% Options for packages loaded elsewhere
\PassOptionsToPackage{unicode}{hyperref}
\PassOptionsToPackage{hyphens}{url}
\documentclass[
  12pt,
  a4paper,
  twocolumn,landscape]{article}
\usepackage{xcolor}
\usepackage[margin=10mm,landscape,includefoot]{geometry}
\usepackage{amsmath,amssymb}
\setcounter{secnumdepth}{-\maxdimen} % remove section numbering
\usepackage{iftex}
\ifPDFTeX
  \usepackage[T1]{fontenc}
  \usepackage[utf8]{inputenc}
  \usepackage{textcomp} % provide euro and other symbols
\else % if luatex or xetex
  \usepackage{unicode-math} % this also loads fontspec
  \defaultfontfeatures{Scale=MatchLowercase}
  \defaultfontfeatures[\rmfamily]{Ligatures=TeX,Scale=1}
\fi
\usepackage{lmodern}
\ifPDFTeX\else
  % xetex/luatex font selection
\fi
% Use upquote if available, for straight quotes in verbatim environments
\IfFileExists{upquote.sty}{\usepackage{upquote}}{}
\IfFileExists{microtype.sty}{% use microtype if available
  \usepackage[]{microtype}
  \UseMicrotypeSet[protrusion]{basicmath} % disable protrusion for tt fonts
}{}
\usepackage{setspace}
\makeatletter
\@ifundefined{KOMAClassName}{% if non-KOMA class
  \IfFileExists{parskip.sty}{%
    \usepackage{parskip}
  }{% else
    \setlength{\parindent}{0pt}
    \setlength{\parskip}{6pt plus 2pt minus 1pt}}
}{% if KOMA class
  \KOMAoptions{parskip=half}}
\makeatother
\setlength{\emergencystretch}{3em} % prevent overfull lines
\providecommand{\tightlist}{%
  \setlength{\itemsep}{0pt}\setlength{\parskip}{0pt}}
% Compact layout and stable page numbers for the private A4 print copy.
\usepackage{microtype}
\usepackage{enumitem}
\setlist[itemize]{leftmargin=1.45em,topsep=2.5pt,itemsep=0.75pt,parsep=0.25pt,partopsep=0pt}
\setlist[enumerate]{leftmargin=1.7em,topsep=2.5pt,itemsep=0.75pt,parsep=0.25pt,partopsep=0pt}

% Supply the CJK glyphs absent from Latin Modern when compiling with XeLaTeX.
\ifXeTeX
  \usepackage{xeCJK}
  \usepackage{newunicodechar}
  \usepackage{seqsplit}
  \usepackage{xstring}
  \setCJKmainfont[AutoFakeSlant=true,AutoFakeBold=true]{PingFang SC}
  \setCJKsansfont[AutoFakeSlant=true,AutoFakeBold=true]{PingFang SC}
  \setCJKmonofont[AutoFakeSlant=true,AutoFakeBold=true]{PingFang SC}
  \newunicodechar{ff}{ff}
\fi

\usepackage{titlesec}
\titleformat{\section}{\normalfont\large\bfseries}{\thesection}{1em}{}
\titleformat{\subsection}{\normalfont\normalsize\bfseries}{\thesubsection}{1em}{}
\titleformat{\subsubsection}{\normalfont\normalsize\bfseries}{\thesubsubsection}{1em}{}
\titlespacing*{\section}{0pt}{8pt plus 1pt minus 1pt}{3pt plus 1pt minus 1pt}
\titlespacing*{\subsection}{0pt}{6pt plus 1pt minus 1pt}{2pt plus 1pt minus 1pt}
\titlespacing*{\subsubsection}{0pt}{4pt plus 1pt minus 1pt}{2pt plus 1pt minus 1pt}

\setlength{\parindent}{0pt}
\setlength{\parskip}{0.12em}
\setlength{\columnsep}{6mm}
\geometry{landscape,includefoot,margin=10mm}
\raggedbottom

\usepackage{fancyhdr}
\setlength{\headheight}{14.5pt}
\setlength{\footskip}{8mm}
\pagestyle{fancy}
\fancyhf{}
\fancyfoot[C]{\thepage}
\renewcommand{\headrulewidth}{0pt}
\renewcommand{\footrulewidth}{0pt}
\fancypagestyle{plain}{%
  \fancyhf{}
  \fancyfoot[C]{\thepage}
  \renewcommand{\headrulewidth}{0pt}
  \renewcommand{\footrulewidth}{0pt}
}

% Make long provenance paths wrap inside the narrow print columns.
\ifXeTeX
  \newcommand{\breakabletexttt}[1]{%
    \IfSubStr{#1}{/}{%
      \path{#1}%
    }{%
      \StrLen{#1}[\breakabletextttlength]%
      \ifnum\breakabletextttlength>60
        \begingroup\ttfamily\seqsplit{#1}\endgroup
      \else
        \begingroup
        \ttfamily
        \CJKfamily{\CJKrmdefault}%
        #1%
        \endgroup
      \fi
    }%
  }
\else
  \newcommand{\breakabletexttt}[1]{\path{#1}}
\fi
\AtBeginDocument{
  \urlstyle{tt}
  \let\texttt\breakabletexttt
  \hypersetup{
    colorlinks=true,
    linkcolor=blue,
    urlcolor=blue,
    citecolor=blue,
    pdftitle={Kajian Pendahuluan: Review dan Tier List}
  }
  \pagenumbering{gobble}
  \begin{titlepage}
    \thispagestyle{empty}
    \pdfbookmark[0]{Sampul}{sampul}
    \vspace*{\fill}
    \begin{center}
      {\Huge\bfseries Kajian Pendahuluan\par}
      \vspace{1.5em}
      {\LARGE Review dan Tier List\par}
      \vspace{2em}
      {\large 70 review per-kode dalam enam domain penilaian\par}
    \end{center}
    \vspace*{\fill}
  \end{titlepage}
  \clearpage
  \pagenumbering{roman}
  \renewcommand{\contentsname}{Daftar Isi}
  \setcounter{tocdepth}{1}
  \pdfbookmark[0]{Daftar Isi}{daftar-isi}
  \tableofcontents
  \clearpage
  \pagenumbering{arabic}
}
\usepackage{bookmark}
\IfFileExists{xurl.sty}{\usepackage{xurl}}{} % add URL line breaks if available
\urlstyle{same}
% fallback for those not using the hyperref driver hyperxmp:
\makeatletter
\@ifundefined{xmpquote}{\newcommand{\xmpquote}[1]{#1}}{}
\makeatother
\hypersetup{
  hidelinks,
  pdfcreator={LaTeX via pandoc}}

\author{}
\date{}

\begin{document}

\setstretch{0.86}
\section{Kajian Pendahuluan: Review dan Tier
List}\label{kajian-pendahuluan-review-dan-tier-list}

Dokumen ini menggabungkan tier list dan 70 review per-kode dalam urutan
daftar kajian pendahuluan. Isi substantif dipertahankan dari file sumber
review; tautan Obsidian pada tier list dinormalisasi menjadi teks untuk
keluaran PDF.

\section{Tier List Kajian
Pendahuluan}\label{tier-list-kajian-pendahuluan}

\subsection{Cara Membaca}\label{cara-membaca}

\begin{itemize}
\tightlist
\item
  Dokumen ini menyusun enam tier list yang berdiri sendiri dari 70
  \texttt{review-\textless{}kode\textgreater{}.md} yang telah tersedia.
\item
  Penilaian tidak dimaksudkan sebagai satu peringkat total. Satu artikel
  dapat berada pada tier berbeda di tiap domain.
\item
  Semua alasan diringkas dari review per-kode. Tautan pada setiap kode
  membuka review yang menjadi basis penilaian.
\item
  \texttt{S} berarti sangat kuat dan selaras; \texttt{A} berarti kuat
  dengan keterbatasan minor; \texttt{B} berarti berguna tetapi memiliki
  keterbatasan penting; \texttt{C} berarti lemah atau hanya dapat
  digunakan secara terbatas; \texttt{D} berarti tidak dapat digunakan
  atau dinilai memadai pada dimensi tersebut.
\item
  \texttt{D} dapat berarti informasi tidak cukup, dimensi tidak berlaku
  untuk desain artikel, atau terdapat ketidakselarasan serius.
  \texttt{D} bukan otomatis vonis bahwa artikel secara keseluruhan
  buruk.
\item
  Ketiadaan \texttt{Future\ Research} tidak dihukum dengan sendirinya.
  Untuk studi nonkuantitatif, \texttt{Dependent\ Variables} dibaca
  secara adaptif sebagai outcome, konstruk, tema, atau objek analisis.
\item
  Dokumen ini adalah sintesis kerja berbasis review TXT, bukan
  verifikasi ulang PDF, data mentah, kode analisis, atau referensi
  artikel secara independen.
\end{itemize}

\subsection{Pemetaan Domain}\label{pemetaan-domain}

\begin{enumerate}
\def\labelenumi{\arabic{enumi}.}
\tightlist
\item
  \textbf{Framing / Argument (F)}: \texttt{Problem\ Statement} dan
  \texttt{Objectives}.
\item
  \textbf{Teori dan Positioning (L)}: \texttt{Literature\ Survey} dan
  \texttt{Research\ Gap}.
\item
  \textbf{Metode dan Evidence (M)}: \texttt{Method\ Used},
  \texttt{Dataset}, \texttt{Population\ /\ Sample}, dan
  \texttt{Dependent\ Variables}.
\item
  \textbf{Hasil dan Kontribusi (E)}: \texttt{Results},
  \texttt{Findings}, \texttt{Conclusions}, dan \texttt{Contributions}.
\item
  \textbf{Transferabilitas (P)}: \texttt{Practical\ Implications} dan
  \texttt{Applications}.
\item
  \textbf{Refleksivitas (X)}: \texttt{Limitations}, \texttt{Challenges},
  dan \texttt{Future\ Research\ (Optional)}.
\end{enumerate}

\subsection{1. Framing / Argument}\label{framing-argument}

\subsubsection{S}\label{s}

\begin{itemize}
\tightlist
\item
  J04 --- Beyond property: masalah tentang peran kawasan industri dalam
  transformasi pascasuburban tajam dan tujuan serta pertanyaan
  penelitian eksplisit.
\item
  J05 --- The Privatization of Metropolitan Jakarta's Urban Fringes:
  masalah privatisasi dan pascasuburbanisasi jelas, dengan tujuan yang
  spesifik dan koheren.
\item
  J07 --- Highway expansion and urban sprawl: masalah kausal, konteks
  JMA, endogenitas, dan tujuan penelitian dirumuskan jelas.
\item
  J08 --- Examining Socio-Economic Inequality Among Commuters: masalah,
  pertanyaan, dan tujuan eksplisit serta selaras dengan fokus
  ketimpangan komuter.
\item
  J11 --- Job Access and Spatial Equity of a Toll Road: masalah,
  konteks, pertanyaan, dan lima tujuan saling terkait dengan baik.
\item
  J26 --- The Dynamics of Jabotabek Development: masalah empiris dan
  kelembagaan serta tujuan analisis dirumuskan jelas untuk studi
  deskriptif-diagnostik.
\item
  A01 --- Stretching the Concept of Borrowed Size: masalah konseptual
  dan tiga pertanyaan empiris dirumuskan jelas serta konsisten.
\item
  A02 --- Borrowing size in networks of cities: masalah dan tujuan
  dirinci menurut skala jaringan, fungsi, dan ukuran kota.
\item
  A03 --- Borrowed Size, Agglomeration Shadows and Cultural Amenities:
  ketegangan borrowed size dan agglomeration shadow diterjemahkan
  menjadi pertanyaan analitis yang jelas.
\item
  A04 --- The Geography of Borrowing Size: masalah, gap
  operasionalisasi, pertanyaan, dan tujuan selaras dengan analisis.
\item
  A05 --- Is there a borrowed size in China's urban agglomerations?:
  problem statement terdiri atas ketidakpastian yang jelas dan tujuan
  dipetakan ke model.
\item
  A10 --- Knowledge innovation network externalities: masalah dan tujuan
  operasional tentang borrowing size versus agglomeration shadow
  spesifik serta selaras.
\item
  A11 --- Do polycentric urban regions promote functional spillovers:
  masalah, tujuan, hipotesis, dan argumen kondisional dirumuskan jelas.
\item
  A12 --- Do large cities contribute to economic growth of small cities:
  masalah, pertanyaan, dan tujuan mencakup hierarki, batas
  administratif, pasar, serta ketahanan.
\item
  A14 --- Does broadband internet allow cities to borrow size: masalah,
  gap, tujuan, mekanisme, dan heterogenitas spasial koheren.
\item
  A17 --- Borrowed size and agglomeration shadows in Japan: masalah dan
  tujuan spesifik, mencakup mekanisme, sektor, dan konteks Jepang.
\item
  A18 --- Improving the Economic Efficiency of Small and Medium-Sized
  Cities: masalah konkret dan hipotesis mencakup kinerja, fungsi,
  heterogenitas, serta mekanisme.
\item
  A19 --- Population agglomeration, borrowed size and urban economic
  growth: masalah, tujuan, dan hipotesis spesifik serta selaras.
\item
  A21 --- The Impact of Borrowing Size on Economic Development: masalah,
  pertanyaan, tujuan, dimensi operasional, dan hipotesis dirumuskan
  jelas.
\item
  A23 --- Shedding Light or Casting Shadows: masalah, konteks, batas
  analisis, dan lima tujuan dijelaskan dengan jelas.
\item
  A25 --- Borrowed Size of Small and Medium Cities: problem dan
  objectives spesifik serta langsung menguji sumber dan skala borrowed
  size.
\item
  A26 --- A three-layer model of borrowed size: masalah, pertanyaan
  penelitian, dan tujuan saling selaras.
\item
  A27 --- Secondary city regionalism: problem statement fokus dan tujuan
  selaras dengan pertanyaan penelitian.
\item
  B01 --- Beyond Polycentricity: masalah, pertanyaan, dan tujuan
  spesifik serta koheren.
\item
  B02 --- Form Follows Function: masalah polisentrisitas dan tujuan
  operasional diterjemahkan langsung ke ukuran morfologis-fungsional.
\item
  B03 --- Functional Polycentrism and Urban Network Development: masalah
  central place versus jaringan dan tujuan lintas skala dirumuskan
  jelas.
\item
  B04 --- On the Economic Foundation of the Urban Network Paradigm:
  masalah dan tiga pertanyaan penelitian dirumuskan eksplisit serta
  operasional.
\item
  B05 --- Assessing Polycentric Urban Systems in the OECD: masalah dan
  tujuan komparatif pada tiga skala jelas serta selaras.
\item
  B07 --- Summing Small Cities Does Not Make a Large City: masalah
  komparatif dan objectives utama maupun operasional spesifik.
\item
  B08 --- Secondary Yet Metropolitan: masalah, konteks, dan tujuan utama
  maupun operasional eksplisit.
\item
  B09 --- The identification of sub-centres: masalah metodologis dan
  tiga tujuan penelitian jelas serta selaras.
\item
  B10 --- Urban Polycentricity and the Costs of Commuting: masalah,
  konteks, dan tujuan empiris selaras dengan pengukuran.
\item
  B11 --- The Rise of Second-Rank Cities: masalah analitis jelas dan
  empat hipotesis spesifik.
\item
  B12 --- City Size and Economic Performance: masalah dan pertanyaan
  utama eksplisit serta tujuan mudah diidentifikasi.
\item
  B13 --- Spatial structure and productivity in US metropolitan areas:
  masalah dan tujuan langsung diterjemahkan ke pengujian struktur
  spasial dan produktivitas.
\item
  B14 --- Polycentric Spatial Structure and Its Economic Performance:
  masalah inkonsistensi bukti dan tujuan meta-analisis dirumuskan jelas.
\item
  B15 --- Spatial structure and productivity in European regions:
  masalah jelas dan empat hubungan yang diuji dirumuskan melalui
  objectives serta hipotesis.
\item
  B16 --- Does Urban Polycentricity Contribute to Regional Economic
  Growth: masalah, tujuan, periode, lokasi, kinerja, desain, dan
  hipotesis eksplisit.
\item
  B17 --- Agglomeration Externalities or Network Externalities: problem,
  gap, pertanyaan, tujuan, dan enam hipotesis jelas.
\item
  B18 --- Big in the neighbourhood: masalah tentang keterbatasan
  peringkat nasional dan tiga tujuan penelitian eksplisit.
\item
  B19 --- Territorial Arrangements of Small and Medium-Sized Towns:
  masalah jelas dengan objectives operasional yang spesifik.
\item
  B20 --- Grid and shake: masalah sensitivitas unit spasial dan tujuan
  metodologis dirumuskan spesifik.
\item
  B21 --- Do New Economic Geography agglomeration shadows: masalah dan
  tujuan empiris jelas serta mencakup mekanisme dan hierarki urban.
\item
  B22 --- Urban growth shadows: masalah dan tujuan menjelaskan perubahan
  hubungan kedekatan pusat besar dengan pertumbuhan serta mekanismenya.
\item
  B23 --- Cities, hinterlands and agglomeration shadows: masalah,
  prediksi NEG, dan tujuan empiris saling terhubung.
\end{itemize}

\subsubsection{A}\label{a}

\begin{itemize}
\tightlist
\item
  J01 --- Identifying post-suburbanization: problem dan tujuan jelas
  serta koheren, tetapi masih memiliki beberapa unsur umum.
\item
  J03 --- Industrial Land Development and Manufacturing Deconcentration:
  masalah, batas wilayah, dan tujuan jelas, meski hubungan lapis masalah
  lebih deskriptif.
\item
  J06 --- Transportation Network and Changes in Urban Structure: masalah
  suburbanisasi dan endogenitas jelas, tetapi pertanyaan formal tidak
  tunggal.
\item
  J09 --- Motorcycle-based ride-hailing among commuters: masalah dan
  tujuan jelas, tetapi framing mobilitas individual dan bentuk spasial
  belum sepenuhnya teroperasionalisasi.
\item
  J10 --- Spatial, mobility, or socio-economic inequity: masalah dan
  tujuan operasional jelas, tetapi pertanyaan formal serta hipotesis
  tidak dilaporkan.
\item
  J22 --- Morphological and Functional Polycentricity in Indonesia:
  masalah, gap, pertanyaan, dan objectives koheren, tetapi mekanisme
  kausal belum diuji.
\item
  J23 --- Urban transformation and inter-suburban commuting: masalah,
  konteks, pertanyaan, dan tujuan jelas, tetapi bahasa kausal melampaui
  desain asosiasional.
\item
  J25 --- Continuity and change in mega-urbanization: masalah dan tujuan
  jelas, tetapi cakupan mega-urbanisasi lebih luas daripada pertanyaan
  desentralisasi.
\item
  J27 --- Peri-urban transformation in the Jakarta metropolitan area:
  masalah dan tujuan transformasi sosial-ekonomi jelas, tetapi
  pertanyaan formal tidak dinyatakan.
\item
  A06 --- Evaluation of Service Function of Small Towns: masalah dan
  tujuan operasional jelas, meski atribusi mekanisme melampaui
  pengukuran.
\item
  A07 --- Urban borrowed size in Guangzhou-Foshan: problem dan
  objectives spesifik, tetapi hipotesis formal dan mekanisme lebih luas
  daripada desain.
\item
  A08 --- Regional Integration Strategies in the Yangtze River Delta:
  masalah, pertanyaan, dan hipotesis jelas, tetapi cakupan analisis
  dipersempit pada pola dua sektor.
\item
  A09 --- Borrowed administrative power and high-speed rail: masalah,
  pertanyaan empiris, tujuan, dan hipotesis jelas, meski mekanisme
  borrowing tidak diukur langsung.
\item
  A13 --- Industrial Composition and Agglomeration Shadow: problem dan
  tujuan jelas, tetapi pertanyaan asosiasi dan klaim shadow agak
  tercampur.
\item
  A15 --- Inter-regional networks and productive efficiency: masalah dan
  tujuan mencakup beberapa jalur analisis, tetapi hipotesis formal tidak
  dirumuskan.
\item
  A16 --- Regional network economies and industrial productivity:
  masalah dan tujuan jelas, tetapi hipotesis formal tidak disajikan.
\item
  A20 --- Borrowing Size and Urban Green Development Efficiency:
  masalah, pertanyaan, dan tujuan selaras, walau sebagian tujuan
  direkonstruksi dari desain.
\item
  A22 --- Small Towns in the South Moravian Region: masalah, gap,
  tujuan, dan hipotesis jelas, tetapi isu pengukuran dan mekanisme
  bergabung.
\item
  A24 --- Border cities: Out of the shadow: masalah dan tujuan spesifik
  serta selaras dengan hipotesis dan estimasi.
\item
  A28 --- Small Firms, Borrowed Size and the Urban-Rural Shift: masalah
  dan tujuan eksplorasi jelas, meski tanpa pertanyaan atau hipotesis
  formal.
\item
  B06 --- The productivity effects of polycentricity: masalah dan tujuan
  terarah, tetapi sebagian tujuan direkonstruksi dari teks.
\item
  B24 --- Identifying Agglomeration Shadows: masalah dan tujuan berlapis
  jelas, walau sebagian uji menjawab pertanyaan yang tidak sepenuhnya
  sama.
\item
  B25 --- Urban Shadow Areas from People Flows: masalah dan tujuan lima
  langkah cukup jelas, tetapi pertanyaan atau hipotesis formal tidak
  ada.
\item
  C02 --- Urban Cannibalism: masalah dan tujuan koheren, tetapi
  objectives direkonstruksi karena tidak ada subbagian formal.
\end{itemize}

\subsubsection{B}\label{b}

\begin{itemize}
\tightlist
\item
  J02 --- Pola Pergerakan dan Dekonsentrasi Pekerjaan: masalah dan
  tujuan cukup jelas, tetapi pertanyaan formal tidak dirumuskan dan dua
  konsep utama tercampur.
\end{itemize}

\subsubsection{C}\label{c}

\begin{itemize}
\tightlist
\item
  Tidak ada artikel pada tier ini.
\end{itemize}

\subsubsection{D}\label{d}

\begin{itemize}
\tightlist
\item
  Tidak ada artikel pada tier ini.
\end{itemize}

\subsection{2. Teori dan Positioning}\label{teori-dan-positioning}

\subsubsection{S}\label{s-1}

\begin{itemize}
\tightlist
\item
  A14 --- Broadband internet and borrowed size: teori borrowed size,
  broadband, pasar kerja, dan beberapa gap diposisikan sangat kuat.
\item
  A17 --- Borrowed size and agglomeration shadows in Japan: landasan
  teori, literatur empiris, saluran mekanisme, dan gap terhubung
  langsung dengan desain.
\item
  A21 --- Borrowing Size and Economic Development: literatur, teori
  network externality, dan empat research gap eksplisit.
\item
  A23 --- Shedding Light or Casting Shadows: lima arus literatur, konsep
  utama, dan gap dirangkum secara terstruktur.
\item
  A25 --- Borrowed Size of Small and Medium Cities: konsep, mekanisme,
  temuan empiris, konteks, dan gap dirumuskan eksplisit.
\item
  B03 --- Functional Polycentrism and Urban Network Development: teori,
  studi empiris, dan konteks kebijakan koheren dengan tiga gap
  eksplisit.
\item
  B10 --- Urban Polycentricity and Commuting Costs: literatur dan temuan
  yang bertentangan dipakai untuk membangun gap pengukuran dan
  pengujian.
\item
  B11 --- The Rise of Second-Rank Cities: agglomeration economies,
  SOUDY, jaringan kota, borrowed size, dan gap terhubung baik.
\item
  B12 --- City Size and Economic Performance: teori-teori yang bersaing
  dibahas dan gap kota tingkat kedua dinyatakan.
\item
  B13 --- Spatial structure and productivity in US metropolitan areas:
  konsep utama, argumen berlawanan, dan gap empiris diposisikan jelas.
\item
  B14 --- Polycentric Spatial Structure and Its Economic Performance:
  perdebatan konsep, pengukuran, skala, dan endogenitas dirangkum lintas
  studi.
\item
  B16 --- Urban Polycentricity and Regional Economic Growth: literatur
  kinerja, mekanisme, konteks China, inkonsistensi hasil, dan gap
  longitudinal jelas.
\item
  B17 --- Agglomeration or Network Externalities: konsep, debat
  pengukuran, dan gap dihubungkan langsung ke enam hipotesis.
\item
  B18 --- Big in the neighbourhood: CPT, mekanisme layanan, studi
  terdahulu, dan gap hinterland dibahas terstruktur.
\item
  B22 --- Urban growth shadows: literatur historis, spasial, dan konteks
  Amerika Serikat membangun gap yang jelas.
\end{itemize}

\subsubsection{A}\label{a-1}

\begin{itemize}
\tightlist
\item
  J01 --- Identifying post-suburbanization: konsep dan gap empiris kuat,
  tetapi tinjauan bersifat naratif.
\item
  J03 --- Industrial Land Development and Manufacturing Deconcentration:
  pustaka lintas konteks dan gap jelas, tanpa strategi tinjauan
  sistematis.
\item
  J04 --- Beyond property: post-suburbia, desakota, dan pembanding
  internasional memberi positioning kuat, tetapi selektif.
\item
  J05 --- The Privatization of Metropolitan Jakarta's Urban Fringes:
  konsep dan pembanding Asia relevan, tetapi gap lanjutan sebagian
  inferensial.
\item
  J06 --- Transportation Network and Urban Structure: literatur dan gap
  empiris, pengukuran, serta konteks jelas, tetapi mekanisme tidak diuji
  langsung.
\item
  J07 --- Highway expansion and urban sprawl: teori, pengukuran,
  identifikasi kausal, konteks, dan gap tercakup, tetapi tinjauan tidak
  sistematis.
\item
  J08 --- Socio-Economic Inequality Among Commuters: literatur relevan
  dan gap terpadu, tetapi mekanisme kausal belum dirumuskan.
\item
  J10 --- Job accessibility evaluation: kerangka tiga dimensi dan gap
  jelas, meski kajian naratif.
\item
  J11 --- Job Access and Spatial Equity: literatur mendukung konsep dan
  metode, tetapi surveinya tidak sistematis.
\item
  J22 --- Morphological and Functional Polycentricity in Indonesia:
  debat, definisi, metode, konteks, dan gap dibahas luas, meski tidak
  sistematis.
\item
  J25 --- Mega-urbanization in the Jakarta-Bandung Region: literatur
  multidimensi dan gap kontekstual dibangun, tetapi strategi pencarian
  tidak dilaporkan.
\item
  J26 --- The Dynamics of Jabotabek Development: literatur relevan dan
  beberapa gap teridentifikasi, walau tidak ada positioning formal
  tersendiri.
\item
  A01 --- Stretching the Concept of Borrowed Size: literatur relevan,
  bukti campuran, dan gap eksplisit, tetapi protokol pencarian tidak
  ada.
\item
  A02 --- Borrowing size in networks of cities: aglomerasi, jaringan,
  shadow, dan gap jelas, meski tinjauan naratif.
\item
  A03 --- Borrowed Size and Cultural Amenities: teori dan literatur
  lintas tradisi terhubung dengan gap, tetapi tidak dibandingkan formal.
\item
  A04 --- Geography of Borrowing Size: konsep, temuan berlawanan, dan
  gap dibangun, meski pencarian literatur tidak dijelaskan.
\item
  A05 --- Borrowed size in China's urban agglomerations: landasan teori
  dan gap luas serta spesifik, tetapi bukan survei sistematis.
\item
  A06 --- Service Function of Small Towns: teori, studi pembanding,
  konteks China, dan gap dibahas luas tanpa protokol penelusuran.
\item
  A07 --- Urban borrowed size in Guangzhou-Foshan: literatur dan gap
  tingkat town relevan, tetapi tinjauan naratif.
\item
  A08 --- Regional Integration Strategies in YRD: teori city region, new
  regionalism, borrowed size, shadow, dan gap terhubung.
\item
  A10 --- Knowledge innovation network externalities: landasan
  konseptual luas dan gap eksternalitas jaringan jelas, tetapi tidak
  sistematis.
\item
  A11 --- Functional spillovers and economic performance: konsep dan gap
  empiris terhubung baik, meski tetap kontekstual.
\item
  A12 --- Large cities and small-city growth: teori dan gap jelas,
  tetapi tinjauan naratif.
\item
  A13 --- Industrial Composition and Agglomeration Shadow: literatur
  NEG, borrowed size, market potential, dan gap eksplisit.
\item
  A15 --- Inter-regional networks and productive efficiency: literatur
  dan gap Jepang spesifik, tetapi survei tidak sistematis.
\item
  A16 --- Regional network economies and productivity: literatur dan gap
  teoretis kuat, tetapi korpus dan seleksi tidak dilaporkan.
\item
  A18 --- Economic Efficiency of Small and Medium-Sized Cities: konsep
  dan gap eksplisit serta terhubung ke analisis, meski naratif.
\item
  A19 --- Population agglomeration and urban economic growth: indikator,
  mekanisme, variasi konteks, dan gap dibahas jelas.
\item
  A20 --- Urban Green Development Efficiency: literatur, hipotesis,
  borrowing size, dan tiga gap diposisikan cukup jelas.
\item
  A22 --- Small Towns in South Moravia: posisi teoritis luas dan gap
  small towns eksplisit.
\item
  A24 --- Border cities: literatur jaringan kota dan efek perbatasan
  terhubung baik dengan konsep serta gap.
\item
  A26 --- Three-layer model of borrowed size: literatur dan gap teoretis
  kuat, tetapi protokol korpus tidak tersedia.
\item
  A27 --- Secondary city regionalism: tinjauan luas dan gap agensi kota
  sekunder jelas, meski naratif.
\item
  A28 --- Small Firms and the Urban-Rural Shift: positioning borrowed
  size dan gap dijelaskan kuat, tetapi selektif.
\item
  B01 --- Beyond Polycentricity: literatur relevan dan gap jelas, meski
  tidak sistematis.
\item
  B02 --- Form Follows Function: literatur historis-konseptual dan gap
  jelas, tetapi protokol korpus tidak ada.
\item
  B04 --- Economic Foundation of the Urban Network Paradigm: kerangka
  teori dan gap kuat, tetapi tinjauan naratif.
\item
  B05 --- Polycentric Urban Systems in the OECD: literatur konseptual
  relevan dan gap eksplisit, meski tidak sistematis.
\item
  B06 --- Productivity effects of polycentricity: agglomeration
  externalities, borrowed size, dan gap kuat, tetapi survei tidak dapat
  diaudit sebagai sistematis.
\item
  B07 --- Summing Small Cities: perdebatan PUR, aglomerasi, borrowed
  size, dan amenitas dipetakan dengan gap jelas.
\item
  B08 --- Secondary Yet Metropolitan: literatur luas dan terintegrasi,
  tetapi gap tidak dirumuskan sebagai bagian khusus.
\item
  B09 --- Identification of sub-centres: literatur morfologis,
  interaksi, dan teori pusat dibahas dengan gap eksplisit.
\item
  B15 --- Spatial structure and productivity in Europe: literatur
  aglomerasi, borrowed size, skala, dan konteks Eropa diposisikan jelas.
\item
  B19 --- Territorial Arrangements of Towns: landasan teori relevan dan
  gap teridentifikasi.
\item
  B20 --- Grid and shake: literatur relevan dengan gap sensitivitas
  agregasi yang eksplisit.
\item
  B21 --- NEG agglomeration shadows and population dynamics: CPT, NEG,
  teori alternatif, dan gap tersusun kuat.
\item
  B23 --- Cities, hinterlands and agglomeration shadows: landasan teori
  dan gap empiris kuat, meski survei naratif.
\item
  B24 --- Agglomeration Shadows from Ancient Ports: rantai literatur dan
  gap identifikasi serta diagnosis pola koheren.
\item
  B25 --- Urban Shadow Areas from People Flows: literatur relevan dan
  gap intra-perkotaan jelas, meski tidak sistematis.
\item
  C02 --- Urban Cannibalism: literatur relevan dan gap konseptual
  spesifik, tetapi tinjauan tidak sistematis.
\end{itemize}

\subsubsection{B}\label{b-1}

\begin{itemize}
\tightlist
\item
  J02 --- Pola Pergerakan dan Dekonsentrasi Pekerjaan: kerangka DUS dan
  literatur relevan, tetapi integrasi literatur dengan gap metodologis
  belum kuat.
\item
  J09 --- Motorcycle-based ride-hailing: literatur luas, tetapi hubungan
  ekonomi berbagi, pascasuburbanisasi, dan rezim perkotaan belum menjadi
  mekanisme operasional.
\item
  J23 --- Urban transformation and inter-suburban commuting: literatur
  relevan, tetapi positioning masih naratif dan mekanisme individu
  diproksikan secara agregat.
\item
  J27 --- Peri-urban transformation: literatur relevan, tetapi gap
  formal dan strategi penelusuran tidak tegas.
\item
  A09 --- Borrowed administrative power and high-speed rail: literatur
  relevan dan gap eksplisit, tetapi survei naratif tanpa penilaian mutu.
\end{itemize}

\subsubsection{C}\label{c-1}

\begin{itemize}
\tightlist
\item
  Tidak ada artikel pada tier ini.
\end{itemize}

\subsubsection{D}\label{d-1}

\begin{itemize}
\tightlist
\item
  Tidak ada artikel pada tier ini.
\end{itemize}

\subsection{3. Metode dan Evidence}\label{metode-dan-evidence}

\subsubsection{S}\label{s-2}

\begin{itemize}
\tightlist
\item
  A18 --- Economic Efficiency of Small and Medium-Sized Cities: desain
  panel, model, operasionalisasi, sumber data, sampel 82 kota, dan uji
  ketahanan dijelaskan memadai.
\item
  B13 --- Spatial structure and productivity in US metropolitan areas:
  metode, variabel, OLS/TSLS, instrumen, data, sampel, dan outcome
  dijelaskan lengkap.
\item
  B16 --- Urban Polycentricity and Regional Economic Growth: FE-IV,
  panel 19 UR, konstruksi indeks, TFP, unit analisis, dan keterbatasan
  data didokumentasikan rinci.
\end{itemize}

\subsubsection{A}\label{a-2}

\begin{itemize}
\tightlist
\item
  A14 --- Broadband internet and borrowed size: data longitudinal, fixed
  effects, lag, kontrol, dan uji sensitivitas kuat, meski identifikasi
  serta pengukuran masih terbatas.
\item
  A15 --- Inter-regional networks and productive efficiency: SFA, fungsi
  translog, panel 47 prefektur, sumber data, unit, dan outcome
  dijelaskan rinci.
\item
  A16 --- Regional network economies and productivity: metode dynamic
  TFP, variabel, sumber data, 47 prefektur, dan periode dijelaskan
  lengkap.
\item
  A17 --- Borrowed size and agglomeration shadows in Japan: panel
  spasial dan RTFP terdokumentasi baik, dengan beberapa inkonsistensi
  pelaporan yang masih terbuka.
\item
  A21 --- Borrowing Size and Economic Development: metode, dataset,
  sampel, pengukuran, dan outcome rinci, meski klasifikasi kota dan
  jarak memiliki masalah.
\item
  A23 --- Shedding Light or Casting Shadows: indikator, dataset, unit,
  klasifikasi, dan model lengkap, walau sampling dan sebagian outcome
  tidak sepenuhnya tersedia.
\item
  A25 --- Borrowed Size of Small and Medium Cities: 152 unit, sumber
  data, eksklusi, delapan outcome, dan OLS dijelaskan baik, dengan
  keterbatasan desain.
\item
  B03 --- Functional Polycentrism and Urban Network Development:
  analisis arus longitudinal, dataset, unit, model, dan outcome
  dijelaskan memadai.
\item
  B05 --- Polycentric Urban Systems in the OECD: indikator, sumber data,
  unit, skala, dan sampel wilayah cukup terjelaskan.
\item
  B06 --- Productivity effects of polycentricity: desain, dataset, TFP,
  regresi spasial, dan GS2SLS dilaporkan rinci.
\item
  B07 --- Summing Small Cities: desain, dataset, unit, indeks
  polisentrisitas, amenitas, dan outcome dijelaskan memadai.
\item
  B09 --- Identification of sub-centres: metode, indikator, ambang,
  sumber data, unit, observasi, dan evaluasi dijelaskan jelas.
\item
  B10 --- Urban Polycentricity and Commuting Costs: metode, dataset 82
  metropolitan areas, outcome, model, dan robustness cukup kuat.
\item
  B11 --- The Rise of Second-Rank Cities: metode, sumber dataset, unit,
  sampel, dan outcome jelas, meski cakupan observasi tidak konsisten.
\item
  B14 --- Polycentric Spatial Structure and Its Economic Performance:
  metode meta-analisis dan 139 estimasi dari 31 studi cukup
  terdokumentasi.
\item
  B17 --- Agglomeration or Network Externalities: metode, data,
  populasi, sampel, dan outcome sangat terperinci, dengan keterbatasan
  panel dan proksi.
\item
  B18 --- Big in the neighbourhood: dataset settlement yang besar,
  subset analisis, metode, dan outcome dijelaskan rinci.
\item
  B21 --- NEG agglomeration shadows and population dynamics: data
  county, sampel, variabel, spesifikasi, dan GMM terdokumentasi kuat.
\item
  B22 --- Urban growth shadows: metode, dataset, unit, sampel, dan
  outcome rinci, meski kode dan data mentah tidak tersedia.
\item
  B23 --- Cities, hinterlands and agglomeration shadows: desain
  longitudinal, dataset, unit, variabel, dan pengujian dijelaskan rinci.
\item
  B24 --- Agglomeration Shadows from Ancient Ports: metode, grid,
  outcome, kontrol, cakupan sampel, dan robustness terdokumentasi baik.
\end{itemize}

\subsubsection{B}\label{b-2}

\begin{itemize}
\tightlist
\item
  J01 --- Identifying post-suburbanization: desain, survei, sampel,
  populasi, dan proksi tersedia, tetapi definisi serta prosedur data
  tidak lengkap.
\item
  J03 --- Industrial Land Development and Manufacturing Deconcentration:
  metode dan indikator longitudinal cukup terlihat, tetapi pembersihan
  data, cakupan sensus, dan kausalitas terbatas.
\item
  J05 --- The Privatization of Metropolitan Jakarta's Urban Fringes:
  pendekatan deskriptif dan sumber sekunder sesuai tujuan, tetapi
  seleksi dan integrasi data kurang dijelaskan.
\item
  J06 --- Transportation Network and Urban Structure: metode, data,
  sampel, dan outcome rinci, tetapi instrumen, proksi, dan observasi
  bermasalah.
\item
  J07 --- Highway expansion and urban sprawl: IV/TSLS, dataset, unit,
  dan outcome rinci, tetapi tahun, nol, dependensi spasial, dan seleksi
  observasi kurang jelas.
\item
  J11 --- Job Access and Spatial Equity: desain skenario, data, unit,
  sampel, dan model tersedia, tetapi sampling, harmonisasi, dan validasi
  terbatas.
\item
  J22 --- Morphological and Functional Polycentricity in Indonesia:
  metode, data, unit, ukuran, dan sampel tersedia, tetapi CBD, jarak,
  waktu, dan prosedur kualitatif bermasalah.
\item
  J26 --- The Dynamics of Jabotabek Development: metode, dataset, unit
  wilayah, dan outcome rinci, tetapi sampling, pengolahan, dan perubahan
  batas tidak tersedia.
\item
  A01 --- Stretching the Concept of Borrowed Size: desain, unit,
  operasionalisasi, dan model cukup rinci, tetapi data, sampel, dan
  kausalitas terbatas.
\item
  A02 --- Borrowing size in networks of cities: dataset 1.114 MUA dan
  model rinci, tetapi desain potong lintang serta waktu dan spasialnya
  tidak sepenuhnya selaras.
\item
  A03 --- Borrowed Size and Cultural Amenities: dataset besar dan
  regresi sesuai dengan indeks, tetapi konstruksi indeks, seleksi, dan
  dependensi FUA kurang lengkap.
\item
  A04 --- Geography of Borrowing Size: unit, indikator, sampel, OLS, dan
  outcome disebutkan, tetapi provenance data, diagnostik, dan
  konsistensi sampel bermasalah.
\item
  A05 --- Borrowed size in China's urban agglomerations: metode, data,
  sampel, dan outcome cukup lengkap, tetapi data hilang, diagnostik, dan
  kausalitas tidak tersedia.
\item
  A09 --- Borrowed administrative power and high-speed rail: DID/DDD,
  panel, 199 county, periode, dan outcome jelas, tetapi PSM, standard
  error, dan proksi terbatas.
\item
  A11 --- Functional spillovers and economic performance: desain,
  dataset, variabel, dan estimasi cukup jelas, tetapi instrumen,
  representativitas, dan robustness terbatas.
\item
  A12 --- Large cities and small-city growth: OLS, 108 kota, periode,
  data, dan outcome dijelaskan, tetapi desain observasional dan proksi
  jarak membatasi.
\item
  A19 --- Population agglomeration and urban economic growth: metode,
  dataset, sampel, dan outcome rinci, tetapi data rahasia, N, bobot
  spasial, dan outcome ambigu.
\item
  A20 --- Urban Green Development Efficiency: variabel, dataset, model,
  pengukuran, dan robustness rinci, tetapi sampel serta observasi tidak
  konsisten.
\item
  A24 --- Border cities: metode, 1.920 MUA, populasi, dan outcome rinci,
  tetapi desain lintas-seksi dan proksi interaksi membatasi.
\item
  A26 --- Three-layer model of borrowed size: desain, data, periode,
  unit, dan indikator jelas, tetapi sampel, formula, serta robustness
  tidak dilaporkan.
\item
  A27 --- Secondary city regionalism: bahan, enam kasus, dan 15
  wawancara cukup dijelaskan, tetapi seleksi kasus, analisis, dan
  saturasi tidak lengkap.
\item
  A28 --- Small Firms and the Urban-Rural Shift: metode survei, dataset
  100 respons, sampel, variabel, dan analisis tersedia, tetapi seleksi
  dan outcome formal kurang jelas.
\item
  B01 --- Beyond Polycentricity: metode, dataset, populasi, sampel, dan
  outcome dijelaskan, tetapi data hilang, bobot, dan diagnostik tidak
  tersedia.
\item
  B02 --- Form Follows Function: kerangka pengukuran, dataset 42 WGR,
  unit, dan indikator tersedia, tetapi survei dan ketidakpastian kurang
  dilaporkan.
\item
  B04 --- Economic Foundation of the Urban Network Paradigm: metode,
  dataset, sampel, unit, dan outcome rinci, tetapi nonresponse dan
  mismatch municipality membatasi.
\item
  B08 --- Secondary Yet Metropolitan: desain, sumber data, sampel, dan
  ukuran perbandingan dijelaskan, tetapi tahun, formula, dan pengukuran
  kelembagaan kurang.
\item
  B12 --- City Size and Economic Performance: dataset, unit, cakupan,
  indikator, dan periode tersedia, tetapi prosedur analisis dan outcome
  formal tidak lengkap.
\item
  B15 --- Spatial structure and productivity in Europe: metode, dataset,
  unit, outcome, OLS, dan TSLS rinci, tetapi agregasi, seleksi, dan
  first-stage bermasalah.
\item
  B19 --- Territorial Arrangements of Towns: prosedur dan unit rinci,
  tetapi provenance dataset, harmonisasi, ambang, dan model luaran
  terbatas.
\item
  B20 --- Grid and shake: metode, dataset, populasi, dan outcome
  tersedia, tetapi detail hilang serta jumlah replikasi dan observasi
  tidak konsisten.
\item
  C02 --- Urban Cannibalism: unit, indeks, data, dan cakupan 149
  heksagon jelas, tetapi kalibrasi, sensitivitas, data mentah, dan kode
  tidak tersedia.
\end{itemize}

\subsubsection{C}\label{c-2}

\begin{itemize}
\tightlist
\item
  J02 --- Pola Pergerakan dan Dekonsentrasi Pekerjaan: desain, 116
  responden, sampling, dan uji tersedia, tetapi representativitas,
  instrumen, dan statistik penting kurang.
\item
  J04 --- Beyond property: dataset dan unit kasus dapat dikenali, tetapi
  metode formal, validasi, harmonisasi, dan outcome tidak dijelaskan.
\item
  J10 --- Job accessibility evaluation: luaran aksesibilitas dan metode
  deskriptif terlihat, tetapi unit, waktu, klasifikasi, sampel, dan
  validasi bermasalah.
\item
  J23 --- Urban transformation and inter-suburban commuting: panel
  spasial dan variabel cukup terlihat, tetapi observasi, data hilang,
  bobot, pemetaan waktu, dan kode tidak tersedia.
\item
  J25 --- Mega-urbanization in the Jakarta-Bandung Region: survei
  regional deskriptif dapat dikenali, tetapi harmonisasi, batas wilayah,
  operasionalisasi, dan kausalitas sangat terbatas.
\item
  J27 --- Peri-urban transformation: desain multisumber dapat dikenali,
  tetapi replikasi, harmonisasi, jumlah responden, dan pengukuran
  segregasi bermasalah.
\item
  A06 --- Service Function of Small Towns: prosedur dan dataset
  tersedia, tetapi observasi hilang, interpolasi, AHP, proksi, dan
  diagnostik kurang transparan.
\item
  A07 --- Urban borrowed size in Guangzhou-Foshan: metode, 66 observasi,
  populasi, dan outcome terlihat, tetapi seleksi, transformasi, proksi,
  dan diagnostik bermasalah.
\item
  A08 --- Regional Integration Strategies in YRD: LQ-DID-DDD
  terstruktur, tetapi unit, simbol, persamaan, tabel, dan fieldwork
  tidak konsisten.
\item
  A10 --- Knowledge innovation network externalities: metode, data,
  cakupan kota, periode, dan outcome dapat diidentifikasi, tetapi
  sampling, atribusi, dan model kurang memadai.
\item
  A13 --- Industrial Composition and Agglomeration Shadow: model dan
  variabel disebutkan, tetapi sumber nilai tambah, konstruksi observasi,
  FE/RE, dan uji spasial tidak tersedia.
\item
  A22 --- Small Towns in South Moravia: dataset, sampel, model, dan
  outcome dijelaskan, tetapi bobot, klasifikasi, pembersihan, waktu, dan
  pengujian kurang.
\item
  B25 --- Urban Shadow Areas from People Flows: metode eksploratif dan
  data dapat dipahami, tetapi parameter, representativitas, validasi,
  dan ketidakpastian tidak dilaporkan.
\end{itemize}

\subsubsection{D}\label{d-2}

\begin{itemize}
\tightlist
\item
  J08 --- Socio-Economic Inequality Among Commuters: informasi analitis,
  pembobotan, pengodean, data hilang, dan rincian instrumen tidak
  memadai untuk menilai atau mereplikasi evidence.
\item
  J09 --- Motorcycle-based ride-hailing: unit analisis, ekspansi data,
  bobot, outcome, diagnostik, dan pelaporan regresi memiliki mismatch
  serius.
\end{itemize}

\subsection{4. Hasil dan Kontribusi}\label{hasil-dan-kontribusi}

\subsubsection{S}\label{s-3}

\begin{itemize}
\tightlist
\item
  J05 --- The Privatization of Metropolitan Jakarta's Urban Fringes:
  hasil indikator dan contoh empiris kaya, temuan serta kesimpulan
  selaras, dan kontribusi dinyatakan jelas.
\item
  A14 --- Broadband internet and borrowed size: hasil rinci,
  heterogenitas diuji, kontribusi eksplisit, dan kesimpulan dikalibrasi
  sebagai bukti sugestif.
\item
  A16 --- Regional network economies and productivity: hasil numerik,
  elastisitas, robustness, perbedaan industri, dan kontribusi terhubung
  baik.
\item
  A18 --- Economic Efficiency of Small and Medium-Sized Cities: hasil,
  heterogenitas, mekanisme, kesimpulan, dan kontribusi dipaparkan
  lengkap.
\item
  A23 --- Shedding Light or Casting Shadows: hasil, statistik, temuan,
  kesimpulan, interpretasi, dan kontribusi dibedakan serta ditelusuri.
\item
  A27 --- Secondary city regionalism: hasil kaya dan tipologi proyek
  konsisten dengan temuan, kesimpulan, serta kontribusi.
\item
  B03 --- Functional Polycentrism and Urban Network Development: hasil
  lintas skala, temuan, kesimpulan, dan kontribusi dilaporkan jelas
  dengan kualifikasi.
\item
  B04 --- Economic Foundation of the Urban Network Paradigm: hasil
  menjawab tujuan, kesimpulan selaras, dan kontribusi
  empiris-metodologis jelas.
\item
  B09 --- Identification of sub-centres: hasil kuantitatif dan
  perbandingan metode terdokumentasi lengkap.
\item
  B11 --- The Rise of Second-Rank Cities: hasil, temuan, kesimpulan, dan
  kontribusi rinci dengan signifikansi statistik.
\item
  B12 --- City Size and Economic Performance: hasil komparatif rinci,
  temuan bernuansa, kesimpulan bersyarat, dan kontribusi jelas.
\item
  B14 --- Polycentric Spatial Structure and Its Economic Performance:
  hasil meta-analisis konsisten dan kontribusi disertai kehati-hatian
  terhadap klaim universal.
\item
  B16 --- Urban Polycentricity and Regional Economic Growth: hasil
  heterogen, kesimpulan hati-hati, dan tiga kontribusi utama
  terdokumentasi lengkap.
\item
  B17 --- Agglomeration or Network Externalities: hasil, temuan,
  kesimpulan, dan tiga kontribusi terhubung, termasuk ketidakselarasan
  yang dilaporkan.
\item
  B21 --- NEG agglomeration shadows and population dynamics: hasil
  kuantitatif dan mekanisme per kelompok dilaporkan lengkap dengan
  kualifikasi.
\item
  B22 --- Urban growth shadows: hasil numerik dan lima stylized facts
  tersusun koheren serta nonkausal.
\end{itemize}

\subsubsection{A}\label{a-3}

\begin{itemize}
\tightlist
\item
  J03 --- Industrial Land Development and Manufacturing Deconcentration:
  hasil longitudinal kaya dan konsisten dengan indikator, dengan
  kontribusi yang terutama diinferensikan.
\item
  J06 --- Transportation Network and Urban Structure: hasil dan
  kontribusi rinci, dibedakan menurut moda dan zona serta diberi
  kualifikasi.
\item
  J07 --- Highway expansion and urban sprawl: hasil, temuan, kesimpulan,
  dan kontribusi terhubung dengan tujuan meski sensitif terhadap
  resolusi.
\item
  J11 --- Job Access and Spatial Equity: hasil, temuan, kesimpulan, dan
  kontribusi konsisten serta didukung angka dan validasi model.
\item
  J22 --- Morphological and Functional Polycentricity in Indonesia:
  hasil, temuan, kesimpulan, dan kontribusi selaras meski ada
  inkonsistensi kecil.
\item
  J26 --- The Dynamics of Jabotabek Development: hasil rinci dan
  kesimpulan konsisten, tetapi kontribusi tidak dinyatakan eksplisit.
\item
  A01 --- Stretching the Concept of Borrowed Size: hasil, temuan,
  kesimpulan, dan kontribusi koheren serta dikualifikasi.
\item
  A02 --- Borrowing size in networks of cities: hasil dibedakan menurut
  fungsi dan ukuran serta terhubung dengan kontribusi.
\item
  A03 --- Borrowed Size and Cultural Amenities: hasil dan kontribusi
  menjawab tujuan, meski beberapa klaim mekanisme tidak langsung.
\item
  A05 --- Borrowed size in China's urban agglomerations: hasil rinci
  dengan heterogenitas dan kesimpulan selaras, walau robustness
  terbatas.
\item
  A10 --- Knowledge innovation network externalities: hasil, temuan,
  kesimpulan, dan kontribusi rinci serta umumnya konsisten.
\item
  A11 --- Functional spillovers and economic performance: hasil
  konsisten dengan hipotesis dan kontribusi empiris teridentifikasi.
\item
  A12 --- Large cities and small-city growth: hasil, temuan, kesimpulan,
  dan kontribusi konsisten dalam beberapa spesifikasi.
\item
  A17 --- Borrowed size and agglomeration shadows in Japan: hasil,
  robustness, temuan, kesimpulan, dan kontribusi tersedia dengan caveat
  internal.
\item
  A20 --- Urban Green Development Efficiency: hasil threshold dan
  kontribusi terdokumentasi serta umumnya konsisten.
\item
  A21 --- Borrowing Size and Economic Development: hasil tabel, temuan,
  hipotesis, kesimpulan, dan kontribusi tersedia.
\item
  A24 --- Border cities: hasil, pengujian hipotesis, kesimpulan, dan
  kontribusi terhubung, dengan caveat kausal.
\item
  A25 --- Borrowed Size of Small and Medium Cities: hasil, temuan,
  kesimpulan, dan kontribusi lengkap, meski ada konflik prosa-tabel.
\item
  A26 --- Three-layer model of borrowed size: hasil numerik, temuan,
  kesimpulan, dan kontribusi koheren dalam batas deskriptif.
\item
  B01 --- Beyond Polycentricity: hasil dan kontribusi lengkap dengan
  nuansa signifikansi serta batas interpretasi.
\item
  B05 --- Polycentric Urban Systems in the OECD: hasil, temuan,
  kesimpulan, dan kontribusi konsisten menurut skala.
\item
  B06 --- Productivity effects of polycentricity: hasil, temuan,
  kesimpulan, dan kontribusi terhubung, dengan satu ambiguitas
  interpretasi.
\item
  B07 --- Summing Small Cities: hasil dan temuan selaras dengan
  analisis, meski ada masalah tanda signifikansi.
\item
  B10 --- Urban Polycentricity and Commuting Costs: hasil dan kontribusi
  rinci, dengan perbedaan EIM-ACT tetap dikualifikasi.
\item
  B13 --- Spatial structure and productivity in US metropolitan areas:
  hasil dan kontribusi kuat, dengan ambiguitas terbatas pada diagnostik
  model.
\item
  B15 --- Spatial structure and productivity in Europe: hasil, temuan,
  kesimpulan, dan kontribusi koheren.
\item
  B18 --- Big in the neighbourhood: hasil, temuan, kesimpulan, dan
  kontribusi saling terhubung.
\item
  B19 --- Territorial Arrangements of Towns: hasil, temuan, kesimpulan,
  dan kontribusi terdokumentasi kuat.
\item
  B23 --- Cities, hinterlands and agglomeration shadows: hasil numerik
  lengkap dan kesimpulan menjawab tujuan.
\item
  B24 --- Agglomeration Shadows from Ancient Ports: hasil utama terukur
  dan selaras dengan temuan serta kontribusi.
\item
  C02 --- Urban Cannibalism: hasil numerik, pola sektoral, kesimpulan,
  dan kontribusi konsisten terhadap indikator.
\end{itemize}

\subsubsection{B}\label{b-3}

\begin{itemize}
\tightlist
\item
  J01 --- Identifying post-suburbanization: hasil, temuan, dan
  kontribusi tersedia, tetapi terdapat konflik tabel-narasi dan klaim
  kurang didukung.
\item
  J02 --- Pola Pergerakan dan Dekonsentrasi Pekerjaan: distribusi pola
  dan asosiasi konkret, tetapi klaim dekonsentrasi melampaui bukti
  terukur.
\item
  J04 --- Beyond property: hasil dan kontribusi cukup kaya, tetapi klaim
  kausal, lingkungan, dan proyek melampaui bukti.
\item
  J08 --- Socio-Economic Inequality Among Commuters: hasil tersedia,
  tetapi interpretasi odds ratio, arah koefisien, dan klaim moda
  bermasalah.
\item
  J10 --- Job accessibility evaluation: hasil numerik dan kesimpulan
  distribusional cukup konsisten, tetapi pekerjaan dan perjalanan
  tercampur.
\item
  J23 --- Urban transformation and inter-suburban commuting: hasil dan
  kontribusi cukup lengkap, tetapi arah arus, koefisien, dan klaim
  pekerjaan bermasalah.
\item
  J25 --- Mega-urbanization in the Jakarta-Bandung Region: sintesis
  empiris kaya dan koheren, tetapi percepatan serta dampak kebijakan
  tidak diuji langsung.
\item
  J27 --- Peri-urban transformation: hasil dan kontribusi dualitas
  pembangunan cukup koheren, tetapi klaim mekanisme melampaui bukti.
\item
  A04 --- Geography of Borrowing Size: hasil menjawab pertanyaan, tetapi
  statistik dan klaim mekanisme melampaui bukti residual.
\item
  A06 --- Service Function of Small Towns: hasil rinci, tetapi
  kesimpulan borrowed size dan arah aksesibilitas masih problematis.
\item
  A07 --- Urban borrowed size in Guangzhou-Foshan: hasil cukup jelas,
  tetapi pola spasial ditafsirkan terlalu jauh sebagai mekanisme dan
  kausalitas.
\item
  A08 --- Regional Integration Strategies in YRD: hasil sektoral jelas,
  tetapi LQ dan beberapa klaim kausal tidak langsung mengukur arus.
\item
  A09 --- Borrowed administrative power and high-speed rail: hasil
  dengan heterogenitas tersedia, tetapi robustness dan mekanisme tidak
  lengkap.
\item
  A13 --- Industrial Composition and Agglomeration Shadow: hasil rinci,
  tetapi bentuk inverse U tidak stabil dan mekanisme tidak langsung.
\item
  A15 --- Inter-regional networks and productive efficiency: hasil cukup
  lengkap, tetapi outcome merupakan proksi dan klaim jaringan tidak
  kausal.
\item
  A19 --- Population agglomeration and urban economic growth: hasil
  rinci, tetapi GDP kurang robust dan klaim shadow tidak langsung.
\item
  A22 --- Small Towns in South Moravia: hasil cukup rinci dan konsisten
  dengan asosiasi, tetapi klaim shadow lebih interpretatif.
\item
  A28 --- Small Firms and the Urban-Rural Shift: hasil dan kesimpulan
  eksploratori rinci, tetapi tidak mengestimasi efek borrowed size.
\item
  B08 --- Secondary Yet Metropolitan: hasil, temuan, kesimpulan, dan
  kontribusi tersedia, tetapi median dan klaim politik belum sepenuhnya
  terverifikasi.
\item
  B20 --- Grid and shake: hasil rinci, tetapi angka narasi-tabel tidak
  konsisten dan statistik KS tidak tersedia.
\item
  B25 --- Urban Shadow Areas from People Flows: hasil konkret dan cukup
  selaras, tetapi klasifikasi serta mekanisme kualitatif berkonflik.
\end{itemize}

\subsubsection{C}\label{c-3}

\begin{itemize}
\tightlist
\item
  J09 --- Motorcycle-based ride-hailing: hasil deskriptif dan regresi
  rinci, tetapi konflik angka-narasi dan klaim asosiasional yang
  melampaui desain serius.
\end{itemize}

\subsubsection{D}\label{d-3}

\begin{itemize}
\tightlist
\item
  B02 --- Form Follows Function: arah temuan utama bertentangan antara
  abstrak, hasil, tabel, kesimpulan, dan subbagian artikel.
\end{itemize}

\subsection{5. Transferabilitas}\label{transferabilitas}

\subsubsection{S}\label{s-4}

\begin{itemize}
\tightlist
\item
  B08 --- Secondary Yet Metropolitan: implikasi sangat konkret dan
  aplikatif, mencakup infrastruktur, koordinasi, tata kelola, dan
  kepemimpinan.
\end{itemize}

\subsubsection{A}\label{a-4}

\begin{itemize}
\tightlist
\item
  J05 --- The Privatization of Metropolitan Jakarta's Urban Fringes:
  implikasi perencanaan konkret dan relevan, meski instrumen
  implementasi belum dirinci.
\item
  J11 --- Job Access and Spatial Equity: implikasi perencanaan dan
  aplikasi evaluasi jalan tol dijelaskan dengan aktor dan tindakan yang
  cukup jelas.
\item
  J25 --- Mega-urbanization in the Jakarta-Bandung Region: implikasi
  tata kelola, pengendalian lahan, kerja sama lintas batas, dan
  transportasi cukup konkret.
\item
  J26 --- The Dynamics of Jabotabek Development: implikasi kebijakan dan
  aplikasi analitis operasional, dengan batas transfer yang dinyatakan.
\item
  A01 --- Stretching the Concept of Borrowed Size: implikasi kebijakan
  dan aplikasi analitis tentang heterogenitas kota serta jaringan
  dijelaskan jelas.
\item
  A17 --- Borrowed size and agglomeration shadows in Japan: implikasi
  sektoral konkret dan terkait temuan, meski implementasi aktual tidak
  dilaporkan.
\item
  A18 --- Economic Efficiency of Small and Medium-Sized Cities:
  rekomendasi transportasi, koordinasi, inovasi, keuangan, dan tenaga
  kerja spesifik serta relevan.
\item
  A21 --- Borrowing Size and Economic Development: rekomendasi
  aksesibilitas, tata letak industri, dan talenta cukup konkret.
\item
  A23 --- Shedding Light or Casting Shadows: implikasi kebijakan dan
  penerapan kerangka tersedia, meski masih umum dan berbasis satu kasus.
\item
  A26 --- Three-layer model of borrowed size: empat pilar kebijakan dan
  contoh penerapan dirumuskan konkret.
\item
  B04 --- Economic Foundation of the Urban Network Paradigm: implikasi
  kebijakan konkret dan prosedur dapat diterapkan, meski belum diuji di
  luar konteks.
\item
  B11 --- The Rise of Second-Rank Cities: implikasi untuk fungsi urban
  dan kerja sama antarkota tersedia, meski rekomendasi umum.
\item
  B12 --- City Size and Economic Performance: implikasi kebijakan
  spesifik dan aplikatif, tanpa evaluasi implementasi.
\item
  B16 --- Urban Polycentricity and Regional Economic Growth: implikasi
  kebijakan konkret dan kontekstual, tetapi terbatas pada UR China.
\item
  B18 --- Big in the neighbourhood: implikasi untuk cohesion,
  perencanaan, zoning, dan regulasi lalu lintas cukup konkret.
\item
  B20 --- Grid and shake: implikasi praktis jelas dan metode diterapkan
  pada dua model empiris.
\end{itemize}

\subsubsection{B}\label{b-4}

\begin{itemize}
\tightlist
\item
  J01 --- Identifying post-suburbanization: implikasi tata kelola dan
  transportasi jelas, tetapi protokol implementasi dan evaluasi tidak
  tersedia.
\item
  J03 --- Industrial Land Development and Manufacturing Deconcentration:
  implikasi tata kelola dan perencanaan relevan, tetapi aplikasi
  spesifik belum diuji.
\item
  J04 --- Beyond property: implikasi diagnosis dan kebijakan ada, tetapi
  pemantauan serta transfer terbatas oleh kasus purposif.
\item
  J06 --- Transportation Network and Urban Structure: relevansi
  kebijakan terlihat, tetapi aplikasi lintas kota dan indikator
  implementasi tidak diuji.
\item
  J07 --- Highway expansion and urban sprawl: rekomendasi transportasi
  cukup konkret, tetapi protokol operasional dan dampak tidak diuji.
\item
  J08 --- Socio-Economic Inequality Among Commuters: rekomendasi
  transportasi dan tata kelola ada, tetapi sebagian besar merupakan
  inferensi asosiasional.
\item
  J10 --- Job accessibility evaluation: rekomendasi transit dan
  konektivitas cukup konkret, tetapi kelayakan dan dampak tidak diuji.
\item
  J22 --- Morphological and Functional Polycentricity in Indonesia:
  implikasi perencanaan relevan, tetapi aplikasi dan dampak kebijakan
  belum diuji.
\item
  J23 --- Urban transformation and inter-suburban commuting: implikasi
  industri, lahan, pekerjaan, dan transportasi tersedia, tetapi transfer
  memerlukan adaptasi konteks.
\item
  J27 --- Peri-urban transformation: relevansi perencanaan dan alat
  diagnosis terlihat, tetapi implikasi masih umum.
\item
  A02 --- Borrowing size in networks of cities: implikasi menurut fungsi
  dan skala berguna, tetapi aplikasi masih diagnostik-konseptual.
\item
  A04 --- Geography of Borrowing Size: tipologi dan implikasi
  perencanaan cukup konkret, tetapi tidak ada evaluasi kebijakan.
\item
  A05 --- Borrowed size in China's urban agglomerations: rekomendasi
  kebijakan konkret, tetapi aplikasi hanya diusulkan sebagai diagnosis.
\item
  A07 --- Urban borrowed size in Guangzhou-Foshan: implikasi
  diferensiasi town dan jaringan cukup berguna, tetapi instrumen dan
  evaluasi tidak tersedia.
\item
  A08 --- Regional Integration Strategies in YRD: rekomendasi kebijakan
  tersedia, tetapi implementasi dan transfer di luar YRD tidak diuji.
\item
  A09 --- Borrowed administrative power and high-speed rail: implikasi
  lahan, layanan, sektor, dan tata kelola relevan, tetapi aplikasi luar
  YRD belum diuji.
\item
  A10 --- Knowledge innovation network externalities: implikasi GBA
  konkret, tetapi masih berupa rekomendasi tanpa evaluasi implementasi.
\item
  A12 --- Large cities and small-city growth: implikasi kontekstual
  cukup jelas, tetapi aplikasi kebijakan dan generalisasi terbatas.
\item
  A13 --- Industrial Composition and Agglomeration Shadow: implikasi
  ketimpangan dan diagnosis dirumuskan hati-hati, tetapi implementasi
  tidak dilaporkan.
\item
  A15 --- Inter-regional networks and productive efficiency: aplikasi
  analitis konkret, tetapi evaluasi proyek dan implementasi kebijakan
  tidak tersedia.
\item
  A16 --- Regional network economies and productivity: rekomendasi
  sektoral relevan, tetapi umum dan tanpa evaluasi aktual.
\item
  A19 --- Population agglomeration and urban economic growth:
  rekomendasi kebijakan relevan, tetapi aplikasi lintas konteks belum
  diuji.
\item
  A25 --- Borrowed Size of Small and Medium Cities: implikasi kebijakan
  konkret, tetapi applications terutama diagnosis komparatif.
\item
  A27 --- Secondary city regionalism: implikasi menurut tipe proyek
  cukup konkret, tetapi keberhasilan dan transfer luar kasus tidak
  ditunjukkan.
\item
  B01 --- Beyond Polycentricity: implikasi dan aplikasi jelas, tetapi
  masih berbasis asosiasi dan belum dievaluasi.
\item
  B02 --- Form Follows Function: berguna sebagai diagnosis
  polisentrisitas, tetapi tidak memberi instrumen atau ambang kebijakan
  spesifik.
\item
  B03 --- Functional Polycentrism and Urban Network Development:
  rekomendasi transportasi dan strategi regional ada, tetapi
  implementasi tidak diuji.
\item
  B05 --- Polycentric Urban Systems in the OECD: benchmarking dan
  implikasi berskala tersedia, tetapi panduan operasional belum ada.
\item
  B07 --- Summing Small Cities: koordinasi regional dan pembiayaan
  relevan, tetapi bukti dampak kebijakan tidak tersedia.
\item
  B09 --- Identification of sub-centres: aplikasi kebijakan dan metode
  ada, tetapi transfer masih terbatas pada Roma dan Milan.
\item
  B10 --- Urban Polycentricity and Commuting Costs: implikasi
  perencanaan cukup konkret, tetapi konteks empiris terbatas pada
  Italia.
\item
  B17 --- Agglomeration or Network Externalities: implikasi perencanaan
  tersedia, tetapi hati-hati dan umum tanpa intervensi lapangan.
\item
  B19 --- Territorial Arrangements of Towns: aplikasi pemetaan dan
  perbandingan jelas, tetapi kebijakan operasional belum tersedia.
\item
  B21 --- NEG agglomeration shadows and population dynamics: aplikasi
  diagnostik tersedia, tetapi instrumen dan evaluasi lintas konteks
  tidak ada.
\item
  B24 --- Agglomeration Shadows from Ancient Ports: evaluasi spasial dan
  koordinasi antarkota relevan, tetapi outcome kesejahteraan serta
  kebijakan tidak diuji.
\item
  B25 --- Urban Shadow Areas from People Flows: alat diagnosis awal
  berguna, tetapi efektivitas dan transfer di luar Nanjing belum
  tervalidasi.
\item
  C02 --- Urban Cannibalism: penyaringan awal kebijakan dan penggunaan
  indeks cukup jelas, tetapi masih normatif tanpa evaluasi dampak.
\end{itemize}

\subsubsection{C}\label{c-4}

\begin{itemize}
\tightlist
\item
  J02 --- Pola Pergerakan dan Dekonsentrasi Pekerjaan: implikasi
  infrastruktur, hunian, dan transportasi disebut, tetapi instrumen
  serta indikator keberhasilan tidak ada.
\item
  J09 --- Motorcycle-based ride-hailing: implikasi koordinasi
  metropolitan tersedia, tetapi transfer dibatasi dua wilayah dan tidak
  ada kebijakan yang diuji.
\item
  A03 --- Borrowed Size and Cultural Amenities: aplikasi diagnostik ada,
  tetapi rekomendasi formal dan konteks fungsi lain tidak diuji.
\item
  A06 --- Service Function of Small Towns: penyaringan fasilitas
  berguna, tetapi belum menjadi instrumen kebijakan operasional.
\item
  A11 --- Functional spillovers and economic performance: implikasi
  praktis ada, tetapi aplikasi masih berupa diagnosis dan transfer belum
  tervalidasi.
\item
  A14 --- Broadband internet and borrowed size: implikasi ada, tetapi
  aplikasi empiris dan evaluasi implementasi tidak dilakukan.
\item
  A20 --- Urban Green Development Efficiency: rekomendasi konkret,
  tetapi tidak ada evaluasi implementasi atau aplikasi empiris.
\item
  A22 --- Small Towns in South Moravia: aplikasi diagnostik jelas,
  tetapi hanya satu region dan tanpa evaluasi intervensi.
\item
  A24 --- Border cities: implikasi ada, tetapi umum dan tanpa instrumen
  atau evaluasi program.
\item
  A28 --- Small Firms and the Urban-Rural Shift: implikasi tersedia,
  tetapi penerapan lintas wilayah, sektor, dan negara memerlukan uji
  ulang.
\item
  B06 --- Productivity effects of polycentricity: transfer dan kebijakan
  terutama bersifat diagnosis komparatif.
\item
  B13 --- Spatial structure and productivity in US metropolitan areas:
  implikasi praktis umum tanpa aplikasi implementatif atau evaluasi
  kebijakan.
\item
  B15 --- Spatial structure and productivity in Europe: implikasi
  terbatas pada agenda UE tanpa prosedur implementasi.
\item
  B23 --- Cities, hinterlands and agglomeration shadows: aplikasi
  tersirat dan terutama inferensi review, tanpa rekomendasi operasional.
\end{itemize}

\subsubsection{D}\label{d-4}

\begin{itemize}
\tightlist
\item
  B14 --- Polycentric Spatial Structure and Its Economic Performance:
  implikasi praktis umum dan tidak ada aplikasi atau evaluasi wilayah
  tertentu yang dilaporkan.
\item
  B22 --- Urban growth shadows: tidak ada rekomendasi kebijakan atau
  aplikasi eksplisit; implikasi praktis hanya inferensi review.
\end{itemize}

\subsection{6. Refleksivitas}\label{refleksivitas}

\subsubsection{S}\label{s-5}

\begin{itemize}
\tightlist
\item
  J11 --- Job Access and Spatial Equity: keterbatasan, tantangan, dan
  tiga agenda penelitian lanjutan dijelaskan rinci.
\item
  J25 --- Mega-urbanization in the Jakarta-Bandung Region: keterbatasan
  serta tantangan diakui dan agenda komparatif tersedia.
\item
  J26 --- The Dynamics of Jabotabek Development: batas desain, data,
  OCR, dan perbedaan inferensi dipaparkan transparan.
\item
  A01 --- Stretching the Concept of Borrowed Size: keterbatasan dan
  tantangan luas serta agenda lanjutan eksplisit.
\item
  A02 --- Borrowing size in networks of cities: isu dinamika,
  kausalitas, pengukuran, dan agenda lanjutan dibahas jelas.
\item
  A06 --- Service Function of Small Towns: keterbatasan, data hilang,
  validasi, dan agenda panel maupun spasial sangat lengkap.
\item
  A08 --- Regional Integration Strategies in YRD: keterbatasan,
  generalisasi, tantangan, dan agenda riset lanjutan diidentifikasi
  luas.
\item
  A11 --- Functional spillovers and economic performance: keterbatasan,
  tantangan, generalisasi, dan agenda lanjutan eksplisit.
\item
  A12 --- Large cities and small-city growth: batas kausalitas, data,
  mekanisme, arah interaksi, dan agenda lanjutan dibahas luas.
\item
  A14 --- Broadband internet and borrowed size: endogenitas, pengukuran,
  generalisasi, dan agenda riset dibahas rinci.
\item
  A17 --- Borrowed size and agglomeration shadows in Japan:
  keterbatasan, tantangan model/data, dan agenda lanjutan beragam serta
  eksplisit.
\item
  A21 --- Borrowing Size and Economic Development: keterbatasan,
  pengolahan data, dan empat arah riset lanjutan dinyatakan jelas.
\item
  A23 --- Shedding Light or Casting Shadows: batas kausalitas,
  delimitasi, ukuran, indikator, dan agenda lanjutan dibahas lengkap.
\item
  A24 --- Border cities: keterbatasan pengukuran, validitas eksternal,
  dan agenda riset komprehensif.
\item
  A26 --- Three-layer model of borrowed size: keterbatasan, tantangan,
  dan agenda lanjutan dinyatakan eksplisit.
\item
  A27 --- Secondary city regionalism: batasan, tantangan, dan agenda
  riset lanjutan dibahas substantif.
\item
  A28 --- Small Firms and the Urban-Rural Shift: keterbatasan dan
  tantangan luas, dengan agenda lanjutan spesifik.
\item
  B01 --- Beyond Polycentricity: keterbatasan, tantangan, dan agenda
  riset dibahas luas serta spesifik.
\item
  B03 --- Functional Polycentrism and Urban Network Development:
  keterbatasan pengukuran, data, tren, jarak, dan definisi wilayah
  dijelaskan.
\item
  B04 --- Economic Foundation of the Urban Network Paradigm:
  keterbatasan, tantangan, dan agenda komparatif-longitudinal
  terdokumentasi jelas.
\item
  B05 --- Polycentric Urban Systems in the OECD: keterbatasan,
  tantangan, dan future research langsung terkait desain.
\item
  B06 --- Productivity effects of polycentricity: keterbatasan, asumsi,
  data, generalisasi, dan agenda lanjutan eksplisit.
\item
  B07 --- Summing Small Cities: keterbatasan dan challenges eksplisit
  serta agenda future research langsung terkait desain.
\item
  B09 --- Identification of sub-centres: keterbatasan dan tantangan
  metodologis serta agenda lanjutan memadai.
\item
  B11 --- The Rise of Second-Rank Cities: keterbatasan dan challenges
  luas serta agenda urban polisentris dan diakronis tersedia.
\item
  B12 --- City Size and Economic Performance: keterbatasan tipologi dan
  bukti serta tantangan tata kelola dibahas.
\item
  B13 --- Spatial structure and productivity in US metropolitan areas:
  keterbatasan pengukuran, kausalitas, dan agenda lanjutan konkret.
\item
  B14 --- Polycentric Spatial Structure and Its Economic Performance:
  keterbatasan heterogenitas, ukuran sampel, dan agenda lanjutan
  eksplisit.
\item
  B15 --- Spatial structure and productivity in Europe: keterbatasan
  skala, data, endogenitas, dan kebutuhan panel dibahas luas.
\item
  B16 --- Urban Polycentricity and Regional Economic Growth:
  keterbatasan empiris, identifikasi, data, dan lima agenda lanjutan
  eksplisit.
\item
  B17 --- Agglomeration or Network Externalities: keterbatasan
  pengukuran, sensitivitas, dan agenda riset lanjut dijelaskan
  substantif.
\item
  B18 --- Big in the neighbourhood: keterbatasan data, linkage,
  populasi, transportasi, batas, dan kausalitas sangat spesifik.
\item
  B19 --- Territorial Arrangements of Towns: keterbatasan dan tantangan
  dibahas luas dengan agenda lanjutan terarah.
\item
  B20 --- Grid and shake: keterbatasan dan tantangan dijelaskan luas,
  meski agenda future research tidak eksplisit.
\item
  B22 --- Urban growth shadows: keterbatasan kausalitas, delineasi,
  jarak, data, dan model dibahas rinci.
\item
  B23 --- Cities, hinterlands and agglomeration shadows: keterbatasan
  rekonstruksi wilayah, periodisasi, pusat, dan lag dijelaskan.
\item
  B24 --- Agglomeration Shadows from Ancient Ports: keterbatasan
  identifikasi, data, model, cakupan, dan sensitivitas dijelaskan sangat
  lengkap.
\item
  C02 --- Urban Cannibalism: keterbatasan dan agenda validasi, data
  arus, kalibrasi, MAUP, dan lintas metropolitan rinci.
\end{itemize}

\subsubsection{A}\label{a-5}

\begin{itemize}
\tightlist
\item
  J02 --- Pola Pergerakan dan Dekonsentrasi Pekerjaan: keterbatasan
  cakupan, representativitas, pengukuran, dan agenda lanjutan diakui
  jelas.
\item
  J03 --- Industrial Land Development and Manufacturing Deconcentration:
  batas desain, sensus, indikator, dan tantangan tata kelola terlihat
  jelas.
\item
  J06 --- Transportation Network and Urban Structure: keterbatasan IV,
  proksi, periode, sampel, dan agenda lanjutan dipetakan.
\item
  J07 --- Highway expansion and urban sprawl: keterbatasan IV, resolusi,
  generalisasi, dan agenda perilaku dibahas.
\item
  J10 --- Job accessibility evaluation: keterbatasan dan agenda
  longitudinal serta variabel sosial-demografis dijelaskan luas.
\item
  J22 --- Morphological and Functional Polycentricity in Indonesia:
  keterbatasan pengukuran, waktu, spasial, dan agenda lanjut
  teridentifikasi.
\item
  J23 --- Urban transformation and inter-suburban commuting:
  keterbatasan, konflik internal, tantangan data, dan agenda lanjut
  eksplisit.
\item
  A03 --- Borrowed Size and Cultural Amenities: keterbatasan, tantangan,
  dan agenda lanjutan cukup ekstensif.
\item
  A04 --- Geography of Borrowing Size: keterbatasan pengukuran,
  agregasi, delineasi, mekanisme, dan agenda lanjut luas.
\item
  A05 --- Borrowed size in China's urban agglomerations: keterbatasan,
  tantangan, dan agenda penelitian lanjut jelas.
\item
  A09 --- Borrowed administrative power and high-speed rail:
  keterbatasan agregasi, endogenitas, transparansi, dan agenda lanjut
  konkret.
\item
  A13 --- Industrial Composition and Agglomeration Shadow: keterbatasan
  konstruk, cakupan, spesifikasi, spasial, dan validitas eksternal
  dipetakan.
\item
  A16 --- Regional network economies and productivity: keterbatasan
  pengukuran, agregasi, kausalitas, dan reproduksibilitas terlihat
  jelas.
\item
  A18 --- Economic Efficiency of Small and Medium-Sized Cities:
  keterbatasan data, omitted variables, lag, pencilan, dan kausalitas
  teridentifikasi.
\item
  A19 --- Population agglomeration and urban economic growth:
  keterbatasan data rahasia, N, bobot, outcome, dan agenda lanjut
  spesifik.
\item
  A25 --- Borrowed Size of Small and Medium Cities: keterbatasan dan
  challenges substantif serta future research spesifik.
\item
  B02 --- Form Follows Function: keterbatasan konsep, skala, arus, dan
  agenda lanjut dijelaskan cukup lengkap.
\item
  B08 --- Secondary Yet Metropolitan: batas konteks, kelembagaan, dan
  agenda lanjut cukup spesifik.
\item
  B10 --- Urban Polycentricity and Commuting Costs: keterbatasan
  konseptual, data, pengukuran, kausalitas, dan spatial challenge
  dicatat.
\item
  B21 --- NEG agglomeration shadows and population dynamics: asumsi,
  desain, pengukuran, dan agenda lanjut terdokumentasi.
\item
  B25 --- Urban Shadow Areas from People Flows: keterbatasan eksplisit,
  tantangan data, dan dua agenda riset disampaikan.
\end{itemize}

\subsubsection{B}\label{b-5}

\begin{itemize}
\tightlist
\item
  J05 --- The Privatization of Metropolitan Jakarta's Urban Fringes:
  tantangan substantif terlihat, tetapi keterbatasan eksplisit dan
  agenda riset terbatas.
\item
  J08 --- Socio-Economic Inequality Among Commuters: tantangan
  metodologis dapat dipetakan, tetapi keterbatasan tidak dinyatakan
  eksplisit.
\item
  J27 --- Peri-urban transformation: tantangan dan keterbatasan dapat
  dipetakan, tetapi sebagian besar berasal dari inferensi review.
\item
  A07 --- Urban borrowed size in Guangzhou-Foshan: tantangan substantif
  dan agenda lanjutan ada, tetapi limitations formal tidak disediakan.
\item
  A10 --- Knowledge innovation network externalities: tantangan dan
  agenda lanjut jelas, tetapi limitations khusus tidak tersedia.
\item
  A15 --- Inter-regional networks and productive efficiency:
  keterbatasan teridentifikasi, tetapi refleksi formal, robustness, dan
  identification terbatas.
\item
  A20 --- Urban Green Development Efficiency: keterbatasan teori dan
  agenda disebut, tetapi isu sampling dan observasi kurang
  direfleksikan.
\end{itemize}

\subsubsection{C}\label{c-5}

\begin{itemize}
\tightlist
\item
  J01 --- Identifying post-suburbanization: keterbatasan penting
  terutama harus diturunkan dari review, bukan diakui eksplisit penulis.
\item
  J04 --- Beyond property: tantangan dibahas, tetapi keterbatasan
  khusus, komparabilitas, dan validitas kasus tidak dirumuskan
  sistematis.
\item
  J09 --- Motorcycle-based ride-hailing: refleksi tidak sistematis dan
  konflik internal tidak direkonsiliasi.
\item
  A22 --- Small Towns in South Moravia: tantangan ada, tetapi
  keterbatasan tidak eksplisit dan sebagian besar inferensial.
\end{itemize}

\subsubsection{D}\label{d-5}

\begin{itemize}
\tightlist
\item
  Tidak ada artikel pada tier ini.
\end{itemize}

\section{Kumpulan Review}\label{kumpulan-review}

\section{Review J01}\label{review-j01}

\textbf{Identitas sumber}

\begin{itemize}
\tightlist
\item
  Kode: J01.
\item
  Judul: \emph{Identifying post-suburbanization: The case of the Jakarta
  metropolitan area (JMA)}.
\item
  Penulis: Adiwan F. Aritenang.
\item
  Afiliasi: Urban and Regional Planning Program, School of Architecture,
  Planning and Policy Development, Institut Teknologi Bandung,
  Indonesia.
\item
  Jenis sumber: artikel jurnal. Status telaah sejawat tidak disebutkan
  dalam file.
\item
  Jurnal: \emph{Habitat International}.
\item
  Volume, tahun, dan nomor artikel: 138 (2023), 102857. Nomor terbitan
  dan rentang halaman tidak disebutkan dalam file.
\item
  DOI: 10.1016/j.habitatint.2023.102857.
\item
  Riwayat naskah: diterima 8 Januari 2023; revisi diterima 1 Juni 2023;
  disetujui 6 Juni 2023; tersedia daring 15 Juni 2023.
\item
  Kata kunci: post-suburbanisation; commuting; polycentric; Jakarta;
  statistics.
\item
  Panjang sumber: 9 halaman berdasarkan metadata PDF.
\item
  Konflik kepentingan: penulis menyatakan tidak ada (bagian
  \emph{Declaration of competing interest}, hlm. 8).
\end{itemize}

\subsection{Summarized Abstract}\label{summarized-abstract}

\textbf{Laporan sumber.} Artikel menguji proses post-suburbanisasi di
Jakarta Metropolitan Area (JMA) melalui pola perjalanan komuter,
asal-tujuan perjalanan, dan karakteristik sosioekonomi komuter. Data
yang digunakan ialah dua survei komuter potong lintang tahun 2014 dan
2019. Penulis menilai tiga dimensi: keterbatasan ketergantungan spasial
kegiatan kerja dan sekolah pada Jakarta; kematangan suburban yang
dikaitkan dengan daya tarik distrik luar bagi kelompok menengah-atas dan
migran yang lebih tua; serta segregasi yang antara lain dibaca dari
besarnya arus komuter dari distrik yang berdekatan (Abstrak, hlm. 1).

\textbf{Interpretasi penulis.} Pola-pola tersebut dipandang menunjukkan
bahwa sebagian kawasan pinggiran bertransformasi menjadi pusat baru
dalam perkembangan metropolitan yang polisentris, meskipun hubungan
dengan kota inti tetap kuat. Artikel menekankan perlunya tata kelola
multitingkat untuk menangani variasi mobilitas, perkembangan spasial,
dan ketimpangan sosioekonomi JMA (Abstrak, hlm. 1; Bagian 5, hlm. 7-8).

\textbf{Inferensi pengolah.} Post-suburbanisasi tidak diukur sebagai
satu keluaran tunggal, tetapi disimpulkan dari sekumpulan proksi
deskriptif pada tiga dimensi. Karena analisisnya bukan pengujian kausal
atau model klasifikasi formal, kesimpulan utama lebih tepat dibaca
sebagai diagnosis pola yang konsisten secara sebagian dengan konsep
post-suburbanisasi, bukan estimasi sebab-akibat atau penetapan ambang
yang tervalidasi.

\subsection{Problem Statement}\label{problem-statement}

\textbf{Laporan sumber.} Pertumbuhan luas wilayah dan kegiatan
sosioekonomi JMA belum diimbangi penelitian yang secara khusus menilai
proses post-suburbanisasi melalui perjalanan komuter dan implikasi
spasialnya. Penulis menyatakan bahwa kajian mengenai kegiatan
sosioekonomi serta mobilitas spasial komuter masih terbatas. Kajian JMA
sebelumnya lebih banyak menggunakan data industri, distribusi tenaga
kerja dan kegiatan, tata guna lahan, atau pembangunan perumahan skala
besar (Bagian 1, hlm. 1).

\textbf{Interpretasi penulis.} Tanpa melihat arus penduduk antara inti
dan pinggiran, perubahan pinggiran dari kota asrama menjadi pusat
kegiatan baru, ketergantungan yang masih bertahan, serta segregasi
antarkawasan tidak dapat dijelaskan secara memadai (Bagian 1, hlm. 1;
Bagian 5, hlm. 7-8).

\textbf{Inferensi pengolah.} Masalah analitis artikel adalah bagaimana
mengenali post-suburbanisasi dari data mobilitas individu yang
diagregasikan menurut distrik, bukan bagaimana mengestimasi dampak satu
faktor tertentu terhadap satu variabel hasil.

\subsection{Objectives}\label{objectives}

\textbf{Laporan sumber.} Tujuan utama artikel ialah memeriksa proses
post-suburbanisasi melalui tiga pola: kemandirian spasial, kematangan
suburban, dan segregasi. Untuk itu, penulis menggunakan karakteristik
sosioekonomi komuter, mobilitas inti-pinggiran, serta implikasi
spasialnya selama 2014-2019 (Bagian 1, hlm. 1; Bagian 4, hlm. 5).

\textbf{Inferensi pengolah atas operasionalisasi tujuan.} Tujuan
tersebut dijabarkan dalam empat pekerjaan analitis:

\begin{itemize}
\tightlist
\item
  Mengukur ketergantungan pinggiran terhadap inti melalui distribusi
  asal-tujuan dan tujuan kegiatan kerja atau sekolah (Bagian 4.1, hlm.
  5-6; Tabel 2-3).
\item
  Menilai kematangan dan daya tarik kawasan suburban melalui umur
  komuter, migrasi lima tahun terakhir, serta proporsi komuter muda dan
  kelas menengah-atas (Bagian 4.2, hlm. 6-7; Tabel 4-5).
\item
  Menilai segregasi sosioekonomi melalui upah, kondisi hunian layak,
  status pekerjaan formal, dan moda transportasi (Bagian 4.3, hlm. 6-7;
  Tabel 6-7).
\item
  Menarik implikasi tata kelola metropolitan dari pola mobilitas dan
  ketimpangan antardistrik (Bagian 5, hlm. 8).
\end{itemize}

\subsection{Literature Survey}\label{literature-survey}

\textbf{Laporan sumber.} Tinjauan pustaka mendefinisikan
post-suburbanisasi sebagai transformasi pinggiran akibat perluasan
infrastruktur, kemajuan teknologi, keterhubungan kota utama dengan kota
yang lebih kecil, desentralisasi pekerjaan, dan pengaruh globalisasi.
Transformasi itu melemahkan sentralitas kota inti dan menghasilkan
kawasan suburban multipusat yang batas urban-suburbannya makin kabur.
Istilah seperti \emph{edge city}, \emph{edgeless city}, dan
\emph{technoburb} ditempatkan sebagai bagian dari diskusi post-suburban
(Bagian 2, hlm. 2).

Pada dimensi sosial-spasial, pustaka yang dirujuk menggambarkan
kematangan suburban melalui penuaan perumahan, penurunan daya tarik bagi
kelompok menengah-atas, pertumbuhan kawasan eksurban, kemungkinan
revitalisasi selektif atau gentrifikasi, spesialisasi ekonomi,
fragmentasi struktural, dan segregasi sosial. Perjalanan penduduk dalam
kondisi tersebut tidak lagi hanya konsentris atau radial, tetapi
cenderung polisentris; pinggiran juga melonggarkan ketergantungannya
pada kota pusat (Bagian 2, hlm. 2).

Artikel membandingkan pengalaman Amerika Serikat, Kanada, Prancis, dan
Tiongkok. Perbandingan tersebut menyoroti peran jalan raya,
desentralisasi, kewenangan pemerintah lokal, industrialisasi suburban,
perumahan skala besar, pusat perbelanjaan, serta perbedaan definisi
ruang urban dan pola permukiman. Penulis menyatakan bahwa penerapan
konsep dari negara Barat tidak dapat dipindahkan begitu saja ke JMA
(Bagian 2, hlm. 2; Bagian 5, hlm. 8).

Untuk JMA, artikel merangkum literatur tentang dekonsentrasi industri,
real estat dan kota baru yang dipimpin sektor swasta, urbanisasi
periurban, investasi asing, privatisasi pengembangan lahan, tata kelola
metropolitan, banjir, serta segregasi sosial dan pekerjaan. Kota-kota
baru sejak 1980-an terutama ditujukan kepada kelompok menengah-atas dan
kerap berdekatan dengan kawasan industri. Penelitian terdahulu juga
menunjukkan bahwa JMA memiliki struktur ketenagakerjaan polisentris,
tetapi tetap memperlihatkan campuran pola suburban tradisional dan baru
(Bagian 3, hlm. 3-5; Gambar 3-4).

\textbf{Interpretasi penulis.} Literatur tersebut digunakan untuk
menurunkan tiga indikator analitis: kemandirian spasial, kematangan, dan
segregasi. Kondisi JMA dibaca sebagai bentuk post-suburbanisasi yang
digerakkan kuat oleh dekonsentrasi industri dan pengembangan lahan
swasta, bukan semata-mata hasil desain pemerintah (Bagian 3, hlm. 4-5).

\textbf{Inferensi pengolah.} Bagian ini merupakan tinjauan naratif,
bukan tinjauan sistematis. File tidak menyebutkan strategi pencarian,
pangkalan data, kriteria inklusi, periode penelusuran, ataupun jumlah
karya yang disaring.

\subsection{Method Used}\label{method-used}

\textbf{Laporan sumber.} Penelitian menggunakan analisis deskriptif dan
spasial atas dua survei komuter potong lintang. Survei diselenggarakan
oleh Central Agency on Statistics atau Badan Pusat Statistik (BPS)
dengan kuesioner dan bobot survei. Strategi sampelnya adalah
stratifikasi spasial dengan pengambilan sampel dua tahap pada tingkat
kecamatan, mempertimbangkan distribusi rumah tangga dan penduduk
berdasarkan Sensus Penduduk 2010 (Bagian 4, hlm. 5).

Analisis membagi wilayah menjadi lima kota administratif Jakarta sebagai
inti, lima wilayah pinggiran yang berbatasan langsung dengan inti, dan
tiga wilayah pinggiran yang tidak berbatasan langsung. Penulis kemudian
membandingkan agregat asal-tujuan, tujuan perjalanan, umur, status
migrasi, kelas sosial-ekonomi, upah, kondisi hunian, status pekerjaan,
dan moda perjalanan menurut distrik dan, bila tersedia, menurut tahun
(Bagian 3-4, hlm. 2 dan 5-7; Gambar 1-2; Tabel 2-7).

Tiga langkah substantifnya adalah:

\begin{itemize}
\tightlist
\item
  Kemandirian spasial dinilai dari proporsi perjalanan inti-ke-inti,
  pinggiran-ke-inti, inti-ke-pinggiran, dan antarpinggiran, lalu
  diperinci menurut perjalanan kerja dan sekolah (Bagian 4.1, hlm. 5-6).
\item
  Kematangan dinilai dari umur rata-rata komuter, kedatangan migran
  dalam lima tahun, dan komposisi komuter muda serta menengah-atas
  sebagai proksi daya tarik dan gentrifikasi (Bagian 4.2, hlm. 6-7).
\item
  Segregasi dinilai dari perbedaan upah, hunian layak, pekerjaan formal,
  dan penggunaan moda transportasi (Bagian 4.3, hlm. 6-7).
\end{itemize}

\textbf{Inferensi pengolah.} File tidak memperlihatkan regresi,
pengujian hipotesis inferensial, estimasi efek, interval kepercayaan,
atau uji signifikansi. Istilah ``analisis spasial'' dalam artikel
diwujudkan terutama melalui klasifikasi lokasi dan tabulasi arus
antardistrik; prosedur statistik spasial formal tidak dijelaskan.

\subsection{Dataset}\label{dataset}

\textbf{Laporan sumber.} Dataset utama adalah JMA Commuting Surveys
tahun 2014 dan 2019. Kedua survei bersifat potong lintang, bukan panel
individu. Informasinya meliputi distrik asal dan tujuan, kegiatan
perjalanan seperti bekerja, bersekolah, dan berbelanja, waktu
perjalanan, jam perjalanan, serta karakteristik sosioekonomi komuter
(Bagian 4, hlm. 5).

Artikel juga menampilkan data kontekstual pertumbuhan Produk Domestik
Regional Bruto harga konstan 2010 untuk 2011-2014, koefisien Gini PDRB,
dan spesialisasi ekonomi distrik (Bagian 3, hlm. 3-4; Tabel 1; Gambar
3-4). Pada pembahasan segregasi, artikel mengutip angka kemiskinan BPS
tahun 2016 untuk Depok, Tangerang Selatan, dan Kota Tangerang, tetapi
angka itu bukan keluaran langsung survei komuter (Bagian 4.3, hlm. 6).

\textbf{Inferensi pengolah.} File tidak menyebutkan nama berkas data,
tautan akses, lisensi, kamus variabel, prosedur pembersihan, penanganan
data hilang, konstruksi bobot, atau kode analisis. Sumber rinci untuk
sebagian data kontekstual pada Tabel 1 dan Gambar 3-4 juga tidak
dijelaskan dalam keterangan yang tersedia.

\subsection{Population / Sample}\label{population-sample}

\textbf{Laporan sumber.} Populasi sasaran yang dinyatakan adalah komuter
harian berusia lebih dari 15 tahun yang tinggal di JMA. Sampel tahun
2014 berjumlah 12.960 komuter dan, setelah pembobotan, disebut mewakili
2,9 juta individu. Sampel tahun 2019 berjumlah 13.120 komuter dan
disebut mewakili 3,2 juta individu (Bagian 4, hlm. 5).

Cakupan wilayah terdiri atas 13 unit tingkat kabupaten/kota di tiga
provinsi:

\begin{itemize}
\tightlist
\item
  Inti DKI Jakarta: Jakarta Pusat, Jakarta Utara, Jakarta Selatan,
  Jakarta Timur, dan Jakarta Barat.
\item
  Pinggiran Jawa Barat: Kota Depok, Kabupaten Bogor, Kota Bogor,
  Kabupaten Bekasi, dan Kota Bekasi.
\item
  Pinggiran Banten: Kabupaten Tangerang, Kota Tangerang, dan Kota
  Tangerang Selatan (Bagian 3-4, hlm. 2 dan 5; Gambar 1).
\end{itemize}

Untuk analisis, kelompok pinggiran yang berbatasan langsung dengan inti
mencakup Depok, Kota Bekasi, Kabupaten Bekasi, Kota Tangerang, dan
Tangerang Selatan. Kelompok tidak berbatasan mencakup Kota Bogor,
Kabupaten Bogor, dan Kabupaten Tangerang (Bagian 3-4, hlm. 2 dan 5).

\textbf{Inferensi pengolah.} Unit observasi survei adalah individu
komuter, sedangkan banyak hasil disajikan sebagai agregat distrik atau
kelompok inti-pinggiran. Karena populasi hanya mencakup komuter berusia
di atas 15 tahun, hasilnya tidak dengan sendirinya merepresentasikan
seluruh penduduk JMA, penduduk yang tidak melakukan perjalanan rutin,
ataupun semua rumah tangga. File tidak menjelaskan tingkat respons,
eksklusi observasi, ukuran sampel per distrik, atau galat desain.

\subsection{Dependent Variables}\label{dependent-variables}

\textbf{Klarifikasi desain.} Studi ini bukan model dengan satu variabel
dependen. Tidak ada persamaan regresi, variabel hasil tunggal, atau
perlakuan kausal. Konstruk utama, yaitu post-suburbanisasi, tidak
diamati secara langsung dan tidak diberi skor komposit; penulis
mengenalinya secara deskriptif melalui beberapa kelompok indikator
(Bagian 4, hlm. 5-7).

Kelompok hasil atau proksi yang paling setara dengan variabel dependen
adalah:

\begin{itemize}
\tightlist
\item
  \textbf{Kemandirian spasial:} jumlah dan proporsi komuter menurut
  pasangan asal-tujuan inti atau pinggiran; proporsi komuter menuju
  Jakarta untuk bekerja dan bersekolah; serta arus antarpinggiran
  (Bagian 4.1, hlm. 5-6; Tabel 2-3).
\item
  \textbf{Kematangan dan daya tarik suburban:} umur rata-rata komuter;
  jumlah serta proporsi migran yang tiba dalam lima tahun; proporsi
  komuter muda; dan proporsi komuter atau migran kelas menengah-atas
  (Bagian 4.2, hlm. 6-7; Tabel 4-5).
\item
  \textbf{Segregasi sosioekonomi:} upah rata-rata, proporsi hunian
  layak, proporsi pekerja formal, dan distribusi moda transportasi
  (Bagian 4.3, hlm. 6-7; Tabel 6-7).
\end{itemize}

Pengelompokan distrik menjadi inti, pinggiran berbatasan, dan pinggiran
tidak berbatasan, serta tahun survei 2014 atau 2019, berfungsi sebagai
dimensi pembanding, bukan variabel bebas dalam model formal. Angka
kemiskinan yang dikutip dari BPS 2016 berfungsi sebagai konteks
pembanding dan bukan variabel hasil yang dihitung dari survei (Bagian
4.3, hlm. 6).

\textbf{Inferensi pengolah.} Istilah ``muda'', ``kelas menengah-atas'',
``hunian layak'', dan ``pekerja formal'' tidak diberi definisi atau
ambang operasional dalam file. Usia perumahan, yang penting dalam
definisi konseptual kematangan suburban, juga tidak diukur secara
langsung. Karena itu, indikator komuter adalah proksi bagi perubahan
wilayah, bukan ukuran langsung seluruh proses kematangan, gentrifikasi,
atau segregasi kawasan.

\subsection{Results}\label{results}

\textbf{Kemandirian spasial.} Tabel 2 melaporkan 3.341.535 komuter
tertimbang pada 2019. Menurut narasi penulis, 40,2\% dari keseluruhan
arus berasal dari pinggiran dan menuju inti; 7,7\% berasal dari inti dan
menuju pinggiran; serta 26,7\%, atau 893.166 komuter, melakukan
perjalanan antarpinggiran. Depok dan Kota Bekasi bersama-sama menyumbang
588.639 komuter menuju inti, disebut lebih dari 43,8\% arus
pinggiran-ke-inti. Lebih dari 61,6\% komuter dari kedua wilayah Bekasi
menuju Jakarta Timur, sedangkan 53,6\% komuter Depok menuju Jakarta
Selatan. Penulis menafsirkan pola ini sebagai akibat kedekatan geografis
tujuan sekaligus tanda tingginya ketergantungan wilayah berbatasan pada
inti (Bagian 4.1, hlm. 5; Tabel 2).

Jakarta Selatan dan Jakarta Pusat menjadi tujuan utama. Artikel
melaporkan bahwa lebih dari 43\% komuter Depok, Kabupaten Bogor, dan
Tangerang Selatan menjalankan kegiatan hariannya di Jakarta Selatan,
sedangkan lebih dari 32\% komuter Kabupaten dan Kota Bogor menuju
Jakarta Pusat (Bagian 4.1, hlm. 5).

Menurut Tabel 3, proporsi komuter Kota Bekasi yang bekerja di inti turun
dari 24,87\% pada 2014 menjadi 21,72\% pada 2019, sedangkan proporsi
untuk sekolah turun dari 31,76\% menjadi 24,03\%. Di Depok, proporsi
bekerja naik dari 19,85\% menjadi 21,73\% dan proporsi sekolah naik dari
23,88\% menjadi 32,48\%. Pada 2019, proporsi perjalanan kerja dan
sekolah menuju inti umumnya lebih rendah pada tiga wilayah yang tidak
berbatasan langsung, khususnya Kota Bogor dan Kabupaten Tangerang,
daripada pada Depok, Kota Bekasi, Kota Tangerang, dan Tangerang Selatan
(Bagian 4.1, hlm. 6; Tabel 3).

\textbf{Kematangan dan gentrifikasi.} Umur rata-rata komuter lebih
tinggi di Kota Bogor (39,5 tahun) dan Kabupaten Tangerang (40,4 tahun)
daripada di Depok (33,2 tahun) dan Kota Bekasi (34,5 tahun). Kota
Tangerang dan Tangerang Selatan masing-masing mencatat 34,7 dan 35,8
tahun (Bagian 4.2, hlm. 6).

Dalam Tabel 4, tujuan migran lima tahun terakhir terbanyak adalah Depok
sebanyak 19.961 orang, Kabupaten Bogor 14.468 orang, dan Kota Tangerang
13.559 orang. Untuk kelompok menengah-atas, proporsi yang ditampilkan
mencapai 12,22\% di Kabupaten Bogor, 8,29\% di Kota Bogor, dan 7,80\% di
Depok. Kota Bogor hanya mencatat 701 migran baru, sedangkan Tangerang
Selatan tercatat nol pada kolom migran lima tahun terakhir (Bagian 4.2,
hlm. 6-7; Tabel 4).

Tabel 5 menunjukkan bahwa pada 2019 proporsi komuter muda tertinggi
terdapat di Jakarta Timur (81,75\%). Di pinggiran, proporsinya relatif
tinggi di Depok (72,04\%), Kota Tangerang (70,94\%), Tangerang Selatan
(70,06\%), dan Kabupaten Bekasi (70,97\%), serta lebih rendah di Kota
Bogor (49,72\%) dan Kabupaten Tangerang (51,34\%). Sebaliknya, proporsi
kelas menengah-atas pada 2019 tertinggi di Kabupaten Tangerang (46,28\%)
dan Kota Bogor (39,33\%); Kota Tangerang dan Tangerang Selatan
masing-masing mencatat 21,60\% dan 22,04\% (Bagian 4.2, hlm. 6-7; Tabel
5).

\textbf{Segregasi sosioekonomi.} Pada 2019, Tabel 6 mencatat upah
rata-rata tertinggi di antara wilayah pinggiran pada Kota Bogor sebesar
Rp8.398.174, diikuti Tangerang Selatan Rp6.720.784. Proporsi hunian
layak tertinggi terdapat di Kabupaten Bekasi sebesar 88,39\%, diikuti
Kota Bekasi sebesar 83,93\% dan Kabupaten Tangerang sebesar 81,82\%.
Proporsi pekerja formal yang ditampilkan sangat tinggi di semua wilayah
pinggiran, berkisar 93,62-100\% pada 2019. Artikel juga mengaitkan
tingginya pekerjaan formal di Depok, Tangerang Selatan, dan Kota
Tangerang dengan angka kemiskinan masing-masing 2,32\%, 1,68\%, dan
4,91\% yang dikutip dari BPS 2016 (Bagian 4.3, hlm. 6-7; Tabel 6).

\textbf{Moda perjalanan.} Narasi hasil menyatakan bahwa sepeda motor dan
mobil mendominasi moda pinggiran-ke-inti pada 2019, masing-masing 60,0\%
dan 8,8\%. Kereta mencakup 16,9\%, transportasi umum lain 5,8\%,
TransJakarta 3,4\%, ojek daring 2,7\%, taksi daring 2,2\%, dan moda
nonmotor 0,1\%. Penulis juga menyatakan penggunaan kereta dan
TransJakarta meningkat, sedangkan penggunaan sepeda motor dan mobil
menurun, tetapi Tabel 7 dalam hasil ekstraksi hanya memperlihatkan satu
set angka per moda (Bagian 4.3, hlm. 7-8; Tabel 7).

\subsection{Findings}\label{findings}

\textbf{Interpretasi penulis.} Temuan inti artikel adalah bahwa
post-suburbanisasi JMA berlangsung melalui perubahan skala dan pola
komuter. Arus antarpinggiran sebesar 26,7\% ditafsirkan sebagai bukti
munculnya kegiatan terspesialisasi dan pusat-pusat baru di pinggiran.
Akan tetapi, arus pinggiran-ke-inti sebesar 40,2\%, termasuk konsentrasi
perjalanan Bekasi ke Jakarta Timur dan Depok ke Jakarta Selatan,
menunjukkan bahwa keterikatan terhadap inti belum hilang (Bagian 4.1,
hlm. 5; Bagian 5, hlm. 8).

Penulis menafsirkan perbedaan usia dan kelas sebagai dua bentuk
gentrifikasi yang berlokasi berbeda: wilayah berbatasan lebih menarik
bagi komuter muda, sedangkan wilayah tidak berbatasan seperti Kota Bogor
dan Kabupaten Tangerang lebih menarik bagi kelas menengah-atas.
Pinggiran secara umum dipandang semakin matang dan menjadi alternatif
tempat tinggal bagi penduduk JMA (Bagian 4.2, hlm. 6-7; Bagian 5, hlm.
8).

Perbedaan upah, hunian, pekerjaan formal, dan moda perjalanan
ditafsirkan sebagai segregasi sosioekonomi. Penulis menghubungkannya
dengan pembangunan kota baru dan kawasan industri oleh sektor swasta,
gentrifikasi, kedekatan ke inti, serta penyediaan infrastruktur yang
tidak seimbang. Wilayah berbatasan dinilai memiliki lebih banyak pilihan
moda, sedangkan komuter berpendapatan lebih rendah dari wilayah tidak
berbatasan disebut lebih bergantung pada bus dan kereta (Bagian 4.3,
hlm. 6-7).

Penulis juga memandang ketersediaan jalan raya dan angkutan umum sebagai
faktor penting dalam arah migrasi dan perjalanan. Rendahnya arus dari
wilayah Banten ke inti dikaitkan dengan keterbatasan kereta dan
transportasi umum, sementara hubungan wilayah Jawa Barat dengan inti
didukung kereta komuter (Bagian 4.2-4.3, hlm. 6-8).

\textbf{Inferensi pengolah.} Bukti paling langsung untuk
post-suburbanisasi dalam file adalah keberadaan arus antarpinggiran yang
besar bersamaan dengan arus kuat menuju inti. Pola tersebut mendukung
diagnosis struktur campuran atau polisentris yang belum sepenuhnya
mandiri. Sebaliknya, klaim tentang gentrifikasi, kematangan properti,
pengaruh infrastruktur, dan peran pembangunan swasta bersandar pada
proksi dan perbandingan deskriptif; data yang ditampilkan tidak
mengisolasi mekanisme penyebabnya.

\textbf{Ketegangan bukti internal.} Penulis menyebut 40,2\% sedikit di
atas ambang 35\% untuk struktur yang dominan monosentris, tetapi
kemudian menyatakan angka itu dapat mengindikasikan kemandirian
pinggiran yang meningkat. Arah hubungan antara ambang tersebut dan
kesimpulan kemandirian tidak dijelaskan secara tuntas (Bagian 5, hlm.
8). Selain itu, narasi bahwa proporsi komuter muda di pinggiran lebih
tinggi daripada inti tidak selaras dengan angka Tabel 5, yang
menunjukkan proporsi muda di semua wilayah inti sekitar 74,61-81,75\%
pada 2019, lebih tinggi daripada angka wilayah pinggiran yang
ditampilkan (Bagian 4.2, hlm. 6-7; Tabel 5).

\subsection{Conclusions}\label{conclusions}

\textbf{Kesimpulan penulis.} JMA memperlihatkan post-suburbanisasi yang
tidak sepenuhnya menyerupai pola kota Barat. Wilayah pinggiran
berkembang dalam jumlah penduduk dan ekonomi, memperoleh campuran guna
lahan serta berbagai kegiatan urban, dan sebagian berfungsi sebagai
pusat baru. Namun, Jakarta tetap menjadi pusat utama kegiatan ekonomi
dan sosial, dan banyak arus komuter pinggiran tetap menuju wilayah inti
terdekat (Bagian 5, hlm. 7-8).

Penulis menyimpulkan bahwa perkembangan polisentris dapat mengurangi
waktu perjalanan dan kemacetan, tetapi perlu dikelola melalui kapasitas
tata kelola lintas pemerintah. Gentrifikasi, revitalisasi, pembangunan
baru, serta konektivitas antara pusat-pusat baru perlu dikendalikan agar
pembangunan JMA lebih inklusif dan berkelanjutan. Kurangnya perencanaan
spasial dan transportasi umum yang efektif dinilai berpotensi
memperdalam pemisahan kawasan pinggiran menjadi real estat besar dan
komunitas berpagar (Bagian 5, hlm. 8).

\textbf{Inferensi pengolah.} Kesimpulan yang paling terdukung oleh angka
yang ditampilkan adalah koeksistensi sentralitas Jakarta dan pertumbuhan
interaksi antarpinggiran. Klaim bahwa polisentrisitas telah mengurangi
waktu perjalanan atau kemacetan lebih lemah dukungannya dalam file
karena hasil rinci mengenai waktu perjalanan dan kemacetan tidak
disajikan.

\subsection{Contributions}\label{contributions}

\textbf{Kontribusi yang diklaim penulis.} Artikel menawarkan data
komuter sebagai sumber alternatif dan inovatif untuk mengenali
post-suburbanisasi. Kombinasi pola asal-tujuan dengan karakteristik
sosioekonomi dipakai untuk menunjukkan pelepasan sebagian fungsi
pinggiran dari inti, pembentukan pusat baru sesuai spesialisasi tiap
wilayah, serta variasi gentrifikasi dan segregasi dalam konteks negara
berkembang (Abstrak, hlm. 1; Bagian 5, hlm. 8).

Artikel juga membawa kontribusi kebijakan dengan menghubungkan diagnosis
mobilitas dan perkembangan spasial pada kebutuhan tata kelola
metropolitan multitingkat. Fokusnya adalah penguatan kapasitas
antarpemerintah untuk mengoordinasikan industri, perumahan, pekerjaan,
pusat belanja, pendidikan, dan infrastruktur di sepanjang koridor
antarkota (Bagian 5, hlm. 8).

\textbf{Inferensi pengolah.} Kontribusi metodologisnya terletak pada
kerangka tiga dimensi yang mengorganisasi indikator komuter, bukan pada
teknik statistik baru. Kontribusi empirisnya adalah menunjukkan secara
bersamaan arus pinggiran-ke-inti dan antarpinggiran sehingga perubahan
JMA tidak direduksi menjadi sepenuhnya monosentris atau sepenuhnya
mandiri.

\subsection{Research Gap}\label{research-gap}

\textbf{Celah yang dinyatakan penulis.} Penelitian post-suburbanisasi
yang menggunakan perjalanan komuter dan implikasi spasialnya masih
kurang. Secara khusus, penelitian tentang kegiatan sosioekonomi komuter
dan mobilitas spasial untuk memahami post-suburbanisasi dinyatakan
terbatas. Kajian JMA sebelumnya lebih sering mengandalkan data industri,
distribusi tenaga kerja, struktur spasial, dan pembangunan perumahan
skala besar daripada data pergerakan penduduk antara inti dan pinggiran
(Bagian 1, hlm. 1).

Penulis juga menempatkan JMA sebagai celah kontekstual. Banyak konsep
post-suburbanisasi dikembangkan dari kota Barat, sedangkan JMA memiliki
peran kuat pengembang swasta, dekonsentrasi industri, kota baru,
ketimpangan infrastruktur, serta perjalanan dari inti menuju kegiatan di
pinggiran yang membuat jalurnya tidak sepenuhnya sama (Bagian 3, hlm. 4;
Bagian 5, hlm. 8).

\textbf{Inferensi pengolah.} Artikel mengisi celah empiris penggunaan
data komuter, tetapi belum mengisi celah validasi pengukuran: file tidak
menguji apakah ketiga kelompok proksi membentuk konstruk
post-suburbanisasi yang konsisten, tidak membandingkannya dengan ukuran
langsung perubahan guna lahan atau lokasi pekerjaan, dan tidak
menetapkan aturan klasifikasi yang dapat direplikasi untuk menyatakan
suatu distrik telah mengalami post-suburbanisasi.

\subsection{Limitations}\label{limitations}

\textbf{Keterbatasan yang dinyatakan secara eksplisit oleh penulis.}
Tidak disebutkan dalam file.

\textbf{Keterbatasan yang diinferensikan pengolah dari desain dan
pelaporan.}

\begin{itemize}
\tightlist
\item
  Dua survei potong lintang membandingkan populasi pada 2014 dan 2019,
  bukan mengikuti individu yang sama. Perubahan komposisi tidak dapat
  langsung dibaca sebagai perubahan perilaku individu.
\item
  Analisis deskriptif dapat memperlihatkan asosiasi spasial, tetapi
  tidak membuktikan bahwa infrastruktur, sektor swasta, gentrifikasi,
  atau kebijakan tertentu menyebabkan pola perjalanan yang diamati.
\item
  Sampel hanya mencakup komuter berusia lebih dari 15 tahun. Perubahan
  wilayah disimpulkan dari perilaku komuter sehingga pengalaman
  nonkomuter dan penduduk di bawah batas usia tidak tercakup.
\item
  Agregasi pada tingkat kabupaten/kota dapat menyembunyikan variasi
  internal yang besar, khususnya karena artikel sendiri menggambarkan
  kompleksitas kota baru, kawasan industri, dan permukiman dalam distrik
  yang sama.
\item
  Umur perumahan, perubahan nilai lahan, perpindahan penduduk lama, atau
  proses penggantian kelas tidak diukur langsung. Oleh sebab itu,
  ``kematangan'' dan ``gentrifikasi'' ditopang oleh proksi usia,
  migrasi, dan kelas komuter.
\item
  Definisi kategori utama, termasuk komuter muda, kelas menengah-atas,
  hunian layak, dan pekerja formal, tidak tersedia. File juga tidak
  menerangkan ketidakpastian estimasi, prosedur pembobotan, data hilang,
  atau keterbandingan instrumen 2014 dan 2019.
\item
  Kerangka sampel disebut menggunakan Sensus Penduduk 2010, tetapi file
  tidak membahas konsekuensinya bagi survei 2019.
\item
  Sejumlah simpulan, terutama penurunan waktu perjalanan dan kemacetan,
  tidak disertai hasil numerik atau tabel yang dapat ditelusuri dalam
  teks yang tersedia (Bagian 5, hlm. 8).
\item
  Beberapa pernyataan naratif tidak selaras dengan angka tabel.
  Rinciannya dipertahankan dalam \texttt{Provenance\ Notes} dan tidak
  diselesaikan melalui sumber eksternal.
\end{itemize}

\subsection{Challenges}\label{challenges}

\textbf{Tantangan yang dilaporkan sumber.} Tata kelola JMA
terfragmentasi di banyak pemerintah dan tiga provinsi, sementara
kegiatan industri, perumahan, pekerjaan, pendidikan, dan transportasi
melintasi batas administratif. Penulis menilai koordinasi kelembagaan
dan komunikasi antara pemerintah, pengembang, serta pengelola kawasan
industri belum memadai (Bagian 3, hlm. 4-5; Bagian 4.3, hlm. 7; Bagian
5, hlm. 8).

Penyediaan infrastruktur juga tidak merata. Jalan raya dan jalan tol
menghubungkan inti dengan banyak wilayah pinggiran, tetapi jaringan
kereta terutama tersedia di arah timur dan selatan serta lebih terbatas
di wilayah barat. Ketimpangan pilihan moda, tingginya penggunaan
kendaraan pribadi, dan kedekatan spasial menghasilkan akses yang berbeda
menurut distrik dan pendapatan (Bagian 3, hlm. 2; Gambar 2; Bagian
4.3-5, hlm. 7-8).

Tantangan sosialnya ialah memastikan pembangunan inklusif ketika kota
baru, kawasan industri, real estat besar, dan komunitas berpagar yang
dipimpin sektor swasta dapat memperkuat segregasi. Tantangan
konseptualnya ialah mengenali bentuk JMA yang sekaligus menunjukkan
pusat baru, ketergantungan kuat pada Jakarta, dan pola yang tidak
sepenuhnya menyerupai pengalaman Barat (Bagian 3, hlm. 4-5; Bagian 5,
hlm. 8).

\textbf{Inferensi pengolah.} Bagi penerapan metode artikel, tantangan
utamanya adalah membedakan kemandirian fungsional dari sekadar kedekatan
tujuan, serta membedakan gentrifikasi dari komposisi komuter tanpa data
longitudinal, data penduduk nonkomuter, dan ukuran langsung perubahan
lingkungan binaan.

\subsection{Practical Implications}\label{practical-implications}

\textbf{Implikasi menurut penulis.} Pemerintah perlu memperkuat
kapasitas kelembagaan antarpemerintah untuk merumuskan,
mengoordinasikan, dan menerapkan regulasi industri dan perumahan di
sepanjang jalan raya antarkota. Tata kelola metropolitan perlu
mengintegrasikan lokasi industri, hunian, pekerjaan, pusat belanja,
kesempatan pendidikan, dan infrastruktur agar perkembangan polisentris
tidak memperdalam ketimpangan (Bagian 5, hlm. 8).

Pengelolaan gentrifikasi, revitalisasi, pembangunan pusat baru, dan
konektivitas perlu dilakukan lintas kabupaten/kota. Penulis menyebut
Badan Pengelola Transportasi Jabodetabek (BPTJ), layanan pengumpan
TransJakarta, dan JR Connexion sebagai bentuk kelembagaan atau layanan
yang dapat mendukung pengelolaan mobilitas metropolitan (Bagian 5, hlm.
8).

Pemerintah juga perlu meningkatkan transportasi umum, khususnya pada
wilayah yang kurang terlayani, dan mengimbangi dominasi pengembang
swasta melalui perencanaan spasial yang efektif. Tujuannya ialah
mencegah segmentasi kawasan pinggiran dan menjaga pembangunan JMA tetap
inklusif serta berkelanjutan (Bagian 5, hlm. 8).

\textbf{Inferensi pengolah.} Karena analisis tidak mengestimasi dampak
suatu intervensi, rekomendasi tersebut menunjukkan arah kebijakan, bukan
bukti kuantitatif bahwa satu layanan atau bentuk kelembagaan tertentu
akan menghasilkan besaran pengurangan waktu perjalanan, kemacetan, atau
segregasi tertentu.

\subsection{Applications}\label{applications}

\textbf{Laporan sumber.} Artikel menawarkan data komuter sebagai data
alternatif untuk mengidentifikasi post-suburbanisasi dan menghubungkan
hasilnya dengan perencanaan transportasi serta koordinasi tata kelola
lintas wilayah (Bagian 5, hlm. 8).

\textbf{Inferensi pengolah.} Kerangka artikel dapat diterapkan sebagai
diagnosis metropolitan dengan menggabungkan matriks asal-tujuan dan
profil sosioekonomi komuter. Penerapannya dapat membantu
mengidentifikasi wilayah yang masih bergantung pada inti, wilayah dengan
interaksi antarpinggiran yang kuat, perbedaan daya tarik kelompok
penduduk, dan kesenjangan akses moda. Artikel tidak menyajikan perangkat
lunak, protokol keputusan, peta intervensi, atau evaluasi implementasi
tertentu. Rincian aplikasi operasional semacam itu: Tidak disebutkan
dalam file.

\subsection{Future Research (Optional)}\label{future-research-optional}

Tidak disebutkan dalam file.

\subsection{Provenance Notes}\label{provenance-notes}

\begin{itemize}
\tightlist
\item
  Review ini hanya menggunakan file TXT
  \texttt{source/journal/J01-aritenang-2023-identifying-post-suburbanization-the-case-of-the-jakarta-metropolitan-area-jma-10.1016-j.habitatint.2023.102857.txt}.
  Tidak ada sumber eksternal yang digunakan untuk melengkapi,
  mengoreksi, atau memverifikasi klaim artikel.
\item
  Metadata ekstraksi menyatakan bahwa teks berasal dari PDF
  \texttt{journal/J01-aritenang-2023-identifying-post-suburbanization-the-case-of-the-jakarta-metropolitan-area-jma-10.1016-j.habitatint.2023.102857.pdf},
  diekstrak pada 31 Agustus 2026 dengan \texttt{pdftotext\ -layout}, dan
  terdiri atas 9 halaman. Review ini membaca seluruh 695 baris TXT,
  termasuk metadata, isi artikel, tabel yang terekstrak, pernyataan
  penulis, deklarasi konflik kepentingan, dan daftar pustaka. PDF asli
  tidak diperiksa secara visual.
\item
  Locator \texttt{hlm.} merujuk pada nomor halaman yang tampak pada
  hasil ekstraksi PDF, bukan rentang halaman bibliografis jurnal. Tabel
  dan gambar dirujuk sesuai nomor yang tercetak dalam artikel.
\item
  File menyebut sampel tertimbang 2019 mewakili 3,2 juta orang,
  sedangkan jumlah pada Tabel 2 adalah 3.341.535. Kedua angka dilaporkan
  apa adanya dan tidak direkonsiliasi (Bagian 4, hlm. 5; Tabel 2).
\item
  Narasi menyebut rerata proporsi seluruh migran dan migran
  menengah-atas sebesar 4,8\% dan 5,4\%, sedangkan baris total Tabel 4
  menampilkan 5,00\% dan 5,69\%. Review mengutamakan angka tabel ketika
  menyebut nilai tertentu, tetapi mencatat perbedaan ini tanpa koreksi
  eksternal (Bagian 4.2, hlm. 6-7; Tabel 4).
\item
  Narasi segregasi menyatakan bahwa pada 2014 wilayah berbatasan
  memiliki upah rata-rata kurang dari Rp5 juta dan wilayah tidak
  berbatasan lebih dari Rp7 juta. Angka Tabel 6 tidak memperlihatkan
  pola itu: seluruh wilayah tidak berbatasan yang dapat diidentifikasi
  berada sekitar Rp2,91-Rp5,01 juta, dan beberapa wilayah berbatasan
  berada di atas Rp5 juta. Narasi berikutnya menyebut Kota Bogor berupah
  tertinggi, sesuai Tabel 6 untuk 2019, tetapi menyebut Kota Bekasi
  memiliki proporsi hunian layak tertinggi meskipun Tabel 6 menampilkan
  Kabupaten Bekasi lebih tinggi pada 2019 (Bagian 4.3, hlm. 6-7; Tabel
  6).
\item
  Bagian diskusi menyebut Tangerang Selatan memiliki upah rata-rata dan
  proporsi hunian layak tertinggi. Tabel 6 justru menampilkan upah 2019
  tertinggi di Kota Bogor dan hunian layak tertinggi di Kabupaten
  Bekasi. Klaim diskusi dan angka tabel dipisahkan, bukan digabungkan
  (Bagian 5, hlm. 8; Tabel 6, hlm. 7).
\item
  Narasi menyatakan proporsi komuter muda di pinggiran lebih tinggi
  daripada inti, sedangkan Tabel 5 memperlihatkan proporsi muda wilayah
  inti yang lebih tinggi daripada semua wilayah pinggiran pada 2019.
  Ketidaksesuaian ini membatasi interpretasi kematangan berdasarkan usia
  (Bagian 4.2, hlm. 6-7; Tabel 5).
\item
  Judul Tabel 7 menyebut 2014 dan 2019, tetapi teks pembahasannya
  menyebut 2019 dan hasil ekstraksi hanya menampilkan satu set angka per
  moda. Karena itu, pernyataan penulis tentang kenaikan penggunaan
  kereta atau TransJakarta dan penurunan sepeda motor atau mobil tidak
  dapat diperiksa ulang dari angka tabel yang tersedia (Bagian 4.3, hlm.
  7-8; Tabel 7).
\item
  Artikel menyatakan adanya penurunan waktu perjalanan dan pengurangan
  kemacetan, tetapi hasil terperinci mengenai kedua ukuran tersebut
  tidak ditemukan dalam teks atau tabel yang tersedia (Bagian 5, hlm.
  8).
\item
  Pemenggalan kata, salah ketik, dan tata letak dua kolom pada TXT dapat
  menimbulkan ambiguitas, khususnya pada tabel. Jika narasi dan angka
  berbeda, review ini menyebut perbedaannya dan tidak mengarang nilai
  pengganti.
\end{itemize}

\section{Review J02}\label{review-j02}

\begin{itemize}
\tightlist
\item
  Kode sumber: J02
\item
  Judul: Pola Pergerakan dan Dekonsentrasi Pekerjaan di Kawasan
  Metropolitan: Studi Kasus Pekerja Industri Cikarang, Bekasi
\item
  Penulis: Putri Sugih Permatasari dan Delik Hudalah
\item
  Afiliasi: Sekolah Arsitektur, Perencanaan dan Pengembangan Kebijakan,
  Institut Teknologi Bandung
\item
  Terbitan: Jurnal Teknik Sipil, Vol. 20, No.~2, Agustus 2013, halaman
  tercetak 97-106
\item
  ISSN: 0853-2982
\item
  DOI: \texttt{10.5614/jts.2013.20.2.3} tercantum dalam nama file
  sumber, tetapi tidak tercetak pada isi artikel atau blok metadata PDF
  di dalam TXT
\item
  Fokus geografis dan sektoral: kawasan industri di Cikarang, Kabupaten
  Bekasi, dalam wilayah metropolitan Jabodetabek; pekerja sektor
  industri pengolahan
\item
  Dasar akses: seluruh teks hasil ekstraksi PDF yang tersedia dalam file
  TXT J02
\end{itemize}

\subsection{Summarized Abstract}\label{summarized-abstract-1}

Laporan sumber: Perkembangan kawasan industri di pinggiran Jabodetabek
meningkatkan jumlah pekerja dan pergerakan harian menuju tempat kerja.
Studi kasus ini memeriksa pola perjalanan pekerja menuju kawasan-kawasan
industri di Cikarang. Data primer berasal dari kuesioner terhadap 116
pekerja industri yang dipilih secara purposif. Hasilnya memuat lima pola
spasial, tetapi sampel sangat didominasi pergerakan lokal di kawasan
perkotaan Cikarang (68,97\%) dan pergerakan intradistrik dari wilayah
lain di Kabupaten Bekasi (21,55\%). Pergerakan dari Kota Bekasi sebesar
6,9\%, dari Kabupaten Bogor 1,72\%, dan dari Kabupaten Karawang 0,86\%.
Usia, status kepemilikan tempat tinggal, bentuk tempat tinggal, jabatan,
serta status kerja tercatat memiliki keterkaitan statistik dengan pola
pergerakan. (Abstrak, hlm. 97; Bagian 3, hlm. 100; Bagian 4 dan Gambar
2, hlm. 100-103; Tabel 2, hlm. 102)

Interpretasi penulis: Kehadiran pergerakan menuju Cikarang dari wilayah
di luar kawasan setempat dipandang sebagai indikasi bahwa orientasi
perjalanan bekerja mulai bergeser dari kota inti ke pusat pekerjaan baru
di pinggiran. Penulis menyebut proses tersebut sebagai dekonsentrasi
pekerjaan, khususnya pada industri pengolahan di Cikarang. (Abstrak,
hlm. 97; Bagian 5, hlm. 103-104)

Inferensi pengolahan: Artikel memberi bukti mengenai keragaman asal
tempat tinggal pekerja yang bekerja di Cikarang, bukan ukuran perubahan
distribusi pekerjaan dari waktu ke waktu. Karena data tidak memiliki
kerangka sampel yang diketahui, tidak mencakup seluruh Jabodetabek, dan
tidak menyajikan seri waktu jumlah pekerjaan, hasilnya paling aman
dibaca sebagai indikasi pola perjalanan menuju pusat pekerjaan pinggiran
pada sampel, sesuai kualifikasi penulis sendiri, bukan sebagai estimasi
tingkat dekonsentrasi metropolitan. (Bagian 3, hlm. 100; Bagian 5, hlm.
104)

\subsection{Problem Statement}\label{problem-statement-1}

Laporan sumber: Pola perjalanan bekerja di Jabodetabek selama ini
dipahami terkonsentrasi menuju kota inti Jakarta. Perkembangan dan
perluasan kegiatan, terutama industri, ke wilayah pinggiran menimbulkan
dugaan bahwa pekerjaan telah mengalami dekonsentrasi. Masalah empiris
artikel adalah apakah pola pergerakan harian pekerja yang tempat
kerjanya berada di kawasan industri Cikarang menunjukkan pergeseran
orientasi tersebut. Pergerakan bekerja dipilih karena berintensitas
tinggi, berlangsung rutin, dan dipakai sebagai indikator struktur
spasial wilayah metropolitan. (Bagian 1, hlm. 97-98)

Cikarang dipilih karena kegiatan industri Kabupaten Bekasi terpusat di
kawasan ini. Artikel melaporkan bahwa sektor industri pengolahan
menyumbang 78,21\% terhadap total PDRB Kabupaten Bekasi dengan laju
pertumbuhan 7,42\%, berdasarkan BPS Kabupaten Bekasi 2010. Pada saat
artikel ditulis, Cikarang disebut memiliki tujuh kawasan industri besar
yang beraglomerasi dan berkembang bersama fasilitas perkotaan,
permukiman, serta fasilitas sosial. (Bagian 1, hlm. 98)

Masalah kedua adalah keterkaitan karakteristik pekerja dengan pola
perjalanan mereka. Artikel menempatkan usia, jenis kelamin, status
perkawinan, asal daerah, pendidikan, pendapatan, kondisi tempat tinggal,
moda transportasi, jabatan, divisi, status kerja, dan lama kerja sebagai
karakteristik yang mungkin berkaitan dengan lokasi tempat tinggal dan
tempat kerja. (Bagian 1, hlm. 98-99)

Inferensi pengolahan: Rumusan masalah tidak disajikan sebagai pertanyaan
penelitian formal. Artikel juga tidak memisahkan secara tegas dua
persoalan yang berbeda, yaitu keberadaan arus perjalanan menuju
pinggiran dan perpindahan stok pekerjaan dari pusat ke pinggiran.
Keduanya dihubungkan melalui konsep Daily Urban System (DUS), tetapi
data yang dikumpulkan secara langsung hanya mengukur pola asal tempat
tinggal pekerja yang bekerja di Cikarang. (Bagian 1, hlm. 97-99; Bagian
3, hlm. 100)

\subsection{Objectives}\label{objectives-1}

Laporan sumber: Tujuan utama penelitian adalah mengidentifikasi dan
memaparkan pola pergerakan harian pekerja industri di wilayah pinggiran
Jabodetabek, dengan kawasan industri Cikarang sebagai studi kasus.
Tujuan tambahan adalah menjelaskan karakteristik demografis dan sosial
ekonomi yang berkaitan dengan pola pergerakan tersebut. (Abstrak, hlm.
97; Bagian 1, hlm. 97 dan 98-99)

Artikel juga menggunakan pola perjalanan untuk menilai ada atau tidaknya
indikasi dekonsentrasi pekerjaan ke Cikarang. Namun, pengukuran besaran,
kecepatan, atau perubahan temporal dekonsentrasi tidak dinyatakan
sebagai tujuan dan kemudian secara eksplisit diakui berada di luar
kemampuan studi. (Bagian 2, hlm. 99-100; Bagian 5, hlm. 104)

Pertanyaan penelitian dan hipotesis statistik formal: Tidak disebutkan
dalam file. Artikel hanya menyajikan lima pola spasial yang
``dihipotesakan'' akan muncul, lalu menguji keterkaitan sejumlah
karakteristik dengan pola pergerakan menggunakan uji Chi-Square. (Bagian
2, hlm. 99-100; keterangan Tabel 2, hlm. 102)

\subsection{Literature Survey}\label{literature-survey-1}

Laporan sumber: Artikel menggunakan konsep DUS dari Van Der Laan (1998)
untuk menggambarkan hubungan pergerakan antara kota inti dan wilayah
pinggiran. Empat tahapan yang dijelaskan ialah orientasi seluruh
aktivitas menuju inti; pergeseran aktivitas menuju pinggiran dan
munculnya perjalanan antarpinggiran; kemandirian relatif inti dan
wilayah fungsional pinggiran; serta pola berkebalikan atau reverse
ketika aktivitas dan perjalanan berlangsung dua arah antara inti dan
pinggiran. Konsep ini divisualkan pada Gambar 1. (Bagian 1 dan Gambar 1,
hlm. 97-98)

Aguilera dkk. (2009) dipakai untuk mendukung gagasan bahwa reverse
commuting dapat menandai tahap lanjut transformasi metropolitan. Ingram
(1998) dirujuk untuk hubungan perkembangan industri pinggiran dengan
indikasi dekonsentrasi pekerjaan. Firman dan Dharmapatni (1995) memberi
konteks perluasan kegiatan ke luar Jakarta, sedangkan Hudalah dan Firman
(2012) menjadi dasar pernyataan bahwa indikasi dekonsentrasi telah
muncul di Jabodetabek serta bahwa kawasan industri Cikarang berkembang
bersama fungsi perkotaan lain. Parnwell (1993) menjadi dasar dimensi
spasial perjalanan, dan Shon (2005) dirujuk untuk penggunaan pola
perjalanan bekerja sebagai indikator struktur spasial metropolitan.
(Bagian 1, hlm. 97-98)

Untuk karakteristik pekerja, Pas (1984) dan Tammaru (2005) menjadi
landasan hubungan faktor demografis, sosial, dan ekonomi dengan perilaku
perjalanan. Punpuing (1993) dilaporkan menemukan hubungan usia, jenis
kelamin, jabatan, waktu di tempat tinggal, status kepemilikan tempat
tinggal, jarak, dan waktu perjalanan dalam kasus Bangkok. Clark, Huang,
dan Withers (2003) dirujuk mengenai jarak perjalanan perempuan yang
lebih pendek serta toleransi waktu dan jarak dalam pilihan tempat
tinggal dan pekerjaan. Modarres (2011) menambahkan lama aktivitas di
tujuan. Ocakci (2000) digunakan untuk jenis industri dan tingkat
profesionalitas pekerja, sedangkan Lee (1966) menyediakan kerangka
faktor asal, tujuan, antara/perjalanan, dan personal. (Bagian 1, hlm.
98-99)

Klasifikasi operasional lima pola perjalanan mengadaptasi dimensi
spasial Parnwell (1993) dan pemetaan metropolitan Maryonoputri (2010).
Jangkauannya bergerak dari perjalanan lokal di Cikarang, perjalanan
dalam Kabupaten Bekasi, perjalanan antara Cikarang dan kota inti beserta
perluasannya, perjalanan antara Cikarang dan pinggiran Jabodetabek lain,
hingga perjalanan antara Cikarang dan wilayah di luar Jabodetabek.
(Bagian 2 dan Tabel 1, hlm. 99)

Inferensi pengolahan: Tinjauan literatur berfungsi sebagai kerangka
konseptual dan sumber daftar variabel, bukan tinjauan sistematis.
Strategi pencarian, basis data, kriteria pemilihan, jumlah publikasi
yang disaring, serta penilaian mutu literatur tidak disebutkan dalam
file. Klaim-klaim dari karya yang dirujuk di atas hanya dilaporkan
sebagaimana diringkas oleh artikel J02 dan tidak diverifikasi terhadap
publikasi aslinya.

\subsection{Method Used}\label{method-used-1}

Laporan sumber: Penelitian memakai studi kasus kawasan industri di
Cikarang dengan data primer dari kuesioner pekerja industri. Penulis
menetapkan lima kategori pola spasial berdasarkan lokasi tempat tinggal
relatif terhadap tempat kerja di Cikarang. Definisi tersebut menerapkan
Parnwell (1993) dan menyesuaikannya dengan pembagian spasial
metropolitan Maryonoputri (2010). (Abstrak, hlm. 97; Bagian 2, hlm. 99;
Bagian 3, hlm. 100)

Sebanyak 116 responden ditentukan secara purposif. Alasan yang diberikan
adalah kerangka sampel tidak tersedia, populasi pekerja berubah
mengikuti sistem kerja industri dan tidak tercatat secara lengkap, data
perusahaan tidak memadai, serta akses izin dan data perusahaan sulit.
Responden adalah pekerja yang menjadi anggota Federasi Serikat Pekerja
Metal Indonesia (FSPMI) Cabang Kabupaten Bekasi. Organisasi tersebut
dipakai untuk mendekatkan peneliti kepada pekerja dan memperoleh
informasi. Responden berasal dari perusahaan-perusahaan di MM2100, BIIE
(Hyundai), Delta Silicon 1 dan 2, Jababeka, dan EJIP. (Bagian 3, hlm.
100)

Analisis yang tampak dalam file terdiri atas statistik deskriptif untuk
proporsi lima pola perjalanan dan profil responden, pemetaan lokasi
tempat tinggal terhadap tempat kerja pada Gambar 2, serta uji Chi-Square
untuk menilai keterkaitan karakteristik pekerja dengan pola perjalanan.
Tabel 2 menandai variabel dengan nilai signifikansi kurang dari 0,05 dan
menyajikan nilai kontingensi Cramer sebagai ukuran keterkaitan. (Bagian
4 dan Gambar 2, hlm. 100-101; Tabel 2, hlm. 102; pembahasan, hlm. 103)

Inferensi pengolahan: Metode tersebut menguji asosiasi bivariat, bukan
pengaruh kausal atau model prediksi multivariat. File tidak menjelaskan
tanggal survei, cara kuesioner disebarkan, instrumen atau butir
pertanyaan, penanganan data hilang, tingkat respons, perangkat lunak
statistik, nilai statistik Chi-Square, derajat kebebasan, nilai p yang
tepat, maupun pemeriksaan asumsi uji. Semua unsur tersebut adalah
\texttt{Tidak\ disebutkan\ dalam\ file}.

\subsection{Dataset}\label{dataset-1}

Laporan sumber: Dataset adalah data primer kuesioner dengan unit terukur
satu jiwa pekerja industri dan 116 observasi responden. Cakupan tempat
kerja berasal dari perusahaan dalam enam nama kawasan atau kompleks yang
disebutkan: MM2100, BIIE (Hyundai), Delta Silicon 1, Delta Silicon 2,
Jababeka, dan EJIP. Cakupan lokasi tempat tinggal dibagi menurut kawasan
perkotaan Cikarang, bagian lain Kabupaten Bekasi, kota inti Jakarta dan
kota perluasannya, wilayah pinggiran Jabodetabek lain, serta wilayah di
luar Jabodetabek. (Bagian 2, Tabel 1, hlm. 99; Bagian 3, hlm. 100)

Bidang data yang dinyatakan dalam uraian meliputi lokasi tempat tinggal
dan tempat kerja; usia; jenis kelamin; status perkawinan; suku atau asal
daerah; tingkat pendidikan; pendapatan; status kepemilikan dan bentuk
tempat tinggal; moda transportasi; jabatan; divisi; status kerja; serta
lama kerja. Artikel juga menyebut jarak dan durasi perjalanan dalam
kerangka konseptual. Akan tetapi, Tabel 2 hanya menyajikan karakteristik
demografis, sosial ekonomi, dan pekerjaan. Hasil rinci tentang moda,
jarak, dan durasi tidak tersedia dalam file. (Bagian 1, hlm. 98-99;
Tabel 2, hlm. 102)

Periode pengumpulan data, kamus data, data mentah, format dataset,
sumber pendanaan khusus pengumpulan data, prosedur validasi jawaban, dan
akses publik dataset: Tidak disebutkan dalam file. Bagian ucapan terima
kasih hanya menyatakan bahwa artikel merupakan bagian penelitian ``Urban
Deconcentration sebagai Model Pengurangan Beban Kota Metropolitan
Jakarta'' yang disponsori Direktorat Pendidikan Tinggi, Kementerian
Pendidikan dan Kebudayaan. (Bagian 6, hlm. 104)

\subsection{Population / Sample}\label{population-sample-1}

Laporan sumber: Populasi sasaran substantif adalah pekerja industri di
kawasan-kawasan industri Cikarang. Besar dan daftar populasi tidak
diketahui; artikel menyatakan populasi tidak terdefinisi dan kerangka
sampel tidak tersedia. Sampel aktual terdiri atas 116 pekerja anggota
FSPMI Cabang Kabupaten Bekasi yang dipilih secara purposif dari
perusahaan-perusahaan di kawasan yang disebutkan. Karena itu, peluang
setiap pekerja untuk terpilih dan tingkat keterwakilan menurut
perusahaan atau kawasan tidak dapat dihitung. (Bagian 3, hlm. 100)

Profil demografis sampel dalam Tabel 2 adalah sebagai berikut. Usia
15-24 tahun berjumlah 19 orang (16,38\%), usia 25-34 tahun 85 orang
(73,28\%), usia 35-44 tahun 11 orang (9,48\%), dan usia di atas 44 tahun
satu orang (0,86\%). Responden laki-laki berjumlah 109 orang (93,97\%)
dan perempuan tujuh orang (6,03\%). Sebanyak 75 orang menikah (64,66\%)
dan 41 orang belum menikah (35,34\%). Asal yang paling banyak adalah
Jawa, 66 orang (56,90\%), dan Sunda, 36 orang (31,03\%). Tabel juga
mencatat Betawi delapan orang dengan persentase tertulis 6,09\%;
Bengkulu, Bugis, Lampung, dan Padang masing-masing satu orang (0,86\%);
serta Palembang dua orang (1,72\%). (Tabel 2, hlm. 102)

Profil sosial ekonomi menunjukkan satu responden berpendidikan SMP
(0,86\%), 106 SMA/SMK (91,38\%), tujuh D3 (6,03\%), dan dua S-1
(1,72\%). Pendapatan Rp1.000.000-Rp2.000.000 dilaporkan untuk 70 orang
(60,34\%), Rp2.000.000-Rp3.000.000 untuk 24 orang (20,69\%),
Rp3.000.000-Rp4.000.000 untuk 19 orang (16,38\%), dan lebih dari
Rp4.000.000 untuk tiga orang (2,59\%). Tempat tinggal milik sendiri
dimiliki 65 orang (56,03\%), kontrak atau sewa 42 orang (36,21\%), dan
menumpang pada saudara atau kerabat sembilan orang (7,76\%). Bentuk
tempat tinggal berupa rumah bagi 79 orang (68,10\%) dan kamar atau
indekos bagi 37 orang (31,90\%). (Tabel 2, hlm. 102)

Profil pekerjaan didominasi operator, yaitu 84 orang (72,41\%), diikuti
leader 19 orang (16,38\%), kategori jabatan lain enam orang (5,17\%),
foreman empat orang (3,45\%), dan supervisor tiga orang (2,59\%). Divisi
produksi mencakup 89 orang (76,72\%), teknisi 11 orang (9,48\%),
personalia tiga orang (2,59\%), keuangan satu orang (0,86\%), dan divisi
lain 12 orang (10,34\%). Pekerja tetap berjumlah 75 orang (64,66\%) dan
pekerja kontrak 41 orang (35,34\%). Lama kerja lebih dari dua tahun
mencakup 100 orang (86,21\%); kurang dari enam bulan tujuh orang
(6,03\%); enam bulan sampai satu tahun empat orang (3,45\%); dan satu
sampai dua tahun lima orang (4,31\%). (Tabel 2, hlm. 102)

Inferensi pengolahan: Sampel sangat terkonsentrasi pada laki-laki,
kelompok usia 25-34 tahun, pendidikan SMA/SMK, operator, divisi
produksi, dan anggota organisasi pekerja tertentu. Komposisi ini perlu
diperlakukan sebagai karakteristik sampel, bukan profil seluruh pekerja
industri Cikarang. Penulis sendiri menyatakan studi tidak dapat
merepresentasikan populasi pekerja industri pinggiran. (Tabel 2, hlm.
102; Bagian 5, hlm. 104)

\subsection{Dependent Variables}\label{dependent-variables-1}

Variabel dependen formal: Tidak disebutkan dalam file. Artikel tidak
menetapkan model regresi, arah sebab-akibat, atau terminologi variabel
dependen dan independen.

Inferensi pengolahan: Jika struktur analisis asosiasi pada Tabel 2
dinyatakan sebagai hubungan variabel, kategori pola pergerakan harian
dapat diperlakukan sebagai keluaran kategorikal utama. Lima kategorinya
adalah:

\begin{itemize}
\tightlist
\item
  Pola 1, lokal: perjalanan antara kawasan industri dan kawasan
  perkotaan Cikarang, yang mencakup Kecamatan Cikarang Pusat, Cikarang
  Timur, Cikarang Utara, dan Cikarang Barat.
\item
  Pola 2, intradistrik: perjalanan antara kawasan industri Cikarang dan
  bagian lain Kabupaten Bekasi.
\item
  Pola 3, antardistrik inti atau perluasan: perjalanan antara Cikarang
  dan DKI Jakarta beserta Kota Depok, Kota Bekasi, atau Kota Tangerang.
  Dalam hasil aktual, asal yang disebut hanya Kota Bekasi.
\item
  Pola 4, antardistrik antarpinggiran: perjalanan antara Cikarang dan
  Kabupaten Tangerang, Kabupaten Bogor, atau Kota Bogor. Dalam hasil
  aktual, asal yang disebut hanya Kabupaten Bogor.
\item
  Pola 5, antarregional: perjalanan antara Cikarang dan kabupaten atau
  kota di luar Jabodetabek. Dalam hasil aktual, asal yang disebut adalah
  Kabupaten Karawang.
\end{itemize}

Definisi pola berasal dari Tabel 1 dan uraian Bagian 2, sedangkan asal
aktual responden dijelaskan pada Bagian 4. Tabel 1 memakai tanda dua
arah antara Cikarang dan wilayah lain, tetapi sampel penelitian terdiri
atas orang yang bekerja di Cikarang. Dengan desain tersebut, data yang
dilaporkan secara langsung menunjukkan perjalanan dari tempat tinggal
menuju pekerjaan di Cikarang, bukan kedua arah arus pekerja. (Tabel 1,
hlm. 99; Bagian 3-4, hlm. 100-101)

\subsection{Results}\label{results-1}

Laporan sumber mengenai distribusi spasial: Kelima pola yang
dihipotesiskan ditemukan dalam sampel. Pola lokal merupakan 68,97\%
responden dan pola intradistrik 21,55\%. Responden pada pola
antardistrik inti atau perluasan berjumlah 6,9\% dan dalam hasil disebut
berasal dari Kota Bekasi. Pola antardistrik antarpinggiran sebesar
1,72\% dan berasal dari Kabupaten Bogor. Pola antarregional sebesar
0,86\% dan berasal dari Kabupaten Karawang. (Bagian 4, hlm. 100-101)

Perhitungan pengolahan: Penjumlahan dua proporsi pertama menunjukkan
bahwa 90,52\% sampel tinggal di kawasan perkotaan Cikarang atau bagian
lain Kabupaten Bekasi. Tiga proporsi dari luar Kabupaten Bekasi
berjumlah 9,48\%. Kedua angka agregat ini dihitung dari persentase yang
dilaporkan sumber, bukan angka agregat yang dicetak oleh penulis.
(Bagian 4, hlm. 100-101)

Untuk pola lokal dan intradistrik, lokasi tempat tinggal lain yang
disebut di Kabupaten Bekasi meliputi Kecamatan Setu, Cibitung, Serang
Baru, Karang Bahagia, Pebayuran, Kedungwaringin, dan Cibarusah. Sebagian
pekerja menyatakan memilih tinggal di sekitar kawasan perkotaan Cikarang
karena dekat dengan pekerjaan dan biaya transportasi terjangkau. Di
antara pekerja yang tinggal lebih jauh, alasan yang dilaporkan mencakup
kedekatan dengan keluarga dan fasilitas tempat tinggal yang dinilai
lebih lengkap daripada fasilitas di kawasan perkotaan Cikarang. Penulis
menduga ketersediaan infrastruktur jalan ikut memungkinkan perjalanan
dari kabupaten pinggiran atau luar Jabodetabek. (Bagian 4 dan Gambar 2,
hlm. 100-101)

Laporan sumber mengenai uji keterkaitan: Lima karakteristik bertanda
nilai signifikansi kurang dari 0,05 pada Tabel 2, yaitu usia, status
kepemilikan tempat tinggal, bentuk tempat tinggal, jabatan, dan status
kerja. Nilai kontingensi Cramer yang dilaporkan masing-masing adalah
0,260; 0,359; 0,441; 0,252; dan 0,311. Artikel tidak memberikan nilai p
yang tepat sehingga hasil hanya dapat dilaporkan menurut tanda
signifikansi pada tabel. Perbandingan langsung oleh pengolahan
menunjukkan bahwa bentuk tempat tinggal mempunyai nilai kontingensi
tertinggi di antara kelima variabel tersebut. (Tabel 2, hlm. 102;
pembahasan, hlm. 103)

Karakteristik yang tidak diberi tanda signifikansi ialah jenis kelamin
dengan nilai kontingensi 0,170; status perkawinan 0,246; suku atau asal
daerah 0,205; tingkat pendidikan 0,217; pendapatan 0,217; divisi kerja
0,141; dan lama kerja 0,115. Penulis mengaitkan ketiadaan pola menurut
jenis kelamin dengan dominasi responden laki-laki dan menyatakan asal
responden tersebar dari berbagai daerah. Untuk status perkawinan dan
pendapatan, penulis menawarkan penjelasan berupa kecenderungan keluarga
berpenghasilan tunggal serta adanya penghasilan sampingan. Dominasi
divisi produksi dan ketidaksesuaian langsung antara lama kerja dengan
status tetap diajukan sebagai penjelasan bagi dua hasil lainnya.
Penjelasan tersebut adalah interpretasi penulis, bukan hasil uji
mekanisme tersendiri. (Tabel 2, hlm. 102; pembahasan, hlm. 103)

Laporan silang karakteristik dan pola menunjukkan bahwa pola lokal,
intradistrik, serta pola dari inti atau perluasannya banyak dilakukan
responden berusia 15-44 tahun, dengan proporsi terbesar pada usia 25-34
tahun. Pola antarpinggiran disebut cenderung dilakukan pekerja berusia
25-44 tahun, sedangkan pola antarregional ditemukan pada pekerja berusia
35-44 tahun. Penulis menafsirkan pekerja muda cenderung mendekati tempat
kerja, sedangkan pekerja lebih tua dapat tinggal lebih jauh. (Bagian 4,
hlm. 101)

Pekerja dengan tempat tinggal sementara, baik menyewa maupun menumpang,
disebut cenderung mengikuti pola lokal dan intradistrik. Responden pada
pola antardistrik dan antarregional cenderung tinggal di rumah milik
sendiri. Jabatan operator mendominasi pola lokal dan intradistrik; pola
antardistrik juga mencakup foreman dan supervisor; satu-satunya kategori
antarregional yang digambarkan ialah operator. Pekerja kontrak cenderung
berada pada pola lokal dan intradistrik, sedangkan pekerja tetap muncul
pada semua pola. Semua pernyataan tersebut adalah deskripsi hubungan
dalam sampel, bukan estimasi populasi atau hubungan sebab-akibat.
(Bagian 4, hlm. 102-103)

\subsection{Findings}\label{findings-1}

Temuan yang dilaporkan: Cikarang menarik pekerja yang tinggal pada lima
lingkup geografis, mulai dari kawasan setempat hingga luar Jabodetabek.
Struktur sampel didominasi pola lokal sebesar 68,97\% dan intradistrik
sebesar 21,55\%. Perhitungan pengolahan atas angka sumber menghasilkan
gabungan 90,52\% untuk kedua pola dalam Kabupaten Bekasi dan 9,48\%
untuk tiga pola dari luar kabupaten tersebut. (Bagian 4, hlm. 100-101)

Interpretasi penulis: Kemunculan Pola 3 dan Pola 4 dipakai sebagai dasar
bahwa orientasi perjalanan bekerja mulai bergeser menuju pinggiran dan
bahwa dekonsentrasi pekerjaan telah terjadi di Cikarang. Cikarang
dinilai telah menjadi pusat pekerjaan baru yang menarik pekerja dari
kota inti atau perluasannya dan dari pinggiran lain. Penulis juga
menafsirkan pergerakan antarregional sebagai tanda bahwa Cikarang
menarik migran dari luar metropolitan. (Bagian 2, hlm. 100; Bagian 4,
hlm. 100-101; Bagian 5, hlm. 103-104)

Inferensi pengolahan: Bukti yang paling langsung adalah keberadaan
perjalanan menuju tempat kerja di Cikarang dari Kota Bekasi, Kabupaten
Bogor, dan Kabupaten Karawang. Uraian hasil tidak melaporkan responden
yang berasal dari DKI Jakarta; seluruh responden Pola 3 yang disebut
berasal dari Kota Bekasi, yang oleh klasifikasi artikel ditempatkan
sebagai kota perluasan. Karena itu, pernyataan pada abstrak dan
kesimpulan bahwa arus berasal dari ``kota inti Jakarta'' lebih luas
daripada rincian hasil yang tersedia dalam TXT. (Abstrak, hlm. 97; Tabel
1, hlm. 99; Bagian 4, hlm. 101; Bagian 5, hlm. 103)

Temuan statistik menunjukkan asosiasi pola perjalanan dengan usia,
kondisi tempat tinggal, jabatan, dan status kerja pada sampel.
Penggunaan kata ``penentu'' oleh penulis tidak boleh dibaca sebagai
kausalitas karena analisis yang tersedia hanya uji Chi-Square dan
kontingensi Cramer, tanpa urutan waktu atau pengendalian serentak
terhadap karakteristik lain. (Abstrak, hlm. 97; Tabel 2, hlm. 102;
Bagian 4-5, hlm. 103-104)

\subsection{Conclusions}\label{conclusions-1}

Kesimpulan penulis terdiri atas tiga pokok. Pertama, semua pola
perjalanan yang dihipotesiskan dinyatakan terjadi. Kehadiran pola
antardistrik dari arah inti atau perluasannya dan dari pinggiran lain
ditafsirkan sebagai indikasi dekonsentrasi. Cikarang dipandang menarik
penduduk setempat, pendatang luar metropolitan, serta pekerja dari
bagian lain wilayah metropolitan. Orientasi bekerja dinilai mulai
bergeser ke pusat pekerjaan industri di pinggiran. (Bagian 5, butir 1,
hlm. 103)

Kedua, penulis menyimpulkan bahwa perkembangan pekerjaan industri
Cikarang telah membantu mengurangi beban Jakarta melalui penciptaan
lapangan pekerjaan di pinggiran. Usia, status dan bentuk tempat tinggal,
jabatan, serta status kerja disebut berkaitan dengan pola perjalanan.
Pengetahuan tentang karakteristik ini dinilai penting bagi penyediaan
infrastruktur pekerja industri. (Bagian 5, butir 2, hlm. 104)

Ketiga, penulis membatasi kesimpulan tersebut sebagai bukti indikatif.
Studi tidak merepresentasikan populasi pekerja industri pinggiran, tidak
dapat menentukan sejauh mana dekonsentrasi telah terjadi, hanya
mengamati Cikarang dan bukan seluruh Jabodetabek, serta terbatas pada
pekerjaan industri. (Bagian 5, butir 3, hlm. 104)

Inferensi pengolahan: Kesimpulan ketiga adalah batas interpretasi yang
paling sesuai dengan desain. Klaim bahwa Cikarang ``mengurangi beban''
Jakarta tidak disertai indikator beban kota, kondisi pembanding, atau
pengukuran perubahan pekerjaan di Jakarta dalam file. Klaim tersebut
perlu diperlakukan sebagai interpretasi penulis, bukan hasil yang diukur
langsung. (Bagian 5, hlm. 104)

\subsection{Contributions}\label{contributions-1}

Kontribusi yang dapat ditelusuri dari isi artikel adalah penerapan
konsep DUS pada kasus pekerja industri Cikarang dan penerjemahannya
menjadi lima kelas perjalanan yang mencakup skala lokal, intradistrik,
antardistrik, dan antarregional. Klasifikasi ini menghubungkan lokasi
tempat tinggal pekerja dengan pusat pekerjaan industri di pinggiran
metropolitan. (Bagian 1-2, Gambar 1 dan Tabel 1, hlm. 97-100)

Artikel juga menyediakan gambaran empiris 116 responden tentang sebaran
asal pekerja dan hubungan pola perjalanan dengan karakteristik
demografis, sosial ekonomi, serta pekerjaan. Hasilnya memperlihatkan
bahwa arus berjarak jauh memang ada, tetapi pergerakan lokal dan dalam
Kabupaten Bekasi jauh lebih dominan. (Bagian 3-4, hlm. 100-103)

Inferensi pengolahan: Nilai utama studi terletak pada bukti kasus dan
operasionalisasi awal, bukan pada estimasi metropolitan. Artikel
menunjukkan cara pola perjalanan pekerja dapat dipakai untuk menyusun
dugaan tentang perubahan struktur spasial, sekaligus secara terbuka
mencatat bahwa desainnya belum dapat mengukur luas dekonsentrasi.
Artikel tidak memiliki bagian yang secara eksplisit berjudul kontribusi;
rumusan kontribusi di atas adalah sintesis pengolahan dari tujuan,
metode, hasil, dan keterbatasan yang dinyatakan penulis.

\subsection{Research Gap}\label{research-gap-1}

Kesenjangan yang dinyatakan sumber: Perkembangan pekerjaan di pinggiran
Jabodetabek diduga telah mengubah pola perjalanan yang sebelumnya
berorientasi ke Jakarta, tetapi pola harian pekerja industri di pusat
pinggiran Cikarang perlu diidentifikasi. Artikel juga berupaya mengisi
kebutuhan penjelasan tentang karakteristik pekerja yang berkaitan dengan
pola tersebut. (Bagian 1, hlm. 97-99)

Kesenjangan yang tetap terbuka menurut penulis meliputi pergerakan dan
mobilitas penduduk di pinggiran lain beserta faktor yang memengaruhinya;
perbandingan pekerja industri antarpinggiran, seperti Kabupaten
Tangerang dan kabupaten industri lain; mobilitas pekerjaan di sektor
perdagangan dan jasa; serta dampak pergerakan pekerja dan penduduk di
wilayah pinggiran metropolitan. (Bagian 5, hlm. 104)

Inferensi pengolahan: Beberapa kesenjangan metodologis belum dijawab
oleh studi ini. Tidak ada data perubahan jumlah atau pangsa pekerjaan
antara inti dan pinggiran dari waktu ke waktu; tidak ada sampel
pembanding pekerja yang bekerja di Jakarta atau pusat lain; tidak ada
estimasi arus dua arah; dan tidak ada kerangka sampel untuk menghasilkan
proporsi populasi. Rincian hasil juga belum memperlihatkan asal DKI
Jakarta walaupun kategori teoretis dan kesimpulan mencakup kota inti.
Data representatif berskala Jabodetabek dan berdimensi waktu diperlukan
jika pertanyaannya adalah besaran serta proses dekonsentrasi, bukan
sekadar keberadaan perjalanan menuju Cikarang. Inferensi ini berasal
dari perbandingan antara tujuan, desain, hasil, dan batasan internal
artikel, tanpa memakai sumber eksternal. (Bagian 2-5, hlm. 99-104)

\subsection{Limitations}\label{limitations-1}

Keterbatasan yang dinyatakan penulis adalah sebagai berikut:

\begin{itemize}
\tightlist
\item
  Kerangka sampel pekerja industri tidak dapat diperoleh sehingga studi
  tidak merepresentasikan populasi pekerja industri di wilayah
  pinggiran.
\item
  Studi hanya membuktikan indikasi dan tidak dapat menggambarkan sejauh
  mana dekonsentrasi telah terjadi.
\item
  Wilayah studi hanya Cikarang, bukan seluruh Jabodetabek.
\item
  Objek hanya pekerja industri sehingga hasil tidak menunjukkan
  dekonsentrasi pekerjaan secara keseluruhan atau pada sektor lain.
\end{itemize}

Semua keterbatasan tersebut dinyatakan pada Bagian 5, butir 3, hlm. 104.

Keterbatasan tambahan berdasarkan pemeriksaan internal TXT adalah
sebagai berikut:

\begin{itemize}
\tightlist
\item
  Pemilihan anggota FSPMI secara purposif dapat menghasilkan bias
  seleksi terhadap pekerja yang terhubung dengan organisasi tersebut.
  Komposisi sampel yang sangat berat pada laki-laki, operator, divisi
  produksi, dan pendidikan SMA/SMK mempersempit dasar generalisasi. Ini
  adalah inferensi pengolahan dari Bagian 3 dan Tabel 2, hlm. 100 dan
  102.
\item
  Kelompok perjalanan luar Kabupaten Bekasi sangat kecil, terutama Pola
  4 sebesar 1,72\% dan Pola 5 sebesar 0,86\%. Tabulasi silang dengan
  banyak kategori dapat bertumpu pada sangat sedikit observasi. File
  tidak melaporkan pemeriksaan frekuensi harapan uji Chi-Square,
  sehingga ketahanan hasil signifikansi tidak dapat dinilai. Ini adalah
  inferensi metodologis dari distribusi pada hlm. 101, ukuran sampel
  pada hlm. 100, dan keterangan Tabel 2 pada hlm. 102.
\item
  Nilai p, statistik uji, derajat kebebasan, interval ketidakpastian,
  dan tabulasi silang lengkap tidak disajikan. Pembaca tidak dapat
  mereplikasi uji atau memeriksa arah setiap asosiasi dari tabel yang
  tersedia. (Tabel 2, hlm. 102)
\item
  Data dijelaskan sebagai satu pengumpulan kuesioner tanpa periode atau
  dimensi waktu yang disebutkan. Secara inferensial, potret ini tidak
  dapat membuktikan proses ``pergeseran'' atau perubahan jumlah
  pekerjaan dari waktu ke waktu. (Bagian 3-4, hlm. 100-103)
\item
  Klasifikasi pada Tabel 1 menggunakan hubungan dua arah, sedangkan
  responden dipilih karena bekerja di Cikarang. Desain hanya mengamati
  asal tempat tinggal menuju tempat kerja Cikarang dan tidak mengukur
  pekerja yang tinggal di Cikarang lalu bekerja di Jakarta atau wilayah
  lain. (Tabel 1, hlm. 99; Bagian 3, hlm. 100)
\item
  Rincian hasil Pola 3 hanya menyebut Kota Bekasi, tetapi abstrak dan
  kesimpulan menyebut pekerja dari kota inti Jakarta. TXT tidak
  menyediakan rincian yang membuktikan keberadaan responden dari DKI
  Jakarta. (Abstrak, hlm. 97; Bagian 4, hlm. 101; Bagian 5, hlm. 103)
\item
  Tabel 2 mencatat delapan responden Betawi dengan persentase 6,09\%.
  Angka orang dan persentase tersebut tidak konsisten secara aritmetis
  untuk total 116 responden. Review ini mempertahankan angka sebagaimana
  tercetak dan tidak mengoreksinya tanpa data asli. (Tabel 2, hlm. 102)
\item
  Uraian pada hlm. 101 merujuk peta lokasi sebagai ``Gambar 1'',
  sedangkan caption pada halaman yang sama menyebut ``Gambar 2. Pola
  pergerakan harian pekerja industri di Cikarang''. Kemungkinan
  kesalahan penomoran dalam artikel atau artefak ekstraksi tidak dapat
  dipastikan hanya dari TXT.
\item
  Moda transportasi, jarak, dan durasi perjalanan dibahas sebagai
  variabel konseptual, tetapi hasil pengukurannya tidak disajikan.
  (Bagian 1, hlm. 98-99; Tabel 2 dan pembahasan, hlm. 102-103)
\end{itemize}

\subsection{Challenges}\label{challenges-1}

Laporan sumber: Tantangan utama penelitian adalah ketiadaan kerangka
sampel. Sistem kerja industri disebut sangat berubah sehingga jumlah
pekerja dapat berubah setiap saat dan tidak tercatat seluruhnya. Data
perusahaan tidak tersedia secara memadai, tidak memuat seluruh
perusahaan di wilayah studi, dan proses memperoleh izin serta data dari
perusahaan sulit. Peneliti mengakses responden melalui FSPMI Cabang
Kabupaten Bekasi sebagai jalan praktis untuk mendekati pekerja dan
memperoleh informasi. (Bagian 3, hlm. 100)

Tantangan analitis yang dapat diinferensikan adalah ketegangan antara
cakupan klaim metropolitan dan cakupan pengamatan yang terbatas pada
satu pusat pekerjaan serta satu saluran rekrutmen responden. Jumlah
responden pada pola jarak jauh juga terlalu kecil untuk mendukung uraian
rinci setiap kombinasi karakteristik secara kuat. Artikel tidak
menjelaskan bagaimana tantangan ini ditangani selain penggunaan sampel
purposif dan pembatasan klaim sebagai indikasi. (Bagian 3-5, hlm.
100-104)

\subsection{Practical Implications}\label{practical-implications-1}

Interpretasi dan rekomendasi penulis: Cikarang dipandang mampu menyerap
pekerjaan industri di pinggiran dan dengan demikian membantu mengurangi
beban Jakarta dalam penyediaan lapangan kerja. Pengetahuan mengenai
usia, kondisi tempat tinggal, jabatan, dan status kerja pekerja dinilai
penting dalam pertimbangan penyediaan infrastruktur bagi pekerja
industri di wilayah pinggiran. Infrastruktur jalan juga diduga
memungkinkan pekerja melakukan perjalanan dari wilayah yang lebih jauh.
(Bagian 4, hlm. 101; Bagian 5, hlm. 104)

Inferensi pengolahan: Temuan dominasi perjalanan lokal dan intradistrik
dapat memberi dasar awal untuk menempatkan kebutuhan transportasi
pekerja dan hunian dekat kawasan industri sebagai persoalan perencanaan
Cikarang dan Kabupaten Bekasi. Namun, artikel tidak menguji kebutuhan
layanan tertentu, tidak membandingkan moda, waktu, jarak, biaya, atau
kapasitas infrastruktur, dan tidak memberikan rancangan kebijakan
operasional. Jenis infrastruktur, lokasi prioritas, pihak pelaksana,
anggaran, serta ukuran keberhasilan: Tidak disebutkan dalam file.

\subsection{Applications}\label{applications-1}

Aplikasi yang sudah dilaksanakan berdasarkan hasil penelitian: Tidak
disebutkan dalam file.

Inferensi pengolahan: Dengan batas sampel yang dinyatakan, klasifikasi
lima pola dapat digunakan secara eksploratif untuk mengelompokkan
wilayah asal pekerja, membaca jangkauan pasar tenaga kerja kawasan
industri, dan merancang survei lanjutan tentang transportasi atau hunian
pekerja. Nilai kontingensi Cramer juga dapat dipakai untuk memilih
karakteristik yang perlu diperiksa lebih lanjut. Penggunaan untuk
proyeksi populasi, penentuan investasi, atau generalisasi seluruh
Jabodetabek belum didukung oleh data artikel. (Tabel 1, hlm. 99; Tabel
2, hlm. 102; Bagian 5, hlm. 104)

\subsection{Future Research
(Optional)}\label{future-research-optional-1}

Agenda yang disebut penulis mencakup:

\begin{itemize}
\tightlist
\item
  Studi pergerakan atau mobilitas penduduk dan aktivitasnya di wilayah
  pinggiran lain, termasuk faktor-faktor yang memengaruhi pergerakan.
\item
  Perbandingan pergerakan dan mobilitas pekerja industri di pinggiran
  lain, misalnya Kabupaten Tangerang atau kabupaten lain dengan
  pertumbuhan industri yang pesat.
\item
  Studi mobilitas pekerjaan di luar industri, terutama perdagangan dan
  jasa, di Jakarta dan wilayah pinggiran.
\item
  Studi dampak pergerakan pekerja dan penduduk di wilayah pinggiran
  metropolitan.
\end{itemize}

Seluruh agenda tersebut terdapat pada Bagian 5, hlm. 104. Variabel
hasil, desain, wilayah rinci selain contoh yang disebut, dan metode
untuk agenda lanjutan: Tidak disebutkan dalam file.

Inferensi pengolahan: Untuk menguji klaim dekonsentrasi dengan lebih
langsung, penelitian lanjutan perlu membedakan bukti keberadaan
perjalanan menuju pinggiran dari bukti perpindahan pekerjaan. Kebutuhan
itu dapat dijawab dengan pengamatan temporal, cakupan inti dan beberapa
pusat pinggiran, arus perjalanan dua arah, serta kerangka sampel yang
diketahui. Ini bukan rekomendasi eksplisit artikel, melainkan inferensi
yang diturunkan dari tujuan penelitian dan keterbatasan yang diakui
penulis. (Bagian 1, hlm. 97-98; Bagian 5, hlm. 104)

\subsection{Provenance Notes}\label{provenance-notes-1}

\begin{itemize}
\tightlist
\item
  Review ini hanya menggunakan seluruh isi file
  \texttt{/Users/mac/Documents/Mac/{[}2{]}\ Obsidian\ Vault/zettelkasten/source/journal/J02-permatasari-hudalah-2013-pola-pergerakan-dan-dekonsentrasi-pekerjaan-di-kawasan-metropolitan-studi-kasus-pekerja-industri-cikarang-bekasi-10.5614-jts.2013.20.2.3.txt},
  termasuk blok metadata dan teks ekstraksi pada baris 1-585. Tidak ada
  sumber eksternal yang digunakan.
\item
  Blok metadata menyatakan PDF sumber mempunyai 10 halaman dan diekstrak
  pada 2026-08-31 menggunakan \texttt{pdftotext\ -layout}. Kolom Title
  dan Author pada metadata PDF kosong. Identitas judul dan penulis
  diambil dari kepala artikel pada halaman tercetak 97, bukan dari kolom
  metadata PDF.
\item
  Locator mengacu pada nomor halaman tercetak artikel yang tertanam
  dalam TXT, yaitu 97-106, serta nama bagian, tabel, atau gambar yang
  tersedia. Halaman 106 dalam TXT hanya memuat kepala atau kaki halaman
  tanpa isi substantif tambahan.
\item
  TXT adalah ekstraksi tata letak dua kolom. Terdapat pemenggalan kata,
  baris kolom yang berdampingan, salah ketik atau ketidakrapian yang
  mungkin berasal dari artikel maupun ekstraksi, dan visual Gambar 1
  serta Gambar 2 tidak dapat diperiksa dari teks saja. Review tidak
  mengklaim verifikasi visual terhadap PDF.
\item
  Isi karya yang dikutip oleh Permatasari dan Hudalah tidak diperiksa
  secara terpisah. Pernyataan tentang literatur terdahulu adalah laporan
  atas penyajian penulis J02, bukan verifikasi terhadap sumber asli
  dalam Daftar Pustaka.
\item
  DOI diambil dari nama file sumber karena tidak tampak pada blok
  metadata atau teks artikel. Status telaah sejawat, tanggal pasti
  survei, data mentah, dan rincian metode yang tidak tersedia tidak
  diinferensikan sebagai fakta.
\item
  Istilah ``laporan sumber'' menandai data atau pernyataan yang
  disajikan artikel; ``interpretasi penulis'' menandai makna yang
  diberikan penulis terhadap hasil; dan ``inferensi pengolahan''
  menandai pembacaan analitis review ini yang hanya bertumpu pada
  informasi internal TXT.
\end{itemize}

\section{Review J03: Industrial Land Development and Manufacturing
Deconcentration in Greater
Jakarta}\label{review-j03-industrial-land-development-and-manufacturing-deconcentration-in-greater-jakarta}

\subsection{Identitas Sumber}\label{identitas-sumber}

\begin{itemize}
\tightlist
\item
  Kode sumber: J03.
\item
  Judul: \emph{Industrial Land Development and Manufacturing
  Deconcentration in Greater Jakarta}.
\item
  Penulis: Delik Hudalah, Dimitra Viantari, Tommy Firman, dan Johan
  Woltjer.
\item
  Publikasi: \emph{Urban Geography}, 2013, Volume 34, Nomor 7, halaman
  950-971.
\item
  DOI: \texttt{10.1080/02723638.2013.783281}.
\item
  Tanggal publikasi daring: 28 Juni 2013.
\item
  Afiliasi yang tercantum: School of Architecture, Planning and Policy
  Development, Institut Teknologi Bandung, Indonesia; serta Faculty of
  Spatial Sciences, University of Groningen, Belanda.
\item
  Status telaah sejawat: Tidak disebutkan dalam file.
\item
  Locator identitas: blok informasi publikasi sebelum artikel dan
  halaman 950.
\end{itemize}

\subsection{Summarized Abstract}\label{summarized-abstract-2}

\textbf{Laporan sumber.} Artikel menelaah apakah perkembangan kawasan
industri swasta di pinggiran Jakarta telah mengurangi dominasi Jakarta
dalam membentuk struktur metropolitan. Latar utamanya ialah dominasi
pengembang swasta dalam pembangunan lahan setelah kebijakan berorientasi
pasar pada akhir 1980-an serta pertumbuhan kawasan industri yang
mengiringi arus investasi asing langsung. Dengan analisis deskriptif dan
eksploratif atas data sekunder, artikel melacak distribusi pekerjaan
manufaktur formal pada 1995, 2001, dan 2010; membedakan konsentrasi
terencana, konsentrasi tidak terencana, dan wilayah nonkonsentrasi;
serta membandingkan spesialisasi sektoral pusat industri. Hasilnya
menunjukkan bahwa pusat-pusat industri pinggiran menangkap sebagian
besar pertumbuhan pekerjaan manufaktur yang keluar dari Jakarta, semakin
terspesialisasi sekaligus terdiversifikasi, dan mendorong pergeseran
dari pola tersebar menuju susunan yang relatif lebih polisentrik.
Meskipun demikian, sebagian besar pekerjaan manufaktur pada 2010 masih
berada di koridor tidak terencana dan kantong yang tersebar. (Locator:
Abstract, hlm. 950; Data and Methods, hlm. 953-954; Tabel 2-5, hlm.
963-966; Conclusion, hlm. 967.)

\textbf{Interpretasi penulis.} Para penulis menafsirkan pergeseran
tersebut sebagai fase lanjut dekonsentrasi industri dan sebagai tanda
bahwa struktur metropolitan polisentrik mungkin muncul apabila tren
berlanjut. Mereka juga menilai sektor swasta sebagai aktor penting,
sementara perencanaan spasial belum tampak memainkan peran krusial dalam
mengarahkan proses itu. (Locator: Scatteration or Polycentricity?, hlm.
963-964; Conclusion, hlm. 967.)

\textbf{Inferensi pengolahan.} Bukti artikel lebih kuat untuk
menunjukkan perubahan distribusi dan komposisi pekerjaan daripada untuk
membuktikan hubungan sebab akibat antara pembangunan kawasan industri
swasta dan dekonsentrasi. Istilah ``menuju polisentrik'' lebih tepat
daripada menyatakan bahwa Greater Jakarta sudah sepenuhnya polisentrik
karena penulis sendiri menyebut 60\% pekerjaan manufaktur masih berada
pada konsentrasi tidak terencana dan wilayah nonkonsentrasi pada 2010.
(Locator dasar inferensi: Data and Methods, hlm. 953-954; Tabel 3, hlm.
964.)

\subsection{Problem Statement}\label{problem-statement-2}

\textbf{Laporan sumber.} Dekonsentrasi penduduk Jakarta telah banyak
didokumentasikan, tetapi restrukturisasi pekerjaan sebagai pembentuk
suburbanisasi Indonesia memperoleh jauh lebih sedikit perhatian.
Literatur Amerika Utara dan Eropa cenderung memusatkan perhatian pada
dekonsentrasi jasa dalam kota pascaindustri, sedangkan perubahan
metropolitan di negara Asia Timur yang sedang berindustrialisasi lebih
banyak didorong manufaktur. Kajian terdahulu pada 1990-an menyimpulkan
bahwa Jakarta masih menjadi pusat manufaktur korporat utama, tetapi
pertumbuhan pesat kawasan industri terencana di pinggiran untuk
mengakomodasi investasi asing langsung menuntut pembaruan bukti tentang
hubungan pusat dan pinggiran. (Locator: Introduction, hlm. 951.)

Perumusan masalah spasialnya juga menyangkut batas wilayah kajian.
Pembangunan kawasan industri telah melampaui Jabodetabek menuju
Kabupaten Serang di barat dan Kabupaten Karawang di timur, sehingga
Jabodetabek saja tidak lagi memadai untuk memahami suburbanisasi
industri Jakarta. Pertanyaan inti artikel ialah sejauh mana pusat
industri swasta mengurangi dominasi pusat Jakarta dalam membentuk
struktur seluruh wilayah metropolitan. (Locator: Implications for
Metropolitan Spatial Structure and Pattern, hlm. 953; From Jabodetabek
to Greater Jakarta, hlm. 955-956; Gambar 1-2.)

\textbf{Inferensi pengolahan.} Masalah analitis artikel terdiri atas dua
lapis yang berkaitan tetapi tidak identik: mengukur dekonsentrasi
pekerjaan manufaktur dan menilai apakah dekonsentrasi itu menghasilkan
penyebaran acak atau reaglomerasi dalam pusat-pusat pinggiran. Artikel
menghubungkan keduanya secara deskriptif, bukan melalui pengujian
kausal. (Locator dasar inferensi: Data and Methods, hlm. 953-954;
Scatteration or Polycentricity?, hlm. 963-964.)

\subsection{Objectives}\label{objectives-2}

\textbf{Laporan sumber.} Tujuan yang dinyatakan ialah:

\begin{itemize}
\tightlist
\item
  Mendeskripsikan kecenderungan suburbanisasi ekonomi mutakhir di
  Greater Jakarta. (Locator: Introduction, hlm. 951.)
\item
  Menilai sejauh mana industri manufaktur telah terdekonsentrasi di
  kawasan metropolitan terbesar Indonesia. (Locator: Introduction, hlm.
  951.)
\item
  Mengeksplorasi pengaruh dekonsentrasi industri terhadap struktur
  spasial metropolitan. (Locator: Introduction, hlm. 951.)
\item
  Mendokumentasikan suburbanisasi industri selama sekitar dua dekade
  serta meninjau faktor sosial-politik dan ekonomi yang berkontribusi
  terhadapnya. (Locator: Introduction, hlm. 951.)
\item
  Menyelidiki sejauh mana pusat industri swasta telah mengurangi
  dominasi Jakarta dalam struktur metropolitan. (Locator: Implications
  for Metropolitan Spatial Structure and Pattern, hlm. 953.)
\item
  Mendeskripsikan bentuk spasial yang dihasilkan dekonsentrasi dan
  membandingkan struktur serta spesialisasi sektor pada pusat industri.
  (Locator: Data and Methods, hlm. 954; Specialization of Industrial
  Centers in Greater Jakarta, hlm. 964-966.)
\end{itemize}

Hipotesis formal: Tidak disebutkan dalam file.

\subsection{Literature Survey}\label{literature-survey-2}

\textbf{Laporan sumber.} Artikel menyusun landasan pustaka dalam empat
gugus berikut.

\begin{itemize}
\tightlist
\item
  \textbf{Definisi dekonsentrasi.} Dekonsentrasi dipahami sebagai
  perubahan kota menjadi wilayah terurbanisasi melalui penyebaran besar
  penduduk dan pekerjaan ke pinggiran. Indikator umumnya ialah turunnya
  bagian penduduk atau pekerjaan di inti metropolitan dan naiknya bagian
  pinggiran. (Locator: Deconcentration, hlm. 951.)
\item
  \textbf{Objek dekonsentrasi pekerjaan.} Literatur dibedakan antara
  dekonsentrasi penduduk dan dekonsentrasi pekerjaan. Kajian masyarakat
  pascaindustri terutama menyoroti jasa bisnis dan suburbanisasi
  perkantoran, walaupun manufaktur juga pernah dan masih berperan dalam
  geografi pekerjaan kota-kota Barat. Di kota global Asia Timur,
  terutama di negara yang sedang berindustrialisasi, manufaktur
  diposisikan sebagai pendorong penting restrukturisasi metropolitan.
  (Locator: Deconcentration, hlm. 952.)
\item
  \textbf{Perbedaan mekanisme Barat dan Asia Timur.} Literatur Amerika
  Utara menekankan keputusan lokasi privat, lahan pinggiran yang lebih
  murah, dan kekuatan pasar lokal yang endogen. Literatur Asia Timur
  menekankan faktor eksternal dan global, khususnya relokasi investasi
  asing langsung, serta peran kebijakan negara. Kasus Tiongkok
  digambarkan lebih terintegrasi dengan perencanaan nasional dan zona
  ekonomi khusus, sedangkan artikel ini memosisikan Indonesia sebagai
  kasus ketika negara memberi panduan umum tetapi pengembang lahan
  swasta menggerakkan bentuk dan geografi suburbanisasi. (Locator:
  Deconcentration, hlm. 952-953.)
\item
  \textbf{Konsekuensi bagi struktur metropolitan.} Perspektif sentris
  berpendapat bahwa ekonomi aglomerasi membuat pekerjaan yang
  terdesentralisasi kembali mengelompok pada sedikit pusat pinggiran
  sehingga membentuk polisentrisitas. Perspektif desentris menekankan
  mobilitas, telekomunikasi, dan globalisasi yang memungkinkan pekerjaan
  menyebar dalam banyak kantong kecil tanpa batas jelas. Artikel
  kemudian menambahkan konteks Asia berupa \emph{desakota}, kota hunian
  atau industri yang direncanakan dan dikelola swasta, pergeseran dari
  ruang yang dikendalikan publik ke swasta, serta zona ekonomi khusus
  yang dipimpin negara di Tiongkok. (Locator: Implications for
  Metropolitan Spatial Structure and Pattern, hlm. 952-953.)
\end{itemize}

\textbf{Interpretasi penulis.} Posisi teoritis artikel lebih dekat pada
argumen reaglomerasi: dekonsentrasi manufaktur dapat tetap menghasilkan
pusat-pusat baru, bukan semata-mata penyebaran acak. Namun, para penulis
mempertahankan kemungkinan bentuk campuran karena pekerjaan di wilayah
tidak terencana dan nonkonsentrasi masih dominan. (Locator:
Introduction, hlm. 951; Scatteration or Polycentricity?, hlm. 963-964.)

\textbf{Inferensi pengolahan.} Nilai pembanding utama survei pustaka ini
terletak pada tipologi mekanisme: pasar lokal Barat, negara nasional
Tiongkok, dan fasilitasi negara yang diikuti dominasi pengembang swasta
di Indonesia. Artikel tidak melakukan tinjauan sistematis, tidak
menjelaskan strategi pencarian pustaka, dan tidak menetapkan kriteria
inklusi atau korpus. (Locator dasar inferensi: bagian Deconcentration
dan Implications for Metropolitan Spatial Structure and Pattern, hlm.
951-953.)

\subsection{Method Used}\label{method-used-2}

\textbf{Laporan sumber.} Penelitian bersifat deskriptif dan eksploratif.
Penulis menghimpun beragam bukti untuk memberi konteks pada evolusi
ekonomi spasial Jakarta selama 15 tahun. Penelitian membahas proses
kausal, khususnya peran pusat industri swasta, tetapi secara eksplisit
tidak mengklaim telah memverifikasi proses tersebut secara empiris;
tujuannya ialah mendeskripsikan dan, sampai taraf tertentu, menguatkan
proses itu. (Locator: Data and Methods, hlm. 953-954.)

Langkah analisis yang dilaporkan meliputi:

\begin{itemize}
\tightlist
\item
  Analisis isi dokumen kebijakan dan penelitian terdahulu untuk
  menjelaskan faktor dan aktor pemicu dekonsentrasi. (Locator: Data and
  Methods, hlm. 954.)
\item
  Perhitungan rasio jumlah pekerja terhadap populasi
  (\emph{employment-population ratio}, E/P) untuk setiap unit spasial
  guna menggambarkan distribusi pekerjaan manufaktur. Indikator relatif
  ini dipilih untuk mengurangi bias akibat distribusi lahan kosong yang
  tidak merata dan penggunaan unit spasial administratif yang arbitrer.
  (Locator: Data and Methods, hlm. 954.)
\item
  Pemetaan E/P pada 1995, 2001, dan 2010 untuk melacak pergeseran pusat
  pekerjaan. Zona dengan E/P lebih dari 0,1 dikategorikan sebagai zona
  konsentrasi manufaktur; zona dengan E/P paling banyak 0,1
  dikategorikan sebagai zona nonkonsentrasi. (Locator: Data and Methods,
  hlm. 954; Scatteration or Polycentricity?, hlm. 963.)
\item
  Penggabungan kompilasi data kawasan industri dengan E/P untuk
  membedakan konsentrasi terencana yang didominasi kawasan industri,
  konsentrasi tidak terencana yang didominasi zona industri, dan wilayah
  nonkonsentrasi. (Locator: Scatteration or Polycentricity?, hlm. 963.)
\item
  Perhitungan \emph{Location Quotient} (LQ), yaitu bagian suatu sektor
  pada satu pusat industri dibagi bagian sektor tersebut pada seluruh
  pusat industri yang dianalisis. LQ yang lebih tinggi ditafsirkan
  sebagai kemungkinan spesialisasi yang lebih tinggi. (Locator: Data and
  Methods, hlm. 954.)
\item
  Penyelarasan klasifikasi sektor 1995 dan 2010 dengan \emph{Table of
  Business Sector Matching 1990-2000}. Data 2010 memuat 16 sektor,
  sedangkan data 1995 disederhanakan menjadi sembilan kelompok karena
  keterbatasan data. (Locator: Specialization of Industrial Centers in
  Greater Jakarta, hlm. 964.)
\end{itemize}

\textbf{Inferensi pengolahan.} Tidak ada model regresi, pengujian
hipotesis, desain kontrafaktual, atau strategi identifikasi kausal yang
dilaporkan. Karena itu, metode mendukung pembacaan pola spasial dan
temporal, tetapi tidak dapat mengisolasi dampak pembangunan kawasan
industri swasta dari kebijakan, jalan tol, harga lahan, investasi asing
langsung, atau faktor lain yang bergerak bersamaan.

\subsection{Dataset}\label{dataset-2}

\textbf{Laporan sumber.} Bahan penelitian seluruhnya berupa data
sekunder, yang mencakup:

\begin{itemize}
\tightlist
\item
  Sensus penduduk dan sensus industri Badan Pusat Statistik untuk
  rentang 1995-2010. Data pekerjaan yang dipakai dalam analisis utama
  berasal dari sensus perusahaan industri formal tahun 1995, 2001, dan
  2010. (Locator: Data and Methods, hlm. 954.)
\item
  Survei pasar tanah dan properti oleh konsultan properti global dari
  1990-an hingga 2010-an. Gambar 4 dan Tabel 1 secara khusus menyebut
  analisis atau perhitungan dari laporan Colliers International tahun
  1997, 2002, 2007, dan 2010. (Locator: Data and Methods, hlm. 954;
  Gambar 4, hlm. 959; Tabel 1, hlm. 959.)
\item
  Dokumen kebijakan dan penelitian terdahulu. Narasi kebijakan secara
  khusus membahas Keputusan Presiden Nomor 53 Tahun 1989 dan Keputusan
  Presiden Nomor 41 Tahun 1996 mengenai kawasan industri. (Locator:
  Major Factors and Actors Triggering Industrial Deconcentration, hlm.
  957-958.)
\item
  Bukti kontekstual lain yang disebut dalam narasi, termasuk data izin
  investasi BKPM tahun 2011 dan laporan tahunan PT Jababeka tahun 2010.
  (Locator: Major Factors and Actors Triggering Industrial
  Deconcentration, hlm. 958.)
\end{itemize}

Cakupan geografis utama ialah Greater Jakarta, yaitu Jabodetabek
ditambah pinggiran luar tambahan ke arah Serang dan Karawang. Cakupan
temporal analisis pekerjaan ialah 1995-2010, sedangkan bukti pasar lahan
dan kebijakan mencakup periode yang lebih luas sejak akhir 1960-an
sampai 2011. (Locator: Gambar 1, hlm. 955; From Jabodetabek to Greater
Jakarta, hlm. 955-956; Data and Methods, hlm. 954.)

Format mikrodata, prosedur pembersihan, penanganan nilai hilang, dan
akses terhadap data mentah: Tidak disebutkan dalam file.

\subsection{Population / Sample}\label{population-sample-2}

\textbf{Laporan sumber.} Populasi observasi sensus industri dibatasi
pada perusahaan industri formal skala menengah dan besar yang
mempekerjakan lebih dari 20 pekerja. Dengan demikian, angka pekerjaan
merujuk pada pekerja di perusahaan yang memenuhi batas tersebut, bukan
seluruh pekerjaan manufaktur. (Locator: Data and Methods, hlm. 954.)

Unit analisis spasial terdiri atas kecamatan di pinggiran dan kotamadya
di Provinsi Daerah Khusus Ibu Kota Jakarta sebagai inti metropolitan.
Untuk menangani perubahan batas kecamatan antara 1995 dan 2010, Greater
Jakarta dipecah menjadi 145 zona, yaitu 140 kecamatan pinggiran dan lima
kotamadya di inti metropolitan. (Locator: Data and Methods, hlm. 954.)

Untuk analisis LQ, pusat yang disertakan harus memiliki E/P lebih dari
0,1. Pada 1995, teks menyebut lima pusat: Jakarta Utara sebagai bagian
dari CBD serta Cikupa-Balaraja, Cikande, Cilegon, dan Cikarang di
pinggiran. Telukjambe di Kabupaten Karawang ditambahkan untuk analisis
2010. Pada klasifikasi bentuk spasial, enam dari tujuh pusat industri
memenuhi kriteria konsentrasi terencana, sedangkan Cikampek dikeluarkan
karena E/P kurang dari 0,1. (Locator: Scatteration or Polycentricity?,
hlm. 963; Specialization of Industrial Centers in Greater Jakarta, hlm.
965.)

Teknik pengambilan sampel probabilitas, tingkat respons sensus, dan
jumlah perusahaan: Tidak disebutkan dalam file. Penelitian menyebut
sensus perusahaan, tetapi tidak memberikan rincian cakupan aktual atau
kehilangan observasi.

\subsection{Dependent Variables}\label{dependent-variables-2}

\textbf{Laporan sumber.} Variabel dependen dalam pengertian model
statistik atau kausal: Tidak disebutkan dalam file. Artikel tidak
menetapkan model dengan variabel terikat dan variabel bebas.

Ukuran hasil yang dianalisis secara deskriptif ialah:

\begin{itemize}
\tightlist
\item
  Jumlah, bagian, dan pertumbuhan pekerjaan manufaktur formal menurut
  zona metropolitan dan bentuk spasial. (Locator: Tabel 2, hlm. 963;
  Tabel 3, hlm. 964.)
\item
  Rasio E/P sebagai indikator relatif konsentrasi pekerjaan manufaktur
  pada unit spasial. (Locator: Data and Methods, hlm. 954; Gambar 5-7,
  hlm. 960-962.)
\item
  LQ sektoral sebagai indikator spesialisasi relatif pusat industri.
  (Locator: Data and Methods, hlm. 954; Tabel 4-5, hlm. 965-966.)
\item
  Pasokan lahan industri dan median harga lahan sebagai bukti
  kontekstual perkembangan pusat industri, bukan sebagai variabel
  dependen dalam suatu model. (Locator: Gambar 4 dan Tabel 1, hlm.
  959-960.)
\end{itemize}

\textbf{Inferensi pengolahan.} Penyebutan ukuran-ukuran tersebut sebagai
``hasil'' tidak berarti artikel mengonstruksi variabel dependen formal.
E/P dan LQ adalah indikator deskriptif yang dipakai untuk klasifikasi
serta perbandingan spasial.

\subsection{Results}\label{results-2}

\textbf{Laporan sumber.} Hasil deskriptif dapat diringkas dalam tujuh
kelompok.

\begin{itemize}
\tightlist
\item
  \textbf{Perluasan wilayah industri metropolitan.} Jabodetabek seluas
  5.898 km2 tidak lagi dianggap cukup untuk menangkap suburbanisasi
  industri karena kawasan industri berkembang menuju Kabupaten Serang
  dan Kabupaten Karawang. Greater Jakarta yang diperluas mencakup
  9.016,43 km2. Artikel mencatat lebih dari 35 kawasan industri dengan
  luas total lebih dari 18.000 hektare; ukuran setiap kawasan berkisar
  50-1.800 hektare dengan rata-rata sekitar 500 hektare. Cikarang
  memiliki hampir 6.000 hektare lahan industri. (Locator: From
  Jabodetabek to Greater Jakarta, hlm. 955-957; Gambar 1-2.)
\item
  \textbf{Keterkaitan koridor dan pusat industri.} Kawasan industri
  mengikuti, atau setidaknya berkorespondensi dengan, pembangunan jalan
  antarkota. Jagorawi yang beroperasi sejak 1978 diikuti zona industri
  tidak terencana di koridor selatan, sedangkan sebagian besar kawasan
  industri terencana sejak 1989 berlokasi di sepanjang Jalan Tol
  Jakarta-Cikampek dan Jakarta-Merak. Enam pusat industri pinggiran
  dengan lahan industri di atas rata-rata 500 hektare berada di
  Cikupa-Balaraja, Cikarang, Cilegon/Ciwandan, Cikande, Telukjambe, dan
  Cikampek. (Locator: From Jabodetabek to Greater Jakarta, hlm. 955-957;
  Gambar 1-2.)
\item
  \textbf{Peralihan pembangunan dari negara dan kota ke swasta dan
  pinggiran.} Kawasan industri awal pada 1970-an dan 1980-an terutama
  dibangun perusahaan negara di pesisir utara Jakarta. Keputusan
  Presiden Nomor 53 Tahun 1989 memfasilitasi partisipasi swasta,
  sedangkan Keputusan Presiden Nomor 41 Tahun 1996 menetapkan pedoman
  kawasan industri beserta infrastruktur dan fasilitas pendukungnya.
  Sejak akhir 1980-an tidak ada kawasan industri baru di Jakarta; sejak
  pertengahan 1990-an pinggiran mengambil alih pasokan lahan industri
  formal. Pada 2011, BKPM memberikan 317 izin tetap investasi asing baru
  terkait manufaktur di Jabodebek, dibandingkan dengan 75 investasi
  domestik. (Locator: Major Factors and Actors Triggering Industrial
  Deconcentration, hlm. 957-959; Gambar 3-4.)
\item
  \textbf{Perbedaan harga dan pasokan lahan.} Sampai puncak ledakan
  properti 1997, median harga lahan industri Jakarta sebesar 228,5 dolar
  AS/m2, sedangkan median pinggiran 74,5 dolar AS/m2, atau sekitar
  sepertiganya. Sejak pertengahan 2000-an Jakarta tidak lagi memasok
  lahan baru untuk industri skala besar. Tabel 1 mencatat median
  pinggiran sebesar 42,7 dolar AS/m2 pada 2002, 61,0 pada 2007, dan
  116,7 pada 2011; nilai untuk Jakarta pada 2007 dan 2011 tidak tersedia
  karena tidak ada lahan skala besar baru. (Locator: Tabel 1, hlm. 959;
  pembahasan hlm. 959-960.)
\item
  \textbf{Pergeseran lokasi 1995-2010.} Pada 1995, zona pekerjaan dengan
  E/P lebih dari 0,1 sudah tersebar, tetapi zona di Jakarta hanya
  terlihat di bagian utara dan total area zona pekerjaan pinggiran lebih
  luas. Pada 2001, pekerjaan manufaktur menurun di Jakarta dan
  konsentrasi bergeser dari zona tidak terencana di selatan serta dari
  kota-kota pinggiran dalam menuju pusat terencana di timur, terutama
  Cikarang dan Telukjambe. Pada 2010, pusat utama di Jakarta relatif
  tidak berubah, yang ditafsirkan penulis sebagai kejenuhan, sementara
  pusat pinggiran lebih dinamis. Periode 1990-an hingga 2000-an ditandai
  perpindahan dari zona industri ke kawasan industri, sedangkan 2000-an
  hingga 2010 ditandai intensifikasi kawasan industri pinggiran yang
  sudah ada. (Locator: The Dynamics of Manufacturing Employment in
  Greater Jakarta, hlm. 960-962; Gambar 5-7.)
\item
  \textbf{Pertumbuhan pekerjaan menurut zona metropolitan.} Total
  pekerjaan manufaktur formal meningkat dari 1.056.399 pada 1995 menjadi
  1.227.709 pada 2010, yaitu bertambah 171.310 atau 16\%. Jakarta turun
  dari 376.386 menjadi 321.270 pekerja, atau 15\%; pinggiran dalam turun
  dari 278.354 menjadi 256.451, atau 8\%; pinggiran luar naik dari
  323.072 menjadi 446.199, atau 38\%; dan pinggiran luar tambahan naik
  dari 78.587 menjadi 203.789, atau 159\%. Jakarta menampung 36\%
  pekerjaan metropolitan pada 1995. Tabel 2 memberi label persentase
  minus 15 sebagai pertumbuhan jumlah pekerjaan Jakarta, walaupun narasi
  menyebutnya sebagai penurunan ``employment share''. (Locator: Tabel 2
  dan pembahasannya, hlm. 962-963.)
\item
  \textbf{Perubahan bentuk spasial.} Pekerjaan pada konsentrasi
  terencana meningkat dari 298.749 menjadi 496.672, atau 66\%, dan
  bagiannya naik dari sekitar 28\% menjadi 40\%. Konsentrasi tidak
  terencana turun dari 375.982 menjadi 351.633, atau 6\%, sedangkan
  wilayah nonkonsentrasi turun dari 381.668 menjadi 379.404, atau 1\%.
  Gabungan konsentrasi tidak terencana dan wilayah nonkonsentrasi turun
  dari 72\% menjadi 60\% pekerjaan metropolitan. (Locator: Tabel 3 dan
  pembahasannya, hlm. 963-964.)
\end{itemize}

Hasil LQ menunjukkan diferensiasi sektoral. Pada 1995, CBD masih
menampung sekitar 170.000 pekerjaan dan menjadi lokasi utama berbagai
sektor, tetapi pada 2010 jumlahnya turun menjadi sekitar 160.000 dan
narasi menyebut keunggulannya terutama bertahan pada barang konsumsi
seperti makanan, pakaian, dan penerbitan. Industri yang memerlukan lahan
luas atau bergantung pada sumber daya, termasuk tekstil, kulit, mesin,
elektronik, dan otomotif, bergeser ke pinggiran. Pada 2010, Karawang
mencatat LQ 3,49 untuk elektronik dan 2,12 untuk mesin, sedangkan
Cikarang mencatat 2,01 untuk elektronik dan 1,70 untuk mesin. Cikande
dan Cikupa-Balaraja menonjol pada kulit serta tekstil, sedangkan Cilegon
menonjol pada energi, kimia, logam, dan kayu. Hasil Cilegon harus dibaca
hati-hati karena jumlah pekerjaannya jauh di bawah rata-rata pusat
industri. (Locator: Tabel 4-5 dan pembahasan, hlm. 965-966.)

\subsection{Findings}\label{findings-2}

\textbf{Temuan yang dilaporkan sumber.}

\begin{itemize}
\tightlist
\item
  Pertumbuhan bersih pekerjaan manufaktur formal selama 1995-2010
  seluruhnya terkonsentrasi di pinggiran luar dan pinggiran luar
  tambahan, sementara Jakarta dan pinggiran dalam kehilangan pekerjaan.
  (Locator: Tabel 2, hlm. 963.)
\item
  Pusat konsentrasi terencana memperoleh tambahan 197.923 pekerjaan,
  lebih besar daripada pertumbuhan bersih Greater Jakarta sebesar
  171.310 pekerjaan, sedangkan dua bentuk spasial lainnya mengalami
  penurunan. Secara aritmetis, kenaikan pusat terencana mencakup seluruh
  pertumbuhan bersih wilayah dan mengimbangi penurunan pada dua bentuk
  lain; data agregat tidak dapat membedakan relokasi perusahaan dari
  penciptaan atau pengurangan pekerjaan. (Locator: Tabel 2-3, hlm.
  963-964. Kalimat kedua merupakan inferensi aritmetis dari angka yang
  dilaporkan.)
\item
  Walaupun pola tersebar masih menaungi 60\% pekerjaan pada 2010,
  bagiannya turun 12 poin persentase dari 72\% pada 1995. Sebaliknya,
  bagian pusat terencana naik dari sekitar 28\% menjadi 40\%. (Locator:
  Scatteration or Polycentricity? dan Tabel 3, hlm. 963-964.)
\item
  Pusat industri berkembang menuju pembagian fungsi sektoral: industri
  berteknologi relatif tinggi terkonsentrasi pada kawasan terencana
  berskala besar di timur, industri padat karya dan mencemari
  terkonsentrasi pada pusat yang lebih tua dan terfragmentasi di barat,
  sedangkan Cilegon mempertahankan peran pada industri berat dan
  ekstraktif. (Locator: Specialization of Industrial Centers in Greater
  Jakarta, Tabel 5 dan pembahasan, hlm. 965-966.)
\end{itemize}

\textbf{Interpretasi penulis.} Penulis menganggap hubungan positif
antara pertumbuhan pekerjaan dan jarak dari CBD sebagai indikasi kuat
bahwa manufaktur formal telah meluas dan semakin tersuburbanisasi.
Mereka juga menilai kawasan industri swasta telah menangkap sebagian
besar pekerjaan yang terdispersi dari Jakarta, mengurangi pola tersebar,
dan mengarahkan wilayah menuju struktur yang relatif lebih polisentrik.
Spesialisasi serta diversifikasi pusat pinggiran dipandang sebagai
pengambilalihan fungsi manufaktur utama yang sebelumnya berada di CBD.
(Locator: The Dynamics of Manufacturing Employment in Greater Jakarta,
hlm. 962-963; Scatteration or Polycentricity?, hlm. 964; Conclusion,
hlm. 967.)

\textbf{Inferensi pengolahan.} Bukti deskriptif menunjukkan
dekonsentrasi dan reaglomerasi pekerjaan, tetapi tidak cukup untuk
menyatakan bahwa pengembang swasta menyebabkan seluruh perubahan. Data
juga memperlihatkan struktur hibrida, yaitu penguatan pusat terencana
bersamaan dengan bertahannya mayoritas pekerjaan di luar pusat tersebut.

\subsection{Conclusions}\label{conclusions-2}

\textbf{Kesimpulan sumber.} Para penulis menyimpulkan bahwa
dekonsentrasi industri Greater Jakarta telah mencapai fase lanjut. Sejak
awal 1990-an, sebagian besar pekerjaan manufaktur telah
tersuburbanisasi; proses itu tampak pada penurunan peran CBD dan
meluasnya batas fungsional metropolitan. Pemerintah dinilai berperan
relatif terbatas sebagai fasilitator investasi serta pembangunan swasta,
sedangkan sektor swasta memainkan peran substansial dalam menyertai
dekonsentrasi pekerjaan manufaktur formal. Penulis tidak menemukan
indikasi bahwa perencanaan spasial memainkan peran krusial dalam
mengarahkan proses tersebut. (Locator: Conclusion, hlm. 967.)

Struktur hasilnya belum sepenuhnya polisentrik. Sebagian besar pekerjaan
formal masih terletak pada koridor zona industri tidak terencana dan
kantong yang tersebar, tetapi kecenderungan itu menurun. Industri
berpindah ke pusat pinggiran tempat kawasan industri swasta
beraglomerasi; pusat tersebut semakin beragam dan mengambil alih sektor
utama CBD. Karena itu, penulis menyimpulkan bahwa susunan spasial
berubah dari tersebar menuju relatif lebih polisentrik. (Locator:
Conclusion, hlm. 967.)

\textbf{Interpretasi penulis.} Perubahan tersebut dipandang sebagai
peluang untuk membentuk metropolitan yang lebih terencana, tetapi
dominasi investor asing dan pengembang besar tanpa kerangka serta
lembaga perencanaan yang memadai dapat memperkuat segregasi sosial,
fragmentasi infrastruktur, dan penurunan mutu lingkungan fisik.
(Locator: Conclusion, hlm. 967.)

\textbf{Inferensi pengolahan.} Kesimpulan tentang perubahan pola spasial
konsisten dengan indikator E/P, pembagian bentuk spasial, dan LQ.
Kesimpulan mengenai peran kausal sektor swasta harus dibaca sebagai
penjelasan yang dikuatkan secara deskriptif, sesuai batas metode yang
dinyatakan penulis sendiri.

\subsection{Contributions}\label{contributions-2}

Kontribusi yang diberi label secara eksplisit oleh penulis: Tidak
disebutkan dalam file.

\textbf{Inferensi pengolahan berdasarkan isi artikel.} Kontribusi
artikel dapat dirumuskan sebagai berikut:

\begin{itemize}
\tightlist
\item
  Memperbarui bukti studi 1990-an dengan gambaran longitudinal pekerjaan
  manufaktur formal pada 1995, 2001, dan 2010. (Locator: Introduction,
  hlm. 951; Data and Methods, hlm. 954; Tabel 2-3, hlm. 963-964.)
\item
  Memperluas unit geografis dari Jabodetabek ke Greater Jakarta agar
  pusat industri di Serang dan Karawang tidak terhapus oleh batas
  metropolitan formal. (Locator: From Jabodetabek to Greater Jakarta,
  hlm. 955-957; Gambar 1-2.)
\item
  Menggabungkan sejarah kebijakan, data pasar lahan, E/P, tipologi
  bentuk spasial, dan LQ untuk menghubungkan pembangunan lahan industri
  dengan dekonsentrasi serta diferensiasi fungsi metropolitan. (Locator:
  Data and Methods, hlm. 953-954; Gambar 3-7; Tabel 1-5.)
\item
  Menawarkan varian kasus Asia Tenggara dalam perdebatan dekonsentrasi:
  negara memfasilitasi investasi, tetapi pengembang swasta lebih dominan
  dalam membentuk geografi pinggiran, berbeda dari gambaran negara yang
  lebih aktif pada kasus Tiongkok. (Locator: Deconcentration, hlm.
  952-953; Conclusion, hlm. 967.)
\item
  Menunjukkan bahwa dikotomi polisentrisitas versus penyebaran tidak
  harus menghasilkan satu kategori mutlak karena Greater Jakarta
  memperlihatkan reaglomerasi pusat terencana sekaligus persistensi
  pekerjaan yang tersebar. (Locator: Tabel 3 dan pembahasan, hlm.
  963-964.)
\end{itemize}

\subsection{Research Gap}\label{research-gap-2}

\textbf{Kesenjangan yang dinyatakan sumber.}

\begin{itemize}
\tightlist
\item
  Peran restrukturisasi pekerjaan dalam suburbanisasi Indonesia kurang
  diperhatikan dibandingkan dekonsentrasi penduduk. (Locator:
  Introduction, hlm. 951.)
\item
  Literatur Barat terlalu berpusat pada jasa dalam masyarakat
  pascaindustri untuk langsung menjelaskan metropolis Asia Timur yang
  perubahan spasialnya banyak digerakkan manufaktur. (Locator:
  Introduction, hlm. 951; Deconcentration, hlm. 952.)
\item
  Kesimpulan penelitian pada 1990-an bahwa Jakarta tetap menjadi pusat
  manufaktur korporat utama perlu diperbarui setelah pertumbuhan kawasan
  industri pinggiran dan arus investasi asing langsung. (Locator:
  Introduction, hlm. 951.)
\item
  Kajian yang hanya memakai batas Jabodetabek tidak lagi menangkap
  perkembangan industri ke Serang dan Karawang. (Locator: From
  Jabodetabek to Greater Jakarta, hlm. 955-956.)
\item
  Belum jelas sejauh mana pusat industri swasta mengurangi dominasi
  pusat Jakarta dan apakah dekonsentrasi menghasilkan penyebaran atau
  pusat-pusat baru. (Locator: Implications for Metropolitan Spatial
  Structure and Pattern, hlm. 953.)
\end{itemize}

\textbf{Kesenjangan tersisa menurut inferensi pengolahan.} Artikel belum
mengidentifikasi efek kausal kawasan industri swasta; belum mencakup
perusahaan kecil dan pekerjaan informal; belum mengukur arus komuter
atau interaksi antarpusat untuk menguji polisentrisitas fungsional;
serta belum menguji dampak sosial, infrastruktur, dan lingkungan yang
disebut pada kesimpulan. Data pekerjaan juga berakhir pada 2010 sehingga
tidak menjelaskan kelanjutan pola setelah periode itu. Ini adalah
inferensi dari cakupan dan batas metode, bukan agenda riset yang
dinyatakan penulis.

\subsection{Limitations}\label{limitations-2}

Bagian keterbatasan khusus: Tidak disebutkan dalam file. Namun, sumber
menyatakan atau memperlihatkan beberapa batas berikut.

\textbf{Batas yang dinyatakan sumber.}

\begin{itemize}
\tightlist
\item
  Analisis bersifat deskriptif dan eksploratif serta tidak bermaksud
  memverifikasi proses kausal secara empiris. (Locator: Data and
  Methods, hlm. 953-954.)
\item
  Sensus industri hanya mencakup perusahaan formal menengah dan besar
  dengan lebih dari 20 pekerja. (Locator: Data and Methods, hlm. 954.)
\item
  LQ yang tinggi pada pusat dengan sedikit pekerjaan dapat melebihkan
  kesan pentingnya spesialisasi. Cilegon hanya memiliki sekitar 5.000
  pekerjaan pada 1995 dan 10.000 pada 2010, jauh di bawah rata-rata
  pusat industri sekitar 56.000-91.000 pekerjaan. (Locator: Data and
  Methods, hlm. 954; catatan Tabel 4-5, hlm. 965-966.)
\item
  Keterbatasan data mengharuskan sektor 1995 digabung menjadi sembilan
  kelompok, sedangkan 2010 memakai 16 sektor setelah penyelarasan
  klasifikasi. (Locator: Specialization of Industrial Centers in Greater
  Jakarta, hlm. 964.)
\end{itemize}

\textbf{Inferensi pengolahan.} Tiga tahun pengamatan hanya memberi
potret berkala, sehingga waktu tepat relokasi perusahaan tidak dapat
dilihat. Penggunaan batas E/P 0,1 menghasilkan klasifikasi yang peka
terhadap ambang, tetapi artikel tidak melaporkan uji sensitivitas.
Harmonisasi 145 zona merespons perubahan batas kecamatan, tetapi rincian
prosedurnya tidak tersedia untuk menilai sepenuhnya keterbandingan
temporal. Artikel juga tidak melaporkan ketidakpastian pengukuran, data
hilang, atau validasi silang antar-sumber. Prediksi polisentrisitas
bergantung pada keberlanjutan tren dan bukan hasil pengujian prospektif.

\subsection{Challenges}\label{challenges-2}

Bagian bernama khusus ``Challenges'': Tidak disebutkan dalam file.

\textbf{Tantangan substantif yang dilaporkan sumber.} Kerangka dan
lembaga perencanaan metropolitan dinilai tidak memadai, perencanaan
spasial terfragmentasi, dan pemerintah daerah pinggiran kekurangan
kapasitas untuk mengantisipasi suburbanisasi cepat. Ekspansi yang banyak
dikendalikan investor asing serta pengembang besar berpotensi memperkuat
segregasi sosial, memecah penyediaan infrastruktur, dan menurunkan
kualitas lingkungan fisik. Kelangkaan lahan kosong, kemacetan Jakarta,
serta pergeseran kegiatan melampaui batas Jabodetabek juga memperumit
koordinasi lintas wilayah. (Locator: Major Factors and Actors Triggering
Industrial Deconcentration, hlm. 957-960; Conclusion, hlm. 967.)

\textbf{Tantangan analitis menurut inferensi pengolahan.} Perubahan
batas kecamatan, perubahan klasifikasi industri, ketimpangan ukuran
pusat, serta keterbatasan ekstraksi pola visual dari peta menyulitkan
perbandingan longitudinal dan interpretasi LQ. Tantangan terakhir
berkaitan dengan file yang ditelaah dalam review ini, bukan tantangan
yang dinyatakan penulis.

\subsection{Practical Implications}\label{practical-implications-2}

\textbf{Implikasi yang dinyatakan sumber.} Penulis menyerukan
peningkatan kapasitas pemerintah daerah pinggiran dan tata kelola
metropolitan agar mampu merumuskan, menyesuaikan, mengoordinasikan,
serta melaksanakan rencana tata guna lahan lokal dan regional. Rencana
tersebut perlu diintegrasikan dengan kebijakan ekonomi regional dan
rencana jangka panjang pemerintah nasional. (Locator: Conclusion, hlm.
967.)

Perubahan menuju pusat industri yang lebih terkonsentrasi dipandang
sebagai peluang bagi metropolitan yang lebih terencana. Namun, peluang
itu memerlukan kerangka institusional yang dapat mengendalikan risiko
segregasi, fragmentasi infrastruktur, dan kerusakan lingkungan, bukan
sekadar mengandalkan penyediaan kawasan oleh swasta. (Locator:
Conclusion, hlm. 967.)

\textbf{Inferensi pengolahan.} Cakupan perencanaan praktis perlu
mengikuti wilayah fungsional Greater Jakarta, termasuk Serang dan
Karawang, bukan berhenti pada batas Jabodetabek. Prioritas infrastruktur
dan layanan juga dapat diarahkan pada pusat yang tumbuh cepat di
pinggiran luar dan tambahan, tetapi artikel tidak menguji intervensi
kebijakan tertentu.

\subsection{Applications}\label{applications-2}

Aplikasi yang dinyatakan atau diuji secara khusus: Tidak disebutkan
dalam file.

\textbf{Inferensi pengolahan.} Pendekatan artikel dapat diterapkan
secara terbatas untuk:

\begin{itemize}
\tightlist
\item
  Menggunakan E/P guna menyaring kecamatan yang berfungsi sebagai pusat
  pekerjaan manufaktur meskipun kepadatan absolutnya rendah.
\item
  Menggunakan LQ untuk membandingkan spesialisasi sektoral antarpusat
  setelah klasifikasi usaha diselaraskan.
\item
  Memadukan peta pusat pekerjaan, lokasi kawasan industri, pasokan
  lahan, dan harga lahan untuk menetapkan cakupan perencanaan
  metropolitan serta kebutuhan koordinasi infrastruktur.
\item
  Memantau apakah pertumbuhan pekerjaan bergerak menuju pusat terencana
  atau tetap menyebar pada koridor tidak terencana.
\end{itemize}

Aplikasi tersebut merupakan kemungkinan penggunaan metode, bukan
efektivitas aplikasi yang dievaluasi dalam artikel. (Locator dasar
inferensi: Data and Methods, hlm. 953-954; Gambar 4-7; Tabel 1-5.)

\subsection{Future Research
(Optional)}\label{future-research-optional-2}

Agenda penelitian lanjutan yang dinyatakan secara eksplisit setelah
hasil penelitian: Tidak disebutkan dalam file. Pernyataan pada
pendahuluan bahwa studi mutakhir perlu dilakukan merupakan motivasi
artikel ini, bukan agenda pascapenelitian. (Locator: Introduction, hlm.
951.)

\textbf{Inferensi pengolahan, bukan usulan penulis.} Penelitian
selanjutnya dapat menguji hubungan kausal antara kebijakan, pembangunan
kawasan industri swasta, jalan tol, harga lahan, investasi asing
langsung, dan relokasi pekerjaan; memperbarui data setelah 2010;
memasukkan perusahaan kecil serta sektor informal; menguji
polisentrisitas fungsional melalui arus komuter dan hubungan antarpusat;
melakukan uji sensitivitas terhadap ambang E/P; dan mengukur dampak
segregasi, fragmentasi infrastruktur, serta kualitas lingkungan yang
baru disebut sebagai risiko.

\subsection{Provenance Notes}\label{provenance-notes-2}

\begin{itemize}
\tightlist
\item
  Review ini hanya menggunakan seluruh isi file TXT
  \texttt{source/journal/J03-hudalah-viantari-firman-woltjer-2013-industrial-land-development-and-manufacturing-deconcentration-in-greater-jakarta-10.1080-02723638.2013.783281.txt},
  termasuk blok metadata, teks artikel, tabel, keterangan gambar,
  kesimpulan, dan daftar referensi. Tidak ada sumber eksternal yang
  digunakan.
\item
  Metadata file menyatakan bahwa TXT diekstrak pada 31 Agustus 2026 dari
  PDF bernama sama dengan metode \texttt{pdftotext\ -layout}. Metadata
  PDF mencatat 23 halaman, sedangkan artikel menggunakan pagination
  jurnal 950-971. Dalam TXT, materi awal platform tampil sebelum halaman
  950.
\item
  Akses review ini adalah terhadap teks lengkap yang tersedia dalam TXT,
  bukan pemeriksaan langsung PDF. Akurasi ekstraksi, tata letak, dan
  kesesuaian dengan PDF asli tidak diverifikasi secara independen.
\item
  Locator \texttt{hlm.} merujuk pada nomor halaman tercetak artikel.
  Locator bagian, tabel, dan gambar mengikuti label dalam TXT.
\item
  Gambar 1-7 tidak tersedia sebagai visual utuh dalam TXT. Karena itu,
  review hanya memakai keterangan gambar dan uraian penulis, bukan
  pembacaan simbol, warna, batas, atau nilai langsung dari peta dan
  grafik.
\item
  Nilai tabel digunakan sebagaimana terbaca dalam ekstraksi. Tabel 2
  mencatat penurunan jumlah pekerjaan Jakarta sebesar 15\% pada kolom
  pertumbuhan, sedangkan narasi menyebut penurunan bagian pekerjaan
  sebesar 15\%. Review tidak menyelaraskan perbedaan istilah tersebut.
\item
  Teks analisis LQ 1995 menyebut Cikarang sebagai salah satu dari lima
  pusat, tetapi kepala kolom Tabel 4 dalam TXT bertuliskan Cibitung.
  Review mempertahankan ketidakselarasan ini dan tidak menganggap
  keduanya identik tanpa penjelasan sumber.
\item
  Rujukan, data lembaga, dan klaim komparatif yang dikutip artikel tidak
  diperiksa terhadap dokumen asal. Status telaah sejawat artikel juga
  tidak dinyatakan dalam file dan tidak diasumsikan.
\item
  Setiap bagian berlabel \textbf{Laporan sumber}, \textbf{Temuan yang
  dilaporkan sumber}, \textbf{Kesimpulan sumber}, atau
  \textbf{Interpretasi penulis} merepresentasikan isi atau penalaran
  artikel. Bagian berlabel \textbf{Inferensi pengolahan} adalah
  pembacaan analitis review ini dan bukan klaim eksplisit penulis.
\end{itemize}

\section{Review J04}\label{review-j04}

\textbf{Identitas sumber}

\begin{itemize}
\tightlist
\item
  Kode sumber: J04.
\item
  Penulis: Delik Hudalah dan Tommy Firman.
\item
  Judul: ``Beyond property: Industrial estates and post-suburban
  transformation in Jakarta Metropolitan Region''.
\item
  Afiliasi penulis: School of Architecture, Planning and Policy
  Development, Bandung Institute of Technology (ITB), Indonesia.
\item
  Publikasi: \emph{Cities}, volume 29, tahun 2012, halaman 40-48.
\item
  Nomor terbitan: Tidak disebutkan dalam file.
\item
  DOI: \texttt{10.1016/j.cities.2011.07.003}.
\item
  Riwayat artikel: diterima 1 Februari 2011, direvisi 20 April 2011,
  diterima untuk publikasi 14 Juli 2011, dan tersedia daring 15 Agustus
  2011 (hlm. 40, blok \emph{article info}).
\item
  Kata kunci: \emph{industrial estate}, Jakarta, \emph{post-suburbia},
  \emph{suburbanization}, dan \emph{urban deconcentration} (hlm. 40,
  blok \emph{article info}).
\item
  Jenis dokumen: artikel jurnal. Status penelaahan sejawat: Tidak
  disebutkan dalam file.
\item
  Dasar akses review: teks lengkap hasil ekstraksi dari PDF sembilan
  halaman, bukan pemeriksaan langsung terhadap PDF atau sumber-sumber
  yang dirujuk artikel.
\end{itemize}

\subsection{Summarized Abstract}\label{summarized-abstract-3}

\textbf{Laporan sumber:} Artikel mengkaji apakah transformasi kawasan
pinggiran Jakarta Metropolitan Region (JMR) telah mengikuti
kecenderungan global menuju \emph{post-suburbia}, yaitu desentralisasi
kehidupan perkotaan ke pinggiran metropolitan, serta sejauh mana
perkembangan industri berkontribusi terhadap transformasi tersebut.
Unsur yang dilaporkan meliputi dekonsentrasi terencana industri
teknologi tinggi dan perusahaan multinasional, perkembangan kawasan
industri menjadi pusat perkotaan, serta proyek yang melampaui penjualan
tanah dan properti, termasuk fasilitas komersial, rekreasi, kesehatan,
teknologi, dan kebudayaan bercorak Barat (hlm. 40, abstrak; hlm. 45-46,
bagian ``Jababeka City: From industrial park to suburban center'').

\textbf{Interpretasi penulis:} Penulis menilai bahwa JMR mulai
memperlihatkan tahap awal \emph{post-suburbia}, tetapi bentuknya tidak
identik dengan pengalaman Amerika Serikat, Eropa, atau China.
Kekhasannya dikaitkan dengan kapasitas inti metropolitan yang menurun,
peran negara yang tetap ada tetapi relatif pasif dan tidak langsung,
hubungan simbiotik pasar-sektor publik, serta privatisasi aturan
perencanaan oleh pengembang (hlm. 40, abstrak; hlm. 47, bagian ``The end
of desakota regions?'').

\textbf{Inferensi review:} Artikel paling tepat dibaca sebagai analisis
konseptual-empiris multiskala: pola pada tingkat JMR dijelaskan melalui
Kabupaten Bekasi dan Cikarang, kemudian diperdalam dengan kasus
Jababeka. Bukti yang disajikan mendukung keberadaan beberapa unsur
\emph{post-suburbia}, tetapi tidak merupakan pengujian kausal atau
pengukuran formal atas tingkat \emph{post-suburbanization}.

\subsection{Problem Statement}\label{problem-statement-3}

\textbf{Laporan sumber:} Kajian terdahulu tentang urbanisasi Asia Timur
dan Indonesia umumnya memahami pinggiran kota besar sebagai perpanjangan
inti metropolitan dan sebagai zona transisi perkotaan-perdesaan. Penulis
menyatakan bahwa kemunculan kota global dan desentralisasi pada beberapa
dekade terakhir menantang konsepsi tersebut (hlm. 40, abstrak dan bagian
``Introduction'').

\textbf{Laporan sumber:} Dampak pembangunan perumahan di pinggiran
metropolitan Indonesia telah banyak diakui, tetapi peran perkembangan
industri dalam membentuk pusat perkotaan baru yang mampu menandingi daya
tarik pusat kota lama masih kurang diteliti. Perubahan politik setelah
jatuhnya rezim Soeharto juga membuka ruang lebih besar bagi pemerintah
lokal dan sektor swasta dalam perkembangan industri (hlm. 40, bagian
``Introduction'').

\textbf{Interpretasi penulis:} Masalah analitisnya bukan hanya apakah
fungsi perkotaan telah berpindah ke pinggiran, melainkan apakah
perubahan itu menandai bentuk \emph{post-suburbia} yang berbeda dari
\emph{desakota} dan apa maknanya dalam wilayah metropolitan yang masih
mengalami industrialisasi (hlm. 40-41, bagian ``Introduction''; hlm. 47,
bagian ``The end of desakota regions?'').

\subsection{Objectives}\label{objectives-3}

\textbf{Laporan sumber:} Artikel menetapkan dua tujuan utama:
mengeksplorasi apakah transformasi perkotaan di pinggiran JMR telah
bergabung dengan kecenderungan global \emph{post-suburbia} dan
menggambarkan sejauh mana perkembangan industri, sebagai sektor ekonomi
utama kawasan, berkontribusi terhadap transformasi itu (hlm. 40, bagian
``Introduction'').

\textbf{Laporan sumber:} Tujuan tersebut dijabarkan menjadi tiga
pertanyaan: apakah \emph{post-suburbia} dapat diidentifikasi di JMR;
apakah perkembangan industri berperan dalam transformasi tersebut; dan
apa makna atau ciri unik \emph{post-suburbia} dalam konteks JMR yang
mengalami industrialisasi (hlm. 40, bagian ``Introduction'').

\subsection{Literature Survey}\label{literature-survey-3}

\textbf{Laporan sumber:} Tinjauan konseptual menelusuri pendorong
\emph{post-suburbia} dari perkembangan mobil, telekomunikasi, dan
komputer; desentralisasi pekerjaan; munculnya fungsi belanja dan
industri teknologi tinggi di pinggiran; serta peran kompleks industri
dalam mendahului atau menopang suburbanisasi perumahan. Globalisasi
neoliberal kemudian memperkuat fleksibilitas lokasi kegiatan ekonomi dan
perubahan gaya hidup (hlm. 41, bagian ``The drivers of post-suburbia'').

\textbf{Laporan sumber:} Literatur yang dirangkum mengidentifikasi
bentuk spasial \emph{post-suburbia} sebagai lanskap antara kota dan desa
yang tidak lagi sekadar zona transisi. Bentuk ini cenderung berkepadatan
rendah, terfragmentasi secara struktural, terspesialisasi secara
ekonomi, tersegregasi secara sosial, dan berpola polisentrik dengan
ikatan yang melemah terhadap inti metropolitan (hlm. 41, bagian
``Elements of post-suburbia'').

\textbf{Laporan sumber:} Artikel membandingkan variasi tata kelola.
Literatur Amerika Serikat menekankan individualisme, kebebasan,
penggunaan mobil, dan berkurangnya visibilitas negara; pengalaman Eropa
masih memberikan peran penting kepada perencanaan kawasan-kota dan
kemitraan publik-swasta; sedangkan pemerintah pusat dan lokal di China
dapat bertindak lebih proaktif sebagai perencana, pengusaha, dan
pengembang pinggiran (hlm. 41-42, bagian ``Elements of post-suburbia'').

\textbf{Laporan sumber:} Untuk Asia Timur, artikel memakai
\emph{desakota} sebagai pembanding utama. \emph{Desakota} digambarkan
sebagai transformasi bertahap, insidental, dan tidak terencana pada
kawasan pertanian di koridor antara kota dan hinterland perdesaan,
termasuk perkembangan industri yang tidak terencana di sawah (hlm. 47,
bagian ``The end of desakota regions?'').

\textbf{Interpretasi penulis:} Kerangka komparatif tersebut digunakan
untuk menilai apakah perkembangan terencana kawasan industri dan fungsi
perkotaan baru di JMR mengikis kekhasan model urbanisasi meluas Asia
atau justru menghasilkan varian lokal \emph{post-suburbia} (hlm. 41,
akhir bagian ``Introduction''; hlm. 47, bagian ``The end of desakota
regions?'').

\textbf{Inferensi review:} Tinjauan pustaka bersifat naratif dan
selektif. File tidak menyebutkan strategi pencarian literatur, kriteria
inklusi, jumlah korpus, atau prosedur penilaian mutu referensi.

\subsection{Method Used}\label{method-used-3}

\textbf{Laporan sumber:} Tidak ada bagian metode formal dalam file.
Artikel menguraikan pola dekonsentrasi perkotaan pada tingkat JMR,
memilih Kabupaten Bekasi untuk analisis lebih terperinci, lalu
memusatkan perhatian pada Jababeka sebagai kawasan industri terbesar dan
paling dinamis di kabupaten tersebut. Kabupaten Bekasi dipilih karena
restrukturisasi ekonominya berlangsung dalam wilayah metropolitan yang
makin urban, tetapi administrasinya masih berupa pemerintahan kabupaten
yang berorientasi perdesaan; kondisi ini dinilai berpotensi
memperlihatkan kompleksitas tata kelola \emph{post-suburbia} (hlm. 41,
bagian ``Introduction'').

\textbf{Laporan sumber:} Analisis menggunakan sensus, statistik tahunan,
data investasi asing langsung, peta penggunaan lahan historis, dokumen
kebijakan, data pemerintah daerah, serta dokumen perusahaan dan
organisasi kawasan industri. Artikel menghitung atau menyesuaikan
beberapa data menjadi Gambar 2 dan Tabel 1-4 (hlm. 42-46, bagian ``Urban
and industrial development in the Jakarta Metropolitan Region'' sampai
``Socio-economic and spatial implications and governance adaptation'').

\textbf{Laporan sumber:} Dalam ucapan terima kasih, artikel dinyatakan
sebagai bagian dari program riset aksi terpadu tentang \emph{Industrial
Linkages} (PHK-I 2010); informasi dan data para peneliti program itu
disebut memungkinkan analisis. Namun, file tidak menjelaskan prosedur
riset aksi, kerja lapangan, wawancara, observasi, atau cara data
tersebut dikumpulkan (hlm. 47, bagian ``Acknowledgments'').

\textbf{Interpretasi penulis:} Pola empiris dibandingkan dengan unsur
\emph{post-suburbia} dari Amerika Serikat, Eropa, dan Asia Timur,
kemudian ditafsirkan terhadap konsep \emph{desakota} dan konteks
politik-ekonomi JMR (hlm. 41-42, bagian ``Post-suburbia: A global
perspective''; hlm. 47, bagian ``The end of desakota regions?'').

\textbf{Inferensi review:} Berdasarkan langkah yang tampak dalam teks,
desainnya dapat diklasifikasikan sebagai studi kasus
deskriptif-interpretatif multiskala yang didukung analisis dokumen dan
statistik sekunder longitudinal. Klasifikasi ini merupakan inferensi
review karena penulis tidak memberi nama eksplisit pada desain
penelitian maupun protokol analisisnya.

\subsection{Dataset}\label{dataset-3}

\textbf{Laporan sumber:} Bahan kuantitatif dan dokumenter yang
disebutkan meliputi:

\begin{itemize}
\tightlist
\item
  Hasil sensus nasional sepuluh tahunan 1961-2000 dari Biro Pusat
  Statistik dan data kependudukan DKI Jakarta untuk menjelaskan
  pergeseran distribusi penduduk JMR (hlm. 42, bagian ``Urban and
  industrial development in the Jakarta Metropolitan Region'').
\item
  Peta penggunaan lahan historis dari Rustiadi (2007), termasuk analisis
  ekspansi lahan terbangun JMR pada 1983-2005 dan angka konversi lahan
  pada 1992-2000 serta 2000-2005 (hlm. 43, Gambar 2).
\item
  Data tahunan BPS 1985-2005 yang dihitung untuk membandingkan
  kontribusi sektoral Jakarta dan wilayah pinggiran terhadap GDRP/PDRB
  JMR (hlm. 43, bagian ``Urban and industrial development in the Jakarta
  Metropolitan Region'').
\item
  Data tahunan Badan Koordinasi Penanaman Modal 1998-2009 tentang
  investasi asing langsung atau FDI, diringkas dalam empat periode tiga
  tahunan untuk sektor sekunder dan tersier (hlm. 43-44, Tabel 1 dan
  Tabel 2).
\item
  Data kawasan industri Cikarang yang disesuaikan dari Shahab (2010) dan
  Provinsi Jawa Barat (2010), mencakup luas industri, luas rencana,
  jumlah penyewa industri, dan pekerja (hlm. 44-45, Tabel 3).
\item
  Data Kabupaten Bekasi (2009) tentang luas, penduduk, kepadatan,
  GDRP/PDRB, dan pendapatan per kapita pada tiga kecamatan industri, 12
  kecamatan pinggiran, enam kecamatan perdesaan, keseluruhan kabupaten,
  dan Jakarta (hlm. 46-47, Tabel 4).
\item
  Keputusan Presiden 53/1989 dan 41/1996; dokumen rencana Provinsi Jawa
  Barat 2009 dan 2010; bahan ZONI 2007; serta publikasi Jababeka
  2008-2010 dan PT Jababeka Infrastruktur 2009 untuk merekonstruksi
  konteks regulasi, rencana zona ekonomi khusus, infrastruktur, dan
  proyek ``beyond property'' (hlm. 44-47, bagian ``Economic and
  political background and policy contexts'' sampai ``Rethinking the
  governance structure'').
\end{itemize}

\textbf{Inferensi review:} Artikel tidak membentuk satu himpunan data
tunggal. Ia menggabungkan seri dengan tahun, unit, dan asal berbeda.
File tidak menyediakan data mentah, kode perhitungan, definisi rinci
setiap sektor, prosedur harmonisasi batas wilayah, ataupun lampiran
replikasi.

\subsection{Population / Sample}\label{population-sample-3}

\textbf{Laporan sumber:} Tidak terdapat populasi partisipan manusia atau
sampel survei dalam file. Unit analisis utamanya bersifat spasial,
ekonomi, dan kelembagaan.

\textbf{Laporan sumber:} Cakupan terluas adalah JMR seluas 5.897 km2,
yang terdiri atas DKI Jakarta dan Bodetabek: tiga kabupaten otonom
(Bogor, Tangerang, dan Bekasi) serta lima kota (Bogor, Depok, Tangerang,
Tangerang Selatan, dan Bekasi). Artikel kemudian mempersempit kajian ke
Kabupaten Bekasi, Cikarang, tujuh kawasan/kota industri, dan akhirnya
Jababeka (hlm. 42, bagian ``Urban and industrial development in the
Jakarta Metropolitan Region'' dan Gambar 1; hlm. 44-45, Tabel 3).

\textbf{Laporan sumber:} Pemilihan kasus bersifat bertujuan. Kabupaten
Bekasi dipilih karena restrukturisasi ekonominya yang progresif
sekaligus kompleksitas status administratifnya sebagai kabupaten;
Jababeka dipilih karena disebut sebagai kota industri terbesar dan
paling dinamis di Kabupaten Bekasi (hlm. 41, bagian ``Introduction'').
Tidak disebutkan dalam file adanya kerangka sampling statistik atau
kriteria pemilihan pembanding.

\textbf{Laporan sumber:} Unit yang dihitung mencakup tujuh perusahaan
kawasan industri di Cikarang dengan 2.288 perusahaan penyewa dan 532.814
pekerja secara keseluruhan. Untuk analisis intra-Kabupaten Bekasi, 21
kecamatan dikelompokkan menjadi tiga kecamatan industri di Cikarang, 12
kecamatan hinterland pinggiran, dan enam kecamatan perdesaan (hlm. 45,
Tabel 3; hlm. 46-47, bagian ``Socio-economic and spatial implications
and governance adaptation'' dan Tabel 4).

\textbf{Laporan sumber:} Angka penduduk yang menjadi konteks, bukan
sampel penelitian, meliputi perkiraan 22,5 juta penduduk JMR pada 2000
dengan 12,8 juta berada di Bodetabek, serta 2.125.960 penduduk Kabupaten
Bekasi dalam Tabel 4 (hlm. 42, bagian ``Urban and industrial development
in the Jakarta Metropolitan Region''; hlm. 46, Tabel 4).

\subsection{Dependent Variables}\label{dependent-variables-3}

\textbf{Laporan sumber:} Tidak disebutkan dalam file. Artikel tidak
merumuskan variabel dependen, variabel independen, hipotesis statistik,
atau model kausal formal.

\textbf{Inferensi review:} Sejumlah indikator berfungsi sebagai hasil
analitis, tetapi bukan variabel dependen yang ditetapkan penulis:
dekonsentrasi penduduk; perluasan lahan terbangun; perubahan kontribusi
sektoral terhadap GDRP/PDRB; distribusi FDI; pertumbuhan kawasan
industri dan fungsi perkotaan; pendapatan per kapita; fragmentasi
infrastruktur; segregasi sosial; serta adaptasi tata kelola (hlm. 42-47,
Gambar 2 dan Tabel 1-4). Mengubah indikator-indikator ini menjadi
variabel dependen formal akan melampaui desain yang dilaporkan.

\subsection{Results}\label{results-3}

\textbf{Laporan sumber:} Distribusi penduduk JMR bergeser dari
konsentrasi pada Jakarta pada 1960-an dan 1970-an menuju dekonsentrasi
sejak 1980-an. Pada 1990-an, penduduk pinggiran telah melampaui penduduk
kota inti, meskipun jumlah penduduk inti tetap bertambah dengan laju
yang menurun. Pada 2000, populasi JMR diperkirakan mencapai 22,5 juta,
dengan 12,8 juta penduduk tersebar di Bodetabek (hlm. 42, bagian ``Urban
and industrial development in the Jakarta Metropolitan Region'').

\textbf{Laporan sumber:} Ekspansi fisik mengikuti pergeseran demografis.
Sekitar 90.760 ha lahan, umumnya pertanian, dikonversi menjadi lahan
terbangun pada masa puncak 1992-2000. Setelah krisis Asia 1997
memperlambat ekspansi, sekitar 22.873,04 ha kembali dikonversi pada
2000-2005 (hlm. 43, Gambar 2 dan uraian yang menyertainya).

\textbf{Laporan sumber:} Kegiatan tersier tetap terkonsentrasi di
Jakarta dalam data 1985-2005: wilayah pinggiran rata-rata menyumbang
8,7\% GDRP/PDRB regional untuk keuangan dan jasa serta 22,5\% untuk
sektor komersial, termasuk hotel dan restoran. Sebaliknya, kontribusi
pinggiran pada manufaktur meningkat dari 24,6\% pada 1985-1990 menjadi
59,8\% pada 2000-2005 (hlm. 43, bagian ``Urban and industrial
development in the Jakarta Metropolitan Region'').

\textbf{Laporan sumber:} Bodetabek menyerap 84,9\%, 83,8\%, 87,0\%, dan
87,3\% FDI sektor sekunder JMR masing-masing pada 1998-2000, 2001-2003,
2004-2006, dan 2007-2009. Kabupaten Bekasi sendiri menyerap 52,4\%,
51,5\%, 45,0\%, dan 46,9\% pada empat periode tersebut. Untuk sektor
tersier, bagian Bodetabek meningkat konsisten dari 5,4\% menjadi 14,9\%,
sementara bagian Kabupaten Bekasi naik dari 2,5\% menjadi 7,0\% (hlm.
43-44, Tabel 1 dan Tabel 2).

\textbf{Laporan sumber:} Mengacu pada sumber Kabupaten Bekasi (2009),
manufaktur dilaporkan menyumbang 79,73\% GDRP/PDRB kabupaten; tahun
observasi angkanya tidak dijelaskan lebih lanjut dalam file. Pertumbuhan
ekonomi daerah dilaporkan 7,42\%, dibandingkan pertumbuhan nasional
5,6\%. Lokasi kawasan industri didukung akses ke Jalan Tol
Jakarta-Cikampek, jarak sekitar 40 km dari pusat Jakarta, kedekatan
relatif dengan Pelabuhan Tanjung Priok, pasokan air Jatiluhur, harga
lahan pertanian yang lebih rendah, serta rencana tata ruang yang
mengantisipasi industri (hlm. 44, bagian ``Industrial estate development
in Cikarang, Kabupaten Bekasi'').

\textbf{Laporan sumber:} Tujuh kawasan/kota industri Cikarang memiliki
total luas industri tercatat 3.866 ha, luas perencanaan 14.620 ha, 2.288
penyewa industri, dan 532.814 pekerja. Empat kawasan berfokus pada
pengembangan dan pengelolaan industri, sedangkan Jababeka, Lippo
Cikarang, dan Deltamas memadukan industri dengan perumahan serta fungsi
perkotaan lain. Data luas industri dan pekerja Deltamas ditandai tidak
tersedia (hlm. 44-45, Tabel 3 dan uraian setelah tabel).

\textbf{Laporan sumber:} Jababeka berkembang dari pengembang kawasan
industri pada 1989, memasuki pembangunan kota pada 1996, dan pada 2007
mengalihkan bisnis intinya menuju kegiatan ``beyond property''. Proyek
yang disebut meliputi Indonesia Movieland 36 ha, Medical City 74 ha,
Cyber City 240 ha, pembangkit gas 130 MW, dan pelabuhan darat 75-150 ha.
Sebagian unsur masih dinyatakan sebagai proyek yang sedang dikembangkan,
direncanakan, diproyeksikan, atau diharapkan, bukan sebagai hasil yang
seluruhnya telah selesai (hlm. 45-46, bagian ``Jababeka City: From
industrial park to suburban center'').

\textbf{Laporan sumber:} Pada saat artikel ditulis, Jababeka dilaporkan
memiliki lebih dari 1.200 penyewa dari sedikitnya 25 negara dan lebih
dari 300.000 pekerja. Fasilitas perkotaan yang disebut telah tersedia
mencakup 30.000 rumah, delapan hotel dan apartemen, 16 universitas dan
lembaga pendidikan tinggi, rumah sakit internasional, 24 mal dan pusat
belanja, lapangan golf internasional, dan kebun raya. Jababeka
dilaporkan menampung lebih dari 958.000 orang, termasuk sekitar 2.450
ekspatriat (hlm. 45, bagian ``Jababeka City: From industrial park to
suburban center'').

\textbf{Laporan sumber:} Ketimpangan intra-Kabupaten Bekasi terlihat
dalam pendapatan per kapita 2007. Tiga kecamatan industri secara agregat
mencapai Rp102.964.624 per kapita, dibandingkan Rp22.059.221 pada 12
kecamatan pinggiran, Rp9.088.036 pada enam kecamatan perdesaan, dan
Rp59.208.934 di Jakarta. Penulis juga melaporkan kemunculan komunitas
berpagar dan fragmentasi fisik antarkawasan industri (hlm. 46-47, Tabel
4 dan bagian ``Socio-economic and spatial implications and governance
adaptation'').

\textbf{Laporan sumber:} Masing-masing kawasan industri membangun jalan,
telekomunikasi, pengolahan air limbah, dan pengolahan air bersih tanpa
koordinasi yang jelas sehingga jaringan antarkawasan tidak tersambung.
Sebagai respons, tujuh kawasan industri menandatangani MOU pada 2006
untuk membentuk ZONI dan mengoordinasikan perencanaan serta
infrastruktur bersama pemerintah; pemerintah Provinsi Jawa Barat
kemudian menyiapkan tim koordinasi dan rencana rinci zona ekonomi khusus
Cikarang-Bekasi (hlm. 46-47, bagian ``Rethinking the governance
structure'').

\subsection{Findings}\label{findings-3}

\textbf{Interpretasi penulis:} Pola JMR dinilai menunjukkan tahap awal
\emph{post-suburbia}: industri teknologi tinggi dan perusahaan
multinasional mengalami dekonsentrasi ke kawasan industri pinggiran yang
direncanakan, lalu sebagian kawasan tersebut berkembang menjadi pusat
perkotaan dengan fungsi komersial, rekreasi, dan kebudayaan bercorak
Barat (hlm. 47, bagian ``The end of desakota regions?'').

\textbf{Interpretasi penulis:} Perkembangan industri bukan sekadar
akibat tambahan dari suburbanisasi perumahan, melainkan salah satu
penggerak utama restrukturisasi ekonomi dan spasial pinggiran JMR.
Kabupaten Bekasi mempertahankan posisi sebagai pusat industri utama
sekaligus mulai memimpin pergeseran ekonomi pinggiran menuju jasa (hlm.
44, setelah Tabel 2; hlm. 47, bagian ``The end of desakota regions?'').

\textbf{Interpretasi penulis:} Transformasi tersebut berbeda dari
gambaran \emph{desakota} yang bertahap, insidental, dan tidak terencana.
Akan tetapi, penulis menggunakan ungkapan hati-hati bahwa pola JMR
``mungkin'' menandai tahap awal \emph{post-suburbia}; artikel tidak
menyatakan bahwa seluruh ciri \emph{desakota} telah hilang (hlm. 47,
bagian ``The end of desakota regions?'').

\textbf{Interpretasi penulis:} Saturasi daya dukung inti Jakarta,
kelangkaan lahan, kekurangan infrastruktur, dan biaya produksi yang
lebih tinggi mendorong FDI baru, terutama manufaktur, ke pinggiran.
Karena itu, penulis menyebut transformasi ini tercipta oleh peluang
akibat keterbatasan inti, bukan terutama oleh kreativitas visi
pembangunan (hlm. 47, bagian ``The end of desakota regions?'').

\textbf{Interpretasi penulis:} Negara tidak merancang
\emph{post-suburbia} JMR secara langsung, tetapi transformasi juga tidak
sepenuhnya spontan. Regulasi pemerintah mendorong perkembangan properti
dan industri, keputusan lokasi serta pelaksanaan didominasi perusahaan
swasta, dan intervensi provinsi yang lebih langsung muncul secara
reaktif ketika daya saing global kawasan industri terancam (hlm. 47,
bagian ``The end of desakota regions?'').

\textbf{Interpretasi penulis:} Berbeda dari orientasi libertarian yang
dilekatkan pada \emph{post-suburbia} Amerika Serikat, kelompok
menengah-atas dan industri teknologi tinggi di JMR membutuhkan kepastian
aturan. Karena kapasitas pemerintah kota dinilai tidak memadai untuk
menerapkan aturan yang disepakati, pengembang swasta mengambil fungsi
merancang dan melaksanakan aturan perencanaan di kawasan pinggiran (hlm.
47, bagian ``The end of desakota regions?'').

\textbf{Interpretasi penulis:} Transformasi meningkatkan kinerja ekonomi
JMR, tetapi disertai ketimpangan regional, ketidakselarasan
infrastruktur, fragmentasi fisik, segregasi sosial, dan degradasi
lingkungan. Dengan demikian, pertumbuhan ekonomi dan pembentukan pusat
baru tidak disamakan dengan pembangunan metropolitan yang kohesif atau
berkelanjutan (hlm. 46, bagian ``Socio-economic and spatial implications
and governance adaptation''; hlm. 47, bagian ``The end of desakota
regions?'').

\textbf{Inferensi review:} Temuan terkuat artikel bukan bukti bahwa JMR
telah sepenuhnya menjadi wilayah \emph{post-suburban}, melainkan
demonstrasi bahwa kawasan industri dapat menjadi mesin pembentuk
sentralitas pinggiran dan tata kelola privat-publik. Status ``tahap
awal'' dan ketergantungan pada indikator deskriptif membatasi kekuatan
klaim transformasi menyeluruh.

\subsection{Conclusions}\label{conclusions-3}

\textbf{Interpretasi penulis:} Penulis menyimpulkan bahwa pola spasial
JMR memperlihatkan tahap awal \emph{post-suburbia} yang dibentuk oleh
dekonsentrasi industri teknologi tinggi dan perusahaan multinasional
serta perubahan beberapa kawasan industri menjadi pusat perkotaan baru.
Perkembangan ini melampaui bisnis lahan dan properti melalui penambahan
fungsi komersial, rekreasi, kesehatan, teknologi, dan kebudayaan (hlm.
47, bagian ``The end of desakota regions?'').

\textbf{Interpretasi penulis:} Pada tahap awal, transformasi dipandang
sebagai implikasi spasial restrukturisasi ekonomi regional dari
pertanian ke manufaktur. Penulis memproyeksikan bahwa dominasi
manufaktur secara bertahap akan diikuti sektor jasa, dengan proyek
Movieland, Medical City, dan Cyber City berpotensi mengubah pinggiran
industri menjadi pusat kebudayaan global regional (hlm. 47, bagian ``The
end of desakota regions?''). Proyeksi ini bukan hasil teramati pada
seluruh proyek.

\textbf{Interpretasi penulis:} Integrasi berkelanjutan dengan
kapitalisme global dan atribut tata kelola yang terdesentralisasi
membuat perkembangan \emph{post-suburban} dinilai tidak terhindarkan di
JMR. Namun, makna lokalnya ditentukan oleh keterbatasan inti Jakarta,
kepemimpinan pasar dengan dukungan regulasi negara, respons pemerintah
yang cenderung reaktif, dan pelaksanaan aturan perencanaan oleh
pengembang swasta (hlm. 47, bagian ``The end of desakota regions?'').

\textbf{Interpretasi penulis:} Manfaat ekonomi disertai masalah
keberlanjutan yang serius. Kesimpulan kebijakannya adalah perlunya
instrumen yang lebih inovatif, koordinasi perencanaan regional yang
lebih baik, dan peningkatan kapasitas pemerintah lokal untuk menangani
ketimpangan, fragmentasi, segregasi, ketidakselarasan infrastruktur, dan
degradasi lingkungan (hlm. 47, bagian ``The end of desakota regions?'').

\textbf{Inferensi review:} Kesimpulan tentang keberadaan unsur
\emph{post-suburbia} konsisten dengan bukti deskriptif yang disajikan,
tetapi istilah ``tidak terhindarkan'' dan hubungan sebab-akibat antara
transformasi industri dengan seluruh dampak regional lebih kuat daripada
yang dapat dipastikan oleh desain tanpa pengujian kausal.

\subsection{Contributions}\label{contributions-3}

\textbf{Interpretasi penulis:} Artikel menunjukkan bahwa perkembangan
pinggiran JMR tidak hanya dapat dibaca melalui konversi lahan,
perumahan, kota baru, atau \emph{desakota}. Kawasan industri terencana
dan perluasan fungsinya menyediakan pintu masuk lain untuk memahami
pembentukan pusat perkotaan baru (hlm. 40-41, bagian ``Introduction'';
hlm. 47, bagian ``The end of desakota regions?'').

\textbf{Inferensi review:} Kontribusi konseptualnya adalah membawa
perdebatan \emph{post-suburbia} dari pengalaman negara maju dan China ke
metropolitan yang masih mengalami industrialisasi, sembari
mempertahankan perbedaan konteks daripada menganggap model Barat berlaku
secara identik.

\textbf{Inferensi review:} Kontribusi empirisnya adalah merangkai bukti
multiskala dari JMR, Kabupaten Bekasi, Cikarang, tujuh kawasan industri,
dan Jababeka. Rangkaian tersebut menghubungkan dekonsentrasi penduduk,
perubahan penggunaan lahan, restrukturisasi sektor ekonomi, distribusi
FDI, fungsi perkotaan baru, ketimpangan lokal, dan adaptasi tata kelola
(hlm. 42-47, Gambar 1-2 dan Tabel 1-4).

\textbf{Inferensi review:} Kontribusi penjelasannya terletak pada
identifikasi konfigurasi lokal: keterbatasan kapasitas inti menciptakan
peluang lokasi di pinggiran; negara membentuk kondisi melalui regulasi
tetapi tidak memimpin dengan visi jangka panjang; pengembang
menginternalisasi penyediaan infrastruktur dan aturan; lalu pemerintah
dan perusahaan membentuk koordinasi baru ketika fragmentasi dan
persaingan meningkat (hlm. 44-47, bagian ``Economic and political
background and policy contexts'', ``Rethinking the governance
structure'', dan ``The end of desakota regions?'').

\textbf{Inferensi review:} Artikel juga menyumbangkan koreksi normatif
bahwa keberhasilan menarik investasi dan membentuk sentralitas baru
dapat berjalan bersama ketimpangan, segregasi, serta jaringan
infrastruktur yang terpecah. Kontribusi ini bersifat interpretatif,
bukan estimasi dampak kausal.

\subsection{Research Gap}\label{research-gap-3}

\textbf{Laporan sumber:} Kesenjangan yang dinyatakan secara eksplisit
adalah kurangnya penelitian tentang apakah perkembangan industri telah
menciptakan pusat perkotaan baru di pinggiran metropolitan Indonesia
yang sama menariknya dengan pusat kota lama. Kajian sebelumnya lebih
banyak mengakui dampak pembangunan perumahan di pinggiran (hlm. 40,
bagian ``Introduction'').

\textbf{Laporan sumber:} Artikel juga menanggapi kesenjangan konseptual
antara pemahaman tradisional pinggiran Asia Timur sebagai perpanjangan
inti atau zona transisi perkotaan-perdesaan dan perubahan yang dipicu
kota global serta desentralisasi (hlm. 40, abstrak dan bagian
``Introduction'').

\textbf{Laporan sumber:} Pertanyaan yang belum mapan dalam literatur
yang dipakai artikel adalah apakah unsur \emph{post-suburbia} dapat
ditemukan di JMR, bagaimana industri berperan, dan ciri apa yang
membedakan bentuk JMR dari pengalaman internasional (hlm. 40, bagian
``Introduction'').

\textbf{Inferensi review:} Artikel mengisi kesenjangan eksploratif
tersebut, tetapi menyisakan kebutuhan akan operasionalisasi konsep
\emph{post-suburbia}, perbandingan sistematis antarkawasan, penelusuran
perubahan setelah proyek yang masih direncanakan selesai atau gagal, dan
pengujian hubungan kausal antara industrialisasi, pembentukan
sentralitas, serta dampak sosial-lingkungan. Kebutuhan ini diturunkan
dari batas bukti dalam file, bukan agenda eksplisit penulis.

\subsection{Limitations}\label{limitations-3}

\textbf{Laporan sumber:} Tidak disebutkan dalam file sebagai bagian
keterbatasan penelitian yang khusus. Pada bagian ucapan terima kasih,
penulis hanya menyatakan tanggung jawab atas kesalahan dan kekurangan
tanpa merincinya (hlm. 47, bagian ``Acknowledgments'').

\textbf{Inferensi review:} Tidak adanya uraian metode formal membuat
prosedur pengumpulan, pembersihan, pemilihan, dan analisis data tidak
dapat direkonstruksi sepenuhnya. Tidak ada penjelasan tentang wawancara,
observasi, validasi silang, reliabilitas data, atau cara konsep
\emph{post-suburbia} dioperasionalkan.

\textbf{Inferensi review:} Seri data berasal dari periode dan sumber
yang berbeda, dari sensus 1961-2000 hingga dokumen 2010. File tidak
menjelaskan harmonisasi batas wilayah, konsistensi definisi sektor,
perubahan nilai nominal, atau ketidakpastian perhitungan. Karena itu,
angka-angka lebih aman dibaca sebagai pola deskriptif daripada satu
rangkaian longitudinal yang sepenuhnya sebanding.

\textbf{Inferensi review:} Sebagian bukti Jababeka berasal dari
publikasi perusahaan dan sebagian status proyek bersifat rencana,
proyeksi, atau harapan. Artikel tidak menjelaskan verifikasi independen
atas klaim perusahaan maupun membedakan secara sistematis fasilitas yang
telah beroperasi dari yang masih direncanakan (hlm. 45-46, bagian
``Jababeka City: From industrial park to suburban center''). Ini
menimbulkan potensi bias sumber, bukan bukti bahwa data tersebut salah.

\textbf{Inferensi review:} Pendalaman empiris berpusat pada Kabupaten
Bekasi dan terutama Jababeka yang sengaja dipilih karena ukuran dan
dinamismenya. Kasus tersebut informatif tetapi berpotensi tidak mewakili
kawasan industri lain, pinggiran JMR bagian lain, atau metropolitan
Indonesia secara umum.

\textbf{Laporan sumber:} Tabel 3 menandai data luas industri dan jumlah
pekerja Deltamas sebagai tidak tersedia (hlm. 45, Tabel 3).
\textbf{Inferensi review:} Kekosongan ini mengurangi kelengkapan
perbandingan tujuh kawasan industri.

\textbf{Inferensi review:} Artikel menyimpulkan adanya degradasi
lingkungan dan mengaitkan transformasi dengan sejumlah dampak
keberlanjutan, tetapi file tidak menyajikan indikator lingkungan
langsung. Demikian pula, hubungan antara industri, pertumbuhan penduduk,
ketimpangan, komunitas berpagar, dan fragmentasi dibahas secara
deskriptif tanpa strategi identifikasi kausal (hlm. 46-47, bagian
``Socio-economic and spatial implications and governance adaptation''
dan ``The end of desakota regions?'').

\textbf{Inferensi review:} Karena data empiris utama berakhir sekitar
2009-2010 dan artikel terbit pada 2012, file tidak dapat memastikan
realisasi jangka panjang proyek, KEK, integrasi infrastruktur, ataupun
proyeksi pergeseran menuju jasa.

\subsection{Challenges}\label{challenges-3}

\textbf{Laporan sumber:} Tantangan substantif pertama adalah menurunnya
daya dukung inti Jakarta akibat kelangkaan lahan, kekurangan
infrastruktur, dan biaya produksi tinggi, yang mengurangi kemampuannya
menampung pertumbuhan dan bersaing menarik FDI (hlm. 47, bagian ``The
end of desakota regions?'').

\textbf{Laporan sumber:} Kawasan industri Cikarang menghadapi persaingan
lebih ketat dari negara-negara Asia lain yang juga membangun kawasan
untuk menarik FDI. Industri manufaktur dipandang mudah berpindah,
sehingga pemerintah dan pengelola berupaya mempertahankan investor serta
memperbaiki iklim usaha (hlm. 46-47, bagian ``Rethinking the governance
structure'').

\textbf{Laporan sumber:} Penyediaan infrastruktur secara sendiri-sendiri
oleh tiap kawasan menciptakan jaringan yang tidak terhubung, sementara
komunikasi antarpengelola sebelumnya terbatas. ZONI dan rencana KEK
muncul sebagai upaya koordinasi, tetapi pada waktu yang sama
memperlihatkan bahwa skala masalah telah melampaui batas masing-masing
kawasan (hlm. 46-47, bagian ``Socio-economic and spatial implications
and governance adaptation'' dan ``Rethinking the governance
structure'').

\textbf{Laporan sumber:} Ketimpangan pendapatan, segregasi sosial,
komunitas berpagar, fragmentasi fisik, dan degradasi lingkungan menjadi
tantangan terhadap kohesi serta keberlanjutan regional. Penulis juga
menyatakan bahwa pemerintah kota di Indonesia belum dibekali kapasitas
kelembagaan yang dibutuhkan untuk mewujudkan aturan perencanaan yang
telah disepakati (hlm. 46-47, Tabel 4 dan bagian ``The end of desakota
regions?'').

\textbf{Interpretasi penulis:} Tantangan tata kelola terletak pada
kebutuhan menggabungkan inisiatif swasta yang dominan dengan koordinasi
dan kapasitas publik yang lebih kuat, tanpa mengandalkan respons
pemerintah yang hanya muncul ketika daya saing terancam (hlm. 47, bagian
``The end of desakota regions?'').

\textbf{Inferensi review:} Tantangan analitis bagi artikel adalah
memisahkan perubahan yang telah teramati dari proyek yang baru
diusulkan, serta menghubungkan indikator pada skala dan tahun berbeda
tanpa metode harmonisasi yang dijelaskan.

\subsection{Practical Implications}\label{practical-implications-3}

\textbf{Interpretasi penulis:} Masalah ketimpangan regional,
ketidakselarasan infrastruktur, fragmentasi fisik, segregasi sosial, dan
degradasi lingkungan menuntut instrumen kebijakan yang lebih inovatif,
koordinasi perencanaan pada tingkat regional yang lebih baik, dan
kapasitas pemerintah lokal yang lebih besar (hlm. 47, bagian ``The end
of desakota regions?'').

\textbf{Laporan sumber:} Pengalaman ZONI menunjukkan satu jalur praktis
yang sedang ditempuh, yakni kerja sama tujuh kawasan industri dengan
pemerintah untuk menyusun rencana terintegrasi dan membangun
infrastruktur secara lebih efisien. Program yang disebut meliputi
jembatan, jalan layang, penambahan lajur Tol Jakarta-Cikampek, dan
perluasan jalan sekitar Cikarang (hlm. 46, bagian ``Rethinking the
governance structure'').

\textbf{Laporan sumber:} Rencana KEK dimaksudkan untuk menerapkan
pengaturan khusus mengenai bea, pajak, perizinan, imigrasi, dan tenaga
kerja, dengan dukungan infrastruktur serta lembaga pengelola berstandar
internasional. Manfaat yang diharapkan adalah mempertahankan investor,
menarik FDI baru, memperluas akses pasar global, merangsang industri
lokal, dan menciptakan pekerjaan (hlm. 46-47, bagian ``Rethinking the
governance structure''). Manfaat tersebut merupakan sasaran kebijakan,
bukan dampak yang telah dievaluasi dalam artikel.

\textbf{Inferensi review:} Karena pengembang telah menjalankan fungsi
perencanaan dan penyediaan layanan yang besar, koordinasi regional perlu
menilai bukan hanya efisiensi investasi, tetapi juga keterhubungan
jaringan, distribusi manfaat, dampak sosial-lingkungan, dan
akuntabilitas aturan privat. Unsur akuntabilitas ini merupakan implikasi
review dari diagnosis artikel, bukan rekomendasi eksplisit penulis.

\subsection{Applications}\label{applications-3}

\textbf{Inferensi review:} Kerangka artikel dapat digunakan sebagai alat
diagnosis untuk menilai apakah suatu kawasan industri tetap berupa
lokasi produksi atau telah berkembang menjadi pusat \emph{post-suburban}
melalui penambahan pekerjaan, hunian, jasa, infrastruktur, rekreasi, dan
kebudayaan. Penerapan ini harus mempertahankan perbedaan antara
indikator keberadaan fungsi dan bukti hubungan sebab-akibat.

\textbf{Inferensi review:} Rangkaian indikator pada Gambar 1-2 dan Tabel
1-4 dapat menjadi titik awal pemantauan dekonsentrasi penduduk,
penggunaan lahan, distribusi investasi, restrukturisasi ekonomi,
konsentrasi pekerja, dan ketimpangan intra-wilayah. Artikel sendiri
tidak menyediakan protokol pemantauan atau ambang klasifikasi.

\textbf{Inferensi review:} Analisis tata kelola dapat diterapkan untuk
mengevaluasi pembagian peran antara pemerintah, pengembang kawasan,
penyedia infrastruktur, dan organisasi antarkawasan seperti ZONI,
khususnya ketika jaringan privat terfragmentasi dan fungsi perkotaan
melampaui batas satu proyek.

\textbf{Inferensi review:} Kasus ini juga dapat dipakai sebagai
pembanding konseptual terhadap \emph{desakota}, \emph{edge city},
\emph{city-edge}, dan bentuk \emph{post-suburbia} lain. Penggunaan
komparatif harus memperlakukan JMR sebagai varian kontekstual, bukan
sebagai bukti bahwa semua pinggiran Asia mengikuti lintasan yang sama.

\subsection{Future Research
(Optional)}\label{future-research-optional-3}

\textbf{Laporan sumber:} Tidak disebutkan dalam file sebagai agenda
penelitian lanjutan yang eksplisit setelah kesimpulan. Pernyataan bahwa
sektor jasa akan mengikuti manufaktur dan proyek ``beyond property''
akan mengubah Cikarang merupakan proyeksi substantif, bukan rancangan
penelitian mendatang (hlm. 47, bagian ``The end of desakota regions?'').

\textbf{Inferensi review:} Penelitian lanjutan dapat memeriksa secara
longitudinal apakah Movieland, Medical City, Cyber City, pelabuhan
darat, ZONI, dan KEK benar-benar direalisasikan serta apakah hasilnya
sesuai dengan proyeksi ekonomi, spasial, sosial, dan budaya dalam
artikel.

\textbf{Inferensi review:} Perbandingan terstruktur antara Jababeka,
enam kawasan lain di Cikarang, bagian Bodetabek lain, dan metropolitan
Indonesia lain diperlukan untuk menilai sejauh mana temuan kasus dapat
digeneralisasi dan untuk membedakan pengaruh karakter lokasi dari model
pengembang.

\textbf{Inferensi review:} Pengukuran formal atas polisentralitas, pola
perjalanan, keterhubungan infrastruktur, segregasi, dampak lingkungan,
kualitas pekerjaan, dan distribusi manfaat dapat menguji klaim
\emph{post-suburbia} serta konsekuensinya. Arah ini tidak diajukan
secara eksplisit oleh penulis; ia diturunkan dari indikator dan
keterbatasan yang tampak dalam file.

\subsection{Provenance Notes}\label{provenance-notes-3}

\begin{itemize}
\tightlist
\item
  \textbf{Cakupan akses:} Review ini hanya menggunakan seluruh isi file
  \texttt{source/journal/J04-hudalah-firman-2012-beyond-property-industrial-estates-and-post-suburban-transformation-in-jakarta-metropolitan-region-10.1016-j.cities.2011.07.003.txt},
  termasuk metadata ekstraksi, isi artikel, tabel, keterangan gambar,
  ucapan terima kasih, dan daftar pustaka. Tidak ada sumber eksternal
  yang digunakan.
\item
  \textbf{Asal teks:} Metadata file menyatakan bahwa teks diekstraksi
  pada 31 Agustus 2026 dengan \texttt{pdftotext\ -layout} dari PDF
  sembilan halaman. PDF sumber bernama
  \texttt{journal/J04-hudalah-firman-2012-beyond-property-industrial-estates-and-post-suburban-transformation-in-jakarta-metropolitan-region-10.1016-j.cities.2011.07.003.pdf}
  (blok \texttt{{[}PDF\ Metadata{]}}). PDF tidak diperiksa langsung
  dalam pekerjaan ini.
\item
  \textbf{Locator:} Nomor halaman mengacu pada paginasi jurnal yang
  tercetak, yaitu halaman 40-48, bukan nomor baris TXT. Klaim penting
  juga diberi locator bagian, gambar, atau tabel ketika tersedia. Daftar
  pustaka berlanjut sampai halaman 48; pembahasan substantif dan ucapan
  terima kasih berakhir pada halaman 47.
\item
  \textbf{Sifat bukti:} Data BPS, BKPM, pemerintah daerah, dokumen
  kebijakan, ZONI, dan Jababeka diketahui hanya melalui pelaporan
  artikel. Dataset dan dokumen asal tidak dibuka atau divalidasi secara
  independen. Karena itu, review tidak mengubah klaim sekunder artikel
  menjadi fakta yang telah diverifikasi langsung.
\item
  \textbf{Kualitas ekstraksi:} Tata letak dua kolom menghasilkan
  pemenggalan kata, ligatur, dan beberapa sambungan kalimat yang
  terpisah dalam TXT. Isi direkonstruksi secara konservatif berdasarkan
  urutan bagian dan paginasi, tanpa menambahkan teks yang tidak ada
  dalam file.
\item
  \textbf{Pembedaan epistemik:} \texttt{Laporan\ sumber} menandai data,
  uraian, atau pernyataan yang disajikan artikel;
  \texttt{Interpretasi\ penulis} menandai penjelasan dan kesimpulan
  analitis penulis; \texttt{Inferensi\ review} menandai klasifikasi,
  evaluasi, implikasi, atau usulan yang disusun dalam review ini.
\item
  \textbf{Batas utama provenance:} Status penelaahan sejawat, prosedur
  metode rinci, data mentah, realisasi proyek setelah periode artikel,
  dan kebenaran klaim pada dokumen sumber sekunder tidak dapat
  dipastikan dari TXT. Tidak ada klaim dalam review ini yang dimaksudkan
  melampaui batas tersebut.
\end{itemize}

\section{Review J05}\label{review-j05}

\textbf{Identitas sumber}

\begin{itemize}
\tightlist
\item
  \textbf{Penulis:} Tommy Firman dan Fikri Zul Fahmi.
\item
  \textbf{Tahun:} 2017.
\item
  \textbf{Judul:} ``The Privatization of Metropolitan Jakarta's
  (Jabodetabek) Urban Fringes: The Early Stages of
  `Post-Suburbanization' in Indonesia.''
\item
  \textbf{Jurnal:} \emph{Journal of the American Planning Association},
  Vol. 83, No.~1, Winter 2017, hlm. 68--79.
\item
  \textbf{DOI:} 10.1080/01944363.2016.1249010.
\item
  \textbf{Tanggal publikasi daring:} 5 Januari 2017.
\item
  \textbf{Objek geografis:} Wilayah metropolitan Jakarta (Jabodetabek),
  Indonesia.
\end{itemize}

\subsection{Summarized Abstract}\label{summarized-abstract-4}

\textbf{Laporan sumber.} Artikel ini menelaah apakah perkembangan
mutakhir di pinggiran Jabodetabek menunjukkan tahap awal
\emph{post-suburbanization}. Dengan pendekatan yang secara eksplisit
disebut deskriptif, penulis memeriksa pola perkembangan fisik serta
perubahan peran sektor publik dan swasta. Pertumbuhan cepat kawasan
pinggiran, perubahan kota-kota komuter menjadi kota yang lebih mandiri,
serta berkembangnya kawasan industri dan kota baru multifungsi dipandang
sebagai ciri yang sejalan dengan \emph{post-suburbia}. Penulis
menghubungkan pola tersebut dengan privatisasi penguasaan, pengembangan,
dan pengelolaan lahan: kebijakan ekonomi pemerintah pusat yang
berorientasi pada pertumbuhan dan otonomi pemerintah daerah memperbesar
kekuasaan pengembang swasta di kawasan pinggiran. Sektor swasta membantu
menyediakan perumahan, pekerjaan, perdagangan, pendidikan,
infrastruktur, dan layanan perkotaan, tetapi orientasi laba dan lemahnya
koordinasi regional berisiko menghasilkan pembangunan yang
terfragmentasi dan bertentangan dengan tujuan kebijakan publik (abstrak,
hlm. 68; bagian ``Privatization of the Urban Fringes and
Post-Suburbanization in Jabodetabek,'' hlm. 76--77).

\textbf{Interpretasi penulis.} Jabodetabek belum menjadi replika penuh
\emph{post-suburbia} Barat. Kawasan ini memperlihatkan tahap awal dengan
campuran bentuk suburban lama dan baru, pertumbuhan pusat-pusat
pinggiran, serta relasi komuter yang masih kuat dengan kota inti (hlm.
76--77).

\subsection{Problem Statement}\label{problem-statement-4}

\textbf{Laporan sumber.} Konsep \emph{post-suburbia} terutama dibentuk
dari pengalaman negara maju: pusat kota tradisional kehilangan dominasi,
sedangkan struktur polisentris, pekerjaan, dan fungsi perkotaan tumbuh
di luar inti. Perkembangan metropolitan Asia juga menunjukkan
desentralisasi penduduk dan kegiatan ekonomi, tetapi berlangsung dalam
konteks globalisasi modal, peran negara, dan perubahan hubungan
kekuasaan publik--swasta yang berbeda. Masalah penelitian artikel ini
ialah seberapa jauh perkembangan Jabodetabek benar-benar mencerminkan
fenomena tersebut, khususnya ketika pengembang swasta memperoleh
kekuasaan yang semakin besar atas keputusan serta layanan pembangunan
lahan (pendahuluan, hlm. 68--69; bagian ``Post-Suburbanization as a
Worldwide Phenomenon of Urban Development,'' hlm. 69--71).

\textbf{Interpretasi penulis.} Persoalannya bukan hanya perubahan
morfologi dari monosentris menuju multisentris. Penulis juga
mempermasalahkan privatisasi tata kelola: pergeseran kewenangan,
kompetensi, dan tanggung jawab dari sektor publik kepada sektor swasta
dapat mempercepat pembangunan sekaligus memperlemah kepatuhan pada
rencana dan koordinasi regional (hlm. 69 dan 76--77).

\textbf{Pertanyaan praktis yang dinyatakan penulis.} Bagaimana perencana
dapat memastikan bahwa pola \emph{post-suburban development} di negara
berkembang tidak menghasilkan masalah seperti yang dialami Amerika
Utara? (pendahuluan, hlm. 69).

\subsection{Objectives}\label{objectives-4}

\textbf{Laporan sumber.} Tujuan utama penelitian adalah mengidentifikasi
sejauh mana fenomena \emph{post-suburbia} yang pertama kali diamati di
negara maju tercermin dalam kondisi metropolitan Jakarta. Untuk itu,
penulis:

\begin{itemize}
\tightlist
\item
  memperbarui kajian mereka sebelumnya mengenai Jabodetabek dengan
  meninjau perubahan dan kesinambungan perkembangan kawasan pinggiran
  selama sekitar empat dekade;
\item
  memeriksa pertumbuhan penduduk, perubahan penggunaan lahan,
  perkembangan kawasan industri, dan perkembangan kota baru;
\item
  menelaah perubahan peran serta hubungan sektor publik dan swasta,
  terutama peran sektor swasta yang difasilitasi oleh meningkatnya
  kewenangan pemerintah daerah; dan
\item
  merefleksikan implikasi pola tersebut bagi ilmu dan praktik
  perencanaan (pendahuluan, hlm. 69; bagian penutup analitis, hlm.
  76--77).
\end{itemize}

\subsection{Literature Survey}\label{literature-survey-4}

\textbf{Laporan sumber.} Telaah pustaka membangun empat kelompok
landasan.

\begin{itemize}
\tightlist
\item
  \textbf{Konsep \emph{post-suburbia}.} Literatur menggambarkannya
  sebagai peralihan dari pola suburban konsentris dan radial menuju
  struktur tambal-sulam, polisentris, dan lebih mandiri. Kawasan luar
  tidak lagi semata-mata berfungsi sebagai tempat tinggal, tetapi juga
  menampung pekerjaan, pusat belanja, manufaktur, kantor, industri
  berteknologi tinggi, dan pendidikan. Istilah terkait yang dibahas
  ialah ``\emph{edgeless city}'' dan ``\emph{technoburbia}'' (hlm.
  68--70).
\item
  \textbf{Pembeda dari suburbanisasi tradisional.} Berdasarkan literatur
  yang dikutip, penulis menyebut tiga ciri: penurunan penduduk serta
  pendapatan rumah tangga suburbia relatif terhadap wilayah;
  desentralisasi pekerjaan jasa dari pusat kota; dan penggunaan lahan
  yang lebih bercampur dengan perkembangan polisentris yang tegas
  (bagian ``Post-Suburbanization as a Worldwide Phenomenon of Urban
  Development,'' hlm. 70).
\item
  \textbf{Pengalaman Asia, terutama Tiongkok.} Beijing dan Shanghai
  digunakan sebagai pembanding karena perkembangan pinggirannya makin
  berorientasi pasar melalui suburbanisasi perumahan, industri, dan
  perdagangan. Namun, hubungan dengan inti kota belum sepenuhnya longgar
  dan sektor publik masih kuat. Literatur Tiongkok membantu penulis
  melihat \emph{post-suburbia} sebagai perubahan tingkat kemandirian
  kawasan, politik dan tata kelola pinggiran, serta susunan kekuasaan
  yang menyatukan wilayah metropolitan, bukan sebagai satu bentuk Barat
  yang seragam (hlm. 70--71).
\item
  \textbf{Mega-urbanisasi dan privatisasi di Asia/Jakarta.} Literatur
  tentang Jabodetabek, Asia Tenggara, desentralisasi, dan privatisasi
  perencanaan digunakan untuk menjelaskan campuran kegiatan ekonomi dan
  permukiman di pinggiran, perluasan kawasan terbangun, investasi real
  estat swasta, klientelisme, serta pergeseran kewenangan pengelolaan
  lahan dari pemerintah kepada pelaku swasta (hlm. 68--69 dan 71--76).
\end{itemize}

\textbf{Interpretasi penulis.} Kasus Tiongkok dipakai sebagai cermin
analitis, bukan sebagai padanan sempurna. Kesamaan utamanya ialah tahap
awal \emph{post-suburbanization} dan bertambahnya pengaruh sektor
swasta; perbedaannya terletak pada konfigurasi peran pemerintah dan
tingkat keterlepasan pinggiran dari inti (hlm. 70--71).

\subsection{Method Used}\label{method-used-4}

\textbf{Laporan sumber.} Penulis menyatakan bahwa pendekatannya bersifat
deskriptif. Artikel disusun dalam tiga bagian analitis: konteks teoretis
\emph{post-suburbanization} sebagai fenomena global; pemeriksaan
perkembangan Jabodetabek melalui pertumbuhan penduduk, perubahan
penggunaan lahan, kota baru, dan kawasan industri; serta refleksi atas
hubungan peran swasta dengan kebijakan pemerintah pusat dan daerah
(abstrak, hlm. 68; pendahuluan, hlm. 69).

\textbf{Inferensi review.} Berdasarkan bahan yang ditampilkan, metode
operasionalnya paling tepat dibaca sebagai kajian kasus deskriptif
berbasis sintesis dokumen dan data sekunder. Artikel menggabungkan
statistik kependudukan dan investasi, hasil studi penginderaan jauh/SIG,
kajian terdahulu, peraturan, laporan pasar properti, laporan perusahaan,
dan berita. Sebutan ``kajian kasus'' dan ``sintesis dokumen'' ini
merupakan klasifikasi review, bukan label metode yang dinyatakan
penulis.

Prosedur pencarian atau pemilihan dokumen, teknik harmonisasi data,
protokol analisis, wawancara, survei primer, observasi lapangan,
pengujian hipotesis, dan model statistik: Tidak disebutkan dalam file.

\subsection{Dataset}\label{dataset-4}

\textbf{Laporan sumber.} Artikel tidak menamai atau menyediakan satu set
data penelitian yang terintegrasi. Bukti kuantitatif dan dokumenter yang
digunakan tersebar dalam sumber-sumber sekunder berikut.

\begin{itemize}
\tightlist
\item
  Statistik penduduk, kepadatan, pertumbuhan, distribusi, migrasi,
  produk domestik bruto, serta investasi dari Badan Pusat Statistik dan
  berbagai studi yang dikutip. Rentang tahun yang tampak dalam
  pembahasan utama mencakup antara lain 1980--1990, 1990--2000,
  1995--2000, 2000--2010, dan kondisi sekitar 2014 (bagian ``Urban
  Development in Jabodetabek'' dan ``Population Growth and Land Use
  Conversion in Jabodetabek,'' hlm. 71--74).
\item
  Hasil studi perubahan penggunaan lahan berbasis penginderaan jauh dan
  sistem informasi geografis untuk 1992--2001, serta estimasi perluasan
  kawasan terbangun pada 2000--2010 (hlm. 74).
\item
  Data perkembangan kawasan industri, harga lahan, luas kawasan,
  investasi, dan perusahaan dari studi terdahulu, Colliers
  International, laporan PT Jababeka, BPS, serta berita bisnis; tahun
  pengamatan yang disebut antara lain pertengahan 2000-an, 2005, 2011,
  2012, 2013, dan 2014 (bagian ``Industrial Estate Development,'' hlm.
  74--75).
\item
  Dokumen kebijakan dan peraturan mengenai kawasan industri, perumahan,
  serta rencana tata ruang, termasuk Keputusan Presiden 41/1996,
  Undang-Undang 3/2014, Peraturan Pemerintah 142/2015, Undang-Undang
  4/1992 yang kemudian diubah oleh Undang-Undang 1/2011, dan Peraturan
  Presiden 54/2008 (hlm. 74--77, catatan 3--6).
\end{itemize}

Ukuran, skema penggabungan, akses, dan pengolahan sebuah dataset
tunggal: Tidak disebutkan dalam file.

\subsection{Population / Sample}\label{population-sample-4}

\textbf{Laporan sumber.} Unit geografis utama ialah Jabodetabek, wilayah
seluas lebih dari 9.000 km² di bagian utara Jawa Barat yang terdiri atas
DKI Jakarta; lima kota, yaitu Bogor, Depok, Tangerang, Tangerang
Selatan, dan Bekasi; serta tiga kabupaten, yaitu Bogor, Tangerang, dan
Bekasi (bagian ``Urban Development in Jabodetabek,'' hlm. 71; Gambar 1,
hlm. 71). Populasi wilayah ini dilaporkan mendekati 30 juta jiwa pada
2014, sedangkan Jakarta sebagai inti memiliki 9,6 juta jiwa pada 2010
(hlm. 73).

Objek empiris tambahan mencakup kawasan pinggiran, kawasan industri dan
kota baru, dengan contoh yang dibahas lebih rinci seperti
Cikarang/Jababeka, Lippo Karawaci, Lippo Cikarang, dan Bumi Serpong
Damai (hlm. 72 dan 74--76).

\textbf{Inferensi review.} Kasus-kasus pengembang dan kota baru
berfungsi sebagai ilustrasi substantif, bukan sampel statistik yang
dinyatakan representatif. Kerangka populasi, metode sampling, jumlah
responden, kriteria inklusi kasus, dan unit observasi mikro: Tidak
disebutkan dalam file.

\subsection{Dependent Variables}\label{dependent-variables-4}

Tidak disebutkan dalam file. Artikel tidak membangun model kausal atau
statistik dan tidak mendefinisikan variabel dependen formal.

\textbf{Inferensi review.} Indikator hasil yang dideskripsikan, tetapi
tidak diperlakukan sebagai variabel dependen, meliputi laju serta
distribusi pertumbuhan penduduk; kepadatan; konversi lahan
pertanian/hutan menjadi lahan terbangun; pertumbuhan kawasan industri
dan kota baru; diferensiasi fungsi pinggiran; arus komuter; penyediaan
layanan perkotaan oleh pengembang; serta pergeseran kewenangan
pembangunan lahan dari sektor publik ke sektor swasta (hlm. 72--76).

\subsection{Results}\label{results-4}

\textbf{Laporan sumber: perubahan demografis dan struktur metropolitan.}

\begin{itemize}
\tightlist
\item
  Populasi Jabodetabek dilaporkan mendekati 30 juta jiwa pada 2014,
  dengan pertumbuhan tahunan 3,6\% selama 2000--2010. Enam dari 12 kota
  Indonesia yang berpenduduk sekurang-kurangnya satu juta jiwa berada di
  Jabodetabek. Jakarta memiliki 9,6 juta penduduk pada 2010 dan menerima
  sekitar dua juta komuter dari wilayah sekitarnya (bagian ``Population
  Growth and Land Use Conversion in Jabodetabek,'' hlm. 73).
\item
  Kepadatan Jabodetabek meningkat dari 37,6 orang per hektare pada 2000
  menjadi 44,6 pada 2010. Kepadatan Jakarta meningkat dari 128,0 menjadi
  145,9 orang per hektare pada periode yang sama (hlm. 73).
\item
  Pangsa Jakarta dalam populasi Jabodetabek turun dari 54,6\% pada 1990
  menjadi 35,5\% pada 2010. Pertumbuhan tahunan Jakarta melambat dari
  3,1\% pada 1980--1990 menjadi 0,4\% pada 1990--2000, lalu naik menjadi
  1,5\% pada 2000--2010. Sebaliknya, pinggiran tumbuh hampir 3\% per
  tahun; Kota Tangerang dan Kota Bekasi masing-masing tumbuh 3,2\% dan
  3,4\% per tahun pada 2000--2010 (hlm. 73).
\item
  Untuk 1995--2000, BPS dipaparkan mengestimasi sekitar 160.000 penduduk
  Jakarta berpindah ke kota dan kabupaten Bogor serta sekitar 190.000
  berpindah ke unit-unit kota/kabupaten Bekasi dan Tangerang yang
  disebut dalam teks (hlm. 73). Redaksi sumber tidak menguraikan lebih
  lanjut pembagian angka 190.000 itu menurut tiap tujuan.
\item
  Bogor, Depok, Serpong, dan Cikarang dipaparkan telah memperoleh basis
  fungsi pendidikan, penelitian, teknologi, industri, perdagangan, atau
  pertemuan yang lebih kuat. Pergerakan komuter tidak hanya menuju
  Jakarta; Cikarang juga menarik perjalanan kerja dari inti dan
  kota-kota kecil Jabodetabek (bagian ``Urban Development in
  Jabodetabek,'' hlm. 72).
\end{itemize}

\textbf{Laporan sumber: konversi lahan.}

\begin{itemize}
\tightlist
\item
  Sekitar 4.000 ha sawah dan 8.000 ha hutan primer di bagian selatan
  Jabodetabek dikonversi menjadi kawasan industri dan permukiman pada
  1994--2001 (hlm. 73).
\item
  Studi penginderaan jauh/SIG yang dikutip melaporkan kenaikan bagian
  lahan terbangun Jabodetabek dari 12\% menjadi 24\% pada 1992--2001,
  sedangkan bagian lahan pertanian turun dari 37\% menjadi 31\% (hlm.
  74).
\item
  Studi lain yang dikutip melaporkan perluasan kawasan
  perkotaan/terbangun di pinggiran Bogor, Tangerang, Bekasi, dan Depok
  dari 544 menjadi 850 km² pada 2000--2010, setara 4,6\% per tahun. Di
  Jakarta, luasnya meningkat dari 560 menjadi 594 km², atau 0,6\% per
  tahun (hlm. 74).
\item
  Konversi lahan pertanian dan ruang terbuka di Bogor--Puncak--Cianjur
  dipandang bermasalah karena kawasan itu ditetapkan sebagai wilayah
  konservasi dan resapan air. Teks menyatakan konversi tersebut
  ``diperkirakan'' menjadi salah satu penyebab utama banjir musiman di
  Jakarta, sehingga hubungan ini disajikan sebagai penilaian dalam
  sumber, bukan hasil uji kausal artikel (hlm. 74).
\end{itemize}

\textbf{Laporan sumber: kawasan industri.}

\begin{itemize}
\tightlist
\item
  Pada 2013 terdapat 35 kawasan industri di pinggiran Jabodetabek,
  masing-masing seluas sekitar 50--1.800 ha; sekitar seperempatnya
  berada di Kabupaten Bekasi. Sekitar 400 ha kawasan industri
  ditambahkan pada 2013--2014, terutama untuk industri otomotif (bagian
  ``Industrial Estate Development,'' hlm. 74).
\item
  Cikarang menjadi konsentrasi utama kawasan industri Bekasi. Jababeka
  digambarkan sebagai kota mandiri dan pusat manufaktur dengan lebih
  dari 1.500 perusahaan nasional dan multinasional dari lebih dari 35
  negara. Luas total kawasan industri di pinggiran Jakarta dilaporkan
  naik dari 11.000 ha pada 2005 menjadi 18.000 ha pada 2012 (hlm. 75).
\item
  Nilai ekspor potensial kawasan industri Cikarang pada pertengahan
  2000-an dilaporkan sebesar US\$30,56 miliar, hampir separuh ekspor
  nonmigas nasional sebesar US\$66,43 miliar pada saat yang sama. Angka
  investasi asing dan domestik lain juga disajikan dari BPS 2006, tetapi
  tidak diuji atau direkonsiliasi oleh artikel (hlm. 75).
\end{itemize}

\textbf{Laporan sumber: kota baru dan layanan privat.}

\begin{itemize}
\tightlist
\item
  Kota-kota baru yang semula berupa komunitas komuter telah memperoleh
  fungsi ekonomi yang lebih beragam. Lippo Karawaci dan Lippo Cikarang
  memadukan hunian dengan universitas, rumah sakit, perdagangan, dan
  kegiatan ekonomi; Jababeka memadukan perumahan, manufaktur, kesehatan,
  serta industri film (bagian ``New Town Residential Development,'' hlm.
  75--76).
\item
  Kota baru terutama merespons permintaan kelompok berpendapatan
  menengah dan atas akan lingkungan yang aman, modern, dan tenang.
  Bentuknya cenderung berupa komunitas berpagar, terpisah dari
  permukiman sekitar, dan tidak bercampur secara sosial maupun budaya
  (hlm. 75--76).
\item
  Pengembang menyediakan bukan hanya infrastruktur dasar, melainkan juga
  jalan, air bersih, pembuangan air limbah, lanskap, keamanan, dan bus
  ulang-alik. Mereka bahkan menunjuk manajer kota sendiri; Lippo
  Karawaci diberikan sebagai contoh (hlm. 76).
\end{itemize}

\subsection{Findings}\label{findings-4}

\textbf{Temuan yang dinyatakan penulis.}

\begin{itemize}
\tightlist
\item
  Jabodetabek menunjukkan tanda tahap awal \emph{post-suburbanization}:
  kawasan pinggiran tumbuh lebih cepat daripada inti, fungsi perkotaan
  makin bercampur, pusat kegiatan ekonomi baru berkembang, dan sejumlah
  kota komuter lama menjadi komunitas dengan basis ekonomi serta layanan
  yang lebih mandiri (hlm. 74 dan 76--77).
\item
  Polanya tetap hibrida, yakni mencampurkan permukiman suburban
  tradisional dengan perkembangan baru. Jabodetabek tidak sepenuhnya
  menyerupai kota-kota Barat karena banyak orang masih tinggal di inti
  tradisional dan pergi ke kawasan pinggiran untuk bekerja atau
  melakukan kegiatan lain (hlm. 77).
\item
  Perubahan fisik tersebut sangat berkaitan, menurut penulis, dengan
  meningkatnya peran sektor swasta dan pergeseran kekuasaan pembangunan
  lahan dari sektor publik ke sektor swasta. Privatisasi berarti
  pengembang lebih agresif membangun proyek serta memperoleh kewenangan
  praktis atas penguasaan, pengembangan, pengelolaan lahan, dan layanan
  perkotaan (hlm. 69 dan 76--77).
\item
  Kebijakan pemerintah pusat yang berorientasi pada pertumbuhan, promosi
  investasi asing, dan pembangunan infrastruktur memungkinkan arus modal
  serta ekspansi swasta. Desentralisasi selanjutnya memberi pemerintah
  daerah kewenangan lebih besar untuk bekerja sama dengan pengembang
  demi pertumbuhan lokal (hlm. 72 dan 76).
\item
  Otonomi lokal juga memperkuat fragmentasi. Pemerintah daerah cenderung
  berfokus pada wilayahnya sendiri, kurang memperhatikan rencana
  nasional atau daerah tetangga, dan dapat mengubah rencana tata ruang
  untuk mengakomodasi kepentingan pengembang (hlm. 75--77).
\end{itemize}

\textbf{Interpretasi penulis.} Kemandirian pinggiran dan privatisasi
tata kelola merupakan dua sisi dari proses yang sama. Kapasitas swasta
mempercepat penyediaan fungsi perkotaan yang tidak mampu atau tidak
segera disediakan pemerintah daerah, tetapi sekaligus menjadikan
rencana, layanan, dan ruang metropolitan lebih terpecah menurut
yurisdiksi serta kawasan privat (hlm. 69, 74, dan 76--77).

\subsection{Conclusions}\label{conclusions-4}

\textbf{Kesimpulan penulis.} Perkembangan mutakhir Jabodetabek mewakili
tahap awal, bukan bentuk matang atau salinan utuh, dari
\emph{post-suburbanization}. Bukti deskriptifnya terlihat pada limpahan
pertumbuhan dari Jakarta ke pinggiran, konversi lahan yang masif,
berkembangnya kawasan industri dan kota baru multifungsi, serta
berubahnya kota asrama menjadi pusat yang lebih mandiri. Ciri itu
diperkuat oleh investasi asing, kebijakan pertumbuhan pemerintah pusat,
otonomi lokal, dan pengalihan kekuasaan pembangunan serta layanan lahan
kepada sektor swasta (bagian ``Privatization of the Urban Fringes and
Post-Suburbanization in Jabodetabek,'' hlm. 76--77).

Penulis menyimpulkan bahwa sektor swasta dapat membantu memenuhi
kebutuhan regional akan perumahan, pekerjaan, perdagangan, pendidikan,
dan infrastruktur. Namun, karena pengembang berorientasi laba dan
pemerintah daerah dapat mengambil tindakan yang menimbulkan masalah
lintas wilayah, perencanaan tetap diperlukan untuk mencegah pelanggaran
tujuan kebijakan, fragmentasi, dan solusi regional yang suboptimal (hlm.
77).

\subsection{Contributions}\label{contributions-4}

\textbf{Kontribusi yang dinyatakan atau langsung ditunjukkan sumber.}

\begin{itemize}
\tightlist
\item
  Artikel memperbarui karya penulis sebelumnya mengenai unsur
  \emph{post-suburban} Jabodetabek dan memperluas rentang pembahasannya
  pada perubahan serta kesinambungan selama sekitar empat dekade
  (pendahuluan, hlm. 69).
\item
  Artikel menggeser perhatian dari bentuk spasial saja menuju hubungan
  antara perubahan fisik pinggiran, privatisasi pengembangan lahan,
  penyediaan layanan oleh pengembang, kebijakan pertumbuhan pusat, dan
  otonomi pemerintah daerah (hlm. 69 dan 76--77).
\item
  Artikel membawa perdebatan \emph{post-suburbia} ke konteks negara
  berkembang melalui perbandingan konseptual dengan Barat dan Tiongkok,
  sambil menolak anggapan bahwa lintasan Jabodetabek harus identik
  dengan salah satunya (hlm. 69--71 dan 77).
\item
  Bagi praktik perencanaan, artikel menunjukkan manfaat sekaligus risiko
  kapasitas swasta: kapasitas tersebut dapat memenuhi kebutuhan
  metropolitan, tetapi memerlukan perencanaan untuk menjaga tujuan
  publik dan koordinasi regional (abstrak, hlm. 68; hlm. 77).
\end{itemize}

\textbf{Inferensi review.} Kontribusi analitis terpentingnya ialah
memperlakukan \emph{post-suburbanization} sebagai perubahan tata kelola
dan distribusi kekuasaan, bukan sekadar desentralisasi penduduk atau
morfologi polisentris. Inferensi ini konsisten dengan argumentasi
artikel, tetapi merupakan rumusan sintesis review.

\subsection{Research Gap}\label{research-gap-4}

\textbf{Kesenjangan yang ditangani sumber.} Penulis menilai masih perlu
dijelaskan derajat kesesuaian perkembangan Jakarta dengan
\emph{post-suburbia} negara maju. Mereka juga menyatakan bahwa kajian
ini melampaui pekerjaan terdahulu mereka dengan memperbarui perkembangan
regional dan memberi fokus khusus pada peran swasta yang difasilitasi
oleh meningkatnya arti pemerintah daerah (pendahuluan, hlm. 69).

\textbf{Kesenjangan tersisa menurut inferensi review, bukan pernyataan
eksplisit penulis.} Artikel belum menguji secara kausal apakah
desentralisasi dan kebijakan pertumbuhan menghasilkan perubahan spasial
yang diamati; belum mengoperasionalkan tingkat
\emph{post-suburbanization} dalam ukuran yang dapat dibandingkan; dan
belum mengukur distribusi manfaat serta kerugian bagi kelompok di dalam
maupun di luar kawasan privat. Hubungan antara perubahan tata kelola,
mobilitas dua arah, segregasi sosial, perubahan lingkungan, dan hasil
regional juga disajikan melalui deskripsi serta contoh, bukan desain
evaluasi yang memungkinkan atribusi sebab-akibat.

\subsection{Limitations}\label{limitations-4}

\textbf{Keterbatasan eksplisit.} Penulis hanya menyebut bahwa pendekatan
penelitian bersifat deskriptif (abstrak, hlm. 68). Tidak ada bagian
khusus yang merinci keterbatasan metode. Dalam ucapan terima kasih,
penulis menyatakan bertanggung jawab atas kesalahan dan kekurangan,
tetapi tidak mengidentifikasi bentuk kekurangan tertentu (hlm. 77).

\textbf{Inferensi review berdasarkan isi file.}

\begin{itemize}
\tightlist
\item
  Artikel bergantung pada kumpulan data dan dokumen sekunder yang
  berasal dari tahun, unit geografis, dan jenis penerbit yang beragam.
  Prosedur pemilihan, pemeriksaan mutu, dan harmonisasi sumber tidak
  dijelaskan.
\item
  Tidak ada definisi operasional atau indeks yang menetapkan ambang
  kapan suatu wilayah dapat disebut berada pada tahap awal
  \emph{post-suburbanization}. Penetapan tahap awal karena itu bersifat
  interpretatif.
\item
  Tidak ada desain kausal, pembanding sistematis, atau analisis
  statistik yang dapat memisahkan pengaruh investasi, jalan tol,
  permintaan perumahan, desentralisasi, kebijakan pusat, dan tindakan
  pengembang.
\item
  Contoh Jababeka, Lippo, dan BSD kaya secara ilustratif, tetapi
  kriteria pemilihannya dan daya wakilnya terhadap seluruh Jabodetabek
  tidak disebutkan.
\item
  Dampak distribusional terhadap kelompok berpendapatan rendah, penduduk
  kampung atau komunitas di sekitar kawasan berpagar, serta mutu dan
  akses layanan di luar kawasan privat tidak diukur secara langsung.
\item
  Artikel menyebut kaitan konversi lahan dengan banjir secara
  berhati-hati, tetapi tidak menyajikan analisis lingkungan yang menguji
  hubungan tersebut (hlm. 74).
\end{itemize}

\subsection{Challenges}\label{challenges-4}

\textbf{Tantangan yang dilaporkan sumber.}

\begin{itemize}
\tightlist
\item
  Otonomi daerah membuat kewenangan pembangunan terpecah
  antaryurisdiksi, sedangkan pemerintah daerah cenderung mengutamakan
  rencana dan pertumbuhan wilayahnya sendiri daripada kepentingan
  metropolitan atau rencana daerah tetangga (hlm. 72 dan 76--77).
\item
  Hubungan patronase/klientelisme antara pemerintah dan pelaku swasta
  serta perlakuan fleksibel terhadap rencana tata ruang memungkinkan
  izin atau perubahan rencana yang tidak selalu sesuai dengan tujuan
  formal (hlm. 72 dan 75--76).
\item
  Pemerintah pusat memiliki kemampuan intervensi yang lebih kecil
  dibandingkan pada 1980-an dan 1990-an, sementara pengembang dapat
  menguasai lahan serta layanan di dalam kawasan privat. Kondisi ini
  menyulitkan koordinasi dan akuntabilitas pada skala regional (hlm.
  76--77).
\item
  Konversi sawah, hutan, dan kawasan resapan; segregasi komunitas
  berpagar; serta layanan yang eksklusif bagi penghuni menunjukkan
  tantangan lingkungan dan pemerataan yang menyertai pertumbuhan
  pinggiran (hlm. 73--76).
\item
  Orientasi laba membuat sektor swasta jarang memperhatikan rencana
  formal dan tujuan kebijakan publik kecuali diwajibkan, sedangkan
  tindakan pemerintah lokal dapat memindahkan biaya atau masalah ke
  wilayah lain (hlm. 77).
\end{itemize}

Tantangan pengumpulan data atau pelaksanaan penelitian yang dialami
penulis: Tidak disebutkan dalam file.

\subsection{Practical Implications}\label{practical-implications-4}

\textbf{Implikasi yang dinyatakan penulis.} Perencana di negara
berkembang perlu menyadari peningkatan cepat peran swasta. Mereka perlu
mengakui kemampuan pengembang membantu pemerintah memenuhi kebutuhan
perumahan, pekerjaan, pusat belanja, pendidikan, infrastruktur, dan
layanan, tetapi pada saat yang sama memastikan tindakan swasta tidak
memperburuk masalah regional atau memicu respons publik yang tidak
terkoordinasi (bagian ``Takeaway for practice,'' hlm. 68; hlm. 77).

Perencanaan guna lahan tetap memiliki fungsi penting untuk mengikat
pembangunan yang berorientasi laba pada rencana dan tujuan kebijakan
publik. Penulis menekankan perlunya perencanaan yang mencegah tindakan
lokal serta privat menghasilkan fragmentasi dan solusi suboptimal
seperti yang dikaitkan dengan \emph{post-suburbia} di Amerika Serikat
dan negara Barat lain (hlm. 77).

\textbf{Inferensi review.} Implikasi tersebut mengarah pada kebutuhan
koordinasi lintas batas administratif, penegakan konsistensi izin dengan
rencana, dan pengaturan kewajiban layanan serta dampak eksternal
pengembang. Instrumen kelembagaan atau regulasi tertentu untuk
melaksanakannya: Tidak disebutkan dalam file.

\subsection{Applications}\label{applications-4}

\textbf{Aplikasi yang didukung oleh isi sumber.} Kerangka artikel dapat
digunakan secara deskriptif untuk membaca transformasi kawasan pinggiran
metropolitan melalui gabungan indikator demografis, penggunaan lahan,
diferensiasi fungsi, arus komuter, perkembangan kawasan industri/kota
baru, dan perubahan pembagian kewenangan publik--swasta. Artikel juga
dapat menjadi bahan diagnosis awal bagi perencana yang menghadapi
ekspansi kawasan privat dan fragmentasi keputusan lintas pemerintah
daerah (hlm. 69--77).

\textbf{Inferensi review.} Pendekatan ini dapat diterapkan sebagai
kerangka pembanding awal pada wilayah metropolitan negara berkembang
lain, asalkan kesamaan tidak diasumsikan dan konfigurasi negara, pasar
tanah, hubungan pusat--daerah, serta keterikatan pinggiran--inti
diperiksa secara kontekstual. Artikel tidak menyediakan alat ukur,
prosedur replikasi, atau pedoman implementasi formal.

\subsection{Future Research
(Optional)}\label{future-research-optional-4}

Agenda penelitian lanjutan yang dinyatakan secara eksplisit oleh
penulis: Tidak disebutkan dalam file.

\textbf{Inferensi review, bukan rekomendasi eksplisit penulis.}
Keterbatasan bukti membuka ruang bagi penelitian longitudinal dan
komparatif yang mengoperasionalkan tingkat kemandirian pinggiran,
menguji mekanisme pergeseran kewenangan publik--swasta, membandingkan
antarkota/kabupaten, serta mengukur dampak sosial, fiskal, mobilitas,
layanan, dan lingkungan di dalam maupun di luar kawasan privat. Usulan
ini diturunkan dari celah metode dan bukti artikel, bukan dari bagian
agenda riset dalam sumber.

\subsection{Provenance Notes}\label{provenance-notes-4}

\begin{itemize}
\tightlist
\item
  Review ini hanya menggunakan seluruh isi file TXT J05 yang ditentukan
  pengguna, termasuk blok metadata PDF, teks artikel, catatan, ucapan
  terima kasih, dan daftar pustaka. Tidak ada sumber eksternal yang
  digunakan untuk menambah atau memverifikasi klaim.
\item
  File TXT mencatat bahwa teks diekstrak pada 31 Agustus 2026 dengan
  \texttt{pdftotext\ -layout} dari PDF 13 halaman. PDF tercatat tidak
  bertanda (\emph{not tagged}). Akses review ini adalah pada ekstraksi
  teks lengkap yang tersedia, bukan pemeriksaan visual langsung terhadap
  PDF.
\item
  Identitas bibliografis diambil dari blok sitasi dan kepala artikel:
  Tommy Firman dan Fikri Zul Fahmi; 2017; \emph{Journal of the American
  Planning Association} 83(1), 68--79; DOI
  10.1080/01944363.2016.1249010. Metadata internal PDF yang hanya
  menampilkan judul \texttt{RJPA\_A\_1249010.indd} dan penulis
  \texttt{user} tidak diperlakukan sebagai identitas bibliografis.
\item
  Locator mengacu pada nomor halaman tercetak artikel, bukan nomor baris
  TXT. Struktur halaman yang tersedia adalah abstrak dan pendahuluan
  pada hlm. 68--69; telaah \emph{post-suburbanization} pada hlm. 69--71;
  perkembangan Jabodetabek pada hlm. 71--76; sintesis privatisasi dan
  \emph{post-suburbanization} pada hlm. 76--77; serta referensi pada
  hlm. 77--79.
\item
  Satu ilustrasi yang dapat diidentifikasi ialah Gambar 1, peta kota dan
  kabupaten Jabodetabek, pada hlm. 71. Tidak ada tabel penelitian yang
  teridentifikasi dalam file. Karena gambar tidak diperiksa secara
  visual, review tidak mengekstraksi informasi kartografis di luar judul
  gambar dan uraian teks.
\item
  Angka, mata uang, nama hukum, dan tahun dipertahankan sesuai teks yang
  diekstrak. Nilainya tidak diverifikasi terhadap dokumen asal yang
  dikutip artikel. Beberapa salah ketik pada ekstraksi/teks referensi
  dan redaksi yang ambigu, termasuk pembagian angka migrasi sekitar
  190.000 orang, tidak diperbaiki melalui dugaan.
\item
  Label \textbf{Laporan sumber} menunjukkan apa yang dilaporkan artikel;
  \textbf{Interpretasi penulis} menunjukkan makna atau hubungan yang
  diajukan oleh penulis artikel; dan \textbf{Inferensi review}
  menunjukkan klasifikasi, evaluasi keterbatasan, atau kemungkinan
  penggunaan yang disusun dalam review ini. Inferensi tidak diperlakukan
  sebagai temuan empiris sumber.
\end{itemize}

\section{Review J06}\label{review-j06}

\subsection{Identitas Sumber}\label{identitas-sumber-1}

\begin{itemize}
\tightlist
\item
  \textbf{Kode:} J06.
\item
  \textbf{Judul:} \emph{Transportation Network and Changes in Urban
  Structure: Evidence from the Jakarta Metropolitan Area}.
\item
  \textbf{Penulis:} Muhammad Halley Yudhistira, Witri Indriyani, Andhika
  Putra Pratama, Yusuf Sofiyandi, dan Yusuf Reza Kurniawan.
\item
  \textbf{Afiliasi:} Institute for Economic and Social Research,
  Department of Economics, Universitas Indonesia; serta Demographic
  Institute, Department of Economics, Universitas Indonesia.
\item
  \textbf{Jurnal:} \emph{Research in Transportation Economics}, volume
  74 (2019), hlm. 52-63.
\item
  \textbf{DOI:} 10.1016/j.retrec.2018.12.003.
\item
  \textbf{Riwayat naskah:} diterima 26 Agustus 2017, direvisi 29
  November 2018, diterima untuk publikasi 22 Desember 2018, dan tersedia
  daring 31 Desember 2018.
\item
  \textbf{Wilayah, periode, dan unit analisis:} Jakarta Metropolitan
  Area (JMA); perubahan hasil utama pada 2000-2010; komunitas atau
  \emph{kelurahan} sebagai unit analisis.
\item
  \textbf{Nomor terbitan:} Tidak disebutkan dalam file.
\item
  \textbf{Status penelaahan sejawat:} Tidak disebutkan dalam file.
\end{itemize}

Locator identitas: metadata PDF dan halaman pertama artikel (hlm. 52).

\subsection{Summarized Abstract}\label{summarized-abstract-5}

\textbf{Laporan sumber.} Artikel ini menguji peran infrastruktur
transportasi dalam proses suburbanisasi di JMA pada tingkat komunitas.
Struktur perkotaan diukur melalui dua hasil, yaitu pertumbuhan kepadatan
penduduk dan pertumbuhan intensitas cahaya malam sebagai proksi
aktivitas ekonomi. Penulis menggunakan Sensus Penduduk 2000 dan 2010,
data cahaya malam pada dua tahun yang sama, ukuran akses jalan tol dan
rel, serta strategi variabel instrumental untuk menangani kemungkinan
hubungan timbal balik antara pembangunan transportasi dan perkembangan
wilayah (Abstrak, hlm. 52; Bagian 3.1-3.2, hlm. 55-56).

\textbf{Laporan hasil sumber.} Perbaikan akses jalan tol mendorong
pertumbuhan kepadatan penduduk terutama di \emph{city suburbs},
sedangkan akses rel mendorong pertumbuhan penduduk di \emph{other
suburbs} atau pinggiran nonkota. Suburbanisasi aktivitas ekonomi yang
ditunjukkan oleh cahaya malam hanya tampak kuat di pinggiran nonkota dan
lebih berkaitan dengan stok jaringan transportasi daripada dengan
perubahan aksesnya (Abstrak, hlm. 52; Bagian 4.2-4.3, hlm. 57-60; Bagian
5, hlm. 61).

\textbf{Interpretasi penulis.} Pola tersebut dipandang sebagai bukti
\emph{transportation-led suburbanization}: moda dan kondisi jaringan
yang berbeda memengaruhi lokasi pertumbuhan penduduk dan aktivitas
ekonomi secara berbeda. Namun, penulis menegaskan bahwa instrumen yang
digunakan tidak sepenuhnya mengatasi endogenitas sehingga estimasi 2SLS
perlu dibaca secara hati-hati (Bagian 4.1, hlm. 56; Bagian 4.5, hlm.
60-61).

\subsection{Problem Statement}\label{problem-statement-5}

\textbf{Laporan sumber.} JMA mengalami urbanisasi cepat sekaligus
perubahan distribusi internal. Walaupun populasi kawasan naik 14,5\%
pada 2000-2010, proporsi penduduk JMA yang tinggal di Jakarta turun dari
39,4\% menjadi 34,0\%. Kepadatan dan aktivitas industri juga berkembang
menuju daerah pinggiran. Pada saat yang sama, jaringan jalan tol dan rel
diperluas sehingga hubungan antara akses transportasi dan perubahan
struktur spasial perlu diuji secara kausal, bukan hanya dideskripsikan
(Bagian 1, hlm. 52-53).

\textbf{Laporan sumber.} Masalah identifikasi muncul karena
infrastruktur dapat memengaruhi kepadatan penduduk dan aktivitas
ekonomi, tetapi pertumbuhan penduduk atau ekonomi juga dapat menciptakan
permintaan atas infrastruktur baru. Selain itu, indikator ekonomi
konvensional tidak tersedia pada tingkat komunitas sehingga perubahan
aktivitas ekonomi lokal sulit diukur secara langsung (Bagian 1, hlm. 53;
Bagian 3.1-3.2, hlm. 55-56).

\textbf{Inferensi pengolahan.} Dengan demikian, persoalan inti artikel
adalah memisahkan pengaruh akses transportasi terhadap suburbanisasi
dari pemilihan lokasi infrastruktur yang mungkin sudah dipengaruhi oleh
perkembangan wilayah, sambil mencari ukuran aktivitas ekonomi yang cukup
rinci secara spasial. Inferensi ini merangkum rumusan masalah yang
tersebar dalam pendahuluan dan bagian identifikasi; artikel tidak
menyajikan satu kalimat pertanyaan penelitian formal.

\subsection{Objectives}\label{objectives-5}

\textbf{Laporan sumber.} Tujuan utama penelitian adalah mengestimasi
hubungan antara perbaikan akses transportasi dan perubahan struktur
perkotaan JMA pada 2000-2010 dengan sampel komunitas atau
\emph{kelurahan} (Bagian 1, hlm. 53).

Tujuan operasional yang dinyatakan atau langsung tercermin dalam desain
penelitian adalah:

\begin{itemize}
\tightlist
\item
  menguji apakah akses yang lebih baik ke ramp jalan tol dan stasiun rel
  meningkatkan pertumbuhan kepadatan penduduk dan aktivitas ekonomi;
\item
  membandingkan respons Jakarta, \emph{city suburbs}, dan
  \emph{other/non-city suburbs} untuk mengidentifikasi suburbanisasi;
\item
  membedakan peran jalan tol dan rel dalam perubahan struktur spasial;
\item
  menggunakan intensitas cahaya malam untuk menangkap dinamika aktivitas
  ekonomi pada tingkat komunitas ketika indikator ekonomi langsung tidak
  tersedia; dan
\item
  menangani potensi kausalitas terbalik melalui pendekatan variabel
  instrumental (Bagian 1, hlm. 53; Bagian 3.1-3.2, hlm. 55-56).
\end{itemize}

\subsection{Literature Survey}\label{literature-survey-5}

\textbf{Laporan sumber.} Literatur yang dibahas dapat dikelompokkan ke
dalam beberapa rumpun berikut.

\begin{itemize}
\tightlist
\item
  Literatur urbanisasi dan perluasan metropolitan menunjukkan bahwa
  perkembangan kota melampaui batas administratif serta terjadi di
  negara maju dan berkembang. Contoh yang dirujuk meliputi Tokyo,
  kota-kota di negara berkembang, dan megakota Asia (Bagian 1, hlm. 52).
\item
  Studi transportasi dan struktur perkotaan menunjukkan hubungan jalan
  atau rel dengan pertumbuhan kota, perubahan penggunaan lahan,
  produktivitas, suburbanisasi, desentralisasi, dan pembentukan
  subpusat. Artikel merujuk, antara lain, Atack et al.~(2009), Baum-Snow
  (2007), Duranton dan Turner (2012), Garcia-Lopez dan Angel (2012),
  Garcia-Lopez et al.~(2015, 2017), Datta (2012), Funderburg et
  al.~(2010), Holl (2016), serta Baum-Snow et al.~(2015) (Bagian 4.5,
  hlm. 60-61).
\item
  Untuk JMA, penelitian terdahulu oleh Firman (2004), Henderson dan
  Kuncoro (1996), serta Hudalah dan Firman (2012) terutama memberikan
  analisis deskriptif mengenai hubungan transportasi dan struktur
  perkotaan. Henderson dan Kuncoro melaporkan suburbanisasi manufaktur
  yang merespons jalan tol keluar Jakarta, sedangkan Hudalah dan Firman
  membahas kawasan industri dan transformasi pascasuburban (Bagian 1,
  hlm. 52; Bagian 4.5, hlm. 60).
\item
  Penggunaan cahaya malam sebagai proksi aktivitas ekonomi didasarkan
  pada Henderson, Storeygard, dan Weil (2012) serta Storeygard (2016).
  Sumber juga mencatat penggunaan data serupa untuk memperkirakan
  konsumsi energi dan populasi (Bagian 2, hlm. 54-55).
\item
  Literatur metodologis mengenai endogenitas dan variabel instrumental
  mencakup Wooldridge (2010), Murray (2006, 2010), Duranton dan Turner
  (2012), Garcia-Lopez et al.~(2015), serta kriteria instrumen lemah
  Stock-Yogo (Bagian 3.2 dan 4.1, hlm. 56).
\end{itemize}

\textbf{Interpretasi penulis.} Dibandingkan penelitian sebelumnya, hasil
jalan tol dalam artikel ini dinilai konsisten arah umumnya dengan bukti
dari Amerika Serikat, Spanyol, dan Tiongkok, meskipun besar estimasinya
relatif kecil dibandingkan rentang dalam Garcia-Lopez dan Angel (2012).
Sebaliknya, peran rel di JMA dinilai lebih lemah daripada temuan
kenaikan kepadatan sebesar 4\% untuk setiap kilometer lebih dekat ke
stasiun di Paris. Penulis juga menyatakan tidak menemukan studi yang
benar-benar sebanding dalam menggabungkan cahaya malam, strategi
estimasi serupa, dan suburbanisasi (Bagian 4.5, hlm. 60-61).

\subsection{Method Used}\label{method-used-5}

\textbf{Laporan sumber.} Penelitian mengadaptasi model regresi
Garcia-Lopez dan Angel (2012). Variabel hasil dinyatakan sebagai
perubahan logaritmik antara 2000 dan 2010 pada komunitas yang sama.
Estimasi awal menggunakan OLS, kemudian 2SLS dengan variabel
instrumental digunakan karena OLS diperkirakan menghasilkan parameter
akses yang tidak konsisten akibat kausalitas terbalik (Bagian 3.1-3.2,
Persamaan 1, hlm. 55-56).

Komponen desain estimasinya adalah:

\begin{itemize}
\tightlist
\item
  ukuran jalan tol berupa perubahan jarak garis lurus dari pusat
  komunitas ke ramp terdekat; teks model menyatakan perubahan akses
  tersebut menggunakan jarak 2010 relatif terhadap 1990;
\item
  ukuran rel berupa jarak ke stasiun rel terdekat pada 1990, yang
  digunakan untuk menangkap peningkatan frekuensi layanan pada jalur
  yang ada;
\item
  kontrol kondisi awal berupa logaritma kepadatan penduduk 2000 atau
  logaritma intensitas cahaya malam 2000;
\item
  kontrol jarak berupa jarak ke pusat Jakarta, pusat distrik, dan garis
  pantai;
\item
  kontrol geografis berupa indeks kekasaran medan, luas lahan, dan
  elevasi;
\item
  kontrol historis berupa populasi distrik pada 1961;
\item
  tiga kandidat instrumen berupa jarak ke Jalan Anyer-Panarukan, jalan
  primer lama, dan stasiun rel lama; serta
\item
  galat baku robust dan uji instrumen lemah menggunakan statistik F
  Kleibergen-Paap terhadap nilai kritis Stock-Yogo (Bagian 3.1-3.2, hlm.
  55-56; catatan Tabel 3-6, hlm. 57-60).
\end{itemize}

\textbf{Laporan sumber.} Instrumen dipilih melalui regresi tahap pertama
dan \emph{reduced form}. Dua instrumen digunakan untuk dua variabel
akses guna menghindari masalah overidentifikasi, tetapi pasangan
instrumen berbeda menurut zona dan hasil. Stasiun rel lama dipertahankan
sebagai kandidat umum; instrumen kedua dipilih antara Jalan
Anyer-Panarukan dan jalan primer lama berdasarkan signifikansi serta
statistik F Kleibergen-Paap. Tabel 4 memakai rel lama dan
Anyer-Panarukan untuk Jakarta serta \emph{city suburbs}, lalu rel lama
dan jalan primer lama untuk kepadatan penduduk di \emph{other suburbs}.
Tabel 5 memakai rel lama dan jalan primer lama untuk cahaya malam di
\emph{city suburbs}, sedangkan Jakarta dan \emph{other suburbs} memakai
rel lama dan Anyer-Panarukan (Bagian 4.1, hlm. 56-57; Tabel 4, hlm. 58;
Tabel 5, hlm. 59).

\textbf{Laporan sumber.} Analisis dilakukan untuk keseluruhan JMA dan
secara terpisah untuk Jakarta, \emph{city suburbs}, dan \emph{other
suburbs}. Untuk cahaya malam, penulis juga menguji spesifikasi
alternatif menggunakan stok akses pada 2000, bukan hanya perubahannya,
karena pertumbuhan aktivitas ekonomi diduga lebih baik dijelaskan oleh
posisi terhadap jaringan yang sudah tersedia (Bagian 4.3, Tabel 6, hlm.
59-60; Lampiran B, hlm. 61-62).

\subsection{Dataset}\label{dataset-5}

\textbf{Laporan sumber.} Data utama dan turunannya meliputi:

\begin{itemize}
\tightlist
\item
  Sensus Penduduk Indonesia 2000 dan 2010 untuk membentuk jumlah serta
  kepadatan penduduk pada tingkat komunitas;
\item
  data \emph{stable light} Defense Meteorological Satellite Program yang
  dikaitkan dalam file dengan National Oceanic and Atmospheric Agency
  (NOAA), sedangkan Gambar 4-5 mencantumkan National Centers for
  Environmental Information (2016), sebagai proksi aktivitas ekonomi
  pada 2000 dan 2010;
\item
  peta digital vektor dari Geospatial Information Bureau of Indonesia
  untuk luas komunitas, posisi jaringan jalan tol dan rel, jarak akses,
  serta variabel jarak dan geografis;
\item
  data jaringan jalan dari Indonesia Highway Authority (2015);
\item
  data jaringan dan layanan rel dari PT Kereta Api Indonesia dan PT KAI
  Commuter Jabodetabek (2016); dan
\item
  populasi distrik 1961 dari Winarso (2010) sebagai kontrol historis
  (Bagian 2, Tabel 1-2 dan Gambar 1, hlm. 53-54; Bagian 3.1, hlm.
  55-56).
\end{itemize}

Data cahaya malam memiliki nilai piksel 0-63. Penulis menumpangtindihkan
data tersebut dengan peta JMA dan menghitung rata-rata nilai piksel di
dalam setiap komunitas. Resolusi piksel disebut berkisar 0,55-2,77 km
persegi, sedangkan rata-rata luas komunitas sekitar 4,55 km persegi;
atas dasar itu, penulis menilai proksi tersebut masih representatif pada
tingkat komunitas (Bagian 3.1, hlm. 55, termasuk catatan kaki 3).

\textbf{Laporan sumber.} Dalam konteks perubahan jaringan 1990-2010,
artikel menyebut pembangunan 45 km jalan tol baru dengan 26 ramp serta
perluasan jaringan rel sepanjang 81 km dengan pembukaan 22 stasiun. Pada
uraian kontekstual 2000-2010, 567 komunitas disebut menjadi rata-rata
8,78 km lebih dekat ke ramp, sedangkan 942 komunitas menjadi rata-rata
12,69 km lebih dekat ke stasiun (Bagian 2, hlm. 53; Bagian 3.1, hlm.
55).

\textbf{Batas akses data.} Lampiran C menyatakan adanya data pelengkap
daring, tetapi isi data pelengkap tersebut tidak terdapat dalam TXT yang
ditinjau (Lampiran C, hlm. 62).

\subsection{Population / Sample}\label{population-sample-5}

\textbf{Laporan sumber.} Sampel yang dinyatakan pada pendahuluan terdiri
atas 1.483 komunitas atau \emph{kelurahan} di JMA untuk 2000 dan 2010.
JMA mencakup 13 wilayah munisipal: enam wilayah Jakarta, yaitu Jakarta
Pusat, Jakarta Utara, Jakarta Timur, Jakarta Selatan, Jakarta Barat, dan
Kepulauan Seribu; serta tujuh wilayah pinggiran, yaitu Kabupaten Bekasi,
Kota Bekasi, Kota Depok, Kabupaten Bogor, Kota Bogor, Kabupaten
Tangerang, dan Kota Tangerang (Bagian 1 dan catatan kaki 1, hlm. 52-53).

Analisis membagi observasi menjadi Jakarta, \emph{city suburbs}, dan
\emph{other suburbs} atau \emph{non-city suburbs}. Tabel 3 dan Tabel 4
melaporkan 261 observasi Jakarta, 340 \emph{city suburbs}, dan 871
\emph{other suburbs}, dengan total 1.472 observasi untuk model kepadatan
penduduk. Tabel 5 melaporkan 261 dan 340 observasi untuk dua kelompok
pertama; pada kelompok \emph{other suburbs}, OLS memuat 871 observasi,
sedangkan spesifikasi IV memuat 882 observasi. Lampiran B juga
melaporkan 1.472 observasi untuk kepadatan penduduk dan 1.483 untuk
cahaya malam (Tabel 3, hlm. 57; Tabel 4, hlm. 58; Tabel 5, hlm. 59;
Tabel B, hlm. 62).

\textbf{Catatan pengolahan.} Alasan berkurangnya 11 observasi dalam
model kepadatan penduduk dan perbedaan 871 versus 882 observasi pada
model cahaya malam tidak dijelaskan dalam file. File juga tidak
memberikan pemetaan eksplisit ketujuh wilayah pinggiran ke kategori
\emph{city suburbs} dan \emph{other/non-city suburbs}. Oleh karena itu,
komposisi kategori tidak ditafsirkan lebih jauh.

\subsection{Dependent Variables}\label{dependent-variables-5}

\textbf{Laporan sumber.} Terdapat dua variabel dependen pada tingkat
komunitas:

\begin{itemize}
\tightlist
\item
  \textbf{Perubahan logaritma kepadatan penduduk}, yaitu
  \texttt{ln(kepadatan\ penduduk\ 2010)\ -\ ln(kepadatan\ penduduk\ 2000)}.
  Kepadatan dibentuk dari jumlah penduduk sensus dan luas lahan
  komunitas; satuan deskriptifnya adalah penduduk per hektare.
\item
  \textbf{Perubahan logaritma intensitas cahaya malam}, yaitu
  \texttt{ln(intensitas\ cahaya\ malam\ 2010)\ -\ ln(intensitas\ cahaya\ malam\ 2000)}.
  Intensitas komunitas merupakan rata-rata nilai piksel \emph{stable
  light} di dalam batas komunitas dan dipakai sebagai proksi aktivitas
  ekonomi (Bagian 3.1, Persamaan 1, hlm. 55; Tabel 4-6, hlm. 58-60;
  Lampiran A, hlm. 61).
\end{itemize}

\textbf{Kualifikasi sumber.} Cahaya malam bukan ukuran ekonomi langsung.
Artikel menggunakannya karena indikator seperti PDB tidak tersedia pada
tingkat komunitas, dan menafsirkan pertumbuhannya sebagai perubahan
aktivitas ekonomi umum (Bagian 1, hlm. 53; Bagian 2-3.1, hlm. 54-55).

\subsection{Results}\label{results-5}

\textbf{Hasil deskriptif yang dilaporkan sumber.} Kepadatan penduduk JMA
mencapai 79,4 penduduk per hektare pada 2010, naik 14,8 penduduk per
hektare atau 23\% dibandingkan 2000. Rata-rata pertumbuhan kepadatan
penduduk pada 2000-2010 lebih tinggi di pinggiran nonkota (39,0\%)
daripada di \emph{city suburbs} (34,3\%). Cakupan cahaya malam pada 2010
juga lebih terang dan lebih luas daripada pada 2000, dengan kenaikan
terbesar terkonsentrasi di pinggiran yang berdekatan dengan Jakarta
(Bagian 2, Gambar 2-5, hlm. 53-55).

Lampiran A melaporkan perubahan rata-rata kepadatan sebesar 14,8
penduduk per hektare untuk seluruh JMA, 18,9 untuk Jakarta, 27,1 untuk
\emph{city suburbs}, dan 8,9 untuk \emph{other suburbs}. Perubahan
rata-rata intensitas cahaya malam masing-masing sebesar 7,8, 0,3, 6,9,
dan 10,5. Angka ini bersifat deskriptif dan tidak dengan sendirinya
menunjukkan pengaruh kausal transportasi (Lampiran A, Tabel A, hlm. 61).

\textbf{Tahap pertama dan \emph{reduced form}.} Sebagian besar instrumen
secara statistik menjelaskan variasi akses jalan tol dan rel pada tahap
pertama. Namun, beberapa koefisien instrumen pada \emph{reduced form}
tidak berbeda secara statistik dari nol. Penulis menyimpulkan bahwa
instrumen bekerja relatif baik dalam banyak spesifikasi, tetapi tidak
secara sempurna menyelesaikan endogenitas (Bagian 4.1 dan Tabel 3, hlm.
56-57).

\textbf{Hasil kepadatan penduduk.} Estimasi 2SLS menunjukkan pola yang
berbeda menurut zona dan moda.

\begin{itemize}
\tightlist
\item
  Di \emph{city suburbs}, setiap perbaikan akses jalan tol sejauh satu
  kilometer dikaitkan dengan pertumbuhan kepadatan penduduk sebesar
  3,7-4,8\%, bergantung pada kontrol. Dalam spesifikasi penuh, koefisien
  perubahan jarak ke ramp adalah -0,048 dan bertanda signifikan menurut
  notasi tabel (Bagian 4.2 dan Tabel 4, Panel B, hlm. 57-58).
\item
  Di \emph{other suburbs}, jarak satu kilometer lebih dekat ke stasiun
  rel dikaitkan dengan pertumbuhan kepadatan penduduk sebesar 1,3-1,7\%.
  Koefisien spesifikasi penuh adalah -0,014 dan bertanda signifikan.
  Sebaliknya, perbaikan akses jalan tol berkaitan dengan pertumbuhan
  kepadatan yang lebih rendah; koefisien perubahan jarak pada
  spesifikasi penuh adalah 0,150 dan bertanda signifikan (Bagian 4.2 dan
  Tabel 4, Panel C, hlm. 57-58).
\item
  Untuk Jakarta, koefisien jalan tol dan rel pada spesifikasi penuh
  tidak berbeda secara statistik dari nol. Koefisien rel di \emph{city
  suburbs} juga tidak signifikan dalam spesifikasi penuh (Tabel 4, Panel
  A-B, hlm. 58).
\item
  Untuk sampel JMA secara keseluruhan, Lampiran B melaporkan bahwa
  perbaikan akses jalan tol berasosiasi dengan pertumbuhan kepadatan
  yang lebih rendah, sedangkan akses rel yang lebih baik berasosiasi
  dengan pertumbuhan kepadatan yang lebih tinggi (Lampiran B dan Tabel
  B, hlm. 61-62).
\end{itemize}

\textbf{Hasil intensitas cahaya malam.} Penulis menyebut hasil 2SLS
untuk cahaya malam kurang intuitif daripada hasil populasi dan
menyatakan bukti yang kuat dari OLS hanya robust pada beberapa
spesifikasi di \emph{other suburbs} (Bagian 4.3, hlm. 58-59).

\begin{itemize}
\tightlist
\item
  Dalam spesifikasi penuh 2SLS untuk \emph{other suburbs}, koefisien
  perubahan jarak ke ramp jalan tol adalah 0,131, sedangkan koefisien
  jarak ke stasiun rel adalah -0,017; keduanya bertanda signifikan
  menurut notasi tabel. Dengan definisi variabel jarak artikel, hasil
  ini berarti perbaikan akses jalan tol berkaitan dengan pertumbuhan
  cahaya malam yang lebih rendah, sedangkan kedekatan ke rel berkaitan
  dengan pertumbuhan yang lebih tinggi (Tabel 5, Panel C, hlm. 59).
\item
  Dalam spesifikasi penuh untuk Jakarta dan \emph{city suburbs},
  koefisien akses jalan tol dan rel tidak berbeda secara statistik dari
  nol (Tabel 5, Panel A-B, hlm. 59).
\item
  Spesifikasi alternatif berbasis stok jaringan pada 2000 menghasilkan
  koefisien jarak ke ramp sebesar -0,009 dan jarak ke stasiun sebesar
  -0,005 untuk \emph{other suburbs} dalam estimasi IV. Penulis membaca
  tanda ini sebagai bukti bahwa komunitas yang sejak awal lebih dekat ke
  jaringan mengalami pertumbuhan cahaya malam lebih cepat (Bagian 4.3
  dan Tabel 6, hlm. 59-60).
\item
  Pada tingkat JMA, Lampiran B melaporkan hubungan negatif antara
  perbaikan akses jalan tol dan pertumbuhan cahaya malam, hubungan
  positif untuk akses rel yang lebih baik, serta pertumbuhan cahaya
  malam yang lebih cepat pada lokasi yang lebih jauh dari pusat Jakarta
  (Lampiran B dan Tabel B, hlm. 61-62).
\end{itemize}

\textbf{Kualifikasi pembacaan angka.} Catatan Tabel 3-6 mendefinisikan
simbol \texttt{***}, \texttt{**}, dan \texttt{*} masing-masing sebagai
signifikansi 10\%, 5\%, dan 1\%, yaitu urutan yang tidak lazim tetapi
dilaporkan demikian dalam file. Review ini mempertahankan status umum
``bertanda signifikan menurut notasi tabel'' dan tidak mengoreksi notasi
tersebut menggunakan asumsi dari luar file.

\subsection{Findings}\label{findings-5}

\textbf{Temuan utama yang dilaporkan sumber.} Infrastruktur transportasi
berkaitan dengan suburbanisasi JMA, tetapi dampaknya tidak seragam
menurut moda, zona, dan hasil yang diukur (Bagian 4.2-4.5, hlm. 57-61).

\begin{itemize}
\tightlist
\item
  Perbaikan akses jalan tol terutama mendorong pertumbuhan penduduk di
  \emph{city suburbs}, bukan di Jakarta atau pinggiran nonkota.
\item
  Akses rel mendorong pertumbuhan penduduk di \emph{other/non-city
  suburbs} dan juga berkaitan dengan pertumbuhan cahaya malam di zona
  tersebut.
\item
  Suburbanisasi aktivitas ekonomi, sebagaimana diproksikan oleh cahaya
  malam, lebih jelas di \emph{other suburbs} daripada di Jakarta atau
  \emph{city suburbs}.
\item
  Untuk cahaya malam, kedekatan terhadap stok jaringan pada 2000
  memberikan pola yang lebih konsisten daripada perubahan akses jalan
  tol selama periode penelitian.
\item
  Perbedaan antara OLS dan 2SLS menunjukkan bahwa penempatan
  infrastruktur bersifat endogen; dalam banyak kasus, besar koefisien
  2SLS melebihi OLS, walaupun beberapa tanda koefisien juga berubah
  (Bagian 4.4, hlm. 60).
\end{itemize}

\textbf{Interpretasi penulis.} Rel dinilai lebih andal dan lebih
menghemat waktu untuk perjalanan jarak jauh sehingga lebih relevan bagi
pertumbuhan penduduk di pinggiran yang lebih jauh. Penulis juga
menghubungkan pertumbuhan cahaya malam dekat stok jalan tol di
\emph{other suburbs} dengan perkembangan kawasan industri, khususnya di
Kabupaten Bekasi. Adapun hubungan negatif antara perbaikan akses jalan
tol dan cahaya malam ditafsirkan sebagai kemungkinan perpindahan
aktivitas baru menjauh dari lokasi ramp yang sudah padat (Bagian
4.2-4.3, hlm. 57-60).

\textbf{Interpretasi penulis.} Perbedaan OLS dan 2SLS dijelaskan melalui
kemungkinan bahwa pemerintah membangun jalan atau stasiun baru di
komunitas yang pada awalnya berkepadatan dan beraktivitas ekonomi rendah
karena biaya pengadaan lahannya lebih murah. Menurut penulis, lokasi
yang semula kurang berkembang kemudian memperoleh efek akses yang lebih
besar daripada wilayah yang sudah berkembang (Bagian 4.4, hlm. 60).

\textbf{Inferensi pengolahan.} Bukti artikel paling kuat mendukung
heterogenitas spasial dan moda, bukan klaim bahwa setiap penambahan
jaringan selalu mempercepat pertumbuhan di lokasi terdekat. Perbedaan
antara variabel perubahan akses dan stok akses, terutama untuk jalan tol
dan cahaya malam, merupakan syarat penting bagi pembacaan temuan
tersebut.

\subsection{Conclusions}\label{conclusions-5}

\textbf{Kesimpulan sumber.} Dengan pendekatan variabel instrumental,
penulis menyimpulkan bahwa jalan tol dan rel membantu menjelaskan
perubahan struktur spasial JMA. Perbaikan akses jalan tol meningkatkan
pertumbuhan penduduk di \emph{city suburbs}. Kedekatan terhadap stok
jalan tol menjelaskan pertumbuhan cahaya malam di \emph{other suburbs}.
Akses rel mendorong pertumbuhan penduduk dan cahaya malam di \emph{other
suburbs} (Bagian 5, hlm. 61).

\textbf{Interpretasi penulis.} Temuan tersebut dipandang sebagai bukti
bahwa perbaikan transportasi mendorong penyebaran perkotaan pada kota
utama yang relatif muda di Asia. Penulis menduga sumber pertumbuhan
menyebar itu adalah migrasi masuk dari provinsi lain, tetapi file tidak
menyajikan pengujian migrasi sebagai bagian dari model (Bagian 5, hlm.
61).

\textbf{Kualifikasi.} Kesimpulan kausal harus dibaca bersama peringatan
penulis bahwa instrumen tidak sempurna dan estimasi 2SLS mungkin masih
mengandung bias (Bagian 4.1, hlm. 56). Karena itu, hasil lebih tepat
dipahami sebagai bukti empiris yang mendukung \emph{transportation-led
suburbanization} dalam konteks JMA, bukan sebagai efek yang seragam atau
bebas dari ketidakpastian identifikasi.

\subsection{Contributions}\label{contributions-5}

\textbf{Kontribusi yang diklaim sumber.} Artikel menyatakan kontribusi
berikut.

\begin{itemize}
\tightlist
\item
  Memperluas pengukuran suburbanisasi dari pertumbuhan kepadatan
  penduduk ke pertumbuhan intensitas cahaya malam sebagai proksi
  aktivitas ekonomi pada tingkat komunitas (Bagian 1, hlm. 53; Bagian 5,
  hlm. 61).
\item
  Menyediakan bukti empiris berbasis IV mengenai hubungan transportasi
  dan perubahan struktur perkotaan JMA, ketika studi JMA terdahulu
  dinilai lebih banyak bersifat deskriptif (Bagian 4.5, hlm. 60-61).
\item
  Menunjukkan bahwa jalan tol dan rel memiliki peran yang berbeda
  menurut lapisan pinggiran: jalan tol untuk pertumbuhan penduduk di
  \emph{city suburbs}, sedangkan rel dan stok jaringan lebih menonjol di
  \emph{other suburbs} (Bagian 4.5 dan 5, hlm. 60-61).
\item
  Menambah bukti mengenai infrastruktur dan suburbanisasi dari kota
  utama di negara berkembang, sehingga melengkapi literatur yang banyak
  membahas kota di negara maju (Bagian 1 dan 4.5, hlm. 53 dan 60-61).
\end{itemize}

\textbf{Kualifikasi pengolahan.} Klaim kebaruan dibatasi pada apa yang
dinyatakan penulis. File menyebut bahwa penulis tidak menemukan studi
yang benar-benar sebanding untuk cahaya malam dan strategi estimasi ini,
tetapi tidak menyediakan prosedur tinjauan sistematis yang dapat
membuktikan keunikan tersebut (Bagian 4.5, hlm. 61).

\subsection{Research Gap}\label{research-gap-5}

\textbf{Kesenjangan yang dinyatakan sumber.} Penelitian sebelumnya telah
menghubungkan transportasi dengan pertumbuhan perkotaan, perubahan
lahan, produktivitas, desentralisasi, dan pembentukan subpusat. Namun,
untuk JMA, bukti mengenai hubungan transportasi dan struktur perkotaan
dinilai terutama deskriptif. Keterbatasan data juga membuat hubungan
tersebut baru relatif belakangan diuji secara empiris (Bagian 1 dan 4.5,
hlm. 52 dan 60-61).

\textbf{Kesenjangan pengukuran.} Indikator kegiatan ekonomi seperti PDB
tidak tersedia pada tingkat komunitas. Artikel mengisi sebagian
kesenjangan itu dengan cahaya malam dan menyatakan tidak menemukan
penelitian yang sangat sebanding dalam memakai proksi tersebut dengan
strategi estimasi yang sama untuk menjelaskan suburbanisasi (Bagian 1,
hlm. 53; Bagian 4.5, hlm. 61).

\textbf{Kesenjangan konteks.} Bukti kausal terdahulu yang dibandingkan
banyak berasal dari Amerika Serikat, Spanyol, Paris, dan Tiongkok.
Artikel memosisikan JMA sebagai tambahan bukti dari kota utama di negara
berkembang dan konteks megakota Asia (Bagian 1 dan 4.5, hlm. 53 dan
60-61).

\textbf{Inferensi pengolahan mengenai gap yang tersisa.} Data hasil
hanya mengamati penduduk dan cahaya malam, bukan migrasi individu,
lokasi perusahaan, perjalanan komuter, atau keluaran ekonomi langsung.
Oleh sebab itu, mekanisme migrasi, relokasi perusahaan, penghematan
waktu, dan biaya lahan yang dipakai untuk menjelaskan koefisien tetap
merupakan interpretasi, bukan mekanisme yang diuji langsung oleh model
dalam file.

\subsection{Limitations}\label{limitations-5}

\textbf{Keterbatasan yang diakui penulis.} Strategi IV tidak sepenuhnya
mengatasi endogenitas. Beberapa instrumen tidak signifikan pada
\emph{reduced form}, pilihan instrumen harus berubah menurut sampel dan
hasil, dan penulis menyatakan estimasi 2SLS mungkin masih bias sehingga
koreksinya perlu ditafsirkan hati-hati (Bagian 4.1, hlm. 56-57).

\textbf{Keterbatasan yang diakui atau diperlihatkan sumber.} Estimasi
rel dinilai relatif lemah dalam menjelaskan suburbanisasi JMA jika
dibandingkan dengan temuan untuk Paris. Hasil cahaya malam 2SLS juga
disebut kurang intuitif, sedangkan bukti OLS hanya robust pada beberapa
spesifikasi di \emph{other suburbs} (Bagian 4.3 dan 4.5, hlm. 58-61).

\textbf{Inferensi pengolahan berbasis data yang dilaporkan.} Cahaya
malam adalah proksi, bukan ukuran langsung PDB atau aktivitas
perusahaan. Nilainya dirata-ratakan pada komunitas dari piksel
beresolusi 0,55-2,77 km persegi; meskipun penulis menilainya
representatif, agregasi ini membatasi ketepatan spasial yang dapat
diklaim (Bagian 3.1, hlm. 55).

\textbf{Inferensi pengolahan berbasis desain.} Dua gelombang hasil, 2000
dan 2010, menangkap perubahan bersih selama satu dekade tetapi tidak
menunjukkan waktu perubahan di dalam periode tersebut. Ukuran akses
berupa jarak garis lurus juga tidak secara langsung mengukur waktu
tempuh, kualitas layanan, biaya perjalanan, atau kemacetan. Variabel
tersebut memang tidak tersedia sebagai hasil model dalam file.

\textbf{Catatan pengolahan.} File tidak menjelaskan kehilangan 11
observasi pada model kepadatan, perbedaan jumlah observasi OLS dan IV
untuk cahaya malam di \emph{other suburbs}, atau pemetaan rinci wilayah
administratif ke dua kategori pinggiran. Hal ini membatasi reproduksi
sampel hanya dari TXT (Tabel 3-5 dan Tabel B, hlm. 57-59 dan 62).

\textbf{Batas generalisasi.} Penelitian hanya menguji JMA. Penulis
mengaitkan relevansinya dengan kota utama Asia, tetapi file tidak
menguji metropolitan lain. Generalisasi di luar JMA karena itu merupakan
potensi aplikasi, bukan hasil empiris langsung.

\subsection{Challenges}\label{challenges-5}

Artikel tidak memiliki subbagian khusus bernama \emph{Challenges}.
Tantangan berikut dilaporkan dalam penjelasan masalah dan metode:

\begin{itemize}
\tightlist
\item
  tidak tersedianya indikator ekonomi tingkat komunitas, yang mendorong
  penggunaan cahaya malam sebagai proksi;
\item
  kausalitas terbalik karena pertumbuhan penduduk dan ekonomi dapat
  memicu pembangunan transportasi, bukan hanya menjadi akibatnya;
\item
  ketergantungan jaringan baru pada jaringan lama, yang membuat lokasi
  akses terkini tidak acak;
\item
  kemungkinan endogenitas jarak stasiun 1990 karena peningkatan
  frekuensi layanan dapat merespons pertumbuhan di sekitar stasiun;
\item
  kebutuhan memilih instrumen berbeda menurut zona dan hasil serta
  menguji instrumen lemah; dan
\item
  penyelarasan data raster cahaya malam dengan batas komunitas yang
  ukurannya relatif kecil (Bagian 1 dan 3.1-3.2, hlm. 53 dan 55-56;
  Bagian 4.1, hlm. 56-57).
\end{itemize}

\textbf{Inferensi pengolahan.} Tanda dan besar koefisien yang berubah
antara OLS, 2SLS, spesifikasi perubahan akses, dan spesifikasi stok
jaringan membuat interpretasi substantif lebih menantang. Artikel
meresponsnya dengan memisahkan zona serta moda dan dengan membahas
kemungkinan lokasi infrastruktur yang endogen (Bagian 4.3-4.4, hlm.
58-60).

\subsection{Practical Implications}\label{practical-implications-5}

\textbf{Implikasi yang dinyatakan penulis.} Penulis menyatakan bahwa
pencarian paket kebijakan untuk menyeimbangkan efek penyebaran perkotaan
dan peningkatan mobilitas penduduk akibat perbaikan transportasi relevan
bagi pemerintah daerah JMA. Relevansi ini didasarkan pada temuan bahwa
investasi transportasi dapat memindahkan pertumbuhan penduduk dan
aktivitas menuju pinggiran, bukan hanya meningkatkan akses di pusat
(Bagian 5, hlm. 61).

\textbf{Inferensi pengolahan.} Karena dampak berbeda menurut moda dan
zona, evaluasi kebijakan sebaiknya tidak memperlakukan semua investasi
transportasi sebagai intervensi yang identik. Temuan artikel menyiratkan
perlunya memperhatikan pertumbuhan penduduk di \emph{city suburbs}
sekitar jalan tol serta pertumbuhan penduduk dan cahaya malam di
\emph{other suburbs} yang terhubung rel atau sudah dekat dengan stok
jaringan. Ini adalah implikasi yang diturunkan dari hasil, bukan
rekomendasi operasional terperinci yang diuji oleh sumber.

\textbf{Batas implikasi.} File tidak menyebut instrumen tata ruang
tertentu, skema pembiayaan, pembagian kewenangan, target layanan, atau
paket kebijakan spesifik. Karena itu, bentuk implementasi praktisnya
tidak dapat dirinci hanya dari sumber ini.

\subsection{Applications}\label{applications-5}

\textbf{Aplikasi yang dinyatakan sumber.} Penulis menyatakan hasilnya
berpotensi membantu pemahaman proses suburbanisasi pada kota utama Asia
seperti Jakarta dan menambah bukti internasional bahwa suburbanisasi
yang dipimpin transportasi juga terjadi di negara berkembang (Bagian 1,
hlm. 53).

\textbf{Aplikasi metodologis yang dinyatakan atau langsung ditunjukkan
sumber.} Penggabungan data sensus, peta jaringan, dan cahaya malam dapat
digunakan untuk menilai perubahan penduduk dan aktivitas ekonomi pada
skala komunitas ketika data ekonomi lokal tidak tersedia. Pemisahan
antara perubahan akses dan stok akses juga dapat membedakan efek
pembukaan atau perbaikan jaringan dari keuntungan lokasi terhadap
jaringan yang sudah ada (Bagian 3.1 dan 4.3, hlm. 55 dan 59-60).

\textbf{Inferensi pengolahan.} Kerangka tersebut dapat menjadi contoh
desain evaluasi untuk wilayah metropolitan lain, tetapi validitas
instrumen historis, definisi zona, kualitas cahaya malam, dan konteks
kelembagaan harus diuji ulang pada setiap penerapan. File tidak
melakukan aplikasi lintas kota.

\subsection{Future Research
(Optional)}\label{future-research-optional-5}

Tidak disebutkan dalam file.

\subsection{Provenance Notes}\label{provenance-notes-5}

\begin{itemize}
\tightlist
\item
  Review ini hanya menggunakan seluruh isi file TXT
  \texttt{source/journal/J06-yudhistira-indriyani-pratama-sofiyandi-kurniawan-2019-transportation-network-and-changes-in-urban-structure-evidence-from-the-jakarta-metropolitan-area-10.1016-j.retrec.2018.12.003.txt},
  yang berjumlah 925 baris dan memuat metadata serta teks artikel dari
  hlm. 52-63.
\item
  Metadata ekstraksi menyatakan bahwa TXT dibuat pada 31 Agustus 2026
  dari PDF 12 halaman dengan \texttt{pdftotext\ -layout}. Review
  dilakukan terhadap TXT, bukan melalui pemeriksaan visual PDF.
\item
  Tidak ada sumber eksternal, laman DOI, berkas PDF, atau data pelengkap
  yang digunakan. Lampiran C hanya memberitahukan keberadaan data
  pelengkap daring; isinya tidak tersedia dalam file yang ditinjau (hlm.
  62).
\item
  Locator halaman mengikuti nomor halaman tercetak jurnal. Locator
  bagian, tabel, gambar, persamaan, dan lampiran mengikuti label yang
  muncul di TXT.
\item
  Tata letak dua kolom menyebabkan beberapa urutan kalimat, rumus,
  simbol, nama sumber gambar, dan judul kolom terpotong atau tersusun
  silang. Review hanya menormalkan pemenggalan yang maknanya jelas dan
  tidak merekonstruksi bagian yang meragukan dengan informasi dari luar.
\item
  Semua angka hasil diparaphrasakan dari teks, tabel, atau lampiran
  dalam TXT. Status \texttt{Laporan\ sumber},
  \texttt{Interpretasi\ penulis}, \texttt{Inferensi\ pengolahan}, dan
  \texttt{Catatan\ pengolahan} digunakan untuk mencegah pencampuran
  hasil empiris, penjelasan penulis, serta evaluasi review.
\item
  Informasi yang tidak dapat dipastikan dari file mencakup nomor
  terbitan, status penelaahan sejawat, alasan perbedaan jumlah
  observasi, pemetaan eksplisit kategori pinggiran, dan agenda riset
  lanjutan.
\end{itemize}

\section{Review J07: Highway Expansion and Urban Sprawl in the Jakarta
Metropolitan
Area}\label{review-j07-highway-expansion-and-urban-sprawl-in-the-jakarta-metropolitan-area}

\subsection{Identitas Sumber}\label{identitas-sumber-2}

\begin{itemize}
\tightlist
\item
  \textbf{Kode:} J07.
\item
  \textbf{Judul:} \emph{Highway expansion and urban sprawl in the
  Jakarta Metropolitan Area}.
\item
  \textbf{Penulis:} Andhika Putra Pratama, Muhammad Halley Yudhistira,
  dan Eric Koomen.
\item
  \textbf{Afiliasi:} Institute for Economic and Social Research serta
  Research Cluster of Urban and Transportation Economics, Universitas
  Indonesia; dan Department of Spatial Economics, Spatial Information
  Laboratory (SPINlab), Vrije Universiteit Amsterdam {[}hlm. 1, blok
  judul dan afiliasi{]}.
\item
  \textbf{Publikasi:} \emph{Land Use Policy}, volume 112, tahun 2022,
  nomor artikel 105856 {[}hlm. 1, header jurnal dan metadata PDF{]}.
  Nomor isu dan rentang halaman bibliografis: Tidak disebutkan dalam
  file. Metadata PDF menyatakan dokumen terdiri atas 11 halaman.
\item
  \textbf{DOI:} 10.1016/j.landusepol.2021.105856 {[}hlm. 1{]}.
\item
  \textbf{Riwayat naskah:} diterima 20 Desember 2020, direvisi 1
  September 2021, diterima untuk publikasi 31 Oktober 2021, dan tersedia
  daring 18 November 2021 {[}hlm. 1{]}.
\item
  \textbf{Kata kunci:} urban development, urban expansion, urban sprawl,
  instrumental variable, transport access, dan the Jakarta Metropolitan
  Area {[}hlm. 1{]}.
\item
  \textbf{Status penelaahan sejawat:} Tidak disebutkan dalam file.
\end{itemize}

\subsection{Summarized Abstract}\label{summarized-abstract-6}

\textbf{Laporan sumber.} Artikel meneliti hubungan kausal yang diklaim
antara perluasan jalan tol dan dua dimensi perkembangan perkotaan di
Jakarta Metropolitan Area (JMA): ekspansi perkotaan, yakni pertambahan
lahan terbangun, dan \emph{urban sprawl}, yakni perkembangan terbangun
yang makin tersebar atau terisolasi. Untuk mengatasi kemungkinan bahwa
pembangunan perkotaan justru mendorong pembangunan jalan tol, penulis
menggunakan kedekatan terhadap infrastruktur transportasi historis
sebagai variabel instrumental. Penulis melaporkan bahwa peningkatan
akses ke jalan tol meningkatkan \emph{sprawl} di dalam JMA, terutama
pada radius 40 hingga 50 km dari pusat Jakarta. Bukti ekspansi lahan
perkotaan yang dipicu jalan tol tidak muncul untuk keseluruhan JMA,
tetapi ditemukan pada subsampel wilayah pinggiran berstatus kota
(\emph{city suburbs}) {[}hlm. 1, Abstract; hlm. 7-9, Bagian 5.2 dan
Tabel 4-7{]}.

\textbf{Interpretasi penulis.} Jalan tol dinilai menghambat perkembangan
menuju bentuk kota yang lebih kompak karena lahan terbangun baru lebih
mungkin muncul terpisah dari lahan perkotaan yang sudah ada. Penulis
juga menafsirkan efek yang lebih kuat di wilayah luar sebagai
konsekuensi kemungkinan dari konsentrasi pembangunan jalan tol di
pinggiran, nilai tanah yang lebih rendah, dan regulasi konversi lahan
yang tidak seketat wilayah pusat {[}hlm. 8-10, Bagian 5.2.2, 5.2.3, dan
6{]}.

\textbf{Inferensi pengolahan.} Bukti artikel paling kuat mendukung
hubungan antara akses jalan tol dan pola spasial pembangunan yang makin
tersebar. Bukti untuk pertambahan total lahan terbangun lebih bersyarat
secara geografis dan tidak mendukung generalisasi bahwa jalan tol
meningkatkan luas terbangun di seluruh JMA.

\subsection{Problem Statement}\label{problem-statement-6}

\textbf{Laporan sumber.} Literatur telah banyak menghubungkan penurunan
biaya perjalanan dan pembangunan jaringan transportasi dengan ekspansi
kota serta suburbanisasi. Namun, pengetahuan kausal tentang pengaruh
jaringan transportasi terhadap pola penggunaan lahan di dalam
metropolitan besar negara berkembang masih terbatas. Metropolitan
seperti JMA berkembang cepat, bersifat luas dan polisentris, serta
memperlihatkan pergantian pembangunan berkepadatan tinggi dan rendah
pada jarak berdekatan, sehingga temuan antarkota dari negara maju belum
langsung menjelaskan dinamika intrametropolitan tersebut {[}hlm. 1-2,
Bagian 1{]}.

Masalah empiris utamanya adalah endogenitas: estimasi OLS dapat keliru
karena pembangunan perkotaan mungkin mendorong perluasan jalan tol,
bukan hanya sebaliknya. Selain itu, ekspansi lahan terbangun dan
\emph{sprawl} bukan konsep yang identik. Ekspansi dapat berlangsung
secara kompak, sedangkan peningkatan \emph{sprawl} dapat terjadi
meskipun tambahan luas terbangun lebih kecil {[}hlm. 4-5, Bagian 3.1,
Gambar 3, dan Bagian 4{]}.

\textbf{Interpretasi penulis.} Kasus JMA relevan karena pertumbuhan
kawasan selama sekitar tiga dekade terjadi bersama perluasan jaringan
jalan tol, suburbanisasi, dan tata ruang yang dipandang kurang mengikat
serta terfragmentasi antarkewenangan {[}hlm. 2-3, Bagian 2.1-2.2{]}.

\subsection{Objectives}\label{objectives-6}

\textbf{Laporan sumber.} Tujuan utama penelitian adalah mengestimasi
efek kausal perluasan atau peningkatan akses jalan tol terhadap
\emph{urban sprawl} di dalam JMA. Tujuan ini dijabarkan menjadi dua
sasaran empiris: menguji pengaruh akses jalan tol terhadap perubahan
indeks \emph{sprawl} dan menguji pengaruhnya terhadap ekspansi lahan
perkotaan {[}hlm. 1-2, Abstract dan Bagian 1{]}.

Penelitian juga bertujuan membedakan dampak transportasi di dalam kota
dari bukti terdahulu yang terutama membandingkan antarkota, menilai
heterogenitas efek pada Jakarta, \emph{city suburbs}, dan \emph{other
suburbs}, serta menguji ketahanan hasil terhadap perubahan wilayah
studi, radius dari pusat Jakarta, pemangkasan observasi ekstrem, ukuran
unit spasial, dan sumber data tutupan lahan {[}hlm. 2, Bagian 1; hlm.
8-9, Bagian 5.2.2-5.2.4{]}.

\subsection{Literature Survey}\label{literature-survey-6}

\begin{itemize}
\tightlist
\item
  \textbf{Landasan teoretis yang dilaporkan sumber:} model kota
  monosentris memprediksi penduduk bergerak menjauhi pusat ketika biaya
  komuter turun, sehingga kawasan perkotaan meluas dan suburbanisasi
  meningkat. Artikel mengaitkan mekanisme ini dengan Glaeser dan Kahn
  (2003), Baum-Snow (2007), Garcia-Lopez (2012), Garcia-Lopez et
  al.~(2015), dan Yudhistira et al.~(2019) {[}hlm. 1, Bagian 1{]}.
\item
  \textbf{Pengukuran awal:} Brueckner dan Fansler (1983) serta McGrath
  (2005) dilaporkan menghubungkan biaya transportasi dengan skala
  spasial metropolitan. Burchfield et al.~(2006) memperkenalkan indeks
  berupa persentase lahan belum terbangun di sekitar lahan permukiman
  untuk menangkap pembangunan yang tersebar {[}hlm. 1-2, Bagian 1{]}.
\item
  \textbf{Peralihan menuju identifikasi kausal:} artikel menyatakan
  studi korelasional terdahulu belum menunjukkan kausalitas.
  Garcia-Lopez (2019) menggunakan jaringan transportasi historis sebagai
  instrumen dan menemukan bahwa stok jalan tol memperluas lahan
  permukiman, meningkatkan fragmentasi petak perkotaan, serta mendorong
  pembangunan pada lahan belum terbangun di sekitar permukiman;
  cakupannya, menurut artikel, masih berupa perbedaan antarkota {[}hlm.
  2, Bagian 1{]}.
\item
  \textbf{Konteks JMA:} studi terdahulu tentang Jakarta dilaporkan
  menggunakan kepadatan penduduk, kegiatan ekonomi, atau data yang
  relatif kasar dan terutama menggambarkan tren dasar. Artikel
  memosisikan GHSL beresolusi tinggi dan dua indikator perkembangan
  perkotaan sebagai kemajuan atas keterbatasan tersebut {[}hlm. 2,
  Bagian 1; hlm. 9-10, Bagian 6{]}.
\item
  \textbf{Dampak lebih luas:} artikel merujuk literatur yang mengaitkan
  \emph{sprawl} dengan lingkungan dan konsumsi energi, produktivitas,
  penyediaan barang publik, kriminalitas, mobilitas sosial-ekonomi, dan
  ketimpangan. Rujukan ini berfungsi sebagai motivasi kebijakan;
  dampak-dampak tersebut tidak diuji dalam penelitian J07 {[}hlm. 1,
  catatan kaki 1; hlm. 10, Bagian 7{]}.
\item
  \textbf{Batas survei:} metode pencarian, kriteria seleksi, jumlah
  korpus, dan penilaian mutu literatur tidak disebutkan dalam file.
  Bagian ini merupakan tinjauan naratif dalam pendahuluan, bukan
  \emph{systematic literature review} yang dijelaskan metodenya.
\end{itemize}

\subsection{Method Used}\label{method-used-6}

\textbf{Laporan sumber.} Penelitian menggunakan analisis regresi atas
unit spasial untuk mengestimasi perubahan antartahun. Model dasar
mengestimasi perubahan logaritmik indikator perkembangan perkotaan
sebagai fungsi perubahan jarak ke pintu jalan tol terdekat, kondisi awal
perkembangan perkotaan, variabel jarak, dan variabel geografis. Variabel
jarak kontrol meliputi jarak ke stasiun kereta terdekat, pusat Jakarta,
pusat distrik, dan garis pantai; variabel geografis meliputi elevasi dan
\emph{terrain ruggedness index} {[}hlm. 4-5, Persamaan 1 dan Bagian 4;
hlm. 7, catatan Tabel 2-4{]}.

Tahapan analisisnya adalah sebagai berikut.

\begin{itemize}
\tightlist
\item
  OLS dipakai sebagai hasil deskriptif awal. Spesifikasi yang disukai
  penulis mengendalikan kondisi awal, seluruh variabel jarak, dan
  variabel geografis {[}hlm. 5-7, Bagian 4 dan 5.1, Tabel 2{]}.
\item
  Untuk menghadapi kausalitas terbalik, penulis menerapkan
  \emph{two-stage least squares} (TSLS). Tahap pertama memprediksi
  perubahan akses jalan tol dengan jarak ke infrastruktur transportasi
  historis; tahap kedua menggunakan perubahan akses terprediksi untuk
  mengestimasi perubahan indeks \emph{sprawl} atau luas terbangun
  {[}hlm. 5, Persamaan 2-3 dan Bagian 4{]}.
\item
  Instrumen utama adalah jarak ke jalan sekunder dan jalan kolektor
  tahun 1924. Untuk subsampel tertentu, penulis juga menggunakan jarak
  ke Jalan Anyer-Panarukan tahun 1815 dan stasiun kereta lama {[}hlm.
  5-8, Gambar 4, Tabel 3-5{]}.
\item
  Relevansi instrumen diperiksa melalui hubungan tahap pertama dan
  statistik F Kleibergen-Paap. Plausibilitas eksogenitas diperiksa
  dengan meregresikan jumlah sekolah dasar, pusat kesehatan distrik,
  pasar, dan rumah tangga miskin tahun 2000 terhadap jarak ke jalan
  historis; tidak ada koefisien yang ditandai signifikan pada Tabel 1
  {[}hlm. 5-6, Bagian 4 dan Tabel 1{]}.
\item
  Penulis memakai galat baku \emph{robust} dan menandai signifikansi
  pada tingkat 10\%, 5\%, dan 1\% {[}hlm. 6-10, catatan Tabel 1-8{]}.
\item
  Uji ketahanan mengubah cakupan wilayah, radius dari pusat Jakarta,
  observasi ekor distribusi, ukuran kisi 2 km menjadi 3 km, dan data
  GHSL menjadi ESA CCI. Uji falsifikasi mengacak penempatan observasi
  instrumen untuk memeriksa kemungkinan efek plasebo {[}hlm. 8-10,
  Bagian 5.2.3-5.2.4 dan Tabel 6-8{]}.
\end{itemize}

\textbf{Interpretasi penulis.} Selisih log digunakan untuk mengendalikan
heterogenitas tak teramati yang tetap sepanjang waktu. Penulis memilih
spesifikasi kontrol penuh dan memprioritaskan kekuatan hubungan tahap
pertama ketika menentukan instrumen {[}hlm. 5, Bagian 4{]}.

\textbf{Pilihan model yang dinyatakan sumber.} Meskipun jarak ke stasiun
kereta dimasukkan sebagai kontrol, penulis mengabaikan kemungkinan efek
kausal kereta terhadap \emph{sprawl} dengan alasan hasil studi
Yudhistira et al.~(2019) hanya sedikit signifikan secara statistik
{[}hlm. 5, Bagian 4{]}.

\subsection{Dataset}\label{dataset-6}

\begin{itemize}
\tightlist
\item
  \textbf{Global Human Settlement Layer (GHSL):} data Joint Research
  Centre, European Commission, yang berasal dari citra Landsat tahun
  1990, 2000, dan 2014 dengan resolusi spasial 38 m. Setiap sel bernilai
  satu jika diklasifikasikan terbangun dan nol jika tidak {[}hlm. 4,
  Bagian 3.1.1{]}.
\item
  \textbf{Jaringan transportasi kontemporer dan kontrol jarak:} data
  topografi OpenStreetMap dipakai sebagai dasar untuk jalan tol atau
  pintu tol, stasiun kereta, pusat distrik, dan garis pantai, kemudian
  diperiksa silang dengan National Highway Authority of Indonesia, KAI
  Commuter Jabodetabek, dan Badan Informasi Geospasial (BIG) {[}hlm. 4,
  Bagian 3.1.1{]}.
\item
  \textbf{Infrastruktur historis:} peta jalan kolonial digunakan untuk
  memperoleh jalan historis, termasuk jaringan tahun 1815 dan 1924;
  Gambar 4 menyebut Dutch Colonial Maps, Leiden University Library
  (2019) sebagai sumber {[}hlm. 5, Gambar 4; hlm. 7-8, Tabel 3-5{]}.
\item
  \textbf{Topografi:} peta elevasi digital resmi BIG digunakan untuk
  elevasi dan \emph{terrain ruggedness index}. File menyebut resolusi
  terbaik sebagai satuan spasial 12 m²; tidak ada penjelasan tambahan
  untuk memastikan apakah notasi tersebut merujuk luas sel atau dimensi
  linear {[}hlm. 4, Bagian 3.1.1{]}.
\item
  \textbf{Uji eksogenitas:} National Village Potential Survey tahun 2000
  menyediakan jumlah sekolah dasar, pusat kesehatan distrik, pasar, dan
  rumah tangga miskin pada tingkat sub-distrik {[}hlm. 6, Bagian 4 dan
  Tabel 1{]}.
\item
  \textbf{Data alternatif untuk uji ketahanan:} ESA Climate Change
  Initiative (ESA CCI) dengan resolusi 300 m dibandingkan dengan GHSL 38
  m {[}hlm. 8-9, Bagian 5.2.3 dan Tabel 7{]}.
\item
  \textbf{Data kontekstual:} sensus penduduk BPS 2000 dan 2010 digunakan
  untuk menghitung pertumbuhan penduduk JMA, tetapi file tidak
  menyatakan bahwa data penduduk menjadi variabel hasil dalam regresi
  utama {[}hlm. 2, Bagian 2.1 dan catatan kaki 2-3{]}.
\item
  \textbf{Statistik data lengkap:} file merujuk deskripsi variabel dan
  statistik utama pada Appendix 1, tetapi isi Appendix 1 tidak terdapat
  dalam TXT. Karena itu, rincian statistik selain yang tampil pada tabel
  utama tidak dapat diverifikasi dari file ini.
\end{itemize}

\subsection{Population / Sample}\label{population-sample-6}

\textbf{Laporan sumber.} Populasi analitis bukan individu atau rumah
tangga, melainkan wilayah JMA yang dibagi menjadi unit spasial seragam
seluas 4 km² atau 2 km × 2 km untuk mengurangi \emph{modifiable areal
unit problem}. Sampel utama terdiri atas 1.421 unit spasial {[}hlm. 4,
Bagian 3.1.1; hlm. 7, Tabel 2-4{]}.

Pembagian subsampel administratif untuk analisis ekspansi lahan adalah
Jakarta sebanyak 155 observasi, \emph{city suburbs} sebanyak 215
observasi, dan \emph{other suburbs} sebanyak 1.051 observasi; totalnya
sama dengan 1.421 {[}hlm. 8, Tabel 5{]}. Untuk uji radius, jumlah
observasi yang terbaca adalah 504 dalam radius 30 km, 866 dalam radius
40 km, dan 1.241 dalam radius 50 km {[}hlm. 9, Tabel 6{]}. Uji ukuran
sel menggunakan 635 observasi untuk kisi 3 km × 3 km, sedangkan uji ESA
CCI menggunakan 703 observasi pada kisi 2 km dan 367 observasi pada kisi
3 km {[}hlm. 9, Tabel 7{]}.

Tabel falsifikasi melaporkan 1.524 observasi untuk indeks \emph{sprawl}
JMA tetapi 1.421 untuk luas terbangun; alasan perbedaan jumlah ini tidak
disebutkan dalam file {[}hlm. 10, Tabel 8{]}. Kriteria rinci inklusi
atau penghapusan sel, daftar batas administratif JMA, dan penanganan sel
tanpa lahan terbangun: Tidak disebutkan dalam file utama yang tersedia.

\subsection{Dependent Variables}\label{dependent-variables-6}

\begin{enumerate}
\def\labelenumi{\arabic{enumi}.}
\tightlist
\item
  \textbf{Perubahan logaritmik indeks \emph{sprawl}.} Indeks dasar
  dihitung dengan mencari persentase ruang terbuka pada delapan sel yang
  langsung bertetangga dengan setiap sel terbangun, lalu
  merata-ratakannya untuk seluruh sel terbangun dalam unit analisis.
  Kenaikan indeks berarti lebih banyak sel terbangun terisolasi dari
  lahan artifisial lain dan dengan demikian menunjukkan pola pembangunan
  yang lebih tersebar {[}hlm. 4, Bagian 3.1 dan Gambar 3; hlm. 4-5,
  Persamaan 1{]}. Tabel utama menampilkan rerata perubahan log indeks
  sebesar -0,282 {[}hlm. 7, Tabel 2 dan 4{]}.
\item
  \textbf{Perubahan logaritmik luas terbangun.} Ekspansi lahan perkotaan
  didefinisikan sebagai pertambahan luas perkotaan di suatu wilayah
  relatif terhadap luas awal, tanpa memperhatikan apakah pola
  tambahannya kompak atau tersebar. Dalam estimasi, indikator ditulis
  sebagai perubahan log luas terbangun dalam km² {[}hlm. 4-5, Bagian 3.1
  dan Persamaan 1; hlm. 7, Tabel 2 dan 4{]}. Tabel utama menampilkan
  rerata sebesar 0,342 {[}hlm. 7, Tabel 2 dan 4{]}.
\end{enumerate}

\textbf{Pembedaan konseptual sumber.} Gambar 3 menunjukkan bahwa
ekspansi besar dapat menurunkan indeks \emph{sprawl} bila sel baru
mengisi ruang secara kompak, sedangkan ekspansi yang lebih kecil dapat
menghasilkan indeks lebih tinggi bila sel baru terpencar {[}hlm. 4,
Gambar 3{]}. Variabel hasil mengenai kepadatan penduduk, perilaku
perjalanan, nilai tanah, penggunaan energi, atau kesejahteraan: Tidak
disebutkan sebagai variabel dependen penelitian ini.

\subsection{Results}\label{results-6}

\textbf{Laporan hasil statistik sumber.} Karena perubahan akses dihitung
sebagai perubahan jarak ke pintu tol terdekat, wilayah yang mengalami
perbaikan lebih besar memiliki nilai perubahan jarak yang lebih negatif.
Oleh sebab itu, koefisien negatif pada indeks \emph{sprawl} ditafsirkan
penulis sebagai kenaikan \emph{sprawl} ketika akses membaik {[}hlm. 7,
Bagian 5.1{]}.

\begin{itemize}
\tightlist
\item
  \textbf{Pemeriksaan eksogenitas:} jarak ke jalan sekunder dan kolektor
  tahun 1924 tidak memiliki hubungan yang ditandai signifikan dengan
  jumlah sekolah dasar, pusat kesehatan distrik, pasar, maupun rumah
  tangga miskin tahun 2000. Penulis menilai hasil ini membuat asumsi
  eksogenitas instrumen masuk akal, bukan membuktikannya secara mutlak
  {[}hlm. 6, Tabel 1 dan Bagian 4{]}.
\item
  \textbf{OLS untuk indeks \emph{sprawl}:} koefisien perubahan jarak
  adalah -0,048, -0,046, dan -0,038 pada spesifikasi bertahap; semuanya
  signifikan pada tingkat 1\%. Spesifikasi kontrol penuh menghasilkan
  -0,038, yang oleh penulis dibaca sebagai kenaikan 3,8\% dalam
  pembangunan tersebar untuk perbaikan akses 1 km {[}hlm. 6-7, Bagian
  5.1 dan Tabel 2 kolom 1-3{]}.
\item
  \textbf{OLS untuk luas terbangun:} koefisiennya 0,054, 0,026, dan
  0,016, seluruhnya signifikan pada tingkat 1\%. Karena perbaikan akses
  berarti perubahan jarak negatif, spesifikasi penuh ditafsirkan sebagai
  pengurangan ekspansi perkotaan sebesar 1,6\% untuk perbaikan 1 km.
  Penulis segera memperingatkan bahwa pola OLS ini mungkin tidak sah
  akibat kausalitas terbalik {[}hlm. 7, Bagian 5.1 dan Tabel 2 kolom
  4-6{]}.
\item
  \textbf{Tahap pertama TSLS:} jarak ke jalan sekunder 1924 memiliki
  koefisien -0,128 dan -0,178, sedangkan jarak ke jalan kolektor 1924
  memiliki koefisien -0,188 dan -0,213; semuanya signifikan pada tingkat
  1\%. Penulis menafsirkan kedekatan terhadap jalan historis sebagai
  terkait dengan perubahan akses tol yang lebih kecil, sehingga jalan
  tol baru cenderung dibangun menjauh dari jaringan jalan lama {[}hlm.
  7, Bagian 5.2.1 dan Tabel 3{]}.
\item
  \textbf{TSLS utama untuk indeks \emph{sprawl}:} koefisiennya -0,066
  dengan instrumen jalan sekunder 1924 dan -0,096 dengan instrumen jalan
  kolektor 1924, masing-masing signifikan pada tingkat 5\% dan 1\%.
  Penulis menerjemahkannya sebagai peningkatan indeks \emph{sprawl}
  sebesar 6,6\% hingga 9,6\% akibat perbaikan akses 1 km. Statistik F
  Kleibergen-Paap sebesar 38,83 dan 38,55 digunakan untuk menyatakan
  instrumen tidak lemah {[}hlm. 7-8, Bagian 5.2.2 dan Tabel 4 kolom
  1-2{]}.
\item
  \textbf{TSLS utama untuk luas terbangun:} koefisien -0,032 dan -0,009
  tidak ditandai signifikan, dengan statistik F 56,02 dan 53,99. Dengan
  demikian, model TSLS utama tidak memberi bukti statistik bahwa
  peningkatan akses tol menyebabkan ekspansi luas terbangun pada
  keseluruhan JMA {[}hlm. 7-8, Bagian 5.2.2 dan Tabel 4 kolom 3-4{]}.
\item
  \textbf{Heterogenitas ekspansi:} pada subsampel Jakarta dan
  \emph{other suburbs}, koefisien TSLS tidak berbeda secara statistik
  dari nol. Pada \emph{city suburbs}, koefisien TSLS -0,016 signifikan
  pada tingkat 1\%, yang ditafsirkan sebagai ekspansi lahan perkotaan
  1,6\% akibat perbaikan akses 1 km. Statistik F untuk Jakarta,
  \emph{city suburbs}, dan \emph{other suburbs} masing-masing 19,43,
  30,08, dan 29,55 {[}hlm. 8, Bagian 5.2.2 dan Tabel 5{]}.
\item
  \textbf{Heterogenitas menurut radius:} efek terhadap indeks
  \emph{sprawl} tidak signifikan dalam radius 30 km dari pusat Jakarta,
  tetapi signifikan sebesar 4,7\% dalam radius 40 km dan 10,2\% dalam
  radius 50 km. Hasil ini menopang klaim bahwa efek lebih kentara di
  wilayah yang lebih jauh dari pusat {[}hlm. 8-9, Bagian 5.2.3 dan Tabel
  6 kolom 6-8{]}.
\item
  \textbf{Perubahan cakupan dan observasi:} penambahan distrik tetangga
  umumnya menghasilkan arah yang konsisten dengan model utama dan
  besaran yang dalam narasi sumber disebut sekitar 8\% hingga 12\%.
  Kombinasi JMA dengan Serang memberi tanda berlawanan, tetapi statistik
  F Kleibergen-Paap-nya rendah sehingga penulis menganggap instrumen
  lemah. Hasil tetap signifikan ketika 5\% atau 10\% observasi teratas
  dikeluarkan, tetapi menjadi tidak signifikan ketika 5\% atau 10\%
  observasi terbawah dikeluarkan {[}hlm. 8-9, Bagian 5.2.3 dan Tabel
  6{]}.
\item
  \textbf{Ukuran sel dan sumber data:} koefisien indeks \emph{sprawl}
  untuk GHSL adalah -0,096 pada sel 2 km dan -0,066 pada sel 3 km; yang
  kedua hanya signifikan pada tingkat 10\%. Dengan ESA CCI, koefisien
  menjadi -0,348 pada sel 2 km dan -0,295 pada sel 3 km, keduanya
  signifikan pada tingkat 1\%. Penulis mengaitkan besaran yang jauh
  lebih tinggi pada ESA CCI dengan resolusi 300 m yang dapat melebihkan
  keberadaan ruang terbuka kecil dan mengabaikan variasi lokal {[}hlm.
  8-9, Bagian 5.2.3 dan Tabel 7{]}.
\item
  \textbf{Ketahanan hasil ekspansi:} penulis menyatakan kebanyakan
  spesifikasi pada Appendix 5 tidak menghasilkan bukti statistik
  ekspansi perkotaan yang dipicu transportasi untuk keseluruhan JMA.
  Pengecualian muncul ketika JMA digabung dengan Lebak atau Serang,
  dengan besaran yang dilaporkan 6,8\% hingga 9,9\%; tabel Appendix 5
  sendiri tidak terdapat dalam TXT sehingga rincian koefisiennya tidak
  dapat diperiksa {[}hlm. 8-9, akhir Bagian 5.2.3{]}.
\item
  \textbf{Uji falsifikasi:} setelah observasi instrumen diacak, tidak
  ada estimasi yang berbeda secara statistik dari nol. Penulis
  menafsirkan ini sebagai penolakan atas kemungkinan bahwa temuan utama
  berasal dari efek plasebo {[}hlm. 9-10, Bagian 5.2.4 dan Tabel 8{]}.
\end{itemize}

\subsection{Findings}\label{findings-6}

\textbf{Temuan yang dilaporkan sumber.} Pertama, akses jalan tol yang
membaik berhubungan secara kausal menurut strategi IV penulis dengan
peningkatan pembangunan terbangun yang tersebar di JMA. Meskipun
rata-rata indeks \emph{sprawl} menurun pada 1990-2014, wilayah dengan
peningkatan akses tol yang lebih besar mengalami konvergensi menuju pola
kompak yang lebih lambat; lahan terbangun baru lebih mungkin terisolasi
dari lahan perkotaan yang telah ada {[}hlm. 7-8, Bagian 5.2.2; hlm. 10,
Bagian 7{]}.

Kedua, efek tersebut heterogen secara spasial. Bukti tidak terlihat
dalam radius 30 km, muncul pada radius 40 km, dan lebih besar pada
radius 50 km. Penulis mengaitkan pola ini dengan pusat yang sudah lebih
kompak dan pembangunan jalan tol baru yang terkonsentrasi di pinggiran
Jakarta {[}hlm. 8-10, Bagian 5.2.3 dan 6{]}.

Ketiga, perluasan jalan tol tidak terbukti menyebabkan pertambahan total
luas terbangun pada sampel JMA secara keseluruhan. Bukti kausal yang
signifikan ditemukan khusus pada \emph{city suburbs}, dengan estimasi
ekspansi 1,6\% untuk perbaikan akses 1 km {[}hlm. 7-10, Tabel 4-5 dan
Bagian 6{]}.

\textbf{Interpretasi penulis.} Kombinasi hasil tersebut berarti
perkembangan JMA berlangsung melalui dua bentuk yang dapat berimpit
tetapi harus dibedakan: penyebaran atau isolasi pembangunan baru pada
skala metropolitan dan ekspansi luas perkotaan pada wilayah pinggiran
tertentu {[}hlm. 10, Bagian 7{]}.

\textbf{Inferensi pengolahan.} Klaim utama tentang \emph{sprawl} lebih
konsisten lintas spesifikasi daripada klaim tentang luas ekspansi.
Namun, konsistensi arah tidak sama dengan kestabilan besaran karena
estimasi berubah cukup besar menurut resolusi data, ukuran sel, radius,
dan observasi yang dikeluarkan.

\subsection{Conclusions}\label{conclusions-6}

\textbf{Kesimpulan penulis.} Dengan menggunakan data spasial runtun
waktu dan infrastruktur transportasi historis sebagai instrumen, artikel
menyimpulkan bahwa peningkatan akses jalan tol menghasilkan pembangunan
lahan terbangun baru yang lebih tersebar di JMA. Ketika jalan tol
bertambah, lahan perkotaan baru lebih mungkin terpisah dari kawasan
terbangun yang sudah ada. Hasil ini dinyatakan \emph{robust} pada
berbagai spesifikasi yang dipilih penulis {[}hlm. 10, Bagian 7{]}.

Penulis juga menyimpulkan adanya ekspansi perkotaan yang dipicu jalan
tol di \emph{city suburbs}. Karena itu, jalan tol dikaitkan dengan dua
perubahan bentuk kota: peningkatan ketersebaran pembangunan pada JMA dan
peningkatan luas lahan terbangun pada bagian pinggiran berstatus kota
{[}hlm. 10, Bagian 7{]}.

\textbf{Kualifikasi pengolahan.} Kesimpulan ekspansi tidak berlaku
seragam untuk seluruh JMA: TSLS agregat, Jakarta, dan \emph{other
suburbs} tidak signifikan. Pernyataan \emph{robust} juga perlu dibaca
bersama hilangnya signifikansi setelah observasi terbawah dikeluarkan
serta perubahan besar besaran pada data ESA CCI {[}hlm. 7-9, Tabel
4-7{]}.

\subsection{Contributions}\label{contributions-6}

\begin{itemize}
\tightlist
\item
  \textbf{Kontribusi empiris yang diklaim penulis:} menyediakan bukti
  tentang bagaimana infrastruktur transportasi membentuk penggunaan
  lahan di dalam metropolitan besar negara berkembang, bukan hanya
  menjelaskan variasi antarkota {[}hlm. 2, Bagian 1{]}.
\item
  \textbf{Kontribusi konseptual-operasional:} memisahkan ekspansi lahan
  perkotaan dari \emph{sprawl} dan mengukur keduanya secara bersamaan,
  sehingga pertambahan luas tidak otomatis dianggap identik dengan
  pembangunan yang tersebar {[}hlm. 2 dan 4, Bagian 1 dan 3.1{]}.
\item
  \textbf{Kontribusi metodologis:} menggunakan variasi jalan dan rel
  historis sebagai calon instrumen untuk mengurangi bias kausalitas
  terbalik antara pembangunan jalan tol dan perkembangan perkotaan
  {[}hlm. 2 dan 5, Bagian 1 dan 4{]}.
\item
  \textbf{Kontribusi data:} memperlihatkan kegunaan data satelit GHSL
  beresolusi tinggi untuk meneliti perubahan penggunaan lahan
  intrametropolitan dan memungkinkan lebih dari satu indikator
  \emph{sprawl} {[}hlm. 2 dan 9-10, Bagian 1 dan 6{]}.
\item
  \textbf{Kontribusi kebijakan yang diklaim:} menambah pemahaman tentang
  konsekuensi tidak langsung pembangunan jalan tol terhadap bentuk kota.
  Artikel tidak mengestimasi secara langsung dampak lanjutan terhadap
  energi, produktivitas, layanan publik, mobilitas sosial, atau
  ketimpangan {[}hlm. 2, Bagian 1; hlm. 10, Bagian 7{]}.
\end{itemize}

\subsection{Research Gap}\label{research-gap-6}

\textbf{Kesenjangan yang dinyatakan sumber.} Bukti terdahulu telah
menjelaskan pengaruh transportasi terhadap suburbanisasi dan
\emph{sprawl} antarkota, khususnya di Eropa dan Amerika Serikat, tetapi
hanya sedikit yang diketahui tentang efek kausalnya di dalam
metropolitan besar negara berkembang. Kondisi metropolitan yang besar,
cepat tumbuh, polisentris, dan memiliki tata ruang kurang mengikat
membuat kesenjangan ini substantif bagi JMA {[}hlm. 1-2, Abstract dan
Bagian 1{]}.

Penulis juga mengidentifikasi kekurangan pengukuran: studi Jakarta
terdahulu menggunakan data relatif kasar dan cenderung mendeskripsikan
tren umum, sedangkan studi \emph{sprawl} sering hanya memakai satu
indikator. Artikel mengisi celah tersebut melalui data GHSL 38 m serta
pengukuran terpisah atas indeks \emph{sprawl} dan ekspansi lahan {[}hlm.
2, Bagian 1; hlm. 9-10, Bagian 6{]}.

\textbf{Kesenjangan yang tetap menurut sumber.} Tidak ada studi yang
secara langsung dapat dibandingkan dengan estimasi intrametropolitan
JMA, dan hubungan antara kebijakan biaya transportasi, perilaku
perjalanan, pilihan lokasi, serta respons kelompok sosial yang berbeda
masih memerlukan penelitian lebih komprehensif {[}hlm. 9-10, Bagian
6-7{]}.

\subsection{Limitations}\label{limitations-6}

Artikel tidak memiliki bagian khusus berjudul keterbatasan. Keterbatasan
dan kualifikasi berikut dinyatakan atau dapat ditelusuri langsung dari
pembahasan metode dan hasil.

\begin{itemize}
\tightlist
\item
  \textbf{Diakui sumber:} OLS rentan terhadap kausalitas terbalik, dan
  perbedaan tanda hasil ekspansi antara OLS dan TSLS menunjukkan
  pentingnya masalah tersebut {[}hlm. 5 dan 7-8, Bagian 4, 5.1, dan
  5.2.2{]}.
\item
  \textbf{Diakui sumber:} besaran koefisien TSLS relatif tinggi
  dibandingkan rerata variabel hasil sehingga harus ditafsirkan dengan
  sangat hati-hati {[}hlm. 8, Bagian 5.2.2{]}.
\item
  \textbf{Diakui sumber:} hasil indeks \emph{sprawl} kehilangan
  signifikansi ketika 5\% atau 10\% observasi terbawah dikeluarkan.
  Kombinasi JMA-Serang juga menghasilkan tanda berlawanan dengan
  instrumen yang lemah {[}hlm. 8-9, Bagian 5.2.3 dan Tabel 6{]}.
\item
  \textbf{Diakui sumber:} besaran efek sensitif terhadap resolusi
  tutupan lahan. ESA CCI 300 m menghasilkan dampak yang jauh lebih besar
  daripada GHSL 38 m, dan penulis menyatakan resolusi kasar dapat gagal
  menangkap ruang hijau, sungai, serta variasi lokal lain {[}hlm. 8-9,
  Bagian 5.2.3 dan Tabel 7{]}.
\item
  \textbf{Pilihan pembatasan model oleh sumber:} efek kausal jaringan
  kereta diabaikan dan hanya jarak ke stasiun yang dikontrol,
  berdasarkan kecilnya signifikansi dalam satu studi terdahulu yang
  dirujuk penulis {[}hlm. 5, Bagian 4{]}.
\item
  \textbf{Inferensi pengolahan:} validitas kausal tetap bergantung pada
  asumsi bahwa jalan historis hanya memengaruhi bentuk kota kontemporer
  melalui akses jalan tol. Uji fasilitas nontransportasi mendukung
  plausibilitas asumsi tersebut, tetapi tidak dapat memeriksa semua
  jalur historis alternatif {[}hlm. 5-6, Bagian 4 dan Tabel 1{]}.
\item
  \textbf{Inferensi pengolahan:} klasifikasi biner GHSL mengukur
  keberadaan lahan terbangun, bukan fungsi lahan, jenis hunian,
  kepadatan bangunan, atau intensitas aktivitas. Karena itu, mekanisme
  perilaku dan jenis pembangunan yang menghasilkan perubahan indeks
  tidak dapat dibedakan dari variabel hasil yang tersedia {[}hlm. 2 dan
  4, Bagian 1 dan 3.1.1{]}.
\item
  \textbf{Inferensi pengolahan:} generalisasi empiris dibatasi pada JMA,
  unit spasial dan rentang waktu yang dipakai, sedangkan hasil
  menunjukkan heterogenitas besar antarzona. Artikel tidak memberikan
  desain untuk menguji apakah besar efek yang sama berlaku pada
  metropolitan lain {[}hlm. 7-10, Bagian 5-7{]}.
\item
  \textbf{Ketidakjelasan dalam file:} GHSL tersedia untuk 1990, 2000,
  dan 2014 serta tabel menyebut perubahan 1990-2014, tetapi definisi
  perubahan jarak di Bagian 5.1 menyebut 2015 dikurangi 1990; Gambar 1
  juga memakai 1990-2015. File tidak menjelaskan rekonsiliasi tahun
  akhir tersebut {[}hlm. 3, Gambar 1; hlm. 4, Bagian 3.1.1; hlm. 7,
  Bagian 5.1; hlm. 9, Tabel 6-7{]}.
\item
  \textbf{Tidak disebutkan dalam file:} penanganan nilai nol sebelum
  transformasi log, bobot spasial atau koreksi dependensi spasial, serta
  kriteria rinci penghapusan observasi.
\end{itemize}

\subsection{Challenges}\label{challenges-6}

\textbf{Tantangan penelitian yang dinyatakan atau diperlihatkan sumber}
meliputi penempatan jalan tol yang tidak acak, kausalitas terbalik
antara infrastruktur dan pembangunan, pemenuhan syarat relevansi serta
eksogenitas instrumen, perbedaan luas unit administrasi yang memunculkan
\emph{modifiable areal unit problem}, potensi bias pengukuran jaringan,
dan sensitivitas pengukuran \emph{sprawl} terhadap resolusi citra
{[}hlm. 4-6, Bagian 3.1.1 dan 4; hlm. 8-9, Bagian 5.2.3{]}.

Penulis merespons tantangan tersebut dengan TSLS berbasis jaringan
historis, kontrol kondisi awal dan geografi, pemeriksaan silang data
transportasi, unit kisi seragam, statistik F Kleibergen-Paap, regresi
fasilitas nontransportasi, berbagai uji ketahanan, dan uji falsifikasi
{[}hlm. 4-10, Bagian 3.1.1-5.2.4{]}.

\textbf{Inferensi pengolahan.} Tantangan interpretasi utama adalah
memisahkan dua hasil yang berbeda: bukti kuat tentang ketersebaran
pembangunan tidak otomatis berarti pertambahan total luas perkotaan pada
semua bagian JMA. Tantangan kedua ialah membandingkan koefisien lintas
spesifikasi karena instrumen, cakupan wilayah, ukuran sel, resolusi
data, dan jumlah observasinya berubah.

\subsection{Practical Implications}\label{practical-implications-6}

\textbf{Implikasi yang dinyatakan penulis.} Pembuat kebijakan perlu
memasukkan dampak tidak langsung investasi jalan tol terhadap bentuk dan
penyebaran pembangunan ke dalam analisis biaya-manfaat, bukan hanya
menghitung manfaat mobilitas langsung. Alasannya, pembangunan tersebar
telah dikaitkan oleh literatur yang dikutip artikel dengan konsumsi
energi, produktivitas, persoalan sosial, dan kesetaraan, walaupun J07
tidak mengukur dampak lanjutan tersebut {[}hlm. 10, Bagian 7{]}.

Penulis menyatakan bahwa pengaruh kuat transportasi privat pada
metropolitan berorientasi mobil seperti JMA membuka kemungkinan
penggunaan pajak bahan bakar atau \emph{road pricing} untuk memengaruhi
pilihan perjalanan dan lokasi. Namun, artikel menegaskan perlunya
penelitian yang lebih komprehensif sebelum menyimpulkan cara kelompok
masyarakat yang berbeda akan merespons kebijakan tersebut {[}hlm. 10,
Bagian 7{]}.

\textbf{Inferensi pengolahan.} Karena efek \emph{sprawl} dilaporkan
lebih jelas pada radius 40 hingga 50 km dan ekspansi signifikan pada
\emph{city suburbs}, pemantauan dampak tata guna lahan dari proyek tol
dapat diprioritaskan pada koridor pinggiran. Ini adalah penerapan yang
diturunkan dari heterogenitas hasil, bukan rekomendasi zonasi spesifik
yang dinyatakan penulis.

\subsection{Applications}\label{applications-6}

\begin{itemize}
\tightlist
\item
  \textbf{Aplikasi yang ditunjukkan sumber:} GHSL dapat dipakai untuk
  memantau pertambahan dan ketersebaran lahan terbangun secara konsisten
  pada skala intrametropolitan; artikel menyatakan ketersediaan data
  spasial terbuka beresolusi rinci menjanjikan perluasan studi empiris
  di kota negara berkembang {[}hlm. 4 dan 9-10, Bagian 3.1.1 dan 6{]}.
\item
  \textbf{Aplikasi yang ditunjukkan sumber:} jarak ke jaringan
  transportasi historis dapat digunakan sebagai instrumen untuk
  mengevaluasi infrastruktur kontemporer ketika endogenitas penempatan
  menjadi masalah, dengan syarat relevansi dan eksogenitas diperiksa
  pada konteks masing-masing {[}hlm. 5-7, Bagian 4 dan Tabel 1-4{]}.
\item
  \textbf{Inferensi pengolahan:} kerangka dua indikator dapat digunakan
  dalam evaluasi rencana jalan tol agar perubahan jumlah lahan terbangun
  tidak keliru diperlakukan sebagai ukuran tunggal \emph{sprawl}.
\item
  \textbf{Inferensi pengolahan:} analisis radius dan subwilayah dapat
  menjadi dasar desain pemantauan spasial yang membedakan pusat,
  pinggiran berstatus kota, dan pinggiran lain. File tidak menyediakan
  perangkat lunak, protokol implementasi, ambang kebijakan, atau
  aplikasi operasional siap pakai.
\end{itemize}

\subsection{Future Research
(Optional)}\label{future-research-optional-6}

\textbf{Agenda eksplisit penulis.} Penelitian lebih komprehensif
diperlukan untuk menghubungkan kebijakan transportasi dengan perilaku
individu di JMA, khususnya untuk memahami bagaimana kelompok masyarakat
yang berbeda menyesuaikan perilaku perjalanan dan pilihan lokasi ketika
biaya transportasi berubah {[}hlm. 10, Bagian 7{]}.

\textbf{Inferensi pengolahan berdasarkan batas studi.} Penelitian
lanjutan dapat menguji kanal kausal kereta yang dalam artikel ini
diabaikan, menjelaskan perbedaan tahun akhir 2014 dan 2015, memeriksa
mengapa hasil sensitif terhadap observasi terbawah, serta membandingkan
estimasi pada data tutupan lahan beresolusi rinci yang dapat membedakan
tipe dan intensitas pembangunan. Agenda ini bukan usulan eksplisit
penulis, melainkan diturunkan dari pilihan model, inkonsistensi
temporal, dan hasil uji ketahanan dalam file.

\subsection{Provenance Notes}\label{provenance-notes-6}

\begin{itemize}
\tightlist
\item
  Review ini disusun hanya dari seluruh isi file
  \texttt{source/journal/J07-pratama-yudhistira-koomen-2022-highway-expansion-and-urban-sprawl-in-the-jakarta-metropolitan-area-10.1016-j.landusepol.2021.105856.txt},
  termasuk metadata PDF dan teks hasil ekstraksi sampai akhir referensi.
\item
  Metadata file menyatakan TXT diekstrak pada 31 Agustus 2026 dari PDF
  11 halaman menggunakan \texttt{pdftotext\ -layout}. PDF asli dan
  lampiran daring tidak dibuka; tidak ada sumber eksternal yang
  digunakan sebagai bukti substantif dalam review ini.
\item
  Locator halaman mengikuti pemisah dan nomor halaman yang tampak dalam
  hasil ekstraksi PDF; locator bagian, persamaan, tabel, dan gambar
  mengikuti label artikel.
\item
  Tata letak Tabel 6 pada halaman 9 terfragmentasi dalam ekstraksi teks.
  Review hanya memakai nilai yang masih dapat dihubungkan secara jelas
  dengan label tabel atau ditegaskan kembali dalam narasi penulis;
  susunan visual tabel asli tidak diasumsikan.
\item
  Artikel merujuk Appendix 1, Appendix 4, Appendix 5, dan
  \emph{supplementary data} daring, tetapi isinya tidak tercakup dalam
  TXT. Karena itu, statistik deskriptif lengkap, hasil \emph{sprawl} per
  subsampel pada Appendix 4, dan rincian \emph{robustness} ekspansi pada
  Appendix 5 tidak dapat diverifikasi.
\item
  Literatur yang disebut dalam artikel diperlakukan hanya sebagai
  laporan penulis J07. Isi dan mutu setiap rujukan tidak diverifikasi
  secara independen.
\item
  Istilah \textbf{Laporan sumber} menunjukkan deskripsi atau hasil yang
  dinyatakan artikel; \textbf{Interpretasi penulis} menunjukkan
  penjelasan mekanisme atau makna yang diberikan penulis;
  \textbf{Inferensi pengolahan} menunjukkan pembacaan analitis review
  ini dan bukan klaim langsung artikel.
\item
  Ketidaksesuaian penyebutan 1990-2014 dan 1990-2015 dipertahankan
  secara eksplisit. Tidak dilakukan koreksi berdasarkan dugaan atau
  informasi di luar TXT.
\end{itemize}

\section{Review J08}\label{review-j08}

\subsection{Identitas Sumber}\label{identitas-sumber-3}

\begin{itemize}
\tightlist
\item
  \textbf{Kode:} J08.
\item
  \textbf{Judul:} ``Examining Socio-Economic Inequality Among Commuters:
  The Case of the Jakarta Metropolitan Area.''
\item
  \textbf{Penulis:} Adiwan Aritenang.
\item
  \textbf{Afiliasi:} Urban and Regional Planning, Institut Teknologi
  Bandung, Indonesia.
\item
  \textbf{Jenis dokumen:} Artikel jurnal. Status telaah sejawat tidak
  dinyatakan secara eksplisit dalam file.
\item
  \textbf{Publikasi:} \emph{Urban Planning}, 2022, Volume 7, Issue 3,
  hlm. 172-184. ISSN 2183-7635.
\item
  \textbf{DOI:} 10.17645/up.v7i3.5271.
\item
  \textbf{Edisi tematik:} ``The Resilient Metropolis: Planning in an Era
  of Decentralization,'' disunting oleh Thomas J. Vicino.
\item
  \textbf{Lisensi:} Creative Commons Attribution 4.0 International (CC
  BY).
\item
  \textbf{Riwayat naskah:} Diserahkan 15 Januari 2022, diterima 28 Mei
  2022, dan diterbitkan 29 Juli 2022.
\item
  \textbf{Kata kunci:} commuters; health; inequality; Jakarta;
  metropolitan area; transportation.
\item
  \textbf{Pendanaan:} PPMI research funding programme, Institut
  Teknologi Bandung, 2021, dengan judul proyek ``New Town Development in
  Jakarta Metropolitan Area (JMA): Reflection of JMA Post-Suburban
  Phenomena'' (hlm. 182, bagian \emph{Acknowledgments}).
\item
  \textbf{Konflik kepentingan:} Penulis menyatakan tidak ada konflik
  kepentingan (hlm. 182, bagian \emph{Conflict of Interests}).
\end{itemize}

\subsection{Summarized Abstract}\label{summarized-abstract-7}

\textbf{Laporan sumber:} Pertumbuhan kawasan perkotaan di sekitar
Jakarta meningkatkan perjalanan komuter di dalam dan antara wilayah
inti-periferi, sementara variasi pekerjaan dan tingkat ekonomi
berpotensi memperlebar ketimpangan sosioekonomi. Artikel membandingkan
perilaku perjalanan, pola spasial, dan kesehatan komuter berpendapatan
di bawah dan di atas upah minimum regional (UMR) dengan menggunakan data
Survei Komuter BPS-Statistics Indonesia 2019. Sampelnya disebut sekitar
atau lebih dari 13.000 responden dari Provinsi Jakarta dan
kabupaten/kota tetangga. Abstrak menyebut statistik deskriptif
kuantitatif dan uji nonparametrik, sedangkan Bagian 3.1 menyebut
statistik ringkasan dan regresi logistik (hlm. 172, abstrak; hlm.
175-176, Bagian 3 dan 3.1).

\textbf{Hasil yang dilaporkan sumber:} Tingkat pendapatan berhubungan
signifikan dengan pilihan kendaraan pribadi, tetapi penulis menyatakan
tidak ada pengaruh pendapatan terhadap pilihan moda angkutan umum.
Penduduk periferi berpendapatan lebih tinggi cenderung bepergian ke
wilayah inti Jakarta, sedangkan komuter berpendapatan lebih rendah lebih
banyak bergerak antarkawasan periferi. Perjalanan berdurasi panjang dan
keberangkatan dini dikaitkan dengan dampak kesehatan fisik yang negatif,
terutama secara deskriptif pada kelompok berpendapatan rendah. Namun,
model kesehatan tidak menemukan pendapatan sebagai prediktor signifikan
setelah kovariat dimasukkan (hlm. 172, abstrak; hlm. 181, Tabel 9 dan
Bagian 5).

\textbf{Interpretasi penulis:} Penyediaan angkutan umum yang nyaman,
aman, andal, dan menjangkau hubungan antarperiferi diperlukan untuk
mengurangi waktu perjalanan, memperbaiki kesehatan komuter, serta
mendukung ketahanan kawasan metropolitan (hlm. 172, abstrak; hlm.
181-182, Bagian 5).

\subsection{Problem Statement}\label{problem-statement-7}

\textbf{Laporan sumber:} Jakarta Metropolitan Area (JMA) berkembang
sebagai kawasan inti-periferi yang melintasi batas administratif.
Desentralisasi sejak 2001, penyebaran industri manufaktur sejak akhir
1980-an, pembangunan perumahan sejak 1990-an, serta ekspansi
infrastruktur dan transportasi meningkatkan populasi dan arus komuter.
Pada saat yang sama, variasi ekonomi dan kapasitas pemerintahan lokal
menghasilkan kondisi sosioekonomi, spasial, transportasi, dan kesehatan
yang tidak merata (hlm. 172-173, Bagian 1).

\textbf{Masalah ilmiah eksplisit:} Walaupun terdapat variasi besar di
dalam JMA, pemahaman mengenai cara ketimpangan sosioekonomi menentukan
pilihan moda, perilaku perjalanan, dan kondisi kesehatan komuter masih
terbatas. Literatur Indonesia yang dirangkum artikel lebih banyak
membahas karakteristik sosioekonomi, pengalaman dan psikologi
perjalanan, pilihan moda, politik dan tata kelola, ekonomi perkotaan,
serta jejaring geografis secara terpisah (hlm. 172-173, Bagian 1).

\textbf{Pertanyaan penelitian:} Artikel bertanya sejauh mana ketimpangan
sosioekonomi terdapat di JMA dan bagaimana kesenjangan tersebut
menentukan perilaku perjalanan serta masalah kesehatan di kawasan
metropolitan (hlm. 173, Bagian 1).

\subsection{Objectives}\label{objectives-7}

\textbf{Tujuan yang dinyatakan atau langsung tercermin dalam sumber:}

\begin{itemize}
\tightlist
\item
  Menguji sejauh mana kota dan kawasan pinggiran yang tidak merata
  memperburuk ketimpangan sosioekonomi dan kesehatan di antara komuter
  JMA (hlm. 173, Bagian 1; hlm. 181, Bagian 5).
\item
  Membandingkan pilihan moda antara komuter berpendapatan di atas dan di
  bawah ambang UMR 4,5 juta IDR (hlm. 177, Bagian 4 dan Tabel 2).
\item
  Menggambarkan hubungan kelompok pendapatan dengan luas hunian, asal
  dan tujuan perjalanan, jarak, durasi, waktu berangkat, jenis
  pekerjaan, serta kondisi kesehatan (hlm. 175-179, Bagian 3-4 dan Tabel
  3-5).
\item
  Mengestimasi faktor sosioekonomi yang berkaitan dengan perjalanan
  antarperiferi, pilihan kendaraan pribadi, dan masalah kesehatan
  melalui regresi logistik (hlm. 176 dan 179-181, Bagian 3.1 serta Tabel
  6-9).
\item
  Menarik implikasi bagi penyediaan transportasi dan ketahanan
  metropolitan, khususnya untuk perjalanan dari dan di antara wilayah
  periferi (hlm. 181-182, Bagian 5).
\end{itemize}

\subsection{Literature Survey}\label{literature-survey-7}

\textbf{Cakupan literatur yang dilaporkan sumber:} Artikel menyusun
kajian pustaka dalam dua lapisan, yaitu ketimpangan sosioekonomi dan
kesehatan pada kawasan metropolitan secara umum serta ketimpangan
tersebut dalam konteks JMA. Uraian berikut adalah ringkasan atas cara
penulis menggunakan pustaka di dalam TXT, bukan verifikasi independen
terhadap sumber-sumber yang dikutip.

\begin{itemize}
\tightlist
\item
  Literatur perkembangan metropolitan menempatkan kawasan metropolitan
  sebagai wilayah perkotaan yang terhubung secara spasial tetapi
  terpisah secara administratif. Akumulasi modal, kemajuan transportasi
  dan komunikasi, serta konsentrasi aktivitas menjadikan metropolitan
  pusat pertumbuhan sekaligus menarik penduduk dan perjalanan (hlm. 173,
  Bagian 2.1).
\item
  Literatur tentang urbanisasi menunjukkan dua sisi pertumbuhan kota.
  Kota menciptakan peluang ekonomi dan dapat mendukung pengurangan
  kemiskinan, tetapi pertumbuhan kendaraan yang lebih cepat daripada
  kapasitas jalan meningkatkan kemacetan, polusi, biaya transportasi,
  dan disparitas pendapatan. Literatur yang dirujuk juga mengusulkan
  identifikasi distribusi pendapatan, konsesi bagi kelompok sasaran, dan
  penyesuaian struktur tarif untuk meningkatkan keterjangkauan
  transportasi (hlm. 174, Bagian 2.1).
\item
  Studi metropolitan yang dirangkum mengaitkan ketimpangan dengan
  mahalnya perumahan bagi kelompok berpendapatan rendah, bertambahnya
  jarak perjalanan ke subpusat pekerjaan, dan biaya sosial perjalanan
  jarak jauh. Artikel juga membedakan kecenderungan pekerja
  berpendidikan rendah atau informal untuk meminimalkan biaya perjalanan
  dengan pasar tenaga kerja terampil yang menjangkau wilayah lebih luas
  (hlm. 174, Bagian 2.1).
\item
  Perspektif polisentris digunakan untuk menjelaskan kemungkinan
  pengurangan jarak perjalanan apabila penduduk dapat tinggal dekat
  subpusat kerja. Pada saat yang sama, ketahanan metropolitan dinyatakan
  memerlukan strategi lintas bidang, termasuk kelembagaan, pendidikan,
  pengembangan tenaga kerja, transportasi, dan upah minimum.
  Ketidakselarasan antara batas pemerintahan dan wilayah ekonomi
  mendasari kebutuhan tata kelola bertingkat (hlm. 174, Bagian 2.1).
\item
  Untuk JMA, artikel mengikuti argumen tentang spasialitas
  neoliberalisme: transformasi lahan yang didominasi pengembang besar
  menghasilkan komunitas kelas menengah berpagar yang tersegregasi,
  tetapi tetap terhubung melalui mobilitas dan aktivitas lain. Literatur
  yang dirangkum juga memperlihatkan koeksistensi apartemen, kantor,
  pusat belanja, dan kampung kota di inti Jakarta serta klaster
  perumahan menengah-atas di periferi (hlm. 174-175, Bagian 2.2).
\item
  Studi terdahulu yang dikutip melaporkan pergeseran komposisi pekerjaan
  JMA, penurunan kemiskinan yang kecil dari 2004 sampai 2014,
  ketimpangan pengeluaran yang lebih tinggi di inti Jakarta, dan
  konsentrasi kelompok pekerjaan atas di South Tangerang City. Artikel
  menggunakan temuan-temuan ini untuk membangun dugaan bahwa ketimpangan
  pendapatan dan ruang tinggal berkaitan dengan pola komuter (hlm. 175,
  Bagian 2.2).
\item
  Literatur kesehatan komuter yang dirujuk menghubungkan moda, jarak,
  dan perjalanan berkepanjangan dengan stres, sakit kepala, sakit
  punggung, kesejahteraan, dan kualitas relasi. Variabel kesehatan fisik
  dan mental kemudian dimasukkan ke kerangka empiris artikel (hlm. 174
  dan 176, Bagian 2.1 dan 3).
\end{itemize}

\textbf{Posisi teoretis penulis:} Penulis berhipotesis secara umum bahwa
ketimpangan sosioekonomi dan pendapatan antara Jakarta dan wilayah
tetangganya dapat menjelaskan perilaku komuter, pilihan moda, dan pada
akhirnya kondisi kesehatan. Hipotesis formal per model, arah hubungan
setiap variabel, dan mekanisme kausal terperinci tidak disebutkan dalam
file (hlm. 173, Bagian 1).

\subsection{Method Used}\label{method-used-7}

\textbf{Desain dan tahap analisis yang dilaporkan sumber:}

\begin{itemize}
\tightlist
\item
  Artikel menggunakan pendekatan kuantitatif pada data Survei Komuter
  2019. \textbf{Inferensi pengolahan:} Karena data berasal dari satu
  tahun survei dan tidak ada pelacakan responden lintas waktu, desain
  empirisnya dapat dibaca sebagai observasional potong lintang; istilah
  desain ini tidak dinyatakan langsung oleh penulis (hlm. 175-176,
  Bagian 3-3.1).
\item
  Tahap pertama adalah statistik ringkasan untuk membandingkan kelompok
  pendapatan menurut moda, wilayah tempat tinggal, tujuan, karakteristik
  perjalanan, jenis pekerjaan, dan kondisi kesehatan (hlm. 176-179,
  Bagian 3.1 dan 4).
\item
  Tahap kedua adalah regresi logistik untuk keluaran biner. Artikel
  menjelaskan tujuan umum regresi sebagai pengujian peluang observasi
  masuk ke kategori tertentu, kemudian menampilkan model perjalanan
  antarperiferi, pilihan moda kendaraan pribadi, dan kesehatan (hlm. 176
  dan 179-181, Bagian 3.1 serta Tabel 6, 7, dan 9).
\item
  Tabel 8 menyajikan efek marginal atau probabilitas prediksi pilihan
  kendaraan pribadi untuk perjalanan periferi-periferi dan periferi-inti
  menurut waktu berangkat, jarak, durasi, dan status pekerjaan formal
  (hlm. 180, Tabel 8).
\item
  Kelayakan model dilaporkan melalui uji likelihood ratio, uji Wald, dan
  uji Hosmer-Lemeshow. Nilai p Hosmer-Lemeshow adalah 0,480 untuk model
  perjalanan antarperiferi, 0,317 untuk model pilihan moda, dan 0,757
  untuk model kesehatan (hlm. 179-181, Tabel 6, 7, dan 9).
\end{itemize}

\textbf{Ketidakselarasan pelaporan:} Abstrak menyebut ``a non-parametric
test'', tetapi Bagian 3.1 hanya menjelaskan statistik ringkasan dan
regresi logistik; nama, prosedur, statistik, dan hasil uji nonparametrik
tidak disebutkan dalam file (hlm. 172, abstrak; hlm. 176, Bagian 3.1).

\textbf{Perangkat lunak, penanganan data hilang, koreksi desain survei
dalam regresi, dan analisis sensitivitas:} Tidak disebutkan dalam file.

\subsection{Dataset}\label{dataset-7}

\begin{itemize}
\tightlist
\item
  \textbf{Dataset utama:} Data mentah Survei Komuter 2019 yang
  dikumpulkan BPS-Statistics Indonesia (hlm. 175, Bagian 3).
\item
  \textbf{Isi data:} Lokasi tempat tinggal, pendapatan, karakteristik
  pekerjaan, pendidikan, kondisi kesehatan, tujuan perjalanan harian,
  moda transportasi, jarak, durasi, dan waktu perjalanan. Jawaban
  individu ditautkan dengan data rumah tangga (hlm. 175-176, Bagian 3).
\item
  \textbf{Dasar distribusi sampel:} Ukuran rumah tangga dan populasi
  berdasarkan Sensus Penduduk Indonesia 2010, data antarsensus 2015, dan
  Survei Angkatan Kerja Nasional tahunan digunakan untuk menentukan
  distribusi sampel (hlm. 175, Bagian 3).
\item
  \textbf{Pembobotan:} Tabel 1 disebut mewakili lebih dari 2,43 juta
  komuter setelah sekitar 13.000 sampel diberi bobot berbasis distrik.
  Rumus bobot, strata, probabilitas pemilihan, dan penerapan bobot pada
  regresi tidak disebutkan dalam file (hlm. 176-177, Bagian 3 dan Tabel
  1).
\item
  \textbf{Data spasial tambahan:} Peta konektivitas jalan raya, jalan
  primer, rel, dan persebaran permukiman pada Gambar 1 diperoleh dari
  OpenStreetMap Indonesia 2021 (hlm. 175-176, Gambar 1). Peta ini
  mendukung deskripsi spasial, bukan dinyatakan sebagai masukan model
  regresi.
\item
  \textbf{Waktu pelaksanaan survei, versi data, akses data, dan nomor
  katalog dataset:} Tidak disebutkan dalam file.
\end{itemize}

\subsection{Population / Sample}\label{population-sample-7}

\begin{itemize}
\tightlist
\item
  \textbf{Populasi sasaran yang tersirat dalam sumber:} Penduduk yang
  melakukan perjalanan komuter dan bertempat tinggal di 13
  kabupaten/kota JMA. Definisi operasional ``komuter'', batas usia,
  frekuensi minimum perjalanan, serta kriteria inklusi dan eksklusi
  tidak disebutkan dalam file (hlm. 175, Bagian 3).
\item
  \textbf{Ukuran sampel:} Abstrak dan pendahuluan menyebut 13.000
  responden, sedangkan Bagian 3 menyebut lebih dari 13.000 responden.
  Jumlah tepat sampel analitis untuk setiap tabel dan model tidak
  disebutkan dalam file (hlm. 172-173, abstrak dan Bagian 1; hlm. 175,
  Bagian 3).
\item
  \textbf{Cakupan inti:} Central Jakarta, North Jakarta, South Jakarta,
  East Jakarta, dan West Jakarta di Provinsi Jakarta (hlm. 175, Bagian
  3).
\item
  \textbf{Cakupan periferi Jawa Barat:} Depok, Bogor, Bogor City,
  Bekasi, dan Bekasi City (hlm. 175, Bagian 3).
\item
  \textbf{Cakupan periferi Banten:} Tangerang, Tangerang City, dan South
  Tangerang (hlm. 175, Bagian 3).
\item
  \textbf{Strategi sampling:} Stratifikasi spasial dan pemilihan sampel
  dua tahap pada tingkat kecamatan. Distribusi rumah tangga dan penduduk
  diperhitungkan dengan basis data sensus dan survei nasional (hlm. 175,
  Bagian 3).
\item
  \textbf{Unit analisis:} \textbf{Inferensi pengolahan:} Unit observasi
  analitis tampaknya komuter individual karena kuesioner individu
  ditautkan ke data rumah tangga dan regresi memodelkan pilihan atau
  kondisi komuter. Penulis tidak memberikan pernyataan formal tentang
  unit analisis (hlm. 175, Bagian 3).
\item
  \textbf{Tingkat respons dan kehilangan observasi:} Tidak disebutkan
  dalam file.
\end{itemize}

\subsection{Dependent Variables}\label{dependent-variables-7}

\begin{itemize}
\tightlist
\item
  \textbf{Perjalanan antarperiferi:} Keluaran biner pada model Tabel 6
  adalah apakah perjalanan berlangsung antara distrik-distrik periferi.
  Prediktor yang ditampilkan ialah pendapatan, sektor ekonomi, pekerjaan
  formal, kendaraan pribadi, dan status menikah. Kategori pembanding
  sektor manufaktur dijelaskan dalam narasi, tetapi pengodean lengkap
  keluaran dan semua kategori referensi tidak disebutkan dalam file
  (hlm. 179-180, Tabel 6 dan pembahasannya).
\item
  \textbf{Pilihan moda kendaraan pribadi:} Model Tabel 7 memperkirakan
  peluang menggunakan kendaraan pribadi. Prediktornya meliputi log
  pendapatan, durasi, jarak, pekerjaan formal, tipe asal-tujuan
  perjalanan, dan status menikah. Definisi moda yang masuk kategori
  ``private vehicle'', kategori referensi tipe perjalanan, satuan atau
  pengodean jarak dan durasi, serta nilai keluaran 0 dan 1 tidak
  disebutkan dalam file (hlm. 180, Tabel 7).
\item
  \textbf{Status kesehatan komuter:} Model Tabel 9 memeriksa adanya
  masalah kesehatan fisik atau terkait stres. Prediktornya ialah log
  pendapatan, tipe perjalanan, durasi, status menikah, dan moda
  transportasi. Instrumen, periode rujukan, ambang kasus, pemisahan
  kesehatan fisik dan mental, serta pengodean variabel kesehatan tidak
  disebutkan dalam file (hlm. 181, Tabel 9 dan Bagian 5).
\item
  \textbf{Keluaran deskriptif tambahan:} Moda, luas lantai rumah per
  kapita, tujuan, jarak, durasi, waktu berangkat, jenis pekerjaan, dan
  persentase masalah kesehatan dibandingkan menurut kelompok pendapatan.
  Luas hunian dibagi menjadi kurang dari 8 m2 per kapita, 8-12,5 m2, dan
  lebih dari 12,5 m2 (hlm. 176-179, Tabel 2-5).
\end{itemize}

\subsection{Results}\label{results-7}

\textbf{Hasil deskriptif yang dilaporkan sumber:}

\begin{itemize}
\tightlist
\item
  Penulis menetapkan 4,5 juta IDR sebagai ambang UMR untuk membedakan
  pendapatan di atas dan di bawah UMR. Narasi menyatakan lebih dari 78\%
  komuter memiliki pendapatan di atas ambang yang disebut berada dalam
  rentang UMR JMA 3,4-4,5 juta IDR (hlm. 176-177, Tabel 1 dan awal
  Bagian 4).
\item
  Sepeda motor mendominasi kedua kelompok, tetapi bagi kelompok di bawah
  UMR porsinya 73,49\%, dibandingkan 56,09\% pada kelompok di atas UMR.
  Mobil pribadi digunakan oleh 17,81\% kelompok di atas UMR dan 3,61\%
  kelompok di bawah UMR. Proporsi kereta komuter adalah 10,13\% dan
  8,51\%, sedangkan bus publik 6,57\% dan 6,70\% untuk kelompok di atas
  dan di bawah UMR (hlm. 177, Tabel 2).
\item
  Penulis menyatakan moda berikutnya bagi kelompok berpendapatan rendah
  ialah kereta komuter, bus, dan TransJakarta; bagi kelompok
  berpendapatan tinggi ialah mobil pribadi, kereta komuter, dan bus.
  Penulis menafsirkan penggunaan mobil dan layanan pemesanan kendaraan
  yang lebih tinggi sebagai indikasi anggaran diskresioner serta
  permintaan atas perjalanan yang lebih hemat waktu (hlm. 177,
  pembahasan Tabel 2).
\item
  Berdasarkan Tabel 3, penulis menyatakan komuter dari Jakarta
  terkonsentrasi pada hunian kecil dan menghubungkannya dengan kampung
  kota. South Tangerang City dan Bogor Regency disebut memiliki bagian
  besar penduduk pada hunian luas, sedangkan Bekasi City dan Depok City
  disebut termasuk asal komuter berpendapatan tinggi (hlm. 177-178,
  Tabel 3 dan pembahasannya).
\item
  Lebih dari 63\% komuter pada kedua kelompok pendapatan bepergian
  menuju wilayah inti. Hampir satu dari tiga bepergian ke wilayah
  periferi, yang ditafsirkan penulis sebagai kemungkinan tanda
  polisentrisitas. Di dalam Jakarta, persentase perjalanan kelompok
  berpendapatan rendah disebut lebih dari 44\%, sedangkan kelompok
  berpendapatan tinggi sekitar 37\%. Sebaliknya, kelompok berpendapatan
  tinggi dari periferi lebih banyak menuju inti, sementara Bogor dan
  Tangerang regencies menjadi kontributor besar perjalanan antarperiferi
  (hlm. 178, Tabel 4 dan pembahasannya).
\item
  Mayoritas komuter menempuh kurang dari 20 km, kurang dari satu jam,
  dan berangkat sebelum pukul 07.00. Untuk kelompok di atas dan di bawah
  UMR, perjalanan kurang dari 20 km masing-masing 55,21\% dan 59,19\%;
  perjalanan lebih dari 40 km 7,94\% dan 7,40\%; durasi lebih dari satu
  jam 27,56\% dan 23,29\%; serta keberangkatan sebelum pukul 07.00
  54,66\% dan 50,50\% (hlm. 178-179, Tabel 5).
\item
  Kelompok berpendapatan rendah dilaporkan lebih banyak bekerja sebagai
  buruh dan staf kantor. Pekerja lepas/informal dan wiraswasta
  berpendapatan tinggi berjumlah 23,50\% dari kelompok tersebut; penulis
  menafsirkannya sebagai kemungkinan cerminan pertumbuhan kerja gig dan
  kewirausahaan di JMA (hlm. 179, Bagian 4).
\item
  Masalah kesehatan fisik dan mental sehari-hari lebih banyak dilaporkan
  pada kelompok berpendapatan rendah, masing-masing 27,13\% dan 24,18\%.
  Sakit kepala dan nyeri lebih tinggi pada kelompok tersebut baik di
  antara pengguna kendaraan pribadi maupun angkutan umum. Variasi
  kesehatan mental antarkelompok disebut relatif kecil, sekitar 32-42\%
  di antara komuter, tetapi denominator dan rincian kategorinya tidak
  ditampilkan (hlm. 179, Bagian 4).
\end{itemize}

\textbf{Hasil regresi yang dilaporkan sumber:}

\begin{itemize}
\tightlist
\item
  Pada model perjalanan antarperiferi, pendapatan memiliki koefisien
  -0,625 dan odds ratio 0,535 dengan p kurang dari 0,001. Sektor jasa
  memiliki odds ratio 0,529 dengan p kurang dari 0,001 dibandingkan
  sektor manufaktur. Kendaraan pribadi memiliki odds ratio 1,372 dengan
  p kurang dari 0,01 dan status menikah 1,251 dengan p kurang dari 0,05.
  Sektor primer dan pekerjaan formal tidak signifikan menurut tanda
  signifikansi tabel. Uji likelihood ratio dan Wald ditampilkan dengan p
  0,000, sedangkan Hosmer-Lemeshow p 0,480 (hlm. 179, Tabel 6).
\item
  \textbf{Interpretasi penulis atas Tabel 6:} Komuter antarperiferi
  cenderung berpendapatan lebih rendah, bekerja di manufaktur,
  menggunakan kendaraan pribadi, dan menikah. Penulis menyebut kenaikan
  pendapatan 1\% mengurangi ``probability'' sebesar 0,535.
  \textbf{Catatan inferensi pengolahan:} Angka 0,535 pada tabel adalah
  odds ratio, bukan perubahan probabilitas langsung; transformasi
  pendapatan dan pengodean outcome tidak cukup dijelaskan untuk
  memverifikasi interpretasi 1\% tersebut (hlm. 179-180, Tabel 6 dan
  pembahasannya).
\item
  Pada model pilihan kendaraan pribadi, log pendapatan memiliki odds
  ratio 1,367 dengan p kurang dari 0,001; durasi 2,658 dengan p kurang
  dari 0,001; jarak 1,918 dengan p kurang dari 0,001; pekerjaan formal
  2,430 dengan p kurang dari 0,05; dan status menikah 1,674 dengan p
  kurang dari 0,001. Kategori perjalanan yang ditampilkan juga
  signifikan: antarperiferi 1,423, inti-periferi 1,598, dan antarkawasan
  inti 1,365. Uji likelihood ratio dan Wald ditampilkan dengan p 0,000;
  Hosmer-Lemeshow p 0,317 (hlm. 180, Tabel 7).
\item
  Penulis menyatakan komuter berdurasi dan berjarak lebih pendek serta
  pekerja formal memiliki peluang lebih tinggi menggunakan kendaraan
  pribadi. \textbf{Catatan inferensi pengolahan:} Arah ini tidak dapat
  direkonstruksi hanya dari tanda koefisien positif karena pengodean
  kategori durasi dan jarak tidak dijelaskan (hlm. 180, pembahasan Tabel
  7).
\item
  Efek marginal Tabel 8 menunjukkan probabilitas kendaraan pribadi yang
  selalu lebih tinggi untuk perjalanan periferi-periferi daripada
  periferi-inti: waktu keberangkatan 0,830 dibandingkan 0,744; jarak
  0,864 dibandingkan 0,787; durasi 0,892 dibandingkan 0,853; dan pekerja
  formal 0,831 dibandingkan 0,745. Semua diberi tanda p kurang dari
  0,001 (hlm. 180, Tabel 8).
\item
  Narasi marginal tambahan menyebut pekerja formal berpendapatan tinggi
  memiliki probabilitas menggunakan kendaraan pribadi sebesar 79,9\%
  pada perjalanan periferi-periferi dan 78,8\% pada perjalanan
  periferi-inti. Untuk pekerja informal berpendapatan rendah,
  probabilitasnya 61,7\% dan 60,1\% (hlm. 180-181, pembahasan Tabel 8).
\item
  Pada model kesehatan, perjalanan inti-periferi memiliki odds ratio
  1,474 dengan p kurang dari 0,01 dan status menikah 1,184 dengan p
  kurang dari 0,05. Durasi memiliki odds ratio 0,614 dengan p kurang
  dari 0,001. Pendapatan 0,999, perjalanan antarperiferi 1,155,
  perjalanan antarkawasan inti 1,076, dan moda 1,043 tidak signifikan
  menurut tabel. Uji likelihood ratio dan Wald ditampilkan dengan p
  0,000; Hosmer-Lemeshow p 0,757 (hlm. 181, Tabel 9).
\item
  \textbf{Interpretasi penulis atas Tabel 9:} Asal-tujuan dan durasi
  perjalanan memengaruhi kesehatan fisik, tetapi tidak ada bukti bahwa
  status sosioekonomi menjelaskan kondisi kesehatan setelah faktor lain
  diperhitungkan. Penggunaan kendaraan pribadi mungkin menurunkan risiko
  secara tidak langsung melalui waktu tunggu dan durasi yang lebih
  pendek; hubungan tidak langsung ini merupakan tafsir penulis dan tidak
  diuji melalui model mediasi yang dilaporkan (hlm. 181, pembahasan
  Tabel 9 dan Bagian 5).
\end{itemize}

\subsection{Findings}\label{findings-7}

\begin{itemize}
\tightlist
\item
  \textbf{Temuan sumber tentang moda:} Perbedaan pendapatan paling jelas
  muncul pada komposisi kendaraan pribadi. Kelompok di bawah UMR sangat
  didominasi sepeda motor, sedangkan kelompok di atas UMR lebih beragam
  dan jauh lebih banyak menggunakan mobil pribadi serta layanan
  pemesanan kendaraan. Penulis menyimpulkan pendapatan berpengaruh
  signifikan terhadap pilihan moda pribadi (hlm. 177 dan 180-181, Tabel
  2, Tabel 7, dan Bagian 5).
\item
  \textbf{Temuan sumber tentang angkutan umum:} Penulis menyatakan
  pendapatan tidak berpengaruh signifikan terhadap pilihan angkutan
  umum; proporsi bus dan kereta relatif dekat antarkelompok. Namun, file
  tidak menampilkan model terpisah untuk pilihan setiap moda angkutan
  umum, sehingga dasar inferensial klaim ini tidak dapat ditelusuri
  sepenuhnya dari tabel regresi yang tersedia (hlm. 177, Tabel 2; hlm.
  181, Bagian 5).
\item
  \textbf{Temuan sumber tentang struktur spasial:} Komuter periferi
  berpendapatan tinggi lebih cenderung menuju inti Jakarta, sedangkan
  kelompok berpendapatan rendah memiliki bagian perjalanan antarperiferi
  yang lebih besar. Penulis menghubungkan hampir sepertiga perjalanan
  menuju periferi dengan kemungkinan struktur JMA yang polisentris (hlm.
  178 dan 181, Tabel 4 dan Bagian 5).
\item
  \textbf{Temuan sumber tentang karakteristik perjalanan:} Distribusi
  jarak, durasi, dan waktu berangkat dinilai relatif serupa
  antarkelompok pendapatan. Kelompok berpendapatan tinggi sedikit lebih
  sering berangkat sebelum pukul 07.00 dan memiliki porsi durasi lebih
  dari satu jam yang lebih tinggi (hlm. 178-179, Tabel 5; hlm. 181,
  Bagian 5).
\item
  \textbf{Temuan sumber tentang kesehatan:} Secara deskriptif, kelompok
  berpendapatan rendah melaporkan masalah fisik dan mental lebih besar.
  Akan tetapi, regresi kesehatan tidak menemukan pendapatan atau moda
  sebagai prediktor signifikan; yang signifikan ialah tipe perjalanan
  inti-periferi, durasi, dan status menikah. Dengan demikian, bukti
  artikel mendukung perbedaan kesehatan deskriptif menurut pendapatan,
  tetapi tidak mendukung pengaruh pendapatan yang independen dalam model
  tersesuaikan (hlm. 179 dan 181, Tabel 9 serta Bagian 5).
\item
  \textbf{Interpretasi penulis:} Perjalanan antarperiferi yang panjang
  oleh kelompok berpendapatan rendah dan keterbatasan alternatif moda
  menjadikan perluasan angkutan umum antarperiferi sebagai isu ketahanan
  metropolitan, bukan hanya efisiensi transportasi (hlm. 181-182, Bagian
  5).
\end{itemize}

\subsection{Conclusions}\label{conclusions-7}

\textbf{Kesimpulan sumber:} Artikel menyimpulkan bahwa komuter JMA
memperlihatkan keragaman sosioekonomi dan kesehatan yang besar.
Pendapatan membedakan pilihan kendaraan pribadi: sepeda motor dominan
pada kelompok berpendapatan rendah, sedangkan moda kelompok
berpendapatan tinggi lebih beragam. Pendapatan tidak dinyatakan
berpengaruh terhadap pilihan angkutan umum. Secara spasial, kelompok
berpendapatan tinggi dari periferi lebih banyak menuju Jakarta,
sementara kelompok berpendapatan rendah lebih banyak bergerak
antarkawasan periferi (hlm. 181, Bagian 5).

\textbf{Kesimpulan kesehatan sumber:} Durasi dan asal-tujuan perjalanan
berkaitan dengan masalah kesehatan, sedangkan model tidak memberikan
bukti bahwa status sosioekonomi menjelaskan kesehatan komuter secara
independen. Pernyataan abstrak tentang dampak kesehatan yang terutama
dialami kelompok berpendapatan rendah sebaiknya dibaca sebagai gabungan
hasil deskriptif dan interpretasi, bukan sebagai efek pendapatan yang
dikonfirmasi model kesehatan (hlm. 172, abstrak; hlm. 181, Tabel 9 dan
Bagian 5).

\textbf{Kesimpulan kebijakan sumber:} Mengembangkan angkutan umum
antarperiferi, mengurangi waktu perjalanan, dan mendorong perjalanan
aktif dinilai dapat memperbaiki kesehatan serta ketahanan JMA.
Pelaksanaannya memerlukan peran pemerintah lokal dan tata kelola
bertingkat karena jaringan metropolitan melintasi batas administratif
(hlm. 181-182, Bagian 5).

\subsection{Contributions}\label{contributions-7}

\begin{itemize}
\tightlist
\item
  \textbf{Kontribusi empiris yang diklaim penulis:} Artikel
  menggabungkan dimensi pendapatan, moda, ruang tinggal, pola
  asal-tujuan, karakteristik perjalanan, pekerjaan, dan kesehatan pada
  satu kajian komuter JMA berbasis Survei Komuter BPS 2019 (hlm.
  175-181, Bagian 3-5).
\item
  \textbf{Kontribusi terhadap literatur:} Kajian memperluas literatur
  ketimpangan metropolitan dengan menunjukkan bahwa posisi inti-periferi
  dan arah perjalanan perlu dipertimbangkan bersama pendapatan ketika
  menilai pilihan moda dan kesehatan komuter (hlm. 173 dan 181, Bagian 1
  dan 5).
\item
  \textbf{Inferensi pengolahan mengenai kontribusi analitis:} Statistik
  deskriptif kelompok pendapatan dilengkapi dengan model logistik untuk
  perjalanan antarperiferi, kendaraan pribadi, dan kesehatan serta
  perbandingan efek marginal periferi-periferi dan periferi-inti.
  Artikel tidak mengklaim regresi logistik itu sendiri sebagai inovasi
  metodologis (hlm. 176 dan 179-181, Bagian 3.1 serta Tabel 6-9).
\item
  \textbf{Kontribusi kebijakan:} Artikel menggeser perhatian dari
  koneksi periferi-inti saja menuju kebutuhan angkutan umum
  antarperiferi, khususnya bagi kelompok berpendapatan rendah yang
  banyak melakukan perjalanan tersebut (hlm. 181-182, Bagian 5).
\end{itemize}

\subsection{Research Gap}\label{research-gap-7}

\textbf{Celah yang dinyatakan sumber:} Studi komuter Indonesia telah
membahas karakteristik sosioekonomi, pengalaman perjalanan, psikologi,
dan pilihan moda, sedangkan studi spasial JMA banyak membahas tata
kelola, kondisi sosioekonomi, ekonomi perkotaan, serta jejaring
geografis. Menurut penulis, masih terbatas pemahaman terpadu mengenai
bagaimana ketimpangan sosioekonomi menentukan pilihan moda, perilaku
perjalanan, pola spasial, dan kesehatan komuter dalam satu kawasan
metropolitan inti-periferi (hlm. 172-173, Bagian 1).

\textbf{Celah residual berdasarkan inferensi pengolahan:} Artikel belum
memisahkan secara terperinci mekanisme antara pendapatan, keterjangkauan
moda, waktu tunggu, durasi, dan kesehatan. Analisis juga belum
menunjukkan apakah hubungan berbeda menurut kabupaten/kota, jenis
kelamin, umur, pendidikan, atau moda angkutan umum tertentu, meskipun
sebagian karakteristik tersebut mungkin tersedia dalam survei. Ini
adalah identifikasi celah dari struktur analisis yang dilaporkan, bukan
pernyataan eksplisit penulis.

\subsection{Limitations}\label{limitations-7}

\textbf{Keterbatasan yang dinyatakan secara eksplisit oleh penulis:}
Tidak disebutkan dalam file.

\textbf{Keterbatasan yang teridentifikasi melalui inferensi pengolahan
atas TXT:}

\begin{itemize}
\tightlist
\item
  Data satu tahun mendukung pemetaan asosiasi, tetapi urutan waktu dan
  hubungan sebab-akibat antara pendapatan, moda, pola perjalanan, dan
  kesehatan tidak dapat dipastikan dari desain yang dilaporkan.
\item
  Jumlah sampel analitis per model, kehilangan data, tingkat respons,
  serta komposisi sampel tanpa bobot tidak diberikan. Karena itu,
  perubahan jumlah observasi dan potensi bias nonrespons tidak dapat
  dinilai.
\item
  Artikel menyebut sampel berbobot mewakili lebih dari 2,43 juta
  komuter, tetapi tidak menjelaskan konstruksi bobot maupun apakah
  stratifikasi, klaster, dan bobot survei diterapkan pada regresi dan
  galat baku (hlm. 176-177, Tabel 1).
\item
  Definisi dan pengodean variabel dependen, kategori referensi, satuan
  prediktor, bentuk variabel durasi dan jarak, serta instrumen kesehatan
  tidak disajikan. Keterbatasan ini menghambat replikasi dan membuat
  arah sejumlah koefisien sulit ditafsirkan (hlm. 176 dan 179-181,
  Bagian 3.1 serta Tabel 6-9).
\item
  Abstrak menyebut uji nonparametrik, tetapi metode dan hasilnya tidak
  muncul dalam Bagian 3.1 atau tabel hasil. Sebaliknya, regresi logistik
  yang menjadi analisis utama tidak disebutkan dalam abstrak (hlm. 172
  dan 176).
\item
  Klaim tidak adanya pengaruh pendapatan pada pilihan angkutan umum
  tidak disertai model khusus moda umum yang dapat dikenali dalam file;
  model Tabel 7 berfokus pada kendaraan pribadi (hlm. 180-181, Tabel 7
  dan Bagian 5).
\item
  Ambang tunggal 4,5 juta IDR digunakan untuk seluruh JMA walaupun
  narasi menyebut rentang UMR regional 3,4-4,5 juta IDR. Dampak pilihan
  ambang terhadap klasifikasi pendapatan tidak diuji (hlm. 176-177,
  Tabel 1-2).
\item
  Label kategori Tabel 5, yaitu kurang dari 30 menit dan kurang dari
  satu jam serta kurang dari 20 km dan kurang dari 40 km, tampak tumpang
  tindih dalam hasil ekstraksi. File tidak menjelaskan apakah kategori
  kedua sebenarnya berarti 30-60 menit dan 20-40 km (hlm. 179, Tabel 5).
\item
  Denominator beberapa persentase tidak selalu jelas. Contohnya, Tabel 3
  tampak menyajikan distribusi kolom, sementara narasi membacanya
  sebagai karakteristik komuter dari suatu wilayah. Tanpa jumlah dasar
  per wilayah, proporsi kondisional tidak dapat diverifikasi (hlm. 178,
  Tabel 3).
\item
  Tidak ada interval kepercayaan, pemeriksaan multikolinearitas,
  diagnostik pengaruh, perbandingan spesifikasi, atau analisis
  sensitivitas yang dilaporkan. Tidak disebutkan dalam file.
\end{itemize}

\subsection{Challenges}\label{challenges-7}

\textbf{Tantangan substantif yang dilaporkan sumber:}

\begin{itemize}
\tightlist
\item
  JMA mencakup banyak pemerintah daerah dengan variasi ekonomi,
  kapasitas, dan kewenangan. Jaringan ekonomi serta perjalanan melampaui
  batas politik, sehingga penyediaan dan pemeliharaan infrastruktur
  metropolitan sulit ditangani oleh satu tingkat pemerintahan (hlm.
  173-175 dan 182, Bagian 1, 2.1, 2.2, dan 5).
\item
  Pertumbuhan kendaraan, keterbatasan ruang jalan, kemacetan, polusi,
  biaya perjalanan, dan waktu tunggu angkutan umum saling berhubungan.
  Kelompok berpendapatan rendah menghadapi keterbatasan pilihan meskipun
  perjalanan hariannya menopang aktivitas metropolitan (hlm. 174 dan
  179-182, Bagian 2.1, 4, dan 5).
\item
  Penulis menempatkan perluasan layanan antarperiferi sebagai kebutuhan
  penting di samping koneksi periferi-inti. Kelompok berpendapatan
  rendah dinyatakan banyak melakukan perjalanan antarperiferi dan
  mengalami perjalanan panjang, sehingga desain jaringan perlu merespons
  pola polisentris (hlm. 178-182, Tabel 4, Tabel 6, dan Bagian 5).
\item
  Segregasi spasial antara kampung kota, pembangunan vertikal mewah, dan
  perumahan berpagar di periferi menyulitkan pemerataan akses tempat
  tinggal, pekerjaan, serta transportasi (hlm. 174-175, Bagian 2.2).
\end{itemize}

\textbf{Tantangan analitis menurut inferensi pengolahan:} Memisahkan
efek pendapatan dari asal-tujuan, jarak, durasi, pekerjaan, kepemilikan
atau akses kendaraan, dan kesehatan memerlukan definisi variabel serta
strategi identifikasi yang lebih lengkap daripada yang dilaporkan.
Perbedaan antara beban kesehatan deskriptif dan hasil regresi
tersesuaikan juga harus dipertahankan agar ketimpangan kelompok tidak
disalahartikan sebagai efek kausal pendapatan.

\subsection{Practical Implications}\label{practical-implications-7}

\begin{itemize}
\tightlist
\item
  \textbf{Rekomendasi penulis:} Pemerintah lokal dan perencana kota
  perlu memperluas angkutan umum antarkawasan periferi, bukan hanya
  koridor menuju inti Jakarta. Layanan harus nyaman, aman, andal, dan
  mampu mengurangi waktu perjalanan (hlm. 172 dan 181-182, abstrak serta
  Bagian 5).
\item
  \textbf{Rekomendasi penulis:} Pengurangan durasi dan waktu tunggu
  diposisikan sebagai intervensi transportasi sekaligus kesehatan,
  terutama bagi komuter berpendapatan rendah yang menghadapi perjalanan
  panjang (hlm. 181-182, Bagian 5).
\item
  \textbf{Rekomendasi penulis:} Perjalanan aktif dapat dipromosikan
  sebagai alternatif yang berpotensi meningkatkan kesehatan dan
  kesejahteraan, dengan dukungan literatur yang dikutip artikel (hlm.
  181-182, Bagian 5).
\item
  \textbf{Rekomendasi penulis:} Tata kelola bertingkat dibutuhkan untuk
  membangun, mengoperasikan, dan memelihara infrastruktur yang melintasi
  batas pemerintah daerah (hlm. 182, Bagian 5).
\item
  \textbf{Kualifikasi inferensial:} Karena artikel tidak mengevaluasi
  intervensi, rekomendasi tersebut merupakan implikasi kebijakan dari
  pola asosiasi, bukan dampak program yang telah diuji.
\end{itemize}

\subsection{Applications}\label{applications-7}

\textbf{Aplikasi yang tersirat oleh sumber dan dirumuskan sebagai
inferensi pengolahan:}

\begin{itemize}
\tightlist
\item
  Hasil asal-tujuan dapat digunakan sebagai dasar diagnosis awal untuk
  memprioritaskan koneksi angkutan umum periferi-periferi, khususnya
  pada arus dari Bogor dan Tangerang regencies yang disebut besar dalam
  kelompok perjalanan tersebut (hlm. 178, Tabel 4).
\item
  Profil moda menurut kelompok pendapatan dapat membantu penilaian
  keterjangkauan dan sasaran layanan, tetapi artikel tidak menguji skema
  tarif, subsidi, ataupun perubahan layanan tertentu (hlm. 177, Tabel
  2).
\item
  Indikator jarak, durasi, waktu berangkat, dan kesehatan dapat dipakai
  bersama untuk menyaring koridor dengan beban perjalanan tinggi. Ini
  merupakan kemungkinan penggunaan hasil, bukan aplikasi yang diuji
  dalam penelitian (hlm. 179 dan 181, Tabel 5 dan 9).
\item
  Kerangka inti-periferi dapat dipakai untuk membedakan kebutuhan
  perjalanan periferi-inti dari periferi-periferi dalam perencanaan
  metropolitan. Validitas penerapannya di kawasan metropolitan lain
  tidak diuji dalam file.
\end{itemize}

\subsection{Future Research
(Optional)}\label{future-research-optional-7}

\textbf{Agenda penelitian lanjutan yang dinyatakan eksplisit oleh
penulis:} Tidak disebutkan dalam file.

\textbf{Inferensi pengolahan, bukan usulan eksplisit penulis:}
Penelitian berikutnya dapat menguji hubungan perjalanan dan kesehatan
secara longitudinal; menjelaskan pengodean serta mengukur kesehatan
fisik dan mental secara terpisah; menggunakan UMR yang sesuai dengan
wilayah asal; mengestimasi pilihan setiap moda angkutan umum; dan
menguji heterogenitas antarwilayah serta kelompok demografis. Usulan ini
diturunkan dari variabel dan keterbatasan pelaporan dalam artikel, bukan
dari sumber eksternal.

\subsection{Provenance Notes}\label{provenance-notes-7}

\begin{itemize}
\tightlist
\item
  Review ini hanya menggunakan file TXT
  \texttt{source/journal/J08-aritenang-2022-examining-socio-economic-inequality-among-commuters-the-case-of-the-jakarta-metropolitan-area-10.17645-up.v7i3.5271.txt}.
  Tidak ada sumber eksternal, Zotero, PDF asli, atau artikel lain yang
  digunakan untuk menambah atau memverifikasi klaim.
\item
  Seluruh 759 baris TXT dibaca, termasuk metadata PDF, teks utama,
  sembilan tabel, satu gambar, ucapan terima kasih, pernyataan konflik
  kepentingan, daftar pustaka, dan profil penulis.
\item
  Metadata ekstraksi menyatakan PDF sumber berjumlah 13 halaman dan
  diekstraksi pada 31 Agustus 2026 dengan \texttt{pdftotext\ -layout}.
  Akses review ini karena itu adalah teks penuh hasil ekstraksi, bukan
  pemeriksaan visual PDF asli.
\item
  Locator halaman merujuk pada nomor halaman tercetak artikel, yaitu
  172-184, bukan nomor baris TXT. Locator bagian, tabel, dan gambar
  ditambahkan untuk klaim penting apabila tersedia.
\item
  Pemenggalan kata, urutan kolom, simbol, dan posisi teks dapat
  terpengaruh ekstraksi dua kolom. Angka dilaporkan mengikuti TXT,
  tetapi ketidakjelasan kategori atau denominator tidak diperbaiki
  dengan dugaan.
\item
  Istilah \textbf{laporan sumber} menandai apa yang dinyatakan atau
  ditampilkan artikel; \textbf{interpretasi penulis} menandai penjelasan
  atau implikasi yang diajukan penulis; \textbf{inferensi pengolahan}
  menandai pembacaan kritis yang diturunkan hanya dari isi TXT dan bukan
  klaim penulis.
\item
  Batas provenance utama ialah tidak adanya verifikasi terhadap PDF
  visual, data mentah BPS, kuesioner, kode analisis, dan pustaka yang
  dikutip. Akibatnya, review tidak mengklaim dapat mengaudit akurasi
  angka, konstruksi variabel, representativitas, atau reproduksibilitas
  model di luar informasi yang tersedia dalam TXT.
\end{itemize}

\section{Review J09}\label{review-j09}

\textbf{Identitas sumber}

\begin{itemize}
\tightlist
\item
  \textbf{Kode:} J09
\item
  \textbf{Judul:} \emph{The crucial role of motorcycle-based
  ride-hailing among commuters: The case of Jakarta and Bandung
  metropolitan areas}
\item
  \textbf{Penulis:} Adiwan Fahlan Aritenang
\item
  \textbf{Afiliasi:} Urban and Regional Planning Program, Transportation
  Program, School of Architecture, Planning and Policy Development,
  Institut Teknologi Bandung, Indonesia
\item
  \textbf{Publikasi:} \emph{Journal of Public Transportation}, volume
  26, artikel 100082, 2024
\item
  \textbf{DOI:} 10.1016/j.jpubtr.2024.100082
\item
  \textbf{Kata kunci sumber:} motorcycle-based ride-hailing;
  metropolitan; commuters; sharing economy; Indonesia
\item
  \textbf{Riwayat naskah:} diterima 29 September 2022, direvisi 21
  Desember 2023, diterima untuk publikasi 20 Februari 2024, dan tersedia
  daring 1 Maret 2024
\item
  \textbf{Status keluaran:} review literatur per kode berdasarkan teks
  lengkap yang tersedia dalam file TXT
\end{itemize}

\subsection{Summarized Abstract}\label{summarized-abstract-8}

\textbf{Laporan sumber:} Artikel ini menilai peran \emph{ride-hailing}
berbasis sepeda motor dalam perjalanan komuter di Jakarta Metropolitan
Area (JMA) dan Bandung Metropolitan Area (BMA). Penelitian memeriksa
proporsi pengguna serta karakteristik sosial, ekonomi, dan spasialnya
dengan statistik deskriptif dan regresi logistik atas Commuter Household
Survey BMA 2017 dan JMA 2019. Pangsa komuter yang menggunakan moda ini
dilaporkan sebesar 4,62\% di JMA dan 3,27\% di BMA. Penulis menyimpulkan
bahwa pengguna terutama berasal dari kelompok muda dan berpendapatan
lebih rendah, termasuk pekerja informal, karena moda tersebut relatif
murah dan mudah diakses. Penulis juga menghubungkan kebutuhan terhadap
\emph{ride-hailing} dengan intensifikasi mobilitas antarkawasan
pinggiran dalam proses pascasuburbanisasi serta menyerukan penguatan
tata kelola metropolitan dan kebijakan angkutan umum lintas daerah.
(Locator: hlm. 1, Abstract; hlm. 6, Tabel 2; hlm. 9, Conclusion.)

\textbf{Interpretasi pengolahan:} Artikel menggabungkan dua persoalan,
yaitu pilihan moda pada tingkat komuter dan bentuk mobilitas
metropolitan. Kontribusi empiris utamanya terletak pada perbandingan dua
survei komuter resmi. Namun, klaim bahwa pengguna muda dan berpendapatan
rendah mendominasi tidak konsisten sepenuhnya dengan seluruh angka yang
disajikan, terutama distribusi usia JMA, distribusi pendapatan terhadap
UMR, dan koefisien pendapatan JMA. Konflik ini tidak direkonsiliasi
dalam review ini. (Locator: hlm. 7, Tabel 3; hlm. 8, Tabel 5; hlm. 9,
Conclusion.)

\subsection{Problem Statement}\label{problem-statement-8}

\textbf{Laporan sumber:} Permasalahan transportasi metropolitan
Indonesia digambarkan berat karena kemacetan, keterbatasan ruang jalan,
pertumbuhan kendaraan pribadi, lemahnya angkutan umum, dan tingginya
mobilitas antara kota inti dan wilayah pinggiran. Di JMA, daerah seperti
Bogor, Depok, Tangerang, dan Bekasi menyumbang mobilitas metropolitan
yang tinggi. Sepeda motor mencakup sekitar 60\% kendaraan di Jakarta dan
70\% di Bandung menurut angka latar yang dikutip artikel. Dalam kondisi
angkutan umum yang kurang baik, ojek menjadi alternatif penting, tetapi
pertumbuhan platform berbagi juga memunculkan persaingan dan konflik
dengan ojek konvensional. (Locator: hlm. 1, Introduction.)

Masalah ilmiah yang dinyatakan penulis ialah kekurangan penelitian
tentang arti penting \emph{ride-hailing} bagi mobilitas komuter di
kawasan metropolitan Indonesia. Penelitian terdahulu yang dibahas
cenderung terpisah antara kajian implikasi ekonomi dan hukum, persaingan
serta harga predator, identitas komunitas pengemudi, karakteristik
komuter, pilihan moda, kesediaan membayar, dan tata kelola
infrastruktur. Sejumlah studi telah menunjukkan bahwa
\emph{ride-sourcing} mendukung keberadaan angkutan umum di JMA, tetapi
belum banyak yang mengukur seberapa penting moda tersebut bagi komuter
metropolitan dan bagaimana karakteristik sosial-ekonomi serta bentuk
spasial metropolitan berkaitan dengan penggunaannya. (Locator: hlm. 1-2,
Introduction.)

\textbf{Interpretasi penulis:} Ketidaksesuaian antara mobilitas lintas
batas administratif, sistem angkutan umum, dan tata kelola yang
terfragmentasi membuat \emph{ride-hailing} bukan hanya persoalan pilihan
individu, melainkan juga persoalan metropolitan. (Locator: hlm. 2,
Introduction; hlm. 9, Conclusion.)

\subsection{Objectives}\label{objectives-8}

\textbf{Tujuan yang dinyatakan sumber:}

\begin{itemize}
\tightlist
\item
  Menilai pentingnya \emph{ride-hailing} berbasis sepeda motor sebagai
  moda perjalanan komuter di JMA dan BMA. (Locator: hlm. 2,
  Introduction.)
\item
  Mengidentifikasi karakteristik sosial, ekonomi, dan spasial
  penggunanya serta proporsinya terhadap mobilitas komuter metropolitan
  secara keseluruhan. (Locator: hlm. 1, Abstract.)
\item
  Menjawab faktor penentu yang mendorong atau menghambat penggunaan
  \emph{ride-hailing} oleh komuter. (Locator: hlm. 2, Introduction.)
\item
  Menilai hubungan bentuk spasial metropolitan dengan arti penting
  \emph{ride-hailing} bagi komuter. (Locator: hlm. 2, Introduction.)
\item
  Merumuskan implikasi kebijakan dari meningkatnya peran
  \emph{ride-hailing} dalam kawasan metropolitan. (Locator: hlm. 2,
  Introduction.)
\end{itemize}

\textbf{Catatan pengolahan:} Artikel tidak merumuskan hipotesis formal.
Pertanyaan tentang ``bentuk spasial'' dibahas melalui perbandingan pola
arus inti-pinggiran dan antarpinggiran, bukan melalui satu variabel
langsung yang diberi nama bentuk metropolitan dalam model. (Locator:
hlm. 5-6, Research methodology dan Gambar 3; hlm. 8, Tabel 5.)

\subsection{Literature Survey}\label{literature-survey-8}

\textbf{Kerangka ekonomi berbagi:} Artikel mendefinisikan ekonomi
berbagi atau konsumsi kolaboratif sebagai akses, penyediaan, atau
pembagian barang dan jasa antarrekan yang dikoordinasikan melalui
layanan daring berbasis komunitas. Kerangka ini diposisikan berlawanan
dengan asumsi kelangkaan karena teknologi memungkinkan pemanfaatan
inventaris yang sebelumnya tidak digunakan. Mengikuti Sundarajan (2016),
artikel merangkum lima ciri: pertukaran berbasis pasar antarrekan,
pemanfaatan modal berdampak tinggi, jejaring berbasis kerumunan,
kaburnya batas personal dan profesional, serta kaburnya batas antara
pekerjaan tetap dan kasual. Artikel menggunakan ciri-ciri itu untuk
menjelaskan peluang sekaligus ketidakjelasan regulasi dan status kerja
dalam \emph{ride-hailing}. (Locator: hlm. 2, Literature framework.)

\textbf{Adopsi dan dampak \emph{ride-hailing}:} Literatur yang dirangkum
menyebut usia, jenis kelamin, pendapatan, dan pendidikan sebagai
prediktor adopsi. Pengguna dalam studi terdahulu umumnya digambarkan
muda, berpendidikan baik, berpendapatan menengah, dan cenderung
melakukan perjalanan pendek. Literatur lain menelaah regulasi,
institusi, konflik sosial, dan hubungan \emph{ride-hailing} dengan
angkutan umum. Artikel juga menunjukkan perdebatan dampak keberlanjutan:
berbagi kendaraan atau sepeda dapat mengurangi kepemilikan mobil,
sedangkan \emph{ride-hailing} dapat menambah perjalanan; moda ini juga
tidak otomatis menggantikan angkutan massal atau terhindar dari
kemacetan. (Locator: hlm. 2-3, Literature framework dan Metropolitan
areas and sharing economy.)

\textbf{Kawasan metropolitan:} Literatur metropolitan digunakan untuk
menjelaskan aglomerasi, urbanisasi, daya saing, perluasan wilayah
terbangun, ketimpangan, dan kebutuhan pengelolaan pertumbuhan. Untuk
JMA, artikel menekankan ketimpangan upah, tingginya kepemilikan sepeda
motor, arus harian dari pinggiran menuju Jakarta, pembangunan kota baru,
segregasi sosial-ekonomi, dan kerugian akibat kemacetan. Untuk BMA,
artikel menekankan ketergantungan wilayah pinggiran kepada Kota Bandung,
perkembangan linear mengikuti jaringan jalan, serta potensi perluasan
dan konurbasi. Peta batas inti dan pinggiran ditampilkan pada Gambar 1
dan Gambar 2. (Locator: hlm. 3-4, Jakarta and Bandung metropolitan
areas; Gambar 1-2.)

\textbf{Konteks Indonesia:} Kajian yang dirangkum menyebut
\emph{ride-hailing} mudah ditemukan, mudah diakses, dan dapat
mempersingkat waktu perjalanan dalam lalu lintas padat. Gojek dan Grab
digambarkan berkembang dari penghubung pengendara dan penumpang menjadi
platform multisisi dengan layanan lain. Alasan penggunaan yang dikutip
dari penelitian terdahulu meliputi waktu tempuh lebih singkat,
penghindaran kemacetan, dan penghindaran desakan di angkutan umum. Pada
saat yang sama, terdapat konflik dengan moda konvensional, pelarangan
penjemputan di lokasi tertentu, serta kebutuhan pengaturan tempat tunggu
dan operasi. (Locator: hlm. 5, The emergence of the sharing economy.)

\textbf{Celah dalam literatur:} Artikel menyoroti kurangnya komunikasi
antara regulasi pemerintah dan ekonomi berbagi, ketidakjelasan dampak
keberlanjutan, pandangan bahwa mobilitas berbagi tidak dapat
menggantikan angkutan massal, kurangnya pengetahuan pemerintah mengenai
keterlibatan masyarakat dengan layanan sesuai permintaan, dan
terbatasnya perhatian pada ruang fisik yang diperlukan oleh kota
berbagi. Kajian pascasuburbanisasi juga disebut jauh lebih berkembang
untuk JMA daripada BMA, sedangkan temuan terdahulu mengenai hubungan
gender dan penggunaan \emph{ride-hailing} di Indonesia dinilai belum
konklusif. (Locator: hlm. 3, Metropolitan areas and sharing economy;
hlm. 5, bagian sebelum Data; hlm. 8, High socio-economic connectivity,
poor infrastructure accessibility.)

\subsection{Method Used}\label{method-used-8}

\textbf{Laporan sumber:} Penelitian menggunakan rancangan kuantitatif
komparatif dengan statistik deskriptif dan regresi logistik. Analisis
pertama mendeskripsikan karakteristik pengguna \emph{ride-hailing}
berbasis sepeda motor, perilaku perjalanan, proporsi pengguna, dan arus
komuter di JMA serta BMA. Analisis kedua mengestimasi faktor penentu
probabilitas penggunaan moda tersebut melalui regresi logistik terpisah
untuk kedua wilayah. (Locator: hlm. 2, Introduction; hlm. 6, Research
methodology.)

Variabel penjelas dikelompokkan menjadi kondisi sosial-ekonomi (usia,
jenis kelamin, jenis aktivitas, gaji, dan jenis pekerjaan), pola komuter
(jarak dan waktu keberangkatan/perjalanan), serta tujuan perjalanan
(inti atau pinggiran). Tabel hasil model juga memasukkan rasio biaya
perjalanan terhadap pengeluaran. Artikel melaporkan koefisien,
\emph{odds ratio}, galat baku, tingkat signifikansi, jumlah observasi,
dan nilai yang diberi label \texttt{Adj\ R2}. (Locator: hlm. 6, Research
methodology; hlm. 8, Tabel 5.)

\textbf{Interpretasi pengolahan:} Desain ini mengidentifikasi asosiasi,
bukan efek kausal. File tidak memberikan persamaan model, metode
pembobotan, kategori referensi seluruh variabel dummy, penanganan data
hilang, pengujian diagnostik, interval kepercayaan, atau penjelasan
apakah \texttt{Adj\ R2} merupakan ukuran \emph{pseudo-R-squared} untuk
model logistik. Karena itu, \emph{odds ratio} harus dibaca sesuai
pelaporan tabel dan tidak boleh diubah menjadi klaim perubahan
probabilitas tanpa informasi tambahan. (Locator: hlm. 5-8, Data and
methodology dan Tabel 4-5.)

\subsection{Dataset}\label{dataset-8}

\textbf{Nama dan pemilik data:} Commuter Household Survey yang dibeli
dari Badan Pusat Statistik (BPS), yakni survei BMA tahun 2017 dan survei
JMA tahun 2019. Survei memberikan informasi lokasi tempat tinggal,
pendapatan, rincian pekerjaan, pendidikan, kondisi kesehatan, tujuan
perjalanan harian, dan moda transportasi komuter. (Locator: hlm. 2,
Introduction; hlm. 5, Data.)

\textbf{Rancangan data yang dijelaskan sumber:} Strategi sampel
memanfaatkan Sensus Penduduk Indonesia 2010 untuk stratifikasi spasial
dengan pemilihan sampel dua tahap pada tingkat kecamatan dan SAKERNAS
tahunan untuk menentukan distribusi sampel. Data disebut diagregasikan
pada tingkat kabupaten/kota untuk analisis deskriptif dan komparatif.
(Locator: hlm. 5, Data.)

\textbf{Cakupan geografis dan waktu:}

\begin{itemize}
\tightlist
\item
  JMA 2019 mencakup 13 wilayah: Jakarta Pusat, Jakarta Utara, Jakarta
  Selatan, Jakarta Timur, dan Jakarta Barat sebagai inti; Depok,
  Kabupaten Bogor, Kota Bogor, Kabupaten Bekasi, dan Kota Bekasi di Jawa
  Barat; serta Kabupaten Tangerang, Kota Tangerang, dan Kota Tangerang
  Selatan di Banten sebagai pinggiran. (Locator: hlm. 5, Data.)
\item
  BMA 2017 mencakup lima wilayah: Kota Bandung sebagai inti serta
  Kabupaten Bandung, Kabupaten Bandung Barat, Kabupaten Sumedang, dan
  Kota Cimahi sebagai pinggiran. (Locator: hlm. 5, Data.)
\end{itemize}

\textbf{Catatan pengolahan:} File menyebut cakupan JMA terdiri atas
``eight municipalities and three regencies'', tetapi daftar setelahnya
berisi 13 wilayah. File juga menyatakan data diagregasikan pada tingkat
kabupaten/kota, sedangkan regresi dan jumlah observasi disajikan
seolah-olah pada tingkat komuter yang telah diekspansi. Tingkat unit
data yang benar-benar masuk ke model dan prosedur ekspansinya tidak
dijelaskan. (Locator: hlm. 5, Data; hlm. 8, Tabel 5.)

\subsection{Population / Sample}\label{population-sample-8}

\textbf{Populasi sasaran:} Komuter yang tinggal dan melakukan perjalanan
di 18 wilayah JMA dan BMA. Definisi operasional ``komuter'', kerangka
populasi, serta kriteria inklusi dan eksklusi responden tidak disebutkan
dalam file. (Locator: hlm. 2, Introduction; hlm. 5, Data.)

\textbf{Sampel keseluruhan:} Artikel menyebut lebih dari 4.000 responden
yang mewakili lebih dari 2,7 juta komuter. Pada bagian pendahuluan dan
Tabel 5, angka populasi atau observasi hasil ekspansi dilaporkan sebagai
2.459.990 untuk JMA dan 403.063 untuk BMA. Jumlah keduanya adalah
2.863.053, tetapi file tidak memberikan total responden mentah yang
tepat untuk setiap metropolitan maupun cara mengubah responden sampel
menjadi jumlah observasi tersebut. (Locator: hlm. 2, Introduction; hlm.
5, Data; hlm. 8, Tabel 5.)

\textbf{Subsampel pinggiran:} Bagian metodologi menyebut 2.856 responden
dari pinggiran JMA yang mewakili 2.236.609 komuter dan 1.112 responden
dari pinggiran BMA yang mewakili 580.921 komuter. Untuk perjalanan
menuju inti, jumlah terbesar di JMA berasal dari Kabupaten Bogor
(483.053), Kota Depok (405.945), dan Kota Bekasi (374.188); di BMA,
jumlah terbesar berasal dari Kabupaten Bandung (301.252) dan Kabupaten
Bandung Barat (141.698). (Locator: hlm. 5, Research methodology; Tabel
1.)

\textbf{Teknik pemilihan sampel:} Sumber menyebut stratifikasi spasial
dan pemilihan dua tahap pada tingkat kecamatan berdasarkan Sensus
Penduduk 2010, dengan distribusi sampel mengacu pada SAKERNAS. Ukuran
klaster, probabilitas pemilihan, bobot survei, tingkat respons,
penanganan nonrespons, dan galat desain tidak disebutkan dalam file.
(Locator: hlm. 5, Data.)

\subsection{Dependent Variables}\label{dependent-variables-8}

\textbf{Variabel dependen regresi:} Penggunaan \emph{ride-hailing}
berbasis sepeda motor sebagai variabel kategoris atau
\texttt{Mode\ (dummy)} dengan nilai minimum 0 dan maksimum 1. Artikel
menyatakan bahwa nilai 0/1 menunjukkan penggunaan \emph{ride-hailing},
tetapi tidak menjelaskan secara eksplisit label mana yang diberi nilai
1, definisi satu kejadian penggunaan, periode rujukan, atau apakah
penggunaan konvensional dan berbasis aplikasi dibedakan. (Locator: hlm.
6, Research methodology; hlm. 8, Tabel 4.)

\textbf{Keluaran deskriptif:} Selain variabel dependen model, artikel
menghitung proporsi pengguna menurut wilayah, usia, jenis kelamin,
aktivitas, alasan memilih moda, kelompok upah, status pekerjaan, jarak,
waktu perjalanan, dan arah tujuan. Ini merupakan keluaran deskriptif,
bukan variabel dependen tambahan yang dirumuskan dalam file. (Locator:
hlm. 6-7, Tabel 2-3.)

\textbf{Prediktor untuk konteks model:} Usia; jenis kelamin; pekerjaan
formal/nonformal; gaji; jarak perjalanan; keberangkatan sebelum pukul
07.00; tujuan menuju inti; serta rasio biaya perjalanan terhadap
pengeluaran. Kategori referensi dan pengodean rinci, selain rentang
dummy yang ditampilkan pada Tabel 4, tidak disebutkan dalam file.
(Locator: hlm. 8, Tabel 4-5.)

\subsection{Results}\label{results-8}

\textbf{Proporsi penggunaan:} Tabel 2 melaporkan bahwa pangsa komuter
pengguna \emph{ride-hailing} berbasis sepeda motor adalah 4,62\% di JMA
dan 3,27\% di BMA. Menurut kelompok upah, proporsi JMA dan BMA
berturut-turut adalah 4,54\% dan 4,03\% untuk upah di bawah Rp2 juta;
3,78\% dan 3,54\% untuk Rp2-2,9 juta; 2,54\% dan 0,28\% untuk Rp3-3,9
juta; serta 4,71\% dan 3,73\% untuk Rp4-4,9 juta. Untuk status kerja,
Tabel 2 mencatat 11,74\% pekerja informal dan 4,05\% pekerja/karyawan di
JMA, serta 2,48\% pekerja/karyawan di BMA; nilai pekerja informal BMA
tidak tersedia. (Locator: hlm. 6, Tabel 2.)

\textbf{Arus spasial:} Gambar 3 menunjukkan bahwa JMA memiliki arus
tujuan yang lebih beragam, termasuk hubungan kuat Kota Bekasi-Kabupaten
Bekasi dan Kota Tangerang-Kabupaten Tangerang. BMA tetap lebih
berorientasi ke Kota Bandung, walaupun terdapat arus antara Kabupaten
Bandung Barat dan Kota Cimahi. Penulis melaporkan 1.218.944 komuter dari
pinggiran menuju Jakarta, dengan 33.233 atau 2,65\% menggunakan
\emph{ride-hailing}, serta 350.055 komuter dari pinggiran menuju Kota
Bandung, dengan 13.549 atau 3,87\% menggunakan moda tersebut. (Locator:
hlm. 6, Gambar 3; hlm. 7, lanjutan Tabel 3.)

\textbf{Usia:} Tabel 3 mencatat distribusi pengguna JMA sebesar 37,19\%
untuk usia di bawah 18 tahun, 17,62\% untuk 18-23 tahun, 24,47\% untuk
23-40 tahun, 20,71\% untuk 40-60 tahun, dan 0,01\% untuk usia di atas 60
tahun. Di BMA, angkanya berturut-turut 25,48\%, 38,02\%, 13,74\%,
22,75\%, dan 0,01\%. Prosa artikel menyatakan bahwa kelompok 23-40 tahun
mendominasi JMA, padahal nilai terbesar pada Tabel 3 adalah kelompok di
bawah 18 tahun. Untuk BMA, prosa dan tabel sama-sama menunjukkan
konsentrasi pengguna yang kuat pada kelompok di bawah 23 tahun.
(Locator: hlm. 7, Tabel 3 dan uraian Ride-hailing commuters'
characteristics.)

\textbf{Jenis kelamin dan aktivitas:} Perempuan mencakup 59,12\%
pengguna di JMA dan 70,17\% di BMA; laki-laki masing-masing 40,88\% dan
29,83\%. Tujuan bekerja mencakup 51,21\% pengguna JMA dan 56,41\%
pengguna BMA, sedangkan tujuan sekolah mencakup 48,79\% dan 43,69\%.
Pendidikan informal/kursus hanya tercatat 0,01\% di kedua wilayah.
(Locator: hlm. 7, Tabel 3.)

\textbf{Alasan memilih moda:} Di JMA, alasan terbesar adalah lebih cepat
(59,91\%), diikuti kenyamanan (19,79\%) dan kemudahan akses atau
kepraktisan (18,70\%); keselamatan tercatat 0,00\% dan biaya 1,59\%. Di
BMA, kemudahan akses atau kepraktisan mendominasi (59,76\%), disusul
kecepatan (31,68\%), kenyamanan (5,14\%), biaya (1,84\%), dan
keselamatan (1,57\%). (Locator: hlm. 7, Tabel 3.)

\textbf{Pendapatan dan status kerja pengguna:} Berdasarkan klasifikasi
terhadap UMR pada Tabel 3, 27,34\% pengguna JMA berada di bawah UMR dan
72,66\% di atas UMR; di BMA, 56,82\% berada di bawah UMR dan 43,18\% di
atas UMR. Distribusi status pengguna JMA adalah 5,86\% informal, 75,24\%
pekerja/karyawan, dan 18,90\% pekerja mandiri/pengusaha. BMA mencatat
59,79\% pekerja/karyawan dan 40,21\% pekerja mandiri/pengusaha, tanpa
sampel pekerja informal. (Locator: hlm. 7, Tabel 3.)

\textbf{Jarak dan waktu:} Di JMA, 52,36\% perjalanan pengguna berjarak
kurang dari 10 km, 20,11\% berjarak 10-19 km, dan 27,54\% melampaui 20
km. Di BMA, proporsinya 55,66\%, 33,12\%, dan 11,21\%. Untuk durasi,
66,38\% pengguna JMA dan 47,21\% pengguna BMA melakukan perjalanan
kurang dari 30 menit; 22,76\% dan 41,23\% selama 30-60 menit; serta
10,86\% dan 11,56\% lebih dari 60 menit. (Locator: hlm. 7, Tabel 3.)

\textbf{Statistik deskriptif model:} Tabel 4 memberikan usia rata-rata
31,49 tahun di BMA (SD 13,69; rentang 6-85) dan 37,63 tahun di JMA (SD
11,58; rentang 17-83). Rata-rata gaji yang tercetak ialah Rp3.536.210 di
BMA dan Rp5.518.431 di JMA. Jarak JMA memiliki rata-rata 21,02 km (SD
15,34; rentang 0-160), sedangkan jarak BMA ditampilkan sebagai dummy
rentang 1-7 tanpa rata-rata. Rata-rata biaya yang tercetak adalah
Rp11.592 untuk BMA dan Rp127.761 untuk JMA. (Locator: hlm. 8, Tabel 4.)

\textbf{Hasil regresi logistik:}

\begin{itemize}
\tightlist
\item
  Usia berasosiasi negatif dan signifikan dengan penggunaan: \emph{odds
  ratio} (OR) 0,980 di JMA dan 0,977 di BMA, keduanya bertanda
  \texttt{***} atau p \textless{} 0,001. Penulis menafsirkannya sebagai
  penurunan kecenderungan penggunaan ketika usia bertambah. (Locator:
  hlm. 7-8, Determinants of using motorcycle-based ride hailing; Tabel
  5.)
\item
  Variabel jenis kelamin memiliki OR 10,118 di JMA dan 12,400 di BMA,
  keduanya p \textless{} 0,001. Prosa sumber menafsirkannya sebagai
  peluang lebih tinggi bagi perempuan, tetapi menyebut OR tersebut
  sebagai ``higher probabilities'' dan tidak menjelaskan kategori
  referensinya. (Locator: hlm. 7-8, Tabel 4-5.)
\item
  Status pekerjaan formal/nonformal memiliki OR 5,006 di JMA dan 0,426
  di BMA, keduanya p \textless{} 0,001. Penulis menafsirkan pekerja
  formal lebih mungkin menggunakan moda ini di JMA, tetapi lebih rendah
  di BMA. (Locator: hlm. 7-8, Determinants; Tabel 5.)
\item
  Gaji memiliki OR 1,357 dan signifikan di JMA (p \textless{} 0,001),
  tetapi OR 1,022 tanpa tanda signifikansi di BMA. Dengan demikian,
  model yang disajikan hanya mendukung asosiasi positif gaji di JMA.
  (Locator: hlm. 8, Tabel 5.)
\item
  Jarak memiliki OR 0,924 di JMA dan 0,454 di BMA, keduanya p
  \textless{} 0,001. Penulis mengaitkannya dengan kecenderungan
  penggunaan pada perjalanan yang lebih pendek. Perbandingan besar OR
  perlu dibatasi karena Tabel 4 menampilkan jarak JMA dalam kilometer,
  tetapi jarak BMA sebagai dummy rentang 1-7. (Locator: hlm. 7-8, Tabel
  4-5.)
\item
  Keberangkatan sebelum pukul 07.00 memiliki OR 0,609 dan signifikan di
  JMA, tetapi OR 1,039 tanpa tanda signifikansi di BMA. (Locator: hlm.
  8, Tabel 5.)
\item
  Tujuan ke inti metropolitan memiliki OR 2,582 di JMA dan 1,787 di BMA;
  rasio biaya perjalanan terhadap pengeluaran memiliki OR 4,325 dan
  4,231. Keempat estimasi bertanda p \textless{} 0,001. (Locator: hlm.
  8, Tabel 5.)
\item
  Jumlah observasi yang dicetak adalah 2.459.990 untuk JMA dan 403.063
  untuk BMA, dengan \texttt{Adj\ R2} masing-masing 0,323 dan 0,309.
  (Locator: hlm. 8, Tabel 5.)
\end{itemize}

\subsection{Findings}\label{findings-8}

\textbf{Temuan yang dinyatakan penulis:}

\begin{itemize}
\tightlist
\item
  \emph{Ride-hailing} berbasis sepeda motor memiliki pangsa kecil
  terhadap seluruh komuter, tetapi diposisikan sebagai moda alternatif
  yang penting karena cepat, praktis, mudah diakses, dan mampu
  menghadapi kemacetan lebih baik daripada pilihan yang tersedia bagi
  sebagian komuter. (Locator: hlm. 6-8, Tabel 2-3 dan Social justice,
  metropolitan areas and the sharing economy.)
\item
  Penulis menyatakan pengguna didominasi komuter muda dan berpendapatan
  lebih rendah, termasuk pekerja informal. Biaya perjalanan yang menurut
  penulis kurang dari 6\% pengeluaran bulanan bagi lebih dari 82\%
  komuter dipakai untuk menjelaskan keterjangkauannya bagi kelompok
  berpendapatan rendah dan menengah. (Locator: hlm. 8, Social justice,
  metropolitan areas and the sharing economy; hlm. 9, Conclusion.)
\item
  Perempuan dilaporkan jauh lebih mungkin menggunakan moda tersebut.
  Penulis menghubungkan pola ini dengan mobilitas dan kemandirian
  perempuan serta profil komuter perempuan yang muda dan memiliki waktu
  perjalanan relatif singkat. (Locator: hlm. 8-9, Discussion.)
\item
  JMA memperlihatkan konektivitas antarpinggiran dan keragaman tujuan
  yang lebih kuat, yang ditafsirkan sebagai pascasuburbanisasi. BMA
  masih lebih bergantung pada Kota Bandung sebagai inti. Penulis
  menyatakan lebih dari 40\% komuter pinggiran JMA bepergian ke
  pinggiran lain, sedangkan proporsinya kurang dari 25\% di BMA.
  (Locator: hlm. 6, Gambar 3; hlm. 9, Conclusion.)
\item
  Penulis menafsirkan perbedaan metropolitan sebagai indikasi bahwa
  ukuran, struktur spasial, jalan raya, alternatif jalan, dan kualitas
  angkutan umum berkaitan dengan penggunaan \emph{ride-hailing}.
  (Locator: hlm. 8-9, Discussion dan Conclusion.)
\end{itemize}

\textbf{Interpretasi pengolahan berdasarkan file:}

\begin{itemize}
\tightlist
\item
  Bukti paling langsung mendukung tiga pola: pangsa penggunaan sedikit
  lebih tinggi di JMA, penggunaan lebih terkonsentrasi pada usia muda,
  dan alasan utama berbeda antara JMA (kecepatan) dan BMA
  (akses/kepraktisan). Pola gender juga kuat dalam tabel, walaupun besar
  persisnya tidak sama dengan klaim pada kesimpulan. (Locator: hlm. 6-7,
  Tabel 2-3; hlm. 9, Conclusion.)
\item
  Klaim dominasi pendapatan rendah hanya jelas untuk distribusi BMA
  terhadap UMR. Di JMA, Tabel 3 justru mencatat 72,66\% pengguna berada
  di atas UMR dan model memberikan OR gaji 1,357 yang positif serta
  signifikan. Klaim kelompok rendah-menengah dapat menggambarkan
  keterjangkauan atau proporsi penggunaan dalam kelompok tertentu,
  tetapi tidak boleh dibaca sebagai pola seragam di kedua metropolitan.
  (Locator: hlm. 7, Tabel 3; hlm. 8, Tabel 5; hlm. 9, Conclusion.)
\item
  Perbandingan JMA-BMA bersifat sugestif, bukan pengujian kausal
  pengaruh bentuk metropolitan. Survei berasal dari tahun berbeda dan
  model tidak memasukkan ukuran langsung bentuk metropolitan. (Locator:
  hlm. 5, Data; hlm. 8, Tabel 5.)
\item
  Pernyataan bahwa perjalanan \emph{ride-hailing} jarak jauh dan
  berdurasi lama memperparah kemacetan merupakan interpretasi penulis;
  tidak ada variabel hasil kemacetan, emisi, perpindahan dari angkutan
  umum, atau pertambahan perjalanan yang diuji dalam model. (Locator:
  hlm. 9, Conclusion.)
\end{itemize}

\subsection{Conclusions}\label{conclusions-8}

\textbf{Kesimpulan sumber:} Penulis menyimpulkan bahwa
\emph{ride-hailing} berbasis sepeda motor memainkan peran penting
sebagai alternatif angkutan komuter di JMA dan BMA. Kelompok muda dan
berpendapatan lebih rendah, termasuk pekerja informal, dinyatakan
sebagai pengguna dominan karena ojek relatif terjangkau dan mudah
diakses. Pada saat yang sama, preferensi dilaporkan meningkat ketika
pendapatan melampaui upah minimum, khususnya untuk aktivitas kerja, dan
dikaitkan dengan keandalan moda dalam menghadapi kemacetan. (Locator:
hlm. 9, Conclusion.)

Penulis juga menyimpulkan bahwa bentuk metropolitan berkaitan dengan
perilaku perjalanan. JMA yang lebih luas memperlihatkan konektivitas
antarpinggiran lebih kuat, sedangkan BMA lebih bergantung pada Kota
Bandung. Pola tersebut dibaca sebagai pascasuburbanisasi pada tingkat
yang berbeda dan sebagai alasan untuk memperkuat kewenangan serta
koordinasi transportasi pada skala metropolitan. (Locator: hlm. 9,
Conclusion.)

\textbf{Batas pembacaan:} Kesimpulan di atas adalah kesimpulan penulis,
bukan fakta kausal yang dibuktikan oleh desain penelitian. Data
mendukung asosiasi dan perbedaan deskriptif, tetapi tidak mengisolasi
pengaruh bentuk metropolitan, kualitas angkutan umum, jalan raya, atau
kemacetan terhadap keputusan menggunakan \emph{ride-hailing}. Beberapa
kesimpulan juga berkonflik dengan angka tabel, sebagaimana dirinci pada
bagian Limitations dan Provenance Notes.

\subsection{Contributions}\label{contributions-8}

\textbf{Kontribusi yang diklaim sumber:}

\begin{itemize}
\tightlist
\item
  Menurut penulis, penelitian ini merupakan penelitian pertama, sejauh
  pengetahuan penulis, yang menggunakan data komuter untuk menelaah
  peran kritis \emph{ride-hailing} di kawasan metropolitan Indonesia.
  Klaim kebaruan ini tidak diverifikasi di luar file. (Locator: hlm. 9,
  Conclusion.)
\item
  Penelitian memperluas literatur rezim perkotaan dan pascasuburbanisasi
  dengan bukti mengenai intensifikasi perjalanan antarpinggiran JMA
  serta membandingkannya dengan ketergantungan pinggiran BMA kepada kota
  inti. (Locator: hlm. 9, Conclusion.)
\item
  Penelitian menambah kajian metropolitan di Global South dengan
  menghubungkan jarak dan durasi penggunaan \emph{ride-hailing},
  kekurangan angkutan umum, serta kemacetan. Hubungan dengan kemacetan
  merupakan interpretasi penulis dan bukan keluaran langsung model.
  (Locator: hlm. 9, Conclusion.)
\item
  Penelitian menambah kajian ekonomi berbagi dengan menunjukkan posisi
  \emph{ride-sourcing} sebagai alternatif angkutan umum, terutama bagi
  komuter yang oleh penulis dikelompokkan sebagai berpendapatan rendah
  dan menengah. (Locator: hlm. 8-9, Discussion dan Conclusion.)
\item
  Penelitian menghasilkan usulan kebijakan mengenai keselamatan,
  keterjangkauan, regulasi lintas daerah, pelayanan angkutan umum, dan
  tata kelola metropolitan. (Locator: hlm. 9, Conclusion.)
\end{itemize}

\textbf{Interpretasi pengolahan:} Nilai tambah yang paling dapat
ditelusuri dari file adalah pemakaian dua survei resmi untuk
menyandingkan profil pengguna, arus inti-pinggiran, dan determinan
statistik di dua metropolitan. Kontribusi teoretis terhadap ``rezim
perkotaan'' dinyatakan dalam kesimpulan, tetapi konsep tersebut tidak
dioperasionalkan secara eksplisit dalam model. Kontribusi mengenai
kemacetan juga lebih tepat dibaca sebagai proposisi interpretatif
daripada bukti dampak.

\subsection{Research Gap}\label{research-gap-8}

\textbf{Celah yang dinyatakan sumber:}

\begin{itemize}
\tightlist
\item
  Masih sedikit kajian yang menilai arti penting ekonomi berbagi,
  khususnya \emph{ride-hailing}, bagi mobilitas komuter di kawasan
  metropolitan Indonesia. (Locator: hlm. 1-2, Introduction.)
\item
  Kajian \emph{ride-hailing} Indonesia lebih banyak membahas ekonomi,
  hukum, kompetisi, harga predator, atau pembentukan identitas komunitas
  pengemudi, sedangkan kajian transportasi metropolitan lebih banyak
  membahas karakteristik komuter, pilihan moda, kesediaan membayar, dan
  tata kelola infrastruktur secara terpisah. (Locator: hlm. 1-2,
  Introduction.)
\item
  Studi yang menghubungkan \emph{ride-sourcing} dengan angkutan umum
  telah tersedia untuk JMA, tetapi arti penting moda tersebut bagi
  komuter metropolitan, faktor penentunya, dan kaitannya dengan bentuk
  spasial belum banyak ditelaah. (Locator: hlm. 2, Introduction.)
\item
  Literatur pascasuburbanisasi lebih banyak berfokus pada JMA;
  penelitian tentang BMA masih terbatas. (Locator: hlm. 8, High
  socio-economic connectivity, poor infrastructure accessibility.)
\item
  Pengetahuan pemerintah dan perencana mengenai keterlibatan masyarakat
  dengan layanan sesuai permintaan, dampak keberlanjutan, serta
  kebutuhan ruang fisik kota berbagi masih kurang. (Locator: hlm. 3,
  Metropolitan areas and sharing economy.)
\item
  Bukti tentang hubungan gender dengan penggunaan \emph{ride-hailing} di
  Indonesia belum konklusif dan pembahasan gender komuter metropolitan
  masih terbatas. (Locator: hlm. 5, bagian sebelum The emergence of the
  sharing economy; hlm. 8-9, Discussion.)
\end{itemize}

\textbf{Interpretasi pengolahan:} Artikel berupaya mengisi celah
deskriptif dan asosiatif tersebut, tetapi belum menutup celah kausal
mengenai apakah \emph{ride-hailing} melengkapi atau menggantikan
angkutan umum, menambah perjalanan, atau memperburuk kemacetan. Artikel
juga belum menguji mekanisme yang menghasilkan perbedaan JMA dan BMA.

\subsection{Limitations}\label{limitations-8}

Tidak disebutkan dalam file sebagai bagian khusus yang berjudul
keterbatasan.

\textbf{Keterbatasan yang diakui atau dinyatakan langsung oleh sumber:}

\begin{itemize}
\tightlist
\item
  Survei BMA tidak memiliki sampel komuter yang bekerja di sektor
  informal sehingga perbandingan kelompok tersebut dengan JMA tidak
  dapat dilakukan. (Locator: hlm. 7, Ride-hailing commuters'
  characteristics.)
\item
  Surabaya sempat dipertimbangkan, tetapi dikeluarkan karena jumlah
  sampel pengguna \emph{ride-hailing} sangat kecil. (Locator: hlm. 2,
  catatan kaki 2.)
\item
  Kesimpulan berasal dari data sampel responden komuter JMA dan BMA;
  sumber tidak menyatakan bahwa hasil berlaku untuk semua kota atau
  semua jenis perjalanan. (Locator: hlm. 9, Conclusion.)
\end{itemize}

\textbf{Keterbatasan yang diidentifikasi melalui pengolahan file:}

\begin{itemize}
\tightlist
\item
  Survei membandingkan BMA 2017 dengan JMA 2019. Perbedaan wilayah tidak
  dapat dipisahkan dari perbedaan tahun pengamatan.
\item
  Desain potong lintang dan regresi logistik menunjukkan asosiasi, bukan
  sebab akibat. Klaim mengenai pengaruh bentuk metropolitan, kualitas
  angkutan umum, jalan raya, dan kemacetan tidak diuji dengan variabel
  hasil langsung.
\item
  Hubungan antara lebih dari 4.000 responden mentah, data yang disebut
  diagregasi per kabupaten/kota, dan 2.459.990 serta 403.063 observasi
  pada regresi tidak dijelaskan. Bobot, desain kompleks survei, galat
  baku berbobot, dan tingkat analisis model juga tidak dilaporkan.
\item
  Pengodean kategori referensi tidak lengkap. Tabel 4 menampilkan jenis
  kelamin sebagai dummy bernilai 1-2, bukan 0-1, dan tidak memberikan
  label kode; artikel kemudian menafsirkan koefisien sebagai efek
  perempuan. Penyebutan OR sebagai probabilitas pada prosa juga tidak
  tepat secara terminologis.
\item
  Pengukuran jarak tidak tampak setara: JMA ditampilkan dalam kilometer
  kontinu, sedangkan BMA ditampilkan sebagai dummy rentang 1-7. Ini
  membatasi perbandingan langsung OR jarak kedua model. (Locator: hlm.
  8, Tabel 4-5.)
\item
  Rentang usia pada Tabel 3 ditulis \texttt{\textless{}18},
  \texttt{18-23}, \texttt{23-40}, \texttt{40-60}, dan
  \texttt{\textgreater{}60}. Jika dibaca secara literal, batas 23 dan 40
  tahun bertumpang tindih; aturan klasifikasi pada titik batas tidak
  disebutkan dalam file. (Locator: hlm. 7, Tabel 3.)
\item
  Definisi hasil tidak cukup rinci untuk memastikan apakah
  ``\emph{ride-hailing}'' hanya berarti layanan berbasis aplikasi atau
  juga ojek konvensional. Bagian diskusi menyebut pengeluaran untuk ojek
  konvensional dan daring, sedangkan judul dan banyak bagian lain
  menggunakan istilah \emph{ride-hailing}. (Locator: hlm. 8, Social
  justice, metropolitan areas and the sharing economy; hlm. 9,
  Conclusion.)
\item
  File tidak menyajikan interval kepercayaan, uji kecocokan model,
  pemeriksaan multikolinearitas, sensitivitas spesifikasi, data hilang,
  atau validasi model.
\item
  Terdapat ketidakkonsistenan internal pada angka dan narasi. Karena
  tidak ada dasar untuk memperbaikinya dari TXT, semua konflik
  dipertahankan sebagai batas bukti dan dirinci pada Provenance Notes.
\end{itemize}

\subsection{Challenges}\label{challenges-8}

\textbf{Tantangan substantif menurut sumber:}

\begin{itemize}
\tightlist
\item
  Sistem transportasi menghadapi kemacetan berat, ruang jalan terbatas,
  tingginya kendaraan pribadi dan sepeda motor, serta akses angkutan
  umum yang belum memadai. (Locator: hlm. 1, Introduction; hlm. 8,
  Discussion.)
\item
  Mobilitas telah melampaui batas kota/kabupaten dan provinsi, sedangkan
  perencanaan serta regulasi masih terfragmentasi. Ini menyulitkan
  penyediaan angkutan umum dan pengaturan \emph{ride-hailing} lintas
  daerah. (Locator: hlm. 3, Metropolitan areas and sharing economy; hlm.
  9, Conclusion.)
\item
  Pemerintah dan perencana belum memiliki kepastian mengenai dampak
  keberlanjutan mobilitas berbagi, apakah layanan tersebut mengurangi
  kepemilikan kendaraan atau justru menambah perjalanan, dan bagaimana
  hubungan layanan tersebut dengan angkutan massal. (Locator: hlm. 3,
  Metropolitan areas and sharing economy.)
\item
  Konflik antara pengemudi daring dan moda konvensional, larangan
  operasi atau penjemputan di tempat tertentu, serta kebutuhan ruang
  tunggu dan parkir merupakan tantangan implementasi. (Locator: hlm. 1,
  Introduction; hlm. 5, The emergence of the sharing economy.)
\item
  Keselamatan dan keterjangkauan tetap menjadi tantangan kebijakan
  meskipun keduanya bukan alasan utama yang dipilih responden dalam
  Tabel 3. (Locator: hlm. 7, Tabel 3; hlm. 9, Conclusion.)
\end{itemize}

\textbf{Tantangan data dan analisis:} Kekosongan sampel pekerja informal
BMA, kecilnya sampel Surabaya, perbedaan tahun survei, ketidakjelasan
bobot, dan inkonsistensi pelaporan menghambat perbandingan yang
sepenuhnya setara. Ini adalah rangkuman keterbatasan yang tampak dalam
file, bukan tantangan tambahan dari sumber eksternal.

\subsection{Practical Implications}\label{practical-implications-8}

\textbf{Usulan penulis:}

\begin{itemize}
\tightlist
\item
  Pemerintah perlu memastikan keselamatan dan keterjangkauan bagi
  komuter yang menggunakan \emph{ride-hailing} setiap hari. (Locator:
  hlm. 9, Conclusion.)
\item
  Diperlukan kerangka hukum yang mengatur layanan \emph{ride-hailing}
  daring lintas kabupaten/kota. (Locator: hlm. 9, Conclusion.)
\item
  Aksesibilitas dan kenyamanan angkutan umum perlu ditingkatkan karena
  kesediaan kelompok berpendapatan rendah membayar perjalanan pendek
  dengan \emph{ride-hailing} dibaca penulis sebagai sinyal masalah mutu
  angkutan umum. (Locator: hlm. 8-9, Social justice, metropolitan areas
  and the sharing economy; Conclusion.)
\item
  Tata kelola skala metropolitan perlu diperkuat agar keragaman arus
  antarpinggiran tidak ditangani hanya sebagai arus menuju kota inti.
  (Locator: hlm. 9, Conclusion.)
\item
  Otoritas metropolitan dan kemitraan dengan sektor swasta, terutama
  pelaku ekonomi berbagi, perlu diperkuat untuk mengoordinasikan
  perencanaan dan penyediaan transportasi yang aman, nyaman, serta hemat
  biaya. (Locator: hlm. 9, Conclusion.)
\end{itemize}

\textbf{Kualifikasi pengolahan:} Usulan tersebut berasal dari
interpretasi penulis atas pola survei. Artikel tidak mengevaluasi biaya
kebijakan, kapasitas kelembagaan, bentuk regulasi tertentu, dampak
distribusional tarif, atau hasil penerapan kebijakan.

\subsection{Applications}\label{applications-8}

\textbf{Aplikasi yang dibahas sumber:} Hasil penelitian ditujukan
sebagai dasar untuk memahami segmen pengguna, arah arus perjalanan, dan
kebutuhan koordinasi transportasi lintas daerah. Profil usia, jenis
kelamin, pendapatan, pekerjaan, jarak, waktu, serta tujuan dapat
digunakan oleh pemerintah metropolitan untuk mengenali kelompok yang
bergantung pada \emph{ride-hailing} dan untuk menempatkan layanan
tersebut dalam perencanaan angkutan umum. (Locator: hlm. 1, Abstract;
hlm. 9, Conclusion.)

\textbf{Inferensi pengolahan:} Dengan batas metodologis yang ada, data
dapat dipakai sebagai garis dasar deskriptif untuk membandingkan JMA dan
BMA, memetakan koridor inti-pinggiran serta antarpinggiran yang
memerlukan konektivitas, dan menyusun pertanyaan evaluasi tentang
keselamatan, keterjangkauan, serta integrasi moda. Model tidak cukup
terdokumentasi untuk langsung dijadikan alat prediksi operasional atau
dasar penetapan tarif tanpa validasi tambahan. Tidak ada aplikasi,
intervensi, atau program kebijakan yang diuji dalam file.

\subsection{Future Research
(Optional)}\label{future-research-optional-8}

Tidak disebutkan dalam file sebagai agenda penelitian lanjutan yang
dirumuskan secara khusus.

\textbf{Inferensi pengolahan berdasarkan celah dan batas file:}

\begin{itemize}
\tightlist
\item
  Mengulang perbandingan dengan tahun survei yang sama, definisi
  variabel yang seragam, bobot survei yang transparan, dan pelaporan
  jumlah responden mentah serta populasi hasil ekspansi.
\item
  Membedakan secara tegas ojek konvensional, \emph{ride-hailing}
  berbasis aplikasi, perjalanan pengumpan angkutan umum, dan perjalanan
  yang menggantikan angkutan umum atau kendaraan pribadi.
\item
  Menguji secara langsung apakah \emph{ride-hailing} menambah
  perjalanan, memperburuk kemacetan, memengaruhi emisi, atau melengkapi
  angkutan umum; artikel sekarang hanya menginterpretasikan hubungan
  tersebut.
\item
  Meneliti mekanisme di balik perbedaan gender, usia, pendapatan, dan
  status pekerjaan, termasuk arti keselamatan serta keterjangkauan bagi
  kelompok pengguna yang berbeda.
\item
  Menguji pengaruh bentuk metropolitan dengan indikator spasial
  eksplisit dan data arus perjalanan yang konsisten, bukan hanya
  perbandingan dua wilayah pada tahun berbeda.
\item
  Memperluas kajian ke Surabaya atau metropolitan lain setelah tersedia
  sampel pengguna yang memadai. Usulan Surabaya diturunkan dari catatan
  sumber bahwa kasus itu dikeluarkan karena sampelnya terlalu kecil.
  (Locator: hlm. 2, catatan kaki 2.)
\end{itemize}

\subsection{Provenance Notes}\label{provenance-notes-8}

\begin{itemize}
\tightlist
\item
  Review ini hanya menggunakan file
  \texttt{/Users/mac/Documents/Mac/{[}2{]}\ Obsidian\ Vault/zettelkasten/source/journal/J09-aritenang-2024-the-crucial-role-of-motorcycle-based-ride-hailing-among-commuters-the-case-of-jakarta-and-bandung-metropolitan-areas-10.1016-j.jpubtr.2024.100082.txt}
  dan tidak menggunakan sumber eksternal.
\item
  Seluruh 753 baris file dibaca, termasuk metadata ekstraksi, teks
  artikel, tabel, keterangan gambar, catatan kaki, deklarasi konflik
  kepentingan, dan daftar pustaka. Metadata menyatakan bahwa TXT
  diekstrak pada 31 Agustus 2026 dari PDF 10 halaman dengan
  \texttt{pdftotext\ -layout}.
\item
  PDF asli tidak dibuka atau diperiksa secara langsung. Status telaah
  sejawat juga tidak dinyatakan dalam file dan tidak diasumsikan hanya
  karena naskah tercantum dalam sebuah jurnal.
\item
  Locator halaman mengikuti pemisah halaman dan nomor halaman artikel
  yang dipertahankan dalam TXT. Isi visual Gambar 1-3 tidak
  diinterpretasikan melampaui keterangan dan uraian penulis yang
  tersedia dalam teks.
\item
  Angka pada tabel ditranskripsikan ke konvensi penulisan bilangan
  Indonesia tanpa mengubah nilainya. Tata letak dua kolom, pemenggalan
  kata, tanda baca angka, dan sejumlah salah ketik dalam TXT dapat
  menghasilkan ambiguitas. Tidak ada angka yang diperbaiki dengan
  merujuk PDF atau sumber lain.
\item
  Semua klaim tentang data BPS, literatur terdahulu, kondisi
  transportasi, dan kebijakan merupakan laporan artikel. Dataset asli
  BPS dan karya-karya dalam daftar pustaka tidak diperiksa sehingga
  akurasi klaim sekunder tersebut tidak diverifikasi secara independen.
\item
  Konflik internal utama yang dipertahankan: UMR disebut Rp4,7 juta
  untuk JMA dan Rp3,7 juta untuk BMA pada hlm. 6, tetapi Rp4,9 juta dan
  Rp2,8 juta pada hlm. 7; pekerja informal JMA tercatat 11,74\% pada
  Tabel 2, 5,86\% pada Tabel 3, dan 4,34\% dalam prosa hlm. 7; prosa
  menyebut usia 23-40 dominan di JMA, tetapi Tabel 3 menempatkan usia di
  bawah 18 tahun sebagai kelompok terbesar; kesimpulan menyebut lebih
  dari 72\% pengguna di kedua wilayah adalah perempuan, tetapi Tabel 3
  melaporkan 59,12\% di JMA dan 70,17\% di BMA.
\item
  Konflik lain: prosa hlm. 6 menyebut sekitar 5\% pengguna JMA dan
  3,46\% pengguna BMA untuk arus pinggiran-inti, sedangkan Tabel 3
  mencatat 2,65\% dan 3,87\%; narasi BMA menyebut kecenderungan lebih
  rendah yang signifikan bagi komuter jam awal ``after 8 am'', tetapi
  Tabel 5 mendefinisikan variabel sebagai sebelum pukul 07.00 dan
  melaporkan OR 1,039 tanpa signifikansi; kesimpulan menyatakan komuter
  \emph{ride-hailing} JMA rata-rata bepergian lebih dari satu jam sejauh
  20 km, sedangkan Tabel 3 hanya mencatat 10,86\% perjalanan di atas 60
  menit dan Tabel 4 tidak melaporkan rata-rata durasi.
\item
  Pemisahan epistemik dalam review ini menggunakan label \textbf{Laporan
  sumber}, \textbf{Interpretasi penulis}, \textbf{Temuan yang dinyatakan
  penulis}, dan \textbf{Inferensi/Interpretasi pengolahan}. Pernyataan
  berlabel pengolahan adalah pembacaan analitis atas TXT, bukan posisi
  yang dinyatakan secara langsung oleh penulis.
\end{itemize}

\section{Review J10}\label{review-j10}

\subsection{Identitas Sumber}\label{identitas-sumber-4}

\begin{itemize}
\tightlist
\item
  \textbf{Kode sumber:} J10.
\item
  \textbf{Judul:} \emph{Spatial, mobility, or socio-economic inequity? A
  district level job accessibility evaluation in Jakarta, Indonesia}.
\item
  \textbf{Penulis:} I Gusti Ayu Andani, Nadia Qamilla, Rakan Pramoe
  Izdihar, Maya Safira, Anjar Dimara Sakti, dan Ibnu Syabri.
\item
  \textbf{Afiliasi:} School of Architecture, Planning and Policy
  Development, Bandung Institute of Technology; Faculty of Earth
  Sciences and Technology, Bandung Institute of Technology; serta Urban
  and Regional Planning Department, Sumatera Institute of Technology.
\item
  \textbf{Jurnal:} \emph{International Journal of Urban Sciences}.
\item
  \textbf{Jenis dokumen:} Artikel jurnal. Status penelaahan sejawat:
  Tidak disebutkan dalam file.
\item
  \textbf{Tahun dan riwayat publikasi:} 2025; diterima redaksi pada 6
  Maret 2025, dinyatakan diterima pada 24 Agustus 2025, dan
  dipublikasikan secara daring pada 15 September 2025 (hlm. 1, blok
  Article History dan informasi sitasi).
\item
  \textbf{DOI:} 10.1080/12265934.2025.2553714.
\item
  \textbf{ISSN:} 1226-5934 (cetak) dan 2161-6779 (daring) (halaman muka
  hasil ekstraksi).
\item
  \textbf{Kata kunci:} job accessibility; equity; Gini coefficient;
  mobility; Palma ratio; Theil index (hlm. 1, blok Keywords).
\item
  \textbf{Volume, nomor terbitan, dan rentang halaman bibliografis:}
  Tidak disebutkan dalam file.
\item
  \textbf{Wilayah dan unit analisis:} Provinsi Jakarta, yang terdiri
  atas lima kota administrasi, dengan Traffic Analysis Zone (TAZ)
  sebagai unit spasial (hlm. 8-9, Bagian 3.1 dan Gambar 3).
\item
  \textbf{Pendanaan dan konflik kepentingan:} Penulis melaporkan
  dukungan Bandung Institute of Technology dengan nomor hibah
  SAPPK.PN-6-05-2023 dan menyatakan tidak ada potensi konflik
  kepentingan (hlm. 24, bagian Funding dan Disclosure statement).
\item
  \textbf{Status output:} Review per-kode selesai dan telah diperiksa
  ulang terhadap TXT J10.
\end{itemize}

\subsection{Summarized Abstract}\label{summarized-abstract-9}

\textbf{Laporan sumber:} Studi ini menilai ketimpangan aksesibilitas
pekerjaan di Jakarta melalui tiga dimensi yang saling terkait, yaitu
spasial, sosial-ekonomi, dan moda/mobilitas. Aksesibilitas dihitung
sebagai jumlah peluang kerja yang dapat dicapai dari setiap TAZ dalam
ambang waktu 41,2 menit dengan mobil, sepeda motor, atau transportasi
publik. Waktu tempuh berasal dari Google Maps API, sedangkan matriks
asal-tujuan dan proksi distribusi pekerjaan berasal dari model
permintaan perjalanan berbasis mikrosimulasi. Distribusi akses kemudian
dinilai dengan koefisien Gini, rasio Palma, dan indeks Theil (hlm. 1,
Abstract; hlm. 7-8, Bagian 3 dan Gambar 2).

\textbf{Laporan sumber:} Hasil utama menunjukkan konsentrasi akses
pekerjaan di inti kota dan akses yang lebih rendah di kawasan periferal
serta berkepadatan rendah. Sepeda motor memberikan jangkauan terluas dan
distribusi paling merata, terutama bagi kelompok berpendapatan rendah
dan menengah, sedangkan transportasi publik memperlihatkan ketimpangan
tertinggi. Koefisien Gini antarkelompok pendapatan hampir sama, dan
dekomposisi Theil menunjukkan bahwa sebagian besar ketimpangan terjadi
di dalam kelompok kepadatan penduduk, bukan di antara kelompok tersebut
(hlm. 1, Abstract; hlm. 16-23, Bagian 4; Tabel 2-4).

\textbf{Interpretasi penulis:} Kesamaan indeks menurut pendapatan
ditafsirkan sebagai tanda bahwa geografi, rancangan jaringan, dan moda
lebih menentukan pola ketimpangan dibandingkan kategori sosial-ekonomi
semata. Dominasi komponen ketimpangan di dalam kelompok kepadatan
dipakai untuk mendukung intervensi setempat pada skala lingkungan,
seperti layanan pengumpan dan konektivitas mil pertama serta terakhir,
alih-alih hanya mengandalkan kebijakan berbasis kategori ruang atau
pendapatan yang luas (hlm. 21-24, Bagian 4.2.1-5).

\textbf{Inferensi pengolahan:} Karena analisis bersifat distribusional,
memakai proksi pekerjaan, dan tidak melaporkan model kausal atau
pengujian inferensial, bukti dalam file paling kuat untuk membandingkan
pola akses antarzona dan antarmoda. Bukti tersebut tidak dengan
sendirinya menetapkan bahwa geografi atau rancangan jaringan secara
kausal lebih berpengaruh daripada pendapatan.

\subsection{Problem Statement}\label{problem-statement-9}

\textbf{Laporan sumber:} Jakarta menghadapi persoalan akses yang tidak
merata terhadap kesempatan kerja di tengah urbanisasi cepat,
heterogenitas moda, dominasi sepeda motor, dan sistem transportasi
publik massal yang masih berkembang. Ketidakmerataan ini penting karena
akses pekerjaan memengaruhi kesejahteraan ekonomi, inklusi sosial,
kebutuhan perjalanan harian, dan peluang mobilitas sosial-ekonomi (hlm.
2-3, Bagian 1).

\textbf{Laporan sumber:} Pemerintah Indonesia telah memprioritaskan BRT,
MRT, dan LRT, tetapi file menyatakan belum ada analisis mengenai
pengaruh perluasan layanan tersebut terhadap akses pekerjaan atau
pengurangan ketimpangan sosial-ekonomi dalam mobilitas. Penelitian
Indonesia terdahulu yang dibahas penulis juga belum mengakomodasi
pilihan moda ketika mengevaluasi akses pekerjaan (hlm. 3, Bagian 1).

\textbf{Laporan sumber:} Penulis membedakan \emph{inequality} sebagai
perbedaan distribusional yang dapat dideskripsikan secara statistik dari
\emph{inequity} sebagai penilaian mengenai adil atau tidak adilnya
perbedaan tersebut. Masalah empiris penelitian adalah mengidentifikasi
wilayah dan kelompok yang tertinggal serta menilai peran sistem
transportasi dalam distribusi kesempatan kerja (hlm. 3, Bagian 1; hlm.
4-5, Bagian 2.1).

\textbf{Inferensi pengolahan:} Rumusan masalah disajikan sebagai narasi,
bukan sebagai pertanyaan penelitian formal. Pertanyaan penelitian dalam
bentuk kalimat tanya: Tidak disebutkan dalam file.

\subsection{Objectives}\label{objectives-9}

\textbf{Laporan sumber:} Tujuan utama penelitian adalah mengevaluasi
ketimpangan aksesibilitas pekerjaan di Jakarta menurut lokasi, tingkat
pendapatan, dan moda transportasi, lalu menawarkan rekomendasi untuk
meningkatkan keadilan mobilitas dan distribusi kesempatan pada tingkat
distrik atau zona (hlm. 3, Bagian 1).

Secara operasional, tujuan tersebut mencakup:

\begin{itemize}
\tightlist
\item
  \textbf{Laporan sumber:} Menghitung jumlah pekerjaan yang dapat
  dicapai dari setiap TAZ dalam 41,2 menit dengan mobil, sepeda motor,
  dan transportasi publik (hlm. 7-8, Bagian 3; hlm. 11, Bagian 3.4).
\item
  \textbf{Laporan sumber:} Memetakan variasi aksesibilitas dari sangat
  rendah hingga sangat tinggi menurut moda dan kelompok pendapatan (hlm.
  12, Bagian 3.4; hlm. 16-17, Bagian 4.1 dan Gambar 8).
\item
  \textbf{Laporan sumber:} Mengukur ketimpangan distribusi secara
  keseluruhan dengan koefisien Gini dan kurva Lorenz, membandingkan
  bagian ekstrem distribusi dengan rasio Palma, serta memisahkan
  ketimpangan di dalam dan di antara kelompok kepadatan penduduk dengan
  indeks Theil (hlm. 14-16, Bagian 3.6).
\item
  \textbf{Laporan sumber:} Menentukan apakah pola ketimpangan lebih
  terkait dengan posisi spasial, kelompok sosial-ekonomi, atau moda,
  serta mengidentifikasi skala intervensi yang relevan (hlm. 1,
  Abstract; hlm. 23-24, Bagian 5).
\end{itemize}

Hipotesis formal atau kriteria keberhasilan yang ditetapkan sebelum
analisis: Tidak disebutkan dalam file.

\subsection{Literature Survey}\label{literature-survey-9}

\textbf{Laporan sumber:} Telaah literatur menempatkan aksesibilitas
sebagai kemudahan mencapai kegiatan melalui hubungan antara sistem
transportasi, tata guna lahan, dan distribusi peluang. Berbeda dari
mobilitas yang menekankan kapasitas bergerak, aksesibilitas dipakai
sebagai ukuran hasil berupa kemampuan berpartisipasi dalam kehidupan
sosial-ekonomi. Akses pekerjaan diprioritaskan karena perjalanan kerja
berlangsung rutin dan berhubungan langsung dengan kesejahteraan ekonomi,
kemiskinan waktu, serta inklusi (hlm. 2, Bagian 1; hlm. 6-7, Bagian
2.2).

\textbf{Laporan sumber:} Kerangka konseptual mengintegrasikan tiga
dimensi. Ketimpangan spasial merujuk pada variasi geografis
infrastruktur, layanan, dan peluang; ketimpangan sosial-ekonomi merujuk
pada perbedaan akses menurut pendapatan, pendidikan, atau pekerjaan;
sedangkan ketimpangan mobilitas fisik merujuk pada perbedaan kapasitas
bergerak karena ketersediaan, keterjangkauan, dan mutu moda. Penulis
memandang ketiganya saling memperkuat secara siklis (hlm. 4-6, Bagian
2.1 dan Gambar 1).

\textbf{Laporan sumber:} Studi-studi yang dirujuk menunjukkan distribusi
akses pekerjaan yang tidak merata di kota-kota Australia, Bogota,
Montevideo, Tokyo, wilayah metropolitan Amerika Serikat, Sao Paulo, dan
kota lainnya. Literatur tersebut menghubungkan akses rendah dengan
\emph{spatial mismatch}, waktu perjalanan panjang, biaya transportasi,
keterbatasan angkutan publik, dan pemusatan pekerjaan atau layanan pada
kawasan tertentu (hlm. 2, Bagian 1; hlm. 6-7, Bagian 2.2).

\textbf{Laporan sumber:} Untuk konteks Global South, terutama Asia
Tenggara, penulis menekankan campuran moda yang khas, termasuk sepeda
motor, angkutan informal, dan layanan \emph{ride-hailing}. Literatur
yang dibahas menyatakan bahwa moda tersebut dapat meningkatkan
fleksibilitas bagi masyarakat berpendapatan rendah, tetapi pola ini
belum banyak masuk ke evaluasi akses pekerjaan multidimensi di Indonesia
(hlm. 3, Bagian 1; hlm. 7, Bagian 2.2).

\textbf{Laporan sumber:} Koefisien Gini dan kurva Lorenz dipandang
berguna untuk merangkum ketimpangan umum, tetapi tidak menunjukkan
kelompok yang dirugikan dan tidak memiliki arah progresif atau regresif.
Rasio Palma ditambahkan untuk membandingkan bagian ekstrem distribusi,
sedangkan indeks Theil dipakai untuk memisahkan variasi di dalam dan di
antara kelompok (hlm. 7, Bagian 2.3).

\textbf{Inferensi pengolahan:} Telaah ini bersifat naratif dan membangun
kerangka analitis, bukan tinjauan sistematis. Basis data pencarian, kata
kunci pencarian, periode penelusuran, kriteria inklusi-eksklusi,
prosedur penilaian mutu, dan jumlah korpus literatur: Tidak disebutkan
dalam file.

\subsection{Method Used}\label{method-used-9}

\textbf{Laporan sumber:} Penelitian menggunakan desain
kuantitatif-spasial yang bersifat deskriptif dan distribusional.
Kerangka analisis memiliki dua tahap: menghitung aksesibilitas pekerjaan
dan mengukur ketimpangan distribusinya menurut moda, pendapatan, dan
kepadatan penduduk (hlm. 7-8, Bagian 3 dan Gambar 2).

\begin{itemize}
\tightlist
\item
  \textbf{Wilayah dan zonasi:} Jakarta dibagi menjadi TAZ berdasarkan
  kerangka JICA (2019). Pembentukan zona dijelaskan sebagai agregasi
  beberapa tingkat administrasi dengan dukungan regresi berganda dan
  model potensi pembangunan berbasis GIS untuk mendistribusikan proyeksi
  penduduk serta indikator sosial-ekonomi (hlm. 8-9, Bagian 3.1 dan
  Gambar 3).
\item
  \textbf{Pengukuran pekerjaan:} Jumlah pekerjaan pada TAZ tujuan
  diproksikan dengan data tarikan perjalanan, yaitu volume perjalanan
  yang menuju zona tersebut. Matriks OD berasal dari model prakiraan
  permintaan perjalanan berbasis mikrosimulasi JABODETABEK Urban
  Transportation Policy Integration Project Phase 2. Sebanyak 65\%
  perjalanan dalam data wawancara OD dikategorikan terkait pekerjaan,
  terdiri atas 33,9\% perjalanan bisnis dan 31,1\% perjalanan kerja
  (hlm. 9-10, Bagian 3.3; Tabel 1; Gambar 5).
\item
  \textbf{Waktu tempuh:} Centroid setiap TAZ menjadi titik asal dan
  tujuan. Google Maps Distance Matrix API diminta menghasilkan waktu
  tempuh untuk mobil (\texttt{driving}), sepeda motor
  (\texttt{two\_wheeler}), dan transportasi publik (\texttt{transit})
  pada Senin, 2 November 2020 pukul 10.00. Model lalu lintas
  \texttt{best\_guess} digunakan untuk moda bermotor. Respons JSON
  diurai menjadi waktu, jarak, dan informasi rute transportasi publik,
  lalu disusun antara lain dalam CSV (hlm. 12-14, Bagian 3.5 dan Gambar
  6).
\item
  \textbf{Aksesibilitas kumulatif:} Untuk setiap zona asal dan moda,
  studi menjumlahkan pekerjaan pada zona tujuan yang dapat dicapai dalam
  waktu paling lama 41,2 menit. Fungsi impedansi bernilai 1 jika waktu
  tempuh memenuhi ambang dan 0 jika melampauinya. Ambang tersebut
  didasarkan pada rata-rata waktu komuter Jakarta dari survei mobilitas
  perkotaan 2019 dan disebut telah dikonfirmasi oleh statistik mobilitas
  Google Maps (hlm. 11, Bagian 3.4, Persamaan 1-2).
\item
  \textbf{Pemetaan:} Nilai aksesibilitas dikelompokkan menjadi lima
  kelas, dari sangat rendah hingga sangat tinggi, menggunakan
  klasifikasi \emph{natural breaks} Jenks pada ArcGIS Pro (hlm. 12,
  Bagian 3.4).
\item
  \textbf{Koefisien Gini dan kurva Lorenz:} Gini merangkum konsentrasi
  aksesibilitas pada skala 0-1, sementara kurva Lorenz memperlihatkan
  distribusi kumulatif akses terhadap populasi (hlm. 14, Bagian 3.6.1,
  Persamaan 3 dan Gambar 7).
\item
  \textbf{Rasio Palma yang diadaptasi:} Karena data pendapatan tingkat
  zona tidak tersedia, pembilangnya adalah rata-rata aksesibilitas 10\%
  TAZ dengan kepadatan tertinggi dan penyebutnya rata-rata aksesibilitas
  40\% TAZ dengan kepadatan terendah. Rasio dihitung untuk moda dan
  kelompok pendapatan yang dilaporkan (hlm. 14-15, Bagian 3.6.2,
  Persamaan 4).
\item
  \textbf{Indeks Theil:} Total ketimpangan didekomposisi menjadi
  komponen di dalam kelompok dan antarkelompok berdasarkan tiga kelas
  kepadatan: rendah, kurang dari 17.090 jiwa/km2; menengah,
  17.090-50.123 jiwa/km2; dan tinggi, sedikitnya 50.124 jiwa/km2 (hlm.
  15-16, Bagian 3.6.3, Persamaan 5-7).
\end{itemize}

\textbf{Inferensi pengolahan:} File tidak melaporkan regresi,
eksperimen, pembobotan kausal, uji hipotesis, interval kepercayaan, atau
prosedur untuk mengestimasi pengaruh relatif setiap dimensi. Metode
karena itu mengukur pola distribusi dan asosiasi deskriptif, bukan efek
kausal.

\subsection{Dataset}\label{dataset-9}

\begin{itemize}
\tightlist
\item
  \textbf{Data waktu tempuh:} Estimasi Google Maps Distance Matrix API
  untuk tiga moda pada satu waktu pengamatan, 2 November 2020 pukul
  10.00. Sumber menyebut data API menggabungkan informasi lalu lintas
  waktu nyata dan historis serta menghasilkan waktu tempuh, jarak, dan
  informasi rute untuk transportasi publik (hlm. 12-14, Bagian 3.5).
\item
  \textbf{Matriks OD:} Matriks asal-tujuan dari model prakiraan
  permintaan perjalanan berbasis mikrosimulasi dalam proyek JICA (2019).
  Data ini menyediakan distribusi perjalanan menurut moda dan tingkat
  pendapatan serta menjadi dasar proksi pekerjaan melalui tarikan
  perjalanan (hlm. 9-10, Bagian 3.3; Tabel 1).
\item
  \textbf{Komposisi perjalanan:} Tabel 1 memuat 14.049.860 perjalanan.
  Menurut moda, 3.965.820 atau 28,23\% menggunakan mobil, 8.781.802 atau
  62,50\% menggunakan sepeda motor, dan 1.302.238 atau 9,27\%
  menggunakan transportasi publik. Menurut pendapatan, 2.608.244 atau
  18,56\% berasal dari kelompok tinggi, 7.669.331 atau 54,59\% dari
  kelompok menengah, dan 3.772.286 atau 26,85\% dari kelompok rendah
  (hlm. 10, Tabel 1).
\item
  \textbf{Data penduduk dan kepadatan:} Jumlah penduduk serta kepadatan
  penduduk per TAZ dipakai untuk karakterisasi zona, pengelompokan
  Theil, dan pengurutan zona dalam rasio Palma. Gambar 4 menampilkan
  distribusinya, tetapi nama himpunan data, tahun observasi khusus,
  serta definisi operasional lengkapnya tidak disebutkan dalam file
  (hlm. 9-10, Bagian 3.2 dan Gambar 4).
\item
  \textbf{Data distribusi pekerjaan:} Tidak berupa hitungan tempat kerja
  langsung; studi mengaproksimasinya dari tarikan perjalanan menuju
  setiap TAZ (hlm. 10-11, Bagian 3.3 dan Gambar 5).
\item
  \textbf{Format antara:} Respons API berformat JSON diproses menjadi
  format terstruktur seperti CSV. Skrip Python digunakan untuk
  permintaan berulang, pemrosesan kelompok, pembatasan laju, dan
  penanganan galat (hlm. 13, Bagian 3.5).
\item
  \textbf{Ketersediaan data mentah, repositori, lisensi, kode analisis,
  dan prosedur akses ulang:} Tidak disebutkan dalam file.
\end{itemize}

\textbf{Inferensi pengolahan:} Dataset menggabungkan sekurang-kurangnya
konteks tahun 2019 untuk matriks OD/ambang komuter dengan waktu tempuh
tahun 2020. File tidak menjelaskan harmonisasi temporal kedua komponen
tersebut atau pengaruh perbedaan waktunya terhadap estimasi
aksesibilitas.

\subsection{Population / Sample}\label{population-sample-9}

\textbf{Laporan sumber:} Populasi analitis adalah penduduk dan
perjalanan dalam wilayah Provinsi Jakarta, meliputi Jakarta Utara,
Jakarta Selatan, Jakarta Timur, Jakarta Barat, dan Jakarta Pusat. Unit
observasi utamanya bukan individu, melainkan TAZ beserta centroid,
jumlah penduduk, kepadatan, tarikan perjalanan, dan aksesibilitasnya
(hlm. 8-10, Bagian 3.1-3.3; Gambar 3-5).

\begin{itemize}
\tightlist
\item
  \textbf{Jumlah unit spasial:} Bagian 3.1 menyatakan bahwa agregasi
  beberapa tingkat administrasi menghasilkan 131 TAZ. Namun, Bagian 3.5
  menyatakan 133 TAZ dan 17.689 kombinasi OD, yakni 133 x 133 (hlm. 9,
  Bagian 3.1; hlm. 12-13, Bagian 3.5).
\item
  \textbf{Kelompok pendapatan:} Tinggi, menengah, dan rendah. Batas
  nominal pendapatan, metode klasifikasi rumah tangga, serta ukuran
  sampel per kelompok dalam data wawancara OD: Tidak disebutkan dalam
  file.
\item
  \textbf{Kelompok kepadatan untuk Theil:} Rendah, kurang dari 17.090
  jiwa/km2; menengah, 17.090-50.123 jiwa/km2; dan tinggi, sedikitnya
  50.124 jiwa/km2 (hlm. 15-16, Bagian 3.6.3).
\item
  \textbf{Cakupan moda:} Mobil, sepeda motor, dan transportasi publik.
  Berjalan kaki dan bersepeda muncul dalam ilustrasi rute, tetapi tidak
  termasuk moda yang dianalisis. Untuk transportasi publik, perjalanan
  akses dan keluar diasumsikan dilakukan dengan berjalan kaki (hlm.
  13-14, Bagian 3.5 dan Gambar 6).
\item
  \textbf{Sampel individu:} Jumlah responden, desain pengambilan sampel,
  tingkat respons, karakteristik demografis responden, dan bobot survei:
  Tidak disebutkan dalam file. Sumber hanya merujuk pada data wawancara
  OD dalam laporan JICA ketika menjelaskan proporsi perjalanan terkait
  pekerjaan (hlm. 9-10, Bagian 3.3).
\end{itemize}

\textbf{Inferensi pengolahan:} Ketidaksesuaian 131 dan 133 TAZ membuat
jumlah observasi spasial yang benar tidak dapat dipastikan hanya dari
file. Karena data dianalisis pada tingkat zona, hasil tidak dapat dibaca
sebagai estimasi pengalaman setiap individu di dalam zona tanpa asumsi
homogenitas tambahan.

\subsection{Dependent Variables}\label{dependent-variables-9}

\textbf{Laporan sumber:} File tidak menetapkan variabel dependen dalam
kerangka regresi atau kausal. Luaran utama yang dianalisis adalah
\textbf{aksesibilitas pekerjaan kumulatif} untuk zona asal \texttt{i},
moda \texttt{m}, aktivitas pekerjaan \texttt{o}, dan ambang waktu
\texttt{s}. Nilainya adalah jumlah proksi pekerjaan pada seluruh zona
tujuan yang dapat dicapai dalam waktu paling lama 41,2 menit (hlm. 11,
Bagian 3.4, Persamaan 1-2).

Luaran turunannya ialah:

\begin{itemize}
\tightlist
\item
  \textbf{Kelas aksesibilitas spasial:} Lima kelas Jenks, dari sangat
  rendah sampai sangat tinggi, untuk peta menurut zona, moda, dan
  tingkat pendapatan (hlm. 12, Bagian 3.4; hlm. 17, Gambar 8).
\item
  \textbf{Koefisien Gini dan kurva Lorenz:} Ukuran dan visualisasi
  konsentrasi distribusi aksesibilitas (hlm. 14, Bagian 3.6.1; hlm. 20,
  Tabel 2 dan Gambar 9).
\item
  \textbf{Rasio Palma:} Perbandingan rata-rata akses antara 10\% TAZ
  terpadat dan 40\% TAZ dengan kepadatan terendah (hlm. 14-15, Bagian
  3.6.2; hlm. 21, Tabel 3).
\item
  \textbf{Indeks Theil:} Ketimpangan total serta komponen di dalam dan
  di antara tiga kelompok kepadatan penduduk (hlm. 15-16, Bagian 3.6.3;
  hlm. 22, Tabel 4 dan Gambar 10).
\end{itemize}

\textbf{Inferensi pengolahan:} Moda, kelompok pendapatan, lokasi TAZ,
dan kelompok kepadatan berfungsi sebagai dimensi disagregasi atau
pembanding, bukan variabel penjelas dalam model statistik. Luaran juga
mengukur peluang kerja yang secara spasial dapat dijangkau menurut
proksi, bukan perolehan pekerjaan aktual, kecocokan
keterampilan-pekerjaan, upah, keterjangkauan biaya, atau probabilitas
diterima bekerja.

\subsection{Results}\label{results-9}

\subsubsection{Pola spasial}\label{pola-spasial}

\begin{itemize}
\tightlist
\item
  \textbf{Laporan sumber:} Akses pekerjaan tertinggi terpusat di inti
  kota, terutama Jakarta Pusat, Jakarta Selatan, dan sebagian Jakarta
  Barat, tempat klaster pekerjaan, pusat komersial, dan wilayah
  administrasi berada. Akses menurun ke arah utara, timur, dan bagian
  selatan yang periferal, khususnya pada zona dengan infrastruktur tidak
  memadai atau koneksi transportasi publik yang lemah (hlm. 16-18,
  Bagian 4.1.1; Gambar 8).
\item
  \textbf{Laporan sumber:} Sepeda motor memiliki cakupan spasial paling
  luas, termasuk ke kawasan semiurban atau permukiman informal yang
  lemah layanan transportasi publiknya. Akses transportasi publik
  terkonsentrasi pada sejumlah koridor BRT dan kereta komuter, kemudian
  menurun di luar jalur utama serta titik perpindahan. Akses dengan
  mobil terkonsentrasi pada kawasan dengan jaringan jalan berkualitas
  lebih baik dan kemacetan lebih rendah (hlm. 16-18, Bagian 4.1.1).
\end{itemize}

\subsubsection{Hasil menurut moda dan
pendapatan}\label{hasil-menurut-moda-dan-pendapatan}

\begin{itemize}
\tightlist
\item
  \textbf{Laporan sumber:} Bagian hasil menyatakan sepeda motor mencakup
  62,5\% seluruh pekerjaan yang dijangkau; kelompok berpendapatan
  menengah dan rendah secara bersama-sama dikatakan mengakses lebih dari
  7 juta pekerjaan melalui moda tersebut. Transportasi publik hanya
  mencakup 9,27\%, termasuk 692.321 atau 4,93\% pada kelompok menengah.
  Mobil berada di antara kedua moda tersebut; kelompok menengah mencatat
  2.163.569 perjalanan/akses pada moda mobil (hlm. 18, Bagian 4.1.2;
  hlm. 10, Tabel 1).
\item
  \textbf{Laporan sumber:} Kelompok berpendapatan menengah memiliki
  nilai semua moda tertinggi, yaitu 7.669.331 atau 54,59\%. Kelompok
  tinggi memiliki 2.608.244 atau 18,56\%, sedangkan kelompok rendah
  memiliki 3.772.286 atau 26,85\% (hlm. 10, Tabel 1; hlm. 18-19, Bagian
  4.1.3).
\item
  \textbf{Inferensi pengolahan:} Angka dan persentase tersebut identik
  dengan tabel yang berjudul jumlah serta proporsi \textbf{perjalanan},
  sementara narasi hasil menyebutnya sebagai \textbf{pekerjaan yang
  dijangkau}. File tidak menyediakan tabel terpisah yang merekonsiliasi
  kedua istilah atau menunjukkan transformasi numeriknya, sehingga angka
  absolut ini perlu dibaca dengan kehati-hatian konseptual.
\end{itemize}

\subsubsection{Koefisien Gini dan kurva
Lorenz}\label{koefisien-gini-dan-kurva-lorenz}

\begin{itemize}
\tightlist
\item
  \textbf{Laporan sumber:} Untuk semua tingkat pendapatan, Gini akses
  transportasi publik adalah 0,417, jauh di atas mobil sebesar 0,140 dan
  sepeda motor sebesar 0,109. Kurva Lorenz transportasi publik paling
  jauh dari garis kesetaraan, sedangkan kurva sepeda motor paling dekat
  (hlm. 19-20, Bagian 4.2.1; Tabel 2 dan Gambar 9).
\item
  \textbf{Laporan sumber:} Variasi Gini antarkelompok pendapatan sangat
  kecil. Pada transportasi publik, nilainya 0,424 untuk kelompok tinggi,
  0,414 untuk menengah, dan 0,421 untuk rendah. Pada mobil, nilainya
  berturut-turut 0,137, 0,142, dan 0,140; pada sepeda motor, 0,108,
  0,110, dan 0,109. Gini semua moda adalah 0,116 untuk kelompok tinggi
  serta 0,117 untuk kelompok menengah dan rendah (hlm. 20-21, Tabel 2
  dan Bagian 4.2.1).
\end{itemize}

\subsubsection{Rasio Palma}\label{rasio-palma}

\begin{itemize}
\tightlist
\item
  \textbf{Laporan sumber:} Untuk semua tingkat pendapatan, rasio Palma
  adalah 1,610 pada mobil, 1,442 pada sepeda motor, 7,361 pada
  transportasi publik, dan 1,508 pada seluruh moda. Penulis membaca
  nilai 7,361 sebagai ketimpangan ekstrem: zona pada kelompok 10\%
  teratas memiliki akses transportasi publik lebih dari tujuh kali
  kelompok 40\% terbawah menurut pemeringkatan kepadatan yang digunakan
  (hlm. 21, Bagian 4.2.2 dan Tabel 3).
\item
  \textbf{Laporan sumber:} Menurut pendapatan, rasio transportasi publik
  adalah 7,616 untuk kelompok tinggi, 7,233 untuk menengah, dan 7,486
  untuk rendah. Nilai menurut pendapatan juga relatif dekat pada mobil,
  yakni 1,584, 1,622, dan 1,609, serta pada sepeda motor, yakni 1,429,
  1,448, dan 1,442 (hlm. 21, Tabel 3).
\end{itemize}

\subsubsection{Indeks Theil}\label{indeks-theil}

\begin{itemize}
\tightlist
\item
  \textbf{Laporan sumber:} Transportasi publik memiliki indeks Theil
  total 0,275, terdiri atas komponen di dalam kelompok sebesar 0,265
  atau 96\% dan komponen antarkelompok sebesar 0,010 atau 4\%. Mobil
  memiliki Theil total 0,036, dengan 91\% di dalam kelompok dan 9\%
  antarkelompok. Sepeda motor memiliki Theil total 0,019, dengan 92\% di
  dalam kelompok dan 8\% antarkelompok (hlm. 22, Bagian 4.2.3 dan Tabel
  4).
\item
  \textbf{Laporan sumber:} Dominasi komponen di dalam kelompok berarti
  variasi besar tetap muncul di antara zona-zona yang berada dalam kelas
  kepadatan yang sama. Beberapa lingkungan padat pun tetap memiliki
  konektivitas atau infrastruktur yang tidak memadai (hlm. 22-23, Bagian
  4.2.3).
\end{itemize}

\textbf{Inferensi pengolahan:} Walaupun abstrak menggunakan ungkapan
akses yang ``significantly lower'', file tidak melaporkan nilai p,
interval kepercayaan, galat baku, atau uji signifikansi. Istilah
tersebut karena itu tidak dapat diverifikasi sebagai signifikansi
statistik dari informasi yang tersedia (hlm. 1, Abstract; hlm. 16-23,
Bagian 4).

\subsection{Findings}\label{findings-9}

\begin{itemize}
\tightlist
\item
  \textbf{Laporan sumber:} Dimensi moda memperlihatkan perbedaan paling
  jelas: sepeda motor merupakan moda dengan distribusi akses paling
  merata, mobil berada di posisi antara, dan transportasi publik
  merupakan moda dengan distribusi paling timpang menurut ketiga indeks
  (hlm. 19-22, Bagian 4.2; Tabel 2-4).
\item
  \textbf{Laporan sumber:} Dimensi spasial ditandai oleh pemusatan akses
  di inti Jakarta dan keterbatasan akses di kawasan periferal atau
  berkepadatan rendah. Namun, Theil memperlihatkan bahwa kategori
  pusat-pinggiran atau padat-renggang saja tidak memadai karena sebagian
  besar variasi berada di dalam kelompok kepadatan (hlm. 16-17, Bagian
  4.1.1; hlm. 22-23, Bagian 4.2.3).
\item
  \textbf{Interpretasi penulis:} Sepeda motor berperan sebagai penyetara
  akses karena efisien, relatif terjangkau, fleksibel di tengah
  kemacetan, dan mampu mengisi kekosongan layanan formal. Penulis
  sekaligus menilai manfaat tersebut memiliki biaya berupa emisi, risiko
  keselamatan, perilaku lalu lintas yang tidak teratur, serta persoalan
  keberlanjutan dan regulasi (hlm. 18, Bagian 4.1.2).
\item
  \textbf{Interpretasi penulis:} Ketimpangan transportasi publik
  dikaitkan dengan cakupan geografis terbatas, integrasi lemah,
  frekuensi rendah, kurangnya sistem pengumpan, dan konektivitas mil
  pertama serta terakhir yang buruk. Karena Gini dan Palma hampir sama
  menurut pendapatan, penulis menyimpulkan bahwa pola eksklusi lebih
  dibentuk oleh geografi dan rancangan jaringan daripada kelas
  pendapatan saja (hlm. 17-18, Bagian 4.1.1-4.1.2; hlm. 20-22, Bagian
  4.2.1-4.2.2).
\item
  \textbf{Interpretasi penulis:} Akses tertinggi kelompok menengah
  dijelaskan dengan fleksibilitas moda dan pilihan bermukim di kawasan
  lama berkepadatan menengah yang dekat dengan pusat pekerjaan. Akses
  lebih rendah kelompok berpendapatan tinggi dijelaskan dengan lokasi
  hunian eksklusif berkepadatan rendah yang cenderung periferal,
  sedangkan keterbatasan kelompok rendah dijelaskan dengan kendala
  finansial, infrastruktur, dan ketergantungan pada moda murah atau
  informal (hlm. 18-19, Bagian 4.1.3).
\item
  \textbf{Inferensi pengolahan:} Penjelasan mengenai pilihan hunian,
  keterjangkauan, frekuensi layanan, atau penyebab akses menurut
  pendapatan tidak diuji langsung dengan model yang dilaporkan.
  Penjelasan tersebut harus diperlakukan sebagai interpretasi penulis
  yang didukung pembacaan spasial dan literatur, bukan sebagai mekanisme
  yang telah diisolasi secara empiris oleh penelitian ini.
\end{itemize}

\subsection{Conclusions}\label{conclusions-9}

\textbf{Laporan sumber:} Penulis menyimpulkan bahwa aksesibilitas
pekerjaan di Jakarta sangat terpusat secara spasial dan bahwa moda
transportasi merupakan pembeda penting. Sepeda motor menghasilkan akses
terluas dan paling merata bagi semua kelompok, terutama kelompok
menengah dan rendah, tetapi ketergantungan tersebut mencerminkan
kelemahan sistem formal serta memunculkan persoalan keselamatan,
keberlanjutan, dan regulasi (hlm. 23-24, Bagian 5).

\textbf{Laporan sumber:} Transportasi publik disimpulkan sebagai moda
paling timpang, dengan Gini 0,417, Palma 7,361, dan Theil 0,275. Lebih
dari 90\% ketimpangan setiap moda terletak di dalam kelompok kepadatan,
sehingga kekurangan layanan pada skala lingkungan dinilai sama
pentingnya dengan struktur kota berskala besar (hlm. 22-24, Tabel 4 dan
Bagian 5).

\textbf{Interpretasi penulis:} Kesamaan Gini menurut kelompok pendapatan
mendukung jawaban bahwa ketimpangan akses Jakarta lebih kuat ditandai
oleh lokasi geografis dan moda daripada kategori sosial-ekonomi semata.
Penulis tidak meniadakan kerugian kelompok berpendapatan rendah, tetapi
menolak penjelasan berbasis pendapatan sebagai satu-satunya atau faktor
dominan (hlm. 20-24, Bagian 4.2.1 dan 5).

\textbf{Inferensi pengolahan:} Kesimpulan komparatif mengenai distribusi
akses konsisten dengan indeks yang dilaporkan. Namun, klaim tentang
faktor yang ``lebih berpengaruh'' sebaiknya dibaca sebagai interpretasi
pola karena desain tidak mengestimasi besaran pengaruh relatif atau
mengendalikan faktor perancu.

\subsection{Contributions}\label{contributions-9}

\textbf{Laporan sumber:} Penulis menyatakan tiga kontribusi utama (hlm.
23, Bagian 5):

\begin{itemize}
\tightlist
\item
  Menawarkan kerangka yang dapat direplikasi untuk mengevaluasi
  aksesibilitas pekerjaan di Global South dengan metrik keadilan formal
  dan multidimensi, serta mengklaimnya sebagai salah satu penerapan awal
  semacam itu pada megakota Asia Tenggara.
\item
  Mengintegrasikan dimensi spasial, sosial-ekonomi, dan moda dalam satu
  pendekatan analitis, bukan menilai masing-masing secara terpisah.
\item
  Menonjolkan pola mobilitas kontekstual yang kerap kurang diperhatikan
  dalam riset aksesibilitas global, khususnya ketergantungan pada sepeda
  motor dan fragmentasi jaringan transit.
\end{itemize}

\textbf{Laporan sumber:} Secara empiris, studi menyediakan perbandingan
Gini dan Palma menurut moda serta pendapatan, lalu Theil menurut moda
dan kelompok kepadatan pada tingkat TAZ di Jakarta. Kombinasi indeks
memungkinkan pembacaan ketimpangan keseluruhan, perbedaan bagian
ekstrem, dan sumber variasi di dalam atau di antara kelompok kepadatan
(hlm. 14-16, Bagian 3.6; hlm. 19-23, Bagian 4.2).

\textbf{Interpretasi penulis:} Temuan diposisikan sebagai dukungan untuk
memasukkan metrik keadilan multidimensi ke dalam praktik perencanaan dan
keputusan investasi transportasi, terutama pada kota Global South dengan
campuran moda yang tidak sepenuhnya diwakili sistem angkutan formal
(hlm. 1, Abstract; hlm. 23-24, Bagian 5).

\textbf{Inferensi pengolahan:} Nilai tambah paling jelas berada pada
integrasi deskriptif tiga lensa dan tiga indeks dalam satu kasus. Klaim
kebaruan sebagai ``salah satu yang pertama'' adalah klaim penulis; file
tidak menyajikan prosedur tinjauan sistematis yang memungkinkan
verifikasi menyeluruh atas prioritas tersebut.

\subsection{Research Gap}\label{research-gap-9}

\textbf{Kesenjangan yang dinyatakan sumber:}

\begin{itemize}
\tightlist
\item
  Penelitian akses pekerjaan di Global South, khususnya Asia Tenggara,
  masih relatif terbatas dibandingkan studi pada kota-kota di wilayah
  lain (hlm. 2-3, Bagian 1; hlm. 7, Bagian 2.2).
\item
  Sistem mobilitas Asia Tenggara memiliki campuran moda khas, termasuk
  dominasi sepeda motor, angkutan informal, dan \emph{ride-hailing},
  tetapi karakteristik tersebut belum banyak diintegrasikan ke evaluasi
  aksesibilitas pekerjaan (hlm. 3, Bagian 1; hlm. 7, Bagian 2.2).
\item
  Hanya sedikit studi yang menggunakan indikator keadilan multidimensi
  untuk menangkap ketimpangan moda dan sosial-ekonomi sekaligus. Studi
  Indonesia terdahulu yang disebut penulis menilai akses pekerjaan dalam
  konteks jalan tol tanpa memasukkan pilihan moda (hlm. 3, Bagian 1).
\item
  Dampak perluasan BRT, MRT, dan LRT Jakarta terhadap akses pekerjaan
  dan pengurangan ketimpangan sosial-ekonomi belum dianalisis dalam
  literatur yang dibahas file (hlm. 3, Bagian 1).
\item
  Gini dan Lorenz saja tidak mengidentifikasi kelompok tertentu yang
  dirugikan atau arah perubahan, sehingga dibutuhkan pelengkap seperti
  Palma dan Theil (hlm. 7, Bagian 2.3).
\end{itemize}

\textbf{Inferensi pengolahan:} Studi menutup sebagian kesenjangan
pengukuran lintas moda dan kelompok, tetapi belum menutup kesenjangan
evaluasi dampak perluasan transportasi massal karena desainnya tidak
membandingkan kondisi sebelum dan sesudah intervensi. Studi juga belum
mengukur pengalaman individual atau perubahan dari waktu ke waktu.

\subsection{Limitations}\label{limitations-9}

\textbf{Keterbatasan yang dinyatakan penulis:}

\begin{itemize}
\tightlist
\item
  Ambang kumulatif 41,2 menit merupakan batas kaku yang menyederhanakan
  penurunan bertahap dalam kesediaan melakukan perjalanan. Toleransi
  waktu sebenarnya dapat berbeda menurut pendapatan, tujuan perjalanan,
  dan ketersediaan moda (hlm. 12, Bagian 3.4).
\item
  Fungsi biner tidak menangkap mutu pengalaman perjalanan seperti
  keandalan, kenyamanan, dan keselamatan (hlm. 12, Bagian 3.4).
\item
  Ambang yang digunakan tidak menangkap variasi waktu dalam sehari atau
  perubahan kemacetan, sehingga akses aktual dapat dinilai terlalu
  rendah atau terlalu tinggi pada jam sibuk (hlm. 12, Bagian 3.4).
\item
  Data pendapatan tingkat zona tidak tersedia. Karena itu, rasio Palma
  diadaptasi dengan mengurutkan TAZ berdasarkan kepadatan penduduk,
  bukan pendapatan (hlm. 14-15, Bagian 3.6.2).
\item
  Akses dan keluar dari transportasi publik diasumsikan dilakukan dengan
  berjalan kaki karena keterbatasan data Google API, meskipun praktiknya
  dapat melibatkan sepeda atau \emph{ride-hailing} (hlm. 13-14, Bagian
  3.5).
\end{itemize}

\textbf{Keterbatasan yang teridentifikasi melalui pengolahan file:}

\begin{itemize}
\tightlist
\item
  \textbf{Jumlah zona tidak konsisten:} Satu bagian menyebut 131 TAZ,
  sedangkan bagian lain menggunakan 133 TAZ dan 17.689 pasangan OD.
  Tidak ada penjelasan mengenai dua zona yang berbeda atau angka mana
  yang dipakai dalam hasil akhir (hlm. 9, Bagian 3.1; hlm. 12-13, Bagian
  3.5).
\item
  \textbf{Dasar ambang tidak sepenuhnya konsisten:} Bagian 3.4
  mengaitkan 41,2 menit dengan rata-rata komuter dari survei mobilitas
  perkotaan 2019 yang dikonfirmasi statistik Google Maps, sedangkan
  Bagian 4.1 menyatakan ambang diperoleh dengan merata-ratakan data
  waktu perjalanan semua moda. Hubungan antara dua penjelasan tersebut
  tidak dirinci (hlm. 11, Bagian 3.4; hlm. 16, Bagian 4.1).
\item
  \textbf{Proksi pekerjaan:} Pekerjaan diperkirakan dari tarikan
  perjalanan, bukan dari inventaris pekerjaan aktual. Formula yang
  ditampilkan tidak membedakan mutu pekerjaan, upah, kecocokan
  keterampilan, persaingan memperoleh pekerjaan, atau kapasitas
  kesempatan (hlm. 10-11, Bagian 3.3-3.4, Persamaan 1-2).
\item
  \textbf{Ambiguitas perjalanan dan pekerjaan:} Tabel 1 berjudul jumlah
  perjalanan, tetapi Bagian 4.1.2-4.1.3 menggunakan angka yang sama
  sebagai jumlah atau proporsi pekerjaan yang dijangkau. Tahap
  transformasi antara kedua konstruk tidak diperlihatkan secara terpisah
  (hlm. 10, Tabel 1; hlm. 18-19, Bagian 4.1.2-4.1.3).
\item
  \textbf{Disagregasi pendapatan tidak sepenuhnya terdokumentasi:} Rumus
  aksesibilitas mengindeks zona, moda, aktivitas, dan ambang waktu,
  tetapi tidak menampilkan indeks pendapatan meskipun hasil dilaporkan
  untuk tiga kelompok pendapatan. Cara memasukkan kelompok pendapatan ke
  perhitungan setelah penggunaan matriks OD tidak dijelaskan secara
  rinci (hlm. 10-12, Bagian 3.3-3.4; Persamaan 1-2).
\item
  \textbf{Representasi waktu tunggal:} Waktu tempuh dikumpulkan hanya
  pada 2 November 2020 pukul 10.00. Penulis menyebut waktu itu mewakili
  kondisi pertengahan pagi yang normal, tetapi analisis sensitivitas
  lintas hari, jam sibuk, musim, atau gangguan jaringan tidak dilaporkan
  (hlm. 12-14, Bagian 3.5).
\item
  \textbf{Ketidakselarasan waktu data:} Matriks OD dan ambang perjalanan
  merujuk pada 2019, sedangkan waktu tempuh API berasal dari 2020.
  Prosedur harmonisasi atau pembaruan data: Tidak disebutkan dalam file.
\item
  \textbf{Definisi kelompok pendapatan:} Batas pendapatan dan metode
  penetapan kategori tinggi, menengah, serta rendah tidak tersedia,
  sehingga reproduksi disagregasi sosial-ekonomi tidak dapat dilakukan
  dari file saja (hlm. 10, Tabel 1).
\item
  \textbf{Agregasi zona:} Penggunaan nilai tingkat TAZ tidak
  memperlihatkan variasi pengalaman di dalam setiap zona. Penafsiran
  terhadap individu berisiko melampaui resolusi data, khususnya ketika
  penulis menjelaskan pilihan hunian atau kendala rumah tangga.
\item
  \textbf{Terminologi unit spasial:} Uraian rumus Theil mendefinisikan
  observasi sebagai \emph{grid}, sedangkan bagian wilayah studi dan
  perhitungan akses menggunakan TAZ. Apakah \emph{grid} hanya istilah
  lain untuk TAZ atau unit yang berbeda: Tidak disebutkan dalam file
  (hlm. 15-16, Bagian 3.6.3).
\item
  \textbf{Kekuatan inferensi:} Tidak ada uji statistik inferensial atau
  model kausal. Oleh sebab itu, kemiripan indeks antarkelompok
  pendapatan mendukung deskripsi pola, tetapi tidak cukup untuk
  menetapkan urutan pengaruh kausal antara pendapatan, geografi, dan
  moda.
\item
  \textbf{Dari ketimpangan ke ketidakadilan:} Bagian teori membedakan
  ketimpangan statistik dari ketidakadilan normatif, tetapi
  operasionalisasi utamanya memakai indeks distribusi. Kriteria etis,
  ambang keadilan, kebutuhan minimum, atau bobot normatif untuk
  menyatakan suatu perbedaan tidak adil: Tidak disebutkan dalam file
  (hlm. 3, Bagian 1; hlm. 4-7, Bagian 2).
\item
  \textbf{Replikasi:} Data mentah, kode Python, rincian kalibrasi atau
  validasi model mikrosimulasi, perlakuan terhadap rute API yang hilang,
  dan parameter lengkap klasifikasi Jenks tidak disediakan dalam file.
\end{itemize}

\subsection{Challenges}\label{challenges-9}

\begin{itemize}
\tightlist
\item
  \textbf{Laporan sumber:} Pengambilan waktu tempuh untuk seluruh
  pasangan OD menghasilkan 17.689 kombinasi dan memerlukan otomatisasi
  Python, pemrosesan dalam kelompok, pembatasan laju permintaan, serta
  penanganan galat agar sesuai dengan kuota API (hlm. 12-13, Bagian
  3.5).
\item
  \textbf{Laporan sumber:} Perbandingan tiga moda membutuhkan parameter
  API yang berbeda dan asumsi khusus untuk akses serta keluar dari
  transportasi publik. Penggabungan respons JSON dengan matriks OD dan
  distribusi pekerjaan juga diperlukan sebelum perhitungan aksesibilitas
  (hlm. 12-14, Bagian 3.5).
\item
  \textbf{Inferensi pengolahan:} Tantangan validitas konstruk muncul
  karena tarikan perjalanan dipakai sebagai proksi pekerjaan dan angka
  perjalanan kemudian dibahas sebagai akses pekerjaan. Dokumentasi
  transformasi yang lebih eksplisit diperlukan untuk membedakan
  permintaan perjalanan, jumlah pekerjaan, dan pekerjaan yang dapat
  dijangkau.
\item
  \textbf{Inferensi pengolahan:} Tantangan komparabilitas timbul dari
  penggunaan satu ambang waktu untuk semua moda dan kelompok, padahal
  penulis sendiri mengakui toleransi waktu dapat berbeda. Nilai yang
  setara secara matematis belum tentu mewakili pengalaman, biaya,
  kenyamanan, atau risiko yang setara.
\item
  \textbf{Inferensi pengolahan:} Tantangan interpretasi terletak pada
  perubahan dari ukuran ketimpangan distribusional menjadi klaim
  ketidakadilan dan mekanisme penyebab. Metrik dapat menunjukkan lokasi
  konsentrasi, tetapi file tidak menyediakan patokan normatif atau
  rancangan identifikasi mekanisme.
\item
  \textbf{Informasi yang tidak tersedia:} Jumlah permintaan API yang
  gagal, proporsi OD tanpa rute transportasi publik, prosedur imputasi,
  pemeriksaan mutu respons, dan validasi hasil terhadap waktu tempuh
  observasi: Tidak disebutkan dalam file.
\end{itemize}

\subsection{Practical Implications}\label{practical-implications-9}

\textbf{Rekomendasi yang dinyatakan penulis:}

\begin{itemize}
\tightlist
\item
  Memperluas jaringan transportasi publik massal yang inklusif dan
  terintegrasi serta memperpanjang cakupan layanan menuju distrik yang
  kurang terlayani (hlm. 1, Abstract; hlm. 23-24, Bagian 5).
\item
  Memperbaiki konektivitas mil pertama dan terakhir melalui layanan
  pengumpan, akses pejalan kaki ke stasiun, dan infrastruktur
  mikromobilitas (hlm. 22-24, Bagian 4.2.3 dan 5).
\item
  Tidak hanya memperluas jaringan pada skala kota atau koridor, tetapi
  juga memperbaiki integrasi stasiun, frekuensi layanan, dan perpindahan
  antarmoda di lingkungan yang secara formal telah dilayani tetapi masih
  menghasilkan akses tidak merata (hlm. 22-23, Bagian 4.2.3).
\item
  Menargetkan perbaikan pada tingkat distrik, subdistrik, atau
  lingkungan karena ketimpangan terbesar berada di dalam kelompok
  kepadatan, bukan sekadar antara kawasan padat dan renggang (hlm.
  22-24, Bagian 4.2.3 dan 5).
\item
  Memasukkan metrik keadilan ke keputusan perencanaan dan investasi agar
  ketimpangan intraurban dapat dikenali serta manfaat peningkatan akses
  dapat dibagikan lebih merata (hlm. 24, Bagian 5).
\item
  Mengakui sepeda motor sebagai solusi adaptif jangka dekat, tetapi
  mengarahkan perencanaan jangka panjang pada moda publik yang lebih
  aman, rendah emisi, andal, dan berkelanjutan (hlm. 23-24, Bagian 5).
\end{itemize}

\textbf{Inferensi pengolahan:} Prioritas lokasi dapat diturunkan dari
TAZ dengan akses rendah dan komponen ketimpangan di dalam kelompok yang
tinggi. Akan tetapi, file tidak melakukan evaluasi biaya, kelayakan
proyek, dampak emisi, atau simulasi manfaat dari rekomendasi tertentu,
sehingga rekomendasi belum merupakan peringkat investasi yang telah
diuji.

\subsection{Applications}\label{applications-9}

\begin{itemize}
\tightlist
\item
  \textbf{Laporan sumber:} Kerangka kumulatif-Gini-Palma-Theil dapat
  digunakan untuk membandingkan distribusi akses pekerjaan menurut moda,
  kelompok pendapatan, dan kelompok kepadatan dalam satu wilayah
  perkotaan (hlm. 7-8, Gambar 2; hlm. 23, Bagian 5).
\item
  \textbf{Laporan sumber:} Peta akses tingkat TAZ dapat digunakan untuk
  mengenali distrik atau lingkungan yang tertinggal dan membedakan
  apakah keterbatasannya terutama tampak pada jaringan transportasi
  publik, mobil, atau sepeda motor (hlm. 16-18, Bagian 4.1 dan Gambar
  8).
\item
  \textbf{Laporan sumber:} Dekomposisi Theil dapat membantu perencana
  memutuskan apakah kebijakan sebaiknya diarahkan pada perbedaan
  antarkategori wilayah atau pada variasi internal di dalam kategori
  yang sama (hlm. 22-23, Bagian 4.2.3 dan Tabel 4).
\item
  \textbf{Laporan sumber:} Penulis memosisikan pendekatan ini sebagai
  kerangka yang dapat direplikasi di kota-kota Global South dan sebagai
  dasar memasukkan pemerataan akses ke perencanaan serta investasi
  transportasi (hlm. 23-24, Bagian 5).
\item
  \textbf{Inferensi pengolahan:} Dengan data lintas waktu yang
  sebanding, kerangka dapat menjadi alat pemantauan sebelum-sesudah
  investasi MRT, LRT, atau BRT. Aplikasi longitudinal ini diusulkan
  penulis sebagai penelitian mendatang dan belum dilakukan dalam studi
  J10 (hlm. 24, Bagian 5).
\end{itemize}

\subsection{Future Research
(Optional)}\label{future-research-optional-9}

\textbf{Laporan sumber:} Penulis mengusulkan agenda penelitian lanjutan
berikut (hlm. 24, Bagian 5):

\begin{itemize}
\tightlist
\item
  Meneliti perubahan akses pekerjaan dan keadilan transportasi dari
  waktu ke waktu, terutama sebagai respons atas investasi MRT, LRT, dan
  BRT Jakarta.
\item
  Membandingkan kondisi sebelum dan sesudah intervensi untuk menguji
  apakah investasi infrastruktur benar-benar mengurangi ketimpangan.
\item
  Memasukkan variabel sosial-demografis yang lebih rinci, seperti
  gender, usia, jenis pekerjaan, dan status disabilitas, agar pengalaman
  eksklusi kelompok yang berbeda dapat dinilai.
\item
  Mengembangkan model aksesibilitas dinamis yang menangkap variasi
  layanan secara waktu nyata.
\item
  Mengkaji hubungan keterjangkauan perumahan, tata guna lahan, dan
  aksesibilitas transportasi, termasuk dampak keadilan dari regenerasi
  perkotaan dan strategi pembangunan berorientasi transit di Jakarta.
\end{itemize}

\subsection{Provenance Notes}\label{provenance-notes-9}

\begin{itemize}
\tightlist
\item
  Review ini disusun hanya dari seluruh isi file TXT
  \texttt{source/journal/J10-andani-qamilla-izdihar-safira-sakti-syabri-2025-spatial-mobility-or-socio-economic-inequity-a-district-level-job-accessibility-evaluation-in-jakarta-indonesia-10.1080-12265934.2025.2553714.txt},
  mulai dari blok metadata PDF sampai akhir daftar pustaka. Tidak ada
  sumber eksternal, Zotero, PDF lain, atau pengetahuan di luar file yang
  digunakan untuk menambah atau memverifikasi klaim.
\item
  Blok metadata menyatakan bahwa TXT diekstrak pada 31 Agustus 2026
  dengan \texttt{pdftotext\ -layout} dari PDF sumber 30 halaman. Dengan
  demikian, basis akses review adalah teks penuh hasil ekstraksi, bukan
  inspeksi visual langsung terhadap PDF asli.
\item
  Locator halaman mengikuti nomor halaman artikel yang tampak pada teks
  hasil ekstraksi. Halaman artikel pertama yang memuat abstrak tidak
  menampilkan angka pada teks, tetapi dirujuk sebagai hlm. 1 mengikuti
  urutan sebelum header hlm. 2. Halaman muka penerbit dibedakan sebagai
  ``halaman muka hasil ekstraksi''.
\item
  Isi peta dan ilustrasi tidak diverifikasi secara visual. Uraian
  mengenai pola pada Gambar 1-10 hanya mengikuti keterangan gambar dan
  narasi penulis yang tersedia dalam TXT. Nilai numerik tabel
  dipertahankan dari hasil ekstraksi Tabel 1-4.
\item
  Simbol beberapa persamaan mengalami distorsi tata letak saat
  diekstrak. Review karena itu menjelaskan fungsi persamaan berdasarkan
  uraian tekstual penulis dan tidak merekonstruksi notasi yang tidak
  terbaca.
\item
  Penanda \textbf{Laporan sumber} merujuk pada data, metode, hasil, atau
  pernyataan eksplisit artikel. Penanda \textbf{Interpretasi penulis}
  merujuk pada makna atau mekanisme yang diajukan penulis atas hasil.
  Penanda \textbf{Inferensi pengolahan} merujuk pada pembacaan kritis
  yang diturunkan dari isi file, bukan klaim eksplisit artikel atau
  posisi pengguna.
\item
  Ketidakkonsistenan 131 dibandingkan 133 TAZ serta pergeseran istilah
  dari jumlah perjalanan ke pekerjaan yang dijangkau dipertahankan
  sebagai batas provenance; angka yang bertentangan tidak diselaraskan
  melalui dugaan.
\item
  File menyatakan bahwa penulis menggunakan Quillbot untuk membantu
  parafrasa dan penyempurnaan bahasa pada beberapa bagian, lalu meninjau
  serta menyunting hasilnya dan bertanggung jawab atas isi akhir (hlm.
  24, Declaration of generative AI and AI-assisted technologies in the
  writing process).
\item
  Status penelaahan sejawat, data mentah, kode analisis, volume, nomor
  terbitan, serta rentang halaman bibliografis tidak dinyatakan secara
  eksplisit dalam file dan tidak diinferensikan dari luar.
\end{itemize}

\section{Review J11}\label{review-j11}

\textbf{Identitas sumber}

\begin{itemize}
\tightlist
\item
  Kode: J11.
\item
  Judul: \emph{Job Access and Spatial Equity of a Toll Road}.
\item
  Penulis: I Gusti Ayu Andani, Lissy La Paix, Shanty Rachmat, Ibnu
  Syabri, dan Karst Geurs.
\item
  Tahun yang terindikasi: 2021. Tahun ini hanya tercantum dalam nama
  file TXT dan nama PDF pada blok metadata ekstraksi; keterangan tahun
  publikasi tidak tampak pada tubuh bab.
\item
  Jenis dokumen menurut isi: bab 13. Isi berulang kali menyebut dokumen
  ini sebagai \emph{chapter}. Status sebagai artikel jurnal dan status
  telaah sejawat: Tidak disebutkan dalam file.
\item
  Wadah publikasi: header halaman berbunyi \emph{applications of
  access}, tetapi rincian judul lengkap wadah, editor, penerbit, tempat
  terbit, ISBN, dan edisi: Tidak disebutkan dalam file.
\item
  Rentang yang tersedia: metadata PDF menyebut 22 halaman. Header cetak
  yang terlihat dalam ekstraksi mencakup halaman 240-260; nomor halaman
  pembuka tidak tercetak dalam teks hasil ekstraksi.
\item
  DOI: Tidak disebutkan dalam file.
\item
  Objek studi: dampak Jalan Tol Cipularang, yang dibangun pada 2005,
  terhadap akses pekerjaan dan keadilan spasial di koridor
  Jakarta-Bandung, Indonesia (bagian Abstract dan Introduction, halaman
  pembuka bab serta hlm. 240).
\end{itemize}

\textbf{Penanda epistemik yang digunakan}

\begin{itemize}
\tightlist
\item
  \textbf{Laporan sumber} berarti pernyataan, metode, angka, atau hasil
  yang disampaikan langsung oleh bab.
\item
  \textbf{Interpretasi penulis} berarti penjelasan sebab, makna, atau
  implikasi yang diberikan oleh penulis bab terhadap hasilnya.
\item
  \textbf{Inferensi review} berarti pembacaan analitis yang dibuat hanya
  dari TXT J11 dan bukan klaim eksplisit penulis. Inferensi semacam ini
  tidak diperlakukan sebagai bukti tambahan.
\end{itemize}

\subsection{Summarized Abstract}\label{summarized-abstract-10}

\textbf{Laporan sumber.} Bab ini mengevaluasi secara \emph{ex post}
dampak Jalan Tol Cipularang terhadap akses pekerjaan dan keadilan
spasial di koridor Jakarta-Bandung. Penulis membangun skenario dengan
dan tanpa jalan tol menggunakan jaringan ArcGIS dan ekstensi model
permintaan transportasi untuk memperoleh estimasi waktu tempuh serta
biaya tergeneralisasi. Akses diukur dengan ukuran akses pekerjaan
potensial dan indeks Shen, sedangkan keadilan spasial dinilai dengan
koefisien Gini, rasio Palma, dan analisis klaster dua tahap (bagian
Abstract; Methods, hlm. 241-248).

\textbf{Laporan sumber.} Pembangunan jalan tol dilaporkan menurunkan
rata-rata waktu tempuh tertimbang di seluruh wilayah sebesar 13\% dan
meningkatkan akses pekerjaan potensial sekitar 5\%. Namun, jalan tol
juga memperbesar jangkauan pekerja terhadap pekerjaan sehingga kompetisi
pekerjaan meningkat, terutama di wilayah miskin pekerjaan di antara
Jakarta dan Bandung. Akibatnya, jumlah pekerjaan yang dapat diakses per
pekerja sedikit menurun. Koefisien Gini dan rasio Palma secara agregat
hampir tidak berubah, walaupun analisis klaster menemukan dampak lokal
yang lebih kuat di wilayah dekat gerbang tol (bagian Abstract; Findings,
hlm. 251-257; Conclusions, hlm. 259).

\textbf{Inferensi review.} Kesimpulan tentang ``tidak ada dampak'' pada
keadilan spasial berlaku untuk perubahan rata-rata pada dua indeks
agregat yang digunakan. Kesimpulan tersebut tidak meniadakan
heterogenitas lokal karena hasil klaster justru menunjukkan penurunan
indeks Shen yang lebih besar pada kelompok distrik dekat gerbang tol
(Table 13.6, hlm. 256; Table 13.7 dan uraian klaster, hlm. 256-257).

\subsection{Problem Statement}\label{problem-statement-10}

\textbf{Laporan sumber.} Ukuran akses yang memperhitungkan kualitas
infrastruktur transportasi dan distribusi spasial aktivitas sering
digunakan dalam analisis keadilan. Akan tetapi, penerapannya dalam
evaluasi \emph{ex post} investasi jalan dinyatakan masih jarang,
khususnya di negara berkembang (bagian Introduction, halaman pembuka
bab).

\textbf{Laporan sumber.} Jalan tol menawarkan rute yang lebih cepat
tetapi berbayar dibandingkan rute yang lebih lambat dan tidak berbayar.
Perubahan waktu dan biaya dapat terdistribusi tidak merata menurut
kondisi sosial ekonomi, bahkan dapat merugikan kelompok tertentu. Bab
memberi contoh dari literatur bahwa komuter berpendapatan tinggi dapat
lebih mungkin memperoleh manfaat, sementara komuter berpendapatan rendah
memilih rute tanpa tol (bagian Introduction, hlm. 240). Pernyataan
contoh tersebut merupakan laporan bab atas penelitian terdahulu, bukan
hasil empiris J11.

\textbf{Laporan sumber.} Dalam konteks Indonesia, evaluasi \emph{ex
post} investasi infrastruktur transportasi dinilai belum mendapat banyak
perhatian. Penulis menempatkan evaluasi \emph{ex post} sebagai bagian
yang sering diabaikan dalam siklus kebijakan, padahal evaluasi tersebut
dapat memperbaiki penilaian proyek berikutnya dan meningkatkan
akuntabilitas. Bab juga menyebut keterbatasan pendanaan khusus dan data
relevan sebagai alasan umum lemahnya pelaksanaan evaluasi \emph{ex post}
(bagian Introduction, hlm. 240-241).

\textbf{Rumusan masalah sumber.} Pertanyaan yang diajukan adalah dampak
jalan tol terhadap akses pekerjaan dan keadilan spasial. Penulis secara
khusus ingin membedakan pengaruh distribusi pekerjaan dan populasi
pekerja serta melihat bagaimana dampak akses tersebar antardistrik,
bukan hanya menghitung perubahan rata-rata wilayah (bagian Questions,
hlm. 241).

\subsection{Objectives}\label{objectives-10}

\begin{itemize}
\tightlist
\item
  \textbf{Laporan sumber:} mengevaluasi secara \emph{ex post} dampak
  Jalan Tol Cipularang terhadap akses pekerjaan dan keadilan spasial di
  koridor Jakarta-Bandung (bagian Introduction dan Questions, hlm.
  240-241).
\item
  \textbf{Laporan sumber:} mengisolasi pengaruh keberadaan jalan tol
  melalui perbandingan skenario jaringan dengan jalan tol dan tanpa
  jalan tol (bagian Methods, hlm. 241-242; Conclusions, hlm. 257).
\item
  \textbf{Laporan sumber:} mengestimasi dampak transportasi berupa arus
  lalu lintas, waktu tempuh, dan biaya tergeneralisasi dengan memasukkan
  hambatan kapasitas atau efek kemacetan (bagian Methods, hlm. 241-243).
\item
  \textbf{Laporan sumber:} mengukur akses pekerjaan dengan ukuran akses
  potensial dan indeks Shen agar distribusi pekerjaan serta kompetisi
  dari populasi pekerja sama-sama diperhitungkan (bagian Access
  measurements, hlm. 242-246).
\item
  \textbf{Laporan sumber:} menilai distribusi keadilan spasial melalui
  koefisien Gini, rasio Palma, dan klaster dua tahap untuk mengenali
  distrik yang paling terdampak beserta karakteristiknya (bagian
  Measuring spatial equity, hlm. 246-248).
\end{itemize}

\subsection{Literature Survey}\label{literature-survey-10}

\textbf{Sifat survei.} Bab menyajikan tinjauan naratif untuk membangun
konsep dan memilih metode. Protokol pencarian, basis data bibliografis,
kriteria inklusi, jumlah studi yang disaring, dan prosedur penilaian
mutu literatur: Tidak disebutkan dalam file. Karena itu, bagian ini
tidak dapat diperlakukan sebagai tinjauan sistematis.

\begin{itemize}
\tightlist
\item
  \textbf{Laporan sumber atas literatur akses:} Hansen (1959) menjadi
  dasar definisi akses sebagai potensi peluang interaksi dan rumus akses
  potensial. Geurs dan Van Wee (2004) dirujuk untuk empat komponen
  akses, yaitu transportasi, tata guna lahan, temporal, dan individual,
  serta empat kategori ukuran akses, yaitu berbasis infrastruktur,
  lokasi, orang, dan utilitas (bagian Access measurements, hlm.
  242-243).
\item
  \textbf{Laporan sumber atas literatur evaluasi jalan:} ukuran berbasis
  infrastruktur menilai kinerja jaringan, kemacetan, dan kecepatan. Bab
  menekankan bahwa estimasi waktu tempuh tanpa hambatan kapasitas serta
  arus lalu lintas dapat meremehkan waktu perjalanan, sehingga model
  lalu lintas diperlukan dalam evaluasi investasi jalan (bagian Access
  measurements, hlm. 243).
\item
  \textbf{Laporan sumber atas literatur kompetisi pekerjaan:} indeks
  Shen memasukkan pasokan pekerjaan yang dapat dijangkau dan permintaan
  potensial dari populasi pekerja. Bab mengaitkannya dengan metode
  \emph{two-step floating catchment area} yang digunakan dalam kajian
  akses pelayanan kesehatan dan pengembangannya dengan peluruhan jarak
  (bagian Potential job access with impedance function and competition,
  hlm. 245-246).
\item
  \textbf{Laporan sumber atas literatur keadilan:} penulis membedakan
  keadilan spasial, yang berhubungan dengan lokasi geografis kelompok
  atau wilayah terdampak, dari keadilan sosial, yang berhubungan dengan
  kondisi ekonomi atau sosial. J11 membatasi analisis utamanya pada
  keadilan spasial (bagian Measuring spatial equity, hlm. 246-247).
\item
  \textbf{Laporan sumber atas pemilihan indeks:} bab menyebut bahwa satu
  ukuran keadilan dapat menghasilkan evaluasi yang bias. Koefisien Gini
  dipilih untuk distribusi keseluruhan, sedangkan rasio Palma
  membandingkan akses rata-rata distrik pada kelompok 10\% terkaya
  dengan kelompok 40\% termiskin. Analisis klaster ditambahkan karena
  kedua indeks agregat tidak menunjukkan lokasi perubahan dan
  karakteristik distrik yang terdampak (bagian Measuring spatial equity,
  hlm. 247-248; bagian Cluster analysis, hlm. 255-256).
\item
  \textbf{Laporan sumber atas heterogenitas dampak:} penelitian
  terdahulu yang dirujuk bab menyatakan bahwa dampak jalan baru dapat
  berbeda antara wilayah perdesaan, pinggiran kota, dan perkotaan serta
  antara segmen pendapatan. Atas dasar itu, pertumbuhan lahan permukiman
  dan pengeluaran bulanan per kapita dimasukkan bersama jumlah pekerjaan
  serta populasi pekerja dalam analisis klaster (bagian Measuring
  spatial equity, hlm. 248).
\end{itemize}

\textbf{Inferensi review.} Kerangka literatur mendukung penggunaan
beberapa ukuran yang saling melengkapi. Namun, file tidak menyajikan
daftar pustaka lengkap, sehingga rincian dan ketepatan seluruh rujukan
yang muncul dalam catatan kaki tidak dapat diperiksa hanya dari TXT ini.

\subsection{Method Used}\label{method-used-10}

\textbf{Laporan sumber.} Penelitian menggunakan evaluasi \emph{ex post}
berbasis model dengan unit analisis distrik. Efek jalan tol
direpresentasikan oleh perbedaan antara skenario jaringan dengan Jalan
Tol Cipularang (\texttt{r\ =\ 1}) dan tanpa jalan tol (\texttt{r\ =\ 0})
(bagian Methods, hlm. 241-242; persamaan 13.1-13.5, hlm. 243-245).

\textbf{Inferensi review atas rancangan.} Perbandingan tersebut
merupakan evaluasi berbasis skenario model, bukan eksperimen atau
perbandingan observasional langsung sebelum-sesudah. File tidak
melaporkan desain eksperimen atau data hasil akses individual sebelum
dan sesudah pembukaan jalan tol.

\textbf{Tahap 1: simulasi transportasi.}

\begin{itemize}
\tightlist
\item
  Penulis membangun jaringan transportasi ArcGIS dari data
  OpenStreetMap, kemudian menggunakan ekstensi model empat tahap
  \texttt{Traffic\ Analyst\ for\ ArcGIS} untuk menghasilkan arus lalu
  lintas, matriks waktu tempuh, dan matriks biaya tergeneralisasi pada
  kedua skenario (bagian Methods, hlm. 241-242).
\item
  Simulasi berlangsung selama 24 jam dan memasukkan hambatan kapasitas
  untuk merepresentasikan kemacetan. Penulis menyatakan bahwa model
  ditujukan untuk memprediksi waktu tempuh wilayah Jakarta-Bandung dan
  terutama arus pada koridor Cipularang; lalu lintas pada jalan lain
  diabaikan menurut uraian metode (bagian Methods, hlm. 241-242).
\item
  Pada tahap bangkitan perjalanan, populasi dan proporsi lahan
  permukiman digunakan sebagai parameter produksi, sedangkan jumlah
  pekerjaan serta proporsi lahan industri dan komersial digunakan
  sebagai parameter tarikan. Parameter diperoleh melalui regresi dengan
  data populasi serta tata guna lahan kawasan metropolitan Jakarta
  (bagian Traffic simulation, hlm. 242).
\item
  Perjalanan seimbang didistribusikan dengan metode Furness. Fungsi
  hambatan menggunakan parameter gravitasi yang dihitung dari survei
  perjalanan angkatan kerja Indonesia 2015 dan matriks biaya lalu lintas
  jalan (bagian Traffic simulation, hlm. 242).
\item
  Biaya tergeneralisasi menggabungkan biaya bahan bakar menurut jarak,
  tarif tol menurut jarak, serta waktu tempuh yang dikalikan nilai
  waktu. Waktu observasi berasal dari Google Maps API, sedangkan nilai
  waktu berasal dari eksperimen pilihan tersurat terhadap 1.600
  responden di wilayah Jakarta-Bandung (Equation 13.2 dan Table 13.3,
  hlm. 243 dan 246; uraian Traffic simulation, hlm. 242).
\item
  Matriks asal-tujuan diestimasi untuk mobil, kendaraan berat, dan
  sepeda motor berdasarkan proporsi moda. Pembebanan jaringan memakai
  algoritma keseimbangan pengguna untuk memperoleh rute, volume lalu
  lintas, waktu, jarak, dan biaya tergeneralisasi antarpasangan zona
  (bagian Traffic simulation, hlm. 242).
\end{itemize}

\textbf{Tahap 2: pengukuran akses.}

\begin{itemize}
\tightlist
\item
  Rata-rata waktu tempuh dan biaya tergeneralisasi dihitung untuk semua
  perjalanan dari setiap asal pada skenario dengan dan tanpa jalan tol
  (Equation 13.1, hlm. 243).
\item
  Ukuran akses pekerjaan potensial menjumlahkan pekerjaan pada setiap
  tujuan setelah dibobot dengan fungsi hambatan waktu atau biaya. Empat
  bentuk fungsi peluruhan diuji, yaitu pangkat negatif, eksponensial
  negatif, log-logistik, dan Gaussian. Fungsi log-logistik dipilih
  karena memberi kecocokan terbaik terhadap data perjalanan; nilai
  \texttt{R2} yang dilaporkan adalah 1,00 untuk waktu dan 1,00 untuk
  biaya tergeneralisasi (Table 13.1, hlm. 244). Parameter log-logistik
  untuk waktu adalah \texttt{a\ =\ -11,55} dan \texttt{b\ =\ 2,50},
  sedangkan untuk biaya tergeneralisasi \texttt{a\ =\ -11,71} dan
  \texttt{b\ =\ 2,10} (Equation 13.4 dan Table 13.2, hlm. 244-245).
\item
  Indeks Shen membagi pasokan potensial pekerjaan yang dapat dijangkau
  dengan permintaan potensial dari populasi pekerja yang bersaing untuk
  pekerjaan tersebut. Nilai 1 ditafsirkan sebagai keseimbangan antara
  pekerjaan dan populasi pekerja; nilai di atas 1 menunjukkan pekerjaan
  relatif lebih banyak (Equation 13.5 dan uraian indeks Shen, hlm.
  245-246).
\end{itemize}

\textbf{Tahap 3: pengukuran keadilan spasial.}

\begin{itemize}
\tightlist
\item
  Koefisien Gini dihitung dengan pendekatan kurva Lorenz trapesium,
  yaitu membandingkan pangsa kumulatif akses dengan pangsa kumulatif
  populasi yang bersesuaian pada distrik-distrik di wilayah studi. Nilai
  0 berarti merata sempurna dan 1 berarti timpang sempurna (bagian
  Measuring spatial equity, hlm. 247).
\item
  Rasio Palma membandingkan akses rata-rata distrik 10\% terkaya dengan
  distrik 40\% termiskin. Karena pendapatan per kapita tidak tersedia,
  peringkat kelompok dibentuk dengan pengeluaran bulanan per kapita
  (bagian Measuring spatial equity, hlm. 247-248).
\item
  Klaster dua tahap menggunakan jumlah klaster yang meminimalkan
  \emph{Bayesian information criterion} (BIC). Karena seluruh jenis
  akses potensial berkorelasi lebih dari 0,9, klaster akhir hanya
  menggunakan indeks Shen dengan peluruhan biaya tergeneralisasi,
  bersama pertumbuhan lahan permukiman, pengeluaran bulanan, jumlah
  pekerjaan, dan populasi pekerja (bagian Cluster analysis, hlm.
  255-256).
\end{itemize}

\textbf{Validasi model.} Kecocokan waktu tempuh simulasi terhadap Google
Maps API diuji dengan RMSPE dan koefisien korelasi. Arus lalu lintas
hasil simulasi pada Jalan Tol Cipularang juga dibandingkan dengan data
arus observasi (bagian Travel time estimation, hlm. 250-251).

\subsection{Dataset}\label{dataset-10}

\begin{itemize}
\tightlist
\item
  \textbf{Jaringan jalan:} data OpenStreetMap sampai tingkat jalan
  lokal, dengan sitasi tahun 2016 pada catatan bab. Tanggal unduh, versi
  jaringan, prosedur pembersihan, atribut kapasitas, dan kelengkapan
  jaringan: Tidak disebutkan dalam file (Methods, hlm. 241-242).
\item
  \textbf{Populasi dan tata guna lahan:} data populasi serta tata guna
  lahan kawasan metropolitan Jakarta dari JICA (2004) untuk mengestimasi
  parameter bangkitan perjalanan (Traffic simulation, hlm. 242).
\item
  \textbf{Mobilitas tenaga kerja:} survei perjalanan angkatan kerja
  Indonesia 2015 dari BPS untuk menghitung parameter gravitasi
  berdasarkan mobilitas angkatan kerja di wilayah studi (Traffic
  simulation, hlm. 242).
\item
  \textbf{Waktu tempuh observasi:} Google Maps API digunakan dalam
  penyusunan biaya lalu lintas dan validasi waktu tempuh simulasi.
  Tanggal serta jam pengambilan, jumlah pasangan asal-tujuan, dan
  kondisi lalu lintas saat pengambilan: Tidak disebutkan dalam file
  (Traffic simulation, hlm. 242; Travel time estimation, hlm. 250-251).
\item
  \textbf{Eksperimen pilihan tersurat:} survei 1.600 responden di
  wilayah Jakarta-Bandung untuk memperoleh nilai waktu. Desain
  eksperimen dirujuk ke Andani et al.~(2021), tetapi rinciannya tidak
  tersedia dalam file J11 (Traffic simulation, hlm. 242).
\item
  \textbf{Pekerjaan dan angkatan kerja:} basis data Indonesia untuk
  riset kebijakan dan ekonomi, yang dirujuk sebagai World Bank (2018),
  digunakan untuk menghitung akses. Data pekerjaan menjadi proksi jumlah
  pekerjaan, sedangkan data angkatan kerja menjadi ukuran pencari
  peluang kerja (bagian indeks Shen, hlm. 245-246).
\item
  \textbf{Skala data tenaga kerja:} data pekerjaan dan angkatan kerja
  hanya tersedia pada tingkat kota/kabupaten, lalu diekstrapolasi ke
  tingkat distrik memakai rasio populasi (bagian indeks Shen, hlm. 246).
\item
  \textbf{Pendapatan:} data pendapatan per kapita tidak tersedia;
  pengeluaran bulanan per kapita digunakan sebagai proksi dalam analisis
  keadilan dan klaster. Sumber, tahun, serta cara konstruksi variabel
  pengeluaran: Tidak disebutkan dalam file (Measuring spatial equity,
  hlm. 247-248).
\item
  \textbf{Perubahan lahan:} bab melaporkan pertumbuhan lahan permukiman
  2,5\% dan lahan industri 1\% selama 2004-2013 serta merujuk Andani et
  al.~(2019). File tidak menjelaskan apakah angka ini dihitung ulang
  dalam J11 atau diambil sepenuhnya dari karya yang dirujuk (Study area,
  Figure 13.3, hlm. 250).
\item
  \textbf{Arus lalu lintas observasi:} digunakan untuk memvalidasi arus
  simulasi di Jalan Tol Cipularang. Sumber data, tahun observasi, jumlah
  titik pencacahan, dan periode pencacahan: Tidak disebutkan dalam file
  (Travel time estimation, hlm. 250-251).
\item
  \textbf{Komponen biaya lain:} tarif tol, biaya bahan bakar, nilai
  waktu, jarak, dan proporsi moda disebut sebagai masukan model. Nilai
  rinci serta tahun acuannya, selain parameter fungsi peluruhan dan
  penjelasan sumber nilai waktu: Tidak disebutkan dalam file (Equation
  13.2 dan Table 13.3, hlm. 243 dan 246).
\end{itemize}

\subsection{Population / Sample}\label{population-sample-10}

\textbf{Unit analisis utama.} Wilayah studi seluas sekitar 15.250 km2
meliputi koridor antara Jakarta dan Bandung, terdiri atas 16 kota atau
kabupaten dan 229 distrik. Jalan Tol Cipularang memiliki lima gerbang;
tiga berada di Kabupaten Bandung Barat dan dua di Kabupaten Purwakarta
(bagian Study area, hlm. 248-249; Figure 13.2, hlm. 249).

\textbf{Populasi substantif.} Pekerjaan didekati melalui jumlah orang
yang bekerja. Orang bekerja didefinisikan sebagai mereka yang bekerja
untuk memperoleh upah atau membantu orang lain memperoleh upah atau
keuntungan setidaknya satu jam pada minggu survei. Angkatan kerja
mencakup penduduk berusia 15 tahun ke atas yang bekerja, sementara tidak
bekerja tetapi memiliki pekerjaan, atau tidak bekerja dan sedang mencari
pekerjaan. Pelajar, pengurus rumah tangga, serta orang yang tidak
melakukan dan tidak mencari kegiatan berpenghasilan dikeluarkan dari
definisi ini (bagian indeks Shen, hlm. 246).

\textbf{Sampel survei nilai waktu.} Eksperimen pilihan tersurat
melibatkan 1.600 responden di wilayah Jakarta-Bandung. Teknik penarikan
sampel, distribusi sosial ekonomi, distribusi geografis, tingkat
respons, dan waktu pelaksanaan: Tidak disebutkan dalam file (Traffic
simulation, hlm. 242).

\textbf{Sampel spasial klaster.} Seluruh 229 distrik dikelompokkan
menjadi tiga klaster. Klaster 1 mencakup 2,6\% distrik, Klaster 2
mencakup 19,2\%, dan Klaster 3 mencakup 78\%. File memberikan
persentase, bukan jumlah distrik bulat pada tiap klaster (uraian
klaster, hlm. 256-257).

\textbf{Ukuran populasi.} Jumlah total penduduk, jumlah total angkatan
kerja, dan jumlah total pekerjaan untuk seluruh wilayah studi: Tidak
disebutkan dalam file. Angka pada Table 13.7 merupakan pusat atau nilai
rata-rata klaster, bukan total populasi wilayah (hlm. 257).

\subsection{Dependent Variables}\label{dependent-variables-10}

\textbf{Keterangan terminologis.} Istilah ``variabel dependen'' tidak
digunakan untuk model evaluasi utama. Jika dimaknai sebagai keluaran
yang dibandingkan antara skenario dengan dan tanpa jalan tol, variabel
hasilnya adalah sebagai berikut.

\begin{itemize}
\tightlist
\item
  Rata-rata waktu tempuh antardistrik dalam menit dan perubahan
  relatifnya (Equation 13.1; Table 13.4, hlm. 243 dan 254).
\item
  Rata-rata biaya tergeneralisasi, dilaporkan dalam ribuan IDR, yang
  menggabungkan bahan bakar, tol, dan nilai waktu (Equation 13.2; Tables
  13.3-13.4, hlm. 243, 246, dan 254).
\item
  Arus lalu lintas hasil pembebanan jaringan serta perbedaannya antara
  skenario (bagian Travel time estimation; Figure 13.5, hlm. 250-252).
\item
  Akses populasi pekerja dan akses pekerjaan potensial dengan hambatan
  waktu maupun biaya tergeneralisasi (Equation 13.3; Table 13.5, hlm.
  244 dan 254-255).
\item
  Indeks Shen sebagai rasio pasokan potensial pekerjaan terhadap
  permintaan potensial pekerja, juga dihitung dengan hambatan waktu dan
  biaya tergeneralisasi (Equation 13.5; Table 13.5, hlm. 245 dan
  254-255).
\item
  Koefisien Gini dan rasio Palma untuk indeks Shen serta akses pekerjaan
  potensial dengan peluruhan biaya tergeneralisasi (Table 13.6, hlm.
  256).
\item
  Keanggotaan dan pusat klaster berdasarkan perubahan relatif indeks
  Shen, pengeluaran per kapita, pertumbuhan lahan permukiman, populasi
  pekerja, dan jumlah pekerjaan (Table 13.7, hlm. 257).
\end{itemize}

\textbf{Variabel pembanding utama.} Keberadaan Jalan Tol Cipularang
dikodekan sebagai skenario \texttt{r\ =\ 1}, sedangkan ketidakhadirannya
sebagai \texttt{r\ =\ 0} (Equations 13.1, 13.3, dan 13.5, hlm. 243-245).
Bab tidak melaporkan model regresi kausal dengan koefisien perlakuan
jalan tol.

\subsection{Results}\label{results-10}

\textbf{Kinerja dan validasi model.}

\begin{itemize}
\tightlist
\item
  \textbf{Laporan sumber:} RMSPE antara waktu tempuh observasi dan
  simulasi adalah 9,7\%, dengan koefisien korelasi 0,90. Penulis menilai
  kesalahan ini cukup kecil dan akurasi model cukup tinggi untuk
  memprediksi waktu tempuh pasangan asal-tujuan (Travel time estimation
  dan Figure 13.4, hlm. 250-251).
\item
  \textbf{Laporan sumber:} rata-rata waktu tempuh model adalah 154
  menit, sedangkan rata-rata Google Maps API adalah 155 menit. Estimasi
  arus lalu lintas pada Jalan Tol Cipularang memiliki galat persentase
  3,7\% terhadap arus observasi (Travel time estimation, hlm. 250-251).
\item
  \textbf{Laporan sumber:} dalam skenario tanpa Jalan Tol Cipularang,
  arus lebih rendah di Bandung dan bagian utara wilayah, tetapi
  meningkat signifikan pada jalan regional nontol lain (Figure 13.5 dan
  uraian, hlm. 251-252).
\end{itemize}

\textbf{Waktu tempuh dan biaya tergeneralisasi.}

\begin{itemize}
\tightlist
\item
  \textbf{Laporan sumber:} setelah dibobot dengan arus lalu lintas,
  Jalan Tol Cipularang menurunkan rata-rata waktu tempuh seluruh wilayah
  sebesar 13\%. Beberapa distrik di Purwakarta mengalami penurunan waktu
  tempuh rata-rata tertimbang hingga 25\% (Travel time estimation, hlm.
  251; Table 13.4, hlm. 254).
\item
  \textbf{Laporan sumber:} rata-rata deskriptif antardistrik pada Table
  13.4 adalah 154,5 menit dengan tol dan 167,5 menit tanpa tol. Nilai
  maksimum, minimum, dan simpangan bakunya masing-masing adalah 332,6;
  109,4; dan 36,5 menit dengan tol serta 335,1; 129,0; dan 35,5 menit
  tanpa tol (Table 13.4, hlm. 254).
\item
  \textbf{Laporan sumber:} biaya tergeneralisasi rata-rata adalah 231,7
  ribu IDR dengan tol dan 233,4 ribu IDR tanpa tol. Nilai maksimum,
  minimum, dan simpangan bakunya masing-masing adalah 453,8; 180,3; dan
  43,6 ribu IDR dengan tol serta 454,3; 180,6; dan 43,6 ribu IDR tanpa
  tol (Table 13.4, hlm. 254).
\item
  \textbf{Laporan sumber:} pengurangan terbesar untuk pasangan
  asal-tujuan ditemukan antara Cipendeuy di Kabupaten Bandung dan Teluk
  Jambe di Karawang; waktu tempuh tanpa tol hampir 1,5 kali waktu tempuh
  dengan tol. Perbedaan biaya tergeneralisasi lebih kecil daripada
  perbedaan waktu karena pengguna harus membayar tol (Travel time
  estimation, hlm. 251-252; Figure 13.6, hlm. 253).
\end{itemize}

\textbf{Ukuran akses.}

\begin{itemize}
\tightlist
\item
  \textbf{Laporan sumber:} ukuran dengan peluruhan waktu menghasilkan
  perubahan yang lebih besar daripada ukuran dengan peluruhan biaya
  tergeneralisasi. Penulis menilai pengukuran berbasis waktu saja
  melebihkan dampak karena tidak memasukkan tarif tol (Potential access
  measures, hlm. 252-253).
\item
  \textbf{Laporan sumber:} akses populasi pekerja berbobot waktu
  bernilai 8.135.674 dengan tol dan 7.804.450 tanpa tol, atau berubah
  4,2\%. Dengan hambatan biaya tergeneralisasi, nilainya 11.304.273
  dengan tol dan 11.264.201 tanpa tol, atau berubah 0,4\% (Table 13.5,
  hlm. 254-255).
\item
  \textbf{Laporan sumber:} indeks Shen berbobot waktu turun dari 0,78
  tanpa tol menjadi 0,73 dengan tol, dengan perubahan relatif -5,6\%.
  Dengan hambatan biaya tergeneralisasi, indeks turun dari 0,74 menjadi
  0,73, dengan perubahan relatif -0,41\% (Table 13.5, hlm. 254-255).
\item
  \textbf{Laporan sumber:} akses pekerjaan potensial berbobot waktu
  bernilai 6.016.274 dengan tol dan 5.755.050 tanpa tol. Dengan hambatan
  biaya tergeneralisasi, nilainya 8.251.967 dengan tol dan 8.221.007
  tanpa tol. Baris perubahan relatif untuk ukuran ini tidak terlihat
  pada Table 13.5 dalam hasil ekstraksi, sedangkan Abstract dan
  Conclusions merangkum peningkatan akses pekerjaan atau
  pekerjaan-populasi pekerja sebagai sekitar 5\% (Table 13.5, hlm.
  254-255; Abstract; Conclusions, hlm. 259).
\item
  \textbf{Laporan sumber:} beberapa distrik di Jakarta dan Purwakarta
  mempunyai indeks Shen sekitar 1. Sebuah distrik di Bekasi mempunyai
  indeks di atas 2, yang berarti pekerjaan lebih banyak daripada pekerja
  dalam ukuran potensial tersebut. Indeks Shen justru lebih rendah pada
  skenario dengan tol, dengan penurunan terkuat di wilayah miskin
  pekerjaan antara Jakarta dan Bandung (Figure 13.7 dan uraian, hlm.
  253-255).
\end{itemize}

\textbf{Indeks keadilan.}

\begin{itemize}
\tightlist
\item
  \textbf{Laporan sumber:} dengan peluruhan biaya tergeneralisasi, rasio
  Palma indeks Shen tetap 1,10 pada kedua skenario, sedangkan koefisien
  Gini naik sedikit dari 0,29 tanpa tol menjadi 0,30 dengan tol (Table
  13.6, hlm. 256).
\item
  \textbf{Laporan sumber:} untuk akses pekerjaan potensial, rasio Palma
  tetap 1,14 dan koefisien Gini tetap 0,38 pada kedua skenario (Table
  13.6, hlm. 256).
\item
  \textbf{Laporan sumber:} ukuran dengan peluruhan waktu tidak
  dimasukkan ke analisis keadilan karena penulis menilainya melebihkan
  dampak jalan tol (bagian Equity indices, hlm. 254-255).
\end{itemize}

\textbf{Klaster spasial.}

\begin{itemize}
\tightlist
\item
  \textbf{Laporan sumber:} tiga klaster akhir diperoleh berdasarkan BIC
  minimum, yaitu distrik terdampak, distrik tidak terdampak yang lebih
  terurbanisasi, serta distrik tidak terdampak yang kurang terurbanisasi
  (bagian Cluster analysis, hlm. 256).
\item
  \textbf{Laporan sumber:} Klaster 1, yang mencakup 2,6\% distrik,
  mempunyai perubahan relatif indeks Shen terendah sebesar -2,02\%.
  Pusat klasternya menunjukkan pengeluaran bulanan per kapita 8.626 ribu
  IDR, pertumbuhan lahan permukiman 5,0\%, populasi pekerja 104.303, dan
  14.526 pekerjaan. Distrik-distrik ini berada dekat gerbang tol, sangat
  mudah diakses, mengalami penurunan waktu terbesar, dan memiliki
  pengeluaran per kapita tertinggi (Table 13.7 dan uraian, hlm.
  256-257).
\item
  \textbf{Laporan sumber:} Klaster 2 mencakup 19,2\% distrik dan
  mempunyai perubahan indeks Shen -0,34\%, pengeluaran 2.473 ribu IDR
  per kapita, pertumbuhan permukiman 5,2\%, populasi pekerja 155.695,
  dan 138.027 pekerjaan. Penulis menggolongkannya sebagai wilayah tidak
  terdampak yang lebih terurbanisasi (Table 13.7 dan uraian, hlm. 257).
\item
  \textbf{Laporan sumber:} Klaster 3 mencakup 78\% distrik dan mempunyai
  perubahan indeks Shen -0,48\%, pengeluaran 1.342 ribu IDR per kapita,
  pertumbuhan permukiman 2,9\%, populasi pekerja 54.766, dan 36.204
  pekerjaan. Klaster ini mempunyai pekerjaan, populasi pekerja,
  pengeluaran rata-rata, dan pertumbuhan permukiman terendah (Table 13.7
  dan uraian, hlm. 257).
\end{itemize}

\subsection{Findings}\label{findings-10}

\begin{itemize}
\tightlist
\item
  \textbf{Interpretasi penulis:} manfaat waktu tempuh terutama
  terkonsentrasi di koridor jalan tol dan dekat gerbang, khususnya
  Purwakarta yang terletak di antara pusat pekerjaan Jakarta dan Bandung
  serta memiliki dua gerbang tol. Bogor dan Cianjur di bagian barat
  termasuk wilayah yang paling sedikit terdampak (Travel time
  estimation, hlm. 251).
\item
  \textbf{Interpretasi penulis:} tarif tol mengurangi besarnya manfaat
  apabila akses dinilai dengan biaya tergeneralisasi. Oleh karena itu,
  perubahan akses berbasis waktu saja lebih besar dan dapat melebihkan
  manfaat dibandingkan ukuran yang memasukkan biaya moneter (Potential
  access measures, hlm. 252-253).
\item
  \textbf{Interpretasi penulis:} jalan tol meningkatkan jangkauan
  terhadap pekerjaan sekaligus memperluas jangkauan pekerja yang
  bersaing. Di wilayah miskin pekerjaan antara Jakarta dan Bandung,
  tambahan populasi pekerja yang dapat menjangkau pekerjaan menurunkan
  jumlah pekerjaan yang dapat diakses per pekerja, sebagaimana
  ditunjukkan oleh penurunan indeks Shen (Figure 13.7 dan uraian, hlm.
  253-255; Conclusions, hlm. 259).
\item
  \textbf{Interpretasi penulis:} kenaikan kecil koefisien Gini untuk
  indeks Shen menunjukkan bahwa pekerja berpendapatan menengah hingga
  tinggi memperoleh manfaat lebih besar daripada pekerja berpendapatan
  rendah. Namun, rasio Palma yang tetap menunjukkan bahwa kelompok
  distrik 10\% teratas tidak memperoleh manfaat lebih besar daripada
  kelompok 40\% terbawah (bagian Equity indices, hlm. 255; Table 13.6,
  hlm. 256).
\item
  \textbf{Laporan sumber dan interpretasi penulis:} kesimpulan
  distribusional tersebut dinyatakan belum sepenuhnya representatif
  karena pilihan moda tidak dimasukkan. Penulis menjelaskan bahwa
  pekerja berpendapatan rendah di Indonesia mengandalkan sepeda motor,
  sedangkan sepeda motor tidak diizinkan di sebagian besar jalan tol;
  dalam praktiknya, pekerja berpendapatan lebih tinggi mungkin lebih
  diuntungkan (bagian Equity indices, hlm. 255-256).
\item
  \textbf{Interpretasi penulis:} ketidakseimbangan antara 104.303
  pekerja dan 14.526 pekerjaan pada pusat Klaster 1 menjelaskan
  penurunan indeks Shen yang lebih besar di wilayah dekat gerbang.
  Penurunan waktu tempuh dan munculnya pekerjaan baru di koridor
  memungkinkan lebih banyak pekerja menjangkau pekerjaan yang tersedia
  (Table 13.7 dan uraian, hlm. 256-257).
\item
  \textbf{Inferensi review:} stabilnya Gini dan Palma secara agregat
  tidak identik dengan ketiadaan dampak di semua lokasi. Hasil klaster
  memperlihatkan bahwa skala pengukuran menentukan kesimpulan: perubahan
  rata-rata wilayah kecil, tetapi subkelompok distrik dekat gerbang
  mengalami perubahan indeks Shen yang lebih negatif (Tables 13.6-13.7,
  hlm. 256-257).
\item
  \textbf{Inferensi review:} karena peringkat ``kaya'' dan ``miskin''
  dibentuk dari pengeluaran per kapita tingkat distrik, hasil utamanya
  menunjukkan pola keadilan spasial antardistrik. Hasil tersebut belum
  langsung membuktikan distribusi manfaat pada individu menurut
  pendapatan, terlebih sumber sendiri mengakui ketiadaan pemodelan
  pilihan moda (Measuring spatial equity, hlm. 247-248; Equity indices,
  hlm. 255-256).
\end{itemize}

\subsection{Conclusions}\label{conclusions-10}

\textbf{Kesimpulan penulis.} Jalan Tol Cipularang menurunkan rata-rata
waktu tempuh wilayah Jakarta-Bandung sebesar 13\%. Untuk pasangan
asal-tujuan tertentu, waktu tempuh tanpa tol dilaporkan hampir 1,5 kali
waktu tempuh dengan tol. Akses pekerjaan dan populasi pekerja meningkat
sekitar 5\%, tetapi kompetisi pekerjaan juga meningkat sehingga
pekerjaan yang dapat diakses per pekerja sedikit menurun (Travel time
estimation, hlm. 251; Conclusions, hlm. 259).

\textbf{Kesimpulan penulis.} Wilayah sepanjang jalan tol dan dekat
gerbang merupakan wilayah yang paling terdampak. Wilayah tersebut
cenderung mempunyai populasi pekerja yang besar dibandingkan jumlah
pekerjaan dan mengalami pertumbuhan lahan permukiman. Penulis
menyimpulkan bahwa jalan tol baru dalam konteks negara berkembang dengan
jaringan jalan yang belum matang dapat memberi dampak positif sekaligus
negatif pada wilayah di dekatnya (Conclusions, hlm. 259).

\textbf{Kesimpulan penulis.} Secara rata-rata, pembangunan jalan tol
dinilai tidak mengubah keadilan spasial menurut rasio Palma dan sebagian
besar hasil koefisien Gini. Pengecualiannya adalah kenaikan kecil Gini
indeks Shen dari 0,29 menjadi 0,30. Analisis klaster sekaligus
menunjukkan adanya dampak spasial yang lebih nyata di wilayah dekat
gerbang (Abstract; Equity indices dan Table 13.6, hlm. 255-256; uraian
klaster, hlm. 257).

\textbf{Inferensi review.} Dukungan empiris dalam file terbatas pada
keluaran model kontrafaktual untuk satu koridor dan satu konfigurasi
data. Generalisasi penulis ke jalan tol baru di negara berkembang perlu
dibaca sebagai interpretasi yang belum diuji lintas proyek, lintas
wilayah, atau lintas kelompok moda dalam J11.

\subsection{Contributions}\label{contributions-10}

\begin{itemize}
\tightlist
\item
  \textbf{Kontribusi yang dinyatakan penulis:} menambah literatur
  tentang dampak jalan tol besar terhadap akses pekerjaan dan keadilan
  spasial dalam konteks negara berkembang (Conclusions, hlm. 259).
\item
  \textbf{Kontribusi yang dinyatakan penulis:} memasukkan efek kemacetan
  dan nilai waktu berbasis eksperimen pilihan tersurat ke simulasi lalu
  lintas untuk mengevaluasi skenario dengan serta tanpa jalan tol
  (Conclusions, hlm. 259).
\item
  \textbf{Kontribusi yang dinyatakan penulis:} memasukkan kompetisi
  pekerjaan melalui indeks Shen sehingga pertambahan jangkauan pekerjaan
  tidak otomatis disamakan dengan pertambahan pekerjaan per pekerja
  (Methods, hlm. 245-246; Conclusions, hlm. 257 dan 259).
\item
  \textbf{Kontribusi yang dinyatakan penulis:} menunjukkan nilai tambah
  analisis akses bagi praktik perencanaan transportasi dan integrasi
  tata guna lahan-transportasi agar masyarakat dapat berpartisipasi
  dalam aktivitas yang diinginkan (Conclusions, hlm. 259).
\item
  \textbf{Inferensi review:} penggunaan bersama Gini, Palma, dan klaster
  memperlihatkan perbedaan antara stabilitas distribusi agregat dan
  dampak lokal. Nilai tambah ini berasal dari kombinasi metode yang
  dilaporkan, meskipun penulis tidak merumuskannya dengan istilah
  ``triangulasi'' (Measuring spatial equity, hlm. 246-248; Tables
  13.6-13.7, hlm. 256-257).
\end{itemize}

\subsection{Research Gap}\label{research-gap-10}

\textbf{Kesenjangan yang memotivasi studi.}

\begin{itemize}
\tightlist
\item
  Evaluasi \emph{ex post} investasi jalan dengan ukuran akses masih
  jarang, terutama di negara berkembang (Introduction, halaman pembuka
  bab).
\item
  Evaluasi \emph{ex post} infrastruktur transportasi di Indonesia belum
  mendapat perhatian memadai, sementara tahap pemantauan, evaluasi, dan
  umpan balik dalam siklus kebijakan sering diabaikan (Introduction,
  hlm. 240-241).
\item
  Penilaian keadilan dengan satu indeks dapat bias, sedangkan Gini dan
  Palma tidak menerangkan lokasi perubahan maupun karakteristik distrik
  yang mendapat manfaat atau kerugian. Bab mengisi bagian ini dengan
  analisis klaster (Measuring spatial equity, hlm. 247-248; Cluster
  analysis, hlm. 255-256).
\item
  Ukuran akses pekerjaan tanpa kompetisi tidak menunjukkan apakah
  pekerjaan yang lebih mudah dijangkau juga dihadapkan pada lebih banyak
  pekerja. Indeks Shen digunakan untuk menutup kesenjangan pengukuran
  tersebut (Access measurements, hlm. 243-246).
\end{itemize}

\textbf{Kesenjangan yang masih tersisa menurut sumber.} Dampak terhadap
distribusi pekerjaan, penduduk, dan biaya perumahan belum dianalisis;
akses menurut kelompok pekerja belum dapat diestimasi; serta perbedaan
perilaku perjalanan dan ketersediaan moda individual belum dimasukkan
(bagian Next, hlm. 259-260).

\textbf{Inferensi review.} File belum menutup kesenjangan antara
keadilan spasial berbasis karakteristik rata-rata distrik dan keadilan
sosial pada tingkat individu. Keterbatasan ini terlihat dari unit
analisis distrik, penggunaan pengeluaran distrik sebagai proksi
pendapatan, dan pengakuan bahwa pilihan moda individual tidak
dimodelkan.

\subsection{Limitations}\label{limitations-10}

\textbf{Keterbatasan yang dinyatakan atau diakui sumber.}

\begin{itemize}
\tightlist
\item
  Penerapan model empat tahap menghadapi keterbatasan data. Model
  berfokus pada prediksi waktu wilayah dan arus koridor Cipularang;
  penulis menyatakan bahwa lalu lintas pada jalan lain diabaikan
  (Methods, hlm. 241-242).
\item
  Data pekerjaan dan angkatan kerja hanya tersedia pada tingkat
  kota/kabupaten, lalu diekstrapolasi ke distrik dengan rasio populasi.
  Penulis tidak dapat membedakan pekerja yang sudah bekerja dari pencari
  kerja secara terperinci dan tidak dapat memasukkan jenis pekerjaan,
  misalnya kelas pekerjaan (bagian indeks Shen, hlm. 246).
\item
  Pendapatan per kapita tidak tersedia sehingga pengeluaran bulanan per
  kapita digunakan sebagai proksi (Measuring spatial equity, hlm. 247).
\item
  Ukuran peluruhan waktu dinilai melebihkan dampak dan dikeluarkan dari
  analisis keadilan karena tidak memasukkan biaya tol (Potential access
  measures dan Equity indices, hlm. 252-255).
\item
  Pilihan moda tidak diperhitungkan dalam hasil keadilan. Hal ini
  penting karena pekerja berpendapatan rendah lebih bergantung pada
  sepeda motor, sedangkan sepeda motor tidak diizinkan di sebagian besar
  jalan tol (Equity indices, hlm. 255-256).
\item
  Dampak jalan tol terhadap perkembangan spasial, termasuk distribusi
  pekerjaan, populasi, dan biaya perumahan, tidak diperiksa (Next, hlm.
  259).
\item
  Akses menurut kelompok pekerja tidak dapat diestimasi karena data
  pekerjaan sektoral yang andal dan data populasi pekerja beresolusi
  spasial tinggi tidak tersedia (Next, hlm. 259).
\item
  Kompleksitas serta perbedaan perilaku perjalanan individual dan
  ketersediaan moda pada tingkat individu tidak dimasukkan (Next, hlm.
  259-260).
\end{itemize}

\textbf{Batas yang terlihat dari file dan diberi label sebagai inferensi
review.}

\begin{itemize}
\tightlist
\item
  Evaluasi membandingkan dua skenario hasil model. File tidak melaporkan
  desain observasional sebelum-sesudah yang secara langsung mengukur
  perubahan akses individu akibat pembukaan tol. Karena itu, angka
  dampak merupakan estimasi berbasis kontrafaktual model, bukan
  perubahan yang seluruhnya diamati.
\item
  Rujukan data dan metode mencantumkan tahun yang berbeda, antara lain
  JICA (2004), BPS (2015), OpenStreetMap (2016), dan World Bank (2018),
  serta periode perubahan lahan 2004-2013, sementara jalan tol disebut
  dibangun pada 2005. Tahun observasi sebenarnya untuk semua variabel
  dan prosedur harmonisasi tahun dasar: Tidak disebutkan dalam file.
\item
  Validasi yang dilaporkan berfokus pada waktu tempuh dan arus Jalan Tol
  Cipularang. Validasi eksternal untuk hasil akses, indeks keadilan, dan
  keanggotaan klaster: Tidak disebutkan dalam file.
\item
  Interval kepercayaan, galat baku, uji signifikansi perubahan Gini atau
  Palma, dan analisis sensitivitas terhadap asumsi utama selain
  perbandingan fungsi peluruhan: Tidak disebutkan dalam file.
\item
  Hasil keadilan menggunakan agregat distrik. Risiko kekeliruan ekologis
  ketika pola distrik ditafsirkan sebagai manfaat bagi pekerja
  berpendapatan tertentu tidak dibahas secara eksplisit dalam file.
\item
  Metode menyebut matriks asal-tujuan untuk tiga moda, termasuk sepeda
  motor, tetapi bagian hasil menyatakan bahwa pilihan moda tidak
  diperhitungkan. Hubungan tepat antara representasi moda dalam simulasi
  lalu lintas dan pengukuran keadilan tidak dijelaskan dalam file.
\end{itemize}

\subsection{Challenges}\label{challenges-10}

\begin{itemize}
\tightlist
\item
  \textbf{Laporan sumber:} keterbatasan pendanaan khusus dan
  ketersediaan data relevan merupakan tantangan umum bagi evaluasi
  \emph{ex post} transportasi (Introduction, hlm. 240-241).
\item
  \textbf{Laporan sumber:} membangun model empat tahap dalam konteks
  data terbatas merupakan tantangan langsung penelitian. Penulis harus
  menggabungkan jaringan OSM, model bangkitan perjalanan, survei
  mobilitas, waktu Google Maps API, nilai waktu hasil survei, serta data
  pekerjaan dan angkatan kerja yang berbeda skala (Methods, hlm.
  241-246).
\item
  \textbf{Laporan sumber:} data tingkat kota/kabupaten harus diturunkan
  secara tidak langsung ke distrik, sementara rincian sektor pekerjaan
  dan kelompok pekerja tidak tersedia (bagian indeks Shen, hlm. 246;
  Next, hlm. 259).
\item
  \textbf{Interpretasi penulis:} pemilihan ukuran merupakan tantangan
  substantif karena akses berbasis waktu, akses berbasis biaya, akses
  potensial, dan indeks Shen menghasilkan besaran serta arah perubahan
  yang berbeda. Penulis memandang hasil tersebut berbeda tetapi saling
  melengkapi (Potential access measures, hlm. 252-253; Conclusions, hlm.
  259).
\item
  \textbf{Inferensi review:} tantangan identifikasi terletak pada upaya
  mengisolasi satu jalan tol melalui skenario model sambil
  mempertahankan kemacetan dan interaksi jaringan. File tidak
  menyediakan semua parameter atau bahan yang diperlukan untuk
  mereplikasi model secara mandiri.
\item
  \textbf{Inferensi review:} tantangan interpretasi keadilan terletak
  pada perbedaan skala. Indeks agregat hampir tetap, sedangkan klaster
  lokal menunjukkan distrik terdampak; kedua temuan harus dibaca bersama
  agar ``tidak ada dampak rata-rata'' tidak berubah menjadi ``tidak ada
  dampak di mana pun''.
\end{itemize}

\subsection{Practical Implications}\label{practical-implications-10}

\begin{itemize}
\tightlist
\item
  \textbf{Implikasi penulis:} analisis akses dapat menambah nilai pada
  praktik perencanaan transportasi dan memperkuat integrasi tata guna
  lahan dengan transportasi (Conclusions, hlm. 259).
\item
  \textbf{Interpretasi penulis dan inferensi review:} evaluasi jalan tol
  sebaiknya tidak hanya memakai penghematan waktu. Biaya tol perlu
  dimasukkan melalui biaya tergeneralisasi karena manfaat yang tampak
  besar dari waktu saja menyusut ketika biaya moneter diperhitungkan
  (Potential access measures, hlm. 252-253; Table 13.5, hlm. 254-255).
\item
  \textbf{Interpretasi penulis:} kenaikan akses pekerjaan perlu dibaca
  bersama perubahan akses populasi pekerja. Di wilayah miskin pekerjaan,
  konektivitas yang lebih tinggi dapat memperbesar kompetisi dan
  menurunkan pekerjaan yang dapat dijangkau per pekerja (Figure 13.7,
  hlm. 253-255; Conclusions, hlm. 259).
\item
  \textbf{Inferensi review berdasarkan hasil:} indeks wilayah perlu
  dilengkapi analisis spasial terpilah. Pemantauan khusus layak
  diarahkan ke distrik dekat gerbang karena perubahan waktu,
  ketidakseimbangan pekerjaan-pekerja, dan perubahan indeks Shen
  terkonsentrasi di sana (Tables 13.6-13.7, hlm. 256-257).
\item
  \textbf{Interpretasi penulis:} penilaian manfaat menurut pendapatan
  perlu memasukkan akses moda individual, khususnya larangan sepeda
  motor di sebagian besar jalan tol. Tanpa itu, ukuran spasial berbasis
  distrik dapat menutupi hambatan yang dialami pekerja berpendapatan
  rendah (Equity indices, hlm. 255-256; Next, hlm. 259-260).
\item
  \textbf{Inferensi review:} keputusan proyek dan evaluasi lanjutan
  perlu membedakan keberhasilan jaringan, perubahan akses, dan
  distribusi manfaat. J11 menunjukkan bahwa ketiga lapisan tersebut
  dapat bergerak berbeda: waktu membaik, akses potensial meningkat,
  tetapi indeks kompetisi pekerjaan menurun.
\end{itemize}

\subsection{Applications}\label{applications-10}

\begin{itemize}
\tightlist
\item
  \textbf{Laporan sumber:} kerangka ini diterapkan untuk evaluasi
  \emph{ex post} Jalan Tol Cipularang pada 229 distrik di koridor
  Jakarta-Bandung (Methods dan Study area, hlm. 241-249).
\item
  \textbf{Aplikasi metodologis yang diperagakan sumber:} simulasi
  jaringan dengan dan tanpa proyek untuk menilai perubahan waktu tempuh,
  biaya tergeneralisasi, dan arus lalu lintas (Methods, hlm. 241-243).
\item
  \textbf{Aplikasi metodologis yang diperagakan sumber:} evaluasi akses
  pekerjaan yang mempertimbangkan distribusi pekerjaan, populasi
  pekerja, kemacetan, biaya tol, dan kompetisi pekerjaan (Access
  measurements, hlm. 242-246).
\item
  \textbf{Aplikasi metodologis yang diperagakan sumber:} pemetaan
  sasaran wilayah melalui klaster distrik berdasarkan perubahan indeks
  Shen dan karakteristik tata guna lahan serta sosial ekonomi (Cluster
  analysis, hlm. 255-257).
\item
  \textbf{Aplikasi kebijakan yang dinyatakan sumber:} memberikan umpan
  balik bagi penilaian proyek berikutnya dan akuntabilitas dalam siklus
  kebijakan transportasi (Introduction, hlm. 240-241).
\item
  \textbf{Inferensi review:} metode serupa berpotensi digunakan pada
  evaluasi jalan berbayar lain apabila data jaringan, waktu, biaya,
  pekerjaan, populasi pekerja, dan karakteristik sosial ekonomi tersedia
  pada skala yang sebanding. Transfer ke konteks lain memerlukan
  kalibrasi ulang; file tidak menguji validitas eksternal di luar
  Cipularang.
\end{itemize}

\subsection{Future Research
(Optional)}\label{future-research-optional-10}

Sumber menyebut tiga perluasan penelitian secara eksplisit pada bagian
Next, hlm. 259-260.

\begin{enumerate}
\def\labelenumi{\arabic{enumi}.}
\tightlist
\item
  Memasukkan dampak jalan tol terhadap perkembangan spasial, khususnya
  distribusi pekerjaan, penduduk, dan biaya perumahan, karena perubahan
  tersebut akan memengaruhi ukuran akses pekerjaan potensial.
\item
  Mengestimasi akses untuk kelompok pekerja yang berbeda dengan data
  pekerjaan tingkat sektor yang andal serta data populasi pekerja pada
  resolusi spasial tinggi.
\item
  Memasukkan kompleksitas dan perbedaan perilaku perjalanan individual,
  termasuk ketersediaan moda pada tingkat individu. Penulis menekankan
  posisi pekerja berpendapatan rendah yang bergantung pada sepeda motor,
  sementara sepeda motor tidak diizinkan di sebagian besar jalan tol.
\end{enumerate}

Usulan masa depan selain tiga butir tersebut: Tidak disebutkan dalam
file. Review ini tidak menambahkan agenda riset eksternal.

\subsection{Provenance Notes}\label{provenance-notes-10}

\begin{itemize}
\tightlist
\item
  Review ini hanya menggunakan seluruh isi file
  \texttt{/Users/mac/Documents/Mac/{[}2{]}\ Obsidian\ Vault/zettelkasten/source/journal/J11-andani-la-paix-rachmat-syabri-geurs-2021-job-access-and-spatial-equity-of-a-toll-road-no-doi.txt},
  sebanyak 888 baris. Tidak ada sumber eksternal, penelusuran web, basis
  data bibliografis, atau dokumen lain yang digunakan untuk menambah
  klaim substantif.
\item
  Blok metadata TXT menyatakan bahwa teks diekstraksi pada 2026-08-31
  dengan \texttt{pdftotext\ -layout} dari
  \texttt{book/J11-andani-la-paix-rachmat-syabri-geurs-2021-job-access-and-spatial-equity-of-a-toll-road-no-doi.pdf}.
  Metadata menyebut 22 halaman, PDF versi 1.5, tidak terenkripsi, dan
  tidak bertanda struktur (\emph{tagged: no}). Review ini membaca TXT,
  bukan PDF asal.
\item
  Terdapat ketidakselarasan klasifikasi: TXT ditempatkan di
  \texttt{source/journal/}, tetapi jalur PDF pada metadata berada di
  \texttt{book/} dan tubuh teks menyebut dokumen sebagai bab 13. Karena
  rincian publikasi tidak tersedia, review ini tidak menetapkannya
  sebagai artikel jurnal.
\item
  Tahun 2021 berasal dari nama file, bukan keterangan bibliografis pada
  tubuh teks. Judul dan nama penulis berasal dari halaman pembuka. DOI,
  editor, penerbit, ISBN, edisi, tanggal publikasi, URL, dan status
  telaah sejawat: Tidak disebutkan dalam file.
\item
  Locator halaman mengikuti nomor header cetak yang masih terlihat dalam
  hasil ekstraksi, yaitu 240-260. Halaman pembuka sebelum header 240
  dirujuk sebagai ``halaman pembuka bab'' dan tidak diberi nomor hasil
  inferensi.
\item
  Teks memuat judul bagian, persamaan 13.1-13.5, Tables 13.1-13.7, dan
  Figures 13.1-13.8. Isi visual peta dan grafik tidak dapat diperiksa
  dari TXT; review hanya menggunakan caption dan uraian naratif yang
  diekstraksi.
\item
  Tata letak beberapa persamaan dan tabel terdistorsi oleh ekstraksi.
  Secara khusus, notasi Equation 13.1 dan Equation 13.5 tidak sepenuhnya
  utuh, sehingga review tidak merekonstruksi simbol yang hilang. Uraian
  konseptual penulis digunakan sebagai dasar penjelasan metode.
\item
  Pada Table 13.5, nilai skenario akses pekerjaan potensial terlihat,
  tetapi baris perubahan relatifnya tidak tampak dalam ekstraksi. Angka
  sekitar 5\% dalam review berasal dari ringkasan eksplisit penulis pada
  Abstract dan Conclusions, bukan dari perhitungan baru.
\item
  Pernyataan penurunan waktu 13\% dinyatakan sumber sebagai rata-rata
  tertimbang arus lalu lintas. Rata-rata deskriptif Table 13.4 adalah
  154,5 dan 167,5 menit. File tidak menguraikan hubungan pembobotan
  keduanya secara rinci, sehingga review tidak menyamakan atau
  menghitung ulang kedua statistik tersebut.
\item
  Figure 13.1 berada di sekitar pembahasan fungsi peluruhan, tetapi
  caption hasil ekstraksi berbunyi ``Net auto vs.~net transit access to
  all jobs via auto within 30 minutes'', yang tidak selaras dengan
  uraian di sekitarnya. Review tidak menggunakan isi figur tersebut
  sebagai bukti.
\item
  Rujukan seperti Hansen (1959), Geurs dan Van Wee (2004), BPS (2015),
  World Bank (2018), serta karya lain hanya diketahui dari sitasi atau
  catatan kaki di dalam J11. Isi sumber-sumber tersebut tidak diakses
  atau diverifikasi secara independen; setiap penyebutannya di sini
  diperlakukan sebagai laporan J11 atas literatur.
\item
  File tidak menyertakan daftar pustaka lengkap. Akibatnya, rincian
  bibliografis rujukan internal dan klaim yang bergantung pada karya
  terdahulu tidak dapat diverifikasi dari file ini saja.
\item
  Batas provenance utama: review dapat merangkum apa yang tertulis dalam
  ekstraksi teks, tetapi tidak dapat memastikan ketepatan OCR/tata letak
  terhadap PDF, membaca informasi visual yang tidak terekstraksi,
  memverifikasi metadata publikasi, menilai sumber data asli, atau
  menguji ulang model dan perhitungannya.
\end{itemize}

\section{Review J22}\label{review-j22}

\subsection{Identitas Sumber}\label{identitas-sumber-5}

\begin{itemize}
\tightlist
\item
  Kode: J22.
\item
  Judul: \emph{Using Morphological and Functional Polycentricity
  Analyses to Study the Indonesian Urban Spatial Structure: The Case of
  Medan, Jakarta, and Denpasar}.
\item
  Penulis: Erie Sadewo, Ibnu Syabri, Anzhelika Antipova, Pradono, dan
  Delik Hudalah.
\item
  Jurnal: \emph{Asian Geographer}.
\item
  DOI: \texttt{10.1080/10225706.2020.1737829}.
\item
  ISSN yang tercantum: 1022-5706 untuk versi cetak dan 2158-1762 untuk
  versi daring.
\item
  Tahun dan status publikasi: blok sitasi menyebut artikel sebagai
  terbitan 2020 dan halaman depan mencatat publikasi daring pada 12
  Maret 2020. Nama file memuat tahun 2021, tetapi volume, nomor, tahun
  edisi final, dan rentang halaman artikel: Tidak disebutkan dalam file.
\item
  Riwayat naskah: diterima redaksi pada 23 Desember 2018 dan diterima
  untuk publikasi pada 9 Desember 2019.
\item
  Afiliasi penulis: School of Architecture, Planning, and Policy
  Development, Institut Teknologi Bandung; Badan Pusat Statistik; dan
  Department of Earth Science, The University of Memphis.
\item
  Objek studi: struktur spasial Metropolitan Medan atau MMA,
  Metropolitan Jakarta atau JMA, dan Metropolitan Denpasar atau DMA,
  dengan fokus pada pusat pekerjaan serta polisentralitas morfologis dan
  fungsional.
\item
  Pernyataan konflik kepentingan: penulis melaporkan tidak ada potensi
  konflik kepentingan.
\item
  Status telaah sejawat: Tidak disebutkan secara eksplisit dalam file.
  Riwayat penerimaan naskah tidak digunakan dalam review ini untuk
  menetapkan status telaah sejawat.
\end{itemize}

\textbf{Penanda epistemik yang digunakan}

\begin{itemize}
\tightlist
\item
  \textbf{Laporan sumber} adalah metode, data, angka, atau pernyataan
  yang disampaikan langsung dalam artikel.
\item
  \textbf{Interpretasi penulis} adalah penjelasan penulis tentang makna,
  mekanisme, sebab, atau implikasi hasil.
\item
  \textbf{Inferensi review} adalah pembacaan analitis yang dibuat dari
  TXT J22. Inferensi ini bukan bukti tambahan dan tidak diperlakukan
  sebagai posisi penulis kecuali dinyatakan demikian.
\end{itemize}

\subsection{Summarized Abstract}\label{summarized-abstract-11}

\textbf{Laporan sumber.} Artikel membandingkan perkembangan spasial MMA,
JMA, dan DMA, tiga wilayah metropolitan yang mempunyai sejarah hubungan
dengan arus modal global tetapi bertumpu pada basis ekonomi yang
berbeda. MMA dipilih untuk basis perkebunan, JMA untuk manufaktur serta
kegiatan komersial, dan DMA untuk jasa pariwisata. Penelitian
mengidentifikasi subpusat pekerjaan, mengukur polisentralitas morfologis
dan fungsional, lalu membaca perbedaannya melalui sejarah
industrialisasi dan konteks pembangunan masing-masing wilayah (Abstract;
Introduction, hlm. 3; An overview of the case studies areas, hlm. 3-6).

\textbf{Laporan sumber.} Identifikasi subpusat menggunakan model
parametrik kepadatan pekerjaan terhadap jarak lurus dari pusat distrik
menuju CBD tradisional. Data utamanya adalah kegiatan usaha menengah dan
besar nonpertanian dari Sensus Ekonomi 2016. Polisentralitas morfologis
dihitung dari kemiringan regresi distribusi \emph{rank-size} jumlah
pekerjaan, sedangkan polisentralitas fungsional dihitung dengan
pendekatan analisis jejaring sosial Green dari survei komuter BPS 2014
untuk JMA dan 2015 untuk MMA serta DMA (Methods and data, hlm. 8-10).

\textbf{Laporan sumber.} Model standar mengidentifikasi 52 subdistrik
pusat ekonomi di JMA, 10 di MMA, dan satu di DMA. Seluruh pusat yang
teridentifikasi di MMA berada dalam radius 10 km dari CBD dan
ditafsirkan sebagai limpahan CBD yang menyatu, bukan subpusat mandiri.
Satu pusat DMA, Kuta, juga berbatasan langsung dengan CBD tradisional.
JMA mempunyai sedikitnya empat gugus pusat, termasuk Cikarang, Bogor,
dan koridor Tangerang-Cikupa-Balaraja di luar inti (Findings, hlm.
10-13).

\textbf{Laporan sumber.} Nilai morfologis JMA adalah -0,44 dan
digolongkan sangat polisentrik. DMA bernilai -1,12 atau rata-rata
polisentrik, sedangkan MMA bernilai -2,21 atau sangat monosentrik.
Indeks polisentralitas fungsional umum (\texttt{Pgf}) adalah 0,11 untuk
JMA dan 0,03 untuk MMA serta DMA. Penulis menilai JMA mempunyai
polisentralitas fungsional tingkat menengah, sementara MMA dan DMA tidak
polisentrik secara fungsional. Kesimpulan keseluruhan artikel adalah
hanya JMA yang mempunyai struktur spasial polisentrik (The comparative
polycentricity, hlm. 13-14; Conclusion, hlm. 19-20).

\textbf{Interpretasi penulis.} Perbedaan struktur tersebut dikaitkan
dengan lintasan industrialisasi, investasi infrastruktur, pengaturan
kelembagaan, dan basis ekonomi regional. Desentralisasi manufaktur serta
jaringan jalan tol dan rel dinilai menopang pembentukan subpusat JMA,
sedangkan skala industrialisasi MMA tetap terbatas dan ekonomi
pariwisata DMA tidak mendorong pusat baru yang jauh dari CBD. Penulis
menyatakan bahwa basis ekonomi mungkin menentukan bentuk kota, tetapi
juga mengakui bahwa bukti tambahan masih diperlukan untuk mendukung
hipotesis ini (Discussion, hlm. 17-19; Conclusion, hlm. 20).

\textbf{Inferensi review.} Hasil kuantitatif mendukung perbedaan
struktur ketiga kasus, tetapi tidak dengan sendirinya mengidentifikasi
industrialisasi sebagai penyebab. Hubungan sebab-akibat tersebut berasal
dari interpretasi komparatif dan narasi historis penulis, bukan model
kausal atau pengamatan longitudinal bentuk spasial.

\subsection{Problem Statement}\label{problem-statement-11}

\textbf{Laporan sumber.} Polisentralitas sering dipromosikan sebagai
bentuk perkotaan yang dapat menyeimbangkan manfaat dan kerugian
aglomerasi. Literatur yang dirangkum artikel mengaitkannya dengan
kepadatan serta pendapatan per kapita yang lebih tinggi, kemiskinan yang
lebih rendah, produktivitas tenaga kerja, efisiensi perjalanan, dan
ketahanan kota. Namun, studi lain melaporkan disparitas regional,
produktivitas yang lebih rendah, serta manfaat kemacetan dan perjalanan
yang bergantung pada ukuran kota dan keseimbangan antara pekerjaan
dengan pekerja. Dengan hasil yang beragam, polisentralitas tetap
bergeser dari alat analisis menjadi agenda normatif kebijakan spasial
(Introduction, halaman pembuka artikel-hlm. 2).

\textbf{Laporan sumber.} Bukti polisentralitas telah banyak dikembangkan
di Eropa, Amerika Serikat, dan Asia Timur, tetapi dinamika
polisentralitas Asia Tenggara masih sedikit diketahui. Kota-kota Asia
Tenggara memiliki sejarah hubungan dengan modal global, sedangkan studi
kontemporer dinilai lebih banyak membahas transisi permukiman daripada
akumulasi modal dan sistem produksi. Akibatnya, penyebab dan konsekuensi
perkembangan spasial polisentrik dalam konteks Asia Tenggara belum
terjelaskan dengan memadai (Introduction, hlm. 2-3).

\textbf{Laporan sumber.} Dalam konteks Indonesia, urbanisasi
mencerminkan kesenjangan regional dan bias kebijakan ekonomi ke wilayah
perkotaan. Kepadatan tinggi yang tidak diimbangi investasi infrastruktur
dikaitkan dengan kemacetan, polusi, dan risiko bencana. Pemahaman
struktur spasial metropolitan diposisikan sebagai masukan bagi strategi
perencanaan lahan untuk menghadapi persoalan tersebut (Introduction,
hlm. 3).

\textbf{Laporan sumber.} Kajian terdahulu tentang JMA menghasilkan
kesimpulan yang berbeda. Hakim dan Parolin (2009) menilai strukturnya
bukan monosentrik ataupun polisentrik, sedangkan Hudalah et al.~(2013)
serta Sadewo, Syabri, dan Pradono (2018) melihat pergeseran menuju
polisentralitas. Artikel mengaitkan perbedaan itu dengan cara pengukuran
dan mencatat bahwa studi tersebut hanya berfokus pada JMA serta tidak
menggabungkan dimensi morfologis dan fungsional (Introduction, hlm. 3).

\textbf{Rumusan masalah sumber.} Hubungan antara fungsi ekonomi dan
bentuk perkotaan JMA maupun metropolitan Indonesia lainnya masih sedikit
diketahui. Artikel menanyakan bagaimana subpusat dan tingkat
polisentralitas berbeda di MMA, JMA, dan DMA, serta bagaimana lintasan
pertumbuhan yang berbeda dalam sistem perencanaan yang serupa dapat
menghasilkan struktur spasial yang berbeda (Introduction, hlm. 3).

\subsection{Objectives}\label{objectives-11}

\begin{itemize}
\tightlist
\item
  Mengeksplorasi perkembangan spasial metropolitan Indonesia melalui
  perbandingan MMA, JMA, dan DMA (Abstract; Introduction, hlm. 3).
\item
  Mengidentifikasi subpusat pada masing-masing wilayah metropolitan
  (Introduction, hlm. 3).
\item
  Membandingkan tingkat polisentralitas morfologis dan fungsional ketiga
  wilayah (Introduction, hlm. 3; Methods and data, hlm. 9-10).
\item
  Membahas bagaimana lintasan pertumbuhan yang berbeda dalam sistem
  perencanaan yang serupa membentuk perkembangan spasial dengan cara
  yang berbeda (Introduction, hlm. 3).
\item
  Menggunakan sejarah serta basis pembangunan masing-masing kasus untuk
  menjelaskan variasi bentuk kontemporer, mengikuti strategi
  perbandingan individualisasi Robinson (2011), bukan mengasumsikan
  model universal yang sama bagi semua metropolis (Introduction, hlm.
  3).
\end{itemize}

\subsection{Literature Survey}\label{literature-survey-11}

\textbf{Sifat survei.} Artikel menyajikan tinjauan naratif untuk
membangun masalah, membandingkan metode identifikasi subpusat, memilih
ukuran polisentralitas, dan menafsirkan sejarah kasus. Protokol
penelusuran, basis data bibliografis, kata kunci, kriteria inklusi,
jumlah studi yang diperiksa, serta penilaian mutu studi: Tidak
disebutkan dalam file. Bagian ini tidak dapat diperlakukan sebagai
tinjauan sistematis.

\begin{itemize}
\tightlist
\item
  Literatur manfaat polisentralitas mengaitkan struktur ini dengan
  keseimbangan aglomerasi, kepadatan dan pendapatan per kapita yang
  lebih tinggi, tingkat kemiskinan yang lebih rendah, produktivitas
  tenaga kerja, efisiensi perjalanan, serta ketahanan kota. Artikel juga
  memuat bukti yang berlawanan tentang disparitas regional,
  produktivitas, dan kemacetan sehingga manfaat polisentralitas
  dinyatakan bergantung pada konteks (Introduction, halaman pembuka
  artikel-hlm. 2).
\item
  Literatur kebijakan menunjukkan perubahan polisentralitas dari konsep
  analitis menjadi agenda normatif. Contoh yang dibahas meliputi
  perencanaan regional dan nasional di Eropa, rencana metropolitan di
  beberapa kota Amerika Serikat, serta strategi pembangunan
  terdesentralisasi Beijing yang belum menghilangkan dominasi pekerjaan
  di pusat kota (Introduction, hlm. 2).
\item
  Definisi yang dipakai menyatakan bahwa struktur polisentrik terbentuk
  ketika satu atau lebih subpusat pekerjaan berkembang di luar CBD.
  Pendorong yang dibahas adalah desentralisasi kegiatan ekonomi,
  peningkatan mobilitas, pola komuter silang, dan distribusi kegiatan
  yang terfragmentasi. Artikel berfokus pada skala intrametropolitan,
  tetapi mengakui penerapan konsep pada skala antarkota dan antarwilayah
  (Identification of subcenters and measuring polycentricity, hlm. 6-7).
\item
  Literatur identifikasi subpusat dikelompokkan ke dalam puncak
  kepadatan pekerjaan, ambang kepadatan dan jumlah pekerjaan, metode
  parametrik, metode nonparametrik, serta varian lain seperti
  \emph{cubic spline}. Pendekatan proporsional dan statistik spasial
  juga dibahas (Identification of subcenters and measuring
  polycentricity, hlm. 6-7).
\item
  Pendekatan puncak kepadatan pekerjaan dikritik karena ambangnya
  arbitrer, dasar teorinya lemah, dan sifatnya statis. Pendekatan
  proporsional memakai rasio pekerjaan terhadap pekerja residen untuk
  mendekati tarikan pekerjaan dan komuter masuk. Artikel melaporkan
  bahwa penggabungan pendekatan dapat menghasilkan pengukuran yang lebih
  efisien (Identification of subcenters and measuring polycentricity,
  hlm. 7).
\item
  Metode parametrik dan nonparametrik mengenali kandidat subpusat
  melalui residual antara kepadatan aktual dan prediksi. Metode
  parametrik lebih sederhana, tetapi mereduksi wilayah menjadi hubungan
  satu dimensi antara kepadatan dan jarak CBD. Metode nonparametrik
  dapat menangkap puncak lokal yang lebih kompleks, tetapi sulit
  membedakan subpusat aktual dari puncak acak. ESDA mempertimbangkan
  autokorelasi serta heterogenitas spasial dan tidak memerlukan nilai
  ambang, tetapi membutuhkan data peta serta data lokasi agregat yang
  lengkap (Identification of subcenters and measuring polycentricity,
  hlm. 7).
\item
  Literatur pengukuran komparatif membedakan polisentralitas morfologis
  dari fungsional. Dimensi morfologis mencakup jumlah pusat, pemisahan,
  dan keseimbangan distribusi ukuran pusat. Kemiringan regresi
  \emph{rank-size} banyak dipakai sebagai proksi hierarki pusat. Nilai
  yang lebih datar atau mendekati nol menunjukkan susunan yang lebih
  polisentrik (Identification of subcenters and measuring
  polycentricity, hlm. 7-8).
\item
  Dimensi fungsional mensyaratkan kepadatan dan keseimbangan interaksi
  antarpusat. Pola komuter banyak digunakan untuk mengukurnya, walaupun
  hubungan fungsional juga dapat berupa arus barang atau informasi.
  Artikel memilih analisis jejaring sosial Green (2007), yang dapat
  digunakan pada beberapa skala tetapi masih mengandung subjektivitas
  dalam menentukan batas nilai maksimum (Identification of subcenters
  and measuring polycentricity, hlm. 8).
\item
  Literatur Asia Tenggara menempatkan bentuk kota dalam hubungan
  historis dengan modal global. Artikel memakai konsep ketergantungan
  lintasan dan strategi perbandingan individualisasi untuk menghubungkan
  bentuk kontemporer dengan basis perkebunan MMA, manufaktur JMA, dan
  pariwisata DMA (Introduction, hlm. 2-3; An overview of the case
  studies areas, hlm. 3-6).
\item
  Literatur kontekstual yang digunakan dalam pembahasan mengaitkan
  restrukturisasi kota dengan perubahan ekonomi, teknologi, dan sosial,
  serta dengan institusi, identitas regional, dan transportasi. Contoh
  lintas negara yang dirangkum penulis menunjukkan hubungan historis
  antara manufaktur, pemindahan pekerjaan, infrastruktur, dan
  pembentukan subpusat. Pernyataan tersebut merupakan laporan J22 atas
  literatur terdahulu, bukan bukti empiris langsung dari tiga kasus
  (Discussion, hlm. 17-19).
\end{itemize}

\textbf{Inferensi review.} Tinjauan memberi dasar yang cukup untuk
menjelaskan pilihan tiga ukuran, yaitu model kepadatan, kemiringan
\emph{rank-size}, dan indeks jejaring komuter. Namun, artikel tidak
menyajikan kerangka formal yang menguji mekanisme sejarah secara seragam
pada ketiga kasus. Literatur historis dipakai sebagai dasar interpretasi
setelah hasil kuantitatif diperoleh.

\subsection{Method Used}\label{method-used-11}

\textbf{Rancangan penelitian.} Penelitian adalah studi komparatif tiga
wilayah metropolitan. Komponen kuantitatifnya mengidentifikasi subpusat
dan membandingkan polisentralitas morfologis serta fungsional. Komponen
kontekstualnya menelusuri sejarah basis ekonomi, industrialisasi,
investasi, perencanaan, dan infrastruktur melalui data deskriptif serta
literatur yang dirujuk artikel. Penulis menyebut strategi
perbandingannya sebagai individualisasi karena tiap metropolis dibaca
menurut lintasan pembangunannya sendiri (Introduction, hlm. 3; Methods
and data, hlm. 8-10; The contextual perspective dan Discussion, hlm.
14-19).

\textbf{Pemilihan kasus.} MMA, JMA, dan DMA dipilih secara bertujuan
karena mewakili hubungan terkuat dengan jaringan akumulasi modal global
dalam tiga sektor berbeda: pertanian atau perkebunan, manufaktur, dan
jasa pariwisata. Dasarnya adalah pemeriksaan \emph{location quotient}
dan nilai ekspor 14 wilayah metropolitan Indonesia (An overview of the
case studies areas dan Table 1, hlm. 3-6).

\textbf{Identifikasi subpusat.}

\begin{itemize}
\tightlist
\item
  Unit model adalah subdistrik. Kepadatan pekerjaan dimodelkan sebagai
  fungsi aksesibilitas terhadap CBD masing-masing wilayah metropolitan
  dengan hipotesis fungsi kepadatan monosentrik klasik (Methods and
  data, hlm. 8-9).
\item
  Jarak aktual antarsubdistrik tidak tersedia. Penulis memakai jarak
  Euklides antara centroid setiap subdistrik dan titik CBD tradisional
  sebagai proksi. Titik CBD dipilih secara arbitrer di alun-alun pusat
  kota: Merdeka Square di Kecamatan Medan Barat, Monas di Gambir, dan
  Puputan-Badung Square di Denpasar Barat (Methods and data, hlm. 8-9).
\item
  Lima spesifikasi parametrik diuji: linear (\texttt{LIN}), eksponensial
  negatif (\texttt{EXP}), bentuk logaritmik dari eksponensial negatif
  atau model standar (\texttt{STD}), potensial terbalik atau gravitasi
  (\texttt{GRAV}), serta bentuk logaritmik gravitasi (\texttt{Log-Log})
  (Methods and data, hlm. 9).
\item
  Model terbaik dipilih dengan koefisien determinasi (\texttt{R2})
  sebagai ukuran kecocokan relatif dan \emph{root mean square error}
  (\texttt{RMSE}) sebagai ukuran absolut. Pemecahan simetri menurut
  sumbu transportasi utama dipertimbangkan apabila kecocokan awal tidak
  memadai (Methods and data, hlm. 9).
\item
  Kandidat subpusat adalah lokasi dengan residual positif dari model
  terbaik yang melampaui ambang kepadatan pekerjaan. Penulis
  melonggarkan ambang rujukan Giuliano dan Small menjadi 1.000 pekerja
  per km2 karena analisis hanya menggunakan usaha menengah dan besar
  atau MLE (Methods and data, hlm. 9).
\item
  Data input mencakup seluruh kegiatan MLE nonpertanian dari Sensus
  Ekonomi 2016 pada tingkat subdistrik. Pilihan ini dimaksudkan agar
  hasil dapat dibandingkan dengan studi JMA terdahulu sekaligus
  menghindari kesimpulan yang hanya bertumpu pada manufaktur (Methods
  and data, hlm. 9).
\end{itemize}

\textbf{Polisentralitas morfologis.}

\begin{itemize}
\tightlist
\item
  Penulis mengikuti Burger dan Meijers (2012) dengan mengestimasi
  kemiringan regresi distribusi \emph{rank-size}. Ukuran node adalah
  jumlah pekerjaan, sedangkan node adalah pusat setiap kota atau
  kabupaten di dalam wilayah metropolitan (Methods and data, hlm. 9-10).
\item
  Karena jumlah node antarkasus berbeda dan lebih banyak node cenderung
  menghasilkan nilai polisentralitas lebih tinggi, regresi dihitung
  untuk dua, tiga, dan empat lokasi dengan jumlah pekerjaan terbesar.
  Rata-rata ketiga kemiringan menjadi nilai morfologis akhir. Kemiringan
  yang lebih datar atau mendekati nol ditafsirkan lebih polisentrik
  (Methods and data, hlm. 9-10).
\item
  MMA semula mempunyai node paling sedikit. Penulis membagi Kabupaten
  Deli Serdang menjadi bagian timur dan barat berdasarkan posisi relatif
  terhadap 98,65 derajat bujur. Masing-masing bagian berisi 12
  subdistrik (Methods and data, hlm. 10).
\end{itemize}

\textbf{Polisentralitas fungsional.}

\begin{itemize}
\tightlist
\item
  Penulis menerapkan analisis jejaring sosial Green (2007) pada matriks
  asal-tujuan komuter untuk menghitung indeks polisentralitas fungsional
  khusus arus masuk dan keluar (\texttt{Psf}) serta indeks umum
  (\texttt{Pgf}) (Methods and data, hlm. 10; Table 4, hlm. 14).
\item
  Secara teoretis, \texttt{Pgf} berada antara 0 dan 1. Artikel memakai
  nilai empiris 0,19 dari kawasan Eropa sebagai pembanding untuk kawasan
  yang dinilai cukup polisentrik secara fungsional (Methods and data,
  hlm. 10).
\item
  Survei komuter dilaksanakan pada 2014 untuk JMA dan 2015 untuk MMA
  serta DMA. Dalam definisi Indonesia yang digunakan, komuter adalah
  perjalanan harian untuk bekerja atau sekolah yang melintasi batas
  administratif kota atau kabupaten tempat tinggal. Karena itu, matriks
  asal-tujuan dibentuk pada batas tingkat kota atau kabupaten, bukan
  seluruh perjalanan antarsubdistrik (Methods and data, hlm. 10).
\end{itemize}

\textbf{Analisis kontekstual.} Penulis membandingkan sejarah hubungan
modal global, waktu pertumbuhan MLE, basis ekonomi, desentralisasi
industri, kebijakan tata ruang, investasi infrastruktur, identitas
regional, dan peran pemerintah serta sektor swasta. Protokol pengumpulan
data kualitatif, teknik pengodean, daftar dokumen sejarah yang dicari
secara sistematis, dan prosedur triangulasi: Tidak disebutkan dalam
file. Dengan demikian, bagian ini adalah interpretasi historis
komparatif, bukan studi kualitatif dengan prosedur yang dilaporkan
lengkap.

\subsection{Dataset}\label{dataset-11}

\textbf{Data utama.} Penelitian menggunakan data agregat kegiatan usaha
menengah dan besar nonpertanian dari Sensus Ekonomi 2016 BPS pada
tingkat subdistrik. Artikel menyatakan bahwa seluruh kegiatan MLE
nonpertanian digunakan, bukan hanya kegiatan manufaktur. Data tersebut
menjadi dasar untuk menghitung kepadatan pekerjaan, mengidentifikasi
subpusat, dan membentuk ukuran polisentralitas morfologis (Methods and
data, hlm. 9-10).

\textbf{Cakupan MLE.} MLE hanya mencakup sebagian kecil unit usaha,
tetapi menyumbang bagian pekerjaan yang cukup besar. Di MMA, 2,61\% unit
MLE mencakup 30,92\% pekerjaan; di JMA, 4,61\% unit MLE mencakup 47,64\%
pekerjaan; dan di DMA, 4,09\% unit MLE mencakup 34,94\% pekerjaan.
Jumlah absolut MLE yang dianalisis: Tidak disebutkan dalam file. Table 2
hanya melaporkan jumlah seluruh kegiatan ekonomi nonpertanian dan
persentase unit MLE (Methods and data dan Table 2, hlm. 9-10).

\textbf{Data komuter.} Data fungsional berasal dari survei komuter BPS
tahun 2014 untuk JMA serta tahun 2015 untuk MMA dan DMA. Lokasi
pekerjaan dalam matriks asal-tujuan dibatasi menurut kota atau kabupaten
karena definisi komuter yang digunakan adalah perjalanan harian untuk
bekerja atau sekolah yang melintasi batas administratif kota atau
kabupaten tempat tinggal (Methods and data, hlm. 10).

\textbf{Data pemilihan dan deskripsi kasus.} Pemilihan tiga kasus
didasarkan pada pengolahan \emph{location quotient} dan nilai ekspor 14
wilayah metropolitan dengan data PDB regional serta ekspor 2016 dari
BPS. Artikel juga menggunakan data populasi, luas, kepadatan, penduduk
bekerja, struktur ekonomi, investasi, ekspor, infrastruktur, dan
pertumbuhan MLE untuk menggambarkan konteks kasus. Daftar lengkap
variabel, berkas data, tautan unduhan, prosedur pembersihan, dan kode
analisis: Tidak disebutkan dalam file (An overview of the case studies
areas, Table 1, dan Table 2, hlm. 3-6 dan 10; The contextual
perspective, hlm. 14-17).

\textbf{Keselarasan waktu.} Data inti tidak berasal dari satu tahun yang
sama: pekerjaan dan kegiatan usaha berasal dari 2016, survei komuter JMA
dari 2014, serta survei MMA dan DMA dari 2015. Artikel melaporkan
perbedaan tahun tersebut, tetapi tidak menjelaskan prosedur harmonisasi
temporal (Methods and data, hlm. 9-10).

\subsection{Population / Sample}\label{population-sample-11}

\textbf{Populasi sasaran kasus.} Unit geografis utama adalah wilayah
metropolitan Medan, Jakarta, dan Denpasar. Ketiganya dipilih dari 14
wilayah metropolitan yang ditetapkan pemerintah karena mewakili hubungan
kuat dengan akumulasi modal global berbasis perkebunan, manufaktur, dan
jasa pariwisata. Pemilihan ini bersifat bertujuan, bukan sampel acak
metropolitan Indonesia (An overview of the case studies areas dan Table
1, hlm. 3-6).

\textbf{Cakupan geografis.} MMA dalam analisis merujuk pada Mebidang,
yaitu Kota Medan, Kota Binjai, dan Kabupaten Deli Serdang; bagian
Kabupaten Karo dikeluarkan karena terletak paling jauh dari inti dan
hanya sebagian wilayahnya termasuk Mebidangro. JMA mencakup 183
subdistrik di DKI Jakarta dan wilayah sekelilingnya di Jawa Barat serta
Banten. DMA mencakup 27 subdistrik di Kota Denpasar serta Kabupaten
Badung, Gianyar, dan Tabanan (An overview of the case studies areas,
hlm. 4-6; Table 2, hlm. 10).

\textbf{Unit observasi kuantitatif.} Untuk model identifikasi subpusat,
unit observasinya adalah subdistrik dengan data kepadatan pekerjaan MLE
nonpertanian. Jumlah subdistrik MMA yang benar-benar masuk ke model
setelah Karo dikeluarkan: Tidak disebutkan secara eksplisit dalam file.
Untuk ukuran morfologis, node adalah pusat kota atau kabupaten; analisis
memakai dua, tiga, dan empat node dengan jumlah pekerjaan terbesar pada
setiap kasus. Deli Serdang dibagi menjadi bagian timur dan barat agar
jumlah node MMA dapat dibandingkan dengan kasus lain (Methods and data,
hlm. 8-10).

\textbf{Sampel survei komuter.} Survei JMA mencakup 13.120 rumah tangga.
Survei MMA dan DMA masing-masing mencakup 2.400 rumah tangga. Kerangka
sampel, metode penarikan sampel, tingkat respons, pembobotan, galat
sampel, dan karakteristik responden: Tidak disebutkan dalam file
(Methods and data, hlm. 10).

\textbf{Besaran populasi wilayah.} Table 2 melaporkan populasi sekitar
4,239 juta untuk MMA, 30,970 juta untuk JMA, dan 2,428 juta untuk DMA.
Catatan tabel menyatakan angka JMA berasal dari 2014, sedangkan MMA dan
DMA dari 2015. Bagian uraian kasus menyebut populasi JMA 30,97 juta pada
2015 dan populasi DMA 2,5 juta tanpa tahun yang dinyatakan pada kalimat
tersebut; perbedaan keterangan tahun dan pembulatan ini tidak
diselesaikan di dalam file (An overview of the case studies areas, hlm.
4-6; Table 2, hlm. 10).

\subsection{Dependent Variables}\label{dependent-variables-11}

\textbf{Variabel dependen yang dinyatakan melalui model.} Kepadatan
pekerjaan MLE nonpertanian pada setiap subdistrik dimodelkan sebagai
fungsi aksesibilitas atau jarak terhadap CBD. Lima bentuk fungsi diuji.
Artikel tidak menyajikan satu model kausal utama; model parametrik
tersebut digunakan untuk memperkirakan pola peluruhan kepadatan dan
mengenali residual positif sebagai kandidat subpusat (Methods and data,
hlm. 8-9; Table 3, hlm. 11).

\textbf{Hasil ukur morfologis.} Kemiringan regresi distribusi
\emph{rank-size} jumlah pekerjaan digunakan sebagai ukuran
polisentralitas morfologis. Nilai akhir adalah rata-rata kemiringan dari
spesifikasi yang memakai dua, tiga, dan empat lokasi terbesar. Artikel
memperlakukan nilai yang lebih dekat ke nol sebagai struktur yang lebih
polisentrik (Methods and data, hlm. 9-10).

\textbf{Hasil ukur fungsional.} Indeks \texttt{Psf} untuk arus masuk dan
keluar serta indeks umum \texttt{Pgf} dihitung dari hubungan asal-tujuan
komuter. Nilai yang lebih tinggi menunjukkan jejaring yang lebih
polisentrik; artikel membandingkan \texttt{Pgf} kasus dengan nilai
empiris maksimum 0,19 yang diambil dari studi wilayah Eropa (Methods and
data, hlm. 10; Table 4, hlm. 14).

\textbf{Hasil lain yang dianalisis.} Jumlah dan lokasi subpusat,
proporsi komuter tradisional, komuter terbalik, dan komuter silang juga
menjadi hasil deskriptif utama. Basis ekonomi, industrialisasi,
infrastruktur, dan institusi tidak dimasukkan sebagai variabel penjelas
dalam model statistik. Hubungannya dengan bentuk spasial dibahas melalui
perbandingan historis dan interpretasi penulis (Findings dan Discussion,
hlm. 10-19).

\subsection{Results}\label{results-11}

\textbf{Kinerja model.} Model standar (\texttt{STD}) dipilih untuk
ketiga wilayah. Untuk JMA, model eksponensial menghasilkan \texttt{R2}
sedikit lebih tinggi daripada model standar, yaitu 55,39 berbanding
54,55, tetapi model standar mempunyai \texttt{RMSE} yang jauh lebih
rendah, yaitu 178,7 pada Table 3 berbanding 190.200 untuk model
eksponensial. Narasi artikel menyebut \texttt{RMSE} standar JMA sebesar
178,8 sehingga terdapat selisih 0,1 dengan tabel. Untuk MMA dan DMA,
model standar memberikan kombinasi \texttt{R2} serta \texttt{RMSE}
terbaik. Nilai model standar yang dilaporkan adalah \texttt{R2} 59,26
dan \texttt{RMSE} 1,95 untuk MMA, \texttt{R2} 54,55 serta \texttt{RMSE}
178,7 pada tabel untuk JMA, dan \texttt{R2} 84,26 serta \texttt{RMSE}
0,93 untuk DMA (Findings dan Table 3, hlm. 10-11).

\textbf{Subpusat MMA.} Model standar menandai 10 subdistrik sebagai
pusat ekonomi. Semuanya berada di Kota Medan, dalam jarak 10 km dari
CBD. Penulis menafsirkan susunan ini sebagai perluasan atau limpahan
CBD, bukan sepuluh subpusat yang berkembang secara mandiri. Karena itu,
perekonomian MMA tetap dinilai didominasi satu pusat (Subcenters in
urban development, hlm. 11-12; Conclusion, hlm. 19).

\textbf{Subpusat DMA.} Model standar hanya mengidentifikasi Kecamatan
Kuta sebagai pusat pekerjaan utama. Kuta bersebelahan dengan CBD
tradisional di Denpasar Barat. Meskipun lokasinya berada di luar CBD
tradisional dan inti Kota Denpasar, kedekatan tersebut membuat penulis
menilai DMA masih mempunyai struktur yang terpusat (Subcenters in urban
development, hlm. 11-12; Conclusion, hlm. 19).

\textbf{Subpusat JMA.} Model standar mengidentifikasi 52 subdistrik
sebagai pusat ekonomi dan sedikitnya empat gugus subpusat. Satu gugus
berada di inti, sedangkan gugus lain berada sekitar 20-60 km dari CBD,
terutama koridor Cikarang di timur, Kota Bogor di selatan, serta koridor
Tangerang-Cikupa-Balaraja di barat. Berbeda dari dua penelitian JMA
terdahulu yang dibahas artikel, metode parametrik ini juga mengenali
Kota Bogor sebagai subpusat (Subcenters in urban development, hlm.
11-13; Figure 2).

\textbf{Polisentralitas morfologis.} Nilai kemiringan rata-rata JMA
sebesar -0,44 diklasifikasikan sebagai sangat polisentrik. Nilai DMA
sebesar -1,12 diklasifikasikan sebagai rata-rata polisentrik karena
pusat pekerjaan bergeser ke Kabupaten Badung di luar inti tradisional.
Nilai MMA sebesar -2,21 diklasifikasikan sebagai sangat monosentrik
karena kesenjangan pekerjaan antara inti dan pinggiran lebih besar (The
comparative polycentricity, hlm. 13).

\textbf{Pola komuter.} Rasio komuter terhadap angkatan kerja adalah 0,18
di MMA, 0,19 di JMA, dan 0,14 di DMA. Jumlah komuter harian yang
dilaporkan mencapai sekitar 315 ribu di MMA, 3,5 juta di JMA, dan 187
ribu di DMA. Komuter tradisional, terbalik, dan silang masing-masing
mencakup 80,63\%, 16,83\%, dan 2,54\% di MMA; 56,41\%, 10,20\%, dan
33,39\% di JMA; serta 40,11\%, 41,71\%, dan 18,18\% di DMA (The
comparative polycentricity dan Table 4, hlm. 13-14).

\textbf{Polisentralitas fungsional.} Nilai \texttt{Psf} arus masuk dan
keluar berturut-turut adalah 0,03 dan 0,04 untuk MMA, 0,10 dan 0,12
untuk JMA, serta 0,01 dan 0,07 untuk DMA. Nilai \texttt{Pgf} adalah 0,03
untuk MMA, 0,11 untuk JMA, dan 0,03 untuk DMA. Berdasarkan pembanding
yang digunakan artikel, MMA dan DMA tidak membentuk jejaring permukiman
yang polisentrik secara fungsional, sedangkan JMA berada pada tingkat
polisentralitas fungsional menengah (The comparative polycentricity dan
Table 4, hlm. 14).

\textbf{Hasil kontekstual deskriptif.} Artikel melaporkan bahwa
pertumbuhan MLE JMA mulai meningkat tajam sejak kemerdekaan, sekitar dua
dekade lebih awal daripada MMA dan DMA, yang peningkatannya lebih
terkait dengan keterbukaan investasi asing pada awal 1970-an. JMA juga
mempunyai jaringan jalan tol dan rel yang lebih luas. MMA mengalami
perluasan jaringan jalan tol yang lambat, sedangkan jalan tol DMA baru
beroperasi pada 2013 dengan jangkauan terbatas dan DMA tidak mempunyai
sistem kereta api (The contextual perspective, hlm. 15-17; Figure 3;
Discussion, hlm. 18-19).

\subsection{Findings}\label{findings-11}

\textbf{Temuan yang dilaporkan sumber.} Hanya JMA yang memenuhi gambaran
struktur spasial polisentrik secara keseluruhan. JMA mempunyai banyak
gugus subpusat, morfologi yang sangat polisentrik, dan jejaring komuter
dengan polisentralitas fungsional menengah. MMA tetap sangat terpusat.
DMA lebih seimbang daripada MMA secara morfologis, tetapi pusat
pekerjaannya berdekatan dengan CBD dan nilai \texttt{Pgf}-nya sama
rendahnya dengan MMA. Karena itu, label morfologis ``rata-rata
polisentrik'' untuk DMA tidak diartikan sebagai bukti bahwa struktur
keseluruhannya polisentrik (The comparative polycentricity, hlm. 13-14;
Conclusion, hlm. 19-20).

\textbf{Temuan yang dilaporkan sumber.} Arah komuter berkaitan dengan
lokasi pusat pekerjaan. MMA didominasi perjalanan tradisional dari
pinggiran ke inti. JMA mempunyai pembagian yang lebih seimbang antara
perjalanan tradisional dan silang karena keberadaan beberapa subpusat.
DMA mempunyai komuter terbalik dan silang sebesar 59,89\%, tetapi
penulis memperingatkan bahwa angka itu tidak menunjukkan kompleksitas
jejaring dengan sendirinya karena CBD DMA terletak di kawasan yang
secara administratif dianggap pinggiran dan bersebelahan dengan inti
(The comparative polycentricity, hlm. 13-14; Conclusion, hlm. 19-20).

\textbf{Interpretasi penulis.} Perbedaan bentuk spasial tidak dianggap
semata-mata sebagai akibat ukuran penduduk atau ekonomi. Penulis
menghubungkan polisentralitas JMA dengan industrialisasi manufaktur,
perpindahan industri ke pinggiran, investasi infrastruktur, dan
pengembangan lahan industri oleh sektor swasta. Industrialisasi MMA
dinilai tidak mencapai skala serupa, sedangkan ketergantungan DMA pada
pariwisata dan keterjangkauan objek wisata dari CBD mengurangi dorongan
bagi pembentukan subpusat yang jauh (The contextual perspective dan
Discussion, hlm. 14-19).

\textbf{Interpretasi penulis.} Struktur polisentrik JMA tidak terbentuk
sebagai hasil langsung strategi perencanaan polisentrik. Penulis
menekankan bias investasi pemerintah ke ibu kota, lokasi strategis JMA,
pembangunan jaringan transportasi, serta kendali sektor swasta atas
kompleks industri dan konversi lahan pinggiran. Dengan demikian,
lintasan industrialisasi dinilai membentuk metropolitan secara berbeda,
terlepas dari strategi tata ruang formal (Discussion, hlm. 18-19;
Conclusion, hlm. 20).

\textbf{Inferensi review.} Perbandingan tiga kasus memperlihatkan
kesesuaian antara pola pusat pekerjaan dan pola komuter, tetapi bukti
yang tersedia tidak mengisolasi pengaruh manufaktur, infrastruktur,
institusi, lokasi strategis, dan ukuran populasi. Oleh sebab itu,
pernyataan tentang mekanisme industrialisasi harus dibaca sebagai
penjelasan historis penulis yang masuk akal untuk kasus-kasus ini, bukan
estimasi efek kausal.

\subsection{Conclusions}\label{conclusions-11}

Penulis menyimpulkan bahwa hanya JMA di antara tiga wilayah studi yang
memiliki struktur spasial polisentrik. Pusat pekerjaan JMA tersebar di
beberapa lokasi pinggiran dan menghasilkan arah komuter yang lebih
seimbang. MMA dan DMA tetap terpusat, tetapi sifat monosentrisnya
berbeda: pusat MMA berada dekat CBD tradisional, sedangkan pusat DMA
bergeser ke luar CBD tradisional dan inti kota tetapi tetap bersebelahan
dengannya (Conclusion, hlm. 19-20).

Penulis juga menyimpulkan bahwa perbedaan morfologi berjalan sejajar
dengan perbedaan polisentralitas fungsional dan lokasi CBD berhubungan
erat dengan pola komuter. Karena dimensi morfologis dan fungsional dapat
memberi gambaran yang tidak identik, keduanya perlu digunakan bersama
dalam studi polisentralitas (Conclusion, hlm. 19-20).

Dalam penjelasan kontekstual, penulis berargumen bahwa polisentralitas
tidak harus bergantung pada jumlah penduduk atau ukuran ekonomi.
Dibandingkan pertanian dan pariwisata, manufaktur dinilai lebih efektif
mendorong transformasi perkotaan Indonesia. Namun, artikel secara
eksplisit menyatakan bahwa bukti tambahan diperlukan untuk mendukung
hipotesis bahwa basis ekonomi regional menentukan bentuk kota
(Conclusion, hlm. 20).

Penulis menegaskan bahwa subpusat kompleks industri JMA bukan hasil
langsung strategi perencanaan pemerintah. Investasi yang berpihak pada
JMA, infrastruktur transportasi, lokasi strategis, serta ekspansi
kompleks industri dan konversi lahan yang dikendalikan sektor swasta
ikut membentuk hasil tersebut. Penyediaan infrastruktur dinilai dapat
membantu mengejar bentuk kota yang diinginkan, tetapi pengalaman MMA dan
DMA menunjukkan bahwa infrastruktur atau strategi formal saja tidak
menjamin lintasan yang sama (Conclusion, hlm. 20).

\subsection{Contributions}\label{contributions-11}

\begin{itemize}
\tightlist
\item
  \textbf{Kontribusi empiris yang dinyatakan penulis.} Artikel
  menyediakan bukti yang disebut langka mengenai perkembangan perkotaan
  polisentrik di Asia Tenggara melalui tiga metropolitan Indonesia
  dengan basis ekonomi dan lintasan globalisasi yang berbeda
  (Introduction, hlm. 3; Conclusion, hlm. 19).
\item
  \textbf{Kontribusi komparatif.} Penelitian memperluas fokus dari JMA
  saja ke perbandingan MMA, JMA, dan DMA serta mengukur spektrum
  polisentralitas, bukan sekadar memberi label biner pada satu wilayah
  (Introduction, hlm. 3; The comparative polycentricity, hlm. 13-14).
\item
  \textbf{Kontribusi metodologis.} Artikel menggabungkan identifikasi
  subpusat, ukuran morfologis \emph{rank-size}, dan indeks jejaring
  komuter. Pendekatan parametriknya mengenali Kota Bogor sebagai
  subpusat yang tidak teridentifikasi oleh studi JMA terdahulu yang
  dibahas artikel (Methods and data, hlm. 8-10; Subcenters in urban
  development, hlm. 11-13).
\item
  \textbf{Kontribusi konseptual.} Hasil DMA menunjukkan bahwa morfologi
  yang relatif seimbang tidak otomatis berarti jejaring fungsional yang
  polisentrik. Artikel menggunakan perbedaan itu untuk menegaskan
  perlunya membaca morfologi dan fungsi secara bersama (The comparative
  polycentricity, hlm. 13-14; Conclusion, hlm. 20).
\item
  \textbf{Kontribusi kontekstual.} Artikel memasukkan sejarah
  industrialisasi, basis ekonomi, infrastruktur, institusi, dan hubungan
  dengan modal global untuk menjelaskan mengapa wilayah yang berada
  dalam sistem perencanaan nasional serupa menghasilkan bentuk yang
  berbeda (Introduction, hlm. 3; Discussion, hlm. 17-19).
\item
  \textbf{Inferensi review.} Kontribusi terkuat terletak pada diagnosis
  komparatif struktur aktual. Kontribusi mengenai penyebab lebih
  tentatif karena mekanisme historis tidak diuji dengan rancangan kausal
  atau prosedur kualitatif yang dilaporkan secara sistematis.
\end{itemize}

\subsection{Research Gap}\label{research-gap-11}

\textbf{Kesenjangan yang hendak diisi.} Dinamika polisentralitas Asia
Tenggara kurang diteliti dibandingkan Eropa, Amerika Serikat, dan Asia
Timur. Studi kontemporer Asia Tenggara dinilai lebih berfokus pada
transisi permukiman daripada akumulasi modal dan sistem produksi
sehingga penyebab serta konsekuensi perkembangan spasial polisentrik
belum cukup dijelaskan (Introduction, hlm. 2-3).

\textbf{Kesenjangan Indonesia.} Pengetahuan tentang hubungan fungsi
ekonomi dan bentuk perkotaan JMA maupun metropolitan Indonesia lainnya
masih terbatas. Kajian JMA terdahulu menghasilkan klasifikasi yang
berbeda, hanya memeriksa JMA, dan tidak menggabungkan aspek morfologis
dengan fungsional. Artikel ini menanggapi kesenjangan tersebut melalui
perbandingan tiga metropolitan dan kombinasi metode (Introduction, hlm.
3).

\textbf{Kesenjangan metode.} Artikel menyatakan bahwa banyak studi
membandingkan polisentralitas antardaerah, tetapi hanya sedikit yang
menggabungkan pendekatan morfologis dan fungsional. Perbedaan potensial
antara kedua dimensi membuat kesimpulan berbasis satu ukuran saja tidak
memadai (Identification of subcenters and measuring polycentricity, hlm.
7-8; Methods and data, hlm. 9).

\textbf{Kesenjangan yang tetap terbuka.} Hubungan antara basis ekonomi
regional dan bentuk kota masih berupa hipotesis yang memerlukan bukti
tambahan. Selain itu, relatif sedikit penelitian menjelaskan bagaimana
wilayah kota polisentrik berkembang dan bagaimana mekanisme tingkat
mikro serta makro saling berhubungan. Cakupan tiga metropolitan juga
belum memungkinkan kesimpulan komprehensif untuk seluruh Indonesia atau
perbandingan yang mengendalikan ukuran kota (Conclusion, hlm. 20).

\subsection{Limitations}\label{limitations-11}

\textbf{Keterbatasan yang dinyatakan penulis.} Identifikasi morfologi
hanya menggunakan MLE. Walaupun MLE menyumbang bagian pekerjaan yang
besar, jumlah unitnya lebih sedikit daripada usaha mikro dan kecil.
Karena itu, penulis menyatakan bahwa gambaran morfologi yang dihasilkan
hanya mencakup sebagian dari struktur ekonomi metropolitan (Conclusion,
hlm. 20).

\textbf{Keterbatasan yang dinyatakan penulis.} Penelitian hanya
membandingkan tiga wilayah metropolitan. Penulis menilai cakupan
tersebut belum cukup untuk menghasilkan kesimpulan yang komprehensif dan
belum mengendalikan variasi ukuran serta jalur pembangunan melalui kasus
pembanding di negara berkembang lain (Conclusion, hlm. 20).

\textbf{Keterbatasan yang dinyatakan penulis.} Pengetahuan mengenai
proses terbentuknya wilayah kota polisentrik masih terbatas. Artikel
menyerukan gabungan metode empiris kuantitatif dan kualitatif untuk
menjelaskan mekanisme tingkat mikro dan makro. Seruan ini juga
menunjukkan bahwa analisis historis dalam artikel belum menguji
mekanisme tersebut secara lengkap (Conclusion, hlm. 20).

\textbf{Batas metode yang diungkapkan dalam file.} Jarak perjalanan
aktual antarsubdistrik tidak tersedia sehingga penelitian memakai jarak
lurus dari centroid ke CBD. Titik CBD dipilih secara arbitrer pada
alun-alun tradisional. Model parametrik juga mereduksi geografi kota
menjadi hubungan kepadatan dan jarak terhadap CBD, suatu kelemahan
pendekatan yang diakui dalam tinjauan metode artikel (Identification of
subcenters and measuring polycentricity, hlm. 7; Methods and data, hlm.
8-9).

\textbf{Batas perbandingan morfologis.} Jumlah node asli berbeda
antarkasus. Penulis mengatasinya dengan memakai dua, tiga, dan empat
lokasi terbesar serta membagi Deli Serdang menjadi bagian timur dan
barat. Prosedur ini membuat perbandingan mungkin dilakukan, tetapi
hasilnya bergantung pada pilihan jumlah node dan pembagian geografis
tersebut (Methods and data, hlm. 9-10). Kalimat terakhir merupakan
inferensi review dari prosedur yang dilaporkan.

\textbf{Batas data fungsional.} Definisi komuter hanya menangkap
perjalanan harian yang melintasi batas kota atau kabupaten. Perjalanan
kerja atau sekolah di dalam batas administratif yang sama tidak masuk ke
matriks asal-tujuan. Implikasinya terhadap nilai \texttt{Pgf}: Tidak
dibahas dalam file. Pernyataan bahwa arus internal tidak tercakup
merupakan inferensi review langsung dari definisi operasional (Methods
and data, hlm. 10).

\textbf{Batas temporal dan tolok ukur.} Data pekerjaan berasal dari
2016, sedangkan data komuter berasal dari 2014 dan 2015. Artikel tidak
melaporkan harmonisasi temporal. Selain itu, interpretasi \texttt{Pgf}
menggunakan tolok ukur empiris 0,19 dari Eropa, sementara tinjauan
artikel sendiri menyebut penentuan batas maksimum dalam pendekatan ini
mengandung unsur subjektif (Identification of subcenters and measuring
polycentricity, hlm. 8; Methods and data, hlm. 10). Potensi pengaruh
kedua hal tersebut pada perbandingan adalah inferensi review dan tidak
dikuantifikasi oleh penulis.

\textbf{Batas penjelasan kausal.} Hubungan antara basis ekonomi,
industrialisasi, infrastruktur, institusi, dan bentuk spasial disusun
dari perbandingan tiga kasus serta narasi historis. Artikel tidak
melaporkan model kausal, desain longitudinal bentuk spasial, protokol
analisis kualitatif, atau uji alternatif mekanisme. Karena itu, hasil
kuantitatif menetapkan perbedaan struktur, tetapi tidak dengan
sendirinya membuktikan penyebab perbedaan tersebut. Ini adalah penilaian
metodologis dalam review, bukan keterbatasan yang dirumuskan dengan
kata-kata yang sama oleh penulis.

\subsection{Challenges}\label{challenges-11}

Artikel tidak menyediakan bagian khusus berjudul tantangan. Tantangan
berikut disarikan hanya dari persoalan yang disebutkan dalam TXT.

\begin{itemize}
\tightlist
\item
  \textbf{Operasionalisasi subpusat.} Literatur memakai beragam definisi
  dan metode dengan kelemahan berbeda. Ambang kepadatan dapat arbitrer,
  metode parametrik menyederhanakan geografi, metode nonparametrik dapat
  menghasilkan puncak acak, dan ESDA membutuhkan data spasial lengkap
  (Identification of subcenters and measuring polycentricity, hlm. 6-8).
\item
  \textbf{Keterbandingan wilayah.} Ketiga metropolitan berbeda dalam
  jumlah node, luas, populasi, basis ekonomi, dan tahap pembangunan.
  Penulis menstandardisasi jumlah node untuk ukuran morfologis, tetapi
  perbedaan ukuran dan konteks tetap menjadi persoalan dalam menafsirkan
  variasi (Methods and data, hlm. 9-10; The contextual perspective, hlm.
  14).
\item
  \textbf{Keterbatasan aksesibilitas.} Tidak tersedianya jarak aktual
  memaksa penggunaan jarak Euklides, sementara letak CBD tradisional
  harus ditentukan secara arbitrer. Hal ini membatasi kemampuan model
  untuk menangkap jaringan transportasi dan aksesibilitas nyata (Methods
  and data, hlm. 8-9). Klausa terakhir adalah inferensi review.
\item
  \textbf{Batas administratif komuter.} Definisi komuter Indonesia
  berbeda dari definisi di Amerika Serikat dan beberapa negara Eropa.
  Ketergantungan pada lintas batas kota atau kabupaten membuat hasil
  sensitif terhadap delineasi administratif, terutama di DMA yang
  CBD-nya berada di kawasan yang dikategorikan sebagai pinggiran
  (Methods and data, hlm. 10; The comparative polycentricity, hlm.
  13-14).
\item
  \textbf{Pemisahan mekanisme.} Artikel menyatakan bahwa sejarah dan
  geografi, waktu dan ruang, sulit dipisahkan. Ukuran populasi,
  industrialisasi, infrastruktur, kebijakan, lokasi strategis, identitas
  regional, dan tindakan sektor swasta saling berkaitan sehingga
  kontribusi masing-masing tidak dapat diisolasi melalui desain
  penelitian ini (The contextual perspective, hlm. 14; Discussion, hlm.
  17-19).
\item
  \textbf{Tantangan kebijakan dalam kasus.} JMA mengalami pembangunan
  pinggiran yang dikendalikan kuat oleh sektor swasta dan memadukan
  ruang terencana dengan tidak terencana. MMA serta DMA mendapat
  dukungan infrastruktur yang lebih lambat atau terbatas. Artikel
  mengaitkan kondisi itu dengan kesulitan menarik manufaktur dan
  membentuk subpusat baru di luar inti (Discussion, hlm. 18-19).
\end{itemize}

\subsection{Practical Implications}\label{practical-implications-11}

\textbf{Implikasi yang dinyatakan atau langsung diturunkan penulis.}
Pengukuran struktur metropolitan sebaiknya tidak hanya menggunakan
persebaran pusat pekerjaan atau arus komuter. Perbedaan antara nilai
morfologis DMA dan klasifikasi fungsionalnya menunjukkan perlunya
menggabungkan kedua dimensi sebelum menyimpulkan bahwa suatu wilayah
polisentrik (Identification of subcenters and measuring polycentricity,
hlm. 7-8; Conclusion, hlm. 20).

\textbf{Implikasi perencanaan.} Penulis menyatakan bahwa pemahaman
organisasi spasial dapat membantu penyusunan strategi perencanaan lahan.
Namun, kasus JMA menunjukkan bahwa bentuk polisentrik dapat muncul
melalui desentralisasi manufaktur, investasi yang terkonsentrasi,
infrastruktur, dan tindakan sektor swasta, bukan karena strategi tata
ruang formal. Oleh karena itu, rencana multipusat perlu dibaca bersama
dengan kebijakan industri, investasi transportasi, tata kelola lahan,
dan lintasan ekonomi setempat (Introduction, hlm. 3; Discussion, hlm.
18-19; Conclusion, hlm. 20). Kalimat terakhir merupakan inferensi review
yang dibatasi pada mekanisme yang dibahas artikel.

\textbf{Implikasi infrastruktur.} Artikel menyatakan bahwa penyediaan
transportasi dapat membantu mengejar bentuk perkotaan yang diinginkan.
Meskipun demikian, pengalaman MMA dan DMA menunjukkan bahwa meniru
infrastruktur atau tujuan manufaktur JMA tidak otomatis menghasilkan
polisentralitas karena kebijakan nasional, lokasi strategis, skala
investasi, dan basis ekonomi berbeda (Discussion, hlm. 18-19;
Conclusion, hlm. 20).

\textbf{Batas implikasi.} Artikel tidak menguji dampak kebijakan
polisentralitas terhadap kemacetan, produktivitas, kemiskinan,
ketahanan, atau kesejahteraan pada ketiga kasus. Karena itu, hasilnya
mendukung diagnosis struktur spasial dan pembacaan lintasan pembangunan,
bukan bukti bahwa penerapan polisentralitas akan menghasilkan manfaat
sosial, ekonomi, atau lingkungan tertentu.

\subsection{Applications}\label{applications-11}

\textbf{Aplikasi analitis.} Rangkaian metode dapat digunakan untuk
mengidentifikasi kandidat subpusat dari kepadatan pekerjaan,
membandingkan keseimbangan ukuran pusat, dan memeriksa keterhubungan
fungsional melalui arus komuter. Aplikasi ini memerlukan data pekerjaan
teragregasi, lokasi centroid serta CBD, batas administratif, dan matriks
asal-tujuan komuter (Methods and data, hlm. 8-10). Pernyataan tentang
penggunaan sebagai rangkaian diagnosis merupakan inferensi review dari
metode yang diterapkan.

\textbf{Aplikasi pada perencanaan metropolitan.} Peta subpusat dapat
menunjukkan apakah pusat aktual berada di inti, menjadi limpahan CBD,
atau membentuk gugus pinggiran. Dalam JMA, metode mengenali Cikarang,
Bogor, dan koridor Tangerang-Cikupa-Balaraja sebagai bagian struktur
multipusat; informasi semacam ini berpotensi digunakan untuk
membandingkan struktur aktual dengan rencana tata ruang dan jaringan
transportasi (Subcenters in urban development, hlm. 11-13; Figure 2).
Potensi penggunaan untuk evaluasi rencana adalah inferensi review;
artikel tidak melaporkan implementasi alat keputusan atau evaluasi
kebijakan.

\textbf{Aplikasi komparatif.} Strategi individualisasi memungkinkan
perbandingan metropolitan tanpa menghapus sejarah dan basis pembangunan
masing-masing. Akan tetapi, aplikasi ke metropolitan lain memerlukan
penyesuaian definisi CBD, unit wilayah, cakupan usaha, dan definisi
komuter. Validasi eksternal di luar tiga kasus: Tidak disebutkan dalam
file.

\subsection{Future Research
(Optional)}\label{future-research-optional-11}

\begin{itemize}
\tightlist
\item
  \textbf{Dinyatakan penulis.} Studi mendatang perlu menggabungkan
  dimensi morfologis dan fungsional dalam analisis polisentralitas
  (Conclusion, hlm. 20).
\item
  \textbf{Dinyatakan penulis.} Metropolitan Indonesia perlu dibandingkan
  dengan metropolitan berukuran serupa di negara berkembang lain agar
  lebih banyak variasi pembangunan dapat dikendalikan dan kesimpulan
  menjadi lebih komprehensif (Conclusion, hlm. 20).
\item
  \textbf{Dinyatakan penulis.} Hipotesis bahwa perbedaan basis ekonomi
  regional menentukan bentuk kota memerlukan bukti tambahan (Conclusion,
  hlm. 20).
\item
  \textbf{Dinyatakan penulis.} Penelitian perlu memberi perhatian lebih
  besar pada gabungan metode empiris kuantitatif dan kualitatif untuk
  menjelaskan mekanisme tingkat mikro dan makro dalam pembentukan
  wilayah kota polisentrik (Conclusion, hlm. 20).
\item
  \textbf{Inferensi review.} Cakupan seluruh jenis usaha, jarak berbasis
  jaringan, data pekerjaan dan komuter pada tahun yang selaras, serta
  pengujian sensitivitas terhadap delineasi administratif dapat
  menanggapi batas data J22. Usulan ini tidak dinyatakan sebagai agenda
  eksplisit oleh penulis dan tidak diperlakukan sebagai temuan sumber.
\end{itemize}

\subsection{Provenance Notes}\label{provenance-notes-11}

\begin{itemize}
\tightlist
\item
  Satu-satunya bahan bukti untuk review ini adalah
  \texttt{/Users/mac/Documents/Mac/{[}2{]}\ Obsidian\ Vault/zettelkasten/source/journal/J22-sadewo-et-al-2021-using-morphological-and-functional-polycentricity-analyses-to-study-the-indonesian-urban-spatial-structure-10.1080-10225706.2020.1737829.txt}.
\item
  Seluruh 1.235 baris TXT dibaca, termasuk halaman depan, abstrak, isi
  artikel, tabel yang terekstrak sebagai teks, keterangan gambar,
  kesimpulan, pernyataan konflik kepentingan, dan daftar pustaka. TXT
  tampak memuat artikel dari halaman pembuka sampai halaman cetak 25,
  tetapi kelengkapannya terhadap berkas penerbit asli tidak diverifikasi
  melalui PDF atau situs web.
\item
  Dasar akses adalah teks penuh yang tersedia dalam bentuk ekstraksi
  TXT. Tidak ada PDF, halaman penerbit, Zotero, DOI, metadata eksternal,
  atau sumber rujukan dalam daftar pustaka yang dibuka untuk menambah
  atau mengoreksi isi review.
\item
  Locator mengacu pada nomor halaman cetak dan nama bagian yang terlihat
  di TXT. Istilah ``halaman pembuka artikel'' digunakan ketika nomor
  halaman pertama tidak terlihat jelas dalam ekstraksi. Nomor tabel dan
  gambar dipertahankan ketika tersedia.
\item
  Pemenggalan baris, ligatur, simbol matematika, spasi ribuan, dan
  susunan tabel dapat berubah selama ekstraksi. Angka \texttt{R2},
  \texttt{RMSE}, kemiringan morfologis, proporsi komuter, \texttt{Psf},
  dan \texttt{Pgf} ditranskripsikan sebagaimana tampil dalam TXT dan
  tidak dikoreksi melalui sumber eksternal.
\item
  Nama file memuat tahun 2021, sedangkan blok sitasi dalam TXT menyebut
  2020 dan halaman depan menyatakan artikel dipublikasikan secara daring
  pada 12 Maret 2020. Review mempertahankan perbedaan ini dan tidak
  menetapkan tahun edisi final tanpa bukti tambahan. Volume, nomor, dan
  rentang halaman bibliografis final: Tidak disebutkan dalam file.
\item
  Terdapat perbedaan internal pada deskripsi luas JMA: uraian menyebut
  hampir 6.392 km2, sedangkan Table 2 mencantumkan 6.935,19 km2. Uraian
  juga menyebut populasi JMA 30,97 juta pada 2015, sedangkan catatan
  Table 2 mengaitkan angka JMA dengan 2014. Perbedaan ini tidak
  diselesaikan dengan sumber luar dan tidak digunakan untuk mengubah
  hasil utama.
\item
  Istilah MMA, JMA, DMA, MLE, CBD, \texttt{R2}, \texttt{RMSE},
  \texttt{Psf}, dan \texttt{Pgf} dipertahankan agar tetap dekat dengan
  terminologi sumber. Uraian dalam bahasa Indonesia merupakan parafrasa,
  bukan kutipan langsung.
\item
  Setiap klaim mekanisme dibedakan sebagai interpretasi penulis atau
  inferensi review. Inferensi tidak diperlakukan sebagai hasil empiris
  artikel.
\end{itemize}

\section{Review J23}\label{review-j23}

\subsection{Identitas Sumber}\label{identitas-sumber-6}

\begin{itemize}
\tightlist
\item
  Kode: J23.
\item
  Judul: \emph{The role of urban transformation on inter-suburban
  commuting: evidence from Jakarta metropolitan area, Indonesia}.
\item
  Penulis: Erie Sadewo, Delik Hudalah, Anzhelika Antipova, Long Cheng,
  dan Ibnu Syabri.
\item
  Afiliasi: Badan Pusat Statistik, Jakarta; School of Architecture,
  Planning, and Policy Development, Institut Teknologi Bandung;
  Department of Earth Sciences, The University of Memphis; dan
  Department of Geography, University of Ghent (halaman pembuka
  artikel).
\item
  Penulis korespondensi: Erie Sadewo, \texttt{erie@bps.go.id} (halaman
  pembuka artikel).
\item
  Jenis sumber menurut isi: artikel jurnal. Status telaah sejawat: Tidak
  disebutkan dalam file. Bagian \emph{Acknowledgements} menyebut
  komentar dari penelaah anonim dan editor, tetapi review ini tidak
  memperlakukannya sebagai verifikasi independen atas status telaah
  sejawat (hlm. 21).
\item
  Jurnal: \emph{Urban Geography}.
\item
  Tahun: 2022.
\item
  DOI: 10.1080/02723638.2022.2125649.
\item
  URL DOI yang tercantum:
  \url{https://doi.org/10.1080/02723638.2022.2125649}.
\item
  Riwayat artikel: diterima 16 Desember 2020, disetujui 8 September
  2022, dan diterbitkan secara daring 6 Oktober 2022 (bagian
  \emph{Article History} pada halaman pembuka artikel dan lembar
  informasi penerbit).
\item
  Volume, nomor terbitan, dan rentang halaman bibliografis: Tidak
  disebutkan dalam file.
\item
  Kata kunci: \emph{post-suburbanization}; \emph{commuting};
  \emph{spatial panel regression}; \emph{Jakarta metropolitan area}
  (halaman pembuka artikel).
\item
  Geografi penelitian: Jakarta Metropolitan Area (JMA), Indonesia,
  khususnya delapan kota atau kabupaten pinggiran yang dibagi menjadi
  pinggiran dalam dan pinggiran luar (bagian \emph{Data and methods},
  hlm. 8-9; Figure 2).
\item
  Periode analisis: 2008-2015 (bagian \emph{Abstract}, halaman pembuka
  artikel; bagian \emph{Data and methods}, hlm. 9).
\item
  Pendanaan dan dukungan yang disebutkan: Bappenas PHRD IV Program,
  Agreement Note No.~35/Ses/NP06/2016, pada bagian
  \emph{Acknowledgements}; ITB Research Program 2020 kategori B pada
  bagian \emph{Funding} (hlm. 21).
\item
  Konflik kepentingan: penulis melaporkan tidak ada potensi konflik
  kepentingan (bagian \emph{Disclosure statement}, hlm. 21).
\end{itemize}

Penanda epistemik dalam review ini adalah sebagai berikut.
\textbf{Laporan sumber} berarti pernyataan, data, metode, atau hasil
yang disampaikan artikel. \textbf{Interpretasi penulis} berarti
penjelasan penulis tentang arti, mekanisme, atau implikasi hasil.
\textbf{Inferensi review} berarti pembacaan analitis yang diturunkan
hanya dari TXT J23 dan tidak diperlakukan sebagai temuan artikel atau
bukti tambahan.

\subsection{Summarized Abstract}\label{summarized-abstract-12}

\textbf{Laporan sumber.} Artikel meneliti pengaruh desentralisasi
pekerjaan dan perbedaan karakteristik pinggiran terhadap perubahan arus
komuter antarpinggiran di Jakarta Metropolitan Area. Data komuter
tahunan dan indikator sosial ekonomi tingkat kota atau kabupaten untuk
2008-2015 dianalisis dengan regresi panel spasial. Penelitian membedakan
arus masuk dari pinggiran lain, arus keluar menuju pinggiran lain, dan
arus tradisional dari pinggiran menuju Jakarta sebagai pembanding
(bagian \emph{Abstract}, halaman pembuka artikel; bagian \emph{Data and
methods}, hlm. 8-12; Table 1).

\textbf{Laporan sumber.} Jumlah komuter antarpinggiran meningkat dari
598.000 menjadi 818.000 antara 2008 dan 2015, dengan pertumbuhan
rata-rata 4,65\%, lebih tinggi daripada pertumbuhan komuter tradisional
sebesar 3,7\%. Perluasan pekerjaan di pinggiran akibat desentralisasi
industri berkaitan dengan kenaikan arus antarpinggiran. Pusat pekerjaan
industri yang monofungsional menarik pekerja dari pinggiran lain, tetapi
pada saat yang sama mendorong sebagian penduduk setempat bepergian
keluar untuk mencari pekerjaan jasa (bagian \emph{Abstract}, halaman
pembuka artikel; bagian \emph{Growing inter-suburban commuting in JMA},
hlm. 13-15; bagian \emph{Discussion}, hlm. 19-20).

\textbf{Interpretasi penulis.} Post-suburbanisasi yang dipimpin
manufaktur tidak otomatis memperbaiki keseimbangan pekerjaan dan
perumahan. Manfaatnya berbeda antarlokasi dan antarkelompok tenaga kerja
karena pekerjaan manufaktur, pekerjaan jasa, upah, biaya perumahan,
kepadatan, dan pengangguran tersebar tidak merata. Penulis menilai kasus
JMA memperluas pemahaman hubungan bentuk kota dan perilaku perjalanan
dari sudut pandang Global South (bagian \emph{Abstract}, halaman pembuka
artikel; bagian \emph{Conclusion}, hlm. 20-21).

\textbf{Inferensi review.} Model menunjukkan hubungan spasial-temporal
dan efek limpahan, tetapi desain yang dijelaskan tidak mengidentifikasi
perubahan bentuk kota sebagai sebab eksogen. Karena itu, istilah
``pengaruh'' dan ``dampak'' dalam artikel lebih aman dibaca sebagai
asosiasi yang diestimasi oleh model panel spasial, bukan efek kausal
yang telah diisolasi.

\subsection{Problem Statement}\label{problem-statement-12}

\textbf{Laporan sumber.} Desentralisasi pekerjaan mengubah pinggiran
yang semula terutama menjadi kawasan hunian menjadi wilayah multifungsi
dengan industri, jasa, dan amenitas. Perubahan ini menantang pola
komuter tradisional dari pinggiran menuju pusat bisnis utama. Literatur
menawarkan dua dugaan yang bertentangan: kedekatan pekerjaan dan hunian
dapat memendekkan perjalanan, sedangkan ketidakseimbangan
pekerjaan-perumahan serta bertambahnya pilihan tujuan dapat
memperpanjang dan memperumit perjalanan antarpinggiran (bagian
\emph{Introduction}, halaman pembuka artikel sampai hlm. 2).

\textbf{Laporan sumber.} Kajian komuter pada era post-suburban lebih
banyak berakar pada pengalaman negara maju dan pada desentralisasi jasa
atau ritel. Komuter antarpinggiran di negara berkembang menerima
perhatian lebih sedikit. Penulis menyatakan bahwa teori dari Global
North mempunyai keterbatasan ketika diterapkan pada Global South, yang
ditandai antara lain oleh deindustrialisasi prematur, pekerjaan industri
skala kecil atau temporer, kemiskinan, informalitas, segregasi spasial,
dan terbatasnya pilihan tempat tinggal pekerja (bagian
\emph{Introduction}, hlm. 2-4).

\textbf{Laporan sumber.} Konteks JMA juga berbeda dari model pinggiran
Barat. Wilayah \emph{desakota} mencampurkan permukiman berkepadatan
tinggi, pertanian, dan manufaktur. Sebagian pekerjaan formal pinggiran
diisi pekerja dari inti, sedangkan rumah tangga pinggiran banyak bekerja
secara lokal atau pada pekerjaan informal. Kombinasi komuter terbalik
dan permukiman campuran tersebut dipandang sebagai tahap awal
post-suburbanisasi Indonesia (bagian \emph{Introduction}, hlm. 3).

\textbf{Laporan sumber.} Penelitian sebelumnya kurang memperhatikan
dimensi spasial dan temporal secara bersamaan. Perbedaan lintasan
pembangunan menciptakan ketimpangan pertumbuhan pekerjaan, penduduk, dan
daya tarik antarpinggiran. Mengabaikan variabel yang khusus lokasi dan
korelasi spasial dapat menimbulkan bias estimasi (bagian
\emph{Introduction}, hlm. 3-4).

\textbf{Rumusan masalah sumber.} Artikel mempertanyakan bagaimana
perubahan lingkungan pinggiran dan dekonsentrasi pekerjaan memengaruhi
dinamika arus masuk dan keluar antarpinggiran, serta apakah penggeraknya
berbeda dari komuter tradisional menuju Jakarta (bagian
\emph{Introduction}, halaman pembuka artikel dan hlm. 3-4; bagian
\emph{Data and methods}, hlm. 8-9).

\subsection{Objectives}\label{objectives-12}

\begin{itemize}
\tightlist
\item
  \textbf{Laporan sumber:} menganalisis dampak transformasi metropolitan
  terhadap dinamika komuter di wilayah pinggiran JMA (bagian
  \emph{Introduction}, hlm. 3-4).
\item
  \textbf{Laporan sumber:} memeriksa penggerak spasial-temporal arus
  komuter tradisional dan antarpinggiran selama 2008-2015 (bagian
  \emph{Introduction}, hlm. 3-4; bagian \emph{Conclusion}, hlm. 20).
\item
  \textbf{Laporan sumber:} menguji pengaruh ketimpangan karakteristik
  pinggiran, yaitu pertumbuhan dan kepadatan penduduk, struktur rumah
  tangga, kesempatan kerja, pengangguran, upah, biaya perumahan, dan
  biaya transportasi, terhadap arus masuk serta arus keluar komuter
  (Figure 1, hlm. 6-7; Table 1, hlm. 10; bagian \emph{Conclusion}, hlm.
  20).
\item
  \textbf{Laporan sumber:} memperhitungkan ketergantungan spasial
  antarkota atau kabupaten dan perubahan antarwaktu dengan model panel
  spasial (bagian \emph{Data and methods}, hlm. 10-12).
\item
  \textbf{Laporan sumber:} menjelaskan hubungan post-suburbanisasi dan
  perubahan pola komuter dari perspektif Global South serta memberi
  bukti tentang implikasi pusat pekerjaan baru di pinggiran bagi
  pembangunan regional berkelanjutan (bagian \emph{Introduction}, hlm.
  3-4).
\end{itemize}

\subsection{Literature Survey}\label{literature-survey-12}

\textbf{Sifat survei.} Artikel menggunakan tinjauan naratif untuk
membangun kerangka konseptual dan memilih variabel serta metode.
Strategi pencarian, pangkalan data, kriteria inklusi, jumlah karya yang
disaring, dan penilaian mutu literatur: Tidak disebutkan dalam file.
Bagian ini bukan tinjauan sistematis berdasarkan informasi yang
tersedia.

\begin{itemize}
\tightlist
\item
  \textbf{Laporan sumber atas konsep post-suburbanisasi:} desentralisasi
  kegiatan ekonomi mengubah pinggiran dari kota asrama menjadi wilayah
  dengan pusat pekerjaan, guna lahan campuran, kepadatan yang meningkat,
  dan batas urban-perdesaan yang kabur. Bentuk di Amerika Serikat
  dikaitkan dengan \emph{edge city} berbasis ruang kantor, bentuk Eropa
  dengan pusat pinggiran yang tetap bergantung pada pusat tradisional,
  dan bentuk Tiongkok dengan kawasan industri yang direncanakan
  pemerintah (bagian \emph{Introduction}, halaman pembuka artikel sampai
  hlm. 2; bagian \emph{The urban form and commuting nexus}, hlm. 4).
\item
  \textbf{Laporan sumber atas keseimbangan pekerjaan-perumahan:}
  pandangan efisiensi polisentris menyatakan bahwa pusat pekerjaan
  tambahan mendekatkan pekerja pada pekerjaan dan mengurangi waktu serta
  biaya perjalanan. Literatur tandingan menunjukkan bahwa pekerjaan di
  subpusat sering diisi nonpenduduk dan bahwa segregasi ras, perbedaan
  pendapatan, serta keterbatasan relokasi dapat meningkatkan
  ketidakseimbangan dan perjalanan berlebih (bagian \emph{The urban form
  and commuting nexus}, hlm. 4-5).
\item
  \textbf{Laporan sumber atas aksesibilitas:} lokasi, kedekatan, ukuran
  pusat pekerjaan, jarak ke CBD, distribusi sumber daya, harga lahan,
  upah, dan biaya perjalanan dipakai dalam literatur untuk menjelaskan
  panjang perjalanan dan keputusan lokasi rumah. Artikel juga merangkum
  argumen bahwa penduduk pinggiran jauh dapat terus menerima biaya
  komuter tinggi ketika relokasi ditunda atau hubungan upah dan biaya
  perjalanan mengompensasinya (bagian \emph{The urban form and commuting
  nexus}, hlm. 5).
\item
  \textbf{Laporan sumber atas komuter lintas kota atau antarpinggiran:}
  perbedaan upah, biaya rumah, jarak, migrasi, rasio dan keragaman
  pekerjaan, ukuran kota, kepadatan, serta struktur pasar kerja
  dilaporkan sebagai penentu arus. Biaya pindah dan ketidaksempurnaan
  pasar perumahan memungkinkan pekerja yang terhubung dengan pusat
  pekerjaan berbeda tetap tinggal berdampingan dan melakukan komuter
  silang (bagian \emph{Commuting between suburbs}, hlm. 5-6).
\item
  \textbf{Laporan sumber atas heterogenitas amenitas:} perbedaan
  kualitas hidup dan amenitas antarkota dapat mendorong komuter.
  Struktur rumah tangga dan tingkat keterampilan anggota rumah tangga
  juga disebut dalam literatur sebagai faktor yang dapat mengubah jarak
  atau pola perjalanan (bagian \emph{Commuting between suburbs}, hlm.
  6).
\item
  \textbf{Laporan sumber atas perbedaan Global South:} pekerjaan baru di
  wilayah yang mengalami deindustrialisasi prematur cenderung ditempati
  pekerja berketerampilan lebih rendah dan dalam pengaturan sementara.
  Di Indonesia, segregasi penduduk perkotaan lebih dikaitkan dengan
  kemiskinan dan informalitas daripada ras. Banyak kawasan pinggiran JMA
  dihuni penduduk berpendapatan rendah, berbeda dari gambaran pinggiran
  kelas menengah Amerika Serikat (bagian \emph{Introduction}, hlm. 2-3).
\item
  \textbf{Laporan sumber atas JMA:} perluasan industri manufaktur,
  perpindahan penduduk ke pinggiran, pertumbuhan beberapa subpusat di
  luar Jakarta, kenaikan jumlah dan panjang perjalanan, serta dominasi
  sepeda motor membentuk konteks empiris penelitian. Artikel melaporkan
  dari karya terdahulu bahwa hampir 750.000 orang bepergian
  antarpinggiran JMA pada 2014, hampir tiga kali komuter terbalik dari
  Jakarta ke wilayah sekitarnya (bagian \emph{Peripheral urban
  development of Jakarta metropolitan area}, hlm. 7-8).
\item
  \textbf{Laporan sumber atas perubahan pekerjaan dan penduduk JMA:}
  artikel mengutip pertumbuhan rata-rata tahunan 93\% untuk jumlah
  perusahaan menengah dan besar di pinggiran selama 1976-2016,
  sebagaimana tercetak dalam TXT. Hampir 40\% industri manufaktur JMA
  dan tujuh kawasan industri besar disebut berada di Kabupaten Bekasi.
  Pertumbuhan penduduk Jakarta turun dari 2,5\% sebelum 1990-an menjadi
  0,11\% satu dekade kemudian, sementara proporsi penduduk JMA yang
  tinggal di Jakarta turun dari 54,6\% menjadi 35,5\% dalam tiga dekade.
  Angka-angka ini berasal dari karya yang dirujuk artikel, bukan hasil
  model J23 (bagian \emph{Peripheral urban development of Jakarta
  metropolitan area}, hlm. 7).
\item
  \textbf{Laporan sumber atas mobilitas JMA sebelum dan selama periode
  studi:} artikel mengutip sekitar 700.000 komuter harian pada
  pertengahan 2000-an, sepuluh kali jumlah dua dekade sebelumnya.
  Rata-rata perjalanan pada 2005 dilaporkan 21,94 km dan 54,78 menit;
  Kabupaten Tangerang mempunyai rata-rata 37,03 km dan Kabupaten Bekasi
  15,33 km. Rata-rata waktu komuter JMA naik menjadi 61,87 menit pada
  2015. Artikel juga mengutip bahwa sepeda motor mencakup hingga 56\%
  perjalanan komuter dan dapat menghemat sekitar 30\% biaya dibandingkan
  angkutan umum. Semua angka pada butir ini adalah laporan J23 atas
  studi terdahulu (bagian \emph{Peripheral urban development of Jakarta
  metropolitan area}, hlm. 7-8).
\item
  \textbf{Kerangka konseptual penulis:} desentralisasi pekerjaan membuat
  sebagian pinggiran menawarkan lebih banyak pekerjaan, upah, dan
  amenitas, sedangkan pinggiran lain mempertahankan atau kehilangan
  fungsi hunian. Komuter antarpinggiran diposisikan sebagai respons
  pekerja terhadap ketimpangan tersebut. Pertumbuhan penduduk yang tidak
  berimpit dengan kesempatan kerja dapat menaikkan pengangguran; harga
  rumah dekat pusat yang lebih menarik dapat memaksa pekerja menukar
  biaya rumah yang lebih rendah dengan biaya perjalanan yang lebih
  tinggi (bagian \emph{Synthesizing inter-suburban commuting: a
  conceptual framework}, Figure 1, hlm. 6-7).
\end{itemize}

\textbf{Inferensi review.} Literatur menghubungkan bentuk kota, pasar
tenaga kerja, pasar perumahan, amenitas, dan mobilitas dalam satu
kerangka. Namun, mekanisme tingkat individu kemudian diuji dengan
indikator agregat kota atau kabupaten. Kesepadanan antara konsep
individual, seperti aspirasi, keterampilan, dan maksimalisasi utilitas,
dengan proksi agregat tidak divalidasi secara khusus dalam file.

\subsection{Method Used}\label{method-used-12}

\textbf{Laporan sumber.} Penelitian menggunakan desain panel spasial
pada tingkat kota atau kabupaten untuk 2008-2015. Perspektif analisisnya
adalah wilayah pinggiran sebagai asal dan tujuan arus. Klasifikasi ruang
mengikuti batas administrasi: lima kota administratif di Jakarta
diperlakukan sebagai inti, sedangkan delapan kota atau kabupaten di luar
Jakarta diperlakukan sebagai pinggiran dan menjadi unit geografis
analisis (bagian \emph{Data and methods}, hlm. 8-9; Figure 2).

Delapan wilayah pinggiran dibagi menjadi:

\begin{itemize}
\tightlist
\item
  pinggiran dalam yang berbatasan dengan Jakarta: Kota Bekasi, Kota
  Depok, Kota Tangerang Selatan, dan Kota Tangerang;
\item
  pinggiran luar: Kabupaten Bekasi, Kabupaten Bogor, Kota Bogor, dan
  Kabupaten Tangerang (bagian \emph{Data and methods}, hlm. 8-9; Figure
  2).
\end{itemize}

\textbf{Operasionalisasi arus.} Arus masuk adalah jumlah orang yang
masuk dari pinggiran lain. Arus keluar antarpinggiran adalah jumlah
orang yang keluar menuju pinggiran lain. Arus komuter tradisional adalah
jumlah orang yang keluar menuju Jakarta. Ketiganya ditransformasi dengan
logaritma natural (Table 1, hlm. 10).

\textbf{Operasionalisasi penggerak.} Model memasukkan pertumbuhan
penduduk tahunan; log kepadatan penduduk; rata-rata jumlah anggota rumah
tangga; log upah minimum kota atau kabupaten; tingkat pengangguran; log
kesempatan kerja baru di manufaktur; log biaya transportasi rata-rata
per kilometer bagi komuter; dan log biaya perumahan rata-rata per meter
persegi (Table 1, hlm. 10).

\textbf{Spesifikasi model.} Artikel memulai uraian dari model panel
nonspasial dengan efek individual yang tidak teramati, lalu menyajikan
model bersarang umum yang dapat memuat lag spasial variabel dependen,
lag spasial variabel penjelas, dan korelasi spasial galat. Parameter
diperkirakan dengan kemungkinan maksimum memakai paket \texttt{spml}
dalam perangkat lunak R (Equations 1-5, hlm. 11).

\textbf{Identifikasi efek panel dan spasial.} Penulis memakai uji
Hausman untuk memilih efek tetap atau efek acak. Uji Lagrange Multiplier
dengan metode gabungan BSK dan perluasannya digunakan untuk memeriksa
kebutuhan model panel spasial, sedangkan metode BSJK menguji efek
wilayah acak ketika galat spasial diizinkan. Komponen galat diperiksa
sebelum estimasi melalui pemusatan dan pemaksimalan iteratif fungsi
kemungkinan (bagian \emph{Data and methods}, hlm. 11-12; Figure 3).

\textbf{Bobot spasial.} Tiga matriks dibandingkan:

\begin{itemize}
\tightlist
\item
  W1, yang disebut bobot kota terdekat, bernilai 1 ketika dua wilayah
  berbagi batas dan 0 jika tidak;
\item
  W2, bobot invers jarak berdasarkan jarak Euclidean antara puncak
  pekerjaan lokal tertinggi pada dua kota atau kabupaten;
\item
  W3, bobot khusus berdasarkan proporsi volume komuter antarpinggiran
  dari suatu asal menuju tujuan selama periode pengamatan.
\end{itemize}

Semua bobot dinormalisasi per baris. Model dipilih dengan membandingkan
Akaike Information Criterion (AIC) dan Bayesian Information Criterion
(BIC). Model komuter tradisional diperkirakan sebagai pembanding bagi
variabel penjelas yang digunakan (Equations 6-8, hlm. 12; Table 2, hlm.
16).

\textbf{Model akhir yang dilaporkan.} Untuk arus keluar antarpinggiran,
model terpilih adalah Spatial Durbin Error Model dengan efek acak
(SDEM-RE). Untuk arus masuk antarpinggiran, model terpilih adalah
General Nesting Spatial Model dengan efek acak (GNSM-RE). Kedua model
antarpinggiran menggunakan W1 menurut uraian dan nilai AIC/BIC dalam
Table 2. Bagian \emph{Discussion} menyatakan bahwa model komuter
tradisional menggunakan efek acak dengan W3 (bagian \emph{The suburban
commuting models}, hlm. 15-17; Table 2; bagian \emph{Discussion}, hlm.
18-19).

\textbf{Inferensi review atas desain.} Unit panel tampaknya merupakan
delapan wilayah pinggiran yang diamati selama delapan tahun, tetapi
jumlah observasi efektif, panel yang seimbang atau tidak seimbang, dan
observasi yang dikeluarkan tidak dinyatakan. Review ini tidak menghitung
ukuran sampel panel dari struktur yang tersirat tersebut.

\subsection{Dataset}\label{dataset-12}

\begin{itemize}
\tightlist
\item
  \textbf{Data komuter:} Survei Angkatan Kerja Nasional (SAKERNAS)
  tahunan dari Badan Pusat Statistik (BPS) menyediakan jumlah komuter
  pada tingkat administrasi kota atau kabupaten. Data tersedia mulai
  2008. Data tingkat kabupaten tidak tersedia dalam SAKERNAS 2016
  sehingga periode analisis dibatasi pada 2008-2015 (bagian \emph{Data
  and methods}, hlm. 9; Table 1, hlm. 10).
\item
  \textbf{Definisi dalam data komuter:} mengikuti konsep SAKERNAS,
  komuter adalah orang yang setiap hari melintasi batas administrasi
  tempat tinggal dan kembali dalam waktu 24 jam (bagian \emph{Data and
  methods}, hlm. 9).
\item
  \textbf{Pertumbuhan dan kepadatan penduduk:} pertumbuhan penduduk
  tahunan dan jumlah penduduk per kilometer persegi berasal dari
  publikasi statistik kota atau kabupaten pada tahun yang bersesuaian
  (Table 1, hlm. 10).
\item
  \textbf{Struktur rumah tangga:} rata-rata jumlah anggota rumah tangga
  berasal dari publikasi statistik kota atau kabupaten pada tahun yang
  bersesuaian (Table 1, hlm. 10).
\item
  \textbf{Upah:} upah minimum kota atau kabupaten dalam rupiah berasal
  dari publikasi statistik kota atau kabupaten pada tahun yang
  bersesuaian (Table 1, hlm. 10).
\item
  \textbf{Pengangguran:} tingkat pengangguran, yang didefinisikan
  sebagai proporsi angkatan kerja aktif yang menganggur, berasal dari
  publikasi statistik kota atau kabupaten pada tahun yang bersesuaian
  (Table 1, hlm. 10).
\item
  \textbf{Kesempatan kerja:} jumlah pekerjaan baru di sektor manufaktur
  berasal dari Sensus Ekonomi 2016 dan digunakan sebagai proksi
  kesempatan kerja (Table 1, hlm. 10). Cara Sensus Ekonomi 2016
  dipetakan ke setiap tahun panel 2008-2015: Tidak disebutkan dalam
  file.
\item
  \textbf{Biaya transportasi:} rata-rata biaya transportasi per
  kilometer yang dikeluarkan komuter, dalam rupiah, berasal dari
  SAKERNAS dan Survei Sosial Ekonomi Nasional (SUSENAS) (Table 1, hlm.
  10).
\item
  \textbf{Biaya perumahan:} rata-rata biaya perumahan per meter persegi,
  dalam rupiah, berasal dari SUSENAS (Table 1, hlm. 10).
\item
  \textbf{Informasi spasial:} W1 memakai ketetanggaan batas wilayah; W2
  memakai jarak Euclidean antarpuncak pekerjaan lokal tertinggi; W3
  dibentuk dari volume relatif komuter antarpinggiran. Sumber geometri
  batas, prosedur penentuan puncak pekerjaan, dan data koordinat: Tidak
  disebutkan dalam file (Equations 6-8, hlm. 12).
\end{itemize}

Nama berkas mikrodata, tautan akses, lisensi, desain dan bobot sampel
SAKERNAS atau SUSENAS, prosedur penggabungan antarsumber, penanganan
nilai hilang, penyesuaian nilai rupiah antarwaktu, serta kode analisis:
Tidak disebutkan dalam file.

\subsection{Population / Sample}\label{population-sample-12}

\textbf{Cakupan wilayah.} Artikel menyatakan bahwa JMA mempunyai luas
6.256 km2 dan terdiri atas delapan kota atau kabupaten serta satu
provinsi. DKI Jakarta, yang mencakup lima kota administratif,
diperlakukan sebagai inti. Delapan kota atau kabupaten di sekelilingnya
menjadi unit analisis pinggiran (bagian \emph{Data and methods}, hlm.
8-9; Figure 2).

\textbf{Populasi substantif.} Populasi yang diukur adalah komuter harian
yang melintasi batas administrasi tempat tinggal dan kembali dalam 24
jam. Analisis arus antarpinggiran mencakup komuter yang bergerak di
antara delapan wilayah pinggiran, sedangkan model pembanding mencakup
komuter dari pinggiran menuju Jakarta (bagian \emph{Data and methods},
hlm. 9; Table 1, hlm. 10).

\textbf{Unit observasi analitis.} Unit geografis adalah kota atau
kabupaten pinggiran pada suatu tahun. Jumlah responden survei, jumlah
catatan individu, jumlah pasangan asal-tujuan, jumlah observasi panel
efektif, proporsi data hilang, dan distribusi sampel per wilayah atau
tahun: Tidak disebutkan dalam file.

\textbf{Inferensi review.} Definisi komuter mengecualikan perjalanan
kerja yang tidak melintasi batas kota atau kabupaten, perjalanan yang
tidak dilakukan setiap hari, dan perjalanan yang tidak kembali dalam 24
jam. Karena hasil disajikan sebagai total arus tingkat wilayah, artikel
menilai dinamika agregat, bukan perubahan perjalanan individu yang sama
dari tahun ke tahun.

\subsection{Dependent Variables}\label{dependent-variables-12}

Artikel memperkirakan tiga variabel dependen dalam model terpisah:

\begin{itemize}
\tightlist
\item
  logaritma natural jumlah komuter masuk dari pinggiran lain, yaitu
  \texttt{Ln(in-commuting)};
\item
  logaritma natural jumlah komuter keluar menuju pinggiran lain, yaitu
  \texttt{Ln(out-commuting)};
\item
  logaritma natural jumlah komuter keluar menuju Jakarta, juga diberi
  label \texttt{Ln(out-commuting)} dalam Table 1, tetapi didefinisikan
  sebagai komuter tradisional (Table 1, hlm. 10).
\end{itemize}

Setiap variabel adalah volume orang pada unit kota atau kabupaten dan
tahun, bukan waktu tempuh, panjang perjalanan, pangsa moda, atau
probabilitas seseorang menjadi komuter. Dalam model spasial umum, lag
spasial variabel dependen (\texttt{WY}) memungkinkan volume di satu
wilayah berhubungan dengan volume wilayah yang bertetangga. Efek endogen
tersebut hanya dilaporkan bermakna secara statistik pada model arus
masuk antarpinggiran (Equations 2-5, hlm. 11; Table 3, hlm. 17-18).

Penanganan arus bernilai nol sebelum transformasi logaritma, pembobotan
total hasil survei, dan alasan memilih volume absolut alih-alih proporsi
terhadap populasi atau angkatan kerja: Tidak disebutkan dalam file.

\subsection{Results}\label{results-12}

\textbf{Perubahan volume dan moda komuter.}

\begin{itemize}
\tightlist
\item
  \textbf{Laporan sumber:} jumlah komuter antarpinggiran meningkat dari
  sekitar 598.000 pada 2008 menjadi 818.000 pada 2015. Pertumbuhan
  rata-ratanya 4,65\%, sedangkan komuter tradisional menuju CBD tumbuh
  3,7\% pada periode yang sama (bagian \emph{Growing inter-suburban
  commuting in JMA}, Figure 4, hlm. 13).
\item
  \textbf{Laporan sumber:} perjalanan menuju CBD tetap mendominasi pola
  JMA, terutama dari wilayah timur. Kota Bekasi adalah asal komuter
  tradisional terbesar di pinggiran dalam, diikuti Kabupaten Bekasi.
  Untuk arus antarpinggiran menurut asal, Kabupaten Bogor mempunyai
  porsi terbesar, disusul Kabupaten Tangerang (bagian \emph{Growing
  inter-suburban commuting in JMA}, hlm. 13).
\item
  \textbf{Laporan sumber:} pengguna kendaraan pribadi pada komuter
  tradisional meningkat dari 58\% menjadi 75\% selama 2008-2015. Pada
  komuter antarpinggiran, proporsinya meningkat dari 48\% menjadi 75\%.
  Penggunaan angkutan umum pada kedua pola turun sekitar separuh menurut
  narasi artikel (bagian \emph{Growing inter-suburban commuting in JMA},
  hlm. 13).
\item
  \textbf{Laporan sumber:} tujuan komuter antarpinggiran lebih
  terkonsentrasi daripada asalnya. Kota Tangerang, Kota Bogor, dan
  Kabupaten Bekasi bersama-sama menerima rata-rata hampir 52\% seluruh
  komuter masuk dari pinggiran lain (Figure 5 dan uraian, hlm. 13-14).
\end{itemize}

\textbf{Restrukturisasi pekerjaan di Kabupaten Bekasi dan Kabupaten
Tangerang.}

\begin{itemize}
\tightlist
\item
  \textbf{Laporan sumber:} selama 2008-2015, perusahaan menengah dan
  besar baru di Kabupaten Bekasi menghasilkan 126.000 pekerjaan,
  terbanyak di antara wilayah pinggiran. Pangsa pekerjaan manufakturnya
  naik dari 30,01\% menjadi 38,47\%, sedangkan upah minimum naik hampir
  tiga kali lipat (bagian \emph{Growing inter-suburban commuting in
  JMA}, hlm. 14).
\item
  \textbf{Laporan sumber:} Kabupaten Tangerang menghasilkan pekerjaan
  19\% lebih sedikit daripada Kabupaten Bekasi dengan karakteristik
  restrukturisasi yang disebut serupa, tetapi komuter tradisionalnya
  turun menjadi 12.000 perjalanan (bagian \emph{Growing inter-suburban
  commuting in JMA}, hlm. 14).
\item
  \textbf{Interpretasi penulis:} pertumbuhan manufaktur Kabupaten Bekasi
  terutama menguntungkan pekerja dari luar wilayah. Pertumbuhan penduduk
  tahunan 4,4\% dan ketidakcocokan kualifikasi membuat pekerjaan baru
  menjadi tarikan bagi pekerja luar, sementara penduduk yang memilih
  pekerjaan jasa atau komersial bepergian ke wilayah lain. Kota Bekasi
  mencatat pertumbuhan arus masuk rata-rata tertinggi dan ditafsirkan
  menarik pekerja jasa dari Kabupaten Bekasi (bagian \emph{Growing
  inter-suburban commuting in JMA}, Figure 6, hlm. 14-15).
\item
  \textbf{Interpretasi penulis:} Kabupaten Tangerang, dengan pertumbuhan
  penduduk tahunan 3,5\%, mengalami pertumbuhan arus masuk negatif dan
  kenaikan arus keluar yang rendah. Penulis menyimpulkan bahwa
  restrukturisasi memberi manfaat yang lebih bersifat regional di
  Kabupaten Bekasi, tetapi lebih lokal di Kabupaten Tangerang (bagian
  \emph{Growing inter-suburban commuting in JMA}, hlm. 14-15).
\end{itemize}

\textbf{Uji model dan pemilihan spesifikasi.}

\begin{itemize}
\tightlist
\item
  \textbf{Laporan sumber:} narasi menyatakan bahwa uji komponen galat
  BSK menolak hipotesis tidak adanya galat spasial dan efek acak, lalu
  uji efek wilayah acak dan BSJK juga mendukung keberadaan efek wilayah
  acak. Namun, Table 2 memuat beberapa hasil tanpa tanda signifikansi,
  termasuk BSK gabungan dan perluasan untuk arus masuk dengan W3 serta
  beberapa uji BSJK. Perbedaan antara ringkasan naratif dan tabel ini
  tidak direkonsiliasi dalam file (bagian \emph{The suburban commuting
  models}, hlm. 15; Table 2, hlm. 16).
\item
  \textbf{Laporan sumber:} uji Hausman memberi hasil campuran. Tiga dari
  sembilan kombinasi model-bobot memilih efek tetap menurut Table 2,
  yaitu arus keluar dengan W3 serta arus masuk dengan W2 dan W3.
  Kombinasi lainnya memakai efek acak (Table 2, hlm. 16).
\item
  \textbf{Laporan sumber:} model efek acak dengan W1 mempunyai AIC dan
  BIC terendah untuk arus keluar dan arus masuk antarpinggiran. Nilainya
  masing-masing AIC 371,06 dan BIC 409,92 untuk arus keluar, serta AIC
  348,55 dan BIC 387,41 untuk arus masuk (Table 2, hlm. 16).
\item
  \textbf{Laporan sumber:} untuk komuter tradisional, Table 2
  menunjukkan AIC 230,22 dan BIC 269,08 pada model efek acak dengan W3,
  lebih rendah daripada W1 dan W2. Bagian \emph{Discussion} juga
  menyebut RE dengan W3 sebagai spesifikasi terbaik (Table 2, hlm. 16;
  bagian \emph{Discussion}, hlm. 18). Pernyataan sebelumnya bahwa W1
  terbaik untuk ``setiap pola komuter'' tidak sesuai dengan angka ini
  (bagian \emph{The suburban commuting models}, hlm. 15).
\item
  \textbf{Laporan sumber:} model arus masuk mempunyai interaksi endogen,
  eksogen, dan galat sehingga diperkirakan sebagai GNSM-RE. Model arus
  keluar mempunyai interaksi eksogen dan galat, tetapi tidak mempunyai
  interaksi endogen yang bermakna, sehingga diperkirakan sebagai SDEM-RE
  (bagian \emph{The suburban commuting models}, hlm. 15-18; Table 3).
\end{itemize}

\textbf{Estimasi arus keluar antarpinggiran.}

\begin{itemize}
\tightlist
\item
  \textbf{Laporan sumber:} koefisien langsung kepadatan adalah -1,02
  dengan tiga tanda bintang. Penulis menafsirkannya sebagai penurunan
  arus keluar sekitar 1\% ketika kepadatan naik 1\%. Lag spasial
  kepadatan bernilai -2,03 dengan tiga tanda bintang, yang ditafsirkan
  sebagai penurunan arus keluar sekitar 2\% di wilayah sekitar ketika
  kepadatan suatu pinggiran naik 1\% (Table 3, hlm. 17; uraian model,
  hlm. 17-18).
\item
  \textbf{Laporan sumber:} koefisien langsung biaya perumahan adalah
  0,73 dengan satu tanda bintang dan lag spasialnya 2,63 dengan dua
  tanda bintang. Penulis menafsirkannya sebagai kenaikan arus keluar
  sekitar 0,7\% di wilayah itu dan 2,6\% di wilayah sekitar ketika biaya
  perumahan naik 1\% (Table 3, hlm. 17; uraian model, hlm. 17-18).
\item
  \textbf{Laporan sumber:} kesempatan kerja manufaktur mempunyai
  koefisien langsung positif 0,07 dengan dua tanda bintang. Penulis
  menafsirkan hasil ini sebagai petunjuk bahwa penduduk yang bepergian
  ke pinggiran lain mencari pekerjaan di luar manufaktur (Table 3, hlm.
  17; uraian model, hlm. 17-18).
\item
  \textbf{Laporan sumber:} parameter galat spasial bernilai 0,55 dengan
  dua tanda bintang. Penulis menyatakan bahwa interaksi galat
  menunjukkan masih adanya pengaruh variabel terlewat yang berkorelasi
  secara spasial (Table 3 dan uraian, hlm. 17-18).
\item
  \textbf{Laporan sumber:} pertumbuhan penduduk, struktur rumah tangga,
  upah, pengangguran, dan biaya transportasi tidak mempunyai tanda
  signifikansi pada efek langsung arus keluar. Lag spasial
  variabel-variabel tersebut juga tidak bertanda signifikan dalam Table
  3 (hlm. 17).
\end{itemize}

\textbf{Estimasi arus masuk antarpinggiran.}

\begin{itemize}
\tightlist
\item
  \textbf{Laporan sumber:} koefisien lag spasial variabel dependen
  adalah -0,40 dengan satu tanda bintang. Penulis menafsirkannya sebagai
  hubungan pertukaran: arus masuk suatu wilayah naik ketika arus masuk
  di wilayah tetangganya turun, sesuai dengan konsentrasi tujuan pada
  beberapa lokasi (Table 3, hlm. 17; uraian model, hlm. 18).
\item
  \textbf{Laporan sumber:} upah mempunyai koefisien langsung -0,75
  dengan satu tanda bintang, sedangkan lag spasial upah bernilai 1,31
  dengan tiga tanda bintang. Penulis menghubungkan pola ini dengan pusat
  manufaktur berupah tinggi di pinggiran luar, arus pekerja
  nonmanufaktur, dan dampak terbalik pada wilayah tetangga (Table 3,
  hlm. 17; uraian model, hlm. 18).
\item
  \textbf{Laporan sumber:} lag spasial tingkat pengangguran bernilai
  -0,05 dengan tiga tanda bintang, sementara efek langsung pengangguran
  tidak bertanda signifikan. Penulis menyatakan bahwa kenaikan arus
  masuk terjadi pada wilayah yang dikelilingi kota atau kabupaten dengan
  tingkat pengangguran lebih rendah (Table 3, hlm. 17; uraian model,
  hlm. 18).
\item
  \textbf{Laporan sumber:} parameter galat spasial bernilai 0,62 dengan
  tiga tanda bintang. Pertumbuhan penduduk, kepadatan, struktur rumah
  tangga, kesempatan kerja manufaktur, biaya transportasi, dan biaya
  perumahan tidak mempunyai tanda signifikansi pada efek langsung arus
  masuk dalam Table 3 (hlm. 17).
\end{itemize}

\textbf{Estimasi komuter tradisional.}

\begin{itemize}
\tightlist
\item
  \textbf{Laporan sumber:} biaya perumahan mempunyai koefisien langsung
  1,35 dengan tiga tanda bintang dan lag spasial 1,94 dengan tiga tanda
  bintang. Lag spasial kepadatan bernilai -1,92 dengan tiga tanda
  bintang. Koefisien langsung kepadatan bernilai -0,38 tanpa tanda
  signifikansi (Table 3, hlm. 17).
\item
  \textbf{Interpretasi penulis:} wilayah yang berdekatan dengan Jakarta,
  yaitu Bekasi, Tangerang, Depok, dan Tangerang Selatan, menjadi lokasi
  hunian penting bagi komuter tradisional. Biaya rumah lebih tinggi
  karena kedekatannya dengan Jakarta, dan dinamika kepadatan serta biaya
  rumah dinilai cukup untuk menjelaskan pola tradisional (bagian
  \emph{Discussion}, hlm. 18-19).
\item
  \textbf{Laporan sumber:} parameter galat spasial (-0,07) dan lag
  spasial variabel dependen (-0,13) tidak mempunyai tanda signifikansi.
  Penulis menyatakan bahwa interaksi endogen dan interaksi galat tidak
  memengaruhi komuter tradisional (Table 3, hlm. 17; bagian
  \emph{Discussion}, hlm. 18-19).
\end{itemize}

\textbf{Arah perjalanan antarpinggiran.} Narasi hasil menyatakan bahwa
perjalanan dari pinggiran dalam menuju pinggiran luar mencapai rata-rata
42\%, sedangkan arah sebaliknya 28\% selama 2008-2015 (Figure 7 dan
uraian, hlm. 17-18). Akan tetapi, bagian \emph{Discussion} menyatakan
bahwa lebih banyak arus keluar berlangsung dari pinggiran luar menuju
pinggiran dalam (hlm. 19). Review ini mempertahankan kedua pernyataan
sebagai ketidaksesuaian internal dan tidak memilih salah satunya tanpa
data visual Figure 7.

Ambang probabilitas yang diwakili satu, dua, dan tiga tanda bintang,
galat baku, nilai p, interval kepercayaan, jumlah observasi, ukuran
kecocokan selain AIC/BIC, dan dekomposisi dampak lengkap: Tidak
disebutkan dalam file.

\subsection{Findings}\label{findings-12}

\begin{itemize}
\tightlist
\item
  \textbf{Interpretasi penulis:} heterogenitas wilayah sekitar
  berhubungan dengan komuter antarpinggiran. Kepadatan, biaya rumah,
  upah, dan pengangguran di wilayah tetangga mempunyai efek limpahan,
  sehingga pinggiran yang berdekatan dipandang saling melengkapi dalam
  pasar kerja metropolitan (bagian \emph{Discussion}, hlm. 19-20).
\item
  \textbf{Interpretasi penulis:} arus utama berlangsung secara lateral
  antara kota atau kabupaten yang berdekatan di sepanjang koridor jalan
  tol. Penulis mengaitkannya dengan kenaikan penggunaan kendaraan
  pribadi selama 2008-2015 dan konsentrasi pekerjaan di sepanjang
  koridor (bagian \emph{Discussion}, hlm. 19).
\item
  \textbf{Interpretasi penulis:} kenaikan kepadatan menurunkan arus
  keluar, sedangkan kenaikan biaya perumahan meningkatkan arus keluar,
  baik secara langsung maupun melalui wilayah sekitar. Hubungan positif
  biaya rumah tidak sesuai dengan kerangka awal yang memperkirakan
  komuter berasal dari lokasi dengan biaya rumah lebih rendah. Penulis
  menjelaskan hasil tersebut melalui besarnya pertukaran perjalanan
  antara pinggiran dalam yang padat dan mahal dengan pinggiran luar
  (uraian model, hlm. 17-18; bagian \emph{Discussion}, hlm. 19-20).
\item
  \textbf{Interpretasi penulis:} post-suburbanisasi manufaktur
  menghasilkan dampak ganda. Pusat manufaktur di Kabupaten Bekasi dan
  Kabupaten Tangerang menarik pekerja dari luar, tetapi keragaman
  pekerjaan lokal yang rendah mendorong penduduk yang menginginkan
  pekerjaan jasa atau komersial untuk bepergian ke Jakarta atau
  pinggiran lain (bagian \emph{Growing inter-suburban commuting in JMA},
  hlm. 14-15; bagian \emph{Discussion}, hlm. 19).
\item
  \textbf{Interpretasi penulis:} upah tinggi di pusat manufaktur
  pinggiran luar tidak otomatis mengurangi arus keluar dan tidak
  menguntungkan semua penduduk. Penulis menghubungkan anomali koefisien
  upah dengan preferensi komuter berpendidikan tinggi terhadap pekerjaan
  jasa dan komersial, serta dengan pemusatan kembali perusahaan jasa di
  Jakarta (bagian \emph{Discussion}, hlm. 19).
\item
  \textbf{Interpretasi penulis:} besarnya efek limpahan kepadatan dan
  biaya perumahan, yang hampir dua kali efek langsung menurut narasi,
  menunjukkan ketergantungan kuat antarlokasi pekerjaan yang
  bertetangga. Perkembangan sepanjang jalan tol dapat menyebarkan
  kesempatan kerja, tetapi komuter antarpinggiran harus berbagi koridor
  dengan komuter tradisional dan komuter terbalik (bagian
  \emph{Discussion}, hlm. 19-20).
\item
  \textbf{Interpretasi penulis:} upah minimum yang tinggi di pinggiran
  manufaktur dapat mengurangi keuntungan biaya bagi perusahaan jasa.
  Kondisi itu menghambat aglomerasi manufaktur dan jasa di lokasi yang
  sama. Perusahaan jasa lebih mungkin memilih pinggiran dalam yang
  berupah sedikit lebih rendah dan mempunyai lebih banyak amenitas
  karena dekat dengan Jakarta (bagian \emph{Discussion}, hlm. 20).
\item
  \textbf{Interpretasi penulis:} pusat industri yang monofungsional
  dapat memperburuk keseimbangan pekerjaan-perumahan. Penambahan
  pekerjaan manufaktur tidak menjamin kecocokan dengan keterampilan atau
  pilihan kerja penduduk setempat, sehingga arus masuk dan arus keluar
  dapat meningkat bersamaan (bagian \emph{Discussion}, hlm. 19-20;
  bagian \emph{Conclusion}, hlm. 20).
\item
  \textbf{Interpretasi penulis:} hubungan positif biaya perumahan dengan
  komuter tradisional dan penurunan arus seiring jarak dari CBD dinilai
  mengonfirmasi model kota monosentris. Namun, frasa sumber menyebut
  ``lag spasial jarak'', sedangkan Table 3 hanya menampilkan lag spasial
  kepadatan. Ketidaksesuaian variabel ini membatasi pemeriksaan atas
  interpretasi tersebut (bagian \emph{Discussion}, hlm. 18-19; Table 3,
  hlm. 17).
\end{itemize}

\textbf{Inferensi review.} Bukti model paling jelas untuk
restrukturisasi manufaktur bukan kenaikan arus masuk akibat variabel
kesempatan kerja. Koefisien langsung kesempatan kerja pada model arus
masuk adalah -0,02 dan tidak bertanda signifikan. Kesimpulan tentang
daya tarik pusat industri bertumpu pada pola deskriptif wilayah,
koefisien upah dan limpahannya, serta penafsiran penulis. Sebaliknya,
kesempatan kerja manufaktur mempunyai hubungan positif yang bertanda
signifikan dengan arus keluar (Table 3, hlm. 17; uraian hasil, hlm.
17-18).

\textbf{Inferensi review.} Banyak variabel kerangka konseptual tidak
bertanda signifikan pada efek langsung, termasuk pertumbuhan penduduk,
struktur rumah tangga, pengangguran, dan biaya transportasi. Hasil tidak
mendukung anggapan bahwa semua unsur heterogenitas pinggiran bekerja
secara langsung. Temuan lebih khusus: beberapa karakteristik tertentu
berhubungan dengan arus, dan sebagian hubungan muncul melalui wilayah
tetangga.

\textbf{Inferensi review.} Parameter galat spasial yang tetap bermakna
pada kedua model antarpinggiran sesuai dengan pengakuan artikel bahwa
variabel terlewat masih mempunyai pola spasial. Hal ini membatasi klaim
bahwa model telah menjelaskan mekanisme transformasi secara lengkap
(bagian \emph{The suburban commuting models}, hlm. 17-18).

\textbf{Ketegangan bukti internal.} Pernyataan arah dominan perjalanan
antara pinggiran dalam dan luar berbeda antara uraian Figure 7 dan
bagian \emph{Discussion}. Kesimpulan substantif tentang pinggiran mana
yang lebih banyak mengirim komuter harus dibaca dengan batas ini (hlm.
17-19).

\subsection{Conclusions}\label{conclusions-12}

\textbf{Kesimpulan penulis.} Komuter antarpinggiran JMA dipengaruhi oleh
ketimpangan karakteristik wilayah, termasuk kepadatan, biaya perumahan,
tingkat upah, dan pengangguran di wilayah sekitar. Regresi panel spasial
menunjukkan bahwa interaksi antarpinggiran penting untuk arus masuk dan
arus keluar, sedangkan komuter tradisional lebih banyak dijelaskan oleh
kepadatan dan biaya perumahan (bagian \emph{The suburban commuting
models}, hlm. 15-18; bagian \emph{Discussion}, hlm. 18-20).

\textbf{Kesimpulan penulis.} Desentralisasi pekerjaan manufaktur
memperluas pekerjaan pinggiran dan meningkatkan arus antarpinggiran,
tetapi hasilnya tidak seragam. Post-suburbanisasi manufaktur tidak
selalu menurunkan arus keluar. Sebagian penduduk harus mencari pekerjaan
di tempat lain karena struktur pekerjaan lokal tidak beragam atau tidak
sesuai, sehingga keseimbangan pekerjaan-perumahan berpotensi memburuk di
JMA dan metropolis lain dengan lintasan serupa (bagian
\emph{Conclusion}, hlm. 20).

\textbf{Kesimpulan penulis.} Perubahan pinggiran JMA tidak dapat
dipahami hanya sebagai penggantian arus menuju CBD dengan perjalanan
lokal yang lebih pendek. Arus antarpinggiran tumbuh lebih cepat daripada
arus tradisional, tujuan terkonsentrasi pada beberapa wilayah, dan
ketergantungan terhadap Jakarta tetap ada untuk pekerjaan dan layanan
tertentu (bagian \emph{Growing inter-suburban commuting in JMA}, hlm.
13-15; bagian \emph{Discussion} dan \emph{Conclusion}, hlm. 19-20).

\textbf{Inferensi review.} Kesimpulan yang paling terdukung adalah
adanya pertumbuhan arus dan asosiasi spasial yang berbeda menurut jenis
komuter. Generalisasi tentang dampak post-suburbanisasi terhadap
keseimbangan pekerjaan-perumahan masih berupa interpretasi karena
keseimbangan tersebut tidak menjadi variabel dependen yang diukur
langsung, kecocokan keterampilan tidak tersedia sebagai variabel, dan
model tidak memakai strategi identifikasi kausal.

\subsection{Contributions}\label{contributions-12}

\begin{itemize}
\tightlist
\item
  \textbf{Kontribusi yang dinyatakan penulis:} menambah bukti tentang
  hubungan post-suburbanisasi dan perubahan pola komuter di pinggiran
  Global South, yang pilihan lokasi huni penduduknya berbeda dari
  konteks Global North (bagian \emph{Introduction}, hlm. 3-4).
\item
  \textbf{Kontribusi yang dinyatakan penulis:} memberi bukti empiris
  tentang dampak subpusat pekerjaan baru di pinggiran terhadap mobilitas
  dan pembangunan regional berkelanjutan (bagian \emph{Introduction},
  hlm. 3-4).
\item
  \textbf{Kontribusi yang dinyatakan penulis:} menjelaskan pengaruh
  desentralisasi pekerjaan dan perubahan lingkungan pinggiran terhadap
  pola komuter intrametropolitan, khususnya arus horizontal
  antarpinggiran (bagian \emph{Conclusion}, hlm. 20).
\item
  \textbf{Kontribusi metodologis yang dilaporkan:} menggabungkan
  struktur panel dengan ketergantungan spasial, membandingkan tiga bobot
  spasial, dan memisahkan efek langsung dari interaksi dengan wilayah
  sekitar (bagian \emph{Data and methods}, hlm. 10-12; Tables 2-3, hlm.
  16-17).
\item
  \textbf{Inferensi review:} perbandingan tiga keluaran, yaitu arus
  masuk antarpinggiran, arus keluar antarpinggiran, dan arus
  tradisional, menunjukkan bahwa satu set karakteristik kota dapat
  mempunyai hubungan berbeda menurut arah perjalanan. Nilai tambah ini
  lebih kuat daripada klaim umum bahwa post-suburbanisasi selalu
  memperpendek atau memperpanjang perjalanan.
\item
  \textbf{Inferensi review:} kasus Kabupaten Bekasi dan Kabupaten
  Tangerang memperlihatkan bahwa pertumbuhan pekerjaan dalam sektor yang
  sama dapat berhubungan dengan hasil pasar kerja lokal yang berbeda.
  Artikel tidak menguji penyebab perbedaan itu secara langsung, tetapi
  perbandingan tersebut memberi dasar empiris bagi pembahasan variasi
  lokal dalam satu metropolis.
\end{itemize}

\subsection{Research Gap}\label{research-gap-12}

\textbf{Kesenjangan yang memotivasi penelitian.}

\begin{itemize}
\tightlist
\item
  Kajian komuter post-suburban terutama berasal dari negara maju dan
  berfokus pada desentralisasi jasa serta ritel. Penjelasan tentang
  komuter antarpinggiran dan restrukturisasi yang dipimpin manufaktur di
  Global South masih terbatas (bagian \emph{Introduction}, hlm. 2-4).
\item
  Teori dari Global North tidak sepenuhnya menjelaskan konteks Global
  South yang mengalami deindustrialisasi prematur, pekerjaan informal
  atau temporer, kemiskinan, segregasi spasial, dan keterbatasan pilihan
  rumah bagi kelas pekerja (bagian \emph{Introduction}, hlm. 2-3).
\item
  Kajian inter-suburban terdahulu lebih sering menjelaskan
  ketidakseimbangan pekerjaan-perumahan di negara maju. Campuran
  permukiman, pertanian, dan manufaktur dalam morfologi \emph{desakota}
  JMA memerlukan pengujian kontekstual tersendiri (bagian
  \emph{Introduction}, hlm. 2-3).
\item
  Dimensi spasial dan temporal komuter jarang diperhitungkan bersama.
  Perbedaan lintasan pembangunan antarpinggiran menghasilkan
  karakteristik khusus lokasi dan ketergantungan spasial yang dapat
  membuat estimasi nonspasial bias (bagian \emph{Introduction}, hlm.
  3-4).
\item
  Model gravitasi dan radiasi dinilai kurang mampu merepresentasikan
  segregasi sosial-geografis, spesialisasi lokasi, ketergantungan
  pilihan tujuan, dan perubahan antarwaktu. Artikel menempatkan panel
  spasial sebagai jawaban metodologis atas sebagian masalah tersebut
  (bagian \emph{Data and methods}, hlm. 10-11).
\end{itemize}

\textbf{Kesenjangan yang tetap terbuka menurut sumber.} Skala kota atau
kabupaten belum menggambarkan guna lahan pusat pekerjaan yang hanya
mencakup sebagian kecil wilayah. Penggunaan volume arus juga belum
mengukur friksi spasial seakurat waktu dan panjang perjalanan. Penulis
meminta penelitian pada skala ruang lebih kecil serta model yang
membedakan waktu dan panjang perjalanan komuter tradisional dari komuter
antarpinggiran (bagian \emph{Conclusion}, hlm. 21).

\textbf{Inferensi review.} Studi belum menutup kesenjangan antara
mekanisme tingkat individu dan hasil agregat wilayah. Keterampilan
pekerja, jenis pekerjaan individu, pilihan tempat tinggal, kualitas
amenitas, serta keputusan moda tidak diobservasi secara langsung,
walaupun semuanya dipakai dalam penjelasan teoretis atau pembahasan
hasil.

\subsection{Limitations}\label{limitations-12}

\textbf{Keterbatasan yang dinyatakan atau diakui penulis.}

\begin{itemize}
\tightlist
\item
  Delineasi administratif pada skala kota atau kabupaten mungkin tidak
  menggambarkan guna lahan pinggiran secara efektif. Pusat pekerjaan
  pinggiran hanya menempati bagian terbatas dari satu wilayah
  administrasi sehingga analisis pada skala lebih kecil diperlukan
  (bagian \emph{Conclusion}, hlm. 21).
\item
  Variabel dependen berupa volume arus komuter tidak menggambarkan
  friksi spasial setepat waktu atau panjang perjalanan. Artikel mengakui
  bahwa banyak studi komuter memakai waktu atau panjang perjalanan
  sebagai keluaran (bagian \emph{Conclusion}, hlm. 21).
\item
  Interaksi galat tetap ada pada model arus masuk dan arus keluar
  antarpinggiran. Penulis menjelaskan bahwa hasil tersebut menunjukkan
  pengaruh variabel terlewat yang berkorelasi spasial antarsatuan
  geografis (bagian \emph{The suburban commuting models}, hlm. 17-18).
\item
  Ketersediaan data komuter skala besar di Indonesia terbatas. Hilangnya
  data tingkat kabupaten dari SAKERNAS 2016 membuat penelitian berhenti
  pada 2015 (bagian \emph{Data and methods}, hlm. 9).
\end{itemize}

\textbf{Keterbatasan yang diinferensikan dari desain dan pelaporan.}

\begin{itemize}
\tightlist
\item
  Model mengestimasi asosiasi panel spasial tanpa intervensi eksogen,
  eksperimen, instrumen, atau rancangan kuasi-eksperimental yang
  dijelaskan. Bahasa sebab-akibat tentang transformasi kota dan
  desentralisasi pekerjaan melampaui identifikasi statistik yang
  tersedia dalam file.
\item
  Data diagregasikan pada kota atau kabupaten. Hubungan tingkat wilayah
  tidak membuktikan bahwa individu dengan karakteristik tertentu
  melakukan perjalanan karena alasan yang dinyatakan. Risiko kekeliruan
  ekologis tidak dibahas secara eksplisit.
\item
  Jumlah responden, jumlah observasi panel, bobot survei, galat
  pengambilan sampel, nilai hilang, dan perubahan desain survei
  antarperiode: Tidak disebutkan dalam file. Ketepatan total arus dan
  keterbandingan tahunan tidak dapat diaudit dari TXT.
\item
  Kesempatan kerja manufaktur bersumber dari Sensus Ekonomi 2016,
  sedangkan panel mencakup 2008-2015. Pemetaan ukuran itu ke tahun
  pengamatan dan kemungkinan penggunaan nilai tetap atau rekonstruksi
  historis: Tidak disebutkan dalam file.
\item
  Penyesuaian inflasi untuk upah, biaya rumah, dan biaya transportasi,
  serta konsistensi harga antarwilayah dan antarwaktu: Tidak disebutkan
  dalam file.
\item
  Transformasi logaritma diterapkan pada ketiga arus dan sebagian besar
  variabel penjelas, tetapi perlakuan terhadap nilai nol dan observasi
  ekstrem: Tidak disebutkan dalam file.
\item
  Ukuran volume absolut tidak dinormalisasi terhadap penduduk atau
  angkatan kerja. Wilayah besar dapat menghasilkan arus besar karena
  ukurannya, walaupun kepadatan dan pertumbuhan penduduk telah
  dimasukkan sebagai kovariat.
\item
  Variabel seperti kualifikasi pekerja, sektor pekerjaan komuter,
  pilihan moda, kualitas amenitas, dan lokasi rumah serta pekerjaan pada
  skala subkota tidak masuk ke model. Beberapa mekanisme utama dalam
  pembahasan karena itu tidak diuji langsung.
\item
  File tidak melaporkan galat baku, interval kepercayaan, nilai p,
  ambang tanda bintang, ukuran sampel, atau uji ketahanan selain
  pergantian matriks bobot. Replikasi dan penilaian ketidakpastian
  menjadi terbatas.
\item
  W2 bergantung pada ``puncak pekerjaan lokal tertinggi'', tetapi cara
  mengidentifikasi puncak dan sumber koordinat: Tidak disebutkan dalam
  file. W3 dibentuk dari arus komuter yang juga menjadi objek
  penelitian; konsekuensi penggunaan informasi arus dalam bobot spasial
  tidak dibahas.
\item
  Model arus masuk yang memuat lag spasial variabel dependen tidak
  disertai penyajian lengkap dampak langsung, tidak langsung, dan total.
  Table 3 hanya menampilkan koefisien dan lag variabel penjelas.
\item
  Arah perjalanan dominan antara pinggiran dalam dan luar tidak
  konsisten antara uraian Figure 7 dan bagian \emph{Discussion}.
  Ketidaksesuaian tersebut membatasi kesimpulan arah tanpa pemeriksaan
  gambar atau data asal.
\item
  Penelitian membahas perubahan pangsa kendaraan pribadi, tetapi moda
  bukan bagian dari model panel. Pengaruh jaringan jalan tol terhadap
  hubungan statistik tidak diisolasi sebagai variabel tersendiri.
\end{itemize}

\subsection{Challenges}\label{challenges-12}

\begin{itemize}
\tightlist
\item
  \textbf{Laporan sumber:} tantangan data adalah sedikitnya sumber yang
  menyediakan informasi komuter berskala besar di Indonesia dan tidak
  tersedianya data tingkat kabupaten untuk 2016 (bagian \emph{Data and
  methods}, hlm. 9).
\item
  \textbf{Laporan sumber:} arus antardaerah mempunyai ketergantungan
  pilihan tujuan, korelasi spasial, dan perubahan waktu. Model gravitasi
  atau pengulangan model potong lintang dinilai tidak memadai untuk
  menangani semuanya sekaligus (bagian \emph{Data and methods}, hlm.
  10-11).
\item
  \textbf{Laporan sumber:} lintasan pembangunan delapan pinggiran
  berbeda. Pekerjaan, penduduk, biaya rumah, upah, dan pengangguran
  tidak tersebar merata, sedangkan sebagian pusat pekerjaan hanya
  mencakup area kecil di dalam batas administrasi yang luas (bagian
  \emph{Introduction}, hlm. 3-4; bagian \emph{Conclusion}, hlm. 21).
\item
  \textbf{Interpretasi penulis:} tantangan pasar kerja adalah
  ketidakcocokan antara struktur pekerjaan manufaktur dan kualifikasi
  atau preferensi penduduk setempat. Pertumbuhan pekerjaan dapat menarik
  pekerja luar sekaligus mendorong penduduk lokal keluar untuk pekerjaan
  jasa (bagian \emph{Growing inter-suburban commuting in JMA}, hlm.
  14-15; bagian \emph{Discussion}, hlm. 19).
\item
  \textbf{Interpretasi penulis:} meningkatnya ketergantungan pada
  kendaraan pribadi dan bertumpuknya komuter antarpinggiran,
  tradisional, serta terbalik pada koridor yang sama dapat menambah
  tekanan transportasi pada masa mendatang (bagian \emph{Growing
  inter-suburban commuting in JMA}, hlm. 13; bagian \emph{Discussion},
  hlm. 19-20).
\item
  \textbf{Inferensi review:} pemilihan model juga tidak sederhana. Uji
  Hausman memberi hasil yang berbeda menurut bobot, beberapa uji
  komponen galat tidak bertanda signifikan dalam Table 2, dan model
  akhir masih mempunyai korelasi galat spasial. Kesimpulan perlu
  mempertahankan ketidakpastian spesifikasi tersebut.
\item
  \textbf{Inferensi review:} reproduksi analisis memerlukan penyelarasan
  beberapa survei dan publikasi tahunan, geometri wilayah, lokasi puncak
  pekerjaan, serta konstruksi tiga bobot spasial. Kode, data olahan, dan
  rincian penyelarasan yang cukup untuk menjalankan ulang estimasi:
  Tidak disebutkan dalam file.
\end{itemize}

\subsection{Practical Implications}\label{practical-implications-12}

\begin{itemize}
\tightlist
\item
  \textbf{Rekomendasi penulis:} kebijakan upah sebagai penarik komuter
  perlu disertai insentif bagi keragaman pekerjaan. Upah manufaktur yang
  tinggi saja tidak memenuhi kebutuhan penduduk yang mencari pekerjaan
  jasa atau komersial (bagian \emph{Conclusion}, hlm. 20).
\item
  \textbf{Rekomendasi penulis:} pengembangan kawasan industri perlu
  diintegrasikan dengan kebijakan tata guna lahan agar pertumbuhan
  klaster manufaktur dan perubahan penggunaan lahan dapat dikelola
  bersama (bagian \emph{Conclusion}, hlm. 20-21).
\item
  \textbf{Rekomendasi penulis:} integrasi tata guna lahan perlu disertai
  perencanaan infrastruktur transportasi yang terdesentralisasi pada
  tingkat lokal untuk memperbaiki hubungan antara pinggiran dalam dan
  luar (bagian \emph{Conclusion}, hlm. 20-21).
\item
  \textbf{Rekomendasi penulis:} infrastruktur transportasi pinggiran
  perlu diperkuat untuk mengantisipasi pertumbuhan komuter
  antarpinggiran dan mobilitas horizontal antara wilayah pinggiran luar
  (bagian \emph{Conclusion}, hlm. 21).
\item
  \textbf{Rekomendasi penulis:} pengembangan berorientasi transit yang
  terintegrasi di sekitar pinggiran JMA perlu dipertimbangkan untuk
  mengurangi ketergantungan pada mobil sekaligus mendukung pertumbuhan
  mobilitas antarpinggiran (bagian \emph{Conclusion}, hlm. 21).
\item
  \textbf{Inferensi review:} pemantauan transportasi sebaiknya
  membedakan arus masuk, arus keluar antarpinggiran, dan arus menuju
  Jakarta. Koefisien serta interaksi spasialnya berbeda sehingga satu
  indikator total komuter dapat menutupi kebutuhan kebijakan yang
  berlawanan.
\item
  \textbf{Inferensi review:} kebijakan kawasan industri memerlukan
  pemantauan komposisi pekerjaan dan kecocokannya dengan tenaga kerja
  lokal, bukan hanya jumlah pekerjaan baru. Artikel menunjukkan bahwa
  pertumbuhan jumlah pekerjaan dapat berjalan bersama kenaikan arus
  keluar, tetapi tidak mengukur program penyesuaian keterampilan
  tertentu.
\end{itemize}

\subsection{Applications}\label{applications-12}

\begin{itemize}
\tightlist
\item
  \textbf{Aplikasi empiris yang dilaporkan:} kerangka diterapkan pada
  delapan kota atau kabupaten pinggiran JMA untuk menjelaskan arus
  antarpinggiran selama 2008-2015 dan membandingkannya dengan arus
  tradisional menuju Jakarta (bagian \emph{Data and methods}, hlm. 8-12;
  bagian \emph{Results}, hlm. 13-18).
\item
  \textbf{Aplikasi analitis yang diperagakan:} panel spasial dapat
  mengenali hubungan suatu wilayah dengan kondisi tetangganya, menguji
  interaksi endogen dan galat, serta mengukur efek limpahan kepadatan,
  biaya rumah, upah, dan pengangguran (Tables 2-3, hlm. 16-17; uraian
  hasil, hlm. 17-18).
\item
  \textbf{Aplikasi perencanaan yang dinyatakan:} hasil dipakai untuk
  mengarahkan integrasi kawasan industri, tata guna lahan, keragaman
  pekerjaan, dan jaringan transportasi antarpinggiran (bagian
  \emph{Conclusion}, hlm. 20-21).
\item
  \textbf{Inferensi review:} metode ini dapat menjadi alat diagnosis
  untuk metropolis lain yang mempunyai data arus tahunan dan
  ketergantungan antarsatuan wilayah. Penerapannya memerlukan definisi
  inti-pinggiran, skala ruang, bobot spasial, dan indikator pasar kerja
  yang sesuai konteks. Artikel tidak menguji validitas eksternal di luar
  JMA.
\item
  \textbf{Inferensi review:} analisis dapat membantu mengenali lokasi
  yang sekaligus menjadi pusat pekerjaan dan asal arus keluar, seperti
  Kabupaten Bekasi dalam uraian artikel. Identifikasi tersebut belum
  sama dengan evaluasi dampak kebijakan tertentu karena model tidak
  membandingkan skenario intervensi.
\end{itemize}

\subsection{Future Research
(Optional)}\label{future-research-optional-12}

Penulis menyebut dua arah penelitian secara eksplisit pada bagian
\emph{Conclusion}, hlm. 21.

\begin{enumerate}
\def\labelenumi{\arabic{enumi}.}
\tightlist
\item
  Menganalisis interaksi guna lahan dan spasial antarpinggiran pada
  skala geografis yang lebih kecil daripada kota atau kabupaten agar
  lokasi pusat pekerjaan dan variasi guna lahan tergambarkan lebih
  akurat.
\item
  Menggunakan waktu dan panjang perjalanan sebagai variabel dependen,
  membedakan keduanya untuk komuter tradisional dan komuter
  antarpinggiran, lalu membandingkan hasil modelnya.
\end{enumerate}

Agenda penelitian tambahan yang dinyatakan penulis: Tidak disebutkan
dalam file. Review ini tidak menambahkan agenda eksternal.

\subsection{Provenance Notes}\label{provenance-notes-12}

\begin{itemize}
\tightlist
\item
  Review ini hanya memakai isi file TXT
  \texttt{/Users/mac/Documents/Mac/{[}2{]}\ Obsidian\ Vault/zettelkasten/source/journal/j23_sadewo_delik_the_role_of_urban_transformation_on_inter_suburban_commuting_10_1080_02723638_2022_2125649.txt}
  untuk seluruh klaim substantif. Seluruh 1.233 baris dibaca, termasuk
  lembar informasi penerbit, isi artikel, tabel yang terekstrak, ucapan
  terima kasih, pernyataan konflik kepentingan, pendanaan, ORCID, dan
  daftar pustaka. Tidak ada sumber eksternal yang dipakai untuk
  melengkapi atau mengoreksi artikel.
\item
  TXT tidak memuat blok metadata ekstraksi yang menyebut nama PDF asal,
  tanggal ekstraksi, perangkat ekstraksi, atau jumlah halaman PDF.
  Informasi tersebut: Tidak disebutkan dalam file. Review ini membaca
  TXT dan tidak memeriksa PDF secara visual.
\item
  Nomor halaman mengikuti header cetak yang tampak dalam TXT. Header
  hlm. 2-26 terlihat; judul, abstrak, dan awal pendahuluan sebelum
  header hlm. 2 dirujuk sebagai ``halaman pembuka artikel'' tanpa
  menetapkan nomor halaman melalui dugaan. Tubuh artikel berakhir pada
  hlm. 21 dan daftar pustaka berlanjut sampai hlm. 26.
\item
  Teks memuat Figures 1-7, Tables 1-3, dan Equations 1-8. Isi gambar
  tidak tersedia sebagai data visual dalam TXT. Review hanya menggunakan
  caption dan uraian naratif. Angka tabel digunakan sejauh susunan baris
  dan kolomnya masih terbaca.
\item
  Tata letak persamaan dan beberapa simbol statistik terdistorsi dalam
  ekstraksi, termasuk notasi galat, varians, dan simbol parameter.
  Review tidak merekonstruksi simbol yang hilang dan memakai penjelasan
  verbal artikel untuk menerangkan fungsi model.
\item
  Table 3 menggunakan satu, dua, dan tiga tanda bintang, tetapi catatan
  yang menjelaskan ambang signifikansinya tidak tampak dalam TXT. Review
  mempertahankan jumlah tanda tersebut dan tidak mengarang nilai
  probabilitas.
\item
  Narasi model menyatakan W1 menghasilkan AIC dan BIC terendah untuk
  ``setiap pola komuter''. Table 2 dan bagian \emph{Discussion}
  menunjukkan bahwa W3 justru terendah untuk komuter tradisional. Review
  mengikuti W1 untuk dua model antarpinggiran dan W3 untuk model
  tradisional, sambil mencatat ketidaksesuaian itu.
\item
  Narasi uji komponen galat menyatakan hasilnya signifikan secara umum,
  tetapi Table 2 menampilkan beberapa statistik tanpa tanda
  signifikansi. Review melaporkan batas tersebut dan tidak menganggap
  semua kombinasi bobot memperoleh dukungan yang sama.
\item
  Uraian Figure 7 menyatakan rata-rata arus pinggiran dalam menuju
  pinggiran luar sebesar 42\% dan arah sebaliknya 28\%. Bagian
  \emph{Discussion} kemudian menyatakan lebih banyak arus keluar
  bergerak dari pinggiran luar menuju pinggiran dalam. Isi visual Figure
  7 tidak dapat dipakai dari TXT untuk menyelesaikan pertentangan ini.
\item
  Pembahasan model tradisional menyebut ``lag spasial jarak'' sebagai
  bernilai negatif, sedangkan Table 3 tidak mencantumkan variabel jarak
  dan menampilkan lag spasial kepadatan sebesar -1,92 dengan tiga tanda
  bintang. Review menggunakan nama variabel dalam Table 3 dan tidak
  menggantinya dengan jarak.
\item
  Table 1 menyatakan bahwa kesempatan kerja manufaktur bersumber dari
  Sensus Ekonomi 2016 untuk panel 2008-2015. Penjelasan konstruksi seri
  tahunannya tidak tersedia. Review tidak mengasumsikan metode
  interpolasi, ekstrapolasi, atau rekonstruksi tertentu.
\item
  Klaim bahwa pusat industri menarik pekerja luar dilaporkan sebagai
  hasil deskriptif dan interpretasi penulis. Koefisien langsung
  kesempatan kerja dalam model arus masuk tidak bertanda signifikan.
  Review menjaga kedua lapisan bukti tersebut tetap terpisah.
\item
  Dua keterangan dukungan dana muncul pada bagian berbeda, yaitu
  Bappenas PHRD IV Program dalam \emph{Acknowledgements} dan ITB
  Research Program 2020 kategori B dalam \emph{Funding}. Keduanya
  dilaporkan apa adanya dan tidak dianggap saling menggantikan.
\item
  Semua karya dalam daftar pustaka dan sitasi internal diperlakukan
  sebagai laporan J23 atas literatur. Isi, metadata, dan temuan
  sumber-sumber tersebut tidak diperiksa secara independen.
\item
  Batas provenance utama adalah ketergantungan pada ekstraksi teks.
  Review dapat memeriksa isi verbal dan tabel yang terbaca, tetapi tidak
  dapat memastikan kesetiaan tata letak terhadap dokumen asal, membaca
  peta atau grafik secara visual, memverifikasi mikrodata, ataupun
  menjalankan ulang estimasi.
\end{itemize}

\section{Review J25}\label{review-j25}

\textbf{Identitas sumber}

\begin{itemize}
\tightlist
\item
  Kode: J25.
\item
  Judul: \emph{The continuity and change in mega-urbanization in
  Indonesia: A survey of Jakarta-Bandung Region (JBR) development}.
\item
  Penulis: Tommy Firman.
\item
  Afiliasi: Research Cluster on Regional and Rural Development, School
  of Architecture, Planning, and Policy Development, Institute of
  Technology, Bandung 41032, Indonesia.
\item
  Jenis sumber menurut isi: artikel jurnal. Status telaah sejawat: Tidak
  disebutkan dalam file.
\item
  Jurnal: \emph{Habitat International}.
\item
  Volume, tahun, dan halaman: 33 (2009), 327-339. Nomor terbitan: Tidak
  disebutkan dalam file.
\item
  ISSN yang tercetak: 0197-3975.
\item
  DOI: 10.1016/j.habitatint.2008.08.005.
\item
  Copyright: 2008 Elsevier Ltd.~Tanggal penerimaan, revisi, persetujuan,
  dan publikasi daring: Tidak disebutkan dalam file.
\item
  Konflik kepentingan: Tidak disebutkan dalam file.
\item
  Kata kunci: mega-urbanization; Indonesia; Jakarta-Bandung Region.
\item
  Objek kajian: perkembangan Jakarta-Bandung Region (JBR), terutama
  Jakarta Metropolitan Area (JMA) dan Bandung Metropolitan Area (BMA),
  dalam konteks krisis ekonomi akhir 1990-an, pemulihan hingga
  pertengahan 2000-an, serta otonomi daerah dan desentralisasi fiskal
  sejak 2001 (Abstract, hlm. 327; Introduction, hlm. 327-329).
\item
  Pendanaan: penelitian didanai sebagian oleh Institute of Research,
  Institute of Technology, Bandung. Penulis menyatakan bertanggung jawab
  sendiri atas kekurangan dan kesalahan artikel (Acknowledgement, hlm.
  337).
\end{itemize}

\textbf{Penanda epistemik yang digunakan}

\begin{itemize}
\tightlist
\item
  \textbf{Laporan sumber} berarti pernyataan, angka, data, atau kondisi
  yang disampaikan di dalam artikel, termasuk laporan artikel atas
  sumber sekundernya.
\item
  \textbf{Interpretasi penulis} berarti penjelasan sebab, makna,
  penilaian, atau implikasi yang diberikan Tommy Firman terhadap bahan
  yang disajikan.
\item
  \textbf{Inferensi pengolah} berarti pembacaan analitis review yang
  hanya diturunkan dari TXT J25. Inferensi ini bukan temuan langsung
  artikel dan bukan bukti tambahan.
\end{itemize}

\subsection{Summarized Abstract}\label{summarized-abstract-13}

\textbf{Laporan sumber.} Artikel menelaah kesinambungan dan perubahan
pembentukan kawasan mega-urban Jakarta-Bandung. Perkembangan fisik JMA
dan BMA dilaporkan telah membentuk sabuk urban sekitar 200 km dari
Jakarta hingga Bandung, dengan campuran kegiatan pertanian, industri,
perdagangan, dan permukiman yang mengaburkan perbedaan desa-kota. Proses
tersebut dinyatakan berlanjut, bahkan lebih cepat daripada pada 1980-an
hingga pertengahan 1990-an, sementara JMA dan BMA berubah dari kota
berinti tunggal menjadi berinti majemuk (Abstract, hlm. 327; Concluding
remarks, hlm. 337).

\textbf{Laporan sumber.} Kajian menggabungkan uraian konseptual mengenai
mega-urbanisasi Asia dengan statistik sekunder tentang penduduk urban,
migrasi, komuter, lokalitas urban, perubahan guna lahan, kota baru,
kawasan industri, investasi, transportasi, dana alokasi umum, dan PDRB.
Artikel juga menelusuri perlambatan pembangunan selama krisis ekonomi
1998 hingga awal 2000-an serta implikasi awal otonomi daerah dan
desentralisasi fiskal setelah 2001 (Introduction, hlm. 327-329; Economic
crises, hlm. 330-331; Recent socioeconomic and physical development,
hlm. 331-335; Early reform and decentralization era, hlm. 335-337).

\textbf{Interpretasi penulis.} Perkembangan yang dipimpin modal swasta
telah memperluas kota baru dan kawasan industri, tetapi juga
menghasilkan kota asrama, segregasi berdasarkan gaya hidup dan
pendapatan, konversi lahan yang tidak terkendali, kemacetan, pencemaran
udara, ekstraksi air tanah berlebihan, serta tekanan terhadap kawasan
resapan. Desentralisasi memberi pemerintah lokal kewenangan dan sumber
fiskal lebih besar, tetapi parokialisme, kapasitas yang lemah, dan
dorongan memaksimalkan pendapatan lokal dinilai memecah pengelolaan
wilayah yang secara fungsional menyatu (Abstract, hlm. 327; Land
conversion, hlm. 332-333; New town, industrial estate and road
development, hlm. 333-335; Early reform and decentralization era, hlm.
335-337).

\textbf{Inferensi pengolah.} Artikel paling tepat dibaca sebagai survei
deskriptif dan sintesis regional berbasis berbagai sumber sekunder,
bukan sebagai survei responden, eksperimen, atau estimasi kausal. Bukti
yang disajikan mendukung dokumentasi koeksistensi integrasi fungsional
dan fragmentasi kelembagaan, tetapi tidak mengisolasi secara statistik
besarnya pengaruh krisis, investasi, infrastruktur, atau desentralisasi
terhadap setiap perubahan spasial.

\subsection{Problem Statement}\label{problem-statement-13}

\textbf{Laporan sumber.} Urbanisasi Asia kontemporer memperlihatkan
ekspansi fisik melewati batas kota dan metropolitan, percampuran
kegiatan urban dengan pertanian, serta kaburnya pembedaan desa-kota.
Upaya banyak negara Asia untuk mengendalikan perluasan metropolitan dan
memperlambat migrasi desa-kota dinilai tidak banyak berhasil. Studi awal
1990-an telah menyimpulkan bahwa Jabotabek dan BMA terintegrasi secara
fisik menjadi sabuk urban, tetapi perkembangan berikutnya berlangsung
dalam kondisi ekonomi dan politik yang berubah (Introduction, hlm.
327-328).

\textbf{Laporan sumber.} Dua kejadian diposisikan sebagai konteks
perubahan utama: krisis ekonomi akhir 1990-an dan awal 2000-an, serta
kebijakan otonomi daerah dan desentralisasi fiskal yang mulai diterapkan
pada 2001. Pertanyaan yang dinyatakan artikel ialah sejauh mana
reformasi otonomi daerah dan desentralisasi fiskal memengaruhi
perkembangan JBR (Introduction, hlm. 328).

\textbf{Interpretasi penulis.} Penulis menolak menjadikan kemiripan atau
perbedaan dengan urbanisasi Barat sebagai pertanyaan utama.
Mega-urbanisasi perlu dipahami melalui kondisi sosial ekonomi, politik,
fisik, aktor, dan institusi pada tempat yang dikaji. Dengan merujuk
Shatkin, artikel mendorong pergeseran dari model kota dunia menuju
pemeriksaan yang lebih membumi atas interaksi aktor dan institusi
global-lokal (Introduction, hlm. 328; Mega-urbanization in Asia, hlm.
329-330).

\textbf{Batas masalah menurut sumber.} Penulis menyatakan bahwa
penilaian dampak desentralisasi masih prematur karena pelaksanaannya
baru berlangsung tujuh tahun pada 2001-2008 dan mencakup dimensi sosial
ekonomi, budaya, dan politik yang luas. Karena itu, artikel hanya
menganggap analisis atas potensi dampaknya terhadap perkembangan urban
dan regional sebagai hal yang mungkin dilakukan (Introduction, hlm.
328).

\textbf{Inferensi pengolah.} Masalah analitis artikel memiliki
ketegangan internal: judul dan kesimpulan membahas kesinambungan serta
perubahan mega-urbanisasi secara luas, sedangkan pertanyaan eksplisit
menyoroti dampak desentralisasi. Tidak ada rancangan yang menghubungkan
kedua fokus itu melalui pengujian kausal; hubungan tersebut dibangun
melalui sintesis deskriptif.

\subsection{Objectives}\label{objectives-13}

\begin{itemize}
\tightlist
\item
  \textbf{Laporan sumber:} memeriksa kesinambungan dan perubahan proses
  pembentukan mega-urban di JBR dalam konteks perkembangan sosial
  ekonomi dan kondisi politik Indonesia mutakhir (Introduction, hlm.
  327-328).
\item
  \textbf{Laporan sumber:} menilai sejauh mana otonomi daerah dan
  desentralisasi fiskal berpotensi memengaruhi perkembangan JBR, sambil
  mengakui bahwa evaluasi dampak definitif masih terlalu dini
  (Introduction, hlm. 328).
\item
  \textbf{Laporan sumber:} menempatkan kasus JBR dalam konteks teori dan
  pengalaman mega-urbanisasi Asia, globalisasi, hubungan lokal-global,
  serta hubungan desa-kota (Mega-urbanization in Asia, hlm. 329-330).
\item
  \textbf{Laporan sumber:} menganalisis dampak krisis ekonomi terhadap
  perkembangan fisik dan ekonomi JBR (Economic crises, hlm. 330-331).
\item
  \textbf{Laporan sumber:} menggambarkan perkembangan sosial ekonomi dan
  fisik mutakhir melalui penduduk urban, migrasi dan komuter, perubahan
  guna lahan, kota baru, kawasan industri, dan infrastruktur jalan
  (Recent socioeconomic and physical development, hlm. 331-335).
\item
  \textbf{Laporan sumber:} membahas implikasi awal reformasi dan
  desentralisasi bagi kapasitas, perilaku, koordinasi, dan pembiayaan
  pemerintah lokal di JBR (Early reform and decentralization era, hlm.
  335-337).
\end{itemize}

\subsection{Literature Survey}\label{literature-survey-13}

\textbf{Sifat survei literatur.} Artikel menyajikan tinjauan naratif dan
kontekstual. Strategi pencarian, pangkalan data, kata kunci penelusuran,
periode pencarian, kriteria inklusi atau eksklusi, jumlah karya yang
disaring, dan penilaian mutu sumber: Tidak disebutkan dalam file. Oleh
karena itu, survei ini tidak dapat diperlakukan sebagai tinjauan
sistematis.

\begin{itemize}
\tightlist
\item
  \textbf{Laporan sumber atas literatur mega-urbanisasi:} Brennan, Hugo,
  Jones, McGee, Lin, dan karya Firman digunakan untuk menjelaskan
  perluasan metropolitan melampaui batas administratif, percampuran guna
  lahan urban dan pertanian, hubungan desa-kota yang intens, serta
  kaburnya batas desa-kota. Globalisasi, investasi domestik dan asing,
  infrastruktur, dan ekonomi aglomerasi ditempatkan sebagai pendorong
  utama (Introduction, hlm. 327; Mega-urbanization in Asia, hlm. 329).
\item
  \textbf{Laporan sumber atas restrukturisasi spasial:} artikel
  merangkum tujuh ciri restrukturisasi mega-urban Asia, yaitu kegiatan
  ekonomi berskala global, pembagian fungsi inti-pinggiran, perubahan
  dari satu inti ke banyak inti, perubahan guna lahan pusat dan konversi
  lahan pertanian pinggiran, pembangunan infrastruktur skala besar,
  peningkatan produksi ruang, serta pertumbuhan komuter dan waktu
  komuter (Mega-urbanization in Asia, hlm. 329).
\item
  \textbf{Laporan sumber atas hubungan lokal-global dan desa-kota:}
  McGee digunakan untuk membaca mega-urbanisasi sebagai pertautan
  hubungan lokal-global dan desa-kota. Ekspansi fisik masuk ke wilayah
  dengan kegiatan perdesaan intensif, sedangkan arus modal, komoditas,
  manusia, dan informasi menghubungkan inti kota dengan pasar global
  sekaligus wilayah metropolitan perluasan yang menyediakan air dan
  pangan (Mega-urbanization in Asia, hlm. 329).
\item
  \textbf{Laporan sumber atas perbandingan Barat-Asia:} Borsdorf,
  Phelps, dan Soja digunakan untuk membahas \emph{post-suburbia},
  struktur polisentris yang terfragmentasi, dan pembalikan relasi
  pusat-pinggiran. McGee menekankan skala luar biasa ekspansi kota Asia,
  sedangkan Dick dan Rimmer, Webster, serta Soja menekankan kesamaan
  proses lintas negara. Firman menerima kemungkinan kesamaan, tetapi
  menegaskan perlunya meletakkan proses dalam masyarakat dan tempatnya
  serta melakukan perbandingan antarsetting (Mega-urbanization in Asia,
  hlm. 329-330).
\item
  \textbf{Laporan sumber atas krisis:} World Bank, McLeod, kajian
  Firman, media properti, dan sumber lain dipakai untuk menghubungkan
  krisis dengan utang tanpa lindung nilai, salah kelola moneter, tata
  kelola buruk, kontraksi ekonomi, masalah perbankan dan properti,
  pengangguran, serta tanah proyek yang menganggur (Economic crises,
  hlm. 330-331).
\item
  \textbf{Laporan sumber atas konversi lahan dan kota baru:} Grigg,
  Leaf, Leisch, Dijkgraaf, Douglass, dan karya Firman mendukung uraian
  tentang rente urban, hak eksklusif izin lokasi, spekulasi, aliran
  modal ke sektor properti, fungsi kota asrama, dan segregasi sukarela
  pada kota baru (Land conversion, hlm. 332-333; New town, industrial
  estate and road development, hlm. 333-334).
\item
  \textbf{Laporan sumber atas desentralisasi:} Alm, Brodjonegoro, Booth,
  Rasyid, Fritzen dan Lim, Hadiz, Kuncoro, serta lainnya digunakan untuk
  merangkum tujuan desentralisasi, yaitu mendekatkan pemerintah kepada
  masyarakat, meningkatkan partisipasi, akuntabilitas, efisiensi,
  responsivitas, dan pelayanan, sekaligus risiko korupsi
  terdesentralisasi, politik uang, kapasitas lokal yang lemah, dan
  lemahnya kewenangan perpajakan (Early reform and decentralization era,
  hlm. 335-336).
\end{itemize}

\textbf{Inferensi pengolah.} Literatur berfungsi sebagai kerangka
interpretasi bagi rangkaian statistik yang heterogen, bukan sebagai
korpus yang dianalisis dengan prosedur eksplisit. Sejumlah rujukan
berupa makalah tidak terbit, surat kabar, majalah properti, laporan
konsultan, tesis sarjana, dan dokumen pemerintah; artikel tidak
melaporkan cara menimbang perbedaan mutu atau kepentingan sumber
tersebut (References, hlm. 337-339).

\subsection{Method Used}\label{method-used-13}

\textbf{Laporan sumber dan inferensi pengolah atas rancangan.} Artikel
tidak memiliki bagian metode tersendiri dan tidak menyebut desain
penelitian dengan terminologi metodologis formal. Berdasarkan pekerjaan
yang tampak dalam file, review mengklasifikasikannya sebagai kajian
deskriptif regional yang menyandingkan literatur, data statistik
sekunder, data tata guna lahan, laporan pemerintah dan konsultan, serta
pemberitaan dan publikasi sektor properti. Pembahasan disusun menurut
konteks teori, krisis ekonomi, perkembangan sosial ekonomi dan fisik,
desentralisasi, lalu kesimpulan (Introduction, hlm. 329; struktur
artikel pada akhir Introduction).

\textbf{Langkah analitis yang tampak dalam file:}

\begin{itemize}
\tightlist
\item
  membandingkan jumlah, proporsi, dan laju pertumbuhan penduduk urban
  antarunit administratif serta antarperiode 1980-1990 dan 1990-2000
  (Table 1 dan uraian Urban population, hlm. 331);
\item
  membandingkan jumlah serta proporsi lokalitas urban dan perdesaan pada
  1999 dan 2005 (Table 2 dan uraian Urban population, hlm. 331-332);
\item
  membandingkan luas kategori guna lahan Bopunjur dan BMA pada 1994 dan
  2001, termasuk perubahan hektare dan persentase yang disediakan tabel
  (Table 3 dan uraian Land conversion, hlm. 332-333);
\item
  menginventarisasi proyek kota baru berluas sekurangnya 1.000 ha pada
  2003 dan kawasan industri pada 2005 (Tables 4-5, hlm. 333-334);
\item
  menyandingkan statistik krisis, investasi, perjalanan, dan
  infrastruktur dari berbagai tahun untuk menggambarkan perubahan
  regional (Economic crises, hlm. 330-331; New town, industrial estate
  and road development, hlm. 333-335);
\item
  membandingkan Dana Alokasi Umum 2001, 2003, dan 2006 serta menampilkan
  PDRB 2000 pemerintah lokal JBR (Table 6, hlm. 336);
\item
  menafsirkan pola tersebut dalam kaitannya dengan kapital swasta,
  perubahan struktur inti, konversi lahan, dampak lingkungan, dan
  fragmentasi tata kelola (Concluding remarks, hlm. 337).
\end{itemize}

\textbf{Tidak disebutkan dalam file:} prosedur pemilihan sumber;
protokol pengumpulan atau harmonisasi data; definisi operasional yang
seragam untuk batas JBR sepanjang seluruh analisis; perangkat lunak;
kode analisis; penanganan data hilang; pembobotan; uji statistik; ukuran
ketidakpastian; model multivariat; strategi identifikasi kausal;
wawancara; observasi lapangan; atau survei responden yang dilakukan
khusus untuk artikel ini.

\textbf{Inferensi pengolah atas desain.} Kata \emph{survey} pada judul
berarti survei atau tinjauan perkembangan regional, bukan bukti adanya
survei sampel primer. Perbandingan sebelum-sesudah yang tersebar di
artikel menunjukkan perubahan waktu, tetapi tanpa kelompok pembanding
atau kontrol atas perubahan batas administratif, sehingga tidak dapat
dengan sendirinya menetapkan penyebab.

\subsection{Dataset}\label{dataset-13}

\begin{itemize}
\tightlist
\item
  \textbf{Penduduk urban:} Table 1 menggunakan \emph{Central Board of
  Statistics} 1991 dan 2001 untuk angka 1990 dan 2000 serta laju
  pertumbuhan 1980-1990 dan 1990-2000 pada unit kota dan kabupaten JBR
  (Table 1, hlm. 331). Nama dataset selain publikasi sensus, berkas,
  tautan, dan prosedur pengolahan: Tidak disebutkan dalam file.
\item
  \textbf{Migrasi:} Sensus Penduduk 2000 digunakan untuk perpindahan
  mantan penduduk Jakarta ke Bekasi, Tangerang, Bogor, dan Depok serta
  migran masuk lima tahun terakhir pada 1995-2000. Angka migrasi BMA
  dilaporkan melalui West Java Office of Central Board of Statistics
  2001 sebagaimana dikutip \emph{Kompas} 2006 (Urban population, hlm.
  331).
\item
  \textbf{Komuter:} angka JICA untuk 2002, sebagaimana dikutip Hidayat
  2007, melaporkan lebih dari tiga juta komuter antara Jakarta dan
  daerah sekitarnya. Statistik koridor Jakarta-Bandung 2004 berasal dari
  Lubis, Dharmowijoyo, dan Armijaya 2005; statistik kereta JMA 2002 dan
  kendaraan 1993-2001 dikutip dari Hata 2003 (Introduction, hlm. 329;
  New town, industrial estate and road development, hlm. 335).
\item
  \textbf{Lokalitas urban-perdesaan:} Table 2 membandingkan 1999 dan
  2005 untuk Jabodetabek dan Bandung Raya, dihitung dari Gardiner dan
  Gardiner 2006 (Table 2, hlm. 332). Kriteria klasifikasi lokalitas dan
  prosedur perhitungannya: Tidak disebutkan dalam file.
\item
  \textbf{Guna lahan:} Table 3 menggunakan data West Java Province
  Planning Board untuk luas 13 kategori guna lahan di Bopunjur dan BMA
  pada 1994 dan 2001 (Table 3, hlm. 333). Resolusi spasial, teknik
  pemetaan, sumber citra, klasifikasi, dan akurasi perubahan: Tidak
  disebutkan dalam file.
\item
  \textbf{Kota baru:} Table 4 menggunakan Real Estate Indonesia untuk
  proyek kota baru seluas 1.000 ha atau lebih pada 2003. Statistik lain
  mengenai operasi proyek, pinjaman properti, penghuni, dan fasilitas
  berasal dari majalah properti serta karya yang dirujuk (Table 4 dan
  uraian, hlm. 333-334).
\item
  \textbf{Kawasan industri dan investasi:} Table 5 menggunakan Collier
  International 2005 dan Investment Coordinating Body 2005. Statistik
  persetujuan investasi asing dan domestik 2000-2004 berasal dari
  Central Board of Statistics 2006 (Table 5, hlm. 334; uraian, hlm.
  335).
\item
  \textbf{Krisis ekonomi dan properti:} statistik PDRB, investasi,
  pinjaman, transaksi properti, proyek terhenti, dan lahan industri
  menganggur berasal dari Central Board of Statistics, World Bank,
  \emph{Properti Indonesia}, \emph{Tempo}, dan karya Firman yang dirujuk
  (Economic crises, hlm. 330-331).
\item
  \textbf{Keuangan pemerintah lokal:} Table 6 menggunakan Ministry of
  Finance dan Central Board of Statistics untuk DAU 2001, 2003, dan 2006
  serta PDRB 2000 (Table 6, hlm. 336).
\item
  \textbf{Data lingkungan dan komersial lain:} perubahan ruang terbuka
  hijau Jakarta dikutip dari \emph{Tempo} 2007, luas pusat perbelanjaan
  dari \emph{Tempo} 2006, dan konversi sawah di pinggiran Bandung dari
  Sarvitri, Rahayu, dan Sari 1997 (Land conversion, hlm. 332; New town,
  industrial estate and road development, hlm. 334).
\end{itemize}

\textbf{Inferensi pengolah.} Dataset artikel bukan satu basis data
terpadu. Tahun, unit, batas wilayah, skala, dan jenis penerbit sumber
berbeda-beda; file tidak menjelaskan proses penyelarasan atau validasi
silang. Karena itu, angka-angka tersebut lebih kuat sebagai potret
deskriptif perindikator daripada sebagai panel regional yang langsung
sebanding.

\subsection{Population / Sample}\label{population-sample-13}

\textbf{Cakupan geografis yang dilaporkan sumber.} JBR terdiri atas JMA
dan BMA serta membentang lebih dari 12.275 km2. Jakarta dan Bandung
hanya menempati masing-masing 164 km2 dan 81 km2 dari wilayah itu. Total
penduduk JBR dilaporkan sekitar 32 juta pada 2004; Jakarta sebagai inti
memiliki sekitar 8,5 juta penduduk dan Bandung sekitar 2 juta. JMA
sendiri disebut hanya mencakup 0,33\% luas daratan nasional, tetapi
menampung 12\% penduduk Indonesia dan menghasilkan hampir seperempat PDB
nasional (Introduction, hlm. 328-329).

\textbf{Cakupan administratif.} Narasi menyebut JMA atau Jabodetabekjur
mencakup DKI Jakarta; Kota Bogor, Depok, Tangerang, dan Bekasi; serta
Kabupaten Bogor, Tangerang, Bekasi, dan Cianjur. BMA dalam narasi
terdiri atas Kota Bandung dan Kabupaten Bandung (Introduction, hlm.
328-329). Namun, cakupan aktual berubah menurut data: Table 1 memasukkan
Kabupaten Karawang ke dalam 12 unit JBR, sedangkan Table 2 menggabungkan
Bandung dan Sumedang dalam Bandung Raya (Tables 1-2, hlm. 331-332).
Perbedaan ini dipertahankan sebagai batas sumber dan tidak
direkonsiliasi oleh review.

\textbf{Populasi substantif.} Penduduk urban JBR dalam Table 1 berjumlah
17.028,5 ribu pada 1990 dan 24.043,9 ribu pada 2000. Artikel juga
membahas penduduk yang berpindah dari Jakarta ke pinggiran, migran masuk
lima tahun terakhir, serta jutaan komuter antara inti dan daerah
sekitarnya (Urban population, hlm. 331; Table 1, hlm. 331).

\textbf{Unit observasi sekunder.} Unit berbeda menurut bahan: kota dan
kabupaten untuk penduduk, investasi, DAU, dan PDRB; lokalitas urban atau
perdesaan untuk Table 2; hektare kategori guna lahan untuk Table 3;
proyek kota baru untuk Table 4; kawasan industri untuk Table 5; serta
arus penumpang atau kendaraan untuk statistik transportasi (hlm.
331-336).

\textbf{Sampel penelitian.} Tidak ada sampel responden atau penarikan
sampel primer yang dilaporkan. Ukuran sampel, kerangka sampel, teknik
sampling, tingkat respons, kriteria inklusi atau eksklusi, dan
karakteristik responden: Tidak disebutkan dalam file. Inventaris Table 4
dibatasi pada proyek kota baru seluas 1.000 ha atau lebih, tetapi file
tidak menjelaskan apakah daftar tersebut lengkap untuk seluruh proyek
yang memenuhi ambang (Table 4, hlm. 333).

\textbf{Inferensi pengolah.} Populasi analitis artikel adalah wilayah
beserta proses pembangunannya, bukan individu. Karena batas JMA, BMA,
dan JBR tidak diterapkan secara identik pada semua tabel, angka dari
bagian yang berbeda tidak selalu dapat diagregasikan menjadi satu
populasi regional yang konsisten.

\subsection{Dependent Variables}\label{dependent-variables-13}

\textbf{Klarifikasi desain.} Artikel tidak menggunakan istilah variabel
dependen, tidak mengestimasi persamaan regresi, dan tidak menetapkan
satu keluaran formal. Krisis ekonomi, investasi, infrastruktur, dan
desentralisasi dibahas sebagai konteks atau mekanisme yang mungkin
memengaruhi perkembangan, bukan sebagai variabel bebas yang diuji dalam
model (Introduction, hlm. 327-329; Concluding remarks, hlm. 337).

Jika dipadankan secara terbatas dengan keluaran analitis, indikator yang
dilaporkan meliputi:

\begin{itemize}
\tightlist
\item
  jumlah, proporsi, kepadatan, dan laju pertumbuhan penduduk urban;
  pangsa Jakarta dalam JBR dan JMA; jumlah lokalitas urban; migrasi;
  serta arus komuter (Introduction dan Urban population, hlm. 328-332;
  Tables 1-2);
\item
  pertumbuhan ekonomi, PDRB dan PDRB per kapita, investasi yang
  disetujui, transaksi properti, proyek yang beroperasi atau terhenti,
  serta lahan industri menganggur selama krisis (Economic crises, hlm.
  330-331);
\item
  luas dan perubahan kategori guna lahan, khususnya hutan, sawah,
  permukiman, dan industri; luas ruang terbuka hijau; serta luas pusat
  perbelanjaan (Land conversion, hlm. 332-333; Table 3);
\item
  jumlah, luas, lokasi, status operasi, fungsi, dan karakter kota baru
  serta kawasan industri (New town, industrial estate and road
  development, hlm. 333-335; Tables 4-5);
\item
  arus penumpang, pembagian moda, pertumbuhan kendaraan, waktu
  perjalanan, dan perkembangan aktivitas ekonomi setelah pembangunan
  jalan tol (New town, industrial estate and road development, hlm. 333
  dan 335);
\item
  DAU, PDRB, kapasitas fiskal, pemanfaatan sumber daya, koordinasi
  lintas batas, dan penyediaan layanan setelah desentralisasi (Early
  reform and decentralization era, hlm. 335-337; Table 6);
\item
  dampak sosial dan lingkungan yang dijelaskan secara naratif, termasuk
  segregasi, kemacetan, pencemaran udara, ekstraksi air tanah, banjir,
  kelangkaan air, dan hilangnya lahan pertanian (Land conversion, hlm.
  332-333; New town, industrial estate and road development, hlm. 334;
  Concluding remarks, hlm. 337).
\end{itemize}

\textbf{Inferensi pengolah.} Konstruk utama, yaitu kesinambungan dan
perubahan mega-urbanisasi, tidak diberi indeks komposit, ambang, atau
pengukuran yang konsisten antarwaktu. Kesimpulan tentangnya diturunkan
dari konvergensi sejumlah indikator deskriptif dan penilaian
kelembagaan.

\subsection{Results}\label{results-13}

\textbf{Skala, konsentrasi, dan pergeseran penduduk urban.}

\begin{itemize}
\tightlist
\item
  \textbf{Laporan sumber:} penduduk urban JBR mencapai lebih dari 24
  juta pada 2000, sekitar 28\% dari seluruh penduduk urban Indonesia,
  dengan laju pertumbuhan sekitar 3,5\% selama 1990-2000. Pangsa Jakarta
  terhadap penduduk urban JBR turun dari 48\% pada 1990 menjadi 34,7\%
  pada 2000 (Urban population dan Table 1, hlm. 331).
\item
  \textbf{Laporan sumber:} laju pertumbuhan penduduk urban Jakarta turun
  dari 3,1\% pada 1980-1990 menjadi 0,16\% pada 1990-2000; Bandung turun
  dari 3,3\% menjadi 0,40\%. Pertumbuhan Kota Bogor naik dari 1,0\%
  menjadi hampir 11\%, tetapi artikel mengaitkannya dengan perluasan
  batas kota. Penurunan pada Kabupaten Bogor, Tangerang, dan Bekasi juga
  dijelaskan bersama pemekaran Depok serta pemisahan Kota Tangerang dan
  Bekasi, sehingga perubahan administratif memengaruhi perbandingan
  (Urban population dan Table 1, hlm. 331).
\item
  \textbf{Laporan sumber:} indeks primasi JMA relatif terhadap empat
  metropolitan terbesar adalah 0,56 pada 1980, 0,58 pada 1990, dan 0,57
  pada 2000. Pangsa JMA terhadap penduduk urban nasional masing-masing
  22,5\%, 23,6\%, dan 21,2\% (Introduction, hlm. 329).
\item
  \textbf{Laporan sumber:} pada 1995-2000 sekitar 190.000 mantan
  penduduk Jakarta berpindah ke Kota dan Kabupaten Bekasi, 192.000 ke
  Kota dan Kabupaten Tangerang, serta 160.000 ke Kota dan Kabupaten
  Bogor serta Kota Depok. Pangsa Jakarta dalam penduduk JMA turun dari
  54,9\% pada 1971 menjadi 39,9\% pada 2000, sedangkan pangsa Bekasi
  naik dari 10,0\% menjadi 15,3\% dan Tangerang dari 12,8\% menjadi
  19,6\% (Urban population, hlm. 331).
\item
  \textbf{Laporan sumber:} jumlah lokalitas urban JMA naik dari 730 pada
  1999 menjadi 1.035 pada 2005, sementara proporsinya meningkat dari
  sekitar dua perlima menjadi hampir tiga perlima. Di BMA, jumlahnya
  naik dari 266 menjadi 412 dan proporsinya dari sekitar sepertiga
  menjadi hampir separuh (Urban population dan Table 2, hlm. 331-332).
\item
  \textbf{Laporan sumber:} Sensus 2000 mencatat lebih dari 1,35 juta
  migran masuk lima tahun terakhir di JMA; sekitar sepertiga berasal
  dari Jawa Tengah, tiga persepuluh dari Jawa Barat, 15\% dari Sumatra,
  dan 9\% dari Jawa Timur. Migran masuk ke Kota dan Kabupaten Bandung
  masing-masing mencapai 315,8 ribu dan 200 ribu, atau 515,8 ribu untuk
  BMA, dengan sekitar dua pertiga berasal dari Jawa (Urban population,
  hlm. 331).
\item
  \textbf{Laporan sumber:} angka JICA yang dikutip artikel mencatat
  lebih dari tiga juta komuter antara Jakarta dan sekitarnya pada 2002,
  termasuk 1,10 juta antara Tangerang-Jakarta, 1,14 juta Bekasi-Jakarta,
  dan hampir satu juta Bogor-Jakarta; daerah yang lebih jauh belum
  termasuk (Introduction, hlm. 329).
\end{itemize}

\textbf{Dampak krisis ekonomi yang dilaporkan.}

\begin{itemize}
\tightlist
\item
  \textbf{Laporan sumber:} JMA mengalami kontraksi dari pertumbuhan
  ekonomi pada kisaran 6,0-8,3\% per tahun selama 1987-1997 menjadi
  -2,74\% pada 1998-1999. PDRB per kapita Jakarta turun dari sekitar
  Rp7,5 juta menjadi Rp6,1 juta pada 1998, lalu dilaporkan mencapai
  Rp7,1 juta pada 2000 atau 3,7 kali rata-rata nasional Rp1,93 juta
  (Economic crises, hlm. 330).
\item
  \textbf{Laporan sumber:} persetujuan investasi domestik di Jakarta
  turun dari hampir Rp63.700 miliar pada 1997 menjadi sekitar Rp18.900
  miliar pada 1998. Transaksi properti Jakarta turun dari Rp2,6 triliun
  pada 1996 menjadi Rp0,8 triliun pada 1998 (Economic crises, hlm. 330).
\item
  \textbf{Laporan sumber:} hanya 50 dari 500 proyek kota baru dan
  perumahan skala besar JMA yang dilaporkan beroperasi dengan baik pada
  1999. Permintaan lahan kawasan industri turun; artikel menyebut
  penurunan kawasan industri sebesar 30\% pada 1997-1998 dan sekitar
  47.000 ha lahan kawasan industri menganggur pada awal 1998 (Economic
  crises, hlm. 330-331).
\item
  \textbf{Laporan sumber:} proyek kota baru, perumahan, kondominium,
  perkantoran, dan kawasan industri melambat atau berhenti, menambah
  tanah yang telah diakuisisi tetapi tidak dimanfaatkan dan meningkatkan
  pengangguran pada sektor properti terkait (Economic crises, hlm.
  330-331).
\end{itemize}

\textbf{Perubahan guna lahan dan lingkungan binaan.}

\begin{itemize}
\tightlist
\item
  \textbf{Laporan sumber:} antara 1994 dan 2001, hutan primer Bopunjur
  turun dari 12.166,2 ha menjadi 5.449,4 ha dan hutan sekunder dari
  1.205,6 ha menjadi 14,0 ha. Di BMA, hutan primer turun dari 57.294,4
  ha menjadi 55.748,7 ha dan hutan sekunder dari 39.349,3 ha menjadi
  5.541,9 ha. Narasi artikel merangkum kehilangan hutan hampir 8.000 ha
  di Bopunjur dan lebih dari 35.000 ha di BMA (Land conversion dan Table
  3, hlm. 332-333).
\item
  \textbf{Laporan sumber:} sawah Bopunjur berkurang 4.388,6 ha, dari
  30.328,2 ha menjadi 25.939,8 ha, sedangkan sawah BMA berkurang
  12.478,7 ha, dari 65.626,1 ha menjadi 53.147,4 ha. Permukiman
  bertambah 2.529,0 ha di Bopunjur dan 3.110,2 ha di BMA (Table 3, hlm.
  333).
\item
  \textbf{Laporan sumber:} artikel menyebut sekitar sepertiga hingga
  separuh area berizin pengembangan di wilayah metropolitan Jakarta pada
  pertengahan 1990-an tidak aktif dikembangkan dan ditahan dari pasar.
  Sistem izin memberi pemegang izin hak eksklusif untuk mengakuisisi
  lahan melalui negosiasi langsung, sedangkan pemilik lahan sering
  dinilai menerima kompensasi yang tidak adil (Land conversion, hlm.
  332).
\item
  \textbf{Laporan sumber:} luas pusat perbelanjaan Jakarta naik dari 1,4
  juta m2 pada 2000 menjadi 2,4 juta m2 pada 2005. Ruang terbuka hijau
  Jakarta turun dari 28,8\% pada 1984 menjadi 9,4\% pada 2000 dan
  diperkirakan 6,2\% pada 2007 (Land conversion, hlm. 332-333).
\end{itemize}

\textbf{Kota baru, segregasi, dan kawasan industri.}

\begin{itemize}
\tightlist
\item
  \textbf{Laporan sumber:} sekitar 25 proyek kota baru berukuran
  500-6.000 ha terdapat di wilayah metropolitan Jakarta pada awal
  2000-an; hanya sekitar tiga perlima yang masih beroperasi pada
  pertengahan 2004. Table 4 mencatat proyek berluas sekurangnya 1.000
  ha, antara lain Bumi Serpong Damai 6.000 ha, Cikarang Baru 5.400 ha,
  Kota Tiga Raksa, Lippo Cikarang, dan Delta Mas masing-masing 3.000 ha,
  serta Kota Baru Parahyangan yang luasnya tidak tersedia pada tabel
  (New town, industrial estate and road development dan Table 4, hlm.
  333).
\item
  \textbf{Laporan sumber:} Lippo Karawaci memiliki sekitar 30.000
  penghuni pada 2000 dan menyediakan jaringan telepon, televisi kabel,
  listrik, air, rumah sakit, hotel, pusat belanja, sekolah, serta pusat
  budaya dan rekreasi. Kota Baru Parahyangan direncanakan seluas 1.250
  ha untuk pengembangan 20 tahun sejak 2000 dan diarahkan menjadi kota
  baru mandiri (New town, industrial estate and road development, hlm.
  334).
\item
  \textbf{Laporan sumber:} sekitar 3.300 ha sawah di pinggiran Bandung
  dikonversi menjadi permukiman baru pada 1990-an. Pada 2000 terdapat
  lebih dari 20 lapangan golf di JBR; salah satunya Karawang
  International Golf Course 27 lubang yang menempati 165 ha (New town,
  industrial estate and road development, hlm. 334).
\item
  \textbf{Laporan sumber:} artikel menyatakan bahwa banyak kawasan
  permukiman berpendapatan rendah di pusat kota telah disingkirkan dan
  dikonversi menjadi kawasan komersial. Jumlah penduduk, luas lahan,
  lokasi rinci, proses pemindahan, dan kompensasi: Tidak disebutkan
  dalam file (Concluding remarks, hlm. 337).
\item
  \textbf{Laporan sumber:} kawasan industri di wilayah metropolitan
  Jakarta menempati hampir 11.000 ha pada akhir 2005; sekitar empat
  persepuluh berada di Kabupaten Bekasi dan sepertiga di Kabupaten
  Karawang. Tingkat okupansi keseluruhannya mencapai 70\% (New town,
  industrial estate and road development, hlm. 334-335; Table 5).
\item
  \textbf{Laporan sumber:} dari persetujuan investasi asing Indonesia
  2000-2004 sebesar US\$64.803,5 juta, hampir 58\% atau US\$37.112,8
  juta berlokasi di JMA. Dari persetujuan investasi domestik sebesar
  Rp265.176,1 juta, hampir sepertiga atau sekitar Rp82.342 juta berada
  di JMA. BMA tidak memiliki kawasan industri besar; artikel menyebut
  zona industri Cibaduyut dan Gedebage yang lebih kecil (New town,
  industrial estate and road development, hlm. 335).
\end{itemize}

\textbf{Transportasi dan hubungan Jakarta-Bandung.}

\begin{itemize}
\tightlist
\item
  \textbf{Laporan sumber:} koridor Jakarta-Bandung diperkirakan melayani
  lebih dari 78 juta penumpang pada 2004, sekitar 95\% melalui jalan.
  Kereta JMA membawa sekitar 118 juta penumpang dan hampir 3.600 juta
  penumpang-km pada 2002, dengan kepadatan penumpang 58.000 penumpang
  per hari per kilometer atau enam kali angka nasional (New town,
  industrial estate and road development, hlm. 335).
\item
  \textbf{Laporan sumber:} jumlah mobil penumpang, truk, bus, dan sepeda
  motor terdaftar di Jakarta tumbuh sekitar 10,6\% per tahun selama
  1993-2001. Artikel melaporkan kecepatan rata-rata kendaraan turun dari
  38,3 km/jam pada 1995 menjadi 34,5 km/jam pada 2006 (New town,
  industrial estate and road development, hlm. 333 dan 335).
\item
  \textbf{Laporan sumber:} penyelesaian Tol Purwakarta-Padalarang
  mempersingkat perjalanan Jakarta-Bandung dari tiga menjadi dua jam.
  Kemudahan perjalanan akhir pekan dari Jakarta mendorong hotel,
  restoran, toko makanan, dan gerai garmen di Bandung; PAD pariwisata
  Bandung mencapai sekitar Rp234 miliar pada 2006, hampir 40\% dari PAD
  total, dengan pertumbuhan 9\% per tahun sejak 2004. Artikel juga
  melaporkan kemacetan akhir pekan yang memburuk (New town, industrial
  estate and road development, hlm. 335).
\end{itemize}

\textbf{Desentralisasi fiskal dan tata kelola regional.}

\begin{itemize}
\tightlist
\item
  \textbf{Laporan sumber:} reformasi mengalihkan tanggung jawab
  pembangunan urban dan regional kepada pemerintah lokal serta memberi
  DAU yang dapat digunakan menurut kebutuhan lokal. Table 6 menunjukkan
  DAU Kabupaten Bandung naik dari Rp654,8 miliar pada 2001 menjadi
  Rp726,2 miliar pada 2003 dan Rp1.168,6 miliar pada 2006; nilai 2006
  yang ditampilkan juga meningkat dibandingkan 2001 untuk sebagian besar
  pemerintah lokal lain (Early reform and decentralization era dan Table
  6, hlm. 336).
\item
  \textbf{Laporan sumber:} kewenangan pemerintah lokal untuk memungut
  pajak tetap terbatas karena kendali utama berada pada pemerintah
  pusat. Artikel melaporkan kecenderungan pemerintah lokal memaksimalkan
  PAD melalui pemanfaatan tanah, air, dan aset fisik serta membangun
  infrastruktur untuk bersaing tanpa selalu mempertimbangkan efisiensi
  regional (Early reform and decentralization era, hlm. 336).
\item
  \textbf{Laporan sumber:} perselisihan lintas batas terjadi pada
  penggunaan lokasi pembuangan antara Kabupaten Bekasi dan Jakarta serta
  pengelolaan sampah antara Kota dan Kabupaten Bandung. Pembangunan vila
  dan permukiman tetap diizinkan di Puncak serta Bandung utara meskipun
  keduanya berfungsi sebagai kawasan resapan (Early reform and
  decentralization era, hlm. 336).
\item
  \textbf{Laporan sumber:} gagasan pengelolaan terpadu muncul melalui
  usulan ``Megapolitan Jabodetabekjur'', tetapi ditentang sebagian pihak
  di wilayah tetangga karena dicurigai sebagai ambisi politik untuk
  menganeksasi wilayah ke Jakarta (Early reform and decentralization
  era, hlm. 336-337).
\end{itemize}

\subsection{Findings}\label{findings-13}

\begin{itemize}
\tightlist
\item
  \textbf{Interpretasi penulis:} indikator pertumbuhan pinggiran,
  perpindahan penduduk keluar dari inti, kenaikan lokalitas urban, dan
  arus komuter menunjukkan limpahan cepat Jakarta dan Bandung serta
  transformasi JMA dan BMA dari struktur berinti tunggal menjadi berinti
  majemuk (Urban population, hlm. 331; Concluding remarks, hlm. 337).
\item
  \textbf{Interpretasi penulis:} krisis ekonomi merupakan interupsi kuat
  tetapi sementara. Pembangunan fisik dan ekonomi merosot pada 1998
  hingga awal 2000-an, kemudian JBR kembali berkembang cepat pada
  pertengahan 2000-an seperti sebelum krisis (Economic crises, hlm.
  330-331; Concluding remarks, hlm. 337).
\item
  \textbf{Interpretasi penulis:} transformasi regional digerakkan kuat
  oleh sektor swasta kapitalistik pada manufaktur, perdagangan, jasa,
  dan properti. Investasi serta pencarian rente menjelaskan perluasan
  kawasan industri, kota baru, spekulasi tanah, dan konversi lahan;
  penulis menegaskan bahwa pengembang kota baru tidak bertujuan
  mendesentralisasikan kegiatan sosial ekonomi dari Jakarta (Land
  conversion, hlm. 332; New town, industrial estate and road
  development, hlm. 333-335; Concluding remarks, hlm. 337).
\item
  \textbf{Interpretasi penulis:} sebagian besar kota baru berfungsi
  sebagai kota asrama karena tidak menyediakan tempat kerja yang memadai
  dan tetap bergantung pada Jakarta atau Bandung. Keterpisahan antarkota
  baru dan lemahnya kaitan dengan sistem infrastruktur regional
  memperbanyak perjalanan komuter, kemacetan, dan pencemaran udara (New
  town, industrial estate and road development, hlm. 333-334; Concluding
  remarks, hlm. 337).
\item
  \textbf{Interpretasi penulis:} kota baru memperkuat segregasi dalam
  tiga bentuk: kantong hunian eksklusif bagi kelompok menengah-atas;
  pemisahan internal kawasan dengan tingkat keamanan berbeda; dan
  pengelolaan fasilitas oleh pengembang swasta yang membatasi pengguna
  dari luar. Penulis menyebutnya segregasi spasial sukarela atau
  segregasi diri berdasarkan gaya hidup dan keamanan (New town,
  industrial estate and road development, hlm. 333-334).
\item
  \textbf{Interpretasi penulis:} pelanggaran rencana tata ruang,
  lemahnya penegakan izin, tekanan investor, dan kepentingan
  politik-ekonomi membuat konversi lahan tidak terkendali. Konversi
  kawasan konservasi dan resapan dipandang berisiko menimbulkan banjir,
  kelangkaan air, hilangnya investasi irigasi, pencemaran, dan kerusakan
  hilir (Land conversion, hlm. 332-333; Early reform and
  decentralization era, hlm. 336; Concluding remarks, hlm. 337).
\item
  \textbf{Interpretasi penulis:} infrastruktur jalan memperkuat
  integrasi fungsional Jakarta-Bandung dengan mempersingkat waktu
  perjalanan dan mendorong pariwisata serta kegiatan jasa di Bandung,
  tetapi ketergantungan pada jalan dan perjalanan akhir pekan juga
  memperburuk kemacetan. Artikel menilai sistem kereta perlu
  dikembangkan untuk mengatasi ketergantungan berlebihan pada jalan (New
  town, industrial estate and road development, hlm. 333 dan 335).
\item
  \textbf{Interpretasi penulis:} desentralisasi menambah kewenangan dan
  sumber fiskal lokal, tetapi pelaksanaannya tidak merata. Kapasitas
  institusional, finansial, dan teknis banyak pemerintah lokal tetap
  lemah; korupsi terdesentralisasi, parokialisme, dan dorongan
  mengeksploitasi sumber lokal menghambat pelayanan serta kerja sama
  lintas batas (Early reform and decentralization era, hlm. 335-337).
\item
  \textbf{Interpretasi penulis:} JBR merupakan wilayah fungsional yang
  tidak dapat dikelola sebagai kumpulan yurisdiksi terpisah. Rencana
  tata ruang terpadu, kepemimpinan lintas pemerintah, koordinasi
  pemangku kepentingan, dan pembangunan institusi regional dinilai
  penting bagi keberlanjutannya (Abstract, hlm. 327; Early reform and
  decentralization era, hlm. 336-337).
\item
  \textbf{Inferensi pengolah:} bukti deskriptif paling kuat menunjukkan
  pertumbuhan pinggiran, perubahan guna lahan, konsentrasi investasi,
  serta masalah koordinasi yang berjalan bersamaan. Bukti tersebut
  konsisten dengan argumentasi mega-urbanisasi, tetapi tidak menunjukkan
  berapa bagian perubahan yang disebabkan secara terpisah oleh
  globalisasi, pasar properti, jalan tol, krisis, atau desentralisasi.
\item
  \textbf{Inferensi pengolah:} klaim bahwa mega-urbanisasi setelah
  krisis berlangsung lebih cepat daripada 1980-an hingga pertengahan
  1990-an tidak disertai satu ukuran regional yang dihitung secara
  konsisten untuk kedua periode. Artikel memberikan banyak indikator
  pendukung dari tahun berbeda, tetapi tidak menyajikan uji langsung
  atas perbandingan laju keseluruhan (Abstract, hlm. 327; Concluding
  remarks, hlm. 337).
\item
  \textbf{Inferensi pengolah:} temuan memperlihatkan ketegangan utama
  antara integrasi fungsional dan fragmentasi kewenangan. Mobilitas,
  pasar lahan, air, sampah, dan dampak lingkungan melampaui batas
  administratif, sementara insentif fiskal serta kewenangan pembangunan
  tetap terlokalisasi (Early reform and decentralization era, hlm.
  336-337).
\end{itemize}

\subsection{Conclusions}\label{conclusions-13}

\textbf{Kesimpulan penulis.} JBR berperan penting dalam ekonomi nasional
dan telah membentuk sabuk urban sekitar 200 km dengan campuran kegiatan
desa-kota. Pertumbuhan penduduk di Jakarta dan Bandung melambat,
sedangkan pertumbuhan di daerah yang berdekatan dan jumlah lokalitas
urban meningkat. Penulis menyimpulkan bahwa pola tersebut mencerminkan
limpahan dari inti dan transformasi JMA serta BMA menuju struktur
berinti majemuk (Concluding remarks, hlm. 337).

\textbf{Kesimpulan penulis.} Krisis ekonomi akhir 1990-an hingga awal
2000-an memperlambat pembangunan karena permintaan lahan kota baru dan
kawasan industri turun, tetapi pembangunan kembali cepat pada
pertengahan 2000-an. Modal swasta pada manufaktur, perdagangan, jasa,
dan properti dipandang sebagai penggerak utama pemulihan serta
transformasi fisik regional (Concluding remarks, hlm. 337).

\textbf{Kesimpulan penulis.} Perkembangan tiga dekade menghasilkan
kemacetan akibat komuter, pencemaran udara di kawasan industri,
ekstraksi air tanah berlebihan, serta konversi lahan yang tidak
terkendali akibat pelanggaran rencana guna lahan oleh pemerintah lokal
dan sektor swasta. Kota baru memperkuat segregasi gaya hidup dan
pendapatan serta tetap bergantung secara sosial ekonomi pada Jakarta
atau Bandung, sedangkan kawasan industri terus berkembang karena
investasi dan akses ke Jakarta (Concluding remarks, hlm. 337).

\textbf{Kesimpulan penulis.} Jalan tol baru memperkuat hubungan kedua
kota dan mendorong pariwisata serta jasa di Bandung, tetapi juga
memperburuk kemacetan. Otonomi daerah dan desentralisasi fiskal
menjadikan pembangunan urusan pemerintah lokal, namun kecenderungan
eksploitasi sumber daya dan parokialisme menghambat pelayanan lintas
batas. Akibatnya, proses mega-urbanisasi terus berlangsung, bahkan
dinilai lebih cepat daripada 1980-an hingga pertengahan 1990-an,
sementara pengelolaan wilayah makin terfragmentasi (Concluding remarks,
hlm. 337).

\textbf{Kesimpulan penulis.} JBR perlu diperlakukan sebagai wilayah
fungsional yang tidak terpisahkan melalui rencana pembangunan spasial
terpadu. Kepemimpinan pemerintah lokal, provinsi, dan pusat, kerja sama
pemangku kepentingan, serta pembangunan institusi regional ditempatkan
sebagai syarat penting keberlanjutan (Abstract, hlm. 327; Early reform
and decentralization era, hlm. 336-337).

\textbf{Inferensi pengolah.} Kesimpulan paling terdukung adalah adanya
perubahan demografis pinggiran, konversi lahan, konsentrasi investasi
dan proyek, serta persoalan koordinasi lintas batas. Kesimpulan mengenai
kecepatan keseluruhan mega-urbanisasi dan dampak spesifik desentralisasi
lebih bersifat interpretatif karena artikel tidak memakai ukuran
longitudinal terpadu atau strategi atribusi kausal.

\subsection{Contributions}\label{contributions-13}

\begin{itemize}
\tightlist
\item
  \textbf{Kontribusi yang dinyatakan penulis:} artikel merupakan salah
  satu dari sedikit kajian yang memeriksa perkembangan mutakhir JBR dan
  secara khusus berusaha menjelaskan kesinambungan serta perubahan
  mega-urbanisasi setelah krisis dan pada awal reformasi (Introduction,
  hlm. 327-328; Concluding remarks, hlm. 337).
\item
  \textbf{Kontribusi konseptual menurut penulis:} kajian menggeser fokus
  dari pertanyaan apakah mega-urbanisasi Asia sama dengan urbanisasi
  Barat menuju pembacaan kontekstual atas kondisi sosial ekonomi,
  politik, fisik, aktor, dan institusi lokal-global (Introduction, hlm.
  328; Mega-urbanization in Asia, hlm. 329-330).
\item
  \textbf{Kontribusi empiris yang diperagakan sumber:} artikel
  memperbarui gambaran integrasi Jakarta-Bandung dari kajian awal
  1990-an dengan menghimpun indikator penduduk, migrasi, lokalitas
  urban, guna lahan, investasi, properti, kawasan industri,
  transportasi, dan fiskal hingga pertengahan 2000-an (hlm. 327-337;
  Tables 1-6).
\item
  \textbf{Kontribusi analitis menurut penulis:} krisis ekonomi dan
  desentralisasi dimasukkan ke dalam penjelasan perubahan regional
  sehingga mega-urbanisasi tidak diperlakukan sebagai proses fisik atau
  demografis yang terlepas dari pergantian ekonomi-politik
  (Introduction, hlm. 327-328; Economic crises, hlm. 330-331; Early
  reform and decentralization era, hlm. 335-337).
\item
  \textbf{Kontribusi kebijakan menurut penulis:} artikel
  mengartikulasikan ketidaksesuaian antara wilayah yang terintegrasi
  secara fungsional dan pemerintahan yang terfragmentasi, lalu
  menempatkan rencana spasial terpadu serta pembangunan institusi
  sebagai kebutuhan regional (Abstract, hlm. 327; Early reform and
  decentralization era, hlm. 336-337).
\item
  \textbf{Inferensi pengolah:} nilai tambah artikel terletak pada
  sintesis lintassektor dan lintasperiode, bukan pada metode statistik
  baru. Artikel tidak menawarkan model, indeks, protokol data, atau
  desain evaluasi yang dapat direplikasi secara langsung.
\end{itemize}

\subsection{Research Gap}\label{research-gap-13}

\textbf{Kesenjangan yang dinyatakan atau dibangun oleh sumber.}

\begin{itemize}
\tightlist
\item
  Studi Jabotabek dan BMA pada awal 1990-an telah mengidentifikasi
  pembentukan sabuk urban, tetapi krisis ekonomi, reformasi politik,
  otonomi daerah, dan desentralisasi fiskal mengubah konteks setelahnya.
  Artikel diposisikan untuk memeriksa kesinambungan dan perubahan dalam
  kondisi baru tersebut (Introduction, hlm. 327-328).
\item
  Penulis menilai perdebatan yang hanya menanyakan kemiripan
  mega-urbanisasi Asia dengan pengalaman Barat kurang memadai.
  Kesenjangan perspektifnya ialah kebutuhan akan pemeriksaan yang
  membumi atas interaksi aktor dan institusi global-lokal dalam setting
  tertentu (Introduction, hlm. 328; Mega-urbanization in Asia, hlm.
  330).
\item
  Dampak reformasi desentralisasi terhadap JBR belum dapat dinilai
  secara matang karena baru berjalan tujuh tahun dan meliputi dimensi
  sosial ekonomi, budaya, serta politik yang luas. Artikel mengisi
  sebagian celah dengan membahas potensi dampak, bukan evaluasi dampak
  definitif (Introduction, hlm. 328).
\item
  Artikel menyatakan bahwa banyak kajian telah membahas pembangunan
  urban di dunia berkembang, tetapi masih sedikit pengetahuan tentang
  pengaruhnya terhadap konversi lahan dari pertanian ke penggunaan urban
  dan dari satu penggunaan urban ke penggunaan lain (Land conversion,
  hlm. 331-332).
\item
  Pada penutup, penulis menyebut kajiannya sebagai salah satu dari
  sedikit studi tentang perkembangan JBR mutakhir, sekaligus mengakui
  sifatnya yang sangat deskriptif (Concluding remarks, hlm. 337).
\end{itemize}

\textbf{Kesenjangan yang masih tersisa sebagai inferensi pengolah.}

\begin{itemize}
\tightlist
\item
  Artikel belum mengisolasi dampak desentralisasi dari pemulihan
  pascakrisis, investasi, pemekaran administratif, pertumbuhan penduduk,
  dan pembangunan infrastruktur.
\item
  Tidak tersedia sistem indikator dengan batas wilayah dan tahun yang
  konsisten untuk menguji klaim perubahan laju mega-urbanisasi secara
  langsung.
\item
  Dampak distribusional terhadap pemilik tanah, penduduk kampung,
  pekerja berpendapatan rendah, dan komunitas yang dipindahkan tidak
  diukur, meskipun segregasi dan perlakuan tidak adil dalam transaksi
  tanah dibahas.
\item
  Dampak lingkungan disebut melalui perubahan tutupan lahan dan contoh
  banjir, air, polusi, serta kemacetan, tetapi hubungan mekanisme dan
  besar dampaknya tidak diuji secara khusus.
\item
  Efektivitas alternatif kelembagaan, rencana terpadu, dan instrumen
  penegakan tidak dibandingkan melalui evaluasi implementasi.
\end{itemize}

\subsection{Limitations}\label{limitations-13}

\textbf{Keterbatasan yang dinyatakan penulis.}

\begin{itemize}
\tightlist
\item
  Evaluasi dampak desentralisasi diakui prematur karena periode
  2001-2008 hanya tujuh tahun dan reformasinya masih berlangsung
  (Introduction, hlm. 328).
\item
  Identifikasi dampak reformasi dinilai sulit karena desentralisasi
  memiliki dimensi sosial ekonomi, budaya, dan politik yang luas.
  Artikel karena itu membatasi diri pada pembahasan potensi dampak
  (Introduction, hlm. 328).
\item
  Penulis secara eksplisit mengakui bahwa kajian ini sangat deskriptif
  (\emph{largely descriptive in character}) (Concluding remarks, hlm.
  337).
\end{itemize}

\textbf{Keterbatasan yang diinferensikan pengolah dari desain dan
pelaporan.}

\begin{itemize}
\tightlist
\item
  Tidak ada bagian metode, protokol pemilihan sumber, atau prosedur
  validasi data. Kelengkapan korpus dan kemungkinan bias pemilihan bukti
  tidak dapat dinilai dari file.
\item
  Data berasal dari sensus, instansi pemerintah, konsultan properti,
  majalah, surat kabar, makalah tidak terbit, dan studi terdahulu.
  Artikel tidak melaporkan penilaian mutu atau rekonsiliasi perbedaan
  definisi di antara sumber-sumber tersebut.
\item
  Indikator menggunakan tahun yang berbeda, antara lain 1980, 1990,
  1994, 1997, 1999, 2000, 2001, 2002, 2003, 2004, 2005, dan 2006, dengan
  beberapa estimasi untuk 2007. File tidak menjelaskan harmonisasi tahun
  dasar atau pembentukan panel yang sebanding.
\item
  Perubahan batas administratif memengaruhi laju pertumbuhan penduduk
  Bogor, Bekasi, Tangerang, dan Depok. Penulis mengakui pengaruh ini
  secara naratif, tetapi tidak menyajikan seri yang diselaraskan pada
  batas tetap (Urban population, hlm. 331).
\item
  Batas wilayah tidak konsisten sepenuhnya. Narasi JMA dan BMA, Table 1,
  serta Table 2 mencakup kombinasi unit yang berbeda; konsekuensinya
  bagi perbandingan tidak dibahas.
\item
  Artikel tidak menyajikan model kausal, kelompok pembanding, uji
  statistik, interval ketidakpastian, atau analisis sensitivitas.
  Hubungan antara krisis, investasi, jalan tol, desentralisasi, dan
  hasil regional tidak dapat dipisahkan secara kuantitatif.
\item
  Konsep mega-urbanisasi, perubahan satu inti ke banyak inti, segregasi,
  dan fragmentasi tidak dioperasionalkan sebagai ukuran yang diuji.
  Beberapa kesimpulan bertumpu pada pembacaan pola dan rujukan
  terdahulu.
\item
  Bukti lingkungan tidak seragam. Table 3 mengukur perubahan luas lahan,
  sedangkan klaim tentang banjir, kelangkaan air, pencemaran, dan
  ekstraksi air tanah sebagian besar dilaporkan secara naratif tanpa
  desain atribusi atau ukuran dampak yang konsisten.
\item
  Kesejahteraan individu dan distribusi biaya-manfaat tidak dianalisis.
  Angka regional tidak dengan sendirinya menunjukkan pengalaman seluruh
  kelompok sosial di dalam wilayah yang sangat heterogen.
\item
  File tidak menyediakan data mentah, metadata variabel, kode, atau
  prosedur perhitungan yang cukup untuk mereplikasi Tables 1-6.
\end{itemize}

\subsection{Challenges}\label{challenges-13}

\textbf{Tantangan penelitian yang diakui atau terlihat dalam sumber.}

\begin{itemize}
\tightlist
\item
  \textbf{Laporan sumber:} memisahkan dampak desentralisasi pada tahap
  awal merupakan tugas sulit karena reformasi masih berlangsung dan
  beroperasi melalui banyak dimensi (Introduction, hlm. 328).
\item
  \textbf{Inferensi pengolah:} tantangan empiris utama ialah menyatukan
  statistik lintassumber, lintasskala, lintasbatas wilayah, dan
  lintastahun tanpa protokol harmonisasi yang dilaporkan.
\item
  \textbf{Inferensi pengolah:} perubahan administratif dapat menyerupai
  pertumbuhan atau penurunan urban yang substantif, sehingga indikator
  demografis harus dibaca bersama riwayat pemekaran (Urban population,
  hlm. 331).
\end{itemize}

\textbf{Tantangan pembangunan regional yang dilaporkan sumber.}

\begin{itemize}
\tightlist
\item
  Lemahnya penegakan rencana guna lahan dan sistem izin, kuatnya tekanan
  investor, spekulasi lahan, serta kapasitas pemerintah lokal yang tidak
  memadai menyulitkan pengendalian konversi lahan (Land conversion, hlm.
  332-333).
\item
  Kota baru tersebar, kurang terhubung satu sama lain dan dengan
  infrastruktur regional, serta tidak menyediakan cukup pekerjaan.
  Ketergantungan pada inti memperbesar perjalanan komuter, kemacetan,
  dan polusi (New town, industrial estate and road development, hlm.
  333-334).
\item
  Ketergantungan tinggi pada moda jalan membebani koridor
  Jakarta-Bandung. Pembangunan jalan tol mempercepat perjalanan, tetapi
  juga mengalihkan permintaan ke jalan dan memperburuk kemacetan akhir
  pekan (New town, industrial estate and road development, hlm. 333 dan
  335).
\item
  Pemerintah lokal memiliki kapasitas institusional, finansial, dan
  teknis yang tidak merata, sementara kewenangan pajaknya terbatas.
  Dorongan memaksimalkan PAD dapat mengalahkan pertimbangan efisiensi
  dan keberlanjutan regional (Early reform and decentralization era,
  hlm. 335-336).
\item
  Korupsi terdesentralisasi, politik uang, parokialisme, dan anggapan
  wilayah sebagai sumber daya eksklusif menghambat tujuan partisipasi,
  akuntabilitas, pelayanan, dan pemerataan (Early reform and
  decentralization era, hlm. 335-336).
\item
  Air, kawasan resapan, sampah, transportasi, polusi, dan banjir
  merupakan persoalan lintas batas, tetapi JBR terbagi ke banyak kota,
  kabupaten, dan provinsi dengan kewenangan masing-masing (Early reform
  and decentralization era, hlm. 336-337).
\item
  Perlindungan kawasan konservasi Bopunjur dan Bandung utara berhadapan
  dengan permintaan perumahan, vila, pariwisata, serta pendapatan izin
  pembangunan (Land conversion, hlm. 332-333; Early reform and
  decentralization era, hlm. 336).
\end{itemize}

\subsection{Practical Implications}\label{practical-implications-13}

\begin{itemize}
\tightlist
\item
  \textbf{Implikasi penulis:} JBR perlu dikelola sebagai satu wilayah
  fungsional melalui rencana tata ruang dan pembangunan yang terpadu,
  meskipun terdiri atas pemerintah lokal dengan kewenangan masing-masing
  (Abstract, hlm. 327; Early reform and decentralization era, hlm.
  336-337).
\item
  \textbf{Implikasi penulis:} pembangunan institusi regional,
  kepemimpinan pemerintah Jakarta, Jawa Barat, Banten, pemerintah lokal,
  dan pemerintah pusat, serta koordinasi seluruh pemangku kepentingan
  diperlukan untuk membuat kerja sama lintas batas berjalan (Early
  reform and decentralization era, hlm. 336-337).
\item
  \textbf{Implikasi penulis:} negara dan pemerintah lokal perlu
  merumuskan kebijakan, strategi, dan institusi pembangunan urban yang
  sesuai dengan kondisi setempat, bukan sekadar mengimpor model
  urbanisasi dari setting lain (Mega-urbanization in Asia, hlm. 330;
  Early reform and decentralization era, hlm. 336).
\item
  \textbf{Inferensi pengolah berdasarkan masalah yang dilaporkan
  sumber:} penegakan rencana guna lahan dan izin perlu diperkuat,
  terutama pada sawah beririgasi, Bopunjur, dan Bandung utara. Pemberian
  izin serta manfaat PAD perlu mempertimbangkan banjir, resapan air,
  kelangkaan air, hilangnya investasi irigasi, dan dampak hilir (Land
  conversion, hlm. 332-333; Early reform and decentralization era, hlm.
  336).
\item
  \textbf{Implikasi penulis:} pengelolaan sampah dan penyediaan air
  memerlukan mekanisme kerja sama lintas batas agar sengketa biaya dan
  lokasi tidak mengganggu pelayanan metropolitan (Early reform and
  decentralization era, hlm. 336).
\item
  \textbf{Implikasi penulis:} ketergantungan berlebihan pada jalan perlu
  dikurangi melalui pengembangan sistem kereta yang komprehensif.
  Artikel juga melaporkan rencana perluasan busway ke Bogor, Depok,
  Tangerang, dan Bekasi serta instrumen tarif parkir dan
  \emph{electronic road pricing} sebagai upaya mengurangi kendaraan
  pribadi dan polusi (New town, industrial estate and road development,
  hlm. 333 dan 335).
\item
  \textbf{Inferensi pengolah:} persetujuan kota baru sebaiknya menilai
  penyediaan pekerjaan, koneksi antarkawasan, akses fasilitas bagi
  komunitas sekitar, dan tambahan perjalanan, bukan hanya luas investasi
  dan penjualan properti. Inferensi ini diturunkan dari fungsi kota
  asrama, segregasi, serta buruknya keterhubungan yang dilaporkan
  artikel (hlm. 333-334).
\item
  \textbf{Inferensi pengolah:} pemantauan regional memerlukan batas
  geografis tetap, tahun yang sebanding, dan indikator bersama untuk
  penduduk, guna lahan, transportasi, investasi, lingkungan, serta
  fiskal agar klaim percepatan dan dampak kebijakan dapat diuji, bukan
  hanya dideskripsikan.
\end{itemize}

\subsection{Applications}\label{applications-13}

\begin{itemize}
\tightlist
\item
  \textbf{Aplikasi yang diperagakan sumber:} kerangka
  kesinambungan-perubahan digunakan untuk membaca perkembangan JBR
  melalui periode pertumbuhan sebelum krisis, kontraksi 1998-awal
  2000-an, pemulihan pertengahan 2000-an, dan awal desentralisasi (hlm.
  327-337).
\item
  \textbf{Aplikasi yang diperagakan sumber:} statistik penduduk,
  lokalitas urban, konversi lahan, kota baru, kawasan industri,
  investasi, transportasi, DAU, dan PDRB dipakai bersama sebagai
  diagnosis regional multidimensi (Tables 1-6, hlm. 331-336).
\item
  \textbf{Aplikasi kebijakan menurut penulis:} hasil kajian dipakai
  untuk mendukung pengelolaan terpadu JBR, perlindungan kawasan resapan,
  koordinasi air dan sampah, serta pembangunan institusi metropolitan
  (Early reform and decentralization era, hlm. 336-337).
\item
  \textbf{Aplikasi konseptual menurut penulis:} kasus JBR menunjukkan
  cara menempatkan mega-urbanisasi dalam hubungan lokal-global,
  desa-kota, kondisi fisik, dan institusi setempat, alih-alih memakai
  kemiripan dengan kota Barat sebagai tolok ukur tunggal (Introduction,
  hlm. 328; Mega-urbanization in Asia, hlm. 329-330).
\item
  \textbf{Inferensi pengolah:} susunan indikator artikel dapat menjadi
  daftar awal untuk audit perkembangan metropolitan lain, tetapi bukan
  perangkat ukur siap pakai. Transfer memerlukan definisi wilayah yang
  konsisten, sumber data yang tervalidasi, dan metode analitis yang
  lebih eksplisit; validitas di luar JBR tidak diuji dalam file.
\end{itemize}

\subsection{Future Research
(Optional)}\label{future-research-optional-13}

\textbf{Agenda yang disebutkan sumber.} Penulis menyatakan perlunya
kajian komparatif pembangunan urban pada setting yang berbeda agar
proses mega-urbanisasi dipahami dalam konteks masyarakat dan tempatnya,
bukan hanya dinilai sama atau berbeda dari pengalaman Barat
(Mega-urbanization in Asia, hlm. 330).

\textbf{Inferensi pengolah, bukan usulan eksplisit penulis.} Pengakuan
bahwa tujuh tahun terlalu singkat untuk menilai desentralisasi
menunjukkan kebutuhan evaluasi longitudinal setelah reformasi lebih
matang, dengan pemisahan pengaruh desentralisasi dari pemulihan ekonomi,
investasi, infrastruktur, dan perubahan administratif (Introduction,
hlm. 328). Artikel tidak menyatakan desain atau jadwal penelitian
lanjutan untuk tujuan tersebut.

\textbf{Inferensi pengolah, bukan usulan eksplisit penulis.} Kesenjangan
pelaporan menunjukkan kegunaan penelitian berikutnya yang memakai batas
regional tetap, seri data yang selaras, ukuran perubahan lingkungan dan
distribusi sosial, serta evaluasi alternatif kelembagaan. Agenda rinci
ini tidak dinyatakan oleh penulis dan tidak diperlakukan sebagai klaim
sumber.

Usulan penelitian masa depan lain yang dinyatakan secara eksplisit:
Tidak disebutkan dalam file.

\subsection{Provenance Notes}\label{provenance-notes-13}

\begin{itemize}
\tightlist
\item
  Review ini hanya menggunakan seluruh isi file
  \texttt{/Users/mac/Documents/Mac/{[}2{]}\ Obsidian\ Vault/zettelkasten/source/journal/J25-firman-2009-the-continuity-and-change-in-mega-urbanization-in-indonesia-a-survey-of-jakarta-bandung-region-(JBR)-development-10.1016:j.habitatint.2008.08.005.txt},
  sebanyak 942 baris. Tidak ada penelusuran web, sumber eksternal, atau
  dokumen lain yang digunakan untuk menambah, mengoreksi, atau
  memverifikasi klaim substantif artikel.
\item
  TXT memuat artikel dari judul dan abstrak sampai
  \emph{Acknowledgement} serta daftar referensi, dengan nomor halaman
  jurnal 327-339. File tidak memuat blok metadata ekstraksi yang
  menjelaskan perangkat, tanggal ekstraksi, nama PDF asal, atau jumlah
  halaman PDF. Review ini membaca TXT dan tidak memeriksa PDF secara
  visual.
\item
  Locator \texttt{hlm.} mengikuti nomor halaman pada header
  \emph{Habitat International} yang terlihat dalam TXT. Nama bagian,
  Table 1-6, dan Fig. 1 dipertahankan sesuai artikel. Klaim dari sumber
  lain di dalam artikel tetap diperlakukan sebagai laporan J25 atas
  sumber sekunder, bukan sebagai hasil verifikasi independen.
\item
  Tata letak dua kolom membuat urutan teks pada beberapa halaman
  berselang antara kolom kiri dan kanan. Halaman 337 terutama
  menggabungkan bagian akhir \emph{Concluding remarks} di kolom kiri
  dengan kelanjutan kesimpulan, \emph{Acknowledgement}, dan awal
  \emph{References} di kolom kanan. Review mengikuti heading serta
  kesinambungan kalimat, tetapi tidak dapat memverifikasi tata letak
  terhadap PDF.
\item
  Fig. 1 hanya tersedia sebagai caption dan posisi referensial dalam
  TXT; isi visual peta tidak dapat diperiksa. Review tidak mengambil
  bukti spasial yang hanya mungkin terlihat pada gambar.
\item
  Tables 1-6 tersedia sebagai teks hasil ekstraksi, tetapi beberapa
  baris dan kolom mengalami pergeseran. Review menggunakan angka hanya
  ketika pasangan label, tahun, dan nilainya dapat dibaca dengan cukup
  jelas serta tidak merekonstruksi sel yang hilang.
\item
  Pada Table 2, susunan tanda pisah dan nilai pada beberapa sel
  lokalitas urban-perdesaan tidak dapat dipetakan secara aman ke semua
  kolom. Pada Table 3, dasar kolom persentase Bopunjur tidak tampak
  konsisten dengan kolom BMA dan tanda kurung pada perubahan \emph{mixed
  garden} Bopunjur tidak selaras dengan kenaikan luas yang tercetak.
  Review karena itu mengutamakan luas hektare dan ringkasan naratif
  penulis, bukan merekonstruksi persentase yang ambigu (hlm. 332-333).
\item
  Cakupan geografis tidak sepenuhnya seragam. Narasi menyebut JMA atau
  Jabodetabekjur dan BMA dengan unit tertentu; Table 1 menambahkan
  Karawang dalam 12 unit JBR; Table 2 menggabungkan Bandung dan Sumedang
  dalam Bandung Raya. Table 4 juga memberi judul ``Jakarta Metropolitan
  Region (JBR)'', padahal JBR digunakan di bagian lain untuk
  Jakarta-Bandung Region. Review tidak menyatukan perbedaan tersebut
  melalui asumsi.
\item
  Narasi menyatakan JMA terdiri atas 12 unit administratif, tetapi
  daftar yang langsung menyusul menyebut DKI Jakarta serta delapan kota
  atau kabupaten, sedangkan BMA kemudian disebut terdiri atas Kota dan
  Kabupaten Bandung (Introduction, hlm. 328-329). Ketidakselarasan
  hitungan ini dipertahankan dan tidak diperbaiki dari luar file.
\item
  Table 6 tidak menunjukkan kenaikan DAU yang sepenuhnya monoton untuk
  setiap pemerintah pada semua titik, meskipun narasi menyatakan semua
  pemerintah lokal JBR menerima jumlah DAU yang meningkat sejak 2001.
  Contohnya, DAU Jakarta tercetak Rp773,0 miliar pada 2001 dan Rp768,1
  miliar pada 2006, sedangkan Kabupaten Bekasi turun dari Rp235,4 miliar
  pada 2001 menjadi Rp231,2 miliar pada 2003 sebelum naik menjadi
  Rp284,9 miliar pada 2006 (Early reform and decentralization era dan
  Table 6, hlm. 336). Review melaporkan tabel dan interpretasi penulis
  secara terpisah.
\item
  Nilai total PDRB Jakarta pada uraian krisis terdistorsi sebagai
  \texttt{Rp.\ 52\ 283\ 1000\ million} untuk 1998. Karena satuannya
  tidak dapat dipastikan dari TXT, review tidak merekonstruksi atau
  menggunakan nilai total tersebut; nilai PDRB per kapita yang tercetak
  lebih jelas dilaporkan apa adanya (Economic crises, hlm. 330).
\item
  Artikel mencetak kepadatan Jakarta sebesar 45,6 orang per km2 dan
  Cianjur 0,59 orang per km2, serta pada uraian kendaraan mencetak
  \texttt{3.8\ motorcycles} di antara 2,5 juta mobil dan 255.000
  kendaraan umum. Kemungkinan skala atau satuan yang hilang tidak dapat
  dipastikan hanya dari TXT, sehingga angka-angka ambigu tersebut tidak
  digunakan sebagai bukti kuantitatif dalam review (Introduction, hlm.
  329; New town, industrial estate and road development, hlm. 333).
\item
  Sejumlah klaim lingkungan memakai bahasa tentatif atau atribusi tidak
  langsung, misalnya pembangunan perumahan Bopunjur ``dicurigai''
  menyebabkan banjir Jakarta 2002. Review tidak mengubah pernyataan
  tersebut menjadi hubungan kausal yang terbukti (New town, industrial
  estate and road development, hlm. 334).
\item
  Tahun jurnal 2009, rentang halaman 327-339, serta DOI berasal dari
  header dan halaman pembuka TXT; copyright tercetak 2008. Perbedaan
  tahun publikasi dan copyright dipertahankan. Nomor terbitan, riwayat
  naskah, tanggal publikasi daring, dan status telaah sejawat: Tidak
  disebutkan dalam file.
\item
  Daftar referensi dibaca sebagai bagian dari seluruh TXT, tetapi isi
  karya yang dirujuk tidak diakses. Rincian metode, data, dan kesimpulan
  karya terdahulu hanya dapat diketahui sejauh J25 melaporkannya.
\item
  Batas provenance utama: review dapat menyatakan apa yang dilaporkan
  dan ditafsirkan artikel serta membuat inferensi yang diberi label,
  tetapi tidak dapat menguji ulang perhitungan, memastikan ketepatan
  ekstraksi terhadap PDF, memeriksa isi visual, menilai mutu sumber
  asli, atau membuktikan hubungan sebab-akibat yang tidak diuji J25.
\end{itemize}

\section{Review J26 - The Dynamics of Jabotabek
Development}\label{review-j26---the-dynamics-of-jabotabek-development}

\subsection{Identitas Sumber}\label{identitas-sumber-7}

\begin{itemize}
\tightlist
\item
  \textbf{Kode:} J26.
\item
  \textbf{Judul:} \emph{The Dynamics of Jabotabek Development}.
\item
  \textbf{Penulis pada halaman judul artikel:} J. Vernon Henderson, Ari
  Kuncoro, dan Damhuri Nasution (hlm. 71; TXT baris 74-90).
\item
  \textbf{Konflik kepengarangan dalam file:} blok metadata dan sitasi
  platform hanya mencantumkan J. Vernon Henderson dan Ari Kuncoro,
  sedangkan halaman judul artikel mencantumkan pula Damhuri Nasution.
  Konflik ini dipertahankan dan tidak diselesaikan dengan sumber luar
  (metadata platform sebelum artikel; hlm. 71).
\item
  \textbf{Afiliasi yang tercetak:} Henderson berafiliasi dengan Brown
  University dan National Bureau of Economic Research, sedangkan Kuncoro
  dan Nasution berafiliasi dengan University of Indonesia (hlm. 71).
  Blok metadata platform hanya menampilkan afiliasi Henderson dan
  Kuncoro.
\item
  \textbf{Jurnal:} \emph{Bulletin of Indonesian Economic Studies}.
\item
  \textbf{Volume, nomor, halaman:} Vol. 32, No.~1, April 1996, hlm.
  71-95 (masthead artikel, hlm. 71; blok sitasi platform).
\item
  \textbf{Tahun dan tanggal:} masthead artikel dan blok sitasi platform
  menyatakan 1996, sedangkan metadata platform menyatakan ``Published
  online: 19 Aug 2006''. Nama file sumber juga memuat angka 2006. Review
  ini tidak menyamakan ketiga penanda tersebut atau mengoreksinya dengan
  informasi eksternal.
\item
  \textbf{DOI sebagaimana tercantum dalam file:}
  10.1080/00074919612331336898.
\item
  \textbf{Wilayah kajian:} kawasan metropolitan Jakarta yang disebut
  Jabotabek, yakni DKI Jakarta dan bagian urban Bogor, Tangerang, serta
  Bekasi (hlm. 72-73; Lampiran A, hlm. 95).
\item
  \textbf{Periode empiris utama:} 1980-1990 untuk perubahan penduduk dan
  komposisi pekerjaan; 1986-1991 untuk perubahan spasial manufaktur;
  kelahiran pabrik 1989-1991; serta sejumlah pengamatan sampai awal
  1990-an, termasuk dokumen penggunaan lahan Agustus 1994 (hlm. 71-72,
  83-86).
\item
  \textbf{Basis akses:} teks lengkap hasil ekstraksi yang mencakup
  artikel cetak hlm. 71-95, metadata platform, catatan kaki, tabel,
  gambar, peta, referensi, dan lampiran.
\end{itemize}

\subsection{Summarized Abstract}\label{summarized-abstract-14}

\textbf{Laporan sumber.} Artikel menelaah dinamika perkembangan spasial
Jabotabek dari 1980 sampai awal 1990-an dengan memusatkan perhatian pada
perubahan lokasi penduduk dan kegiatan usaha. Penduduk dan manufaktur
sama-sama mengalami desentralisasi, tetapi dengan pola yang tidak
seragam. Kepadatan penduduk di pusat menurun sementara wilayah yang
lebih jauh mengalami densifikasi. Manufaktur formal bersuburbanisasi
dengan cepat ke Botabek setelah pembangunan jalan tol ke arah timur dan
barat; antara 1986 dan 1991, struktur manufaktur berubah dari pola
monosentris yang berorientasi pada Jakarta bagian utara-tengah menjadi
pola multisenter berbentuk sejumlah klaster (abstrak, hlm. 71; bagian
``Population Dispersion'', hlm. 73-82; bagian ``Industrial Composition
and Spatial Dispersion'', hlm. 82-93).

\textbf{Interpretasi penulis.} Penulis menilai pola desentralisasi
tersebut secara umum sesuai dengan pola perkembangan metropolitan yang
diperkirakan teori perkotaan, tetapi Jabotabek menghadapi masalah
khusus: kepadatan pusat masih sangat tinggi, investasi transportasi dan
layanan dasar rendah, tata guna lahan bersifat reaktif, pembangunan
pinggiran terpencar di sepanjang jalan utama, dan hak atas tanah
penduduk kampung lemah. Penulis juga menganggap kebijakan modernisasi
Jakarta Utara berisiko merusak industri skala kecil yang menopang
manufaktur formal serta memperlebar ketidakcocokan spasial antara tempat
tinggal pekerja berpendapatan rendah dan pekerjaan pinggiran (abstrak,
hlm. 71; hlm. 76-81, 90-93).

\textbf{Inferensi pengolahan.} Nilai utama artikel terletak pada
pembacaan terpadu antara perubahan bentuk kota, infrastruktur
transportasi, kelembagaan pasar tanah, dan restrukturisasi manufaktur.
Namun, hubungan kausal yang dikemukakan terutama bertumpu pada
kesesuaian waktu, pola spasial, teori, dan bukti deskriptif; artikel
tidak menyajikan rancangan identifikasi kausal untuk mengisolasi dampak
jalan tol, perencanaan, atau hak atas tanah.

\subsection{Problem Statement}\label{problem-statement-14}

\textbf{Laporan sumber.} Jabotabek mengalami pertumbuhan penduduk urban
rata-rata 5,9\% per tahun pada 1980-1990. Arus penduduk dan investasi
infrastruktur yang besar menimbulkan pertanyaan tentang bagaimana lokasi
penduduk, pekerjaan, dan penggunaan lahan berubah di dalam kawasan
metropolitan tersebut (pendahuluan, hlm. 72).

Masalah empiris pertama adalah kombinasi yang tampak bertentangan antara
kepadatan sangat tinggi di DKI Jakarta dan \emph{urban sprawl} di
pinggiran Botabek. Kepadatan pusat memang menurun, tetapi gradien
kepadatan masih curam, sementara proyek-proyek pinggiran berkembang
sebagai kantong atau pita yang tidak terhubung baik oleh jalan
pengumpan, air perpipaan, dan sewerage (hlm. 73-79).

Masalah kedua adalah perubahan sangat cepat dalam lokasi manufaktur.
Penulis ingin mengetahui apakah perpindahan manufaktur dari pusat
berarti dispersi sederhana atau pembentukan pusat-pusat industri baru,
serta apa konsekuensinya bagi industri kecil yang lebih sukar berpindah
(hlm. 82-91).

Masalah ketiga menyangkut institusi dan kebijakan: perencanaan
penggunaan lahan hampir tidak proaktif, investasi infrastruktur relatif
rendah, dan banyak penghuni kampung tidak memiliki hak terdaftar yang
cukup kuat untuk menjual, mengembangkan, atau memperoleh kompensasi
penuh atas tanahnya (hlm. 78-81).

\textbf{Interpretasi penulis.} Tanpa respons kebijakan, kombinasi
suburbanisasi pekerjaan, lemahnya mobilitas industri kecil, dan
penggusuran atau gentrifikasi kawasan pusat dapat menciptakan
ketidakcocokan spasial, pengangguran atau setengah pengangguran
perkotaan, serta ketidakpuasan sosial (abstrak, hlm. 71; hlm. 90-91).

\subsection{Objectives}\label{objectives-14}

\textbf{Laporan sumber.} Tujuan yang dinyatakan atau langsung tercermin
dalam pendahuluan adalah:

\begin{itemize}
\tightlist
\item
  Menelaah perkembangan spasial Jabotabek sebagai respons terhadap
  pertumbuhan penduduk dan investasi infrastruktur yang besar (hlm. 72).
\item
  Menggambarkan dan menganalisis desentralisasi serta suburbanisasi
  penduduk antara 1980 dan 1990 melalui perubahan gradien kepadatan
  (hlm. 73-77; Gambar 1-2; Tabel 1).
\item
  Mengidentifikasi ciri yang dianggap tidak biasa, terutama kepadatan
  pusat yang tinggi bersamaan dengan sprawl pinggiran, suburbanisasi
  cepat manufaktur formal, dan perubahan penggunaan lahan (hlm. 72).
\item
  Menilai perubahan distribusi spasial kegiatan industri pada 1986-1991
  dan menentukan apakah struktur manufaktur Jabotabek tetap monosentris
  atau telah menjadi multisenter (hlm. 82-89; Peta 2; Tabel 4-7).
\item
  Membahas persoalan kebijakan dan kelembagaan yang ditimbulkan oleh
  perubahan tersebut, termasuk tata guna lahan, hak atas tanah, layanan
  sanitasi, dan posisi industri skala kecil (hlm. 78-81, 90-93).
\end{itemize}

\textbf{Inferensi pengolahan.} Artikel juga secara implisit menguji
apakah model gradien kepadatan perkotaan yang lazim dapat menjelaskan
perubahan Jabotabek, kemudian menggunakan penyimpangan atau ciri khusus
lokal untuk merumuskan diagnosis kebijakan. Tujuan implisit ini
disimpulkan dari susunan analisis; file tidak menyatakannya sebagai
hipotesis formal.

\subsection{Literature Survey}\label{literature-survey-14}

\textbf{Laporan sumber.} Artikel tidak memuat bagian tinjauan pustaka
yang berdiri sendiri. Literatur dipakai secara tematik di sepanjang
analisis:

\begin{itemize}
\tightlist
\item
  Ingram dan Carroll (1981) digunakan sebagai rujukan pola
  desentralisasi penduduk dan industri di negara lain serta sebagai
  pembanding gradien kepadatan kota-kota Amerika Latin (hlm. 72, 74-75).
\item
  Henderson dan Kuncoro (1996) dirujuk untuk konteks sentralisasi sumber
  daya perkotaan dan subsidi implisit bagi lokasi industri di
  metropolitan terbesar Indonesia (hlm. 72).
\item
  World Bank (1993) digunakan untuk perbandingan laju pertumbuhan
  metropolitan, pertumbuhan penduduk urban, dan konteks layanan
  infrastruktur (hlm. 72, 77-78; daftar referensi, hlm. 95). Pada
  pembahasan kontaminasi air, teks hasil ekstraksi mencetak ``World Bank
  1953'', sedangkan daftar referensi hanya memuat World Bank (1993);
  konflik internal ini tidak diperbaiki (hlm. 78).
\item
  GOI (1993) digunakan untuk batas kawasan urban, analisis kepadatan
  tingkat desa, dan angka perjalanan lintas batas administratif (catatan
  kaki 2, hlm. 72; hlm. 82; daftar referensi, hlm. 94).
\item
  Struyk dkk. (1990), Aksoro (1994), serta Dowell dan Leaf (1991)
  menjadi dasar pembahasan sprawl, izin lokasi, kerumitan hak atas
  tanah, gradien harga tanah, dan selisih nilai antara klaim tanah kuat
  dan lemah (hlm. 78-80).
\item
  Newby, sebagaimana tercetak pada pembahasan septic system, dipakai
  untuk menghubungkan kepadatan rumah tapak dengan kegagalan sistem
  septik. Tahun yang tercetak dalam tubuh teks tampak berbeda dari
  daftar referensi dan tidak dinormalisasi dalam review ini (hlm. 78;
  referensi, hlm. 94).
\item
  Henderson (1991) dipakai untuk menjelaskan perbedaan biaya upah
  antarkota dan insentif desentralisasi industri (hlm. 85-86).
\item
  Mukherjee, Sarker dan Mukherjee, Lee, serta Hamer memberikan bukti
  dari Calcutta, Rio de Janeiro, Seoul, dan São Paulo mengenai hubungan
  pemasok industri kecil dengan perusahaan besar dan fungsi inkubasi
  industri kecil. Bukti tersebut berasal dari kota lain, bukan pengujian
  langsung atas hubungan produktivitas industri kecil di Jabotabek (hlm.
  90).
\end{itemize}

\textbf{Interpretasi penulis.} Literatur perkotaan dipakai sebagai tolok
ukur ``pola yang diharapkan'', sedangkan literatur kelembagaan tanah dan
industri kecil dipakai untuk menafsirkan mekanisme serta konsekuensi
kebijakan yang tidak semuanya diukur langsung oleh data Jabotabek.

\textbf{Inferensi pengolahan.} Tinjauan ini bersifat naratif dan
argumentatif, bukan tinjauan sistematis. Protokol pencarian, kriteria
inklusi, ukuran korpus, dan penilaian mutu studi: \textbf{Tidak
disebutkan dalam file}.

\subsection{Method Used}\label{method-used-14}

\textbf{Laporan sumber.} Artikel menggunakan kombinasi analisis
deskriptif, pemetaan, regresi gradien kepadatan, indeks konsentrasi,
perbandingan lintas waktu, dan pembahasan institusional.

\begin{itemize}
\tightlist
\item
  Untuk penduduk, penulis mengestimasi gradien kepadatan 1980 dan 1990
  pada kecamatan urban. Bentuk dasarnya adalah logaritma kepadatan
  penduduk per hektare sebagai fungsi jarak garis lurus dari Monas,
  kuadrat jarak, indikator kecamatan Bogor, dan jarak ke Kota Bogor bagi
  observasi Bogor. Jarak dari Monas dinyatakan dalam satuan 3,4 km.
  Model diestimasi dengan memasukkan dan mengeluarkan komponen Bogor
  untuk memeriksa kepekaan pola terhadap pusat kedua tersebut (hlm.
  75-76; Persamaan 1; Tabel 1; Gambar 1-2).
\item
  Penulis membandingkan rotasi gradien akibat perbaikan transportasi
  dengan pergeseran gradien akibat pertumbuhan penduduk. Sebuah
  kalkulasi kasar mengintegrasikan gradien untuk memperkirakan kepadatan
  pusat kontrafaktual apabila pertumbuhan penduduk tidak terjadi.
  Penulis sendiri menyebutnya \emph{back-of-the-envelope calculation}
  (hlm. 76-77 dan catatan kaki 3).
\item
  Untuk manufaktur, penulis memetakan kecamatan berdasarkan ambang 8.000
  pekerja manufaktur formal pada 1986 dan 1991 serta membandingkan
  perubahan pekerjaan di dalamnya (Peta 2, hlm. 84-85).
\item
  Penulis mengestimasi gradien logaritma kepadatan pekerjaan manufaktur
  per hektare terhadap jarak dari Monas, indikator Bogor, dan jarak ke
  Bogor. Berbeda dari model penduduk, suku kuadrat dikeluarkan. Estimasi
  dilakukan untuk seluruh manufaktur dan secara terpisah untuk tekstil
  dan pakaian jadi, kimia, serta mesin pada 1986 dan 1991, dengan
  pemeriksaan tambahan yang mengeluarkan Bogor (hlm. 86-88; Tabel 5-6).
\item
  Untuk membedakan desentralisasi dari dispersi, penulis menghitung
  Hirschman-Herfindahl Index (HHI), yakni jumlah kuadrat pangsa
  pekerjaan setiap kecamatan, dengan 61 kecamatan yang diseragamkan pada
  batas 1986 agar indeks 1986 dan 1991 dapat dibandingkan (hlm. 88-89;
  Tabel 7).
\item
  Tabel deskriptif digunakan untuk membandingkan komposisi
  sosial-ekonomi, struktur pekerjaan, tingkat suburbanisasi manufaktur,
  kelahiran pabrik, dan pekerjaan penduduk di sektor keuangan serta jasa
  bisnis (Tabel 2-4 dan 8, hlm. 81-83, 92).
\item
  Analisis kualitatif atas sprawl, izin lokasi, hak atas tanah, rencana
  waterfront, dan penggunaan lahan melengkapi hasil kuantitatif (hlm.
  78-81, 85, 90-93).
\end{itemize}

\textbf{Inferensi pengolahan.} Desain tersebut efektif untuk menunjukkan
perubahan pola spasial, tetapi bukan eksperimen atau kuasi-eksperimen.
Penafsiran pembangunan jalan tol sebagai pemicu suburbanisasi didasarkan
pada urutan waktu, orientasi koridor, besarnya perubahan, dan teori
biaya lokasi; tidak ada kelompok kontrol, instrumen,
\emph{difference-in-differences}, atau pengujian kausal lain yang
disebutkan dalam file.

\subsection{Dataset}\label{dataset-14}

\textbf{Laporan sumber.} Data dan bahan yang disebutkan meliputi:

\begin{itemize}
\tightlist
\item
  Rangkaian data penduduk 1980 dan 1990 digunakan untuk kepadatan,
  pendidikan, komposisi pekerjaan, dan industri tempat penduduk bekerja.
  Sensus Penduduk disebut secara eksplisit dalam pembahasan lokasi
  kegiatan nonmanufaktur, tetapi hanya menunjukkan tempat tinggal
  pekerja, bukan tempat kerjanya. Rincian sumber mentah setiap tabel
  penduduk tidak selalu dicantumkan dalam ekstraksi (hlm. 75-76, 81-83,
  91-92; Tabel 1-3 dan 8).
\item
  Analisis kepadatan tingkat desa, penetapan batas urban 1990, dan angka
  perjalanan lintas batas administratif berasal dari \emph{Jabotabek
  Metropolitan Development Plan Review} (GOI 1993). Data kepadatan
  kemudian diagregasi atau ditampilkan pada tingkat kecamatan (catatan
  kaki 2, hlm. 72-73; hlm. 82).
\item
  BPS \emph{Annual Survey of Medium and Large Scale Manufacturing},
  dalam berkas yang direvisi, diproses ulang, dan diperbarui oleh
  Development Studies Project (DSP), untuk pekerjaan manufaktur formal
  1986 dan 1991 serta kelahiran pabrik 1989-1991 (Tabel 4, Peta 2, dan
  sumber tabel/peta, hlm. 83-85).
\item
  Sensus Ekonomi 1986 untuk pekerjaan perusahaan kecil dan keseluruhan
  kegiatan industri, serta untuk membandingkan besaran sektor informal
  (Tabel 4, hlm. 83; hlm. 90 dan catatan kaki 10).
\item
  Dokumen penggunaan lahan Agustus 1994 yang tersimpan di Ministry of
  Public Works untuk mengamati perluasan industri ke Serang dan Karawang
  setelah 1991 (hlm. 85).
\item
  Grafik kenaikan upah menurut pekerjaan yang tidak dipublikasikan dan
  disiapkan DSP untuk mendukung pembahasan perbedaan biaya tenaga kerja
  (catatan kaki 9, hlm. 86).
\end{itemize}

Resolusi mikrodata, prosedur pembersihan, definisi seluruh variabel
sensus, bobot sampel, dan akses publik berkas DSP: \textbf{Tidak
disebutkan dalam file}.

\subsection{Population / Sample}\label{population-sample-14}

\textbf{Laporan sumber.} Unit analisis utama berbeda menurut bagian:

\begin{itemize}
\tightlist
\item
  Analisis gradien penduduk mencakup 78 kecamatan urban Jabotabek
  berdasarkan definisi 1990. Delapan belas observasi pada bagian Bogor
  dikeluarkan dalam spesifikasi pembanding, sehingga tersisa 60
  kecamatan (hlm. 75-76; Tabel 1).
\item
  Lampiran A merinci cakupan sebagai 42 kecamatan DKI Jakarta tanpa
  Kepulauan Seribu; tiga kecamatan Kotip Depok; lima Kecamatan Kotip
  Tangerang; empat Kecamatan Kotip Bekasi; lima kecamatan Kotamadya
  Bogor; serta kecamatan tambahan pada kabupaten Bogor, Tangerang, dan
  Bekasi. Lampiran mengingatkan bahwa sejumlah kecamatan kemudian
  dipecah atau digambar ulang batasnya (hlm. 95).
\item
  Analisis konsentrasi manufaktur lintas waktu menggunakan 61 kecamatan
  yang dibakukan pada batas 1986. Dalam regresi gradien pekerjaan,
  jumlah observasi berubah mengikuti definisi wilayah: Tabel 5
  menampilkan 61 dan 43 observasi untuk 1986 serta 78 dan 60 untuk 1991
  pada spesifikasi dengan dan tanpa Bogor (hlm. 86-89; Tabel 5 dan 7).
\item
  Manufaktur formal didefinisikan sebagai perusahaan menengah dan besar
  dengan lebih dari 20 pekerja. Perusahaan kecil memiliki 5-20 pekerja,
  sedangkan industri rumah tangga memiliki kurang dari lima pekerja; dua
  kategori terakhir diperlakukan sebagai sektor informal dalam estimasi
  untuk DKI Jakarta (hlm. 85, 90).
\item
  Tabel komposisi sosial-ekonomi mencakup penduduk berusia sepuluh tahun
  atau lebih. Tabel pangsa industri dan jasa mencakup pekerja berusia
  sepuluh tahun atau lebih yang bekerja pada minggu sebelumnya (catatan
  Tabel 2, 3, dan 8, hlm. 81, 83, 92).
\end{itemize}

Jumlah individu, rumah tangga, dan perusahaan di balik agregat
kecamatan: \textbf{Tidak disebutkan dalam file}. Teknik pengambilan
sampel pada tingkat individu atau perusahaan: \textbf{Tidak disebutkan
dalam file}.

\subsection{Dependent Variables}\label{dependent-variables-14}

\textbf{Laporan sumber.} Variabel terikat atau keluaran utama adalah:

\begin{itemize}
\tightlist
\item
  Logaritma kepadatan penduduk residensial per hektare pada 1980 dan
  1990 untuk regresi gradien kepadatan (Persamaan 1 dan Tabel 1, hlm.
  75-76).
\item
  Logaritma kepadatan pekerjaan manufaktur formal per hektare pada 1986
  dan 1991, baik untuk seluruh manufaktur maupun untuk tekstil dan
  pakaian jadi, kimia, serta mesin (Tabel 5-6, hlm. 86-88).
\item
  HHI geografis pekerjaan manufaktur sebagai ukuran konsentrasi spasial,
  dihitung untuk seluruh manufaktur dan tiga industri utama (Tabel 7,
  hlm. 88-89).
\item
  Ukuran deskriptif berupa pangsa pekerjaan menurut sektor, pangsa
  pekerjaan atau kelahiran pabrik di Botabek, rasio komposisi pendidikan
  wilayah terhadap komposisi metropolitan, persentase komuter lintas
  batas, dan pangsa pekerja penduduk dalam keuangan serta jasa bisnis
  (Tabel 2-4 dan 8, hlm. 81-83, 92).
\end{itemize}

Variabel penjelas regresi kepadatan terutama berupa jarak dari Monas,
kuadrat jarak pada model penduduk, indikator Bogor, dan jarak ke Kota
Bogor bagi kecamatan Bogor. Tidak ada variabel terikat tunggal untuk
pembahasan kualitatif tentang sprawl, hak atas tanah, sanitasi, atau
industri kecil; untuk bagian tersebut, konsep variabel terikat formal:
\textbf{Tidak berlaku}.

\subsection{Results}\label{results-14}

\textbf{Laporan hasil empiris: penduduk dan kepadatan.}

\begin{itemize}
\tightlist
\item
  Penduduk urban Jabotabek tumbuh rata-rata 5,9\% per tahun pada
  1980-1990; sumber juga menyebut kenaikan total sekitar 75\% selama
  dekade tersebut (hlm. 72, 77).
\item
  Kepadatan dekat pusat pada 1980 sekitar 42.000 penduduk per km
  persegi. Angka ini ditempatkan dekat kepadatan Mexico City dalam
  pembanding yang digunakan penulis (hlm. 73-75).
\item
  Model gradien penduduk memiliki daya penjelas yang dinyatakan tinggi,
  dengan R kuadrat sekitar 0,71 pada 1980 dan 0,73 pada 1990. Pola umum
  serupa ketika observasi Bogor dikeluarkan (hlm. 75-76; Tabel 1).
\item
  Pada jarak 3,4 km dari Monas, laju penurunan kepadatan per kilometer
  berkurang dari 17\% pada 1980 menjadi 12\% pada 1990. Kepadatan
  menurun di pusat, tetapi meningkat di sebagian besar kawasan
  metropolitan yang lebih jauh, termasuk bagian DKI Jakarta (hlm. 76-77;
  Gambar 1).
\item
  Kalkulasi kasar penulis menyatakan bahwa tanpa pertumbuhan penduduk,
  rotasi gradien yang dikaitkan dengan investasi transportasi akan
  menurunkan kepadatan pusat dari sekitar 42.000 menjadi sekitar 20.000
  penduduk per km persegi pada 1990. Ini adalah hasil kontrafaktual
  kasar, bukan pengamatan langsung (hlm. 77 dan catatan kaki 3).
\item
  Bagian gradien menuju Bogor tidak berpotongan antara 1980 dan 1990;
  penulis membacanya sebagai stagnasi perkembangan Bogor pada periode
  tersebut (hlm. 77; Gambar 2).
\end{itemize}

\textbf{Laporan hasil empiris: penggunaan lahan, sosial-ekonomi, dan
mobilitas.}

\begin{itemize}
\tightlist
\item
  Pembangunan pinggiran berlangsung sebagai kantong-kantong proyek
  privat yang dangkal di sepanjang jalan utama, dengan koneksi jalan
  sekunder, air, dan sewerage yang lemah. Ruang hijau di antaranya
  digambarkan lebih sebagai akibat sementara daripada hasil rencana
  jangka panjang (hlm. 78-79).
\item
  Tabel 2 menunjukkan konsentrasi relatif penduduk berpendidikan tinggi
  di Jakarta Selatan, Timur, dan Pusat, sedangkan Jakarta Utara dan
  Botabek memiliki representasi relatif lebih besar dari penduduk
  berpendidikan rendah. Antara 1980 dan 1990, baik kelompok berstatus
  sosial-ekonomi tinggi maupun rendah menunjukkan suburbanisasi (hlm.
  80-82; Tabel 2).
\item
  Di DKI Jakarta pada 1990, 27-43\% tenaga kerja penduduk suatu
  kotamadya bepergian ke luar kotamadyanya. Untuk wilayah Botabek
  terdekat, teks menyatakan persentasenya jauh lebih tinggi, tetapi
  rentang angkanya rusak dalam ekstraksi menjadi \texttt{5M2\%}; review
  ini tidak menebak nilai aslinya (hlm. 82).
\end{itemize}

\textbf{Laporan hasil empiris: komposisi dan lokasi industri.}

\begin{itemize}
\tightlist
\item
  Pangsa manufaktur dalam seluruh pekerjaan Jabotabek meningkat dari
  13\% pada 1980 menjadi 22\% pada 1990, sedangkan keuangan dan jasa
  bisnis naik dari 3\% menjadi 5\%. Pangsa jasa kemasyarakatan, sosial,
  dan personal turun dari 34\% menjadi 25\% (hlm. 82-83; Tabel 3).
\item
  Pangsa pekerjaan manufaktur formal Jabotabek yang berlokasi di Botabek
  naik dari 43\% pada 1986 menjadi 56\% pada 1991. Pada 1989-1991, 73\%
  kelahiran pabrik sektor korporasi terjadi di Botabek, dibandingkan
  32\% pada sektor formal yang tidak berbadan hukum
  (\emph{unincorporated}) (hlm. 83-85; Tabel 4).
\item
  Peta 2 menunjukkan pertumbuhan pekerjaan manufaktur 40-120\% pada
  kecamatan yang telah memiliki lebih dari 8.000 pekerja pada kedua
  tahun. Di kecamatan yang baru melampaui ambang tersebut pada 1991,
  kenaikan paling lambat mencapai 210\%; satu kecamatan tumbuh empat
  belas kali lipat menjadi 93.000 pekerja, sedangkan tiga kecamatan lain
  bergerak dari nol menjadi 13.000, 16.000, dan 22.000 pekerja. Sejumlah
  kecamatan DKI dekat pusat justru turun dari atas menjadi bawah ambang
  8.000 (hlm. 84-85; Peta 2).
\item
  Penulis melaporkan selisih upah sekitar 25\% antara manufaktur formal
  Jakarta dan Botabek sebagai insentif relokasi. Klaim ini disajikan
  dalam pembahasan biaya lokasi, bukan sebagai koefisien dari model upah
  dalam artikel (hlm. 85-86).
\item
  Gradien kepadatan pekerjaan manufaktur bersifat signifikan dan
  monosentris pada 1986, dengan R kuadrat tersesuaikan sekitar 0,382.
  Pada 1991, kepadatan pekerjaan pusat berkurang separuh, koefisien
  jarak tidak signifikan, dan R kuadrat tersesuaikan sekitar 0,078.
  Penulis menafsirkan hilangnya hubungan dengan jarak pusat sebagai
  hilangnya monosentrisitas manufaktur; hasil serupa diperoleh ketika
  Bogor dikeluarkan (hlm. 86-87; Tabel 5).
\item
  Tekstil dan pakaian jadi, kimia, serta mesin masing-masing mencakup
  sekitar 41\%, 17\%, dan 19\% pekerjaan manufaktur Jabotabek. Ketiganya
  menunjukkan gradien monosentris yang signifikan pada 1986, tetapi
  tidak pada 1991 (hlm. 87-88; Tabel 6).
\item
  HHI meningkat untuk seluruh manufaktur dan ketiga industri utama,
  sehingga kehilangan monosentrisitas tidak sama dengan penyebaran
  merata. Nilai yang terbaca jelas naik dari 0,0467 menjadi 0,0567 untuk
  tekstil dan pakaian jadi, dari 0,0789 menjadi 0,0863 untuk kimia,
  serta dari 0,0925 menjadi 0,0929 untuk mesin. Nilai 1991 bagi seluruh
  manufaktur terdistorsi dalam TXT menjadi \texttt{4585}; kesimpulan
  arah kenaikannya mengikuti pernyataan eksplisit penulis, bukan koreksi
  angka oleh pengolah (hlm. 88-89; Tabel 7).
\item
  Aktivitas manufaktur membentuk klaster di Jakarta Utara, koridor
  Bekasi, koridor Tangerang, dan jalur Kabupaten Bogor. Tiga pusat
  pertama pada praktiknya membentuk pita timur-barat yang tersambung
  melalui jalan tol dan jalan utama, sedangkan jalur Bogor lebih berat
  pada tekstil dan pakaian jadi (hlm. 89; Peta 2).
\end{itemize}

\textbf{Laporan hasil empiris: sektor informal dan jasa.}

\begin{itemize}
\tightlist
\item
  Penulis memperkirakan sekitar 40\% pekerjaan manufaktur DKI Jakarta
  pada 1990 berada pada perusahaan dengan kurang dari 20 pekerja.
  Perkiraan diperoleh dengan membandingkan Sensus Penduduk dan survei
  perusahaan manufaktur menengah-besar, serta disebut sebanding dengan
  angka Sensus Ekonomi 1986 (hlm. 90 dan catatan kaki 10).
\item
  Pada 1986, hanya 24\% pekerjaan perusahaan kecil berada di Botabek,
  dibandingkan 43\% pekerjaan manufaktur formal. Perbedaan kelahiran
  pabrik korporasi dan formal tidak berbadan hukum pada 1989-1991 juga
  menunjukkan mobilitas lebih tinggi pada perusahaan besar (hlm. 90;
  Tabel 4).
\item
  Jakarta Utara masih menampung 25\% pekerjaan manufaktur formal
  metropolitan dan konsentrasi besar industri kecil, sekalipun
  manufaktur secara keseluruhan bersuburbanisasi (hlm. 90-91).
\item
  Pangsa pekerja penduduk dalam keuangan dan jasa bisnis meningkat di
  semua wilayah antara 1980 dan 1990. Angkanya naik dari 4,8\% menjadi
  8,4\% di Jakarta Selatan; 3,0\% menjadi 6,6\% di Jakarta Timur; 4,4\%
  menjadi 9,7\% di Jakarta Pusat; 3,5\% menjadi 6,5\% di Jakarta Barat;
  2,0\% menjadi 5,9\% di Jakarta Utara; 1,0\% menjadi 1,2\% di Bogor;
  1,2\% menjadi 2,8\% di Bekasi; dan 2,0\% menjadi 3,0\% di Tangerang
  (hlm. 91-92; Tabel 8).
\end{itemize}

\subsection{Findings}\label{findings-14}

\textbf{Temuan yang dilaporkan sumber.}

\begin{itemize}
\tightlist
\item
  Jabotabek mengikuti kecenderungan umum metropolitan berupa penurunan
  kepadatan pusat, densifikasi pinggiran, suburbanisasi penduduk, dan
  suburbanisasi pekerjaan. Namun, kepadatan pusat dan kecuraman
  gradiennya tetap tinggi (hlm. 71-77).
\item
  Dalam lima tahun yang mencakup perluasan jalan tol ke Bekasi dan
  Tangerang, manufaktur formal berubah dari pola monosentris menjadi
  multisenter. Perubahan itu berbentuk pembentukan klaster baru, bukan
  dispersi seragam, karena konsentrasi geografis justru meningkat (hlm.
  83-89; Tabel 4-7; Peta 2).
\item
  Manufaktur korporasi jauh lebih mudah berpindah ke Botabek daripada
  perusahaan kecil dan perusahaan formal tidak berbadan hukum. Karena
  itu, Jakarta Utara tetap penting sebagai pusat manufaktur formal
  sekaligus industri kecil (hlm. 89-91).
\item
  Jasa keuangan dan bisnis bertumbuh tetapi tetap relatif
  tersentralisasi di ``Golden Triangle'' Jakarta Pusat-Jakarta Selatan;
  keterbatasan data tempat kerja membuat hasil untuk sektor
  nonmanufaktur hanya merupakan gambaran kasar (hlm. 91-92; Tabel 8).
\end{itemize}

\textbf{Interpretasi penulis.}

\begin{itemize}
\tightlist
\item
  Penulis mengaitkan gradien penduduk yang masih curam dengan kurangnya
  investasi transportasi, ketiadaan angkutan massal, buruknya
  pengelolaan lalu lintas kendaraan, dan lemahnya fasilitas nonkendaraan
  (hlm. 77).
\item
  Suburbanisasi manufaktur ditafsirkan sebagai pelepasan permintaan
  terpendam untuk keluar dari biaya tanah dan upah yang tinggi di DKI
  setelah koridor transportasi tersedia (hlm. 85-86, 93).
\item
  Peningkatan HHI ditafsirkan sebagai bukti manfaat kedekatan
  antarpabrik, limpahan informasi, dan penghematan biaya pengangkutan
  dalam rantai produksi. Mekanisme eksternalitas ini dijelaskan secara
  teoretis, tetapi tidak diestimasi langsung (hlm. 89).
\item
  Hak tanah yang lemah dianggap menghambat konsolidasi, konversi, dan
  peningkatan kualitas lahan serta mengurangi insentif penghuni untuk
  berinvestasi pada rumah dan lingkungan (hlm. 79-80).
\item
  Penulis menilai penghilangan industri kecil dari Jakarta Utara akan
  menurunkan produktivitas manufaktur formal dan menimbulkan
  ketidakcocokan spasial bagi pekerja kampung. Peran produktivitas
  tersebut didukung analogi dari kota lain, bukan pengukuran hubungan
  input-output atau inovasi industri kecil Jabotabek dalam artikel ini
  (hlm. 90-91).
\end{itemize}

\textbf{Inferensi pengolahan.} Artikel memperlihatkan dua proses yang
harus dibedakan: desentralisasi terhadap pusat historis dan konsentrasi
kembali pada pusat-pusat baru. Pembedaan ini didukung langsung oleh
kombinasi gradien dan HHI. Sebaliknya, klaim bahwa jalan tol menyebabkan
perpindahan, bahwa pemberian hak akan menghasilkan konversi yang
efisien, dan bahwa industri kecil menaikkan produktivitas korporasi
lebih tepat dibaca sebagai interpretasi mekanisme yang masuk akal
daripada efek kausal yang telah diukur.

\subsection{Conclusions}\label{conclusions-14}

\textbf{Kesimpulan penulis.} Jabotabek memiliki investasi privat yang
kuat, ruang hijau yang masih tersebar, fondasi jaringan jalan arteri,
penurunan kepadatan pusat, dan struktur metropolitan yang makin
multisenter. Pada saat yang sama, pertumbuhannya menghadirkan persoalan
besar yang membutuhkan respons kelembagaan dan infrastruktur (bagian
``Conclusions'', hlm. 92-93).

\begin{itemize}
\tightlist
\item
  Perbaikan transportasi dan kenaikan pendapatan diperkirakan akan terus
  menurunkan kepadatan dekat pusat serta mendorong densifikasi dan
  suburbanisasi lebih jauh. Urutan perubahan kawasan seharusnya
  memengaruhi waktu investasi sewerage dan infrastruktur lain (hlm.
  92-93).
\item
  Perencanaan tata guna lahan yang reaktif menghasilkan pembangunan
  pita, sprawl, dan layanan air, jalan lokal, serta sewerage yang tidak
  memadai. Hak tanah yang tidak jelas menghambat pembangunan kembali
  kawasan kampung (hlm. 93).
\item
  Manufaktur telah menjadi multisenter melalui klaster-klaster baru,
  bukan melalui penyebaran acak. Suburbanisasi lebih lanjut dan
  kemungkinan desentralisasi ke kota kecil-menengah di bagian lain Jawa
  diproyeksikan berlanjut (hlm. 93).
\item
  Jakarta Utara tetap merupakan pusat manufaktur utama. Rencana
  waterfront dan modernisasi wilayah tersebut seharusnya
  mempertimbangkan keberadaan serta fungsi industri skala kecil (hlm.
  90-93).
\end{itemize}

\textbf{Inferensi pengolahan.} Kesimpulan artikel bersifat diagnosis
perkembangan dan rekomendasi kebijakan metropolitan, bukan kesimpulan
evaluasi dampak. Kekuatan buktinya paling tinggi untuk perubahan
distribusi spasial yang terukur dan lebih rendah untuk mekanisme
produktivitas, kesejahteraan, serta keresahan sosial yang diproyeksikan.

\subsection{Contributions}\label{contributions-14}

Bagian kontribusi eksplisit: \textbf{Tidak disebutkan dalam file}.
Kontribusi berikut merupakan \textbf{inferensi pengolahan} berdasarkan
isi artikel:

\begin{itemize}
\tightlist
\item
  Menyediakan gambaran lintas waktu yang menghubungkan perubahan
  kepadatan penduduk 1980-1990 dengan restrukturisasi lokasi manufaktur
  1986-1991 dalam satu kerangka metropolitan Jabotabek.
\item
  Menunjukkan secara empiris bahwa hilangnya monosentrisitas tidak
  identik dengan berkurangnya konsentrasi. Gradien jarak melemah, tetapi
  HHI meningkat karena kegiatan terkumpul dalam pusat-pusat baru (hlm.
  86-89; Tabel 5-7).
\item
  Menggabungkan ukuran kuantitatif dengan analisis kelembagaan atas izin
  lokasi, hak tanah, infrastruktur, dan tata guna lahan sehingga
  perubahan bentuk kota tidak dijelaskan hanya oleh jarak atau biaya
  transportasi (hlm. 78-81).
\item
  Mengangkat perbedaan mobilitas antara perusahaan korporasi, perusahaan
  formal tidak berbadan hukum, dan industri kecil sebagai dimensi
  penting suburbanisasi industri (hlm. 83, 89-91; Tabel 4).
\item
  Menghubungkan restrukturisasi ruang produksi dengan risiko distributif
  bagi penghuni kampung dan pekerja berpendapatan rendah, walaupun
  konsekuensi kesejahteraan tersebut tidak diukur secara langsung (hlm.
  79-81, 90-91).
\end{itemize}

\subsection{Research Gap}\label{research-gap-14}

Bagian bernama \emph{research gap}: \textbf{Tidak disebutkan dalam
file}. Namun, sumber menyatakan dua kekosongan pengetahuan atau data
secara langsung:

\begin{itemize}
\tightlist
\item
  Penulis menyatakan bahwa pengaruh perubahan kepadatan penduduk yang
  cepat terhadap kontaminasi, pembangunan kembali kawasan, dan penentuan
  waktu infrastruktur belum dipahami dengan baik (hlm. 78).
\item
  Perkembangan spasial sektor nonmanufaktur sukar dilacak karena Sensus
  Ekonomi hanya berlangsung sepuluh tahunan, data 1986 tidak mudah
  tersedia dan sebagian besar belum dianalisis, sedangkan Sensus
  Penduduk mencatat tempat tinggal, bukan tempat kerja (hlm. 91).
\end{itemize}

\textbf{Inferensi pengolahan.} Kesenjangan yang masih tersisa dari
desain artikel adalah:

\begin{itemize}
\tightlist
\item
  Belum ada estimasi kausal terpisah atas pengaruh jalan tol terhadap
  relokasi industri dibandingkan pengaruh harga tanah, upah, kebijakan
  zona ekspor, dan tren pertumbuhan umum.
\item
  Fungsi pemasok, daur ulang, perbaikan, inovasi, dan inkubasi industri
  kecil di Jabotabek belum diuji langsung; argumen artikel terutama
  memakai bukti dari kota lain (hlm. 90).
\item
  Dampak aktual suburbanisasi terhadap waktu tempuh, pengangguran,
  setengah pengangguran, pendapatan, dan keresahan sosial belum diukur.
\item
  Dampak pemberian hak terdaftar kepada penghuni lama terhadap
  efisiensi, kompensasi, investasi rumah, penggusuran, dan distribusi
  keuntungan pembangunan kembali belum dievaluasi.
\item
  Perubahan setelah 1991 hanya ditunjukkan secara terbatas melalui
  dokumen lahan 1994 dan proyeksi, bukan seri data pekerjaan yang
  sebanding.
\end{itemize}

\subsection{Limitations}\label{limitations-14}

Bagian keterbatasan formal: \textbf{Tidak disebutkan dalam file}.
Keterbatasan berikut terdiri atas peringatan yang dilaporkan sumber dan
inferensi metodologis pengolahan.

\textbf{Keterbatasan yang dilaporkan sumber.}

\begin{itemize}
\tightlist
\item
  Data lokasi nonmanufaktur sangat terbatas. Sensus Penduduk
  menghubungkan industri dengan tempat tinggal pekerja, bukan tempat
  kerja, sehingga pola keuangan dan jasa bisnis pada Tabel 8 hanya
  memberi gambaran kasar (hlm. 91-92).
\item
  Analisis utama suburbanisasi industri menggunakan perusahaan menengah
  dan besar; gambar tersebut secara eksplisit disebut tidak lengkap
  karena mengecualikan sebagian besar perusahaan kecil dan industri
  rumah tangga (hlm. 89-90).
\item
  Batas kecamatan berubah antara 1986 dan 1990/1991. Penulis menggunakan
  batas 1986 untuk perbandingan HHI dan peta tertentu, tetapi jumlah
  serta definisi observasi tetap berbeda pada beberapa analisis (catatan
  kaki 2, hlm. 72; catatan Tabel 5, hlm. 87; Lampiran A, hlm. 95).
\item
  Perkiraan kepadatan pusat tanpa pertumbuhan penduduk merupakan
  kalkulasi kasar yang bergantung pada integrasi gradien, kemiringan
  1990, dan total penduduk 1980 (catatan kaki 3, hlm. 77).
\item
  Perkiraan sektor informal sekitar 40\% diperoleh dari perbandingan dua
  sumber agregat, bukan pencacahan khusus sektor informal dalam desain
  artikel (hlm. 90 dan catatan kaki 10).
\end{itemize}

\textbf{Inferensi pengolahan.}

\begin{itemize}
\tightlist
\item
  Regresi gradien memakai sedikit prediktor spasial dan tidak
  mengendalikan seluruh karakteristik lokasi yang mungkin berubah
  bersamaan dengan pembangunan jalan tol.
\item
  Perbandingan sebelum-sesudah dan pola koridor tidak cukup untuk
  mengisolasi sebab-akibat; klaim mekanisme harus dibaca dengan
  kehati-hatian.
\item
  Analisis pada agregat kecamatan berpotensi menyembunyikan variasi
  intrakecamatan dan tidak dapat langsung diterjemahkan menjadi perilaku
  rumah tangga atau perusahaan individual.
\item
  Bukti mengenai industri kecil, kesehatan, hak tanah, dan
  ketidakcocokan spasial tidak setara kedalamannya dengan bukti gradien
  serta konsentrasi manufaktur.
\item
  Tidak ada ukuran ketidakpastian untuk sejumlah persentase deskriptif,
  HHI, proyeksi, atau rekomendasi kebijakan yang disebutkan.
\end{itemize}

\subsection{Challenges}\label{challenges-14}

\textbf{Tantangan pembangunan yang dilaporkan dan ditafsirkan penulis.}

\begin{itemize}
\tightlist
\item
  Menurunkan kepadatan pusat dan eksternalitasnya tanpa mendorong sprawl
  tidak terkoordinasi di pinggiran (hlm. 77-79).
\item
  Mengatasi pencemaran akibat kombinasi sumur air tanah, sistem septik,
  dan kepadatan rumah tapak, sambil menentukan waktu sewerage agar tidak
  segera usang akibat pembangunan kembali (hlm. 78).
\item
  Mengubah perencanaan dari pemberian izin lokasi yang reaktif menjadi
  perkiraan permintaan lahan, alokasi penggunaan lahan, dan penyediaan
  infrastruktur sebelum pembangunan (hlm. 79).
\item
  Memperkuat hak penduduk kampung tanpa mengabaikan persoalan
  pemerataan, kompensasi, spekulasi, dan penanganan pendudukan baru.
  Penulis mengusulkan hak bagi penghuni lama dengan klaim penguasaan
  kuat, tetapi catatan kaki juga menyatakan penghuni liar baru
  seharusnya secara rutin digusur; keduanya dilaporkan tanpa dukungan
  normatif dari review ini (hlm. 80 dan catatan kaki 5).
\item
  Mengelola perjalanan lintas wilayah yang meningkat akibat
  suburbanisasi penduduk dan pekerjaan (hlm. 81-82).
\item
  Menjaga fungsi industri kecil di Jakarta Utara di tengah modernisasi,
  gentrifikasi, dan rencana pembangunan waterfront (hlm. 90-91).
\item
  Mencegah ketidakcocokan spasial bagi pekerja berkeahlian rendah ketika
  pekerjaan bergerak ke pinggiran tetapi sumber daya perjalanan mereka
  terbatas (hlm. 90-91).
\end{itemize}

\textbf{Tantangan analitis yang dilaporkan sumber.} Keterbatasan data
tempat kerja nonmanufaktur, perubahan batas kecamatan, dan cakupan
survei manufaktur formal menyulitkan penyusunan gambaran metropolitan
yang sepenuhnya sebanding (hlm. 72, 87, 89-92).

\subsection{Practical Implications}\label{practical-implications-14}

\textbf{Implikasi dan rekomendasi penulis.}

\begin{itemize}
\tightlist
\item
  Pemerintah lokal perlu memperkirakan permintaan lahan sebelum atau
  bersamaan dengan pembangunan jalan tol, merevisi rencana penggunaan
  lahan dan kebutuhan infrastruktur, lalu menilai proposal pengembang
  berdasarkan rencana tersebut (hlm. 79).
\item
  Jaringan jalan pengumpan, jalan sekunder, air, dan sewerage perlu
  direncanakan lintas proyek agar pembangunan tidak hanya membentuk pita
  dangkal di sepanjang jalan raya (hlm. 78-79, 93).
\item
  Investasi jalan lokal, sewerage, dan pipa air di lingkungan lama perlu
  diselaraskan dengan urutan pembangunan kembali kawasan. Penulis
  menyarankan investasi besar mengikuti konfigurasi pembangunan kembali
  agar tidak segera menjadi tidak sesuai (hlm. 78).
\item
  Mekanisme pemberian klaim terdaftar kepada penghuni lama dengan hak
  penguasaan kuat dipandang dapat menaikkan efisiensi transaksi,
  insentif perbaikan rumah, nilai kompensasi, dan bagian keuntungan
  pembangunan kembali yang diterima penduduk kampung (hlm. 79-80, 93).
\item
  Investasi transportasi massal, manajemen lalu lintas, serta fasilitas
  bagi perjalanan nonkendaraan dipandang diperlukan untuk meratakan
  gradien kepadatan dan mendukung akses metropolitan (hlm. 77).
\item
  Rencana waterfront Jakarta Utara seharusnya mengakui bahwa wilayah
  tersebut masih menjadi pusat manufaktur dan tidak menghapus industri
  kecil tanpa mempertimbangkan hubungan produksi serta akses kerja
  penduduk kampung (hlm. 90-93).
\end{itemize}

\textbf{Inferensi pengolahan.} Rekomendasi tersebut perlu diterapkan
sebagai hipotesis kebijakan yang diuji, bukan sebagai efek yang telah
pasti, karena artikel tidak mengestimasi biaya-manfaat, dampak
distribusi, atau respons perilaku dari intervensi yang diusulkan.

\subsection{Applications}\label{applications-14}

\textbf{Aplikasi yang langsung tercermin dalam sumber.}

\begin{itemize}
\tightlist
\item
  Gradien kepadatan dapat dipakai untuk memantau rotasi struktur
  penduduk antara pusat dan pinggiran serta untuk mengantisipasi lokasi
  kebutuhan transportasi dan sanitasi (hlm. 75-78).
\item
  Kombinasi gradien pekerjaan dan HHI dapat membedakan hilangnya
  dominasi pusat lama dari dispersi kegiatan secara merata; dalam kasus
  Jabotabek, kombinasi itu mengidentifikasi pembentukan klaster
  multisenter (hlm. 86-89).
\item
  Peta ambang pekerjaan dapat menandai kecamatan yang menjadi pusat
  baru, tetap kuat, atau mengalami penurunan manufaktur (Peta 2, hlm.
  84-85).
\item
  Analisis perbedaan mobilitas menurut ukuran dan bentuk usaha dapat
  digunakan dalam penilaian rencana regenerasi industri dan perlindungan
  akses kerja (Tabel 4, hlm. 83, 89-91).
\item
  Diagnosis izin lokasi dan hak tanah dapat digunakan untuk menilai
  apakah rencana wilayah mendahului pembangunan atau sekadar disesuaikan
  setelah izin dan konstruksi terjadi (hlm. 79-80).
\end{itemize}

\textbf{Inferensi pengolahan.} Kerangka ini dapat menjadi alat diagnosis
awal bagi perencanaan Jabotabek, tetapi penerapan pada periode atau
metropolitan lain memerlukan data batas wilayah yang seragam, lokasi
tempat kerja, sektor informal, jaringan transportasi, dan pengujian
mekanisme lokal. File tidak mengklaim validitas universal.

\subsection{Future Research
(Optional)}\label{future-research-optional-14}

Agenda penelitian masa depan yang dirumuskan secara eksplisit:
\textbf{Tidak disebutkan dalam file}. Sumber lebih banyak memberikan
proyeksi: densifikasi pinggiran, suburbanisasi manufaktur lanjutan,
kemungkinan relokasi ke kota kecil-menengah di Jawa, dan peningkatan
perjalanan silang wilayah (hlm. 77-78, 82, 85-86, 93). Proyeksi tersebut
bukan rancangan penelitian.

\textbf{Inferensi pengolahan untuk penelitian lanjutan.}

\begin{itemize}
\tightlist
\item
  Mengestimasi dampak kausal pembukaan jalan tol terhadap kelahiran,
  perpindahan, pekerjaan, dan produktivitas pabrik dengan data panel
  perusahaan dan pembanding koridor yang memadai.
\item
  Membangun data tempat kerja untuk sektor jasa agar sentralisasi
  ``Golden Triangle'' dapat diukur dari lokasi pekerjaan, bukan lokasi
  tempat tinggal pekerja.
\item
  Menguji secara langsung hubungan pemasok, inovasi, daur ulang, dan
  inkubasi antara industri kecil dan perusahaan korporasi Jabotabek.
\item
  Mengukur ketidakcocokan spasial melalui waktu tempuh, biaya
  perjalanan, akses pekerjaan, pengangguran, dan hasil sosial rumah
  tangga berpendapatan rendah.
\item
  Mengevaluasi dampak pemberian hak tanah terdaftar terhadap investasi
  rumah, transaksi, kompensasi, perpindahan paksa, dan distribusi
  keuntungan pembangunan kembali.
\item
  Memperbarui gradien, HHI, dan peta klaster setelah 1991 dengan unit
  geografis yang diharmonisasi untuk menguji proyeksi desentralisasi
  lebih jauh.
\item
  Menilai bersama urutan pembangunan kembali kawasan dan investasi air
  limbah agar rekomendasi waktu infrastruktur dapat diuji secara teknis,
  fiskal, dan distributif.
\end{itemize}

\subsection{Provenance Notes}\label{provenance-notes-14}

\begin{itemize}
\tightlist
\item
  Review ini disusun hanya dari seluruh file TXT
  \texttt{source/journal/J26-henderson-kuncoro-nasution-2006-the-dynamics-of-jabotabek-development-10.1080:00074919612331336898.txt},
  yang dibaca sampai akhir pada baris 1198. Tidak ada sumber luar yang
  digunakan untuk menambah, mengoreksi, atau merekonsiliasi isi.
\item
  Teks yang tersedia mencakup metadata unduhan dan platform, isi artikel
  cetak hlm. 71-95, delapan tabel, dua gambar gradien, dua peta, catatan
  kaki, daftar referensi, serta Lampiran A. Oleh karena itu, basis
  review adalah teks lengkap hasil ekstraksi, bukan sekadar abstrak.
  Tata letak visual PDF asli tidak diperiksa.
\item
  Locator utama mengikuti nomor halaman cetak yang tampak pada
  ekstraksi. Locator tabel, gambar, peta, persamaan, bagian, catatan
  kaki, dan lampiran ditambahkan ketika tersedia. Nomor baris TXT
  digunakan khusus untuk konflik identitas pada bagian awal dan cakupan
  berkas.
\item
  Semua pernyataan berlabel \textbf{Laporan sumber} atau \textbf{Temuan
  yang dilaporkan sumber} merupakan parafrasa atas apa yang dinyatakan
  atau ditampilkan artikel. Pernyataan berlabel \textbf{Interpretasi
  penulis} mempertahankan penjelasan mekanisme, penilaian, proyeksi,
  atau rekomendasi penulis. Pernyataan berlabel \textbf{Inferensi
  pengolahan} adalah penilaian analitis review dan bukan posisi yang
  diklaim telah diuji oleh artikel.
\item
  Identitas dalam file tidak sepenuhnya konsisten. Artikel cetak dan
  sitasi platform memberi tahun 1996, metadata platform memberi tanggal
  terbit daring 19 Agustus 2006, dan nama file memuat 2006. Halaman
  judul artikel memuat Henderson, Kuncoro, dan Nasution, sedangkan
  metadata/sitasi platform hanya memuat Henderson dan Kuncoro. Semua
  versi dicatat tanpa memilih koreksi eksternal.
\item
  Hasil ekstraksi mengandung korupsi OCR. Contoh yang material adalah
  rentang komuter Botabek yang tampil sebagai \texttt{5M2\%}, nilai HHI
  seluruh manufaktur 1991 yang tampil sebagai \texttt{4585}, beberapa
  koefisien tabel, ejaan nama tempat, serta konflik tahun pada sitasi
  World Bank/Newby di tubuh teks dan daftar referensi. Angka yang tidak
  dapat dibaca secara andal tidak direkonstruksi. Arah temuan hanya
  dipertahankan bila dijelaskan eksplisit dalam prosa artikel.
\item
  Referensi sekunder yang dikutip artikel tidak dibuka atau
  diverifikasi. Klaim yang bersandar pada kota lain, laporan pemerintah,
  dokumen penggunaan lahan, atau grafik DSP diperlakukan sebagai laporan
  artikel mengenai bahan tersebut, bukan sebagai hasil verifikasi
  mandiri.
\item
  Status telaah sejawat artikel tidak diverifikasi dari luar file.
  Ucapan terima kasih menyebut komentar dua referee anonim (hlm. 71),
  tetapi review ini tidak menggunakan keterangan itu untuk membuat klaim
  tambahan tentang status penerbitan.
\end{itemize}

\section{Review J27}\label{review-j27}

\subsection{Identitas Sumber}\label{identitas-sumber-8}

\begin{itemize}
\tightlist
\item
  \textbf{Kode sumber:} J27.
\item
  \textbf{Judul:} \emph{Peri-urban transformation in the Jakarta
  metropolitan area}.
\item
  \textbf{Penulis:} Haryo Winarso, Delik Hudalah, dan Tommy Firman.
\item
  \textbf{Afiliasi:} School of Architecture, Planning and Policy
  Development, Bandung Institute of Technology, Indonesia.
\item
  \textbf{Penulis korespondensi:} Haryo Winarso.
\item
  \textbf{Jenis dokumen:} Artikel jurnal. Status penelaahan sejawat:
  Tidak disebutkan dalam file.
\item
  \textbf{Jurnal:} \emph{Habitat International}.
\item
  \textbf{Volume, tahun, dan halaman:} Volume 49, 2015, hlm. 221-229.
  Nomor terbitan: Tidak disebutkan dalam file.
\item
  \textbf{DOI:} 10.1016/j.habitatint.2015.05.024.
\item
  \textbf{ISSN:} 0197-3975.
\item
  \textbf{Riwayat artikel:} Diterima 25 Februari 2015; naskah revisi
  diterima 4 Mei 2015; disetujui 22 Mei 2015; tersedia secara daring 7
  Juni 2015 (hlm. 221, blok \emph{Article history}).
\item
  \textbf{Kata kunci:} \emph{peri-urbanisation}; Jakarta; \emph{large
  scale land development}; \emph{segregation}; \emph{socio-economic
  change} (hlm. 221, blok \emph{Keywords}).
\item
  \textbf{Wilayah studi:} Jakarta Metropolitan Area (JMA/Jabodetabek),
  dengan analisis kasus pada empat kelurahan di sekitar Bumi Serpong
  Damai (BSD), Kecamatan Serpong, Kabupaten Tangerang (hlm. 221-223 dan
  225-226, Bagian 1 dan 3.1-3.2).
\item
  \textbf{Pendanaan:} Sebagian artikel didasarkan pada studi tahun 2007
  yang didanai ITB Research Grant (hlm. 228, bagian
  \emph{Acknowledgement}).
\item
  \textbf{Konflik kepentingan:} Tidak disebutkan dalam file.
\end{itemize}

\subsection{Penanda Epistemik}\label{penanda-epistemik}

\begin{itemize}
\tightlist
\item
  \textbf{Laporan sumber} adalah data, metode, hasil, atau pernyataan
  yang disampaikan secara eksplisit dalam artikel.
\item
  \textbf{Interpretasi penulis} adalah penjelasan makna, mekanisme, atau
  implikasi yang diberikan penulis atas data dan hasilnya.
\item
  \textbf{Inferensi review} adalah pembacaan analitis yang diturunkan
  hanya dari TXT J27. Inferensi ini bukan klaim eksplisit artikel dan
  bukan bukti tambahan.
\end{itemize}

\subsection{Summarized Abstract}\label{summarized-abstract-15}

\textbf{Laporan sumber.} Artikel membahas transformasi kawasan periurban
JMA yang terutama berlangsung melalui pembangunan lahan skala besar
sejak awal 1990-an. Penulis menelusuri perpindahan batas periurban ke
arah luar serta perubahan demografi dan ekonomi melalui data statistik
deret waktu. Untuk melihat segregasi secara lebih tajam, artikel memakai
BSD sebagai kasus ekstrem dan membandingkan kondisi kota baru tersebut
dengan empat kelurahan di sekitarnya melalui survei rumah tangga
(Abstract, hlm. 221; Bagian 1, hlm. 222-223; Bagian 3-3.2, hlm.
224-226).

\textbf{Laporan sumber.} Transformasi yang dilaporkan meliputi kenaikan
proporsi migran, peralihan pekerjaan dari sektor primer menuju sektor
sekunder dan terutama tersier, penurunan pengangguran, kemunculan usaha
kecil kreatif sebagai strategi penghidupan, serta kenaikan pendapatan
nominal rumah tangga. Pembangunan kota baru juga membentuk pusat
kegiatan di luar Jakarta melalui penyediaan perumahan, perkantoran,
pusat belanja, dan fasilitas lainnya (Abstract, hlm. 221; Bagian
3-3.2.3, hlm. 224-227; Bagian 3.3, hlm. 227-228).

\textbf{Interpretasi penulis.} Periurbanisasi JMA dipahami sebagai
proses yang menghasilkan pekerjaan dan inovasi akar rumput, tetapi
sekaligus membentuk pembangunan yang tidak merata, enklave berfasilitas
tinggi, dan segregasi sosial-spasial. Kota-kota baru skala besar dinilai
terutama menguntungkan kelompok kaya, sedangkan kelompok miskin
menghadapi kesenjangan pendapatan dan eksternalitas negatif. Penulis
menggolongkan pola akhirnya sebagai \emph{self-segregation} atau
\emph{voluntary spatial segregation} dan memperingatkan potensi konflik
antarkomunitas (Abstract, hlm. 221; Bagian 3.3, hlm. 227-228; Bagian 4,
hlm. 228).

\textbf{Inferensi review.} Bukti paling langsung dalam file menunjukkan
koeksistensi perubahan struktur sosial-ekonomi dan pemisahan ruang pada
satu kasus terpilih. Bukti tersebut tidak cukup untuk mengestimasi
dampak kausal pembangunan BSD atau untuk mewakili seluruh JMA karena
penulis sendiri menyebut BSD sebagai kasus ekstrem yang tidak ditujukan
bagi generalisasi (Bagian 1, hlm. 223).

\subsection{Problem Statement}\label{problem-statement-15}

\textbf{Laporan sumber.} Periurbanisasi merupakan isu penting di negara
berkembang karena urbanisasi di perdesaan mempertemukan migrasi keluar
dari inti kota dengan migrasi penduduk perdesaan menuju kawasan
pinggiran. Kawasan transisi ini menghadapi bentang ruang dan kelembagaan
yang terfragmentasi, tekanan konversi lahan, kekurangan infrastruktur,
degradasi lingkungan, dan potensi konflik sosial sehingga menjadi
tantangan besar bagi perencana (Bagian 1, hlm. 221; Bagian 2.1, hlm.
223).

\textbf{Laporan sumber.} Di JMA, lebih dari 300.000 ha lahan perdesaan
disebut telah dikonversi sektor swasta menjadi kota-kota baru dalam tiga
dekade. Sejak 1990-an, kebijakan ekonomi yang lebih liberal, lemahnya
perencanaan dan regulasi, investasi jalan, serta pencarian lahan murah
mendorong pembangunan terencana dan tidak terencana semakin jauh ke
kawasan perdesaan, bahkan ke lokasi yang belum memiliki jalan, drainase,
dan listrik (Bagian 1, hlm. 221-222).

\textbf{Laporan sumber.} Masalah analitis artikel adalah bahwa
pergeseran spasial tersebut tidak hanya mengubah penggunaan lahan dan
lokasi penduduk, tetapi juga struktur migrasi, pekerjaan, pendapatan,
serta hubungan antara kota baru dan permukiman sekitar. Manfaat berupa
pusat kegiatan dan lapangan kerja dapat berlangsung bersamaan dengan
ketimpangan dan segregasi (Bagian 1, hlm. 222-223; Bagian 3.3, hlm.
227-228).

\textbf{Rumusan masalah eksplisit dalam bentuk kalimat tanya:} Tidak
disebutkan dalam file. \textbf{Inferensi review:} Dari pernyataan tujuan
dan argumen, persoalan yang diuji adalah bagaimana kawasan periurban JMA
bergerak keluar, bagaimana struktur sosial-ekonominya berubah, dan
bagaimana pembangunan lahan skala besar berkaitan dengan segregasi
sosial-spasial (Bagian 1, hlm. 222-223; Bagian 4, hlm. 228).

\subsection{Objectives}\label{objectives-15}

\begin{itemize}
\tightlist
\item
  \textbf{Laporan sumber:} Menunjukkan bahwa kawasan periurban JMA
  bergerak ke arah luar dan bahwa struktur sosial-ekonominya berubah
  dari kegiatan yang didominasi karakter perdesaan menjadi kegiatan yang
  lebih urban (Bagian 1, hlm. 222; Bagian 4, hlm. 228).
\item
  \textbf{Laporan sumber:} Membahas fenomena periurbanisasi JMA terkini
  dengan fokus pada transformasi sosial-ekonomi dan segregasi spasial
  yang menyertainya (Bagian 1, hlm. 222-223).
\item
  \textbf{Laporan sumber:} Menggunakan data statistik deret waktu untuk
  memperlihatkan perubahan demografi dan ekonomi yang memungkinkan
  periurbanisasi berlangsung selama beberapa dekade (Bagian 1, hlm.
  223).
\item
  \textbf{Laporan sumber:} Menggunakan empat kelurahan di sekitar BSD
  untuk memperlihatkan perubahan pola migrasi, struktur pekerjaan, dan
  pendapatan rumah tangga serta membandingkan kondisi permukiman sekitar
  dengan kondisi sosial-ekonomi kota baru (Bagian 1, hlm. 223; Bagian
  3.1-3.2.3, hlm. 225-227).
\item
  \textbf{Laporan sumber:} Mengungkap bagaimana segregasi spasial
  berlangsung pada kasus yang sengaja dipilih sebagai kasus ekstrem,
  bukan menghasilkan generalisasi statistik bagi seluruh periurban JMA
  (Bagian 1, hlm. 223).
\item
  \textbf{Hipotesis formal dan kriteria keberhasilan yang ditetapkan
  sebelum analisis:} Tidak disebutkan dalam file.
\end{itemize}

\subsection{Literature Survey}\label{literature-survey-15}

\textbf{Sifat survei.} Artikel menggunakan tinjauan naratif untuk
membangun pengertian periurban dan segregasi serta menempatkan
transformasi JMA dalam literatur metropolitanisasi. Strategi pencarian,
pangkalan data, kata kunci, rentang penelusuran, kriteria
inklusi-eksklusi, jumlah studi yang disaring, dan penilaian mutu
literatur: Tidak disebutkan dalam file. Oleh sebab itu, bagian ini tidak
dapat diperlakukan sebagai tinjauan sistematis.

\begin{itemize}
\tightlist
\item
  \textbf{Laporan sumber atas literatur periurban:} Tidak ada
  kesepakatan tunggal mengenai definisi kawasan periurban. Ciri yang
  umumnya disepakati ialah kedekatan dengan kawasan terbangun, pengaruh
  kegiatan urban yang kuat, akses ke pasar dan fasilitas kota, arus
  komuter dua arah, batas yang tidak kaku, campuran penggunaan lahan,
  pasar tanah yang dinamis, serta campuran tenaga kerja dan penduduk.
  Artikel juga mencatat bahwa periurban negara berkembang kerap dihuni
  migran miskin yang menjalankan beberapa kegiatan pendapatan informal
  dan memiliki karakter berbeda dari periurban negara maju (Bagian 2.1,
  hlm. 223).
\item
  \textbf{Laporan sumber atas literatur transisi:} Dalam negara
  berkembang yang kekurangan perencanaan dan regulasi, kawasan periurban
  mengalami tekanan pembangunan serta peralihan dari kegiatan perdesaan
  menuju campuran kegiatan dan penggunaan lahan perdesaan-perkotaan.
  Sejumlah karya memahaminya sebagai metropolitanisasi, yakni perluasan
  pasar komoditas dan tenaga kerja metropolitan ke rumah tangga
  pertanian (Bagian 2.1, hlm. 223).
\item
  \textbf{Laporan sumber atas literatur skala ruang:} Pertumbuhan kota
  mendorong periurban bergerak keluar secara bergelombang dan
  meninggalkan sprawl. Artikel melaporkan bahwa jangkauan periurban Asia
  dapat mencapai 300 km pada kasus pesisir Tiongkok, lebih jauh daripada
  rentang 30-50 km yang disebut untuk Afrika (Bagian 2.1, hlm. 223).
  Angka tersebut merupakan laporan J27 atas literatur, bukan hasil
  pengukuran kasus Jakarta.
\item
  \textbf{Laporan sumber atas literatur perubahan ekonomi:} Kajian Asia
  Timur menggambarkan periurban sebagai tempat transformasi dari ekonomi
  berbasis pertanian menuju industri dan jasa, didorong manufaktur dan
  modal asing, disertai pertumbuhan penduduk cepat, peningkatan
  kebutuhan tanah urban, dan kenaikan harga lahan (Bagian 2.1, hlm.
  223).
\item
  \textbf{Laporan sumber atas literatur segregasi:} Segregasi spasial
  dijelaskan sebagai distribusi kelompok sosial yang tidak homogen dan
  cenderung terkonsentrasi menurut status, etnisitas, atau asal. Konsep
  tersebut mencakup diferensiasi hunian, pembagian sosial ruang,
  kedekatan dan jarak antarkelas, serta lima dimensi yang dirujuk dari
  Massey dan Denton: kesetaraan/ketidaksetaraan, isolasi/eksposur,
  konsentrasi, sentralisasi, dan pengelompokan. Kondisi fisik hunian dan
  lingkungan juga dipakai dalam literatur yang dirujuk sebagai penanda
  selera, gaya hidup, dan posisi sosial (Bagian 2.2, hlm. 223).
\item
  \textbf{Laporan sumber atas literatur dampak segregasi:} Segregasi
  dapat menghambat akses kelompok berpendapatan rendah terhadap tempat
  kerja, layanan, dan peluang. Dampaknya dinilai lebih berat jika
  berlangsung intensif pada skala geografis yang melampaui jarak
  berjalan kaki dan membatasi interaksi antarkelas; pada skala kecil,
  dampaknya dapat lebih terbatas (Bagian 2.2, hlm. 223-224).
\item
  \textbf{Laporan sumber atas konteks JMA:} Literatur yang dirujuk
  menghubungkan pembangunan JMA dengan pertumbuhan ekonomi, investasi
  asing, liberalisasi era Suharto, peningkatan peran pengembang swasta,
  jalan antarkota, konversi lahan, kota baru, bentuk kota yang sekaligus
  terencana dan tidak terencana, serta segregasi pada pengembangan
  berpagar (Bagian 1, hlm. 221-222; Bagian 3-3.3, hlm. 224-228).
\end{itemize}

\textbf{Interpretasi penulis.} Kerangka tersebut mengarahkan pembacaan
JMA sebagai periurbanisasi bercorak kapitalistik: akumulasi modal
melalui pembangunan skala besar menghasilkan pusat kegiatan baru dan
pekerjaan, tetapi akses terhadap hunian serta fasilitas bermutu
terkonsentrasi pada kelompok yang mampu membayar (Bagian 1, hlm. 222;
Bagian 3.3-4, hlm. 227-228).

\textbf{Inferensi review.} Artikel menggabungkan beberapa tradisi
konsep, tetapi tidak mengoperasionalkan lima dimensi segregasi yang
ditinjau menjadi indeks formal. Literatur berfungsi sebagai lensa
penafsiran atas perbedaan sosial-ekonomi, bentuk enklave, dan contoh
konflik, bukan sebagai dasar pengukuran segregasi multidimensi yang
lengkap.

\subsection{Method Used}\label{method-used-15}

\textbf{Laporan sumber.} Artikel menggabungkan analisis statistik
deskriptif deret waktu dengan studi kasus rumah tangga. Data statistik
digunakan untuk menunjukkan perubahan demografi dan ekonomi JMA selama
beberapa dekade. Empat kelurahan di sebelah BSD dianalisis untuk
menjelaskan transformasi sosial-ekonomi dan membandingkan permukiman
sekitar dengan kota baru (Bagian 1, hlm. 223; Bagian 3-3.2.3, hlm.
224-227).

\textbf{Inferensi review atas rancangan.} Kombinasi survei rumah tangga,
data sekunder, observasi koridor, dan wawancara dapat disebut rancangan
multisumber yang dominan deskriptif. Artikel tidak menyebut desainnya
sebagai \emph{mixed methods}, tidak menyajikan model identifikasi
kausal, dan tidak membentuk kelompok kontrol sebelum-sesudah.

\textbf{Tahap analisis makro JMA:}

\begin{itemize}
\tightlist
\item
  Populasi inti DKI Jakarta dan Bodetabek dibandingkan pada sensus 1961,
  1971, 1980, 1990, 2000, dan 2010 (Tabel 1, hlm. 224).
\item
  Perubahan kepadatan penduduk inti, zona dalam, dan zona luar dibaca
  dari Gambar 2. Tahun-tahun pengamatan rinci pada gambar tidak terbaca
  sebagai angka dalam TXT, selain rujukan naratif pada 1987 dan
  perkembangan sejak awal 1990-an (Bagian 3, hlm. 224; Gambar 2).
\item
  Kontribusi sektor primer, sekunder, dan tersier dalam produk domestik
  regional bruto dibandingkan untuk 1985, 1990, 2000, dan 2005 (Gambar
  3, hlm. 225). Nilai numerik setiap sektor tidak tersedia dalam teks
  hasil ekstraksi.
\item
  Uraian historis digunakan untuk menelusuri Jabotabek Development Plan
  1974, pembangunan jalan menuju Tangerang, Bekasi, dan Bogor,
  meningkatnya peran pengembang swasta, perluasan kota baru, serta
  pergeseran jarak kawasan periurban dari pusat Jakarta (Bagian 3, hlm.
  224).
\end{itemize}

\textbf{Pemilihan kasus dan sampel rumah tangga:}

\begin{itemize}
\tightlist
\item
  BSD dipilih karena disebut sebagai kota baru terbesar dan paling
  dinamis yang pernah dibangun di kawasan periurban Jakarta. Penulis
  menegaskan bahwa BSD tidak mewakili seluruh periurban JMA dan
  diperlakukan sebagai kasus ekstrem untuk memperjelas proses segregasi
  (Bagian 1, hlm. 223).
\item
  Empat kelurahan yang dipilih ialah Rawa Mekar Jaya, Rawabuntu,
  Jelupang, dan Cilenggang. Alasan pemilihannya ialah masih memiliki
  ciri perdesaan, berjarak sekitar 15-25 menit dari BSD City, dan berada
  di kawasan periurban Jakarta (Bagian 3.2, hlm. 225).
\item
  Penulis menerapkan \emph{proportionate random sampling} kepada 201
  rumah tangga. Di lapangan, 220 nama dipilih secara acak dari kerangka
  sampel berupa nama rumah tangga yang diperoleh dari setiap kelurahan.
  Tabel 2 melaporkan 201 responden: 52 di Rawa Mekar Jaya, 61 di
  Rawabuntu, 56 di Jelupang, dan 32 pada baris yang tertulis
  ``Cilempang'' (Bagian 3.2 dan Tabel 2, hlm. 225). Hubungan antara 220
  nama dan 201 responden akhir serta perbedaan nama Cilenggang/Cilempang
  tidak dijelaskan dalam file.
\item
  Kuesioner digunakan untuk memperoleh informasi transformasi
  sosial-ekonomi yang mencakup pola migrasi, struktur pekerjaan, dan
  pendapatan (Bagian 3.2, hlm. 225).
\item
  Responden dikelompokkan menurut status sosial-ekonomi, kemudian data
  survei permukiman informal dibandingkan dengan informasi
  sosial-ekonomi kota baru untuk menggambarkan segregasi antara
  permukiman sekitar dan BSD (Bagian 1, hlm. 223). Definisi kategori
  status sosial-ekonomi dan prosedur pencocokan kedua kelompok: Tidak
  disebutkan dalam file.
\end{itemize}

\textbf{Observasi dan wawancara pelengkap:}

\begin{itemize}
\tightlist
\item
  Survei pengamatan dari kendaraan (\emph{windshield survey}) dilakukan
  pada 2010 di sepanjang jalan utama dekat BSD. Peneliti menghitung 191
  kegiatan perdagangan dan jasa formal maupun informal (Bagian 3.2.2,
  hlm. 226).
\item
  Setelah inspeksi dan wawancara mendalam terhadap pelaku perdagangan
  dan jasa terpilih, 20 usaha digolongkan sebagai usaha kecil kreatif
  dan inovatif yang dijalankan individu atau keluarga untuk
  mempertahankan penghidupan (Bagian 3.2.2, hlm. 226). Jumlah pihak yang
  diwawancarai, pedoman wawancara, cara pemilihan, kriteria klasifikasi,
  perekaman, dan prosedur analisis kualitatif: Tidak disebutkan dalam
  file.
\end{itemize}

\textbf{Teknik analisis statistik.} Artikel menampilkan jumlah,
proporsi, laju pertumbuhan, kategori pekerjaan, kelompok pendapatan, dan
perbandingan lintas tahun. Uji hipotesis, regresi, galat baku, interval
kepercayaan, uji signifikansi, pembobotan survei, penanganan data
hilang, serta perangkat lunak analisis: Tidak disebutkan dalam file.

\subsection{Dataset}\label{dataset-15}

\begin{itemize}
\tightlist
\item
  \textbf{Statistik penduduk JMA:} Tabel 1 memuat populasi DKI Jakarta,
  Bodetabek, dan total wilayah metropolitan untuk enam tahun sensus,
  1961-2010, dengan sumber BPS (2010) (Tabel 1, hlm. 224).
\item
  \textbf{Kepadatan penduduk:} Gambar 2 membandingkan kepadatan inti,
  zona dalam, dan kawasan periurban. Nama himpunan data, tahun lengkap,
  unit kepadatan, serta metode pembentukan zona: Tidak disebutkan dalam
  teks yang tersedia (Gambar 2 dan Bagian 3, hlm. 224).
\item
  \textbf{Struktur ekonomi regional:} Gambar 3 menggunakan analisis
  berdasarkan Central Statistics Bureau (2009) untuk membandingkan
  kontribusi sektor primer, sekunder, dan tersier dalam GDRP JMA pada
  1985, 1990, 2000, dan 2005 (Gambar 3, hlm. 225). Nilai angka di dalam
  grafik tidak terekstraksi ke TXT.
\item
  \textbf{Data administrasi lokal:} BPS (2007) digunakan untuk
  menyatakan jumlah penduduk gabungan empat kelurahan kasus pada 2006,
  sedangkan BPS (2010) dipakai untuk populasi dan laju pertumbuhan JMA
  serta wilayah periurban (Bagian 3 dan 3.2, hlm. 224-225).
\item
  \textbf{Survei rumah tangga lapangan:} Data kuesioner tahun 2006
  mencakup migrasi, jenis pekerjaan, pekerjaan utama dan sampingan,
  serta pendapatan rumah tangga pada empat kelurahan sekitar BSD. Gambar
  5 secara eksplisit menyebut \emph{Field survey 2006} sebagai sumber
  (Bagian 3.2-3.2.3 dan Gambar 5, hlm. 225-227).
\item
  \textbf{Data pembanding 1991:} Gambar 5-7 membandingkan pekerjaan dan
  pendapatan 1991 dengan 2006. Sumber khusus data 1991, cara memperoleh
  ingatan historis, kesetaraan instrumen, dan apakah responden yang sama
  diamati pada kedua tahun: Tidak disebutkan dalam file (Bagian
  3.2.2-3.2.3 dan Gambar 5-7, hlm. 226-227).
\item
  \textbf{Observasi usaha 2010:} Survei koridor menginventarisasi 191
  kegiatan perdagangan dan jasa formal maupun informal, dengan inspeksi
  serta wawancara lanjutan pada usaha terpilih (Bagian 3.2.2, hlm. 226).
\item
  \textbf{Data BSD:} Artikel menggunakan Harmanurjeni (2006) untuk
  keterangan pendapatan penghuni BSD, BSD City (2009) untuk luas
  terbangun, unit hunian, dan jumlah penduduk, serta Leisch (2002) dan
  Winarso dan Sari (2007) untuk profil sosial-ekonomi penghuni kota baru
  (Bagian 3.1, hlm. 225). Data tersebut diperlakukan sebagai data
  sekunder yang dilaporkan J27; sumber aslinya tidak diperiksa dalam
  review ini.
\item
  \textbf{Data pendapatan pembanding:} Artikel merujuk BPS (1994) untuk
  angka rata-rata pendapatan pekerja di Jakarta dan Jawa Barat serta BPS
  (2006) untuk pendapatan rumah tangga urban Indonesia (catatan kaki
  4-5, hlm. 225). Catatan kaki 4 menyatakan inflasi tidak dihitung pada
  perbandingan yang memakai angka BPS (1994).
\item
  \textbf{Harga hunian dan fasilitas kota baru:} Tabel 3 memuat lima
  pengembangan di Kabupaten Tangerang pada 2006, berikut tipe
  bangunan/lahan, harga, dan fasilitas. Sumber tabel, cara pemilihan
  proyek, tanggal harga, serta apakah harga merupakan harga penawaran
  atau transaksi: Tidak disebutkan dalam file (Tabel 3, hlm. 228).
\item
  \textbf{Bukti kejadian konflik:} Artikel mengutip laporan \emph{Tempo
  Interaktif} (2003) mengenai tenaga bongkar-muat berbayar ilegal dan
  kolom jurnalisme warga Barutu (2012) mengenai pemblokiran akses
  kompleks perumahan berpendapatan tinggi (Bagian 3.3, hlm. 228).
  Keduanya merupakan contoh yang dilaporkan artikel, bukan hasil survei
  rumah tangga J27.
\item
  \textbf{Ketersediaan data:} Nama berkas, kamus variabel, data mentah,
  transkrip, repositori, lisensi, kode analisis, dan prosedur akses
  ulang: Tidak disebutkan dalam file.
\end{itemize}

\subsection{Population / Sample}\label{population-sample-15}

\textbf{Cakupan makro.} JMA atau Jabodetabek dijelaskan sebagai Jakarta
beserta wilayah Bogor, Depok, Tangerang, dan Bekasi. Tabel 1
mengoperasionalkan inti sebagai DKI Jakarta dan periurban sebagai
Bodetabek (Bagian 1, hlm. 221; Tabel 1 dan Gambar 2, hlm. 224). Artikel
juga memakai singkatan JMR pada Tabel 1 dan Gambar 2 tanpa menjelaskan
perbedaannya dari JMA.

\textbf{Populasi regional.} Tabel 1 melaporkan total 27.447.027 penduduk
pada 2010, terdiri atas 9.607.787 di DKI Jakarta dan 17.839.240 di
Bodetabek (Tabel 1, hlm. 224). Bagian pendahuluan menyebut populasi JMA
lebih dari 30 juta dengan sitasi BPS (2010), tetapi tidak menjelaskan
perbedaannya dari total pada tabel (Bagian 1, hlm. 221).

\textbf{Populasi kasus.} Kecamatan Serpong disebut memiliki luas
4.599,798 ha dan terdiri atas 16 desa/kelurahan. Empat kelurahan kasus
memiliki jumlah penduduk gabungan 54.250 pada 2006 (Bagian 3.2, hlm.
225).

\textbf{Sampel rumah tangga.} Sampel analitis yang ditampilkan berjumlah
201 rumah tangga, dengan alokasi 52, 61, 56, dan 32 responden. Penulis
juga menyatakan bahwa 220 nama dipilih secara acak dari kerangka nama
rumah tangga. Hubungan antara 220 nama dan 201 responden, tingkat
respons, kriteria inklusi-eksklusi, jumlah rumah tangga dalam kerangka,
bobot, karakteristik nonresponden, dan perhitungan ukuran sampel: Tidak
disebutkan dalam file (Bagian 3.2 dan Tabel 2, hlm. 225).

\textbf{Populasi pembanding BSD.} BSD dilaporkan memiliki sekitar
100.000 penduduk dan 25.000 unit rumah pada area 1.500 ha yang telah
terbangun pada 2009. Angka ini adalah profil kota baru, bukan sampel
rumah tangga yang diwawancarai dalam J27 (Bagian 3.1, hlm. 225). Jumlah
observasi individu atau rumah tangga BSD yang menjadi dasar data
pembanding: Tidak disebutkan dalam file.

\textbf{Sampel usaha.} Sebanyak 191 kegiatan perdagangan dan jasa
ditemukan melalui pengamatan 2010. Setelah inspeksi dan wawancara
terhadap pelaku terpilih, 20 usaha digolongkan sebagai usaha kecil
kreatif dan inovatif. Jumlah keseluruhan usaha yang diwawancarai, cara
memilihnya dari kegiatan yang diamati, dan tingkat partisipasinya: Tidak
disebutkan dalam file (Bagian 3.2.2, hlm. 226).

\subsection{Dependent Variables}\label{dependent-variables-15}

\textbf{Klarifikasi desain.} Artikel tidak menetapkan satu variabel
dependen dalam model regresi atau eksperimen. Transformasi periurban
merupakan konstruk luas yang dibaca melalui beberapa keluaran deskriptif
berikut.

\begin{itemize}
\tightlist
\item
  \textbf{Perluasan spasial periurban:} jarak perkembangan dari pusat
  Jakarta dan urutan fase pembangunan permukiman, koridor jalan, serta
  proyek lahan skala besar (Bagian 1 dan 3, hlm. 222 dan 224; Bagian 4,
  hlm. 228).
\item
  \textbf{Perubahan demografi:} jumlah penduduk, laju pertumbuhan,
  kepadatan menurut zona, proporsi migran, tahun kedatangan, dan daerah
  asal migran (Tabel 1 dan Gambar 2, hlm. 224; Bagian 3.2.1 dan Gambar
  4, hlm. 225-226).
\item
  \textbf{Perubahan struktur ekonomi:} kontribusi sektor primer,
  sekunder, dan tersier terhadap GDRP JMA (Gambar 3, hlm. 225).
\item
  \textbf{Perubahan pekerjaan rumah tangga:} jumlah atau proporsi
  pencari nafkah menurut jenis pekerjaan, peralihan pekerjaan
  antarsektor, pengangguran, usaha sewa kamar, serta kegiatan
  perdagangan dan jasa (Bagian 3.2.2 dan Gambar 5, hlm. 226-227).
\item
  \textbf{Perubahan pendapatan:} rentang dan modus pendapatan utama,
  kelompok pendapatan utama, serta kelompok pendapatan sampingan pada
  1991 dan 2006 (Bagian 3.2.3 dan Gambar 6-7, hlm. 226-227).
\item
  \textbf{Pemisahan sosial-spasial:} perbedaan pendapatan, tipe dan
  harga hunian, fasilitas, serta kondisi lingkungan antara kota baru dan
  permukiman di sekitarnya. Artikel tidak menghitung indeks segregasi,
  jarak antarkelompok, atau ukuran konsentrasi spasial formal (Bagian 1,
  hlm. 223; Bagian 3.1 dan 3.3, hlm. 225 dan 227-228).
\item
  \textbf{Inovasi penghidupan:} jumlah kegiatan perdagangan/jasa dan
  jumlah usaha yang diklasifikasikan kreatif serta inovatif (Bagian
  3.2.2, hlm. 226).
\item
  \textbf{Potensi konflik:} contoh pemaksaan penggunaan tenaga kerja dan
  pemblokiran akses kompleks dipakai sebagai indikasi ketegangan.
  Frekuensi konflik, korban, kerugian, persepsi penduduk yang diukur
  sistematis, atau variabel konflik formal: Tidak disebutkan dalam file
  (Bagian 3.3, hlm. 227-228).
\end{itemize}

\textbf{Inferensi review.} Pembangunan lahan skala besar, pertumbuhan
jalan, dan waktu berfungsi sebagai konteks penjelas dalam narasi, bukan
sebagai variabel perlakuan yang efeknya diestimasi. Karena itu,
perubahan keluaran tidak dapat langsung dibaca sebagai dampak kausal BSD
atau pengembang swasta.

\subsection{Results}\label{results-15}

\subsubsection{Perluasan JMA dan perubahan
regional}\label{perluasan-jma-dan-perubahan-regional}

\begin{itemize}
\tightlist
\item
  \textbf{Laporan sumber:} Populasi DKI Jakarta/Bodetabek/total wilayah
  metropolitan pada 1961 adalah 2.904.533/2.794.712/5.699.245; pada 1971
  sebesar 4.546.492/3.483.537/8.030.029; pada 1980 sebesar
  6.503.449/5.413.271/11.916.720; pada 1990 sebesar
  7.259.257/8.878.256/16.137.513; pada 2000 sebesar
  8.347.083/12.842.626/21.189.709; dan pada 2010 sebesar
  9.607.787/17.839.240/27.447.027 (Tabel 1, hlm. 224).
\item
  \textbf{Laporan sumber:} Jakarta bertambah dari sekitar 2,9 juta
  penduduk pada 1961 menjadi 4,6 juta pada 1971, dengan laju pertumbuhan
  tahunan 5,8\%. Untuk 2010, JMA dilaporkan tumbuh 2,05\% per tahun;
  Kabupaten Bogor 2,45\%; Kabupaten Tangerang 4,02\%; dan Kabupaten
  Bekasi 3,37\%. Angka Tangerang dihitung dengan menggabungkan Kota
  Tangerang Selatan dan Kabupaten Tangerang setelah pemekaran 2009
  (Bagian 3, hlm. 224; catatan kaki 2).
\item
  \textbf{Interpretasi penulis:} Pertumbuhan Bodetabek yang lebih cepat
  daripada DKI Jakarta, penurunan rata-rata kepadatan inti, serta
  kenaikan kepadatan zona dalam dan luar dibaca sebagai perpindahan
  penduduk dari inti menuju pinggiran dan sebagai pengaruh pembangunan
  kota baru sejak 1990 (Bagian 3 dan Gambar 2, hlm. 224).
\item
  \textbf{Laporan sumber:} Nama Jabotabek disebut berasal dari konsep
  yang dirumuskan konsultan Belanda pada 1970. Jabotabek Development
  Plan kemudian diperkenalkan pada 1974 dan berfokus pada pembangunan
  jalan yang menghubungkan Tangerang di barat, Bekasi di timur, dan
  Bogor di selatan. Kawasan periurban disebut berjarak 10-15 km dari
  pusat pada masa tersebut, sekitar 20 km pada 1980, dan pembangunan
  fisik mencapai 30-45 km pada akhir 1980-an dan awal 1990-an (catatan
  kaki 1, hlm. 221; Bagian 3, hlm. 224).
\item
  \textbf{Laporan sumber:} Pada era Orde Baru, kebijakan perumahan
  bergeser ke penyediaan berbasis pasar oleh sektor swasta. Pada 1987,
  jumlah rumah formal yang dibangun pengembang swasta disebut telah
  melampaui produksi perusahaan perumahan milik pemerintah (Bagian 3,
  hlm. 224).
\item
  \textbf{Laporan sumber:} Kontribusi sektor primer terhadap GDRP
  menurun, sedangkan sektor sekunder dan tersier tumbuh dalam
  perbandingan 1985, 1990, 2000, dan 2005. Sektor jasa dalam catatan
  artikel mencakup perdagangan, restoran dan hotel, transportasi dan
  komunikasi, penyewaan dan jasa perusahaan, serta jasa lain. Penulis
  menyatakan puncak pertumbuhan metropolitan terjadi pada 1990-1995,
  ketika investasi asing bertambah dan perusahaan asing membuka kantor
  cabang di JMA (Bagian 3 dan Gambar 3, hlm. 224-225; catatan kaki 3).
\item
  \textbf{Interpretasi penulis:} Pergeseran sektoral tersebut dipandang
  sebagai fase kematangan metropolitan. Penurunan sektor primer terutama
  dijelaskan melalui penyusutan lahan perdesaan akibat konversi oleh
  pengembang, sedangkan kota-kota baru membentuk pusat kegiatan di luar
  Jakarta (Bagian 3, hlm. 224-225).
\end{itemize}

\subsubsection{Profil BSD dan permukiman
sekitar}\label{profil-bsd-dan-permukiman-sekitar}

\begin{itemize}
\tightlist
\item
  \textbf{Laporan sumber:} BSD terletak sekitar 30 km di barat daya
  pusat Jakarta. Kota baru ini direncanakan dalam tiga tahap pada area
  6.000 ha untuk menampung 600.000 orang pada 2035. Pada 2009, tahap
  pertama seluas 1.500 ha telah dikembangkan dengan 25.000 rumah dan
  sekitar 100.000 penghuni (Bagian 3.1, hlm. 225).
\item
  \textbf{Laporan sumber:} Artikel menyatakan lebih dari 90\% penghuni
  BSD berpendapatan harian di atas Rp33.000 atau ``US\$3300'', yang
  kemudian ditulis setara sekitar Rp825.000 atau ``US\$82,500'' per
  bulan. Artikel juga mencantumkan rata-rata pendapatan pekerja Jakarta
  Rp255.463 (US\$25,5), Jawa Barat Rp151.618 (US\$15,1), modus bulanan
  BSD Rp1.000.000 (US\$100), dan pendapatan bulanan seluruh rumah tangga
  urban Indonesia Rp1.250.000 (US\$137) (Bagian 3.1 dan catatan kaki
  4-5, hlm. 225).
\item
  \textbf{Batas pelaporan:} Angka dolar, periodisasi, dan klaim bahwa
  Rp1.000.000 ``jauh di atas'' Rp1.250.000 tidak konsisten secara
  internal. Review ini mempertahankan angka sesuai TXT dan tidak
  membetulkannya melalui dugaan atau sumber luar.
\item
  \textbf{Laporan sumber atas data sekunder:} Komposisi usia, ukuran
  keluarga, pendidikan, pekerjaan, dan pendapatan penghuni kota baru
  ditafsirkan sebagai profil profesional muda di sektor swasta. Konsumen
  perumahan disebut terutama berasal dari kelompok menengah dan
  berpendapatan tinggi serta keturunan Tionghoa Indonesia (Bagian 3.1,
  hlm. 225). Angka rinci untuk karakteristik tersebut: Tidak disebutkan
  dalam file.
\end{itemize}

\subsubsection{Migrasi}\label{migrasi}

\begin{itemize}
\tightlist
\item
  \textbf{Laporan sumber:} Data yang dikumpulkan pada 2006 menunjukkan
  bahwa 52\% populasi dalam wilayah studi dikategorikan sebagai migran
  yang datang setelah pengembangan BSD pada awal 1990-an, sedangkan 48\%
  dikategorikan ``indigenous'' dan disebut telah tinggal di bekas
  kawasan perkebunan karet sejak 1967 (Bagian 3.2.1, hlm. 225-226).
  Proporsi ini berasal dari analisis kasus berbasis sampel rumah tangga,
  bukan pencacahan seluruh 54.250 penduduk.
\item
  \textbf{Laporan sumber:} Asal migran yang diberi angka ialah Kotamadya
  Jakarta 24,7\%, Kabupaten Tangerang 27,1\%, Provinsi Jawa Barat
  10,3\%, Provinsi Jawa Tengah 14,0\%, dan Sumatra 4,7\%. Migran dari
  Kalimantan dan Sulawesi juga disebut, tetapi proporsinya tidak
  tertulis dalam narasi TXT (Bagian 3.2.1 dan Gambar 4, hlm. 226).
\item
  \textbf{Interpretasi penulis:} Arus dari inti dan zona dalam JMA
  dipahami sebagai suburbanisasi, sedangkan arus dari wilayah perdesaan
  terdekat dipahami sebagai migrasi desa-kota menuju pinggiran.
  Mayoritas migran disebut hanya pernah berpindah sekali, terutama dari
  Jakarta dan Tangerang, sehingga kawasan periurban dinilai menarik
  penduduk inti sekaligus, dalam jumlah lebih kecil, penduduk perdesaan
  sekitar (Bagian 3.2.1, hlm. 226).
\end{itemize}

\subsubsection{Pekerjaan dan usaha
lokal}\label{pekerjaan-dan-usaha-lokal}

\begin{itemize}
\tightlist
\item
  \textbf{Laporan sumber:} Dalam rentang 15 tahun, jumlah pencari nafkah
  rumah tangga yang bekerja sebagai buruh, pegawai swasta, pengusaha,
  dan pedagang meningkat, sedangkan petani berkurang. Usaha penyewaan
  kamar atau akomodasi sebagai pendapatan utama juga bertambah (Bagian
  3.2.2, hlm. 226).
\item
  \textbf{Laporan sumber:} Gambar 5 menunjukkan penurunan pekerjaan
  sektor primer serta pertumbuhan sektor sekunder dan tersier pada
  1991-2006. Pengangguran juga dilaporkan menurun (Bagian 3.2.2 dan
  Gambar 5, hlm. 226-227). Nilai angka tiap sektor dan pengangguran
  tidak tersedia dalam teks hasil ekstraksi.
\item
  \textbf{Interpretasi penulis:} Penurunan pengangguran dijelaskan oleh
  tersedianya pekerjaan tersier, terutama pekerjaan informal seperti
  pedagang kaki lima dan jasa perbaikan sepeda motor atau peralatan
  rumah tangga. Kegiatan tersebut dipandang sebagai inovasi kota dari
  akar rumput dan strategi penghidupan (Bagian 3.2.2, hlm. 226).
\item
  \textbf{Laporan sumber:} Pengamatan 2010 menemukan 191 kegiatan
  perdagangan dan jasa formal maupun informal di jalan utama dekat BSD.
  Inspeksi dan wawancara terhadap pelaku terpilih menghasilkan
  klasifikasi 20 usaha kecil kreatif dan inovatif yang dijalankan
  individu atau keluarga (Bagian 3.2.2, hlm. 226).
\end{itemize}

\subsubsection{Pendapatan rumah tangga sekitar
BSD}\label{pendapatan-rumah-tangga-sekitar-bsd}

\begin{itemize}
\tightlist
\item
  \textbf{Laporan sumber:} Pada 2006, pendapatan pekerjaan utama
  berkisar dari Rp500.000 sampai lebih dari Rp15.000.000 per bulan.
  Modusnya Rp1.250.000 per bulan dan disebut sebanding dengan rata-rata
  pendapatan rumah tangga urban Indonesia, sedangkan modus 1991 adalah
  Rp500.000 (Bagian 3.2.3, hlm. 226).
\item
  \textbf{Laporan sumber:} Lima kelompok pendapatan utama yang digunakan
  adalah kurang dari Rp1.250.000; Rp1.250.000-Rp2.500.000;
  Rp2.500.000-Rp5.000.000; Rp5.000.000-Rp10.000.000; dan lebih dari
  Rp10.000.000. Kelompok terbawah menurun, sedangkan kelompok lainnya
  meningkat, terutama kelompok kedua. Seluruh kelompok pendapatan
  sampingan juga meningkat dan kelompok kedua tumbuh paling besar
  (Bagian 3.2.3 dan Gambar 6-7, hlm. 226-227).
\item
  \textbf{Batas perbandingan yang dinyatakan sumber:} Inflasi rata-rata
  7,7\% per tahun tidak dimasukkan ke perbandingan pendapatan 1991-2006
  tersebut (catatan kaki 6, hlm. 226).
\item
  \textbf{Interpretasi penulis:} Meskipun pendapatan nominal berubah,
  tingkat pendapatan permukiman sekitar masih dinilai jauh di bawah BSD.
  Penulis menyatakan tidak ada perbaikan kesejahteraan yang berarti bagi
  penduduk lebih miskin, sedangkan penduduk lebih kaya meningkat cepat
  secara sosial-ekonomi (Bagian 3.2.3, hlm. 226).
\item
  \textbf{Inferensi review:} Karena perbandingan 1991-2006 tidak
  mengoreksi inflasi dan angka rinci grafik tidak tersedia, arah
  perubahan nominal dapat dilaporkan, tetapi besarnya perubahan
  pendapatan riil dan klaim kesejahteraan tidak dapat diuji ulang dari
  TXT.
\end{itemize}

\subsubsection{Fasilitas, segregasi, dan
konflik}\label{fasilitas-segregasi-dan-konflik}

\begin{itemize}
\tightlist
\item
  \textbf{Laporan sumber:} Kota-kota baru menyediakan lingkungan
  berstandar tinggi berupa lapangan golf, sekolah swasta mahal, rumah
  sakit mahal, ruang hijau, kantor, pusat belanja, dan amenitas utama.
  Ruang hijau disebut melebihi 10 m2 per orang dan dibandingkan dengan
  ``standar Eropa Barat'' 6 m2 per orang (Bagian 3.3, hlm. 227). Catatan
  kaki 7 justru mengutip pedoman Skotlandia sebesar 60 m2 ruang terbuka
  per rumah tangga; hubungan kedua standar tidak dijelaskan.
\item
  \textbf{Laporan sumber:} Tabel 3 menunjukkan rentang harga 2006
  berikut: Bintaro Jaya Rp96.184.880-Rp147.996.000; BSD
  Rp62.600.000-Rp373.200.000; Citra Raya Tangerang
  Rp147.280.000-Rp307.500.000; TigaRaksa Rp38.800.000-Rp128.800.000; dan
  LippoKarawaci Rp131.595.000-Rp167.000.000. Angka tipe bangunan/lahan
  yang dicantumkan berkisar dari komponen bangunan 37 sampai 232 dan
  komponen lahan 84 sampai 332; satuannya tidak ditulis dalam tabel
  hasil ekstraksi (Tabel 3, hlm. 228).
\item
  \textbf{Laporan sumber:} Fasilitas yang dicantumkan meliputi pusat
  komersial dan bisnis, gedung multifungsi, fasilitas olahraga, lapangan
  golf, \emph{country club}, sekolah dan universitas, fasilitas
  keagamaan, klinik dan rumah sakit, rekreasi, transportasi, pusat
  belanja, hotel, kantor, serta apartemen (Tabel 3, hlm. 228).
\item
  \textbf{Interpretasi penulis:} Pusat baru tersebut mengurangi tarikan
  sentripetal Jakarta dan perjalanan harian sebagian anggota rumah
  tangga. Namun, hanya kelompok mampu yang dapat membayar perjalanan,
  hunian, dan fasilitas tersebut, sedangkan kelompok miskin menanggung
  kemacetan serta polusi sebagai eksternalitas pembangunan yang tidak
  merata (Bagian 3.3, hlm. 227).
\item
  \textbf{Laporan sumber:} Artikel menyebut 180 unit rumah seharga Rp1
  miliar per unit terjual dalam satu tahun pada 2011 di salah satu kota
  baru, sementara kelompok miskin kesulitan mendapatkan tempat tinggal
  di dalam pengembangan tersebut (Bagian 3.3, hlm. 227-228). Nama kota
  baru dan sumber khusus angka penjualan: Tidak disebutkan dalam file.
\item
  \textbf{Laporan sumber:} Dua contoh ketegangan ialah tuntutan agar
  penghuni memakai tenaga bongkar-muat berbayar dengan tarif di atas
  normal dan pemblokiran akses menuju kompleks perumahan berpendapatan
  tinggi oleh kelompok miskin. Penulis menyatakan perilaku semacam itu
  tumbuh di sejumlah kawasan hunian kaya dan menimbulkan rasa tidak aman
  (Bagian 3.3, hlm. 227-228).
\item
  \textbf{Interpretasi penulis:} Pembangunan bercorak akumulasi modal
  dipandang memperkuat segregasi di kawasan periurban dan di dalam kota
  baru. Perubahan guna lahan menciptakan pemenang serta pihak yang
  menanggung kerugian, sehingga antagonisme kelas dapat berkembang
  menjadi konflik lokal, terutama ketika kelompok miskin tidak memiliki
  akses terhadap bantuan hukum (Bagian 3.3, hlm. 227-228).
\end{itemize}

\subsection{Findings}\label{findings-15}

\begin{itemize}
\tightlist
\item
  \textbf{Laporan sumber:} Kawasan periurban Jakarta bergerak semakin
  jauh dari inti. Perkembangannya digambarkan dalam tiga fase:
  penyebaran permukiman kecil tanpa pola yang tegas; pertumbuhan
  perdagangan dan jasa mengikuti investasi jalan antarkota pada 1980-an;
  serta pembangunan lahan skala besar yang menembus kawasan perdesaan
  sejak akhir 1980-an dan terutama 1990-an (Bagian 1, hlm. 221-222;
  Bagian 3 dan 4, hlm. 224 dan 228).
\item
  \textbf{Laporan sumber:} Periurban JMA tumbuh lebih cepat daripada
  inti, kepadatan zona luar meningkat, dan struktur GDRP bergeser dari
  sektor primer ke sektor sekunder serta tersier (Tabel 1 dan Gambar 2,
  hlm. 224; Gambar 3, hlm. 225).
\item
  \textbf{Interpretasi penulis:} Pengembangan kota baru menarik penduduk
  dari Jakarta dan zona dalam sebagai suburbanisasi sekaligus menarik
  penduduk wilayah perdesaan sekitar sebagai migrasi desa-kota. Kedua
  arus tersebut bersama-sama membentuk periurbanisasi (Bagian 3.2.1,
  hlm. 225-226; Bagian 4, hlm. 228).
\item
  \textbf{Laporan sumber:} Data sampel empat kelurahan menunjukkan
  proporsi migran sebesar 52\%; pekerjaan sektor primer menurun;
  pekerjaan sekunder dan tersier meningkat; pengangguran menurun; usaha
  perdagangan dan jasa bertambah; serta kelompok pendapatan nominal di
  atas kategori terendah meningkat (Bagian 3.2.1-3.2.3, hlm. 225-227;
  Gambar 4-7).
\item
  \textbf{Interpretasi penulis:} Pekerjaan informal dan usaha keluarga
  yang tumbuh di sekitar BSD dipahami sebagai inovasi akar rumput yang
  membantu keberlanjutan penghidupan. Dengan demikian, pembangunan skala
  besar bukan hanya memisahkan ruang, tetapi juga menciptakan pasar dan
  pekerjaan bagi masyarakat sekitar (Bagian 3.2.2, hlm. 226; Bagian 4,
  hlm. 228).
\item
  \textbf{Interpretasi penulis:} Perubahan pendapatan tidak menghapus
  kesenjangan. Hunian dan fasilitas bermutu tinggi di kota baru
  berorientasi pada kelompok menengah dan kaya, sedangkan kelompok
  miskin tetap berada di luar pengembangan dan menanggung sebagian
  eksternalitasnya (Bagian 3.1 dan 3.3, hlm. 225 dan 227-228).
\item
  \textbf{Interpretasi penulis:} Segregasi JMA diklasifikasikan sebagai
  \emph{self-segregation} atau \emph{voluntary spatial segregation}.
  Penulis menyamakannya secara konseptual dengan komunitas berpagar yang
  menjadi ``pulau kekayaan di lautan kemiskinan'' dan menghubungkannya
  dengan potensi konflik kelas (Bagian 4, hlm. 228).
\item
  \textbf{Inferensi review:} Temuan perubahan demografi, pekerjaan, dan
  pendapatan didukung oleh statistik serta survei deskriptif, sedangkan
  label segregasi sukarela dan mekanisme konflik terutama merupakan
  interpretasi teoretis atas perbandingan kondisi ruang, data sekunder
  kota baru, dan sejumlah contoh kejadian. Artikel tidak mengukur
  keputusan segregasi sebagai pilihan sukarela pada tingkat rumah
  tangga.
\item
  \textbf{Inferensi review:} Pernyataan bahwa pengembangan lahan
  menghasilkan transformasi sosial-ekonomi selaras dengan urutan waktu
  yang diceritakan, tetapi desain tidak memisahkan pengaruh BSD dari
  pertumbuhan metropolitan, jalan, kebijakan nasional, migrasi, atau
  perubahan ekonomi lain. Hubungan tersebut bersifat deskriptif dan
  interpretatif, bukan estimasi kausal.
\item
  \textbf{Inferensi review:} Potensi konflik merupakan peringatan yang
  masuk akal dalam kerangka artikel, tetapi dua contoh media tidak
  memberikan estimasi prevalensi, tren, atau perbandingan risiko konflik
  antara kawasan dengan dan tanpa kota baru.
\end{itemize}

\subsection{Conclusions}\label{conclusions-15}

\textbf{Kesimpulan penulis.} Periurban JMA telah bergerak dari sekitar
10 km dari pusat pada akhir 1970-an menjadi sekitar 40-45 km pada saat
artikel ditulis. Perluasan itu berlangsung dari permukiman kecil, menuju
kegiatan perdagangan dan jasa di sepanjang jalan Tangerang-Bogor-Bekasi
pada 1980-an, lalu menuju pembangunan lahan skala besar pada 1990-an
(Bagian 4, hlm. 228).

\textbf{Kesimpulan penulis.} Pembangunan skala besar menarik penduduk
inti dan membentuk suburbanisasi, sekaligus menarik migran dari
perdesaan sekitar dan mempercepat urbanisasi kawasan pinggiran. Pola
tersebut menghasilkan banyak subpusat dengan campuran penggunaan lahan,
menciptakan pekerjaan, serta mendorong usaha kecil kreatif dan inovatif
(Bagian 4, hlm. 228).

\textbf{Kesimpulan penulis.} Manfaat transformasi tidak terdistribusi
merata. Kota baru periurban terutama dikembangkan untuk kelompok kaya;
meskipun sebagian rumah tangga sekitar mengalami perubahan
sosial-ekonomi, kesenjangan kaya-miskin tetap besar. Ketimpangan dan
antagonisme kelas dinilai berpotensi menimbulkan konflik, sementara
bentuk ruangnya diklasifikasikan sebagai segregasi mandiri atau sukarela
(Bagian 4, hlm. 228).

\textbf{Inferensi review.} Kesimpulan mengenai arah umum perubahan JMA
didukung oleh beberapa rangkaian data dan sejarah pembangunan. Namun,
kesimpulan tentang hubungan pembangunan skala besar, perbaikan
kesejahteraan, segregasi, dan konflik harus dibatasi pada bukti
deskriptif serta kasus ekstrem BSD; file tidak menyediakan desain yang
dapat menggeneralisasi atau menguji hubungan sebab-akibat tersebut.

\subsection{Contributions}\label{contributions-15}

Bagian khusus yang merumuskan daftar kontribusi penelitian: Tidak
disebutkan dalam file. Butir berikut merangkum kontribusi yang
dinyatakan dalam tujuan atau diperagakan oleh analisis.

\textbf{Kontribusi yang dinyatakan atau diperagakan sumber:}

\begin{itemize}
\tightlist
\item
  Artikel melengkapi kajian perubahan periurban dengan menempatkan
  transformasi sosial-ekonomi dan segregasi spasial sebagai dua hasil
  yang berlangsung bersamaan dalam JMA (Bagian 1, hlm. 222-223).
\item
  Artikel menyusun lintasan historis-spasial perluasan periurban Jakarta
  dari rencana regional dan jalan antarkota menuju kota baru skala
  besar, lalu menghubungkannya dengan perubahan populasi, kepadatan, dan
  struktur ekonomi (Bagian 3, hlm. 224-225; Bagian 4, hlm. 228).
\item
  Artikel memberi bukti rumah tangga tentang komposisi migran, perubahan
  jenis pekerjaan, pengembangan usaha, dan pergeseran pendapatan di
  empat kelurahan sekitar BSD (Bagian 3.2-3.2.3, hlm. 225-227).
\item
  Pemilihan BSD sebagai kasus ekstrem digunakan untuk memperlihatkan
  secara tajam kontras antara permukiman sekitar dan enklave kota baru
  berfasilitas tinggi, tanpa mengklaim keterwakilan statistik (Bagian 1,
  hlm. 223; Bagian 3.1 dan 3.3, hlm. 225 dan 227-228).
\item
  Artikel menunjukkan dualitas substantif: periurbanisasi dapat
  membentuk subpusat, pekerjaan, dan inovasi penghidupan sekaligus
  memperbesar pemisahan sosial-spasial dan potensi konflik (Abstract,
  hlm. 221; Bagian 4, hlm. 228).
\end{itemize}

\textbf{Inferensi review.} Kontribusi metodologis artikel terletak pada
penggabungan skala regional dan rumah tangga serta pada penggunaan kasus
ekstrem sebagai alat penjelas. Artikel tidak menawarkan ukuran baru,
model kausal, atau prosedur replikasi lengkap untuk mengukur segregasi.

\subsection{Research Gap}\label{research-gap-15}

\textbf{Kesenjangan yang dinyatakan atau dibangun sumber:}

\begin{itemize}
\tightlist
\item
  Definisi kawasan periurban belum memiliki kesepakatan umum meskipun
  terdapat ciri bersama berupa transisi penggunaan lahan, pengaruh
  urban, pasar, fasilitas, dan arus komuter (Bagian 2.1, hlm. 223).
\item
  Studi periurban yang dibahas artikel sebagian besar berfokus pada
  implikasi perencanaan dan pengelolaan serta hubungan antara kegiatan
  urban dan perdesaan. J27 melengkapinya dengan pembahasan fenomena JMA
  terkini yang menitikberatkan transformasi sosial-ekonomi dan segregasi
  spasial (Bagian 1 dan 2.1, hlm. 222-223).
\item
  Pertumbuhan besar JMA sejak 1990-an memerlukan penjelasan yang
  menghubungkan pergeseran batas periurban, konversi lahan skala besar,
  migrasi, pekerjaan, pendapatan, pembentukan subpusat, dan pemisahan
  sosial (Bagian 1, hlm. 221-223).
\end{itemize}

\textbf{Pernyataan formal bahwa studi terdahulu gagal menguji pertanyaan
tertentu, beserta tinjauan sistematis yang membuktikannya:} Tidak
disebutkan dalam file.

\textbf{Kesenjangan yang masih tersisa sebagai inferensi review:}
Artikel belum menyediakan pengukuran segregasi yang dapat dibandingkan
lintas ruang dan waktu, sampel yang mewakili seluruh JMA, pembanding
kawasan tanpa pembangunan skala besar, atau identifikasi kausal dampak
kota baru. Artikel juga belum mengukur konflik secara sistematis maupun
memisahkan perubahan pendapatan nominal dari perubahan riil. Kesenjangan
ini diturunkan dari desain dan batas data J27, bukan agenda riset yang
dinyatakan penulis.

\subsection{Limitations}\label{limitations-15}

\textbf{Keterbatasan yang dinyatakan sumber:}

\begin{itemize}
\tightlist
\item
  BSD tidak mewakili kawasan periurban JMA. Kasus itu sengaja
  diperlakukan sebagai kasus ekstrem dan analisisnya tidak ditujukan
  untuk generalisasi, melainkan untuk memperlihatkan bagaimana segregasi
  berlangsung (Bagian 1, hlm. 223).
\item
  Perbandingan pendapatan 1991-2006 tidak memperhitungkan inflasi yang
  disebut rata-rata 7,7\% per tahun (catatan kaki 6, hlm. 226).
  Perbandingan pendapatan BSD dengan angka BPS (1994) juga secara
  terpisah dinyatakan tidak menghitung inflasi (catatan kaki 4, hlm.
  225).
\item
  Bagian khusus yang berjudul keterbatasan penelitian: Tidak disebutkan
  dalam file.
\end{itemize}

\textbf{Keterbatasan yang diinferensikan dari desain dan pelaporan:}

\begin{itemize}
\tightlist
\item
  Empat kelurahan dipilih karena masih berciri perdesaan, dekat dengan
  BSD, dan berada di kawasan periurban. Pemilihan lokasi bersifat
  bertujuan, sedangkan pengambilan acak proporsional baru dilakukan pada
  tingkat rumah tangga di dalam lokasi tersebut; keterwakilan bagi
  Serpong atau JMA tidak dapat diasumsikan.
\item
  Artikel menyebut sampel 201 rumah tangga dan pemilihan 220 nama,
  tetapi tidak menjelaskan selisih 19 antara keduanya, tingkat respons,
  nonrespons, atau bobot sampel (Bagian 3.2 dan Tabel 2, hlm. 225).
\item
  Nama kelurahan keempat tertulis Cilenggang dalam uraian dan
  ``Cilempang'' pada Tabel 2. Identitas yang benar tidak dapat
  dipastikan hanya dari file (Bagian 3.2 dan Tabel 2, hlm. 225).
\item
  Data rumah tangga sekitar BSD dibandingkan dengan informasi penghuni
  BSD dari sumber sekunder. Prosedur sampling, tahun, instrumen, dan
  definisi variabel kedua kelompok tidak diselaraskan secara rinci,
  sehingga kekuatan perbandingannya terbatas (Bagian 1, hlm. 223; Bagian
  3.1, hlm. 225).
\item
  Sumber dan cara memperoleh nilai 1991 pada Gambar 5-7 tidak
  dijelaskan. File tidak memastikan apakah perbandingan berasal dari
  panel, pertanyaan retrospektif, sampel berbeda, atau himpunan data
  terdahulu.
\item
  Data regional, rumah tangga, kota baru, dan contoh konflik berasal
  dari berbagai tahun dan sumber. Prosedur harmonisasi waktu dan
  definisi: Tidak disebutkan dalam file.
\item
  Analisis bersifat deskriptif tanpa uji ketidakpastian atau
  pengendalian faktor perancu. Pertumbuhan penduduk, pergeseran
  pekerjaan, dan pendapatan karena itu tidak dapat diatribusikan secara
  unik kepada BSD atau pembangunan lahan skala besar.
\item
  Segregasi tidak diukur dengan indeks spasial atau sampel rumah tangga
  yang dilaporkan secara sebanding pada kedua sisi pemisahan. Pola
  interaksi antarkelompok, alasan pemilihan tempat tinggal, dan derajat
  kesukarelaan segregasi tidak diukur secara langsung dalam hasil yang
  disajikan.
\item
  Konflik sosial didukung oleh potensi teoretis dan contoh media, bukan
  inventaris kejadian atau survei persepsi yang representatif.
  Kesimpulan mengenai pertumbuhan perilaku koersif tidak disertai angka
  deret waktu.
\item
  Angka pendapatan BSD, konversi rupiah-dolar, dan perbandingan
  Rp1.000.000 dengan Rp1.250.000 tidak konsisten secara internal (Bagian
  3.1, hlm. 225). Koreksi atau nilai pengganti tidak dapat ditentukan
  dari file.
\item
  Angka grafik untuk kepadatan, kontribusi sektoral, perubahan
  pekerjaan, dan distribusi pendapatan tidak muncul dalam TXT. Review
  hanya dapat memakai arah perubahan yang dijelaskan narasi dan caption.
\item
  Batas geografis periurban berubah dalam narasi, tetapi kriteria
  spasial yang dipakai untuk menentukan garis 10, 20, 30, atau 45 km dan
  tahun acuannya tidak dijelaskan secara operasional.
\item
  Validitas, reliabilitas, uji coba kuesioner, definisi kategori
  pekerjaan dan pendapatan, penanganan data hilang, serta bahan
  replikasi: Tidak disebutkan dalam file.
\end{itemize}

\subsection{Challenges}\label{challenges-15}

\textbf{Tantangan substantif yang dilaporkan sumber:}

\begin{itemize}
\tightlist
\item
  Pertumbuhan periurban dalam kondisi lemahnya perencanaan dan regulasi
  mendorong \emph{sprawl} serta dualitas pembangunan terencana dan tidak
  terencana (Bagian 1, hlm. 221-222).
\item
  Pengembangan dapat mendahului jalan, drainase, dan jaringan listrik
  karena harga lahan lebih menentukan keputusan pengembang daripada
  ketersediaan infrastruktur (Bagian 1, hlm. 222).
\item
  Bentang spasial dan kelembagaan periurban terfragmentasi, sementara
  campuran penggunaan lahan dapat menimbulkan degradasi lingkungan,
  kekurangan infrastruktur, dan konflik sosial (Bagian 1 dan 2.1, hlm.
  221 dan 223).
\item
  Pembentukan subpusat mengurangi sebagian perjalanan ke Jakarta, tetapi
  akses terhadap hunian dan fasilitas bermutu tetap ditentukan kemampuan
  membayar. Kemacetan dan polusi dapat dialihkan kepada penduduk yang
  lebih miskin (Bagian 3.3, hlm. 227).
\item
  Kesenjangan ruang dan kelas dapat memunculkan konflik ketika perubahan
  guna lahan menciptakan manfaat bagi investor dan kerugian bagi
  kelompok yang lemah secara ekonomi maupun hukum (Bagian 3.3, hlm.
  227-228).
\end{itemize}

\textbf{Tantangan analitis sebagai inferensi review:}

\begin{itemize}
\tightlist
\item
  Menentukan batas periurban yang bergerak memerlukan definisi spasial
  konsisten, sedangkan artikel menekankan bahwa batasnya cair dan tidak
  menyediakan algoritme delineasi.
\item
  Memisahkan transformasi akibat satu kota baru dari perubahan ekonomi
  metropolitan memerlukan pembanding atau strategi identifikasi yang
  tidak tersedia dalam desain ini.
\item
  Membandingkan kesejahteraan 1991 dan 2006 memerlukan koreksi inflasi,
  definisi pendapatan yang sama, serta dokumentasi sumber tahun awal.
\item
  Menilai segregasi membutuhkan data lokasi dan interaksi kedua kelompok
  pada skala yang sebanding; J27 menggabungkan survei permukiman sekitar
  dengan informasi sekunder kota baru dan contoh media.
\item
  Sejumlah ketidakkonsistenan sampel, nama lokasi, satuan, serta
  konversi mata uang membatasi reproduksi dan menuntut agar angka tidak
  diselaraskan melalui asumsi.
\end{itemize}

\subsection{Practical Implications}\label{practical-implications-15}

\textbf{Makna praktis dalam interpretasi penulis:}

\begin{itemize}
\tightlist
\item
  \textbf{Interpretasi penulis:} Penciptaan subpusat, pekerjaan, dan
  usaha lokal berlangsung bersama segregasi serta distribusi
  eksternalitas yang tidak merata. Transformasi ekonomi karena itu
  memiliki manfaat dan beban sosial-spasial yang perlu dibaca bersamaan
  (Abstract, hlm. 221; Bagian 3.3-4, hlm. 227-228).
\item
  \textbf{Interpretasi penulis:} Pembangunan yang terutama melayani
  kelompok berpendapatan tinggi menciptakan potensi konflik ketika
  kelompok miskin kesulitan memperoleh hunian dan menanggung dampak
  perubahan guna lahan (Bagian 3.3-4, hlm. 227-228).
\item
  \textbf{Interpretasi penulis:} Usaha perdagangan dan jasa skala kecil
  di sekitar kota baru berfungsi sebagai strategi penghidupan dan
  inovasi akar rumput dalam ekonomi periurban (Bagian 3.2.2, hlm. 226;
  Bagian 4, hlm. 228).
\end{itemize}

\textbf{Inferensi review untuk praktik:}

\begin{itemize}
\tightlist
\item
  Pemantauan pembangunan skala besar sebaiknya mencakup perubahan lahan,
  migrasi, pekerjaan, pendapatan riil, keterjangkauan hunian, akses
  terhadap fasilitas, dan eksternalitas pada permukiman sekitar agar
  manfaat dan beban dapat dibaca bersama.
\item
  Penyediaan infrastruktur dasar perlu mendahului atau menyertai
  perluasan, khususnya ketika pengembang bergerak ke lahan murah yang
  belum memiliki jalan, drainase, dan listrik.
\item
  Pencegahan konflik memerlukan mekanisme pengaduan, akses hukum, dan
  forum antarkomunitas. Artikel mendukung kebutuhan umum tersebut
  melalui diagnosisnya, tetapi tidak menguji rancangan kelembagaan
  tertentu.
\item
  Rekomendasi program, pembagian kewenangan, biaya, jadwal implementasi,
  indikator evaluasi, dan instrumen regulasi khusus: Tidak disebutkan
  dalam file. Karena itu, implikasi di atas merupakan arah penggunaan
  temuan, bukan evaluasi kebijakan yang telah diuji.
\end{itemize}

\subsection{Applications}\label{applications-15}

\begin{itemize}
\tightlist
\item
  \textbf{Aplikasi yang diperagakan sumber:} Menggunakan deret waktu
  populasi, kepadatan, dan struktur ekonomi untuk mengenali pergeseran
  kawasan periurban pada skala metropolitan (Bagian 3, hlm. 224-225).
\item
  \textbf{Aplikasi yang diperagakan sumber:} Menggunakan survei rumah
  tangga di sekitar kota baru untuk memetakan komposisi migran,
  perubahan pekerjaan, sumber penghidupan, dan distribusi pendapatan
  (Bagian 3.2-3.2.3, hlm. 225-227).
\item
  \textbf{Aplikasi yang diperagakan sumber:} Memilih kasus ekstrem untuk
  mengungkap kontras sosial-spasial yang mungkin kurang terlihat dalam
  statistik regional, dengan syarat hasil tidak diperlakukan sebagai
  representasi seluruh wilayah (Bagian 1, hlm. 223).
\item
  \textbf{Aplikasi yang diperagakan sumber:} Menggabungkan harga dan
  fasilitas kota baru, kondisi permukiman sekitar, serta contoh konflik
  untuk membangun diagnosis segregasi dan pembangunan tidak merata
  (Bagian 3.1 dan 3.3, hlm. 225 dan 227-228).
\item
  \textbf{Inferensi review:} Kerangka indikator J27 dapat dipakai
  sebagai daftar awal pemantauan transformasi di kawasan periurban lain.
  Namun, penerapannya memerlukan definisi batas, data pembanding,
  koreksi harga, dan ukuran segregasi yang lebih konsisten; validitas
  transfer di luar kasus BSD tidak diuji dalam file.
\end{itemize}

\subsection{Future Research
(Optional)}\label{future-research-optional-15}

Tidak disebutkan dalam file.

\subsection{Provenance Notes}\label{provenance-notes-15}

\begin{itemize}
\tightlist
\item
  Review ini hanya menggunakan seluruh isi file TXT
  \texttt{/Users/mac/Documents/Mac/{[}2{]}\ Obsidian\ Vault/zettelkasten/source/journal/J27-winarso-hudalah-firman-2015-peri-urban-transformation-in-the-jakarta-metropolitan-area-10.1016:j.habitatint.2015.05.024.txt},
  sebanyak 529 baris. Tidak ada sumber eksternal, penelusuran web,
  Zotero, PDF lain, atau pengetahuan di luar file yang digunakan untuk
  menambah maupun mengoreksi klaim substantif.
\item
  TXT dimulai pada halaman artikel 221 dan berakhir pada halaman 229
  setelah daftar pustaka. File tidak memuat blok metadata proses
  ekstraksi, sehingga nama PDF asal, tanggal ekstraksi, perangkat
  ekstraksi, jumlah halaman menurut metadata, status enkripsi, dan
  status OCR: Tidak disebutkan dalam file.
\item
  Identitas bibliografis diambil dari halaman 221. DOI pada isi artikel
  ditulis \texttt{10.1016/j.habitatint.2015.05.024}; tanda titik dua
  pada nama file tidak diperlakukan sebagai bentuk DOI bibliografis.
\item
  Locator \texttt{hlm.} mengikuti nomor halaman cetak yang tampak pada
  header artikel. Bagian dan nomor tabel/gambar mengikuti label yang
  terdapat dalam TXT.
\item
  Tata letak artikel dua kolom menyebabkan beberapa alinea dan daftar
  pustaka tersusun bersisian dalam hasil teks. Review merekonstruksi
  urutan hanya ketika hubungan kalimat masih jelas dan tidak mengisi
  bagian yang hilang melalui dugaan.
\item
  Isi visual Gambar 1-7 tidak dapat diperiksa langsung dari TXT. Review
  hanya menggunakan caption, angka yang tertulis dalam teks, dan
  interpretasi eksplisit penulis. Nilai yang hanya mungkin terdapat di
  dalam grafik tidak direkonstruksi.
\item
  Tabel 1-3 tersedia sebagai teks, tetapi hasil ekstraksi memuat
  ketidakselarasan nama dan format. Tabel 2 menulis ``Cilempang'',
  sedangkan narasi menulis Cilenggang. Keduanya dipertahankan sebagai
  ketidakkonsistenan sumber.
\item
  Artikel menyebut 220 nama dipilih dari kerangka sampel, sedangkan
  sampel dan total Tabel 2 berjumlah 201. File tidak menjelaskan apakah
  perbedaannya merupakan nonrespons, eksklusi, atau hal lain.
\item
  Pendahuluan menyebut populasi JMA lebih dari 30 juta dengan sitasi BPS
  (2010), sedangkan Tabel 1 memberikan total 27.447.027 untuk 2010.
  Kedua pernyataan dicatat tanpa rekonsiliasi eksternal.
\item
  Angka pendapatan BSD pada hlm. 225 memuat konversi yang tampak tidak
  konsisten, termasuk Rp33.000 menjadi ``US\$3300'', Rp825.000 menjadi
  ``US\$82,500'', dan klaim bahwa modus Rp1.000.000 berada jauh di atas
  Rp1.250.000. Review melaporkan bentuk TXT apa adanya dan tidak menebak
  desimal atau periode yang dimaksud.
\item
  Uraian ruang hijau menyebut lebih dari 10 m2 per orang dibandingkan
  standar 6 m2 per orang, sedangkan catatan kaki 7 menguraikan pedoman
  60 m2 per rumah tangga. File tidak menjelaskan konversi di antara
  kedua satuan.
\item
  Artikel bergantian memakai JMA dan JMR untuk wilayah metropolitan.
  Perbedaan konseptual atau batas antara keduanya: Tidak disebutkan
  dalam file.
\item
  Bentuk rentang seperti \texttt{221e229}, pemenggalan kata, ligatur,
  dan salah ketik merupakan artefak atau bentuk yang terdapat dalam TXT.
  Review menormalkan ejaan untuk keterbacaan, tetapi tidak mengubah
  angka atau substansi yang meragukan.
\item
  Seluruh rujukan dalam daftar pustaka dan sitasi internal diperlakukan
  sebagai laporan artikel atas literatur. Isi karya Adell, Webster, BPS,
  Harmanurjeni, Leisch, Winarso, maupun sumber lain tidak diakses atau
  diverifikasi secara independen.
\item
  Bagian \emph{Acknowledgement} menyatakan bahwa Maulien Karina Sari
  melakukan studi lapangan dan menyediakan hasil statistik, Silvania Dwi
  Utami menyediakan sebagian informasi, dan sebagian artikel berasal
  dari studi tahun 2007 yang didanai ITB Research Grant (hlm. 228).
  Peran analitis lebih lanjut dan hubungan tepat studi 2007 dengan
  survei 2006 serta pengamatan 2010: Tidak disebutkan dalam file.
\item
  Batas provenance utama: review dapat menjelaskan apa yang tertulis
  dalam TXT, tetapi tidak dapat memastikan kesetiaan tata letak terhadap
  PDF, membaca angka grafik yang tidak terekstraksi, memvalidasi data
  mentah, memeriksa sumber sekunder, atau mereplikasi perhitungan
  penelitian.
\end{itemize}

\section{Review A01}\label{review-a01}

\textbf{Identitas sumber}

\begin{itemize}
\tightlist
\item
  Kode: A01.
\item
  Judul: \emph{Stretching the Concept of `Borrowed Size'}.
\item
  Penulis: Evert J. Meijers dan Martijn J. Burger.
\item
  Afiliasi: Delft University of Technology, Belanda, untuk Evert J.
  Meijers; Erasmus University Rotterdam, Belanda, untuk Martijn J.
  Burger (halaman pembuka, hlm. 1).
\item
  Jenis dokumen menurut halaman pembuka: artikel dalam \emph{Urban
  Studies}. Status telaah sejawat: Tidak disebutkan dalam file.
\item
  Tahun: nama file mencantumkan 2017. Tahun publikasi final dalam tubuh
  artikel: Tidak disebutkan dalam file. Halaman pembuka mencantumkan hak
  cipta Urban Studies Journal Limited 2015; naskah masuk pada November
  2014 dan disetujui pada Juni 2015 (halaman pembuka, hlm. 1).
\item
  Volume dan nomor terbitan: Tidak disebutkan dalam file.
\item
  Rentang halaman yang tersedia: 1-23 menurut header artikel dan
  metadata PDF (halaman pembuka; blok \texttt{{[}PDF\ Metadata{]}}).
\item
  DOI: \texttt{10.1177/0042098015597642} (halaman pembuka, hlm. 1).
\item
  Kata kunci: \emph{agglomeration economies}, \emph{agglomeration
  spillovers}, \emph{network externalities}, \emph{territorial
  development}, dan \emph{urban system} (bagian Keywords, hlm. 1).
\item
  Objek kajian: konsep \emph{borrowed size}, bukti empiris terdahulu
  yang berkaitan dengan konsep tersebut, dan eksplorasi empiris terhadap
  fungsi metropolitan pada 1.967 kota di EU25, Norwegia, dan Swiss
  (bagian Introduction, hlm. 2; Exploring the borrowed size concept dan
  Research approach, hlm. 8-10).
\item
  Dukungan pendanaan: Netherlands Organisation for Scientific Research
  (NWO) dan Platform31 (bagian Funding, hlm. 21).
\end{itemize}

\textbf{Penanda epistemik yang digunakan}

\begin{itemize}
\tightlist
\item
  \textbf{Laporan sumber} berarti konsep, metode, data, angka, atau
  hasil yang dinyatakan langsung oleh artikel.
\item
  \textbf{Laporan sumber atas literatur} berarti cara Meijers dan Burger
  merangkum karya yang mereka rujuk; review ini tidak memperlakukannya
  sebagai verifikasi mandiri atas karya asli tersebut.
\item
  \textbf{Interpretasi penulis} berarti penjelasan Meijers dan Burger
  tentang arti, mekanisme, atau implikasi hasil mereka.
\item
  \textbf{Inferensi review} berarti pembacaan analitis yang diturunkan
  hanya dari TXT A01. Inferensi ini tidak diperlakukan sebagai temuan
  langsung artikel atau sebagai bukti tambahan.
\end{itemize}

\subsection{Summarized Abstract}\label{summarized-abstract-16}

\textbf{Laporan sumber.} Artikel ini menelaah dan memperjelas konsep
\emph{borrowed size} yang mulai digunakan sebagai konsep kebijakan di
beberapa negara Eropa, tetapi dinilai masih tidak terdefinisi dengan
baik dan belum mempunyai landasan empiris yang meyakinkan. Dari uraian
singkat Alonso (1973), penulis merekonstruksi lima premis: konsep ini
bersifat analitis; terutama menyangkut kota kecil; berlangsung dalam
kompleks metropolitan multisenter; menghasilkan keseimbangan manfaat dan
biaya aglomerasi yang lebih menguntungkan; serta dimungkinkan oleh
aksesibilitas dan konektivitas jaringan. Artikel kemudian meninjau bukti
terdahulu dan melakukan analisis eksploratif atas fungsi metropolitan di
1.967 kota Eropa (Abstract, hlm. 1; Origins and meaning of the `borrowed
size' concept, hlm. 3; Exploring the borrowed size concept, hlm. 8-10).

\textbf{Laporan sumber.} Analisis menguji apakah peminjaman ukuran lebih
sering terjadi dalam wilayah metropolitan multisenter atau polisentris,
apakah fenomena tersebut terbatas pada kota kecil, dan apakah fenomena
itu berkaitan dengan integrasi spasial. Fungsi metropolitan diukur dalam
lima ranah, yaitu politik-administratif, ilmu pengetahuan, korporasi,
budaya, dan olahraga. Penulis menggunakan korelasi Pearson, regresi OLS
dengan galat baku robust, residual model fungsi metropolitan, regresi
logistik, serta perbandingan residual menurut tingkat integrasi yang
direpresentasikan oleh FUA, polyFUA, dan PIA (Research approach, hlm.
8-10; Tables 1-3, hlm. 11-15; Borrowed size and integration, hlm.
16-18).

\textbf{Laporan sumber.} Hubungan ukuran kota dan fungsi metropolitan
jauh lebih lemah di wilayah metropolitan polisentris daripada
monosentris. Dalam wilayah polisentris, populasi kota-kota tetangga
berasosiasi positif dengan fungsi metropolitan suatu kota, sedangkan
ketersediaan fungsi di kota-kota tetangga berasosiasi negatif. Peluang
diklasifikasikan sebagai meminjam ukuran justru lebih tinggi pada kota
menengah, besar, dan sangat besar daripada kota kecil. Integrasi
fungsional yang kuat juga berkaitan dengan hasil yang lebih
menguntungkan daripada hubungan yang hanya sedang atau lemah, meskipun
bayang-bayang aglomerasi masih sering muncul (Tables 1-3, hlm. 11-15;
Figures 2 dan 4 serta uraian, hlm. 15-18).

\textbf{Interpretasi penulis.} Konsep itu perlu diperluas dalam lingkup
dan skala. \emph{Borrowed size} didefinisikan ulang sebagai keadaan
ketika sebuah kota memiliki fungsi perkotaan dan/atau tingkat kinerja
yang lazimnya berkaitan dengan kota yang lebih besar karena interaksi
dalam jaringan kota pada berbagai skala spasial menyediakan pengganti
manfaat aglomerasi. Konsep tidak terbatas pada kota kecil atau hubungan
regional yang dekat, tidak semua manfaat aglomerasi dapat dipinjam, dan
hasil pada tingkat jaringan perlu dinilai melalui konsep eksternalitas
jaringan kota (Conclusion, hlm. 18-21).

\textbf{Inferensi review.} Bukti empiris A01 secara langsung hanya
mengoperasionalkan dimensi peminjaman fungsi, bukan peminjaman kinerja
atau keseluruhan definisi baru tersebut. Karena klasifikasi
\emph{borrowed size} dibentuk dari residual model potong lintang, hasil
menunjukkan penyimpangan fungsi dari pola yang diharapkan berdasarkan
ukuran dan negara, bukan identifikasi kausal bahwa interaksi jaringan
telah menghasilkan penyimpangan itu (Research approach, hlm. 10-12;
Borrowed size and city size, hlm. 12-16; The missing link, hlm. 20).

\subsection{Problem Statement}\label{problem-statement-16}

\textbf{Laporan sumber.} Kebijakan pembangunan spasial-ekonomi di
Belanda dan Flanders mengadopsi \emph{borrowed size} sebagai cara
mengatasi kekurangan massa kritis kota-kota kecil dan menengah.
Penggunaan ini muncul di tengah literatur, terutama \emph{New Economic
Geography}, yang menempatkan manfaat aglomerasi kota besar sebagai
penggerak utama pertumbuhan. Negara-negara Benelux dinilai memiliki
konsentrasi perkotaan yang relatif rendah akibat pilihan politik dan
sosial yang historis, tetapi pembuat kebijakan berharap manfaat
metropolis besar dapat diorganisasi melalui jaringan kota-kota yang
lebih kecil (Introduction, hlm. 1-2).

\textbf{Laporan sumber.} Masalah utamanya ialah konsep perencanaan dapat
diadopsi tanpa elaborasi konseptual dan dukungan empiris yang kuat.
Kelenturan semacam itu membantu pemangku kepentingan menyesuaikan konsep
dengan situasi masing-masing, tetapi juga menimbulkan pertanyaan apakah
konsep yang kabur dapat mengarahkan tujuan dan prioritas investasi
secara efektif. Artikel mengutip penilaian Phelps bahwa \emph{borrowed
size} menarik secara intuitif, tetapi sukar ditentukan secara analitis
dan kegunaannya belum terbukti (Introduction, hlm. 2).

\textbf{Laporan sumber.} Bukti yang tersedia sebelum analisis A01
dinilai campuran dan jauh dari konklusif. Penulis secara khusus
meragukan anggapan bahwa peminjaman ukuran terkait secara eksklusif
dengan kota kecil, bahwa keberadaan kota tetangga dengan sendirinya
menguntungkan, dan bahwa jarak, aksesibilitas, atau interaksi
benar-benar mendorong proses tersebut. Pada saat yang sama,
\emph{agglomeration shadow} menawarkan mekanisme berlawanan: kedekatan
dengan konsentrasi kegiatan dapat menekan perkembangan wilayah sekitar
melalui kompetisi (Existing evidence of `borrowed size', hlm. 4-8;
Exploring the borrowed size concept, hlm. 8).

\textbf{Rumusan masalah sumber.} Artikel menanyakan apakah premis asli
Alonso cukup tepat untuk menjelaskan sistem perkotaan kontemporer dan
cukup kuat untuk menjadi prinsip kebijakan. Pertanyaan empiris
eksplisitnya adalah: (1) apakah \emph{borrowed size} lebih sering muncul
di wilayah metropolitan multisenter atau polisentris; (2) apakah
fenomena itu terbatas pada kota kecil di dalam sistem tersebut; dan (3)
apakah fenomena itu berkaitan dengan integrasi spasial (Exploring the
borrowed size concept, hlm. 8).

\subsection{Objectives}\label{objectives-16}

\begin{itemize}
\tightlist
\item
  \textbf{Laporan sumber:} memperjelas konsep \emph{borrowed size} dan
  menilai justifikasi empirisnya pada tahap awal penggunaan konsep
  tersebut dalam kebijakan dan perdebatan (Introduction, hlm. 2).
\item
  \textbf{Laporan sumber:} menyintesis literatur yang secara langsung
  ataupun tidak langsung menguji lima premis konseptual Alonso
  (Introduction, hlm. 2; Existing evidence of `borrowed size', hlm.
  4-8).
\item
  \textbf{Laporan sumber:} menghasilkan informasi baru melalui analisis
  eksploratif tentang organisasi sistem perkotaan, ukuran kota, dan
  integrasi spasial sebagai kondisi yang berkaitan dengan peminjaman
  ukuran (Introduction, hlm. 2; Exploring the borrowed size concept,
  hlm. 8).
\item
  \textbf{Laporan sumber:} merumuskan ulang lingkup dan skala konsep
  agar lebih bernilai sebagai kerangka teoretis untuk sistem perkotaan
  kontemporer dan sebagai prinsip perencanaan (Introduction, hlm. 2;
  Conclusion, hlm. 18-21).
\item
  \textbf{Interpretasi penulis:} mengarahkan perhatian dari ukuran dan
  kedekatan semata menuju interaksi dalam jaringan kota serta dari hasil
  lokal menuju eksternalitas positif atau negatif pada tingkat jaringan
  (Borrowed size: Scale dan From borrowed size to city network
  externalities, hlm. 19-20).
\end{itemize}

\subsection{Literature Survey}\label{literature-survey-16}

\textbf{Sifat survei.} Artikel melakukan tinjauan naratif dan konseptual
terhadap studi yang membahas seluruh atau sebagian premis \emph{borrowed
size}. Basis data bibliografis, kata kunci pencarian, rentang pencarian,
kriteria inklusi-eksklusi, jumlah rekaman yang disaring, serta prosedur
penilaian mutu studi: Tidak disebutkan dalam file. Karena itu, survei
ini tidak dapat diperlakukan sebagai tinjauan sistematis (Introduction,
hlm. 2; Existing evidence of `borrowed size', hlm. 4).

\textbf{Rekonstruksi konsep Alonso menurut sumber.} Penulis menyarikan
lima premis dari kurang dari satu halaman tulisan Alonso (1973):

\begin{itemize}
\tightlist
\item
  \textbf{Laporan sumber atas literatur:} \emph{borrowed size} semula
  merupakan konsep analitis yang menggambarkan keadaan yang dapat
  diamati, bukan konsep normatif kebijakan (Origins and meaning, butir
  A, hlm. 3).
\item
  \textbf{Laporan sumber atas literatur:} prosesnya dirumuskan terutama
  sebagai kota kecil yang memperoleh keuntungan dari kota lebih besar,
  walaupun contoh Alonso mengenai Jerman dan Low Countries membuka
  kemungkinan kota-kota kecil yang setara saling meminjam ukuran (butir
  B, hlm. 3).
\item
  \textbf{Laporan sumber atas literatur:} kota harus dekat dengan kota
  lain dan menjadi bagian dari entitas fungsional yang sama; tidak semua
  kota kecil dapat meminjam ukuran (butir C, hlm. 3).
\item
  \textbf{Laporan sumber atas literatur:} kota kecil dalam kompleks
  multisenter dapat mempertahankan keuntungan ukuran kecil, seperti
  kemacetan yang lebih rendah, sekaligus mengakses fasilitas, layanan
  bisnis, pergudangan, serta pasar tenaga kerja yang lebih luas milik
  pusat lain (butir D, hlm. 3).
\item
  \textbf{Laporan sumber atas literatur:} aksesibilitas dan konektivitas
  jaringan merupakan syarat. Alonso mengusulkan \emph{population
  potential}, yaitu jumlah penduduk yang mudah dicapai dan peluang
  interaksi yang tersedia, sebagai proksi konektivitas (butir E, hlm.
  3).
\end{itemize}

\textbf{Operasionalisasi dalam literatur menurut sumber.} Peminjaman
ukuran menuntut diskoneksi positif antara ukuran tempat dan
karakteristik yang biasanya terkait dengan ukuran, dan diskoneksi itu
harus berhubungan dengan posisi tempat terhadap kota lain. Karakteristik
yang pernah dipakai meliputi pendapatan per kapita, pertumbuhan
penduduk, biaya perumahan, spesialisasi jasa, lapangan kerja jasa, ruang
kantor dan ritel, perusahaan kecil, fasilitas budaya, serta fungsi
metropolitan. Posisi kota dapat dinilai melalui jarak, waktu tempuh,
keanggotaan wilayah metropolitan, aksesibilitas, atau keterlekatan
jaringan. Artikel menegaskan bahwa hubungan dasar antara ukuran dan
karakteristik harus terlebih dahulu terbukti agar penyimpangannya
bermakna (Analytical concept: Operationalising borrowed size, hlm. 4).

\textbf{Kota kecil dan arah peminjaman.} Sebagian besar literatur
terdahulu berfokus pada kota kecil-menengah, pinggiran kota, atau
\emph{edge city}. Penulis tidak menemukan studi yang secara eksplisit
menguji pengaruh kedekatan dengan kota yang lebih kecil atau bertingkat
lebih rendah. Mereka mengemukakan bahwa kota besar juga dapat meminjam
basis dukungan dari kota hinterland sehingga ukuran pusat dan fungsi
pusat menjadi tidak selaras; perbedaan antara kota kecil dan kawasan
nonperkotaan sekitarnya pun dapat menjadi semakin kabur (Small cities,
hlm. 5).

\textbf{Multisentrisitas, polisentrisitas, dan bukti campuran.} Menurut
rangkuman artikel, penelitian tingkat regional menemukan bahwa
multisenter tidak berpengaruh signifikan terhadap PDB per kapita,
sedangkan polisentrisitas dapat berkaitan negatif dengan kinerja akibat
kekurangan manfaat aglomerasi dan duplikasi fungsi tingkat rendah. Pada
tingkat lokal, studi-studi terdahulu menghasilkan temuan berbeda:
sebagian tidak menemukan kinerja kota kecil dekat metropolis yang lebih
baik, sebagian menemukan pertumbuhan pekerjaan dan penduduk yang lebih
cepat dekat kota besar, dan sebagian menunjukkan penalti akibat
keterpencilan atau pengaruh pertumbuhan kota tetangga. Penulis
menekankan bahwa hasil positif satu kota dapat tetap merupakan permainan
berjumlah nol pada skala regional (Multicentric metropolitan areas, hlm.
5-6).

\textbf{Laporan sumber atas mekanisme tandingan.} \emph{Agglomeration
shadow} dijelaskan sebagai efek persaingan konsentrasi perusahaan
terhadap wilayah sekitar. Dalam kerangka ini, tempat di dalam
bayang-bayang pusat hanya mampu menopang barang dan jasa dasar dan
menghadapi peluang perkembangan yang lebih sempit. Mekanisme ini
merupakan kebalikan konseptual dari \emph{borrowed size} (Multicentric
metropolitan areas, hlm. 6).

\textbf{Keseimbangan manfaat dan biaya.} Literatur yang dirangkum
menyatakan bahwa perusahaan dan rumah tangga menimbang manfaat kota
besar dengan kemacetan, pencemaran, kejahatan, upah, dan biaya lahan.
Kota kecil dekat metropolis dapat menerima relokasi industri yang peka
terhadap biaya sekaligus mengakses pasar, jasa bisnis, pengetahuan,
tenaga kerja, dan amenitas kota besar. Manfaat dapat menjangkau wilayah
yang lebih luas daripada biaya aglomerasi yang lebih terkonsentrasi di
dalam batas kota; artikel melaporkan bukti bahwa keseimbangan tersebut
lebih baik di wilayah metropolitan polisentris Amerika Serikat daripada
wilayah yang didominasi satu kota besar (Balance of agglomeration
benefits and costs, hlm. 6-7).

\textbf{Konektivitas dan interaksi.} Literatur yang dirangkum mengaitkan
setiap tambahan jarak dari pusat berorde lebih tinggi dengan penalti
pertumbuhan, tetapi penulis berpendapat jarak mungkin penting karena
menghambat interaksi. Phelps dan koleganya dipakai untuk menempatkan
interaksi sekaligus sebagai ukuran dan mekanisme \emph{borrowed size}.
Literatur lain menyarankan jangkauan perjalanan satu jam untuk
mempertahankan kontak tatap muka, keterlekatan dalam jaringan korporasi,
pasar, informasi, dan komunikasi melampaui kedekatan fisik, serta
integrasi pasar tenaga kerja sebagai penyebab penting. Namun,
konektivitas tidak menggantikan kedekatan untuk semua fungsi dan juga
dapat meningkatkan kompetisi; isolasi dapat melindungi pusat kecil dan
menopang layanan publik institusional meskipun massa kritis lokal tidak
mencukupi (Network connectivity, hlm. 7-8).

\textbf{Inferensi review.} Tinjauan membangun alasan yang kuat untuk
membedakan hasil positif dan negatif pada tingkat kota dari hasil
agregat jaringan. Akan tetapi, pilihan dan kelengkapan korpus tidak
dapat diaudit dari TXT karena metode penelusuran literatur tidak
dilaporkan. Semua klaim tentang studi terdahulu dalam review A01 ini
tetap merupakan laporan Meijers dan Burger atas sumber-sumber tersebut,
bukan hasil pemeriksaan karya aslinya (Introduction, hlm. 2; Existing
evidence of `borrowed size', hlm. 4-8).

\subsection{Method Used}\label{method-used-16}

\textbf{Laporan sumber.} Penelitian menggabungkan rekonstruksi
konseptual, tinjauan naratif bukti terdahulu, dan analisis kuantitatif
eksploratif potong lintang. Bagian empiris mengadopsi pandangan
fungsional: peminjaman ukuran dicari melalui kelebihan fungsi
metropolitan relatif terhadap fungsi yang lazim diharapkan dari ukuran
kota dan posisi negara (Introduction, hlm. 2; Research approach, hlm.
8-10; Borrowed size and city size, hlm. 12).

\textbf{Tahap 1: penyusunan unit spasial dan lingkungan kota.}

\begin{itemize}
\tightlist
\item
  Data fungsi metropolitan pada tingkat LAU2 diagregasikan menjadi
  \emph{Morphological Urban Areas} (MUA), yaitu kawasan terbangun yang
  bersambungan. MUA diperlakukan sebagai kota (Research approach, hlm.
  9).
\item
  Kota beserta cekungan komuternya membentuk \emph{Functional Urban
  Area} (FUA). Tingkat berikutnya ialah polyFUA yang menggabungkan FUA
  bersebelahan berdasarkan ukuran kota dan jarak. Kota berpenduduk lebih
  dari 500.000 digabung bila berjarak kurang dari 60 km dan cekungan
  tenaga kerjanya bersentuhan; ambang jarak untuk kota lebih kecil
  adalah 30 km (Research approach, hlm. 9).
\item
  \emph{Potential Integration Area} (PIA) dibentuk ketika wilayah yang
  dapat dicapai dalam 45 menit dari beberapa FUA bertumpang tindih
  sekurang-kurangnya 33\%. FUA, polyFUA, dan PIA diperlakukan sebagai
  urutan dari integrasi fungsional lebih kuat menuju lebih lemah,
  walaupun nama PIA sendiri menunjukkan potensi integrasi, bukan
  integrasi yang telah terwujud (Research approach, hlm. 9; Borrowed
  size and integration, hlm. 16-17).
\item
  Kota dianggap mempunyai tetangga atau berada dalam lingkungan
  multisenter apabila FUA, polyFUA, atau PIA-nya mencakup lebih dari
  satu kota. Kota lainnya dianggap terisolasi (Research approach, hlm.
  9).
\item
  Di antara wilayah multisenter, polisentrisitas morfologis dihitung
  dengan \emph{Herfindahl-Hirschman Index} (HHI): pangsa penduduk setiap
  kota terhadap total penduduk kota-kota dalam wilayah dikuadratkan lalu
  dijumlahkan. Nilai mendekati 0 menunjukkan ketiadaan inti dominan,
  sedangkan nilai mendekati 1 menunjukkan dominasi satu kota. Wilayah
  dengan HHI di bawah 0,56 diklasifikasikan polisentris; penulis
  menyebut ambang ini cukup baik untuk mengeluarkan contoh monosentris
  yang jelas dan memasukkan kasus polisentris yang lazim dikenal
  (Research approach, hlm. 9-10).
\end{itemize}

\textbf{Tahap 2: pembentukan indeks fungsi metropolitan.}

\begin{itemize}
\tightlist
\item
  Fungsi metropolitan dibagi ke dalam lima ranah: politik-administratif,
  ilmu pengetahuan, korporasi, budaya, dan olahraga. Setiap indikator
  dinormalisasi dari 0, yang berarti tidak ada fungsi, sampai 100, yang
  berarti nilai maksimum. Skor indikator dijumlahkan lalu dibagi dengan
  jumlah indikator untuk membentuk indeks ranah dan indeks agregat
  (Research approach, hlm. 8-9).
\item
  Fungsi yang tidak dianggap berkaitan kuat dengan ukuran kota, yaitu
  pariwisata, status ibu kota, dan transportasi, dikeluarkan dari indeks
  fungsi metropolitan. Bandar udara dapat berlokasi di kota kecil di
  luar kota besar dan pelabuhan hanya terdapat di pesisir; alasan
  tersebut diberikan dalam Note 2 (hlm. 21).
\end{itemize}

\textbf{Tahap 3: korelasi ukuran dan fungsi menurut sistem perkotaan.}

\begin{itemize}
\tightlist
\item
  Korelasi Pearson orde nol dihitung antara ukuran kota dan indeks
  fungsi agregat maupun kelima indeks ranah. Korelasi dibandingkan untuk
  seluruh kota, kota tanpa tetangga, kota dengan tetangga, kota dalam
  wilayah metropolitan monosentris, dan kota dalam wilayah metropolitan
  polisentris. Korelasi yang lebih rendah ditafsirkan sebagai diskoneksi
  ukuran dan fungsi yang lebih besar, tetapi penulis mengakui bahwa
  diskoneksi dapat mencakup peminjaman ukuran maupun bayang-bayang
  aglomerasi (Borrowed size, multicentricity and polycentricity; Table
  1, hlm. 10-11).
\end{itemize}

\textbf{Tahap 4: model OLS fungsi metropolitan dalam wilayah
polisentris.}

\begin{itemize}
\tightlist
\item
  Untuk 606 kota dengan data memadai di wilayah metropolitan
  polisentris, penulis mengestimasi model OLS bentuk tereduksi. Variabel
  hasilnya adalah indeks fungsi metropolitan kota. Prediktor utamanya
  ialah penduduk lokal, jumlah penduduk di bagian lain polyFUA, dan
  jumlah indeks fungsi metropolitan di kota-kota lain dalam polyFUA
  (persamaan model dan Table 2, hlm. 11-13).
\item
  Model mengendalikan indeks pariwisata, indeks transportasi, PDB per
  kapita tingkat NUTS-3, dan variabel semu negara. Indeks pariwisata
  berasal dari keberadaan situs warisan dunia dalam BBSR; indeks
  transportasi didasarkan pada bandar udara, pelabuhan, dan hubungan
  kereta internasional. Sumber serta tahun PDB per kapita: Tidak
  disebutkan dalam file (Note 4, hlm. 21).
\item
  Seluruh model menggunakan galat baku robust. Penulis melaporkan bahwa
  statistik \emph{Variance Inflation Factor} tidak menunjukkan masalah
  multikolinearitas. Model diestimasi untuk indeks agregat dan untuk
  setiap ranah fungsi (Borrowed size, multicentricity and
  polycentricity; Table 2, hlm. 12-13).
\end{itemize}

\textbf{Tahap 5: klasifikasi residual dan regresi logistik.}

\begin{itemize}
\tightlist
\item
  Penulis meregresikan indeks total fungsi metropolitan terhadap ukuran
  kota dan variabel semu negara. Residual positif menandai kota dengan
  fungsi lebih banyak daripada yang diharapkan dan diklasifikasikan
  sebagai meminjam ukuran; residual negatif menandai fungsi lebih
  sedikit daripada yang diharapkan dan ditafsirkan sebagai
  \emph{agglomeration shadow}. Kota yang terlalu kecil untuk diharapkan
  mempunyai fungsi dan memang tidak mempunyai fungsi ditempatkan dalam
  kategori ketiga (Borrowed size and city size, hlm. 12).
\item
  Jika kota begitu kecil sehingga prediksi fungsi berada di bawah nol
  dan tidak ada fungsi yang ditemukan, residual positifnya diubah
  menjadi nol. Penulis juga menguji spesifikasi residual dengan kontrol
  pariwisata, transportasi, dan PDB per kapita dan menyatakan
  kesimpulannya sangat serupa (Notes 5-7, hlm. 21).
\item
  Regresi logistik multivariat terhadap 1.967 kota mengestimasi peluang
  klasifikasi \emph{borrowed size}. Prediktornya adalah kelas ukuran
  kota, jenis lingkungan kota, dan interaksi keduanya. Hasil dilaporkan
  sebagai \emph{odds ratio} dengan interval kepercayaan 95\% (Table 3,
  hlm. 14; uraian, hlm. 12 dan 15).
\end{itemize}

\textbf{Tahap 6: perbandingan tingkat integrasi.}

\begin{itemize}
\tightlist
\item
  Penulis kembali memakai residual dari model fungsi metropolitan
  terhadap ukuran dan negara, lalu membandingkan residual rata-rata kota
  menurut kombinasi hubungan kuat dalam FUA, sedang dalam polyFUA, dan
  lemah dalam PIA. Residual rata-rata dipilih agar besarnya peminjaman
  fungsi atau bayang-bayang aglomerasi ikut diperhitungkan, bukan hanya
  kategori biner (Borrowed size and integration, hlm. 16-18; Figure 4,
  hlm. 17).
\item
  Banyak kota masuk ke beberapa delimitasi sekaligus. Analisis pertama
  membandingkan kategori murni kuat, sedang, dan lemah; analisis kedua
  membandingkan kombinasi hubungan kuat-sedang dengan sedang-lemah;
  analisis ketiga membandingkan kota dengan hubungan pada ketiga tingkat
  dengan kota tanpa tetangga. Kombinasi kuat-lemah yang mencakup 200
  kota dikeluarkan karena penulis menyatakan hasilnya tidak dapat
  ditafsirkan secara bermakna (Figure 4 dan uraian, hlm. 17-18).
\end{itemize}

\textbf{Inferensi review atas desain.} Analisis ini bersifat
observasional, eksploratif, dan asosiatif. Tidak ada perlakuan
eksperimental, urutan waktu sebab-akibat, instrumen, atau strategi
identifikasi kausal yang dilaporkan. FUA, polyFUA, dan PIA adalah proksi
bertingkat untuk integrasi; TXT tidak melaporkan arus interaksi aktual
yang seragam untuk seluruh tingkatan. Karena itu, istilah ``dimungkinkan
oleh interaksi'' dalam definisi akhir lebih luas daripada hal yang dapat
diidentifikasi langsung oleh model potong lintang A01 (Research
approach, hlm. 8-10; Borrowed size and integration, hlm. 16-18; Borrowed
size: Scale, hlm. 20).

\subsection{Dataset}\label{dataset-16}

\textbf{Laporan sumber.} Bagian empiris menggabungkan data fungsi
metropolitan, delimitasi spasial, jumlah penduduk, dan variabel kontrol
sebagai berikut.

\begin{itemize}
\tightlist
\item
  \textbf{Basis data fungsi metropolitan:} database German Federal
  Institute for Research on Building, Urban Affairs and Spatial
  Development atau BBSR (2011). Data fungsi individual dikumpulkan untuk
  2004-2009 dan mayoritas observasi merujuk 2008. Artikel mengarahkan
  pembaca ke BBSR (2011) untuk uraian rinci cara pengumpulan; prosedur
  tersebut tidak direproduksi dalam file A01 (Research approach, hlm.
  8).
\item
  \textbf{Politik-administratif:} mencakup keberadaan institusi Uni
  Eropa, organisasi internasional, dan organisasi nonpemerintah
  (Research approach, hlm. 8).
\item
  \textbf{Ilmu pengetahuan:} mencakup universitas besar dan organisasi
  penelitian internasional (Research approach, hlm. 8).
\item
  \textbf{Korporasi:} mencakup perusahaan 500 teratas berdasarkan omzet
  dan jumlah staf, jasa produsen maju, bank, serta pameran dagang
  (Research approach, hlm. 8).
\item
  \textbf{Budaya:} mencakup acara musik, pameran seni, festival film,
  teater tingkat atas, gedung opera, galeri, dan museum (Research
  approach, hlm. 8).
\item
  \textbf{Olahraga:} mencakup stadion, tempat Olimpiade Musim Panas, dan
  acara olahraga besar (Research approach, hlm. 8).
\item
  \textbf{Skala asal dan agregasi:} data tersedia pada tingkat LAU2
  untuk seluruh Eropa lalu diagregasikan ke MUA. Delimitasi MUA, FUA,
  dan polyFUA berasal dari proyek ESPON 1.4.3 yang memperbarui ESPON
  1.1.1; PIA berasal dari ESPON 1.1.1 (Research approach, hlm. 9;
  References, hlm. 22).
\item
  \textbf{Penduduk kota:} dipakai sebagai ukuran kota dan sebagai dasar
  HHI serta model statistik. Sumber, tahun referensi, prosedur
  harmonisasi dengan data fungsi, dan definisi rinci jumlah penduduk:
  Tidak disebutkan dalam file.
\item
  \textbf{Variabel kontrol:} indeks pariwisata dari BBSR (2011), indeks
  transportasi, dan PDB per kapita NUTS-3. Sumber data transportasi dan
  PDB, tahun observasi keduanya, nilai hilang, serta cara pencocokan
  spasial: Tidak disebutkan dalam file (Note 4, hlm. 21).
\item
  \textbf{Ketersediaan dan replikasi:} tautan unduh data olahan, kode
  analisis, daftar seluruh kota, dan kamus variabel: Tidak disebutkan
  dalam file. Daftar lengkap wilayah metropolitan polisentris dan
  rincian identifikasinya dinyatakan tersedia dalam Meijers et
  al.~(2015), bukan di dalam TXT A01 (Note 3, hlm. 21).
\end{itemize}

\subsection{Population / Sample}\label{population-sample-16}

\textbf{Populasi spasial utama.} Analisis mencakup 1.967 aglomerasi
perkotaan atau MUA yang selanjutnya disebut kota, tersebar di EU25,
Norwegia, dan Swiss (Research approach, hlm. 9).

\textbf{Klasifikasi ukuran kota.} Kota kecil berpenduduk kurang dari
76.000; kota menengah 76.000-200.000; kota besar 201.000-500.000; dan
kota sangat besar lebih dari 500.000. Jumlah kota dalam setiap kelas
ukuran: Tidak disebutkan dalam file (Borrowed size and city size, hlm.
12; Table 3, hlm. 14).

\textbf{Sampel menurut lingkungan kota.} Angka dalam TXT tidak
sepenuhnya konsisten:

\begin{itemize}
\tightlist
\item
  Uraian inventaris pada hlm. 10 menyebut 651 kota tanpa tetangga dan
  1.316 kota dengan tetangga; 609 dari kelompok bertetangga berada dalam
  wilayah polisentris.
\item
  Table 1 pada hlm. 11 menyebut 648 kota tanpa tetangga dan 1.319 kota
  dengan tetangga, terdiri atas 710 kota dalam wilayah monosentris dan
  609 dalam wilayah polisentris.
\item
  Caption Figure 4 pada hlm. 17 menyebut 649 kota tanpa tetangga.
\item
  Penjelasan yang merekonsiliasi selisih 648, 649, dan 651 kota
  terisolasi serta 1.316 dan 1.319 kota bertetangga: Tidak disebutkan
  dalam file. Review ini tidak memilih salah satu angka sebagai koreksi
  atas angka lainnya.
\end{itemize}

\textbf{Sampel model.} Table 2 memakai 606 kota dalam wilayah
polisentris karena tiga dari 609 kota kekurangan data. Table 3 memakai
seluruh 1.967 kota. Hanya lima kota sangat besar yang terisolasi;
penulis menyatakan jumlah ini mengurangi daya statistik perbandingan
lingkungan untuk kelas tersebut (Table 2, Note, hlm. 13; Borrowed size
and city size, hlm. 15).

\textbf{Keanggotaan delimitasi integrasi.} Sebanyak 586 kota berada
dalam FUA yang memuat beberapa kota, 327 berada dalam polyFUA, dan 1.035
berada dalam PIA; kategori ini bertumpang tindih. Untuk Figure 4,
kategori yang dilaporkan terdiri atas 180 kota hanya dengan hubungan
kuat, 18 hanya sedang, 611 hanya lemah, 85 kuat dan sedang, 103 sedang
dan lemah, 121 kuat-sedang-lemah, 200 kuat dan lemah yang dikeluarkan
dari interpretasi, serta 649 tanpa tetangga (Borrowed size and
integration, hlm. 16-17; caption Figure 4, hlm. 17).

\textbf{Populasi manusia atau sampel responden individual.} Tidak
berlaku. Unit observasi penelitian adalah kota, bukan individu, rumah
tangga, atau perusahaan. Teknik penarikan sampel probabilitas individu:
Tidak berlaku.

\subsection{Dependent Variables}\label{dependent-variables-16}

\textbf{Keterangan terminologis.} Artikel memakai beberapa keluaran
analitis, bukan satu variabel dependen yang sama untuk seluruh tahapan.

\begin{itemize}
\tightlist
\item
  Dalam korelasi Pearson, keluaran yang dikaitkan dengan ukuran kota
  ialah indeks agregat fungsi metropolitan serta indeks
  politik-administratif, ilmu pengetahuan, korporasi, budaya, dan
  olahraga (Table 1, hlm. 11).
\item
  Dalam model OLS, variabel dependen ialah nilai indeks fungsi
  metropolitan kota, baik agregat maupun per ranah (persamaan model dan
  Table 2, hlm. 11-13).
\item
  Dalam regresi logistik, variabel dependen biner ialah klasifikasi
  peminjaman ukuran: \texttt{1\ =\ ya}, berdasarkan residual positif
  dari estimasi fungsi terhadap ukuran dan variabel semu negara.
  Kelompok pembanding mencakup kota dalam bayang-bayang aglomerasi dan
  kota tanpa fungsi metropolitan (Table 3, hlm. 14; uraian, hlm. 12).
\item
  Dalam analisis integrasi, keluaran yang dibandingkan ialah residual
  rata-rata fungsi metropolitan menurut kategori hubungan kota. Residual
  positif menunjukkan surplus fungsi terhadap nilai yang diperkirakan,
  sedangkan residual negatif menunjukkan defisit (Borrowed size and
  integration, hlm. 16-18; Figure 4, hlm. 17).
\item
  Tingkat kinerja kota, pertumbuhan, produktivitas, pendapatan, atau
  ukuran kesejahteraan bukan variabel dependen empiris A01. Penulis
  menyebut peminjaman kinerja sebagai dimensi konseptual yang belum
  diuji dalam artikel (Borrowed size: Scope, hlm. 18-19; The missing
  link, hlm. 20).
\end{itemize}

\textbf{Prediktor dan pembeda utama.} Ukuran penduduk lokal, penduduk
kota lain dalam polyFUA, fungsi kota lain dalam polyFUA, kelas ukuran
kota, lingkungan terisolasi/monosentris/polisentris, interaksi
ukuran-lingkungan, tingkat integrasi, kontrol pariwisata, kontrol
transportasi, PDB per kapita, dan variabel semu negara digunakan pada
tahapan yang berlainan (Tables 1-3, hlm. 11-14; Notes 4-6, hlm. 21).

\subsection{Results}\label{results-16}

\textbf{Hubungan ukuran dan fungsi menurut sistem perkotaan.}

\begin{itemize}
\tightlist
\item
  \textbf{Laporan sumber:} untuk seluruh 1.967 kota, korelasi ukuran
  dengan indeks fungsi agregat adalah 0,886. Korelasi per ranah ialah
  0,475 untuk politik-administratif, 0,771 untuk ilmu pengetahuan, 0,850
  untuk korporasi, 0,833 untuk budaya, dan 0,807 untuk olahraga. Seluruh
  korelasi dalam Table 1 dinyatakan signifikan pada taraf 0,01, dua sisi
  (Table 1, hlm. 11).
\item
  \textbf{Laporan sumber:} pada kota tanpa tetangga, korelasi fungsi
  agregat adalah 0,749; nilai per ranah berturut-turut 0,402, 0,352,
  0,678, 0,684, dan 0,741. Pada kota dengan tetangga, korelasi agregat
  justru 0,887; nilai per ranah 0,474, 0,777, 0,853, 0,835, dan 0,806
  (Table 1, hlm. 11).
\item
  \textbf{Laporan sumber:} perbedaan terbesar muncul ketika kota
  bertetangga dibagi menurut struktur wilayah. Dalam wilayah
  monosentris, korelasi agregat mencapai 0,943 dan korelasi per ranah
  adalah 0,869, 0,867, 0,929, 0,886, dan 0,809. Dalam wilayah
  polisentris, korelasi agregat turun menjadi 0,686 dan nilai per ranah
  menjadi 0,252, 0,514, 0,697, 0,637, dan 0,803 (Table 1, hlm. 11).
\item
  \textbf{Interpretasi penulis:} keberadaan tetangga saja tidak otomatis
  memicu peminjaman ukuran. Diskoneksi ukuran-fungsi lebih sering
  terdapat dalam wilayah polisentris daripada monosentris, kecuali pada
  fungsi olahraga. Di wilayah monosentris, kota dominan tampak menyerap
  fungsi sehingga ukuran dan fungsi sangat erat; proses peminjaman
  ukuran dan bayang-bayang aglomerasi lebih sering muncul di antara kota
  tetangga yang relatif seimbang (Borrowed size, multicentricity and
  polycentricity, hlm. 10-11).
\end{itemize}

\textbf{Model OLS pada kota dalam wilayah polisentris.}

\begin{itemize}
\tightlist
\item
  \textbf{Laporan sumber:} untuk indeks fungsi agregat, koefisien
  penduduk lokal per 1.000 orang adalah 0,006 dengan galat baku 0,001;
  koefisien penduduk di bagian lain wilayah metropolitan per 10.000
  orang adalah 0,004 dengan galat baku 0,001. Keduanya ditandai
  signifikan pada taraf 0,01 dalam Table 2 (hlm. 13).
\item
  \textbf{Laporan sumber dengan catatan ekstraksi:} indeks fungsi di
  bagian lain wilayah berasosiasi negatif dengan fungsi kota sendiri.
  Angka koefisien agregat tercetak sebagai \texttt{20.063} dalam TXT
  dengan galat baku 0,021 dan tanda signifikansi 0,01; uraian penulis
  secara eksplisit menyebut arah hubungan negatif. Review ini tidak
  mengubah bentuk angka yang terdistorsi menjadi nilai baru (Table 2,
  hlm. 13; uraian hasil, hlm. 12).
\item
  \textbf{Laporan sumber:} model indeks agregat mencakup 606 observasi
  dan mempunyai \texttt{R2\ =\ 0,612}. \texttt{R2} model ranah adalah
  0,212 untuk politik-administratif, 0,431 untuk ilmu pengetahuan, 0,579
  untuk korporasi, 0,705 untuk budaya, dan 0,553 untuk olahraga (Table
  2, hlm. 13).
\item
  \textbf{Interpretasi penulis:} tambahan penduduk kota tetangga
  menyediakan permintaan eksternal yang berkaitan positif dengan fungsi
  di suatu kota, sedangkan fungsi yang sudah tersedia di tetangga
  menjadi pasokan pesaing. Pola gabungan itu mendukung pembacaan bahwa
  diskoneksi ukuran-fungsi di wilayah polisentris mencerminkan
  \emph{borrowed size} sekaligus \emph{agglomeration shadow}. Bukti per
  ranah dinilai paling jelas untuk fungsi korporasi, budaya, dan
  olahraga; ketidakpastian parameter terlalu besar untuk kesimpulan
  serupa pada politik-administratif dan ilmu pengetahuan (Borrowed size,
  multicentricity and polycentricity, hlm. 11-12).
\end{itemize}

\textbf{Residual dan peluang peminjaman ukuran.}

\begin{itemize}
\tightlist
\item
  \textbf{Laporan sumber:} dibandingkan kota kecil, kota menengah
  mempunyai \emph{odds ratio} (OR) 1,86 dengan interval kepercayaan 95\%
  1,44-2,42; kota besar OR 1,87 dengan interval 1,28-2,72; dan kota
  sangat besar OR 2,70 dengan interval 1,68-4,36. Ketiganya ditandai
  signifikan pada taraf 0,01. Penulis menyatakan tidak ada perbedaan
  signifikan antara kota menengah, besar, dan sangat besar (Table 3,
  hlm. 14; uraian, hlm. 12).
\item
  \textbf{Laporan sumber:} regresi ukuran kota dan variabel semu negara
  yang digunakan untuk membentuk residual mempunyai
  \texttt{adjusted\ R2\ =\ 0,800}. Spesifikasi tambahan yang juga
  memasukkan pariwisata, transportasi, dan PDB per kapita dinyatakan
  menghasilkan kesimpulan yang sangat serupa (Notes 5-6, hlm. 21).
\item
  \textbf{Laporan sumber:} ketika ukuran kota belum diinteraksikan
  dengan lingkungan, kota dalam wilayah monosentris mempunyai OR 1,06
  dengan interval 0,81-1,38 dan kota dalam wilayah polisentris OR 1,00
  dengan interval 0,76-1,32 dibandingkan kota terisolasi. Keduanya tidak
  signifikan (Table 3, hlm. 14; uraian, hlm. 12).
\item
  \textbf{Laporan sumber:} setelah interaksi, peluang kota menengah di
  wilayah polisentris lebih rendah daripada kota menengah terisolasi, OR
  0,55 dengan interval 0,32-0,92 dan \texttt{p\ \textless{}\ 0,05}.
  Perbandingan kota menengah monosentris terhadap kota menengah
  terisolasi menghasilkan OR 0,67 dengan interval 0,39-1,16 dan tidak
  signifikan pada 5\% (Table 3, hlm. 14; uraian, hlm. 15).
\item
  \textbf{Laporan sumber:} kota besar dalam wilayah polisentris
  mempunyai peluang lebih tinggi daripada kota besar terisolasi, OR 4,23
  dengan interval 1,33-13,4 dan \texttt{p\ \textless{}\ 0,05}. Kota
  besar monosentris mempunyai OR 1,88 dengan interval 0,64-6,01
  dibandingkan kota besar terisolasi dan perbedaannya tidak signifikan.
  Perbedaan lingkungan pada kota sangat besar juga tidak signifikan,
  terutama karena hanya lima kota sangat besar yang terisolasi (Table 3,
  hlm. 14; uraian, hlm. 15).
\item
  \textbf{Laporan sumber:} Figure 2 menunjukkan bahwa jauh lebih banyak
  kota tidak meminjam ukuran. Kota yang lebih besar lebih sering
  meminjam daripada kota kecil. Ketiadaan tetangga mengurangi peluang,
  terutama bagi kota lebih besar, tetapi kota menengah terisolasi lebih
  jarang menghadapi bayang-bayang aglomerasi dan dapat menjadi pusat
  bagi hinterland perdesaan. Kota kecil bertetangga hanya sedikit lebih
  sering meminjam ukuran daripada kota kecil terisolasi. Di antara kota
  bertetangga, kota lebih besar dalam wilayah polisentris lebih sering
  meminjam daripada padanannya dalam wilayah monosentris (Figure 2 dan
  uraian, hlm. 15-16).
\end{itemize}

\textbf{Integrasi dan residual fungsi metropolitan.}

\begin{itemize}
\tightlist
\item
  \textbf{Laporan sumber:} di antara kategori yang hanya mempunyai satu
  tingkat hubungan, kota dengan integrasi kuat dalam FUA berkinerja
  lebih baik dalam residual fungsi daripada kota yang hanya mempunyai
  hubungan sedang dalam polyFUA atau lemah dalam PIA. Penulis
  menafsirkan bahwa integrasi lebih kuat dapat mengatasi sebagian efek
  persaingan dan mengurangi bayang-bayang aglomerasi (Figure 4a dan
  uraian, hlm. 17-18).
\item
  \textbf{Laporan sumber:} ketika kota dengan hubungan kuat-sedang
  dibandingkan dengan kota berhubungan sedang-lemah, pola umumnya
  serupa. Bayang-bayang aglomerasi masih dominan, tetapi lebih kuat
  ketika hubungan kuat tidak ada (Figure 4b dan uraian, hlm. 18).
\item
  \textbf{Laporan sumber:} penulis meminta kehati-hatian karena rasio
  tepat antara hubungan kuat dan sedang atau antara hubungan sedang dan
  lemah dalam setiap kategori tidak dapat ditentukan (Borrowed size and
  integration, hlm. 18).
\item
  \textbf{Laporan sumber:} kota tanpa tetangga tidak selalu lebih buruk
  karena isolasi melindungi dari kompetisi regional dan menciptakan
  fungsi pusat untuk hinterland, yang disebut \emph{institutional
  borrowed size}. Namun, kota tersebut tidak dapat meminjam fungsi dari
  kota lain. Dibandingkan kelompok terisolasi, kota yang sekaligus
  mempunyai banyak hubungan kuat, sedang, dan lemah lebih didominasi
  efek peminjaman ukuran (Figure 4c dan uraian, hlm. 18).
\item
  Nilai numerik residual pada Figure 4 dan proporsi setiap kategori yang
  benar-benar meminjam ukuran: Tidak disebutkan dalam file. TXT
  menyediakan caption, jumlah kota, dan uraian arah hasil, tetapi tidak
  memuat label numerik grafik sebagai teks yang dapat diperiksa (Figure
  4, hlm. 17).
\end{itemize}

\subsection{Findings}\label{findings-16}

\begin{itemize}
\tightlist
\item
  \textbf{Interpretasi penulis:} \emph{borrowed size} tidak terbukti
  sebagai konsekuensi otomatis dari mempunyai kota tetangga. Korelasi
  ukuran-fungsi justru lebih rendah pada kota terisolasi daripada kota
  bertetangga secara keseluruhan. Di antara kota bertetangga, korelasi
  jauh lebih rendah dalam wilayah polisentris daripada monosentris.
  Dengan demikian, mempunyai tetangga saja tidak memadai untuk
  menjelaskan diskoneksi ukuran-fungsi (Table 1 dan uraian, hlm. 10-11).
\item
  \textbf{Interpretasi penulis:} wilayah polisentris memuat dua proses
  yang berjalan bersamaan. Populasi kota lain memperbesar basis
  permintaan dan mendukung fungsi, tetapi fungsi yang sudah tersedia di
  kota lain menjadi pesaing. Peminjaman ukuran dan bayang-bayang
  aglomerasi karena itu bukan klasifikasi wilayah yang saling
  meniadakan; kota-kota berbeda dalam jaringan yang sama dapat mengalami
  hasil berlawanan (model dan Table 2, hlm. 11-13; From borrowed size to
  city network externalities, hlm. 20).
\item
  \textbf{Interpretasi penulis:} peminjaman fungsi tidak terbatas pada
  kota kecil. Peluang klasifikasi \emph{borrowed size} meningkat pada
  kelas kota yang lebih besar, dan kota besar dalam wilayah polisentris
  mempunyai keunggulan paling jelas terhadap kota besar terisolasi.
  Penulis menyimpulkan bahwa fenomena ini lebih sering muncul pada kota
  lebih besar, terutama yang menjadi bagian dari entitas metropolitan
  polisentris (Table 3, hlm. 14; Borrowed size and city size, hlm.
  15-16).
\item
  \textbf{Interpretasi penulis:} integrasi fungsional yang kuat lebih
  mendukung peminjaman fungsi daripada kedekatan atau potensi integrasi
  yang lemah. Isolasi dapat melindungi kota dari kompetisi dan
  menghasilkan fungsi pusat bagi hinterland, tetapi tidak memungkinkan
  peminjaman fungsi melalui jaringan kota (Borrowed size and
  integration, hlm. 16-18).
\item
  \textbf{Interpretasi penulis:} ``ukuran'' yang dipinjam dapat berupa
  basis dukungan bagi fungsi atau tingkat kinerja yang lazimnya dimiliki
  kota lebih besar. Sebuah kota dapat menopang fungsi bagi kota lain,
  sedangkan kota pemberi basis dukungan memperoleh akses dan kinerja
  dari fungsi tersebut. Jika peminjaman fungsi dan kinerja terjadi
  bersama, penulis menyebutnya \emph{genuine borrowed size} (Borrowed
  size: Scope, hlm. 18-19).
\item
  \textbf{Interpretasi penulis:} tidak semua eksternalitas aglomerasi
  sama mudahnya dipinjam. Keberagaman input dan proses pencocokan
  diperkirakan lebih mudah dipinjam daripada limpahan pengetahuan;
  layanan berorde tinggi, lembaga serta infrastruktur peka-skala juga
  mungkin dibagi. Artikel menegaskan bahwa tidak semua manfaat ukuran
  dapat dipinjam dan bahwa ukuran juga membawa biaya, seperti kemacetan,
  pencemaran, dan kejahatan (Borrowed size: Scope, hlm. 19).
\item
  \textbf{Interpretasi penulis:} peminjaman ukuran dapat berlangsung
  dari skala metropolitan sampai global karena interaksi lebih penting
  daripada jarak atau akses belaka. Eksternalitas yang berlainan dapat
  dibagi pada skala berbeda; ekonomi urbanisasi khususnya dapat dibagi
  pada skala regional atau metropolitan (Borrowed size: Scale, hlm.
  19-20).
\item
  \textbf{Inferensi review:} korelasi yang lebih rendah hanya
  menunjukkan bahwa ukuran kurang menentukan fungsi; ukuran itu sendiri
  tidak menunjukkan apakah penyimpangan bersifat positif atau negatif.
  Penulis menangani persoalan ini melalui arah residual dan model fungsi
  kota tetangga, tetapi klasifikasi residual tetap bergantung pada
  spesifikasi nilai yang ``diharapkan'' oleh model (Table 1, hlm. 11;
  uraian residual, hlm. 12).
\item
  \textbf{Inferensi review:} hasil tidak mendukung proposisi sederhana
  bahwa polisentrisitas selalu menguntungkan semua kelas kota. Kota
  besar polisentris menunjukkan peluang peminjaman yang lebih tinggi,
  tetapi kota menengah polisentris justru mempunyai peluang lebih rendah
  daripada kota menengah terisolasi. Efek yang dilaporkan bergantung
  pada ukuran kota dan lingkungan sekaligus (Table 3, hlm. 14-15).
\item
  \textbf{Inferensi review:} bukti A01 paling kuat dibaca sebagai
  pemetaan empiris ketidakselarasan fungsi dan ukuran dalam satu sistem
  kota Eropa. Mekanisme interaksi, peminjaman kinerja, dan hasil bersih
  pada tingkat jaringan masih merupakan penjelasan teoretis yang belum
  diukur langsung dalam artikel (Research approach, hlm. 8-10; The
  missing link, hlm. 20).
\end{itemize}

\subsection{Conclusions}\label{conclusions-16}

\textbf{Kesimpulan penulis.} Lima premis asli Alonso terlalu sempit
untuk sistem perkotaan kontemporer. Berdasarkan tinjauan dan analisis
eksploratif, konsep perlu diperjelas dan diperluas pada dimensi lingkup
serta skala. Peminjaman ukuran tidak terbatas pada kota kecil atau
kawasan multisenter/polisentris, walaupun dimensinya tampak paling
menonjol dalam kawasan polisentris (Conclusion, hlm. 18-20).

\textbf{Definisi ulang penulis.} Peminjaman ukuran terjadi ketika sebuah
kota menunjukkan fungsi perkotaan dan/atau tingkat kinerja yang biasanya
dikaitkan dengan kota lebih besar sebagai hasil interaksi dalam jaringan
kota pada berbagai skala spasial yang menyediakan pengganti bagi manfaat
aglomerasi (Borrowed size: Scale, hlm. 20).

\textbf{Kesimpulan penulis.} Hubungan peminjaman dapat berjalan dua
arah. Kota kecil dapat mengakses manfaat kota besar, tetapi kota kecil
juga membantu kota besar mempertahankan lebih banyak fungsi metropolitan
daripada yang dapat ditopang kota besar itu sendiri. Penulis
memperkirakan kota besar lebih sering meminjam manfaat fungsional,
sedangkan kota kecil lebih lazim meminjam dalam bentuk kinerja; hubungan
empiris kedua dimensi masih harus ditentukan (Borrowed size: Scale dan
The missing link, hlm. 20).

\textbf{Kesimpulan penulis.} Pada skala sistem, hasil lokal perlu
dijumlahkan sebagai eksternalitas jaringan kota. Eksternalitas jaringan
positif muncul ketika peminjaman ukuran lebih dominan daripada
bayang-bayang aglomerasi, sedangkan eksternalitas negatif muncul ketika
bayang-bayang lebih dominan. Penulis menilai eksternalitas jaringan kota
positif sebagai tujuan kebijakan regional yang lebih tepat daripada
sekadar mengejar \emph{borrowed size} suatu kota (From borrowed size to
city network externalities, hlm. 20).

\textbf{Kesimpulan teoretis penulis.} Konsep tersebut tidak membuang
mekanisme manfaat aglomerasi, tetapi menantang landasan geografis teori
aglomerasi dengan menyatakan bahwa eksternalitas tidak harus berhenti di
batas aglomerasi dan dapat dibagi melalui jaringan kota. Penulis
mengajukannya sebagai bagian dari penjelasan atas dinamika kota Eropa
yang menurut literatur mereka tidak sepenuhnya sesuai dengan ekspektasi
bahwa kota besar menjadi penggerak pertumbuhan utama (The missing link
dan paragraf penutup, hlm. 20-21).

\textbf{Inferensi review.} Definisi akhir mencakup fungsi, kinerja,
interaksi, banyak skala, dan substitusi manfaat aglomerasi, tetapi
pengujian A01 hanya mencakup fungsi dan memakai posisi dalam delimitasi
spasial sebagai proksi integrasi. Oleh karena itu, artikel memberikan
dasar eksploratif bagi perluasan konsep, bukan pengujian lengkap atas
seluruh unsur definisi barunya (Research approach, hlm. 8-10; Borrowed
size: Scope dan Scale, hlm. 18-20; The missing link, hlm. 20).

\subsection{Contributions}\label{contributions-16}

\begin{itemize}
\tightlist
\item
  \textbf{Kontribusi yang dinyatakan dan diperagakan sumber:} membongkar
  uraian singkat Alonso menjadi lima premis yang dapat diperiksa secara
  konseptual dan empiris (Origins and meaning, hlm. 3).
\item
  \textbf{Kontribusi yang dinyatakan sumber:} menyatukan literatur yang
  sebelumnya membahas peminjaman ukuran secara langsung atau menyentuh
  bagian-bagian premisnya secara tidak langsung, sekaligus
  memperlihatkan bahwa bukti yang tersedia campuran (Existing evidence,
  hlm. 4-8; Exploring the borrowed size concept, hlm. 8).
\item
  \textbf{Kontribusi empiris sumber:} mengoperasionalkan peminjaman
  fungsi sebagai diskoneksi positif antara ukuran dan fungsi
  metropolitan serta membedakannya dari residual negatif yang dibaca
  sebagai bayang-bayang aglomerasi (Research approach, hlm. 10-12).
\item
  \textbf{Kontribusi empiris sumber:} membandingkan 1.967 kota di
  beberapa jenis sistem perkotaan, menguji heterogenitas menurut kelas
  ukuran, dan menghubungkan surplus atau defisit fungsi dengan beberapa
  tingkat integrasi spasial (Tables 1-3 dan Figures 2-4, hlm. 11-18).
\item
  \textbf{Kontribusi konseptual sumber:} membedakan peminjaman fungsi
  dari peminjaman kinerja, menjelaskan kemungkinan hubungan timbal balik
  antara kota penopang dan kota penyedia fungsi, serta membatasi istilah
  \emph{genuine borrowed size} pada keadaan ketika kedua dimensi muncul
  bersama (Borrowed size: Scope, hlm. 18-19).
\item
  \textbf{Kontribusi konseptual sumber:} memperluas skala dari hubungan
  kota kecil-kota besar yang berdekatan menuju interaksi jaringan kota
  dari skala metropolitan sampai global, lalu menghubungkan hasil kota
  individual dengan eksternalitas jaringan kota positif dan negatif
  (Borrowed size: Scale dan From borrowed size to city network
  externalities, hlm. 19-20).
\item
  \textbf{Kontribusi kebijakan sumber:} mengingatkan bahwa koneksi,
  kedekatan, atau struktur polisentris tidak otomatis menghasilkan
  manfaat, bahwa kompetisi dapat mengalahkan peminjaman, dan bahwa tidak
  semua manfaat ukuran dapat dibagi (Network connectivity, hlm. 7-8;
  Conclusion, hlm. 18-20).
\item
  \textbf{Inferensi review:} artikel mengubah \emph{borrowed size} dari
  label deskriptif tentang kota kecil menjadi agenda pengukuran
  berlapis: hasil fungsi dan kinerja pada tingkat kota, mekanisme
  interaksi pada berbagai skala, serta saldo positif-negatif pada
  tingkat jaringan. Nilai tambah ini bersifat arsitektural dan teoretis;
  hanya lapisan fungsi kota yang diuji secara langsung dalam A01
  (Origins and meaning, hlm. 3; Conclusion, hlm. 18-21).
\end{itemize}

\subsection{Research Gap}\label{research-gap-16}

\textbf{Kesenjangan yang memotivasi studi.}

\begin{itemize}
\tightlist
\item
  Konsep yang diperkenalkan Alonso pada 1973 hanya mendapat sedikit
  perhatian selama beberapa dekade dan kemudian digunakan dalam
  penelitian serta kebijakan tanpa definisi analitis yang mapan (Origins
  and meaning, hlm. 3; Existing evidence, hlm. 4).
\item
  Bukti empiris terdahulu mengenai keuntungan kota kecil yang dekat
  dengan kota besar bersifat campuran dan jauh dari konklusif
  (Multicentric metropolitan areas, hlm. 5-6; Exploring the borrowed
  size concept, hlm. 8).
\item
  Penulis tidak menemukan penelitian yang secara eksplisit menguji
  pengaruh kedekatan dengan kota lebih kecil atau berorde lebih rendah
  sehingga kemungkinan kota besar meminjam basis dukungan dari kota
  kecil kurang diperhatikan (Small cities, hlm. 5).
\item
  Hubungan antara struktur multisenter atau polisentris dan hasil lokal
  belum dibedakan secara memadai dari hasil regional. Sebuah kota dapat
  melampaui ekspektasi ukurannya tanpa meningkatkan hasil agregat
  wilayah (Multicentric metropolitan areas, hlm. 5).
\item
  Literatur sering menekankan kedekatan atau aksesibilitas, sedangkan
  peran integrasi fungsional dan interaksi sebagai mekanisme belum cukup
  dijelaskan (Network connectivity, hlm. 7-8; Borrowed size and
  integration, hlm. 16).
\end{itemize}

\textbf{Kesenjangan yang tetap terbuka menurut sumber.}

\begin{itemize}
\tightlist
\item
  Pengujian empiris peminjaman kinerja masih belum dilakukan dalam A01;
  analisis hanya mengukur peminjaman fungsi (The missing link, hlm. 20).
\item
  Hubungan antara peminjaman fungsi, peminjaman kinerja, dan
  eksternalitas jaringan kota belum ditentukan (The missing link, hlm.
  20).
\item
  Tren peminjaman ukuran dari waktu ke waktu belum dieksplorasi (The
  missing link, hlm. 20).
\item
  Jenis eksternalitas aglomerasi yang dapat dipinjam atau dibagi serta
  skala spasial masing-masing belum dipetakan secara memadai (Borrowed
  size: Scope, hlm. 19; The missing link, hlm. 20).
\item
  Penggerak proses peminjaman ukuran, kemampuan eksternalitas menyebar
  dalam jarak lebih luas, dan kapasitas kota memanfaatkan integrasi
  serta konektivitas masih memerlukan analisis mendalam (The missing
  link, hlm. 20).
\item
  Saldo agregat fungsi dan kinerja pada tingkat jaringan kota masih
  diajukan sebagai pertanyaan, belum dijawab oleh hasil lokal A01 (From
  borrowed size to city network externalities, hlm. 20).
\end{itemize}

\textbf{Inferensi review.} A01 belum menutup kesenjangan identifikasi
antara ``kota mempunyai fungsi melebihi ekspektasi ukuran'' dan ``kota
memperoleh fungsi itu karena interaksi jaringan''. Residual menyediakan
indikator hasil, sedangkan data interaksi dan urutan kausal yang
diperlukan untuk menguji mekanisme belum disajikan (Borrowed size and
city size, hlm. 12; Borrowed size and integration, hlm. 16-18; The
missing link, hlm. 20).

\subsection{Limitations}\label{limitations-16}

\textbf{Keterbatasan yang dinyatakan atau diakui sumber.}

\begin{itemize}
\tightlist
\item
  Penulis menyebut analisisnya eksploratif dan bukti keseluruhan
  mengenai konsep masih jauh dari konklusif (Introduction, hlm. 2;
  Exploring the borrowed size concept, hlm. 8; Conclusion, hlm. 18).
\item
  Perlakuan empiris hanya berfokus pada peminjaman fungsi, bukan
  peminjaman kinerja (The missing link, hlm. 20).
\item
  Diskoneksi ukuran-fungsi tidak dengan sendirinya membuktikan bahwa
  kekurangan fungsi di satu kota dikompensasi oleh kelebihan di kota
  tetangga; alasan inilah yang mendorong model tambahan pada Table 2
  (Borrowed size, multicentricity and polycentricity, hlm. 11).
\item
  Keberadaan fungsi yang melampaui ekspektasi dapat dipengaruhi pola
  historis, politik, institusional, atau penggerak lain di luar proses
  ekonomi yang dicari. Penulis secara eksplisit menyinggung kemungkinan
  \emph{institutional borrowed size} dan kekuatan pendorong lain (hlm.
  10; Network connectivity, hlm. 8).
\item
  Tidak semua keuntungan aglomerasi dapat dipinjam. Kemudahan peminjaman
  bergantung pada jenis dan mobilitas eksternalitas, sedangkan ukuran
  kota juga membawa biaya yang tidak dimaksudkan pembuat kebijakan untuk
  dipinjam (Borrowed size: Scope, hlm. 19).
\item
  Perbandingan integrasi campuran pada Figure 4 harus dibaca hati-hati
  karena rasio tepat antara hubungan kuat, sedang, dan lemah dalam
  setiap kategori tidak diketahui (Borrowed size and integration, hlm.
  18).
\item
  Perbandingan lingkungan untuk kota sangat besar kekurangan daya
  statistik karena hanya lima kota sangat besar yang terisolasi
  (Borrowed size and city size, hlm. 15).
\item
  Tiga kota polisentris tidak masuk model OLS Table 2 karena data tidak
  memadai (Note pada Table 2, hlm. 13).
\item
  Metode rinci pengumpulan indikator fungsi diserahkan kepada BBSR
  (2011): Tidak disebutkan dalam file (Research approach, hlm. 8).
\end{itemize}

\textbf{Batas yang terlihat dari file dan diberi label sebagai inferensi
review.}

\begin{itemize}
\tightlist
\item
  Data fungsi mencakup 2004-2009 dan mayoritas merujuk 2008, tetapi
  model tidak memanfaatkan dimensi waktu. Analisis potong lintang tidak
  menunjukkan apakah integrasi mendahului surplus fungsi atau sebaliknya
  (Research approach, hlm. 8-10).
\item
  Residual positif merupakan hasil relatif terhadap spesifikasi model.
  Perubahan bentuk fungsi, variabel kontrol, ambang klasifikasi, atau
  kualitas pengukuran dapat mengubah kota yang masuk kategori
  \emph{borrowed size}. Penulis melaporkan satu uji kontrol alternatif
  dengan hasil serupa, tetapi analisis sensitivitas lain: Tidak
  disebutkan dalam file (Notes 5-7, hlm. 21).
\item
  Kategori integrasi dibangun dari delimitasi yang sebagian berbasis
  morfologi, jarak, cekungan komuter, dan tumpang tindih jangkauan 45
  menit. Ukuran arus aktual antarkota yang sebanding untuk semua
  kategori: Tidak disebutkan dalam file (Research approach, hlm. 9-10;
  Borrowed size and integration, hlm. 16-18).
\item
  Ambang HHI 0,56 dipilih sebagai batas yang dianggap cukup baik. Uji
  sensitivitas hasil terhadap ambang polisentrisitas lain: Tidak
  disebutkan dalam file (Research approach, hlm. 10).
\item
  Model menyertakan variabel semu negara dan beberapa kontrol, tetapi
  strategi untuk menyingkirkan seluruh perancu atau ketergantungan
  spasial antarkota: Tidak disebutkan dalam file. Koefisien karenanya
  tidak dapat dibaca sebagai efek kausal (Tables 2-3, hlm. 13-15; Notes
  4-6, hlm. 21).
\item
  Tahun dan sumber penduduk kota, transportasi, serta PDB per kapita:
  Tidak disebutkan dalam file. Keselarasan temporal variabel-variabel
  itu dengan data fungsi 2004-2009 tidak dapat diperiksa dari TXT
  (Research approach, hlm. 8-10; Note 4, hlm. 21).
\item
  Hitungan kota terisolasi dan bertetangga berbeda antara narasi, Table
  1, dan Figure 4. Alasan perbedaan tersebut: Tidak disebutkan dalam
  file. Karena itu, ukuran subkelompok harus dilaporkan sesuai locator
  masing-masing, bukan disatukan secara paksa (hlm. 10; Table 1, hlm.
  11; Figure 4, hlm. 17).
\item
  Sampel hanya mencakup kota di EU25, Norwegia, dan Swiss. Validitas
  empiris definisi untuk sistem perkotaan di luar cakupan tersebut tidak
  diuji dalam A01 (Research approach, hlm. 9).
\item
  Tinjauan literatur tidak melaporkan protokol pencarian atau penilaian
  mutu sehingga cakupan dan potensi bias seleksi literatur tidak dapat
  dinilai dari file (Introduction, hlm. 2; Existing evidence of
  `borrowed size', hlm. 4-8).
\item
  Figure 2, Figure 4, dan Figure 5 tidak terekstraksi sebagai data
  numerik lengkap. Review dapat memakai uraian penulis dan caption,
  tetapi tidak dapat memeriksa semua elemen visualnya dari TXT (Figures
  2, 4, dan 5, hlm. 15, 17, dan 19).
\end{itemize}

\subsection{Challenges}\label{challenges-16}

\begin{itemize}
\tightlist
\item
  \textbf{Laporan sumber:} tantangan konseptual pertama ialah menentukan
  karakteristik yang biasanya berkaitan dengan ukuran dan ukuran posisi
  kota terhadap kota lain. Tanpa hubungan dasar ukuran-karakteristik,
  diskoneksi yang menjadi inti operasionalisasi tidak dapat dikenali
  (Analytical concept: Operationalising borrowed size, hlm. 4).
\item
  \textbf{Laporan sumber:} tantangan analitis berikutnya ialah
  memisahkan peminjaman ukuran dari bayang-bayang aglomerasi. Keduanya
  sama-sama melemahkan hubungan ukuran-fungsi, tetapi mempunyai arah
  residual dan implikasi yang berlawanan (Borrowed size, multicentricity
  and polycentricity, hlm. 10-12).
\item
  \textbf{Interpretasi penulis:} kedekatan bukan mekanisme yang memadai
  jika tidak menghasilkan interaksi. Pada saat yang sama, interaksi
  dapat membawa akses terhadap manfaat maupun paparan terhadap kompetisi
  sehingga konektivitas tidak selalu menghasilkan hasil positif (Network
  connectivity, hlm. 7-8).
\item
  \textbf{Laporan sumber:} delimitasi integrasi saling bertumpang
  tindih; sebuah kota dapat mempunyai hubungan pada beberapa skala
  sekaligus. Hal ini menyulitkan isolasi pengaruh satu tingkat integrasi
  dan menyebabkan kategori kuat-lemah dikeluarkan dari interpretasi
  (Borrowed size and integration, hlm. 16-18; Figure 4, hlm. 17).
\item
  \textbf{Interpretasi penulis:} skala hasil merupakan tantangan
  penting. Surplus fungsi di satu kota dapat diimbangi defisit kota lain
  sehingga keberhasilan lokal belum tentu menjadi eksternalitas positif
  bagi jaringan secara keseluruhan (Multicentric metropolitan areas,
  hlm. 5; From borrowed size to city network externalities, hlm. 20).
\item
  \textbf{Inferensi review:} tantangan pengukuran terletak pada
  penggunaan fungsi yang mudah dilokasikan sebagai indikator hasil,
  sementara interaksi yang diduga memungkinkannya tidak diobservasi
  secara langsung dalam satu ukuran jaringan yang konsisten (Research
  approach, hlm. 8-10; Borrowed size and integration, hlm. 16-18).
\item
  \textbf{Inferensi review:} tantangan komunikasi kebijakan ialah
  mencegah definisi yang diperluas menjadi terlalu elastis. Jika fungsi,
  kinerja, berbagai arah peminjaman, banyak skala, dan hasil jaringan
  digabungkan tanpa indikator terpisah, konsep berisiko kembali sulit
  diuji, persis seperti masalah awal yang dikritik artikel
  (Introduction, hlm. 2; Conclusion, hlm. 18-20).
\end{itemize}

\subsection{Practical Implications}\label{practical-implications-16}

\begin{itemize}
\tightlist
\item
  \textbf{Implikasi penulis:} pembuat kebijakan tidak dapat menganggap
  jaringan kota kecil secara otomatis setara dengan satu kota besar.
  Kedekatan, multisenter, atau label polisentris tidak cukup; manfaat
  bergantung pada interaksi dan tingkat integrasi, sedangkan persaingan
  dapat menghasilkan bayang-bayang aglomerasi (hlm. 10-18).
\item
  \textbf{Implikasi penulis:} kebijakan tidak semestinya hanya diarahkan
  kepada kota kecil. Kota besar juga dapat meminjam basis dukungan dan
  fungsi melalui kota lain, dan hasil regresi menunjukkan peluang
  peminjaman fungsi yang lebih tinggi pada kelas kota lebih besar (Table
  3, hlm. 14-16; Borrowed size: Scale, hlm. 20).
\item
  \textbf{Implikasi penulis:} fungsi dan kinerja perlu diukur terpisah.
  Keberadaan fungsi yang melampaui ekspektasi ukuran tidak otomatis
  membuktikan kinerja yang lebih tinggi, sedangkan akses terhadap fungsi
  kota lain dapat meningkatkan kinerja tanpa memindahkan fungsi ke kota
  peminjam (Borrowed size: Scope, hlm. 18-19).
\item
  \textbf{Implikasi penulis:} kebijakan regional sebaiknya menilai saldo
  pada tingkat jaringan, yaitu apakah peminjaman ukuran secara agregat
  lebih besar daripada bayang-bayang aglomerasi. Penulis menawarkan
  eksternalitas jaringan kota positif sebagai sasaran yang lebih tepat
  (From borrowed size to city network externalities, hlm. 20).
\item
  \textbf{Implikasi penulis:} jenis fungsi dan eksternalitas perlu
  dibedakan karena tidak semuanya dapat dipinjam pada skala yang sama.
  Pembuat kebijakan juga harus mengingat bahwa ukuran membawa biaya,
  bukan hanya manfaat (Borrowed size: Scope dan Scale, hlm. 19-20).
\item
  \textbf{Inferensi review berdasarkan hasil:} evaluasi kebijakan perlu
  ditabulasi silang menurut kelas ukuran dan lingkungan kota. Koefisien
  agregat lingkungan yang tidak signifikan menutupi perbedaan antara
  kota menengah polisentris dan kota besar polisentris pada model
  interaksi (Table 3, hlm. 14-15).
\item
  \textbf{Inferensi review:} indikator surplus fungsi sebaiknya
  diperlakukan sebagai diagnosis awal, lalu diuji dengan data arus
  antarkota dan perubahan kinerja sebelum dipakai untuk menyimpulkan
  bahwa kebijakan konektivitas menyebabkan \emph{borrowed size}. Langkah
  lanjutan ini konsisten dengan kesenjangan yang diakui penulis, tetapi
  tidak dilaksanakan dalam A01 (uraian residual, hlm. 12; The missing
  link, hlm. 20).
\end{itemize}

\subsection{Applications}\label{applications-16}

\begin{itemize}
\tightlist
\item
  \textbf{Aplikasi yang diperagakan sumber:} membandingkan kuat-lemahnya
  hubungan ukuran-fungsi pada kota terisolasi, kota bertetangga, serta
  kota dalam wilayah monosentris dan polisentris (Table 1, hlm. 11).
\item
  \textbf{Aplikasi yang diperagakan sumber:} menguji apakah penduduk dan
  fungsi di kota-kota lain dalam polyFUA berkaitan dengan pasokan fungsi
  suatu kota setelah karakteristik lokal dan negara dikendalikan (Table
  2, hlm. 11-13).
\item
  \textbf{Aplikasi yang diperagakan sumber:} mengklasifikasikan kota
  sebagai meminjam fungsi, berada dalam bayang-bayang aglomerasi, atau
  tidak mempunyai fungsi yang diharapkan berdasarkan residual model
  (Borrowed size and city size, hlm. 12).
\item
  \textbf{Aplikasi yang diperagakan sumber:} mengidentifikasi
  heterogenitas peluang peminjaman ukuran menurut kelas kota dan
  lingkungan metropolitan melalui interaksi regresi logistik (Table 3,
  hlm. 14-15).
\item
  \textbf{Aplikasi yang diperagakan sumber:} membandingkan surplus atau
  defisit fungsi menurut hubungan kuat, sedang, dan lemah pada beberapa
  skala spasial (Figure 4 dan uraian, hlm. 16-18).
\item
  \textbf{Aplikasi konseptual yang diajukan sumber:} memakai kerangka
  fungsi-kinerja dan peminjaman-bayang-bayang untuk mengevaluasi
  eksternalitas positif maupun negatif pada tingkat jaringan kota, bukan
  hanya keberhasilan satu kota (Figure 5 dan uraian, hlm. 18-20).
\item
  \textbf{Inferensi review:} kerangka dapat diterapkan pada sistem kota
  lain hanya jika hubungan normal antara ukuran dan hasil ditetapkan
  terlebih dahulu, fungsi serta kinerja tidak dicampur, interaksi
  jaringan dapat diukur, dan klasifikasi spasial dikalibrasi terhadap
  konteks setempat. A01 sendiri tidak menguji pemindahan metode ke luar
  cakupan Eropa yang ditelitinya (Analytical concept: Operationalising
  borrowed size, hlm. 4; Research approach, hlm. 8-10; The missing link,
  hlm. 20).
\end{itemize}

\subsection{Future Research
(Optional)}\label{future-research-optional-16}

Sumber menyatakan agenda penelitian berikut secara eksplisit pada bagian
\emph{The missing link}, hlm. 20:

\begin{enumerate}
\def\labelenumi{\arabic{enumi}.}
\tightlist
\item
  Menguji peminjaman kinerja karena analisis empiris A01 hanya berfokus
  pada peminjaman fungsi.
\item
  Menentukan hubungan antara peminjaman fungsi, peminjaman kinerja, dan
  eksternalitas jaringan kota.
\item
  Menelusuri tren peminjaman ukuran dari waktu ke waktu.
\item
  Menentukan jenis eksternalitas aglomerasi yang dapat dipinjam atau
  dibagi dalam jaringan serta skala spasial tempat setiap proses
  berlangsung.
\item
  Mengidentifikasi penggerak peminjaman ukuran melalui analisis mendalam
  mengenai kemampuan berbagai eksternalitas menyebar pada wilayah yang
  lebih luas.
\item
  Menganalisis faktor yang menentukan kapasitas kota untuk memanfaatkan
  integrasi dan konektivitas.
\item
  Mengukur agregat hasil lokal pada tingkat jaringan kota, baik dalam
  fungsi maupun kinerja, agar dominasi peminjaman ukuran atau
  bayang-bayang aglomerasi dapat ditentukan (From borrowed size to city
  network externalities, hlm. 20).
\end{enumerate}

Agenda di luar butir-butir yang dinyatakan artikel: Tidak ditambahkan
dalam review ini.

\subsection{Provenance Notes}\label{provenance-notes-16}

\begin{itemize}
\tightlist
\item
  Review ini hanya menggunakan seluruh isi file
  \texttt{/Users/mac/Documents/Mac/{[}2{]}\ Obsidian\ Vault/zettelkasten/source/journal/A01-meijers-burger-2017-stretching-the-concept-of-borrowed-size-10.1177-0042098015597642.txt},
  sebanyak 1.214 baris. Tidak ada penelusuran web, basis data
  bibliografis, PDF lain, catatan argumentasi, atau pengetahuan luar
  yang digunakan untuk menambah klaim substantif.
\item
  Blok \texttt{{[}PDF\ Metadata{]}} menyatakan bahwa TXT diekstraksi
  pada 2026-08-31 dengan \texttt{pdftotext\ -layout} dari
  \texttt{journal/A01-meijers-burger-2017-stretching-the-concept-of-borrowed-size-10.1177-0042098015597642.pdf}.
  PDF mempunyai 23 halaman, tidak terenkripsi, dan tidak bertanda
  struktur (\texttt{Tagged:\ no}). Review ini membaca TXT, bukan PDF
  asal.
\item
  Judul, penulis, afiliasi, jurnal, DOI, hak cipta 2015, serta tanggal
  penerimaan naskah berasal dari halaman pembuka. Tahun 2017 hanya
  tampak pada nama file. Volume, nomor terbitan, tahun publikasi final
  yang tertulis dalam artikel, dan status telaah sejawat: Tidak
  disebutkan dalam file.
\item
  Locator halaman mengikuti nomor 1-23 yang muncul pada header hasil
  ekstraksi. Nama bagian, Table 1-3, Figure 1-5, dan Notes 1-7
  dipertahankan sebagai locator tambahan ketika tersedia.
\item
  Teks menyertakan daftar pustaka pada hlm. 21-23, tetapi isi semua
  karya yang dirujuk tidak diakses secara mandiri. Pernyataan tentang
  Alonso, Phelps, BBSR, ESPON, dan penelitian lain diperlakukan sebagai
  laporan A01 atas literatur.
\item
  Tata letak dua kolom dan tabel lanskap menimbulkan distorsi. Pada
  Table 2, tanda minus tampil sebagai awalan \texttt{2}, misalnya
  \texttt{20.063}, dan operator signifikansi pada beberapa tempat tampil
  sebagai garis miring terbalik. Review menggunakan arah hubungan dan
  tingkat signifikansi hanya ketika didukung pula oleh uraian naratif
  atau penanda tabel yang jelas; angka yang terdistorsi tidak
  direkonstruksi secara diam-diam.
\item
  Catatan Table 3 hasil ekstraksi menggunakan simbol yang sama untuk dua
  taraf signifikansi (\texttt{**p} untuk 0,01 dan kembali \texttt{**p}
  untuk 0,05), sedangkan tubuh teks membedakan
  \texttt{p\ \textless{}\ 0,01} dan \texttt{p\ \textless{}\ 0,05}.
  Review mengikuti taraf yang dinyatakan dalam tubuh teks untuk
  perbandingan yang dibahas penulis dan tidak memperbaiki legenda tabel
  berdasarkan pengetahuan luar.
\item
  Figure 1 menyediakan skema delimitasi, Figure 2 berupa diagram
  lingkaran klasifikasi kota, Figure 3 menggambarkan tingkat integrasi,
  Figure 4 memuat residual rata-rata, dan Figure 5 memuat dimensi
  teoretis \emph{borrowed size}. Grafik dan label numerik tidak
  seluruhnya terekstraksi; review tidak menebak nilai visual yang
  hilang.
\item
  Jumlah kota menurut lingkungan tidak konsisten di dalam TXT: narasi
  hlm. 10 melaporkan 651 kota terisolasi dan 1.316 bertetangga, Table 1
  melaporkan 648 dan 1.319, sedangkan Figure 4 melaporkan 649 tanpa
  tetangga. Review mempertahankan ketiganya beserta locator dan tidak
  mengarang rekonsiliasi.
\item
  Nilai \texttt{adjusted\ R2\ =\ 0,800} pada Note 5 berkaitan dengan
  regresi pembentuk residual yang hanya memasukkan ukuran dan variabel
  semu negara, sedangkan \texttt{R2\ =\ 0,612} pada Table 2 berkaitan
  dengan model OLS fungsi agregat untuk 606 kota polisentris. Review
  tidak menyamakan kedua statistik tersebut.
\item
  Batas provenance utama: review dapat memeriksa pernyataan yang
  tersedia dalam ekstraksi, tetapi tidak dapat memastikan kesetiaan tata
  letak terhadap PDF, memulihkan nilai yang hanya terdapat dalam grafik,
  memverifikasi sumber data asli, mengaudit korpus tinjauan, atau
  mereplikasi model tanpa data dan kode yang tidak disertakan.
\end{itemize}

\section{Review A02}\label{review-a02}

\textbf{Identitas sumber}

\begin{itemize}
\tightlist
\item
  Kode: A02.
\item
  Judul: \emph{Borrowing size in networks of cities: City size, network
  connectivity and metropolitan functions in Europe}.
\item
  Penulis: Evert J. Meijers, Martijn J. Burger, dan Marloes M.
  Hoogerbrugge.
\item
  Afiliasi: Meijers dan Hoogerbrugge berafiliasi dengan Faculty of
  Architecture and the Built Environment, Delft University of
  Technology; Burger berafiliasi dengan Department of Applied Economics
  dan Erasmus Happiness Economics Research Organisation, Erasmus
  University Rotterdam, serta Tinbergen Institute; Hoogerbrugge juga
  berafiliasi dengan Erasmus Happiness Economics Research Organisation,
  Erasmus University Rotterdam (halaman judul, hlm. 181).
\item
  Jenis sumber menurut tubuh teks: artikel jurnal. Status telaah
  sejawat: Tidak disebutkan dalam file.
\item
  Jurnal: \emph{Papers in Regional Science}.
\item
  Volume, nomor, dan waktu terbit: Volume 95, Nomor 1, Maret 2016
  (footer artikel, hlm. 181-198).
\item
  Rentang halaman artikel yang terindikasi: 181-198. Nomor 181 tidak
  tercetak pada hasil ekstraksi halaman judul, tetapi halaman berikutnya
  bernomor 182 dan daftar pustaka berakhir pada hlm. 198.
\item
  DOI: 10.1111/pirs.12181.
\item
  Riwayat naskah: diterima 20 Oktober 2014 dan disetujui 23 Juni 2015
  (halaman judul, hlm. 181).
\item
  Klasifikasi JEL: R12 dan R11.
\item
  Kata kunci: \emph{urbanization economies}, \emph{network
  externalities}, \emph{borrowed size}, \emph{agglomeration shadow}, dan
  \emph{connectivity}.
\item
  Pendanaan yang diakui penulis: Netherlands Organisation for Scientific
  Research (NWO) dan Platform31 (catatan kaki halaman judul, hlm. 181).
\item
  Cakupan studi: 1.114 aglomerasi perkotaan pada 16 negara, yaitu
  negara-negara bekas EU15 selain UK, ditambah Norwegia dan Swiss
  (Bagian 3.2, hlm. 187).
\end{itemize}

\textbf{Penanda epistemik yang digunakan}

\begin{itemize}
\tightlist
\item
  \textbf{Laporan sumber} berarti pertanyaan, konsep, metode, data,
  angka, atau hasil yang disampaikan langsung dalam artikel.
\item
  \textbf{Interpretasi penulis} berarti penjelasan mekanisme, makna,
  atau implikasi yang diberikan penulis terhadap konteks atau hasil
  empirisnya.
\item
  \textbf{Inferensi review} berarti pembacaan analitis yang dibuat hanya
  dari TXT A02. Inferensi ini tidak diperlakukan sebagai pernyataan
  penulis ataupun bukti tambahan.
\end{itemize}

\subsection{Summarized Abstract}\label{summarized-abstract-17}

\textbf{Laporan sumber.} Artikel berangkat dari ketegangan antara teori
yang menempatkan manfaat aglomerasi kota besar sebagai penggerak
pertumbuhan dan dinamika Eropa Barat yang relatif stabil tingkat
urbanisasinya, tidak mengalami kenaikan pangsa penduduk kota besar,
serta memperlihatkan kinerja kota kecil dan menengah yang semakin baik.
Penulis menguji apakah eksternalitas aglomerasi tidak lagi terbatas pada
batas kota, tetapi dapat dibagi melalui jaringan kota dalam bentuk
\emph{borrowed size}, sekaligus dapat menimbulkan \emph{agglomeration
shadow} akibat persaingan. Pertanyaan utamanya ialah sejauh mana
konektivitas jaringan kota berkaitan dengan tingkat fungsi metropolitan
yang lebih tinggi dan seberapa penting konektivitas tersebut
dibandingkan faktor lokal (Abstract dan Bagian 1, hlm. 181-183).

\textbf{Laporan sumber.} Penelitian menggabungkan data fungsi
metropolitan, delimitasi perkotaan ESPON, ukuran populasi, aksesibilitas
jaringan, dan variabel kontrol menjadi dataset 1.114 \emph{morphological
urban areas} (MUA). Kehadiran fungsi metropolitan agregat dan per jenis
dianalisis dengan regresi beta terinflasi nol dan satu. Ukuran utama
penjelas ialah populasi lokal, populasi yang dapat dicapai dalam 45
menit dengan mobil sebagai konektivitas regional, serta potensi
aksesibilitas populasi melalui jaringan jalan, rel, dan udara sebagai
konektivitas (inter)nasional. Analisis juga membandingkan efek marginal,
menguji fungsi metropolitan satu per satu, dan memasukkan interaksi
ukuran-konektivitas (Bagian 3, hlm. 186-190; Bagian 4.3-4.4, hlm.
191-195).

\textbf{Laporan sumber.} Ukuran lokal dan konektivitas (inter)nasional
dilaporkan berhubungan positif dengan kehadiran fungsi metropolitan,
tetapi ukuran lokal secara agregat kira-kira 2,5 kali lebih penting.
Konektivitas regional secara umum tidak meningkatkan fungsi metropolitan
dan untuk beberapa fungsi menunjukkan hubungan negatif. Ukuran lokal
signifikan bagi semua jenis fungsi, sedangkan konektivitas
(inter)nasional terutama penting bagi fungsi perusahaan, institusi
internasional, dan sains, tetapi umumnya tidak bagi fungsi budaya dan
olahraga. Kota kecil lebih memperoleh tambahan fungsi dari pertambahan
ukuran, sedangkan kota besar lebih mampu memanfaatkan konektivitas
regional maupun (inter)nasional (Bagian 4, hlm. 189-195; Bagian 5, hlm.
195-196).

\textbf{Interpretasi penulis.} Hubungan positif pada skala
(inter)nasional ditafsirkan sebagai bukti bahwa kota dapat meminjam
ukuran melalui jaringan. Sebaliknya, hubungan regional yang tidak
signifikan atau negatif ditafsirkan sebagai kemungkinan dominasi
bayang-bayang aglomerasi akibat persaingan antarkota yang berdekatan.
Penulis menyimpulkan bahwa ekonomi jaringan kota melengkapi, bukan
menghapus, peran massa lokal (Bagian 4.1-4.4, hlm. 189-195; Bagian 5,
hlm. 195-196).

\textbf{Inferensi review.} Bukti empiris artikel menunjukkan asosiasi
bersyarat antara proksi konektivitas dan lokasi fungsi metropolitan.
Bukti tersebut belum mengidentifikasi aliran manfaat jaringan secara
langsung dan tidak membuktikan bahwa konektivitas menyebabkan
\emph{borrowed size} atau \emph{agglomeration shadow}. Batas ini sejalan
dengan pengakuan penulis mengenai kausalitas terbalik (Bagian 5, hlm.
196).

\subsection{Problem Statement}\label{problem-statement-17}

\textbf{Laporan sumber atas konteks.} Eropa, terutama bagian padat dari
bekas EU15, tidak mengikuti logika pertumbuhan kota besar secara linear.
Lebih dari setengah penduduk perkotaan Uni Eropa dilaporkan tinggal di
kota kecil dan menengah berpenduduk 5.000-100.000 jiwa, sedangkan hanya
7\% warga Uni Eropa tinggal di megakota berpenduduk setidaknya 5 juta
jiwa. Tingkat urbanisasi Eropa Barat sudah lama tinggi; kota besar tidak
tumbuh lebih cepat daripada kota kecil; kontribusi kawasan metropolitan
besar terhadap PDB hanya meningkat tipis; dan kota lapis kedua
dilaporkan mengungguli ibu kota dalam pertumbuhan PDB pada sebagian
besar bekas negara EU15 selama 2000-2007 (Bagian 1, hlm. 181-182).
Angka-angka tersebut dilaporkan artikel dari literatur yang dirujuk,
bukan diestimasi ulang dalam studi A02.

\textbf{Interpretasi penulis atas masalah teoretis.} Pola tersebut
dianggap bertentangan dengan teori aglomerasi dan model \emph{new
economic geography} yang cenderung memprediksi konsentrasi pada kota
yang lebih sedikit dan lebih besar. Penulis tidak menolak manfaat
aglomerasi, tetapi mempertanyakan anggapan bahwa manfaat tersebut harus
terkurung di dalam batas kota. Keterhubungan regional, nasional, dan
internasional mungkin memungkinkan kota yang lebih kecil mengompensasi
kekurangan massa lokal (Bagian 1, hlm. 182-183).

\textbf{Laporan sumber.} Pengetahuan tentang jangkauan spasial manfaat
dan biaya aglomerasi masih terbatas. Literatur telah menyatakan bahwa
efek tersebut melemah terhadap jarak dan jangkauannya mungkin meluas,
tetapi belum jelas apakah ekonomi urbanisasi dapat dibagi antarkota.
Kajian terdahulu juga lebih banyak menaruh perhatian pada konektivitas
internasional, sedangkan konsep \emph{borrowed size} pada skala regional
mendapat perhatian terbatas dan kemungkinan efek persaingan jaringan
perlu dipertimbangkan (Bagian 2, hlm. 184-186).

\textbf{Rumusan masalah sumber.} Pertanyaan penelitian dinyatakan secara
eksplisit: sejauh mana konektivitas jaringan kota diterjemahkan menjadi
tingkat fungsi metropolitan yang lebih tinggi, dan seberapa penting
konektivitas tersebut dibandingkan faktor lokal? Pertanyaan ini
diperinci menurut skala jaringan, jenis fungsi metropolitan, dan kelas
ukuran kota (Bagian 1, hlm. 183).

\textbf{Inferensi review.} Masalah operasionalnya ialah memisahkan peran
massa yang berada di dalam MUA dari potensi populasi yang dapat dicapai
melalui jaringan. Karena desainnya potong lintang dan konektivitas tidak
diinstrumentasi, penelitian memisahkan keduanya secara statistik, bukan
mengidentifikasi proses kausal atau pertukaran fungsi antarkota secara
langsung (Bagian 3.1-3.3, hlm. 186-189; Bagian 5, hlm. 196).

\subsection{Objectives}\label{objectives-17}

\begin{itemize}
\tightlist
\item
  \textbf{Laporan sumber:} menilai sejauh mana konektivitas jaringan
  kota berhubungan dengan tingkat fungsi metropolitan yang lebih tinggi
  dan membandingkan kepentingannya dengan faktor lokal, terutama ukuran
  populasi kota (Bagian 1, hlm. 183; Bagian 3.1, hlm. 186).
\item
  \textbf{Laporan sumber:} membedakan konektivitas jaringan regional
  dari konektivitas nasional dan internasional untuk mengetahui apakah
  skala jaringan mengubah hasil (Bagian 1, hlm. 183; Bagian 3.1-3.2,
  hlm. 186-188).
\item
  \textbf{Laporan sumber:} menguji apakah pengaruh ukuran dan
  konektivitas berbeda antara fungsi institusi internasional,
  perusahaan, sains, budaya, dan olahraga (Bagian 1, hlm. 183; Bagian
  3.1-3.2, hlm. 186-187).
\item
  \textbf{Laporan sumber:} menguji apakah kota dari ukuran yang berbeda
  memperoleh manfaat yang berbeda dari pertambahan ukuran lokal dan
  konektivitas jaringan (Bagian 1, hlm. 183; Bagian 4.4, hlm. 193-195).
\item
  \textbf{Laporan sumber:} mengeksplorasi \emph{borrowed size} sebagai
  pengaruh positif konektivitas dan \emph{agglomeration shadow} sebagai
  pengaruh negatif konektivitas terhadap fungsi metropolitan (Bagian 2,
  hlm. 185-186; Bagian 3.1, hlm. 186).
\item
  \textbf{Inferensi review:} tujuan empiris artikel adalah menilai
  \emph{borrowed function}, yaitu kehadiran fungsi yang melampaui apa
  yang diasosiasikan dengan ukuran lokal, bukan mengukur \emph{borrowed
  performance} seperti pertumbuhan, produktivitas, atau kesejahteraan.
  Pembedaan fungsi dan kinerja dibahas dalam tinjauan konsep, sedangkan
  variabel hasil studi hanya mengukur fungsi (Bagian 2, hlm. 185; Bagian
  3.1, hlm. 186).
\end{itemize}

\subsection{Literature Survey}\label{literature-survey-17}

\textbf{Inferensi review atas sifat survei.} Artikel menyajikan tinjauan
naratif untuk membangun kerangka konseptual dan operasionalisasi.
Strategi pencarian, basis data bibliografis, periode pencarian, kriteria
inklusi-eksklusi, jumlah studi yang disaring, dan penilaian mutu
literatur: Tidak disebutkan dalam file. Karena itu, bagian ini tidak
dapat diperlakukan sebagai tinjauan sistematis.

\begin{itemize}
\tightlist
\item
  \textbf{Laporan sumber atas ekonomi urbanisasi:} kota yang lebih
  besar, padat, dan beragam memberi pengurangan biaya, peningkatan
  keluaran, serta manfaat utilitas melalui pasar input dan tenaga kerja
  yang lebih besar, infrastruktur dan fasilitas yang lebih baik, jasa
  bisnis khusus, institusi pengetahuan, transmisi informasi, konsumsi,
  kontak tatap muka, dan keragaman industri. Artikel berfokus pada
  ekonomi urbanisasi, bukan klaster perusahaan terspesialisasi (Bagian
  2, hlm. 183-184).
\item
  \textbf{Laporan sumber atas mekanisme aglomerasi:} Duranton dan Puga
  serta Puga dirujuk untuk tiga saluran berupa \emph{sharing},
  \emph{matching}, dan \emph{learning}. A02 terutama menguji
  \emph{sharing} melalui keberadaan fasilitas yang tidak terbagi,
  pemasok, dan fungsi metropolitan; \emph{matching} dan \emph{learning}
  tidak menjadi fokus empiris utama (Bagian 1, hlm. 183; Bagian 2, hlm.
  184).
\item
  \textbf{Laporan sumber atas biaya aglomerasi:} konsentrasi pekerjaan
  dan penduduk dapat meningkatkan kemacetan, persaingan lokasi pusat,
  polusi, kriminalitas, dan skala masalah sosial. Artikel melaporkan
  argumen literatur bahwa kota lebih kecil mungkin lebih mampu
  mengendalikan biaya sosial, ekonomi, dan lingkungan tersebut (Bagian
  2, hlm. 184).
\item
  \textbf{Laporan sumber atas perspektif jaringan:} konsep eksternalitas
  jaringan perkotaan menekankan bahwa interaksi antarkota dapat
  memengaruhi kinerja tempat yang terhubung. Jaringan perusahaan, modal,
  pengetahuan, manusia, dan barang dapat melengkapi, bahkan dalam
  beberapa keadaan menggantikan, ekonomi aglomerasi lokal dalam
  memfasilitasi \emph{sharing}, \emph{matching}, dan \emph{learning}
  (Bagian 2, hlm. 184-185).
\item
  \textbf{Laporan sumber atas \emph{borrowed size}:} Alonso dirujuk
  untuk gagasan bahwa kota kecil di dekat konsentrasi penduduk lain
  dapat memperlihatkan karakteristik kota yang lebih besar, memperoleh
  sebagian manfaat aglomerasi tetangganya, dan menghindari sebagian
  biayanya. Artikel memperluas kedekatan regional tersebut ke jaringan
  kota nasional dan internasional (Bagian 2, hlm. 185).
\item
  \textbf{Laporan sumber atas pembedaan konsep:} Meijers dan Burger
  dirujuk untuk membedakan \emph{borrowed function} dari \emph{borrowed
  performance}. Artikel lain yang dirangkum menyatakan bahwa kota yang
  relatif lebih besar dalam suatu wilayah lebih mampu menampung fungsi,
  sedangkan akses ke fungsi dan jaringan kota peringkat pertama dapat
  meningkatkan manfaat lokasi kota lapis kedua meskipun fungsi atau
  jaringan itu tidak hadir secara lokal (Bagian 2, hlm. 185).
\item
  \textbf{Laporan sumber atas \emph{agglomeration shadow}:} model
  \emph{new economic geography} memprediksi bahwa pertumbuhan dekat kota
  bertingkat lebih tinggi dapat dibatasi oleh persaingan. Artikel juga
  merangkum bukti yang mempertanyakan efek tersebut serta bukti bahwa
  kota besar dapat mengeksploitasi basis dukungan kota tetangga yang
  lebih kecil, sekalipun kota kecil secara bersamaan memperoleh akses ke
  manfaat kota besar (Bagian 2, hlm. 185-186).
\item
  \textbf{Kerangka konseptual penulis:} \emph{borrowed size} dan
  \emph{agglomeration shadow} diposisikan sebagai dua sisi eksternalitas
  jaringan. Pengaruh positif konektivitas terhadap kehadiran ekonomi
  urbanisasi disebut \emph{borrowed size}, sedangkan pengaruh negatif
  disebut \emph{agglomeration shadow} (Bagian 2, hlm. 185-186).
\end{itemize}

\textbf{Inferensi review.} Definisi operasional tersebut membuat tanda
asosiasi konektivitas menjadi penentu klasifikasi: positif dibaca
sebagai peminjaman ukuran dan negatif sebagai bayang-bayang. Mekanisme
berupa pembagian sumber daya atau persaingan tidak diamati secara
terpisah. Karena itu, istilah konseptual dalam hasil bergantung pada
interpretasi tanda model, bukan pengukuran langsung transaksi, arus
pengguna, atau perpindahan fungsi (Bagian 2, hlm. 185-186; Bagian 3.1,
hlm. 186).

\subsection{Method Used}\label{method-used-17}

\textbf{Laporan sumber.} Penelitian menggunakan dataset tingkat MUA
untuk menjelaskan kehadiran fungsi metropolitan melalui ukuran lokal,
konektivitas jaringan pada dua skala, dan variabel kontrol. Penulis
menyebut analisisnya sebagai eksplorasi empiris dan pada bagian akhir
menegaskan bahwa hasil harus dibaca sebagai asosiasi bersyarat, bukan
hubungan kausal (Bagian 3, hlm. 186-189; Bagian 5, hlm. 196).

\textbf{Inferensi review atas rancangan.} Susunan data dan model
merupakan analisis kuantitatif potong lintang: fungsi metropolitan dari
satu tahun per indikator digabungkan dengan ukuran aksesibilitas dan
kovariat untuk membandingkan kota, bukan mengikuti perubahan kota dari
waktu ke waktu (Bagian 3.2, hlm. 186-188; Bagian 5, hlm. 196).

\textbf{Pengukuran konstruk utama.}

\begin{itemize}
\tightlist
\item
  Ukuran lokal diukur dengan jumlah penduduk kota (Bagian 3.1, hlm.
  186).
\item
  Konektivitas regional diukur dengan jumlah penduduk di luar kota itu
  sendiri yang dapat dicapai dalam waktu 45 menit dengan mobil. Wilayah
  isokron ini disebut \emph{potential urban strategic horizon} (PUSH)
  (Bagian 3.1-3.2, hlm. 186 dan 187).
\item
  Konektivitas (inter)nasional diukur dengan potensi aksesibilitas
  terhadap populasi di luar FUA melalui jalan, rel, dan udara, dengan
  populasi tujuan dibobot menurut waktu perjalanan (Bagian 3.2, hlm.
  187-188).
\item
  Ekonomi urbanisasi diproksikan dengan kehadiran fungsi metropolitan
  dalam domain institusi internasional, perusahaan, sains, budaya, dan
  olahraga. Fokusnya terutama pada saluran \emph{sharing} (Bagian 1,
  hlm. 183; Bagian 3.1, hlm. 186).
\item
  Variabel kontrol meliputi PDB per kapita pada skala FUA, variabel
  boneka ibu kota, indeks pariwisata, dan variabel boneka negara. Indeks
  pariwisata menggabungkan bintang dalam panduan Michelin dan keberadaan
  situs warisan UNESCO (Bagian 3.2, hlm. 188).
\end{itemize}

\textbf{Model utama.} Artikel memakai regresi beta terinflasi nol dan
satu karena indeks fungsi metropolitan dibatasi antara 0 dan 1 serta
memuat observasi bernilai tepat 0 dan 1. Model terdiri atas regresi
logistik untuk peluang nilai 0, regresi logistik untuk peluang nilai 1,
dan regresi beta untuk proporsi di antara 0 dan 1. Penulis terutama
menaruh perhatian pada bagian terinflasi nol karena banyak kota tidak
memiliki fungsi metropolitan dan mungkin terdapat ambang populasi untuk
menopang fungsi tersebut (Bagian 3.3, hlm. 188-189).

\textbf{Pelaporan dan pengujian heterogenitas.}

\begin{itemize}
\tightlist
\item
  Tabel 3 menyajikan koefisien, galat baku kokoh, dan \emph{average
  marginal effects} (AME). Bagian terinflasi satu tidak ditampilkan
  karena hanya memuat intersep; semua model memuat intersep dan variabel
  boneka negara (Tabel 3, hlm. 189).
\item
  Kepentingan substantif ukuran dan konektivitas dibandingkan melalui
  AME per unit serta perubahan prediksi ketika masing-masing variabel
  berpindah dari persentil ke-25 ke persentil ke-75 (Bagian 4.1-4.2,
  hlm. 189-191).
\item
  Model dasar diestimasi ulang untuk setiap fungsi metropolitan guna
  menguji heterogenitas jenis fungsi. Tabel 4 hanya menampilkan AME dan
  pergeseran persentil; model lengkap dinyatakan tersedia atas
  permintaan (Bagian 4.3 dan Tabel 4, hlm. 191-192).
\item
  Interaksi antara ukuran dan kedua ukuran konektivitas ditambahkan
  untuk menilai heterogenitas antarkota. Karena interaksi dalam model
  nonlinear sulit ditafsirkan, penulis menampilkan secara grafis
  perubahan AME sepanjang ukuran kota 15.000-1.500.000 jiwa, yaitu
  observasi yang kira-kira berada dalam tiga simpangan baku dari rerata
  setelah pencilan dikeluarkan (Bagian 4.4 dan Gambar 1-3, hlm.
  193-194).
\item
  Estimasi ulang dengan OLS dinyatakan menghasilkan simpulan serupa,
  tetapi hasilnya hanya tersedia dari penulis atas permintaan.
  \emph{Fractional logit} tidak dipakai karena data mengalami
  overdispersi (catatan kaki 2-3, hlm. 189-190).
\end{itemize}

\textbf{Inferensi review atas identifikasi.} Model mengendalikan
kovariat teramati dan perbedaan negara, tetapi tidak memakai instrumen,
data panel, eksperimen, atau urutan temporal untuk mengidentifikasi
sebab-akibat. Penulis sendiri menyatakan kemungkinan bahwa ukuran dan
konektivitas justru didorong oleh fungsi metropolitan (Bagian 5, hlm.
196).

\textbf{Informasi metode yang tidak tersedia.} Perangkat lunak dan
versinya, kode analisis, prosedur penanganan data hilang, pemeriksaan
ketergantungan spasial, diagnostik pengaruh pencilan selain pembatasan
rentang grafik interaksi, spesifikasi lengkap model per fungsi, serta
uji sensitivitas terhadap batas 45 menit: Tidak disebutkan dalam file.

\subsection{Dataset}\label{dataset-17}

\begin{itemize}
\tightlist
\item
  \textbf{Fungsi metropolitan:} lokasi fungsi diperoleh dari Federal
  Institute for Research on Building, Urban Affairs and Spatial
  Development di Jerman (BBSR 2011). Data tiap indikator berasal dari
  satu tahun dalam rentang 2004-2009 dan sebagian besar dari 2008
  (Bagian 3.2, hlm. 186-187).
\item
  \textbf{Cakupan fungsi:} institusi internasional meliputi kantor PBB,
  institusi Uni Eropa, organisasi internasional, dan LSM. Domain
  perusahaan meliputi perusahaan Fortune 500 menurut omzet dan pekerja,
  jasa produsen maju, bank, dan pameran. Domain sains meliputi
  universitas besar dan organisasi penelitian internasional. Budaya
  dibagi menjadi acara dan tempat budaya, sedangkan olahraga mencakup
  stadion serta acara tingkat tinggi, termasuk Olimpiade Musim Panas
  (Bagian 3.2, hlm. 186-187).
\item
  \textbf{Konstruksi indeks:} setiap indikator dinormalisasi pada skala
  0-100, dijumlahkan dalam lima domain, lalu dibagi dengan jumlah
  indikator domain. Skor kelima domain diproses dengan prosedur yang
  sama menjadi indeks umum fungsi metropolitan. Dalam model, indeks
  dependen dinyatakan berada pada skala 0-1 (Bagian 3.2-3.3, hlm.
  187-188).
\item
  \textbf{Indikator per fungsi:} Tabel 4 menganalisis organisasi
  internasional, kantor pusat, bank, jasa produsen maju, pameran,
  universitas, konferensi, paten, tempat budaya, acara budaya, tempat
  olahraga, dan acara olahraga. Indikator kantor pusat berasal dari
  Csomos dan Derudder dan dinyatakan memberi hasil serupa dengan data
  perusahaan BBSR (Tabel 4 dan catatan b, hlm. 192).
\item
  \textbf{Delimitasi spasial:} data awal berada terutama pada LAU 2 atau
  skala munisipal, kemudian ditautkan ke MUA, FUA, dan PUSH dari proyek
  ESPON. MUA adalah kawasan terbangun yang bersambung; FUA menggabungkan
  MUA dengan hinterland komuter; PUSH merepresentasikan wilayah dalam
  jangkauan 45 menit (Bagian 3.2, hlm. 187).
\item
  \textbf{Cakupan geografis:} ketidakselarasan delimitasi antarsumber
  membuat penulis membatasi data pada bekas EU15 selain UK, ditambah
  Norwegia dan Swiss. Dataset akhir memuat seluruh 1.114 MUA di 16
  negara tersebut (Bagian 3.2, hlm. 187).
\item
  \textbf{Konektivitas regional:} populasi yang dapat dicapai dalam 45
  menit dengan mobil, tidak termasuk penduduk kota asal. Delimitasi PUSH
  berasal dari pekerjaan ESPON yang dirujuk dalam artikel (Bagian 3.2,
  hlm. 187).
\item
  \textbf{Konektivitas (inter)nasional:} potensi aksesibilitas multimoda
  jalan, rel, dan udara berasal dari penelitian ESPON oleh University of
  Tours serta Spiekermann dan Wegener, dengan pembaruan 2009 untuk data
  tahun 2006. Indikator awal pada NUTS 3 diadaptasi ke FUA dalam
  penelitian ini (Bagian 3.2 dan catatan kaki 1, hlm. 187).
\item
  \textbf{Kontrol:} PDB per kapita diukur pada FUA; status ibu kota
  dikodekan sebagai variabel boneka; pariwisata berasal dari BBSR dan
  menggabungkan bintang destinasi Michelin dengan situs UNESCO; variabel
  boneka negara menangkap antara lain perbedaan institusional dalam
  produksi fungsi metropolitan (Bagian 3.2, hlm. 188).
\item
  \textbf{Sumber ukuran populasi lokal dan PDB per kapita:} nama
  dataset, tahun referensi, serta prosedur pengolahannya secara khusus:
  Tidak disebutkan dalam file.
\item
  \textbf{Ketersediaan data:} tautan unduh dataset gabungan, lisensi,
  kamus variabel, kode konstruksi indeks, dan repositori replikasi:
  Tidak disebutkan dalam file.
\end{itemize}

\textbf{Inferensi review.} Dataset menyatukan indikator dengan tahun dan
satuan spasial yang berbeda: fungsi terutama 2008 tetapi mencakup
2004-2009, aksesibilitas memakai data 2006, dan beberapa variabel
memakai MUA, FUA, atau PUSH. Artikel menjelaskan penautan geografis,
tetapi tidak melaporkan uji sensitivitas terhadap ketidakselarasan waktu
dan unit tersebut (Bagian 3.2, hlm. 186-188).

\subsection{Population / Sample}\label{population-sample-17}

\textbf{Unit analisis.} Unit observasi adalah MUA, yaitu kawasan
terbangun yang bersambung dan dianggap paling mendekati konsep
aglomerasi perkotaan. Penulis menyebut unit-unit tersebut sebagai
``kota'' untuk memudahkan pembahasan, sekalipun banyak MUA
terfragmentasi secara kelembagaan (Bagian 3.2, hlm. 187).

\textbf{Cakupan populasi analitis.} Dataset mencakup seluruh 1.114
aglomerasi perkotaan yang berhasil dibentuk pada 16 negara dalam cakupan
studi, bukan sampel probabilitas dari kota-kota tersebut. UK dikeluarkan
dari bekas EU15, sedangkan Norwegia dan Swiss dimasukkan. Negara Eropa
lain di luar cakupan tersebut tidak dianalisis (Bagian 3.2, hlm. 187).

\textbf{Ukuran kota.} Tabel 1 melaporkan populasi lokal dalam satuan
100.000 jiwa dengan rerata 1,530, simpangan baku 4,497, minimum 0,051,
dan maksimum 95,910. Dalam perbandingan persentil, persentil ke-25
berjumlah 34.000 penduduk dan persentil ke-75 berjumlah 112.000 penduduk
(Tabel 1, hlm. 188; Bagian 4.2, hlm. 190-191).

\textbf{Subcakupan grafik interaksi.} Gambar 1-3 memvisualisasikan kota
berpenduduk 15.000-1.500.000 jiwa, yaitu observasi yang menurut penulis
kira-kira berada dalam tiga simpangan baku dari rerata setelah pencilan
dikeluarkan. File tidak menyebutkan jumlah tepat observasi yang tersisa
dalam rentang grafik tersebut (Bagian 4.4, hlm. 193-194).

\textbf{Klarifikasi populasi manusia.} Penelitian tidak merekrut
responden individu, rumah tangga, perusahaan, atau pengguna fasilitas.
Ukuran penduduk berfungsi sebagai atribut kota dan sebagai potensi
populasi yang dapat dicapai. Teknik penarikan sampel individu,
karakteristik responden, dan tingkat respons: Tidak berlaku.

\textbf{Inferensi review.} Generalisasi empiris langsung terbatas pada
sistem perkotaan matang dalam 16 negara yang dapat ditautkan secara
memadai ke delimitasi ESPON. Cakupan 1.114 MUA luas di dalam wilayah
tersebut, tetapi bukan sampel yang secara desain mewakili seluruh Eropa
atau sistem kota di luar Eropa Barat (Bagian 1, hlm. 182; Bagian 3.2,
hlm. 187).

\subsection{Dependent Variables}\label{dependent-variables-17}

\textbf{Variabel dependen utama.} Indeks agregat kehadiran fungsi
metropolitan berada antara 0 dan 1. Nilai 0 berarti tidak ada fungsi
metropolitan dalam indeks, sedangkan nilai 1 diberikan kepada MUA dengan
skor tertinggi. Rerata indeks adalah 0,019, simpangan baku 0,066,
minimum 0, dan maksimum 1 (Bagian 3.3 dan Tabel 1, hlm. 188-189).

\textbf{Variabel dependen terpilah.} Model heterogenitas mengganti
indeks agregat dengan skor untuk 12 fungsi: organisasi internasional,
kantor pusat, bank, jasa produsen maju, pameran, universitas,
konferensi, paten, tempat budaya, acara budaya, tempat olahraga, dan
acara olahraga (Tabel 4, hlm. 192).

\textbf{Massa pada nilai nol.} Tabel 4 melaporkan proporsi MUA tanpa
tiap fungsi: 96,0\% untuk organisasi internasional; 93,5\% kantor pusat;
76,5\% bank; 73,4\% jasa produsen maju; 90,9\% pameran; 89,3\%
universitas; 89,2\% konferensi; 45,1\% paten; 48,3\% tempat budaya;
91,0\% acara budaya; 64,9\% tempat olahraga; dan 91,2\% acara olahraga
(Tabel 4, hlm. 192). Distribusi dengan banyak nilai nol ini menjadi
alasan substantif penggunaan model terinflasi nol.

\textbf{Variabel penjelas utama.} Populasi lokal dinyatakan per 100.000
jiwa, konektivitas regional per satu juta penduduk yang dapat dicapai,
dan konektivitas (inter)nasional per sepuluh juta penduduk potensial
yang dapat dicapai. Kontrolnya ialah PDB per kapita FUA dalam ribuan,
variabel boneka ibu kota, indeks pariwisata, dan variabel boneka negara
(Tabel 1 dan Tabel 3, hlm. 188-189).

\textbf{Batas konstruk.} Variabel hasil mengukur lokasi atau kehadiran
fungsi metropolitan, bukan penggunaan fungsi oleh penduduk kota lain,
kualitas fungsi, produktivitas kota, pertumbuhan ekonomi, pertumbuhan
penduduk, atau kesejahteraan. Variabel hasil untuk \emph{borrowed
performance}: Tidak disebutkan dalam file sebagai bagian analisis
empiris A02.

\textbf{Inferensi review.} Penyebutan \emph{borrowed size} dalam model
sebenarnya merujuk pada \emph{borrowed function}: kota yang lebih
terkoneksi memiliki skor fungsi lebih tinggi setelah ukuran dan kontrol
diperhitungkan. Penelitian tidak mengamati apakah penduduk atau
perusahaan benar-benar mengakses fungsi di kota lain (Bagian 2, hlm.
185; Bagian 3.1, hlm. 186).

\subsection{Results}\label{results-17}

\textbf{Deskripsi variabel dan korelasi.}

\begin{itemize}
\tightlist
\item
  \textbf{Laporan sumber:} rerata konektivitas regional adalah 1,986
  juta penduduk dengan simpangan baku 2,322, rentang 0-12,436 juta.
  Rerata konektivitas (inter)nasional adalah 5,630 unit sepuluh juta
  penduduk dengan simpangan baku 1,887 dan rentang 1,392-11,368. Rerata
  PDB per kapita FUA adalah 27,356 ribu, sedangkan rerata variabel
  boneka ibu kota dan indeks pariwisata masing-masing 0,015 dan 0,062
  (Tabel 1, hlm. 188).
\item
  \textbf{Laporan sumber:} korelasi terbesar yang ditampilkan ialah 0,64
  antara konektivitas regional dan (inter)nasional. Konektivitas
  (inter)nasional berkorelasi 0,44 dengan PDB per kapita dan 0,18 dengan
  populasi lokal. Populasi lokal berkorelasi 0,55 dengan status ibu kota
  dan 0,50 dengan pariwisata. Korelasi populasi lokal dengan
  konektivitas regional hanya 0,03 (Tabel 2, hlm. 188).
\end{itemize}

\textbf{Model agregat pada bagian proporsi.}

\begin{itemize}
\tightlist
\item
  \textbf{Laporan sumber:} populasi lokal mempunyai koefisien 0,081
  dengan galat baku 0,014 dan AME 0,0077 dengan galat baku 0,0009;
  keduanya ditandai signifikan pada \texttt{p\ \textless{}\ 0,01} (Tabel
  3, hlm. 189).
\item
  \textbf{Laporan sumber:} konektivitas (inter)nasional mempunyai
  koefisien 0,092 dengan galat baku 0,035, ditandai
  \texttt{p\ \textless{}\ 0,01}, dan AME 0,0013 dengan galat baku
  0,0006, ditandai \texttt{p\ \textless{}\ 0,05} (Tabel 3, hlm. 189).
\item
  \textbf{Laporan sumber:} baris konektivitas regional pada hasil
  ekstraksi Tabel 3 menampilkan koefisien 0,024 dengan galat baku 0,025
  dan AME 0,0005 dengan galat baku 0,0004, tanpa tanda signifikansi.
  Akan tetapi, narasi hasil kemudian menyebut ``tanda negatif'' untuk
  konektivitas regional. Ketidakselarasan ini tidak dapat diselesaikan
  hanya dari TXT (Tabel 3 dan Bagian 4.1, hlm. 189-190).
\item
  \textbf{Laporan sumber:} PDB per kapita mempunyai koefisien 0,009 dan
  AME 0,0002, keduanya tanpa tanda signifikansi. Status ibu kota
  mempunyai koefisien 0,436 tanpa tanda signifikansi dan AME 0,0399 yang
  ditandai \texttt{p\ \textless{}\ 0,01}. Pariwisata mempunyai koefisien
  2,004 dan AME 0,0403, keduanya ditandai \texttt{p\ \textless{}\ 0,01}
  (Tabel 3, hlm. 189).
\end{itemize}

\textbf{Bagian terinflasi nol dan statistik model.}

\begin{itemize}
\tightlist
\item
  \textbf{Laporan sumber:} koefisien yang ditampilkan ialah 2,873 untuk
  populasi lokal, 0,068 untuk konektivitas regional, 0,062 untuk
  konektivitas (inter)nasional, 3,616 untuk PDB per kapita, 14,748 untuk
  status ibu kota, dan 0,018 untuk pariwisata. Populasi lokal, PDB per
  kapita, dan status ibu kota ditandai signifikan pada
  \texttt{p\ \textless{}\ 0,01}; tiga variabel lainnya tidak ditandai
  signifikan (Tabel 3, hlm. 189).
\item
  \textbf{Laporan sumber:} model memakai 1.114 observasi, menghasilkan
  Wald chi-square 450,33 yang ditandai \texttt{p\ \textless{}\ 0,01},
  serta \texttt{ln\ phi} 3,46 dengan galat baku 0,104 dan tanda
  \texttt{p\ \textless{}\ 0,01}. Variabel boneka negara disertakan
  (Tabel 3, hlm. 189).
\item
  \textbf{Inferensi review atas batas interpretasi:} file tidak
  menjelaskan secara eksplisit apakah parameter logit bagian nol
  dimodelkan sebagai peluang berada pada nilai nol atau peluang tidak
  berada pada nilai nol. Review ini karena itu tidak memberi makna arah
  pada koefisien bagian tersebut (Bagian 3.3 dan Tabel 3, hlm. 188-189).
\end{itemize}

\textbf{Perbandingan kepentingan ukuran dan konektivitas.}

\begin{itemize}
\tightlist
\item
  \textbf{Laporan sumber:} tambahan 100.000 penduduk lokal dikaitkan
  dengan kenaikan rata-rata indeks fungsi metropolitan sebesar 0,0077,
  sedangkan tambahan sepuluh juta penduduk dalam potensi konektivitas
  (inter)nasional dikaitkan dengan kenaikan 0,0013 (Bagian 4.2 dan Tabel
  3, hlm. 189-190).
\item
  \textbf{Laporan sumber:} perpindahan populasi lokal dari persentil
  ke-25, yaitu 34.000 jiwa, ke persentil ke-75, yaitu 112.000 jiwa,
  dikaitkan dengan kenaikan indeks rata-rata 0,0080. Perpindahan
  konektivitas (inter)nasional dari 42,3 juta ke 69,9 juta penduduk
  potensial dikaitkan dengan kenaikan 0,0032 (Bagian 4.2, hlm. 190-191).
\item
  \textbf{Interpretasi penulis:} perbandingan tersebut menunjukkan bahwa
  ukuran lokal secara agregat lebih penting daripada konektivitas
  jaringan. Bagian simpulan menyatakannya sebagai efek yang kira-kira
  2,5 kali lebih besar (Bagian 4.2, hlm. 190-191; Bagian 5, hlm. 195).
\end{itemize}

\textbf{Heterogenitas antarjenis fungsi.}

\begin{itemize}
\tightlist
\item
  \textbf{Laporan sumber:} ukuran lokal merupakan determinan signifikan
  untuk seluruh fungsi metropolitan yang diuji. Konektivitas
  (inter)nasional signifikan untuk sebagian besar, tetapi tidak semua,
  fungsi (Bagian 4.3, hlm. 191; Tabel 4, hlm. 192).
\item
  \textbf{Laporan sumber:} dari fungsi budaya dan olahraga, hanya tempat
  budaya yang dikaitkan secara signifikan dengan konektivitas
  (inter)nasional. Acara budaya, tempat olahraga, dan acara olahraga
  tidak dilaporkan mempunyai hubungan signifikan tersebut (Bagian 4.3,
  hlm. 191-193; Tabel 4, hlm. 192).
\item
  \textbf{Laporan sumber:} konektivitas (inter)nasional lebih penting
  bagi fungsi perusahaan dan organisasi internasional. Berdasarkan
  pergeseran persentil, kehadiran pameran dan organisasi internasional
  dinilai lebih bergantung pada konektivitas (inter)nasional daripada
  ukuran lokal (Bagian 4.3, hlm. 191).
\item
  \textbf{Laporan sumber:} dalam domain sains, konektivitas jaringan
  dinilai penting untuk paten dan konferensi internasional, tetapi tidak
  menjelaskan kehadiran universitas. Penulis mengaitkan pengecualian
  universitas dengan waktu pendiriannya yang sudah lama dan sifatnya
  yang tidak mudah berpindah (Bagian 4.3, hlm. 191).
\item
  \textbf{Laporan sumber:} konektivitas regional secara umum tidak
  berpengaruh. Narasi menyebut AME negatif dan signifikan untuk jasa
  produsen maju, konferensi, paten, tempat budaya, dan acara olahraga;
  mayoritas tanda fungsi lain juga disebut negatif tetapi tidak
  signifikan (Bagian 4.3, hlm. 191-193). Sejumlah tanda minus pada Tabel
  4 tidak terbaca dengan andal dalam ekstraksi, sehingga arah ini
  dilaporkan dari narasi penulis, bukan direkonstruksi dari susunan
  tabel.
\end{itemize}

\textbf{Heterogenitas menurut ukuran kota.}

\begin{itemize}
\tightlist
\item
  \textbf{Laporan sumber:} tambahan ukuran lokal terutama meningkatkan
  fungsi metropolitan kota kecil, sedangkan efek marginal ukuran
  dinyatakan kecil pada kota yang lebih besar (Bagian 4.4 dan Gambar 1,
  hlm. 193).
\item
  \textbf{Laporan sumber:} AME konektivitas regional menjadi positif
  bagi kota berpenduduk lebih dari 500.000 jiwa dan secara umum negatif
  bagi kota berpenduduk kurang dari 250.000 jiwa (Bagian 4.4 dan Gambar
  2, hlm. 193-194).
\item
  \textbf{Laporan sumber:} AME konektivitas (inter)nasional secara
  signifikan lebih besar pada kota yang lebih besar. Penulis menilai
  hasil ini menunjukkan bahwa fungsi metropolitan ditopang oleh
  kombinasi skala dan konektivitas (Bagian 4.4 dan Gambar 3, hlm.
  193-195).
\item
  \textbf{Inferensi review atas masalah internal:} narasi menyatakan
  bahwa efek tambahan ukuran kecil pada kota berpenduduk ``lebih dari
  150.000'' dan memberi keterangan ``85\% dari kota''. Angka persentil
  pada bagian sebelumnya menunjukkan persentil ke-75 hanya 112.000 jiwa,
  sehingga keterangan 85\% tidak selaras jika arah ``lebih dari'' dibaca
  secara harfiah. Review mempertahankan pernyataan sumber tanpa
  mengganti tanda atau persentasenya (Bagian 4.2, hlm. 190-191; Bagian
  4.4, hlm. 193).
\end{itemize}

\subsection{Findings}\label{findings-17}

\begin{itemize}
\tightlist
\item
  \textbf{Interpretasi penulis:} asosiasi positif konektivitas
  (inter)nasional setelah ukuran lokal dikendalikan menunjukkan bahwa
  kota dapat meminjam ukuran melalui keterhubungan jalan, rel, dan udara
  pada skala nasional serta internasional (Bagian 4.1, hlm. 189-190).
\item
  \textbf{Interpretasi penulis:} tidak signifikannya konektivitas
  regional pada model agregat mungkin terjadi karena manfaat
  \emph{borrowed size} bagi sebagian kota dikompensasi oleh kerugian
  persaingan atau \emph{agglomeration shadow} bagi kota lain dalam PUSH
  yang sama. Penulis menyajikan mekanisme ini sebagai penjelasan yang
  mungkin, bukan hasil pengukuran persaingan langsung (Bagian 4.1, hlm.
  190).
\item
  \textbf{Interpretasi penulis:} hasil fungsi terpilah memperkuat dugaan
  bahwa bayang-bayang aglomerasi lebih dominan pada skala regional,
  khususnya untuk jasa produsen maju, konferensi, paten, tempat budaya,
  dan acara olahraga (Bagian 4.3, hlm. 191-193).
\item
  \textbf{Interpretasi penulis:} skala jaringan menentukan jenis
  eksternalitas yang tampak. Konektivitas (inter)nasional cenderung
  memberi manfaat, sedangkan konektivitas regional hanya menguntungkan
  sebagian kota dan rata-rata dapat membawa persaingan. Karena itu,
  konektivitas (inter)nasional dinilai sebagai pengganti ukuran lokal
  yang lebih penting daripada konektivitas regional (Bagian 4.3, hlm.
  193).
\item
  \textbf{Interpretasi penulis:} fungsi yang memiliki jangkauan
  internasional lebih bergantung pada konektivitas, sedangkan budaya dan
  olahraga sering bertumpu pada basis dukungan lokal. Universitas
  menjadi pengecualian dalam sains karena lokasinya dinilai historis dan
  tidak mudah berpindah (Bagian 4.3, hlm. 191).
\item
  \textbf{Interpretasi penulis:} kota besar lebih mampu memanfaatkan
  konektivitas karena beberapa fungsi membutuhkan sekaligus massa
  minimum dan jaringan (inter)nasional. Hasil tersebut dikaitkan dengan
  pembedaan \emph{town-ness}, yang menghubungkan kota dan hinterland
  secara hierarkis, dari \emph{city-ness}, yang mencakup hubungan
  horizontal antarkota di luar hinterland (Bagian 4.4, hlm. 193-195).
\item
  \textbf{Inferensi review:} temuan tidak mendukung proposisi bahwa
  jaringan secara umum menggantikan ukuran. Dukungan yang lebih tepat
  ialah substitusi parsial dan bersyarat: berlaku terutama pada
  konektivitas (inter)nasional, jenis fungsi tertentu, serta kota yang
  sudah memiliki skala lebih besar (Bagian 4.2-4.4, hlm. 190-195; Bagian
  5, hlm. 195-196).
\item
  \textbf{Inferensi review:} istilah \emph{agglomeration shadow} pada
  hasil regional bersandar pada hubungan negatif dan penjelasan
  persaingan. Karena model tidak mengamati redistribusi fungsi dari satu
  kota ke kota lain, hasil belum membedakan persaingan dari faktor
  regional tak teramati atau mekanisme lain yang dapat menghasilkan
  asosiasi serupa (Bagian 3.1, hlm. 186; Bagian 4.1 dan 4.3, hlm.
  190-193).
\item
  \textbf{Inferensi review:} fungsi metropolitan yang jarang hadir
  menghasilkan banyak nilai nol. Model menangani bentuk distribusinya,
  tetapi kelangkaan tersebut juga berarti hasil untuk beberapa fungsi
  menggambarkan perbedaan antara segelintir pusat dan mayoritas kota
  tanpa fungsi, bukan perubahan intensitas yang merata di seluruh sistem
  (Bagian 3.3, hlm. 188-189; Tabel 4, hlm. 192).
\end{itemize}

\subsection{Conclusions}\label{conclusions-17}

\textbf{Kesimpulan penulis.} Artikel merangkum enam hasil utama (Bagian
5, hlm. 195-196):

\begin{enumerate}
\def\labelenumi{\arabic{enumi}.}
\tightlist
\item
  Ukuran lokal dan konektivitas dalam jaringan (inter)nasional sama-sama
  berkontribusi positif terhadap kehadiran fungsi metropolitan.
\item
  Kota meminjam ukuran melalui jaringan (inter)nasional, sedangkan
  keterhubungan regional pada umumnya tidak menghasilkan tingkat fungsi
  metropolitan yang lebih tinggi.
\item
  Efek ukuran lokal secara umum jauh lebih besar, kira-kira 2,5 kali
  efek konektivitas jaringan.
\item
  Kepentingan ukuran dan konektivitas berbeda menurut fungsi. Ukuran
  selalu penting; konektivitas tidak berperan bagi sebagian besar fungsi
  budaya dan olahraga, tetapi penting dan kadang lebih besar daripada
  ukuran bagi perusahaan, institusi internasional, dan sains.
\item
  Konektivitas regional yang lebih kuat menghasilkan efek persaingan
  bagi beberapa fungsi, sehingga \emph{agglomeration shadow} lebih
  dominan daripada \emph{borrowed size} pada skala regional.
\item
  Kota kecil terutama memperoleh fungsi dari pertambahan ukuran,
  sedangkan kota besar lebih memperoleh fungsi dari peningkatan
  konektivitas regional dan (inter)nasional.
\end{enumerate}

\textbf{Kesimpulan teoretis penulis.} Eksternalitas jaringan kota
dipandang sebagai mata rantai yang melengkapi teori aglomerasi dalam
menjelaskan dinamika Eropa Barat. Karena manfaat yang biasanya disebut
ekonomi aglomerasi dapat dibagi melalui jaringan, penulis meminta agar
fondasi geografis teori aglomerasi ditinjau kembali, bahkan
mempertanyakan ketepatan istilah ``ekonomi aglomerasi'' ketika
manfaatnya tidak terbatas pada aglomerasi itu sendiri (Bagian 5, hlm.
195).

\textbf{Kualifikasi penulis.} Hasil harus ditafsirkan sebagai asosiasi
bersyarat, bukan hubungan kausal, karena mungkin terjadi kausalitas
terbalik: fungsi metropolitan dapat mendorong ukuran dan konektivitas,
bukan hanya dipengaruhi keduanya (Bagian 5, hlm. 196).

\textbf{Inferensi review.} Simpulan terkuat yang didukung desain ialah
bahwa ukuran lokal dan konektivitas (inter)nasional mengandung informasi
statistik yang berbeda tentang distribusi fungsi metropolitan. Label
\emph{borrowed size}, \emph{agglomeration shadow}, dan penggantian
ukuran adalah interpretasi konseptual atas pola tersebut; kekuatan bukti
untuk mekanisme dan kausalitas lebih rendah daripada kekuatan bukti
untuk asosiasi (Bagian 4, hlm. 189-195; Bagian 5, hlm. 195-196).

\subsection{Contributions}\label{contributions-17}

\begin{itemize}
\tightlist
\item
  \textbf{Kontribusi yang diklaim penulis:} menawarkan penjelasan
  alternatif bagi ketegangan antara dinamika kota-kota Eropa Barat dan
  teori aglomerasi dengan menempatkan eksternalitas jaringan kota
  sebagai pelengkap ekonomi aglomerasi lokal (Bagian 1, hlm. 182-183;
  Bagian 5, hlm. 195).
\item
  \textbf{Kontribusi yang diklaim penulis:} memperluas analogi ekonomi
  jaringan dari jaringan antarperusahaan ke jaringan antarkota dan
  menghubungkannya secara eksplisit dengan ekonomi urbanisasi (Bagian 5,
  hlm. 195).
\item
  \textbf{Kontribusi konseptual penulis:} menyatukan \emph{borrowed
  size} dan \emph{agglomeration shadow} sebagai dua kemungkinan tanda
  eksternalitas jaringan. Pengaruh positif konektivitas terhadap fungsi
  metropolitan diposisikan sebagai \emph{borrowed size}, sedangkan
  pengaruh negatif diposisikan sebagai \emph{agglomeration shadow}
  (Bagian 2, hlm. 185-186).
\item
  \textbf{Inferensi review atas kontribusi empiris:} menguji ukuran
  lokal bersama dua skala konektivitas pada 1.114 MUA di 16 negara,
  dengan fungsi metropolitan pada domain institusi internasional,
  perusahaan, sains, budaya, dan olahraga (Bagian 3, hlm. 186-188).
\item
  \textbf{Inferensi review atas kontribusi analitis:} tidak berhenti
  pada indeks agregat, tetapi membandingkan 12 fungsi dan menguji
  interaksi ukuran-konektivitas untuk menunjukkan bahwa hasil berubah
  menurut jenis fungsi serta ukuran kota (Bagian 4.3-4.4, hlm. 191-195).
\item
  \textbf{Inferensi review atas kontribusi pengukuran:} membandingkan
  AME dalam satuan asli dengan perubahan dari persentil ke-25 ke
  persentil ke-75 agar signifikansi statistik tidak disamakan dengan
  kepentingan substantif (Bagian 4.1-4.3, hlm. 189-193).
\item
  \textbf{Inferensi review:} nilai tambah utama studi bukan teknik
  regresi yang baru, melainkan pengoperasian kerangka
  ukuran-jaringan-fungsi secara multiskala dan heterogen. Kerangka
  tersebut mempersempit klaim ``jaringan menggantikan ukuran'' menjadi
  klaim yang bergantung pada skala jaringan, fungsi, dan ukuran kota
  (Bagian 3-4, hlm. 186-195).
\end{itemize}

\subsection{Research Gap}\label{research-gap-17}

\textbf{Kesenjangan yang memotivasi studi.}

\begin{itemize}
\tightlist
\item
  Teori aglomerasi dan model \emph{new economic geography} dinilai belum
  memadai untuk menjelaskan stabilitas urbanisasi, terbatasnya
  pertumbuhan kota besar, serta kinerja kota kecil dan menengah dalam
  sistem perkotaan Eropa Barat yang matang (Bagian 1, hlm. 181-183).
\item
  Jangkauan spasial manfaat dan biaya aglomerasi belum banyak
  dieksplorasi. Literatur menunjukkan efek melemah terhadap jarak dan
  mungkin meluas, tetapi belum jelas apakah ekonomi urbanisasi
  benar-benar dapat dibagi melalui jaringan kota (Bagian 2, hlm. 184).
\item
  Literatur jaringan kota lebih banyak berfokus pada konektivitas
  internasional, sedangkan perbandingan jaringan regional dan
  (inter)nasional dengan faktor lokal masih kurang (Bagian 2, hlm.
  184-185; pertanyaan penelitian, hlm. 183).
\item
  Konsep \emph{borrowed size} mendapat perhatian empiris terbatas
  meskipun dinilai berpotensi menjelaskan dinamika Eropa Barat. Konsep
  tersebut juga perlu diperluas dari kedekatan regional ke jaringan
  nasional dan internasional (Bagian 2, hlm. 185).
\item
  Heterogenitas menurut jenis fungsi dan ukuran kota perlu
  diperhitungkan karena ukuran serta konektivitas mungkin tidak
  mempunyai kepentingan yang sama untuk semua fungsi dan kota (Bagian
  3.1, hlm. 186).
\end{itemize}

\textbf{Kesenjangan yang dinyatakan masih tersisa.} Artikel belum
menilai perubahan kepentingan konektivitas dari waktu ke waktu, fungsi
atau indikator kinerja lain, tipologi kota yang lebih luas, kausalitas,
saluran \emph{matching} dan \emph{learning}, ragam eksternalitas
jaringan, ataupun cara mengatasi dominasi bayang-bayang aglomerasi
regional (Bagian 5, hlm. 196).

\textbf{Inferensi review.} A02 mengisi celah tentang \emph{borrowed
function} pada data lintas kota, tetapi belum mengisi celah pengukuran
mekanisme. Asosiasi antara potensi populasi dan fungsi belum menunjukkan
kota mana yang meminjam dari kota mana, jenis aliran yang
menghubungkannya, siapa pengguna lintas batasnya, atau fungsi mana yang
berpindah akibat persaingan (Bagian 2, hlm. 185-186; Bagian 3.1, hlm.
186).

\subsection{Limitations}\label{limitations-17}

\textbf{Keterbatasan yang dinyatakan atau diakui penulis.}

\begin{itemize}
\tightlist
\item
  Ketidakselarasan delimitasi LAU 2, MUA, FUA, dan PUSH tidak dapat
  diselesaikan secara memuaskan untuk semua negara Eropa. Akibatnya,
  studi dibatasi pada bekas EU15 selain UK, ditambah Norwegia dan Swiss
  (Bagian 3.2, hlm. 187).
\item
  Ukuran potensi populasi dipilih karena penelitian berfokus pada
  ekonomi aglomerasi melalui \emph{sharing}. Penulis menyatakan bahwa
  data jaringan lain, seperti ko-publikasi, ko-inventor, kerja sama
  riset, kolaborasi R\&D, perdagangan antarkota, dan investasi asing,
  lebih tepat untuk menilai \emph{learning} atau \emph{matching} (Bagian
  3.1, hlm. 186).
\item
  Analisis hanya menunjukkan asosiasi bersyarat. Kausalitas terbalik
  mungkin terjadi jika fungsi metropolitan mendorong ukuran serta
  konektivitas. Penulis menyatakan bahwa ketiadaan instrumen kredibel
  untuk konektivitas dan kurangnya data historis membuat masalah ini
  sulit diatasi (Bagian 5, hlm. 196).
\item
  Analisis belum memiliki perspektif dinamis untuk menguji apakah
  kepentingan konektivitas meningkat dari waktu ke waktu dibandingkan
  faktor lokal (Bagian 5, hlm. 196).
\item
  Fokus hasil terbatas pada fungsi metropolitan yang dipilih. Indikator
  kinerja ekonomi lain dan klasifikasi kota selain ukuran belum diuji
  (Bagian 5, hlm. 196).
\item
  Artikel terutama menilai saluran \emph{sharing} dan belum memvalidasi
  eksternalitas jaringan melalui \emph{matching} atau \emph{learning}
  (Bagian 3.1, hlm. 186; Bagian 5, hlm. 196).
\end{itemize}

\textbf{Batas desain dan pelaporan yang diinferensikan dari TXT.}

\begin{itemize}
\tightlist
\item
  Fungsi metropolitan berasal dari tahun-tahun dalam rentang 2004-2009,
  sebagian besar 2008, sedangkan konektivitas (inter)nasional memakai
  data 2006. Uji kepekaan terhadap ketidakselarasan tahun: Tidak
  disebutkan dalam file.
\item
  Potensi populasi mengukur akses yang dimungkinkan infrastruktur, bukan
  arus interaksi aktual. Intensitas perjalanan, penggunaan fasilitas
  lintas kota, hubungan organisasi, dan arah jaringan antarpasangan kota
  tidak diamati dalam variabel konektivitas yang dilaporkan.
\item
  Konektivitas regional diukur dengan ambang tetap 45 menit berkendara.
  Alasan substantif pemilihan tepat 45 menit dan uji sensitivitas
  terhadap ambang lain: Tidak disebutkan dalam file.
\item
  Variabel dibangun pada MUA, FUA, dan PUSH yang berbeda. Artikel
  menjelaskan fungsi tiap skala, tetapi konsekuensi agregasi spasial
  atau perubahan hasil jika delimitasi diganti: Tidak disebutkan dalam
  file.
\item
  Variabel hasil adalah proksi kehadiran fungsi. Frekuensi penggunaan,
  kapasitas, mutu, jangkauan pelayanan, dan manfaat ekonomi aktual tiap
  fungsi tidak diukur.
\item
  Model mengendalikan negara, tetapi pengujian ketergantungan spasial,
  galat spasial, atau pengelompokan galat baku menurut negara atau
  wilayah: Tidak disebutkan dalam file.
\item
  Model lengkap per fungsi dalam Tabel 4, hasil OLS pembanding, kode
  analisis, dan dataset replikasi tidak tersedia dalam TXT; artikel
  menyatakan sebagian hasil tersedia atas permintaan.
\item
  Sumber dan tahun populasi lokal serta PDB per kapita tidak dirinci.
  Penanganan observasi hilang, kriteria pencilan di luar grafik
  interaksi, dan validasi indeks komposit juga tidak dijelaskan dalam
  file.
\item
  Banyak fungsi sangat jarang, dengan 89,2\%-96,0\% MUA tidak memiliki
  sejumlah fungsi. Model terinflasi menangani distribusi tersebut,
  tetapi file tidak melaporkan uji kestabilan koefisien untuk
  fungsi-fungsi paling langka.
\end{itemize}

\textbf{Inferensi review.} Karena desain tidak mengamati urutan waktu
dan aliran antarkota, hasil tidak dapat memastikan apakah jaringan
menciptakan fungsi, fungsi menciptakan jaringan, atau keduanya terbentuk
oleh faktor lain. Hal ini membatasi penggunaan tanda koefisien sebagai
bukti tunggal mekanisme \emph{borrowed size} atau \emph{agglomeration
shadow} (Bagian 3.1-3.3, hlm. 186-189; Bagian 5, hlm. 196).

\subsection{Challenges}\label{challenges-17}

\begin{itemize}
\tightlist
\item
  \textbf{Laporan sumber:} tantangan data geografis adalah menautkan
  lokasi fungsi pada LAU 2 dengan delimitasi MUA, FUA, dan PUSH yang
  berasal dari proyek berbeda. Ketidakselarasan ini langsung
  mempersempit cakupan negara (Bagian 3.2, hlm. 187).
\item
  \textbf{Laporan sumber:} tantangan pengukuran adalah memilih bentuk
  jaringan yang sesuai dengan mekanisme. Potensi populasi dinilai
  memadai untuk \emph{sharing}, tetapi pengujian \emph{matching} dan
  \emph{learning} membutuhkan jaringan pengetahuan, penelitian,
  perdagangan, atau investasi yang berbeda (Bagian 3.1, hlm. 186; Bagian
  5, hlm. 196).
\item
  \textbf{Laporan sumber:} tantangan identifikasi adalah kausalitas
  terbalik dan ketiadaan instrumen konektivitas yang kredibel, terutama
  karena data historis tidak tersedia (Bagian 5, hlm. 196).
\item
  \textbf{Interpretasi penulis:} manfaat dan kerugian jaringan dapat
  muncul bersamaan pada skala regional. Satu kota mungkin meminjam
  ukuran sementara kota lain kehilangan fungsi akibat persaingan,
  sehingga efek agregat dapat tidak signifikan (Bagian 4.1 dan 4.3, hlm.
  190-193).
\item
  \textbf{Interpretasi penulis:} kota kecil menghadapi dilema kebijakan.
  Konektivitas (inter)nasional menawarkan peluang mengganti sebagian
  kekurangan massa, tetapi kota kecil justru kurang mampu
  mengeksploitasi jaringan dan lebih rentan terhadap bayang-bayang
  aglomerasi regional (Bagian 5, hlm. 195-196).
\item
  \textbf{Laporan sumber:} tantangan kebijakan yang belum terjawab ialah
  mengatasi dominasi \emph{agglomeration shadow} pada skala regional.
  Integrasi fungsional, institusional, dan kultural antarkota diusulkan
  sebagai kemungkinan, tetapi penulis menyatakan usulan ini masih perlu
  landasan lebih kuat (Bagian 5, hlm. 196).
\item
  \textbf{Inferensi review:} tantangan replikasi berasal dari tidak
  tersedianya model lengkap per fungsi, hasil OLS, sumber rinci beberapa
  variabel, kode, dan data gabungan. TXT cukup untuk memahami rancangan,
  tetapi tidak cukup untuk menjalankan ulang estimasi (Bagian 3.2-3.3,
  hlm. 186-190; Tabel 4, hlm. 192).
\item
  \textbf{Inferensi review:} tantangan interpretasi terletak pada
  penggunaan tanda asosiasi sebagai label mekanisme. Hubungan negatif
  regional dapat konsisten dengan persaingan, tetapi model yang
  dilaporkan tidak menunjukkan pemindahan fungsi atau pasangan kota yang
  memperoleh dan kehilangan fungsi (Bagian 3.1, hlm. 186; Bagian 4.1 dan
  4.3, hlm. 190-193).
\end{itemize}

\subsection{Practical Implications}\label{practical-implications-17}

\begin{itemize}
\tightlist
\item
  \textbf{Implikasi penulis:} kebijakan pembangunan perkotaan yang
  mengabaikan ekonomi jaringan kota berisiko memakai skala spasial
  perkotaan yang terlalu sempit. Fungsi kota tidak hanya ditentukan
  massa lokal, tetapi juga posisi dalam jaringan (Bagian 5, hlm.
  195-196).
\item
  \textbf{Implikasi penulis:} keterhubungan (inter)nasional dapat
  membantu kota kecil dan menengah mengompensasi sebagian kekurangan
  massa. Namun, peluang ini tidak boleh dibaca sebagai pengganti penuh
  karena ukuran lokal tetap penting, jaringan tidak mendukung semua
  fungsi, dan kota kecil kurang mampu mengeksploitasi konektivitas
  (Bagian 5, hlm. 195-196).
\item
  \textbf{Implikasi penulis:} konektivitas regional tidak otomatis
  menguntungkan setiap kota. Kedekatan dengan kota lain dapat
  menciptakan persaingan dan bayang-bayang aglomerasi, terutama bagi
  kota yang lebih kecil (Bagian 4.3, hlm. 191-193; Bagian 5, hlm. 196).
\item
  \textbf{Inferensi review berdasarkan hasil:} strategi perlu dibedakan
  menurut fungsi. Konektivitas lebih relevan bagi organisasi
  internasional, perusahaan, konferensi, dan paten, sedangkan banyak
  fungsi budaya dan olahraga lebih bergantung pada basis lokal (Bagian
  4.3, hlm. 191-193).
\item
  \textbf{Implikasi penulis:} integrasi fungsional, institusional, dan
  kultural antarkota dalam wilayah yang sama mungkin mengurangi dominasi
  bayang-bayang aglomerasi, tetapi efektivitas strategi tersebut belum
  diuji dalam artikel (Bagian 5, hlm. 196).
\item
  \textbf{Inferensi review:} kebijakan konektivitas sebaiknya dievaluasi
  bersama kesiapan massa lokal, jenis fungsi yang dituju, posisi kota
  dalam hierarki regional, dan risiko konsentrasi fungsi pada pusat yang
  sudah besar. Hasil A02 tidak mendukung rekomendasi seragam berupa
  ``tingkatkan konektivitas'' untuk semua kota (Bagian 4.3-4.4, hlm.
  191-195; Bagian 5, hlm. 195-196).
\item
  \textbf{Inferensi review:} karena estimasi bersifat asosiasional,
  artikel belum memberi besaran dampak kausal suatu investasi jalan,
  rel, bandara, atau program integrasi tertentu. Angka AME tidak dapat
  langsung dipakai sebagai prakiraan manfaat proyek tanpa desain
  evaluasi tambahan (Bagian 4.1-4.2, hlm. 189-191; Bagian 5, hlm. 196).
\end{itemize}

\subsection{Applications}\label{applications-17}

\begin{itemize}
\tightlist
\item
  \textbf{Aplikasi empiris sumber:} kerangka diterapkan untuk
  membandingkan kehadiran fungsi metropolitan pada 1.114 MUA menurut
  ukuran lokal, konektivitas regional, dan konektivitas (inter)nasional
  (Bagian 3-4, hlm. 186-195).
\item
  \textbf{Aplikasi konseptual sumber:} arah hubungan konektivitas
  digunakan untuk mengeksplorasi \emph{borrowed size} ketika positif dan
  \emph{agglomeration shadow} ketika negatif, dengan tetap
  membandingkannya terhadap peran ukuran lokal (Bagian 2, hlm. 185-186;
  Bagian 3.1, hlm. 186).
\item
  \textbf{Aplikasi analitis sumber:} indeks agregat dapat diuraikan per
  fungsi untuk mengetahui fungsi yang lebih ditentukan pasar lokal dan
  fungsi yang lebih bergantung pada jaringan (Bagian 4.3 dan Tabel 4,
  hlm. 191-193).
\item
  \textbf{Aplikasi analitis sumber:} interaksi ukuran-konektivitas dapat
  dipakai untuk mengenali apakah kota kecil dan besar merespons posisi
  jaringan secara berbeda (Bagian 4.4 dan Gambar 1-3, hlm. 193-195).
\item
  \textbf{Inferensi review:} dengan data yang sebanding, rancangan dapat
  digunakan sebagai diagnosis awal sistem kota lain untuk memetakan
  asosiasi massa-jaringan-fungsi. Penerapan di luar 16 negara studi
  memerlukan delimitasi kota, fungsi, jaringan, dan kontrol yang sesuai
  konteks; validitas eksternal tersebut tidak diuji dalam A02 (Bagian 3,
  hlm. 186-189).
\item
  \textbf{Inferensi review:} kerangka dapat membantu penyaringan awal
  fungsi yang layak dianalisis lebih lanjut dalam kebijakan jaringan
  regional. Artikel tidak menyediakan perangkat lunak, ambang keputusan
  kebijakan, metode identifikasi pasangan pemberi-penerima manfaat, atau
  evaluasi proyek operasional. Rincian aplikasi tersebut: Tidak
  disebutkan dalam file (Bagian 3-4, hlm. 186-195).
\end{itemize}

\subsection{Future Research
(Optional)}\label{future-research-optional-17}

Sumber menyebut agenda penelitian berikut secara eksplisit pada Bagian
5, hlm. 196.

\begin{enumerate}
\def\labelenumi{\arabic{enumi}.}
\tightlist
\item
  Menambahkan perspektif dinamis untuk menguji apakah konektivitas
  jaringan makin penting dari waktu ke waktu dibandingkan faktor lokal.
\item
  Memasukkan jenis fungsi metropolitan lain, indikator kinerja ekonomi,
  serta tipologi kota yang lebih luas daripada kelas ukuran.
\item
  Menangani kemungkinan kausalitas terbalik antara fungsi metropolitan,
  ukuran, dan konektivitas. Penulis sekaligus mengakui bahwa ketiadaan
  instrumen konektivitas yang kredibel dan data historis membuat agenda
  ini sulit.
\item
  Memperkuat dasar konseptual dan validasi empiris ekonomi jaringan kota
  serta menempatkan \emph{borrowed size} lebih kokoh dalam teori
  aglomerasi.
\item
  Membedakan ragam eksternalitas jaringan dan menguji saluran
  \emph{matching} serta \emph{learning} memakai data lain, seperti
  kolaborasi R\&D, jaringan ko-inventor, dan jaringan perdagangan.
\item
  Meneliti cara mengatasi dominasi \emph{agglomeration shadow} regional,
  termasuk menguji apakah integrasi fungsional, institusional, dan
  kultural antarkota dapat menguranginya.
\end{enumerate}

Review ini tidak menambahkan agenda riset eksternal di luar yang
dinyatakan artikel.

\subsection{Provenance Notes}\label{provenance-notes-17}

\begin{itemize}
\tightlist
\item
  Review ini hanya menggunakan seluruh isi file
  \texttt{/Users/mac/Documents/Mac/{[}2{]}\ Obsidian\ Vault/zettelkasten/source/journal/A02-meijers-burger-hoogerbrugge-2016-borrowing-size-in-networks-of-cities-city-size-network-connectivity-and-metropolitan-functions-in-europe-10.1111-pirs.12181.txt},
  sebanyak 1.033 baris. Tidak ada sumber eksternal, penelusuran web,
  basis data bibliografis, PDF lain, atau catatan kajian lain yang
  digunakan untuk menambah atau mengoreksi klaim substantif.
\item
  Blok metadata TXT menyatakan bahwa teks diekstrak pada 31 Agustus 2026
  dengan \texttt{pdftotext\ -layout} dari PDF
  \texttt{journal/A02-meijers-burger-hoogerbrugge-2016-borrowing-size-in-networks-of-cities-city-size-network-connectivity-and-metropolitan-functions-in-europe-10.1111-pirs.12181.pdf}.
  PDF terdiri atas 19 halaman, tidak terenkripsi, tidak bertanda
  struktur, dan menggunakan PDF versi 1.4. Review ini membaca TXT, bukan
  memeriksa PDF secara visual.
\item
  Judul dan tiga nama penulis berasal dari halaman judul artikel. Kolom
  \texttt{Author} pada metadata PDF hanya mencantumkan Evert Meijers dan
  Martijn Burger serta menghilangkan Marloes Hoogerbrugge. Review
  mempertahankan tiga penulis yang terlihat pada tubuh artikel dan
  mencatat ketidakselarasan metadata ini tanpa verifikasi eksternal.
\item
  Footer artikel menyebut \emph{Papers in Regional Science}, Volume 95,
  Nomor 1, Maret 2016. Halaman tubuh yang terlihat berurutan dari 182
  sampai 198; halaman judul yang tepat mendahului 182 diperlakukan
  sebagai hlm. 181. Rentang 181-198 karena itu merupakan pembacaan
  urutan halaman dalam TXT, bukan metadata rentang halaman yang
  terpisah.
\item
  Halaman artikel utama mencantumkan hak cipta 2015, sedangkan halaman
  terakhir yang berisi abstrak bahasa Spanyol dan Jepang mencantumkan
  hak cipta 2016. Review tidak menafsirkan perbedaan ini sebagai
  perubahan tahun publikasi; identitas terbit mengikuti footer Volume
  95, Nomor 1, Maret 2016.
\item
  Locator \texttt{hlm.} mengikuti nomor halaman cetak yang tampak pada
  hasil ekstraksi. Rujukan bagian, Tabel 1-4, dan Gambar 1-3 mengikuti
  penomoran artikel. Nomor baris TXT tidak dipakai sebagai pengganti
  locator substantif.
\item
  Tata letak Tabel 4 sangat terdistorsi oleh ekstraksi: sebagian tanda
  minus, batas interval kepercayaan, pemisah desimal, dan kesejajaran
  kolom tidak dapat dibaca dengan aman. Review menggunakan tabel hanya
  ketika nilai atau labelnya cukup jelas dan memakai narasi hasil untuk
  arah hubungan regional per fungsi. Nilai yang tidak dapat dipastikan
  tidak direkonstruksi.
\item
  Tabel 3 hasil ekstraksi menampilkan konektivitas regional dengan
  koefisien 0,024 dan AME 0,0005, sedangkan narasi menyebut ``tanda
  negatif''. Review melaporkan kedua bentuk tersebut dan tidak memilih
  atau memperbaiki salah satunya tanpa melihat PDF asli (Tabel 3 dan
  Bagian 4.1, hlm. 189-190).
\item
  Bagian 4.4 menyebut kota berpenduduk lebih dari 150.000 jiwa sebagai
  85\% kota, tetapi Bagian 4.2 menyatakan persentil ke-75 hanya 112.000
  jiwa. Kedua informasi tidak mungkin selaras jika kata ``lebih dari''
  dibaca secara harfiah. Review tidak mengubah arah ketimpangan atau
  persentasenya.
\item
  Gambar 1-3 tidak menyertakan nilai titik atau kurva yang dapat dibaca
  dari TXT. Review hanya menggunakan caption dan uraian naratif penulis,
  bukan menaksir nilai dari gambar.
\item
  Abstrak bahasa Spanyol terbaca dalam ekstraksi, sedangkan sebagian
  besar abstrak bahasa Jepang menjadi karakter rusak. Keduanya mengulang
  isi abstrak dan tidak digunakan untuk menambah hasil.
\item
  Daftar pustaka tersedia pada hlm. 196-198, tetapi isi karya yang
  dirujuk tidak diperiksa secara independen. Semua uraian penelitian
  terdahulu dalam review diperlakukan sebagai laporan A02 atas
  literatur, bukan bukti langsung dari karya-karya tersebut.
\item
  Status telaah sejawat, sumber dan tahun khusus populasi lokal serta
  PDB per kapita, perangkat lunak analisis, kode, data replikasi, dan
  spesifikasi lengkap seluruh model: Tidak disebutkan dalam file.
\item
  Batas provenance utama: review dapat menyatakan apa yang dilaporkan
  dalam ekstraksi dan mengaudit konsistensi internalnya, tetapi tidak
  dapat memastikan kesetiaan tata letak terhadap PDF, membaca detail
  visual grafik, memverifikasi metadata jurnal secara eksternal, menguji
  ulang model, atau memastikan mekanisme kausal.
\end{itemize}

\section{Review A03}\label{review-a03}

\textbf{Identitas sumber}

\begin{itemize}
\tightlist
\item
  Kode: A03.
\item
  Judul: \emph{Borrowed Size, Agglomeration Shadows and Cultural
  Amenities in North-West Europe}.
\item
  Penulis: Martijn J. Burger, Evert J. Meijers, Marloes M. Hoogerbrugge,
  dan Jaume Masip Tresserra.
\item
  Afiliasi: Department of Applied Economics, Tinbergen Institute,
  Erasmus University Rotterdam, Belanda; Faculty of Architecture and the
  Built Environment, Delft University of Technology, Belanda; serta
  Centre of Land Policy and Valuations, Polytechnic University of
  Catalonia, Barcelona, Spanyol (halaman judul, hlm. 1090).
\item
  Jenis sumber: artikel empiris dalam jurnal. Status telaah sejawat
  secara eksplisit: Tidak disebutkan dalam file. Bagian Acknowledgements
  menyatakan bahwa penulis menerima masukan dari dua penelaah anonim,
  tetapi file tidak memberi keterangan formal tentang status proses
  telaah (hlm. 1104-1105).
\item
  Jurnal: \emph{European Planning Studies}.
\item
  Tahun, volume, nomor, dan halaman: 2015, Volume 23, Nomor 6, hlm.
  1090-1109 (header publikasi, hlm. 1090).
\item
  DOI dan URL DOI: 10.1080/09654313.2014.905002;
  http://dx.doi.org/10.1080/09654313.2014.905002 (header publikasi, hlm.
  1090).
\item
  Riwayat naskah: diterima Juli 2013 dan disetujui Januari 2014 (halaman
  judul, hlm. 1090). Tanggal yang lebih terperinci: Tidak disebutkan
  dalam file.
\item
  Korespondensi: Martijn J. Burger, Erasmus University Rotterdam; alamat
  surel \texttt{mburger@ese.eur.nl} (halaman judul, hlm. 1090).
\item
  Cakupan empiris: 2.794 kota, kota kecil, dan desa di tujuh negara
  Eropa Barat, yang ditempatkan dalam 417 \emph{functional urban areas}
  (FUA); data amenitas budaya terutama merujuk 2006-2007 (Bagian 2.1,
  hlm. 1094-1097).
\item
  Pendanaan: Dutch Science Council NWO dan Platform31. Riset merupakan
  bagian dari proyek \emph{Networks, Agglomeration and Polycentric
  Metropolitan Areas: New Perspectives for Improved Economic
  Performance} dalam program \emph{Knowledge for Strong Cities} (KKS)
  (Acknowledgements dan Funding, hlm. 1104-1105).
\item
  Pernyataan konflik kepentingan dan kata kunci: Tidak disebutkan dalam
  file.
\item
  Panjang berkas sumber: metadata PDF menyebut 21 halaman; rentang
  halaman artikel yang tercetak adalah 1090-1109.
\end{itemize}

\textbf{Penanda epistemik yang digunakan}

\begin{itemize}
\tightlist
\item
  \textbf{Laporan sumber} berarti pertanyaan, data, metode, angka, atau
  hasil yang dinyatakan langsung dalam artikel.
\item
  \textbf{Laporan sumber atas literatur} berarti ringkasan artikel
  terhadap karya yang dirujuk, bukan hasil empiris A03 dan bukan hasil
  pemeriksaan independen terhadap rujukan tersebut.
\item
  \textbf{Interpretasi penulis} berarti makna, mekanisme, generalisasi,
  atau implikasi yang ditarik oleh penulis dari hasilnya.
\item
  \textbf{Inferensi review} berarti penilaian analitis yang dibuat hanya
  dari TXT A03. Inferensi ini tidak diperlakukan sebagai bukti tambahan
  dan tidak disandarkan pada sumber eksternal.
\end{itemize}

\subsection{Summarized Abstract}\label{summarized-abstract-18}

\textbf{Laporan sumber.} Artikel mengonseptualisasikan dan
mengeksplorasi secara empiris \emph{borrowed size} di Eropa Barat Laut.
Sebuah tempat disebut meminjam ukuran apabila menampung lebih banyak
fungsi perkotaan daripada yang biasanya dapat ditopang oleh ukurannya
sendiri. Sebaliknya, tempat yang menampung lebih sedikit fungsi daripada
yang biasanya dapat ditopangnya berada dalam \emph{agglomeration
shadow}. Kehadiran fungsi perkotaan diproksikan dengan indeks amenitas
budaya kelas atas, sedangkan kemampuan memanfaatkan massa lokal,
regional, serta nasional dan internasional diperiksa melalui ukuran
penduduk dan posisi tempat dalam sistem perkotaan regional (Abstract,
hlm. 1090; Bagian 2, hlm. 1093-1094).

\textbf{Laporan sumber.} Analisis mencakup 2.794 tempat yang tersarang
dalam 417 FUA. Variabel hasilnya adalah \emph{Cultural Amenities Index}
berskala 0-1. Penulis mengestimasi regresi beta dengan inflasi nol dan
satu, lalu memeriksa peran jumlah penduduk lokal, penduduk di sisa FUA,
aksesibilitas nasional dan internasional, peringkat tempat dalam FUA,
serta interaksi antara posisi peringkat dan ketiga ukuran potensi
penduduk (Bagian 2.1-2.2, hlm. 1094-1097; Tabel 3, hlm. 1098-1099;
Lampiran 1, hlm. 1109).

\textbf{Interpretasi penulis.} Ukuran penduduk lokal merupakan penjelas
terpenting bagi keberadaan amenitas budaya kelas atas. Tempat terbesar
dalam FUA lebih mampu memanfaatkan massanya sendiri, penduduk FUA di
sekitarnya, dan aksesibilitas nasional-internasional. Amenitas secara
tidak proporsional terkonsentrasi di tempat berperingkat pertama,
sedangkan tempat berperingkat lebih rendah cenderung berada dalam
bayang-bayang aglomerasi. Penulis menilai pola ini lebih dekat dengan
logika pusat pelayanan Christaller dan prediksi \emph{New Economic
Geography} daripada dengan dugaan bahwa kota kecil umumnya meminjam
ukuran dari tetangga besarnya (Bagian 3.1, hlm. 1100-1103; Bagian 4,
hlm. 1103-1104).

\textbf{Inferensi review.} Bukti A03 menunjukkan asosiasi spasial potong
lintang antara ukuran, posisi hierarkis, aksesibilitas, dan amenitas.
Artikel tidak mengamati perpindahan konsumen antartempat atau
mengidentifikasi secara langsung pasangan tempat yang ``meminjamkan''
dan ``meminjam'' ukuran. Karena itu, \emph{borrowed size} dan
\emph{agglomeration shadow} dalam studi ini merupakan interpretasi atas
ketidakselarasan ukuran-fungsi dan pola koefisien model, bukan
pengukuran langsung aliran manfaat antarkota.

\subsection{Problem Statement}\label{problem-statement-18}

\textbf{Laporan sumber.} Amenitas konsumsi, terutama fungsi kelas atas,
umumnya bergantung kuat pada jumlah penduduk lokal dan secara
tradisional terkonsentrasi di kota besar. Pada saat yang sama, perluasan
skala proses sosial-ekonomi, integrasi metropolitan, dan pembentukan
jaringan kota memungkinkan massa kritis untuk suatu fungsi berasal dari
luar batas tempat yang menampungnya (Bagian 1.1-1.2, hlm. 1090-1091).

\textbf{Laporan sumber.} Interdependensi spasial dapat bekerja dalam dua
arah. Permintaan eksternal dari tempat lain dapat memungkinkan suatu
tempat menampung fungsi yang tidak dapat ditopang sendiri, tetapi
pasokan amenitas di tempat lain juga dapat menarik konsumen lokal dan
mengurangi fungsi yang dapat dipertahankan. Artikel menyatakan bahwa
masih belum jelas tempat mana yang memperoleh manfaat melalui
\emph{borrowed size} dan tempat mana yang kehilangan fungsi melalui
\emph{agglomeration shadow} (Bagian 1.3-1.4, hlm. 1091-1093).

\textbf{Rumusan masalah sumber.} Pertanyaan umumnya ialah: tempat mana
yang meminjam ukuran dan tempat mana yang menghadapi bayang-bayang
aglomerasi? Pertanyaan khusus artikel ialah apakah posisi suatu tempat
dalam sistem perkotaan regional memengaruhi kapasitasnya untuk meminjam
ukuran pada skala lokal, regional, serta nasional dan internasional.
Pertanyaan ini diperiksa melalui hubungan antara penduduk lokal, potensi
penduduk di luar tempat, dan keberadaan amenitas budaya kelas atas
(Bagian 1.4, hlm. 1093).

\textbf{Inferensi review.} Masalah empirisnya bukan hanya apakah
jaringan perkotaan menyediakan massa tambahan, melainkan siapa yang
dapat menangkap massa tersebut. Peringkat dalam FUA digunakan sebagai
proksi kapasitas penangkapan itu, sehingga artikel menguji heterogenitas
menurut hierarki regional alih-alih mengasumsikan seluruh tempat dalam
jaringan mendapat manfaat yang sama.

\subsection{Objectives}\label{objectives-18}

\begin{itemize}
\tightlist
\item
  \textbf{Laporan sumber:} mengonseptualisasikan \emph{borrowed size}
  dan \emph{agglomeration shadow} sebagai ketidakselarasan antara ukuran
  tempat dan fungsi perkotaan yang ditampungnya (Abstract, hlm. 1090;
  Bagian 2, hlm. 1093-1094).
\item
  \textbf{Laporan sumber:} mengeksplorasi seberapa kuat ukuran dan
  fungsi berhubungan pada tempat-tempat di Eropa Barat Laut dengan
  amenitas budaya kelas atas sebagai proksi fungsi perkotaan (Abstract,
  hlm. 1090; Bagian 1.4, hlm. 1093).
\item
  \textbf{Laporan sumber:} menilai pengaruh jumlah penduduk lokal,
  jumlah penduduk di bagian FUA selain tempat itu, dan aksesibilitas
  nasional-internasional terhadap keberadaan amenitas budaya (Bagian
  1.4, hlm. 1093; Bagian 2.1, hlm. 1094-1097).
\item
  \textbf{Laporan sumber:} menguji apakah peringkat suatu tempat dalam
  FUA memengaruhi keberadaan amenitas dan memoderasi kemampuan tempat
  tersebut untuk memanfaatkan potensi penduduk pada ketiga skala (Bagian
  2.1, hlm. 1096-1097).
\item
  \textbf{Laporan sumber:} menjelaskan mengapa satu tempat tampak
  meminjam ukuran sedangkan tempat lain menghadapi bayang-bayang
  aglomerasi melalui posisi masing-masing dalam sistem perkotaan
  regional (Abstract, hlm. 1090).
\end{itemize}

\subsection{Literature Survey}\label{literature-survey-18}

\textbf{Inferensi review tentang sifat survei.} Artikel menggunakan
tinjauan naratif untuk membangun kerangka konseptual dan ekspektasi
empiris. Strategi pencarian, pangkalan data, rentang penelusuran,
kriteria inklusi-eksklusi, jumlah karya yang disaring, serta penilaian
mutu literatur: Tidak disebutkan dalam file. Oleh karena itu, bagian
pustaka A03 tidak dapat diperlakukan sebagai tinjauan sistematis.

\begin{itemize}
\tightlist
\item
  \textbf{Laporan sumber atas literatur kota konsumsi:} kota tidak lagi
  hanya dipahami sebagai pusat produksi, tetapi juga sebagai pusat
  konsumsi. Upah nominal meningkat bersama ukuran kota, tetapi biaya
  hidup meningkat lebih besar, yang oleh literatur ditafsirkan sebagai
  kesediaan penduduk membayar untuk amenitas kota besar. Amenitas
  meliputi aset estetis dan warisan, barang dan jasa terspesialisasi,
  serta layanan publik khusus (Bagian 1.1, hlm. 1090-1091).
\item
  \textbf{Laporan sumber atas literatur skala dan fungsi:} amenitas
  konsumsi, terutama fungsi kelas atas, berkorelasi kuat dengan jumlah
  penduduk lokal dan mencerminkan ekonomi urbanisasi melalui konsumsi.
  Universitas, rumah sakit, stadion, gedung konser, teater, galeri,
  museum, dan ritel berorde tinggi secara tradisional berada di kota
  besar (Bagian 1.1, hlm. 1090-1091).
\item
  \textbf{Laporan sumber atas literatur jaringan kota:} perkembangan
  pusat suburban, keterhubungan fungsional antarkawasan yang dahulu
  terpisah, wilayah perkotaan polisentris, dan kawasan megakota
  memperluas ruang interaksi. Tempat-tempat dalam jaringan dapat
  memiliki spesialisasi sektoral dan fungsional yang saling melengkapi,
  sehingga hubungan ukuran-fungsi tidak harus mengikuti hierarki
  sederhana (Bagian 1.2, hlm. 1091).
\item
  \textbf{Laporan sumber atas konsep \emph{borrowed size}:} Alonso
  (1973) digunakan sebagai sumber konsep bahwa kota kecil atau wilayah
  metropolitan dapat memperlihatkan karakteristik kota lebih besar
  apabila berdekatan dengan konsentrasi penduduk lain. Menurut ringkasan
  A03, kota kecil berpotensi memperoleh sebagian manfaat aglomerasi
  tetangga besar sambil menghindari biaya aglomerasi, dan pola Jerman
  serta Negara-Negara Rendah pernah diajukan sebagai contoh penting
  (Bagian 1.3, hlm. 1091-1092).
\item
  \textbf{Laporan sumber atas konsep \emph{agglomeration shadow}:} Teori
  Tempat Sentral, teori sistem perkotaan, dan \emph{New Economic
  Geography} digunakan untuk menjelaskan bagaimana pusat berorde tinggi
  membatasi pertumbuhan atau fungsi wilayah di sekitarnya melalui
  persaingan. Kedekatan dengan kota besar dapat membuat tempat berukuran
  serupa memiliki lebih sedikit amenitas daripada tempat yang terisolasi
  (Bagian 1.3, hlm. 1092-1093).
\item
  \textbf{Laporan sumber atas mekanisme aglomerasi amenitas:} fungsi
  kelas atas memilih lokasi dekat basis konsumen terbesar. Pengelompokan
  amenitas juga menghemat biaya perjalanan, pencarian, dan transaksi
  dalam perjalanan multitujuan serta meningkatkan daya tarik bagi
  produsen dan konsumen. Mekanisme ini memberi keunggulan kumulatif
  kepada pusat yang sudah memiliki banyak amenitas (Bagian 1.3, hlm.
  1092-1093).
\item
  \textbf{Laporan sumber atas dua gaya jaringan:} tambahan permintaan
  eksternal dapat menopang amenitas di suatu tempat, sedangkan tambahan
  pasokan eksternal dapat menarik konsumennya. Dengan demikian,
  \emph{borrowed size} dan bayang-bayang aglomerasi diposisikan sebagai
  dua kemungkinan yang berlawanan dari interdependensi spasial yang sama
  (Bagian 1.4, hlm. 1093).
\item
  \textbf{Laporan sumber atas skala integrasi:} catatan artikel merujuk
  pembagian polisentralitas mikro sebagai integrasi internal wilayah
  metropolitan, meso sebagai jaringan perkotaan nasional yang rapat, dan
  makro sebagai keterlekatan wilayah metropolitan dalam jaringan
  metropolis transnasional atau global (Catatan 3, hlm. 1105).
\item
  \textbf{Laporan sumber atas kritik:} artikel mengarahkan pembaca
  kepada Peck (2005) serta Storper dan Scott (2009) untuk kritik
  terhadap hubungan amenitas dan pertumbuhan, tetapi isi kritik tersebut
  tidak diringkas dalam file (Catatan 2, hlm. 1105).
\end{itemize}

\textbf{Interpretasi penulis.} Penulis menggabungkan paradigma jaringan
kota dengan logika hierarki pusat pelayanan. Kedekatan dan konektivitas
tidak otomatis menghasilkan pembagian manfaat yang merata karena
keuntungan kolokasi dan besarnya basis konsumen dapat membuat pusat
berperingkat pertama menangkap permintaan dari seluruh wilayah (Bagian
1.3-1.4, hlm. 1091-1093).

\textbf{Inferensi review.} Kerangka tersebut menghasilkan ketegangan
teoretis yang dapat diuji: Alonso mengarahkan perhatian pada manfaat
jaringan bagi kota kecil, sedangkan Teori Tempat Sentral dan \emph{New
Economic Geography} mengarahkan perhatian pada dominasi pusat besar. A03
mengoperasionalkan ketegangan ini melalui interaksi peringkat FUA dengan
potensi penduduk, tetapi tidak melakukan pengujian formal yang
membandingkan keseluruhan teori sebagai model alternatif.

\subsection{Method Used}\label{method-used-18}

\textbf{Laporan sumber.} Unit analisis berada pada tingkat tempat, yang
sebagian besar didefinisikan sebagai unit LAU-2. Tempat-tempat tersebut
tersarang dalam FUA yang mewakili wilayah pasar tenaga kerja dengan
interdependensi sosial dan ekonomi tinggi serta radius maksimum 60
kilometer (Bagian 2.1, hlm. 1094-1096).

\textbf{Inferensi review atas bentuk data.} Karena variabel terutama
merujuk satu potret 2006 dan data amenitas mencakup 2006-2007, review
ini mengklasifikasikan desainnya sebagai kuantitatif potong lintang,
bukan longitudinal (Bagian 2.1, hlm. 1094 dan 1096).

\textbf{Konstruksi variabel hasil.} \emph{Cultural Amenities Index}
dibangun dari enam indikator: jumlah teater; jumlah gedung opera dan
teater musik; jumlah acara musik besar; jumlah dan kualitas lembaga seni
publik; jumlah pameran seni dan festival film; serta jumlah galeri seni.
Tempat dengan skor rata-rata tertinggi pada enam dimensi diberi nilai 1,
sedangkan tempat lain diberi skor relatif terhadap tempat tertinggi,
dengan nilai minimum 0. Rincian pengukuran setiap indikator dirujuk
penulis ke BBSR (2011) dan tidak direproduksi lengkap dalam file (Bagian
2.1, hlm. 1094).

\textbf{Variabel penjelas utama.} Jumlah penduduk lokal diukur pada
tingkat LAU-2 per 100.000 penduduk. Potensi regional diukur sebagai
jumlah penduduk, dalam juta orang, di bagian FUA selain tempat yang
diamati. Potensi nasional dan internasional diukur dengan aksesibilitas
absolut multimoda terhadap penduduk di luar FUA, dibobot menurut waktu
perjalanan melalui jalan, kereta, dan udara, serta dinyatakan per 10
juta penduduk (Bagian 2.1, hlm. 1094-1096; Tabel 2, hlm. 1096).

\textbf{Posisi dalam FUA.} Penulis mula-mula menggunakan variabel
indikator apakah tempat bukan yang terbesar dalam FUA. Model berikutnya
membedakan tempat terbesar, terbesar kedua, terbesar ketiga, terbesar
keempat, dan tempat lain yang berperingkat lebih rendah. Interaksi
antara status bukan tempat terbesar dan masing-masing variabel penduduk
digunakan untuk menilai apakah posisi hierarkis memoderasi kemampuan
mengeksploitasi massa lokal, regional, dan nasional-internasional
(Bagian 2.1, hlm. 1096-1097; Tabel 3, hlm. 1098-1099; Lampiran 1, hlm.
1109).

\textbf{Variabel kontrol.} Model memasukkan indikator keberadaan situs
warisan yang dalam TXT disebut ``UNICEF heritage site'', PDB per kapita
pada tingkat FUA, dan indikator negara. Indikator negara dimaksudkan
menangkap, antara lain, perbedaan kelembagaan dalam penyediaan barang
dan jasa publik, termasuk fasilitas budaya (Bagian 2.1, hlm. 1096).
Istilah ``UNICEF'' dipertahankan sesuai file dan tidak dikoreksi dengan
pengetahuan eksternal.

\textbf{Strategi estimasi.} Karena indeks hasil dibatasi pada 0 dan 1
serta massa probabilitas sangat terkonsentrasi pada nol, penulis menilai
regresi linear dapat menghasilkan estimasi tidak efisien dan bias,
sedangkan logit fraksional dapat menghadapi overdispersi. Mereka
menggunakan regresi beta dengan inflasi nol dan satu. Model terdiri atas
regresi logistik untuk peluang nilai 0, regresi logistik untuk peluang
nilai 1, dan regresi beta untuk proporsi di antara 0 dan 1; parameter
\texttt{phi} digunakan untuk menyesuaikan varians kondisional (Bagian
2.2, hlm. 1097).

\textbf{Spesifikasi dan pelaporan.} Model 1 merupakan spesifikasi dasar;
Model 2 menambahkan indikator bukan tempat terbesar; Model 3 memerinci
peringkat kedua, ketiga, keempat, dan peringkat lain. Model 4-6
masing-masing menambahkan interaksi posisi dengan penduduk lokal,
penduduk di sisa FUA, dan aksesibilitas nasional-internasional. Seluruh
model memakai intersep, indikator negara, dan galat baku robust. Bagian
inflasi satu hanya berupa intersep sehingga tidak ditampilkan. Penulis
melaporkan koefisien dan \emph{average marginal effects} (AME) untuk
Model 1-3; AME interaksi Model 4-6 disajikan melalui Gambar 4-6,
sedangkan Lampiran 1 hanya menampilkan koefisien (Tabel 3, hlm.
1098-1099; Bagian 3.1, hlm. 1101-1103; Lampiran 1, hlm. 1109).

\textbf{Pemeriksaan spesifikasi tambahan.} Penulis menyatakan telah
menjalankan regresi serupa yang menginteraksikan urutan peringkat
terperinci dengan penduduk lokal, penduduk di sisa FUA, dan
aksesibilitas nasional-internasional. Grafiknya tidak ditampilkan demi
keringkasan, tetapi penulis menyatakan hasilnya tidak mengubah
kesimpulan tentang \emph{borrowed size} dan persaingan spasial (Catatan
10, hlm. 1105).

\textbf{Inferensi review atas desain.} Model mengestimasi asosiasi
kondisional pada satu potret waktu. Tidak ada perlakuan eksperimental,
urutan waktu sebelum-sesudah, instrumen, atau strategi identifikasi
kausal yang dilaporkan. Bahasa mengenai ``kenaikan'' penduduk dan
``dampak'' aksesibilitas dalam uraian AME harus dibaca sebagai perubahan
prediksi model dengan kovariat lain tetap, bukan efek kausal perubahan
aktual.

\subsection{Dataset}\label{dataset-18}

\begin{itemize}
\tightlist
\item
  \textbf{Data fungsi metropolitan BBSR:} basis data German Federal
  Institute for Research on Building, Urban Affairs and Spatial
  Development menyediakan informasi fungsi metropolitan pada
  tempat-tempat di bagian Eropa yang terurbanisasi. A03 menggunakan data
  acara dan fasilitas budaya 2006-2007 untuk 2.794 tempat (Bagian 2.1,
  hlm. 1094).
\item
  \textbf{Enam komponen amenitas budaya:} data mencakup teater, opera
  dan teater musik, acara musik besar, lembaga seni publik, pameran seni
  dan festival film, serta galeri seni. Indikator acara musik besar
  menggunakan tanggal tur The Rolling Stones, Madonna, Sting, Bon Jovi,
  Anna Netrebko, Anne-Sophie Mutter, Vienna Philharmonic, New York
  Philharmonic Orchestra, pertunjukan \emph{Cats}, serta festival jazz
  internasional. Lembaga seni publik mencakup museum dan galeri publik,
  klub seni privat atau nirlaba, yayasan dan asosiasi, arsip dan
  perpustakaan terkait, serta perguruan tinggi seni (Bagian 2.1, hlm.
  1094; Catatan 5-6, hlm. 1105).
\item
  \textbf{Data FUA:} definisi dan batas FUA mengikuti proyek ESPON 1.4.3
  yang dirujuk sebagai Peeters (2011). FUA dipahami sebagai wilayah
  pasar tenaga kerja yang menyatukan kota, kota kecil, dan desa yang
  saling bergantung secara ekonomi dan sosial (Bagian 2.1, hlm.
  1095-1096).
\item
  \textbf{Data aksesibilitas:} ukuran potensi aksesibilitas absolut
  multimoda berasal dari riset ESPON yang dirujuk sebagai Spiekermann
  dan Wegener (2006). Ukuran awalnya dikembangkan pada tingkat NUTS-3,
  kemudian diadaptasi ke FUA untuk A03 (Bagian 2.1, hlm. 1096; Catatan
  7, hlm. 1105).
\item
  \textbf{Profil aksesibilitas:} wilayah yang paling mudah diakses
  disebut berada di megaregion dengan bandar udara internasional besar,
  seperti Rhine-Main di sekitar Frankfurt am Main dan Randstad di
  sekitar Amsterdam. Wilayah yang paling sulit diakses berada di bagian
  periferal Prancis, Jerman, dan Austria (Bagian 2.1, hlm. 1096).
\item
  \textbf{Data penduduk lokal:} artikel melaporkan ukuran dan
  statistiknya, tetapi nama sumber data penduduk lokal secara khusus:
  Tidak disebutkan dalam file (Bagian 2.1, hlm. 1094-1096; Tabel 2, hlm.
  1096).
\item
  \textbf{Data kontrol:} PDB per kapita pada tingkat FUA, indikator
  situs warisan, dan negara dimasukkan ke model. Sumber khusus, definisi
  rinci, mata uang atau paritas PDB, dan prosedur pencocokan data
  kontrol: Tidak disebutkan dalam file (Bagian 2.1, hlm. 1096-1097).
\item
  \textbf{Waktu pengukuran:} penulis menyatakan seluruh variabel diukur
  untuk 2006 kecuali jika disebutkan lain, sedangkan data acara dan
  fasilitas budaya disebut mencakup 2006-2007 (Bagian 2.1, hlm. 1094 dan
  1096).
\item
  \textbf{Statistik deskriptif:} indeks amenitas memiliki rerata 0,0047,
  simpangan baku 0,037, minimum 0, dan maksimum 1. Penduduk lokal dalam
  satuan 100.000 memiliki rerata 0,32, simpangan baku 1,13, minimum
  0,00061, dan maksimum 34,61. Penduduk di sisa FUA dalam juta memiliki
  rerata 2,16, simpangan baku 3,07, minimum 0, dan maksimum 11,17 (Tabel
  2, hlm. 1096).
\item
  \textbf{Statistik deskriptif lanjutan:} aksesibilitas per 10 juta
  penduduk memiliki rerata 6,85, simpangan baku 1,51, minimum 2,96, dan
  maksimum 10,01. PDB per kapita FUA dalam satuan ribuan memiliki rerata
  32,39, simpangan baku 8,47, minimum 13, dan maksimum 54. Proporsi
  observasi dengan situs warisan adalah sekitar 0,02, sedangkan proporsi
  tempat yang bukan terbesar dalam FUA adalah sekitar 0,85 (Tabel 2,
  hlm. 1096).
\item
  \textbf{Ketersediaan dan reproduksibilitas:} tautan unduh data
  analisis, lisensi, versi berkas, kamus variabel, prosedur pembersihan,
  penanganan data hilang, dan kode analisis: Tidak disebutkan dalam
  file. BBSR disebut telah membagikan data kepada penulis
  (Acknowledgements, hlm. 1104).
\end{itemize}

\subsection{Population / Sample}\label{population-sample-18}

\textbf{Unit analisis.} Sampel analitis terdiri atas 2.794 kota, kota
kecil, dan desa, sebagian besar pada tingkat LAU-2, yang berada dalam
417 FUA. Dengan demikian, observasi utama adalah tempat, bukan individu,
rumah tangga, pengunjung amenitas, atau organisasi budaya (Bagian 2.1,
hlm. 1094-1096).

\textbf{Cakupan negara.} Sebanyak 35\% unit berada di Prancis, 29\% di
Jerman, 14\% di Swiss, 8\% di Belanda, 7\% di Austria, 6\% di Belgia,
dan 1\% di Luksemburg. Artikel bergantian menyebut cakupan ini sebagai
Eropa Barat dan Eropa Barat Laut (Bagian 2.1, hlm. 1094; judul dan
Gambar 1, hlm. 1090 dan 1095).

\textbf{Rentang ukuran tempat.} Jumlah penduduk tempat berkisar dari 61
orang di Gerstengrund hingga hampir 3,5 juta di Berlin. Rerata ukuran
tempat sedikit di atas 30.000 penduduk, sedangkan rerata tempat yang
terbesar dalam FUA-nya hampir 130.000 penduduk (Bagian 2.1, hlm.
1094-1096).

\textbf{Struktur FUA.} Setiap FUA rata-rata mencakup 6-7 unit LAU-2.
Ukurannya disebut berkisar dari wilayah megakota seperti Paris dengan
11,1 juta penduduk, Milano dengan 6,0 juta, dan Dortmund dengan 5,3 juta
hingga desa mandiri di kawasan lebih perdesaan dengan sedikit di atas
10.000 penduduk. FUA memiliki radius maksimum 60 kilometer (Bagian 2.1,
hlm. 1095-1096).

\textbf{Komposisi hierarkis.} Sekitar 85\% observasi bukan tempat
terbesar dalam FUA. Tabel 2 melaporkan sekitar 11\% sebagai tempat
terbesar kedua, 8\% terbesar ketiga, 7\% terbesar keempat, dan 58\%
sebagai tempat bukan terbesar pada kategori peringkat lainnya (Tabel 2,
hlm. 1096).

\textbf{Batas informasi sampel.} Kerangka pemilihan tempat, aturan
inklusi tempat tanpa amenitas, kelengkapan cakupan tiap negara, jumlah
observasi yang dikeluarkan, dan penanganan perubahan batas LAU-2 atau
FUA: Tidak disebutkan dalam file. Contoh Milano muncul dalam uraian
rentang FUA meskipun Italia tidak muncul dalam komposisi tujuh negara;
file tidak menjelaskan apakah contoh itu berasal dari basis ESPON yang
lebih luas atau sampel analitis (Bagian 2.1, hlm. 1094-1096).

\subsection{Dependent Variables}\label{dependent-variables-18}

\textbf{Laporan sumber.} Variabel dependen tunggal adalah \emph{Cultural
Amenities Index}, sebuah ukuran gabungan amenitas budaya kelas atas pada
tingkat tempat. Nilainya dibatasi dari 0 untuk tempat tanpa amenitas
yang tercakup hingga 1 untuk tempat dengan skor rata-rata tertinggi pada
enam dimensi (Bagian 2.1-2.2, hlm. 1094 dan 1097).

\textbf{Komponen distribusi model.} Spesifikasi inflasi nol memodelkan
apakah indeks sama dengan 0, spesifikasi inflasi satu memodelkan apakah
indeks sama dengan 1, dan bagian proporsi memodelkan nilai indeks yang
berada di antara 0 dan 1. Ketiganya merupakan bagian pemodelan atas
variabel hasil yang sama, bukan tiga variabel dependen substantif yang
berbeda (Bagian 2.2, hlm. 1097).

\textbf{Klarifikasi operasional.} \emph{Borrowed size} dan
\emph{agglomeration shadow} tidak menjadi variabel biner yang diamati
langsung. Keduanya disimpulkan dari apakah fungsi tampak lebih banyak
atau lebih sedikit daripada yang dapat ditopang ukuran tempat serta dari
hubungan indeks dengan massa eksternal dan posisi FUA (Bagian 2, hlm.
1093-1094).

\textbf{Variabel bukan dependen.} Penduduk lokal, penduduk di sisa FUA,
aksesibilitas nasional-internasional, peringkat dalam FUA, interaksi
peringkat, PDB per kapita, situs warisan, dan indikator negara adalah
prediktor atau kontrol, bukan keluaran penelitian (Bagian 2.1, hlm.
1094-1097).

\subsection{Results}\label{results-18}

\textbf{Distribusi amenitas.}

\begin{itemize}
\tightlist
\item
  \textbf{Laporan sumber:} 78\% tempat memiliki nilai nol pada indeks.
  Rerata seluruh sampel hanya 0,0047, yang menunjukkan konsentrasi kuat
  amenitas kelas atas pada sejumlah kecil tempat (Bagian 2.1, hlm. 1094;
  Tabel 2, hlm. 1096).
\item
  \textbf{Laporan sumber:} tempat berpenduduk kurang dari 50.000
  memiliki skor rata-rata 0,001; tempat berpenduduk 100.000-350.000
  memiliki skor rata-rata 0,029; dan 25 tempat terbesar, yang
  masing-masing berpenduduk lebih dari 350.000, memiliki skor rata-rata
  0,257 (Bagian 2.1, hlm. 1094).
\item
  \textbf{Laporan sumber:} Berlin menempati skor tertinggi, yaitu 1,00,
  diikuti Paris 0,96 dan Vienna 0,65. Tabel 1 juga menunjukkan variasi
  di antara kota besar: Duisburg, misalnya, memiliki penduduk 0,49 juta
  tetapi skor 0,01, sedangkan Zurich berpenduduk 0,37 juta dan berskor
  0,17 (Tabel 1, hlm. 1095).
\end{itemize}

\textbf{Model dasar dan kontrol.}

\begin{itemize}
\tightlist
\item
  \textbf{Laporan sumber:} jumlah penduduk lokal merupakan determinan
  terpenting dalam Model 1. AME sebesar sekitar 0,007 berarti tambahan
  100.000 penduduk berasosiasi dengan kenaikan prediksi indeks rata-rata
  sekitar 0,007, dengan variabel lain tetap (Bagian 3.1, hlm. 1097-1100;
  Tabel 3, hlm. 1098-1099).
\item
  \textbf{Laporan sumber:} AME penduduk di sisa FUA dan aksesibilitas
  nasional-internasional tidak signifikan dalam spesifikasi dasar. PDB
  per kapita juga tidak menunjukkan efek dalam Model 1, sedangkan
  indikator situs warisan berhubungan positif dengan amenitas; AME
  indikator warisan adalah sekitar 0,0037 (Bagian 3.1, hlm. 1100; Tabel
  3, hlm. 1098-1099).
\item
  \textbf{Laporan sumber:} setelah posisi tempat dalam FUA dikontrol,
  PDB per kapita memiliki AME positif yang kecil. Kenaikan PDB per
  kapita FUA sebesar 1.000 euro berasosiasi dengan kenaikan indeks
  sekitar 0,0001 (Catatan 8, hlm. 1105; Tabel 3, hlm. 1098-1099).
\end{itemize}

\textbf{Ketaklinearan ukuran lokal.}

\begin{itemize}
\tightlist
\item
  \textbf{Laporan sumber:} efek marginal penduduk lokal kecil pada
  tempat berpenduduk kurang dari sekitar 100.000-200.000, kemudian
  meningkat setelah ambang tersebut. Untuk tempat berpenduduk
  150.000-200.000, tambahan 100.000 penduduk berasosiasi dengan kenaikan
  indeks sekitar 0,006 (Gambar 3 dan uraian, hlm. 1100).
\item
  \textbf{Laporan sumber:} untuk kota berpenduduk 500.000, 750.000, dan
  1.000.000, tambahan 100.000 penduduk masing-masing berasosiasi dengan
  kenaikan indeks sekitar 0,013, 0,025, dan 0,042. Lonjakan lain pada
  rentang 10.000-100.000 ditafsirkan sebagai tanda bahwa sebagian
  komponen indeks memiliki ambang penduduk lebih rendah (Gambar 3 dan
  uraian, hlm. 1100).
\end{itemize}

\textbf{Peringkat dalam FUA.}

\begin{itemize}
\tightlist
\item
  \textbf{Laporan sumber:} bukan menjadi tempat terbesar dalam FUA
  berasosiasi dengan skor indeks sekitar 0,003 lebih rendah secara
  rata-rata, dengan variabel lain tetap. Model 3 menunjukkan gradien
  hierarkis: dibandingkan tempat terbesar, tempat terbesar kedua
  memiliki AME sekitar -0,0022, terbesar ketiga -0,0031, terbesar
  keempat -0,0032, dan tempat berperingkat kelima atau lebih rendah
  sekitar -0,0040 (Bagian 3.1, hlm. 1100-1101; Tabel 3, hlm. 1098-1099).
\item
  \textbf{Laporan sumber:} seluruh Model 1-3 menggunakan 2.794 observasi
  dan indikator negara. Statistik Wald masing-masing 246,0, 521,0, dan
  548,3 serta dilaporkan signifikan pada \texttt{p\ \textless{}\ 0,01};
  nilai \texttt{ln\ phi} masing-masing 3,86, 4,06, dan 4,08 (Tabel 3,
  hlm. 1098-1099).
\end{itemize}

\textbf{Interaksi posisi dan potensi penduduk.}

\begin{itemize}
\tightlist
\item
  \textbf{Laporan sumber:} AME penduduk lokal pada indeks lebih kuat
  bagi tempat yang terbesar dalam FUA. Keuntungan moderasi dari status
  terbesar paling menonjol pada tempat berpenduduk sekitar 40.000-80.000
  (Gambar 4 dan uraian, hlm. 1101-1102).
\item
  \textbf{Laporan sumber:} bagi tempat terbesar dalam FUA, tambahan satu
  juta penduduk di sisa FUA berasosiasi dengan kenaikan indeks rata-rata
  sekitar 0,05. Ukuran sisa FUA tidak memengaruhi keberadaan amenitas
  kelas atas pada tempat lain yang bukan terbesar (Gambar 5 dan uraian,
  hlm. 1102-1103).
\item
  \textbf{Laporan sumber:} aksesibilitas nasional-internasional
  berdampak lebih positif pada tempat terbesar dalam FUA. Namun,
  besarannya terbatas: tambahan aksesibilitas 10 juta penduduk hanya
  berasosiasi dengan kenaikan indeks sekitar 0,001-0,002 bahkan pada
  tempat terbesar. Tempat paling mudah diakses diprediksi berskor
  sekitar 0,007-0,014 lebih tinggi daripada tempat paling tidak mudah
  diakses, dengan variabel lain tetap (Gambar 6 dan uraian, hlm.
  1102-1103).
\item
  \textbf{Laporan sumber:} koefisien interaksi Model 4-6 tersedia dalam
  Lampiran 1, tetapi AME interaksi tidak ditampilkan dalam lampiran.
  Penulis mengandalkan Gambar 4-6 untuk interpretasi substantif (Bagian
  3.1, hlm. 1101; Lampiran 1, hlm. 1109).
\end{itemize}

\textbf{Batas pelaporan komponen inflasi.} Tabel 3 juga memuat koefisien
bagian inflasi nol, sedangkan bagian inflasi satu tidak ditampilkan
karena hanya mengandung intersep. Teks utama tidak memberi interpretasi
substantif terperinci atas koefisien inflasi nol. Karena tanda koefisien
pada tabel TXT terdistorsi, review ini tidak merekonstruksi arah atau
nilai yang tidak dikonfirmasi uraian naratif (Tabel 3, hlm. 1098-1099).

\subsection{Findings}\label{findings-18}

\begin{itemize}
\tightlist
\item
  \textbf{Interpretasi penulis:} ukuran tempat sendiri tetap merupakan
  faktor utama bagi amenitas budaya kelas atas; hasil ini menunjukkan
  bahwa ruang lingkup geografis konsumsi masih sangat lokal (Bagian 4,
  hlm. 1103).
\item
  \textbf{Interpretasi penulis:} tempat terbesar dalam FUA menjalankan
  fungsi pusat pelayanan dan lebih mampu mengeksploitasi massa lokal.
  Pusat tersebut juga menangkap permintaan dari sisa FUA, sehingga
  meminjam ukuran dari wilayah sekitarnya sekaligus menimbulkan
  bayang-bayang aglomerasi bagi tempat berperingkat lebih rendah (Bagian
  3.1, hlm. 1100-1103).
\item
  \textbf{Interpretasi penulis:} tempat terbesar juga paling mampu
  memanfaatkan aksesibilitas nasional-internasional, tetapi efek
  aksesibilitas tersebut kecil secara substantif dibandingkan
  konsentrasi yang berhubungan dengan ukuran dan peringkat (Bagian 3.1,
  hlm. 1103).
\item
  \textbf{Interpretasi penulis:} distribusi amenitas kelas atas
  mengikuti logika Christaller, yaitu konsentrasi layanan di tempat
  terbesar dalam FUA karena permintaan lokal yang lebih tinggi dan akses
  terhadap konsumen dari tempat tetangga. Rata-rata, kota besar menaungi
  kota tetangga yang lebih kecil, bukan kota kecil meminjam ukuran dari
  tetangga besarnya sebagaimana dugaan Alonso (Bagian 4, hlm. 1104).
\item
  \textbf{Interpretasi penulis:} kota besar memperoleh manfaat terbesar
  dari \emph{borrowed size}, sedangkan tempat kecil dan menengah yang
  bukan pusat peringkat pertama lebih mungkin mengalami persaingan
  spasial. Penulis merangkum pola ini dengan pernyataan bahwa tempat
  kecil dan menengah di Eropa biasanya berkembang ketika terisolasi
  (Bagian 4, hlm. 1104).
\item
  \textbf{Inferensi review:} hasil dasar yang tidak signifikan untuk
  penduduk sisa FUA menyembunyikan heterogenitas posisi. Setelah
  interaksi diperiksa, massa regional berasosiasi positif hanya bagi
  pusat terbesar. Karena itu, kesimpulan tentang ekonomi jaringan
  bergantung pada siapa yang berada pada puncak hierarki, bukan sekadar
  pada ukuran jaringan secara agregat.
\item
  \textbf{Inferensi review:} istilah ``bayang-bayang'' didukung oleh
  gradien amenitas menurut peringkat dan perbedaan respons terhadap
  massa eksternal, tetapi model tidak menunjukkan bahwa amenitas
  tertentu berpindah dari satu tempat ke tempat lain. Mekanisme
  persaingan konsumen tetap merupakan interpretasi penulis yang
  konsisten dengan pola, bukan aliran yang diukur langsung.
\item
  \textbf{Inferensi review:} klaim bahwa tempat kecil dan menengah
  ``berkembang ketika terisolasi'' lebih luas daripada keluaran yang
  diukur. Variabel hasil hanya menangkap amenitas budaya kelas atas,
  bukan pertumbuhan ekonomi, kesejahteraan, atau keseluruhan kemampuan
  tempat untuk berkembang.
\end{itemize}

\subsection{Conclusions}\label{conclusions-18}

\textbf{Kesimpulan penulis.} A03 menempatkan \emph{borrowed size} dan
\emph{agglomeration shadow} sebagai konsep berlawanan untuk menerangkan
tempat yang masing-masing menampung lebih banyak atau lebih sedikit
fungsi daripada yang lazim ditopang ukurannya. Dengan amenitas budaya
kelas atas sebagai ukuran fungsi, jumlah penduduk lokal menjadi faktor
penjelas utama, sedangkan ukuran FUA dan aksesibilitas
nasional-internasional secara umum lebih kecil perannya (Bagian 4, hlm.
1103-1104).

\textbf{Kesimpulan penulis.} Posisi dalam FUA mengubah hasil tersebut.
Amenitas terkonsentrasi secara tidak proporsional di tempat terbesar;
tempat itu lebih mampu menggunakan massanya sendiri, penduduk sisa FUA,
dan aksesibilitas nasional-internasional. Tempat berperingkat lebih
rendah lebih kecil kemungkinannya meminjam ukuran dan lebih besar
kemungkinannya menghadapi persaingan serta bayang-bayang pusat peringkat
pertama (Bagian 4, hlm. 1103-1104).

\textbf{Kesimpulan penulis.} Untuk amenitas kelas atas, pola rata-rata
Eropa Barat Laut lebih mendukung konsentrasi pusat daripada pembagian
fungsi yang menguntungkan kota kecil. Namun, penulis membatasi
kesimpulan itu dengan menyatakan bahwa amenitas budaya berorde lebih
rendah, acara berskala kecil, atau fungsi perkotaan lain dapat mengikuti
geografi berbeda (Bagian 4, hlm. 1104).

\textbf{Inferensi review.} Kesimpulan yang paling langsung ditopang
analisis adalah asosiasi antara posisi peringkat pertama dan kelebihan
amenitas budaya kelas atas. Generalisasi ke seluruh fungsi kota, semua
bentuk \emph{borrowed size}, atau dinamika pertumbuhan perlu ditahan
karena desainnya potong lintang, proksinya khusus, dan interaksi
konsumen antartempat tidak diamati.

\subsection{Contributions}\label{contributions-18}

\begin{itemize}
\tightlist
\item
  \textbf{Kontribusi yang dinyatakan penulis:} menyediakan salah satu
  eksplorasi empiris awal atas konsep \emph{borrowed size}, yang
  sebelumnya dinilai menjanjikan untuk menjelaskan urbanisasi Eropa
  Barat Laut tetapi hanya mendapat perhatian terbatas dalam ekonomi dan
  geografi (Bagian 1.3, hlm. 1092; Bagian 4, hlm. 1103).
\item
  \textbf{Kontribusi yang dinyatakan penulis:} mengonseptualisasikan
  \emph{agglomeration shadow} sebagai antonim \emph{borrowed size} dan
  menempatkan keduanya sebagai dua sisi dari ketidakselarasan
  ukuran-fungsi akibat interdependensi spasial (Bagian 2, hlm.
  1093-1094; Bagian 4, hlm. 1103).
\item
  \textbf{Kontribusi yang dinyatakan penulis:} menunjukkan bahwa posisi
  relatif dalam FUA membantu menjelaskan tempat mana yang dapat
  menangkap massa regional dan nasional-internasional, bukan hanya
  apakah massa eksternal tersedia (Abstract, hlm. 1090; Bagian 3.1, hlm.
  1101-1103).
\item
  \textbf{Kontribusi empiris sumber:} menyediakan perbandingan lintas
  2.794 tempat dan 417 FUA di tujuh negara dengan satu indeks amenitas
  budaya kelas atas yang sama (Bagian 2.1, hlm. 1094-1096).
\item
  \textbf{Inferensi review:} kontribusi analitis terkuat ialah
  memperlakukan hierarki sebagai moderator ekonomi jaringan. Pendekatan
  ini mencegah efek positif bagi pusat dan efek nol atau negatif bagi
  tempat bawahan saling meniadakan dalam satu koefisien rata-rata.
\item
  \textbf{Inferensi review:} pemakaian model beta terinflasi nol dan
  satu menyesuaikan bentuk distribusi indeks yang sangat jarang dan
  terbatas. File menyajikannya sebagai pilihan metodologis yang sesuai,
  tetapi tidak mengklaim bahwa model tersebut merupakan inovasi
  statistik baru (Bagian 2.2, hlm. 1097).
\end{itemize}

\subsection{Research Gap}\label{research-gap-18}

\textbf{Kesenjangan yang memotivasi studi.}

\begin{itemize}
\tightlist
\item
  Konsep \emph{borrowed size} telah lama diajukan sebagai kunci
  pemahaman pola urbanisasi Jerman dan Negara-Negara Rendah, tetapi
  menurut penulis hanya memperoleh perhatian terbatas dalam ekonomi dan
  geografi (Bagian 1.3, hlm. 1091-1092).
\item
  Literatur jaringan kota menjelaskan bahwa massa kritis dapat berasal
  dari luar tempat, sedangkan teori pusat pelayanan dan \emph{New
  Economic Geography} menjelaskan persaingan serta bayang-bayang pusat
  besar. Sebelum A03, masih belum jelas tempat mana yang memperoleh
  manfaat dan tempat mana yang kehilangan fungsi dalam interdependensi
  tersebut (Bagian 1.3-1.4, hlm. 1091-1093).
\item
  Hubungan ukuran-fungsi perlu dijelaskan bersama posisi regional,
  nasional, dan internasional. Artikel mengisi celah ini dengan menguji
  peringkat FUA sebagai pemoderasi tiga ukuran potensi penduduk (Bagian
  1.4, hlm. 1093; Bagian 2.1, hlm. 1096-1097).
\end{itemize}

\textbf{Kesenjangan yang masih tersisa menurut sumber.} Penelitian belum
mencakup amenitas berorde lebih rendah, layanan konsumen selain budaya
kelas atas, fungsi perkotaan yang lebih luas, komplementaritas fungsi
antartempat, ukuran posisi jaringan yang lebih terperinci, atau seluruh
faktor yang memungkinkan suatu tempat mengambil massa kritis dari
wilayah lain (Bagian 4, hlm. 1104).

\textbf{Inferensi review.} A03 mengisi celah asosiasi
ukuran-posisi-fungsi, tetapi belum mengisi celah mekanisme. Data tidak
mengidentifikasi asal pengunjung, arus belanja budaya, substitusi
antaramenitas, perubahan sebelum-sesudah, atau pasangan pemberi-penerima
manfaat. Dengan demikian, siapa yang benar-benar ``meminjamkan'' ukuran
kepada siapa masih belum terukur langsung.

\subsection{Limitations}\label{limitations-18}

\textbf{Keterbatasan yang dinyatakan atau diakui sumber.}

\begin{itemize}
\tightlist
\item
  Fokus penelitian terutama pada amenitas budaya kelas atas. Penulis
  mengakui bahwa kegiatan budaya lokal atau regional, acara berskala
  kecil, layanan konsumen lain, dan fungsi perkotaan lain dapat memiliki
  pola lokasi berbeda (Bagian 4, hlm. 1104).
\item
  Ukuran PDB per kapita berada pada tingkat FUA, sedangkan indeks
  amenitas berada pada tingkat tempat. Akibatnya, variabel PDB tidak
  menangkap heterogenitas pendapatan yang dapat kuat antara pusat dan
  kawasan suburban yang berdekatan (Catatan 8, hlm. 1105).
\item
  Ukuran aksesibilitas multimoda awalnya dikembangkan untuk wilayah
  NUTS-3 dan diadaptasi ke FUA. Prosedur adaptasi secara terperinci:
  Tidak disebutkan dalam file (Catatan 7, hlm. 1105).
\item
  Gambar interaksi penduduk lokal tidak menampilkan ukuran di atas
  sekitar 200.000 bagi tempat yang bukan terbesar karena observasi
  terbesar dalam kelompok tersebut adalah Oberhausen dengan sedikit di
  atas 200.000 penduduk (Catatan 11, hlm. 1105).
\item
  Analisis grafis potensi regional dibatasi hingga satu juta penduduk di
  sisa FUA karena kurang dari 0,5\% tempat terbesar dalam FUA memiliki
  hinterland di atas ambang itu (Catatan 12, hlm. 1105).
\item
  Sebagian hasil estimasi ketaklinearan ukuran lokal dinyatakan tersedia
  atas permintaan dan tidak tercantum dalam file (Catatan 9, hlm. 1105).
\end{itemize}

\textbf{Batas yang diinferensikan review dari desain dan pelaporan.}

\begin{itemize}
\tightlist
\item
  Desain potong lintang tidak menetapkan urutan waktu dan tidak
  memisahkan kemungkinan bahwa amenitas menarik penduduk dari
  kemungkinan bahwa penduduk menopang amenitas. Hasil tidak dapat dengan
  sendirinya dibaca sebagai efek kausal ukuran atau aksesibilitas.
\item
  \emph{Borrowed size} dan bayang-bayang aglomerasi tidak diukur melalui
  aliran konsumen atau pertukaran fungsi. Penetapannya bergantung pada
  proksi amenitas, peringkat FUA, dan asosiasi model.
\item
  Sebanyak 78\% observasi bernilai nol. Model secara khusus menangani
  distribusi ini, tetapi informasi empiris tentang variasi positif
  berasal dari minoritas tempat yang memiliki amenitas tercakup (Bagian
  2.1-2.2, hlm. 1094 dan 1097).
\item
  Indeks menggabungkan enam bentuk amenitas yang berbeda. File
  menyatakan bahwa skor rata-rata enam dimensi digunakan, tetapi
  prosedur standardisasi setiap dimensi sebelum perataan, sumber rinci
  setiap komponen, penanganan data hilang, dan uji kepekaan konstruksi
  indeks: Tidak disebutkan dalam file; rincian dialihkan ke BBSR (2011)
  (Bagian 2.1, hlm. 1094).
\item
  Ukuran aksesibilitas merepresentasikan potensi penduduk berbobot
  waktu, bukan penggunaan aktual kereta, jalan, atau udara untuk
  mengunjungi amenitas. Kualitas layanan, harga, preferensi, bahasa, dan
  hambatan perbatasan tidak dimodelkan secara eksplisit.
\item
  Peringkat ukuran dalam FUA merupakan ukuran posisi yang kasar. Artikel
  sendiri menyebut perlunya ukuran jaringan dan posisi yang lebih rinci,
  termasuk posisi dalam wilayah atau negara yang lebih luas (Bagian 4,
  hlm. 1104).
\item
  Sampel sangat terkonsentrasi di Prancis dan Jerman, yang masing-masing
  mencakup 35\% dan 29\% observasi. Validitas eksternal untuk bagian
  Eropa lain atau sistem perkotaan di luar Eropa: Tidak diuji dalam file
  (Bagian 2.1, hlm. 1094).
\item
  Model memasukkan indikator negara dan galat baku robust, tetapi cara
  menangani ketergantungan observasi yang tersarang dalam 417 FUA,
  misalnya pengelompokan galat baku atau efek FUA, tidak dijelaskan
  dalam file (Tabel 3, hlm. 1098-1099; Lampiran 1, hlm. 1109).
\item
  Uji validasi eksternal terhadap indeks, pemeriksaan spesifikasi
  alternatif selain keluarga model yang dibahas, ukuran kecocokan
  prediksi, dan analisis sensitivitas terhadap definisi FUA: Tidak
  disebutkan dalam file.
\item
  Data fungsi budaya merujuk 2006-2007, sehingga hasil tidak menunjukkan
  perubahan dinamis \emph{borrowed size} atau bayang-bayang aglomerasi
  dari waktu ke waktu (Bagian 2.1, hlm. 1094 dan 1096).
\end{itemize}

\subsection{Challenges}\label{challenges-18}

\textbf{Tantangan yang dinyatakan atau ditangani sumber.} Artikel tidak
memiliki bagian khusus bernama \emph{Challenges}, tetapi
mengidentifikasi beberapa persoalan analitis berikut.

\begin{itemize}
\tightlist
\item
  \textbf{Laporan sumber:} efek jaringan dapat berlawanan: satu tempat
  memanfaatkan permintaan eksternal, sedangkan tempat lain kehilangan
  permintaan karena pasokan eksternal. Penulis memperkirakan koefisien
  rata-rata massa eksternal dapat menjadi tidak signifikan apabila kedua
  proses saling meniadakan, lalu menggunakan posisi FUA dan interaksi
  untuk membedakannya (Bagian 1.4, hlm. 1093; Bagian 2.1, hlm.
  1096-1097).
\item
  \textbf{Laporan sumber:} variabel hasil dibatasi pada 0-1 dan sangat
  terkonsentrasi pada nol. Tantangan distribusi ini ditangani dengan
  regresi beta terinflasi nol dan satu serta parameter overdispersi
  (Bagian 2.2, hlm. 1097).
\item
  \textbf{Laporan sumber:} efek interaksi dalam model nonlinier sulit
  ditafsirkan. Penulis karena itu menyajikan AME secara grafis pada
  Gambar 4-6 dan menempatkan koefisien regresi di Lampiran 1 (Bagian
  3.1, hlm. 1101; Lampiran 1, hlm. 1109).
\item
  \textbf{Laporan sumber:} skala geografis proses dapat meliputi FUA,
  negara, dan jaringan internasional. Penelitian harus membedakan massa
  lokal, penduduk regional di luar tempat, dan aksesibilitas di luar FUA
  (Bagian 2, hlm. 1093-1094; Bagian 2.1, hlm. 1094-1096).
\item
  \textbf{Interpretasi penulis:} konsentrasi ekonomi aglomerasi dalam
  konsumsi menantang kemampuan kota kecil dan menengah untuk bersaing
  dengan kota besar melalui amenitas budaya kelas atas (Bagian 4, hlm.
  1104).
\item
  \textbf{Inferensi review:} tantangan pengukuran utamanya ialah
  membedakan fungsi yang benar-benar ditopang permintaan eksternal dari
  fungsi yang sekadar berkorelasi dengan status pusat. Tanpa data asal
  pengunjung atau perubahan temporal, model tidak dapat mengamati proses
  peminjaman secara langsung.
\item
  \textbf{Inferensi review:} tantangan interpretasi muncul karena
  ``lebih banyak daripada yang diharapkan berdasarkan ukuran''
  didefinisikan secara konseptual, tetapi artikel tidak menerbitkan
  klasifikasi residual setiap tempat sebagai peminjam atau sebagai
  tempat yang terbayangi.
\end{itemize}

\subsection{Practical Implications}\label{practical-implications-18}

\textbf{Implikasi praktis yang dinyatakan sebagai rekomendasi kebijakan
formal:} Tidak disebutkan dalam file. Artikel lebih banyak membahas
makna hasil bagi persaingan dan pembagian fungsi perkotaan daripada
merumuskan instrumen kebijakan tertentu.

\begin{itemize}
\tightlist
\item
  \textbf{Interpretasi penulis:} kota kecil dan menengah lebih kecil
  kemungkinannya meminjam ukuran dari tetangga besar untuk menopang
  amenitas budaya kelas atas. Dalam rata-rata sampel, kedekatan dalam
  FUA lebih sering menguntungkan pusat terbesar dan menempatkan pusat
  bawahan dalam bayang-bayang aglomerasi (Bagian 4, hlm. 1104).
\item
  \textbf{Inferensi review berdasarkan caveat penulis:} pola untuk
  kegiatan budaya lokal, regional, atau berskala kecil dapat berbeda.
  Contoh penulis tentang musisi yang lebih berorientasi lokal dan acara
  budaya lebih kecil menunjukkan kemungkinan strategi amenitas yang
  tidak dibatasi pada fungsi kelas atas, tetapi strategi tersebut tidak
  diuji dalam A03 (Bagian 4, hlm. 1104).
\item
  \textbf{Inferensi review:} perencanaan budaya metropolitan perlu
  memperhitungkan hierarki internal FUA. Menambah aksesibilitas atau
  basis penduduk regional saja tidak menjamin penyebaran amenitas karena
  pusat peringkat pertama lebih mampu menangkap potensi tersebut.
\item
  \textbf{Inferensi review:} evaluasi struktur polisentris sebaiknya
  tidak menyamakan keberadaan banyak pusat dengan pembagian fungsi yang
  seimbang. A03 menunjukkan bahwa integrasi regional dapat berdampingan
  dengan konsentrasi amenitas pada satu pusat.
\item
  \textbf{Inferensi review:} karena efek aksesibilitas
  nasional-internasional dilaporkan kecil dan terutama menguntungkan
  pusat terbesar, investasi konektivitas tidak boleh dianggap otomatis
  memperkuat amenitas semua kota dalam jaringan. Artikel tidak menguji
  intervensi transportasi atau kebijakan redistribusi tertentu.
\item
  \textbf{Inferensi review:} pembagian fungsi yang saling melengkapi
  dapat menjadi alternatif terhadap persaingan meniru amenitas yang
  sama, tetapi artikel mengajukannya sebagai kemungkinan untuk
  penelitian mendatang, bukan sebagai kebijakan yang telah diuji (Bagian
  4, hlm. 1104).
\end{itemize}

\subsection{Applications}\label{applications-18}

\begin{itemize}
\tightlist
\item
  \textbf{Aplikasi yang diperagakan sumber:} menggunakan hubungan
  ukuran-fungsi dan posisi dalam jaringan regional untuk memeriksa
  \emph{borrowed size} dan \emph{agglomeration shadow} pada amenitas
  budaya kelas atas (Bagian 2, hlm. 1093-1094; Bagian 2.1, hlm.
  1096-1097).
\item
  \textbf{Aplikasi yang diperagakan sumber:} membandingkan kemampuan
  tempat berperingkat pertama dan tempat bawahan dalam memanfaatkan
  massa lokal, penduduk sisa FUA, serta aksesibilitas
  nasional-internasional (Bagian 2.1, hlm. 1096-1097; Gambar 4-6, hlm.
  1101-1103).
\item
  \textbf{Aplikasi yang diperagakan sumber:} memodelkan indeks fungsi
  perkotaan yang terbatas dan mengandung banyak nilai nol melalui
  regresi beta terinflasi nol dan satu (Bagian 2.2, hlm. 1097).
\item
  \textbf{Inferensi review:} kerangka dapat digunakan sebagai alat
  diagnosis awal untuk menilai apakah suatu fungsi terkonsentrasi
  melebihi pola yang diperkirakan dari ukuran tempat dan apakah
  konsentrasi itu berbeda menurut hierarki regional. Aplikasi tersebut
  memerlukan indikator fungsi yang sebanding dan definisi FUA yang
  jelas.
\item
  \textbf{Inferensi review:} penerapan pada layanan kesehatan,
  pendidikan, olahraga, ritel, atau fungsi ekonomi lain belum diuji
  dalam A03. Artikel sendiri memperingatkan bahwa satu tempat dapat
  menghadapi bayang-bayang pada satu fungsi tetapi meminjam ukuran pada
  fungsi lain (Bagian 4, hlm. 1104).
\item
  \textbf{Batas aplikasi:} perangkat lunak, kode estimasi, protokol
  klasifikasi tempat, ambang keputusan perencanaan, dan studi
  implementasi kebijakan: Tidak disebutkan dalam file.
\end{itemize}

\subsection{Future Research
(Optional)}\label{future-research-optional-18}

Sumber menyebut agenda penelitian berikut secara eksplisit pada Bagian
4, hlm. 1104.

\begin{enumerate}
\def\labelenumi{\arabic{enumi}.}
\tightlist
\item
  Memasukkan jenis layanan konsumen lain, khususnya amenitas budaya
  berorde lebih rendah dan acara berskala lebih kecil yang mungkin
  mengikuti pola lokasi berbeda dari amenitas kelas atas.
\item
  Memperluas variabel hasil ke ragam fungsi perkotaan yang jauh lebih
  luas, bukan hanya amenitas budaya kelas atas.
\item
  Menguji komplementaritas fungsi antartempat dalam wilayah metropolitan
  sebagai kemungkinan kunci untuk memahami \emph{borrowed size} dan
  bayang-bayang aglomerasi.
\item
  Memeriksa kemungkinan bahwa tempat yang sama mengalami bayang-bayang
  untuk satu fungsi tetapi meminjam ukuran untuk fungsi lain.
\item
  Menggunakan ukuran posisi tempat yang lebih terperinci dalam wilayah
  dan jaringan, termasuk apakah suatu tempat merupakan yang terbesar
  dalam wilayah atau negara yang lebih luas.
\item
  Menganalisis lebih banyak faktor yang menentukan kemampuan suatu
  tempat memperoleh massa kritis dari wilayah lain, termasuk kualitas
  infrastruktur, kemudahan akses, kualitas umum layanan, serta
  keseluruhan tingkat permintaan dan penawaran.
\end{enumerate}

Agenda penelitian lain di luar yang dinyatakan penulis: Tidak
ditambahkan oleh review ini.

\subsection{Provenance Notes}\label{provenance-notes-18}

\begin{itemize}
\tightlist
\item
  Review ini hanya menggunakan seluruh isi file
  \texttt{/Users/mac/Documents/Mac/{[}2{]}\ Obsidian\ Vault/zettelkasten/source/journal/A03-burger-meijers-hoogerbrugge-masip-tresserra-2015-borrowed-size-agglomeration-shadows-and-cultural-amenities-in-north-west-europe-10.1080-09654313.2014.905002.txt},
  sebanyak 894 baris. Tidak ada situs web, basis data bibliografis, PDF
  lain, atau sumber eksternal yang digunakan untuk menambah, mengoreksi,
  atau memverifikasi klaim substantif.
\item
  Blok metadata menyatakan bahwa TXT diekstrak pada 31 Agustus 2026
  dengan \texttt{pdftotext\ -layout} dari
  \texttt{journal/A03-burger-meijers-hoogerbrugge-masip-tresserra-2015-borrowed-size-agglomeration-shadows-and-cultural-amenities-in-north-west-europe-10.1080-09654313.2014.905002.pdf}.
  PDF memiliki 21 halaman, tidak terenkripsi, dan tidak bertanda
  struktur (\texttt{Tagged:\ no}). Review membaca TXT, bukan memeriksa
  PDF secara visual.
\item
  Locator halaman mengikuti nomor halaman jurnal yang masih tampak dalam
  ekstraksi, yaitu 1090-1109. Nomor bagian, tabel, gambar, catatan, dan
  lampiran mengikuti label yang tercetak dalam artikel.
\item
  Isi visual Gambar 1-6 tidak tersedia sebagai data grafis yang dapat
  diperiksa dari TXT. Review hanya menggunakan caption dan uraian
  naratif penulis. Karena itu, bentuk kurva, rentang kepercayaan,
  geografi peta, dan detail visual lain tidak dinilai secara independen.
\item
  Tata letak Tabel 3 dan Lampiran 1 terdistorsi dalam ekstraksi.
  Sejumlah tanda negatif muncul sebagai karakter \texttt{2}, dan
  beberapa baris interaksi tidak sejajar dengan jelas. Review tidak
  merekonstruksi seluruh koefisien dari karakter yang rusak; angka AME
  dan arah hubungan diambil ketika didukung oleh uraian naratif atau
  bagian tabel yang terbaca jelas.
\item
  Tabel 1 dan Tabel 2 terbaca sebagai teks, tetapi review tidak
  menghitung ulang rerata, simpangan baku, persentase, AME, atau
  perbedaan prediksi. Seluruh angka dilaporkan sebagaimana tertulis.
\item
  File memakai istilah ``UNICEF heritage site'' untuk variabel kontrol.
  Karena sumber eksternal dilarang dan tidak ada klarifikasi lain dalam
  TXT, review mempertahankan istilah itu tanpa menggantinya dengan nama
  lembaga lain (Bagian 2.1, hlm. 1096).
\item
  Judul dan pembahasan memakai ``North-West Europe'', sedangkan uraian
  data menyebut tempat-tempat ``West-European''. Sampel yang dirinci
  mencakup Prancis, Jerman, Swiss, Belanda, Austria, Belgia, dan
  Luksemburg. Review mempertahankan variasi istilah geografis tersebut
  dan tidak menetapkan batas baru.
\item
  Uraian FUA menyebut Milano sebagai contoh meskipun Italia tidak
  tercantum dalam distribusi negara sampel. File tidak menjelaskan
  apakah contoh itu merujuk basis ESPON yang lebih luas atau sampel
  analitis; ketidakjelasan ini tidak diselesaikan melalui dugaan.
\item
  Header jurnal menetapkan tahun publikasi 2015, sedangkan metadata PDF
  pada kolom Subject menyebut \emph{European Planning Studies, 2014} dan
  halaman hak cipta menyebut 2014. Review menggunakan 2015 sebagai tahun
  edisi berdasarkan header bibliografis, sambil mempertahankan
  ketidakselarasan ini sebagai batas metadata.
\item
  Status telaah sejawat tidak dinyatakan secara formal. Ucapan terima
  kasih kepada dua penelaah anonim dilaporkan sebagai isi file, bukan
  dianggap verifikasi independen atas status telaah.
\item
  Rujukan seperti Alonso (1973), BBSR (2011), Peeters (2011),
  Spiekermann dan Wegener (2006), serta karya lain hanya diketahui
  melalui sitasi dan daftar pustaka A03. Isi rujukan tersebut tidak
  dibaca atau diverifikasi; setiap penyebutannya diperlakukan sebagai
  laporan A03 atas literatur.
\item
  Batas provenance utama: review dapat memeriksa apa yang tertulis dalam
  ekstraksi teks lengkap, tetapi tidak dapat memastikan kesetiaan tata
  letak terhadap PDF, membaca detail grafis, memeriksa data mentah,
  mereplikasi model, atau memverifikasi klaim kausal dan bibliografis
  secara independen.
\end{itemize}

\section{Review A04}\label{review-a04}

\textbf{Identitas sumber}

\begin{itemize}
\tightlist
\item
  Kode: A04.
\item
  Judul: \emph{The Geography of Borrowing Size: Exploring Spatial
  Distributions for German Urban Regions}.
\item
  Penulis: Kati Volgmann dan Karsten Rusche.
\item
  Afiliasi: Research Institute for Regional and Urban Development GmbH
  (ILS), Dortmund, Jerman.
\item
  Jenis sumber menurut isi: artikel jurnal. Status telaah sejawat: Tidak
  disebutkan dalam file.
\item
  Jurnal: \emph{Tijdschrift voor Economische en Sociale Geografie}.
\item
  Keterangan bibliografis yang tercetak: 2019, Volume 0, Nomor 0, hlm.
  1-20. Nama file menyebut tahun 2020; ketidakselarasan ini tidak
  diselesaikan dengan sumber eksternal.
\item
  DOI: 10.1111/tesg.12362.
\item
  Riwayat naskah: diterima Juni 2018 dan disetujui Januari 2019.
\item
  Kata kunci: \emph{borrowed size}; \emph{borrowed function};
  \emph{borrowed performance}; pertumbuhan perkotaan; fungsi
  metropolitan; kawasan perkotaan.
\item
  Wilayah dan unit studi: munisipalitas yang memiliki fungsi
  metropolitan di dalam 20 kawasan perkotaan fungsional Jerman, dengan
  analisis utama terhadap 60 munisipalitas (bagian Research Design, hlm.
  5-10; Table 3-4, hlm. 15-16).
\item
  Panjang sumber: 20 halaman berdasarkan metadata PDF dan penomoran
  halaman artikel.
\end{itemize}

\textbf{Penanda epistemik yang digunakan}

\begin{itemize}
\tightlist
\item
  \textbf{Laporan sumber} berarti tujuan, konsep, data, prosedur, angka,
  atau hasil yang disampaikan langsung oleh artikel. Laporan artikel
  mengenai karya terdahulu tetap diperlakukan sebagai laporan artikel,
  bukan verifikasi independen atas karya tersebut.
\item
  \textbf{Interpretasi penulis} berarti penjelasan mekanisme, makna,
  implikasi, atau generalisasi yang diberikan Volgmann dan Rusche
  terhadap hasil mereka.
\item
  \textbf{Inferensi review} berarti pembacaan analitis yang dibuat hanya
  dari TXT A04. Inferensi ini tidak diperlakukan sebagai temuan empiris
  tambahan.
\end{itemize}

\subsection{Summarized Abstract}\label{summarized-abstract-19}

\textbf{Laporan sumber.} Artikel ini mengoperasionalkan konsep
\emph{borrowed size} yang diperluas menjadi empat tipe: \emph{borrowed
size}, \emph{borrowed performance}, \emph{borrowed function}, dan
\emph{agglomeration shadow}. Kasus empirisnya adalah sistem perkotaan
Jerman yang polisentris. Pada skala munisipalitas, penulis menyusun
indeks fungsi metropolitan dari 47 indikator, mengukur kinerja relatif
melalui perubahan pangsa pekerjaan 2008-2015, lalu menggunakan residual
dua regresi OLS untuk menempatkan 60 munisipalitas ke dalam empat tipe
tersebut (Abstract, hlm. 1; bagian Operationalisation, hlm. 6-10; Table
2, hlm. 10).

\textbf{Laporan sumber.} Klasifikasi menghasilkan 20 munisipalitas dalam
\emph{agglomeration shadow}, 7 dalam \emph{borrowed function}, 16 dalam
\emph{borrowed performance}, dan 17 dalam \emph{borrowed size}.
Distribusinya berbeda menurut tipe kota dan struktur kawasan perkotaan.
Dari lima kota besar, Munich diklasifikasikan sebagai \emph{borrowed
size}, Berlin sebagai \emph{borrowed performance}, serta Hamburg,
Cologne, dan Frankfurt am Main sebagai \emph{borrowed function}. Enam
belas dari 17 kasus \emph{borrowed size} berada pada kota kecil,
sedangkan 15 dari 20 kasus \emph{agglomeration shadow} berada pada kota
menengah (Table 2-3, hlm. 10 dan 15; uraian hasil, hlm. 13-16).

\textbf{Interpretasi penulis.} Penulis menafsirkan perbedaan geografi
tersebut sebagai bukti bahwa ukuran kota, derajat sentralitas, dan
konteks spasial kawasan perkotaan memengaruhi kemampuan kota untuk
memperoleh manfaat kinerja atau fungsi dari kota lain maupun untuk
berada dalam bayang-bayang aglomerasi. Mereka juga menilai bahwa ekonomi
aglomerasi tidak berhenti pada batas fisik kota, tetapi dapat melimpah
ke kota-kota di sekitarnya (Abstract, hlm. 1; Conclusion, hlm. 17).

\textbf{Inferensi review.} Bukti langsung artikel adalah penyimpangan
fungsi dan kinerja dari nilai yang diprediksi berdasarkan ukuran dan
kepadatan, bukan pengamatan langsung atas proses ``meminjam'' fungsi,
arus manfaat, atau hubungan kausal antarkota. Oleh karena itu, empat
label paling aman dibaca sebagai tipologi empiris berbasis residual yang
konsisten dengan konsep \emph{borrowed size}, bukan sebagai identifikasi
kausal atas mekanisme peminjaman (persamaan 1-2 dan Table 2, hlm. 9-10).

\subsection{Problem Statement}\label{problem-statement-19}

\textbf{Laporan sumber.} Artikel berangkat dari dua gejala perkembangan
perkotaan Eropa. Pertama, kemampuan menarik fungsi tingkat tinggi atau
fungsi metropolitan dan perkembangan ekonomi tidak selalu sejalan dengan
ukuran kota, sehingga hubungan antara skala atau kepadatan perkotaan dan
produktivitas dinilai tidak jelas. Kedua, pola pertumbuhan kawasan
perkotaan Eropa sangat heterogen karena ketergantungan lintasan,
perbedaan endowment faktor, dan saling ketergantungan antardaerah
(Introduction, hlm. 1).

\textbf{Laporan sumber.} Konsep \emph{borrowed size} menawarkan
penjelasan bahwa kota kecil dan menengah dalam kawasan perkotaan
fungsional yang lebih besar dapat memanfaatkan fungsi metropolitan,
jaringan, dan eksternalitas ekonomi kota lain. Sebaliknya, kota yang
gagal memperoleh manfaat dari tetangga yang lebih besar dapat mengalami
\emph{backwash} atau \emph{agglomeration shadow}. Namun, temuan
terdahulu berbeda-beda mengenai lokasi, kawasan, dan tipe kota tempat
efek tersebut muncul (Introduction, hlm. 1-2; Theoretical Foundation,
hlm. 2-4).

\textbf{Laporan sumber.} Celah utama yang dinyatakan penulis adalah
kurangnya operasionalisasi empiris yang dapat menganalisis dan
mengklasifikasikan distribusi spasial efek \emph{borrowed size} pada
skala kecil dengan kumpulan indikator yang kuat dan metode yang
konsisten. Kerangka empat dimensi Meijers dan Burger (2017) tersedia
secara konseptual, tetapi masih perlu diterjemahkan menjadi indikator
fungsi dan kinerja yang terukur (Introduction, hlm. 2; Research gap,
hlm. 5).

\textbf{Inferensi review.} Masalah pengukuran artikel terdiri atas dua
keputusan: menentukan apa yang dimaksud dengan fungsi metropolitan dan
kinerja kota, lalu menentukan kapan nilai keduanya lebih tinggi atau
lebih rendah daripada yang ``diharapkan berdasarkan ukuran''. Residual
regresi menjadi penghubung antara konsep teoretis dan klasifikasi
empiris tersebut (bagian Operationalisation, hlm. 6-10).

\subsection{Objectives}\label{objectives-19}

\textbf{Laporan sumber.} Artikel mengajukan tiga pertanyaan penelitian:

\begin{enumerate}
\def\labelenumi{\arabic{enumi}.}
\tightlist
\item
  Bagaimana kategori \emph{borrowed size} dapat diukur secara statistik?
  Tujuan pertama ialah mengeksplorasi dan mengoperasionalkan hubungan
  antara \emph{borrowed size}, \emph{borrowed performance},
  \emph{borrowed function}, dan \emph{agglomeration shadow} melalui
  indikator yang dapat mengategorikan kota-kota Jerman (Research gap,
  hlm. 5).
\item
  Sejauh mana kota-kota dalam kawasan perkotaan Jerman dipengaruhi efek
  spasial \emph{borrowed size}? Analisis dilakukan pada skala kecil
  dengan memasukkan kota besar, menengah, dan kecil untuk memetakan
  variasi keempat tipe (Research gap, hlm. 5).
\item
  Apakah ukuran kota dan struktur spasial memberikan efek yang berbeda
  terhadap \emph{borrowing function} dan \emph{borrowing performance}?
  Penulis secara khusus ingin membandingkan perilaku kota besar dengan
  kota menengah dan kecil pada struktur kawasan yang berbeda (Research
  gap, hlm. 5).
\end{enumerate}

\textbf{Laporan sumber.} Tujuan analitis turunannya adalah menggambarkan
distribusi fungsi metropolitan dan kinerja ekonomi, mengklasifikasikan
kota berdasarkan kombinasi residual kedua dimensi, serta membandingkan
klasifikasi menurut tipe kota dan struktur kawasan perkotaan monosentris
atau polisentris (Analysis and Empirical Findings, hlm. 10-16).

\subsection{Literature Survey}\label{literature-survey-19}

\textbf{Sifat survei.} Artikel menyajikan tinjauan naratif untuk
membangun konsep dan operasionalisasi. Strategi pencarian, pangkalan
data bibliografis, kriteria inklusi, jumlah karya yang disaring, dan
prosedur penilaian mutu: Tidak disebutkan dalam file. Bagian ini tidak
dapat diperlakukan sebagai tinjauan sistematis.

\begin{itemize}
\tightlist
\item
  \textbf{Laporan sumber atas asal konsep:} Alonso (1973) mengemukakan
  bahwa kota kecil dapat menunjukkan karakteristik kota yang lebih besar
  jika dekat dengan konsentrasi penduduk lain. Phelps dan kolega
  kemudian menghubungkannya dengan pertumbuhan suburban, \emph{edge
  cities}, subpusat ekonomi baru, serta ketegangan antara aglomerasi dan
  desentralisasi (Theoretical Foundation, hlm. 2).
\item
  \textbf{Laporan sumber atas manfaat yang dipinjam:} penduduk dan
  perusahaan di permukiman kecil dapat mempertahankan keunggulan
  konsentrasi serta sewa yang rendah sambil mengakses pasar, layanan,
  infrastruktur, dan amenitas kota besar, serta menghindari sebagian
  biaya aglomerasi. Studi Greater London oleh Phelps et al.~(2001)
  dilaporkan menemukan perusahaan kecil di kota tepi memperoleh
  keuntungan dari kawasan perkotaan besar di dekatnya, bukan hanya dari
  ekonomi lokal (Theoretical Foundation, hlm. 2).
\item
  \textbf{Laporan sumber atas efek negatif:} \emph{agglomeration shadow}
  merujuk pada keterbatasan pertumbuhan akibat kompetisi di dekat
  konsentrasi perusahaan. Gagasan ini disejajarkan dengan teori tempat
  sentral dan sistem perkotaan, yang memperkirakan kota kecil di dekat
  kota besar memiliki lebih sedikit fungsi sentral daripada kota
  terisolasi dengan ukuran serupa, serta dengan konsep \emph{spread} dan
  \emph{backwash} Myrdal (Theoretical Foundation, hlm. 2-3).
\item
  \textbf{Laporan sumber atas bukti yang tidak seragam:} Brezzi dan
  Veneri (2015), Burger et al.~(2015), serta Meijers (2008) dilaporkan
  menemukan kinerja negatif atau bayang-bayang fungsi di sekitar kota
  besar pada beberapa kawasan polisentris. Sebaliknya, penelitian lain
  melaporkan pertumbuhan pekerjaan atau penduduk yang positif bagi kota
  kecil dalam kawasan besar. Besar manfaat dikaitkan dengan fungsi
  tingkat tinggi, jaringan, infrastruktur, konektivitas, mobilitas,
  pendidikan, jarak, dan struktur demografi (Theoretical Foundation,
  hlm. 3).
\item
  \textbf{Laporan sumber atas kedekatan dan interaksi:} jarak perjalanan
  dan waktu tempuh digunakan dalam literatur sebagai proksi akses
  terhadap eksternalitas aglomerasi. Akan tetapi, artikel juga mengutip
  argumen bahwa hubungan kooperatif, pertukaran informasi, dan
  keterhubungan dalam jaringan nasional atau internasional lebih penting
  daripada kedekatan fisik semata; \emph{borrowed size} dinyatakan
  sebagai produk interaksi nyata (Theoretical Foundation, hlm. 3-4).
\item
  \textbf{Laporan sumber atas fungsi metropolitan:} fungsi tingkat
  tinggi memerlukan ambang massa perkotaan tertentu. Literatur
  mengelompokkannya ke dalam fungsi pengambilan keputusan dan kontrol,
  inovasi dan kompetisi, gerbang, serta simbolik. Fungsi ini dipandang
  sebagai ekspresi kekuasaan geografis sekaligus posisi hierarkis dalam
  sistem tempat sentral (Theoretical Foundation, hlm. 4).
\item
  \textbf{Laporan sumber atas kerangka empat dimensi:} Meijers dan
  Burger (2017) mendefinisikan \emph{borrowed performance} sebagai
  kinerja yang lebih baik daripada yang diharapkan berdasarkan ukuran
  dan \emph{borrowed function} sebagai jumlah fungsi yang lebih banyak
  daripada yang diharapkan. Ketika keduanya positif secara bersamaan,
  kategorinya adalah \emph{borrowed size}. Kombinasi kinerja dan fungsi
  yang sama-sama lebih rendah menjadi \emph{agglomeration shadow} dalam
  operasionalisasi artikel (Figure 1, hlm. 4; uraian, hlm. 5; Table 2,
  hlm. 10).
\item
  \textbf{Interpretasi penulis atas ekspektasi teoretis:} kota kecil
  diperkirakan lebih sering memperoleh \emph{borrowed performance},
  sedangkan kota besar lebih sering memperoleh \emph{borrowed function}
  dan dapat membayangi hinterland-nya. Kawasan polisentris diperkirakan
  lebih memungkinkan kota berbagi eksternalitas, meskipun kota dominan
  juga dapat membatasi fungsi kota sekitar (Research gap, hlm. 5).
\end{itemize}

\textbf{Inferensi review.} Tinjauan menunjukkan ketegangan yang tidak
sepenuhnya diselesaikan oleh ukuran kedekatan: teori yang dikutip
menekankan ``interaksi nyata'', tetapi rancangan empiris A04 terutama
memakai lokasi dalam kawasan, ukuran, kepadatan, fungsi, dan perubahan
pekerjaan. Arus interaksi antarkota tidak menjadi variabel klasifikasi
utama (Theoretical Foundation, hlm. 3-4; persamaan 1-2, hlm. 10).

\subsection{Method Used}\label{method-used-19}

\textbf{Laporan sumber.} Penelitian merupakan studi kuantitatif
eksploratif dan klasifikatif pada sistem perkotaan Jerman. Tahap
utamanya meliputi delineasi kawasan perkotaan fungsional, konstruksi dua
ukuran empiris, estimasi dua regresi OLS, klasifikasi berdasarkan tanda
residual, serta pemetaan dan tabulasi silang tipe hasil menurut ukuran
kota dan struktur kawasan (Research Design, hlm. 5-10; Analysis and
Empirical Findings, hlm. 10-16).

\textbf{Tahap 1: delineasi kawasan dan tipe kota.}

\begin{itemize}
\tightlist
\item
  Kota inti ditetapkan memiliki sedikitnya 200.000 penduduk dan 100.000
  pekerja atau pekerjaan yang tercakup asuransi sosial. Kriteria ini
  menghasilkan 32 kota inti (Research Design, hlm. 6).
\item
  Lima kota inti dengan lebih dari 600.000 penduduk, yaitu Berlin,
  Munich, Hamburg, Frankfurt, dan Cologne, diklasifikasikan sebagai kota
  besar. Dua puluh tujuh kota inti lainnya disebut kota menengah
  (Research Design, hlm. 6).
\item
  Hinterland fungsional ditentukan melalui analisis jaringan dari kota
  inti ke titik fokus setiap munisipalitas dengan ambang waktu tempuh
  mobil. Besar buffer mengikuti kurva gravitasi yang diturunkan dari
  sentralisasi pekerjaan; ambangnya berkisar dari 30 menit untuk Erfurt
  hingga 60 menit untuk Berlin (Research Design, hlm. 6; Figure 2, hlm.
  7).
\item
  Karena kawasan beberapa inti bertumpang tindih, seperti Duisburg,
  Essen, dan Dortmund, setiap munisipalitas ditempatkan ke inti terdekat
  berdasarkan waktu perjalanan untuk menghasilkan 20 kawasan yang tidak
  saling tumpang tindih. Kawasan dibedakan menjadi monosentris dan
  polisentris menurut jumlah kota inti (Research Design, hlm. 6).
\item
  Munisipalitas yang memiliki fungsi metropolitan tetapi bukan kota
  besar dan bukan inti kawasan diklasifikasikan sebagai kota kecil dalam
  analisis ini (Research Design, hlm. 6).
\end{itemize}

\textbf{Tahap 2: pengukuran fungsi metropolitan.}

\begin{itemize}
\tightlist
\item
  Penulis menggunakan 47 indikator atribut fungsional terlokalisasi
  untuk 2008-2010. Nilai mentah distandardisasi dengan skor-z, kemudian
  dihubungkan secara aditif menjadi satu indeks metropolitan statis
  (bagian Metropolitan functions, hlm. 6-8; Table 1, hlm. 9).
\item
  Karena analisis dilakukan pada tingkat munisipalitas, indeks tingkat
  county dipindahkan ke munisipalitas terbesar yang termasuk dalam
  county tersebut. Setiap munisipalitas perkotaan kemudian memiliki satu
  nilai indeks (bagian Metropolitan functions, hlm. 8).
\end{itemize}

\textbf{Tahap 3: pengukuran kinerja perkotaan.}

\begin{itemize}
\tightlist
\item
  Kinerja dibatasi pada dimensi ekonomi melalui pekerjaan yang tercakup
  asuransi sosial. Untuk 60 munisipalitas yang memiliki fungsi
  metropolitan, penulis menghitung pangsa pekerjaan setiap kota dari
  total pekerjaan seluruh kelompok, lalu menghitung perubahan pangsa
  antara 2008 dan 2015. Kenaikan pangsa diperlakukan sebagai kinerja
  relatif yang baik dan penurunan sebagai kinerja relatif yang buruk
  (bagian Urban performance, hlm. 8).
\item
  Penulis sengaja tidak memakai perubahan pekerjaan absolut karena kota
  dengan basis awal besar cenderung mencatat perubahan besar, dan tidak
  memakai persentase pertumbuhan murni karena perubahan absolut kecil
  pada kota berbasis kecil dapat menghasilkan persentase tinggi (bagian
  Urban performance, hlm. 8).
\end{itemize}

\textbf{Tahap 4: ekspektasi berdasarkan ukuran dan klasifikasi
residual.}

\begin{itemize}
\tightlist
\item
  Regresi pertama memodelkan indeks metropolitan 2008-2010 (\texttt{MF})
  dengan populasi (\texttt{POP}) dan kepadatan penduduk
  (\texttt{POP.DENSITY}) tahun 2008. Regresi kedua memodelkan perubahan
  pangsa pekerjaan 2008-2015 (\texttt{EMP.SHARE}) dengan stok pekerjaan
  (\texttt{EMP}) dan kepadatan pekerjaan (\texttt{EMP.DENSITY}) tahun
  2008. Kepadatan dihitung terhadap luas kawasan permukiman (bagian
  Metropolitan functions, hlm. 8; persamaan 1-2, hlm. 10).
\item
  Residual positif berarti fungsi atau kinerja lebih tinggi daripada
  prediksi model berdasarkan ukuran dan kepadatan; residual negatif
  berarti lebih rendah. Kombinasi dua tanda menghasilkan empat kelompok:
  negatif-negatif sebagai \emph{agglomeration shadow}, positif
  fungsi-negatif kinerja sebagai \emph{borrowed function}, negatif
  fungsi-positif kinerja sebagai \emph{borrowed performance}, dan
  positif-positif sebagai \emph{borrowed size} (Table 2, hlm. 10).
\item
  Daya jelas tersesuaikan yang dilaporkan adalah \texttt{R2\ =\ 0,86}
  untuk model indeks metropolitan dan \texttt{R2\ =\ 0,51} untuk model
  pangsa pekerjaan. Penulis menyebut model kedua lebih lemah tetapi
  masih dapat diterima (bagian Expectations ``given size'', hlm. 9-10).
\item
  Hasil residual ditampilkan dalam diagram pencar, dipetakan, lalu
  dibandingkan menurut tiga tipe kota serta struktur kawasan monosentris
  atau polisentris (Figure 5-6, hlm. 13-14; Table 3-4, hlm. 15-16).
\end{itemize}

\textbf{Inferensi review atas rancangan.} Metode mengukur deviasi
bersyarat dari dua regresi dan kemudian memberi label konseptual pada
kuadrannya. File tidak melaporkan desain kausal, pengamatan langsung
atas aliran fungsi, atau pengujian apakah satu kota benar-benar
memanfaatkan fungsi kota lain. File juga tidak menerangkan transformasi
variabel, koefisien, galat baku, nilai-p, pemeriksaan dependensi
spasial, atau uji ketahanan spesifikasi (Research Design, hlm. 5-10;
Analysis and Empirical Findings, hlm. 10-16).

\subsection{Dataset}\label{dataset-19}

\begin{itemize}
\tightlist
\item
  \textbf{Batas spasial dan data kawasan:} referensi delineasi berasal
  dari Siedentop dan Kaup (2017), sedangkan Figure 2 menyebut \emph{The
  Regional Database Germany} sebagai sumber data. Jenis jaringan jalan,
  tahun jaringan, perangkat lunak perutean, dan prosedur validasi waktu
  tempuh: Tidak disebutkan dalam file (Research Design dan Figure 2,
  hlm. 6-7).
\item
  \textbf{Fungsi metropolitan:} 47 indikator terlokalisasi untuk periode
  2008-2010, dirujuk kepada Volgmann (2014). Empat kelompok besarnya
  adalah fungsi kontrol, inovasi, gerbang, dan simbolik (bagian
  Metropolitan functions, hlm. 6-8; Table 1, hlm. 9).
\item
  \textbf{Fungsi kontrol:} indikator mencakup kedudukan dan omzet
  perusahaan besar, aset bank, pendapatan perusahaan asuransi, bursa
  efek, perusahaan ritel pangan, kementerian federal, pegawai pemerintah
  federal dan negara bagian, pengadilan, institusi Uni Eropa atau
  Perserikatan Bangsa-Bangsa, kedutaan dan konsulat, asosiasi,
  organisasi pembangunan, serta yayasan (Table 1, hlm. 9).
\item
  \textbf{Fungsi inovasi:} indikator mencakup lokasi perusahaan
  inovatif, insinyur, jasa berorientasi bisnis, pekerja bergelar
  akademik, paten perusahaan, lokasi riset DFG dan organisasi riset
  utama, universitas, perpustakaan ilmiah, paten lembaga riset, serta
  pekerja riset (Table 1, hlm. 9).
\item
  \textbf{Fungsi gerbang:} indikator mencakup aktivitas dan penumpang
  bandara internasional, stasiun ICE atau TGV, volume bongkar muat
  pelabuhan laut dan darat, omzet perusahaan logistik, angkutan udara,
  serta lokasi dan luas pameran dagang (Table 1, hlm. 9).
\item
  \textbf{Fungsi simbolik:} indikator mencakup penyiaran, studio film,
  penerbit, surat kabar nasional, domain internet, pekerjaan budaya,
  bangunan arsitek terkenal, gedung tinggi, kompetisi pembangunan
  perkotaan, penonton opera, teater, dan konser, kapasitas stadion sepak
  bola, serta tamu hotel (Table 1, hlm. 9).
\item
  \textbf{Pekerjaan:} jumlah pekerjaan atau pekerja dalam skema asuransi
  sosial diperoleh dari sumber data regional publik
  \texttt{regionalstatistik.de}. Stok 2008 dipakai sebagai prediktor
  ukuran ekonomi, sedangkan perubahan pangsa pekerjaan dihitung antara
  2008 dan 2015 (bagian Urban performance, hlm. 8; persamaan 2, hlm.
  10).
\item
  \textbf{Populasi dan kepadatan:} populasi, kepadatan penduduk, dan
  luas kawasan permukiman tahun 2008 digunakan dalam regresi fungsi;
  pekerjaan dan kepadatan pekerjaan terhadap luas permukiman tahun 2008
  digunakan dalam regresi kinerja (persamaan 1-2, hlm. 10).
\item
  \textbf{Ketersediaan data:} data pendapatan pada skala lokal
  dinyatakan tidak tersedia sehingga penulis memakai pekerjaan sebagai
  ukuran kinerja. Nama berkas, tautan unduh khusus, lisensi, data
  mentah, kode analisis, prosedur data hilang, rincian tahun untuk
  setiap indikator fungsi, dan sumber primer masing-masing dari 47
  indikator: Tidak disebutkan dalam file (bagian Metropolitan functions
  dan Urban performance, hlm. 8; Table 1, hlm. 9).
\end{itemize}

\subsection{Population / Sample}\label{population-sample-19}

\begin{itemize}
\tightlist
\item
  \textbf{Populasi konseptual dan geografi kasus:} sistem perkotaan
  Jerman dipilih karena tidak memiliki satu kota primat seperti sistem
  Prancis atau Inggris, tetapi memiliki struktur polisentris dengan
  sekitar 10-12 kota inti terkemuka dan variasi kondisi struktural serta
  fungsional. Unit spasial dasar penelitian adalah munisipalitas di
  dalam kawasan perkotaan Jerman (bagian Study region, hlm. 5-6).
\item
  \textbf{Kerangka kawasan:} penulis mula-mula mengidentifikasi 32 kota
  inti berdasarkan ambang sedikitnya 200.000 penduduk dan 100.000
  pekerja atau pekerjaan dalam skema asuransi sosial. Lima inti
  berpenduduk lebih dari 600.000 dikategorikan sebagai kota besar dan 27
  inti lain sebagai kota menengah. Analisis jaringan dan penetapan
  setiap munisipalitas ke inti terdekat menghasilkan 20 kawasan
  perkotaan yang tidak bertumpang tindih (bagian Study region, hlm. 6;
  Figure 2, hlm. 7).
\item
  \textbf{Sampel analitis utama:} regresi residual dan klasifikasi empat
  tipe menggunakan 60 munisipalitas yang memiliki nilai indeks atau
  fungsi metropolitan. Menurut Table 3, sampel ini terdiri atas 5 kota
  besar, 26 kota menengah, dan 29 kota kecil. Kota kecil dalam rancangan
  ini adalah munisipalitas yang memiliki fungsi metropolitan, tetapi
  bukan kota besar dan bukan kota inti kawasan (bagian Study region dan
  Urban performance, hlm. 6 dan 8; Table 2-3, hlm. 10 dan 15).
\item
  \textbf{Struktur spasial sampel:} 12 dari 60 munisipalitas berada
  dalam kawasan monosentris dan 48 berada dalam kawasan polisentris.
  Pembagian tersebut adalah hasil delineasi fungsional, bukan
  pengambilan sampel acak (Table 4, hlm. 16).
\item
  \textbf{Periode observasi:} indikator fungsi metropolitan berasal dari
  2008-2010; ukuran awal populasi, pekerjaan, dan kepadatan berasal dari
  2008; kinerja dihitung dari perubahan pangsa pekerjaan 2008-2015
  (bagian Metropolitan functions dan Urban performance, hlm. 8;
  persamaan 1-2, hlm. 10).
\item
  \textbf{Cakupan yang perlu dibedakan:} delineasi kawasan dan peta
  perubahan pekerjaan membahas seluruh munisipalitas dalam kawasan yang
  dianalisis, tetapi klasifikasi residual hanya diterapkan kepada 60
  munisipalitas yang memiliki fungsi metropolitan. Jumlah seluruh
  munisipalitas dalam 20 kawasan sebelum pembatasan tersebut: Tidak
  disebutkan dalam file (Abstract, hlm. 1; bagian Urban performance,
  hlm. 8; uraian Figure 4 dan Table 2, hlm. 10).
\item
  \textbf{Ketidakselarasan jumlah:} rancangan mengidentifikasi 27 kota
  inti menengah, sedangkan Table 3 hanya memuat 26 kota menengah dalam
  sampel klasifikasi. Alasan satu kota inti menengah tidak masuk ke
  dalam 60 kasus: Tidak disebutkan dalam file (bagian Study region, hlm.
  6; Table 3, hlm. 15).
\item
  \textbf{Sifat sampel:} tidak ada partisipan manusia atau unit
  individual. Prosedur sampling probabilistik, kriteria eksklusi selain
  ketersediaan fungsi metropolitan, penanganan data hilang, dan analisis
  daya statistik: Tidak disebutkan dalam file (Research Design, hlm.
  5-10).
\end{itemize}

\subsection{Dependent Variables}\label{dependent-variables-19}

\textbf{Laporan sumber.} Dua variabel pada sisi kiri persamaan OLS
menjadi keluaran empiris yang residualnya dipakai untuk klasifikasi:

\begin{enumerate}
\def\labelenumi{\arabic{enumi}.}
\tightlist
\item
  \textbf{Indeks fungsi metropolitan (\texttt{MF}).} Variabel statis ini
  merupakan hasil standardisasi-z dan penggabungan aditif 47 indikator
  fungsi kontrol, inovasi, gerbang, dan simbolik untuk 2008-2010. Dalam
  regresi pertama, indeks dijelaskan oleh populasi dan kepadatan
  penduduk tahun 2008. Residual positif menunjukkan fungsi lebih banyak
  daripada yang diprediksi berdasarkan ukuran dan kepadatan, sedangkan
  residual negatif menunjukkan fungsi lebih sedikit (bagian Metropolitan
  functions, hlm. 8; Table 1, hlm. 9; persamaan 1 dan uraian residual,
  hlm. 9-10).
\item
  \textbf{Perubahan pangsa pekerjaan (\texttt{EMP.SHARE}).} Penulis
  menghitung pangsa pekerjaan setiap kota dari jumlah pekerjaan seluruh
  60 kota, lalu menghitung perubahannya antara 2008 dan 2015. Dalam
  regresi kedua, perubahan tersebut dijelaskan oleh stok pekerjaan dan
  kepadatan pekerjaan tahun 2008. Residual positif menunjukkan kinerja
  relatif di atas prediksi, sedangkan residual negatif menunjukkan
  kinerja relatif di bawah prediksi (bagian Urban performance, hlm. 8;
  persamaan 2, hlm. 10).
\end{enumerate}

\textbf{Laporan sumber.} Kombinasi tanda kedua residual menghasilkan
keluaran kategoris turunan: fungsi negatif-kinerja negatif menjadi
\emph{agglomeration shadow}; fungsi positif-kinerja negatif menjadi
\emph{borrowed function}; fungsi negatif-kinerja positif menjadi
\emph{borrowed performance}; dan fungsi positif-kinerja positif menjadi
\emph{borrowed size} (Table 2, hlm. 10).

\textbf{Inferensi review.} Label tersebut mengukur posisi relatif
terhadap nilai prediksi dua model. Label bukan ukuran langsung atas
akses, peminjaman, atau transfer suatu fungsi antarkota. Pekerjaan
absolut, pertumbuhan penduduk, pendapatan, produktivitas, PDB per
kapita, arus komuter, dan konektivitas jaringan tidak menjadi variabel
dependen dalam dua regresi A04 (bagian Urban performance dan
Expectations ``given size'', hlm. 8-10).

\subsection{Results}\label{results-19}

\textbf{Daya jelas model yang dilaporkan.} Regresi fungsi metropolitan
menghasilkan \texttt{R2} tersesuaikan sebesar 0,86. Regresi perubahan
pangsa pekerjaan menghasilkan \texttt{R2} tersesuaikan sebesar 0,51;
penulis menilai hasil kedua lebih lemah, tetapi masih dapat diterima.
Koefisien, galat baku, interval kepercayaan, nilai-p, jumlah derajat
kebebasan, dan diagnostik residual selain pembacaan diagram pencar:
Tidak disebutkan dalam file (bagian Expectations ``given size'', hlm.
9-10; Figure 5, hlm. 13).

\textbf{Distribusi fungsi metropolitan dan kinerja.}

\begin{itemize}
\tightlist
\item
  Indeks 2008-2010 menunjukkan hierarki yang didominasi kira-kira oleh
  lima kota inti: Berlin 414,76; Munich 260,54; Hamburg 221,80;
  Frankfurt am Main 206,12; dan Cologne 141,03. Penulis menghubungkan
  keunggulan Berlin dengan relokasi pemerintahan dan institusi publik
  serta budaya ke ibu kota (Analysis and Empirical Findings, hlm. 10;
  Figure 3, hlm. 11).
\item
  Untuk perubahan pangsa pekerjaan 2008-2015, Berlin, Munich, Hamburg,
  dan Cologne memimpin peningkatan pangsa di antara munisipalitas yang
  dibahas; Leipzig, sebuah kota menengah, juga disebut sebagai salah
  satu kota yang berkembang lebih kuat. Tidak ada satu kawasan perkotaan
  dengan pertumbuhan pangsa yang sepenuhnya dominan karena setiap
  kawasan fungsional memuat campuran hinterland berkinerja relatif
  rendah dan tinggi (uraian Figure 4, hlm. 10 dan 12-13).
\end{itemize}

\textbf{Klasifikasi residual keseluruhan.} Dari 60 munisipalitas, 20
masuk \emph{agglomeration shadow}, 7 masuk \emph{borrowed function}, 16
masuk \emph{borrowed performance}, dan 17 masuk \emph{borrowed size}.
Sebagian besar pasangan residual berada dekat titik nol. Dua kasus pada
ujung diagram adalah Berlin, yang berkinerja jauh di atas prediksi
tetapi memiliki sedikit lebih sedikit fungsi daripada prediksi sehingga
masuk \emph{borrowed performance}, dan Frankfurt am Main, yang
berkinerja sedikit di bawah prediksi tetapi memiliki jauh lebih banyak
fungsi sehingga masuk \emph{borrowed function} (Table 2 dan Figure 5,
hlm. 10 dan 13).

\textbf{Distribusi menurut tipe kota.}

\begin{itemize}
\tightlist
\item
  Lima kota besar terbagi menjadi tiga \emph{borrowed function}, satu
  \emph{borrowed performance}, dan satu \emph{borrowed size}; tidak ada
  kota besar dalam \emph{agglomeration shadow} (Table 3, hlm. 15).
\item
  Dari 26 kota menengah, 15 masuk \emph{agglomeration shadow}, 4
  \emph{borrowed function}, dan 7 \emph{borrowed performance}; tidak ada
  kota menengah dalam \emph{borrowed size} (Table 3, hlm. 15).
\item
  Dari 29 kota kecil, 5 masuk \emph{agglomeration shadow}, 8
  \emph{borrowed performance}, dan 16 \emph{borrowed size}; tidak ada
  kota kecil dalam \emph{borrowed function} (Table 3, hlm. 15).
\item
  Munich adalah satu-satunya kota besar dalam \emph{borrowed size}.
  Berlin adalah satu-satunya kota besar dalam \emph{borrowed
  performance}. Hamburg, Cologne, dan Frankfurt am Main masuk
  \emph{borrowed function}; empat kota menengah dalam kategori yang sama
  adalah Hanover, Dusseldorf, Bonn, dan Stuttgart (uraian hasil dan
  Table 3, hlm. 15-16).
\item
  Keenam belas kasus kota kecil dalam \emph{borrowed size} disebut
  berada di sekitar kota besar seperti Munich, Frankfurt am Main,
  Cologne, dan Berlin. Penulis juga membahas Esslingen, Ludwigsburg,
  Waiblingen, Darmstadt, Mainz, Bad Homburg, dan Neuss sebagai kota
  sekitar yang memperoleh efek pekerjaan positif tanpa menjalankan
  fungsi metropolitan penting, yaitu pola yang sesuai dengan
  \emph{borrowed performance} (uraian hasil, hlm. 15).
\end{itemize}

\textbf{Distribusi menurut struktur kawasan.}

\begin{itemize}
\tightlist
\item
  Pada 12 kasus di kawasan monosentris, terdapat 5 \emph{agglomeration
  shadow}, tidak ada \emph{borrowed function}, 5 \emph{borrowed
  performance}, dan 2 \emph{borrowed size} (Table 4, hlm. 16).
\item
  Pada 48 kasus di kawasan polisentris, terdapat 15 \emph{agglomeration
  shadow}, 7 \emph{borrowed function}, 11 \emph{borrowed performance},
  dan 15 \emph{borrowed size} (Table 4, hlm. 16).
\item
  Delapan kasus \emph{agglomeration shadow} berada di Rhine-Ruhr.
  Penulis menempatkan Dortmund, Essen, Duisburg, Bochum, Wuppertal, dan
  Monchengladbach dalam konteks Ruhr dan Bergisches Land yang lemah
  secara struktural, serta menafsirkan Dusseldorf, Cologne, dan Bonn
  sebagai pusat besar yang membayangi kota-kota Ruhr. Kasus
  bayang-bayang lain muncul pada kawasan yang kurang sentral dengan inti
  menengah, seperti Dresden, Aachen, Bremen, Bielefeld, dan Magdeburg,
  serta kawasan Erfurt dan Nuremberg (uraian hasil dan Table 4, hlm.
  15-16).
\end{itemize}

\subsection{Findings}\label{findings-19}

\textbf{Jawaban atas pertanyaan pertama.} Penulis menunjukkan bahwa
kategori empat dimensi dapat dioperasionalkan dengan dua regresi OLS:
satu untuk fungsi metropolitan dan satu untuk perubahan pangsa
pekerjaan. Tanda kedua residual menyediakan aturan klasifikasi statistik
yang seragam bagi 60 munisipalitas (bagian Expectations ``given size'',
hlm. 9-10; Table 2, hlm. 10).

\textbf{Jawaban atas pertanyaan kedua.} Keempat tipe ditemukan dalam
sistem perkotaan Jerman, tetapi geografinya tidak seragam. Kasus positif
terdapat pada kota besar maupun kota sekitar, sedangkan
\emph{agglomeration shadow} muncul baik di kawasan Rhine-Ruhr yang
polisentris maupun di sejumlah kawasan yang lebih terisolasi. Karena
itu, kedekatan dengan aglomerasi besar tidak menghasilkan satu jenis
efek yang sama pada semua kota (Figure 6 dan uraian hasil, hlm. 14-16).

\textbf{Jawaban atas pertanyaan ketiga.} Ukuran kota berkaitan kuat
dengan tipe hasil: tiga dari lima kota besar memiliki fungsi lebih
tinggi daripada prediksi tetapi kinerja lebih rendah; 16 dari 17 kasus
\emph{borrowed size} adalah kota kecil; dan 15 dari 20 kasus
\emph{agglomeration shadow} adalah kota menengah. Struktur kawasan juga
membedakan distribusi: semua tujuh kasus \emph{borrowed function} dan 15
dari 17 kasus \emph{borrowed size} berada di kawasan polisentris,
sedangkan lima dari 12 kasus monosentris berupa \emph{borrowed
performance} (Table 3-4, hlm. 15-16).

\textbf{Interpretasi penulis.} Penulis menilai kota besar berfungsi
sebagai pusat pertumbuhan dan memperoleh dukungan dari kota menengah
serta kecil dalam jaringan kawasan. Mereka juga menafsirkan kota-kota
polisentris sebagai lebih mampu mentransfer keunggulan fungsional, kota
monosentris besar sebagai lebih efektif dalam \emph{borrowed
performance}, dan kota atau kawasan yang relatif kecil serta terisolasi
sebagai lebih sulit mempertahankan tingkat minimum metropolitanitas dan
kinerja dalam sistem nasional (uraian hasil, hlm. 15-16).

\textbf{Inferensi review.} Table 4 tidak mendukung anggapan sederhana
bahwa polisentrisme selalu menguntungkan: selain memuat seluruh kasus
\emph{borrowed function} dan mayoritas \emph{borrowed size}, kawasan
polisentris juga memuat 15 dari 20 kasus \emph{agglomeration shadow}.
Karena 48 dari 60 kasus memang berada di kawasan polisentris,
perbandingan jumlah mentah juga perlu dibaca bersama ketidakseimbangan
ukuran kelompok. Bukti artikel menunjukkan asosiasi spasial dan deviasi
dari ekspektasi ukuran, bukan efek kausal ukuran kota, sentralitas, atau
struktur kawasan (Table 4, hlm. 16; Conclusion, hlm. 17).

\subsection{Conclusions}\label{conclusions-19}

\textbf{Laporan sumber.} Penulis menyimpulkan bahwa kerangka empat
dimensi Meijers dan Burger dapat dioperasionalkan untuk
mengklasifikasikan kota dan menggambarkan distribusi spasial
\emph{borrowed size}, \emph{borrowed performance}, \emph{borrowed
function}, dan \emph{agglomeration shadow} dalam sistem perkotaan
Jerman. Pertumbuhan fungsi dan ekonomi tidak hanya terkonsentrasi di
kota besar, meskipun ukuran kota tetap dipandang sebagai faktor sentral
bagi keberadaan fungsi metropolitan (Conclusion, hlm. 17).

\textbf{Interpretasi penulis.} Ukuran relatif kota dan lokasinya dalam
kawasan monosentris atau polisentris dinilai berdampak penting terhadap
tipe yang muncul. Konteks kawasan perlu diperhitungkan karena hubungan
antara kota besar, pusat dalam kota, dan diferensiasi fungsi suburban
membentuk pembagian kerja spasial yang terintegrasi. Ekonomi aglomerasi
karena itu tidak dipandang berhenti pada batas fisik kota, tetapi dapat
melimpah ke kota atau kawasan sekitarnya (Conclusion, hlm. 17).

\textbf{Interpretasi penulis.} Penulis menghubungkan hasil tersebut
dengan proses konsentrasi dan dekonsentrasi dalam kawasan perkotaan
serta konsep \emph{interplaces}. Mereka menyatakan bahwa infrastruktur
transportasi regional dan teknologi informasi serta komunikasi dapat
mendukung \emph{borrowed size}, dan bahwa identifikasi bentuk
perkembangan spasial diperlukan agar kebijakan disesuaikan dengan
struktur kawasan (Conclusion, hlm. 17-18).

\textbf{Inferensi review.} Kesimpulan yang paling kuat didukung oleh
analisis adalah bahwa empat pola residual memiliki distribusi berbeda
menurut tipe kota dan konteks kawasan. Kesimpulan bahwa fungsi atau
manfaat benar-benar dipinjam, ditransfer, atau disebabkan oleh kota
tetangga tetap bersifat interpretatif karena penelitian tidak mengukur
arus interaksi atau mengidentifikasi hubungan kausal (persamaan 1-2 dan
Table 2, hlm. 10; Conclusion, hlm. 17-18).

\subsection{Contributions}\label{contributions-19}

\begin{itemize}
\tightlist
\item
  \textbf{Kontribusi konseptual:} artikel menerjemahkan kerangka empat
  dimensi yang semula analitis menjadi tipologi empiris bersama antara
  fungsi dan kinerja. Dengan demikian, \emph{borrowed function} dan
  \emph{borrowed performance} tidak dianalisis secara terpisah dari
  kombinasi \emph{borrowed size} dan \emph{agglomeration shadow}
  (Research gap, hlm. 5; bagian Expectations ``given size'', hlm. 9-10).
\item
  \textbf{Kontribusi pengukuran:} artikel menggabungkan indeks 47
  indikator fungsi metropolitan dengan perubahan pangsa pekerjaan,
  kemudian mengendalikan ekspektasi ukuran dan kepadatan melalui dua
  regresi residual. Penambahan kepadatan serta pilar kinerja merupakan
  perluasan yang dinyatakan penulis terhadap pendekatan regresi
  sebelumnya (bagian Metropolitan functions sampai Expectations ``given
  size'', hlm. 8-10).
\item
  \textbf{Kontribusi empiris-spasial:} operasionalisasi diterapkan pada
  skala munisipalitas terhadap 60 kasus dalam kerangka 20 kawasan
  perkotaan Jerman, mencakup kota besar, menengah, dan kecil serta
  konteks monosentris dan polisentris. Hasil tidak hanya disajikan
  sebagai model, tetapi juga sebagai diagram residual, peta, dan
  distribusi menurut tipe kota serta struktur kawasan (Figure 5-6 dan
  Table 3-4, hlm. 13-16).
\item
  \textbf{Kontribusi substantif:} artikel menunjukkan pemisahan yang
  berguna antara fungsi dan kinerja. Sebuah kota dapat memiliki fungsi
  di atas prediksi tetapi kinerja di bawah prediksi, seperti Frankfurt
  am Main, atau memperlihatkan pola sebaliknya, seperti Berlin. Temuan
  ini mencegah fungsi metropolitan diperlakukan otomatis setara dengan
  kinerja ekonomi relatif (uraian Figure 5, hlm. 13).
\item
  \textbf{Kontribusi terhadap perdebatan polisentrisitas:} hasil
  memperlihatkan bahwa tipe positif dan negatif dapat hidup berdampingan
  dalam sistem serta kawasan polisentris yang sama. Lokasi dalam kawasan
  besar bukan jaminan tunggal atas manfaat, dan ukuran kota perlu dibaca
  bersama konteks spasial dalam interpretasi penulis (uraian hasil, hlm.
  15-16; Conclusion, hlm. 17).
\item
  \textbf{Inferensi review atas nilai tambah:} nilai utama A04 terletak
  pada perangkat diagnosis komparatif yang transparan pada tingkat
  aturan klasifikasi. Nilai tambahnya bukan pembuktian mekanisme kausal,
  melainkan cara yang konsisten untuk menemukan kasus yang menyimpang
  dari pola fungsi dan kinerja yang diprediksi oleh ukuran serta
  kepadatan (persamaan 1-2 dan Table 2, hlm. 10).
\end{itemize}

\subsection{Research Gap}\label{research-gap-19}

\textbf{Celah yang dinyatakan sebelum penelitian.}

\begin{itemize}
\tightlist
\item
  Literatur empiris dinilai belum menyediakan operasionalisasi, analisis
  distribusi spasial, dan klasifikasi kota pada skala kecil dengan
  kumpulan indikator yang kuat serta metode yang konsisten
  (Introduction, hlm. 2).
\item
  Kerangka empat dimensi tersedia, tetapi matriks teoretisnya perlu
  diterjemahkan menjadi indikator fungsi dan kinerja yang dapat diukur
  serta aturan untuk menentukan hasil yang lebih tinggi atau lebih
  rendah daripada ekspektasi berdasarkan ukuran (Research gap dan awal
  Research Design, hlm. 5-6).
\item
  Literatur mengenai kinerja perkotaan belum memberi jawaban empiris
  rinci tentang cara mengukur kinerja dalam kerangka \emph{borrowed
  size}; Meijers dan Burger disebut menggambarkan dimensinya tanpa
  mengoperasionalkannya secara terperinci (bagian Urban performance,
  hlm. 8).
\item
  Analisis diferensiasi spasial dan fungsional dalam sistem perkotaan
  disebut masih langka, khususnya mengenai hubungan antara fungsi
  perkotaan dan saling ketergantungan pertumbuhan (Theoretical
  Foundation, hlm. 4).
\end{itemize}

\textbf{Celah yang tetap terbuka menurut penulis.}

\begin{itemize}
\tightlist
\item
  Indeks metropolitan agregat belum membedakan fungsi tipikal bagi
  setiap dimensi \emph{borrowed size}, distribusi dan spesialisasi
  fungsi antarkota, kaitannya dengan koaglomerasi spasial, serta
  perubahan saling ketergantungan dari waktu ke waktu (Conclusion, hlm.
  17-18).
\item
  Faktor pendorong kemampuan suatu kota untuk meminjamkan atau
  menggunakan massa kritis kota di dekatnya masih memerlukan penjelasan
  lebih dalam. Penulis menunjuk pendapatan, koneksi transportasi,
  aksesibilitas, dan arus komuter sebagai indikator yang dibutuhkan
  (Conclusion, hlm. 17-18).
\end{itemize}

\textbf{Inferensi review atas celah tersisa.} Hubungan interaksi yang
dianggap sentral dalam teori belum diuji secara langsung. Rancangan
belum membedakan efek kedekatan, konektivitas, komplementaritas fungsi,
seleksi lokasi, dan kondisi awal; belum menguji apakah hasil bertahan
terhadap definisi kawasan, ambang kota, spesifikasi regresi, atau
dependensi spasial alternatif. Karena itu, celah antara label tipologi
dan mekanisme \emph{borrowing} masih terbuka (Theoretical Foundation,
hlm. 3-4; Research Design, hlm. 6-10; Conclusion, hlm. 17-18).

\subsection{Limitations}\label{limitations-19}

\textbf{Keterbatasan yang dinyatakan atau diakui penulis.}

\begin{itemize}
\tightlist
\item
  Kinerja dipersempit menjadi dimensi ekonomi berdasarkan pekerjaan
  dalam skema asuransi sosial. Penulis menyatakan pilihan ini mungkin
  tidak optimal dan menyebut indikator demografi atau pendapatan sebagai
  alternatif; pendapatan tidak dipakai karena tidak tersedia pada skala
  lokal yang dibutuhkan (bagian Urban performance, hlm. 8).
\item
  Indeks fungsi diperlakukan statis. Penulis memilih indikator statis
  dengan alasan fungsi metropolitan tidak berubah secara nyata antara
  2008 dan 2015, tetapi analisis ini tidak mengukur perubahan fungsi
  sebagai keluaran dinamis (bagian Metropolitan functions, hlm. 8;
  persamaan 1 dan penjelasannya, hlm. 10).
\item
  Penulis menyebut pendekatan ini sebagai titik awal dan mengakui
  kebutuhan akan pemisahan jenis fungsi, variabel tambahan, analisis
  perubahan temporal, dan pemahaman lebih dalam atas faktor pendorong
  (Conclusion, hlm. 17-18).
\end{itemize}

\textbf{Keterbatasan yang teridentifikasi dalam review TXT.}

\begin{itemize}
\tightlist
\item
  \textbf{Cakupan kasus:} klasifikasi dibatasi pada 60 munisipalitas
  dengan fungsi metropolitan. Munisipalitas tanpa nilai indeks tidak
  memperoleh kategori, sehingga klaim tentang seluruh munisipalitas
  dalam kawasan perlu dibedakan dari cakupan model residual (bagian
  Urban performance, hlm. 8; Table 2, hlm. 10).
\item
  \textbf{Sifat relatif ukuran kinerja:} pangsa setiap kota dihitung
  terhadap total pekerjaan 60 kota. Secara konstruksi, kenaikan pangsa
  satu atau beberapa kota diimbangi perubahan pangsa kota lain;
  perubahan ini mengukur posisi relatif dan tidak menunjukkan apakah
  pekerjaan absolut suatu kota tumbuh atau menyusut (bagian Urban
  performance, hlm. 8).
\item
  \textbf{Atribusi spasial indeks:} nilai indeks pada tingkat county
  dipindahkan ke munisipalitas terbesar di county tersebut. Prosedur ini
  dapat memusatkan atribut yang sebenarnya tersebar ke satu
  munisipalitas; pengujian sensitivitas terhadap aturan atribusi
  tersebut tidak dilaporkan (bagian Metropolitan functions, hlm. 8).
\item
  \textbf{Agregasi pengukuran:} 47 indikator distandardisasi dan
  digabungkan secara aditif menjadi satu indeks. Bobot rinci, korelasi
  antarindikator, validasi konstruk, uji reliabilitas, pengaruh
  indikator ekstrem, dan hasil menurut empat kelompok fungsi secara
  terpisah: Tidak disebutkan dalam file (bagian Metropolitan functions,
  hlm. 8; Table 1, hlm. 9).
\item
  \textbf{Pelaporan regresi:} selain \texttt{R2} tersesuaikan dan
  diagram residual, file tidak melaporkan koefisien, ketidakpastian
  estimasi, nilai-p, transformasi variabel, pemeriksaan
  heteroskedastisitas, multikolinearitas, normalitas, pengamatan
  berpengaruh, atau uji dependensi spasial. Karena banyak residual
  disebut dekat nol, klasifikasi berdasarkan tanda dapat peka terhadap
  kesalahan ukur atau spesifikasi, tetapi ketidakpastian kategori tidak
  dihitung (bagian Expectations ``given size'', hlm. 9-10; Figure 5 dan
  uraiannya, hlm. 13).
\item
  \textbf{Identifikasi mekanisme:} akses terhadap fungsi kota lain, arus
  komuter, hubungan kooperatif, pertukaran informasi, dan konektivitas
  disebut penting oleh teori, tetapi tidak menjadi variabel dalam dua
  model. Kedekatan atau keanggotaan kawasan tidak dengan sendirinya
  membuktikan proses peminjaman (Theoretical Foundation, hlm. 3-4;
  persamaan 1-2, hlm. 10).
\item
  \textbf{Delineasi dan kategori kota:} file tidak melaporkan uji
  sensitivitas terhadap ambang 200.000 penduduk, 100.000 pekerjaan,
  batas 600.000 penduduk, buffer 30-60 menit, atau aturan penempatan ke
  inti terdekat. Penyederhanaan kawasan yang bertumpang tindih menjadi
  kawasan diskret dapat menghilangkan hubungan lintas-batas (bagian
  Study region, hlm. 6).
\item
  \textbf{Konsistensi pelaporan:} 27 inti menengah didefinisikan dalam
  rancangan, tetapi Table 3 memuat 26 kota menengah. Selain itu, kawasan
  polisentris didefinisikan sebagai memiliki sedikitnya dua kota inti
  pada halaman 6, tetapi sebagai memiliki lebih dari dua kota inti pada
  halaman 16. Kedua ketidakselarasan tidak dijelaskan dalam file (bagian
  Study region, hlm. 6; uraian Table 4, hlm. 16).
\item
  \textbf{Klaim signifikansi:} Abstract menyebut ukuran dan derajat
  sentralitas memiliki pengaruh signifikan, dan Conclusion menyebut
  dampak signifikan. Statistik inferensial yang dapat memverifikasi
  penggunaan istilah ``signifikan'' tidak disajikan dalam TXT; hasil
  yang tersedia terutama berupa regresi untuk menghasilkan residual dan
  tabulasi deskriptif tipe (Abstract, hlm. 1; Analysis and Empirical
  Findings, hlm. 10-16; Conclusion, hlm. 17).
\item
  \textbf{Keterterapan eksternal:} kasus Jerman dipilih karena struktur
  historis dan polisentrisitasnya yang khas. Pengujian lintas negara
  atau sistem perkotaan lain tidak dilakukan, sehingga generalisasi di
  luar konteks Jerman belum diuji dalam artikel (bagian Study region,
  hlm. 5-6).
\end{itemize}

\subsection{Challenges}\label{challenges-19}

\textbf{Status bagian.} Artikel tidak memiliki bagian khusus bernama
\emph{Challenges}. Poin berikut merangkum persoalan pengukuran dan
analisis yang dinyatakan atau diperlihatkan dalam keputusan rancangan
sumber.

\begin{itemize}
\tightlist
\item
  \textbf{Menerjemahkan konsep ke ukuran:} tantangan inti penelitian
  adalah mengubah matriks teoretis empat bidang menjadi dua ukuran yang
  dapat menentukan kapan kota memiliki fungsi atau kinerja di atas dan
  di bawah ekspektasi berdasarkan ukuran (Research gap dan
  Operationalisation, hlm. 5-6).
\item
  \textbf{Memilih indikator fungsi nontrivial:} penulis perlu memilih
  atribut yang menangkap karakter metropolitan, bukan sekadar populasi
  atau kepadatan, sekaligus menjaga ketersediaan, keterbandingan waktu,
  objektivitas, dan keandalan statistik pada skala county serta kota
  (bagian Metropolitan functions, hlm. 6 dan 8).
\item
  \textbf{Mendefinisikan kinerja yang adil antarskala:} perubahan
  absolut menguntungkan kota dengan basis pekerjaan besar, sedangkan
  persentase pertumbuhan dapat membesar akibat basis kecil. Penulis
  merespons masalah ini dengan perubahan pangsa pekerjaan, tetapi ukuran
  yang dipilih tetap hanya menggambarkan pergeseran kepentingan relatif
  (bagian Urban performance, hlm. 8).
\item
  \textbf{Memisahkan ukuran dari metropolitanitas:} karena fungsi
  tingkat tinggi sering berkaitan dengan massa kritis, model harus
  menentukan fungsi dan kinerja yang melampaui nilai yang diharapkan
  dari ukuran. Penulis menambahkan kepadatan ke dalam regresi yang
  ringkas untuk menangkap dimensi aglomerasi yang tidak
  direpresentasikan ukuran saja (bagian Expectations ``given size'',
  hlm. 9-10).
\item
  \textbf{Mendelineasi kawasan yang tumpang tindih:} jaringan perjalanan
  dari 32 inti menghasilkan tumpang tindih, khususnya Duisburg, Essen,
  dan Dortmund. Penulis harus menetapkan setiap munisipalitas kepada
  inti terdekat agar terbentuk 20 kawasan diskret yang dapat
  dibandingkan (bagian Study region, hlm. 6).
\item
  \textbf{Menafsirkan heterogenitas spasial:} pola positif dan negatif
  muncul di dalam kawasan yang sama, sementara mekanisme yang
  dikemukakan teori meliputi jarak, aksesibilitas, jaringan, hubungan
  kooperatif, dan arus informasi. Dengan variabel yang tersedia, artikel
  menghadapi tantangan untuk membedakan mekanisme tersebut dari konteks
  ukuran dan lokasi (Theoretical Foundation, hlm. 3-4; Analysis and
  Empirical Findings, hlm. 13-16).
\item
  \textbf{Keterbatasan data lokal:} tidak tersedianya pendapatan pada
  skala lokal mendorong penggunaan pekerjaan. Data koneksi transportasi,
  aksesibilitas, dan arus komuter baru diajukan sebagai kebutuhan
  penelitian berikutnya, bukan tersedia dalam analisis ini (bagian Urban
  performance, hlm. 8; Conclusion, hlm. 18).
\end{itemize}

\subsection{Practical Implications}\label{practical-implications-19}

\textbf{Implikasi yang dinyatakan penulis.}

\begin{itemize}
\tightlist
\item
  Pembuat kebijakan dan perencana regional perlu mengidentifikasi serta
  mengklasifikasikan fase dan bentuk spasial perkembangan kawasan
  sebelum memilih kebijakan. Pertumbuhan monosentris dan perkembangan
  kota polisentris memerlukan konsep perencanaan infrastruktur yang
  berbeda (Introduction, hlm. 1).
\item
  Penulis menyatakan bahwa ekonomi aglomerasi tidak terbatas pada batas
  fisik kota dan dapat melimpah ke kota sekitar. Mereka menghubungkan
  dukungan terhadap \emph{borrowed size} dengan perbaikan transportasi
  regional serta teknologi informasi dan komunikasi (Conclusion, hlm.
  17).
\item
  Kota kecil dan menengah di sekitar pusat besar dinilai perlu memiliki
  peran khusus dalam pembangunan kembali perkotaan dan kebijakan
  perencanaan pada masa depan (Conclusion, hlm. 18).
\item
  \textbf{Inferensi review:} pemetaan empat tipe dapat membantu
  membedakan kota yang menunjukkan kelebihan fungsi, kelebihan kinerja,
  keduanya, atau justru kekurangan relatif pada kedua dimensi. Perbedaan
  ini relevan agar wilayah \emph{agglomeration shadow} tidak
  diperlakukan sama dengan kota yang mengalami \emph{borrowed
  performance} atau \emph{borrowed function} (Table 2, hlm. 10; uraian
  hasil, hlm. 15-16).
\end{itemize}

\textbf{Inferensi review untuk penggunaan kebijakan.} Tipologi sebaiknya
dipakai sebagai diagnosis awal, bukan dasar tunggal untuk memilih
intervensi. Sebelum menyimpulkan bahwa konektivitas atau kedekatan
menyebabkan suatu kategori, pengguna kebijakan perlu memeriksa
pertumbuhan pekerjaan absolut, distribusi fungsi tertentu, arus komuter,
aksesibilitas, hubungan jaringan, dan kondisi lokal yang tidak masuk ke
dalam model A04 (bagian Urban performance, hlm. 8; Conclusion, hlm.
17-18).

\subsection{Applications}\label{applications-19}

\begin{itemize}
\tightlist
\item
  \textbf{Aplikasi dalam artikel:} dua residual dipakai untuk
  mengklasifikasikan dan memetakan 60 munisipalitas Jerman ke dalam
  empat tipe, kemudian membandingkannya menurut ukuran kota dan struktur
  kawasan. Ini adalah aplikasi deskriptif-komparatif pada skala
  munisipalitas (Table 2, Figure 5-6, dan Table 3-4, hlm. 10 dan 13-16).
\item
  \textbf{Aplikasi sebagai penyaringan kasus:} sebagai inferensi review,
  kota dengan kombinasi residual yang berbeda dapat dipilih untuk studi
  kasus lanjutan. Misalnya, Berlin dan Frankfurt am Main menyediakan
  pola fungsi-kinerja yang berlawanan, sedangkan Rhine-Ruhr memuat
  konsentrasi kasus \emph{agglomeration shadow} (uraian Figure 5, hlm.
  13; Table 4 dan uraian hasil, hlm. 15-16).
\item
  \textbf{Aplikasi perbandingan kawasan:} kerangka dapat digunakan untuk
  membandingkan susunan fungsi dan kinerja antara kota besar, menengah,
  dan kecil serta antara kawasan monosentris dan polisentris. Penulis
  menyebut pendekatannya sebagai titik awal bagi penelitian lanjutan
  menggunakan konsep yang telah dioperasionalkan (Conclusion, hlm. 17).
\item
  \textbf{Aplikasi dengan ukuran yang diperluas:} rancangan dapat
  dikembangkan dengan memisahkan kelompok fungsi metropolitan atau
  menambahkan pendapatan, koneksi transportasi, aksesibilitas, dan arus
  komuter untuk menguji penjelasan yang lebih terperinci (Conclusion,
  hlm. 17-18).
\item
  \textbf{Batas aplikasi:} artikel tidak menguji alat prediksi, evaluasi
  program, simulasi skenario, atau penarikan kesimpulan kausal. Data
  mentah, kode, dan paket implementasi juga tidak disebutkan dalam file,
  sehingga penerapan ulang memerlukan rekonstruksi dataset dan keputusan
  pengukuran dari uraian metode (bagian Research Design, hlm. 5-10).
\end{itemize}

\subsection{Future Research
(Optional)}\label{future-research-optional-19}

\textbf{Agenda yang disebutkan langsung oleh penulis.}

\begin{itemize}
\tightlist
\item
  Memisahkan indeks metropolitan ke dalam ragam fungsi yang lebih luas
  agar dapat diketahui fungsi tipikal bagi setiap dimensi \emph{borrowed
  size} dan kekayaan data tidak hilang dalam indeks agregat (Conclusion,
  hlm. 17).
\item
  Meneliti hubungan fungsi tersebut dengan koaglomerasi spasial dan
  menguji apakah saling ketergantungan berubah dari waktu ke waktu
  (Conclusion, hlm. 17).
\item
  Mengkaji distribusi serta spesialisasi fungsi metropolitan antarkota
  dalam satu kawasan untuk mengidentifikasi fungsi mana yang dapat
  dipinjam dari kota kecil dan menengah di sekitarnya (Conclusion, hlm.
  17).
\item
  Memperdalam penjelasan mengenai faktor yang memungkinkan kota
  meminjamkan atau menggunakan massa kritis dan ukuran kota yang
  berdekatan (Conclusion, hlm. 17-18).
\item
  Menambahkan indikator pendapatan, koneksi transportasi, aksesibilitas,
  dan arus komuter dalam kawasan perkotaan untuk menjelaskan geografi
  empat dimensi dengan lebih baik (Conclusion, hlm. 18).
\item
  Mengembangkan pemahaman mengenai peran khusus kota kecil dan menengah
  di sekitar pusat dalam pembangunan kembali perkotaan dan kebijakan
  perencanaan (Conclusion, hlm. 18).
\end{itemize}

\textbf{Inferensi review yang mengikuti batas bukti A04.} Replikasi
dapat menguji ketahanan kategori terhadap ambang delineasi, spesifikasi
fungsional dan spasial, aturan atribusi county-ke-munisipalitas, serta
ketidakpastian residual. Desain longitudinal atau jaringan juga
diperlukan untuk membedakan limpahan kausal, komplementaritas fungsi,
dan seleksi lokasi. Agenda ini tidak dinyatakan dengan rincian tersebut
oleh penulis, tetapi mengikuti keterbatasan metode dan kebutuhan akan
analisis perubahan serta interaksi yang mereka akui (bagian Metropolitan
functions dan Expectations ``given size'', hlm. 8-10; Conclusion, hlm.
17-18).

\subsection{Provenance Notes}\label{provenance-notes-19}

\begin{itemize}
\tightlist
\item
  Review ini hanya menggunakan TXT A04:
  \texttt{A04-volgmann-rusche-2020-the-geography-of-borrowing-size-exploring-spatial-distributions-for-german-urban-regions-10.1111-tesg.12362.txt}.
  Seluruh 903 baris yang tersedia diperiksa, termasuk metadata
  ekstraksi, isi artikel, tabel, keterangan gambar, kesimpulan, dan
  daftar pustaka. Tidak ada informasi substantif dari sumber eksternal
  yang ditambahkan.
\item
  Blok metadata menyatakan bahwa TXT diekstrak dari PDF pada 31 Agustus
  2026 dengan \texttt{pdftotext\ -layout}. PDF memiliki 20 halaman,
  tidak bertag, dan tidak dienkripsi (blok {[}PDF Metadata{]}, sebelum
  teks artikel).
\item
  Locator \texttt{hlm.} dalam review mengikuti nomor halaman tercetak
  artikel, bukan nomor baris TXT atau urutan halaman pembaca PDF.
  Locator Figure dan Table mengikuti label asli sumber.
\item
  Figure 2-6 terutama hadir sebagai gambar dengan keterangan, sedangkan
  detail visual peta dan diagram tidak seluruhnya terwakili dalam teks
  hasil ekstraksi. Review hanya memakai angka dan interpretasi visual
  yang dinyatakan dalam prosa atau tabel; pola simbol, batas peta, dan
  posisi titik tidak dibaca sebagai bukti tambahan.
\item
  Tata letak dua kolom menyebabkan beberapa kalimat pada halaman 15-16
  tersusun berselang dalam TXT. Nama kota dan kategori hanya digunakan
  ketika hubungan tersebut dapat dipulihkan dari definisi kategori,
  Table 3-4, dan uraian yang berdekatan. Ejaan seperti
  \texttt{Muenster}, \texttt{Dusseldorf}, dan \texttt{Monchengladbach}
  dipertahankan dari TXT dan tidak dinormalisasi melalui sumber
  eksternal.
\item
  Artikel tercetak menyatakan tahun 2019, Volume 0, Nomor 0, dan halaman
  1-20, sedangkan nama file memakai tahun 2020. Review mempertahankan
  ketidakselarasan tersebut dan tidak menyelesaikannya melalui DOI atau
  katalog eksternal (baris bibliografis pada hlm. 1; nama file sumber).
\item
  Status telaah sejawat, versi publikasi akhir, lisensi dataset, sumber
  primer terperinci untuk setiap indikator, data mentah, dan kode
  analisis: Tidak disebutkan dalam file. Daftar pustaka dan laporan
  artikel mengenai penelitian terdahulu tidak diverifikasi secara
  independen.
\item
  Rancangan menyebut 27 kota inti menengah, sedangkan Table 3 memuat 26
  kota menengah. Table 4 menampilkan baris Chemnitz tanpa kasus, tetapi
  file tidak menyatakan apakah baris kosong itu menjelaskan selisih
  jumlah (bagian Study region, hlm. 6; Table 3-4, hlm. 15-16).
\item
  Definisi kawasan polisentris berubah dari sedikitnya dua inti pada
  halaman 6 menjadi lebih dari dua inti pada halaman 16. File tidak
  merekonsiliasi kedua rumusan tersebut (bagian Study region, hlm. 6;
  uraian Table 4, hlm. 16).
\item
  Klaim Abstract bahwa analisis mengintegrasikan seluruh munisipalitas
  perlu dibedakan dari klasifikasi residual yang hanya mencakup 60
  munisipalitas dengan fungsi metropolitan. Jumlah seluruh munisipalitas
  yang masuk ke dalam tahap delineasi dan pemetaan tidak diberikan
  (Abstract, hlm. 1; bagian Urban performance, hlm. 8; Table 2, hlm.
  10).
\item
  Indeks metropolitan berulang kali diberi periode 2008-2010, tetapi
  uraian setelah persamaan 1 menyatakan bahwa semua nilai berasal dari
  2008. Review memakai 2008-2010 untuk indeks dan 2008 untuk prediktor
  populasi serta kepadatan; batas temporal yang lebih rinci tidak dapat
  diselesaikan dari file (bagian Metropolitan functions, hlm. 8;
  persamaan 1 dan uraian Analysis and Empirical Findings, hlm. 10).
\item
  Istilah ``signifikan'' dilaporkan sebagai pilihan kata penulis dalam
  Abstract dan Conclusion. Karena TXT tidak memuat koefisien, galat
  baku, nilai-p, atau pengujian formal hubungan kategori dengan tipe
  kota dan struktur kawasan, review tidak mengubahnya menjadi klaim
  signifikansi statistik yang telah terverifikasi (Abstract, hlm. 1;
  Conclusion, hlm. 17).
\item
  Seluruh label konseptual dipertahankan dalam bahasa asli dan dicetak
  miring agar tidak disamakan dengan observasi langsung. Pernyataan
  mengenai perpindahan, peminjaman, limpahan, atau bayang-bayang yang
  melampaui keluaran residual diberi status interpretasi penulis atau
  inferensi review.
\end{itemize}

\section{Review A05 - Is there a borrowed size in China's urban
agglomerations?}\label{review-a05---is-there-a-borrowed-size-in-chinas-urban-agglomerations}

\subsection{Identitas Sumber}\label{identitas-sumber-9}

\begin{itemize}
\tightlist
\item
  \textbf{Kode sumber:} A05.
\item
  \textbf{Judul:} \emph{Is there a borrowed size in China's urban
  agglomerations?}
\item
  \textbf{Penulis:} Yang Tongbin, Zhu Yingming, dan Du Jiazhen.
\item
  \textbf{Afiliasi:} School of Economics and Management, Nanjing
  University of Science and Technology; serta Jiangsu Industrial Cluster
  Decision-Making Consulting Research Base, Nanjing University of
  Science and Technology, Nanjing, China.
\item
  \textbf{Tahun:} 2022.
\item
  \textbf{Jurnal:} \emph{Progress in Geography}, volume 41, nomor 7,
  halaman 1156-1167.
\item
  \textbf{DOI:} 10.18306/dlkxjz.2022.07.002.
\item
  \textbf{Tanggal naskah:} diterima 17 Januari 2022 dan direvisi 17
  Maret 2022.
\item
  \textbf{Jenis dan status publikasi:} artikel jurnal empiris. Status
  penelaahan sejawat: Tidak disebutkan dalam file.
\item
  \textbf{Kata kunci:} \emph{borrowed size}, \emph{borrowed function},
  \emph{borrowed performance}, \emph{agglomeration shadow}, dan
  aglomerasi perkotaan.
\item
  \textbf{Pendanaan yang dicantumkan:} National Social Science
  Foundation of China (nomor 20BJL106 dan 21CJY075) serta Cultural
  Experts and ``Four batch'' Talents Independently Selected Topic
  Project (nomor ZXGZ(2018)86).
\item
  \textbf{Bahasa sumber:} isi utama berbahasa Mandarin; judul, identitas
  penulis, abstrak, dan kata kunci juga tersedia dalam bahasa Inggris
  pada halaman 1167. Abstrak Inggris merupakan padanan dari abstrak
  Mandarin, bukan bukti atau hasil tambahan.
\item
  \textbf{Cakupan akses:} teks lengkap hasil ekstraksi 12 halaman PDF
  dengan \texttt{pdftotext\ -layout}, termasuk metadata ekstraksi, tabel
  yang terbaca sebagai teks, keterangan gambar, daftar pustaka, dan
  abstrak dwibahasa.
\end{itemize}

\subsection{Summarized Abstract}\label{summarized-abstract-20}

\textbf{Laporan sumber.} Artikel ini meneliti apakah \emph{borrowed
size} terdapat dalam 14 aglomerasi perkotaan China, kondisi spasial yang
mendukungnya, dan mekanisme pembentukannya. Analisis menggunakan data
kota setingkat prefektur atau lebih tinggi selama 2008-2019.
\emph{Borrowed size} dinilai melalui dua dimensi, yaitu \emph{borrowed
function} dan \emph{borrowed performance}. Penulis memakai metode
residual untuk mengidentifikasi kedua dimensi tersebut, menguji pengaruh
struktur polisentris dengan model panel, lalu menguji kedekatan
geografis, keterhubungan jaringan, dan interaksi keduanya dengan data
penampang lintang 2019 (hlm. 1156; bagian 1.1-1.3, hlm. 1158-1160).

\textbf{Laporan sumber.} Pada tingkat aglomerasi, tujuh wilayah
menunjukkan \emph{borrowed function} dan \emph{borrowed performance}
sekaligus pada 2019: Chang-Zhu-Tan, Delta Sungai Mutiara, Wuhan,
Liaoning Tengah dan Selatan, Harbin-Changchun, Teluk Beibu, serta
Beijing-Tianjin-Hebei. Kota besar lebih mungkin menunjukkan kedua
dimensi, sedangkan kota kecil dan menengah lebih sering hanya
menunjukkan \emph{borrowed function}. Lima aglomerasi -
Chengdu-Chongqing, Dataran Tengah, Semenanjung Shandong, Pesisir
Guangdong-Fujian-Zhejiang, dan Delta Sungai Yangtze - menunjukkan
\emph{agglomeration shadow}. Struktur polisentris berasosiasi positif
dengan kedua dimensi; semua kelompok ukuran kota memperoleh
\emph{borrowed function}, tetapi hanya kota menengah dan besar yang
memperoleh \emph{borrowed performance} secara signifikan. Kedekatan
geografis berasosiasi negatif dengan \emph{borrowed function} tetapi
positif dengan \emph{borrowed performance}. Keterhubungan jaringan
berasosiasi positif dengan keduanya pada sampel penuh, tetapi efek
promosi yang signifikan pada analisis heterogenitas hanya tampak pada
kota besar (bagian 2.1-2.3 dan 3.1, hlm. 1160-1165).

\textbf{Interpretasi penulis.} Penulis memaknai hasil tersebut sebagai
bukti bahwa manfaat aglomerasi tidak selalu terkunci di dalam batas
kota: eksternalitas jaringan memungkinkan kota memperoleh dukungan skala
penduduk dari sistem perkotaan yang lebih luas. Struktur polisentris
memperluas peluang itu, tetapi posisi kota dalam jaringan tetap tidak
setara. Kedekatan fisik juga membawa dua proses yang berlawanan, yakni
kompetisi fungsi dan kerja sama ekonomi (bagian 2.2-3.2, hlm.
1164-1166).

\textbf{Inferensi pengolahan.} Temuan empiris menunjukkan asosiasi
kondisional berdasarkan residual dan regresi, bukan pengamatan langsung
atas perpindahan skala penduduk ataupun identifikasi kausal yang
eksperimental. Karena itu, istilah ``mekanisme'' dalam artikel sebaiknya
dibaca sebagai mekanisme yang konsisten dengan pola estimasi, bukan
sebagai mekanisme kausal yang telah dibuktikan secara tuntas (model
(1)-(6), hlm. 1158-1160).

\subsection{Problem Statement}\label{problem-statement-20}

\textbf{Laporan sumber.} Kajian yang ada dinilai belum memberi perhatian
memadai pada \emph{borrowed size} pada skala aglomerasi perkotaan.
Padahal, China sedang mengalami urbanisasi cepat, aglomerasi menjadi
bentuk utama urbanisasi baru, dan keterhubungan kota besar, menengah,
serta kecil terus menguat. Kondisi tersebut memunculkan pertanyaan
empiris apakah kota-kota dalam aglomerasi China benar-benar dapat
memperoleh fungsi atau kinerja yang melampaui apa yang diperkirakan dari
ukuran internalnya (abstrak dan pendahuluan, hlm. 1156-1157).

Masalah tersebut dijabarkan menjadi tiga ketidakpastian. Pertama, arah
\emph{borrowed size} tidak harus mengikuti gagasan awal Alonso bahwa
kota kecil meminjam dari kota besar; dalam konteks China, kota besar
yang berada di pusat aglomerasi mungkin justru lebih mampu memanfaatkan
skala kota kecil dan sekaligus menimbulkan \emph{agglomeration shadow}.
Kedua, belum jelas apakah \emph{borrowed size} lebih sering muncul dalam
struktur monosentris atau polisentris. Ketiga, belum jelas apakah
kedekatan geografis masih menjadi mekanisme utama atau telah digantikan
maupun diperkuat oleh keterhubungan jaringan di tengah peningkatan
transportasi dan kompresi ruang-waktu (pendahuluan, hlm. 1157).

\textbf{Interpretasi penulis.} Artikel menempatkan \emph{borrowed size}
dan \emph{agglomeration shadow} sebagai dua sisi eksternalitas jaringan.
Kota dapat memperoleh manfaat dari skala sistem, tetapi persaingan
antarkota juga dapat mengalihkan pengguna, fungsi, dan peluang
pertumbuhan ke kota yang lebih dominan (pendahuluan, hlm. 1157).

\subsection{Objectives}\label{objectives-20}

\textbf{Laporan sumber.} Penelitian memiliki tiga tujuan utama:

\begin{itemize}
\tightlist
\item
  menguji keberadaan \emph{borrowed size} pada 14 aglomerasi perkotaan
  China dari dimensi fungsi dan kinerja, sekaligus memeriksa kemungkinan
  \emph{agglomeration shadow} dan arah peminjaman antarkelas ukuran
  kota;
\item
  menentukan apakah \emph{borrowed size} lebih sering terjadi dalam
  aglomerasi berstruktur polisentris atau monosentris serta menguji
  heterogenitas pengaruh struktur tersebut pada kota kecil, menengah,
  dan besar; dan
\item
  menguji kedekatan geografis dan keterhubungan jaringan sebagai
  mekanisme pembentukan \emph{borrowed function} dan \emph{borrowed
  performance}, termasuk apakah keduanya bersifat substitutif atau
  komplementer dan apakah efeknya berbeda menurut ukuran kota.
\end{itemize}

Ketiga tujuan dinyatakan dalam rangkaian pertanyaan pendahuluan dan
diwujudkan berturut-turut melalui model (1)-(2), model (3)-(4), serta
model (5)-(6) (hlm. 1157-1159).

\subsection{Literature Survey}\label{literature-survey-20}

\textbf{Laporan sumber.} Survei pustaka artikel membangun landasan
berikut:

\begin{itemize}
\tightlist
\item
  Alonso {[}1{]} memperkenalkan \emph{borrowed size} pada 1973. Gagasan
  intinya ialah kota kecil dapat memanfaatkan ekonomi aglomerasi kota
  besar yang berdekatan, mempertahankan keuntungan ukuran kecil, dan
  menunjukkan karakter ekonomi yang biasanya diasosiasikan dengan kota
  besar. Contoh awal yang disebutkan mencakup kota kecil di kawasan
  metropolitan pesisir Atlantik serta kota-kota di Jerman dan
  Negara-Negara Rendah (pendahuluan, hlm. 1156).
\item
  Phelps, Fallon, dan Williams {[}2{]} dilaporkan menemukan bukti pada
  kawasan London Raya bahwa perusahaan kecil di kota pinggiran
  memperoleh keuntungan bukan dari ekonomi lokal saja, melainkan dari
  kota besar di dekatnya. Dijkstra, Garcilazo, dan McCann {[}3{]}
  menunjukkan kinerja pertumbuhan kota lapis kedua di Eropa, sedangkan
  Meijers dan Burger {[}4{]} serta Volgmann dan Rusche {[}5{]}
  menggunakan \emph{borrowed size} untuk menjelaskan kinerja kota kecil
  dan menengah di kawasan metropolitan Eropa (pendahuluan, hlm. 1156).
\item
  Literatur membedakan perspektif fungsi dan kinerja. Fungsi merujuk
  pada artefak, aktivitas, atau fasilitas khas kota besar yang dapat
  ditopang oleh basis pengguna dari luar kota; kinerja merujuk pada
  tingkat capaian yang lebih tinggi daripada yang diperkirakan dari
  ukuran kota sendiri. Artikel mengikuti perluasan konsep oleh Meijers
  dan Burger {[}4{]} serta metode residual yang juga digunakan Volgmann
  dan Rusche {[}5{]} (bagian 1.1, hlm. 1158).
\item
  \emph{Agglomeration shadow} digunakan untuk menerangkan sisi merugikan
  hubungan pusat-pinggiran. Tervo {[}8{]} melaporkan kemunculannya dalam
  perkembangan spasial jangka panjang Finlandia setelah Perang Dunia II.
  Burger dkk. {[}9{]} menemukan fasilitas budaya di Eropa barat laut
  terkonsentrasi secara tidak proporsional di kota terbesar. Meijers dan
  Burger {[}4{]} melaporkan bahwa kota menengah yang terisolasi
  menghadapi bayangan lebih kecil daripada kota dengan tetangga dekat,
  sedangkan Volgmann dan Rusche {[}5{]} menemukan sekitar sepertiga kota
  di Jerman menunjukkan bayangan, paling jelas pada kota menengah
  (pendahuluan, hlm. 1157).
\item
  Teori tempat sentral, sistem kota, dan geografi ekonomi baru
  menerangkan tarikan faktor produksi oleh pusat serta ambang jarak
  untuk menghindari efek isap. Sebaliknya, gagasan efek
  balik-penyebaran, polarisasi-\emph{trickle-down}, dan pusat-pinggiran
  yang dirujuk artikel menekankan bahwa aglomerasi dan difusi berubah
  menurut tahap perkembangan. Artikel karena itu meminta pembacaan yang
  tidak menganggap kedekatan dengan kota besar selalu menguntungkan atau
  selalu merugikan (pendahuluan, hlm. 1157).
\item
  Paradigma jaringan kota menyatakan bahwa eksternalitas dapat melampaui
  batas kota sehingga fungsi tingkat tinggi tidak lagi melekat secara
  deterministik pada ukuran kota. Penulis merujuk Capello {[}6{]},
  Batten {[}7{]}, ulasan eksternalitas jaringan {[}11{]}, dan studi
  Huang, Hong, dan Ma {[}12{]} untuk mengajukan bahwa koneksi jaringan
  dapat menggantikan manfaat kedekatan fisik (pendahuluan, hlm.
  1156-1157).
\item
  Untuk China, penelitian sebelumnya disebut terutama berfokus pada
  Beijing-Tianjin-Hebei dan Delta Sungai Yangtze {[}13-14{]} serta pada
  efek ekonomi \emph{borrowed size} {[}15{]}. Artikel ini menggunakan
  kekurangan cakupan wilayah, kondisi spasial, mekanisme, dan
  heterogenitas ukuran kota tersebut sebagai dasar perluasan kajian
  (akhir pendahuluan, hlm. 1158).
\end{itemize}

\textbf{Inferensi pengolahan.} Bagian ini merupakan survei naratif untuk
membangun model dan pertanyaan penelitian. File tidak menyebutkan
strategi pencarian, kriteria inklusi, jumlah literatur yang disaring,
atau prosedur telaah sistematis; oleh sebab itu, bagian tersebut tidak
dapat diperlakukan sebagai tinjauan sistematis (pendahuluan dan daftar
pustaka, hlm. 1156-1158 dan 1166-1167).

\subsection{Method Used}\label{method-used-20}

\textbf{Laporan sumber.} Penelitian menjalankan tiga tahap estimasi:

\begin{itemize}
\tightlist
\item
  \textbf{Identifikasi \emph{borrowed function} dan \emph{borrowed
  performance}.} Model (1) meregresikan indeks fungsi kota
  (\texttt{FUN}) terhadap ukuran kota (\texttt{SIZ}), jumlah penduduk
  kota-kota lain dalam aglomerasi (\texttt{RSI}), dan rerata fungsi kota
  lain (\texttt{RFU}), dengan efek tetap kota dan tahun. Model (2)
  meregresikan kinerja kota (\texttt{PER}) terhadap \texttt{SIZ},
  \texttt{SIZ\^{}2}, \texttt{RSI}, dan rerata kinerja kota lain
  (\texttt{RPE}), juga dengan efek tetap kota dan tahun. Kuadrat ukuran
  dimasukkan untuk menangkap kemungkinan pertukaran antara efek skala
  dan kemacetan. Residual positif model (1) berarti fungsi aktual
  melebihi fungsi yang diperkirakan dan diklasifikasikan sebagai
  \emph{borrowed function}; residual positif model (2) berarti kinerja
  aktual melebihi kinerja yang diperkirakan dan diklasifikasikan sebagai
  \emph{borrowed performance}. Kedua residual positif sekaligus
  didefinisikan sebagai \emph{borrowed size}, sedangkan keduanya negatif
  sekaligus didefinisikan sebagai \emph{agglomeration shadow} (bagian
  1.1 dan Gambar 1, hlm. 1158-1159). Untuk membaca hasil pada tingkat
  aglomerasi, penulis merata-ratakan residual kota-kota anggota (bagian
  2.1, hlm. 1161).
\item
  \textbf{Pengujian kondisi struktur spasial.} Residual \emph{borrowed
  function} dan \emph{borrowed performance} menjadi variabel hasil dalam
  model (3) dan (4), dengan indeks struktur polisentris (\texttt{POL})
  sebagai variabel penjelas utama, seperangkat kontrol, efek tetap kota,
  dan efek tetap tahun. \texttt{POL} dihitung setiap tahun untuk setiap
  aglomerasi melalui koefisien Pareto dari regresi peringkat-ukuran.
  Koefisien kurang dari 1 menunjukkan kecenderungan monosentris, lebih
  dari 1 menunjukkan kecenderungan polisentris, dan sama dengan 1 sesuai
  dengan hukum Zipf (bagian 1.1-1.2, hlm. 1158-1159).
\item
  \textbf{Pengujian mekanisme kedekatan dan jaringan.} Model (5) dan (6)
  menggunakan residual fungsi dan kinerja sebagai variabel hasil serta
  kedekatan geografis (\texttt{PRO}), keterhubungan jaringan
  (\texttt{NET}), interaksi \texttt{NET\ x\ PRO}, dan kontrol sebagai
  variabel penjelas. Tanda interaksi yang sama dengan tanda efek
  kedekatan ditafsirkan sebagai komplementaritas; tanda yang berlawanan
  ditafsirkan sebagai substitusi. Karena data \texttt{PRO} dan
  \texttt{NET} yang berubah dari waktu ke waktu sulit diperoleh, tahap
  ini memakai penampang lintang 2019 (bagian 1.1 dan 1.3, hlm.
  1158-1160).
\end{itemize}

Indeks fungsi kota dibentuk dari sembilan indikator positif dalam tiga
kelompok. Fungsi komersial mencakup pangsa aset perusahaan A-share
Shanghai-Shenzhen, pangsa penjualan eceran barang konsumsi, dan pangsa
kumulatif lahan komersial/jasa dalam aglomerasi. Fungsi sains-teknologi
mencakup pangsa belanja fiskal iptek, pangsa jumlah dosen perguruan
tinggi, dan pangsa permohonan paten. Fungsi budaya-olahraga mencakup
pangsa koleksi perpustakaan umum, pangsa jumlah objek wisata kelas 5A,
serta pangsa jumlah medali emas Olimpiade Musim Panas dan Musim Dingin.
Penulis memberi bobot objektif dengan metode deviasi standar untuk
menghasilkan indeks komposit (bagian 1.2 dan Tabel 1, hlm. 1159-1160).
Kinerja kota diukur dengan PDB per kapita (bagian 1.2, hlm. 1159).

Kedekatan geografis dioperasionalkan sebagai jumlah kota lain dalam
aglomerasi yang berbatasan langsung dengan suatu kota. Keterhubungan
jaringan dioperasionalkan sebagai rerata kebalikan waktu perjalanan
kereta api antara kota tersebut dan kota-kota lain dalam aglomerasi;
waktu yang lebih singkat menghasilkan nilai koneksi lebih tinggi.
Variabel kontrol mencakup rasio investasi langsung asing terhadap PDB,
rasio nilai tambah sektor tersier terhadap sektor sekunder, rasio
mahasiswa perguruan tinggi terhadap jumlah penduduk akhir tahun, dan
tingkat administratif. Tingkat administratif dikodekan 1 untuk kota di
atas tingkat prefektur dan 0 untuk kota tingkat prefektur biasa, lalu
diinteraksikan dengan tren waktu linear (bagian 1.2, hlm. 1159-1160).

Untuk analisis heterogenitas, kota dikelompokkan berdasarkan standar
Dewan Negara China 2014: kota kecil berpenduduk kurang dari 500.000,
kota menengah 500.000-1.000.000, dan kota besar lebih dari 1.000.000
penduduk (bagian 2.2, hlm. 1164).

\textbf{Inferensi pengolahan.} Rancangan ini mendeteksi kelebihan fungsi
atau kinerja relatif terhadap prediksi model. Ia tidak mengukur arus
peminjaman penduduk secara langsung. Model efek tetap memperkuat
pengendalian atas heterogenitas kota yang tidak berubah dan guncangan
waktu bersama, tetapi file tidak melaporkan desain instrumen, eksperimen
alam, atau prosedur lain untuk mengidentifikasi kausalitas (bagian 1.1,
hlm. 1158-1159).

\subsection{Dataset}\label{dataset-20}

\textbf{Laporan sumber.} Model (1)-(4) menggunakan panel 2008-2019.
Model (5)-(6) menggunakan penampang lintang 2019 karena data kedekatan
dan jaringan yang berubah dari waktu ke waktu dinyatakan sulit diperoleh
(bagian 1.3, hlm. 1160). Data dasarnya berasal dari:

\begin{itemize}
\tightlist
\item
  \emph{China City Statistical Yearbook} untuk sebagian besar variabel;
\item
  perhitungan tidak langsung penduduk tetap wilayah perkotaan dari PDB
  dan PDB per kapita;
\item
  China Land Market Network untuk luas lahan komersial dan jasa yang
  ditawarkan sejak 2000 dan dijumlahkan secara kumulatif per tahun;
\item
  basis data Guotai'an untuk aset perusahaan A-share Shanghai dan
  Shenzhen;
\item
  Chinese Research Data Services Platform untuk jumlah permohonan paten;
  dan
\item
  kompilasi manual dari Baidu Baike, Baidu Maps, dan situs 12306 untuk
  jumlah objek wisata 5A, medali emas Olimpiade, jumlah kota yang
  berbatasan, serta waktu perjalanan kereta api antarkota.
\end{itemize}

PDB, penjualan eceran barang konsumsi, belanja fiskal iptek, dan
indikator nominal lain dikonversi menjadi nilai riil dengan tahun dasar
2001 menggunakan indeks harga yang relevan (bagian 1.3, hlm. 1160). File
tidak menyebutkan nama atau bentuk indeks harga untuk masing-masing
seri, prosedur penanganan data hilang, kriteria pembersihan data, maupun
akses ke data replikasi.

\subsection{Population / Sample}\label{population-sample-20}

\textbf{Laporan sumber.} Populasi sasaran adalah kota setingkat
prefektur atau lebih tinggi di dalam 14 aglomerasi perkotaan: Delta
Sungai Mutiara, Delta Sungai Yangtze, Beijing-Tianjin-Hebei,
Harbin-Changchun, Semenanjung Shandong, Liaoning Tengah dan Selatan,
Chengdu-Chongqing, Wuhan, Chang-Zhu-Tan, Danau Poyang, Dataran Tengah,
Teluk Beibu, Pesisir Guangdong-Fujian-Zhejiang, dan Guanzhong. Batas
setiap aglomerasi disesuaikan dengan rencana pengembangan aglomerasi
yang diterbitkan pemerintah pusat atau pemerintah daerah (bagian 1.3,
hlm. 1160).

Tabel 2 dan 3 melaporkan 2.568 observasi kota-tahun untuk panel
2008-2019. Tabel 4 melaporkan 214 observasi untuk penampang lintang
2019. Dalam regresi struktur polisentris menurut ukuran, terdapat 233
observasi kota-tahun kecil, 719 menengah, dan 1.616 besar. Dalam regresi
penampang lintang 2019 menurut ukuran, terdapat 10 kota kecil, 52 kota
menengah, dan 152 kota besar (Tabel 2-4, hlm. 1160 dan 1164).

\textbf{Inferensi pengolahan.} Angka 214 pada penampang lintang dan
2.568 pada panel konsisten dengan 214 kota dikali 12 tahun. Namun, file
tidak menyatakan secara eksplisit bahwa panel seimbang, tidak memberikan
daftar lengkap 214 kota dalam teks yang terekstraksi, dan tidak
menjelaskan penanganan kota yang batas administratif atau kelas
ukurannya berubah. Oleh sebab itu, konsistensi hitungan tersebut tidak
boleh menggantikan dokumentasi sampel yang tidak tersedia.

\subsection{Dependent Variables}\label{dependent-variables-20}

\textbf{Laporan sumber.} Variabel dependen berubah menurut tahap
analisis:

\begin{itemize}
\tightlist
\item
  Dalam model (1), variabel dependen adalah indeks komposit fungsi kota
  (\texttt{FUN}). Indeks ini merangkum fungsi komersial,
  sains-teknologi, serta budaya-olahraga melalui sembilan indikator pada
  Tabel 1 dan pembobotan metode deviasi standar (bagian 1.1-1.2 dan
  Tabel 1, hlm. 1158-1160).
\item
  Dalam model (2), variabel dependen adalah kinerja kota (\texttt{PER})
  yang diproksikan dengan PDB per kapita (bagian 1.1-1.2, hlm.
  1158-1159).
\item
  Dalam model (3) dan (5), variabel dependen substantif adalah residual
  model fungsi, yaitu indeks \emph{borrowed function}. Nilai residual di
  atas nol berarti fungsi aktual lebih tinggi daripada fungsi yang
  diprediksi berdasarkan spesifikasi model; nilai di bawah nol berarti
  tidak ada \emph{borrowed function} menurut aturan klasifikasi artikel
  (bagian 1.1 dan Gambar 1, hlm. 1158-1159).
\item
  Dalam model (4) dan (6), variabel dependen substantif adalah residual
  model kinerja, yaitu indeks \emph{borrowed performance}. Nilai di atas
  nol berarti kinerja aktual lebih tinggi daripada kinerja yang
  diprediksi; nilai di bawah nol berarti tidak ada \emph{borrowed
  performance} menurut aturan artikel (bagian 1.1 dan Gambar 1, hlm.
  1158-1159).
\end{itemize}

Kombinasi residual fungsi dan kinerja menghasilkan klasifikasi utama:
keduanya positif berarti \emph{borrowed size}, sedangkan keduanya
negatif berarti \emph{agglomeration shadow}. Kombinasi satu positif dan
satu negatif menunjukkan hanya salah satu dimensi yang dipinjam (Gambar
1, hlm. 1159).

\textbf{Catatan terminologis.} Tabel 3 dan 4 memberi label kolom
\texttt{FUN} dan \texttt{PER}, tetapi uraian model menyatakan bahwa yang
diregresikan adalah residual fungsi dan residual kinerja. Review ini
mengikuti definisi substantif dalam model, yakni indeks \emph{borrowed
function} dan \emph{borrowed performance}, tanpa menganggap label tabel
sebagai ukuran mentah kedua variabel.

\subsection{Results}\label{results-20}

\textbf{Laporan sumber: estimasi dasar.} Pada model fungsi, ukuran kota
berkoefisien positif (\texttt{0,1778}, \texttt{t\ =\ 3,76}, signifikan
1\%), sedangkan jumlah penduduk kota lain dalam aglomerasi berkoefisien
negatif (\texttt{-0,0463}, \texttt{t\ =\ -6,20}, signifikan 1\%) dan
rerata fungsi kota lain juga negatif (\texttt{-4,3710},
\texttt{t\ =\ -38,18}, signifikan 1\%). Pada model kinerja, ukuran kota
positif (\texttt{35,2497}, \texttt{t\ =\ 16,82}, signifikan 1\%) dan
kuadratnya negatif (\texttt{-90,1004}, \texttt{t\ =\ -11,71}, signifikan
1\%), sehingga hubungan yang diestimasi berbentuk U terbalik. Jumlah
penduduk kota lain (\texttt{0,9530}, \texttt{t\ =\ 4,65}) dan rerata
kinerja kota lain (\texttt{0,5792}, \texttt{t\ =\ 14,11}) sama-sama
positif dan signifikan 1\%. Nilai \texttt{R\^{}2} masing-masing model
adalah \texttt{0,3886} dan \texttt{0,2764}, dengan \texttt{N\ =\ 2.568}
serta efek tetap kota dan tahun (Tabel 2, hlm. 1160).

\textbf{Interpretasi penulis atas estimasi dasar.} Koefisien fungsi
ditafsirkan sebagai konsentrasi fungsi pada kota berpenduduk besar,
adanya ambang pengguna, dan kompetisi fungsi antarkota. Bentuk U
terbalik pada kinerja ditafsirkan sebagai perimbangan manfaat aglomerasi
dengan biaya kemacetan. Koefisien positif \texttt{RSI} dan \texttt{RPE}
pada model kinerja dipandang konsisten dengan peminjaman skala penduduk
dan limpahan pertumbuhan ekonomi dari lingkungan aglomerasi yang lebih
maju (bagian 2.1, hlm. 1160-1161).

\textbf{Laporan sumber: keberadaan pada tingkat aglomerasi.} Selama
2008-2019, tren rerata \emph{borrowed function} relatif datar. Sembilan
aglomerasi selalu menunjukkan ciri \emph{borrowed function}:
Chang-Zhu-Tan, Delta Sungai Mutiara, Wuhan, Liaoning Tengah dan Selatan,
Harbin-Changchun, Teluk Beibu, Danau Poyang, Guanzhong, serta
Beijing-Tianjin-Hebei. Untuk \emph{borrowed performance}, Wuhan, Pesisir
Guangdong-Fujian-Zhejiang, dan Teluk Beibu mengalami kenaikan yang
jelas. Pada 2019, sembilan aglomerasi menunjukkan \emph{borrowed
performance}: Delta Sungai Mutiara, Beijing-Tianjin-Hebei,
Chang-Zhu-Tan, Wuhan, Liaoning Tengah dan Selatan, Semenanjung Shandong,
Pesisir Guangdong-Fujian-Zhejiang, Teluk Beibu, serta Harbin-Changchun.
Irisan kedua dimensi menghasilkan tujuh aglomerasi dengan \emph{borrowed
size}: Chang-Zhu-Tan, Delta Sungai Mutiara, Wuhan, Liaoning Tengah dan
Selatan, Harbin-Changchun, Teluk Beibu, serta Beijing-Tianjin-Hebei
(bagian 2.1 dan Gambar 2, hlm. 1161-1162).

\textbf{Laporan sumber: pola tingkat kota.} Pada 2019, kota yang disebut
paling jelas berada pada kuadran \emph{borrowed function} dan
\emph{borrowed performance} sekaligus ialah Xi'an, Wuhan, Changsha,
Nanchang, Shenzhen, Beijing, Dalian, Nanning, Zhengzhou, Qingdao, dan
Fuzhou. Semuanya dinyatakan sebagai kota utama dalam aglomerasi
masing-masing. Kota besar cenderung berada pada kelompok yang meminjam
kinerja, sedangkan kota kecil dan menengah umumnya memiliki
\emph{borrowed function} tanpa \emph{borrowed performance}. Proporsi
kota yang memiliki kedua residual negatif dan karena itu
diklasifikasikan berada dalam \emph{agglomeration shadow} mencapai 81\%
di Chengdu-Chongqing, 93\% di Dataran Tengah, 65\% di Semenanjung
Shandong, 65\% di Pesisir Guangdong-Fujian-Zhejiang, dan 66\% di Delta
Sungai Yangtze (bagian 2.1 dan Gambar 3, hlm. 1162-1164).

\textbf{Laporan sumber: struktur polisentris.} Untuk sampel penuh,
\texttt{POL} berasosiasi positif dengan \emph{borrowed function}
(\texttt{0,0192}, \texttt{t\ =\ 7,28}, signifikan 1\%) dan
\emph{borrowed performance} (\texttt{0,0275}, \texttt{t\ =\ 3,39},
signifikan 1\%). Pada fungsi, koefisien positif muncul untuk kota kecil
(\texttt{0,0128}, signifikan 10\%), menengah (\texttt{0,0397},
signifikan 1\%), dan besar (\texttt{0,0177}, signifikan 1\%). Pada
kinerja, koefisien untuk kota kecil (\texttt{0,2091},
\texttt{t\ =\ 0,76}) tidak signifikan, sedangkan kota menengah
(\texttt{0,4764}, signifikan 1\%) dan besar (\texttt{0,2029}, signifikan
5\%) positif signifikan. Koefisien terbesar untuk kedua dimensi terdapat
pada kelompok kota menengah. Nilai \texttt{R\^{}2} sampel penuh hanya
\texttt{0,0380} untuk fungsi dan \texttt{0,0450} untuk kinerja; model
kelompok berkisar \texttt{0,0271-0,2470} (Tabel 3 dan bagian 2.2, hlm.
1164).

\textbf{Laporan sumber: kedekatan dan jaringan.} Pada sampel penuh 2019,
kedekatan geografis (\texttt{PRO}) berasosiasi negatif dengan
\emph{borrowed function} (\texttt{-0,0520}, signifikan 1\%) tetapi
positif dengan \emph{borrowed performance} (\texttt{0,0199}, signifikan
5\%). Keterhubungan jaringan (\texttt{NET}) positif untuk fungsi
(\texttt{0,1728}, signifikan 5\%) dan kinerja (\texttt{1,0089},
signifikan 10\%). Interaksi \texttt{NET\ x\ PRO} positif pada fungsi
(\texttt{0,0093}, signifikan 5\%) dan kinerja (\texttt{0,1048},
signifikan 5\%). Karena efek utama \texttt{PRO} pada fungsi negatif
sementara interaksi positif, penulis mengidentifikasi substitusi:
jaringan mengurangi efek negatif kedekatan. Karena efek utama dan
interaksi pada kinerja sama-sama positif, penulis mengidentifikasi
komplementaritas (Tabel 4 dan bagian 2.3, hlm. 1164-1165).

\textbf{Laporan sumber: heterogenitas kedekatan dan jaringan.}
\texttt{PRO} negatif signifikan untuk \emph{borrowed function} di semua
kelas ukuran, tetapi positif signifikan untuk \emph{borrowed
performance} di semua kelas. \texttt{NET} justru negatif signifikan
terhadap fungsi kota kecil (\texttt{-0,5910}, signifikan 5\%), positif
tetapi tidak signifikan untuk fungsi kota menengah (\texttt{0,1998}),
dan positif signifikan untuk fungsi kota besar (\texttt{0,2535},
signifikan 5\%). Pada kinerja, efek \texttt{NET} tidak signifikan untuk
kota kecil (\texttt{0,9218}) dan menengah (\texttt{-0,2782}), tetapi
positif signifikan untuk kota besar (\texttt{2,1584}, signifikan 1\%).
Interaksi pada fungsi positif signifikan untuk ketiga kelas; pada
kinerja, interaksi hanya positif signifikan pada kota besar
(\texttt{0,3264}, signifikan 10\%). Berdasarkan pola ini, penulis
menyimpulkan bahwa daya promosi langsung jaringan terhadap kedua dimensi
hanya tampak pada kota besar (Tabel 4 dan bagian 2.3, hlm. 1164-1165).

\subsection{Findings}\label{findings-20}

\textbf{Temuan yang dilaporkan sumber.} Bukti utama bukan bahwa semua
aglomerasi China memiliki \emph{borrowed size}, melainkan bahwa
keberadaannya selektif menurut wilayah, dimensi, dan ukuran kota. Tujuh
dari 14 aglomerasi memenuhi dua syarat residual positif pada tingkat
agregat 2019. Sementara itu, lima aglomerasi memiliki proporsi besar
kota dengan dua residual negatif. Dengan demikian, \emph{borrowed size}
dan \emph{agglomeration shadow} dapat muncul dalam sistem perkotaan
nasional yang sama (bagian 2.1, hlm. 1161-1164).

\textbf{Temuan yang dilaporkan sumber.} Posisi dalam hierarki kota
penting. Kota-kota utama lebih sering memperoleh fungsi dan kinerja
sekaligus; kota kecil serta menengah lebih sering memperoleh fungsi
tanpa kenaikan kinerja yang sebanding (bagian 2.1 dan 3.1, hlm.
1163-1165).

\textbf{Inferensi pengolahan.} Pola tersebut berbeda dari pembacaan satu
arah atas konsep awal yang hanya menempatkan kota kecil sebagai peminjam
dari kota besar. Artikel sendiri menjadikan kemungkinan arah yang
berlawanan sebagai pertanyaan empiris, tetapi perbandingan ini merupakan
pembacaan konseptual atas hasil, bukan koefisien tersendiri
(pendahuluan, hlm. 1157; bagian 2.1, hlm. 1163-1164).

\textbf{Interpretasi penulis.} Penulis menganggap adanya perbedaan
fungsi atau kinerja di dalam aglomerasi sebagai prasyarat bagi
peminjaman: aglomerasi dengan varians fungsi atau kinerja lebih kecil
cenderung memiliki tingkat peminjaman yang lebih rendah. Pada kota kecil
dan menengah, basis industri yang relatif lemah dinilai membatasi
kemampuan menarik penduduk untuk mendukung kinerja ekonomi, meskipun
diferensiasi fungsi masih memungkinkan mereka mengembangkan fungsi yang
memiliki keunggulan komparatif (bagian 2.1, hlm. 1162-1164).

\textbf{Temuan yang dilaporkan sumber.} Struktur polisentris memperbesar
peluang \emph{borrowed function} pada seluruh kelas ukuran, tetapi
manfaat kinerja terukur hanya signifikan pada kota menengah dan besar.
Kota menengah memperoleh koefisien terbesar pada kedua dimensi (Tabel 3,
hlm. 1164).

\textbf{Interpretasi penulis.} Struktur polisentris dipandang
menyebarkan faktor produksi secara lebih merata, memperdalam pembagian
kerja, dan mempererat hubungan pasar. Kota menengah dianggap paling
mungkin berkembang menjadi simpul baru sehingga paling banyak memperoleh
manfaat. Kota kecil yang berada di tepian jaringan dinilai belum mampu
mengubah dukungan struktur polisentris menjadi kinerja ekonomi (bagian
2.2, hlm. 1164).

\textbf{Temuan yang dilaporkan sumber.} Kedekatan fisik memiliki arah
yang berbeda pada fungsi dan kinerja, sedangkan efek langsung positif
jaringan tidak merata antarkelas kota. Hubungan jaringan mengurangi
hambatan kedekatan terhadap fungsi dan memperkuat manfaat kedekatan
terhadap kinerja pada sampel penuh (Tabel 4, hlm. 1164).

\textbf{Interpretasi penulis.} Banyak tetangga berbatasan meningkatkan
kompetisi memperebutkan basis pengguna fungsi, sehingga ``isolasi''
dapat melindungi fungsi lokal. Namun, kedekatan juga memungkinkan
pembagian kerja industri dan kerja sama ekonomi yang menopang kinerja.
Jaringan kereta memperluas jangkauan peminjaman melampaui tetangga
langsung. Dominasi dan aksesibilitas kota besar membuat manfaat langsung
jaringan terutama terkonsentrasi pada kota besar, sementara kota kecil
dan menengah masih dibatasi ruang geografis (bagian 2.3, hlm. 1165).

\textbf{Inferensi pengolahan.} Koefisien jaringan kota kecil yang
negatif signifikan untuk fungsi memperumit pernyataan umum bahwa
jaringan selalu meningkatkan kedua dimensi. Rumusan yang paling sesuai
dengan seluruh hasil ialah: asosiasi positif terdapat pada sampel penuh,
tetapi arah dan signifikansinya heterogen; efek promosi langsung yang
konsisten hanya ditemukan pada kota besar (Tabel 4 dan bagian 2.3, hlm.
1164-1165).

\subsection{Conclusions}\label{conclusions-20}

\textbf{Kesimpulan penulis.} Pertama, tujuh aglomerasi memiliki
\emph{borrowed size} pada tingkat agregat, kota besar lebih mungkin
meminjam fungsi dan kinerja sekaligus, kota kecil serta menengah lebih
mungkin hanya meminjam fungsi, dan lima aglomerasi menunjukkan
\emph{agglomeration shadow} (bagian 3.1, hlm. 1165).

Kedua, \emph{borrowed size} lebih sering terjadi pada aglomerasi
polisentris. Semua kelas ukuran dapat memperoleh \emph{borrowed
function} dalam struktur tersebut, tetapi hanya kota menengah dan besar
yang memperoleh \emph{borrowed performance} secara signifikan. Pengaruh
terbesar terdapat pada kota menengah (bagian 3.1, hlm. 1165).

Ketiga, kedekatan geografis menurunkan \emph{borrowed function} tetapi
meningkatkan \emph{borrowed performance}. Keterhubungan jaringan
meningkatkan keduanya pada sampel penuh; jaringan dan kedekatan bersifat
substitutif untuk fungsi serta komplementer untuk kinerja. Dalam
pengujian menurut ukuran, efek promosi langsung jaringan pada kedua
hasil hanya terlihat pada kota besar (bagian 3.1, hlm. 1165).

\textbf{Batas interpretasi.} Kesimpulan tersebut berlaku bagi kota
setingkat prefektur atau lebih tinggi dalam 14 batas aglomerasi resmi
dan periode yang dianalisis. File tidak menyediakan dasar untuk
menggeneralisasikannya ke semua kota China, kota tingkat county, negara
lain, atau periode setelah 2019 (bagian 1.3, hlm. 1160; akhir bagian
3.2, hlm. 1166-1167).

\subsection{Contributions}\label{contributions-20}

\textbf{Kontribusi yang diklaim penulis.} Artikel memperluas penelitian
dari konsentrasi sebelumnya pada Beijing-Tianjin-Hebei dan Delta Sungai
Yangtze menjadi 14 aglomerasi dengan tingkat perkembangan yang beragam,
sehingga hasilnya dimaksudkan turut memberi rujukan bagi pengaturan
spasial aglomerasi yang kurang berkembang (akhir pendahuluan, hlm.
1158).

Artikel juga memperluas fokus dari efek ekonomi \emph{borrowed size}
menuju dua pertanyaan pembentukan: struktur polisentris sebagai kondisi
spasial serta kedekatan geografis dan keterhubungan jaringan sebagai
mekanisme. Analisis fungsi dan kinerja ditempatkan dalam satu kerangka
dengan \emph{agglomeration shadow}, bukan diperlakukan sebagai fenomena
terpisah (pendahuluan dan bagian 1.1, hlm. 1157-1159).

Kontribusi ketiga ialah analisis heterogenitas berdasarkan kota kecil,
menengah, dan besar. Hasilnya menunjukkan bahwa kesempatan memperoleh
fungsi tidak identik dengan kesempatan memperoleh kinerja dan bahwa
manfaat jaringan tidak terbagi merata menurut ukuran kota (akhir
pendahuluan, hlm. 1158; bagian 2.2-2.3, hlm. 1164-1165).

\textbf{Inferensi pengolahan.} Nilai tambah analitis terpenting terletak
pada pemisahan dua dimensi hasil. Tanpa pemisahan tersebut, efek
kompetitif kedekatan pada fungsi dapat tertutup oleh efek kooperatifnya
pada kinerja. Namun, kontribusi ini tetap bergantung pada validitas
residual sebagai indikator \emph{borrowed size} dan pada proksi spasial
serta jaringan yang digunakan (bagian 1.1-1.2, hlm. 1158-1159; Tabel 4,
hlm. 1164).

\subsection{Research Gap}\label{research-gap-20}

\textbf{Kesenjangan yang dinyatakan sumber.} Artikel menyatakan satu
kesenjangan umum dan tiga perluasan khusus yang ingin diisi:

\begin{itemize}
\tightlist
\item
  perhatian terhadap \emph{borrowed size} pada skala aglomerasi
  perkotaan masih kurang;
\item
  bukti China sebelumnya terutama berasal dari wilayah maju, khususnya
  Beijing-Tianjin-Hebei dan Delta Sungai Yangtze, sehingga variasi
  lintas 14 aglomerasi belum cukup dikaji;
\item
  penelitian lebih banyak berfokus pada dampak ekonomi daripada kondisi
  spasial dan mekanisme terbentuknya \emph{borrowed size}; dan
\item
  perbedaan kapasitas peminjaman menurut ukuran kota belum cukup
  diperhitungkan.
\end{itemize}

Kesenjangan tersebut dinyatakan dalam abstrak dan akhir pendahuluan
(hlm. 1156 dan 1158). Penulis juga menyebut bahwa metode untuk menguji
keberadaan \emph{borrowed size} masih terbatas, lalu memilih metode
residual karena dinilai memiliki daya penjelas yang baik dan telah
diterima dalam studi terdahulu (bagian 1.1, hlm. 1158).

\textbf{Kesenjangan yang masih tersisa menurut penulis.} Pengukuran
fungsi belum mencakup fungsi administratif secara memadai; dimensi
kedekatan dan jaringan yang berubah dari waktu ke waktu belum tersedia
secara memadai; serta cakupan spasial masih mengikuti batas resmi dan
belum memasukkan kota tingkat county (akhir bagian 3.2, hlm. 1166-1167).

\textbf{Inferensi pengolahan.} Penelitian ini belum menutup kesenjangan
antara indikator residual dan proses peminjaman aktual. Masih diperlukan
bukti arus penduduk, pengguna, pengetahuan, perdagangan, atau faktor
produksi yang dapat menunjukkan dari kota mana skala dipinjam dan
melalui relasi apa. Kebutuhan tersebut merupakan inferensi dari
operasionalisasi model, bukan agenda yang dirumuskan dengan kata-kata
yang sama oleh penulis (bagian 1.1-1.2, hlm. 1158-1159).

\subsection{Limitations}\label{limitations-20}

\textbf{Keterbatasan yang dinyatakan penulis.} Pertama, data panel untuk
beberapa fungsi, termasuk fungsi administratif, langka sehingga sistem
indikator fungsi kota masih dapat dioptimalkan. Kedua, data kedekatan
dan jaringan yang berubah dari waktu ke waktu juga langka; akibatnya,
pengujian kedua mekanisme hanya menggunakan data 2019. Ketiga, wilayah
studi mengikuti batas aglomerasi dalam dokumen resmi dan hanya mencakup
kota setingkat prefektur atau lebih tinggi; penulis mengakui perlunya
delineasi aglomerasi yang lebih ilmiah dan pengujian kota tingkat county
(akhir bagian 3.2, hlm. 1166-1167).

\textbf{Keterbatasan berdasarkan informasi yang dilaporkan sumber.}
Kinerja kota direpresentasikan hanya oleh PDB per kapita, sedangkan
jaringan direpresentasikan hanya oleh waktu perjalanan kereta api dan
kedekatan hanya oleh jumlah kota yang berbatasan. Proksi tersebut
dinyatakan dalam bagian metode, tetapi file tidak melaporkan uji dengan
ukuran kinerja, ukuran jaringan, atau definisi kedekatan alternatif
(bagian 1.2, hlm. 1159).

\textbf{Inferensi pengolahan.} Klasifikasi \emph{borrowed size}
bergantung pada tanda residual relatif terhadap nol. File tidak
melaporkan interval ketidakpastian untuk klasifikasi tiap kota atau uji
kepekaan terhadap spesifikasi model, sehingga kota dengan residual dekat
nol berpotensi berpindah kategori jika model berubah. Selain itu,
residual positif menunjukkan hasil di atas prediksi, bukan dengan
sendirinya membuktikan bahwa kelebihan tersebut berasal dari skala kota
lain (bagian 1.1 dan Gambar 1, hlm. 1158-1159).

\textbf{Inferensi pengolahan.} Analisis mekanisme adalah penampang
lintang satu tahun dan tidak menyediakan urutan waktu antara
konektivitas dan hasil. Hubungan dua arah tetap mungkin terjadi,
misalnya kota berkinerja tinggi dapat memperoleh jaringan kereta yang
lebih baik. File tidak menyebutkan instrumen atau strategi identifikasi
lain untuk mengatasi kemungkinan tersebut (bagian 1.1 dan 1.3, hlm.
1158-1160).

\textbf{Inferensi pengolahan.} Kelompok kota kecil pada regresi 2019
hanya berisi 10 observasi, jauh lebih sedikit daripada 52 kota menengah
dan 152 kota besar. Estimasi serta ketidaksignifikanan pada kelompok
kecil karena itu perlu dibaca dengan hati-hati. Nilai \texttt{R\^{}2}
model struktur polisentris juga rendah pada sampel penuh
(\texttt{0,0380} dan \texttt{0,0450}), meskipun rendahnya daya penjelas
tidak meniadakan signifikansi koefisien \texttt{POL} (Tabel 3-4, hlm.
1164).

File tidak menyebutkan dalam teks yang tersedia: prosedur khusus untuk
data hilang dan pencilan, bentuk galat baku atau pengelompokan galat
baku, uji diagnostik model, uji ketahanan alternatif, serta ketersediaan
kode dan data replikasi.

\subsection{Challenges}\label{challenges-20}

Tidak terdapat bagian khusus bernama tantangan dalam file. Tiga kelompok
berikut merupakan pengelompokan review atas masalah yang dilaporkan
sumber, bukan kategori formal dari penulis:

\begin{itemize}
\tightlist
\item
  \textbf{Tantangan data:} panel untuk fungsi administratif serta ukuran
  kedekatan dan jaringan yang berubah dari waktu ke waktu sulit
  diperoleh. Hal ini membatasi indeks fungsi dan memaksa analisis
  mekanisme menggunakan satu tahun (bagian 1.3 dan akhir bagian 3.2,
  hlm. 1160 dan 1166-1167).
\item
  \textbf{Tantangan distribusi manfaat:} fungsi tingkat tinggi masih
  terkonsentrasi pada kota inti, sedangkan kota kecil dan menengah
  memiliki kelemahan dalam kualitas tenaga kerja, akumulasi modal, dan
  inovasi. Kondisi ini membuat mereka rentan berpartisipasi secara pasif
  dalam pembagian kerja dan jatuh dalam \emph{agglomeration shadow}
  (bagian 3.2, hlm. 1166).
\item
  \textbf{Tantangan integrasi spasial:} struktur polisentris dan
  jaringan tidak otomatis menghasilkan manfaat yang sama bagi semua
  kota. Temuan bahwa pengaruh langsung positif jaringan hanya signifikan
  pada kota besar menunjukkan kesulitan kota kecil dan menengah untuk
  menembus kendala ruang serta memanfaatkan posisi jaringan (bagian 2.3,
  hlm. 1165).
\end{itemize}

\textbf{Inferensi pengolahan.} Tantangan implementasinya ialah mendorong
konektivitas tanpa memperkuat penyerapan fungsi oleh kota dominan.
Inferensi ini mengikuti perbedaan antara efek jaringan sampel penuh,
koefisien negatif fungsi pada kota kecil, dan keberadaan
\emph{agglomeration shadow}; file tidak menguji intervensi kebijakan
tertentu (Gambar 3 dan Tabel 4, hlm. 1163-1165).

\subsection{Practical Implications}\label{practical-implications-20}

\textbf{Rekomendasi penulis.} Pemerintah pada berbagai tingkat
disarankan menjadikan \emph{borrowed size} sebagai salah satu jalur
koordinasi aglomerasi, dengan tiga arah kebijakan (bagian 3.2, hlm.
1166):

\begin{itemize}
\tightlist
\item
  \textbf{Menata fungsi kota secara rasional.} Karena fungsi masih
  terkonsentrasi di kota inti, kota kecil dan menengah disarankan
  membangun fungsi tingkat tinggi berdasarkan keunggulan komparatif,
  menargetkan pengguna pada skala seluruh aglomerasi, dan menggunakan
  basis pengguna lintas kota untuk mencapai ambang layanan.
\item
  \textbf{Menjaga kinerja kota kecil dan menengah.} Untuk menghindari
  \emph{agglomeration shadow}, kota-kota ini disarankan menonjolkan
  kelebihan lingkungan perkotaan, biaya transformasi, dan potensi pasar
  agar lebih menarik serta dapat menggunakan skala penduduk aglomerasi
  untuk menciptakan kinerja ekonomi yang lebih tinggi.
\item
  \textbf{Mendorong strategi polisentris dan jaringan kota.} Aglomerasi
  yang masih sangat monosentris disarankan memperkuat efek penyebaran
  kota pusat, menaikkan kapasitas simpul sekunder, dan memasukkan kota
  kecil serta menengah ke dalam integrasi aglomerasi. Infrastruktur
  transportasi perlu diperkuat untuk menurunkan biaya pergerakan faktor
  produksi, memperdalam koneksi jaringan, dan membuka jalur
  \emph{borrowed size}.
\end{itemize}

\textbf{Kualifikasi dari hasil.} Rekomendasi konektivitas perlu dibaca
bersama heterogenitas hasil: jaringan berasosiasi positif pada sampel
penuh, tetapi efek promosi langsung terhadap fungsi dan kinerja hanya
signifikan pada kota besar, bahkan koefisien fungsi kota kecil negatif.
Artikel tetap menyatakan bahwa interaksi jaringan-kedekatan yang positif
bagi fungsi kota kecil dan menengah membuat penguatan jaringan penting
bagi kelompok tersebut (bagian 2.3, hlm. 1165).

\subsection{Applications}\label{applications-20}

\textbf{Aplikasi yang dinyatakan sumber.} Hasil dimaksudkan sebagai
rujukan bagi penataan struktur spasial aglomerasi yang kurang berkembang
dan bagi kebijakan pembangunan terkoordinasi antarkota. Kerangka
fungsi-kinerja dapat digunakan untuk membedakan kota yang memperoleh
fungsi saja, kinerja saja, keduanya sebagai \emph{borrowed size}, atau
tidak memperoleh keduanya sebagai \emph{agglomeration shadow} (akhir
pendahuluan dan Gambar 1, hlm. 1158-1159).

\textbf{Inferensi pengolahan.} Dalam batas data dan model artikel,
pendekatan tersebut dapat berfungsi sebagai alat diagnostik untuk:
membandingkan pola antarkota dan antaraglomerasi; mengenali kota yang
fungsi relatifnya tidak diikuti kinerja; serta memeriksa apakah strategi
polisentris dan peningkatan akses jaringan berkorelasi dengan perubahan
posisi kota. Ini merupakan kemungkinan penggunaan hasil, bukan aplikasi
yang diuji melalui program kebijakan atau evaluasi dampak di dalam file
(model (1)-(6), hlm. 1158-1159; bagian 2, hlm. 1160-1165).

Tidak disebutkan dalam file adanya perangkat lunak, instrumen
operasional, studi implementasi, atau evaluasi kebijakan yang telah
menerapkan hasil penelitian ini.

\subsection{Future Research
(Optional)}\label{future-research-optional-20}

\textbf{Agenda yang dinyatakan penulis.} Penulis mengusulkan tiga arah
penelitian lanjutan (akhir bagian 3.2, hlm. 1166-1167):

\begin{itemize}
\tightlist
\item
  mengoptimalkan sistem indikator fungsi kota, termasuk memasukkan
  fungsi administratif ketika data panel yang memadai tersedia, untuk
  memperbaiki pengukuran \emph{borrowed function};
\item
  memperdalam kajian mekanisme melalui kedekatan multidimensi, termasuk
  kedekatan kognitif dan institusional, serta jaringan multidimensi,
  termasuk jaringan pengetahuan dan perdagangan, terutama dengan data
  yang berubah dari waktu ke waktu; dan
\item
  mengidentifikasi batas aglomerasi secara lebih ilmiah, lalu menilai
  apakah kota tingkat county juga dapat meminjam skala dari kota lain.
\end{itemize}

\textbf{Inferensi pengolahan.} Pengembangan yang juga logis, tetapi
tidak dinyatakan secara eksplisit oleh penulis, ialah menguji
sensitivitas klasifikasi residual, menggunakan ukuran kinerja
alternatif, melacak arus antarkota secara langsung, dan menerapkan
rancangan longitudinal atau kuasi-eksperimental untuk membedakan
pengaruh jaringan dari seleksi kota ke dalam jaringan (diturunkan dari
bagian 1.1-1.3, hlm. 1158-1160, serta keterbatasan pada akhir bagian
3.2, hlm. 1166-1167).

\subsection{Provenance Notes}\label{provenance-notes-20}

\begin{itemize}
\tightlist
\item
  Review ini hanya menggunakan file
  \texttt{source/journal/A05-yang-zhu-du-2022-is-there-a-borrowed-size-in-chinas-urban-agglomerations-10.18306-dlkxjz.2022.07.002.txt}.
  Tidak ada sumber luar, berkas PDF, catatan Zotero, maupun artikel
  dalam daftar pustaka yang dibuka untuk menambah atau memverifikasi
  klaim substantif.
\item
  Metadata ekstraksi menyatakan bahwa TXT berasal dari
  \texttt{journal/A05-yang-zhu-du-2022-is-there-a-borrowed-size-in-chinas-urban-agglomerations-10.18306-dlkxjz.2022.07.002.pdf},
  diekstraksi pada 31 Agustus 2026 dengan \texttt{pdftotext\ -layout},
  dan PDF berjumlah 12 halaman. Dasar akses review ini adalah teks
  lengkap yang tersedia dalam TXT.
\item
  Identitas bibliografis diambil dari kepala artikel, baris sitasi, dan
  abstrak Inggris pada halaman cetak 1156 serta 1167. Metadata PDF
  memberi label produksi internal yang diawali
  \texttt{2-\/-\/-2022-0039} dan memuat nama penulis pertama dalam
  aksara Mandarin, serta mencantumkan \texttt{Founder\ magtech2011}
  sebagai \texttt{Author}. Entri mesin tersebut tidak diperlakukan
  sebagai kepengarangan artikel karena kepala artikel secara eksplisit
  mencantumkan Yang Tongbin, Zhu Yingming, dan Du Jiazhen.
\item
  Tubuh substantif artikel berbahasa Mandarin terdapat pada halaman
  1156-1166, sedangkan halaman 1167 memuat lanjutan daftar pustaka dan
  abstrak Inggris dengan hasil yang sama. Informasi Inggris digunakan
  untuk mengonfirmasi istilah dan identitas, tidak dihitung sebagai
  temuan kedua.
\item
  Locator \texttt{hlm.} merujuk pada nomor halaman cetak jurnal
  1156-1167, bukan nomor baris TXT. Locator bagian, model, tabel, dan
  gambar dipertahankan jika tersedia.
\item
  Tabel 1-4 terbaca sebagai teks dan menjadi dasar angka yang
  dilaporkan. Gambar 2 dan 3 tidak menyajikan seluruh nilai grafis dalam
  ekstraksi TXT; review hanya memakai pola dan angka yang dijelaskan
  dalam narasi penulis, bukan memperkirakan nilai dari gambar yang tidak
  terlihat. Gambar 1 juga dirujuk hanya untuk aturan klasifikasi yang
  dijelaskan kembali dalam teks.
\item
  Tata letak dua kolom menyebabkan beberapa kalimat halaman 1166
  terpisah dalam urutan ekstraksi. Makna bagian diskusi dan agenda
  penelitian direkonstruksi hanya dari fragmen kalimat yang bersebelahan
  pada halaman yang sama; tidak ada materi yang ditambahkan dari luar
  TXT.
\item
  Semua pernyataan evaluatif yang tidak dinyatakan oleh artikel diberi
  label \textbf{Inferensi pengolahan}. Penjelasan kausal yang diberikan
  artikel diberi label \textbf{Interpretasi penulis}, sedangkan angka,
  spesifikasi, pola, dan rekomendasi eksplisit diberi label
  \textbf{Laporan sumber}, \textbf{Temuan yang dilaporkan sumber}, atau
  \textbf{Rekomendasi penulis}.
\item
  Status penelaahan sejawat, daftar lengkap kota sampel, prosedur data
  hilang, dan artefak replikasi: Tidak disebutkan dalam file.
\end{itemize}

\section{Review A06 - Sun et
al.~(2022)}\label{review-a06---sun-et-al.-2022}

Status: Review literatur berbasis teks lengkap hasil ekstraksi PDF.

\subsection{Identitas Artikel}\label{identitas-artikel}

\begin{itemize}
\tightlist
\item
  \textbf{Kode:} A06.
\item
  \textbf{Judul:} \emph{Evaluation of Service Function of Small Towns in
  China from the Perspective of Borrowed Size and Agglomeration Shadow:
  A Case Study of Southern Jiangsu Province}.
\item
  \textbf{Penulis:} Sun Dongqi, Lu Jiayi, Zhang Mingdou, dan Niu Fangqu
  (penulis korespondensi).
\item
  \textbf{Tahun:} 2022.
\item
  \textbf{Jurnal:} \emph{Progress in Geography}, volume 41, nomor 2,
  hlm. 199-213.
\item
  \textbf{DOI:} 10.18306/dlkxjz.2022.02.002.
\item
  \textbf{Tanggal naskah:} diterima 16 Maret 2021 dan direvisi 25
  Oktober 2021.
\item
  \textbf{Pendanaan:} National Natural Science Foundation of China,
  proyek nomor 41971162.
\item
  \textbf{Kata kunci:} borrowed size; agglomeration shadow; evaluasi
  fungsi pelayanan; kota kecil; Jiangsu selatan.
\item
  \textbf{Wilayah dan periode data utama:} Suzhou, Wuxi, dan Changzhou
  di Jiangsu selatan; data POI tahun 2019.
\item
  \textbf{Locator identitas:} hlm. 199, blok judul, abstrak, catatan
  penulis, dan format sitasi; hlm. 213, judul dan abstrak berbahasa
  Inggris.
\end{itemize}

\subsection{Summarized Abstract}\label{summarized-abstract-21}

\textbf{Laporan sumber.} Artikel ini menanggapi kekurangan penelitian
geografis tentang fungsi pelayanan kota kecil pada skala wilayah
setingkat kota praja. Penulis menilai fungsi pelayanan 248 unit yang
terdiri atas kota praja dan subdistrik di Suzhou, Wuxi, dan Changzhou,
yaitu wilayah Jiangsu selatan yang berdekatan dengan Shanghai. Bahan
utamanya ialah 741.209 titik fasilitas pelayanan dari 14 kelompok
industri yang diseleksi dari 1.964.591 data \emph{point of interest}
(POI) Amap tahun 2019. Analisis menggabungkan estimasi kepadatan kernel,
\emph{analytic hierarchy process} (AHP), klasifikasi spasial, dan
regresi linear berganda yang kemudian diestimasi dengan \emph{feasible
generalized least squares} (FGLS) karena heteroskedastisitas. Hasil
deskriptif menunjukkan konsentrasi fasilitas yang tinggi di pusat kota
dan pusat kota setingkat kabupaten serta sepanjang koridor transportasi;
indeks fungsi pelayanan cenderung lebih tinggi pada zona terdekat dengan
Shanghai; dan 55 unit diidentifikasi sebagai area bayangan pelayanan.
Regresi menunjukkan hubungan positif antara jarak ke Shanghai dan indeks
pelayanan, yang berarti indeks cenderung meningkat ketika lokasi makin
jauh dari Shanghai setelah kovariat model diperhitungkan. (Hlm. 199,
abstrak; hlm. 201, Bagian 1.1-1.2; hlm. 203-205, Bagian 1.3 dan 2.1-2.2;
hlm. 207-210, Bagian 2.3 dan 3.)

\textbf{Interpretasi penulis.} Penulis membaca pola tersebut sebagai
bukti bahwa pengaruh Shanghai terhadap fungsi pelayanan kota kecil di
Jiangsu selatan signifikan, paling menonjol pada zona 100 km, serta
bahwa proses borrowed size dan agglomeration shadow berlangsung secara
bersamaan. Mereka mengaitkan kepadatan tinggi di dekat Shanghai dengan
hubungan fungsional regional, sedangkan koefisien jarak yang positif dan
keberadaan kota berindeks rendah di antara kota berindeks lebih tinggi
dibaca sebagai bayangan aglomerasi. (Hlm. 208-210, Bagian 2.3 dan 3.)

\textbf{Inferensi pengolahan.} Bukti paling langsung dalam artikel
mendukung adanya pola spasial pelayanan dan asosiasi yang konsisten
dengan agglomeration shadow. Bukti borrowed size lebih tidak langsung
karena berasal dari indeks POI lintas bagian, kedekatan geografis,
signifikansi regresi, dan narasi keterkaitan historis, bukan dari
pengukuran aliran pemanfaatan layanan atau mekanisme peminjaman fungsi.
Karena desainnya potong lintang, istilah ``proses'' sebaiknya dipahami
sebagai interpretasi penulis atas pola asosiasi, bukan identifikasi
kausal atau perubahan dari waktu ke waktu.

\subsection{Problem Statement}\label{problem-statement-21}

\textbf{Laporan sumber.} Penelitian tentang kota kecil di Tiongkok
relatif banyak, tetapi keterbatasan data membuat kajian perencanaan
lebih sering berfokus pada tata kelola, sedangkan kajian geografi pada
skala kota praja masih sedikit. Penelitian yang tersedia terutama
membahas distribusi spasial, penggunaan lahan, infrastruktur, atau
perilaku belanja. Fungsi pelayanan kota kecil pada skala kota praja,
terutama dalam relasinya dengan metropolitan terdekat, belum dinilai
secara memadai. Padahal, kota kecil berfungsi sebagai jembatan antara
wilayah perkotaan makro dan hinterland perdesaan serta sebagai penyedia
pelayanan bagi kawasan perdesaan. (Hlm. 200-201, pendahuluan.)

\textbf{Laporan sumber.} Literatur borrowed size dan agglomeration
shadow lebih sering menguji dampak kota besar terhadap pertumbuhan
penduduk atau ekonomi wilayah sekitarnya. Bukti empiris bahwa kota kecil
meminjam fungsi kota besar masih terbatas, dan penelitian khusus sektor
pelayanan juga sedikit. Karena itu, belum jelas apakah kedekatan dengan
Shanghai memperkuat fungsi pelayanan kota kecil melalui borrowed size
atau justru menekannya melalui kompetisi dan agglomeration shadow. (Hlm.
199-201, pendahuluan.)

\textbf{Inferensi pengolahan.} Masalah empiris artikel mempunyai dua
lapisan yang tidak sepenuhnya sama: mengukur banyaknya fungsi pelayanan
pada tingkat kota praja dan mengatribusikan pola itu kepada borrowed
size atau agglomeration shadow. Metode pertama relatif operasional
melalui POI, sedangkan atribusi mekanisme menuntut bukti yang lebih kuat
daripada kedekatan dan korelasi spasial.

\subsection{Objectives}\label{objectives-21}

\textbf{Laporan sumber.} Tujuan penelitian adalah:

\begin{itemize}
\tightlist
\item
  menilai secara kuantitatif fungsi pelayanan kota kecil di Jiangsu
  selatan pada skala wilayah kota praja;
\item
  memetakan diferensiasi spasial dan tingkat pemusatan fasilitas
  pelayanan;
\item
  membangun indeks fungsi pelayanan untuk setiap kota kecil;
\item
  mengenali area berfungsi pelayanan rendah yang dikelilingi unit dengan
  fungsi lebih tinggi sebagai kandidat ``shadow area'';
\item
  menguji pengaruh potensi internal, posisi dalam sistem perkotaan, dan
  aksesibilitas infrastruktur transportasi terhadap indeks fungsi
  pelayanan; dan
\item
  memeriksa apakah relasi kota kecil dengan Shanghai memperlihatkan
  borrowed size, agglomeration shadow, atau keduanya.
\end{itemize}

Tujuan tersebut dinyatakan dalam pendahuluan dan dijabarkan menjadi lima
langkah penelitian. (Hlm. 201, akhir pendahuluan dan Bagian 1.2.)

\subsection{Literature Survey}\label{literature-survey-21}

\textbf{Laporan sumber.} Tinjauan pustaka bergerak dari teori tempat
sentral menuju teori jaringan kota. Teori tempat sentral melahirkan
gagasan pusat-pinggiran serta polarisasi-difusi, tetapi globalisasi,
informatisasi, \emph{space of flows}, dan \emph{central flow theory}
menekankan bahwa fungsi sebuah tempat juga bergantung pada hubungan
nonlokal. Dalam wilayah polisentris, kota-kota dapat memiliki
spesialisasi yang saling melengkapi dan kota kecil dapat menanggung
fungsi yang melampaui kapasitas lokalnya. (Hlm. 199, pendahuluan.)

\textbf{Laporan sumber.} Borrowed size ditelusuri ke Alonso (1973): kota
kecil yang dekat dengan pusat populasi metropolitan dapat menunjukkan
ciri kota besar karena memanfaatkan keuntungan aglomerasi sembari
menghindari sebagian biayanya. Meijers dan Burger kemudian memperluasnya
sebagai kemampuan kota kecil memanfaatkan fungsi metropolitan dan
meningkatkan posisi dalam jaringan. Studi Phelps dan kolega tentang
London serta studi lain mengenai interaksi hierarki perkotaan
menyediakan dukungan empiris, tetapi bukti pada fungsi pelayanan kota
kecil masih terbatas. (Hlm. 199-200, pendahuluan.)

\textbf{Laporan sumber.} Agglomeration shadow dijelaskan sebagai sisi
kompetitif atau efek balik dari aglomerasi: kota kecil dekat pusat besar
dapat menanggung fungsi lebih sedikit daripada kota kecil yang lebih
jauh karena aktivitas tersedot ke pusat. Konsep ini dikaitkan dengan
teori ekonomi geografi baru serta efek difusi dan \emph{backwash}
Myrdal. Sebagian besar studi terdahulu mengukurnya dengan pertumbuhan
penduduk atau ekonomi. (Hlm. 200, pendahuluan.)

\textbf{Laporan sumber.} Malý menjadi acuan metodologis terdekat karena
menilai fungsi pelayanan kota kecil di Moravia Selatan dan
meregresikannya terhadap aksesibilitas transportasi, sejarah, budaya,
serta daya tarik wisata; studi tersebut menemukan borrowed size dan
agglomeration shadow dapat berdampingan. De Vos dan kolega menemukan
borrowed size pada industri internet pita lebar di Swedia. Penulis
menempatkan kedua studi itu sebagai pengecualian dalam literatur yang
masih sedikit membahas pelayanan. (Hlm. 200, pendahuluan.)

\textbf{Laporan sumber.} Dalam konteks Tiongkok, Zhang Jingxiang dan
kolega memperkenalkan konsep ``metropolitan shadow area'' melalui
interaksi polarisasi kota besar, daya tarik kota kecil, dan kebijakan.
Kajian lanjutan membahas hubungan industri, strategi aktor, pertumbuhan
ekonomi, polisentrisitas, dan jasa kehidupan di Shanghai. Temuan
terdahulu tidak tunggal: ada bukti kota besar mendukung pertumbuhan kota
kecil di Delta Sungai Yangtze, tetapi ada pula bukti agglomeration
shadow di Beijing-Tianjin-Hebei. (Hlm. 200-201, pendahuluan.)

\textbf{Interpretasi penulis.} Literatur tersebut dipakai untuk
menjustifikasi bahwa borrowed size dan agglomeration shadow bukan
pilihan yang saling meniadakan. Jenis fungsi yang diukur dan posisi
suatu kota dalam jaringan dapat membuat manfaat serta tekanan aglomerasi
muncul bersamaan. Artikel lalu memindahkan perdebatan yang biasanya
berada pada skala kota atau kabupaten ke skala kota praja dan sektor
pelayanan. (Hlm. 200-201, pendahuluan.)

\textbf{Inferensi pengolahan.} Tinjauan ini memberi landasan konseptual
yang luas dan menyebut studi pembanding utama, tetapi tidak disajikan
sebagai tinjauan sistematis. File tidak menjelaskan strategi pencarian,
kriteria inklusi, rentang waktu korpus, atau penilaian mutu studi
terdahulu; karena itu, kelengkapan survei literatur tidak dapat dinilai
dari TXT.

\subsection{Method Used}\label{method-used-21}

\textbf{Laporan sumber.} Penelitian memakai pendekatan kuantitatif
dengan data utama tahun 2019 dan unit analisis administratif berupa 248
kota praja dan subdistrik. Alurnya ialah sebagai berikut. (Hlm. 201,
Bagian 1.1-1.2.)

\begin{itemize}
\tightlist
\item
  Sebanyak 1.964.591 POI Amap tahun 2019 disaring menurut klasifikasi
  industri nasional Tiongkok dan klasifikasi POI Amap. Penulis
  mempertahankan 741.209 POI dalam 14 kelompok pelayanan. (Hlm. 201 dan
  203, Bagian 1.2 dan Tabel 1.)
\item
  Estimasi kepadatan kernel digunakan untuk memetakan konsentrasi
  kontinu POI. Analisis dilakukan dengan ArcGIS 10.3 dan bandwidth
  bawaan yang dihitung melalui aturan empiris Silverman. (Hlm. 203,
  Bagian 1.3.1, Persamaan 1.)
\item
  AHP satu tingkat dibangun dengan Yaahp 10.3. Empat belas kelompok
  industri dibandingkan berpasangan untuk memperoleh bobot nilai
  pelayanan. Bobot terbesar diberikan kepada perdagangan besar dan
  eceran (0,2278), diikuti jasa penduduk, reparasi, dan jasa lain
  (0,1551), serta akomodasi dan makanan (0,1520). Bobot terkecil
  diberikan kepada real estat (0,0151) dan keuangan (0,0184). (Hlm. 204,
  Bagian 1.3.2 dan Tabel 3.)
\item
  Indeks fungsi pelayanan setiap unit dihitung sebagai jumlah 14 hasil
  kali antara jumlah fasilitas industri dan bobot industri, yaitu
  \texttt{I\ =\ sum(F\_i\ V\_i)}. Indeks lalu dikelompokkan dengan
  \emph{natural breaks} ArcGIS 10.3 menjadi lima kelas: tinggi, agak
  tinggi, sedang, agak rendah, dan rendah. (Hlm. 204-205, Persamaan 2
  dan Bagian 2.2.)
\item
  Sebuah unit kelas rendah ditetapkan sebagai ``shadow area'' apabila
  dikelilingi unit kelas tinggi, agak tinggi, atau sedang. Penulis
  menolak menyamakan semua unit berkelas rendah dengan bayangan
  aglomerasi karena ketimpangan regional dapat mempunyai penyebab lain.
  (Hlm. 201, Bagian 1.2, langkah 3.)
\item
  Posisi terhadap Shanghai dibagi dengan radius 100, 150, dan 200 km
  dari koordinat Pemerintah Kota Shanghai. Unit pada garis batas
  dimasukkan ke zona yang memuat bagian wilayahnya paling luas,
  sedangkan sejumlah kecil unit di luar 200 km digabungkan ke kelompok
  200 km. (Hlm. 205 dan 207, Bagian 2.2 serta catatan Tabel 4.)
\item
  Regresi linear berganda menghubungkan logaritma natural indeks
  pelayanan dengan enam prediktor: penduduk tetap, PDB, jarak pemerintah
  kota praja ke pemerintah kabupaten, jarak ke pemerintah kota tingkat
  prefektur, jarak ke Pemerintah Kota Shanghai, dan waktu berkendara
  terpendek ke pintu jalan tol. Penulis menguji multikolinearitas dengan
  \emph{variance inflation factor} (VIF) dan heteroskedastisitas dengan
  uji White. Karena hanya heteroskedastisitas yang dilaporkan
  terdeteksi, estimasi dilanjutkan dengan FGLS melalui pembobotan WLS.
  (Hlm. 203-204, Bagian 1.2 dan 1.3.3, Persamaan 3.)
\item
  Model diestimasi untuk sampel penuh dan secara terpisah untuk tiga
  zona jarak dari Shanghai. (Hlm. 207-209, Bagian 2.3 dan Tabel 6.)
\end{itemize}

\textbf{Inferensi pengolahan.} Berdasarkan penggunaan data utama dari
satu tahun, rancangan ini bersifat potong lintang. Rancangan
menggabungkan pemetaan eksploratif, pembentukan indeks normatif, dan
model asosiasional. Ia tidak merupakan eksperimen, kuasi-eksperimen,
analisis panel, atau model kausal; istilah ``pengaruh'' dalam laporan
hasil karenanya lebih aman dibaca sebagai hubungan statistik bersyarat.

\subsection{Dataset}\label{dataset-21}

\textbf{Laporan sumber.} Dataset utama dan pendukung terdiri atas:

\begin{itemize}
\tightlist
\item
  \textbf{POI Amap:} 1.964.591 POI tahun 2019 untuk Suzhou, Wuxi, dan
  Changzhou dikumpulkan melalui platform API terbuka Amap. Berdasarkan
  klasifikasi industri nasional GB/4754-2011 dan kode klasifikasi Amap,
  741.209 POI dipilih sebagai fasilitas yang berkaitan dengan fungsi
  pelayanan. (Hlm. 201, Bagian 1.2.)
\item
  \textbf{Empat belas kelompok pelayanan:} perdagangan besar dan eceran
  (311.723 POI); transportasi, pergudangan, dan pos (3.109); akomodasi
  dan makanan (135.003); transmisi informasi, perangkat lunak, dan
  layanan teknologi informasi (3.924); keuangan (6.785); real estat
  (24.988); persewaan dan layanan bisnis (11.154); penelitian ilmiah dan
  layanan teknis (5.478); pengelolaan air, lingkungan, dan fasilitas
  publik (23.452); jasa penduduk, reparasi, dan jasa lain (73.688);
  pendidikan (25.326); kesehatan dan pekerjaan sosial (17.907);
  kebudayaan, olahraga, dan hiburan (53.494); serta administrasi publik,
  jaminan sosial, dan organisasi sosial (45.178). (Hlm. 203, Tabel 1.)
\item
  \textbf{Penduduk tetap:} data dari Jiangsu Government Service Network.
  Dari 248 unit, 20 nilai hilang dan dilengkapi dengan interpolasi.
  Jenis atau prosedur interpolasi tidak diterangkan dalam file. (Hlm.
  201, Bagian 1.2, langkah 4.)
\item
  \textbf{PDB:} data berasal dari buku statistik atau buku tahunan tahun
  2020 milik distrik atau kota setingkat kabupaten di Suzhou, Wuxi, dan
  Changzhou. File tidak menegaskan tahun observasi setiap nilai PDB,
  seri atau tabel spesifik yang dipakai, ataupun penanganan nilai yang
  mungkin hilang. (Hlm. 201, Bagian 1.2, langkah 4.)
\item
  \textbf{Jarak dan waktu tempuh:} posisi dalam sistem perkotaan
  diwakili oleh jarak dari pemerintah kota praja ke pemerintah
  kabupaten, pemerintah kota tingkat prefektur, dan Pemerintah Kota
  Shanghai. Akses jalan tol diperoleh dengan memasukkan titik asal dan
  tujuan satu per satu di Amap untuk mendapatkan waktu berkendara
  terpendek ke pintu tol terdekat. (Hlm. 201, Bagian 1.2; hlm. 203,
  Tabel 2.)
\item
  \textbf{Data spasial administratif:} batas kota praja dan subdistrik
  Suzhou, Wuxi, dan Changzhou berasal dari basis data sistem informasi
  geografis Departemen Sumber Daya Alam Provinsi Jiangsu; batas
  administratif Shanghai berasal dari basis data sistem informasi
  geografis dasar nasional. (Hlm. 201, Bagian 1.2.)
\end{itemize}

\textbf{Inferensi pengolahan.} Dataset POI mengukur keberadaan atau
jumlah lokasi yang terdaftar, bukan kapasitas fasilitas, mutu pelayanan,
jumlah tenaga kerja, volume transaksi, jam operasi, keterjangkauan, atau
pemanfaatan oleh penduduk. File tidak menyediakan tanggal pengunduhan
POI, parameter permintaan API, aturan deduplikasi, pemeriksaan
koordinat, tingkat kelengkapan Amap, tanggal batas administratif,
ataupun uji bias cakupan antarkategori dan antarwilayah. Oleh sebab itu,
indeks paling tepat dibaca sebagai proksi kuantitas fasilitas yang
tercatat.

\subsection{Population / Sample}\label{population-sample-21}

\textbf{Laporan sumber.} Populasi wilayah studi dalam arti administratif
adalah seluruh 248 kota praja dan subdistrik di wilayah tradisional
Jiangsu selatan, yaitu Suzhou, Wuxi, dan Changzhou. Wilayah ini seluas
17.669,79 km2, berpenduduk 22,026 juta jiwa pada 2019, dan berdekatan
dengan Shanghai. Karena tidak ada definisi skala kota kecil yang tegas
pada tingkat nasional, penulis memperlakukan semua kota praja dan
subdistrik sebagai ``small towns''. (Hlm. 201, Bagian 1.1.)

\textbf{Laporan sumber.} Analisis indeks dan pemetaan kelas pelayanan
menggunakan 248 unit. Distribusinya ialah 81 unit pada kolom zona 100
km, 93 pada zona 150 km, dan 74 pada zona 200 km, dengan unit di luar
200 km digabungkan ke kelompok terakhir. Regresi hanya memakai 216
observasi: 65 pada model 100 km, 88 pada model 150 km, dan 63 pada model
200 km. (Hlm. 207, Tabel 4; hlm. 208-209, Tabel 5-6.)

\textbf{Interpretasi penulis.} Penulis menyatakan statistik sampel yang
dekat dengan parameter populasi menunjukkan representativitas dan
reliabilitas inferensi yang tinggi. Dasarnya ialah penyebaran variabel
yang mereka nilai relatif kecil. (Hlm. 207-208, Bagian 2.3.)

\textbf{Inferensi pengolahan.} Tidak ada pengambilan sampel acak; desain
awalnya lebih menyerupai sensus administratif wilayah studi. File tidak
menjelaskan mengapa 248 unit pada analisis spasial menyusut menjadi 216
dalam regresi, padahal 20 nilai penduduk yang hilang disebut telah
diinterpolasi. Alur eksklusi 32 unit dan pola data hilangnya tidak
tersedia, sehingga representativitas sampel regresi belum dapat
diverifikasi. Selain itu, penggabungan subdistrik pusat kota, pusat
pemerintahan kabupaten, zona pengembangan, dan kota praja perdesaan
dalam satu label ``small towns'' menghasilkan populasi analitis yang
sangat heterogen.

\subsection{Dependent Variables}\label{dependent-variables-21}

\textbf{Laporan sumber.} Variabel terikat utama ialah \textbf{indeks
fungsi pelayanan kota kecil} (\texttt{y} atau \texttt{I}). Indeks
tingkat unit dihitung dengan menjumlahkan jumlah POI pada setiap
kelompok industri (\texttt{F\_i}) yang dikalikan bobot AHP kelompok
tersebut (\texttt{V\_i}) untuk 14 industri. Dengan demikian, indeks
merupakan skor aditif berbobot berbasis hitungan fasilitas. (Hlm. 204,
Bagian 1.3.2, Persamaan 2 dan Tabel 3.)

\textbf{Laporan sumber.} Untuk pemetaan, indeks mentah dibagi menjadi
lima kelas dengan \emph{natural breaks}: tinggi (2.000,01-4.360,87),
agak tinggi (1.000,01-2.000,00), sedang (500,01-1.000,00), agak rendah
(250,01-500,00), dan rendah (0,10-250,00). Klasifikasi ini juga menjadi
dasar identifikasi area bayangan. (Hlm. 205, Bagian 2.2.)

\textbf{Laporan sumber.} Untuk regresi, metode menyatakan bahwa variabel
terikat ditransformasikan dengan logaritma natural, ditulis
\texttt{ln\ y}. Tabel statistik deskriptif melaporkan rerata
\texttt{ln\ y} sebesar 5,742, simpangan baku 0,831, minimum 3,913, dan
maksimum 7,960 untuk 216 observasi. (Hlm. 204, Bagian 1.3.3, Persamaan
3; hlm. 208, Tabel 5.)

\textbf{Inferensi pengolahan.} File hanya menyatakan secara eksplisit
bahwa variabel terikat dilogaritmakan, tetapi narasi hasil menafsirkan
kenaikan semua prediktor dalam persentase. Rerata prediktor pada Tabel 5
juga tampak berada pada suatu skala transformasi. Bentuk transformasi
ruas kanan, standardisasi yang mungkin dilakukan, dan alasan mengapa
rerata variabel terikat pada Tabel 6 (0,331 untuk sampel penuh) berbeda
dari rerata \texttt{ln\ y} pada Tabel 5 (5,742) tidak dijelaskan.
Akibatnya, makna numerik koefisien tidak sepenuhnya dapat direplikasi
dari TXT.

\textbf{Inferensi pengolahan.} Karena indeks memakai jumlah POI mentah
tanpa normalisasi luas, penduduk, kapasitas, atau permintaan, skor
menggabungkan ukuran unit dengan fungsi pelayanannya. Sebutan
``centrality'' yang diberikan penulis karenanya merupakan interpretasi
atas skor, bukan ukuran aliran pelayanan atau surplus fungsi yang
dibuktikan secara terpisah.

\subsection{Results}\label{results-21}

\textbf{Laporan sumber: kepadatan fasilitas.} Peta kepadatan kernel
menghasilkan lima pola utama. Pertama, kepadatan fasilitas kota kecil
dekat Shanghai umumnya tinggi. Kedua, kepadatan tertinggi terkonsentrasi
di kawasan pusat Suzhou, Wuxi, dan Changzhou beserta hinterland
terdekat, membentuk pola pusat-pinggiran. Ketiga, pusat kota setingkat
kabupaten seperti Kunshan, Changshu, Zhangjiagang, Jiangyin, Yixing, dan
Liyang membentuk inti sekunder. Keempat, inti tersebut mengikuti jalan
tol, jalan nasional, dan rel, sehingga membentuk struktur titik-sumbu.
Kelima, unit yang jauh dari pusat relatif lebih seimbang, kecuali
beberapa unit berpenduduk sedikit; bagian barat mempunyai kepadatan
lebih rendah. (Hlm. 204-205, Bagian 2.1 dan Gambar 2.)

\textbf{Laporan sumber: perbedaan sektor.} Perdagangan besar dan eceran
mencapai kepadatan kernel maksimum 496, sedangkan akomodasi dan makanan
mencapai 182. Jasa penduduk, reparasi, dan jasa lain mencapai 85;
kebudayaan, olahraga, dan hiburan mencapai 78. Sejumlah sektor lain
berada pada rentang 10-50. Penelitian ilmiah dan teknis, keuangan,
transportasi-pergudangan-pos, serta informasi dan teknologi mempunyai
nilai maksimum di bawah 10, tetapi persebarannya disebut lebih merata.
Penulis menilai perbedaan antarsektor ada, tetapi pola umumnya serupa.
(Hlm. 205-206, Bagian 2.1 dan Gambar 3.)

\textbf{Laporan sumber: kelas indeks.} Dari 248 unit, terdapat 6 unit
berindeks tinggi, 18 agak tinggi, 43 sedang, 68 agak rendah, dan 113
rendah. Dengan demikian, gabungan kelas agak rendah dan rendah berjumlah
181 unit. Di antara 24 unit pada gabungan kelas tinggi dan agak tinggi,
14 berada pada zona 100 km, 4 pada zona 150 km, dan 6 pada zona 200 km.
Dari 43 unit kelas sedang, distribusinya masing-masing 15, 16, dan 12.
Dari 181 unit kelas agak rendah dan rendah, distribusinya masing-masing
52, 73, dan 56. (Hlm. 205 dan 207, Bagian 2.2 dan Tabel 4.)

\textbf{Laporan sumber: pola jarak dan bayangan.} Indeks pelayanan
secara keseluruhan dinyatakan lebih tinggi dekat Shanghai, terutama pada
zona 100 km, lalu menurun pada zona 150-200 km. Berdasarkan aturan
spasial penulis, 55 unit diklasifikasikan sebagai ``shadow areas''. Area
itu muncul pada setiap zona, termasuk pusat kota utama dan sekitar
Shanghai. (Hlm. 205 dan 208, Bagian 2.2, Gambar 4-5.)

\textbf{Laporan sumber: regresi sampel penuh.} Model FGLS sampel penuh
(\texttt{n\ =\ 216}) mempunyai \texttt{R2\ =\ 0,876} dan uji F 245,132
dengan \texttt{Prob\ \textgreater{}\ F\ =\ 0}. Koefisien penduduk
sebesar 0,658 (\texttt{p\ \textless{}\ 0,001}), PDB 0,567
(\texttt{p\ \textless{}\ 0,001}), jarak ke pemerintah kabupaten 0,123
(\texttt{p\ =\ 0,085}), jarak ke pemerintah kota -0,038
(\texttt{p\ =\ 0,565}), jarak ke Shanghai 0,198 (\texttt{p\ =\ 0,005}),
dan waktu ke pintu tol 0,237 (\texttt{p\ \textless{}\ 0,001}).
Berdasarkan ambang penulis, semua kecuali jarak ke pemerintah kota
signifikan setidaknya pada taraf 10\%. (Hlm. 209, Tabel 6.)

\textbf{Interpretasi penulis: regresi sampel penuh.} Penulis menafsirkan
koefisien penduduk dan PDB sebagai dukungan terhadap peran potensi
internal. Koefisien positif jarak ke pemerintah kabupaten dan Shanghai
ditafsirkan berarti bahwa semakin dekat suatu kota kecil kepada pusat
tersebut, indeksnya justru semakin rendah, sehingga mendukung
agglomeration shadow. Jarak ke pemerintah kota tingkat prefektur tidak
signifikan. Koefisien positif waktu ke pintu tol disebut menunjukkan
pengaruh positif aksesibilitas transportasi. (Hlm. 207-208, Bagian 2.3.)

\textbf{Laporan sumber: regresi zona 100 km.} Model mempunyai 65
observasi dan \texttt{R2\ =\ 0,888}. Hanya jarak ke Shanghai yang
signifikan pada taraf 5\%, dengan koefisien 0,572 dan
\texttt{p\ =\ 0,011}; penduduk, PDB, dua jarak pusat administratif lain,
dan waktu ke jalan tol tidak signifikan. (Hlm. 209, Tabel 6.)

\textbf{Laporan sumber: regresi zona 150 km.} Model mempunyai 88
observasi dan \texttt{R2\ =\ 0,889}. Penduduk berkoefisien 0,893
(\texttt{p\ \textless{}\ 0,001}), PDB 0,515
(\texttt{p\ \textless{}\ 0,001}), jarak ke Shanghai 0,252
(\texttt{p\ =\ 0,035}), dan waktu ke pintu tol 0,171
(\texttt{p\ =\ 0,093}). Dua jarak pusat administratif tidak signifikan.
(Hlm. 209, Tabel 6.)

\textbf{Laporan sumber: regresi zona 200 km.} Model mempunyai 63
observasi dan \texttt{R2\ =\ 0,891}. Jarak ke Shanghai berkoefisien
0,656 (\texttt{p\ \textless{}\ 0,001}) dan waktu ke pintu tol 0,236
(\texttt{p\ =\ 0,035}); prediktor substantif lain tidak signifikan pada
taraf 10\%. (Hlm. 209, Tabel 6.)

\textbf{Interpretasi penulis: perbandingan zona.} Penulis menyimpulkan
Shanghai paling dominan pada zona 100 km karena jarak ke Shanghai
menjadi satu-satunya variabel signifikan di sana. Pada zona 150 dan 200
km, faktor lain juga signifikan: potensi internal terutama penting pada
150 km dan waktu ke jalan tol penting pada 150 serta 200 km. Penulis
menyebut hal ini sebagai agglomeration shadow Shanghai yang paling nyata
dalam 100 km. (Hlm. 208-209, Bagian 2.3 dan Tabel 6.)

\subsection{Findings}\label{findings-21}

\textbf{Temuan yang dilaporkan sumber.} Fungsi pelayanan kota kecil di
Jiangsu selatan tidak tersebar acak. Pusat kota utama dan pusat kota
setingkat kabupaten menjadi inti, sedangkan jaringan transportasi
membentuk sumbu konsentrasi. Pola tersebut sekaligus memperlihatkan
logika tempat sentral dan struktur regional berbasis jaringan. (Hlm.
204-206, Bagian 2.1, Gambar 2-3.)

\textbf{Temuan yang dilaporkan sumber.} Secara deskriptif, unit dekat
Shanghai memiliki indeks lebih tinggi. Namun, setelah penduduk, PDB,
posisi administratif, dan akses jalan tol dimasukkan ke model, koefisien
jarak ke Shanghai tetap positif pada sampel penuh dan ketiga zona: unit
yang lebih jauh justru berkaitan dengan indeks lebih tinggi. (Hlm. 205
dan 207-209, Bagian 2.2-2.3, Gambar 4 dan Tabel 6.)

\textbf{Interpretasi penulis.} Kombinasi kedua pola itu dibaca sebagai
koeksistensi borrowed size dan agglomeration shadow. Kedekatan dan
integrasi fungsional dengan Shanghai dipandang memungkinkan kota-kota
Jiangsu selatan meminjam sumber daya, teknologi, industri, dan pasar
metropolitan. Pada saat yang sama, kompetisi dengan pusat besar menekan
fungsi pelayanan sejumlah kota yang sangat dekat atau terbenam di antara
pusat lain. (Hlm. 208-210, Bagian 2.3 dan 3.)

\textbf{Inferensi pengolahan.} Hasil lebih kuat menunjukkan paradoks
antara pola kasar dan pola kondisional daripada membuktikan dua
mekanisme secara langsung: konsentrasi mentah tinggi di dekat Shanghai,
tetapi asosiasi parsial jarak menunjukkan gradien yang konsisten dengan
bayangan aglomerasi. Borrowed size belum diukur melalui aliran pengguna,
akses lintas batas, jaringan perusahaan, komplementaritas fungsi, atau
kinerja yang melampaui ukuran lokal.

\textbf{Inferensi pengolahan.} Pernyataan bahwa efek Shanghai ``paling
signifikan'' pada 100 km memerlukan kualifikasi. Jarak ke Shanghai
memang satu-satunya prediktor signifikan pada model 100 km, tetapi
koefisiennya lebih besar dan nilai \texttt{p}-nya lebih kecil pada model
200 km (0,656; \texttt{p\ \textless{}\ 0,001}) daripada pada model 100
km (0,572; \texttt{p\ =\ 0,011}). Jadi, ``paling signifikan'' tampaknya
merujuk pada dominasi relatif terhadap kovariat lain, bukan pada besaran
koefisien atau tingkat signifikansi statistik tertinggi.

\textbf{Inferensi pengolahan.} Koefisien positif waktu tempuh ke pintu
tol secara literal berarti indeks meningkat ketika waktu ke pintu tol
bertambah, yaitu ketika aksesibilitas memburuk. Ini tidak sejalan dengan
narasi penulis bahwa aksesibilitas transportasi memberi pengaruh
positif, kecuali variabel dikodekan terbalik atau ditransformasikan
dengan cara yang tidak dijelaskan. File tidak menyediakan keterangan
yang dapat menyelesaikan persoalan arah ini.

\subsection{Conclusions}\label{conclusions-21}

\textbf{Kesimpulan yang dinyatakan penulis.} Penulis merumuskan tiga
kesimpulan utama. (Hlm. 209-210, Bagian 3.1.)

\begin{itemize}
\tightlist
\item
  Kota kecil dekat Shanghai secara umum mempunyai kepadatan fasilitas
  pelayanan lebih tinggi. Area berkepadatan tinggi dan agak tinggi
  terkonsentrasi pada pusat kota utama, pusat distrik atau kota
  setingkat kabupaten, dan sekitarnya. Persebarannya mengikuti jalan
  tol, jalan nasional, dan rel sehingga membentuk pola titik-sumbu. Unit
  yang jauh dari pusat relatif seimbang kecuali beberapa unit
  berpenduduk sedikit; perbedaan pola antarsektor ada, tetapi tidak
  besar.
\item
  Indeks fungsi pelayanan cenderung lebih tinggi dekat Shanghai,
  terutama pada zona 100 km, dan menurun pada zona yang lebih jauh
  hingga 200 km.
\item
  Shanghai mempunyai hubungan signifikan dengan indeks fungsi pelayanan
  kota kecil Jiangsu selatan. Penulis menyatakan pengaruhnya paling
  menonjol dalam 100 km dan menyimpulkan borrowed size serta
  agglomeration shadow berlangsung berdampingan.
\end{itemize}

\textbf{Inferensi pengolahan.} Dua kesimpulan pertama didukung langsung
oleh distribusi POI, kelas indeks, dan peta. Bagian agglomeration shadow
pada kesimpulan ketiga memperoleh dukungan dari koefisien jarak yang
positif dan identifikasi 55 area berindeks rendah. Bagian borrowed size
masih merupakan interpretasi awal karena tidak mempunyai indikator
tersendiri yang mengukur fungsi yang dipinjam dari Shanghai. Kesimpulan
juga tidak boleh digeneralisasi menjadi efek kausal Shanghai atau pola
universal kota kecil Tiongkok.

\subsection{Contributions}\label{contributions-21}

\textbf{Kontribusi yang diklaim atau ditunjukkan sumber.}

\begin{itemize}
\tightlist
\item
  \textbf{Kontribusi skala analisis:} artikel membawa perdebatan
  borrowed size dan agglomeration shadow ke skala kota praja, skala yang
  oleh penulis dinilai kurang diteliti karena kelangkaan data. (Hlm.
  201, pendahuluan; hlm. 210, Bagian 3.2.)
\item
  \textbf{Kontribusi substantif:} alih-alih hanya memakai pertumbuhan
  ekonomi atau penduduk, artikel menjadikan fungsi pelayanan lintas 14
  sektor sebagai objek empiris. (Hlm. 199-201, pendahuluan; hlm. 203,
  Tabel 1.)
\item
  \textbf{Kontribusi data:} pemanfaatan 741.209 POI memungkinkan
  perbandingan fasilitas pada 248 unit administratif dengan resolusi
  lebih rinci daripada statistik agregat kota. (Hlm. 201 dan 203, Bagian
  1.2 dan Tabel 1.)
\item
  \textbf{Kontribusi metode:} artikel menyusun alur dari kepadatan
  kernel, pembobotan AHP, indeks pelayanan, klasifikasi \emph{natural
  breaks}, identifikasi lembah spasial, hingga regresi FGLS. Penulis
  menawarkan aturan ``unit rendah yang dikelilingi unit lebih tinggi''
  sebagai cara awal mengenali area bayangan pelayanan. (Hlm. 201 dan
  203-205, Bagian 1.2-1.3; hlm. 210, Bagian 3.2.)
\item
  \textbf{Kontribusi empiris:} artikel mendokumentasikan bahwa pola
  konsentrasi dekat Shanghai dan pola kondisional yang konsisten dengan
  bayangan aglomerasi dapat muncul dalam wilayah yang sama. (Hlm.
  205-210, Bagian 2.2-3.1.)
\item
  \textbf{Kontribusi kebijakan:} hasil memunculkan kelompok kota
  berfungsi rendah yang dapat diprioritaskan dalam agenda pemerataan
  pelayanan publik. (Hlm. 210, Bagian 3.2.)
\end{itemize}

\textbf{Inferensi pengolahan.} Kontribusi terkuat ialah operasionalisasi
fungsi pelayanan pada resolusi kota praja dan pengungkapan perbedaan
antara pola spasial mentah dengan hubungan regresi bersyarat. Kontribusi
teoritis terhadap borrowed size masih bersifat eksploratif karena
konstruk tersebut belum diidentifikasi secara langsung.

\subsection{Research Gap}\label{research-gap-21}

\textbf{Kesenjangan yang dinyatakan sumber.}

\begin{itemize}
\tightlist
\item
  Kajian geografis tentang fungsi pelayanan kota kecil pada skala kota
  praja masih kurang. (Hlm. 199, abstrak; hlm. 201, pendahuluan.)
\item
  Studi borrowed size lebih sering menunjukkan dampak kota besar
  terhadap wilayah sekitarnya daripada menunjukkan secara empiris bahwa
  kota kecil meminjam fungsi kota besar. (Hlm. 200, pendahuluan.)
\item
  Kajian borrowed size dan agglomeration shadow yang secara khusus
  memakai fungsi pelayanan masih sedikit. (Hlm. 200, pendahuluan.)
\item
  Di Tiongkok, keterbatasan data membuat riset kota kecil lebih banyak
  berfokus pada tata kelola atau memakai skala yang lebih agregat; riset
  skala kota praja tentang fungsi pelayanan jarang dilakukan. (Hlm. 201,
  pendahuluan.)
\item
  Belum ada metode spasial yang mapan untuk mengidentifikasi
  agglomeration shadow; regresi dapat menunjukkan asosiasi, tetapi tidak
  dengan sendirinya menentukan lokasi bayangan. (Hlm. 210, Bagian 3.2.)
\item
  Mekanisme borrowed size di wilayah polisentris, ukuran efeknya,
  kelompok kota yang mengalaminya, dan implikasinya bagi integrasi serta
  koordinasi spasial masih belum terjawab. (Hlm. 210, Bagian 3.2.)
\end{itemize}

\textbf{Inferensi pengolahan: kesenjangan yang tersisa setelah studi.}
Artikel belum memisahkan secara terukur manfaat borrowed size dari daya
tarik internal kota, limpahan umum metropolitan, dan tekanan
agglomeration shadow. Belum ada pengamatan arus pengguna atau penyedia
layanan lintas batas, dimensi waktu, pembanding kontrafaktual,
pengukuran kapasitas serta mutu layanan, maupun rancangan yang menangani
endogenitas lokasi. Kesenjangan tersebut membatasi kemampuan hasil untuk
menjelaskan mekanisme, bukan sekadar pola.

\subsection{Limitations}\label{limitations-21}

\textbf{Keterbatasan yang dinyatakan penulis.} (Hlm. 210, Bagian 3.2.)

\begin{itemize}
\tightlist
\item
  Data skala kota praja sulit diperoleh, sehingga indikator potensi
  internal, posisi dalam sistem perkotaan, dan aksesibilitas
  transportasi disederhanakan.
\item
  Potensi internal hanya diwakili terutama oleh penduduk dan PDB.
  Penulis menyebut pekerjaan, tingkat urbanisasi, struktur industri,
  serta modal asing sebagai variabel yang seharusnya dipertimbangkan.
  Modal Jerman di Taicang dan modal Jepang di Wuxi diberikan sebagai
  contoh faktor relevan yang tidak masuk model.
\item
  Borrowed size hanya dinilai secara awal dari signifikansi hasil
  regresi. Penulis mengakui masih perlu dijawab apakah kota besar
  benar-benar menyediakan borrowed size, apakah efek itu hanya terjadi
  pada kota kecil, bagaimana ukurannya ditentukan, dan apakah efek
  tersebut memperkuat integrasi serta pembangunan terkoordinasi.
\item
  Analisis hanya menjadikan Shanghai sebagai pusat. Pengaruh Nanjing,
  Hangzhou, dan kota lain yang setingkat tidak dimodelkan, padahal
  Jiangsu selatan berada dalam sistem polisentris.
\item
  Metode ilmiah untuk mengenali agglomeration shadow secara spasial
  belum mapan. Aturan lembah indeks yang digunakan artikel sendiri
  dinyatakan sebagai percobaan awal.
\end{itemize}

\textbf{Keterbatasan yang diturunkan dari isi file sebagai inferensi
pengolahan.}

\begin{itemize}
\tightlist
\item
  \textbf{Validitas konstruk:} jumlah POI berbobot tidak menunjukkan
  kapasitas, mutu, jangkauan, permintaan, penggunaan aktual, atau apakah
  sebuah layanan melayani penduduk dari kota lain. Indeks karena itu
  belum mengukur borrowed size secara langsung.
\item
  \textbf{Normalisasi:} penggunaan hitungan POI mentah membuat indeks
  sensitif terhadap jumlah penduduk, luas unit, kepadatan, dan perbedaan
  pencatatan Amap. Hasil tinggi di pusat kota sebagian dapat
  mencerminkan ukuran serta kepadatan unit.
\item
  \textbf{Definisi populasi:} penyamaan semua kota praja, subdistrik,
  pusat kabupaten, dan zona pengembangan sebagai kota kecil berpotensi
  mencampur unit dengan fungsi administratif dan tingkat urbanisasi yang
  sangat berbeda. Hal ini juga membuat temuan tingginya fasilitas di
  ``small towns'' dalam pusat kota perlu ditafsirkan hati-hati.
\item
  \textbf{Kehilangan observasi:} jumlah unit turun dari 248 menjadi 216
  dalam regresi tanpa penjelasan mengenai 32 unit yang tidak dianalisis,
  meskipun 20 nilai penduduk hilang disebut sudah diinterpolasi.
\item
  \textbf{Interpolasi:} bentuk interpolasi, variabel pembentuk,
  validasi, dan ketidakpastiannya tidak tersedia.
\item
  \textbf{AHP:} file tidak memuat matriks perbandingan berpasangan,
  identitas atau jumlah penilai, rasio konsistensi, proses konsensus,
  dan analisis sensitivitas bobot. Replikasi bobot tidak mungkin
  dilakukan dari TXT saja.
\item
  \textbf{Klasifikasi spasial:} batas kelas \emph{natural breaks}
  bergantung pada distribusi sampel. Tidak ada uji sensitivitas terhadap
  jumlah kelas, ambang alternatif, bandwidth KDE, radius buffer, ataupun
  definisi ketetanggaan pada kriteria ``dikelilingi''.
\item
  \textbf{Zona jarak:} Tabel 4 membagi seluruh 248 unit ke tiga kolom
  yang jumlahnya saling melengkapi, tetapi label ``dalam buffer 150 km''
  dan ``dalam buffer 200 km'' dapat pula dibaca secara kumulatif. File
  tidak memberi definisi formal yang sepenuhnya menghilangkan ambiguitas
  antara buffer kumulatif dan sabuk jarak.
\item
  \textbf{Model regresi:} penelitian potong lintang tidak membuktikan
  sebab-akibat dan rentan terhadap variabel terlewat, kausalitas
  terbalik, serta seleksi lokasi. File tidak melaporkan pemeriksaan
  autokorelasi spasial residual walaupun unit dan hasilnya bersifat
  spasial.
\item
  \textbf{Transparansi estimasi:} statistik uji White dan VIF, bentuk
  fungsi heteroskedastisitas, bobot WLS, transformasi seluruh prediktor,
  serta diagnostik residual tidak ditampilkan. Perbedaan skala rerata
  variabel terikat antara Tabel 5 dan Tabel 6 juga tidak dijelaskan.
\item
  \textbf{Arah aksesibilitas:} koefisien positif untuk waktu ke pintu
  tol menunjukkan hubungan dengan waktu yang lebih panjang, tetapi
  dibahas sebagai manfaat aksesibilitas. Tanpa definisi pengodean
  tambahan, arah substantif hasil ini tidak pasti.
\item
  \textbf{Representativitas statistik:} Tabel 5 memberi kolom simpangan
  baku, tetapi narasi menyebutnya galat baku dan memakai perbandingannya
  dengan rerata untuk menyimpulkan representativitas. Perbandingan
  tersebut tidak cukup untuk menunjukkan bahwa sampel regresi mewakili
  248 unit, terlebih mekanisme kehilangan datanya tidak dilaporkan.
\item
  \textbf{Validitas eksternal:} Jiangsu selatan merupakan wilayah maju
  dengan sejarah industrialisasi kota praja dan integrasi kuat dengan
  Shanghai. Polanya tidak otomatis berlaku bagi kota kecil di wilayah
  Tiongkok yang lebih perifer atau bagi negara lain.
\end{itemize}

\subsection{Challenges}\label{challenges-21}

\textbf{Tantangan yang dilaporkan sumber.} Tantangan utama ialah
kelangkaan data sosial-ekonomi yang konsisten pada skala kota praja.
Penulis harus menginterpolasi 20 nilai penduduk, menyederhanakan enam
prediktor, mengumpulkan waktu tempuh ke pintu tol melalui pemasukan
asal-tujuan satu per satu, dan mengandalkan POI sebagai pengganti
informasi pelayanan konvensional. Tidak adanya definisi nasional tentang
ukuran kota kecil juga diatasi secara operasional dengan memasukkan
semua kota praja dan subdistrik. (Hlm. 201, Bagian 1.1-1.2; hlm. 210,
Bagian 3.2.)

\textbf{Inferensi pengolahan.} Tantangan analitis terbesarnya ialah
membedakan empat penjelasan yang dapat menghasilkan pola serupa: ukuran
lokal, status sebagai pusat administratif, borrowed size dari Shanghai,
dan penarikan fungsi oleh aglomerasi. Konsentrasi POI di pusat dan
asosiasi positif jarak ke Shanghai tidak dengan sendirinya menentukan
aliran mekanisme. Tantangan replikasi juga besar karena keputusan
pembersihan POI, pembobotan AHP, pembentukan zona, hubungan
ketetanggaan, serta pembobotan FGLS tidak didokumentasikan secara
lengkap.

\subsection{Practical Implications}\label{practical-implications-21}

\textbf{Implikasi yang dinyatakan penulis.} Kota kecil merupakan
penghubung antara wilayah perkotaan dan hinterland perdesaan. Karena
itu, kota dengan indeks sangat rendah yang terbenam di antara unit
berindeks lebih tinggi perlu mendapat perhatian khusus dalam upaya
pemerataan pelayanan publik. Peta indeks dan area bayangan dimaksudkan
sebagai informasi latar ilmiah bagi pengambilan keputusan serta sebagai
pengalaman untuk wilayah serupa. (Hlm. 201, akhir pendahuluan; hlm. 210,
Bagian 3.2.)

\textbf{Inferensi pengolahan.} Bagi perencana, hasil dapat dipakai
sebagai alat penyaringan awal untuk menemukan lokasi yang mungkin
mengalami kekurangan fasilitas meskipun berada dekat pusat pertumbuhan.
Namun, penetapan prioritas investasi tidak seharusnya hanya memakai skor
POI. Verifikasi kebutuhan penduduk, kapasitas fasilitas, akses
perjalanan, mutu layanan, batas administratif, dan kemungkinan warga
memakai layanan di kota tetangga masih diperlukan sebelum intervensi.

\textbf{Inferensi pengolahan.} Temuan koeksistensi menyiratkan bahwa
kebijakan kedekatan atau konektivitas ke metropolitan tidak selalu
cukup. Sebagian kota mungkin memperoleh manfaat jaringan, sedangkan kota
lain kehilangan fungsi lokal karena kompetisi. Implikasi ini konsisten
dengan interpretasi penulis, tetapi artikel tidak menguji instrumen
kebijakan tertentu atau dampaknya.

\subsection{Applications}\label{applications-21}

\textbf{Aplikasi yang didukung langsung oleh sumber.}

\begin{itemize}
\tightlist
\item
  memetakan konsentrasi fasilitas pelayanan lintas sektor pada skala
  kota praja;
\item
  menghitung dan membandingkan indeks kuantitas pelayanan antarkota;
\item
  menyaring lokasi yang menjadi lembah fungsi pelayanan di antara
  pusat-pusat yang lebih kuat;
\item
  menguji asosiasi indeks pelayanan dengan kapasitas internal, posisi
  hierarkis, kedekatan metropolitan, dan akses jalan tol; serta
\item
  menyediakan bahan awal untuk evaluasi pemerataan pelayanan publik di
  Jiangsu selatan. (Hlm. 201-210, Bagian 1.2-3.2.)
\end{itemize}

\textbf{Inferensi pengolahan.} Alur POI-KDE-AHP-regresi dapat diadaptasi
untuk wilayah metropolitan lain hanya sebagai diagnosis eksploratif.
Transfer aplikasi memerlukan kalibrasi ulang klasifikasi POI dan bobot,
definisi unit yang konsisten, pengukuran akses aktual, uji sensitivitas,
serta validasi terhadap data pelayanan lapangan. File tidak melaporkan
aplikasi operasional, implementasi kebijakan, perangkat lunak yang
dibagikan, atau evaluasi pascapenerapan; informasi tersebut
\textbf{tidak disebutkan dalam file}.

\subsection{Future Research
(Optional)}\label{future-research-optional-21}

\textbf{Agenda yang dinyatakan penulis.} (Hlm. 210, Bagian 3.2.)

\begin{itemize}
\tightlist
\item
  memperkaya variabel potensi internal dengan pekerjaan, urbanisasi,
  struktur industri, dan investasi asing;
\item
  menguji secara lebih langsung apakah borrowed size benar-benar
  terjadi, kota mana yang mengalaminya, seberapa besar efeknya, serta
  apakah efek itu mendukung integrasi dan koordinasi spasial;
\item
  mengganti kerangka Shanghai tunggal dengan analisis polisentris yang
  memasukkan Nanjing, Hangzhou, dan kota relevan lain;
\item
  mengembangkan metode spasial yang lebih ilmiah untuk mengidentifikasi
  agglomeration shadow;
\item
  memberi perhatian lebih besar kepada kota berindeks sangat rendah yang
  tertanam di antara kota dengan fungsi lebih kuat; dan
\item
  menelusuri penyebab serta mekanisme borrowed size dan agglomeration
  shadow secara kualitatif untuk memberi umpan balik kepada teori.
\end{itemize}

\textbf{Inferensi pengolahan.} Agenda tersebut dapat diperkuat dengan
data panel dan perubahan fasilitas dari waktu ke waktu, data mobilitas
atau asal-tujuan pengguna layanan, ukuran kapasitas dan mutu, model
dependensi spasial, desain pembanding, serta pengujian kepekaan terhadap
definisi kota kecil, radius pengaruh, bobot AHP, dan kriteria area
bayangan. Usulan ini merupakan inferensi review, bukan agenda yang
tertulis eksplisit dalam artikel.

\subsection{Provenance Notes}\label{provenance-notes-21}

\begin{itemize}
\tightlist
\item
  \textbf{Basis akses:} review ini hanya menggunakan file TXT
  \texttt{source/journal/A06-sun-lu-zhang-niu-2022-evaluation-of-service-function-of-small-towns-in-china-from-the-perspective-of-borrowed-size-and-agglomeration-shadow-a-case-study-of-southern-jiangsu-province-10.18306-dlkxjz.2022.02.002.txt}.
  Tidak ada PDF, sumber eksternal, metadata daring, catatan review lain,
  atau pustaka yang dikutip artikel yang dibuka untuk menambah
  substansi.
\item
  \textbf{Cakupan baca:} seluruh 780 baris TXT dibaca, meliputi metadata
  ekstraksi, teks artikel hlm. 199-213, persamaan, enam tabel, lima
  gambar beserta keterangan tekstualnya, daftar pustaka, dan abstrak
  berbahasa Inggris.
\item
  \textbf{Asal teks:} blok metadata menyatakan TXT diekstrak pada 31
  Agustus 2026 dari PDF 15 halaman dengan \texttt{pdftotext\ -layout}.
  PDF asli bernama sama dengan kode A06 dan tercatat tidak terenkripsi
  serta tidak bertanda struktur (\emph{not tagged}). (TXT, blok
  \texttt{{[}PDF\ Metadata{]}}, baris 1-29.)
\item
  \textbf{Dasar identitas bibliografis:} metadata mesin PDF memberi
  judul internal \texttt{2-\/-\/-2021-0204} diikuti nama Sun Dongqi dan
  mencantumkan \texttt{Founder\ magtech2011} sebagai author. Karena
  nilai itu tampak sebagai metadata produksi, identitas artikel dalam
  review diambil dari blok judul, format sitasi, DOI, dan abstrak
  Inggris yang semuanya tersedia di TXT. (Hlm. 199 dan 213; TXT, baris
  7-13.)
\item
  \textbf{Locator:} nomor halaman dalam review merujuk pada nomor
  halaman jurnal yang tercetak, bukan nomor baris TXT. Tabel, gambar,
  bagian, dan persamaan disebut ketika tersedia agar klaim dapat
  ditelusuri.
\item
  \textbf{Keterbatasan ekstraksi:} tata letak dua kolom menyebabkan
  beberapa kalimat pada TXT terselang dengan kolom sebelah, dan rumus
  serta tabel tidak selalu tersusun linear. Angka yang dilaporkan di
  sini diperiksa terhadap label tabel dan narasi yang tersedia, tetapi
  isi visual rinci Gambar 1-5 tidak dapat diverifikasi melampaui
  keterangan dan pembahasan tekstual dalam TXT.
\item
  \textbf{Informasi tidak tersedia:} file tidak menyebut tanggal
  pengambilan POI, protokol pembersihan data, kode atau data terbuka,
  rincian interpolasi, alasan berkurangnya sampel regresi, matriks dan
  uji konsistensi AHP, statistik lengkap uji White dan VIF, bobot FGLS,
  pemeriksaan dependensi spasial, konflik kepentingan, ataupun
  pernyataan ketersediaan data. Status telaah sejawat juga tidak
  diverifikasi dari luar TXT.
\item
  \textbf{Batas epistemik:} semua bagian berlabel \textbf{Laporan
  sumber}, \textbf{Temuan yang dilaporkan sumber}, atau
  \textbf{Kesimpulan yang dinyatakan penulis} adalah parafrasa isi
  artikel. Bagian \textbf{Interpretasi penulis} merekam makna yang
  diberikan para penulis terhadap hasilnya. Bagian \textbf{Inferensi
  pengolahan} adalah evaluasi analitis berdasarkan isi dan
  ketidakkonsistenan internal TXT, bukan klaim yang dibuat artikel atau
  informasi eksternal.
\end{itemize}

\section{Review A07 - Zhao et
al.~(2022)}\label{review-a07---zhao-et-al.-2022}

Status: Review literatur berbasis pembacaan seluruh TXT hasil ekstraksi
PDF.

\subsection{Identitas Sumber}\label{identitas-sumber-10}

\begin{itemize}
\tightlist
\item
  \textbf{Kode:} A07.
\item
  \textbf{Judul asli:}
  \emph{广佛都市圈网络外部性的城镇借用规模绩效检验}.
\item
  \textbf{Judul bahasa Inggris:} \emph{Examining performance of urban
  borrowed size based on the towns' network externalities of
  Guangzhou-Foshan metropolitan areas}.
\item
  \textbf{Penulis:} Zhao Miaoxi, Wang Yankai, Hu Yuke, Guo Zhensong, dan
  Wei Zhaobin. Blok afiliasi mengaitkan tiga penulis pertama dengan
  South China University of Technology, Guo Zhensong dengan Urban
  Planning \& Design Institute of Shenzhen, dan Wei Zhaobin dengan
  Huizhou University. Guo Zhensong ditandai sebagai penulis
  korespondensi. (Lokator: hlm. 2367, blok judul dan catatan penulis;
  hlm. 2384, judul dan abstrak bahasa Inggris; TXT baris 36-39, 76-82,
  753-759.)
\item
  \textbf{Jurnal:} \emph{Geographical Research} atau \emph{地理研究},
  volume 41, nomor 9, September 2022, hlm. 2367-2384. (Lokator: TXT
  baris 30-31, 82-83, 749.)
\item
  \textbf{DOI:} \texttt{10.11821/dlyj020211036}. (Lokator: hlm. 2367;
  TXT baris 57.)
\item
  \textbf{Tanggal naskah:} diterima 5 November 2021 dan diterima untuk
  publikasi 11 April 2022. (Lokator: TXT baris 76.)
\item
  \textbf{Pendanaan:} Guangdong Science and Technology Plan Soft Science
  Project \texttt{2019A101002076} dan National Natural Science
  Foundation of China \texttt{51711530711}. (Lokator: TXT baris 77.)
\item
  \textbf{Kata kunci:} borrowed size, agglomeration shadow, network
  externalities, urban network, dan metropolitan area. (Lokator: hlm.
  2367 dan 2384; TXT baris 55-56, 789-790.)
\item
  \textbf{Wilayah studi:} metropolitan Guangzhou-Foshan di kawasan Delta
  Sungai Mutiara, dengan Guangzhou dan Foshan sebagai inti yang
  berdekatan. (Lokator: hlm. 2371-2372, Bagian 3; TXT baris 232-273.)
\item
  \textbf{Periode data utama:} data perusahaan baru terdaftar dan
  jaringan perusahaan untuk 2017-2019; populasi yang dipakai sebagai
  ukuran awal berasal dari sensus penduduk keenam nasional. (Lokator:
  hlm. 2373-2374, Bagian 4.2; TXT baris 332-347.)
\item
  \textbf{Status telaah sejawat:} Tidak disebutkan dalam file. Artikel
  memang menyampaikan ucapan terima kasih kepada penelaah anonim, tetapi
  TXT tidak memberi pernyataan status publikasi yang dapat diverifikasi
  secara terpisah. (Lokator: TXT baris 629-631.)
\end{itemize}

\subsection{Summarized Abstract}\label{summarized-abstract-22}

\textbf{Laporan sumber.} Artikel menguji apakah jaringan eksternalitas
kota memungkinkan town atau kota praja memanfaatkan borrowed size untuk
melampaui batas geografis dan hierarki ukuran, lalu meningkatkan
pertumbuhan ekonomi bersama di metropolitan Guangzhou-Foshan. Penelitian
menggunakan ukuran jumlah perusahaan baru terdaftar dan titik jaringan
perusahaan sebagai indikator kinerja relatif terhadap ukuran penduduk
yang sudah ada. Tahap pertama menggunakan deviasi terhadap rerata dan
kelipatan simpangan baku untuk membaca regional performance serta
agglomeration shadow. Tahap kedua menggunakan regresi berganda untuk
menguji keterkaitan ukuran town, borrowed size, aksesibilitas jaringan
transportasi, dan limpahan inovasi teknologi dengan pertumbuhan
perusahaan. (Lokator: hlm. 2367, abstrak; TXT baris 42-57.)

\textbf{Laporan sumber.} Tiga hasil utama yang diringkas abstrak ialah:
ambang 100.000 penduduk membedakan karakter kinerja, dengan town di atas
ambang tersebut menunjukkan borrowed-size performance dan town di
bawahnya menunjukkan agglomeration shadow; kinerja jumlah perusahaan
baru dan titik jaringan perusahaan membentuk pola ruang bercincin,
dengan pola titik jaringan lebih menonjol; dan kinerja jumlah perusahaan
baru paling erat berkaitan dengan ukuran town yang telah ada, disusul
secara berurutan oleh aksesibilitas hub transportasi regional, borrowed
size, borrowed performance, dan paten kerja sama lintas-town. (Lokator:
hlm. 2367, abstrak; hlm. 2375-2377, Bagian 5.1; TXT baris 47-54,
373-442.)

\textbf{Interpretasi penulis.} Penulis membaca peningkatan jaringan
transportasi regional dan saluran kerja sama teknologi sebagai kondisi
yang membantu pemusatan faktor dan memengaruhi pembangunan urbanisasi
melalui efek eksternalitas jaringan. Mereka memosisikan borrowed size
dan agglomeration shadow sebagai dua kemungkinan hasil dari hubungan
antartown, bukan sebagai proses yang harus selalu menghasilkan manfaat
bagi semua unit. (Lokator: hlm. 2367, abstrak; hlm. 2368-2371, Bagian
2.1-2.3; TXT baris 122-130, 198-229.)

\textbf{Inferensi review.} Ringkasan tersebut lebih tepat dipahami
sebagai bukti pola relatif dan hubungan statistik daripada identifikasi
kausal borrowed size. \texttt{PER} dibangun dari rasio indikator
perusahaan terhadap populasi yang sudah ada dan penyimpangan dari
rerata; ukuran itu belum secara langsung mengamati fungsi yang
benar-benar dipinjam, aliran pengguna, atau keputusan relasional
perusahaan. (Lokator dasar: hlm. 2372-2373, Persamaan 1-2; hlm.
2378-2380, Tabel 3-6; TXT baris 279-330, 462-579.)

\subsection{Problem Statement}\label{problem-statement-22}

\textbf{Laporan sumber.} Kebijakan pembangunan metropolitan dan jalur
urbanisasi baru menjadikan wilayah padat town sebagai fokus. Teori
aglomerasi menjelaskan manfaat konsentrasi, tetapi teori pusat-pinggiran
dan ketergantungan jalur sulit menjelaskan mengapa sebagian town kecil
di wilayah maju tetap tumbuh. Artikel menyoroti kemungkinan bahwa town
dapat memanfaatkan koneksi dengan kota besar, sementara town lain dapat
mengalami penarikan sumber daya atau agglomeration shadow. (Lokator:
hlm. 2367-2368, Bagian 1; TXT baris 62-98.)

\textbf{Laporan sumber.} Penulis menyatakan bahwa riset empiris yang ada
cenderung melihat performa regional secara agregat dan kurang
menganalisis interaksi spasial internal, khususnya borrowed-size effect
pada tingkat town. Dalam konteks metropolitan Guangzhou-Foshan, mereka
menilai belum jelas bagaimana ukuran town yang telah ada, borrowed size,
lokasi transportasi, dan jaringan kerja sama teknologi saling
berhubungan dengan organisasi spasial perusahaan. (Lokator: hlm.
2368-2371, Bagian 1 dan 2.3; TXT baris 99-107, 198-206.)

\textbf{Laporan sumber.} Masalah penelitian dirumuskan dalam dua
pertanyaan atau proposisi kerja:

\begin{itemize}
\tightlist
\item
  Apakah konektivitas jaringan yang lebih kuat, ketika dipadukan dengan
  ukuran awal town, borrowed size, aksesibilitas transportasi, dan
  jaringan kerja sama teknologi, berkaitan dengan pertumbuhan perusahaan
  baru serta perbedaan regional performance dan agglomeration shadow?
  (Lokator: hlm. 2370-2371, Bagian 2.3, Gambar 1; TXT baris 207-211.)
\item
  Apakah jaringan transportasi yang memudahkan perpindahan tenaga kerja,
  modal, dan teknologi meningkatkan kemungkinan kantor pusat menempatkan
  cabang di lokasi lain, sehingga hubungan fungsional perusahaan
  membentuk kinerja pertumbuhan atau bayangan aglomerasi antartown?
  (Lokator: hlm. 2370-2371, Bagian 2.3, Gambar 2; TXT baris 212-229.)
\end{itemize}

\textbf{Inferensi review.} Problem statement memiliki dua lapisan
analitis yang perlu dipisahkan. Lapisan pertama adalah pengukuran apakah
pertumbuhan perusahaan suatu town lebih tinggi atau lebih rendah
daripada yang diharapkan berdasarkan ukuran penduduknya. Lapisan kedua
adalah penjelasan apakah perbedaan itu berasal dari borrowed size,
aglomerasi lokal, kedekatan ke hub, kerja sama teknologi, atau tekanan
kompetisi. Desain yang ditampilkan artikel lebih langsung menjawab
lapisan pertama daripada membuktikan mekanisme lapisan kedua. (Lokator
dasar: hlm. 2372-2374, Bagian 4.1-4.2; TXT baris 276-369.)

\subsection{Objectives}\label{objectives-22}

\textbf{Laporan sumber.} Tujuan penelitian yang dapat ditarik dari
abstrak, pertanyaan penelitian, dan rancangan analisis ialah:

\begin{itemize}
\tightlist
\item
  menguji kinerja borrowed size town dalam kerangka eksternalitas
  jaringan kota di metropolitan Guangzhou-Foshan;
\item
  mengukur regional performance berdasarkan jumlah perusahaan baru
  terdaftar dan titik jaringan perusahaan relatif terhadap populasi yang
  telah ada;
\item
  mengidentifikasi perbedaan kinerja dan agglomeration shadow menurut
  kelas ukuran penduduk serta lokasi spasial;
\item
  menguji peran gabungan ukuran town, borrowed size, borrowed
  performance, aksesibilitas jaringan transportasi lokal dan hub
  regional, serta limpahan inovasi teknologi;
\item
  membandingkan pola perusahaan baru dengan pola hubungan kantor
  pusat-cabang lintas-town; dan
\item
  menjalankan analisis tambahan untuk manufaktur, perdagangan dan
  sirkulasi komersial, persewaan dan jasa bisnis, serta jasa teknologi
  dan perangkat lunak. (Lokator: hlm. 2367, abstrak; hlm. 2370-2373,
  Bagian 2.3 dan 4.1; hlm. 2374, Bagian 4.2; hlm. 2378-2380, Bagian 5.2;
  TXT baris 42-54, 223-229, 276-330, 364-369, 443-579.)
\end{itemize}

\textbf{Interpretasi penulis.} Tujuan substantifnya bukan hanya
menentukan town yang tumbuh cepat, melainkan menjelaskan bagaimana dua
bentuk jaringan, yaitu jaringan transportasi fisik dan jaringan kerja
sama pengetahuan, berkopel dengan ukuran town untuk menghasilkan manfaat
jaringan atau bayangan aglomerasi. (Lokator: hlm. 2369-2371, Bagian
2.1-2.3; TXT baris 122-178, 198-229.)

\textbf{Inferensi review.} Artikel tidak menyatakan hipotesis statistik
formal dengan arah koefisien untuk setiap variabel. Istilah
\texttt{假说模型} atau model hipotesis dipakai untuk kerangka
konseptual, sedangkan pertanyaan yang ditampilkan menjadi dasar
pengujian regresi. Karena itu, tujuan artikel paling aman dirumuskan
sebagai pengujian asosiasi dan mekanisme yang diusulkan, bukan pengujian
hipotesis kausal yang telah dipraregistrasi atau dinyatakan secara
formal. (Lokator: hlm. 2370-2371, Bagian 2.3; TXT baris 198-229.)

\subsection{Literature Survey}\label{literature-survey-22}

\textbf{Laporan sumber.} Tinjauan pustaka dimulai dari teori ekonomi
aglomerasi. Aglomerasi dipaparkan sebagai sumber eksternalitas positif,
increasing returns, pasar yang tebal, efek keterkaitan, limpahan
pengetahuan, dan ekonomi eksternal murni. New Economic Geography
menjelaskan konsentrasi pusat-periferi, tetapi ketergantungan jalur yang
menyertainya dapat menyiratkan keterkuncian town kecil. Artikel memakai
fakta pertumbuhan sebagian town kecil di wilayah metropolitan Eropa
sebagai alasan untuk mencari penjelasan alternatif. (Lokator: hlm.
2367-2368, Bagian 1; TXT baris 68-86.)

\textbf{Laporan sumber.} Perubahan dari masyarakat berbasis hinterland
menuju masyarakat jaringan dipakai untuk memperkenalkan pergeseran dari
teori \texttt{central\ place} ke \texttt{central\ flow}. Dalam pandangan
ini, perkembangan town tidak hanya bergantung pada sumber daya lokal,
tetapi juga pada kemudahan menghubungkan diri dengan dunia luar. Artikel
mengutip pembahasan tentang jaringan transportasi, jaringan perusahaan,
jaringan sosial, dan jaringan inovasi untuk menunjukkan bahwa fungsi
kota dapat dibentuk oleh hubungan lintas batas. (Lokator: hlm.
2368-2370, Bagian 1 dan 2.1; TXT baris 87-124.)

\textbf{Laporan sumber.} Eksternalitas jaringan kota dirujukkan kepada
Capello sebagai manfaat koordinasi dan komplementaritas yang timbul dari
jaringan fungsional antarkota. Artikel juga memakai gagasan network
value yang meningkat dengan jumlah partisipan, lalu membedakan jaringan
transportasi sebagai jaringan infrastruktur yang masih dibatasi kondisi
fisik dari jaringan inovasi sebagai hubungan informasi yang dapat
melintasi batas geografis. (Lokator: hlm. 2369-2370, Bagian 2.1; TXT
baris 112-178.)

\textbf{Laporan sumber.} Literatur transportasi yang dirangkum tidak
sepakat secara tunggal. Aksesibilitas dapat memperluas pasar dan memicu
aglomerasi, tetapi penurunan biaya perjalanan juga dapat memperkuat
pusat dan menarik faktor dari town nonpusat. Studi yang dirujuk artikel
memperlihatkan hasil yang berbeda menurut skala: efek jaringan
transportasi dapat mendorong ekonomi kota tingkat prefektur, tetapi
kereta cepat dapat memperburuk efek hisap pusat dalam jangka pendek pada
sebagian kota kecil dan menengah. (Lokator: hlm. 2369, Bagian 2.1.1; TXT
baris 131-155.)

\textbf{Laporan sumber.} Untuk jaringan inovasi, artikel membedakan
eksternalitas spesialisasi, eksternalitas diversifikasi, dan kerja sama
lintas wilayah. Paten lintas-town dianggap dapat merekam limpahan
pengetahuan yang menyeberangi batas geografis dan lebih merefleksikan
hubungan berbasis kontrak atau kepentingan bersama. Paten dengan tiga
atau lebih kategori aplikasi dipakai untuk menangkap teknologi umum atau
hubungan lintas bidang di dalam town. (Lokator: hlm. 2369-2370 dan
2373-2374, Bagian 2.1.2 dan 4.2; TXT baris 156-178, 346-363.)

\textbf{Laporan sumber.} Borrowed size ditelusuri ke Alonso: town kecil
yang dekat dengan pusat populasi metropolitan dapat memperlihatkan
sebagian karakter kota besar dengan memanfaatkan sumber daya dan fungsi
kota tersebut. Meijers dan kolega dipakai untuk membedakan borrowed size
dari agglomeration shadow. Borrowed size menggambarkan manfaat yang
diterima town kecil dari skala kota besar, sedangkan agglomeration
shadow menggambarkan pertumbuhan yang lebih rendah daripada ekspektasi
untuk ukuran town tersebut karena kedekatan dengan kota besar. (Lokator:
hlm. 2370, Bagian 2.2; TXT baris 179-197.)

\textbf{Laporan sumber.} Pengukuran borrowed performance dalam studi
terdahulu yang dirangkum memakai pendapatan, pertumbuhan pekerjaan,
pendapatan per kapita, atau upah. Artikel menyatakan bahwa penelitian
sebelumnya lebih banyak memakai wilayah perkotaan atau city-region di
Eropa dan analisis kota atau urban agglomeration di Tiongkok, sementara
kinerja relatif pada tingkat town dengan ukuran dan lokasi yang berbeda
belum cukup diperiksa. (Lokator: hlm. 2370, Bagian 2.2; TXT baris
186-197.)

\textbf{Interpretasi penulis.} Penulis menggabungkan teori Capello dan
Camagni tentang eksternalitas jaringan dengan konsep borrowed size untuk
membaca town sebagai bagian dari sistem relasional. Mereka menilai bahwa
kerangka tersebut lebih sesuai daripada menjelaskan pertumbuhan hanya
melalui Marshallian agglomeration atau Jacobian diversity, karena
perusahaan dapat terhubung secara bersamaan melalui kedekatan fisik dan
hubungan lintas wilayah. (Lokator: hlm. 2369-2371, Bagian 2.1-2.3; TXT
baris 144-178, 198-229.)

\textbf{Inferensi review.} Tinjauan pustaka dalam TXT bersifat naratif
dan argumentatif. Tidak ada keterangan tentang strategi pencarian,
kriteria inklusi, rentang waktu korpus, atau penilaian mutu studi yang
dirujuk. Kelengkapan survei literatur karena itu \textbf{Tidak
disebutkan dalam file}; daftar pustaka yang tampil hanya menjadi bukti
tentang karya yang dikutip artikel, bukan hasil verifikasi independen
atas temuan karya-karya tersebut. (Lokator: TXT baris 633-748.)

\subsection{Method Used}\label{method-used-22}

\textbf{Laporan sumber.} Penelitian menggunakan pendekatan kuantitatif
dengan dua tahap utama. Unit kawasan yang disebutkan berjumlah 73,
sedangkan tabel variabel dan analisis town menggunakan 66 unit. Tahap
pertama mengukur regional relative performance dan memetakan
persebarannya. Tahap kedua menguji kombinasi faktor yang diasumsikan
membentuk pertumbuhan perusahaan melalui regresi berganda. (Lokator:
hlm. 2373-2374, Bagian 4; hlm. 2375-2380, Bagian 5; TXT baris 276-369,
371-579.)

\textbf{Laporan sumber: pengukuran kinerja relatif.} Untuk setiap town,
jumlah perusahaan baru atau titik jaringan perusahaan tahun 2017-2019
dibagi dengan populasi yang telah ada sehingga menghasilkan
\texttt{E\_i}. Nilai tersebut dibandingkan dengan rerata wilayah dan
dinyatakan sebagai residual yang dinormalisasi dengan simpangan baku,
disebut \texttt{PER\_i} atau regional performance. Persamaan (1) dan (2)
digunakan untuk menghitung ukuran tersebut. Nilai positif dipakai untuk
menunjukkan kinerja di atas ekspektasi regional dan nilai negatif untuk
kinerja di bawahnya, yang dikaitkan dengan agglomeration shadow.
(Lokator: hlm. 2372-2373, Bagian 4.1, Persamaan 1-2; TXT baris 279-304.)

\textbf{Laporan sumber: analisis spasial.} Town dikelompokkan
berdasarkan populasi: kurang dari 50.000, 50.000-100.000,
100.000-200.000, 200.000-500.000, dan lebih dari 500.000. Nilai kinerja
dibandingkan antarkelompok, kemudian GIS dipakai untuk membaca pola
spasial dan pola berlapis atau bercincin. (Lokator: hlm. 2372-2373 dan
2375-2377, Bagian 4.1 dan 5.1, Tabel 2 serta Gambar 6-7; TXT baris
299-304, 373-442.)

\textbf{Laporan sumber: model regresi.} Penulis membandingkan regresi
berbobot geografis dan regresi berganda berdasarkan AICC. AICC regresi
berbobot geografis adalah 105,22, sedangkan AICC regresi berganda 98,36;
karena nilai yang lebih rendah dianggap lebih baik, regresi berganda
dipilih. Persamaan (3) memuat \texttt{ln\ PER\_i} sebagai variabel
terikat dan ukuran penduduk, borrowed size, borrowed performance,
aksesibilitas transportasi lokal, aksesibilitas hub regional, paten
teknologi umum multi-kategori, serta paten kerja sama lintas-town
sebagai kandidat penjelas. (Lokator: hlm. 2373, Bagian 4.1, Persamaan 3;
TXT baris 305-317.)

\textbf{Laporan sumber: konstruksi borrowed size dan borrowed
performance.} Borrowed size town \texttt{i} dihitung sebagai jumlah
populasi town lain \texttt{j} dibagi jarak ke town tersebut, sesuai
Persamaan (4). Borrowed performance dihitung secara analog dengan
mengganti populasi town lain dengan jumlah perusahaan terdaftar atau
titik jaringan perusahaan, sesuai Persamaan (5). (Lokator: hlm. 2373,
Bagian 4.1, Persamaan 4-5; TXT baris 318-330.)

\textbf{Laporan sumber: pemilihan model.} Semua variabel diubah ke
logaritma natural untuk standardisasi, lalu regresi berganda dan regresi
bertahap digunakan untuk mengekstrak kombinasi variabel yang mempunyai
peran penjelas. Variabel kepadatan penduduk semula dipertimbangkan
tetapi dikeluarkan karena masalah multikolinearitas dengan nilai VIF
lebih dari 5. Regresi tambahan dijalankan untuk empat kelompok industri.
(Lokator: hlm. 2373-2374, Bagian 4.1-4.2; TXT baris 305-317, 364-369.)

\textbf{Inferensi review.} Berdasarkan satu rangkaian periode perusahaan
2017-2019 dan ukuran penduduk sensus yang disebut tanpa seri waktu
pembanding, desain empiris ini bersifat potong lintang atau setidaknya
tidak disajikan sebagai analisis panel. Regresi yang dipilih berdasarkan
AICC dapat membandingkan kecocokan model, tetapi tidak dengan sendirinya
membuktikan arah sebab-akibat. Artikel tidak menyebut eksperimen,
kuasi-eksperimen, variabel instrumental, atau strategi kontrafaktual.
(Lokator dasar: hlm. 2372-2374, Bagian 4.1-4.2; TXT baris 279-369.)

\subsection{Dataset}\label{dataset-22}

\textbf{Laporan sumber.} Dataset yang dipakai untuk membentuk variabel
penelitian meliputi:

\begin{itemize}
\tightlist
\item
  \textbf{Data perusahaan:} basis data pendaftaran perusahaan industri
  dan komersial baru untuk 2017-2019 dari Qichacha atau
  \texttt{qcc.com}. Dari data tersebut penulis membentuk jumlah
  perusahaan baru dan jaringan kantor pusat-cabang sebagai indikator
  hubungan antartown. (Lokator: hlm. 2373, Bagian 4.2; TXT baris
  332-336.)
\item
  \textbf{Data penduduk:} populasi penduduk menetap dari sensus penduduk
  keenam nasional, sebagai ukuran populasi awal town. (Lokator: hlm.
  2373, Bagian 4.2; TXT baris 334-336.)
\item
  \textbf{Data jaringan transportasi lokal:} waktu perjalanan rata-rata
  tertimbang untuk menggambarkan aksesibilitas dari satu town ke town
  lain. Jumlah penduduk town menjadi bobot; perhitungan menggunakan
  penduduk town lain dan tidak memasukkan penduduk town itu sendiri
  dalam komponen tujuannya. (Lokator: hlm. 2373, Bagian 4.2; TXT baris
  337-342.)
\item
  \textbf{Data jaringan hub regional:} waktu tempuh dari tiap town
  menuju Bandara Baiyun dan Stasiun Guangzhou Selatan. Artikel
  menjelaskan bahwa kapasitas fasilitas dapat dipakai sebagai bobot
  dalam jaringan hub, tetapi untuk indikator penelitian yang disebutkan
  secara spesifik, dua simpul tersebut menjadi tujuan pengukuran.
  (Lokator: hlm. 2373, Bagian 4.2; TXT baris 342-345.)
\item
  \textbf{Data inovasi:} data paten penemuan baru dari Wanfang
  (\texttt{wanfangdata.com.cn}). Paten dibedakan menurut kategori
  aplikasi dan lokasi institusi pemohon; paten dengan tiga atau lebih
  kategori aplikasi dipakai untuk teknologi umum multi-kategori,
  sedangkan kerja sama antarlembaga yang melintasi unit spasial town
  dipakai untuk paten kerja sama lintas-town. (Lokator: hlm. 2373-2374,
  Bagian 4.2; TXT baris 346-363.)
\item
  \textbf{Data spasial:} seluruh variabel ditempatkan pada unit town
  atau street menggunakan ArcGIS untuk menghasilkan peta dan analisis
  spasial. (Lokator: hlm. 2374, Bagian 4.2; TXT baris 364-369.)
\end{itemize}

\textbf{Laporan sumber.} Tabel 1 mencatat 66 observasi. Statistik
deskriptif yang ditampilkan adalah: variabel terikat jumlah perusahaan
baru \texttt{ln\ PER\_i} memiliki rerata 8,908 dan simpangan baku 1,309;
titik jaringan perusahaan \texttt{ln\ PER\_i} memiliki rerata 5,852 dan
simpangan baku 1,412; \texttt{ln\ Pop} 11,643 dan 0,962;
\texttt{ln\ Bor\_Size\_ij} 12,576 dan 0,308; borrowed performance
berbasis jumlah perusahaan terdaftar 10,185 dan 0,340; borrowed
performance berbasis titik jaringan 7,310 dan 0,354; \texttt{ln\ Trans}
-0,159 dan 0,742; \texttt{ln\ Hub} 0,778 dan 0,200;
\texttt{ln\ Mul\_pat} 3,293 dan 2,231; serta \texttt{ln\ Cross\_pat}
5,449 dan 1,942. (Lokator: hlm. 2376, Tabel 1; TXT baris 397-409.)

\textbf{Inferensi review.} Rincian berikut tidak tersedia dalam TXT:
tanggal pengunduhan data Qichacha dan Wanfang, parameter pencarian,
aturan pembersihan atau deduplikasi, validasi koordinat, tingkat
kelengkapan basis data, penanganan nilai hilang, serta kode atau data
replikasi. Karena itu, angka observasi dan definisi umum dapat
dilaporkan, tetapi proses pembentukan dataset tidak sepenuhnya dapat
diulang hanya dari file ini. (Lokator dasar: hlm. 2373-2374, Bagian 4.2;
TXT baris 332-369.)

\textbf{Inferensi review: ketidakkonsistenan label paten.} Uraian model
pada Bagian 4.1 menyebut \texttt{ln\ Mul\_pat} sebagai jumlah paten
teknologi umum multi-kategori dan \texttt{ln\ Cross\_pat} sebagai jumlah
paten kerja sama lintas-town. Namun, deskripsi indikator pada Tabel 1
tampak membalik keduanya: \texttt{ln\ Mul\_pat} dijelaskan sebagai paten
dengan pemohon dari town berbeda, sedangkan \texttt{ln\ Cross\_pat}
dijelaskan sebagai paten dengan tiga atau lebih kategori aplikasi. TXT
tidak menyediakan koreksi atau keterangan bahwa label tersebut tertukar.
Review ini mempertahankan nama variabel sesuai bagian yang dikutip dan
tidak mengklaim kepastian pemetaan nama ke konstruk. (Lokator
pembanding: hlm. 2373, Bagian 4.1; hlm. 2376, Tabel 1; TXT baris 308-320
dan 397-409.)

\subsection{Population / Sample}\label{population-sample-22}

\textbf{Laporan sumber.} Populasi wilayah studi adalah metropolitan
Guangzhou-Foshan. Guangzhou dan Foshan digambarkan sebagai metropolitan
yang bersebelahan secara spasial, memiliki lingkungan alam dan budaya
yang serupa, serta memiliki jaringan transportasi dan informasi yang
berkembang. Pada 2019, total GDP metropolitan tersebut dilaporkan
sebesar 35770 亿元, dengan satuan sebagaimana tertulis dalam TXT, dan
jumlah penduduknya 24,17 juta. (Lokator: hlm. 2371-2372, Bagian 3; TXT
baris 232-256.)

\textbf{Laporan sumber.} Artikel menyebut 73 unit spasial dalam
metropolitan Guangzhou-Foshan, dengan 66 unit berupa town periferal di
luar wilayah inti kota. Analisis indikator pada Tabel 1 dan analisis
perbandingan ukuran pada Tabel 2 menggunakan 66 observasi. Kelompok
populasi 66 town tersebut terdiri atas 15 unit berpenduduk kurang dari
50.000, 15 unit berpenduduk 50.000-100.000, 19 unit berpenduduk
100.000-200.000, 13 unit berpenduduk 200.000-500.000, dan 4 unit
berpenduduk lebih dari 500.000. (Lokator: hlm. 2373, Bagian 4.2; hlm.
2375-2376, Tabel 1-2; TXT baris 332-336, 397-427.)

\textbf{Laporan sumber.} Unit analisis disebut sebagai town atau street.
Town yang dianalisis mencakup ukuran dan lokasi yang berbeda, termasuk
town di sekitar pusat kota, kawasan pengembangan, dan town periferal
yang lebih jauh. Artikel menggunakan 66 unit itu sebagai observasi dalam
regresi jumlah perusahaan baru dan titik jaringan perusahaan. (Lokator:
hlm. 2374-2376, Bagian 4.2 dan 5.1; TXT baris 364-369, 373-442.)

\textbf{Inferensi review.} Prosedur pengambilan sampel probabilistik,
kriteria rinci pemilihan 66 unit dari 73 unit, dan perlakuan terhadap
tujuh unit lainnya \textbf{Tidak disebutkan dalam file}. Dengan
demikian, angka 66 paling aman dibaca sebagai himpunan unit town yang
dipakai dalam analisis, bukan sebagai sampel acak yang mewakili semua
metropolitan atau semua town di Tiongkok. (Lokator dasar: hlm.
2373-2376, Bagian 4.2 dan Tabel 1-2; TXT baris 332-336, 397-427.)

\textbf{Inferensi review.} Cakupan eksternal penelitian terbatas pada
satu metropolitan yang secara historis telah terindustrialisasi dan
memiliki konektivitas kuat. Artikel tidak melaporkan pembanding
metropolitan lain, sehingga transfer langsung temuan ke wilayah dengan
struktur ekonomi, jaringan, atau tingkat urbanisasi berbeda
\textbf{Tidak disebutkan dalam file} dan tidak dapat dipastikan dari
desain ini. (Lokator dasar: hlm. 2371-2372, Bagian 3; TXT baris
232-273.)

\subsection{Dependent Variables}\label{dependent-variables-22}

\textbf{Laporan sumber.} Penelitian memakai dua variabel terikat yang
sama-sama ditujukan untuk menangkap performa ekonomi town, tetapi
melalui objek perusahaan yang berbeda:

\begin{itemize}
\tightlist
\item
  \textbf{Jumlah perusahaan baru terdaftar:} jumlah perusahaan industri
  dan komersial yang baru terdaftar pada 2017-2019 di tiap town. Artikel
  menyebut jumlah perusahaan sebagai indikator umum dalam penelitian
  eksternalitas dan menggunakannya untuk membaca pertumbuhan perusahaan
  town. (Lokator: hlm. 2372-2373, Bagian 4.1; hlm. 2376, Tabel 1; TXT
  baris 284-304, 397-401.)
\item
  \textbf{Titik jaringan perusahaan atau jaringan perusahaan
  lintas-town:} jumlah perusahaan baru terdaftar yang kantor pusat dan
  cabangnya berada di town yang berbeda. Indikator ini dimaksudkan untuk
  merepresentasikan intensitas hubungan fungsional kantor pusat-cabang
  antartown. (Lokator: hlm. 2373-2374, Bagian 4.2; hlm. 2376, Tabel 1;
  TXT baris 332-336, 397-405.)
\end{itemize}

\textbf{Laporan sumber.} Untuk kedua objek tersebut, penulis terlebih
dahulu membentuk \texttt{E\_i} sebagai rasio indikator perusahaan
terhadap populasi yang sudah ada, lalu menghitung penyimpangan
terstandardisasi dari rerata regional sebagai \texttt{PER\_i}. Regresi
memakai bentuk yang diberi label \texttt{ln\ PER\_i}. Nilai kinerja
kelompok pada Tabel 2 menunjukkan angka negatif untuk town di bawah
kinerja relatif regional dan angka positif untuk sebagian kelompok di
atasnya. (Lokator: hlm. 2372-2373, Persamaan 1-2; hlm. 2375-2376, Tabel
2; TXT baris 279-304, 413-427.)

\textbf{Laporan sumber.} Pada Tabel 2, nilai kinerja jumlah perusahaan
baru untuk kelompok populasi berturut-turut adalah -0,429; -0,327;
0,471; 0,060; dan 0,398. Nilai kinerja titik jaringan perusahaan
berturut-turut adalah -0,676; -0,302; 0,236; 0,301; dan 1,572. Artikel
memakai pola ini untuk menyatakan bahwa 100.000 penduduk merupakan titik
pembeda utama, meskipun kelompok 200.000-500.000 memperlihatkan lembah
relatif pada ukuran tertentu. (Lokator: hlm. 2375-2376, Tabel 2 dan
Bagian 5.1.1; TXT baris 388-394, 413-427.)

\textbf{Inferensi review.} Konstruk dependent variable tidak sepenuhnya
transparan. TXT menjelaskan \texttt{PER\_i} sebagai residual
terstandardisasi dari rasio, tetapi Tabel 1 mendeskripsikan
\texttt{ln\ PER\_i} dengan definisi hitungan perusahaan dan melaporkan
rerata positif; file tidak menjelaskan secara eksplisit apakah log
diterapkan pada hitungan mentah sebelum pembentukan rasio, pada rasio,
atau pada residual setelah pembentukan \texttt{PER\_i}. Karena residual
dapat bernilai negatif pada Tabel 2, urutan transformasi dan penanganan
nilai nonpositif perlu diklarifikasi sebelum koefisien dianggap dapat
direplikasi. (Lokator pembanding: hlm. 2372-2373, Persamaan 1-3; hlm.
2376, Tabel 1-2; TXT baris 279-330, 397-427.)

\textbf{Inferensi review.} Kedua variabel terikat merupakan proksi
aktivitas perusahaan yang tercatat. Keduanya tidak secara langsung
mengukur produktivitas, pendapatan, pekerjaan, nilai tambah, kualitas
jaringan, arus tenaga kerja, volume transaksi, atau fungsi metropolitan
yang benar-benar dikonsumsi town. Variabel titik jaringan lebih dekat
dengan hubungan organisasi perusahaan daripada ukuran umum kesejahteraan
ekonomi. (Lokator dasar: hlm. 2372-2374, Bagian 4.1-4.2; TXT baris
301-304, 332-369.)

\subsection{Results}\label{results-22}

\textbf{Laporan sumber: perbedaan kinerja menurut ukuran penduduk.} Pada
66 unit town, semakin kecil populasi yang telah ada, semakin rendah
nilai regional performance. Town dengan populasi kurang dari 100.000
disebut memiliki karakter agglomeration shadow, sedangkan town dengan
populasi lebih dari 100.000 disebut memperlihatkan borrowed-size
performance. Tabel 2 menunjukkan bahwa kelompok kurang dari 50.000
memiliki kinerja jumlah perusahaan baru -0,429 dan kinerja titik
jaringan -0,676; kelompok 50.000-100.000 masing-masing -0,327 dan
-0,302; kelompok 100.000-200.000 masing-masing 0,471 dan 0,236; kelompok
200.000-500.000 masing-masing 0,060 dan 0,301; dan kelompok lebih dari
500.000 masing-masing 0,398 dan 1,572. (Lokator: hlm. 2375-2376, Bagian
5.1 dan Tabel 2; TXT baris 373-394, 413-427.)

\textbf{Interpretasi penulis.} Penulis menjelaskan nilai yang relatif
rendah pada kelompok 200.000-500.000 melalui contoh Lecong dan
Nanzhuang, yaitu town dengan produksi khusus seperti furnitur dan bahan
bangunan. Struktur industri yang tunggal, kebutuhan transformasi,
perpindahan perusahaan, kemunduran atau penutupan perusahaan lama, serta
lemahnya kesinambungan perusahaan baru dipandang menurunkan hubungan
fungsional perusahaan dan menciptakan masalah pengosongan industri.
(Lokator: hlm. 2375, Bagian 5.1; TXT baris 388-394.)

\textbf{Laporan sumber: pola spasial jumlah perusahaan baru.} Jumlah
perusahaan baru pada 2017-2019 memperlihatkan pola pemusatan. Nansha
disebut mempunyai jumlah perusahaan baru yang sangat menonjol
dibandingkan populasi awalnya, sedangkan Lecong dan Huangpu juga
merupakan unit dengan residual terstandardisasi besar. Nancun, Shiling,
dan Taihe di sekitar pusat kota memiliki regional performance yang
menonjol; Genghe dan Liangkou di wilayah lebih jauh menunjukkan
agglomeration shadow. Distribusi keseluruhan digambarkan sebagai pola
bercincin, dan kedekatan antartown yang menunjukkan rekursi spasial
dibaca sebagai indikasi bahwa agglomeration externality lebih kuat
daripada network externality pada sebagian pola tersebut. (Lokator: hlm.
2375-2377, Bagian 5.1.1 dan Gambar 6; TXT baris 418-434.)

\textbf{Laporan sumber: pola spasial titik jaringan perusahaan.} Kinerja
yang dibangun dari titik jaringan perusahaan memperlihatkan pola
bercincin yang lebih jelas. Sejumlah town di cincin suburban dekat pusat
Guangzhou dan Foshan, seperti Huangpu, Huadu, Taihe, Renhe, Panyu,
Nancun, Lecong, Chencun, Shawan, dan Shunde, disebut memiliki
borrowed-size performance yang menonjol. Sebaliknya, town pada cincin
suburban yang lebih jauh, seperti Aotou, Liangkou, Lü tian, Paitan,
Zhengguo, Genghe, Baini, dan Mingcheng, disebut memperlihatkan
agglomeration shadow yang nyata. (Lokator: hlm. 2377, Bagian 5.1.2 dan
Gambar 7; TXT baris 435-442.)

\textbf{Laporan sumber: regresi jumlah perusahaan baru.} Tabel 3
menampilkan lima spesifikasi regresi bertahap. Tanda \texttt{**} berarti
\texttt{p\ \textless{}\ 0,01} dan \texttt{*} berarti
\texttt{p\ \textless{}\ 0,05}. Angka yang ditampilkan sumber adalah
sebagai berikut:

\begin{itemize}
\tightlist
\item
  Model 1 memuat koefisien \texttt{ln\ Pop} 1,113**, \texttt{ln\ Hub}
  0,649\emph{, dan \texttt{ln\ Mul\_pat} 0,083}, dengan \texttt{R2}
  0,886 dan adjusted \texttt{R2} 0,881.
\item
  Model 2 memuat \texttt{ln\ Pop} 1,089**, \texttt{ln\ Bor\_Size\_ij}
  0,429\emph{, dan \texttt{ln\ Mul\_pat} 0,035}, dengan \texttt{R2}
  0,885 dan adjusted \texttt{R2} 0,879.
\item
  Model 3 memuat \texttt{ln\ Pop} 1,102**, \texttt{ln\ Bor\_Per\_ij}
  berbasis jumlah perusahaan terdaftar 0,372\emph{, dan
  \texttt{ln\ Mul\_pat} 0,127}, dengan \texttt{R2} 0,884 dan adjusted
  \texttt{R2} 0,879.
\item
  Model 4 memuat \texttt{ln\ Pop} 1,167** dan \texttt{ln\ Trans} 0,227*,
  dengan \texttt{R2} 0,876 dan adjusted \texttt{R2} 0,872.
\item
  Model 5 memuat \texttt{ln\ Trans} 0,246\textbf{, \texttt{ln\ Mul\_pat}
  0,186}, dan \texttt{ln\ Cross\_pat} 0,246*, dengan \texttt{R2} 0,641
  dan adjusted \texttt{R2} 0,624. (Lokator: hlm. 2378, Tabel 3; TXT
  baris 464-479.)
\end{itemize}

\textbf{Laporan sumber.} Narasi artikel menyatakan bahwa jumlah
perusahaan terdaftar berhubungan dengan ukuran penduduk yang sudah ada,
aksesibilitas hub regional, borrowed size, borrowed performance,
aksesibilitas jaringan transportasi lokal, dan paten kerja sama
lintas-town. Penulis menyatakan koefisien elastisitasnya menurun menurut
urutan ukuran awal town, hub transportasi regional, borrowed size,
borrowed performance, dan paten kerja sama lintas-town. (Lokator: hlm.
2378, Bagian 5.2.1; TXT baris 482-493.)

\textbf{Interpretasi penulis.} Ukuran penduduk awal dipandang sebagai
fondasi pertumbuhan karena konsentrasi perusahaan menciptakan lapangan
kerja dan menarik penduduk. Borrowed size dibaca sebagai kecenderungan
perusahaan memilih town yang dekat dengan populasi padat dan memiliki
aksesibilitas baik. Paten kerja sama lintas-town ditafsirkan sebagai
saluran berbagi pengetahuan, pencocokan, dan pembelajaran yang mendorong
pertumbuhan jumlah perusahaan. (Lokator: hlm. 2378, Bagian 5.2.1; TXT
baris 482-493.)

\textbf{Laporan sumber: perbedaan sektor untuk jumlah perusahaan baru.}
Regresi bertahap juga dijalankan untuk manufaktur, perdagangan dan
sirkulasi komersial, persewaan dan jasa bisnis, serta jasa teknologi dan
perangkat lunak. Narasi sumber menyatakan bahwa manufaktur dijelaskan
oleh ukuran town yang telah ada, aksesibilitas transportasi lokal, dan
paten teknologi umum multi-kategori; borrowed size tidak masuk ke
persamaan manufaktur. Pada tiga sektor lain, borrowed size dan paten
lintas-town disebut sebagai variabel yang signifikan. Tabel 4 melaporkan
\texttt{R2} dan adjusted \texttt{R2} masing-masing 0,796 dan 0,786 untuk
manufaktur, 0,776 dan 0,765 untuk perdagangan, 0,747 dan 0,735 untuk
persewaan atau jasa bisnis, serta 0,721 dan 0,707 untuk teknologi atau
perangkat lunak. (Lokator: hlm. 2378, Bagian 5.2.1 dan Tabel 4; TXT
baris 494-526.)

\textbf{Laporan sumber: regresi titik jaringan perusahaan.} Tabel 5
memperlihatkan lima spesifikasi untuk pertumbuhan titik jaringan
perusahaan. \texttt{ln\ Pop} bernilai positif dan signifikan pada empat
spesifikasi pertama, yaitu 1,151\textbf{, 1,144}, 1,274\textbf{, dan
1,282}. \texttt{ln\ Bor\_Size\_ij} tampil sebesar 0,440* dan 0,516* pada
dua spesifikasi; borrowed performance berbasis titik jaringan tampil
sebesar 0,382* pada Model 1 dan 0,449* pada Model 5; \texttt{ln\ Hub}
tampil sebesar 0,796\textbf{, 0,822}, 0,821\textbf{, dan 0,789* pada
empat spesifikasi; \texttt{ln\ Trans} sebesar 0,596} pada Model 5;
\texttt{ln\ Mul\_pat} sebesar 0,089** pada dua spesifikasi; dan
\texttt{ln\ Cross\_pat} sebesar 0,421** pada Model 5. \texttt{R2}
masing-masing adalah 0,915, 0,915, 0,904, 0,904, dan 0,641, sedangkan
adjusted \texttt{R2} masing-masing 0,909, 0,909, 0,900, 0,900, dan
0,624. (Lokator: hlm. 2379, Tabel 5; TXT baris 529-543.)

\textbf{Laporan sumber.} Penulis merangkum hasil Tabel 5 sebagai
hubungan signifikan antara titik jaringan perusahaan dengan ukuran town
yang telah ada, borrowed size, aksesibilitas hub regional, dan paten
kerja sama lintas-town. Model dengan empat variabel penjelas disebut
memiliki adjusted \texttt{R2} sebesar 0,909. (Lokator: hlm. 2379, Bagian
5.2.2 dan Tabel 5; TXT baris 503-507.)

\textbf{Interpretasi penulis.} Aksesibilitas hub regional, terutama
lingkungan Bandara Baiyun dan Stasiun Guangzhou Selatan, dipandang
penting bagi pembentukan jaringan kantor pusat-cabang. Infrastruktur
transportasi dan fasilitas pertemuan di sekitar hub dianggap mempermudah
penyebaran pengetahuan tacit; hubungan sosial perusahaan tersebut lalu
dipahami sebagai hasil kopling antara bentuk jaringan kota,
aksesibilitas, dan aliran pengetahuan teknologi. (Lokator: hlm. 2379,
Bagian 5.2.2; TXT baris 508-551.)

\textbf{Laporan sumber: perbedaan sektor untuk titik jaringan.} Untuk
empat sektor yang sama, sebagian besar model titik jaringan dapat
memasukkan ukuran town yang telah ada, aksesibilitas jaringan hub
regional, dan paten kerja sama lintas-town. Penulis menyatakan bahwa
elastisitas borrowed size tertinggi muncul pada titik jaringan
perusahaan manufaktur; perdagangan dan jasa teknologi atau perangkat
lunak masing-masing menambahkan borrowed performance dan paten teknologi
umum multi-kategori sebagai variabel penjelas. Tabel 6 melaporkan
\texttt{R2} dan adjusted \texttt{R2} sebesar 0,800 dan 0,787 untuk
manufaktur, 0,881 dan 0,873 untuk perdagangan, 0,862 dan 0,855 untuk
persewaan atau jasa bisnis, serta 0,740 dan 0,731 untuk teknologi atau
perangkat lunak. (Lokator: hlm. 2380, Bagian 5.2.2 dan Tabel 6; TXT
baris 552-579.)

\textbf{Inferensi review.} Hasil regresi menunjukkan kekuatan asosiasi
yang relatif tinggi dalam spesifikasi yang ditampilkan, tetapi nilai
\texttt{R2} atau adjusted \texttt{R2} tidak mengisolasi kontribusi
kausal setiap jaringan. Perbandingan model juga perlu dibaca sebagai
perbandingan spesifikasi bertahap, bukan bukti bahwa seluruh variabel
selalu signifikan dalam satu persamaan final. Selain itu, pemetaan nama
paten yang terbalik pada Tabel 1 membuat interpretasi
\texttt{ln\ Mul\_pat} dan \texttt{ln\ Cross\_pat} perlu mengikuti narasi
masing-masing tabel dengan kehati-hatian. (Lokator dasar: hlm. 2373-2376
dan 2378-2380, Bagian 4.1-5.2 serta Tabel 1, 3, 5, dan 6; TXT baris
305-320, 397-409, 464-579.)

\subsection{Findings}\label{findings-22}

\textbf{Temuan yang dilaporkan sumber.} Ukuran penduduk awal berkorelasi
positif dengan kinerja perusahaan relatif. Batas 100.000 penduduk
dipakai penulis untuk membedakan kelompok yang cenderung menunjukkan
borrowed-size performance dari kelompok yang cenderung menunjukkan
agglomeration shadow. Namun, kelompok 200.000-500.000 tidak mengikuti
peningkatan sederhana pada seluruh indikator karena memiliki nilai
jumlah perusahaan baru yang hanya 0,060, jauh di bawah kelompok
100.000-200.000 yang bernilai 0,471. (Lokator: hlm. 2375-2376, Bagian
5.1 dan Tabel 2; TXT baris 373-394, 413-427.)

\textbf{Temuan yang dilaporkan sumber.} Kinerja jumlah perusahaan baru
dan kinerja titik jaringan sama-sama membentuk pola ruang bercincin,
tetapi pola titik jaringan disebut lebih kuat. Town pada suburban dekat
pusat Guangzhou-Foshan mempunyai kinerja relatif yang lebih menonjol,
sedangkan town pada suburban yang lebih jauh lebih sering menjadi lokasi
agglomeration shadow. (Lokator: hlm. 2375-2377, Bagian 5.1.1-5.1.2 dan
Gambar 6-7; TXT baris 418-442.)

\textbf{Temuan yang dilaporkan sumber.} Dalam regresi jumlah perusahaan
baru, ukuran town yang telah ada merupakan penjelas yang paling
konsisten. Aksesibilitas hub regional, borrowed size, borrowed
performance, aksesibilitas transportasi lokal, serta indikator paten
juga tampil dalam spesifikasi yang berbeda. Dalam regresi titik
jaringan, ukuran town, borrowed size, aksesibilitas hub regional,
aksesibilitas transportasi, borrowed performance, dan paten tampil dalam
kombinasi model bertahap; model dengan empat variabel penjelas mempunyai
adjusted \texttt{R2} 0,909 menurut narasi artikel. (Lokator: hlm.
2378-2379, Bagian 5.2 dan Tabel 3-5; TXT baris 443-507, 529-543.)

\textbf{Temuan yang dilaporkan sumber.} Hasil menurut sektor tidak
seragam. Untuk jumlah perusahaan baru, manufaktur dikaitkan dengan
ukuran town, aksesibilitas transportasi lokal, dan teknologi umum
multi-kategori, tanpa borrowed size masuk ke persamaan; tiga sektor jasa
atau perdagangan lainnya disebut memasukkan borrowed size dan paten
lintas-town. Untuk titik jaringan, ukuran town relatif stabil sebagai
penjelas, borrowed-size elasticity disebut paling tinggi pada
manufaktur, dan beberapa sektor menambahkan borrowed performance atau
teknologi umum multi-kategori. (Lokator: hlm. 2378-2380, Bagian
5.2.1-5.2.2 dan Tabel 4 serta 6; TXT baris 494-500, 552-579.)

\textbf{Interpretasi penulis.} Penulis menyimpulkan bahwa pertumbuhan
perusahaan di metropolitan Guangzhou-Foshan berasal dari kopling antara
akumulasi ukuran town, kedekatan dan aksesibilitas jaringan
transportasi, borrowed size, serta aliran pengetahuan lintas-town.
Mereka menilai bahwa paten kerja sama lintas-town muncul lebih sering
sebagai penjelas daripada teknologi umum multi-kategori, sehingga
eksternalitas jaringan inovasi antardaerah dianggap lebih kuat daripada
inovasi diversifikasi internal. (Lokator: hlm. 2380-2381, Bagian 6; TXT
baris 584-615.)

\textbf{Interpretasi penulis.} Dalam pembacaan teori yang dinyatakan
tidak ketat, penulis mengurutkan ciri regional Guangzhou-Foshan sebagai
\texttt{agglomeration\ externality\ \textgreater{}\ network\ externality\ \textgreater{}\ diversity\ externality}.
Mereka menghubungkannya dengan lintasan historis town spesialisasi:
konsentrasi spesialisasi pada tahap awal, ekonomi jalan nasional dan
provinsi, lalu kerja sama inovasi serta konektivitas ke pusat
metropolitan pada tahap berikutnya. (Lokator: hlm. 2380-2381, Bagian 6;
TXT baris 608-623.)

\textbf{Inferensi review.} Bukti paling langsung dalam artikel mendukung
adanya gradien ukuran, perbedaan spasial cincin, serta asosiasi antara
indikator perusahaan dan variabel jaringan. Bukti tersebut belum
memisahkan secara empiris apakah kinerja di town dekat pusat merupakan
fungsi yang dipinjam, aglomerasi lokal, kebijakan khusus, atau
keunggulan lokasi. Kesimpulan tentang urutan tiga jenis eksternalitas
karena itu merupakan interpretasi penulis atas pola koefisien dan
frekuensi masuknya variabel, bukan dekomposisi eksperimental atas
besaran masing-masing eksternalitas. (Lokator dasar: hlm. 2375-2381,
Bagian 5-6; TXT baris 373-434, 443-623.)

\textbf{Inferensi review.} Ambang 100.000 penduduk berasal dari
pengelompokan dan perbandingan nilai kinerja, bukan dari estimasi titik
patah atau uji nonlinier yang dijelaskan dalam TXT. Demikian pula, pola
bercincin dari peta tidak dengan sendirinya membuktikan bahwa kedekatan
menghasilkan borrowed size atau bahwa jarak menghasilkan shadow; peta
tersebut menunjukkan konfigurasi spasial yang kemudian ditafsirkan
melalui teori. (Lokator dasar: hlm. 2372-2373, Bagian 4.1; hlm.
2375-2377, Bagian 5.1 dan Gambar 6-7; TXT baris 279-304, 373-442.)

\subsection{Conclusions}\label{conclusions-22}

\textbf{Kesimpulan yang dinyatakan penulis.} Artikel menyimpulkan bahwa
eksternalitas jaringan yang berkopel dengan banyak faktor dapat
menghasilkan regional performance maupun agglomeration shadow dalam
perkembangan town. Ukuran penduduk yang telah ada harus mencapai tingkat
tertentu agar efek skala terbentuk; town dengan populasi di bawah
100.000 cenderung memiliki kinerja relatif rendah, sedangkan town di
atas 100.000 menunjukkan kinerja borrowed-size dalam pembagian yang
dibuat penulis. (Lokator: hlm. 2380-2381, Bagian 6; TXT baris 582-591.)

\textbf{Kesimpulan yang dinyatakan penulis.} Regional performance jumlah
perusahaan baru dan titik jaringan perusahaan mempunyai pola spasial
bercincin. Town di suburban dekat pusat kota memperlihatkan kinerja yang
lebih jelas, sedangkan town di suburban jauh memperlihatkan
agglomeration shadow yang lebih signifikan. (Lokator: hlm. 2380-2381,
Bagian 6; TXT baris 584-591.)

\textbf{Kesimpulan yang dinyatakan penulis.} Titik jaringan perusahaan
berhubungan dengan ukuran town yang telah ada, borrowed size,
aksesibilitas jaringan hub regional, dan paten kerja sama lintas-town.
Secara lebih luas, penulis menyatakan bahwa peningkatan jaringan
transportasi regional dan jaringan kerja sama teknologi membantu
pemusatan faktor serta memengaruhi pembangunan urbanisasi metropolitan
melalui eksternalitas jaringan. (Lokator: hlm. 2381, Bagian 6; TXT baris
592-600.)

\textbf{Kesimpulan yang dinyatakan penulis.} Penulis menyatakan bahwa
eksternalitas aglomerasi, jaringan, dan diversifikasi perlu dibaca
secara terpadu. Mereka menilai teori Marshall dan Jacobs tidak dapat
secara mandiri menjelaskan pertumbuhan town spesialisasi di
Guangzhou-Foshan, sedangkan pendekatan gabungan Capello-Camagni lebih
sesuai dengan hasil yang mereka baca. (Lokator: hlm. 2381, Bagian 6; TXT
baris 608-615.)

\textbf{Interpretasi penulis.} Diskusi akhir menggambarkan lintasan town
spesialisasi dari ketergantungan pada pusat melalui istilah
``星期六工程师'' pada 1980-an menuju kerja sama inovasi dengan pusat
metropolitan pada periode baru. Beijiao di Shunde dipakai sebagai contoh
town yang, melalui akses ke Guangzhou South Railway Station dan hubungan
dengan perusahaan besar, berpotensi menjadi bagian penting dari sistem
polisentris metropolitan. Penulis juga memperkirakan kawasan inovasi
seperti Songshan Lake European Town dan Beijiao dapat menjadi kelompok
fungsi penting melalui borrowed size. (Lokator: hlm. 2381, Bagian 6; TXT
baris 615-627.)

\textbf{Inferensi review.} Kesimpulan empiris sebaiknya dibatasi pada
town di metropolitan Guangzhou-Foshan, indikator perusahaan yang
tersedia, dan spesifikasi regresi yang ditampilkan. Kata
``mempengaruhi'' atau ``mendorong'' dalam kesimpulan penulis tidak dapat
dibaca sebagai bukti kausal yang berdiri sendiri karena desain, urutan
waktu variabel, dan strategi identifikasi sebab-akibat tidak dijelaskan
secara memadai dalam TXT. (Lokator dasar: hlm. 2373-2381, Bagian 4-6;
TXT baris 332-369, 443-623.)

\subsection{Contributions}\label{contributions-22}

\textbf{Kontribusi yang diklaim atau ditunjukkan sumber.} Artikel
memperluas pembahasan borrowed size dan agglomeration shadow ke skala
town di dalam metropolitan, bukan hanya ke kota atau wilayah
metropolitan yang lebih agregat. Penulis secara eksplisit menempatkan
ketiadaan pemeriksaan empiris pada tingkat town sebagai alasan
penelitian. (Lokator: hlm. 2368 dan 2370, Bagian 1 dan 2.2; TXT baris
99-107, 191-197.)

\textbf{Kontribusi yang diklaim atau ditunjukkan sumber.} Artikel
mengoperasionalkan kinerja ekonomi melalui dua indikator aktivitas
perusahaan: jumlah perusahaan baru terdaftar dan titik jaringan kantor
pusat-cabang lintas-town. Indikator kedua dipakai untuk membawa hubungan
fungsional antarkota ke dalam analisis borrowed size. (Lokator: hlm.
2372-2374, Bagian 4.1-4.2; TXT baris 279-304, 332-336.)

\textbf{Kontribusi yang diklaim atau ditunjukkan sumber.} Kerangka
empiris menggabungkan ukuran awal town, borrowed size berbasis populasi
town lain dan jarak, borrowed performance, aksesibilitas transportasi
lokal, aksesibilitas hub regional, serta jaringan paten. Dengan
demikian, artikel membedakan jaringan fisik yang dibatasi jarak dari
jaringan pengetahuan yang melintasi batas geografis. (Lokator: hlm.
2369-2374, Bagian 2.1 dan 4.1-4.2; TXT baris 122-178, 305-363.)

\textbf{Kontribusi yang diklaim atau ditunjukkan sumber.} Dari sisi
metode, artikel menggunakan residual terstandardisasi untuk
membandingkan performa terhadap ukuran penduduk, GIS untuk membaca pola
spasial, perbandingan AICC untuk memilih antara regresi berbobot
geografis dan regresi berganda, serta regresi bertahap menurut sektor.
(Lokator: hlm. 2372-2374, Bagian 4.1-4.2; hlm. 2375-2380, Bagian 5; TXT
baris 279-369, 371-579.)

\textbf{Kontribusi empiris yang diklaim atau ditunjukkan sumber.}
Artikel melaporkan ambang pembeda 100.000 penduduk, pola cincin yang
lebih kuat pada titik jaringan perusahaan, perbedaan antara suburban
dekat dan jauh, serta heterogenitas faktor menurut sektor. Temuan
tersebut memberi pembacaan mikro terhadap struktur internal metropolitan
Guangzhou-Foshan. (Lokator: hlm. 2375-2381, Bagian 5-6; TXT baris
373-623.)

\textbf{Kontribusi praktis yang diklaim atau ditunjukkan sumber.} Hasil
artikel diarahkan untuk membantu memahami bagaimana pembangunan jaringan
transportasi, kerja sama teknologi, dan organisasi fungsi perusahaan
dapat mendukung perkembangan town dalam strategi metropolitan dan
urbanisasi baru. (Lokator: hlm. 2367, abstrak; hlm. 2371-2372 dan 2381,
Bagian 3 dan 6; TXT baris 99-107, 232-273, 592-600.)

\textbf{Inferensi review.} Kontribusi paling kuat yang dapat dipastikan
dari TXT adalah pemetaan dan pengukuran komparatif aktivitas perusahaan
pada unit town serta penggabungan dua proksi jaringan. Kontribusi
teoritis terhadap pembuktian borrowed size masih eksploratif, sebab
model tidak mengobservasi secara langsung pertukaran fungsi, aliran
pengguna, atau hubungan sebab-akibat antara kedekatan jaringan dan
performa. (Lokator dasar: hlm. 2372-2374, Bagian 4; TXT baris 279-369.)

\subsection{Research Gap}\label{research-gap-22}

\textbf{Kesenjangan yang dinyatakan sumber sebelum penelitian.} Artikel
menyebut beberapa kekosongan penelitian:

\begin{itemize}
\tightlist
\item
  Riset empiris jaringan eksternalitas sering melihat agregat regional
  dan kurang menelaah interaksi spasial internal metropolitan, terutama
  pada tingkat town. (Lokator: hlm. 2368, Bagian 1; TXT baris 99-107.)
\item
  Riset borrowed size lebih sering memeriksa dampak kota besar terhadap
  pertumbuhan wilayah sekitarnya daripada membuktikan bagaimana town
  kecil meminjam skala atau fungsi kota besar. (Lokator: hlm. 2370,
  Bagian 2.2; TXT baris 179-197.)
\item
  Penelitian sebelumnya yang dirangkum lebih banyak menggunakan ukuran
  kota, wilayah kota, pekerjaan, pendapatan per kapita, atau GDP;
  performa relatif town dengan ukuran dan lokasi berbeda belum cukup
  diuji. (Lokator: hlm. 2370, Bagian 2.2; TXT baris 186-197.)
\item
  Penelitian jaringan eksternalitas domestik dinilai belum cukup
  memperhatikan ukuran awal town, lokasi transportasi, dan organisasi
  spasial perusahaan secara bersamaan. (Lokator: hlm. 2370-2371, Bagian
  2.3; TXT baris 198-229.)
\end{itemize}

\textbf{Inferensi review: kesenjangan yang masih tersisa.} Berdasarkan
desain dan hasil yang tersedia, beberapa gap tetap terbuka:

\begin{itemize}
\tightlist
\item
  Belum ada pengukuran langsung tentang fungsi yang dipinjam dari kota
  besar, seperti aliran pengguna, hubungan kantor pusat-cabang yang
  lengkap, pertukaran faktor, atau komplementaritas fungsi; yang
  tersedia terutama adalah jumlah perusahaan dan proksi jaringan.
\item
  Belum ada pemisahan kausal antara aglomerasi lokal, kebijakan,
  aksesibilitas hub, borrowed size, dan agglomeration shadow. Model
  regresi potong lintang dengan spesifikasi bertahap belum menyediakan
  kontrafaktual.
\item
  Belum ada pembuktian dinamika dari waktu ke waktu. Periode perusahaan
  2017-2019 disebut sebagai sumber data, tetapi TXT tidak menampilkan
  analisis perubahan tahunan atau panel.
\item
  Belum ada uji lintas-metropolitan untuk mengetahui apakah pola 100.000
  penduduk dan suburban bercincin berlaku di luar Guangzhou-Foshan.
\item
  Pemetaan konstruk paten, transformasi \texttt{PER}, arah variabel
  waktu tempuh, serta pemilihan 66 unit masih memerlukan dokumentasi
  agar hasil dapat direplikasi.
\end{itemize}

Semua butir pada daftar kedua adalah inferensi review dari isi dan
ketidakjelasan TXT, bukan agenda eksplisit yang telah dinyatakan
penulis. (Lokator dasar: hlm. 2372-2381, Bagian 4-6; TXT baris 279-369,
373-623.)

\subsection{Limitations}\label{limitations-22}

\textbf{Laporan sumber.} Artikel tidak memiliki subbagian keterbatasan
yang berdiri sendiri dan tidak menyajikan daftar keterbatasan formal.
Keterbatasan atau kendala yang disebut secara langsung adalah pemilihan
data perusahaan 2017-2019 karena ketersediaan data dan penghapusan
variabel kepadatan penduduk setelah diagnosis multikolinearitas
menunjukkan VIF lebih dari 5. (Lokator: hlm. 2373-2374, Bagian 4.2; TXT
baris 332-369.)

\textbf{Laporan sumber.} Artikel memakai satu kasus metropolitan
Guangzhou-Foshan, 73 unit spasial dalam konteks kawasan, dan 66 unit
dalam tabel analisis. Data yang dipakai sebagai ukuran awal penduduk
adalah sensus keenam nasional, sedangkan indikator perusahaan berasal
dari 2017-2019. Artikel tidak menawarkan replikasi pada metropolitan
lain atau seri pembanding lain. (Lokator: hlm. 2371-2374, Bagian 3-4.2;
TXT baris 232-273, 332-369.)

\textbf{Inferensi review: validitas konstruk.} Jumlah perusahaan baru
dan jumlah perusahaan yang memiliki kantor pusat serta cabang
lintas-town mengukur aktivitas perusahaan yang tercatat, bukan
keseluruhan pertumbuhan ekonomi atau kesejahteraan town. Indikator
tersebut tidak memuat pendapatan, pekerjaan, produktivitas, nilai
tambah, kapasitas produksi, umur perusahaan, pemanfaatan layanan, atau
volume hubungan antarunit. (Lokator dasar: hlm. 2372-2374, Bagian
4.1-4.2; TXT baris 284-304, 332-369.)

\textbf{Inferensi review: borrowed size.} Rumus borrowed size dan
borrowed performance menjumlahkan ukuran atau performa town lain yang
dibagi jarak. Bentuk tersebut memberi proksi potensi kedekatan dan skala
eksternal, tetapi tidak membuktikan bahwa town i benar-benar menerima
fungsi dari town j. Variabel jarak juga tidak dengan sendirinya
membedakan manfaat koneksi dari kompetisi atau tarikan pusat. (Lokator
dasar: hlm. 2373, Persamaan 4-5; TXT baris 318-330.)

\textbf{Inferensi review: ukuran dan seleksi unit.} File tidak
menjelaskan secara rinci kriteria yang mengubah 73 unit kawasan menjadi
66 observasi tabel, maupun apakah tujuh unit lainnya dikeluarkan karena
status inti kota atau alasan data. Prosedur sampling, potensi seleksi,
dan uji sensitivitas terhadap definisi town \textbf{Tidak disebutkan
dalam file}. (Lokator: hlm. 2373-2376, Bagian 4.2 dan Tabel 1-2; TXT
baris 332-336, 397-427.)

\textbf{Inferensi review: ketidakselarasan waktu.} Sumber menyebut data
perusahaan 2017-2019 dan sensus penduduk keenam, tetapi tidak menyebut
tahun observasi sensus itu atau bagaimana perbedaan waktu tersebut
ditangani. Karena itu, rasio perusahaan terhadap populasi awal mungkin
mencampur periode pengukuran yang berbeda; besarnya persoalan tersebut
\textbf{Tidak disebutkan dalam file}. (Lokator: hlm. 2373, Bagian 4.2;
TXT baris 332-336.)

\textbf{Inferensi review: transformasi variabel.} Bagian metode
menyatakan bahwa \texttt{E\_i} dibentuk dari rasio indikator terhadap
populasi lalu diubah menjadi residual terstandardisasi \texttt{PER\_i},
sementara persamaan regresi dan Tabel 1 memakai label
\texttt{ln\ PER\_i}. TXT tidak menjelaskan urutan logaritma dan
pembentukan residual, atau bagaimana nilai nonpositif ditangani. Ini
membatasi interpretasi elastisitas dan replikasi angka. (Lokator: hlm.
2372-2373, Persamaan 1-3; hlm. 2376, Tabel 1; TXT baris 279-330,
397-409.)

\textbf{Inferensi review: arah aksesibilitas.} \texttt{ln\ Trans}
dijelaskan sebagai aksesibilitas jaringan transportasi lokal yang
dihitung melalui waktu perjalanan rata-rata tertimbang. \texttt{ln\ Hub}
juga dijelaskan melalui waktu tempuh menuju hub, tetapi koefisien yang
dilaporkan bernilai positif dan narasi penulis menafsirkan aksesibilitas
sebagai faktor yang menguntungkan. TXT tidak menyatakan apakah waktu
dibalik, dinormalisasi, atau dikodekan ulang, sehingga arti substantif
tanda koefisien tidak dapat dipastikan hanya dari file. (Lokator: hlm.
2373-2374, Bagian 4.1-4.2; hlm. 2378-2379, Tabel 3 dan 5; TXT baris
318-345, 464-479, 529-543.)

\textbf{Inferensi review: label paten.} Uraian model memetakan
\texttt{ln\ Mul\_pat} pada paten teknologi umum multi-kategori dan
\texttt{ln\ Cross\_pat} pada paten kerja sama lintas-town, sedangkan
Tabel 1 menampilkan uraian operasional yang tampak terbalik. Tanpa
klarifikasi, interpretasi koefisien kedua variabel itu tidak sepenuhnya
aman. (Lokator: hlm. 2373, Bagian 4.1; hlm. 2376, Tabel 1; TXT baris
308-320, 397-409.)

\textbf{Inferensi review: model statistik.} Artikel membandingkan AICC
dan menyebut VIF, tetapi TXT tidak menyajikan statistik diagnostik
lengkap, standard error, interval kepercayaan, bentuk pembobotan,
pemeriksaan residual, uji autokorelasi spasial, atau uji sensitivitas
terhadap ambang populasi, jarak, dan formula kinerja. Analisis spasial
yang menghasilkan pola cincin tidak disertai model eksplisit untuk
ketergantungan spasial dalam residual. (Lokator: hlm. 2373-2374, Bagian
4.1-4.2; hlm. 2378-2380, Tabel 3-6; TXT baris 305-369, 464-579.)

\textbf{Inferensi review: generalisasi.} Karena kasus hanya satu
metropolitan yang terhubung kuat dan tidak ada pembanding eksternal,
temuan tentang threshold, pola cincin, atau hierarki eksternalitas tidak
dapat digeneralisasi menjadi hukum umum kota kecil tanpa pengujian
tambahan. (Lokator dasar: hlm. 2371-2372, Bagian 3; TXT baris 232-273.)

\subsection{Challenges}\label{challenges-22}

\textbf{Tantangan yang dilaporkan atau tampak dari rancangan sumber.}
Penelitian harus menerjemahkan konsep borrowed size dan eksternalitas
jaringan yang abstrak ke unit town yang sangat kecil dibandingkan skala
kota metropolitan. Artikel membedakan jaringan transportasi yang masih
dibatasi jarak dan biaya perjalanan dari jaringan informasi atau paten
yang dapat melintasi batas geografis, tetapi juga mengakui bahwa kedua
eksternalitas sulit dipisahkan sepenuhnya dalam skala metropolitan.
(Lokator: hlm. 2369-2370, Bagian 2.1.1-2.1.2; TXT baris 144-178.)

\textbf{Tantangan yang dilaporkan atau tampak dari rancangan sumber.}
Data pada tingkat town dikumpulkan dari beberapa sumber dan bentuk:
pendaftaran perusahaan, hubungan kantor pusat-cabang, sensus penduduk,
waktu tempuh, serta paten penemuan. Penulis menyatakan bahwa pemilihan
data perusahaan 2017-2019 berkaitan dengan ketersediaan data.
Perhitungan aksesibilitas juga memerlukan pembobotan penduduk, jarak
atau waktu perjalanan, dan tujuan hub tertentu. (Lokator: hlm.
2373-2374, Bagian 4.1-4.2; TXT baris 332-369.)

\textbf{Tantangan yang dilaporkan atau tampak dari rancangan sumber.}
Pengukuran regional performance harus mengubah jumlah perusahaan menjadi
rasio terhadap populasi lalu membandingkannya dengan rerata wilayah.
Penulis kemudian harus menentukan kelompok ukuran populasi, membaca peta
GIS, dan menyusun regresi bertahap yang berbeda untuk jumlah perusahaan,
titik jaringan, dan empat sektor. Masalah multikolinearitas membuat
kepadatan penduduk dikeluarkan dari kandidat model. (Lokator: hlm.
2372-2374, Bagian 4.1-4.2; TXT baris 279-330, 364-369.)

\textbf{Tantangan substantif yang dibahas penulis.} Town spesialisasi
berukuran 200.000-500.000 menghadapi persoalan transformasi: perpindahan
atau kematian perusahaan lama, berkurangnya ruang untuk manufaktur baru,
dan belum tercapainya tingkat kota yang diperlukan bagi pertumbuhan
jasa. Heterogenitas ini membuat ukuran penduduk saja tidak cukup untuk
menjelaskan jalur pertumbuhan setiap town. (Lokator: hlm. 2375 dan 2381,
Bagian 5.1 dan 6; TXT baris 388-394, 601-607.)

\textbf{Inferensi review.} Tantangan analitis terbesar adalah membedakan
ukuran internal, status administratif, aglomerasi lokal, borrowed size,
akses ke hub, kebijakan khusus, dan kompetisi pusat, karena semuanya
dapat menghasilkan jumlah perusahaan yang tinggi atau rendah. Tantangan
replikasi bertambah karena prosedur pembersihan data, pembentukan
jaringan, transformasi, dan keputusan spasial tidak dijelaskan secara
lengkap. (Lokator dasar: hlm. 2371-2374 dan 2375-2380, Bagian 3-5; TXT
baris 257-273, 332-369, 373-579.)

\textbf{Inferensi review.} Ketidakkonsistenan nama \texttt{ln\ Mul\_pat}
dan \texttt{ln\ Cross\_pat}, serta ketidakjelasan arah ukuran berbasis
waktu tempuh, merupakan tantangan khusus dalam membaca hasil. Review
tidak menyelesaikan ambiguitas tersebut dengan dugaan; statusnya adalah
\textbf{Tidak disebutkan dalam file}. (Lokator: hlm. 2373-2374 dan 2376,
Bagian 4.1-4.2 serta Tabel 1; TXT baris 308-320, 342-345, 397-409.)

\subsection{Practical Implications}\label{practical-implications-22}

\textbf{Interpretasi penulis.} Hasil penelitian mendukung perhatian pada
jaringan transportasi regional dan kerja sama teknologi sebagai bagian
dari strategi urbanisasi metropolitan. Peningkatan konektivitas
dipandang dapat mengurangi biaya hubungan, memperkuat pertukaran faktor,
dan memperluas kesempatan town untuk menerima fungsi atau aktivitas dari
pusat. (Lokator: hlm. 2367, abstrak; hlm. 2370-2371 dan 2381, Bagian 2.3
dan 6; TXT baris 42-54, 198-229, 592-600.)

\textbf{Inferensi review.} Perencanaan tidak seharusnya memperlakukan
semua town sebagai unit yang sama. Town di suburban dekat pusat
menunjukkan potensi kinerja borrowed-size, sedangkan town di suburban
jauh lebih rentan terhadap agglomeration shadow. Town kecil di bawah
100.000 penduduk dan town spesialisasi berukuran menengah memerlukan
pembacaan yang berbeda karena basis populasi dan struktur industrinya
tidak sama. (Lokator: hlm. 2375-2377 dan 2380-2381, Bagian 5.1 dan 6;
TXT baris 373-442, 584-607.)

\textbf{Laporan sumber.} Artikel menempatkan analisis kinerja regional
dan hubungan fungsional perusahaan sebagai bahan untuk memahami
pembagian fungsi metropolitan, menerima limpahan industri, dan membangun
sistem polisentris. Contoh yang dibahas adalah town periferal seperti
Huangpu, Nansha, Lecong, dan Beijiao yang tumbuh melalui kombinasi
posisi, sumber daya, kerja sama industri, dan jaringan infrastruktur.
(Lokator: hlm. 2371-2372 dan 2381, Bagian 3 dan 6; TXT baris 246-273,
615-627.)

\textbf{Inferensi review.} Dalam penggunaan perencanaan, indeks
\texttt{PER} dan peta cincin paling aman dijadikan alat diagnosis awal
untuk menyaring town yang performanya menyimpang dari ukuran penduduk,
bukan sebagai dasar tunggal untuk menetapkan bahwa suatu lokasi pasti
menerima borrowed size atau mengalami shadow. Verifikasi tambahan atas
struktur industri, kapasitas lahan, perusahaan yang keluar atau masuk,
perjalanan aktual, dan hubungan antartown diperlukan sebelum keputusan
investasi atau redistribusi fungsi dibuat. (Lokator dasar: hlm.
2372-2374, Bagian 4; hlm. 2375-2381, Bagian 5-6; TXT baris 279-369,
373-623.)

\textbf{Inferensi review.} Temuan positif tentang jaringan tidak boleh
diterjemahkan secara otomatis menjadi kebijakan mempercepat semua
hubungan. Artikel sendiri membedakan manfaat limpahan dari risiko hisap
atau bayangan aglomerasi. Karena arah substantif variabel waktu tempuh
tidak dijelaskan sepenuhnya, implikasi kebijakan tentang memperpendek
waktu perjalanan perlu diuji ulang dengan definisi aksesibilitas yang
eksplisit. (Lokator dasar: hlm. 2369, Bagian 2.1.1; hlm. 2373-2374 dan
2378-2379, Bagian 4.1-4.2 dan 5.2; TXT baris 131-155, 318-345, 464-479,
529-543.)

\textbf{Batas implikasi praktis.} Instrumen kebijakan spesifik, biaya
intervensi, mekanisme implementasi, indikator pemantauan, dan evaluasi
dampak kebijakan \textbf{Tidak disebutkan dalam file}. Artikel memberi
arah interpretasi untuk metropolitan planning, bukan evaluasi kebijakan
yang sudah diterapkan.

\subsection{Applications}\label{applications-22}

\textbf{Aplikasi yang dilakukan dalam artikel.} Kerangka penelitian
diterapkan untuk:

\begin{itemize}
\tightlist
\item
  memetakan regional performance jumlah perusahaan baru di 66 unit town;
\item
  memetakan regional performance titik jaringan perusahaan kantor
  pusat-cabang lintas-town;
\item
  membandingkan nilai kinerja menurut lima kelompok ukuran populasi;
\item
  membaca perbedaan spasial antara suburban dekat dan suburban jauh
  melalui GIS;
\item
  menguji hubungan kinerja dengan ukuran town, borrowed size, borrowed
  performance, aksesibilitas lokal, aksesibilitas hub, dan indikator
  paten; serta
\item
  membandingkan faktor penjelas pada manufaktur, perdagangan dan
  sirkulasi komersial, persewaan dan jasa bisnis, serta teknologi dan
  perangkat lunak. (Lokator: hlm. 2372-2380, Bagian 4-5 dan Tabel 1-6;
  TXT baris 279-579.)
\end{itemize}

\textbf{Aplikasi yang didukung langsung oleh hasil sumber.} Jumlah
perusahaan baru dapat dipakai sebagai indikator awal kecenderungan
pertumbuhan perusahaan town, sedangkan titik jaringan dapat dipakai
untuk membaca hubungan fungsional kantor pusat-cabang lintas-town. Kedua
indikator dapat dipakai bersama agar analisis tidak hanya menangkap
banyaknya perusahaan, tetapi juga keterhubungan organisasi perusahaan.
(Lokator: hlm. 2372-2374, Bagian 4.1-4.2; TXT baris 301-304, 332-336.)

\textbf{Inferensi review.} Prosedur rasio-populasi, residual
terstandardisasi, pemetaan GIS, dan regresi dapat diadaptasi sebagai
diagnosis eksploratif untuk metropolitan lain. Adaptasi tersebut
memerlukan definisi unit yang sebanding, dokumentasi data perusahaan dan
jaringan, validasi indikator paten, kejelasan transformasi log, serta
pengujian sensitivitas terhadap ambang 100.000 dan formula jarak. Ini
adalah kemungkinan penggunaan yang diturunkan dari metode, bukan
aplikasi lintas wilayah yang dilaporkan artikel.

\textbf{Batas aplikasi.} TXT tidak menyebut perangkat lunak atau kode
analisis selain penyebutan ArcGIS dalam penempatan spasial, tidak
menyediakan evaluasi penerapan di lapangan, dan tidak melaporkan
perubahan kebijakan setelah hasil penelitian. Informasi tentang aplikasi
operasional pascapenelitian \textbf{Tidak disebutkan dalam file}.
(Lokator: hlm. 2374, Bagian 4.2; TXT baris 364-369.)

\subsection{Future Research
(Optional)}\label{future-research-optional-22}

\textbf{Laporan sumber.} Artikel tidak memiliki bagian khusus yang
diberi label future research atau agenda penelitian. Pernyataan agenda
terstruktur \textbf{Tidak disebutkan dalam file}. Meski demikian,
penulis menyebut bahwa analisis eksternalitas jaringan, borrowed size,
dan agglomeration shadow masih merupakan isu yang perlu diperdalam,
terutama untuk menghubungkan pola positif atau negatif dengan praktik
kebijakan metropolitan. (Lokator: hlm. 2369, Bagian 2.1; TXT baris
122-130.)

\textbf{Laporan sumber.} Pada bagian akhir, penulis mendorong analisis
yang lebih berorientasi ke depan mengenai pola dan kecepatan pertumbuhan
ekonomi. Mereka memperkirakan bahwa kawasan inovasi seperti Songshan
Lake European Town dan Beijiao di Shunde dapat menjadi kelompok fungsi
penting dalam sistem polisentris metropolitan melalui keunggulan
borrowed size. Pernyataan ini merupakan pembacaan prospektif penulis,
bukan hasil pengujian masa depan. (Lokator: hlm. 2381, Bagian 6; TXT
baris 624-627.)

\textbf{Inferensi review.} Berdasarkan celah dan keterbatasan yang
terlihat dalam TXT, agenda lanjutan yang relevan adalah:

\begin{itemize}
\tightlist
\item
  membangun data panel agar perubahan perusahaan, jaringan, penduduk,
  dan aksesibilitas dapat diamati dari waktu ke waktu;
\item
  mengukur arus atau hubungan kantor pusat-cabang secara langsung, bukan
  hanya jumlah perusahaan yang memenuhi kriteria lintas-town;
\item
  menguji fungsi borrowed size dengan data pertukaran faktor, pengguna,
  pasar, atau layanan yang dapat membedakannya dari aglomerasi lokal;
\item
  memasukkan beberapa pusat metropolitan dan membandingkan metropolitan
  dengan struktur ekonomi serta jaringan yang berbeda;
\item
  menguji ketergantungan spasial, seleksi lokasi, endogenitas, dan
  alternatif kontrafaktual;
\item
  melakukan uji sensitivitas atas ambang populasi, definisi jarak,
  ukuran aksesibilitas, urutan transformasi log, dan pemetaan
  \texttt{ln\ Mul\_pat} atau \texttt{ln\ Cross\_pat}; serta
\item
  mendokumentasikan proses akuisisi, pembersihan, deduplikasi, dan
  pembentukan jaringan agar analisis dapat direplikasi.
\end{itemize}

Daftar ini adalah inferensi review berbasis desain artikel, bukan
rekomendasi future research yang dinyatakan eksplisit oleh penulis.
(Lokator dasar: hlm. 2372-2381, Bagian 4-6; TXT baris 279-623.)

\subsection{Provenance Notes}\label{provenance-notes-22}

\begin{itemize}
\tightlist
\item
  \textbf{Basis sumber:} Review ini hanya menggunakan file TXT
  \texttt{source/journal/A07-zhao-wang-hu-guo-wei-2022-examining-performance-of-urban-borrowed-size-based-on-the-towns-network-externalities-of-guangzhou-foshan-metropolitan-areas-10.11821-dlyj020211036.txt}.
  PDF asli, sumber daring, dan karya yang tercantum dalam daftar pustaka
  tidak dibuka untuk menambah substansi.
\item
  \textbf{Cakupan baca:} seluruh 791 baris TXT dibaca, termasuk blok
  \texttt{{[}PDF\ Metadata{]}}, teks artikel berbahasa Mandarin pada
  hlm. 2367-2384, enam tabel, tujuh keterangan gambar yang tersedia
  dalam ekstraksi, persamaan, daftar pustaka, serta abstrak dan kata
  kunci bahasa Inggris pada bagian akhir.
\item
  \textbf{Metadata ekstraksi:} TXT menyatakan bahwa teks diekstrak pada
  31 Agustus 2026 dengan \texttt{pdftotext\ -layout} dari PDF 18
  halaman. Metadata mesin menyebut judul internal
  \texttt{4-\/-\/-2021-1036\ 赵渺希}, \texttt{Founder\ magtech2011}
  sebagai author, serta PDF tidak terenkripsi dan tidak bertanda
  struktur. Nilai metadata produksi tersebut tidak dipakai untuk
  menggantikan identitas penulis pada blok judul artikel. (Lokator: TXT
  baris 1-29.)
\item
  \textbf{Dasar identitas:} nama penulis, judul, jurnal, volume, nomor,
  halaman, DOI, afiliasi, tanggal naskah, pendanaan, dan kata kunci
  diambil dari blok judul, catatan penulis, kepala halaman, DOI, dan
  abstrak Inggris yang tersedia di TXT. (Lokator: TXT baris 30-57,
  76-82, 749-790.)
\item
  \textbf{Konvensi locator:} locator utama menggunakan nomor halaman
  jurnal yang tercetak dan bagian, tabel, gambar, atau persamaan
  artikel. Rentang baris TXT ditambahkan agar klaim dapat ditelusuri
  pada file yang direview.
\item
  \textbf{Keterbatasan ekstraksi:} format dua kolom dan pemisahan
  halaman membuat beberapa rumus serta tabel tidak tersusun linear.
  Karena itu, pemetaan label paten pada Tabel 1 versus uraian model,
  urutan transformasi \texttt{PER}, dan arti tanda ukuran berbasis waktu
  tempuh ditandai sebagai ambiguitas, bukan diselesaikan dengan dugaan.
  (Lokator: TXT baris 284-330, 397-409, 464-479, 529-575.)
\item
  \textbf{Informasi yang tidak tersedia:} prosedur sampling rinci,
  alasan lengkap penggunaan 66 dari 73 unit, tanggal dan protokol unduh
  data, pembersihan atau deduplikasi dataset, penanganan data hilang,
  rincian pembobotan jarak dan waktu, diagnostik residual atau
  autokorelasi spasial, urutan transformasi log, koreksi label paten,
  kode replikasi, evaluasi kebijakan, dan agenda future research formal
  \textbf{Tidak disebutkan dalam file}.
\item
  \textbf{Status peer review:} TXT menyebut ucapan terima kasih kepada
  penelaah anonim, tetapi status telaah sejawat formal \textbf{Tidak
  disebutkan dalam file} dan tidak diverifikasi dari luar sumber.
\item
  \textbf{Pemisahan epistemik:} bagian berlabel \textbf{Laporan sumber}
  memparafrasekan data, metode, hasil, atau pernyataan artikel.
  \textbf{Interpretasi penulis} merekam penjelasan dan atribusi
  mekanisme yang diberikan penulis. \textbf{Inferensi review} adalah
  evaluasi yang diturunkan hanya dari isi, struktur, angka, dan
  ketidakjelasan internal TXT; bukan klaim penulis dan bukan informasi
  eksternal.
\end{itemize}

\section{Review A08 - Zhen, Shi, dan Lu
(2023)}\label{review-a08---zhen-shi-dan-lu-2023}

Status: Review literatur berbasis seluruh teks TXT yang tersedia dari
ekstraksi PDF. Review ini hanya menggunakan sumber A08.

\subsection{Identitas Artikel}\label{identitas-artikel-1}

\begin{itemize}
\tightlist
\item
  \textbf{Kode:} A08.
\item
  \textbf{Judul:} \emph{The Impact of Regional Integration Strategies on
  the Formation of City Regions and Its Agglomeration Shadow: Evidence
  from the Yangtze River Delta, China}.
\item
  \textbf{Penulis:} Yanlin Zhen, Dehao Shi, dan Yanan Lu.
\item
  \textbf{Afiliasi yang tercantum:} Spatial Planning Center, Yangtze
  Delta Region Institute of Tsinghua University, Jiaxing; Department of
  Geography and Resource Management, The Chinese University of Hong
  Kong; serta School of Public Affairs, Zhejiang University, Hangzhou.
\item
  \textbf{Tahun terbit:} 2023.
\item
  \textbf{Jurnal:} \emph{Land}, volume 12, artikel 1053. Nomor edisi
  tidak disebutkan dalam file; PDF yang tersedia terdiri atas 16
  halaman.
\item
  \textbf{DOI:} 10.3390/land12051053.
\item
  \textbf{Editor akademik:} Thomas Panagopoulos.
\item
  \textbf{Riwayat naskah:} diterima 2 April 2023, direvisi 9 Mei 2023,
  diterima 10 Mei 2023, dan diterbitkan 12 Mei 2023.
\item
  \textbf{Jenis sumber:} artikel jurnal empiris. Status penelaahan
  sejawat tidak dapat diverifikasi secara mandiri dari TXT; bagian
  penghargaan menyebut anonymous reviewers dan editors.
\item
  \textbf{Lisensi:} artikel menyatakan akses terbuka dengan lisensi
  Creative Commons Attribution (CC BY).
\item
  \textbf{Pendanaan:} bagian pendanaan menyatakan tidak ada pendanaan
  eksternal.
\item
  \textbf{Pernyataan konflik kepentingan:} penulis menyatakan tidak ada
  konflik kepentingan dan funder tidak berperan dalam desain,
  pengumpulan, analisis, interpretasi, penulisan, atau keputusan
  publikasi.
\item
  \textbf{Ketersediaan data:} data dinyatakan akan tersedia atas
  permintaan yang relevan kepada penulis korespondensi; data atau kode
  tidak disertakan dalam TXT.
\item
  \textbf{Kata kunci:} \emph{city region}, \emph{regional integration
  strategy}, \emph{industrial integration}, \emph{DDD model}, dan
  \emph{Yangtze River Delta}.
\item
  \textbf{Objek, wilayah, dan periode:} pengaruh \emph{regional
  integration strategy} (RIS) terhadap integrasi industri dan
  pembentukan \emph{city region} di Delta Sungai Yangtze (YRD), dengan
  periode analisis 2000-2017 dan kebijakan utama tahun 2010.
\item
  \textbf{Catatan jumlah unit:} abstrak menyebut 122 unit tingkat
  kabupaten, sedangkan uraian studi dan kesimpulan menyebut 121 entitas
  perkotaan. Tabel regresi melaporkan \texttt{N\ =\ 2178}. File tidak
  menjelaskan perbedaan angka tersebut. (Locator: metadata PDF dan
  abstrak, TXT baris 7-10 dan 45-56; Bagian 3.1, hlm. 5, TXT baris
  157-166 dan 431-456; Bagian 6, hlm. 13, TXT baris 953-958; Tabel 1-3,
  hlm. 8-11, TXT baris 624-651, 677-719, dan 812-858.)
\item
  \textbf{Locator identitas:} kepala artikel, abstrak, dan sitasi pada
  hlm. 1; riwayat naskah pada hlm. 1; metadata ekstraksi TXT baris 1-28.
\end{itemize}

\subsection{Summarized Abstract}\label{summarized-abstract-23}

\textbf{Laporan sumber.} Artikel menguji apakah RIS yang dipandu negara
memengaruhi pembentukan \emph{city region} di YRD dan apakah pengaruh
tersebut memiliki heterogenitas industri serta spasial. Abstrak menyebut
sampel 122 unit tingkat kabupaten pada 2000-2017. Dalam analisis, YRDRP
yang disetujui pemerintah pusat pada 2010 diperlakukan sebagai
eksperimen kuasi-alamiah. Model \emph{difference-in-differences} (DID)
digunakan untuk mengestimasi pengaruh RIS terhadap integrasi industri,
sedangkan model \emph{difference-in-difference-in-difference} (DDD)
digunakan untuk menguji heterogenitas spasial. Artikel melaporkan bahwa
RIS berdampak positif dan signifikan terhadap integrasi industri, dengan
pola yang berbeda antara industri sekunder dan tersier. Industri
sekunder disebut memperoleh tata letak yang lebih seimbang, sedangkan
industri tersier cenderung mengelompok menuju kota pusat. Pengaruh
negatif RIS terhadap integrasi kedua sektor disebut lebih kuat di kota
yang berdekatan dengan metropolitan, sehingga RIS dipandang memperbesar
\emph{agglomeration shadow}. Fieldwork mendalam di Jiaxing digunakan
untuk mengurai mekanisme bayangan tersebut. (Locator: abstrak, hlm. 1,
TXT baris 45-56.)

\textbf{Interpretasi penulis.} Penulis memandang RIS sebagai intervensi
negara yang mengurangi hambatan institusional, mempercepat arus lintas
batas, serta mengubah hubungan antarkota dari struktur hierarkis menuju
struktur berjaringan. Dalam kerangka ini, penyebaran industri sekunder
ke wilayah perifer dan konsentrasi industri tersier di metropolitan
dianggap sebagai bentuk integrasi industri dalam pembentukan \emph{city
region}. Namun, kedekatan dengan metropolitan juga dapat menekan kota
yang berada di antaranya. (Locator: pendahuluan, hlm. 1-3, TXT baris
99-151; Bagian 2.1-2.2, hlm. 3-4, TXT baris 168-207 dan 208-360.)

\textbf{Inferensi pengolahan.} Bukti utama artikel adalah perubahan dan
perbedaan indeks lokasi industri, bukan pengukuran langsung atas arus
pekerja, modal, konsumsi, perusahaan, atau interaksi antarkota. Karena
itu, klaim tentang pembentukan \emph{city region}, \emph{borrowed size},
dan \emph{agglomeration shadow} perlu dibaca sebagai interpretasi atas
pola LQ dan koefisien DID/DDD. Ketidaksesuaian jumlah unit antara
abstrak dan bagian metode juga membuat ukuran sampel perlu diperlakukan
sebagai informasi yang belum terselesaikan, bukan dipilih secara
sepihak.

\subsection{Problem Statement}\label{problem-statement-23}

\textbf{Laporan sumber.} Literatur tentang pembangunan perkotaan dan
wilayah telah banyak membahas aglomerasi, eksternalitas ekonomi, dan
pembentukan \emph{city region}. Studi terdahulu menelaah urban sprawl,
peningkatan industri, akumulasi modal, mobilitas penduduk,
restrukturisasi industri, penggunaan lahan, dan tata kelola. Akan
tetapi, penulis menilai sebagian besar literatur mengenai pembentukan
\emph{city region} berfokus pada redistribusi spontan yang dipimpin
pasar dan kurang memperhatikan peran kekuasaan negara. (Locator:
pendahuluan, hlm. 1-2, TXT baris 65-127.)

\textbf{Laporan sumber.} Penulis juga menyatakan bahwa penelitian
tentang eksternalitas \emph{city region} lebih banyak
mengonseptualisasikan dan mengukur eksternalitas ekonomi, tetapi hanya
sedikit yang menilai heterogenitas spasial eksternalitas tersebut.
Menurut artikel, belum ada literatur relevan yang secara langsung
menguji dampak kekuasaan negara terhadap pembentukan \emph{city region}
sekaligus heterogenitas spasialnya. (Locator: pendahuluan, hlm. 2, TXT
baris 114-127.)

\textbf{Laporan sumber.} Dalam konteks China, pemerintah secara bertahap
mengadopsi pendekatan \emph{new regionalism} dan menggunakan RIS untuk
mengurangi hambatan institusional, mempercepat aliran sumber daya lintas
batas, serta meningkatkan interaksi ekonomi dan pembagian kerja
industri. YRDRP 2010 mengarahkan industri jasa modern untuk
terkonsentrasi di metropolitan seperti Shanghai, Nanjing, dan Hangzhou,
sementara manufaktur diproyeksikan berpindah ke kota lain di wilayah
tersebut. Namun, belum jelas apakah strategi negara ini benar-benar
mendorong relokasi aktivitas sosial-ekonomi, dan belum jelas pula
bagaimana pengaruhnya berbeda menurut sektor dan kedekatan terhadap
metropolitan. (Locator: pendahuluan, hlm. 2-3, TXT baris 128-151 dan
194-207; Bagian 2.1, hlm. 3, TXT baris 168-193.)

\textbf{Interpretasi penulis.} Artikel menyusun masalah tersebut sebagai
ketegangan antara dua proses. RIS dapat memperkuat integrasi dengan
memperlancar arus faktor produksi dan pembagian kerja regional, tetapi
juga dapat memperkuat konsentrasi fungsi di kota pusat dan menekan kota
yang berdekatan dengannya. \emph{Borrowed size} dan \emph{agglomeration
shadow} dipakai untuk membaca dua kemungkinan hasil dari jaringan dan
kedekatan metropolitan. (Locator: Bagian 2.2, hlm. 4, TXT baris
208-351.)

\textbf{Inferensi pengolahan.} Secara operasional, masalah besar
pembentukan \emph{city region} dipersempit menjadi dampak RIS terhadap
pola lokasi industri sekunder dan tersier. Penyempitan ini sesuai dengan
keterbatasan data yang dinyatakan artikel, tetapi berarti bahwa
integrasi institusional, budaya, mobilitas, dan interaksi antarkota
tidak diuji sebagai variabel hasil utama.

\subsection{Objectives}\label{objectives-23}

\textbf{Laporan sumber.} Pertanyaan penelitian yang dinyatakan artikel
adalah:

\begin{itemize}
\tightlist
\item
  apakah strategi integrasi regional yang dipandu negara, sebagai simbol
  penting pembentukan \emph{city region}, memajukan pembentukan
  \emph{city region} dan sejauh mana pengaruhnya; serta
\item
  bentuk heterogenitas apa yang muncul dalam \emph{city region} yang
  diorkestrasi negara, khususnya dari sisi struktur industri dan
  spasial. (Locator: pendahuluan, hlm. 2-3, TXT baris 145-151.)
\end{itemize}

\textbf{Laporan sumber.} Tujuan tersebut diterjemahkan menjadi dua
hipotesis:

\begin{itemize}
\tightlist
\item
  \textbf{Hipotesis 1:} implementasi RIS secara signifikan mendorong
  integrasi industri dalam proses pembentukan \emph{city region}.
  Industri manufaktur atau sekunder dan industri jasa atau tersier
  diperkirakan memperlihatkan pola spasial yang berbeda, yaitu
  desentralisasi industri sekunder dan aglomerasi industri tersier.
\item
  \textbf{Hipotesis 2:} RIS meningkatkan heterogenitas spasial dalam
  proses pembentukan \emph{city region}. (Locator: Bagian 2.1-2.2, hlm.
  3-4, TXT baris 194-207 dan 322-360.)
\end{itemize}

\textbf{Laporan sumber.} Secara empiris, tujuan pertama diuji melalui
pengukuran \emph{location quotient} dan regresi DID. Tujuan kedua diuji
melalui regresi DDD dengan kelompok yang ditetapkan sebagai wilayah
berpotensi mengalami bayangan di sekitar Shanghai, Hangzhou, dan
Nanjing. Fieldwork di Jiaxing dipakai untuk menelusuri mekanisme yang
mungkin menjelaskan hasil DDD. (Locator: Bagian 3.2, hlm. 5-7, TXT baris
507-521 dan 570-585; Bagian 4.3, hlm. 10-11, TXT baris 761-806.)

\textbf{Inferensi pengolahan.} File tidak menyebutkan tujuan eksplisit
untuk mengestimasi arus aktual antarkota, mengukur besaran konsumsi
lintas batas, atau menguji efek RIS pada kategori industri yang lebih
rinci. Hal-hal tersebut tidak boleh dianggap sebagai tujuan yang telah
dicapai oleh desain penelitian ini.

\subsection{Literature Survey}\label{literature-survey-23}

\textbf{Laporan sumber.} Survei pustaka dimulai dari teori aglomerasi,
\emph{city region}, dan \emph{new regionalism}. Pembentukan \emph{city
region} dipaparkan sebagai proses integrasi ekonomi, industri,
institusional, dan budaya antara kota pusat dan wilayah sekitarnya.
Literatur ekonomi klasik menekankan diferensial sewa lahan dan perubahan
struktur inti-periferi, sedangkan ekonomi geografi baru menekankan pasar
yang berkompetisi tidak sempurna, rasionalitas terbatas, dan
\emph{increasing returns to scale}. Literatur tentang jaringan kota
kemudian menyoroti eksternalitas yang dapat melampaui batas fisik kota.
(Locator: pendahuluan, hlm. 1-2, TXT baris 65-127.)

\textbf{Laporan sumber.} Artikel membedakan dua pendekatan atas
integrasi \emph{city region}. Pendekatan pertama menekankan dampak
sosial-ekonomi dan pergerakan lintas batas tenaga kerja, modal, serta
faktor lain. Pendekatan kedua menekankan agenda negara,
\emph{reterritorialization}, \emph{rescaling}, dan respons pemerintah
terhadap persoalan reproduksi sosial serta keberlanjutan. Penulis
menyatakan bahwa kedua perspektif tersebut belum cukup sering
diintegrasikan untuk menguji hubungan kausal antara intervensi negara
dan perpindahan aktivitas sosial-ekonomi. (Locator: Bagian 2.1, hlm. 3,
TXT baris 170-193.)

\textbf{Laporan sumber.} Dalam pembahasan \emph{borrowed size}, artikel
merujuk Alonso untuk gagasan bahwa kota kecil dapat memperoleh ekonomi
aglomerasi dari kota besar yang berdekatan. Kota kecil dan menengah di
dalam jaringan perkotaan dapat menampung fungsi perkotaan yang melampaui
ukuran penduduknya dengan menarik populasi yang diperlukan dari jaringan
yang lebih luas. Sebaliknya, \emph{agglomeration shadow} merujuk pada
kemungkinan penduduk kota perifer bekerja dan mengonsumsi di kota pusat,
sehingga kota perifer memperoleh efek yang berkurang atau berlawanan
dari jaringan spasial. (Locator: Bagian 2.2, hlm. 4, TXT baris 208-321.)

\textbf{Laporan sumber.} Artikel menggunakan pengalaman Jiaxing yang
berdekatan dengan Shanghai sebagai konteks untuk membahas bayangan pada
imigrasi penduduk, peningkatan industri, dan aspek lain. Penulis
mencatat adanya saran dalam literatur \emph{new regionalism} agar
pemerintah daerah merumuskan kebijakan integrasi yang lebih kuat untuk
menghilangkan bayangan, tetapi menyatakan bahwa sedikit studi yang
memeriksa apakah dan sejauh mana bayangan tersebut berakar pada RIS.
(Locator: Bagian 2.2, hlm. 4, TXT baris 322-351.)

\textbf{Interpretasi penulis.} Survei pustaka mendukung posisi bahwa
pembentukan \emph{city region} di China tidak cukup dijelaskan oleh
kekuatan pasar. RIS diposisikan sebagai instrumen top-down yang dapat
membentuk pola industri dan jaringan, sedangkan \emph{borrowed size} dan
\emph{agglomeration shadow} diperlakukan sebagai eksternalitas spasial
yang mungkin muncul secara bersamaan. (Locator: pendahuluan, hlm. 2-3,
dan Bagian 2.1-2.2, hlm. 3-4.)

\textbf{Inferensi pengolahan.} Survei ini bersifat naratif dan digunakan
untuk membangun hipotesis. TXT tidak menyebutkan strategi pencarian,
basis data yang ditelusuri, kriteria inklusi atau eksklusi, jumlah
literatur yang dipertimbangkan, rentang waktu korpus, atau prosedur
penilaian mutu. Karena itu, kelengkapan survei pustaka dan statusnya
sebagai tinjauan sistematis tidak dapat dinilai dari file.

\subsection{Method Used}\label{method-used-23}

\subsubsection{Desain umum}\label{desain-umum}

\textbf{Laporan sumber.} Proses empiris memiliki tiga tahap. Pertama,
\emph{location quotient} digunakan untuk menggambarkan integrasi
industri YRD. Kedua, DID digunakan untuk mengestimasi dampak bersih RIS
terhadap integrasi industri. Ketiga, DDD digunakan untuk menguji
heterogenitas spasial dan menentukan apakah pengaruh RIS konsisten
dengan \emph{agglomeration shadow} atau \emph{borrowed size}. (Locator:
Bagian 3.2, hlm. 5-7, TXT baris 507-521.)

\textbf{Laporan sumber.} YRDRP yang diterapkan pada 2010 diperlakukan
sebagai eksperimen kuasi-alamiah. Kelompok perlakuan terdiri atas unit
tingkat kabupaten di dalam 16 kota tingkat prefektur inti YRD. Unit
tingkat kabupaten dari kota lain menjadi kelompok kontrol. Distrik
munisipal di dalam satu kota tingkat prefektur diperlakukan sebagai satu
unit statistik. (Locator: Bagian 3.1, hlm. 5, TXT baris 373-478.)

\subsubsection{Location quotient}\label{location-quotient}

\textbf{Laporan sumber.} Artikel menghitung \emph{location quotient}
dari output industri untuk mengukur konsentrasi atau spesialisasi
relatif suatu industri di kota dibandingkan dengan keseluruhan wilayah.
\texttt{LQSI} adalah \emph{location quotient} industri sekunder dan
\texttt{LQTI} adalah \emph{location quotient} industri tersier. Nilai
yang lebih tinggi menunjukkan konsentrasi atau spesialisasi yang lebih
besar dalam pengertian ukuran lokasi. (Locator: Bagian 3.2.1, hlm. 6-7,
TXT baris 523-536.)

\textbf{Laporan sumber.} Dalam kerangka pembentukan \emph{city region},
penurunan LQ industri sekunder diperlakukan sebagai kecenderungan yang
positif karena menunjukkan distribusi sekunder yang lebih seimbang,
sedangkan kenaikan LQ industri tersier diperlakukan sebagai
kecenderungan positif karena menunjukkan konsentrasi dan spesialisasi
jasa di kota pusat. Dengan demikian, integrasi industri tidak diwakili
oleh satu indeks gabungan, tetapi oleh arah yang berbeda dari dua LQ.
(Locator: Bagian 3.2 dan Bagian 3.3, hlm. 6-8, TXT baris 513-521 dan
602-609.)

\subsubsection{Difference-in-differences}\label{difference-in-differences}

\textbf{Laporan sumber.} Dalam model DID, \texttt{Treat\_i} bernilai 1
untuk unit yang termasuk dalam 16 kota inti dan 0 untuk unit lain.
\texttt{T\_t} bernilai 1 setelah implementasi RIS pada 2010 dan 0
sebelum periode tersebut. \texttt{RIS\_it} didefinisikan sebagai
interaksi \texttt{Treat\_i\ x\ T\_t}. Koefisien RIS digunakan untuk
mengestimasi pengaruh implementasi RIS terhadap integrasi industri. Efek
tetap kota mengendalikan heterogenitas kota dan efek tetap tahun
mengendalikan tren waktu. (Locator: Bagian 3.2.2, hlm. 6-7, TXT baris
540-568.)

\textbf{Laporan sumber.} Model DID menggunakan \texttt{LQSI} dan
\texttt{LQTI} sebagai keluaran. Artikel menyatakan bahwa koefisien RIS
yang positif berarti peningkatan spesialisasi industri dalam \emph{city
region}, sedangkan koefisien yang negatif menunjukkan kecenderungan
homogenisasi. Dalam interpretasi sektoralnya, koefisien negatif untuk
LQSI dibaca sebagai penyebaran industri sekunder dan koefisien positif
untuk LQTI dibaca sebagai konsentrasi industri tersier. (Locator: Bagian
3.2.2, hlm. 6-7, dan Bagian 4.2, hlm. 8-10, TXT baris 558-571 dan
662-759.)

\subsubsection{Difference-in-difference-in-differences}\label{difference-in-difference-in-differences}

\textbf{Laporan sumber.} Model DDD menambahkan \texttt{Sha} dan
interaksinya dengan RIS, perlakuan, serta waktu. \texttt{Sha} bernilai 1
untuk unit yang berdekatan dengan area munisipal Shanghai, Hangzhou, dan
Nanjing, yang ditetapkan sebagai wilayah berpotensi mengalami bayangan,
dan bernilai 0 untuk unit lain. Koefisien interaksi \texttt{RIS\ x\ Sha}
dipakai untuk mengestimasi heterogenitas pengaruh RIS. Nilai positif
ditafsirkan sebagai \emph{borrowed size}, sedangkan nilai negatif
ditafsirkan sebagai \emph{agglomeration shadow}. (Locator: Bagian 3.1
dan 3.2.3, hlm. 5 dan 7, TXT baris 458-478 dan 570-585.)

\textbf{Laporan sumber.} Penetapan kelompok potensial bayangan
didasarkan pada penelitian terdahulu dan investigasi lapangan di YRD.
File tidak menyatakan bahwa \texttt{Sha} adalah ukuran bayangan yang
diamati setelah perlakuan; variabel tersebut ditentukan sebelum estimasi
sebagai kategori lokasi yang berpotensi terdampak. (Locator: Bagian 3.1,
hlm. 5, TXT baris 458-478.)

\subsubsection{Data dan pengujian
tambahan}\label{data-dan-pengujian-tambahan}

\textbf{Laporan sumber.} Karena perbedaan satuan, variabel industri
sekunder, industri tersier, penduduk, PDB per kapita, investasi aset
tetap, investasi real estat, pengeluaran fiskal publik, paten, dan
struktur industri diproses ke dalam bentuk logaritmik untuk analisis
berikutnya. File tidak menjelaskan secara lengkap apakah dan bagaimana
\texttt{LQSI} serta \texttt{LQTI} ikut ditransformasikan. (Locator:
Bagian 3.3, hlm. 7-8, TXT baris 587-609.)

\textbf{Laporan sumber.} Artikel menggunakan galat baku robust pada
regresi benchmark. Uji tren paralel digunakan sebagai pemeriksaan
prasyarat DID dan DDD. Artikel melaporkan bahwa koefisien sebelum
implementasi tidak signifikan, sedangkan pengaruh setelah implementasi
menjadi signifikan dan perubahan LQ tidak tampak signifikan sampai tahun
kedua setelah RIS. (Locator: Bagian 4.1, hlm. 8, TXT baris 615-660;
Bagian 4.4, hlm. 11-12, TXT baris 874-914.)

\textbf{Laporan sumber.} Fieldwork mendalam di Jiaxing, termasuk
wawancara dan penelitian lapangan yang disebut dalam bagian hasil dan
kesimpulan, digunakan untuk memberi penjelasan mekanisme bayangan.
Artikel tidak memberikan protokol pengumpulan data kualitatif secara
rinci. (Locator: Bagian 4.3, hlm. 10-11, TXT baris 789-806; Bagian 6,
hlm. 13, TXT baris 969-977.)

\textbf{Inferensi pengolahan.} DID dan DDD bukan eksperimen acak.
Interpretasi kausal bergantung pada asumsi tren paralel, ketepatan
pembentukan kelompok, tidak adanya gangguan spasial yang mengubah
kontrol, dan spesifikasi kovariat. Artikel sendiri menyatakan bahwa
regresi benchmark menunjukkan korelasi dan membutuhkan bukti tambahan
untuk hubungan sebab-akibat. (Locator: Bagian 4.1, hlm. 8, TXT baris
653-660.)

\textbf{Inferensi pengolahan.} Persamaan DDD yang diekstrak memuat
\texttt{alpha4} dua kali, masing-masing sebelum \texttt{RIS\_it} dan
\texttt{Sha\_i}. File tidak menjelaskan apakah ini kesalahan penomoran
atau merepresentasikan parameter yang berbeda. Karena itu, spesifikasi
persamaan tidak dapat direkonstruksi lebih jauh dari TXT.

\subsection{Dataset}\label{dataset-23}

\textbf{Laporan sumber.} Sumber data utama adalah buku tahunan statistik
Provinsi Zhejiang, Provinsi Jiangsu, dan Munisipalitas Shanghai untuk
output bruto industri sekunder, output bruto industri tersier, dan PDB
per kapita. Investasi aset tetap, investasi real estat, pengeluaran
fiskal publik, dan paten invensi yang diberikan diambil dari buku
tahunan statistik kota tingkat prefektur dan kabupaten. (Locator: Bagian
3.3, hlm. 7-8, TXT baris 587-601.)

\textbf{Laporan sumber.} Variabel dan satuannya dalam Lampiran A adalah:

\begin{itemize}
\tightlist
\item
  \texttt{SecI}, output industri sekunder, dalam seratus juta RMB.
  Rerata 259,97; simpangan baku 466,17; minimum 2,49; maksimum 4.454,87.
\item
  \texttt{TerI}, output industri tersier, dalam seratus juta RMB. Rerata
  232,24; simpangan baku 545,51; minimum 2,39; maksimum 7.500,59.
\item
  \texttt{LQSI}, \emph{location quotient} industri sekunder. Rerata
  0,96; simpangan baku 0,17; minimum 0,26; maksimum 1,36.
\item
  \texttt{LQTI}, \emph{location quotient} industri tersier. Rerata 0,92;
  simpangan baku 0,15; minimum 0,60; maksimum 1,52.
\item
  \texttt{Treat}, indikator kota dalam 16 kota inti, dengan 1 untuk unit
  yang terkena RIS dan 0 untuk unit lain. Rerata 0,68.
\item
  \texttt{Sha}, indikator unit yang berdekatan dengan area munisipal
  metropolitan dan berpotensi mengalami bayangan. Rerata 0,12.
\item
  \texttt{T}, indikator waktu setelah penerapan RIS, dengan 1 untuk
  tahun setelah RIS dan 0 untuk tahun lain. Rerata 0,44.
\item
  \texttt{RIS}, interaksi wilayah inti dan waktu, ditulis sebagai
  \texttt{Core\ x\ T}. Rerata 0,30.
\item
  \texttt{Pop}, penduduk tetap, dalam sepuluh ribu orang. Rerata 93,65;
  simpangan baku 76,71; minimum 7,62; maksimum 680,67.
\item
  \texttt{pGDP}, PDB per kapita dalam RMB. Rerata 43.483,35; simpangan
  baku 36.006,68; minimum 2.900,00; maksimum 218.984,10.
\item
  \texttt{FAI}, investasi aset tetap, dalam miliar RMB. Rerata 220,94;
  simpangan baku 410,60; minimum 0,32; maksimum 5.176,24.
\item
  \texttt{REI}, investasi real estat, dalam miliar RMB. Rerata 61,42;
  simpangan baku 176,09; minimum 0,01; maksimum 2.566,90.
\item
  \texttt{PFE}, pengeluaran fiskal publik, dalam miliar RMB. Rerata
  55,13; simpangan baku 113,50; minimum 0,94; maksimum 1.396,89.
\item
  \texttt{PAT}, jumlah paten yang diberikan. Rerata 2.348,27; simpangan
  baku 4.737,55; minimum 2,00; maksimum 49.720,00.
\item
  \texttt{INDS}, struktur industri, yaitu industri tersier dibagi
  industri sekunder. Rerata 0,82; simpangan baku 0,33; minimum 0,32;
  maksimum 4,49. (Locator: Lampiran A, Tabel A1, hlm. 14-15, TXT baris
  1010-1072.)
\end{itemize}

\textbf{Laporan sumber.} Variabel RIS dibentuk dari status 16 kota inti
dan indikator waktu. Artikel menyatakan bahwa keputusan memilih RIS 2010
terkait dengan kebutuhan mempercepat integrasi kota, aliran faktor, dan
pembagian kerja industri. Periode 2000-2017 dipilih untuk mencakup
periode keterlibatan China dalam globalisasi sebelum sengketa
perdagangan China-Amerika Serikat dimulai pada 2018. (Locator: Bagian
3.1, hlm. 5-6, TXT baris 486-505.)

\textbf{Catatan simbol dataset.} Tabel regresi menggunakan simbol
\texttt{RSI} sebagai kovariat pada Tabel 1-3, sedangkan Lampiran A
mendefinisikan investasi real estat sebagai \texttt{REI}. File tidak
menjelaskan apakah \texttt{RSI} adalah salah ketik untuk \texttt{REI}
atau variabel yang berbeda. \texttt{GDP} juga muncul pada beberapa tabel
regresi, sementara definisi yang eksplisit di Lampiran lebih jelas untuk
\texttt{pGDP}. (Locator: Tabel 1-3, hlm. 8-11, TXT baris 624-651,
677-719, dan 812-858; Lampiran A, hlm. 14-15, TXT baris 1052-1060.)

\textbf{Informasi tidak tersedia dalam file.} TXT tidak memberikan
prosedur pembersihan data, aturan untuk pencilan, penanganan data
hilang, tanggal pengumpulan setiap seri, daftar unit yang masuk, akses
aktual ke data mentah, kode analisis, atau prosedur yang dapat
mereplikasi penggabungan buku tahunan. Pernyataan data hanya menyebutkan
ketersediaan atas permintaan kepada penulis korespondensi. (Locator:
Data Availability Statement, hlm. 13, TXT baris 990-998.)

\subsection{Population / Sample}\label{population-sample-23}

\textbf{Laporan sumber.} Wilayah penelitian mencakup unit tingkat
kabupaten di Munisipalitas Shanghai, Provinsi Zhejiang, dan Provinsi
Jiangsu dalam YRD. YRDRP 2010 menetapkan 16 kota inti: Shanghai; delapan
kota di Jiangsu, yaitu Nanjing, Suzhou, Wuxi, Changzhou, Zhenjiang,
Yangzhou, Taizhou, dan Nantong; serta tujuh kota di Zhejiang, yaitu
Hangzhou, Ningbo, Huzhou, Jiaxing, Shaoxing, Zhoushan, dan Taizhou. Unit
tingkat kabupaten di dalam kota-kota inti menjadi kelompok perlakuan,
sedangkan unit dari kota lain menjadi kelompok kontrol. (Locator: Bagian
3.1, hlm. 5, TXT baris 373-478.)

\textbf{Laporan sumber.} Artikel menyatakan bahwa distrik munisipal di
dalam satu kota tingkat prefektur digabung sebagai satu unit statistik
karena fokusnya adalah interaksi entitas perkotaan. Bagian studi area
menyebut total 121 unit entitas perkotaan, sedangkan abstrak menyebut
122 unit tingkat kabupaten. Model regresi pada Tabel 1-3 menggunakan
2.178 observasi. (Locator: abstrak, hlm. 1, TXT baris 45-56; Bagian 3.1,
hlm. 5, TXT baris 431-456; Tabel 1-3, hlm. 8-11, TXT baris 624-651,
677-719, dan 812-858.)

\textbf{Laporan sumber.} Untuk DDD, unit yang berdekatan dengan area
munisipal Shanghai, Nanjing, dan Hangzhou ditetapkan sebagai kelompok
\texttt{Shadow}; seluruh unit lain ditetapkan sebagai kelompok
non-shadow. Lampiran melaporkan rerata \texttt{Sha} sebesar 0,12, tetapi
jumlah pasti unit shadow dan daftar unitnya tidak disebutkan dalam teks.
(Locator: Bagian 3.1, hlm. 5, TXT baris 458-478; Lampiran A, hlm. 14,
TXT baris 1040-1047.)

\textbf{Laporan sumber.} Fieldwork berfokus pada Jiaxing, kota tingkat
prefektur yang disebut berdekatan dengan Shanghai. Artikel menyebut
investigasi lapangan, wawancara, serta contoh kabupaten Tongxiang,
Haining, Haiyan, Pinghu, dan Jiashan. Ukuran dan cara pemilihan informan
tidak diberikan. (Locator: Bagian 2.2, hlm. 4, TXT baris 322-351; Bagian
4.3, hlm. 10-11, TXT baris 789-806.)

\textbf{Inferensi pengolahan.} Desain utama lebih dekat pada sensus unit
yang ditetapkan dalam wilayah studi daripada sampel acak. Namun,
penyatuan distrik munisipal dan penggunaan label unit tingkat kabupaten
atau entitas perkotaan membuat unit analisis tidak sepenuhnya seragam.
Perbedaan 122 dan 121 juga belum memungkinkan verifikasi ukuran populasi
analitis dari TXT saja.

\textbf{Inferensi pengolahan.} Kelompok shadow ditetapkan berdasarkan
kedekatan dengan tiga metropolitan sebelum hasil DDD dipaparkan. Karena
itu, hasil DDD menguji apakah efek RIS berbeda pada lokasi yang dianggap
berpotensi mengalami bayangan, bukan mengidentifikasi lokasi bayangan
sepenuhnya dari ukuran hasil tanpa klasifikasi awal.

\subsection{Dependent Variables}\label{dependent-variables-23}

\textbf{Laporan sumber.} Variabel dependen utama adalah \texttt{LQSI}
dan \texttt{LQTI}:

\begin{itemize}
\tightlist
\item
  \texttt{LQSI} mengukur lokasi relatif output industri sekunder pada
  suatu unit terhadap lokasi industri sekunder di seluruh YRD.
\item
  \texttt{LQTI} mengukur lokasi relatif output industri tersier pada
  suatu unit terhadap lokasi industri tersier di seluruh YRD.
\end{itemize}

Keduanya dihitung dari proporsi output industri suatu unit terhadap
total wilayah dan dipakai sebagai indikator konsentrasi atau
spesialisasi industri. Persamaan LQ tercantum sebagai Persamaan (1),
tetapi tata letak persamaan pada TXT tidak sepenuhnya linear. (Locator:
Bagian 3.2.1, hlm. 6-7, TXT baris 523-536.)

\textbf{Laporan sumber.} Dalam pembacaan substantif penulis, penurunan
\texttt{LQSI} menunjukkan penyebaran atau distribusi industri sekunder
yang lebih seimbang di wilayah, sedangkan kenaikan \texttt{LQTI}
menunjukkan konsentrasi dan spesialisasi industri tersier di kota pusat.
Kedua arah tersebut secara bersama-sama diperlakukan sebagai indikator
integrasi industri dan pembentukan \emph{city region}. (Locator: Bagian
3.3, hlm. 7-8, TXT baris 602-609; Bagian 4.2, hlm. 8-10, TXT baris
662-759.)

\textbf{Laporan sumber.} Tabel A1 melaporkan rentang \texttt{LQSI}
0,26-1,36 dan rentang \texttt{LQTI} 0,60-1,52. Artikel menyatakan bahwa
seluruh variabel tertentu diproses dalam bentuk logaritmik, tetapi tidak
memberi keterangan lengkap tentang transformasi LQ. (Locator: Bagian
3.3, hlm. 7-8, TXT baris 587-609; Lampiran A, hlm. 14, TXT baris
1027-1033.)

\textbf{Inferensi pengolahan.} LQ mengukur posisi relatif atau
spesialisasi output, bukan arus relokasi industri, jumlah perusahaan
yang berpindah, lapangan kerja, produktivitas, konsumsi lintas batas,
atau interaksi antarkota. Karena itu, istilah ``integrasi industri''
dalam artikel perlu dibatasi pada perubahan pola lokasi dua sektor yang
dioperasionalkan melalui LQ.

\textbf{Inferensi pengolahan.} Tidak adanya indeks gabungan membuat arah
normatif integrasi bersifat sektoral dan asimetris: penurunan LQSI
dianggap baik, sedangkan kenaikan LQTI juga dianggap baik. File tidak
menjelaskan uji sensitivitas terhadap indeks gabungan atau ukuran
integrasi alternatif.

\subsection{Results}\label{results-23}

\subsubsection{Statistik deskriptif dan
benchmark}\label{statistik-deskriptif-dan-benchmark}

\textbf{Laporan sumber.} LQ industri sekunder dan tersier memiliki
rentang sebagaimana disebutkan dalam Lampiran A. Pada uraian deskriptif,
artikel juga menyebut rentang \texttt{LEST} 0,26-1,36 dan \texttt{LETI}
0,60-1,52, sementara simbol yang digunakan di bagian lain adalah
\texttt{LQSI} dan \texttt{LQTI}. Penulis menyatakan bahwa angka tersebut
menunjukkan aglomerasi dan spesialisasi industri tersier lebih tinggi
daripada industri sekunder. (Locator: Bagian 3.3, hlm. 7-8, TXT baris
602-609; Lampiran A, hlm. 14, TXT baris 1027-1033.)

\textbf{Laporan sumber.} Pada regresi benchmark Tabel 1, koefisien RIS
untuk \texttt{LQSI} adalah \texttt{-0,016} dengan tanda \texttt{***} dan
nilai t \texttt{-6,31}. Untuk \texttt{LQTI}, koefisien RIS dicetak
\texttt{0,002} dengan tanda \texttt{*} dan nilai t \texttt{0,71}. Kedua
model memiliki efek tetap kota dan tahun, \texttt{N\ =\ 2178}, dan
\texttt{R\_square\ =\ 0,98}. (Locator: Tabel 1, hlm. 8, TXT baris
615-651.)

\textbf{Interpretasi penulis.} Penulis membaca koefisien negatif pada
LQSI sebagai tanda bahwa distribusi industri sekunder makin seimbang,
dan koefisien positif pada LQTI sebagai tanda bahwa industri tersier
makin terpolarisasi. Hasil benchmark dinyatakan menunjukkan tren
integrasi industri yang substansial sejak RIS diterapkan. (Locator:
Bagian 4.1, hlm. 8, TXT baris 653-660.)

\textbf{Inferensi pengolahan.} Pada kolom LQTI Tabel 1, tanda
signifikansi \texttt{*} tidak tampak selaras dengan nilai t
\texttt{0,71} jika dibaca menggunakan catatan tabel yang sama. Selain
itu, file tidak memberikan p-value. Klaim signifikansi untuk hasil
benchmark LQTI karena itu bersumber dari narasi dan tanda tabel yang
tidak sepenuhnya dapat direkonsiliasi dari TXT.

\subsubsection{Estimasi DID}\label{estimasi-did}

\textbf{Laporan sumber.} Tabel 2 menunjukkan bahwa sebelum RIS,
perbedaan perlakuan dan kontrol adalah \texttt{-0,004} untuk LQSI dengan
t \texttt{-0,72} dan \texttt{-0,007} untuk LQTI dengan t
\texttt{-0,840}. Setelah RIS, perbedaan tersebut menjadi \texttt{-0,022}
untuk LQSI dengan tanda \texttt{***} dan \texttt{0,016} untuk LQTI
dengan tanda \texttt{***}. Koefisien RIS adalah \texttt{-0,019} untuk
LQSI dengan tanda \texttt{***} dan \texttt{0,024} untuk LQTI dengan
tanda \texttt{***}. (Locator: Tabel 2, hlm. 8-9, TXT baris 677-719.)

\textbf{Laporan sumber.} Model DID pada Tabel 2 menggunakan
\texttt{N\ =\ 2178}, dengan \texttt{R\_square\ =\ 0,85} untuk LQSI dan
\texttt{0,79} untuk LQTI. Pada model LQSI, kovariat yang dicetak
meliputi GDP, PDB per kapita, pengeluaran fiskal publik, investasi aset
tetap, \texttt{RSI}, dan struktur industri. Pada model LQTI, kovariat
yang dicetak meliputi penduduk, PDB per kapita, pengeluaran fiskal
publik, paten, \texttt{RSI}, dan struktur industri. (Locator: Tabel 2,
hlm. 8-9, TXT baris 677-719.)

\textbf{Interpretasi penulis.} Koefisien RIS yang negatif pada LQSI
ditafsirkan sebagai pengurangan kecenderungan aglomerasi industri
sekunder di YRD. Koefisien RIS yang positif pada LQTI ditafsirkan
sebagai peningkatan spesialisasi industri tersier. Penulis menyimpulkan
bahwa RIS memfasilitasi penyebaran industri sekunder dan konsentrasi
industri tersier, yang dianggap sebagai dampak positif terhadap
integrasi industri. (Locator: Bagian 4.2, hlm. 8-10, TXT baris 662-759.)

\textbf{Laporan sumber.} Penulis membandingkan pola tersebut dengan
kecenderungan evolusi spontan. Pemerintah lokal yang bersaing untuk
menarik investasi dan perusahaan manufaktur disebut telah mendorong
konsentrasi serta homogenisasi industri sekunder. RIS dinilai mengurangi
kompetisi tersebut. Untuk sektor tersier, artikel menyebut jasa keuangan
Shanghai, industri internet Hangzhou, dan jasa produsen Nanjing sebagai
contoh konsentrasi di metropolitan. (Locator: Bagian 4.2, hlm. 9-10, TXT
baris 730-759.)

\textbf{Inferensi pengolahan.} Estimasi DID mendukung perubahan pola
lokasi yang diasosiasikan dengan RIS, tetapi tidak dengan sendirinya
membuktikan bahwa RIS menyebabkan perpindahan faktor atau bahwa LQ yang
berubah berasal dari interaksi antarkota. Artikel mengandalkan
interpretasi kebijakan dan asumsi DID untuk melangkah dari perubahan LQ
menuju pembentukan \emph{city region}.

\subsubsection{Estimasi DDD dan heterogenitas
spasial}\label{estimasi-ddd-dan-heterogenitas-spasial}

\textbf{Laporan sumber.} Pada Tabel 3, interaksi \texttt{RIS\ x\ Sha}
untuk LQSI bernilai \texttt{-0,035} dan signifikan pada taraf 1 persen,
dengan t \texttt{-2,39}. Pada LQTI, interaksi yang sama bernilai
\texttt{-0,048} dan signifikan pada taraf 1 persen, dengan t
\texttt{-2,78}. Interaksi sebelum RIS tidak signifikan, yaitu
\texttt{0,013} untuk LQSI dan \texttt{-0,001} untuk LQTI. (Locator:
Bagian 4.3, hlm. 10-11, TXT baris 761-806; Tabel 3, hlm. 11, TXT baris
812-858.)

\textbf{Interpretasi penulis untuk industri sekunder.} Penurunan LQSI di
kota yang berbatasan dengan metropolitan dibaca sebagai efek RIS yang
lebih terlihat di wilayah tersebut dibandingkan kota yang tidak
berbatasan. Penulis mengaitkannya dengan peningkatan transportasi
antarkota dan limpahan harga lahan yang membuat relokasi manufaktur ke
kota berdekatan kurang menarik. Aksesibilitas yang membaik di seluruh
YRD kemudian memungkinkan kota perifer yang sebelumnya dirugikan oleh
transportasi menjadi tujuan relokasi manufaktur. (Locator: Bagian 4.3,
hlm. 10-11, TXT baris 776-788.)

\textbf{Interpretasi penulis untuk industri tersier.} Koefisien negatif
LQTI sebesar \texttt{-0,048} menunjukkan bahwa keunggulan relatif
industri tersier di kota yang berbatasan dengan metropolitan menurun
setelah RIS. Hal ini berjalan berlawanan dengan koefisien DID
keseluruhan yang positif \texttt{0,024}: konsentrasi industri tersier
meningkat secara keseluruhan, tetapi keunggulan di kota yang berdekatan
dengan metropolitan menurun. Penulis menafsirkan kondisi ini sebagai
\emph{agglomeration shadow} pada wilayah perifer metropolitan. (Locator:
Bagian 4.3, hlm. 10-11, TXT baris 789-806.)

\textbf{Laporan sumber dari Jiaxing.} Artikel menyebut kenaikan cepat
harga lahan dan kekurangan tenaga kerja berbakat sebagai hambatan bagi
perusahaan teknologi bernilai tambah tinggi untuk menetap di Jiaxing.
Institusi riset yang disebut adalah Jiaxing Science and Technology City,
Xiuzhou National High-tech Zone, dan Yangtze River Delta Research
Institute; semuanya dikatakan menghadapi kekurangan talenta dan
kapasitas ekspansi yang terbatas. Artikel juga melaporkan bahwa penduduk
Tongxiang, Haining, dan Haiyan cenderung pergi ke Hangzhou untuk
konsumsi akhir pekan, sedangkan penduduk Pinghu dan Jiashan cenderung
membelanjakan tabungan di Shanghai. (Locator: Bagian 4.3, hlm. 10-11,
TXT baris 795-806.)

\textbf{Interpretasi penulis.} Penulis menyatakan bahwa RIS memperkuat
\emph{agglomeration shadow} industri tersier di perifer metropolitan.
Untuk industri sekunder, artikel menggunakan dua pembacaan yang
berdekatan: penyebaran dari metropolitan ke wilayah perifer dipandang
memiliki ciri \emph{borrowed size}, tetapi kenaikan biaya lahan dan
keluarnya industri juga dapat menghasilkan kemerosotan industri sekunder
di sekitar metropolitan yang dipandang sebagai \emph{agglomeration
shadow}. Penulis menyimpulkan bahwa mekanisme kedua konsep tersebut pada
dasarnya identik. (Locator: Bagian 4.3, hlm. 11, TXT baris 860-872.)

\subsubsection{Pemeriksaan tren paralel}\label{pemeriksaan-tren-paralel}

\textbf{Laporan sumber.} Artikel menyatakan bahwa koefisien
\texttt{RIS\ x\ Before\ n} secara statistik tidak signifikan, sedangkan
koefisien setelah implementasi menjadi signifikan. Artikel menggunakan
hasil tersebut sebagai dukungan terhadap asumsi tren paralel. Artikel
juga menyatakan bahwa perubahan LQSI dan LQTI baru signifikan pada tahun
kedua setelah implementasi RIS, sehingga terdapat jeda waktu dampak.
(Locator: Bagian 4.4, hlm. 11-12, TXT baris 874-894.)

\textbf{Inferensi pengolahan.} Tabel 4 yang terekstrak menampilkan tanda
\texttt{*} pada \texttt{RIS\ x\ Before\ 1} untuk LQTI, dengan angka
\texttt{0,007} dan t \texttt{-1,24}, meskipun narasi menyatakan semua
koefisien sebelum implementasi tidak signifikan. Tabel juga menampilkan
beberapa pasangan koefisien dan t yang sulit dibaca karena tata letak.
Dari TXT saja, klaim bahwa seluruh uji pra-perlakuan tidak signifikan
perlu diberi kualifikasi.

\subsection{Findings}\label{findings-23}

\textbf{Temuan yang dilaporkan sumber.} RIS mempunyai pengaruh berbeda
menurut sektor industri. Koefisien DID negatif pada LQSI dikaitkan
dengan penyebaran atau tata letak industri sekunder yang lebih seimbang.
Koefisien positif pada LQTI dikaitkan dengan peningkatan konsentrasi
industri tersier di kota pusat. (Locator: Tabel 2 dan Bagian 4.2, hlm.
8-10, TXT baris 677-759.)

\textbf{Temuan yang dilaporkan sumber.} Heterogenitas spasial tampak di
unit yang berdekatan dengan Shanghai, Hangzhou, dan Nanjing. Interaksi
DDD \texttt{RIS\ x\ Sha} bernilai negatif dan signifikan untuk kedua
sektor. Untuk industri tersier, hal ini dibaca sebagai penurunan
keunggulan relatif di wilayah yang berbatasan dengan metropolitan. Untuk
industri sekunder, pengaruh RIS lebih menonjol di wilayah sekitar
metropolitan dalam bentuk relokasi atau penyebaran, tetapi tekanan harga
lahan juga dapat memicu keluarnya industri dan kemerosotan lokal.
(Locator: Bagian 4.3, hlm. 10-11, TXT baris 761-872.)

\textbf{Temuan yang dilaporkan sumber.} Fieldwork Jiaxing memberikan
contoh mekanisme yang dikaitkan dengan bayangan: kenaikan harga lahan,
lemahnya daya tarik talenta, keterbatasan ekspansi institusi riset, dan
kebocoran konsumsi penduduk ke Hangzhou atau Shanghai. (Locator: Bagian
4.3, hlm. 10-11, TXT baris 789-806; Bagian 6, hlm. 13, TXT baris
969-977.)

\textbf{Interpretasi penulis.} RIS dipandang memfasilitasi perpindahan
dan arus faktor melalui infrastruktur transportasi, sehingga struktur
spasial bergeser dari hierarkis menuju berjaringan. Kota pusat menjadi
lokasi ``headquarters'' dan jasa produsen, sedangkan wilayah perifer
menjadi lokasi ``branches'' dan manufaktur. (Locator: Bagian 5, hlm.
12-13, dan Bagian 6, hlm. 13, TXT baris 917-944 dan 953-968.)

\textbf{Inferensi pengolahan.} Temuan terkuat secara langsung adalah
perubahan relatif konsentrasi output dua sektor setelah kebijakan dan
perbedaan perubahan di wilayah yang telah didefinisikan sebagai
berpotensi mengalami bayangan. Bukti tersebut belum membedakan secara
langsung apakah perubahan terjadi melalui perpindahan perusahaan, arus
tenaga kerja, relokasi modal, perubahan permintaan, atau perubahan
pencatatan output.

\textbf{Inferensi pengolahan.} Hasil DDD negatif untuk LQSI dan LQTI
tidak sepenuhnya bersinonim dengan satu mekanisme. Untuk LQTI, arah
interpretasi bayangan relatif jelas dalam kerangka artikel. Untuk LQSI,
penulis sekaligus menyebut penyebaran sebagai ciri \emph{borrowed size}
dan kemerosotan akibat harga lahan sebagai \emph{agglomeration shadow}.
Karena itu, klaim bahwa RIS memperbesar satu proses tunggal perlu dibaca
sebagai kesimpulan teoritis penulis, bukan pemisahan empiris yang
tuntas.

\textbf{Inferensi pengolahan.} Fieldwork Jiaxing memperkaya penjelasan
mekanisme, tetapi karena ukuran, prosedur, dan strategi analisis data
kualitatif tidak tersedia, bukti tersebut tidak dapat dipakai untuk
mengukur prevalensi mekanisme pada seluruh YRD.

\subsection{Conclusions}\label{conclusions-23}

\textbf{Kesimpulan penulis.} Artikel menyimpulkan bahwa RIS membentuk
\emph{city region} YRD melalui restrukturisasi industri. Industri
tersier bernilai tambah tinggi cenderung mengelompok di kota pusat,
sedangkan industri sekunder menyebar ke wilayah perifer. Kota pusat
digambarkan sebagai hub yang didominasi kantor pusat dan jasa produsen,
sedangkan wilayah perifer didominasi cabang dan manufaktur. (Locator:
Bagian 6, hlm. 13, TXT baris 953-968.)

\textbf{Kesimpulan penulis.} DDD digunakan untuk menyimpulkan bahwa RIS
meningkatkan konsentrasi industri tersier di metropolitan dan
menghasilkan penurunan yang berarti di wilayah perifer. Untuk industri
sekunder, kebijakan memfasilitasi penyebaran dari metropolitan, tetapi
harga lahan dan perumahan yang tinggi di wilayah sekitar metropolitan
mendorong industri menuju lokasi yang lebih jauh. Kondisi ini dipakai
untuk menjelaskan pembentukan \emph{agglomeration shadow}. (Locator:
Bagian 6, hlm. 13, TXT baris 969-977.)

\textbf{Kesimpulan penulis.} Harga lahan yang meningkat cepat dan
ketidakcukupan daya tarik talenta diidentifikasi melalui penelitian di
Jiaxing sebagai faktor utama di balik bayangan \emph{city region}.
Artikel menyatakan bahwa kedua faktor tersebut dapat diteliti lebih
lanjut sebagai variabel perantara. (Locator: Bagian 6, hlm. 13, TXT
baris 969-977.)

\textbf{Interpretasi penulis.} Penulis mengusulkan bahwa kebijakan
integrasi regional mengubah model sewa yang berpusat pada kota
individual menjadi model yang mencakup keseluruhan wilayah dengan
memfasilitasi mobilitas faktor. Penjelasan ini dinyatakan sebagai
potensi penjelasan kausal, bukan sebagai estimasi terpisah atas setiap
rantai mekanisme. (Locator: Bagian 6, hlm. 13, TXT baris 959-968.)

\textbf{Inferensi pengolahan.} Kesimpulan empiris paling aman adalah
bahwa RIS berasosiasi dengan pola perubahan LQSI dan LQTI yang berbeda
secara sektoral dan spasial dalam unit studi YRD. TXT belum memberi
dasar untuk menggeneralisasi efek tersebut ke seluruh China, wilayah
lain, periode setelah 2017, atau semua bentuk pembentukan \emph{city
region}.

\subsection{Contributions}\label{contributions-23}

\textbf{Kontribusi yang diklaim atau ditunjukkan sumber.}

\begin{itemize}
\tightlist
\item
  \textbf{Kontribusi teoritis:} artikel membawa peran kekuasaan negara
  dan RIS ke dalam pembahasan pembentukan \emph{city region}, yang
  sebelumnya lebih sering menekankan redistribusi spontan berbasis
  pasar. (Locator: pendahuluan, hlm. 2, TXT baris 103-127.)
\item
  \textbf{Kontribusi substantif:} artikel menilai integrasi industri
  dengan membedakan pola industri sekunder dan tersier, bukan
  memperlakukan integrasi sebagai perubahan ekonomi yang homogen.
  (Locator: Bagian 2.1 dan 3.2, hlm. 3 dan 5-7, TXT baris 194-207 dan
  507-521.)
\item
  \textbf{Kontribusi skala analisis:} penggunaan unit tingkat kabupaten
  atau entitas perkotaan dimaksudkan untuk menangkap interaksi yang
  lebih rinci dibandingkan skala kota tingkat prefektur yang lebih
  agregat. (Locator: Bagian 5, hlm. 12-13, TXT baris 930-944.)
\item
  \textbf{Kontribusi metode:} artikel menggabungkan LQ, DID, DDD, uji
  tren paralel, dan fieldwork Jiaxing untuk menghubungkan dampak
  kebijakan dengan heterogenitas spasial dan penjelasan mekanisme.
  (Locator: Bagian 3.2-3.3, hlm. 5-8, dan Bagian 4.4, hlm. 11-12.)
\item
  \textbf{Kontribusi empiris:} artikel melaporkan bahwa RIS dapat
  mendorong penyebaran industri sekunder dan konsentrasi industri
  tersier secara bersamaan, dengan efek yang berbeda di dekat
  metropolitan. (Locator: Bagian 4.2-4.3, hlm. 8-11, TXT baris 662-872.)
\item
  \textbf{Kontribusi pada studi bayangan aglomerasi:} artikel
  menghubungkan \emph{agglomeration shadow} dengan kebijakan integrasi
  regional, harga lahan, kekurangan talenta, dan kebocoran konsumsi di
  wilayah antara metropolitan. (Locator: Bagian 4.3 dan 6, hlm. 10-13,
  TXT baris 789-806 dan 969-977.)
\end{itemize}

\textbf{Inferensi pengolahan.} Nilai tambah paling kuat artikel adalah
memperlihatkan bahwa integrasi regional yang diukur melalui industri
dapat bersifat asimetris menurut sektor dan lokasi. Klaim mekanisme
\emph{borrowed size} dan \emph{agglomeration shadow} lebih tepat
dipandang sebagai kontribusi eksploratif karena proksi hasilnya berbasis
LQ dan kelompok shadow ditentukan berdasarkan kedekatan awal.

\subsection{Research Gap}\label{research-gap-23}

\textbf{Kesenjangan yang dinyatakan sumber.}

\begin{itemize}
\tightlist
\item
  Peran kekuasaan negara dalam pembentukan \emph{city region} dan
  heterogenitas spasialnya belum diteliti secara langsung dalam
  literatur yang dirujuk artikel. (Locator: pendahuluan, hlm. 2, TXT
  baris 114-127.)
\item
  Hubungan kausal antara implementasi RIS sebagai intervensi negara dan
  perpindahan aktivitas sosial-ekonomi antarkota belum jelas. (Locator:
  Bagian 2.1, hlm. 3, TXT baris 186-193.)
\item
  Sedikit studi yang menguji apakah bayangan spasial dan sejauh mana
  bayangan tersebut berakar pada RIS. (Locator: Bagian 2.2, hlm. 4, TXT
  baris 322-351.)
\item
  Perspektif yang mengukur interaksi antarkota dan perspektif yang
  menjelaskan kegunaan kebijakan negara belum cukup terintegrasi.
  (Locator: Bagian 2.1, hlm. 3, TXT baris 174-193.)
\item
  Pengaruh bottom-up dalam pembentukan \emph{city region} belum diteliti
  secara memadai karena hipotesis artikel berfokus pada strategi
  top-down. (Locator: Bagian 5, hlm. 12-13, TXT baris 917-929.)
\item
  Keterbatasan data menghalangi penilaian langsung atas interaksi
  antarunit tingkat kabupaten. (Locator: Bagian 5, hlm. 12-13, TXT baris
  930-951.)
\item
  Data industri tingkat kabupaten hanya memungkinkan klasifikasi ke
  dalam sektor sekunder dan tersier, bukan kategori industri yang lebih
  rinci. (Locator: Bagian 3.2, hlm. 5-6, TXT baris 507-521.)
\end{itemize}

\textbf{Kesenjangan yang masih tersisa menurut artikel.} Artikel
menyebut perlunya mengeksplorasi pengaruh bottom-up, menjembatani
kekurangan pengukuran interaksi tingkat kabupaten, serta meneliti harga
lahan dan daya tarik talenta sebagai variabel perantara. (Locator:
Bagian 5, hlm. 12-13, TXT baris 917-951; Bagian 6, hlm. 13, TXT baris
969-977.)

\textbf{Inferensi pengolahan.} Artikel belum menutup kesenjangan antara
perubahan LQ dan proses peminjaman atau bayangan yang dialami aktor.
Masih belum tersedia pengukuran langsung atas arus komuter, migrasi,
konsumsi, modal, perusahaan, talenta, atau relokasi industri.
Kesenjangan ini diturunkan dari definisi variabel dan tidak dinyatakan
sebagai hasil tambahan oleh penulis.

\textbf{Inferensi pengolahan.} Perlu juga diuji apakah hasil bertahan
ketika batas YRD, definisi unit, kelompok shadow, tahun kebijakan,
kategori sektor, dan ukuran integrasi diubah. Usulan ini merupakan
inferensi review, bukan agenda yang secara eksplisit dijabarkan dalam
file.

\subsection{Limitations}\label{limitations-23}

\textbf{Keterbatasan yang dinyatakan penulis.}

\begin{itemize}
\tightlist
\item
  Keterbatasan data tingkat kabupaten membuat industri hanya dapat
  dibagi menjadi sektor sekunder dan tersier, tanpa kategori industri
  yang lebih rinci. (Locator: Bagian 3.2, hlm. 5-6, TXT baris 513-521.)
\item
  Keterbatasan data menghalangi penilaian langsung atas interaksi
  antarunit tingkat kabupaten. Artikel menyebut penelitian ini sebagai
  eksplorasi awal yang membutuhkan penelitian lanjutan. (Locator: Bagian
  5, hlm. 12-13, TXT baris 930-951.)
\item
  Kerangka hipotesis menekankan pengaruh strategi negara dari atas ke
  bawah, sedangkan pengaruh bottom-up masih perlu diteliti. (Locator:
  Bagian 5, hlm. 12-13, TXT baris 917-929.)
\item
  Harga lahan dan daya tarik talenta baru diidentifikasi sebagai
  kemungkinan variabel perantara melalui fieldwork Jiaxing dan belum
  diuji sebagai bagian terpisah dari model utama. (Locator: Bagian 6,
  hlm. 13, TXT baris 969-977.)
\end{itemize}

\textbf{Keterbatasan yang diturunkan dari isi file sebagai inferensi
pengolahan.}

\begin{itemize}
\tightlist
\item
  \textbf{Ketidakselarasan sampel:} abstrak menyebut 122 unit, sedangkan
  bagian studi area dan kesimpulan menyebut 121 entitas perkotaan. Tabel
  regresi melaporkan \texttt{N\ =\ 2178}. File tidak menyelesaikan
  perbedaan tersebut. (Locator: abstrak, hlm. 1; Bagian 3.1, hlm. 5;
  Bagian 6, hlm. 13; TXT baris 45-56, 431-456, dan 953-958.)
\item
  \textbf{Validitas konstruk:} LQSI dan LQTI menunjukkan konsentrasi
  atau spesialisasi output relatif. Keduanya tidak menunjukkan arus
  aktual, mutu, kapasitas, produktivitas, jumlah pekerja, relokasi
  perusahaan, atau pemanfaatan lintas batas.
\item
  \textbf{Definisi integrasi yang asimetris:} penurunan LQSI dan
  kenaikan LQTI sama-sama diberi makna positif, tetapi file tidak
  menawarkan indeks gabungan atau pengujian alternatif untuk memastikan
  bahwa kedua arah tersebut mewakili satu konstruk integrasi.
\item
  \textbf{Klasifikasi shadow:} kelompok \texttt{Sha} ditentukan dari
  kedekatan dengan tiga metropolitan dan investigasi terdahulu atau
  lapangan. Alasan pemilihan hanya tiga metropolitan, ambang kedekatan,
  dan daftar lengkap unit tidak disebutkan dalam file. (Locator: Bagian
  3.1, hlm. 5.)
\item
  \textbf{Identifikasi kebijakan:} RIS tidak ditentukan melalui
  penugasan acak. Kota inti berasal dari rencana regional dan berpotensi
  memiliki karakteristik awal yang berbeda dari unit kontrol. File
  menyajikan uji tren paralel, tetapi tidak menyebut placebo, alternatif
  tahun kebijakan, pencocokan, atau strategi lain untuk menguji seleksi
  kebijakan.
\item
  \textbf{Interferensi spasial:} unit kontrol berada di wilayah yang
  sama dengan unit perlakuan dan dapat ikut menerima pengaruh
  transportasi atau arus industri. File tidak menyebut uji spillover,
  autokorelasi spasial, atau model dependensi spasial.
\item
  \textbf{Transformasi variabel:} artikel menyatakan banyak variabel
  dilogaritmakan, tetapi transformasi yang diterapkan pada LQSI dan LQTI
  serta formula lengkap variabel regresi tidak dijelaskan secara
  memadai.
\item
  \textbf{Ketidakjelasan simbol:} \texttt{RSI} muncul dalam Tabel 1-3,
  sedangkan Lampiran mendefinisikan \texttt{REI}; GDP muncul di tabel
  tanpa definisi seterang pGDP. File tidak memungkinkan pemastian apakah
  ini salah ketik atau variabel berbeda.
\item
  \textbf{Ketidakkonsistenan persamaan dan tabel:} Persamaan DDD
  menampilkan \texttt{alpha4} dua kali. Tabel 1 mencetak tanda
  signifikansi pada LQTI dengan t \texttt{0,71}, dan Tabel 4 tampak
  mencetak tanda \texttt{*} pada koefisien pra-perlakuan LQTI meskipun
  narasinya menyatakan tidak signifikan. Beberapa tanda t dalam Tabel 2
  juga tidak mengikuti tanda koefisien secara jelas. (Locator: Persamaan
  3, hlm. 7; Tabel 1-4, hlm. 8-12.)
\item
  \textbf{Pemeriksaan tren paralel:} artikel melaporkan hasil yang
  mendukung tren paralel, tetapi detail lengkap definisi periode
  referensi, spesifikasi galat, dan rekonstruksi seluruh statistik tidak
  tersedia dalam ekstraksi TXT.
\item
  \textbf{Data kualitatif:} artikel menyebut fieldwork dan wawancara di
  Jiaxing, tetapi tidak memberikan jumlah informan, kriteria pemilihan,
  tanggal, pertanyaan, prosedur pencatatan, metode pengodean, atau cara
  menghubungkan bukti kualitatif dengan estimasi kuantitatif.
\item
  \textbf{Kelengkapan data:} prosedur penanganan data hilang, pencilan,
  batas administratif yang berubah, dan perbedaan tahun pengamatan antar
  buku tahunan tidak disebutkan.
\item
  \textbf{Replikasi:} kode dan data mentah tidak tersedia dalam TXT.
  Data hanya dinyatakan tersedia atas permintaan kepada penulis
  korespondensi.
\item
  \textbf{Validitas eksternal:} periode 2000-2017 dan wilayah YRD
  mencerminkan satu konteks kebijakan dan ekonomi yang spesifik. File
  tidak menyediakan dasar untuk menggeneralisasikan hasil ke wilayah
  China lain, negara lain, atau periode setelah 2017.
\end{itemize}

\subsection{Challenges}\label{challenges-23}

\textbf{Keterangan sumber.} File tidak memiliki bagian khusus bernama
\texttt{Challenges}. Tantangan berikut diringkas dari bagian metode,
diskusi, dan kesimpulan; statusnya tetap dibedakan dari keterbatasan
formal.

\textbf{Tantangan yang dilaporkan sumber.}

\begin{itemize}
\tightlist
\item
  memperoleh data sosial-ekonomi yang cukup panjang dan konsisten pada
  tingkat kabupaten untuk periode 2000-2017;
\item
  mengukur pembentukan \emph{city region} ketika data hanya memungkinkan
  pembedaan sektor sekunder dan tersier;
\item
  membedakan integrasi yang dipengaruhi kebijakan negara dari
  perpindahan aktivitas sosial-ekonomi dan interaksi antarkota;
\item
  menilai heterogenitas spasial di wilayah yang berdekatan dengan
  Shanghai, Hangzhou, dan Nanjing; serta
\item
  menelusuri mekanisme di Jiaxing ketika harga lahan, kekurangan
  talenta, dan konsumsi lintas metropolitan berlangsung bersamaan.
  (Locator: Bagian 2.1-2.2, hlm. 3-4; Bagian 3.2-3.3, hlm. 5-8; Bagian
  4.3, hlm. 10-11.)
\end{itemize}

\textbf{Inferensi pengolahan.} Tantangan analitis utama adalah
memisahkan beberapa proses yang menghasilkan arah LQ serupa: penyebaran
manufaktur, keluarnya manufaktur dari wilayah mahal, konsentrasi jasa di
pusat, kebocoran konsumsi, dan perubahan relatif output. Desain
LQ-DID-DDD membantu mendeteksi pola kebijakan, tetapi tidak
mengidentifikasi setiap saluran tersebut secara mandiri.

\textbf{Inferensi pengolahan.} Tantangan replikasi juga besar karena
jumlah unit tidak konsisten, simbol kovariat tidak seragam, detail
transformasi dan pengolahan data tidak lengkap, serta teks tabel
mengalami gangguan tata letak. Informasi tersebut tidak boleh diisi
dengan asumsi dari sumber lain.

\subsection{Practical Implications}\label{practical-implications-23}

\textbf{Rekomendasi penulis.} Pemerintah daerah disarankan
memperhitungkan perbedaan yang mendasari perkembangan industri, wilayah,
dan ekonomi perkotaan ketika menerapkan integrasi regional. Penulis
menekankan bahwa dampak RIS yang berbeda terhadap industri sekunder dan
tersier harus dipertimbangkan secara bersamaan. (Locator: Bagian 6, hlm.
13, TXT baris 978-985.)

\textbf{Rekomendasi penulis.} Kota dengan orientasi pembangunan industri
yang berbeda disarankan menggunakan kerangka kebijakan yang berbeda
pula, sehingga keunggulan masing-masing sektor dapat dipakai untuk
mencapai pembangunan regional. Pemerintah lokal juga disarankan
mengenali eksternalitas khusus dari pembentukan \emph{city region} dan
memperkuat kerja sama dengan unit tetangga secara terarah. (Locator:
Bagian 6, hlm. 13, TXT baris 978-988.)

\textbf{Implikasi dari hasil yang dilaporkan.} RIS tidak seharusnya
dinilai hanya dari kenaikan integrasi agregat. Evaluasi perlu memeriksa
apakah penyebaran industri sekunder benar-benar memberi manfaat kepada
kota perifer dan apakah konsentrasi industri tersier memperlemah kota
yang berada di antara metropolitan. (Locator: Bagian 4.2-4.3, hlm.
8-11.)

\textbf{Inferensi pengolahan.} Untuk perencanaan, LQ dapat dipakai
sebagai alat penyaringan awal untuk memantau perubahan relatif sektor,
bukan sebagai satu-satunya dasar investasi atau klaim tentang
kesejahteraan. Verifikasi harga lahan, kapasitas talenta, arus konsumsi,
pekerjaan, dan kebutuhan penduduk tetap diperlukan sebelum kebijakan
diarahkan kepada kota yang dianggap mengalami bayangan.

\textbf{Inferensi pengolahan.} Peningkatan konektivitas perlu dibarengi
pemantauan distribusi manfaat. Artikel menunjukkan bahwa akses dan
integrasi dapat memperkuat hubungan dengan metropolitan sekaligus
memperbesar konsentrasi pusat atau kebocoran konsumsi. Namun, file tidak
menguji instrumen kebijakan tertentu sehingga implikasi tersebut belum
merupakan evaluasi dampak kebijakan.

\subsection{Applications}\label{applications-23}

\textbf{Aplikasi yang ditunjukkan oleh desain studi.}

\begin{itemize}
\tightlist
\item
  menggunakan LQ untuk membandingkan perubahan konsentrasi dan
  spesialisasi industri sekunder serta tersier di dalam suatu wilayah;
\item
  menggunakan DID untuk mengevaluasi hubungan antara kebijakan integrasi
  regional dan perubahan pola industri ketika tersedia kelompok
  perlakuan, kontrol, serta periode sebelum dan sesudah kebijakan;
\item
  menggunakan DDD untuk menguji apakah pengaruh kebijakan berbeda antara
  unit yang berdekatan dengan metropolitan dan unit lain;
\item
  menggabungkan analisis kuantitatif dengan fieldwork untuk menelusuri
  kemungkinan mekanisme, seperti harga lahan, talenta, dan konsumsi
  lintas batas; serta
\item
  menyediakan bahan awal bagi penataan fungsi industri dan koordinasi
  antarkota di YRD. (Locator: Bagian 3.2-3.3, hlm. 5-8; Bagian 4.2-4.4,
  hlm. 8-12; Bagian 6, hlm. 13.)
\end{itemize}

\textbf{Batas aplikasi.} Artikel tidak melaporkan implementasi hasil di
wilayah lain, perangkat lunak atau kode yang dibagikan, instrumen
kebijakan yang dievaluasi setelah penelitian, ataupun studi evaluasi
pascapenerapan. Informasi tersebut tidak disebutkan dalam file.

\textbf{Inferensi pengolahan.} Kerangka ini dapat diadaptasi sebagai
diagnosis eksploratif di wilayah metropolitan lain hanya jika definisi
unit, periode kebijakan, klasifikasi industri, ketersediaan data, asumsi
tren paralel, dan kelompok pembanding dapat dipertanggungjawabkan.
Adaptasi tersebut bukan validasi otomatis bahwa hasil YRD berlaku di
tempat lain.

\subsection{Future Research
(Optional)}\label{future-research-optional-23}

\textbf{Agenda yang dinyatakan penulis.}

\begin{itemize}
\tightlist
\item
  mengeksplorasi pengaruh bottom-up dalam pembentukan \emph{city region}
  sebagai pelengkap terhadap fokus top-down artikel;
\item
  menjembatani kekurangan data tentang interaksi antarunit tingkat
  kabupaten;
\item
  meneliti harga lahan dan daya tarik talenta sebagai variabel perantara
  yang mungkin menjelaskan \emph{agglomeration shadow};
\item
  melanjutkan eksplorasi awal atas hubungan antara RIS, \emph{borrowed
  size}, dan \emph{agglomeration shadow}. (Locator: Bagian 5, hlm.
  12-13, TXT baris 917-951; Bagian 6, hlm. 13, TXT baris 969-977.)
\end{itemize}

\textbf{Inferensi pengolahan.} Agenda tambahan yang logis dari
keterbatasan TXT adalah:

\begin{itemize}
\tightlist
\item
  memperluas pengukuran industri ketika data tingkat kabupaten
  memungkinkan klasifikasi yang lebih rinci;
\item
  mengumpulkan data langsung mengenai relokasi perusahaan, lapangan
  kerja, komuter, migrasi, konsumsi, modal, dan arus pengetahuan;
\item
  menggunakan desain longitudinal atau kuasi-eksperimental tambahan
  untuk membedakan efek RIS dari seleksi kota inti dan tren ekonomi yang
  sudah ada;
\item
  menguji spillover dan dependensi spasial antarkota, termasuk
  kemungkinan bahwa unit kontrol ikut terkena pengaruh RIS;
\item
  memeriksa kepekaan terhadap definisi kelompok shadow, ambang
  kedekatan, batas YRD, tahun perlakuan, pembobotan sektor, dan
  transformasi LQ; serta
\item
  mendokumentasikan protokol fieldwork Jiaxing dan membandingkannya
  dengan kota yang berada di luar kelompok shadow.
\end{itemize}

Usulan ini adalah inferensi review, bukan agenda eksplisit yang
seluruhnya ditulis oleh penulis.

\subsection{Provenance Notes}\label{provenance-notes-23}

\begin{itemize}
\tightlist
\item
  \textbf{Basis akses:} review ini hanya menggunakan file TXT
  \texttt{source/journal/A08-zhen-shi-lu-2023-the-impact-of-regional-integration-strategies-on-the-formation-of-city-regions-and-its-agglomeration-shadow-evidence-from-the-yangtze-river-delta-china-10.3390-land12051053.txt}.
  Tidak ada PDF, sumber eksternal, metadata daring, catatan Zotero, atau
  artikel dalam daftar pustaka yang dibuka untuk menambah atau
  memverifikasi substansi.
\item
  \textbf{Cakupan baca:} seluruh isi TXT dibaca, mulai dari blok
  metadata PDF, kepala artikel, abstrak, pendahuluan, tinjauan pustaka,
  metode, Tabel 1-4, hasil, diskusi, kesimpulan, pernyataan penulis,
  Lampiran A Tabel A1, daftar pustaka, hingga disclaimer penerbit.
  Pembacaan mencakup seluruh baris yang tersedia sampai akhir file.
\item
  \textbf{Asal teks:} blok metadata menyatakan bahwa PDF berjumlah 16
  halaman dan diekstrak pada 31 Agustus 2026 dengan
  \texttt{pdftotext\ -layout}. Metadata juga menyatakan PDF tidak
  terenkripsi dan tidak bertag. (Locator: TXT baris 1-28.)
\item
  \textbf{Locator:} \texttt{hlm.} merujuk pada nomor halaman artikel
  yang tercetak dalam ekstraksi, sedangkan \texttt{Bagian},
  \texttt{Persamaan}, \texttt{Tabel}, dan \texttt{Lampiran} merujuk pada
  label sumber. Nomor baris TXT ditambahkan untuk klaim penting dan
  untuk bagian yang mengalami ketidakselarasan.
\item
  \textbf{Keterbatasan ekstraksi:} teks dua kolom mengandung pengulangan
  dan interleaving pada bagian sekitar hlm. 4-5, terutama TXT baris
  208-360 dan 382-478. Tabel dan persamaan juga tidak selalu tersusun
  linear. Review mempertahankan isi yang dapat dibaca dan menandai
  bagian yang tidak dapat direkonstruksi dengan pasti.
\item
  \textbf{Visual yang tidak tersedia:} Figure 1 dan Figure 2 memiliki
  keterangan dan uraian tekstual dalam TXT, tetapi bentuk grafis
  lengkapnya tidak tersedia sebagai data terstruktur. Review tidak
  memperkirakan nilai atau batas spasial dari gambar.
\item
  \textbf{Ketidakkonsistenan yang dipertahankan:} perbedaan 122 unit
  dalam abstrak dan 121 unit dalam bagian metode, pengulangan
  \texttt{alpha4} dalam Persamaan 3, perbedaan simbol \texttt{RSI} dan
  \texttt{REI}, serta beberapa ketidakselarasan tanda signifikansi dan
  statistik t tidak diperbaiki berdasarkan dugaan.
\item
  \textbf{Informasi tidak tersedia:} daftar lengkap unit, jumlah pasti
  kelompok \texttt{Sha}, ambang kedekatan, prosedur pembersihan dan
  imputasi data, penanganan pencilan atau data hilang, formula
  transformasi lengkap, rincian galat dan pengelompokan galat, uji
  spillover spasial, protokol fieldwork, kode, dan data mentah tidak
  tersedia dalam TXT.
\item
  \textbf{Status penelaahan:} file memuat pernyataan bahwa penulis
  berterima kasih kepada anonymous reviewers dan editors, tetapi status
  penelaahan sejawat tidak diverifikasi dari luar TXT dan tidak
  diperlakukan sebagai fakta tambahan.
\item
  \textbf{Pemisahan tingkat klaim:} bagian berlabel \textbf{Laporan
  sumber}, \textbf{Temuan yang dilaporkan sumber}, \textbf{Kesimpulan
  penulis}, dan \textbf{Rekomendasi penulis} merupakan parafrasa atau
  pelaporan isi artikel. \textbf{Interpretasi penulis} merekam makna
  yang diberikan artikel terhadap hasilnya. \textbf{Inferensi
  pengolahan} adalah evaluasi review yang diturunkan hanya dari isi dan
  ketidaklengkapan internal TXT, bukan klaim tambahan dari artikel.
\item
  \textbf{Batas epistemik:} angka, definisi, dan hasil statistik hanya
  dilaporkan sejauh dapat ditelusuri ke TXT. Tidak ada informasi dari
  sumber luar yang digunakan untuk mengisi informasi yang tidak
  tersedia.
\end{itemize}

\section{Review A09 - Shi et
al.~(2025)}\label{review-a09---shi-et-al.-2025}

Status: Review literatur berbasis seluruh teks TXT hasil ekstraksi PDF.

\subsection{Identitas Artikel}\label{identitas-artikel-2}

\begin{itemize}
\tightlist
\item
  \textbf{Kode:} A09.
\item
  \textbf{Judul:} \emph{Borrowed size and borrowed administrative power:
  Effects of high-speed rail network on industrial upgrading and
  variegated externalities in the Yangtze River Delta, China}.
\item
  \textbf{Penulis:} Dehao Shi, Lei Wang, Xianchun Zhang, dan Tao Yu. Lei
  Wang dan Xianchun Zhang dicantumkan sebagai penulis korespondensi.
\item
  \textbf{Afiliasi:} School of Architecture and Urban Planning, Nanjing
  University; Department of Geography and Resource Management, The
  Chinese University of Hong Kong; Key Laboratory of Lake and Watershed
  Science for Water Security, Nanjing Institute of Geography and
  Limnology, Chinese Academy of Sciences; University of Chinese Academy
  of Sciences; Department of Planning, Property and Environmental
  Management, The University of Manchester; dan School of Public
  Affairs, Zhejiang University.
\item
  \textbf{Jurnal:} \emph{Journal of Transport Geography}, volume 123,
  tahun 2025, artikel 104113.
\item
  \textbf{DOI:} \texttt{10.1016/j.jtrangeo.2025.104113}.
\item
  \textbf{Tanggal editorial yang tersedia:} diterima 14 Januari 2024,
  direvisi 20 November 2024, diterima 3 Januari 2025, dan tersedia
  daring 8 Januari 2025.
\item
  \textbf{Kata kunci:} high-speed rail, industrial upgrading,
  agglomeration externality, administrative externality, dan
  difference-in-difference-in-differences (DDD) model.
\item
  \textbf{Wilayah dan unit analisis:} Yangtze River Delta (YRD), dengan
  199 unit tingkat county yang dianalisis selama 2000-2017.
\item
  \textbf{Jenis berdasarkan isi TXT:} artikel penelitian empiris. Status
  penelaahan sejawat tidak disebutkan dalam file.
\item
  \textbf{Ketersediaan data:} artikel menyatakan bahwa data akan
  tersedia berdasarkan permintaan.
\item
  \textbf{Locator identitas:} TXT baris 1-28 untuk metadata PDF; baris
  42-107 untuk kepala artikel, abstrak, DOI, dan tanggal editorial;
  baris 787-790 untuk pernyataan ketersediaan data.
\end{itemize}

\subsection{Summarized Abstract}\label{summarized-abstract-24}

\textbf{Laporan sumber.} Artikel mengkaji bagaimana pembangunan jaringan
kereta api berkecepatan tinggi atau high-speed rail (HSR) memengaruhi
peningkatan struktur industri dan perkembangan regional di YRD. Analisis
tidak hanya memakai kerangka eksternalitas aglomerasi berdasarkan ukuran
kota, tetapi juga memasukkan kekuatan administratif sebagai sumber
eksternalitas yang berbeda. Data yang digunakan berupa panel 199 unit
tingkat county selama 2000-2017. Penulis menerapkan
difference-in-differences (DID) bertahap dan
difference-in-difference-in-differences (DDD). Dalam kerangka tersebut,
\emph{borrowed size} dan \emph{agglomeration shadow} dipakai untuk
membaca perbedaan efek HSR di sekitar pusat ekonomi besar, sedangkan
\emph{borrowed administrative power} dan \emph{administrative shadow}
dipakai untuk membaca perbedaan di sekitar pusat administratif
berpangkat tinggi. (TXT baris 62-80 dan 149-161.)

\textbf{Laporan sumber.} Abstrak melaporkan bahwa HSR secara signifikan
mendorong peningkatan industri di YRD. Output sektor sekunder berkurang
4,4\%, sedangkan output sektor tersier meningkat 7,4\%. Abstrak juga
menyatakan adanya hubungan positif antara ukuran kota dan keluarnya
sektor manufaktur, serta bahwa kekuatan administratif meningkatkan
konsentrasi sektor jasa di wilayah yang terhubung HSR. (TXT baris
64-80.)

\textbf{Interpretasi penulis.} Penulis membaca HSR sebagai mekanisme
perantara yang menghubungkan geografi transportasi dengan ekonomi
politik perkotaan. Efek HSR dipandang tidak seragam karena dipengaruhi
oleh skala ekonomi kota, posisi administratif pemerintah kota, dan
karakter sektor industri. Dengan demikian, HSR dapat menghasilkan
eksternalitas jaringan yang menguntungkan sebagian wilayah melalui
\emph{borrowed size}, tetapi juga menghasilkan bayangan aglomerasi atau
bayangan administratif pada wilayah lain. (TXT baris 68-80, 130-161, dan
737-740.)

\textbf{Inferensi pengolahan.} Bukti yang tersedia mengukur efek
rata-rata HSR dan perbedaan efeknya pada kelompok county yang telah
ditetapkan berdasarkan kedekatan dengan pusat ekonomi atau
administratif. File tidak menunjukkan arus penduduk, pasar, layanan
publik, kewenangan, negosiasi, frekuensi kereta, atau perpindahan fungsi
administratif secara langsung. Oleh karena itu, istilah \emph{borrowed
size} dan \emph{borrowed administrative power} dalam review ini dipahami
sebagai interpretasi atas koefisien interaksi DDD dan definisi kelompok
potensial, bukan observasi langsung atas proses peminjaman. (TXT baris
281-296 dan 311-330.)

\subsection{Problem Statement}\label{problem-statement-24}

\textbf{Laporan sumber.} Literatur HSR yang dibahas artikel telah
meneliti urbanisasi, pertumbuhan ekonomi, perubahan industri, migrasi,
investasi aset tetap, perpindahan manufaktur, dan konsentrasi jasa.
Penjelasan yang umum dipakai menekankan pengurangan waktu dan biaya
perjalanan, mobilitas faktor, efek limpahan, efek koridor, kompresi
ruang-waktu, dan ekonomi aglomerasi. Namun, penelitian terdahulu juga
menunjukkan hasil yang berlawanan: kota besar China dapat memperoleh
perbaikan aksesibilitas dan peningkatan industri yang lebih besar,
sementara sebagian kota kecil dan menengah tidak mengalami perbaikan
atau bahkan mengalami efek merugikan. (TXT baris 111-147.)

\textbf{Laporan sumber.} Artikel mengidentifikasi dua persoalan utama.
Pertama, evolusi spasial regional tidak hanya dibentuk oleh aglomerasi
yang berpusat pada kota besar, tetapi juga oleh interaksi jaringan
antarkota kecil dan menengah. Peran kota pinggiran dinilai kurang
diperhatikan dalam pendekatan yang berpusat pada kota inti. Kedua,
terdapat persoalan endogenitas antara pembangunan HSR dan peningkatan
industri: wilayah dengan potensi pertumbuhan yang lebih besar mungkin
memang lebih mungkin dipilih untuk pembangunan HSR. (TXT baris 149-161.)

\textbf{Laporan sumber.} Persoalan lain adalah belum jelasnya perbedaan
antara pengaruh ukuran ekonomi dan pengaruh kekuatan administratif.
Pemerintah kota berpangkat tinggi dapat memengaruhi rute, lokasi
stasiun, dan frekuensi kereta melalui keputusan serta negosiasi
antartingkat pemerintahan. Pada saat yang sama, beberapa kota yang lebih
rendah tingkatnya dapat memakai infrastruktur lintas batas dan integrasi
regional untuk memanfaatkan kapasitas atau sumber daya dari kota yang
lebih tinggi. Artikel menyatakan bahwa sedikit penelitian yang
membandingkan secara rinci kesamaan dan perbedaan kedua sumber pengaruh
tersebut. (TXT baris 111-137 dan 222-250.)

\textbf{Laporan sumber.} Artikel juga memandang peningkatan industri
sebagai proses yang berbeda menurut sektor. Literatur yang dirujuk
penulis mengarah pada kemungkinan bahwa sektor sekunder lebih mudah
terdorong keluar oleh kenaikan biaya lahan dan tenaga kerja, sedangkan
sektor tersier dapat lebih terkonsentrasi di pusat atau kota yang
memiliki fasilitas dan kewenangan lebih tinggi. (TXT baris 195-221 dan
242-250.)

\textbf{Inferensi pengolahan.} Masalah penelitian artikel terdiri atas
tiga pertanyaan empiris yang tidak identik: apakah HSR mengubah struktur
industri, apakah perubahan itu berbeda di sekitar pusat ekonomi atau
administratif tertentu, dan apakah perbedaan tersebut dapat disebut
sebagai \emph{borrowed size} atau bayangan administratif. Dua pertanyaan
pertama dioperasionalkan dengan DID dan DDD, sedangkan pertanyaan ketiga
bergantung pada proksi kedekatan kelompok. Karena itu, identifikasi
mekanisme lebih lemah daripada estimasi perubahan output yang
dilaporkan. (TXT baris 311-330.)

\subsection{Objectives}\label{objectives-24}

\textbf{Laporan sumber.} Tujuan penelitian yang dapat direkonstruksi
dari pendahuluan, hipotesis, dan rancangan metode adalah:

\begin{itemize}
\tightlist
\item
  menguji pengaruh pembangunan HSR terhadap peningkatan struktur
  industri di YRD;
\item
  menguraikan peningkatan tersebut ke dalam perubahan output sektor
  sekunder dan sektor tersier;
\item
  menguji apakah efek HSR berbeda menurut eksternalitas ukuran kota atau
  \emph{borrowed size};
\item
  menguji apakah efek HSR berbeda menurut kekuatan administratif atau
  \emph{borrowed administrative power};
\item
  membandingkan eksternalitas aglomerasi dan administratif pada wilayah
  yang berada di sekitar pusat ekonomi atau administratif;
\item
  menangani waktu pembukaan HSR yang bertahap melalui model multi-stage
  DID;
\item
  mengurangi potensi bias seleksi melalui propensity score matching
  (PSM); dan
\item
  menguji ketahanan hasil dengan uji tren paralel dan placebo.
\end{itemize}

\textbf{Laporan sumber.} Tujuan tersebut dirumuskan melalui tiga
hipotesis. Hipotesis 1 menyatakan bahwa peningkatan industri yang dipicu
HSR memiliki eksternalitas jaringan yang dapat tampil sebagai
\emph{borrowed size} atau \emph{agglomeration shadow}. Hipotesis 2
menyatakan bahwa efek HSR pada peningkatan industri memiliki
eksternalitas yang beragam, yang bersumber dari ukuran kota dan tingkat
administratif. Hipotesis 3 menyatakan bahwa efek HSR berbeda antar
sektor; sektor sekunder berpotensi mengalami efek \emph{borrowed size}
ekonomi, sedangkan sektor tersier dapat mengalami \emph{agglomeration
shadow} yang bersumber dari kekuatan administratif. (TXT baris 163-171
dan 182-250.)

\textbf{Inferensi pengolahan.} Tujuan artikel mencampurkan evaluasi
dampak kebijakan dengan pengujian teori eksternalitas. DID terutama
menjawab apakah terdapat perubahan yang berbeda antara wilayah HSR dan
non-HSR, sedangkan DDD menjawab apakah perbedaan itu lebih besar atau
lebih kecil di kelompok potensial. DDD tidak dengan sendirinya
memverifikasi bahwa kelompok tersebut benar-benar meminjam skala
penduduk atau kekuatan administratif. (TXT baris 311-330.)

\subsection{Literature Survey}\label{literature-survey-24}

\textbf{Laporan sumber.} Survei pustaka artikel dimulai dari penelitian
mengenai dampak HSR terhadap transformasi sosial-ekonomi. Dampak yang
disebut meliputi migrasi, investasi aset tetap, perpindahan manufaktur,
konsentrasi industri jasa, urbanisasi, pertumbuhan ekonomi, efek
limpahan, efek koridor, dan kompresi ruang-waktu. Teori lokasi
tradisional menjelaskan dampak tersebut melalui pengurangan waktu serta
biaya perjalanan dan peningkatan mobilitas faktor. New Economic
Geography menambahkan ekonomi aglomerasi, \emph{increasing returns to
scale}, dan efek \emph{crowding-out} dari konsentrasi berlebih. (TXT
baris 111-147.)

\textbf{Laporan sumber.} Artikel kemudian memindahkan pembahasan dari
struktur hierarkis menuju struktur jaringan. \emph{Borrowed size}, yang
diperkenalkan melalui Alonso dan dikembangkan dalam literatur regional
yang dirujuk artikel, menggambarkan ketidaksesuaian antara ukuran kota
dan kinerjanya dalam sistem regional. Kota yang lebih kecil dapat
memperoleh manfaat dari skala pasar dan ekonomi aglomerasi kota besar di
dekatnya. Sebaliknya, \emph{agglomeration shadow} menggambarkan kondisi
ketika kota besar membatasi pertumbuhan kota tetangga yang lebih kecil
sehingga fungsi dan kinerjanya lebih rendah daripada yang diharapkan.
(TXT baris 163-219.)

\textbf{Laporan sumber.} Dalam pembacaan jaringan, HSR dipandang dapat
mengubah sistem perkotaan hierarkis menjadi jaringan. Eksternalitas
jaringan dapat bersifat positif ketika kota pinggiran memperoleh akses
terhadap pasar, tenaga kerja, dan fungsi dari kota besar, tetapi dapat
pula bersifat negatif ketika kota dominan menyerap kegiatan dari wilayah
sekitarnya. Artikel memakai dualitas ini untuk membangun Hipotesis 1.
(TXT baris 196-219.)

\textbf{Laporan sumber.} Literatur ekonomi politik yang dirujuk artikel
memperluas analisis dengan menempatkan keputusan pemerintah, niat lokal,
dan negosiasi antartingkat pemerintahan sebagai bagian dari proses
pembangunan HSR. Pemerintah kota berpangkat tinggi disebut memiliki
pengaruh terhadap rute, lokasi stasiun, dan frekuensi kereta. Namun,
penelitian mengenai Pearl River Delta yang juga dirujuk artikel
menunjukkan bahwa hasil pembangunan infrastruktur dapat dipengaruhi oleh
bargaining regional dan tidak selalu mengikuti ukuran ekonomi semata.
(TXT baris 222-250.)

\textbf{Laporan sumber.} Artikel menggunakan contoh perbandingan Beijing
dan Shanghai untuk menunjukkan bahwa wilayah dengan pola ekonomi yang
mirip dapat mengalami lintasan berbeda setelah HSR. Beijing disebut
memiliki kapasitas intervensi administratif untuk menarik konsentrasi
sektor tersier, sedangkan Shanghai disebut memiliki kapasitas intervensi
yang berbeda. Kota sekitar Shanghai, seperti Suzhou, dipakai sebagai
contoh kota yang dapat menggunakan integrasi regional untuk menarik
infrastruktur dan investasi manufaktur. (TXT baris 225-235.)

\textbf{Laporan sumber.} Literatur yang dibahas juga mendukung
heterogenitas sektoral. Kenaikan aksesibilitas dan harga lahan dapat
mendorong manufaktur yang sensitif terhadap biaya untuk berpindah,
sedangkan sektor tersier yang lebih intensif modal pengetahuan dapat
terkonsentrasi di pusat. Karena intensitas aglomerasi dan
\emph{crowding-out} berbeda antar sektor, perubahan rasio industri tidak
harus berarti ekspansi output semua sektor. (TXT baris 201-221.)

\textbf{Interpretasi penulis.} Dari survei tersebut, penulis
menyimpulkan bahwa efek HSR perlu dibaca melalui kombinasi struktur
jaringan, ukuran ekonomi, dan kekuatan administratif. Mereka menganggap
kerangka \emph{borrowed size} dapat diperluas menjadi \emph{borrowed
administrative power} agar pengaruh politik dan tata kelola tidak
disamakan dengan eksternalitas pasar. (TXT baris 222-250 dan 735-740.)

\textbf{Inferensi pengolahan.} Survei ini adalah tinjauan naratif yang
digunakan untuk membangun hipotesis, bukan tinjauan sistematis. TXT
tidak menyebutkan strategi pencarian, kriteria inklusi, jumlah studi
yang disaring, rentang korpus, atau prosedur penilaian mutu. Temuan
karya yang dirujuk hanya diketahui melalui cara artikel A09
merangkumnya; karya-karya tersebut tidak dibuka atau diverifikasi secara
terpisah dalam review ini. (TXT baris 163-250 dan 795-884.)

\subsection{Method Used}\label{method-used-24}

\textbf{Laporan sumber.} Penelitian memakai desain kuantitatif dengan
panel 199 unit tingkat county di YRD selama 2000-2017. DID digunakan
karena membandingkan unit yang memiliki stasiun HSR dengan unit yang
tidak memiliki stasiun HSR sebelum dan sesudah pembukaan. Penulis
menyebut pendekatan tersebut sebagai quasi-experimental design dan
menyatakan bahwa HSR diperlakukan sebagai guncangan kebijakan yang
berlangsung bertahap. (TXT baris 149-161 dan 311-324.)

\textbf{Laporan sumber.} Pada DID dasar, \texttt{Treat\_i} adalah dummy
yang bernilai 1 jika kota atau county memiliki sedikitnya satu stasiun
HSR dan 0 jika tidak. \texttt{T\_t} dipakai untuk menunjukkan waktu
operasi HSR. Variabel efek utama dibentuk sebagai
\texttt{HSR\_it\ =\ Treat\_i\ x\ T\_t}. Persamaan yang ditampilkan
adalah
\texttt{gamma\_it\ =\ alpha\ +\ tau\ HSR\_it\ +\ beta1\ Treat\_i\ +\ beta2\ T\_t\ +\ beta3\ Cov\_it\ +\ mu\_i\ +\ delta\_t\ +\ epsilon\_it},
dengan efek tetap unit \texttt{mu\_i} dan efek tetap tahun
\texttt{delta\_t}. \texttt{Cov\_it} berisi kovariat sosial-ekonomi dan
administratif. (TXT baris 333-360.)

\textbf{Laporan sumber.} Karena stasiun HSR dibuka pada waktu yang
berbeda, penulis menerapkan multi-stage DID. \texttt{HSR1\_it}
menunjukkan tahun pertama HSR beroperasi di unit i, sedangkan
\texttt{HSR2\_it} sampai \texttt{HSRy\_it} menggunakan standar penetapan
yang sama untuk tahun-tahun berikutnya. Koefisien masing-masing tahap
dipakai untuk membaca efek bersih HSR pada perubahan industri menurut
waktu operasi. (TXT baris 362-375.)

\textbf{Laporan sumber.} DDD menambahkan interaksi
\texttt{HSR\_it\ x\ BOR\_i}. \texttt{BOR\_i} adalah dummy yang bernilai
1 jika unit i berada dalam kelompok potensial yang berdekatan dengan
pusat eksternalitas, dan 0 jika tidak. Model dipisahkan menjadi kelompok
\emph{borrowed size} dan kelompok \emph{borrowed administrative power}.
Koefisien interaksi yang positif dan signifikan ditafsirkan sebagai
\emph{borrowed size} atau \emph{borrowed administrative power},
sedangkan koefisien negatif dan signifikan ditafsirkan sebagai
\emph{agglomeration shadow} atau \emph{administrative shadow}. (TXT
baris 305-330.)

\textbf{Laporan sumber.} Kelompok \emph{borrowed size} dibentuk dari
county yang berdekatan dengan yurisdiksi Shanghai, Nanjing, Hangzhou,
Suzhou, dan Hefei. Lima kota tersebut dipilih karena disebut memiliki
populasi lebih dari 3 juta pada 2013. Kelompok \emph{borrowed
administrative power} dibentuk dari county yang berbatasan dengan
Shanghai, Hangzhou, Nanjing, dan Ningbo, yang disebut sebagai kota
berpangkat administratif vice-provincial atau lebih tinggi. Dengan
demikian, kelompok ekonomi memiliki lima pusat, sedangkan kelompok
administratif memiliki empat pusat. (TXT baris 271-296.)

\textbf{Laporan sumber.} Kovariat yang dicantumkan meliputi tingkat
administratif, ukuran penduduk, tingkat pembangunan ekonomi yang
diwakili PDB per kapita, investasi real estat, pengeluaran keuangan
publik, penyediaan layanan publik, dan tingkat urbanisasi. Artikel
menyatakan bahwa unit kota dikelompokkan dalam lima tingkatan hierarki,
dari subdistrik munisipal hingga county umum. Analisis statistik
dilakukan dengan Stata 15. (TXT baris 333-358 dan 381-401.)

\textbf{Laporan sumber.} Penulis menggunakan PSM untuk mengatasi potensi
endogenitas akibat pemilihan lokasi HSR dan perbedaan karakteristik
antara kelompok perlakuan dan kontrol. Pencocokan dilakukan
satu-lawan-satu dengan \emph{nearest neighbour matching} dan
penggantian. Sampel di luar dukungan (\emph{off-support}) dihapus.
Gambar 2 dan uraian penulis menunjukkan bahwa sebagian besar sampel
berada dalam dukungan serta bahwa perbedaan karakteristik berkurang
setelah pencocokan. (TXT baris 359-419.)

\textbf{Laporan sumber.} Penulis mengubah variabel dependen dan kovariat
ke bentuk logaritmik untuk mengurangi heteroskedastisitas. Mereka juga
menyatakan bahwa korelasi Pearson antara variabel dependen dan indikator
lain signifikan, serta bahwa seluruh nilai VIF kurang dari 5 untuk
menghindari multikolinearitas. (TXT baris 297-309.)

\textbf{Laporan sumber.} Ketahanan hasil diuji melalui tren paralel dan
placebo. Uji tren paralel memakai koefisien \texttt{HSR-B4} sampai
\texttt{HSR-B1} untuk perubahan sebelum pembukaan HSR. Placebo diulang
1.000 kali dengan tahun implementasi fiktif yang dipilih secara acak
karena pembukaan HSR merupakan peristiwa yang tidak dapat dibalik dan
berlangsung bertahap. (TXT baris 616-680.)

\textbf{Inferensi pengolahan.} Desain tersebut lebih tepat disebut
estimasi kuasi-eksperimental bersyarat daripada eksperimen. Efek tetap,
PSM, tren paralel, dan placebo memperkuat strategi identifikasi yang
dinyatakan penulis, tetapi TXT tidak menyediakan model propensity secara
rinci, jumlah observasi sebelum dan sesudah penghapusan
\emph{off-support}, definisi periode acuan untuk multi-stage DID, jenis
galat baku, pengelompokan galat baku, atau rincian bobot pencocokan.
Informasi yang tidak tersedia itu membatasi replikasi dan penilaian
penuh atas klaim kausal. (TXT baris 311-419 dan 626-680.)

\subsection{Dataset}\label{dataset-24}

\textbf{Laporan sumber.} Wilayah studi adalah YRD yang mencakup Shanghai
Municipality, Jiangsu, Zhejiang, dan Anhui. Artikel menyebut luas
wilayah 358.000 km2, GDP 2019 sebesar 23,72 triliun yuan, dan populasi
227 juta jiwa. Angka tersebut dilaporkan sebagai karakteristik regional,
bukan sebagai ukuran langsung sampel county 2000-2017. (TXT baris
138-148 dan 254-280.)

\textbf{Laporan sumber.} Data HSR, termasuk tahun pembukaan dan kota
yang dilayani, dikumpulkan dari situs resmi HSR China 12306 dan situs
statistik gaotie.cn. Data sosial-ekonomi berasal dari buku tahunan
statistik county dan kota yang tersedia melalui platform CNKI dan EPS.
Artikel menyatakan bahwa penyesuaian kesalahan yang berkaitan dengan
perubahan pembagian administratif juga dipertimbangkan. (TXT baris
292-309.)

\textbf{Laporan sumber.} Variabel yang tersedia dalam Tabel 1 adalah:

\begin{itemize}
\tightlist
\item
  \texttt{IndS}, yaitu rasio output sektor tersier terhadap sektor
  sekunder, dengan rerata 0,84, simpangan baku 0,35, minimum 0,36, dan
  maksimum 4,49;
\item
  \texttt{Ind2} (dicetak sebagai \texttt{Imd2} pada Tabel 1), yaitu
  output sektor sekunder dalam 100 juta RMB, dengan rerata 218,09,
  simpangan baku 580,86, minimum 0,52, dan maksimum 9.251,4;
\item
  \texttt{Ind3}, yaitu output sektor tersier dalam 100 juta RMB, dengan
  rerata 208,16, simpangan baku 884,67, minimum 0,85, dan maksimum
  20.783,47;
\item
  \texttt{Treat}, yaitu status kota atau county HSR, dengan rerata 0,41;
\item
  \texttt{T}, yaitu waktu atau status pembukaan HSR sebagaimana
  didefinisikan dalam artikel, dengan rerata yang dicetak 2008,5,
  simpangan baku 5,19, minimum 2000, dan maksimum 2017;
\item
  \texttt{HSR}, yaitu interaksi \texttt{Treat\ x\ T}, dengan rerata
  0,13;
\item
  \texttt{BS} dan \texttt{BAP}, yaitu dummy lokasi kelompok
  \emph{borrowed size} dan \emph{borrowed administrative power}, dengan
  rerata masing-masing 0,14 dan 0,07; dan
\item
  \texttt{Admin}, \texttt{Pop}, \texttt{pGDP}, \texttt{REI},
  \texttt{PFE}, \texttt{PS}, dan \texttt{UR}, yang masing-masing
  mewakili hierarki administratif, populasi, GDP per kapita, investasi
  real estat, pengeluaran keuangan publik, penyediaan layanan publik,
  dan tingkat urbanisasi.
\end{itemize}

(TXT baris 333-401.)

\textbf{Laporan sumber.} \texttt{Admin} diukur pada hierarki kota 1-5.
\texttt{Pop} diukur dalam 10.000 penduduk, \texttt{pGDP} dalam RMB,
\texttt{REI}, \texttt{PFE}, dan \texttt{PS} dalam miliar RMB, sedangkan
\texttt{UR} dijelaskan sebagai proporsi industri non-pertanian. Artikel
juga menyebut populasi, pembangunan ekonomi, kepadatan penduduk,
ketahanan ekonomi, investasi real estat, urbanisasi, pengeluaran publik,
dan layanan publik sebagai pertimbangan pemilihan kovariat. Namun, Tabel
1 secara eksplisit mencantumkan \texttt{Pop} sebagai ukuran kota, bukan
variabel terpisah bernama kepadatan penduduk. (TXT baris 341-355 dan
381-401.)

\textbf{Inferensi pengolahan.} TXT tidak menjelaskan daftar tahun buku
statistik yang dipakai untuk setiap variabel, prosedur pembersihan,
penanganan nilai hilang, deflator atau indeks harga, cara mengoreksi
perubahan batas administratif, dan protokol penggabungan data HSR dengan
data county. TXT juga tidak memuat kode analisis atau data replikasi.
Pernyataan ketersediaan data hanya berbunyi bahwa data akan tersedia
berdasarkan permintaan. (TXT baris 292-309 dan 787-790.)

\textbf{Inferensi pengolahan.} Ada ambiguitas internal pada \texttt{T}.
Uraian metode memperlakukannya sebagai dummy yang menunjukkan apakah HSR
telah beroperasi, sedangkan Tabel 1 memberi deskripsi
\texttt{1\ =\ post-HSR,\ 0\ =\ otherwise} tetapi statistiknya memiliki
rentang tahun 2000-2017 dan rerata 2008,5. Review ini tidak memilih
salah satu definisi tersebut karena TXT tidak menyelesaikan
perbedaannya. (TXT baris 333-350 dan 381-394.)

\subsection{Population / Sample}\label{population-sample-24}

\textbf{Laporan sumber.} Populasi analitis adalah 199 unit tingkat
county di YRD yang diamati selama 2000-2017. Penulis memilih skala
county, bukan prefecture, karena pembangunan HSR menggunakan kota dan
county sebagai unit dasar, serta pemerintah county China disebut
memiliki sistem fiskal dan kapasitas untuk merumuskan strategi ruang
serta bernegosiasi dengan pemerintah kota dan regional. (TXT baris
271-296.)

\textbf{Laporan sumber.} Periode 2000-2017 dipilih untuk menyediakan
tahun sebelum dan sesudah awal HSR pada 2008. Penulis sengaja tidak
memasukkan tahun setelah 2018 karena strategi integrasi YRD yang
diperkenalkan pemerintah pusat setelah 2018 dipandang sebagai guncangan
tambahan terhadap struktur industri dan dapat menjadi faktor pembaur.
(TXT baris 271-280.)

\textbf{Laporan sumber.} Lima pusat yang dipakai untuk kelompok
\emph{borrowed size} adalah Shanghai, Nanjing, Hangzhou, Suzhou, dan
Hefei. Empat pusat untuk kelompok \emph{borrowed administrative power}
adalah Shanghai, Hangzhou, Nanjing, dan Ningbo. Unit yang berdekatan
atau berbatasan dengan pusat-pusat tersebut dimasukkan ke kelompok
potensial masing-masing. TXT tidak memberikan daftar lengkap county
anggota atau jumlah absolut county untuk setiap kelompok. (TXT baris
281-296.)

\textbf{Laporan sumber.} Statistik deskriptif mencatat rerata
\texttt{Treat} sebesar 0,41, yang menunjukkan proporsi status HSR dalam
data menurut pengodean Tabel 1, serta rerata variabel \texttt{HSR}
sebesar 0,13. Rerata dummy \texttt{BS} adalah 0,14 dan rerata dummy
\texttt{BAP} adalah 0,07. PSM menghapus observasi \emph{off-support},
tetapi artikel tidak mencantumkan jumlah observasi yang tersisa setelah
proses tersebut. (TXT baris 381-419.)

\textbf{Inferensi pengolahan.} Artikel tidak menyebutkan N regresi
secara eksplisit dalam potongan Tabel 3-5 yang tersedia, jumlah unit
perlakuan dan kontrol secara absolut, pola data hilang, atau apakah
panel 199 unit seimbang. Karena itu, 199 unit adalah ukuran sampel yang
dinyatakan pada desain, bukan jaminan bahwa jumlah unit tersebut tetap
sama pada setiap estimasi. Generalisasi juga dibatasi oleh karakter YRD
sebagai wilayah maju, terurbanisasi, dan polisentris. (TXT baris 271-280
dan 381-419.)

\textbf{Inferensi pengolahan.} Pemilihan kelompok berdasarkan kedekatan
dengan pusat yang besar secara ekonomi atau tinggi secara administratif
mengidentifikasi wilayah yang \emph{potentially affected}, bukan sampel
yang terbukti memiliki relasi peminjaman. Kedekatan geografis juga dapat
berisi beberapa mekanisme sekaligus, seperti akses pasar, komuter,
kompetisi lahan, seleksi lokasi HSR, dan keputusan administratif. (TXT
baris 281-296 dan 311-330.)

\subsection{Dependent Variables}\label{dependent-variables-24}

\textbf{Laporan sumber.} Variabel dependen utama adalah transformasi
industri \texttt{gamma\_it} yang diwakili oleh tiga ukuran:

\begin{itemize}
\tightlist
\item
  \texttt{IndS}, yaitu rasio output sektor tersier terhadap output
  sektor sekunder;
\item
  \texttt{Ind2}, yaitu besaran output sektor sekunder; dan
\item
  \texttt{Ind3}, yaitu besaran output sektor tersier.
\end{itemize}

Ketiga ukuran tersebut bersumber dari buku tahunan statistik resmi.
\texttt{IndS} dipakai untuk membaca perubahan struktur industri,
sedangkan \texttt{Ind2} dan \texttt{Ind3} dipakai untuk mengidentifikasi
heterogenitas sektoral. (TXT baris 333-341 dan 381-387.)

\textbf{Laporan sumber.} Dalam Tabel 3, koefisien HSR untuk
\texttt{IndS} pada spesifikasi kolom 8 adalah 0,075, koefisien untuk
\texttt{Ind2} pada kolom 9 adalah -0,044, dan koefisien untuk
\texttt{Ind3} pada kolom 10 adalah 0,074. Artikel menyatakan bahwa
variabel dependen dan kovariat ditransformasikan ke bentuk logaritmik.
(TXT baris 297-309 dan 500-531.)

\textbf{Laporan sumber.} Dalam DDD, ketiga ukuran yang sama dipakai
sebagai hasil untuk dua jenis kelompok. Kolom 11-13 menganalisis
kelompok \emph{borrowed size}, sedangkan kolom 14-16 menganalisis
kelompok \emph{borrowed administrative power}. Tanda interaksi
\texttt{HSR\ x\ BOR} digunakan untuk memberi label \emph{borrowed size}
atau \emph{shadow} secara relatif terhadap kelompok pembanding. (TXT
baris 537-582.)

\textbf{Inferensi pengolahan.} Rasio output tersier-sekunder dapat
berubah karena kenaikan sektor tersier, penurunan sektor sekunder, atau
keduanya. Karena itu, peningkatan \texttt{IndS} tidak identik dengan
peningkatan kesejahteraan, produktivitas, pekerjaan, atau kualitas
struktur industri. Output sektor juga tidak menunjukkan distribusi
internal antar subindustri. Artikel sendiri mengakui bahwa dampak HSR
dapat berbeda di dalam subkategori sektor sekunder dan tersier. (TXT
baris 735-758.)

\textbf{Inferensi pengolahan.} Interpretasi persentase koefisien perlu
berhati-hati karena TXT tidak menjelaskan secara terpisah apakah
statistik deskriptif pada Tabel 1 ditampilkan sebelum atau sesudah log,
bagaimana rasio \texttt{IndS} ditransformasikan, serta apakah seluruh
komponen persamaan menggunakan transformasi yang sama. Artikel
menafsirkan koefisien HSR 0,075 sebagai kenaikan 7,5\%, tetapi rincian
transformasi lengkap tidak tersedia. (TXT baris 297-309, 381-401, dan
467-490.)

\subsection{Results}\label{results-24}

\textbf{Laporan sumber: perkembangan HSR.} Perkembangan HSR di YRD
dibagi secara naratif menjadi dua tahap. Pada 2008-2013, jaringan
terutama menghubungkan pusat regional seperti Shanghai, Nanjing, dan
Hangzhou. Sejak 2013, ekspansi lebih banyak menghubungkan kota-kota
sekunder dan membuka rute khusus, dengan Huangshan-Hangzhou Tourism
Railway sebagai contoh. Pada 2017, 17 jalur HSR dan 116 stasiun disebut
telah beroperasi di YRD. Tabel 2 melaporkan jumlah unit tingkat
prefecture dan county yang baru dilayani setiap tahun 2008-2017. (Bagian
4.1, Tabel 2, TXT baris 409-451.)

\textbf{Laporan sumber: pola spasial dan industri.} Gambar 3
menggambarkan pembangunan HSR dan perubahan struktur industri pada 2000,
2008, 2013, dan 2017. Narasi artikel menyebut wilayah yang dekat dengan
Hangzhou dan Nanjing secara konsisten menunjukkan peningkatan struktur
industri, sedangkan Hefei mengalami optimasi komposisi internal dan
beberapa wilayah periferal mengalami penurunan proporsi kontribusi
sektor tersier. (Bagian 4.1, TXT baris 421-466.)

\textbf{Laporan sumber: hasil DID utama.} Koefisien HSR dinyatakan
positif dan signifikan pada seluruh iterasi regresi untuk struktur
industri. Pada spesifikasi dengan kontrol lengkap, koefisien HSR pada
\texttt{IndS} adalah 0,075 dengan tanda signifikansi 1\%, yang
ditafsirkan penulis sebagai peningkatan struktur industri sebesar 7,5\%
di kota yang terhubung HSR. Nilai R2 spesifikasi tersebut adalah 0,7214,
dengan efek tetap county dan tahun. (Bagian 4.2, Tabel 3, TXT baris
467-531.)

\textbf{Laporan sumber: hasil sektoral.} Pada kolom output sektoral,
koefisien HSR adalah -0,044 dan signifikan untuk output sektor sekunder,
serta 0,074 dan signifikan untuk output sektor tersier. Dengan demikian,
peningkatan rasio struktur industri dalam model berlangsung bersamaan
dengan penurunan output sekunder dan kenaikan output tersier, bukan
melalui ekspansi sektor tersier saja. (Bagian 4.2, Tabel 3, TXT baris
491-531.)

\textbf{Interpretasi penulis: hasil DID.} Penulis menilai HSR memperkuat
aktivitas jasa, informasi, dan keuangan melalui konektivitas, tetapi
kenaikan harga lahan dan biaya tenaga kerja menekan manufaktur serta
sektor padat karya. Penulis menyatakan bahwa koefisien \texttt{Treat}
konsisten positif dan memakai pembacaan itu untuk menyatakan bahwa kota
HSR memiliki lintasan peningkatan industri yang lebih kuat daripada kota
non-HSR. Penulis juga membaca koefisien positif populasi sebagai
dukungan bagi ekonomi aglomerasi, sedangkan koefisien negatif GDP per
kapita dan tingkat administratif sebagai tanda bahwa kota yang lebih
maju atau berpangkat lebih tinggi dapat lebih kurang responsif terhadap
perubahan yang dipicu HSR karena memiliki jalur pembangunan dan proyek
lain. (TXT baris 460-490.)

\textbf{Inferensi pengolahan.} Pernyataan tentang \texttt{Treat} tidak
sepenuhnya selaras dengan Tabel 3: tabel mencetak koefisien
\texttt{Treat} negatif pada kolom 1-4, yaitu -0,09, -0,135, -0,061, dan
-0,027, lalu positif pada kolom 5-10. Karena itu, kata ``konsisten''
merupakan ketidakkonsistenan pelaporan internal yang perlu dipertahankan
sebagai kualifikasi; review ini tidak memilih narasi atau angka tabel
sebagai satu-satunya versi yang benar. (TXT baris 467-490 dan 500-531.)

\textbf{Interpretasi penulis: layanan publik sebagai kondisi.} Ketika
pengeluaran layanan publik dimasukkan sebagai kontrol, besaran dan
signifikansi koefisien HSR meningkat dari spesifikasi sebelumnya.
Penulis menafsirkan hal ini sebagai indikasi bahwa pengaruh HSR terhadap
peningkatan industri bergantung pada investasi layanan publik pemerintah
lokal dan tidak bersifat linear. (TXT baris 474-490.)

\textbf{Laporan sumber: DDD kelompok borrowed size.} Pada kelompok yang
berdekatan dengan pusat ekonomi besar, interaksi \texttt{HSR\ x\ BOR}
untuk \texttt{IndS} bernilai 0,033 dengan statistik t 0,62 dan tidak
bertanda signifikan. Interaksi untuk \texttt{Ind2} bernilai -0,082
dengan tanda \texttt{**} dan statistik t -2,34. Interaksi untuk
\texttt{Ind3} bernilai 0,034 dengan statistik t 0,74 dan tidak bertanda
signifikan. Selisih perlakuan-kontrol sebelum dan sesudah untuk
\texttt{IndS} berubah dari 0,133 menjadi 0,166; untuk \texttt{Ind2} dari
-0,129 menjadi -0,212; dan untuk \texttt{Ind3} dari 0,032 menjadi 0,065.
(Bagian 4.3, Tabel 4, TXT baris 537-614.)

\textbf{Interpretasi penulis: borrowed size.} Penulis menyatakan bahwa
kelompok kota kecil dan menengah menunjukkan tanda eksternalitas
aglomerasi atau \emph{borrowed size} dalam penyesuaian industri, tetapi
efek tersebut tidak menjadi bagian dari efek HSR pada struktur industri
secara keseluruhan karena interaksi pada \texttt{IndS} tidak signifikan.
Pada output sekunder, interaksi negatif dibaca sebagai penguatan efek
\emph{crowding-out} setelah HSR. Penulis menyatakan bahwa perubahan
tersebut dapat dipicu oleh kemakmuran dan aktivitas konsumen kota besar
yang meningkatkan biaya di kota sekitar. Output tersier dinyatakan tidak
terpengaruh secara signifikan oleh eksternalitas aglomerasi dalam model
ini. (TXT baris 585-614.)

\textbf{Laporan sumber: DDD kelompok borrowed administrative power.}
Pada kelompok yang berdekatan dengan pusat administratif berpangkat
tinggi, interaksi \texttt{HSR\ x\ BOR} untuk \texttt{IndS} bernilai
-0,09 dengan tanda \texttt{**} dan statistik t -1,65. Interaksi untuk
\texttt{Ind2} bernilai -0,002 dengan statistik t -0,04 dan tidak
bertanda signifikan. Interaksi untuk \texttt{Ind3} bernilai -0,14 dengan
tanda \texttt{***} dan statistik t -3,31. Selisih perlakuan-kontrol
sebelum dan sesudah untuk \texttt{IndS} berubah dari 0,083 menjadi
-0,007; untuk \texttt{Ind2} dari -0,139 menjadi -0,141; dan untuk
\texttt{Ind3} dari -0,004 menjadi -0,145. (Bagian 4.3, Tabel 4, TXT
baris 537-582.)

\textbf{Interpretasi penulis: borrowed administrative power.} Interaksi
negatif pada \texttt{IndS} dan terutama \texttt{Ind3} dibaca sebagai
\emph{administrative shadow}. Penulis menyatakan bahwa wilayah tingkat
lebih rendah di sekitar pusat administratif yang lebih tinggi mengalami
peningkatan industri yang lebih lemah setelah HSR, terutama pada sektor
tersier. Kekuatan administratif dianggap memengaruhi sektor tersier
melalui penyediaan layanan publik, perencanaan industri, dan subsidi,
sedangkan sektor sekunder lebih banyak dipengaruhi biaya lahan dan
pertimbangan ekonomi. (TXT baris 615-621 dan 699-722.)

\textbf{Inferensi pengolahan: makna DDD.} Hasil DDD tidak menunjukkan
bahwa semua kelompok \emph{borrowed size} menerima manfaat ekonomi
tambahan. Untuk \texttt{IndS}, efek diferensial kelompok ekonomi tidak
signifikan; untuk \texttt{Ind2}, efeknya negatif. Demikian pula, istilah
\emph{borrowed administrative power} menggambarkan nama kelompok
potensial, sementara koefisien yang signifikan terutama menunjukkan
bayangan administratif pada struktur dan sektor tersier. Pembacaan yang
paling hati-hati adalah bahwa kedekatan dengan pusat besar mengubah
distribusi relatif efek HSR antar sektor dan kelompok, bukan bahwa
proses peminjaman telah diamati langsung. (TXT baris 305-330 dan
537-614.)

\textbf{Laporan sumber: tren paralel.} Koefisien sebelum intervensi
\texttt{HSR-B4}, \texttt{HSR-B3}, \texttt{HSR-B2}, dan \texttt{HSR-B1}
tidak bertanda signifikan untuk struktur industri, sektor sekunder,
maupun sektor tersier. Penulis menggunakan hasil tersebut untuk
menyatakan bahwa tidak terdapat perbedaan tren awal yang dapat dibedakan
antara kelompok perlakuan dan kontrol. Tabel 5 juga memuat koefisien
\texttt{HSR-1st}, \texttt{HSR-2nd}, dan \texttt{HSR-3rd} untuk periode
operasi setelah pembukaan. (Bagian 5.1, Tabel 5, TXT baris 616-710.)

\textbf{Laporan sumber: placebo.} Placebo dilakukan 1.000 kali dengan
tahun implementasi fiktif yang dipilih secara acak. Dalam uraian Gambar
4, penulis menyatakan bahwa koefisien placebo terutama berada di sekitar
nol dan bahwa distribusi tersebut mendukung ketahanan hasil regresi.
(Bagian 5.2, Gambar 4, TXT baris 656-680.)

\textbf{Inferensi pengolahan: batas hasil robustness.} TXT tidak
menyediakan angka numerik lengkap dari seluruh pengulangan placebo atau
visual Gambar 4 selain deskripsi umum. Pernyataan bahwa sebagian besar
nilai placebo tidak signifikan tidak dapat diperiksa secara rinci dari
TXT saja. Selain itu, uji tren paralel yang ditampilkan adalah untuk
model DID utama; file tidak menjelaskan secara terpisah apakah asumsi
tren paralel juga diuji untuk setiap definisi kelompok DDD \texttt{BS}
dan \texttt{BAP}. (TXT baris 626-638 dan 656-680.)

\subsection{Findings}\label{findings-24}

\textbf{Temuan yang dilaporkan sumber.} HSR berkaitan dengan peningkatan
struktur industri di county YRD yang terhubung. Efek tersebut bersifat
sektoral: output sekunder berkurang, sedangkan output tersier meningkat.
Karena itu, peningkatan rasio tersier-sekunder merupakan gabungan dari
dua perubahan yang berlawanan arah. (TXT baris 467-498 dan 750-760.)

\textbf{Temuan yang dilaporkan sumber.} Efek HSR tidak sama di semua
wilayah. Di kelompok \emph{borrowed size}, interaksi struktur industri
tidak signifikan, interaksi output sekunder negatif dan signifikan,
serta interaksi output tersier tidak signifikan. Di kelompok
\emph{borrowed administrative power}, interaksi struktur industri dan
output tersier negatif serta bertanda signifikan, sedangkan interaksi
output sekunder tidak signifikan. (TXT baris 537-614.)

\textbf{Interpretasi penulis.} Penulis menganggap pola tersebut sebagai
eksternalitas yang beragam: ekonomi aglomerasi dan skala kota terutama
membentuk perpindahan atau pengurangan sektor sekunder, sedangkan
kekuatan administratif terutama membentuk konsentrasi atau bayangan
sektor tersier. Kota pusat dengan fasilitas dan kewenangan tinggi dapat
menarik fungsi tersier, sementara kota yang lebih rendah di sekitarnya
terhubung ke jaringan tetapi tidak selalu memperoleh pertumbuhan relatif
yang sama. (TXT baris 585-614 dan 699-722.)

\textbf{Temuan yang dilaporkan sumber.} Artikel juga melaporkan dua
gelombang peningkatan industri setelah HSR, yaitu pada tahun pertama dan
tahun ketiga operasi. Aktivitas sektor tersier disebut merespons pada
tahun pertama dan tahun berikutnya, sedangkan sektor sekunder baru
tampak terdampak pada tahun ketiga. (TXT baris 639-655.)

\textbf{Interpretasi penulis.} Perbedaan jeda waktu tersebut dijelaskan
melalui karakter aset dan biaya relokasi. Sektor tersier yang lebih
ringan aset dipandang lebih cepat merespons perubahan aktivitas
komersial, sedangkan sektor sekunder yang lebih padat modal memerlukan
waktu untuk menyesuaikan biaya dan kelayakan relokasi. (TXT baris
639-655.)

\textbf{Inferensi pengolahan.} Ada perbedaan penting antara efek
rata-rata dan efek diferensial. HSR secara rata-rata meningkatkan output
tersier sebesar 7,4\%, tetapi interaksi negatif pada kelompok
administratif berarti efek relatif kelompok tersebut lebih rendah
dibanding pembanding dalam spesifikasi DDD. Angka interaksi negatif
tidak boleh dibaca sebagai bukti bahwa output tersier absolut selalu
turun di setiap county yang berdekatan dengan pusat administratif.
Sebaliknya, efek rata-rata positif juga tidak membuktikan bahwa semua
county memperoleh pertumbuhan tersier yang sama.

\textbf{Inferensi pengolahan.} Klaim \emph{borrowed size} paling kuat
dalam artikel berada pada tingkat kerangka dan interpretasi, bukan pada
koefisien positif DDD yang konsisten. Koefisien utama yang signifikan
untuk kelompok ekonomi justru negatif pada output sekunder, sedangkan
koefisien positif pada struktur industri tidak signifikan. Dengan
demikian, bukti TXT lebih langsung mendukung heterogenitas dan
kemungkinan \emph{crowding-out} daripada pengukuran peminjaman skala
yang bersifat positif.

\textbf{Inferensi pengolahan.} Contoh Jiaxing dan Nanjing membantu
narasi mekanisme, tetapi dalam TXT keduanya disajikan sebagai ilustrasi
diskusi, bukan sebagai estimasi kasus terpisah. Jiaxing disebut menarik
komuter Shanghai melalui pengembangan di sekitar stasiun HSR, tetapi
kenaikan harga lahan menghambat relokasi manufaktur dari Shanghai.
Nanjing disebut memanfaatkan fasilitas pendidikan dan medis untuk
menarik tenaga kerja berkeahlian tinggi dan industri teknologi tinggi
tanpa secara signifikan menyedot manufaktur kota kecil dan menengah
sekitar. (TXT baris 680-722.)

\subsection{Conclusions}\label{conclusions-24}

\textbf{Kesimpulan yang dinyatakan penulis.} HSR sebagai infrastruktur
regional yang dipimpin pemerintah menghasilkan dua jenis konsekuensi
yang saling berkaitan, yaitu eksternalitas ekonomi berdasarkan ukuran
kota dan eksternalitas administratif berdasarkan kekuatan pemerintah
kota. Artikel menyatakan bahwa penggunaan infrastruktur sebagai lensa
perantara mempertemukan geografi ekonomi, geografi transportasi, dan
ekonomi politik perkotaan. (TXT baris 735-750.)

\textbf{Kesimpulan yang dinyatakan penulis.} HSR mendorong peningkatan
industri di kota-kota sepanjang jaringan. Pusat kota dapat mengalami
relokasi sektor sekunder ke county non-HSR, bukan selalu ke kota HSR
lain. Kota non-pusat dapat menarik sektor tersier, tetapi pada saat yang
sama menghadapi efek penyedotan dari pusat. Karena itu, peningkatan
industri pada kota periferal tidak menjamin bahwa total volume sektor
tersier akan meningkat; hasilnya bergantung pada kemampuan kota tersebut
menarik kegiatan dari wilayah non-HSR dan mengurangi penyedotan oleh
pusat. (TXT baris 750-760.)

\textbf{Kesimpulan yang dinyatakan penulis.} Wilayah di sekitar pusat
dengan kekuatan administratif lebih tinggi mengalami bayangan sektor
tersier yang lebih nyata, sedangkan dampak pada sektor sekunder lebih
lemah. Penulis mengaitkan perbedaan tersebut dengan layanan publik,
perencanaan, dan subsidi untuk sektor tersier, serta biaya lahan dan
kalkulasi biaya-manfaat untuk sektor sekunder. (TXT baris 699-722.)

\textbf{Kesimpulan yang dinyatakan penulis.} Dampak infrastruktur
transportasi di negara berkembang dengan intervensi negara yang kuat
tidak hanya ditentukan geografi ekonomi yang sudah ada. Kota dengan
peringkat administratif tinggi dapat memengaruhi perencanaan jaringan
transportasi walaupun tidak selalu menjadi kota terbesar secara ekonomi.
Penulis memberi contoh bahwa Suzhou, meskipun memiliki output ekonomi
terbesar kedua setelah Shanghai di YRD menurut uraian artikel, disebut
tidak memperoleh izin membangun bandara karena kedekatannya dengan
bandara Shanghai dan bandara Pukou Nanjing. (TXT baris 716-740.)

\textbf{Batas interpretasi.} Kesimpulan tersebut berlaku pada 199 unit
county di YRD selama 2000-2017 dan tidak secara otomatis berlaku untuk
wilayah pedalaman China, seluruh China, atau negara lain. Artikel
sendiri menyatakan bahwa struktur YRD yang maju, terurbanisasi, dan
polisentris membatasi penerapan langsung hasil ke wilayah lain. (TXT
baris 741-758.)

\textbf{Inferensi pengolahan.} Kesimpulan kausal artikel didukung oleh
desain DID bertahap, PSM, tren paralel, dan placebo sebagaimana
dilaporkan, tetapi kekuatan kausal tetap bergantung pada rincian seleksi
lokasi, definisi waktu, komposisi kelompok, dan spesifikasi yang tidak
sepenuhnya tersedia dalam TXT. Kesimpulan paling aman dari bukti yang
dapat diperiksa adalah adanya perubahan output yang berbeda menurut
sektor dan perbedaan efek relatif di sekitar pusat ekonomi serta pusat
administratif.

\subsection{Contributions}\label{contributions-24}

\textbf{Kontribusi yang diklaim sumber.} Artikel menyatakan atau
menunjukkan beberapa kontribusi berikut:

\begin{itemize}
\tightlist
\item
  memperluas analisis HSR dari aksesibilitas dan pertumbuhan ekonomi
  menuju hubungan antara infrastruktur, peningkatan industri, dan
  eksternalitas regional;
\item
  membawa konsep \emph{borrowed size} ke analisis tingkat county dalam
  YRD;
\item
  memperkenalkan \emph{borrowed administrative power} sebagai cara untuk
  memasukkan peran tingkat administratif dan intervensi pemerintah ke
  dalam analisis eksternalitas HSR;
\item
  membedakan dampak pada struktur industri, output sektor sekunder, dan
  output sektor tersier;
\item
  membandingkan pusat ekonomi berdasarkan ukuran penduduk dengan pusat
  administratif berdasarkan tingkat pemerintahan;
\item
  menggabungkan DID, multi-stage DID, PSM, dan DDD untuk memeriksa efek
  rata-rata, dinamika waktu, dan heterogenitas eksternalitas; dan
\item
  menjembatani literatur geografi transportasi dengan ekonomi politik
  perkotaan melalui infrastruktur transportasi sebagai mekanisme
  perantara.
\end{itemize}

(TXT baris 149-161, 311-419, 537-614, dan 735-740.)

\textbf{Kontribusi kebijakan yang diklaim sumber.} Artikel berupaya
memberi wawasan bagi wilayah yang sedang menggunakan transportasi
berkecepatan tinggi untuk menata struktur ekonomi, aktivitas industri,
dan struktur regional. Penulis menyebut relevansi potensial bagi kawasan
seperti Asia Tenggara, Afrika, dan India, tetapi sekaligus
memperingatkan bahwa perbedaan kondisi politik dan ekonomi membatasi
penerapan langsung. (TXT baris 723-740.)

\textbf{Inferensi pengolahan.} Nilai tambah analitis terkuat terletak
pada pemisahan antara efek DID rata-rata dan perbedaan DDD menurut
kedekatan dengan pusat ekonomi atau administratif. Namun, kontribusi
teoritis untuk membuktikan proses \emph{borrowed size} atau
\emph{borrowed administrative power} masih eksploratif karena konstruk
tersebut diproksikan melalui dummy kedekatan dan tidak diukur melalui
aliran, keputusan administratif, atau kapasitas layanan yang diamati.

\subsection{Research Gap}\label{research-gap-24}

\textbf{Kesenjangan yang dinyatakan sumber.} Artikel menyebut beberapa
kekurangan literatur sebelumnya:

\begin{itemize}
\tightlist
\item
  pendekatan yang berpusat pada kota inti kurang memperhatikan agensi
  dan interaksi kota kecil serta menengah di jaringan regional;
\item
  dampak jaringan antarkota periferal sering diremehkan dibandingkan
  efek aglomerasi yang berpusat pada kota besar;
\item
  hubungan endogen antara pembangunan HSR dan pertumbuhan atau
  peningkatan industri belum mudah dipisahkan;
\item
  penelitian yang membandingkan pengaruh ukuran kota dengan kekuatan
  administratif masih sedikit;
\item
  studi tentang HSR sering membahas hasil agregat dan belum cukup
  menguraikan perbedaan antara sektor sekunder dan tersier; dan
\item
  literatur sebelumnya belum cukup menjelaskan apakah wilayah sekitar
  pusat besar memperoleh manfaat \emph{borrowed size}, mengalami
  \emph{agglomeration shadow}, atau mengalami kombinasi keduanya.
\end{itemize}

(TXT baris 149-161, 171-219, dan 222-250.)

\textbf{Kesenjangan yang masih tersisa menurut isi artikel.} Artikel
mengakui perlunya penelitian yang membandingkan eksternalitas jaringan
antarwilayah dengan kondisi politik dan ekonomi yang berbeda. Artikel
juga menyebut perlunya metode kualitatif untuk memeriksa mekanisme
eksternalitas administratif. Keterbatasan eksplisit mengenai subkategori
industri, frekuensi kereta, dan penempatan stasiun menunjukkan bahwa
rincian jaringan dan sektoral juga masih menjadi celah. (TXT baris
741-758.)

\textbf{Inferensi pengolahan.} Setelah studi ini, masih belum terjawab
apakah county yang diklasifikasikan sebagai kelompok potensial
benar-benar menerima skala penduduk, pasar, investasi, atau kapasitas
administratif dari pusat tertentu. Belum tersedia pengukuran langsung
tentang asal-tujuan komuter, relokasi perusahaan, penggunaan layanan,
bargaining antarpemerintah, frekuensi kereta, kualitas stasiun, atau
aliran sumber daya. Kesenjangan ini berasal dari operasionalisasi dan
informasi yang tersedia dalam TXT, bukan dari penambahan sumber luar.

\subsection{Limitations}\label{limitations-24}

\textbf{Keterbatasan yang dinyatakan penulis.} Artikel secara eksplisit
mengakui bahwa:

\begin{itemize}
\tightlist
\item
  analisis hanya membahas sektor sekunder dan tersier, padahal efek HSR
  dapat sangat berbeda di dalam subkategori masing-masing;
\item
  frekuensi layanan kereta dan penempatan strategis stasiun dapat
  menghasilkan heterogenitas efek HSR, tetapi faktor tersebut tidak
  dimasukkan dalam analisis utama; dan
\item
  wilayah studi hanya YRD yang sangat maju. Karakter urbanisasi maju dan
  struktur polisentris YRD membuat kesimpulan tidak dapat langsung
  diterapkan pada wilayah pedalaman China atau negara lain di Global
  South.
\end{itemize}

Penulis mengusulkan perbandingan antarwilayah dengan kondisi politik dan
ekonomi berbeda serta metode kualitatif untuk mempelajari mekanisme
eksternalitas administratif. (TXT baris 741-758.)

\textbf{Keterbatasan berdasarkan informasi yang dilaporkan sumber.}
Kelompok \texttt{BS} dan \texttt{BAP} bukan ukuran langsung peminjaman.
\texttt{BS} dibentuk dari kedekatan dengan lima kota berpopulasi lebih
dari 3 juta, sedangkan \texttt{BAP} dibentuk dari kedekatan dengan empat
kota berpangkat administratif tinggi. Definisi ini tidak memperlihatkan
apakah penduduk, pasar, fungsi, layanan, kewenangan, atau investasi
benar-benar bergerak melintasi batas county. (TXT baris 271-296.)

\textbf{Inferensi pengolahan.} Klaim kausal menghadapi keterbatasan
transparansi. Artikel menyebut endogenitas pemilihan lokasi HSR, PSM,
efek tetap, tren paralel, dan placebo, tetapi TXT tidak mencantumkan
model propensity, seluruh kovariat yang dipakai untuk matching secara
eksplisit, ukuran sampel setelah matching, statistik balance, jenis
galat baku, pengelompokan galat baku, atau strategi untuk ketergantungan
spasial. (TXT baris 149-161 dan 359-419.)

\textbf{Inferensi pengolahan.} Pembukaan HSR berlangsung bertahap,
tetapi TXT tidak menjelaskan kategori waktu yang dihilangkan sebagai
referensi, batas jendela event-time, atau bagaimana seluruh efek tahap
dibandingkan. Tabel 1 juga tidak konsisten sepenuhnya dengan uraian
\texttt{T}: deskripsi dummy pasca-HSR berdampingan dengan statistik yang
tampak seperti tahun pembukaan. Perbedaan ini dapat memengaruhi
interpretasi koefisien jika tidak diselesaikan dalam data atau kode
asli. (TXT baris 362-375 dan 381-394.)

\textbf{Inferensi pengolahan.} Variabel \texttt{IndS} dan output
sektoral adalah ukuran berbasis output, bukan ukuran produktivitas,
pekerjaan, kualitas layanan, nilai tambah per pekerja, komposisi
subindustri, atau kesejahteraan. Rasio \texttt{IndS} juga dapat
meningkat karena sektor sekunder turun, sehingga label ``upgrading''
perlu dibaca sesuai definisi operasional artikel. (TXT baris 333-341 dan
735-758.)

\textbf{Inferensi pengolahan.} Efek publikasi layanan publik tidak diuji
sebagai mekanisme mediasi. Fakta bahwa koefisien HSR berubah ketika
\texttt{PS} dimasukkan menunjukkan sensitivitas spesifikasi atau
kemungkinan kondisi pendukung, tetapi tidak membuktikan urutan kausal
HSR lalu layanan publik lalu peningkatan industri. (TXT baris 474-490.)

\textbf{Inferensi pengolahan.} Data dan proses konstruksi variabel belum
cukup terdokumentasi untuk direplikasi dari TXT saja. Tidak disebutkan
daftar county, tanggal dan prosedur pengumpulan data, penanganan nilai
hilang, pemrosesan perubahan batas administratif secara rinci, aturan
deflasi, kode analisis, atau data replikasi. Data availability hanya
dinyatakan tersedia berdasarkan permintaan. (TXT baris 292-309 dan
787-790.)

\textbf{Inferensi pengolahan.} Bukti tren paralel dan placebo memberi
dukungan pada robustness model utama menurut penulis, tetapi TXT hanya
menyajikan ringkasan dan bukan keluaran numerik lengkap placebo. File
juga tidak menjelaskan apakah hasil tersebut tahan terhadap definisi
pusat, radius atau batas kedekatan alternatif, pengukuran kualitas
jaringan, dan spesifikasi sektor yang lebih rinci. (TXT baris 626-680.)

\textbf{Inferensi pengolahan.} Ada dua unit klasifikasi yang berbeda,
yaitu lima pusat ekonomi dan empat pusat administratif. Perbedaan
tersebut substantif dan sesuai tujuan perbandingan, tetapi artikel tidak
menyediakan jumlah anggota, tumpang tindih anggota, atau hasil
sensitivitas ketika pusat yang tumpang tindih diperlakukan secara
berbeda. (TXT baris 281-296.)

\subsection{Challenges}\label{challenges-24}

\textbf{Tantangan yang dilaporkan sumber.} Artikel menghadapi persoalan
pemilihan lokasi HSR yang tidak acak. Wilayah yang berpotensi tumbuh
lebih cepat dapat lebih mungkin memperoleh HSR, sehingga penulis memakai
PSM, efek tetap, DID bertahap, tren paralel, dan placebo untuk
mengurangi persoalan tersebut. (TXT baris 149-161 dan 359-419.)

\textbf{Tantangan yang dilaporkan sumber.} Pembangunan jaringan
berlangsung berurutan dari pusat regional ke kota sekunder dan rute
khusus. Kondisi ini menyulitkan pemisahan efek waktu pembukaan, efek
jaringan yang terus berkembang, dan perubahan sosial-ekonomi lain
sepanjang 2000-2017. (TXT baris 409-451 dan 639-680.)

\textbf{Tantangan yang dilaporkan sumber.} Pemerintah dan tingkat
administratif berperan dalam rute, lokasi stasiun, layanan kereta, serta
pembangunan regional. Karena itu, analisis HSR tidak cukup jika hanya
menggunakan ukuran pasar atau ukuran penduduk. Artikel berusaha
menangkap persoalan tersebut melalui kelompok \texttt{BAP} dan model
DDD. (TXT baris 222-250 dan 281-330.)

\textbf{Tantangan yang dilaporkan sumber.} Efek HSR berbeda antar sektor
dan dapat terjadi dengan jeda waktu yang berbeda. Sektor tersier
merespons lebih cepat, sementara sektor sekunder lebih lambat karena
biaya relokasi dan karakter aset. Hal ini menuntut analisis sektoral dan
tidak cukup dijawab hanya dengan satu indeks struktur industri. (TXT
baris 639-655.)

\textbf{Inferensi pengolahan.} Tantangan utama interpretasi adalah
memisahkan setidaknya lima mekanisme yang dapat menghasilkan pola
serupa: ekonomi lokal, efek aglomerasi pusat, konektivitas HSR, kenaikan
biaya lahan, dan keputusan administratif. Dummy kedekatan \texttt{BS}
atau \texttt{BAP} menggabungkan mekanisme tersebut, sehingga label
eksternalitas sebaiknya dipahami sebagai kategori analitis awal.

\textbf{Inferensi pengolahan.} Tantangan penerapan juga bersifat
distribusional. Koneksi HSR dapat menaikkan output tersier secara
agregat sekaligus memperlemah posisi relatif county yang berada dekat
pusat administratif atau ekonomi. Kebijakan yang hanya mengejar
konektivitas tanpa memeriksa sektor, harga lahan, layanan publik, dan
distribusi kewenangan berpotensi membaca manfaat agregat sebagai manfaat
yang merata. Inferensi ini diturunkan dari perbedaan hasil DID dan DDD,
bukan dari evaluasi kebijakan tambahan. (TXT baris 537-614 dan 750-760.)

\subsection{Practical Implications}\label{practical-implications-24}

\textbf{Implikasi yang dinyatakan atau dibahas penulis.} Pembangunan HSR
tidak menjamin peningkatan ekonomi lokal yang seragam. Kota periferal
perlu menyiapkan strategi pembangunan sebelum terintegrasi ke jaringan.
Wilayah yang berfokus pada sektor tersier, transisi dari manufaktur ke
jasa, atau pengembangan real estat dapat memperoleh manfaat dari HSR,
sedangkan wilayah yang ingin menghidupkan kembali sektor sekunder
mungkin tidak memperoleh manfaat maksimum dari HSR penumpang. (TXT baris
723-733.)

\textbf{Implikasi yang dinyatakan atau dibahas penulis.} Pemerintah
perlu memperhatikan bahwa konektivitas dapat membawa kenaikan harga
lahan dan biaya tenaga kerja, yang mendorong manufaktur keluar dari
pusat atau county yang berdekatan. Pada saat yang sama, sektor tersier
dapat terkonsentrasi di kota berpangkat lebih tinggi. Karena itu,
perencanaan HSR perlu dibaca bersama strategi lahan, investasi layanan
publik, dan kemampuan kota periferal untuk menarik fungsi dari wilayah
non-HSR. (TXT baris 460-466, 585-614, dan 750-760.)

\textbf{Implikasi yang dinyatakan atau dibahas penulis.} Peringkat
administratif perlu dipertimbangkan dalam tata kelola jaringan
transportasi. Kota yang memiliki kewenangan lebih tinggi dapat
memengaruhi hasil pembangunan jaringan walaupun tidak selalu merupakan
kota terbesar secara ekonomi. Perspektif tersebut dipakai artikel untuk
mendorong analisis pembangunan transportasi yang memasukkan politik
skala dan hubungan antarpemerintah. (TXT baris 716-740.)

\textbf{Inferensi pengolahan.} Hasil A09 dapat digunakan sebagai
peringatan agar indikator keberhasilan HSR tidak hanya berupa rasio
sektor tersier-sekunder atau output regional. Evaluasi praktis
seharusnya memeriksa output absolut kedua sektor, pekerjaan, harga
lahan, frekuensi kereta, lokasi stasiun, penggunaan layanan, relokasi
perusahaan, dan distribusi antar county. Daftar ini adalah inferensi
dari batas pengukuran TXT, bukan rekomendasi empiris yang diuji oleh
artikel.

\subsection{Applications}\label{applications-24}

\textbf{Aplikasi yang didukung langsung oleh sumber.} Kerangka artikel
dapat digunakan untuk:

\begin{itemize}
\tightlist
\item
  mengevaluasi perubahan struktur industri di tingkat county setelah
  pembukaan HSR;
\item
  memisahkan perubahan rasio struktur industri dari perubahan output
  sektor sekunder dan tersier;
\item
  membandingkan efek HSR pada wilayah yang berdekatan dengan pusat
  ekonomi besar dan wilayah pembanding;
\item
  membandingkan efek HSR pada wilayah yang berdekatan dengan pusat
  administratif berpangkat tinggi dan wilayah pembanding;
\item
  memeriksa dinamika tahun pertama, tahun kedua, dan tahun ketiga
  setelah HSR melalui multi-stage DID; dan
\item
  menyediakan bahan bagi perencanaan regional yang memperhitungkan
  ekonomi aglomerasi serta kekuatan administratif.
\end{itemize}

(TXT baris 311-419, 500-582, dan 639-680.)

\textbf{Aplikasi yang diklaim penulis.} Temuan artikel dimaksudkan
memberi wawasan bagi wilayah yang menggunakan transportasi berkecepatan
tinggi untuk menata pembangunan ekonomi, aktivitas industri, dan
struktur regional. Penulis menyebut kemungkinan relevansi bagi Asia
Tenggara, Afrika, dan India, tetapi tidak melaporkan penerapan empiris
di wilayah-wilayah tersebut. (TXT baris 723-740.)

\textbf{Inferensi pengolahan.} Di luar YRD, alur DID-PSM-DDD hanya layak
dipakai sebagai kerangka diagnosis awal setelah definisi unit, kualitas
data, pembukaan stasiun, pusat ekonomi, pusat administratif, dan
struktur politik disesuaikan. Transfer langsung bobot, ambang populasi 3
juta, atau definisi kedekatan tidak didukung oleh TXT.

\textbf{Informasi tidak tersedia.} File tidak melaporkan adanya
perangkat lunak yang dibagikan, dashboard, instrumen operasional,
program kebijakan yang sudah menerapkan hasil, atau evaluasi
pascapenerapan. Pernyataan yang tersedia hanya bahwa data akan diberikan
berdasarkan permintaan. (TXT baris 787-790.)

\subsection{Future Research
(Optional)}\label{future-research-optional-24}

\textbf{Agenda yang dinyatakan penulis.} Penulis mengusulkan penelitian
lanjutan untuk:

\begin{itemize}
\tightlist
\item
  membandingkan eksternalitas jaringan antarwilayah dengan kondisi
  politik dan ekonomi yang berbeda;
\item
  memakai metode kualitatif untuk mengeksplorasi mekanisme eksternalitas
  administratif dalam pembangunan transportasi; dan
\item
  menguji implikasi HSR di luar YRD dengan memperhatikan perbedaan
  tingkat pembangunan, struktur politik, dan ekonomi.
\end{itemize}

(TXT baris 741-758.)

\textbf{Arah lanjutan yang tersirat dari keterbatasan sumber.} Penulis
menyebut heterogenitas subkategori sektor, frekuensi layanan kereta, dan
penempatan strategis stasiun sebagai keterbatasan. Memasukkan
faktor-faktor tersebut dalam studi berikutnya merupakan arah yang masuk
akal dari isi artikel, tetapi bukan agenda penelitian yang dirumuskan
secara eksplisit dalam TXT. (TXT baris 741-758.)

\textbf{Inferensi pengolahan.} Berdasarkan keterbatasan
operasionalisasi, penelitian berikutnya juga perlu mengukur arus komuter
dan pengguna layanan, relokasi perusahaan, investasi, aliran
pengetahuan, kapasitas layanan publik, serta proses negosiasi
administratif secara langsung. Desain panel yang lebih terperinci,
variasi kualitas jaringan, dan pengujian sensitivitas kelompok kedekatan
dapat membantu membedakan efek pusat ekonomi, jaringan HSR, dan kekuatan
administratif. Usulan ini adalah inferensi review, bukan agenda yang
dinyatakan secara eksplisit dalam artikel.

\subsection{Provenance Notes}\label{provenance-notes-24}

\begin{itemize}
\tightlist
\item
  \textbf{Basis akses:} review ini hanya menggunakan file TXT
  \texttt{source/journal/A09-shi-wang-zhang-yu-2025-borrowed-size-and-borrowed-administrative-power-effects-of-high-speed-rail-network-on-industrial-upgrading-and-variegated-externalities-in-the-yangtze-river-delta-china-10.1016-j.jtrangeo.2025.104113.txt}.
  Tidak ada PDF, sumber daring, catatan Zotero, atau karya dalam daftar
  pustaka yang dibuka untuk menambah atau memverifikasi substansi.
\item
  \textbf{Cakupan baca:} seluruh 890 baris TXT dibaca, termasuk blok
  metadata PDF, teks artikel, abstrak, persamaan, Tabel 1-5, keterangan
  Gambar 1-4, diskusi, kesimpulan, CRediT authorship contribution
  statement, declaration of competing interest, acknowledgment, data
  availability, dan daftar pustaka.
\item
  \textbf{Asal teks:} blok metadata menyatakan bahwa TXT diekstrak pada
  31 Agustus 2026 dari PDF 13 halaman dengan
  \texttt{pdftotext\ -layout}. Metadata menyatakan PDF tidak terenkripsi
  dan tidak bertag. (TXT baris 1-28.)
\item
  \textbf{Dasar identitas bibliografis:} identitas penulis, judul,
  jurnal, DOI, tanggal editorial, dan afiliasi diambil dari kepala
  artikel serta bagian \texttt{ARTICLE\ INFO}, bukan dari metadata
  \texttt{pdfinfo} saja. Metadata PDF hanya mencantumkan Dehao Shi
  sebagai author, sedangkan kepala artikel mencantumkan empat penulis.
  (TXT baris 7-13 dan 42-107.)
\item
  \textbf{Locator:} karena TXT tidak menyediakan nomor halaman jurnal
  terpisah untuk setiap bagian, locator utama dalam review ini
  menggunakan nama bagian, nomor tabel atau gambar, dan rentang baris
  TXT. Pemisah halaman PDF tetap muncul sebagai karakter form-feed dan
  nomor halaman internal pada TXT.
\item
  \textbf{Tabel dan gambar:} Tabel 1-5 terbaca sebagai teks dan angka
  yang dilaporkan di review diambil dari baris tabel yang tersedia.
  Gambar 1-4 terutama tersedia melalui keterangan dan uraian naratif;
  review tidak memperkirakan nilai visual yang tidak tercetak dalam TXT.
  (TXT baris 268, 381-451, 537-582, dan 682-710.)
\item
  \textbf{Pemisahan epistemik:} bagian berlabel \textbf{Laporan sumber}
  memparafrasekan isi artikel; \textbf{Interpretasi penulis} merekam
  penjelasan atau makna kausal yang diberikan penulis; \textbf{Inferensi
  pengolahan} adalah evaluasi berbasis isi dan ketidaklengkapan TXT,
  bukan klaim artikel atau informasi dari luar.
\item
  \textbf{Informasi tidak tersedia:} TXT tidak memberikan daftar lengkap
  199 county, jumlah absolut unit perlakuan dan kontrol, jumlah sampel
  setelah PSM, rincian propensity score dan balance test, penanganan
  data hilang, prosedur koreksi batas administratif, rincian deflasi,
  definisi lengkap variabel waktu \texttt{T}, jenis galat baku dan
  pengelompokannya, pemeriksaan dependensi spasial, frekuensi kereta,
  kualitas stasiun, daftar subindustri, kode, atau data replikasi.
\item
  \textbf{Ambiguitas internal yang dipertahankan:} definisi dan
  statistik \texttt{T} tidak sepenuhnya selaras; jumlah pusat kelompok
  ekonomi dan administratif berbeda; serta label output sekunder pada
  Tabel 1 tampak sebagai \texttt{Imd2} sementara tabel hasil menggunakan
  \texttt{Ind2}. Review tidak memperbaiki atau menebak maksud asli dari
  perbedaan tersebut. (TXT baris 333-350 dan 381-401.)
\item
  \textbf{Status publikasi:} artikel mencantumkan tanggal diterima dan
  hak cipta Elsevier, tetapi status penelaahan sejawat tidak
  diverifikasi dari luar TXT dan karena itu dinyatakan tidak disebutkan
  dalam file.
\end{itemize}

\section{Review A10}\label{review-a10}

\subsection{Identitas Sumber}\label{identitas-sumber-11}

\begin{itemize}
\tightlist
\item
  Judul: \emph{Knowledge innovation network externalities in the
  Guangdong-Hong-Kong-Macao Greater Bay Area: borrowing size or
  agglomeration shadow?} {[}TXT lines 40-49, 85-87{]}.
\item
  Penulis: Wenyi Yang, Fei Fan, Xueli Wang, dan Haichao Yu {[}TXT lines
  44, 88-93{]}.
\item
  Jurnal: \emph{Technology Analysis \& Strategic Management} {[}TXT
  lines 30, 79-87{]}.
\item
  Tahun yang tercantum dalam sitasi: 2021; artikel dinyatakan terbit
  daring pada 15 Juni 2021 {[}TXT lines 46-56{]}. Nama file TXT memuat
  \texttt{2022}, tetapi tahun pada sitasi dan tanggal publikasi di dalam
  TXT adalah 2021.
\item
  DOI: \texttt{10.1080/09537325.2021.1940922}; tautan DOI juga
  dicantumkan dalam TXT {[}TXT lines 48-51, 80{]}.
\item
  Afiliasi: School of Public Administration, Zhongnan University of
  Economics and Law, serta Institute of Regional and Urban-Rural
  Development, Wuhan University {[}TXT lines 88-93{]}.
\item
  Volume, nomor terbitan, dan halaman artikel pada informasi
  bibliografis yang terbaca: Tidak disebutkan dalam file. Metadata PDF
  hanya menyebutkan 19 halaman {[}TXT lines 22-24{]}.
\item
  Status peer-review: Tidak disebutkan dalam file.
\end{itemize}

\subsection{Summarized Abstract}\label{summarized-abstract-25}

\textbf{Laporan sumber:} Artikel menganalisis pola spasial dan
eksternalitas jaringan inovasi pengetahuan di Guangdong-Hong Kong-Macao
Greater Bay Area (GBA) berdasarkan data 2001-2019. Metodenya
menggabungkan analisis jaringan sosial dan model ekonometrika spasial.
Guangzhou dan Hong Kong disebut selalu menjadi inti jaringan; Shenzhen
muncul sebagai pusat inovasi setelah 2012; kota-kota lain cenderung
berada di posisi periferal. Kota kecil dan menengah disebut tidak
memperoleh manfaat dari perkembangan inovatif kota inti, melainkan
terperangkap dalam \emph{agglomeration shadow}. Perbedaan institusional
dan budaya disebut sebagai hambatan utama kerja sama inovasi, sedangkan
jarak memiliki keterbatasan yang lebih kecil. Penulis menyimpulkan bahwa
eksternalitas jaringan inovasi di wilayah ini bersifat negatif dan bahwa
kesenjangan spasial perlu dipersempit {[}TXT lines 96-114{]}.

\textbf{Interpretasi penulis:} Jaringan inovasi GBA belum menghasilkan
\emph{borrowing size} bagi kota kecil dan menengah. Posisi pusat yang
kuat justru berkaitan dengan dampak negatif terhadap kota sekitar,
terutama ketika terdapat perbedaan institusional {[}TXT lines 501-509,
662-665{]}.

\textbf{Inferensi pengolahan:} Pertanyaan utama artikel bukan hanya
apakah konektivitas antarkota meningkat, melainkan apakah pola
konektivitas tersebut menyebarkan kapasitas inovasi atau memperkuat
konsentrasi di kota inti. Bukti yang tersedia dalam TXT mendukung
kesimpulan penulis pada konteks kota-kota GBA dan periode 2001-2019;
generalisasi ke wilayah, periode, atau ukuran inovasi lain tidak
dinyatakan dalam file.

\subsection{Problem Statement}\label{problem-statement-25}

\textbf{Laporan sumber:} Perubahan dari \emph{space of place} ke
\emph{space of flows} membuat kota dipahami sebagai simpul jaringan,
bukan hanya titik yang berdiri sendiri. Interaksi antarkota dipandang
penting bagi kinerja inovasi {[}TXT lines 119-125{]}. GBA diposisikan
sebagai wilayah yang terbuka dan vital secara ekonomi. Pemerintah China
melalui GBA Planning Outline 2019 menargetkan wilayah tersebut sebagai
aglomerasi perkotaan kelas dunia dan kawasan inovasi internasional,
dengan Guangzhou, Shenzhen, Hong Kong, dan Macao sebagai empat pusat
inovasi {[}TXT lines 126-130{]}.

Masalah empirisnya adalah belum jelas apakah jaringan inovasi GBA telah
membentuk pola multisentra dan eksternalitas seperti apa yang
dihasilkannya. Penulis menanyakan apakah kota di sekitar pusat inovasi
memperoleh \emph{borrowing size} atau justru mengalami
\emph{agglomeration shadow} {[}TXT lines 131-148{]}. Kompleksitas dua
sistem politik dan tiga wilayah pabean juga membuat kerja sama inovasi
lebih sulit {[}TXT lines 286-298, 304-309{]}.

\textbf{Interpretasi penulis:} Studi tentang jaringan perkotaan GBA
sebelumnya lebih banyak membahas arus faktor ekonomi, jalur perkembangan
inovasi, struktur spasial jaringan, kinerja inovasi, dan faktor yang
memengaruhinya. Menurut penulis, studi-studi tersebut belum menjelaskan
pengaruh eksternal struktur jaringan terhadap inovasi regional {[}TXT
lines 131-145, 292-298{]}.

\subsection{Objectives}\label{objectives-25}

\textbf{Laporan sumber:} Tujuan yang dinyatakan adalah mengukur pola
spasial dan eksternalitas jaringan inovasi GBA, menjelaskan pengaruh
kota besar terhadap inovasi regional, dan menyediakan rujukan untuk
perbaikan jaringan {[}TXT lines 147-155{]}. Secara operasional, artikel
juga bertujuan untuk:

\begin{itemize}
\tightlist
\item
  Mengidentifikasi posisi dan perkembangan pusat-pusat jaringan inovasi
  pengetahuan.
\item
  Menguji pengaruh sentralitas jaringan terhadap inovasi keseluruhan,
  inovasi kerja sama, dan inovasi independen kota.
\item
  Menguji apakah perbedaan sentralitas berkaitan dengan kesenjangan
  inovasi antarkota.
\item
  Membedakan pengaruh jarak geografis, perbedaan budaya, dan perbedaan
  institusional terhadap eksternalitas jaringan {[}TXT lines 371-417{]}.
\item
  Menilai apakah jaringan yang ada menghasilkan eksternalitas positif
  berupa \emph{borrowing size} atau eksternalitas negatif berupa
  \emph{agglomeration shadow} {[}TXT lines 221-225, 286-298{]}.
\end{itemize}

\textbf{Catatan:} Tidak disebutkan dalam file adanya hipotesis formal
bernomor atau target prediksi kuantitatif sebelum estimasi.

\subsection{Literature Survey}\label{literature-survey-25}

\textbf{Laporan sumber:} Tinjauan pustaka dibangun dari beberapa
rangkaian gagasan. Pertama, teori \emph{backwash-spread}, \emph{trickle
down}, dan \emph{core-periphery} menjelaskan kemungkinan perkembangan
yang tidak seimbang antara pusat dan perifer. Kota besar dengan kondisi
inovasi yang lebih baik dapat menjadi inti, sedangkan wilayah sekitarnya
mengambil posisi subordinat {[}TXT lines 158-170{]}.

Kedua, konsep \emph{borrowed size} menjelaskan bahwa kota kecil yang
dekat dengan metropolitan dapat menikmati sebagian keuntungan kota besar
sambil mempertahankan keunggulan lokalnya, misalnya sewa lebih rendah
dan kemacetan lebih kecil. Sebaliknya, \emph{agglomeration shadow}
menggambarkan keadaan ketika kota kecil tidak memperoleh manfaat kota
besar di dekatnya dan justru dirugikan {[}TXT lines 171-177, 181-189{]}.

Ketiga, perspektif eksternalitas jaringan dan \emph{space of flows}
digunakan untuk melampaui penjelasan yang hanya bertumpu pada kedekatan
geografis. Capello dirujuk untuk gagasan bahwa semakin banyak simpul
yang berpartisipasi, semakin besar eksternalitas jaringan; partisipasi
tersebut dapat menghasilkan sinergi horizontal melalui kerja sama
sederhana dan sinergi vertikal melalui pembagian kerja serta
spesialisasi {[}TXT lines 197-220{]}. Dalam kerangka ini,
\emph{borrowing size} dapat dicapai melalui keterlekatan dalam jaringan
dan tidak semata-mata melalui kedekatan geografis {[}TXT lines
212-220{]}.

Keempat, literatur jaringan inovasi membahas faktor internal dan
eksternal. Kepadatan serta sentralitas dikaitkan dengan kinerja inovasi;
hubungan kuat dapat meningkatkan kepercayaan dan berbagi pengetahuan,
sedangkan hubungan lemah dapat memperluas keterbukaan dan mendukung
gagasan baru. Artikel juga merujuk pada argumen bahwa jaringan yang
efektif dapat memerlukan kombinasi struktur internal yang kuat dan
koneksi eksternal yang lebih lemah {[}TXT lines 228-264{]}.

Kelima, kedekatan kognitif, sosial, organisasional, institusional, dan
geografis dipaparkan sebagai faktor yang memengaruhi evolusi jaringan
inovasi. Jarak dapat membatasi efek jaringan, tetapi keterikatan
geografis yang terlalu tertutup juga dapat menghasilkan \emph{lock-in}.
Kedekatan budaya dan institusional dapat memudahkan kepercayaan,
mengurangi ketidakpastian, dan meredakan konflik, tetapi keseragaman
atau kedekatan yang berlebihan juga dipaparkan berpotensi mengurangi
keragaman dan insentif inovasi {[}TXT lines 265-282{]}.

\textbf{Interpretasi penulis:} Konteks GBA penting karena kombinasi
keterbukaan jaringan dan batas institusional-budaya membuat arah
eksternalitas tidak dapat diasumsikan positif. Artikel menempatkan
jarak, budaya, dan institusi sebagai lingkungan eksternal yang perlu
diuji bersama struktur jaringan {[}TXT lines 283-298{]}.

\textbf{Inferensi pengolahan:} Tinjauan ini berbentuk sintesis
konseptual dan pengajuan masalah, bukan tinjauan sistematis yang dapat
direplikasi. Strategi pencarian, kriteria inklusi, jumlah korpus, dan
prosedur penilaian literatur tidak disebutkan dalam file.

\subsection{Method Used}\label{method-used-25}

\textbf{Laporan sumber:} Penelitian menggabungkan dua metode utama.
Analisis jaringan sosial dengan perangkat lunak UCINET digunakan untuk
menghitung sentralitas dan kepadatan. Jaringan inovasi pengetahuan
diperlakukan sebagai jaringan tidak berarah. Sentralitas derajat
dihitung sebagai \texttt{d(i)/(n-1)}, dengan \texttt{d(i)} sebagai
jumlah hubungan langsung simpul \texttt{i} dan \texttt{n-1} sebagai
jumlah hubungan maksimum. Kepadatan dihitung dari jumlah hubungan aktual
\texttt{m} dibandingkan jumlah hubungan maksimum teoretis
\texttt{n(n-1)/2} {[}TXT lines 355-368{]}.

Untuk mengukur eksternalitas, penulis membentuk tiga matriks bobot
spasial: matriks jarak geografis, matriks perbedaan institusional, dan
matriks perbedaan budaya. Persamaan menempatkan nilai jarak atau nilai
perbedaan antarkota pada pasangan kota dan menetapkan diagonal
\texttt{i=j} sebesar 0 {[}TXT lines 371-389{]}. Jarak geografis,
perbedaan institusional, dan perbedaan budaya digunakan sebagai dasar
bobot tanpa penjelasan tambahan dalam TXT mengenai standardisasi atau
normalisasi matriks.

Model spasial yang digunakan disebut sebagai spatial Dubin model dalam
TXT. Penulis menyatakan bahwa hasil uji LM dan robust LM melewati
signifikansi 1\%, lalu menggunakan model tersebut untuk mengestimasi
hubungan antara jumlah paper kota, sentralitas, variabel kontrol, dan
lag spasial sentralitas \texttt{W*centrality} {[}TXT lines 393-402{]}.
Variabel \texttt{paper} diestimasi dalam tiga bentuk: total paper, paper
kolaboratif, dan paper lokal atau paper tanpa kerja sama dengan kota
lain.

Model tambahan menghubungkan kesenjangan jumlah paper antara setiap
pasangan kota dengan kesenjangan sentralitas serta perbedaan variabel
kontrol. Model ini digunakan untuk menguji apakah konsentrasi jaringan
memperlebar kesenjangan inovasi antarkota {[}TXT lines 403-411{]}.

Penulis juga mengidentifikasi kemungkinan endogenitas dua arah: kerja
sama pengetahuan dapat meningkatkan keluaran inovasi, tetapi kemampuan
inovasi yang lebih kuat juga dapat menarik lebih banyak hubungan kerja
sama. Untuk merevisi model, penulis menggunakan lag sentralitas kerja
sama sebagai variabel instrumental dan melaporkan hasilnya pada Table 3
{[}TXT lines 582-593{]}.

\textbf{Interpretasi penulis:} Penggabungan analisis jaringan dengan
model spasial dimaksudkan agar penelitian tidak berhenti pada deskripsi
struktur jaringan, tetapi juga menguji dampak spasial posisi kota dalam
jaringan {[}TXT lines 412-417{]}.

\textbf{Inferensi pengolahan:} TXT tidak menjelaskan model efek tetap
atau acak, struktur error secara rinci, prosedur pembentukan hubungan
jaringan, uji validitas instrumen, atau diagnostik lain di luar
pernyataan bahwa LM dan robust LM signifikan. Karena itu, rincian
identifikasi statistik yang dapat diverifikasi dari file terbatas.

\subsection{Dataset}\label{dataset-25}

\textbf{Laporan sumber:} Artikel membedakan paper dan paten sebagai
indikator inovasi, dengan paper untuk merefleksikan inovasi pengetahuan
dan paten untuk lebih merepresentasikan inovasi teknologi. Namun,
operasionalisasi penelitian ini menggunakan jumlah seluruh paper yang
dipublikasikan di setiap kota, paper kolaboratif antarkota, dan paper
tanpa kerja sama dengan kota lain {[}TXT lines 313-320{]}. Ketiga jenis
data paper tersebut berasal dari Web of Science dan mencakup 2001-2019
{[}TXT lines 313-325{]}.

Data sosial-ekonomi untuk variabel kontrol berasal dari buku tahunan
statistik Provinsi Guangdong dan setiap kota {[}TXT lines 349-351{]}.
Table 1 memperlihatkan kontrol \texttt{Lnresearcher}, \texttt{Lnexpend},
\texttt{lngdp}, \texttt{Popuflow}, dan \texttt{Internet} {[}TXT lines
536-548{]}. Dalam uraian hasil, variabel-variabel tersebut dibaca
sebagai kualitas peneliti, pengeluaran atau pendanaan, tingkat ekonomi,
mobilitas penduduk, dan penetrasi Internet {[}TXT lines 523-555{]}.

Perbedaan budaya diproksikan dengan dialek. Dialek dibagi menjadi tiga
tingkat, yaitu wilayah dialek, distrik dialek, dan zona dialek, lalu
dummy perbedaan budaya diberi nilai 0 sampai 3 untuk setiap pasangan
kota {[}TXT lines 326-338{]}. Perbedaan institusional menggabungkan
perbedaan sistem politik, wilayah pabean, dan hierarki administratif;
nilainya untuk setiap pasangan kota ditetapkan dari 1 sampai 4 {[}TXT
lines 339-348{]}.

Jumlah total paper, jumlah paper per kota dan tahun, kriteria
pengambilan paper dari Web of Science, prosedur atribusi afiliasi ke
kota, penanganan paper dengan banyak afiliasi, data yang hilang, serta
statistik deskriptif tidak disebutkan dalam file. Catatan kaki artikel
secara eksplisit menyatakan bahwa deskripsi data dan karakteristik
statistik tidak ditampilkan untuk menjaga panjang artikel {[}TXT lines
720-722{]}.

\subsection{Population / Sample}\label{population-sample-25}

\textbf{Laporan sumber:} Ruang lingkup GBA terdiri dari dua wilayah
administratif khusus, Hong Kong dan Macao, serta sembilan kota di
Provinsi Guangdong: Guangzhou, Foshan, Zhaoqing, Jiangmen, Zhongshan,
Zhuhai, Shenzhen, Dongguan, dan Huizhou {[}TXT lines 302-309{]}.

\textbf{Inferensi pengolahan:} Daftar tersebut berjumlah 11 kota. Unit
analisis utama adalah kota pada tahun tertentu untuk model inovasi dan
pasangan antarkota untuk model kesenjangan. Jaringan diperlakukan
sebagai jaringan tidak berarah, sehingga hubungan kota digunakan sebagai
pasangan tanpa arah {[}TXT lines 355-368, 403-411{]}.

\textbf{Inferensi pengolahan:} Prosedur sampling tidak disebutkan dalam
file; sumber tampak menggunakan cakupan seluruh kota yang didefinisikan
sebagai GBA, bukan sampel sebagian kota. Jumlah observasi panel, jumlah
pasangan kota yang benar-benar masuk estimasi, dan jumlah paper yang
dianalisis tidak disebutkan dalam file.

\subsection{Dependent Variables}\label{dependent-variables-25}

\textbf{Laporan sumber:} Variabel dependen utama dalam model spasial
adalah \texttt{paper\_i,t}, yang terdiri dari tiga kategori {[}TXT lines
395-402{]}.

\begin{itemize}
\tightlist
\item
  Inovasi keseluruhan atau kemampuan inovasi keseluruhan: jumlah total
  paper yang dipublikasikan kota.
\item
  Kemampuan inovasi kolaboratif: paper hasil kerja sama antara kota.
\item
  Kemampuan inovasi independen: paper yang tidak melibatkan kerja sama
  dengan kota lain.
\end{itemize}

Model kesenjangan menggunakan perbedaan jumlah total paper antara kota
\texttt{i} dan kota \texttt{j} sebagai variabel dependen, yaitu
kesenjangan pengetahuan atau kemampuan inovasi antarkota {[}TXT lines
571-579, 613-626{]}.

Sentralitas bukan variabel dependen utama, melainkan variabel penjelas
inti. \texttt{W*centrality} digunakan untuk menangkap pengaruh posisi
jaringan kota terhadap wilayah sekitarnya. Variabel kontrol yang
ditampilkan meliputi peneliti, pengeluaran, GDP, arus penduduk, dan
Internet {[}TXT lines 397-402, 536-548{]}.

\subsection{Results}\label{results-25}

\textbf{Pola jaringan yang dilaporkan sumber:} Jumlah total paper GBA
meningkat sepanjang 2001-2019. Sebelum 2011, Hong Kong berada di
peringkat pertama untuk total paper. Pada 2013, total paper dan paper
lokal Guangzhou melampaui Hong Kong. Guangzhou mempertahankan posisi
pertama untuk paper kerja sama sepanjang periode. Setelah 2012, Shenzhen
berada di urutan ketiga untuk total paper dan kesenjangannya dengan Hong
Kong terus menyempit. Jumlah paper Macao relatif kecil {[}TXT lines
421-431{]}.

Kepadatan keseluruhan jaringan meningkat dari 2001 sampai 2019, yang
oleh penulis dibaca sebagai peningkatan hubungan kerja sama inovasi
antarkota. Hong Kong memiliki sentralitas keseluruhan tertinggi sebelum
2012, kemudian dilampaui Guangzhou setelah 2013. Shenzhen melampaui
Zhongshan pada 2013 dan berada di peringkat ketiga; pada 2019
sentralitas keseluruhannya mendekati Hong Kong {[}TXT lines 432-439{]}.

Untuk sentralitas kerja sama, Guangzhou selalu berada di posisi
tertinggi. Zhongshan berada di posisi kedua sebelum 2018, tetapi berada
di belakang Hong Kong pada 2019; Shenzhen berada di posisi keempat
{[}TXT lines 440-442{]}. Penulis merangkum bahwa Guangzhou menjadi kota
inti sejak 2013, Hong Kong dan Shenzhen memiliki kapasitas inovasi kuat
tetapi kerja sama dengan kota lain relatif terbatas, dan Macao lemah
dalam inovasi keseluruhan maupun kerja sama antarkota {[}TXT lines
456-459{]}.

Figure 8 menggambarkan perubahan dari struktur segitiga yang berpusat
pada Hong Kong, Guangzhou, dan Zhongshan pada 2001, menjadi struktur
empat tautan yang menampilkan Guangzhou, Hong Kong, Shenzhen, dan
Zhongshan pada 2010. Untuk 2019, TXT secara literal menyebut struktur
berlian dengan ``Guangzhou, Shenzhen and Guangzhou'' sebagai tiga pusat,
serta Zhongshan, Dongguan, dan Foshan sebagai subpusat {[}TXT lines
460-464{]}. Pengulangan nama Guangzhou tersebut tidak dikoreksi dalam
review ini karena nama pusat ketiga yang dimaksud tidak dapat dipastikan
dari TXT.

\textbf{Eksternalitas inovasi keseluruhan dan kerja sama:} Pada
spesifikasi berbasis jarak, sentralitas kerja sama tidak berpengaruh
signifikan terhadap inovasi keseluruhan kota sendiri, tetapi
\texttt{W*centrality} berdampak negatif dan signifikan terhadap wilayah
sekitar, dengan koefisien -0.91**. Pada spesifikasi berbasis perbedaan
institusional, sentralitas sendiri berdampak negatif dan signifikan
(-0.14\emph{), sedangkan lag spasialnya juga negatif dan signifikan
(-0.89}). Pada spesifikasi budaya, pengaruh sentralitas (-0.03) dan lag
spasial (0.44) tidak signifikan {[}TXT lines 501-509, 536-548{]}.

Untuk kemampuan inovasi kolaboratif, sentralitas tidak signifikan pada
matriks jarak (-0.15), negatif dan signifikan 10\% pada matriks
institusi (-0.17*), serta negatif dan signifikan 5\% pada matriks budaya
(-0.57\textbf{). Lag spasialnya negatif dan signifikan pada ketiga
spesifikasi: -0.83* untuk jarak, -0.91* untuk institusi, dan -4.19}
untuk budaya {[}TXT lines 536-548{]}.

Untuk kemampuan inovasi independen, sentralitas sendiri bernilai negatif
dan tidak signifikan pada jarak (-0.10), negatif dan signifikan pada
institusi (-0.11\emph{), serta negatif dan signifikan pada budaya
(-0.14\textbf{). \texttt{W*centrality} negatif dan signifikan pada semua
spesifikasi: -0.93} untuk jarak, -0.82\textbf{ untuk institusi, dan
-1.72}} untuk budaya {[}TXT lines 565-569, 536-548{]}. Tanda
signifikansi dalam Table 1 adalah *** 1\%, ** 5\%, dan * 10\%; angka
dalam tanda kurung mengikuti catatan tabel sebagai standard errors
{[}TXT lines 536-548{]}.

Penulis menyatakan bahwa hasil tersebut menunjukkan jaringan pengetahuan
GBA tidak memiliki eksternalitas positif: kota dengan sentralitas tinggi
berdampak negatif pada inovasi kota sekitar, dan dampaknya lebih nyata
pada kota dengan perbedaan institusional besar {[}TXT lines 501-509{]}.

\textbf{Inferensi pengolahan:} Uraian penulis yang menyebut pengaruh
budaya ``relatively weak'' perlu dibaca per keluaran, karena Table 1
tetap menunjukkan beberapa koefisien budaya yang negatif dan signifikan
untuk inovasi kolaboratif serta independen {[}TXT lines 504-506,
536-548{]}.

\textbf{Kontrol:} Kualitas peneliti berpengaruh positif terhadap inovasi
lokal, tetapi menjadi tidak signifikan ketika perbedaan budaya
dimasukkan. Pendanaan atau pengeluaran, tingkat ekonomi, mobilitas
penduduk, dan penetrasi Internet masing-masing menunjukkan pola yang
berbeda antarspesifikasi; secara umum penulis melaporkan pengaruh
positif pendanaan dan tingkat ekonomi, hubungan negatif mobilitas
penduduk dengan inovasi pengetahuan, serta pengaruh positif penetrasi
Internet {[}TXT lines 523-555{]}.

\textbf{Kesenjangan antarkota:} Kesenjangan sentralitas berhubungan
positif dan sangat signifikan dengan kesenjangan kemampuan inovasi;
koefisien \texttt{gap\_centrality} pada lima model berada sekitar 0.9958
sampai 1.0007 dan semuanya diberi tanda signifikansi 1\% {[}TXT lines
571-577, 613-626{]}. Interaksi perbedaan jarak dengan sentralitas
signifikan pada Model 2, tetapi tidak pada Model 5. Interaksi perbedaan
institusional signifikan pada Model 3 dan Model 5, sedangkan interaksi
budaya tidak signifikan pada model yang ditampilkan {[}TXT lines
616-625{]}. Penulis menyatakan jarak dan perbedaan institusional
memperparah eksternalitas negatif pusat inovasi, sementara perbedaan
budaya relatif kecil pengaruhnya. Kesenjangan tingkat Internet juga
dilaporkan memperlebar kesenjangan inovasi {[}TXT lines 574-579{]}.

\textbf{Uji endogenitas dan robustness:} Lag sentralitas kerja sama
memiliki pengaruh negatif pada kemampuan inovasi. Dalam matriks jarak,
pengaruhnya tidak signifikan untuk inovasi keseluruhan (-0.0578),
kolaboratif (-0.0996), maupun independen (-0.0393). Dalam matriks
institusi, pengaruhnya negatif dan signifikan untuk ketiga keluaran
(-0.1360**, -0.1432*, dan -0.1299**). Dalam matriks budaya, pengaruhnya
negatif dan signifikan untuk inovasi keseluruhan (-0.0437*) dan
kolaboratif (-0.0763**), tetapi tidak signifikan untuk inovasi
independen (-0.0361) {[}TXT lines 582-593, 630-638{]}. Penulis menilai
hasil ini konsisten dengan hasil sebelumnya sehingga hasilnya robust dan
reliable {[}TXT lines 588-593{]}.

\textbf{Mekanisme yang diajukan penulis:} Penulis menyatakan bahwa
perbedaan tingkat inovasi awal antarkota dan dukungan kebijakan yang
tidak sama, khususnya dalam perkembangan Shenzhen, dapat memperkuat
ketidakseimbangan jaringan, penjarahan sumber daya, dan kesenjangan.
Kota dengan sentralitas tinggi juga disebut lebih toleran terhadap
perbedaan budaya dan institusi serta lebih unggul dalam keragaman budaya
dan reformasi institusional, sehingga lebih mampu menarik talenta dan
sumber daya inovasi dari kota sekitar {[}TXT lines 594-644{]}. Formulasi
penulis menyebut alasan-alasan ini sebagai kemungkinan, bukan sebagai
mekanisme yang diuji secara terpisah dalam TXT.

\subsection{Findings}\label{findings-25}

\textbf{Laporan sumber:} Jaringan inovasi pengetahuan GBA bertumbuh
dalam jumlah paper, kepadatan, dan beberapa indikator sentralitas,
tetapi peningkatan konektivitas tersebut tidak menghasilkan
eksternalitas positif bagi kota sekitar. Guangzhou dan Hong Kong
merupakan pusat historis, Shenzhen menguat setelah 2012, dan Macao belum
menjadi pusat. Kerja sama paling banyak terjadi di antara kota-kota
Guangdong, terutama Guangzhou dan Shenzhen {[}TXT lines 648-661{]}.

\textbf{Interpretasi penulis:} Kota kecil dan menengah berada dalam
\emph{agglomeration shadow}, bukan memperoleh \emph{borrowing size}.
Perbedaan institusional dan budaya dianggap lebih menghambat kerja sama
inovasi daripada jarak geografis, sehingga perbaikan lingkungan lunak
diposisikan sebagai prasyarat untuk mengurangi kesenjangan spasial
{[}TXT lines 662-668{]}.

\textbf{Inferensi pengolahan:} Temuan utama memiliki dua lapis. Pada
lapis struktur, jaringan menjadi lebih padat dan memiliki beberapa
pusat. Pada lapis distribusi manfaat, sentralitas yang lebih tinggi
berkorelasi dengan kesenjangan yang lebih besar dan dengan kinerja
inovasi yang lebih rendah di wilayah sekitar dalam spesifikasi yang
dilaporkan. Dengan demikian, ``multisentra'' secara struktural tidak
otomatis berarti pemerataan manfaat inovasi. Kesimpulan inferensial ini
tetap dibatasi pada indikator paper, desain model, wilayah, dan periode
yang tersedia dalam TXT.

\subsection{Conclusions}\label{conclusions-25}

\textbf{Laporan sumber:} Penulis menyimpulkan bahwa kemampuan inovasi
pengetahuan dan kepadatan jaringan GBA terus meningkat, tetapi pola
empat pusat yang direncanakan pemerintah belum terbentuk. Guangzhou dan
Hong Kong selalu menjadi pusat, Shenzhen menjadi pusat tambahan sejak
2012, dan Macao belum cukup kuat. Kerja sama Hong Kong dengan kota
daratan relatif terbatas selain dengan Guangzhou dan Shenzhen, sedangkan
sebagian besar kerja sama antarkota berlangsung di Guangdong {[}TXT
lines 648-661{]}.

Model spasial dipakai untuk menyimpulkan bahwa sentralitas inovasi kota
berdampak negatif pada wilayah sekitar. Karena itu, kota kecil dan
menengah tidak memperoleh \emph{borrowing size}, tetapi berada dalam
\emph{agglomeration shadow}. Penulis juga menyimpulkan bahwa perbedaan
budaya dan institusional jauh lebih negatif bagi inovasi daripada jarak
{[}TXT lines 662-665{]}.

\textbf{Interpretasi penulis:} Pengoptimalan jaringan memerlukan
pengurangan hambatan institusional, lingkungan sosial yang lebih
toleran, infrastruktur transportasi dan Internet yang memadai, penguatan
basis riset Macao, serta perlindungan hak kekayaan intelektual {[}TXT
lines 666-707{]}. Saran tersebut merupakan agenda kebijakan yang
diajukan penulis, bukan evaluasi atas kebijakan yang sudah dilaksanakan.

\subsection{Contributions}\label{contributions-25}

\textbf{Laporan sumber:} Artikel menyatakan dua pembeda utama dari
penelitian sebelumnya. Pertama, artikel tidak hanya menganalisis
struktur jaringan perkotaan, tetapi juga efek spasialnya. Kedua, artikel
memusatkan perhatian pada kerja sama inovasi antarkota dan dampak
spasial kota pusat {[}TXT lines 147-151{]}. Penulis juga menyatakan
bahwa penggabungan metode jaringan sosial dan model pengukuran spasial
memungkinkan analisis atribut jaringan sekaligus efek spasialnya {[}TXT
lines 412-417{]}.

\textbf{Inferensi pengolahan:} Kontribusi operasional artikel adalah
menjadikan sentralitas sebagai indikator posisi jaringan, lalu
menghubungkannya dengan tiga keluaran paper dan kesenjangan antarkota
melalui tiga jenis lingkungan spasial. Hal ini memberi cara yang lebih
spesifik untuk membedakan ``jaringan yang lebih terkoneksi'' dari
``jaringan yang manfaatnya menyebar''. Klaim kontribusi ini adalah
pembacaan atas rancangan artikel; TXT tidak menyediakan pembandingan
kuantitatif dengan semua studi terdahulu.

\subsection{Research Gap}\label{research-gap-25}

\textbf{Gap yang dinyatakan sumber:} Penelitian terdahulu tentang GBA
disebut telah membahas evolusi struktur jaringan, kinerja inovasi, dan
faktor yang memengaruhinya, tetapi mengabaikan eksternalitas jaringan.
Akibatnya, penelitian tersebut belum dapat menunjukkan efek spasial
sebenarnya dari jaringan inovasi {[}TXT lines 286-298{]}. Artikel ini
mengisi gap itu dengan menguji dampak jaringan terhadap kota sendiri,
kota sekitar, dan kesenjangan antarkota.

\textbf{Gap lanjutan yang dinyatakan sumber:} Penulis menyebut bahwa
pengukuran struktur jaringan masih dapat diperbaiki, efek jaringan
terhadap ekonomi, masyarakat, dan tata ruang kota perlu dianalisis lebih
dalam, serta eksternalitas jaringan inovasi teknologi perlu dikaji
karena situasi dan mekanismenya berbeda dari inovasi pengetahuan {[}TXT
lines 708-717{]}.

\textbf{Batas gap dari TXT:} Teks tidak menyebutkan adanya analisis pada
tingkat perusahaan, kualitas paper, kutipan, jenis teknologi, atau
periode setelah 2019. Hal ini tidak otomatis berarti aspek-aspek
tersebut tidak pernah diteliti; hanya saja aspek tersebut tidak tersedia
sebagai cakupan artikel dalam file yang dibaca.

\subsection{Limitations}\label{limitations-25}

\textbf{Keterbatasan yang dinyatakan atau tersirat langsung oleh
sumber:} Artikel tidak memiliki bagian berjudul limitations yang
terpisah. Catatan kaki menyatakan bahwa deskripsi data dan karakteristik
statistik tidak ditampilkan {[}TXT lines 720-722{]}. Pada bagian future
research, penulis sendiri mengakui perlunya perbaikan pengukuran
struktur jaringan, pendalaman efek sosial-ekonomi dan tata ruang, serta
pemisahan analisis inovasi pengetahuan dari inovasi teknologi {[}TXT
lines 708-717{]}.

\textbf{Keterbatasan yang dapat diidentifikasi dari TXT (inferensi
pengolahan):}

\begin{itemize}
\tightlist
\item
  Dataset hanya mengoperasionalkan inovasi pengetahuan melalui jumlah
  paper. TXT tidak melaporkan ukuran kualitas paper, sitasi, bidang
  ilmu, atau indikator paten dalam estimasi utama.
\item
  Jumlah paper, jumlah observasi, statistik deskriptif, kriteria inklusi
  Web of Science, dan prosedur penempatan afiliasi paper ke kota tidak
  disebutkan dalam file.
\item
  Sentralitas derajat didefinisikan berdasarkan jumlah hubungan, tetapi
  TXT tidak menjelaskan apakah hubungan dibobot oleh jumlah paper
  kolaboratif, memakai ambang tertentu, atau dibentuk sebagai hubungan
  biner.
\item
  Proksi budaya berupa perbedaan dialek dan proksi institusi berupa
  gabungan sistem politik, wilayah pabean, serta hierarki administratif.
  TXT tidak melaporkan uji validitas konstruk atau alternatif pengukuran
  untuk kedua proksi tersebut.
\item
  Persamaan matriks bobot menampilkan jarak dan perbedaan secara
  langsung, tetapi TXT tidak menjelaskan transformasi, normalisasi, atau
  alasan pemilihan bentuk bobot tersebut. Karena itu, perbandingan besar
  koefisien lintas matriks perlu kehati-hatian.
\item
  Penulis menyebut lag sentralitas sebagai variabel instrumental untuk
  menangani endogenitas, tetapi hasil uji tahap pertama, kekuatan
  instrumen, pembatasan eksklusi, dan diagnostik instrumen lain tidak
  disebutkan dalam file.
\item
  Uji LM dan robust LM hanya dilaporkan melalui pernyataan bahwa
  hasilnya melewati 1\%; nilai statistik, derajat kebebasan, dan rincian
  keputusan model tidak ditampilkan dalam TXT.
\item
  Angka pada Figure 2 sampai Figure 8 tidak tersedia sebagai deret
  numerik dalam ekstraksi teks. Struktur yang diringkas dari gambar
  bergantung pada uraian dan caption yang terbaca.
\item
  Deskripsi Figure 8 pada 2019 mengulang nama Guangzhou sebagai dua dari
  tiga pusat. TXT tidak menyediakan koreksi yang dapat diverifikasi.
\end{itemize}

Keterbatasan-keterbatasan inferensial tersebut tidak membatalkan hasil
yang dilaporkan, tetapi membatasi pemeriksaan ulang dan penilaian atas
mekanisme kausalnya hanya dari TXT.

\subsection{Challenges}\label{challenges-25}

\textbf{Laporan sumber:} Tantangan konteks utama adalah GBA menjalankan
dua sistem politik dan tiga wilayah pabean. Sembilan kota Guangdong
berada dalam satu wilayah pabean, sedangkan Hong Kong dan Macao memiliki
wilayah pabean masing-masing; sistem politik dan tingkat
administratifnya juga berbeda {[}TXT lines 304-309, 339-348{]}.
Perbedaan tersebut membuat kerja sama inovasi dan arus faktor lintas
kota sulit.

Penulis melaporkan bahwa mekanisme kerja sama lintas batas belum
memadai, sehingga menghambat mobilitas personel, berbagi sumber daya
inovasi, dan proyek bersama {[}TXT lines 669-677{]}. Perbedaan budaya,
yang diproksikan dengan dialek, juga menjadi tantangan; jarak geografis
tetap membatasi arus faktor, tetapi disebut lebih kecil pengaruhnya
dibanding hambatan institusional dan budaya {[}TXT lines 552-564,
687-695{]}.

Tantangan lain adalah ketimpangan kapasitas awal dan dukungan kebijakan.
Penulis mengaitkannya dengan konsentrasi jaringan, daya tarik kota pusat
terhadap talenta dan sumber daya, serta kesenjangan perkembangan kota
kecil dan menengah {[}TXT lines 594-644{]}. Dalam pelaksanaan kebijakan,
pembangunan transportasi perbatasan, konektivitas Internet, pendidikan,
inklusi sosial, riset Macao, dan perlindungan kekayaan intelektual juga
diposisikan sebagai pekerjaan yang belum selesai {[}TXT lines
678-707{]}.

\textbf{Tantangan analitis:} TXT menyebut hubungan timbal balik antara
sentralitas kerja sama dan kemampuan inovasi sebagai masalah endogenitas
{[}TXT lines 582-588{]}. Cara mengatasinya dilaporkan melalui lag
sentralitas, tetapi rincian pengujian validitasnya tidak disebutkan
dalam file.

\subsection{Practical Implications}\label{practical-implications-25}

\textbf{Laporan sumber dan rekomendasi penulis:}

\begin{itemize}
\tightlist
\item
  Pemerintah kota perlu mengurangi hambatan institusional antara kota
  daratan dan Hong Kong/Macao melalui kebijakan yang memudahkan
  mobilitas personel, berbagi sumber daya inovasi, dan pengembangan
  proyek kerja sama {[}TXT lines 669-675{]}.
\item
  Kota kecil dan menengah perlu meningkatkan investasi inovasi
  pengetahuan dan reformasi institusional agar kesenjangan dengan kota
  pusat menyempit {[}TXT lines 675-677{]}.
\item
  Lingkungan sosial perlu dibuat lebih toleran terhadap keragaman budaya
  dan kepribadian. Penulis mengusulkan perlindungan kepentingan imigran,
  perlakuan yang setara, dan peningkatan pendidikan untuk memperkuat
  inklusivitas {[}TXT lines 678-683{]}.
\item
  Kota-kota perlu bekerja sama membangun transportasi di wilayah
  perbatasan dan memperbaiki fasilitas Internet. Hubungan lalu lintas
  antara kota daratan dan Hong Kong/Macao serta akses Internet bagi kota
  kecil dan menengah diberi perhatian khusus {[}TXT lines 687-695{]}.
\item
  Macao disarankan memperkuat investasi riset ilmiah dan pertukaran
  talenta dengan Hong Kong serta Guangzhou untuk memperkuat riset dasar
  {[}TXT lines 696-701{]}.
\item
  Perlindungan hak kekayaan intelektual perlu diperkuat melalui
  pemanfaatan pengalaman Hong Kong, pembentukan pusat perlindungan dan
  transaksi regional, serta mekanisme pertukaran informasi dan platform
  berbagi {[}TXT lines 702-707{]}.
\end{itemize}

\textbf{Batas evidensi:} Saran tersebut merupakan implikasi kebijakan
yang diajukan penulis dari hasil model. TXT tidak melaporkan evaluasi
dampak atau uji implementasi atas rekomendasi tersebut.

\subsection{Applications}\label{applications-25}

\textbf{Aplikasi yang disebut atau tersirat langsung oleh sumber:} Hasil
artikel dapat digunakan sebagai dasar untuk mengoptimalkan pola jaringan
eksternalitas inovasi GBA, mengevaluasi dampak kota pusat terhadap kota
sekitar, dan mengarahkan kebijakan pada pengurangan kesenjangan spasial
{[}TXT lines 147-155, 662-668{]}.

\textbf{Inferensi pengolahan:} Indikator sentralitas, kepadatan, dan
kesenjangan antarkota yang digunakan dalam artikel dapat menjadi
perangkat diagnosis untuk memantau apakah pertumbuhan jaringan diikuti
penyebaran manfaat.

\textbf{Inferensi pengolahan:} Dalam batas artikel, aplikasi paling
dekat adalah perencanaan kerja sama lintas batas dan penentuan prioritas
intervensi pada hambatan institusional, budaya, transportasi, Internet,
kapasitas riset, serta hak kekayaan intelektual. Pernyataan bahwa
perangkat tersebut cocok untuk wilayah lain, sektor lain, atau inovasi
teknologi tidak disebutkan dalam file dan tidak diasumsikan di sini.

\subsection{Future Research
(Optional)}\label{future-research-optional-25}

\textbf{Laporan sumber:} Penulis mengusulkan tiga arah penelitian:

\begin{itemize}
\tightlist
\item
  Memperbaiki pengukuran struktur jaringan. Contoh yang diberikan adalah
  \emph{locally weighted regression} untuk mengidentifikasi pusat
  jaringan secara lebih presisi melalui pemodelan permukaan kepadatan
  kota {[}TXT lines 708-711{]}.
\item
  Menganalisis lebih dalam efek spasial struktur jaringan terhadap aspek
  ekonomi, sosial, dan tata ruang kota, bukan hanya terhadap inovasi
  {[}TXT lines 712-715{]}.
\item
  Menganalisis eksternalitas jaringan dari perspektif inovasi teknologi,
  karena situasi dan mekanismenya disebut berbeda dari inovasi
  pengetahuan {[}TXT lines 715-717{]}.
\end{itemize}

Arah lain di luar tiga usulan tersebut: Tidak disebutkan dalam file.

\subsection{Provenance Notes}\label{provenance-notes-25}

\begin{itemize}
\tightlist
\item
  Kode review: A10.
\item
  Sumber tunggal yang digunakan:
  \texttt{/Users/mac/Documents/Mac/{[}2{]}\ Obsidian\ Vault/zettelkasten/source/journal/A10-yang-fan-wang-yu-2022-knowledge-innovation-network-externalities-in-the-guangdong-hong-kong-macao-greater-bay-area-borrowing-size-or-agglomeration-shadow-10.1080-09537325.2021.1940922.txt}.
\item
  TXT menyatakan bahwa sumber PDF berada pada path
  \texttt{journal/A10-yang-fan-wang-yu-2022-knowledge-innovation-network-externalities-in-the-guangdong-hong-kong-macao-greater-bay-area-borrowing-size-or-agglomeration-shadow-10.1080-09537325.2021.1940922.pdf},
  diekstrak pada 2026-08-31 dengan \texttt{pdftotext\ -layout} {[}TXT
  lines 1-4{]}.
\item
  Seluruh 863 baris TXT dibaca, termasuk blok metadata PDF, teks
  artikel, catatan kaki, disclosure statement, funding, daftar
  kontributor, ORCID, dan referensi {[}TXT lines 1-863{]}. Review ini
  hanya memakai isi TXT tersebut; referensi yang tercantum tidak
  diperlakukan sebagai bukti yang dibaca secara terpisah.
\item
  Locator utama menggunakan nomor baris TXT. Locator tambahan berupa
  bagian artikel, Table 1, Table 2, Table 3, dan Figure 8 sebagaimana
  ditandai dalam ekstraksi.
\item
  Basis akses: full-text dalam bentuk ekstraksi TXT. Gambar tersedia
  dalam TXT sebagai caption dan deskripsi terbatas, bukan sebagai data
  visual yang dapat dibaca ulang secara numerik.
\item
  Disclosure statement menyatakan tidak ada potensi konflik kepentingan
  yang dilaporkan. Sumber menyebut dukungan NSFC Grant No.~42071154,
  Ministry of Education Humanities and Social Sciences Planning Fund
  Grant No.~20YJA790010, dan Natural Science Foundation of Shanghai
  Grant No.~21ZR1421100 {[}TXT lines 726-734{]}.
\item
  Elemen metadata jurnal yang tidak tercantum lengkap, termasuk ISSN
  yang tampil kosong pada header, volume, issue, dan halaman
  bibliografis: Tidak disebutkan dalam file.
\item
  Tidak ada informasi dari luar TXT yang digunakan untuk memperbaiki
  tahun pada nama file, mengoreksi Figure 8, memvalidasi status
  peer-review, atau menambahkan jumlah observasi maupun paper.
\end{itemize}

\section{Review A11}\label{review-a11}

\textbf{Judul artikel:} \emph{Do polycentric urban regions promote
functional spillovers and economic performance? Evidence from China}\\
\textbf{Penulis:} Yixiao Wang, Bindong Sun, dan Tinglin Zhang\\
\textbf{Jurnal:} \emph{Regional Studies}\\
\textbf{Tahun dan tanggal terbit:} 2020; diterbitkan daring pada 17
Desember 2020\\
\textbf{DOI:} 10.1080/00343404.2020.1841147\\
\textbf{Kode kajian:} A11\\
\textbf{Jenis sumber:} Artikel jurnal empiris\\
\textbf{Wilayah kajian:} 13 klaster kota di China\\
\textbf{Periode data:} 2004-2014, dengan jeda dua tahun\\
\textbf{Status peer-review:} Tidak dapat diverifikasi dari file TXT.\\
\textbf{Tidak disebutkan dalam file:} Volume, nomor terbitan, dan
rentang halaman jurnal tidak disebutkan secara eksplisit dalam file.
Metadata PDF hanya menyebutkan 13 halaman; nomor halaman artikel yang
tampil dalam ekstraksi mencapai halaman 12.

Nama file A11 memuat tahun 2022, tetapi metadata PDF dan sitasi pada
artikel menyebut tahun 2020, dengan penerbitan daring pada 17 Desember
2020. Review ini mengikuti tahun yang dinyatakan di dalam artikel dan
metadata sumber, bukan tahun pada nama file. (Laporan sumber; locator:
TXT baris 1-8, 40-55.)

\subsection{Identitas}\label{identitas}

Artikel ini menguji apakah struktur spasial polisentris pada wilayah
perkotaan polisentris (polycentric urban regions/PURs) menghasilkan
spillover fungsional lintas batas klaster dan, melalui spillover
tersebut, meningkatkan produktivitas tenaga kerja. Fokus empirisnya
adalah fungsi sektor jasa produsen dan relasinya dengan sektor
manufaktur pada klaster kota di China. Afiliasi yang tampil dalam file
menghubungkan ketiga penulis dengan beberapa unit East China Normal
University di Shanghai, China, tetapi pemetaan rinci tiap unit kepada
tiap penulis tidak sepenuhnya jelas karena tata letak ekstraksi.
(Laporan sumber; locator: TXT baris 40-55 dan 121-136, judul, penulis,
afiliasi, serta keterangan penerbitan.)

Pendanaan yang dilaporkan berasal dari Major Program of the National
Social Science Foundation of China (grant 17ZDA068), National Natural
Science Foundation of China (grant 42001183), dan Humanity and Social
Science Youth Foundation of the Ministry of Education of China (grant
20YJCZH233). Pernyataan konflik menyebutkan bahwa tidak ada potensi
konflik kepentingan yang dilaporkan. (Laporan sumber; locator: TXT baris
786-803, bagian ``DISCLOSURE STATEMENT'' dan ``FUNDING''.)

Artikel mendefinisikan PUR sebagai sekumpulan pusat perkotaan tanpa satu
pusat yang dominan dalam jumlah penduduk. Dalam artikel ini,
polycentricity yang diukur adalah polycentricity morfologis berdasarkan
distribusi rank-size penduduk kota, sedangkan outcome fungsionalnya
dibangun dari fungsi jasa produsen. (Laporan sumber; locator: TXT baris
114-119 dan 290-316, bagian ``INTRODUCTION'' serta ``The measurement of
the mono/polycentricity of city clusters''.)

Penanda yang digunakan dalam review ini adalah sebagai berikut:

\begin{itemize}
\tightlist
\item
  \textbf{Laporan sumber} berarti apa yang secara eksplisit dilaporkan,
  diukur, atau ditampilkan oleh artikel.
\item
  \textbf{Interpretasi penulis} berarti penjelasan teoritis, pembacaan
  hasil, kesimpulan, atau implikasi yang dinyatakan oleh penulis
  artikel.
\item
  \textbf{Inferensi pengolahan} berarti kesimpulan terbatas yang ditarik
  dari isi TXT untuk menilai cakupan atau kehati-hatian interpretasi;
  bagian ini bukan pernyataan penulis.
\end{itemize}

\subsection{Summarized Abstract}\label{summarized-abstract-26}

\textbf{Laporan sumber:} Artikel berangkat dari klaim bahwa literatur
sering mengaitkan kekuatan ekonomi PUR dengan kemampuan kota-kota
penyusunnya menghasilkan spillover fungsional melalui efek
\emph{borrowed size}, tetapi bukti empirisnya masih terbatas. Dengan
data panel 2004-2014 untuk 13 klaster kota China, artikel melaporkan
bahwa polycentricity meningkatkan spillover fungsional positif sektor
jasa produsen ke luar batas klaster hanya ketika klaster memiliki
populasi lebih besar dan konektivitas infrastruktur perkotaan yang lebih
baik. Spillover fungsional tersebut kemudian dikaitkan dengan
produktivitas tenaga kerja yang lebih tinggi. Artikel mengusulkan
pembagian kerja antarkota sebagai mekanisme penting yang menghubungkan
polycentricity dengan spillover fungsional. (Locator: TXT baris 84-90
dan 169-179, bagian ``ABSTRACT'' dan ``INTRODUCTION'', PDF hlm. 1-2.)

\textbf{Interpretasi penulis:} Kekuatan ekonomi PUR melalui efek
\emph{borrowed size} bukanlah efek otomatis dari bentuk polisentris.
Manfaatnya bersyarat pada skala populasi dan konektivitas infrastruktur,
setidaknya untuk klaster kota China dan fungsi jasa produsen yang
dianalisis. (Locator: TXT baris 174-179, 717-726, dan 732-737, bagian
``INTRODUCTION'' dan ``CONCLUSIONS'', PDF hlm. 2 dan 10.)

\textbf{Inferensi pengolahan:} Ringkasan artikel paling tepat dibaca
sebagai argumen kondisional: polycentricity -\textgreater{} pembagian
kerja/spillover fungsional -\textgreater{} produktivitas, dengan
populasi dan konektivitas sebagai kondisi penguat. TXT tidak mendukung
penyederhanaan bahwa semua PUR atau semua bentuk polycentricity pasti
meningkatkan kinerja ekonomi.

\subsection{Problem Statement}\label{problem-statement-26}

\textbf{Laporan sumber:} Artikel menempatkan PUR dalam ketegangan antara
dua kemungkinan. Dibandingkan wilayah monosen-trik, kota-kota dalam PUR
berukuran lebih kecil secara individual dan berjarak lebih jauh,
sehingga dapat mengalami kerugian aglomerasi, keterbatasan massa kritis,
dan sirkulasi arus yang kurang nyaman. Di sisi lain, interaksi antarkota
dapat menghasilkan eksternalitas jaringan melalui efek \emph{borrowed
size}. Sebagian kota dapat memiliki fungsi lebih besar daripada yang
diperkirakan dari ukurannya dan menikmati \emph{borrowed size},
sedangkan kota lain dapat memiliki fungsi lebih kecil dan mengalami
\emph{agglomeration shadow}. (Locator: TXT baris 101-118 dan 143-163,
bagian ``INTRODUCTION'', PDF hlm. 1-2.)

\textbf{Laporan sumber:} Masalah empiris yang diidentifikasi adalah
kelangkaan bukti tentang apakah PUR memperoleh keuntungan ekonomi
melalui spillover fungsional regional. Penelitian terdahulu yang
diringkas artikel lebih sering menguji \emph{borrowed size} atau
\emph{agglomeration shadow} pada kota individual, bukan spillover
fungsional dan kinerja ekonomi PUR sebagai keseluruhan. Artikel juga
menyatakan bahwa mekanisme di balik eksternalitas jaringan kota belum
cukup dieksplorasi. (Locator: TXT baris 164-174 dan 236-260, bagian
``INTRODUCTION'' serta ``Literature review'', PDF hlm. 2-3.)

\textbf{Interpretasi penulis:} Persoalan yang hendak dijawab bukan
sekadar apakah struktur kota lebih polisentris, melainkan apakah
polycentricity dapat menghasilkan spillover fungsional positif yang
totalnya mengalahkan \emph{agglomeration shadow}, lalu memberi manfaat
ekonomi pada tingkat klaster. Karena keterbatasan data dan metode,
artikel tidak menguji hipotesis besar bahwa keunggulan ekonomi PUR
mengungguli ekonomi aglomerasi wilayah monosen-trik yang memiliki pusat
utama lebih besar. Artikel mempersempit persoalan menjadi dua hipotesis
konkret tentang hubungan polycentricity, spillover, populasi,
konektivitas, dan produktivitas. (Locator: TXT baris 208-220, PDF hlm.
3.)

\textbf{Inferensi pengolahan:} Batas persoalan tersebut penting untuk
review: artikel memberikan uji empiris atas rantai hubungan tertentu,
bukan perbandingan menyeluruh antara seluruh manfaat dan biaya PUR
versus wilayah monosen-trik.

\subsection{Objectives}\label{objectives-26}

\textbf{Laporan sumber:} Tujuan empiris yang dinyatakan adalah menguji
dampak polycentricity 13 klaster kota China terhadap spillover
fungsional regional di luar batas klaster pada sektor jasa produsen,
serta outcome ekonomi yang mengikutinya. Strategi estimasi terlebih
dahulu menguji hubungan polycentricity-spillover dan spillover-kinerja
ekonomi, kemudian mengeksplorasi pembagian kerja antarkota sebagai
mekanisme penghubung. (Locator: TXT baris 169-179 dan 330-337, bagian
``INTRODUCTION'' dan ``DATA AND METHOD'', PDF hlm. 2 dan 4.)

\textbf{Laporan sumber:} Artikel merumuskan dua hipotesis:

\begin{itemize}
\tightlist
\item
  \textbf{Hipotesis 1:} PUR dapat memperdalam pembagian kerja dan kerja
  sama antarkota dengan menyediakan lokasi produksi yang lebih
  menguntungkan daripada wilayah monosen-trik, sehingga menghasilkan
  lebih banyak spillover fungsional regional yang mendukung kinerja
  ekonomi.
\item
  \textbf{Hipotesis 2:} Pengaruh polycentricity terhadap pembagian kerja
  spasial dan spillover fungsional bergantung pada peluang interaksi
  antarkota. Populasi yang lebih besar dan konektivitas infrastruktur
  yang lebih baik dipandang penting untuk memfasilitasi interaksi dan
  pembagian kerja.
\end{itemize}

(Locator: TXT baris 222-228 dan 273-287, bagian ``Our hypotheses'', PDF
hlm. 3-4.)

\textbf{Interpretasi penulis:} Tujuan teoritisnya adalah menjelaskan
bagaimana pembagian kerja antara sektor jasa produsen dan manufaktur
bergeser dari berlangsung di dalam kota besar menjadi berlangsung
antarkota seiring meningkatnya polycentricity. (Locator: TXT baris
230-266 dan 317-328, bagian ``Our hypotheses'', PDF hlm. 3-4.)

\subsection{Literature Survey}\label{literature-survey-26}

\textbf{Laporan sumber:} Tinjauan literatur menelusuri asal-usul dan
perluasan konsep \emph{borrowed size}. Alonso memperkenalkannya sebagai
keadaan ketika kota kecil atau wilayah metropolitan kecil dapat
menunjukkan karakteristik kota yang lebih besar karena berada dekat
konsentrasi penduduk dan memanfaatkan infrastruktur serta fasilitas
tingkat tinggi kota besar. Parr kemudian diringkas sebagai pembedaan
bahwa biaya aglomerasi lebih terikat pada batas kota besar daripada
manfaat aglomerasinya. (Locator: TXT baris 114-119 dan 143-163, bagian
``INTRODUCTION'', PDF hlm. 1-2.)

\textbf{Laporan sumber:} Artikel juga membedakan \emph{borrowed size},
\emph{borrowed function}, dan \emph{borrowed performance} berdasarkan
literatur yang dikutip. \emph{Borrowed function} menekankan akses kota
terhadap fungsi tingkat tinggi kota lain, sedangkan \emph{borrowed
performance} menekankan kinerja kota yang lebih baik karena dukungan
kota di sekitarnya. Lawannya adalah \emph{agglomeration shadow}, yaitu
ketika kota yang lebih besar menaungi kota tetangga sehingga fungsi kota
tetangga lebih rendah daripada yang diharapkan dari ukurannya. (Locator:
TXT baris 157-179, bagian ``INTRODUCTION'', PDF hlm. 2.)

\textbf{Laporan sumber:} Teori jaringan kota dipakai untuk menjelaskan
bahwa jaringan dibangun oleh koneksi transportasi serta teknologi
informasi dan komunikasi, interaksi fungsional, komplementaritas, dan
sinergi. Interaksi antarfirma memungkinkan pembagian kerja spasial dalam
PUR. Melalui jaringan, kota-kota dapat mengeksploitasi ekonomi
aglomerasi melalui hubungan komplementer dan aktivitas kooperatif
sehingga mencapai massa kritis dalam bentuk eksternalitas jaringan kota.
(Locator: TXT baris 190-220, bagian ``Literature review'', PDF hlm.
2-3.)

\textbf{Laporan sumber:} Artikel merangkum literatur empiris sebagai
belum memadai untuk menilai spillover fungsional regional dan manfaat
ekonomi PUR secara keseluruhan. Studi pada kota individual dinilai tidak
cukup untuk menunjukkan arti eksternalitas jaringan bagi wilayah.
Pengukuran fungsi dan kinerja pada tingkat kota, jaringan kota, dan
wilayah perkotaan dipandang diperlukan. Artikel menyebut temuan bahwa
korelasi antara ukuran dan fungsi kota lebih rendah dalam PUR
dibandingkan wilayah monosen-trik, tetapi menyatakan bahwa isu tersebut
belum diuji secara langsung dengan pemeriksaan yang ketat. (Locator: TXT
baris 222-260, bagian ``Literature review'', PDF hlm. 3.)

\textbf{Interpretasi penulis:} Untuk konteks China, meningkatnya
polycentricity diperkirakan menaikkan ukuran relatif kota kecil.
Manufaktur yang sensitif terhadap biaya dapat berpindah ke kota kecil,
sementara jasa produsen cenderung tetap mengelompok di kota besar.
Pembagian kerja antara manufaktur di kota kecil dan jasa produsen di
kota besar diperkirakan memperkuat permintaan, kapasitas penyediaan
jasa, \emph{borrowed size}, dan interaksi fungsional, sekaligus
mengurangi \emph{agglomeration shadow}. (Locator: TXT baris 230-266 dan
243-260, PDF hlm. 3-4.)

\textbf{Laporan sumber:} Sektor jasa produsen dalam artikel mencakup
transportasi, pergudangan dan layanan pos; transmisi informasi, layanan
komputer dan industri perangkat lunak; jasa keuangan; industri real
estat; leasing dan jasa bisnis; serta penelitian ilmiah, layanan teknis,
dan eksplorasi geologi. (Locator: TXT baris 317-333, bagian ``The
measurement of the functional spillovers of a city cluster in the
producer service sector'', PDF hlm. 4.)

\textbf{Batas tinjauan:} Artikel bukan artikel tinjauan sistematis.
Metode pencarian literatur, kriteria inklusi, jumlah corpus literatur,
dan prosedur sintesis literatur tidak disebutkan dalam file. Ringkasan
di atas hanya melaporkan cara artikel menyajikan literatur yang
dirujuknya, bukan verifikasi independen atas setiap rujukan.

\subsection{Method Used}\label{method-used-26}

\textbf{Laporan sumber:} Strategi estimasi terdiri atas dua tahap
substantif. Tahap pertama menguji hubungan polycentricity dengan
spillover fungsional dan hubungan spillover fungsional dengan
produktivitas. Tahap kedua menguji mekanisme pembagian kerja antarkota.
Karena spillover fungsional menjadi variabel dependen pada persamaan
spillover dan variabel penjelas pada persamaan produktivitas, artikel
memperlakukannya sebagai endogen dan menggunakan model persamaan
simultan. Estimasi utama memakai \emph{three-stage least squares}
(3SLS), sedangkan OLS dan 2SLS dipertahankan sebagai pembanding.
(Locator: TXT baris 364-420, bagian ``Estimation model for the effect of
polycentricity on functional spillovers and economic performance'', PDF
hlm. 5-6.)

\textbf{Laporan sumber:} Polycentricity morfologis diukur dari
distribusi rank-size penduduk kota dalam klaster dengan persamaan
\texttt{ln(Rank\ -\ 0.5)\ =\ a\ +\ Pareto\ x\ ln(Size)\ +\ error}.
Pengurangan 0.5 pada rank dimaksudkan untuk mengurangi bias ukuran
sampel potensial pada eksponen Pareto. Nilai absolut koefisien Pareto
dipakai sebagai indeks; nilai yang lebih tinggi diartikan sebagai
klaster yang lebih polisentris. (Locator: TXT baris 290-313, persamaan
(1), PDF hlm. 4.)

\textbf{Laporan sumber:} Spillover fungsional sektor jasa produsen
dibangun melalui empat langkah:

\begin{itemize}
\tightlist
\item
  Pangsa tenaga kerja jasa produsen tiap kota dibandingkan dengan total
  tenaga kerja jasa produsen klaster digunakan sebagai tingkat aktual
  fungsi jasa produsen kota.
\item
  Pangsa penduduk kota dalam klaster diregresikan terhadap pangsa jasa
  produsen yang diharapkan, dengan kontrol atas pangsa tenaga kerja
  manufaktur, modal manusia, inovasi, intervensi pemerintah, dan
  informatization.
\item
  Residual regresi dianggap sebagai selisih antara fungsi aktual dan
  fungsi yang diharapkan dari ukuran penduduk. Residual positif
  ditafsirkan sebagai pusat jasa produsen yang menikmati \emph{borrowed
  size}; residual negatif sebagai \emph{agglomeration shadow}.
\item
  Residual kota-kota dalam setiap klaster dijumlahkan menjadi spillover
  fungsional klaster. Jumlah positif diartikan sebagai \emph{borrowed
  size} yang lebih besar daripada \emph{agglomeration shadow}.
\end{itemize}

(Laporan sumber; locator: TXT baris 334-363 dan 387-431, bagian
pengukuran spillover fungsional, PDF hlm. 4-6.)

\textbf{Laporan sumber:} Persamaan (2) menjelaskan \texttt{lnFS} dengan
polycentricity, populasi, konektivitas infrastruktur, jumlah provinsi,
pengguna internet, diversitas industri, dummy wilayah central dan
western, serta interaksi polycentricity dengan populasi dan
konektivitas. Persamaan (3) menjelaskan GDP per pekerja dengan spillover
fungsional, modal manusia, investasi, populasi, dan dummy wilayah. Tahun
dikontrol dalam kedua persamaan. (Locator: TXT baris 422-462, persamaan
(2)-(3), PDF hlm. 5-6.)

\textbf{Laporan sumber:} Konektivitas infrastruktur diukur dengan
frekuensi rata-rata kereta pada kota-kota dalam klaster. Penulis
menyatakan bahwa frekuensi kereta yang lebih tinggi dapat meningkatkan
pemanfaatan infrastruktur, memperbaiki konektivitas, dan menurunkan
biaya transaksi antarkota. Jumlah provinsi digunakan untuk menangkap
kemungkinan friksi batas provinsi, sedangkan pengguna internet per
kapita merefleksikan komunikasi antarkota. (Locator: TXT baris 463-505,
bagian penjelasan variabel dan persamaan (4), PDF hlm. 6.)

\textbf{Catatan reproduksibilitas:} Rumus persamaan (4) hadir dalam TXT,
tetapi simbol penjumlahan dan tata letak penyebutnya terdistorsi oleh
ekstraksi \texttt{pdftotext\ -layout}. Review ini tidak merekonstruksi
rumus tersebut; definisi verbal frekuensi kereta rata-rata digunakan
sebagai bukti yang dapat dibaca. (Locator: TXT baris 487-497, PDF hlm.
6.)

\textbf{Laporan sumber:} Untuk mekanisme, pembagian kerja diukur dari
ketidakmiripan struktur industri kota. Grey relational analysis
menghasilkan \emph{grey relational grade} (GRG), yang lebih tinggi
berarti struktur industri kota lebih mirip dengan kota referensi.
Koefisien variasi (CV) GRG antarkota dipakai sebagai indikator pembagian
kerja; CV yang lebih tinggi berarti struktur industri kota lebih tidak
mirip. Persamaan (6) menguji polycentricity terhadap CV, sedangkan
persamaan (7) menguji CV terhadap spillover fungsional dengan kontrol
dan interaksi. Model mekanisme juga diestimasi menggunakan persamaan
simultan 3SLS. (Locator: TXT baris 498-526 dan 556-575, bagian
``Estimation model for the effect of polycentricity on functional
spillovers through labour division among cities'', PDF hlm. 6-7.)

\textbf{Laporan sumber:} Variabel dalam interaksi dipusatkan pada nilai
rata-rata sebelum interaksi dihitung agar hasil lebih mudah ditafsirkan.
Penulis tidak mengontrol efek tetap klaster kota karena variasi utama
polycentricity dan spillover berasal dari variasi antarklaster, bukan
variasi dalam klaster sepanjang waktu. (Locator: TXT baris 599-603,
659-661, dan 478-492, PDF hlm. 6 dan 8-9.)

\textbf{Tidak disebutkan dalam file:} Teks yang tersedia tidak
menjelaskan instrumen spesifik untuk 2SLS/3SLS, bentuk lengkap prosedur
identifikasi, atau rincian pengujian diagnostik lain di luar keterangan
over-identification dan pengujian residual yang dilaporkan.

\subsection{Dataset}\label{dataset-26}

\textbf{Laporan sumber:} Dataset merupakan panel klaster kota China
untuk 2004-2014 dengan jeda setiap dua tahun. Sumber data yang
disebutkan adalah sebagai berikut:

\begin{itemize}
\tightlist
\item
  Data penduduk tiap kota berasal dari LandScan.
\item
  Frekuensi kereta tiap kota dalam klaster dikumpulkan dari jadwal
  kereta China yang diterbitkan China Railway Corporation.
\item
  Jumlah pengguna internet, jumlah pekerja, GDP, investasi total pada
  aset tetap, jumlah mahasiswa, dan konsumsi pemerintah dikumpulkan dari
  \emph{China city statistical yearbook}.
\end{itemize}

(Locator: TXT baris 273-287, bagian ``Research areas and data'', PDF
hlm. 4.)

Data tenaga kerja sektor jasa produsen dan manufaktur dipakai dalam
pembentukan indikator spillover dan pembagian kerja. Variabel kontrol
mencakup modal manusia, inovasi, intervensi pemerintah, informatization,
diversitas industri, jumlah provinsi, serta dummy wilayah. (Locator: TXT
baris 357-387, 504-519, dan 540-553, PDF hlm. 5-7.)

\textbf{Statistik deskriptif yang dilaporkan:} Untuk 78 observasi,
\texttt{lnFS} memiliki mean 0.074, simpangan baku 0.136, minimum -0.125,
dan maksimum 0.324. \texttt{lnGDP} atau GDP per pekerja memiliki mean
2.453, simpangan baku 0.376, minimum 1.510, dan maksimum 3.320, dalam
satuan 10.000 yuan per pekerja. \texttt{lnPOP} memiliki mean 7.918 dan
rentang 7.103-9.240; \texttt{lnPARETO} mean 0.850 dan rentang
0.494-1.231; dan \texttt{lnINF\_CONNECT} mean 3.915 dan rentang
0.336-6.237. (Laporan sumber; locator: TXT baris 533-545 dan 649-651,
Table 2 dan Table 4, PDF hlm. 7-8.)

\texttt{lnINTERNET} memiliki mean 0.150 dan rentang 0.038-0.523;
\texttt{ADMIN} mean 1.385 dan rentang 1-3; dummy \texttt{WESTERN} mean
0.154; dummy \texttt{CENTRAL} mean 0.231; \texttt{lnDIVERSITY} mean
1.524 dan rentang 0.190-3.921; \texttt{lnHUMAN} mean 6.567 dan rentang
5.458-7.416; serta \texttt{lnINVEST} mean 2.144 dan rentang 0.762-3.066.
Untuk analisis mekanisme, \texttt{lnGOV} memiliki mean 0.084, minimum
0.042, maksimum 0.152, dan simpangan baku 0.024; \texttt{lnCV} memiliki
mean 0.226, minimum 0.067, maksimum 0.390, dan simpangan baku 0.067.
(Laporan sumber; locator: TXT baris 546-553 dan 556-579, Table 2 serta
uraian variabel mekanisme, PDF hlm. 7.)

\textbf{Inferensi pengolahan:} Dataset memungkinkan pengujian panel pada
tingkat klaster-tahun dan pembentukan indikator berbasis kota di dalam
tiap klaster, tetapi TXT tidak menunjukkan data arus aktual antarkota.
Karena itu, konektivitas kereta dan residual fungsi menjadi proksi utama
untuk interaksi serta \emph{borrowed size}, bukan pengukuran langsung
seluruh arus ekonomi atau mobilitas.

\subsection{Population / Sample}\label{population-sample-26}

\textbf{Laporan sumber:} Sampel terdiri atas 13 klaster kota China yang
disebut relatif terkenal di China dan diidentifikasi oleh Ning (2016).
Table 1 melaporkan jumlah kota penyusun, populasi tahun 2014 dalam
ribuan penduduk, dan indeks polycentricity tahun 2014 sebagai berikut:

\begin{itemize}
\tightlist
\item
  Central Plains: 5 kota; populasi 29,576.292 (satuan x1000);
  polycentricity 2.422.
\item
  Changchun-Jilin: 2 kota; populasi 12,355.282 (satuan x1000);
  polycentricity 1.971.
\item
  Chang-zhu-tan: 3 kota; populasi 13,864.109 (satuan x1000);
  polycentricity 1.717.
\item
  Mindongnan: 5 kota; populasi 26,039.169 (satuan x1000); polycentricity
  1.662.
\item
  Shandong Peninsula: 6 kota; populasi 39,217.211 (satuan x1000);
  polycentricity 1.556.
\item
  Liaozhongnan: 7 kota; populasi 26,871.16 (satuan x1000);
  polycentricity 1.393.
\item
  Ha-da-qi: 3 kota; populasi 19,211.908 (satuan x1000); polycentricity
  1.245.
\item
  Pearl River Delta: 9 kota; populasi 56,841.152 (satuan x1000);
  polycentricity 1.216.
\item
  Yangtze River Delta: 15 kota; populasi 102,938.639 (satuan x1000);
  polycentricity 1.190.
\item
  Jing-jin-ji: 5 kota; populasi 48,242.533 (satuan x1000);
  polycentricity 1.074.
\item
  Chengdu-Chongqing: 5 kota; populasi 53,875.884 (satuan x1000);
  polycentricity 0.815.
\item
  Guanzhong: 4 kota; populasi 18,412.486 (satuan x1000); polycentricity
  0.722.
\item
  Wuhan: 3 kota; populasi 13,495.167 (satuan x1000); polycentricity
  0.721.
\end{itemize}

(Locator: TXT baris 273-287 dan 344-354, bagian ``Research areas and
data'' dan Table 1, PDF hlm. 4-5.)

\textbf{Laporan sumber:} Persamaan menggunakan \texttt{i} untuk klaster
kota dan \texttt{t} untuk tahun. Dengan 13 klaster dan periode
pengamatan yang berjarak dua tahun, artikel melaporkan 78 observasi pada
Table 4 dan Table 5. Data kota penyusun digunakan untuk membentuk
residual fungsi jasa produsen dan CV struktur industri sebelum
diagregasikan atau dimasukkan ke persamaan tingkat klaster. (Locator:
TXT baris 447-452, 556-579, 649-651, dan 774-779, PDF hlm. 6-10.)

\textbf{Tidak disebutkan dalam file:} Nama seluruh kota penyusun tiap
klaster, kriteria administratif atau fungsional untuk menetapkan batas
klaster, prosedur pemilihan 13 klaster selain rujukan pada Ning (2016),
penanganan data hilang, pembobotan observasi, dan representativitas
sampel terhadap seluruh klaster kota China tidak disebutkan dalam file.

\subsection{Dependent Variables}\label{dependent-variables-26}

\textbf{Spillover fungsional regional (\texttt{lnFS}):} Variabel ini
merupakan jumlah residual kota-kota penyusun dalam tiap klaster untuk
fungsi sektor jasa produsen. Residual positif pada kota diperlakukan
sebagai indikasi fungsi jasa produsen yang melampaui permintaan
berdasarkan ukuran penduduk dan sebagai \emph{borrowed size}; residual
negatif diperlakukan sebagai \emph{agglomeration shadow}. Jumlah
residual klaster menjadi ukuran trade-off antara kedua efek. Table 2
melaporkan mean 0.074, simpangan baku 0.136, minimum -0.125, dan
maksimum 0.324. (Laporan sumber; locator: TXT baris 357-431 dan 533-538,
PDF hlm. 5 dan 7.)

\textbf{Produktivitas tenaga kerja (\texttt{lnGDP} atau GDP per
worker):} Variabel dependen pada persamaan (3) adalah GDP per pekerja.
Table 2 menyatakan satuannya 10.000 yuan per pekerja, dengan mean 2.453,
simpangan baku 0.376, minimum 1.510, dan maksimum 3.320. Artikel
mengharapkan hubungan positif antara spillover fungsional dan
produktivitas tenaga kerja. (Laporan sumber; locator: TXT baris 439-462
dan 533-539, persamaan (3) dan Table 2, PDF hlm. 6-7.)

\textbf{Pembagian kerja (\texttt{lnCV}) dalam analisis mekanisme:} Pada
persamaan (6), variabel yang dijelaskan adalah CV dari GRG struktur
industri antarkota. CV lebih tinggi diartikan sebagai struktur industri
yang lebih tidak mirip dan, karenanya, pembagian kerja antarkota yang
lebih besar. Table 5 melaporkan mean \texttt{lnCV} 0.226, minimum 0.067,
maksimum 0.390, dan simpangan baku 0.067. (Laporan sumber; locator: TXT
baris 498-520 dan 556-579, PDF hlm. 6-7.)

Pada persamaan (7), \texttt{lnFS} kembali menjadi variabel dependen,
sedangkan \texttt{lnCV} menjadi variabel penjelas utama untuk menguji
apakah pembagian kerja berhubungan dengan spillover fungsional. (Laporan
sumber; locator: TXT baris 556-575, persamaan (6)-(7), PDF hlm. 7.)

\textbf{Catatan pengukuran:} Notasi \texttt{ln} digunakan dalam
persamaan dan tabel, tetapi definisi substantif \texttt{FS} yang
dijelaskan penulis adalah jumlah residual kota. TXT tidak memberikan
uraian tambahan tentang transformasi log untuk setiap variabel di luar
notasi persamaan dan definisi Table 2; review ini mempertahankan notasi
artikel tanpa mengasumsikan transformasi yang tidak dijelaskan.

\subsection{Results}\label{results-26}

\textbf{Hasil hubungan polycentricity dan spillover fungsional:} Pada
model tanpa interaksi, koefisien polycentricity dalam 3SLS adalah 0.017
dengan statistik 0.29 dan tidak menunjukkan hubungan langsung yang
signifikan dalam tabel. OLS dan 2SLS sama-sama melaporkan koefisien
0.082 dengan statistik 1.23. Setelah interaksi dengan populasi dan
konektivitas dimasukkan, koefisien utama polycentricity dalam model 3SLS
menjadi 0.265 dengan \texttt{***} dan statistik 3.38. Koefisien
interaksi polycentricity-populasi adalah 0.281 dengan \texttt{**} dan
statistik 2.15, sedangkan interaksi polycentricity-konektivitas adalah
0.127 dengan \texttt{**} dan statistik 2.08. (Laporan sumber; locator:
TXT baris 599-632, Table 4, PDF hlm. 8.)

\textbf{Interpretasi penulis:} Polycentricity tidak cukup untuk
menghasilkan spillover fungsional jika efek moderasi tidak
diperhitungkan. Pada nilai populasi dan konektivitas rata-rata, struktur
polisentris dapat meningkatkan spillover; populasi yang lebih besar
memperluas peluang interaksi, sedangkan konektivitas yang lebih baik
menurunkan friksi jarak dan biaya transaksi. (Locator: TXT baris
659-690, bagian ``Spatial structure, regional functional spillovers and
regional labour productivity'', PDF hlm. 9.)

\textbf{Hasil ambang marginal:} Berdasarkan Figure 1 yang menggunakan
model 4 pada Table 4, penulis melaporkan bahwa polycentricity lebih
kondusif bagi spillover fungsional hanya pada klaster berpenduduk lebih
dari 18.08 juta orang dan dengan konektivitas lebih dari 24 kereta per
hari. (Laporan sumber dan interpretasi penulis; locator: TXT baris
691-703, Figure 1, PDF hlm. 9.)

\textbf{Hasil spillover dan produktivitas:} Koefisien \texttt{FS} pada
persamaan GDP per pekerja positif pada semua model di Table 4. Dalam
model 3SLS tanpa moderasi, koefisiennya 0.571 dengan \texttt{***} dan
statistik 4.25. Dalam model 3SLS dengan moderasi, koefisiennya 0.408
dengan \texttt{***} dan statistik 3.46. Koefisien OLS masing-masing
0.318 dengan \texttt{***}, sedangkan 2SLS masing-masing 0.570 dan 0.405
dengan tanda signifikansi yang dilaporkan. (Laporan sumber; locator: TXT
baris 633-650, Table 4, PDF hlm. 8.)

\textbf{Hasil variabel kontrol:} Populasi berkoefisien positif pada
persamaan spillover dan produktivitas. Jumlah provinsi (\texttt{ADMIN})
berkoefisien negatif pada persamaan spillover. Dummy central dan western
bernilai negatif pada spesifikasi yang ditampilkan, dan artikel
menginterpretasikan klaster central serta western memiliki spillover
lebih sedikit daripada klaster eastern. Modal manusia dan investasi
berkoefisien positif pada produktivitas. (Laporan sumber dan
interpretasi penulis; locator: TXT baris 611-624, 636-645, dan 663-683,
Table 4 serta uraian hasil, PDF hlm. 8-9.)

\textbf{Hasil uji model simultan:} Penulis melaporkan bahwa residual
pada persamaan (2) berpengaruh signifikan terhadap persamaan (3) pada
tingkat 10 persen dengan \texttt{p\ =\ 0.058}, sehingga penggunaan model
persamaan simultan dianggap perlu. Hasil 3SLS disebut konsisten dengan
2SLS; OLS dipertahankan sebagai pembanding. (Laporan sumber dan
interpretasi penulis; locator: TXT baris 464-476 dan 659-663, PDF hlm. 6
dan 9.)

\textbf{Hasil mekanisme pembagian kerja:} Pada persamaan (6), model
dengan interaksi melaporkan koefisien polycentricity 0.198 dengan
\texttt{***} dan statistik 4.19. Interaksi polycentricity-populasi
bernilai 0.369 dengan \texttt{***} dan statistik 4.01. Koefisien utama
populasi bernilai -0.032 dengan \texttt{**} dan statistik -2.11.
Interaksi polycentricity-konektivitas bernilai -0.052 tanpa tanda
signifikansi dalam tabel. Pada persamaan (7), koefisien CV adalah 0.817
dengan \texttt{**} dan statistik 2.24 pada model dengan interaksi,
sementara interaksi CV-konektivitas adalah 0.614 dengan \texttt{***} dan
statistik 3.40; interaksi CV-populasi bernilai -0.006 tanpa tanda
signifikansi. (Laporan sumber; locator: TXT baris 727-779, Table 5, PDF
hlm. 10.)

\textbf{Interpretasi penulis:} Perbandingan model Table 5 dipakai untuk
menyatakan bahwa pertumbuhan ukuran pasar atau populasi, bersama dengan
polycentricity, mendorong ketidakmiripan industri antarkota dan
memperdalam pembagian kerja. Pembagian kerja meningkatkan spillover
fungsional ketika infrastruktur lebih terhubung. (Locator: TXT baris
710-722, bagian ``Further analysis'', PDF hlm. 10.)

\textbf{Kinerja model yang dilaporkan:} Table 4 memiliki 78 observasi
pada setiap model; R2 untuk persamaan spillover berada pada 0.614-0.693
dan R2 untuk persamaan GDP per pekerja pada 0.923-0.929. Table 5
memiliki 78 observasi; R2 model pembagian kerja adalah 0.2119 dan
0.3557, sedangkan R2 persamaan spillover mekanisme adalah 0.0899 dan
0.5515. (Laporan sumber; locator: TXT baris 629-650 dan 752-779, Table
4-5, PDF hlm. 8 dan 10.)

\subsection{Findings}\label{findings-26}

\textbf{Temuan yang dilaporkan:}

\begin{itemize}
\tightlist
\item
  Polycentricity tidak mempunyai hubungan langsung yang signifikan
  dengan spillover fungsional pada spesifikasi 3SLS tanpa moderator.
\item
  Hubungan positif muncul pada model yang memperhitungkan interaksi
  dengan populasi dan konektivitas; kedua interaksi tersebut bertanda
  positif dan signifikan dalam model 3SLS yang dilaporkan.
\item
  Artikel memberikan ambang marginal lebih dari 18.08 juta penduduk dan
  lebih dari 24 kereta per hari untuk kondisi ketika polycentricity
  lebih mendukung spillover fungsional.
\item
  Spillover fungsional berhubungan positif dengan GDP per pekerja pada
  semua model yang ditampilkan.
\item
  Pembagian kerja antarkota, yang diproksikan dengan ketidakmiripan
  struktur industri, dipakai sebagai mekanisme yang menghubungkan
  polycentricity dengan spillover; hubungan pembagian kerja-spillover
  menjadi lebih kuat ketika konektivitas infrastruktur lebih baik.
\item
  Batas provinsi berhubungan negatif dengan spillover fungsional dalam
  hasil yang diringkas penulis.
\end{itemize}

(Locator: TXT baris 667-690, 691-703, 710-722, dan 717-737, PDF hlm.
9-10.)

\textbf{Interpretasi penulis:} Temuan tersebut dianggap mendukung kedua
hipotesis dan menunjukkan bahwa pembagian kerja antarkota dapat menjadi
kekuatan utama efek \emph{borrowed size} untuk fungsi jasa produsen.
(Locator: TXT baris 717-737, bagian ``CONCLUSIONS'', PDF hlm. 10.)

\textbf{Inferensi pengolahan:} Bukti paling kuat dalam TXT adalah bukti
hubungan kondisional pada sampel klaster kota China, bukan bukti efek
universal atau bukti bahwa PUR secara keseluruhan selalu lebih unggul
daripada wilayah monosen-trik. Kesimpulan kausal yang lebih luas perlu
ditahan karena ukuran spillover dibentuk dari residual, instrumen 3SLS
tidak dirinci, dan artikel sendiri menyatakan bahwa perbandingan besar
dengan wilayah monosen-trik tidak dapat diuji.

\subsection{Conclusions}\label{conclusions-26}

\textbf{Interpretasi penulis:} Dengan menggunakan 13 klaster kota China
dan 3SLS untuk hubungan simultan, artikel menyimpulkan bahwa
polycentricity lebih bermanfaat bagi spillover fungsional regional
ketika klaster memiliki populasi lebih besar dan konektivitas
infrastruktur yang lebih baik. Spillover tersebut berkaitan dengan
kenaikan produktivitas tenaga kerja. Pembagian kerja antara sektor jasa
produsen dan manufaktur dipandang memainkan peran positif dalam pengaruh
polycentricity terhadap spillover, serta menjadi kekuatan penting efek
\emph{borrowed size} setidaknya untuk fungsi jasa produsen. (Locator:
TXT baris 710-737, bagian ``CONCLUSIONS'', PDF hlm. 10.)

Penulis menyatakan bahwa temuan mendukung hipotesis mereka, tetapi
membatasi klaim pada konteks klaster kota China. Mereka tidak menyatakan
bahwa penelitian ini telah membuktikan keunggulan PUR atas wilayah
monosen-trik secara menyeluruh. (Locator: TXT baris 208-220 dan 717-737,
PDF hlm. 3 dan 10.)

\textbf{Inferensi pengolahan:} Rumusan kesimpulan yang paling aman
adalah bahwa konfigurasi polisentris dapat menjadi produktif apabila
ditopang massa penduduk, konektivitas, dan diferensiasi fungsi
antarkota. Variabel-variabel tersebut merupakan kondisi analitis dalam
model, bukan jaminan hasil perencanaan di luar populasi, periode,
sektor, dan negara yang diteliti.

\subsection{Contributions}\label{contributions-26}

\textbf{Inferensi pengolahan berdasarkan gap, metode, dan hasil
artikel:} Sumbangan yang dapat diidentifikasi dari artikel mencakup
beberapa hal:

\begin{itemize}
\tightlist
\item
  Mengalihkan perhatian dari pengukuran \emph{borrowed size} atau
  \emph{agglomeration shadow} pada kota individual ke spillover
  fungsional pada tingkat klaster dan dampaknya terhadap produktivitas.
\item
  Menghubungkan polycentricity morfologis, spillover fungsional sektor
  jasa produsen, dan GDP per pekerja dalam satu kerangka persamaan
  simultan.
\item
  Menunjukkan bahwa hubungan polycentricity-spillover dimoderasi oleh
  ukuran populasi dan konektivitas infrastruktur, bukan bersifat
  konstan.
\item
  Menguji pembagian kerja antarkota sebagai mekanisme yang menghubungkan
  polycentricity dengan spillover fungsional.
\item
  Memberikan bukti dari 13 klaster kota China dengan panel 2004-2014.
\end{itemize}

(Locator: TXT baris 164-179, 208-220, 236-260, 330-337, dan 710-737, PDF
hlm. 2-4 dan 10.)

\textbf{Inferensi pengolahan:} Kontribusi artikel terutama bersifat
empiris-operasional dan kontekstual. Artikel memperjelas kondisi yang
menyertai efek \emph{borrowed size}, tetapi tidak menyelesaikan
perdebatan umum tentang apakah polycentricity lebih baik daripada
monocentricity, karena hipotesis perbandingan tersebut secara eksplisit
tidak diuji.

\subsection{Research Gap}\label{research-gap-26}

\textbf{Gap yang dinyatakan artikel:}

\begin{itemize}
\tightlist
\item
  Bukti empiris tentang manfaat ekonomi PUR melalui spillover fungsional
  regional dinilai masih sedikit dibandingkan literatur teoritis.
\item
  Studi terdahulu lebih banyak mengamati kota individual sehingga belum
  cukup menjelaskan eksternalitas jaringan pada tingkat PUR sebagai
  keseluruhan.
\item
  Mekanisme yang menghubungkan struktur jaringan kota dengan spillover
  fungsional belum dieksplorasi secara memadai.
\item
  Pengukuran fungsi dan kinerja pada tingkat kota, jaringan kota, dan
  wilayah diperlukan agar manfaat regional tidak disimpulkan dari
  performa satu kota saja.
\item
  Karena keterbatasan data dan metode, perbandingan besar antara
  keunggulan ekonomi PUR dan ekonomi aglomerasi wilayah monosen-trik
  belum diuji.
\end{itemize}

(Locator: TXT baris 164-174, 208-220, dan 236-260, bagian
``INTRODUCTION'' serta ``Literature review'', PDF hlm. 2-3.)

\textbf{Gap yang masih terbuka menurut artikel:} Generalisasi ke ekonomi
maju atau pasca-industri belum diketahui; pengukuran residual belum sama
dengan ukuran aktual \emph{borrowed size}; dan mekanisme lain di luar
pembagian kerja jasa produsen-manufaktur mungkin turut menjelaskan
spillover atau kinerja regional. (Locator: TXT baris 750-779, bagian
akhir ``CONCLUSIONS'', PDF hlm. 10-11.)

\textbf{Inferensi pengolahan:} Untuk mengisi gap tersebut, bukti
lanjutan perlu mempertahankan desain kondisional dan menguji apakah
ambang serta mekanisme yang dilaporkan tetap berlaku pada konteks,
sektor, ukuran wilayah, dan bentuk jaringan yang berbeda. Ini adalah
inferensi review, bukan agenda eksplisit selain yang tercantum pada
bagian future research artikel.

\subsection{Limitations}\label{limitations-26}

\textbf{Keterbatasan yang dinyatakan penulis:}

\begin{itemize}
\tightlist
\item
  Artikel tidak dapat menguji hipotesis besar tentang apakah kekuatan
  ekonomi PUR mengimbangi atau melampaui ekonomi aglomerasi wilayah
  monosen-trik karena keterbatasan data dan metode.
\item
  Mekanisme promosi polycentricity diturunkan dari konteks China sebagai
  negara berkembang/emerging country. Kota kecil dalam wilayah
  monosen-trik China dijelaskan memiliki daya tarik lebih rendah bagi
  manufaktur karena keterbatasan tenaga kerja, skala pasar, dan
  infrastruktur. Generalisasi ke ekonomi maju atau pasca-industri masih
  perlu diverifikasi.
\item
  Metode residual untuk mengukur \emph{borrowed size} tidak sempurna.
  Masalah variabel yang dihilangkan membuat residual hanya mendekati
  derajat \emph{borrowed size}, bukan derajat aktualnya, walaupun
  variabel kontrol telah digunakan.
\item
  Data arus akan lebih tepat untuk mengukur \emph{borrowed size} apabila
  tersedia, karena arus merepresentasikan proses yang benar-benar
  berlangsung dalam \emph{borrowed size}.
\item
  Penelitian hanya menganalisis mekanisme pembagian kerja antara jasa
  produsen dan manufaktur; mekanisme lain yang dapat meningkatkan
  spillover atau kinerja jasa produsen mungkin ada.
\end{itemize}

(Locator: TXT baris 208-220 dan 750-779, bagian ``INTRODUCTION'' dan
``CONCLUSIONS'', PDF hlm. 3 dan 10-11.)

\textbf{Keterbatasan yang diinferensikan dari desain:}

\begin{itemize}
\tightlist
\item
  Sampel terdiri atas 13 klaster yang disebut relatif terkenal, dengan
  78 observasi dan interval waktu dua tahun. TXT tidak memberi dasar
  untuk menganggapnya sebagai sampel acak atau representatif dari semua
  klaster kota China.
\item
  Penulis tidak mengontrol efek tetap klaster karena variasi utama
  berada antarklaster. Inferensinya, perbedaan tetap antarklaster yang
  tidak teramati tidak dihilangkan oleh efek tetap klaster dalam
  spesifikasi yang dijelaskan. Ini adalah konsekuensi desain yang perlu
  diperhitungkan, bukan klaim bahwa estimasinya pasti bias.
\item
  Polycentricity diukur secara morfologis, sedangkan spillover diukur
  melalui residual fungsi jasa produsen dan konektivitas melalui
  frekuensi kereta. Ketiganya tidak identik dengan pengukuran langsung
  struktur fungsional, arus, atau seluruh dimensi jaringan kota.
\item
  Tidak adanya rincian instrumen 2SLS/3SLS dalam TXT membatasi penilaian
  eksternal atas strategi identifikasi simultan, walaupun artikel
  melaporkan over-identification dan uji residual sebagai dukungan
  penggunaan model tersebut.
\end{itemize}

\textbf{Tidak disebutkan dalam file:} Perlakuan terhadap missing data,
uji sensitivitas terhadap definisi alternatif polycentricity atau
spillover, pengujian autokorelasi/heteroskedastisitas, rincian
instrumen, dan dasar pemilihan bentuk kontrol tidak dijelaskan secara
lengkap dalam file.

\subsection{Challenges}\label{challenges-26}

\textbf{Tantangan substantif yang dihadapi studi:}

\begin{itemize}
\tightlist
\item
  Membedakan struktur morfologis dari fungsi ekonomi. Pareto mengukur
  distribusi ukuran penduduk, sedangkan spillover dibangun dari fungsi
  jasa produsen; hubungan keduanya harus dijembatani secara teoritis dan
  statistik.
\item
  Mengubah konsep \emph{borrowed size} dan \emph{agglomeration shadow}
  menjadi residual yang dapat diagregasikan. Residual harus dibaca
  sebagai selisih fungsi aktual dan fungsi yang diharapkan, bukan
  sebagai observasi arus yang langsung terjadi.
\item
  Menangani kemungkinan endogenitas karena spillover menjadi outcome
  dalam satu persamaan dan prediktor pada persamaan produktivitas.
  Artikel meresponsnya dengan persamaan simultan 3SLS, tetapi rincian
  instrumen tidak tersedia dalam TXT.
\item
  Menafsirkan interaksi. Koefisien utama polycentricity yang tidak
  signifikan tanpa moderator tidak boleh dibaca terpisah dari interaksi
  populasi dan konektivitas; penulis memusatkan variabel sebelum
  membentuk interaksi.
\item
  Mengukur konektivitas dengan frekuensi kereta. Ukuran ini tersedia
  untuk klaster, tetapi TXT tidak menunjukkan bahwa ia mencakup seluruh
  bentuk transportasi, komunikasi, atau arus antarfirma.
\item
  Menghubungkan ketidakmiripan industri dengan pembagian kerja. CV GRG
  dipakai sebagai indikator, sehingga interpretasi pembagian kerja
  bergantung pada kesesuaian proksi tersebut dengan pembagian fungsi
  ekonomi antarkota.
\item
  Menjaga batas konteks. Artikel sendiri menempatkan mekanisme dalam
  konteks industrialisasi regional China dan menyatakan bahwa
  penerapannya pada ekonomi pasca-industri perlu diverifikasi.
\end{itemize}

(Laporan sumber dan inferensi pengolahan; locator: TXT baris 290-316,
334-431, 464-526, 659-690, dan 750-779, PDF hlm. 4-7 dan 9-11.)

\textbf{Tantangan teknis pembacaan:} Beberapa persamaan, terutama
persamaan (4), mengalami distorsi tata letak pada TXT. Simbol yang tidak
terbaca tidak direkonstruksi dalam review ini agar tidak mengubah bukti
sumber. (Locator: TXT baris 487-497, PDF hlm. 6.)

\subsection{Practical Implications}\label{practical-implications-26}

\textbf{Interpretasi penulis:} Implikasi kebijakan yang dinyatakan ada
dua. Pertama, konsep perencanaan spasial polisentris lebih sesuai untuk
PUR dengan populasi besar dan konektivitas infrastruktur yang baik,
setidaknya dalam pengertian efek \emph{borrowed size}. Kedua,
rasionalitas ekonomi PUR dapat diperkuat dengan meningkatkan
konektivitas jaringan infrastruktur transportasi antarkota dalam PUR.
(Locator: TXT baris 739-747, bagian ``CONCLUSIONS'', PDF hlm. 10.)

\textbf{Laporan sumber yang relevan bagi kebijakan:} Batas provinsi
memiliki pengaruh negatif terhadap spillover fungsional dalam model, dan
penulis menginterpretasikannya sebagai hambatan administratif yang
mengganggu spillover sektor jasa produsen. Namun, penulis tidak
merumuskan instrumen kebijakan tata kelola lintas-provinsi secara
terpisah. (Locator: TXT baris 667-674, PDF hlm. 9.)

\textbf{Inferensi pengolahan:} Temuan ini menyarankan agar perencanaan
tidak memakai label ``polisentris'' sebagai tujuan yang cukup. Penilaian
perlu memeriksa apakah massa penduduk, jaringan transportasi, dan
pembagian fungsi antarkota memungkinkan interaksi yang mendukung
spillover. Saran ini dibatasi pada logika model artikel dan tidak boleh
dipindahkan langsung ke konteks non-China tanpa pengujian.

\subsection{Applications}\label{applications-26}

\textbf{Laporan sumber:} Artikel menerapkan modelnya pada 13 klaster
kota China dan pada sektor jasa produsen, tetapi tidak melaporkan
implementasi kebijakan, proyek perencanaan, atau evaluasi intervensi
tertentu setelah hasil empiris diperoleh. Penerapan operasional di luar
sampel tidak disebutkan dalam file.

\textbf{Inferensi pengolahan yang dapat diturunkan secara terbatas:}
Metode artikel dapat digunakan sebagai kerangka diagnosis untuk:

\begin{itemize}
\tightlist
\item
  memeriksa apakah perubahan distribusi ukuran penduduk menunjukkan
  polycentricity morfologis;
\item
  mengidentifikasi kota yang fungsi jasa produsennya berada di atas atau
  di bawah tingkat yang diharapkan dari ukuran penduduk melalui
  residual;
\item
  menilai ketidakmiripan struktur industri antarkota sebagai indikasi
  pembagian kerja;
\item
  menguji apakah konektivitas transportasi memperkuat hubungan
  polycentricity, pembagian kerja, dan spillover; serta
\item
  menghubungkan ukuran spillover dengan GDP per pekerja.
\end{itemize}

Ambang lebih dari 18.08 juta penduduk dan lebih dari 24 kereta per hari
dapat dipakai hanya sebagai titik rujukan dari Figure 1, bukan sebagai
standar universal. TXT tidak menyatakan bahwa ambang tersebut telah
divalidasi di luar 13 klaster dan periode penelitian. (Locator: TXT
baris 691-703, 702-703, dan 739-747, Figure 1 serta implikasi kebijakan,
PDF hlm. 9-10.)

\subsection{Future Research
(Optional)}\label{future-research-optional-26}

\textbf{Agenda yang dinyatakan penulis:}

\begin{itemize}
\tightlist
\item
  Memverifikasi apakah kesimpulan dan mekanisme berlaku pada ekonomi
  maju atau pasca-industri, bukan hanya pada konteks China sebagai
  negara berkembang/emerging country.
\item
  Menggunakan data arus jika tersedia, karena data arus dipandang lebih
  tepat untuk mengukur derajat \emph{borrowed size} yang aktual.
\item
  Menguji mekanisme lain di luar pembagian kerja antara sektor jasa
  produsen dan manufaktur yang mungkin meningkatkan spillover fungsional
  atau kinerja jasa produsen pada skala regional.
\end{itemize}

(Locator: TXT baris 750-779, bagian akhir ``CONCLUSIONS'', PDF hlm.
10-11.)

\textbf{Inferensi pengolahan:} Validasi ambang populasi dan frekuensi
kereta pada sampel atau negara lain merupakan tindak lanjut logis dari
hasil moderasi, tetapi tidak dinyatakan sebagai agenda tersendiri oleh
penulis.

\subsection{Provenance Notes}\label{provenance-notes-26}

Review ini hanya menggunakan file
\texttt{source/journal/A11-wang-sun-zhang-2022-do-polycentric-urban-regions-promote-functional-spillovers-and-economic-performance-evidence-from-china-10.1080-00343404.2020.1841147.txt}.
Seluruh isi yang tersedia dibaca, dari TXT baris 1 sampai 889, termasuk
blok \texttt{{[}PDF\ Metadata{]}} pada baris 1-28 dan
\texttt{{[}Extracted\ Text{]}} mulai baris 29. Metadata menyatakan bahwa
sumber PDF asal adalah file dengan prefiks A11 pada folder
\texttt{journal/}, ekstraksi dilakukan pada 2026-08-31 dengan
\texttt{pdftotext\ -layout}, dan PDF memiliki 13 halaman. (Laporan
sumber; locator: TXT baris 1-29.)

Nomor baris dalam review merujuk pada baris TXT tersebut. Rujukan PDF
hlm. mengikuti nomor halaman artikel yang tampak dalam teks hasil
ekstraksi, sedangkan judul bagian, tabel, dan gambar dicantumkan bila
tersedia. Artikel menyebut Table 1 untuk karakteristik 13 klaster, Table
2 untuk statistik deskriptif, Table 3 untuk sumber variasi, Table 4
untuk hubungan polycentricity-spillover-produktivitas, Table 5 untuk
mekanisme pembagian kerja, dan Figure 1 untuk efek marginal. (Locator:
TXT baris 344-354, 533-579, 605-651, 702-703, dan 727-779.)

Tidak ada PDF, artikel lain, metadata eksternal, atau rujukan dalam
daftar pustaka yang digunakan sebagai bukti tambahan. Rujukan yang
disebut dalam review hanya dipakai sejauh artikel A11 sendiri menyajikan
atau menggunakannya. Informasi yang tidak tersedia dalam TXT dinyatakan
secara eksplisit; simbol persamaan yang rusak karena ekstraksi tidak
diisi dengan dugaan.

Review membedakan laporan sumber, interpretasi penulis, dan inferensi
pengolahan. Inferensi tentang batas generalisasi, proksi, ketiadaan efek
tetap, keterbatasan identifikasi, dan kegunaan aplikasi bukan kutipan
atau klaim eksplisit penulis; semuanya ditandai sebagai inferensi dan
diturunkan hanya dari desain serta keterangan yang tersedia dalam TXT.

\section{Review A12 - Sun dan Ding
(2016)}\label{review-a12---sun-dan-ding-2016}

Status: Review literatur berbasis seluruh teks TXT hasil ekstraksi PDF.

\subsection{Identitas Artikel}\label{identitas-artikel-3}

\begin{itemize}
\tightlist
\item
  \textbf{Kode:} A12.
\item
  \textbf{Judul:} \emph{Do large cities contribute to economic growth of
  small cities? Evidence from Yangtze River Delta in China}; judul
  Mandarin pada teks ialah
  ``大城市有利于小城市的经济增长吗？------来自长三角城市群的证据''.
\item
  \textbf{Penulis:} Sun Bindong dan Ding Song.
\item
  \textbf{Tahun:} 2016.
\item
  \textbf{Jurnal:} \emph{Geographical Research} (\emph{地理研究}),
  volume 35, nomor 9.
\item
  \textbf{Halaman:} 1615-1625.
\item
  \textbf{DOI:} 10.11821/dlyj201609002.
\item
  \textbf{Afiliasi yang tercantum:} Center for Modern Chinese City
  Studies dan School of Urban and Regional Science, East China Normal
  University, Shanghai, China.
\item
  \textbf{Tanggal naskah:} diterima 23 Februari 2016 dan direvisi 22
  Juni 2016.
\item
  \textbf{Pendanaan:} National Natural Science Foundation of China,
  Shanghai Philosophy and Social Science Planning Project, serta East
  China Normal University Excellent Doctoral Dissertation Cultivation
  Project; nomor proyek tercantum di TXT.
\item
  \textbf{Kata kunci:} kota besar; kota kecil; interaksi spasial;
  limpahan pertumbuhan; agglomeration shadow; Delta Sungai Yangtze.
\item
  \textbf{Locator identitas:} TXT baris 30-62 dan 78-85; judul/abstrak
  Inggris pada baris 548-585.
\end{itemize}

\subsection{Summarized Abstract}\label{summarized-abstract-27}

\textbf{Laporan sumber.} Artikel menanyakan apakah kedekatan dengan kota
besar membantu atau merugikan pertumbuhan ekonomi kota kecil. Penulis
menggunakan 108 kota kecil di Delta Sungai Yangtze dan menambahkan jarak
geografis ke kota besar, hubungan dengan batas administratif, serta
potensi pasar ke dalam model pertumbuhan ekonomi tradisional. Model
tersebut digunakan untuk mengestimasi pengaruh spasial kota-kota besar
dengan tingkat hierarki berbeda terhadap pertumbuhan kota kecil dan
menguji heterogenitasnya. (TXT, hlm. 1615 dan 1625, abstrak Mandarin dan
Inggris, baris 52-60 dan 555-585.)

\textbf{Laporan sumber.} Hasil utama yang dilaporkan ialah bahwa
kedekatan dengan kota besar membantu pertumbuhan ekonomi kota kecil di
Delta Sungai Yangtze dan artikel tidak menemukan bukti langsung
agglomeration shadow. Interaksi spasial lebih banyak berlangsung dari
kota berhierarki tinggi ke kota berhierarki rendah; limpahan pertumbuhan
dari kota subprovinsial disebut paling signifikan, sedangkan hubungan
antarkota kecil dengan hierarki yang sama relatif lemah. Batas
administratif disebut menciptakan segmentasi pasar yang melemahkan
limpahan spasial. (TXT, hlm. 1615 dan 1625, abstrak, baris 52-61 dan
567-583.)

\textbf{Interpretasi penulis.} Penulis membaca hasil tersebut sebagai
dukungan empiris bagi strategi pengembangan metropolitan dan penguatan
fungsi radiasi kota besar, khususnya di wilayah Delta Sungai Yangtze.
Mereka menekankan bahwa kesimpulan tersebut tidak otomatis berlaku bagi
seluruh Tiongkok karena tahap perkembangan, struktur industri, dan
lingkungan kelembagaan antardaerah berbeda. (TXT, hlm. 1616-1617, baris
120-133; hlm. 1622-1623, baris 416-433.)

\textbf{Inferensi pengolahan.} Bukti yang disajikan lebih tepat dibaca
sebagai hubungan statistik bersyarat antara jarak atau potensi pasar dan
laju pertumbuhan selama satu interval 2000-2010. Kedekatan tidak diukur
melalui arus tenaga kerja, perdagangan, pengetahuan, perusahaan, atau
pengguna layanan yang diamati secara langsung. Karena itu, hasil
mendukung adanya pola yang konsisten dengan limpahan pertumbuhan dari
kota besar, tetapi tidak dengan sendirinya mengidentifikasi mekanisme
atau membuktikan bahwa kedekatan menyebabkan pertumbuhan. (Inferensi
dari TXT, hlm. 1617-1618, bagian 2.1-2.2, baris 151-175 dan 178-215.)

\subsection{Problem Statement}\label{problem-statement-27}

\textbf{Laporan sumber.} Strategi pembangunan wilayah Tiongkok
menekankan koordinasi antarkota besar, menengah, kecil, dan kota kecil.
Namun, artikel menunjukkan dua gejala yang tampak bertentangan.
Kota-kota kecil di kawasan pesisir tenggara dapat memperoleh dorongan
pertumbuhan karena dekat dengan pusat utama, sedangkan gejala seperti
``sabuk kemiskinan di sekitar Beijing-Tianjin'' dan ``gelap di bawah
lampu'' tampak menunjukkan persaingan serta penyerapan faktor oleh kota
besar. Pertanyaan yang hendak dijawab ialah apakah kota besar menjadi
``berkah'' atau ``bencana'' bagi pertumbuhan kota kecil. (TXT, hlm.
1615-1616, bagian 1, baris 67-94.)

\textbf{Laporan sumber.} Kerangka teoritis yang dirujuk memberi jawaban
yang tidak tunggal. Teori Myrdal tentang backwash-spread, Hirschman
tentang polarization-trickle-down, dan Friedmann tentang core-periphery
memandang kota besar dapat sekaligus mengonsentrasikan faktor dan
menyebarkan manfaat, dengan hasil yang bergantung pada tahap
perkembangan. New Economic Geography memperkenalkan agglomeration
shadow: kedekatan dengan pusat dapat membuat pendirian perusahaan baru
tidak menguntungkan karena faktor terserap oleh kota besar, sedangkan
kota yang melampaui ambang jarak tertentu dapat menghindari efek hisap
tersebut. (TXT, hlm. 1616, bagian 1, baris 95-105.)

\textbf{Laporan sumber.} Penelitian terdahulu yang dirangkum artikel
umumnya membahas interaksi spasial antarkota dengan tingkat hierarki
sama atau menggunakan skala besar seperti zona ekonomi, provinsi, kota
tingkat prefektur, dan wilayah kabupaten. Pendekatan tersebut dinilai
kurang menjelaskan perbedaan interaksi antara kota berhierarki tinggi
dan rendah, serta kurang menangkap hubungan internal di dalam satu
aglomerasi perkotaan. (TXT, hlm. 1616-1617, bagian 1, baris 106-119.)

\textbf{Interpretasi penulis.} Artikel memosisikan Delta Sungai Yangtze
sebagai kasus penting untuk menguji apakah efek limpahan pertumbuhan
atau agglomeration shadow lebih dominan pada kawasan perkotaan yang
relatif maju dan terintegrasi. Fokusnya bukan hanya arah pengaruh,
tetapi juga heterogenitas menurut hierarki kota, keberadaan batas
administratif, dan potensi pasar. (TXT, hlm. 1616-1617, bagian 1, baris
120-143.)

\textbf{Inferensi pengolahan.} Masalah penelitian memiliki dua tingkat
yang perlu dibedakan: pertama, apakah pertumbuhan kota kecil berkorelasi
dengan kedekatan atau potensi pasar kota lain; kedua, apakah korelasi
tersebut dapat disebut limpahan pertumbuhan atau agglomeration shadow.
Model jarak dan potensi pasar mengoperasionalkan tingkat pertama,
sedangkan atribusi mekanisme pada tingkat kedua masih bergantung pada
interpretasi penulis. (Inferensi dari TXT, hlm. 1617-1621, bagian 2-3,
baris 151-175 dan 218-388.)

\subsection{Objectives}\label{objectives-27}

\textbf{Laporan sumber.} Tujuan penelitian yang dinyatakan atau
dijabarkan dalam artikel ialah:

\begin{itemize}
\tightlist
\item
  menguji apakah kota besar membantu atau menghambat pertumbuhan ekonomi
  kota kecil di dalam skala aglomerasi perkotaan;
\item
  mengestimasi pengaruh spasial kota dengan hierarki berbeda terhadap
  pertumbuhan 108 kota kecil pada 2000-2010;
\item
  membedakan limpahan dari kota tingkat tinggi dan interaksi antarkota
  kecil yang memiliki tingkat hierarki sama;
\item
  menguji heterogenitas pengaruh antara kota tingkat prefektur umum,
  kota subprovinsial, dan kota setingkat munisipalitas langsung;
\item
  menilai apakah batas administratif menciptakan segmentasi pasar dan
  mengubah limpahan spasial;
\item
  menggabungkan jarak geografis dan potensi pasar sebagai ukuran
  interaksi spasial; dan
\item
  menggunakan serangkaian uji ketahanan untuk memeriksa reliabilitas
  hasil dasar. (TXT, hlm. 1617-1618, bagian 1 dan 2, baris 134-175 dan
  178-215.)
\end{itemize}

Pertanyaan pokoknya dirumuskan secara eksplisit sebagai apakah kota
besar berkontribusi terhadap pertumbuhan ekonomi kota kecil. Hipotesis
teoretis yang diuji tidak menghasilkan satu arah yang dianggap
universal; artikel menyatakan bahwa hasil bergantung pada tahap
perkembangan wilayah dan jarak kota kecil ke kota besar. (TXT, hlm.
1616-1617, bagian 1, baris 95-105 dan 120-133.)

\subsection{Literature Survey}\label{literature-survey-27}

\textbf{Laporan sumber.} Artikel memulai dari teori interaksi spasial
dan teori pembangunan wilayah. Myrdal, Hirschman, dan Friedmann
digunakan untuk menjelaskan kemungkinan simultan antara konsentrasi atau
penyerapan di satu sisi dan difusi atau limpahan di sisi lain. Perbedaan
antara tahap awal dan tahap lanjut pembangunan menjadi dimensi waktu
dalam penalaran artikel. (TXT, hlm. 1616, bagian 1, baris 95-105.)

\textbf{Laporan sumber.} New Economic Geography digunakan untuk
menjelaskan agglomeration shadow secara spasial. Dalam ringkasan
artikel, perusahaan baru yang dekat dengan pusat kota dapat menghadapi
kondisi tidak menguntungkan karena pusat menyerap faktor, sehingga kota
kecil yang cukup jauh dari pusat dapat mengalami pertumbuhan relatif
lebih baik. Kerangka ini dibandingkan dengan kemungkinan limpahan
positif ketika kedekatan memudahkan akses terhadap eksternalitas
aglomerasi. (TXT, hlm. 1617, bagian 2.1, baris 151-155; hlm. 1616,
bagian 1, baris 95-105.)

\textbf{Laporan sumber.} Kajian empiris yang dirangkum menyatakan bahwa
perkembangan ekonomi suatu wilayah dipengaruhi oleh faktor internal dan
perkembangan wilayah sekitarnya. Artikel menilai banyak studi sebelumnya
berfokus pada kota dengan hierarki sama dan belum cukup mempertimbangkan
bahwa eksternalitas aglomerasi dapat berbeda menurut hierarki kota serta
menurun secara nonlinier terhadap jarak. (TXT, hlm. 1616-1617, bagian 1,
baris 106-114.)

\textbf{Laporan sumber.} Artikel juga menyoroti ketidakseimbangan skala
penelitian. Ada studi pada zona ekonomi, provinsi, kota tingkat
prefektur, dan wilayah kabupaten, tetapi studi kasus mengenai interaksi
internal dalam aglomerasi perkotaan dinilai lebih terbatas. Penulis
menilai skala aglomerasi lebih relevan bagi implementasi kebijakan
urbanisasi berbasis kota-kota yang saling terhubung. (TXT, hlm.
1616-1617, bagian 1, baris 114-119.)

\textbf{Laporan sumber.} Untuk mengisi celah tersebut, artikel
menggabungkan tiga jenis ukuran: sifat pertama ruang berupa jarak
geografis, serta sifat kedua ruang berupa segmentasi pasar dan potensi
pasar. Artikel menyatakan bahwa kombinasi ini digunakan agar limpahan
spasial tidak diperlakukan hanya sebagai fungsi jarak, melainkan juga
mempertimbangkan biaya atau hubungan ekonomi yang berada di balik jarak.
(TXT, hlm. 1617, bagian 1, baris 139-143; hlm. 1621, bagian 3.2.3, baris
375-388.)

\textbf{Interpretasi penulis.} Literatur dipakai untuk membangun matriks
konseptual empat kemungkinan berdasarkan dua dimensi: tahap perkembangan
dan kedekatan dengan kota besar. Pada tahap awal, efek penyerapan atau
backwash diperkirakan lebih kuat; pada tahap lanjut, kota yang dekat
diperkirakan memperoleh lebih banyak limpahan dan radiasi. Delta Sungai
Yangtze dipilih untuk menguji sebagian dari argumen tersebut karena
keterbatasan data membuat perbandingan lintas tahap perkembangan belum
dapat dilakukan. (TXT, hlm. 1616-1617, bagian 1, baris 120-133.)

\textbf{Inferensi pengolahan.} Tinjauan pustaka ini bersifat naratif dan
berbasis referensi yang dipilih penulis. TXT tidak menjelaskan strategi
pencarian, kriteria inklusi, rentang waktu korpus, prosedur penilaian
mutu, atau cara menyelesaikan perbedaan temuan antarstudi. Karena itu,
kelengkapan tinjauan literatur tidak dapat diverifikasi dari TXT; yang
dapat diverifikasi ialah posisi konseptual dan celah penelitian
sebagaimana dirumuskan artikel. (Inferensi berdasarkan tidak adanya
keterangan tersebut dalam bagian 1 dan daftar referensi, TXT baris
65-143 dan 473-544.)

\subsection{Method Used}\label{method-used-27}

\textbf{Laporan sumber.} Penelitian menggunakan model pertumbuhan
ekonomi kuantitatif dan estimasi Ordinary Least Squares (OLS) untuk
model dasar. Persamaan umum yang ditampilkan ialah:

\texttt{Delta\ pgdp\_i,t\ =\ alpha\ +\ beta1\ dis\_i,0\ +\ beta2\ pgdp\_i,0\ +\ beta3\ den\_i,0\ +\ beta4\ inv\_i,0\ +\ beta5\ lab\_i,0\ +\ beta6\ urb\_i,0\ +\ beta7\ fdi\_i,0\ +\ beta8\ gov\_i,0\ +\ beta9\ edu\_i,0\ +\ beta10\ ame\_i,0\ +\ epsilon\_i.}

Dalam persamaan, \texttt{i} menunjukkan kota kecil, \texttt{0} tahun
awal, dan \texttt{t} rentang waktu. Variabel penjelas memakai
karakteristik awal tahun 2000, sedangkan variabel yang dijelaskan ialah
pertumbuhan tahunan rata-rata PDB per kapita pada 2000-2010. Penulis
menyatakan bahwa penggunaan karakteristik awal dimaksudkan untuk relatif
mengurangi persoalan pertumbuhan yang memengaruhi karakteristik kota dan
untuk menangkap pengaruh jangka panjang karakteristik awal. (TXT, hlm.
1617, bagian 2.1, baris 146-175.)

\textbf{Laporan sumber.} Jarak geografis digunakan untuk menangkap arah
dan heterogenitas interaksi spasial. Artikel membedakan jarak antarkota
kecil, jarak ke kota tingkat prefektur atau lebih tinggi, jarak ke kota
subprovinsial, jarak ke munisipalitas langsung, serta beberapa kombinasi
jarak menurut hierarki. Variabel dummy menunjukkan apakah kota kecil
terdekat berada dalam kota tingkat prefektur yang sama atau memiliki
hubungan administratif vertikal. (TXT, hlm. 1617-1618, bagian 2.1-2.2,
baris 156-206.)

\textbf{Laporan sumber.} Uji ketahanan dilakukan dalam tiga bentuk
utama:

\begin{itemize}
\tightlist
\item
  suku kuadrat jarak ditambahkan untuk menguji kemungkinan hubungan
  nonlinier; seluruh suku linear dan kuadrat dilaporkan tidak
  signifikan;
\item
  uji ketergantungan spasial dilakukan terhadap residual model dasar
  dengan matriks tetangga terdekat K=1 sampai K=5, menggunakan Moran's I
  serta statistik LM-LAG, Robust LM-LAG, LM-ERR, dan Robust LM-ERR;
\item
  ukuran jarak inkremental dan potensi pasar digunakan sebagai pengganti
  atau tambahan ukuran jarak geografis. (TXT, hlm. 1619-1621, bagian 3.1
  dan 3.2, baris 300-314 dan 324-388.)
\end{itemize}

\textbf{Laporan sumber.} Potensi pasar dihitung dengan rumus
\texttt{MP\_i\ =\ sum(GDP\_j\ /\ Dis\_ij)}. \texttt{GDP\_j} adalah PDB
tahun 2000 dari kota kecil terdekat atau wilayah perkotaan terdekat,
sedangkan \texttt{Dis\_ij} ialah jarak geografis yang dihitung
sebelumnya. Model potensi pasar tetap mengontrol variabel lain dan
diestimasi untuk 108 observasi. (TXT, hlm. 1621-1622, bagian 3.2.3 dan
Tabel 4, baris 375-411.)

\textbf{Interpretasi penulis.} Penulis menyatakan bahwa tidak
ditemukannya ketergantungan spasial residual membuat hasil OLS model
dasar dapat dianggap andal. Mereka juga menyatakan bahwa hasil dengan
jarak inkremental dan potensi pasar relatif konsisten dengan model
dasar. (TXT, hlm. 1620-1621, bagian 3.2.1-3.2.3, baris 324-388.)

\textbf{Inferensi pengolahan.} Rancangan ini merupakan estimasi
observasional atas satu interval pertumbuhan dengan karakteristik awal,
bukan eksperimen, kuasi-eksperimen, atau panel pertumbuhan berulang. Uji
residual dan beberapa spesifikasi alternatif meningkatkan pemeriksaan
ketahanan, tetapi tidak mengubah fakta bahwa koefisien jarak dapat
menangkap faktor lain yang berkorelasi dengan lokasi kota kecil.
(Inferensi dari struktur persamaan dan periode pengamatan dalam TXT,
hlm. 1617-1618, baris 165-175 dan 178-187.)

\subsection{Dataset}\label{dataset-27}

\textbf{Laporan sumber.} Data ekonomi terutama berasal dari buku
statistik tahun 2001 dan 2011 milik provinsi atau kota terkait serta
\emph{China County (City) Social and Economic Statistical Yearbook}
tahun 2001. Data penduduk berasal dari sensus penduduk nasional kelima
dan keenam. Batas administratif diseragamkan mengikuti batas tahun 2010.
Karena data hilang, Chongming County dikeluarkan dari sampel. (TXT, hlm.
1618, bagian 2.2, baris 178-184.)

\textbf{Laporan sumber.} Variabel pertumbuhan menggunakan PDB per kapita
2000-2010 dengan mengeluarkan sektor primer. Indeks harga konsumen (CPI)
dari provinsi atau kota terkait, dengan tahun 2000 sebagai harga
konstan, digunakan untuk mengeluarkan pengaruh inflasi. Populasi yang
dipakai untuk indikator penduduk ialah penduduk perkotaan dari data
sensus. (TXT, hlm. 1618, bagian 2.2, baris 185-187.)

\textbf{Laporan sumber.} Jarak dihitung sebagai jarak Euclidean
menggunakan peta dasar basis data topografi Tiongkok skala 1:4.000.000
dan perangkat lunak ArcGIS. Potensi pasar menggunakan PDB wilayah
terkait pada 2000 dan jarak antardaerah yang sama dengan perhitungan
sebelumnya. (TXT, hlm. 1618 dan 1621, bagian 2.2 dan 3.2.3, baris
178-184 dan 375-384.)

\textbf{Laporan sumber.} Variabel kontrol terdiri atas logaritma PDB per
kapita awal, kepadatan penduduk, rasio investasi aset tetap terhadap
PDB, rasio tenaga kerja nonpertanian terhadap total penduduk, rasio
penduduk nonpertanian, rasio investasi asing aktual terhadap PDB, rasio
belanja pemerintah tertentu terhadap PDB, proporsi penduduk
berpendidikan setingkat SMA atau lebih, dan jumlah ranjang rumah sakit
per 10.000 penduduk. (TXT, hlm. 1618, bagian 2.2, baris 207-215.)

\textbf{Inferensi pengolahan.} TXT tidak menyediakan data mentah, daftar
lengkap 108 kota, nilai setiap variabel, kode pengolahan, atau prosedur
yang cukup untuk mereplikasi seluruh dataset. TXT juga tidak menjelaskan
rincian penanganan data hilang selain pengeluaran Chongming County,
sehingga kelengkapan sampel setelah keputusan tersebut tidak dapat
diperiksa lebih jauh dari file. (Inferensi berdasarkan isi TXT, baris
178-215.)

\subsection{Population / Sample}\label{population-sample-27}

\textbf{Laporan sumber.} Populasi analitis terdiri atas 108 kota kecil
di Delta Sungai Yangtze, yang dijelaskan sebagai kabupaten dan kota
tingkat kabupaten; untuk definisi operasional kota kecil, artikel
menggunakan bagian penduduk perkotaan dari kabupaten dan kota tingkat
kabupaten. Wilayah penelitian dirujuk sebagai dua provinsi dan satu
kota. Batas administrasi seluruh unit disesuaikan dengan kondisi tahun
2010. (TXT, hlm. 1617-1618, bagian 1 dan 2.2, baris 134-137 dan
178-184.)

\textbf{Laporan sumber.} Artikel membandingkan kota kecil dengan kota
tingkat prefektur, kota subprovinsial, dan munisipalitas langsung
sebagai kota berhierarki tinggi. Untuk memberi konteks ukuran, teks
menyebut populasi perkotaan 2010 Shanghai sekitar 20 juta, Nanjing 5,83
juta, Hangzhou 5,16 juta, Ningbo 2,58 juta, rata-rata 21 kota tingkat
prefektur sekitar 1,27 juta, dan rata-rata 108 kabupaten atau kota
tingkat kabupaten sekitar 370 ribu. (TXT, hlm. 1617, bagian 2.1, baris
156-164.)

\textbf{Laporan sumber.} Seluruh estimasi model dasar dan model
ketahanan memakai jumlah sampel 108. Chongming County dikeluarkan karena
data hilang. Tidak ada uraian tentang pembagian sampel
pelatihan-pengujian atau pengambilan sampel acak. (TXT, hlm. 1618-1622,
Tabel 1, Tabel 3, dan Tabel 4, baris 178-184, 294-296, 360-373, dan
396-411.)

\textbf{Inferensi pengolahan.} Sampel ini lebih tepat dipahami sebagai
kasus kawasan terpilih daripada sampel acak yang mewakili seluruh kota
kecil Tiongkok. Artikel sendiri membatasi penerapan hasil pada wilayah
domestik yang berada pada tahap perkembangan sebanding dengan Delta
Sungai Yangtze. Generalisasi ke wilayah dengan tahap perkembangan,
struktur industri, atau institusi berbeda tidak didukung langsung oleh
desain ini. (TXT, hlm. 1616-1617 dan 1622, baris 129-133 dan 423-433.)

\subsection{Dependent Variables}\label{dependent-variables-27}

\textbf{Laporan sumber.} Variabel dependen utama ialah
\texttt{Delta\ pgdp}, yaitu laju pertumbuhan tahunan rata-rata PDB per
kapita kota kecil selama 2000-2010, dengan sektor primer dikeluarkan dan
pengaruh inflasi dikoreksi menggunakan CPI berharga konstan 2000. Tidak
ada variabel dependen kedua yang dilaporkan. (TXT, hlm. 1617-1618,
bagian 2.1-2.2, baris 165-175 dan 185-187.)

\textbf{Laporan sumber.} Variabel dependen tersebut dimodelkan terhadap
jarak geografis, karakteristik ekonomi awal, faktor aglomerasi, input
modal dan tenaga kerja, urbanisasi, investasi asing, intervensi
pemerintah, modal manusia, dan kualitas hidup. Jarak dan potensi pasar
bukan variabel dependen; keduanya merupakan ukuran interaksi atau
limpahan spasial yang dipakai sebagai penjelas pertumbuhan. (TXT, hlm.
1617-1618, bagian 2.1-2.2, baris 165-175 dan 188-215.)

\textbf{Interpretasi penulis.} Koefisien jarak yang negatif ditafsirkan
sebagai tanda bahwa kota kecil yang lebih dekat dengan kota besar
mengalami pertumbuhan lebih cepat, yaitu adanya limpahan pertumbuhan
positif. Koefisien dummy hubungan administratif dan ukuran potensi pasar
digunakan untuk menafsirkan segmentasi pasar serta interaksi ekonomi.
(TXT, hlm. 1619-1621, bagian 3.1 dan 3.2.3, baris 229-240, 305-314, dan
375-388.)

\textbf{Inferensi pengolahan.} Agglomeration shadow tidak diukur sebagai
variabel dependen atau indikator ambang yang berdiri sendiri. Ia
disimpulkan dari tanda dan signifikansi jarak, bentuk nonlinier yang
diuji, dummy administratif, dan potensi pasar. Oleh karena itu, ``tidak
ada bukti langsung'' dalam artikel tidak identik dengan pembuktian bahwa
agglomeration shadow tidak ada. (Inferensi dari TXT, hlm. 1616-1621,
baris 95-105, 300-304, dan 324-388.)

\subsection{Results}\label{results-27}

\textbf{Laporan sumber: model dasar.} Enam model pada Tabel 1 semuanya
memakai 108 observasi. Adjusted R2 berada pada rentang 0,6529-0,6894.
Model (1), yang tidak membedakan hierarki kota besar, tidak menemukan
pengaruh signifikan dari jarak ke kota tingkat prefektur atau lebih
tinggi. Namun, dummy hierarki kota besar bernilai positif dan
signifikan, sehingga penulis menilai pembedaan hierarki diperlukan.
(TXT, hlm. 1619, bagian 3.1 dan Tabel 1, baris 229-240 dan 242-296.)

\textbf{Laporan sumber: heterogenitas hierarki.} Ketika tingkat kota
dibedakan, jarak ke kota subprovinsial dan jarak ke Shanghai dalam
spesifikasi terkait berkoefisien negatif dan signifikan. Artikel
menafsirkan hasil ini sebagai limpahan pertumbuhan positif, dengan
limpahan dari kota subprovinsial lebih kuat daripada dari Shanghai;
pengaruh kota tingkat prefektur umum dinilai kecil. Dalam model yang
mempertimbangkan tiga tingkat sekaligus, hanya jarak ke kota
subprovinsial yang tetap signifikan; koefisiennya -0,0074 dan penulis
melaporkan bahwa pemendekan jarak 100 km menaikkan pertumbuhan kota
kecil sekitar 0,74\%. (TXT, hlm. 1619, Tabel 1 dan pembahasan model
(2)-(6), baris 229-240 dan 245-296.)

\textbf{Laporan sumber: karakteristik kontrol.} Pada model (6), log PDB
per kapita awal berkoefisien -7,1251 dan signifikan; tenaga kerja,
urbanisasi, investasi asing, belanja pemerintah, modal manusia, dan
kualitas hidup berkoefisien positif serta signifikan pada tingkat yang
berbeda. Kepadatan penduduk dan rasio investasi tidak signifikan. Kota
tingkat kabupaten juga mempunyai dummy positif dan signifikan
dibandingkan kabupaten. (TXT, hlm. 1619-1620, Tabel 1 dan pembahasan,
baris 260-323.)

\textbf{Interpretasi penulis: model dasar.} Koefisien negatif PDB awal
dibaca sebagai konvergensi kondisional: kota kecil dengan tingkat awal
lebih rendah tumbuh lebih cepat. Koefisien tenaga kerja, urbanisasi,
investasi asing, belanja pemerintah, modal manusia, dan kualitas hidup
dibaca sebagai dukungan terhadap peran faktor-faktor tersebut dalam
pertumbuhan. Ketidaksignifikanan kepadatan dan investasi dibaca sebagai
tidak adanya pengaruh yang jelas dalam spesifikasi tersebut. (TXT, hlm.
1619-1620, bagian 3.1, baris 315-323.)

\textbf{Laporan sumber: nonlinieritas jarak.} Penambahan suku kuadrat
untuk jarak antarkota kecil, jarak ke kota tingkat prefektur atau lebih
tinggi, kota tingkat prefektur umum, kota subprovinsial, dan Shanghai
menghasilkan suku linear serta kuadrat yang tidak signifikan. Penulis
mengusulkan bahwa efek nonlinier mungkin memerlukan ruang geografis yang
lebih luas, sedangkan jarak di dalam sampel Delta Sungai Yangtze relatif
kecil dan seluruh kawasan masih berada dalam jangkauan radiasi
metropolitan. (TXT, hlm. 1619, bagian 3.1, baris 300-304.)

\textbf{Laporan sumber: batas administratif.} Dummy kota tingkat
prefektur umum bernilai negatif dan signifikan pada beberapa model.
Interaksi antara jarak ke kota tingkat prefektur umum dan dummy tersebut
juga bernilai negatif dan signifikan. Penulis menyatakan bahwa kota
kecil yang berada dalam administrasi kota tingkat prefektur terdekat
tumbuh lebih lambat daripada kota kecil di luar batas tersebut, yang
dibaca sebagai indikasi penyerapan sumber daya oleh kota atas dan
segmentasi pasar akibat batas administratif. (TXT, hlm. 1619, bagian 3.1
dan Tabel 1, baris 305-314.)

\textbf{Laporan sumber: uji ketergantungan spasial.} Untuk residual
model (6), Moran's I pada K=1 sampai K=5 masing-masing dilaporkan
0,1094, 0,1120, 0,0571, 0,0480, dan 0,0064, dengan nilai P dalam tanda
kurung 0,1690, 0,1070, 0,1640, 0,1870, dan 0,3390. Statistik LM-LAG,
Robust LM-LAG, LM-ERR, dan Robust LM-ERR juga tidak signifikan pada
seluruh K yang diuji. Penulis menyimpulkan residual tidak memiliki
ketergantungan spasial yang signifikan dan hasil OLS model dasar dapat
dipercaya. (TXT, hlm. 1620-1621, bagian 3.2.1 dan Tabel 2, baris
324-347.)

\textbf{Laporan sumber: jarak inkremental.} Tabel 3 menunjukkan bahwa
hasil dasar tetap relatif stabil setelah memasukkan jarak inkremental.
Dalam model (6), jarak inkremental antara kota subprovinsial dan kota
tingkat prefektur umum berkoefisien -0,0096 dan signifikan pada tingkat
1\%, sedangkan jarak inkremental antara munisipalitas langsung dan kota
subprovinsial berkoefisien -0,0019 dan tidak signifikan. Dummy hubungan
administratif tetap negatif dan signifikan. Adjusted R2 model-model
Tabel 3 berada pada rentang 0,6556-0,6858. (TXT, hlm. 1621, bagian 3.2.2
dan Tabel 3, baris 348-373.)

\textbf{Laporan sumber: potensi pasar.} Dalam Tabel 4, potensi pasar
keseluruhan mempunyai koefisien 0,4251 dan signifikan pada model yang
mempertimbangkan 25 wilayah perkotaan. Ketika potensi pasar dibedakan
menurut hierarki, potensi pasar kota subprovinsial merupakan ukuran
hierarki yang dilaporkan signifikan dan positif, sekitar 0,0141-0,0144
pada spesifikasi terkait; potensi pasar kota tingkat prefektur umum dan
munisipalitas langsung tidak dilaporkan signifikan. Penulis menyatakan
hasil ini mendukung ketahanan kesimpulan model dasar. (TXT, hlm.
1621-1622, bagian 3.2.3 dan Tabel 4, baris 375-411.)

\textbf{Inferensi pengolahan.} Hasil deskriptif dan hasil kondisional
perlu dibaca terpisah. Model menunjukkan hubungan positif antara
kedekatan ke kota besar dan pertumbuhan, tetapi ukuran yang digunakan
adalah jarak serta proksi potensi pasar, bukan observasi limpahan
aktual. Selain itu, artikel menyebut pengujian OLS, tetapi TXT tidak
menampilkan seluruh koefisien ketidakpastian, prosedur penanganan semua
kemungkinan faktor terlewat, atau desain identifikasi kausal. (Inferensi
dari TXT, hlm. 1617-1622, baris 165-175 dan 221-411.)

\subsection{Findings}\label{findings-27}

\textbf{Temuan yang dilaporkan sumber.} Di Delta Sungai Yangtze,
kedekatan dengan kota besar cenderung meningkatkan pertumbuhan ekonomi
kota kecil. Artikel tidak menemukan bukti langsung agglomeration shadow
dalam sampel tersebut. (TXT, hlm. 1615 dan 1622, bagian 4, baris 52-60
dan 416-429.)

\textbf{Temuan yang dilaporkan sumber.} Arah interaksi lebih banyak
berasal dari kota berhierarki tinggi menuju kota berhierarki rendah.
Interaksi antarkota kecil dengan hierarki sama relatif lemah, sedangkan
limpahan pertumbuhan kota subprovinsial paling konsisten dan paling
signifikan dalam spesifikasi yang dilaporkan. (TXT, hlm. 1619-1622,
bagian 3 dan 4, baris 229-240, 347-358, dan 434-450.)

\textbf{Temuan yang dilaporkan sumber.} Batas administratif mengurangi
limpahan dari kota tingkat prefektur umum ke kota kecil di luar wilayah
administrasinya dan berkaitan dengan pertumbuhan yang lebih rendah bagi
kota kecil yang berada di bawah administrasi kota terdekat. Artikel
menyebut kondisi ini sebagai segmentasi pasar yang lahir dari ``ekonomi
administratif''. (TXT, hlm. 1619-1623, bagian 3.1 dan 4, baris 305-314
dan 451-462.)

\textbf{Interpretasi penulis.} Penulis menyimpulkan bahwa Delta Sungai
Yangtze secara keseluruhan masih berada dalam jangkauan radiasi
metropolitan. Oleh sebab itu, efek penyebaran dari kota besar lebih
dominan daripada efek bayangan aglomerasi dalam konteks wilayah dan
periode penelitian. (TXT, hlm. 1622, bagian 4, baris 416-429.)

\textbf{Interpretasi penulis.} Pengaruh Shanghai terhadap kota kecil
mungkin tidak bekerja langsung, melainkan melalui peran perantara kota
subprovinsial. Penulis juga menduga bahwa perbedaan kewenangan ekonomi
tingkat provinsi, kompetisi antarpemerintah kota kecil, dan struktur
ekonomi menjelaskan lemahnya interaksi antarkota kecil, tetapi
menyatakan dugaan tersebut masih membutuhkan pengujian empiris. (TXT,
hlm. 1622-1623, bagian 4, baris 434-450.)

\textbf{Inferensi pengolahan.} Temuan utama paling kuat mendukung
pernyataan kontekstual bahwa kota kecil yang lebih dekat atau lebih
terhubung secara potensial dengan kota berhierarki tinggi memiliki
pertumbuhan lebih tinggi setelah kontrol tertentu. Temuan ini tidak
memisahkan secara empiris apakah hubungan tersebut bekerja melalui
tenaga kerja, input-output, pengetahuan, investasi, akses pasar, atau
kebijakan. (Inferensi dari TXT, hlm. 1617-1621, baris 151-175 dan
375-411.)

\textbf{Inferensi pengolahan.} ``Tidak ada agglomeration shadow'' perlu
dipahami sebagai tidak ditemukannya bukti langsung dalam model dan
kawasan ini, bukan sebagai penolakan universal terhadap teori tersebut.
Artikel sendiri menyatakan bahwa hasil dapat berbeda menurut tahap
perkembangan, struktur industri, dan lingkungan kelembagaan. (TXT, hlm.
1616-1617 dan 1622, baris 120-133 dan 423-433.)

\subsection{Conclusions}\label{conclusions-27}

\textbf{Kesimpulan yang dinyatakan penulis.} Artikel merumuskan tiga
kesimpulan utama:

\begin{itemize}
\tightlist
\item
  Di kawasan Delta Sungai Yangtze, kedekatan dengan kota besar membantu
  pertumbuhan ekonomi kota kecil. Agglomeration shadow yang diprediksi
  New Economic Geography tidak didukung oleh bukti langsung dalam sampel
  ini, sehingga pembangunan ekonomi kota kecil masih berada dalam
  pengaruh kawasan metropolitan. (TXT, hlm. 1622, bagian 4, baris
  416-429.)
\item
  Interaksi spasial lebih banyak menunjukkan pengaruh kota berhierarki
  tinggi terhadap kota berhierarki rendah. Interaksi antarkota kecil
  dengan hierarki sama tidak signifikan, dan kota subprovinsial
  menunjukkan limpahan pertumbuhan positif yang paling signifikan. (TXT,
  hlm. 1622-1623, bagian 4, baris 434-450.)
\item
  Batas administratif menyebabkan segmentasi pasar dan melemahkan
  limpahan spasial, khususnya limpahan kota tingkat prefektur umum
  kepada kota kecil yang tidak berada di bawah administrasinya. (TXT,
  hlm. 1623, bagian 4, baris 451-462.)
\end{itemize}

\textbf{Interpretasi penulis.} Kesimpulan tersebut digunakan untuk
mendukung penguatan fungsi radiasi dan dorongan kota besar, pelepasan
limpahan spasial secara menyeluruh antarkota, serta pengurangan hambatan
yang timbul dari batas administratif. Penulis menekankan bahwa kebijakan
harus dibedakan menurut tahap perkembangan aglomerasi. (TXT, hlm.
1622-1623, bagian 4, baris 416-462.)

\textbf{Inferensi pengolahan.} Kesimpulan empiris artikel valid terutama
sebagai temuan kasus untuk 108 kota kecil di Delta Sungai Yangtze selama
2000-2010. Ia tidak cukup untuk menyatakan efek kausal kota besar
terhadap semua kota kecil, tidak membuktikan bahwa limpahan selalu satu
arah, dan tidak mengukur secara langsung mekanisme yang membuat
kedekatan produktif atau merugikan. (Inferensi dari desain, sampel, dan
batasan generalisasi yang dinyatakan TXT, baris 134-175 dan 423-433.)

\subsection{Contributions}\label{contributions-27}

\textbf{Kontribusi yang diklaim sumber.} Artikel menyatakan tiga
kontribusi marginal utama:

\begin{itemize}
\tightlist
\item
  \textbf{Arah interaksi:} interaksi spasial dibedakan antara kota
  berhierarki tinggi terhadap kota berhierarki rendah dan interaksi
  antarkota dengan hierarki sama; artikel juga mempertimbangkan
  heterogenitas pengaruh munisipalitas langsung, kota subprovinsial, dan
  kota tingkat prefektur umum.
\item
  \textbf{Skala aglomerasi:} penelitian memusatkan perhatian pada
  limpahan spasial di dalam skala aglomerasi perkotaan, dengan
  mempertimbangkan bahwa interaksi mempunyai batas geografis.
\item
  \textbf{Ukuran limpahan:} penelitian menggabungkan sifat pertama ruang
  berupa jarak geografis dengan sifat kedua ruang berupa segmentasi
  pasar dan potensi pasar, lalu menggunakan beberapa uji ketahanan untuk
  memeriksa konsistensi estimasi. (TXT, hlm. 1617, bagian 1, baris
  134-143.)
\end{itemize}

\textbf{Kontribusi empiris yang ditunjukkan oleh isi artikel.} Dengan
108 unit kota kecil dan periode 2000-2010, artikel menyediakan satu
kasus yang membandingkan secara serentak hierarki kota besar, interaksi
antarkota kecil, batas administratif, dan potensi pasar dalam model
pertumbuhan. Uji jarak inkremental, potensi pasar, dan residual spasial
digunakan untuk memperkuat klaim bahwa hasil dasar tidak bergantung pada
satu ukuran saja. (TXT, hlm. 1618-1622, bagian 2-3, baris 178-215 dan
324-411.)

\textbf{Inferensi pengolahan.} Kontribusi paling jelas ialah pengalihan
fokus dari limpahan antarkota yang homogen ke struktur hierarkis di
dalam aglomerasi. Namun, kontribusi terhadap penjelasan mekanisme
agglomeration shadow atau limpahan belum bersifat final karena konstruk
tersebut diproksikan melalui jarak, dummy administratif, dan potensi
pasar, bukan melalui data arus atau mekanisme langsung. (Inferensi dari
TXT, hlm. 1617-1621, baris 139-143, 188-206, dan 375-388.)

\subsection{Research Gap}\label{research-gap-27}

\textbf{Kesenjangan yang dinyatakan sumber sebelum penelitian.} Artikel
mengidentifikasi beberapa kekosongan:

\begin{itemize}
\tightlist
\item
  studi terdahulu lebih sering menganalisis interaksi antarkota dengan
  hierarki sama dan kurang membedakan pengaruh kota dengan hierarki
  berbeda;
\item
  studi pada skala nasional, zona ekonomi, provinsi, kota tingkat
  prefektur, atau kabupaten belum cukup menjelaskan interaksi internal
  di dalam aglomerasi perkotaan;
\item
  pengaruh kota besar terhadap kota kecil perlu diuji dengan
  mempertimbangkan heterogenitas antara munisipalitas langsung, kota
  subprovinsial, dan kota tingkat prefektur umum;
\item
  jarak geografis saja tidak cukup untuk mewakili hubungan spasial,
  sehingga segmentasi pasar dan potensi pasar perlu dimasukkan; dan
\item
  perlu dibedakan kondisi tahap awal dan tahap lanjut pembangunan untuk
  memahami kapan efek penyerapan atau limpahan lebih dominan. (TXT, hlm.
  1616-1617, bagian 1, baris 95-143.)
\end{itemize}

\textbf{Kesenjangan yang masih tersisa menurut agenda sumber.}
Penelitian belum mencakup pengaruh kota berhierarki rendah terhadap kota
berhierarki tinggi, belum membandingkan secara dinamis wilayah dengan
tahap perkembangan berbeda, dan belum cukup mengidentifikasi hubungan
sebab-akibat serta mekanisme internal limpahan spasial. Artikel menyebut
kebutuhan untuk mengukur kolam pasar tenaga kerja, hubungan
input-output, dan limpahan pengetahuan. (TXT, hlm. 1623, bagian 4, baris
463-471.)

\textbf{Inferensi pengolahan.} Artikel menjawab apakah terdapat hubungan
pertumbuhan yang konsisten dengan limpahan, tetapi belum sepenuhnya
menjawab mengapa hubungan tersebut terjadi atau kapan efek bayangan
muncul. Gap residual juga mencakup pengujian dengan ukuran akses selain
jarak Euclidean, observasi perubahan sepanjang waktu, dan pengukuran
arus ekonomi yang dapat membedakan korelasi lokasi dari mekanisme
limpahan. Usulan terakhir merupakan inferensi review, bukan agenda
eksplisit seluruhnya dalam artikel. (Inferensi dari TXT, hlm. 1618-1623,
baris 178-215 dan 463-471.)

\subsection{Limitations}\label{limitations-27}

\textbf{Keterbatasan yang dinyatakan atau diakui sumber.}

\begin{itemize}
\tightlist
\item
  Data yang tersedia membuat penelitian hanya menguji satu kawasan,
  yaitu Delta Sungai Yangtze, sehingga kesimpulan terutama relevan bagi
  wilayah domestik pada tahap perkembangan yang sebanding.
\item
  Penalaran lengkap tentang perbedaan tahap perkembangan membutuhkan
  data dari aglomerasi yang matang, relatif matang, dan masih
  berkembang; data semacam itu belum digunakan dalam artikel.
\item
  Artikel belum mempertimbangkan pengaruh kota berhierarki rendah
  terhadap kota berhierarki tinggi secara komprehensif.
\item
  Pertimbangan terhadap identifikasi kausal dan mekanisme internal
  limpahan spasial masih terbatas; tenaga kerja, input-output, dan
  limpahan pengetahuan disebut sebagai agenda lanjutan.
\item
  Sumber data yang hilang menyebabkan Chongming County dikeluarkan dari
  sampel. (TXT, hlm. 1616-1618 dan 1622-1623, baris 129-133, 178-184,
  dan 463-471.)
\end{itemize}

\textbf{Keterbatasan metodologis yang merupakan inferensi pengolahan.}

\begin{itemize}
\tightlist
\item
  \textbf{Identifikasi kausal:} penggunaan karakteristik awal membantu
  mengurangi kemungkinan pertumbuhan memengaruhi karakteristik kota,
  tetapi satu estimasi atas pertumbuhan 2000-2010 tetap rentan terhadap
  faktor terlewat, seleksi lokasi, dan hubungan sebab-akibat yang tidak
  ditangani secara eksplisit.
\item
  \textbf{Pengukuran jarak:} jarak Euclidean tidak sama dengan waktu
  perjalanan, biaya angkutan, arus komuter, akses pasar aktual, atau
  intensitas hubungan antarkota. TXT tidak melaporkan validasi ukuran
  jarak terhadap ukuran aliran tersebut.
\item
  \textbf{Pengukuran mekanisme:} jarak ke kota terdekat, dummy batas
  administratif, dan potensi pasar merupakan proksi. TXT tidak
  memberikan data langsung tentang perusahaan, tenaga kerja,
  input-output, pengetahuan, perdagangan, atau mobilitas yang dapat
  menunjukkan proses limpahan.
\item
  \textbf{Agregasi:} outcome PDB per kapita dan variabel kontrol bekerja
  pada unit administratif; variasi internal di dalam kabupaten atau kota
  tingkat kabupaten dapat tersamarkan. TXT tidak menjelaskan rincian
  spasial yang lebih halus.
\item
  \textbf{Nonlinieritas:} semua suku jarak linear dan kuadrat dilaporkan
  tidak signifikan, tetapi model tidak ditampilkan dalam bentuk tabel
  lengkap. Penjelasan bahwa wilayah terlalu kecil untuk menampilkan
  nonlinieritas adalah interpretasi penulis, bukan pengujian terhadap
  seluruh bentuk ambang yang mungkin.
\item
  \textbf{Sampel:} daftar lengkap 108 kota, aturan pembentukan populasi
  awal, dan dampak pengeluaran Chongming County tidak tersedia dalam
  TXT. Karena itu, kemungkinan pengaruh keputusan sampel tidak dapat
  dihitung dari file.
\item
  \textbf{Replikasi:} TXT tidak menyediakan data mentah, kode, matriks
  jarak, matriks bobot spasial, atau rincian proses pembersihan dan
  penyelarasan data administratif.
\item
  \textbf{Inferensi spasial:} uji Moran's I dan LM dilakukan terhadap
  residual dengan K=1 sampai K=5, tetapi TXT tidak melaporkan estimasi
  model spasial alternatif. Kesimpulan bahwa OLS memadai mengikuti uji
  residual yang dilaporkan penulis.
\item
  \textbf{Ketidakpastian statistik:} tabel memberikan koefisien, tanda
  signifikansi, adjusted R2, dan jumlah observasi, tetapi TXT tidak
  memuat seluruh standard error atau interval kepercayaan setiap
  koefisien.
\end{itemize}

\subsection{Challenges}\label{challenges-27}

\textbf{Tantangan yang dihadapi menurut sumber.} Tantangan substantif
pertama ialah menjelaskan gejala empiris yang berlawanan: kota kecil
dekat pusat dapat tumbuh cepat, tetapi sebagian wilayah dekat kota besar
juga tampak tertekan. Tantangan kedua ialah memisahkan pengaruh tingkat
hierarki kota besar dari pengaruh jarak semata. (TXT, hlm. 1615-1617,
bagian 1, baris 74-143.)

\textbf{Tantangan data dan pengukuran menurut sumber.} Artikel harus
menggabungkan buku statistik, buku tahunan, sensus penduduk, basis data
topografi, dan perhitungan GIS untuk membentuk variabel pertumbuhan,
populasi, jarak, dan potensi pasar. Ketidaktersediaan data juga
menyebabkan satu unit dikeluarkan dari sampel. (TXT, hlm. 1618, bagian
2.2, baris 178-215.)

\textbf{Tantangan analitis menurut sumber.} Jarak geografis dapat
mengabaikan mekanisme ekonomi yang membuat kota saling berinteraksi,
sehingga artikel menambahkan ukuran segmentasi administratif, jarak
inkremental, dan potensi pasar. Penulis juga memeriksa nonlinieritas
serta ketergantungan spasial agar hasil model dasar lebih dapat
dipercaya. (TXT, hlm. 1617-1622, bagian 2.1 dan 3.1-3.2, baris 151-175
dan 300-411.)

\textbf{Inferensi pengolahan.} Tantangan interpretasi terbesar ialah
bahwa satu pola koefisien negatif pada jarak dapat konsisten dengan
limpahan, tetapi tidak menunjukkan saluran kausalnya. Dummy
administratif juga dapat menangkap banyak perbedaan antara unit di dalam
dan di luar wilayah administrasi, bukan hanya segmentasi pasar.
Distingsi tersebut perlu dipertahankan ketika hasil artikel dipakai
untuk kajian lain. (Inferensi dari TXT, hlm. 1618-1623, baris 188-215
dan 305-314.)

\subsection{Practical Implications}\label{practical-implications-27}

\textbf{Implikasi yang dinyatakan sumber.} Artikel mendukung penguatan
fungsi radiasi dan dorongan kota besar, karena dalam kasus Delta Sungai
Yangtze kedekatan dengan kota besar berkaitan dengan pertumbuhan kota
kecil. Strategi pengembangan metropolitan dan urbanisasi berbasis
aglomerasi dipandang memiliki manfaat ekonomi dalam konteks wilayah
tersebut. (TXT, hlm. 1622, bagian 4, baris 416-422.)

\textbf{Implikasi yang dinyatakan sumber.} Limpahan tidak seharusnya
hanya dipahami sebagai hubungan kota besar dengan kota kecil. Karena
interaksi antarkota kecil masih lemah dan arah limpahan lebih banyak
satu arah, koordinasi perlu diarahkan pada pelepasan limpahan spasial
secara menyeluruh di antara kota dengan hierarki sama maupun berbeda.
(TXT, hlm. 1622-1623, bagian 4, baris 434-450.)

\textbf{Implikasi yang dinyatakan sumber.} Batas administratif perlu
dikurangi sebagai penghambat koordinasi ekonomi. Artikel mengusulkan
pembaruan sistem evaluasi pejabat, eksplorasi reformasi pembagian
administratif pada wilayah yang memenuhi syarat, perencanaan pembangunan
ekonomi regional yang terpadu, pencegahan jalan yang terputus di batas
wilayah, dan uji coba integrasi berbagai rencana. (TXT, hlm. 1623,
bagian 4, baris 451-462.)

\textbf{Implikasi yang dinyatakan sumber.} Kebijakan perlu dibedakan
menurut tahap perkembangan aglomerasi. Untuk aglomerasi yang masih
berada pada tahap awal, artikel menyarankan kewaspadaan terhadap
penyebaran berlebihan dan pembangunan yang terlalu tersebar, serta
penguatan pusat regional melalui peningkatan kepadatan, pemendekan
jarak, dan pengurangan fragmentasi. Untuk wilayah Beijing-Tianjin-Hebei,
artikel menyebut perlunya relokasi fungsi noninti ibu kota dan
pengurangan intervensi administratif melalui konektivitas infrastruktur,
koordinasi industri, dan identitas budaya. (TXT, hlm. 1622-1623, bagian
4, baris 423-433.)

\textbf{Inferensi pengolahan.} Hasil ini dapat dipakai sebagai
peringatan agar kebijakan kedekatan metropolitan tidak digeneralisasi
secara mekanis. Sebelum mengubah batas administratif atau memusatkan
investasi, diperlukan verifikasi mekanisme, karakteristik wilayah, serta
dampak distribusional pada kota kecil yang berada di bawah administrasi
kota besar. Ini merupakan implikasi kehati-hatian dari review, bukan
evaluasi kebijakan yang dilakukan artikel.

\subsection{Applications}\label{applications-27}

\textbf{Aplikasi yang didukung langsung oleh sumber.} Kerangka artikel
dapat digunakan untuk:

\begin{itemize}
\tightlist
\item
  menguji hubungan pertumbuhan kota kecil dengan jarak ke kota besar
  pada skala aglomerasi;
\item
  membedakan limpahan dari kota tingkat prefektur umum, kota
  subprovinsial, dan munisipalitas langsung;
\item
  menilai apakah hubungan administratif melemahkan limpahan spasial;
\item
  mengganti atau melengkapi jarak geografis dengan ukuran potensi pasar;
\item
  memeriksa kestabilan hasil dengan jarak inkremental, suku kuadrat
  jarak, dan uji residual spasial; serta
\item
  menyediakan bukti awal untuk kebijakan pengembangan metropolitan dan
  koordinasi antarkota. (TXT, hlm. 1617-1623, bagian 1-4, baris 134-143,
  300-411, dan 416-462.)
\end{itemize}

\textbf{Batas aplikasi.} Artikel tidak melaporkan implementasi
kebijakan, evaluasi pascapenerapan, perangkat lunak atau kode yang
dibagikan, maupun replikasi di aglomerasi lain. Informasi tersebut tidak
disebutkan dalam file. Adaptasi ke wilayah lain karena itu hanya dapat
diperlakukan sebagai aplikasi eksploratif dan memerlukan pemeriksaan
kesetaraan tahap perkembangan, data, batas administrasi, serta ukuran
interaksi. (Inferensi dari TXT, baris 129-143 dan 416-471.)

\subsection{Future Research
(Optional)}\label{future-research-optional-27}

\textbf{Agenda yang dinyatakan sumber.} Artikel mengusulkan:

\begin{itemize}
\tightlist
\item
  perbandingan dinamis antarwilayah dan antarperiode, termasuk
  perbandingan efek jangka pendek, menengah, dan panjang di Delta Sungai
  Yangtze serta perbandingan dengan kawasan saudara pada tahap
  perkembangan berbeda;
\item
  penelitian yang lebih kuat mengenai arah interaksi spasial, termasuk
  pengaruh kota berhierarki rendah terhadap kota berhierarki tinggi dan
  bentuk interaksi yang lebih beragam dalam satu kerangka;
\item
  penguatan identifikasi kausal dan mekanisme limpahan spasial melalui
  pengukuran kolam tenaga kerja, hubungan input-output, dan limpahan
  pengetahuan. (TXT, hlm. 1623, bagian 4, baris 463-471.)
\end{itemize}

\textbf{Inferensi pengolahan.} Agenda tersebut dapat dilengkapi dengan
data arus komuter, perdagangan, jaringan perusahaan, mobilitas, dan
perubahan fasilitas atau PDB dari waktu ke waktu untuk menguji saluran
limpahan secara langsung. Pengujian kepekaan terhadap definisi kota
kecil, ukuran jarak, batas administratif, dan bentuk hubungan nonlinier
juga relevan. Usulan ini adalah inferensi review dari keterbatasan
pengukuran dalam TXT, bukan agenda eksplisit yang seluruhnya ditulis
penulis.

\subsection{Provenance Notes}\label{provenance-notes-27}

\begin{itemize}
\tightlist
\item
  \textbf{Basis akses:} Review ini hanya menggunakan file TXT A12 yang
  diminta, yaitu
  \texttt{source/journal/A12-sun-ding-2016-do-large-cities-contribute-to-economic-growth-of-small-cities-evidence-from-yangtze-river-delta-in-china-10.11821-dlyj201609002.txt}.
  Tidak ada sumber eksternal, pencarian daring, PDF, atau artikel
  rujukan yang digunakan untuk menambahkan bukti.
\item
  \textbf{Cakupan baca:} Seluruh 586 baris TXT dibaca, termasuk blok
  metadata PDF, teks artikel berbahasa Mandarin pada hlm. 1615-1625,
  persamaan, Tabel 1-4, Gambar 1, abstrak berbahasa Inggris, dan daftar
  referensi yang tercantum di dalam TXT.
\item
  \textbf{Asal teks:} Blok metadata menyatakan TXT diekstrak pada 31
  Agustus 2026 dengan \texttt{pdftotext\ -layout} dari PDF 11 halaman.
  Metadata juga menyatakan PDF tidak terenkripsi dan tidak bertanda
  struktur (\emph{not tagged}). Informasi ini dilaporkan sebagai
  metadata TXT, bukan hasil pemeriksaan PDF terpisah. (TXT, baris 1-28.)
\item
  \textbf{Locator:} Nomor halaman dalam review merujuk nomor halaman
  jurnal yang tercetak, sedangkan rentang ``baris'' merujuk nomor baris
  TXT A12. Bagian, persamaan, tabel, dan gambar dicantumkan ketika
  tersedia. Karena tata letak hasil ekstraksi mempertahankan kolom dan
  pemisah halaman, beberapa angka tabel dibaca bersama label dan narasi
  artikel yang berdekatan.
\item
  \textbf{Identitas:} Nama penulis, judul, jurnal, volume, nomor,
  halaman, tanggal naskah, pendanaan, kata kunci, dan DOI diambil dari
  blok judul, catatan penulis, format jurnal, dan abstrak Inggris dalam
  TXT, bukan dari metadata produksi PDF yang mencantumkan
  \texttt{Founder\ magtech2011} sebagai author. (TXT, baris 7-14 dan
  41-62.)
\item
  \textbf{Informasi tidak tersedia dalam file:} status telaah sejawat,
  daftar lengkap unit sampel, data mentah, kode replikasi, protokol
  pembersihan data, rincian keputusan ketika data hilang selain
  Chongming County, standard error atau interval kepercayaan lengkap,
  serta rincian identifikasi kausal dan mekanisme limpahan tidak
  tersedia atau tidak dijelaskan secara lengkap.
\item
  \textbf{Pemisahan epistemik:} Label \textbf{Laporan sumber} merangkum
  apa yang dinyatakan atau ditampilkan artikel. Label
  \textbf{Interpretasi penulis} merekam makna yang diberikan Sun dan
  Ding terhadap hasil. Label \textbf{Inferensi pengolahan} menunjukkan
  evaluasi atau kehati-hatian analitis yang diturunkan dari isi TXT;
  bagian tersebut bukan klaim eksplisit artikel dan tidak memakai bukti
  di luar TXT.
\item
  \textbf{Batas generalisasi:} Artikel menyatakan bahwa kesimpulannya
  terutama berlaku bagi wilayah yang tahap perkembangannya sebanding
  dengan Delta Sungai Yangtze dan tidak otomatis dapat memandu seluruh
  Tiongkok. Review ini mempertahankan batas tersebut. (TXT, hlm.
  1616-1617 dan 1622, baris 129-133 dan 423-433.)
\end{itemize}

\section{Review A13 - Li dan Sun
(2020)}\label{review-a13---li-dan-sun-2020}

Status: Review literatur berbasis teks lengkap hasil ekstraksi PDF.

\subsection{Identitas Artikel}\label{identitas-artikel-4}

\begin{itemize}
\tightlist
\item
  \textbf{Kode:} A13.
\item
  \textbf{Judul:} \emph{Industrial Composition and Agglomeration Shadow:
  Evidence from China's Large Urban Systems}.
\item
  \textbf{Penulis:} Jiaming Li dan Dongqi Sun.
\item
  \textbf{Jenis yang tercantum dalam file:} Research Article.
\item
  \textbf{Jurnal:} \emph{Complexity}.
\item
  \textbf{Volume dan identitas artikel:} Volume 2020, Article ID
  5717803, 11 halaman. Nomor edisi dan nomor halaman artikel yang
  terpisah tidak disebutkan dalam file.
\item
  \textbf{DOI:} 10.1155/2020/5717803.
\item
  \textbf{Afiliasi:} Institute of Geographic Sciences and Natural
  Resources Research, Chinese Academy of Sciences, Beijing 100101,
  China.
\item
  \textbf{Penulis korespondensi:} Jiaming Li; alamat surel tercantum
  dalam file.
\item
  \textbf{Riwayat naskah:} diterima 14 April 2020, direvisi 12 Mei 2020,
  diterima 18 Mei 2020, dan diterbitkan 30 Juni 2020.
\item
  \textbf{Guest editor:} Wen-Ze Yue.
\item
  \textbf{Pendanaan:} Second Scientific Investigation of the Tibetan
  Plateau, Grant no. 2019QZKK0406, dan National Natural Science
  Foundation of China, Grant no. 41701128.
\item
  \textbf{Konflik kepentingan:} penulis menyatakan tidak ada konflik
  kepentingan.
\item
  \textbf{Wilayah empiris:} sistem perkotaan Tiongkok dengan enam kota
  inti, yaitu Beijing, Shanghai, Guangzhou, Tianjin, Chongqing, dan
  Shenzhen, serta kota-kota noninti yang berada dalam wilayah
  pengaruhnya.
\item
  \textbf{Periode utama:} pertumbuhan penduduk dianalisis untuk blok
  lima tahun dan sepuluh tahun dalam rentang 1980-2015; blok sepuluh
  tahun terakhir hanya mencakup 2010-2015.
\item
  \textbf{Locator identitas:} TXT baris 1-51; judul, abstrak, dan
  informasi penerbitan tercetak pada hlm. 1.
\item
  \textbf{Status telaah sejawat:} Tidak disebutkan dalam file.
\end{itemize}

\subsection{Summarized Abstract}\label{summarized-abstract-28}

\textbf{Laporan sumber.} Artikel ini berangkat dari pertanyaan dalam
\emph{new economic geography} (NEG) tentang mengapa efek
\emph{agglomeration shadow} tampak signifikan di sebagian wilayah
perkotaan, tetapi tidak di wilayah lain. Penelitian menguji dampak
komposisi industri kota inti terhadap sistem perkotaan regional di
Tiongkok dengan memakai model lokasi perkotaan dari NEG. Fokus
empirisnya adalah pertumbuhan penduduk enam kota inti dan sistem
perkotaan di sekitarnya. Hasil abstrak menyatakan bahwa jasa mempunyai
efek negatif yang signifikan terhadap \emph{market potential}, sedangkan
manufaktur mempunyai efek positif. Hasil itu disebut tetap muncul pada
skala spasial dan rentang waktu yang berbeda. (Hlm. 1, abstrak; TXT
baris 52-61.)

\textbf{Laporan sumber.} Abstrak mengaitkan hasil tersebut dengan efek
pemusatan dan aglomerasi yang kuat pada jasa tingkat tinggi. Menurut
penulis, efek itu lebih mungkin menciptakan \emph{agglomeration shadow}
pada kota noninti. Perbedaan komposisi industri dipakai untuk
menjelaskan mengapa sistem perkotaan yang berpusat pada Beijing lebih
mungkin membentuk bayangan atas kota sekitarnya dibandingkan sistem yang
berpusat pada Shanghai. (Hlm. 1, abstrak; TXT baris 55-61.)

\textbf{Interpretasi penulis.} Artikel membaca manufaktur sebagai
saluran limpahan yang cenderung mendukung pertumbuhan kota sekitarnya,
sedangkan jasa sebagai sektor yang lebih memusatkan kegiatan di kota
inti dan lebih berpotensi menekan kota noninti. Implikasi yang
dinyatakan berkaitan dengan pembangunan yang tidak merata dalam sistem
perkotaan regional dan pembangunan kawasan metropolitan di Tiongkok.
(Hlm. 1, abstrak; TXT baris 57-61.)

\textbf{Inferensi pengolahan.} Abstrak tidak mengukur
\emph{agglomeration shadow} sebagai variabel terikat yang berdiri
sendiri. Variabel hasil dalam model adalah pertumbuhan penduduk kota,
sedangkan \emph{market potential} dan interaksinya dengan komposisi
industri menjadi bagian dari penjelasan. Karena itu, klaim tentang
bayangan aglomerasi merupakan interpretasi atas hubungan model dan pola
spasial, bukan identifikasi langsung atau bukti kausal mengenai
mekanisme bayangan.

\subsection{Problem Statement}\label{problem-statement-28}

\textbf{Laporan sumber.} Tiongkok mengalami urbanisasi cepat dan
mempunyai banyak kota besar. Dalam pengantar, file menyebut 392 kota
dengan penduduk sedikitnya 300.000 jiwa pada 2014. Perkembangan ini
membentuk sistem regional seperti Beijing-Tianjin-Hebei, Delta Sungai
Yangtze, dan Delta Sungai Mutiara. Namun, perkembangan kota-kota sekitar
tidak seragam. Konsentrasi sumber daya di Beijing, termasuk
infrastruktur komunikasi, sekolah, dan rumah sakit, dikaitkan dengan
terbentuknya ``poverty belt'' di sekitar ibu kota. Sebaliknya, kota-kota
tingkat kedua dan ketiga di sekitar Shanghai dan Guangzhou disebut
berkembang lebih baik. (Hlm. 1, Bagian 1; TXT baris 66-85.)

\textbf{Laporan sumber.} Dalam kerangka NEG, kota inti dapat mendominasi
lanskap perkotaan dan menghambat pertumbuhan kota yang lebih kecil
melalui efek kompetisi atau \emph{agglomeration shadow}. Akan tetapi,
bukti empiris tentang keberadaan dan kekuatan bayangan tersebut
bercampur. Studi tentang Paris, Amerika Serikat, Eropa, dan Jepang yang
dirangkum artikel menemukan bahwa kedekatan dengan kota besar kadang
mendukung pertumbuhan, kadang mengikis pasar kota kecil, dan kadang
tidak menghasilkan hubungan yang stabil. (Hlm. 2-3, Bagian 2; TXT baris
89-145.)

\textbf{Laporan sumber.} Ukuran penduduk dan biaya transportasi saja
dipandang belum cukup untuk menjelaskan variasi tersebut. Artikel
membandingkan Beijing dan Shanghai yang pada 2015 disebut memiliki
populasi yang sama-sama sangat besar, sekitar 20 juta jiwa, tetapi
menunjukkan pola pengaruh yang berbeda terhadap kota sekitarnya. Artikel
lalu mengusulkan komposisi industri kota inti sebagai faktor yang
mungkin menjelaskan variasi efek limpahan dan \emph{backwash}. (Hlm.
2-3, Bagian 2; TXT baris 137-160.)

\textbf{Laporan sumber.} Penelitian terdahulu sudah membahas perbedaan
limpahan sektor manufaktur dan jasa, tetapi hubungan langsung antara
komposisi industri kota inti dan perkembangan sistem perkotaan regional
masih dinilai kurang diteliti, khususnya di negara berkembang seperti
Tiongkok. Literatur tentang \emph{agglomeration shadow} lebih banyak
berfokus pada Amerika Serikat dan Eropa; penelitian di Asia dan Tiongkok
disebut masih terbatas. (Hlm. 3, Bagian 2; TXT baris 161-192.)

\textbf{Interpretasi penulis.} Artikel mengajukan hipotesis bahwa
komposisi industri kota inti memengaruhi \emph{agglomeration shadow}
dan, melalui itu, sistem perkotaan regional. Jasa yang bersifat
informasional atau intensif secara sosial, misalnya keuangan dan ritel,
dianggap mempunyai biaya komunikasi yang lebih tinggi dalam arti luas.
Manufaktur diperkirakan lebih mampu menghasilkan jaringan produksi
regional melalui limpahan spasial. (Hlm. 3-4, akhir Bagian 2 dan awal
Bagian 3; TXT baris 193-209.)

\textbf{Inferensi pengolahan.} Problem artikel terdiri atas dua
pertanyaan yang berbeda: apakah komposisi industri berkaitan dengan
pertumbuhan kota sekitar, dan apakah hubungan itu cukup untuk
menyimpulkan adanya bayangan aglomerasi. Model dapat menguji asosiasi
pertumbuhan penduduk dengan komposisi industri, \emph{market potential},
dan interaksi keduanya. Namun, hubungan tersebut tidak dengan sendirinya
memisahkan kompetisi, limpahan produksi, pemusatan jasa, kebijakan kota,
atau faktor lokasi lain.

\subsection{Objectives}\label{objectives-28}

\textbf{Laporan sumber.} Tujuan yang dapat diidentifikasi dari
pendahuluan, metodologi, dan rancangan analisis adalah:

\begin{itemize}
\tightlist
\item
  menjelaskan mengapa \emph{agglomeration shadow} berbeda antarwilayah
  perkotaan melalui komposisi industri kota inti;
\item
  menguji hubungan komposisi industri kota inti dengan pertumbuhan
  penduduk kota-kota dalam sistem perkotaan Tiongkok;
\item
  membandingkan pengaruh pangsa nilai tambah manufaktur dan jasa
  terhadap hubungan antara kota inti dan kota sekitarnya;
\item
  menguji hubungan jangka pendek melalui blok lima tahun dan hubungan
  jangka lebih panjang melalui blok sepuluh tahun;
\item
  memeriksa apakah hasil berbeda pada batas spasial 100 km, 300 km, dan
  500 km;
\item
  menggunakan interaksi antara \emph{market potential} dan komposisi
  industri untuk menilai apakah sektor tertentu memperkuat atau
  melemahkan pengaruh kota inti; dan
\item
  membandingkan sistem perkotaan yang berpusat pada enam kota inti,
  dengan perhatian khusus pada perbedaan Beijing dan Shanghai.
\end{itemize}

Tujuan-tujuan tersebut dirumuskan dari pernyataan bahwa artikel ingin
menguji pengaruh komposisi industri terhadap \emph{agglomeration
shadow}, uraian model, serta pembagian analisis menurut waktu dan skala
spasial. Artikel tidak menyajikan daftar pertanyaan penelitian bernomor
yang terpisah. (Hlm. 1-4, Bagian 1-4; TXT baris 95-103, 193-209,
320-366.)

\subsection{Literature Survey}\label{literature-survey-28}

\textbf{Laporan sumber.} Artikel menempatkan gagasan \emph{agglomeration
shadow} dalam tradisi Alonso, yang menekankan bahwa kota kecil di
sekitar kota besar perlu dipahami sebagai bagian dari sistem perkotaan.
Kota kecil dapat memperoleh manfaat kedekatan dengan kota besar melalui
\emph{borrowed size}, tetapi manfaat itu tidak selalu muncul. Studi
tentang amenitas konsumsi di Eropa Barat Laut, misalnya, dirangkum
sebagai bukti bahwa kekuatan aglomerasi kota dominan dapat menghambat
pertumbuhan amenitas di kota kecil. (Hlm. 2, Bagian 2; TXT baris
105-126.)

\textbf{Laporan sumber.} Dalam NEG, bayangan aglomerasi berkembang
ketika pusat pertumbuhan mengindustrialisasi dan dominasi kota tertentu
mengorbankan kota lain di hinterland. Krugman menggunakan gagasan
\emph{market potential} untuk merepresentasikan daya tarik lokasi
manufaktur berdasarkan populasi dan jarak. Artikel juga merangkum
hubungan nonlinier antara jarak ke kota inti dan \emph{market
potential}, termasuk kemungkinan bentuk ``S'' dan zona bayangan pada
lokasi yang tidak cukup jauh dari pusat. (Hlm. 2, Bagian 2; TXT baris
127-160.)

\textbf{Laporan sumber.} Literatur yang diringkas memberi bukti yang
tidak seragam. Studi tentang kota-kota Amerika Serikat menemukan
bayangan di sekitar wilayah metropolitan tingkat tertinggi atau
menunjukkan pentingnya jarak bagi kota baru, tetapi tidak selalu bagi
kota yang sudah ada. Simulasi pembentukan megalopolis dan penelitian
tentang biaya transportasi menghasilkan kesimpulan bahwa bayangan dapat
kecil atau sulit ditemukan. Penelitian tentang Prancis, Eropa, dan
Jepang juga menunjukkan bahwa efek kedekatan dan bayangan bergantung
pada konteks serta industri. (Hlm. 2-3, Bagian 2; TXT baris 89-145,
177-187.)

\textbf{Laporan sumber.} Artikel menggunakan kerangka ``first nature''
dan ``second nature'' dari Krugman. \emph{First nature} merujuk pada
karakteristik intrinsik lokasi, sedangkan \emph{second nature}
bergantung pada struktur ekonomi dan hubungan antarkota. Dalam kerangka
itu, komposisi industri kota inti diperkirakan mengubah hubungan
antarkota melalui perbedaan biaya transportasi, biaya komunikasi,
jaringan produksi, dan pola penyebaran industri. (Hlm. 2 dan 10, Bagian
2 dan 8; TXT baris 134-160, 639-650.)

\textbf{Laporan sumber.} Untuk perbedaan sektor, artikel merangkum studi
yang menemukan bahwa manufaktur dan jasa menghasilkan limpahan yang
berbeda. Jasa produsen dianggap memperkuat konsentrasi dan fungsi
komando kota besar, sementara manufaktur dapat membentuk jaringan
produksi. Penulis menghubungkan jasa dengan \emph{hierarchical
diffusion}: jasa cenderung berpindah dari kota tingkat tertinggi ke kota
tingkat kedua, bukan kepada kota kecil yang sangat dekat. Manufaktur
dikaitkan dengan \emph{contagious diffusion} melalui jaringan produksi
regional. (Hlm. 3 dan 10-11, Bagian 2 dan 8; TXT baris 151-182,
651-690.)

\textbf{Laporan sumber.} Dalam konteks Tiongkok, artikel menyebut
penelitian tentang perubahan industri, konsentrasi jasa produsen,
perbedaan sistem Beijing-Tianjin-Hebei dan Delta Sungai Yangtze, serta
aglomerasi manufaktur. Penelitian sebelumnya disebut menunjukkan bahwa
kota dengan jasa produsen yang berkembang dapat tersebar berbeda dari
kota yang berfokus pada manufaktur, tetapi penelitian yang dirujuk
tentang jasa produsen lebih berfokus pada konsentrasi perusahaan
daripada dampak komposisi industri terhadap \emph{agglomeration shadow}.
(Hlm. 3, Bagian 2; TXT baris 161-192.)

\textbf{Interpretasi penulis.} Tinjauan pustaka dipakai untuk
menempatkan komposisi industri sebagai penjelasan atas perbedaan antara
limpahan positif dan efek \emph{backwash}. Manufaktur yang memerlukan
jaringan produksi diperkirakan menghasilkan manfaat bagi kota sekitar,
sedangkan jasa tingkat tinggi yang intensif dalam komunikasi dan
interaksi tatap muka diperkirakan memperdalam hierarki kota. (Hlm. 3-4
dan 10-11, Bagian 2 dan 8.)

\textbf{Inferensi pengolahan.} Tinjauan ini merupakan tinjauan naratif,
bukan tinjauan sistematis. File tidak menyebut strategi pencarian,
kriteria inklusi, rentang waktu korpus, jumlah studi yang ditelusuri,
atau prosedur penilaian mutu literatur. Karena itu, kelengkapan survei
dan kemungkinan bias pemilihan referensi tidak dapat dinilai dari TXT.

\subsection{Method Used}\label{method-used-28}

\textbf{Laporan sumber.} Penelitian menggunakan pendekatan kuantitatif
yang berangkat dari model lokasi perkotaan NEG. Unit analisis dinyatakan
sebagai kota \texttt{i}, kota inti \texttt{j}, dan periode \texttt{t}.
Variabel hasilnya adalah pertumbuhan penduduk kota antara periode
\texttt{t} dan \texttt{t-1}. Analisis dirancang untuk melihat pengaruh
komposisi industri kota inti melalui \emph{market potential}. (Hlm. 5-6,
Bagian 4; TXT baris 342-374.)

\textbf{Laporan sumber.} \emph{Market potential} disesuaikan dari ukuran
tradisional dengan hanya memasukkan kota inti yang dianggap paling
berpengaruh terhadap kota-kota menengah dan kecil di sekitarnya. Untuk
kota \texttt{i} dan kota inti \texttt{j}, rumus yang ditampilkan dalam
TXT berbentuk \texttt{MP\ =\ POP\ /\ Distance}, dengan \texttt{POP}
sebagai ukuran rata-rata populasi kota inti pada periode terkait dan
\texttt{Distance} sebagai jarak Euclidean antara pusat kota. Jarak yang
digunakan adalah jarak garis lurus, bukan biaya transportasi aktual.
(Hlm. 5-6, Bagian 4.1, Persamaan 1; TXT baris 320-338.)

\textbf{Laporan sumber.} Populasi rata-rata periode dipakai karena
menggunakan populasi tahun awal saja dianggap dapat meremehkan
\emph{market potential} ketika populasi kota inti meningkat cepat.
Artikel menyebut bahwa populasi enam kota inti meningkat rata-rata 11,9
juta jiwa antara 1980 dan 2015, atau lebih dari 1,7 juta jiwa setiap
lima tahun. Penetapan kota inti untuk suatu kota didasarkan pada kota
yang memberikan pengaruh terbesar, bukan semata-mata kota yang paling
dekat. (Hlm. 5-6, Bagian 4.1 dan 4.2; TXT baris 320-340, 383-400.)

\textbf{Laporan sumber.} Model regresi memasukkan populasi awal 1980,
\emph{market potential}, komposisi industri kota inti pada periode
sebelumnya, kuadrat komposisi industri, serta interaksi
\texttt{MP\ x\ IC}. Komposisi industri diukur secara terpisah melalui
pangsa nilai tambah industri/manufaktur dan pangsa nilai tambah yang
berkaitan dengan jasa. Kuadrat dimasukkan karena literatur yang dirujuk
mengindikasikan dampak struktur industri dapat nonlinier. (Hlm. 6,
Bagian 4.2, Persamaan 2; TXT baris 346-381.)

\textbf{Laporan sumber.} Analisis dilakukan untuk blok waktu lima tahun
dan sepuluh tahun. Untuk skala spasial, artikel menetapkan batas 100 km,
300 km, dan 500 km untuk memeriksa sistem perkotaan Tiongkok. Hasil juga
diestimasi dengan model \emph{fixed effects} (FE), \emph{random effects}
(RE), dan, untuk perbandingan sepuluh tahun, OLS. (Hlm. 6-9, Bagian 4.2
dan 6; TXT baris 342-367, 437-470, 474-597.)

\textbf{Laporan sumber.} Artikel berfokus pada manufaktur dan jasa,
bukan pertanian. Alasannya adalah bobot industri Tiongkok dianggap
besar, kota inti didorong menjadi kota dunia melalui jasa terutama jasa
produsen, dan pangsa pertanian di kota inti disebut sangat kecil.
Penulis juga menyebut bahwa modal manusia penting secara teoritis,
tetapi pangsanya dianggap kecil pada 1980 dan sebagian efeknya mungkin
tertangkap oleh populasi karena jumlah pekerja terampil berkorelasi
dengan ukuran penduduk kota. (Hlm. 6, Bagian 4.2; TXT baris 346-362.)

\textbf{Interpretasi penulis.} \emph{Market potential} diperlakukan
sebagai saluran utama yang menghubungkan komposisi industri kota inti
dengan pertumbuhan kota sekitar. Interaksi positif untuk manufaktur
dibaca sebagai penguatan efek kota inti, sedangkan interaksi negatif
untuk jasa dibaca sebagai pelemahan limpahan dan kemungkinan
terbentuknya bayangan. (Hlm. 7-10, Bagian 6-8; TXT baris 425-470,
600-690.)

\textbf{Inferensi pengolahan.} Desain ini merupakan analisis
observasional atas perubahan populasi lintas kota dan periode, bukan
eksperimen atau kuasi-eksperimen. Penggunaan FE/RE tidak dengan
sendirinya mengidentifikasi sebab-akibat. File tidak menyajikan
penjelasan tentang pemilihan antara FE dan RE, pemeriksaan endogenitas,
instrumen, atau desain kontrafaktual. Karena itu, istilah ``pengaruh''
dalam artikel perlu dibaca sebagai hubungan model dengan kualifikasi,
kecuali ketika bagian ini secara eksplisit menyebutnya sebagai
interpretasi penulis.

\subsection{Dataset}\label{dataset-28}

\textbf{Laporan sumber.} Data penduduk berasal dari Department of
Economic and Social Affairs Perserikatan Bangsa-Bangsa. Kriteria
pemilihan kota adalah sedikitnya 300.000 penduduk pada 2015, yang
menurut artikel menghasilkan sekitar 392 kota di Tiongkok. Data berisi
angka penduduk tahunan kota yang dikumpulkan setiap lima tahun dari 1980
sampai 2015. (Hlm. 4, Bagian 3.3; TXT baris 239-247.)

\textbf{Laporan sumber.} Tabel 1 melaporkan bahwa pada 2015 rerata
ukuran kota adalah 1.248 ribu jiwa dan mediananya 607 ribu jiwa. Tabel
tersebut juga melaporkan 153 kota berpenduduk 0,5-1 juta, 86 kota
berpenduduk 1-5 juta, 10 kota berpenduduk 5-10 juta, dan 6 kota
berpenduduk lebih dari 10 juta. Deret pada tabel mencakup 1980, 1990,
2000, 2010, dan 2015. (Hlm. 5, Tabel 1; TXT baris 300-307.)

\textbf{Laporan sumber.} Gambar 1 memperlihatkan perubahan populasi
total, populasi perkotaan, dan rasio urbanisasi Tiongkok pada 1978-2015.
Gambar 2 menampilkan pola spasial perkotaan Tiongkok pada 2015. Artikel
menyatakan bahwa keenam kota inti memiliki populasi lebih dari 10 juta
pada 2015. (Hlm. 3-5, Bagian 3.1-3.2, Gambar 1-2; TXT baris 193-231,
248-296.)

\textbf{Laporan sumber.} Tabel 3 memberikan statistik deskriptif untuk
populasi kota inti, populasi kota noninti, pangsa nilai tambah industri
di kota inti pada 2010, pangsa nilai tambah jasa di kota inti pada 2010,
pertumbuhan populasi kota inti, pertumbuhan populasi kota noninti, dan
jarak terpendek ke kota inti. Rerata yang ditampilkan adalah 15.312,33
ribu jiwa untuk populasi kota inti, 1.029,23 ribu jiwa untuk populasi
kota noninti, 38,37\% untuk pangsa industri, 54,74\% untuk pangsa jasa,
11.983,83 ribu jiwa untuk pertumbuhan kota inti, 800,32 ribu jiwa untuk
pertumbuhan kota noninti, dan 418,23 km untuk jarak terpendek. (Hlm. 7,
Tabel 3; TXT baris 407-422.)

\textbf{Laporan sumber.} Data ketersediaan pada bagian akhir menyatakan
bahwa penelitian memakai data populasi kota Tiongkok dari 1950 sampai
2015 dan bahwa data tersebut disimpan di repositori \emph{World
Urbanization Prospects} milik United Nations Department of Economic and
Social Affairs, Population Division. (Hlm. 10, Data Availability; TXT
baris 631-639.)

\textbf{Tidak disebutkan dalam file.} Sumber rinci, seri statistik, dan
prosedur pengumpulan untuk data pangsa nilai tambah industri dan jasa
tidak dijelaskan dalam bagian pengumpulan data. File juga tidak
menjelaskan sumber koordinat pusat kota, basis perhitungan pembagian
wilayah pengaruh, aturan pembersihan data penduduk, atau apakah data
penduduk 1950-1975 hanya tersedia untuk konteks dan bukan untuk estimasi
utama.

\textbf{Inferensi pengolahan.} Ada perbedaan cakupan waktu internal:
Bagian 3.3 menyebut data penelitian 1980-2015, sedangkan bagian Data
Availability menyebut data 1950-2015. Perbedaan itu tidak dapat
diselesaikan hanya dari TXT. Dalam review ini, 1980-2015 dipakai sebagai
periode analisis utama karena itulah yang dijelaskan pada bagian
pemilihan data dan model; pernyataan 1950-2015 dipertahankan sebagai
keterangan ketersediaan data.

\subsection{Population / Sample}\label{population-sample-28}

\textbf{Laporan sumber.} Populasi kajian dibentuk oleh kota-kota di
Tiongkok yang memenuhi ambang sedikitnya 300.000 penduduk pada 2015.
Dari sistem tersebut, enam kota dengan populasi di atas 10 juta pada
2015 dipilih sebagai kota inti: Beijing, Shanghai, Tianjin, Chongqing,
Guangzhou, dan Shenzhen. Empat kota pertama berstatus munisipalitas di
bawah pemerintah pusat dan dua kota terakhir dijelaskan sebagai pusat
penting perdagangan luar negeri atau zona ekonomi khusus. (Hlm. 4-5,
Bagian 3.2-3.3; TXT baris 248-269.)

\textbf{Laporan sumber.} Artikel menyatakan bahwa enam kota itu dipilih
karena dominan dalam sistem regionalnya dan karena banyak kota kecil
serta menengah tumbuh di sekitarnya. Pengaruhnya dianggap menjangkau
Tiongkok utara, selatan, timur, dan barat. Shenzhen diperlakukan berbeda
menurut waktu karena pada awal reformasi belum lama menjadi kota inti
dan pada periode awal hanya sedikit memengaruhi kota sekitarnya. (Hlm. 5
dan 7, Bagian 3.2 dan 5.1; TXT baris 262-269, 390-400.)

\textbf{Laporan sumber.} Jumlah kota sekitar tiap kota inti berubah
menurut blok waktu. Pada 1980-1990, jumlahnya adalah Beijing 138,
Shanghai 161, Guangzhou 44, Tianjin 5, Chongqing 38, dan Shenzhen 0.
Pada 1990-2000, jumlahnya adalah 129, 156, 50, 3, 45, dan 3. Pada
2000-2010, jumlahnya adalah 126, 159, 42, 1, 47, dan 11. Pada 2010-2015,
jumlahnya adalah 132, 155, 41, 1, 45, dan 12. (Hlm. 5, Tabel 2; TXT
baris 311-317.)

\textbf{Laporan sumber.} Tabel 4, Tabel 5, dan Tabel 8 masing-masing
melaporkan 2.702 observasi. Tabel 6 dan Tabel 7, yang membagi analisis
lima tahun menurut skala, melaporkan 162 observasi untuk 100 km, 814
untuk 300 km, dan 1.493 untuk 500 km. File tidak menjelaskan secara
eksplisit bagaimana jumlah observasi itu dipetakan menjadi jumlah kota,
periode, dan kemungkinan pengulangan kota. (Hlm. 8-10, Tabel 4-8; TXT
baris 474-597.)

\textbf{Inferensi pengolahan.} Ini bukan sampel acak. Pemilihan kota
inti berbasis ambang populasi dan pertimbangan dominasi sistem
perkotaan, sedangkan kota noninti berasal dari cakupan kota yang
memenuhi kriteria data. Karena kota inti serta sistem pengaruhnya
dipilih dari konteks Tiongkok, hasil terutama berlaku untuk populasi
urban yang tercakup dalam desain itu, bukan otomatis untuk semua kota
atau sistem metropolitan di luar konteks tersebut.

\textbf{Inferensi pengolahan.} Perubahan jumlah kota sekitar antara blok
waktu menunjukkan bahwa pembagian kota ke dalam sistem pengaruh bersifat
dinamis. Namun, file tidak memberi algoritme atau nilai ambang yang
digunakan untuk menentukan kota yang paling dipengaruhi oleh suatu kota
inti. Hal ini membatasi replikasi dan penilaian apakah perubahan jumlah
kota mencerminkan perubahan pengaruh, perubahan data, atau keduanya.

\subsection{Dependent Variables}\label{dependent-variables-28}

\textbf{Laporan sumber.} Variabel terikat utama dalam Persamaan 2 adalah
pertumbuhan populasi kota \texttt{DeltaPOP\_it}, yaitu perubahan
populasi kota \texttt{i} antara periode \texttt{t} dan \texttt{t-1}.
Model menjelaskan pertumbuhan kota tersebut dengan populasi awal 1980,
\emph{market potential}, komposisi industri kota inti pada periode
sebelumnya, kuadrat komposisi industri, dan interaksi \emph{market
potential} dengan komposisi industri. (Hlm. 6, Bagian 4.2, Persamaan 2;
TXT baris 363-381.)

\textbf{Laporan sumber.} Persamaan dan tabel membedakan dua ukuran
komposisi industri: \texttt{ICm} untuk pangsa manufaktur/industri dan
\texttt{ICs} untuk pangsa jasa. Masing-masing dipakai sebagai variabel
linear, kuadrat, serta komponen dalam interaksi dengan \texttt{MP}.
\emph{Market potential} sendiri merupakan variabel penjelas yang
dibangun dari populasi rata-rata kota inti dibagi jarak Euclidean. (Hlm.
5-6, Bagian 4.1-4.2, Tabel 4-8; TXT baris 323-338, 364-381, 474-597.)

\textbf{Tidak disebutkan dalam file.} Satuan numerik
\texttt{DeltaPOP\_it} dalam setiap regresi tidak dijelaskan secara
konsisten. Tabel 3 menggunakan satuan ribu jiwa untuk statistik populasi
dan pertumbuhan, tetapi rumus regresi tidak menyatakan secara eksplisit
apakah seluruh variabel estimasi memakai jiwa, ribu jiwa, transformasi,
atau skala lain.

\textbf{Inferensi pengolahan.} Tidak ada variabel terikat yang secara
langsung mengukur intensitas \emph{agglomeration shadow}. Bayangan
merupakan interpretasi atas tanda dan signifikansi interaksi, bentuk
nonlinier, pola jarak, serta perbandingan pertumbuhan kota sekitar.
Dengan demikian, hasil regresi lebih tepat disebut menjelaskan
pertumbuhan penduduk bersyarat pada hubungan kota inti-kota sekitar
daripada mengestimasi ukuran bayangan secara langsung.

\textbf{Inferensi pengolahan.} Karena komposisi industri diukur dengan
pangsa nilai tambah dan bukan ukuran limpahan aktual, variabel tersebut
tidak secara langsung mengobservasi jaringan produksi, perjalanan
bisnis, arus komuter, hubungan pemasok-pelanggan, konsumsi jasa, atau
penggunaan pasar kota inti oleh penduduk kota noninti. Kesesuaian
konstruk antara ukuran komposisi dan mekanisme \emph{borrowed size} atau
bayangan karena itu terbatas.

\subsection{Results}\label{results-28}

\textbf{Laporan sumber: perubahan ukuran dan komposisi kota inti.}
Populasi enam kota inti tetap besar setelah reformasi, sementara
komposisi industrinya berubah. Rerata pangsa nilai tambah jasa meningkat
dari 28,1\% pada 1980 menjadi 54,7\% pada 2010. Pada 2010, pangsa jasa
melampaui 50\% di empat dari enam kota inti; Beijing mencapai 75,1\%,
dan narasi sumber menyebut angka itu hampir 20\% lebih tinggi daripada
Shanghai. Pada tahun yang sama, pangsa manufaktur hanya melampaui pangsa
jasa di Tianjin dan Chongqing, dan keduanya berada dekat 50\%. (Hlm. 7,
Bagian 5.1 dan Tabel 3; TXT baris 368-389.)

\textbf{Laporan sumber: konteks jarak dan sektor.} Jarak rata-rata
antara kota inti yang paling dekat adalah 418,2 km, dengan rentang jarak
terpendek ke kota inti dari 18,24 km sampai 2.724,71 km. Artikel memberi
Foshan sebagai contoh kota kurang dari 50 km dari Guangzhou dan Yining
sebagai contoh kota lebih dari 1.000 km dari Chongqing. Pertambangan
dimasukkan dalam sektor industri, tetapi kehadirannya di kota inti
disebut kecil. (Hlm. 7, Bagian 5.1 dan Tabel 3; TXT baris 390-403,
417-422.)

\textbf{Laporan sumber: manufaktur pada blok lima tahun.} Pada Tabel 4,
\texttt{MP} positif dan signifikan pada Model 1 untuk FE dan RE, dengan
koefisien masing-masing 0,8926 dan 0,9821 serta nilai dalam tanda kurung
0,000. Pangsa manufaktur \texttt{ICm} negatif dan signifikan pada Model
1, masing-masing -1,4962 untuk FE dan -1,2569 untuk RE. Setelah kuadrat
\texttt{ICm} dimasukkan pada Model 2, koefisien linear menjadi positif
dan koefisien kuadrat negatif: 3,6674 dan -0,0545 pada FE, serta 4,2983
dan -0,0586 pada RE; nilai dalam tanda kurung untuk komponen tersebut
ditampilkan 0,000. (Hlm. 8, Tabel 4; TXT baris 474-495.)

\textbf{Laporan sumber: manufaktur dan interaksi.} Pada Model 3 Tabel 4,
interaksi \texttt{MP\ x\ ICm} positif dan ditampilkan signifikan, dengan
koefisien 0,0441 pada FE dan 0,0226 pada RE serta nilai dalam tanda
kurung 0,000. Pada model yang sama, \texttt{MP} tidak lagi signifikan:
-0,3714 pada FE dengan nilai 0,146 dan 0,1799 pada RE dengan nilai
0,423. Narasi artikel menyatakan bahwa interaksi tersebut menunjukkan
pangsa manufaktur memperkuat efek positif \emph{market potential}
terhadap pertumbuhan penduduk kota sekitar. (Hlm. 7-8, Bagian 6.1 dan
Tabel 4; TXT baris 437-455, 480-492.)

\textbf{Laporan sumber: jasa pada blok lima tahun.} Tabel 5 menunjukkan
\texttt{MP} positif dan signifikan pada model-model jasa. Pada Model 3,
koefisien \texttt{MP} adalah 4,1945 pada FE dan 2,3634 pada RE, keduanya
dengan nilai dalam tanda kurung 0,000. Pangsa jasa \texttt{ICs}
mempunyai koefisien positif, sedangkan \texttt{ICs2} mempunyai koefisien
negatif dan ditampilkan signifikan. Interaksi \texttt{MP\ x\ ICs}
negatif dan ditampilkan signifikan, dengan koefisien -0,0488 pada FE dan
-0,0243 pada RE, keduanya dengan nilai dalam tanda kurung 0,000. (Hlm.
8, Tabel 5; TXT baris 467-521.)

\textbf{Interpretasi penulis: perbandingan sektor.} Penulis menafsirkan
hasil manufaktur sebagai bukti bahwa manufaktur memperkuat hubungan
\emph{market potential} dengan pertumbuhan kota sekitarnya, sehingga
bayangan aglomerasi mungkin tidak muncul atau lebih lemah. Hasil jasa
ditafsirkan sebagai bukti bahwa ekspansi jasa menekan efek \emph{market
potential} kota inti terhadap pertumbuhan kota sekitar. (Hlm. 7-8,
Bagian 6.1-6.2; TXT baris 425-455 dan 467-470.)

\textbf{Laporan sumber: skala spasial lima tahun.} Pada Tabel 6,
interaksi \texttt{MP\ x\ ICm} positif dan ditampilkan signifikan pada
100 km, 300 km, dan 500 km, baik dalam FE maupun RE. Koefisien FE yang
ditampilkan adalah 0,1191, 0,0587, dan 0,0503; koefisien RE adalah
0,0967, 0,0295, dan 0,0270. Nilai \texttt{R2} pada model-model tersebut
adalah 0,0460, 0,0799, dan 0,0850 untuk FE, serta 0,2878, 0,4105, dan
0,4110 untuk RE. (Hlm. 8, Tabel 6; TXT baris 526-547.)

\textbf{Laporan sumber: skala spasial jasa.} Pada Tabel 7 yang berjudul
regresi pangsa jasa, interaksi yang ditampilkan bernilai negatif dan
bertanda signifikansi pada ketiga skala 100 km, 300 km, dan 500 km, baik
untuk FE maupun RE. Untuk FE, baris interaksi terbaca \texttt{-01023},
\texttt{-0.0604}, dan \texttt{-0.0517}; untuk RE, terbaca
\texttt{-0.0732}, \texttt{-0.0298}, dan \texttt{-0.0257}. Teks ekstraksi
tidak memiliki titik desimal pada angka pertama, sehingga angka tersebut
tidak dinormalisasi dalam review ini. \texttt{MP} ditampilkan positif
dan signifikan pada keenam model. (Hlm. 9, Tabel 7; TXT baris 551-572.)

\textbf{Inferensi pengolahan: anomali Tabel 7.} Judul Tabel 7 menyebut
pangsa jasa, tetapi baris variabel dalam hasil ekstraksi terbaca
\texttt{ICm} dan \texttt{ICm2}, bukan \texttt{ICs} dan \texttt{ICs2}.
Narasi Bagian 6.2 menyebut hasil itu sebagai hasil jasa. Review
mengikuti judul dan narasi sumber, tetapi mempertahankan anomali notasi
karena TXT tidak memberi dasar untuk memperbaikinya.

\textbf{Laporan sumber: kestabilan hubungan nonlinier menurut skala.}
Artikel menyatakan bahwa pengaruh komposisi industri pada \emph{market
potential} tetap terlihat pada skala yang berbeda, dan efeknya menjadi
lebih signifikan pada skala yang lebih kecil. Namun, hubungan ``inverse
U'' antara pangsa manufaktur atau jasa dan pertumbuhan kota sekitar
tidak stabil, terutama pada skala kecil. Untuk manufaktur, hubungan
tersebut tidak tetap signifikan; untuk jasa, hubungan itu juga tidak
cukup stabil bahkan sampai 300 km. (Hlm. 7, Bagian 6.1-6.2; TXT baris
437-450.)

\textbf{Laporan sumber: blok sepuluh tahun.} Tabel 8 menunjukkan bahwa
interaksi manufaktur \texttt{MP\ x\ IC} positif dan signifikan hanya
pada RE, dengan koefisien 0,0317 dan nilai 0,009. Pada OLS, koefisiennya
0,0132 dengan nilai 0,281; pada FE, -0,0102 dengan nilai 0,593.
Interaksi jasa negatif pada ketiga model: -0,0283 pada OLS dengan nilai
0,010, -0,0344 pada FE dengan nilai 0,045, dan -0,0439 pada RE dengan
nilai 0,000. (Hlm. 9, Tabel 8; TXT baris 576-597.)

\textbf{Laporan sumber: blok sepuluh tahun dan populasi awal.} Pada
Tabel 8, koefisien populasi awal 1980 positif dan ditampilkan signifikan
untuk model manufaktur maupun jasa. Artikel menyatakan bahwa pengaruh
populasi awal lebih kuat dalam rentang sepuluh tahun untuk kedua
industri dan menyimpulkan bahwa hukum Gibrat ditolak untuk Tiongkok.
Penulis juga menyatakan bahwa dalam model RE koefisien interaksi pada
periode sepuluh tahun lebih besar daripada pada periode lima tahun,
sehingga dampak komposisi industri terhadap \emph{market potential}
tidak menghilang dan mungkin menguat sepanjang waktu. (Hlm. 9-10, Tabel
8 dan Bagian 6.3; TXT baris 453-470, 600-602.)

\textbf{Laporan sumber: bentuk puncak dan jarak.} Artikel menyatakan
bahwa pangsa manufaktur dan jasa sama-sama memperlihatkan hubungan
\emph{inverse U}, sehingga pangsa rendah maupun tinggi kurang positif
bagi pertumbuhan kota sekitar. Puncak jasa disebut berada sekitar
100-300 km dari kota inti, sedangkan puncak industri sekitar 300 km atau
lebih. Ketika suku kubik ditambahkan, efeknya masih signifikan untuk
jasa tetapi tidak signifikan untuk industri. (Hlm. 10-11, Bagian 8; TXT
baris 675-690.)

\textbf{Inferensi pengolahan.} Klaim robustitas perlu dibatasi pada
tanda interaksi dan perbandingan sektor yang berulang pada beberapa
skala dan periode. Bukti untuk bentuk \emph{inverse U} lebih lemah
karena artikel sendiri menyebut hubungan itu tidak stabil pada skala
kecil, dan hasil interaksi manufaktur sepuluh tahun berubah antara OLS,
FE, dan RE. Nilai \texttt{R2} FE yang rendah pada sejumlah model dan
perbedaan hasil antarpendekatan estimasi juga menunjukkan bahwa
kesimpulan substantif bergantung pada spesifikasi model.

\subsection{Findings}\label{findings-28}

\textbf{Temuan yang dilaporkan sumber.} Komposisi industri kota inti
berkaitan dengan pertumbuhan penduduk kota sekitar melalui \emph{market
potential}. Pada lima tahun, interaksi pangsa manufaktur dengan
\emph{market potential} positif, sedangkan interaksi pangsa jasa
negatif. Pola tanda interaksi itu diulang pada batas 100 km, 300 km, dan
500 km. (Hlm. 7-9, Bagian 6.1-6.3, Tabel 4-8.)

\textbf{Temuan yang dilaporkan sumber.} Pangsa industri dan jasa tidak
hanya diperlakukan sebagai hubungan linear. Model keseluruhan dan model
sepuluh tahun memperlihatkan komponen linear positif serta kuadrat
negatif, yang oleh penulis dibaca sebagai hubungan \emph{inverse U}.
Akan tetapi, stabilitas hubungan itu berkurang pada skala spasial yang
lebih kecil. (Hlm. 7-10, Bagian 6.1-6.3; TXT baris 431-470.)

\textbf{Interpretasi penulis.} Manufaktur dipandang mengikuti
\emph{contagious diffusion} karena jaringan produksi regional berkembang
di sekitar kota inti untuk menghemat biaya transportasi. Jasa dipandang
mengikuti \emph{hierarchical diffusion} karena waktu dan komunikasi
lebih penting, sehingga jasa tingkat tinggi cenderung berpindah dari
kota teratas ke kota tingkat kedua, bukan ke kota kecil yang berdekatan.
(Hlm. 10-11, Bagian 8; TXT baris 651-673.)

\textbf{Interpretasi penulis.} Perbedaan komposisi industri dipakai
untuk menjelaskan mengapa Beijing lebih mungkin menciptakan bayangan
aglomerasi dibandingkan Shanghai. Beijing disebut lebih berfokus pada
pembangunan berbasis jasa, sedangkan pangsa manufaktur Shanghai dan
Guangzhou pada 2015 masih lebih dari 10 poin persentase lebih rendah
daripada Beijing. Artikel mengaitkan fokus jasa Beijing dengan limpahan
geografis yang lebih rendah. (Hlm. 10, Bagian 7; TXT baris 605-621.)

\textbf{Temuan yang dilaporkan sumber.} Pertumbuhan penduduk kota inti
dan kota noninti tidak mengikuti implikasi hukum Gibrat menurut
pembacaan penulis. Populasi awal mempunyai hubungan positif dan lebih
kuat pada rentang waktu sepuluh tahun, sehingga ukuran awal kota tetap
penting dalam pertumbuhan berikutnya. (Hlm. 9-10, Bagian 6.3; TXT baris
600-602.)

\textbf{Inferensi pengolahan.} Bukti empiris artikel paling langsung
mendukung pernyataan bahwa komposisi sektor memoderasi asosiasi antara
\emph{market potential} kota inti dan pertumbuhan kota sekitar. Bukti
untuk mekanisme \emph{agglomeration shadow} lebih tidak langsung, karena
tidak ada ukuran arus ekonomi, penggunaan layanan, kompetisi pasar,
jaringan perusahaan, atau perubahan lokasi industri yang menghubungkan
tanda koefisien dengan mekanisme secara langsung.

\textbf{Inferensi pengolahan.} Pernyataan bahwa jasa mempunyai ``efek
negatif'' harus dibaca sebagai efek interaksi dan bukan sebagai klaim
bahwa semua jasa secara absolut mengurangi pertumbuhan. Dalam Tabel 5,
\texttt{MP} dan \texttt{ICs} sendiri ditampilkan positif pada beberapa
spesifikasi, sementara \texttt{MP\ x\ ICs} negatif. Artinya, yang
didukung oleh model adalah pelemahan hubungan \emph{market
potential}-pertumbuhan ketika pangsa jasa meningkat, bukan penurunan
pertumbuhan dalam setiap kondisi.

\textbf{Inferensi pengolahan.} Perbandingan Beijing-Shanghai bersifat
penjelasan berbasis komposisi dan narasi model. TXT tidak menampilkan
estimasi terpisah yang secara langsung menguji perbedaan kausal kedua
kota tersebut dengan kontrol yang sama. Karena itu, perbandingan
tersebut tidak boleh diperlakukan sebagai eksperimen komparatif atau
atribusi tunggal kepada pangsa jasa.

\subsection{Conclusions}\label{conclusions-28}

\textbf{Kesimpulan yang dinyatakan penulis.} Artikel menyimpulkan bahwa
analisis empiris memberi bukti bahwa manufaktur, secara keseluruhan,
lebih mendorong pertumbuhan kota sekitar daripada jasa melalui limpahan.
Komposisi industri tidak berdampak langsung pada kota sekitar, tetapi
bekerja melalui \emph{market potential}. Penulis menempatkan mekanisme
ini dalam gagasan \emph{first nature} dan \emph{second nature}, dengan
hubungan ekonomi antarkota sebagai bagian dari \emph{second nature}.
(Hlm. 10, Bagian 8; TXT baris 626-650.)

\textbf{Kesimpulan yang dinyatakan penulis.} Untuk manufaktur, jaringan
produksi regional dianggap berkembang di sekitar kota inti dan mengikuti
difusi menular. Untuk jasa, waktu komunikasi dianggap lebih penting
daripada biaya transportasi. Kota kecil dekat kota inti dinilai tidak
selalu mempunyai keunggulan dibandingkan kota tingkat kedua yang lebih
jauh, sehingga limpahan jasa kepada kota yang dekat dengan inti menjadi
lemah. (Hlm. 10-11, Bagian 8; TXT baris 651-673.)

\textbf{Kesimpulan yang dinyatakan penulis.} Pangsa manufaktur dan jasa
sama-sama mempunyai hubungan \emph{inverse U} dengan pertumbuhan kota
sekitar, tetapi lokasi puncaknya berbeda. Puncak jasa berada sekitar
100-300 km dari kota inti, sedangkan puncak industri berada sekitar 300
km atau lebih. Suku kubik tetap signifikan untuk jasa dan tidak
signifikan untuk industri. (Hlm. 11, Bagian 8; TXT baris 675-690.)

\textbf{Kesimpulan yang dinyatakan penulis.} Temuan tersebut dianggap
penting untuk memahami pembangunan sistem perkotaan regional yang tidak
merata dan untuk pembangunan kawasan metropolitan di Tiongkok. Artikel
menekankan bahwa sistem yang berpusat pada kota dengan komposisi jasa
yang kuat lebih mungkin menciptakan bayangan dibandingkan sistem yang
mempunyai peran manufaktur lebih besar. (Hlm. 1 dan 10-11, abstrak dan
Bagian 8.)

\textbf{Inferensi pengolahan.} Kesimpulan tentang arah interaksi cukup
konsisten di beberapa tabel, tetapi kekuatan kesimpulan tentang
mekanisme dan bentuk kurva perlu dibatasi oleh perbedaan FE, RE, dan
OLS, ketidakstabilan \emph{inverse U} pada skala kecil, serta tidak
adanya ukuran langsung untuk bayangan. Artikel lebih meyakinkan sebagai
studi asosiasi sektor dan pertumbuhan daripada sebagai pengujian kausal
lengkap atas \emph{agglomeration shadow}.

\subsection{Contributions}\label{contributions-28}

\textbf{Kontribusi yang diklaim atau ditunjukkan sumber.}

\begin{itemize}
\tightlist
\item
  \textbf{Kontribusi substantif:} artikel memindahkan perhatian dari
  ukuran kota dan biaya transportasi saja kepada komposisi industri kota
  inti sebagai penjelasan variasi \emph{agglomeration shadow}. (Hlm.
  2-4, Bagian 2-3.)
\item
  \textbf{Kontribusi teoritis:} artikel menghubungkan NEG, \emph{market
  potential}, \emph{first nature/second nature}, dan perbedaan jaringan
  manufaktur serta jasa dalam satu kerangka empiris. (Hlm. 2 dan 10,
  Bagian 2 dan 8.)
\item
  \textbf{Kontribusi empiris:} artikel membandingkan enam kota inti di
  Tiongkok dan kota noninti yang berada dalam sistem pengaruhnya, bukan
  hanya satu wilayah metropolitan. (Hlm. 4-7, Bagian 3 dan 5.)
\item
  \textbf{Kontribusi komparatif:} artikel membedakan hasil manufaktur
  dan jasa, menguji jangka lima dan sepuluh tahun, serta membandingkan
  tiga skala spasial. (Hlm. 6-9, Bagian 4 dan 6.)
\item
  \textbf{Kontribusi metodologis:} artikel memasukkan pangsa industri,
  kuadrat pangsa, dan interaksi dengan \emph{market potential} ke dalam
  model pertumbuhan penduduk; \emph{market potential} juga memakai
  populasi rata-rata periode dan kota inti yang paling berpengaruh.
  (Hlm. 5-6, Bagian 4.)
\item
  \textbf{Kontribusi pada perdebatan Beijing-Shanghai:} artikel
  menawarkan komposisi industri sebagai penjelasan mengapa dua kota inti
  dengan ukuran populasi yang sama-sama besar dapat mempunyai limpahan
  regional yang berbeda. (Hlm. 2-3 dan 10, Bagian 2 dan 7.)
\item
  \textbf{Kontribusi kebijakan yang dinyatakan:} artikel menyediakan
  dasar untuk mempertimbangkan ketimpangan perkembangan sistem perkotaan
  dan pembangunan kawasan metropolitan berdasarkan perbedaan struktur
  industri. (Hlm. 1 dan 10-11, abstrak dan Bagian 8.)
\end{itemize}

\textbf{Inferensi pengolahan.} Kontribusi paling kuat dari artikel
adalah menunjukkan bahwa analisis kota inti-kota sekitar dapat
menghasilkan kesimpulan berbeda ketika komposisi sektor dan interaksi
dengan \emph{market potential} diperhitungkan. Klaim kontribusi
mekanisme masih bersifat eksploratif karena indikator manufaktur dan
jasa tidak mengukur arus limpahan atau bayangan secara langsung.

\subsection{Research Gap}\label{research-gap-28}

\textbf{Kesenjangan yang dinyatakan sumber sebelum penelitian.}

\begin{itemize}
\tightlist
\item
  Belum jelas mengapa \emph{agglomeration shadow} kuat pada sebagian
  sistem perkotaan dan kecil atau tidak ada pada sistem lain. (Hlm. 2-3,
  Bagian 2.)
\item
  Hubungan antara komposisi industri kota inti dan perkembangan sistem
  perkotaan regional masih kurang diteliti. (Hlm. 3, Bagian 2.)
\item
  Bukti tentang perbedaan limpahan manufaktur dan jasa telah ada, tetapi
  belum banyak dipakai untuk menjelaskan variasi bayangan aglomerasi.
  (Hlm. 3, Bagian 2.)
\item
  Kajian \emph{agglomeration shadow} lebih terkonsentrasi pada Amerika
  Serikat dan Eropa; bukti dari negara berkembang seperti Tiongkok masih
  terbatas. (Hlm. 3, Bagian 2.)
\item
  Studi Tiongkok yang dirujuk tentang jasa produsen lebih menelaah
  konsentrasi perusahaan daripada pengaruh komposisi industri kota inti
  terhadap pertumbuhan kota sekitar. (Hlm. 3, Bagian 2.)
\item
  Perubahan sektor perlu diperhitungkan untuk menjelaskan pergerakan
  penduduk perkotaan kontemporer, karena kesenjangan antara teori dan
  realitas empiris dianggap meningkat. (Hlm. 3, Bagian 2.)
\end{itemize}

\textbf{Inferensi pengolahan: kesenjangan yang tersisa setelah studi.}

\begin{itemize}
\tightlist
\item
  Belum ada pengukuran langsung tentang arus barang, jasa, komuter,
  investasi, jaringan pemasok, atau penggunaan pasar yang dapat menguji
  mekanisme limpahan dan \emph{borrowed size}.
\item
  Belum ada desain kontrafaktual atau strategi identifikasi yang dapat
  memisahkan pengaruh komposisi industri dari kebijakan, ukuran awal
  kota, aksesibilitas, keunggulan lokasi, dan seleksi lokasi industri.
\item
  Analisis belum memperlakukan seluruh struktur polisentris secara
  eksplisit; hanya kota inti yang dianggap paling berpengaruh dimasukkan
  dalam ukuran \emph{market potential}.
\item
  Belum jelas bagaimana hasil berubah jika manufaktur, jasa, komunikasi,
  keuangan, ritel, dan jasa produsen dipecah menjadi subsektor yang
  lebih rinci.
\item
  Belum ada uji yang cukup untuk kestabilan ambang 100, 300, dan 500 km,
  definisi kota inti, atau aturan penugasan kota noninti kepada kota
  inti.
\item
  Belum ada verifikasi apakah pangsa nilai tambah mewakili kapasitas,
  produktivitas, atau limpahan aktual ke kota sekitarnya.
\end{itemize}

\subsection{Limitations}\label{limitations-28}

\textbf{Keterbatasan yang dinyatakan melalui keputusan desain sumber.}
Artikel menyatakan bahwa ukuran populasi dan jarak dipandang cukup untuk
tujuan pengukuran \emph{market potential} yang bukan penilaian
konseptual. Upah tidak dimasukkan karena berubah dramatis setelah
reformasi sehingga tingkat pendapatan tidak dianggap sebanding. Artikel
juga memusatkan perhatian pada manufaktur dan jasa karena pertanian
mempunyai pangsa kecil di kota inti, serta mengakui secara teoritis
pentingnya modal manusia tetapi menggunakan populasi sebagai pendekatan
yang mungkin menangkap sebagian efeknya. (Hlm. 5-6, Bagian 4.1-4.2; TXT
baris 262-269, 346-362.)

\textbf{Keterbatasan yang diturunkan dari isi TXT sebagai inferensi
pengolahan.}

\begin{itemize}
\tightlist
\item
  \textbf{Validitas konstruk:} pertumbuhan penduduk merupakan proksi
  hasil sistem perkotaan, bukan ukuran langsung \emph{agglomeration
  shadow}. Tanda interaksi tidak menunjukkan besaran penduduk atau
  fungsi ekonomi yang benar-benar ``terbayang''.
\item
  \textbf{Pengukuran komposisi:} pangsa nilai tambah industri dan jasa
  tidak sama dengan jumlah perusahaan, tenaga kerja, produktivitas,
  kapasitas, mutu jasa, jaringan produksi, atau aliran transaksi.
  Pertambangan juga digabung ke dalam industri, sementara rincian
  subsektor tidak tersedia.
\item
  \textbf{Pengukuran jarak:} jarak Euclidean antara pusat kota tidak
  menangkap waktu tempuh, konektivitas aktual, biaya logistik, hambatan
  administratif, atau jaringan transportasi. Artikel sendiri menyebut
  jarak garis lurus sebagai pendekatan yang dapat diterima pada skala
  makro, bukan ukuran biaya transportasi yang paling akurat. (Hlm. 5,
  Bagian 4.1.)
\item
  \textbf{Cakupan kota inti:} enam kota dipilih karena ukuran dan
  dominasi sistemnya. Pemilihan ini dapat membuat hasil sensitif
  terhadap ambang lebih dari 10 juta jiwa dan tidak menguji kota inti
  lain yang mungkin relevan, khususnya dalam sistem polisentris.
\item
  \textbf{Penugasan wilayah pengaruh:} file tidak memaparkan algoritme,
  nilai ambang formal, atau prosedur reproduktif untuk menentukan kota
  inti yang paling berpengaruh terhadap setiap kota noninti. Jumlah kota
  sekitar juga berubah antarperiode.
\item
  \textbf{Batas spasial:} batas 100 km, 300 km, dan 500 km dipilih untuk
  analisis, tetapi TXT tidak menyajikan uji sensitivitas terhadap batas
  lain atau menjelaskan secara formal apakah semua batas bersifat
  kumulatif atau merupakan sabuk yang saling terpisah.
\item
  \textbf{Cakupan waktu:} blok sepuluh tahun 2010-2015 sebenarnya hanya
  mencakup lima tahun. Perbandingan lima dan sepuluh tahun karena itu
  tidak sepenuhnya simetris.
\item
  \textbf{Kehilangan informasi model:} TXT tidak menjelaskan seluruh
  kovariat potensial seperti kebijakan, modal asing, amenitas, modal
  manusia terukur, kualitas infrastruktur, dan struktur pemerintahan.
  Sebagian variabel yang disebut secara teoritis tidak masuk ke
  estimasi.
\item
  \textbf{Spesifikasi estimasi:} hasil manufaktur sepuluh tahun berbeda
  antara OLS, FE, dan RE; interaksi manufaktur hanya signifikan pada RE.
  File tidak melaporkan alasan pemilihan model, uji Hausman, atau
  pemeriksaan sensitivitas terhadap spesifikasi alternatif.
\item
  \textbf{Dependensi spasial:} unit observasi berada dalam sistem
  geografis yang saling berdekatan, tetapi TXT tidak melaporkan uji
  autokorelasi spasial, model spasial, atau dampak tetangga terhadap
  residual.
\item
  \textbf{Replikasi dataset:} sumber rinci pangsa nilai tambah,
  koordinat, aturan pembersihan, dan prosedur pembentukan sampel tidak
  tersedia dalam TXT. Tabel hasil memberikan angka koefisien, tetapi
  tidak seluruh keputusan pengolahan yang menghasilkan angka tersebut.
\item
  \textbf{Cakupan waktu dataset:} Bagian pengumpulan data menyebut
  1980-2015, sedangkan Data Availability menyebut 1950-2015.
  Ketidakselarasan ini membatasi kepastian tentang seri yang benar-benar
  masuk ke model.
\item
  \textbf{Anomali ekstraksi hasil:} Tabel 7 mempunyai notasi
  \texttt{ICm} pada tabel yang berjudul jasa dan satu koefisien
  interaksi terbaca \texttt{-01023}. Catatan kaki Tabel 4 juga terbaca
  memberi ambang \texttt{***\ p\ \textless{}\ 0.1}, sedangkan catatan
  kaki tabel lain memakai ambang berbeda. Anomali tersebut tidak dapat
  diperbaiki tanpa sumber selain TXT.
\item
  \textbf{Validitas eksternal:} hasil berasal dari sistem perkotaan
  Tiongkok pada masa urbanisasi dan transformasi industri 1980-2015.
  File tidak memberi dasar untuk menggeneralisasikan hasil ke negara
  lain atau periode setelah 2015.
\end{itemize}

\textbf{Tidak disebutkan dalam file.} Artikel tidak memuat bagian
keterbatasan formal yang merinci ukuran efek, ketidakpastian prediksi,
diagnostik residual, atau analisis sensitivitas lengkap. Keterangan di
atas membedakan keputusan desain yang disebut penulis dari keterbatasan
yang diturunkan melalui pemeriksaan isi TXT.

\subsection{Challenges}\label{challenges-28}

\textbf{Tantangan yang dilaporkan atau terlihat dalam desain sumber.}
Artikel menghadapi perubahan cepat ukuran populasi kota inti sehingga
populasi tahun awal dianggap tidak memadai untuk menghitung \emph{market
potential}. Penulis juga harus menentukan batas sistem perkotaan di
negara yang memiliki banyak kota inti dan kedekatan antarkota yang
berbeda. Perubahan peran Shenzhen dari kota yang baru menjadi inti ke
kota yang memengaruhi lebih banyak kota menambah persoalan penetapan
pengaruh menurut waktu. (Hlm. 5-7, Bagian 4.1 dan 5.1; TXT baris
320-340, 390-400.)

\textbf{Tantangan yang dilaporkan atau terlihat dalam desain sumber.}
Perbedaan mekanisme sektor merupakan tantangan substantif: biaya
transportasi, biaya waktu, biaya komunikasi, jaringan produksi, dan
fungsi komando tidak bergerak dengan cara yang sama. Artikel mengatasi
sebagian persoalan itu dengan memisahkan model manufaktur dan jasa serta
melakukan perbandingan lintas skala dan waktu. (Hlm. 3, 6-10, Bagian 2
dan 6; TXT baris 151-182, 342-366.)

\textbf{Tantangan yang dilaporkan atau terlihat dalam desain sumber.}
Ketersediaan variabel juga terbatas. Artikel menilai bahwa ukuran
penduduk dan jarak merupakan dua unsur paling penting untuk \emph{market
potential}, tidak memasukkan upah karena perubahan tingkat pendapatan,
dan tidak memasukkan pertanian karena pangsanya kecil di kota inti.
Modal manusia disebut penting tetapi diperkirakan sebagian tertangkap
oleh populasi. (Hlm. 5-6, Bagian 4.1-4.2; TXT baris 262-269, 346-362.)

\textbf{Inferensi pengolahan.} Tantangan analitis utama adalah
membedakan empat penjelasan yang dapat menghasilkan tanda hasil serupa:
kapasitas internal kota, posisi administratif atau hierarkis, limpahan
manufaktur, dan penarikan fungsi oleh jasa kota inti. Model yang
tersedia menguji hubungan statistik, tetapi tidak mengobservasi seluruh
proses yang membedakan penjelasan tersebut.

\textbf{Inferensi pengolahan.} Tantangan replikasi juga muncul dari tata
letak TXT dua kolom dan beberapa kerusakan karakter hasil ekstraksi.
Rumus, notasi tabel, dan catatan kaki tidak selalu terbaca linear. Hal
ini memerlukan kehati-hatian agar angka atau tanda yang ambigu tidak
dinormalisasi menjadi fakta baru.

\subsection{Practical Implications}\label{practical-implications-28}

\textbf{Implikasi yang dinyatakan sumber.} Temuan artikel dipandang
relevan untuk memahami perkembangan regional yang tidak merata dan
pembangunan kawasan metropolitan di Tiongkok. Komposisi industri kota
inti perlu dipertimbangkan bersama ukuran kota dan \emph{market
potential}, karena kota inti dengan jasa tingkat tinggi yang kuat dapat
mempunyai efek pemusatan lebih besar, sedangkan manufaktur dapat
memberikan limpahan melalui jaringan produksi. (Hlm. 1, abstrak; hlm.
10-11, Bagian 8.)

\textbf{Interpretasi penulis.} Perbandingan Beijing dan Shanghai dipakai
untuk menyiratkan bahwa kebijakan pembangunan yang terlalu berfokus pada
jasa di kota inti dapat mengurangi limpahan geografis ke kota
sekitarnya. Sebaliknya, hubungan produksi manufaktur yang lebih tersebar
dapat membantu menghubungkan kota inti dengan kota noninti. (Hlm. 10,
Bagian 7; TXT baris 605-621.)

\textbf{Inferensi pengolahan.} Bagi perencana, kerangka artikel dapat
digunakan sebagai alat diagnosis awal untuk menanyakan apakah
pertumbuhan kota sekitar berbeda menurut struktur ekonomi kota inti.
Namun, hasilnya tidak cukup untuk menetapkan proyek atau mengalokasikan
investasi tanpa verifikasi kapasitas fasilitas, arus komuter dan barang,
kualitas layanan, kebijakan lokal, biaya perjalanan, serta karakteristik
penduduk.

\textbf{Inferensi pengolahan.} Implikasi kebijakan tidak boleh
disederhanakan menjadi ``manufaktur selalu baik'' atau ``jasa selalu
buruk''. Artikel sendiri menunjukkan hubungan nonlinier dan puncak yang
berbeda menurut jarak. Sektor, skala, fase urbanisasi, dan posisi kota
dalam hierarki perlu diperiksa sebelum hasil diterjemahkan menjadi
strategi pembangunan.

\subsection{Applications}\label{applications-28}

\textbf{Aplikasi yang didukung langsung oleh sumber.}

\begin{itemize}
\tightlist
\item
  membandingkan sistem perkotaan regional yang berpusat pada kota inti
  dengan ukuran dan komposisi industri yang berbeda;
\item
  mengestimasi \emph{market potential} berbasis populasi rata-rata dan
  jarak untuk menilai hubungan kota inti-kota sekitar;
\item
  menguji apakah pangsa manufaktur atau jasa memperkuat atau melemahkan
  hubungan \emph{market potential} dengan pertumbuhan penduduk;
\item
  membandingkan hasil pada blok lima tahun dan sepuluh tahun; dan
\item
  menguji sensitivitas temuan terhadap batas 100 km, 300 km, dan 500 km
  untuk kajian ketimpangan sistem perkotaan serta pembangunan kawasan
  metropolitan.
\end{itemize}

\textbf{Inferensi pengolahan.} Alur tersebut dapat dipakai sebagai
kerangka eksploratif untuk menyusun hipotesis mengenai limpahan
sektoral, bukan sebagai alat penetapan sebab atau kebijakan secara
otomatis. Penerapan di wilayah lain memerlukan data industri dan
penduduk yang sebanding, definisi kota inti yang transparan, ukuran
jarak yang sesuai, pemeriksaan dependensi spasial, dan validasi
mekanisme melalui data aliran.

\textbf{Tidak disebutkan dalam file.} File tidak melaporkan implementasi
kebijakan yang menggunakan model ini, penggunaan oleh pemerintah,
perangkat lunak yang dibagikan, atau evaluasi aplikasi setelah
penelitian.

\subsection{Future Research
(Optional)}\label{future-research-optional-28}

\textbf{Laporan sumber.} Artikel tidak menyediakan bagian atau daftar
agenda penelitian lanjutan yang diberi label \emph{future research}.
Artikel hanya menyatakan bahwa hasil jasa menimbulkan pertanyaan apakah
pertumbuhan tersebut akan bertahan ketika ekspansi jasa berlanjut, serta
mendiskusikan kebutuhan memahami perbedaan sektor dan pergerakan
penduduk perkotaan. Agenda penelitian formal yang lebih rinci:
\textbf{Tidak disebutkan dalam file}. (Hlm. 10, Bagian 7; TXT baris
605-621.)

\textbf{Inferensi pengolahan.} Berdasarkan batas bukti dalam TXT,
penelitian lanjutan yang relevan dapat mencakup:

\begin{itemize}
\tightlist
\item
  membangun data panel yang lebih rinci untuk menelusuri perubahan
  komposisi industri, lokasi perusahaan, dan pertumbuhan penduduk dari
  waktu ke waktu;
\item
  mengukur arus barang, jasa, komuter, perjalanan bisnis, investasi, dan
  hubungan pemasok-pelanggan agar mekanisme limpahan serta bayangan
  tidak hanya disimpulkan dari interaksi;
\item
  memasukkan beberapa kota inti sekaligus untuk merepresentasikan sistem
  polisentris, bukan hanya kota inti yang dipilih sebagai pemberi
  \emph{market potential} terbesar;
\item
  menguji ambang jarak, definisi kota inti, penugasan kota noninti, dan
  ukuran jarak alternatif melalui analisis sensitivitas;
\item
  memisahkan subsektor jasa tingkat tinggi, jasa konsumen, keuangan,
  ritel, dan subsektor manufaktur untuk menguji apakah tanda agregat
  menyembunyikan mekanisme yang berbeda;
\item
  menguji dependensi spasial, variabel yang terlewat, seleksi lokasi,
  dan kemungkinan kausalitas terbalik dengan desain identifikasi yang
  lebih kuat; dan
\item
  memvalidasi pangsa nilai tambah dan pertumbuhan penduduk terhadap
  ukuran produktivitas, kapasitas, kualitas, serta pemanfaatan aktual.
\end{itemize}

Semua butir di atas adalah inferensi pengolahan, bukan agenda yang
secara eksplisit ditulis oleh penulis artikel.

\subsection{Provenance Notes}\label{provenance-notes-28}

\begin{itemize}
\tightlist
\item
  \textbf{Basis akses:} review ini hanya menggunakan file TXT
  \texttt{source/journal/A13-li-sun-2020-industrial-composition-and-agglomeration-shadow-evidence-from-chinas-large-urban-systems-10.1155-2020-5717803.txt}.
  PDF asli, sumber daring, metadata eksternal, dan artikel yang
  tercantum dalam daftar pustaka tidak dibuka untuk menambah substansi.
\item
  \textbf{Cakupan baca:} seluruh 766 baris TXT dibaca, termasuk blok
  metadata PDF, teks artikel, abstrak, persamaan, Gambar 1-2, Tabel 1-8,
  bagian Data Availability, Conflicts of Interest, Acknowledgments, dan
  seluruh daftar pustaka.
\item
  \textbf{Metadata ekstraksi:} blok \texttt{{[}PDF\ Metadata{]}}
  menyatakan TXT diekstrak pada 31 Agustus 2026 dengan
  \texttt{pdftotext\ -layout} dari PDF 11 halaman. Metadata itu dicatat
  sebagai isi TXT, bukan diverifikasi melalui PDF atau sumber penerbit.
\item
  \textbf{Locator:} nomor halaman dalam review mengikuti nomor halaman
  artikel yang tercetak pada TXT, sedangkan locator \texttt{TXT\ baris}
  mengikuti nomor baris file. Tabel, gambar, bagian, dan persamaan
  disebut ketika tersedia.
\item
  \textbf{Keterbatasan ekstraksi:} tata letak dua kolom membuat beberapa
  kalimat terpotong dan karakter seperti \texttt{This} atau \texttt{The}
  terbaca sebagai \texttt{.is}. Notasi Persamaan 2 dan beberapa angka
  tabel juga tidak selalu tersusun linear. Tanda atau angka yang ambigu
  dipertahankan dan diberi keterangan, bukan diperbaiki berdasarkan
  sumber luar.
\item
  \textbf{Anomali internal yang dipertahankan:} Bagian 3.3 menyebut
  periode data 1980-2015, sedangkan Data Availability menyebut
  1950-2015. Tabel 7 berjudul jasa tetapi memuat label \texttt{ICm}, dan
  satu angka interaksi terbaca \texttt{-01023}. Catatan kaki Tabel 4
  juga memiliki ambang signifikansi yang tampak tidak konsisten dengan
  catatan kaki tabel lain.
\item
  \textbf{Informasi yang tidak tersedia:} sumber rinci pangsa nilai
  tambah, protokol pembersihan data, koordinat dan algoritme pembagian
  pengaruh, satuan tepat variabel regresi, alasan pemilihan FE/RE, uji
  kausalitas, diagnostik residual, uji dependensi spasial, dan
  implementasi kebijakan tidak disebutkan dalam file.
\item
  \textbf{Pembedaan epistemik:} bagian berlabel \textbf{Laporan sumber}
  memparafrasakan apa yang dinyatakan atau ditampilkan artikel;
  \textbf{Interpretasi penulis} merekam makna yang diberikan penulis
  terhadap hasil; \textbf{Inferensi pengolahan} adalah evaluasi
  berdasarkan isi dan keterbatasan internal TXT, bukan klaim artikel
  atau informasi eksternal.
\end{itemize}

\section{Review A14}\label{review-a14}

\subsection{Identitas}\label{identitas-1}

\begin{itemize}
\tightlist
\item
  \textbf{Kode kajian:} A14. Kode ini mengikuti permintaan review dan
  prefiks nama file sumber.
\item
  \textbf{Judul:} \emph{Does broadband internet allow cities to `borrow
  size'? Evidence from the Swedish labour market} (TXT baris 40--47 dan
  91--92, judul artikel, PDF pp.~1 dan 2).
\item
  \textbf{Penulis:} Duco de Vos, Urban Lindgren, Maarten van Ham, dan
  Evert Meijers (TXT baris 43 dan 93, baris penulis, PDF pp.~1--2).
\item
  \textbf{Jurnal:} \emph{Regional Studies} (TXT baris 30, 35, dan
  85--92, kepala jurnal dan judul artikel, PDF pp.~1--2).
\item
  \textbf{Tahun dan publikasi:} Teks sitasi menyebut tahun 2020; artikel
  tercatat terbit daring pada 9 Januari 2020. Metadata PDF juga memuat
  subjek ``Regional Studies, 2019'', sehingga tahun yang muncul di
  metadata tidak sepenuhnya seragam (TXT baris 8, 45--47, dan 62, blok
  metadata dan informasi publikasi).
\item
  \textbf{DOI:} \texttt{10.1080/00343404.2019.1699238} (TXT baris 45--49
  dan 86, informasi sitasi dan DOI).
\item
  \textbf{Tautan artikel:}
  \texttt{https://doi.org/10.1080/00343404.2019.1699238} (TXT baris 49,
  informasi tautan artikel).
\item
  \textbf{Kata kunci:} broadband internet, agglomeration economies,
  borrowed size, commuting, employment (TXT baris 9 dan 102--103,
  metadata dan ``KEYWORD'').
\item
  \textbf{JEL:} L96, O3, O18, R11 (TXT baris 106, ``JEL'').
\item
  \textbf{Riwayat naskah:} diterima 8 Maret 2019 dan dalam bentuk revisi
  11 November 2019 (TXT baris 107, ``HISTORY'').
\item
  \textbf{Volume, nomor, dan halaman jurnal:} Tidak disebutkan dalam
  file. File hanya memuat 13 halaman PDF hasil ekstraksi, bukan
  informasi volume/nomor/paginasi jurnal yang lengkap (jumlah halaman
  PDF: TXT baris 22; ketiadaan volume/nomor/paginasi lengkap terlihat
  dari informasi sitasi pada baris 45--49).
\item
  \textbf{Status peer-review:} Tidak dapat diverifikasi dari file TXT;
  status tersebut tidak diasumsikan.
\end{itemize}

\subsection{Summarized Abstract}\label{summarized-abstract-29}

{[}Laporan sumber{]} Artikel mengeksplorasi apakah broadband internet
membantu kota-kota kecil memperoleh manfaat pasar kerja kota yang lebih
besar, yaitu proses yang dibahas sebagai \emph{borrowed size}. Dengan
micro-data Swedia 2007--2015 dan data broadband yang disebut unik,
penulis melaporkan bukti yang bersifat sugestif bahwa broadband membantu
kota-kota kecil memperoleh manfaat tersebut. Pola itu terutama didorong
oleh penetrasi broadband keseluruhan di tempat tinggal, bukan sekadar
ketersediaan broadband di lokasi hunian individual (TXT baris 95--101,
bagian ``ABSTRACT'', PDF p.~1).

{[}Interpretasi penulis{]} Kata ``suggestive'' membatasi kekuatan klaim
abstrak. Abstrak tidak menyajikan hasil sebagai pembuktian universal
atau sebagai estimasi kausal yang tanpa syarat (TXT baris 98--101,
bagian ``ABSTRACT'', PDF p.~1).

\subsection{Problem Statement}\label{problem-statement-29}

{[}Laporan sumber{]} \emph{Borrowed size} merujuk pada keadaan ketika
kota kecil atau wilayah metropolitan kecil dekat konsentrasi penduduk
yang lebih besar dapat menunjukkan sebagian karakteristik kota besar.
Literatur yang dirujuk artikel menekankan kedekatan, tetapi kurang
memperhatikan mekanisme integrasi dan transfer informasi yang mungkin
membuat manfaat aglomerasi menjangkau tempat yang lebih jauh. Artikel
menyatakan bahwa hubungan broadband internet dan \emph{borrowed size}
banyak diasumsikan, tetapi belum ada kerja empiris yang menyelidiki
hubungan tersebut (TXT baris 110--130, bagian ``INTRODUCTION'', PDF
pp.~1--2).

{[}Laporan sumber{]} Masalah substantif dipersempit pada pasar kerja.
Pasar kerja kota kecil dalam wilayah metropolitan besar disebut memiliki
sedikit perhatian dibandingkan layanan bisnis serta akses ke fungsi
perkotaan seperti belanja dan rekreasi. Karena itu, artikel menanyakan
apakah broadband mengubah aksesibilitas pasar kerja dan tingkat
pekerjaan secara berbeda menurut posisi tempat dalam hierarki perkotaan
(TXT baris 156--197, bagian ``LITERATURE REVIEW'', PDF p.~2).

{[}Inferensi pengolahan{]} Unit masalahnya bukan hanya apakah rumah
tangga memiliki koneksi. Artikel membedakan ketersediaan broadband di
lokasi hunian sebagai kemungkinan saluran keputusan individual dari
penetrasi broadband pada tingkat tempat sebagai kemungkinan saluran
lingkungan ekonomi lokal (TXT baris 160--175 dan 193--198, bagian
``LITERATURE REVIEW'').

\subsection{Objectives}\label{objectives-29}

{[}Laporan sumber{]} Tujuan utama penulis adalah mengisi kekosongan
empiris tentang hubungan internet dan \emph{borrowed size}. Tujuan
operasionalnya adalah:

\begin{itemize}
\tightlist
\item
  Menilai apakah broadband memperbaiki aksesibilitas tenaga kerja di
  luar kota-kota besar Swedia (TXT baris 182--192, bagian
  ``INTRODUCTION'', PDF p.~2).
\item
  Menilai apakah pengenalan broadband secara lokal meningkatkan tingkat
  pekerjaan di luar kota-kota terbesar (TXT baris 188--192, bagian
  ``INTRODUCTION'', PDF p.~2).
\item
  Mengestimasi hubungan perubahan broadband dengan perubahan jarak
  komuter dan tingkat pekerjaan melalui model fixed effects, dengan
  fokus pada efek dua tahun setelah pengenalan broadband (TXT baris
  193--206, bagian ``INTRODUCTION'', PDF p.~2).
\item
  Membedakan efek ketersediaan broadband di tempat tinggal dari efek
  penetrasi broadband pada tingkat tempat, serta menguji heterogenitas
  antara inti perkotaan, wilayah perkotaan, wilayah rural, dan jarak
  dari inti perkotaan terdekat (TXT baris 193--206, bagian
  ``INTRODUCTION'', PDF p.~2; baris 338--375, bagian ``METHODOLOGY'').
\end{itemize}

\subsection{Literature Survey}\label{literature-survey-29}

{[}Laporan sumber{]} Artikel membangun kerangka \emph{borrowed size}
dari Alonso, yang menggambarkannya sebagai keadaan ketika kota kecil
atau wilayah metropolitan kecil memiliki sebagian karakteristik kota
besar karena dekat dengan konsentrasi penduduk lain. Artikel juga
merangkum lima ciri dari uraian tersebut: konsep analitis; berkenaan
dengan kota kecil; berada dekat wilayah metropolitan yang lebih besar;
mempertahankan keunggulan kota kecil sambil menikmati manfaat
aglomerasi; dan memperoleh manfaat melalui aksesibilitas serta
konektivitas jaringan (TXT baris 112--128 dan 156--175, bagian
``INTRODUCTION'' dan ``Borrowed size in labour markets'', PDF pp.~1--2).

{[}Laporan sumber{]} Broadband ditempatkan sebagai infrastruktur yang
dapat menyediakan aksesibilitas sekaligus konektivitas jaringan. Penulis
membedakan efek aksesibilitas bagi individu, yang diproksikan dengan
ketersediaan broadband di lokasi hunian, dari efek konektivitas jaringan
terhadap ekonomi lokal, yang diproksikan dengan pangsa koneksi broadband
di tempat tinggal atau tempat tersebut (TXT baris 160--175, bagian
``Borrowed size in labour markets'', PDF p.~2). Artikel merujuk argumen
bahwa perbaikan akses ke layanan, termasuk broadband, dapat terkait
dengan pertumbuhan pusat kecil dan wilayah rural, serta bahwa jaringan
informasi dapat mengurangi ketergantungan pada kedekatan fisik (TXT
baris 116--130, bagian ``INTRODUCTION'').

{[}Laporan sumber{]} Dalam literatur pasar kerja, artikel mengaitkan
broadband dengan teleworking, pencarian pekerjaan, pencocokan
pekerja--perusahaan, dan lokasi perusahaan. Broadband dapat memungkinkan
kerja dari rumah dan komuter yang lebih panjang pada hari lain; dapat
menurunkan biaya pencarian kerja; serta dapat memberi perusahaan di kota
kecil akses pada tenaga kerja dan eksternalitas informasi kota besar di
dekatnya. Artikel juga menjelaskan kemungkinan berlawanan pada tingkat
tempat: broadband dapat menyebarkan manfaat aglomerasi dan mengurangi
komuter berlebih jika pekerjaan lokal bertambah (TXT baris 160--180 dan
193--228, bagian ``Borrowed size in labour markets'').

{[}Laporan sumber{]} Literatur tentang kinerja lokal dan teknologi
informasi dalam artikel mencakup hasil yang tidak seragam. Artikel
menyebut studi tentang pencarian kerja daring yang menemukan durasi
menganggur tidak lebih pendek atau mungkin lebih panjang, studi lanjutan
yang menemukan pencari kerja daring memperoleh pekerjaan sekitar 25\%
lebih cepat, studi Craigslist yang tidak menemukan pengaruh pada tingkat
pengangguran, dan studi Jerman berbasis instrumental variable yang tidak
menemukan efek signifikan broadband terhadap pengangguran. Artikel juga
melaporkan bahwa riset lokasi perusahaan mengaitkan broadband dengan
pendirian perusahaan baru, dengan kekuatan hubungan yang berbeda menurut
ruang dan industri (TXT baris 233--272, bagian ``Local performance and
IT'', PDF p.~3).

{[}Laporan sumber{]} Studi Atasoy disebut paling mirip dengan artikel
ini: studi tersebut mengestimasi ekspansi broadband dan hasil pasar
kerja di Amerika Serikat pada 1999--2007, mengukur penetrasi sebagai
pangsa penduduk county yang tinggal di area dengan broadband, dan
melaporkan tingkat pekerjaan 1,8 poin persentase lebih tinggi dengan
efek lebih besar di wilayah yang lebih terpencil. Artikel ini menyatakan
menambah literatur dengan heterogenitas urban--suburban--rural, jarak ke
inti kota, skala geografis yang lebih rinci, serta pemisahan penetrasi
tingkat tempat dan ketersediaan di lokasi hunian (TXT baris 229--248,
bagian ``Local performance and IT'', PDF p.~3).

{[}Interpretasi penulis{]} Hipotesis teoretisnya adalah: ketersediaan
broadband di hunian meningkatkan jarak komuter dan pekerjaan melalui
kerja dari rumah serta akses pada pekerjaan yang jauh; penetrasi
broadband pada tingkat tempat menarik perusahaan dan meningkatkan
pekerjaan lokal sehingga dapat menurunkan komuter; dan pola
\emph{borrowed size} terlihat jika efek terbesar muncul pada jarak
tertentu dari kota besar, terutama wilayah yang sebelumnya berada di
tepi pasar kerja lalu menjadi lebih terintegrasi (TXT baris 250--272,
bagian ``Hypotheses'', PDF p.~3).

\subsection{Method Used}\label{method-used-29}

{[}Laporan sumber{]} Data diagregasi ke grid berukuran 250 × 250 meter.
Penulis menggunakan model fixed effects untuk mengestimasi hubungan
variasi waktu broadband dengan variasi waktu jarak komuter dan tingkat
pekerjaan. Spesifikasi mencakup ketersediaan broadband di tempat
tinggal, penetrasi broadband pada tingkat tempat, kontrol pada tingkat
grid/tempat/wilayah, fixed effects grid, fixed effects tahun, dan galat
(TXT baris 338--358, bagian ``METHODOLOGY'', PDF p.~5).

{[}Laporan sumber{]} Model spasial menginteraksikan broadband dengan
kategori inti FUA, bagian FUA di luar inti, dan wilayah di luar FUA.
Model lain memungkinkan efek broadband berubah menurut jarak dari FUA
terdekat. FUA didasarkan pada definisi OECD; penulis mengakui definisi
tersebut berperspektif monosentris, beberapa zona saling bersentuhan,
dan mungkin tumpang tindih dalam kenyataan, tetapi memilihnya karena
tersedia dan konsisten sehingga hasil dapat direproduksi (TXT baris
359--375, bagian ``METHODOLOGY'', PDF p.~5).

{[}Laporan sumber{]} Penulis memasukkan kontrol populasi dan pendapatan
pada tingkat grid; populasi, pendapatan, struktur pendidikan dan
demografi, jumlah perusahaan, ukuran perusahaan, pendapatan agregat,
pangsa penduduk urban, serta 21 pangsa industri pada tingkat wilayah
pasar kerja. Mereka melaporkan masalah multikolinearitas yang
memengaruhi koefisien kontrol dan karena itu berfokus pada variabel
independen utama yang memiliki VIF konsisten rendah (TXT baris 372--398
dan 497--501, bagian ``Summary statistics'', ``METHODOLOGY'', dan
``RESULTS'').

{[}Laporan sumber{]} Untuk mengurangi kekhawatiran bahwa pekerjaan dan
perilaku komuter mendorong adopsi broadband, analisis utama menggunakan
broadband dan penetrasi yang dilag dua tahun. Penulis menyatakan bahwa
timing kedatangan broadband mungkin memiliki unsur acak, tetapi
inisiatif komunitas dapat endogen; jika upaya lokal endogen, estimasi
sebaiknya dibaca sebagai batas atas efek pengenalan broadband (TXT baris
442--450, bagian ``METHODOLOGY'', PDF pp.~6--7).

{[}Laporan sumber{]} Analisis sensitivitas membedakan teknologi tetap
tercepat dan membandingkan DSL dengan kabel/serat optik, sambil tetap
memasukkan mobile 3G dan penetrasi tingkat tempat. Pola umum dinilai
serupa; efek mobile 3G lebih kecil dan paling sedikit bervariasi menurut
ruang (TXT baris 631--671, bagian ``Sensitivity analysis: internet
technologies'', PDF pp.~10--11).

\subsection{Dataset}\label{dataset-29}

{[}Laporan sumber{]} Sumber pertama adalah ASTRID, data register
longitudinal di Umeå University yang mencakup semua individu di Swedia
berusia 18--64 tahun pada 2007--2015. Variabelnya mencakup pendapatan
kerja, pendidikan, demografi, lokasi tempat tinggal, lokasi tempat
kerja, identitas perusahaan dan establishment, lokasi, industri, sektor,
serta ukuran perusahaan (TXT baris 276--305, bagian ``ASTRID data'', PDF
pp.~3--4).

{[}Laporan sumber{]} Lokasi hunian tersedia dengan ketelitian 100 meter
lalu diubah menjadi grid 250 × 250 meter agar dapat dipadankan dengan
data broadband. Lokasi tempat kerja yang tepat digunakan untuk
menghitung jarak komuter Euclidean. Informasi industri memakai SNI 2007
pada dua digit sehingga menghasilkan 21 industri (TXT baris 295--307,
bagian ``ASTRID data'', PDF p.~4).

{[}Laporan sumber{]} Sumber broadband berasal dari survei tahunan
pemasok internet yang dilakukan Swedish Post and Telecom Authority sejak
2007. Pemasok diminta mencatat seluruh alamat yang memiliki koneksi
menurut jenis koneksi; PTS mengubahnya menjadi data bereferensi
geografis pada grid 250 × 250 meter. Broadband dianggap tersedia jika
sedikitnya satu alamat dalam grid memiliki koneksi, dan jumlah alamat
digunakan untuk membangun ukuran penetrasi grid (TXT baris 308--326,
bagian ``Broadband data'', PDF p.~4).

{[}Laporan sumber{]} Ukuran utama broadband hunian adalah dummy koneksi
tetap apa pun, yaitu DSL, kabel, atau serat optik. Analisis sensitivitas
memakai koneksi tetap tercepat, sedangkan mobile 3G juga dimodelkan.
Teks membandingkan kecepatan DSL sekitar 14 Mbit/s dengan kabel/serat
optik sekitar 100 Mbit/s serta menyebut kecepatan 3G sekitar 6,5 Mbit/s
(TXT baris 356--369, bagian ``Broadband data'', PDF p.~5).

{[}Laporan sumber{]} Statistik ringkas menunjukkan tingkat pekerjaan
rata-rata grid 0,639, jarak komuter rata-rata 23,60 km, broadband hunian
0,975, dan broadband tempat 0,986. Nilai broadband yang tinggi tetap
memiliki variasi, terutama pada tingkat grid. Jumlah observasi statistik
ringkas adalah 871.688; observasi jarak komuter hanya 837.368 karena
grid tanpa pekerjaan dikeluarkan dari ukuran tersebut (TXT baris
372--438, Tabel 1, PDF p.~6).

{[}Inferensi pengolahan{]} File TXT tidak memuat data mentah, kode
analisis, atau isi lengkap data suplementer. Karena itu, review ini
hanya menggunakan definisi, statistik, koefisien, dan keterangan
analitis yang tercetak atau terekstrak dalam TXT.

\subsection{Population / Sample}\label{population-sample-29}

{[}Laporan sumber{]} Populasi yang tercakup dalam ASTRID adalah seluruh
individu di Swedia berusia 18--64 tahun selama 2007--2015. Namun, unit
observasi analisis bukan individu secara langsung: informasi tersebut
diagregasi terutama ke grid 250 × 250 meter, kemudian indikator
sosial-ekonomi dan struktur industri juga dibentuk pada 60 wilayah pasar
kerja lokal yang mencakup seluruh Swedia (TXT baris 276--305, bagian
``ASTRID data'', PDF pp.~3--4).

{[}Laporan sumber{]} Setiap wilayah pasar kerja lokal memiliki rata-rata
sekitar 100.000 penduduk usia kerja. Untuk klasifikasi spasial, grid
dibedakan menurut inti FUA, bagian wilayah urban di luar inti, dan
wilayah di luar FUA. Tempat didefinisikan menggunakan \emph{tätorter}
dengan sedikitnya 200 penduduk dan jarak antargedung kurang dari 200
meter, serta \emph{småorter} dengan 50--200 penduduk dan jarak
antargedung tidak lebih dari 150 meter (TXT baris 295--305 dan 736--742,
bagian ``ASTRID data'' dan ``NOTES'', PDF pp.~4 dan 12).

{[}Laporan sumber{]} Tabel 1 mencatat N 871.688 untuk statistik ringkas.
Karena grid tanpa pekerjaan tidak memiliki observasi jarak komuter,
ukuran jarak komuter memiliki 837.368 observasi. Pada regresi, jumlah
observasi berbeda menurut outcome dan spesifikasi: 825.353 pada Tabel 2
kolom (1), 634.033 pada kolom (2)--(3), 844.449 pada kolom (4), dan
646.345 pada kolom (5)--(6). Teks tidak menjelaskan secara lengkap
prosedur penyusutan sampel dari cakupan ASTRID ke setiap N regresi (TXT
baris 405--438, Tabel 1, PDF p.~6; baris 464--485, Tabel 2, PDF
pp.~7--8).

\subsection{Dependent Variables}\label{dependent-variables-29}

{[}Laporan sumber{]} Ada dua variabel dependen utama:

\begin{itemize}
\tightlist
\item
  \textbf{Jarak komuter:} jarak Euclidean antara lokasi tempat tinggal
  dan tempat kerja untuk setiap pekerja, lalu dirata-ratakan pada
  seluruh individu bekerja dalam grid. Penulis mencatat bahwa ukuran
  Euclidean dapat mengandung error pengukuran; arah error tidak jelas,
  walaupun artikel mengharapkan error terbatas berdasarkan hubungan yang
  disebutnya dengan biaya transportasi umum (TXT baris 308--321, bagian
  ``ASTRID data'', PDF p.~4).
\item
  \textbf{Tingkat pekerjaan:} proporsi individu pada grid yang memiliki
  pendapatan tahunan dari bekerja lebih tinggi dari 160.000 SEK,
  kira-kira €15.000. Penulis memperingatkan bahwa mahasiswa atau pekerja
  paruh waktu dapat dikategorikan sebagai tidak bekerja, tetapi
  memperkirakan error terkait perubahan antar-tahun terbatas karena
  ambang tersebut lebih rendah daripada pendapatan terendah yang
  disepakati secara kolektif (TXT baris 323--331, bagian ``ASTRID
  data'', PDF p.~4).
\end{itemize}

{[}Laporan sumber{]} Variabel utama yang menjelaskan kedua outcome
adalah dummy broadband tetap di lokasi hunian, penetrasi broadband pada
tingkat tempat, dan pengenalan mobile 3G. Spesifikasi utama menggunakan
broadband yang dilag dua tahun untuk menilai perubahan outcome setelah
jeda adaptasi (TXT baris 193--206, 356--369, dan 442--450, bagian
``INTRODUCTION'', ``Broadband data'', dan ``METHODOLOGY'').

\subsection{Results}\label{results-29}

{[}Laporan sumber{]} Tabel 2 menggunakan kontrol, fixed effects tahun,
dan fixed effects grid pada seluruh kolom. Catatan tabel menyatakan
bahwa \texttt{***}, \texttt{**}, dan \texttt{*} masing-masing berarti p
\textless{} 0,01, p \textless{} 0,05, dan p \textless{} 0,1 (TXT baris
464--485, Tabel 2, PDF pp.~7--8). Hasil utamanya adalah:

\begin{itemize}
\tightlist
\item
  Pada spesifikasi dengan broadband hunian tahun berjalan, broadband
  berkorelasi dengan penurunan jarak komuter sebesar 1,721 km (SE 0,302;
  p \textless{} 0,01), sedangkan mobile 3G berkorelasi dengan kenaikan
  0,480 km (SE 0,140; p \textless{} 0,01). Penulis merangkum koefisien
  broadband itu sebagai penurunan sekitar 1,7 km yang berbarengan dengan
  pengenalan broadband (TXT baris 467--474 dan 497--511, Tabel 2 kolom
  (1), PDF pp.~7--8).
\item
  Pada spesifikasi dengan lag dua tahun, broadband hunian berkaitan
  dengan penurunan jarak komuter 1,149 km (SE 0,377; p \textless{}
  0,01), sedangkan mobile 3G tidak signifikan. Penulis menafsirkan hasil
  ini sebagai broadband pada tahun \texttt{t\ −\ 2} yang berbarengan
  dengan penurunan jarak komuter pada tahun \texttt{t}, bukan sebagai
  penghapusan seluruh masalah kausalitas (TXT baris 469--476 dan
  512--521, Tabel 2 kolom (2), PDF pp.~7--8).
\item
  Untuk tingkat pekerjaan, broadband hunian tahun berjalan memiliki
  koefisien 0,00843 (SE 0,00141; p \textless{} 0,01), yang oleh penulis
  dibaca sebagai kenaikan 0,8 poin persentase. Dengan lag dua tahun,
  koefisiennya 0,00344 (SE 0,00177; p \textless{} 0,1), dibaca sebagai
  kenaikan 0,3 poin persentase. Mobile 3G positif dan signifikan pada
  model pekerjaan, tetapi penulis menyebut efeknya kecil (TXT baris
  467--476, 497--521, Tabel 2 kolom (4)--(5), PDF pp.~7--8).
\item
  Setelah penetrasi broadband tingkat tempat dimasukkan, koefisien
  broadband hunian untuk jarak komuter menjadi −0,708 dan tidak
  signifikan, sedangkan penetrasi broadband tempat menjadi −1,750 (SE
  0,826; p \textless{} 0,05). Penulis menghitung bahwa kenaikan 0,1 pada
  pangsa koneksi tempat berhubungan dengan penurunan rata-rata sekitar
  0,18 km. Untuk tingkat pekerjaan, broadband hunian menjadi 0,00117 dan
  tidak signifikan, sedangkan broadband tempat menjadi 0,00920 (SE
  0,00389; p \textless{} 0,05); kenaikan 0,1 pada penetrasi tempat
  dibaca sebagai kenaikan 0,09 poin persentase. Mobile 3G pada model ini
  tidak signifikan untuk jarak komuter dan tetap kecil tetapi signifikan
  untuk pekerjaan (TXT baris 469--476, 535--555, Tabel 2 kolom (3) dan
  (6), PDF pp.~7--9).
\item
  Model Tabel 2 memiliki R² 0,530 pada kolom komuter tahun berjalan,
  0,579 pada dua kolom komuter berlag, 0,732 pada pekerjaan tahun
  berjalan, dan 0,764 pada dua kolom pekerjaan berlag. N-nya
  masing-masing 825.353, 634.033, 634.033, 844.449, 646.345, dan 646.345
  (TXT baris 479--485, Tabel 2, PDF pp.~7--8). {[}Laporan sumber{]}
  Untuk R² 0,530 pada kolom (1), penulis menyatakan bahwa nilainya
  terutama berasal dari fixed effects (TXT baris 507--511, bagian
  ``RESULTS'', PDF p.~8). {[}Inferensi pengolahan{]} Karena itu, review
  ini tidak memperlakukan R² tersebut sebagai bukti tersendiri atas
  mekanisme \emph{borrowed size}.
\end{itemize}

{[}Laporan sumber{]} Tabel 3 dan analisis jarak menunjukkan
heterogenitas spasial:

\begin{itemize}
\tightlist
\item
  Untuk jarak komuter berdasarkan kategori FUA, tidak ada efek broadband
  yang signifikan di inti FUA. Di luar inti tetapi masih dalam FUA,
  broadband hunian bernilai −2,175 (SE 0,615; p \textless{} 0,01),
  sedangkan broadband tempat bernilai 0,316 dan tidak signifikan. Di
  luar FUA, broadband hunian bernilai 1,304 (SE 0,699; p \textless{}
  0,1) dan broadband tempat bernilai −4,325 (SE 1,202; p \textless{}
  0,01) (TXT baris 562--579, Tabel 3 kolom (1), PDF p.~9).
\item
  Untuk tingkat pekerjaan berdasarkan kategori FUA, efek broadband yang
  signifikan hanya berasal dari penetrasi tingkat tempat. Di inti FUA,
  koefisiennya 0,143 (SE 0,0741; p \textless{} 0,1); di luar inti tetapi
  masih dalam FUA, 0,0140 (SE 0,00564; p \textless{} 0,05). Di luar FUA,
  koefisien broadband hunian −0,000732 dan broadband tempat 0,00529,
  keduanya tidak signifikan. Broadband hunian di inti dan di luar inti
  FUA juga tidak signifikan (TXT baris 565--579, Tabel 3 kolom (3), PDF
  p.~9).
\item
  Pada model jarak dari inti FUA, ketersediaan broadband hunian kecil
  dan negatif sampai sekitar 30 km, lalu pada tempat berjarak 50 km atau
  lebih menjadi positif sekitar 4,3 km. Penulis mengaitkan pola wilayah
  jauh ini dengan kerja dari rumah, tetapi menyebutnya sebagai
  kemungkinan. Penetrasi broadband tempat dikaitkan dengan komuter yang
  lebih panjang pada jarak menengah, khususnya 20--30 km; pada tempat
  sampai 20 km atau berjarak 50 km dan lebih, perubahan umum penetrasi
  tidak banyak mengubah situasi menurut uraian penulis (TXT baris
  605--620, Gambar 3(a), PDF pp.~9--10).
\item
  Untuk tingkat pekerjaan menurut jarak, ketersediaan broadband hunian
  signifikan hanya dekat dan di dalam inti urban, sedangkan penetrasi
  broadband tempat signifikan pada jarak menengah 20--50 km dari inti
  FUA. Penulis menyebut wilayah 20--50 km tersebut sebagai ``zona
  ekspansi'' pasar kerja urban (TXT baris 636--650, Gambar 3(b), PDF
  p.~10).
\end{itemize}

{[}Laporan sumber{]} Analisis sensitivitas yang membedakan DSL dari
kabel/serat optik menghasilkan pola yang disebut serupa dengan hasil
utama. Efek positif broadband hunian hanya tampak di luar FUA dan
tampaknya didorong kabel/serat optik; efek pekerjaan didorong penetrasi
tempat dan hanya hadir di area menengah dalam FUA di luar inti. Mobile
3G memiliki efek kecil tetapi signifikan pada pekerjaan, sedangkan pola
jarak teknologi pada Gambar 4 disebut mirip dengan estimasi sebelumnya
(TXT baris 631--671, bagian ``Sensitivity analysis: internet
technologies'' dan Gambar 4, PDF pp.~10--11). Nilai lengkap Tabel A4
tidak tersedia dalam TXT; file hanya memuat ringkasan naratifnya.

\subsection{Findings}\label{findings-29}

{[}Laporan sumber{]} Hasil agregat tidak mendukung ekspektasi awal bahwa
broadband hunian akan memperpanjang komuter. Dengan lag dua tahun,
broadband hunian justru berkaitan dengan komuter yang lebih pendek; efek
pekerjaan positif pada model sederhana tetapi menjadi tidak signifikan
setelah penetrasi tempat dikendalikan. Artikel menyatakan bahwa hasil
negatif pada komuter dan positif pada pekerjaan terutama didorong
penetrasi broadband di tingkat tempat, bukan ketersediaan broadband di
hunian (TXT baris 502--521 dan 539--555, bagian ``RESULTS'', PDF
pp.~8--9).

{[}Interpretasi penulis{]} Penulis menilai ketersediaan broadband di
hunian kemungkinan berkorelasi dengan penetrasi broadband keseluruhan
dan struktur ekonomi atau iklim bisnis lokal. Di luar kota pusat,
kondisi bisnis yang lebih baik dapat menarik lebih banyak perusahaan dan
pekerjaan lokal; akibatnya, lebih banyak penduduk mungkin bekerja dekat
rumah dan tidak perlu menjadi komuter jarak jauh. Ini adalah penjelasan
penulis atas hasil, bukan mekanisme yang diukur langsung dalam model
(TXT baris 522--538, bagian ``Broadband and labour market outcomes'',
PDF p.~8).

{[}Laporan sumber{]} Pola yang dianggap paling mendukung \emph{borrowed
size} berada pada jarak menengah dari inti urban: penetrasi broadband
tempat berkaitan dengan pekerjaan pada jarak 20--50 km dan dengan jarak
komuter pada 20--30 km. Penulis menyebut area tersebut sebagai wilayah
yang ditambahkan atau lebih terintegrasi ke dalam pasar kerja urban
melalui broadband. Di lokasi yang lebih dekat, pasar kerja dipandang
sudah lebih terintegrasi; di lokasi 50 km atau lebih, jaraknya dinilai
terlalu jauh untuk integrasi yang sama (TXT baris 598--620, Gambar 3,
PDF pp.~9--10).

{[}Interpretasi penulis{]} Penulis menafsirkan penetrasi broadband
tempat sebagai faktor yang dapat membuat lokasi menengah lebih menarik
bagi perusahaan dan penduduk, sehingga mereka dapat mengakses manfaat
pasar kerja urban di dekatnya. Mereka menyebut ini sebagai tanda jelas
\emph{borrowed size} dan sebagai indikasi bahwa broadband mendorong
ekspansi spasial pasar kerja urban (TXT baris 598--604 dan 698--709,
bagian ``Broadband and borrowed size'' dan ``CONCLUSIONS'', PDF
pp.~9--11).

{[}Interpretasi penulis{]} Untuk broadband hunian, penulis membedakan
pola kerja dari rumah di wilayah remote dari pola \emph{borrowed size}.
Kenaikan komuter di wilayah jauh dinilai konsisten dengan kerja dari
rumah, sedangkan tidak adanya hasil signifikan dekat inti membuat
penulis menyatakan belum ada hubungan kausal antara koneksi broadband
hunian dan \emph{borrowed size} di pasar kerja (TXT baris 610--620 dan
687--697, bagian ``Broadband and borrowed size'' dan ``CONCLUSIONS'',
PDF pp.~9--11).

{[}Inferensi pengolahan{]} Berdasarkan pembandingan hipotesis dan
spesifikasi, bukti artikel lebih kuat untuk mekanisme tingkat tempat
pada zona menengah daripada untuk klaim bahwa setiap koneksi rumah
tangga secara langsung menghasilkan \emph{borrowed size}. Namun,
inferensi ini tetap dibatasi oleh peringatan artikel sendiri tentang
confounder, endogenitas, dan ketiadaan pengukuran kerja dari rumah (TXT
baris 442--450, 663--675, dan 713--721, bagian ``METHODOLOGY'',
``CONCLUSIONS'', dan ``LIMITATIONS'' yang dibahas penulis).

\subsection{Conclusions}\label{conclusions-29}

{[}Laporan sumber{]} Penulis menyimpulkan bahwa artikel mengestimasi
hubungan penetrasi broadband dengan perilaku komuter dan pekerjaan,
dengan memisahkan inti urban, wilayah urban, dan rural serta
memungkinkan efek berubah menurut jarak dari inti FUA. Inovasi yang
mereka nyatakan adalah memodelkan mekanisme potensial \emph{borrowed
size} melalui variasi waktu broadband dan membedakan efek di hunian dari
efek pada tempat (TXT baris 673--686, bagian ``CONCLUSIONS'', PDF
p.~11).

{[}Laporan sumber{]} Kesimpulan substantifnya adalah bahwa broadband
hunian berkaitan dengan komuter lebih panjang di wilayah remote, yang
dapat dijelaskan oleh kerja dari rumah, serta dengan pekerjaan yang
lebih tinggi di dalam inti urban. Sebaliknya, penetrasi broadband tempat
memperlihatkan pola \emph{borrowed size}: efek pekerjaan berada pada
jarak 20--50 km dan efek komuter pada 20--30 km dari inti urban. Penulis
menyatakan bahwa konektivitas broadband tampaknya mendorong ekspansi
spasial pasar kerja urban (TXT baris 687--712, bagian ``CONCLUSIONS'',
PDF pp.~11--12).

{[}Interpretasi penulis{]} Kesimpulan tersebut tidak setara dengan klaim
kausal tanpa syarat. Dalam bagian metodologi, penulis menyatakan bahwa
endogenitas upaya lokal dapat membuat estimasi menjadi batas atas; dalam
kesimpulan, mereka secara khusus menolak menyatakan hubungan kausal
untuk broadband hunian dan \emph{borrowed size} (TXT baris 442--450 dan
690--697, bagian ``METHODOLOGY'' dan ``CONCLUSIONS'', PDF pp.~6--7 dan
11).

\subsection{Contributions}\label{contributions-29}

{[}Laporan sumber{]} Kontribusi yang dinyatakan atau ditunjukkan artikel
adalah:

\begin{itemize}
\tightlist
\item
  Mengisi kekosongan empiris mengenai hubungan broadband internet dan
  \emph{borrowed size}, yang menurut pendahuluan belum pernah diselidiki
  secara empiris (TXT baris 156--159, bagian ``INTRODUCTION'', PDF
  p.~2).
\item
  Memperluas pembahasan \emph{borrowed size} ke outcome pasar kerja,
  yang disebut kurang mendapat perhatian dibandingkan layanan bisnis dan
  fungsi urban (TXT baris 182--192, bagian ``Borrowed size in labour
  markets'', PDF p.~2).
\item
  Menggabungkan data register Swedia dengan data broadband yang
  menelusuri difusi teknologi pada 2007--2015 (TXT baris 182--187,
  bagian ``INTRODUCTION'', PDF p.~2).
\item
  Memisahkan broadband pada tingkat hunian dari penetrasi pada tingkat
  tempat untuk membedakan keputusan individual dan lingkungan ekonomi
  lokal (TXT baris 193--198 dan 675--686, bagian ``INTRODUCTION'' dan
  ``CONCLUSIONS'').
\item
  Menguji heterogenitas antara inti urban, wilayah di luar inti, rural,
  dan jarak ke inti terdekat, bukan hanya melaporkan satu efek rata-rata
  (TXT baris 203--206 dan 522--538, bagian ``INTRODUCTION'' dan
  ``RESULTS'').
\item
  Menganalisis efek pada skala geografis yang disebut lebih rinci
  daripada Atasoy serta memisahkan broadband hunian dari penetrasi
  tempat, sebagaimana dinyatakan dalam tinjauan pustaka artikel (TXT
  baris 238--248, bagian ``Local performance and IT'', PDF p.~3).
\end{itemize}

{[}Interpretasi penulis{]} Artikel menyebut dirinya sebagai salah satu
studi awal yang secara eksplisit memodelkan mekanisme yang mendasari
proses \emph{borrowed size}. Pernyataan ini adalah klaim kontribusi
penulis dalam artikel dan tidak diverifikasi terhadap literatur di luar
TXT (TXT baris 680--686, bagian ``CONCLUSIONS'', PDF p.~11).

\subsection{Research Gap}\label{research-gap-29}

{[}Laporan sumber{]} Gap awal yang secara eksplisit dirumuskan adalah
tidak adanya penelitian empiris yang menguji hubungan broadband internet
dan \emph{borrowed size}. Gap kedua adalah kurangnya perhatian pada
outcome pasar kerja dalam literatur \emph{borrowed size}, yang lebih
sering membahas layanan bisnis dan akses ke fungsi perkotaan (TXT baris
156--159 dan 182--192, bagian ``INTRODUCTION'' dan ``Borrowed size in
labour markets'', PDF p.~2).

{[}Laporan sumber{]} Artikel juga menyatakan gap pengukuran dan
mekanisme: studi ini perlu memisahkan akses broadband individual dari
konektivitas jaringan tempat, serta menguji apakah efek berbeda menurut
jarak dan posisi dalam hierarki urban. Artikel Atasoy disebut belum
menguji heterogenitas dengan cara yang sama dan menggunakan skala
geografis yang lebih kasar (TXT baris 160--175 dan 229--248, bagian
``Borrowed size in labour markets'' dan ``Local performance and IT'',
PDF p.~3).

{[}Interpretasi penulis{]} Setelah studi ini, gap yang masih terbuka
meliputi status kausal hubungan tersebut, generalisasi di luar Swedia,
mekanisme kerja dari rumah, ukuran \emph{matching} yang lebih baik,
dampak 4G, faktor lokal yang menentukan kapasitas \emph{borrowed size},
jenis manfaat aglomerasi yang dapat dipinjam, dan mekanisme alternatif
(TXT baris 663--686 dan 713--721, bagian ``CONCLUSIONS'', PDF
pp.~11--12).

\subsection{Limitations}\label{limitations-29}

{[}Laporan sumber{]} Keterbatasan yang dinyatakan atau diakui dalam
artikel meliputi:

\begin{itemize}
\tightlist
\item
  \textbf{Generalisasi geografis:} Swedia memiliki kombinasi khusus
  berupa tingkat urbanisasi tinggi dan jarak panjang antara tempat
  sentral dan remote; artikel meminta pengujian di negara lain (TXT
  baris 713--718, bagian ``CONCLUSIONS'', PDF p.~11).
\item
  \textbf{Confounding dan kausalitas:} pertumbuhan broadband dapat
  berkorelasi dengan proyek infrastruktur, perumahan, dan dinamika lain
  yang sekaligus mendorong pekerjaan. Kontrol dinilai menangkap banyak
  confounder, tetapi penulis menyatakan riset lebih lanjut tentang sifat
  kausal masih diperlukan (TXT baris 442--450 dan 718--721, bagian
  ``METHODOLOGY'' dan ``CONCLUSIONS'', PDF pp.~6--7 dan 12).
\item
  \textbf{Endogenitas upaya lokal:} inisiatif komunitas dan timing
  kedatangan broadband tidak sepenuhnya acak. Jika usaha lokal endogen,
  penulis meminta hasil dibaca sebagai batas atas efek pengenalan
  broadband (TXT baris 442--450, bagian ``METHODOLOGY'', PDF pp.~6--7).
\item
  \textbf{Proksi pencocokan pasar kerja:} tingkat pekerjaan dipakai
  untuk mengukur \emph{better matching}, tetapi tidak mengukur
  kesesuaian pekerja--perusahaan seperti overeducation atau gaji (TXT
  baris 663--669, bagian ``CONCLUSIONS'', PDF p.~11).
\item
  \textbf{Tidak ada ukuran kerja dari rumah:} kerja dari rumah dipakai
  sebagai penjelasan yang masuk akal untuk efek komuter di wilayah
  remote, tetapi data tingkat kerja dari rumah dan jumlah komuter tidak
  tersedia sehingga bukti mekanismenya belum konklusif (TXT baris
  670--675, bagian ``CONCLUSIONS'', PDF p.~11).
\item
  \textbf{Perubahan teknologi:} efek kecil 3G diamati sebelum 4G;
  penulis menyatakan hasilnya mungkin berubah setelah 4G hadir (TXT
  baris 667--671, bagian ``CONCLUSIONS'', PDF p.~11).
\item
  \textbf{Konstruksi konsep:} \emph{Borrowed size} masih memerlukan
  substansiasi empiris, terutama mengenai faktor lokal yang memengaruhi
  kapasitas meminjam ukuran dan manfaat aglomerasi yang dapat dipinjam
  (TXT baris 675--686, bagian ``CONCLUSIONS'', PDF p.~11).
\item
  \textbf{Jarak komuter:} jarak Euclidean dapat mengandung error
  pengukuran dan penulis tidak menentukan arah error tersebut (TXT baris
  308--319, bagian ``ASTRID data'', PDF p.~4).
\item
  \textbf{Ukuran pekerjaan:} ambang pendapatan dapat membuat mahasiswa
  atau pekerja paruh waktu dikategorikan sebagai tidak bekerja, walaupun
  penulis mengharapkan error perubahan antar-tahun terbatas (TXT baris
  323--326 dan 286--294, bagian ``ASTRID data'', PDF p.~4).
\item
  \textbf{Bobot spasial:} jarak komuter rata-rata relatif tinggi karena
  grid berpenduduk rendah yang memiliki komuter lebih jauh diberi bobot
  sama dengan grid berpenduduk tinggi (TXT baris 372--380, ``Summary
  statistics'', PDF p.~6).
\item
  \textbf{Definisi FUA:} definisi OECD yang monosentris dapat tidak
  mencerminkan tumpang tindih wilayah pasar kerja antarkota, walaupun
  dipilih karena konsisten dan dapat direproduksi (TXT baris 362--375,
  bagian ``METHODOLOGY'', PDF p.~5).
\item
  \textbf{Definisi ketersediaan broadband:} broadband dinyatakan
  tersedia apabila sedikitnya satu alamat dalam grid memiliki koneksi.
  File tidak menjelaskan lebih lanjut bagaimana ukuran alamat pada grid
  diagregasi menjadi seluruh ukuran penetrasi tingkat tempat, sehingga
  rincian denominator ukuran tersebut tidak tersedia (TXT baris 308--326
  dan 535--542, bagian ``Broadband data'' dan ``RESULTS'', PDF pp.~4 dan
  8--9).
\end{itemize}

\subsection{Challenges}\label{challenges-29}

{[}Inferensi pengolahan{]} TXT tidak menyediakan bagian terpisah
berjudul ``Challenges''. Heading ini merangkum kesulitan analitis yang
dilaporkan penulis, bukan menambahkan tantangan dari luar sumber:

\begin{itemize}
\tightlist
\item
  Memisahkan arah hubungan broadband dan outcome karena pekerjaan serta
  perilaku komuter dapat mendorong adopsi broadband; penulis
  meresponsnya dengan lag dua tahun, kontrol, dan fixed effects, tetapi
  tidak menganggap masalah selesai (TXT baris 376--398 dan 442--453,
  bagian ``METHODOLOGY'').
\item
  Mengisolasi efek individual dari efek lingkungan tempat karena
  ketersediaan broadband hunian dapat berkorelasi dengan penetrasi
  broadband tingkat tempat; artikel memasukkan kedua ukuran sekaligus,
  dengan peringatan bahwa efek tempat harus dibaca hati-hati (TXT baris
  384--398 dan 522--539, bagian ``METHODOLOGY'' dan ``RESULTS'').
\item
  Menghadapi multikolinearitas pada kontrol tingkat wilayah pasar kerja,
  sehingga penulis berfokus pada variabel utama yang memiliki VIF rendah
  dan tidak membahas seluruh koefisien kontrol sebagai temuan substantif
  (TXT baris 497--501, bagian ``RESULTS'', PDF p.~8).
\item
  Menyamakan unit geografis dengan proses urban yang tidak sepenuhnya
  monosentris; batas FUA yang bersentuhan belum tentu menggambarkan
  tumpang tindih pasar kerja sebenarnya (TXT baris 359--375, bagian
  ``METHODOLOGY'', PDF p.~5).
\item
  Mengukur akses pasar kerja melalui jarak Euclidean dan tingkat
  pekerjaan yang masing-masing memiliki keterbatasan pengukuran
  sebagaimana dijelaskan pada bagian data (TXT baris 308--326, bagian
  ``ASTRID data'', PDF p.~4).
\item
  Menguji mekanisme kerja dari rumah tanpa variabel langsung tentang
  frekuensi kerja dari rumah atau jumlah komuter (TXT baris 670--675,
  bagian ``CONCLUSIONS'', PDF p.~11).
\end{itemize}

\subsection{Practical Implications}\label{practical-implications-29}

{[}Laporan sumber{]} Penulis menarik tiga implikasi kebijakan:

\begin{itemize}
\tightlist
\item
  Pemerintah yang ingin meningkatkan kondisi ekonomi tempat melalui
  broadband perlu menyadari bahwa dalam jangka pendek efek terhadap
  pekerjaan kecil dan sangat bergantung pada kedekatan dengan pusat
  pekerjaan lain (TXT baris 687--694, bagian ``CONCLUSIONS'', PDF
  p.~11).
\item
  Pemerintah sebaiknya mempertimbangkan prioritas pada koneksi tetap
  karena efek mobile jauh lebih kecil, dengan catatan bahwa hasil
  tersebut berasal dari periode sebelum 4G. Artikel tidak menemukan efek
  dari peningkatan DSL ke kabel atau serat optik (TXT baris 692--697,
  bagian ``CONCLUSIONS'', PDF p.~11).
\item
  Untuk wilayah di luar kota, ketersediaan broadband di hunian tidak
  menunjukkan efek pada prospek pekerjaan penduduk, sedangkan penetrasi
  broadband tingkat tempat menunjukkan efek signifikan. Karena itu,
  kebijakan konektivitas cepat di wilayah rural menurut penulis
  sebaiknya pada tahap awal memprioritaskan koneksi tempat, bukan
  menghubungkan setiap penduduk secara individual (TXT baris 698--705,
  bagian ``CONCLUSIONS'', PDF pp.~11--12).
\end{itemize}

\subsection{Applications}\label{applications-29}

{[}Laporan sumber{]} File tidak memuat bagian terpisah berjudul
``Applications'' dan tidak melaporkan studi implementasi kebijakan
tertentu. Penerapan yang secara eksplisit dibahas berada dalam bentuk
implikasi kebijakan pada bagian kesimpulan, bukan evaluasi program
setelah kebijakan diterapkan (TXT baris 687--705, bagian
``CONCLUSIONS'', PDF pp.~11--12).

{[}Inferensi pengolahan{]} Berdasarkan desain yang benar-benar dipakai,
aplikasi analitis yang paling langsung adalah mengevaluasi intervensi
broadband dengan memisahkan indikator koneksi hunian dan penetrasi
tingkat tempat, menggunakan jeda waktu, serta membandingkan inti FUA,
wilayah di luar inti, wilayah di luar FUA, dan kategori jarak. Ini
adalah kemungkinan penggunaan kerangka penelitian, bukan hasil aplikasi
baru yang dilaporkan artikel (TXT baris 338--375 dan 522--538, bagian
``METHODOLOGY'' dan ``RESULTS'').

\subsection{Future Research
(Optional)}\label{future-research-optional-29}

{[}Laporan sumber{]} Arah penelitian lanjutan yang disebutkan penulis
adalah:

\begin{itemize}
\tightlist
\item
  Menguji apakah hasil Swedia juga berlaku di negara lain (TXT baris
  713--718, bagian ``CONCLUSIONS'', PDF p.~11).
\item
  Meneliti lebih lanjut sifat kausal pengaruh broadband dengan
  menghadapi korelasi bersama infrastruktur, perumahan, dan dinamika
  ekonomi (TXT baris 718--721, bagian ``CONCLUSIONS'', PDF p.~12).
\item
  Menggunakan ukuran \emph{matching} pasar kerja selain tingkat
  pekerjaan, termasuk overeducation dan gaji (TXT baris 663--669, bagian
  ``CONCLUSIONS'', PDF p.~11).
\item
  Menambahkan tingkat kerja dari rumah dan jumlah komuter agar
  penjelasan mekanisme kerja dari rumah dapat diuji lebih konklusif (TXT
  baris 670--675, bagian ``CONCLUSIONS'', PDF p.~11).
\item
  Menilai apakah hadirnya 4G mengubah peran mobile internet dibandingkan
  broadband tetap (TXT baris 667--671, bagian ``CONCLUSIONS'', PDF
  p.~11).
\item
  Memberi substansiasi empiris lebih lanjut pada konsep \emph{borrowed
  size}, termasuk faktor lokal yang menentukan kapasitas meminjam
  ukuran, manfaat aglomerasi yang dapat dipinjam, dan mekanisme
  alternatif yang mungkin mendasarinya (TXT baris 675--686, bagian
  ``CONCLUSIONS'', PDF p.~11).
\end{itemize}

\subsection{Provenance Notes}\label{provenance-notes-29}

\begin{itemize}
\tightlist
\item
  \textbf{Bahan yang digunakan:} hanya file
  \texttt{source/journal/A14-de-vos-lindgren-van-ham-meijers-2020-does-broadband-internet-allow-cities-to-borrow-size-evidence-from-the-swedish-labour-market-10.1080-00343404.2019.1699238.txt}.
  Seluruh 800 baris yang tersedia dibaca, termasuk blok metadata PDF
  pada baris 1--28, teks artikel pada baris 29--721, catatan, ORCID, dan
  daftar referensi pada baris 728--799.
\item
  \textbf{Provenans ekstraksi:} metadata menyatakan sumber PDF berada di
  folder \texttt{journal/}, diekstrak pada 31 Agustus 2026 dengan
  \texttt{pdftotext\ -layout}; PDF memiliki 13 halaman dan tidak
  terenkripsi (TXT baris 1--28, blok ``{[}PDF Metadata{]}'').
\item
  \textbf{Aturan locator:} rujukan utama memakai nomor baris TXT dan,
  bila tersedia, halaman PDF, bagian, tabel, atau gambar. Nomor halaman
  dalam review merujuk penanda halaman artikel yang muncul di hasil
  ekstraksi, bukan paginasi jurnal lengkap.
\item
  \textbf{Pembedaan evidensi:} \texttt{{[}Laporan\ sumber{]}} berarti
  artikel secara eksplisit melaporkan data, metode, hasil, atau
  pernyataan literatur; \texttt{{[}Interpretasi\ penulis{]}} berarti
  mekanisme, makna, atau klaim kausal yang dinyatakan penulis;
  \texttt{{[}Inferensi\ pengolahan{]}} berarti penyusunan atau pembacaan
  analitis review ini dan tidak ditulis sebagai temuan baru artikel.
\item
  \textbf{Materi yang tidak tersedia:} isi lengkap data suplementer,
  Tabel A1, Tabel A2, Tabel A4, dan Gambar A1 tidak ada dalam TXT;
  artikel hanya merujuknya. Data mentah, kode analisis, rincian
  pemilihan observasi untuk setiap N regresi, denominator lengkap ukuran
  penetrasi tingkat tempat, volume/nomor/halaman jurnal, dan status
  peer-review juga tidak disebutkan dalam file.
\item
  \textbf{Batas penggunaan referensi:} daftar referensi dan ringkasan
  studi terdahulu diperlakukan hanya sebagai bagian dari apa yang
  dilaporkan artikel A14. Referensi tersebut tidak dibaca atau
  diverifikasi sebagai bukti terpisah.
\item
  \textbf{Perubahan file:} review ini dibuat sebagai
  \texttt{output/kajian/review-A14.md}. Tidak ada perubahan pada TXT
  sumber, \texttt{Todo.md}, daftar kode, atau file lain.
\end{itemize}

\section{Review A15 - Inter-regional networks and productive efficiency
in
Japan}\label{review-a15---inter-regional-networks-and-productive-efficiency-in-japan}

Status: review literatur berbasis teks lengkap yang tersedia dalam TXT.
Seluruh substansi di bawah ini hanya berasal dari file sumber A15.

\subsection{Identitas Sumber}\label{identitas-sumber-12}

\begin{itemize}
\tightlist
\item
  \textbf{Kode sumber:} A15.
\item
  \textbf{Judul:} \emph{Inter-regional networks and productive
  efficiency in Japan}.
\item
  \textbf{Penulis:} Akihiro Otsuka.
\item
  \textbf{Afiliasi:} Association of International Arts and Science,
  Yokohama City University, Yokohama, Japan.
\item
  \textbf{Tahun publikasi yang tercantum dalam TXT:} 2019. Baris sitasi
  menyebut \emph{Papers in Regional Science}, 2019, 99, 115-133; nama
  berkas memuat angka 2020, tetapi angka tersebut tidak dipakai untuk
  mengubah identitas yang dinyatakan di dalam artikel.
\item
  \textbf{Jurnal:} \emph{Papers in Regional Science}.
\item
  \textbf{Volume dan halaman:} Volume 99, halaman 115-133.
\item
  \textbf{Nomor terbitan:} Tidak disebutkan dalam file.
\item
  \textbf{DOI:} 10.1111/pirs.12474.
\item
  \textbf{Riwayat naskah:} diterima 22 April 2019, direvisi 19 Juli
  2019, dan diterima 7 Agustus 2019.
\item
  \textbf{Jenis sumber:} artikel jurnal empiris dengan analisis ekonomi
  regional.
\item
  \textbf{Status penelaahan sejawat:} Tidak diverifikasi secara formal
  dari file. Bagian ucapan terima kasih menyebut reviewers, tetapi file
  tidak menyatakan status peer review dengan formula khusus.
\item
  \textbf{Pendanaan:} JSPS KAKENHI, Grant/Award Number 18K01614.
\item
  \textbf{Pernyataan konflik kepentingan:} Penulis menyatakan tidak ada
  konflik kepentingan dan sponsor tidak berperan dalam desain,
  pengumpulan, analisis, interpretasi data, penulisan, atau keputusan
  publikasi.
\item
  \textbf{Kata kunci:} agglomeration, inter-regional network, Japan,
  regional disparities, dan stochastic frontier analysis.
\item
  \textbf{Klasifikasi JEL:} R10 dan R11.
\item
  \textbf{Bahasa:} teks utama berbahasa Inggris. TXT juga memuat abstrak
  berbahasa Spanyol dan Jepang pada bagian akhir.
\item
  \textbf{Locator identitas:} hlm. 115 untuk judul, penulis, afiliasi,
  abstrak, kata kunci, dan klasifikasi JEL; hlm. 126 untuk pernyataan
  konflik kepentingan; hlm. 129-133 untuk materi lampiran; baris 780-781
  TXT untuk sitasi artikel.
\end{itemize}

\subsection{Summarized Abstract}\label{summarized-abstract-30}

\textbf{Laporan sumber.} Artikel berangkat dari ketidakmampuan teori
aglomerasi tradisional menjelaskan penyusutan disparitas regional Jepang
yang berkelanjutan. Dengan stochastic frontier analysis, penelitian ini
menguji peran jaringan antarwilayah dalam meningkatkan efisiensi
produktif melalui konsep \emph{borrowed size}. Hasil yang dilaporkan
menyatakan bahwa perbaikan jaringan antarwilayah memperkuat ekonomi
lokal dan mendorong konvergensi disparitas regional. \emph{Borrowed
size} dari aglomerasi kota besar dinyatakan sangat efektif di wilayah
lokal, dan pembangunan jaringan transportasi memungkinkan wilayah lokal
mengejar wilayah metropolitan besar. Penulis menyimpulkan bahwa jaringan
transportasi berkecepatan tinggi dapat membantu mengurangi disparitas
regional dalam ekonomi Jepang. (Hlm. 115, abstrak.)

\textbf{Laporan sumber.} Analisis menggunakan panel 47 prefektur Jepang
selama 1980-2014. Efisiensi produktif diestimasi untuk seluruh industri,
industri manufaktur, dan industri non-manufaktur. Kepadatan penduduk
dipakai untuk merepresentasikan kekuatan aglomerasi lokal, sedangkan
indeks akses pasar yang menggabungkan skala penduduk wilayah lain dan
waktu tempuh dipakai untuk merepresentasikan jaringan antarwilayah.
(Hlm. 118-121, bagian 3; hlm. 131-133, Appendix C.)

\textbf{Interpretasi penulis.} Penulis membaca koefisien akses pasar
yang positif sebagai bukti bahwa jaringan memperluas manfaat aglomerasi
kota besar ke wilayah lokal. Efek jaringan dinilai lebih besar pada
manufaktur daripada non-manufaktur, dan pengaruhnya terhadap efisiensi
produktif dianggap lebih besar daripada manfaat aglomerasi penduduk
lokal. Karena kontribusi akses pasar mengurangi disparitas lebih besar
daripada kontribusi kepadatan penduduk yang memperlebar disparitas,
penulis menyatakan bahwa jaringan transportasi membantu konvergensi
regional. (Hlm. 123-127, bagian 4-5.)

\textbf{Inferensi pengolahan.} Bukti yang tersedia adalah hasil estimasi
SFA pada data panel wilayah, bukan pengamatan langsung atas arus
penduduk, perjalanan bisnis, pengetahuan, atau transaksi antardaerah.
Indeks akses pasar juga menjumlahkan populasi wilayah lain yang
didiskontokan oleh waktu tempuh, sehingga merupakan proksi potensi pasar
dan jaringan, bukan ukuran langsung atas skala yang benar-benar dipinjam
oleh setiap wilayah. Kesimpulan tentang \emph{borrowed size} dan
pengurangan disparitas karena itu sebaiknya dibaca sebagai interpretasi
atas hubungan model dan dekomposisi, bukan sebagai identifikasi kausal
langsung. (Hlm. 119-121, persamaan 3-5; hlm. 124-126, Tabel 3-4.)

\subsection{Problem Statement}\label{problem-statement-30}

\textbf{Laporan sumber.} Teori aglomerasi tradisional menekankan skala,
kedekatan geografis, biaya transportasi, dan limpahan pengetahuan
sebagai sumber produktivitas. Teori tersebut memprediksi bahwa
pertumbuhan kota besar akan berlangsung dengan mengorbankan wilayah
perifer dan bahwa efek \emph{lock-in} aglomerasi akan memperlebar
disparitas. Namun, di Jepang, konsentrasi kegiatan ekonomi di Tokyo
meningkat sejak pertengahan 1990-an, sementara disparitas regional tidak
meningkat secara terus-menerus. Kekuatan ekonomi wilayah lokal justru
dilaporkan mengejar dan kesenjangan dengan kota besar menyempit. (Hlm.
115-116, bagian 1.)

\textbf{Laporan sumber.} Jepang juga mengembangkan jaringan transportasi
berkecepatan tinggi dengan Tokyo sebagai hub. Jaringan tersebut
menurunkan waktu tempuh, memperbesar peluang kontak, dan memungkinkan
perusahaan di luar aglomerasi mengakses ekonomi aglomerasi. Penelitian
mengenai perluasan cakupan geografis manfaat aglomerasi melalui jaringan
transportasi telah tersedia untuk banyak negara, tetapi menurut artikel
penelitian untuk Jepang masih tidak memadai. (Hlm. 116, bagian 1.)

\textbf{Laporan sumber.} Studi sebelumnya di Jepang menunjukkan hubungan
positif antara aglomerasi, jaringan antarwilayah, dan efisiensi
produktif pada industri tertentu. Akan tetapi, semakin dekatnya jarak ke
kota besar juga dapat meningkatkan kompetisi antarwilayah dan
menghasilkan \emph{agglomeration shadows}, yang dalam konteks Jepang
disebut fenomena \emph{pull-in}. Artikel karena itu mempertanyakan
apakah efek positif jaringan atau efek negatif kompetisi yang lebih
dominan. (Hlm. 116-117, bagian 1.)

\textbf{Laporan sumber.} Masalah tersebut memiliki dimensi tambahan.
Wilayah metropolitan besar seperti Greater Tokyo dan Kansai memiliki
kepadatan penduduk serta kekuatan aglomerasi tinggi, sedangkan banyak
wilayah lokal mengalami penurunan kepadatan dan pelemahan aglomerasi.
Pada saat yang sama, perbaikan infrastruktur transportasi menarik
investasi pabrik ke wilayah lokal dan mendukung pembagian kerja vertikal
manufaktur di seluruh Jepang. Belum ada pemeriksaan yang memadai
mengenai apakah jaringan antarwilayah menyumbang pada penyempitan
disparitas efisiensi produktif. (Hlm. 116-117, bagian 1.)

\textbf{Interpretasi penulis.} \emph{Borrowed size} ditempatkan sebagai
penjelasan yang dapat menjembatani kontradiksi antara meningkatnya
konsentrasi Tokyo dan menyempitnya disparitas regional. Wilayah lokal
dapat meminjam manfaat aglomerasi kota besar melalui waktu tempuh yang
lebih pendek tanpa harus memiliki ukuran aglomerasi yang sama. Namun,
manfaat tersebut perlu dibedakan dari kemungkinan kompetisi dan
penarikan kegiatan menuju pusat. (Hlm. 116-118, bagian 1-2.)

\textbf{Inferensi pengolahan.} Pertanyaan inti artikel bukan hanya
apakah aksesibilitas berkaitan dengan efisiensi, melainkan apakah
hubungan tersebut cukup untuk menjelaskan konvergensi antarprefektur.
Ini membutuhkan dua tahap pembacaan: mengestimasi efisiensi dan
kontribusi akses pasar, lalu menguji apakah kontribusi tersebut
berkaitan dengan perubahan disparitas. File tidak menunjukkan desain
kontrafaktual yang memisahkan secara tegas investasi jaringan dari
karakteristik wilayah yang mendorong investasi tersebut.

\subsection{Objectives}\label{objectives-30}

\textbf{Laporan sumber.} Tujuan penelitian yang dapat direkonstruksi
dari pendahuluan adalah sebagai berikut:

\begin{itemize}
\tightlist
\item
  menganalisis pengaruh aglomerasi dan jaringan antarwilayah terhadap
  efisiensi produktif ekonomi regional Jepang;
\item
  menguji apakah jaringan antarwilayah menghasilkan efek positif
  \emph{borrowed size} atau efek negatif \emph{agglomeration shadow};
\item
  menguji apakah perbaikan jaringan berhubungan dengan penyempitan
  disparitas efisiensi produktif antardaerah;
\item
  membandingkan efek jaringan pada industri manufaktur dan
  non-manufaktur;
\item
  menilai apakah pengaruh jaringan lebih besar daripada pengaruh
  aglomerasi penduduk lokal; dan
\item
  menyediakan bukti empiris untuk menilai relevansi kebijakan jaringan
  antarwilayah dan gagasan pembentukan \emph{super-megaregion} dalam
  rencana tata ruang nasional Jepang. (Hlm. 116-117, bagian 1.)
\end{itemize}

\textbf{Inferensi pengolahan.} Tujuan tersebut tidak disajikan sebagai
daftar hipotesis formal dengan ambang atau prediksi koefisien terpisah.
Struktur model menunjukkan ekspektasi bahwa \texttt{deltaDEN} dan
\texttt{deltaACC} positif berarti kepadatan penduduk dan akses pasar
mengurangi inefisiensi, sedangkan dekomposisi beta menguji arah
kontribusi keduanya terhadap disparitas. (Hlm. 119-121, persamaan 3-8.)

\subsection{Literature Survey}\label{literature-survey-30}

\textbf{Laporan sumber.} Artikel memulai tinjauan dari teori aglomerasi.
Konsentrasi penduduk dan perusahaan, ekonomi skala, serta kedekatan
geografis menjelaskan pertukaran informasi, pengetahuan, inovasi, dan
peningkatan daya saing. Teori tersebut menghubungkan ukuran perkotaan
dengan efisiensi dan pertumbuhan, tetapi tidak cukup menjelaskan wilayah
lokal yang tumbuh lebih cepat atau mengejar kota besar. (Hlm. 115-116,
bagian 1; hlm. 117, bagian 2.)

\textbf{Laporan sumber.} Teori jaringan kota dipakai untuk menjelaskan
kontradiksi tersebut. Kota dapat meminjam fungsi kota lain bukan hanya
melalui kedekatan fisik, tetapi juga melalui hubungan horizontal dan
non-hierarkis. Konsep \emph{urban network externalities} memperluas
cakupan manfaat perkotaan melampaui batas fisik kota. Artikel mengaitkan
gagasan ini dengan Camagni, Capello, Hall dan Pain, serta literatur
jaringan kota lain yang disebut dalam teks. (Hlm. 117, bagian 2.)

\textbf{Laporan sumber.} Titik awal konsep \emph{borrowed size}
ditelusuri ke Alonso. Kota kecil dapat menggunakan fasilitas belanja dan
hiburan, gudang dan jasa bisnis, serta pasar tenaga kerja kota lain
untuk melengkapi kapasitasnya sendiri. Dengan demikian, kota kecil dapat
memperoleh sebagian manfaat aglomerasi sembari menghindari sebagian
biaya kemacetan. Artikel juga menyebut konsep \emph{externality fields}
dan \emph{regional externalities} untuk menegaskan bahwa keuntungan
urban dapat meluas jauh di luar batas kota. (Hlm. 117, bagian 2.)

\textbf{Laporan sumber.} Ukuran tradisional \emph{borrowed size}
dijelaskan sebagai jumlah populasi kota lain yang didiskontokan oleh
jarak: \texttt{borrowed\ size\_c\ =\ jumlah(pop\_j\ /\ w\_cj)} untuk
kota \texttt{j} yang berbeda dari kota \texttt{c}. Ukuran ini meningkat
ketika populasi kota tetangga bertambah dan ketika jarak ke kota besar
memendek. Artikel menilai ukuran tersebut terutama menangkap ekonomi
urbanisasi positif melalui konektivitas. (Hlm. 118, persamaan 1.)

\textbf{Laporan sumber.} Literatur geografi ekonomi baru dipakai untuk
menjelaskan sisi negatif konektivitas. Jaringan dapat meningkatkan
kompetisi dan menciptakan \emph{agglomeration shadows}. Dalam konteks
Jepang, efek negatif kompetisi disebut \emph{pull-in phenomenon}. Karena
itu, jaringan antarwilayah dipahami memiliki efek positif dan negatif,
dan pertanyaan empirisnya adalah dampak bersih mana yang dominan. (Hlm.
118, bagian 2.)

\textbf{Laporan sumber.} Literatur empiris yang dirujuk mencakup studi
tentang efisiensi produktif pada industri atau wilayah tertentu di
berbagai negara. Untuk Jepang, Otsuka, Goto, dan Sueyoshi serta Otsuka
dan Goto dilaporkan menemukan pengaruh positif aglomerasi dan perbaikan
jaringan terhadap efisiensi produktif industri. Artikel membedakan
manufaktur sebagai industri dasar atau ekspor dengan transaksi lintas
wilayah yang relatif banyak, sedangkan non-manufaktur dipandang lebih
berpusat pada jasa lokal dan kontak antarpersonal. (Hlm. 116-117, bagian
1; hlm. 117-118, bagian 2.)

\textbf{Interpretasi penulis.} Tinjauan pustaka mengarahkan artikel pada
perluasan unit analisis dari kota atau industri tertentu menuju ekonomi
regional nasional, serta pada pengujian apakah jaringan melengkapi
aglomerasi lokal. Pemisahan manufaktur dan non-manufaktur dipakai untuk
menguji dugaan bahwa aktivitas manufaktur lebih sensitif terhadap
jaringan perjalanan lintas wilayah. (Hlm. 116-117, bagian 1.)

\textbf{Inferensi pengolahan.} Tinjauan ini bersifat naratif dan
teoritis. File tidak menyebutkan strategi pencarian, kriteria inklusi,
jumlah studi yang diseleksi, penilaian mutu, atau prosedur telaah
sistematis. Daftar pustaka yang panjang tidak dengan sendirinya
membuktikan bahwa seluruh literatur relevan telah dicakup.

\subsection{Method Used}\label{method-used-30}

\textbf{Laporan sumber.} Penelitian menggunakan stochastic frontier
analysis dengan pendekatan satu tahap dari Battese dan Coelli.
Pendekatan ini mengestimasi fungsi frontier produksi dan determinan
inefisiensi secara simultan. Artikel menolak pendekatan dua tahap karena
distribusi inefisiensi yang diasumsikan pada fungsi frontier dapat tidak
konsisten dengan persamaan regresi determinan inefisiensi. Estimasi
dilakukan dengan maximum likelihood menggunakan Frontier 4.1. (Hlm.
118-119 dan 121, bagian 3.1 dan 4.)

\textbf{Laporan sumber.} Fungsi frontier produksi ditulis sebagai fungsi
produksi translog. Dalam bentuk ringkas, \texttt{ln\ Y\_jt} dijelaskan
oleh log tenaga kerja \texttt{ln\ L\_jt}, log modal \texttt{ln\ K\_jt},
tren waktu \texttt{T}, kuadrat dan interaksi variabel tersebut, efek
tetap prefektur \texttt{alpha\_j}, serta error acak
\texttt{v\_jt\ -\ u\_jt}. \texttt{Y} adalah output riil, \texttt{L}
input tenaga kerja, \texttt{K} input modal, dan \texttt{T} tren waktu
yang merepresentasikan kemajuan teknologi. Efek tetap prefektur dipakai
untuk mengontrol karakteristik setiap prefektur yang tidak tertangkap
oleh variabel penjelas. (Hlm. 119, persamaan 2.)

\textbf{Laporan sumber.} Error terdiri atas \texttt{v\_jt}, yang
diasumsikan berdistribusi normal dan independen dari \texttt{u\_jt}
serta variabel penjelas, dan \texttt{u\_jt}, yaitu komponen inefisiensi
yang bersifat non-negatif. Efisiensi produktif \texttt{TEF\_jt}
didefinisikan sebagai rasio output yang diamati terhadap output pada
frontier, sehingga \texttt{TEF\_jt\ =\ exp(-u\_jt)} dan berada di atas
nol serta tidak lebih besar dari satu. (Hlm. 119, persamaan 2 dan
definisi TEF.)

\textbf{Laporan sumber.} Rata-rata inefisiensi dirumuskan sebagai
\texttt{mu\_jt\ =\ delta0\ -\ deltaDEN\ ln\ DEN\_jt\ -\ deltaACC\ ln\ ACC\_jt}.
\texttt{DEN} adalah kepadatan penduduk, yaitu penduduk per unit lahan
yang dapat dihuni, dan merepresentasikan kekuatan aglomerasi wilayah.
\texttt{ACC} adalah indeks akses pasar yang merepresentasikan kekuatan
jaringan antarwilayah. Karena determinan ditulis dengan tanda negatif,
koefisien \texttt{deltaDEN} dan \texttt{deltaACC} yang positif berarti
variabel tersebut mengurangi inefisiensi dan meningkatkan efisiensi.
(Hlm. 119, persamaan 3 dan catatan kaki 3.)

\textbf{Laporan sumber.} Indeks akses pasar dirumuskan sebagai jumlah,
untuk semua wilayah \texttt{k} selain \texttt{j}, dari populasi wilayah
\texttt{k} yang dibagi resistensi jarak \texttt{T\_jkt}. \texttt{T\_jkt}
adalah waktu tempuh antarwilayah. Waktu tersebut dihitung dari kombinasi
perjalanan udara, kereta api, dan mobil; untuk setiap kombinasi
asal-tujuan, waktu terpendek di antara moda yang tersedia dimasukkan,
lalu digunakan dalam perhitungan berbobot. Indeks meningkat ketika waktu
tempuh menurun dan ketika pasar besar lebih mudah diakses. (Hlm.
119-120, persamaan 4 dan catatan kaki 4.)

\textbf{Laporan sumber.} Perubahan efisiensi didekomposisi dengan
persamaan 5 menjadi kontribusi perubahan kepadatan penduduk dan
kontribusi perubahan akses pasar:
\texttt{dotTEF\_jt\ =\ deltaDEN\ dotDEN\_jt\ +\ deltaACC\ dotACC\_jt}.
Dekomposisi ini dipakai untuk menilai sumber perubahan efisiensi pada
masing-masing wilayah dan industri. (Hlm. 120, persamaan 5.)

\textbf{Laporan sumber.} Konvergensi diuji dengan meregresikan laju
perubahan tahunan rata-rata efisiensi dari periode awal ke periode
\texttt{t} terhadap tingkat efisiensi awal. Koefisien konvergensi
\texttt{beta} negatif berarti disparitas efisiensi menurun, sedangkan
koefisien positif berarti disparitas meningkat. Koefisien tersebut
dipecah menjadi \texttt{beta\ =\ betaDEN\ +\ betaACC}; kontribusi
positif memperlebar disparitas dan kontribusi negatif menguranginya.
(Hlm. 120-121, persamaan 6-8.)

\textbf{Laporan sumber.} Enam spesifikasi diestimasi untuk tiga cakupan
industri. Model A, C, dan E tidak memasukkan efek waktu, sedangkan model
B, D, dan F memasukkan efek waktu dalam determinan kepadatan penduduk
dan akses pasar. Artikel memilih B, D, dan F berdasarkan nilai
log-likelihood, setelah menyatakan tidak ada perbedaan signifikan antara
hasil dengan dan tanpa efek waktu. (Hlm. 121, bagian 4; Tabel 1.)

\textbf{Inferensi pengolahan.} Rancangan ini menggunakan data panel dan
efek tetap prefektur, tetapi file tidak melaporkan instrumen, eksperimen
alam, atau strategi kuasi-eksperimental untuk menangani kemungkinan
bahwa wilayah yang lebih produktif lebih mungkin memperoleh jaringan
transportasi. Estimasi SFA dapat mengendalikan komponen inefisiensi
sesuai asumsi model, tetapi tidak dengan sendirinya membuktikan bahwa
investasi jaringan menyebabkan peningkatan efisiensi.

\subsection{Dataset}\label{dataset-30}

\textbf{Laporan sumber.} Dataset dan konstruksi variabel dijelaskan
terutama pada Appendix C:

\begin{itemize}
\tightlist
\item
  \textbf{Output riil atau nilai tambah (\texttt{Y}):} berasal dari
  \emph{Annual Report on Prefectural Accounts} milik Cabinet Office dan
  dideflasi dengan GDP deflator.
\item
  \textbf{Input tenaga kerja (\texttt{L}):} dihitung dengan mengalikan
  jumlah pekerja dari \emph{Annual Report on Prefectural Accounts}
  dengan jumlah jam kerja dari \emph{Monthly Labour Survey} milik
  Ministry of Health, Labour and Welfare.
\item
  \textbf{Input modal (\texttt{K}):} stok modal swasta dikalikan rasio
  utilisasi kapasitas. Data stok modal swasta diestimasi oleh Central
  Research Institute of Electric Power Industry. Rasio utilisasi
  kapasitas manufaktur berasal dari indeks yang diterbitkan Ministry of
  Economy, Trade and Industry.
\item
  \textbf{Rasio utilisasi kapasitas non-manufaktur:} karena statistik
  resmi tidak tersedia, penulis mengasumsikan perubahan \texttt{ln(Y/K)}
  berasal dari perubahan utilisasi modal dan tren waktu. Proksi
  utilisasi dibentuk dari residual \texttt{epsilon} dalam regresi
  \texttt{ln(Y/K)\ =\ alpha\ +\ betaT\ +\ epsilon}, dengan kemiringan
  \texttt{beta} yang tidak berubah menurut waktu.
\item
  \textbf{Kepadatan penduduk (\texttt{DEN}):} penduduk dibagi luas lahan
  pemukiman. Data penduduk berasal dari \emph{Basic Resident Register
  Population} milik Ministry of Internal Affairs and Communications.
  Data luas pemukiman berasal dari \emph{A System of Social and
  Demographic Statistics} milik kementerian yang sama.
\item
  \textbf{Akses pasar (\texttt{ACC}):} matriks waktu tempuh antardaerah
  berasal dari \emph{National Integrated Transport Analysis System}
  milik Ministry of Land, Infrastructure, Transport and Tourism.
\end{itemize}

\textbf{Laporan sumber.} Tabel C1 menunjukkan bahwa pada seluruh periode
1980-2014, rerata output/value-added adalah 10.341.909 juta yen untuk
semua industri, 1.964.520 juta yen untuk manufaktur, dan 8.355.152 juta
yen untuk non-manufaktur. Rerata input tenaga kerja masing-masing adalah
1.364.108, 262.363, dan 1.078.642 person-hour; rerata input modal
masing-masing adalah 18.940.610, 6.133.024, dan 12.795.671 juta yen.
Rerata nasional kepadatan penduduk adalah 1.343 dan rerata indeks akses
pasar adalah 24,16. (Hlm. 132-133, Tabel C1.)

\textbf{Laporan sumber.} Output riil meningkat sepanjang periode pada
manufaktur dan non-manufaktur, walaupun laju kenaikannya lebih moderat
pada 2000-an. Input tenaga kerja menurun secara stabil sejak 1990,
sedangkan input modal meningkat. Kepadatan penduduk naik moderat sampai
2010 lalu sedikit turun pada 2014. Indeks akses pasar meningkat tajam
sejak 1990, yang dikaitkan artikel dengan pembukaan Yamagata Shinkansen,
Akita Shinkansen, Nagano Shinkansen, dan Kyushu Shinkansen. (Hlm.
132-133, Appendix C.)

\textbf{Laporan sumber.} Rerata regional pada Tabel C2 memperlihatkan
bahwa Greater Tokyo memiliki kepadatan penduduk 4.619 dan akses pasar
35,93; Kansai 2.407 dan 32,13; Chubu 1.246 dan 25,82; sedangkan Hokkaido
257 dan 17,88. Nilai regional lain adalah Tohoku 493 dan 19,25, North
Kanto 798 dan 25,64, Hokuriku 738 dan 20,64, Chugoku 864 dan 20,33,
Shikoku 864 dan 21,08, Kyushu 846 dan 20,34, serta Okinawa 1.158 dan
24,46. (Hlm. 133, Tabel C2.)

\textbf{Inferensi pengolahan.} File tidak menjelaskan secara rinci tahun
dasar dan prosedur penerapan GDP deflator, prosedur pembersihan data,
penanganan nilai hilang atau pencilan, pembentukan bobot waktu tempuh,
maupun ketersediaan data mentah. Proksi utilisasi modal non-manufaktur
juga bergantung pada asumsi mengenai perubahan \texttt{ln(Y/K)},
sehingga pengukuran input modal pada sektor tersebut tidak sepenuhnya
bersumber dari statistik utilisasi resmi.

\subsection{Population / Sample}\label{population-sample-30}

\textbf{Laporan sumber.} Populasi analisis terdiri atas 47 prefektur
Jepang yang diamati selama 1980-2014. Unit analisis utama adalah
prefektur-tahun. Tabel 1 melaporkan 1.645 observasi untuk setiap model
dan setiap cakupan industri. Analisis juga mengagregasikan skor
efisiensi dan kontribusi pertumbuhan ke 11 wilayah Jepang untuk
keperluan penyajian deskriptif. (Hlm. 121-123, Tabel 1-2.)

\textbf{Laporan sumber.} Sebelas wilayah agregat dan prefekturnya
adalah: Hokkaido; Tohoku yang terdiri atas Aomori, Iwate, Miyagi, Akita,
Yamagata, Fukushima, dan Niigata; North Kanto yang terdiri atas Ibaraki,
Tochigi, Gunma, dan Yamanashi; Greater Tokyo area yang terdiri atas
Saitama, Chiba, Tokyo, dan Kanagawa; Chubu yang terdiri atas Nagano,
Gifu, Shizuoka, Aichi, dan Mie; Hokuriku yang terdiri atas Toyama,
Ishikawa, dan Fukui; Kansai yang terdiri atas Shiga, Kyoto, Osaka,
Hyogo, Nara, dan Wakayama; Chugoku yang terdiri atas Tottori, Shimane,
Okayama, Hiroshima, dan Yamaguchi; Shikoku yang terdiri atas Tokushima,
Kagawa, Ehime, dan Kochi; Kyushu yang terdiri atas Fukuoka, Saga,
Nagasaki, Kumamoto, Oita, Miyazaki, dan Kagoshima; serta Okinawa. (Hlm.
123 dan 124, catatan Tabel 2-3.)

\textbf{Laporan sumber.} Tabel 2 menyatakan bahwa skor regional adalah
rerata prefektur dalam setiap wilayah agregat. Tidak ada perusahaan,
rumah tangga, pekerja individual, atau unit kota yang digunakan sebagai
sampel langsung. (Hlm. 121-123, Tabel 2.)

\textbf{Inferensi pengolahan.} Perkalian 47 prefektur dengan 35 tahun,
jika kedua batas tahun dihitung inklusif, menghasilkan 1.645 dan cocok
dengan jumlah observasi Tabel 1. Ini konsisten dengan panel lengkap,
tetapi file tidak menyebutkan secara eksplisit istilah \emph{balanced
panel} atau prosedur seleksi observasi. Karena unit analisis adalah
prefektur, hasil tidak otomatis menggambarkan variasi di dalam prefektur
atau mekanisme pada tingkat perusahaan.

\subsection{Dependent Variables}\label{dependent-variables-30}

\textbf{Laporan sumber.} Artikel memiliki beberapa variabel hasil yang
tersusun dalam tahap estimasi:

\begin{itemize}
\tightlist
\item
  \textbf{Log output riil (\texttt{ln\ Y\_jt}):} variabel sisi kiri
  fungsi frontier produksi translog. Output dianalisis untuk semua
  industri, manufaktur, dan non-manufaktur.
\item
  \textbf{Efisiensi produktif (\texttt{TEF\_jt}):} skor yang diturunkan
  dari komponen inefisiensi SFA, dengan
  \texttt{TEF\_jt\ =\ exp(-u\_jt)}. Skor ini berada di atas nol dan
  paling tinggi bernilai satu.
\item
  \textbf{Laju perubahan efisiensi:} variabel sisi kiri persamaan
  konvergensi, yaitu laju perubahan tahunan rata-rata \texttt{TEF} dari
  periode awal ke periode pengamatan.
\item
  \textbf{Kontribusi determinan efisiensi:} perubahan efisiensi dipecah
  menjadi kontribusi perubahan kepadatan penduduk (\texttt{DEN}) dan
  perubahan akses pasar (\texttt{ACC}) melalui persamaan 5. Kontribusi
  tersebut kemudian dipakai dalam dekomposisi koefisien konvergensi,
  bukan sebagai variabel hasil yang diukur secara terpisah.
\end{itemize}

\textbf{Laporan sumber.} Kepadatan penduduk dan akses pasar adalah
determinan inefisiensi dalam persamaan 3. Dengan konvensi tanda artikel,
koefisien positif pada keduanya berarti inefisiensi lebih rendah dan
efisiensi lebih tinggi. (Hlm. 119-120, persamaan 3-5.)

\textbf{Inferensi pengolahan.} \emph{Borrowed size} bukan variabel
dependen yang diukur secara mandiri. Konstruk tersebut dioperasionalkan
melalui hubungan antara \texttt{ACC}, perubahan efisiensi, dan
konvergensi. Karena \texttt{ACC} sekaligus memuat populasi wilayah lain
dan waktu tempuh, model tidak memisahkan secara langsung apakah
peningkatan efisiensi berasal dari kota besar tertentu, jaringan
antarkota secara umum, atau potensi pasar yang lebih besar.

\subsection{Results}\label{results-30}

\textbf{Laporan sumber: estimasi SFA.} Uji terhadap hipotesis nol
\texttt{gamma\ =\ 0} melawan alternatif
\texttt{gamma\ \textgreater{}\ 0} menolak hipotesis nol pada taraf
signifikansi 1\%. Artikel menafsirkan hal ini sebagai bukti adanya efek
efisiensi dan sebagai dasar untuk mengevaluasi determinan efisiensi.
Hasil dengan dan tanpa efek waktu tidak menunjukkan perbedaan signifikan
menurut artikel. Model B, D, dan F dipilih berdasarkan log-likelihood.
(Hlm. 121, bagian 4; Tabel 1.)

\textbf{Laporan sumber: determinan efisiensi.} Koefisien kepadatan
penduduk dan akses pasar positif serta signifikan pada seluruh model
yang dilaporkan. Untuk model tanpa efek waktu dan dengan efek waktu,
koefisien \texttt{deltaDEN} dan \texttt{deltaACC} berturut-turut adalah
0,028 dan 0,031 serta 0,065 dan 0,048 untuk semua industri; 0,045 dan
0,036 serta 0,076 dan 0,083 untuk manufaktur; dan 0,033 dan 0,021 serta
0,053 dan 0,063 untuk non-manufaktur. Semua nilai tersebut ditandai
signifikan pada taraf 1\%. (Hlm. 121-122, Tabel 1.)

\textbf{Interpretasi penulis.} Parameter akses pasar lebih besar
daripada parameter kepadatan penduduk pada manufaktur maupun
non-manufaktur. Parameter akses pasar juga lebih besar pada manufaktur
daripada non-manufaktur. Penulis membaca pola ini sebagai bukti bahwa
dampak \emph{borrowed size} melalui jaringan antarwilayah lebih kuat
pada manufaktur, karena siklus perencanaan produk, pengembangan,
pengadaan, dan pengiriman lebih sering melintasi batas wilayah. (Hlm.
121, bagian 4.)

\textbf{Laporan sumber: komponen frontier dan skor rata-rata.} Koefisien
tenaga kerja dan modal pada fungsi produksi frontier bernilai positif
dan signifikan pada keenam spesifikasi. Untuk model dengan efek waktu
yang dipilih, koefisien tenaga kerja adalah 0,712 pada semua industri,
0,407 pada manufaktur, dan 0,419 pada non-manufaktur; koefisien modal
masing-masing 0,195, 0,495, dan 0,602. Rerata skor efisiensi pada model
yang dipilih adalah 0,898 untuk semua industri, 0,847 untuk manufaktur,
dan 0,839 untuk non-manufaktur. (Hlm. 121-122, Tabel 1.)

\textbf{Laporan sumber: skor efisiensi menurut wilayah.} Tabel 2
memberikan rerata selama 1980-2014. Untuk semua industri, manufaktur,
dan non-manufaktur secara berurutan, nilainya adalah: Hokkaido 0,809,
0,729, dan 0,755; Tohoku 0,842, 0,770, dan 0,782; North Kanto 0,911,
0,872, dan 0,854; Greater Tokyo 0,989, 0,971, dan 0,938; Chubu 0,923,
0,882, dan 0,863; Hokuriku 0,871, 0,806, dan 0,808; Kansai 0,969, 0,944,
dan 0,916; Chugoku 0,873, 0,808, dan 0,807; Shikoku 0,879, 0,819, dan
0,815; Kyushu 0,869, 0,806, dan 0,805; serta Okinawa 0,918, 0,870, dan
0,853. Rerata nasionalnya adalah 0,898, 0,847, dan 0,839. (Hlm. 123,
Tabel 2.)

\textbf{Laporan sumber.} Greater Tokyo memiliki skor efisiensi tertinggi
untuk ketiga cakupan industri, sedangkan Hokkaido memiliki skor
terendah. Skor manufaktur umumnya lebih tinggi daripada skor
non-manufaktur. Artikel mengaitkan skor tinggi manufaktur dengan rantai
pasok yang menjangkau berbagai wilayah dan manfaat pemendekan waktu
perjalanan; pada non-manufaktur, kontak antarpersonal lebih lokal dan
sebagian kegiatan lebih mudah dilakukan secara telecommuting. (Hlm.
121-123, bagian 4.)

\textbf{Laporan sumber: pertumbuhan dan kontribusi.} Tabel 3 memberikan
laju pertumbuhan skor efisiensi serta kontribusi \texttt{DEN} dan
\texttt{ACC}. Untuk semua industri, nilainya adalah: Hokkaido 0,187,
dengan kontribusi \texttt{DEN} -0,008 dan \texttt{ACC} 0,194; Tohoku
0,152, -0,012, dan 0,163; North Kanto 0,148, 0,005, dan 0,142; Greater
Tokyo 0,122, 0,024, dan 0,098; Chubu 0,152, 0,004, dan 0,148; Hokuriku
0,073, -0,001, dan 0,074; Kansai 0,113, 0,005, dan 0,108; Chugoku 0,107,
-0,006, dan 0,113; Shikoku 0,192, -0,006, dan 0,198; Kyushu 0,169,
-0,001, dan 0,170; serta Okinawa 0,098, 0,011, dan 0,087. Rerata
nasionalnya 0,138, dengan kontribusi \texttt{DEN} 0,001 dan \texttt{ACC}
0,136. (Hlm. 124, Tabel 3.)

\textbf{Laporan sumber.} Untuk manufaktur, angka pertumbuhan dan
kontribusi \texttt{DEN} serta \texttt{ACC} berturut-turut adalah:
Hokkaido 0,239, -0,010, 0,249; Tohoku 0,194, -0,015, 0,210; North Kanto
0,190, 0,007, 0,183; Greater Tokyo 0,157, 0,031, 0,125; Chubu 0,195,
0,005, 0,189; Hokuriku 0,094, -0,001, 0,096; Kansai 0,145, 0,007, 0,138;
Chugoku 0,137, -0,008, 0,145; Shikoku 0,246, -0,008, 0,254; Kyushu
0,217, -0,001, 0,218; dan Okinawa 0,126, 0,014, 0,112. Rerata
nasionalnya 0,176, 0,002, dan 0,174. (Hlm. 124, Tabel 3.)

\textbf{Laporan sumber.} Untuk non-manufaktur, angka pertumbuhan dan
kontribusi \texttt{DEN} serta \texttt{ACC} berturut-turut adalah:
Hokkaido 0,182, -0,006, 0,187; Tohoku 0,149, -0,009, 0,158; North Kanto
0,142, 0,004, 0,138; Greater Tokyo 0,113, 0,018, 0,094; Chubu 0,146,
0,003, 0,143; Hokuriku 0,071, -0,001, 0,072; Kansai 0,108, 0,004, 0,104;
Chugoku 0,105, -0,005, 0,109; Shikoku 0,186, -0,005, 0,191; Kyushu
0,164, -0,001, 0,164; dan Okinawa 0,093, 0,008, 0,084. Rerata
nasionalnya 0,132, 0,001, dan 0,131. (Hlm. 124, Tabel 3.)

\textbf{Interpretasi penulis.} Hokkaido, Shikoku, dan Kyushu dilaporkan
mengalami peningkatan efisiensi yang signifikan dan melebihi rerata
nasional. Peningkatan akses pasar dianggap sangat berpengaruh pada
ketiga wilayah tersebut: penerbangan memperbaiki akses Hokkaido ke
daratan utama, Jembatan Honshu-Shikoku memperbaiki akses Shikoku, dan
Kyushu Shinkansen yang dibuka pada 2004 dan 2011 memperbaiki akses
Kyushu. Tohoku, North Kanto, dan Chubu juga disebut terdampak kuat oleh
perbaikan akses. (Hlm. 123-124, bagian 4.)

\textbf{Laporan sumber: hubungan dinamis.} Gambar 1 menunjukkan hubungan
menanjak antara peningkatan akses pasar dan peningkatan efisiensi untuk
semua industri. Nagano dan Tokushima berada di bagian kanan atas gambar,
yang dalam narasi artikel berarti keduanya mengalami perbaikan akses dan
efisiensi yang tinggi. Gambar 2 menunjukkan kecenderungan menurun antara
tingkat efisiensi awal dan laju perubahannya: wilayah dengan efisiensi
awal lebih rendah cenderung mengalami pertumbuhan efisiensi lebih
tinggi. Nagano, Tokushima, dan Hokkaido disebut membaik secara
signifikan, sedangkan Saitama, Osaka, Kanagawa, dan Tokyo hanya
mengalami sedikit perbaikan. (Hlm. 124-125, Gambar 1-2.)

\textbf{Laporan sumber: konvergensi.} Koefisien konvergensi untuk semua
industri adalah -0,0029 dengan standard error 0,0010; kontribusi
kepadatan penduduk 0,0012 dengan standard error 0,0002; dan kontribusi
akses pasar -0,0042 dengan standard error 0,0010. Untuk manufaktur,
ketiganya adalah -0,0025 (0,0009), 0,0011 (0,0001), dan -0,0036
(0,0008). Untuk non-manufaktur, ketiganya adalah -0,0030 (0,0009),
0,0009 (0,0001), dan -0,0039 (0,0009). Semua koefisien pada Tabel 4
ditandai signifikan pada taraf 1\%. (Hlm. 126, Tabel 4.)

\textbf{Interpretasi penulis.} Koefisien konvergensi yang negatif
menunjukkan disparitas efisiensi produktif berkurang. Kontribusi
kepadatan penduduk yang positif menunjukkan aglomerasi penduduk
memperlebar disparitas, sedangkan kontribusi akses pasar yang negatif
menunjukkan jaringan menguranginya. Besarnya kontribusi akses pasar
melebihi kontribusi kepadatan penduduk, sehingga efek bersihnya adalah
penurunan disparitas. (Hlm. 124-126, bagian 4.)

\textbf{Inferensi pengolahan.} Kesimpulan bahwa \emph{agglomeration
shadows} tidak tampak berarti tidak ditemukannya efek negatif bersih
akses pasar pada spesifikasi dan periode penelitian, bukan pembuktian
bahwa tidak ada wilayah atau sektor yang mengalami penarikan kegiatan
oleh aglomerasi. Model mengestimasi efek rata-rata pada 47 prefektur dan
tidak melaporkan pengujian terpisah atas lokasi bayangan, arus kegiatan,
atau heterogenitas yang dapat membuat efek negatif tertutup oleh efek
positif di wilayah lain.

\subsection{Findings}\label{findings-30}

\textbf{Temuan yang dilaporkan sumber.} Temuan utama artikel dapat
dirangkum sebagai berikut:

\begin{itemize}
\tightlist
\item
  Kepadatan penduduk dan akses pasar sama-sama berhubungan dengan
  efisiensi produktif yang lebih tinggi dalam model SFA.
\item
  Efek akses pasar lebih besar daripada efek kepadatan penduduk, baik
  pada manufaktur maupun non-manufaktur.
\item
  Efek akses pasar lebih besar pada manufaktur daripada non-manufaktur.
\item
  Greater Tokyo, Kansai, dan Chubu memiliki tingkat efisiensi rata-rata
  tinggi, sedangkan Hokkaido terendah; wilayah lokal tertentu mengalami
  pertumbuhan efisiensi yang kuat selama periode pengamatan.
\item
  Hubungan antara peningkatan akses pasar dan peningkatan efisiensi
  bersifat menanjak, sedangkan hubungan antara efisiensi awal dan
  pertumbuhan efisiensi bersifat menurun.
\item
  Disparitas efisiensi produktif menurun pada semua industri,
  manufaktur, dan non-manufaktur menurut koefisien konvergensi yang
  negatif.
\item
  Kontribusi kepadatan penduduk memperlebar disparitas, sementara
  kontribusi akses pasar menguranginya dengan magnitudo yang lebih
  besar. (Hlm. 121-126, Tabel 1-4 dan Gambar 1-2.)
\end{itemize}

\textbf{Interpretasi penulis.} Penulis menyimpulkan bahwa jaringan
transportasi memperluas manfaat aglomerasi kota besar ke wilayah lokal
melalui \emph{borrowed size}. Manufaktur lebih diuntungkan karena
hubungan produksinya lebih lintas wilayah. Kota metropolitan besar tetap
memperoleh manfaat aglomerasi berbasis ukuran, tetapi wilayah lokal
dapat mengejar melalui pemendekan waktu tempuh. (Hlm. 126-127, bagian
5.)

\textbf{Inferensi pengolahan.} Pola yang paling kuat adalah perbedaan
antara dua mekanisme agregat: aglomerasi penduduk lokal meningkatkan
efisiensi tetapi cenderung memperlebar ketimpangan, sedangkan akses
pasar berkaitan dengan peningkatan efisiensi sekaligus konvergensi.
Namun, \texttt{ACC} tidak hanya merepresentasikan koneksi transportasi;
ia juga memasukkan populasi wilayah lain. Dengan demikian, hasil tidak
dapat menunjukkan bagian mana dari efek yang berasal dari perjalanan,
ukuran pasar, atau karakteristik wilayah yang berkorelasi dengan
keduanya.

\subsection{Conclusions}\label{conclusions-30}

\textbf{Kesimpulan yang dinyatakan penulis.} Bagian diskusi dan
kesimpulan merumuskan lima kesimpulan utama:

\begin{itemize}
\tightlist
\item
  Perbaikan jaringan antarwilayah meningkatkan efek peminjaman, yaitu
  manfaat aglomerasi kota besar yang melimpah ke ekonomi lokal dan
  meningkatkan efisiensi produktif.
\item
  Dampak \emph{borrowed size} melalui jaringan lebih besar daripada
  manfaat aglomerasi berbasis ukuran lokal. Efek positif jaringan tampak
  kuat, sehingga efek \emph{agglomeration shadow} tidak terlihat dalam
  hasil agregat.
\item
  Peningkatan efisiensi akibat jaringan lebih besar pada manufaktur
  daripada non-manufaktur karena manufaktur berperan sebagai industri
  dasar atau ekspor dan bergantung pada jaringan eksternal.
\item
  Greater Tokyo dan Kansai memperoleh manfaat besar dari ukuran wilayah
  dan aglomerasi penduduk, sedangkan wilayah lokal di luar metropolitan
  besar sangat efektif menggunakan \emph{borrowed size} dari aglomerasi
  kota besar.
\item
  Penguatan jaringan transportasi membantu wilayah lokal mengejar
  wilayah metropolitan besar dan berkontribusi pada konvergensi
  disparitas efisiensi produktif. (Hlm. 126-127, bagian 5.)
\end{itemize}

\textbf{Interpretasi penulis.} Temuan tersebut dipakai untuk menyatakan
bahwa basis geografis teori aglomerasi perlu dibangun ulang agar tidak
berhenti pada batas aglomerasi. Jika manfaat ekonomi dapat dibagi
melalui jaringan regional, istilah dan kebijakan yang hanya berfokus
pada aglomerasi internal dianggap terlalu sempit. (Hlm. 126-127, bagian
5.)

\textbf{Batas interpretasi.} Kesimpulan berlaku pada 47 prefektur
Jepang, periode 1980-2014, definisi output-input yang digunakan, dan
indeks akses pasar yang dibangun penulis. File tidak menyediakan dasar
untuk menggeneralisasikan hasil ke tingkat perusahaan, kota di negara
lain, periode setelah 2014, atau proyek transportasi tertentu tanpa
pengujian tambahan.

\subsection{Contributions}\label{contributions-30}

\textbf{Kontribusi yang diklaim penulis.} Artikel menyatakan beberapa
kontribusi:

\begin{itemize}
\tightlist
\item
  menghubungkan teori aglomerasi, \emph{borrowed size}, dan jaringan
  antarwilayah dengan analisis efisiensi produktif ekonomi regional
  Jepang;
\item
  menguji dampak jaringan terhadap disparitas efisiensi produktif, yang
  menurut artikel belum diperiksa dalam studi sebelumnya untuk Jepang;
\item
  menggunakan indeks akses pasar yang memperhitungkan waktu tempuh,
  bukan hanya jarak fisik, sehingga jaringan transportasi udara, kereta
  api, dan mobil masuk ke dalam ukuran potensi pasar;
\item
  membandingkan efek jaringan pada manufaktur dan non-manufaktur;
\item
  menguraikan perubahan efisiensi menjadi kontribusi kepadatan penduduk
  dan akses pasar, lalu menghubungkannya dengan konvergensi disparitas;
  dan
\item
  memberi dasar empiris untuk mendiskusikan kebijakan
  \emph{super-megaregion} dan investasi transportasi berkecepatan tinggi
  dalam rencana tata ruang nasional Jepang. (Hlm. 116-117, 119-121, dan
  126-127.)
\end{itemize}

\textbf{Inferensi pengolahan.} Nilai tambah analitis terkuat artikel
adalah memperlakukan jaringan bukan sekadar kondisi akses, melainkan
sebagai komponen yang dapat mengubah distribusi manfaat aglomerasi
antardaerah. Namun, kontribusi tersebut bergantung pada validitas indeks
\texttt{ACC}, asumsi SFA, dan interpretasi efisiensi sebagai hasil yang
dapat dipengaruhi oleh jaringan. Artikel belum mengukur secara langsung
proses mikro yang membuat akses pasar menjadi produktivitas atau
\emph{borrowed size}.

\subsection{Research Gap}\label{research-gap-30}

\textbf{Kesenjangan yang dinyatakan sumber sebelum penelitian.} Artikel
menyebutkan kesenjangan berikut:

\begin{itemize}
\tightlist
\item
  teori aglomerasi tradisional belum menjelaskan mengapa disparitas
  regional Jepang menyempit ketika konsentrasi ekonomi Tokyo meningkat;
\item
  penelitian tentang perluasan manfaat aglomerasi melalui jaringan
  transportasi untuk Jepang masih tidak memadai;
\item
  studi sebelumnya cenderung berfokus pada aktivitas ekonomi tertentu,
  bukan seluruh ekonomi nasional;
\item
  pengaruh jaringan antarwilayah terhadap konvergensi disparitas
  efisiensi produktif belum diperiksa secara memadai;
\item
  perbedaan antara manufaktur dan non-manufaktur dalam dampak jaringan
  belum dijelaskan secara cukup; dan
\item
  belum ada bukti empiris yang mendukung secara langsung kebijakan
  pembentukan \emph{super-megaregion} melalui penguatan jaringan
  antarwilayah. (Hlm. 115-117, bagian 1.)
\end{itemize}

\textbf{Kesenjangan yang masih tersisa menurut artikel.} Penulis
menyatakan bahwa mekanisme \emph{borrowed size} perlu diteliti dari
perspektif mikro. Mereka secara khusus menyebut perlunya memahami apakah
penurunan waktu perjalanan meningkatkan jumlah perjalanan bisnis dan
bagaimana perjalanan bisnis tersebut meningkatkan efisiensi produktif.
Artikel juga menyatakan bahwa jaringan antarwilayah internasional dan
hubungan antardaerah lintas negara memerlukan penelitian lebih lanjut.
(Hlm. 126-127, bagian 5.)

\textbf{Inferensi pengolahan.} Studi ini belum menutup kesenjangan
antara jaringan yang diukur sebagai potensi akses dan peminjaman manfaat
yang berlangsung secara nyata. Belum ada pengukuran arus perjalanan
bisnis, asal-tujuan transaksi, jaringan perusahaan, aliran pengetahuan,
atau perubahan produktivitas perusahaan yang dapat menjelaskan mekanisme
di antara waktu tempuh dan efisiensi regional.

\subsection{Limitations}\label{limitations-30}

\textbf{Keterbatasan yang dinyatakan atau diakui dalam file.} Artikel
tidak memiliki bagian khusus berjudul limitations. Namun, agenda
penelitian pada bagian akhir secara eksplisit mengakui bahwa mekanisme
efek \emph{borrowed size} belum dielaborasi dari perspektif mikro,
dampak perjalanan bisnis terhadap efisiensi belum diteliti secara
memadai, dan jaringan antardaerah internasional belum dianalisis. (Hlm.
126-127, bagian 5.)

\textbf{Keterbatasan pengukuran sebagai inferensi pengolahan.} Efisiensi
produktif diturunkan dari komponen inefisiensi model frontier, bukan
diamati langsung. \texttt{ACC} menggabungkan populasi wilayah lain dan
waktu tempuh sehingga tidak memisahkan ukuran pasar, konektivitas, dan
kedekatan ke aglomerasi kota besar. Kepadatan penduduk juga merupakan
proksi aglomerasi yang tidak menangkap struktur industri, kualitas
tenaga kerja, inovasi, atau kepadatan perusahaan secara terpisah.

\textbf{Keterbatasan data sebagai inferensi pengolahan.} Proksi
utilisasi modal non-manufaktur dibentuk dari residual karena statistik
resmi tidak tersedia. File tidak memberikan rincian validasi proksi
tersebut. File juga tidak melaporkan prosedur penanganan data hilang dan
pencilan, dokumentasi lengkap matriks waktu tempuh, bobot setiap moda,
aturan agregasi waktu perjalanan, atau akses terhadap data replikasi.

\textbf{Keterbatasan desain sebagai inferensi pengolahan.} Panel 47
prefektur memungkinkan analisis perubahan sepanjang waktu dan efek tetap
prefektur, tetapi file tidak menyebutkan strategi untuk mengatasi
kausalitas terbalik, variabel terlewat, atau seleksi lokasi investasi.
Wilayah yang sudah produktif dapat lebih mudah menarik infrastruktur,
sehingga hubungan positif antara akses pasar dan efisiensi tidak secara
otomatis menunjukkan arah sebab-akibat.

\textbf{Keterbatasan cakupan sebagai inferensi pengolahan.} Observasi
berhenti pada 2014, unit utama adalah prefektur, dan kategori industri
hanya dibagi menjadi semua industri, manufaktur, dan non-manufaktur.
Hasil tidak menjelaskan variasi dalam prefektur, perbedaan
antarperusahaan, mekanisme pada sektor yang lebih rinci, atau perubahan
setelah jaringan dan kebijakan transportasi berikutnya.

\textbf{Informasi metodologis yang tidak tersedia dalam file.} Tidak
disebutkan analisis sensitivitas terhadap ukuran akses pasar alternatif,
uji placebo, uji ketahanan terhadap spesifikasi frontier lain, uji
dependensi spasial, instrumen, desain kuasi-eksperimental, kode
analisis, atau data replikasi. File juga tidak menyajikan penjelasan
lebih lanjut tentang asumsi distribusi inefisiensi selain formula yang
digunakan.

\subsection{Challenges}\label{challenges-30}

\textbf{Tantangan yang dilaporkan sumber.} Artikel menghadapi tantangan
menjelaskan penyempitan disparitas dalam situasi ketika Tokyo semakin
terkonsentrasi. Tantangan empiris lainnya adalah mengukur jaringan
antarwilayah dengan waktu tempuh yang berubah, menggabungkan udara,
kereta api, dan mobil, serta membedakan manfaat akses dari kompetisi
yang dapat menghasilkan \emph{agglomeration shadow}. (Hlm. 115-120,
bagian 1-3.)

\textbf{Tantangan data yang dilaporkan sumber.} Statistik resmi untuk
rasio utilisasi kapasitas non-manufaktur tidak tersedia, sehingga
penulis membangun proksi berbasis residual. Wilayah kepulauan atau
wilayah yang tidak berada di daratan utama juga memiliki kondisi akses
yang khas; artikel secara khusus membahas perbaikan penerbangan Hokkaido
dan akses Okinawa yang dipengaruhi jumlah penerbangan ke Tokyo. (Hlm.
123 dan 131-133, bagian 4 dan Appendix C.)

\textbf{Tantangan substantif yang dilaporkan sumber.} Perbedaan struktur
produksi membuat efek jaringan mungkin tidak seragam. Manufaktur
memiliki rantai pasok dan transaksi eksternal yang lebih luas, sedangkan
non-manufaktur lebih berpusat pada layanan lokal dan kontak
antarpersonal. Karena itu, satu kebijakan jaringan dapat menghasilkan
besaran manfaat yang berbeda antarindustri. (Hlm. 116-117 dan 123,
bagian 1 dan 4.)

\textbf{Inferensi pengolahan.} Tantangan utama dalam membaca hasil
adalah memisahkan tiga hal yang dapat menghasilkan korelasi positif yang
sama: aglomerasi lokal, akses ke pasar besar, dan investasi
infrastruktur yang dipilih oleh wilayah berpotensi tinggi. Artikel
berhasil mendekomposisi kepadatan dan akses dalam kerangka model, tetapi
belum membuktikan bahwa dekomposisi tersebut identik dengan arus
peminjaman manfaat secara aktual.

\subsection{Practical Implications}\label{practical-implications-30}

\textbf{Implikasi yang dinyatakan penulis.} Hasil mendukung strategi
nasional Jepang yang menggabungkan pemadatan atau \emph{compactization}
wilayah dengan pembentukan jaringan antarwilayah untuk mempertahankan
ukuran pasar ketika populasi menurun. Rencana tata ruang nasional
bertujuan membentuk negara yang mendorong interaksi melalui pergerakan
orang, barang, uang, dan informasi, serta menggunakan jaringan
transportasi berkecepatan tinggi untuk menciptakan
\emph{super-megaregion}. (Hlm. 126-127 dan 131-132, bagian 5 dan
Appendix B.)

\textbf{Rekomendasi penulis.} Investasi yang mempersingkat waktu tempuh
antara wilayah metropolitan besar dan wilayah lokal dipandang dapat
memperbesar efek \emph{borrowed size}, meningkatkan efisiensi lokal, dan
mengurangi disparitas. Linear Chuo Shinkansen disebut sebagai contoh
kebijakan yang dirancang menghubungkan Tokyo, Nagoya, dan Osaka serta
memperkuat interaksi ekonomi. (Hlm. 126-127 dan 131-132, bagian 5 dan
Appendix B.)

\textbf{Kualifikasi inferensial.} Bukti artikel mendukung prioritas pada
akses jaringan sebagai strategi pemerataan dalam batas model, terutama
untuk manufaktur. Namun, rekomendasi tersebut tidak boleh dibaca sebagai
evaluasi dampak proyek Linear Chuo Shinkansen tertentu, karena artikel
tidak mengestimasi proyek itu sebagai intervensi dan observasi berakhir
pada 2014. Pemeriksaan manfaat dan kemungkinan efek kompetitif di
tingkat wilayah atau sektor tetap diperlukan.

\subsection{Applications}\label{applications-30}

\textbf{Aplikasi yang digunakan dalam sumber.} Kerangka artikel
diterapkan untuk:

\begin{itemize}
\tightlist
\item
  mengestimasi skor efisiensi produktif prefektur Jepang untuk semua
  industri, manufaktur, dan non-manufaktur;
\item
  membandingkan tingkat efisiensi antara 11 wilayah agregat Jepang;
\item
  menguraikan perubahan efisiensi menjadi kontribusi kepadatan penduduk
  dan akses pasar;
\item
  menguji konvergensi atau divergensi disparitas efisiensi produktif;
  dan
\item
  membaca kemungkinan relevansi jaringan transportasi bagi rencana
  pembangunan wilayah nasional. (Hlm. 121-127, Tabel 1-4, Gambar 1-2,
  dan bagian 5.)
\end{itemize}

\textbf{Aplikasi potensial sebagai inferensi pengolahan.} Dengan data
yang sebanding, metode tersebut dapat digunakan sebagai alat diagnostik
untuk membandingkan efisiensi antardaerah, memeriksa apakah peningkatan
akses pasar berkaitan dengan pengejaran wilayah ber-efisiensi rendah,
dan membedakan kontribusi aglomerasi lokal dari kontribusi akses
jaringan. Ini adalah kemungkinan penggunaan analitis, bukan aplikasi
kebijakan yang diuji dalam file.

\textbf{Batas aplikasi.} File menyebut perangkat lunak Frontier 4.1,
tetapi tidak melaporkan kode, paket replikasi, atau implementasi
operasional bagi pemerintah. Tidak disebutkan adanya evaluasi
pascapembangunan, uji proyek tertentu, atau studi implementasi
kebijakan.

\subsection{Future Research
(Optional)}\label{future-research-optional-30}

\textbf{Agenda yang dinyatakan penulis.} Penulis mengusulkan:

\begin{itemize}
\tightlist
\item
  memperdalam temuan dari perspektif mikro;
\item
  menjelaskan mekanisme efek \emph{borrowed size};
\item
  menguji apakah penurunan waktu perjalanan meningkatkan jumlah
  perjalanan bisnis;
\item
  menjelaskan bagaimana perjalanan bisnis dapat meningkatkan efisiensi
  produktif;
\item
  meneliti pemanfaatan konektivitas jaringan antarnegara seiring
  globalisasi; dan
\item
  menganalisis serta memahami jaringan antardaerah internasional. (Hlm.
  126-127, bagian 5.)
\end{itemize}

\textbf{Inferensi pengolahan.} Berdasarkan keterbatasan yang tampak
dalam TXT, penelitian lanjutan juga dapat menguji ukuran akses pasar
alternatif, mengikuti arus asal-tujuan dan jaringan perusahaan, memakai
data tingkat perusahaan atau sektor, memeriksa efek heterogen
antarwilayah, serta menggunakan desain identifikasi yang lebih kuat
untuk membedakan pengaruh jaringan dari pemilihan lokasi investasi.
Usulan ini adalah inferensi review, bukan agenda yang dinyatakan dengan
rumusan yang sama oleh penulis.

\subsection{Provenance Notes}\label{provenance-notes-30}

\begin{itemize}
\tightlist
\item
  Review ini hanya menggunakan file
  \texttt{source/journal/A15-otsuka-2020-inter-regional-networks-and-productive-efficiency-in-japan-10.1111-pirs.12474.txt}.
  Tidak ada PDF, situs web, catatan Zotero, atau referensi yang dikutip
  artikel yang dibuka untuk menambah atau memverifikasi substansi.
\item
  Seluruh 994 baris TXT yang tersedia dibaca, termasuk blok
  \texttt{{[}PDF\ Metadata{]}}, teks utama, persamaan, Tabel 1-4, Gambar
  1-2 beserta keterangan tekstual, Appendix A-C, Tabel C1-C2, daftar
  pustaka, dan abstrak berbahasa Spanyol serta Jepang.
\item
  Metadata TXT menyatakan sumber berasal dari PDF dengan nama yang sama,
  diekstraksi pada 31 Agustus 2026 menggunakan
  \texttt{pdftotext\ -layout}. Metadata PDF menyebut 20 halaman, tidak
  terenkripsi, dan tidak bertag. Informasi tersebut diperlakukan sebagai
  metadata file, bukan sebagai bukti tambahan dari luar TXT.
\item
  Judul pada metadata PDF mengalami kerusakan karakter
  (\texttt{Inter�regional}), sedangkan judul pada tubuh artikel
  dan baris sitasi terbaca sebagai \emph{Inter-regional networks and
  productive efficiency in Japan}. Identitas judul mengikuti tubuh
  artikel dan sitasi yang tersedia dalam TXT.
\item
  Nomor halaman dalam review merujuk pada nomor halaman cetak jurnal,
  terutama hlm. 115-133, bukan nomor baris TXT. Locator bagian,
  persamaan, tabel, dan gambar dipertahankan ketika tersedia.
\item
  Tabel 1-4 dan Tabel C1-C2 terbaca sebagai teks dan menjadi dasar angka
  yang dilaporkan. Nilai Gambar 1-2 tidak diperkirakan dari visual;
  review hanya menggunakan kecenderungan dan contoh lokasi yang
  dijelaskan dalam narasi atau keterangan gambar.
\item
  Label \textbf{Laporan sumber} atau \textbf{Temuan yang dilaporkan
  sumber} digunakan untuk parafrasa atas hasil, metode, data, dan
  pernyataan eksplisit artikel. Label \textbf{Interpretasi penulis}
  digunakan untuk makna yang diberikan penulis terhadap koefisien atau
  pola. Label \textbf{Inferensi pengolahan} digunakan untuk evaluasi
  yang diturunkan dari isi TXT dan tidak diposisikan sebagai klaim
  penulis.
\item
  Nomor terbitan jurnal, metode pengambilan atau seleksi sampel, rincian
  pembersihan data, prosedur data hilang dan pencilan, bobot lengkap
  waktu tempuh, kode analisis, data replikasi, uji ketahanan, instrumen,
  dan desain kuasi-eksperimental tidak disebutkan dalam file.
\item
  Status penelaahan sejawat tidak dipastikan secara formal. Ucapan
  terima kasih kepada reviewers dicatat sebagai isi sumber, bukan
  sebagai dasar untuk mengklaim status publikasi tertentu.
\end{itemize}

\section{Review A16 - Otsuka (2024)}\label{review-a16---otsuka-2024}

Status: Review literatur berbasis seluruh teks TXT hasil ekstraksi PDF.
Tidak ada verifikasi dengan sumber di luar TXT.

\subsection{Identitas Artikel}\label{identitas-artikel-5}

\textbf{Laporan sumber.}

\begin{itemize}
\tightlist
\item
  \textbf{Kode:} A16, berdasarkan prefiks nama file sumber.
\item
  \textbf{Judul:} \emph{Effects of regional network economies on
  industrial productivity in Japan: dynamic total factor productivity
  function approach}.
\item
  \textbf{Penulis:} Akihiro Otsuka.
\item
  \textbf{Jurnal:} \emph{Asia-Pacific Journal of Regional Science},
  2024, volume 8, hlm. 1111-1134.
\item
  \textbf{DOI:} 10.1007/s41685-024-00358-2.
\item
  \textbf{Afiliasi yang dicantumkan:} Association of International Arts
  and Science, Yokohama City University.
\item
  \textbf{Riwayat naskah:} diterima 16 Juli 2024, diterima untuk
  publikasi 30 Oktober 2024, dan dipublikasikan secara daring 13
  November 2024.
\item
  \textbf{Kata kunci:} inter-regional network, borrowed size, total
  factor productivity, panel data, dan Japan.
\item
  \textbf{Klasifikasi JEL:} R10 dan R11.
\end{itemize}

Locator: TXT baris 1-28 dan 30-79; PDF hlm. 1111.

\textbf{Informasi tidak tersedia dalam TXT.} Nomor isu tidak disebutkan
dalam file. TXT menyebut naskah sebagai \texttt{ARTICLE} dan memuat
riwayat penerimaan, tetapi status peer-review tidak diverifikasi secara
terpisah dalam review ini.

\subsection{Summarized Abstract}\label{summarized-abstract-31}

\textbf{Laporan sumber.} Artikel berangkat dari pentingnya peningkatan
kinerja produktivitas bagi keberlanjutan ekonomi regional. Penulis
menyatakan bahwa agen ekonomi regional memperoleh manfaat dari interaksi
lokal di dalam wilayah sekaligus dari ekonomi eksternal di luar wilayah
melalui masyarakat yang terhubung dalam jaringan. Pertanyaan utamanya
ialah apakah ekonomi eksternal yang terlokalisasi atau limpahan
eksternalitas melalui jaringan antarwilayah lebih penting bagi
produktivitas. Studi ini menggunakan data ekonomi regional Jepang dari
perspektif industri dan pendekatan fungsi total factor productivity
(TFP) dinamis. Studi juga menilai dampak peningkatan infrastruktur
transportasi berkualitas tinggi terhadap pengejaran ketertinggalan
produktivitas. Hasil abstrak menyatakan bahwa jaringan antarwilayah
mempunyai efek jangka panjang terhadap kinerja produktivitas industri;
efek borrowed size berkontribusi signifikan pada perbaikan produktivitas
industri manufaktur dan non-manufaktur; serta penguatan infrastruktur
transportasi dan jaringan transportasi antarwilayah diperkirakan dapat
memperbaiki produktivitas industri.

Locator: TXT baris 48-68; PDF hlm. 1111, bagian Abstract.

\textbf{Interpretasi penulis.} Hasil tersebut dibaca sebagai dukungan
bahwa eksternalitas jaringan bukan hanya fenomena umum antardaerah,
tetapi juga relevan bagi produktivitas industri dan dapat memperluas
cakupan geografis aktivitas produksi melalui transportasi berkualitas
tinggi.

\textbf{Inferensi pengolahan.} Abstrak secara eksplisit menyebut efek
jangka panjang; keterangan tentang efek jangka pendek dan jangka panjang
secara bersamaan dijelaskan dalam pendahuluan serta metode. Kata-kata
seperti ``likely to improve'' dalam abstrak adalah implikasi kebijakan
penulis, bukan hasil eksperimen atau estimasi efek kausal kebijakan
transportasi.

\subsection{Problem Statement}\label{problem-statement-31}

\textbf{Laporan sumber.} Penelitian terdahulu tentang ekonomi aglomerasi
umumnya menekankan bahwa kedekatan geografis, pengelompokan agen
ekonomi, dan limpahan pengetahuan lokal memengaruhi produktivitas.
Artikel mencatat bahwa bentuk limpahan pengetahuan yang tidak sepenuhnya
terikat geografis juga mungkin terjadi melalui kedekatan kognitif,
teknologi, dan sosial. Karena itu, wilayah dapat memperoleh kapabilitas
komplementer dari lokasi yang jauh melalui kolaborasi dan jaringan
antarwilayah. Masalah yang hendak dijawab adalah bagaimana eksternalitas
geografis tersebut menentukan kinerja produktivitas ekonomi regional,
khususnya pada aktivitas industri.

Locator: TXT baris 86-130; PDF hlm. 1112-1113, bagian 1 Introduction.

\textbf{Laporan sumber.} Penulis menyatakan bahwa bukti empiris
sebelumnya belum cukup untuk aktivitas industri regional. Industri
manufaktur diposisikan sebagai industri berorientasi ekspor yang
cenderung membangun sistem produksi dan rantai pasok dalam jaringan
wilayah luas. Sebaliknya, industri non-manufaktur, terutama jasa,
diposisikan sebagai industri intra-regional yang lebih dibatasi oleh
kebutuhan permintaan lokal. Perbedaan ini membuat efek jaringan
antarwilayah dan eksternalitas lokal mungkin berbeda antarindustri.
Artikel juga menempatkan isu ini dalam konteks kebijakan Jepang pada
2000-an yang mendorong lokasi industri, klaster, dan hubungan antara
perusahaan, institusi akademik, serta organisasi riset lintas batas
administratif.

Locator: TXT baris 134-176; PDF hlm. 1113-1114, bagian 1 Introduction.

\textbf{Laporan sumber.} Kerangka masalah diperjelas melalui dua
kemungkinan eksternalitas. Borrowed size memungkinkan wilayah yang lebih
kecil memperoleh manfaat skala dari pusat atau wilayah lain.
Agglomeration shadow sebaliknya dapat terjadi ketika pusat besar
menyerap sumber daya dari wilayah tetangga. Artikel menyatakan bahwa
kedua proses dapat bersifat komplementer atau kompetitif, tetapi studi
ini tidak memodelkan agglomeration shadow secara eksplisit. Penulis
memilih mengestimasi efek bersih eksternalitas geografis dari perspektif
makroregional dengan fokus statistik pada borrowed size.

Locator: TXT baris 259-335; PDF hlm. 1115-1117, bagian 2.2-2.3.

\textbf{Inferensi pengolahan.} Problem statement artikel memiliki dua
lapisan: mengukur apakah jaringan antarwilayah berkaitan dengan
produktivitas, dan menentukan apakah pola tersebut berbeda antara
manufaktur dan non-manufaktur serta berhubungan dengan pengejaran
ketertinggalan produktivitas. Lapisan kedua memerlukan interpretasi
mekanisme yang lebih kuat daripada koefisien borrowed size saja karena
agglomeration shadow tidak diestimasi sebagai variabel terpisah.

\subsection{Objectives}\label{objectives-31}

\textbf{Laporan sumber.} Tujuan penelitian yang dapat dirangkum dari
abstrak, pendahuluan, dan bagian fokus penelitian adalah:

\begin{itemize}
\tightlist
\item
  mengidentifikasi efek eksternalitas geografis terhadap kinerja
  produktivitas ekonomi regional;
\item
  mengklarifikasi efek pembentukan jaringan antarwilayah dibandingkan
  dengan eksternalitas geografis konvensional;
\item
  mengevaluasi efek borrowed size terhadap peningkatan produktivitas
  dari perspektif industri;
\item
  membandingkan efek tersebut antara seluruh industri, industri
  manufaktur, dan industri non-manufaktur;
\item
  menangkap efek jangka pendek dan jangka panjang melalui fungsi TFP
  dinamis; dan
\item
  menilai apakah pembentukan jaringan antarwilayah serta transportasi
  berkualitas tinggi membantu wilayah dengan produktivitas rendah
  mengejar wilayah dengan produktivitas tinggi.
\end{itemize}

Locator: TXT baris 150-176 dan 337-356; PDF hlm. 1113-1114 dan 1117,
bagian 1 dan 2.4.

\textbf{Inferensi pengolahan.} Artikel tidak menyajikan subbagian
hipotesis formal atau daftar tujuan bernomor. Daftar di atas adalah
perumusan ulang dari fokus penelitian yang dinyatakan penulis, bukan
tambahan tujuan baru.

\subsection{Literature Survey}\label{literature-survey-31}

\textbf{Laporan sumber.} Survei literatur dimulai dari konsep
produktivitas dan aglomerasi. Produktivitas didefinisikan sebagai rasio
output terhadap input. Penulis membedakan produktivitas tenaga kerja
dari TFP karena produktivitas tenaga kerja hanya menggunakan tenaga
kerja sebagai input dan dapat terpengaruh oleh perbedaan rasio
modal-tenaga kerja antardaerah. TFP diposisikan sebagai ukuran yang
mencakup produktivitas tenaga kerja dan modal. Tabel 1 membagi kemajuan
teknologi dalam TFP menjadi kemajuan langsung, seperti investasi R\&D,
stok pengetahuan, dan limpahan pengetahuan, serta kemajuan tidak
langsung, seperti skala ekonomi, kematangan modal, modal manusia,
infrastruktur, dan eksternalitas geografis berupa aglomerasi serta
ekonomi jaringan.

Locator: TXT baris 188-236; PDF hlm. 1114-1115, bagian 2.1 dan Tabel 1.

\textbf{Laporan sumber.} Literatur aglomerasi yang dirujuk penulis
menyatakan bahwa inovasi dan limpahan pengetahuan cenderung mengelompok
serta terlokalisasi. Tacit knowledge dianggap sulit ditransfer dalam
jarak jauh. Namun, artikel juga merujuk penelitian yang menekankan
kedekatan kognitif, teknologi, dan sosial serta pentingnya menggabungkan
pengetahuan lokal dengan pengetahuan dari lokasi yang jauh. Dari sini,
kedekatan geografis diperlakukan sebagai bukan kondisi yang selalu perlu
atau cukup untuk peningkatan produktivitas.

Locator: TXT baris 97-116; PDF hlm. 1112, bagian 1 Introduction.

\textbf{Laporan sumber.} Borrowed size ditelusuri kepada Alonso dan
dijelaskan sebagai kemampuan kota atau wilayah yang lebih kecil untuk
memanfaatkan ekonomi aglomerasi dari aglomerasi yang lebih besar tanpa
menanggung seluruh kerugiannya, seperti kemacetan. Literatur jaringan
kota yang dirujuk menyatakan bahwa jaringan antarwilayah memungkinkan
wilayah memperoleh kapabilitas komplementer, memanfaatkan skala dari
wilayah lain, serta meningkatkan produktivitas meskipun ukuran lokal
terbatas. Artikel juga menjelaskan bahwa transportasi yang lebih baik
dapat menurunkan waktu perjalanan dan biaya interaksi, memperluas
jangkauan limpahan pengetahuan, dan memperkuat eksternalitas jaringan.

Locator: TXT baris 117-130 dan 259-291; PDF hlm. 1112-1113 dan
1115-1116, bagian 1 dan 2.2.

\textbf{Laporan sumber.} Agglomeration shadow dijelaskan melalui karya
Fujita dan kolega sebagai penekanan terhadap pertumbuhan kota kecil atau
wilayah tetangga ketika pusat besar menyerap sumber daya, termasuk
tenaga kerja. Penulis menyatakan bahwa borrowed size dan agglomeration
shadow dapat hadir bersamaan: wilayah kecil dapat ditekan dalam beberapa
aktivitas tetapi memperoleh manfaat dari sumber daya dan infrastruktur
pusat dalam aktivitas lain. Literatur empiris yang dirujuk menghasilkan
temuan yang berlawanan tentang apakah eksternalitas jaringan melengkapi
atau menggantikan eksternalitas geografis tradisional.

Locator: TXT baris 294-335 dan 342-351; PDF hlm. 1116-1117, bagian
2.3-2.4.

\textbf{Laporan sumber.} Artikel menempatkan studi sebelumnya yang
dikutip sebagai bukti bahwa hubungan jaringan-lokal belum konsisten.
Sebagian studi yang dirujuk mendukung hubungan komplementer, sedangkan
studi lain mendukung hubungan substitusi. Artikel juga menyebut hasil
dari Tiongkok yang menunjukkan pengaruh eksternalitas jaringan pada
kegiatan inovasi dapat lebih besar daripada eksternalitas geografis
tradisional dalam konteks tertentu.

Locator: TXT baris 337-356; PDF hlm. 1117, bagian 2.4.

\textbf{Interpretasi penulis.} Literatur tersebut dipakai untuk
membangun posisi bahwa efek jaringan harus diuji bersama faktor lokal,
tidak hanya melalui ukuran atau kepadatan wilayah. Pembagian manufaktur
sebagai industri berorientasi ekspor dan non-manufaktur sebagai industri
intra-regional menjadi alasan teoritis untuk mengharapkan perbedaan
koefisien antarindustri.

\textbf{Inferensi pengolahan.} Tinjauan ini adalah narasi konseptual dan
bukan tinjauan sistematis. TXT tidak menjelaskan strategi pencarian,
kriteria pemilihan referensi, cakupan korpus, atau prosedur penilaian
kualitas studi terdahulu. Isi dan kesimpulan penelitian-penelitian yang
dikutip di sini hanya dilaporkan sebagaimana diringkas oleh artikel A16
dan tidak diverifikasi terhadap publikasi asalnya.

\subsection{Method Used}\label{method-used-31}

\textbf{Laporan sumber: pengukuran TFP.} TFP regional dihitung dengan
metodologi indeks yang dirujuk kepada Good et al., Aw et al., dan
Otsuka. Untuk setiap sektor industri, TFP wilayah j pada waktu t
dihitung sebagai penyimpangan relatif dari agen ekonomi representatif
yang memiliki tingkat input dan output rata-rata pada waktu tertentu.
Perubahan TFP agen representatif dari waktu ke waktu turut
diperhitungkan sehingga tingkat TFP antardaerah dan antarperiode dapat
dibandingkan. Output \texttt{Y} adalah output wilayah, input \texttt{X}
adalah faktor produksi, dan \texttt{S} adalah pangsa biaya input. Dalam
data penelitian, input terdiri atas tenaga kerja dan modal.

Locator: TXT baris 359-391 dan 567-576; PDF hlm. 1117-1118 dan 1121,
bagian 3.1 dan 3.4.

\textbf{Laporan sumber: model determinan.} Persamaan (2) memodelkan
\texttt{ln\ TFP} sebagai fungsi dari \texttt{ln\ BS}, \texttt{ln\ DS},
\texttt{ln\ TC}, dan \texttt{ln\ HC}, dengan efek individual wilayah dan
error. \texttt{BS} adalah indeks borrowed size, \texttt{DS} adalah
kepadatan penduduk, \texttt{TC} adalah stok modal transportasi, dan
\texttt{HC} adalah kumpulan tenaga kerja terampil. Dalam kerangka
penulis, BS mewakili limpahan eksternalitas lintas wilayah, sedangkan
DS, TC, dan HC melengkapi eksternalitas yang terlokalisasi.

Locator: TXT baris 394-408; PDF hlm. 1118, bagian 3.2 dan Persamaan (2).

\textbf{Laporan sumber: operasionalisasi variabel.} Indeks borrowed size
dihitung sebagai \texttt{BS\_jt\ =\ jumlah\ GRP\_kt\ /\ d\_jkt} untuk
seluruh wilayah k selain j. GRP atau nilai tambah regional dipakai
sebagai proksi ukuran klaster. \texttt{d\_jkt} adalah resistansi jarak
berbasis waktu tempuh. Waktu tempuh tersebut merupakan rata-rata
berbobot dari waktu perjalanan udara, kereta api, dan mobil, dengan
bobot berupa proporsi penumpang tiap moda untuk pasangan asal-tujuan dan
tahun tertentu. Penulis memilih GRP daripada populasi karena populasi
Kanagawa, misalnya, dapat besar akibat penduduk yang bepergian ke Tokyo
sehingga berpotensi melebihkan ukuran ekonomi yang dipinjam.

Locator: TXT baris 409-450; PDF hlm. 1118-1119, bagian 3.2.

\textbf{Laporan sumber: variabel lokal.} DS dihitung sebagai penduduk
dibagi luas lahan layak huni. TC mencakup agregat stok modal sosial
transportasi berupa jalan, pelabuhan, dan bandar udara. HC dibentuk
karena tidak ada statistik resmi tentang kumpulan tenaga kerja terampil.
Nilainya menggunakan jumlah tenaga kerja menurut tingkat pendidikan dan
jenis kelamin yang diberi bobot rasio upah terhadap upah rata-rata
menurut tingkat pendidikan, lalu dibagi jumlah tenaga kerja wilayah.

Locator: TXT baris 451-493; PDF hlm. 1119-1120, bagian 3.2.

\textbf{Laporan sumber: estimasi dinamis.} Karena terdapat kekhawatiran
endogenitas, penulis menggunakan model panel dinamis dalam bentuk beda
pertama:

\texttt{Delta\ lnTFP\_jt\ =\ alpha\ Delta\ lnTFP\_j,t-1\ +\ beta\ Delta\ lnBS\_jt\ +\ gamma\ Delta\ lnDS\_jt\ +\ delta\ Delta\ lnTC\_jt\ +\ lambda\ Delta\ lnHC\_jt\ +\ Delta\ epsilon\_jt}.

Estimasi dilakukan dengan metode two-step first-difference GMM.
Instrumen yang dicantumkan adalah \texttt{Delta\ TFP(t-2)} dan lag satu
periode dari variabel penjelas. Penulis menjelaskan bahwa model panel
dinamis sesuai untuk banyak unit individual dan rentang waktu yang
relatif singkat, yaitu data 47 prefektur dengan periode 2000-2018.

Locator: TXT baris 496-525 dan 695-703; PDF hlm. 1120-1121 dan 1124,
bagian 3.3 serta catatan Tabel 3.

\textbf{Laporan sumber: tanda koefisien dan elastisitas.} Penulis
mengharapkan koefisien BS, TC, dan HC bertanda positif. DS diperkirakan
positif jika ekonomi kepadatan dominan, tetapi dapat bertanda negatif
jika eksternalitas negatif kepadatan atau kepadatan berlebih dominan. TC
juga dapat bertanda negatif jika efek negatif seperti kemacetan lebih
besar daripada manfaat modal transportasi. Model dinamis menghasilkan
elastisitas jangka pendek dari koefisien estimasi dan elastisitas jangka
panjang dengan memperhitungkan koefisien lag TFP; untuk BS, rumus
kontribusi menggunakan \texttt{beta\ /\ (1\ -\ alpha)}.

Locator: TXT baris 527-558 dan 770-776; PDF hlm. 1120-1121 dan 1126,
bagian 3.3 serta rumus kontribusi.

\textbf{Laporan sumber: pemeriksaan dan robustness.} Penulis menggunakan
uji Arellano-Bond untuk autokorelasi serial dan mensyaratkan
autokorelasi orde pertama dapat signifikan tetapi orde kedua tidak. Uji
Sargan-Hansen digunakan untuk restriksi overidentifikasi instrumen.
Robustness diperiksa dengan membatasi periode observasi menjadi
2000-2015. Kontribusi BS kemudian dihitung untuk perubahan
\texttt{ln\ BS} antara 2000 dan 2018 dan dikorelasikan dengan TFP awal
tahun 2000.

Locator: TXT baris 656-672 dan 764-818; PDF hlm. 1124-1127, bagian 4,
Tabel 3, Tabel 5, Tabel 6, dan Gambar 2-4.

\textbf{Inferensi pengolahan.} Desain ini merupakan analisis panel
observasional dengan estimasi dinamis dan GMM, bukan eksperimen atau
kuasi-eksperimen. Penggunaan lag dan uji overidentifikasi menunjukkan
strategi penanganan endogenitas yang dilaporkan penulis, tetapi TXT
tidak memberi dasar untuk menyimpulkan bahwa seluruh sumber endogenitas
telah terselesaikan. Penulis sendiri menyatakan bahwa autokorelasi
spasial belum dapat ditangani secara memadai.

\subsection{Dataset}\label{dataset-31}

\textbf{Laporan sumber.} Dataset penelitian adalah panel regional Jepang
untuk 2000-2018, mencakup seluruh industri, manufaktur, dan
non-manufaktur di 47 prefektur. Data TFP berasal dari
\emph{Regional-Level Japan Industrial Productivity Database} (R-JIP)
yang diterbitkan Research Institute of Economy, Trade and Industry
(RIETI). Output diukur dengan nilai tambah riil, input tenaga kerja
dengan jam kerja, input modal dengan stok modal riil, dan pangsa biaya
dengan biaya tenaga kerja serta modal yang tercatat dalam basis data.

Locator: TXT baris 561-576; PDF hlm. 1121, bagian 3.4.

\textbf{Laporan sumber.} Data waktu tempuh untuk BS berasal dari
National Integrated Traffic Analysis System (NITAS) yang dikelola
Ministry of Land, Infrastructure, Transport and Tourism. Data penduduk
dan luas untuk DS berasal dari Social and Demographic Statistics System
milik Statistics Bureau. Data stok modal sosial transportasi untuk TC
berasal dari Cabinet Office. Data jumlah tenaga kerja menurut pendidikan
dan data upah menurut pendidikan untuk HC berasal dari Basic Survey on
Employment Structure dan Basic Survey on Wage Structure.

Locator: TXT baris 585-595; PDF hlm. 1122, bagian 3.4.

\textbf{Laporan sumber.} Artikel menyatakan bahwa seluruh data asli
dapat diakses publik tanpa biaya. Bagian Data availability juga
menyatakan bahwa dataset yang digunakan atau dianalisis tersedia dari
penulis korespondensi atas permintaan dan mencantumkan alamat sumber
untuk R-JIP, NITAS, Social and Demographic Statistics System, Social
Capital Stock Estimation Data, Basic Survey on Employment Structure, dan
Basic Survey on Wage Structure.

Locator: TXT baris 563-566 dan 912-923; PDF hlm. 1121 dan 1130, bagian
3.4 serta Data availability.

\textbf{Inferensi pengolahan.} Dataset menggabungkan pengukuran
produktivitas, nilai tambah regional, kepadatan, waktu tempuh lintas
moda, stok infrastruktur, pendidikan, dan upah. TXT tidak memuat berkas
data mentah, kode pemrosesan, daftar industri rinci yang membentuk
agregat manufaktur/non-manufaktur, prosedur harmonisasi semua seri, atau
rincian pembentukan setiap bobot moda transportasi.

\subsection{Population / Sample}\label{population-sample-31}

\textbf{Laporan sumber.} Unit analisis regional adalah 47 prefektur di
Jepang. Estimasi dilakukan untuk tiga cakupan industri: ALL atau seluruh
industri, MANU atau industri manufaktur, dan NMANU atau industri
non-manufaktur. Artikel tidak melaporkan prosedur pengambilan sampel
karena analisis dinyatakan menggunakan panel seluruh prefektur.

Locator: TXT baris 170-176, 561-565, dan 697-703; PDF hlm. 1113-1114,
1121, dan 1124, bagian 1, 3.4, serta catatan Tabel 3.

\textbf{Laporan sumber.} Model utama pada Tabel 3 memiliki 799 observasi
untuk masing-masing dari tiga cakupan industri. Robustness dengan
periode 2000-2015 pada Tabel 5 memiliki 658 observasi untuk
masing-masing cakupan. Tabel 3 menyebut regresi diestimasi menggunakan
data panel seluruh prefektur.

Locator: TXT baris 680-703 dan 780-803; PDF hlm. 1124 dan 1126, Tabel 3
dan Tabel 5.

\textbf{Inferensi pengolahan.} Sampel analitis berada pada tingkat
prefektur-industri-tahun, bukan tingkat perusahaan atau pekerja. Jumlah
industri individual yang berada di dalam agregat MANU dan NMANU, rincian
observasi yang hilang, dan alasan perbedaan jumlah observasi antara
kemungkinan kombinasi unit-waktu tidak dijelaskan secara terpisah dalam
TXT. Karena unit analisisnya agregat, hasil tidak langsung menunjukkan
perilaku perusahaan individual.

\subsection{Dependent Variables}\label{dependent-variables-31}

\textbf{Laporan sumber.} Variabel terikat substantif adalah tingkat TFP
regional menurut industri. TFP mencakup output nilai tambah riil dan
input tenaga kerja serta modal, dengan pangsa biaya input dalam
konstruksi indeks. TFP dihitung sebagai tingkat relatif terhadap agen
representatif dan memperhitungkan perubahan agen representatif sepanjang
waktu.

Locator: TXT baris 192-211 dan 363-391; PDF hlm. 1114 dan 1117-1118,
bagian 2.1 dan 3.1.

\textbf{Laporan sumber.} Dalam estimasi utama, variabel terikat ditulis
sebagai \texttt{Delta\ lnTFP(t)}, yaitu perubahan pertama log TFP.
\texttt{Delta\ lnTFP(t-1)} menjadi variabel dinamis di sisi kanan.
Estimasi dipisahkan untuk ALL, MANU, dan NMANU. Analisis kontribusi BS
terhadap pengejaran ketertinggalan memakai kontribusi jangka panjang BS
dan hubungan korelasionalnya dengan TFP awal 2000, bukan variabel
terikat baru dalam regresi utama.

Locator: TXT baris 522-538, 680-703, dan 770-818; PDF hlm. 1121, 1124,
dan 1126-1127, Persamaan (3), Tabel 3, rumus kontribusi, dan Tabel 6.

\textbf{Inferensi pengolahan.} TFP adalah indeks relatif, bukan ukuran
produktivitas tenaga kerja sederhana. TXT menjelaskan konstruksi
konseptual dan variabel input, tetapi tidak memberikan kode, nilai agen
representatif, atau langkah numerik lengkap yang memungkinkan reproduksi
independen hanya dari TXT. Istilah ``produktivitas'' dalam kesimpulan
harus dibaca sebagai perubahan atau tingkat TFP yang didefinisikan
artikel.

\subsection{Results}\label{results-31}

\textbf{Laporan sumber: statistik deskriptif.} Rata-rata TFP untuk
seluruh industri meningkat dari 0,035 pada 2000-2009 menjadi 0,075 pada
2010-2018. Pada saat yang sama, rentang TFP seluruh industri meningkat
dari 0,727 menjadi 0,761, sehingga penulis melaporkan peningkatan
produktivitas bersama pelebaran disparitas regional. Rata-rata TFP
manufaktur meningkat dari 0,179 menjadi 0,366 dan rentangnya menyempit
dari 1,302 menjadi 1,149. Rata-rata TFP non-manufaktur berubah dari
-0,007 menjadi -0,009 dan rentangnya melebar dari 0,767 menjadi 0,831.
Penulis menyimpulkan bahwa manufaktur secara konsisten lebih produktif
dan tumbuh lebih kuat, sedangkan non-manufaktur hanya mengalami
perbaikan minimal serta disparitas yang lebih besar.

Locator: TXT baris 596-602 dan 634-641; PDF hlm. 1122-1123, bagian 3.4
dan Tabel 2.

\textbf{Laporan sumber: perubahan variabel eksternalitas.} Rata-rata BS
hampir tidak berubah antara 2000-2009 dan 2010-2018, tetapi nilai
maksimum naik dan minimum turun, yang oleh penulis dibaca sebagai
pelebaran disparitas borrowed size. Rata-rata DS relatif tidak berubah,
sementara nilai maksimum dan rentangnya meningkat. Rata-rata TC dan HC
meningkat. Rentang TC melebar, sedangkan rentang HC menyempit, yang oleh
penulis dikaitkan dengan kemajuan keterampilan tenaga kerja yang lebih
tersebar secara nasional.

Locator: TXT baris 642-652; PDF hlm. 1123, bagian 3.4 dan Tabel 2.

\textbf{Laporan sumber: estimasi model utama.} Tabel 3 melaporkan hasil
two-step first-difference GMM dengan 799 observasi per model. Untuk ALL,
koefisien lag \texttt{Delta\ lnTFP(t-1)} adalah 0,424 dengan standar
error 0,034 dan signifikan pada 1\%; BS 0,153 dengan standar error 0,066
dan signifikan pada 5\%; DS 0,402 dengan standar error 0,048 dan
signifikan pada 1\%; TC 0,244 dengan standar error 0,037 dan signifikan
pada 1\%; serta HC 0,187 dengan standar error 0,001 dan signifikan pada
1\%.

Locator: TXT baris 680-703; PDF hlm. 1124, Tabel 3.

\textbf{Laporan sumber: estimasi manufaktur.} Untuk MANU, koefisien lag
TFP adalah 0,446 dengan standar error 0,018 dan signifikan pada 1\%.
Koefisien BS adalah 0,190 dengan standar error 0,085 dan signifikan pada
5\%. DS bernilai -0,015 dengan standar error 0,152 dan tidak diberi
tanda signifikansi. TC bernilai 0,909 dengan standar error 0,031 dan
signifikan pada 1\%. HC bernilai 0,132 dengan standar error 0,039 dan
signifikan pada 1\%. Dengan demikian, di antara variabel penjelas, TC
memiliki koefisien terbesar pada manufaktur.

Locator: TXT baris 680-703; PDF hlm. 1124, Tabel 3.

\textbf{Laporan sumber: estimasi non-manufaktur.} Untuk NMANU, koefisien
lag TFP adalah 0,381 dengan standar error 0,051 dan signifikan pada 1\%.
Koefisien BS adalah 0,169 dengan standar error 0,022 dan signifikan pada
1\%. DS bernilai 0,151 dengan standar error 0,061 dan signifikan pada
5\%. TC bernilai -0,130 dengan standar error 0,015 dan signifikan pada
1\%. HC bernilai 0,130 dengan standar error 0,016 dan signifikan pada
1\%.

Locator: TXT baris 680-703; PDF hlm. 1124, Tabel 3.

\textbf{Interpretasi penulis atas model utama.} Untuk ALL, semua
koefisien bertanda positif dan signifikan. Penulis menyatakan DS
mempunyai pengaruh terbesar, sedangkan TC dan HC juga lebih besar
daripada efek borrowed size, sehingga eksternalitas lokal dan jaringan
dipandang bersifat komplementer. Untuk MANU dan NMANU, BS positif dan
signifikan pada keduanya, dengan koefisien manufaktur lebih besar. DS
hanya signifikan pada NMANU, yang ditafsirkan sebagai bukti bahwa
aglomerasi penduduk lebih penting bagi industri intra-regional. TC
sangat besar dan positif pada MANU, sedangkan pengaruh HC positif tetapi
paling kecil di antara variabel penjelas pada dua kelompok industri.

Locator: TXT baris 707-735; PDF hlm. 1124-1125, bagian 4.

\textbf{Laporan sumber: elastisitas jangka pendek dan jangka panjang.}
Tabel 4 melaporkan nilai berikut: untuk ALL, BS 0,153 dan 0,265; DS
0,402 dan 0,697; TC 0,244 dan 0,424; HC 0,187 dan 0,324. Untuk MANU, BS
0,190 dan 0,343; DS -0,015 dan -0,028; TC 0,909 dan 1,639; HC 0,132 dan
0,237. Untuk NMANU, BS 0,169 dan 0,274; DS 0,151 dan 0,244; TC -0,130
dan -0,210; HC 0,130 dan 0,209. Penulis menyatakan bahwa efek jangka
panjang lebih besar daripada efek jangka pendek dan bahwa dampak jangka
panjang BS paling besar pada manufaktur, sementara dampak jangka panjang
TC juga sangat besar pada manufaktur.

Locator: TXT baris 736-763; PDF hlm. 1125-1126, bagian 4 dan Tabel 4.

\textbf{Interpretasi penulis atas elastisitas.} Penulis mengaitkan
besarnya efek jangka panjang TC pada manufaktur dengan peran
infrastruktur transportasi berkualitas tinggi dalam memperkuat
eksternalitas intra- dan antarwilayah. Efek jangka panjang DS yang
positif pada non-manufaktur dikaitkan dengan ukuran permintaan lokal dan
pembentukan saluran penjualan yang lebih efisien.

Locator: TXT baris 736-763; PDF hlm. 1125-1126, bagian 4.

\textbf{Laporan sumber: robustness.} Estimasi periode 2000-2015 pada
Tabel 5 memiliki 658 observasi per model. Koefisien BS masing-masing
adalah 0,248 untuk ALL, 0,401 untuk MANU, dan 0,212 untuk NMANU;
ketiganya diberi tanda signifikan pada 1\%. Koefisien DS adalah 0,398
untuk ALL dan 0,227 untuk NMANU, keduanya signifikan pada 1\% dan 5\%;
koefisien DS MANU adalah -0,344 tanpa tanda signifikansi. Koefisien TC
adalah 0,320 untuk ALL, 1,151 untuk MANU, dan -0,100 untuk NMANU,
semuanya diberi tanda signifikan pada 1\%. Koefisien HC adalah 0,224,
0,197, dan 0,146 untuk ALL, MANU, dan NMANU, semuanya diberi tanda
signifikan pada 1\%. Penulis menyatakan bahwa tidak ada perubahan besar
pada urutan relatif koefisien maupun signifikansinya.

Locator: TXT baris 764-803; PDF hlm. 1126-1127, bagian 4 dan Tabel 5.

\textbf{Laporan sumber: diagnostik estimasi.} Pada model utama,
probabilitas J-statistic adalah 0,369 untuk ALL, 0,370 untuk MANU, dan
0,157 untuk NMANU. Statistik AR(1) masing-masing -5,098, -3,039, dan
-3,987 dengan tanda signifikansi; statistik AR(2) masing-masing -0,510,
1,025, dan -1,011 tanpa tanda signifikansi. Penulis menyatakan kondisi
AR(2) dan uji Sargan-Hansen memenuhi persyaratan yang mereka gunakan
untuk estimasi panel dinamis. Pada robustness, probabilitas J-statistic
adalah 0,387, 0,345, dan 0,165; statistik AR(2) adalah -1,350, 0,581,
dan -0,935 tanpa tanda signifikansi.

Locator: TXT baris 658-668, 680-703, dan 780-803; PDF hlm. 1124 dan
1126, Tabel 3 dan Tabel 5.

\textbf{Laporan sumber: kontribusi BS dan pengejaran ketertinggalan.}
Kontribusi BS dihitung menggunakan elastisitas jangka panjang
\texttt{beta\ /\ (1\ -\ alpha)} dikalikan perubahan \texttt{ln\ BS}
antara 2000 dan 2018. Penulis menyatakan bahwa hubungan antara
kontribusi BS dan TFP awal 2000 menurun ke kanan pada semua grafik.
Koefisien korelasi pada Tabel 6 adalah -0,427 untuk ALL dengan t-value
3,164 dan p-value 0,003; -0,368 untuk MANU dengan t-value 2,657 dan
p-value 0,011; serta -0,369 untuk NMANU dengan t-value 2,663 dan p-value
0,011. Penulis membaca korelasi negatif tersebut sebagai indikasi bahwa
kontribusi BS lebih besar di wilayah yang semula berproduktivitas rendah
dan mungkin mengurangi disparitas TFP regional.

Locator: TXT baris 770-818 dan 867-875; PDF hlm. 1126-1129, rumus
kontribusi, Gambar 2-4, dan Tabel 6.

\textbf{Interpretasi penulis.} Penulis mengaitkan perluasan Shinkansen
di Tohoku dan Kyushu dengan berkurangnya waktu tempuh menuju kawasan
metropolitan besar. Karena wilayah-wilayah tersebut disebut memiliki
produktivitas awal yang rendah, penulis menyatakan bahwa peningkatan
transportasi berkualitas tinggi sangat mungkin membantu perbaikan
produktivitas industri di sana.

Locator: TXT baris 820-841; PDF hlm. 1127-1128, bagian 4.

\textbf{Inferensi pengolahan.} Bukti regresi mendukung adanya asosiasi
positif BS dengan perubahan TFP dalam tiga agregat industri dan pola
korelasi yang konsisten dengan konvergensi kontribusi BS. Namun,
korelasi kontribusi dengan TFP awal tidak dengan sendirinya membuktikan
bahwa jaringan transportasi menyebabkan pengejaran ketertinggalan.
Selain itu, Tabel 2 melaporkan bahwa disparitas TFP seluruh industri dan
non-manufaktur melebar, sehingga kontribusi BS yang berpotensi
mengurangi disparitas tidak boleh disamakan dengan penurunan disparitas
TFP total.

\textbf{Inferensi pengolahan.} Hasil menunjukkan heterogenitas yang
penting. Kesimpulan umum tentang komplementaritas eksternalitas lokal
dan jaringan paling jelas pada model ALL; pada MANU, DS tidak signifikan
dan TC sangat dominan, sedangkan pada NMANU, TC bertanda negatif dan DS
positif. Tanda negatif TC NMANU secara umum sesuai dengan kemungkinan
eksternalitas negatif transportasi yang dijelaskan dalam metode, tetapi
TXT tidak menguji secara terpisah apakah kemacetan atau mekanisme lain
benar-benar menjadi penyebabnya.

\subsection{Findings}\label{findings-31}

\textbf{Temuan yang dilaporkan sumber.} Borrowed size berhubungan
positif dan signifikan dengan TFP pada seluruh industri, manufaktur, dan
non-manufaktur. Efek tersebut bertahan dalam estimasi jangka panjang dan
robustness periode 2000-2015. Koefisien serta elastisitas jangka panjang
BS lebih besar pada manufaktur daripada non-manufaktur dalam hasil
utama.

Locator: TXT baris 707-721, 736-751, dan 780-803; PDF hlm. 1124-1126,
bagian 4 dan Tabel 3-5.

\textbf{Temuan yang dilaporkan sumber.} Eksternalitas lokal juga
berkaitan dengan TFP, tetapi bentuknya berbeda menurut industri.
Kepadatan penduduk paling besar pada ALL dan signifikan pada NMANU,
sedangkan modal transportasi paling kuat serta positif pada MANU. HC
bertanda positif pada semua cakupan industri. Pada NMANU, TC bertanda
negatif dan signifikan.

Locator: TXT baris 673-735; PDF hlm. 1124-1125, Tabel 3 dan pembahasan
bagian 4.

\textbf{Temuan yang dilaporkan sumber.} Kontribusi BS berkorelasi
negatif dengan TFP awal di tiga cakupan industri. Artikel menafsirkan
pola ini sebagai indikasi bahwa wilayah berproduktivitas rendah menerima
kontribusi BS lebih besar, sehingga jaringan antarwilayah mungkin
mendukung pengejaran ketertinggalan.

Locator: TXT baris 807-818 dan 867-875; PDF hlm. 1127-1129, Gambar 2-4
dan Tabel 6.

\textbf{Interpretasi penulis.} Secara keseluruhan, penulis menyimpulkan
bahwa borrowed size melengkapi ekonomi eksternal lokal. Perluasan
jaringan transportasi berkualitas tinggi dipandang memperluas cakupan
geografis kegiatan produksi serta memungkinkan limpahan yang sebelumnya
terikat wilayah mengalir antarwilayah.

Locator: TXT baris 844-896; PDF hlm. 1128-1130, bagian 5 Conclusion.

\textbf{Inferensi pengolahan.} Temuan paling kuat adalah asosiasi
diferensial menurut industri dan perbedaan antara efek jaringan dalam
model dengan dinamika disparitas TFP keseluruhan. Temuan tersebut tidak
mengisolasi borrowed size dari faktor lain yang berubah bersamaan, dan
tidak mengukur secara langsung aliran pengetahuan, perjalanan, rantai
pasok, atau transaksi antarwilayah.

\subsection{Conclusions}\label{conclusions-31}

\textbf{Kesimpulan yang dinyatakan penulis.} Studi menyimpulkan bahwa
borrowed size memengaruhi TFP industri manufaktur dan non-manufaktur.
Efek jangka panjangnya lebih besar pada manufaktur, yang dipandang
mencerminkan interaksi eksternalitas geografis lintas batas ketika
jaringan transportasi berkualitas tinggi memperluas ruang produksi.
Kontribusi borrowed size juga lebih besar di wilayah dengan
produktivitas awal yang rendah.

Locator: TXT baris 844-880; PDF hlm. 1128-1129, bagian 5 Conclusion.

\textbf{Kesimpulan yang dinyatakan penulis.} Eksternalitas lokal
dianggap penting bagi kedua kelompok industri dan terutama besar pada
non-manufaktur. Penulis menyatakan bahwa efek borrowed size melengkapi
ekonomi eksternal lokal. Untuk meningkatkan produktivitas, penulis
merekomendasikan penguatan strategis wilayah individual sekaligus
pembangunan jaringan transportasi berkualitas tinggi yang menghubungkan
wilayah dan klaster.

Locator: TXT baris 881-896; PDF hlm. 1129-1130, bagian 5 Conclusion.

\textbf{Kesimpulan yang dinyatakan penulis.} Penulis juga menyatakan
bahwa peningkatan keterampilan tenaga kerja berkontribusi terhadap
produktivitas di semua industri. Dukungan kebijakan kepada universitas
dan lembaga riset regional dipandang sebagai cara untuk meningkatkan
produktivitas industri.

Locator: TXT baris 897-903; PDF hlm. 1130, bagian 5 Conclusion.

\textbf{Inferensi pengolahan.} Kesimpulan tersebut layak dibaca sebagai
kesimpulan berbasis model panel observasional dan bukan estimasi dampak
kausal program transportasi, kebijakan klaster, atau dukungan
universitas. Artikel tidak menunjukkan hasil evaluasi sebelum-sesudah
dengan kelompok pembanding untuk kebijakan tertentu.

\subsection{Contributions}\label{contributions-31}

\textbf{Kontribusi yang diklaim atau ditunjukkan sumber.}

\begin{itemize}
\tightlist
\item
  \textbf{Kontribusi empiris:} menggunakan data industri regional Jepang
  2000-2018 dari 47 prefektur untuk menguji eksternalitas geografis pada
  TFP, bukan hanya pada pertumbuhan ekonomi regional umum.
\item
  \textbf{Kontribusi perbandingan industri:} memisahkan hasil ALL, MANU,
  dan NMANU sehingga perbedaan antara industri berorientasi ekspor dan
  industri intra-regional dapat ditunjukkan secara statistik.
\item
  \textbf{Kontribusi temporal:} menggunakan fungsi TFP dinamis untuk
  melaporkan elastisitas jangka pendek dan jangka panjang.
\item
  \textbf{Kontribusi pengukuran:} membangun BS dengan GRP sebagai proksi
  ukuran klaster dan resistansi jarak berbasis waktu tempuh lintas moda,
  bukan hanya ukuran populasi atau jarak fisik.
\item
  \textbf{Kontribusi terhadap isu konvergensi:} menghitung kontribusi BS
  selama 2000-2018 dan menghubungkannya dengan TFP awal untuk menilai
  kemungkinan pengejaran ketertinggalan wilayah.
\item
  \textbf{Kontribusi kebijakan:} menghubungkan hasil produktivitas
  dengan penguatan infrastruktur transportasi, jaringan antarwilayah,
  keterampilan tenaga kerja, universitas, dan lembaga riset regional.
\end{itemize}

Locator: TXT baris 150-176, 363-450, 539-558, 736-763, dan 844-903; PDF
hlm. 1113-1114, 1117-1119, 1121, 1125-1130.

\textbf{Inferensi pengolahan.} Kontribusi terkuat artikel adalah
menggabungkan ukuran TFP industri, indeks jaringan berbasis GRP dan
waktu tempuh, serta dinamika jangka pendek-panjang dalam satu analisis
prefektur. Klaim kebaruan atau keunggulan relatif terhadap seluruh
literatur tidak dapat diverifikasi lebih jauh karena TXT tidak
menyediakan tinjauan sistematis atau perbandingan metodologis
menyeluruh.

\subsection{Research Gap}\label{research-gap-31}

\textbf{Kesenjangan yang dinyatakan sumber sebelum penelitian.}

\begin{itemize}
\tightlist
\item
  Bukti empiris mengenai eksternalitas regional pada aktivitas industri
  dinilai belum memadai.
\item
  Perbedaan efek jaringan antarwilayah dan eksternalitas geografis
  konvensional menurut industri belum cukup dijelaskan.
\item
  Temuan sebelumnya tentang hubungan komplementer atau substitutif
  antara eksternalitas jaringan dan eksternalitas lokal masih
  bertentangan.
\item
  Efek borrowed size terhadap kinerja produktivitas industri, termasuk
  perbedaan jangka pendek dan jangka panjang, belum cukup diuji.
\item
  Keadaan yang membuat agglomeration shadow dominan belum sepenuhnya
  jelas karena kompleksitas analisis.
\end{itemize}

Locator: TXT baris 134-149, 322-335, dan 337-356; PDF hlm. 1113,
1116-1117, dan 1117, bagian 1 dan 2.3-2.4.

\textbf{Kesenjangan yang masih tersisa setelah studi.} Penulis secara
eksplisit menyatakan bahwa kemungkinan autokorelasi spasial belum dapat
dibahas sepenuhnya dan bahwa metode panel dinamis yang kredibel serta
mempertimbangkan autokorelasi spasial belum berkembang secara memadai.

Locator: TXT baris 904-909; PDF hlm. 1130, bagian 5 Conclusion.

\textbf{Inferensi pengolahan.} Dari desain dan variabel yang dilaporkan,
masih terbuka kesenjangan mengenai mekanisme langsung borrowed size: TXT
tidak menunjukkan pengukuran aliran pengetahuan, perjalanan, hubungan
perusahaan, rantai pasok, atau penggunaan kapasitas eksternal secara
langsung. Efek kebijakan transportasi tertentu juga belum dipisahkan
dari perubahan ekonomi regional lain. Ini adalah inferensi dari cakupan
data dan model, bukan daftar agenda yang dinyatakan penulis.

\subsection{Limitations}\label{limitations-31}

\textbf{Keterbatasan yang dinyatakan sumber.} Artikel mengakui bahwa
kemungkinan autokorelasi spasial tidak dapat dibahas sepenuhnya.
Interaksi spasial ada dalam aktivitas produksi, tetapi metode panel
dinamis yang kredibel dan sekaligus memperhitungkan autokorelasi spasial
dinyatakan belum tersedia atau belum cukup berkembang untuk digunakan.
Penulis menyebut pengembangan metode analitis untuk isu tersebut sebagai
kebutuhan masa depan.

Locator: TXT baris 904-909; PDF hlm. 1130, bagian 5 Conclusion.

\textbf{Keterbatasan yang diturunkan dari isi TXT sebagai inferensi
pengolahan.}

\begin{itemize}
\tightlist
\item
  \textbf{Identifikasi kausal:} model menggunakan data observasional dan
  lag sebagai instrumen dalam GMM, tetapi TXT tidak menunjukkan desain
  kontrafaktual atau variasi eksogen yang mengisolasi dampak pembangunan
  transportasi terhadap TFP.
\item
  \textbf{Ketergantungan spasial:} keterbatasan yang diakui penulis
  berarti standard error, koefisien, atau interpretasi spillover dapat
  dipengaruhi hubungan antarprefektur yang belum dimodelkan secara
  spasial. Besaran dampaknya tidak tersedia dalam TXT.
\item
  \textbf{Validitas konstruk BS:} BS dibentuk dari GRP wilayah lain dan
  resistansi waktu tempuh. Ukuran ini menangkap potensi akses ekonomi
  berbobot jarak, tetapi TXT tidak melaporkan observasi langsung tentang
  interaksi, kolaborasi, arus barang, arus pekerja, atau limpahan
  pengetahuan aktual.
\item
  \textbf{Agregasi:} unit analisis adalah prefektur dan tiga kelompok
  industri besar. Hasilnya tidak langsung dapat dipindahkan ke
  perusahaan, kota, atau subsektor individual.
\item
  \textbf{Cakupan waktu:} pengamatan berakhir pada 2018. TXT tidak
  menyediakan pengujian periode setelah itu atau pemisahan dampak setiap
  program kebijakan yang disebut di pendahuluan.
\item
  \textbf{Spesifikasi:} model utama hanya menampilkan BS, DS, TC, HC,
  efek individual, dan dinamika TFP. TXT tidak menjelaskan secara
  lengkap apakah efek waktu, struktur industri yang lebih rinci, atau
  kovariat regional lain dimasukkan.
\item
  \textbf{Reproduksibilitas:} TXT tidak menyertakan data mentah, kode,
  prosedur pembersihan, atau rincian lengkap konstruksi indeks TFP dan
  bobot waktu tempuh.
\end{itemize}

\textbf{Informasi tidak tersedia dalam TXT.} Daftar keterbatasan
tambahan yang disusun penulis selain isu autokorelasi spasial tidak
disebutkan dalam file. Penilaian tambahan di atas tidak boleh dibaca
sebagai pengakuan penulis.

\subsection{Challenges}\label{challenges-31}

\textbf{Tantangan yang dilaporkan sumber.} Pengukuran TFP menurut
industri di Jepang menghadapi keterbatasan karena data untuk mengukur
TFP industri belum berkembang dengan baik. Artikel menggunakan R-JIP
yang baru dihimpun untuk mengatasi kebutuhan tersebut. Penulis juga
menghadapi rentang waktu yang relatif singkat dan jumlah unit lintas
seksi yang besar, sehingga memilih estimasi panel dinamis.

Locator: TXT baris 514-521 dan 567-572; PDF hlm. 1120-1121, bagian
3.3-3.4.

\textbf{Tantangan yang dilaporkan sumber.} Tidak ada statistik resmi
untuk pool tenaga kerja terampil, sehingga HC harus dibangun dari data
tenaga kerja menurut pendidikan dan data upah. Pengukuran kedekatan
klaster juga memerlukan keputusan metodologis: GRP dipilih untuk
menghindari bias populasi akibat komuter, dan jarak fisik diganti dengan
waktu tempuh berbobot udara, kereta api, serta mobil.

Locator: TXT baris 415-450 dan 468-485; PDF hlm. 1118-1120, bagian 3.2.

\textbf{Tantangan yang dilaporkan sumber.} Hubungan borrowed size dan
agglomeration shadow kompleks. Karena sulit menafsirkan agglomeration
shadow secara eksplisit, penulis memilih fokus pada efek bersih
eksternalitas geografis. Artikel juga mengakui kesulitan memperlakukan
autokorelasi spasial dalam model panel dinamis.

Locator: TXT baris 316-335 dan 904-909; PDF hlm. 1116-1117 dan 1130,
bagian 2.3 dan 5.

\textbf{Inferensi pengolahan.} Tantangan analitis utama adalah
membedakan potensi jaringan yang dihitung dari GRP dan waktu tempuh dari
mekanisme interaksi aktual. Tantangan lain ialah menafsirkan koefisien
TC yang berlawanan tanda antara MANU dan NMANU tanpa observasi terpisah
tentang kemacetan, kapasitas terpakai, atau jenis infrastruktur yang
bekerja melalui masing-masing industri.

\subsection{Practical Implications}\label{practical-implications-31}

\textbf{Implikasi yang dinyatakan sumber.} Peningkatan produktivitas
industri memerlukan dua arah kebijakan sekaligus: penguatan strategis
setiap wilayah dan pembangunan jaringan transportasi berkualitas tinggi
yang menghubungkan wilayah. Menghubungkan klaster melalui jaringan
transit cepat dan memperkuat hubungan antarwilayah dipandang dapat
memperluas limpahan eksternalitas yang sebelumnya terikat pada wilayah
asal.

Locator: TXT baris 887-896; PDF hlm. 1129-1130, bagian 5 Conclusion.

\textbf{Implikasi yang dinyatakan sumber.} Karena HC bertanda positif
pada semua kelompok industri, dukungan terhadap peningkatan keterampilan
tenaga kerja dipandang relevan untuk daya saing industri regional.
Penulis secara khusus menyebut dukungan kebijakan bagi universitas dan
lembaga riset regional sebagai sarana meningkatkan produktivitas
industri.

Locator: TXT baris 897-903; PDF hlm. 1130, bagian 5 Conclusion.

\textbf{Inferensi pengolahan.} Kebijakan tidak sebaiknya diterapkan
seragam tanpa memperhatikan industri. Hasil mengindikasikan alasan untuk
memberi perhatian lebih besar pada konektivitas transportasi bagi
manufaktur, sementara kepadatan atau ukuran permintaan lokal tampak
lebih relevan bagi non-manufaktur. Ini adalah implikasi analitis dari
perbedaan koefisien, bukan rekomendasi sektoral yang dirinci secara
formal oleh penulis.

\textbf{Inferensi pengolahan.} Rekomendasi kebijakan perlu disertai
kehati-hatian karena autokorelasi spasial belum ditangani dan desain
tidak mengisolasi dampak kausal proyek transportasi tertentu. Hasil
lebih tepat dipakai sebagai dasar diagnosis dan penentuan hipotesis
kebijakan daripada bukti tunggal bahwa investasi tertentu pasti
meningkatkan TFP.

\subsection{Applications}\label{applications-31}

\textbf{Aplikasi yang didukung oleh desain penelitian.} Kerangka artikel
dapat digunakan untuk diagnosis produktivitas regional dengan
menggabungkan TFP menurut industri, ukuran potensi ekonomi wilayah lain,
waktu tempuh lintas moda, kepadatan penduduk, stok modal transportasi,
dan keterampilan tenaga kerja. Aplikasi tersebut masih berupa penggunaan
kerangka pengukuran; artikel tidak melaporkan implementasi kebijakan
berdasarkan hasilnya.

Locator: TXT baris 363-493 dan 561-595; PDF hlm. 1117-1122, bagian 3.

\textbf{Aplikasi yang disarankan oleh temuan sumber.} Hasil dapat
diterapkan sebagai masukan untuk memprioritaskan hubungan transportasi
antarwilayah dan koneksi antarklaster, terutama bila tujuan kebijakan
adalah memperluas manfaat ke wilayah yang produktivitas awalnya rendah.
Hasil juga dapat dipakai untuk membedakan kebutuhan diagnosis manufaktur
dan non-manufaktur.

Locator: TXT baris 844-903; PDF hlm. 1128-1130, bagian 5 Conclusion.

\textbf{Inferensi pengolahan.} Aplikasi pada wilayah di luar 47
prefektur Jepang, pada periode lain, atau pada skala perusahaan
memerlukan pengujian ulang karena definisi BS, sumber data, struktur
industri, dan pola spasial dapat berbeda. TXT tidak menyediakan validasi
eksternal atau demonstrasi bahwa model telah diterapkan di luar konteks
tersebut.

\subsection{Future Research
(Optional)}\label{future-research-optional-31}

\textbf{Agenda yang dinyatakan sumber.} Penelitian masa depan perlu
mengembangkan metode analitis yang kredibel untuk panel dinamis dan
autokorelasi spasial, agar interaksi spasial dalam aktivitas produksi
dapat dipertimbangkan secara statistik.

Locator: TXT baris 904-909; PDF hlm. 1130, bagian 5 Conclusion.

\textbf{Inferensi pengolahan, bukan agenda eksplisit penulis.}
Berdasarkan keterbatasan konstruk dan cakupan model, pengembangan
lanjutan secara logis dapat menguji ukuran interaksi antarwilayah yang
lebih langsung, memisahkan kebijakan transportasi tertentu dari
perubahan regional lain, dan memeriksa heterogenitas pada subsektor yang
lebih rinci. TXT tidak menyatakan rincian agenda tersebut secara
eksplisit.

\subsection{Provenance Notes}\label{provenance-notes-31}

\begin{itemize}
\tightlist
\item
  \textbf{Sumber tunggal:}
  \texttt{/Users/mac/Documents/Mac/{[}2{]}\ Obsidian\ Vault/zettelkasten/source/journal/A16-otsuka-2024-effects-of-regional-network-economies-on-industrial-productivity-in-japan-dynamic-total-factor-productivity-function-approach-10.1007-s41685-024-00358-2.txt}.
\item
  \textbf{Cakupan pembacaan:} seluruh TXT, baris 1-1.139, termasuk
  metadata PDF, teks artikel, tabel, keterangan gambar, bagian Data
  availability, deklarasi, dan referensi.
\item
  \textbf{Metadata ekstraksi:} TXT menyatakan sumber PDF berada di
  folder \texttt{journal/}, diekstrak pada 2026-08-31 dengan
  \texttt{pdftotext\ -layout}; metadata PDF menyatakan 24 halaman.
  Locator utama dalam review menggunakan nomor halaman PDF yang
  tercetak, bagian, persamaan, tabel, atau gambar; rentang baris TXT
  ditambahkan untuk audit.
\item
  \textbf{Batasan akses:} review ini tidak membuka atau memeriksa PDF,
  data mentah, kode estimasi, atau publikasi yang ada di daftar
  referensi. Semua laporan tentang referensi terdahulu berasal dari
  uraian A16 dalam TXT.
\item
  \textbf{Pemisahan epistemik:} label \texttt{Laporan\ sumber} merangkum
  pernyataan, angka, atau prosedur yang tercantum dalam TXT;
  \texttt{Interpretasi\ penulis} merangkum cara artikel membaca
  hasilnya; \texttt{Inferensi\ pengolahan} adalah evaluasi yang
  diturunkan dari isi TXT dan bukan klaim penulis.
\item
  \textbf{Informasi yang tidak tersedia:} TXT tidak memberikan nomor isu
  jurnal, kode/data replikasi, rincian lengkap pembersihan data dan
  konstruksi indeks, maupun verifikasi peer-review. TXT juga memiliki
  ketidakkonsistenan penomoran gambar: narasi pada baris 807-809
  menyebut Gambar 1-3, sedangkan keterangan yang tampil adalah Gambar 2
  untuk ALL, Gambar 3 untuk MANU, dan Gambar 4 untuk NMANU; judul Tabel
  6 kembali merujuk Gambar 1-3. Review mengikuti isi keterangan gambar
  dan angka Tabel 6, bukan mengoreksi penomoran tersebut.
\item
  \textbf{Integritas sumber:} tidak ada informasi dari luar TXT yang
  digunakan dan file TXT sumber tidak diubah.
\end{itemize}

\section{Review A17}\label{review-a17}

Sumber yang ditinjau: Akihiro Otsuka, ``Borrowed size and agglomeration
shadows in Japan: a spatial econometric analysis of regional
productivity''. Artikel ini tercantum sebagai artikel \emph{Asia-Pacific
Journal of Regional Science} (2025) 9:1115-1144 dengan DOI
\texttt{10.1007/s41685-025-00398-2}.

\subsection{Summarized Abstract}\label{summarized-abstract-32}

\begin{itemize}
\tightlist
\item
  \textbf{Dilaporkan sumber:} Studi ini menguji pengaruh
  saling-ketergantungan spasial terhadap \emph{regional total factor
  productivity} (RTFP) di Jepang dengan memusatkan perhatian pada dua
  efek yang berlawanan, yaitu \emph{borrowed size} dan
  \emph{agglomeration shadow}. Analisis memakai data panel 47 prefektur
  Jepang pada 2000-2018 dan membandingkan sektor manufaktur dengan
  nonmanufaktur melalui \emph{Spatial Durbin Model} (SDM) serta
  \emph{Spatial Durbin-SARAR Model} (SDM-SARAR). Model menguraikan efek
  menjadi efek langsung di dalam wilayah dan efek tidak langsung atau
  limpahan antardaerah.
\item
  \textbf{Dilaporkan sumber:} Aglomerasi industri dan modal transportasi
  dilaporkan berpengaruh positif terhadap produktivitas lokal di kedua
  sektor. Limpahan negatif dari konsentrasi di wilayah tetangga juga
  ditemukan dan ditafsirkan sebagai \emph{agglomeration shadow}. Modal
  manusia berpengaruh positif secara konsisten pada manufaktur, dengan
  bukti limpahan antardaerah. Abstrak menyarankan respons kebijakan yang
  dibedakan menurut sektor: jejaring klaster industri dan mobilitas
  tenaga kerja untuk manufaktur, serta ekonomi jasa yang lebih
  terdesentralisasi dan investasi infrastruktur yang berkeadilan spasial
  untuk nonmanufaktur. {[}Abstract, hlm. 1115{]}
\end{itemize}

\subsection{Problem Statement}\label{problem-statement-32}

\begin{itemize}
\tightlist
\item
  \textbf{Dilaporkan sumber:} Masalah utama penelitian adalah bagaimana
  efek limpahan antardaerah memengaruhi RTFP melalui dua mekanisme yang
  berlawanan. \emph{Borrowed size} menggambarkan keuntungan
  produktivitas dari pemanfaatan sumber daya ekonomi di wilayah
  tetangga, sedangkan \emph{agglomeration shadow} menggambarkan
  penyerapan talenta dan sumber daya oleh pusat aglomerasi besar
  sehingga menghambat perkembangan wilayah sekitar. {[}§1, hlm.
  1115-1117{]}
\item
  \textbf{Dilaporkan sumber:} Bukti empiris mengenai operasi simultan
  kedua mekanisme masih terbatas. Studi terdahulu cenderung menganalisis
  eksternalitas positif jejaring perkotaan atau limpahan negatif
  aglomerasi secara terpisah. Operasionalisasi \emph{borrowed size} juga
  disebut belum konsisten dan ambigu. Selain itu, literatur yang dirujuk
  banyak berfokus pada sistem perkotaan Eropa, sementara bukti dari
  sistem regional non-Barat yang matang, termasuk Jepang, masih
  terbatas. {[}§1, hlm. 1116-1118{]}
\end{itemize}

\subsection{Objectives}\label{objectives-32}

\begin{itemize}
\tightlist
\item
  \textbf{Dilaporkan sumber:} Mengidentifikasi jalur yang menghubungkan
  spesialisasi industri, modal manusia, dan modal transportasi dengan
  limpahan spasial pada RTFP. {[}§1, hlm. 1117-1118{]}
\item
  \textbf{Dilaporkan sumber:} Mengestimasi secara bersama-sama efek
  \emph{borrowed size} dan \emph{agglomeration shadow} menggunakan
  kerangka ekonometrika spasial, lalu membandingkan dinamika sektor
  manufaktur dan nonmanufaktur. {[}§1, hlm. 1117-1119; Abstract, hlm.
  1115{]}
\item
  \textbf{Dilaporkan sumber:} Memeriksa apakah konsep dan teori ekonomi
  spasial yang telah dikembangkan dalam konteks lain berlaku, serta apa
  keterbatasannya, dalam konteks Jepang yang memiliki tingkat urbanisasi
  tinggi dan struktur regional polisentris. {[}§1, hlm. 1118-1119{]}
\end{itemize}

\subsection{Literature Survey}\label{literature-survey-32}

\begin{itemize}
\tightlist
\item
  \textbf{Posisi teoretis yang dilaporkan sumber:} Konsep \emph{borrowed
  size} berakar pada teori ekonomi eksternal Alonso (1973) dan
  dikembangkan lebih lanjut oleh Burger et al.~(2015). Wilayah kecil
  dapat melampaui keterbatasan fisiknya dengan memanfaatkan ukuran
  populasi, massa pasar, dan keragaman fungsi dari kota besar di
  dekatnya. Dengan demikian, ekonomi aglomerasi dapat melampaui batas
  administratif ke ruang regional yang lebih luas. {[}§1, hlm. 1115{]}
\item
  \textbf{Posisi teoretis yang dilaporkan sumber:} Jaringan antardaerah,
  infrastruktur, kereta cepat, dan bandar udara dapat memperluas
  pertukaran pengetahuan serta menghasilkan keuntungan produktivitas dan
  limpahan positif. Namun, jaringan juga dapat memusatkan sumber daya di
  kota besar dan menciptakan efek penyedotan dari wilayah pinggiran.
  Konsep \emph{agglomeration shadow} yang diperkenalkan Fujita et
  al.~(1999) menjelaskan hambatan terhadap perkembangan mandiri kota
  kecil dan menengah di sekitar pusat aglomerasi. Kedua mekanisme dapat
  bekerja secara komplementer ataupun kompetitif. {[}§1, hlm.
  1115-1117{]}
\item
  \textbf{Literatur empiris yang dirangkum artikel:} Sejumlah studi
  menggunakan konektivitas fungsional atau kinerja perkotaan untuk
  mengkaji \emph{borrowed size}, sedangkan studi lain menekankan saluran
  seperti pola komuter atau jejaring infrastruktur. Artikel menyatakan
  bahwa kajian tersebut sering memakai kerangka satu dimensi dan
  berpusat pada sistem perkotaan Eropa. {[}§1, hlm. 1117{]}
\item
  \textbf{Literatur tentang tiga saluran yang dirangkum artikel:} Untuk
  spesialisasi industri, artikel merujuk pada studi tentang spesialisasi
  okupasi, premi upah, pencocokan keterampilan, dan distrik Marshallian.
  Untuk modal manusia, artikel membahas limpahan lintas batas, investasi
  kebijakan kohesi, pembangunan kapasitas, serta akar institusional
  kesenjangan. Untuk modal transportasi, artikel merujuk pada
  ketidakmerataan waktu dan lokasi investasi modal publik serta
  perbedaan keuntungan menurut posisi jaringan dan aksesibilitas. {[}§1,
  hlm. 1118{]}
\item
  \textbf{Inferensi pengolahan:} Survei literatur dalam TXT mengarahkan
  desain penelitian dari ukuran atau konektivitas perkotaan yang
  bersifat umum menuju penguraian tiga saluran sumber daya yang dapat
  diamati, sekaligus mempertahankan perhatian pada limpahan positif dan
  negatif. Inferensi ini disusun dari uraian literatur dan kontribusi
  penulis, bukan dari sumber di luar TXT.
\end{itemize}

\subsection{Method Used}\label{method-used-32}

\begin{itemize}
\tightlist
\item
  \textbf{Dilaporkan sumber:} Penelitian menggunakan analisis
  ekonometrika spasial atas data panel prefektur dengan efek tetap
  regional. Model utama adalah SDM dan SDM-SARAR. SDM memasukkan
  \emph{spatial lag} dari variabel dependen dan variabel penjelas;
  SDM-SARAR menambahkan autokorelasi spasial pada komponen galat untuk
  menangkap ketergantungan spasial yang tidak teramati. {[}§2.2-§2.3,
  hlm. 1120-1122{]}
\item
  \textbf{Dilaporkan sumber:} RTFP menjadi variabel dependen. Variabel
  penjelasnya adalah \emph{location quotient} (LQ) sebagai ukuran
  spesialisasi industri, \emph{human capital} (HC), dan \emph{transport
  capital} (TC). Koefisien variabel penjelas yang dilag secara spasial
  dipakai artikel untuk membaca limpahan \emph{borrowed size} jika
  positif dan \emph{agglomeration shadow} jika negatif. {[}§2.1 dan
  §2.3, hlm. 1119-1122{]}
\item
  \textbf{Dilaporkan sumber:} Estimasi SDM dan SDM-SARAR dilakukan
  dengan metode \emph{maximum likelihood} (ML) pada panel seimbang dan
  efek tetap regional. Artikel menyatakan bahwa ML dirancang untuk
  menangani endogenitas dari \emph{spatial lag} variabel dependen yang
  dapat menimbulkan bias simultan. {[}§2.3, hlm. 1122{]}
\item
  \textbf{Dilaporkan sumber:} Matriks bobot spasial utama adalah matriks
  kontiguitas biner tipe \emph{queen} yang menganggap dua prefektur
  bertetangga jika berbagi batas atau satu titik. Matriks simetris,
  diagonalnya nol, lalu distandardisasi menurut baris. Karena Okinawa
  terisolasi dan menghasilkan baris nol, artikel menghubungkannya dengan
  Kagoshima sebagai daratan terdekat. {[}§2.4, hlm. 1124{]}
\item
  \textbf{Dilaporkan sumber:} Analisis dampak memakai dekomposisi LeSage
  dan Pace menjadi efek langsung, tidak langsung, dan total berdasarkan
  Persamaan (6). Efek langsung adalah dampak rata-rata perubahan
  variabel di wilayah yang sama; efek tidak langsung adalah pengaruh
  perubahan di wilayah tetangga; efek total adalah jumlah keduanya.
  {[}§3, hlm. 1122-1123{]}
\item
  \textbf{Dilaporkan sumber:} Analisis pendahuluan memakai Moran's I
  untuk autokorelasi spasial, F test untuk efek individual, dan Hausman
  test untuk memilih efek tetap atau acak. Robustness check mengestimasi
  ulang model dengan matriks \emph{K-nearest neighbors} (KNN), K = 4.
  Seluruh estimasi dilakukan di R; untuk spesifikasi SDM-SARAR penuh
  artikel menyebut fungsi \texttt{sacsarlm()} dari paket
  \texttt{spatialreg}. {[}§3.1-§3.2, hlm. 1128-1131; §4.2, hlm.
  1134-1136{]}
\item
  \textbf{Interpretasi penulis:} Kombinasi SDM dan SDM-SARAR dinilai
  memungkinkan pemodelan limpahan yang dapat diamati sekaligus
  ketergantungan laten, sedangkan struktur efek tetap dipilih karena
  karakteristik regional yang tidak teramati berkorelasi dengan variabel
  penjelas. {[}§2.2-§2.3 dan §3.2, hlm. 1120-1122, 1131{]}
\item
  \textbf{Inferensi pengolahan:} Desain ini memisahkan pengaruh lokal
  dari pengaruh yang berasal dari wilayah tetangga pada dua sektor,
  tetapi TXT tidak menyebut eksperimen atau randomisasi. Karena itu,
  review ini tidak menyebut hasilnya sebagai bukti eksperimental.
\end{itemize}

\subsection{Dataset}\label{dataset-32}

\begin{itemize}
\tightlist
\item
  \textbf{Dilaporkan sumber:} Data mencakup periode 2000-2018 dan
  bersumber dari statistik nasional Jepang yang tersedia untuk publik.
  RTFP dihitung dari \emph{Regional-Level Japan Industrial Productivity
  Database} (R-JIP) milik Research Institute of Economy, Trade, and
  Industry (RIETI). LQ juga dibangun dari nilai tambah riil menurut
  industri dan prefektur dalam R-JIP. {[}§2.5, hlm. 1124-1125{]}
\item
  \textbf{Dilaporkan sumber:} HC dihitung dari data pekerjaan dan upah
  yang dipilah menurut tingkat pendidikan, jenis kelamin, serta wilayah,
  dengan menggunakan \emph{Basic Survey on Employment Structure} dari
  Statistics Bureau of Japan dan \emph{Basic Survey on Wage Structure}
  dari Ministry of Health, Labour and Welfare. TC, sebagai proksi
  investasi infrastruktur, berasal dari \emph{Statistics on Social
  Capital Stock} milik Cabinet Office of Japan. {[}§2.1 dan §2.5, hlm.
  1119-1120, 1124-1125{]}
\item
  \textbf{Dilaporkan sumber:} R-JIP menjadi pembatas utama periode
  analisis karena 2018 adalah tahun terbaru yang tersedia di dalam basis
  data tersebut. Data modal publik lain memiliki tahun terbaru 2020,
  tetapi penelitian tetap memakai periode analisis sampai 2018. Artikel
  mengaitkan ketiadaan observasi pascapandemi dengan keterbatasan
  institusional dan pilihan memakai dataset terbuka yang dapat
  direproduksi. {[}§2.5, hlm. 1124-1125{]}
\item
  \textbf{Dilaporkan sumber:} Dalam RTFP, output diukur sebagai nilai
  tambah riil; input terdiri atas jam kerja dan stok modal riil;
  \emph{cost share} diperoleh dari data biaya tenaga kerja dan modal
  dalam R-JIP. {[}§2.5, hlm. 1125{]}
\item
  \textbf{Ringkasan deskriptif sumber:} Pada seluruh periode 2000-2018,
  rerata RTFP manufaktur adalah 0,267 dengan simpangan baku 0,221 dan
  rentang -0,436 sampai 0,902; rerata RTFP nonmanufaktur adalah -0,008
  dengan simpangan baku 0,126 dan rentang -0,212 sampai 0,619. Rerata LQ
  manufaktur adalah 1,062, sedangkan nonmanufaktur 0,982. Rerata HC
  adalah 0,776 dan rerata TC adalah 5.041.015 juta yen. {[}Table 1, hlm.
  1126{]}
\item
  \textbf{Ringkasan deskriptif sumber:} Rerata RTFP manufaktur meningkat
  dari 0,140 pada 2000-2005 menjadi 0,381 pada 2012-2018; rerata RTFP
  nonmanufaktur berada dekat nol pada ketiga subperiode. Gambar 1
  menggambarkan pola manufaktur yang tersebar dengan beberapa wilayah
  RTFP tinggi di Jepang barat dan tengah, sedangkan Gambar 2
  menggambarkan sentralisasi nonmanufaktur di sekitar Tokyo dan Osaka;
  artikel menyatakan wilayah pinggiran cenderung memiliki RTFP lebih
  rendah. {[}Table 1 dan Fig. 1-2, hlm. 1125-1127{]}
\item
  \textbf{Dilaporkan sumber:} Bagian \emph{Data sources} mencantumkan
  keempat sumber tersebut sebagai bahan yang dipublikasikan dalam bahasa
  Jepang. \emph{Data availability statement} menyatakan dataset yang
  digunakan dan/atau dianalisis tersedia dari penulis korespondensi
  berdasarkan permintaan. {[}§6 dan Data availability, hlm. 1141{]}
\end{itemize}

\subsection{Population / Sample}\label{population-sample-32}

\begin{itemize}
\tightlist
\item
  \textbf{Dilaporkan sumber:} Unit analisis adalah seluruh 47 prefektur
  Jepang, bukan subset prefektur, pada periode 2000-2018. Analisis
  dilakukan terpisah untuk sektor manufaktur dan nonmanufaktur dengan
  panel yang disebut seimbang. {[}Abstract, hlm. 1115; §2.5 dan §2.3,
  hlm. 1122, 1124{]}
\item
  \textbf{Dilaporkan sumber:} Variasi yang dianalisis berada pada
  tingkat prefektur dan industri menurut waktu. TXT tidak menyebutkan
  jumlah observasi numerik selain cakupan 47 prefektur, periode
  2000-2018, serta dua kelompok sektor.
\end{itemize}

\subsection{Dependent Variables}\label{dependent-variables-32}

\begin{itemize}
\tightlist
\item
  \textbf{Dilaporkan sumber:} Variabel dependen adalah RTFP menurut
  industri dan wilayah, yang dalam model dinyatakan sebagai vektor RTFP
  untuk sektor manufaktur atau nonmanufaktur. RTFP dihitung mengikuti
  Good et al.~(1997) dan penerapan regional oleh Otsuka (2024a, b).
  {[}§2.1 dan §2.3, hlm. 1119-1121{]}
\item
  \textbf{Dilaporkan sumber:} Formulasi RTFP mengukur tingkat
  produktivitas relatif terhadap rerata antarwilayah sekaligus perubahan
  lintas waktu. Outputnya adalah nilai tambah riil; faktor inputnya
  adalah tenaga kerja yang diukur melalui jam kerja dan modal yang
  diukur melalui stok modal riil. {[}Persamaan (1) dan §2.5, hlm. 1119,
  1125{]}
\item
  \textbf{Inferensi pengolahan:} Karena model diestimasi terpisah
  menurut sektor, hasil RTFP yang dilaporkan tidak dimaksudkan sebagai
  satu ukuran gabungan manufaktur dan nonmanufaktur.
\end{itemize}

\subsection{Results}\label{results-32}

\begin{itemize}
\tightlist
\item
  \textbf{Autokorelasi spasial, dilaporkan sumber:} Moran's I RTFP
  manufaktur pada Tabel 2 bernilai 0,45 pada 2000 dan 0,13 pada 2018,
  dengan rerata waktu 0,30 (z = 3,00; p = 0,00; \textbf{\emph{). Untuk
  nonmanufaktur, nilainya 0,18 pada 2000 dan 0,22 pada 2018, dengan
  rerata waktu 0,26 (z = 2,89; p = 0,00; }}). Uraian artikel membaca
  manufaktur sebagai pengelompokan yang lebih kuat dan nonmanufaktur
  sebagai pengelompokan yang lebih lemah, tetapi tetap memiliki
  ketergantungan spasial. {[}Table 2 dan §3.1, hlm. 1128-1129{]}
\item
  \textbf{Autokorelasi spasial, interpretasi penulis:} Artikel
  mengaitkan penurunan Moran's I manufaktur pada awal 2000-an dengan
  kemungkinan \emph{offshoring}, lalu mengaitkan pemulihan 2017-2018
  dengan kembalinya aktivitas manufaktur bernilai tambah tinggi ke
  Jepang. Kenaikan Moran's I nonmanufaktur dikaitkan dengan limpahan
  dari aglomerasi jasa maju di kawasan metropolitan Tokyo. Artikel juga
  menyebut contoh kemungkinan limpahan melalui hubungan rantai pasok
  Aichi dan Hiroshima serta penyebaran jaringan bisnis dari Tokyo ke
  Saitama, Chiba, dan Kanagawa. {[}§3.1, hlm. 1129-1130{]}
\item
  \textbf{Pemilihan efek, dilaporkan sumber:} F test menolak hipotesis
  homogenitas efek regional untuk manufaktur (72,59; p = 0,00) dan
  nonmanufaktur (125,83; p = 0,00). Hausman test juga menolak
  konsistensi model efek acak untuk manufaktur (3.812,47; p = 0,00) dan
  nonmanufaktur (102,33; p = 0,00). Model yang dipilih untuk kedua
  sektor adalah efek tetap. {[}Table 3 dan §3.2, hlm. 1131{]}
\item
  \textbf{Perbandingan model, dilaporkan sumber:} SDM-SARAR-FE
  mengungguli SDM-FE menurut AIC dan \emph{log-likelihood}. Untuk
  manufaktur, AIC dan LogL masing-masing berubah dari -1.414,4 dan 716,2
  pada SDM-FE menjadi -1.417,9 dan 718,9 pada SDM-SARAR-FE. Untuk
  nonmanufaktur, nilainya berubah dari -2.967,4 dan 1.492,7 menjadi
  -2.980,9 dan 1.500,4. LR test menolak hipotesis kesesuaian model yang
  sama pada taraf 1\% untuk kedua sektor; LM test juga menemukan
  autokorelasi spasial residual yang signifikan pada SDM. {[}Table 4-5
  dan §4.1, hlm. 1132-1133{]}
\item
  \textbf{Koefisien SDM-SARAR-FE, dilaporkan sumber:} Pada manufaktur,
  LQ lokal bernilai 0,461 (\textbf{\emph{), HC 0,752 (}}), dan TC 0,158
  (\textbf{\emph{). LQ tetangga (WLQ) bernilai -0,144 (}}), HC tetangga
  (WHC) 0,082 dan tidak signifikan, serta TC tetangga (WTC) -0,038 dan
  tidak signifikan. Rho bernilai -0,203 (\textbf{), sedangkan Lambda
  0,807 (}*). {[}Table 5, hlm. 1133{]}
\item
  \textbf{Koefisien SDM-SARAR-FE, dilaporkan sumber:} Pada
  nonmanufaktur, LQ lokal bernilai 0,374 (\textbf{\emph{), HC 0,045 dan
  tidak signifikan, serta TC 0,092 (}}). WLQ bernilai -0,376
  (\textbf{\emph{), WHC 0,051 dan tidak signifikan, serta WTC -0,079
  (}}). Rho bernilai 0,891 (\textbf{\emph{), sedangkan Lambda -0,269
  (}}). {[}Table 5, hlm. 1133{]}
\item
  \textbf{Dampak manufaktur, dilaporkan sumber:} Dari dekomposisi
  SDM-SARAR pada Tabel 8, dampak langsung LQ, HC, dan TC masing-masing
  adalah 0,455, 0,748, dan 0,157. Dampak tidak langsung masing-masing
  adalah -0,147, 0,076, dan -0,035. Dampak totalnya adalah 0,308, 0,824,
  dan 0,122. Uraian §4.3 menyatakan dampak langsung ketiga variabel
  positif dan signifikan; limpahan LQ negatif serta TC negatif dan
  signifikan; limpahan HC positif dan signifikan. {[}Table 8 dan §4.3,
  hlm. 1136-1137{]}
\item
  \textbf{Dampak nonmanufaktur, dilaporkan sumber:} Dampak langsung LQ,
  HC, dan TC masing-masing adalah 0,373, 0,045, dan 0,092. Dampak tidak
  langsungnya adalah -0,079, 0,049, dan -0,078, sedangkan dampak
  totalnya 0,294, 0,094, dan 0,014. Uraian §4.3 menyatakan dampak
  langsung LQ positif tetapi lebih lemah daripada manufaktur, dampak
  langsung HC tidak signifikan, dampak langsung TC kuat, limpahan LQ
  negatif, limpahan HC positif tetapi tidak signifikan, dan limpahan TC
  negatif serta signifikan. {[}Table 8 dan §4.3, hlm. 1136-1138{]}
\item
  \textbf{Robustness check, dilaporkan sumber:} Dengan matriks KNN, K =
  4, tanda dan arah hubungan utama pada SDM-FE tetap serupa dengan
  matriks \emph{queen} (Table 6). LM test untuk autokorelasi residual
  pada SDM-FE-KNN tidak signifikan untuk manufaktur (0,914) maupun
  nonmanufaktur (0,834). Pada SDM-SARAR-FE-KNN, manufaktur menunjukkan
  Lambda signifikan tetapi Rho tidak, sedangkan nonmanufaktur
  menunjukkan Rho signifikan tetapi Lambda tidak; variabel penjelas
  utama dan arah limpahannya tetap stabil. {[}Table 6-7 dan §4.2, hlm.
  1135-1136{]}
\item
  \textbf{Interpretasi penulis:} Penulis menilai SDM-SARAR lebih tepat
  karena menggabungkan limpahan yang dapat diamati dengan ketergantungan
  spasial laten. LQ lokal yang positif dan WLQ yang negatif dibaca
  sebagai keberadaan manfaat aglomerasi lokal bersama bayangan
  aglomerasi dari konsentrasi wilayah tetangga. {[}§4.1 dan §4.3, hlm.
  1133-1138{]}
\end{itemize}

\subsection{Findings}\label{findings-32}

\begin{itemize}
\tightlist
\item
  \textbf{Sintesis pengolahan:} Temuan utama menunjukkan koeksistensi
  keuntungan lokal dan tekanan antardaerah. Aglomerasi industri
  meningkatkan RTFP di wilayah yang sama pada kedua sektor, tetapi
  aglomerasi industri di wilayah tetangga memiliki dampak negatif
  menurut koefisien WLQ dan dampak tidak langsung LQ. Pola ini adalah
  dasar empiris yang dipakai artikel untuk mengidentifikasi
  \emph{borrowed size} dan \emph{agglomeration shadow} secara bersamaan.
  {[}Table 5 dan Table 8, hlm. 1133, 1136{]}
\item
  \textbf{Sintesis pengolahan:} Manufaktur lebih terkait dengan HC
  daripada nonmanufaktur: HC lokal dan dampak totalnya paling besar di
  manufaktur, sementara artikel melaporkan limpahan HC positif. Pada
  nonmanufaktur, dampak HC lokal tidak signifikan dan dampak totalnya
  kecil. TC meningkatkan produktivitas lokal di kedua sektor, tetapi
  limpahan TC negatif, terutama pada nonmanufaktur. {[}Table 5 dan Table
  8, hlm. 1133, 1136-1138{]}
\item
  \textbf{Interpretasi penulis:} Penulis mengaitkan hasil manufaktur
  dengan peran tenaga kerja terampil, integrasi pasar kerja, difusi
  pengetahuan, serta risiko persaingan atas input bergerak. Untuk
  nonmanufaktur, penulis menilai ketergantungan pada permintaan internal
  dan infrastruktur fisik membuat sektor ini lebih rentan terhadap
  limpahan negatif dan kompetisi antardaerah. {[}§4.3, hlm. 1137-1138{]}
\item
  \textbf{Catatan pengolahan atas TXT:} Tabel 5 melaporkan WHC tidak
  signifikan di kedua sektor, sedangkan uraian §4.3 menyatakan dampak
  tidak langsung HC pada manufaktur positif dan signifikan. Review ini
  mempertahankan kedua laporan tersebut dan tidak menggabungkannya
  menjadi satu klaim yang lebih kuat. {[}Table 5 dan §4.3, hlm. 1133,
  1136-1137{]}
\end{itemize}

\subsection{Conclusions}\label{conclusions-32}

\begin{itemize}
\tightlist
\item
  \textbf{Dilaporkan sumber:} Penelitian menyimpulkan bahwa
  interdependensi spasial memengaruhi RTFP secara berbeda pada
  manufaktur dan nonmanufaktur. Aglomerasi industri berpengaruh positif
  pada RTFP kedua sektor, sementara limpahan negatif dari wilayah
  tetangga yang ditafsirkan sebagai \emph{agglomeration shadow} juga
  signifikan. HC berkontribusi positif, terutama pada manufaktur. Dampak
  TC disebut moderat dan kadang negatif, khususnya melalui limpahan
  spasial. {[}§5, hlm. 1139-1140{]}
\item
  \textbf{Interpretasi penulis:} SDM-SARAR memberi kerangka terpadu
  untuk mengoperasionalkan \emph{borrowed size} dan \emph{agglomeration
  shadow}, mengestimasi limpahan sekaligus ketergantungan antardaerah
  yang tidak teramati, dan memasukkan perbedaan sektor dalam analisis
  eksternalitas spasial. {[}§5, hlm. 1139-1140{]}
\item
  \textbf{Inferensi pengolahan:} Dalam batas data dan spesifikasi yang
  digunakan, kesimpulan artikel lebih tepat dibaca sebagai pola hubungan
  spasial yang berbeda menurut sektor daripada sebagai ukuran tunggal
  keuntungan aglomerasi untuk semua wilayah.
\end{itemize}

\subsection{Contributions}\label{contributions-32}

\begin{itemize}
\tightlist
\item
  \textbf{Kontribusi yang diklaim penulis:} Studi ini termasuk yang
  pertama, menurut penulis, yang mengidentifikasi dan mengukur bersama
  efek \emph{borrowed size} dan \emph{agglomeration shadow} dalam satu
  kerangka ekonometrika spasial dengan SDM dan SDM-SARAR. {[}§1, hlm.
  1118{]}
\item
  \textbf{Kontribusi yang diklaim penulis:} Studi memecah eksternalitas
  spasial ke dalam tiga saluran, yaitu spesialisasi industri, HC, dan
  TC, alih-alih mengoperasionalkan \emph{borrowed size} hanya melalui
  konektivitas regional atau sentralitas fungsional. {[}§1, hlm. 1118{]}
\item
  \textbf{Kontribusi yang diklaim penulis:} Studi menyediakan bukti dari
  Jepang sebagai konteks non-Barat dengan tingkat urbanisasi tinggi dan
  struktur regional polisentris, lalu menilai secara kritis penerapan
  dan keterbatasan teori ekonomi spasial dalam kondisi institusional dan
  geografis yang berbeda. {[}§1, hlm. 1118-1119{]}
\item
  \textbf{Kontribusi metodologis yang dilaporkan penulis:} Dekomposisi
  efek lokal dan limpahan serta pemodelan ketergantungan spasial laten
  diposisikan sebagai kemajuan untuk analisis produktivitas regional
  sektoral. {[}§5, hlm. 1139-1140{]}
\end{itemize}

\subsection{Research Gap}\label{research-gap-32}

\begin{itemize}
\tightlist
\item
  \textbf{Gap yang dilaporkan sumber sebelum penelitian:} Kajian empiris
  tentang operasi simultan \emph{borrowed size} dan \emph{agglomeration
  shadow} masih terbatas; literatur sering memisahkan manfaat jejaring
  perkotaan dari limpahan negatif aglomerasi. Operasionalisasi
  \emph{borrowed size} disebut tidak konsisten dan ambigu, sedangkan
  bukti dari sistem regional non-Barat yang matang seperti Jepang masih
  kurang. Kondisi ketika \emph{borrowed size} mengungguli
  \emph{agglomeration shadow} juga disebut belum terselesaikan. {[}§1,
  hlm. 1116-1118{]}
\item
  \textbf{Cara artikel merespons gap:} Artikel mengestimasi kedua
  mekanisme bersama-sama, memisahkan tiga saluran limpahan, dan menguji
  perbedaan manufaktur serta nonmanufaktur pada panel 47 prefektur
  Jepang. {[}§1 dan §2.5, hlm. 1118-1119, 1124{]}
\item
  \textbf{Gap yang masih terbuka menurut sumber:} Mekanisme transmisi
  seperti kompetisi tenaga kerja, sentralisasi pasar, ketergantungan
  teknologi, difusi pengetahuan, mobilitas tenaga kerja, dan interaksi
  institusional belum dapat diidentifikasi sepenuhnya. Generalisasi ke
  negara atau wilayah lain dengan struktur institusional dan hierarki
  spasial berbeda juga belum diuji. {[}§5, hlm. 1140-1141{]}
\end{itemize}

\subsection{Limitations}\label{limitations-32}

\begin{itemize}
\tightlist
\item
  \textbf{Dilaporkan sumber:} Analisis terbatas pada periode prapandemi
  sampai 2018 karena ketersediaan R-JIP; akibatnya, respons
  produktivitas terhadap gangguan dan pemulihan COVID-19 belum tercakup.
  {[}§2.5 dan §5, hlm. 1124-1125, 1140-1141{]}
\item
  \textbf{Dilaporkan sumber:} Kerangka ekonometrika spasial tingkat
  makro tidak dapat mengidentifikasi sepenuhnya saluran transmisi
  spesifik, termasuk persaingan tenaga kerja, pemusatan pasar, dan
  ketergantungan teknologi. Studi kasus serta analisis institusional
  kualitatif berada di luar cakupan penelitian. {[}§5, hlm. 1140{]}
\item
  \textbf{Dilaporkan sumber:} Kasus Jepang bersifat kontekstual,
  sehingga apakah hasilnya dapat digeneralisasi ke pengaturan nasional
  atau subnasional lain masih perlu diperiksa. {[}§5, hlm. 1140-1141{]}
\item
  \textbf{Dilaporkan sumber:} SDM-SARAR dapat menghadapi masalah
  identifikasi ketika variabel penjelas yang dilag secara spasial dan
  galat spasial dibangun dengan matriks bobot yang sama, karena
  kemungkinan multikolinearitas. Penulis menyatakan mitigasi dilakukan
  dengan spesifikasi Durbin penuh, panel dengan variasi ruang-waktu, dan
  pemeriksaan stabilitas estimasi. {[}§2.3 dan §4.1, hlm. 1122, 1134{]}
\item
  \textbf{Tidak disebutkan dalam file:} Keterbatasan lain di luar
  butir-butir di atas.
\end{itemize}

\subsection{Challenges}\label{challenges-32}

\begin{itemize}
\tightlist
\item
  \textbf{Dilaporkan sumber:} Penyusunan indeks HC menghadapi ketiadaan
  indeks komposit resmi, sehingga artikel membangun ukuran berbobot upah
  dari pekerja menurut jenis kelamin dan tingkat pendidikan. {[}§2.1,
  hlm. 1119-1120{]}
\item
  \textbf{Dilaporkan sumber:} Ketersediaan data regional dan konsisten
  menurut industri menjadi alasan penggunaan R-JIP dan membatasi periode
  analisis. Isolasi geografis Okinawa juga menimbulkan baris nol dalam
  matriks kontiguitas, sehingga Okinawa dihubungkan dengan Kagoshima
  untuk menghindari masalah numerik. {[}§2.4-§2.5, hlm. 1124-1125{]}
\item
  \textbf{Dilaporkan sumber:} Struktur model menghadapi endogenitas dari
  \emph{spatial lag} variabel dependen, ketergantungan spasial yang
  tidak teramati, dan potensi kesulitan membedakan pengaruh
  \emph{spatially lagged regressors} dari galat spasial. Artikel
  menggunakan ML, SDM-SARAR, Durbin penuh, dan robustness check untuk
  menangani atau mengevaluasi persoalan tersebut. {[}§2.2-§2.3 dan
  §4.1-§4.2, hlm. 1120-1122, 1133-1136{]}
\item
  \textbf{Inferensi pengolahan:} Tantangan tersebut membuat interpretasi
  harus mempertahankan perbedaan antara koefisien lokal, koefisien
  limpahan, dan dampak hasil dekomposisi; ketiganya tidak dapat
  diperlakukan sebagai ukuran yang sama.
\end{itemize}

\subsection{Practical Implications}\label{practical-implications-32}

\begin{itemize}
\tightlist
\item
  \textbf{Interpretasi penulis:} Kebijakan spasial sebaiknya spesifik
  sektor dan adaptif terhadap konteks ruang; pendekatan seragam dinilai
  tidak memadai untuk pembangunan regional yang seimbang. {[}§4.4, hlm.
  1138{]}
\item
  \textbf{Interpretasi penulis untuk manufaktur:} Kebijakan dapat
  memperkuat klaster lokal, akumulasi keterampilan, kolaborasi
  geografis, inovasi, mobilitas tenaga kerja, dan integrasi pasar kerja.
  Namun, kebijakan diarahkan pada jejaring klaster yang tersebar dan
  bukan satu kutub industri, karena konsentrasi berlebihan serta
  limpahan negatif dari pusat tetangga dapat menekan wilayah lain.
  {[}Abstract dan §4.4, hlm. 1115, 1138-1139{]}
\item
  \textbf{Interpretasi penulis untuk TC:} Perencanaan infrastruktur
  manufaktur sebaiknya menekankan konektivitas, komplementaritas,
  aksesibilitas regional, integrasi multimoda, dan pengurangan
  pengalihan sumber daya, bukan hanya memperbesar skala investasi.
  {[}§4.4, hlm. 1138-1139{]}
\item
  \textbf{Interpretasi penulis untuk nonmanufaktur:} Kebijakan diarahkan
  pada desentralisasi kegiatan ekonomi dan penguatan ekonomi jasa
  berbasis tempat, termasuk pariwisata, budaya, dan kewirausahaan lokal.
  Penulis juga menyarankan dukungan bagi identitas ekonomi lokal,
  kapasitas institusional, kualitas hidup, layanan publik, perluasan
  pasar, investasi infrastruktur strategis, dan intervensi sisi
  permintaan. {[}§4.4, hlm. 1139{]}
\item
  \textbf{Interpretasi penulis:} Investasi di kota regional berukuran
  menengah dapat membantu berfungsi sebagai pusat sekunder dan
  mempertahankan HC. Dampak spasial yang tidak diinginkan dari
  infrastruktur metropolitan juga perlu dievaluasi agar aktivitas
  ekonomi tidak terdorong keluar dari wilayah noninti. {[}§4.4, hlm.
  1138-1139{]}
\end{itemize}

\subsection{Applications}\label{applications-32}

\begin{itemize}
\tightlist
\item
  \textbf{Dilaporkan sumber:} Temuan ditujukan untuk merancang strategi
  pembangunan regional yang dibedakan menurut sektor dan kondisi
  spasial. {[}Abstract dan §1, hlm. 1115, 1118-1119{]}
\item
  \textbf{Dilaporkan sumber:} Kerangka dampak langsung-tidak
  langsung-total digunakan artikel untuk menilai bagaimana perubahan LQ,
  HC, dan TC di suatu wilayah serta wilayah tetangga berkaitan dengan
  RTFP. Artikel memosisikan pendekatan ini sebagai dasar analitis untuk
  mengevaluasi distribusi spasial sumber daya dan kebijakan pembangunan
  regional. {[}§3, hlm. 1122-1123{]}
\item
  \textbf{Inferensi pengolahan:} Dalam batas yang dinyatakan artikel,
  aplikasi yang didukung TXT adalah perbandingan kebijakan klaster,
  tenaga kerja, dan infrastruktur menurut sektor; TXT tidak melaporkan
  implementasi kebijakan tertentu atau evaluasi dampak kebijakan setelah
  penerapan.
\end{itemize}

\subsection{Future Research
(Optional)}\label{future-research-optional-32}

\begin{itemize}
\tightlist
\item
  \textbf{Agenda yang disebut sumber:} Menelusuri mekanisme limpahan
  antardaerah melalui difusi pengetahuan, mobilitas tenaga kerja, dan
  interaksi institusional, termasuk kondisi munculnya
  \emph{agglomeration shadow} seperti kompetisi atas tenaga kerja
  terampil dan kapasitas inovasi yang tidak simetris. {[}§5, hlm.
  1140{]}
\item
  \textbf{Agenda yang disebut sumber:} Memasukkan dinamika temporal
  melalui model spasial dinamis untuk menangkap penyesuaian tertunda dan
  limpahan antarkurun waktu. {[}§5, hlm. 1140{]}
\item
  \textbf{Agenda yang disebut sumber:} Menggunakan analisis kebijakan
  berbasis simulasi untuk menghasilkan skenario kuantitatif mengenai
  perubahan struktur industri, capaian pendidikan, dan investasi
  infrastruktur. {[}§5, hlm. 1140{]}
\item
  \textbf{Agenda yang disebut sumber:} Melengkapi analisis makro dengan
  studi kasus dan analisis institusional kualitatif, serta menguji
  generalisasi pada konteks nasional atau subnasional lain. {[}§5, hlm.
  1140-1141{]}
\item
  \textbf{Agenda yang disebut sumber:} Memperluas kerangka ke periode
  setelah 2018 untuk memeriksa respons produktivitas regional terhadap
  gangguan dan pemulihan akibat COVID-19. {[}§5, hlm. 1141{]}
\end{itemize}

\subsection{Provenance Notes}\label{provenance-notes-32}

\begin{itemize}
\tightlist
\item
  \textbf{File TXT yang dipakai:}
  \texttt{/Users/mac/Documents/Mac/{[}2{]}\ Obsidian\ Vault/zettelkasten/source/journal/A17-otsuka-2025-borrowed-size-and-agglomeration-shadows-in-japan-a-spatial-econometric-analysis-of-regional-productivity-10.1007-s41685-025-00398-2.txt}.
  Tidak ada TXT lain yang dipakai atau diproses untuk assignment ini.
\item
  \textbf{Cakupan pembacaan:} Seluruh 1.415 baris TXT dibaca, mulai dari
  blok \texttt{{[}PDF\ Metadata{]}} sampai
  \texttt{Publisher\textquotesingle{}s\ Note}, termasuk teks artikel,
  Table 1-8, Fig. 1-3, bagian \emph{Data sources}, \emph{Data
  availability}, deklarasi, dan daftar referensi. Blok metadata menyebut
  PDF 30 halaman, diekstrak pada 2026-08-31 dengan
  \texttt{pdftotext\ -layout}.
\item
  \textbf{Identitas dan locator:} Judul, penulis, tahun, jurnal, halaman
  1115-1144, DOI, periode, dan cakupan wilayah diambil dari TXT. Locator
  dalam review mengikuti nomor halaman jurnal yang tercetak serta nomor
  bagian, persamaan, tabel, dan gambar sebagaimana tersedia dalam
  ekstraksi; karakter pemisah halaman dan pemenggalan kata dari
  \texttt{pdftotext\ -layout} tidak diperlakukan sebagai isi substantif.
\item
  \textbf{Batas sumber:} Review ini hanya memproses TXT A17 dan tidak
  memeriksa dataset asli, kode R, atau PDF secara independen. Pernyataan
  tentang hasil, signifikansi, interpretasi, dan implikasi kebijakan
  adalah parafrasa dari TXT, bukan verifikasi ulang perhitungan.
\item
  \textbf{Pemisahan epistemik:} Label \textbf{Dilaporkan sumber} merujuk
  pada hasil, metode, atau pernyataan eksplisit artikel;
  \textbf{Interpretasi penulis} merujuk pada pembacaan mekanisme atau
  rekomendasi yang dinyatakan penulis; \textbf{Inferensi pengolahan}
  hanya menandai kesimpulan yang disusun saat mengorganisasi isi TXT.
\item
  \textbf{Catatan konsistensi internal:} Selain perbedaan antara
  ketidaksignifikanan WHC pada Table 5 dan pernyataan limpahan HC
  manufaktur yang positif-signifikan pada §4.3, uraian §3.1
  menggambarkan Moran's I manufaktur sebagai konsisten dan signifikan,
  sedangkan Table 2 menampilkan nilai yang menurun dan p \textgreater=
  0,10 pada 2014-2017. Kedua bentuk pelaporan dipertahankan sebagai
  batas interpretasi, tanpa menambahkan penyelesaian dari luar TXT.
\item
  \textbf{Status provenance:} TXT tersedia dan memuat bahan yang
  diperlukan untuk review per-kode; tidak ada indikasi sumber hilang
  dalam file yang dibaca. Kelengkapan ekstraksi terhadap PDF asli tetap
  tidak diverifikasi secara independen.
\end{itemize}

\section{Review A18 - Improving the Economic Efficiency of Small and
Medium-Sized Cities in the Yangtze River delta: What Role for Borrowed
Size?}\label{review-a18---improving-the-economic-efficiency-of-small-and-medium-sized-cities-in-the-yangtze-river-delta-what-role-for-borrowed-size}

\subsection{Summarized Abstract}\label{summarized-abstract-33}

\begin{itemize}
\tightlist
\item
  \textbf{{[}Laporan sumber{]}} Artikel ini meneliti peran
  \texttt{borrowed\ size} dalam meningkatkan efisiensi ekonomi kota
  kecil dan menengah di aglomerasi perkotaan, dengan Yangtze River Delta
  sebagai kasus. Penulis menyatakan bahwa riset tentang aglomerasi
  perkotaan meningkat, tetapi kajian tentang kota kecil dan menengah
  masih lebih sedikit dan kurang diteorikan (Abstract, hlm. 1).
\item
  \textbf{{[}Laporan sumber{]}} Hasil utama menunjukkan bahwa
  \texttt{borrowed\ size} berpengaruh positif dan signifikan terhadap
  efisiensi ekonomi kota kecil dan menengah. Perbaikan salah alokasi
  sumber daya disebut sebagai mekanisme penting, dan artikel menyajikan
  bukti empiris multidimensi untuk pembangunan aglomerasi perkotaan yang
  terkoordinasi (Abstract, hlm. 1).
\item
  \textbf{{[}Interpretasi penulis{]}} \texttt{Borrowed\ size} dianalisis
  melalui dua dimensi, yaitu \texttt{borrowed\ performance} dan
  \texttt{borrowed\ functions}, bukan hanya melalui ukuran atau skala
  kota (Introduction, hlm. 1-2; Theoretical background, hlm. 2-3).
\item
  \textbf{{[}Inferensi pengolahan{]}} Kesimpulan empiris utama perlu
  dibaca dalam cakupan sampel 82 kota kecil dan menengah di Yangtze
  River Delta pada 2003-2020, bukan sebagai bukti tanpa batas untuk
  semua kota.
\end{itemize}

\subsection{Problem Statement}\label{problem-statement-33}

\begin{itemize}
\tightlist
\item
  \textbf{{[}Laporan sumber{]}} Kota kecil dan menengah dalam banyak
  aglomerasi menghadapi struktur industri yang tunggal, pembangunan
  infrastruktur yang tertinggal, dan kehilangan penduduk. Kondisi
  tersebut membatasi peningkatan efisiensi ekonomi dan memperlebar
  kesenjangan dengan kota-kota pusat (Introduction, hlm. 1).
\item
  \textbf{{[}Laporan sumber{]}} Sebagian kota kecil dan menengah di
  Yangtze River Delta, termasuk Jiangyin, Zhangjiagang, dan Changshu,
  disebut dapat melampaui beberapa ibu kota provinsi dalam skala ekonomi
  total dengan memanfaatkan sistem manufaktur maju dan layanan produktif
  di kota pusat seperti Shanghai. Namun, tidak semua kota kecil dan
  menengah mencapai sasaran pertumbuhan; sebagian justru masuk ke dalam
  bayang-bayang aglomerasi tanpa pertumbuhan yang berarti (Introduction,
  hlm. 1).
\item
  \textbf{{[}Laporan sumber{]}} Literatur \texttt{borrowed\ size} yang
  diringkas penulis terutama membahas distribusi spasial aktivitas
  ekonomi dan faktor-faktor yang memengaruhinya. Masih terdapat
  kekurangan kajian tentang peran \texttt{borrowed\ size} dalam alokasi
  faktor antarkota, khususnya dari perspektif kota kecil dan menengah
  (Introduction, hlm. 2).
\item
  \textbf{{[}Interpretasi penulis{]}} Masalah penelitian dirumuskan
  sebagai kebutuhan untuk menemukan cara meningkatkan efisiensi ekonomi
  kota kecil dan menengah dengan mempertimbangkan hubungan mereka dengan
  kota pusat, bukan dengan memandang setiap kota secara terisolasi
  (Introduction, hlm. 1-2).
\end{itemize}

\subsection{Objectives}\label{objectives-33}

\begin{itemize}
\tightlist
\item
  \textbf{{[}Laporan sumber{]}} Tujuan utama artikel adalah
  mengeksplorasi peran \texttt{borrowed\ size} dalam meningkatkan
  efisiensi ekonomi kota kecil dan menengah di aglomerasi perkotaan,
  menggunakan Yangtze River Delta sebagai contoh (Abstract, hlm. 1;
  Introduction, hlm. 2).
\item
  \textbf{{[}Laporan sumber{]}} Penulis menguji pengaruh
  \texttt{borrowed\ size} dari dimensi \texttt{borrowed\ performance}
  dan \texttt{borrowed\ functions}, menganalisis heterogenitas
  berdasarkan kawasan metropolitan, percepatan transportasi, dan G60
  Science and Innovation Corridor, serta menguji mekanisme melalui salah
  alokasi tenaga kerja dan modal (Introduction, hlm. 2).
\item
  \textbf{{[}Laporan sumber{]}} Hipotesis 1 menyatakan bahwa
  \texttt{borrowed\ size} dapat meningkatkan efisiensi ekonomi kota
  kecil dan menengah melalui kedua dimensinya (Research hypothesis, hlm.
  2-3). Hipotesis 2 menyatakan bahwa \texttt{borrowed\ size} dapat
  memperbaiki salah alokasi sumber daya melalui alokasi multi-pusat,
  pembangunan kota baru, dan mekanisme pasar dalam kerangka perencanaan
  neoliberal (Mechanism analysis, hlm. 3).
\item
  \textbf{{[}Laporan sumber{]}} Artikel juga menyatakan bahwa hasilnya
  dimaksudkan untuk memberi referensi bagi pembangunan ekonomi kota
  kecil dan menengah di berbagai tempat (Introduction, hlm. 1-2).
\end{itemize}

\subsection{Literature Survey}\label{literature-survey-33}

\begin{itemize}
\tightlist
\item
  \textbf{{[}Laporan sumber{]}} Landasan konsep berasal dari Alonso
  (1973), yang menggambarkan kota kecil atau kawasan metropolitan kecil
  yang dekat dengan konsentrasi penduduk lain dapat menunjukkan sebagian
  karakteristik kota yang lebih besar. Meijers dan Burger (2017) membagi
  \texttt{borrowed\ size} menjadi \texttt{borrowed\ performance} dan
  \texttt{borrowed\ functions}, sedangkan Camagni et al.~(2016)
  menekankan perlunya membedakan fungsi yang dipinjam dari kota pusat
  (Theoretical background, hlm. 2).
\item
  \textbf{{[}Laporan sumber{]}} Kelompok literatur pertama yang
  diringkas artikel memakai \texttt{borrowed\ size} untuk menjelaskan
  distribusi spasial kegiatan ekonomi, termasuk kemajuan kota periferal
  atau suburb dan lokasi bisnis dekat kawasan metropolitan. Kelompok ini
  terutama menggunakan observasi deskriptif dan studi kasus. Kelompok
  kedua mengkaji faktor yang berkaitan dengan \texttt{borrowed\ size},
  seperti fungsi perkotaan, ukuran kota, dan kedekatan geografis
  (Introduction, hlm. 2).
\item
  \textbf{{[}Laporan sumber{]}} Dalam penjelasan teoritis artikel,
  interaksi antarkota, permintaan pasar kota pusat, relokasi tenaga
  kerja, dan akses ke layanan keuangan, informasi, riset, serta
  komersial diposisikan sebagai jalur yang dapat membantu kota kecil dan
  menengah memperoleh manfaat aglomerasi. Struktur polisentris,
  infrastruktur transportasi, teknologi informasi, spesialisasi
  fungsional, pembangunan kota baru, dan mekanisme pasar juga dibahas
  sebagai konteks pendukung (Research hypothesis, hlm. 2-3).
\item
  \textbf{{[}Interpretasi penulis{]}} Artikel memosisikan
  \texttt{borrowed\ size} sebagai penghubung antara manfaat kota pusat
  dan efisiensi kota kecil dan menengah, dengan alokasi faktor antarkota
  sebagai mekanisme yang masih kurang dikaji dalam literatur yang mereka
  telaah (Introduction, hlm. 2; Mechanism analysis, hlm. 3).
\item
  \textbf{{[}Inferensi pengolahan{]}} File tidak menjelaskan prosedur
  tinjauan sistematis, kriteria pemilihan korpus, atau strategi
  pencarian literatur. Karena itu, survei literatur di sini dibaca
  sebagai telaah naratif yang disajikan dalam artikel, bukan sebagai
  hasil tinjauan sistematis yang dapat direkonstruksi dari file.
\end{itemize}

\subsection{Method Used}\label{method-used-33}

\begin{itemize}
\tightlist
\item
  \textbf{{[}Laporan sumber{]}} Model utama adalah model efek tetap
  untuk menguji pengaruh \texttt{borrowed\ size} terhadap efisiensi
  ekonomi. Persamaan (1) menempatkan DTFP sebagai variabel dependen,
  \texttt{borr\_size} sebagai variabel penjelas utama, \texttt{X}
  sebagai kumpulan variabel kontrol, serta efek tetap waktu dan kota.
  Indeks \texttt{i} menunjukkan wilayah dan \texttt{t} menunjukkan
  tahun. Koefisien \texttt{beta\_1} yang positif dan signifikan
  ditafsirkan sebagai peningkatan efisiensi; koefisien negatif dikaitkan
  dengan efek \texttt{agglomeration\ shadow} (Empirical strategy, hlm.
  3-4).
\item
  \textbf{{[}Laporan sumber{]}} \texttt{Borrowed\ performance} dihitung
  sebagai aksesibilitas populasi kota pusat yang didiskontokan oleh
  jarak, yaitu jumlah populasi kota pusat dibagi bobot jarak
  \texttt{w\_i,m} untuk setiap kota kecil dan menengah (persamaan (2),
  hlm. 4). Bobot jarak bernilai nol untuk pasangan dalam kategori yang
  sama dan berupa kebalikan jarak geografis dalam kilometer untuk
  pasangan kota pusat dan kota kecil/menengah (Variable selection, hlm.
  4).
\item
  \textbf{{[}Laporan sumber{]}} \texttt{Borrowed\ functions} dihitung
  melalui aksesibilitas fungsi tingkat tinggi kota pusat yang
  didiskontokan oleh jarak (persamaan (4), hlm. 4). Fungsi tingkat
  tinggi kota \texttt{m} diringkas dari empat kelompok industri, yaitu
  informasi, transmisi komputer dan perangkat lunak; keuangan; penyewaan
  dan layanan bisnis; serta riset ilmiah, layanan teknis, dan eksplorasi
  geologi (persamaan (3), hlm. 4).
\item
  \textbf{{[}Laporan sumber{]}} Variabel kontrol meliputi tingkat
  pembangunan ekonomi, struktur industri, tingkat upah, penanaman modal
  asing, investasi tetap, intervensi pemerintah, infrastruktur
  transportasi, dan infrastruktur informasi. Definisi operasionalnya
  dijelaskan pada bagian 3.2.3, hlm. 5.
\item
  \textbf{{[}Laporan sumber{]}} Salah alokasi modal dan tenaga kerja
  diukur dengan \texttt{MisK\_i,t} dan \texttt{MisL\_i,t} melalui
  persamaan (7) dan (8). Karena salah alokasi dapat berupa kekurangan
  maupun kelebihan alokasi, artikel menggunakan nilai absolut indeks
  tersebut; nilai absolut yang lebih besar menunjukkan salah alokasi
  yang lebih parah (Mechanism variables, hlm. 5-6).
\item
  \textbf{{[}Laporan sumber{]}} Uji ketahanan mencakup penggantian
  ukuran efisiensi dengan DEA-Malmquist dan Stochastic Frontier Analysis
  (SFA), model panel dinamis System GMM untuk menangani endogenitas,
  winsorisasi dua sisi 1\% untuk pencilan, dan penggunaan lag pertama
  \texttt{borrowed\ size}. Heterogenitas diuji melalui pemisahan kota di
  dalam dan di luar kawasan metropolitan serta interaksi dengan
  pembukaan kereta api berkecepatan tinggi dan keanggotaan G60 Science
  and Innovation Corridor (Robustness test dan Heterogeneity analysis,
  hlm. 6-9).
\item
  \textbf{{[}Interpretasi penulis{]}} Rangkaian metode tersebut
  dimaksudkan untuk menguji bukan hanya hubungan utama, tetapi juga
  dimensi fungsi dan kinerja, kondisi wilayah yang berbeda, serta jalur
  alokasi faktor yang menjelaskan hubungan tersebut.
\end{itemize}

\subsection{Dataset}\label{dataset-33}

\begin{itemize}
\tightlist
\item
  \textbf{{[}Laporan sumber{]}} Dataset berbentuk panel kota tingkat
  kabupaten dan di atasnya, mencakup 82 kota kecil dan menengah di
  Yangtze River Delta pada 2003-2020. Jumlah observasi yang dilaporkan
  adalah 1.476 (bagian 3.3 dan Table 1, hlm. 5-6).
\item
  \textbf{{[}Laporan sumber{]}} Data yang diperlukan untuk menghitung
  variabel independen bersumber dari \emph{China City Statistical
  Yearbook}. Data jarak antara kota pusat dan kota kecil/menengah
  diperoleh dengan mengolah \emph{China Terrain Database} menggunakan
  jarak garis lurus Euclidean. Data untuk variabel dependen, kontrol,
  dan mekanisme berasal dari \emph{China City Statistical Yearbook} dan
  \emph{China County Statistical Yearbook} (Data sources and descriptive
  statistics, hlm. 5).
\item
  \textbf{{[}Laporan sumber{]}} Kota pusat yang digunakan untuk konteks
  Yangtze River Delta adalah Shanghai, Nanjing, Hangzhou, Hefei, dan
  Ningbo. Kota kecil dan menengah didefinisikan menggunakan penduduk
  perkotaan permanen kurang dari 1 juta pada tahun awal 2003, mengikuti
  standar yang dirujuk artikel. Sampel terdiri atas 33 distrik kota
  tingkat prefektur dan 49 kota tingkat kabupaten (Data sources and
  descriptive statistics, hlm. 5).
\item
  \textbf{{[}Laporan sumber{]}} Table 1 melaporkan 1.476 observasi untuk
  DTFP, \texttt{borrowed\ performance}, \texttt{borrowed\ functions},
  dan seluruh variabel kontrol. Rata-rata yang tercantum untuk DTFP
  adalah 11.266, untuk \texttt{borrowed\ performance} 17.650, dan untuk
  \texttt{borrowed\ functions} 0.353; statistik minimum, maksimum, dan
  simpangan baku tersedia pada Table 1, hlm. 5-6.
\item
  \textbf{{[}Laporan sumber{]}} Artikel menyatakan bahwa data akan
  disediakan atas permintaan (Data availability, hlm. 10).
\end{itemize}

\subsection{Population / Sample}\label{population-sample-33}

\begin{itemize}
\tightlist
\item
  \textbf{{[}Laporan sumber{]}} Unit analisis adalah wilayah kota pada
  tingkat kabupaten dan di atasnya. Sampel empiris terdiri atas 82 kota
  kecil dan menengah di Yangtze River Delta, dengan 33 distrik kota
  tingkat prefektur dan 49 kota tingkat kabupaten, selama 2003-2020
  (Data sources and descriptive statistics, hlm. 5).
\item
  \textbf{{[}Laporan sumber{]}} Artikel menetapkan Shanghai, Nanjing,
  Hangzhou, Hefei, dan Ningbo sebagai kota pusat dalam penghitungan
  aksesibilitas \texttt{borrowed\ performance} dan
  \texttt{borrowed\ functions}, sedangkan sampel yang disebutkan terdiri
  atas 82 kota kecil dan menengah (Variable selection dan Data sources,
  hlm. 4-5).
\item
  \textbf{{[}Laporan sumber{]}} Kriteria kota kecil dan menengah adalah
  penduduk perkotaan permanen kurang dari 1 juta pada tahun awal 2003,
  mengikuti standar yang dirujuk artikel (Data sources and descriptive
  statistics, hlm. 5).
\item
  \textbf{{[}Inferensi pengolahan{]}} Populasi manusia, responden, atau
  peserta penelitian tidak digunakan sebagai populasi survei dalam
  desain panel kota ini. Untuk konteks tersebut: \textbf{Tidak berlaku}.
\item
  \textbf{{[}Inferensi pengolahan{]}} Klaim generalisasi paling langsung
  dari desain ini berlaku pada panel kota yang dipilih di Yangtze River
  Delta; file tidak menyediakan prosedur sampling probabilistik atau
  cakupan populasi kota di luar sampel.
\end{itemize}

\subsection{Dependent Variables}\label{dependent-variables-33}

\begin{itemize}
\tightlist
\item
  \textbf{{[}Laporan sumber{]}} Variabel dependen utama adalah efisiensi
  ekonomi perkotaan (\texttt{DTFP}). Artikel mengukurnya dengan total
  factor productivity melalui metode residual Solow (Dependent variable:
  Economic efficiency, hlm. 4-5; Table 1, hlm. 5-6).
\item
  \textbf{{[}Laporan sumber{]}} Dalam penghitungan tersebut, output
  \texttt{Y} adalah GDP aktual yang diubah menjadi GDP riil dengan tahun
  2001 sebagai tahun dasar; input tenaga kerja \texttt{L} adalah jumlah
  pekerja; input modal \texttt{K} diestimasi melalui perpetual inventory
  method; dan elastisitas output modal \texttt{alpha} diperoleh melalui
  regresi dengan merujuk Jiang dan Feng (2012) (Dependent variable, hlm.
  4-5).
\item
  \textbf{{[}Laporan sumber{]}} Persamaan (5) menghubungkan \texttt{TFP}
  dengan \texttt{Y}, \texttt{K}, \texttt{L}, dan elastisitas modal
  \texttt{alpha}. Tata letak rumus pada ekstraksi TXT tidak sepenuhnya
  terbaca, tetapi definisi komponen dan metode residual Solow dijelaskan
  dalam file (hlm. 5).
\item
  \textbf{{[}Laporan sumber{]}} Dalam regresi mekanisme, indeks salah
  alokasi modal (\texttt{MisK}) dan salah alokasi tenaga kerja
  (\texttt{MisL}) menjadi variabel yang dijelaskan untuk menilai apakah
  \texttt{borrowed\ size} mengurangi ketidaksesuaian alokasi (Mechanism
  variables dan Table 7, hlm. 5-6 dan 9).
\end{itemize}

\subsection{Results}\label{results-33}

\begin{itemize}
\tightlist
\item
  \textbf{{[}Laporan sumber{]}} Pada baseline regression (Table 2, hlm.
  6), koefisien \texttt{Borr\_size} untuk dimensi kinerja adalah 0.024
  dengan standard error 0.007 pada kolom (1), dan 0.027 dengan standard
  error 0.007 pada kolom (2). Untuk dimensi fungsi, koefisiennya 0.570
  dengan standard error 0.289 pada kolom (3), dan 0.568 dengan standard
  error 0.246 pada kolom (4). Seluruh model memuat city fixed effects
  dan year fixed effects; kolom (2) dan (4) memuat seluruh variabel
  kontrol. Jumlah sampel tiap model adalah 1.476.
\item
  \textbf{{[}Laporan sumber{]}} Teks artikel menyatakan bahwa koefisien
  \texttt{borrowed\ performance} tetap positif dan signifikan pada
  tingkat 1\% setelah kontrol dimasukkan, sedangkan hasil dimensi fungsi
  pada kolom (4) signifikan pada tingkat 5\% (Baseline regression, hlm.
  6). Nilai R2 pada Table 2 berturut-turut adalah 0.783, 0.847, 0.779,
  dan 0.841.
\item
  \textbf{{[}Laporan sumber{]}} Pada perubahan ukuran efisiensi (Table
  3, hlm. 6), koefisien \texttt{Borr\_size} tetap positif: DEA-Malmquist
  menghasilkan 0.004 dengan standard error 0.001 untuk kinerja dan 0.084
  dengan standard error 0.038 untuk fungsi; SFA menghasilkan 0.003
  dengan standard error 0.001 untuk kinerja dan 0.066 dengan standard
  error 0.037 untuk fungsi.
\item
  \textbf{{[}Laporan sumber{]}} Pada uji endogenitas dengan System GMM
  (Table 4, hlm. 7), koefisien \texttt{Borr\_size} adalah 0.004 dengan
  standard error 0.002 untuk kinerja dan 0.189 dengan standard error
  0.081 untuk fungsi. Uji AR(1) memiliki p-value 0.009 dan 0.007, AR(2)
  0.468 dan 0.500, sedangkan Hansen 0.499 dan 0.837. Artikel menyatakan
  bahwa hasil ini menunjukkan validitas instrumen dan pengaruh positif
  pada tingkat 5\%; jumlah sampel adalah 1.394.
\item
  \textbf{{[}Laporan sumber{]}} Setelah winsorisasi dua sisi 1\% pada
  seluruh variabel, koefisien \texttt{Borr\_size} pada Table 4 kolom (3)
  dan (4) masing-masing 0.025 dengan standard error 0.003 dan 0.393
  dengan standard error 0.121, keduanya dilaporkan positif dan
  signifikan pada tingkat 1\%. Jumlah sampelnya 1.476.
\item
  \textbf{{[}Laporan sumber{]}} Dengan mengganti variabel penjelas utama
  menjadi lag pertama \texttt{L.Borr\_size}, koefisiennya adalah 0.025
  dengan standard error 0.003 untuk kinerja dan 0.751 dengan standard
  error 0.132 untuk fungsi, keduanya positif dan signifikan (Table 4
  kolom (5)-(6), hlm. 7). Jumlah sampelnya 1.394.
\item
  \textbf{{[}Laporan sumber{]}} Pada heterogeneity test kawasan
  metropolitan (Table 5, hlm. 7-8), efek \texttt{Borr\_size} pada kota
  metropolitan adalah 0.035 dengan standard error 0.009 untuk kinerja
  dan 0.733 dengan standard error 0.263 untuk fungsi; keduanya
  signifikan menurut penjelasan artikel. Pada kota non-metropolitan,
  koefisiennya -0.040 dengan standard error 0.041 untuk kinerja dan
  -1.100 dengan standard error 0.768 untuk fungsi, dan tidak melewati
  uji signifikansi. Uji beda antarkelompok dilaporkan memiliki p-value
  0.000 untuk kinerja dan 0.001 untuk fungsi. Sampel subkelompok
  masing-masing 900 dan 576.
\item
  \textbf{{[}Laporan sumber{]}} Pada interaksi dengan pembukaan
  high-speed rail (Table 6, hlm. 8), interaksi \texttt{Borr\_size*HSR}
  bernilai 0.005 dengan standard error 0.002 untuk kinerja dan 0.118
  dengan standard error 0.180 untuk fungsi. Artikel melaporkan interaksi
  kinerja signifikan positif, sedangkan interaksi fungsi tidak
  signifikan.
\item
  \textbf{{[}Laporan sumber{]}} Pada interaksi dengan G60 Science and
  Innovation Corridor (Table 6, hlm. 8), \texttt{Borr\_size*SIC}
  bernilai -0.010 dengan standard error 0.004 untuk kinerja dan 0.058
  dengan standard error 0.133 untuk fungsi. Artikel melaporkan efek
  interaksi kinerja signifikan negatif, sedangkan efek interaksi fungsi
  tidak signifikan.
\item
  \textbf{{[}Laporan sumber{]}} Pada mechanism test (Table 7, hlm. 9),
  koefisien \texttt{Borr\_size} terhadap capital mismatch adalah -0.011
  dengan standard error 0.002 untuk kinerja dan -0.174 dengan standard
  error 0.086 untuk fungsi. Terhadap labor mismatch, koefisiennya -0.009
  dengan standard error 0.003 untuk kinerja dan -0.196 dengan standard
  error 0.115 untuk fungsi. Artikel menyatakan seluruh arah tersebut
  menunjukkan pengurangan salah alokasi dengan tingkat signifikansi yang
  tercantum pada tabel.
\item
  \textbf{{[}Interpretasi penulis{]}} Hasil baseline dan berbagai uji
  ketahanan dipakai penulis untuk menyatakan bahwa
  \texttt{borrowed\ size} meningkatkan efisiensi ekonomi melalui kinerja
  dan fungsi, sementara heterogenitas menunjukkan bahwa efek tersebut
  bergantung pada konteks metropolitan, transportasi, dan koridor
  inovasi (hlm. 6-10).
\end{itemize}

\subsection{Findings}\label{findings-33}

\begin{itemize}
\tightlist
\item
  \textbf{{[}Laporan sumber{]}} Penulis melaporkan tiga temuan utama:
  \texttt{borrowed\ size} meningkatkan efisiensi ekonomi kota kecil dan
  menengah pada dimensi kinerja maupun fungsi; hasil tersebut tetap
  bertahan setelah pengukuran alternatif, penanganan endogenitas,
  penghapusan pencilan, dan pengujian lag; serta dampak lebih kuat di
  kota dalam kawasan metropolitan dibandingkan kota di luar kawasan
  metropolitan (Conclusions, hlm. 9-10).
\item
  \textbf{{[}Laporan sumber{]}} Percepatan infrastruktur transportasi
  memperkuat efek \texttt{borrowed\ performance}, tetapi tidak
  menunjukkan pengaruh signifikan terhadap \texttt{borrowed\ functions}.
  Keberadaan G60 Science and Innovation Corridor memperlemah efek
  \texttt{borrowed\ performance}, sementara pengaruhnya terhadap
  \texttt{borrowed\ functions} tidak signifikan (Conclusions, hlm.
  9-10).
\item
  \textbf{{[}Laporan sumber{]}} Uji mekanisme menunjukkan bahwa
  pengurangan salah alokasi modal dan tenaga kerja merupakan jalur
  penting bagi peningkatan efisiensi ekonomi, baik pada dimensi kinerja
  maupun fungsi (Mechanism test dan Conclusions, hlm. 9-10).
\item
  \textbf{{[}Interpretasi penulis{]}} Penulis mengaitkan temuan positif
  tersebut dengan kemampuan kota kecil dan menengah untuk meminjam
  permintaan pasar, faktor produksi, layanan antara, fasilitas, dan
  fungsi dari kota pusat. Mereka mengaitkan efek negatif G60 pada
  kinerja dengan konsentrasi sumber daya inovasi dan daya tarik kota
  pusat yang dapat membentuk \texttt{concentration\ shadow} (Discussion,
  hlm. 10).
\item
  \textbf{{[}Inferensi pengolahan{]}} Temuan tidak menunjukkan efek yang
  seragam untuk semua dimensi dan konteks: arah positif paling konsisten
  pada hubungan utama, sedangkan interaksi fungsi dengan HSR dan G60
  tidak signifikan dalam hasil yang dilaporkan.
\end{itemize}

\subsection{Conclusions}\label{conclusions-33}

\begin{itemize}
\tightlist
\item
  \textbf{{[}Interpretasi penulis{]}} Berdasarkan sampel kota kecil dan
  menengah tingkat kabupaten atau lebih tinggi di Yangtze River Delta,
  penulis menyimpulkan bahwa \texttt{borrowed\ size} memfasilitasi
  peningkatan efisiensi ekonomi pada dimensi
  \texttt{borrowed\ performance} dan \texttt{borrowed\ functions}.
  Kesimpulan ini disebut tetap berlaku setelah serangkaian uji ketahanan
  (Conclusions, hlm. 9-10).
\item
  \textbf{{[}Interpretasi penulis{]}} Penulis menyimpulkan bahwa
  pengaruh lebih menonjol di kota dalam kawasan metropolitan;
  infrastruktur transportasi membantu efek kinerja, sedangkan G60
  Science and Innovation Corridor melemahkannya. Kedua faktor tersebut
  tidak dilaporkan berpengaruh signifikan terhadap efek fungsi
  (Conclusions, hlm. 9-10).
\item
  \textbf{{[}Interpretasi penulis{]}} Penulis menyimpulkan bahwa
  perbaikan salah alokasi modal dan tenaga kerja adalah mekanisme
  penting yang menghubungkan \texttt{borrowed\ size} dengan efisiensi
  ekonomi (Conclusions, hlm. 9-10).
\item
  \textbf{{[}Inferensi pengolahan{]}} Kesimpulan di atas adalah
  kesimpulan internal artikel untuk desain, periode, dan wilayah yang
  diteliti; file tidak menyediakan dasar untuk memperluasnya melampaui
  cakupan tersebut.
\end{itemize}

\subsection{Contributions}\label{contributions-33}

\begin{itemize}
\tightlist
\item
  \textbf{{[}Interpretasi penulis{]}} Kontribusi pertama yang diklaim
  adalah menggeser perhatian dari sekadar efek limpahan antarkota ke
  mekanisme alokasi faktor antarkota, dengan menganalisis
  \texttt{borrowed\ size} melalui dimensi kinerja dan fungsi
  (Introduction, hlm. 2).
\item
  \textbf{{[}Interpretasi penulis{]}} Kontribusi kedua adalah menguji
  heterogenitas berdasarkan kawasan metropolitan, percepatan
  transportasi, dan G60 Science and Innovation Corridor, lalu memeriksa
  mekanisme melalui salah alokasi tenaga kerja dan modal (Introduction,
  hlm. 2).
\item
  \textbf{{[}Interpretasi penulis{]}} Dalam Discussion, artikel
  menyatakan bahwa hasilnya memperluas kajian tentang skala pinjaman,
  spesialisasi perkotaan, dan mobilitas faktor dengan menyoroti masalah
  praktis salah alokasi modal dan tenaga kerja (Discussion, hlm. 10).
\item
  \textbf{{[}Inferensi pengolahan{]}} Kontribusi empiris yang dapat
  diidentifikasi dari file adalah penggabungan pengukuran aksesibilitas
  populasi dan fungsi tingkat tinggi dengan outcome efisiensi ekonomi
  dalam satu panel kota; ini merupakan ringkasan atas desain yang
  dilaporkan, bukan klaim tambahan dari luar file.
\end{itemize}

\subsection{Research Gap}\label{research-gap-33}

\begin{itemize}
\tightlist
\item
  \textbf{{[}Laporan sumber{]}} Artikel menyatakan bahwa kajian tentang
  kota kecil dan menengah di dalam aglomerasi perkotaan masih lebih
  sedikit dan kurang diteorikan dibandingkan kajian aglomerasi perkotaan
  secara umum (Abstract dan Introduction, hlm. 1).
\item
  \textbf{{[}Laporan sumber{]}} Penulis secara khusus mengidentifikasi
  kekurangan literatur mengenai peran \texttt{borrowed\ size} dalam
  alokasi faktor antarkota, terutama dari perspektif kota kecil dan
  menengah (Introduction, hlm. 2).
\item
  \textbf{{[}Interpretasi penulis{]}} Artikel menempatkan pengujian
  pengaruh \texttt{borrowed\ performance}, \texttt{borrowed\ functions},
  dan mekanisme salah alokasi sebagai respons terhadap celah tersebut.
\item
  \textbf{{[}Laporan sumber{]}} Celah penelitian lanjutan yang
  dinyatakan secara eksplisit setelah studi ini: \textbf{Tidak
  disebutkan dalam file}.
\end{itemize}

\subsection{Limitations}\label{limitations-33}

\begin{itemize}
\tightlist
\item
  \textbf{{[}Laporan sumber{]}} Tidak terdapat bagian khusus berjudul
  \texttt{Limitations} dalam file. Batasan yang disebutkan secara
  langsung adalah bahwa sampel 82 kota dipilih dengan mempertimbangkan
  ketersediaan dan konsistensi data (Data sources and descriptive
  statistics, hlm. 5).
\item
  \textbf{{[}Laporan sumber{]}} Penulis menyatakan bahwa banyak faktor
  memengaruhi efisiensi ekonomi kota dan sebagian sulit diukur karena
  keterbatasan ketersediaan data. Untuk mengurangi potensi bias variabel
  yang dihilangkan, artikel memasukkan sejumlah variabel kontrol dan
  melakukan uji endogenitas, tetapi file tidak menyatakan bahwa seluruh
  faktor telah terukur (Processing endogeneity, hlm. 7).
\item
  \textbf{{[}Laporan sumber{]}} Data penelitian dinyatakan tersedia atas
  permintaan, bukan disertakan sebagai data terbuka di dalam file review
  sumber (Data availability, hlm. 10).
\item
  \textbf{{[}Inferensi pengolahan{]}} Penggunaan jarak garis lurus
  Euclidean, klasifikasi kota berdasarkan tahun awal 2003, dan batas
  wilayah sampel merupakan pilihan desain yang perlu dipertahankan saat
  membaca hasil; file tidak menguji secara eksplisit dampak pilihan
  tersebut sebagai keterbatasan.
\end{itemize}

\subsection{Challenges}\label{challenges-33}

\begin{itemize}
\tightlist
\item
  \textbf{{[}Laporan sumber{]}} Konteks yang dihadapi kota kecil dan
  menengah mencakup struktur industri yang tunggal, infrastruktur yang
  tertinggal, dan kehilangan penduduk. Kota-kota tersebut juga berisiko
  mengalami bayang-bayang aglomerasi ketika manfaat kota pusat tidak
  menghasilkan pertumbuhan yang cukup (Introduction, hlm. 1).
\item
  \textbf{{[}Laporan sumber{]}} Penulis mengidentifikasi salah alokasi
  modal dan tenaga kerja sebagai kendala terhadap efisiensi ekonomi.
  Mereka juga membahas kemungkinan hubungan kausal dua arah antara
  \texttt{borrowed\ size} dan efisiensi ekonomi, faktor yang sulit
  diukur karena keterbatasan data, pencilan, dan kemungkinan efek waktu
  tertunda (Mechanism analysis dan Robustness test, hlm. 3 dan 7).
\item
  \textbf{{[}Interpretasi penulis{]}} Konsentrasi sumber daya inovasi di
  kota yang lebih besar dan daya tarik kota pusat di dalam G60 Science
  and Innovation Corridor dipandang dapat menciptakan
  \texttt{concentration\ shadow} dan membatasi efek
  \texttt{borrowed\ performance} kota kecil dan menengah (Heterogeneity
  analysis dan Discussion, hlm. 8 dan 10).
\item
  \textbf{{[}Inferensi pengolahan{]}} Tantangan analitis artikel bukan
  hanya mengukur kedekatan dengan kota pusat, tetapi juga membedakan
  efek kinerja dari efek fungsi dan menguji apakah hubungan tersebut
  bertahan setelah koreksi yang dilaporkan penulis.
\end{itemize}

\subsection{Practical Implications}\label{practical-implications-33}

\begin{itemize}
\tightlist
\item
  \textbf{{[}Interpretasi penulis{]}} Pemerintah disarankan memperkuat
  dukungan kebijakan untuk efek \texttt{borrowed\ size}, termasuk
  kebijakan fiskal, dana dukungan industri, subsidi fiskal, kolaborasi
  industri lintas wilayah, serta platform kerja sama dan berbagi
  informasi antarperusahaan (Policy recommendations, hlm. 10).
\item
  \textbf{{[}Interpretasi penulis{]}} Penulis menyarankan perluasan efek
  radiasi kawasan metropolitan dan perbaikan transportasi: integrasi
  infrastruktur di kawasan metropolitan, peningkatan konektivitas antara
  kota kecil/menengah dan kota pusat, serta investasi transportasi
  antarkota.
\item
  \textbf{{[}Interpretasi penulis{]}} Dalam pembangunan G60 Science and
  Innovation Corridor, penulis menyarankan agar konsentrasi sumber daya
  inovasi yang berlebihan di kota inti dihindari supaya efek
  \texttt{borrowed\ size} kota kecil dan menengah tidak melemah.
\item
  \textbf{{[}Interpretasi penulis{]}} Untuk mengoptimalkan salah
  alokasi, artikel merekomendasikan penguatan pasar keuangan regional
  dan inovasi keuangan untuk mengatasi kesulitan serta biaya pembiayaan,
  disertai perbaikan mekanisme mobilitas tenaga kerja dan reformasi
  sistem registrasi rumah tangga untuk mengurangi hambatan perpindahan
  lintas kota (Policy recommendations, hlm. 10).
\item
  \textbf{{[}Inferensi pengolahan{]}} Implikasi praktis tersebut
  sebaiknya dipahami sebagai rekomendasi penulis yang diturunkan dari
  hasil sampel Yangtze River Delta; file tidak menguji dampak
  implementasi kebijakan-kebijakan itu secara langsung.
\end{itemize}

\subsection{Applications}\label{applications-33}

\begin{itemize}
\tightlist
\item
  \textbf{{[}Laporan sumber{]}} Artikel menerapkan ukuran
  \texttt{borrowed\ performance} dan \texttt{borrowed\ functions} untuk
  mengevaluasi hubungan kota kecil dan menengah dengan lima kota pusat
  dalam panel Yangtze River Delta, serta menerapkan uji heterogenitas
  pada kawasan metropolitan, HSR, dan G60 Science and Innovation
  Corridor (Methodology and data serta Heterogeneity analysis, hlm.
  3-9).
\item
  \textbf{{[}Laporan sumber{]}} Penulis menggunakan hasil tersebut
  sebagai dasar rekomendasi untuk perencanaan aglomerasi perkotaan,
  perencanaan metropolitan, pembangunan infrastruktur transportasi,
  kolaborasi industri, dan pengurangan salah alokasi modal serta tenaga
  kerja (Conclusions and policy recommendations, hlm. 9-10).
\item
  \textbf{{[}Laporan sumber{]}} Tujuan penerapan yang dinyatakan penulis
  adalah menyediakan referensi bagi pembangunan ekonomi kota kecil dan
  menengah; artikel tidak melaporkan implementasi kebijakan atau
  evaluasi program setelah penelitian (Introduction, hlm. 1-2; Policy
  recommendations, hlm. 10).
\item
  \textbf{{[}Inferensi pengolahan{]}} Aplikasi di luar wilayah, periode,
  definisi kota, dan sumber data yang digunakan artikel: \textbf{Tidak
  disebutkan dalam file}.
\end{itemize}

\subsection{Future Research
(Optional)}\label{future-research-optional-33}

\begin{itemize}
\tightlist
\item
  \textbf{{[}Laporan sumber{]}} Arah penelitian masa depan atau agenda
  riset lanjutan yang dinyatakan secara eksplisit: \textbf{Tidak
  disebutkan dalam file}.
\end{itemize}

\subsection{Provenance Notes}\label{provenance-notes-33}

\begin{itemize}
\tightlist
\item
  \textbf{Sumber utama:}
  \texttt{source/journal/A18-zhao-wang-li-2025-improving-the-economic-efficiency-of-small-and-medium-sized-cities-in-the-yangtze-river-delta-what-role-for-borrowed-size-10.1016-j.cities.2025.105954.txt}.
\item
  \textbf{Verifikasi nama file:} penelusuran prefiks \texttt{A18} di
  bawah \texttt{source/} menemukan satu TXT aktual, yaitu file di atas.
  TXT tidak diubah.
\item
  \textbf{Basis evidence:} seluruh 1.021 baris TXT dibaca, termasuk blok
  \texttt{{[}PDF\ Metadata{]}} pada baris 1-28 dan blok
  \texttt{{[}Extracted\ Text{]}} mulai baris 29. TXT menyatakan sumber
  PDF memiliki 12 halaman, diekstrak pada 2026-08-31 dengan
  \texttt{pdftotext\ -layout}, dan memuat artikel \emph{Cities} 162
  (2025) 105954 dengan DOI 10.1016/j.cities.2025.105954.
\item
  \textbf{Identitas yang tersedia dalam file:} judul artikel, Weishu
  Zhao, Zhe Wang, dan Jiani Li, jurnal \emph{Cities}, volume 162, tahun
  2025, nomor artikel 105954, serta DOI tercantum pada blok metadata dan
  halaman awal.
\item
  \textbf{Locator:} review mempertahankan locator halaman artikel,
  bagian, persamaan, dan Table 1-7 apabila tersedia. Nomor halaman
  merujuk nomor halaman artikel yang muncul di footer ekstraksi, bukan
  nomor baris TXT.
\item
  \textbf{Label evidence:} \texttt{{[}Laporan\ sumber{]}} merangkum atau
  memparafrasekan isi yang dilaporkan artikel;
  \texttt{{[}Interpretasi\ penulis{]}} merujuk pada penjelasan, klaim
  kontribusi, kesimpulan, atau rekomendasi yang dinyatakan penulis;
  \texttt{{[}Inferensi\ pengolahan{]}} adalah pembacaan pengolah atas
  batas cakupan atau struktur bukti dan bukan klaim tambahan dari
  sumber.
\item
  \textbf{Batas evidence:} ekstraksi \texttt{pdftotext\ -layout} memuat
  artefak tata letak pada rumus, pemenggalan kolom tabel, dan simbol
  signifikansi. Rumus (5) tidak seluruhnya terbaca secara utuh dalam
  TXT; karena itu review hanya menyatakan komponen yang dijelaskan dan
  tidak merekonstruksi bagian rumus yang hilang. Table 5 juga
  menampilkan notasi \texttt{0.035****} dalam ekstraksi, sementara
  catatan tabel menjelaskan tiga tingkat tanda signifikansi; notasi
  tersebut tidak dinormalisasi dalam review.
\item
  \textbf{Batas akses:} evidence yang dibaca adalah seluruh teks turunan
  TXT dari PDF. Review ini tidak menambahkan isi dari sumber eksternal,
  tidak menganggap data tersedia atas permintaan sebagai data yang telah
  diperiksa, dan tidak mengisi agenda future research dengan dugaan.
\end{itemize}

\section{Review A19 - Chen et
al.~(2024)}\label{review-a19---chen-et-al.-2024}

Status: Review selesai dan terverifikasi berbasis teks lengkap hasil
ekstraksi PDF.

\subsection{Identitas Artikel}\label{identitas-artikel-6}

\begin{itemize}
\tightlist
\item
  \textbf{Kode:} A19.
\item
  \textbf{Judul:} \emph{Population agglomeration, borrowed size and
  urban economic growth in China}.
\item
  \textbf{Penulis:} Le Chen, Xi Wang, Yuanwei Hao, Jiangbin Yin, dan
  Dongsheng Chen.
\item
  \textbf{Jurnal:} \emph{Cities}, volume 155 (2024), artikel 105491.
\item
  \textbf{DOI:} 10.1016/j.cities.2024.105491.
\item
  \textbf{Unit dan periode utama:} 285 kota tingkat prefektur di China,
  2005-2020.
\item
  \textbf{Locator identitas:} hlm. 1, blok judul dan abstrak; hlm. 20,
  daftar pustaka.
\end{itemize}

\subsection{Summarized Abstract}\label{summarized-abstract-34}

\textbf{Laporan sumber.} Artikel membangun kerangka teoritis dan model
pengukuran untuk menguji hubungan aglomerasi penduduk dengan pertumbuhan
ekonomi perkotaan. Kerangka itu menggabungkan aglomerasi intrakota
dengan aglomerasi antarkota, terutama borrowed size dan agglomeration
shadow. Pengujian memakai data 285 kota tingkat prefektur di China pada
2005-2020. Data saling diverifikasi dengan lebih dari 500 buku tahunan
statistik provinsi dan kota, 4.560 buletin statistik, unsur alam, serta
citra cahaya malam DMSP/OLS dan NPP/VIIRS yang dipadankan. (Hlm. 1,
abstrak.)

\textbf{Temuan yang dilaporkan sumber.} Aglomerasi penduduk intrakota
berpengaruh positif terhadap pertumbuhan ekonomi perkotaan, dan pengaruh
ini dimoderasi oleh faktor spasial serta kebijakan. Efek borrowed size
lebih kuat daripada agglomeration shadow di antara kota-kota China.
Aglomerasi penduduk intrakota juga mempunyai limpahan spasial positif;
tingkat aglomerasi yang tinggi di kota tetangga meningkatkan pertumbuhan
ekonomi kota yang diamati. Koefisien aglomerasi intrakota turun sekitar
0,1 setelah unsur aglomerasi antarkota dimasukkan. (Hlm. 1, abstrak.)

\textbf{Interpretasi penulis.} Hasil tersebut dipakai untuk menyatakan
bahwa borrowed size mendominasi agglomeration shadow, meskipun efek
bayangan tetap ada. Penulis juga membaca kerangka gabungan
intrakota-antarkota sebagai penjelasan yang lebih lengkap atas hubungan
aglomerasi penduduk dan pertumbuhan ekonomi. (Hlm. 4-5, Bagian 2.2-3.2;
hlm. 18, Bagian 5.2.)

\textbf{Inferensi pengolahan.} Bukti paling konsisten dalam TXT adalah
asosiasi positif aglomerasi intrakota dengan proksi pendapatan lahan
perkotaan per kapita. Hasil untuk PDB per kapita lebih lemah dan pada
beberapa model berubah tanda, sehingga klaim positif perlu dibaca
bersama perbedaan antarproksi dan antarspesifikasi model. Panel
observasional, GMM, dan IV memperkuat pemeriksaan ketahanan, tetapi TXT
tidak menunjukkan eksperimen atau kontrafaktual yang mengubah hasil
menjadi bukti kausal langsung.

\subsection{Problem Statement}\label{problem-statement-34}

\textbf{Laporan sumber.} Studi terdahulu sering hanya membahas
aglomerasi intrakota atau salah satu bentuk aglomerasi antarkota.
Kerangka dan model empiris yang sekaligus mempertimbangkan aglomerasi
intrakota, borrowed size, dan agglomeration shadow dinilai masih kurang.
Artikel juga menyebut dua persoalan metodologis: endogenitas antara
aglomerasi penduduk dan pertumbuhan ekonomi, serta persoalan pengumpulan
dan verifikasi data statistik China. (Hlm. 2-4, pendahuluan dan Bagian
2.2.)

\textbf{Laporan sumber.} Literatur belum mencapai konsensus mengenai
apakah pengaruh antarkota lebih banyak berupa manfaat borrowed size atau
tekanan agglomeration shadow. Persoalan itu dianggap penting untuk
kota-kota China yang padat dan mengalami urbanisasi cepat, termasuk
karena pertumbuhan dapat melemah ketika ukuran atau kepadatan kota
melewati titik tertentu. (Hlm. 1-4, pendahuluan dan Bagian 2.1-2.2.)

\textbf{Interpretasi penulis.} Penulis merumuskan masalah substantif
sebagai pertanyaan apakah aglomerasi penduduk mendorong atau menghambat
pertumbuhan, bagaimana limpahan spasial bekerja, dan bentuk efek
antarkota mana yang lebih dominan. (Hlm. 2, pendahuluan.)

\textbf{Inferensi pengolahan.} Masalah artikel terdiri atas pengukuran
efek lokal dan atribusi mekanisme antarkota. Model dapat menguji
hubungan statistik dengan indikator borrowed size dan shadow, tetapi
pengukuran mekanisme secara langsung memerlukan bukti arus, jaringan,
atau proses yang tidak sepenuhnya tersedia dalam TXT.

\subsection{Objectives}\label{objectives-34}

\textbf{Laporan sumber.} Artikel bertujuan untuk:

\begin{itemize}
\tightlist
\item
  menguji apakah aglomerasi penduduk intrakota berkontribusi secara
  signifikan terhadap pertumbuhan ekonomi perkotaan;
\item
  menguji apakah aglomerasi penduduk memiliki limpahan spasial;
\item
  menilai apakah aglomerasi penduduk kota tetangga lebih kuat
  menghasilkan borrowed size atau agglomeration shadow;
\item
  menyusun kerangka teoritis dan model pengukuran yang menggabungkan
  aglomerasi intrakota serta antarkota;
\item
  mengukur aglomerasi ekonomi dan pertumbuhan ekonomi dengan seperangkat
  indikator, bukan satu indikator tunggal; dan
\item
  memeriksa ketahanan hasil terhadap endogenitas, perbedaan indikator
  statistik, ketidakakuratan GDP, objek studi, dan periode 2020. (Hlm.
  2-4, pendahuluan dan Bagian 3; hlm. 14-18, Bagian 4.3-5.2.)
\end{itemize}

\textbf{Hipotesis yang dilaporkan sumber.} Hipotesis 1 menyatakan bahwa
aglomerasi penduduk dapat mendorong pertumbuhan ekonomi, walaupun
kemacetan dapat menghambatnya. Hipotesis 2 menyatakan bahwa aglomerasi
penduduk memiliki limpahan spasial positif, aglomerasi tinggi di kota
tetangga meningkatkan pertumbuhan kota, dan borrowed size lebih kuat
daripada agglomeration shadow. (Hlm. 4, Bagian 3.1.)

\subsection{Literature Survey}\label{literature-survey-34}

\textbf{Laporan sumber.} Tinjauan pustaka dimulai dari indikator
aglomerasi intrakota. Spesialisasi industri dan keragaman disebut
sebagai indikator klasik, sedangkan ukuran penduduk, kepadatan penduduk,
dan keanggotaan dalam aglomerasi perkotaan digunakan untuk menangkap
tingkat aglomerasi keseluruhan. Literatur yang ditinjau banyak
melaporkan hubungan positif aglomerasi dengan pertumbuhan, tetapi
sebagian studi menemukan hubungan berbentuk U terbalik atau menunjukkan
ukuran serta kepadatan optimal. (Hlm. 2-4, Bagian 2.1.)

\textbf{Laporan sumber.} Artikel meninjau bukti bahwa lokasi, kebijakan,
modal, sumber daya alam, dan kondisi wilayah dapat mengubah hubungan
aglomerasi-pertumbuhan. Contoh yang dibahas meliputi perbedaan wilayah
timur, tengah, dan barat China; perbedaan kota ibu kota dan nonibu kota;
serta perbedaan kota pesisir dan nonpesisir. (Hlm. 3-4, Bagian 2.1.)

\textbf{Laporan sumber.} Pada aglomerasi antarkota, borrowed size
dijelaskan sebagai kondisi ketika kota kecil memperoleh akses terhadap
pasar tenaga kerja, modal, layanan, dan keuntungan aglomerasi kota besar
tanpa menanggung seluruh masalah kepadatan. Agglomeration shadow
dijelaskan sebagai kondisi ketika kota besar menarik tenaga kerja,
modal, dan aktivitas dari wilayah tetangga. Studi di Eropa, Amerika
Serikat, Swedia, Polandia, dan berbagai wilayah China dilaporkan
menghasilkan bukti yang tidak seragam. (Hlm. 1-4, Bagian 1 dan 2.2.)

\textbf{Laporan sumber.} Sejumlah studi yang dikutip menemukan borrowed
size lebih besar daripada shadow, sedangkan studi lain menemukan shadow
pada wilayah sekitar Beijing-Tianjin-Hebei, Beijing, dan
Guangzhou-Foshan. Artikel juga mengutip temuan bahwa efek positif atau
negatif dapat berbeda menurut waktu, ruang, ukuran kota, dan
konektivitas. (Hlm. 1-4, Bagian 1 dan 2.2.)

\textbf{Interpretasi penulis.} Perbedaan hasil terdahulu menjadi alasan
untuk menguji kedua mekanisme secara bersamaan di seluruh kota tingkat
prefektur China, bukan hanya pada satu wilayah metropolitan atau satu
bentuk aglomerasi. (Hlm. 3-5, Bagian 2.2-3.1.)

\textbf{Inferensi pengolahan.} Tinjauan pustaka dalam TXT adalah
tinjauan naratif yang disusun untuk mendukung model penelitian. File
tidak menyebutkan strategi pencarian, kriteria inklusi, rentang korpus,
atau prosedur penilaian mutu studi terdahulu. Karena itu, kelengkapan
survei literatur tidak dapat diverifikasi dari TXT saja.

\subsection{Method Used}\label{method-used-34}

\textbf{Laporan sumber.} Penelitian menggunakan pendekatan kuantitatif
dengan panel kota tingkat prefektur China pada 2005-2020. Kerangka
teoritisnya berangkat dari fungsi produksi Cobb-Douglas dan membedakan:

\begin{itemize}
\tightlist
\item
  aglomerasi intrakota yang diukur melalui log kepadatan penduduk, log
  ukuran kota awal, dan dummy keanggotaan dalam aglomerasi perkotaan;
\item
  aglomerasi antarkota yang diwakili aglomerasi kota tetangga dan dummy
  kedekatan dengan kota berpopulasi tinggi;
\item
  dua proksi pertumbuhan ekonomi, yaitu pendapatan lahan perkotaan per
  kapita dan PDB per kapita; serta
\item
  variabel kendali berupa modal, modal manusia, struktur pekerjaan
  manufaktur, dan suhu tahunan rata-rata. (Hlm. 4-7, Bagian 3.1-3.3.)
\end{itemize}

\textbf{Laporan sumber.} Persamaan dasar diestimasi dengan regresi panel
efek tetap atau regresi \texttt{xi:reg} sesuai karakter indikator.
Hausman test dengan \texttt{p\ =\ 0,000} mendukung efek tetap pada model
yang diuji. Model panel dinamis dengan lag orde pertama dari tingkat
pendapatan diestimasi menggunakan GMM untuk menangani endogenitas. (Hlm.
5-9, Bagian 3.2-4.1 dan Tabel 2.)

\textbf{Laporan sumber.} Limpahan spasial diuji melalui spatial
autoregressive model (SAR), spatial error model (SEM), dan spatial
Durbin model (SDM). Karena ukuran kota awal dan dummy keanggotaan
aglomerasi tidak berubah sepanjang waktu, regresi panel spasial hanya
menggunakan kepadatan penduduk sebagai indikator aglomerasi. Persamaan
tambahan memakai dummy kedekatan dengan kota yang masuk 5 persen teratas
berdasarkan populasi dan interaksi dummy tersebut dengan indikator
aglomerasi lokal untuk membedakan borrowed size dan shadow. (Hlm. 5-6
dan 11-13, Bagian 3.2, Bagian 4.2, dan Tabel 5-7.)

\textbf{Laporan sumber.} Heterogenitas diuji melalui interaksi dengan
lokasi timur, tengah, atau barat; status ibu kota provinsi atau kota
terencana; status kota pesisir terbuka; serta zona iklim tropis atau
temperate. Uji robustitas memakai 2SLS dengan RDLS dan kepadatan
jaringan sungai sebagai instrumen kepadatan penduduk, serta kemiringan
sebagai instrumen ukuran kota awal. Karena ukuran kota awal tidak
berubah sepanjang waktu, uji IV menggantinya dengan penduduk tetap pada
2005-2020. Instrumen dibentuk bersama interaksi unsur geografi fisik dan
kebalikan nilai tukar sebagai guncangan eksternal. (Hlm. 9-10 dan 13-14,
Bagian 3.3 dan 4.3.)

\textbf{Laporan sumber.} Pemeriksaan tambahan dilakukan pada data
distrik municipal untuk tahun 2000, 2010, dan 2020; GDP yang
dimodifikasi dengan cahaya malam; 35 ibu kota provinsi dan kota
terencana; serta 285 kota dengan tahun 2020 dikeluarkan. Tujuannya ialah
menguji kemungkinan ketidakcocokan antara statistik tingkat kota,
penduduk, dan area terbangun, serta kemungkinan distorsi tahun COVID-19
dan pengaruh kota besar. (Hlm. 14-16, Bagian 4.3 dan Tabel 9-11.)

\textbf{Inferensi pengolahan.} Rancangan ini adalah studi observasional
panel dengan beberapa strategi koreksi dan uji ketahanan, bukan
eksperimen. TXT tidak menjelaskan matriks bobot spasial secara rinci,
sehingga konstruksi kedekatan antarkota dan reproduksi penuh model
spasial tidak dapat dipastikan dari file saja.

\subsection{Dataset}\label{dataset-34}

\textbf{Laporan sumber.} Dataset utama dan penggunaannya adalah sebagai
berikut:

\begin{itemize}
\tightlist
\item
  \textbf{Statistik GDP:} data GDP berasal dari \emph{China City
  Statistical Yearbook} 2006-2021. Data tingkat kota tahun 2017 yang
  tidak tersedia sebagai data seluruh kota diganti dengan data dari
  buletin statistik masing-masing kota. GDP nominal diubah menjadi GDP
  riil dengan deflator dari \emph{China Statistical Yearbook} 2006-2021.
  (Hlm. 6-7, Bagian 3.3.)
\item
  \textbf{Penduduk:} ukuran penduduk berupa penduduk tetap yang
  dikumpulkan dari lebih dari 500 buku tahunan statistik pada 35
  provinsi untuk 2006-2021 dan diverifikasi dengan 4.560 buletin
  statistik kota. Data sensus 2010 dan 2020 digunakan sebagai tolok
  ukur. Untuk Liaoning, Jilin, dan Heilongjiang, angka penduduk tetap
  dihitung dari penduduk rumah tangga ditambah penduduk sementara. (Hlm.
  6-7, Bagian 3.3.)
\item
  \textbf{Area terbangun:} data DMSP/OLS digunakan untuk 2005-2012 dan
  NPP/VIIRS untuk 2013-2020. Data dikoreksi lintas sensor, dinormalisasi
  secara temporal dan spasial, disaring dari sumber cahaya nonperkotaan,
  lalu dicocokkan dengan data penggunaan lahan 30 meter tahun 2010 dan
  2015. (Hlm. 5-7, Bagian 3.2-3.3 dan Gambar 2.)
\item
  \textbf{Cahaya malam untuk koreksi GDP:} data cahaya malam 285 kota
  tahun 2005-2020 digunakan untuk memodifikasi data GDP statistik
  melalui fungsi kuadratik antara GDP dan nilai cahaya malam. (Hlm. 7
  dan 15-16, Bagian 3.3 dan Tabel 10.)
\item
  \textbf{Modal:} stok modal riil dihitung dengan metode perpetual
  inventory menggunakan tingkat depresiasi seragam 9,6 persen. (Hlm. 6,
  Bagian 3.3.)
\item
  \textbf{Modal manusia:} jumlah aplikasi paten per 10.000 penduduk
  diperoleh dari buletin statistik kota dan China National Intellectual
  Property Administration, lalu diverifikasi dengan buku statistik
  provinsi dan buku sains-teknologi. (Hlm. 6, Bagian 3.3.)
\item
  \textbf{Suhu:} suhu tahunan rata-rata diperoleh dari Shanghai Bancheng
  Data Technology Company dan dibandingkan dengan data lingkungan China
  pada 10 persen titik data yang dipilih acak; perbedaan yang dilaporkan
  sebesar 2 derajat C-1,2 derajat C. (Hlm. 6-7, Bagian 3.3.)
\item
  \textbf{Instrumen:} RDLS dan kemiringan berasal dari data model
  elevasi digital 30 meter Geospatial Data Cloud, sedangkan kepadatan
  jaringan sungai berasal dari Chinese Scientific Database. (Hlm. 7,
  Bagian 3.3.)
\end{itemize}

\textbf{Laporan sumber.} Statistik deskriptif pada Tabel 1 menggunakan
4.560 observasi untuk sebagian besar variabel. Nilai minimum-maksimum,
rerata, dan simpangan baku masing-masing variabel dicantumkan pada tabel
tersebut, termasuk log pendapatan lahan perkotaan per kapita, log GDP
per kapita, log kepadatan penduduk, log ukuran kota awal, dummy
aglomerasi, dan dummy kedekatan dengan kota berpopulasi tinggi. (Hlm. 7,
Tabel 1.)

\textbf{Inferensi pengolahan.} Data availability menyatakan bahwa data
yang digunakan bersifat rahasia. TXT juga tidak memberi daftar lengkap
285 kota, protokol pembersihan dan deduplikasi setiap sumber, atau
berkas kode. Karena itu, dataset tidak dapat direkonstruksi sepenuhnya
dari TXT saja.

\subsection{Population / Sample}\label{population-sample-34}

\textbf{Laporan sumber.} Sampel utama terdiri atas 285 kota tingkat
prefektur di China yang diamati dari 2005 sampai 2020, sehingga
statistik deskriptif utama berjumlah 4.560 observasi kota-tahun. Artikel
menyebut 285 kota tersebut sebagai contoh untuk memvalidasi kerangka dan
model pengukuran. (Hlm. 2 dan 6-7, pendahuluan serta Bagian 3.3.)

\textbf{Laporan sumber.} Ukuran sampel berubah menurut model. Model
dengan ukuran kota awal mengeluarkan Chongqing karena penduduk tetapnya
dianggap pencilan, sehingga \texttt{N\ =\ 4.544}. Tabel 2 mencantumkan
\texttt{N\ =\ 4.275} untuk sebagian model GMM dan \texttt{N\ =\ 4.260}
untuk model GMM dengan ukuran kota awal. Tabel 5 memakai 4.560
observasi; Tabel 6 memakai 4.544 untuk model ukuran kota awal dan 4.560
untuk indikator lain. (Hlm. 8, 11-12, Tabel 2, 5, dan 6.)

\textbf{Laporan sumber.} Analisis khusus mencakup 115 kota berbasis
sumber daya pada 2005-2020, panel distrik municipal untuk 2000, 2010,
dan 2020 dengan \texttt{N\ =\ 854} atau \texttt{N\ =\ 851}, 35 ibu kota
provinsi dan kota terencana pada 2005-2020 dengan \texttt{N\ =\ 560}
atau \texttt{N\ =\ 544}, serta 285 kota pada 2005-2019 dengan
\texttt{N\ =\ 4.275} atau \texttt{N\ =\ 4.260}. (Hlm. 8-9 dan 15-16,
Bagian 4.1, Bagian 4.3, serta Tabel 9-11.)

\textbf{Inferensi pengolahan.} File tidak menjelaskan kriteria pemilihan
285 kota, apakah cakupannya sensus atas seluruh kota tingkat prefektur
yang memenuhi syarat, atau bagaimana observasi yang hilang selain kasus
Chongqing ditangani. Karena itu, representativitas untuk seluruh China
tidak dapat dinilai hanya dari TXT. Perbedaan \texttt{N} yang
dicantumkan pada Tabel 2 juga tidak seluruhnya direkonsiliasi oleh
catatan tabel yang menyebut ukuran sampel GMM 4.260.

\subsection{Dependent Variables}\label{dependent-variables-34}

\textbf{Laporan sumber.} Pertumbuhan ekonomi perkotaan dioperasionalkan
dengan dua variabel terikat:

\begin{itemize}
\tightlist
\item
  \textbf{Pendapatan lahan perkotaan per kapita:} ditulis sebagai
  \texttt{Qit/Cit}, dengan \texttt{Qit} sebagai tingkat pendapatan dan
  \texttt{Cit} sebagai area terbangun kota pada rumus produksi yang
  dijelaskan penulis. Variabel ini menjadi outcome utama pada banyak
  model panel, spasial, IV, dan robustitas.
\item
  \textbf{GDP per kapita:} ditulis sebagai \texttt{Qit/Nit}, dengan
  \texttt{Nit} sebagai ukuran penduduk. Variabel ini dipakai sebagai
  outcome pembanding untuk menguji ketahanan hasil terhadap proksi
  pertumbuhan yang berbeda. (Hlm. 5-6, Bagian 3.2.)
\end{itemize}

\textbf{Laporan sumber.} Data GDP yang dimodifikasi dengan cahaya malam
digunakan sebagai outcome alternatif dalam Tabel 10. Bentuk logaritmik
digunakan pada model, sebagaimana terlihat dari label seperti
\texttt{Ln(per\ capita\ urban\ land\ revenue)} dan
\texttt{Ln(GDP\ per\ capita)}. (Hlm. 7 dan 15, Bagian 3.3 serta Tabel
1-2 dan 10.)

\textbf{Inferensi pengolahan.} Kedua outcome adalah proksi ekonomi
agregat, bukan ukuran langsung atas kesejahteraan, produktivitas total,
distribusi pendapatan, atau kualitas pembangunan. Makna label ``per
kapita'' pada \texttt{Qit/Cit} mengikuti penjelasan dan notasi sumber;
TXT tidak memberikan definisi operasional tambahan yang dapat
menyelesaikan hubungan antara pendapatan lahan dan area terbangun.

\subsection{Results}\label{results-34}

\textbf{Laporan sumber: aglomerasi intrakota.} Tabel 2 menunjukkan
hubungan positif dan signifikan antara kepadatan penduduk dengan
pendapatan lahan perkotaan per kapita pada model efek tetap, dengan
koefisien 0,897 dan 0,869. Ukuran kota awal juga positif, dengan
koefisien 0,029 dan 0,040, sedangkan dummy keanggotaan aglomerasi
bernilai 0,649. Pada model GMM, koefisiennya masing-masing 0,473 untuk
kepadatan penduduk, 0,031 untuk ukuran kota awal, dan 0,060 untuk dummy
keanggotaan aglomerasi. (Hlm. 8, Tabel 2.)

\textbf{Interpretasi penulis.} Penulis menerjemahkan koefisien tersebut
sebagai kenaikan 8,7 persen pendapatan lahan perkotaan per kapita untuk
setiap kenaikan 10 persen kepadatan penduduk, kenaikan 2,9-4 persen
ketika ukuran penduduk awal menjadi dua kali lipat, dan kenaikan
rata-rata 0,65 bagi kota yang menjadi bagian dari aglomerasi perkotaan.
Kesamaan hasil model dasar dan GMM dipakai untuk menyatakan ketahanan
temuan. (Hlm. 8, Bagian 4.1.)

\textbf{Laporan sumber: perbedaan proksi GDP.} Untuk GDP per kapita,
model efek tetap menghasilkan koefisien positif untuk kepadatan
penduduk, yaitu 0,040 dan 0,013, untuk ukuran kota awal, yaitu 0,033 dan
0,072, serta untuk dummy keanggotaan aglomerasi, yaitu 0,524. Pada model
GMM, koefisien kepadatan penduduk turun menjadi 0,002, ukuran kota awal
menjadi -0,001 dan tidak signifikan, sedangkan dummy keanggotaan menjadi
-0,010. Penulis menyatakan hasil GDP kurang robust dibandingkan hasil
pendapatan lahan per kapita. (Hlm. 8-9, Tabel 2.)

\textbf{Laporan sumber: kota berbasis sumber daya dan variabel kendali.}
Pada 115 kota berbasis sumber daya, pengaruh kepadatan penduduk terhadap
pendapatan lahan per kapita tetap sekitar 0,87, tetapi pengaruhnya
terhadap GDP per kapita tidak signifikan pada taraf 1 persen. Ukuran
kota awal berdampak negatif pada kedua outcome. Stok modal lahan per
kapita dan pangsa pekerjaan manufaktur umumnya positif terhadap
pendapatan lahan, sedangkan hasil aplikasi paten tidak sejalan dengan
ekspektasi konvensional dan tanda suhu tahunan tidak robust. Penambahan
kuadrat suhu menghasilkan hubungan U terbalik dengan koefisien kuadrat
-0,07 yang signifikan pada taraf 1 persen. (Hlm. 8-9, Bagian 4.1 dan
Tabel 2.)

\textbf{Laporan sumber: heterogenitas ruang, kebijakan, dan alam.}
Dengan outcome pendapatan lahan per kapita, interaksi kepadatan penduduk
dengan wilayah timur dan tengah bernilai positif, masing-masing 0,019
dan 0,032 pada model efek tetap; estimasi GMM juga positif,
masing-masing sekitar 0,027 dan 0,011. Penulis menyatakan efek
aglomerasi-pertumbuhan lebih kuat di wilayah timur dan tengah daripada
barat. Penguatan untuk ibu kota provinsi dan kota terencana lebih lemah
pada sebagian model efek tetap tetapi positif pada GMM. Hasil untuk kota
pesisir terbuka dan faktor alam tidak konsisten. (Hlm. 9-10, Tabel 3.)

\textbf{Laporan sumber.} Untuk GDP per kapita, penulis kembali
melaporkan perbedaan menurut lokasi dan kebijakan, tetapi hasilnya
kurang robust terutama pada kota-kota barat. Penulis mengaitkan
ketidakrobustan itu dengan proporsi kota berbasis sumber daya yang besar
di wilayah barat dan perbedaan antara urbanisasi lahan serta urbanisasi
penduduk. (Hlm. 10-11, Tabel 4.)

\textbf{Laporan sumber: limpahan spasial.} Moran's I untuk pendapatan
lahan perkotaan per kapita dan GDP per kapita berada pada rentang
0,26-0,53 dan signifikan pada taraf 1 persen. Tabel 5 menunjukkan
koefisien kepadatan penduduk positif pada model SAR, SEM, dan SDM untuk
outcome pendapatan lahan, masing-masing 0,757, 0,845, dan 0,767. Dalam
SDM, efek tidak langsung kepadatan penduduk kota tetangga bernilai 0,106
dan signifikan pada taraf 1 persen. Parameter spasial juga positif dan
signifikan, yaitu \texttt{rho\ =\ 0,178} pada SAR,
\texttt{lambda\ =\ 0,579} pada SEM, dan \texttt{rho\ =\ 0,571} pada SDM.
(Hlm. 11, Tabel 5.)

\textbf{Laporan sumber.} Untuk GDP per kapita, koefisien langsung
kepadatan penduduk pada Tabel 5 adalah -0,004 pada SAR, -0,013 pada SEM,
dan -0,029 pada SDM; efek tidak langsung pada SDM adalah 0,060 dan
signifikan. Penulis tetap menyimpulkan adanya limpahan spasial positif
dan pola hubungan ekonomi kuat-kuat antarkota berdasarkan parameter
spasial, efek tidak langsung, serta ekonomi tetangga. Dalam model SDM
dengan pendapatan lahan, penulis menghitung kontribusi aglomerasi
intrakota sekitar 88 persen dan aglomerasi antarkota sekitar 12 persen.
(Hlm. 11-12, Bagian 4.2 dan Tabel 5.)

\textbf{Laporan sumber: borrowed size dan agglomeration shadow.} Pada
Tabel 6, dummy kedekatan dengan kota berpopulasi tinggi, yang
didefinisikan sebagai kota yang berdampingan dengan kota dalam 5 persen
teratas populasi, bernilai positif pada model pendapatan lahan untuk
kepadatan penduduk (\texttt{0,538}) dan ukuran kota awal
(\texttt{0,201}). Efeknya tidak signifikan pada model dummy keanggotaan
aglomerasi (\texttt{0,068}). Interaksi ukuran kota awal dengan kedekatan
bernilai positif dan signifikan (\texttt{0,037}), sedangkan interaksi
kepadatan penduduk tidak signifikan (\texttt{-0,006}) dan interaksi
dummy keanggotaan aglomerasi bernilai negatif serta signifikan
(\texttt{-0,253}). Hasil GDP memperlihatkan pola kedekatan yang positif
pada model kepadatan (\texttt{0,538}), ukuran awal (\texttt{0,744}), dan
dummy aglomerasi (\texttt{0,218} pada taraf 10 persen), dengan interaksi
ukuran awal \texttt{0,103} dan interaksi dummy aglomerasi \texttt{0,135}
yang signifikan. (Hlm. 12, Tabel 6.)

\textbf{Interpretasi penulis.} Koefisien positif kedekatan dengan kota
yang memiliki aglomerasi tinggi dibaca sebagai borrowed size, sedangkan
pengurangan kontribusi aglomerasi lokal ketika unsur tetangga dimasukkan
dan beberapa interaksi negatif dibaca sebagai bukti bahwa agglomeration
shadow juga ada. Karena efek positif tetangga lebih besar secara
keseluruhan, penulis menyatakan borrowed size lebih kuat daripada shadow
dan kota dengan ukuran awal lebih besar lebih mungkin meminjam skala
kota tetangga. (Hlm. 12-13, Bagian 4.2.)

\textbf{Laporan sumber: variasi antarkota.} Tabel 7 menunjukkan bahwa
efek kedekatan dengan kota yang memiliki aglomerasi tinggi berbeda
menurut lokasi dan kebijakan. Efeknya lebih kuat di wilayah timur serta
pada ibu kota provinsi dan kota terencana, sedangkan koefisien untuk
kota pesisir terbuka tetap negatif. Penulis juga melaporkan pola yang
searah untuk GDP per kapita, dengan pengaruh kedekatan yang relatif
lebih kuat di wilayah timur, ibu kota, dan zona temperate. (Hlm. 13,
Tabel 7.)

\textbf{Laporan sumber: uji ketahanan.} Estimasi IV/2SLS, data distrik
municipal, GDP yang dimodifikasi dengan cahaya malam, perubahan objek
studi, dan pengeluaran tahun 2020 umumnya mempertahankan arah utama
hasil. Narasi penulis menyebut koefisien kepadatan penduduk sekitar 0,78
untuk pendapatan lahan per kapita dan 0,52 untuk GDP per kapita dalam
uji IV, serta koefisien aglomerasi antarkota sekitar 0,04. Tabel 8
menampilkan beberapa koefisien IV yang berbeda menurut instrumen,
termasuk 0,806 dan 0,230 pada dua outcome untuk instrumen RDLS dikalikan
nilai tukar, serta 0,731 dan 0,666 untuk instrumen kepadatan jaringan
sungai dikalikan nilai tukar. (Hlm. 13-14, Bagian 4.3 dan Tabel 8.)

\textbf{Laporan sumber.} Analisis distrik municipal menghasilkan
koefisien sekitar 0,85 untuk kepadatan penduduk dan 0,77 untuk ukuran
kota awal terhadap pendapatan lahan per kapita; koefisien kedekatan
antarkota sekitar 0,5. Ketika unsur aglomerasi antarkota dimasukkan,
koefisien kepadatan turun sekitar 0,2 untuk pendapatan lahan per kapita
dan 0,1 untuk GDP per kapita. Koreksi GDP dengan cahaya malam serta
analisis 35 ibu kota provinsi dan kota terencana juga disebut mendukung
hasil utama. (Hlm. 15-16, Bagian 4.3 dan Tabel 9-11.)

\textbf{Inferensi pengolahan.} Hasil robustitas mendukung kestabilan
arah hubungan dalam spesifikasi yang dipilih, tetapi tidak menghapus
perbedaan besar antara outcome pendapatan lahan dan GDP, terutama pada
GMM dan subkelompok kota berbasis sumber daya. Karena itu, kesimpulan
substantif paling kuat adalah hubungan positif pada proksi pendapatan
lahan dengan limpahan antarkota positif; generalisasi yang sama ke semua
ukuran pertumbuhan ekonomi memerlukan kehati-hatian.

\subsection{Findings}\label{findings-34}

\textbf{Temuan yang dilaporkan sumber.} Empat temuan utama artikel
adalah:

\begin{itemize}
\tightlist
\item
  aglomerasi penduduk intrakota berkaitan positif dengan pertumbuhan
  ekonomi perkotaan, terutama ketika pertumbuhan diukur dengan
  pendapatan lahan perkotaan per kapita;
\item
  efek aglomerasi intrakota lebih kuat di wilayah timur dan tengah,
  serta lebih kuat pada ibu kota provinsi dan kota terencana dalam
  spesifikasi yang mendukungnya;
\item
  aglomerasi penduduk kota tetangga mempunyai limpahan positif dan
  borrowed size lebih kuat daripada agglomeration shadow;
\item
  setelah aglomerasi antarkota dimasukkan, koefisien aglomerasi
  intrakota turun sekitar 0,1, yang dibaca sebagai keberadaan efek
  shadow yang kecil; meskipun demikian, aglomerasi intrakota tetap lebih
  kuat sebagai pendorong pertumbuhan daripada aglomerasi antarkota.
  (Hlm. 1, abstrak; hlm. 11-18, Bagian 4-5.)
\end{itemize}

\textbf{Interpretasi penulis.} Penulis menafsirkan pola tersebut sebagai
pola pembangunan saling meminjam di antara kota-kota China. Kota
tetangga dengan aglomerasi tinggi dapat mendorong pertumbuhan kota yang
diamati, sementara sebagian sumber daya lokal dapat terserap oleh
wilayah tetangga. (Hlm. 12-18, Bagian 4.2-5.2.)

\textbf{Inferensi pengolahan.} Bukti artikel lebih langsung menunjukkan
korelasi spasial dan perubahan koefisien setelah kovariat antarkota
ditambahkan daripada mengamati proses peminjaman secara langsung.
Istilah borrowed size dan agglomeration shadow dalam review ini karena
itu dipertahankan sebagai interpretasi penulis atas proksi dan model,
bukan sebagai observasi langsung atas arus tenaga kerja, modal, layanan,
atau pengetahuan.

\subsection{Conclusions}\label{conclusions-34}

\textbf{Kesimpulan yang dinyatakan penulis.} Penulis menyimpulkan bahwa:

\begin{itemize}
\tightlist
\item
  aglomerasi penduduk intrakota berdampak positif terhadap pertumbuhan
  ekonomi perkotaan, dengan efek yang dimoderasi oleh lokasi dan faktor
  kebijakan; efeknya lebih kuat di wilayah timur dan tengah, serta pada
  ibu kota provinsi dan kota terencana;
\item
  borrowed size lebih kuat daripada agglomeration shadow di antara
  kota-kota China, aglomerasi intrakota memiliki limpahan spasial
  positif, dan aglomerasi tinggi di kota tetangga meningkatkan
  pertumbuhan kota;
\item
  efek positif aglomerasi tetangga mempunyai variasi ruang dan
  kebijakan, sementara kontribusi aglomerasi intrakota terhadap
  pertumbuhan turun sekitar 0,1 setelah unsur antarkota dimasukkan.
  (Hlm. 18, Bagian 5.2.)
\end{itemize}

\textbf{Inferensi pengolahan.} Kesimpulan pertama dan ketiga paling kuat
untuk outcome pendapatan lahan perkotaan per kapita dan spesifikasi yang
mempertahankan arah hasil. Kesimpulan tentang GDP per kapita memerlukan
kualifikasi karena beberapa estimasi GMM tidak signifikan atau berubah
tanda. Kesimpulan borrowed size juga sebaiknya tidak dibaca sebagai
identifikasi kausal atau bukti mekanisme lengkap.

\subsection{Contributions}\label{contributions-34}

\textbf{Kontribusi yang diklaim sumber.} Artikel menyatakan beberapa
kontribusi:

\begin{itemize}
\tightlist
\item
  menyusun kerangka teoritis dan model pengukuran yang menggabungkan
  aglomerasi intrakota, borrowed size, dan agglomeration shadow;
\item
  memvalidasi kerangka tersebut pada 285 kota tingkat prefektur China;
\item
  mengumpulkan dan saling memverifikasi dataset berskala besar dari
  lebih dari 500 buku tahunan statistik, 4.560 buletin statistik, data
  CNIPA, data Geospatial Data Cloud, serta citra cahaya malam;
\item
  menggunakan seperangkat indikator aglomerasi dan pertumbuhan, yaitu
  kepadatan penduduk, ukuran kota awal, dummy aglomerasi perkotaan, GDP
  per kapita, dan pendapatan lahan perkotaan per kapita, bukan satu
  indikator tunggal; dan
\item
  memperluas perdebatan yang sebelumnya sering dipisahkan antara
  aglomerasi intrakota dan antarkota dalam konteks kota-kota China.
  (Hlm. 2-4, pendahuluan.)
\end{itemize}

\textbf{Inferensi pengolahan.} Kontribusi empiris terkuat adalah menguji
indikator lokal dan tetangga dalam satu kerangka panel serta memeriksa
hasil dengan beberapa sumber data. Kontribusi terhadap pengukuran
mekanisme borrowed size masih bersifat tidak langsung karena variabel
utama antarkota berupa indikator kedekatan, aglomerasi tetangga, dan
interaksi model.

\subsection{Research Gap}\label{research-gap-34}

\textbf{Kesenjangan yang dinyatakan sumber sebelum penelitian.} Artikel
mengidentifikasi:

\begin{itemize}
\tightlist
\item
  belum adanya kerangka teoritis dan model empiris yang komprehensif
  untuk aglomerasi intrakota dan antarkota secara bersamaan;
\item
  belum adanya konsensus apakah borrowed size atau agglomeration shadow
  lebih dominan dalam hubungan antarkota;
\item
  keterbatasan pengujian limpahan spasial aglomerasi penduduk pada
  kota-kota China;
\item
  persoalan endogenitas antara aglomerasi penduduk dan pertumbuhan
  ekonomi; dan
\item
  persoalan akuisisi serta verifikasi data statistik, terutama untuk
  data kepadatan penduduk dan area terbangun. (Hlm. 2-5, pendahuluan dan
  Bagian 2.2.)
\end{itemize}

\textbf{Kesenjangan yang masih terbuka menurut sumber.} Pada bagian
keterbatasan, artikel menyatakan perlunya membandingkan aglomerasi
industri tradisional dengan aglomerasi perkotaan umum, menganalisis
hubungan kota sebagai jaringan alih-alih hanya hierarki, dan
mengeksplorasi proses serta mekanisme penjelas dampak aglomerasi. (Hlm.
17-18, Bagian 5.1-5.2.)

\textbf{Inferensi pengolahan.} Setelah studi ini, masih belum ada
pengukuran langsung mengenai arus atau fungsi yang dipinjam,
pembandingan kontrafaktual kota yang tidak berdekatan dengan kota besar,
dan pemisahan yang kuat antara efek ukuran lokal, kedekatan, dan
mekanisme transfer antarkota. Ini adalah inferensi dari variabel serta
batasan yang dijelaskan TXT, bukan gap tambahan yang dinyatakan
eksplisit oleh penulis.

\subsection{Limitations}\label{limitations-34}

\textbf{Keterbatasan yang dinyatakan sumber.} Penulis menyebut tiga
keterbatasan utama:

\begin{itemize}
\tightlist
\item
  penelitian hanya berfokus pada aglomerasi penduduk dan tidak
  membandingkannya dengan aglomerasi industri tradisional;
\item
  analisis belum memperhitungkan hubungan antarkota yang semakin
  menyerupai jaringan, termasuk hubungan yang dibentuk oleh transportasi
  dan informasi; dan
\item
  proses dampak aglomerasi serta mekanisme penjelasnya belum
  dieksplorasi secara mendalam. (Hlm. 17-18, Bagian 5.1.)
\end{itemize}

\textbf{Keterbatasan yang diturunkan dari isi file sebagai inferensi
pengolahan.}

\begin{itemize}
\tightlist
\item
  \textbf{Validitas konstruk:} kedekatan dengan kota berpopulasi tinggi
  dan aglomerasi penduduk tetangga adalah proksi untuk borrowed size
  atau shadow, bukan pengukuran langsung aliran tenaga kerja, modal,
  layanan, teknologi, atau pengetahuan.
\item
  \textbf{Perbedaan outcome:} hasil pendapatan lahan perkotaan per
  kapita lebih konsisten daripada hasil GDP per kapita. Perbedaan ini
  membatasi klaim bahwa aglomerasi selalu meningkatkan semua ukuran
  pertumbuhan ekonomi.
\item
  \textbf{Identifikasi:} panel observasional dan strategi IV/GMM tidak
  dengan sendirinya membuktikan kausalitas. Asumsi eksogenitas dan jalur
  pengecualian instrumen tidak dijelaskan secara lengkap dalam TXT.
\item
  \textbf{Cakupan sampel:} kriteria pemilihan 285 kota dan daftar kota:
  \textbf{Tidak disebutkan dalam file}. Karena itu, batas generalisasi
  tidak dapat diperiksa sepenuhnya.
\item
  \textbf{Ketersediaan data:} pernyataan ketersediaan data menyebut data
  yang digunakan bersifat rahasia. File juga tidak menyebutkan apakah
  kode analisis dibagikan.
\item
  \textbf{Ukuran sampel:} Tabel 2 mencantumkan \texttt{N\ =\ 4.275} pada
  sebagian model GMM dan \texttt{N\ =\ 4.260} pada model lainnya,
  sedangkan catatan tabel menyebut ukuran sampel GMM 4.260 secara umum.
  TXT tidak menjelaskan perbedaan itu.
\item
  \textbf{Kekuatan instrumen:} narasi Bagian 4.3 menyatakan statistik
  Kleibergen-Paap Wald rk F melebihi nilai kritis 16,38, tetapi Tabel 8
  menampilkan nilai 10,66 dan 10,64 untuk model berbasis instrumen
  slope, lebih rendah daripada 16,38. Ketidaksesuaian ini tidak
  diselesaikan dalam file.
\item
  \textbf{Transparansi model:} matriks bobot spasial, konstruksi
  ketetanggaan secara rinci, pembersihan data, penanganan seluruh nilai
  hilang, dan kode estimasi tidak tersedia dalam TXT.
\item
  \textbf{Notasi outcome:} rumus dan label \texttt{Qit/Cit} mengikuti
  penjelasan penulis, tetapi file tidak memberi klarifikasi operasional
  tambahan atas penggunaan area terbangun \texttt{Cit} dalam outcome
  yang dinamai pendapatan lahan per kapita.
\end{itemize}

\subsection{Challenges}\label{challenges-34}

\textbf{Tantangan yang dilaporkan sumber.} Penelitian menghadapi
endogenitas antara kepadatan penduduk dan pertumbuhan ekonomi, serta
hubungan timbal balik antara kepadatan penduduk dan pendapatan lahan per
kapita. Data area terbangun pada tingkat kota disebut langka, data
statistik tahun 2017 tidak sepenuhnya sesuai dengan unit kota yang
dibutuhkan, dan data paten pada \emph{China City Statistical Yearbook}
hanya tersedia mulai 2017. Penulis juga perlu menggabungkan dan
mengoreksi dua seri cahaya malam yang berbeda, menangani ukuran kota
awal yang tidak berubah sepanjang waktu, serta memeriksa pengaruh kota
berbasis sumber daya dan tahun 2020. (Hlm. 3-7 dan 13-16, Bagian 2.2-3.3
dan 4.3.)

\textbf{Tantangan yang dilaporkan sumber.} Penulis menyebut adanya
hambatan regional yang telah membatasi arus tenaga kerja dan modal
antarkota. Hambatan tersebut menjadi alasan bagi rekomendasi untuk
membangun konektivitas antarkota. (Hlm. 17, Bagian 5.1.)

\textbf{Inferensi pengolahan.} Tantangan replikasi terbesar adalah
menghubungkan banyak sumber statistik, data cahaya malam, dan unsur
geografi ketika data akhir dinyatakan rahasia. Perbedaan \texttt{N} dan
statistik kekuatan instrumen antara narasi dan tabel juga perlu
diselesaikan sebelum hasil dapat direplikasi atau dibandingkan secara
kuantitatif.

\subsection{Practical Implications}\label{practical-implications-34}

\textbf{Implikasi yang dinyatakan penulis.} Artikel memberi tiga
rekomendasi kebijakan:

\begin{itemize}
\tightlist
\item
  mendorong perencanaan perkotaan dan perdesaan berbasis aglomerasi,
  termasuk perencanaan stok lahan dan peningkatan aglomerasi penduduk,
  industri, serta jasa;
\item
  membangun sistem konektivitas antarkota yang lebih dalam melalui
  mekanisme lintas kota, termasuk jalan berkecepatan tinggi dan
  jembatan, agar tenaga kerja dan modal dapat bergerak; dan
\item
  menyebarkan fungsi kota besar di dalam aglomerasi perkotaan melalui
  mobilitas tenaga kerja, modal, dan teknologi serta transfer industri
  yang rasional ke kota kecil dan menengah di sekitarnya. (Hlm. 17-18,
  Bagian 5.2.)
\end{itemize}

\textbf{Interpretasi penulis.} Penyebaran fungsi kota besar dipandang
layak untuk mengurangi homogenisasi ekonomi dan kesenjangan antarkota,
sementara konektivitas dipandang dapat memperkuat manfaat borrowed size.
(Hlm. 17, Bagian 5.1-5.2.)

\textbf{Inferensi pengolahan.} Rekomendasi tersebut tidak boleh
diterapkan hanya berdasarkan koefisien kedekatan atau kepadatan
penduduk. Validasi tambahan atas arus, kapasitas infrastruktur, karakter
kota berbasis sumber daya, dan perbedaan antara GDP serta pendapatan
lahan diperlukan sebelum menyimpulkan bahwa intervensi tertentu akan
menghasilkan borrowed size.

\subsection{Applications}\label{applications-34}

\textbf{Laporan sumber.} File tidak menyediakan bagian terpisah bernama
``Applications''. Penggunaan yang ditunjukkan langsung adalah memakai
kerangka dan indikator gabungan untuk mengevaluasi pengaruh aglomerasi
penduduk intrakota serta antarkota terhadap pertumbuhan ekonomi 285 kota
tingkat prefektur, menguji limpahan spasial, dan menyusun rekomendasi
perencanaan serta konektivitas antarkota. (Hlm. 2-7 dan 11-18, Bagian
1-5.)

\textbf{Inferensi pengolahan.} Kerangka tersebut dapat dipakai sebagai
diagnosis awal untuk membandingkan kota lokal dan kota tetangga, tetapi
transfer ke wilayah lain merupakan inferensi, bukan aplikasi yang diuji
dalam TXT. Penggunaan semacam itu memerlukan definisi unit, indikator,
data, dan hubungan spasial yang setara.

\subsection{Future Research
(Optional)}\label{future-research-optional-34}

\textbf{Agenda yang dinyatakan penulis.} Penulis menyarankan penelitian
berikutnya untuk:

\begin{itemize}
\tightlist
\item
  membandingkan dampak aglomerasi industri dengan aglomerasi perkotaan
  umum terhadap pertumbuhan ekonomi;
\item
  menggunakan statistik jaringan untuk menganalisis efek aglomerasi
  ketika hubungan antarkota dipahami sebagai jaringan transportasi atau
  informasi; dan
\item
  mengeksplorasi proses dampak aglomerasi serta mekanisme penjelasnya
  secara lebih mendalam. (Hlm. 17-18, Bagian 5.1.)
\end{itemize}

\textbf{Inferensi pengolahan.} Berdasarkan celah yang terlihat dalam
TXT, pengukuran arus antarkota dan pengujian sensitivitas terhadap
indikator pertumbuhan dapat menjadi pengembangan lanjutan. Usulan ini
adalah inferensi review, bukan agenda yang dinyatakan eksplisit oleh
penulis.

\subsection{Provenance Notes}\label{provenance-notes-34}

\begin{itemize}
\tightlist
\item
  \textbf{Sumber yang digunakan:} tepat satu file TXT aktual dengan
  prefiks A19 ditemukan dan diverifikasi melalui pencarian
  \texttt{source/**/A19*.txt}:
  \texttt{source/journal/A19-chen-wang-hao-yin-chen-2024-population-agglomeration-borrowed-size-and-urban-economic-growth-in-china-10.1016-j.cities.2024.105491.txt}.
\item
  \textbf{Cakupan baca:} seluruh 1.679 baris TXT dibaca, dari blok
  \texttt{{[}PDF\ Metadata{]}} pada baris 1-28, blok
  \texttt{{[}Extracted\ Text{]}} mulai baris 29, teks artikel, seluruh
  tabel dan bagian hasil, hingga daftar pustaka pada baris 1458-1678.
\item
  \textbf{Metadata ekstraksi:} TXT menyatakan PDF diekstrak pada
  2026-08-31 dengan \texttt{pdftotext\ -layout}. Metadata mencatat PDF
  20 halaman, tidak terenkripsi, ukuran 595,276 x 793,701 points, dan
  ukuran file 2.210.356 bytes. Creator yang tercatat adalah Elsevier dan
  producer-nya Acrobat Distiller 8.1.0 (Windows); PDF tidak ditandai
  (\emph{tagged: no}) dan tidak memuat JavaScript. Path PDF yang dicatat
  adalah file A19 dengan nama dasar yang sama. (TXT, baris 1-28.)
\item
  \textbf{Identitas:} metadata PDF hanya mencantumkan \texttt{Le\ Chen}
  sebagai author, sedangkan blok judul artikel mencantumkan lima
  penulis. Review menggunakan identitas dari blok judul, afiliasi, DOI,
  abstrak, dan keterangan artikel dalam TXT, bukan memperluas metadata
  mesin PDF. (TXT, baris 7-13 dan 42-50.)
\item
  \textbf{Basis akses:} review ini hanya memakai TXT A19. PDF tidak
  dibuka sebagai sumber tambahan, dan tidak ada sumber eksternal,
  metadata daring, atau artikel yang tercantum di daftar pustaka yang
  dibuka untuk menambah substansi.
\item
  \textbf{Locator:} nomor halaman dalam review merujuk pada nomor
  halaman artikel yang tampil dalam hasil ekstraksi; bagian, persamaan,
  tabel, dan gambar disebut bila tersedia. Nomor baris TXT digunakan
  terutama untuk menjelaskan metadata dan cakupan baca.
\item
  \textbf{Keterbatasan ekstraksi:} tata letak dua kolom membuat sebagian
  tabel, persamaan, dan kalimat tersusun berselang dalam TXT. Angka pada
  review diambil dari label tabel dan narasi yang tersedia;
  ketidaksesuaian internal tidak disamarkan dan dicatat pada bagian
  Limitations.
\item
  \textbf{Data dan informasi yang tidak tersedia:} data penelitian
  dinyatakan rahasia. Kriteria pemilihan kota, daftar lengkap kota, kode
  analisis, matriks bobot spasial, rincian pembersihan data, dan
  beberapa penjelasan tentang perbedaan ukuran sampel: \textbf{Tidak
  disebutkan dalam file}.
\item
  \textbf{Pembedaan epistemik:} label \textbf{Laporan sumber},
  \textbf{Temuan yang dilaporkan sumber}, dan \textbf{Kesimpulan yang
  dinyatakan penulis} adalah parafrasa isi artikel. Label
  \textbf{Interpretasi penulis} merekam makna yang diberikan penulis
  terhadap hasil. Label \textbf{Inferensi pengolahan} adalah pembacaan
  analitis atas isi, proksi, dan ketidakkonsistenan internal TXT; label
  itu bukan klaim tambahan dari artikel.
\end{itemize}

\section{Review A20}\label{review-a20}

\textbf{Identitas sumber}

\begin{itemize}
\tightlist
\item
  Guo, Aijun; Liu, Peixian; Zhong, Fanglei; Yang, Chunlin; dan Luo,
  Xijing. 2022. ``Borrowing Size and Urban Green Development Efficiency
  in the City Network of China: Impact Measures and Size Thresholds.''
  \emph{Land}, 11, 493. DOI: \texttt{10.3390/land11040493}.
\item
  Review ini disusun dari TXT ekstraksi sumber A20, tanpa menambahkan
  pengetahuan dari luar file.
\end{itemize}

\subsection{Summarized Abstract}\label{summarized-abstract-35}

\textbf{Temuan sumber:} Artikel menguji pengaruh tiga dimensi borrowing
size, yaitu borrowing population size (Bpop), borrowing economic
activity density (Bden), dan borrowing advanced functions (Bfun),
terhadap urban green development efficiency (GDE) pada 280 kota tingkat
prefektur di China selama 2009-2019. GDE diukur dengan super-efficiency
slacks-based measure (SBM) yang memasukkan output tidak diinginkan.
Pengaruh linear diuji dengan model efek tetap ganda, sedangkan hubungan
nonlinier diuji dengan panel threshold regression menggunakan city size
sebagai variabel ambang. Ketiga dimensi borrowing size dilaporkan
berpengaruh positif terhadap GDE. Penulis juga melaporkan
double-threshold effect: Bpop dan Bfun memiliki pola U-shaped terhadap
GDE menurut rentang city size, sedangkan Bden memiliki pola inverted
U-shaped. Penulis merekomendasikan perhatian pada borrowing size dan
perluasan city size secara wajar. (Locator: Abstract, PDF p.~1; blok
metadata \texttt{Subject}.)

\textbf{Interpretasi penulis:} Borrowing size diposisikan sebagai sarana
bagi kota untuk memperoleh manfaat jaringan, sementara besarnya kota
membatasi atau mengubah efek hijau tersebut. (Locator: Abstract, PDF
p.~1.)

\subsection{Problem Statement}\label{problem-statement-35}

\textbf{Temuan sumber:} Penulis menempatkan masalah pada transisi
urbanisasi China dari ekspansi ukuran menuju peningkatan kualitas,
dengan konflik antara peningkatan kualitas pembangunan dan persoalan
sumber daya serta lingkungan. GDE dipakai sebagai indikator komprehensif
yang memperhatikan input sumber daya dan polusi selain efisiensi
ekonomi. Pada saat yang sama, konektivitas antarkota dan pembagian kerja
meningkatkan peran jaringan kota, sehingga muncul pertanyaan apakah
borrowing size berbasis jaringan kota memengaruhi kualitas pembangunan
dan bagaimana ukuran kota sendiri membatasi efek tersebut. (Locator:
Introduction, PDF pp.~1-2.)

\textbf{Konteks perdebatan sumber:} Artikel mengaitkan pertanyaan itu
dengan perdebatan antara pembangunan kota besar dan pembangunan kota
kecil-menengah dalam jaringan kota polisentris. Penulis menyatakan bahwa
ekspansi ukuran dapat menghasilkan agglomeration economy sampai titik
tertentu, tetapi setelah melewati titik kritis dapat memunculkan
diseconomy, penyakit kota besar, serta tekanan sumber daya dan
lingkungan. (Locator: Introduction, PDF pp.~2-3.)

\textbf{Pertanyaan yang dirumuskan penulis:} Apakah borrowing size
berbasis jaringan kota memengaruhi GDE, dan bagaimana pengaruh tersebut
berubah di bawah batasan city size? (Locator: Introduction, PDF p.~2.)

\subsection{Objectives}\label{objectives-35}

\textbf{Tujuan yang dinyatakan atau dibangun dari rancangan penelitian
penulis:}

\begin{itemize}
\tightlist
\item
  Mengukur borrowing size dalam tiga dimensi: Bpop, Bden, dan Bfun.
  (Locator: Introduction, PDF p.~3; Section 3.2, PDF pp.~7-8.)
\item
  Mengukur GDE untuk 280 kota tingkat prefektur di China menggunakan
  super-efficiency SBM dengan output tidak diinginkan. (Locator:
  Introduction, PDF p.~3; Section 3.2, PDF pp.~6-7.)
\item
  Menguji pengaruh borrowing size terhadap GDE dengan model efek tetap
  ganda. (Locator: Introduction, PDF p.~3; Section 3.1, Equation (1),
  PDF p.~5.)
\item
  Menguji hubungan nonlinier dengan city size sebagai variabel threshold
  melalui panel threshold model. (Locator: Introduction, PDF p.~3;
  Section 3.1, Equation (2), PDF pp.~5-6.)
\end{itemize}

\textbf{Inferensi pengolahan:} Empat butir di atas adalah pengelompokan
tujuan berdasarkan pertanyaan, ukuran variabel, dan model yang
dijelaskan sumber; file tidak menyajikannya sebagai daftar heading
``Objectives'' tersendiri.

\subsection{Literature Survey}\label{literature-survey-35}

\textbf{Posisi literatur yang dilaporkan penulis:}

\begin{itemize}
\tightlist
\item
  GDE dipahami sebagai ukuran yang menggabungkan pembangunan ekonomi,
  konservasi sumber daya, dan polusi lingkungan. Urban network
  externality dipahami melalui pembagian kerja, komplementaritas fungsi,
  dan sinergi antarkota. (Locator: Introduction, PDF pp.~2-3; Hypothesis
  Development, PDF p.~4.)
\item
  Borrowing size didefinisikan sebagai akses atau pemanfaatan potensi
  ukuran kota lain melalui posisi nodal dan hubungan fungsional dalam
  jaringan kota. Penulis menggunakan tiga dimensi borrowing size yang
  juga dipakai dalam studi terdahulu yang mereka rujuk. (Locator:
  Hypothesis Development, PDF p.~4; References {[}24{]}, {[}31{]}.)
\item
  Literatur yang dibahas membedakan efek jaringan positif berupa
  borrowing size dari efek negatif kompetisi, siphoning effect, dan
  agglomeration shadow. Penulis menyatakan bahwa borrowing size dan
  agglomeration shadow dapat hidup bersamaan, tetapi besar efeknya dapat
  berbeda. (Locator: Introduction, PDF p.~3.)
\item
  Penulis merangkum literatur sebelumnya sebagai telah banyak
  mengonfirmasi fenomena borrowing size, tetapi menyatakan belum ada
  studi yang secara langsung mengonfirmasi keberadaan agglomeration
  shadow. Ini adalah klaim tinjauan penulis, bukan verifikasi independen
  dalam review ini. (Locator: Introduction, PDF p.~3.)
\end{itemize}

\textbf{Hipotesis sumber:} H1 menyatakan bahwa penggunaan borrowing size
merupakan cara efektif untuk meningkatkan GDE. H2 menyatakan bahwa
peningkatan Bpop, Bden, dan Bfun mendukung pembangunan hijau yang
efisien. H3 menyatakan bahwa efek borrowing size dibatasi city size:
ekspansi ukuran dapat memperkuat efek hijau sampai titik kritis, lalu
dapat menurunkan GDE. (Locator: Hypothesis Development, PDF pp.~4-5.)

\subsection{Method Used}\label{method-used-35}

\textbf{Metode yang dilaporkan sumber:}

\begin{itemize}
\tightlist
\item
  GDE dihitung dengan super-efficiency SBM yang memasukkan non-desired
  outputs. Sistem indikatornya memakai input tenaga kerja, modal, dan
  sumber daya; desired outputs ekonomi, sosial, dan lingkungan; serta
  non-desired outputs berupa polusi industri. (Locator: Section 3.2, PDF
  pp.~6-7; Table 1.)
\item
  Model dasar adalah panel double-fixed-effect: GDE menjadi variabel di
  sisi kiri, Bpop, Bden, dan Bfun menjadi variabel inti, X menjadi
  variabel kontrol, serta terdapat city-specific effects dan time-fixed
  effects. (Locator: Section 3.1, Equation (1), PDF p.~5.)
\item
  Model threshold menggunakan city size (CZ) sebagai threshold variable
  dalam double-threshold panel regression untuk menguji mekanisme
  pembatasan ukuran kota. (Locator: Section 3.1, Equation (2), PDF
  pp.~5-6.)
\item
  Sebelum pengujian threshold, penulis melaporkan bahwa data panel lolos
  LLC test dan IPS test. Jumlah threshold diuji dengan bootstrap 300
  kali menggunakan Stata 16.0. (Locator: Section 4.2, PDF pp.~11-12;
  Table 5.)
\item
  Robustness check pertama memakai borrowing size yang dilag satu
  periode sebagai instrumental variable dan metode two-stage least
  squares (2SLS) untuk menguji endogeneity pada fixed-effect model.
  (Locator: Section 4.3, PDF p.~14; Table 8.)
\item
  Robustness check kedua memakai system generalized method of moments
  dynamic panel model dengan lagged GDE serta variabel borrowing size
  dan kontrol sebagaimana dirumuskan dalam Equation (7). (Locator:
  Section 4.3, Equation (7), PDF pp.~14-15; Table 9.)
\item
  Bpop, Bden, dan Bfun dihitung sebagai jumlah nilai kota lain yang
  dibobot dengan jarak transportasi; jarak antarkota diperoleh melalui
  crawling driving path planning Gaode Map menggunakan R programming
  language. (Locator: Section 3.2, Equations (3)-(6), PDF pp.~7-8;
  Section 3.3, PDF p.~9.)
\end{itemize}

\subsection{Dataset}\label{dataset-35}

\textbf{Temuan sumber:} Dataset utama berupa panel 280 kota tingkat
prefektur di China pada periode 2009-2019. Data makro terutama berasal
dari \emph{China City Statistical Yearbook} (2010-2020). Matriks jarak
transportasi antarkota dihitung dari driving path planning Gaode Map
dengan R; perencanaan rute dilakukan 78.400 kali. Kesenjangan data untuk
sebagian kota dan tahun diisi dengan average interpolation. (Locator:
Abstract, PDF p.~1; Section 3.3, PDF p.~9.)

\textbf{Cakupan indikator:} Table 1 memuat 3.080 observasi untuk setiap
indikator input-output GDE. Table 3 memuat 2.800 observasi untuk GDE,
Bpop, Bden, Bfun, CZ, Fdi, Road, Invest, Gov, Tec, dan lnTec. Nilai
ringkas yang dilaporkan untuk GDE adalah mean 1.0689, standar deviasi
0.3287, minimum 0.0715, dan maksimum 12.1295; untuk CZ adalah mean
450.167, standar deviasi 319.3643, minimum 19.5, dan maksimum 3416.
(Locator: Tables 1 and 3, PDF pp.~7 and 9.)

\textbf{Catatan kesesuaian observasi:} Jumlah observasi berbeda menurut
tabel dan model: 3.080 pada Table 1, 2.800 pada Table 3 serta beberapa
estimasi, dan 2.080 pada beberapa model di Table 4 dan Table 9. TXT
tidak menjelaskan secara khusus penyebab setiap perbedaan jumlah
observasi. (Locator: Tables 1, 3, 4, and 9, PDF pp.~7, 9-10, and 15.)

\textbf{Ketersediaan data menurut sumber:} Sumber menyatakan bahwa
sumber data yang tersedia untuk publik dijelaskan dalam artikel; untuk
data lain, pembaca diminta menghubungi corresponding author dengan
alasan yang wajar. (Locator: Data Availability Statement, PDF p.~17.)

\subsection{Population / Sample}\label{population-sample-35}

\textbf{Temuan sumber:} Unit analisis adalah kota-tahun; cakupan
analisis yang dinyatakan adalah 280 kota tingkat prefektur di China.
Periode observasi yang dinyatakan adalah 2009-2019. File tidak
menjelaskan daftar kota, kriteria inklusi-eksklusi, atau teknik sampling
terpisah; karena itu, teknik sampling: \textbf{Tidak disebutkan dalam
file}. (Locator: Abstract, PDF p.~1; Section 3.2, PDF pp.~6-7.)

\textbf{Inferensi pengolahan:} Review ini menyebut 280 kota sebagai
cakupan dataset yang dilaporkan, bukan sebagai klaim bahwa seluruh kota
tingkat prefektur di China tercakup, karena TXT tidak memberikan daftar
atau aturan pemilihannya.

\subsection{Dependent Variables}\label{dependent-variables-35}

\textbf{Variabel dependen menurut struktur model:} GDE adalah outcome
pada sisi kiri Equation (1) dan Equation (2), serta disebut ``Explained
Variable: GDE'' pada Table 4. GDE diukur dengan super-efficiency SBM
yang memasukkan output tidak diinginkan. (Locator: Section 3.1,
Equations (1)-(2), PDF pp.~5-6; Table 4, PDF p.~10.)

\textbf{Komponen pengukuran GDE:}

\begin{itemize}
\tightlist
\item
  Input: city employment sebagai ukuran tenaga kerja; capital stock;
  urban construction land area; total water supply; dan total
  electricity consumption of the whole society. (Locator: Section 3.2,
  PDF pp.~6-7; Table 1.)
\item
  Desired outputs: real GDP, average wage of urban residents, dan park
  green area. (Locator: Section 3.2, PDF p.~7; Table 1.)
\item
  Non-desired outputs: industrial wastewater discharge, industrial SO2
  emissions, dan industrial fume and dust emissions. (Locator: Section
  3.2, PDF p.~7; Table 1.)
\end{itemize}

\textbf{Inferensi pengolahan:} Section 3.2 dan Table 2 menyebut GDE
sebagai ``Explanatory variable'', tetapi persamaan dan Table 4
memperlakukannya sebagai variabel yang dijelaskan. Review ini
mengklasifikasikan GDE sebagai dependent variable berdasarkan posisi
persamaan dan label Table 4, sambil mempertahankan catatan inkonsistensi
terminologi sumber.

\subsection{Results}\label{results-35}

\textbf{Baseline regression, temuan sumber:} Dalam model tanpa kontrol
pada Table 4 model (1), koefisien Bpop = 0.0087 (\textbf{\emph{; t =
5.96), Bden = 0.0872 (}}; t = 5.33), dan Bfun = 0.7227 (\textbf{\emph{;
t = 6.27). Setelah kontrol ditambahkan pada model (2), masing-masing
tetap positif: Bpop = 0.0024 (}; t = 1.74), Bden = 0.0641 (}\emph{; t =
4.32), dan Bfun = 0.3848 (}\textbf{; t = 3.60). Model (1) dan (2)
memiliki 2.080 observasi dengan R-squared masing-masing 0.2780 dan
0.2832. Model (3)-(5), yang memasukkan satu dimensi borrowing size
secara terpisah, juga melaporkan koefisien positif dan signifikan: Bpop
= 0.0028 (}), Bden = 0.0639 (\textbf{\emph{), dan Bfun = 0.3933 (}});
masing-masing memakai 2.800 observasi. (Locator: Section 4.1, PDF
pp.~10-11; Table 4.)

\textbf{Koefisien kontrol dan interpretasi penulis:} Fdi berkoefisien
negatif dan signifikan, Road negatif dan signifikan, sedangkan Invest,
Gov, dan lnTec positif dan signifikan dalam spesifikasi baseline
berkontrol. Penulis mengaitkan tanda negatif Fdi dengan kemungkinan
dominannya pollution paradise dibanding pollution halo; tanda negatif
Road dengan perluasan jalan, kenaikan kepemilikan kendaraan, dan emisi;
serta tanda positif Invest, Gov, dan lnTec dengan peran modal fisik,
intervensi pemerintah, serta inovasi dan spillover
pengetahuan-teknologi. Ini adalah interpretasi atau kemungkinan penulis,
bukan pengujian mekanisme tambahan yang dilaporkan di bagian tersebut.
(Locator: Section 4.1, PDF pp.~10-11; Table 4.)

\textbf{Threshold test, temuan sumber:} Bootstrap 300 iterasi
menghasilkan single-threshold effect yang signifikan pada 5\% (F =
708.39; p = 0.0133) dan double-threshold effect yang signifikan pada 1\%
(F = 1880.78; p = 0.0033). Triple-threshold effect tidak signifikan (F =
84.43; p = 0.2867), sehingga penulis memilih double-threshold model. Dua
threshold city size dilaporkan sebesar 343.1000 dengan 95\% confidence
interval {[}341.9000, 353.0000{]} dan 358.9600 dengan interval
{[}347.2000, 362.0000{]}. (Locator: Section 4.2, PDF pp.~11-12; Tables
5-6; Figure 3.)

\textbf{Piecewise threshold regression, temuan sumber:} Pada Table 7,
rentang city size dibagi menurut dua threshold tersebut. Bpop memiliki
koefisien 0.0062 (\textbf{\emph{; t = 4.23) pada stage 1, -0.0130 (}}; t
= -8.90) pada stage 2, dan 0.0028 (\textbf{\emph{; t = 2.83) pada stage
3. Bden memiliki koefisien 0.0012 (tidak signifikan; t = 0.11), 4.1659
(}}; t = 51.83), dan 0.2048 (\textbf{\emph{; t = 4.41). Bfun memiliki
koefisien 0.4812 (}}; t = 5.39), 0.3803 (\textbf{; t = 2.02), dan 0.4336
(}*; t = 2.84). N = 2.800 dan R = 0.6382. (Locator: Section 4.2, PDF
pp.~12-14; Table 7.)

\textbf{Pola yang diberi nama oleh penulis:} Penulis menyebut hubungan
Bpop-GDE dan Bfun-GDE di bawah threshold city size sebagai U-shaped,
sedangkan hubungan Bden-GDE sebagai inverted U-shaped. Untuk Bpop,
penulis menawarkan agglomeration shadow atau siphoning effect, serta
keterbatasan kota kecil-menengah dalam mengubah potensi ukuran pasar
menjadi manfaat pembangunan, sebagai dua kemungkinan penjelasan untuk
stage 2. Untuk Bfun, penulis mengaitkan penurunan intensitas pada stage
2 dengan agglomeration shadow atau siphoning, kompetisi dan homogenisasi
fungsi/produk, segmentasi pasar, serta hambatan tidak terlihat.
(Locator: Section 4.2, PDF pp.~12-14.)

\textbf{Robustness results, temuan sumber:} Pada 2SLS, lag borrowing
size satu periode berpengaruh positif signifikan pada tahap pertama;
CD-Wald F-statistic untuk Bpop, Bden, dan Bfun masing-masing 44.63,
14.80, dan 15.55. Tahap kedua tetap menghasilkan efek positif: Bpop =
0.0036 (*), Bden = 0.0025 (\textbf{), dan Bfun = 0.7293 (}), dengan
2.800 observasi. Penulis menyatakan F-values tersebut menunjukkan
instrumental variables valid dan hasilnya konsisten dengan baseline.
(Locator: Section 4.3, PDF p.~14; Table 8.)

\textbf{Dynamic robustness, temuan sumber:} System GMM melaporkan
koefisien positif untuk borrowing size. Dalam model terpisah, Bpop =
0.0026 (\textbf{), Bden = 0.0477 (}\emph{), dan Bfun = 0.0056
(}\textbf{); dalam model gabungan, Bpop = 0.0426 (}\emph{), Bden =
0.0435 (}**), dan Bfun = 0.0529 (*). Penulis menyatakan Sargan serta
AR(1) dan AR(2) tests lolos, dan menyimpulkan hasil utama robust.
Observasi pada tiga model pertama adalah 2.800 dan model gabungan 2.080.
(Locator: Section 4.3, PDF pp.~14-15; Table 9.)

\subsection{Findings}\label{findings-35}

\textbf{Temuan sumber yang diringkas:}

\begin{itemize}
\tightlist
\item
  Ketiga dimensi borrowing size berhubungan positif dengan GDE dalam
  baseline regression, setelah kontrol, dan robustness checks. (Locator:
  Section 4.1 and Section 4.3, PDF pp.~10-11 and 14-15; Tables 4, 8, and
  9.)
\item
  City size memoderasi hubungan tersebut secara nonlinier melalui double
  threshold pada 343.1000 dan 358.9600. (Locator: Section 4.2, PDF
  pp.~11-12; Tables 5-6.)
\item
  Menurut klasifikasi penulis, Bpop dan Bfun mengikuti pola U-shaped,
  sedangkan Bden mengikuti pola inverted U-shaped ketika city size
  berubah lintas rentang threshold. (Locator: Section 4.2, PDF
  pp.~12-14; Table 7.)
\item
  Efek Bden paling besar dilaporkan pada stage 2, sedangkan intensitas
  Bfun paling tinggi pada stage 1, menurun pada stage 2, lalu meningkat
  pada stage 3. (Locator: Section 4.2, PDF pp.~12-14; Table 7.)
\end{itemize}

\textbf{Interpretasi penulis:} Hasil tersebut dipakai penulis untuk
menyatakan bahwa kota dapat memanfaatkan jaringan kota untuk
meningkatkan GDE, tetapi kapasitas kota, kompetisi, agglomeration shadow
atau siphoning, dan homogenisasi fungsi dapat mengubah besarnya manfaat.
(Locator: Sections 2 and 4.2, PDF pp.~4-5 and 12-14.)

\textbf{Inferensi pengolahan:} Review ini merangkum ``positif'' sebagai
arah koefisien yang dilaporkan, bukan sebagai bukti bahwa hubungan
tersebut pasti kausal. TXT menyebut robustness dan pengujian
endogeneity, tetapi tidak memberi dasar bagi review ini untuk menilai
validitas kausal di luar laporan penulis.

\subsection{Conclusions}\label{conclusions-35}

\textbf{Kesimpulan penulis:}

\begin{enumerate}
\def\labelenumi{\arabic{enumi}.}
\tightlist
\item
  Ketiga dimensi borrowing size berpengaruh positif signifikan terhadap
  GDE; kota dapat meningkatkan GDE melalui borrowing size dalam jaringan
  kota.
\item
  Terdapat nonlinear panel threshold effect yang signifikan dan
  berbentuk double threshold antara city size, borrowing size, dan GDE.
\item
  Di bawah batasan city size, Bpop dan Bfun memiliki hubungan U-shaped
  dengan GDE, sedangkan Bden memiliki hubungan inverted U-shaped.
\end{enumerate}

(Locator: Section 5.1, PDF p.~15.)

\textbf{Interpretasi penulis:} Borrowing size dipresentasikan sebagai
gagasan baru untuk mendorong pembangunan kota dan sebagai dasar
informasi bagi peningkatan GDE di China. Penulis juga menyatakan bahwa
studi ini dapat menjadi referensi bagi penelitian di negara lain dan
memperkaya teori tentang dampak borrowing size terhadap pembangunan
kota. (Locator: Section 5.1, PDF pp.~15-16.)

\subsection{Contributions}\label{contributions-35}

\textbf{Kontribusi yang diklaim penulis:}

\begin{itemize}
\tightlist
\item
  Menghitung matriks jarak lalu lintas antarkota dengan crawling driving
  path planning Gaode Map melalui R, bukan hanya memakai jarak Euclidean
  atau matriks waktu kereta yang dibahas memiliki keterbatasan. Penulis
  menyatakan hal ini membuat bobot lebih realistis dan pengukuran
  borrowing size lebih akurat. (Locator: Introduction, PDF p.~3; Section
  3.3, PDF p.~9.)
\item
  Menempatkan borrowing size, GDE, dan city size dalam satu kerangka
  penelitian untuk 280 kota tingkat prefektur di China, lalu menguji
  hubungan nonliniernya dengan panel threshold model. (Locator:
  Introduction, PDF p.~3.)
\item
  Memberikan informasi untuk promosi GDE di China dari perspektif urban
  network externality dan menyatakan relevansi referensial bagi studi
  negara lain. (Locator: Section 5.1, PDF pp.~15-16.)
\end{itemize}

\textbf{Batas status klaim:} Butir-butir ini adalah kontribusi yang
dinyatakan artikel. Review ini tidak memverifikasi akurasi perbandingan
metode jarak, kebaruan literatur, atau generalisasinya ke negara lain.

\subsection{Research Gap}\label{research-gap-35}

\textbf{Kesenjangan eksplisit yang diidentifikasi penulis:}

\begin{enumerate}
\def\labelenumi{\arabic{enumi}.}
\tightlist
\item
  Walaupun faktor-faktor yang memengaruhi GDE telah diuji dari berbagai
  perspektif, masih sedikit studi yang menghubungkan borrowing size
  dengan GDE dari perspektif jaringan kota.
\item
  Belum ada konsensus atas temuan mengenai kemungkinan efek
  agglomeration economy dan congestion dari ekspansi city size terhadap
  pembangunan hijau; masih sedikit studi yang memasukkan borrowing size,
  GDE, dan city size dalam satu kerangka.
\item
  Pengukuran borrowing size membutuhkan distance weights karena
  aksesibilitas lalu lintas antarkota memengaruhinya. Penulis menilai
  jarak Euclidean dan matriks waktu perjalanan kereta yang lazim dipakai
  memiliki keterbatasan, sehingga diperlukan pengaturan bobot berbasis
  jarak transportasi yang lebih realistis.
\end{enumerate}

(Locator: Introduction, PDF pp.~3-4.)

\subsection{Limitations}\label{limitations-35}

\textbf{Keterbatasan yang dinyatakan penulis:}

\begin{itemize}
\tightlist
\item
  Pada analisis mekanisme teoretis, artikel belum membangun model
  matematis yang sesuai dari teori-teori klasik dan baru melakukan
  analisis kualitatif pada tingkat teori. (Locator: Section 5.1, PDF
  p.~16.)
\item
  Pada objek penelitian, efek positif borrowing size bergantung pada
  aktor mikro seperti individu dan perusahaan, tetapi penelitian ini
  tidak masuk lebih dalam ke level mikro. (Locator: Section 5.1, PDF
  p.~16.)
\item
  Penulis mengaitkan kekurangan tersebut dengan ketersediaan data dan
  keterbatasan tingkat penelitian, lalu menyatakan perlunya integrasi
  data kota makro, perusahaan mikro, dan industri meso bersama analisis
  teoretis serta model kuantitatif. (Locator: Section 5.1, PDF p.~16.)
\end{itemize}

\subsection{Challenges}\label{challenges-35}

\textbf{Status bagian dalam sumber:} Tidak ada bagian tersendiri
berjudul ``Challenges''. Tantangan yang tampak dari uraian metode dan
diskusi sumber adalah:

\begin{itemize}
\tightlist
\item
  Menangani perbedaan antarkota dan faktor tak teramati yang berubah
  menurut waktu; penulis meresponsnya dengan individual dan year fixed
  effects. (Locator: Section 3.1, PDF p.~5.)
\item
  Menetapkan bobot aksesibilitas yang lebih realistis; penulis melakukan
  78.400 crawling rute Gaode Map dan mengisi sebagian data yang hilang
  dengan average interpolation. (Locator: Section 3.3, PDF p.~9.)
\item
  Menguji kemungkinan endogeneity dan persistence atau stickiness GDE;
  penulis menambahkan 2SLS dengan lag borrowing size dan system GMM
  dynamic panel. (Locator: Section 4.3, PDF pp.~14-15.)
\item
  Membedakan manfaat jaringan dari kemungkinan kompetisi, siphoning
  effect, dan agglomeration shadow dalam jaringan kota. (Locator:
  Introduction and Hypothesis Development, PDF pp.~2-5.)
\end{itemize}

\subsection{Practical Implications}\label{practical-implications-35}

\textbf{Rekomendasi kebijakan penulis:}

\begin{itemize}
\tightlist
\item
  Memperbaiki infrastruktur jaringan kota dan aksesibilitas
  transportasi, termasuk proyek transportasi komprehensif seperti jalan
  raya dan rel, agar arus faktor produksi antarkota lebih efisien.
  (Locator: Section 5.2, PDF p.~16.)
\item
  Memandang city size secara tepat dan mengejar pembangunan antarkota
  yang terkoordinasi: menghindari crowding effect dari sedikit mega atau
  super cities, serta memperluas ukuran kota kecil-menengah secara
  moderat agar manfaat agglomeration dan eksternalitas penduduk dapat
  digunakan. (Locator: Section 5.2, PDF pp.~16-17.)
\item
  Membangun jaringan kota polisentris dengan fungsi yang komplementer,
  memperkuat status kota sebagai node, dan menggunakan posisi serta
  fungsi tiap kota untuk memperoleh borrowing population, advanced
  functions, dan economic activity density. (Locator: Section 5.2, PDF
  p.~17.)
\end{itemize}

\textbf{Status evidensi:} Butir di atas adalah policy recommendations
atau interpretasi penulis berdasarkan hasil model; TXT tidak melaporkan
evaluasi implementasi kebijakan tersebut.

\subsection{Applications}\label{applications-35}

\textbf{Aplikasi yang dinyatakan sumber:} Penulis menyebut hasilnya
sebagai informasi untuk mendorong GDE di China dan sebagai referensi
bagi studi dampak borrowing size di negara lain. (Locator: Section 5.1,
PDF pp.~15-16.)

\textbf{Batas aplikasi:} Tidak ada aplikasi empiris, studi implementasi,
atau evaluasi kebijakan yang dilaporkan di luar rekomendasi tersebut.
Rincian aplikasi operasional: \textbf{Tidak disebutkan dalam file}.

\subsection{Future Research
(Optional)}\label{future-research-optional-35}

\textbf{Agenda yang disebutkan eksplisit oleh penulis:} Menggabungkan
data makro kota, data mikro perusahaan, dan data meso industri dengan
analisis teoretis serta model kuantitatif untuk mencapai integrasi
lintas skala dan menghasilkan kesimpulan penelitian yang lebih bermakna.
(Locator: Section 5.1, PDF p.~16.)

\subsection{Provenance Notes}\label{provenance-notes-35}

\begin{itemize}
\tightlist
\item
  Kode dan status kerja dicocokkan dengan \texttt{Todo.md} serta
  \texttt{output/kajian/kode-kajian-pendahuluan.md}. Filename aktual
  yang diverifikasi sebelum pembacaan adalah
  \texttt{source/journal/A20-guo-liu-zhong-yang-luo-2022-borrowing-size-and-urban-green-development-efficiency-in-the-city-network-of-china-impact-measures-and-size-thresholds-10.3390-land11040493.txt};
  hanya file ini yang menjadi bahan sumber A20.
\item
  TXT menyatakan sumber PDF relatif sebagai
  \texttt{journal/A20-guo-liu-zhong-yang-luo-2022-borrowing-size-and-urban-green-development-efficiency-in-the-city-network-of-china-impact-measures-and-size-thresholds-10.3390-land11040493.pdf},
  diekstrak pada 2026-08-31 dengan \texttt{pdftotext\ -layout}. Blok
  metadata juga menyebut judul, lima penulis, DOI, dan
  \texttt{Pages:\ 20}; teks artikel sendiri menampilkan footer sampai
  \texttt{19\ of\ 19}. Locator review mengikuti nomor halaman artikel
  yang tercetak, yaitu PDF pp.~1-19, dan perbedaan jumlah halaman ini
  tidak diselesaikan dengan asumsi.
\item
  Seluruh teks yang tersedia dalam TXT dibaca, termasuk blok
  \texttt{{[}PDF\ Metadata{]}}, metadata \texttt{pdfinfo}, abstrak, isi
  Sections 1-5, Tables 1-9, Figures 1-3, Author Contributions, Funding,
  Data Availability Statement, Acknowledgments, Conflicts of Interest,
  dan References {[}1{]}-{[}66{]}.
\item
  Dalam batas TXT, teks berakhir setelah References {[}66{]} dan tidak
  menunjukkan bagian artikel yang terpotong; artefak layout yang
  terlihat dicatat di bawah ini dan tidak diisi dengan dugaan.
\item
  Ekstraksi layout memiliki gangguan yang membatasi pembacaan presisi:
  Equation (2) pada PDF pp.~5-6 tampil terduplikasi atau terpecah;
  caption Figure 2 pada PDF p.~8 berulang dan sebagian labelnya tidak
  konsisten; dan Table 8 diberi judul ``Estimation results of
  double-threshold regression'' walaupun Section 4.3 menjelaskannya
  sebagai hasil instrumental-variable. Review mempertahankan isi yang
  dapat dibaca dan menandai gangguan ini, tanpa mengisi bagian yang
  rusak dari luar TXT.
\item
  Ada inkonsistensi terminologi lain: Section 3.2 menyebut GDE sebagai
  ``Explanatory variable'', sedangkan Equation (1), Equation (2), dan
  Table 4 menempatkannya sebagai variabel yang dijelaskan. Klasifikasi
  GDE sebagai dependent variable di review adalah inferensi pengolahan
  berbasis struktur persamaan dan label Table 4.
\item
  Angka, tanda signifikansi, confidence interval, jumlah observasi,
  serta penjelasan mekanisme di review diperlakukan sebagai laporan
  sumber pada locator yang disebutkan. Penjelasan seperti agglomeration
  shadow, siphoning, pollution paradise, dan homogenisasi ditandai
  sebagai interpretasi atau kemungkinan penulis, bukan temuan mekanisme
  yang diverifikasi terpisah.
\item
  Tidak ada sumber eksternal, pembacaan PDF langsung, atau klaim tentang
  studi lain di luar bibliografi yang tercantum dalam TXT yang digunakan
  untuk menulis review ini. Future research tidak ditambahkan di luar
  agenda eksplisit Section 5.1.
\end{itemize}

\section{Review A21}\label{review-a21}

Kode: A21

Label evidensi dalam review ini: \textbf{Laporan sumber} merangkum hal
yang dinyatakan atau ditampilkan dalam TXT; \textbf{Interpretasi
penulis} merangkum penjelasan dan penafsiran Gong dan Zhong;
\textbf{Inferensi pengolahan} menandai kesimpulan tentang batas atau
keterbacaan bukti yang dibuat saat mengolah TXT.

\subsection{Summarized Abstract}\label{summarized-abstract-36}

\textbf{Laporan sumber:} Artikel menguji pengaruh borrowing size
terhadap perkembangan ekonomi kota kecil dan menengah di Tiongkok.
Penelitian menggunakan tiga dimensi borrowing size, data 285 kota
tingkat prefektur, dan data cahaya malam untuk 2004--2013 dalam model
fixed-effects dengan data panel. Pengaruh borrowing size terhadap
perkembangan ekonomi kota kecil dan menengah dilaporkan positif dan
signifikan, sedangkan pengaruhnya terhadap kota besar tidak signifikan
secara umum. Besarnya pengaruh berbeda antarwilayah: di Tiongkok tengah
dan barat, borrowing economic activity density dan borrowing advanced
functions paling positif; di timur, borrowing population paling positif
sementara borrowing advanced functions negatif; di timur laut, borrowing
economic activity density positif. Artikel merekomendasikan peningkatan
aksesibilitas antarkota, penataan lokasi industri, dan dukungan terhadap
pendidikan (Abstract, p.~1; PDF metadata Subject).

\textbf{Interpretasi penulis:} Borrowing size diposisikan sebagai cara
bagi kota kecil dan menengah untuk berbagi manfaat aglomerasi kota besar
melalui jaringan kota, sekaligus sebagai dasar kebijakan koordinasi
regional.

\subsection{Problem Statement}\label{problem-statement-36}

\textbf{Laporan sumber:} Pengembangan kota kecil dan menengah dipandang
penting untuk urbanisasi sehat, integrasi kota--desa, keseimbangan
efisiensi dan keadilan, serta pembangunan regional terkoordinasi. Namun,
perkembangannya tidak seragam di Tiongkok. Kota-kota kecil di Delta
Sungai Yangtze dan Delta Sungai Mutiara disebut mendominasi kelompok
county terbaik, sedangkan ``poverty belt around Beijing and Tianjin''
dan ``lamp shadow area'' masih ada (Section 1, pp.~1--2). Teori
tradisional terutama membahas peran dominan kota besar dalam aglomerasi
dan difusi ke kota kecil, serta kurang memperhatikan peran aktif kota
kecil dan menengah dalam sistem perkotaan (Section 1, p.~2).

\textbf{Laporan sumber:} Artikel merumuskan dua pertanyaan: bagaimana
meningkatkan pertumbuhan ekonomi kota kecil dan menengah, dan apa yang
menyebabkan perbedaan tingkat perkembangan ekonomi kota-kota tersebut
(Section 1, p.~2). Abstrak menyatakan masih kurang penelitian empiris
kuantitatif yang relevan mengenai pembagian manfaat aglomerasi kota
besar melalui borrowing size.

\subsection{Objectives}\label{objectives-36}

\textbf{Laporan sumber:} Tujuan penelitian adalah memperkenalkan konsep
borrowing size dari perspektif jaringan kota nasional dan menganalisis
pengaruhnya terhadap perkembangan ekonomi kota kecil dan menengah.
Penelitian juga bertujuan mengidentifikasi faktor yang menguntungkan dan
merugikan, menemukan kekurangan, serta menyediakan rekomendasi kebijakan
untuk mendukung perkembangan kota, koordinasi urbanisasi baru Tiongkok,
pembangunan regional, efisiensi, dan keadilan (Section 1, pp.~2--3).

\textbf{Laporan sumber:} Secara operasional, penelitian mengukur tiga
dimensi borrowing size---borrowing population, borrowing economic
activity density, dan borrowing advanced functions---kemudian menguji
hipotesis mengenai pengaruh ketiganya serta heterogenitas menurut empat
wilayah utama dan pembagian utara--selatan (Sections 2 dan 4.3,
pp.~3--5, 10--13).

\subsection{Literature Survey}\label{literature-survey-36}

\textbf{Laporan sumber:} Tinjauan pustaka memaparkan perdebatan antara
prioritas pengembangan kota besar dan kota kecil-menengah. Pendukung
kota besar menggunakan argumen aglomerasi dan economies of scale, tetapi
sumber juga mencatat diseconomies of scale seperti diminishing marginal
returns dan polusi lingkungan. Pendukung kota kecil-menengah menilai
pengembangan kota-kota tersebut dapat mempercepat urbanisasi, menurunkan
biaya psikologis tenaga kerja desa yang masuk kota, dan mempertemukan
pasar kota--desa dengan lebih baik (Section 2, pp.~2--3).

\textbf{Laporan sumber:} Literatur tentang urban network dan network
externality digunakan untuk menjelaskan manfaat sinergi dan
komplementaritas antarkota. Borrowing size diperkenalkan oleh Alonso
sebagai perhatian terhadap ukuran kota yang saling berinteraksi, bukan
hanya ukuran optimal satu kota. Kajian terdahulu juga menghubungkan
borrowing size dengan kemampuan kota kecil di sekitar London dan
perkembangan perusahaan di kota kecil. Sumber membedakan tiga bentuk
kota kecil yang dapat meminjam ukuran: yang berdekatan dengan kota
besar, kelompok kota kecil yang saling berdekatan, dan kota kecil yang
tidak berdekatan dengan kota besar maupun kelompok kota kecil tetapi
dapat meminjam melalui jaringan kota (Section 2, pp.~3--4).

\textbf{Laporan sumber:} Tinjauan juga menekankan bahwa eksternalitas
jaringan kota dapat negatif. Agglomeration shadows muncul terutama dari
kompetisi antarkota; kota kecil menghadapi tekanan dari kota besar dan
kota kecil tetangga. Karena itu borrowing size dan agglomeration shadows
dapat hadir bersamaan, dengan intensitas yang berbeda untuk tiap kota
(Section 2, p.~3).

\textbf{Interpretasi penulis:} Artikel mendefinisikan borrowing size
sebagai proses ketika suatu kota memperoleh manfaat skala dengan
berpartisipasi dalam jaringan dengan kota atau wilayah lain. Kerangka
teorinya menggabungkan spatial interaction theory, absolute advantage
theory, comparative advantage theory, flow space theory, dan network
externality theory. Tiga hipotesisnya adalah: borrowing population
berpengaruh positif tetapi relatif tidak signifikan; borrowing economic
activity density berpengaruh positif; dan borrowing advanced functions
memiliki pengaruh ganda, yaitu dapat positif maupun negatif (Sections 2,
pp.~3--5).

\subsection{Method Used}\label{method-used-36}

\textbf{Laporan sumber:} Penelitian menggunakan analisis empiris
kuantitatif dengan model fixed-effects dan data panel untuk data
sosial-ekonomi serta cahaya malam kota tingkat prefektur pada
2004--2013. Dari 285 kota, 254 kota kecil dan menengah menjadi objek
penelitian dan 31 kota besar menjadi pembanding (Section 4.2,
pp.~9--10).

\textbf{Laporan sumber:} Tiga ukuran borrowing size dihitung sebagai
jumlah nilai suatu dimensi di kota lain yang dibobot dengan jarak jalan
terpendek: borrowing population (\texttt{bpop}), borrowing economic
activity density (\texttt{bden}), dan borrowing advanced functions
(\texttt{bfun}). Matriks jarak yang digunakan berukuran 285 × 285,
dengan 81.225 nilai jarak dari rute jalan antarkota (Section 3.1, p.~6;
Section 3.2, pp.~6--7).

\textbf{Laporan sumber:} Economic activity density dihitung dari total
tahunan penumpang kereta api, jalan raya, air, dan udara di dalam kota,
dibagi luas kawasan terbangun dan 365 hari. Advanced functions diukur
dengan rasio jumlah pekerja pada layanan tingkat tinggi terhadap jumlah
pekerja manufaktur; industri layanan tingkat tinggi yang dicantumkan
meliputi transportasi, pergudangan dan pengiriman, transmisi informasi,
layanan komputasi dan perangkat lunak, keuangan, real estat, leasing dan
jasa bisnis, riset ilmiah, layanan teknis, serta prospecting geologi
(Section 3.2, pp.~6--7).

\textbf{Laporan sumber:} Persamaan benchmark (Equation 4, Section 3.3,
p.~6) memasukkan lag satu periode untuk bpop, bfun, bden, log populasi
(\texttt{lnpop}), kuadrat log populasi (\texttt{lnpop2}), investasi aset
tetap (\texttt{pinvest}), foreign direct investment (\texttt{pfdi}), dan
modal manusia lokal (\texttt{pedu}). Semua variabel independen dilag
untuk mengurangi kemungkinan kausalitas timbal balik dengan
produktivitas perkotaan; variabel populasi dilogaritmakan untuk
mengurangi heteroskedastisitas (Section 3.3, p.~6).

\textbf{Laporan sumber:} Bagian 3.4 juga mendefinisikan \texttt{fun}
sebagai advanced functions lokal, \texttt{den} sebagai economic activity
density lokal, \texttt{pinvest} sebagai tingkat investasi aset tetap,
\texttt{pfdi} sebagai proporsi foreign direct investment terhadap GDP
tahun berjalan, \texttt{pedu} sebagai proporsi mahasiswa perguruan
tinggi terhadap total penduduk, dan \texttt{pgov} sebagai proporsi
pengeluaran pemerintah dalam GDP tahun berjalan (Section 3.4, pp.~6--7).
Table 1 menampilkan \texttt{Lpgov} bersama variabel kontrol.

\textbf{Inferensi pengolahan:} Equation 4 tidak menampilkan
\texttt{pgov}, meskipun \texttt{pgov} dijelaskan di Section 3.4 dan
\texttt{Lpgov} muncul di Table 1. TXT tidak menjelaskan perbedaan ini,
sehingga review tidak menganggap \texttt{pgov} sebagai bagian dari
Equation 4 tanpa catatan.

\textbf{Laporan sumber:} Perubahan spasial-temporal borrowing size dan
cahaya malam ditampilkan dan dianalisis menggunakan ArcGIS, sedangkan
tiga dimensi borrowing size diukur menggunakan MATLAB (Section 3.1,
p.~6). Uji heterogenitas membagi kota menurut timur laut, timur, tengah,
barat, serta utara dan selatan (Section 4.3, pp.~10--13).

\subsection{Dataset}\label{dataset-36}

\textbf{Laporan sumber:} Data cahaya malam berasal dari 34 laporan
selama periode 22 tahun 1992--2013 yang diterbitkan National Geophysical
Data Center dalam National Oceanic and Atmospheric Administration. Data
mencakup cahaya kota, town, lokasi dengan penerangan kontinu, dan gas
flares; kebakaran, pantulan sinar matahari, aurora, dan noise latar
diproses atau dihapus. Rentang DN asli adalah 1--63 dengan background
noise bernilai nol (Section 3.1, p.~6).

\textbf{Laporan sumber:} Data cahaya malam untuk 2004--2013 diproses
melalui continuity correction, pixel saturation, dan abnormal value
processing. Nilai grid dijumlahkan berdasarkan batas kota tingkat
prefektur, kemudian dihitung nilai DN per kapita. Penulis menyatakan
hasil olahan ini lebih sesuai untuk perbandingan lintas tahun dan lintas
wilayah (Section 3.1, p.~6).

\textbf{Laporan sumber:} Data ekonomi terutama berasal dari China City
Statistical Yearbook edisi 2005--2014. Kota Haidong, Tongren, Bijie,
Lhasa, dan Sansha dikeluarkan karena data tidak lengkap. Nama asli kota
yang mengalami perubahan administratif dipertahankan untuk menjaga
integritas sampel (Section 3.1, p.~6).

\textbf{Laporan sumber:} Jarak antarkota diperoleh dari 2345 Intercity
Highway Distance Search Tool dan digunakan untuk membentuk matriks jarak
jalan terpendek. Pernyataan ketersediaan data menyebut sumber data
publik telah dijelaskan dalam artikel, sedangkan data lain dapat diminta
kepada penulis korespondensi dengan alasan yang wajar (Section 3.1,
p.~6; Data Availability Statement, p.~16).

\subsection{Population / Sample}\label{population-sample-36}

\textbf{Laporan sumber:} Unit analisis adalah kota tingkat prefektur,
bukan municipal district, di Tiongkok. Sampel akhir berisi 285 kota
tingkat prefektur setelah penghapusan kota dengan data ekonomi tidak
lengkap. Dalam analisis regresi, 254 kota kecil dan menengah dianalisis,
sedangkan 31 kota besar digunakan sebagai pembanding (Section 3.1, p.~6;
Section 4.2, pp.~9--10).

\textbf{Laporan sumber:} Kota kecil dan menengah dalam studi tidak
diklasifikasikan dengan ambang populasi. Sumber menyatakan bahwa kota
selain ibu kota provinsi, munisipalitas, dan Shenzhen diklasifikasikan
sebagai kota kecil dan menengah untuk penelitian ini. Untuk uji
heterogenitas, Table 2 memuat 73 kota barat, 74 kota tengah, 76 kota
timur, dan 31 kota timur laut; pembagian kota kecil-menengah utara
berjumlah 116 dan selatan 138. Jumlah observasinya berturut-turut adalah
657, 666, 684, 279, 1.044, dan 1.242 (Sections 4.3 dan 5.2, pp.~10--15;
Table 2, p.~11).

\textbf{Laporan sumber:} Untuk pembagian geografis atas 285 kota, kota
ditempatkan di wilayah selatan jika ibu kota provinsinya berada di bawah
32° N dan di wilayah utara jika berada di atas 32° N. Daftar yang
diberikan sumber menghasilkan 154 kota selatan dan 131 kota utara,
sedangkan Table 2 melaporkan subset kota kecil-menengah untuk regresi,
yaitu 138 kota selatan dan 116 kota utara (Section 4.3, pp.~10--11;
Table 2, p.~11).

\textbf{Laporan sumber:} Populasi statistik di luar 285 kota tingkat
prefektur yang diteliti tidak disebutkan dalam file.

\subsection{Dependent Variables}\label{dependent-variables-36}

\textbf{Laporan sumber:} Variabel dependen atau explained variable
adalah \texttt{lightp}, yaitu kecerahan cahaya malam per kapita, yang
digunakan untuk menunjukkan tingkat perkembangan ekonomi kota. Dalam
penjelasan Equation 4, \texttt{lightp} juga disebut sebagai urban
productivity (Sections 3.3--3.4, pp.~6--7). Tidak ada variabel dependen
lain yang dinyatakan untuk analisis regresi utama.

\subsection{Results}\label{results-36}

\textbf{Laporan sumber:} Figure 1 menunjukkan distribusi
spasial-temporal kecerahan cahaya malam per kapita pada 2004, 2009, dan
2013. Kota dengan nilai tinggi terutama berada di wilayah maju atau
wilayah dengan perkembangan ekonomi rata-rata tetapi populasi relatif
kecil. Kota dengan DN per kapita di atas 20 berjumlah empat pada 2004,
lima yang disebutkan pada 2009, dan 15 pada 2013. Kota bernilai tinggi
cenderung berada di kawasan pesisir, bagian tengah dan hilir Sungai
Yangtze, serta kawasan Bohai Bay (Section 4.1, p.~7; Figure 1).

\textbf{Laporan sumber:} Figure 2 menunjukkan borrowing population
terutama terkonsentrasi di North China Plain dan Middle and Lower
Yangtze Plains. Nilai kurang dari 50 unit terutama berada di wilayah
perbatasan barat laut, termasuk Urumqi dan Karamay. Distribusinya
relatif stabil dengan sedikit peningkatan sepanjang waktu (Section 4.1,
p.~8; Figure 2).

\textbf{Interpretasi penulis:} Naik-turunnya borrowing advanced
functions pada Figure 3 dijelaskan melalui perubahan relatif kekuatan
layanan tingkat tinggi dan manufaktur. Penulis mengaitkan kondisi 2009
dengan dampak krisis keuangan global terhadap manufaktur, serta kondisi
2013 dengan pemulihan ekonomi global dan perhatian Tiongkok terhadap
ekonomi riil. Borrowing advanced functions pada 2013 tetap lebih rendah
daripada 2009 karena layanan tingkat tinggi dinilai masih relatif lemah
terhadap manufaktur (Section 4.1, pp.~8--9; Figure 3).

\textbf{Laporan sumber:} Figure 4 menunjukkan jumlah kota dengan
borrowing economic activity density di atas 0,14 unit meningkat dari
tiga pada 2004 menjadi 12 pada 2009 dan 46 pada 2013. Nilai tinggi
berada di bagian tengah wilayah studi, sedangkan nilai rendah umumnya
berada di tepi wilayah (Section 4.1, p.~9; Figure 4).

\textbf{Laporan sumber:} Pada Table 1, untuk 254 kota kecil dan
menengah, model terpisah melaporkan koefisien borrowing population
sebesar 0,341 (fe1), borrowing advanced functions sebesar 47,77 (fe2),
dan borrowing economic activity density sebesar 70,50 (fe3), semuanya
bertanda \texttt{***}. Dalam model bersama fe4, koefisien ketiganya
masing-masing 0,164, 19,93, dan 41,17, semuanya bertanda \texttt{***}.
Goodness of fit untuk model kota kecil-menengah berada pada R-squared
0,326--0,401, dengan 2.286 observasi dan 254 kota. Untuk 31 kota besar
pada fe5, borrowing population sebesar 0,255 bertanda \texttt{**},
sedangkan borrowing advanced functions sebesar 4,845 dan borrowing
economic activity density sebesar 28,39 tidak memiliki tanda
signifikansi. Model kota besar memiliki 279 observasi dan R-squared
0,652. Catatan tabel menetapkan \texttt{***} untuk p \textless{} 0,01,
\texttt{**} untuk p \textless{} 0,05, dan \texttt{*} untuk p \textless{}
0,1 (Section 4.2, pp.~9--10; Table 1, p.~9).

\textbf{Laporan sumber:} Untuk kota kecil dan menengah, koefisien lnpop
positif dan lnpop2 negatif, keduanya signifikan pada tingkat 1\%. Sumber
menafsirkan pola ini sebagai hubungan berbentuk U terbalik antara
populasi lokal dan perkembangan ekonomi: ekspansi populasi membantu pada
tahap awal, tetapi dapat merugikan pada tahap lanjut (Section 4.2,
p.~10; Table 1, p.~9).

\textbf{Laporan sumber:} Table 2 melaporkan pola berikut untuk kota
kecil dan menengah. Di barat, bfun = 20,66 (\texttt{**}) dan bden =
43,78 (\texttt{***}), sedangkan bpop = 0,0552 tidak signifikan. Di
wilayah tengah, bfun = 19,51 (\texttt{***}) dan bden = 42,33
(\texttt{***}), sedangkan bpop = 0,0009 tidak signifikan. Di timur, bpop
= 0,482 (\texttt{***}), bfun = −19,04 (\texttt{*}), dan bden = 9,66
tidak signifikan. Di timur laut, bden = 117,9 (\texttt{***}), sedangkan
bpop = 0,16 dan bfun = −1 tidak signifikan. Di utara, bfun = 29,58
(\texttt{***}) dan bden = 50,72 (\texttt{***}), sedangkan bpop = 0,0557
tidak signifikan. Di selatan, bpop = 0,273 (\texttt{***}) dan bfun =
−9,947 (\texttt{***}), sedangkan bden = 7,371 tidak signifikan (Section
4.3, pp.~10--13; Table 2, p.~11).

\textbf{Interpretasi penulis:} Hasil Table 2 dikaitkan dengan perbedaan
kemampuan inovasi, tingkat industrialisasi, aksesibilitas, kepadatan
penduduk, marketization, struktur ekonomi, dan kondisi geografis
antarwilayah. Penulis menilai bahwa pusat dan barat terutama membutuhkan
hubungan fungsi maju dan limpahan pengetahuan; wilayah timur lebih mampu
mengubah populasi yang dipinjam menjadi keuntungan ekonomi tetapi
menghadapi kompetisi fungsi maju; timur laut lebih memperoleh manfaat
dari kepadatan aktivitas ekonomi; dan perbedaan utara--selatan
mencerminkan perbedaan struktur ekonomi serta perkembangan regional
(Section 4.3, pp.~10--13).

\subsection{Findings}\label{findings-36}

\textbf{Interpretasi penulis:} Temuan utama artikel adalah bahwa ketiga
dimensi borrowing size dapat membantu perkembangan ekonomi kota kecil
dan menengah secara agregat. Manfaatnya tidak seragam: dimensi yang
efektif bergantung pada konteks regional. Bagi kota besar, sumber
menyatakan manfaat yang terdeteksi terutama berasal dari borrowing
population, bukan borrowing economic activity density atau borrowing
advanced functions (Sections 4.2--4.3, pp.~9--13).

\textbf{Interpretasi penulis:} Mekanisme yang diajukan untuk borrowing
population adalah perluasan pasar tenaga kerja dan pasar konsumen.
Borrowing economic activity density dipahami sebagai saluran limpahan
pengetahuan, teknologi, dan informasi. Borrowing advanced functions
dipahami sebagai pembagian kerja antarruang dan akses pada fungsi
budaya, pendidikan, hiburan, serta layanan; tetapi fungsi maju juga
dapat meningkatkan kompetisi dan homogenisasi industri (Section 2,
pp.~4--5).

\textbf{Interpretasi penulis:} Dalam penjelasan hasil, kota kecil dan
menengah di tengah dan barat dinilai lemah dalam kapasitas inovasi dan
industrialisasi, kota timur menghadapi persaingan fungsi maju, timur
laut menghadapi industri tradisional yang menua dan overcapacity,
sedangkan kota utara lebih didominasi sumber daya, energi, dan industri
berat. Ini adalah penjelasan penulis atas pola regresi, bukan pengujian
mikro terhadap individu atau perusahaan (Sections 4.3 dan 5.2,
pp.~10--15).

\subsection{Conclusions}\label{conclusions-36}

\textbf{Laporan sumber:} Section 5.1 merumuskan tiga kesimpulan.
Pertama, borrowing population lebih terkonsentrasi di North China Plain
dan Middle and Lower Yangtze Plain serta hanya sedikit meningkat; jumlah
kota dengan tingkat borrowing advanced functions yang lebih tinggi
mengalami pola turun--naik--turun; borrowing economic activity density
bergantung pada lokasi dan jumlah kota dengan nilai tinggi meningkat.
Kedua, pengaruh agregat borrowing population, borrowing advanced
functions, dan borrowing economic activity density terhadap perkembangan
ekonomi kota kecil dan menengah signifikan positif, sementara pada kota
besar pengaruh yang terutama terlihat adalah borrowing population.
Ketiga, pengaruh ketiga dimensi berbeda secara signifikan menurut
wilayah (Sections 5.1, pp.~13--14).

\textbf{Laporan sumber:} Untuk pengujian hipotesis, penelitian mendukung
bagian positif Hypothesis 1 tetapi tidak membuktikan bahwa pengaruhnya
relatif tidak signifikan. Hypothesis 2 dikonfirmasi. Hypothesis 3
memperoleh dukungan untuk pengaruh positif borrowing advanced functions,
tetapi hubungan negatifnya tidak dikonfirmasi (Section 5.1, pp.~13--14).

\textbf{Inferensi pengolahan:} Ada ketidaksesuaian internal yang perlu
dipertahankan. Uraian Section 4.3 dan Table 2 menyatakan bden di timur
tidak signifikan dan bden di timur laut positif signifikan, tetapi satu
kalimat pada kesimpulan menyebut pola paling positif itu berada di
``eastern and northern regions'' (p.~14). Review ini tidak mengganti
istilah sumber atau menyimpulkan mana yang benar.

\subsection{Contributions}\label{contributions-36}

\textbf{Interpretasi penulis:} Artikel menyatakan bahwa pengukuran
borrowing size secara multidimensi dan pengujian empiris dengan
fixed-effects model pada data panel memperkaya penelitian kuantitatif
tentang pengaruh borrowing size terhadap perkembangan kota kecil dan
menengah. Artikel juga menyatakan bahwa hasilnya dapat menjadi rujukan
bagi penelitian di negara berkembang lain dan memperkaya teori terkait
hubungan borrowing size dan perkembangan kota kecil-menengah (Section
5.2, pp.~14--15).

\textbf{Laporan sumber:} Kontribusi empiris yang benar-benar dilakukan
dalam artikel adalah membangun jaringan kota dari 285 kota tingkat
prefektur, mengukur tiga dimensi borrowing size, memetakan
distribusinya, dan membandingkan pengaruhnya antarwilayah (Sections
3--5, pp.~6--15).

\subsection{Research Gap}\label{research-gap-36}

\textbf{Laporan sumber:} Berdasarkan tinjauan penulis, terdapat empat
kesenjangan penelitian. Pertama, studi sebelumnya lebih banyak berobjek
negara maju di Eropa dan Amerika pada tahap akhir urbanisasi, sedangkan
studi mengenai negara berkembang, khususnya Tiongkok, masih sedikit.
Kedua, belum ada definisi borrowing size yang diterima secara universal
karena peneliti cenderung mendefinisikannya sesuai kebutuhan studi
masing-masing. Ketiga, studi di Tiongkok terutama berfokus pada urban
agglomeration tingkat nasional dan masih sedikit yang menganalisis empat
wilayah utama serta perbedaan tingkat perkembangan utara--selatan.
Keempat, penentuan bobot jarak di Tiongkok terutama menggunakan jarak
Euclidean atau waktu tempuh kereta api, dan kedua metode tersebut
dinilai memiliki keterbatasan (Section 2, pp.~3--4).

\subsection{Limitations}\label{limitations-36}

\textbf{Laporan sumber:} Penulis menyatakan empat keterbatasan. Pertama,
bagian teori belum membangun model matematis yang sesuai dari teori
klasik dan baru melakukan analisis kualitatif. Kedua, pengaruh borrowing
size pada dasarnya bergantung pada subjek mikro seperti individu dan
perusahaan, tetapi penelitian ini tidak berada pada tingkat mikro.
Ketiga, cakupan terutama kota tingkat prefektur; klasifikasi kota
kecil-menengah tidak memakai standar populasi, dan kota tingkat
county---terutama kota kecil---belum diteliti secara rinci. Keempat,
bobot jarak terutama memakai jarak jalan terpendek yang tersedia pada
September 2019, sehingga tidak mencerminkan peningkatan jalan,
densifikasi jaringan jalan, penambahan penerbangan dan kereta api, serta
jaringan komunikasi antarkota (Section 5.2, pp.~14--15).

\subsection{Challenges}\label{challenges-36}

\textbf{Laporan sumber:} Tidak disebutkan dalam file sebagai subbagian
bernama ``Challenges''. Kendala yang dilaporkan meliputi data kota yang
tidak lengkap sehingga lima kota dikeluarkan, kebutuhan melakukan
continuity correction, pixel saturation, dan abnormal value processing
pada data cahaya malam, serta perubahan nama administratif yang harus
dipertahankan agar sampel tetap konsisten (Section 3.1, p.~6). Sumber
juga menghadapi keterbatasan dalam merepresentasikan jaringan kota yang
berubah karena matriks jarak hanya memakai jarak jalan terpendek dan
tidak memasukkan jaringan komunikasi (Section 5.2, pp.~14--15).

\textbf{Inferensi pengolahan:} Karena sumber tidak memisahkan
``Challenges'' dari keterbatasan metodologis dan pengolahan data, bagian
ini hanya mengulang kendala yang dilaporkan dan tidak menambahkan
tantangan baru.

\subsection{Practical Implications}\label{practical-implications-36}

\textbf{Laporan sumber:} Artikel mengusulkan tiga kelompok rekomendasi
kebijakan. Pertama, pemerintah perlu meningkatkan jaringan jalan dan
aksesibilitas antarkota agar kota kecil-menengah dapat meminjam ukuran,
mengakses pasar produk dan faktor, serta memperlancar pertukaran bahan
baku, modal, tenaga kerja, pengetahuan, teknologi, dan informasi
(Section 5.2, pp.~14--15).

\textbf{Laporan sumber:} Kedua, tata letak industri perlu dirancang
melalui pembagian kerja yang saling melengkapi. Kota pusat disarankan
memperkuat fungsi layanan tingkat tinggi dan teknologi inti. Kota
kecil-menengah disarankan memanfaatkan biaya lahan dan tenaga kerja yang
rendah, memperbaiki lingkungan usaha, mengembangkan industri yang sesuai
kondisi lokal, dan mempertahankan basis manufaktur. Kota dengan tata
letak industri yang tidak efisien dapat menggunakan jaringan kota untuk
memperoleh fungsi komplementer (Section 5.2, pp.~14--15).

\textbf{Laporan sumber:} Ketiga, pemerintah perlu mengembangkan dan
menarik talenta. Bentuk yang disebutkan meliputi konsultasi, kerja sama
proyek, reemployment setelah pensiun, talent leasing, perekrutan lulusan
perguruan tinggi, dan penyediaan akomodasi yang sesuai bagi talenta yang
direkrut (Section 5.2, pp.~14--15).

\subsection{Applications}\label{applications-36}

\textbf{Laporan sumber:} Penerapan yang dilakukan dalam file adalah
penggunaan jaringan 285 kota tingkat prefektur untuk mengukur tiga
dimensi borrowing size, memetakan distribusi spasial-temporalnya pada
2004, 2009, dan 2013, menguji pengaruhnya terhadap \texttt{lightp},
serta membandingkan wilayah barat, tengah, timur, timur laut, utara, dan
selatan (Sections 3--4, pp.~6--13). Hasilnya kemudian digunakan sebagai
dasar rekomendasi mengenai aksesibilitas, tata letak industri, dan
talenta (Section 5.2, pp.~14--15).

\textbf{Laporan sumber:} Penerapan di luar konteks Tiongkok, di luar
rentang 2004--2013, atau di luar kota tingkat prefektur: Tidak
disebutkan dalam file sebagai hasil empiris.

\subsection{Future Research
(Optional)}\label{future-research-optional-36}

\textbf{Laporan sumber:} Penulis menyebutkan empat arah penelitian
lanjutan. Pertama, teori klasik dapat dimodelkan secara matematis untuk
membuat argumen lebih rigor. Kedua, penelitian dapat dilakukan dari
perspektif individu dan perusahaan. Ketiga, pengumpulan data tingkat
county dapat membuat penelitian lebih rinci. Keempat, informasi
pembangunan jaringan transportasi antarkota dan koneksi jaringan
informasi perlu dikumpulkan untuk membangun tabel bobot jarak yang
komprehensif dan dinamis (Section 5.2, pp.~15--16).

\subsection{Provenance Notes}\label{provenance-notes-36}

\textbf{Laporan sumber:} Verifikasi prefiks A21 menemukan tepat satu TXT
aktual dengan path
\texttt{source/journal/A21-gong-zhong-2021-the-impact-of-borrowing-size-on-the-economic-development-of-small-and-medium-sized-cities-in-china-10.3390-land10020134.txt}.
Blok metadata menyebut Source PDF
\texttt{journal/A21-gong-zhong-2021-the-impact-of-borrowing-size-on-the-economic-development-of-small-and-medium-sized-cities-in-china-10.3390-land10020134.pdf},
tanggal ekstraksi 2026-08-31, dan metode \texttt{pdftotext\ -layout}.
Metadata PDF mencantumkan judul, penulis Xiaoxia Gong dan Fanglei Zhong,
DOI \texttt{10.3390/land10020134}, serta \texttt{Pages:\ 19}; teks
artikel memiliki footer sampai \texttt{18\ of\ 18} (PDF Metadata; title
page; pp.~1--18).

\textbf{Laporan sumber:} Teks juga memuat metadata publikasi: Land 2021,
10, 134; diterima 28 Januari 2021; diterbitkan 31 Januari 2021; serta
lisensi CC BY (title page, p.~1). Informasi pendanaan, kontribusi
penulis, ketersediaan data, ucapan terima kasih, dan konflik kepentingan
tercantum pada pp.~15--16.

\textbf{Inferensi pengolahan:} Seluruh teks yang tersedia dalam TXT,
termasuk blok \texttt{{[}PDF\ Metadata{]}}, bagian artikel, tabel,
gambar yang berbentuk caption, dan daftar referensi, telah dibaca. TXT
menggunakan pemisah form-feed dan memiliki beberapa karakter kontrol
serta perataan whitespace akibat ekstraksi layout. Isi substantif dan
locator utama terbaca, tetapi perataan tabel dan satu konflik istilah
wilayah pada bagian kesimpulan dipertahankan sebagai batas evidence.
Rujukan bibliografis hanya dibaca sebagai bagian dari TXT dan tidak
dibuka untuk menambah atau memverifikasi bukti. PDF tidak digunakan
sebagai sumber bukti tambahan; review ini hanya menggunakan TXT.

\section{Review A22 - Malý (2016)}\label{review-a22---maluxfd-2016}

Status: Review literatur berbasis seluruh teks TXT hasil ekstraksi PDF.

\subsection{Identitas Artikel}\label{identitas-artikel-7}

\begin{itemize}
\tightlist
\item
  \textbf{Kode:} A22.
\item
  \textbf{Judul:} \emph{Small Towns in the Context of ``Borrowed Size''
  and ``Agglomeration Shadow'' Debates: The Case of the South Moravian
  Region (Czech Republic)}.
\item
  \textbf{Penulis:} Jiří Malý.
\item
  \textbf{Tahun:} 2016.
\item
  \textbf{Jurnal dan lokasi:} \emph{European Countryside}, 4 (2016),
  hlm. 333-350. Header TXT tidak merinci apakah angka 4 merupakan volume
  atau nomor.
\item
  \textbf{DOI:} 10.1515/euco-2016-0024.
\item
  \textbf{Tanggal naskah:} diterima 22 Februari 2016 dan diterima untuk
  publikasi 19 Juli 2016.
\item
  \textbf{Kata kunci:} small towns, borrowed size, agglomeration shadow,
  service function, periphery, metropolitan area.
\item
  \textbf{Pendanaan:} Specific research project MUNI/A/1315/2015,
  ``Integrovaný výzkum environmentálních změn v krajinné sféře Země''.
\item
  \textbf{Wilayah dan objek:} small centres di South Moravian Region,
  Czech Republic, dalam sistem perkotaan yang dipengaruhi Brno.
\item
  \textbf{Locator identitas:} hlm. 333, blok judul, abstrak, afiliasi,
  dan DOI; hlm. 345-350, bagian penutup dan referensi.
\end{itemize}

\subsection{Summarized Abstract}\label{summarized-abstract-37}

\textbf{Laporan sumber.} Artikel menguji sejauh mana penyediaan layanan
pada small towns ditentukan oleh lokasi mereka dalam sistem perkotaan
regional yang kuat dipengaruhi kawasan metropolitan. Kasus yang dipilih
ialah South Moravian Region di Republik Ceko dengan Brno sebagai pusat
utama. Penulis menyusun indeks fungsi pelayanan dari 30 jenis layanan
pada 672 municipality, memilih 113 small centres, lalu menghubungkan
variasi fungsi pelayanan dengan ukuran penduduk, posisi dalam sistem
perkotaan, akses transportasi, sejarah, daya tarik wisata, dan karakter
wilayah. Hubungan tersebut dievaluasi dengan regresi logistik
multinomial. (Hlm. 333, abstrak; hlm. 337-340, bagian 3.)

Hasil menunjukkan bahwa layanan mengikuti hierarki threshold dan range,
tetapi small centres berukuran penduduk serupa dapat memiliki fungsi
pelayanan yang sangat berbeda. Fungsi pelayanan lebih tinggi pada
sejumlah pusat yang lebih perifer dan menarik secara wisata, sedangkan
kedekatan dengan pusat yang lebih besar dapat berkaitan dengan fungsi
yang lebih rendah. Di Brno Metropolitan Area (BMA), koneksi dengan pusat
besar dan kemampuan menarik komuter atau kegiatan komersial berkaitan
dengan fungsi pelayanan yang lebih tinggi pada sebagian pusat. (Hlm.
341-345, bagian 4; hlm. 345-346, bagian 5.)

\textbf{Interpretasi penulis.} Penulis menyimpulkan bahwa borrowed size
dan agglomeration shadow dapat berlangsung bersamaan. Sebuah small
centre dapat meminjam fungsi atau kinerja dari jaringan metropolitan
dalam satu aspek, tetapi mengalami kekurangan fungsi pelayanan karena
persaingan dengan pusat yang lebih besar dalam aspek lain. Daya tarik
wisata dan komersial tempat tertentu juga dipandang penting bagi hasil
akhir fungsi pelayanan. (Hlm. 333, abstrak; hlm. 345-346, bagian 5.)

\textbf{Inferensi pengolahan.} Bukti paling langsung dalam TXT ialah
variasi indeks fasilitas dan asosiasi statistik antara kategori fungsi
pelayanan dengan karakter spasial. Kategori borrowed size dan
agglomeration shadow merupakan interpretasi atas pola tersebut; desain
yang dilaporkan tidak mengukur secara langsung peminjaman fungsi oleh
pengguna layanan atau mengidentifikasi efek kausal. (Hlm. 340, bagian 3;
hlm. 343-345, bagian 4.)

\subsection{Problem Statement}\label{problem-statement-37}

\textbf{Laporan sumber.} Teori sistem perkotaan yang menekankan hierarki
dan kota besar, serta banyak teori urban kontemporer yang berfokus pada
kawasan padat dan aglomerasi sebagai kutub pertumbuhan, cenderung
mengabaikan peran small towns walaupun kota-kota tersebut penting bagi
wilayah perdesaan. Perubahan akibat globalisasi, mobilitas, jaringan
kota, dan hubungan horizontal juga membuat ukuran penduduk tidak selalu
sejalan dengan fungsi urban yang diharapkan. (Hlm. 334-335, bagian 1.)

Artikel menyatakan bahwa konsep borrowed size dan agglomeration shadow
telah membantu menjelaskan pengaruh eksternalitas jaringan terhadap
fungsi kota, tetapi hanya sedikit memberi perhatian pada fungsi
pelayanan small towns. Banyak studi terdahulu bekerja pada skala
internasional, menekankan fungsi berorde tinggi, dan memakai unit
spasial yang lebih besar. Karena itu, fungsi pelayanan small towns pada
skala regional dan spektrum layanan yang lebih luas belum cukup
dipahami. (Hlm. 333-334, abstrak dan bagian 1.)

\textbf{Interpretasi penulis.} Persoalan empirisnya ialah menentukan
apakah small towns yang berada dekat atau terhubung dengan pusat
metropolitan memperoleh keuntungan jaringan, atau justru kehilangan
fungsi karena kompetisi dan agglomeration shadow. Posisi dalam sistem
perkotaan, jaringan transportasi, sejarah, dan daya tarik wisata
diperlakukan sebagai kemungkinan penjelas di luar ukuran penduduk. (Hlm.
334, bagian 1; hlm. 336-337, bagian 2.)

\textbf{Inferensi pengolahan.} Problem artikel menggabungkan dua
pertanyaan yang berbeda: bagaimana mengukur fungsi pelayanan sebuah
municipality dan bagaimana menafsirkan perbedaan skor sebagai borrowed
size atau agglomeration shadow. Pertanyaan pengukuran dijawab dengan
hitungan fasilitas berbobot; pertanyaan mekanisme dijawab melalui
asosiasi regresi dan pembacaan pola spasial, sehingga tingkat kepastian
keduanya tidak sama.

\subsection{Objectives}\label{objectives-37}

\textbf{Laporan sumber.} Tujuan utama penelitian adalah mengevaluasi
fungsi pelayanan small towns di South Moravian Region dengan
memperhatikan:

\begin{itemize}
\tightlist
\item
  lokasi dalam jaringan perkotaan;
\item
  lokasi dan keterhubungan dalam jaringan transportasi;
\item
  kepentingan historis;
\item
  daya tarik wisata; dan
\item
  hubungan dengan pusat metropolitan Brno.
\end{itemize}

Penulis mengajukan hipotesis bahwa fungsi pelayanan dipengaruhi bukan
hanya oleh ukuran penduduknya sendiri, melainkan juga oleh posisi town
terhadap kota inti kawasan metropolitan dan infrastruktur transportasi.
Tujuan operasionalnya diterjemahkan menjadi empat langkah: menetapkan
layanan dan nilainya, menentukan small centres, menetapkan variabel yang
berpotensi memengaruhi fungsi pelayanan, serta mengevaluasi hubungan
antarvariabel dengan regresi. (Hlm. 334 dan 337, bagian 1 dan 3.)

\textbf{Inferensi pengolahan.} Artikel juga secara implisit bertujuan
menguji apakah dua proses yang sering dipertentangkan dapat muncul dalam
satu sistem regional. Namun, tujuan tersebut lebih tepat dipahami
sebagai pengujian pola konseptual melalui indikator pelayanan, bukan
pengujian eksperimental atas mekanisme jaringan.

\subsection{Literature Survey}\label{literature-survey-37}

\textbf{Laporan sumber.} Tinjauan pustaka artikel bergerak melalui
beberapa rangkaian konsep:

\begin{itemize}
\tightlist
\item
  \emph{Central Place Theory} Christaller dan model hierarkis
  menghubungkan ukuran, tingkat pusat, threshold, range, dan fungsi
  pelayanan. Kerangka ini dipakai untuk menjelaskan mengapa layanan
  dasar tersebar lebih luas, sedangkan layanan berorde tinggi lebih
  terkonsentrasi.
\item
  Model jaringan kota dan polisentrisitas menekankan interaksi
  antarkota, spesialisasi, hubungan horizontal, serta dimensi fungsional
  sistem perkotaan. Dalam kerangka ini, posisi sebuah tempat dan ukuran
  hinterland-nya dapat sama pentingnya dengan ukuran tempat itu sendiri.
\item
  Borrowed size dijelaskan melalui gagasan bahwa tempat yang lebih kecil
  dapat memanfaatkan kepadatan, sumber daya, atau hubungan dengan
  kawasan urban yang lebih besar dan menampung fungsi yang tidak akan
  muncul jika berdiri sendiri. Artikel mengaitkan pembahasan ini dengan
  Alonso, Phelps dan kolega, serta Meijers dan kolega.
\item
  Agglomeration shadow dipaparkan sebagai kemungkinan sisi kompetitif
  dari kedekatan dengan pusat besar: tempat yang dekat pusat lebih besar
  dapat memiliki lebih sedikit fungsi dibanding tempat berukuran sama
  yang lebih terisolasi. Artikel membandingkan bukti tentang amenitas
  budaya, pertumbuhan penduduk, interaksi komuter, dan fungsi perkotaan.
\item
  Literatur tentang small towns dan pembangunan perdesaan menempatkan
  kota kecil sebagai penyedia pekerjaan dan layanan bagi hinterland.
  Perdebatan territorial cohesion dan dokumen perencanaan Eropa serta
  Ceko digunakan untuk menggarisbawahi arti aksesibilitas pusat urban
  dan layanan kepentingan umum.
\item
  Artikel menekankan bahwa efek borrowed size atau agglomeration shadow
  bergantung pada jenis fungsi yang diukur, skala geografis,
  konektivitas transportasi, sejarah, perbedaan budaya, dan daya tarik
  wisata. Malý menyebut studi terdahulu tentang Moravia Selatan, Eropa
  Barat Laut, Swedia, dan sistem perkotaan Amerika sebagai konteks
  empiris perdebatan.
\end{itemize}

Rujukan-rujukan tersebut juga menunjukkan bahwa hasil penelitian
sebelumnya tidak seragam: beberapa studi menemukan interaksi positif
atau manfaat kedekatan, sedangkan studi lain menemukan fungsi berorde
tinggi yang lebih rendah di bawah bayangan kota besar. (Hlm. 334-337,
bagian 1-2.)

\textbf{Interpretasi penulis.} Literatur dipakai untuk menolak pemisahan
biner antara hierarki dan jaringan, serta antara borrowed size dan
agglomeration shadow. Penulis memandang fungsi pelayanan sebagai arena
yang memungkinkan manfaat jaringan dan kompetisi metropolitan terjadi
secara bersamaan, dengan hasil yang berbeda menurut tempat dan jenis
layanan. (Hlm. 336-337, bagian 2; hlm. 345-346, bagian 5.)

\textbf{Inferensi pengolahan.} File menyediakan daftar referensi dan
narasi tinjauan, tetapi tidak menjelaskan strategi pencarian, kriteria
inklusi, rentang korpus, atau penilaian mutu literatur. Karena itu,
survei ini tidak dapat diperlakukan sebagai tinjauan sistematis; cakupan
dan keseimbangan literatur tidak dapat diverifikasi dari TXT saja.

\subsection{Method Used}\label{method-used-37}

\textbf{Laporan sumber.} Penelitian merupakan studi kasus kuantitatif
pada South Moravian Region, wilayah yang sistem perkotaannya relatif
monosentris dan sangat dipengaruhi Brno. Metode terdiri atas empat
langkah:

\begin{itemize}
\tightlist
\item
  mendefinisikan layanan dan nilai masing-masing layanan;
\item
  menentukan small centres dan menghitung centrality fungsi pelayanan;
\item
  menentukan variabel yang berpotensi menjelaskan centrality; dan
\item
  mengevaluasi hubungan antara variabel spasial dan skor centrality
  dengan regresi.
\end{itemize}

Penulis menyusun nilai layanan berdasarkan jumlah municipality yang
memiliki fasilitas tersebut: lebih dari 500 municipality diberi nilai 1;
250-499 diberi 5; 100-249 diberi 10; 50-99 diberi 15; 10-49 diberi 20;
dan kurang dari 10 diberi 25. Indeks fungsi pelayanan setiap
municipality dihitung sebagai jumlah hasil kali jumlah fasilitas setiap
jenis layanan dan nilai layanan, ditulis sebagai
\texttt{I\ =\ sum(F\_i\ V\_i)}. Contoh dalam artikel: dua pub, satu toko
bahan pangan, dan dua restoran menghasilkan
\texttt{(2\ x\ 1)\ +\ (1\ x\ 1)\ +\ (2\ x\ 5)\ =\ 13}. (Hlm. 337-339,
Tabel 1 dan bagian 3.)

Small centres dibatasi pada municipality dengan populasi maksimum 15.000
jiwa dan indeks fungsi pelayanan minimum 80. Brno serta lima permukiman
berpenduduk 20.000-35.000 diperlakukan sebagai large centres. Tidak ada
batas minimum penduduk; ambang indeks dipakai agar municipality yang
hampir tidak memiliki fungsi pelayanan tidak masuk analisis. Karena
sebagian selected centres tidak memiliki status administratif ``town'',
artikel menggunakan istilah ``small centres''. Nilai absolut kemudian
direlatifkan terhadap jumlah penduduk untuk membandingkan pusat dengan
ukuran berbeda. (Hlm. 339, bagian 3.)

Variabel penjelas dikelompokkan ke dalam inner potential, position
within the urban system, dan access to transport infrastructure.
Variabelnya mencakup status town, status pusat pada masa perencanaan
sosialis, fasilitas akomodasi per 1.000 penduduk, proporsi pekerjaan di
pertanian, kepadatan penduduk hinterland yang dapat dijangkau dalam 20
menit, waktu berkendara ke large centre terdekat, proporsi komuter kerja
dari large centres, crude net migration rate 1990-2015, waktu berkendara
ke pintu masuk motorway terdekat, variasi koneksi angkutan umum ke tiga
large centres yang paling terhubung, dan jumlah koneksi tersebut. (Hlm.
339-340, Tabel 2.)

Regresi yang digunakan ialah regresi logistik multinomial. Nilai relatif
indeks diklasifikasikan menjadi tiga kategori fungsi pelayanan, yaitu
high, medium, dan low centrality, sebagai variabel terikat. Variabel
spasial menjadi prediktor. Model diestimasi untuk seluruh small centres
di region dan secara terpisah untuk BMA. Penulis menyatakan bahwa
variabel diperiksa dan dikoreksi bila diperlukan agar memenuhi asumsi
regresi multinomial. Perangkat lunak dan rincian prosedur koreksi
tersebut \textbf{tidak disebutkan dalam file}. (Hlm. 340, bagian 3; hlm.
343-344, Tabel 4.)

\textbf{Interpretasi penulis.} Rancangan ini dipakai untuk mengetahui
variabel yang meningkatkan peluang suatu small centre memiliki fungsi
pelayanan lebih tinggi dibanding centre lain yang berukuran serupa.
Pemisahan model region dan BMA dimaksudkan untuk menangkap dinamika
khusus hinterland terdekat Brno. (Hlm. 340, bagian 3.)

\textbf{Inferensi pengolahan.} Berdasarkan periode observasi layanan
2012-2014 dan penggunaan satu set unit wilayah, desainnya bersifat
potong lintang untuk outcome pelayanan. Model menguji asosiasi
bersyarat, bukan eksperimen, kuasi-eksperimen, panel, atau identifikasi
kausal. Ambang indeks 80 juga membuat pemilihan unit analisis bergantung
pada indikator yang kemudian dijelaskan oleh regresi.

\subsection{Dataset}\label{dataset-37}

\textbf{Laporan sumber.} Dataset pelayanan mencakup 30 layanan dasar dan
berorde lebih tinggi pada 672 municipality di South Moravian Region
untuk periode 2012-2014. Jenis layanan dikelompokkan menurut nilai
layanan sebagai berikut:

\begin{itemize}
\tightlist
\item
  \textbf{Nilai 1:} public library, grocery shop, pub, CzechPoint, dan
  cultural centre.
\item
  \textbf{Nilai 5:} early childhood education, primary education,
  restaurant, post office, dan bank.
\item
  \textbf{Nilai 10:} general practitioner, paediatrician, low-secondary
  education, dentist, dan drugstore.
\item
  \textbf{Nilai 15:} museum and gallery, nursing home, cinema hall, ATM,
  dan insurance company.
\item
  \textbf{Nilai 20:} upper secondary education, retirement home, grammar
  school, emergency medical service, hospital, dan facility for
  physically disabled people.
\item
  \textbf{Nilai 25:} shelter, low-threshold facility for children and
  youth, shopping centre, dan theatre.
\end{itemize}

Nilai tersebut diturunkan dari jumlah municipality yang memiliki
fasilitas, dengan rentang dari lebih dari 500 municipality untuk nilai 1
sampai kurang dari 10 municipality untuk nilai 25. Sumber data pelayanan
yang dicantumkan pada Tabel 1 ialah The South Moravian Region (2012),
Czech Statistical Office (2013), Banky (2015), Czech Post (2015), dan
CzechPoint (2015). (Hlm. 338, Tabel 1.)

Data untuk variabel penjelas dihimpun dari TERPLAN (1985), MFČR (2013),
IDOS (2015), dan Czech Statistical Office (2011, 2016a, 2016b). Data
tersebut dipakai untuk merepresentasikan status historis, daya tarik
wisata, karakter perdesaan, posisi dalam sistem perkotaan, migrasi,
akses motorway, dan koneksi angkutan umum. Peta sistem perkotaan juga
menyebut Czech Statistical Office (2016b) serta delimitasi BMA oleh
Mulíček et al.~(2013). (Hlm. 337, keterangan Gambar 1; hlm. 339-340,
Tabel 2.)

\textbf{Inferensi pengolahan.} Dataset yang tersedia dalam TXT adalah
uraian sumber dan agregat pada tabel, bukan data mentah per
municipality. Hitungan fasilitas mengukur keberadaan fasilitas yang
tercatat, bukan kapasitas, kualitas, volume penggunaan, atau asal
pengguna. File tidak menjelaskan prosedur pembersihan, deduplikasi,
kelengkapan pencatatan, format data mentah, atau kode analisis. Tahun
data juga tidak sepenuhnya seragam: layanan dirangkum untuk 2012-2014,
sementara beberapa prediktor merujuk 1985, 1990-2015, 2011, 2013, 2015,
dan 2016.

\subsection{Population / Sample}\label{population-sample-37}

\textbf{Laporan sumber.} Populasi wilayah studi mencakup 672
municipality di South Moravian Region. Brno, dengan hampir 400.000
penduduk, dan lima settlement berpenduduk 20.000-35.000 diperlakukan
sebagai large centres. Small centres dibatasi pada municipality
berpenduduk paling banyak 15.000 jiwa dan mempunyai indeks fungsi
pelayanan minimal 80. Batas minimum penduduk tidak ditetapkan karena
penulis ingin mempertimbangkan municipality berpenduduk sekitar 3.000
jiwa atau kurang. (Hlm. 337-339, bagian 3.)

Dengan aturan tersebut, terdapat 113 small centres. Penulis menyebut
kelompok ini hanya sekitar seperenam dari seluruh municipality,
sedangkan municipality non-central berjumlah 559 pada Tabel 3. Sebagian
selected centres tidak memiliki status administratif ``town''; karena
itu istilah small centres digunakan untuk mencakup pusat yang melayani
hinterland tetapi bukan town secara formal. (Hlm. 341-342, Tabel 3; hlm.
345, bagian 4.)

Analisis regresi dilakukan untuk seluruh small centres dan secara
terpisah untuk BMA. Tabel 4 menampilkan \texttt{N\ (cases)\ =\ 113} pada
kolom all small centres. Sel pada kolom BMA tampil sebagai \texttt{1668}
dalam hasil ekstraksi TXT; catatan kaki 8 menyatakan bahwa regresi BMA
memasukkan semua municipality untuk memastikan ukuran sampel memadai.
Penanda catatan kaki tampak melekat pada angka tersebut, sehingga nilai
BMA tidak dinormalisasi lebih lanjut dalam review ini. (Hlm. 343-344,
Tabel 4 dan catatan kaki 8.)

\textbf{Interpretasi penulis.} Pemilihan small centres dimaksudkan untuk
menangkap pusat lokal atau mikroregional yang benar-benar memiliki
fungsi pelayanan, bukan seluruh desa yang hampir tidak menyediakan
layanan. Perbandingan berdasarkan nilai indeks relatif dipakai untuk
menunjukkan bahwa pusat dengan ukuran penduduk serupa dapat memiliki
centrality yang berbeda. (Hlm. 339 dan 341-342, bagian 3-4.)

\textbf{Inferensi pengolahan.} Ini bukan sampel acak. Inventarisasi awal
menyerupai sensus wilayah, tetapi analisis small centres menerapkan
seleksi berdasarkan indeks outcome. Karena itu, hasil regresi berlaku
pada definisi analitis small centres yang dipakai artikel, bukan
otomatis pada seluruh municipality atau seluruh small towns di luar
South Moravian Region. Ukuran BMA juga memiliki ambiguitas ekstraksi
yang membatasi replikasi langsung dari TXT.

\subsection{Dependent Variables}\label{dependent-variables-37}

\textbf{Laporan sumber.} Variabel terikat ialah indeks service function
centrality municipality. Indeks dihitung sebagai jumlah berbobot
\texttt{I\ =\ sum(F\_i\ V\_i)}, yaitu jumlah fasilitas pada setiap jenis
layanan (\texttt{F\_i}) dikalikan nilai layanan (\texttt{V\_i}) yang
ditetapkan dari kelangkaan layanan di antara 672 municipality. Nilai
indeks juga direlatifkan terhadap jumlah penduduk, dengan satuan indeks
per 100 inhabitants pada Tabel 3. (Hlm. 338-339, Tabel 1 dan bagian 3;
hlm. 342, catatan kaki 6.)

Untuk regresi, nilai relatif indeks dikategorikan menjadi \textbf{high
centrality}, \textbf{medium centrality}, dan \textbf{low centrality}.
Kategori tersebut menjadi outcome multinomial; kategori low centrality
digunakan sebagai reference category pada Tabel 4. TXT tidak menyebutkan
batas numerik atau prosedur klasifikasi tiga kategori itu. (Hlm. 340,
bagian 3; hlm. 343-344, Tabel 4.)

\textbf{Inferensi pengolahan.} Indeks adalah proksi aditif atas jumlah
fasilitas yang tercatat dan diberi bobot normatif berdasarkan
kelangkaan, bukan ukuran langsung atas arus layanan, kapasitas, mutu,
akses aktual, atau surplus fungsi dibanding kebutuhan lokal.
Relativisasi per 100 penduduk membantu perbandingan antarukuran
municipality, tetapi tidak dengan sendirinya menghilangkan perbedaan
luas wilayah, struktur permintaan, kapasitas fasilitas, atau kualitas
pencatatan.

\subsection{Results}\label{results-37}

\textbf{Laporan sumber: distribusi layanan dan variasi centrality.}
Distribusi 30 layanan mengikuti pola hierarkis yang konsisten dengan
threshold dan range. Layanan umum seperti public library, grocery shop,
dan pub tersedia di banyak municipality; layanan dengan sedikit lokasi
seperti theatre, shopping centre, hospital, dan layanan sosial atau
pendidikan berorde tinggi lebih terkonsentrasi pada pusat yang lebih
besar. (Hlm. 338, Tabel 1 dan bagian 3.)

Tabel 3 menunjukkan variasi besar di antara small centres dengan ukuran
penduduk serupa. Pada kategori 9.000-11.999 penduduk, rerata relative
centrality adalah 7,9 dengan maksimum 10,3 dan koefisien variasi 27,0\%.
Pada kategori 600-999 penduduk, rerata relative centrality adalah 14,2
dengan maksimum 30,6 dan koefisien variasi 38,8\%. Koefisien variasi
kategori small centres lain berada pada 31,3\%-44,6\%. (Hlm. 341-342,
Tabel 3.)

\textbf{Laporan sumber: regresi seluruh small centres.} Model pada Tabel
4 menggunakan low centrality sebagai referensi. Nilai kecocokan model
yang tercantum ialah \texttt{-2LL\ =\ 175.396}, Nagelkerke
\texttt{R\ Square\ =\ 0.529}, dan
\texttt{Percentage\ Correctly\ Estimate\ =\ 64.6\%}, dengan
\texttt{N\ (cases)\ =\ 113}. Variabel yang signifikan pada perbandingan
berikut dilaporkan dengan \texttt{Exp(B)} dan \texttt{Sig.}:

\begin{itemize}
\tightlist
\item
  \textbf{Medium dibanding low centrality:} tourist attractiveness
  (\texttt{Exp(B)\ =\ 2.202}, \texttt{p\ =\ 0.007}), peripherality yang
  diukur melalui kepadatan hinterland (\texttt{0.579},
  \texttt{p\ =\ 0.038}), dan jarak berkendara ke large centre terdekat
  (\texttt{1.190}, \texttt{p\ =\ 0.011}).
\item
  \textbf{High dibanding low centrality:} tourist attractiveness
  (\texttt{3.075}, \texttt{p\ =\ 0.000}), kepadatan hinterland
  (\texttt{0.522}, \texttt{p\ =\ 0.033}), jarak ke large centre terdekat
  (\texttt{1.288}, \texttt{p\ =\ 0.002}), importance for higher-ranked
  centres (\texttt{3.936}, \texttt{p\ =\ 0.001}), dan access to the
  nearest motorway (\texttt{1.077}, \texttt{p\ =\ 0.027}).
\end{itemize}

Status town, status sebagai centre pada masa perencanaan sosialis, rural
character, suburban character, functional polycentricity, dan network
density tidak dilaporkan signifikan dalam model seluruh small centres.
(Hlm. 343, Tabel 4; hlm. 343-344, bagian 4.)

\textbf{Laporan sumber: regresi BMA.} Tabel 4 mencantumkan
\texttt{-2LL\ =\ 247.803}, Nagelkerke \texttt{R\ Square\ =\ 0.479}, dan
\texttt{Percentage\ Correctly\ Estimate\ =\ 64.5\%} untuk kolom BMA.
Variabel historical importance dan rural character dikeluarkan dari
model BMA karena dinyatakan tidak relevan. Variabel signifikan yang
tersisa ialah:

\begin{itemize}
\tightlist
\item
  \textbf{Medium dibanding low centrality:} distance ke large centre
  (\texttt{Exp(B)\ =\ 1.270}, \texttt{p\ =\ 0.024}), importance for
  higher-ranked centres (\texttt{2.190}, \texttt{p\ =\ 0.000}), suburban
  character atau net migration (\texttt{0.678}, \texttt{p\ =\ 0.016}),
  dan network density (\texttt{1.470}, \texttt{p\ =\ 0.038}).
\item
  \textbf{High dibanding low centrality:} tourist attractiveness
  (\texttt{2.159}, \texttt{p\ =\ 0.001}), distance ke large centre
  (\texttt{1.489}, \texttt{p\ =\ 0.020}), importance for higher-ranked
  centres (\texttt{2.279}, \texttt{p\ =\ 0.008}), dan network density
  (\texttt{2.537}, \texttt{p\ =\ 0.001}).
\end{itemize}

\textbf{Laporan sumber: pola spasial dan contoh tempat.} Small centres
di bagian barat daya dan utara region, serta di dekat batas selatan BMA,
mempunyai centrality tinggi ketika berperan sebagai pusat lokal yang
lebih jauh dari large centres. Tempat yang sangat menarik bagi
wisatawan, seperti Mikulov, Valtice, Lednice, dan Vranov nad Dyjí,
memiliki centrality tinggi meskipun posisi dalam jaringan tidak selalu
dominan. Modřice dan Rajhrad di selatan Brno serta Chvaletice di selatan
Znojmo disebut sebagai tempat yang memperoleh centrality lebih tinggi
melalui lokasi strategis dan suburbanisasi komersial. (Hlm. 343, bagian
4; hlm. 344-345, Gambar 4 dan bagian 5.)

Sebaliknya, fungsi pelayanan small centres di sekitar Brno, Hodonín, dan
Vyškov disebut lebih rendah karena kedekatan dengan higher-ranked
centres. Di BMA, keterhubungan angkutan umum dengan large centres,
terutama Brno, dan kemampuan menarik komuter dari core city berkaitan
dengan centrality yang lebih tinggi. Namun, beberapa municipality di
hinterland Brno mengalami pertumbuhan penduduk karena suburbanisasi
hunian tanpa pertambahan fasilitas yang sebanding; Kuřim, Šlapanice, dan
Mokrá-Horákov diberikan sebagai contoh. (Hlm. 343-344, bagian 4.)

\textbf{Interpretasi penulis.} Pada skala seluruh region, centrality
yang lebih tinggi dikaitkan dengan daya tarik wisata, posisi yang lebih
perifer, kepadatan hinterland yang lebih rendah, dan hubungan kerja
dengan large centres. Pada BMA, network density dan arus komuter dari
pusat lebih menonjol. Penulis membaca pola tersebut sebagai borrowed
size bagi pusat yang terhubung dan mampu menarik kegiatan, serta
agglomeration shadow bagi pusat dekat aglomerasi yang kehilangan fungsi
pelayanan. (Hlm. 343-346, bagian 4-5.)

\textbf{Inferensi pengolahan.} Angka \texttt{Exp(B)} dan \texttt{Sig.}
mendukung asosiasi kategori centrality dengan prediktor dalam model,
tetapi tidak cukup untuk menentukan bahwa Brno atau large centres secara
kausal menimbulkan borrowed size atau agglomeration shadow. Koefisien
jarak yang positif berarti peluang kategori centrality lebih tinggi
meningkat ketika jarak ke large centre bertambah, sesuai pembacaan
shadow oleh penulis; arah substantif variabel waktu atau koneksi
transportasi tetap harus mengikuti definisi pengukuran yang diberikan
artikel.

\subsection{Findings}\label{findings-37}

\textbf{Temuan yang dilaporkan sumber.} Ukuran penduduk hanya menjadi
indikator umum fungsi urban. Di dalam kategori ukuran yang sama, small
centres mempunyai tingkat centrality dan ragam layanan yang berbeda
besar. Fungsi pelayanan terutama berasosiasi dengan daya tarik wisata,
kepentingan dalam jaringan large centres, posisi terhadap pusat yang
lebih besar, dan kepadatan hinterland. (Hlm. 341-345, bagian 4; hlm.
345-346, bagian 5.)

\textbf{Interpretasi penulis.} Penulis menyatakan bahwa kedua proses
tidak perlu dipisahkan secara dikotomis:

\begin{itemize}
\tightlist
\item
  small centres yang lebih perifer dan tidak berhadapan dengan centre
  yang lebih kuat dapat mempertahankan fungsi pelayanan sebagai pusat
  lokal;
\item
  sebagian pusat di suburban BMA memperoleh commercial activities dan
  komuter dari large centres sehingga dapat ``borrow functions'';
\item
  sebagian pusat dekat Brno mengalami kekurangan fungsi pelayanan karena
  competitive advantage core city, tetapi tetap dapat ``borrow
  performance'' melalui pertumbuhan penduduk; dan
\item
  kekurangan fungsi pelayanan di bawah agglomeration shadow dapat
  dikompensasi oleh fungsi urban Brno serta perjalanan penduduk ke core
  city.
\end{itemize}

Penulis menyebut hubungan tersebut sebagai interdependence dan mutual
conditionality antara borrowed size dan agglomeration shadow, bukan dua
kondisi yang saling meniadakan. Di kawasan perdesaan dengan banyak small
centres berukuran serupa, hubungan antarpusat dipandang lebih mungkin
berbentuk kompetisi yang dipengaruhi kekhususan lokal. (Hlm. 345-346,
bagian 5.)

\textbf{Inferensi pengolahan.} Temuan empiris paling kuat ialah
heterogenitas fungsi pelayanan dan pola asosiasi spasial yang berbeda
antara region secara keseluruhan dan BMA. Istilah borrowed size lebih
tepat dibaca sebagai interpretasi konseptual atas konektivitas, komuter,
dan aktivitas komersial; TXT tidak menyediakan pengukuran terpisah yang
membuktikan fungsi metropolitan benar-benar dipakai atau dipindahkan
kepada small centre.

\subsection{Conclusions}\label{conclusions-37}

\textbf{Kesimpulan yang dinyatakan penulis.} Penulis merangkum
kesimpulan utama sebagai berikut:

\begin{itemize}
\tightlist
\item
  South Moravian Region memiliki pola yang kuat dipengaruhi Brno sebagai
  pusat regional yang relatif monosentris, sementara small centres tetap
  penting sebagai pusat bagi settlement kecil dan hinterland perdesaan.
\item
  Distribusi layanan dasar dan layanan berorde tinggi mengikuti
  threshold, range, dan hierarki sistem perkotaan, tetapi small centres
  dengan ukuran penduduk serupa berbeda jauh dalam fungsi pelayanannya.
\item
  Fungsi pelayanan terutama berhubungan dengan importance for
  higher-ranked centres yang diukur melalui arus komuter kerja, tourist
  attractiveness yang diukur melalui fasilitas akomodasi, posisi
  terhadap large centres, dan kepadatan penduduk hinterland.
\item
  Centre yang lebih perifer dan tidak berhadapan dengan pusat kuat
  cenderung memiliki fungsi pelayanan lebih tinggi; centre di wilayah
  urban dan dekat aglomerasi lebih mungkin mengalami kekurangan fungsi
  pelayanan yang dibaca sebagai agglomeration shadow.
\item
  Di BMA, centre yang berada pada sumbu perkembangan zona komersial atau
  terhubung dengan large centres dapat memperlihatkan borrowing size,
  sedangkan centre yang dekat Brno tidak selalu meminjam fungsi
  pelayanan dan dapat tetap underserviced meskipun mengalami pertumbuhan
  penduduk.
\end{itemize}

Penulis menutup dengan pandangan bahwa borrowed size dan agglomeration
shadow saling bergantung dan dapat hadir pada aspek yang berbeda dari
tempat yang sama. (Hlm. 345-346, bagian 5.)

\textbf{Inferensi pengolahan.} Kesimpulan tentang variasi indeks dan
asosiasi prediktor konsisten dengan hasil yang dilaporkan. Namun,
kesimpulan bahwa terjadi borrowed size sebaiknya dipahami sebagai
interpretasi atas keterhubungan, arus komuter, aktivitas komersial, dan
perbedaan fungsi antarpusat; file tidak menunjukkan pengukuran langsung
atas fungsi metropolitan yang dipakai oleh penduduk atau perusahaan
small centre.

\subsection{Contributions}\label{contributions-37}

\textbf{Laporan sumber.} Artikel tidak menyediakan subbagian yang secara
khusus berjudul ``Contributions''. Kontribusi yang dapat diidentifikasi
dari tujuan, metode, hasil, dan kesimpulannya ialah:

\begin{itemize}
\tightlist
\item
  membawa perdebatan borrowed size dan agglomeration shadow ke fungsi
  pelayanan small centres pada skala regional;
\item
  mengoperasionalkan fungsi pelayanan melalui 30 jenis layanan dan
  indeks berbobot pada tingkat municipality;
\item
  membandingkan centre dengan ukuran penduduk berbeda melalui nilai
  relatif per 100 inhabitants;
\item
  menghubungkan centrality dengan karakter internal, posisi dalam sistem
  perkotaan, konektivitas transportasi, dan relasi dengan large centres
  melalui regresi logistik multinomial; serta
\item
  menunjukkan secara empiris bahwa manfaat jaringan dan kekurangan
  fungsi akibat kedekatan dengan aglomerasi dapat muncul bersamaan,
  khususnya dengan membedakan ``borrow functions'' dari ``borrow
  performance''.
\end{itemize}

\textbf{Inferensi pengolahan.} Sumbangan paling jelas dari artikel ialah
pemindahan unit perhatian dari kota besar dan fungsi tingkat tinggi ke
small centres sebagai simpul pelayanan hinterland, serta pembacaan yang
tidak biner atas borrowed size dan agglomeration shadow. Klaim kebaruan
absolut, kemampuan generalisasi, dan keunggulan indeks dibanding ukuran
lain \textbf{tidak disebutkan dalam file}.

\subsection{Research Gap}\label{research-gap-37}

\textbf{Kesenjangan yang dinyatakan sumber.} Artikel menempatkan
penelitiannya pada beberapa kekosongan berikut:

\begin{itemize}
\tightlist
\item
  Peran small towns kurang mendapat perhatian dalam teori dan penelitian
  sistem perkotaan, terutama dibandingkan aglomerasi besar dan wilayah
  padat.
\item
  Konsep borrowed size dan agglomeration shadow masih memberi sedikit
  perhatian pada fungsi pelayanan small towns.
\item
  Banyak penelitian terdahulu berada pada skala internasional,
  menggunakan unit spasial yang lebih besar, atau menilai fungsi berorde
  tinggi; dimensi regional dan spektrum layanan yang lebih luas belum
  cukup dieksplorasi.
\item
  Mekanisme yang mendasari kerja sama dan kompetisi antarpusat masih
  perlu diteliti, terutama pada sistem perkotaan dengan kondisi yang
  mendukung susunan polisentris. (Hlm. 334-337, bagian 1-2; hlm. 346,
  bagian 5.)
\end{itemize}

\textbf{Inferensi pengolahan: kesenjangan yang tersisa.} Dari isi TXT,
artikel belum memisahkan secara langsung kontribusi borrowed size dari
daya tarik internal, status administratif, suburbanisasi, dan kompetisi
pusat besar. Artikel juga tidak menyediakan outcome pelayanan
longitudinal, ukuran kapasitas atau mutu, ukuran penggunaan layanan
lintas batas, atau pembanding wilayah lain. Ini adalah kesenjangan
analitis yang diturunkan dari desain file, bukan agenda yang dinyatakan
secara eksplisit oleh penulis.

\subsection{Limitations}\label{limitations-37}

\textbf{Keterbatasan yang dinyatakan sumber.} File tidak memuat
subbagian atau daftar keterbatasan penelitian yang berdiri sendiri.
Penulis menyatakan bahwa definisi small towns merupakan persoalan
penting, tidak menetapkan batas minimum penduduk, dan memilih indeks
pelayanan minimum untuk menghindari masuknya municipality yang hampir
tidak memiliki fungsi pelayanan. Penulis juga menyatakan bahwa mekanisme
di balik bentuk kerja sama dan kompetisi perlu diteliti lebih lanjut.
(Hlm. 339, bagian 3; hlm. 346, bagian 5.) Rincian keterbatasan lain
dalam format eksplisit \textbf{tidak disebutkan dalam file}.

\textbf{Inferensi pengolahan.} Batasan yang terlihat dari isi TXT
meliputi:

\begin{itemize}
\tightlist
\item
  \textbf{Cakupan kasus:} satu region di Republik Ceko dengan sistem
  yang dipengaruhi Brno. Validitas eksternal untuk sistem polisentris
  lain, wilayah yang lebih perifer, atau negara lain tidak diuji.
\item
  \textbf{Definisi unit:} small centres dipilih dengan ambang indeks
  pelayanan dan mencakup municipality yang sebagian tidak berstatus
  town. Hal ini membuat populasi analitis heterogen dan berpotensi
  mengaitkan prediktor dengan outcome yang ikut menentukan seleksi.
\item
  \textbf{Cakupan hitungan:} narasi menyebut 672 municipality, sedangkan
  Tabel 3 menampilkan Brno dan lima large centres di samping 113 small
  centres serta 559 non-central. TXT tidak menjelaskan denominator atau
  kemungkinan tumpang tindih baris-baris tersebut; review mempertahankan
  angka sesuai locator masing-masing tanpa memaksakan penjumlahan.
\item
  \textbf{Validitas konstruk:} jumlah fasilitas berbobot tidak sama
  dengan kapasitas, kualitas, tingkat pemakaian, keterjangkauan, atau
  kecukupan layanan. Arus komuter yang diukur adalah arus kerja dari
  large centres, bukan arus pengguna layanan.
\item
  \textbf{Pembobotan:} nilai layanan ditentukan dari kelangkaan
  fasilitas; matriks pembandingan, penguji, uji sensitivitas, atau
  justifikasi alternatif bobot tidak tersedia dalam TXT.
\item
  \textbf{Klasifikasi:} batas numerik untuk kategori high, medium, dan
  low centrality tidak disebutkan. Dampak perubahan ambang indeks 80,
  relativisasi per penduduk, atau aturan klasifikasi tidak diuji.
\item
  \textbf{Waktu:} outcome dirangkum untuk 2012-2014 dan prediktor
  berasal dari beberapa tahun atau rentang waktu. File tidak melaporkan
  pengamatan panel atau perubahan fungsi pelayanan antarwaktu.
\item
  \textbf{Model:} regresi logistik multinomial menunjukkan asosiasi,
  bukan sebab-akibat. TXT tidak melaporkan secara lengkap data mentah,
  pemeriksaan residual, dependensi spasial, bentuk koreksi asumsi, atau
  analisis sensitivitas.
\item
  \textbf{Replikasi:} jumlah kasus BMA dirender sebagai \texttt{1668}
  dengan penanda catatan kaki yang melekat, dan file tidak menyediakan
  data atau kode untuk menyelesaikan ambiguitas tersebut.
\end{itemize}

\subsection{Challenges}\label{challenges-37}

\textbf{Tantangan yang dilaporkan sumber.} Tantangan utama penelitian
ialah menetapkan definisi operasional small towns dalam hierarki
regional ketika tidak ada batas minimum penduduk yang digunakan. Penulis
juga harus membedakan status town formal dari peran sebuah municipality
sebagai centre, menyusun ukuran fungsi dari banyak jenis layanan, dan
memasukkan faktor sejarah, wisata, suburbanisasi, posisi jaringan, serta
transportasi ke dalam satu kerangka analisis. (Hlm. 337-340, bagian 3.)

Tantangan substantif lain ialah menjelaskan mengapa centre yang dekat
dengan metropolitan dapat sekaligus terhubung dengan jaringan besar dan
kehilangan sebagian fungsi pelayanan. Penulis menangani persoalan ini
dengan memisahkan analisis seluruh region dan BMA serta membandingkan
indikator borrowed size dengan agglomeration shadow. (Hlm. 340 dan
343-346, bagian 3-5.)

\textbf{Inferensi pengolahan.} Tantangan analitis terbesar adalah
membedakan efek ukuran penduduk, status administratif, daya tarik
internal, manfaat konektivitas, dan kompetisi metropolitan ketika semua
dapat menghasilkan skor fasilitas yang tinggi atau rendah. Tantangan
replikasi juga besar karena rincian data per municipality, aturan
klasifikasi, pembobotan layanan, dan beberapa detail sampel BMA tidak
tersedia secara lengkap dalam TXT.

\subsection{Practical Implications}\label{practical-implications-37}

\textbf{Implikasi yang dinyatakan atau ditunjukkan oleh pembahasan
penulis.} Small centres dipandang penting untuk menjaga penyediaan
pekerjaan dan layanan bagi hinterland perdesaan serta mengurangi
kesenjangan spasial. Hasil artikel menyarankan perlunya memperhatikan
centre yang berada dekat aglomerasi tetapi memiliki fungsi pelayanan
rendah, centre perifer yang menjadi pusat lokal, serta centre pada sumbu
komersial yang memperoleh manfaat dari hubungan metropolitan. Penilaian
hubungan core-hinterland juga perlu memperhatikan mobilitas penduduk dan
kemampuan mengakses transportasi, telekomunikasi, serta broadband. (Hlm.
335-337, bagian 2; hlm. 345-346, bagian 5.)

\textbf{Inferensi pengolahan.} Indeks dan hasil regresi dapat dipakai
sebagai diagnosis awal untuk membedakan wilayah yang kekurangan
fasilitas dari wilayah yang mungkin mengandalkan fungsi core city. Skor
tersebut tidak cukup sebagai dasar tunggal penentuan investasi;
verifikasi kapasitas, kebutuhan penduduk, penggunaan aktual, akses
perjalanan, dan fungsi yang tersedia di municipality tetangga tetap
diperlukan. Artikel tidak menguji instrumen kebijakan atau dampak
intervensi tertentu.

\subsection{Applications}\label{applications-37}

\textbf{Aplikasi yang didukung langsung oleh penelitian.} Prosedur dalam
artikel dapat digunakan untuk:

\begin{itemize}
\tightlist
\item
  memetakan distribusi layanan dasar dan berorde tinggi pada tingkat
  municipality;
\item
  menyusun indeks perbandingan fungsi pelayanan small centres;
\item
  mengenali perbedaan fungsi di antara centre dengan ukuran penduduk
  serupa;
\item
  menguji hubungan fungsi pelayanan dengan posisi terhadap large
  centres, hinterland, transportasi, sejarah, wisata, dan suburbanisasi;
  serta
\item
  membedakan secara eksploratif centre yang tampak berfungsi sebagai
  local centre dari centre yang berada dekat aglomerasi tetapi memiliki
  layanan lebih rendah.
\end{itemize}

\textbf{Batas aplikasi.} Implementasi kebijakan, penggunaan oleh lembaga
perencana, perangkat lunak atau data yang dibagikan, dan evaluasi
pascapenerapan \textbf{tidak disebutkan dalam file}. Transfer prosedur
ke wilayah lain merupakan inferensi pengolahan dan memerlukan
penyesuaian klasifikasi layanan, bobot, unit wilayah, serta validasi
lokal.

\subsection{Future Research
(Optional)}\label{future-research-optional-37}

\textbf{Agenda yang disebut sumber.} Penulis menyatakan bahwa mekanisme
yang mendasari bentuk kerja sama dan kompetisi antarpusat perlu
diselidiki lebih lanjut, terutama dengan melibatkan sistem perkotaan
yang memiliki kondisi untuk susunan polisentris. Penulis juga menekankan
bahwa evaluasi kelebihan dan kekurangan konfigurasi metropolitan perlu
memperhatikan tingkat mobilitas individual serta akses dan kemampuan
menggunakan transportasi, telekomunikasi, dan broadband. (Hlm. 346,
bagian 5.)

Agenda tambahan yang tidak dinyatakan sumber: \textbf{Tidak disebutkan
dalam file}.

\subsection{Provenance Notes}\label{provenance-notes-37}

\begin{itemize}
\tightlist
\item
  \textbf{Verifikasi filename dan pasangan sumber:} Pemeriksaan rekursif
  pada \texttt{source/} menemukan satu TXT dengan prefiks A22, yaitu
  \texttt{source/journal/A22-maly-2016-small-towns-in-the-context-of-borrowed-size-and-agglomeration-shadow-debates-the-case-of-the-south-moravian-region-czech-republic-10.1515-euco-2016-0024.txt}.
  TXT menunjuk pada PDF relatif
  \texttt{journal/A22-maly-2016-small-towns-in-the-context-of-borrowed-size-and-agglomeration-shadow-debates-the-case-of-the-south-moravian-region-czech-republic-10.1515-euco-2016-0024.pdf};
  DOI, judul, penulis, tahun, dan rentang halaman pada blok teksnya
  cocok dengan pasangan tersebut.
\item
  \textbf{Basis akses:} Review ini hanya menggunakan TXT A22 yang
  disebut di atas. Tidak ada sumber luar, metadata daring, PDF pasangan,
  atau referensi yang dikutip artikel yang dibuka untuk menambah
  substansi. Daftar referensi dalam artikel hanya dilaporkan sebagai
  bagian dari isi A22.
\item
  \textbf{Cakupan baca:} Seluruh 896 baris TXT dibaca, termasuk blok
  \texttt{{[}PDF\ Metadata{]}} pada baris 1-23, blok
  \texttt{{[}Extracted\ Text{]}}, artikel dari hlm. 333-350, abstrak
  Inggris dan Czech, bagian 1-5, Tabel 1-4, Gambar 1-4 beserta
  keterangan tekstualnya, acknowledgement, dan daftar referensi sampai
  baris 895.
\item
  \textbf{Metadata ekstraksi:} Blok metadata menyatakan ekstraksi
  dilakukan 31 Agustus 2026 dengan \texttt{pdftotext\ -layout}. PDF
  tercatat memiliki 18 halaman, ukuran A4, tidak terenkripsi, berukuran
  1.333.929 byte, versi 1.5, dan diproduksi dengan iTextSharp 5.1.3.
  CreationDate dan ModDate sama-sama 4 Januari 2017; custom metadata,
  metadata stream, tagging, user properties, form, JavaScript, dan
  optimasi tercatat tidak ada.
\item
  \textbf{Locator:} Nomor halaman dalam review adalah nomor halaman
  jurnal yang tercetak, hlm. 333-350, bukan nomor baris TXT. Bagian,
  persamaan, tabel, gambar, dan catatan kaki disebut ketika tersedia.
\item
  \textbf{Keterbatasan ekstraksi:} Teks berasal dari tata letak PDF dua
  kolom sehingga beberapa baris terputus oleh form-feed dan penanda
  catatan kaki dapat menempel pada angka. Contoh penting ialah sel
  \texttt{N\ (cases)} BMA yang tampil sebagai \texttt{1668}; review
  mempertahankan bentuk yang terbaca dan menandai ambiguitasnya. Detail
  visual gambar tidak dinilai melampaui keterangan dan pembahasan
  tekstual yang tersedia dalam TXT.
\item
  \textbf{Ketidakselarasan jumlah:} Narasi metode menyebut 672
  municipality, sementara Tabel 3 melaporkan 113 small centres dan 559
  non-central sekaligus menampilkan Brno serta lima large centres
  sebagai baris tersendiri. TXT tidak menjelaskan apakah baris-baris
  tersebut memakai denominator yang sama atau sebagian tumpang tindih;
  review tidak mengoreksi angka sumber.
\item
  \textbf{Informasi yang tidak tersedia:} File tidak menyebutkan
  strategi tinjauan pustaka, data mentah, kode, protokol pembersihan,
  batas kategori centrality, rincian pembobotan dan validasi, ukuran BMA
  yang tidak ambigu dalam ekstraksi, atau implementasi kebijakan. Status
  telaah sejawat \textbf{tidak disebutkan dalam file} dan tidak
  diverifikasi dari luar.
\item
  \textbf{Pemisahan epistemik:} Bagian berlabel \textbf{Laporan sumber}
  atau \textbf{Temuan yang dilaporkan sumber} memparafrasakan isi
  artikel; \textbf{Interpretasi penulis} merekam makna yang diberikan
  Malý terhadap hasil; \textbf{Inferensi pengolahan} adalah evaluasi
  yang diturunkan dari isi dan keterbatasan internal TXT, bukan klaim
  artikel atau informasi eksternal.
\end{itemize}

\section{Review A23}\label{review-a23}

\subsection{Summarized Abstract}\label{summarized-abstract-38}

File TXT tidak memuat bagian yang diberi judul ``Abstract''. Bagian ini
merupakan ringkasan berbasis teks yang tersedia.

\textbf{Temuan dilaporkan:} Meijers dan Cardoso mengeksplorasi dan
menguji secara empiris kerangka \emph{borrowed size/agglomeration
shadows} pada 120 municipality di kawasan Randstad Holland. Mereka
memperluas matriks empat kuadran menjadi matriks sembilan kotak dengan
tujuh kemungkinan hasil lokal, lalu mengaitkan posisi municipality
dengan enam faktor: kepentingan historis, ukuran, komposisi sektoral
ekonomi lokal, status sosial, \emph{network embeddedness}, dan
pariwisata (Introduction; Classifying cities in the Randstad;
Conclusion).

\textbf{Interpretasi penulis:} Integrasi metropolitan dapat menghasilkan
manfaat pada tingkat kawasan, tetapi manfaat dan biayanya tidak tersebar
merata di antara kota-kota. Hasil eksploratif menunjukkan bahwa kota
sekunder di Randstad sangat heterogen, sehingga strategi pengembangan
kota tidak seharusnya seragam. Penulis secara eksplisit tidak bermaksud
menetapkan kausalitas pada tahap ini (Introduction; Conclusion).

\subsection{Problem Statement}\label{problem-statement-38}

\textbf{Masalah yang dilaporkan dalam sumber:} Konsepsi tradisional
tentang kota dengan hinterland yang terutama rural perlu ditantang
karena banyak kawasan fungsional atau metropolitan kini terdiri atas
banyak kota. Banyak teori perkotaan masih berangkat dari model kota
tunggal dan hinterland, sementara praktik kebijakan menghadapi kawasan
multicentric dengan batas dan kewenangan yang tumpang tindih serta
tantangan bersama (Introduction).

\textbf{Interpretasi penulis:} Integrasi antarkota menimbulkan manfaat
sosioekonomi yang generatif sekaligus efek distributif di dalam kawasan.
Belum jelas mengapa satu kota memperoleh manfaat berupa \emph{borrowed
size}, sedangkan kota lain berada dalam \emph{agglomeration shadow}
(Introduction; Theory).

\textbf{Batas masalah:} Bab ini menyatakan bahwa analisisnya bersifat
eksploratif dan tidak dimaksudkan untuk menetapkan hubungan sebab-akibat
atau menghasilkan rekomendasi kebijakan yang kuat pada tahap awal
(Introduction; Conclusion).

\subsection{Objectives}\label{objectives-38}

\textbf{Tujuan yang dinyatakan penulis:} - Mengeksplorasi dan memberi
dukungan empiris awal terhadap kerangka \emph{borrowed
size/agglomeration shadows} dengan menempatkan kota-kota di Randstad
Holland ke dalam kerangka tersebut (Introduction). - Mengklasifikasikan
municipality berdasarkan hubungan antara ukuran dan fungsi perkotaan,
serta antara ukuran dan kinerja atau dinamika perkotaan (Classifying
cities in the Randstad). - Mengidentifikasi faktor-faktor yang
berasosiasi dengan posisi dalam kuadran atau kategori
\emph{agglomeration shadow} dan \emph{borrowed size} (Introduction;
Research approach). - Menggambarkan karakteristik umum yang membedakan
kota primer dan sekunder, termasuk heterogenitas di antara kota
sekunder, untuk menghindari strategi pengembangan yang
\emph{one-size-fits-all} (Introduction). - Menguji dan mengembangkan
lebih lanjut kerangka \emph{borrowed size/agglomeration shadow}
(Introduction).

\subsection{Literature Survey}\label{literature-survey-38}

\textbf{Tinjauan yang disajikan penulis:} Literatur tentang hubungan
antarkota dibingkai melalui konsep \emph{urban network externalities},
yaitu manfaat yang berasal dari relasi antartempat dan jaringan yang
menghubungkan aglomerasi. Sumber menyebut lima arus literatur: skala
regional atau intra-metropolitan dengan konsep polycentricity; skala
global dengan jaringan perusahaan dan komunikasi kota dunia; ekonomi
perkotaan atau \emph{new economic geography}; teori simulasi dan
kompleksitas; serta kajian distribusi ukuran kota (Theory, mengutip
Peris et al., 2018).

Penulis menghubungkan integrasi kota dengan manfaat dan biaya yang
mungkin memiliki efek generatif pada jaringan kota, tetapi juga efek
redistributif di antara kota-kota. Literatur transportasi, \emph{new
economic geography}, dan argumen \emph{two-way road} digunakan untuk
menjelaskan bahwa konektivitas dapat mendorong investasi ke luar maupun
ke dalam pusat. Hasil integrasi bergantung pada keseimbangan gaya
sentripetal dan sentrifugal serta biaya transportasi (Theory, mengutip
Preston, 2001; Krugman, 1991; Fujita et al., 1999).

Kerangka utama dibangun dari konsep \emph{borrowed size} dan
\emph{agglomeration shadows}. \emph{Borrowed size} terjadi ketika suatu
kota memiliki fungsi perkotaan atau tingkat kinerja yang biasanya
diasosiasikan dengan kota yang lebih besar sebagai konsekuensi integrasi
dengan kota-kota tersebut. \emph{Agglomeration shadow} adalah kondisi
ketika kota memiliki fungsi atau kinerja lebih rendah daripada yang
diharapkan berdasarkan ukurannya akibat efek kompetisi dari kota-kota
yang terhubung melalui jaringan dan relasi (Theory, mengutip Alonso,
1973; Krugman, 1993; Meijers and Burger, 2017).

Sumber juga menempatkan network embeddedness dalam perdebatan tentang
apakah jaringan menggantikan kedekatan atau melengkapinya. Studi yang
dirujuk menunjukkan pentingnya konektivitas untuk fungsi metropolitan,
tetapi interaksi antara faktor lokal dan faktor jaringan masih belum
diketahui (Theory, mengutip Johansson and Quigley, 2004; Meijers et al.,
2016). Kajian terdahulu tentang integrasi fungsional dan kelembagaan di
kawasan polycentric digunakan sebagai konteks bagi pengujian pada
tingkat kota individual (Theory, mengutip Maly, 2016; Meijers et al.,
2018).

\subsection{Method Used}\label{method-used-38}

\textbf{Metode yang dilaporkan:} Penelitian menggunakan studi kasus
empiris eksploratif atas Randstad Holland, yang diposisikan sebagai
kawasan metropolitan multicentric dan polycentric. Unit analisisnya
adalah 120 municipality di Randstad (Classifying cities in the Randstad,
terutama pernyataan unit analisis sebelum Figure 2.1).

Kota ditempatkan dalam matriks berdasarkan dua hubungan: ukuran dengan
fungsi perkotaan, dan ukuran dengan kinerja atau dinamika perkotaan.
Fungsi diproksikan dengan jumlah pekerjaan pada 2016 dan keberadaan atau
jumlah restoran pada 2011. Kinerja diproksikan dengan pertumbuhan
penduduk pada 2002-2018 dan perkembangan jumlah restoran pada 1996-2011
(Classifying cities in the Randstad).

Kerangka klasifikasi diperluas dari empat kuadran menjadi matriks
sembilan kotak untuk memasukkan municipality yang memiliki jumlah fungsi
dan dinamika sesuai harapan. Sumber memberi tujuh label hasil lokal:
\emph{agglomeration shadow}, \emph{performance shadow}, \emph{functional
shadow}, \emph{as expected}, \emph{borrowed functions}, \emph{borrowed
performance}, dan \emph{borrowed size} (Classifying cities in the
Randstad; Figure 2.2; Conclusion).

Untuk klasifikasi, penulis menghitung pekerjaan dan restoran per kapita,
kemudian memakai penyimpangan lebih dari 0.75 standar deviasi di atas
atau di bawah mean sebagai batas fungsi lebih banyak atau lebih sedikit
daripada yang diharapkan. Prosedur serupa diterapkan pada pertumbuhan
penduduk dan perkembangan restoran untuk dimensi kinerja. Jika dua
indikator pada dimensi yang sama berlawanan, penulis memakai aturan
kategorisasi berbasis indikator yang menyimpang dan, pada kasus yang
tidak biasa, rata-rata skor terstandar (Classifying cities in the
Randstad; Figure 2.3).

Asosiasi dieksplorasi dengan box plot untuk enam faktor pada Figures
2.4-2.9. Perbandingan kategori menggunakan uji-t pada dimensi fungsional
dan kinerja, lalu model ordinal regression dengan enam covariate
digunakan untuk menjelaskan skor pada masing-masing dimensi (Results:
aggregate patterns; Results: functional dimension; Results: performance
dimension).

\subsection{Dataset}\label{dataset-38}

\textbf{Data yang dilaporkan digunakan dalam latihan empiris:} -
Populasi municipality dan pertumbuhan penduduk; ukuran populasi yang
dipakai sebagai faktor merujuk pada 2016, sedangkan pertumbuhan kinerja
merujuk pada 2002-2018 (Classifying cities in the Randstad). - Jumlah
pekerjaan per 1.000 penduduk pada 2016 dari Statistics Netherlands. -
Jumlah restoran per 1.000 penduduk pada 2011 dari Bedrijfschap Horeca en
Catering. Tahun 2011 dipilih karena memungkinkan perbandingan
perkembangan restoran pada 1996-2011. - Komposisi sektoral ekonomi
municipality pada 2016 dari Statistics Netherlands, dihitung sebagai
rasio pekerja kerah putih di sektor layanan publik dan privat terhadap
pekerja kerah biru di konstruksi dan manufaktur. - Status sosial pada
2016 dari Netherlands Institute for Social Research. Ukuran ini
menggabungkan pendidikan, posisi dalam pasar kerja, dan pendapatan; data
pada skala neighbourhood dirata-ratakan dengan pembobotan populasi ke
skala municipality. - Kepentingan historis yang dioperasionalkan dari
peringkat kota berdasarkan ukuran populasi pada 1849, tahun sensus
pertama di Belanda. - \emph{Network embeddedness} dari ko-kemunculan
nama tempat Belanda pada seluruh 25 juta situs web dengan ekstensi
\texttt{.nl}, dikoreksi menggunakan gravity model terhadap ukuran dan
jarak. - Pariwisata yang diukur dengan jumlah tempat tidur hotel per
kapita pada 2011 dari Bedrijfschap Horeca en Catering.

Untuk \emph{network embeddedness}, nilai hubungan aktual dibandingkan
dengan nilai prediksi untuk hampir 1.700 tempat di Belanda. Nilai
prediksi kurang dari 0 ditetapkan menjadi 0. Pada skala municipality,
penulis menggunakan hubungan tempat yang memberi nama municipality,
tempat terbesar jika nama tempat municipality tidak digunakan, atau
rata-rata jika nama municipality terdiri atas dua nama (Research
approach).

File TXT tidak menyediakan berkas data mentah, daftar observasi lengkap
di luar daftar klasifikasi pada Figure 2.3, atau format pengunduhan
dataset. Informasi tersebut: Tidak disebutkan dalam file.

\subsection{Population / Sample}\label{population-sample-38}

\textbf{Unit yang diteliti:} Populasi atau unit analisis yang dinyatakan
penulis adalah 120 municipality dalam delimitasi operasional Randstad
yang mereka gunakan. Randstad tidak memiliki delimitasi formal yang
presisi; estimasi penduduk kawasan dapat berkisar antara 7 dan 8 juta.
Delimitasi yang dipakai mengecualikan bekas pulau Goeree-Overflakkee,
Noordoostpolder, sebagian besar utara Holland termasuk Alkmaar, dan
Noordoostpolder di provinsi Flevoland. Kawasan operasional tersebut
memiliki 7.46 juta penduduk pada Januari 2018, dan sekitar sepertiganya
tinggal di empat kota utama: Amsterdam, Rotterdam, The Hague, dan
Utrecht (Classifying cities in the Randstad; Figure 2.1).

Jenis pengambilan sampel probabilistik, sampel individu, dan
representasi di luar 120 municipality: Tidak disebutkan dalam file.
Perbandingan pada dimensi fungsional mencakup 30 municipality dengan
fungsi lebih banyak daripada yang diharapkan dan 38 dengan fungsi lebih
sedikit. Pada dimensi kinerja, kategori yang tumbuh atau berkinerja
lebih baik berisi 27 municipality dan kategori yang lebih buruk berisi
26 (Results: functional dimension; Results: performance dimension).

Sebagai konteks pendahuluan, perhitungan berbasis data ESPON 1.1.1 dan
1.4.3 menyebut 1,001 dari 1,967 kota Eropa berada dalam kawasan
metropolitan multicentric. Dari 265 juta penduduk kota-kota tersebut,
hampir 192 juta atau 72.4\% tinggal di kawasan urban multicentric
(Introduction). Angka ini adalah konteks Eropa, bukan populasi analisis
utama Randstad.

\subsection{Dependent Variables}\label{dependent-variables-38}

File tidak menggunakan istilah ``dependent variable'' secara eksplisit.
\textbf{Inferensi pengolahan, terbatas pada struktur analisis yang
dijelaskan sumber:} skor atau kategori pada dua dimensi berikut
diperlakukan sebagai hasil yang dijelaskan oleh enam faktor dalam
ordinal regression.

\begin{itemize}
\tightlist
\item
  \textbf{Dimensi fungsional:} apakah municipality memiliki lebih
  banyak, lebih sedikit, atau jumlah fungsi yang sesuai harapan
  berdasarkan ukurannya. Fungsi diukur melalui pekerjaan per kapita dan
  restoran per kapita; hasil klasifikasinya tampak dalam Figure 2.3 dan
  hubungan faktornya dianalisis pada Figure 2.10.
\item
  \textbf{Dimensi kinerja:} apakah municipality berkembang atau
  berkinerja lebih baik, lebih buruk, atau sesuai harapan berdasarkan
  ukurannya. Kinerja diukur melalui pertumbuhan penduduk dan
  perkembangan jumlah restoran; hubungan faktornya dianalisis pada
  Figure 2.11.
\item
  \textbf{Keluaran klasifikasi gabungan:} kategori lokal dalam matriks
  \emph{borrowed size/agglomeration shadow}, termasuk \emph{borrowed
  size}, \emph{borrowed functions}, \emph{borrowed performance},
  \emph{as expected}, \emph{functional shadow}, \emph{performance
  shadow}, dan \emph{agglomeration shadow} (Figure 2.2; Conclusion).
\end{itemize}

Enam faktor yang dimasukkan sebagai covariate dalam model tersebut
adalah kepentingan historis, ukuran populasi, komposisi sektoral, status
sosial, \emph{network embeddedness}, dan pariwisata (Research approach).
Penetapan label dependent variable di atas adalah inferensi pengolahan;
sumber hanya menyatakan bahwa model menjelaskan skor pada dimensi
fungsional dan kinerja.

\subsection{Results}\label{results-38}

\textbf{Temuan dilaporkan dalam klasifikasi dan statistik deskriptif:} -
Untuk 120 municipality, rata-rata restoran adalah 0.55 per 1,000
penduduk dengan SD 0.30 dan rentang 0.17-2.47. Rata-rata pekerjaan
adalah 40.11 per 1,000 penduduk dengan SD 15.55 dan rentang 17-107
(Classifying cities in the Randstad). - Rata-rata pertumbuhan penduduk
Randstad adalah 8.4\% pada 2002-2018 dengan SD 10.79. Rata-rata
peningkatan restoran adalah 9.69\% pada 1996-2011 dengan SD 29.32
(Classifying cities in the Randstad). - Figure 2.3 mengklasifikasikan
municipality ke dalam kategori fungsi dan kinerja berdasarkan ukuran.
Kategori yang ditampilkan mencakup \emph{agglomeration shadow},
\emph{performance shadow}, \emph{functional shadow}, \emph{as expected},
\emph{borrowed functions}, \emph{borrowed performance}, dan
\emph{borrowed size}.

\textbf{Temuan dilaporkan pada pola agregat, Figures 2.4-2.9:} -
Municipality yang penting secara historis lebih banyak terwakili dalam
kategori \emph{borrowed size} dan \emph{borrowed function}. Tempat yang
meminjam kinerja cenderung kurang penting secara historis; beberapa
growth centre yang disebut adalah Almere, Houten, Lansingerland, dan
Pijnacker-Nootdorp (Figure 2.4). - Kota yang lebih besar lebih sering
berada dalam kategori \emph{borrowed size} atau \emph{borrowed
functions}, sedangkan kota yang lebih kecil lebih sering berada dalam
\emph{functional shadow} atau \emph{agglomeration shadow}. Almere berada
dalam \emph{borrowed performance} tetapi belum memiliki kepentingan
fungsional; Nissewaard disebut tidak memenuhi ekspektasi yang dibentuk
oleh ukurannya dan mengalami stagnasi dalam periode yang lebih baru
(Figure 2.5). - Municipality dengan manufaktur yang lebih kuat dan
layanan yang relatif lebih sedikit lebih sering mengalami
\emph{agglomeration shadow}. Ekonomi lokal tempat yang meminjam lebih
berciri layanan (Figure 2.6). - Pola status sosial dilaporkan campuran
atau tidak ada. Rotterdam, yang disebut memiliki status sosial terendah,
berada dalam kategori \emph{borrowed size}. Tempat yang tumbuh lebih
banyak cenderung memiliki status sosial agak lebih tinggi, sedangkan
beberapa municipality dalam \emph{performance shadow} termasuk yang
secara sosial lebih lemah (Figure 2.7). - Tempat dalam \emph{borrowed
size} lebih kuat hubungannya dengan tempat lain, sedangkan tempat dalam
\emph{agglomeration shadow} memiliki lebih sedikit hubungan. Tempat
dalam \emph{borrowed functions} juga memiliki hubungan yang lebih kuat;
tempat dalam \emph{performance shadow} sedikit lebih terhubung,
sementara tempat dalam \emph{borrowed performance} lebih lemah
hubungannya. Dugaan bahwa \emph{functional shadow} akan memiliki
hubungan lebih kuat tidak terdukung oleh pola box plot (Figure 2.8). -
Kategori \emph{borrowed size} dan \emph{borrowed functions} lebih
turistik, sedangkan kategori bayangan cenderung kurang turistik.
Pariwisata tidak tampak terkait dengan \emph{borrowed performance} dalam
pola yang dilaporkan. Indikator tempat tidur hotel juga menangkap tempat
yang sentral dalam jaringan bisnis; Haarlemmermeer dan Schiphol disebut
sebagai contoh paling jelas (Figure 2.9).

\textbf{Temuan dilaporkan pada dimensi fungsional:} Municipality dengan
fungsi lebih banyak (N = 30) dibandingkan yang lebih sedikit (N = 38)
secara signifikan lebih penting secara historis,
\texttt{t(66)\ =\ -4.0,\ p\ =\ .000}; lebih besar,
\texttt{t(30)\ =\ 2.5,\ p\ =\ .017}; memiliki status sosial lebih
rendah, \texttt{t(66)\ =\ -2.2,\ p\ =\ .033}; dan menarik lebih banyak
wisatawan, \texttt{t(32)\ =\ 2.4,\ p\ =\ .022}. Hubungan dengan tempat
lain juga lebih kuat, \texttt{t(66)\ =\ 4.9,\ p\ =\ .000}. Perbedaan
orientasi layanan antara kategori tersebut tidak signifikan secara
statistik (Results: functional dimension; Figure 2.10).

Dalam ordinal regression untuk skor dimensi fungsional dengan keenam
covariate sekaligus, pariwisata dan \emph{network embeddedness}
signifikan pada \texttt{p\ \textless{}\ 0.05}, sedangkan kepentingan
historis signifikan hanya pada tingkat \texttt{p\ \textless{}\ 0.10}
(Results: functional dimension).

\textbf{Temuan dilaporkan pada dimensi kinerja:} Dibandingkan
municipality yang berkinerja lebih buruk, municipality yang tumbuh lebih
cepat tampak lebih penting secara historis, lebih besar, memiliki status
sosial lebih tinggi, memiliki lebih banyak pekerjaan sektor layanan,
menerima lebih banyak wisatawan, dan lebih kuat hubungannya dengan
tempat lain. Namun, pada perbandingan mean, hanya perbedaan ukuran yang
signifikan, \texttt{t(28)\ =\ 2.1,\ p\ =\ .046}; perbedaan pariwisata
signifikan hanya pada tingkat \texttt{p\ =\ .094},
\texttt{t(50)\ =\ 1.7}. Kelompok lebih baik berisi 27 municipality dan
kelompok lebih buruk 26, sementara sumber menyebut setidaknya N = 30
direkomendasikan untuk uji-t sampel independen (Results: performance
dimension; Figure 2.11).

Dalam ordinal regression untuk skor kinerja, status sosial adalah faktor
paling signifikan (\texttt{p\ =\ 0.028}), diikuti ukuran populasi
(\texttt{p\ =\ 0.052}) dan \emph{network embeddedness}
(\texttt{p\ =\ 0.091}). Sumber menyatakan bahwa perbedaan pada dimensi
kinerja lebih lemah daripada pada dimensi fungsional (Results:
performance dimension).

\subsection{Findings}\label{findings-38}

\textbf{Temuan dilaporkan:} - Kota sekunder dalam pengertian yang
dipakai buku berada pada posisi yang beragam: Dordrecht dan Vlaardingen
menghadapi \emph{performance shadow}; Schiedam dan Zaanstad berkinerja
sesuai harapan; Delft, Haarlem, dan Hilversum meminjam fungsi; sedangkan
Amersfoort dan Leiden meminjam ukuran. Tidak satu pun kota sekunder yang
disebut tersebut memiliki fungsi lebih sedikit daripada yang diharapkan
(Conclusion). - Kota primer yang disebut, yaitu Amsterdam dan The Hague,
berada dalam kategori \emph{borrowed functions}, sedangkan Utrecht dan
Rotterdam berada dalam \emph{borrowed size} (Conclusion). - Pada dimensi
fungsional, kepentingan historis, ukuran, status sosial yang lebih
rendah, pariwisata, dan terutama hubungan yang lebih kuat dengan kota
lain berasosiasi dengan fungsi yang lebih banyak daripada yang
diharapkan. Pada model multivariat, pariwisata dan \emph{network
embeddedness} mencapai \texttt{p\ \textless{}\ 0.05} (Results:
functional dimension). - Pada dimensi kinerja, ukuran merupakan
perbedaan mean yang signifikan pada perbandingan kategori. Dalam model
ordinal, status sosial memiliki \texttt{p\ =\ 0.028}, diikuti ukuran
populasi \texttt{p\ =\ 0.052} dan \emph{network embeddedness}
\texttt{p\ =\ 0.091} (Results: performance dimension).

\textbf{Interpretasi penulis:} Hasil tersebut mendukung proposisi bahwa
jaringan dapat menggantikan ukuran, khususnya karena tempat yang
memiliki hubungan lebih kuat cenderung memiliki lebih banyak fungsi.
Penulis juga menafsirkan variasi kategori sebagai tanda bahwa kota
sekunder merupakan kelompok yang heterogen dan bahwa keadaan lokal serta
\emph{path dependency} dapat memengaruhi kemampuan kota memanfaatkan
integrasi metropolitan (Results: functional dimension; Conclusion).

\textbf{Inferensi pengolahan:} Bukti yang tersedia dalam file
membandingkan asosiasi dan model ordinal eksploratif; file tidak
menyediakan dasar untuk mengubah temuan tersebut menjadi klaim kausal.
Batas ini konsisten dengan penolakan kausalitas oleh penulis sendiri
(Introduction; Conclusion).

\subsection{Conclusions}\label{conclusions-38}

\textbf{Kesimpulan penulis:} Bab ini memperluas matriks awal
\emph{borrowed size/agglomeration shadow} menjadi kerangka yang
membedakan tujuh hasil lokal: \emph{agglomeration shadow},
\emph{performance shadow}, \emph{functional shadow}, \emph{as expected},
\emph{borrowed functions}, \emph{borrowed performance}, dan
\emph{borrowed size} (Conclusion).

Penerapan pada Randstad menunjukkan bahwa 120 municipality dalam kawasan
metropolitan dengan sekitar 7.5 juta penduduk menempati posisi yang
berbeda. Semua enam faktor yang dipertimbangkan disebut tampak berperan
dalam positioning, meskipun kekuatan asosiasi berbeda ketika dimensi
fungsi dan kinerja dipisahkan. Tempat dengan lebih banyak fungsi
daripada yang diharapkan berkaitan dengan kepentingan historis, populasi
lebih besar, status sosial rata-rata lebih rendah, lebih banyak
wisatawan, dan hubungan yang lebih kuat dengan kota lain. Tempat yang
berkembang lebih kuat daripada yang lain dicirikan terutama oleh
populasi lebih besar, status sosial lebih tinggi, lebih banyak
wisatawan, dan lebih banyak relasi dengan kota lain; box plot
mengisyaratkan bukti untuk dua faktor yang tersisa (Conclusion).

\textbf{Interpretasi penulis:} Integrasi metropolitan dapat positif bagi
kawasan secara keseluruhan tetapi tidak membagi manfaat dan biaya secara
merata di antara kota primer dan berbagai jenis kota sekunder. Karena
jaringan memungkinkan spesialisasi, diperlukan pemahaman tentang
interdependensi baru dan kerangka strategis perencanaan regional untuk
memanfaatkan integrasi sekaligus menyebarkan biaya dan manfaat secara
lebih merata (Conclusion).

\subsection{Contributions}\label{contributions-38}

File tidak memiliki bagian yang secara eksplisit berjudul
``Contributions''. \textbf{Inferensi pengolahan atas kontribusi yang
dapat ditelusuri dari pernyataan penulis:} - Menyempurnakan kerangka
empat kuadran menjadi matriks sembilan kotak dan tujuh label hasil lokal
untuk menangkap municipality yang berkinerja sesuai harapan serta
variasi borrowing dan shadowing (Classifying cities in the Randstad;
Figure 2.2; Conclusion). - Menerapkan dan menguji kerangka tersebut pada
120 municipality di Randstad Holland (Classifying cities in the
Randstad; Conclusion). - Menghubungkan positioning dengan enam faktor
dan mengoperasionalkan \emph{network embeddedness} melalui ko-kemunculan
toponim yang dikoreksi dengan gravity model (Research approach). -
Menghasilkan agenda untuk mengembangkan analisis dinamis, kausal,
historis, dan multi-skala atas hubungan antarkota (Conclusion).

\subsection{Research Gap}\label{research-gap-38}

File tidak memiliki heading ``Research Gap'', tetapi penulis menyebut
beberapa hal yang belum jelas atau perlu dianalisis lebih lanjut: -
Mengapa satu kota memperoleh \emph{borrowed size}, sedangkan kota lain
mengalami \emph{agglomeration shadow} (Introduction). - Bagaimana
pentingnya jaringan dibandingkan faktor lokal, serta apakah faktor lokal
dan faktor jaringan saling menggantikan atau melengkapi (Theory). -
Bagaimana manfaat jaringan pada tingkat kawasan didistribusikan di
antara kota-kota individual, karena efek generatif pada jaringan dapat
menyembunyikan efek distributif intra-kawasan (Theory). - Apakah relasi
antarpasangan kota muncul karena banyaknya fungsi yang sudah dimiliki
kota, atau justru relasi tersebut menggantikan ukuran dan memungkinkan
kota mempertahankan lebih banyak fungsi (Conclusion). - Bagaimana
menjelaskan pergeseran posisi kota dari waktu ke waktu dengan perspektif
dinamis, proksi fungsi dan kinerja yang lebih komprehensif, serta faktor
yang lebih banyak. Penulis juga menyebut kemungkinan pengalaman
borrowing dan shadowing yang simultan serta tumpang tindih (Conclusion).
- Bagaimana mengembangkan ukuran \emph{network embeddedness} pada skala
spasial yang berbeda dan mengukur relasi kota sekunder dengan kota
primer secara khusus (Conclusion).

\textbf{Inferensi pengolahan:} Kesenjangan paling langsung dari desain
yang dilaporkan adalah belum ditetapkannya kausalitas, tetapi ini bukan
klaim baru karena penulis menyatakannya sebagai agenda penelitian pada
Conclusion.

\subsection{Limitations}\label{limitations-38}

\textbf{Keterbatasan yang dinyatakan atau terlihat langsung dalam file:}
- Analisis bersifat eksploratif dan menggunakan asosiasi, bukan hubungan
kausal. Karena itu penulis menyatakan belum dapat menghasilkan
rekomendasi kebijakan yang kuat (Introduction; Conclusion). - Delimitasi
Randstad tidak presisi atau formal, sehingga estimasi penduduknya
bervariasi antara 7 dan 8 juta (Classifying cities in the Randstad). -
Kerangka yang dipakai semula tidak dapat mengklasifikasikan tempat yang
memiliki jumlah fungsi dan lintasan perkembangan sesuai harapan; masalah
ini ditangani dengan perluasan ke matriks sembilan kotak (Classifying
cities in the Randstad). - Snapshot empiris menggunakan jumlah indikator
yang oleh penulis disebut agak terbatas untuk menangkap dimensi fungsi
dan kinerja (Conclusion). - Kelompok perbandingan dimensi kinerja hanya
berisi 27 dan 26 municipality; sumber menyebut setidaknya N = 30
direkomendasikan untuk uji-t sampel independen (Results: performance
dimension). - Ukuran \emph{network embeddedness} didasarkan pada
ko-kemunculan nama tempat di situs \texttt{.nl}, yang oleh penulis
dinyatakan lebih luas daripada koneksi atau arus fisik. Pada tingkat
municipality, ukuran tersebut juga diproksikan melalui satu nama tempat,
tempat terbesar, atau rata-rata dua nama (Research approach).

\textbf{Inferensi pengolahan, bukan keterbatasan yang diberi label
eksplisit oleh penulis:} Indikator fungsi dan kinerja menggunakan tahun
atau rentang waktu yang berbeda: pekerjaan 2016, tingkat restoran 2011,
pertumbuhan penduduk 2002-2018, dan perkembangan restoran 1996-2011.
File tidak menyatakan dampak perbedaan waktu ini terhadap hasil,
sehingga dampaknya tidak dapat ditentukan dari sumber.

\subsection{Challenges}\label{challenges-38}

\textbf{Tantangan yang dibahas penulis:} - Mengatur kawasan metropolitan
multicentric yang memiliki batas dan kewenangan tumpang tindih, tetapi
juga tantangan dan prioritas bersama (Introduction). - Menentukan
bagaimana fungsi dan kinerja yang ``diharapkan'' berdasarkan ukuran
dapat diukur ketika dua indikator dalam satu dimensi memberikan arah
yang berbeda. Penulis menerapkan aturan kategorisasi dan rata-rata skor
terstandar untuk kasus tersebut (Classifying cities in the Randstad). -
Mengukur \emph{network embeddedness} secara presisi. Ko-kemunculan
toponim menangkap keterkaitan yang lebih luas daripada arus fisik, dan
pengembangan metode tersebut masih memerlukan analisis tentang jenis
serta arah relasi (Research approach; Results: aggregate patterns). -
Menangani kompetisi antarkota ketika jumlah fungsi dan kegiatan di
kawasan multicentric terbatas dan distribusinya pada dasarnya merupakan
permainan zero-sum, terutama ketika pendapatan pajak terutama berasal
dari pajak lokal (Theory). - Menjelaskan kausalitas dan perubahan
historis. Penulis menyatakan bahwa penjelasan evolusi kota sekunder
dalam hubungannya dengan kota primer mungkin memerlukan narasi historis,
bukan hanya analisis statistik kuantitatif (Conclusion).

\textbf{Interpretasi penulis:} Heterogenitas kota sekunder dan
distribusi manfaat-biaya integrasi membuat strategi pembangunan yang
seragam menjadi tidak memadai. Tantangan kebijakan bukan hanya mendorong
integrasi, tetapi juga memahami interdependensi dan mengelola
konsekuensi distributifnya (Introduction; Conclusion).

\subsection{Practical Implications}\label{practical-implications-38}

\textbf{Implikasi yang dinyatakan atau dibahas penulis:} - Kebijakan
pengembangan perkotaan perlu menghindari strategi
\emph{one-size-fits-all}, bukan hanya karena perbedaan antara kota
primer dan sekunder, tetapi juga karena heterogenitas di antara kota
sekunder (Introduction). - Perencana dan pembuat kebijakan perlu
memperhatikan efek distributif integrasi antarkota, bukan hanya manfaat
generatif bagi kawasan secara keseluruhan (Theory; Conclusion). -
Kerangka strategis perencanaan regional diperlukan untuk memahami
interdependensi baru, memanfaatkan spesialisasi jaringan, dan
mengupayakan penyebaran biaya serta manfaat integrasi yang lebih merata
(Conclusion).

\textbf{Batas implikasi:} Penulis menyatakan bahwa analisis ini tidak
dimaksudkan untuk menghasilkan rekomendasi kebijakan yang kuat karena
belum menetapkan kausalitas (Introduction). Implikasi di atas adalah
arah yang dibahas sumber, bukan evaluasi dampak kebijakan.

\subsection{Applications}\label{applications-38}

\textbf{Penerapan yang dilaporkan:} Kerangka digunakan pada studi kasus
Randstad Holland untuk memposisikan 120 municipality ke dalam kategori
\emph{borrowed size/agglomeration shadow}. Figure 2.1 menyajikan
Randstad Holland, Figure 2.2 menyajikan kerangka yang diperluas, Figure
2.3 menyajikan klasifikasi municipality, dan Figures 2.4-2.11 menyajikan
asosiasi faktor dengan dimensi fungsional atau kinerja.

Penerapan di luar latihan empiris Randstad atau penggunaan operasional
lainnya: Tidak disebutkan dalam file.

\subsection{Future Research
(Optional)}\label{future-research-optional-38}

\textbf{Agenda yang disebutkan penulis:} - Menggunakan perspektif
dinamis untuk menjelaskan pergeseran posisi kota, dengan proksi fungsi
dan kinerja yang lebih komprehensif serta himpunan faktor yang lebih
besar (Conclusion). - Menyelidiki kemungkinan bahwa kota mengalami
borrowing dan shadowing secara simultan dan tumpang tindih, sehingga
positioning perlu dicatat secara lebih hati-hati (Conclusion). -
Menganalisis kausalitas secara lebih hati-hati, termasuk kemungkinan
menggunakan pendekatan ``instrumented variables'' sebagaimana istilah
dalam file, dan melengkapi analisis kuantitatif dengan kajian evolusi
historis kota sekunder dan kota primer (Conclusion). - Mengembangkan
variabel \emph{network embeddedness} pada skala spasial yang berbeda dan
mengukur tingkat hubungan kota sekunder dengan kota primer secara khusus
(Conclusion). - Menelaah apakah relasi antarkota menyebabkan kota
mempertahankan lebih banyak fungsi atau justru terbentuk karena kota
telah memiliki banyak fungsi (Conclusion).

Semua butir di atas bersumber dari agenda yang dinyatakan penulis;
tambahan agenda di luar daftar tersebut: Tidak disebutkan dalam file.

\subsection{Provenance Notes}\label{provenance-notes-38}

\begin{itemize}
\tightlist
\item
  Sumber yang dibaca seluruhnya adalah
  \texttt{source/journal/A23-meijers-cardoso-2021-shedding-light-or-casting-shadows-relations-between-primary-and-secondary-cities-10.1332-policypress-9781529212075.003.0002.txt},
  setelah nama file diverifikasi berdasarkan prefiks \texttt{A23} di
  bawah \texttt{source/}.
\item
  TXT berisi 1,206 baris yang mencakup blok
  \texttt{{[}PDF\ Metadata{]}}, \texttt{{[}Extracted\ Text{]}}, isi bab,
  dan References. Seluruh teks sampai akhir References dibaca. Tidak ada
  perubahan pada TXT.
\item
  Metadata ekstraksi menyebut PDF diekstrak pada 2026-08-31 dengan
  \texttt{pdftotext\ -layout}. Metadata PDF menyebut 35 halaman, tidak
  terenkripsi, ukuran halaman 612 x 792 points, dan ukuran berkas
  1,028,869 bytes. Creator dan Producer tercatat sebagai Microsoft Word
  for Microsoft 365; field Author tercatat \texttt{pendras}.
\item
  Isi file mengidentifikasi karya ini sebagai Chapter 2, ``Shedding
  Light or Casting Shadows? Relations between Primary and Secondary
  Cities'', oleh Evert Meijers dan Rodrigo Cardoso. Citation block
  menyebut versi buku dalam M. Pendras dan C. Williams (Eds.),
  \emph{Secondary Cities: Exploring Uneven Development in Dynamic Urban
  Regions of the Global North}, halaman 25-54, Bristol University Press
  (bagian pembuka file).
\item
  Metadata \texttt{Source\ PDF} menunjuk ke
  \texttt{book/A23-meijers-cardoso-2021-shedding-light-or-casting-shadows-relations-between-primary-and-secondary-cities-10.1332-policypress-9781529212075.003.0002.pdf},
  sedangkan TXT yang diverifikasi berada di \texttt{source/journal/}.
  Perbedaan jalur metadata dan lokasi TXT dicatat sebagai batas
  provenance. Isi TXT tetap dapat dibaca lengkap, sehingga perbedaan ini
  tidak dipakai untuk menambah atau menebak informasi bibliografis.
\item
  Locator yang dipertahankan dalam review adalah nama bagian sumber,
  Figures 2.1-2.11, dan angka statistik atau model yang disebut pada
  bagian Results. Nomor halaman cetak per klaim tidak terlihat di teks
  hasil ekstraksi; yang tersedia adalah rentang halaman 25-54 pada blok
  citation dan metadata 35 halaman PDF. Tabel atau lampiran: Tidak
  disebutkan dalam file.
\item
  Klaim substantif review ini dibatasi pada TXT A23. References dibaca
  sebagai bagian dari teks, tetapi tidak digunakan untuk menambahkan
  bukti atau klaim dari sumber lain. Pemisahan ``Temuan dilaporkan'',
  ``Interpretasi penulis'', dan ``Inferensi pengolahan'' menunjukkan
  status klaim masing-masing.
\end{itemize}

\section{Review A24 - Border cities: Out of the
shadow}\label{review-a24---border-cities-out-of-the-shadow}

Status: Review literatur berbasis teks lengkap hasil ekstraksi PDF.

\subsection{Identitas Artikel}\label{identitas-artikel-8}

\begin{itemize}
\tightlist
\item
  \textbf{Kode:} A24.
\item
  \textbf{Judul:} \emph{Border cities: Out of the shadow}.
\item
  \textbf{Penulis:} Christophe Sohn, Julien Licheron, dan Evert Meijers.
\item
  \textbf{Afiliasi:} Luxembourg Institute of Socio-Economic Research,
  Luxembourg; dan Utrecht University, Netherlands.
\item
  \textbf{Jurnal:} \emph{Papers in Regional Science}, 101(2), 417-438,
  tahun 2022.
\item
  \textbf{DOI:} 10.1111/pirs.12653.
\item
  \textbf{Riwayat naskah:} diterima 10 Juni 2021, direvisi 29 September
  2021, dan diterima 7 Desember 2021.
\item
  \textbf{Kata kunci:} \emph{agglomeration shadow}, \emph{border
  effects}, \emph{borrowed size}, \emph{European border cities}, dan
  \emph{potential urban interaction}.
\item
  \textbf{Jenis sumber:} artikel jurnal empiris. Status telaah sejawat:
  Tidak disebutkan dalam file.
\item
  \textbf{Wilayah dan unit analisis:} 1.920 morphological urban areas
  (MUA) yang disebut sebagai kota, di EU 27 ditambah Norwegia dan Swiss;
  355 di antaranya dikategorikan sebagai kota perbatasan.
\item
  \textbf{Periode data utama:} data fungsi metropolitan dikumpulkan pada
  2003-2009 dengan 2006 sebagai tahun rujukan utama; variabel lain
  memakai tahun yang dijelaskan dalam Bagian Dataset.
\item
  \textbf{Locator identitas:} hlm. 417, judul, abstrak, afiliasi, dan
  kata kunci; hlm. 438, format sitasi.
\end{itemize}

\subsection{Summarized Abstract}\label{summarized-abstract-39}

\textbf{Laporan sumber.} Artikel menguji bagaimana keberadaan perbatasan
nasional mengubah pola \emph{borrowed size} dan \emph{agglomeration
shadow} yang biasanya muncul dari kedekatan antarkota. Analisis memakai
distribusi fungsi metropolitan pada 1.920 kota Eropa, termasuk 355 kota
perbatasan, serta menghitung interaksi potensial berdasarkan ukuran
penduduk dan jarak melalui persamaan gravitasi. Dengan model
\emph{spatial autoregressive} (SAR), penulis membandingkan interaksi
domestik dan lintas batas, membedakan perbatasan terkendali dan terbuka
serta durasi keterbukaannya, lalu memisahkan fungsi yang digerakkan
pasar dan fungsi yang digerakkan publik. (Hlm. 417, abstrak; hlm.
418-419, Bagian 1; hlm. 424-426, Bagian 3.)

\textbf{Temuan yang dilaporkan sumber.} Dalam konteks domestik,
kota-kota berukuran serupa memperoleh lebih banyak fungsi metropolitan,
sedangkan kota yang lebih kecil mengalami \emph{agglomeration shadow}
ketika berdekatan dengan kota yang lebih besar. Di konteks lintas batas,
tidak ditemukan efek bayangan dari kota besar terhadap kota kecil; kota
kecil justru dapat memperoleh manfaat dari kedekatan dengan kota besar
di seberang perbatasan. Efek \emph{borrowed size} antarkota berukuran
serupa lebih kuat lintas batas daripada domestik. Efek tersebut tidak
muncul pada perbatasan terkendali dan menguat seiring durasi keterbukaan
perbatasan. Manfaat lintas batas dari kota besar terhadap kota kecil
hanya ditemukan untuk fungsi yang digerakkan pasar, sedangkan fungsi
publik lebih sensitif terhadap efek penghalang dan perlindungan
perbatasan. (Hlm. 417, abstrak; hlm. 427-430, Bagian 4.)

\textbf{Interpretasi penulis.} Penulis menyimpulkan bahwa perbatasan
memoderasi regularitas normal sistem permukiman: peran perlindungan
perbatasan dapat menahan bayangan aglomerasi dan, ketika integrasi
berlangsung lebih lama, memperbesar peluang peminjaman ukuran. Karena
keunggulan ini, kota perbatasan dinilai perlu keluar dari bayang-bayang.
(Hlm. 431-432, Bagian 5.)

\textbf{Inferensi pengolahan.} Hasil artikel menunjukkan asosiasi
bersyarat antara proksi interaksi potensial dan indeks fungsi
metropolitan. Karena interaksi aktual tidak tersedia dan model
menggunakan data fungsi terutama pada tahun rujukan 2006, hasil tersebut
tidak dengan sendirinya mengidentifikasi arus pemanfaatan fungsi atau
efek kausal perbatasan. (Hlm. 424-426, Bagian 3.2-3.3.)

\subsection{Problem Statement}\label{problem-statement-39}

\textbf{Laporan sumber.} Literatur tentang eksternalitas jaringan kota
mengakui bahwa hubungan antarkota dapat memengaruhi pertumbuhan dan
kinerja metropolitan, tetapi keterikatan teritorial kota dan
kedekatannya dengan perbatasan nasional sebagian besar diabaikan.
Perhatian juga lebih banyak diarahkan pada keterhubungan kota besar
dalam jaringan global, sementara relasi regional antara kota kecil dan
menengah serta kasus polisentris lintas batas jarang dipertimbangkan.
(Hlm. 418, Bagian 1.)

\textbf{Laporan sumber.} Masalah pengukuran juga penting karena data
interaksi antarkota pada skala kota, lintas banyak negara, dan untuk
berbagai jenis fungsi sulit diperoleh. Pada skala regional, penelitian
biasanya mengasumsikan integrasi dari kedekatan atau tumpang tindih
wilayah fungsional, padahal tingkat pengaruh kota tetangga dapat berbeda
menurut ukuran dan jarak. (Hlm. 418-419, Bagian 1.)

\textbf{Laporan sumber.} Artikel menanyakan bagaimana perbatasan
memoderasi dampak keberadaan kota tetangga terhadap kinerja
metropolitan. Penulis menguji apakah perbatasan mengurangi peluang kota
besar meminjam ukuran dari kota kecil, melindungi kota kecil dari
bayangan kota besar, dan mengurangi manfaat hubungan antarkota berukuran
serupa. Mereka juga mempertanyakan apakah hasil tersebut berubah menurut
tingkat keterbukaan dan durasi pembukaan perbatasan serta menurut jenis
fungsi metropolitan. (Hlm. 418-419 dan 420-424, Bagian 1-2.)

\textbf{Inferensi pengolahan.} Problem statement artikel memiliki dua
lapisan: menjelaskan perbedaan kinerja fungsi metropolitan menurut
konfigurasi tetangga, dan menentukan apakah perbedaan tersebut
dimoderasi oleh institusi perbatasan. Lapisan kedua sulit dipisahkan
secara langsung karena ukuran interaksi lintas batas dibentuk sebagai
potensi berdasarkan populasi dan jarak, bukan arus aktual. (Hlm.
424-426, Bagian 3.2-3.3.)

\subsection{Objectives}\label{objectives-39}

\textbf{Laporan sumber.} Tujuan utama penelitian adalah menguji
bagaimana perbatasan nasional memoderasi pola \emph{borrowed size} dan
\emph{agglomeration shadow} yang dihasilkan oleh interaksi dengan kota
tetangga. Tujuan tersebut diwujudkan dalam beberapa pengujian:

\begin{itemize}
\tightlist
\item
  menguji efek kedekatan dengan kota yang lebih kecil, berukuran serupa,
  dan lebih besar terhadap jumlah fungsi metropolitan;
\item
  membandingkan efek interaksi dengan kota dalam negara yang sama dan
  dengan kota di seberang perbatasan;
\item
  menguji peran status perbatasan, yaitu terkendali atau terbuka, serta
  durasi keterbukaan perbatasan;
\item
  membandingkan efek perbatasan pada fungsi metropolitan yang digerakkan
  pasar dan yang digerakkan publik; dan
\item
  memeriksa ketahanan hasil terhadap definisi kota perbatasan, koefisien
  peluruhan jarak, dan kelompok ukuran kota.
\end{itemize}

Tujuan dan hipotesis dinyatakan dalam Bagian 1-2 dan diterjemahkan ke
dalam tiga tahap estimasi pada Bagian 3. (Hlm. 418-419, 420-424, dan
427-431.)

\subsection{Literature Survey}\label{literature-survey-39}

\textbf{Laporan sumber.} Tinjauan pustaka menghubungkan dua rumpun
literatur, yaitu eksternalitas jaringan kota dan efek perbatasan.
Kerangka jaringan kota dikembangkan untuk melampaui model tempat sentral
yang menekankan hierarki serta kontrol wilayah pasar. Eksternalitas
jaringan dipahami sebagai manfaat ekonomi yang dapat diperoleh
perusahaan dan rumah tangga dari lokasi di kota yang tertanam dalam
jaringan kota lain. Pada wilayah perkotaan multisentris, integrasi yang
lebih dekat dapat menghasilkan eksternalitas positif, tetapi manfaatnya
tidak harus terbagi sama kepada semua kota. (Hlm. 419, Bagian 2.1.)

\textbf{Laporan sumber.} \emph{Borrowed size} ditelusuri kepada Alonso
(1973) dan diperluas oleh Meijers dan Burger sebagai situasi ketika
sebuah kota mempunyai fungsi urban yang biasanya diasosiasikan dengan
kota yang lebih besar karena berinteraksi dengan kota lain pada beberapa
skala spasial. \emph{Agglomeration shadow} dipakai untuk menjelaskan
sisi kompetitifnya, yaitu ketika sebuah kota mempunyai fungsi lebih
sedikit atau pertumbuhan lebih rendah daripada yang diharapkan karena
persaingan dengan kota tetangga. Kedua konsep tersebut diposisikan
sebagai dua kemungkinan hasil dari interaksi antarkota. (Hlm. 419-420,
Bagian 2.1.)

\textbf{Laporan sumber.} Dari kerangka tersebut, penulis mengajukan tiga
hipotesis untuk konteks tanpa pembedaan perbatasan: kota besar yang
dekat dengan kota lebih kecil mempunyai lebih banyak fungsi daripada
yang diharapkan (H1); kota kecil yang dekat dengan kota lebih besar
mempunyai lebih sedikit fungsi daripada yang diharapkan (H2); dan
kota-kota yang berukuran kurang lebih sama memperoleh lebih banyak
fungsi daripada yang diharapkan (H3). (Hlm. 420, Bagian 2.1.)

\textbf{Laporan sumber.} Literatur efek perbatasan banyak menggunakan
model gravitasi untuk mengukur pengaruh hambatan politik dan
administratif terhadap arus perdagangan. Studi yang dirujuk menunjukkan
bahwa batas negara dapat tetap mengurangi arus meskipun hambatan formal
dihapus, tetapi temuan mengenai manfaat integrasi bagi wilayah
perbatasan tidak seragam. Literatur tersebut juga menekankan bahwa
dampak dapat berubah menurut waktu pembukaan perbatasan dan bahwa
hambatan fisik, institusional, serta budaya tidak hilang secara
otomatis. (Hlm. 421-422, Bagian 2.2.)

\textbf{Laporan sumber.} Penulis lalu merumuskan hipotesis moderasi
perbatasan. H4 menyatakan kota perbatasan besar lebih sedikit memperoleh
manfaat dari kota kecil di seberang perbatasan dibandingkan konteks
domestik. H5 menyatakan kota perbatasan kecil lebih sedikit terkena
bayangan kota besar lintas batas. H6 menyatakan kota perbatasan
berukuran serupa mengalami \emph{borrowed size} yang lebih lemah. H7
menyatakan kota yang dipisahkan perbatasan terkendali tidak mengalami
efek interaksi yang signifikan. H8 menyatakan durasi pembukaan yang
lebih lama memperkuat efek \emph{borrowed size} dan \emph{agglomeration
shadow} lintas batas. (Hlm. 422-423, Bagian 2.2.)

\textbf{Laporan sumber.} H9 menyatakan efek moderasi perbatasan lebih
kuat untuk fungsi yang digerakkan publik daripada fungsi yang digerakkan
pasar. Dasarnya adalah dugaan bahwa fungsi pasar mengikuti dukungan
basis lintas batas secara lebih rasional, sedangkan fungsi publik lebih
dipengaruhi hambatan bahasa, budaya, institusi, kedaulatan wilayah, dan
keinginan otoritas untuk mempertahankan fasilitas di masing-masing sisi
perbatasan. (Hlm. 423-424, Bagian 2.2.)

\textbf{Interpretasi penulis.} Literatur tersebut digunakan untuk
mengharapkan dua peran perbatasan yang berlawanan: perbatasan dapat
menjadi penghalang yang mengurangi integrasi dan sekaligus pelindung
yang mengurangi kompetisi atau duplikasi rasionalisasi fungsi. Peran
aktualnya diperkirakan bergantung pada ukuran kota tetangga, durasi
keterbukaan, dan jenis fungsi. (Hlm. 421-424, Bagian 2.2.)

\textbf{Inferensi pengolahan.} Bagian ini merupakan tinjauan naratif
untuk membangun hipotesis, bukan tinjauan sistematis. Strategi
pencarian, kriteria inklusi, jumlah literatur yang disaring, dan
prosedur penilaian mutu literatur terdahulu: Tidak disebutkan dalam
file.

\subsection{Method Used}\label{method-used-39}

\textbf{Laporan sumber.} Penelitian menggunakan model \emph{spatial
autoregressive} (SAR) dalam tiga tahap. Tahap pertama mengestimasi
hubungan antara indeks fungsi metropolitan, variabel kota, dan pengaruh
kota tetangga tanpa membedakan perbatasan. Matriks bobot spasial
digunakan untuk mengukur derajat pengaruh tetangga berdasarkan interaksi
potensial. (Hlm. 424, Bagian 3.1, Persamaan 1.)

\textbf{Laporan sumber.} Tahap kedua memisahkan pengaruh domestik dan
lintas batas. Matriks \texttt{W\_D} hanya mempertahankan pengaruh
potensial kota di negara yang sama, sedangkan matriks \texttt{W\_B}
hanya mempertahankan pengaruh potensial kota perbatasan di negara
berbeda. Variabel dummy \texttt{B} menandai kota perbatasan. Tahap
ketiga menggunakan struktur model tersebut dengan dua keluaran khusus,
yaitu fungsi metropolitan yang digerakkan pasar dan fungsi metropolitan
yang digerakkan publik. (Hlm. 424, Bagian 3.1, Persamaan 2; hlm.
424-425, Bagian 3.2.)

\textbf{Laporan sumber.} Interaksi potensial pasangan kota dihitung
dengan persamaan gravitasi berdasarkan populasi kedua kota dan jarak
Euclidian, dengan koefisien proporsionalitas ditetapkan 1 dan koefisien
peluruhan jarak \texttt{alpha\ =\ 2}. Penulis tidak menganggap nilai
tersebut sebagai arus tertentu; interaksi yang dihitung adalah potensi
umum. Interaksi dengan jarak lebih dari 75 km dikeluarkan. Robustness
check memakai \texttt{alpha\ =\ 1,5} dan \texttt{alpha\ =\ 2,5}. (Hlm.
425-426, Bagian 3.3.)

\textbf{Laporan sumber.} Kota tetangga dikelompokkan menurut
perbandingan populasinya dengan kota yang dikaji: lebih kecil jika
populasinya di bawah 50 persen, berukuran serupa jika lebih dari 50
persen tetapi kurang dari 200 persen, dan lebih besar jika lebih dari
200 persen. Variabel kontrol mencakup populasi kota, pertumbuhan
penduduk, status ibu kota, keberadaan organisasi supranasional,
konektivitas global, status negara federal, keragaman fungsi, dan dummy
negara sebagaimana dijelaskan dalam teks serta Lampiran Tabel A1. (Hlm.
425-427, Bagian 3.3-3.5; hlm. 435, Tabel A1.)

\textbf{Laporan sumber.} Kota perbatasan ditentukan melalui waktu tempuh
mobil kurang dari 45 menit menuju perbatasan terdekat. Perbatasan
terkendali adalah perbatasan eksternal Uni Eropa pada 2006, dengan
pengecualian perbatasan EU-EFTA yang memiliki perjanjian khusus;
perbatasan terbuka adalah perbatasan internal EU dan perbatasan bersama
dengan negara EFTA. Durasi pembukaan dikelompokkan menjadi masih
terkendali, kurang dari 10 tahun, 10-20 tahun, dan lebih dari 20 tahun.
(Hlm. 426, Bagian 3.4; hlm. 435, Tabel A1.)

\textbf{Laporan sumber.} Ketahanan hasil diuji dengan tiga perubahan
asumsi: definisi kota perbatasan kurang dari 25 km dari perbatasan
berdasarkan jarak garis lurus; koefisien peluruhan jarak 1,5 dan 2,5;
serta estimasi terpisah untuk kota kecil, menengah, dan besar. Hasil
tambahan disebut tersedia dalam lampiran. (Hlm. 430-431, Bagian 4.4;
Lampiran Tabel A2-A4.)

\textbf{Inferensi pengolahan.} Rancangan ini adalah studi observasional
dengan model spasial dan indeks fungsi metropolitan sebagai keluaran.
Persamaan gravitasi memberikan proksi interaksi potensial, bukan
pengukuran arus, sedangkan kategori durasi perbatasan membandingkan
kondisi pada satu kerangka tahun 2006. Dengan demikian, hasil lebih
tepat dibaca sebagai asosiasi bersyarat daripada identifikasi kausal
atau pengamatan perubahan kinerja dari waktu ke waktu. (Hlm. 424-426,
Bagian 3.)

Perangkat lunak estimasi dan kode replikasi: Tidak disebutkan dalam
file.

\subsection{Dataset}\label{dataset-39}

\textbf{Laporan sumber.} Basis data fungsi metropolitan dibuat oleh
German Federal Spatial Planning Agency pada 2010 dan diterbitkan sebagai
BBSR (2011), lalu diolah kembali oleh penulis. Data awal diperoleh pada
skala munisipalitas LAU 2 di seluruh benua Eropa dan diagregasikan ke
skala \emph{morphological urban area} (MUA) yang didefinisikan ESPON
agar kota yang terdiri atas beberapa yurisdiksi kecil dapat dianalisis
sebagai satu unit. (Hlm. 424-425, Bagian 3.2.)

\textbf{Laporan sumber.} Indeks fungsi metropolitan mencakup 28 variabel
dalam empat domain:

\begin{itemize}
\tightlist
\item
  \textbf{Ekonomi:} kantor pusat perusahaan top-500, jasa produsen
  tingkat lanjut, bank, dan pameran dagang.
\item
  \textbf{Sains:} universitas top-500, asosiasi ilmiah, dan kongres
  internasional.
\item
  \textbf{Transportasi:} transportasi penumpang udara dan kereta api,
  transportasi barang maritim, dan lalu lintas data.
\item
  \textbf{Budaya dan olahraga:} teater, gedung opera, acara musik,
  galeri, stadion olahraga, dan acara olahraga penting.
\end{itemize}

Data fungsi dikumpulkan selama 2003-2009, dengan 2006 sebagai tahun
rujukan utama. Nilai indeks berkisar dari 0, yang berarti tidak ada
fungsi metropolitan, sampai 100, yang berarti maksimum. (Hlm. 424-425,
Bagian 3.2.)

\textbf{Laporan sumber.} Data variabel kontrol dan perbatasan pada
Lampiran Tabel A1 bersumber dari beberapa basis: IGEAT (2007) untuk
populasi MUA dan dummy negara; ESPON (2012) untuk waktu tempuh menuju
perbatasan; ESPON (2016) untuk pertumbuhan penduduk dan basis data MUA;
BBSR (2011) untuk ibu kota, organisasi supranasional, dan keragaman
fungsi; Taylor (2000) atau dataset GaWC 12 untuk konektivitas global;
Forum of Federations (2015) untuk status federal; European Commission
(2021) dan berbagai sumber untuk status serta durasi pembukaan
perbatasan. Tahun rujukan setiap variabel dijelaskan dalam Lampiran
Tabel A1. (Hlm. 435, Tabel A1.)

\textbf{Laporan sumber.} Karena data komposit mengenai arus perdagangan,
komuter, komunikasi, korporasi, dan kolaborasi ilmiah tidak tersedia
pada granularitas 1.920 kota dan cakupan 29 negara, penulis membentuk
interaksi potensial dari populasi dan jarak. Data pekerja lintas batas
disebut tidak tersedia untuk seluruh Eropa pada tingkat geografis yang
diperlukan. (Hlm. 425-426, Bagian 3.3.)

\textbf{Inferensi pengolahan.} Dataset tidak mengukur satu jenis arus
atau penggunaan fungsi tertentu. Ia menggabungkan atribut fungsi
metropolitan dengan proksi hubungan spasial, sehingga hasil tidak dapat
langsung dibaca sebagai volume perdagangan, mobilitas, atau pertukaran
pengetahuan. Tanggal pengunduhan, prosedur pembersihan, aturan
deduplikasi, dan ketersediaan data mentah: Tidak disebutkan dalam file.

\subsection{Population / Sample}\label{population-sample-39}

\textbf{Laporan sumber.} Populasi analitis mencakup 1.920 MUA di 29
negara, yaitu EU 27 ditambah Norwegia dan Swiss. MUA merupakan
pengelompokan munisipalitas yang memiliki sedikitnya 20.000 penduduk.
Penulis menggunakan istilah ``kota'' untuk MUA demi kesederhanaan.
Sebanyak 355 dari 1.920 unit dikategorikan sebagai kota perbatasan
berdasarkan waktu tempuh mobil kurang dari 45 menit ke perbatasan
terdekat. (Hlm. 424-425, Bagian 3.2; hlm. 426, Bagian 3.4.)

\textbf{Laporan sumber.} Observasi analisis utama berjumlah 1.920. Tidak
ada partisipan individual atau sampel manusia; unit analisisnya adalah
MUA atau kota pada tingkat wilayah. Analisis subkelompok menggunakan
definisi alternatif kota perbatasan dan kelas ukuran kota sebagaimana
dijelaskan dalam Bagian Method Used. (Hlm. 424-431, Bagian 3-4.)

\textbf{Inferensi pengolahan.} File menyajikan cakupan unit analisis
sebagai keseluruhan 1.920 kota dalam wilayah studi, bukan prosedur
pengambilan sampel probabilistik. Karena data fungsi memakai 2006
sebagai tahun utama, generalisasi empiris terutama berlaku pada sistem
kota Eropa dan kerangka klasifikasi perbatasan yang dipakai artikel,
bukan otomatis pada kota atau periode di luar cakupan tersebut.

\subsection{Dependent Variables}\label{dependent-variables-39}

\textbf{Laporan sumber.} Variabel terikat utama adalah indeks fungsi
metropolitan (\texttt{Y}). Indeks ini dibangun dari 28 variabel pada
domain ekonomi, sains, transportasi, budaya, dan olahraga, dengan
rentang nilai 0-100. Model pertama dan kedua memakai indeks seluruh
fungsi metropolitan. (Hlm. 424-425, Bagian 3.1-3.2.)

\textbf{Laporan sumber.} Pada tahap ketiga, indeks dipisahkan menjadi
dua variabel terikat: \textbf{fungsi yang digerakkan pasar}, yang
mencakup domain ekonomi dan transportasi dan terutama dikendalikan aktor
privat; serta \textbf{fungsi yang digerakkan publik}, yang mencakup
domain sains dan aktivitas budaya yang dikendalikan pemangku kepentingan
publik. Dengan demikian, penelitian menggunakan tiga keluaran: seluruh
fungsi metropolitan, fungsi pasar, dan fungsi publik. (Hlm. 424-425,
Bagian 3.1-3.2.)

\textbf{Inferensi pengolahan.} Variabel terikat mengukur keberadaan atau
indeks fungsi metropolitan pada tingkat MUA, bukan arus manfaat, jumlah
pengguna, kualitas layanan, produktivitas, pertumbuhan penduduk, atau
pertumbuhan pekerjaan. Pemisahan pasar-publik merupakan klasifikasi
domain komposit artikel, bukan pengukuran terpisah atas seluruh dimensi
institusional perbatasan. (Hlm. 424-425, Bagian 3.2.)

\subsection{Results}\label{results-39}

\textbf{Laporan sumber: model dasar.} Model 1 pada Tabel 1 memiliki
\texttt{pseudo\ R2\ =\ 0,4780} dan \texttt{n\ =\ 1.920}. Di antara
variabel kontrol, populasi kota, pertumbuhan penduduk, keberadaan
organisasi supranasional, konektivitas global, dan keragaman fungsi
berpengaruh positif serta signifikan; status ibu kota dan status negara
federal tidak signifikan. Interaksi dengan kota yang lebih kecil tidak
signifikan (\texttt{0,0009}), interaksi dengan kota berukuran serupa
positif dan signifikan (\texttt{0,0494}, taraf 5\%), sedangkan interaksi
dengan kota yang lebih besar negatif dan signifikan (\texttt{-0,3738},
taraf 5\%). (Hlm. 427, Tabel 1.)

\textbf{Interpretasi penulis: model dasar.} Tidak adanya efek signifikan
dari kota yang lebih kecil membuat H1 tidak terkonfirmasi. Efek positif
dari kota berukuran serupa mendukung H3 dan dipahami sebagai
\emph{borrowed size} yang menguntungkan kota-kota dalam wilayah
polisentris. Efek negatif dari kota yang lebih besar mendukung H2 dan
dipahami sebagai \emph{agglomeration shadow} terhadap kota yang lebih
kecil. (Hlm. 427-428, Bagian 4.1.)

\textbf{Laporan sumber: perbandingan domestik dan lintas batas.} Model 2
pada Tabel 2 memiliki \texttt{pseudo\ R2\ =\ 0,4828} dan
\texttt{n\ =\ 1.920}. Untuk interaksi domestik, koefisien kota yang
lebih kecil tidak signifikan (\texttt{0,0012}), kota berukuran serupa
positif pada taraf 10\% (\texttt{0,0512}), dan kota yang lebih besar
negatif pada taraf 1\% (\texttt{-0,4316}). Untuk interaksi lintas batas,
kota yang lebih kecil tidak signifikan (\texttt{0,0012}), sedangkan kota
berukuran serupa (\texttt{0,0815}) dan kota yang lebih besar
(\texttt{0,0326}) positif dan signifikan pada taraf 5\%. Estimasi
variabel kontrol dan pengaruh domestik tidak berubah secara berarti
dibandingkan model dasar. (Hlm. 428-429, Tabel 2.)

\textbf{Interpretasi penulis: perbandingan domestik dan lintas batas.}
Tidak adanya efek signifikan kota yang lebih kecil terhadap kota besar
dalam konteks domestik maupun lintas batas membuat H4 tidak
terkonfirmasi. Berbeda dari bayangan yang ditemukan dalam konteks
domestik, kota besar di seberang perbatasan tidak menimbulkan
\emph{agglomeration shadow}; kota kecil bahkan memperoleh manfaat
kedekatan tersebut. H5 terverifikasi, tetapi efek perlindungannya lebih
kuat daripada dugaan awal karena bayangan bukan hanya berkurang,
melainkan tidak ditemukan dan digantikan peminjaman ukuran. Efek positif
kota berukuran serupa lintas batas lebih kuat daripada efek domestik,
sehingga H6 tidak terverifikasi. (Hlm. 428-429, Bagian 4.2.)

\textbf{Laporan sumber: status dan durasi keterbukaan perbatasan.} Model
3 memiliki \texttt{pseudo\ R2\ =\ 0,4964} dan \texttt{n\ =\ 1.920}. Pada
perbatasan terkendali, pengaruh kota yang lebih kecil
(\texttt{-0,0005}), berukuran serupa (\texttt{-0,0014}), dan lebih besar
(\texttt{-0,0004}) tidak signifikan. Pada perbatasan yang terbuka kurang
dari 10 tahun, hanya pengaruh kota yang lebih besar yang signifikan pada
taraf 10\% (\texttt{0,0261}). Pada perbatasan yang terbuka 10-20 tahun,
pengaruh kota berukuran serupa (\texttt{0,0801}) dan lebih besar
(\texttt{0,0274}) signifikan; pada perbatasan yang terbuka lebih dari 20
tahun, keduanya tetap positif dan lebih kuat (\texttt{0,1261} dan
\texttt{0,0284}). Pengaruh kota yang lebih kecil tidak signifikan pada
semua kategori durasi. (Hlm. 429, Tabel 2.)

\textbf{Interpretasi penulis: status dan durasi.} Tidak adanya efek
interaksi pada perbatasan terkendali memverifikasi H7. Kenaikan efek
peminjaman dari kota berukuran serupa atau lebih besar seiring durasi
keterbukaan sejalan dengan H8, tetapi H8 tidak terkonfirmasi sepenuhnya
karena durasi tidak mengubah pengaruh kota yang lebih kecil. (Hlm.
428-429, Bagian 4.2.)

\textbf{Laporan sumber: fungsi pasar dan fungsi publik.} Model 4 untuk
fungsi pasar memiliki \texttt{pseudo\ R2\ =\ 0,5187}, sedangkan model 5
untuk fungsi publik memiliki \texttt{pseudo\ R2\ =\ 0,4842}; keduanya
menggunakan \texttt{n\ =\ 1.920}. Dalam konteks domestik, bayangan dari
kota yang lebih besar lebih kuat untuk fungsi pasar (\texttt{-0,7321},
taraf 1\%) daripada fungsi publik (\texttt{-0,3235}, taraf 5\%).
Peminjaman dari kota berukuran serupa domestik signifikan untuk fungsi
pasar (\texttt{0,0812}, taraf 10\%) tetapi tidak untuk fungsi publik
(\texttt{0,0203}). (Hlm. 429-430, Tabel 3.)

\textbf{Laporan sumber.} Dalam konteks lintas batas, manfaat kedekatan
kota yang lebih besar bagi kota kecil hanya signifikan untuk fungsi
pasar (\texttt{0,0302}, taraf 10\%), bukan fungsi publik
(\texttt{0,0105}, tidak signifikan). Pengaruh kota berukuran serupa
positif untuk fungsi pasar (\texttt{0,1022}, taraf 10\%) dan publik
(\texttt{0,0510}, taraf 10\%), tetapi lebih besar untuk fungsi pasar.
(Hlm. 430, Tabel 3.)

\textbf{Interpretasi penulis: jenis fungsi.} Fungsi publik lebih
sensitif terhadap efek penghalang dan perlindungan perbatasan daripada
fungsi pasar. Pola ini dinyatakan konsisten dengan H9. (Hlm. 430, Bagian
4.3.)

\textbf{Laporan sumber: robustness check.} Mengganti definisi kota
perbatasan dari waktu tempuh mobil 45 menit menjadi jarak garis lurus
kurang dari 25 km tidak mengubah estimasi secara signifikan (Tabel A2).
Mengganti koefisien peluruhan jarak dari 2 menjadi 1,5 dan 2,5 juga
tidak menghasilkan perubahan signifikan (Tabel A3). Analisis subkelompok
menunjukkan efek peminjaman antara kota perbatasan berukuran serupa
terutama berkaitan dengan kota berukuran menengah, yaitu 100.000-250.000
penduduk, dan pada tingkat lebih rendah kota berpenduduk lebih dari
250.000; manfaat dari kota yang lebih besar di seberang perbatasan
terutama berkaitan dengan kota berukuran menengah (Tabel A4). (Hlm.
430-431, Bagian 4.4; hlm. 436-438, Tabel A2-A4.)

\subsection{Findings}\label{findings-39}

\textbf{Temuan yang dilaporkan sumber.} Dalam sistem domestik Eropa yang
dianalisis, konfigurasi tetangga berukuran serupa berkaitan dengan lebih
banyak fungsi metropolitan, sedangkan kedekatan dengan kota lebih besar
berkaitan dengan lebih sedikit fungsi pada kota yang lebih kecil. Ini
adalah pola dasar \emph{borrowed size} dan \emph{agglomeration shadow}
yang menjadi pembanding bagi analisis lintas batas. (Hlm. 427-428,
Bagian 4.1.)

\textbf{Temuan yang dilaporkan sumber.} Perbatasan tidak bekerja hanya
sebagai penghalang. Pada konteks lintas batas, kota kecil tidak
menunjukkan bayangan dari kota besar di seberang perbatasan dan bahkan
memperoleh manfaat kedekatan; kota berukuran serupa juga menunjukkan
efek \emph{borrowed size} yang lebih kuat daripada pasangan domestik.
(Hlm. 428-429, Bagian 4.2.)

\textbf{Interpretasi penulis.} Penulis menjelaskan keunggulan kota kecil
lintas batas melalui kombinasi diferensial dan komplementaritas lintas
batas, perluasan basis pasar, perbedaan budaya serta institusional, dan
bertahannya kerangka keputusan nasional. Perbatasan menjaga agar proses
rasionalisasi fungsi tidak sepenuhnya mengosongkan kota kecil, sementara
integrasi pasar memberi akses kepada basis dukungan yang lebih luas.
(Hlm. 431-432, Bagian 5.)

\textbf{Temuan yang dilaporkan sumber.} Efek lintas batas tergantung
pada rezim perbatasan. Tidak ada efek interaksi yang signifikan pada
perbatasan terkendali. Pengaruh positif dari kota berukuran serupa dan
lebih besar muncul atau menguat pada perbatasan yang telah terbuka lebih
lama, tetapi durasi pembukaan tidak memberi pengaruh berarti untuk kota
tetangga yang lebih kecil. (Hlm. 429, Bagian 4.2.)

\textbf{Temuan yang dilaporkan sumber.} Efek juga berbeda menurut jenis
fungsi. Manfaat dari kota besar di seberang perbatasan hanya tampak pada
fungsi pasar. Peminjaman dari kota berukuran serupa tampak pada fungsi
pasar dan publik, tetapi lebih kuat pada fungsi pasar; bayangan domestik
dari kota besar juga lebih kuat untuk fungsi pasar. (Hlm. 430, Bagian
4.3.)

\textbf{Interpretasi penulis.} Peran perlindungan perbatasan dinilai
lebih kuat daripada peran penghalangnya untuk kasus kota berukuran
serupa. Penulis merangkum gambaran besarnya sebagai lebih banyak
\emph{borrowed size} seiring waktu, sementara \emph{agglomeration
shadow} lintas batas tetap tidak muncul. (Hlm. 431-432, Bagian 5.)

\textbf{Inferensi pengolahan.} Bukti paling kuat artikel adalah
perbedaan koefisien dan signifikansi antara interaksi domestik serta
lintas batas dalam model SAR. Label ``meminjam ukuran'' dan ``dilindungi
perbatasan'' tetap merupakan interpretasi atas proksi interaksi
potensial dan indeks fungsi, bukan observasi langsung atas arus fungsi
atau mekanisme institusional tertentu. (Hlm. 424-426 dan 428-430.)

\subsection{Conclusions}\label{conclusions-39}

\textbf{Kesimpulan yang dinyatakan penulis.} Perbatasan nasional
memoderasi regularitas normal sistem permukiman Eropa. Dalam pengaturan
domestik, kota yang lebih besar dapat membayangi kota yang lebih kecil,
sedangkan kota berukuran serupa saling memperoleh fungsi. Dalam
pengaturan lintas batas, kota besar tidak memproyeksikan bayangan kepada
kota kecil; kota kecil dapat memperoleh manfaat dari kota besar di
negara tetangga, dan kota berukuran serupa dapat meminjam ukuran lebih
banyak daripada pasangan domestik. (Hlm. 431-432, Bagian 5.)

\textbf{Kesimpulan yang dinyatakan penulis.} Efek tersebut dipengaruhi
oleh durasi pembukaan perbatasan dan jenis fungsi metropolitan.
Keterbukaan yang lebih lama berkaitan dengan lebih banyak fungsi pada
wilayah polisentris lintas batas, sedangkan efek peminjaman positif
lebih besar pada fungsi yang digerakkan pasar daripada fungsi yang
digerakkan publik. (Hlm. 431-432, Bagian 5.)

\textbf{Interpretasi penulis.} Hasil dipandang sebagai indikasi bahwa
integrasi Eropa lebih lanjut, khususnya integrasi metropolitan lintas
batas, bermanfaat bagi kota perbatasan. Para pemangku kepentingan di
kota dan wilayah perbatasan dianjurkan menyadari keunggulan tersebut dan
mengembangkan strategi untuk memanfaatkannya. (Hlm. 432, Bagian 5.)

\textbf{Inferensi pengolahan.} Kesimpulan empiris paling aman dibatasi
pada hubungan statistik dalam 1.920 MUA Eropa, definisi kota perbatasan,
kategori keterbukaan, dan tahun data yang digunakan. Kesimpulan tersebut
tidak membuktikan bahwa pembukaan perbatasan secara kausal menciptakan
fungsi atau bahwa mekanisme yang sama berlaku di luar konteks studi.

\subsection{Contributions}\label{contributions-39}

\textbf{Kontribusi yang diklaim atau ditunjukkan sumber.} Artikel
membawa literatur eksternalitas jaringan kota ke dalam percakapan dengan
literatur efek perbatasan dan secara empiris memasukkan keterikatan
teritorial kota ke dalam analisis pengaruh kota tetangga. (Hlm. 418-419,
Bagian 1.)

\textbf{Kontribusi metode.} Penelitian menggunakan model SAR dan proksi
interaksi potensial berbasis gravitasi untuk membandingkan pengaruh kota
yang lebih kecil, serupa, dan lebih besar. Kerangka tersebut diperluas
dengan pemisahan efek domestik-lintas batas, status perbatasan
terkendali-terbuka, durasi pembukaan, serta fungsi pasar-publik. (Hlm.
424-427, Bagian 3.)

\textbf{Kontribusi empiris.} Dengan 1.920 kota Eropa dan 355 kota
perbatasan, artikel menunjukkan bahwa pola lintas batas tidak sama
dengan pola domestik: bayangan kota besar tidak melintasi perbatasan,
sedangkan hubungan kota berukuran serupa lintas batas dapat menghasilkan
peminjaman ukuran yang lebih kuat. (Hlm. 424-425 dan 428-432.)

\textbf{Kontribusi substantif.} Artikel memperlihatkan bahwa efek
perbatasan bersifat multidimensional dalam arti yang digunakan penulis:
status dan lamanya keterbukaan serta jenis fungsi mengubah arah atau
kekuatan hubungan. Perbedaan antara fungsi pasar dan publik menunjukkan
bahwa manfaat integrasi dan perlindungan perbatasan tidak terbagi sama
pada semua fungsi metropolitan. (Hlm. 423-424 dan 429-432.)

\textbf{Inferensi pengolahan.} Nilai tambah utama artikel adalah desain
perbandingan yang menempatkan efek domestik sebagai pembanding bagi
variasi lintas batas dan menguji beberapa kondisi moderasi dalam satu
analisis. Kontribusi mekanistiknya tetap terbatas karena komponen fisik,
institusional, dan budaya perbatasan tidak diestimasi secara terpisah
serta interaksi aktual tidak tersedia. (Hlm. 422 dan 425-426.)

\subsection{Research Gap}\label{research-gap-39}

\textbf{Kesenjangan yang dinyatakan sumber sebelum penelitian.}
Literatur eksternalitas jaringan kota dan studi pola polisentris sering
mengabaikan lokasi kota terhadap batas negara. Kasus lintas batas,
terutama hubungan regional antara kota kecil dan menengah, jarang
dipertimbangkan. Artikel mengisi kesenjangan tersebut dengan menanyakan
bagaimana perbatasan memoderasi dampak kota tetangga terhadap kinerja
metropolitan. (Hlm. 418-419, Bagian 1.)

\textbf{Kesenjangan yang dinyatakan sumber.} Data interaksi antarkota
pada skala rinci dan lintas negara sulit diperoleh, sehingga penelitian
sebelumnya sering memakai kedekatan atau tumpang tindih wilayah sebagai
asumsi integrasi. Artikel ini mencoba membangun proksi yang lebih
terukur melalui ukuran penduduk dan jarak, sambil membedakan jenis
hubungan dan kondisi perbatasan. (Hlm. 418-419 dan 425-426.)

\textbf{Kesenjangan yang masih terbuka menurut sumber.} Penulis
menyarankan studi kasus wilayah metropolitan lintas batas untuk
mengeksplorasi pola umum secara lebih rinci, penguraian perbatasan
menjadi hambatan yang berbeda, dan penggunaan ukuran kinerja lain di
luar kehadiran fungsi metropolitan. (Hlm. 431-432, Bagian 5.)

\textbf{Inferensi pengolahan: gap yang tersisa.} Karena interaksi
dibentuk dari populasi dan jarak, studi belum menguji arus aktual,
mekanisme institusional spesifik, atau urutan sebab-akibat antara
keterbukaan perbatasan dan kinerja. Ia juga belum menunjukkan apakah
pola tersebut bertahan untuk ukuran kinerja longitudinal, meskipun
artikel mengusulkan produktivitas, pertumbuhan pekerjaan dan penduduk,
serta indikator firma dan rumah tangga sebagai arah lanjutan. (Hlm.
425-426 dan 431-432.)

\subsection{Limitations}\label{limitations-39}

\textbf{Keterbatasan yang dilaporkan sumber.} Data arus komposit
antarkota tidak tersedia pada granularitas 1.920 kota dan cakupan 29
negara. Data pekerja lintas batas juga tidak tersedia di seluruh Eropa
pada tingkat geografis yang diperlukan. Karena itu, penelitian memakai
interaksi potensial, bukan arus perdagangan, komuter, komunikasi,
korporasi, atau kolaborasi ilmiah tertentu. (Hlm. 425-426, Bagian 3.3.)

\textbf{Keterbatasan yang dilaporkan sumber.} Koefisien peluruhan jarak
\texttt{alpha\ =\ 2} dipilih dari studi terdahulu dan diakui sebagai
pilihan yang agak arbitrer. Ambang interaksi 75 km ditetapkan untuk
fokus pada skala regional. Definisi kota perbatasan 45 menit juga
didasarkan pada data perjalanan tahun 2008. Penulis menguji
alternatif-alternatif tersebut melalui Lampiran Tabel A2-A3, tetapi
pengujian ketahanan tidak mengubah fakta bahwa ukuran awalnya adalah
keputusan operasional. (Hlm. 425-426 dan 430-431.)

\textbf{Keterbatasan yang dilaporkan sumber.} Artikel tidak menguraikan
efek perbatasan menjadi hambatan fisik, institusional, dan budaya secara
terpisah karena hal itu dinyatakan berada di luar cakupan penelitian.
Sebagai gantinya, artikel membandingkan status, durasi keterbukaan, dan
jenis fungsi metropolitan. (Hlm. 422 dan 423-424.)

\textbf{Inferensi pengolahan.} Fungsi metropolitan diukur dengan indeks
komposit 28 variabel pada tingkat MUA dan terutama memakai 2006 sebagai
tahun rujukan. Ukuran tersebut tidak menunjukkan volume penggunaan,
kualitas, kapasitas, atau arus manfaat antar kota. Kategori
\emph{borrowed size} karenanya merupakan pembacaan atas hubungan indeks
dengan interaksi potensial, bukan pengukuran langsung atas peminjaman
fungsi.

\textbf{Inferensi pengolahan.} Rancangan observasional ini tidak
menyediakan pembanding kontrafaktual atau identifikasi sebab-akibat.
Kategori durasi keterbukaan membandingkan pasangan perbatasan dalam
kerangka data yang sama, tetapi tidak melacak perubahan fungsi kota
sebelum dan sesudah pembukaan. Variabel yang tidak tercantum dalam model
dan perbedaan historis antarnegara masih dapat ikut menjelaskan hasil.

\textbf{Inferensi pengolahan.} Validitas eksternal dibatasi oleh cakupan
EU 27, Norwegia, dan Swiss, definisi MUA, ambang minimal 20.000
penduduk, serta sistem fungsi metropolitan yang datanya dikumpulkan pada
2003-2009. Penerapan pada kota di luar cakupan geografis atau periode
tersebut: Tidak disebutkan dalam file.

\textbf{Informasi metodologis yang tidak tersedia.} Prosedur pembersihan
data, daftar lengkap pasangan kota, rincian pembentukan semua matriks
bobot, statistik diagnostik model, data mentah, dan kode replikasi:
Tidak disebutkan dalam file.

\subsection{Challenges}\label{challenges-39}

\textbf{Status bagian.} Artikel tidak memiliki bagian khusus bernama
\emph{Challenges}. Tantangan berikut diringkas dari persoalan data dan
pengukuran yang dijelaskan penulis, bukan dari kategori formal artikel.

\textbf{Tantangan yang dilaporkan sumber.} Tantangan utama adalah
mengukur interaksi aktual antara banyak kota di berbagai negara.
Interaksi yang relevan mencakup berbagai arus material dan imaterial,
tetapi data kompositnya tidak tersedia pada tingkat geografis dan
cakupan yang diperlukan. Data komuter lintas batas disebut sebagai salah
satu kemungkinan, namun tidak tersedia untuk seluruh Eropa pada
granularitas tersebut. (Hlm. 425-426, Bagian 3.3.)

\textbf{Tantangan yang dilaporkan sumber.} Penulis juga harus menangani
sifat perbatasan yang tidak seragam. Perbatasan berbeda menurut tingkat
keterbukaan dan lama integrasi, serta mengandung hambatan fisik,
institusional, budaya, dan bahasa yang dapat memengaruhi fungsi
metropolitan secara berbeda. (Hlm. 421-424, Bagian 2.2.)

\textbf{Tantangan yang dilaporkan sumber.} Tantangan operasional lainnya
ialah menetapkan kota perbatasan secara realistis melalui waktu tempuh,
memilih koefisien peluruhan jarak ketika arus aktual tidak tersedia,
membatasi interaksi pada skala regional, dan membedakan fungsi pasar
dari fungsi publik. (Hlm. 425-427, Bagian 3.3-3.5.)

\textbf{Inferensi pengolahan.} Tantangan analitis artikel adalah
membedakan efek penghalang, efek perlindungan, diferensial ekonomi,
komplementaritas, dan kompetisi dari satu proksi interaksi potensial.
Karena komponen perbatasan tidak diestimasi satu per satu, mekanisme
yang menjelaskan perbedaan koefisien tetap tidak terurai secara
langsung. (Hlm. 422-424 dan 425-426.)

\textbf{Inferensi pengolahan.} Replikasi juga menghadapi tantangan
karena keputusan pembentukan indeks 28 variabel, matriks spasial, batas
75 km, kategori ukuran kota, dan pengelompokan durasi perlu
direkonstruksi dari uraian TXT. Rincian prosedur replikasi: Tidak
disebutkan dalam file.

\subsection{Practical Implications}\label{practical-implications-39}

\textbf{Implikasi yang dinyatakan penulis.} Integrasi Eropa lebih
lanjut, terutama integrasi metropolitan lintas batas, dipandang
bermanfaat bagi kota perbatasan. Pemangku kepentingan di kota dan
wilayah perbatasan disarankan lebih menyadari keunggulan unik tersebut
dan mengembangkan strategi untuk memanfaatkannya. (Hlm. 432, Bagian 5.)

\textbf{Implikasi yang dinyatakan penulis.} Kota perbatasan tidak perlu
dipahami hanya sebagai lokasi yang dirugikan oleh kedekatan dengan kota
besar. Perlindungan perbatasan dapat mencegah kota kecil kehilangan
fungsi, sedangkan integrasi pasar lintas batas dapat memperluas basis
pasar. Manfaat ini terutama berkaitan dengan fungsi yang digerakkan
pasar; fungsi publik lebih dipengaruhi hambatan budaya, bahasa,
institusional, dan kedaulatan. (Hlm. 431-432, Bagian 5.)

\textbf{Inferensi pengolahan.} Hasil mendukung penggunaan kebijakan yang
membedakan kota kecil lintas batas, kota berukuran serupa, fungsi pasar,
dan fungsi publik, tetapi TXT tidak menguji instrumen kebijakan
tertentu. Indeks fungsi dan klasifikasi efek sebaiknya diperlakukan
sebagai diagnosis komparatif, bukan sebagai dasar tunggal untuk
menetapkan intervensi atau menyimpulkan dampak pembukaan perbatasan.

\subsection{Applications}\label{applications-39}

\textbf{Aplikasi analitis yang dilakukan sumber.} Kerangka artikel
diterapkan untuk membandingkan pengaruh kota tetangga dalam negara yang
sama dan di seberang perbatasan, menguji variasi menurut status serta
durasi keterbukaan perbatasan, dan membedakan hasil pada fungsi pasar
serta publik. Aplikasi empirisnya mencakup 1.920 MUA Eropa dan 355 kota
perbatasan melalui model SAR. (Hlm. 424-431, Bagian 3-4.)

\textbf{Aplikasi potensial menurut inferensi pengolahan.} Kerangka ini
dapat dipakai sebagai penyaringan awal untuk membandingkan pola
\emph{borrowed size} dan \emph{agglomeration shadow} antarwilayah
perbatasan, selama definisi fungsi, arus interaksi, dan kondisi
perbatasan dikalibrasi ulang. Penggunaan tersebut merupakan inferensi
dari desain penelitian, bukan aplikasi kebijakan yang diuji dalam
artikel.

\textbf{Batas aplikasi.} Implementasi kebijakan, perangkat operasional,
evaluasi program, dan hasil penerapan pada wilayah tertentu: Tidak
disebutkan dalam file.

\subsection{Future Research
(Optional)}\label{future-research-optional-39}

\textbf{Agenda yang dinyatakan penulis.} Penulis mengusulkan:

\begin{itemize}
\tightlist
\item
  studi kasus wilayah metropolitan lintas batas untuk mengeksplorasi
  pola umum artikel secara lebih rinci;
\item
  penguraian perbatasan menjadi hambatan yang berbeda agar efek
  masing-masing dapat dinilai secara terpisah;
\item
  penggunaan ukuran kinerja selain kehadiran fungsi metropolitan,
  termasuk produktivitas, pertumbuhan pekerjaan, dan pertumbuhan
  penduduk;
\item
  pelacakan tren dari waktu ke waktu dengan ukuran kinerja tersebut; dan
\item
  pelengkapan ukuran agregat dengan indikator hasil pada tingkat mikro
  untuk perusahaan dan rumah tangga.
\end{itemize}

Agenda tersebut disebutkan dalam bagian kesimpulan dan diskusi. (Hlm.
431-432, Bagian 5.)

\subsection{Provenance Notes}\label{provenance-notes-39}

\begin{itemize}
\tightlist
\item
  \textbf{Sumber aktual yang diverifikasi:}
  \texttt{source/journal/A24-sohn-licheron-meijers-2022-border-cities-out-of-the-shadow-10.1111-pirs.12653.txt}.
  Tidak ada TXT lain berprefiks \texttt{A24} yang dipakai.
\item
  \textbf{Basis akses:} review ini hanya menggunakan satu file TXT
  tersebut. PDF tidak dibuka terpisah; metadata PDF yang dilaporkan di
  dalam TXT tetap dibaca dan dicatat.
\item
  \textbf{Cakupan baca:} seluruh 1.197 baris tersedia dibaca, termasuk
  blok \texttt{{[}PDF\ Metadata{]}}, teks artikel, abstrak, persamaan,
  Tabel 1-3, Lampiran Tabel A1-A4, daftar pustaka, dan abstrak tambahan
  berbahasa Spanyol serta Jepang.
\item
  \textbf{Metadata ekstraksi:} TXT menyatakan berasal dari
  \texttt{journal/A24-sohn-licheron-meijers-2022-border-cities-out-of-the-shadow-10.1111-pirs.12653.pdf},
  diekstrak pada 31 Agustus 2026 dengan \texttt{pdftotext\ -layout}.
  Metadata PDF menyatakan 23 halaman, tidak terenkripsi, dan tidak
  bertag. (TXT, baris 1-29.)
\item
  \textbf{Dasar identitas bibliografis:} judul dan tiga penulis diambil
  dari kepala artikel, sedangkan jurnal, volume, nomor, halaman, dan DOI
  dikonfirmasi dari format sitasi pada akhir TXT. Metadata PDF hanya
  mencantumkan Christophe Sohn pada field \texttt{Author}, sehingga
  field tersebut tidak dipakai sebagai daftar penulis lengkap. (TXT,
  baris 41-72 dan 1026-1027.)
\item
  \textbf{Locator:} \texttt{Hlm.} merujuk pada nomor halaman tercetak
  jurnal, bukan nomor baris TXT. Bagian, persamaan, tabel, dan lampiran
  dipertahankan jika tersedia.
\item
  \textbf{Keterbatasan ekstraksi:} tata letak dua kolom, simbol
  persamaan, dan beberapa tabel tidak selalu tersusun linear dalam TXT.
  Angka dan makna hasil diperiksa terhadap label tabel serta narasi
  artikel yang berdekatan; tidak ada nilai yang diperkirakan dari sumber
  di luar TXT.
\item
  \textbf{Informasi yang tidak tersedia:} status telaah sejawat yang
  terverifikasi, prosedur pembersihan data, data interaksi aktual, data
  mentah, kode replikasi, dan rincian diagnostik model: Tidak disebutkan
  dalam file.
\item
  \textbf{Penanda epistemik:} \textbf{Laporan sumber} dan \textbf{Temuan
  yang dilaporkan sumber} adalah parafrasa atas isi artikel;
  \textbf{Interpretasi penulis} merekam penjelasan yang diberikan
  penulis; \textbf{Inferensi pengolahan} adalah pembacaan analitis yang
  diturunkan hanya dari teks dan metadata TXT, bukan klaim tambahan dari
  luar sumber.
\item
  \textbf{Batas kerja:} TXT sumber dan file review kode lain tidak
  diubah. Dokumen induk tidak disusun.
\end{itemize}

\section{Review A25 - Borrowed Size of Small and Medium Cities in a
Hierarchical Urban
System}\label{review-a25---borrowed-size-of-small-and-medium-cities-in-a-hierarchical-urban-system}

Catatan pembacaan: - \textbf{Laporan sumber} berarti informasi yang
dinyatakan atau ditampilkan dalam TXT. - \textbf{Interpretasi penulis}
berarti penjelasan yang diberikan Kwon dan Sohn atas hasil atau kerangka
mereka. - \textbf{Inferensi pengolahan} berarti simpulan yang ditarik
dari isi TXT saat menyusun review ini, bukan pernyataan eksplisit
sumber.

\subsection{Summarized Abstract}\label{summarized-abstract-40}

\begin{itemize}
\tightlist
\item
  \textbf{Laporan sumber:} Kota kecil dan menengah dapat memperoleh
  manfaat dari kedekatan dengan kota besar melalui borrowed size dan
  efek spillover. Kedekatan geografis diperkirakan memfasilitasi difusi
  spasial, sedangkan kedekatan hierarkis, yaitu status yang lebih dekat
  dalam sistem urban, dapat menghasilkan manfaat tanpa bergantung
  semata-mata pada ukuran kota. Studi ini menganalisis 152 kota kecil
  dan menengah di Korea dengan delapan indeks dan model regresi. Hasil
  abstrak menyatakan bahwa kedekatan geografis dengan kota primat dan
  posisi hierarkis yang lebih tinggi merupakan pendorong utama borrowed
  size. (Abstract, p.~305.)
\item
  \textbf{Interpretasi penulis:} Kedua jenis efek tersebut diasosiasikan
  masing-masing dengan contagious diffusion dan hierarchical diffusion.
  Studi ditempatkan dalam perdebatan tentang agglomeration economies dan
  network effects dalam pembangunan ekonomi urban. (Abstract, p.~305;
  Introduction, pp.~305-306.)
\end{itemize}

\subsection{Problem Statement}\label{problem-statement-40}

\begin{itemize}
\tightlist
\item
  \textbf{Laporan sumber:} Hubungan geografis dengan kota besar dapat
  menjadi penalti pertumbuhan melalui kompetisi, tetapi sejumlah studi,
  terutama tentang kota-kota Eropa, menunjukkan bahwa kota kecil dapat
  memperoleh manfaat dari kota besar di dekatnya. Dalam sistem urban
  yang hierarkinya kuat, belum terselesaikan apakah sumber utama
  borrowed size berasal dari kota besar dalam functional urban region
  yang sama atau dari kota primat/ibu kota. (Introduction, pp.~305-306.)
\item
  \textbf{Laporan sumber:} Belum ada konsensus apakah kedekatan
  hierarkis pada tingkat nasional atau regional lebih penting, dan
  apakah kedekatan geografis, kedekatan hierarkis, atau konektivitas
  jaringan paling efektif menghasilkan borrowed size. Literatur yang
  dirangkum sumber juga lebih banyak berfokus pada kota Eropa yang lebih
  polisentris, sementara sistem urban Asia Timur dengan hierarki kuat
  kurang diperiksa. (Introduction, pp.~306, 310.)
\item
  \textbf{Inferensi pengolahan:} Masalah penelitian dapat diringkas
  sebagai pengujian sumber dan skala spillover dalam sistem urban yang
  didominasi kota primat, bukan sebagai asumsi bahwa kedekatan dengan
  kota besar selalu menguntungkan. (Inferensi dari Introduction,
  pp.~305-306, 310.)
\end{itemize}

\subsection{Objectives}\label{objectives-40}

\begin{itemize}
\tightlist
\item
  \textbf{Laporan sumber:} Studi mengeksplorasi keberadaan dan dinamika
  borrowed size dalam sistem urban Korea yang sangat hierarkis. Tujuan
  khususnya adalah menguji apakah dua faktor kedekatan, yaitu geografis
  dan hierarkis, memengaruhi borrowed size pada kota kecil dan menengah,
  serta faktor mana yang lebih kuat. (Introduction, p.~306.)
\item
  \textbf{Laporan sumber:} Studi juga membandingkan kedekatan geografis
  dan hierarkis pada skala nasional dan regional, dengan menilai dua
  aspek efek tersebut: borrowed functions dan performance.
  (Introduction, p.~306; Research Methodology, pp.~310-313.)
\end{itemize}

\subsection{Literature Survey}\label{literature-survey-40}

\begin{itemize}
\tightlist
\item
  \textbf{Laporan sumber mengenai literatur:} Agglomeration shadow
  berakar pada New Economic Geography dan menjelaskan mengapa kota yang
  dekat dengan kota lebih besar dapat terhambat oleh efek kompetisi.
  Dalam ringkasan sumber, agglomeration shadow terjadi ketika kota
  perifer memiliki lebih sedikit fungsi daripada yang diharapkan
  berdasarkan ukurannya. (Agglomeration Shadows and Borrowed Sizes on
  the Network of Cities, pp.~306-307.)
\item
  \textbf{Laporan sumber mengenai literatur:} Borrowed size dipaparkan
  sebagai kebalikan dari agglomeration shadow: kota kecil atau
  metropolitan area dapat menunjukkan sebagian karakteristik kota yang
  lebih besar ketika dekat dengan konsentrasi penduduk lain. Borrowed
  size dikaitkan dengan sinergi dan network economy, sehingga kota kecil
  dapat mempertahankan ukuran kecil sambil memperoleh manfaat aglomerasi
  dan menghindari sebagian biaya aglomerasi. (Agglomeration Shadows and
  Borrowed Sizes on the Network of Cities, pp.~306-307.)
\item
  \textbf{Laporan sumber mengenai literatur:} Konsep borrowed size
  kemudian dibedakan menjadi borrowed function dan borrowed performance.
  Jika fungsi dan performance sama-sama muncul, sumber menyebutnya
  borrowed size. Interaksi jaringan pada beberapa skala dan functional
  complementarity dipaparkan sebagai kondisi yang membedakan hubungan
  kolaboratif dari hubungan kompetitif. (Agglomeration Shadows and
  Borrowed Sizes on the Network of Cities, pp.~307-308.)
\item
  \textbf{Laporan sumber mengenai literatur:} Sumber menguraikan tiga
  cara kedekatan hierarkis dapat bekerja: difusi hierarkis fungsi
  budaya, layanan khusus, pengetahuan, dan informasi; jarak geografis
  antarkota yang hierarkisnya berdekatan sehingga mengurangi kompetisi;
  serta institutional borrowed size melalui alokasi sumber daya dan
  layanan publik. Figure 1 menggambarkan hubungan kedekatan geografis,
  kedekatan hierarkis, dan konektivitas jaringan terhadap borrowed size.
  (Agglomeration Shadows and Borrowed Sizes on the Network of Cities,
  pp.~307-308; Figure 1, p.~308.)
\item
  \textbf{Laporan sumber mengenai literatur:} Faktor yang dirangkum
  meliputi jarak, konektivitas jaringan, ukuran kota dan kota mitra,
  complementarity fungsi, serta struktur spasial monosentrik atau
  polisentris. Sumber mencatat hasil yang tidak seragam: borrowed size
  lebih sering dikaitkan dengan wilayah polisentris, tetapi satu studi
  yang diringkas melaporkan bahwa monocentricity dapat mengaktifkan
  borrowed size melalui mobilitas tenaga kerja. (Factors Influencing
  Borrowed Size Effects, pp.~309-310.)
\item
  \textbf{Laporan sumber mengenai literatur:} Studi-studi Eropa, China,
  dan Jepang yang dirangkum memperluas bukti empiris mengenai borrowed
  size, termasuk peran infrastruktur transportasi antardaerah. Namun,
  menurut penulis, sebagian besar studi Asia masih menekankan kedekatan
  geografis dan konektivitas dalam region yang sama, dengan perhatian
  yang lebih kecil terhadap hierarki urban nasional dan regional.
  (Factors Influencing Borrowed Size Effects, pp.~309-310.)
\end{itemize}

\subsection{Method Used}\label{method-used-40}

\begin{itemize}
\tightlist
\item
  \textbf{Laporan sumber:} Studi menggunakan analisis empiris
  kuantitatif dengan delapan indeks fungsi dan performance sebagai
  response variables dalam model regresi untuk 152 kota kecil dan
  menengah di Korea. (Introduction, p.~306; Research Methodology,
  pp.~310-313.)
\item
  \textbf{Laporan sumber:} Model dibangun dalam dua skema klasifikasi
  hierarki: dua kelompok, yaitu kota tingkat kedua dan tingkat ketiga;
  serta empat kelompok, yaitu upper dan lower second-tier serta upper
  dan lower third-tier. Masing-masing skema memiliki model pada skala
  nasional, regional, dan gabungan, yang ditampilkan sebagai Models 1-6
  dan Equations (1)-(6). (Research Methodology, pp.~312-314.)
\item
  \textbf{Laporan sumber:} Pada skala nasional, jarak diukur dari Seoul
  melalui DIST\_SEOUL; pada skala regional, jarak diukur dari pusat
  ekonomi regional melalui DIST\_CENTER. Variabel hierarki nasional dan
  regional dibentuk sebagai dummy. Pemisahan upper dan lower dilakukan
  dengan k-means clustering berdasarkan ukuran populasi. (Research
  Methodology, pp.~312-313.)
\item
  \textbf{Laporan sumber:} Estimasi dilakukan dengan ordinary least
  squares (OLS) menggunakan cluster-robust standard errors berdasarkan
  provinsi untuk menghindari heteroscedasticity dan autocorrelation.
  Koefisien distandardisasi untuk perbandingan antarmodel. (Research
  Methodology, pp.~313-314; Tables 2-7.)
\item
  \textbf{Interpretasi penulis:} Koefisien positif untuk dummy kota yang
  lebih dekat ke tingkat pertama dipakai untuk menilai apakah manfaat
  hierarkis terkonsentrasi pada posisi tertentu. Perbandingan
  signifikansi variabel nasional dan regional dipakai untuk menilai
  skala yang lebih berpengaruh. (Research Methodology, pp.~312-313.)
\end{itemize}

\subsection{Dataset}\label{dataset-40}

\begin{itemize}
\tightlist
\item
  \textbf{Laporan sumber:} Data fungsi dan performance dikumpulkan pada
  skala kota dan county. Semua data dikumpulkan untuk 2022, kecuali GRDP
  yang menggunakan 2020. (Research Methodology, pp.~311-312.)
\item
  \textbf{Laporan sumber:} Korean Statistical Information Service (2022)
  menjadi sumber untuk jumlah karyawan advanced producer services,
  department stores, karyawan sektor creative and art-related services,
  jumlah establishment sektor information technology, dan per capita
  GRDP. (Table 1, p.~312.)
\item
  \textbf{Laporan sumber:} Korean Railway Industry Information Center
  (2022) dan Public Data Portal (2022) menjadi sumber untuk jumlah
  penumpang yang menggunakan stasiun kereta berkecepatan tinggi.
  Intellectual Property Statistics Service (2022) menjadi sumber jumlah
  paten, sedangkan Organisation for Venture Enterprise Certification
  (2022) menjadi sumber jumlah startup. (Table 1, p.~312.)
\item
  \textbf{Laporan sumber:} Variabel posisi hierarkis dan jarak geografis
  dihitung oleh penulis. Variabel tersebut mencakup status
  administratif/hierarkis nasional dan regional, Euclidean distance ke
  Seoul, serta Euclidean distance ke pusat ekonomi regional. (Table 1,
  p.~312.)
\end{itemize}

\subsection{Population / Sample}\label{population-sample-40}

\begin{itemize}
\tightlist
\item
  \textbf{Laporan sumber:} Unit analisis adalah 152 kota kecil dan
  menengah di Korea kontinental. Tujuh kota besar, yaitu Seoul dan enam
  Metropolitan Cities dengan populasi lebih dari satu juta, dipisahkan
  dari 152 kota tersebut. (Research Methodology, pp.~310-311.)
\item
  \textbf{Laporan sumber:} Kota terbesar di tiap provinsi dipakai
  sebagai pusat regional, kecuali Incheon dan Ulsan karena terlalu dekat
  dengan Seoul dan Busan. Karena tidak ada Metropolitan City di Gangwon,
  Chungcheongbuk, dan Jeollabuk, Wonju, Cheongju, dan Jeonju
  diklasifikasikan sebagai pusat regional. Dua kota di Jeju dikeluarkan
  karena studi hanya mencakup Korea kontinental, dan Sejong dikeluarkan
  karena statusnya sebagai special self-governing city dengan
  konsentrasi fungsi pemerintah pusat. (Research Methodology,
  pp.~310-311.)
\item
  \textbf{Laporan sumber:} Data dikumpulkan pada tingkat administratif
  kedua, yaitu kota dan county; district tidak disertakan karena
  merupakan bagian dari Metropolitan Cities, bukan kota yang berdiri
  administratif secara mandiri untuk tujuan studi. (Endnote 1, p.~323.)
\item
  \textbf{Tidak disebutkan dalam file:} Prosedur sampling probabilistik,
  kriteria ukuran sampel selain klasifikasi unit di atas, dan tingkat
  respons tidak dilaporkan. (Research Methodology, pp.~310-311.)
\item
  \textbf{Inferensi pengolahan:} Sampel analitis tampak merupakan
  himpunan unit yang tersisa setelah aturan klasifikasi dan eksklusi
  geografis/administratif, bukan sampel survei; TXT tidak menyebut
  proses sampling survei. (Inferensi dari Research Methodology,
  pp.~310-311; Endnote 1, p.~323.)
\end{itemize}

\subsection{Dependent Variables}\label{dependent-variables-40}

\begin{itemize}
\tightlist
\item
  \textbf{Laporan sumber:} Delapan dependent variables dipakai untuk
  merepresentasikan metropolitan functions dan urban performance.
  BUSINESS adalah jumlah karyawan di sektor advanced producer services;
  CONSUME adalah jumlah department stores; CULTURE adalah jumlah
  karyawan di sektor creative and art-related services; TRANSPORT adalah
  jumlah penumpang yang menggunakan stasiun kereta berkecepatan tinggi.
  (Table 1, p.~312.)
\item
  \textbf{Laporan sumber:} IT adalah jumlah establishment di sektor
  information technology; INCOME adalah per capita GRDP; INNOVATE adalah
  jumlah paten; VITALITY adalah jumlah startup. Delapan indeks tersebut
  digunakan sebagai ukuran untuk menilai keberadaan borrowed size
  effects. (Table 1, p.~312; Conclusion, pp.~322-323.)
\item
  \textbf{Laporan sumber:} Lima kategori metropolitan functions yang
  dipilih adalah business, consumption, cultural amenities,
  transportation, dan information technology. Urban performance
  diringkas dalam tiga kategori, yaitu income, innovation, dan vitality.
  (Research Methodology, pp.~310-312.)
\end{itemize}

\subsection{Results}\label{results-40}

\begin{itemize}
\tightlist
\item
  \textbf{Laporan sumber:} Pada Model 1, yang menguji skala nasional
  dengan SECOND dan DIST\_SEOUL, koefisien SECOND positif dan signifikan
  untuk tujuh dari delapan indeks; pengecualiannya adalah INCOME.
  DIST\_SEOUL memiliki koefisien negatif dan signifikan untuk BUSINESS,
  CULTURE, IT, INNOVATE, dan VITALITY, sehingga lima indeks menunjukkan
  distance decay. (Table 2, p.~316; Estimated Impacts of Urban Hierarchy
  and Geography on Borrowed Size Effects, pp.~314-316.)
\item
  \textbf{Laporan sumber:} Pada Model 2, REG\_SECOND positif dan
  signifikan untuk enam indeks, yaitu BUSINESS, CONSUME, CULTURE,
  TRANSPORT, INNOVATE, dan VITALITY; variabel ini tidak signifikan untuk
  INCOME dan IT. DIST\_CENTER tidak signifikan pada kedelapan indeks.
  (Table 3, p.~316.)
\item
  \textbf{Laporan sumber:} Pada Model 3, ketika hierarki nasional,
  hierarki regional, DIST\_SEOUL, dan DIST\_CENTER dimasukkan bersama,
  signifikansi SECOND melemah dan hanya tetap signifikan untuk
  TRANSPORT. REG\_SECOND signifikan untuk BUSINESS, CONSUME, TRANSPORT,
  dan VITALITY. DIST\_SEOUL tetap negatif dan signifikan untuk lima
  indeks, sedangkan DIST\_CENTER tetap tidak signifikan. (Table 4,
  p.~317; Discussion, pp.~319-320.)
\item
  \textbf{Laporan sumber:} Pada Model 4, klasifikasi nasional empat
  tingkat menunjukkan bahwa UP\_SECOND positif dan signifikan pada tujuh
  indeks selain INCOME. LOW\_SECOND dan UP\_THIRD menunjukkan pengaruh
  yang terbatas dan bercampur; DIST\_SEOUL mempertahankan pola distance
  decay pada lima indeks. (Table 5, p.~318; Estimated Impacts\ldots,
  pp.~317-318.)
\item
  \textbf{Laporan sumber:} Pada Model 5, REG\_UP\_SECOND berpengaruh
  positif pada lima indeks, sedangkan REG\_LOW\_SECOND dan
  REG\_UP\_THIRD menunjukkan pengaruh yang lebih terbatas. DIST\_CENTER
  tidak menunjukkan signifikansi yang berarti. Penulis menyatakan bahwa
  pada skala regional hierarki lebih berperan daripada kedekatan
  geografis dengan pusat regional. (Table 6, p.~318; Estimated
  Impacts\ldots, pp.~318-319.)
\item
  \textbf{Laporan sumber:} Pada Model 6, yang menggabungkan seluruh
  variabel hierarki dan geografis, UP\_SECOND nasional memiliki pengaruh
  positif dan signifikan terutama pada CULTURE, TRANSPORT, IT, dan
  VITALITY. Variabel hierarki regional pada umumnya kehilangan
  signifikansi ketika variabel nasional dimasukkan, sedangkan
  DIST\_CENTER tetap tidak berarti. Tabel 7 juga menampilkan koefisien
  REG\_UP\_SECOND yang positif dan signifikan untuk INCOME serta
  koefisien DIST\_SEOUL yang negatif dan signifikan untuk CULTURE, IT,
  INCOME, dan VITALITY. (Table 7, p.~319; Estimated Impacts\ldots,
  pp.~319-320.)
\item
  \textbf{Laporan sumber:} Figure 4 merangkum jumlah signifikansi di
  antara delapan dependent variables pada model gabungan. Dalam
  perbandingan dua tingkat, REG\_SECOND tercatat empat kali dan SECOND
  nasional satu kali; DIST\_SEOUL tercatat lima kali dan DIST\_CENTER
  nol kali. Dalam perbandingan empat tingkat, UP\_SECOND nasional dan
  DIST\_SEOUL masing-masing tercatat empat kali, REG\_UP\_SECOND satu
  kali, dan DIST\_CENTER nol kali. (Figure 4, p.~320; catatan Figure 4
  menyatakan bahwa perbandingan berasal dari Models 3 dan 6, Tables 4
  dan 7.)
\item
  \textbf{Inferensi pengolahan:} Bukti empiris dalam TXT disajikan
  sebagai pola asosiasi OLS antara posisi/jarak dan delapan indikator,
  bukan sebagai hasil eksperimen atau estimasi yang memisahkan mekanisme
  kausal secara langsung. (Inferensi dari Research Methodology,
  pp.~312-314, dan Tables 2-7.)
\end{itemize}

\subsection{Findings}\label{findings-40}

\begin{itemize}
\tightlist
\item
  \textbf{Interpretasi penulis:} Posisi dalam hierarki penting, tetapi
  hasilnya bergantung pada skema klasifikasi. Dalam klasifikasi dua
  tingkat, kota tingkat kedua nasional dan regional cenderung memiliki
  fungsi serta performance urban yang lebih tinggi; dalam klasifikasi
  empat tingkat, manfaat nasional yang bermakna terutama terkonsentrasi
  pada upper second-tier cities. Pada skala regional, kota dengan
  tingkat lebih rendah masih menunjukkan sebagian efek positif, tetapi
  lebih lemah. (Discussion, pp.~319-321.)
\item
  \textbf{Interpretasi penulis:} Kedekatan dengan Seoul merupakan faktor
  geografis yang konsisten. Jarak ke Seoul berasosiasi negatif dengan
  borrowed size effects, sedangkan jarak ke pusat regional sebagian
  besar tidak signifikan. Penulis mengaitkan asimetri ini dengan
  dominasi Seoul dalam ukuran, pengaruh ekonomi, politik, dan budaya,
  serta jaringan transportasi Korea yang berorientasi ke ibu kota.
  (Discussion, pp.~320-321.)
\item
  \textbf{Interpretasi penulis:} Dalam model dua tingkat, kedekatan
  spasial dengan Seoul lebih berpengaruh daripada hierarki regional;
  dalam model empat tingkat, pengaruh upper second-tier nasional dan
  jarak ke Seoul menjadi sebanding. Klasifikasi yang lebih rinci
  dipandang membuat peran hierarki lebih terlihat. (Discussion,
  pp.~320-321; Figure 4, p.~320.)
\item
  \textbf{Interpretasi penulis:} Secara keseluruhan, kota kecil dan
  menengah di Korea dapat memperoleh borrowed size terutama melalui
  salah satu dari dua kondisi: berada pada posisi teratas di antara kota
  kecil-menengah dalam sistem urban atau berdekatan secara geografis
  dengan Seoul. (Discussion, pp.~321-322; Conclusion, pp.~322-323.)
\item
  \textbf{Inferensi pengolahan:} Karena pola signifikansi berubah antara
  klasifikasi dua dan empat tingkat, kesimpulan tentang faktor yang
  paling berpengaruh sensitif terhadap cara posisi hierarkis
  dioperasionalkan. (Inferensi dari Tables 4, 6, dan 7 serta Figure 4,
  p.~320.)
\end{itemize}

\subsection{Conclusions}\label{conclusions-40}

\begin{itemize}
\tightlist
\item
  \textbf{Laporan sumber:} Studi menyimpulkan bahwa posisi hierarkis
  berpengaruh, tetapi kekuatannya berubah menurut klasifikasi. Pada
  klasifikasi dua tingkat, kota tingkat kedua nasional dan regional
  memperoleh borrowed size, dengan hierarki regional lebih berpengaruh;
  pada klasifikasi empat tingkat, hanya upper second-tier cities pada
  tingkat nasional yang menunjukkan efek positif yang berarti.
  (Conclusion, pp.~322-323.)
\item
  \textbf{Laporan sumber:} Jarak ke Seoul menunjukkan asosiasi negatif
  yang konsisten dengan borrowed size effects, sedangkan jarak ke pusat
  regional tidak memiliki pengaruh berarti. Ketika hierarki dan jarak
  dibandingkan, kedekatan dengan Seoul paling berpengaruh dalam
  klasifikasi sederhana; dengan klasifikasi lebih rinci, pengaruh
  hierarki urban nasional menjadi sebanding dengan kedekatan geografis.
  (Conclusion, pp.~322-323.)
\item
  \textbf{Interpretasi penulis:} Dominasi Seoul dan struktur hierarkis
  yang kaku membuat syarat borrowed size sulit dicapai oleh banyak kota
  kecil dan menengah. Perbaikan konektivitas transportasi untuk
  mengurangi jarak relatif dinilai sebagai pilihan praktis, tetapi harus
  disertai complementarity fungsi agar aksesibilitas tidak memperkuat
  pemusatan fungsi di Seoul dan agglomeration shadow. (Conclusion,
  p.~323.)
\end{itemize}

\subsection{Contributions}\label{contributions-40}

\begin{itemize}
\tightlist
\item
  \textbf{Laporan sumber:} Penulis menyatakan bahwa studi berkontribusi
  pada diskusi lebih luas tentang agglomeration dan network effects
  dengan memperjelas kondisi struktural ketika borrowed size muncul.
  (Conclusion, p.~323.)
\item
  \textbf{Laporan sumber:} Studi memperluas pembahasan yang sebelumnya
  banyak berfokus pada kota Eropa dengan menguji sistem urban Korea yang
  hierarkinya kuat, membandingkan kedekatan geografis dan hierarkis pada
  skala nasional serta regional, dan menggabungkan delapan indikator
  fungsi serta performance. (Introduction, pp.~306, 310; Research
  Methodology, pp.~310-313.)
\end{itemize}

\subsection{Research Gap}\label{research-gap-40}

\begin{itemize}
\tightlist
\item
  \textbf{Laporan sumber:} Sebelum studi ini, belum jelas apakah
  borrowed size dalam sistem hierarkis terutama berasal dari kota besar
  dalam functional urban region atau dari kota primat, apakah kedekatan
  hierarkis nasional atau regional lebih penting, dan apakah kedekatan
  geografis, kedekatan hierarkis, atau konektivitas jaringan merupakan
  faktor yang paling efektif. (Introduction, pp.~305-306, 310.)
\item
  \textbf{Laporan sumber:} Penulis menilai studi sebelumnya, termasuk
  studi Asia, lebih sering membahas kedekatan geografis dan konektivitas
  antarkota dalam region yang sama daripada peran hierarki urban
  nasional dan regional dalam sistem yang hierarkinya kuat. (Factors
  Influencing Borrowed Size Effects, pp.~309-310.)
\item
  \textbf{Laporan sumber:} Gap yang masih tersisa adalah pemisahan
  proses individual kedekatan hierarkis dan geografis, perubahan efek
  dari waktu ke waktu, kapasitas tiap kota untuk memanfaatkan efek
  tersebut, dan pengaruh jenis koneksi jaringan yang berbeda.
  (Conclusion, p.~323.)
\item
  \textbf{Inferensi pengolahan:} TXT tidak melaporkan pengujian lintas
  negara selain rujukan literatur yang dibahas, sehingga bukti empiris
  langsung dalam studi ini tetap berpusat pada sistem urban Korea.
  (Inferensi dari Abstract, p.~305; Research Methodology, pp.~310-313.)
\end{itemize}

\subsection{Limitations}\label{limitations-40}

\begin{itemize}
\tightlist
\item
  \textbf{Laporan sumber:} Penulis tidak dapat membedakan proses
  individual kedekatan hierarkis dan kedekatan geografis sebagai
  variabel penjelas yang terpisah; karena itu, studi tidak dapat
  menentukan secara definitif mana yang paling kuat memengaruhi borrowed
  size. (Conclusion, p.~323.)
\item
  \textbf{Laporan sumber:} Model tidak dapat memisahkan pengaruh
  institusional yang terkait dengan kedekatan hierarkis. Hasilnya dapat
  menangkap gabungan kecenderungan pemerintah pusat untuk
  memprioritaskan kota tingkat lebih tinggi dan kebijakan pembangunan
  berimbang. (Discussion, pp.~320-321.)
\item
  \textbf{Inferensi pengolahan:} Rentang waktu tidak dianalisis sebagai
  dinamika temporal dalam model yang dilaporkan, karena data fungsi dan
  performance terutama merujuk pada 2022 dan GRDP pada 2020; hal ini
  konsisten dengan future research penulis tentang perubahan temporal.
  (Inferensi dari Research Methodology, p.~311; Conclusion, p.~323.)
\item
  \textbf{Tidak disebutkan dalam file:} Keterbatasan lain di luar
  batasan identifikasi mekanisme, faktor institusional, dan cakupan
  waktu yang dinyatakan atau dapat diturunkan secara langsung dari TXT.
\end{itemize}

\subsection{Challenges}\label{challenges-40}

\begin{itemize}
\tightlist
\item
  \textbf{Laporan sumber:} Tantangan struktural utama adalah dominasi
  Seoul dan hierarki urban Korea yang kaku. Kota kecil dan menengah yang
  jauh dari Seoul menghadapi kesulitan memperoleh borrowed size;
  perpindahan ke tingkat hierarki yang lebih tinggi juga sulit kecuali
  kota telah memiliki fungsi ekonomi atau administratif utama.
  (Discussion, pp.~320-321; Conclusion, p.~323.)
\item
  \textbf{Laporan sumber:} Aksesibilitas yang lebih baik tidak otomatis
  menghasilkan manfaat. Tanpa complementarity fungsi, peningkatan akses
  dapat memusatkan fungsi di Seoul dan memperluas agglomeration shadow.
  Sumber juga menyatakan bahwa pemerintah lokal masih sangat bergantung
  pada proyek dan anggaran pemerintah nasional serta belum cukup
  mengembangkan complementarity dengan kota besar. (Discussion,
  pp.~321-322.)
\item
  \textbf{Laporan sumber:} Penurunan populasi, urban shrinkage, dan
  tingkat fertilitas Korea yang rendah membuat strategi sekadar menambah
  penduduk perlu dipertimbangkan ulang; bahkan penyediaan layanan dasar
  di kota yang menyusut dapat menjadi tidak berkelanjutan tanpa
  intervensi kebijakan. (Discussion, p.~322.)
\end{itemize}

\subsection{Practical Implications}\label{practical-implications-40}

\begin{itemize}
\tightlist
\item
  \textbf{Laporan sumber:} Pemerintah dapat mengurangi jarak relatif
  melalui peningkatan konektivitas transportasi, karena lokasi absolut
  kota tidak dapat diubah tetapi hubungan fungsionalnya dapat
  diperbaiki. (Discussion, pp.~321-322; Conclusion, p.~323.)
\item
  \textbf{Laporan sumber:} Peningkatan konektivitas perlu disertai
  kebijakan yang memperkuat functional complementarity, spesialisasi
  industri regional, dan interaksi antarkota agar manfaat aglomerasi
  tidak berubah menjadi pemusatan fungsi di Seoul. (Discussion,
  pp.~321-322.)
\item
  \textbf{Laporan sumber:} Dalam kebijakan megaregion Korea, pemerintah
  lokal didorong membentuk collaborative zones yang menghubungkan kota
  besar dan kota kecil di sekitarnya serta menyusun joint regional
  master plans. (Discussion, p.~322.)
\item
  \textbf{Laporan sumber:} Untuk kota kecil dan menengah yang jauh dari
  Seoul, alternatif mempertahankan urban vitality yang disebutkan adalah
  menarik imigran dan memanfaatkan populasi bergerak, termasuk
  commuters, part-time residents, business travellers, dan workation
  participants. (Discussion, p.~322.)
\end{itemize}

\subsection{Applications}\label{applications-40}

\begin{itemize}
\tightlist
\item
  \textbf{Laporan sumber:} Kerangka delapan indeks, klasifikasi posisi
  hierarkis, dan ukuran jarak ke Seoul serta pusat regional diterapkan
  untuk menilai fungsi dan performance kota kecil-menengah dalam sistem
  urban Korea. (Research Methodology, pp.~310-314; Tables 1-7.)
\item
  \textbf{Interpretasi penulis:} Hasil dapat dipakai untuk mengarahkan
  kebijakan konektivitas, complementarity fungsi, spesialisasi regional,
  dan perencanaan lintas kota, bukan hanya untuk mengidentifikasi
  kekurangan fungsi di kota yang lebih kecil. (Discussion, pp.~321-322.)
\item
  \textbf{Inferensi pengolahan:} Dalam batas bukti TXT, aplikasi yang
  paling jelas adalah diagnosis komparatif pada skala nasional dan
  regional; TXT tidak melaporkan implementasi atau evaluasi kebijakan
  tertentu setelah rekomendasi tersebut. (Inferensi dari Research
  Methodology, pp.~310-314, dan Discussion, pp.~321-322.)
\end{itemize}

\subsection{Future Research
(Optional)}\label{future-research-optional-40}

\begin{itemize}
\tightlist
\item
  \textbf{Laporan sumber:} Penulis menyarankan penelitian lanjutan
  tentang perubahan efek dari waktu ke waktu, kapasitas masing-masing
  kota untuk memanfaatkan borrowed size, serta bagaimana jenis koneksi
  jaringan yang berbeda, seperti commuting, logistics, dan
  accessibility, membentuk geografi borrowed size effects di Korea dan
  di luar Korea. (Conclusion, p.~323.)
\end{itemize}

\subsection{Provenance Notes}\label{provenance-notes-40}

\begin{itemize}
\tightlist
\item
  \textbf{Sumber utama:}
  \texttt{source/journal/A25-kwon-sohn-2026-borrowed-size-of-small-and-medium-cities-in-a-hierarchical-urban-system-10.1111-tesg.70061.txt}.
\item
  \textbf{Verifikasi filename:} Pencarian rekursif
  \texttt{source/**/A25*.txt} menemukan tepat satu TXT aktual, yaitu
  path di atas. Tidak ada TXT A25 kedua yang dipakai.
\item
  \textbf{Batas pembacaan:} Seluruh TXT dibaca, termasuk
  \texttt{{[}PDF\ Metadata{]}} pada baris 1-26 dan
  \texttt{{[}Extracted\ Text{]}} pada baris 27-1.173. TXT menyatakan
  sumber PDF
  \texttt{journal/A25-kwon-sohn-2026-borrowed-size-of-small-and-medium-cities-in-a-hierarchical-urban-system-10.1111-tesg.70061.pdf},
  diekstrak pada 2026-08-31 dengan \texttt{pdftotext\ -layout}, dan
  memiliki 21 halaman. (TXT, baris 1-26.)
\item
  \textbf{Identitas yang tersedia dalam file:} Judul, penulis Kyusang
  Kwon dan Jungyul Sohn, jurnal Tijdschrift voor Economische en Sociale
  Geografie, 2026, volume 117 nomor 2, halaman 305-325, dan DOI
  10.1111/tesg.70061 tercantum di bagian awal/footer teks ekstraksi.
  (TXT, baris 28-37, 77-81.)
\item
  \textbf{Aturan locator:} Nomor halaman dalam review mengikuti
  pagination artikel yang tercetak di TXT, bukan nomor baris TXT.
  Locator bagian, tabel, gambar, dan endnote dipertahankan ketika
  tersedia.
\item
  \textbf{Basis evidence:} Review hanya memakai isi TXT A25. PDF asli,
  data mentah, kode estimasi, dan karya dalam daftar referensi tidak
  dibuka untuk menambah atau memverifikasi substansi.
\item
  \textbf{Batasan ekstraksi:} TXT memuat form-feed, header/footer
  penerbit, pemenggalan kata dengan hyphen, dan tata letak tabel yang
  terpecah ke beberapa baris. Review tidak mengubah TXT dan tidak
  menganggap artefak pemformatan sebagai isi substantif.
\item
  \textbf{Ketidaksesuaian internal yang dipertahankan:} Prosa pada
  bagian pembahasan menyatakan bahwa variabel hierarki regional
  kehilangan signifikansi dalam Model 6, tetapi Table 7 menampilkan
  REG\_UP\_SECOND untuk INCOME sebesar 0.229 dengan \texttt{***}. Prosa
  yang sama menyebut pengaruh DIST\_SEOUL sebagai ``significant and
  positive'', sedangkan koefisien signifikan DIST\_SEOUL yang
  ditampilkan Table 7 untuk CULTURE, IT, INCOME, dan VITALITY bernilai
  negatif. Review melaporkan kedua lapisan bukti ini tanpa memilih atau
  memperbaiki salah satunya. (Table 7, p.~319; Estimated Impacts\ldots,
  pp.~319-320.)
\item
  \textbf{Nuansa generalisasi:} Kalimat kesimpulan tentang asosiasi
  negatif DIST\_SEOUL yang ``konsisten'' dibaca bersama tabel rinci;
  beberapa indeks pada Tables 2-4 tidak signifikan meskipun pola negatif
  atau distance decay dilaporkan untuk sejumlah indeks. (Tables 2-4,
  pp.~316-317; Conclusion, pp.~322-323.)
\item
  \textbf{Audit review:} File ini memuat seluruh 19 heading wajib, tidak
  menggunakan tabel, memberi label laporan sumber/interpretasi
  penulis/inferensi pengolahan, dan menempatkan locator pada klaim
  substantif utama. Audit tidak mengubah atau membuat file review lain.
\end{itemize}

\section{Review A26}\label{review-a26}

\subsection{Summarized Abstract}\label{summarized-abstract-41}

\textbf{{[}Dilaporkan sumber{]}} Otsuka meninjau kembali konsep
\emph{borrowed size} untuk menjelaskan pertumbuhan wilayah yang
melampaui manfaat aglomerasi tradisional. Artikel mengusulkan tiga
lapisan: eksternalitas terkait skala, konektivitas jaringan sebagai
mediator, serta fondasi institusional dan relasional. Dengan data
prefektur Jepang, artikel mengaitkan aksesibilitas yang membaik,
spesialisasi industri, perencanaan spasial nasional, dan transfer fiskal
dengan kemungkinan difusinya manfaat aglomerasi dari kawasan inti ke
wilayah periferal (Abstrak, hlm. 975).

\textbf{{[}Interpretasi penulis{]}} Dalam konteks Jepang yang ditandai
sentralisasi pemerintahan, konsentrasi metropolitan, dan penuaan
penduduk, \emph{borrowed size} diposisikan sebagai proses aktif yang
dimediasi institusi, bukan limpahan pasif. Artikel menyatakan bahwa
kerangka ini menantang dikotomi aglomerasi-dispersi dan memiliki
pelajaran yang dapat ditransfer untuk ekonomi menua yang
terdesentralisasi (Abstrak, hlm. 975).

\textbf{{[}Inferensi pengolahan{]}} Bukti empiris yang disajikan
bersifat deskriptif dan digunakan untuk mengilustrasikan kecocokan
struktural dengan model; bukti tersebut tidak memvalidasi hubungan
sebab-akibat formal karena artikel secara eksplisit tidak bertujuan
melakukan inferensi kausal (Section 1, hlm. 976; Section 4, hlm. 983).

\subsection{Problem Statement}\label{problem-statement-41}

\textbf{{[}Dilaporkan sumber{]}} Konsentrasi aktivitas ekonomi di
Greater Tokyo tetap kuat, sementara penurunan penduduk dan pengosongan
industri menjadi masalah akut bagi wilayah periferal. Perubahan menuju
telework setelah pandemi COVID-19 membuka kembali pertanyaan tentang
kelayakan hunian di wilayah periferal, tetapi tidak menghapus persoalan
struktural tersebut (Section 1, hlm. 975).

\textbf{{[}Dilaporkan sumber{]}} Pertanyaan penelitian utamanya adalah
bagaimana wilayah urban yang lebih kecil atau periferal dapat mengakses
manfaat ekonomi skala tanpa mengalami aglomerasi fisik, serta mekanisme
institusional dan relasional apa yang memungkinkan akses itu (Section 1,
hlm. 976).

\textbf{{[}Interpretasi penulis{]}} Kerangka aglomerasi yang menyamakan
kapasitas fungsional dengan ukuran, kepadatan, dan kedekatan fisik
dinilai tidak cukup untuk menjelaskan koeksistensi konsentrasi
metropolitan dengan penyempitan kesenjangan GVA per kapita antardaerah
di Jepang (Section 1, hlm. 976; Section 4.1, hlm. 984-985).

\subsection{Objectives}\label{objectives-41}

\textbf{{[}Dilaporkan sumber{]}} Tujuan yang dinyatakan atau dijelaskan
dalam artikel adalah:

\begin{itemize}
\tightlist
\item
  merekonstruksi konsep \emph{borrowed size} sebagai struktur tiga
  lapisan yang menghubungkan eksternalitas skala, konektivitas jaringan,
  dan fondasi institusional-relasional (Sections 1 dan 3, hlm. 976,
  979-983);
\item
  menjelaskan \emph{borrowed size} sekaligus sebagai keadaan, mekanisme,
  dan proses yang saling bertingkat (Section 1, hlm. 976; Section 3,
  hlm. 979-980);
\item
  memberi landasan empiris deskriptif melalui struktur ekonomi prefektur
  Jepang, khususnya disparitas GVA riil, kedekatan spasial terhadap
  Greater Tokyo, dan konfigurasi atau spesialisasi industri (Section 1,
  hlm. 976-977; Section 4, hlm. 983);
\item
  menilai relevansi kerangka tersebut untuk pertumbuhan wilayah yang
  terdistribusi dan merumuskan implikasi kebijakan berbasis jaringan
  serta desain institusional (Section 1, hlm. 977; Section 5, hlm.
  990-992).
\end{itemize}

\textbf{{[}Dilaporkan sumber{]}} Artikel menyatakan bahwa komponen
empirisnya bukan untuk membangun inferensi kausal, melainkan melengkapi
model teoretis dengan keragaman struktural pertumbuhan wilayah (Section
1, hlm. 976-977).

\subsection{Literature Survey}\label{literature-survey-41}

\textbf{{[}Dilaporkan sumber{]}} Artikel menelusuri asal konsep pada
Alonso (1973), rekonstruksi Meijers dan Burger (2010) dalam teori sistem
perkotaan, polisentrik, dan kawasan urban fungsional, serta penerapan
empiris oleh Phelps et al.~(2001) dan Meijers et al.~(2016). Parr
(2007), Burger et al.~(2015), dan Meijers dan Burger (2017) digunakan
untuk memperluas pembahasan menuju jangkauan fungsional, pasar
perumahan, jaringan transportasi, struktur industri, dan peran kota
kecil-menengah (Section 2, hlm. 977-978).

\textbf{{[}Dilaporkan sumber{]}} Untuk dimensi institusional, artikel
mengacu pada Rodríguez-Pose (2013), Rodríguez-Pose dan Storper (2006),
Storper dan Scott (2016), North (1990), Rodrik (2007), serta Aoki
(2001). Untuk sisi negatif eksternalitas, artikel membahas
\emph{agglomeration shadow} melalui Krugman (1991), Fujita et
al.~(1999), Tabuchi (1998), dan Brakman et al.~(2004). Camagni et
al.~(2015) dan Boschma (2005) dipakai sebagai landasan yang kemudian
diperluas ke mekanisme tata kelola teritorial (Sections 2-3, hlm.
977-982).

\textbf{{[}Interpretasi penulis{]}} Kesenjangan teoretis utama terletak
pada kurangnya penjelasan tentang fondasi institusional dan relasional
yang memediasi \emph{borrowed size}: mengapa sebagian wilayah mampu
memanfaatkannya, bagaimana struktur mediator muncul dan bertahan, serta
bagaimana \emph{borrowed size} dapat muncul, berubah, atau menghilang.
Kedekatan spasial saja tidak dianggap menjamin akses terhadap
eksternalitas (Section 2, hlm. 977-978).

\textbf{{[}Inferensi pengolahan{]}} Section 2 merupakan tinjauan
literatur naratif dalam artikel. Protokol penelusuran, kriteria inklusi,
jumlah korpus, dan prosedur \emph{systematic review} tidak disebutkan
dalam file.

\subsection{Method Used}\label{method-used-41}

\textbf{{[}Dilaporkan sumber{]}} Metode artikel terdiri atas
rekonstruksi teoretis dan analisis empiris deskriptif pada tingkat
prefektur Jepang. Kerangka konseptual menyusun tiga lapisan: sumber
eksternalitas skala, struktur mediasi jaringan, dan basis penerimaan
berupa institusi serta relasi (Sections 3 dan 3.4, hlm. 979-983).

\textbf{{[}Dilaporkan sumber{]}} Analisis empiris membandingkan atau
menggambarkan: distribusi GVA riil dan penduduk; perubahan koefisien
variasi GVA per kapita; hubungan GVA per kapita awal dengan pertumbuhan
berikutnya; waktu tempuh ke Tokyo dan pertumbuhan produktivitas tenaga
kerja; pengurangan waktu tempuh dan pertumbuhan produktivitas;
spesialisasi manufaktur dan pertumbuhan produktivitas; serta koefisien
variasi sebelum dan sesudah transfer fiskal (Tables 1-2 dan Figures 3-8,
hlm. 983-990).

\textbf{{[}Dilaporkan sumber{]}} Prefektur digunakan sebagai unit
administratif dan statistik karena ketersediaan statistik jangka panjang
yang konsisten serta perannya dalam tata kelola ekonomi dan perencanaan
infrastruktur. Secara teoretis, cakupan \emph{borrowed size} tetap
melampaui batas administratif untuk mencakup kawasan ekonomi fungsional
dan jaringan arus manusia, barang, serta informasi (Section 1, hlm.
976-977).

\textbf{{[}Dilaporkan sumber{]}} Artikel menegaskan bahwa analisisnya
tidak dimaksudkan untuk menetapkan hubungan kausal atau melakukan
validasi formal, melainkan menunjukkan karakteristik struktural yang
relevan bagi model (Section 1, hlm. 976-977; Section 4, hlm. 983).

\subsection{Dataset}\label{dataset-41}

\textbf{{[}Dilaporkan sumber{]}} Data utama berasal dari
\emph{Prefectural-level Economic Database} milik Cabinet Office,
Government of Japan. Data tersebut digunakan untuk GVA riil, penduduk,
GVA per kapita, produktivitas tenaga kerja, dan spesialisasi manufaktur
pada tabel serta gambar empiris (Tables 1-2 dan Figures 3-7, hlm.
983-990).

\textbf{{[}Dilaporkan sumber{]}} Data waktu tempuh atau jarak temporal
ke Tokyo bersumber dari \emph{National Integrated Transport Analysis
System} (NITAS) milik Ministry of Land, Infrastructure, Transport and
Tourism (Figure 5-6, hlm. 985-987). Perhitungan Figure 8 menggunakan
\emph{Annual Report on Prefectural Accounts} dari Cabinet Office,
\emph{Basic Resident Registers} dari Ministry of Internal Affairs and
Communications, dan \emph{Nikkei Economic Electronic Databank System}
(Figure 8, hlm. 989).

\textbf{{[}Dilaporkan sumber{]}} Cakupan waktu yang tampak dalam
analisis adalah 1990-2019 untuk tren dan hubungan pertumbuhan, 2019
untuk distribusi utama, serta pembanding 2014 pada Table 2. Data
pendukung disebut berasal dari sumber publik Cabinet Office dan
Statistics Bureau of Japan; dataset olahan dan kode analisis tersedia
dari penulis atas permintaan yang wajar (Data availability statement,
hlm. 994).

\textbf{{[}Inferensi pengolahan{]}} Jumlah observasi, definisi
operasional lengkap indeks spesialisasi manufaktur, formula koefisien
variasi, dan spesifikasi statistik: \textbf{Tidak disebutkan dalam
file}.

\subsection{Population / Sample}\label{population-sample-41}

\textbf{{[}Dilaporkan sumber{]}} Populasi atau unit analisis adalah
prefektur-prefektur Jepang (Section 1, hlm. 976-977; Section 4, hlm.
983).

\textbf{{[}Inferensi pengolahan{]}} Prosedur pengambilan sampel, jumlah
prefektur yang masuk dalam setiap perhitungan, dan kriteria pengecualian
tidak disebutkan dalam file; karena itu ukuran sampel adalah
\textbf{Tidak disebutkan dalam file}.

\textbf{{[}Dilaporkan sumber{]}} Istilah populasi pada Table 1 dan Table
2 merujuk pada jumlah penduduk sebagai indikator agregat wilayah, bukan
responden survei. Artikel menyatakan bahwa penelitian tidak melibatkan
manusia sebagai partisipan atau data personal teridentifikasi (Research
ethics and consent, hlm. 994).

\subsection{Dependent Variables}\label{dependent-variables-41}

\textbf{{[}Inferensi pengolahan{]}} Artikel tidak menetapkan variabel
dependen formal dalam model kausal atau regresi. Untuk desain deskriptif
tanpa inferensi kausal, variabel dependen formal adalah \textbf{Tidak
berlaku}.

\textbf{{[}Dilaporkan sumber{]}} Beberapa indikator hasil yang dipetakan
atau dibandingkan adalah:

\begin{itemize}
\tightlist
\item
  GVA riil dan GVA per kapita, termasuk perubahan koefisien variasinya,
  sebagai indikator kinerja dan disparitas regional (Table 1 dan Figure
  3, hlm. 983-985);
\item
  pertumbuhan GVA per kapita setelah tingkat awal 1990, untuk
  menunjukkan pola konvergensi beta (Figure 4, hlm. 985);
\item
  pertumbuhan produktivitas tenaga kerja dalam kaitannya dengan waktu
  tempuh ke Tokyo, pengurangan waktu tempuh, dan spesialisasi manufaktur
  (Figures 5-7, hlm. 985-988);
\item
  koefisien variasi GVA per kapita sebelum dan sesudah transfer fiskal
  (Figure 8, hlm. 989-990).
\end{itemize}

\textbf{{[}Inferensi pengolahan{]}} Indikator-indikator tersebut
berfungsi sebagai outcome deskriptif dalam review ini, bukan sebagai
variabel dependen yang telah diestimasi dalam spesifikasi statistik yang
dijelaskan lengkap. Persamaan, kovariat, dan prosedur estimasi tidak
disebutkan dalam file.

\subsection{Results}\label{results-41}

\textbf{{[}Dilaporkan sumber{]}} Hasil numerik dan pola utama yang
dilaporkan artikel adalah sebagai berikut:

\begin{itemize}
\tightlist
\item
  Greater Tokyo Area menyumbang 33,53\% GVA riil nasional dan 28,89\%
  penduduk pada 2019; selisih pangsa tersebut adalah 4,64 poin
  persentase. Kansai menyumbang 15,28\% GVA dan 16,28\% penduduk,
  sedangkan Chubu menyumbang 14,55\% GVA dan 13,48\% penduduk (Table 1,
  hlm. 983).
\item
  Tokyo memiliki 19,75\% pangsa GVA nasional dan 10,89\% pangsa penduduk
  pada 2019. Lima prefektur teratas mencakup 44,12\% GVA dan 36,88\%
  penduduk; nilai pembanding 2014 masing-masing 42,79\% dan 35,50\%
  (Table 2, hlm. 984).
\item
  Koefisien variasi GVA per kapita turun tajam pada 1990-1995, sempat
  meningkat pada 2005, kemudian kembali turun hingga 2019 dan berada di
  bawah tingkat 1990 (Figure 3, hlm. 984-985).
\item
  Hubungan antara GVA per kapita pada 1990 dan laju pertumbuhan sampai
  2019 menunjukkan korelasi negatif yang oleh artikel disebut
  mengindikasikan konvergensi beta (Figure 4, hlm. 985).
\item
  Wilayah yang lebih jauh dari Tokyo menunjukkan pertumbuhan
  produktivitas tenaga kerja yang relatif lebih tinggi. Pengurangan
  waktu tempuh ke Tokyo berkorelasi positif dan signifikan dengan
  pertumbuhan produktivitas tenaga kerja, dengan koefisien korelasi
  0,35; pola tersebut paling menonjol pada Tokushima, Yamaguchi,
  Shimane, Aomori, dan Iwate (Figures 5-6, hlm. 985-987).
\item
  Spesialisasi manufaktur berkorelasi positif kuat dengan pertumbuhan
  produktivitas tenaga kerja. Shiga, Tokushima, Gunma, dan Ibaraki
  disebut sebagai prefektur dengan aglomerasi manufaktur tinggi dan
  pertumbuhan produktivitas di atas rata-rata (Figure 7, hlm. 988).
\item
  Koefisien variasi GVA per kapita menunjukkan pergeseran turun setelah
  transfer fiskal (Figure 8, hlm. 989-990).
\end{itemize}

\textbf{{[}Inferensi pengolahan{]}} Ukuran sampel, p-value, interval
ketidakpastian, persamaan, dan uji robustness: \textbf{Tidak disebutkan
dalam file}. Karena itu, pola-pola tersebut tidak dapat dibaca secara
statistik lebih rinci.

\subsection{Findings}\label{findings-41}

\textbf{{[}Interpretasi penulis{]}} Artikel membaca konvergensi GVA
regional sebagai indikasi bahwa eksternalitas ekonomi yang tetap
terkonsentrasi di inti dapat didistribusikan melalui jaringan
transportasi dan keterkaitan industri. Penurunan jarak waktu diposisikan
sebagai bukti relevansi lapisan konektivitas jaringan, sedangkan
keterkaitan industri dan dukungan institusional diposisikan sebagai
fondasi yang membuat manfaat tersebut dapat diserap dan dipertahankan
(Section 4.4, hlm. 990).

\textbf{{[}Interpretasi penulis{]}} Spesialisasi manufaktur diperlakukan
sebagai \emph{absorptive interface}: keunggulan sektoral membantu
wilayah periferal menyalurkan dan memperkuat limpahan metropolitan.
Pembagian kerja yang saling melengkapi antara wilayah manufaktur dan
metropolitan dipandang berbeda dari ketergantungan satu arah (Section
4.3, hlm. 987-990).

\textbf{{[}Interpretasi penulis{]}} National Spatial Strategy (NSS),
Local Allocation Tax, infrastruktur, koordinasi, dan modal sosial
diposisikan sebagai contoh lapisan institusional-relasional. Transfer
fiskal dan perencanaan spasial nasional disebut membantu memperkuat
kapasitas wilayah periferal untuk berpartisipasi dalam jaringan produksi
serta mempertahankan manfaat eksternalitas (Section 4.3, hlm. 988-990).

\textbf{{[}Inferensi pengolahan{]}} Secara keseluruhan, hasil deskriptif
konsisten dengan mekanisme tiga lapisan sebagaimana dirumuskan penulis,
tetapi tidak dapat membedakan secara kausal apakah konektivitas,
spesialisasi, atau institusi menghasilkan pertumbuhan karena artikel
tidak melakukan inferensi kausal atau validasi formal.

\subsection{Conclusions}\label{conclusions-41}

\textbf{{[}Interpretasi penulis{]}} Kesimpulan artikel adalah bahwa
\emph{borrowed size} sebaiknya direkonseptualisasi sebagai konstruksi
tiga lapisan yang mencakup eksternalitas terkait skala, struktur
jaringan mediasi, dan fondasi institusional-relasional. Kerangka
multisakala dan nonlinier ini dimaksudkan untuk menjelaskan koeksistensi
aglomerasi yang bertahan dengan penurunan disparitas antardaerah di
Jepang (Section 6, hlm. 992-993).

\textbf{{[}Interpretasi penulis{]}} \emph{Borrowed size} dipahami
sebagai proses berkelanjutan dari mediasi institusional dan adaptasi
jaringan, bukan kondisi statis. Karena itu, kebijakan regional
disarankan berfokus pada kerja sama institusional, koordinasi tata
kelola, dan sirkulasi pengetahuan, bukan sekadar meniru aglomerasi
metropolitan (Section 6, hlm. 992-993).

\textbf{{[}Dilaporkan sumber{]}} Artikel menyatakan bahwa ilustrasi
prefektur Jepang mendukung interpretasi berbasis proses, tetapi juga
menegaskan bahwa kajian ini merupakan fondasi konseptual dan ilustrasi
relevansi kontekstual, bukan inferensi kausal atau validasi formal
(Section 6, hlm. 993).

\subsection{Contributions}\label{contributions-41}

\textbf{{[}Interpretasi penulis{]}} Artikel mengklaim tiga kontribusi
utama:

\begin{itemize}
\tightlist
\item
  mereframe \emph{borrowed size} sebagai teori yang memasukkan ketebalan
  institusional dan konektivitas relasional;
\item
  menempatkan kerangka tersebut dalam konteks empiris ekonomi regional
  Jepang, dengan potensi generalisasi ke konteks Asia lain;
\item
  menawarkan alternatif terhadap kebijakan regional yang berpusat pada
  aglomerasi melalui desain institusional berorientasi jaringan dan
  kemungkinan pertumbuhan regional terdistribusi (Section 1, hlm. 977).
\end{itemize}

\textbf{{[}Interpretasi penulis{]}} Artikel juga menyatukan pemahaman
\emph{borrowed size} sebagai keadaan, mekanisme, dan proses dalam satu
hierarki konseptual, serta menghubungkan institusi dengan konfigurasi
spasial (Sections 3 dan 6, hlm. 979-980, 992-994).

\subsection{Research Gap}\label{research-gap-41}

\textbf{{[}Dilaporkan sumber{]}} Sebelum artikel ini, literatur yang
ditinjau dinilai lebih sering menjelaskan \emph{borrowed size} melalui
kedekatan spasial dan jarak fisik. Fondasi institusional-relasional,
alasan keberhasilan yang berbeda antarwilayah, serta kondisi yang
memungkinkan struktur mediator muncul dan bertahan disebut masih kurang
dieksplorasi (Section 2, hlm. 977-978).

\textbf{{[}Dilaporkan sumber{]}} Artikel juga menyebut kurangnya
pemahaman tentang kemunculan, evolusi, dan kemungkinan pembubaran
\emph{borrowed size}. Pada bagian penutup, pengukuran komponen
institusional dan relasional masih dinilai belum berkembang secara
metodologis, sementara indikator tingkat prefektur belum menangkap
interaksi mikroskala dan dinamika perubahan institusi (Section 6, hlm.
993).

\textbf{{[}Inferensi pengolahan{]}} Dengan demikian, gap empiris yang
tetap terbuka adalah menguji mekanisme tiga lapisan melalui ukuran
relasional dan institusional yang teroperasionalisasi, desain
longitudinal atau multisakala, serta strategi identifikasi yang
melampaui korelasi deskriptif. Kebutuhan ini diturunkan dari batasan dan
agenda riset yang dinyatakan sumber, bukan dari sumber eksternal.

\subsection{Limitations}\label{limitations-41}

\textbf{{[}Dilaporkan sumber{]}} Artikel menyebut tiga keterbatasan
utama:

\begin{itemize}
\tightlist
\item
  tingkat abstraksi teoretis yang tinggi menimbulkan pertanyaan tentang
  keseimbangan antara kejelasan konseptual dan ketepatan empiris;
\item
  operasionalisasi dan pengukuran komponen institusional-relasional
  \emph{borrowed size} masih belum berkembang secara metodologis;
\item
  cakupan indikator lintas-seksi pada tingkat prefektur mungkin tidak
  menangkap interaksi mikroskala atau dinamika temporal perubahan
  institusi (Section 6, hlm. 993).
\end{itemize}

\textbf{{[}Dilaporkan sumber{]}} Artikel tidak bertujuan menghasilkan
inferensi kausal atau validasi formal. Karena itu, bukti yang disajikan
berfungsi sebagai fondasi konseptual dan ilustrasi relevansi dalam
konteks Jepang, bukan pengujian formal atas model (Section 6, hlm. 993).

\subsection{Challenges}\label{challenges-41}

\textbf{{[}Dilaporkan sumber{]}} Tantangan yang dibahas dalam file
meliputi:

\begin{itemize}
\tightlist
\item
  penurunan penduduk dan pengosongan industri yang mengancam otonomi
  serta ketahanan ekonomi wilayah regional (Section 1, hlm. 975);
\item
  menerjemahkan kedekatan temporal menjadi keuntungan ekonomi, karena
  kedekatan merupakan kondisi yang mungkin perlu tetapi tidak cukup;
  hasilnya bergantung pada kepadatan, redundansi, dan ketahanan jaringan
  serta dukungan kebijakan dan tata kelola (Section 4.2, hlm. 987);
\item
  membangun kapasitas fiskal, koordinasi administratif, dan kemampuan
  \emph{network brokerage} sambil mengelola konflik, termasuk kompetisi
  pajak antarkewenangan (Section 3.3, hlm. 981-982);
\item
  menjaga agar manfaat tidak bocor kembali ke inti melalui arus keuangan
  dan sumber daya yang tidak seimbang, atau melalui masalah \emph{leaky
  bucket} dan \emph{agglomeration shadow} (Sections 3.2-3.3, hlm.
  981-982).
\end{itemize}

\subsection{Practical Implications}\label{practical-implications-41}

\textbf{{[}Interpretasi penulis{]}} Artikel merumuskan empat pilar
desain kebijakan yang dipetakan ke lapisan model:

\begin{itemize}
\tightlist
\item
  memastikan interoperabilitas institusional melalui harmonisasi
  regulasi, perbaikan tata kelola antarkewenangan, dan koordinasi
  administratif (Layer 3; Section 5.4.1, hlm. 991-992);
\item
  menumbuhkan modal relasional melalui mobilitas manusia, jaringan
  akademia-industri lintas wilayah, dan kemitraan berbasis kepercayaan
  (Layers 2-3; Section 5.4.2, hlm. 992);
\item
  merancang ulang skala fungsional agar kebijakan mengikuti ruang
  ekonomi dan ruang hidup aktual, bukan semata batas administratif
  (Layers 1-2; Section 5.4.3, hlm. 992);
\item
  memperkuat struktur mediasi melalui organisasi perantara, hub
  regional, dan kemitraan publik-swasta untuk inovasi serta koordinasi
  kerja sama lintas wilayah (Layer 2; Section 5.4.4, hlm. 992).
\end{itemize}

\textbf{{[}Interpretasi penulis{]}} Transfer fiskal, pembagian
pendapatan, reinvestasi pendapatan lokal, infrastruktur
transportasi-komunikasi-logistik, dan platform perencanaan nasional
diposisikan sebagai cara memperkuat kapasitas penyerapan wilayah
periferal. Arah kebijakannya bergeser dari redistribusi reaktif menuju
integrasi proaktif yang dapat mengubah disparitas menjadi
komplementaritas (Sections 3.2-3.3 dan 5.3-5.4, hlm. 981-982, 991-992).

\subsection{Applications}\label{applications-41}

\textbf{{[}Dilaporkan sumber{]}} Penerapan empiris artikel berada pada
konteks prefektur Jepang. National Spatial Strategy menjadi contoh
platform formal untuk mengatur kebijakan spasial, infrastruktur, dan
hubungan antarprefektur; rencana Chuo Shinkansen sepanjang 438 km antara
Tokyo, Nagoya, dan Osaka digunakan sebagai ilustrasi kerja bersama Layer
3 dan Layer 2 (Section 4.3, hlm. 988-989).

\textbf{{[}Interpretasi penulis{]}} Kerangka ini ditawarkan untuk
menganalisis pertumbuhan terdistribusi dan eksternalitas yang dimediasi
jaringan di negara atau wilayah lain, terutama yang mengalami transisi
demografis, konsentrasi metropolitan, dan restrukturisasi spasial
(Section 5.4, hlm. 991-992; Section 6, hlm. 993-994).

\textbf{{[}Inferensi pengolahan{]}} Implementasi kebijakan yang telah
dievaluasi atau penerapan kasus di luar Jepang: \textbf{Tidak disebutkan
dalam file}. Kerangka tersebut karena itu belum dapat diperlakukan
sebagai alat prediksi yang tervalidasi.

\subsection{Future Research
(Optional)}\label{future-research-optional-41}

\textbf{{[}Dilaporkan sumber{]}} Agenda riset yang dinyatakan artikel
adalah:

\begin{itemize}
\tightlist
\item
  mengembangkan pemodelan dinamis melalui simulasi jaringan spasial
  evolusioner dan dataset longitudinal untuk menangkap pembentukan,
  pendalaman, dan pembubaran \emph{borrowed size};
\item
  melakukan analisis multisakala yang mencakup munisipalitas, kawasan
  metropolitan, dan kawasan fungsional yang lebih luas untuk menguji
  cakupan serta ketergantungan skala mekanisme;
\item
  melakukan studi institusional komparatif lintas negara dan wilayah
  untuk membedakan universalitas dari kekhususan kontekstual
  \emph{borrowed size};
\item
  mengembangkan indikator relasional baru, termasuk lintasan mobilitas,
  keterkaitan rantai pasok, dan kolaborasi universitas-industri (Section
  6, hlm. 993).
\end{itemize}

\subsection{Provenance Notes}\label{provenance-notes-41}

\begin{itemize}
\tightlist
\item
  \textbf{Bahan utama:}
  \texttt{source/journal/A26-otsuka-2025-a-three-layer-model-of-borrowed-size-empirical-insights-from-japan-10.1080-21681376.2025.2590896.txt}.
  Nama file TXT dengan prefiks \texttt{A26} diverifikasi di bawah
  \texttt{source/} sebelum pembacaan; hanya pasangan sumber-review A26
  yang diproses.
\item
  \textbf{Cakupan pembacaan:} seluruh isi TXT dibaca dari baris 1 sampai
  akhir, dengan pembacaan mencapai baris bernomor 1089, termasuk blok
  \texttt{{[}PDF\ Metadata{]}} dan \texttt{{[}Extracted\ Text{]}},
  tabel, keterangan gambar, pernyataan data, dan referensi. TXT memuat
  teks artikel pada hlm. 975-995.
\item
  \textbf{Metadata ekstraksi:} TXT mencatat
  \texttt{Extracted\ on:\ 2026-08-31} dan
  \texttt{Extraction\ method:\ pdftotext\ -layout}. Metadata PDF
  mencatat judul \emph{A three-layer model of borrowed size: empirical
  insights from Japan}, 22 halaman, PDF versi 1.5, serta sumber PDF
  dengan stem nama yang sama. DOI yang tercetak dalam artikel adalah
  \texttt{10.1080/21681376.2025.2590896} (PDF Metadata dan halaman awal
  artikel, hlm. 975).
\item
  \textbf{Batas provenance:} pembacaan didasarkan pada TXT ekstraksi,
  bukan pengambilan sumber eksternal. Tata letak TXT mempertahankan
  pemenggalan halaman dan beberapa pemenggalan kata; detail visual atau
  angka yang hanya berada di dalam grafik tidak dapat diverifikasi
  melampaui narasi dan caption yang tersedia. \textbf{{[}Inferensi
  pengolahan{]}} Ini membatasi pemeriksaan ulang visual atas Figures
  3-8.
\item
  \textbf{Ketersediaan data:} artikel menyatakan data pendukung berasal
  dari sumber publik Cabinet Office dan Statistics Bureau of Japan,
  sedangkan dataset olahan dan kode analisis tersedia dari penulis atas
  permintaan yang wajar (Data availability statement, hlm. 994). Dataset
  dan kode tersebut tidak diakses dalam review ini.
\item
  \textbf{Batas sumber:} nama dan karya pada Literature Survey
  dilaporkan sebagai rujukan yang tercantum dalam artikel; tidak ada
  sumber eksternal yang ditambahkan atau digunakan untuk melengkapi
  review.
\item
  \textbf{Perubahan file:} dalam pekerjaan ini file yang dibuat adalah
  \texttt{output/kajian/review-A26.md}; TXT sumber tidak diubah. Setelah
  audit, hanya checkbox \texttt{A26} pada \texttt{Todo.md} dan
  \texttt{output/kajian/kode-kajian-pendahuluan.md} yang dicentang.
\end{itemize}

\section{Review A27}\label{review-a27}

\textbf{Kode kajian:} A27

\textbf{Identitas sumber}

\begin{itemize}
\tightlist
\item
  \textbf{Judul:} \emph{Secondary city regionalism: China's cross
  boundary development projects}
\item
  \textbf{Penulis:} Yi Li, Nicholas A. Phelps, dan Guoliang Xu
\item
  \textbf{Jurnal:} \emph{Cities}, volume 171, tahun 2026, artikel 106795
\item
  \textbf{DOI:} \texttt{10.1016/j.cities.2026.106795}
\item
  \textbf{Jenis sumber:} Artikel jurnal
\item
  \textbf{Kata kunci:} City regionalism; Secondary cities; Urban and
  economic strategizing; Borrowed size and function; Cross-boundary
  development projects; China
\item
  \textbf{Riwayat publikasi yang tercantum:} Diterima 25 Desember 2024;
  revisi diterima 3 November 2025; diterima untuk publikasi 8 Januari
  2026; tersedia daring 20 Januari 2026
\item
  \textbf{Bahan yang ditinjau:} File TXT hasil ekstraksi PDF A27, 10
  halaman, dengan metode ekstraksi \texttt{pdftotext\ -layout}
\end{itemize}

\subsection{Summarized Abstract}\label{summarized-abstract-42}

\textbf{{[}Dilaporkan sumber{]}} Artikel mengkaji bentuk-bentuk
regionalisme kota mutakhir yang diprakarsai oleh kota sekunder di
Tiongkok pascareformasi. Wacana akademik dan kebijakan selama ini banyak
berpusat pada kemunculan kawasan kota global dan penskalaan ulang peran
negara, sedangkan peran serta strategi kota sekunder dalam kawasan
kotanya kurang diperhatikan. Artikel mengisi celah tersebut melalui
perspektif \emph{outside-in}, yaitu melihat strategi otoritas kota
sekunder untuk memperoleh manfaat aglomerasi dan membatasi dampak
marginalisasi dengan memobilisasi peluang pembangunan lintas batas.
Berdasarkan kerja lapangan, artikel mengidentifikasi tiga jenis
pembangunan lintas batas: kota baru, kawasan industri, dan enklave
inovasi (Abstract, hal. PDF 1).

\textbf{{[}Interpretasi penulis{]}} Strategi regionalisasi dan program
kemitraan yang diluncurkan kota sekunder dibaca sebagai cara untuk
menemukan serta mengeksploitasi ceruk di dalam kawasan kota. Artikel
menekankan bahwa kota sekunder bukan penerima pasif limpahan dari kota
primer; otoritasnya secara aktif membentuk proyek lintas batas untuk
meminjam ukuran dan fungsi, mengakses aglomerasi, serta menghindari
bayang-bayang marginalisasi (Abstract; Sections 1 dan 5, hal. PDF 1-2
dan 7-8).

\textbf{{[}Inferensi review{]}} Kontribusi empiris utama artikel adalah
penjelasan proses dan mekanisme strategisasi, bukan pengukuran dampak
ekonomi bersih dari proyek-proyek tersebut. Desain evaluasi kausal,
ukuran efek, atau pembanding sistematis antarkasus adalah \textbf{Tidak
disebutkan dalam file}.

\subsection{Problem Statement}\label{problem-statement-42}

\textbf{{[}Dilaporkan sumber{]}} Literatur regionalisme kota sering
memperlakukan kawasan kota sebagai unit yang tidak terdiferensiasi dan
menaikkan analisis dari kota ke kawasan kota atau kawasan kota
polisentris. Pendekatan tersebut jarang membaca sistem urban, lanskap
ekonomi, dan lanskap kebijakan dari perspektif \emph{outside-in}.
Literatur juga belum cukup menempatkan kota sekunder dalam pembagian
kerja intra-kawasan kota serta imajinasi kebijakan yang menyertainya
(Section 1, hal. PDF 1-2).

\textbf{{[}Dilaporkan sumber{]}} Logika regionalisme kota yang berpusat
pada aglomerasi kota primer mengandalkan pertumbuhan, limpahan, dan
\emph{trickle-down} ke wilayah lain. Dalam konteks Tiongkok, ketimpangan
yang membesar antara Beijing dan daerah sekitarnya disebut sebagai
konsekuensi pendekatan yang mengutamakan kota primer. Kritik terhadap
pendekatan ini menunjukkan bahwa limpahan yang diasumsikan tidak selalu
terjadi dan bahwa kota provinsi, kota kecil, serta wilayah perdesaan di
luar atau di dalam hinterland kawasan kota dapat menerima dampak negatif
(Section 2, hal. PDF 2-3).

\textbf{{[}Dilaporkan sumber{]}} Masalah penelitian dirumuskan melalui
pertanyaan: bagaimana kota sekunder merespons proses yang digerakkan
aglomerasi, serta imajinasi dan praktik perencanaan strategis
urban-ekonomi apa yang tampak dalam regionalisme kota. Artikel
memusatkan perhatian pada perjuangan strategis otoritas kota sekunder
untuk menemukan ceruk dalam organisasi ekonomi dan pengaturan tata
kelola kawasan kota, terutama melalui peluang pembangunan lintas batas
(Sections 1 dan 2.2, hal. PDF 1-3).

\textbf{{[}Interpretasi penulis{]}} Regionalisme kota dipahami bukan
sebagai hasil ekonomi yang otomatis, melainkan sebagai proyek
politik-ekonomi yang terus berubah. Karena itu, batas yurisdiksi
diperlakukan sekaligus sebagai hambatan, ruang konflik, dan objek yang
dapat dimobilisasi untuk menghasilkan hubungan antarkota atau
\emph{interplace economy} (Sections 2 dan 2.2, hal. PDF 2-3).

\textbf{{[}Inferensi review{]}} Problem statement artikel menggeser
fokus analitis dari pertanyaan apakah aglomerasi kota primer menetes ke
bawah menuju pertanyaan tentang siapa yang mengorkestrasi aliran,
melalui proyek apa, dan dengan pembagian kerja lintas batas seperti apa.

\subsection{Objectives}\label{objectives-42}

\textbf{{[}Dilaporkan sumber{]}} Tujuan yang dinyatakan atau dijelaskan
dalam artikel adalah:

\begin{itemize}
\tightlist
\item
  mengkaji modalitas regionalisme kota yang diprakarsai kota sekunder di
  Tiongkok pascareformasi;
\item
  mengisi celah konseptual dan empiris dengan memeriksa agensi serta
  strategi otoritas kota sekunder;
\item
  menjelaskan bagaimana kota sekunder mencari ceruk dalam organisasi
  ekonomi dan tata kelola kawasan kota;
\item
  menelaah mobilisasi peluang pembangunan lintas batas sebagai bagian
  dari strategi kota sekunder;
\item
  mengidentifikasi dan membandingkan tiga jenis objek batas: kota baru
  lintas batas, kawasan industri lintas batas, dan enklave inovasi
  lintas batas;
\item
  menunjukkan bagaimana proyek-proyek tersebut digunakan untuk meminjam
  fungsi, ukuran industri, talenta, layanan, merek, dan pembagian kerja
  dari kota primer atau kawasan metropolitan (Sections 1-4, hal. PDF
  1-7).
\end{itemize}

\textbf{{[}Dilaporkan sumber{]}} Artikel menyatakan dua kontribusi
langsung sejak pendahuluan: memusatkan perhatian pada perjuangan
strategis otoritas kota sekunder untuk menemukan ceruk, serta menarik
perhatian pada mobilisasi peluang pembangunan lintas batas dalam konteks
Tiongkok (Section 1, hal. PDF 1-2).

\textbf{{[}Inferensi review{]}} Hipotesis formal, pertanyaan
subhipotesis, atau target pengujian statistik adalah \textbf{Tidak
disebutkan dalam file}.

\textbf{{[}Dilaporkan sumber{]}} Artikel menyajikan pemahaman konseptual
dan empiris berbasis kasus.

\subsection{Literature Survey}\label{literature-survey-42}

\textbf{{[}Dilaporkan sumber{]}} Artikel memetakan tiga rumpun literatur
tentang pembentukan kawasan kota. Rumpun pertama membahas kekuatan
ekonomi yang menghasilkan raksasa urban global, dengan penekanan pada
aglomerasi, pertumbuhan, dan interaksi spasial. Rumpun kedua memetakan
polisentrisitas kawasan kota global melalui ekonomi jasa produsen dan
konektivitas informasi. Rumpun ketiga menguraikan dinamika politik dan
tata kelola kawasan kota (Section 1, hal. PDF 1-2).

\textbf{{[}Dilaporkan sumber{]}} Pembahasan kemudian membedakan
perencanaan urban-ekonomi setelah kawasan kota terbentuk dari
pembayangan awal yang berkontribusi pada pembentukan kawasan kota.
Literatur yang ada dinilai lebih sering berfokus pada strategi negara
dan otoritas kota primer daripada otoritas kota sekunder. Literatur
polisentrisitas dipakai untuk menegaskan keberadaan banyak pusat, mode
pembentukan seperti \emph{centrifugal}, \emph{incorporation}, dan
\emph{fusion}, serta kemungkinan pembagian kerja fungsional di antara
aglomerasi ekonomi yang berbeda (Sections 1 dan 2.1, hal. PDF 1-3).

\textbf{{[}Dilaporkan sumber{]}} Konsep \emph{borrowed size} ditelusuri
ke Alonso, lalu dikembangkan dalam kajian tentang sistem perkotaan dan
kawasan kota polisentris. Konsep tersebut menjelaskan peluang perusahaan
untuk memperoleh biaya kemacetan dan sewa yang lebih rendah di
permukiman kecil sambil mengakses pasar besar, layanan bisnis, dan pasar
tenaga kerja yang beragam di permukiman besar. Artikel membedakan
\emph{borrowed size} dari \emph{borrowed function}, serta dari efek
bayangan ketika permukiman di orbit kota besar mengalami eksternalitas
negatif (Section 2.1, hal. PDF 3).

\textbf{{[}Dilaporkan sumber{]}} Literatur tentang perifer urban,
suburbanisasi, dan kritik terhadap \emph{methodological cityism}
digunakan untuk menunjukkan bahwa batas serta wilayah perifer tidak
cukup dipahami dari perspektif kota utama. Literatur tentang
regionalisme kota, perencanaan spasial strategis, politik kawasan, dan
geopolitik kapitalisme mendukung argumen bahwa kawasan kota dibentuk
oleh banyak aktor, bukan hanya kekuatan pasar atau negara nasional
(Sections 2 dan 2.2, hal. PDF 2-3).

\textbf{{[}Dilaporkan sumber{]}} Untuk konteks Tiongkok, artikel
menggunakan literatur tentang hierarki administratif urban, aneksasi
administratif, tata kelola regional, kewirausahaan negara, pembangunan
kawasan suburban, kawasan ekonomi khusus, dan hubungan antarkota di
Delta Sungai Yangtze. Konsep \emph{boundary object} dipakai untuk
membaca kota baru, kawasan industri, dan enklave inovasi sebagai objek
konkret tempat aktor negara, pengembang, sektor privat, dan warga
menghubungkan kepentingan yang berbeda (Sections 2.2 dan 3, hal. PDF
3-4).

\textbf{{[}Interpretasi penulis{]}} Sintesis literatur mengarah pada
kerangka bahwa limpahan aglomerasi tidak berlangsung sebagai hukum alam.
Limpahan bergantung pada strategisasi aktor, pembagian kerja,
infrastruktur, hubungan lintas yurisdiksi, dan pengaturan kelembagaan.
Kota sekunder karena itu diposisikan sebagai promotor imajinasi kawasan
kota bersama negara dan sektor privat (Sections 2 dan 2.2, hal. PDF
2-3).

\textbf{{[}Inferensi review{]}} Tinjauan pustaka dalam file berbentuk
naratif. Protokol pencarian, basis data yang ditelusuri, kriteria
inklusi-eksklusi, jumlah korpus, dan prosedur penilaian kualitas
literatur adalah \textbf{Tidak disebutkan dalam file}. Rujukan yang
disebut di sini hanya dilaporkan sebagaimana tercantum dalam artikel,
bukan diverifikasi melalui sumber luar.

\subsection{Method Used}\label{method-used-42}

\textbf{{[}Dilaporkan sumber{]}} Penelitian didasarkan pada lebih dari
20 kunjungan lapangan ke kawasan kota di Tiongkok, terutama
Beijing-Tianjin-Hebei atau BTH dan Delta Sungai Yangtze atau YRD, pada
periode 2008-2023. Kunjungan berulang dipakai untuk membangun
sensitivitas terhadap jaringan aktor yang mendasari pembangunan kawasan
kota, membangun kepercayaan, melakukan wawancara, dan mengumpulkan
dokumen internal serta bahan sekunder (Section 3, hal. PDF 3-4).

\textbf{{[}Dilaporkan sumber{]}} Lebih dari 10 lokasi dikunjungi,
sebagian lebih dari sekali. Kasus kemudian dipilih untuk
merepresentasikan tiga jenis perencanaan \emph{boundary object} yang
ditekankan artikel. Lokasi contoh kasus ditunjukkan dalam Fig. 1 dan
tipologi, logika strategisasi, serta latar belakang contoh diringkas
dalam Table 2 (Sections 3 dan 4, hal. PDF 3-4).

\textbf{{[}Dilaporkan sumber{]}} Wawancara menggunakan \emph{snowball
sampling}. Peneliti berupaya memperluas jangkauan informan untuk
mencakup pejabat provinsi, kota, kabupaten atau distrik; perencana;
akademisi universitas; pengusaha; dan penduduk. Sebanyak 15 wawancara
dilakukan dengan pejabat pemerintah, perusahaan privat, perencana, dan
informan lain. Enam simposium pemerintah mengenai pembangunan kota kecil
dan kota berlangsung juga dihadiri. Wawancara berdurasi 30 menit sampai
lebih dari satu jam; seluruh wawancara dan simposium berjumlah sekitar
50 jam. Transkrip dibuat segera setelah wawancara untuk menjaga
ketepatan ingatan (Section 3 dan Table 1, hal. PDF 3-4).

\textbf{{[}Dilaporkan sumber{]}} Data wawancara ditriangulasi dengan
upaya memperluas tipe informan dan dengan konsultasi terhadap situs web
pemerintah lokal, surat kabar lokal, serta karya yang telah
dipublikasikan. Nama informan dan individu yang disebut informan
dianonimkan (Section 3, hal. PDF 3-4).

\textbf{{[}Inferensi review{]}} Secara operasional, desain ini dapat
dibaca sebagai studi kasus kualitatif multi-kasus dengan triangulasi
sumber dan variasi tipe proyek. Teknik analisis seperti pengodean
transkrip, analisis tematik, proses perbandingan kasus, atau prosedur
penelusuran bukti adalah \textbf{Tidak disebutkan dalam file}. Karena
itu, prosedur analisis yang dapat direplikasi tidak tersedia secara
lengkap.

\textbf{{[}Dilaporkan sumber{]}} Artikel menggunakan dua atau tiga
contoh dari BTH atau YRD untuk setiap jenis proyek guna menunjukkan
logika dan langkah konkret strategisasi (Section 4, hal. PDF 4).

\textbf{{[}Inferensi review{]}} Tujuan pemilihan kasus untuk estimasi
representatif populasi kota sekunder adalah \textbf{Tidak disebutkan
dalam file}.

\subsection{Dataset}\label{dataset-42}

\textbf{{[}Dilaporkan sumber{]}} Bahan empiris utama adalah hasil
kunjungan lapangan, wawancara, simposium pemerintah, dokumen internal,
situs web pemerintah lokal, surat kabar lokal, dan karya yang telah
dipublikasikan. Artikel menyebutkan bahwa bahan-bahan ini digunakan
untuk menguraikan karakteristik dasar, aktor, konteks, dan motivasi tiga
jenis pembangunan lintas batas (Section 3, hal. PDF 3-4).

\textbf{{[}Dilaporkan sumber{]}} Table 1 memuat 15 wawancara, tanggal,
contoh proyek, serta relevansi informan. Kasus yang tercantum adalah
Gu'an New Town, Huaqiao New Town, Jiangyin-Jinjiang Industrial Park,
Suzhou-Chuzhou Industrial Park, Changxing innovation enclave, dan Quzhou
innovation enclave. Table 2 mengelompokkan kasus tersebut ke dalam tiga
tipe proyek dan mencantumkan fungsi strategis serta latar belakang
masing-masing (Tables 1-2, hal. PDF 4).

\textbf{{[}Dilaporkan sumber{]}} Angka tambahan di dalam narasi berasal
dari dokumen atau survei yang dirujuk artikel. Laporan komuter Beijing
tahun 2022 digunakan untuk menyebut sekitar 330.000 orang yang setiap
hari melakukan perjalanan ke Beijing dari wilayah sekitar, setara 3\%
dari seluruh komuter Beijing. Laporan komuter lintas batas YRD tahun
2020 digunakan untuk menyebut Kunshan dan Taicang sebagai sumber utama
komuter ke Shanghai, masing-masing 74,8\% dan 12,7\% dalam uraian
artikel (Section 4.1, hal. PDF 5).

\textbf{{[}Dilaporkan sumber{]}} Penulis memperoleh akses selama kerja
lapangan terhadap hasil survei Taicang tahun 2022 yang tidak tersedia
untuk publik. Dari 953 perusahaan yang berpartisipasi, 671 disebut
terhubung dengan Shanghai di sisi hulu atau hilir rantai industri, atau
70,4\%. Untuk manufaktur, proporsinya 76,9\%; untuk sektor jasa, artikel
menyebut 50\% terkait dengan Shanghai dan 60\% dari kelompok yang
terkait itu berada pada perdagangan dan logistik. Angka lain menyebut
623 dari 953 perusahaan memiliki karyawan yang melakukan komuter, atau
65\%; teks menyebut 54\% melakukan komuter harian dari Shanghai ke
Taicang, dengan 29,6\% manajer dan 37,3\% teknisi (Section 4.1, hal. PDF
5-6).

\textbf{{[}Dilaporkan sumber{]}} File memiliki pernyataan \emph{Data
availability} yang berbunyi: ``No data was used for the research
described in the article.'' Pernyataan tersebut tercantum setelah bagian
Etika, meskipun bagian metode dan hasil menerangkan wawancara, dokumen,
kunjungan lapangan, serta hasil survei yang diakses penulis (Data
availability dan Section 3, hal. PDF 8 serta 3-4).

\textbf{{[}Inferensi review{]}} Dataset terstruktur, transkrip, matriks
pengodean, instrumen wawancara, lampiran dokumen, atau kode analisis
adalah \textbf{Tidak disebutkan dalam file}. Status akses terhadap data
mentah selain uraian survei Taicang juga \textbf{Tidak disebutkan dalam
file}. Ketidaksesuaian antara pernyataan ketersediaan data dan uraian
bahan empiris tidak dapat diselesaikan hanya dari TXT.

\subsection{Population / Sample}\label{population-sample-42}

\textbf{{[}Inferensi review berbasis definisi sumber{]}} Cakupan
konseptual penelitian adalah kota sekunder dalam kawasan kota global di
Tiongkok. Pada bagian pendahuluan, kota sekunder dijelaskan secara
kontekstual sebagai: kota tingkat prefektur yang lebih rendah dan bukan
munisipalitas yang dikelola pusat atau ibu kota provinsi; kota tingkat
kabupaten yang signifikan secara ekonomi dan memiliki skala ekonomi atau
peran khusus dalam rantai nilai regional, seperti Kunshan; serta kota
dengan status administratif sama tetapi berskala jauh lebih kecil
daripada kota tingkat prefektur pembanding, seperti Chuzhou dibandingkan
Suzhou (Section 1, hal. PDF 1-2).

\textbf{{[}Dilaporkan sumber{]}} Pada bagian definisi, istilah kota
sekunder juga digunakan secara dasar untuk merujuk pada kota nonprimer
di dalam kawasan kota global tertentu. Klasifikasi ini sengaja longgar
dan bukan kerangka teliti untuk mengategorikan kota kecil, menengah, dan
besar dalam sistem urban regional, nasional, atau internasional. Kota
sekunder dipilih melalui lensa relasional untuk menunjukkan dilema
pembangunan dan strategi otoritas dalam konteks pembangunan tidak merata
(Section 3, hal. PDF 3-4).

\textbf{{[}Dilaporkan sumber{]}} Sampel kasus terdiri atas dua contoh
kota baru lintas batas, dua contoh kawasan industri lintas batas, dan
dua contoh enklave inovasi lintas batas. Contohnya adalah Gu'an New Town
dan Huaqiao New Town; Jiangyin-Jingjiang Industrial Park dan
Suzhou-Chuzhou Industrial Park; serta Changxing innovation enclave di
Shanghai dan Quzhou innovation enclave di Hangzhou (Table 2 dan Sections
4.1-4.3, hal. PDF 4-7).

\textbf{{[}Dilaporkan sumber{]}} Unit informan berjumlah 15 wawancara.
Table 1 mengaitkan wawancara dengan proyek sebagai berikut: Gu'an dengan
profesional real estat dan perusahaan; Huaqiao dengan kepala perencana
dan pejabat pemerintah; Jiangyin-Jinjiang dengan perusahaan dan pejabat
pemerintah; Suzhou-Chuzhou dengan dua informan perusahaan dan satu
pejabat pemerintah; Changxing dengan pejabat pemerintah, perusahaan, dan
pejabat enklave; serta Quzhou dengan pejabat enklave dan perusahaan
(Table 1, hal. PDF 4).

\textbf{{[}Dilaporkan sumber{]}} Enam simposium pemerintah dan lebih
dari 10 lokasi kunjungan juga menjadi bagian dari proses pengumpulan
bahan (Section 3, hal. PDF 3-4).

\textbf{{[}Inferensi review{]}} Jumlah peserta simposium, jumlah
kunjungan per lokasi, dan distribusi informan menurut kategori adalah
\textbf{Tidak disebutkan dalam file}.

\textbf{{[}Inferensi review{]}} Ukuran populasi kota sekunder di
Tiongkok, jumlah kota yang disaring, kriteria eksklusi kasus selain
keterwakilan tiga tipe, dan dasar penentuan kecukupan atau saturasi
wawancara adalah \textbf{Tidak disebutkan dalam file}. Sampel ini karena
itu lebih tepat dipahami sebagai sampel kasus bertujuan dan informan
berbasis jejaring, bukan sampel probabilitas.

\textbf{{[}Dilaporkan sumber{]}} Artikel menyatakan bahwa informan
memberi izin untuk direkam dan dikutip, menyetujui penyebutan peran
profesional, dan nama mereka dianonimkan. Tidak ada dokumen persetujuan
etika formal yang ditandatangani karena penulis menyatakan Tiongkok
tidak memiliki sistem aplikasi etika yang ekuivalen (Ethical statement,
hal. PDF 8).

\subsection{Dependent Variables}\label{dependent-variables-42}

\textbf{{[}Inferensi review{]}} Variabel dependen formal adalah
\textbf{Tidak berlaku}. Artikel tidak membangun model kausal, regresi,
atau eksperimen, dan tidak menetapkan variabel dependen atau independen
formal pada uraian survei yang dikutip.

\textbf{{[}Dilaporkan sumber{]}} Objek analisis yang dipelajari adalah
modalitas dan strategisasi kota sekunder dalam regionalisme kota,
terutama pembentukan proyek lintas batas. Artikel mengamati bagaimana
otoritas kota sekunder berupaya meminjam: fungsi dan standar layanan
melalui kota baru; ukuran atau massa industri melalui kawasan industri;
serta ukuran kumpulan talenta dan fungsi teknologi melalui enklave
inovasi (Table 2 dan Sections 4.1-4.3, hal. PDF 4-7).

\textbf{{[}Dilaporkan sumber{]}} Konsekuensi atau proses yang ditelusuri
secara naratif mencakup konektivitas transportasi, aliran komuter dan
penduduk, hubungan rantai pasok, pertukaran merek dan layanan
pengelolaan, relokasi atau perluasan produksi, serta kemungkinan
perpindahan talenta dan produksi di kemudian hari (Sections 4.1-4.3,
hal. PDF 5-7).

\textbf{{[}Inferensi review{]}} Angka komuter dan keterkaitan perusahaan
berfungsi sebagai bukti kontekstual dalam TXT, bukan outcome yang
diestimasi melalui spesifikasi statistik.

\textbf{{[}Inferensi review{]}} Definisi operasional bersama untuk
keberhasilan proyek, kinerja ekonomi, peningkatan kesejahteraan, atau
pengurangan marginalisasi adalah \textbf{Tidak disebutkan dalam file}.
Artikel sendiri menyatakan bahwa penelitian berfokus pada kehadiran,
logika, dan manifestasi strategisasi, bukan pada penimbangan efeknya
(Conclusion, hal. PDF 8).

\subsection{Results}\label{results-42}

\textbf{{[}Dilaporkan sumber{]}} Artikel mengelompokkan proyek
pembangunan lintas batas ke dalam tiga tipe. Kota baru lintas batas
dikaitkan dengan \emph{borrowing function}, merek kota primer, dan
standar layanan. Kawasan industri lintas batas dikaitkan dengan
peminjaman ukuran atau massa industri. Enklave inovasi lintas batas
dikaitkan dengan peminjaman ukuran talenta berteknologi tinggi serta
peminjaman fungsi melalui pembagian kerja industri (Table 2, hal. PDF
4).

\textbf{{[}Dilaporkan sumber{]}} Pada tipe kota baru, Gu'an New Town di
tepi Langfang diposisikan dekat Beijing dan dipasarkan melalui kedekatan
dengan Tian'anmen Square, termasuk melalui papan iklan yang
didokumentasikan dalam Fig. 2. Strateginya ialah menggunakan
transportasi regional dan citra kedekatan dengan Beijing untuk menarik
investasi dan industri limpahan. Huaqiao New Town di Kunshan, dekat
Shanghai, digunakan sebagai contoh lain; jalur metro lintas provinsi
pertama di Tiongkok dibangun untuk mengintegrasikan Huaqiao dengan
Shanghai (Section 4.1 dan Figs. 1-2, hal. PDF 4-5).

\textbf{{[}Dilaporkan sumber{]}} Kota baru tersebut tidak hanya dibangun
sebagai ruang hunian. Gu'an berupaya meningkatkan layanan kolektif dan
standar agar menyamai Beijing atau suburbnya; langkah yang disebut
mencakup perencanaan, kesehatan, dan pendidikan. Artikel juga menyatakan
bahwa investor diharapkan memperoleh akses pada infrastruktur,
fasilitas, pendidikan, kesehatan, riset ilmiah, dan talenta yang
diasosiasikan dengan Beijing, dengan biaya produksi dan reproduksi yang
lebih rendah. Penulis menyebut empat dimensi yang dituju, yaitu karisma,
daya tarik, kapasitas, dan daya saing. Strategi ini diharapkan
menyamakan limpahan penduduk dan industri, tetapi artikel menyatakan
bahwa keduanya sering tidak cocok dan kota baru tetap menghadapi
persoalan swasembada (Section 4.1, hal. PDF 4-6).

\textbf{{[}Dilaporkan sumber{]}} Bukti aliran menuju kota primer
menunjukkan sekitar 330.000 orang melakukan komuter ke Beijing setiap
hari dari wilayah sekitar, setara 3\% total komuter Beijing; sumber yang
dikutip menyebut Langfang, Gu'an, Zhuozhou, dan tempat lain sebagai asal
utama. Dalam hubungan Shanghai, Kunshan dan Taicang disebut sebagai
sumber utama komuter lintas batas, masing-masing 74,8\% dan 12,7\%
menurut laporan tahun 2020 yang dirujuk artikel (Section 4.1, hal. PDF
5).

\textbf{{[}Dilaporkan sumber{]}} Hasil survei Taicang yang diakses
penulis memperlihatkan keterkaitan industri lintas batas: 671 dari 953
perusahaan, atau 70,4\%, terhubung dengan Shanghai di sisi hulu atau
hilir rantai industri. Proporsi manufaktur yang terkait disebut 76,9\%.
Untuk sektor jasa, artikel menyebut 50\% perusahaan terkait dengan
Shanghai dan 60\% dari kelompok yang terkait itu berada di perdagangan
dan logistik. Teks juga menyebut 623 dari 953 perusahaan memiliki
karyawan komuter, atau 65\%; 54\% komuter harian berasal dari Shanghai;
29,6\% komuter adalah manajer; dan 37,3\% teknisi. Denominator untuk
angka 54\% tidak diperjelas lebih lanjut dalam file (Section 4.1, hal.
PDF 5-6).

\textbf{{[}Dilaporkan sumber{]}} Pada tipe kawasan industri,
Jiangyin-Jingjiang Industrial Park didirikan pada 2003 di Jingjiang,
kota tetangga Jiangyin yang dipisahkan Sungai Yangtze. Jembatan dan
jalan raya nasional menjadi prasyarat fisik. Jiangyin memiliki
perkembangan ekonomi lebih awal melalui \emph{town and village
enterprises} dan reformasi kepemilikan saham, sedangkan Jingjiang lebih
terbelakang meskipun memiliki kondisi pelabuhan alami. Pemerintah
tingkat lebih tinggi, dengan panduan pemerintah Provinsi Jiangsu,
diperlukan untuk menjembatani ketimpangan kekuasaan antarotoritas
(Section 4.2, hal. PDF 6).

\textbf{{[}Dilaporkan sumber{]}} Modal terdaftar fase pertama kawasan
Jiangyin-Jingjiang adalah 100 juta yuan, dengan Jiangyin menyumbang 90\%
dan Jingjiang 10\%; pendapatan investasi selama 10 tahun tetap digunakan
untuk pengembangan bergulir kawasan. Strateginya ialah menggunakan lahan
dan pelabuhan Jingjiang serta memindahkan industri dari Jiangyin. Namun,
kerja sama kemudian mengalami dampak negatif: ekonomi Jiangyin terdampak
arus investasi asing, sedangkan Jingjiang menuduh Jiangyin sekadar
memindahkan industri yang tidak diinginkan dan menggunakan terlalu
banyak sumber daya lahannya. Inisiatif terhenti setelah dekade pertama
dan hanya mengalami kemajuan terbatas (Section 4.2, hal. PDF 6).

\textbf{{[}Dilaporkan sumber{]}} Artikel mencatat biaya relokasi
perusahaan yang tinggi, peningkatan biaya administrasi, lambatnya
pertumbuhan pendapatan, dan tekanan pengelolaan ketika tiba waktunya
membagi keuntungan. Setelah kegagalan tersebut, kota sekunder lebih
berfokus pada perluasan produksi yang sudah ada daripada relokasi
perusahaan. Suzhou-Chuzhou Industrial Park menjadi contoh perubahan ini
(Section 4.2, hal. PDF 6-7).

\textbf{{[}Dilaporkan sumber{]}} Suzhou-Chuzhou Industrial Park
dikembangkan ketika Sino-Singapore Suzhou Industrial Park Development
Group bekerja sama dengan pemerintah Chuzhou. Strateginya dijalankan
secara tidak langsung melalui otoritas zona pembangunan, bukan secara
langsung melalui pemerintah lokal kota maju. Informan menyatakan bahwa
70-80\% perusahaan yang ditarik berasal dari Delta Sungai Yangtze,
banyak di antaranya perusahaan privat kecil dan perusahaan asing. Sejak
2015, Komite Manajemen Chuzhou menarik sekitar dua pertiga perusahaan
tersebut, tetapi biaya agen pencarian investasi dibagi dengan
Sino-Singapore Park karena merek dan pengelolaannya tetap dipakai
(Interview, 19 November 2021; Section 4.2, hal. PDF 6-7).

\textbf{{[}Dilaporkan sumber{]}} Untuk perusahaan berteknologi tinggi,
informan menyatakan empat perlima investasi ditarik oleh Sino-Singapore
Park. Model ini meminjam merek dan layanan pengelolaan dari Suzhou,
termasuk pelatihan inkubator dan layanan inovasi. Aliran yang dihasilkan
mencakup komunikasi sektor publik, informasi, kader, dan sumber daya
bisnis, serta aliran rantai pasok antara perusahaan di kawasan gabungan
dan kota asal. Banyak perusahaan di kawasan asal disubsidi untuk membuka
pabrik cabang ketika perlu memperluas kapasitas produksi, sehingga
limpahan berlangsung melalui aliran yang tidak tampak dan bukan terutama
relokasi fisik (Section 4.2, hal. PDF 6-7).

\textbf{{[}Dilaporkan sumber{]}} Pada tipe enklave inovasi, 22 enklave
didirikan di Shanghai dan kota-kota lain di YRD dari 2015 sampai akhir
2023 menurut wawancara yang dirujuk artikel. Jumlahnya lebih tinggi di
Hangzhou karena pemerintah Provinsi Zhejiang lebih aktif mendorong kota
sekunder mendirikan enklave di ibu kota provinsi. Enklave biasanya hanya
menempati beberapa lantai ruang kantor, berfungsi sebagai inkubator dan
platform inovasi untuk sejumlah industri teknologi tinggi, berbeda dari
kawasan industri yang memerlukan lahan luas (Section 4.3, hal. PDF 7).

\textbf{{[}Dilaporkan sumber{]}} Industri yang ditargetkan mencakup
kecerdasan buatan, mesin dan peralatan, energi baru, manufaktur
berteknologi tinggi, serta industri medis dan farmasi. Harapannya,
ketika teknologi dan bisnis matang, pabrik akan berada di kota sekunder
dan mempekerjakan tenaga kerja dari kota tersebut. Dengan demikian, kota
sekunder meminjam skala talenta, pengetahuan teknologi, dan layanan dari
metropolitan (Section 4.3, hal. PDF 7).

\textbf{{[}Dilaporkan sumber{]}} Kasus Changxing menunjukkan tekanan
untuk menarik industri yang lebih maju dan keterbatasan pemerintah
kabupaten dalam menandingi sumber daya politik, finansial, kebijakan,
dan regulasi kota terkemuka. Penelusuran pejabat lokal mengarahkan
Changxing pada ruang dekat Zhejiang University di Hangzhou; hubungan
sebelumnya antara pimpinan Changxing dan Hangzhou disebut memudahkan
penyewaan lahan terlantar serta pembentukan inkubator. Kasus Quzhou
menunjukkan penggunaan ruang di Hangzhou karena talenta tingkat tinggi
lebih sulit menetap di Quzhou dan diharapkan tetap dapat berkontribusi
pada Quzhou (Interviews, 23 Juli 2024 dan 13 Desember 2021; Section 4.3,
hal. PDF 7).

\textbf{{[}Dilaporkan sumber{]}} Enklave inovasi dibingkai melalui
metafora menarik ``phoenix'' dari teknologi, manufaktur, talenta, dan
pengusaha yang lebih maju di metropolitan ke ``nest'' kota sekunder.
Namun, kota primer dapat bersikap konservatif atau memusuhi enklave
tersebut karena melihatnya sebagai \emph{waqiangjiao}, yaitu pencurian
sumber daya teknologi dan talenta (Section 4.3, hal. PDF 7).

\textbf{{[}Interpretasi penulis{]}} Tiga tipe kasus menunjukkan bahwa
strategi kota sekunder memiliki lintasan berbeda: kota baru
mengintegrasikan kota sekunder secara fisik dan fungsional dengan kota
primer; kawasan industri mengatur kerja sama, ekspansi produksi, dan
aliran rantai pasok; enklave inovasi mengakses talenta, teknologi, dan
fungsi inovasi dari metropolitan. Limpahan dipahami sebagai sesuatu yang
diorkestrasi melalui proyek, aktor, merek, infrastruktur, dan hubungan
antarkewenangan, bukan sesuatu yang muncul spontan (Sections 4 dan 5,
hal. PDF 4-8).

\subsection{Findings}\label{findings-42}

\textbf{{[}Interpretasi penulis{]}} Temuan inti artikel ialah bahwa kota
sekunder adalah agen aktif dalam regionalisme kota. Kota-kota tersebut
mengidentifikasi ceruk yang dapat dieksploitasi untuk memperoleh
\emph{borrowed size} atau \emph{borrowed function}, meredam efek
bayangan kota primer, dan membentuk ekonomi antartempat melampaui batas
administratif (Sections 1, 4, dan 5, hal. PDF 1-2 dan 4-8).

\textbf{{[}Interpretasi penulis{]}} \emph{Borrowed function} paling
jelas pada kota baru dan enklave inovasi. Kota baru meminjam citra,
kedekatan, standar layanan, dan konektivitas kota primer. Enklave
inovasi meminjam talenta, pengetahuan teknologi, dan fungsi industri
dengan harapan produksi matang kelak berpindah atau terhubung ke kota
sekunder (Sections 4.1 dan 4.3, hal. PDF 4-7).

\textbf{{[}Interpretasi penulis{]}} \emph{Borrowed size} pada kawasan
industri tidak hanya berarti memindahkan massa industri dari inti ke
perifer. Pada Suzhou-Chuzhou, ukuran dan fungsi dipinjam melalui merek,
layanan pengelolaan, jaringan perusahaan, perluasan produksi, dan
hubungan rantai pasok. Pada Jiangyin-Jingjiang, strategi tersebut
memperlihatkan bahwa ketimpangan kepentingan, pembagian keuntungan,
biaya relokasi, dan tuduhan pemindahan industri yang tidak diinginkan
dapat menghentikan kerja sama (Section 4.2, hal. PDF 6-7).

\textbf{{[}Interpretasi penulis{]}} Batas yurisdiksi bukan sekadar
residu dari skema klasifikasi. Batas menjadi lokasi tempat pemerintah,
perusahaan, dan warga menegosiasikan akses terhadap lahan, layanan,
talenta, industri, pendapatan, dan mobilitas. Proyek-proyek lintas batas
karena itu membentuk ruang antarkota yang dapat berupa koridor atau
jaringan node yang bersaing sekaligus bekerja sama (Sections 2.2 dan 5,
hal. PDF 3 dan 7-8).

\textbf{{[}Interpretasi penulis{]}} Bukti dari Tiongkok diposisikan
berbeda dari praktik pembuangan penggunaan lahan yang tidak diinginkan
kepada yurisdiksi tetangga yang sering dibahas dalam konteks Barat.
Dalam kasus yang diteliti, otoritas kota sekunder justru secara aktif
membangun objek batas untuk menarik atau mengarahkan aliran penduduk,
industri, talenta, serta pengetahuan. Artikel tetap mengakui bahwa
kepentingan kota primer, sektor privat, dan warga tidak selalu sejalan
(Conclusion, hal. PDF 8).

\textbf{{[}Inferensi review{]}} Kasus-kasus tersebut mendukung
plausibilitas mekanisme strategisasi, tetapi tidak memungkinkan
penilaian tentang apakah setiap proyek meningkatkan produktivitas,
kesejahteraan, atau pemerataan secara bersih. Hasil yang terlihat juga
dapat dipengaruhi oleh dorongan kota primer, perusahaan, kebijakan
tingkat lebih tinggi, dan kondisi rantai nilai; artikel sendiri
menyatakan pengaruh relatif strategi yang disengaja dan limpahan yang
tidak disengaja masih belum jelas (Conclusion, hal. PDF 8).

\subsection{Conclusions}\label{conclusions-42}

\textbf{{[}Interpretasi penulis{]}} Artikel menyimpulkan bahwa studi
regionalisme kota perlu memberi perhatian lebih besar pada agensi kota
sekunder dan pada imajinasi kawasan kota yang mereka orkestrasi. Tiga
mode pembangunan lintas batas memperlihatkan cara berbeda untuk
merespons regionalisme yang berpusat pada metropolitan dan aglomerasi:
integrasi fungsional melalui kota baru, peminjaman ukuran dan fungsi
industri melalui kawasan industri, serta peminjaman talenta dan fungsi
teknologi melalui enklave inovasi (Section 5, hal. PDF 7-8).

\textbf{{[}Interpretasi penulis{]}} Kawasan kota di Tiongkok digambarkan
berkembang menjadi ruang koridor megalopolitan atau urban-regional yang
melampaui garis administratif dan semakin menyerupai mosaik atau
arkipelago node yang saling terhubung, bersaing, dan bekerja sama.
Kesimpulan ini berasal dari pembacaan atas kasus, bukan dari pengukuran
seluruh kawasan kota di Tiongkok (Section 5, hal. PDF 7-8).

\textbf{{[}Dilaporkan sumber{]}} Studi kasus menunjukkan bahwa limpahan
tidak terjadi secara otonom seperti yang diasumsikan teori aglomerasi
dan \emph{trickle-down}. Kota sekunder tidak menunggu secara pasif,
tetapi berupaya membangun jalur perkembangan melalui perencanaan
urban-ekonomi dan proyek lintas batas. Penulis menyatakan bahwa
penelitian ini berfokus pada keberadaan, logika, dan manifestasi
strategisasi, bukan pada penimbangan efeknya (Conclusion, hal. PDF 8).

\textbf{{[}Inferensi review{]}} Kesimpulan yang dapat
dipertanggungjawabkan dari TXT adalah kesimpulan mengenai proses,
agensi, dan variasi strategi. Tidak ada dasar dalam file untuk
menyimpulkan bahwa ketiga tipe proyek secara umum berhasil, mengurangi
ketimpangan, atau menghasilkan manfaat yang lebih besar daripada biaya
dan konflik.

\subsection{Contributions}\label{contributions-42}

\textbf{{[}Dilaporkan sumber{]}} Artikel secara eksplisit mengklaim dua
kontribusi. Pertama, artikel memusatkan perhatian pada perjuangan
strategis otoritas kota sekunder untuk menemukan ceruk dalam organisasi
ekonomi dan tata kelola kawasan kota. Kedua, dari konteks Tiongkok,
artikel menunjukkan mobilisasi peluang pembangunan lintas batas sebagai
bagian dari strategisasi otoritas kota sekunder (Section 1, hal. PDF
1-2).

\textbf{{[}Interpretasi penulis{]}} Secara konseptual, perspektif
\emph{outside-in} memperluas kajian regionalisme kota yang sebelumnya
lebih sering dibaca dari kota primer, negara, atau aktor yang bersekutu
dengan kota primer. Artikel menghubungkan \emph{borrowed size},
\emph{borrowed function}, pembagian kerja urban, batas yurisdiksi, dan
agensi perencanaan dalam satu pembacaan strategis (Sections 1-3, hal.
PDF 1-4).

\textbf{{[}Dilaporkan sumber{]}} Secara empiris, artikel menyediakan
tipologi tiga jenis proyek dan narasi enam contoh dari BTH serta YRD.
Tipologi itu mengaitkan masing-masing proyek dengan motif dan mekanisme
berbeda, sehingga memperlihatkan keragaman profil, lintasan, tantangan,
dan peluang kota sekunder (Table 2 dan Conclusion, hal. PDF 4 dan 8).

\textbf{{[}Inferensi review{]}} Sumbangan paling kuat artikel adalah
pemetaan mekanisme dan aktor yang sering hilang jika kawasan kota
diperlakukan sebagai unit homogen. Pengukuran komparatif dampak dan
prosedur analisis kualitatif yang rinci adalah \textbf{Tidak disebutkan
dalam file}.

\subsection{Research Gap}\label{research-gap-42}

\textbf{{[}Dilaporkan sumber{]}} Celah yang hendak diisi adalah
kurangnya perhatian terhadap agensi kota sekunder, strategi mereka dalam
pembentukan kawasan kota, dan cara mereka merespons efek aglomerasi
serta marginalisasi. Literatur sebelumnya dinilai lebih sering
menempatkan kota primer, negara, atau sektor privat sebagai penggerak
utama regionalisme kota (Sections 1 dan 2, hal. PDF 1-3).

\textbf{{[}Dilaporkan sumber{]}} Artikel juga menyatakan bahwa
konsekuensi skala dan kompleksitas perifer urban, bentuk strategisasi
politik-birokratis serta ekonomi yang berpusat pada batas yurisdiksi,
dan kemunculan ruang antarkota belum cukup dielaborasi. Dalam konteks
Tiongkok, kesulitan menetapkan batas perifer diperbesar oleh pembedaan
administratif yang terutama membedakan ruang urban dan rural (Section
2.2 dan Section 3, hal. PDF 3-4).

\textbf{{[}Interpretasi penulis{]}} Dengan menghadirkan kasus kota
sekunder, artikel berupaya memperluas teori regionalisme kota dari
perspektif \emph{outside-in}. Fokusnya bukan hanya pada dampak limpahan,
tetapi pada proses politik, kelembagaan, relasional, dan spasial yang
membuat limpahan dapat dimobilisasi melalui objek batas (Sections 2.2,
4, dan 5, hal. PDF 3-8).

\textbf{{[}Dilaporkan sumber{]}} Gap yang masih terbuka menurut penutup
artikel mencakup pengaruh strategi kota sekunder terhadap hasil
pembangunan, perbedaan antara efek strategi yang disengaja dan limpahan
kota primer yang tidak disengaja, serta perbandingan strategi kota
sekunder secara keseluruhan dan strategi yang berpusat pada batas
yurisdiksi (Conclusion, hal. PDF 8).

\textbf{{[}Inferensi review{]}} Gap empiris yang tersisa juga mencakup
pengujian lintas kasus dengan indikator hasil yang konsisten, penjelasan
mekanisme seleksi dan kegagalan proyek, serta evaluasi siapa yang
memperoleh atau menanggung manfaat. Kebutuhan ini diturunkan dari
batasan yang dinyatakan artikel dan dari absennya ukuran outcome dalam
TXT, bukan dari literatur eksternal.

\subsection{Limitations}\label{limitations-42}

\textbf{{[}Dilaporkan sumber{]}} Artikel mengakui bahwa fokusnya adalah
mengelusidasi kehadiran, logika, dan manifestasi strategisasi kota
sekunder, bukan menimbang efek strategisasi tersebut. Karena itu, temuan
tidak dimaksudkan sebagai evaluasi dampak pembangunan, pemerataan,
kesejahteraan, atau efisiensi ekonomi (Conclusion, hal. PDF 8).

\textbf{{[}Dilaporkan sumber{]}} Definisi kota sekunder sengaja longgar
dan bukan kerangka kategorisasi yang teliti. Konsep kawasan kota sendiri
disebut kabur, sedangkan peripherality urban dan objek batas sulit
didefinisikan dalam sistem urban Tiongkok. Kesulitan tersebut
memengaruhi ketegasan batas populasi dan klasifikasi kasus (Section 3,
hal. PDF 3-4).

\textbf{{[}Dilaporkan sumber{]}} Cakupan empiris terutama berada di BTH
dan YRD, dengan lebih dari 10 lokasi yang dikunjungi dan enam contoh
yang ditampilkan (Sections 3-4, hal. PDF 3-4).

\textbf{{[}Inferensi review{]}} Pernyataan bahwa kasus tersebut mencakup
seluruh jenis kota sekunder atau seluruh variasi regionalisme kota di
Tiongkok adalah \textbf{Tidak disebutkan dalam file}.

\textbf{{[}Inferensi review{]}} Pemilihan kasus untuk merepresentasikan
tiga tipe dan penggunaan \emph{snowball sampling} untuk informan
membatasi keterwakilan statistik. Jumlah kota sekunder yang menjadi
kerangka populasi, kriteria pemilihan dari lebih dari 10 lokasi,
distribusi informan per kasus, dan dasar saturasi adalah \textbf{Tidak
disebutkan dalam file}.

\textbf{{[}Inferensi review{]}} Artikel tidak menjelaskan prosedur
pengodean, audit transkrip, aturan perbandingan kasus, atau cara
menggabungkan bukti wawancara dengan dokumen. Hasilnya kaya secara
naratif, tetapi keterlacakan analisis dari bahan mentah menuju
kesimpulan tidak dapat diperiksa sepenuhnya dari TXT.

\textbf{{[}Dilaporkan sumber{]}} Jenis proyek dapat beririsan secara
material: enklave inovasi dapat ditumpangkan pada atau dipahat dari
kawasan industri lintas batas dan kota baru sebagai bagian dari evolusi
kebijakan kawasan ekonomi khusus (Section 2.2, hal. PDF 3).

\textbf{{[}Inferensi review{]}} TXT tidak menyajikan uji klasifikasi
yang menentukan apakah setiap proyek hanya termasuk satu tipe.

\textbf{{[}Inferensi review{]}} Ada ketidakselarasan temporal yang perlu
dicatat. Section 3 menyebut lebih dari 20 kunjungan lapangan pada
2008-2023, sedangkan Table 1 memuat wawancara Changxing pada 22 dan 23
Juli 2024. TXT tidak menjelaskan apakah wawancara 2024 berada di luar
rentang kunjungan lapangan atau merupakan pembaruan terpisah.

\textbf{{[}Inferensi review{]}} Pernyataan \emph{Data availability} yang
menyebut tidak ada data digunakan tidak mudah dipadukan dengan uraian 15
wawancara, enam simposium, dokumen internal, dan hasil survei Taicang.
Review tidak mengasumsikan salah satu bagian lebih benar; keduanya
dipertahankan sebagai laporan sumber yang belum direkonsiliasi.

\subsection{Challenges}\label{challenges-42}

\textbf{{[}Dilaporkan sumber{]}} Kota sekunder menghadapi posisi kurang
menguntungkan ketika harus menarik perusahaan, industri maju, talenta,
dan penduduk dari kota primer yang memiliki sumber daya politik,
finansial, kebijakan, serta regulasi lebih besar. Tekanan evaluasi
pembangunan dan kesulitan meningkatkan industri tradisional mendorong
pencarian jalur alternatif seperti enklave inovasi (Section 4.3, hal.
PDF 7).

\textbf{{[}Dilaporkan sumber{]}} Proyek kota baru menghadapi tantangan
untuk menyelaraskan limpahan penduduk dengan limpahan industri. Penduduk
dapat tertarik oleh perumahan dan layanan, tetapi tidak selalu bekerja
pada industri yang ditarik. Ketergantungan pemerintah lokal pada
penjualan lahan untuk membiayai penarikan industri juga membuat ledakan
properti tidak otomatis menghasilkan aliran pendapatan jangka panjang
(Section 4.1, hal. PDF 5-6).

\textbf{{[}Dilaporkan sumber{]}} Kawasan industri lintas batas
menghadapi perbedaan insentif dan kekuasaan antarotoritas, kebutuhan
intervensi pemerintah tingkat lebih tinggi, biaya relokasi, pembagian
biaya serta keuntungan, dan konflik tentang apakah industri yang
dipindahkan benar-benar menguntungkan kota penerima. Dalam kasus
Jiangyin-Jingjiang, konflik tersebut menyebabkan inisiatif terhenti
(Section 4.2, hal. PDF 6).

\textbf{{[}Dilaporkan sumber{]}} Kerja sama berbasis merek dan layanan
juga menghadapi tantangan atribusi dan ketergantungan. Pada
Suzhou-Chuzhou, pengaruh pemerintah Chuzhou sulit dipisahkan dari merek
Sino-Singapore Park, pengelola zona, penggerak perusahaan, dan hubungan
investasi antarzona. Artikel menyebut hubungan korporasi dan inovasi di
seluruh YRD sebagai faktor yang juga penting (Section 4.2, hal. PDF
6-7).

\textbf{{[}Dilaporkan sumber{]}} Enklave inovasi menghadapi kesulitan
membuat talenta tingkat tinggi menetap di kota sekunder, keterbatasan
pengetahuan lokal, dan sikap kota primer yang dapat melihat enklave
sebagai pencurian sumber daya teknologi serta talenta. Strategi enklave
juga bergantung pada harapan bahwa perusahaan, teknologi, dan produksi
kelak terhubung kembali dengan kota sekunder (Section 4.3, hal. PDF 7).

\textbf{{[}Dilaporkan sumber{]}} Tantangan konseptual dan tata kelola
meliputi penetapan batas perifer, pengurangan hambatan
antaradministrasi, pembangunan hubungan yang tidak setara, dan
pengelolaan kawasan kota yang mencakup banyak tingkat pemerintahan.
Batas yang sama dapat menjadi penghalang aliran atau peluang untuk
mengatur hubungan antartempat (Sections 2.2 dan 3, hal. PDF 3-4).

\textbf{{[}Inferensi review{]}} Tantangan penelitian yang paling
menonjol adalah membedakan efek strategi kota sekunder dari dorongan
kota primer, kepentingan perusahaan, kebijakan tingkat lebih tinggi, dan
limpahan yang tidak direncanakan. Artikel mengidentifikasi masalah ini,
tetapi strategi identifikasi untuk menyelesaikannya adalah \textbf{Tidak
disebutkan dalam file}.

\subsection{Practical Implications}\label{practical-implications-42}

\textbf{{[}Interpretasi penulis{]}} Regionalisme kota dan kebijakan
pembangunan tidak seharusnya hanya dipahami sebagai ekspansi kota primer
atau penantian limpahan. Perencanaan perlu mengenali kota sekunder
sebagai aktor yang dapat membentuk jalur pengembangan, membangun ceruk,
dan mengorkestrasi aliran lintas batas (Sections 1, 4, dan 5, hal. PDF
1-2 dan 4-8).

\textbf{{[}Inferensi review berbasis bukti sumber{]}} Untuk kota baru,
implikasi praktisnya ialah menggabungkan konektivitas transportasi
dengan peningkatan layanan, fasilitas, pendidikan, kesehatan, dan
kapasitas ekonomi. Infrastruktur saja tidak cukup jika aliran penduduk
tidak diikuti kesempatan kerja dan basis industri yang sesuai (Section
4.1, hal. PDF 4-6).

\textbf{{[}Inferensi review berbasis bukti sumber{]}} Untuk kawasan
industri, kerja sama perlu memperjelas pembagian peran, pembagian
pendapatan, biaya relokasi, tanggung jawab pengelolaan, dan standar
layanan. Kasus Jiangyin-Jingjiang menunjukkan bahwa pembangunan fisik
dan pemindahan industri tanpa keselarasan insentif dapat memperbesar
konflik, sedangkan Suzhou-Chuzhou menunjukkan pentingnya aliran merek,
layanan, dan ekspansi produksi (Section 4.2, hal. PDF 6-7).

\textbf{{[}Inferensi review berbasis bukti sumber{]}} Untuk enklave
inovasi, desain kebijakan dapat memusatkan perhatian pada akses talenta,
kedekatan dengan universitas, inkubasi, hubungan perusahaan, dan jalur
yang menghubungkan kematangan teknologi dengan produksi di kota
sekunder. Namun, bukti bahwa desain tersebut telah menghasilkan
perpindahan talenta atau produksi secara luas adalah \textbf{Tidak
disebutkan dalam file} (Section 4.3, hal. PDF 7).

\textbf{{[}Interpretasi penulis{]}} Batas yurisdiksi perlu diperlakukan
sebagai ruang strategis untuk kerja sama dan pembentukan ekonomi
antartempat, bukan hanya garis administratif. Pencapaian tersebut
menuntut keterlibatan berbagai tingkat pemerintahan, aktor privat, dan
warga serta perhatian pada ketimpangan kekuasaan di antara kota primer
dan sekunder (Sections 2.2, 3, dan 5, hal. PDF 3-4 dan 7-8).

\subsection{Applications}\label{applications-42}

\textbf{{[}Dilaporkan sumber{]}} Kerangka artikel diterapkan secara
empiris pada enam proyek di dua kawasan utama Tiongkok. Gu'an New Town
merepresentasikan kota sekunder di tepi Langfang yang mengaitkan diri
dengan Beijing; Huaqiao New Town merepresentasikan hubungan
Kunshan-Shanghai; Jiangyin-Jingjiang merepresentasikan kawasan industri
yang dipisahkan Sungai Yangtze; Suzhou-Chuzhou merepresentasikan
peminjaman merek dan layanan kawasan industri; Changxing
merepresentasikan enklave inovasi di Shanghai; dan Quzhou
merepresentasikan enklave inovasi di Hangzhou (Fig. 1, Table 2, dan
Sections 4.1-4.3, hal. PDF 4-7).

\textbf{{[}Dilaporkan sumber{]}} Penerapan kerangka juga tampak dalam
pembedaan tiga objek batas dan tiga logika strategisasi: meminjam fungsi
pada kota baru, meminjam ukuran atau massa industri pada kawasan
industri, serta meminjam ukuran talenta dan fungsi industri pada enklave
inovasi. Contoh-contoh tersebut digunakan untuk menerangkan hubungan
antara perencanaan, aliran ekonomi, konektivitas, dan batas
administratif (Table 2 dan Section 4, hal. PDF 4-7).

\textbf{{[}Dilaporkan sumber{]}} Artikel menyebut bahwa kondisi Tiongkok
tidak harus diperlakukan sebagai sesuatu yang tak dapat dibandingkan
dengan Barat, melainkan dapat memperluas spektrum dinamika sistem kota,
pembentukan kawasan kota global dan megalopolitan, strategisasi suburban
serta kota sekunder, dan politik batas yurisdiksi (Conclusion, hal. PDF
8).

\textbf{{[}Inferensi review{]}} Penerapan kebijakan di luar kasus
Tiongkok dan evaluasi implementasi yang mengukur keberhasilan tiap
proyek adalah \textbf{Tidak disebutkan dalam file}. Penerapan lintas
negara, penggunaan kerangka sebagai alat prediksi, dan replikasi
kuantitatif juga \textbf{Tidak disebutkan dalam file}.

\subsection{Future Research
(Optional)}\label{future-research-optional-42}

\textbf{{[}Dilaporkan sumber{]}} Agenda riset yang dinyatakan dalam
penutup artikel meliputi:

\begin{itemize}
\tightlist
\item
  meneliti agensi kota sekunder dalam pembentukan kawasan kota dan
  mengembangkan lensa relasional yang tidak hanya melihat hubungan
  global-lokal, tetapi juga dinamika kompleks di dalam kawasan kota;
\item
  menempatkan perencanaan kawasan kota sebagai agenda utama untuk
  memahami hasil potensial pembangunan sebagai proses ekonomisasi yang
  dibentuk geopolitik global, hubungan internasional, dan politik
  antartingkat kota serta aktor lain dalam negara;
\item
  meneliti lebih lanjut perbedaan antara strategi kota sekunder yang
  disengaja dan efek limpahan kota primer yang tidak disengaja;
\item
  melakukan analisis komparatif atas strategisasi kota sekunder secara
  keseluruhan dan strategisasi yang berpusat pada batas yurisdiksi;
\item
  membandingkan mode pembentukan kota, kawasan kota global, dan kawasan
  megalopolitan, serta kedalaman strategisasi suburban dan kota sekunder
  dalam kondisi Tiongkok dan Barat;
\item
  mengkaji potensi kota kecil dan menengah sebagai penemu jalur
  pembangunan dalam bidang kesejahteraan dan \emph{catch-up
  development}, sekaligus sebagai aktor politik yang mendorong
  representasi, suara, dan kekuasaan lobi yang relevan (Conclusion, hal.
  PDF 8).
\end{itemize}

\textbf{{[}Inferensi review{]}} Rencana kerja, periode, desain sampel,
atau metode khusus untuk menjalankan agenda tersebut adalah
\textbf{Tidak disebutkan dalam file}.

\subsection{Provenance Notes}\label{provenance-notes-42}

\begin{itemize}
\tightlist
\item
  \textbf{Bahan utama:}
  \texttt{/Users/mac/Documents/Mac/{[}2{]}\ Obsidian\ Vault/zettelkasten/source/journal/A27-li-phelps-xu-2026-secondary-city-regionalism-chinas-cross-boundary-development-projects-10.1016-j.cities.2026.106795.txt}.
\item
  \textbf{Verifikasi kode dan nama file:} Prefiks \texttt{A27} dan judul
  pada nama file cocok dengan judul artikel di PDF Metadata, yaitu
  \emph{Secondary city regionalism: China's cross boundary development
  projects}. Hanya file TXT A27 ini yang digunakan sebagai bukti
  substantif; PDF hanya diidentifikasi melalui metadata yang tercantum
  dalam TXT.
\item
  \textbf{Cakupan pembacaan:} Seluruh TXT dibaca dari baris 1 sampai
  764, termasuk \texttt{{[}PDF\ Metadata{]}},
  \texttt{{[}Extracted\ Text{]}}, abstrak, Sections 1-5, Tables 1-2,
  Figs. 1-2, kutipan wawancara, pernyataan etika, pendanaan,
  ketersediaan data, dan seluruh referensi.
\item
  \textbf{Metadata ekstraksi:} TXT mencatat
  \texttt{Extracted\ on:\ 2026-08-31} dan
  \texttt{Extraction\ method:\ pdftotext\ -layout}. Metadata PDF
  mencatat 10 halaman, PDF versi 1.7, ukuran halaman 595.276 x 793.701
  points, serta DOI \texttt{10.1016/j.cities.2026.106795}.
\item
  \textbf{Locator:} Penyebutan \texttt{hal.\ PDF} mengikuti halaman
  PDF/ekstraksi yang terlihat dalam TXT; locator Section, Table, Fig.,
  tanggal wawancara, dan kutipan laporan dipertahankan agar klaim dapat
  ditelusuri dalam file sumber. Artikel tidak memiliki nomor halaman
  jurnal berkelanjutan karena menggunakan nomor artikel 106795.
\item
  \textbf{Batas akses:} Review hanya menggunakan teks TXT hasil
  ekstraksi, bukan sumber web, PDF lain, dataset mentah, transkrip, atau
  referensi eksternal. Pemenggalan kolom dan baris dari
  \texttt{pdftotext\ -layout} dipertahankan sebagai batas teknis; detail
  visual yang tidak diterangkan caption atau narasi tidak ditafsirkan
  ulang.
\item
  \textbf{Keterbatasan provenance penting:} Section 3 menyebut kunjungan
  lapangan 2008-2023, sedangkan Table 1 memuat wawancara Changxing pada
  22 dan 23 Juli 2024. TXT tidak menjelaskan hubungan kedua rentang
  waktu tersebut.
\item
  \textbf{Keterbatasan provenance penting:} Bagian \emph{Data
  availability} menyatakan ``No data was used for the research described
  in the article'', sementara isi artikel melaporkan 15 wawancara, enam
  simposium, kunjungan lapangan, dokumen, dan hasil survei Taicang yang
  tidak tersedia untuk publik. Review mempertahankan kedua laporan
  tersebut dan tidak mengarang resolusinya.
\item
  \textbf{Etika dan anonimisasi:} Artikel menyatakan tidak ada dokumen
  persetujuan etika formal yang ditandatangani, tetapi tujuan penelitian
  disampaikan, izin merekam dan mengutip diminta, peran profesional
  boleh disebut, dan nama informan dianonimkan (Ethical statement, hal.
  PDF 8).
\item
  \textbf{Pemisahan epistemik:} Pernyataan berlabel
  \texttt{{[}Dilaporkan\ sumber{]}} merangkum isi eksplisit TXT;
  \texttt{{[}Interpretasi\ penulis{]}} menyajikan cara artikel membaca
  bukti; \texttt{{[}Inferensi\ review{]}} adalah penilaian terbatas yang
  diturunkan dari TXT dan tidak diperlakukan sebagai temuan penulis.
\item
  \textbf{Informasi yang tidak tersedia:} Ukuran populasi kota sekunder,
  kriteria eksklusi kasus yang lengkap, prosedur pengodean, data mentah,
  instrumen wawancara, ukuran efek, dan evaluasi outcome adalah
  \textbf{Tidak disebutkan dalam file}; variabel dependen formal adalah
  \textbf{Tidak berlaku}.
\item
  \textbf{Audit status:} Review ini memuat tepat 19 heading wajib, tanpa
  tabel, dan klaim substantifnya diverifikasi terhadap seluruh 764 baris
  TXT.
\item
  \textbf{Perubahan file:} Review A27 diaudit dan diperbarui. TXT
  sumber, dokumen induk, serta file untuk kode lain tidak diubah; hanya
  checkbox A27 di \texttt{Todo.md} dan
  \texttt{output/kajian/kode-kajian-pendahuluan.md} yang disinkronkan.
\end{itemize}

\section{Review Kode A28}\label{review-kode-a28}

\textbf{Identitas sumber}

\begin{itemize}
\tightlist
\item
  Judul: \emph{Small Firms, Borrowed Size and the Urban-Rural Shift}.
\item
  Penulis: N. A. Phelps, R. J. Fallon, dan C. L. Williams.
\item
  Publikasi: \emph{Regional Studies}, volume 35, nomor 7, halaman
  613-624, tahun 2001.
\item
  DOI: 10.1080/00343400120075885.
\item
  Afiliasi yang tercantum: School of Geography, University of Leeds;
  Environment and Regeneration Department, Planning Division, Luton
  Borough Council; dan CB Hillier Parker.
\item
  Riwayat naskah yang tercantum: diterima Februari 1997 dan direvisi
  September 2000. Blok penerbit juga mencatat ``Published online: 18 Aug
  2010''; TXT tidak menjelaskan hubungan antara tanggal daring tersebut
  dan terbitan jurnal tahun 2001.
\item
  Jenis sumber: artikel jurnal. Status peer-review formal: Tidak
  disebutkan dalam file; acknowledgement menyebut tiga anonymous
  referees (p.~623, TXT lines 802-804).
\item
  Dasar review: seluruh teks yang tersedia dalam TXT A28, termasuk blok
  metadata hasil ekstraksi PDF.
\item
  Status audit: Lulus. Review dicocokkan ulang dengan satu TXT aktual
  A28; batas sumber/ekstraksi tetap dicatat dan dikualifikasi pada
  \texttt{Provenance\ Notes}, tetapi tidak menghalangi penyelesaian
  review.
\end{itemize}

\subsection{Summarized Abstract}\label{summarized-abstract-43}

Laporan sumber dalam abstrak menyatakan bahwa sebagian unsur pergeseran
pembentukan dan pertumbuhan perusahaan dari wilayah urban ke rural
mungkin berkaitan dengan fenomena \emph{borrowed size}. Konsep ini
menghubungkan literatur tentang geografi pembentukan dan pertumbuhan
perusahaan kecil, teori aglomerasi, serta hubungan antarpemukiman dalam
sistem urban. Perusahaan kecil dapat berlokasi di pemukiman yang lebih
kecil, tetapi tetap mengakses tenaga kerja khusus dan eksternalitas
informasi dari kawasan urban yang lebih besar di dekatnya (abstrak,
p.~613; TXT lines 95-108).

Laporan sumber menyebutkan bahwa artikel menguji relevansi sebagian
aspek \emph{borrowed size} pada hubungan perusahaan kecil di Leighton
Linslade dan Lichfield dengan ekonomi lokal serta pusat urban yang lebih
besar di sekitarnya. Artikel secara khusus mempertanyakan kontribusi
dukungan kelembagaan lokal terhadap pertumbuhan perusahaan rural, lalu
membahas implikasinya bagi pemahaman tentang ekonomi lokal dan kebijakan
ekonomi lokal (abstrak, p.~613; TXT lines 105-108).

\textbf{Interpretasi penulis:} \emph{Borrowed size} membantu menjelaskan
bagaimana perusahaan memperoleh manfaat eksternalitas urban tanpa harus
berlokasi di inti urban. \textbf{Inferensi review:} artikel memosisikan
pergeseran urban-rural sebagai proses relasional antarlokasi, bukan
sekadar perpindahan aktivitas dari kota ke desa atau keberhasilan yang
hanya ditentukan oleh kapasitas ekonomi lokal.

\subsection{Problem Statement}\label{problem-statement-43}

Laporan sumber menggambarkan masalah awal berupa ketegangan antara
pandangan konvensional dan bukti empiris. Pandangan konvensional
mengaitkan perusahaan kecil dengan ekonomi eksternal yang terlokalisasi
dan menganggap kota besar sebagai lingkungan yang paling sesuai untuk
pembentukan serta pertumbuhan perusahaan kecil. Namun, pergeseran
urban-rural dalam pembentukan dan pertumbuhan perusahaan kecil, serta
temuan tentang lemahnya integrasi perusahaan kecil dengan ekonomi lokal,
menantang anggapan tersebut (pendahuluan, pp.~613-614).

Artikel juga menempatkan kebijakan pengembangan ekonomi lokal sebagai
bagian dari masalah. Kebijakan yang dirujuk dalam artikel menekankan
pembangunan jaringan antarfirma dan penguatan atau koordinasi institusi
lokal. Penulis berpendapat bahwa pendekatan tersebut dapat gagal
mengenali sifat relasional ekonomi eksternal dalam sistem urban, karena
manfaat perusahaan tidak selalu berasal dari pemukiman tempat perusahaan
berada (pendahuluan, pp.~613-614).

\textbf{Inferensi review:} Tidak disebutkan dalam file satu pertanyaan
penelitian formal yang ditulis dalam bentuk kalimat tunggal. Secara
operasional, masalahnya adalah apakah perusahaan kecil di pemukiman
rural yang mudah diakses dapat memanfaatkan ukuran ekonomi pusat urban
yang lebih besar, serta seberapa penting ekonomi lokal dan dukungan
institusi lokal dibandingkan hubungan dengan pusat-pusat di luar lokasi
langsung.

\subsection{Objectives}\label{objectives-43}

Laporan sumber menyatakan bahwa penelitian ini bertujuan mengeksplorasi
relevansi \emph{borrowed size} bagi dinamika lokasi perusahaan kecil di
wilayah rural yang mudah diakses (bagian \emph{Research Methodology},
p.~616). Tujuan yang dijalankan dalam artikel meliputi:

\begin{itemize}
\tightlist
\item
  Menghubungkan geografi pembentukan dan pertumbuhan perusahaan kecil
  dengan teori aglomerasi serta hubungan antarpemukiman dalam sistem
  urban (abstrak, p.~613).
\item
  Mengkaji hubungan perusahaan kecil di Leighton Linslade dan Lichfield
  dengan ekonomi lokal serta pusat urban besar di sekitarnya (abstrak,
  p.~613).
\item
  Menelusuri lokasi berbagai eksternalitas melalui hubungan
  input-output, tenaga kerja, dan keunggulan lokasi (metodologi,
  pp.~616-619).
\item
  Menelusuri sumber eksternalitas teknologis melalui informasi teknis,
  informasi pasar, informasi pemasok, kontak bisnis informal, dan lokasi
  pesaing (pp.~621-622, Tabel 7).
\item
  Memeriksa relevansi serta penggunaan penyedia layanan dan institusi
  lokal oleh perusahaan kecil (pp.~620-621, Tabel 6).
\item
  Menguji apakah pola di luar South East England, khususnya di
  Lichfield, juga konsisten dengan gagasan \emph{borrowed size} (catatan
  5, p.~623).
\end{itemize}

Tidak disebutkan dalam file sebagai hipotesis formal, tujuan prediksi
kausal, atau tujuan estimasi pengaruh. Inferensi review: ekspektasi
substantif penulis adalah bahwa aksesibilitas ke konurbasi besar dapat
lebih penting daripada kedekatan administratif dengan ekonomi lokal,
terutama untuk perusahaan jasa bisnis.

\subsection{Literature Survey}\label{literature-survey-43}

Laporan sumber memulai tinjauan dari gagasan bahwa perusahaan kecil
dengan sumber daya internal terbatas bergantung pada pasar properti,
tenaga kerja, dan input antara yang berfungsi baik, serta pada arus
pengetahuan dan keahlian yang terlokalisasi. Hoover dan Vernon digunakan
untuk mendukung pandangan historis bahwa kota besar menyediakan
lingkungan yang sesuai bagi pembentukan dan pertumbuhan perusahaan
kecil. Literatur produksi fleksibel dan aglomerasi, termasuk Piore dan
Sabel serta Scott, dikaitkan dengan perumusan kebijakan mengenai
perusahaan kecil dan pembangunan ekonomi lokal (pendahuluan,
pp.~613-614).

Artikel lalu menyajikan bukti yang menantang kebijakan yang terlalu
berpusat pada ekonomi lokal. Curran dan Blackburn dirujuk karena
berpendapat bahwa integrasi perusahaan kecil dengan ekonomi lokalnya
hampir tidak ada selain dalam arti simbolis. Keeble, Keeble dan Tyler,
serta karya terkait digunakan untuk menunjukkan adanya pergeseran
urban-rural dalam pembentukan dan pertumbuhan perusahaan kecil. North
dirujuk untuk pola konsentrasi aktivitas usaha kecil dan menengah di
wilayah rural yang berdekatan dengan kawasan metropolitan; Allen, Coe
dan Townsend, serta Hitchens dan rekan-rekan dipakai untuk menunjukkan
keterkaitan perusahaan di pemukiman kecil dengan London, Birmingham, dan
pusat urban lain (pendahuluan, pp.~614-615).

Konsep utama artikel berasal dari Alonso. Dalam kutipan yang dipaparkan
penulis, kota kecil di kawasan metropolitan dapat menunjukkan sebagian
karakter kota besar jika berada dekat dengan pusat populasi lain.
Keuntungan yang dapat dipertahankan di pemukiman kecil meliputi
kemacetan yang lebih rendah dan sewa yang lebih rendah, sementara
keuntungan kota besar meliputi pasar, layanan bisnis, keahlian, pasar
tenaga kerja yang lebih besar dan beragam, serta amenitas budaya
(pp.~614-615). Penulis mengaitkan gagasan ini dengan \emph{non-place
urban realm} Weber, teori sistem kota, dan literatur jaringan.

Posisi teoretis artikel membedakan eksternalitas pecuniary atau yang
berkaitan dengan ukuran pasar dari eksternalitas teknologis. Pembedaan
itu dipertautkan dengan tiga serangkai eksternalitas Marshallian:
transaksi input-output material, pengumpulan tenaga kerja, dan
eksternalitas teknologis. Penulis menyatakan bahwa peningkatan
transportasi, infrastruktur, dan proses produksi yang makin berputar
membuat sebagian eksternalitas pecuniary bekerja pada ruang yang luas,
sedangkan kontak antarpersonal yang menjadi dasar eksternalitas
teknologis tetap lebih terbatas secara geografis (pp.~615-616).

\textbf{Interpretasi penulis:} \emph{Borrowed size} paling mungkin
terlihat pada industri yang dapat memperoleh eksternalitas pecuniary dan
teknologis pada skala konurbasi atau kawasan urban yang sangat
terurbanisasi. Layanan bisnis dipilih karena dipandang bergantung pada
pasar, layanan keuangan dan kantor pusat, serta jaringan informasi yang
dapat berasal dari kota besar (pp.~615-616).

Inferensi review: tinjauan ini adalah landasan konseptual dan selektif
untuk studi kasus, bukan tinjauan sistematis. Protokol pencarian,
kriteria inklusi, dan korpus literatur yang lengkap tidak disebutkan
dalam file.

\subsection{Method Used}\label{method-used-43}

Laporan sumber menyebut desain penelitian sebagai studi empiris berskala
kecil, eksploratori, dan berbasis studi kasus. Studi pertama dilakukan
di Leighton Linslade pada September 1996 sebagai bagian dari disertasi
mahasiswa dan mencakup perusahaan manufaktur serta jasa bisnis. Studi
kedua dilakukan pada awal hingga pertengahan 1998 di Lichfield untuk
menguji relevansi \emph{borrowed size} di luar South East England;
replikasinya terhadap studi pertama hanya dapat dilakukan secara tidak
sempurna dan pada praktiknya berfokus pada perusahaan jasa bisnis
(catatan 5, p.~623).

Metode pengumpulan data utama adalah survei melalui telepon. Sebagian
perusahaan meminta kuesioner dikirim dan dikembalikan melalui pos atau
faks, dan pada studi Lichfield respons telepon serta respons pos/faks
kurang lebih berimbang. Pemilihan survei telepon dijelaskan oleh
kebutuhan untuk memperoleh informasi kuantitatif sekaligus menangkap
unsur kualitatif mengenai relevansi \emph{borrowed size} dengan sumber
daya yang terbatas (metodologi, pp.~616-617).

Kerangka sampel disusun dari direktori bisnis lokal dan nasional serta
Yellow Pages. Perusahaan dan tempat usaha didefinisikan sebagai
perusahaan kecil jika memiliki 100 atau lebih sedikit karyawan.
Kuesioner meminta informasi tentang lokasi geografis tenaga kerja, input
material dan teknis, penjualan, keunggulan dan kelemahan kedua lokasi,
relevansi serta penggunaan institusi lokal, sumber informasi bisnis, dan
pesaing (pp.~616-617, Tabel 1; pp.~619-622, Tabel 3-Tabel 7).

Analisis terdiri atas statistik deskriptif, perbandingan bivariat, uji
chi-kuadrat, analisis varians, dan korelasi peringkat Spearman. Penulis
menganalisis jenis eksternalitas satu per satu. Tidak ada regresi atau
teknik multivariat dalam hasil yang dilaporkan; penulis sendiri
menyatakan bahwa teknik multivariat mungkin akan memberi pemahaman lebih
baik tentang komplementaritas eksternalitas pecuniary dan teknologis
(kesimpulan, p.~622).

\textbf{Inferensi review:} Unit analisis yang tampak dari metode adalah
perusahaan atau tempat usaha kecil dalam sektor jasa bisnis di dua
pemukiman. Tidak disebutkan dalam file adanya panel, pelacakan
perusahaan yang sama sepanjang waktu, eksperimen, atau strategi
identifikasi kausal.

\subsection{Dataset}\label{dataset-43}

Laporan sumber menunjukkan bahwa dataset terdiri atas respons kuesioner
dari 100 perusahaan atau tempat usaha yang valid: 64 dari Leighton
Linslade dan 36 dari Lichfield. Kerangka valid yang dikirimi survei
berjumlah 65 di Leighton Linslade dan 53 di Lichfield, sehingga Tabel 1
mencatat 118 kiriman dan 100 respons. Kerangka awal sebelum penyaringan
berjumlah 71 dan 82, tetapi banyak entitas kemudian diketahui telah
tutup, pindah, atau tidak memenuhi definisi industri (pp.~616-617, Tabel
1).

Klasifikasi sektor pada catatan 4 menggunakan kode aktivitas SIC 1980:
8370 professional and technical services n.e.s.; 8380 advertising; 8394
computer services; dan 8395 miscellaneous business services (p.~623).
Tabel 1 menunjukkan bahwa komposisi kerangka dan respons berbeda antara
lokasi: Leighton Linslade memiliki proporsi perusahaan jasa komputer
yang lebih besar, sedangkan Lichfield memiliki lebih banyak jasa bisnis
umum, terutama agen rekrutmen tenaga kerja (p.~617, Tabel 1).

Variabel yang tersedia dalam bentuk agregat meliputi karakteristik asal,
status single-site atau bagian dari grup, periode pendirian, pertumbuhan
omzet dan tenaga kerja, rata-rata persentase geografis penjualan, input,
dan tenaga kerja, keunggulan lokasi, relevansi serta penggunaan lima
penyedia layanan lokal, dan sumber geografis informasi serta pesaing
(pp.~618-622, Tabel 2-Tabel 7 dan Gambar 2).

\textbf{Batas evidence:} Data mentah tingkat perusahaan, kuesioner
lengkap, codebook, bobot sampel, prosedur penanganan data hilang, dan
dokumentasi jawaban ganda: Tidak disebutkan dalam file. Tidak disebutkan
dalam file pula apakah persentase pada Tabel 7 menggunakan penyebut yang
sama untuk semua jenis informasi atau bagaimana jawaban lebih dari satu
sumber diperlakukan.

\subsection{Population / Sample}\label{population-sample-43}

Populasi sasaran yang dinyatakan secara operasional adalah perusahaan
dan tempat usaha kecil di sektor jasa bisnis pada dua pemukiman yang
mudah diakses dan dekat dengan konurbasi besar. Leighton Linslade
terdiri atas Leighton Buzzard dan Linslade yang menyatu melalui
pertumbuhan penduduk cepat; Lichfield adalah kota dengan sejarah dan
identitas tunggal yang lebih panjang. Pada 1991, populasinya
masing-masing 31.855 dan 27.577 (p.~616).

Penulis menekankan bahwa kedua pemukiman tersebut jauh lebih besar
daripada definisi pemukiman rural Keeble dan Tyler yang berpenduduk
kurang dari 10.000. Karena itu, penggunaan istilah pemukiman rural yang
mudah diakses memiliki batasan sejak awal (p.~616). Gambar 1 menempatkan
kedua lokasi terhadap jalan tol, jalur kereta utama, batas perkiraan
kawasan urban besar, London, dan kawasan West Midlands (p.~617, Gambar
1).

Kerangka awal mengidentifikasi 71 entitas di Leighton Linslade dan 82 di
Lichfield. Setelah perusahaan yang tutup, pindah, atau tidak sesuai
definisi dikeluarkan, kerangka sah menjadi 65 dan 53. Respons lengkap
diperoleh dari 64 perusahaan di Leighton Linslade, atau 98,5 persen, dan
36 di Lichfield, atau 67,4 persen (p.~617, Tabel 1). Penolakan langsung
dan kuesioner pos/faks yang tidak dikembalikan menjadi masalah khusus di
Lichfield. Uji chi-kuadrat yang dilaporkan bernilai 3,54 dan dianggap
menunjukkan tidak ada perbedaan signifikan pada tingkat respons, tetapi
hasilnya disebut marginal (p.~617).

Dalam respons yang terkumpul, 61,6 persen merupakan perusahaan mikro
dengan lima atau lebih sedikit pekerja, 86,0 persen merupakan perusahaan
single-site, dan 83,8 persen dibentuk di dua pemukiman studi. Hampir
semua perusahaan, kecuali satu, berdiri sejak 1970. Perusahaan yang
pindah ke Lichfield lebih banyak daripada yang pindah ke Leighton
Linslade, tetapi perbedaan tersebut tidak signifikan secara statistik.
Di kedua lokasi, mayoritas perusahaan yang pindah berasal dari pemukiman
sekitar, bukan langsung dari London atau West Midlands; sekitar tiga
perempat dari perusahaan yang pindah berasal dari lokasi seperti
Bedford, Hemel Hempstead, Luton, Watford, Cannock, Stafford, Tamworth,
dan Derby (p.~618, Tabel 2).

\textbf{Inferensi review:} Tidak disebutkan dalam file apakah pemilihan
dari kerangka sah dilakukan secara acak, sensus, atau dengan bobot
tertentu. Perbedaan ukuran kerangka, nonrespons Lichfield, perbedaan
komposisi sektor, dan karakter kedua lokasi membatasi keterwakilan
sampel terhadap perusahaan kecil rural secara umum.

\subsection{Dependent Variables}\label{dependent-variables-43}

Dependent variables formal: Tidak disebutkan dalam file. Artikel tidak
menyajikan model yang menetapkan satu variabel hasil dan seperangkat
prediktor secara eksplisit.

Indikator hasil atau respons yang dilaporkan meliputi:

\begin{itemize}
\tightlist
\item
  Pertumbuhan omzet selama tiga tahun terakhir. Secara keseluruhan, 80,2
  persen perusahaan melaporkan pertumbuhan omzet (p.~618).
\item
  Pertumbuhan tenaga kerja selama tiga tahun terakhir. Secara
  keseluruhan, angkanya 40,4 persen; Lichfield mencapai 51 persen,
  sedangkan Leighton Linslade 34 persen (p.~618).
\item
  Distribusi geografis nilai penjualan, sumber input, dan asal tenaga
  kerja (pp.~618-620, Tabel 3-Tabel 5).
\item
  Keunggulan lokasi yang dipilih responden, serta penilaian relevansi
  dan penggunaan layanan institusi lokal (pp.~619-621, Gambar 2 dan
  Tabel 6).
\item
  Sumber geografis informasi pasar, teknis, pemasok, kontak bisnis
  informal, dan pesaing (pp.~621-622, Tabel 7).
\end{itemize}

\textbf{Inferensi review:} Karakteristik seperti asal perusahaan, status
single-site atau grup, kode industri, dan tahun pendirian berfungsi
sebagai deskripsi sampel, bukan dependent variable yang dimodelkan.
Variabel-variabel di atas lebih tepat dibaca sebagai indikator
deskriptif tentang akses eksternalitas dan kinerja yang dilaporkan
sendiri, bukan sebagai bukti pengaruh \emph{borrowed size} terhadap
pertumbuhan.

\subsection{Results}\label{results-43}

Laporan sumber pertama-tama menunjukkan bahwa perusahaan yang disurvei
relatif kecil, muda, dan sebagian besar berakar di lokasi studi.
Mayoritas 61,6 persen adalah perusahaan mikro, 86,0 persen single-site,
dan 83,8 persen dibentuk di Leighton Linslade atau Lichfield. Sebanyak
80,2 persen mengalami pertumbuhan omzet dalam tiga tahun terakhir,
tetapi hanya 40,4 persen mengalami pertumbuhan tenaga kerja. Pertumbuhan
tenaga kerja lebih sering dilaporkan di Lichfield, yaitu 51 persen,
dibandingkan Leighton Linslade, yaitu 34 persen (p.~618, Tabel 2).

\textbf{Penjualan dan input.} Menurut Tabel 3, rata-rata penjualan
Leighton Linslade terdiri atas 12,2 persen ke lokalitas segera, 58,2
persen ke konurbasi dan kawasan UK, 22,8 persen ke bagian UK lainnya,
dan 7,1 persen ke luar negeri. Untuk Lichfield, angkanya 17,5 persen,
22,1 persen, 51,1 persen, dan 9,7 persen. Total kedua lokasi adalah 14,0
persen, 45,6 persen, 32,6 persen, dan 8,0 persen. Tidak ada perbedaan
signifikan untuk penjualan lokal dan luar negeri, tetapi terdapat
perbedaan signifikan pada penjualan ke konurbasi atau kawasan UK serta
bagian UK lainnya; nilai F yang dilaporkan untuk dua kategori terakhir
adalah 53,58 dan 27,97 dengan p \textless{} 0,01 (p.~618, Tabel 3).

Tabel 4 menunjukkan rata-rata sumber input Leighton Linslade sebesar
22,4 persen dari lokalitas segera, 34,4 persen dari lokalitas lain, 11,8
persen dari konurbasi, 30,4 persen dari bagian UK lainnya, dan 0,9
persen dari luar negeri. Untuk Lichfield, angkanya 43,1 persen, 12,7
persen, 16,1 persen, 22,9 persen, dan 5,2 persen. Perbedaan signifikan
terutama berkaitan dengan sumber input lokal: perusahaan Lichfield
mengambil proporsi input yang lebih besar dari lokalitas segera dan
lokalitas lain daripada perusahaan Leighton Linslade (p.~619, Tabel 4).
Untuk Leighton Linslade, Milton Keynes dan Luton disebut lebih penting
sebagai sumber input daripada London (p.~619).

\textbf{Keunggulan lokasi dan tenaga kerja.} Gambar 2 meminta responden
menyebutkan tiga keunggulan utama lokasi. Penulis melaporkan bahwa
preferensi tempat tinggal pemilik atau manajer, lingkungan yang menarik,
serta ketersediaan dan biaya tempat usaha merupakan atribut penting.
Sekitar sepertiga responden Lichfield menyebut kualitas lingkungan.
Namun, ketersediaan dan biaya tenaga kerja serta sebagian faktor lain,
termasuk infrastruktur dukungan bisnis lokal, hanya disebut oleh
sebagian kecil perusahaan. Komunikasi dan sentralitas adalah keunggulan
yang paling sering disebut di kedua lokasi. Kedekatan dengan London
disebut oleh sekitar seperlima perusahaan jasa bisnis di Leighton
Linslade (p.~619-620, Gambar 2).

Perbedaan keunggulan lokasi antara dua pemukiman tidak signifikan;
korelasi peringkat Spearman yang dilaporkan adalah 0,780 (p.~620, Gambar
2). Tabel 5 menunjukkan bahwa tenaga kerja berasal terutama dari
lokalitas segera: 51,9 persen di Leighton Linslade dan 54,5 persen di
Lichfield. Sumber dari lokalitas lain adalah 37,5 persen dan 21,3
persen; dari konurbasi 5,6 persen dan 5,8 persen; dan dari bagian UK
lainnya 5,0 persen dan 18,5 persen. Lokalitas sekitar tetap penting,
sedangkan konurbasi bukan sumber langsung terbesar tenaga kerja (p.~620,
Tabel 5).

\textbf{Institusi lokal.} Tabel 6 menunjukkan bahwa mayoritas perusahaan
di kedua lokasi tidak menggunakan layanan Training and Enterprise
Councils, pemerintah lokal, Business Link, dan perguruan tinggi lokal.
Penggunaan relatif lebih tinggi terlihat pada Chamber of Commerce di
kedua lokasi dan Business Link di Lichfield, tetapi tabel tetap mencatat
60,6 persen perusahaan Lichfield tidak menggunakan Business Link.
Proporsi penggunaan yang dilaporkan adalah sebagai berikut: TEC 6,3
persen di Leighton Linslade dan 25,8 persen di Lichfield; pemerintah
lokal 17,2 persen dan 24,2 persen; Business Link 10,9 persen dan 39,4
persen; Chamber of Commerce 51,6 persen dan 65,6 persen; serta perguruan
tinggi lokal 23,4 persen dan 35,3 persen.

Dalam penilaian relevansi, TEC, pemerintah lokal, dan perguruan tinggi
lokal termasuk yang paling sering dianggap tidak relevan. Chamber of
Commerce dan Business Link dinilai relatif lebih relevan, tetapi
proporsi yang menganggap layanan itu tidak relevan tetap besar. Uji
tabel kontingensi dua kali dua atas penggunaan dan tidak penggunaan
menunjukkan hasil signifikan pada batas kepercayaan 95 persen, kecuali
untuk Chamber of Commerce. Setelah kategori relevan dan sangat relevan
digabung, serta kategori tidak relevan dan tidak tahu digabung, uji
chi-kuadrat tidak menunjukkan perbedaan signifikan antarlokasi dalam
persepsi relevansi secara keseluruhan (pp.~620-621, Tabel 6).

\textbf{Informasi dan pesaing.} Tabel 7 memperlihatkan bahwa sumber
nasional, termasuk jurnal dan pameran yang dimasukkan ke kategori bagian
UK lainnya, merupakan sumber utama informasi bisnis secara keseluruhan.
Ada perbedaan orientasi antara lokasi, tetapi korelasi peringkat
Spearman tidak signifikan. Di Leighton Linslade, ekonomi lokal segera
hampir tidak penting sebagai sumber informasi bisnis dan persaingan;
lokalitas lain seperti Milton Keynes dan Luton lebih sering menjadi
sumber informasi pasar dan kontak bisnis informal, serta penting sebagai
sumber informasi pemasok. Di Lichfield, ekonomi lokal segera lebih
menonjol untuk informasi pasar, informasi pemasok, kontak informal, dan
persaingan (pp.~621-622, Tabel 7).

Beberapa angka Tabel 7 memperlihatkan pola tersebut. Untuk informasi
teknis, bagian UK lainnya dicatat sebagai sumber bagi 26 respons
Lichfield atau 83,9 persen dan 29 respons Leighton Linslade atau 45,3
persen. Untuk kontak bisnis informal, bagian UK lainnya dicatat sebesar
57,6 persen di Lichfield, sedangkan lokalitas lain sebesar 50,0 persen
di Leighton Linslade. Untuk pesaing, bagian UK lainnya mencapai 60,6
persen di Lichfield dan 35,9 persen di Leighton Linslade; konurbasi
mencapai 32,8 persen di Leighton Linslade dan 27,3 persen di Lichfield
(p.~622, Tabel 7). Persentase tersebut dilaporkan sebagaimana tertulis
dalam TXT; file tidak menjelaskan penyebut atau aturan jawaban ganda
secara lebih lanjut.

\subsection{Findings}\label{findings-43}

Interpretasi penulis adalah bahwa perusahaan kecil di kedua lokasi
memperoleh manfaat dari hubungan geografis yang melampaui lokalitas
langsung. Penjualan, input, tenaga kerja, dan informasi tidak seluruhnya
terlokalisasi di pemukiman tempat perusahaan berada. Keunggulan utama
kedua lokasi lebih banyak berkaitan dengan aksesibilitas dan sentralitas
ke pemukiman lain sebagai pasar dan sumber eksternalitas pecuniary
daripada dengan keunggulan internal ekonomi lokal (pp.~619-622).

Penulis menilai ada dukungan terhadap dugaan bahwa London dan West
Midlands menjadi sumber eksternalitas teknologis dan informasi bagi
perusahaan jasa bisnis di dua lokasi. Lichfield tampak terhubung
sekaligus dengan ekonomi regional South East yang lebih luas dan dengan
konurbasi West Midlands. Leighton Linslade terhubung dengan Greater
London serta pusat provinsi lain di South East dan sekitarnya. Penulis
juga menekankan bahwa perusahaan di kedua lokasi mengakses eksternalitas
teknologis dari konurbasi besar dengan tingkat kepentingan yang kurang
lebih sama, meskipun sumber lokal dan sumber UK lainnya berbeda
(pp.~620-622).

Temuan tentang institusi lokal mengarah pada interpretasi penulis bahwa
peran kapasitas kelembagaan formal dalam wilayah rural yang mudah
diakses dapat terlalu dibesar-besarkan. Mayoritas perusahaan tidak
menggunakan layanan institusi lokal atau menganggapnya tidak relevan.
Penulis tidak menyimpulkan bahwa semua institusi lokal tidak berguna;
Chamber of Commerce dan Business Link, terutama di Lichfield,
menunjukkan penggunaan yang lebih tinggi (p.~621, Tabel 6).

Inferensi review: bukti artikel mendukung adanya pola \emph{borrowed
size} sebagai hubungan akses dan jaringan, tetapi tidak menguji efek
kausalnya terhadap pertumbuhan. Hasil paling kuat berlaku untuk
perusahaan jasa bisnis kecil di dua pemukiman yang relatif besar dan
mudah diakses, bukan untuk seluruh perusahaan rural atau seluruh bentuk
pergeseran urban-rural.

\subsection{Conclusions}\label{conclusions-43}

Interpretasi penulis menyatakan bahwa \emph{borrowed size} berguna untuk
memahami sebagian pola geografis pembentukan dan pertumbuhan perusahaan
kecil di UK. Pemikiran ini menggeser perhatian dari karakteristik satu
lokalitas ke interaksi antarkelompok lokalitas. Definisi ``lokal'' yang
berbeda dapat relevan untuk jenis eksternalitas yang berbeda dan
menentukan titik pengungkit kebijakan yang sesuai (pendahuluan,
pp.~614-615).

Penulis menyimpulkan bahwa pemukiman rural yang mudah diakses mungkin
menawarkan hubungan yang bersifat komplementer: perusahaan memperoleh
eksternalitas pecuniary dan teknologis pada skala konurbasi atau ekonomi
regional yang terurbanisasi tanpa harus menanggung seluruh disekonomi
ukuran urban. Pada saat yang sama, lokasi tersebut mungkin hanya
memiliki sedikit eksternalitas teknologis atau fondasi kelembagaan yang
bersifat lokal (kesimpulan, pp.~622-623).

\textbf{Interpretasi penulis:} Penulis memberi kualifikasi kuat terhadap
kesimpulan itu. UK mungkin terlalu kecil dan terlalu didominasi oleh
satu pusat utama untuk menguji \emph{borrowed size} secara empiris.
Studi yang lebih jauh dari London, analisis multivariat, dan penelitian
tentang karakteristik fisik wilayah rural diperlukan sebelum
generalisasi lebih luas dibuat (kesimpulan, p.~622).

Inferensi review: kesimpulan artikel bersifat eksploratori dan
kontekstual. Ia menawarkan mekanisme relasional yang masuk akal dan
didukung oleh pola jaringan yang dilaporkan, tetapi tidak membuktikan
bahwa akses ke kota besar menyebabkan pertumbuhan omzet atau tenaga
kerja.

\subsection{Contributions}\label{contributions-43}

Bagian kontribusi formal yang terpisah tidak disebutkan dalam file. Poin
berikut adalah \textbf{inferensi review dari isi TXT}:

\begin{itemize}
\tightlist
\item
  \textbf{Inferensi review:} Artikel menghidupkan kembali penggunaan
  eksplisit konsep \emph{borrowed size} Alonso dalam studi perusahaan
  kecil dan menghubungkannya dengan literatur aglomerasi, perusahaan
  kecil, serta sistem urban (pp.~614-616).
\item
  \textbf{Inferensi review:} Artikel memberi bukti eksploratori dari dua
  kasus yang kontras: Leighton Linslade di South East England dan
  Lichfield di luar South East, sehingga menguji apakah gagasan tersebut
  memiliki relevansi di luar kasus yang berorientasi London (pp.~616-617
  dan catatan 5, p.~623).
\item
  \textbf{Inferensi review:} Artikel memecah akses eksternalitas menjadi
  dimensi input-output, tenaga kerja, lokasi, informasi, persaingan, dan
  institusi, sehingga menunjukkan bahwa ``lokal'' tidak identik untuk
  semua fungsi (Tabel 3-Tabel 7).
\item
  \textbf{Inferensi review:} Artikel menantang asumsi bahwa penguatan
  dukungan institusi lokal selalu menjadi basis utama pertumbuhan
  perusahaan kecil di wilayah rural yang mudah diakses (pp.~620-623).
\item
  \textbf{Inferensi review:} Artikel menghubungkan temuan perusahaan
  dengan implikasi kebijakan lintas-lokalitas, regional, dan nasional,
  bukan hanya kebijakan yang dibatasi oleh batas pemukiman (p.~623).
\end{itemize}

Inferensi review: kontribusi empirisnya lebih tepat disebut pemetaan
pola dan pengajuan agenda penelitian daripada estimasi pengaruh
\emph{borrowed size}. Tidak ada dasar dalam file untuk menyebut artikel
ini sebagai verifikasi final atau pengukuran umum bagi semua perusahaan
kecil.

\subsection{Research Gap}\label{research-gap-43}

Gap yang dinyatakan atau ditandai oleh artikel adalah sebagai berikut.
Poin-poin ini merupakan \textbf{laporan sumber atau interpretasi
penulis} kecuali jika diberi label lain:

\begin{itemize}
\tightlist
\item
  \textbf{Laporan sumber:} Penulis menyatakan bahwa, sejauh pengetahuan
  penulis, konsep \emph{borrowed size} belum dikembangkan secara
  eksplisit sejak penjelasan singkat Alonso pada 1973 (p.~614).
\item
  \textbf{Laporan sumber:} Bukti mengenai pengaruh \emph{borrowed size}
  di luar South East England disebut lebih terbatas dibandingkan bukti
  tentang hubungan pemukiman kecil dengan Greater London (pp.~614-615).
\item
  \textbf{Laporan sumber:} Peran aktual dukungan institusi lokal
  dibandingkan akses ke eksternalitas dari kota besar masih
  diperdebatkan. Catatan 6 memperbandingkan Keeble dan Tyler, Bryson dan
  Daniels, serta Curran dan Blackburn dengan posisi yang berbeda tentang
  eksternalitas teknologis lokal dan pentingnya ekonomi lokal (p.~623).
\item
  \textbf{Interpretasi penulis:} Hubungan komplementer antara
  eksternalitas pecuniary dan teknologis belum ditangkap secara memadai
  oleh analisis satu per satu. Penulis mengusulkan analisis multivariat
  untuk mengisi celah ini (p.~622).
\item
  \textbf{Laporan sumber:} Karakteristik fisik wilayah rural belum
  dieksplorasi secara eksplisit dalam artikel (pp.~622-623).
  \textbf{Inferensi review:} Peran amenitas rural dalam mekanisme
  \emph{borrowed size} juga tidak diuji secara terpisah.
\item
  \textbf{Interpretasi penulis:} Belum jelas apakah eksternalitas
  teknologis yang nyata, atau sekurang-kurangnya fondasi kelembagaannya,
  tersedia bagi perusahaan di pemukiman rural dan rural yang mudah
  diakses (p.~623).
\item
  \textbf{Interpretasi penulis:} Dibutuhkan kasus rural yang lebih jauh
  dari London serta basis kasus yang lebih luas karena penulis mengakui
  keterbatasan generalisasi dari dua studi kasus (pp.~622-623).
\end{itemize}

Inferensi review: artikel ini membuka gap pengukuran dan perbandingan,
tetapi tidak menutup gap kausal, temporal, dan keterwakilan geografis.

\subsection{Limitations}\label{limitations-43}

\textbf{Laporan sumber:} Keterbatasan yang dinyatakan oleh penulis
meliputi:

\begin{itemize}
\tightlist
\item
  Bahan empiris berskala kecil dan eksploratori, berbasis dua studi
  kasus, sehingga generalisasi terbatas (p.~616 dan kesimpulan, p.~622).
\item
  Leighton Linslade dan Lichfield lebih besar daripada ambang 10.000
  penduduk yang dipakai dalam definisi rural Keeble dan Tyler. Karena
  itu, temuan tidak dapat langsung dianggap mewakili semua pemukiman
  rural (p.~616).
\item
  UK mungkin terlalu kecil dengan tingkat primasi terlalu tinggi untuk
  menguji \emph{borrowed size} secara memadai. Penulis mengusulkan kasus
  yang lebih jauh dari London (p.~622).
\item
  Studi Lichfield merupakan replikasi yang tidak sempurna terhadap studi
  Leighton Linslade dan dilakukan pada periode berbeda, yaitu 1998
  dibandingkan 1996 (catatan 5, p.~623).
\item
  Kerangka sampel menghadapi perusahaan yang tutup, pindah, atau tidak
  sesuai definisi; Lichfield juga menghadapi perpindahan perusahaan dari
  kota ke kawasan industri di wilayah otoritas lokal (pp.~616-617).
\item
  Respons Lichfield lebih rendah, dengan penolakan langsung dan
  kuesioner pos/faks yang tidak dikembalikan; penulis menyebut hasil uji
  perbedaan tingkat respons sebagai marginal (p.~617).
\item
  Analisis terbatas pada statistik deskriptif dan bivariat atas jenis
  eksternalitas secara terpisah. Komplementaritas antareksternalitas
  belum diuji dengan teknik multivariat (p.~622).
\item
  Karakteristik fisik wilayah rural tidak diteliti secara eksplisit
  (pp.~622-623). \textbf{Inferensi review:} Hubungan rinci antara
  amenitas rural dan kemampuan meminjam manfaat ukuran urban juga tidak
  diuji secara terpisah.
\item
  Cakupan sektoral difokuskan pada jasa bisnis dan empat kelompok SIC
  1980, sehingga keterterapan pada sektor lain tidak ditunjukkan (p.~616
  dan catatan 4, p.~623).
\end{itemize}

\textbf{Inferensi review:} Desain yang dilaporkan tidak menetapkan
hubungan sebab-akibat, pertumbuhan hanya dilaporkan sebagai kondisi tiga
tahun sebelumnya, dan tidak ada analisis longitudinal tingkat
perusahaan. Perlakuan terhadap data hilang, jawaban ganda, serta
penyebut persentase pada Tabel 7: \textbf{Tidak disebutkan dalam file}.
Perbedaan waktu pengumpulan serta mode respons dapat mengurangi
keterbandingan dua kasus, tetapi kemungkinan tersebut adalah inferensi
review, bukan klaim kausal penulis.

\subsection{Challenges}\label{challenges-43}

\textbf{Laporan sumber:} Tantangan operasional dalam membangun sampel
terlihat dari banyaknya entitas yang semula teridentifikasi tetapi
ternyata sudah tutup, pindah, atau tidak termasuk definisi industri. Di
Lichfield, perpindahan dari kota ke kawasan industri di tempat lain
dalam wilayah otoritas lokal membuat kerangka berbasis lokasi kota
semakin sulit ditetapkan (pp.~616-617).

\textbf{Laporan sumber:} Pengumpulan respons juga menantang. Sebagian
responden meminta kuesioner melalui pos atau faks; Lichfield mengalami
penolakan langsung dan pengembalian yang tidak lengkap. Hal ini
menyebabkan tingkat respons dua lokasi berbeda jauh meskipun uji
statistik yang dilakukan tidak menemukan perbedaan signifikan pada batas
yang digunakan (p.~617).

\textbf{Interpretasi penulis:} Penelitian ini menggabungkan informasi
kuantitatif dan unsur kualitatif untuk menilai relevansi \emph{borrowed
size}, dengan perhatian pada transaksi, tenaga kerja, informasi, dan
kontak antarpemukiman (pp.~616-617). \textbf{Inferensi review:}
Keterbatasan sumber daya disebut sebagai konteks pemilihan survei
telepon, sementara hasil akhirnya disajikan terutama melalui statistik
deskriptif dan bivariat (pp.~616-617 dan p.~622).

\textbf{Laporan sumber:} Perbedaan komposisi sektor juga menjadi
tantangan perbandingan: Leighton Linslade memiliki lebih banyak
perusahaan jasa komputer, sedangkan Lichfield memiliki lebih banyak jasa
bisnis umum. Artikel juga membandingkan studi yang dimulai pada waktu
berbeda dan tidak sepenuhnya identik (p.~617, Tabel 1; catatan 5,
p.~623). \textbf{Inferensi review:} Tantangan tersebut membuat perbedaan
antarlokasi sulit dipisahkan secara bersih dari perbedaan sektor, waktu,
dan cara pengumpulan data.

\subsection{Practical Implications}\label{practical-implications-43}

\textbf{Interpretasi penulis:} Kebijakan ekonomi yang murni lokal
mungkin memiliki batas ketika perusahaan memperoleh manfaat penting dari
ekonomi eksternal di pemukiman lain. Karena \emph{borrowed size} bekerja
melalui jaringan antarlokasi dan tidak selalu melalui ekonomi lokal
langsung, titik pengungkit kebijakan perlu disesuaikan dengan jenis
eksternalitas yang hendak diperkuat (pendahuluan, pp.~614-615).

\textbf{Interpretasi penulis:} Kebijakan relevan pada tingkat regional,
termasuk kolaborasi antar-institusi di pemukiman yang terhubung secara
fungsional. Penulis juga menyatakan bahwa kurangnya massa kritis pada
cabang industri tertentu dapat meningkatkan relevansi kebijakan industri
nasional, bukan hanya dukungan bisnis lokal (kesimpulan, p.~623).

\textbf{Interpretasi penulis:} Temuan institusi lokal menyarankan agar
dukungan formal tidak otomatis dianggap sebagai penjelasan utama
keberhasilan perusahaan rural yang mudah diakses. Sebaliknya,
aksesibilitas, sentralitas, hubungan dengan pasar, dan jaringan
informasi dari kota atau pemukiman sekitar perlu dipertimbangkan bersama
amenitas lokal dan biaya lokasi (pp.~619-623).

Inferensi review: evaluasi kebijakan sebaiknya memisahkan setidaknya
empat saluran, yaitu penjualan, input, tenaga kerja, serta informasi dan
kontak. Menggabungkannya sebagai satu ukuran ``ekonomi lokal'' dapat
menyembunyikan bahwa perusahaan menggunakan skala geografis yang berbeda
untuk fungsi yang berbeda.

\subsection{Applications}\label{applications-43}

Tidak disebutkan dalam file sebagai bagian aplikasi yang terpisah atau
sebagai evaluasi terhadap program kebijakan tertentu. Aplikasi berikut
adalah inferensi review yang dibatasi oleh desain artikel:

\begin{itemize}
\tightlist
\item
  Kerangka \emph{borrowed size} dapat digunakan sebagai alat diagnosis
  untuk memetakan asal penjualan, input, tenaga kerja, informasi, dan
  pesaing perusahaan kecil pada beberapa skala geografis.
\item
  Survei serupa dapat dipakai untuk membandingkan pemukiman kecil yang
  dekat dengan konurbasi berbeda, dengan memperlakukan aksesibilitas dan
  jaringan sebagai bagian dari ekonomi lokal fungsional.
\item
  Tabel relevansi dan penggunaan institusi dapat digunakan untuk
  membedakan keberadaan penyedia layanan dari penggunaan aktual dan
  persepsi kegunaannya.
\item
  Temuan dapat diterapkan sebagai peringatan metodologis dalam membaca
  pertumbuhan perusahaan di wilayah rural yang mudah diakses:
  pertumbuhan lokal tidak dengan sendirinya membuktikan adanya basis
  eksternalitas teknologis lokal.
\end{itemize}

Aplikasi pada wilayah, sektor, atau negara lain: Tidak disebutkan dalam
file. Inferensi review untuk penerapan di luar konteks artikel
memerlukan pengujian ulang, bukan pemindahan kesimpulan secara langsung.

\subsection{Future Research
(Optional)}\label{future-research-optional-43}

\textbf{Laporan sumber:} Artikel secara eksplisit menyarankan penelitian
tentang apakah eksternalitas teknologis yang nyata, atau setidaknya
fondasi kelembagaannya, tersedia bagi perusahaan di pemukiman rural dan
rural yang mudah diakses. Saran ini ditempatkan dalam konteks
keterbatasan generalisasi dua studi kasus dan perdebatan literatur yang
belum selesai (kesimpulan, p.~623).

Agenda lain yang disebutkan atau ditandai dalam TXT meliputi:

\begin{itemize}
\tightlist
\item
  \textbf{Laporan sumber:} Meneliti pemukiman rural yang mudah diakses
  tetapi berada lebih jauh dari London daripada Lichfield (p.~622).
\item
  \textbf{Laporan sumber:} Menggunakan teknik statistik multivariat
  untuk menangkap komplementaritas eksternalitas pecuniary dan
  teknologis (p.~622).
\item
  \textbf{Laporan sumber:} Mengeksplorasi karakteristik fisik pemukiman
  rural dan peran amenitas rural dalam memperoleh manfaat ukuran urban
  tanpa seluruh disekonominya (pp.~622-623).
\item
  \textbf{Inferensi review:} Memperluas perbandingan kasus untuk menguji
  generalisasi dan membedakan konteks South East England dari wilayah UK
  lain (pp.~622-623).
\end{itemize}

Urutan prioritas, ukuran sampel yang disarankan, periode pengamatan, dan
desain operasional penelitian lanjutan: Tidak disebutkan dalam file.

\subsection{Provenance Notes}\label{provenance-notes-43}

\begin{itemize}
\tightlist
\item
  Bahan yang ditinjau hanya
  \texttt{/Users/mac/Documents/Mac/{[}2{]}\ Obsidian\ Vault/zettelkasten/source/journal/A28-phelps-fallon-williams-2001-small-firms-borrowed-size-and-the-urban-rural-shift-10.1080-00343400120075885.txt};
  tidak ada TXT A28 lain yang ditemukan di bawah \texttt{source/}.
\item
  TXT memuat blok \texttt{{[}PDF\ Metadata{]}} pada lines 1-26 dan blok
  \texttt{{[}Extracted\ Text{]}} pada lines 27-892. Metadata mencatat
  judul \emph{Small Firms, Borrowed Size and the Urban-Rural Shift};
  field \texttt{Author} kosong; \texttt{Creator}
  \texttt{RealPage\ PDF\ Generator\ 1.0}; \texttt{Producer}
  \texttt{iText\ 4.2.0\ by\ 1T3XT}; \texttt{CreationDate} 10 January
  2002 dan \texttt{ModDate} 10 October 2014;
  \texttt{Custom\ Metadata:\ no}; \texttt{Metadata\ Stream:\ yes};
  \texttt{Tagged:\ no}; \texttt{UserProperties:\ no};
  \texttt{Suspects:\ no}; \texttt{Form:\ none};
  \texttt{JavaScript:\ no}; 13 halaman; tidak terenkripsi; ukuran
  halaman \texttt{555.024\ x\ 773.008\ pts}; rotasi 0; ukuran file
  417140 bytes; tidak dioptimalkan; PDF versi 1.4 (TXT lines 1-26).
  Metadata juga mencatat ekstraksi pada 2026-08-31 dengan
  \texttt{pdftotext\ -layout}.
\item
  Bagian awal hasil ekstraksi mencatat artikel diunduh oleh
  \texttt{{[}The\ Aga\ Khan\ University{]}} pada 10 October 2014,
  penerbit Routledge, serta tanggal ``Published online: 18 Aug 2010''
  (TXT lines 28-58). Header artikel yang sama mencantumkan terbitan
  \emph{Regional Studies} 35:7, pp.~613-624, tahun 2001 dan DOI
  10.1080/00343400120075885 (TXT lines 79-145).
\item
  Locator dalam review memakai nomor halaman jurnal tercetak 613-624 dan
  locator internal seperti Tabel 1-Tabel 7, Gambar 1-Gambar 2, serta
  catatan bernomor. Rentang lines TXT dicantumkan pada bagian yang
  merangkum abstrak atau metadata bila membantu pelacakan.
\item
  Seluruh angka, lokasi, periode, istilah, dan interpretasi yang
  diatribusikan kepada artikel berasal dari TXT. Tidak dilakukan
  penelusuran atau verifikasi eksternal terhadap DOI, jurnal, penerbit,
  atau daftar referensi.
\item
  \textbf{Masalah sumber/ekstraksi:} hasil \texttt{pdftotext\ -layout}
  memiliki artefak tata letak dan encoding pada sejumlah kata serta
  simbol statistik. Review ini terutama menggunakan parafrasa; tidak ada
  kutipan panjang yang direkonstruksi dari teks yang artefaknya tidak
  pasti. Tabel 7 juga tidak menjelaskan penyebut persentase atau
  perlakuan jawaban ganda (TXT lines 692-715). Batas ini dicatat sebagai
  keterbatasan provenance, tetapi tidak dengan sendirinya menghalangi
  penyelesaian review karena klaim dikualifikasi dan informasi yang
  tidak tersedia dinyatakan secara eksplisit.
\item
  Pernyataan berlabel \texttt{Laporan\ sumber} merujuk pada data atau
  klaim yang dilaporkan artikel. \texttt{Interpretasi\ penulis} merujuk
  pada penjelasan dan kesimpulan penulis. \texttt{Inferensi\ review}
  adalah pembacaan analitis dari TXT dan bukan pernyataan baru dari
  artikel.
\item
  Informasi yang tidak tersedia atau tidak relevan dinyatakan dengan
  \texttt{Tidak\ disebutkan\ dalam\ file} atau \texttt{Tidak\ berlaku}
  sesuai kebutuhan. Tidak ada perubahan pada TXT sumber.
\end{itemize}

\section{Review B01}\label{review-b01}

\subsection{Identitas Sumber}\label{identitas-sumber-13}

\begin{itemize}
\tightlist
\item
  Kode kajian: B01.
\item
  Judul: \emph{Beyond Polycentricity: Does Stronger Integration Between
  Cities in Polycentric Urban Regions Improve Performance?}
\item
  Penulis: Evert Meijers, Marloes Hoogerbrugge, dan Rodrigo Cardoso.
\item
  Publikasi: \emph{Tijdschrift voor Economische en Sociale Geografie},
  2018, volume 109, nomor 1, halaman 1-21.
\item
  DOI: 10.1111/tesg.12292.
\item
  Afiliasi yang dicantumkan: Delft University of Technology, Faculty of
  Architecture and the Built Environment; serta Erasmus Happiness
  Economics Research Organisation.
\item
  Riwayat naskah: diterima Desember 2016 dan disetujui Juni 2017.
\item
  Jenis sumber: artikel jurnal. Status peer-review: Tidak disebutkan
  dalam file dan tidak diasumsikan.
\item
  Basis bahan review: seluruh teks yang tersedia dalam TXT hasil
  ekstraksi PDF, termasuk blok metadata ekstraksi, tabel, gambar,
  catatan, dan referensi.
\end{itemize}

\textbf{Penanda epistemik:} \texttt{{[}Laporan\ sumber{]}} berarti isi
yang dilaporkan atau dihitung dalam artikel;
\texttt{{[}Interpretasi\ penulis{]}} berarti penjelasan teoretis,
substantif, atau kebijakan dari penulis artikel;
\texttt{{[}Inferensi\ pengolahan{]}} berarti pembacaan hati-hati yang
dibuat dalam review ini dan bukan klaim langsung penulis.

\subsection{Summarized Abstract}\label{summarized-abstract-44}

{[}Laporan sumber{]} Artikel menyatakan bahwa sekitar seperempat
penduduk Eropa tinggal di \emph{polycentric urban regions} (PUR), yaitu
gugus kota yang secara historis dan administratif berbeda, tetapi
berdekatan, terhubung baik, dan berukuran relatif serupa. Tujuan
analisisnya adalah mengeksplorasi apakah integrasi yang lebih erat dapat
meningkatkan manfaat aglomerasi pada tingkat PUR. Penulis menyusun
daftar komprehensif pertama yang mereka laporkan mengenai PUR Eropa,
berjumlah 117, mengukur integrasi fungsional, institusional, dan
kultural, lalu menguji hubungannya dengan kinerja.

{[}Laporan sumber{]} Kinerja didefinisikan sebagai sejauh mana ekonomi
urbanisasi berkembang, dengan kehadiran fungsi-fungsi metropolitan
sebagai proksi. Analisis potong lintang pertama terhadap seluruh PUR
tersebut menemukan bukti bahwa ketiga dimensi integrasi berarah positif,
dengan integrasi fungsional paling kuat signifikansinya. Untuk integrasi
institusional, keberadaan suatu bentuk kerja sama metropolitan dinilai
lebih penting daripada bentuk persisnya. Artikel menyimpulkan bahwa
hasilnya mendukung asumsi teoretis bahwa jaringan dapat menggantikan
sebagian fungsi kedekatan fisik. (Locator: PDF p.~1, bagian ABSTRACT;
TXT baris 40-55.)

\subsection{Problem Statement}\label{problem-statement-44}

{[}Laporan sumber{]} Literatur aglomerasi yang diringkas penulis sering
mengaitkan ekonomi urbanisasi dengan ukuran atau kepadatan kota. Namun,
PUR yang terdiri atas beberapa kota menengah tidak selalu menghasilkan
tingkat manfaat aglomerasi yang sebanding dengan satu kota besar dengan
jumlah penduduk gabungan yang sama. Artikel mengutip temuan bahwa dua
kota berdekatan yang masing-masing berpenduduk setengah juta tidak
otomatis dapat menyediakan manfaat seperti satu kota berpenduduk satu
juta, baik dalam fungsi budaya, rekreasi, olahraga, ritel
terspesialisasi, maupun ekonomi urbanisasi secara umum. Pada saat yang
sama, kinerja antar-PUR berbeda dan perbedaan itu memerlukan penjelasan.
(Locator: PDF pp.~1-3, INTRODUCTION; TXT baris 59-130.)

{[}Interpretasi penulis{]} Penulis mengajukan integrasi antar-kota
sebagai salah satu penjelasan penting. Interaksi diperlukan untuk
berbagi, mencocokkan, dan belajar; kota-kota yang secara fisik terpisah
tetapi terikat secara fungsional, institusional, dan kultural
diperkirakan lebih menyerupai satu aglomerasi besar. Dengan demikian,
integrasi diposisikan sebagai mekanisme yang mungkin mengatasi kerugian
PUR yang terfragmentasi. (Locator: PDF pp.~2-3, INTRODUCTION; TXT baris
112-159.)

{[}Laporan sumber{]} Pertanyaan penelitian yang memandu artikel adalah:
``does stronger integration between cities in Polycentric Urban Regions
enable them to organise more urbanisation economies?'' Artikel secara
sengaja memusatkan perhatian pada PUR yang polisentris secara
morfologis, tanpa mensyaratkan bahwa PUR tersebut telah diidentifikasi
sebelumnya sebagai entitas metropolitan yang koheren atau memiliki batas
administratif tertentu. (Locator: PDF p.~3, INTRODUCTION; TXT baris
137-163.)

{[}Inferensi pengolahan{]} Masalah yang diuji bukan sekadar apakah suatu
wilayah memiliki struktur banyak-pusat, melainkan apakah struktur
tersebut dapat dikonversi menjadi manfaat aglomerasi melalui hubungan
antar-kota. Karena variabel hasilnya adalah proksi fungsi metropolitan,
masalah kesejahteraan atau produktivitas langsung tidak diuji secara
langsung dalam laporan ini. (Locator: PDF pp.~3 dan 8-10; TXT baris
137-159, 482-563.)

\subsection{Objectives}\label{objectives-44}

{[}Laporan sumber{]} Tujuan utama yang dinyatakan atau dijelaskan dalam
artikel adalah:

\begin{itemize}
\tightlist
\item
  mengidentifikasi dan mendefinisikan secara konsisten PUR di Eropa dari
  perspektif morfologis;
\item
  mengukur tiga dimensi integrasi antar-kota, yaitu fungsional,
  institusional, dan kultural;
\item
  membangun ukuran kinerja PUR berdasarkan tingkat fungsi metropolitan
  yang diamati dibandingkan dengan tingkat yang diharapkan jika
  kota-kotanya bekerja sebagai satu aglomerasi;
\item
  menguji apakah integrasi yang lebih kuat berkaitan dengan kinerja PUR
  yang lebih baik, baik untuk indikator individual maupun indeks
  gabungan;
\item
  memeriksa apakah ada hubungan penguatan antar-dimensi integrasi serta
  apakah hubungan integrasi-kinerja berbeda menurut ukuran PUR, lokasi
  di Eropa Timur, atau status lintas batas.
\end{itemize}

{[}Inferensi pengolahan{]} Lima tujuan di atas dirangkum dari pertanyaan
penelitian, uraian pendekatan, dan urutan model yang dijelaskan penulis;
artikel tidak menyajikannya sebagai daftar tujuan formal tersendiri.
(Locator: PDF pp.~3-5 dan 12-18, INTRODUCTION, RESEARCH APPROACH,
RESULTS; TXT baris 151-187, 264-318, 702-946.)

\subsection{Literature Survey}\label{literature-survey-44}

\begin{itemize}
\tightlist
\item
  \textbf{Ekonomi aglomerasi.} {[}Laporan sumber{]} Artikel membedakan
  ekonomi lokalisasi, yang muncul dari konsentrasi perusahaan dalam
  industri yang sama, dari ekonomi urbanisasi, yang diperoleh dari kota
  yang besar dan beragam. Ekonomi urbanisasi diringkas melalui arus
  pengetahuan dan informasi lintas industri, pasar tenaga kerja yang
  terdiversifikasi dan terspesialisasi, infrastruktur bersama, jasa
  bisnis terspesialisasi, dan amenitas konsumen. Penulis juga merangkum
  temuan bahwa ukuran atau kepadatan sering berkaitan dengan
  produktivitas dan premi upah perkotaan, meskipun efeknya bervariasi
  menurut sektor, negara, dan pilihan pemodelan. (Locator: PDF pp.~1-2,
  INTRODUCTION; TXT baris 59-97.)
\item
  \textbf{Jaringan dan kedekatan.} {[}Laporan sumber{]} Penulis
  menggunakan literatur yang menyatakan bahwa ekonomi jaringan dapat
  menggantikan sebagian ekonomi aglomerasi atau kedekatan. Literatur
  tersebut menjadi dasar bagi gagasan bahwa jaringan kota-kota kecil dan
  menengah dapat menjadi alternatif terhadap konsentrasi lebih lanjut di
  megakota, terutama ketika pasokan perumahan tidak elastis atau
  terdapat hambatan historis dan institusional terhadap pertumbuhan kota
  besar. (Locator: PDF pp.~2-3, INTRODUCTION; TXT baris 98-130.)
\item
  \textbf{PUR dan kinerja.} {[}Laporan sumber{]} PUR ditempatkan dalam
  literatur tentang gugus kota yang berbeda secara administratif tetapi
  berdekatan. Studi terdahulu yang diringkas meliputi studi kasus PUR
  tertentu dan pemodelan kuantitatif yang mengukur mono-polisentrisitas
  atau polisentrisitas pada wilayah fungsional atau administratif.
  Penulis menilai bahwa belum ada daftar komprehensif PUR Eropa yang
  sesuai dengan kepentingan mereka pada bentuk morfologis sekaligus
  memungkinkan pengujian hubungan dengan integrasi. (Locator: PDF
  pp.~3-4, INTRODUCTION dan RESEARCH APPROACH; TXT baris 144-187,
  252-318.)
\item
  \textbf{Integrasi fungsional.} {[}Laporan sumber{]} Literatur yang
  dibahas menekankan dampak fragmentasi ketika PUR hanya menjadi
  himpunan kota menengah yang terpisah. Fragmentasi dapat menghambat
  terbentuknya lingkungan metropolitan, pasar perumahan dan tenaga kerja
  yang efisien, serta hubungan transportasi yang terjangkau. Sebaliknya,
  hubungan komplementer antarkota dapat memperluas pasar, masukan
  pengetahuan, infrastruktur, dan peluang subkontrak. Fungsi yang
  tersebar di banyak kota juga menciptakan permintaan transportasi dan
  interaksi. (Locator: PDF pp.~3-4, LINKING INTEGRATION TO PERFORMANCE;
  TXT baris 160-226.)
\item
  \textbf{Integrasi institusional.} {[}Laporan sumber{]} Kerja sama
  metropolitan, penggabungan munisipalitas, atau kolaborasi
  antarmunisipalitas dibahas sebagai cara untuk berbagi sumber daya,
  mengoordinasikan infrastruktur dan tata guna lahan, mengurangi
  kompetisi redundan, serta memperoleh posisi tawar yang lebih kuat
  terhadap tingkat pemerintahan yang lebih tinggi. Literatur yang
  diringkas juga mengaitkan fragmentasi tata kelola dengan produktivitas
  yang lebih rendah. (Locator: PDF pp.~4-5, LINKING INTEGRATION TO
  PERFORMANCE; TXT baris 227-258.)
\item
  \textbf{Integrasi kultural.} {[}Laporan sumber{]} Penulis
  menghubungkan kerja sama historis dan mobilitas dengan terbentuknya
  identitas metropolitan. Identitas bersama dipahami bukan hanya sebagai
  kesadaran fungsional atas ruang ekonomi yang saling terkait, tetapi
  juga sebagai keterikatan emosional. Sebaliknya, jarak mental,
  perbedaan budaya, perbedaan orientasi politik, dan hambatan bahasa
  dapat mengurangi kemungkinan membangun hubungan ekonomi atau strategi
  bersama. (Locator: PDF p.~4, LINKING INTEGRATION TO PERFORMANCE; TXT
  baris 201-226.)
\item
  \textbf{Kerangka konseptual.} {[}Laporan sumber{]} Tiga dimensi
  tersebut diturunkan dari kerangka Kloosterman dan Musterd (2001).
  Penulis mengakui bahwa ketiganya saling terkait dan dapat saling
  memperkuat atau membatasi, lalu menguji bukan hanya hubungan langsung
  masing-masing dimensi dengan kinerja, tetapi juga kemungkinan umpan
  balik antardimensi. (Locator: PDF pp.~4-5 dan 10-12; TXT baris
  231-249, 580-597, 853-907.)
\end{itemize}

{[}Inferensi pengolahan{]} Survei literatur membangun argumen teoritis
yang konsisten untuk arah hubungan positif, tetapi artikel juga
menyatakan bahwa asumsi jaringan sebagai pengganti kedekatan tersebut
sebelumnya belum dibuktikan secara komprehensif pada seluruh PUR Eropa.
Literatur yang dipakai terutama berfungsi sebagai landasan mekanisme dan
perumusan hipotesis, bukan sebagai meta-analisis baru. (Locator: PDF
pp.~2-4; TXT baris 123-159, 190-258.)

\subsection{Method Used}\label{method-used-44}

\begin{itemize}
\tightlist
\item
  \textbf{Desain dan unit analisis.} {[}Laporan sumber{]} Penelitian
  menggunakan pendekatan kuantitatif potong lintang terhadap 117 PUR
  Eropa. Hubungan antara indikator integrasi dan kategori kinerja
  dianalisis dengan regresi logistik berurutan (\emph{ordered logistic
  regression}). Artikel menyajikan model untuk indikator individual,
  indeks dimensi, indeks gabungan, variabel kontrol, dan interaksi.
  (Locator: PDF pp.~3, 12-18; TXT baris 160-187, 702-718, 927-953.)
\item
  \textbf{Identifikasi PUR.} {[}Laporan sumber{]} Kota dibatasi sebagai
  \emph{morphological urban areas} (MUA), yaitu seluruh munisipalitas
  yang membentuk kawasan terbangun kontigu menurut ESPON 1.4.3. PUR
  hanya mencakup wilayah dengan sedikitnya dua aglomerasi dan setiap
  aglomerasi memiliki paling sedikit 40.000 penduduk. Beberapa batas
  wilayah digunakan: FUA dan hinterland berbasis basin komuter,
  \emph{polyFUA} yang menggabungkan FUA kontigu berdasarkan ukuran dan
  jarak, \emph{suprapolyFUA} untuk contoh Rhein-Ruhr dan Randstad, serta
  \emph{potential integration areas} (PIA) yang menggabungkan FUA ketika
  isokron 45 menitnya tumpang tindih sedikitnya 33 persen. Untuk
  mencegah PIA yang terlalu luas, semua kota inti harus berjarak
  maksimal 60 menit dan sedikitnya dua kota inti harus berjarak maksimal
  45 menit. (Locator: PDF pp.~5-6, RESEARCH APPROACH; TXT baris
  264-318.)
\item
  \textbf{Ukuran polisentrisitas.} {[}Laporan sumber{]} Penulis
  memperkenalkan indeks Herfindahl sebagai ukuran distribusi ukuran
  kota. Rumusnya adalah H = jumlah kuadrat pangsa penduduk setiap kota,
  dengan pangsa dihitung dari total penduduk kota-kota di wilayah dan N
  sebagai jumlah kota. Nilai lebih rendah berarti lebih polisentris.
  Ambang inklusi ditetapkan pada 0,56 melalui pertimbangan praktis dan
  ``common sense'', agar kasus yang dianggap jelas monosentris tidak
  masuk. (Locator: PDF p.~5, RESEARCH APPROACH; TXT baris 284-304.)
\item
  \textbf{Pembentukan ukuran kinerja.} {[}Laporan sumber{]} Kinerja
  diukur dari kehadiran fungsi metropolitan memakai basis data BBSR.
  Skor fungsi kota-kota dalam setiap PUR diagregasikan, tingkat yang
  diharapkan dihitung berdasarkan persamaan regresi untuk kota tunggal,
  lalu tingkat aktual dibandingkan dengan tingkat yang diharapkan.
  Kategori hasil disusun menjadi empat tingkat berurutan: secara
  signifikan lebih rendah dari prediksi, lebih rendah tetapi tidak
  signifikan, lebih tinggi dari prediksi, dan secara signifikan lebih
  tinggi dari prediksi. (Locator: PDF pp.~8-10, Measuring the
  performance of PURs; TXT baris 482-563.)
\item
  \textbf{Kalibrasi model kota tunggal.} {[}Laporan sumber{]} Karena
  sebagian MUA tidak mempunyai fungsi metropolitan sama sekali, penulis
  memakai regresi beta yang diinflasi nol, yang terdiri atas bagian
  logistik untuk probabilitas skor nol dan bagian beta untuk proporsi
  antara 0 dan 1. Model pada 1.947 MUA memasukkan ukuran kota, ukuran
  hinterland, status ibu kota, PDB per kapita, keterlekatan dalam
  jaringan politik internasional, konektivitas bandara, pariwisata, dan
  dummy negara. (Locator: PDF p.~9, Table 2; TXT baris 509-563.)
\item
  \textbf{Pengukuran integrasi.} {[}Laporan sumber{]} Integrasi
  fungsional diwakili efisiensi koneksi jalan, efisiensi koneksi kereta,
  dan frekuensi kereta antar setiap pasangan kota. Integrasi
  institusional diwakili keberadaan badan metropolitan yang mencakup
  lebih dari 50 persen PUR, jumlah tahun badan aktif, dan tipe kemitraan
  dari kesepakatan informal sampai otoritas metropolitan terpilih.
  Integrasi kultural diwakili homogenitas preferensi politik wali kota
  dan homogenitas bahasa. Semua indikator dinormalisasi dan
  diagregasikan menjadi indeks dengan skor lebih tinggi berarti
  integrasi lebih besar. (Locator: PDF pp.~10-12, Table 3; TXT baris
  580-702.)
\item
  \textbf{Pengujian tambahan.} {[}Laporan sumber{]} Model agregat
  memasukkan ukuran PUR, dummy Eropa Timur, dan derajat lintas batas.
  Penulis juga menghitung fragmentasi administratif-teritorial pada
  tingkat NUTS 2 dan NUTS 3, serta menguji interaksi integrasi dengan
  ukuran, lokasi Eropa Timur, dan derajat lintas batas. Interaksi
  antardimensi fungsional, institusional, dan kultural juga diuji;
  variabel untuk model interaksi dipusatkan pada rata-rata, kecuali
  dummy Eropa Timur. (Locator: PDF pp.~14-17; TXT baris 811-841,
  853-953.)
\end{itemize}

\subsection{Dataset}\label{dataset-44}

\begin{itemize}
\tightlist
\item
  {[}Laporan sumber{]} Dataset fungsi metropolitan berasal dari basis
  data \emph{Metropolitan Areas in Europe} milik BBSR (2011). Domainnya
  mencakup sains, ekonomi, budaya, dan olahraga. Data fungsi individual
  dikumpulkan untuk periode 2004-2009, dengan mayoritas observasi
  bersesuaian dengan 2008. (Locator: PDF pp.~8-9; TXT baris 482-505.)
\item
  {[}Laporan sumber{]} Variabel konektivitas fungsional memakai Google
  Maps untuk koneksi jalan dan, bersama Deutsche Bahn, untuk koneksi
  kereta. Frekuensi kereta dihitung dari jumlah rata-rata kereta antar
  setiap pasangan kota pada pukul 08.00-20.00. (Locator: PDF pp.~10-11,
  Table 3; TXT baris 586-619.)
\item
  {[}Laporan sumber{]} Informasi kerja sama metropolitan dikumpulkan
  melalui pencarian internet. Informasi homogenitas preferensi politik
  berasal dari situs web, ikhtisar per negara, dan basis data pemilu;
  homogenitas bahasa didasarkan pada peta ``Languages of Europe''.
  Populasi, MUA, dan beberapa batas wilayah mengacu pada proyek ESPON
  1.4.3 dan ESPON 1.1.1. (Locator: PDF pp.~5-8 dan 10-12; TXT baris
  264-318, 621-656.)
\item
  {[}Laporan sumber{]} Daftar Table 1 memuat 117 PUR, negara yang
  terlibat, jumlah MUA, populasi, dan indeks polisentrisitas. Penulis
  melaporkan hampir 122 juta orang tinggal di PUR dan mengaitkan angka
  itu dengan 25 persen populasi EU; frasa TXT juga menyebut ``Norway and
  Switzerland'' tanpa menjelaskan kembali relasi statistiknya. Catatan
  sumber menyebut data ESPON tidak tersedia untuk sejumlah negara atau
  wilayah tertentu, tetapi rincian cakupan lengkap dataset tidak
  diberikan di luar catatan tersebut. (Locator: PDF pp.~6-8, Table 1 dan
  catatan; TXT baris 322-479, 482.)
\item
  {[}Laporan sumber{]} Penulis menyatakan dataset disediakan agar
  analisis dapat diuji dengan indikator lain, tetapi lokasi, URL,
  format, dan isi berkas dataset tidak disebutkan dalam TXT. (Locator:
  PDF p.~18, DISCUSSION AND CONCLUSION; TXT baris 1020-1046.)
\item
  {[}Inferensi pengolahan{]} Terdapat ketidaksesuaian internal yang
  perlu diverifikasi pada dataset asli: catatan 2 menyebut beberapa
  pengecualian termasuk ``the Baltic states'', sementara Table 1
  menampilkan kode negara ``EE'' pada baris Narva-Kohtla Jarve. Review
  ini tidak menafsirkan cakupan geografis tersebut di luar apa yang
  tertulis dalam TXT. (Locator: PDF pp.~7 dan 19, Table 1 dan Notes; TXT
  baris 367, 1079-1082.)
\end{itemize}

\subsection{Population / Sample}\label{population-sample-44}

\begin{itemize}
\tightlist
\item
  {[}Laporan sumber{]} Populasi analitis artikel adalah PUR yang
  diidentifikasi secara morfologis di Eropa. Seluruh 117 PUR dalam
  daftar dipakai dalam model integrasi-kinerja; Table 7 sampai Table 10
  masing-masing menunjukkan 117 observasi. (Locator: PDF pp.~6-8 dan
  13-17; TXT baris 322-479, 775, 802, 883, 948.)
\item
  {[}Laporan sumber{]} Kriteria pembentukan mencakup sedikitnya dua MUA,
  setiap MUA sekurang-kurangnya 40.000 penduduk, dan indeks Herfindahl
  paling tinggi 0,56. PIA juga dibatasi oleh ambang waktu perjalanan
  agar tidak menjadi wilayah yang berlebihan ukurannya. PUR lintas batas
  tetap dimasukkan. (Locator: PDF pp.~5-8; TXT baris 264-318, 322-479.)
\item
  {[}Laporan sumber{]} Sebanyak 1.947 MUA digunakan untuk membangun
  model dasar prediksi fungsi metropolitan pada kota tunggal. MUA
  tersebut adalah basis kalibrasi pembanding, bukan jumlah PUR dalam
  analisis utama. (Locator: PDF p.~9, Table 2; TXT baris 509-563.)
\item
  {[}Inferensi pengolahan{]} Desain ini lebih dekat dengan cakupan
  seluruh kasus yang berhasil didefinisikan daripada sampel acak. Karena
  itu, tidak ada dasar dalam TXT untuk menyebutnya sebagai sampel
  probabilistik atau untuk menghitung galat representasi populasi PUR
  Eropa yang lebih luas. (Locator: PDF pp.~5-8 dan 13-17; TXT baris
  264-479, 775, 802, 883, 948.)
\item
  {[}Batas provenance{]} Jumlah kandidat PUR sebelum penyaringan, jumlah
  observasi yang hilang, aturan penanganan data hilang, dan apakah ada
  PUR yang dikeluarkan selain berdasarkan kriteria definisi yang
  dijelaskan: Tidak disebutkan dalam file. (Locator: PDF pp.~5-8 dan
  13-17; TXT baris 264-479, 775, 802, 883, 948.)
\end{itemize}

\subsection{Dependent Variables}\label{dependent-variables-44}

{[}Laporan sumber{]} Variabel dependen adalah ``performance'' PUR, yaitu
kategori berurutan yang menunjukkan seberapa dekat atau jauh kehadiran
fungsi metropolitan aktual dari tingkat yang diprediksi jika seluruh PUR
berfungsi sebagai satu aglomerasi. Prediksi berasal dari model kota
tunggal dan perbandingan dilakukan setelah skor fungsi kota-kota dalam
PUR diagregasikan. Empat kategori performa adalah: signifikan di bawah
prediksi pada p \textless{} 0,05; di bawah prediksi tetapi tidak
signifikan; di atas prediksi; dan signifikan di atas prediksi. (Locator:
PDF pp.~8-10, Measuring the performance of PURs; TXT baris 482-563.)

{[}Laporan sumber{]} Indeks fungsi metropolitan menggabungkan fungsi
dalam empat domain: sains, ekonomi, budaya, dan olahraga. Contohnya
meliputi universitas besar dan organisasi riset internasional; kantor
pusat perusahaan Fortune-500, jasa produsen, bank, dan pameran; konser,
pameran seni, festival film, teater, opera, galeri, dan museum; serta
stadion, lokasi Olimpiade, dan acara olahraga besar. (Locator: PDF
pp.~8-9; TXT baris 482-505.)

{[}Inferensi pengolahan{]} Outcome ini mengukur kinerja relatif terhadap
ekspektasi fungsi metropolitan, bukan produktivitas, pendapatan,
kesejahteraan, atau efisiensi ekonomi secara langsung. Karena bentuknya
ordinal dan berbasis perbandingan dengan model pembanding, koefisien
ordered logit harus dibaca sebagai hubungan dengan kategori kinerja
tersebut, bukan sebagai kenaikan langsung pada produktivitas atau jumlah
fungsi. (Locator: PDF pp.~8-10 dan 12-18; TXT baris 482-563, 702-718.)

\subsection{Results}\label{results-44}

\begin{itemize}
\tightlist
\item
  \textbf{Deskripsi awal.} {[}Laporan sumber{]} Tingkat kinerja berbeda
  di sebagian besar negara pada Figure 1 dan tidak memperlihatkan pola
  spasial yang jelas, kecuali kecenderungan kinerja yang agak lebih
  lemah di Eropa Timur. (Locator: PDF p.~10, Figure 1; TXT baris
  571-578.)
\item
  \textbf{Statistik deskriptif integrasi.} {[}Laporan sumber{]} Table 4
  melaporkan rata-rata efisiensi jalan 1,12, efisiensi kereta 0,65,
  frekuensi kereta 23,56, keberadaan badan metropolitan 0,48, tahun
  aktif 4,74, tipe kemitraan 1,18, homogenitas preferensi politik 0,60,
  dan homogenitas bahasa 0,86. Rentang yang dilaporkan berturut-turut
  adalah 0,30-1,59; 0-2,63; 0-98; 0-1; 0-41; 0-4; 0,31-1; dan 0-1.
  (Locator: PDF p.~12, Table 4; TXT baris 666-677.)
\item
  \textbf{Integrasi fungsional.} {[}Laporan sumber{]} Pada Table 7,
  efisiensi jalan bernilai 1,204 dengan galat baku 0,901 dan tidak
  ditandai signifikan. Efisiensi kereta bernilai 0,840 dengan galat baku
  0,440 dan signifikan pada p \textless{} 0,10. Frekuensi koneksi kereta
  bernilai 0,028 dengan galat baku 0,009 dan signifikan pada p
  \textless{} 0,01; dalam model gabungan indikator, frekuensi tetap
  signifikan dengan koefisien 0,027 dan galat baku 0,009. Indeks
  integrasi fungsional bernilai 0,931 dengan galat baku 0,293 dan
  signifikan pada p \textless{} 0,01. Model gabungan indikator memiliki
  LR chi-square 13,46 dan pseudo-R2 0,0512, sedangkan model indeks
  memiliki LR chi-square 11,01 dan pseudo-R2 0,0419. (Locator: PDF
  p.~13, Table 7; TXT baris 752-780.)
\item
  \textbf{Integrasi institusional.} {[}Laporan sumber{]} Pada model
  individual Table 8, keberadaan badan metropolitan memiliki koefisien
  0,816 dengan galat baku 0,358 dan signifikan pada p \textless{} 0,05.
  Jumlah tahun aktif bernilai 0,039 dengan galat baku 0,022 dan
  signifikan pada p \textless{} 0,10. Tipe kemitraan bernilai 0,222
  dengan galat baku 0,130 dan signifikan pada p \textless{} 0,10. Indeks
  integrasi institusional bernilai 0,423 dengan galat baku 0,198 dan
  signifikan pada p \textless{} 0,05. (Locator: PDF p.~14, Table 8; TXT
  baris 787-807.)
\item
  \textbf{Model institusional simultan.} {[}Laporan sumber{]} Presence
  of a metropolitan body dan type of partnership sangat berkorelasi pada
  Table 5, dengan korelasi 0,91 dan signifikansi p \textless{} 0,01.
  Karena itu, penulis menyebut model simultan yang memasukkan keduanya
  kurang relevan. Pada model tersebut, koefisien keberadaan badan tidak
  lagi ditandai signifikan, jumlah tahun aktif bernilai 0,010 dan tidak
  signifikan, sedangkan tipe kemitraan ditampilkan dalam TXT sebagai
  ``20.249'' dengan galat baku 0,303 dan tidak signifikan. (Locator: PDF
  pp.~13-14, Tables 5 dan 8; TXT baris 724-737, 793-807.)
\item
  \textbf{Integrasi kultural.} {[}Laporan sumber{]} Homogenitas
  preferensi politik, homogenitas bahasa, dan indeks integrasi kultural
  tidak signifikan dalam model-model Table 9. Koefisien yang terlihat
  dalam TXT antara lain 1,304 dengan galat baku 0,813 untuk preferensi
  politik, 0,439 dengan galat baku 0,540 untuk bahasa, dan 0,379 dengan
  galat baku 0,239 untuk indeks kultural; tidak satu pun diberi tanda
  signifikansi. LR chi-square model-model tersebut adalah 2,58, 0,68,
  2,90, dan 2,61, dengan pseudo-R2 antara 0,003 dan 0,011. (Locator: PDF
  p.~15, Table 9; TXT baris 871-888.)
\item
  \textbf{Indeks gabungan dan kontrol.} {[}Laporan sumber{]} Pada Table
  10, indeks fungsional tetap positif dan signifikan pada kelima model:
  0,782 (p \textless{} 0,01) pada Model 16, lalu 0,771, 0,776, 0,769,
  dan 0,785 pada Model 17 sampai 20, masing-masing tetap signifikan
  sekurang-kurangnya pada p \textless{} 0,05. Indeks institusional
  bernilai positif tetapi tidak signifikan pada kelima model. Indeks
  kultural tidak signifikan pada Model 16, menjadi signifikan lemah pada
  p \textless{} 0,10 setelah kontrol ditambahkan pada Model 17, dan
  kembali tidak konsisten signifikan pada model interaksi berikutnya.
  (Locator: PDF pp.~16-17, Table 10; TXT baris 927-953.)
\item
  \textbf{Kontrol wilayah.} {[}Laporan sumber{]} Model dengan kontrol
  menunjukkan arah hubungan negatif antara ukuran PUR dan kinerja,
  dengan beberapa koefisien signifikan pada p \textless{} 0,10. Dummy
  Eropa Timur juga berarah negatif dan signifikan lemah, atau signifikan
  pada p \textless{} 0,05 dalam sebagian model. Derajat lintas batas
  berarah positif tetapi tidak signifikan. Penulis menafsirkan ini
  sebagai indikasi bahwa PUR yang lebih besar lebih sulit memanfaatkan
  massa kritis, PUR Eropa Timur cenderung berkinerja lebih rendah, dan
  kompleksitas lintas batas tidak menghasilkan kinerja yang lebih lemah
  secara statistik. (Locator: PDF pp.~16-17; TXT baris 892-910,
  934-950.)
\item
  \textbf{Kecocokan model dan interaksi.} {[}Laporan sumber{]}
  Penambahan tiga indeks integrasi dan kontrol meningkatkan pseudo-R2
  dari 0,056 pada Model 16 menjadi 0,102 pada Model 17 dan sekitar
  0,095-0,103 pada Model 18-20; LR chi-square seluruh model Table 10
  signifikan pada p \textless{} 0,01. Namun, tidak ada interaksi yang
  signifikan antara integrasi fungsional-institusional,
  fungsional-kultural, atau institusional-kultural. Model tambahan yang
  tidak dilaporkan juga tidak menemukan interaksi signifikan antara
  integrasi dan ukuran PUR, lokasi Eropa Timur, atau derajat lintas
  batas. (Locator: PDF pp.~16-17, Table 10 dan pembahasan; TXT baris
  853-953.)
\item
  \textbf{Fragmentasi administratif.} {[}Laporan sumber{]} Penulis
  menemukan hubungan bahwa fragmentasi administratif yang lebih besar
  berkaitan dengan lebih banyak kerja sama pada Figure 2. Akan tetapi,
  model yang menjelaskan kinerja PUR dengan fragmentasi
  administratif-teritorial tidak dilaporkan secara rinci dan indikator
  fragmentasi tersebut disebut tetap jauh dari signifikan. (Locator: PDF
  pp.~14-15, Figure 2; TXT baris 811-841, 849-850.)
\end{itemize}

\subsection{Findings}\label{findings-44}

{[}Interpretasi penulis{]} Temuan inti artikel adalah bahwa integrasi
fungsional merupakan dimensi yang paling konsisten terkait dengan
kemampuan PUR mengorganisasikan ekonomi urbanisasi. Mobilitas dan
koneksi antar-kota, terutama frekuensi kereta, dipandang lebih dekat
dengan permintaan dan arus aktual daripada sekadar efisiensi koneksi.
(Locator: PDF pp.~12-18, RESULTS dan DISCUSSION AND CONCLUSION; TXT
baris 702-780, 970-982.)

{[}Interpretasi penulis{]} Integrasi institusional juga dikaitkan dengan
kinerja yang lebih baik, tetapi ukuran paling penting tampaknya adalah
ada atau tidaknya badan kerja sama metropolitan yang bekerja. Bentuk
formal dan tingkat kewenangan yang persis tidak tampak sepenting
keberadaan kerja sama; durasi kerja sama memberi indikasi positif yang
lebih lemah. (Locator: PDF pp.~14 dan 18; TXT baris 811-828, 984-1004.)

{[}Interpretasi penulis{]} Bukti untuk integrasi kultural lebih lemah
dan bergantung pada spesifikasi model. Indikator individual tidak
signifikan, sedangkan beberapa model agregat dengan kontrol memberikan
petunjuk positif pada tingkat p \textless{} 0,10. Tidak ada bukti
kuantitatif bahwa ketiga dimensi secara otomatis memperkuat satu sama
lain dalam suatu ``upward spiral''. (Locator: PDF pp.~15-18; TXT baris
823-828, 892-907, 1006-1018.)

{[}Interpretasi penulis{]} Penulis menyatakan bahwa hubungan
integrasi-kinerja tidak berbeda secara signifikan menurut ukuran PUR,
lokasi di Eropa, atau status lintas batas. Mereka menghubungkan hasil
keseluruhan dengan gagasan bahwa jaringan dapat menggantikan kedekatan
dalam mengorganisasikan manfaat aglomerasi dan dengan proses
``metropolisation'', yaitu terbentuknya kualitas ekonomi, fungsional,
administratif, dan sosio-spasial pada skala PUR. (Locator: PDF
pp.~17-18; TXT baris 918-930, 970-975, 1020-1030.)

{[}Inferensi pengolahan{]} Karena hasil utama berasal dari desain potong
lintang dan outcome berupa kategori kinerja relatif, temuan ini
mendukung asosiasi yang dilaporkan, bukan pembuktian bahwa integrasi
secara kausal menghasilkan fungsi metropolitan. TXT tidak melaporkan
strategi identifikasi kausal. (Locator: PDF pp.~3 dan 12-18; TXT baris
160-187, 702-953.)

\subsection{Conclusions}\label{conclusions-44}

{[}Laporan sumber{]} Dalam bagian diskusi dan kesimpulan, penulis
merangkum bahwa:

\begin{itemize}
\tightlist
\item
  semakin kuat integrasi fungsional kota-kota dalam PUR, semakin baik
  kinerjanya dalam mengorganisasikan ekonomi urbanisasi;
\item
  integrasi institusional atau tata kelola metropolitan berdampak
  positif, tetapi efeknya lebih kecil daripada integrasi fungsional;
  keberadaan suatu bentuk kerja sama lebih penting daripada bentuk dan
  cakupannya, sementara durasi kerja sama hanya memiliki dukungan
  indikatif;
\item
  beberapa model memberi petunjuk bahwa integrasi kultural juga
  berdampak positif;
\item
  peningkatan timbal balik antarberbagai bentuk integrasi tidak berhasil
  dibuktikan secara empiris;
\item
  hubungan integrasi-kinerja tidak terbukti berbeda menurut ukuran,
  lokasi Eropa, atau status lintas batas PUR.
\end{itemize}

{[}Interpretasi penulis{]} Penulis menyimpulkan bahwa tantangan utama
PUR adalah berpindah dari fragmentasi menuju integrasi. Mereka mendorong
integrasi fungsional, investasi konektivitas dan angkutan umum
antar-kota, serta proses \emph{region-building} yang lebih luas.
Investasi pada metropolisasi PUR diposisikan sebagai alternatif yang
layak dan menguntungkan terhadap konsentrasi pembangunan di kawasan ibu
kota. (Locator: PDF p.~18, DISCUSSION AND CONCLUSION; TXT baris
970-1046.)

\subsection{Contributions}\label{contributions-44}

\begin{itemize}
\tightlist
\item
  {[}Laporan sumber{]} Kontribusi empiris yang diklaim penulis adalah
  daftar komprehensif pertama berisi 117 PUR Eropa dengan definisi
  morfologis dan pengukuran yang konsisten. Daftar tersebut mencantumkan
  negara, kota/MUA, ukuran penduduk, dan indeks polisentrisitas.
  (Locator: PDF pp.~3 dan 6-8; TXT baris 151-160, 299-479.)
\item
  {[}Inferensi pengolahan{]} Secara metodologis, artikel menghubungkan
  ukuran polisentrisitas morfologis, ukuran integrasi tiga dimensi, dan
  ukuran kinerja berbasis fungsi metropolitan dalam satu analisis
  komparatif lintas Eropa. Artikel juga menggunakan model kota tunggal
  sebagai dasar menghitung fungsi yang diharapkan pada PUR. (Locator:
  PDF pp.~5-18; TXT baris 264-563, 580-702, 702-953.)
\item
  {[}Interpretasi penulis{]} Kontribusi teoretisnya adalah memberikan
  dukungan empiris pada asumsi bahwa ekonomi jaringan dapat menggantikan
  kedekatan, setidaknya dalam arti kemampuan mengorganisasikan fungsi
  metropolitan. Hasilnya juga memperluas pembahasan dari sekadar
  polisentrisitas menuju proses integrasi atau metropolisasi. (Locator:
  PDF pp.~3 dan 18; TXT baris 137-159, 970-975.)
\item
  {[}Laporan sumber{]} Kontribusi kebijakan yang dikemukakan adalah
  menunjukkan bahwa penguatan hubungan antar-kota dan kerja sama
  metropolitan dapat menjadi alternatif terhadap strategi konsentrasi
  pada kawasan ibu kota, terutama karena banyak penduduk Eropa tinggal
  di PUR. (Locator: PDF p.~18; TXT baris 1032-1046.)
\item
  {[}Catatan verifikasi{]} Klaim ``first'' dan ``first-ever'' merupakan
  klaim kebaruan yang dibuat artikel. Klaim tersebut tidak diverifikasi
  dengan sumber di luar TXT. (Locator: PDF pp.~1, 3, dan 18; TXT baris
  40-52, 151-160, 987-1005.)
\end{itemize}

\subsection{Research Gap}\label{research-gap-44}

{[}Laporan sumber{]} Gap sebelum penelitian ini adalah belum adanya
daftar komprehensif PUR Eropa. Penulis mengaitkannya dengan
ketidakjelasan konseptual dan perdebatan apakah polisentrisitas hanya
menyangkut aspek morfologis atau juga relasi antar-pusat. Studi yang
tersedia lebih sering berupa studi kasus PUR tertentu atau pemodelan
kuantitatif pada wilayah fungsional dan administratif; asumsi bahwa
integrasi dapat mengubah jaringan kota menjadi sumber manfaat aglomerasi
disebut masih luas tetapi belum tersubstansiasi. (Locator: PDF pp.~3-5;
TXT baris 144-187, 252-318.)

{[}Inferensi pengolahan{]} Artikel menanggapi gap tersebut dengan
membatasi identifikasi awal pada polisentrisitas morfologis, kemudian
menguji secara komparatif apakah integrasi fungsional, institusional,
dan kultural menjelaskan variasi kinerja 117 PUR. Strategi ini menjawab
kebutuhan akan perbandingan yang luas, tetapi tidak menyelesaikan
kebutuhan akan pengukuran yang sepenuhnya menangkap relasi dan konteks
tiap wilayah. (Locator: PDF pp.~3-5 dan 10-18; TXT baris 151-187,
264-318, 580-702, 702-953.)

{[}Inferensi pengolahan{]} Gap yang masih terbuka adalah arah kausal,
dinamika waktu, serta validitas proksi integrasi dan kinerja. Kesimpulan
ini ditarik dari desain potong lintang, perbedaan tahun data, dan agenda
riset yang disebutkan penulis; TXT tidak menyatakan bahwa studi ini
telah mengidentifikasi mekanisme kausal. (Locator: PDF pp.~3, 10, dan
18; TXT baris 160-187, 571-585, 1006-1046.)

\subsection{Limitations}\label{limitations-44}

\begin{itemize}
\tightlist
\item
  {[}Laporan sumber{]} Keterbatasan data mengharuskan penggunaan
  indikator yang tersedia pada tingkat kota dan dapat dikumpulkan dengan
  cakupan seluruh Eropa. Penulis mengakui indikator tersebut tidak
  menangkap seluruh kompleksitas integrasi fungsional, institusional,
  dan kultural. (Locator: PDF pp.~10-12; TXT baris 580-597, 680-718.)
\item
  {[}Laporan sumber{]} Data fungsi metropolitan berasal dari 2004-2009,
  mayoritas 2008, sedangkan indikator integrasi hanya dapat dikumpulkan
  untuk suatu tahun dalam rentang 2014-2016. Penulis menyebut
  ketidakselarasan waktu ini tidak ideal, meskipun mereka menilai
  sebagian besar indikator integrasi cukup konstan dari waktu ke waktu.
  (Locator: PDF p.~10; TXT baris 571-585.)
\item
  {[}Laporan sumber{]} Keterbatasan kelayakan mencegah penulis
  mengeksplorasi kondisi integrasi kultural dalam periode waktu yang
  lebih panjang. Ukuran kultural karena itu bersifat seperti snapshot,
  bukan sejarah integrasi yang lengkap. (Locator: PDF pp.~10-12; TXT
  baris 680-718.)
\item
  {[}Laporan sumber{]} Homogenitas politik digunakan sebagai pendekatan
  terhadap kedekatan kultural, dengan asumsi bahwa warna politik wali
  kota mencerminkan tanda tangan kultural-politik kota. Penulis sendiri
  menyatakan asumsi ini lebih mungkin berlaku di negara yang wali
  kotanya dipilih langsung atau melalui dewan lokal daripada di negara
  yang wali kotanya ditunjuk pemerintah pusat. Homogenitas bahasa
  direduksi menjadi variabel dummy: nilainya 1 berarti tidak ada
  hambatan bahasa, sedangkan nilainya 0 diberikan bila sedikitnya 10
  persen penduduk PUR berbicara bahasa yang berbeda. (Locator: PDF
  pp.~11-12, Table 3 dan pembahasan; TXT baris 642-656, 691-702.)
\item
  {[}Laporan sumber{]} Identifikasi PUR menggunakan perspektif
  morfologis untuk sengaja tidak memilih hanya wilayah yang sudah
  terintegrasi secara fungsional. Namun, pilihan ini berarti dimensi
  relasional tidak dipakai untuk mendefinisikan PUR sejak awal. Ambang
  Herfindahl 0,56 juga ditentukan secara pragmatis melalui ``common
  sense'', bukan melalui prosedur ambang yang dirinci lebih lanjut.
  (Locator: PDF pp.~3 dan 5; TXT baris 151-159, 264-304.)
\item
  {[}Laporan sumber{]} Ukuran kinerja adalah kehadiran fungsi
  metropolitan sebagai proksi ekonomi urbanisasi. Artikel tidak
  menyamakan ukuran tersebut dengan seluruh aspek produktivitas,
  kesejahteraan, atau kualitas hidup. (Locator: PDF pp.~8-9; TXT baris
  482-505.)
\item
  {[}Laporan sumber{]} Korelasi tinggi antara keberadaan badan
  metropolitan dan tipe kemitraan membuat model institusional simultan
  kurang relevan. Interaksi dengan ukuran, lokasi, dan status lintas
  batas juga tidak dilaporkan lengkap di dalam tabel, hanya hasil
  signifikansinya yang diringkas. (Locator: PDF pp.~13-17; TXT baris
  680-699, 811-828, 913-930.)
\item
  {[}Inferensi pengolahan{]} Desain potong lintang, tanpa strategi
  identifikasi kausal yang disebutkan dalam TXT, membatasi interpretasi
  pada hubungan statistik. Arah sebab-akibat, kemungkinan sebab bersama,
  dan perubahan integrasi sebelum atau sesudah peningkatan fungsi
  metropolitan tidak dapat ditentukan dari laporan ini saja. (Locator:
  PDF pp.~3, 12-18; TXT baris 160-187, 702-953.)
\item
  {[}Batas provenance{]} Bobot atau formula agregasi setelah normalisasi
  indeks, pemeriksaan asumsi ordered logit, penanganan data hilang,
  analisis sensitivitas terhadap ambang 0,56, dan lokasi dataset yang
  diklaim tersedia: Tidak disebutkan dalam file. (Locator: PDF pp.~5,
  10-12, dan 18; TXT baris 284-304, 580-702, 1020-1046.)
\end{itemize}

\subsection{Challenges}\label{challenges-44}

\begin{itemize}
\tightlist
\item
  {[}Laporan sumber{]} Tantangan konseptual pertama adalah
  ketidakjelasan definisi PUR dan perdebatan apakah relasi antar-kota
  harus menjadi bagian dari definisi polisentrisitas. Penulis mengatasi
  hal ini dengan memisahkan identifikasi morfologis dari pengukuran
  integrasi. (Locator: PDF pp.~5-6; TXT baris 252-318.)
\item
  {[}Laporan sumber{]} Tantangan delimitation mencakup penentuan batas
  kota, hinterland, FUA, polyFUA, PIA, dan wilayah lintas negara,
  sekaligus mencegah ukuran wilayah yang terlalu besar atau seleksi
  hanya pada wilayah yang sudah terintegrasi. (Locator: PDF pp.~5-6; TXT
  baris 264-318.)
\item
  {[}Laporan sumber{]} Analisis atas lebih dari 100 PUR memungkinkan
  gambaran luas yang konsisten, tetapi tidak dapat memberikan detail
  mendalam seperti studi kasus. Penulis menyebut pengumpulan indikator
  regional dengan cakupan Eropa sebagai pekerjaan yang sulit dan kadang
  melelahkan. (Locator: PDF p.~3 dan pp.~10-12; TXT baris 160-187,
  580-597.)
\item
  {[}Laporan sumber{]} Pengukuran institusi metropolitan harus
  membandingkan bentuk kerja sama yang sangat berbeda, dari kesepakatan
  informal hingga otoritas terpilih, di bawah konteks nasional yang
  berbeda. Pengukuran kultural juga menghadapi perbedaan sistem
  pemilihan wali kota, bahasa, dan PUR lintas batas. (Locator: PDF
  pp.~10-12; TXT baris 621-656, 680-718.)
\item
  {[}Laporan sumber{]} Kompleksitas lintas batas diakui sebagai kondisi
  yang menantang, walaupun model tidak menemukan hubungan
  integrasi-kinerja yang berbeda secara signifikan untuk PUR lintas
  batas. (Locator: PDF pp.~16-18; TXT baris 853-910, 918-930.)
\item
  {[}Laporan sumber{]} Multikolinearitas antara keberadaan badan
  metropolitan dan tipe kemitraan menyulitkan interpretasi model
  simultan. (Locator: PDF pp.~13-14; TXT baris 724-737, 811-822.)
\end{itemize}

\subsection{Practical Implications}\label{practical-implications-44}

{[}Interpretasi penulis{]} Implikasi praktis paling langsung adalah
bahwa kebijakan PUR sebaiknya tidak berhenti pada keberadaan beberapa
pusat kota atau pada penjumlahan ukuran penduduk. Pemerintah dan aktor
metropolitan perlu membangun koneksi fungsional, khususnya infrastruktur
penghubung dan angkutan umum antar-kota, agar hubungan antar-pusat
memungkinkan arus dan interaksi. (Locator: PDF p.~18; TXT baris
1032-1046.)

{[}Interpretasi penulis{]} Pada sisi kelembagaan, temuan mendukung
pembentukan atau pemeliharaan suatu badan kerja sama metropolitan yang
bekerja, tanpa mengharuskan satu bentuk kelembagaan yang seragam.
Koordinasi sumber daya, infrastruktur, tata guna lahan, dan
komplementaritas antar-pusat merupakan arah tindakan yang konsisten
dengan mekanisme yang dibahas. (Locator: PDF pp.~4-5 dan 18; TXT baris
227-258, 984-1004, 1035-1046.)

{[}Interpretasi penulis{]} Pada tingkat nasional, investasi dalam
metropolisasi PUR disebut sebagai alternatif yang layak terhadap
konsentrasi investasi lebih lanjut di kawasan ibu kota. Penulis juga
menyatakan bahwa tindakan yang mendorong integrasi akan ``pay off''.
(Locator: PDF p.~18; TXT baris 1032-1046.)

{[}Inferensi pengolahan{]} Pernyataan ``pay off'' sebaiknya dibaca
sebagai implikasi kebijakan dari hubungan statistik artikel, bukan
sebagai evaluasi kausal atas program tertentu. (Locator: PDF p.~18; TXT
baris 1032-1046.)

\subsection{Applications}\label{applications-44}

\begin{itemize}
\tightlist
\item
  {[}Laporan sumber{]} Penulis menyatakan bahwa daftar 117 PUR yang
  konsisten membuka peluang penelitian komparatif dan dapat membantu
  menjelaskan kontras kinerja antar-PUR. (Locator: PDF pp.~6-8 dan 18;
  TXT baris 322-479, 987-1005.)
\item
  {[}Inferensi pengolahan{]} Kerangka indikator dapat diterapkan sebagai
  alat diagnosis awal: konektivitas jalan dan kereta serta frekuensi
  kereta untuk dimensi fungsional; badan, durasi, dan tipe kerja sama
  untuk dimensi institusional; serta politik dan bahasa untuk dimensi
  kultural. Hasil diagnosis tidak boleh diperlakukan sebagai ukuran
  lengkap tanpa validasi konteks, karena keterbatasan proksi diakui
  artikel. (Locator: PDF pp.~10-12; TXT baris 580-702.)
\item
  {[}Interpretasi penulis{]} Dalam perencanaan, aplikasi yang disarankan
  adalah mengarahkan investasi pada hubungan antar-kota, transit
  antarkota, dan proses region-building agar PUR berfungsi sebagai
  entitas metropolitan yang lebih terintegrasi. (Locator: PDF p.~18; TXT
  baris 1032-1046.)
\item
  {[}Batas provenance{]} Implementasi kebijakan tertentu, evaluasi
  program, studi kasus pascakebijakan, atau bukti bahwa suatu pemerintah
  telah menerapkan indeks artikel: Tidak disebutkan dalam file.
  (Locator: PDF p.~18; TXT baris 970-1046.)
\end{itemize}

\subsection{Future Research
(Optional)}\label{future-research-optional-44}

{[}Laporan sumber{]} Agenda riset yang secara eksplisit disebut penulis
meliputi: (Locator: PDF p.~18, DISCUSSION AND CONCLUSION; TXT baris
1006-1046.)

\begin{itemize}
\tightlist
\item
  menggunakan indikator lain untuk mengukur ketiga dimensi integrasi dan
  menguji ketahanan analisis;
\item
  menilai relevansi pembagian kerja tertentu antar-kota terhadap
  kinerja;
\item
  menguji polisentrisitas fungsional wilayah, bukan hanya
  polisentrisitas morfologis;
\item
  membuat pembedaan yang lebih halus di antara entitas tata kelola
  metropolitan, dengan memperhitungkan kekhususan regulasi nasional dan
  warisan sejarah;
\item
  mengembangkan indikator integrasi kultural yang lebih komprehensif,
  termasuk perbedaan etnis, agama, dan persepsi pembentukan identitas
  pada skala PUR;
\item
  menggunakan proksi kinerja lain; dan
\item
  menambahkan dimensi waktu untuk memahami evolusi integrasi dan
  hubungannya dengan kinerja.
\end{itemize}

{[}Laporan sumber{]} Penulis juga menyatakan bahwa dataset disediakan
untuk membantu pengujian dengan indikator lain. Lokasi atau tautan
dataset tersebut tidak disebutkan dalam TXT, sehingga detail aksesnya
tetap Tidak disebutkan dalam file. (Locator: PDF p.~18; TXT baris
1006-1046.)

\subsection{Provenance Notes}\label{provenance-notes-44}

\begin{itemize}
\tightlist
\item
  Review ini hanya menggunakan file
  \texttt{source/journal/B01-meijers-hoogerbrugge-cardoso-2018-beyond-polycentricity-does-stronger-integration-between-cities-in-polycentric-urban-regions-improve-performance-10.1111-tesg.12292.txt}.
  Verifikasi filename dengan pola \texttt{source/**/B01*.txt}
  menghasilkan tepat satu TXT aktual, yaitu file tersebut. Tidak ada
  informasi dari referensi, web, PDF lain, atau file kajian lain yang
  ditambahkan sebagai bukti.
\item
  Seluruh 1.196 baris logis TXT dibaca, dari baris 1 sampai 1.196,
  termasuk blok metadata, isi artikel, Table 1-10, Figure 1-2, catatan,
  acknowledgements, dan referensi yang tercantum dalam ekstraksi.
  Perintah penghitung baris berbasis karakter baris baru menghasilkan
  1.195 karena baris terakhir tidak diakhiri karakter baris baru.
  Referensi di dalam TXT tidak diverifikasi secara independen.
\item
  TXT memuat metadata bahwa sumber PDF memiliki 21 halaman, diekstrak
  pada 2026-08-31 dengan \texttt{pdftotext\ -layout}, dan footer PDF
  menunjukkan unduhan dari Wiley pada 30/08/2026. Locator utama dalam
  review memakai nomor halaman cetak PDF, bagian, tabel, atau gambar,
  disertai rentang baris TXT bila relevan. (Locator: TXT baris 1-26.)
\item
  Ekstraksi mempertahankan header/footer, pemisah halaman, dan beberapa
  artefak karakter. Dalam sejumlah tabel, tanda minus tampak sebagai
  prefiks atau karakter ``2'', misalnya \texttt{20.249},
  \texttt{20.00036}, atau \texttt{21.151}. Arah negatif hanya diringkas
  ketika juga ditegaskan oleh prosa artikel; nilai tabel yang ambigu
  tidak dikoreksi secara diam-diam.
\item
  Judul Table 9 pada ekstraksi menyebut pengaruh integrasi
  institusional, tetapi isi baris tabel dan uraian hasil membahas
  indikator serta indeks integrasi kultural. Review mempertahankan
  ketidaksesuaian ini dan mengikuti isi hasil ketika merangkum Table 9;
  penyelesaiannya melalui PDF atau sumber lain: Tidak dilakukan.
\item
  Klaim artikel seperti ``first comprehensive list'' dan ``first-ever
  cross-sectional analysis'' diperlakukan sebagai klaim penulis, bukan
  hasil verifikasi eksternal. Status peer-review, URL dataset, formula
  bobot agregasi indeks, rincian data hilang, diagnostik model, dan
  keluaran lengkap model yang tidak dilaporkan adalah Tidak disebutkan
  dalam file.
\item
  Bagian yang diberi label \texttt{{[}Inferensi\ pengolahan{]}} tidak
  ditulis sebagai temuan penulis. Batas utama inferensi adalah desain
  potong lintang, ketidaksesuaian periode data, proksi fungsi
  metropolitan, dan keterbatasan indikator integrasi sebagaimana
  dijelaskan di dalam TXT.
\end{itemize}

\section{Review B02}\label{review-b02}

Kode kajian: B02

Identitas sumber: - Judul: \emph{Form Follows Function? Linking
Morphological and Functional Polycentricity} - Penulis: Martijn Burger
dan Evert Meijers. - Jurnal: \emph{Urban Studies}, volume 49, nomor 5,
April 2012, halaman 1127-1149. - DOI: 10.1177/0042098011407095. -
Riwayat versi yang tercantum dalam TXT: OnlineFirst 6 Juni 2011; Version
of Record 9 Maret 2012. - Jenis sumber: artikel jurnal. - Status
peer-review: Tidak disebutkan dalam file.

Penanda evidensi: {[}Laporan sumber{]} berarti isi yang dilaporkan atau
didefinisikan dalam artikel; {[}Interpretasi penulis{]} berarti
penjelasan atau penarikan makna yang dibuat oleh Burger dan Meijers;
{[}Inferensi pengolahan{]} berarti pembacaan analitis saya atas isi TXT
dan bukan klaim eksplisit penulis.

\subsection{Summarized Abstract}\label{summarized-abstract-45}

{[}Laporan sumber{]} Artikel berangkat dari masih tipisnya penelitian
empiris tentang biaya dan manfaat sistem urban polisentris. Penulis
mengaitkan persoalan ini dengan bertahannya dua pendekatan yang secara
analitis berbeda: pendekatan morfologis yang berpusat pada ciri nodal,
serta pendekatan fungsional yang berfokus pada hubungan antarpusat.
Artikel menyusun kerangka teoretis umum yang menghubungkan keduanya,
menjelaskan cara pengukuran dan perbandingan yang koheren, lalu
mengujinya pada 42 region WGR di Belanda. Abstrak menyatakan bahwa
sebagian besar region lebih polisentris secara morfologis daripada
secara fungsional. Perbedaan tersebut terutama dikaitkan dengan ukuran,
konektivitas eksternal, dan tingkat swasembada pusat utama suatu region
(locator: TXT baris 83-95; Abstract).

{[}Laporan sumber{]} Pernyataan abstrak tersebut bertentangan dengan
uraian hasil dan kesimpulan utama artikel, yang menyatakan bahwa hampir
semua region relatif lebih polisentris secara fungsional daripada secara
morfologis. Kontradiksi ini tidak diselesaikan oleh TXT, sehingga kedua
versi dipertahankan dalam review ini (locator: TXT baris 642-648 dan
856-866; hlm. 1138 dan 1144).

{[}Interpretasi penulis{]} Morfologi dan fungsi tidak perlu diperlakukan
sebagai satu dimensi yang sama. Keduanya berbagi prinsip mengenai
keseimbangan kepentingan pusat, tetapi memakai konsep kepentingan yang
berbeda: nodalitas untuk kepentingan absolut dan sentralitas untuk
kepentingan relatif (locator: TXT baris 174-186 dan 406-426; hlm.
1128-1129 dan 1133).

\subsection{Problem Statement}\label{problem-statement-45}

{[}Laporan sumber{]} Literatur kebijakan dan akademik tentang
polisentrisme berkembang luas, tetapi polisentrisitas tetap menjadi
konsep yang ``fuzzy'' dan digunakan dengan beragam arti. Ketidakjelasan
ini menghambat kemajuan empiris untuk menentukan manfaat aktual
pembangunan polisentris serta konsekuensi lingkungannya, ekonominya, dan
sosialnya. Artikel juga melaporkan bahwa perdebatan utama berkisar pada
apakah polisentrisitas hanya menyangkut bentuk dan ukuran pusat kota
atau juga mencakup relasi di antara pusat-pusat tersebut (locator: TXT
baris 97-165; hlm. 1127-1129).

{[}Laporan sumber{]} Dua pendekatan tersebut belum banyak dihubungkan.
Distribusi ukuran pusat yang seimbang tidak dengan sendirinya
membuktikan adanya hubungan fungsional, hubungan yang setara, atau pola
arus multidireksional. Sebaliknya, ketidakseimbangan ukuran tidak
otomatis menunjukkan bahwa relasi fungsionalnya hierarkis. Artikel
menyatakan bahwa masih belum jelas mengapa sebagian sistem morfologisnya
polisentris tetapi tidak fungsional, atau sebaliknya (locator: TXT baris
148-165 dan 174-186; hlm. 1128-1129).

{[}Inferensi pengolahan{]} Masalah penelitian dapat dirumuskan sebagai
masalah pengukuran dan penjelasan, bukan sebagai pengujian dampak
kebijakan. Artikel hendak menjelaskan hubungan serta perbedaan dua
ukuran polisentrisitas; ia tidak langsung mengestimasi apakah
polisentrisitas menyebabkan kinerja regional yang lebih baik (locator:
TXT baris 174-208 dan 845-925; hlm. 1129 dan 1144-1145).

\subsection{Objectives}\label{objectives-45}

{[}Laporan sumber{]} Tujuan yang dinyatakan adalah mengeksplorasi
hubungan antara polisentrisitas morfologis dan fungsional; menyajikan
kerangka teoretis umum yang berakar pada penelitian sistem urban;
menunjukkan saling ketergantungan antara keseimbangan distribusi
kepentingan absolut dan relatif pusat; serta menjelaskan mengapa tingkat
kedua bentuk polisentrisitas berbeda di dalam wilayah. Penulis juga
bertujuan menguji kerangka tersebut pada region-region di Belanda dan
membandingkan kedua ukuran dengan cara yang koheren (locator: TXT baris
174-208 dan 340-365; hlm. 1129 dan 1132).

{[}Laporan sumber{]} Tujuan operasionalnya adalah mengukur
polisentrisitas morfologis melalui distribusi rank-size skor nodalitas,
mengukur polisentrisitas fungsional melalui distribusi rank-size skor
sentralitas internal, dan memisahkan keseimbangan arah relasi dari
kepadatan jaringan (locator: TXT baris 406-465 dan 497-526; hlm.
1133-1135).

\subsection{Literature Survey}\label{literature-survey-45}

{[}Inferensi pengolahan{]} Dari cara telaah disajikan, bagian pustaka
terbaca sebagai sintesis konseptual dan historis, bukan tinjauan
sistematis yang metodenya dijelaskan secara eksplisit. File tidak
menyebutkan strategi pencarian, kriteria inklusi-eksklusi, atau jumlah
korpus literatur yang ditelaah. {[}Laporan sumber{]} Penulis
mengelompokkan literatur polisentrisitas ke dalam tiga pemakaian yang
berkaitan tetapi berbeda: strategi perencanaan normatif pada skala
metropolitan, nasional, dan transnasional; proses spasial berupa difusi
fungsi urban dari pusat besar ke pusat yang lebih kecil; serta
konfigurasi spasial yang dihasilkan (locator: TXT baris 132-165; hlm.
1128-1129).

{[}Laporan sumber{]} Landasan teoretis utama berasal dari teori tempat
sentral dan tradisi sistem urban. Teori tempat sentral Christaller dan
Losch membahas distribusi, ukuran, dan jumlah kota serta hirarki pusat.
Dalam gambaran hierarkis klasik, pusat tingkat rendah bergantung pada
pusat tingkat tinggi dan pola interaksi menyerupai jaringan bintang yang
terpusat. Tradisi sistem urban, yang dipengaruhi teori sistem umum,
melihat kota sebagai simpul yang saling bergantung dan interaksi di
antara simpul sebagai bagian dari sistem. Penulis mencatat bahwa hirarki
semakin sulit menjelaskan organisasi metropolitan yang lebih polisentris
dalam era pascaindustri dan globalisasi (locator: TXT baris 193-253;
hlm. 1129-1130).

{[}Laporan sumber{]} Konsep Preston tentang nodalitas dan sentralitas
menjadi jembatan operasional. Nodalitas adalah kepentingan absolut
pusat, sedangkan sentralitas adalah bagian kepentingan yang berasal dari
penyediaan barang, jasa, atau pekerjaan melebihi kebutuhan penduduk
pusat itu sendiri. Dalam sistem tertutup, sentralitas ditulis sebagai
\texttt{Cc\ =\ Nc\ -\ Lc}, dengan \texttt{Nc} sebagai kepentingan
absolut dan \texttt{Lc} sebagai kepentingan lokal dari arus internal.
Untuk sistem terbuka, sentralitas internal ditulis sebagai
\texttt{Cci\ =\ Nc\ -\ Cce\ -\ Lc}, dengan \texttt{Cce} sebagai
sentralitas eksternal (locator: TXT baris 257-308 dan 327-365; hlm.
1130-1132).

{[}Laporan sumber{]} Literatur morfologis menyamakan polisentrisitas
dengan distribusi pusat yang lebih seimbang dalam ukuran atau
kepentingan absolut. Literatur fungsional memperluas perhatian itu ke
pola interaksi, terutama keseimbangan arus masuk dan arah yang
multidireksional. Green menambahkan dimensi kepadatan jaringan, yang
dalam artikel dipisahkan dari keseimbangan arah arus karena jaringan
yang padat masih dapat bersifat hierarkis, sedangkan jaringan yang
jarang dapat memiliki pusat-pusat dengan konektivitas relatif seimbang
(locator: TXT baris 380-426 dan 437-465; hlm. 1132-1134).

{[}Interpretasi penulis{]} Artikel menempatkan perbedaan dengan studi
Polynet sebagai persoalan definisi dan operasionalisasi: studi tersebut
menggabungkan keseimbangan arus dan kepadatan jaringan dalam indeks
fungsional, sedangkan artikel ini memisahkan keduanya agar perbandingan
dengan ukuran morfologis lebih jelas (locator: TXT baris 642-661 dan
432-465; hlm. 1138 dan 1133-1134).

\subsection{Method Used}\label{method-used-45}

{[}Laporan sumber{]} Penelitian menggunakan kerangka teoritis yang
menghubungkan nodalitas sebagai kepentingan absolut dan sentralitas
sebagai kepentingan relatif. Polisentrisitas morfologis diukur dari
keseimbangan distribusi nodalitas, sedangkan polisentrisitas fungsional
diukur dari keseimbangan distribusi sentralitas internal. Kepadatan
jaringan dihitung terpisah sebagai rasio sentralitas internal terhadap
nodalitas; karena itu, kepadatan jaringan tidak dimasukkan ke dalam
ukuran polisentrisitas fungsional (locator: TXT baris 347-365, 406-465,
dan 612-631; hlm. 1132-1134 dan 1138).

{[}Laporan sumber{]} Secara empiris, penulis mengestimasi kemiringan
garis regresi yang paling sesuai dengan distribusi rank-size. Skor
nodalitas menghasilkan ukuran polisentrisitas morfologis, sedangkan skor
sentralitas internal menghasilkan ukuran polisentrisitas fungsional.
Garis yang lebih datar diperlakukan sebagai lebih polisentris dan garis
yang lebih curam sebagai lebih monosentris. Untuk menghindari
ketergantungan pada satu pilihan ukuran sampel pusat, penulis memakai 2,
3, dan 4 tempat per region lalu merata-ratakan ketiga skor tersebut
(locator: TXT baris 497-526 dan 559-586; hlm. 1135-1137).

{[}Laporan sumber{]} Analisis dijalankan secara terpisah untuk arus
perjalanan kerja dan perjalanan belanja. Penulis membandingkan skor
morfologis dan fungsional, menyusun matriks korelasi untuk
polisentrisitas, kepadatan jaringan, serta distribusi sentralitas
eksternal, kemudian memeriksa perbedaan antarregion melalui Gambar 5-9.
Parameter rank-size diperkirakan dengan pendekatan Gabaix dan Ibragimov
yang, menurut catatan artikel, mengoreksi bias sampel kecil (locator:
TXT baris 581-621, 726-769, dan 978-985; hlm. 1137-1141 serta catatan
8-9).

{[}Inferensi pengolahan{]} Desain ini adalah perbandingan deskriptif dan
korelasional atas struktur spasial. TXT tidak menyebutkan model kausal,
estimasi efek kebijakan, uji signifikansi, atau spesifikasi regresi
multivariat untuk menjelaskan perbedaan antara dua skor. Karena itu,
hubungan yang disajikan tidak cukup untuk dibaca sebagai bukti
sebab-akibat (locator: TXT baris 581-621 dan 726-769; hlm. 1137-1141).

\subsection{Dataset}\label{dataset-45}

{[}Laporan sumber{]} Data arus berasal dari \emph{Dutch National Travel
Survey 2004-08} (\emph{Mobiliteitsonderzoek Nederland}). Penulis
menghitung skor tahunan rata-rata dan memberi bobot pada skor agar
representatif bagi seluruh populasi Belanda. Dua jenis arus yang dipakai
adalah perjalanan rumah-kerja dan perjalanan belanja (locator: TXT baris
587-621 dan 978-985; hlm. 1137 serta catatan 8-9).

{[}Laporan sumber{]} Untuk dimensi pekerjaan, nodalitas tempat
ditentukan dari pekerjaan, yaitu total arus masuk perjalanan kerja yang
mencakup arus dari pusat tempat itu sendiri dan dari luar region WGR.
Sentralitas internal ditentukan dari total arus masuk perjalanan kerja
yang berasal dari tempat-tempat dalam WGR yang sama. Untuk dimensi
belanja, nodalitas memakai jumlah total pembelanja dan sentralitas
internal memakai total arus belanja masuk dari tempat-tempat dalam WGR
yang sama (locator: TXT baris 587-621; hlm. 1137).

{[}Laporan sumber{]} Dataset spasialnya adalah 42 region WGR yang
bersama-sama mencakup seluruh Belanda. Batas WGR berasal dari
pengelompokan kelembagaan berdasarkan \emph{Intermunicipal Statutory
Regulations Act} dan persepsi administrator serta anggota dewan
municipal/provinsial tentang skala kerja sama regional. CBS dicantumkan
sebagai sumber peta WGR pada Gambar 3 (locator: TXT baris 470-557; hlm.
1134-1136).

{[}Inferensi pengolahan{]} TXT tidak menyediakan ukuran mikrodata
survei, daftar pertanyaan, jumlah observasi per arus, tingkat
nonrespons, penanganan data hilang, atau rincian prosedur pembobotan.
Informasi tersebut tidak boleh diisi dengan dugaan (locator: TXT baris
587-621 dan 978-985).

\subsection{Population / Sample}\label{population-sample-45}

{[}Laporan sumber{]} Populasi spasial penelitian adalah 42 region WGR
yang diperlakukan sebagai region yang secara fungsional koheren dan
mencakup seluruh wilayah Belanda. Penulis menyebut batas tersebut
sebagai proksi tidak langsung bagi region yang koheren secara
fungsional; batasnya dinilai cukup dapat dipertahankan karena cukup
bertepatan dengan area perjalanan kerja (locator: TXT baris 470-495;
hlm. 1134-1135).

{[}Laporan sumber{]} Di dalam setiap WGR, data nodalitas dan sentralitas
dikumpulkan untuk empat kota atau tempat terbesar. Prosedur pengukuran
kemudian menggunakan 2, 3, dan 4 tempat terbesar secara bergantian dan
merata-ratakan skornya. Gambar 4 mengilustrasikan empat tempat terbesar
berdasarkan pekerjaan untuk region Maastricht dan Sittard-Geleen
(locator: TXT baris 487-495 dan 579-596; hlm. 1135 dan 1137).

{[}Inferensi pengolahan{]} Unit analisis yang tampak dalam analisis
adalah distribusi pusat dalam region, bukan individu responden.
{[}Laporan sumber{]} Data survei diberi bobot agar mewakili populasi
Belanda, tetapi ukuran sampel responden, kerangka sampling, dan jumlah
observasi yang masuk ke setiap region tidak disebutkan dalam file
(locator: TXT baris 978-985; catatan 8-9).

\subsection{Dependent Variables}\label{dependent-variables-45}

{[}Inferensi pengolahan{]} Variabel dependen kausal: Tidak berlaku,
karena artikel tidak menggunakan istilah variabel dependen dalam
formulasi kausal. Untuk memenuhi kategori ini tanpa memaksakan desain
kausal, ukuran hasil utama yang dibandingkan adalah sebagai berikut: -
Skor polisentrisitas morfologis, yaitu kemiringan regresi rank-size atas
skor nodalitas pusat-pusat terbesar; garis yang lebih datar berarti
distribusi kepentingan absolut lebih seimbang. - Skor polisentrisitas
fungsional, yaitu kemiringan regresi rank-size atas skor sentralitas
internal pusat-pusat terbesar; garis yang lebih datar berarti distribusi
kepentingan relatif atau arus internal lebih seimbang. - Selisih MP-FP,
yaitu perbedaan kedua skor dalam satuan perbedaan poin persentase yang
dipakai untuk membandingkan apakah suatu region relatif lebih morfologis
atau lebih fungsional polisentris.

Definisi dan arah interpretasi ukuran tersebut dijelaskan pada TXT baris
497-526, 579-586, dan 717-721 (hlm. 1135-1137 serta 1139-1140).

{[}Laporan sumber{]} Kepadatan jaringan, distribusi sentralitas
eksternal, ukuran pusat utama, orientasi lokal, dan tingkat swasembada
dipakai sebagai dimensi pembanding atau penjelas perbedaan, bukan
sebagai bagian dari skor polisentrisitas fungsional. Kepadatan jaringan
secara khusus dihitung sebagai rasio sentralitas internal terhadap
nodalitas (locator: TXT baris 437-465 dan 612-631; hlm. 1133-1134 dan
1138).

{[}Inferensi pengolahan{]} Tidak ada variabel dependen berupa
produktivitas, kinerja regional, manfaat sosial, dampak lingkungan, atau
keberhasilan kebijakan. Maka artikel mengukur konfigurasi spasial dan
relasi arus, bukan konsekuensi normatif dari pembangunan polisentris
(locator: TXT baris 845-925; hlm. 1144-1145).

\subsection{Results}\label{results-45}

{[}Laporan sumber{]} Struktur spasial berbeda antarregion. Sebagian WGR
cenderung monosentris, sebagian cenderung polisentris, dan sebagian
besar berada di antara kedua kutub tersebut. Pola umumnya serupa untuk
arus perjalanan kerja dan belanja (locator: TXT baris 630-641 dan
845-866; hlm. 1138 dan 1144).

{[}Laporan sumber{]} Korelasi antara skor polisentrisitas morfologis dan
fungsional adalah 0,84 untuk pekerjaan/perjalanan kerja dan 0,57 untuk
belanja. Walaupun korelasinya cukup besar, hampir semua region relatif
lebih polisentris secara fungsional daripada morfologis. Dalam kedua
jenis arus, distribusi arus masuk di dalam WGR lebih merata daripada
distribusi ukuran pusat (locator: TXT baris 642-651; Tabel 1 dan Tabel 2
pada hlm. 1138-1139; TXT baris 666-740).

{[}Laporan sumber{]} Terdapat ketidaksesuaian internal yang tidak boleh
diabaikan: abstrak menyatakan mayoritas region lebih morfologis daripada
fungsional, sedangkan bagian hasil, Tabel 1, dan kesimpulan menyatakan
hampir semua lebih fungsional daripada morfologis. Tabel 1
memperlihatkan mayoritas nilai MP-FP negatif, tetapi TXT tidak memberi
erratum atau penjelasan untuk perbedaan dengan abstrak (locator: TXT
baris 83-95, 642-648, 666-721, dan 856-866; Abstract, Tabel 1, hlm.
1127, 1138-1140, dan 1144).

{[}Laporan sumber{]} Tabel 1 memperlihatkan arah perbedaan tersebut
melalui nilai MP-FP. Contoh yang relatif lebih fungsional adalah
Groningen, dengan MP-FP -0,82 untuk pekerjaan dan -0,93 untuk belanja,
serta Maastricht, dengan -0,67 dan -0,86. Contoh region yang relatif
lebih morfologis adalah Veendam, dengan 0,16 untuk pekerjaan dan 0,31
untuk belanja. Utrecht memiliki perbedaan yang kecil, yaitu -0,13 untuk
pekerjaan dan -0,23 untuk belanja. Nilai negatif berarti skor FP lebih
datar daripada MP menurut konvensi artikel, sedangkan nilai positif
berarti sebaliknya (locator: TXT baris 666-714 dan 717-721; Tabel 1 hlm.
1139 serta penjelasan hlm. 1139-1140).

{[}Laporan sumber{]} Hubungan polisentrisitas morfologis dengan
kepadatan jaringan relatif lemah tetapi positif: 0,30 untuk pekerjaan
dan 0,41 untuk belanja. Kepadatan jaringan hampir tidak berkorelasi
dengan keseimbangan fungsional: 0,10 untuk pekerjaan dan -0,01 untuk
belanja. Hasil ini dipakai penulis untuk mendukung pemisahan antara
kepadatan jaringan dan arah atau keseimbangan arus. Korelasi distribusi
sentralitas eksternal dengan MP adalah 0,90 untuk pekerjaan dan 0,55
untuk belanja; dengan FP masing-masing 0,78 dan 0,47; dengan kepadatan
jaringan masing-masing 0,18 dan 0,14 (locator: TXT baris 612-641 dan
726-740; Tabel 2 hlm. 1138-1139).

{[}Laporan sumber{]} Region yang lebih fungsional daripada morfologis
cenderung memiliki pusat utama yang lebih berorientasi pada pasar lokal
dan menarik arus lebih kuat dari luar WGR. Perbedaan besar juga
berasosiasi dengan pusat utama yang besar secara absolut, seperti
Amsterdam, Den Haag, Groningen, dan Maastricht. Sebaliknya, Roermond dan
Veendam disebut sebagai contoh region dengan pusat utama yang relatif
kecil dan subordinat dalam sistem urban supraregional, serta relatif
lebih morfologis daripada fungsional (locator: TXT baris 748-792; Gambar
6-8 pada hlm. 1140-1142).

{[}Laporan sumber{]} Perbandingan rinci atas Groningen, Utrecht, dan
Veendam memperlihatkan variasi regional. Di Groningen, distribusi
nodalitas lebih tidak seimbang daripada distribusi sentralitas internal,
sementara distribusi sentralitas eksternal sedikit lebih miring daripada
distribusi nodalitas. Sebanyak 54 persen pekerjaan di kota Groningen
diisi oleh penduduk yang juga tinggal di Groningen, dibandingkan dengan
24 persen pada kota-kota berperingkat lebih rendah di region tersebut.
Utrecht memiliki distribusi lokal, sentralitas internal, dan sentralitas
eksternal yang lebih serupa. Veendam sangat polisentris secara
morfologis, tetapi arus komuter internalnya tidak otomatis sama
meratanya (locator: TXT baris 775-842; Gambar 9 hlm. 1142-1143).

{[}Laporan sumber{]} Subbagian 5.3 juga menyebut Groningen sebagai
region yang jauh lebih morfologis daripada fungsional dan Veendam
sebagai salah satu yang lebih fungsional daripada morfologis. Pernyataan
ini berlawanan dengan nilai MP-FP pada Tabel 1, yang menunjukkan
Groningen negatif dan Veendam positif, serta dengan uraian sesudahnya
yang menyebut Veendam sangat polisentris secara morfologis. {[}Inferensi
pengolahan{]} TXT tidak memberi penjelasan untuk perbedaan ini; konflik
dipertahankan tanpa memilih salah satu versi (locator: TXT baris 775-792
dan 837-842; Tabel 1, TXT baris 666-714; hlm. 1142-1143).

\subsection{Findings}\label{findings-45}

{[}Interpretasi penulis{]} Temuan inti artikel adalah bahwa keseimbangan
ukuran pusat dan keseimbangan relasi pusat saling berkaitan tetapi tidak
identik. Sebuah region dapat mempunyai beberapa pusat dengan ukuran yang
relatif seimbang, namun arus yang menghubungkan pusat tersebut tetap
memiliki orientasi tertentu; sebaliknya, distribusi ukuran yang lebih
hierarkis dapat disertai distribusi sentralitas internal yang lebih
seimbang (locator: TXT baris 406-465, 630-661, dan 824-842; hlm.
1133-1134, 1138, dan 1143).

{[}Interpretasi penulis{]} Perbedaan utama antara MP dan FP dijelaskan
melalui komposisi nodalitas. Nodalitas mencakup kepentingan lokal dan
eksternal, sedangkan sentralitas internal hanya menangkap surplus arus
yang datang dari tempat lain di WGR. Pusat utama yang besar dapat lebih
swasembada karena memiliki pasar tenaga kerja dan fungsi tingkat tinggi
yang lebih beragam, sekaligus menarik komuter dan pembelanja dari luar
region. Akibatnya, ukuran absolutnya dapat lebih dominan daripada
distribusi arus internalnya (locator: TXT baris 649-661 dan 894-925;
hlm. 1138-1140 dan 1144-1145).

{[}Interpretasi penulis{]} Kepadatan jaringan tidak boleh disamakan
dengan polisentrisitas fungsional. Relasi yang banyak tidak selalu
terdistribusi secara seimbang, dan relasi yang sedikit dapat tetap
relatif seimbang di antara pusat-pusat. Menurut penulis, pemisahan ini
penting secara konseptual dan membuat perbandingan dengan
polisentrisitas morfologis lebih konsisten (locator: TXT baris 437-465;
hlm. 1133-1134).

{[}Inferensi pengolahan{]} Bukti artikel mendukung jawaban bersyarat
terhadap pertanyaan pada judul: fungsi tidak sekadar mengikuti bentuk
secara satu-ke-satu. Korelasi 0,84 dan 0,57 menunjukkan hubungan
empiris, tetapi perbedaan skor, variasi antarregion, dan hasil Veendam
menunjukkan bahwa bentuk morfologis tidak cukup sebagai proksi tunggal
bagi struktur fungsional (locator: TXT baris 630-661, 726-740, dan
837-842). Ini adalah inferensi dari pola yang dilaporkan, bukan
formulasi kausal penulis.

\subsection{Conclusions}\label{conclusions-45}

{[}Laporan sumber{]} Kesimpulan eksplisit artikel dapat diringkas
sebagai berikut (locator: TXT baris 845-884; hlm. 1144): - Tidak ada
satu tipe organisasi spasial yang dominan di seluruh region; region
berbeda-beda dari monosentris sampai polisentris, dengan banyak region
berada di antaranya. - Polisentrisitas morfologis dan fungsional
berkorelasi cukup besar, tetapi hampir semua region relatif lebih
polisentris secara fungsional daripada morfologis. - Semakin besar
dominasi FP atas MP, semakin besar kecenderungan pusat utama untuk
bersandar pada pasar tenaga kerja dan konsumen lokalnya sendiri, menarik
arus dari luar region, dan berukuran lebih besar.

{[}Laporan sumber{]} Kesimpulan pada bagian utama bertentangan dengan
kalimat ringkasan abstrak: bagian utama menyebut hampir semua region
lebih FP daripada MP, sedangkan abstrak menyebut sebagian besar lebih MP
daripada FP. TXT tidak menjelaskan apakah perbedaan ini merupakan
koreksi, perbedaan definisi, atau kesalahan editorial. Karena itu, arah
hasil harus dibaca dengan kualifikasi ini (locator: TXT baris 83-95 dan
848-884; Abstract dan hlm. 1144).

{[}Interpretasi penulis{]} Hubungan tersebut dapat dijelaskan oleh
asosiasi ukuran pusat dengan keragaman sektor, keragaman pekerjaan,
pasar tenaga kerja lokal yang lebih besar, serta fungsi tingkat tinggi
termasuk ritel khusus. Penulis memperkirakan mekanisme serupa mungkin
muncul di negara atau skala lain, tetapi menegaskan bahwa kepadatan
penduduk Belanda dan dominasi kota kecil-menengah dapat membuat
generalisasi itu tidak aman tanpa pengujian tambahan (locator: TXT baris
868-925; hlm. 1144-1145).

{[}Laporan sumber{]} Artikel tidak menyimpulkan bahwa pembangunan
polisentris pasti menghasilkan manfaat kebijakan. Sebaliknya, artikel
menyerukan basis bukti untuk menguji klaim manfaat polisentrisitas dan
keberlanjutan kebijakan tersebut, serta menyarankan agar manfaat dikaji
bersama kepadatan jaringan dan kapasitas region menarik arus dari jarak
yang lebih jauh (locator: TXT baris 894-925; hlm. 1144-1145).

\subsection{Contributions}\label{contributions-45}

{[}Laporan sumber{]} Kontribusi teoretis yang ditawarkan adalah model
yang menghubungkan polisentrisitas morfologis dan fungsional melalui
pembedaan nodalitas, sentralitas, kepentingan lokal, sentralitas
internal, dan sentralitas eksternal. Kerangka tersebut menghidupkan
kembali hubungan dengan teori tempat sentral dan penelitian sistem
urban, yang menurut penulis kurang menonjol dalam dua dekade sebelumnya
(locator: TXT baris 174-208 dan 257-365; hlm. 1129-1132).

{[}Laporan sumber{]} Kontribusi metodologisnya adalah memakai prinsip
rank-size yang sama untuk mengukur dua dimensi, memakai beberapa jumlah
pusat terbesar, mengoreksi bias sampel kecil pada estimasi parameter,
serta memisahkan keseimbangan arah arus dari kepadatan jaringan.
Kontribusi empirisnya adalah perbandingan dua ukuran pada 42 WGR dengan
dua domain arus, yaitu perjalanan kerja dan belanja (locator: TXT baris
497-526, 579-621, dan 824-845; hlm. 1135-1137 dan 1143-1144).

{[}Inferensi pengolahan{]} Nilai tambah analitis artikel terletak pada
pencegahan dua kekeliruan pembacaan: menganggap pusat yang ukurannya
seimbang pasti membentuk jaringan fungsional, atau menganggap jaringan
yang padat pasti polisentris secara fungsional. Namun, kontribusi ini
bersifat kerangka pengukuran dan diagnosis struktur; TXT tidak
menunjukkan bahwa model tersebut sudah membuktikan efek terhadap kinerja
regional (locator: TXT baris 437-465 dan 845-925).

\subsection{Research Gap}\label{research-gap-45}

{[}Laporan sumber{]} Gap yang hendak diisi artikel adalah kurangnya
bukti empiris tentang biaya dan manfaat sistem urban polisentris,
ketidakjelasan konseptual polisentrisitas, dan minimnya usaha untuk
menghubungkan pendekatan morfologis dengan fungsional. Penulis juga
mengidentifikasi ketidakjelasan mengenai alasan suatu sistem berbeda
dalam derajat MP dan FP serta kesulitan memperoleh data arus antarpusat
(locator: TXT baris 83-95, 120-165, dan 275-315; hlm. 1127-1131).

{[}Laporan sumber{]} Setelah perbandingan Belanda dilakukan, artikel
sendiri menyisakan gap generalisasi dan gap konsekuensi. Belum diketahui
apakah pola yang sama berlaku di negara lain atau skala lain; belum
diuji secara memadai bagaimana arus selain komuter dan belanja mengubah
perbedaan tersebut; dan belum dijawab apakah bentuk-bentuk
polisentrisitas itu menghasilkan kinerja regional, manfaat, atau
keberlanjutan kebijakan yang dijanjikan (locator: TXT baris 920-925 dan
934-942; hlm. 1145).

{[}Inferensi pengolahan{]} Karena tinjauan pustaka artikel tidak
dilaporkan sebagai tinjauan sistematis dan tidak menyediakan kriteria
korpus, review ini tidak dapat menyatakan bahwa gap tersebut mencakup
seluruh literatur polisentrisitas. Gap di atas harus dibaca sebagai gap
yang dirumuskan dan dibatasi oleh artikel ini (locator: TXT baris 97-186
dan 947-1134).

\subsection{Limitations}\label{limitations-45}

{[}Laporan sumber{]} Keterbatasan generalisasi yang dinyatakan penulis
berhubungan dengan konteks Belanda: negara ini relatif padat penduduk,
sebagian besar kotanya kecil atau menengah, dan tingkat polisentrisitas
umumnya mungkin relatif tinggi. Penulis juga menduga sentralitas
eksternal pusat utama mungkin relatif rendah. Karena itu, penelitian di
negara lain diperlukan untuk menilai apakah hasilnya dapat
digeneralisasi (locator: TXT baris 920-932; hlm. 1145).

{[}Laporan sumber{]} Cakupan arus dibatasi pada perjalanan kerja dan
belanja. Penulis menyatakan bahwa pengaruh distribusi sentralitas
eksternal terhadap selisih MP-FP relatif terbatas dalam analisis ini;
untuk perdagangan antarperusahaan atau hubungan pemegang saham, cakupan
geografis yang lebih luas mungkin membuat sentralitas eksternal lebih
penting. Pembahasan berikutnya juga berfokus terutama pada pekerjaan dan
komuter, sementara hasil belanja yang lebih luas disebut tersedia atas
permintaan (locator: TXT baris 934-942 dan 986-988; hlm. 1145 serta
catatan 10).

{[}Laporan sumber{]} Ukuran fungsional sengaja hanya memakai sentralitas
internal dan tidak menggabungkan kepadatan jaringan. Ini menjaga
pemisahan konsep, tetapi berarti skor FP tidak boleh dibaca sebagai
ukuran tunggal keterhubungan atau integrasi jaringan. Selain itu,
pengukuran memakai empat pusat terbesar dengan sensitivitas 2, 3, dan 4
pusat; pusat yang lebih kecil tidak menjadi bagian utama peringkat yang
dianalisis (locator: TXT baris 559-586 dan 437-465; hlm. 1136-1137 dan
1133-1134).

{[}Keterbatasan yang terlihat dari prosedur dan pelaporan{]} WGR
merupakan proksi tidak langsung bagi region yang koheren secara
fungsional, dengan batas yang dibentuk dari persepsi aktor pemerintahan
dan bukan dari satu kriteria arus yang dijelaskan secara numerik. TXT
juga tidak menyebutkan jumlah responden, tingkat nonrespons, data
hilang, interval kepercayaan, nilai-p, atau spesifikasi lengkap garis
pada Gambar 6-8. Informasi tersebut tidak tersedia dalam file dan tidak
boleh diasumsikan (locator: TXT baris 483-495, 511-520, 587-621, dan
717-792).

{[}Inferensi pengolahan{]} Artikel belum menguji variabel hasil yang
diperlukan untuk menjawab apakah pembangunan polisentris menghasilkan
manfaat lingkungan, ekonomi, atau sosial. Ini merupakan batas cakupan
substantif: artikel memperjelas cara mengukur struktur, tetapi belum
menghubungkan ukuran tersebut dengan outcome kebijakan atau kinerja
regional (locator: TXT baris 120-131 dan 894-925).

\subsection{Challenges}\label{challenges-45}

{[}Laporan sumber{]} Tantangan konseptual pertama adalah membedakan tiga
penggunaan polisentrisitas sebagai strategi normatif, proses spasial,
dan konfigurasi spasial, lalu membedakan lagi bentuk morfologis dari
relasi fungsional. Variasi istilah dan pemakaian berisiko menghasilkan
pengukuran yang tidak sebanding (locator: TXT baris 97-165 dan 132-163;
hlm. 1127-1129).

{[}Laporan sumber{]} Tantangan pengukuran meliputi pemisahan kepentingan
lokal, internal, dan eksternal dalam sistem yang terbuka; memperoleh
data arus; memilih jumlah pusat yang masuk ke distribusi rank-size;
serta menangani bias sampel kecil. Artikel menjelaskan bahwa ambang
ukuran atau proporsi total dapat membuat ukuran sampel itu sendiri
menjadi indikator polisentrisitas, sehingga penulis memilih jumlah pusat
yang tetap dan beberapa spesifikasi (locator: TXT baris 327-365 dan
527-586; hlm. 1131-1137).

{[}Interpretasi penulis{]} Tantangan lain adalah menjaga agar kepadatan
jaringan tidak tercampur dengan keseimbangan arus. Sistem dengan banyak
hubungan tetapi terpusat dan sistem dengan sedikit hubungan tetapi
seimbang dapat menerima penilaian berbeda pada kedua dimensi tersebut.
Pemisahan analitis diperlukan agar konsep dan perbandingan tidak
terdistorsi (locator: TXT baris 437-465; hlm. 1133-1134).

{[}Inferensi pengolahan{]} Tantangan replikasi terutama terletak pada
transparansi data dan keputusan delimitasi. Karena TXT tidak memuat
rincian mikrodata dan ketidakpastian statistik, peneliti lain akan sulit
menilai sensitivitas hasil hanya dari file ini; ini adalah konsekuensi
keterbatasan informasi yang tersedia, bukan tuduhan tentang data yang
tidak dilakukan (locator: TXT baris 587-621 dan 717-792).

\subsection{Practical Implications}\label{practical-implications-45}

{[}Interpretasi penulis{]} Untuk penelitian dan evaluasi kebijakan,
tingkat polisentrisitas sebaiknya tidak dilaporkan sebagai satu atribut
tunggal. Ukuran distribusi pusat secara morfologis perlu dibedakan dari
keseimbangan relasi fungsional, dan keduanya perlu dibaca bersama
kepadatan jaringan serta keterbukaan terhadap arus eksternal (locator:
TXT baris 432-465 dan 894-925; hlm. 1133-1134 dan 1144-1145).

{[}Interpretasi penulis{]} Temuan memperingatkan bahwa pusat utama yang
besar dapat membuat suatu region tampak lebih hierarkis secara
morfologis, sementara arus internalnya relatif seimbang. Sebaliknya,
keberadaan beberapa pusat dengan ukuran lebih seimbang tidak menjamin
arus komuter atau belanja yang multidireksional. Pembacaan seperti ini
penting ketika konsep polisentris digunakan untuk mendukung tujuan
kohesi dan daya saing, karena penulis menyatakan kombinasi kedua tujuan
tersebut tidak otomatis terbukti (locator: TXT baris 380-404 dan
824-842; hlm. 1132-1133 dan 1143).

{[}Laporan sumber{]} Artikel menganjurkan basis bukti sebelum manfaat
pembangunan polisentris atau keberlanjutannya dijadikan dasar kebijakan
dan investasi publik. Manfaat yang diklaim perlu diuji dengan membedakan
MP dan FP serta mempertimbangkan kepadatan jaringan dan kemampuan
menarik arus dari jarak yang lebih jauh (locator: TXT baris 901-925;
hlm. 1144-1145).

{[}Inferensi pengolahan{]} TXT tidak memberikan instrumen kebijakan,
urutan intervensi, atau ambang nilai untuk menentukan region
polisentris. Maka implikasi praktis yang dapat dipertanggungjawabkan
terbatas pada perbaikan diagnosis dan evaluasi, bukan rekomendasi
intervensi spesifik (locator: TXT baris 845-925).

\subsection{Applications}\label{applications-45}

{[}Laporan sumber{]} Aplikasi empiris aktual artikel adalah pengukuran
struktur spasial di 42 WGR yang mencakup seluruh Belanda, menggunakan
data pekerjaan/perjalanan kerja dan belanja. Kerangka yang sama dipakai
untuk membandingkan region, mengidentifikasi perbedaan MP-FP, dan
memeriksa kaitan perbedaan itu dengan orientasi lokal, orientasi
eksternal, serta ukuran pusat utama (locator: TXT baris 470-526 dan
581-621; hlm. 1134-1137).

{[}Laporan sumber{]} Penulis menyatakan kerangka tersebut pada dasarnya
bebas skala dan dapat diterapkan pada entitas spasial mulai dari kota
individual sampai benua. Namun, dalam file ini pengujian empiris yang
benar-benar dilaporkan hanya mencakup WGR di Belanda; penerapan pada
skala atau negara lain belum dilakukan (locator: TXT baris 507-520 dan
920-932; hlm. 1135 dan 1145).

{[}Inferensi pengolahan{]} Aplikasi yang aman dari artikel ini adalah
sebagai perangkat klasifikasi dan perbandingan struktur urban, bukan
sebagai alat untuk menyimpulkan kinerja atau menetapkan kebijakan secara
langsung. Setiap aplikasi perlu mempertahankan pemisahan antara
nodalitas, sentralitas internal, sentralitas eksternal, dan kepadatan
jaringan (locator: TXT baris 406-465 dan 612-631).

\subsection{Future Research
(Optional)}\label{future-research-optional-45}

{[}Laporan sumber{]} Agenda penelitian yang dinyatakan penulis mencakup
pengujian pada negara lain dan skala lain, termasuk negara atau
makroregion lintas batas, untuk menilai generalisasi hasil Belanda.
Penulis secara khusus mengingatkan bahwa karakteristik kepadatan dan
ukuran kota di Belanda dapat membuat hasilnya tidak langsung berlaku di
tempat lain (locator: TXT baris 920-932; hlm. 1145).

{[}Laporan sumber{]} Penulis menyarankan pemeriksaan arus lain, terutama
perdagangan antarperusahaan dan hubungan pemegang saham. Karena arus
tersebut biasanya memiliki jangkauan geografis lebih luas, sentralitas
eksternal dapat berperan lebih besar daripada dalam data komuter dan
belanja. Agenda lain adalah menguji apakah manfaat polisentrisitas dan
kebijakan pembangunan polisentris benar-benar ada, apakah kebijakannya
berkelanjutan, dan bagaimana hubungannya dengan kinerja regional
(locator: TXT baris 934-942 dan 901-925; hlm. 1145).

{[}Laporan sumber{]} Penelitian lanjutan tersebut perlu tetap membedakan
MP dan FP, serta mempertimbangkan kepadatan jaringan dan kapasitas
region menarik arus dari jauh. TXT tidak menyebutkan desain sampel,
teknik estimasi, atau hipotesis spesifik untuk agenda tersebut (locator:
TXT baris 918-925).

\subsection{Provenance Notes}\label{provenance-notes-45}

\begin{itemize}
\tightlist
\item
  Bahan substantif review ini hanya file
  \texttt{/Users/mac/Documents/Mac/{[}2{]}\ Obsidian\ Vault/zettelkasten/source/journal/B02-burger-meijers-2012-form-follows-function-linking-morphological-and-functional-polycentricity-10.1177-0042098011407095.txt}.
  File dibaca dari baris 1 sampai 1140, termasuk blok metadata
  ekstraksi, teks artikel, catatan, ucapan terima kasih, dan daftar
  referensi.
\item
  Blok metadata TXT mencatat sumber PDF, tanggal ekstraksi 2026-08-31,
  metode \texttt{pdftotext\ -layout}, dan 24 halaman. Review ini tidak
  membuka PDF, URL jurnal, atau sumber yang tercantum di daftar
  referensi; semua klaim substantif dibatasi pada TXT (locator: TXT
  baris 1-24 dan 25-69).
\item
  Identitas bibliografis utama berasal dari judul, penulis, jurnal,
  tahun, halaman awal, dan DOI yang tercantum pada TXT (locator: TXT
  baris 32-38 dan 70-79). Status peer-review tidak dinyatakan dalam
  file, sehingga tidak disimpulkan.
\item
  Locator utama memakai nomor halaman cetak artikel yang muncul di TXT,
  bagian, tabel, atau gambar; locator tambahan memakai nomor baris TXT
  agar klaim dapat ditelusuri langsung. Tabel 1 dan Tabel 2 hanya
  diringkas dalam paragraf atau poin; tidak ada tabel baru dalam review
  ini.
\item
  {[}Laporan sumber{]} digunakan untuk parafrasa isi, definisi,
  prosedur, dan hasil yang dilaporkan artikel. {[}Interpretasi
  penulis{]} digunakan untuk penjelasan mekanisme dan implikasi yang
  dinyatakan Burger dan Meijers. {[}Inferensi pengolahan{]} digunakan
  untuk kesimpulan kehati-hatian yang saya tarik dari teks; label
  tersebut tidak mengubahnya menjadi klaim penulis.
\item
  Informasi yang tidak tersedia dalam TXT dinyatakan secara eksplisit,
  terutama ukuran sampel survei, rincian sampling dan nonrespons,
  penanganan data hilang, nilai-p atau interval kepercayaan, serta
  spesifikasi statistik lengkap untuk beberapa gambar.
\item
  TXT memuat kontradiksi internal mengenai arah perbandingan MP-FP:
  abstrak menyatakan mayoritas region lebih morfologis, sedangkan hasil
  utama dan kesimpulan menyatakan hampir semua lebih fungsional.
  Subbagian 5.3 juga membalik contoh Groningen dan Veendam dibandingkan
  Tabel 1 serta uraian sesudahnya. Review tidak memilih salah satunya
  secara diam-diam; setiap versi dilaporkan bersama locator
  masing-masing (locator: TXT baris 83-95, 642-648, 666-714, 775-792,
  837-842, dan 856-866).
\item
  File keluaran ini adalah review B02 saja. Tidak ada perubahan pada
  file TXT sumber, dokumen induk, atau review kode lain. Setelah audit
  struktur dan evidence, checkbox B02 dicentang pada \texttt{Todo.md}
  dan daftar kode; checkbox kode lain tidak diubah.
\end{itemize}

\section{Review B03}\label{review-b03}

\textbf{Identitas sumber}

\begin{itemize}
\tightlist
\item
  Kode kajian: B03.
\item
  Penulis: Bastiaan De Goei, Martijn J. Burger, Frank G. Van Oort, dan
  Michael Kitson.
\item
  Tahun: 2010.
\item
  Judul: \emph{Functional Polycentrism and Urban Network Development in
  the Greater South East, United Kingdom: Evidence from Commuting
  Patterns, 1981-2001}.
\item
  Jurnal: \emph{Regional Studies}, volume 44, nomor 9, hlm. 1149-1170.
\item
  DOI: \texttt{10.1080/00343400903365102}.
\item
  Publikasi daring: 8 Februari 2010; artikel tercetak sebagai edisi
  November 2010.
\item
  File yang ditelaah:
  \texttt{source/journal/B03-de-goei-burger-van-oort-kitson-2010-functional-polycentrism-and-urban-network-development-in-the-greater-south-east-united-kingdom-evidence-from-commuting-patterns-1981-2001-10.1080-00343400903365102.txt}.
\item
  Cakupan bahan: seluruh teks yang tersedia pada TXT, yaitu baris
  1-1273, termasuk metadata ekstraksi PDF, isi artikel, tabel, gambar
  yang memiliki keterangan, lampiran, catatan, dan daftar referensi.
\item
  Status publikasi: TXT mencatat artikel diterima pada April 2008 dan
  direvisi pada April 2009 serta menyebut komentar dua anonymous
  referees, tetapi status peer-review tidak dinyatakan secara eksplisit
  dan tidak diverifikasi di luar TXT.
\end{itemize}

\subsection{Summarized Abstract}\label{summarized-abstract-46}

\textbf{Laporan penulis.} Artikel menguji bagaimana jaringan perkotaan
berkembang di Greater South East, Inggris, dengan menggunakan pola
komuter pada 1981, 1991, dan 2001. Titik berangkatnya adalah argumen
dalam literatur kontemporer bahwa konseptualisasi \emph{central place}
yang tradisional perlu dilengkapi atau digantikan oleh pandangan
jaringan yang menekankan pola ketergantungan silang antarsatuan ruang.
Hasil empiris menunjukkan bahwa Greater South East belum dapat dicirikan
sebagai wilayah perkotaan polisentris atau jaringan perkotaan yang
terintegrasi. Namun, terdapat bukti perkembangan jaringan pada tingkat
lokal atau intra-perkotaan dan desentralisasi sistem pada tingkat
regional atau inter-perkotaan. {[}Laporan penulis; TXT baris 113-122,
PDF hlm. 1149{]}

\textbf{Inferensi pengolahan.} Kesimpulan abstrak bersifat bersyarat
pada ukuran fungsional yang dipakai: polisentrisme fungsional
diperlakukan sebagai integrasi jaringan perkotaan yang terlihat dari
kekuatan relatif arus komuter setelah ukuran distrik dan jarak
dikendalikan. Dengan demikian, keberadaan banyak pusat secara morfologis
saja tidak dianggap cukup sebagai bukti jaringan. {[}TXT baris 233-248,
331-352, dan 671-701, PDF hlm. 1151 dan 1160{]}

\textbf{Inferensi pengolahan dengan batasan cakupan.} Ringkasan tidak
berarti bahwa sama sekali tidak ada hubungan antarkota. Klaim yang
didukung TXT adalah bahwa belum ada jaringan yang sepenuhnya
terintegrasi; hubungan lokal dan sebagian pergeseran menuju
desentralisasi tetap ditemukan. {[}TXT baris 893-915, PDF hlm. 1163{]}

\subsection{Problem Statement}\label{problem-statement-46}

\textbf{Laporan penulis.} Masalah penelitian muncul dari ketegangan
antara dua cara membaca sistem perkotaan. Model \emph{central place}
memandang hubungan kota-hinterland sebagai hierarkis dan terpusat,
sedangkan model jaringan menekankan interdependensi yang saling silang
pada skala intra-perkotaan dan inter-perkotaan. Penulis menyatakan bahwa
riset empiris yang menguji seberapa cocok model jaringan dengan
kenyataan masih terbatas. Banyak studi mengandalkan karakteristik
simpul, seperti \emph{location quotient}, hubungan pangkat-peringkat,
indeks kecukupan, atau rasio pekerjaan terhadap penduduk bekerja, bukan
karakteristik arus. Studi berbasis arus yang ada juga sering hanya
menggambarkan satu titik waktu, sehingga arah perubahan sistem belum
jelas. {[}Laporan penulis; TXT baris 219-248, PDF hlm. 1151{]}

\textbf{Laporan penulis.} Greater South East menjadi kasus problematis
karena dipromosikan oleh peneliti dan lembaga kebijakan sebagai jaringan
perkotaan atau wilayah polisentris yang terintegrasi, sementara ukuran
London dapat secara otomatis menghasilkan dugaan sistem monosentris.
Penulis juga menunjukkan bahwa visualisasi arus mentah dapat bias oleh
ukuran populasi dan pekerjaan serta kedekatan fisik London. Karena itu,
pertanyaan tidak cukup dijawab dengan melihat banyaknya pusat atau
menggambar garis arus; pengujian perlu mengendalikan massa dan jarak.
{[}TXT baris 438-469 dan 485-539, PDF hlm. 1155-1157{]}

\textbf{Inferensi pengolahan.} Masalah inti artikel adalah apakah label
kebijakan ``Greater South East yang terintegrasi'' benar-benar
ditunjukkan oleh interaksi fungsional, dan apakah hubungan tersebut
berubah sepanjang waktu. Artikel dengan sengaja membedakan perkembangan
lokal dari integrasi regional, sehingga ``desentralisasi'' tidak
otomatis diperlakukan sebagai ``jaringan terintegrasi''. {[}TXT baris
463-469, 514-532, dan 855-865, PDF hlm. 1155-1157 dan 1162{]}

\textbf{Inferensi pengolahan.} Rumusan hipotesis bernomor atau daftar
\emph{problem statement} formal: \textbf{Tidak disebutkan dalam file}.
Kriteria pengujian dijelaskan dalam bagian metode melalui perbandingan
kekuatan relatif tiga belas jenis interdependensi. {[}TXT baris 629-718,
PDF hlm. 1159-1160{]}

\subsection{Objectives}\label{objectives-46}

\textbf{Laporan penulis.} Tujuan eksplisitnya adalah menginvestigasi
apakah Greater South East terdiri atas \emph{network-cities} pada skala
intra-perkotaan dan apakah distrik-distrik yang bersebelahan membentuk
jaringan perkotaan berskala lebih besar yang sepenuhnya terintegrasi.
Pengujian dilakukan dengan mengendalikan massa distrik dan jarak fisik.
{[}Laporan penulis; TXT baris 510-521, PDF hlm. 1156{]}

\textbf{Inferensi pengolahan.} Tujuan empiris tersebut dapat diringkas
menjadi dua pertanyaan analitis. Pertama, bagaimana struktur rata-rata
hubungan komuter pada skala intra-perkotaan dan inter-perkotaan? Kedua,
bagaimana kekuatan relatif berbagai rezim interdependensi itu berubah
antara 1981 dan 2001? Perubahan diuji melalui \emph{time dummies} dan
interaksi atau \emph{slope dummies} antara waktu dan rezim spasial.
{[}TXT baris 629-718, PDF hlm. 1159-1160{]}

\textbf{Inferensi pengolahan.} Desain juga berfungsi menguji apakah
perpindahan dari sistem terpusat menuju sistem polisentris dapat diamati
sebagai tren bertahap, bukan hanya sebagai kondisi pada satu tahun. Ini
merupakan pembacaan atas tujuan yang dinyatakan, bukan tujuan tambahan
yang ditulis sebagai daftar tersendiri. {[}TXT baris 233-248 dan
707-718, PDF hlm. 1151 dan 1160{]}

\textbf{Inferensi pengolahan.} Tujuan yang berkaitan dengan evaluasi
kausal atas kebijakan tertentu, pengukuran kesejahteraan rumah tangga,
atau pengujian dampak ekonomi regional secara langsung: \textbf{Tidak
disebutkan dalam file}. {[}TXT keseluruhan; tidak ada tujuan tersebut
dalam teks artikel{]}

\subsection{Literature Survey}\label{literature-survey-46}

\textbf{Laporan penulis.} Posisi teoretisnya menempatkan sistem
perkotaan sebagai himpunan kota yang saling bergantung secara
fungsional. Struktur sistem dapat bergerak dari monosentris ke
polisentris dan dominasi struktur dapat berbeda menurut skala. Model
monosentris yang berakar pada Burgess, lalu diperluas oleh Alonso dan
Muth, menggambarkan pusat bisnis sebagai inti kegiatan ekonomi dengan
suburb sebagai pemasok tenaga kerja. Perubahan harga lahan, amenitas,
biaya transportasi, dan penyebaran pekerjaan kemudian dapat mendorong
terbentuknya pusat-pusat lokal, arus komuter timbal balik, dan suburb
yang lebih mandiri. {[}TXT baris 273-315, PDF hlm. 1152{]}

\textbf{Laporan penulis.} Pada skala inter-perkotaan, literatur yang
dibahas mengarah pada gagasan \emph{polycentric urban region}, yaitu
beberapa wilayah metropolitan yang secara historis dan spasial terpisah
tetapi berpotensi membentuk jaringan. Penulis membedakan jaringan nodal,
yang terutama menghubungkan pusat-pusat lama, dari jaringan yang
sepenuhnya terintegrasi, yang juga menghubungkan pusat dengan suburb
lintas wilayah dan suburb dengan suburb lintas wilayah. Penulis
mengingatkan bahwa jaringan inter-perkotaan tidak harus menjadi tahap
evolusi berikutnya setelah \emph{network-city} di tingkat
intra-perkotaan. {[}TXT baris 326-352 dan 669-718, PDF hlm. 1153-1160{]}

\textbf{Laporan penulis.} Istilah juga dibatasi secara konseptual. Dalam
artikel ini, ``functional polycentrism'' disamakan dengan ``urban
network integration'', walaupun catatan penulis mengakui bahwa sebagian
sarjana ingin menyimpan ``polycentricity'' untuk polisentrisme
morfologis dan menggunakan ``urban networks'' untuk hubungan fungsional.
Penulis memilih fokus pada jaringan fungsional antarkota, bukan sekadar
jumlah pusat. {[}TXT baris 233-248 dan 1058-1074, PDF hlm. 1151 dan
catatan 2-3{]}

\textbf{Laporan penulis.} Survei literatur mengelompokkan penggerak
perubahan menjadi fleksibilitas dan mobilitas perusahaan, fleksibilitas
dan mobilitas rumah tangga, serta kebijakan lokal dan regional.
Hipotesis restrukturisasi mengaitkan penyebaran pekerjaan dengan
teknologi informasi dan komunikasi serta pergeseran dari manufaktur ke
jasa. Hipotesis dekonsentrasi mengaitkan perubahan dengan preferensi
tempat tinggal, mobilitas, demografi, dan fleksibilitas tempat kerja.
Kebijakan pertumbuhan, pengendalian lahan, Green Belt, serta kerja sama
antarkota dapat mendorong atau, sebagai produk samping yang tidak
disengaja, membentuk jaringan. {[}TXT baris 355-404, PDF hlm.
1154-1155{]}

\textbf{Laporan penulis.} Literatur dan dokumen kebijakan yang dikutip
menggambarkan Greater South East sebagai kawasan yang berpusat pada
London tetapi memiliki pusat kegiatan seperti Cambridge dan Oxford serta
berbagai koridor pertumbuhan. Pada saat yang sama, Hall dan Green serta
Hall dan Pain dilaporkan menyatakan bahwa relasi fungsional yang tidak
terkait London masih terbatas. Penulis menggunakan ketegangan ini
sebagai konteks untuk pengujian kuantitatif, bukan sebagai bukti awal
yang diterima begitu saja. {[}TXT baris 405-469, PDF hlm. 1155-1156{]}

\textbf{Inferensi pengolahan.} Bagian ini merupakan tinjauan naratif
atas teori, studi empiris, dan wacana kebijakan. Protokol pencarian,
kriteria inklusi-eksklusi, ukuran korpus, atau prosedur \emph{systematic
review}: \textbf{Tidak disebutkan dalam file}. Referensi yang
dicantumkan adalah sumber yang dipakai oleh artikel; tidak ada yang
diverifikasi secara terpisah dalam review ini. {[}TXT keseluruhan bagian
literatur dan daftar referensi{]}

\subsection{Method Used}\label{method-used-46}

\textbf{Inferensi pengolahan.} Desain ini dapat diklasifikasikan sebagai
analisis longitudinal berbasis arus pada distrik-distrik di Greater
South East untuk tiga tahun sensus: 1981, 1991, dan 2001. {[}TXT baris
485-521, PDF hlm. 1156{]}

\textbf{Laporan penulis.} Unit analisisnya adalah interaksi komuter
antara atau di dalam distrik asal dan tujuan, dengan dua skala:
intra-perkotaan, yaitu arus yang tetap berada dalam satu wilayah
perkotaan, dan inter-perkotaan, yaitu arus antarkawasan perkotaan.
{[}TXT baris 485-521 dan catatan 4 pada baris 1042-1057, PDF hlm. 1156
dan 1167{]}

\textbf{Laporan penulis.} Interaksi dimodelkan sebagai fungsi dari massa
distrik asal, massa distrik tujuan, dan jarak fisik. Dalam formulasi
awal, intensitas interaksi berbanding lurus dengan massa kedua lokasi
dan berbanding terbalik dengan jarak. Jarak diukur sebagai jarak jalan
aktual antara distrik, bukan jarak udara. Karena jarak antara pusat
gravitasi distrik tetangga dapat terlalu besar, variabel
\emph{contiguity} dimasukkan sebagai koreksi. {[}TXT baris 553-595 dan
961-977, PDF hlm. 1158 dan 1164-1165{]}

\textbf{Laporan penulis.} Dalam spesifikasi estimasi yang paling dasar,
efek tetap asal dan tujuan menggantikan variabel massa serta
mengendalikan karakteristik spesifik distrik. Efek tetap waktu
ditambahkan untuk mengendalikan perbedaan antarperiode. Penulis
menyatakan bahwa spesifikasi efek tetap asal-tujuan dimaksudkan agar
jumlah komuter yang diprediksi konsisten dengan jumlah yang diamati pada
distrik asal dan tujuan. {[}TXT baris 559-595 dan 760-762, PDF hlm.
1158-1159 dan tabel 2{]}

\textbf{Laporan penulis.} Karena variabel hasil berupa jumlah komuter,
penulis tidak menjadikan OLS sebagai pilihan utama dan menggunakan
keluarga model Poisson. Pengujian dan distribusi yang dipilih berbeda
menurut skala: regresi binomial negatif digunakan untuk model
intra-perkotaan karena \emph{over-dispersion}, sedangkan regresi Poisson
zero-inflated digunakan untuk model inter-perkotaan karena adanya banyak
nol. Tabel 3 menyatakan bahwa bagian inflasi juga diestimasi menggunakan
log jarak jalan, \emph{contiguity}, dan interdependensi. {[}TXT baris
595-623 dan 790-805, PDF hlm. 1158-1159 dan 1161{]}

\textbf{Laporan penulis.} Tiga belas rezim interdependensi spasial
dimasukkan sebagai variabel indikator. Pada skala intra-perkotaan,
rezimnya mencakup arus di dalam pusat kota, di dalam suburb, suburb
-\textgreater{} pusat kota, pusat kota -\textgreater{} suburb, dan
suburb -\textgreater{} suburb. Pada skala inter-perkotaan, rezimnya
mencakup hubungan antarpusat, pusat -\textgreater{} suburb, suburb
-\textgreater{} pusat, dan suburb -\textgreater{} suburb antarkawasan,
serta empat arah hubungan antara wilayah Greater South East dan London.
{[}TXT baris 624-667, PDF hlm. 1159{]}

\textbf{Laporan penulis.} Kategori referensi pada analisis
intra-perkotaan adalah interdependensi yang seluruhnya berada di dalam
distrik pusat. Kategori referensi pada analisis inter-perkotaan adalah
arus yang berada di dalam wilayah perkotaan. Polisentrisme atau
integrasi jaringan dinilai dari apakah hubungan yang menjadi referensi
masih lebih kuat dan apakah hierarki antarkategori tetap terlihat
setelah massa, jarak, serta efek tetap dikendalikan. {[}TXT baris
669-718, PDF hlm. 1160{]}

\textbf{Laporan penulis.} Model 1 dan Model 3 menguji struktur umum
selama 1981-2001, sedangkan Model 2 dan Model 4 menambahkan tren linear
waktu serta interaksi waktu dengan rezim spasial. Penulis sendiri
mencatat bahwa hanya tren linear yang diuji. {[}TXT baris 720-724,
790-808, dan catatan 12 pada baris 1105-1107, PDF hlm. 1160-1162 dan
1168{]}

\textbf{Inferensi pengolahan.} Metode ini menguji asosiasi antara tipe
hubungan spasial dan intensitas komuter, bukan efek kausal dari Green
Belt, infrastruktur, relokasi pekerjaan, atau kebijakan tertentu.
Perlakuan eksperimental, strategi identifikasi kausal, atau variabel
kejutan eksternal sebagai bagian dari estimasi utama: \textbf{Tidak
disebutkan dalam file}. {[}TXT baris 701-718, 899-905, dan 916-930, PDF
hlm. 1160-1163{]}

\subsection{Dataset}\label{dataset-46}

\textbf{Laporan penulis.} Data interaksi komuter berasal dari
\emph{Special Workplace Statistics (Set C)} dalam sensus Inggris untuk
1981, 1991, dan 2001. Catatan sumber menyebutkan bahwa data sensus 1981
diestimasi ulang untuk batas 1991, data 1991 berasal dari \emph{Special
Workplace Statistics (Set C)}, dan data 2001 berasal dari \emph{Special
Workplace Statistics (Level 1)}. {[}TXT baris 505-521 dan 1075-1082, PDF
hlm. 1156 dan catatan 6 pada hlm. 1167-1168{]}

\textbf{Laporan penulis.} Penulis menggunakan \emph{CIDS 1991/2001
Common Geography} untuk mengurangi masalah perubahan batas distrik.
Dengan batas tersebut, Greater South East terdiri atas 146 distrik dan
27 wilayah perkotaan, masing-masing dengan satu distrik inti. Wilayah
perkotaan terutama dibangun dari NUTS-III yang sedikit diadaptasi; TTWA
digunakan ketika NUTS-III menghasilkan pembagian yang tidak sesuai.
Definisi kawasan yang dipilih mengikuti tiga \emph{Governmental Office
Regions}: London, South East, dan East of England, karena dianggap
paling sesuai dengan realitas politik yang dibahas artikel. {[}TXT baris
506-521 dan 938-977, PDF hlm. 1156 dan 1164-1165{]}

\textbf{Laporan penulis.} Jarak jalan dihitung dari jaringan jalan
Greater South East tahun 2005 melalui matriks biaya asal-tujuan, dengan
biaya ditetapkan sebagai jarak. Jarak tersebut diterapkan pada ketiga
periode dengan alasan bahwa sebagian besar jalan umum telah dibangun
sejak awal abad ke-20 dan peningkatan jalan tidak mengubah jarak aktual.
Jarak intra-distrik dihitung dari luas distrik dengan asumsi wilayah
berbentuk lingkaran, sedangkan \emph{contiguity} dipakai untuk menangani
kemungkinan jarak pusat gravitasi yang berlebihan pada distrik yang
berbatasan. {[}TXT baris 580-595 dan 1039-1046, PDF hlm. 1158-1159 dan
1167{]}

\textbf{Laporan penulis.} Gambar 3 hanya menampilkan \emph{net flows} di
atas ambang 50 komuter. Visualisasi tersebut tidak mengendalikan ukuran
absolut distrik dan jarak; keterbatasan ini menjadi alasan penulis
beralih ke model gravitasi. {[}TXT baris 522-539 dan 550-552, PDF hlm.
1156-1157{]}

\textbf{Inferensi pengolahan.} Ambang 50 komuter digunakan untuk
visualisasi, bukan dinyatakan sebagai kriteria seleksi observasi model.
{[}TXT baris 522-523 dan 550-552, PDF hlm. 1156-1157{]}

\textbf{Laporan penulis.} Nilai sekitar 21 juta penduduk dan GDP sekitar
US\$900 miliar adalah deskripsi konteks Greater South East, bukan ukuran
observasi dataset yang dirinci sebagai sampel. {[}TXT baris 438-462, PDF
hlm. 1155{]}

\textbf{Inferensi pengolahan.} Jumlah total komuter dalam seluruh
wilayah, matriks arus mentah, kamus data lengkap, bobot atau prosedur
penanganan data hilang, serta kode estimasi: \textbf{Tidak disebutkan
dalam file}. {[}TXT bagian Dataset dan Methods{]}

\subsection{Population / Sample}\label{population-sample-46}

\textbf{Inferensi pengolahan.} Studi tidak menggunakan sampel responden
dalam pengertian survei yang dilaporkan secara terpisah.

\textbf{Laporan penulis.} Populasi empiris yang tersedia dalam desain
adalah arus komuter yang tercatat antara dan di dalam 146 distrik pada
geografi bersama di Greater South East, untuk tiga tahun sensus. Dua
puluh tujuh wilayah perkotaan digunakan untuk mengelompokkan distrik
menjadi inti dan distrik perifer. {[}TXT baris 506-521 dan 943-977, PDF
hlm. 1156 dan 1164-1165{]}

\textbf{Laporan penulis.} Tabel 2 melaporkan 2.052 observasi untuk model
intra-perkotaan. Tabel 3 melaporkan 60.681 observasi untuk model
inter-perkotaan. Angka tersebut adalah jumlah observasi estimasi, bukan
jumlah orang, rumah tangga, atau jumlah distrik unik. {[}TXT baris
734-768 dan 804-850, PDF hlm. 1161-1162{]}

\textbf{Inferensi pengolahan.} Jumlah pasangan distrik unik, perlakuan
terhadap arus nol pada setiap tahap, atau apakah setiap arus merupakan
hitungan satu arah yang diduplikasi sebagai pasangan arah: \textbf{Tidak
disebutkan dalam file}. {[}TXT baris 734-768 dan 804-850, PDF hlm.
1161-1162{]}

\textbf{Inferensi pengolahan.} Secara geografis, cakupan mengikuti
wilayah dan klasifikasi yang ditetapkan penulis, bukan pengambilan
sampel probabilitas. Karena unitnya adalah arus antar-distrik, hasil
paling langsung berlaku pada konfigurasi distrik dan periode tersebut.
{[}TXT baris 510-521, 943-977, dan 995-1032, PDF hlm. 1156, 1164-1166{]}

\textbf{Inferensi pengolahan.} Ukuran sampel individual, tingkat
respons, karakteristik demografis responden, atau aturan
inklusi-eksklusi di tingkat pekerja: \textbf{Tidak disebutkan dalam
file}. TXT hanya menyediakan sumber statistik sensus dan struktur
agregat wilayah. {[}TXT baris 505-521 dan 1075-1082, PDF hlm. 1156 dan
catatan 6{]}

\subsection{Dependent Variables}\label{dependent-variables-46}

\textbf{Laporan penulis.} Variabel dependen utama adalah intensitas
interaksi \texttt{Iij}, yang dijelaskan sebagai jumlah komuter antara
wilayah asal \texttt{i} dan wilayah tujuan \texttt{j}, termasuk
interaksi di dalam distrik dan antardistrik. Variabel ini diperlakukan
sebagai data hitungan dan dimodelkan melalui nilai rata-rata kondisional
\texttt{mij}. {[}TXT baris 553-578 dan 597-623, PDF hlm. 1158-1159{]}

\textbf{Inferensi pengolahan.} Karena asal dan tujuan dibedakan, arah
arus penting dalam kategori seperti suburb -\textgreater{} pusat kota
dan pusat kota -\textgreater{} suburb. {[}TXT baris 624-667, PDF hlm.
1159{]}

\textbf{Laporan penulis.} Analisis memiliki dua penerapan variabel hasil
yang sama pada skala berbeda. Analisis intra-perkotaan melihat arus yang
berada dalam satu wilayah perkotaan; analisis inter-perkotaan melihat
arus antarwilayah perkotaan, termasuk hubungan dengan London. Tahun
1981, 1991, dan 2001 ditangani melalui efek tetap waktu dan, pada Model
2 serta Model 4, tren interaksi dengan rezim spasial. {[}TXT baris
629-718 dan 720-848, PDF hlm. 1159-1162{]}

\textbf{Inferensi pengolahan.} \emph{Net flow} dengan ambang lebih dari
50 komuter pada Gambar 3 bukan variabel dependen yang terpisah,
melainkan cara visualisasi arus. Variabel jarak jalan,
\emph{contiguity}, rezim interdependensi, serta efek tetap asal, tujuan,
dan waktu berperan sebagai prediktor atau pengendali. {[}TXT baris
522-539, 559-595, dan 760-762, PDF hlm. 1156-1159 dan tabel 2{]}

\textbf{Inferensi pengolahan.} Ukuran variabel dependen pada tingkat
individu, distribusi lengkap hitungan arus, atau transformasi tambahan
selain formulasi gravitasi dan spesifikasi model hitungan yang
dijelaskan: \textbf{Tidak disebutkan dalam file}. {[}TXT baris 553-623
dan 734-850, PDF hlm. 1158-1162{]}

\subsection{Results}\label{results-46}

\textbf{Laporan penulis (gambaran awal).} Gambar 3 menunjukkan kenaikan
absolut jumlah komuter antara 1981 dan 2001, tetapi arus relatif masih
banyak terjadi di dalam wilayah perkotaan masing-masing. Arus
antarkawasan terutama mengarah ke London, membentuk pola hub-and-spoke,
sedangkan aktivitas komuter antara wilayah South East dan East of
England tampak terbatas. Perubahan besar tidak terlihat secara visual
selain perubahan intensitas lokal, misalnya peningkatan di sekitar
Norwich dan Crawley serta penurunan di sekitar Oxford dan Portsmouth.
{[}TXT baris 522-539, PDF hlm. 1156-1157{]}

\textbf{Laporan penulis (model intra-perkotaan).} Model 1 pada Tabel 2
menggunakan regresi binomial negatif. Koefisien log jarak jalan adalah
-1,362 dengan standard error 0,049 dan signifikan pada p \textless{}
0,01; penulis menafsirkannya sebagai penurunan prediksi komuter sebesar
1,36\% untuk kenaikan jarak fisik 1\%. Koefisien \emph{contiguity}
adalah sekitar 1,34 dan signifikan pada p \textless{} 0,01; teks
menafsirkan distrik yang berbatasan memiliki arus sekitar 3,85 kali
lebih tinggi daripada pasangan yang tidak berbatasan, dengan massa dan
jarak dikendalikan. {[}TXT baris 724-733 dan 734-770, PDF hlm.
1160-1161{]}

\textbf{Laporan penulis.} Dengan kategori referensi berupa arus di dalam
pusat kota, arus di dalam suburb tidak berbeda signifikan dari kategori
referensi pada Model 1. Sebaliknya, koefisien suburb -\textgreater{}
pusat kota adalah -1,630, pusat kota -\textgreater{} suburb -1,557, dan
suburb -\textgreater{} suburb -2,224; ketiganya signifikan pada p
\textless{} 0,01. Dalam istilah prediksi penulis, arus di dalam pusat
kota sekitar lima kali lebih tinggi daripada hubungan suburb
-\textgreater{} pusat kota dan pusat kota -\textgreater{} suburb, serta
sekitar sembilan kali lebih tinggi daripada arus suburb -\textgreater{}
suburb. Hubungan suburb -\textgreater{} pusat kota juga secara
signifikan lebih kuat daripada hubungan suburb -\textgreater{} suburb,
berdasarkan uji Wald chi-square 12,05 dengan 1 derajat kebebasan dan p
\textless{} 0,01. {[}TXT baris 688-705 dan 740-748, PDF hlm.
1160-1161{]}

\textbf{Laporan penulis.} Model 2 menambahkan tren linear. Dibandingkan
hubungan di dalam pusat kota, rasio kekuatan pusat kota tetap lebih
besar tetapi menyempit: sekitar 6,4 kali terhadap pusat kota
-\textgreater{} suburb pada 1981 menjadi 4,1 kali pada 2001, dan sekitar
12,2 kali terhadap suburb -\textgreater{} pusat kota pada 1981 menjadi
8,2 kali pada 2001. Interaksi waktu untuk pusat kota -\textgreater{}
suburb sebesar 0,216 dan untuk suburb -\textgreater{} suburb sebesar
0,199, keduanya signifikan pada p \textless{} 0,05. Tren untuk arus di
dalam suburb dan suburb -\textgreater{} pusat kota tidak signifikan pada
tingkat yang ditandai dalam tabel; penulis tetap melaporkan kenaikan
relatif suburb -\textgreater{} pusat kota, tetapi tidak signifikan.
{[}TXT baris 723-732, 773-789, dan 754-759, PDF hlm. 1160-1161{]}

\textbf{Laporan penulis.} Dengan mengekstrapolasi tren yang diamati
tanpa kejutan eksternal, penulis memperkirakan bahwa diperlukan
setidaknya sekitar 80 tahun lagi sebelum dapat disebut
\emph{network-city} yang tepat pada skala intra-perkotaan. {[}TXT baris
773-779 dan 899-905, PDF hlm. 1160 dan 1163{]}

\textbf{Inferensi pengolahan.} Ekstrapolasi tersebut berbasis perbedaan
magnitudo, bukan prediksi yang memasukkan pembukaan infrastruktur,
penciptaan atau relokasi pekerjaan, atau pembangunan perumahan baru.
{[}TXT baris 899-905, PDF hlm. 1163{]}

\textbf{Laporan penulis (model inter-perkotaan).} Model 3 pada Tabel 3
menggunakan regresi Poisson zero-inflated. Koefisien log jarak jalan
adalah -1,931 dengan standard error 0,018 dan \emph{contiguity} sebesar
0,636 dengan standard error 0,039; keduanya signifikan pada p
\textless{} 0,01. Dengan arus di dalam wilayah perkotaan sebagai
referensi, hubungan antarpusat antarwilayah South East dan East of
England memiliki koefisien -0,461, suburb -\textgreater{} pusat kota
-0,766, pusat kota -\textgreater{} suburb -0,781, dan suburb
-\textgreater{} suburb -0,665; seluruhnya signifikan pada p \textless{}
0,01. Penulis menyatakan bahwa, dengan massa dan jarak konstan, arus di
dalam wilayah perkotaan rata-rata sekitar 100\% lebih besar daripada
arus antardistrik yang melampaui batas wilayah perkotaan. {[}TXT baris
773-777, 790-802, dan 810-824, PDF hlm. 1160-1162{]}

\textbf{Interpretasi penulis.} Di antara tipe hubungan inter-perkotaan
South East dan East of England, hubungan antarpusat lebih kuat daripada
hubungan pusat-suburb lintas wilayah, suburb-pusat lintas wilayah, dan
suburb-suburb. Penulis menafsirkan bahwa jaringan antarkawasan yang
muncul terutama berupa hubungan antarwilayah inti, bukan jaringan penuh
yang juga menghubungkan seluruh suburb. {[}TXT baris 773-789, PDF hlm.
1160-1161{]}

\textbf{Laporan penulis.} Hubungan dari pusat kota dan suburb Greater
South East menuju London tidak secara signifikan lebih kecil daripada
arus di dalam wilayah perkotaan, dengan koefisien masing-masing -0,076
dan -0,150 yang tidak signifikan. Arah sebaliknya, London
-\textgreater{} pusat kota Greater South East sebesar -1,882 dan London
-\textgreater{} suburb Greater South East sebesar -1,622, signifikan
pada p \textless{} 0,01. {[}TXT baris 795-829, PDF hlm. 1161-1162{]}

\textbf{Interpretasi penulis.} Penulis menafsirkan asimetri ini sebagai
banyak pekerja bergerak dari South East dan East of England ke London,
sedangkan relatif sedikit pekerja bergerak dari London ke dua wilayah
tersebut. {[}TXT baris 855-860, PDF hlm. 1162{]}

\textbf{Laporan penulis.} Model 4 menunjukkan beberapa perubahan waktu.
Hubungan pusat kota -\textgreater{} suburb antarwilayah South East dan
East of England meningkat relatif terhadap referensi, dengan interaksi
waktu 0,125 dan p \textless{} 0,05. Untuk hubungan dengan London,
Greater South East suburb -\textgreater{} London memiliki interaksi
-0,118 dan London -\textgreater{} Greater South East central cities
memiliki interaksi 0,290; London -\textgreater{} Greater South East
suburbs memiliki interaksi 0,121. Ketiga interaksi tersebut signifikan
pada tingkat yang ditandai dalam Tabel 3. Interaksi Greater South East
central cities -\textgreater{} London sebesar -0,057 tidak signifikan.
{[}TXT baris 830-841, PDF hlm. 1162{]}

\textbf{Interpretasi penulis.} Penulis merangkum hasil Model 4 sebagai
bukti desentralisasi Greater South East pada tingkat inter-perkotaan.
Namun, hubungan antarkawasan South East dan East of England tidak
memperoleh kekuatan relatif terhadap hubungan di dalam wilayah
perkotaan, kecuali hubungan pusat kota -\textgreater{} suburb. Karena
itu, bukti perkembangan jaringan antarkawasan selama dua puluh tahun
dinilai masih kurang. {[}TXT baris 855-865, PDF hlm. 1162{]}

\textbf{Laporan penulis (diagnostik model).} Tabel 2 melaporkan
parameter \emph{over-dispersion} 0,298 dan 0,294 untuk Model 1 dan Model
2, keduanya signifikan pada p \textless{} 0,01. Tabel 3 melaporkan
statistik Vuong 13,17 dan 12,77 untuk Model 3 dan Model 4, keduanya
signifikan pada p \textless{} 0,01; catatan metode menyatakan statistik
tersebut mendukung model zero-inflated dibandingkan pasangan
non-zero-inflated. Jumlah observasi yang dilaporkan adalah 2.052 untuk
masing-masing model intra-perkotaan dan 60.681 untuk masing-masing model
inter-perkotaan. {[}TXT baris 760-768, 842-853, dan 1078-1082, PDF hlm.
1161-1162 dan catatan 10{]}

\subsection{Findings}\label{findings-46}

\textbf{Interpretasi penulis pada skala intra-perkotaan.} Greater South
East tidak sepenuhnya monosentris maupun polisentris. Hubungan di dalam
pusat kota masih kuat, tetapi hubungan di dalam suburb tidak lebih lemah
secara signifikan. Hal ini ditafsirkan sebagai tanda bahwa banyak
distrik yang diklasifikasikan sebagai suburb telah mandiri secara
ekonomi dan tidak semata-mata berfungsi sebagai hinterland pusat. Dengan
demikian, wilayah perkotaan menunjukkan campuran model kota monosentris
dan polisentris serta perkembangan menuju bentuk yang lebih polisentris.
{[}Interpretasi penulis; TXT baris 671-721 dan 773-789, PDF hlm.
1160-1161{]}

\textbf{Interpretasi penulis pada skala inter-perkotaan.} Greater South
East belum merupakan jaringan perkotaan yang terintegrasi.
Ketergantungan di dalam wilayah perkotaan masih lebih kuat daripada
sebagian besar ketergantungan antarkawasan. Hubungan dengan London
bersifat tidak seimbang dan secara umum mendukung pembacaan Greater
South East sebagai sistem yang masih monosentris atau hub-and-spoke,
walaupun terdapat tanda desentralisasi. {[}Interpretasi penulis; TXT
baris 790-865 dan 893-915, PDF hlm. 1161-1163{]}

\textbf{Sintesis lintas skala.} Temuan utama bukan satu klasifikasi
tunggal untuk seluruh kawasan. Perubahan lokal dapat berlangsung menuju
kemandirian suburb sementara hubungan regional belum menjadi jaringan
yang saling silang. Desentralisasi arus menuju atau dari London juga
tidak sama dengan integrasi penuh, karena syarat integrasi menuntut
hilangnya hierarki dan hubungan yang lebih merata antarjenis distrik.
{[}Inferensi pengolahan berdasarkan TXT baris 669-718, 855-865, dan
893-915, PDF hlm. 1160-1163{]}

\subsection{Conclusions}\label{conclusions-46}

\textbf{Kesimpulan yang dilaporkan penulis.} Artikel menyimpulkan bahwa
Greater South East belum merupakan jaringan perkotaan yang sepenuhnya
terintegrasi. Kekuatan interdependensi di antara distrik suburb
menunjukkan bahwa simpul suburb semakin beroperasi secara independen
dari pusat kota dan mendukung perkembangan wilayah polisentris pada
skala intra-perkotaan. Analisis tren mengarah ke bentuk yang lebih
polisentris, tetapi ekstrapolasi sederhana dari dua puluh tahun
perubahan menghasilkan waktu sekitar 80 tahun lagi untuk mencapai
\emph{network-city} penuh pada skala tersebut jika tidak ada kejutan
eksternal. {[}Laporan penulis; TXT baris 880-905, PDF hlm. 1162-1163{]}

\textbf{Laporan penulis.} Pada skala inter-perkotaan, hasil umumnya
mendukung interpretasi monosentris yang relatif stabil, terutama karena
peran London dan lemahnya perkembangan jaringan antarkawasan. Ketiadaan
hierarki pada sebagian hubungan regional inter-perkotaan dipandang
membuka kemungkinan menuju polisentrisme, tetapi penulis menyatakan
bahwa hal itu memerlukan investasi yang terarah pada infrastruktur,
lokasi pekerjaan, dan perencanaan. {[}TXT baris 905-915, PDF hlm.
1163{]}

\textbf{Laporan penulis.} Penulis menutup dengan kualifikasi bahwa
komuter hanya satu cara untuk menilai struktur sistem perkotaan.
Kesimpulan tentang jaringan ekonomi yang lebih luas tidak boleh dibaca
seolah-olah telah dibuktikan oleh arus perjalanan kerja saja. {[}TXT
baris 916-936, PDF hlm. 1163{]}

\subsection{Contributions}\label{contributions-46}

\textbf{Kontribusi yang diklaim penulis.} Kontribusi utama artikel
adalah memberikan penilaian empiris terhadap perubahan sistem perkotaan
menggunakan data arus komuter, bukan hanya karakteristik simpul, dan
menganalisis perkembangan jaringan sepanjang waktu 1981-2001. Artikel
secara eksplisit menempatkan pendekatan tersebut sebagai respons
terhadap keterbatasan studi yang hanya menguji kondisi pada satu titik
waktu. {[}Laporan penulis; TXT baris 219-248 dan 233-248, PDF hlm.
1151{]}

\textbf{Laporan penulis.} Kontribusi substantifnya adalah memperlakukan
Greater South East sebagai kasus yang harus diuji, bukan sebagai contoh
jaringan terintegrasi yang sudah terbukti. Dengan membedakan skala
intra-perkotaan dan inter-perkotaan, artikel menghasilkan kesimpulan
yang tidak seragam: terdapat perkembangan lokal dan desentralisasi
regional, tetapi belum ada integrasi penuh. {[}TXT baris 233-248,
510-521, dan 893-915, PDF hlm. 1151, 1156, dan 1163{]}

\textbf{Inferensi pengolahan.} Perluasan model gravitasi dengan tiga
belas rezim interdependensi, jarak jalan, \emph{contiguity}, efek tetap
asal-tujuan-waktu, serta pemilihan model binomial negatif dan
zero-inflated Poisson memungkinkan perbandingan arus antarjenis hubungan
setelah faktor ukuran dan kedekatan diperhitungkan. Ini adalah
kontribusi analitis yang dapat dikaitkan langsung dengan desain artikel,
tetapi tidak boleh dibaca sebagai bukti bahwa metode tersebut
membuktikan hubungan kausal. {[}TXT baris 553-623 dan 629-718, PDF hlm.
1158-1160{]}

\textbf{Inferensi pengolahan.} Klaim kontribusi berupa perangkat lunak,
dataset terbuka, replikasi kode, atau evaluasi langsung keberhasilan
kebijakan: \textbf{Tidak disebutkan dalam file}. {[}TXT keseluruhan{]}

\subsection{Research Gap}\label{research-gap-46}

\textbf{Laporan penulis (kesenjangan yang ditangani).} Artikel
menargetkan tiga kekurangan yang dinyatakan penulis: sedikitnya
pengujian empiris kuantitatif yang ketat atas model jaringan perkotaan;
dominasi ukuran berbasis simpul yang hanya menjadi proksi interaksi; dan
keterbatasan studi berbasis arus yang membandingkan perubahan lintas
waktu. Penggunaan arus komuter, pengendalian massa dan jarak, serta tiga
tahun sensus dirancang untuk menutup sebagian kesenjangan tersebut.
{[}TXT baris 219-248 dan 485-521, PDF hlm. 1151 dan 1156{]}

\textbf{Laporan penulis (kesenjangan yang masih tersisa).} Pertama,
komuter perlu dilengkapi data interaksi ekonomi lain, seperti hubungan
pembeli-pemasok dan kolaborasi inovasi. Kedua, perjalanan yang tidak
terjadi setiap hari, termasuk perjalanan rekreasi dan bisnis, serta
perdagangan antarkota belum dimasukkan. Ketiga, komuter mingguan jarak
jauh atau \emph{super-commutes} dapat mengganggu gambaran ``ruang
perkotaan harian'' karena pekerja mungkin memiliki dua tempat tinggal.
{[}TXT baris 874-900 dan 916-934, PDF hlm. 1162-1163{]}

\textbf{Laporan penulis.} Keempat, hubungan waktu non-linear perlu
diteliti lebih hati-hati karena artikel hanya menguji tren linear.
Kelima, pembacaan masa depan tidak hanya bergantung pada tren arus yang
telah diamati, tetapi juga pada kejutan eksternal seperti pembukaan
infrastruktur, penciptaan atau relokasi pekerjaan, dan pembangunan
fasilitas perumahan. {[}TXT baris 899-905 dan catatan 12 pada baris
1105-1107, PDF hlm. 1163 dan 1168{]}

\textbf{Inferensi pengolahan.} Pembaruan kesenjangan berdasarkan
literatur setelah artikel ini: \textbf{Tidak disebutkan dalam file}.
Review ini tidak menyatakan apakah kesenjangan tersebut telah ditutup
oleh penelitian lain. {[}TXT keseluruhan{]}

\subsection{Limitations}\label{limitations-46}

\textbf{Batasan yang dinyatakan penulis.} Data komuter hanya menangkap
satu jenis pergerakan manusia dan bukan ukuran sempurna untuk interaksi
ekonomi. Kelompok masyarakat memiliki kesediaan bepergian yang berbeda,
dan banyak orang tetap memilih tinggal relatif dekat dengan tempat kerja
utama. Data hubungan antarfirma atau kolaborasi inovasi disebut sulit
diperoleh atau terlalu lemah untuk menghasilkan inferensi statistik yang
kuat. {[}Laporan penulis; TXT baris 916-934, PDF hlm. 1163{]}

\textbf{Laporan penulis.} Data yang tersedia terutama berasal dari
pertanyaan tentang perilaku komuter harian. \emph{Super-commutes}
mingguan jarak jauh dapat tidak tercatat dengan baik atau mengganggu
interpretasi karena pekerja dapat mempunyai dua tempat tinggal. Penulis
mengutip perkiraan bahwa jumlah orang yang melakukan jenis komuter
tersebut di Inggris sedikit di atas 1\% penduduk bekerja. {[}TXT baris
874-900, PDF hlm. 1162-1163{]}

\textbf{Inferensi pengolahan.} Perincian perhitungan perkiraan tersebut
untuk Greater South East: \textbf{Tidak disebutkan dalam file}. {[}TXT
baris 874-900, PDF hlm. 1162-1163{]}

\textbf{Inferensi pengolahan.} Hanya tiga titik waktu yang dianalisis
dan perubahan dimodelkan sebagai tren linear. Ekstrapolasi 80 tahun
sangat bergantung pada asumsi bahwa pola relatif dua puluh tahun tetap
berlanjut; penulis sendiri memperingatkan perlunya memasukkan kejutan
eksternal dan menguji hubungan non-linear. {[}TXT baris 723-724,
899-905, dan 1105-1107, PDF hlm. 1160, 1163, dan 1168{]}

\textbf{Laporan penulis.} Jarak jalan tahun 2005 digunakan untuk ketiga
periode, berdasarkan alasan historis bahwa jarak aktual sebagian besar
tidak berubah. {[}TXT baris 580-595 dan 1039-1046, PDF hlm. 1158-1159
dan 1167{]}

\textbf{Inferensi pengolahan.} Uji sensitivitas terhadap jaringan jalan
historis, perubahan waktu tempuh, atau perubahan aksesibilitas:
\textbf{Tidak disebutkan dalam file}. Ini merupakan keterbatasan
pengukuran yang diturunkan dari spesifikasi data, bukan kritik yang
dinyatakan sebagai temuan terpisah oleh penulis. {[}TXT baris 580-595
dan 1039-1046, PDF hlm. 1158-1159 dan 1167{]}

\textbf{Laporan penulis.} Definisi wilayah perkotaan bersifat
operasional. Penulis mengakui adanya perdebatan mengenai dimensi Greater
South East dan memilih tiga wilayah administratif karena alasan
kesesuaian dengan realitas kebijakan. {[}TXT baris 943-977 dan
1058-1074, PDF hlm. 1164-1165 dan catatan 2-5{]}

\textbf{Inferensi pengolahan.} Pengelompokan NUTS-III dan TTWA,
pemilihan distrik dengan kota terbesar sebagai inti, serta perlakuan
khusus untuk London dapat memengaruhi kategori intra- dan
inter-perkotaan. {[}TXT baris 943-977 dan 1058-1074, PDF hlm. 1164-1165
dan catatan 2-5{]}

\textbf{Inferensi pengolahan.} Analisis sensitivitas terhadap batas
alternatif, pengujian pada data tingkat individu, pemeriksaan kausal,
rincian data hilang, atau kode replikasi: \textbf{Tidak disebutkan dalam
file}. Karena itu, validitas kesimpulan harus dipahami sebagai validitas
bersyarat pada geografi, arus komuter, periode, dan kriteria integrasi
yang dipakai artikel. {[}TXT baris 916-934, 943-977, dan 1058-1074{]}

\subsection{Challenges}\label{challenges-46}

\textbf{Laporan penulis (tantangan pengukuran ruang).} Perubahan batas
distrik selama dua puluh tahun mengharuskan penulis menggunakan
\emph{CIDS 1991/2001 Common Geography} dan menyesuaikan data 1981.
Penentuan 27 wilayah perkotaan juga tidak sepenuhnya mekanis karena
NUTS-III kadang berdiri sebagai satu unit sendiri, sehingga TTWA
digunakan dalam beberapa kasus. {[}TXT baris 506-521 dan 948-977, PDF
hlm. 1156 dan 1164-1165{]}

\textbf{Laporan penulis (tantangan membedakan ukuran dari hubungan).}
London memiliki populasi besar, lapangan kerja banyak, kawasan yang
relatif padat, dan jarak jalan yang lebih pendek dengan beberapa
distrik. Penulis menilai bahwa visualisasi sederhana dapat salah
menganggap arus besar sebagai bukti integrasi. Model gravitasi dipakai
untuk memisahkan pengaruh ukuran, jarak, dan tipe hubungan. {[}TXT baris
471-505, PDF hlm. 1156-1157{]}

\textbf{Laporan penulis (tantangan distribusi data).} Arus komuter
berupa data hitungan dan menghadirkan \emph{over-dispersion} serta
banyak nilai nol. Penulis harus memilih binomial negatif untuk skala
intra-perkotaan dan zero-inflated Poisson untuk skala inter-perkotaan,
lalu menggunakan pengujian \emph{over-dispersion} dan statistik Vuong.
{[}TXT baris 597-623, 734-768, dan 804-853, PDF hlm. 1158-1159 dan
1161-1162{]}

\textbf{Laporan penulis (tantangan konseptual dan kebijakan).} Greater
South East memiliki definisi yang diperdebatkan dan berbagai lembaga
mengusulkan pusat serta koridor pertumbuhan yang tidak selalu sama.
Selain itu, literatur membedakan polisentrisme morfologis, jaringan
nodal, dan jaringan yang sepenuhnya terintegrasi. Artikel harus
menerjemahkan perbedaan tersebut ke dalam tiga belas kategori arus dan
dua skala analisis. {[}TXT baris 434-480, 624-718, dan 1058-1074, PDF
hlm. 1155-1156 dan catatan 2-5{]}

\textbf{Laporan penulis.} Data pembeli-pemasok, kolaborasi inovasi,
perjalanan mingguan, dan interaksi yang lebih jarang sulit diperoleh
atau tidak cocok dengan data yang tersedia. {[}TXT baris 928-936, PDF
hlm. 1163{]}

\textbf{Inferensi pengolahan.} Uraian teknis lebih jauh tentang proses
memperoleh, membersihkan, atau menggabungkan data tersebut:
\textbf{Tidak disebutkan dalam file}. {[}TXT baris 928-936, PDF hlm.
1163{]}

\subsection{Practical Implications}\label{practical-implications-46}

\textbf{Laporan penulis (implikasi transportasi).} Karena bukti
perkembangan jaringan lebih kuat pada skala intra-perkotaan daripada
inter-perkotaan, penulis menyarankan agar perencanaan transportasi
memberi perhatian khusus pada jaringan di dalam kawasan perkotaan.
Prioritas yang disebutkan adalah peningkatan jalan sekunder,
transportasi publik, dan rute sepeda dalam struktur perkotaan yang sudah
ada, bukan secara otomatis memprioritaskan jalan berkecepatan tinggi
yang menghubungkan wilayah perkotaan. {[}TXT baris 900-910, PDF hlm.
1163{]}

\textbf{Laporan penulis (implikasi untuk strategi regional).} Jika
tujuan kebijakan adalah mendorong polisentrisme inter-perkotaan, penulis
menyebut perlunya investasi yang terarah pada infrastruktur, lokasi
pekerjaan, dan perencanaan. Namun, hasil artikel tidak mendukung asumsi
bahwa kawasan tersebut sudah berfungsi sebagai jaringan yang
terintegrasi atau bahwa kerja sama antarkota dengan sendirinya
menghasilkan pertumbuhan di setiap kota. {[}TXT baris 395-404 dan
905-915, PDF hlm. 1155 dan 1163{]}

\textbf{Laporan penulis (implikasi tata ruang).} Jika persaingan
antarlokasi tetap menghambat pendekatan jaringan, perencanaan lokasi
perumahan dan kegiatan usaha disarankan terutama melayani permintaan
lokal dan meningkatkan efisiensi sistem perkotaan harian. Penulis juga
menilai bahwa perencanaan perkembangan penduduk menjadi lebih penting
karena lokasi pekerjaan semakin endogen terhadap lokasi penduduk, atau
dalam ungkapan artikel, ``jobs follow people'' lebih daripada ``people
follow jobs''. {[}TXT baris 916-930, PDF hlm. 1163{]}

\textbf{Inferensi pengolahan.} Implikasi tersebut bersifat rekomendasi
berbasis pola arus, bukan hasil evaluasi biaya-manfaat atau evaluasi
dampak kebijakan. Urutan prioritas investasi, besaran biaya, target
implementasi, atau indikator evaluasi kebijakan: \textbf{Tidak
disebutkan dalam file}. {[}TXT baris 900-930{]}

\subsection{Applications}\label{applications-46}

\textbf{Inferensi pengolahan (aplikasi analitis yang diturunkan dari
studi).} Kerangka ini dapat digunakan sebagai alat diagnosis untuk
membedakan apakah arus komuter terutama bertahan di dalam wilayah
perkotaan, menghubungkan pusat-pusat antarkawasan, atau meluas ke
hubungan pusat-suburb dan suburb-suburb. Model juga dapat dipakai untuk
memeriksa klaim bahwa suatu kawasan merupakan jaringan perkotaan, dengan
syarat data arus, geografi, jarak, dan kategori hubungan tersedia.
{[}TXT baris 629-718 dan 790-865, PDF hlm. 1159-1162{]}

\textbf{Inferensi pengolahan (aplikasi perencanaan).} Hasilnya dapat
diterapkan secara terbatas untuk mengarahkan perhatian pada konektivitas
intra-perkotaan, menguji apakah koridor atau pusat pertumbuhan yang
diusulkan benar-benar memiliki hubungan fungsional, dan membandingkan
hubungan London dengan hubungan antarkawasan lain. Ini merupakan
penggunaan analitis yang diturunkan dari metode dan implikasi penulis,
bukan penerapan yang diuji sebagai bagian dari artikel. {[}TXT baris
471-480, 485-505, dan 900-930, PDF hlm. 1156-1157 dan 1163{]}

\textbf{Laporan penulis.} Pendekatan yang sama dapat dipadukan dengan
data pembeli-pemasok, kolaborasi inovasi, perjalanan rekreasi,
perjalanan bisnis, atau perdagangan antarkota untuk memperoleh gambaran
fungsional yang lebih luas. Saran tersebut muncul dalam diskusi artikel.
{[}TXT baris 916-934, PDF hlm. 1163{]}

\textbf{Inferensi pengolahan.} Hasil empiris dari data tambahan itu:
\textbf{Tidak disebutkan dalam file}. {[}TXT baris 916-934, PDF hlm.
1163{]}

\textbf{Inferensi pengolahan.} Contoh implementasi nyata, evaluasi
sebelum-sesudah kebijakan, atau bukti bahwa rekomendasi transportasi
telah diterapkan dan menghasilkan dampak tertentu: \textbf{Tidak
disebutkan dalam file}. {[}TXT keseluruhan{]}

\subsection{Future Research
(Optional)}\label{future-research-optional-46}

\textbf{Laporan penulis.} Agenda masa depan yang dinyatakan atau
langsung disarankan dalam TXT meliputi:

\begin{itemize}
\tightlist
\item
  Menguji hubungan waktu non-linear, karena artikel hanya menggunakan
  tren linear antara 1981 dan 2001. {[}Catatan penulis; TXT baris
  1105-1107, PDF hlm. 1168{]}
\item
  Menggabungkan data interaksi pembeli-pemasok dan kolaborasi inovasi
  untuk melengkapi indikator komuter. {[}Laporan penulis; TXT baris
  928-934, PDF hlm. 1163{]}
\item
  Meneliti perjalanan yang lebih jarang, terutama perjalanan rekreasi
  dan bisnis, serta hubungan fungsional lain seperti perdagangan
  antarkota. {[}Laporan penulis; TXT baris 893-900, PDF hlm.
  1162-1163{]}
\item
  Memperhitungkan \emph{super-commutes} mingguan dan kemungkinan ``ruang
  perkotaan mingguan'' atau ``ruang perkotaan bulanan'', bukan hanya
  ruang perkotaan harian. {[}Laporan penulis; TXT baris 874-900, PDF
  hlm. 1162-1163{]}
\item
  Menguji perubahan sistem dengan memperhitungkan kejutan eksternal
  seperti pembukaan infrastruktur, penciptaan dan relokasi pekerjaan,
  serta pembangunan fasilitas perumahan, alih-alih mengekstrapolasi
  perbedaan tren yang telah diamati. {[}Inferensi pengolahan; agenda
  tersirat, bukan agenda eksplisit, berdasarkan TXT baris 899-905, PDF
  hlm. 1163{]}
\end{itemize}

\textbf{Inferensi pengolahan.} Agenda yang lebih spesifik mengenai
periode setelah 2001, desain sampel individual, atau model kausal:
\textbf{Tidak disebutkan dalam file}. {[}TXT keseluruhan{]}

\subsection{Provenance Notes}\label{provenance-notes-46}

\begin{itemize}
\tightlist
\item
  Review ini hanya menggunakan file TXT B03 yang diminta. Daftar
  referensi di dalam artikel dibaca sebagai bagian dari TXT, tetapi
  sumber-sumber yang dirujuk tidak dibuka atau digunakan sebagai bukti
  tambahan.
\item
  TXT menyatakan bahwa teks diekstrak dari PDF dengan
  \texttt{pdftotext\ -layout} pada 31 Agustus 2026. Metadata PDF
  menyebut 23 halaman, ukuran A4, dan judul metadata \texttt{untitled};
  judul artikel dan identitas bibliografis di atas diambil dari judul,
  byline, dan blok sitasi yang tercetak di dalam teks artikel. {[}TXT
  baris 1-25 dan 35-64{]}
\item
  Locator dalam review menggunakan \texttt{TXT\ baris\ x-y} sebagai
  penunjuk utama. \texttt{PDF\ hlm.} merujuk pada nomor halaman artikel
  yang tercetak di dalam hasil ekstraksi, bukan indeks halaman PDF
  keseluruhan.
\item
  Ekstraksi mempertahankan header, footer, nomor halaman, dan beberapa
  artefak tata letak. Simbol panah pada kategori hubungan tampil sebagai
  \texttt{!} di sejumlah baris TXT; review menuliskannya sebagai
  \texttt{-\textgreater{}} hanya untuk menjelaskan arah yang sudah
  dijelaskan oleh teks.
\item
  Detail visual gambar di luar caption dan uraian naratif: \textbf{Tidak
  disebutkan dalam file}. Review tidak menilai bentuk grafis, peta, atau
  ukuran panah secara independen. {[}TXT baris 217-218, 326-329,
  550-552, dan 976-977{]}
\item
  Ada ketidakselarasan angka kecil di dalam bahan: Tabel 2 menampilkan
  koefisien \emph{contiguity} Model 1 sebesar 1,343, sedangkan catatan
  11 menyebut 1,349 dan menggunakan nilai itu untuk memperoleh faktor
  3,85. Review mempertahankan interpretasi tekstual 3,85 kali dan tidak
  merekonstruksi koefisien yang berbeda. {[}TXT baris 740-742 dan
  1085-1104, PDF hlm. 1161 dan catatan 11 pada hlm. 1167{]}
\item
  Ada ketegangan internal yang dipertahankan: hasil Model 3 menyatakan
  terdapat hierarki antartipe interdependensi di antara wilayah South
  East dan East of England, sedangkan diskusi akhir menyebut absennya
  hierarki regional membuka kemungkinan menuju polisentrisme. Review
  tidak menyelesaikan perbedaan formulasi ini di luar TXT. {[}TXT baris
  778-789 dan 905-915, PDF hlm. 1161 dan 1163{]}
\item
  Ada pula ketidakselarasan tahun pada Appendix C: uraian utama
  konsisten menyebut 1981, 1991, dan 2001, sedangkan satu kalimat
  menyebut periode \texttt{1980,\ 1991,\ and\ 2001}. Review tidak
  mengoreksi kalimat tersebut secara spekulatif dan menggunakan tiga
  tahun yang berulang pada abstrak, metode, tabel, dan kesimpulan.
  {[}TXT baris 506-507, 734-768, 804-848, dan 1042-1046{]}
\item
  Data mentah, kode replikasi, rincian data hilang, jumlah komuter
  total, jumlah pasangan distrik unik, bobot survei, serta validasi
  terhadap batas alternatif atau jaringan jalan historis: \textbf{Tidak
  disebutkan dalam file}.
\item
  Label ``Laporan penulis'' dipakai untuk isi yang dinyatakan atau
  dilaporkan artikel; ``Interpretasi penulis'' untuk cara artikel
  mengartikan hasil; dan ``Inferensi pengolahan'' untuk simpulan
  analitis yang ditarik review secara terbatas dari teks. Tidak ada
  klaim dari luar TXT yang ditambahkan.
\end{itemize}

\section{Review B04}\label{review-b04}

\textbf{Identitas sumber}

\begin{itemize}
\tightlist
\item
  Kode kajian: B04.
\item
  Judul: \emph{On the Economic Foundation of the Urban Network Paradigm:
  Spatial Integration, Functional Integration and Economic
  Complementarities within the Dutch Randstad}.
\item
  Penulis: Frank van Oort, Martijn Burger, dan Otto Raspe.
\item
  Publikasi: \emph{Urban Studies}, volume 47, nomor 4, halaman 725-748,
  April 2010.
\item
  DOI: 10.1177/0042098009352362.
\item
  Riwayat naskah yang tercantum: pertama diterima Agustus 2007 dan dalam
  bentuk final Februari 2009.
\item
  Status peer-review: Tidak disebutkan dalam file; TXT hanya menunjukkan
  publikasi ini sebagai artikel di \emph{Urban Studies} (TXT baris 26,
  68-70).
\item
  Basis akses: seluruh teks hasil ekstraksi tersedia dalam file
  \texttt{source/journal/B04-van-oort-burger-raspe-2010-on-the-economic-foundation-of-the-urban-network-paradigm-spatial-integration-functional-integration-and-economic-complementarities-within-the-dutch-randstad-10.1177-0042098009352362.txt}.
  Metadata menyatakan sumber PDF 24 halaman, diekstrak pada 2026-08-31
  dengan \texttt{pdftotext\ -layout} (TXT, baris 1-24).
\item
  Penanda epistemik: \textbf{Laporan sumber} berarti isi yang dinyatakan
  atau ditampilkan oleh artikel; \textbf{interpretasi penulis} berarti
  penjelasan atau kesimpulan Van Oort, Burger, dan Raspe;
  \textbf{inferensi pengolahan} berarti pembacaan analitis review ini
  yang tidak dinyatakan eksplisit oleh penulis.
\end{itemize}

\subsection{Summarized Abstract}\label{summarized-abstract-47}

\textbf{Laporan sumber:} Artikel menguji apakah integrasi spasial,
integrasi fungsional, dan komplementaritas ekonomi benar-benar hadir
dalam jaringan ekonomi Dutch Randstad. Data yang digunakan adalah relasi
antarperusahaan di Randstad. Hasil utamanya menunjukkan hierarki yang
jelas dalam berbagai jenis ketergantungan spasial, dengan model central
place tetap dominan. Tidak ditemukan bukti untuk integrasi fungsional
antarmunicipality secara umum. Artikel menyimpulkan bahwa Randstad pada
saat penelitian belum berfungsi sebagai wilayah yang terintegrasi secara
spasial dan fungsional; kebijakan ekonomi spasial lebih baik diarahkan
pada wilayah-wilayah yang lebih kecil di dalam Randstad. Temuan ini juga
mempertanyakan penerapan konsep urban network secara umum karena
Randstad lazim dipandang sebagai contoh utama sistem urban polisentris
yang berhasil (hlm. 725, TXT baris 41-52).

\textbf{Interpretasi penulis:} Keberadaan banyak kota dan hubungan
ekonomi antarkota tidak cukup untuk membuktikan bahwa seluruh Randstad
merupakan satu jaringan urban dengan subekonomi yang saling melengkapi.
Komplementaritas harus terbentuk dari kemunculan bersama integrasi
spasial dan integrasi fungsional, bukan dari interaksi saja (hlm.
731-732, TXT baris 328-355).

\textbf{Inferensi pengolahan:} Ringkasan abstrak mendukung pembacaan
bahwa klaim artikel bersifat kondisional pada definisi operasional dan
data relasi antarperusahaan tahun 2005. Ia tidak membuktikan bahwa
Randstad tidak pernah atau tidak mungkin menjadi jaringan urban pada
periode lain, sektor lain, atau definisi integrasi yang lebih longgar
(berdasarkan hlm. 725 dan 732-744, TXT baris 41-52, 391-415, dan
922-967).

\subsection{Problem Statement}\label{problem-statement-47}

\textbf{Laporan sumber:} Randstad adalah konurbasi di Belanda bagian
barat yang terdiri atas Amsterdam, Rotterdam, The Hague, Utrecht, dan
banyak kota yang lebih kecil. Dalam perencanaan dan kebijakan ekonomi,
nama Randstad sering diperlakukan sebagai satu kawasan urban dan ekonomi
yang terintegrasi. Literatur serta dokumen kebijakan mengaitkannya
dengan spesialisasi yang saling bergantung, daya saing, dan potensi
pertumbuhan ekonomi (hlm. 725-726, TXT baris 55-101).

\textbf{Laporan sumber:} Artikel menyoroti bahwa pemberian nama kepada
kumpulan kota tidak otomatis menjadikannya satu kota yang terintegrasi
secara spasial dan fungsional. Agar kota-kota dalam urban network
bersifat komplementer, kota-kota tersebut harus memiliki spesialisasi
industri yang berbeda sekaligus menunjukkan interaksi spasial yang kuat.
Relasi antarperusahaan diposisikan sebagai dasar penting pembentukan dan
evolusi jaringan urban (hlm. 726-727, TXT baris 102-123).

\textbf{Interpretasi penulis:} Masalah empirisnya adalah apakah citra
Randstad sebagai model jaringan kota didukung oleh relasi ekonomi
antarperusahaan yang aktual. Konsep spatial network juga berisiko hanya
menjadi nama baru untuk interaksi biasa jika pertukaran masih terutama
dijelaskan oleh ukuran node, jarak, dan hubungan hierarkis yang sudah
dijelaskan oleh paradigma spatial interaction tradisional (hlm. 727, TXT
baris 124-145).

\textbf{Inferensi pengolahan:} Problem statement mengandung dua
pengujian yang saling terkait: pertama, apakah pola aliran melampaui
pengaruh massa dan jarak; kedua, apakah perbedaan spesialisasi
benar-benar berkaitan dengan aliran antarwilayah. Karena itu, label
polisentris atau urban network tidak diperlakukan sebagai bukti awal,
melainkan sebagai klaim yang harus diuji (berdasarkan hlm. 727 dan
731-732, TXT baris 124-145 dan 328-355).

\subsection{Objectives}\label{objectives-47}

\textbf{Laporan sumber:} Artikel merumuskan tiga pertanyaan penelitian:

\begin{enumerate}
\def\labelenumi{\arabic{enumi}.}
\tightlist
\item
  Sejauh mana Randstad beroperasi sebagai entitas yang terintegrasi
  secara spasial?
\item
  Sejauh mana Randstad beroperasi sebagai entitas yang terintegrasi
  secara fungsional?
\item
  Apakah komplementaritas urban antarmunicipality, yang didefinisikan
  sebagai kemunculan bersama integrasi spasial dan fungsional, terdapat
  di Randstad?
\end{enumerate}

\textbf{Laporan sumber:} Tujuan operasionalnya adalah menggunakan data
relasi antarperusahaan dari 1.676 establishment di Randstad untuk
menguji kondisi integrasi dan komplementaritas pada tingkat
municipality. Relasi penjualan dan pembelian yang paling penting
dikumpulkan untuk barang, jasa, dan informasi, kemudian diagregasikan ke
municipality (hlm. 727 dan 732, TXT baris 147-158 dan 391-415).

\textbf{Interpretasi penulis:} Tujuan tersebut secara langsung menguji
dasar ekonomi paradigma urban network, bukan hanya morfologi kota,
jumlah penduduk, atau distribusi pekerjaan. Keberhasilan konsep jaringan
mensyaratkan integrasi arus dan pembagian fungsi ekonomi secara
bersamaan (hlm. 726-727 dan 731, TXT baris 110-145 dan 328-355).

\subsection{Literature Survey}\label{literature-survey-47}

\textbf{Inferensi pengolahan:} Teks menyajikan tinjauan konseptual dan
naratif, bukan tinjauan sistematis. TXT tidak memuat protokol
penelusuran, kriteria inklusi, jumlah korpus, atau prosedur penilaian
kualitas literatur. Pembahasan bergerak dari model monocentric city dan
central place, ke daily urban systems (DUS), polycentric urban region
(PUR), lalu ke fondasi ekonomi jaringan urban (hlm. 728-731, TXT baris
163-370).

\textbf{Laporan sumber:} Dalam model monocentric, kegiatan ekonomi
terkonsentrasi pada urban core, suburb terutama berfungsi sebagai tempat
tinggal, dan arus komuter maupun relasi ekonomi membentuk hubungan
hierarchical-nodal. Perkembangan transportasi, komunikasi,
suburbanisasi, dan perluasan wilayah mobilitas kemudian mendorong
terbentuknya metropolitan area atau DUS (hlm. 728-729, TXT baris
170-213).

\textbf{Laporan sumber:} PUR dipahami sebagai jaringan wilayah
metropolitan yang secara historis dan spasial terpisah tetapi berada
dalam kawasan yang lebih luas. Relasi dapat berupa periphery-core,
inter-core, criss-cross, dan intraurban. PUR tidak selalu merupakan
network-city pada semua skala; sistem dapat polisentris pada tingkat
antarurban tetapi tetap monosentris di dalam setiap wilayah
metropolitan, atau sebaliknya (hlm. 729, TXT baris 219-249).

\textbf{Laporan sumber:} Literatur ekonomi menempatkan mobilitas dan
fleksibilitas perusahaan, agglomeration economies, pembagian kerja
spasial, outsourcing, kerja sama vertikal, ICT, globalisasi, dan
perubahan menuju ekonomi jasa sebagai kemungkinan fondasi jaringan
urban. Artikel juga menyajikan perdebatan antara pandangan substitution
atau death of distance dan pandangan complementary yang menekankan
konsentrasi, ketergantungan historis, kontak tatap muka, serta relasi
kepercayaan (hlm. 730-731, TXT baris 271-337).

\textbf{Laporan sumber:} Tiga konsep kunci didefinisikan sebagai
berikut. Integrasi spasial terkait arus mobilitas dan komunikasi yang
menguat antarlokasi. Integrasi fungsional mensyaratkan diferensiasi
municipality dalam fungsi ekonomi sekaligus penggunaan spesialisasi satu
sama lain. Komplementaritas ekonomi muncul apabila perbedaan
spesialisasi tersebut bertepatan dengan interaksi ekonomi yang besar
(hlm. 731, TXT baris 328-355).

\textbf{Interpretasi penulis:} Literatur kebijakan dan akademik banyak
memakai ketiga konsep tersebut, tetapi hanya sedikit studi empiris yang
menilai seberapa baik model urban network menjelaskan sistem ekonomi
kontemporer. Bukti yang tersedia sering berangkat dari karakteristik
node dan menjelaskan interaksi melalui jarak serta ukuran node, bukan
melalui karakteristik arus antarperusahaan (hlm. 732, TXT baris
391-405).

\textbf{Inferensi pengolahan:} Tinjauan literatur berfungsi terutama
untuk menurunkan definisi dan kondisi pengujian, bukan untuk membangun
sintesis literatur yang independen dari kasus Randstad. Karena itu,
research gap yang diklaim artikel harus dibaca sebagai gap yang
dirumuskan dari cakupan referensi dan persoalan data dalam artikel ini,
bukan sebagai verifikasi menyeluruh terhadap seluruh literatur urban
network (berdasarkan hlm. 728-732, TXT baris 163-370 dan 391-405).

\subsection{Method Used}\label{method-used-47}

\textbf{Laporan sumber:} Analisis dilakukan dalam beberapa tahap.
Pertama, pola relasi divisualisasikan pada Figure 2. Kedua, Figure 3
membandingkan intensitas relasi aktual dengan intensitas yang diharapkan
setelah mempertimbangkan ukuran municipality asal dan tujuan serta
kedekatan fisik; perbandingan menggunakan likelihood ratio chi-squared
statistic. Ketiga, penulis menggunakan model spatial interaction untuk
menguji integrasi spasial dengan lebih akurat (hlm. 733-734, TXT baris
428-473).

\textbf{Laporan sumber:} Model utamanya adalah gravity model untuk data
hitungan. Intensitas interaksi \texttt{Iij} adalah jumlah linkage antar
atau di dalam municipality. Massa municipality didefinisikan sebagai
jumlah total relasi antarperusahaan yang tertanam dalam jaringan,
sedangkan jarak diukur sebagai jarak garis lurus antarkota. Untuk relasi
intra-municipality, jarak dihitung menggunakan dua pertiga radius dari
area lingkaran yang diasumsikan. Karena relasi diperlakukan sebagai
undirected, tidak ada pembedaan akhir antara municipality asal dan
tujuan; logaritma produk massa kedua municipality digunakan sebagai satu
variabel (hlm. 735-736, TXT baris 515-551).

\textbf{Laporan sumber:} Penulis memilih zero-inflated negative binomial
regression karena data menunjukkan overdispersion dan jumlah nol yang
berlebih. Komponen negative binomial memodelkan jumlah interaksi untuk
pasangan yang memiliki probabilitas interaksi tidak nol, sedangkan
komponen zero-inflated atau binary memodelkan peluang pasangan masuk ke
kelompok strictly zero. Likelihood ratio test digunakan untuk menilai
overdispersion dan Vuong statistic untuk membandingkan model
zero-inflated dengan model yang tidak zero-inflated (hlm. 736-737, TXT
baris 565-610).

\textbf{Laporan sumber:} Enam jenis ketergantungan spasial
diklasifikasikan: hubungan intra-municipality; core-periphery dan
criss-cross di dalam subwilayah urban yang sama; intercore;
core-periphery antara core dan suburb non-asalnya; serta criss-cross
antarsuburb dari subwilayah yang berbeda. Tiga kondisi integrasi spasial
diuji: hubungan intraurban tidak lebih kuat daripada hubungan antarkota;
hubungan dalam subwilayah tidak lebih kuat daripada hubungan lintas
subwilayah; dan tidak ada hierarki yang dapat diamati di antara jenis
hubungan tersebut (hlm. 737-738, TXT baris 601-632).

\textbf{Laporan sumber:} Integrasi fungsional diukur melalui location
quotient pada 27 basic sectors, dengan dua ukuran spesialisasi: jumlah
employment dan jumlah firms. Nilai location quotient di atas satu
berarti sektor relatif terwakili secara berlebih dan nilai di bawah satu
berarti relatif kurang terwakili. Perbedaan spesialisasi
antarmunicipality pada kedua ukuran tersebut digabungkan dengan bobot
yang sama menjadi indikator functional integration untuk setiap
interaksi ekonomi (hlm. 741-742, TXT baris 769-837).

\textbf{Laporan sumber:} Model 5 dan 6 memasukkan indikator integrasi
fungsional bersama massa serta jarak. Model 7 dan 8 memasukkan integrasi
spasial dan fungsional secara bersamaan. Tidak ada interaksi antara
integrasi fungsional dan jarak yang signifikan pada Model 6 dan 8,
sehingga term tersebut dihilangkan dari model final. Korelasi
antarvariabel penjelas dilaporkan tidak lebih tinggi dari 0,4 (catatan
7-8, TXT baris 990-1003).

\textbf{Inferensi pengolahan:} Desain ini adalah desain observasional
lintas bagian pada jaringan relasi tahun 2005. Gravity model menguji
pola asosiasi setelah kontrol massa dan jarak; ia tidak, berdasarkan
informasi yang tersedia dalam TXT, mengidentifikasi efek kausal
kebijakan, ICT, spesialisasi, atau perubahan kelembagaan (berdasarkan
hlm. 732-743, TXT baris 391-415 dan 515-914).

\subsection{Dataset}\label{dataset-47}

\textbf{Laporan sumber:} Dataset berasal dari survei tahun 2005 terhadap
establishment yang aktif di sektor manufaktur, wholesale, dan business
services di Dutch Randstad. Sampling frame-nya adalah LISA, yang
dijelaskan sebagai employment register seluruh establishment ekonomi di
Belanda. Sampel acak berstrata mempertimbangkan ukuran perusahaan,
sektor, dan wilayah (hlm. 732, TXT baris 391-405).

\textbf{Laporan sumber:} Kuesioner meminta setiap establishment
mengidentifikasi sumber dan tujuan dari sepuluh relasi penjualan dan
pembelian yang paling penting. Relasi tersebut berkaitan dengan
subcontracting serta pembelian dan penjualan produk, jasa, dan
pengetahuan. Jawaban kemudian diagregasikan menjadi arus
antarmunicipality (hlm. 727 dan 732, TXT baris 110-123 dan 410-420).

\textbf{Laporan sumber:} Sektor yang dipilih adalah sektor dasar yang
tidak secara langsung terikat pada konsumen sebagai penentu lokasi,
berbeda dari retail. Catatan artikel menyebut sektor non-basic antara
lain retail, pendidikan dasar dan sekolah, layanan publik lokal seperti
polisi, pemadam kebakaran, dan layanan kesehatan, serta pemerintah lokal
dan regional (hlm. 732 dan catatan 5, TXT baris 395-407 dan 980-984).

\textbf{Tidak disebutkan dalam file:} TXT tidak memberikan matriks
relasi mentah, daftar lengkap pasangan municipality, salinan kuesioner,
bobot setiap relasi, atau prosedur pembersihan dan pembobotan data di
luar penjelasan sampling dan agregasi tersebut.

\subsection{Population / Sample}\label{population-sample-47}

\textbf{Laporan sumber:} Table 1 memberikan populasi 41.398
establishment, sampel terpilih 20.301, dan 1.676 respons dengan response
rate total 8,3 persen. Rinciannya berdasarkan wilayah adalah Amsterdam:
populasi 19.045, sampel 7.035, respons 574, response rate 8,2 persen;
Rotterdam: 10.789, 5.668, 514, dan 9,1 persen; The Hague: 5.468, 3.655,
291, dan 8,0 persen; Utrecht: 6.096, 3.943, 297, dan 7,5 persen (Table
1, hlm. 732, TXT baris 372-382).

\textbf{Laporan sumber:} Berdasarkan sektor, manufaktur memiliki
populasi 5.322, sampel 5.307, respons 367, dan response rate 6,9 persen;
wholesale 10.991, 4.807, 376, dan 7,8 persen; business services 25.085,
10.186, 933, dan 9,2 persen. Totalnya tetap 41.398, 20.301, 1.676, dan
8,3 persen (Table 1, hlm. 732, TXT baris 383-387).

\textbf{Laporan sumber:} Unit observasi awal adalah establishment. Dalam
analisis, relasi diubah menjadi interaksi pada tingkat municipality
LAU-2. Artikel menyebut terdapat 69 municipality di Randstad dan
mengidentifikasi empat core city, yaitu Amsterdam, Rotterdam, The Hague,
dan Utrecht, beserta municipality suburban mereka (hlm. 732 dan 737, TXT
baris 391-405 dan 588-613).

\textbf{Laporan sumber:} Penulis menyatakan bahwa 1.676 respons, sekitar
8 persen, representatif terhadap stratifikasi menurut wilayah, ukuran,
dan sektor (hlm. 732, TXT baris 395-410).

\textbf{Inferensi pengolahan:} Pernyataan tersebut adalah klaim
representativitas terhadap skema stratifikasi yang digunakan, bukan
validasi bahwa seluruh relasi ekonomi Randstad teramati (berdasarkan
hlm. 732-734, TXT baris 395-415 dan 458-462).

\textbf{Tidak disebutkan dalam file:} Profil establishment yang tidak
merespons, penyesuaian nonresponse, bobot survei, tanggal pelaksanaan
yang lebih rinci dari tahun 2005, serta daftar pasti municipality tidak
diberikan dalam TXT.

\textbf{Catatan konsistensi sumber:} Caption Figure 1 menyebut ``67
municipalities in the survey'', sedangkan uraian desain dan indikator
fungsional menyebut 69 municipality. Nilai \texttt{N\ =\ 2415} pada
Model 1 dan 2 serta \texttt{N\ =\ 2346} pada Model 3 dan 4 juga
konsisten secara aritmetis dengan suatu konstruksi 69 municipality,
tetapi TXT tidak menjelaskan perbedaan dengan caption 67. Kecocokan
aritmetis tersebut adalah inferensi pengolahan dari angka yang
ditampilkan, bukan penjelasan yang diberikan penulis. Ketidaksesuaian
ini tidak dapat diselesaikan hanya dari TXT (Figure 1, hlm. 726, TXT
baris 77; hlm. 732, 741, dan Table 2-3, TXT baris 391-405, 655-658, dan
869-874).

\subsection{Dependent Variables}\label{dependent-variables-47}

\textbf{Laporan sumber:} Variabel dependen utama adalah intensitas
interaksi ekonomi \texttt{Iij}, yaitu frekuensi atau jumlah linkage
antarperusahaan antara municipality \texttt{i} dan \texttt{j}, termasuk
relasi yang seluruhnya berada di dalam satu municipality. Hubungan
diperlakukan sebagai undirected. Dalam model spatial integration,
variabel ini menjelaskan jumlah interaksi yang diamati untuk setiap
pasangan atau unit municipality (hlm. 735-736, TXT baris 526-551).

\textbf{Laporan sumber:} Karena menggunakan model zero-inflated negative
binomial, outcome memiliki dua komponen terkait: indikator probabilitas
suatu pasangan termasuk kelompok strictly zero dan jumlah interaksi pada
bagian yang memiliki probabilitas non-nol. Dengan demikian, nol tidak
diperlakukan semata-mata sebagai nilai rendah dari variabel hitungan
(hlm. 736-737, TXT baris 565-600).

\textbf{Laporan sumber:} Untuk pengujian integrasi fungsional, outcome
tetap berupa flow frequency antar-municipality. Indikator perbedaan
spesialisasi berperan sebagai variabel penjelas, bersama joint mass,
distance, dan dummy jenis ketergantungan spasial pada model gabungan
(Table 3, hlm. 743, TXT baris 869-914).

\textbf{Inferensi pengolahan:} Studi ini mengukur keterhubungan ekonomi
antarwilayah, bukan outcome langsung seperti pertumbuhan, produktivitas,
pendapatan, kesejahteraan, atau daya saing. Oleh sebab itu, hubungan
yang lebih banyak tidak dapat disamakan begitu saja dengan kinerja
ekonomi yang lebih baik (berdasarkan definisi outcome dan tujuan
pengujian pada hlm. 727, 735-743, TXT baris 147-158 dan 515-914).

\subsection{Results}\label{results-47}

\textbf{Laporan sumber:} Visualisasi jaringan menunjukkan banyak relasi
berada di dalam empat kota besar, tetapi terdapat pula arus antarpusat
dan relasi criss-cross di luar kota pusat. Dalam jumlah absolut, empat
kota besar menjadi pusat jaringan; Amsterdam disebut memiliki posisi
kunci dalam jaringan keseluruhan. Sekitar 50 persen relasi
antarperusahaan establishment yang berlokasi di Randstad melampaui batas
Randstad, sehingga Figure 2 hanya memperlihatkan sekitar setengah cerita
jaringan tersebut (Figure 2, hlm. 733-734, TXT baris 428-462).

\textbf{Laporan sumber:} Figure 3 menunjukkan bahwa relasi dalam
municipality yang sama dan relasi antara central city dengan hinterland
langsungnya umumnya lebih besar dari yang diproyeksikan setelah
memperhitungkan massa dan jarak. Sebaliknya, relasi antarempat central
city lebih kecil daripada yang diharapkan (Figure 3, hlm. 734, TXT baris
443-459).

\textbf{Laporan sumber:} Uji likelihood ratio untuk overdispersion dan
Vuong statistic mendukung zero-inflated negative binomial sebagai model
yang paling sesuai di Table 2. Pada komponen negative binomial Model 1,
joint mass berkoefisien positif 0,91 dan distance berkoefisien negatif
-0,85, keduanya signifikan pada p\textless0,01. Penulis menafsirkan
koefisien tersebut sebagai elastisitas: kenaikan 1 persen jarak
memprediksi penurunan 0,85 persen intensitas interaksi, sedangkan
kenaikan 1 persen joint mass memprediksi kenaikan 0,91 persen (Model 1,
Table 2, hlm. 738-740, TXT baris 655-693 dan 712-733).

\textbf{Laporan sumber:} Komponen zero-inflated menunjukkan bahwa
municipality dengan joint mass rendah lebih mungkin masuk kelompok tanpa
interaksi. Penulis menghubungkan bagian terbesar non-interaction dengan
banyak locality kecil yang tidak memiliki critical mass untuk
menghasilkan arus (Table 2, hlm. 739-740, TXT baris 697-704).

\textbf{Laporan sumber:} Model 2 menguji kondisi pertama. Setelah massa
dan jarak dikontrol, interaksi intra-municipality secara signifikan
lebih kuat daripada interaksi antarmunicipality dalam subwilayah yang
sama maupun antarmunicipality di luar subwilayah. Intensitas yang
diprediksi di dalam municipality lebih dari 75 persen lebih tinggi
daripada interaksi antarmunicipality dalam wilayah yang sama dan sekitar
110 persen lebih tinggi daripada interaksi antarmunicipality lintas
wilayah (Table 2, hlm. 740, TXT baris 734-754).

\textbf{Laporan sumber:} Model 3 menguji kondisi kedua dan menemukan
bahwa interaksi antarmunicipality di dalam subwilayah yang sama
diprediksi lebih dari 30 persen lebih tinggi daripada interaksi
antarmunicipality lintas subwilayah, dengan massa dan jarak tetap (Table
2, hlm. 740, TXT baris 713-729).

\textbf{Laporan sumber:} Model 4 menguji kondisi ketiga dengan
interregional criss-cross sebagai kategori acuan. Intercore dan
intraregional core-periphery adalah jenis hubungan relatif terkuat.
Interregional core-periphery dan intraregional criss-cross disebut
sebagai jenis yang paling lemah, sehingga terdapat hierarki dan bukan
kesetaraan antarjenis hubungan. Koefisien intercore adalah 0,52 dengan
p\textless0,05, intraregional core-periphery 0,46 dengan p\textless0,01,
sedangkan beberapa koefisien lain tidak signifikan (Table 2, hlm.
740-741, TXT baris 734-755 dan 759-767).

\textbf{Interpretasi penulis:} Karena ketiga kondisi integrasi spasial
tidak terpenuhi, Randstad belum berfungsi sebagai cluster yang
terintegrasi secara spasial. Relasi perusahaan bersifat regionalized di
dalam empat subwilayah dan interregionalized terutama di antara empat
kota besar (hlm. 741, TXT baris 759-767).

\textbf{Laporan sumber:} Pada integrasi fungsional, Model 5 tidak
menemukan pengaruh signifikan indikator functional integration terhadap
flow frequency. Model 6 menunjukkan interaksi negatif yang signifikan
antara indikator perbedaan spesialisasi dan massa: bukti lemah integrasi
fungsional muncul pada kelompok municipality yang relatif kecil. Dalam
Table 3, koefisien perbedaan spesialisasi adalah -0,10 pada Model 5 dan
1,50 pada Model 6; interaksinya dengan massa adalah -0,18 pada Model 6,
dengan tingkat signifikansi masing-masing tercantum pada tabel (Table 3,
hlm. 742-743, TXT baris 838-914).

\textbf{Laporan sumber:} Model 7 dan 8 memasukkan variabel spasial dan
fungsional sekaligus. Hubungan variabel spasial tetap serupa dengan
Model 2-4. Indikator functional integration hanya sedikit signifikan dan
kembali terutama pada municipality kecil. Pada Model 8, koefisien
perbedaan spesialisasi adalah 1,15 dengan p\textless0,10 dan
interaksinya dengan massa -0,15 dengan p\textless0,10; Model 7 tidak
menunjukkan signifikansi langsung pada indikator tersebut (Table 3, hlm.
743, TXT baris 875-914).

\subsection{Findings}\label{findings-47}

\textbf{Laporan sumber:} Temuan empiris utama adalah bahwa struktur
relasi antarperusahaan Randstad tidak memenuhi tiga kondisi yang dipakai
artikel untuk menyebut suatu sistem sebagai terintegrasi secara spasial.
Efek massa dan jarak tetap kuat, hubungan intra-municipality dan
intra-subwilayah lebih kuat, serta hubungan antarjenis menunjukkan
hierarki central place.

\textbf{Laporan sumber:} Tidak ada bukti jelas bahwa keseluruhan
municipality Randstad terintegrasi secara fungsional. Bukti yang relatif
lemah hanya muncul pada municipality terkecil, yang tidak mampu
menampung seluruh fungsi ekonomi di dalam batasnya sendiri dan karena
itu dapat berkomplementer dengan municipality terdekat yang memiliki
spesialisasi lain. Kota-kota besar dipandang telah memiliki sumber daya
dan fungsi khusus yang cukup dalam wilayah fungsionalnya masing-masing
(hlm. 741-744, TXT baris 820-837 dan 922-939).

\textbf{Interpretasi penulis:} Amsterdam, Rotterdam, The Hague, dan
Utrecht adalah mature nodes dalam wilayah ekonomi fungsionalnya sendiri.
Hubungan antarkota besar memang ada, tetapi pola tersebut tidak cukup
kuat untuk menggantikan struktur DUS dan central place dengan satu urban
network Randstad yang terpadu (hlm. 741-744, TXT baris 759-767 dan
922-939).

\textbf{Inferensi pengolahan:} Hasil mengarah pada polisentrisitas yang
tidak otomatis berarti integrasi jaringan. Empat node besar dapat hadir
bersamaan dengan ketergantungan yang tetap hierarkis dan
terregionalisasi. Ini membedakan keberadaan banyak pusat dari fungsi
kolektif satu sistem ekonomi (berdasarkan Figure 2, Table 2, dan
kesimpulan artikel, hlm. 733-744, TXT baris 428-462, 655-767, dan
922-939).

\subsection{Conclusions}\label{conclusions-47}

\textbf{Interpretasi penulis:} Jawaban terhadap pertanyaan pertama
adalah bahwa Randstad belum merupakan entitas yang terintegrasi secara
spasial menurut tiga kondisi yang diuji. Jawaban terhadap pertanyaan
kedua adalah tidak terdapat bukti jelas untuk integrasi fungsional
lintas municipality. Jawaban terhadap pertanyaan ketiga adalah
komplementaritas urban, yang didefinisikan sebagai kombinasi kedua
bentuk integrasi tersebut, saat ini tidak hadir di Randstad;
pengecualian yang ditemukan hanya bukti lemah pada municipality terkecil
(hlm. 742-744, TXT baris 922-939).

\textbf{Interpretasi penulis:} Penulis menyatakan bahwa DUS dari empat
kota terbesar lebih sesuai dengan realitas spasial-ekonomi daripada
memperlakukan seluruh Randstad sebagai satu sistem. Sekitar 50 persen
relasi berhubungan dengan wilayah nasional atau internasional di luar
Randstad, sehingga ketergantungan spasial dan fungsional melampaui DUS,
tetapi skala polycentric urban region justru tampak terlewati dalam
trajektori local-global yang dibahas (hlm. 744-745, TXT baris 940-967).

\textbf{Inferensi pengolahan:} Kesimpulan tidak mendukung klaim
universal bahwa semua urban network gagal. Kesimpulan yang dapat
dipertahankan dari TXT lebih sempit: pada sektor yang disurvei, wilayah
yang diteliti, relasi yang dinominasikan, dan observasi tahun 2005,
Randstad lebih menyerupai kumpulan DUS yang saling berhubungan daripada
satu jaringan urban yang terintegrasi (berdasarkan hlm. 732-745, TXT
baris 391-415 dan 922-984).

\subsection{Contributions}\label{contributions-47}

\textbf{Laporan sumber:} Kontribusi pertama adalah menyediakan pengujian
empiris terhadap dimensi ekonomi urban network menggunakan data
relasional antarperusahaan, ketika penelitian sebelumnya disebut
kekurangan data interfirm. Dataset berasal dari survei besar terhadap
1.676 establishment (hlm. 727 dan 744, TXT baris 147-158 dan 940-952).

\textbf{Laporan sumber:} Kontribusi kedua adalah memperkenalkan
pengujian berbasis karakteristik arus yang membedakan efek massa dan
jarak dari konteks spasial-fungsional hubungan. Tiga kondisi integrasi
spasial dan enam tipe ketergantungan memberi aturan pembanding yang
eksplisit, bukan hanya deskripsi peta jaringan (hlm. 732-738, TXT baris
391-405 dan 479-632).

\textbf{Laporan sumber tentang operasionalisasi:} Pengujian integrasi
fungsional menggabungkan ukuran spesialisasi employment dan jumlah firms
pada 27 sektor basic untuk membangun indikator perbedaan spesialisasi
yang digunakan dalam pengujian integrasi fungsional dan komplementaritas
(hlm. 741-742, TXT baris 774-837). TXT tidak menyebutnya sebagai
kontribusi yang terpisah.

\textbf{Laporan sumber:} Penulis menyatakan bahwa sistem pengujian
tersebut dirumuskan secara independen dari skala wilayah dan kota
sehingga dapat diterapkan pada konteks lain (hlm. 738, TXT baris
625-632).

\textbf{Inferensi pengolahan:} Artikel menyebut pengujian pembentukan
urban network di Greater South East, UK, sebagai contoh penggunaan
sejenis, tetapi TXT B04 tidak melaporkan hasil wilayah tersebut sebagai
bagian dari analisis artikel ini (hlm. 738 dan 744-745, TXT baris
625-632 dan 962-973).

\textbf{Inferensi pengolahan:} Nilai tambah utama artikel berada pada
operasionalisasi konsep yang biasanya normatif atau deskriptif menjadi
serangkaian kondisi yang dapat ditolak. Namun, kontribusi metodologis
itu tetap bergantung pada kategori hubungan, definisi massa, dan
indikator functional integration yang dipilih penulis (berdasarkan hlm.
732-742, TXT baris 391-405, 478-632, dan 769-837).

\subsection{Research Gap}\label{research-gap-47}

\textbf{Laporan sumber:} Gap awal yang hendak diisi adalah sedikitnya
penelitian empiris tentang kohesi ekonomi spasial dan fungsional
urban-networked regions. Penyebab yang disebut adalah kurangnya data
relasi antarperusahaan. Bukti empiris yang tersedia lebih sering memakai
karakteristik node seperti location quotients, rank-size relations,
sufficiency indices, dan employment-to-work ratios; pendekatan tersebut
belum secara langsung menguji pola arus interfirm (hlm. 732, TXT baris
391-405).

\textbf{Laporan sumber:} Artikel juga menyatakan bahwa dimensi ekonomi
urban network belum diuji secara sistematis dan belum ada penelitian
empiris sebelumnya yang menguji apakah citra Randstad sebagai contoh
polisentrisitas sesuai dengan relasi antarperusahaan masa kini (hlm.
744, TXT baris 940-948).

\textbf{Gap yang masih tersisa menurut sumber:} Penulis meminta
penelitian komparatif di wilayah lain untuk mengonfirmasi atau menantang
hasil Randstad. Mereka juga menekankan perlunya pengukuran berulang
karena relasi antarperusahaan dapat berkembang menuju pembentukan
jaringan antardaerah pada masa depan (hlm. 744-745, TXT baris 959-984).

\textbf{Inferensi pengolahan:} Gap yang tersisa bukan hanya replikasi
geografis, tetapi juga pengujian apakah hasil berubah jika cakupan
sektor, arah relasi, batas wilayah, atau periode pengamatan berubah.
Poin-poin ini adalah implikasi analitis dari desain yang dilaporkan,
bukan agenda eksplisit yang seluruhnya ditulis oleh penulis (berdasarkan
hlm. 732-745, TXT baris 391-415, 548-552, 769-837, dan 954-984).

\subsection{Limitations}\label{limitations-47}

\textbf{Keterbatasan yang dinyatakan atau tersirat langsung oleh
sumber:}

\begin{itemize}
\tightlist
\item
  Cakupan data dibatasi pada manufaktur, wholesale, dan business
  services yang dipilih sebagai basic sectors; sektor konsumen dan
  layanan publik lokal tidak masuk analisis. Karena itu, temuan tidak
  mewakili seluruh ekonomi Randstad (hlm. 732 dan catatan 5, TXT baris
  395-407 dan 980-984).
\item
  Data relasional dibentuk dari sepuluh relasi penjualan dan pembelian
  yang paling penting, bukan sensus seluruh relasi. Hubungan kemudian
  diagregasikan ke municipality, sehingga detail establishment dan
  hubungan antaraktor hilang (hlm. 727 dan 732, TXT baris 110-123 dan
  410-420).
\item
  Relasi diperlakukan sebagai undirected. Pembedaan supplier-customer
  atau arah arus tidak dipertahankan dalam model akhir, sehingga arti
  fungsional dari arah pertukaran tidak dapat dinilai dari hasil ini
  (hlm. 736, TXT baris 548-552).
\item
  Hanya hubungan yang diamati pada survei tahun 2005 yang digunakan.
  Penulis sendiri menyatakan bahwa relasi antarperusahaan perlu diukur
  berulang kali karena dapat berkembang dari waktu ke waktu (hlm.
  744-745, TXT baris 978-984).
\item
  Sekitar setengah relasi establishment di Randstad melampaui batas
  Randstad. Analisis regional karenanya tidak memuat keseluruhan
  jaringan eksternal yang diakui artikel sebagai bagian dari cerita
  (hlm. 734 dan 744, TXT baris 458-462 dan 954-957).
\item
  Uji sektor-spesifik untuk manufaktur, distribusi, dan business
  services dilakukan sebagai pemeriksaan robustness, tetapi tidak
  dieksplorasi lebih lanjut karena artikel berfokus pada interdependensi
  lintas sektor untuk mengukur integrasi fungsional (catatan 3, TXT
  baris 1004-1017).
\item
  Definisi integrasi spasial sengaja ketat: konteks spasial-fungsional
  tidak boleh memberi efek tambahan setelah massa dan jarak
  diperhitungkan (hlm. 734-738, TXT baris 478-513).
\end{itemize}

\textbf{Inferensi pengolahan:} Kesimpulan ``tidak terintegrasi'' karena
itu bergantung pada ambang konseptual yang dipilih (berdasarkan hlm.
734-738, TXT baris 478-513).

\textbf{Inferensi pengolahan:} Response rate 8,3 persen dan tidak
tersedianya informasi profil nonresponden membuat potensi bias
nonresponse tidak dapat dinilai dari TXT. Klaim representativitas
penulis terbatas pada stratifikasi wilayah, ukuran, dan sektor
(berdasarkan Table 1, hlm. 732, TXT baris 372-410).

\textbf{Inferensi pengolahan:} Joint mass didefinisikan dari jumlah
relasi antarperusahaan yang tertanam dalam jaringan, sementara
intensitas interaksi juga menjadi outcome. Kedekatan konseptual antara
kontrol massa dan outcome perlu diperhatikan ketika menafsirkan
koefisien, meskipun TXT tidak membahasnya sebagai masalah endogenitas
(berdasarkan hlm. 735-740, TXT baris 526-551 dan 655-733).

\subsection{Challenges}\label{challenges-47}

\textbf{Laporan sumber:} Tantangan utama penelitian adalah ketiadaan
data relasional antarperusahaan yang memadai dalam penelitian urban
network sebelumnya. Artikel harus mengubah jawaban establishment tentang
relasi penting menjadi jaringan dan kemudian mengagregasikannya ke 69
municipality (hlm. 727 dan 732, TXT baris 147-158 dan 391-415).

\textbf{Laporan sumber:} Perbedaan ukuran municipality dan jarak fisik
dapat membuat kota besar atau kota yang berdekatan tampak lebih
terhubung tanpa menunjukkan integrasi jaringan. Karena itu, visualisasi
mentah dianggap belum cukup dan perlu dikontrol melalui gravity model
(hlm. 733-734, TXT baris 463-488).

\textbf{Laporan sumber:} Data memiliki overdispersion dan excessive zero
counts. Hal tersebut memaksa pemilihan model yang mampu menangani dua
proses, yaitu pasangan yang selalu nol dan pasangan yang memiliki
peluang menghasilkan hitungan interaksi (hlm. 736-737, TXT baris
527-600).

\textbf{Laporan sumber:} Tantangan konseptual lainnya adalah membedakan
hubungan tradisional central place dari urban network. Penulis perlu
mengoperasionalisasikan perbedaan tersebut melalui enam tipe hubungan
dan tiga kondisi integrasi, lalu menghubungkannya dengan spesialisasi
pada 27 sektor (hlm. 737-742, TXT baris 601-837).

\textbf{Inferensi pengolahan:} Batas wilayah juga merupakan tantangan
analitis karena separuh relasi melampaui Randstad. Model yang hanya
melihat pasangan di dalam wilayah dapat mengukur embeddedness internal,
tetapi tidak secara lengkap mengukur posisi Randstad dalam jaringan
nasional dan internasional (berdasarkan Figure 2 dan kesimpulan artikel,
hlm. 733-745, TXT baris 428-462 dan 940-967).

\subsection{Practical Implications}\label{practical-implications-47}

\textbf{Interpretasi penulis:} Kebijakan ekonomi spasial sebaiknya tidak
mengasumsikan bahwa seluruh Randstad sudah menjadi satu kawasan urban
network yang terintegrasi. Fokus yang lebih sesuai adalah daily urban
systems dari empat kota terbesar atau wilayah yang lebih kecil di dalam
Randstad (hlm. 725 dan 744, TXT baris 47-52 dan 948-967).

\textbf{Interpretasi penulis:} Upaya kebijakan untuk menciptakan dan
mempertahankan pembangunan polisentris pada skala lebih tinggi perlu
dipertanyakan ketika komplementaritas berbasis relasi antarperusahaan
tidak melebihi skala city-region monosentris. Tanpa bukti interaksi
lintas wilayah yang memadai, target interaksi tingkat tinggi berisiko
lebih dekat pada wishful thinking daripada kondisi ekonomi yang teramati
(hlm. 744, TXT baris 940-952).

\textbf{Inferensi pengolahan:} Perencanaan dapat menggunakan diagnosis
berbasis arus sebelum menetapkan kawasan metropolitan sebagai satu unit
kebijakan. Diagnosis tersebut sebaiknya memeriksa apakah arus aktual
melampaui ekspektasi berdasarkan ukuran dan jarak, apakah hubungan masih
berhierarki, serta apakah perbedaan spesialisasi benar-benar disertai
pertukaran (berdasarkan Figure 3, Table 2-3, hlm. 734-743, TXT baris
443-488 dan 655-914).

\subsection{Applications}\label{applications-47}

\textbf{Laporan sumber:} Sistem pengujian artikel dapat diterapkan pada
wilayah lain karena penulis menyatakan bahwa sistem itu tidak bergantung
pada skala wilayah dan kota (hlm. 738 dan 744-745, TXT baris 625-632 dan
962-973).

\textbf{Inferensi pengolahan:} Fungsi aplikatif yang paling jelas adalah
membandingkan apakah pola wilayah lain juga mengikuti central place,
membentuk integrasi spasial, atau menunjukkan integrasi fungsional
(berdasarkan hlm. 737-738 dan 744-745, TXT baris 601-632 dan 962-973).

\textbf{Inferensi pengolahan:} Dalam kajian perkotaan, prosedur ini
dapat dipakai sebagai audit klaim urban network pada tiga tingkat:
pemeriksaan jaringan antar-establishment, pengelompokan relasi ke
subwilayah dan core-periphery, serta pengujian perbedaan spesialisasi
antarlokasi. Hasilnya dapat membantu membedakan label polisentris dari
integrasi ekonomi yang benar-benar teramati (berdasarkan hlm. 732-742,
TXT baris 391-415, 601-632, dan 769-837).

\textbf{Inferensi pengolahan:} Untuk pemantauan kebijakan, desain survei
yang sama dapat diulang pada periode berbeda dan dibandingkan dengan
kondisi 2005. Aplikasi ini harus disebut sebagai usulan pengolahan yang
diturunkan dari anjuran penulis untuk mengukur relasi antarperusahaan
secara berulang, bukan sebagai aplikasi yang sudah dilakukan dalam
artikel (hlm. 744-745, TXT baris 978-984).

\subsection{Future Research
(Optional)}\label{future-research-optional-47}

\textbf{Laporan sumber:} Penulis secara eksplisit mengusulkan penelitian
komparatif di wilayah lain untuk mengonfirmasi atau menantang hasil
Randstad. Sistem pengujian yang diperkenalkan dinilai dapat digunakan
untuk tujuan tersebut (hlm. 744-745, TXT baris 959-973).

\textbf{Laporan sumber:} Penulis juga menyatakan pentingnya pengukuran
berulang terhadap relasi antarperusahaan. Relasi tersebut mungkin
berkembang menuju jaringan antardaerah walaupun hubungan komuter dan
belanja yang dirujuk artikel tidak menunjukkan perkembangan
interregional yang mendalam selama 20 tahun sebelumnya (hlm. 744-745,
TXT baris 978-984 dan catatan 1, TXT baris 990-996).

\textbf{Inferensi pengolahan:} Agenda lanjutan yang konsisten dengan
batasan data adalah menguji arah supplier-customer secara terpisah,
mencakup sektor non-basic, dan menilai dampak batas eksternal Randstad.
Ketiganya bukan hasil penelitian ini dan tidak dinyatakan sebagai
rangkaian agenda lengkap oleh penulis.

\subsection{Provenance Notes}\label{provenance-notes-47}

\begin{itemize}
\tightlist
\item
  Review ini hanya menggunakan TXT B04 yang ditentukan pengguna. Seluruh
  isi TXT dibaca, dari blok metadata ekstraksi pada baris 1 sampai akhir
  referensi pada baris 1162. Tidak ada fakta yang ditambahkan dari web,
  PDF lain, Zotero, atau file review kode lain.
\item
  Locator utama menggunakan nomor halaman artikel yang tercetak dalam
  hasil ekstraksi, nomor bagian, nomor tabel atau gambar, dan rentang
  baris TXT. Nomor halaman artikel mencakup hlm. 725-748; nomor baris
  merujuk langsung pada file TXT.
\item
  Blok metadata menyatakan PDF tidak terenkripsi, berjumlah 24 halaman,
  dan diekstrak dengan \texttt{pdftotext\ -layout} pada 2026-08-31 (TXT
  baris 1-24). Review ini menilai teks hasil ekstraksi, bukan melakukan
  verifikasi eksternal terhadap PDF atau sumber referensi di dalam
  daftar pustaka.
\item
  Beberapa persamaan dan simbol pada TXT mengalami distorsi tata letak
  ekstraksi. Karena itu, metode dijelaskan secara substantif sesuai teks
  yang terbaca dan tidak merekonstruksi simbol persamaan yang tidak
  tampil utuh.
\item
  Informasi yang tidak diberikan atau tidak dapat dipastikan dari TXT
  telah ditandai sebagai ``Tidak disebutkan dalam file'' atau dijelaskan
  sebagai keterbatasan. Status peer-review, daftar pasti 67 atau 69
  municipality, bobot survei, profil nonresponden, matriks relasi
  mentah, serta prosedur causal identification tidak disebutkan dalam
  file.
\item
  Terdapat dua ketidaksesuaian internal yang dipertahankan apa adanya:
  Figure 1 menyebut 67 municipality, sementara narasi dan konstruksi
  jumlah observasi menunjuk 69; judul Table 3 menyebut
  \texttt{N\ =\ 2346}, sedangkan Model 5 dan Model 6 di dalam tabel
  mencantumkan \texttt{N\ =\ 2415} (TXT baris 77, 391-405, 655-658, dan
  869-874). Review tidak memilih salah satu angka sebagai fakta
  definitif.
\item
  TXT sumber tidak diubah. Setelah audit review selesai, sinkronisasi
  status B04 hanya dilakukan pada \texttt{Todo.md} dan
  \texttt{output/kajian/kode-kajian-pendahuluan.md} sesuai instruksi
  kerja; tidak ada dokumen induk atau review kode lain yang dibuat.
\end{itemize}

\section{Review B05}\label{review-b05}

\textbf{Identitas sumber}

\begin{itemize}
\tightlist
\item
  \textbf{Kode:} B05, berdasarkan prefiks nama file yang ditugaskan.
\item
  \textbf{Judul:} \emph{Assessing Polycentric Urban Systems in the OECD:
  Country, Regional and Metropolitan Perspectives}. Judul pada badan
  artikel terbaca ``in the''; metadata PDF pada TXT baris 7 memuat
  bentuk salah-urai ``in'the''. (Locator: TXT baris 7, 43-45, 97-99.)
\item
  \textbf{Penulis:} Monica Brezzi dan Paolo Veneri. (Locator: TXT baris
  47, 101-102.)
\item
  \textbf{Jurnal:} \emph{European Planning Studies}. (Locator: TXT baris
  38, 91-92.)
\item
  \textbf{Tahun dan publikasi:} badan artikel dan blok sitasi menyebut
  2014; dipublikasikan daring pada 4 Juni 2014. Nama file sumber memuat
  2015, sehingga perbedaan tahun tersebut tidak diselesaikan di luar
  informasi TXT. (Locator: TXT baris 2, 51, 59-61.)
\item
  \textbf{DOI:} 10.1080/09654313.2014.905005. (Locator: TXT baris 59-63,
  92.)
\item
  \textbf{Afiliasi yang tercantum:} Organisation for Economic
  Co-operation and Development, Regional Development Policy Division,
  GOV/RDP, Paris, France. (Locator: TXT baris 49-50, 107-108.)
\item
  \textbf{Riwayat naskah:} diterima Agustus 2013 dan disetujui Januari
  2014. (Locator: TXT baris 112.)
\item
  \textbf{Volume, nomor, dan halaman jurnal:} Tidak disebutkan dalam
  file. Nomor halaman artikel yang dipakai sebagai locator di bawah
  berasal dari penanda halaman dalam ekstraksi, bukan metadata
  volume/nomor.
\item
  \textbf{Bahan yang ditelaah:} seluruh TXT B05, termasuk blok metadata
  ekstraksi, teks artikel, catatan, dan daftar pustaka; tidak ada sumber
  eksternal yang digunakan.
\end{itemize}

\subsection{Summarized Abstract}\label{summarized-abstract-48}

\textbf{Laporan sumber:} Artikel memandang sistem perkotaan kontemporer
di negara-negara OECD sebagai sistem yang tersusun di sekitar functional
regions yang sering melampaui batas administratif kota. Dengan definisi
functional urban areas (FUAs) yang diharmonisasi di negara-negara OECD,
penulis mendefinisikan polycentricity pada tiga skala: metropolitan,
regional, dan nasional. Artikel kemudian mengukur polycentricity dan
mengeksplorasi implikasi ekonominya. Hasil ringkasnya adalah bahwa
region yang relatif lebih monocentric memiliki GDP per capita lebih
tinggi daripada region yang lebih polycentric, sedangkan pada tingkat
negara polycentricity berasosiasi dengan GDP per capita yang lebih
tinggi. (Locator: TXT baris 116-125; artikel p.~1.)

\textbf{Interpretasi penulis:} Penulis menekankan bahwa makna dan
implikasi kebijakan polycentricity sangat dipengaruhi oleh skala spasial
pengamatannya. (Locator: TXT baris 738-787; artikel p.~15.)

\textbf{Inferensi pengolahan:} Karena artikel menggunakan bahasa
asosiasi dan menyebut analisis ekonomi tersebut sebagai preliminary atau
descriptive, klaim ekonomi dalam abstrak tidak boleh dibaca sebagai
pernyataan bahwa polycentricity secara kausal menaikkan atau menurunkan
GDP per capita. (Locator: TXT baris 186-205, 640-647, 695-706, 778-787;
artikel pp.~2-3, 12-13, dan 15.)

\textbf{Inferensi pengolahan:} Pesan utama artikel adalah
ketergantungan-skala: struktur yang lebih hierarkis pada tingkat
regional muncul bersama tingkat GDP per capita yang lebih tinggi, tetapi
struktur nasional yang lebih polycentric muncul bersama GDP per capita
yang lebih tinggi. Perbedaan ini tidak boleh diringkas menjadi satu
resep spasial yang berlaku untuk semua tingkat wilayah. (Locator: TXT
baris 616-639, 695-705, 764-787; artikel pp.~11-15.)

\subsection{Problem Statement}\label{problem-statement-48}

\textbf{Laporan sumber:} Perubahan teknologi komunikasi, perpindahan
orang dan barang, serta perkembangan ekonomi memperluas ruang tempat
orang tinggal dan bekerja. Suburbanisasi dan integrasi kota dengan
hinterland membuat interdependensi ekonomi berlangsung dalam functional
regions, bukan semata-mata dalam batas administratif. Penulis menekankan
relevansi kebijakan karena 275 large FUAs disebut menampung sekitar
setengah populasi OECD dan berkontribusi lebih dari setengah GDP serta
pekerjaan OECD. (Locator: TXT baris 130-161; artikel pp.~1-2.)

\textbf{Laporan sumber:} Polycentricity dapat ditelaah melampaui skala
metropolitan, karena kota dan wilayah perkotaan dalam satu region atau
satu negara dapat saling terhubung. Pengetahuan tentang fungsi dan
efisiensi hubungan tersebut diperlukan untuk memahami kaitan urbanisasi
dan pembangunan ekonomi serta untuk menargetkan kebijakan, merencanakan
layanan publik, dan merancang tata kelola. Namun, bukti empiris tentang
hubungan polycentricity dan pembangunan ekonomi belum konklusif;
hasilnya bergantung pada negara yang dipilih, definisi konseptual, dan
indikator pengukuran. (Locator: TXT baris 167-185; artikel p.~2.)

\textbf{Interpretasi pengolahan:} Masalah penelitian terdiri atas dua
persoalan yang saling terkait: bagaimana membandingkan polycentricity
lintas negara secara konsisten, dan bagaimana menilai apakah hubungan
polycentricity dengan kemakmuran berubah menurut skala metropolitan,
regional, dan nasional. Artikel tidak menyajikan problem statement
sebagai subjudul formal tersendiri, tetapi kedua persoalan tersebut
dirumuskan dalam pendahuluan. (Locator: TXT baris 167-205; artikel
pp.~2-3.)

\subsection{Objectives}\label{objectives-48}

\textbf{Laporan sumber:} Tujuan pertama adalah menilai struktur
polycentric sistem perkotaan OECD pada tiga skala geografis, yaitu
metropolitan, regional, dan nasional, dengan definisi FUA yang sama yang
diterapkan pada 29 negara OECD. Tujuan kedua adalah mengeksplorasi
implikasi polycentricity pada tingkat nasional dan regional terhadap
tingkat kemakmuran ekonomi secara keseluruhan. (Locator: TXT baris
186-193; artikel p.~2.)

\textbf{Laporan sumber:} Struktur artikel menunjukkan fungsi tambahan
tiap bagian: bagian metropolitan menggambarkan organisasi spasial dan
perubahan menuju sprawl atau compact development; bagian regional dan
nasional mengukur polycentricity serta melakukan analisis awal
hubungannya dengan GDP per capita. (Locator: TXT baris 194-205; artikel
p.~2.)

\textbf{Interpretasi pengolahan:} Tujuan artikel bersifat pengukuran
komparatif dan eksplorasi asosiasi. Tidak ada tujuan yang dinyatakan
untuk mengidentifikasi efek kausal, mengevaluasi satu intervensi
kebijakan, atau menguji outcome ekonomi dan lingkungan pada skala
metropolitan. (Locator: TXT baris 186-205; artikel p.~2.)

\subsection{Literature Survey}\label{literature-survey-48}

\textbf{Laporan sumber:} Literatur yang dirangkum membedakan
polycentricity morfologis dan fungsional. Dimensi morfologis berfokus
pada ukuran dan distribusi populasi, pekerjaan, atau penggunaan lahan di
antara pusat-pusat kota. Dimensi fungsional berfokus pada fungsi yang
dijalankan kota dan hubungan antarkota, termasuk arus komuter. Penulis
menyatakan kedua dimensi terkait, tetapi tidak identik. (Locator: TXT
baris 230-257; artikel pp.~3-4.)

\textbf{Laporan sumber:} Mengikuti klasifikasi ESPON 3.1, artikel
menggunakan perspektif metropolitan, regional, dan nasional. Perspektif
metropolitan merujuk pada organisasi dalam satu atau beberapa labour
market areas yang tumpang tindih; perspektif regional merujuk pada
jaringan sedikitnya dua FUA yang terhubung secara fungsional dalam
region administratif yang lebih besar; perspektif nasional merujuk pada
sistem perkotaan seluruh negara. (Locator: TXT baris 258-280; artikel
p.~4.)

\textbf{Laporan sumber:} Untuk skala metropolitan, artikel mengaitkan
model monocentric tradisional dengan central business district dan
penurunan kepadatan menurut jarak, lalu menempatkan ekspansi wilayah
metropolitan, integrasi pusat baru, dan desentralisasi pekerjaan sebagai
tantangan terhadap model tersebut. Rujukan yang disebut di bagian ini
meliputi Alonso, Muth, Mills, Anas, White, dan Champion. (Locator: TXT
baris 317-346; artikel p.~5.)

\textbf{Laporan sumber:} Pembahasan metropolitan juga menghubungkan
polycentricity dengan decentralized concentration, sprawl, efisiensi
layanan publik, transportasi, energi, ruang hijau, dan penggunaan lahan.
Artikel menyebut bahwa bukti mengenai pengaruh polycentricity terhadap
waktu dan jarak perjalanan masih ambigu, sedangkan pekerjaan untuk
menilai apakah polycentricity memperbaiki efisiensi penggunaan lahan
atau kondisi lingkungan masih lebih terbatas dibandingkan kajian tentang
biaya sprawl. (Locator: TXT baris 305-316, 349-400; artikel pp.~5-7.)

\textbf{Laporan sumber:} Pada tingkat regional, literatur tentang
Polynucleated Urban Field dan polycentric urban region (PUR) digunakan
untuk menjelaskan wilayah dengan beberapa pusat yang terpisah secara
morfologis, berdekatan secara fisik, serta terhubung secara fungsional
dan/atau memiliki spesialisasi yang saling melengkapi. Gagasan borrowed
size dari Alonso dipakai untuk mengajukan bahwa sinergi dan
komplementaritas antarpusat dapat mengompensasi ketiadaan satu
aglomerasi besar. Parr juga dirujuk untuk konsep regional externalities
yang mencakup infrastruktur bersama, layanan khusus, serta faktor
pembangunan keras dan lunak. (Locator: TXT baris 413-448; artikel
pp.~7-8.)

\textbf{Laporan sumber:} Pada tingkat nasional, artikel meninjau bukti
bahwa distribusi ukuran kota umumnya mengikuti power law atau Pareto
distribution, lalu menggunakan hubungan rank-size untuk membandingkan
hierarki perkotaan antarnegara. Penjelasan teoretis yang disebut
meliputi guncangan acak terkait migrasi, produktivitas, dan inovasi,
serta pendekatan Christallerian mengenai fungsi kota dengan ukuran
berbeda. (Locator: TXT baris 649-687; artikel pp.~12-13.)

\textbf{Interpretasi pengolahan:} Literatur dipakai untuk membangun
definisi, skala analisis, dan ekspektasi borrowed size, bukan sebagai
systematic review. Prosedur pencarian literatur, kriteria inklusi,
ukuran corpus, dan penilaian kualitas studi terdahulu tidak disebutkan
dalam file. Artikel ini sebaiknya dipahami sebagai studi empiris
komparatif dengan tinjauan konseptual, bukan artikel review sistematis.
(Locator: TXT baris 1-935, khususnya 230-280, 413-448, dan 649-687;
artikel pp.~3-4, 7-8, dan 12-13.)

\subsection{Method Used}\label{method-used-48}

\textbf{Laporan sumber:} Identifikasi FUA menggabungkan sumber geografis
berbasis GIS dengan sumber administratif dan survei agar dapat menangkap
urban cores yang padat serta arus komuter tanpa terikat batas
administratif. Prosedurnya terdiri atas tiga langkah: mengidentifikasi
urban cores yang berdekatan atau sangat terhubung memakai data grid
populasi 1 km2; menggabungkan urban cores yang tidak berdekatan tetapi
terintegrasi bila lebih dari 15\% populasi bekerja dari salah satu core
berkomuter ke core lain; dan menentukan commuting shed atau hinterland
dengan memilih municipality yang mengirim sedikitnya 15\% angkatan
kerjanya ke cores. (Locator: TXT baris 208-227; artikel p.~3.)

\textbf{Laporan sumber:} Pada skala metropolitan, artikel menggunakan
pola kepadatan sebagai ilustrasi dan Sprawl Index (SI) untuk mengukur
pertumbuhan built-up area yang disesuaikan dengan pertumbuhan populasi.
Rumus yang tercantum adalah
\texttt{SI\_i\ =\ {[}urb\_i,t+n\ -\ (urb\_i,t\ *\ (pop\_i,t+n\ /\ pop\_i,t)){]}\ /\ urb\_i,t\ *\ 100},
dengan \texttt{urb} sebagai luas built-up area dan \texttt{pop} sebagai
populasi total metropolitan pada tahun awal dan akhir. Share penduduk
atau pekerjaan di urban centres disebut sebagai ukuran alternatif
polycentricity metropolitan. (Locator: TXT baris 336-382; artikel
pp.~5-6.)

\textbf{Laporan sumber:} Pada skala regional, unit analisisnya adalah
large administrative regions tingkat TL2. Penulis memakai urban primacy,
yaitu share populasi dalam FUA terbesar terhadap total populasi region,
dan koefisien beta dari regresi rank-size
\texttt{ln(rank)\ =\ a\ +\ b\ ln(size)}. Beta dipakai sebagai indikator
hierarki antarfungsi urban; karena beta secara definisi negatif, nilai
absolut beta yang lebih tinggi ditafsirkan sebagai polycentricity yang
lebih tinggi. Estimasi beta menggunakan empat FUA terbesar di setiap
region untuk menjaga konsistensi jumlah observasi. (Locator: TXT baris
449-532; artikel pp.~8-10.)

\textbf{Laporan sumber:} Untuk mengaitkan struktur regional dengan
ekonomi, penulis menjalankan empat spesifikasi OLS dengan dependent
variable GDP per capita. Model 1 memasukkan primacy; Model 2 memasukkan
polycentricity; Model 3 memasukkan primacy dan polycentricity; Model 4
mengganti primacy dengan metro\_share. Variabel kontrolnya meliputi
ukuran populasi region, pendidikan tersier, share penduduk di FUA, dan
dummy lokasi Eropa Timur serta Eropa Barat sesuai spesifikasi. Semua
variabel merujuk pada 2010 dan standard error yang dilaporkan bersifat
robust. (Locator: TXT baris 545-576, 601-608; artikel pp.~10-11.)

\textbf{Laporan sumber:} Pada skala nasional, beta diestimasi dari empat
FUA terbesar di setiap negara. Hubungan beta dan log GDP per capita 2010
ditampilkan melalui partial residual plot setelah mengontrol urban
primacy, pendidikan tersier, dan dummy Eropa Timur serta Eropa Barat.
Penulis menyebut analisis ini sebagai deskripsi awal karena jumlah
observasi terbatas dan desainnya lintas negara. (Locator: TXT baris
675-705, 718-735; artikel pp.~13-14.)

\textbf{Interpretasi pengolahan:} Strategi ini mengutamakan
keterbandingan lintas negara dan kesederhanaan indikator, dengan
mengorbankan dimensi konektivitas dan sebagian kompleksitas struktural.
OLS dan partial residual plot mengukur asosiasi bersyarat pada kontrol
yang digunakan; dari TXT tidak ada dasar untuk menyebut hasil tersebut
sebagai efek kausal. (Locator: TXT baris 465-505, 545-576, 640-647,
675-705; artikel pp.~8-13.)

\subsection{Dataset}\label{dataset-48}

\textbf{Laporan sumber:} Data utama berasal dari OECD regional and
metropolitan database. Definisi FUA dibangun dari population grid 1 km2,
informasi GIS, sumber administratif dan survei, sedangkan analisis
metropolitan menggunakan data built-up area dan populasi. Figure 2
menggunakan data National Census untuk membandingkan pola kepadatan
Paris dan San Francisco. Analisis nasional merujuk pada variabel dari
OECD database, termasuk URL database OECD yang dicantumkan dalam
catatan. (Locator: TXT baris 208-227, 242-245, 343-346, 553-569,
799-806; artikel pp.~3, 5-6, 10, 16.)

\textbf{Laporan sumber:} Tahun rujukan regresi regional dan nasional
adalah 2010. Sprawl Index membandingkan 2000 dan 2006. Data yang
digunakan mencakup populasi, luas built-up area, GDP per capita dalam
PPP US\$, pendidikan tersier, lokasi makrogeografis, jumlah dan ukuran
FUA, serta ukuran primacy dan beta. (Locator: TXT baris 367-382,
553-569, 695-704; artikel pp.~6, 10, 13.)

\textbf{Informasi tidak tersedia:} File data mentah, versi database,
definisi rinci setiap field, prosedur pembersihan atau imputasi, kode
analisis, dan daftar lengkap negara serta region pada tiap analisis:
\textbf{Tidak disebutkan dalam file}. TXT menyebutkan beberapa
pengecualian cakupan, termasuk ketiadaan land-use layer untuk SI dan
pengecualian negara atau region karena aturan pengukuran. (Locator: TXT
baris 383-409, 510-520, 688-694, 799-801; artikel pp.~6-7, 9, 13, dan
16.)

\textbf{Inferensi pengolahan:} Dataset ini lebih tepat dipahami sebagai
basis data spasial-ekonomi komparatif pada beberapa tingkat agregasi,
bukan dataset survei individu. Perbandingan antarhasil harus
memperhatikan bahwa cakupan dan unit analisis berubah menurut skala.
(Locator: TXT baris 208-227, 367-382, 449-569, dan 688-704; artikel
pp.~3, 6, 8-10, dan 13.)

\subsection{Population / Sample}\label{population-sample-48}

\textbf{Laporan sumber:} Cakupan konseptual artikel adalah sistem
perkotaan di 29 negara OECD. Pada tingkat regional, unit administratif
yang dipakai adalah region TL2, sedangkan node urban dibentuk oleh FUA
dan indikatornya dihitung dengan data seluruh FUA, termasuk core dan
hinterland. Region tanpa FUA diperlakukan sebagai rural, region dengan
satu FUA sebagai monocentric, dan region dengan dua atau lebih FUA
sebagai polycentric. (Locator: TXT baris 449-452, 521-528; artikel
pp.~8-9.)

\textbf{Laporan sumber:} Analisis OLS regional memiliki 206 observasi.
Untuk perbandingan koefisien beta, penulis menyebut 147 OECD regions
yang masing-masing memiliki sedikitnya empat FUA. Pada tingkat nasional,
analisis mencakup 26 negara OECD karena Luxembourg, Estonia, dan
Slovenia dikeluarkan; masing-masing hanya memiliki 1, 3, dan 2 FUA
sehingga tidak memenuhi pilihan empat FUA terbesar. (Locator: TXT baris
521-532, 588-599, 688-694; artikel pp.~9-11 dan 13.)

\textbf{Laporan sumber:} Untuk analisis SI metropolitan, cakupan yang
disebutkan dalam Figure 3 adalah Eropa, Jepang, dan Amerika Serikat.
Jumlah total metropolitan area yang dianalisis pada skala ini:
\textbf{Tidak disebutkan dalam file}. SI tidak dapat dihitung untuk
Canada, Chile, Korea, dan Mexico karena tidak tersedia land-use layer
pada dua titik waktu. (Locator: TXT baris 383-409, 799-801; artikel
pp.~6-7 dan 16.)

\textbf{Informasi tidak tersedia:} Sampel individual, responden, rumah
tangga, atau prosedur probability sampling: \textbf{Tidak disebutkan
dalam file}. Jumlah FUA total, jumlah TL2 yang dikeluarkan selain
informasi tentang primacy lebih besar dari 1, serta daftar unit
metropolitan lengkap juga \textbf{Tidak disebutkan dalam file}.
(Locator: TXT baris 1-935; cakupan yang tersedia diringkas pada baris
383-409, 449-532, dan 688-694.)

\textbf{Inferensi pengolahan:} ``Sample'' di sini adalah himpunan unit
wilayah yang tersedia dan memenuhi aturan pengukuran, bukan sampel
manusia. Pengurangan unit karena ketersediaan FUA, land-use layer, dan
kecocokan batas dapat menghasilkan cakupan analisis yang berbeda pada
tiga skala. (Locator: TXT baris 383-409, 449-532, dan 688-694; artikel
pp.~6-7, 8-11, dan 13.)

\subsection{Dependent Variables}\label{dependent-variables-48}

\textbf{Laporan sumber:} Pada regresi regional, dependent variable
adalah natural logarithm GDP per capita dalam Purchasing Power Parity
US\$ untuk 2010. Catatan Tabel 2 menyebutnya sebagai GDP per capita 2010
dalam US\$ PPP, sementara uraian metode menyatakan bahwa bentuk yang
digunakan adalah logaritma natural. (Locator: TXT baris 545-555,
578-608; artikel pp.~10-11.)

\textbf{Laporan sumber:} Pada analisis nasional, outcome yang dipetakan
adalah natural logarithm GDP per capita dalam PPP US\$ pada 2010. Figure
7 menunjukkan hubungan antara ukuran polycentricity nasional dan log GDP
per capita setelah kontrol terhadap primacy, pendidikan, serta dummy
Eropa Timur dan Barat. (Locator: TXT baris 695-705, 718-735; artikel
pp.~13-14.)

\textbf{Laporan sumber:} Pada skala metropolitan, artikel menyatakan
bahwa bagian ini tidak mengeksplorasi secara empiris hubungan struktur
spasial dengan outcome ekonomi, lingkungan, atau sosial. Sprawl Index
serta share penduduk atau pekerjaan di urban centres disajikan sebagai
ukuran struktur atau polycentricity metropolitan. (Locator: TXT baris
349-400; artikel pp.~5-7.)

\textbf{Inferensi pengolahan:} Dependent variable metropolitan dalam
model outcome: \textbf{Tidak berlaku} dalam desain yang dilaporkan.
Karena itu, Sprawl Index dan share penduduk atau pekerjaan tersebut
tidak boleh diperlakukan sebagai dependent variable dari regresi yang
menguji konsekuensi polycentricity. (Locator: TXT baris 349-400; artikel
pp.~5-7.)

\textbf{Informasi tidak tersedia:} Rincian sumber GDP, tahun dasar PPP,
transformasi GDP selain penyebutan logaritma natural, dan alasan formal
pemilihan dependent variable: \textbf{Tidak disebutkan dalam file}.
(Locator: TXT baris 545-569, 695-704; artikel pp.~10-13.)

\subsection{Results}\label{results-48}

\textbf{Laporan sumber - skala metropolitan:} Figure 2 memperlihatkan
bahwa kepadatan populasi di FUA Paris dan San Francisco tidak turun
secara monoton ketika jarak dari pusat utama bertambah; terdapat
beberapa local peaks yang dibaca sebagai indikasi sub-centres. (Locator:
TXT baris 336-346; artikel p.~5.)

\textbf{Laporan sumber - skala metropolitan:} SI menunjukkan
heterogenitas tinggi pada pola perkembangan metropolitan di Eropa,
Jepang, dan Amerika Serikat, tanpa pola sprawl keseluruhan. Pada
2000-2006, sekitar sepertiga metropolitan area mengalami variasi SI
positif, dengan nilai rata-rata OECD 0,8\%. Variasi positif berarti
built-up area tumbuh lebih cepat daripada populasi sehingga built-up
area per orang meningkat. Beberapa metropolitan area di Jepang, Las
Palmas, Zaragoza, dan Tallin memiliki nilai lebih dari 10\%;
metropolitan area tersebut pada 2000 memiliki built-up area per orang
yang relatif lebih rendah dibandingkan metropolitan area di Amerika
Serikat. (Locator: TXT baris 367-392; artikel pp.~6-7.)

\textbf{Laporan sumber - skala metropolitan:} Bagian ini tidak
mengestimasi hubungan empiris antara struktur metropolitan,
polycentricity, dan outcome ekonomi, lingkungan, atau sosial. Penulis
hanya menyatakan kemungkinan implikasinya dan menilai kajian tentang
apakah polycentricity lebih efisien atau lebih baik bagi lingkungan
daripada sprawl masih terbatas. (Locator: TXT baris 393-400; artikel
p.~7.)

\textbf{Laporan sumber - skala regional:} Tabel 2 memuat empat model OLS
dengan 206 observasi. Adjusted R-squared masing-masing adalah 0,376,
0,395, 0,402, dan 0,393; mean VIF masing-masing adalah 1,17, 1,39, 1,38,
dan 1,40. Catatan tabel menyatakan standard error robust dan tingkat
signifikansi 10\%, 5\%, serta 1\%. (Locator: TXT baris 570-608; artikel
pp.~10-11.)

\textbf{Laporan sumber - skala regional:} Primacy berkoefisien positif
dan signifikan pada Model 1, yaitu 0,372 dengan standard error 0,134,
serta pada Model 3, yaitu 0,256 dengan standard error 0,137. Ukuran
populasi region dan share tenaga kerja berpendidikan tersier berasosiasi
positif dengan GDP per capita pada seluruh model. Dummy Eropa Barat juga
positif dan signifikan pada seluruh model, sedangkan dummy Eropa Timur
tidak signifikan dalam tabel yang diekstrak. (Locator: TXT baris
582-596, 616-639; artikel pp.~10-12.)

\textbf{Laporan sumber - skala regional:} Baris
\texttt{beta*\ polycentricity} pada Tabel 2 terbaca sebagai
\texttt{20.170\ (0.057)***} pada Model 2 dan \texttt{20.149\ (0.058)**}
pada Model 3; karakter awal angka pada ekstraksi ini tidak stabil untuk
dibaca sebagai tanda aritmetika. Penulis tetap menyatakan bahwa tanda
koefisien tersebut tidak mendukung hipotesis bahwa region dengan
pusat-pusat berbeda ukuran dapat memperoleh GDP per capita lebih tinggi
melalui mekanisme borrowed size. Share total penduduk urban
(\texttt{metro\_share}) tidak signifikan dalam model yang memasukkannya.
(Locator: TXT baris 593-608, 621-639; artikel pp.~11-12.)

\textbf{Catatan audit ekstraksi:} Narasi penulis menyatakan bahwa
\texttt{d\_polycentricity} tidak signifikan. Namun, pada transkripsi
Tabel 2 di TXT, baris tersebut terbaca memuat \texttt{20.055\ (0.113)},
\texttt{20.052\ (0.112)}, dan \texttt{20.167\ (0.061)***} pada tiga
spesifikasi yang relevan. Karena tata letak tabel terpecah oleh
ekstraksi PDF, karakter awal angka tidak dinormalisasi di sini dan
narasi penulis tidak menguraikan perbedaan Model 4. Klaim yang aman
adalah narasi penulis menilai dummy tersebut tidak berkorelasi dengan
dependent variable, sedangkan pembacaan baris tabel perlu diverifikasi
pada dokumen asal. (Locator: TXT baris 578-608, 621-627; artikel
pp.~10-12.)

\textbf{Laporan sumber - skala nasional:} Pada 26 negara OECD, beta
nasional tampak berkorelasi positif dengan rata-rata kemakmuran ekonomi.
Figure 7 menunjukkan bahwa negara yang lebih polycentric rata-rata
memiliki GDP per capita lebih tinggi setelah beberapa kontrol. Penulis
menyatakan bahwa urban primacy tidak berasosiasi dengan GDP per capita
pada tingkat negara. Koefisien, standard error, p-value, dan
goodness-of-fit numerik untuk analisis nasional: \textbf{Tidak
disebutkan dalam file}. (Locator: TXT baris 688-705, 718-735; artikel
pp.~13-14.)

\textbf{Interpretasi penulis:} Hasil regional dan nasional bergerak
berlawanan. Penulis mengajukan bahwa polycentricity nasional mungkin
tidak merepresentasikan borrowed size karena area urban dan metropolitan
di negara polycentric belum tentu berdekatan atau membentuk network
relationship. Kemungkinan penjelasan lain adalah biaya aglomerasi lebih
rendah karena tersebar di beberapa FUA dan bagian wilayah nasional yang
lebih besar berada dekat setidaknya satu FUA besar. (Locator: TXT baris
723-735; artikel p.~14.)

\subsection{Findings}\label{findings-48}

\textbf{Laporan sumber:} Temuan lintas skala adalah bahwa struktur
spatial tidak memiliki implikasi ekonomi yang seragam. Pada tingkat
regional, primacy lebih tinggi dan tingkat polycentricity lebih rendah
berasosiasi dengan GDP per capita lebih tinggi. Pada tingkat nasional,
polycentricity lebih tinggi berasosiasi dengan GDP per capita lebih
tinggi. Pada tingkat metropolitan, hasil yang tersedia terutama
menggambarkan variasi kepadatan dan sprawl, bukan hubungan ekonomi.
(Locator: TXT baris 616-639, 695-705, 764-780; artikel pp.~11-15.)

\textbf{Interpretasi penulis:} Hasil regional dianggap mungkin
mencerminkan pentingnya kedekatan fisik dan aglomerasi orang serta
pekerja bagi kondisi sosial-ekonomi region. Hasil nasional mungkin
berkaitan dengan penyebaran manfaat kedekatan terhadap FUA besar dan
pengurangan biaya aglomerasi, tetapi penjelasan ini disajikan sebagai
kemungkinan, bukan mekanisme yang diuji langsung. (Locator: TXT baris
621-639, 723-735; artikel pp.~11-14.)

\textbf{Inferensi pengolahan:} Perbedaan hasil tidak membuktikan bahwa
monocentricity selalu lebih baik di region atau polycentricity selalu
lebih baik di negara. Perbedaan dapat dibaca sebagai indikasi bahwa unit
agregasi, definisi beta, kedekatan antar-FUA, dan mekanisme ekonomi yang
relevan berubah menurut skala. Karena konektivitas tidak diukur dan
desainnya observasional lintas wilayah, inferensi mekanistik harus tetap
terbatas. (Locator: TXT baris 510-520, 621-647, 688-735, dan 748-758;
artikel pp.~9, 11-15.)

\textbf{Inferensi pengolahan:} Artikel menghasilkan temuan morfologis
yang cukup kuat untuk perbandingan deskriptif, tetapi belum memberikan
bukti tentang bagaimana pusat-pusat tersebut berinteraksi atau apakah
perubahan struktur menyebabkan perubahan GDP. Kesimpulan tentang
network, borrowed size, agglomeration costs, dan territorial access
masih berada pada tingkat interpretasi atau hipotesis penjelas.
(Locator: TXT baris 465-485, 621-639, 695-735, dan 748-758; artikel
pp.~8-15.)

\subsection{Conclusions}\label{conclusions-48}

\textbf{Laporan sumber:} Penulis menyimpulkan bahwa polycentricity perlu
didefinisikan dan diperlakukan secara terpisah pada skala metropolitan,
regional, dan nasional. Definisi FUA yang konsisten digunakan sebagai
building block untuk membandingkan struktur metropolitan, regional, dan
nasional pada cakupan internasional yang luas. (Locator: TXT baris
738-747; artikel pp.~15.)

\textbf{Laporan sumber:} Perbandingan internasional memaksa penggunaan
ukuran polycentricity yang sangat aproksimatif. Artikel terutama
mengukur dimensi morfologis dan tidak memasukkan konektivitas
antarpusat, padahal konektivitas merupakan dimensi penting khususnya
pada tingkat regional. Penulis menyebut perlunya informasi lintas negara
yang lebih harmonis dan time series tahunan untuk menghasilkan bukti
yang lebih kuat. (Locator: TXT baris 748-758; artikel p.~15.)

\textbf{Laporan sumber:} Implikasi kebijakan berbeda menurut skala. Pada
tingkat metropolitan, perhatian mencakup efisiensi layanan publik,
transportasi, interaksi tatap muka, dan lingkungan. Pada tingkat
regional, aglomerasi besar dapat memberi manfaat ekonomi dan konsumsi,
sedangkan network kota dipandang berpotensi mengompensasinya. Namun,
eksplorasi artikel justru menemukan bahwa polycentricity regional yang
lebih rendah dan primacy lebih tinggi berasosiasi dengan GDP per capita
lebih tinggi. Pada tingkat nasional, struktur yang lebih polycentric
berasosiasi positif dengan GDP per capita. (Locator: TXT baris 764-784;
artikel p.~15.)

\textbf{Interpretasi penulis:} Artikel mendorong agar kebijakan yang
bertujuan mendorong polycentricity dirujukkan pada skala spasial
tertentu dan diinformasikan oleh kemungkinan implikasi ekonomi serta
lingkungan, bukan diterapkan sebagai tujuan umum tanpa diferensiasi
skala. (Locator: TXT baris 779-787; artikel p.~15.)

\textbf{Inferensi pengolahan:} Kesimpulan yang dapat dipertahankan dari
TXT adalah kesimpulan komparatif dan bersyarat, bukan rekomendasi
universal mengenai bentuk urban system yang optimal. (Locator: TXT baris
738-787; artikel pp.~15.)

\subsection{Contributions}\label{contributions-48}

\textbf{Laporan sumber:} Penulis menyatakan bahwa kontribusi artikel
adalah memperkuat pemahaman tentang struktur spasial sistem perkotaan
dan relevansinya bagi kebijakan. Kontribusi operasionalnya adalah
memakai definisi FUA yang konsisten untuk mengukur polycentricity pada
tiga skala dan pada basis lintas negara OECD yang luas. (Locator: TXT
baris 738-747; artikel p.~15.)

\textbf{Laporan sumber:} Artikel juga menyediakan ukuran sederhana yang
dapat dibandingkan, yaitu Sprawl Index pada skala metropolitan serta
primacy dan beta rank-size pada skala regional dan nasional. Ia
menambahkan eksplorasi awal hubungan ukuran tersebut dengan GDP per
capita untuk 206 observasi regional dan 26 negara pada analisis
nasional. (Locator: TXT baris 367-382, 465-490, 521-532, 688-705;
artikel pp.~6, 8-10, 13.)

\textbf{Inferensi pengolahan:} Nilai tambah utama yang tampak dari TXT
adalah pemisahan analitis antara tiga skala dan penolakan terhadap
penggunaan satu indikator polycentricity untuk semua skala. Nilai tambah
ini adalah kontribusi kerangka pengukuran dan benchmarking awal; bukan
validasi kausal atas kebijakan polycentric development. (Locator: TXT
baris 230-280, 367-400, 465-505, 675-705, dan 738-787; artikel pp.~3-6,
8-13, dan 15.)

\subsection{Research Gap}\label{research-gap-48}

\textbf{Gap yang disebutkan sumber:} Pengukuran perlu diperbaiki agar
memasukkan konektivitas antarpusat dan dimensi fungsional yang lebih
komprehensif. Penulis menyarankan harmonisasi informasi lintas negara
yang lebih lanjut untuk memenuhi kebutuhan tersebut. (Locator: TXT baris
640-647, 748-755; artikel pp.~12 dan 15.)

\textbf{Gap yang disebutkan sumber:} Diperlukan time series tahunan
untuk variabel FUA agar analisis longitudinal dapat dilakukan. Menurut
penulis, hal ini akan memberikan bukti empiris yang lebih kuat mengenai
relevansi struktur spasial dalam membentuk proses sosial-ekonomi negara,
region, dan metropolitan area. (Locator: TXT baris 753-758; artikel
p.~15.)

\textbf{Gap yang disebutkan sumber:} Pada skala metropolitan, masih
diperlukan lebih banyak kajian mengenai apakah polycentricity dapat
meningkatkan efisiensi penggunaan lahan atau kondisi lingkungan
dibandingkan sprawl. (Locator: TXT baris 393-400; artikel p.~7.)

\textbf{Inferensi pengolahan:} TXT juga menunjukkan celah untuk menguji
mekanisme yang menjelaskan hasil regional dan nasional yang berlawanan,
terutama peran jarak fisik, konektivitas, biaya aglomerasi, dan akses ke
FUA besar. Ini adalah inferensi dari batas desain dan penjelasan
penulis, bukan agenda penelitian tambahan yang dirumuskan secara
terpisah oleh sumber. (Locator: TXT baris 621-639, 723-735, dan 748-758;
artikel pp.~11-15.)

\subsection{Limitations}\label{limitations-48}

\textbf{Keterbatasan yang dinyatakan sumber:} Perbandingan lintas banyak
negara membuat ukuran polycentricity harus sangat aproksimatif. Ukuran
yang dipakai terutama morfologis dan mengabaikan konektivitas
antarpusat. (Locator: TXT baris 748-755; artikel p.~15.)

\textbf{Keterbatasan yang dinyatakan sumber:} Penyelarasan FUA yang
ditentukan secara fungsional dengan region TL2 yang administratif tidak
selalu mudah. FUA besar dapat melampaui region administratif, dan
beberapa FUA dapat berada dalam satu region. Untuk mengurangi
inkonsistensi, FUA dialokasikan berdasarkan lokasi urban core, tetapi
indikator spasial dihitung dengan seluruh FUA. Region dengan urban
primacy lebih besar dari 1 dikeluarkan dari analisis. (Locator: TXT
baris 510-520, 640-647; artikel pp.~9 dan 12.)

\textbf{Keterbatasan yang dinyatakan sumber:} Sprawl Index tidak dapat
dihitung untuk Canada, Chile, Korea, dan Mexico karena tidak tersedia
land-use layer pada dua titik waktu. Penggunaan empat FUA terbesar juga
mengecualikan Luxembourg, Estonia, dan Slovenia dari perbandingan
nasional karena jumlah FUA mereka terlalu sedikit untuk aturan tersebut.
(Locator: TXT baris 688-694, 799-801; artikel pp.~13 dan 16.)

\textbf{Keterbatasan yang dinyatakan sumber:} Analisis nasional memiliki
jumlah observasi terbatas dan bersifat cross-country, sehingga hanya
memungkinkan deskripsi pertama tentang hubungan yang diteliti. Analisis
regional juga disebut sebagai preliminary exploration atas hubungan yang
kompleks dan multidimensi. (Locator: TXT baris 640-647, 695-706; artikel
pp.~12-13.)

\textbf{Keterbatasan yang terlihat dari metode:} Bagian metropolitan
tidak menguji outcome ekonomi, sosial, atau lingkungan secara empiris;
ukuran beta memakai empat FUA terbesar; dan analisis regresi yang
dilaporkan berbasis asosiasi pada satu tahun, bukan desain longitudinal.
Poin terakhir adalah inferensi pengolahan dari periode dan metode yang
dijelaskan, bukan klaim bahwa penulis menyebut istilah ``tidak kausal''
secara eksplisit. (Locator: TXT baris 349-400, 521-532, 545-569,
688-705, dan 748-758; artikel pp.~5-7, 9-10, 13, dan 15.)

\textbf{Informasi tidak tersedia:} Analisis sensitivitas terhadap jumlah
FUA yang dipilih, model alternatif di luar empat spesifikasi regional,
validasi indikator konektivitas, atau pengujian kausal: \textbf{Tidak
disebutkan dalam file}. (Locator: TXT baris 1-935; keterbatasan dan
spesifikasi yang tersedia diringkas pada baris 465-505, 570-608, dan
748-758.)

\subsection{Challenges}\label{challenges-48}

\textbf{Laporan sumber:} Tantangan utama adalah mendefinisikan wilayah
fungsional secara konsisten ketika batas kehidupan dan kerja melampaui
batas administratif. Tantangan berikutnya adalah menyatukan sumber grid,
GIS, administratif, survei, dan land-use lintas negara dengan tingkat
ketersediaan yang berbeda. (Locator: TXT baris 208-227, 510-520; artikel
pp.~3 dan 9.)

\textbf{Laporan sumber:} Polycentricity memiliki dimensi ukuran,
distribusi, dan konektivitas. Ukuran komposit yang memuat semua dimensi
akan lebih lengkap, tetapi lebih diskresional dan membutuhkan lebih
banyak data. Penulis memilih ukuran sederhana dan langsung agar lebih
dapat dibandingkan. (Locator: TXT baris 465-485; artikel pp.~8-9.)

\textbf{Laporan sumber:} Menghubungkan self-organization melalui FUA
dengan struktur administratif TL2 tetap menjadi persoalan karena wilayah
fungsional dapat melintasi beberapa batas region. Hal ini menyulitkan
alokasi FUA ke entitas administratif yang dianggap tepat. (Locator: TXT
baris 640-647; artikel p.~12.)

\textbf{Interpretasi pengolahan:} Tantangan tersebut membatasi bukan
hanya presisi indikator, tetapi juga interpretasi hasil. Nilai beta atau
primacy yang berbeda dapat mencerminkan konfigurasi morfologis yang
tidak sepenuhnya menangkap hubungan fungsional antarkota. (Locator: TXT
baris 465-505, 510-520, dan 640-647; artikel pp.~8-9 dan 12.)

\subsection{Practical Implications}\label{practical-implications-48}

\textbf{Laporan sumber:} Pada skala metropolitan, distribusi penduduk
dan kegiatan ekonomi relevan untuk efisiensi penyediaan layanan publik,
transportasi, interaksi tatap muka, konsumsi energi, ruang hijau, dan
dampak lingkungan dari pola penggunaan lahan seperti sprawl. (Locator:
TXT baris 305-316, 764-769; artikel pp.~5 dan 15.)

\textbf{Laporan sumber:} Pada skala regional, keberadaan metropolitan
area besar dapat mendukung agglomeration economies dan manfaat konsumsi,
sedangkan jaringan beberapa kota dapat menjadi dasar kerja sama,
infrastruktur bersama, layanan khusus, dan komplementaritas regional.
Akan tetapi, hasil empiris awal artikel berasosiasi dengan GDP per
capita yang lebih tinggi pada region yang lebih hierarkis dan lebih
primacy-driven. (Locator: TXT baris 426-448, 764-777; artikel pp.~7-8
dan 15.)

\textbf{Laporan sumber:} Pada skala nasional, kebijakan urban perlu
mempertimbangkan kontribusi banyak kota, koherensi kebijakan, potensi
aglomerasi, dan ketimpangan teritorial. Artikel menyatakan bahwa tujuan
polycentric development perlu dirujukkan pada skala spasial spesifik dan
diinformasikan oleh kemungkinan konsekuensi ekonomi serta lingkungan.
(Locator: TXT baris 283-302, 649-659, 778-787; artikel pp.~4, 12, dan
15.)

\textbf{Interpretasi pengolahan:} Implikasi praktis yang paling aman
adalah memakai skala sebagai penyaring awal kebijakan, bukan menjadikan
koefisien regresi sebagai dasar tunggal untuk memindahkan investasi dari
kota besar ke kota kecil atau sebaliknya. Hasil yang disebut preliminary
dan keterbatasan konektivitas mengharuskan validasi kontekstual sebelum
kebijakan diterapkan. (Locator: TXT baris 621-647, 695-706, dan 764-787;
artikel pp.~11-15.)

\subsection{Applications}\label{applications-48}

\textbf{Laporan sumber:} Kerangka FUA dapat digunakan untuk membaca
struktur urban tanpa membatasi analisis pada batas administratif. Pada
skala metropolitan, SI dapat dipakai untuk membandingkan pertumbuhan
built-up area terhadap pertumbuhan populasi. Pada skala regional dan
nasional, primacy serta distribusi rank-size dapat digunakan untuk
menggambarkan hierarki dan tingkat polycentricity. (Locator: TXT baris
208-227, 349-382, 482-505, 675-694; artikel pp.~3, 5-6, 9, dan 13.)

\textbf{Inferensi pengolahan:} Aplikasi yang didukung oleh desain ini
adalah benchmarking deskriptif lintas wilayah, pemetaan awal pola
hierarki urban, dan penyaringan pertanyaan kebijakan yang perlu diuji
lebih lanjut. Contoh Paris-San Francisco, Aragon-Brittany, serta
Korea-Jerman berfungsi sebagai ilustrasi struktur atau rank-size, bukan
evaluasi program kebijakan. (Locator: TXT baris 336-346, 458-479,
679-687; artikel pp.~5 dan 8-9, 13.)

\textbf{Informasi tidak tersedia:} Studi implementasi, evaluasi
intervensi, panduan operasional pemerintah, atau pengujian dampak
aplikasi kebijakan: \textbf{Tidak disebutkan dalam file}. (Locator: TXT
baris 1-935; aplikasi yang dilaporkan terbatas pada ukuran dan
perbandingan deskriptif pada baris 208-227, 349-382, 482-505, dan
675-694.)

\subsection{Future Research
(Optional)}\label{future-research-optional-48}

\textbf{Laporan sumber:} Penulis menyarankan dua penyempurnaan utama:
pertama, mengharmonisasi informasi lintas negara agar pengukuran dapat
memasukkan konektivitas dan dimensi fungsional secara lebih
komprehensif; kedua, menyusun time series tahunan variabel FUA untuk
memungkinkan analisis longitudinal. (Locator: TXT baris 748-758; artikel
p.~15.)

\textbf{Laporan sumber:} Sumber juga menunjukkan kebutuhan riset lebih
lanjut tentang hubungan polycentricity metropolitan dengan efisiensi
penggunaan lahan dan kondisi lingkungan dibandingkan sprawl. (Locator:
TXT baris 393-400; artikel p.~7.)

\textbf{Inferensi pengolahan:} Berdasarkan keterbatasan yang dijelaskan
sumber, agenda yang paling langsung adalah menguji apakah hubungan
regional dan nasional tetap bertahan ketika konektivitas, jarak
antar-FUA, dan perubahan waktu tersedia. Ini merupakan inferensi yang
ditandai, bukan rekomendasi tambahan dari luar TXT. (Locator: TXT baris
621-647, 723-735, dan 748-758; artikel pp.~11-15.)

\subsection{Provenance Notes}\label{provenance-notes-48}

\begin{itemize}
\tightlist
\item
  \textbf{File sumber yang digunakan:}
  \texttt{/Users/mac/Documents/Mac/{[}2{]}\ Obsidian\ Vault/zettelkasten/source/journal/B05-brezzi-veneri-2015-assessing-polycentric-urban-systems-in-the-oecd-country-regional-and-metropolitan-perspectives-10.1080-09654313.2014.905005.txt}.
\item
  \textbf{Cakupan pembacaan:} TXT dibaca dari baris 1 sampai 935,
  termasuk metadata PDF pada baris 1-28, seluruh teks artikel, catatan,
  dan daftar pustaka.
\item
  \textbf{Metadata ekstraksi:} TXT mencatat ekstraksi pada 2026-08-31
  dengan metode \texttt{pdftotext\ -layout}; metadata PDF mencatat 20
  pages. File TXT juga memuat pemberitahuan unduhan oleh Selcuk
  Universitesi pada 7 Februari 2015. (Locator: TXT baris 1-3, 22, 30-34,
  90.)
\item
  \textbf{Konvensi locator:} Locator utama memakai nomor baris TXT;
  nomor halaman artikel dicantumkan bila penanda halaman tersedia dalam
  ekstraksi. Locator tidak dimaksudkan sebagai nomor halaman jurnal
  karena volume, nomor, dan pagination bibliografis tidak disebutkan.
\item
  \textbf{Batas akses:} Review ini hanya memproses TXT yang diberikan.
  PDF asli, halaman penerbit, database OECD, dan referensi yang
  tercantum tidak dibuka atau digunakan untuk menambah informasi.
\item
  \textbf{Pembedaan klaim:} ``Laporan sumber'' berarti isi yang
  dinyatakan atau ditampilkan penulis; ``Interpretasi penulis'' berarti
  pembacaan atau penjelasan yang diberikan artikel; ``Interpretasi
  pengolahan'' dan ``Inferensi pengolahan'' berarti sintesis terbatas
  dari TXT oleh proses review.
\item
  \textbf{Ketidakpastian penting:} Ekstraksi layout membuat sebagian
  tanda negatif dan penjajaran Tabel 2 tidak sepenuhnya jelas. Review
  mempertahankan hasil yang dapat dibaca, mencatat ketidaksesuaian
  antara narasi dan baris \texttt{d\_polycentricity}, dan tidak mengisi
  koefisien nasional yang tidak tersedia.
\item
  \textbf{Perubahan workspace:} TXT sumber tidak diubah. Audit ini hanya
  memperbarui \texttt{review-B05.md} dan status B05 pada
  \texttt{Todo.md} serta daftar kode; kode lain dan dokumen induk tidak
  ditangani.
\end{itemize}

\section{Review B06}\label{review-b06}

\textbf{Laporan sumber:} Kode review: B06.

\textbf{Laporan sumber:} Judul sumber: \emph{The productivity effects of
polycentricity: A systematic analysis of urban regions in Europe}.

\textbf{Laporan sumber:} Penulis: Freke Caset, Yuting Yang, Ben
Derudder, dan Krasen Samardzhiev.

\textbf{Laporan sumber:} Publikasi: \emph{Papers in Regional Science},
2023, volume 102, issue 6, halaman 1193-1213 menurut keterangan sitasi
pada TXT. Metadata PDF pada bagian awal menyebut rentang 1193-1214;
perbedaan ini dipertahankan sebagai catatan provenance dan tidak
diselesaikan melalui dugaan.

\textbf{Laporan sumber:} Identifier: DOI 10.1111/pirs.12765.

\textbf{Laporan sumber:} Jenis sumber: Artikel jurnal. Status peer
review tidak diverifikasi dalam TXT.

\textbf{Laporan sumber:} Afiliasi dan pendanaan: Afiliasi penulis
mencakup Ghent University, Vrije Universiteit Brussel, KU Leuven,
Nicolaus Copernicus University, dan University of Liverpool. Pendanaan
disebut berasal dari Fonds Wetenschappelijk Onderzoek, grant G014119N.

\textbf{Catatan provenance:} Hanya file TXT B06 yang diminta. TXT
tersebut memuat metadata ekstraksi, teks artikel, tabel, catatan kaki,
referensi, pemberitahuan supporting information, serta ringkasan bahasa
Spanyol dan Jepang. Penanda \texttt{Laporan\ sumber},
\texttt{Interpretasi\ penulis}, dan \texttt{Inferensi\ pengolahan}
digunakan untuk menjaga pemisahan epistemik.

\subsection{Summarized Abstract}\label{summarized-abstract-49}

\textbf{Laporan sumber:} Artikel menguji sejauh mana polycentric urban
regions (PURs) dapat menggantikan manfaat ekonomi aglomerasi yang
biasanya disediakan kota besar. Penulis menggunakan perangkat lunak
sumber terbuka PURban untuk mengidentifikasi perkembangan polisentris,
lalu mengukur struktur spasial berdasarkan ukuran, dispersi, dan
polycentricity pada 94 urban regions (URs) di dalam cakupan 34 negara
Eropa. Struktur tersebut dihubungkan dengan total factor productivity
(TFP). Abstrak melaporkan bahwa wilayah yang lebih polisentris dan
wilayah yang lebih tersebar berasosiasi dengan tingkat produktivitas
yang lebih rendah. (Artikel, abstrak, TXT baris 52-77; ringkasan
Spanyol, TXT baris 1107-1112.)

\textbf{Interpretasi penulis:} Hasil utama dipakai untuk menolak
dukungan empiris terhadap dugaan umum bahwa polycentricity meningkatkan
produktivitas regional di Eropa. Dalam pendahuluan, penulis merumuskan
hasilnya sebagai efek POLY yang negatif atau tidak ada, bergantung pada
apakah ukuran populasi dan dispersi dimasukkan, serta sebagai indikasi
berlanjutnya pentingnya manfaat aglomerasi berpusat tunggal. (Artikel,
Bagian 1, hlm. 1193-1194; TXT baris 122-137.)

\textbf{Inferensi pengolahan:} Temuan ini harus dibaca sebagai hubungan
statistik dalam kumpulan UR yang dibentuk oleh kriteria tertentu, bukan
sebagai bukti bahwa semua bentuk polycentricity selalu menurunkan
produktivitas atau bahwa struktur monosentris pasti lebih baik. Artikel
sendiri membedakan polycentricity morfologis dari kemungkinan pendekatan
fungsional, dan beberapa hasil POLY menjadi tidak signifikan ketika DISP
diperhitungkan. (TXT baris 753-776 dan 911-926.)

\subsection{Problem Statement}\label{problem-statement-49}

\textbf{Laporan sumber:} Penulis menempatkan masalah dalam perubahan
kota menjadi wilayah urban yang terintegrasi, dengan hubungan antara
struktur spasial regional dan produktivitas ekonomi sebagai pertanyaan
penting dalam studi urban dan regional. PUR dipahami sebagai wilayah
yang memiliki sinergi antara pusat-pusat kota yang berdekatan dan
terhubung secara padat. Di Eropa, polycentric development juga telah
menjadi paradigma kebijakan yang dikaitkan sekaligus dengan cohesion
teritorial dan competitiveness. Namun, bukti empiris mengenai manfaat
produktivitasnya disebut langka, terfragmentasi, tidak merata, dan
kontradiktif. (Artikel, Bagian 1 dan 2.1, hlm. 1193-1195; TXT baris
81-121 dan 141-188.)

\textbf{Laporan sumber:} Masalahnya bukan hanya apakah polycentricity
berkaitan dengan produktivitas, tetapi juga bagaimana hubungan tersebut
diukur. Studi terdahulu memakai konsepsi, indikator, unit wilayah, skala
analisis, serta ukuran kinerja ekonomi yang berbeda. Penulis mencatat
bahwa perbedaan kecil dalam definisi dan operasionalisasi dapat
menghasilkan keluaran empiris yang sangat berbeda. Ukuran produktivitas
juga sensitif terhadap tingkat disagregasi wilayah dan pilihan antara
GDP per capita dan TFP. (TXT baris 293-312.)

\textbf{Interpretasi penulis:} Narasi kebijakan bahwa pusat-pusat kota
dapat menggabungkan manfaat aglomerasi melalui borrowed size belum
memperoleh dasar bukti yang konsisten pada skala regional Eropa.
Penjelasan alternatif yang dibahas mencakup tetap pentingnya kedekatan
fisik, biaya perjalanan dan arus komoditas yang lebih panjang, duplikasi
fungsi perkotaan, persaingan antarpusat, agglomeration shadow, serta
keunggulan kota tunggal sebagai node integrasi ekonomi global. (TXT
baris 320-383.)

\textbf{Inferensi pengolahan:} Problem statement artikel dapat diringkas
sebagai persoalan pengujian empiris yang sekaligus bersifat konseptual,
spasial, dan pengukuran: apakah struktur polisentris benar-benar
menggantikan aglomerasi kota besar, setelah ukuran dan dispersi
diperhitungkan, dengan unit wilayah yang tidak dipaksakan oleh batas
administratif. Ini adalah rekonstruksi dari abstrak dan pendahuluan,
karena TXT tidak memberikan satu kalimat pertanyaan penelitian yang
diformat secara eksplisit.

\subsection{Objectives}\label{objectives-49}

\textbf{Laporan sumber:} Tujuan yang dinyatakan melalui pendahuluan dan
rancangan penelitian adalah memperluas analisis polycentric development,
produktivitas ekonomi, dan hubungan keduanya dengan cara yang lebih
komparatif, rinci, dan dapat direproduksi. Penulis secara khusus
menyatakan kontribusi pada tiga hal: memakai data yang lebih
fine-grained dan menghubungkannya dengan unit geografis UR; mencakup
wilayah Eropa yang lebih luas dengan lebih dari dua kali jumlah negara
dibanding Ouwehand et al.~(2022); dan memeriksa ketergantungan spasial
antarvariabel. (TXT baris 122-130 dan 417-427.)

\textbf{Inferensi pengolahan:} Dari tujuan tersebut, sasaran operasional
studi adalah:

\begin{itemize}
\tightlist
\item
  Mengidentifikasi UR berdasarkan kondisi morfologis yang mendukung
  kemungkinan interaksi dan aglomerasi regional melalui PURban.
\item
  Mengukur derajat polycentricity, ukuran populasi, dan dispersi pada UR
  yang teridentifikasi.
\item
  Mengukur produktivitas regional dengan TFP, bukan hanya GDP per
  capita.
\item
  Mengestimasi hubungan struktur spasial dan TFP dengan regresi spasial,
  sambil menguji apakah ketergantungan spasial dan endogeneity mengubah
  hasil.
\item
  Menilai apakah polycentricity dapat berperan sebagai pengganti manfaat
  aglomerasi kota besar dalam konteks Eropa.
\end{itemize}

\textbf{Inferensi pengolahan:} Tujuan operasional di atas merupakan
sintesis dari teks sumber, bukan daftar tujuan bernomor yang ditulis
penulis.

\subsection{Literature Survey}\label{literature-survey-49}

\textbf{Laporan sumber:} Literatur dibahas melalui dua jalur. Jalur
pertama menjelaskan landasan teoritis dan kebijakan: ESDP 1999,
Territorial Agenda 2030, Pact of Amsterdam, serta kebijakan polycentric
development di berbagai negara Eropa. Landasan ekonominya adalah
agglomeration externalities, meliputi urbanization externalities dari
pasar dan tenaga kerja yang besar, localization externalities dari
ko-lokasi sektor, dan complexity externalities dari ko-presensi sektor
yang berbeda. (TXT baris 141-172.)

\textbf{Laporan sumber:} Jalur kedua membahas borrowed size. Mengikuti
Alonso (1973) sebagaimana dipaparkan artikel, kota yang lebih kecil
dalam kompleks metropolitan dapat mengakses manfaat ukuran kota yang
lebih besar melalui akses mudah ke pusat lain. Karena itu, PUR
diperkirakan dapat mempertahankan keuntungan kota yang lebih kecil
sembari memperoleh critical mass dan spillover dari jejaring kota di
sekitarnya. Artikel juga menyebut studi yang menemukan asosiasi positif,
negatif, maupun tidak signifikan antara polycentricity dan
produktivitas. (TXT baris 162-181.)

\textbf{Laporan sumber:} Table 1 merangkum studi pan-Eropa yang dipilih
penulis. ESPON 1.1.1 (2004) melaporkan hubungan positif antara
polycentricity negara dan GDP per capita. ESPON 1.4.3 (2007) tidak
menemukan korelasi signifikan dan melihat kecenderungan kecil ke arah
kinerja yang lebih baik pada negara yang lebih monosentris. Vandermotten
et al.~(2007) juga mengaitkan keberhasilan yang lebih tinggi dengan
wilayah yang lebih monosentris. Meijers dan Sandberg (2008)
menghubungkan sistem urban yang lebih polisentris dengan disparitas
regional GDP per capita yang lebih besar. Brezzi dan Veneri (2015)
menemukan hasil yang berbeda menurut skala: relatif lebih monosentris
pada tingkat regional berkaitan dengan GDP per capita yang lebih tinggi,
sementara pada tingkat negara polycentricity dikaitkan dengan GDP per
capita yang lebih tinggi. Meijers dan Burger (2017) menemukan borrowed
size lebih sering pada wilayah metropolitan polisentris. Ouwehand et
al.~(2022) tidak menemukan dampak langsung polycentricity terhadap TFP,
tetapi menemukan pengaruh yang lebih negatif ketika ukuran wilayah
meningkat. (Artikel, Bagian 2.2, Table 1, hlm. 1195-1197; TXT baris
192-292.)

\textbf{Interpretasi penulis:} Perbedaan hasil terutama dipahami sebagai
akibat dari ketidaksamaan analytical-operational scheme: definisi
polycentricity, komposisi indeks, unit dan skala analisis, serta ukuran
produktivitas. Perbedaan itu membuat temuan antarstudi tidak dapat
dibandingkan secara sederhana. Penulis juga menekankan bahwa mayoritas
studi terdahulu memakai indikator morfologis yang mengukur keseimbangan
ukuran pusat, bukan relasi aktual antar pusat. (TXT baris 293-316.)

\textbf{Laporan sumber:} Untuk menjelaskan hasil negatif atau tidak
signifikan, artikel membahas beberapa mekanisme dalam literatur:
globalisasi dan tersierisasi yang menguntungkan node layanan tingkat
tinggi; pentingnya proximity fisik dibanding relational proximity; biaya
perjalanan, arus komoditas, dan informasi; duplikasi fungsi tingkat
rendah; kompetisi dan rivalry; kepadatan serta interaksi tatap muka;
economic stress; serta growth shadow atau agglomeration shadow. (TXT
baris 320-363.)

\textbf{Interpretasi penulis:} Literatur memberi alasan untuk tidak
memodelkan POLY secara terisolasi. Ukuran kota dan tingkat dispersi
dapat memoderasi atau mengubah hubungan POLY-TFP, tetapi temuan mengenai
interaksi ukuran dan polycentricity masih bertentangan. Artikel
menggunakan ketidakselarasan ini untuk membenarkan model yang memasukkan
POLY, POP, DISP, dan interaksinya. (TXT baris 364-383.)

\textbf{Inferensi pengolahan / batas survey:} TXT tidak memuat protokol
pencarian basis data, kriteria inklusi-eksklusi, tanggal pencarian, atau
klaim bahwa seluruh publikasi yang relevan telah dikumpulkan. Karena
itu, review ini memperlakukan bagian tersebut sebagai tinjauan naratif
dan perbandingan studi yang dirangkum dalam Table 1; corpus systematic
literature review yang dapat diaudit tidak dapat diverifikasi dari TXT.
Istilah ``systematic analysis'' pada judul dipertahankan sebagai judul
sumber, tanpa menggunakannya untuk menyimpulkan protokol survei
literatur yang tidak dijelaskan.

\subsection{Method Used}\label{method-used-49}

\textbf{Laporan sumber:} Studi menggunakan desain kuantitatif komparatif
pada skala interurban regional, dengan perspektif morfologis yang
mengikuti degree of urbanization. PURban digunakan untuk membangun unit
analisis dari urban centres (UCs), sedangkan TFP digunakan sebagai
ukuran produktivitas. Regresi spasial dan instrumental variables
analysis digunakan untuk menguji hubungan struktur spasial dan
produktivitas. (Artikel, Bagian 3, hlm. 1200-1201; TXT baris 431-449.)

\textbf{Laporan sumber:} PURban mendefinisikan UC berdasarkan GHS Urban
Centre Database (GHS-UCDB). UC adalah klaster sel grid bersebelahan
seluas 1 km2 dengan kepadatan penduduk sekurang-kurangnya 1.500 penduduk
per km2 atau built-up area sekurang-kurangnya 50\%, serta jumlah
penduduk minimum 50.000. Basis ini menyediakan 734 UCs untuk cakupan 34
negara yang dianalisis. (TXT baris 451-460.)

\textbf{Laporan sumber:} Catatan kaki 3 menyebut bahwa UC yang berada di
negara EU dipilih kecuali yang berada di Iceland, Corsica, Sardinia,
Ibiza, Mallorca, dan Crete; UC non-EU yang secara geografis berada di
antara negara-negara EU, termasuk Kaliningrad yang berada di Rusia serta
UC di Switzerland, Serbia, Albania, dan negara Balkan non-EU lainnya,
juga dimasukkan. (TXT baris 465-468.)

\textbf{Laporan sumber:} PURban menyaring PUR/UR menggunakan empat
kondisi morfologis dan spasial:

\begin{itemize}
\tightlist
\item
  Jarak geografis maksimum antar pasangan UC adalah waktu tempuh 45
  menit dengan mobil melalui jalan, termasuk keterlambatan kemacetan.
\item
  Setiap PUR harus memiliki sedikitnya dua UC.
\item
  Total penduduk PUR sekurang-kurangnya 1.5 juta orang.
\item
  Pusat-pusatnya harus relatif seimbang dalam kepentingan atau ukuran,
  yang diukur melalui indeks gabungan dari primacy index, Herfindahl
  index, rank-size slope, dan indeks berbasis standard deviation.
\end{itemize}

\textbf{Laporan sumber:} Indeks komposit dibentuk melalui min-max
normalization pada setiap indikator, dirata-ratakan, lalu dinormalisasi
lagi sehingga bernilai 0 sampai 1. Perhitungan hanya memakai empat UC
terbesar menurut total penduduk. (TXT baris 461-495.)

\textbf{Laporan sumber:} Algoritme menyaring area di sekitar setiap UC,
mengidentifikasi kondisi 45 menit, jumlah pusat, dan critical mass,
kemudian menghasilkan 94 UR unik yang berada di 10 negara. UR yang
terdiri dari himpunan UC persis sama hanya dihitung satu kali. Karena
algoritme dapat menghasilkan UR yang sebagian tumpang tindih, penulis
memeriksa autocorrelation spasial dan menggunakan model ekonometrika
spasial. (TXT baris 497-535.)

\textbf{Laporan sumber:} TFP dihitung melalui dua persamaan. Persamaan
pertama memprediksi GVA dari modal dan tenaga kerja, untuk setiap PUR
dan seluruh sektor ekonomi:

\begin{itemize}
\tightlist
\item
  \texttt{ln\ GVA\_j\ =\ rho\_1\ W\ ln\ GVA\_j\ +\ kappa\ ln\ K\_j\ +\ delta\ ln\ L\_j\ +\ epsilon\_j}.
\end{itemize}

Residual dari persamaan tersebut diperlakukan sebagai TFP dan menjadi
variabel dependen pada persamaan utama:

\begin{itemize}
\tightlist
\item
  \texttt{TFP\_j\ =\ alpha\_0\ +\ rho\_2\ W\ TFP\_j\ +\ alpha\_1\ POLY\_j\ +\ alpha\_2\ ln\ POP\_j\ +\ alpha\_3\ DISP\_j\ +\ alpha\_4\ X\_j\ +\ epsilon\_j}.
\end{itemize}

\textbf{Laporan sumber:} W adalah matriks bobot spasial: untuk pasangan
PUR i dan j, nilainya 1 jika ada sedikitnya satu UC yang sama-sama
dimiliki, dan 0 jika tidak. Variabel X mencakup human capital,
infrastructure connectivity, dan innovation. (TXT baris 516-577.)

\textbf{Laporan sumber:} Penulis memulai dari spatial Durbin model
sebagai titik awal, menguji beberapa spesifikasi, dan menyatakan bahwa
SAR paling sesuai. Lagrange multiplier tests menunjukkan adanya efek
spasial; likelihood ratio test memilih SAR dibandingkan alternatif yang
diuji. Koefisien spasial pada TFP digunakan untuk menangkap
ketergantungan produktivitas antar-UR. (TXT baris 530-535 dan catatan
kaki 6 pada baris 581-587.)

\textbf{Laporan sumber:} Karena arah hubungan dapat bersifat dua arah,
penulis juga menerapkan generalized spatial two-stage least square
(GS2SLS). Instrumen dibangun dari kondisi historis tahun 1975:
historical polycentricity (POLY\_HIST), historical population
(POP\_HIST), dan historical dispersion (DISP\_HIST). Asumsi instrumen
adalah kondisi historis membentuk struktur spasial saat ini tetapi tidak
berkaitan dengan produktivitas regional saat ini. Diagnostics menyatakan
instrumen relevan dan seluruh variabel independen utama eksogen kecuali
POP. Karena itu, hanya POP yang diinstrumentasikan dalam model GS2SLS,
sedangkan POLY dan DISP diestimasi dalam SAR. (TXT baris 592-612.)

\textbf{Laporan sumber:} Table 3 memakai pendekatan stepwise: variabel
utama ditambahkan sendiri-sendiri, berpasangan, atau melalui interaksi
POLY\emph{DISP, POLY}POP, dan DISP*POP. Robust standard errors
dilaporkan. Seluruh model diberi label ``Region Fixed'' dalam tabel,
tetapi TXT tidak menjelaskan lebih lanjut spesifikasi teknis label
tersebut. (TXT baris 739-752 dan 791-822.)

\textbf{Inferensi pengolahan:} Strategi ini meningkatkan perhatian
terhadap skala, tumpang tindih wilayah, autocorrelation, dan potensi
endogeneity, tetapi tidak mengubah desain menjadi eksperimen. Koefisien
tetap perlu dibaca sebagai asosiasi model yang bergantung pada definisi
UR, pemilihan sampel, konstruksi TFP, validitas instrumen, dan asumsi
alokasi data.

\subsection{Dataset}\label{dataset-49}

\textbf{Laporan sumber:} Data utama berasal dari beberapa sumber dan
tidak dikumpulkan pada satu tahun yang seragam. Table 2 mencantumkan
hal-hal berikut:

\begin{itemize}
\tightlist
\item
  GVA untuk EU27 bersumber dari Eurostat tahun 2018; untuk UK bersumber
  dari Office for National Statistics tahun 2018. Variabel ini tersedia
  pada NUTS 3, dengan keterangan klasifikasi NUTS 3 (2016) dalam tabel.
\item
  LAB berupa total employment berbasis tempat kerja untuk EU27 dari
  Eurostat/Ardeco, dengan rentang tahun yang dicantumkan 1980-2023; data
  UK berupa jumlah orang bekerja menurut occupation code dari Office for
  National Statistics, tahun 2015-2017. Keduanya dipakai pada NUTS 3
  (2016).
\item
  CAP berupa gross fixed capital formation tahun 2015 pada NUTS 2 dari
  Eurostat, kemudian didisagregasi ke NUTS 3 menggunakan active
  enterprises tahun 2014 dari database ESPON. Table 2 mencantumkan unit
  teritorial NUTS 3 (2013) untuk CAP.
\item
  POLY memakai GHS-UCDB dan HERE routing API, tahun 2015, pada unit
  urban centre.
\item
  POP dan DISP berasal dari GHS-UCDB; POP adalah total penduduk UC dalam
  UR, sedangkan DISP adalah jumlah UC dalam UR.
\item
  HUM berasal dari database ESPON tahun 2016 dan mengukur jumlah orang
  yang bekerja di sektor intensif pengetahuan pada NUTS 3 (2013).
\item
  ACC berasal dari MRS ESPON tahun 2014 dan menggabungkan potential
  accessibility melalui jalan, kereta api, dan udara pada NUTS 3 (2013).
\item
  INNO berasal dari database ESPON, menggunakan jumlah paten 4.0 per
  1.000 penduduk pada rentang 2010-2015, diindeks terhadap rata-rata
  negara ESPON, pada NUTS 3 (2013).
\item
  ACTI adalah jumlah perusahaan aktif setidaknya selama sebagian periode
  rujukan, dari database ESPON tahun 2014 pada NUTS 3 (2013); variabel
  ini dipakai sebagai covariate untuk disaggregasi modal.
\item
  POLY\_HIST, POP\_HIST, dan DISP\_HIST dibangun untuk tahun 1975.
  Penulis membuat set UC historis sendiri dengan mereplikasi metodologi
  yang disebut dalam artikel, menggunakan GHS-BUILT dan GHS-POP
  historis.
\end{itemize}

\textbf{Catatan provenance:} Rincian sumber, tahun, unit, statistik
deskriptif, dan missing values tersebut berasal dari Table 2, bukan dari
pengambilan data baru. (Artikel, Table 2, hlm. 1204-1206; TXT baris
613-734.)

\textbf{Laporan sumber:} Karena CAP hanya tersedia pada NUTS 2, data
tersebut didisagregasi ke NUTS 3 menggunakan active enterprises dan
sebuah R package untuk Bayesian Spatial Disaggregation Modelling.
Setelah itu, sebagian besar variabel pada NUTS 3 dialokasikan ke UC.
Jika UC melintasi beberapa NUTS 3 atau beberapa UC berada dalam satu
NUTS 3, penulis menggunakan proporsi luas permukaan. Prosedur alokasi
dan agregasi dibedakan menurut apakah variabel berupa nilai absolut atau
relatif, lalu nilai UC dikumpulkan kembali menjadi skor UR. (TXT baris
613-633.)

\textbf{Laporan sumber:} Missing values ditangani dengan beberapa cara.
Nilai GVA Ticino diambil dari statistik nasional Swiss. Missing LAB di
wilayah sekitar Glasgow, Edinburgh, dan Ticino diestimasi dengan
expectation maximization di SPSS menggunakan POP sebagai covariate.
Nilai HUM Ticino dan Zuid-Limburg, serta missing CAP Ticino, diestimasi
menggunakan rata-rata wilayah NUTS tetangga. Table 2 mencantumkan
missing GVA dan CAP sekitar 0,5 persen, LAB 5 persen, dan HUM 1 persen;
ACC, INNO, dan ACTI tidak memiliki missing values menurut tabel. (TXT
baris 642-734.)

\textbf{Inferensi pengolahan:} Perbedaan tahun, tingkat wilayah,
prosedur disaggregasi, alokasi berbasis luas, dan imputasi berarti TFP
tidak berasal dari satu himpunan observasi yang dikumpulkan pada satu
tahun seragam. TXT tidak menyediakan raw data, kode pemrosesan, nilai
observasi per UR, atau rincian algoritme statistik di luar uraian di
atas.

\subsection{Population / Sample}\label{population-sample-49}

\textbf{Laporan sumber:} Populasi analitis berangkat dari 734 UCs yang
berada dalam cakupan 34 negara Eropa. UCs dipilih dari GHS-UCDB dengan
kriteria degree of urbanization dan kemudian disaring oleh PURban
berdasarkan jarak tempuh, jumlah minimum pusat, critical mass, dan
keseimbangan ukuran. Hasil penyaringan adalah 94 UR unik yang berlokasi
di 10 negara. (TXT baris 451-468 dan 497-512.)

\textbf{Laporan sumber:} Analisis SAR menggunakan 94 observasi UR.
Analisis GS2SLS menggunakan 79 observasi, karena sejumlah UC tidak dapat
diidentifikasi berdasarkan data tahun 1975 sehingga sebagian UR tidak
dapat masuk ke model berinstrumen. Table 3 juga mencatat kemungkinan
hubungan tumpang tindih antar-UR melalui UC yang sama. (TXT baris
741-745 dan 807-819.)

\textbf{Inferensi pengolahan:} Ini bukan sampel acak penduduk atau
survei. Sampelnya adalah himpunan wilayah yang lolos algoritme dan
ambang batas substantif, sehingga generalisasi paling aman dibatasi pada
UR yang memenuhi kondisi tersebut. Cakupan basisnya 34 negara, tetapi UR
yang memenuhi syarat hanya ditemukan di 10 negara. TXT tidak
mencantumkan daftar tekstual 10 negara tersebut, tidak menjelaskan
probabilitas pemilihan, dan tidak memberikan kerangka sampling lain.

\textbf{Batas informasi:} Sampel individu, rumah tangga, perusahaan,
atau responden survei: Tidak berlaku. Daftar atau ukuran populasi
penduduk di luar 94 UR yang gagal memenuhi kriteria: Tidak disebutkan
dalam file; Table 2 hanya menyediakan statistik deskriptif untuk
variabel yang dianalisis.

\subsection{Dependent Variables}\label{dependent-variables-49}

\textbf{Laporan sumber:} Variabel dependen utama pada model
struktur-produktivitas adalah TFP untuk UR. TFP bukan pengukuran yang
langsung tercantum sebagai data observasi; TFP merupakan residual dari
persamaan produksi yang memprediksi log GVA menggunakan log capital (K)
dan log labour (L). GVA mencakup seluruh sektor ekonomi yang ada di
wilayah PUR. (TXT baris 536-553 dan 648-650.)

\textbf{Laporan sumber:} GVA berfungsi sebagai dependent variable pada
tahap perhitungan TFP, sedangkan TFP menjadi dependent variable pada
persamaan utama. Persamaan utama juga memasukkan spatial lag
\texttt{W\ TFP}, sehingga produktivitas UR diperlakukan sebagai memiliki
ketergantungan terhadap produktivitas UR lain yang berbagi UC. (TXT
baris 537-565.)

\textbf{Interpretasi penulis:} TFP dipilih karena dianggap menangkap
technological sophistication dan production efficiency dengan lebih
tepat untuk membedakan perkembangan ekonomi regional Eropa dibandingkan
hanya mengandalkan GDP per capita. Artikel mengaitkan pilihan ini dengan
literatur yang menekankan peran technological progress, innovation, dan
knowledge externalities dalam produktivitas. (TXT baris 407-416.)

\textbf{Inferensi pengolahan:} Karena TFP adalah residual yang
bergantung pada GVA, K, L, disaggregasi, imputasi, dan spesifikasi
persamaan pertama, koefisien POLY atau DISP tidak dapat dibaca sebagai
perubahan langsung pada output yang terukur. Unit, tahun data, nilai
missing, dan statistik deskriptif khusus TFP ditandai \texttt{/} dalam
Table 2 dan tidak tersedia dalam TXT. (TXT baris 642-650.)

\subsection{Results}\label{results-49}

\textbf{Laporan sumber:} Diagnostics menyatakan bahwa instrumen relevan
dan variabel independen utama eksogen kecuali POP. POP karena itu
diinstrumentasikan hanya pada GS2SLS. POLY dan DISP tidak menunjukkan
masalah yang sama dan dimasukkan dalam SAR. (TXT baris 739-745.)

\textbf{Laporan sumber:} Koefisien spasial \texttt{Rho} bernilai positif
dan signifikan pada seluruh model menurut uraian penulis, dengan nilai
sekitar 0,74 pada model SAR dan sekitar 0,68-0,69 pada model GS2SLS.
(TXT baris 746-752; Table 3, TXT baris 806-812.)

\textbf{Interpretasi penulis:} Penulis menafsirkan \texttt{Rho} sebagai
bukti bahwa produktivitas UR saling bergantung secara spasial dan
memiliki efek eksternalitas positif. (TXT baris 746-752.)

\textbf{Laporan sumber:} Hasil untuk variabel struktur spasial adalah
sebagai berikut:

\begin{itemize}
\tightlist
\item
  Pada SAR model (1), POLY memiliki koefisien negatif sekitar -0,22
  dengan standard error 0,08 dan signifikan pada tingkat 5 persen.
\item
  Pada SAR model (2), DISP memiliki koefisien negatif sekitar -0,02
  dengan standard error 0,00 dan signifikan pada tingkat 1 persen.
\item
  Ketika DISP ditambahkan pada SAR model (3) dan (4), koefisien POLY
  tetap negatif sekitar -0,10 tetapi tidak signifikan, sedangkan DISP
  tetap negatif dan signifikan pada tingkat yang lebih lemah sebagaimana
  tanda pada tabel.
\item
  Dalam GS2SLS model (6) dan (8), POLY kembali memiliki koefisien
  negatif dan signifikan, masing-masing sekitar -0,48 dan -0,56.
  Perbedaan spesifikasi dan jumlah observasi perlu diperhatikan saat
  membandingkan angka ini dengan SAR.
\item
  Dalam GS2SLS model (7) dan (9), DISP memiliki koefisien negatif
  sekitar -0,04 dan signifikan.
\item
  POP tidak signifikan pada semua model menurut uraian penulis.
\end{itemize}

\textbf{Catatan provenance:} Angka dan tanda signifikansi diringkas dari
Table 3 dan uraian Bagian 4; simbol minus pada ekstraksi tabel beberapa
kali muncul sebagai karakter nonstandar. (TXT baris 753-776 dan Table 3,
baris 791-822.)

\textbf{Laporan sumber:} Tidak ada interaksi yang signifikan kecuali
DISP\emph{POP pada model (9). Koefisien interaksi tersebut mendekati nol
dan negatif, sekitar -0,03 dengan standard error 0,01. POLY}POP dan
POLY*DISP tidak signifikan. (TXT baris 769-776 dan Table 3, baris
803-805.)

\textbf{Interpretasi penulis:} Koefisien DISP*POP dipahami berarti efek
negatif DISP sedikit lebih kuat pada UR dengan populasi lebih besar.
(TXT baris 769-772.)

\textbf{Laporan sumber:} HUM dan ACC memiliki pengaruh positif dalam
uraian penulis. ACC konsisten positif dan signifikan dalam Table 3,
dengan koefisien sekitar 0,25-0,43 menurut model. HUM disebut berdampak
positif, tetapi tingkat signifikansinya tidak konsisten
antarspesifikasi. INNO dinyatakan penulis tidak memengaruhi TFP dalam
model mana pun. (TXT baris 777-787 dan Table 3, baris 800-802.)

\textbf{Catatan ketidakselarasan internal:} Baris INNO pada hasil
ekstraksi Table 3 tampak menampilkan koefisien -0,03 dengan standard
error 0,01 dan dua tanda bintang pada model (2), sedangkan narasi Bagian
4 menyatakan INNO tidak berdampak pada TFP di semua model. TXT juga
mengandung artefak layout dan simbol. Karena hanya TXT yang boleh
digunakan, ketidakselarasan ini tidak diputuskan secara sepihak; hasil
INNO perlu dianggap tidak pasti pada level tabel-versus-narasi.

\textbf{Laporan sumber:} GS2SLS memiliki 79 observasi dan nilai R2
sekitar 0,49-0,52. Uji alternatif DISP berupa median jarak antara UC
pusat dan UC lain tidak signifikan pada model mana pun. (TXT baris
809-822 dan 858-864.)

\textbf{Interpretasi penulis:} Penulis menghubungkan hasil alternatif
tersebut dengan praseleksi UR yang sudah dibatasi oleh waktu tempuh
jalan 45 menit. (TXT baris 858-864.)

\subsection{Findings}\label{findings-49}

\textbf{Interpretasi penulis:} Temuan utama adalah tidak adanya bukti
empiris untuk dampak positif polycentricity terhadap produktivitas
regional Eropa. POLY berasosiasi negatif ketika berdiri sendiri atau
dikombinasikan dengan POP pada spesifikasi tertentu, tetapi pengaruhnya
menjadi tidak signifikan ketika DISP dimasukkan. Karena DISP merupakan
prediktor negatif yang lebih konsisten, penulis menganggap tingkat
persebaran populasi antar-UC sangat penting dalam menjelaskan TFP. (TXT
baris 753-761 dan 839-852.)

\textbf{Interpretasi penulis:} POP yang tidak signifikan dipahami
sebagian sebagai konsekuensi praseleksi. Semua UR sudah memiliki minimum
1.5 juta penduduk, yang menurut logika penulis menempatkan UR dalam
kondisi yang cukup mendukung agglomeration externalities; perbedaan
antara 1,6 dan 1,9 juta mungkin tidak cukup membedakan hasil dalam
sampel ini. (TXT baris 762-768.)

\textbf{Interpretasi penulis:} Hasil interaksi DISP\emph{POP menunjukkan
bahwa dampak negatif dispersi sedikit lebih kuat pada wilayah yang lebih
berpenduduk, sementara tidak ditemukannya interaksi POP}POLY berbeda
dari temuan Ouwehand et al.~(2022) dan Meijers dan Burger (2010).
Spatial \texttt{Rho} yang positif diinterpretasikan sebagai externality
produktivitas antar-UR. (TXT baris 769-776 dan 746-752.)

\textbf{Interpretasi penulis:} Dampak positif HUM dianggap masuk akal
karena sektor intensif pengetahuan berkaitan dengan industri bernilai
dan bertumbuh tinggi. Dampak positif ACC dipahami melalui aksesibilitas
eksternal Eropa, meskipun indikator tersebut lebih banyak
merepresentasikan kemudahan mencapai wilayah Eropa lain daripada
efisiensi logistik internal. (TXT baris 777-787.)

\textbf{Inferensi pengolahan:} Bukti yang paling stabil dalam artikel
adalah asosiasi negatif DISP-TFP, bukan bukti interaksi positif atau
negatif yang kuat antara POLY dan DISP. Oleh karena itu, kalimat abstrak
bahwa wilayah lebih polisentris dan lebih tersebar memiliki
produktivitas lebih rendah harus dibaca bersama hasil bahwa efek POLY
dapat hilang setelah DISP dimasukkan dan bahwa interaksi POLY*DISP tidak
signifikan.

\subsection{Conclusions}\label{conclusions-49}

\textbf{Laporan sumber:} Penulis menyimpulkan bahwa analisis 94 UR Eropa
tidak menemukan bukti empiris untuk dampak positif polycentricity
terhadap produktivitas regional. POLY berpengaruh negatif atau tidak
berpengaruh, bergantung pada variabel struktur spasial tambahan,
sedangkan DISP berpengaruh negatif pada berbagai model. Gabungan hasil
itu dipakai untuk menunjukkan pentingnya manfaat aglomerasi berpusat
tunggal. (Artikel, Bagian 5, hlm. 1209; TXT baris 837-864.)

\textbf{Interpretasi penulis:} Penulis menilai pendekatan yang
memodelkan beberapa karakteristik struktur spasial secara bersama-sama
lebih memadai daripada menjadikan polycentricity sebagai prediktor
tunggal. Mereka juga melihat kelebihan delineasi bottom-up berbasis UC
karena dapat menangkap hubungan lintas batas administratif, misalnya
hubungan Brussels dengan pusat-pusat di wilayah Flanders dan Wallonia.
(TXT baris 865-876.)

\textbf{Laporan sumber:} Penulis sekaligus menyatakan bahwa delineasi
bottom-up menyulitkan rekomendasi kebijakan yang sangat spesifik, karena
UR yang terbentuk tidak selalu bertepatan dengan unit pengambilan
keputusan ekonomi regional. Model spasial kuantitatif dapat
menggambarkan pola umum skala Eropa, tetapi bukan pengganti studi kasus
mendalam. (TXT baris 881-910.)

\textbf{Interpretasi penulis:} Fokus morfologis tidak menutup
kemungkinan hasil berbeda bila polycentricity didekati secara
fungsional. Penulis mengutip kritik bahwa ukuran demografis dan
rank-size terutama menggambarkan hierarki kota, bukan relasi efektif
antar kota, tetapi tetap menggunakan ukuran morfologis karena data
relasional Eropa yang sesuai belum tersedia. (TXT baris 911-926.)

\textbf{Inferensi pengolahan:} Kesimpulan artikel bersifat kontekstual:
ia melemahkan narasi bahwa polycentricity secara otomatis meningkatkan
TFP pada UR yang memenuhi kriteria studi, tetapi tidak menyelesaikan
pertanyaan apakah integrasi fungsional, komplementaritas, atau desain
kelembagaan tertentu dapat menghasilkan efek yang berbeda.

\subsection{Contributions}\label{contributions-49}

\textbf{Laporan sumber:} Kontribusi yang dinyatakan penulis meliputi:

\begin{itemize}
\tightlist
\item
  Menerapkan penilaian yang lebih robust, komparatif, dan reproducible
  terhadap kondisi spasial PUR di Eropa dengan membangun analisis pada
  PURban, perangkat lunak sumber terbuka yang menggabungkan beberapa
  kerangka pengukuran dan dataset polycentricity morfologis.
\item
  Menggunakan TFP sebagai ukuran produktivitas ekonomi dan bukan hanya
  GDP per capita.
\item
  Menggunakan territorial level yang lebih fine-grained, terutama NUTS
  3, lalu menghubungkannya dengan UC dan UR melalui prosedur yang dapat
  menangani UC yang melintasi batas administratif.
\item
  Memperluas cakupan geografis hingga 34 negara dalam basis analisis,
  dengan hasil UR berada di 10 negara, dan menyatakan cakupan tersebut
  lebih dari dua kali jumlah negara dalam Ouwehand et al.~(2022).
\item
  Memasukkan ukuran, dispersi, polycentricity, spatial regression, dan
  instrumental variables untuk memeriksa ketergantungan spasial serta
  potensi bidirectionality.
\end{itemize}

\textbf{Catatan provenance:} Dasar untuk daftar ini terdapat pada
pendahuluan, proposed approach, dan ringkasan rancangan metode. (TXT
baris 122-130, 388-427, dan 431-449.)

\textbf{Inferensi pengolahan:} Kontribusi artikel terutama berada pada
integrasi pengukuran morfologis, penyelarasan skala, dan pemodelan
spasial. Kontribusi tersebut memperbaiki desain perbandingan, tetapi
tidak menghapus ketergantungan pada proxy morfologis, alokasi data, dan
asumsi instrumen; artikel tidak mengklaim telah membuktikan mekanisme
kausal secara definitif.

\subsection{Research Gap}\label{research-gap-49}

\textbf{Laporan sumber:} Gap yang diidentifikasi pada awal artikel
adalah:

\begin{itemize}
\tightlist
\item
  Bukti hubungan polycentricity-produktivitas masih fragmented, uneven,
  scarce, dan contradictory meskipun polycentric development kuat dalam
  kebijakan Eropa.
\item
  Masih sedikit studi yang menguji hubungan tersebut secara sistematis
  pada skala seluruh Eropa.
\item
  Literatur lebih sering bekerja pada skala intra-urban daripada
  interurban atau regional, padahal skala regional relevan bagi agenda
  pembangunan kota dan produktivitas Uni Eropa.
\item
  Tidak ada common and shared understanding tentang definisi dan
  pengukuran polycentricity, sehingga perbedaan subindikator, indeks,
  unit, dan skala memengaruhi hasil.
\item
  Ukuran produktivitas terdahulu sering memakai GDP per capita karena
  ketersediaan data; penggunaan TFP pada skala regional yang lebih halus
  masih perlu dikembangkan.
\item
  Pendekatan yang membatasi kota di dalam batas administratif berpotensi
  mengabaikan koneksi lintas batas yang relevan bagi produktivitas dan
  polycentric development.
\end{itemize}

\textbf{Catatan provenance:} Pernyataan tersebut bersumber dari Bagian
1, Table 1, dan proposed approach. (TXT baris 107-130, 182-188, 293-312,
dan 394-427.)

\textbf{Gap yang tersisa menurut sumber:} Penulis tetap mengidentifikasi
kebutuhan akan ukuran DISP yang berbeda, data sosial-ekonomi langsung
pada tingkat UC, dan pendekatan fungsional yang memuat konektivitas,
interaksi, komplementaritas, serta tindakan institusional. Data arus
komuter atau hubungan perusahaan pada skala seluruh Eropa disebut belum
memadai dalam TXT. (TXT baris 858-864 dan 894-926.)

\textbf{Inferensi pengolahan:} Artikel menutup sebagian gap
komparabilitas dan spasial, tetapi gap fungsional-relasional dan gap
penyelarasan data belum selesai. Kesimpulan tentang produktivitas karena
itu masih melekat pada polycentricity morfologis dan UR yang dihasilkan
algoritme.

\subsection{Limitations}\label{limitations-49}

\textbf{Laporan sumber:} Keterbatasan yang dinyatakan atau dijelaskan
langsung oleh artikel meliputi:

\begin{itemize}
\tightlist
\item
  Waktu tempuh untuk menentukan kedekatan antar-UC hanya menggunakan
  perjalanan dengan mobil. Catatan kaki penulis menyatakan hal ini dapat
  menimbulkan bias di wilayah yang memiliki proporsi komuter kereta api
  yang signifikan. (TXT baris 477-481 dan catatan kaki baris 521-525.)
\item
  Data ekonomi dan kontrol pada NUTS 3 harus dialokasikan ke UC yang
  bentuknya tidak selalu berimpit dengan NUTS 3. Alokasi berbasis
  proporsi luas mengasumsikan bahwa seluruh aktivitas ekonomi yang
  terukur di NUTS 3 terkonsentrasi di dalam UC.
\item
  Asumsi alokasi tersebut dapat menyebabkan upward bias pada estimasi
  tingkat UC karena rata-rata 40 persen penduduk NUTS 3 dalam analisis
  berada di luar UC. Penulis mengusulkan nighttime light atau population
  shares sebagai cara untuk mengurangi masalah ini. (TXT baris 881-893.)
\item
  UR bottom-up dan lintas batas tidak selalu sesuai dengan unit
  administratif tempat keputusan ekonomi regional dibuat, sehingga
  rekomendasi kebijakan yang tailored sulit diturunkan.
\item
  Polycentricity yang digunakan bersifat morfologis. Ukuran ini tidak
  secara langsung mengukur arus komuter, hubungan perusahaan, interaksi,
  atau fungsi kelembagaan. Penulis menyebut kritik bahwa rank-size hanya
  indikator tidak langsung dari sistem relasional. (TXT baris 894-926.)
\item
  Analisis spasial kuantitatif memberi gambaran luas, bukan detail studi
  kasus untuk wilayah tertentu. (TXT baris 906-910.)
\end{itemize}

\textbf{Laporan sumber:} Artikel juga mengakui potensi bidirectionality
antara produktivitas dan struktur spasial. Diagnostics memperlakukan POP
berbeda dari POLY dan DISP karena POP tidak eksogen dalam model SAR;
hanya POP yang diinstrumentasikan dalam GS2SLS. Validitas instrumen
bergantung pada relevansi dan eksogenitas historical population,
polycentricity, dan dispersion. (TXT baris 594-612 dan 739-745.)

\textbf{Inferensi pengolahan:} Ada batasan desain tambahan yang terlihat
dari TXT tetapi tidak diberi label sebagai daftar limitations formal: 94
UR berasal dari praseleksi minimum 1.5 juta penduduk dan 45 menit; hasil
berasal dari 10 negara meskipun basisnya 34 negara; tahun data berbeda
antarvariabel; missing values diimputasi; dan model GS2SLS berkurang
menjadi 79 observasi. Hal-hal ini membatasi komparabilitas dan
generalisasi, tetapi tidak boleh diperlakukan sebagai ukuran bias yang
sudah dihitung penulis.

\textbf{Batas informasi:} TXT tidak memuat nilai diagnostics lengkap di
Appendix S2-S3, analisis sensitivitas penuh, raw data, kode PURban atau
kode estimasi, sehingga kekuatan numerik setiap uji robustness tidak
dapat dinilai lebih jauh dari yang dilaporkan pada teks dan Table 3.

\subsection{Challenges}\label{challenges-49}

\textbf{Laporan sumber:} Artikel menghadapi tantangan konseptual karena
polycentricity memiliki banyak definisi dan indikator. Studi terdahulu
yang sama-sama memakai label morfologis masih menghasilkan kesimpulan
berbeda karena mengubah unit, skala, subindikator, atau ukuran
produktivitas. (TXT baris 293-312.)

\textbf{Laporan sumber:} Tantangan data dan spasialnya mencakup:

\begin{itemize}
\tightlist
\item
  Mengidentifikasi UR yang tidak dibatasi wilayah administratif, dapat
  tumpang tindih, dan dapat melintasi batas negara atau wilayah
  administratif.
\item
  Menghubungkan UC dengan NUTS 3 ketika satu UC berada di beberapa NUTS
  3 atau satu NUTS 3 berisi beberapa UC.
\item
  Mendisagregasi CAP dari NUTS 2 ke NUTS 3 dan mengimputasi missing
  values pada beberapa wilayah.
\item
  Membangun instrumen historis ketika UC tahun 1975 tidak selalu dapat
  diidentifikasi, yang menyebabkan berkurangnya sampel GS2SLS.
\item
  Menangani spatial autocorrelation dan memilih SAR setelah pengujian
  beberapa spesifikasi.
\item
  Memisahkan kemungkinan arah kausal struktur-produktivitas dari
  endogeneity POP.
\end{itemize}

\textbf{Catatan provenance:} Sumber tantangan ini terdapat pada Bagian
3, Table 2, catatan kaki, dan Bagian 4. (TXT baris 518-535, 594-633,
638-641, 732-734, dan 741-745.)

\textbf{Interpretasi penulis:} Tantangan paling substantif adalah
mengukur kondisi yang memungkinkan aglomerasi regional tanpa data
relasional Eropa yang memadai. Karena itu, artikel memilih proxy
morfologis yang lebih konsisten secara geografis, tetapi mengakui bahwa
proxy tersebut tidak identik dengan hubungan aktual antar kota. (TXT
baris 394-406 dan 911-926.)

\textbf{Inferensi pengolahan:} Tantangan teknis dan konseptual saling
terkait. Semakin bottom-up unit wilayahnya, semakin besar kebutuhan
untuk melakukan alokasi data dari unit administratif; alokasi itu justru
dapat memengaruhi TFP yang hendak dijelaskan. Ini adalah inferensi dari
prosedur dan asumsi sumber, bukan pernyataan kausal tambahan dari
penulis.

\subsection{Practical Implications}\label{practical-implications-49}

\textbf{Interpretasi penulis:} Hasil tidak mendukung penggunaan
polycentric development sebagai asumsi umum bahwa integrasi beberapa
pusat otomatis meningkatkan competitiveness regional. Dalam terjemahan
kebijakan yang eksplisit, penulis menyatakan ada bukti berlanjutnya
pentingnya manfaat aglomerasi single-centre. (TXT baris 837-847 dan
907-910.)

\textbf{Interpretasi penulis:} Penilaian kebijakan sebaiknya tidak
memperlakukan POLY sebagai satu-satunya karakteristik struktur spasial.
Ukuran populasi dan terutama dispersi perlu dipertimbangkan bersama,
karena DISP muncul sebagai prediktor negatif yang lebih kuat atau lebih
konsisten daripada POLY dalam spesifikasi artikel. (TXT baris 848-864.)

\textbf{Laporan sumber:} Delineasi bottom-up dapat membantu mengenali
geografi urban lintas batas yang mungkin terlewat bila analisis hanya
memakai OECD TL2 atau batas administratif lain. Namun, unit bottom-up
tersebut tidak selalu bertepatan dengan unit tempat kebijakan diambil,
sehingga penulis membatasi implikasi kebijakannya pada gambaran umum
Eropa dan tidak menawarkan rekomendasi tailored untuk setiap UR. (TXT
baris 865-910.)

\textbf{Inferensi pengolahan:} Implikasi praktis yang aman adalah
kehati-hatian dalam mengubah indeks morfologis menjadi resep kebijakan.
Hasil ini dapat dipakai untuk menguji kembali klaim bahwa polycentricity
merupakan panacea ekonomi, tetapi tidak cukup untuk menentukan apakah
suatu wilayah harus memilih struktur monosentris atau polisentris tanpa
data fungsi, tata kelola, dan konteks lokal.

\subsection{Applications}\label{applications-49}

\textbf{Laporan sumber:} Aplikasi analitis yang benar-benar dilakukan
adalah penggunaan PURban untuk mengidentifikasi, memetakan, dan
menganalisis derajat polycentricity morfologis; membangun 94 UR;
menghubungkan struktur UR dengan data NUTS 3; serta mengestimasi
hubungan dengan TFP melalui SAR dan GS2SLS. PURban disebut sebagai
open-source software tool, tetapi TXT tidak mencantumkan alamat
repositori atau instruksi penggunaannya. (TXT baris 388-406, 451-512,
dan 516-603.)

\textbf{Laporan sumber:} Penulis melihat pendekatan degree of
urbanization dan delineasi non-administratif berpotensi berguna untuk
perbandingan antarnegara dan kebijakan yang menargetkan geografi urban
tertentu, karena pendekatan tersebut memungkinkan perbandingan ``likes
with likes'' serta menangkap koneksi lintas batas. (TXT baris 894-903.)

\textbf{Tidak disebutkan dalam file:} TXT tidak melaporkan implementasi
kebijakan yang menggunakan hasil studi, evaluasi program, perubahan
regulasi, atau dampak operasional pada salah satu dari 94 UR. Tidak ada
aplikasi eksperimental atau studi kasus kebijakan individual yang dapat
ditambahkan tanpa sumber lain.

\textbf{Inferensi pengolahan:} Aplikasi paling langsung adalah sebagai
kerangka diagnosis komparatif, bukan sebagai alat prediksi kebijakan
wilayah tertentu. Penggunaan untuk keputusan spesifik memerlukan
kehati-hatian karena unit analisis, data fungsi, dan unit pemerintahan
tidak selalu sama.

\subsection{Future Research
(Optional)}\label{future-research-optional-49}

\textbf{Laporan sumber:} Agenda riset lanjut yang disebut eksplisit
adalah:

\begin{itemize}
\tightlist
\item
  Menguji ukuran dispersi yang berbeda, termasuk ukuran yang
  memperhitungkan jarak geografis antar-UC. Artikel melaporkan bahwa
  median jarak dari UC pusat ke UC lain tidak signifikan dalam modelnya,
  tetapi tetap menganjurkan eksplorasi lebih lanjut terhadap ukuran
  DISP.
\item
  Mengalokasikan data NUTS 3 dengan nighttime light atau population
  shares, bukan hanya surface area, untuk mengurangi bias dari aktivitas
  dan penduduk yang berada di luar UC.
\item
  Mengembangkan dataset sosial-ekonomi yang lebih luas langsung pada
  tingkat wilayah yang konsisten dengan GHSL/UC sehingga post hoc
  imputation dari NUTS 3 dapat dikurangi.
\item
  Menguji pendekatan fungsional terhadap polycentricity yang mencakup
  connectivity, interaction, complementarity, dan institutional action;
  data commuting flows atau hubungan antarperusahaan disebut sebagai
  kebutuhan yang belum memadai pada skala pan-Eropa.
\end{itemize}

\textbf{Catatan provenance:} Sumber agenda tersebut adalah TXT baris
858-864, 881-903, dan 911-926.

\textbf{Tidak disebutkan dalam file:} Tidak ada agenda eksplisit
mengenai tahun observasi tertentu, negara tambahan, desain longitudinal,
eksperimen kebijakan, atau metode selain agenda yang dirangkum di atas.

\subsection{Provenance Notes}\label{provenance-notes-49}

\textbf{Bahan dan batas akses:} Review ini hanya memakai
\texttt{/Users/mac/Documents/Mac/{[}2{]}\ Obsidian\ Vault/zettelkasten/source/journal/B06-caset-yang-derudder-samardzhiev-2023-the-productivity-effects-of-polycentricity-a-systematic-analysis-of-urban-regions-in-europe-10.1111-pirs.12765.txt}.
Tidak ada informasi yang ditambahkan dari PDF, situs DOI, referensi yang
dikutip, supporting information, atau sumber eksternal. Referensi pada
Bagian 2 diringkas hanya sebagaimana dipresentasikan oleh artikel.

\textbf{Verifikasi file aktual:} Pencocokan prefiks \texttt{B06}
terhadap TXT di bawah \texttt{source/} menemukan tepat satu file, yaitu
file pada path di atas. File tersebut berasosiasi dengan PDF yang
disebut pada baris 2 TXT dan dibaca seluruhnya.

\textbf{Metadata ekstraksi:} Blok metadata TXT menyebut PDF sumber,
tanggal ekstraksi 2026-08-31, metode \texttt{pdftotext\ -layout}, 22
halaman, tidak terenkripsi, dan ukuran PDF 679734 bytes. Teks yang
dibaca mencakup baris 1-1125, termasuk metadata, isi artikel, Table 1-3,
catatan kaki, referensi, supporting information notice, serta ringkasan
bahasa Spanyol dan Jepang. (TXT baris 1-29 dan 1093-1125.)

\textbf{Konvensi locator:} Locator utama memakai nomor baris TXT hasil
ekstraksi. Jika tersedia, locator juga menyebut nomor halaman artikel
yang tercetak, bagian, table, atau figure. Nomor halaman PDF metadata
tidak dipakai untuk mengoreksi pagination artikel. Keterangan sitasi
artikel menyebut halaman 1193-1213, sedangkan field Subject pada
metadata awal menyebut 1193-1214; konflik ini tetap dinyatakan.

\textbf{Artefak ekstraksi:} Layout dua kolom dan tabel menghasilkan
pemisahan kolom yang tidak selalu stabil. Simbol minus pada beberapa
koefisien muncul sebagai karakter nonstandar, beberapa kata terbelah
oleh tanda hubung, dan Table 3 memiliki artefak pada simbol
signifikansi. Karena itu, hasil utama diringkas dengan mengutamakan
tanda dan uraian naratif yang terbaca, sementara ketidakselarasan INNO
antara Table 3 dan narasi dicatat terbuka pada bagian Results.

\textbf{Informasi yang tidak tersedia:} TXT tidak memuat daftar lengkap
34 negara dan daftar tekstual 10 negara yang menghasilkan UR; raw data;
kode PURban atau kode estimasi; isi rinci Appendix S1-S3; angka lengkap
diagnostics instrumen, LM, dan likelihood ratio; rincian implementasi
label ``Region Fixed''; nilai observasi setiap UR; protokol systematic
literature search; atau verifikasi status peer review. Tidak ada dugaan
yang digunakan untuk mengisi kekosongan tersebut.

\textbf{Status evidensi:} Pernyataan bertanda \texttt{Laporan\ sumber}
adalah isi atau hasil yang dilaporkan artikel.
\texttt{Interpretasi\ penulis} merujuk pada penjelasan atau implikasi
yang secara eksplisit diberikan penulis. \texttt{Inferensi\ pengolahan}
adalah pembacaan hati-hati dalam review ini dan tidak boleh diperlakukan
sebagai kutipan atau klaim penulis.

\section{Review Kode B07}\label{review-kode-b07}

\textbf{Identitas sumber}

\begin{itemize}
\tightlist
\item
  Kode: B07.
\item
  Judul: ``Summing Small Cities Does Not Make a Large City: Polycentric
  Urban Regions and the Provision of Cultural, Leisure and Sports
  Amenities''.
\item
  Penulis: Evert Meijers.
\item
  Jurnal: \emph{Urban Studies}, volume 45, nomor 11, halaman 2323-2342,
  October 2008.
\item
  DOI: \texttt{10.1177/0042098008095870}.
\item
  URL yang tercetak dalam TXT:
  \texttt{http://usj.sagepub.com/content/45/11/2323}.
\item
  Basis review: seluruh TXT B07 dibaca, termasuk blok metadata PDF dan
  seluruh teks hasil ekstraksi. Tidak ada sumber di luar TXT yang
  digunakan.
\end{itemize}

\textbf{Penanda evidensi:} \texttt{{[}Laporan\ sumber{]}} berarti isi
yang dinyatakan, diukur, atau dicetak dalam artikel;
\texttt{{[}Interpretasi\ penulis{]}} berarti penjelasan atau penilaian
yang diberikan Meijers; \texttt{{[}Inferensi\ pengolahan{]}} berarti
pembacaan analitis dari isi TXT, bukan klaim yang secara eksplisit
dinyatakan penulis.

\subsection{Summarized Abstract}\label{summarized-abstract-50}

{[}Laporan sumber{]} Artikel menguji apakah \emph{polycentric urban
region} (PUR), meskipun terdiri atas beberapa kota yang terpisah secara
spasial, dapat memperoleh keuntungan dari ukuran gabungannya dengan
tingkat yang serupa dengan wilayah kota yang berukuran sama tetapi
monosentris. Pertanyaan tersebut dibatasi pada penyediaan amenitas
budaya, rekreasi, dan olahraga. Kehadiran amenitas di 42 wilayah Belanda
dinyatakan dalam sebuah indeks yang digunakan sebagai variabel dependen
dalam model regresi berganda. Tingkat polisentrisitas menjadi salah satu
variabel penjelas, dengan pengendalian terhadap ukuran penduduk, jumlah
pengunjung, dan pendapatan rata-rata. Hasil yang diringkas abstrak
menunjukkan hubungan negatif: semakin polisentris suatu wilayah, semakin
sedikit amenitas budaya, rekreasi, dan olahraga yang tersedia; wilayah
yang lebih monosentris memiliki lebih banyak amenitas. {[}TXT, baris
88-98; artikel, hlm. 2323{]}

{[}Interpretasi penulis{]} Temuan tersebut dipakai untuk menolak asumsi
sederhana bahwa menjumlahkan beberapa kota kecil atau menengah otomatis
menghasilkan kapasitas amenitas seperti satu kota besar. {[}TXT, baris
815-831 dan 834-921; artikel, hlm. 2337-2340{]}

{[}Inferensi pengolahan{]} Abstrak dan analisis menguji hubungan
antarwilayah, bukan dampak kausal dari polisentrisitas terhadap
amenitas. Karena itu, kesimpulan paling aman adalah bahwa bentuk sistem
perkotaan berasosiasi dengan indeks amenitas dalam kasus dan cakupan
yang diteliti, bukan bahwa polisentrisitas secara mandiri menyebabkan
berkurangnya amenitas.

\subsection{Problem Statement}\label{problem-statement-50}

{[}Laporan sumber{]} Kebijakan pembangunan regional dan perencanaan
spasial di Eropa sering mempromosikan kumpulan kota yang berdekatan dan
saling terhubung sebagai \emph{city networks}, \emph{urban networks},
atau PUR. Pembuat kebijakan mengasumsikan bahwa kota-kota kecil dan
menengah yang bertindak bersama dapat membuka peluang pertumbuhan
regional dan ``bermain di liga utama'' kompetisi nasional atau
internasional. Asumsi tersebut kadang dikaitkan dengan peningkatan
integrasi fungsional dan perluasan ruang hidup perkotaan. Namun,
rasionalitas fungsional dan realitas ekonominya dinyatakan tidak selalu
diterima. {[}TXT, baris 100-141; artikel, hlm. 2323-2324{]}

{[}Laporan sumber{]} Masalah substantifnya adalah apakah ukuran penduduk
yang tersebar pada beberapa kota menghasilkan basis dukungan yang sama
dengan ukuran penduduk yang terkonsentrasi pada satu kota. Amenitas
tingkat tinggi memerlukan massa kritis, baik dari sisi permintaan maupun
pasokan modal manusia. Artikel merumuskan masalahnya dengan perbandingan
tiga kota berdekatan yang masing-masing berpenduduk sekitar 100.000
dengan satu kota berpenduduk 300.000: apakah ketiganya dapat
mengorganisasi dukungan yang sebanding bagi amenitas perkotaan? {[}TXT,
baris 167-200 dan 202-228; artikel, hlm. 2324-2326{]}

{[}Interpretasi penulis{]} Di satu sisi, PUR diperkirakan dapat
menikmati ekonomi skala, cakupan, dan kompleksitas tanpa menanggung
seluruh kemacetan, kelangkaan lahan, harga tanah tinggi, kriminalitas,
dan polusi yang terkait dengan aglomerasi besar. Di sisi lain, manfaat
tertentu mungkin bergantung pada kepadatan, kedekatan, kontak tatap
muka, interaksi informal, dan lingkungan metropolitan yang tidak
otomatis tersedia dalam struktur kota yang tersebar. {[}TXT, baris
124-141 dan 219-240; artikel, hlm. 2324-2326{]}

\subsection{Objectives}\label{objectives-50}

{[}Laporan sumber{]} Tujuan utama artikel adalah mengembangkan pengujian
empiris atas klaim bahwa keuntungan aglomerasi dapat terbentuk pada
skala PUR, dengan membandingkan wilayah perkotaan yang relatif
polisentris dan relatif monosentris. Fokus operasionalnya adalah
keberadaan amenitas budaya, rekreasi, dan olahraga yang membutuhkan
dukungan supralokal atau massa kritis. {[}TXT, baris 142-166 dan
272-280; artikel, hlm. 2324-2326{]}

{[}Laporan sumber{]} Tujuan operasional yang terlihat dari rancangan
penelitian adalah sebagai berikut:

\begin{itemize}
\tightlist
\item
  Menginventarisasi kuantitas dan kualitas amenitas pada 42 WGR-region
  di Belanda.
\item
  Membentuk indeks kehadiran amenitas untuk enam kategori dan satu
  indeks total.
\item
  Mengukur derajat mono/polisentrisitas melalui kemiringan garis regresi
  distribusi rank-size lima kota terbesar di tiap wilayah.
\item
  Menghubungkan indeks amenitas dengan indikator polisentrisitas melalui
  regresi berganda sambil mengendalikan penduduk, populasi nonpermanen
  yang diproksikan dengan tempat tidur hotel terisi, dan pendapatan
  rumah tangga rata-rata.
\item
  Memeriksa ketahanan hasil pengukuran polisentrisitas dengan indikator
  yang dihitung dari tiga kota terbesar.
\item
  Menguji korelasi antara polisentrisitas dan jenis amenitas individual
  setelah jumlahnya dinyatakan relatif terhadap penduduk.
\end{itemize}

{[}Inferensi pengolahan{]} TXT tidak menyajikan hipotesis formal
bernomor. Ekspektasi arah yang dinyatakan secara eksplisit terutama
berlaku untuk ukuran penduduk, populasi pengunjung, dan pendapatan, yang
diperkirakan berhubungan positif dengan amenitas; arah hubungan
polisentrisitas justru menjadi pertanyaan empiris utama. {[}TXT, baris
654-691; artikel, hlm. 2334{]}

\subsection{Literature Survey}\label{literature-survey-50}

{[}Laporan sumber{]} Telaah literatur artikel berbentuk narasi teoritis
dan konseptual. Artikel menempatkan PUR dalam sejarah gagasan tentang
kumpulan kota yang berdekatan, dari pembahasan sejak 1960-an hingga
meningkatnya perhatian kebijakan dan akademik pada dekade sebelum
artikel terbit. Literatur yang diringkas meliputi asal-usul dan karakter
PUR, agenda penelitian, penerapan pada wilayah tertentu, kapasitas
kelembagaan, pemerintahan, serta pembangunan spasial polisentris.
{[}TXT, baris 136-166; artikel, hlm. 2324{]}

{[}Laporan sumber{]} Poros teoritis pertama adalah massa kritis dan
ekonomi aglomerasi. Amenitas langka atau tingkat tinggi memerlukan
ukuran pasar minimum, variasi permintaan, dan sumber daya manusia yang
memadai. Literatur yang dikutip dipakai untuk menjelaskan mengapa
aktivitas ekonomi bernilai tinggi dan fungsi budaya yang beragam
biasanya terkonsentrasi di aglomerasi terbesar. {[}TXT, baris 202-228;
artikel, hlm. 2325-2326{]}

{[}Laporan sumber{]} Poros kedua adalah gagasan \emph{borrowed size} dan
eksternalitas regional atau jaringan perkotaan. Alonso dipakai untuk
merujuk pada situasi ketika kota-kota yang berdekatan dan terhubung
dapat menampung fungsi yang biasanya hanya ada di kota besar karena
basis dukungannya dianggap lebih luas secara gabungan. Capello dan Parr
digunakan untuk membedakan eksternalitas jaringan atau regional dari
keuntungan yang melekat pada satu aglomerasi. {[}TXT, baris 174-218;
artikel, hlm. 2325-2326{]}

{[}Laporan sumber{]} Literatur yang sama juga memuat keberatan terhadap
penyamaan sederhana antara integrasi beberapa kota dan aglomerasi
monosentris. Parr dikutip mengenai arus perjalanan, komoditas, dan
informasi yang lebih panjang atau kurang nyaman, serta ketiadaan
lingkungan metropolitan. Henderson dipakai untuk menunjukkan bahwa
konsentrasi yang terlalu tinggi maupun penyebaran yang terlalu besar
dapat bermasalah, sehingga bentuk ``konsentrasi yang terdesentralisasi''
mungkin menjadi kompromi. {[}TXT, baris 219-270; artikel, hlm.
2325-2326{]}

{[}Laporan sumber{]} Artikel membedakan PUR sebagai karakter morfologis,
yaitu pengelompokan kota-kota yang relatif sama besar dan berdekatan,
dari \emph{urban network} sebagai PUR yang juga memiliki integrasi
fungsional, interaksi, spesialisasi, dan komplementaritas. Karena
istilah polisentrisitas dapat dipakai pada banyak skala dan dalam makna
analitis maupun normatif, artikel memilih menggunakan polisentrisitas
terutama untuk persoalan morfologis pada skala regional. {[}TXT, baris
433-588; artikel, hlm. 2331-2332{]}

{[}Laporan sumber{]} Poros terakhir adalah literatur tentang amenitas
perkotaan. Artikel mengaitkan amenitas budaya, rekreasi, dan olahraga
dengan daya tarik kota, populasi permanen maupun sementara, serta arti
penting kota bagi pekerja pengetahuan. {[}TXT, baris 142-200; artikel,
hlm. 2324-2325{]}

{[}Inferensi pengolahan{]} TXT tidak menjelaskan prosedur pencarian
literatur, kriteria inklusi, jumlah korpus, atau metode tinjauan
sistematis. Karena itu, bagian ini harus dibaca sebagai landasan teori
dan pemetaan perdebatan yang dipilih penulis, bukan sebagai tinjauan
sistematis yang dapat direplikasi.

\subsection{Method Used}\label{method-used-50}

{[}Laporan sumber{]} Penelitian menggunakan desain komparatif
lintas-seksi pada 42 WGR-region. Penulis menginventarisasi amenitas,
menormalkan 20 indikator, membangun indeks, mengukur bentuk sistem
perkotaan dengan rank-size, lalu menggunakan regresi berganda untuk
menghubungkan bentuk sistem perkotaan dengan indeks amenitas. {[}TXT,
baris 272-338 dan 706-752; artikel, hlm. 2326-2329 dan 2336{]}

{[}Laporan sumber{]} Delimitasi WGR dipilih karena data perjalanan yang
secara khusus menunjukkan penggunaan amenitas tidak cukup andal. Batas
stadsgewest dianggap terlalu kecil dan secara definisional cenderung
monosentris. WGR-samenwerkingsgebieden dipakai sebagai proksi tidak
langsung untuk wilayah yang secara fungsional koheren, karena batasnya
dibentuk dari persepsi administrator dan anggota dewan mengenai skala
masalah yang memerlukan koordinasi antarmunicipal. {[}TXT, baris
240-318; artikel, hlm. 2326-2327{]}

{[}Laporan sumber{]} Polisentrisitas dihitung dari kemiringan garis
regresi distribusi ukuran lima kota atau tempat terbesar di setiap
wilayah. Garis yang lebih datar diperlakukan sebagai tanda hierarki yang
lebih lemah dan polisentrisitas yang lebih tinggi. Pengukuran tiga kota
digunakan sebagai pemeriksaan ketahanan. Artikel menyatakan bahwa hasil
dari lima dan tiga kota sangat berkorelasi, lalu menggunakan indikator
lima kota sebagai ukuran utama. {[}TXT, baris 589-691; artikel, hlm.
2332-2335{]}

{[}Laporan sumber{]} Dalam model regresi, keempat variabel independen
dimasukkan secara simultan: populasi WGR-region, jumlah tempat tidur
hotel terisi, pendapatan rumah tangga rata-rata, dan indikator
polisentrisitas berbasis lima kota. Semua 42 wilayah menjadi kasus
analisis. {[}TXT, baris 706-752; artikel, hlm. 2336{]}

{[}Laporan sumber{]} Analisis tambahan menyatakan jumlah amenitas
individual sebagai frekuensi per 1.000 penduduk, kemudian menghitung
korelasinya dengan indikator polisentrisitas. Catatan metodologis
artikel juga menyatakan tidak ada multikolinearitas; residual terstandar
berdistribusi normal dan varians residual konstan berdasarkan plot
diagnostik. {[}TXT, baris 755-800 dan 928-934; artikel, hlm.
2337-2340{]}

{[}Inferensi pengolahan{]} Rancangan ini memungkinkan perbandingan
statistik antarwilayah dan pengendalian beberapa perbedaan dasar, tetapi
TXT tidak menunjukkan eksperimen, perubahan kebijakan yang
dieksploitasi, panel waktu, atau strategi identifikasi kausal. Maka
hasilnya paling tepat diperlakukan sebagai hubungan asosiasional pada
satu periode pengamatan.

\subsection{Dataset}\label{dataset-50}

{[}Laporan sumber{]} Dataset penelitian berupa inventaris 20 variabel
amenitas pada enam kategori. Data mencakup kuantitas dan kualitas,
dengan pemilihan amenitas yang dianggap memiliki signifikansi supralokal
dan memerlukan dukungan atau massa kritis wilayah. Amenitas yang hanya
berfungsi pada skala kota kecil dan sebagian besar amenitas berskala
nasional dikeluarkan. {[}TXT, baris 319-338; artikel, hlm. 2327{]}

{[}Laporan sumber{]} Rincian data yang tercetak dalam Tabel 1 adalah
sebagai berikut:

\begin{itemize}
\tightlist
\item
  Cinema: jumlah layar dan kursi bioskop, jumlah layar dan kursi bioskop
  art-house, serta jumlah layar pada bioskop terbesar; seluruhnya tahun
  2005 dari NVB atau Dutch Association of Film Exhibitors.
\item
  Theatre: jumlah teater besar dengan lebih dari 400 kursi, teater
  menengah dengan 200-400 kursi, dan teater kecil dengan kurang dari 200
  kursi; tahun 2005 dari Theater Instituut Nederland.
\item
  Museum: jumlah museum tahun 2006 dari Stichting Museumkaart dan jumlah
  museum yang termasuk 50 museum paling banyak dikunjungi di Belanda,
  dengan data tahun 2004 dari The Netherlands Museum Association.
\item
  Concert: jumlah gedung konser tahun 2006, kapasitas tempat pertunjukan
  pop/konser, dan kapasitas tempat pop/konser terbesar; sumbernya
  organisasi asosiasi gedung konser dan tempat pop/festival yang
  tercetak dalam Tabel 1.
\item
  Sports: jumlah gedung ski dalam ruang, arena seluncur, lintasan
  karting, luas meter persegi dinding panjat dalam ruang, dan kapasitas
  stadion; sumbernya kombinasi inventaris penulis, NKBV, ESAC, dan World
  Stadiums, seluruhnya bertahun 2006.
\item
  Leisure: jumlah kasino yang terbuka untuk publik dan jumlah bintang
  Michelin restoran; tahun 2006, dengan sumber Holland Casino dan
  Michelin.
\end{itemize}

{[}Laporan sumber{]} Setiap variabel diubah menjadi z-score. Nilai z
sebesar 0 dikonversi menjadi skor 100, sedangkan satu simpangan baku
diberi nilai 20 sehingga, misalnya, z sebesar 1 menjadi skor 120.
Variabel dalam setiap kategori diberi bobot sama, dan keenam kategori
juga diberi bobot sama dalam indeks total. {[}TXT, baris 304-330;
artikel, hlm. 2327{]}

{[}Laporan sumber{]} Variabel penjelas diambil dari sumber dan tahun
yang berbeda: populasi dari Statistics Netherlands (CBS) pada 1 Januari
2005; tempat tidur hotel dan tingkat okupansi rata-rata tahun 2004 dari
Bedrijfsschap Horeca en Catering; serta pendapatan rumah tangga
rata-rata tahun 2002 dari CBS. {[}TXT, baris 654-691; artikel, hlm.
2334{]}

{[}Inferensi pengolahan{]} TXT tidak menyediakan berkas data mentah,
kode analisis, prosedur penanganan nilai hilang, atau dokumentasi
replikasi yang terpisah. Tabel 2 hanya mencetak skor wilayah, sedangkan
Tabel 1, Tabel 3, dan Tabel 4 mencetak definisi indikator serta keluaran
analisis.

\subsection{Population / Sample}\label{population-sample-50}

{[}Laporan sumber{]} Unit analisis adalah 42 WGR-region di Belanda.
Seluruh 42 wilayah tersebut bersama-sama mencakup seluruh wilayah
Belanda, bukan sekadar beberapa kota yang dipilih sebagai contoh.
Rata-rata populasi per wilayah pada 2005 adalah 387.034 penduduk; nilai
minimum adalah 107.620 di Zeeuwsch-Vlaanderen dan maksimum 1.352.680 di
Agglomeration Amsterdam. {[}TXT, baris 282-318; artikel, hlm.
2326-2327{]}

{[}Laporan sumber{]} WGR-region umumnya terletak dalam satu provinsi,
meskipun beberapa melintasi lebih dari satu provinsi. Penulis menyatakan
bahwa perjalanan dari satu tempat ke tempat lain dalam wilayah yang sama
umumnya dapat dilakukan dalam waktu setengah jam. Populasi untuk
pengukuran rank-size menggunakan kota dan desa aktual, bukan
municipality tempat kota atau desa tersebut berada. Rata-rata satu
WGR-region memuat hampir 60 kota dan desa pada 2005. {[}TXT, baris
299-318 dan 589-605; artikel, hlm. 2327 dan 2331-2332{]}

{[}Laporan sumber{]} Delimitasi wilayah ditentukan melalui pengaturan
kerja sama antarmunicipal WGR, yang batasnya mencerminkan persepsi
administrator dan anggota dewan mengenai skala persoalan yang
membutuhkan koordinasi regional. Artikel menyebutnya sebagai proksi
tidak langsung bagi koherensi fungsional dan menyatakan bahwa hasil
delimitasinya secara umum dapat dipertahankan. {[}TXT, baris 254-318;
artikel, hlm. 2326-2327{]}

{[}Inferensi pengolahan{]} ``Sample'' dalam studi ini bukan sampel
responden atau rumah tangga. TXT tidak menyebut metode sampling
probabilistik, jumlah responden, tingkat respons, atau karakteristik
individu. Empat puluh dua wilayah tersebut lebih tepat dipahami sebagai
sensus kasus WGR yang mencakup seluruh Belanda menurut delimitasi yang
dipilih, dengan generalisasi langsung terbatas pada konteks wilayah yang
sebanding.

\subsection{Dependent Variables}\label{dependent-variables-50}

{[}Laporan sumber{]} Variabel dependen utama adalah indeks total
kehadiran amenitas budaya, rekreasi, dan olahraga pada setiap
WGR-region. Indeks total merupakan rata-rata skor enam kategori: cinema,
theatre, museums, concerts, leisure, dan sports. Indeks dibangun dari 20
indikator yang telah dinormalisasi dan diberi bobot sama sesuai prosedur
yang dijelaskan penulis. {[}TXT, baris 274-338; artikel, hlm.
2326-2329{]}

{[}Laporan sumber{]} Enam indeks kategori juga tersedia dalam Tabel 2
sebagai gambaran kehadiran amenitas di masing-masing wilayah. Indeks
memadukan ukuran kuantitas dan kualitas, misalnya kapasitas gedung,
ukuran tempat pertunjukan, museum yang termasuk 50 paling banyak
dikunjungi, dan bintang Michelin. {[}TXT, baris 322-330 dan 359-390;
artikel, hlm. 2327-2329{]}

{[}Laporan sumber{]} Analisis tambahan memakai jumlah amenitas
individual per 1.000 penduduk untuk mengoreksi ukuran populasi ketika
memeriksa hubungan dengan polisentrisitas. Ukuran per 1.000 tersebut
berfungsi sebagai keluaran amenitas spesifik dalam analisis korelasi,
bukan menggantikan indeks total dalam model regresi utama. {[}TXT, baris
755-800; artikel, hlm. 2337{]}

{[}Inferensi pengolahan{]} Skor yang lebih tinggi dimaksudkan untuk
menunjukkan kehadiran amenitas yang lebih kuat, tetapi TXT tidak memberi
uji validasi eksternal atas indeks, uji sensitivitas bobot, atau alasan
statistik mengapa keenam kategori harus memiliki bobot identik. Hal
tersebut merupakan pilihan operasional penelitian, bukan fakta alamiah
tentang kualitas suatu wilayah.

\subsection{Results}\label{results-50}

{[}Laporan sumber{]} Model regresi utama menghasilkan N = 42, R2 =
0.953, adjusted R2 = 0.948, F = 187.764, dan Sign. F = 0.000. Artikel
menyatakan bahwa keempat variabel independen secara signifikan
berkontribusi pada kehadiran amenitas. {[}TXT, baris 706-752; artikel,
hlm. 2336{]}

{[}Laporan sumber{]} Koefisien dan koefisien beta terstandar yang
dicetak dalam Tabel 3 adalah sebagai berikut:

\begin{itemize}
\tightlist
\item
  Number of occupied hotel beds: B = 0.001 dan Beta = 0.264, dengan p
  yang dicetak 0.00.
\item
  Household income: B = -0.642 dan Beta = -0.078, dengan p yang dicetak
  0.05.
\item
  Polycentricity: B = -5.381 dan Beta = -0.109, dengan p yang dicetak
  0.01.
\item
  Population size: B = 3.991E-05 dan Beta = 0.761, dengan p yang dicetak
  0.00.
\item
  Constant: B = 95.591, dengan standard error 10.129.
\end{itemize}

{[}Laporan sumber{]} Populasi adalah penjelas terkuat; artikel menyebut
adjusted R2 akan menjadi 0.874 bila hanya variabel populasi yang
digunakan. Jumlah penduduk nonpermanen yang diproksikan dengan tempat
tidur hotel terisi memiliki beta tertinggi kedua. Bentuk sistem
perkotaan kemudian menjelaskan sedikit lebih banyak daripada pendapatan
rumah tangga. Hubungan pendapatan dengan indeks amenitas justru negatif,
berlawanan dengan ekspektasi awal penulis. {[}TXT, baris 715-733;
artikel, hlm. 2336{]}

{[}Laporan sumber{]} Dengan nilai rata-rata populasi 387.034, tempat
tidur hotel terisi 1.956,8, pendapatan rumah tangga 30.188, dan nilai
polisentrisitas minimum -1,94 yang diperlakukan sebagai wilayah paling
monosentris, model memprediksi indeks total 104,28. Ketika nilai
polisentrisitas maksimum digunakan dengan kontrol lain tetap pada
rata-ratanya, prediksi indeks menjadi 97,12. Selisihnya lebih dari tujuh
poin indeks, hampir setengah simpangan baku indeks total menurut
perbandingan penulis. {[}TXT, baris 755-771; artikel, hlm. 2337{]}

{[}Laporan sumber{]} Penggunaan indikator polisentrisitas yang dihitung
dari tiga kota tetap menghasilkan hubungan yang signifikan, dengan p =
0.04. Beta-nya sedikit lebih kecil daripada beta indikator lima kota.
{[}TXT, baris 772-775; artikel, hlm. 2337{]}

{[}Laporan sumber{]} Setelah jumlah amenitas dinyatakan per 1.000
penduduk, delapan korelasi yang signifikan dengan polisentrisitas
semuanya negatif:

\begin{itemize}
\tightlist
\item
  Kursi bioskop per 1.000 penduduk: r = -0.533.
\item
  Layar bioskop per 1.000 penduduk: r = -0.473.
\item
  Kursi bioskop art-house per 1.000 penduduk: r = -0.384.
\item
  Layar bioskop art-house per 1.000 penduduk: r = -0.427.
\item
  Kasino terbuka untuk publik per 1.000 penduduk: r = -0.431.
\item
  Gedung konser per 1.000 penduduk: r = -0.343.
\item
  Kapasitas penonton tempat pop/konser per 1.000 penduduk: r = -0.317.
\item
  Luas dinding panjat dalam ruang per 1.000 penduduk: r = -0.377.
\end{itemize}

{[}TXT, baris 777-800; artikel, hlm. 2337{]}

{[}Laporan sumber{]} Hubungan yang tidak signifikan tercetak untuk
sebagian besar jenis teater, museum, museum yang termasuk 50 paling
banyak dikunjungi, bintang Michelin, arena seluncur, gedung ski dalam
ruang, lintasan karting, dan kapasitas stadion. Teater kecil berkorelasi
negatif dan teater menengah, museum, serta gedung ski dalam ruang
berkorelasi positif, tetapi tanpa tanda signifikansi pada Tabel 4.
{[}TXT, baris 785-800; artikel, hlm. 2337{]}

{[}Provenance penting{]} Tabel 3 mencetak tanda bintang pada sejumlah
p-value, sedangkan catatannya menyatakan
\texttt{***\ Significant\ at\ the\ 0.001\ level}. Tanda bintang yang
tercetak bersama p = 0.05 dan p = 0.01 tidak sepenuhnya konsisten dengan
catatan tersebut. Review ini mempertahankan angka dan catatan
sebagaimana muncul dalam TXT dan tidak memperbaiki ambang signifikansi
berdasarkan dugaan.

\subsection{Findings}\label{findings-50}

{[}Interpretasi penulis{]} Hasil utama ditafsirkan sebagai kerugian PUR
dibandingkan wilayah kota yang lebih monosentris setelah perbedaan
populasi, jumlah pengunjung, dan pendapatan dikendalikan. Semakin
polisentris suatu wilayah, semakin sedikit amenitas budaya, rekreasi,
dan olahraga yang tersedia. Penulis mengaitkan korelasi negatif pada
bioskop, kasino, gedung konser, tempat pop/konser, dan dinding panjat
dengan kemungkinan bahwa pelaku usaha memandang basis dukungan wilayah
polisentris lebih lemah. {[}TXT, baris 806-831; artikel, hlm.
2337-2338{]}

{[}Interpretasi penulis{]} Amenitas yang memerlukan basis dukungan besar
karena menyediakan layanan khusus tampak lebih jarang pada wilayah
polisentris. Sebaliknya, bentuk sistem perkotaan tidak tampak menentukan
bagi sebagian amenitas olahraga dan sektor teater secara keseluruhan.
Teater menengah tampak lebih umum dan museum serta gedung ski dalam
ruang memiliki korelasi positif, tetapi hubungan-hubungan tersebut tidak
signifikan dalam Tabel 4. {[}TXT, baris 821-831; artikel, hlm. 2338{]}

{[}Interpretasi penulis{]} Penulis memandang hasil ini mendukung
kekhawatiran Parr serta Turok dan Bailey bahwa PUR tidak memperoleh
eksternalitas aglomerasi atau eksternalitas regional dengan cara yang
sama seperti wilayah monosentris. Susunan ruang yang tersebar membuat
ukuran gabungan tidak otomatis berfungsi sebagai satu massa perkotaan.
{[}TXT, baris 815-831; artikel, hlm. 2338{]}

{[}Interpretasi penulis{]} Empat mekanisme yang diajukan untuk
menjelaskan temuan adalah sebagai berikut:

\begin{itemize}
\tightlist
\item
  Jarak antarkota dapat menghasilkan arus perjalanan, komoditas, dan
  informasi yang lebih panjang atau kurang nyaman, sekaligus tidak
  menyediakan lingkungan metropolitan yang padat.
\item
  Rivalitas dan kompetisi antarkota dapat mendorong duplikasi fasilitas,
  citra kota, proyek pembangunan kembali, inovasi, dan investasi, bukan
  spesialisasi serta komplementaritas. Duplikasi tersebut mengurangi
  keragaman amenitas pada skala PUR.
\item
  Amenitas tertanam dalam konteks perkotaan dan jalur perkembangan
  historis. Amenitas masa kini sering dibangun di atas infrastruktur
  lama, pola hubungan lokal, dan pola sentralitas yang diwariskan.
\item
  Pembiayaan amenitas banyak berlangsung pada basis municipality. Hibah
  umum pemerintah pusat, menurut uraian artikel, menggunakan kriteria
  seperti jumlah unit hunian, penduduk, dan area terbangun yang
  cenderung menguntungkan kota yang lebih besar. Akibatnya, satu
  municipality berpenduduk 300.000 dianggap memiliki basis finansial
  lebih besar daripada tiga kota berpenduduk 100.000 secara
  bersama-sama.
\end{itemize}

{[}TXT, baris 832-879; artikel, hlm. 2338-2339{]}

{[}Inferensi pengolahan{]} Empat mekanisme tersebut adalah penjelasan
yang masuk akal menurut penulis, tetapi tidak diuji satu per satu dalam
model utama. Model mengukur bentuk wilayah dan amenitas, bukan
kompetisi, duplikasi, kepadatan interaksi, struktur subsidi, atau arus
perjalanan secara langsung. Oleh karena itu, mekanisme tersebut tidak
boleh diperlakukan sebagai hasil kausal yang telah teridentifikasi.

\subsection{Conclusions}\label{conclusions-50}

{[}Laporan sumber{]} Jawaban artikel terhadap pertanyaan utamanya adalah
bahwa kumpulan kota yang terpisah secara spasial tidak dapat memperoleh
keuntungan ukuran urban dengan tingkat yang sama seperti satu wilayah
kota monosentris, setidaknya untuk amenitas budaya, rekreasi, dan
olahraga yang dipilih. Setelah pengendalian terhadap populasi,
pengunjung, dan pendapatan rata-rata, PUR memiliki indeks amenitas lebih
rendah daripada wilayah yang lebih monosentris. {[}TXT, baris 806-831
dan 834-921; artikel, hlm. 2337-2340{]}

{[}Interpretasi penulis{]} Kesimpulan ini diringkas melalui judul dan
kalimat akhir bahwa sekumpulan kota yang berdekatan, dalam hal basis
dukungan amenitas urban, bukan hanya tidak menjadi kota besar; mereka
bahkan tidak mencapai jumlah sederhana dari bagian-bagiannya. Namun,
penulis secara eksplisit membatasi klaim tersebut pada WGR-region dan
jenis amenitas yang diteliti. {[}TXT, baris 880-921; artikel, hlm.
2339-2340{]}

{[}Laporan sumber{]} Penulis tidak menyimpulkan bahwa seluruh jenis
keuntungan aglomerasi pasti lebih rendah dalam PUR. Justru, ia
menyatakan bahwa keunggulan perusahaan dan pekerja, indikator lain
termasuk diseconomies aglomerasi, tahap integrasi fungsional, dan
wilayah di negara lain masih dapat menghasilkan pola berbeda. {[}TXT,
baris 880-937; artikel, hlm. 2339-2340{]}

{[}Inferensi pengolahan{]} Kesimpulan yang dapat dipertahankan dari TXT
adalah kesimpulan kontekstual dan berbasis indikator, bukan hukum umum
tentang semua PUR. Generalisasi di luar Belanda, WGR-region, periode
data, dan amenitas yang dipilih memerlukan bukti tambahan.

\subsection{Contributions}\label{contributions-50}

{[}Laporan sumber{]} Penulis menyatakan bahwa artikel ini bertujuan
mengembangkan perdebatan tentang PUR dengan menguji secara empiris
apakah keuntungan ukuran urban berkembang pada skala PUR. {[}TXT, baris
142-166; artikel, hlm. 2324{]}

{[}Inferensi pengolahan{]} Berdasarkan desain dan hasil yang tercetak,
kontribusi substantif dan metodologis artikel dapat dirumuskan sebagai
berikut:

\begin{itemize}
\tightlist
\item
  Mengubah klaim kebijakan tentang ``ukuran gabungan'' PUR menjadi
  perbandingan empiris antara sistem perkotaan polisentris dan
  monosentris.
\item
  Menyediakan operasionalisasi indeks amenitas yang memadukan 20
  indikator kuantitas dan kualitas ke dalam enam kategori dan indeks
  total.
\item
  Menggunakan kemiringan distribusi rank-size lima kota terbesar sebagai
  ukuran kontinu bentuk morfologis, sehingga polisentrisitas tidak
  diperlakukan hanya sebagai kategori biner.
\item
  Menunjukkan bahwa hubungan negatif tetap muncul ketika ukuran
  polisentrisitas dihitung dari tiga kota, meskipun indikatornya sedikit
  lebih lemah.
\item
  Memisahkan secara konseptual morfologi PUR dari \emph{urban network}
  yang membutuhkan integrasi fungsional, spesialisasi, interaksi, dan
  komplementaritas.
\item
  Menempatkan pembiayaan, koordinasi regional, dan kompetisi antarkota
  sebagai kemungkinan saluran yang perlu diteliti untuk memahami mengapa
  massa penduduk gabungan tidak setara dengan massa kota besar.
\end{itemize}

{[}Inferensi pengolahan{]} Daftar di atas adalah sintesis atas isi
artikel, bukan daftar kontribusi resmi yang ditulis dalam subbagian
khusus. TXT tidak menyediakan pernyataan kontribusi penulis dalam format
deklaratif atau pembandingan sistematis dengan semua studi terdahulu.

\subsection{Research Gap}\label{research-gap-50}

{[}Laporan sumber{]} Artikel mengidentifikasi bahwa banyak penelitian
sebelumnya berfokus pada pembentukan konsep PUR, karakteristik
definisional, penyusunan agenda, penerapan pada wilayah tertentu,
kapasitas kelembagaan, dan tata kelola. Penulis menyatakan bahwa
perhatian perlu dialihkan pada pembuktian empiris dan validasi terhadap
berbagai klaim tentang PUR. Artikel ini mengisi sebagian celah tersebut
dengan menguji hubungan bentuk sistem perkotaan dan penyediaan amenitas.
{[}TXT, baris 142-166; artikel, hlm. 2324{]}

{[}Laporan sumber{]} Artikel juga menilai bahwa bentuk sistem perkotaan
sering dianggap sudah ada begitu saja dalam analisis sistem perkotaan,
padahal bentuk tersebut mungkin memengaruhi keuntungan dan kerugian
aglomerasi. Pengukuran morfologis berbasis rank-size digunakan untuk
mengatasi pengabaian tersebut pada skala regional. {[}TXT, baris 240-270
dan 589-628; artikel, hlm. 2326 dan 2332{]}

{[}Laporan sumber{]} Celah yang masih terbuka menurut bagian akhir
artikel meliputi hal-hal berikut:

\begin{itemize}
\tightlist
\item
  Pengujian keuntungan yang lebih langsung diterima perusahaan dan
  pekerja, bukan hanya amenitas internal kota yang eksternal bagi
  perusahaan.
\item
  Analisis tren untuk mengetahui apakah kesenjangan PUR dan wilayah
  monosentris menyempit atau melebar.
\item
  Pengujian indikator lain, termasuk indikator yang biasanya dirangkum
  sebagai diseconomies aglomerasi.
\item
  Perbandingan yang lebih luas di negara dan wilayah lain.
\item
  Perbandingan PUR yang telah berkembang menjadi \emph{urban network}
  dengan PUR yang baru memulai proses tersebut.
\item
  Pengujian peran interaksi, komplementaritas, integrasi fungsional,
  aksesibilitas internal, kapasitas pengorganisasian regional, kerja
  sama, dan koordinasi.
\item
  Pemahaman yang lebih mendalam tentang proses yang mendorong atau
  menghambat keuntungan aglomerasi melalui metode lain, termasuk
  wawancara pemangku kepentingan.
\end{itemize}

{[}TXT, baris 880-948; artikel, hlm. 2339-2340{]}

{[}Inferensi pengolahan{]} Celah metodologis yang juga terlihat dari TXT
adalah belum adanya pengujian mekanisme yang diajukan, validasi
alternatif terhadap bobot indeks, dan desain temporal yang dapat
membedakan urutan sebab-akibat. Tiga hal ini merupakan inferensi dari
rancangan yang dilaporkan, bukan agenda tambahan yang dinyatakan sebagai
daftar khusus oleh penulis.

\subsection{Limitations}\label{limitations-50}

{[}Laporan sumber{]} Keterbatasan yang diakui secara langsung adalah
keterbatasan ruang lingkup dan skala. Penulis menegaskan bahwa temuan
berlaku untuk WGR-region dan amenitas yang dipelajari, serta bahwa
persoalan eksternalitas regional dalam sistem polisentris memerlukan
eksplorasi lebih lanjut. {[}TXT, baris 880-900; artikel, hlm. 2339{]}

{[}Laporan sumber{]} Data perjalanan yang memetakan penggunaan amenitas
tidak cukup andal untuk dijadikan dasar delimitasi ideal. Karena itu,
WGR-region merupakan proksi tidak langsung koherensi fungsional yang
berasal dari persepsi administrator dan anggota dewan, bukan batas yang
dibangun langsung dari arus perjalanan menuju amenitas. {[}TXT, baris
240-318; artikel, hlm. 2326-2327{]}

{[}Laporan sumber{]} Pengukuran polisentrisitas terutama menangkap
distribusi ukuran kota secara morfologis. Artikel sendiri membedakan hal
tersebut dari integrasi fungsional \emph{urban network} dan kemudian
meminta riset lebih lanjut tentang interaksi, komplementaritas,
aksesibilitas internal, serta kapasitas koordinasi. {[}TXT, baris
547-588 dan 909-937; artikel, hlm. 2331-2332 dan 2340{]}

{[}Laporan sumber{]} Cakupan substantif dibatasi pada amenitas budaya,
rekreasi, dan olahraga yang dipilih karena signifikansi supralokal.
Keuntungan yang langsung diterima perusahaan dan pekerja belum diukur.
Negara yang diteliti hanya Belanda, dan pengamatan utama menggunakan
satu gambaran tahun untuk ukuran kota serta populasi. {[}TXT, baris
319-338, 589-605, dan 880-925; artikel, hlm. 2327, 2331-2332, dan
2339{]}

{[}Inferensi pengolahan{]} Beberapa batas analitis dapat ditarik dari
prosedur yang dilaporkan:

\begin{itemize}
\tightlist
\item
  Tahun data tidak seragam: sebagian besar amenitas berasal dari
  2005-2006, tempat tidur hotel dan okupansi dari 2004, serta pendapatan
  dari 2002. TXT mencantumkan perbedaan ini tetapi tidak melaporkan
  analisis sensitivitas terhadap ketidakserentakan waktu.
\item
  Bobot sama untuk semua indikator dalam kategori dan semua kategori
  dalam indeks merupakan keputusan konstruksi yang tidak dibandingkan
  dengan skema bobot alternatif dalam TXT.
\item
  Pemilihan lima kota terbesar dibahas sebagai pilihan yang tidak
  sepenuhnya arbitrer, tetapi TXT hanya melaporkan pemeriksaan
  alternatif dengan tiga kota, bukan berbagai ukuran sampel atau ambang
  ukuran lain.
\item
  Model menjelaskan 94,8 persen varians indeks total, tetapi angka
  tersebut tidak dengan sendirinya membuktikan kausalitas atau
  memastikan bahwa semua mekanisme substantif telah terukur.
\item
  TXT tidak menyebut prosedur untuk nilai hilang, audit reliabilitas
  antar-sumber, atau validasi independen atas inventaris amenitas.
\end{itemize}

{[}Provenance{]} Tabel 3 memiliki ketidakselarasan tampilan antara tanda
bintang dan catatan ambang signifikansi, sebagaimana dicatat pada bagian
Results. Hal tersebut membatasi kepastian dalam membaca tingkat
signifikansi persis dari baris-baris tabel, meskipun arah koefisien dan
angka p yang dicetak tetap dapat dilaporkan.

\subsection{Challenges}\label{challenges-50}

{[}Laporan sumber{]} Tantangan utama yang secara eksplisit disebut
penulis adalah bagaimana memanfaatkan potensi penduduk yang ada di PUR
melalui koordinasi regional agar kompetisi yang tidak efisien dan
duplikasi dapat dihindari. Penulis menilai pembiayaan amenitas pada
basis regional, bukan hanya lokal, diperlukan; metode untuk
menyeimbangkan keuntungan dan beban antarkota juga dianggap penting.
{[}TXT, baris 945-961; artikel, hlm. 2340{]}

{[}Laporan sumber{]} Pada tahap penelitian, tantangannya adalah
menemukan unit wilayah yang cukup fungsional tetapi tidak secara
definisional monosentris, memilih ukuran sampel kota yang masuk akal
untuk rank-size, serta membedakan pengaruh jumlah penduduk dari bentuk
sistem perkotaan. Penulis membahas ketidakandalan data perjalanan,
masalah ambang ukuran tetap, dan sifat arbitrer dari jumlah kota yang
dipilih. {[}TXT, baris 240-270 dan 614-667; artikel, hlm. 2326-2327 dan
2332-2334{]}

{[}Interpretasi penulis{]} Tantangan kelembagaan dan politik juga muncul
dari sejarah rivalitas antarkota, kecenderungan meniru proyek dan citra
kota, serta sistem hibah yang menurut uraian artikel lebih menguntungkan
municipality besar. Tantangan ini menghambat pergeseran dari duplikasi
lokal menuju komplementaritas regional. {[}TXT, baris 832-879; artikel,
hlm. 2338-2339{]}

{[}Inferensi pengolahan{]} Artikel belum menyelesaikan tantangan
pengukuran integrasi fungsional, mekanisme koordinasi, atau distribusi
manfaat dan beban; ia terutama mengidentifikasi tantangan tersebut
sebagai penjelasan dan agenda kerja.

\subsection{Practical Implications}\label{practical-implications-50}

{[}Interpretasi penulis{]} Pemerintah dan perencana regional tidak
seharusnya menganggap bahwa pengelompokan beberapa kota kecil atau
menengah akan otomatis menyediakan amenitas tingkat tinggi seperti satu
kota besar dengan jumlah penduduk gabungan yang sama. Indeks amenitas
dan bentuk sistem perlu dibaca bersama dengan ukuran penduduk, populasi
pengunjung, pendapatan, serta kondisi pembiayaan. {[}TXT, baris 806-831
dan 880-899; artikel, hlm. 2338-2339{]}

{[}Interpretasi penulis{]} Implikasi kebijakan yang diajukan adalah
memperkuat koordinasi regional, mengurangi kompetisi yang menghasilkan
pengulangan fasilitas, dan mengarahkan kota-kota pada pembagian fungsi
atau komplementaritas. Penulis juga menekankan pentingnya pembiayaan
amenitas pada skala regional serta mekanisme untuk membagi keuntungan
dan beban antarkota. {[}TXT, baris 909-961; artikel, hlm. 2340{]}

{[}Interpretasi penulis{]} Integrasi fungsional yang lebih besar,
identitas regional, cara pandang regional, aksesibilitas internal,
kapasitas pengorganisasian, dan koordinasi disebut sebagai kondisi yang
mungkin membantu PUR memanfaatkan potensi penduduknya. {[}TXT, baris
909-937; artikel, hlm. 2340{]}

{[}Inferensi pengolahan{]} Implikasi di atas adalah rekomendasi
konseptual yang diturunkan dari hasil observasional dan mekanisme yang
diusulkan. TXT tidak melaporkan evaluasi program, perubahan anggaran,
atau bukti bahwa intervensi koordinasi dan pembiayaan regional telah
meningkatkan amenitas.

\subsection{Applications}\label{applications-50}

{[}Inferensi pengolahan{]} Berdasarkan prosedur yang dilaporkan,
pendekatan artikel dapat diterapkan sebagai alat diagnosis komparatif
untuk menilai apakah wilayah dengan ukuran penduduk gabungan memiliki
kapasitas amenitas yang sebanding dengan wilayah monosentris. Penerapan
semacam itu dapat menggunakan tiga komponen yang memang dicontohkan
dalam TXT: inventaris amenitas, ukuran rank-size sistem perkotaan, dan
model yang mengendalikan populasi serta pengunjung.

{[}Interpretasi penulis{]} Penggunaan kebijakan yang paling dekat dengan
isi artikel adalah perencanaan pembangunan regional, koordinasi
fasilitas budaya, rekreasi, dan olahraga, serta pengaturan pembiayaan
agar tidak seluruh keputusan dan beban berada pada municipality
individual. {[}TXT, baris 945-961; artikel, hlm. 2340{]}

{[}Provenance{]} Tidak ada studi implementasi, evaluasi kebijakan, atau
contoh penerapan model pada wilayah di luar 42 WGR-region yang
dilaporkan dalam TXT. Karena itu, keberhasilan aplikasi praktis tidak
tersedia dan tidak boleh diasumsikan.

\subsection{Future Research
(Optional)}\label{future-research-optional-50}

{[}Laporan sumber{]} Penulis secara eksplisit mengusulkan agenda riset
berikut:

\begin{itemize}
\tightlist
\item
  Menguji keuntungan yang lebih langsung diterima perusahaan dan
  pekerja, bukan hanya amenitas yang internal bagi kota tetapi eksternal
  bagi perusahaan.
\item
  Mengikuti perubahan dari waktu ke waktu untuk melihat apakah
  kesenjangan amenitas antara PUR dan wilayah monosentris menyempit atau
  melebar.
\item
  Menilai indikator lain, termasuk kemungkinan keuntungan PUR dalam
  mengurangi atau menghindari diseconomies aglomerasi.
\item
  Memperluas bukti ke lebih banyak negara dan wilayah.
\item
  Membandingkan PUR yang telah berkembang menjadi \emph{urban network}
  dengan PUR yang baru memulai proses pengembangan tersebut.
\item
  Menguji apakah interaksi, komplementaritas, integrasi fungsional,
  aksesibilitas internal, identitas regional, kapasitas
  pengorganisasian, kerja sama, dan koordinasi yang lebih kuat berkaitan
  dengan kinerja yang lebih baik.
\item
  Menggunakan metode lain, seperti wawancara pemangku kepentingan, untuk
  memahami proses yang mendorong atau menghambat keuntungan aglomerasi.
\end{itemize}

{[}TXT, baris 880-948; artikel, hlm. 2339-2340{]}

{[}Inferensi pengolahan{]} Agenda tambahan yang tersirat dari desain
adalah uji sensitivitas bobot indeks, pengukuran mekanisme pembiayaan
dan duplikasi secara langsung, serta replikasi dengan batas wilayah dan
indikator polisentrisitas alternatif. Ini bukan usulan eksplisit yang
dirumuskan sebagai daftar oleh penulis.

\subsection{Provenance Notes}\label{provenance-notes-50}

\begin{itemize}
\tightlist
\item
  File sumber yang digunakan hanya
  \texttt{/Users/mac/Documents/Mac/{[}2{]}\ Obsidian\ Vault/zettelkasten/source/journal/B07-meijers-2008-summing-small-cities-does-not-make-a-large-city-polycentric-urban-regions-and-the-provision-of-cultural-leisure-and-sports-amenities-10.1177-0042098008095870.txt}.
\item
  Verifikasi filename di bawah \texttt{source/} menemukan tepat satu TXT
  yang cocok dengan prefiks \texttt{B07}: file jurnal yang tercantum di
  atas. Tidak ada TXT kode lain yang diproses dalam audit ini.
\item
  Seluruh isi TXT dibaca dari baris 1 sampai 1.049. TXT terdiri atas
  blok \texttt{{[}PDF\ Metadata{]}} pada baris 1-26, blok
  \texttt{{[}Extracted\ Text{]}} mulai baris 27, dan karakter form-feed
  penutup pada baris 1.049.
\item
  Metadata ekstraksi menyatakan bahwa teks diambil pada 2026-08-31
  dengan \texttt{pdftotext\ -layout} dari PDF 21 halaman. Tanggal
  tersebut adalah tanggal ekstraksi yang tercetak dalam TXT, bukan
  tanggal akses yang diverifikasi secara terpisah.
\item
  Identitas bibliografis dalam review diambil dari byline dan header
  artikel pada TXT, terutama baris 34-38 dan 74-119. Metadata teknis PDF
  mencetak judul \texttt{2323-2342\ USJ\_095870.indd} dan author
  \texttt{Administrator}; keduanya diperlakukan sebagai metadata teknis
  yang tidak menggantikan judul serta penulis artikel yang tercetak pada
  naskah.
\item
  Locator utama menggunakan nomor halaman artikel, nomor bagian, nomor
  tabel/gambar, dan rentang baris TXT. Nomor halaman artikel tercetak di
  dalam hasil ekstraksi, sedangkan nomor baris merujuk pada posisi dalam
  TXT yang dibaca.
\item
  Beberapa bagian hasil ekstraksi dengan tata letak dua kolom, terutama
  Tabel 2 dan halaman yang memuat perpindahan kolom, tidak tersusun
  seperti tampilan PDF. Angka dan klaim yang digunakan dalam review
  hanya diambil ketika konteksnya dapat dibaca dari TXT; tidak ada
  pengurutan ulang atau pengisian angka yang tidak tercetak.
\item
  Tabel 1 mencetak tahun dan sumber berbeda untuk indikator amenitas,
  tetapi TXT tidak menyediakan data mentah, kode analisis, definisi
  rinci seluruh sumber, atau prosedur validasi inventaris.
\item
  Status peer-review tidak disebutkan dalam file. TXT memang mencetak
  nama jurnal dan keterangan ``Version of Record'', tetapi status
  peer-review tidak diasumsikan.
\item
  Kutipan literal dalam review dibatasi pada judul, istilah, atau frasa
  singkat yang memang tercetak; seluruh uraian substantif lainnya adalah
  parafrasa yang diberi label laporan, interpretasi penulis, atau
  inferensi pengolahan.
\end{itemize}

\section{Review B08}\label{review-b08}

\textbf{Identitas sumber}

\begin{itemize}
\tightlist
\item
  Kode kajian: B08. Kode ini diturunkan dari prefiks nama file TXT; kode
  tersebut tidak tercantum di dalam artikel.
\item
  Judul: \emph{Secondary Yet Metropolitan? The Challenges of
  Metropolitan Integration for Second-Tier Cities}.
\item
  Penulis: Rodrigo V. Cardoso dan Evert J. Meijers.
\item
  Jurnal: \emph{Planning Theory \& Practice}, 2017.
\item
  DOI: \texttt{10.1080/14649357.2017.1371789}.
\item
  Publikasi online: 10 Oktober 2017. Riwayat artikel yang tercantum
  dalam TXT: diterima 19 Desember 2016 dan disetujui 22 Agustus 2017.
\item
  Fokus geografis dan substantif: kota second-tier di Eropa, hubungan
  kota dengan wilayah metropolitan, manfaat aglomerasi, kapasitas
  kelembagaan, dan integrasi politik.
\item
  Status akses untuk review ini: seluruh teks hasil ekstraksi TXT
  dibaca; PDF hanya disebut oleh blok metadata TXT dan tidak digunakan
  sebagai bahan tambahan.
\end{itemize}

Penanda epistemik yang digunakan: \textbf{Laporan sumber} berarti isi
yang dinyatakan atau dilaporkan artikel; \textbf{Interpretasi penulis}
berarti posisi analitis yang dibangun penulis; \textbf{Inferensi
pengolahan} berarti pembacaan kritis yang ditarik dari struktur
teks/metode, bukan klaim eksplisit artikel.

\subsection{Summarized Abstract}\label{summarized-abstract-51}

\begin{itemize}
\tightlist
\item
  \textbf{Laporan sumber:} Artikel menanyakan apakah wilayah tempat
  integrasi metropolitan secara umum menguntungkan kota beririsan dengan
  wilayah kerugian yang lazim dialami kota second-tier di Eropa, serta
  apa implikasi dari pertemuan kedua hal tersebut. Integrasi
  metropolitan didefinisikan melalui dimensi spasial-fungsional,
  politik-kelembagaan, dan simbolik-kultural. Potensi manfaatnya
  mencakup pemanfaatan manfaat aglomerasi pada skala metropolitan,
  penggunaan sumber daya bersama secara efisien, dan peningkatan
  pengaruh politik-institusional pada tingkat pemerintahan yang lebih
  tinggi (TXT lines 99-119; PDF p.~1).
\item
  \textbf{Laporan sumber:} Artikel menilai secara empiris pertambahan
  massa demografis dan fungsional yang dapat diperoleh kota second-tier
  ketika skala metropolitan digabungkan. Untuk kapasitas kelembagaan dan
  politik, artikel menggunakan literatur yang ada serta contoh-contoh
  wilayah metropolitan sebagai ilustrasi (TXT lines 109-114; PDF p.~1).
\item
  \textbf{Interpretasi penulis:} Pembentukan wilayah metropolitan
  dipandang sebagai strategi yang menjanjikan terutama bagi kota
  second-tier yang tertanam dalam wilayah perkotaan yang besar dan padat
  atau berada di negara dengan kota first-tier yang dominan. Potensi itu
  tidak otomatis terwujud, sehingga perlu dikelola melalui strategi
  perencanaan dan tata kelola yang sesuai konteks (TXT lines 115-119;
  PDF p.~1).
\end{itemize}

\subsection{Problem Statement}\label{problem-statement-51}

\begin{itemize}
\tightlist
\item
  \textbf{Laporan sumber:} Kota second-tier didefinisikan sebagai kota
  non-ibu kota, umumnya berukuran menengah, yang kinerja ekonomi dan
  sosialnya cukup penting untuk memengaruhi kinerja ekonomi nasional.
  Kota-kota ini merupakan simpul penting sistem perkotaan Eropa, tetapi
  biasanya tidak memiliki fungsi komando dan kontrol pada skala yang
  lebih tinggi (TXT lines 122-131; PDF pp.~1-2).
\item
  \textbf{Laporan sumber:} Kerugian yang dibahas dibandingkan dengan
  kota first-tier meliputi kemampuan yang lebih lemah untuk memperoleh
  manfaat aglomerasi karena ukuran dan fungsi perkotaan yang lebih
  kecil, pengabaian relatif dalam penerimaan investasi publik untuk daya
  saing, serta pengaruh politik yang terbatas dalam pembuatan kebijakan
  nasional atau Eropa. Kerugian ini lebih mungkin menonjol dalam sistem
  perkotaan yang tersentralisasi dan didominasi ibu kota, atau di luar
  koridor ekonomi yang diuntungkan (TXT lines 169-182, 291-304; PDF
  pp.~2, 4-5).
\item
  \textbf{Interpretasi penulis:} Masalah penelitian bukan sekadar apakah
  kota-kota di sekitar suatu pusat dapat dijumlahkan, melainkan apakah
  integrasi fungsional, kelembagaan, dan simbolik dapat mengubah massa
  yang tersebar menjadi kapasitas metropolitan yang berguna bagi kota
  second-tier. Artikel juga menegaskan sejak awal bahwa kebutuhan dan
  urgensi integrasi berbeda antarwilayah (TXT lines 255-288, 302-324;
  PDF pp.~4-5).
\end{itemize}

\subsection{Objectives}\label{objectives-51}

\begin{itemize}
\tightlist
\item
  \textbf{Laporan sumber:} Tujuan utama artikel adalah membahas tumpang
  tindih antara bidang manfaat integrasi metropolitan dan bidang
  kerugian yang menetap pada banyak kota second-tier Eropa dibandingkan
  kota first-tier (TXT lines 195-204; PDF p.~3).
\item
  \textbf{Laporan sumber:} Tujuan operasionalnya adalah mendefinisikan
  integrasi metropolitan dan manfaat potensialnya; mengidentifikasi
  kerugian khas kota second-tier; membandingkan pertambahan relatif
  populasi dan fungsi perkotaan pada skala metropolitan; serta menilai
  potensi berbagi sumber daya dan peningkatan suara politik melalui
  literatur dan contoh yang relevan (TXT lines 195-205; PDF p.~3).
\item
  \textbf{Laporan sumber:} Tujuan lanjutan adalah membahas pendekatan
  perencanaan dan tata kelola untuk mengelola peluang, hambatan, relasi
  pusat-periferi, dan jalur menuju integrasi. Artikel ditutup dengan
  kesimpulan dan saran penelitian (TXT lines 205-214; PDF p.~3).
\item
  \textbf{Inferensi pengolahan:} TXT tidak merumuskan hipotesis
  statistik atau pertanyaan kausal formal. Tujuan artikel lebih tepat
  dibaca sebagai pengujian potensi dan penyusunan argumen kebijakan
  daripada pengukuran efek integrasi sebagai perlakuan.
\end{itemize}

\subsection{Literature Survey}\label{literature-survey-51}

\begin{itemize}
\tightlist
\item
  \textbf{Laporan sumber:} Kerangka konseptual artikel mendefinisikan
  integrasi metropolitan sebagai keterjalinan dan interaksi
  tempat-tempat dalam lingkup pengaruh potensial satu sama lain. Dimensi
  spasial-fungsional mencakup spesialisasi ekonomi yang saling
  melengkapi dan jaringan yang efisien; dimensi politik-kelembagaan
  mencakup struktur tata kelola metropolitan dan kerja sama
  antarpemerintah kota; dimensi simbolik-kultural mencakup pembangunan
  identitas metropolitan dan perluasan keterikatan tempat (TXT lines
  217-225; PDF p.~3).
\item
  \textbf{Laporan sumber:} Dalam literatur ekonomi aglomerasi, artikel
  menghubungkan ukuran dan kepadatan dengan mekanisme \emph{sharing,
  matching, and learning}, tetapi juga membahas biaya konsentrasi
  berlebihan. Literatur jaringan kota dan \emph{agglomeration
  externality fields} dipakai untuk menjelaskan gagasan bahwa
  tempat-tempat dapat meminjam ukuran (\emph{borrow size}) dalam ruang
  metropolitan yang lebih luas. Namun, penjumlahan ukuran beberapa pusat
  saja tidak dianggap setara dengan manfaat satu kota besar; jaringan,
  koordinasi fungsi, konektivitas, dan kerja sama antarlembaga
  diperlukan untuk mengubah agregat menjadi integrasi (TXT lines
  226-266; PDF pp.~3-4).
\item
  \textbf{Laporan sumber:} Literatur kota second-tier digunakan untuk
  membahas bias first-city, akumulasi fungsi dan investasi di ibu kota,
  variasi struktur sistem perkotaan nasional, serta ketimpangan
  kapasitas ekonomi, fungsional, dan kelembagaan. Artikel juga
  menempatkan isu ini dalam literatur perencanaan strategis, kapasitas
  pengorganisasian regional, pemerintahan metropolitan, kompleksitas
  tata kelola, identitas metropolitan, dan relasi pusat-periferi (TXT
  lines 169-191, 291-324, 633-647; PDF pp.~2-5, 12).
\item
  \textbf{Interpretasi penulis:} Tinjauan literatur diarahkan untuk
  menggabungkan dua percakapan yang sering diperlakukan terpisah:
  keuntungan integrasi pada skala metropolitan dan kerugian struktural
  kota second-tier. Dari sana, artikel membangun dugaan bahwa insentif
  untuk berintegrasi dapat lebih besar pada kota second-tier, sekaligus
  mengakui bahwa kemampuan melaksanakannya dapat lebih lemah.
\item
  \textbf{Batas informasi:} Artikel tidak melaporkan protokol pencarian
  literatur, basis data yang ditelusuri, kata kunci, kriteria
  inklusi-eksklusi, jumlah korpus, atau prosedur telaah sistematis.
  Karena itu, bagian literaturnya adalah sintesis naratif dan
  argumentatif, bukan systematic review yang dapat direplikasi dari TXT.
\end{itemize}

\subsection{Method Used}\label{method-used-51}

\begin{itemize}
\tightlist
\item
  \textbf{Laporan sumber:} Desain artikel menggabungkan pembahasan
  konseptual, perbandingan deskriptif lintas kota berbasis data
  sekunder, dan pembacaan contoh-contoh kebijakan atau tata kelola.
  Analisis empiris berfokus pada pertambahan relatif populasi dan indeks
  fungsi perkotaan; kapasitas sumber daya, kelembagaan, dan politik
  dibahas melalui literatur serta contoh (TXT lines 195-214, 487-493;
  PDF pp.~3, 8-9).
\item
  \textbf{Laporan sumber:} Kerangka sampel dimulai dari 26 kota
  first-tier dan 124 kota second-tier dalam daftar ESPON/SGPTD (2012).
  Kota second-tier yang tidak tercakup studi fungsi perkotaan BBSR
  dikeluarkan; kota second-tier yang berada dalam wilayah metropolitan
  BBSR kota first-tier yang lebih besar beserta pasangan first-tier-nya
  juga dikeluarkan. Bern diganti dengan Zurich, ibu kota negara tanpa
  kota second-tier dikeluarkan, dan kota yang GDP-nya lebih tinggi
  daripada ibu kotanya tidak diperlakukan sebagai second-tier. Hasil
  akhirnya adalah 22 first-tier dan 71 second-tier (TXT lines 357-368;
  PDF p.~6).
\item
  \textbf{Laporan sumber:} Untuk populasi, artikel membandingkan
  populasi \emph{morphological urban area} (MUA), yaitu kawasan
  terbangun yang menyatu dari kota inti, dengan populasi
  \emph{metropolitan area} (MA) BBSR yang dibangun berdasarkan akses 60
  menit menuju klaster fungsi metropolitan tingkat atas. Ukuran
  pertambahan dihitung sebagai populasi MA dibagi populasi MUA (TXT
  lines 415-424; PDF pp.~6-7).
\item
  \textbf{Laporan sumber:} Uji pembanding menggunakan \emph{functional
  urban area} (FUA) ESPON yang berbasis wilayah perjalanan kerja dan
  membandingkannya dengan MUA. Untuk fungsi, indeks BBSR pada unit
  \emph{Local Administrative Unit} tingkat LAU2 diagregasikan pada skala
  kota dan wilayah metropolitan dengan perhitungan analog terhadap
  Gambar 1 (TXT lines 444-451, 468-473; PDF pp.~8-9).
\item
  \textbf{Inferensi pengolahan:} TXT tidak menunjukkan wawancara,
  observasi lapangan, eksperimen, model regresi, uji signifikansi, atau
  estimasi efek kausal. Karena itu, kata ``gains'' dan ``potential''
  dalam artikel harus dipahami sebagai perbedaan agregasi spasial dan
  potensi kebijakan, bukan bukti bahwa integrasi menyebabkan pertumbuhan
  atau peningkatan kapasitas.
\end{itemize}

\subsection{Dataset}\label{dataset-51}

\begin{itemize}
\tightlist
\item
  \textbf{Laporan sumber:} Dataset utama berasal dari data sekunder
  ESPON dan BBSR. BBSR (2011) mencakup studi 125 wilayah metropolitan
  Eropa dan menyediakan indeks fungsi perkotaan pada unit LAU2; fungsi
  tingkat atas yang disebut artikel mencakup politik, ekonomi, ilmu
  pengetahuan, transportasi, dan kebudayaan (TXT lines 336-343, 468-473;
  PDF pp.~5, 9).
\item
  \textbf{Laporan sumber:} Delimitasi spasial MUA dan FUA bersumber dari
  ESPON (2007), sedangkan MA BBSR menggunakan akses 60 menit ke klaster
  fungsi metropolitan tingkat atas. Catatan artikel juga menjelaskan
  PUSH sebagai area yang dapat dicapai dengan mobil dalam 45 menit dan
  PIA sebagai penggabungan FUA dengan tumpang tindih PUSH minimum 33
  persen (TXT lines 419-423, 870-882; PDF pp.~7-8, 16).
\item
  \textbf{Laporan sumber:} Variabel yang dipakai dalam perbandingan
  adalah populasi pada delimitasi berbeda, rasio populasi MA terhadap
  MUA atau FUA terhadap MUA, indeks kinerja fungsional pada skala MUA
  dan MA, serta proporsi potensi MA yang berada di luar FUA untuk Gambar
  3 (TXT lines 415-424, 444-451, 657-666; PDF pp.~7-8, 12-13).
\item
  \textbf{Batas informasi:} TXT tidak menyediakan baris data mentah,
  tahun observasi tiap variabel, kamus variabel, formula pembobotan
  indeks BBSR, aturan penanganan data hilang, kode pengolahan, atau
  ketidakpastian statistik. Tahun-tahun yang muncul dalam nama laporan
  dan referensi, termasuk 2005, 2007, 2011, 2012, dan 2016, tidak
  menetapkan satu periode observasi bersama untuk seluruh variabel;
  bagian Funding juga menyebut 2012-2015 sebagai periode riset asli,
  bukan sebagai tahun observasi variabel.
\end{itemize}

\subsection{Population / Sample}\label{population-sample-51}

\begin{itemize}
\tightlist
\item
  \textbf{Inferensi pengolahan:} Unit analisis yang tampak dari desain
  perbandingan adalah kota inti dan wilayah metropolitan yang terkait
  dengannya; artikel tidak menggunakan individu, rumah tangga,
  perusahaan, atau responden survei. Sampel analitis terdiri atas 22
  kota first-tier dan 71 kota second-tier, total 93 kota; daftar nama
  lengkap disajikan dalam Table 1 (TXT lines 357-368, 388-411; PDF
  pp.~6-7).
\item
  \textbf{Laporan sumber:} Kota yang dipilih mencakup beragam ukuran dan
  prospek sosial-ekonomi dari daftar ESPON/SGPTD, tetapi hanya kota yang
  memenuhi ketersediaan data BBSR dan aturan pemisahan
  first-tier-second-tier yang dipakai dalam perbandingan. Contoh
  pasangan yang dikeluarkan karena keterkaitan metropolitan adalah
  Rotterdam-Amsterdam, Malmo-Copenhagen, dan Antwerp-Brussels (TXT lines
  291-304, 357-368; PDF pp.~4-6).
\item
  \textbf{Interpretasi penulis:} Perbandingan dimaksudkan untuk membaca
  pola lintas sistem perkotaan Eropa, khususnya kondisi yang lebih
  merugikan kota second-tier di negara dengan ibu kota dominan atau
  dalam wilayah perkotaan padat. Artikel tidak mengklaim bahwa semua
  kota second-tier memiliki urgensi integrasi yang sama (TXT lines
  294-304, 845-850; PDF pp.~5, 16-17).
\item
  \textbf{Batas informasi:} Tidak ada prosedur sampling probabilistik,
  tingkat respons, kriteria pemilihan kasus kualitatif yang formal, atau
  analisis representativitas sampel terhadap seluruh kota Eropa yang
  dilaporkan dalam TXT.
\end{itemize}

\subsection{Dependent Variables}\label{dependent-variables-51}

\begin{itemize}
\tightlist
\item
  \textbf{Batas konseptual:} Artikel tidak menetapkan variabel dependen
  dalam arti model kausal atau regresi. Integrasi metropolitan juga
  tidak dioperasionalkan sebagai satu variabel perlakuan tunggal,
  melainkan sebagai proses dengan dimensi spasial-fungsional,
  politik-kelembagaan, dan simbolik-kultural (TXT lines 217-225; PDF
  p.~3).
\item
  \textbf{Laporan sumber:} Ukuran hasil yang paling dekat dengan
  variabel dependen adalah rasio potensi populasi, yaitu populasi MA
  dibagi populasi MUA, dan rasio pembanding FUA terhadap MUA. Nilai yang
  lebih besar dibaca sebagai pertambahan relatif massa demografis ketika
  wilayah sekitar diintegrasikan pada delimitasi yang lebih luas (TXT
  lines 419-424, 444-447; PDF pp.~7-8).
\item
  \textbf{Laporan sumber:} Ukuran hasil kedua adalah rasio indeks
  kinerja fungsi metropolitan pada skala yang lebih luas terhadap indeks
  fungsi MUA. Artikel memakai indeks fungsi BBSR yang diagregasikan dari
  LAU2, bukan ukuran langsung keberhasilan tata kelola atau
  kesejahteraan (TXT lines 468-484; PDF pp.~9-10).
\item
  \textbf{Laporan sumber:} Gambar 3 memakai proporsi populasi dan fungsi
  metropolitan yang berada di luar FUA sebagai ukuran potensi yang belum
  terintegrasi secara fungsional; artikel melaporkan sekitar 44 persen
  untuk populasi dan 38 persen untuk fungsi pada kota second-tier (TXT
  lines 657-666; PDF pp.~12-13).
\item
  \textbf{Inferensi pengolahan:} Kapasitas pengorganisasian regional,
  efisiensi berbagi sumber daya, pengaruh politik, identitas, dan kerja
  sama pusat-periferi merupakan konstruk argumentatif dalam artikel,
  tetapi tidak memiliki variabel dependen numerik yang dilaporkan untuk
  mengujinya secara langsung.
\end{itemize}

\subsection{Results}\label{results-51}

\begin{itemize}
\tightlist
\item
  \textbf{Laporan sumber:} Pada ukuran populasi MUA-ke-MA, kota
  second-tier secara konsisten menunjukkan pertambahan relatif yang
  lebih besar daripada kota first-tier. Pencilan yang disebut pada
  kelompok second-tier adalah Dusseldorf, Cardiff, dan Maribor, yang
  berada dalam wilayah polisentris dengan kota second-tier lain di
  dekatnya; Ljubljana dan Zurich disebut sebagai pencilan relatif di
  kelompok first-tier. Salerno memiliki rasio 28,6 kali dan tidak
  digambarkan (TXT lines 419-440; PDF pp.~7-8).
\item
  \textbf{Laporan sumber:} Untuk delimitasi FUA-ke-MUA, median yang
  dilaporkan adalah 1,77 untuk second-tier dan 1,22 untuk first-tier.
  Artikel menyebut hasil ini mengonfirmasi pola populasi pada delimitasi
  lain, dengan catatan bahwa skala integrasi tetap perlu dibatasi agar
  integrasi fungsional, kelembagaan, dan kultural tetap layak (TXT lines
  444-451; PDF p.~8).
\item
  \textbf{Batas verifikasi angka:} Ada ketidaksesuaian internal pada
  median populasi. Uraian hasil menyebut median first-tier 1,53 dan
  second-tier 3,35, sedangkan caption Gambar 1 menyebut first-tier 1,78
  dan second-tier 3,35. Kedua angka dicatat sebagaimana muncul dalam
  TXT; review ini tidak memilih salah satunya tanpa data atau sumber
  lain (TXT lines 432-440; PDF pp.~7-8).
\item
  \textbf{Laporan sumber:} Pada indeks fungsi metropolitan, artikel
  kembali melaporkan pertambahan proporsional yang lebih tinggi pada
  kota second-tier. Pencilan yang disebut adalah Toulon di dekat
  Marseilles dan Lens di dekat Lille; Maribor memiliki rasio 43 kali dan
  Toulon 22 kali, keduanya tidak digambarkan. Artikel memperingatkan
  bahwa pada kasus tersebut keuntungan agregasi dapat lebih banyak
  memperkuat kota second-tier yang lebih besar daripada mitra yang lebih
  kecil (TXT lines 468-484, 500-503; PDF pp.~9-10).
\item
  \textbf{Batas verifikasi angka:} Caption Gambar 2 menyebut median
  first-tier 1,24 dan second-tier 1,78. Uraian hasil menyebut ``median
  values (1.78 versus 1.07 in first-tier cities)'', sehingga angka
  first-tier yang tersurat dalam uraian adalah 1,07, berbeda dari 1,24
  pada caption. TXT tidak menjelaskan perbedaan tersebut; review ini
  mempertahankan kedua angka tanpa memilih salah satunya (TXT lines
  474-484, 500-503; PDF pp.~9-10).
\item
  \textbf{Laporan sumber:} Gambar 3 menunjukkan bahwa sekitar 44 persen
  potensi populasi dan 38 persen potensi fungsi pada wilayah
  metropolitan second-tier berada di luar FUA yang sudah terintegrasi
  secara fungsional. Kota first-tier disebut menyisakan sedikit potensi
  di luar FUA yang telah mereka kuasai (TXT lines 657-666; PDF
  pp.~12-13).
\item
  \textbf{Laporan sumber:} Untuk kapasitas sumber daya dan politik,
  artikel tidak menyajikan estimasi numerik baru. Artikel mengutip
  penelitian dan menggunakan contoh Manchester, West of England,
  Lund-Malmo, Prancis, serta proses devolusi Manchester untuk
  menunjukkan kemungkinan berbagi sumber daya, koordinasi, serta
  penguatan suara politik (TXT lines 505-570, 593-630; PDF pp.~9-12).
\end{itemize}

\subsection{Findings}\label{findings-51}

\begin{itemize}
\tightlist
\item
  \textbf{Interpretasi penulis:} Potensi kenaikan massa demografis dan
  fungsional melalui penggabungan skala metropolitan relatif lebih besar
  bagi kota second-tier daripada kota first-tier. Penjelasan yang
  diajukan adalah bahwa wilayah sekitar kota second-tier sering
  menyimpan penduduk dan fungsi yang tersebar, sehingga pusat-pusatnya
  dapat ``meminjam ukuran'' melalui integrasi; sebaliknya, kota
  first-tier cenderung mengonsentrasikan fungsi tingkat atas di inti dan
  menimbulkan bayangan aglomerasi pada sekitarnya (TXT lines 454-484;
  PDF pp.~8-10).
\item
  \textbf{Interpretasi penulis:} Integrasi dapat membantu kota
  second-tier mengatasi kelangkaan investasi publik dengan
  mengorganisasi dan berbagi sumber daya yang sudah ada, serta dapat
  memperbesar bobot demografis-ekonominya dalam negosiasi dengan tingkat
  pemerintahan yang lebih tinggi. Namun, manfaat ini menuntut kapasitas
  pengorganisasian regional dan rekam jejak koordinasi, bukan hanya
  pembentukan batas administratif baru (TXT lines 517-560, 603-630; PDF
  pp.~10-12).
\item
  \textbf{Interpretasi penulis:} Kondisi kota second-tier memberi dua
  kemungkinan yang berlawanan dalam relasi pusat-periferi. Kesamaan
  ukuran dan ketergantungan timbal balik dapat memudahkan kerja sama,
  tetapi tidak adanya kota pemimpin yang cukup kuat dapat mengurangi
  kemampuan memobilisasi sumber daya, memediasi konflik, dan membangun
  identitas metropolitan bersama (TXT lines 754-797; PDF pp.~15-16).
\item
  \textbf{Inferensi pengolahan:} Bukti paling langsung dalam artikel
  menyangkut potensi demografis dan fungsional yang dihitung dari
  agregasi wilayah. Klaim tentang kapasitas politik, kelembagaan, dan
  simbolik didukung oleh teori, literatur, dan ilustrasi kasus, tetapi
  tidak diuji dengan desain pengukuran yang sebanding dengan dua ukuran
  kuantitatif tersebut.
\end{itemize}

\subsection{Conclusions}\label{conclusions-51}

\begin{itemize}
\tightlist
\item
  \textbf{Laporan sumber:} Artikel menyimpulkan bahwa manfaat relatif
  dari massa demografis dan fungsional lebih besar pada kota second-tier
  dalam perbandingan yang dilakukan. Untuk manfaat bertindak sebagai
  entitas besar dan terkoordinasi, artikel merumuskan kondisi yang dapat
  mendukung berbagi sumber daya dan peningkatan pengaruh politik
  berdasarkan sejarah perkotaan Eropa, tren kebijakan, dan contoh yang
  relevan (TXT lines 812-830; PDF p.~16).
\item
  \textbf{Laporan sumber:} Penulis membatasi klaimnya: tidak ada dasar
  untuk menyatakan bahwa integrasi selalu menguntungkan wilayah
  metropolitan second-tier. Potensi harus dieksplorasi dan disesuaikan
  dengan konteks sejarah, politik, geografis, dan sosial-ekonomi
  setempat (TXT lines 831-844; PDF pp.~16-17).
\item
  \textbf{Laporan sumber:} Konsep ``second-tier metropolitan region''
  dinilai berguna untuk kebijakan dan penelitian, terutama pada negara
  dengan ibu kota dominan dan sistem perkotaan tersentralisasi serta
  pada kota yang berada dalam wilayah urbanisasi besar dan padat.
  Kerangka kelembagaan nasional yang berbeda membuat kemampuan integrasi
  dan pengaruh kebijakan antarwilayah tidak setara (TXT lines 845-863;
  PDF pp.~17-18).
\item
  \textbf{Laporan sumber:} Kesimpulan terakhir menyerukan penelitian
  studi kasus yang rinci dan kerja komparatif yang luas untuk menilai
  perkembangan relevansi kebijakan konsep tersebut dari waktu ke waktu
  serta menghadapkan potensi integrasi pada tantangan sosial, ekonomi,
  dan lingkungan kota-kota Eropa (TXT lines 864-868; PDF p.~17).
\end{itemize}

\subsection{Contributions}\label{contributions-51}

\begin{itemize}
\tightlist
\item
  \textbf{Laporan sumber:} Artikel menyumbang pada diskusi strategi
  pengembangan kota second-tier dengan menghubungkan kerugian khas
  kota-kota tersebut dengan potensi integrasi metropolitan. Abstrak
  secara khusus menyebut penilaian empiris terhadap pertambahan massa
  demografis dan fungsional serta penggunaan contoh untuk kapasitas
  kelembagaan dan politik (TXT lines 109-119; PDF p.~1).
\item
  \textbf{Interpretasi penulis:} Kontribusi substantifnya adalah
  menunjukkan bahwa metropolitanisasi dapat menjadi strategi untuk
  mengompensasi keterbatasan ukuran dan fungsi kota second-tier, sambil
  memindahkan perhatian perencanaan dari kota inti saja ke wilayah
  perkotaan yang saling bergantung.
\item
  \textbf{Laporan sumber:} Artikel menautkan argumen tersebut pada tiga
  perdebatan perencanaan dan tata kelola: bagaimana mengelola
  heterogenitas wilayah perkotaan dan sumber daya fungsional agregat;
  bagaimana kompleksitas membatasi ruang lingkup pemerintahan serta
  pembangunan imajinasi bersama; dan bagaimana menyeimbangkan relasi
  pusat-periferi untuk kerja sama metropolitan (TXT lines 633-647,
  838-844; PDF pp.~12-14, 16-17).
\item
  \textbf{Inferensi pengolahan:} Kontribusi artikel lebih kuat sebagai
  sintesis konseptual, perbandingan deskriptif, dan agenda perencanaan
  daripada sebagai estimasi kausal tentang efek integrasi atau evaluasi
  kelembagaan yang terukur.
\end{itemize}

\subsection{Research Gap}\label{research-gap-51}

\begin{itemize}
\tightlist
\item
  \textbf{Batas struktur sumber:} Artikel tidak memiliki heading khusus
  ``research gap''. Poin di bawah ini diringkas dari pertanyaan, batas
  klaim, dan saran penelitian yang dinyatakan dalam isi serta
  kesimpulan.
\item
  \textbf{Laporan sumber:} Penulis menyatakan perlunya penelitian lebih
  lanjut mengenai kondisi dan tipe wilayah tempat integrasi lebih
  penting, karena pertanyaan tentang besaran manfaat dan bentuk
  kelembagaan yang tepat tidak memiliki jawaban generik (TXT lines
  276-285, 812-818; PDF pp.~4, 16).
\item
  \textbf{Laporan sumber:} Artikel mengidentifikasi kebutuhan akan studi
  kasus rinci dan kerja komparatif luas untuk menilai bagaimana
  relevansi kebijakan konsep berkembang dari waktu ke waktu serta
  bagaimana potensi integrasi menghadapi tantangan sosial, ekonomi, dan
  lingkungan (TXT lines 864-868; PDF p.~17).
\item
  \textbf{Inferensi pengolahan:} Dari desain artikel, masih terbuka
  kebutuhan untuk mengoperasionalkan dan menguji secara langsung jalur
  kelembagaan, politik, serta simbolik yang hanya dibahas melalui
  literatur dan contoh. Artikel juga belum memberikan dasar empiris yang
  lengkap untuk memilih batas atas skala integrasi atau menentukan
  bentuk organisasi metropolitan yang paling efektif.
\end{itemize}

\subsection{Limitations}\label{limitations-51}

\begin{itemize}
\tightlist
\item
  \textbf{Laporan sumber:} Penulis menyatakan bahwa pengukuran spektrum
  tingkat kerugian atau urgensi antarwilayah berada di luar cakupan
  artikel. Mereka juga mengingatkan perlunya batas atas skala teritorial
  agar integrasi fungsional dan kelembagaan tetap layak serta faktor
  kultural tetap bermakna (TXT lines 291-304, 444-451; PDF pp.~5, 8).
\item
  \textbf{Laporan sumber:} Argumen dibatasi oleh keragaman sejarah,
  politik, ekonomi, kerangka kelembagaan, sumber daya keuangan, otonomi
  fiskal, dan legitimasi demokratis di negara-negara Eropa. Penulis
  secara eksplisit menolak klaim universal mengenai manfaat integrasi
  (TXT lines 831-863; PDF pp.~16-18).
\item
  \textbf{Inferensi pengolahan:} Sampel analitis bergantung pada irisan
  daftar ESPON/SGPTD dan ketersediaan data BBSR, lalu dipersempit dengan
  aturan eksklusi. TXT tidak menyajikan analisis sensitivitas terhadap
  seleksi kota, standardisasi indeks, atau perbedaan delimitasi selain
  pembanding MUA-MA dan MUA-FUA.
\item
  \textbf{Inferensi pengolahan:} Tidak adanya model kausal, periode
  observasi yang jelas, ukuran ketidakpastian, pengukuran langsung
  kapasitas kelembagaan/politik/simbolik, dan prosedur komparasi kasus
  yang formal membatasi interpretasi hasil sebagai potensi, bukan dampak
  aktual integrasi.
\item
  \textbf{Batas provenance:} Angka median pada uraian dan caption Gambar
  1 serta Gambar 2 tidak sepenuhnya konsisten. Selain itu, teks TXT
  tidak memuat visual lengkap grafik, sehingga distribusi seluruh
  pengamatan tidak dapat diaudit hanya dari hasil ekstraksi (TXT lines
  432-440, 474-484, 500-503; PDF pp.~7-10).
\end{itemize}

\subsection{Challenges}\label{challenges-51}

\begin{itemize}
\tightlist
\item
  \textbf{Laporan sumber:} Kota second-tier sering menghadapi kelangkaan
  sumber daya metropolitan, pengabaian dalam investasi transportasi,
  fungsi tingkat atas, dan layanan kolektif, serta pengaruh politik yang
  lemah pada tingkat pemerintahan yang lebih tinggi (TXT lines 487-528;
  PDF pp.~8-10).
\item
  \textbf{Laporan sumber:} Integrasi fungsional justru lebih mendesak
  bagi second-tier karena proporsi potensi populasi dan fungsi yang
  belum terintegrasi lebih besar, tetapi kota-kota ini secara historis
  lebih lemah dalam mengarahkan organisasi spasial-fungsional wilayahnya
  (TXT lines 650-677; PDF pp.~12-13).
\item
  \textbf{Laporan sumber:} Fragmentasi kelembagaan, perbedaan
  kepentingan, kompleksitas jaringan, dan kesulitan membangun identitas
  metropolitan bersama dapat menghambat perencanaan. Wilayah yang lebih
  sederhana mungkin memiliki ruang untuk visi yang lebih proaktif,
  tetapi tetap perlu menguji ambang kompleksitas yang dapat dikelola
  (TXT lines 720-749; PDF pp.~14-15).
\item
  \textbf{Laporan sumber:} Hubungan yang terlalu asimetris dengan kota
  inti dapat membuat mitra enggan bekerja sama. Sebaliknya, wilayah
  second-tier mungkin tidak memiliki kota pemimpin yang mampu
  menyediakan sumber daya teknis dan keuangan, memediasi konflik, dan
  memobilisasi identitas bersama. Tantangannya adalah menjaga
  keseimbangan antara kepemimpinan dan dominasi (TXT lines 754-797; PDF
  pp.~15-16).
\item
  \textbf{Laporan sumber:} Perbedaan kerangka nasional, kekuasaan
  politik, sumber daya keuangan, otonomi fiskal, dan legitimasi
  demokratis menentukan apakah wilayah metropolitan dapat bertindak
  sebagai aktor kebijakan yang relevan (TXT lines 851-863; PDF
  pp.~17-18).
\end{itemize}

\subsection{Practical Implications}\label{practical-implications-51}

\begin{itemize}
\tightlist
\item
  \textbf{Interpretasi penulis:} Perencana dan pembuat kebijakan perlu
  memperlakukan integrasi sebagai proses yang harus memberi kualitas
  integratif pada agregasi ukuran melalui jaringan yang berfungsi,
  koordinasi fungsi, konektivitas, dan kerja sama antarlembaga;
  menjumlahkan populasi kota-kota yang berdekatan saja tidak cukup (TXT
  lines 255-275; PDF pp.~4-5).
\item
  \textbf{Laporan sumber:} Strategi praktis yang dibahas mencakup
  pembangunan infrastruktur penghubung, pencarian sinergi antar tempat,
  dan pemanfaatan sumber daya yang sudah ada secara bersama, misalnya
  universitas atau bandara, terutama ketika investasi nasional baru
  sulit diperoleh (TXT lines 517-547, 650-666; PDF pp.~10, 12-13).
\item
  \textbf{Laporan sumber:} Koordinasi tidak harus selalu berbentuk
  tingkat pemerintahan metropolitan formal. Namun, organisasi yang
  terlalu longgar atau hanya berfokus pada layanan sempit dapat kurang
  mampu melaksanakan kebijakan dan mengatasi konflik; cakupan,
  kedalaman, kapasitas implementasi, dan koordinasi kebijakan lebih
  penting daripada label formal semata (TXT lines 548-560; PDF p.~10).
\item
  \textbf{Laporan sumber:} Perencanaan dapat memandang seluruh wilayah
  sebagai wilayah perkotaan heterogen dan multisentris, bukan sebagai
  inti yang memerintah periferi secara hierarkis. Artikel mengaitkan
  pendekatan ini dengan penyediaan kualitas lingkungan binaan, amenitas,
  transportasi publik, infrastruktur, dan kelayakhunian pada skala
  metropolitan (TXT lines 685-714; PDF pp.~13-14).
\item
  \textbf{Laporan sumber:} Tata kelola perlu mengelola kompleksitas
  secara pragmatis, memperluas akses warga kepada pemerintah lokal,
  membuat jaringan dan kemitraan metropolitan terlihat serta dapat
  dihubungkan, mendukung mobilitas lintas wilayah, dan membangun narasi
  identitas yang dapat dipahami warga (TXT lines 720-749; PDF
  pp.~14-15).
\item
  \textbf{Laporan sumber:} Kota inti perlu mengambil peran memimpin
  tanpa mendominasi. Konsesi terhadap bobot ekonomi-demografisnya dan
  pembatasan kekuasaan pengambilan keputusan dipandang dapat membantu
  membangun persepsi keuntungan yang simetris (TXT lines 782-797; PDF
  pp.~15-16).
\end{itemize}

\subsection{Applications}\label{applications-51}

\begin{itemize}
\tightlist
\item
  \textbf{Inferensi pengolahan:} Prosedur kuantitatif artikel dapat
  dipakai sebagai diagnosis awal: bandingkan MUA dengan MA atau FUA
  untuk mengidentifikasi massa penduduk dan fungsi yang dapat
  dimobilisasi, lalu ukur bagian potensi yang belum terintegrasi secara
  fungsional. Artikel menggunakan prosedur ini pada 22 kota first-tier
  dan 71 kota second-tier (TXT lines 357-368, 415-424, 657-666; PDF
  pp.~6-7, 12-13).
\item
  \textbf{Laporan sumber:} Manchester Airport dipakai sebagai contoh
  sumber daya ekonomi dan simbolik yang mendorong skala city-region
  sebagai pusat kekuatan alternatif terhadap London. West of England
  dipakai sebagai contoh promosi bersama Bristol, Bath, universitas, dan
  simpul ekonomi-industrial dalam strategi ekonomi-spasial terintegrasi
  (TXT lines 529-547; PDF pp.~10-11).
\item
  \textbf{Laporan sumber:} Kerja sama Lund-Malmo dipaparkan sebagai
  upaya membangun citra wilayah pengetahuan dan mengembangkan lahan
  perkotaan melalui kemitraan pemerintah kota dan otoritas perencanaan
  regional, dengan memanfaatkan dua laboratorium riset yang dibiayai Uni
  Eropa (TXT lines 539-547; PDF p.~10-11).
\item
  \textbf{Laporan sumber:} SCOT di Prancis digunakan sebagai contoh
  kerangka perencanaan antarpemerintah kota untuk infrastruktur, ruang
  alam, perumahan, pekerjaan, dan layanan kolektif. Artikel mencatat
  bahwa semua kota besar Prancis tercakup metropolitan SCOT dengan
  pengecualian Paris, yang menghadapi kompleksitas politik tersendiri
  (TXT lines 561-570; PDF p.~11).
\item
  \textbf{Laporan sumber:} Contoh lain meliputi strategi ``City of
  Cities'' untuk Milan, perkembangan perencanaan metropolitan di Turin,
  Lyon, dan Hanover, proses devolusi Manchester, serta perubahan posisi
  periferi Toulouse. Contoh-contoh ini berfungsi sebagai ilustrasi jalur
  integrasi dan tata kelola, bukan sebagai evaluasi komparatif dengan
  ukuran keberhasilan yang seragam (TXT lines 695-714, 782-809; PDF
  pp.~13-16).
\end{itemize}

\subsection{Future Research
(Optional)}\label{future-research-optional-51}

\begin{itemize}
\tightlist
\item
  \textbf{Laporan sumber:} Penulis menyerukan studi kasus yang rinci dan
  penelitian komparatif luas untuk menguji kondisi, tipe wilayah, dan
  perkembangan relevansi kebijakan integrasi second-tier dari waktu ke
  waktu (TXT lines 812-818, 864-868; PDF pp.~16-18).
\item
  \textbf{Laporan sumber:} Arah penelitian yang disebut langsung dalam
  pembahasan meliputi penilaian ambang kompleksitas yang masih
  memungkinkan visi dan identitas metropolitan, penentuan bentuk
  organisasi yang memiliki cakupan serta kapasitas implementasi memadai,
  dan pengujian keseimbangan relasi pusat-periferi (TXT lines 276-285,
  548-560, 730-749; PDF pp.~4, 10, 14-15).
\item
  \textbf{Inferensi pengolahan:} Replikasi yang transparan dengan data
  baris, definisi indeks yang konsisten, pemeriksaan angka median, dan
  pelaporan sensitivitas antar-delimitasi akan diperlukan untuk
  memperkuat hasil komparatif sebelum potensi demografis-fungsional
  diterjemahkan sebagai dampak integrasi.
\end{itemize}

\subsection{Provenance Notes}\label{provenance-notes-51}

\begin{itemize}
\tightlist
\item
  Review ini hanya menggunakan file
  \texttt{source/journal/B08-cardoso-meijers-2017-secondary-yet-metropolitan-the-challenges-of-metropolitan-integration-for-second-tier-cities-10.1080-14649357.2017.1371789.txt}.
  Seluruh 1.076 baris yang tersedia dibaca, termasuk blok metadata
  ekstraksi, teks artikel, catatan kaki, disclosure statement, funding,
  notes on contributors, dan daftar referensi.
\item
  Blok metadata TXT menyatakan bahwa sumber asalnya adalah PDF 21
  halaman, diekstraksi pada 31 Agustus 2026 dengan
  \texttt{pdftotext\ -layout}. Review ini tidak mengambil informasi dari
  PDF, situs jurnal, atau karya-karya yang tercantum di daftar referensi
  secara terpisah (TXT lines 1-28).
\item
  Locator utama memakai nomor baris TXT. Nomor halaman PDF dicantumkan
  hanya ketika penanda halaman muncul di dalam TXT. Caption Gambar 1-3
  dan uraian tekstualnya tersedia, tetapi visual grafik lengkap tidak
  tersedia sebagai data yang dapat diperiksa di dalam TXT.
\item
  Laporan sumber, interpretasi penulis, dan inferensi pengolahan diberi
  label secara eksplisit. Klaim tentang institusi, kapasitas politik,
  dan contoh wilayah diperlakukan sebagai argumen atau ilustrasi
  artikel, bukan sebagai hasil pengukuran primer yang dibuat oleh review
  ini.
\item
  Ketidaksesuaian angka median dicatat tanpa rekonsiliasi: median
  first-tier pada populasi dilaporkan 1,53 di uraian dan 1,78 di caption
  Gambar 1; angka fungsi diuraikan sebagai 1,78 dan 1,07 dengan pemetaan
  yang ambigu, sedangkan caption Gambar 2 menyebut first-tier 1,24 dan
  second-tier 1,78 (TXT lines 432-440, 474-484, 500-503).
\item
  Informasi yang tidak tersedia dalam TXT tidak diisi dengan dugaan,
  terutama tahun observasi yang seragam, data mentah, formula indeks
  lengkap, uji statistik, prosedur telaah literatur, representativitas
  sampel, dan hasil kelembagaan-politik yang terukur.
\end{itemize}

\section{Review B09}\label{review-b09}

\subsection{Summarized Abstract}\label{summarized-abstract-52}

\textbf{Temuan sumber:} Artikel ini memperkenalkan metode untuk
mengidentifikasi sub-pusat metropolitan berdasarkan pendekatan interaksi
yang menggunakan arus komuter. Hasilnya dibandingkan dengan pendekatan
berbasis kepadatan pekerjaan dan/atau ambang ukuran pekerjaan. Metode
diterapkan pada kawasan metropolitan (MA) Roma dan Milan; menurut hasil
artikel, metode yang diajukan memiliki kesesuaian yang lebih baik
daripada pendekatan berbasis kepadatan pekerjaan. Hasil identifikasi
sangat sensitif terhadap metode dan, sampai tingkat tertentu, terhadap
konsep pusat yang digunakan. {[}Abstrak, hlm. 177{]}

\textbf{Interpretasi penulis:} Identifikasi sub-pusat perlu
memperhatikan dimensi fungsional, bukan hanya konsentrasi morfologis
pekerjaan. {[}Abstrak; Introduction, hlm. 177{]}

\subsection{Problem Statement}\label{problem-statement-52}

\textbf{Temuan sumber:} Struktur spasial kawasan metropolitan semakin
kompleks; pekerjaan tersebar ke seluruh wilayah, sementara eksternalitas
aglomerasi melampaui satu inti perkotaan. Identifikasi sub-pusat disebut
sebagai langkah awal untuk memahami organisasi spasial metropolitan,
tetapi riset empiris terutama masih menggunakan kepadatan pekerjaan,
meskipun perhatian terhadap polisentrisitas fungsional meningkat.
Penulis juga menyatakan bahwa pengukuran dampak polisentrisitas terhadap
kinerja ekonomi metropolitan masih kurang diteliti secara empiris.
{[}Introduction, hlm. 177{]}

\textbf{Interpretasi penulis:} Masalah metodologisnya adalah kurangnya
alat identifikasi sentralitas metropolitan yang menggunakan ukuran
interaksi untuk melengkapi atau menggantikan ukuran berbasis kepadatan.
Perbandingan langsung antara pendekatan interaksi dan kepadatan juga
dinilai masih kurang. {[}Introduction, hlm. 177{]}

\subsection{Objectives}\label{objectives-52}

\textbf{Temuan sumber:} Tujuan artikel adalah mengusulkan pendekatan
metodologis untuk mengidentifikasi sub-pusat metropolitan, menerapkannya
pada MA Roma dan Milan, serta membandingkan hasilnya dengan dua metode
lain yang berbasis kepadatan dan ukuran pekerjaan. {[}Introduction, hlm.
177{]}

\textbf{Interpretasi penulis:} Pendekatan tersebut dimaksudkan untuk
mengisi kesenjangan dalam deskripsi struktur perkotaan dengan
mengidentifikasi sentralitas melalui relasi fungsional antarnode.
{[}Introduction, hlm. 177{]}

\subsection{Literature Survey}\label{literature-survey-52}

\textbf{Temuan sumber:} Tinjauan membedakan metodologi statis/morfologis
dari metodologi interaksi. Metode morfologis yang dibahas meliputi
ambang pekerjaan dan kepadatan Giuliano dan Small, rasio pekerjaan
terhadap pekerja penduduk (E/R) Shearmur dan Coffey, puncak kepadatan
lokal, residu model ekonometrik kepadatan, kernel density, serta Local
Indicators of Spatial Autocorrelation (LISA). Literatur interaksi
menggunakan data arus, terutama arus komuter, untuk membaca hubungan dan
posisi node dalam jaringan metropolitan. {[}Literature review, hlm.
178--180{]}

\textbf{Temuan sumber:} Kerangka teoretis artikel terutama terkait teori
central place, yang menekankan hierarki fungsi antarnode. Literatur
network of cities yang menekankan relasi horizontal dan spesialisasi
antarpusat juga dibahas, tetapi penulis menilai teori central place
tetap relevan untuk satu MA. Pendekatan morfologis dan fungsional tidak
selalu bertentangan; studi ESPON yang dirujuk menggabungkan ukuran,
lokasi, dan konektivitas. {[}Literature review, hlm. 179--180{]}

\textbf{Interpretasi penulis:} Perbedaan metode berakar pada perbedaan
konsep pusat: kepadatan dan ukuran pekerjaan menangkap nodalitas
absolut, sedangkan penyediaan fungsi dan pekerjaan melebihi permintaan
penduduk menangkap sentralitas relatif terhadap wilayah sekitar.
{[}Literature review, hlm. 179--180{]}

\subsection{Method Used}\label{method-used-52}

\textbf{Temuan sumber:} Unit analisis yang dipilih adalah tingkat
perkotaan; sub-pusat dipilih dari seluruh munisipalitas yang berada
dalam MA. Italia belum memiliki definisi resmi MA, sehingga batas MA
diambil dari Boix dan Veneri (2009), yang menggunakan algoritme yang
terinspirasi metodologi US Federal Register 1990. {[}Methods and data,
hlm. 180{]}

\textbf{Temuan sumber:} MA dimodelkan sebagai jaringan dengan
munisipalitas sebagai node dan arus komuter sebagai link. Untuk setiap
munisipalitas dihitung flow-centrality ratio (FCi), yaitu perbandingan
in-degree terhadap out-degree, dan directional dominance index (DIIi),
yaitu ukuran keterlibatan node yang membandingkan perjalanan masuk
dengan rata-rata perjalanan masuk dalam MA. Munisipalitas dengan FCi dan
DIIi lebih besar dari 1 diperlakukan sebagai kandidat sub-pusat tingkat
kedua. {[}Methods and data, hlm. 180--181{]}

\textbf{Temuan sumber:} Selanjutnya dihitung productive completeness
(PCi), yaitu rasio jumlah sektor lima digit dengan sedikitnya satu
pekerjaan pada 2001 di suatu munisipalitas terhadap rata-rata jumlah
sektor lima digit dalam MA. Munisipalitas dengan FCi \textgreater{} 1,
DIIi \textgreater{} 1, dan PCi \textgreater{} 1 dipilih sebagai
sub-pusat metropolitan tingkat pertama. Node yang terpilih dan
bersebelahan digabungkan berdasarkan kontiguitas dan interaksi;
agregatnya harus memiliki self-containment arus komuter
sekurang-kurangnya 15\%. {[}Methods and data, hlm. 180--181; catatan
kaki 4, hlm. 181{]}

\textbf{Temuan sumber:} Hasil tersebut dibandingkan dengan metode
Giuliano--Small (GS), yang memakai sedikitnya 10.000 pekerjaan dan
kepadatan sedikitnya 10 pekerjaan per acre, serta metode
Shearmur--Coffey (SC), yang memakai rasio E/R \textgreater{} 1 dan
sedikitnya 5.000 pekerjaan. Evaluasi kesesuaian dilakukan dengan
mengestimasi fungsi kepadatan pekerjaan polisentris: bentuk eksponensial
negatif logaritmik dan model gravitasi logaritmik, kemudian
membandingkan adjusted R2. {[}Empirical findings, hlm. 181--183;
Evaluation of methods, hlm. 183--184{]}

\subsection{Dataset}\label{dataset-52}

\textbf{Temuan sumber:} Data arus komuter antarmunisipalitas berasal
dari Istituto Nazionale di Statistica (Istat) dan menggunakan data 2001.
Untuk setiap MA tersedia matriks yang memuat jumlah komuter dari setiap
munisipalitas ke semua munisipalitas lain dalam MA. PCi menggunakan data
Sensus Istat dengan klasifikasi sektoral NACE, pada tingkat sektor lima
digit dan tahun 2001. {[}Methods and data, hlm. 180--181{]}

\textbf{Temuan sumber:} Evaluasi kepadatan menggunakan kepadatan
pekerjaan bruto per kilometer persegi pada setiap munisipalitas serta
jarak dari CBD dan sub-pusat terkonsolidasi; sumber tidak memberikan
uraian dataset terpisah untuk setiap variabel evaluasi tersebut. Nama
sumber dan periode tambahan untuk variabel evaluasi: Tidak disebutkan
dalam file. {[}Evaluation of methods, hlm. 183--184{]}

\subsection{Population / Sample}\label{population-sample-52}

\textbf{Temuan sumber:} Populasi analisis adalah seluruh munisipalitas
dalam dua MA yang dipilih. Tabel 1 mencatat 200 observasi untuk Roma dan
597 observasi untuk Milan pada masing-masing indikator FCi, DIIi, dan
PCi. {[}Table 1, hlm. 181{]}

\textbf{Temuan sumber:} Desain sampel probabilistik, responden individu,
atau survei sampel: Tidak disebutkan dalam file. Unit observasinya
adalah munisipalitas dan arus di antara munisipalitas. {[}Methods and
data, hlm. 180--181{]}

\subsection{Dependent Variables}\label{dependent-variables-52}

\textbf{Temuan sumber:} Variabel dependen tunggal konvensional pada
tahap identifikasi: Tidak disebutkan dalam file. FCi, DIIi, dan PCi
adalah indikator seleksi sub-pusat. Pada tahap evaluasi, variabel yang
dijelaskan dalam persamaan adalah D, yaitu kepadatan pekerjaan bruto per
kilometer persegi di setiap munisipalitas; jarak dari CBD dan sub-pusat
digunakan dalam fungsi kepadatan eksponensial dan gravitasi. {[}Methods
and data, hlm. 180--181; Evaluation of methods, hlm. 183{]}

\textbf{Interpretasi pengolahan:} Dengan demikian, keluaran utama tahap
identifikasi adalah status/kumpulan sub-pusat, sedangkan adjusted R2
pada tahap evaluasi dipakai untuk membandingkan seberapa besar metode
identifikasi menjelaskan variasi kepadatan pekerjaan. Ini adalah
pembacaan struktur analisis dari persamaan dan uraian sumber, bukan
istilah variabel dependen yang dipakai sebagai judul analisis oleh
penulis. {[}Evaluation of methods, hlm. 183--184{]}

\subsection{Results}\label{results-52}

\textbf{Temuan sumber:} Statistik deskriptif pada Tabel 1 mencatat hal
berikut. Untuk Roma, FCi memiliki 200 observasi, minimum 0,02, maksimum
4,62, mean 0,38, dan simpangan baku 0,51; DIIi memiliki rentang
0,00--195,94, mean 1,49, dan simpangan baku 13,88; PCi memiliki rentang
0,02--2,02, mean 0,33, dan simpangan baku 0,31. Untuk Milan, FCi
memiliki 597 observasi, rentang 0,00--4,63, mean 0,58, dan simpangan
baku 0,47; DIIi memiliki rentang 0,00--235,07, mean 1,11, dan simpangan
baku 9,71; PCi memiliki rentang 0,01--2,01, mean 0,37, dan simpangan
baku 0,28. {[}Table 1, hlm. 181{]}

\textbf{Temuan sumber:} Pendekatan fungsional mengidentifikasi empat
munisipalitas pusat di Roma dan mengonsolidasikannya menjadi tiga
sub-pusat berdasarkan kontiguitas dan interaksi. Tabel A.1 mencantumkan
Rome, Pomezia, Latina, dan Civitavecchia; Tabel 2 menampilkan hasil
pusat terkonsolidasi sebagai Rome, Latina, dan Civitavecchia. Di Milan,
11 munisipalitas diidentifikasi dan diagregasikan menjadi sembilan
sub-pusat terkonsolidasi. Tabel A.1 mencantumkan Milan, Como, Pavia,
Crema, Lodi, Segrate, Saronno, Trezzano sul Naviglio, Cantù, Legnano,
dan Monza; Tabel 2 menampilkan agregasi tersebut dengan Milan terdiri
atas tiga munisipalitas dan pusat lain masing-masing satu. {[}Empirical
findings, hlm. 181; Table 2, hlm. 181; Table A.1, hlm. 184{]}

\textbf{Temuan sumber:} Dengan metode SC, 35 munisipalitas di MA Milan
dan enam munisipalitas di MA Roma memenuhi syarat awal, lalu
diagregasikan masing-masing menjadi 15 dan tiga sub-pusat
terkonsolidasi. Dengan metode GS, tidak ada sub-pusat yang
teridentifikasi di Roma, sedangkan di Milan ditemukan satu sub-pusat
yang terdiri atas tiga munisipalitas. {[}Evaluation of methods, hlm.
182--183; Table 2, hlm. 181; Table A.1, hlm. 184{]}

\textbf{Temuan sumber:} Adjusted R2 untuk bentuk log eksponensial di
Roma adalah 0,525 dengan metode fungsional dan 0,393 dengan SC; di Milan
nilainya 0,671 dengan metode fungsional, 0,674 dengan SC, dan 0,426
dengan GS. Untuk bentuk log gravitasi, nilainya di Roma adalah 0,546
dengan metode fungsional dan 0,396 dengan SC; di Milan 0,703 dengan
metode fungsional, 0,711 dengan SC, dan 0,414 dengan GS. {[}Table 3,
hlm. 183{]}

\subsection{Findings}\label{findings-52}

\textbf{Temuan sumber:} Peta dan uraian artikel menggambarkan Milan
sebagai struktur yang lebih polisentris, dengan beberapa node pusat
mengelilingi munisipalitas pivot, sedangkan Roma memiliki lebih sedikit
node pusat dan peran munisipalitas pivot yang lebih menonjol.
{[}Empirical findings, hlm. 181; Fig. 1--2, hlm. 181--182{]}

\textbf{Interpretasi penulis:} Pendekatan fungsional lebih restriktif
daripada SC karena mencari node dengan sentralitas neto positif dan
kapasitas sebagai pemasok fungsi, bukan sekadar node berkepadatan
pekerjaan tinggi di sekitar pivot. Untuk Roma, pendekatan fungsional
menjelaskan varians kepadatan pekerjaan jauh lebih besar daripada SC;
pada kedua spesifikasi, selisih adjusted R2 sekurang-kurangnya 13,2 poin
persentase. Untuk Milan, metode fungsional dan SC menjelaskan bagian
varians yang secara substantif serupa, sedangkan GS menghasilkan nilai
yang lebih rendah. {[}Evaluation of methods, hlm. 182--183{]}

\textbf{Interpretasi penulis:} Hasil identifikasi sub-pusat bergantung
kuat pada metode yang dipakai dan konsep pusat yang hendak diukur.
Perbandingan adjusted R2 juga dapat dipengaruhi oleh struktur data,
termasuk ukuran dan jumlah unit spasial, fungsi yang dipilih, serta
struktur perkotaan aktual. {[}Evaluation of methods, hlm. 183--184{]}

\subsection{Conclusions}\label{conclusions-52}

\textbf{Interpretasi penulis:} Penulis menyimpulkan bahwa metode
fungsional yang diusulkan bekerja baik secara teoretis dan empiris serta
memberikan kesesuaian yang lebih baik ketika fungsi kepadatan pekerjaan
polisentris diestimasi. Pendekatan ini dinilai lebih tepat untuk
pusat-pusat dalam MA Italia yang terbentuk melalui koalesensi, sedangkan
ukuran kepadatan lebih dekat dengan gagasan desentralisasi pekerjaan
dari CBD yang padat. {[}Concluding remarks, hlm. 183--184{]}

\textbf{Interpretasi penulis:} Metode GS dinilai tidak memadai dalam
konteks metropolitan Italia, terutama ketika ambangnya tidak diperbarui
berdasarkan pengetahuan lokal. Pilihan metode harus bergantung pada
apakah analisis diarahkan pada polisentrisitas morfologis atau
fungsional. {[}Concluding remarks, hlm. 184{]}

\subsection{Contributions}\label{contributions-52}

\textbf{Temuan sumber:} Artikel mengusulkan prosedur yang menggabungkan
unsur interaksi dan fungsi: FCi serta DIIi untuk relasi dan dominasi
arus, lalu PCi untuk kelengkapan struktur sektoral. Prosedur itu
mengidentifikasi node tanpa menetapkan ambang absolut populasi atau
pekerjaan. Penulis menyatakan bahwa karya ini berkontribusi pada
literatur tentang pendekatan yang menilai sentralitas melalui fungsi dan
pekerjaan yang disediakan melebihi jumlah yang diminta penduduk sekitar.
{[}Methods and data, hlm. 180--181; Concluding remarks, hlm. 184{]}

\textbf{Interpretasi pengolahan:} Kontribusi metodologis yang dapat
diidentifikasi dari TXT adalah menjadikan konsep pusat fungsional dapat
dibandingkan secara eksplisit dengan dua prosedur berbasis ambang dalam
dua MA Italia. {[}Introduction, hlm. 177; Empirical findings, hlm.
181--183{]}

\subsection{Research Gap}\label{research-gap-52}

\textbf{Temuan sumber:} Penulis menyatakan bahwa perhatian terhadap
relasi fungsional antarnode di dalam MA telah meningkat, tetapi jauh
lebih sedikit penelitian yang mengidentifikasi sentralitas metropolitan
memakai ukuran interaksi sebagai pengganti pendekatan berbasis
kepadatan. Studi komparatif antara hasil kedua pendekatan tersebut juga
disebut masih kurang. Selain itu, topik pengaruh polisentrisitas
terhadap kinerja ekonomi metropolitan dinyatakan masih kurang diteliti
secara empiris. {[}Introduction, hlm. 177{]}

\textbf{Inferensi pengolahan/batas klaim:} Ini adalah kesenjangan yang
dirumuskan penulis dalam artikel; TXT tidak menyediakan tinjauan
sistematis dengan protokol pencarian yang memungkinkan verifikasi
independen atas kelengkapan klaim tersebut. {[}Introduction; Literature
review, hlm. 177--180{]}

\subsection{Limitations}\label{limitations-52}

\textbf{Temuan sumber:} Penulis menyebut ketersediaan data sebagai
batasan serius: data arus biasanya tidak seluas dan tidak sesering
diperbarui seperti data stok. Masalah unit areal yang dapat dimodifikasi
(MAUP) dapat muncul karena tingkat spasial analisis bergantung pada
tingkat tempat data fungsional seperti komuter tersedia. {[}Concluding
remarks, hlm. 184{]}

\textbf{Temuan sumber:} Jumlah sektor lima digit hanya merupakan
pendekatan terhadap fungsi ekonomi yang disediakan pusat; hubungan
antara klasifikasi NACE dan fungsi ekonomi perlu diteliti lebih lanjut,
walaupun hal itu bukan tujuan artikel. Penulis juga mencatat bahwa
metode berbasis ambang sulit diterapkan lintas konteks karena perbedaan
struktur permukiman dan unit spasial yang tersedia. {[}Methods and data,
hlm. 181; Concluding remarks, hlm. 184{]}

\textbf{Temuan sumber:} Arus komuter pekerjaan tidak mencakup seluruh
pergerakan untuk konsumsi, studi, dan rekreasi. Karena itu, kandidat
awal dari FCi dan DIIi disebut sub-pusat tingkat kedua, dan PCi
diperlukan untuk menangkap seperangkat fungsi perkotaan yang lebih luas.
{[}Methods and data, hlm. 180--181{]}

\subsection{Challenges}\label{challenges-52}

\textbf{Temuan sumber:} Tantangan metodologis yang dibahas dalam TXT
meliputi pemilihan tingkat spasial ketika data arus hanya tersedia pada
unit tertentu, penentuan ambang yang dapat bersifat diskresioner, serta
penggabungan munisipalitas yang bersebelahan menjadi pusat
terkonsolidasi dengan syarat interaksi dan self-containment minimal
15\%. {[}Literature review, hlm. 178; Methods and data, hlm. 180--181{]}

\textbf{Temuan sumber:} Perbandingan antarhasil juga menghadapi
persoalan bahwa ambang GS berasal dari struktur wilayah Los Angeles dan
diterapkan pada unit yang berbeda dari munisipalitas Italia. SC dapat
memilih munisipalitas di sabuk pertama dan kedua yang memiliki banyak
pekerjaan karena kedekatannya dengan pivot, meskipun belum tentu
merupakan pusat fungsional. {[}Empirical findings, hlm. 181--183{]}

\subsection{Practical Implications}\label{practical-implications-52}

\textbf{Interpretasi penulis:} Identifikasi sub-pusat memberi
pengetahuan tentang organisasi spasial metropolitan yang diperlukan
untuk kebijakan tata ruang. Hasil identifikasi juga dapat membantu
mengenali tempat yang menjadi prioritas investasi publik dan tempat yang
membentuk efek aglomerasi di dalam wilayah. {[}Introduction, hlm. 177{]}

\textbf{Interpretasi pengolahan:} Alat pengukuran struktur perkotaan
yang lebih andal diperlukan untuk menguji efek polisentrisitas terhadap
kinerja ekonomi metropolitan, tetapi artikel ini terutama mengusulkan
dan mengevaluasi metode identifikasi, bukan mengestimasi efek kebijakan
tersebut. {[}Introduction, hlm. 177; Concluding remarks, hlm. 184{]}

\subsection{Applications}\label{applications-52}

\textbf{Temuan sumber:} Metode diterapkan pada seluruh munisipalitas
dalam MA Roma dan Milan menggunakan arus komuter 2001, lalu hasilnya
dibandingkan dengan GS dan SC. Aplikasi lain di luar dua MA tersebut
tidak disebutkan dalam file. {[}Methods and data, hlm. 180--181;
Empirical findings, hlm. 181--183{]}

\subsection{Future Research
(Optional)}\label{future-research-optional-52}

\textbf{Temuan sumber:} Penulis mengusulkan penelitian lanjutan untuk
menguji apakah terdapat jenis sub-pusat yang berbeda menurut asal
pembentukannya, yaitu desentralisasi atau koalesensi, dan apakah jenis
tersebut menghasilkan dampak yang berbeda terhadap struktur
metropolitan. Penulis juga mengusulkan pengkajian hubungan
antarsub-pusat serta apakah fungsi mereka bersifat alternatif atau
komplementer, dengan mengeksplorasi koneksi konsumsi, rekreasi, dan
pertukaran antarperusahaan. {[}Concluding remarks, hlm. 184{]}

\subsection{Provenance Notes}\label{provenance-notes-52}

\textbf{Catatan verifikasi:} Hanya satu TXT aktual berprefiks B09
ditemukan di bawah \texttt{source/}, yaitu
\texttt{source/journal/B09-veneri-2013-the-identification-of-sub-centres-in-two-italian-metropolitan-areas-a-functional-approach-10.1016-j.cities.2012.04.006.txt}.
Nama tersebut cocok dengan artikel Paolo Veneri, ``The identification of
sub-centres in two Italian metropolitan areas: A functional approach,''
\emph{Cities} 31 (2013), hlm. 177--185, DOI
\texttt{10.1016/j.cities.2012.04.006}. {[}Metadata artikel dan halaman
177--185{]}

\textbf{Catatan ekstraksi:} TXT memuat seluruh teks yang tersedia dari
617 baris, termasuk blok \texttt{{[}PDF\ Metadata{]}}, teks artikel,
Appendix A, dan References. Metadata menyebut PDF sumber
\texttt{journal/B09-veneri-2013-the-identification-of-sub-centres-in-two-italian-metropolitan-areas-a-functional-approach-10.1016-j.cities.2012.04.006.pdf},
diekstrak pada 2026-08-31 dengan \texttt{pdftotext\ -layout}; PDF
berjumlah 9 halaman, tidak terenkripsi, berukuran 595.276 x 793.701 pts,
berukuran 617632 bytes, dan versi PDF 1.7. Metadata juga menyatakan
Custom Metadata: no, Metadata Stream: yes, Tagged: no, UserProperties:
no, Suspects: no, Form: none, JavaScript: no, dan Optimized: no. {[}PDF
Metadata, baris 1--20{]}

\textbf{Batasan evidence:} Klaim review ini hanya diproses dari TXT B09
tersebut. Ekstraksi memiliki artefak kecil berupa ligatur/karakter rumus
dan susunan dua kolom pada beberapa bagian, tetapi teks naratif, angka
tabel, locator halaman, appendix, dan kesimpulan terbaca. TXT dan PDF
sumber tidak diubah. Catatan kaki menyatakan pandangan yang disampaikan
adalah pandangan penulis dan tidak selalu mencerminkan pandangan OECD.
{[}PDF Metadata; catatan kaki, hlm. 177{]}

\section{Review B10}\label{review-b10}

\textbf{Sumber:} Paolo Veneri, ``Urban Polycentricity and the Costs of
Commuting: Evidence from Italian Metropolitan Areas,'' \emph{Growth and
Change}, Vol. 41 No.~3 (September 2010), hlm. 403--429.

\subsection{Summarized Abstract}\label{summarized-abstract-53}

\textbf{Laporan sumber:} Artikel menguji peran polisentrisitas,
kekompakan, diversifikasi fungsional, dan ukuran kawasan metropolitan
terhadap biaya komuter. Biaya eksternal diproksikan dengan emisi CO2 per
kapita, sedangkan biaya privat diproksikan dengan waktu perjalanan.
Derajat polisentrisitas diukur secara dinamis menggunakan arus komuter
dan alat analisis jejaring sosial pada basis data 82 metropolitan areas
(MAs) di Italia. Abstrak melaporkan bahwa MAs yang lebih polisentris
lebih baik dalam biaya privat dan eksternal; kekompakan berkaitan dengan
biaya lingkungan yang lebih rendah tetapi biaya privat yang lebih
tinggi; ukuran berkaitan dengan kedua biaya yang lebih tinggi; dan
diversifikasi fungsional tidak signifikan secara statistik.
Sosiodemografi juga berperan (hlm. 403, Abstract).

\textbf{Laporan sumber:} Pada bagian hasil dan kesimpulan, koefisien
polisentrisitas terhadap waktu perjalanan memang negatif tetapi tidak
signifikan, sehingga sumber menyatakan bahwa polisentrisitas tidak
terbukti memengaruhi durasi perjalanan secara signifikan (hlm. 420--423,
Table 7 dan ``Concluding Remarks'').

\textbf{Inferensi pengolahan:} Ringkasan ini mempertahankan klaim
abstrak, tetapi memberi kualifikasi hasil utama: dukungan empiris yang
paling kuat untuk polisentrisitas berada pada biaya eksternal CO2, bukan
pada pengurangan waktu perjalanan.

\subsection{Problem Statement}\label{problem-statement-53}

\textbf{Laporan sumber:} Organisasi spasial kota memengaruhi pola
mobilitas dan menimbulkan biaya sosial privat maupun eksternal.
Polisentrisitas telah memperoleh peran normatif sebagai alat perencanaan
untuk daya saing, kohesi sosial, dan keberlanjutan lingkungan, tetapi
hipotesis tersebut belum cukup diuji secara empiris. Konsep
polisentrisitas, khususnya pada skala metropolitan, juga belum
didefinisikan dan diukur secara memadai. Banyak pengukuran sebelumnya
bersifat statis dan morfologis, dengan fokus pada distribusi pekerjaan
atau penduduk, tanpa menangkap interaksi antar-node perkotaan (hlm.
403--405).

\textbf{Interpretasi penulis:} Hasil empiris yang saling bertentangan
mengenai jarak, waktu, dan penggunaan transportasi publik dapat sebagian
dijelaskan oleh perbedaan definisi dan pengukuran polisentrisitas.
Desentralisasi pekerjaan tidak selalu menghasilkan polisentrisitas; ia
juga dapat menghasilkan kawasan perkotaan yang tersebar. Karena itu,
untuk konteks Eropa, pusat perlu dipahami sebagai tempat yang
menjalankan fungsi sentral dan mengorganisasi wilayah sekitarnya, bukan
hanya sebagai lokasi dengan kepadatan pekerjaan tinggi (hlm. 406--409).

\subsection{Objectives}\label{objectives-53}

\textbf{Laporan sumber:} Tujuan utama artikel adalah menguji apakah
derajat polisentrisitas perkotaan yang tinggi bersifat baik dalam
kaitannya dengan biaya mobilitas dan karena itu dapat dipandang sebagai
tujuan perencanaan yang diinginkan. Artikel juga menilai peran ukuran
perkotaan, dispersi/kekompakan, dan keragaman fungsional sebagai
determinan biaya komuter (hlm. 404--405).

\textbf{Laporan sumber:} Untuk menjalankan tujuan tersebut, penulis
mendefinisikan dan mengidentifikasi sub-centres dari perspektif
interaksi, lalu mengukur polisentrisitas dengan pendekatan dinamis
berbasis analisis jejaring sosial. Dua indikator polisentrisitas
digunakan untuk robustness: ukuran relatif jumlah sub-centres dan
Special Functional Polycentricity (PSF) (hlm. 404, 408--411).

\subsection{Literature Survey}\label{literature-survey-53}

\textbf{Laporan sumber:} Literatur yang diringkas mengelompokkan
karakteristik spasial yang memengaruhi komuter ke dalam polisentrisitas,
dispersi, pencampuran fungsi, dan ukuran (hlm. 405). Argumen yang
mendukung polisentrisitas mencakup hipotesis co-location yang
mendekatkan rumah dan pekerjaan; kemungkinan turunnya jarak karena sewa
dekat sub-centre lebih rendah daripada kawasan monosen\-tris; dukungan
terhadap transportasi massal karena sub-centre memberi skala minimum;
serta pengurangan kemacetan melalui distribusi penduduk, pekerjaan, dan
fungsi ke beberapa sub-centre (hlm. 405--406).

\textbf{Interpretasi penulis:} Bukti terdahulu tentang polisentrisitas
belum konsisten. Sebagian studi melaporkan desentralisasi pekerjaan
berkaitan dengan jarak komuter dan penggunaan mobil yang lebih tinggi,
atau perjalanan lebih panjang; studi lain melaporkan durasi komuter yang
lebih pendek. Hubungan dengan penggunaan transportasi publik juga
berlawanan antar-studi (hlm. 406). Penulis mengaitkan sebagian
ketidakkonsistenan tersebut dengan penggunaan ukuran berbasis kepadatan
dan penyamaan desentralisasi dengan polisentrisitas.

\textbf{Laporan sumber:} Untuk dispersi, literatur yang diringkas
cenderung mengaitkan kepadatan rendah dan pola terpencar dengan jarak
perjalanan lebih panjang, ketidakefisienan transportasi massal, dan
konsumsi energi yang lebih tinggi, tetapi hubungan kepadatan dengan
waktu perjalanan masih tidak jelas. Diversifikasi fungsional atau
keseimbangan pekerjaan-perumahan dipaparkan sebagai mekanisme yang dapat
memperpendek jarak dan meningkatkan penggunaan transit. Deregulasi
pembangunan rumah, perubahan preferensi, dan meluasnya transportasi
privat juga dipaparkan sebagai faktor yang dapat menghasilkan perjalanan
lebih panjang; stok perumahan yang lebih baru dikaitkan dengan emisi
komuter yang lebih tinggi dalam studi sistem tenaga kerja lokal Italia
(hlm. 407--408).

\textbf{Laporan sumber:} Penulis membedakan evolusi polisentrisitas
melalui desentralisasi dari CBD ke sub-centre baru, yang disebut lebih
intens di kota-kota Amerika Serikat, dan melalui integrasi pusat-pusat
yang sebelumnya mandiri ke dalam satu MA melalui coalescence atau
incorporation, yang disebut lebih lazim di kota-kota Eropa. Pendekatan
morfologis melihat distribusi pekerjaan di dua atau lebih pusat,
sedangkan pendekatan fungsional memperhatikan distribusi fungsi dan
hierarki node (hlm. 408--409).

\subsection{Method Used}\label{method-used-53}

\textbf{Laporan sumber:} Unit analisis ekonometrik adalah 82 MAs Italia
dalam desain cross-section. Jaringan internal setiap MA dimodelkan
sebagai jaringan municipality: municipality adalah node dan arus komuter
adalah link (hlm. 409--410, 420). Data arus komuter tahun 2001 digunakan
untuk menghitung in-degree setiap municipality. In-degree mengukur
jumlah link yang menuju node tertentu dan digunakan sebagai ukuran
sentralitas.

\textbf{Laporan sumber:} Untuk ukuran pertama, penulis menggunakan
matriks empat nearest neighbours yang didikotomikan; catatan 1
menetapkan ambang minimum lima komuter, dengan nilai link 1 untuk empat
municipality tetangga terdekat dan 0 untuk kasus lain. Municipality pada
persentil ke-95 in-degree dipilih sebagai sub-centres. Jumlah
sub-centres kemudian dibandingkan dengan populasi MA dalam scatter plot;
jarak vertikal setiap observasi dari garis interpolasi ditafsirkan
sebagai derajat polisentrisitas relatif terhadap ukuran MA (hlm.
409--411, Figure 1, Note 1).

\textbf{Laporan sumber:} Ukuran kedua adalah PSF, indeks 0--1 dari Green
(2007), yang dihitung dari dispersi in-degree nodal dan kepadatan
jaringan. Nilai yang lebih tinggi menunjukkan derajat polisentrisitas
fungsional yang lebih tinggi. Penulis menyatakan bahwa PSF menangkap
hierarki jaringan sekaligus kekuatan interaksi melalui kepadatan
jaringan (hlm. 410--411, Equation 2, Figure 2).

\textbf{Laporan sumber:} Biaya eksternal dihitung melalui Environmental
Impact of Mobility (EIM), yang menggabungkan jumlah komuter menurut
pasangan asal-tujuan dan moda, faktor emisi CO2 per penumpang-kilometer,
serta jarak jalan. Biaya privat dihitung melalui Average Commuting Time
(ACT), yaitu rata-rata waktu perjalanan komuter berdasarkan kelas durasi
dan nilai representatif 7,5; 22,5; 45; dan 75 menit (hlm. 413--416,
Equations 3--4, Table 5).

\textbf{Laporan sumber:} Dua model cross-section diestimasi. Model EIM
menggunakan spatial lag model dengan maximum likelihood dan matriks
spasial contiguity; model ACT menggunakan OLS dengan robust standard
errors. Seluruh variabel distandardisasi dengan mean nol dan unit
variance. Uji robust LM dan pemeriksaan pada Appendix Tables A3--A4
digunakan untuk menangani autokorelasi spasial, sedangkan korelasi dan
VIF pada Appendix Tables A1--A2 digunakan untuk memeriksa
multikolinearitas (hlm. 420, Table 7; hlm. 428--429, Appendix).

\textbf{Inferensi pengolahan:} Karena regresi menggunakan satu potongan
waktu dan observasi agregat pada tingkat MA, hasil yang diringkas di
sini paling aman dibaca sebagai hubungan statistik antar-karakteristik
spasial dan biaya komuter, bukan bukti kausal individual.

\subsection{Dataset}\label{dataset-53}

\textbf{Laporan sumber:} Analisis menggunakan data Istat Population
Census tahun 2001 untuk arus komuter, moda, waktu perjalanan, penduduk,
dan variabel terkait; data Istat Population Census tahun 1981 dan 2001
serta Istat Industry and Services Census tahun 2001 digunakan untuk
variabel independen (hlm. 411--416, 420, Table 6). Jarak dihitung dengan
GIS. Tabel 1 dan Tabel 3 mencatat sumber sebagai elaborasi atas data
sensus Istat; faktor emisi pada Table 4 bersumber dari Amici della Terra
(2005).

\textbf{Laporan sumber:} Tabel 1 mencatat 10.933.774 komuter: 1.118.017
pengguna transportasi publik (10,23 persen) dan 9.815.757 pengguna
transportasi privat (89,77 persen). Mobil sebagai pengemudi mencakup
6.952.851 orang (63,59 persen), mobil sebagai penumpang 511.400 (4,68
persen), motor atau skuter 689.287 (6,30 persen), dan sepeda, berjalan
kaki, atau moda lain 1.662.219 (15,20 persen) (hlm. 413, Table 1).

\textbf{Laporan sumber:} Table 2a menggunakan 9.271.555 komuter karena
tidak memasukkan komuter dengan sepeda atau berjalan kaki dan
mengeluarkan observasi yang moda atau waktunya tidak tercatat. Jarak
pada tabel ini adalah jarak garis lurus dari centroid ke centroid. Table
3 melaporkan jarak jalan yang dihitung dengan GIS; rata-rata waktu
seluruh komuter adalah 22 menit dan rata-rata jarak 7,1 kilometer.
Perbedaan definisi jarak dan cakupan observasi tersebut perlu
dipertahankan ketika membandingkan kedua tabel (hlm. 414, Tables 2a--3).

\textbf{Laporan sumber:} Faktor emisi yang digunakan adalah 35 gram CO2
per passenger-kilometre untuk kereta, 32 untuk trem, 21,3 untuk kereta
bawah tanah, 72 untuk bus kota, 26 untuk bus luar kota, 31 untuk bus
sekolah/perusahaan, 105 untuk mobil, 80 untuk motor, dan 0 untuk sepeda,
berjalan kaki, atau moda lain (hlm. 415, Table 4).

\textbf{Laporan sumber:} Table 6 mencantumkan variabel penjelas berikut:
\texttt{polyc\_green} sebagai PSF; \texttt{polyc\_nsr} sebagai indeks
relative sub-centre number; \texttt{gini\_area\_job} sebagai perbedaan
absolut antara share area dan share pekerjaan tiap municipality;
\texttt{ldensity} sebagai log kepadatan bruto penduduk per km2;
\texttt{gini\_pop\_job} sebagai perbedaan absolut antara share penduduk
dan share pekerjaan; \texttt{house\_age} sebagai proporsi rumah yang
dibangun setelah 1982; \texttt{rel\_pub\_comp} sebagai rasio kecepatan
rata-rata transportasi publik terhadap moda privat; \texttt{area}
sebagai luas MA; dummy Central dan Southern Italy;
\texttt{demo\_structure} sebagai rasio penduduk usia 40-64 terhadap usia
15-39; serta \texttt{graduates} sebagai share penduduk bergelar (hlm.
418, Table 6).

\textbf{Laporan sumber:} \texttt{gini\_area\_job} merepresentasikan
konsentrasi spasial pekerjaan, sedangkan distribusi pekerjaan dan
penduduk yang lebih merata pada \texttt{gini\_pop\_job} dipakai sebagai
pendekatan diversifikasi fungsional atau mixed land-use. Variabel
kontrol \texttt{house\_age}, daya saing transportasi publik, lokasi
geografis, struktur umur, dan pendidikan dimasukkan untuk menangkap
determinan non-spasial atau kontekstual biaya komuter (hlm. 417-419,
Table 6).

\subsection{Population / Sample}\label{population-sample-53}

\textbf{Laporan sumber:} Unit analisis penelitian adalah 82 Italian
metropolitan areas (MAs). MAs tersebut diidentifikasi dengan algoritme
fungsional yang didasarkan pada metodologi SMA dari Federal Register,
dengan adaptasi Clusa dan Roca; catatan 3 menyebut unit analisis ini
diidentifikasi dalam Boix dan Veneri (2009) (hlm. 411-412, Note 3). Di
dalam jaringan tiap MA, municipality diperlakukan sebagai node dan arus
komuter sebagai link (hlm. 409-410).

\textbf{Laporan sumber:} Tabel 1 mencakup 10.933.774 komuter dalam 82
MAs. Tabel 2a menggunakan 9.271.555 komuter setelah mengecualikan
perjalanan dengan sepeda atau berjalan kaki serta observasi yang moda
atau waktunya tidak tercatat (hlm. 413-414, Tables 1-2a).

\textbf{Tidak disebutkan dalam file:} Desain pengambilan sampel
probabilistik, ukuran populasi individu di luar cakupan 82 MAs, dan
kriteria sampling individual.

\textbf{Inferensi pengolahan:} Karena unit ekonometriknya adalah MA,
observasi regresi merupakan agregat kawasan, bukan sampel individu
komuter.

\subsection{Dependent Variables}\label{dependent-variables-53}

\textbf{Laporan sumber:} Penelitian menggunakan dua variabel dependen
untuk mewakili dua komponen biaya mobilitas:

\begin{itemize}
\tightlist
\item
  \textbf{Environmental Impact of Mobility (EIM):} indikator biaya
  eksternal yang mengestimasi rata-rata emisi CO2 per komuter
  berdasarkan arus dari municipality asal ke tujuan, moda transportasi,
  faktor emisi per passenger-kilometre, dan jarak jalan. Perhitungannya
  diberikan pada Equation 3 (hlm. 413-415).
\item
  \textbf{Average Commuting Time (ACT):} indikator biaya privat berupa
  rata-rata waktu perjalanan komuter. Nilai ini dihitung dari kelas
  durasi 0-15, 15-30, 30-60, dan lebih dari 60 menit, dengan nilai
  representatif berturut-turut 7,5; 22,5; 45; dan 75 menit (hlm.
  415-416, Equation 4 dan Table 5).
\end{itemize}

\textbf{Laporan sumber:} EIM diestimasi dengan spatial lag model dan ACT
dengan OLS ber-robust standard errors. Keterbatasan data membuat biaya
privat hanya diwakili ACT; biaya tiket dan bahan bakar tidak dihitung
dalam penelitian ini (hlm. 415-416, 420).

\textbf{Tidak disebutkan dalam file:} Satuan keluaran EIM yang
dilaporkan sebagai satu angka kawasan selain penjelasan bahwa indikator
tersebut mengukur jumlah rata-rata gas rumah kaca per komuter.

\subsection{Results}\label{results-53}

\textbf{Laporan sumber:} Statistik deskriptif menunjukkan dominasi
transportasi privat. Dari 10.933.774 komuter, 10,23 persen menggunakan
transportasi publik dan 89,77 persen menggunakan transportasi privat;
mobil sebagai pengemudi mencakup 63,59 persen (hlm. 413, Table 1). Pada
Table 3, rata-rata durasi transportasi publik adalah 38 menit dengan
jarak 10,4 km, sedangkan transportasi privat bermotor rata-rata 19 menit
dengan jarak 8,3 km; seluruh moda rata-rata 22 menit dan 7,1 km. Table
2a menggunakan cakupan dan definisi jarak yang berbeda, sehingga tidak
dapat dibaca sebagai pengukuran yang identik dengan Table 3 (hlm. 414,
Tables 2a-3).

\textbf{Laporan sumber:} Dalam dua spesifikasi EIM, kedua ukuran
polisentrisitas berkoefisien negatif dan signifikan:
\texttt{polyc\_green} = -0,299 (p = 0,008) dan \texttt{polyc\_nsr} =
-0,184 (p = 0,006). Dalam dua spesifikasi ACT, koefisiennya juga negatif
tetapi tidak signifikan: \texttt{polyc\_green} = -0,092 (p = 0,341) dan
\texttt{polyc\_nsr} = -0,077 (p = 0,224) (hlm. 420-421, Table 7). Setiap
spesifikasi menggunakan 82 observasi; adjusted R-squared berada pada
0,745-0,746 untuk EIM dan 0,786-0,789 untuk ACT (hlm. 421, Table 7).

\textbf{Laporan sumber:} Variabel konsentrasi pekerjaan
\texttt{gini\_area\_job} berkoefisien negatif dan signifikan pada EIM
(-0,332 dan -0,355; keduanya p = 0,000), tetapi positif dan signifikan
lemah pada ACT (0,122, p = 0,096; 0,117, p = 0,078). \texttt{ldensity}
berkoefisien negatif dan signifikan pada EIM (-0,413 dan -0,546; p =
0,000) serta positif dan signifikan pada ACT (0,488 dan 0,445; p =
0,000). \texttt{gini\_pop\_job} tidak signifikan pada spesifikasi
pertama EIM (-0,162, p = 0,115), signifikan negatif pada spesifikasi
kedua (-0,338, p = 0,000), dan tidak signifikan pada kedua spesifikasi
ACT (-0,116, p = 0,131; -0,132, p = 0,123) (hlm. 421, Table 7).

\textbf{Laporan sumber:} \texttt{rel\_public\_comp} tidak signifikan
pada EIM, tetapi negatif dan signifikan pada ACT di kedua spesifikasi
(-0,115; p = 0,047 dan p = 0,037). \texttt{area} positif dan signifikan
pada seluruh spesifikasi, dengan koefisien 0,542-0,595 untuk EIM dan
0,578-0,588 untuk ACT. Share \texttt{graduates} positif dan signifikan
pada seluruh spesifikasi, sedangkan \texttt{demo\_structure} tidak
signifikan (hlm. 421, Table 7).

\textbf{Laporan sumber:} Dummy Southern Italy positif pada kedua
spesifikasi EIM dan signifikan pada tingkat yang berbeda (0,537, p =
0,062; 0,705, p = 0,007), sementara dummy South tidak signifikan pada
ACT. Dummy Central Italy tidak signifikan pada tiga spesifikasi pertama,
tetapi pada spesifikasi ACT kedua tercatat -0,306 dengan tanda
\texttt{*} dan p = 0,064. \texttt{house\_age} umumnya tidak signifikan
dalam tabel, tetapi pada spesifikasi ACT kedua tercatat 0,129 dengan
tanda \texttt{*} dan p = 0,090. Narasi penulis menyatakan bahwa dummy
geografis tidak signifikan kecuali Southern MAs pada biaya eksternal dan
bahwa umur stok perumahan tidak memiliki efek yang signifikan (hlm.
421-422, Table 7; lihat juga hlm. 422).

\textbf{Laporan sumber:} Parameter spasial \texttt{rho} positif dan
signifikan dalam dua model spatial lag EIM (0,267, p = 0,008; 0,287, p =
0,004). Catatan 4 menyatakan uji robust LM mengarahkan penggunaan
spatial lag model dan residualnya tidak lagi menunjukkan autokorelasi
spasial; catatan 5 menyatakan Moran I untuk model ACT tidak signifikan
dengan beberapa matriks bobot (hlm. 420, 423, Notes 4-5; Appendix Tables
A3-A4).

\subsection{Findings}\label{findings-53}

\textbf{Interpretasi penulis:} Polisentrisitas berkaitan dengan biaya
eksternal komuter yang lebih rendah, terutama emisi CO2 rata-rata per
komuter. Hubungan ini didukung oleh dua ukuran polisentrisitas, dan
penulis menyatakan hasilnya juga konsisten dengan indikator morfologis
konsentrasi pekerjaan dalam MA (hlm. 420-423, Table 7 dan ``Concluding
Remarks'').

\textbf{Interpretasi penulis:} Polisentrisitas tidak terbukti
memengaruhi waktu perjalanan secara signifikan. Dengan demikian, sumber
tidak mendukung kesimpulan bahwa MA yang lebih polisentris pasti
memiliki perjalanan yang lebih singkat, meskipun arah koefisien ACT
negatif (hlm. 420-423, Table 7 dan ``Concluding Remarks'').

\textbf{Interpretasi penulis:} Kepadatan yang lebih tinggi berkaitan
dengan emisi CO2 yang lebih rendah tetapi durasi perjalanan yang lebih
panjang. Penulis membaca pola ini sebagai trade-off antara biaya
eksternal dan biaya internal, sehingga efek bersih kepadatan masih tidak
diketahui (hlm. 422).

\textbf{Interpretasi penulis:} Diversifikasi fungsional tidak memiliki
hubungan yang dapat ditegaskan secara statistik dengan biaya mobilitas.
Koefisiennya negatif pada model biaya eksternal tetapi kehilangan
signifikansi dalam satu dari dua spesifikasi, sedangkan pada model biaya
privat koefisien selalu tidak signifikan dan bertanda positif (hlm.
422).

\textbf{Interpretasi penulis:} Ukuran MA yang lebih besar berkaitan
dengan emisi CO2 dan waktu perjalanan yang lebih tinggi, yang dijelaskan
penulis melalui jarak rata-rata yang lebih panjang di MA berukuran besar
(hlm. 422).

\textbf{Interpretasi penulis:} Daya saing relatif transportasi publik
berkaitan dengan durasi komuter yang lebih singkat. Share lulusan
berkaitan dengan emisi CO2 dan waktu perjalanan yang lebih tinggi,
sedangkan struktur umur penduduk tidak menunjukkan peran yang
signifikan. MA di wilayah Selatan menunjukkan dampak CO2 yang lebih
tinggi dalam spesifikasi EIM (hlm. 422).

\textbf{Inferensi pengolahan:} Bukti paling kuat dalam hasil rinci
berada pada hubungan polisentrisitas dengan EIM, bukan pada pengurangan
ACT. Pernyataan abstrak bahwa MA yang lebih polisentris lebih baik dalam
biaya privat perlu dibaca bersama hasil Table 7 dan kesimpulan yang
menyatakan efek ACT tidak signifikan.

\subsection{Conclusions}\label{conclusions-53}

\textbf{Laporan sumber:} Sumber menyimpulkan bahwa MA kontemporer
memiliki struktur yang semakin polisentris, sementara paradigma
polisentris mempromosikan polisentrisitas sebagai alat perencanaan untuk
meningkatkan daya saing lokal dan keberlanjutan. Namun, konsep
polisentrisitas dan dampaknya terhadap daya saing serta keberlanjutan
masih dinyatakan belum cukup terdefinisi dan diteliti secara teoretis
maupun empiris (hlm. 423, ``Concluding Remarks'').

\textbf{Laporan sumber:} Berdasarkan 82 Italian MAs, polisentrisitas
dinilai sebagai model struktur perkotaan yang berkelanjutan terutama
dari perspektif lingkungan karena berkaitan dengan penurunan emisi CO2
akibat komuter. Hasil tersebut robust untuk dua ukuran polisentrisitas
dan konsisten dengan indikator morfologis konsentrasi pekerjaan, tetapi
polisentrisitas tidak tampak memengaruhi waktu perjalanan secara
signifikan (hlm. 423).

\textbf{Laporan sumber:} Dispersi perkotaan berkaitan dengan biaya
eksternal yang lebih tinggi sekaligus biaya privat yang lebih rendah.
Sumber menyatakan adanya trade-off antara penahanan biaya kolektif
mobilitas, seperti polusi dan gas rumah kaca, dan pemenuhan preferensi
perumahan di area atau lingkungan ber-kepadatan rendah (hlm. 423).

\subsection{Contributions}\label{contributions-53}

\textbf{Laporan sumber:} Sumber menyebut dua kontribusi utama. Pertama,
artikel memberikan definisi formal polisentrisitas pada skala
metropolitan dan pengukuran kuantitatif untuk analisis komparatif.
Kedua, artikel menjelaskan dampak derajat polisentrisitas perkotaan
terhadap dua komponen utama biaya komuter (hlm. 423, ``Concluding
Remarks'').

\textbf{Laporan sumber:} Konsep polisentrisitas dibangun dari perspektif
interaksi dan diukur dengan indikator berbasis social network analysis.
Artikel memperkenalkan indikator polisentrisitas baru serta menghitung
indikator lain untuk robustness dan kelengkapan analisis (hlm. 423).
Pada tingkat operasional, pendekatan tersebut mendefinisikan sub-centres
sebagai tempat yang menjalankan fungsi sentral dan mengorganisasi
wilayah di sekitarnya, bukan sekadar tempat dengan kepadatan pekerjaan
tinggi (hlm. 404, 408-410).

\subsection{Research Gap}\label{research-gap-53}

\textbf{Laporan sumber:} Sebelum analisis dilakukan, sumber
mengidentifikasi bahwa hipotesis mengenai polisentrisitas sebagai alat
untuk daya saing, kohesi sosial, dan keberlanjutan lingkungan belum
cukup diuji secara empiris. Konsep itu juga belum didefinisikan dan
diukur secara memadai pada skala metropolitan; banyak studi sebelumnya
memakai pendekatan statis dan morfologis yang berfokus pada sebaran
pekerjaan atau penduduk tanpa menangkap interaksi antar-node (hlm.
404-405).

\textbf{Interpretasi penulis:} Literatur tentang hubungan
polisentrisitas dan biaya mobilitas menghasilkan temuan yang kontras.
Sumber mengaitkan sebagian ketidakkonsistenan itu dengan definisi dan
pengukuran yang menyamakan desentralisasi dengan polisentrisitas, serta
dengan perbedaan antara evolusi melalui desentralisasi dan melalui
coalescence atau incorporation yang dianggap lebih lazim di kota-kota
Eropa (hlm. 406-409).

\textbf{Laporan sumber:} Pada bagian kesimpulan, sumber masih menyatakan
bahwa konsep polisentrisitas dan dampaknya terhadap daya saing serta
keberlanjutan perlu diteliti lebih lanjut secara teoretis dan empiris
(hlm. 423).

\textbf{Inferensi pengolahan:} Review ini tidak menetapkan research gap
tambahan di luar gap yang dinyatakan sumber.

\subsection{Limitations}\label{limitations-53}

\textbf{Laporan sumber:} Penulis mengakui bahwa data longitudinal lebih
disukai untuk menetapkan kausalitas, tetapi keterbatasan ketersediaan
data membuat penelitian menggunakan desain cross-section (hlm. 420).
Data biaya privat yang tersedia juga terbatas, sehingga penelitian hanya
menghitung ACT dan tidak memasukkan biaya tiket maupun bahan bakar (hlm.
415-416).

\textbf{Laporan sumber:} ACT merupakan pendekatan terhadap waktu
perjalanan yang dibangun dari kelas durasi dan nilai representatif 7,5;
22,5; 45; dan 75 menit, bukan dari durasi kontinu tiap komuter (hlm.
415-416, Table 5). Cakupan komuter pada Table 2a juga berbeda dari Table
1 dan Table 3 karena pengecualian perjalanan berjalan kaki atau sepeda
serta data moda atau waktu yang tidak tercatat (hlm. 414, Notes to Table
2a).

\textbf{Laporan sumber:} Penulis menyebut penggunaan dummy wilayah
sebagai cara untuk menangkap perbedaan ekonomi, budaya, dan teknologi
antarmakro-wilayah Italia secara agak kasar (hlm. 419).

\textbf{Inferensi pengolahan:} Desain cross-section dengan 82 unit
agregat membatasi pembacaan hasil sebagai hubungan statistik tingkat
kawasan. TXT tidak melaporkan desain sampling individu, sehingga
keterwakilan komuter individual dan generalisasi di luar 82 MAs tidak
dapat dinilai dari file ini.

\textbf{Tidak disebutkan dalam file:} Bagian tersendiri yang merangkum
seluruh keterbatasan penelitian secara formal.

\subsection{Challenges}\label{challenges-53}

\textbf{Laporan sumber:} Tantangan konseptual utama adalah
mendefinisikan dan mengukur polisentrisitas pada skala metropolitan.
Sumber harus membedakan pusat morfologis berbasis kepadatan pekerjaan
dari pusat fungsional yang menjalankan fungsi sentral dan mengorganisasi
wilayah sekitarnya (hlm. 404, 408-410).

\textbf{Laporan sumber:} Identifikasi sub-centres juga menghadapi
perbedaan evolusi kawasan. Desentralisasi dari CBD ke sub-centre baru
disebut lebih menonjol di kota-kota Amerika Serikat, sedangkan integrasi
pusat-pusat yang sebelumnya mandiri melalui coalescence atau
incorporation disebut lebih lazim di kota-kota Eropa. Hal ini membuat
pendekatan kepadatan sederhana kurang memadai untuk konteks Eropa (hlm.
408-409).

\textbf{Laporan sumber:} Bukti empiris sebelumnya yang kontras mengenai
jarak, durasi, dan penggunaan transportasi publik menjadi tantangan
untuk menguji hubungan polisentrisitas dan biaya mobilitas. Penulis
menanganinya dengan dua ukuran polisentrisitas, yaitu relative
sub-centre number dan PSF, serta dengan indikator morfologis konsentrasi
pekerjaan sebagai pembanding (hlm. 406-410, 420-423).

\textbf{Laporan sumber:} Ketergantungan spasial merupakan tantangan
teknis pada model EIM. Uji robust LM mengarahkan penggunaan spatial lag
model, sedangkan untuk model ACT catatan penulis menyatakan Moran I
tidak signifikan dengan beberapa matriks bobot (hlm. 420, Notes 4-5;
Appendix Tables A3-A4). Penulis juga melaporkan pemeriksaan korelasi dan
VIF; mean VIF adalah 2,96 dan nilai di atas 10 disebut dapat
mengindikasikan masalah kolinearitas (hlm. 428-429, Appendix Table A2).

\subsection{Practical Implications}\label{practical-implications-53}

\textbf{Interpretasi penulis:} Polisentrisitas dapat dipandang sebagai
tujuan perencanaan yang diinginkan untuk menekan biaya eksternal
mobilitas, khususnya emisi CO2. Namun, hasil penelitian tidak memberi
dasar untuk menjanjikan durasi perjalanan yang lebih singkat hanya
karena struktur kawasan lebih polisentris (hlm. 404-405, 420-423).

\textbf{Interpretasi penulis:} Perencanaan perlu mempertimbangkan
trade-off kepadatan. Kepadatan lebih tinggi berkaitan dengan biaya
lingkungan yang lebih rendah tetapi waktu perjalanan yang lebih tinggi,
sedangkan dispersi atau area berkepadatan rendah berkaitan dengan biaya
eksternal lebih tinggi dan biaya privat lebih rendah dalam kesimpulan
sumber (hlm. 422-423).

\textbf{Interpretasi penulis:} Daya saing relatif transportasi publik
relevan bagi biaya privat karena koefisiennya berkaitan dengan durasi
komuter yang lebih pendek. Sumber juga menempatkan sub-centres dan
keseimbangan antara keuntungan aglomerasi serta pengurangan kemacetan
sebagai alasan perencanaan yang mendukung polisentrisitas (hlm. 405,
416-419, 422).

\subsection{Applications}\label{applications-53}

\textbf{Laporan sumber:} Penerapan empiris yang ditunjukkan adalah
penggunaan indikator berbasis arus komuter dan social network analysis
untuk mengidentifikasi sub-centres serta membandingkan derajat
polisentrisitas pada 82 Italian MAs. EIM dan ACT kemudian digunakan
untuk menguji dua komponen biaya mobilitas (hlm. 409-416, 420-423).

\textbf{Inferensi pengolahan:} Kerangka indikator tersebut dapat dipakai
sebagai perangkat analisis komparatif kawasan, tetapi TXT tidak
mendokumentasikan penerapan pada negara, kota, atau program kebijakan di
luar 82 Italian MAs.

\subsection{Future Research
(Optional)}\label{future-research-optional-53}

\textbf{Catatan pengolahan:} Agenda future research yang ditulis sebagai
bagian tersendiri tidak disebutkan dalam file.

\textbf{Laporan sumber:} Dua kebutuhan penelitian yang dinyatakan secara
langsung adalah penggunaan data longitudinal untuk membantu penetapan
kausalitas (hlm. 420) dan penelitian teoretis serta empiris lebih lanjut
tentang konsep polisentrisitas serta dampaknya pada daya saing dan
keberlanjutan (hlm. 423).

\textbf{Catatan pengolahan:} Pernyataan di atas dicatat sebagai
kebutuhan atau arah yang tersirat dalam teks sumber, bukan usulan future
research baru dari reviewer.

\subsection{Provenance Notes}\label{provenance-notes-53}

\begin{itemize}
\tightlist
\item
  \textbf{TXT aktual yang diverifikasi:}
  \texttt{source/journal/B10-veneri-2010-urban-polycentricity-and-the-costs-of-commuting-evidence-from-italian-metropolitan-areas-10.1111-j.1468-2257.2010.00531.x.txt}.
  Pencarian prefiks \texttt{B10} di bawah \texttt{source/} menemukan
  satu file TXT ini.
\item
  \textbf{Metadata ekstraksi yang dibaca:} blok
  \texttt{{[}PDF\ Metadata{]}} menyatakan PDF sumber berada pada path
  relatif
  \texttt{journal/B10-veneri-2010-urban-polycentricity-and-the-costs-of-commuting-evidence-from-italian-metropolitan-areas-10.1111-j.1468-2257.2010.00531.x.pdf},
  diekstrak pada 2026-08-31 dengan \texttt{pdftotext\ -layout}; metadata
  PDF mencatat judul artikel, 27 halaman, PDF tidak terenkripsi, dan
  ukuran halaman 432 x 648 pts. Creator tercatat
  \texttt{XyEnterprise\ XPP\ 8.0C.1\ Patch\ \#3}, producer
  \texttt{PDFlib\ PLOP\ 2.0.0p6\ (SunOS)/Acrobat\ Distiller\ 7.0.5\ (Windows)},
  PDF version 1.3, dan ukuran file 385221 bytes (TXT, baris 1-25).
\item
  \textbf{Batas akses:} Review ini disusun hanya dari seluruh isi TXT,
  termasuk metadata ekstraksi dan teks artikel. Daftar pustaka dan karya
  yang dikutip di dalam artikel tidak diverifikasi secara terpisah;
  tidak ada sumber luar yang ditambahkan.
\item
  \textbf{Locator:} Klaim substantif menggunakan nomor halaman tercetak
  artikel 403-429 serta nomor tabel, gambar, persamaan, dan catatan
  sebagaimana muncul dalam TXT. Nomor halaman PDF pada metadata
  diperlakukan sebagai informasi teknis, bukan pengganti locator
  artikel.
\item
  \textbf{Perbedaan internal sumber yang dipertahankan:} Abstrak
  menyatakan MA yang lebih polisentris lebih baik dalam biaya privat dan
  eksternal (hlm. 403), tetapi hasil rinci menunjukkan koefisien ACT
  negatif namun tidak signifikan (hlm. 420-423, Table 7). Selain itu,
  narasi menyatakan \texttt{house\_age} tidak signifikan, sedangkan
  Table 7 memberi \texttt{0,129*} dengan p = 0,090 pada spesifikasi ACT
  kedua; footnote tabel menetapkan \texttt{*\ p\ \textless{}\ .10} (hlm.
  421-422).
\item
  \textbf{Perbedaan internal sumber yang dipertahankan:} Abstrak
  menyatakan MA yang lebih polisentris lebih baik dalam biaya privat dan
  eksternal (hlm. 403), tetapi hasil rinci menunjukkan koefisien ACT
  negatif namun tidak signifikan (hlm. 420-423, Table 7). Selain itu,
  narasi menyatakan \texttt{house\_age} dan dummy geografis tidak
  signifikan (kecuali Southern MAs pada EIM), sedangkan Table 7 memberi
  \texttt{house\_age\ =\ 0,129*} dengan p = 0,090 dan
  \texttt{dummy\_centre\ =\ -0,306*} dengan p = 0,064 pada spesifikasi
  ACT kedua; footnote tabel menetapkan \texttt{*\ p\ \textless{}\ .10}
  (hlm. 421-422).
\item
  \textbf{Perbedaan internal sumber yang dipertahankan:} Abstrak
  menyatakan MA yang lebih polisentris lebih baik dalam biaya privat dan
  eksternal (hlm. 403), tetapi hasil rinci menunjukkan koefisien ACT
  negatif namun tidak signifikan (hlm. 420-423, Table 7). Selain itu,
  narasi menyatakan \texttt{house\_age} dan dummy geografis tidak
  signifikan (kecuali Southern MAs pada EIM), sedangkan Table 7 memberi
  \texttt{house\_age\ =\ 0,129*} dengan p = 0,090 dan
  \texttt{dummy\_centre\ =\ -0,306*} dengan p = 0,064 pada spesifikasi
  ACT kedua; footnote tabel menetapkan \texttt{*\ p\ \textless{}\ .10}
  (hlm. 421-422). Table 6 mendefinisikan \texttt{gini\_area\_job}
  sebagai konsentrasi pekerjaan, sedangkan narasi hasil menyebut
  konsentrasi populasi ketika menjelaskan koefisiennya (hlm. 418, 421).
\item
  \textbf{Penamaan variabel:} Table 6 mencetak \texttt{rel\_pub\_comp},
  sedangkan uraian teks dan Table 7 menggunakan
  \texttt{rel\_public\_comp}; review mempertahankan nama yang muncul
  pada bagian hasil dan tidak menganggapnya sebagai dua variabel berbeda
  (hlm. 418, 420-421).
\item
  \textbf{Batas inferensi:} Label \textbf{Laporan sumber} merangkum
  pernyataan, metode, hasil, atau interpretasi yang ditulis artikel;
  \textbf{Interpretasi penulis} merujuk pada penafsiran Veneri;
  \textbf{Inferensi pengolahan} adalah kesimpulan terbatas dari penataan
  ulang isi TXT dan tidak dimaksudkan sebagai klaim kausal atau
  generalisasi baru.
\end{itemize}

\section{Review B11}\label{review-b11}

\textbf{Kode:} B11\\
\textbf{Sumber yang direview:}
\texttt{source/journal/B11-camagni-capello-caragliu-2015-the-rise-of-second-rank-cities-what-role-for-agglomeration-economies-10.1080-09654313.2014.904999.txt}

\textbf{Penanda epistemik:} \texttt{{[}Laporan\ sumber{]}} menyatakan
temuan atau isi yang dilaporkan artikel;
\texttt{{[}Interpretasi\ penulis{]}} menyatakan penjelasan atau posisi
substantif penulis artikel; \texttt{{[}Inferensi\ pengolahan{]}}
menyatakan pembacaan yang dibuat dalam review ini, bukan klaim langsung
sumber.

\subsection{Summarized Abstract}\label{summarized-abstract-54}

{[}Laporan sumber{]} Artikel membahas mengapa kota peringkat kedua di
Eropa pada beberapa periode dapat menyamai atau melampaui kinerja kota
peringkat pertama. Sumber merangkum bukti bahwa dalam sekitar 15 tahun
terakhir kota peringkat pertama tidak selalu memimpin kinerja ekonomi
nasional; ketika memimpin, perbedaannya dengan kota peringkat kedua juga
tidak selalu signifikan. Sumber juga menyebut klaim terdahulu bahwa kota
peringkat kedua telah menjadi penggerak utama kinerja nasional (TXT
baris 109-115).

{[}Interpretasi penulis{]} Penulis menolak penjelasan yang
menyederhanakan kinerja kota menjadi eksploitasi ekonomi aglomerasi yang
hanya bergantung pada ukuran kota. Mereka mengusulkan bahwa fungsi
tingkat tinggi, efisiensi internal, dan jaringan kerja sama antarkota
merupakan kondisi yang memungkinkan ekonomi aglomerasi dimanfaatkan dan
imbal hasil urban yang menurun diatasi (TXT baris 116-126, 185-194).

{[}Inferensi pengolahan{]} Artikel menggabungkan argumen teoretis dengan
analisis deskriptif dan estimasi ekonometrik. Outcome ekonometrik
utamanya adalah \emph{average location benefits} yang diproksikan dengan
log sewa urban, bukan GDP secara langsung; GDP digunakan pada tahap
deskriptif untuk mengidentifikasi kota yang \emph{outperforming} (TXT
baris 418-481, 515-578).

\subsection{Problem Statement}\label{problem-statement-54}

{[}Laporan sumber{]} Perdebatan tentang kebangkitan kota peringkat kedua
sering menjelaskan kinerja relatif mereka dengan ekonomi aglomerasi yang
dianggap tersedia di kota kecil dan menengah, sementara disekonomi skala
dianggap terutama muncul di kota besar. Penulis menilai penjelasan ini
tidak cukup karena menggunakan batas struktural untuk menjelaskan gejala
siklis dan tidak menjelaskan mengapa perbedaan kinerja muncul pada waktu
dan tempat tertentu (TXT baris 231-254).

{[}Interpretasi penulis{]} Masalah analitis terdiri atas dua pertanyaan
yang berbeda: bagaimana setiap kelas ukuran kota dapat memanfaatkan
imbal hasil urban berskala meningkat, dan kapan suatu kelas kota lebih
cocok untuk ekspansi atau masa krisis. Pertanyaan pertama dipandang
lebih mendasar karena kota-kota dengan ukuran sama dapat memiliki
kinerja berbeda (TXT baris 243-260).

{[}Laporan sumber{]} Artikel juga menyatakan bahwa literatur yang ada
belum cukup menguraikan isi ``black box'' ekonomi aglomerasi dan terlalu
menempatkan ukuran sebagai determinan tunggal dinamika urban. Karena
itu, kemampuan menarik atau mengembangkan fungsi tingkat tinggi serta
membangun jaringan kerja sama perlu diuji secara empiris (TXT baris
185-229).

\subsection{Objectives}\label{objectives-54}

{[}Laporan sumber{]} Tujuan utama artikel adalah memeriksa kondisi yang
membuat kota peringkat kedua di Uni Eropa pada beberapa periode
mengungguli kawasan metropolitan yang lebih besar. Peringkat ditentukan
menurut ukuran penduduk fisik wilayah \emph{larger urban zone} (LUZ),
dengan kota peringkat pertama berpenduduk lebih dari 1 juta dan kota
peringkat kedua berpenduduk 200.000 sampai kurang dari 1 juta
berdasarkan data 2011 (TXT baris 195-203).

{[}Laporan sumber{]} Penulis merumuskan empat hipotesis: kota kecil dan
menengah dapat berhenti tumbuh atau menurun jika tidak naik ke peringkat
urban yang lebih tinggi; kota kecil dan menengah terbaik dapat menunda
atau mengatasi imbal hasil yang menurun; kemampuan itu bergantung pada
fungsi bernilai lebih tinggi dan jaringan kerja sama eksternal; dan kota
peringkat kedua dapat meminjam faktor tersebut dari kawasan urban besar
yang berdekatan (TXT baris 387-401).

{[}Inferensi pengolahan{]} Tujuan empirisnya mencakup tiga pengujian
yang saling terkait: hubungan ukuran dengan manfaat lokasi, peran fungsi
tingkat tinggi dan jaringan kota, serta apakah \emph{borrowed size} dari
kota peringkat pertama membantu kota peringkat kedua yang berkinerja di
atas rata-rata. Pembagian ini diringkas dari persamaan (1)-(3) dan
susunan Section 5, bukan daftar tujuan formal terpisah dalam TXT (TXT
baris 515-581, 665-907).

\subsection{Literature Survey}\label{literature-survey-54}

\begin{itemize}
\tightlist
\item
  {[}Laporan sumber{]} Literatur yang diringkas mengaitkan ekonomi
  aglomerasi dengan ukuran kota, tetapi penulis menempatkan perdebatan
  Eropa tentang kinerja kota peringkat pertama dan kedua sebagai bukti
  bahwa ukuran saja tidak cukup. Penjelasan berbasis disekonomi skala
  dipandang berisiko menjadi penalaran melingkar jika tidak menjelaskan
  ambang strukturalnya (TXT baris 231-260).
\item
  {[}Laporan sumber{]} Model SOUDY (\emph{Supply Oriented Urban
  DYnamics}) menyatakan bahwa setiap fungsi ekonomi dan peringkat urban
  memiliki rentang ukuran efisien, dengan batas minimum dan maksimum
  ketika manfaat produksi melebihi biaya lokasi. Kota yang mendekati
  ukuran maksimum untuk peringkatnya memasuki area tidak stabil dan
  dapat naik melalui pengembangan atau penarikan fungsi tingkat lebih
  tinggi (TXT baris 261-294).
\item
  {[}Interpretasi penulis{]} Dari SOUDY, ukuran optimal tunggal
  sebaiknya diganti dengan rentang ukuran efisien; rentang itu berbeda
  menurut fungsi yang dijalankan, dan peringkat urban dapat dipisahkan
  dari ukuran fisik. Dua kota dengan ukuran sama dapat berada pada
  peringkat berbeda jika salah satunya memiliki fungsi tingkat lebih
  tinggi (TXT baris 295-310).
\item
  {[}Laporan sumber{]} Teori jaringan kota menjelaskan kerja sama
  sebagai bentuk organisasi antara pertumbuhan internal dan eksternal.
  Sumber membahas spesialisasi, fungsi tingkat tinggi pada pusat
  berperingkat lebih rendah, hubungan horizontal antarkota, jaringan
  komplementer seperti Randstad dan Veneto, integrasi
  vertikal-horizontal, serta eksternalitas jaringan (TXT baris 312-335).
\item
  {[}Laporan sumber{]} Konsep \emph{borrowed size} dari Alonso
  menjelaskan bahwa kota kecil atau kawasan metropolitan dapat
  memperlihatkan sebagian karakter kota lebih besar karena dekat dengan
  konsentrasi penduduk lain. Sumber menyebut lokasi dalam sistem urban
  polisentris dapat menggantikan sebagian ukuran fisik kota, tetapi juga
  mengingatkan bahwa layanan tertentu dapat kehilangan skala jika
  ambangnya tidak dapat disebarkan (TXT baris 336-356).
\item
  {[}Interpretasi penulis{]} \emph{City-network theory} memperluas
  \emph{borrowed size}: ukuran dapat diperoleh bukan hanya dari
  kedekatan fisik dengan pusat besar, tetapi juga dari hubungan
  horizontal nonhierarkis dengan kota komplementer atau serupa yang
  lebih jauh. Dengan demikian, konektivitas, fungsi tingkat tinggi, dan
  integrasi jaringan diposisikan sebagai sumber keunggulan selain ukuran
  (TXT baris 357-386).
\end{itemize}

\subsection{Method Used}\label{method-used-54}

{[}Laporan sumber{]} Penelitian menggabungkan analisis deskriptif tren
GDP dengan estimasi fungsi produksi urban agregat. Model menjelaskan
\emph{average location benefits} (ALB) melalui populasi metropolitan,
kuadrat populasi, interaksi populasi dengan fungsi tingkat tinggi dan
jaringan, serta indikator kota peringkat kedua yang mengungguli
rata-rata nasional (TXT baris 404-481, 484-578).

{[}Laporan sumber{]} ALB diproksikan dengan sewa urban yang diukur
melalui harga rumah per meter persegi untuk apartemen kualitas rata-rata
di kawasan \emph{Central Business District} (CBD). Persamaan dasar
memasukkan Pop dan Pop2 untuk menangkap imbal hasil meningkat atau
menurun, serta interaksi antara populasi dan fungsi atau jaringan
setelah tercapai massa kritis. Persamaan lanjutan menambahkan interaksi
dengan dummy kota yang \emph{outperforming} dan persamaan alternatif
yang mengganti fungsi serta jaringan dengan \emph{borrowed size} (TXT
baris 515-578).

{[}Laporan sumber{]} Estimasi pada Table 3 dan Table 5 memakai
\emph{random effects}, galat baku robust, dan 158 observasi. Penulis
menolak \emph{fixed effects} kota secara konseptual, sedangkan uji
Breusch-Pagan dilaporkan mendukung \emph{random effects} dibanding OLS.
Table 3 memuat empat spesifikasi; Table 5 membandingkan model dasar,
\emph{borrowed size}, dan interaksi \emph{borrowed size} dengan dummy
kota yang \emph{outperforming} (TXT baris 665-717, 832-870).

{[}Laporan sumber{]} Kinerja relatif pada tahap deskriptif dihitung dari
pertumbuhan GDP dan GDP per kapita pada periode 1996-2009. Fokus empiris
kemudian diletakkan pada periode 2001-2004 dan 2008-2009, ketika kota
peringkat kedua tumbuh lebih cepat; dummy bernilai 1 jika pertumbuhan
wilayah metropolitan melebihi rata-rata pertumbuhan GDP negara (TXT
baris 418-481, 503-511).

{[}Laporan sumber{]} \emph{Borrowed size} dihitung sebagai populasi kota
peringkat pertama yang didiskontokan oleh jarak dari setiap kota
peringkat kedua. Matriks bobot bernilai nol untuk pasangan dengan
peringkat sama dan menggunakan kebalikan jarak geografis dalam kilometer
untuk pasangan kota peringkat pertama-kedua (TXT baris 648-664).

{[}Inferensi pengolahan{]} Karena appendix menyebut dua observasi per
kota dan tabel menggunakan \emph{random effects}, desainnya dapat dibaca
sebagai panel pendek pada tingkat wilayah metropolitan. Istilah panel:
Tidak disebutkan dalam file. Strategi identifikasi kausal: Tidak
disebutkan dalam file. Hasil paling aman dibaca sebagai hubungan
statistik pada struktur data yang digunakan (TXT baris 729-734,
1102-1123).

\subsection{Dataset}\label{dataset-54}

{[}Laporan sumber{]} Dataset baru disusun dari data wilayah metropolitan
EUROSTAT. Wilayah metropolitan dibentuk dari agregasi wilayah
administratif NUTS3 dengan sedikitnya 250.000 penduduk dan berbasis LUZ,
dengan tujuan memasukkan sabuk komuter di sekitar kota (TXT baris
580-600).

{[}Laporan sumber{]} GDP dalam harga berjalan dan populasi dikumpulkan
untuk 1995-2009. GDP dideflasi menggunakan deflator negara EUROSTAT
menjadi harga konstan 2005; sumber menyatakan variasi harga di dalam
negara belum sepenuhnya dapat dihilangkan. Untuk tahun awal ketika data
populasi tidak tersedia, populasi wilayah NUTS3 pembentuk kawasan
metropolitan diagregasikan mengikuti klasifikasi resmi EUROSTAT (TXT
baris 593-600).

{[}Laporan sumber{]} Data sewa lahan berupa harga per meter persegi
apartemen kualitas rata-rata di CBD dikumpulkan untuk 2004 dan 2011,
lalu dideflasi dengan deflator negara yang sama. Fungsi urban tingkat
tinggi diukur dari proporsi profesional berketerampilan tinggi terhadap
total penduduk, menggunakan mikrodata \emph{Labour Force Survey}
EUROSTAT yang diagregasikan pada NUTS2 dan ditetapkan untuk kawasan
metropolitan. Data tersedia 1995-2007; rata-rata 1998-2002 dan 2002-2006
dihitung untuk analisis (TXT baris 601-624, 636-638).

{[}Laporan sumber{]} Jaringan kota diukur dari jumlah proyek Framework
Programme 5 dan 6 yang diikuti institusi di setiap kota. CORDIS menjadi
sumber data, dengan periode FP5 1998-2002 dan FP6 2002-2006. Sumber
memposisikannya sebagai proksi kolaborasi ilmiah, pertukaran
pengetahuan, dan jaringan eksternal jarak jauh (TXT baris 625-647).

{[}Laporan sumber{]} Table 2 merangkum lima kelompok data: GDP konstan
2005 dari EUROSTAT (1995-2009), populasi tahunan rata-rata dari EUROSTAT
(1995-2009), sewa lahan dari EUROSTAT/Urban Audit dan integrasi penulis
(2004 dan 2011), fungsi tingkat tinggi dari \emph{Labour Force Survey}
(1995-2007), serta jaringan kota dari CORDIS untuk FP5 dan FP6 (TXT
baris 612-638).

\subsection{Population / Sample}\label{population-sample-54}

{[}Laporan sumber{]} Analisis deskriptif menggunakan sampel 136 kawasan
metropolitan Uni Eropa. Kota peringkat pertama adalah LUZ dengan lebih
dari 1 juta penduduk, sedangkan kota peringkat kedua memiliki 200.000
sampai kurang dari 1 juta penduduk; kelas ukuran tersebut ditetapkan
menggunakan data 2011 (TXT baris 404-417).

{[}Laporan sumber{]} Database untuk analisis empiris disebut dalam
appendix terdiri atas 100 kawasan metropolitan, dengan dua observasi per
kota dan cakupan waktu sebagaimana Section 4. Table 3 dan Table 5
masing-masing melaporkan 158 observasi. Estimasi utama pada Table 3
dijelaskan untuk kota peringkat kedua di negara-negara Barat; catatan 7
menyatakan bahwa analisis Section 5 hanya menggunakan kota peringkat
kedua di negara EU15 karena dinamika kota peringkat kedua di Eropa Timur
dipandang khas (TXT baris 665-734, 971-973, 1102-1123).

{[}Inferensi pengolahan{]} Angka 136 kawasan metropolitan, 100 kawasan
dalam database appendix, dan 158 observasi merujuk pada cakupan tahap
analisis yang berbeda, sehingga tidak disamakan sebagai satu ukuran
sampel. TXT tidak merinci jumlah kota peringkat kedua yang tersisa
setelah pembatasan EU15, aturan sampling probabilistik, atau daftar
observasi yang dikeluarkan selain keterangan tentang Eropa Timur.
Informasi tersebut: Tidak disebutkan dalam file (TXT baris 404-417,
665-734, 1102-1123).

\subsection{Dependent Variables}\label{dependent-variables-54}

{[}Laporan sumber{]} Variabel dependen pada persamaan ekonometrik adalah
ALB, yaitu \emph{average location benefits}, yang dioperasionalkan
sebagai log sewa urban atau log harga apartemen kualitas rata-rata per
meter persegi di CBD. Sumber mendasarkan indikator ini pada gagasan
bahwa harga rumah mencerminkan evaluasi pasar atas daya tarik relatif
dan keunggulan lokasi bersih suatu kota (TXT baris 535-551, 674-716).

{[}Laporan sumber{]} Pada analisis deskriptif, pertumbuhan GDP tahunan
dan pertumbuhan GDP per kapita digunakan untuk membandingkan kinerja
kota peringkat pertama dan kedua selama 1996-2009. Dummy D yang bernilai
1 ketika pertumbuhan wilayah metropolitan melebihi rata-rata negara
kemudian dipakai sebagai variabel penjelas dalam persamaan interaksi,
bukan sebagai variabel dependen utama (TXT baris 418-481, 552-578).

{[}Inferensi pengolahan{]} ALB dan GDP growth menangkap keluaran yang
berbeda: ALB merepresentasikan manfaat lokasi yang dinilai melalui pasar
perumahan, sedangkan GDP growth merepresentasikan dinamika kinerja
ekonomi. Ukuran dependen lain untuk estimasi utama: Tidak disebutkan
dalam file (TXT baris 418-481, 535-578, 674-716).

\subsection{Results}\label{results-54}

\begin{itemize}
\tightlist
\item
  {[}Laporan sumber{]} Analisis deskriptif menemukan pola siklis pada
  ekonomi Uni Eropa, dengan fase pertumbuhan cepat sekitar 1995-2000 dan
  2003-2007, perlambatan 2000-2003, lalu perlambatan dan kontraksi
  terkait krisis setelah 2007. Kota peringkat pertama tampak mendorong
  pertumbuhan pada fase ekspansi, sedangkan kota peringkat kedua tumbuh
  lebih cepat pada periode 2001-2004 dan mengalami penurunan GDP yang
  lebih kecil pada 2008-2009 (TXT baris 418-481; Figure 2; Table 1).
\item
  {[}Laporan sumber{]} Dalam pertumbuhan GDP per kapita, kota peringkat
  kedua mengungguli kota peringkat pertama pada 8 dari 10 tahun terakhir
  yang diperiksa. Pola itu terlihat pada EU27 dan EU15, terutama pada
  periode kedua dan keempat. Di New Member States, kota peringkat
  pertama lebih sering berkinerja lebih baik dan tidak tampak pola
  pemulihan yang jelas bagi kota peringkat kedua (TXT baris 455-469;
  Figure 3).
\item
  {[}Interpretasi penulis{]} Ketidakmunculan pola pemulihan kota
  peringkat kedua di New Member States mungkin berkaitan dengan dominasi
  kota peringkat pertama dalam penciptaan kekayaan atau belum
  terbentuknya jaringan kota peringkat kedua yang terstruktur dan saling
  terhubung. Sumber menyajikan ini sebagai kemungkinan penjelasan, bukan
  pengujian kausal terpisah (TXT baris 460-469).
\item
  {[}Laporan sumber{]} Perbedaan rata-rata antar-kelompok tidak
  memperoleh dukungan statistik yang tegas dari \emph{t-tests} klasik
  pada keempat periode. Hasil ini menjadi alasan bagi penulis untuk
  melanjutkan pada analisis empiris mengenai kondisi yang membuat
  sebagian kota peringkat kedua mengungguli kota peringkat pertama (TXT
  baris 472-481).
\item
  {[}Laporan sumber{]} Pada Table 3, bentuk inverted U untuk hubungan
  ukuran dan ALB tidak berlaku bagi seluruh sampel kota peringkat kedua
  di negara Barat dalam spesifikasi awal, bahkan setelah fungsi dan
  jaringan ditambahkan. Fungsi signifikan dalam nilai absolut pada
  spesifikasi terkait, sedangkan jaringan tidak signifikan pada model
  yang disebutkan (TXT baris 726-747; Table 3).
\item
  {[}Laporan sumber{]} Ketika interaksi dengan dummy kota yang
  \emph{outperforming} dimasukkan, kota yang tidak berkinerja lebih baik
  menunjukkan pola inverted U tradisional. Kota yang tumbuh cepat
  menunjukkan imbal hasil menurun sampai ukuran kritis tertentu, lalu
  beralih ke imbal hasil meningkat. Pada spesifikasi konseptual (2),
  koefisien interaksi umumnya memiliki tanda yang diharapkan dan
  signifikan, kecuali jaringan pada kota \emph{outperforming} (TXT baris
  737-747; Table 3).
\item
  {[}Laporan sumber{]} Figure 4 menggambarkan kurva inverted U untuk
  sampel umum dan kurva U untuk kota \emph{outperforming}. Sumber
  melaporkan ukuran optimal kota peringkat kedua sekitar 714.973
  penduduk, ukuran kritis kota \emph{outperforming} sekitar 296.559
  penduduk, serta kombinasi 7,5\% pekerjaan pada fungsi tingkat lanjut
  dan ukuran 545.000 penduduk untuk mencapai tingkat manfaat urban yang
  sama dengan kota lain (TXT baris 748-807; Figure 4; Table 4).
\item
  {[}Laporan sumber{]} Figure 5 menunjukkan bahwa peran fungsi tingkat
  tinggi dalam manfaat marginal lokasi besar bagi kota peringkat kedua
  yang \emph{outperforming}, tetapi kemiringan pengaruhnya tetap konstan
  ketika intensitas fungsi meningkat (TXT baris 819-874; Figure 5).
\item
  {[}Laporan sumber{]} Pada Table 5, \emph{borrowed size} tidak cukup
  untuk menghasilkan manfaat yang lebih tinggi pada model yang
  memasukkannya tanpa interaksi. Pada model yang memasukkan interaksi
  dengan dummy kota peringkat kedua yang \emph{outperforming}, interaksi
  tersebut bernilai positif dan ditandai signifikan; koefisien
  \emph{borrowed size} langsung tetap tidak ditandai signifikan. Table 5
  menggunakan 158 observasi, estimasi \emph{random effects}, dan galat
  baku robust (TXT baris 878-896; Table 5, baris 832-870).
\item
  {[}Laporan sumber{]} Figure 6 memperlihatkan, menurut perhitungan dan
  interval kepercayaan berbasis kuintil distribusi, bahwa ALB meningkat
  ketika akses kota peringkat kedua terhadap fungsi dan jaringan di
  kawasan metropolitan besar meningkat, meskipun fungsi dan jaringan itu
  tidak berada di dalam kota tersebut (TXT baris 897-914; Figure 6).
\end{itemize}

\subsection{Findings}\label{findings-54}

{[}Interpretasi penulis{]} Temuan inti artikel adalah bahwa kota
peringkat kedua yang berkinerja baik tidak sekadar memperoleh keuntungan
dari ukuran. Mereka mampu mengubah imbal hasil menurun menjadi ekonomi
aglomerasi melalui fungsi tingkat tinggi, jaringan kerja sama, atau
akses terhadap fungsi dan jaringan di kota besar (TXT baris 814-817,
918-941).

{[}Interpretasi penulis{]} Penulis menganggap hipotesis kedua
terkonfirmasi: kota yang \emph{outperforming} dapat mengatasi disekonomi
skala. Fungsi tingkat tinggi meningkatkan manfaat lokasi, dan
\emph{borrowed size} berperan sebagai katalis bagi kota peringkat kedua
yang sudah mengungguli rata-rata negara, bukan sebagai pengganti yang
otomatis menguntungkan semua kota (TXT baris 810-817, 885-896).

{[}Inferensi pengolahan{]} Bukti untuk peran jaringan perlu dibaca lebih
hati-hati daripada kesimpulan umum artikel: jaringan tidak signifikan
pada bagian fungsi dan tidak menghasilkan efek langsung signifikan pada
model \emph{borrowed size}, sedangkan interaksi \emph{borrowed size}
dengan status \emph{outperforming} dilaporkan signifikan. Jadi, TXT
lebih kuat mendukung hubungan kondisional daripada manfaat seragam bagi
semua kota (TXT baris 695-704, 848-851, 878-896).

{[}Inferensi pengolahan{]} Hasil deskriptif dan ekonometrik tidak
identik: keunggulan pertumbuhan kota peringkat kedua pada periode
tertentu tidak dengan sendirinya membuktikan bahwa fungsi atau
\emph{borrowed size} menyebabkan ALB lebih tinggi. Model memakai agregat
wilayah metropolitan dan dua observasi per kota, sementara strategi
identifikasi kausal tidak disebutkan dalam TXT (TXT baris 418-481,
706-713, 1102-1123).

\subsection{Conclusions}\label{conclusions-54}

{[}Laporan sumber{]} Penulis menyimpulkan bahwa anggapan umum tentang
determinan pertumbuhan urban terlalu sederhana. Ekonomi aglomerasi
memang ada dan disekonomi aglomerasi dapat muncul, tetapi kota dapat
mengatasi disekonomi skala dengan berinovasi pada fungsi yang dijalankan
atau pada pengorganisasian kegiatan bersama kota lain melalui jaringan.
Sumber juga menyatakan bahwa kota dapat ``borrow size'' dari kawasan
metropolitan besar di sekitarnya untuk memperoleh akses ke fungsi dan
jaringan tanpa menanggung seluruh kerugian lokasi kota besar (TXT baris
918-941).

{[}Laporan sumber{]} Studi empiris menyatakan bahwa kota yang
\emph{outperforming} dicirikan oleh ekonomi skala dan bahwa ekonomi
tersebut berkaitan dengan tingkat fungsi dan jaringan yang dimiliki
kota, atau dengan kapasitasnya untuk meminjam ukuran dari kawasan
metropolitan besar (TXT baris 938-941).

{[}Interpretasi penulis{]} Dalam pilihan kebijakan antara memusatkan
sumber daya pada kota besar atau menyebarkan investasi ke lebih banyak
kota, penulis mendukung investasi pada kota-kota yang memungkinkan kota
dengan ukuran apa pun mengubah risiko imbal hasil menurun menjadi
ekonomi aglomerasi. Instrumen yang disebut adalah pembaruan fungsi dan
cara kerja sama antarkota (TXT baris 942-949).

\subsection{Contributions}\label{contributions-54}

{[}Laporan sumber{]} Artikel secara eksplisit menyatakan bahwa
kontribusinya terhadap literatur adalah pandangan yang lebih kompleks:
kinerja urban tidak hanya bergantung pada ekonomi aglomerasi, dan
ekonomi aglomerasi tidak hanya bergantung pada ukuran kota. Artikel
menghubungkan gagasan ini dengan model SOUDY dan teori jaringan kota
(TXT baris 185-194).

{[}Inferensi pengolahan{]} Kontribusi substantif yang dapat
diidentifikasi dari TXT adalah:

\begin{itemize}
\tightlist
\item
  kontribusi konseptual, yaitu menguraikan determinan ekonomi aglomerasi
  ke dalam ukuran, fungsi tingkat tinggi, jaringan, dan \emph{borrowed
  size};
\item
  kontribusi empiris, yaitu membandingkan tren kinerja kota peringkat
  pertama dan kedua serta menguji kondisi ALB pada database wilayah
  metropolitan Eropa;
\item
  kontribusi operasional, yaitu menerjemahkan fungsi tingkat tinggi,
  partisipasi jaringan FP, dan akses berbobot jarak ke dalam variabel
  model; dan
\item
  kontribusi kebijakan, yaitu menawarkan pembaruan fungsi dan kerja sama
  sebagai alternatif terhadap konsentrasi investasi semata pada kota
  besar.
\end{itemize}

Pengelompokan empat jenis kontribusi tersebut adalah sintesis review;
TXT tidak menyediakan subbagian formal bernama ``Contributions''. Dasar
faktualnya adalah uraian tujuan, dataset baru, persamaan, hasil, dan
implikasi kebijakan (TXT baris 219-229, 484-664, 918-949).

\subsection{Research Gap}\label{research-gap-54}

{[}Laporan sumber{]} Gap yang hendak diisi adalah penjelasan yang belum
memadai tentang mengapa sebagian kota peringkat kedua berkinerja lebih
baik, selain jawaban umum berupa ukuran dan disekonomi kota besar.
Penulis menyatakan bahwa teori yang tersedia tentang rentang ukuran
efisien dan jaringan kota perlu dibuktikan secara empiris, serta bahwa
ekonomi aglomerasi perlu diuraikan menurut fungsi dan hubungan antarkota
(TXT baris 231-260, 261-335).

{[}Interpretasi penulis{]} Artikel mengisi gap tersebut dengan menguji
fungsi tingkat tinggi, jaringan kerja sama, dan \emph{borrowed size}
sebagai kondisi yang dapat membantu kota kecil dan menengah mengatasi
imbal hasil menurun. Analisis juga membedakan kedekatan dengan kota
besar dari hubungan jaringan horizontal, sehingga \emph{borrowed size}
tidak dipahami hanya sebagai efek jarak fisik (TXT baris 336-386,
515-578, 878-914).

{[}Laporan sumber{]} Gap lanjutan yang disebut secara eksplisit adalah
perlunya menjelaskan kondisi yang membuat struktur urban polisentris
muncul dan berhasil, serta karakteristik sistem urban lain yang
menentukan kemampuan kota peringkat kedua memperoleh manfaat dari
kedekatan dengan kawasan metropolitan besar (TXT baris 908-914).

{[}Inferensi pengolahan{]} Keterbatasan data dan desain juga menyisakan
kebutuhan untuk menguji apakah hubungan yang dilaporkan tetap berlaku
pada cakupan kota dan periode yang lebih luas. Ini merupakan inferensi
dari pembatasan sampel, indikator, dan dua observasi per kota, bukan
agenda eksplisit tambahan di luar yang disebutkan sumber (TXT baris
665-734, 989-1001, 1102-1123).

\subsection{Limitations}\label{limitations-54}

\begin{itemize}
\tightlist
\item
  {[}Laporan sumber{]} Analisis ekonometrik tidak mencakup seluruh kota
  peringkat kedua Eropa. Kota peringkat kedua di Eropa Timur dikeluarkan
  dari analisis Section 5 karena sumber menilai dinamikanya khas; model
  berfokus pada kota-kota peringkat kedua di negara EU15/negara Barat
  (TXT baris 726-734, 971-973).
\item
  {[}Laporan sumber{]} ALB memakai harga rumah sebagai proksi sewa atau
  manfaat lokasi dan bergantung pada asumsi bahwa perbedaan harga
  mencerminkan daya tarik relatif. Untuk perbandingan dinamis, sumber
  juga mengasumsikan kemiringan kurva pasokan rumah sama pada dimensi
  relatif yang dibandingkan; sumber memperingatkan bahwa asumsi ini
  dapat menimbulkan bias berat terutama pada perbandingan area inti dan
  cincin (TXT baris 535-551, 974-988).
\item
  {[}Laporan sumber{]} Ukuran fungsi tingkat tinggi dibatasi pada
  proporsi kategori ISCO 88 kategori 1, yaitu legislator, pejabat
  senior, dan manajer. Data mikro diagregasikan di tingkat NUTS2 dan
  nilainya ditetapkan ke kawasan metropolitan (TXT baris 607-624,
  636-638).
\item
  {[}Laporan sumber{]} Ukuran jaringan kota hanya menghitung partisipasi
  dalam proyek FP5 dan FP6 dan terutama merepresentasikan konektivitas
  ilmiah atau pertukaran pengetahuan berorientasi inovasi. Sumber
  menyebutnya hanya sebagian dari kemungkinan jaringan transnasional dan
  bukan ukuran seluruh jaringan atau aksesibilitas fisik (TXT baris
  640-647, 989-996).
\item
  {[}Laporan sumber{]} Ketika \emph{borrowed size} dimasukkan, sumber
  mencatat multikolinearitas antara \emph{borrowed size} dan variabel
  fungsi serta jaringan (TXT baris 564-570, 989-990).
\item
  {[}Laporan sumber{]} Perbandingan deskriptif antarkelompok tidak
  mendapat dukungan tegas dari \emph{t-tests} klasik pada seluruh
  periode, walaupun pola grafik digunakan untuk memilih periode fokus
  (TXT baris 472-481).
\item
  {[}Laporan sumber{]} Ketersediaan data tidak seragam: GDP dan populasi
  mencakup 1995-2009, sewa urban hanya 2004 dan 2011, fungsi tingkat
  tinggi 1995-2007, serta jaringan FP5/FP6 1998-2006. Rata-rata periode
  untuk sejumlah variabel dihitung dari tahun yang tersedia (TXT baris
  503-511, 612-638).
\item
  {[}Inferensi pengolahan{]} Karena hasil utama diestimasi pada data
  agregat wilayah metropolitan dengan dua observasi per kota dan
  strategi identifikasi kausal: Tidak disebutkan dalam file, arah
  sebab-akibat dan pengaruh kebijakan tidak dapat dipastikan dari
  artikel ini saja (TXT baris 706-713, 729-734, 1102-1123).
\item
  {[}Informasi tidak tersedia{]} Prosedur umum untuk menangani data
  hilang, rincian penyesuaian bobot selain formula jarak, dan uji
  sensitivitas terhadap definisi kelas kota: Tidak disebutkan dalam
  file.
\end{itemize}

\subsection{Challenges}\label{challenges-54}

\begin{itemize}
\tightlist
\item
  {[}Laporan sumber{]} Tantangan konseptualnya adalah memisahkan
  persoalan struktural tentang cara kota memperoleh imbal hasil skala
  dari persoalan siklis tentang kapan kota tertentu tumbuh atau
  mengalami krisis (TXT baris 243-260).
\item
  {[}Laporan sumber{]} Tantangan teoritis dan pengukuran adalah
  menjelaskan perbedaan kinerja kota berukuran sama tanpa menganggap
  ukuran sebagai satu-satunya determinan. Penulis meresponsnya dengan
  menggabungkan rentang ukuran efisien, fungsi tingkat tinggi, jaringan
  kota, dan \emph{borrowed size} (TXT baris 261-335, 357-386).
\item
  {[}Laporan sumber{]} Penyusunan dataset menghadapi ketersediaan data
  metropolitan yang baru dan tidak seragam, termasuk kebutuhan
  mengagregasikan NUTS3 untuk tahun awal serta mengintegrasikan data
  harga rumah dari sumber nasional (TXT baris 584-606, 612-638).
\item
  {[}Laporan sumber{]} Operasionalisasi jaringan menghadapi pilihan
  indikator yang dapat dibandingkan lintas kota. Penulis memilih jumlah
  proyek FP5 dan FP6 karena dapat menangkap kolaborasi dan pertukaran
  pengetahuan tanpa bergantung pada jarak geografis, tetapi membatasi
  cakupan jaringan pada konektivitas ilmiah (TXT baris 640-647,
  989-996).
\item
  {[}Laporan sumber{]} Estimasi harus membedakan variasi antarkota dari
  pengaruh ukuran, fungsi, dan jaringan. Penulis menolak \emph{fixed
  effects} kota secara konseptual dan menggunakan hasil uji
  Breusch-Pagan untuk memilih \emph{random effects}; multikolinearitas
  muncul ketika \emph{borrowed size} dimasukkan bersama fungsi dan
  jaringan (TXT baris 726-734, 832-870, 989-990).
\item
  {[}Inferensi pengolahan{]} Perbedaan cakupan deskriptif EU27, analisis
  fokus EU15, dan database appendix menuntut kehati-hatian ketika hasil
  digeneralisasi ke seluruh Eropa. TXT melaporkan cakupan tersebut,
  tetapi tidak menyediakan rekonsiliasi lengkap jumlah unit pada setiap
  tahap (TXT baris 404-417, 460-481, 665-734, 1102-1123).
\end{itemize}

\subsection{Practical Implications}\label{practical-implications-54}

{[}Interpretasi penulis{]} Implikasi praktis utama adalah bahwa
kebijakan pertumbuhan tidak perlu memusatkan investasi hanya pada kota
terbesar. Investasi dapat diarahkan untuk memperbarui fungsi kota dan
cara kota bekerja sama agar kota, terlepas dari ukurannya, memiliki
peluang mengubah disekonomi skala menjadi ekonomi aglomerasi (TXT baris
942-949).

{[}Interpretasi penulis{]} Bagi kota peringkat kedua, pengembangan
fungsi bernilai tinggi dan hubungan dengan kawasan metropolitan besar
dipandang sebagai cara memperoleh manfaat aglomerasi tanpa harus
menanggung seluruh kerugian lokasi kota besar. Untuk jaringan, sumber
menekankan kerja sama dan organisasi kegiatan antarkota, bukan hanya
kedekatan fisik (TXT baris 331-386, 891-896).

{[}Inferensi pengolahan{]} Rekomendasi tersebut perlu dibaca sebagai
implikasi kebijakan dari hubungan statistik yang dilaporkan, bukan
sebagai bukti bahwa investasi tertentu pasti menghasilkan pertumbuhan.
Kesesuaian indikator harga rumah, fungsi, dan jaringan dengan konteks
lain harus diperiksa terlebih dahulu (TXT baris 535-551, 640-647,
706-713).

\subsection{Applications}\label{applications-54}

{[}Laporan sumber{]} Artikel menerapkan kerangka tersebut untuk
membandingkan pertumbuhan GDP kota peringkat pertama dan kedua,
mengidentifikasi kota yang \emph{outperforming}, mengestimasi kurva ALB,
dan menguji peran fungsi, jaringan, serta \emph{borrowed size} pada
wilayah metropolitan Eropa (TXT baris 404-481, 484-664, 665-914).

{[}Interpretasi penulis{]} Hasilnya dapat digunakan untuk mendukung
keputusan apakah sumber daya publik perlu disebarkan ke lebih banyak
kota dan untuk merancang investasi pada fungsi urban serta kerja sama
antarkota (TXT baris 942-949).

{[}Inferensi pengolahan{]} Di luar analisis Eropa yang dilaporkan,
aplikasi empiris pada negara, wilayah, atau kebijakan tertentu lainnya:
Tidak disebutkan dalam file. Penggunaan indikator artikel di luar
cakupan tersebut tidak diasumsikan dalam review ini.

\subsection{Future Research
(Optional)}\label{future-research-optional-54}

{[}Laporan sumber{]} Sumber secara eksplisit membuka penelitian lanjutan
untuk menjelaskan kondisi yang memungkinkan struktur urban polisentris
muncul dan berhasil, termasuk karakteristik sistem urban yang
memengaruhi kemampuan kota peringkat kedua memanfaatkan kedekatan dengan
kawasan metropolitan besar (TXT baris 908-914).

{[}Laporan sumber{]} Catatan 15 menyebut bahwa peran jaringan akan
disoroti ketika analisis diakronis dikembangkan. TXT tidak menyajikan
analisis diakronis tersebut sebagai bagian tersendiri dari artikel ini
(TXT baris 998-1001).

{[}Inferensi pengolahan{]} Agenda future research lain yang tidak
dinyatakan dalam TXT: Tidak disebutkan dalam file. Review ini tidak
menambahkan dugaan agenda riset.

\subsection{Provenance Notes}\label{provenance-notes-54}

\begin{itemize}
\tightlist
\item
  Review ini hanya menggunakan TXT aktual dengan prefiks B11:
  \texttt{source/journal/B11-camagni-capello-caragliu-2015-the-rise-of-second-rank-cities-what-role-for-agglomeration-economies-10.1080-09654313.2014.904999.txt}.
  Seluruh teks yang tersedia, dari baris 1 sampai 1124, dibaca, termasuk
  metadata ekstraksi, teks artikel, catatan, referensi, dan appendix.
\item
  Blok metadata menyatakan sumber PDF berjudul \emph{The Rise of
  Second-Rank Cities: What Role for Agglomeration Economies?}, memiliki
  23 halaman, diekstrak pada 2026-08-31 dengan
  \texttt{pdftotext\ -layout}, dan memiliki DOI
  \texttt{10.1080/09654313.2014.904999} (TXT baris 1-28). Review ini
  tidak menggunakan PDF terpisah atau sumber luar.
\item
  TXT memuat identitas artikel pada header, tiga penulis, \emph{European
  Planning Studies}, publikasi online 04 Jun 2014, DOI, afiliasi
  Politecnico di Milano, serta riwayat diterima Juni 2013 dan disetujui
  Januari 2014 (TXT baris 38-61, 95-105). Volume, nomor, dan rentang
  halaman jurnal: Tidak disebutkan dalam file. Metadata \texttt{pdfinfo}
  hanya mencantumkan Roberto Camagni sebagai Author, sedangkan header
  artikel mencantumkan Roberto Camagni, Roberta Capello, dan Andrea
  Caragliu; review mengikuti identitas tiga penulis pada header artikel.
\item
  Locator utama menggunakan rentang baris TXT dan nama bagian, tabel,
  atau gambar yang terlihat dalam ekstraksi. Nomor halaman PDF hanya
  digunakan bila ditampilkan secara jelas dalam teks ekstraksi; nomor
  halaman bukan hasil penelusuran PDF terpisah.
\item
  Ekstraksi \texttt{-layout} mempertahankan header, footer, pemisah
  halaman, dan artefak karakter. Sejumlah tanda minus pada tabel tampil
  seperti prefiks \texttt{2}, misalnya \texttt{20.04}, \texttt{20.28},
  \texttt{20.004}, atau \texttt{20.30}; nilai ambigu tidak dikoreksi
  diam-diam. Ringkasan numerik karena itu mengutamakan prosa sumber,
  pola signifikansi, dan angka yang terbaca jelas (TXT baris 672-723,
  832-870).
\item
  Artikel menyatakan \emph{random effects}, uji Breusch-Pagan, 158
  observasi pada tabel estimasi, database appendix 100 kawasan
  metropolitan dengan dua observasi per kota, dan sampel deskriptif 136
  kawasan metropolitan. Perbedaan cakupan tersebut dipertahankan sebagai
  laporan tahap analisis yang berbeda, bukan dipaksa menjadi satu ukuran
  sampel (TXT baris 404-417, 706-713, 729-734, 1102-1123).
\item
  Label \texttt{{[}Inferensi\ pengolahan{]}} tidak mewakili temuan atau
  interpretasi penulis. Batas utama inferensi adalah cakupan EU15 pada
  estimasi utama, proksi ALB berupa harga rumah, proksi jaringan berupa
  proyek FP5/FP6, ketidakseragaman periode data, dan strategi
  identifikasi kausal: Tidak disebutkan dalam file.
\end{itemize}

\section{Review B12}\label{review-b12}

Status: selesai setelah audit isi, struktur, locator, dan batasan
sumber.

Sumber kerja:
\texttt{source/journal/B12-parkinson-meegan-karecha-2015-city-size-and-economic-performance-is-bigger-better-small-more-beautiful-or-middling-marvellous-10.1080-09654313.2014.904998.txt}

\subsection{Summarized Abstract}\label{summarized-abstract-55}

{[}Laporan sumber{]} Artikel membahas kontribusi kota tingkat kedua
terhadap kinerja ekonomi nasional di Eropa. Artikel meninjau teori yang
bersaing tentang ukuran kota, investasi, dan kinerja ekonomi; menyajikan
bukti tentang lebih dari 150 kota ibu kota dan kota tingkat kedua di 31
negara; lalu merumuskan pesan kebijakan untuk pembuat kebijakan lokal,
nasional, dan Eropa. Bukti yang dilaporkan mendukung desentralisasi
tanggung jawab, kewenangan, dan sumber daya, penyebaran investasi, serta
peningkatan kinerja di banyak kota, bukan konsentrasi investasi hanya di
ibu kota. (Locator: TXT baris 49-62; halaman cetak PDF 1054.)

{[}Interpretasi penulis{]} Pemerintah nasional pada masa austerity
sebaiknya menahan tekanan untuk memusatkan investasi di ibu kota dan
lebih banyak berinvestasi di kota tingkat kedua apabila kesenjangan
dengan ibu kota besar dan membesar, infrastruktur bisnis kota tingkat
kedua lemah karena kurang investasi nasional, dan terdapat bukti
eksternalitas negatif pertumbuhan ibu kota. Artikel juga berpendapat
bahwa Komisi Eropa perlu mengembalikan isu kota tingkat kedua ke agenda
ekonominya. (Locator: TXT baris 56-62 dan 514-548; halaman cetak PDF
1054 dan 1066-1067.)

\subsection{Problem Statement}\label{problem-statement-55}

{[}Laporan sumber{]} Artikel menempatkan masalah pada dominasi ekonomi
dan kebijakan ibu kota, terutama ketika resesi, krisis Eurozone, dan
austerity meningkatkan persaingan atas sumber daya publik dan swasta.
Pertanyaan kunci yang dirumuskan adalah: ``Why should policy-makers
invest beyond the capital cities in an age of austerity?'' Pertanyaan
terkaitnya ialah apakah ukuran kota berpengaruh terhadap kinerja kota
dan ekonomi nasional, serta apakah investasi sebaiknya dikonsentrasikan
di ibu kota atau disebarkan ke lebih banyak kota. (Locator: TXT baris
66-104; halaman cetak PDF 1054-1055.)

{[}Interpretasi penulis{]} Risiko yang hendak diuji bukan sekadar
ketertinggalan kota tingkat kedua, melainkan kemungkinan bahwa
konsentrasi investasi di ibu kota memperlebar kesenjangan antarkota dan
mengurangi kinerja ekonomi nasional secara keseluruhan. (Locator: TXT
baris 88-103 dan 202-218; halaman cetak PDF 1055-1057.)

\subsection{Objectives}\label{objectives-55}

{[}Laporan sumber{]} Artikel menyatakan bahwa pembahasannya
mengeksplorasi pertanyaan kebijakan dan penelitian yang muncul dari
perdebatan tersebut, meninjau teori yang bersaing, menyajikan bukti
komparatif, dan mengidentifikasi pesan kebijakan. (Locator: TXT baris
104-105 dan 49-62; halaman cetak PDF 1055 dan 1054.)

{[}Inferensi pengolahan{]} Dari tujuan yang dinyatakan itu, sasaran
kerja artikel dapat diringkas sebagai: membandingkan kontribusi ibu kota
dan kota tingkat kedua; memeriksa kaitan antara desentralisasi,
penyebaran investasi, dan kinerja ekonomi; menantang asumsi bahwa
investasi ibu kota selalu menghasilkan keuntungan nasional terbesar;
serta menerjemahkan bukti tersebut menjadi arahan kebijakan. Ini adalah
perumusan ulang tujuan dari teks, bukan daftar tujuan formal yang
ditulis terpisah oleh penulis. (Locator: TXT baris 184-218 dan 220-263;
halaman cetak PDF 1057-1058.)

\subsection{Literature Survey}\label{literature-survey-55}

{[}Laporan sumber{]} Bagian kajian membahas beberapa kerangka teoritis:
teori pertumbuhan endogen berbasis neoklasik, ekonomi geografis atau New
Economic Geography, serta teori institusional dan evolusioner. Ekonomi
aglomerasi dijelaskan melalui spesialisasi, increasing returns, ekonomi
urbanisasi, jaringan, kepercayaan, dan modal sosial. (Locator: TXT baris
108-127; halaman cetak PDF 1055-1056.)

{[}Laporan sumber{]} Perbedaan teori tersebut dipasangkan dengan dua
paradigma kebijakan. Posisi neoklasik mendukung aglomerasi ibu kota dan
pendekatan place-neutral atau space-blind dengan layanan publik
universal serta intervensi spasial yang terbatas. Perspektif
institusional dan evolusioner menyoroti diseconomies, eksternalitas
negatif, diminishing marginal returns, pentingnya middle regions, dan
kebutuhan kebijakan place-based yang menggabungkan pendekatan bottom-up
dan top-down. (Locator: TXT baris 130-182; halaman cetak PDF 1056-1057.)

{[}Interpretasi penulis{]} Penulis mendukung perspektif place-based dan
desentralisasi, tetapi tidak menyatakan bahwa aglomerasi tidak
bermanfaat. Argumennya adalah bahwa manfaat aglomerasi memiliki batas
dan kontribusi kolektif wilayah atau kota di luar puncak hierarki perlu
diperhitungkan. Rujukan Markusen et al.~disebut telah mengkaji fenomena
kota tingkat kedua, terutama di luar Eropa, sedangkan artikel ini
menyatakan menambahkan informasi dan mengembangkan konsep tersebut untuk
konteks Eropa. (Locator: TXT baris 156-182 dan 202-218; halaman cetak
PDF 1056-1057.)

\subsection{Method Used}\label{method-used-55}

{[}Laporan sumber{]} Artikel bertumpu pada studi ESPON yang mencakup 124
kota tingkat kedua dan 31 ibu kota di Eropa. Bukti yang dipakai dalam
artikel berasal dari analisis data kuantitatif, tinjauan kebijakan, dan
studi kota individual. Artikel menegaskan bahwa hanya sebagian terbatas
dari bukti kuantitatif itu yang dapat ditampilkan. (Locator: TXT baris
184-218 dan 254-263; halaman cetak PDF 1057-1058.)

{[}Laporan sumber{]} Analisis yang tampak dalam teks membandingkan
kinerja kota dan kota-wilayah pada masa boom 2000-2007, perubahan
pertumbuhan 2000-2008, dan masa resesi 2008-2010, dengan contoh lintas
negara serta pembahasan khusus Jerman. (Locator: TXT baris 265-340 dan
368-381; halaman cetak PDF 1058-1063.)

{[}Inferensi pengolahan{]} Berdasarkan bentuk bukti yang dijelaskan,
metode artikel ini paling tepat dibaca sebagai analisis komparatif
deskriptif yang dipadukan dengan sintesis kebijakan dan studi kasus
ilustratif. Model statistik, prosedur estimasi, strategi identifikasi
kausal, dan prosedur pemilihan kota tidak diuraikan dalam TXT. (Locator:
TXT baris 184-218 dan 254-263; halaman cetak PDF 1057-1058.)

\subsection{Dataset}\label{dataset-55}

{[}Laporan sumber{]} Bahan empiris utama adalah studi ESPON tentang 124
kota tingkat kedua dan 31 ibu kota di 31 negara Eropa. Indikator yang
disebutkan meliputi populasi, pekerjaan, output atau GDP, pertumbuhan
GDP, produktivitas, inovasi, serta kontribusi terhadap ekonomi nasional.
Periode yang digunakan dalam bagian hasil mencakup boom 2000-2007,
pertumbuhan kota-wilayah 2000-2008, dan resesi 2008-2010. (Locator: TXT
baris 184-201, 284-340, dan 368-381; halaman cetak PDF 1057-1063.)

{[}Laporan sumber{]} Gambar 1 berisi total GDP dalam PPS tahun 2007 dan
sumbernya dicantumkan sebagai Parkinson et al.~(2012). Gambar 2 berisi
perubahan persentase riil tahunan rata-rata total GDP kota tingkat kedua
pada 2000-2007 dan sumbernya Eurostat. Map 1 bersumber dari Parkinson et
al.~(2014) untuk pertumbuhan ekonomi kota-wilayah 2000-2008, sedangkan
Map 2 bersumber dari Parkinson et al.~(2014) untuk resesi 2008-2010.
(Locator: TXT baris 323-347 dan 368-381; halaman cetak PDF 1060-1063.)

{[}Inferensi pengolahan{]} Teks tidak memberikan kamus variabel, format
data mentah, unit setiap nilai, atau prosedur penggabungan sumber data
secara rinci. Untuk rincian tersebut: Tidak disebutkan dalam file.
(Locator: TXT baris 184-263 dan keterangan Gambar 1-2 serta Map 1-2.)

\subsection{Population / Sample}\label{population-sample-55}

{[}Laporan sumber{]} Unit kajian adalah kota atau kota-wilayah, bukan
responden individu. Cakupannya terdiri atas 124 kota tingkat kedua dan
31 ibu kota di 31 negara Eropa. Kota tingkat kedua didefinisikan sebagai
kota di luar ibu kota yang kinerja ekonomi dan sosialnya cukup penting
untuk memengaruhi potensi kinerja ekonomi nasional; istilah itu tidak
berarti kota tersebut selalu kota terbesar kedua atau kurang penting.
Kota-kota tersebut mencakup hampir 80\% populasi urban metropolitan
Eropa menurut teks. (Locator: TXT baris 184-201; halaman cetak PDF
1057.)

{[}Laporan sumber{]} Penulis mengakui bahwa pembedaan ibu kota dan kota
tingkat kedua bersifat administratif, bukan fungsional, meskipun dalam
banyak negara ibu kota juga merupakan kota terbesar secara ekonomi dan
demografis. (Locator: TXT baris 202-218; halaman cetak PDF 1057.)

{[}Inferensi pengolahan{]} Prosedur sampling atau kriteria inklusi
setiap kota selain definisi operasional tersebut tidak dipaparkan dalam
TXT. Sampel responden manusia: Tidak berlaku. (Locator: TXT baris
184-218; halaman cetak PDF 1057.)

\subsection{Dependent Variables}\label{dependent-variables-55}

{[}Inferensi pengolahan{]} Tidak disebutkan dalam file adanya model
statistik dengan variabel dependen formal. Artikel melaporkan ukuran
kinerja berupa total GDP, pertumbuhan GDP tahunan, produktivitas,
populasi, pekerjaan, output, dan inovasi. (Locator: TXT baris 184-201,
265-305, dan 331-339; halaman cetak PDF 1057-1061.)

{[}Inferensi pengolahan{]} Ukuran-ukuran tersebut dapat dibaca sebagai
indikator hasil yang dibandingkan antar kota, tetapi tidak boleh
diperlakukan sebagai variabel dependen yang diestimasi dalam model
tertentu. Definisi operasional, transformasi, dan unit pengukurannya:
Tidak disebutkan dalam file. (Locator: TXT baris 184-263 dan 284-340;
halaman cetak PDF 1057-1061.)

\subsection{Results}\label{results-55}

{[}Laporan sumber{]} Untuk periode 2000-2007, dalam negara federal
seluruh kota tingkat kedua di Jerman dan Austria serta setengah kota
tingkat kedua di Belgia mengungguli ibu kotanya. Dalam negara regional,
seluruh kota tingkat kedua di Spanyol dan sepertiga di Italia tumbuh
lebih cepat daripada ibu kota. Di negara-negara Nordik, seluruh kota
tingkat kedua yang dikaji tumbuh lebih cepat daripada ibu kotanya.
Sebaliknya, di negara kesatuan yang terpusat, seluruh kota tingkat kedua
di Hongaria, Slovakia, Slovenia, Estonia, Lituania, dan Bulgaria, serta
semua kecuali satu di Republik Ceko, memiliki pertumbuhan lebih rendah
daripada ibu kota; hanya di Rumania, Latvia, dan Kroasia beberapa kota
tingkat kedua mengungguli ibu kota. (Locator: TXT baris 220-242; halaman
cetak PDF 1057-1058.)

{[}Laporan sumber{]} Ibu kota tetap mendominasi sistem urban Eropa dalam
populasi, pekerjaan, dan output. Pada 2007, total GDP ibu kota lebih
besar daripada kota non-ibu kota terdepan di semua negara kecuali Jerman
dan Italia. Di 19 negara, total GDP ibu kota lebih dari dua kali GDP
kota non-ibu kota terbesar, dan di Inggris, Prancis, Hongaria, serta
Latvia perbandingannya mencapai delapan kali. Namun, di 16 dari 26
negara, sedikitnya satu kota tingkat kedua memiliki pertumbuhan GDP
tahunan lebih tinggi daripada ibu kotanya pada masa boom. (Locator: TXT
baris 284-305; halaman cetak PDF 1059-1060.)

{[}Laporan sumber{]} Kinerja masa boom berbeda menurut wilayah. Kelompok
kota-wilayah di Eropa Utara, Tengah, dan Barat menampung lebih dari
empat perlima kota-wilayah terkemuka. Di negara kesatuan bekas sosialis
yang disebutkan, hanya tiga dari 27 kota-wilayah yang terkemuka dan tiga
berada pada tingkat menengah, sementara mayoritas tertinggal.
Kota-wilayah yang tertinggal di Eropa Timur dan Eropa Tenggara mengalami
lonjakan pertumbuhan; teks menyatakan laju rata-ratanya enam dan tujuh
kali laju wilayah mapan di Barat dan Utara. (Locator: TXT baris 307-328;
halaman cetak PDF 1059-1060.)

{[}Laporan sumber{]} Pada masa krisis, lebih dari 75\% kota mengalami
penurunan GDP pada 2008-2010 dan ibu kota berkinerja lebih baik daripada
kota tingkat kedua. Sembilan belas tempat dengan pertumbuhan tercepat
berada di Eropa Timur, termasuk 12 tempat di Polandia. Di Jerman hanya
Berlin yang tumbuh; semua kota Jerman lainnya menurun. Seluruh 14 kota
di Inggris dan seluruh 12 kota di Italia menurun, sedangkan delapan dari
sembilan kota di Spanyol menurun. (Locator: TXT baris 331-339; halaman
cetak PDF 1060.)

{[}Laporan sumber{]} Jerman menjadi contoh khusus: pada 2000-2007,
populasi meningkat lebih cepat daripada Berlin di enam kota tingkat
kedua, produktivitas tumbuh lebih cepat daripada Berlin di seluruh 14
kota tingkat kedua, lima dari sepuluh kota teratas dalam pertumbuhan GDP
berada di Jerman, dan lima dari sepuluh kota paling inovatif juga berada
di Jerman. (Locator: TXT baris 265-282; halaman cetak PDF 1058-1059.)

\subsection{Findings}\label{findings-55}

{[}Interpretasi penulis{]} Temuan utama penulis adalah bahwa
desentralisasi tanggung jawab, kewenangan, dan sumber daya, bersama
penyebaran investasi ke berbagai kota, dapat menghasilkan manfaat
nasional. Kota tingkat kedua secara individual dapat tertinggal dari ibu
kota, tetapi kontribusi kolektifnya terhadap kinerja ekonomi nasional
tetap sangat penting dan dalam banyak kasus lebih besar daripada
kontribusi ibu kota. Penulis secara eksplisit tidak mengklaim bahwa
setiap kota tingkat kedua selalu berkinerja baik atau mengungguli ibu
kota. (Locator: TXT baris 220-263; halaman cetak PDF 1057-1058.)

{[}Interpretasi penulis{]} Penulis menyimpulkan bahwa hubungan ibu kota
dan kota tingkat kedua bersifat win-win, bukan zero-sum. Ibu kota tetap
diperlukan untuk posisi global dan daya saing nasional, tetapi dominasi
berlebihan dapat menghasilkan pembangunan yang timpang. Manfaat
aglomerasi juga memiliki batas; kemacetan, kelangkaan lahan, sprawl,
marginalisasi modal manusia, dan penurunan kualitas infrastruktur
disebut sebagai eksternalitas atau diseconomies yang dapat mengurangi
daya saing ibu kota. (Locator: TXT baris 371-420; halaman cetak PDF
1062-1064.)

{[}Inferensi pengolahan{]} Bukti yang disajikan lebih mendukung
kesimpulan bersyarat menurut konteks tata kelola, periode, dan wilayah
daripada aturan bahwa kota besar atau kota tingkat kedua selalu
berkinerja lebih baik. Pembacaan ini mengikuti variasi lintas negara dan
peringatan penulis tentang keterbatasan generalisasi. (Locator: TXT
baris 202-218, 254-263, 307-340, dan 503-527; halaman cetak PDF
1057-1067.)

\subsection{Conclusions}\label{conclusions-55}

{[}Interpretasi penulis{]} Kesimpulan kebijakan artikel adalah bahwa ibu
kota dan kota tingkat kedua sama-sama perlu didukung. Pemerintah
sebaiknya membantu kota tingkat kedua menjadi investment-ready agar
dapat menampung sebagian pertumbuhan ibu kota ketika kapasitas ibu kota
mencapai batas dan biaya pertumbuhan mulai melampaui manfaatnya.
Investasi yang lebih besar di kota tingkat kedua terutama dibenarkan
apabila tiga kondisi terpenuhi: kesenjangan dengan ibu kota besar dan
membesar; infrastruktur bisnis kota tingkat kedua lemah karena kurang
investasi nasional; dan terdapat bukti jelas tentang eksternalitas
negatif pertumbuhan ibu kota. (Locator: TXT baris 417-425 dan 503-527;
halaman cetak PDF 1064-1067.)

{[}Interpretasi penulis{]} Kesimpulan tambahannya ialah bahwa
desentralisasi harus disertai kewenangan dan sumber daya yang sepadan,
tata kelola perlu bekerja pada skala ekonomi fungsional kota-wilayah,
strategi investasi teritorial perlu transparan, dan dampak belanja arus
utama terhadap kota tingkat kedua perlu dipantau. Untuk Eropa, Komisi
Eropa diminta mengakui potensi kota tingkat kedua dan mengintegrasikan
dampak teritorial kebijakan sektoral serta pendanaan arus utama.
(Locator: TXT baris 351-362, 441-500, dan 530-548; halaman cetak PDF
1061-1067.)

\subsection{Contributions}\label{contributions-55}

{[}Laporan sumber{]} Penulis menyatakan bahwa artikel atau kerja yang
mendasarinya menambahkan informasi baru dan mengembangkan konsep kota
tingkat kedua yang sebelumnya dibahas oleh Markusen et al.~Artikel juga
menyatukan analisis data kuantitatif, tinjauan kebijakan, dan studi kota
individual untuk membahas hubungan ukuran kota, investasi, tata kelola,
dan kinerja ekonomi di Eropa. (Locator: TXT baris 184-218 dan 254-263;
halaman cetak PDF 1057-1058.)

{[}Interpretasi penulis{]} Kontribusi argumentatif yang diklaim adalah
menantang asumsi bahwa ibu kota harus menjadi pilihan pertama investasi
demi keberhasilan nasional, tanpa menghapus pentingnya ibu kota. Artikel
menawarkan pesan kebijakan yang menempatkan penguatan ibu kota dan kota
tingkat kedua sebagai strategi yang saling melengkapi. (Locator: TXT
baris 243-263 dan 397-425; halaman cetak PDF 1058 dan 1064.)

\subsection{Research Gap}\label{research-gap-55}

{[}Laporan sumber{]} Artikel menyatakan bahwa penelitian akademik
sebelumnya banyak berfokus pada kota global yang lebih besar, sehingga
pembuat kebijakan perlu memberi perhatian lebih besar pada kontribusi
potensial dan aktual kota tingkat kedua. Artikel juga menyatakan bahwa
banyak negara belum memiliki kebijakan eksplisit untuk kota tingkat
kedua, sehingga kepentingan kolektifnya terabaikan. (Locator: TXT baris
82-87 dan 397-415; halaman cetak PDF 1055 dan 1064.)

{[}Laporan sumber{]} Artikel meminta pembuat kebijakan mengkaji secara
sistematis bagaimana keputusan mereka memengaruhi kota tingkat kedua dan
meminta pemerintah memantau serta mempublikasikan dampak teritorial
program belanja. (Locator: TXT baris 254-263 dan 482-500; halaman cetak
PDF 1058 dan 1065-1066.)

{[}Inferensi pengolahan{]} TXT tidak memuat bagian formal berjudul
research gap atau agenda metodologis yang lebih rinci. Gap yang dapat
dicatat tanpa informasi eksternal adalah kurangnya perhatian sistematis
terhadap kota tingkat kedua dan dampak teritorial keputusan investasi,
sebagaimana dinyatakan atau diminta oleh penulis. (Locator: TXT baris
82-87, 254-263, dan 482-500.)

\subsection{Limitations}\label{limitations-55}

{[}Laporan sumber{]} Penulis mengakui bahwa semua tipologi kota memiliki
keterbatasan. Pembedaan ibu kota dan kota tingkat kedua bersifat
administratif, bukan fungsional; kota tingkat kedua sangat beragam; dan
istilah itu tidak menunjuk satu kategori ukuran yang seragam. (Locator:
TXT baris 184-218; halaman cetak PDF 1057.)

{[}Laporan sumber{]} Penulis menegaskan bahwa Jerman bersifat unik
karena sistem federal, perpindahan ibu kota, sejarah pembagian dan
penyatuan, kewenangan kota tingkat kedua, sistem perbankan regional, dan
perusahaan menengahnya. Karena itu, negara lain tidak dapat begitu saja
meniru karakteristik struktural Jerman. Artikel juga hanya menyajikan
sebagian terbatas dari bukti kuantitatif studi yang lebih besar.
(Locator: TXT baris 265-282 dan 254-263; halaman cetak PDF 1058-1059.)

{[}Laporan sumber{]} Penulis membatasi klaimnya dengan menyatakan bahwa
bukti tidak menunjukkan semua kota tingkat kedua di semua negara
berkinerja baik atau mengungguli ibu kota. (Locator: TXT baris 254-263;
halaman cetak PDF 1058.)

{[}Inferensi pengolahan{]} Karena TXT tidak menguraikan model, prosedur
estimasi, definisi operasional lengkap, atau strategi identifikasi
kausal, hubungan antara desentralisasi dan kinerja tidak dapat dibaca
sebagai estimasi kausal formal dari artikel ini. Keterbatasan rincian
teknis tersebut adalah batas informasi TXT, bukan keterbatasan yang
secara eksplisit diklaim penulis. (Locator: TXT baris 184-218 dan
254-263.)

\subsection{Challenges}\label{challenges-55}

{[}Laporan sumber{]} Tantangan konteksnya adalah resesi, krisis
Eurozone, austerity, persoalan ekonomi dan sosial yang membesar, serta
kelangkaan sumber daya publik dan swasta. Kondisi itu berisiko
memperlebar kesenjangan antara kota tingkat kedua dan ibu kota.
(Locator: TXT baris 88-103; halaman cetak PDF 1055.)

{[}Laporan sumber{]} Tantangan tata kelola meliputi pendelegasian
tanggung jawab tanpa kewenangan dan sumber daya yang sepadan, pembagian
peran yang tidak jelas, ketergantungan kapasitas fiskal pada transfer,
lemahnya pemerintah regional yang dipilih secara demokratis, fragmentasi
tata kelola metropolitan, dan batas administratif yang tidak sesuai
dengan wilayah ekonomi fungsional. (Locator: TXT baris 351-362, 441-479;
halaman cetak PDF 1061 dan 1065.)

{[}Interpretasi penulis{]} Tantangan kebijakan Eropa adalah memastikan
bukan hanya dana yang ditargetkan secara eksplisit, tetapi juga seluruh
pendanaan arus utama Komisi Eropa, memperhitungkan dampaknya pada kota
tingkat kedua secara koheren. (Locator: TXT baris 530-548; halaman cetak
PDF 1066-1067.)

\subsection{Practical Implications}\label{practical-implications-55}

{[}Interpretasi penulis{]} Pemerintah nasional perlu mempertahankan ibu
kota yang kuat, tetapi tidak menjadikan ibu kota satu-satunya tujuan
investasi. Investasi di kota tingkat kedua perlu diprioritaskan dengan
mempertimbangkan tiga kondisi yang dinyatakan artikel: kesenjangan yang
besar dan membesar, infrastruktur bisnis yang lemah karena kurang
investasi nasional, dan eksternalitas negatif pertumbuhan ibu kota.
(Locator: TXT baris 503-527; halaman cetak PDF 1066-1067.)

{[}Interpretasi penulis{]} Desentralisasi perlu memasangkan tanggung
jawab dengan kewenangan dan sumber daya. Kebijakan akan lebih efektif
bila memberi ruang bagi kondisi lokal, membangun kapasitas manusia dan
fiskal serta otonomi kota, menjaga konsistensi dan transparansi, dan
didukung kepercayaan serta tanggung jawab bersama. (Locator: TXT baris
351-362 dan 456-466; halaman cetak PDF 1061 dan 1064-1065.)

{[}Interpretasi penulis{]} Pemerintah sebaiknya menaikkan skala tata
kelola ke tingkat kota-wilayah fungsional, mendorong kolaborasi,
menjelaskan kriteria investasi teritorial, serta memantau dan
mempublikasikan dampak belanja terhadap kota-wilayah. Kebijakan inovasi,
riset dan pengembangan, pendidikan dan keterampilan, transportasi,
konektivitas, dan infrastruktur disebut sebagai saluran penting.
(Locator: TXT baris 469-500; halaman cetak PDF 1065-1066.)

\subsection{Applications}\label{applications-55}

{[}Laporan sumber{]} Artikel menerapkan perbandingan ibu kota dan kota
tingkat kedua untuk membaca kontribusi kota terhadap daya saing
nasional, termasuk melalui analisis lintas negara, contoh Jerman, angka
pertumbuhan GDP, serta peta pertumbuhan dan resesi kota-wilayah.
(Locator: TXT baris 184-218, 265-340, dan 323-381; halaman cetak PDF
1057-1063.)

{[}Inferensi pengolahan{]} Berdasarkan rekomendasi artikel, keluaran ini
dapat digunakan untuk menilai strategi investasi teritorial nasional,
menguji dampak program arus utama pada kota-wilayah, dan merancang
koordinasi tata kelola pada skala fungsional. Ini adalah penggunaan
analitis yang diturunkan dari rekomendasi sumber, bukan laporan bahwa
aplikasi tersebut telah diimplementasikan dan dievaluasi. Evaluasi
implementasi: Tidak disebutkan dalam file. (Locator: TXT baris 469-500
dan 530-548; halaman cetak PDF 1065-1067.)

\subsection{Future Research
(Optional)}\label{future-research-optional-55}

{[}Laporan sumber{]} Artikel menyatakan bahwa pembuat kebijakan dan
peneliti perlu memberi perhatian lebih besar pada kontribusi kota
tingkat kedua pada masa mendatang dan bahwa dampak keputusan terhadap
kota tingkat kedua perlu diperiksa secara sistematis. Artikel juga
mendorong Komisi Eropa untuk mengakui potensi ekonomi kota tingkat kedua
dalam strategi masa depan. (Locator: TXT baris 82-87, 254-263, dan
530-548; halaman cetak PDF 1055, 1058, dan 1066-1067.)

{[}Inferensi pengolahan{]} Agenda riset minimal yang dapat diturunkan
secara eksplisit dari seruan tersebut adalah pemeriksaan sistematis atas
dampak keputusan investasi dan pemantauan dampak teritorialnya.
Rancangan penelitian, hipotesis, dan metode untuk agenda itu: Tidak
disebutkan dalam file. (Locator: TXT baris 254-263 dan 482-500; halaman
cetak PDF 1058 dan 1065-1066.)

\subsection{Provenance Notes}\label{provenance-notes-55}

\begin{itemize}
\tightlist
\item
  {[}Inferensi pengolahan{]} Kode yang diproses hanya B12. Verifikasi
  nama menemukan satu TXT aktual berprefiks \texttt{B12} di bawah
  \texttt{source/}, yaitu file sumber yang tercantum di awal review.
  Baris 1-2 menghubungkannya dengan PDF
  \texttt{journal/B12-parkinson-meegan-karecha-2015-city-size-and-economic-performance-is-bigger-better-small-more-beautiful-or-middling-marvellous-10.1080-09654313.2014.904998.pdf}.
\item
  {[}Inferensi pengolahan{]} Seluruh TXT dibaca, dari baris 1 sampai
  625, termasuk blok \texttt{{[}PDF\ Metadata{]}} pada baris 1-27,
  \texttt{{[}Extracted\ Text{]}} pada baris 28-625, halaman artikel,
  referensi, dan pemberitahuan hak cipta pada baris 621-625. Teks
  menyatakan metode ekstraksi \texttt{pdftotext\ -layout} dan tanggal
  ekstraksi 2026-08-31. (Locator: TXT baris 1-5 dan 621-625.)
\item
  {[}Laporan sumber{]} Metadata PDF mencatat judul
  \texttt{City\ Size\ and\ Economic\ Performance:\ Is\ Bigger\ Better,\ Small\ More\ Beautiful\ or\ Middling\ Marvellous?};
  Subject
  \texttt{European\ Planning\ Studies,\ 2014.\ doi:10.1080/09654313.2014.904998};
  Author \texttt{MICHAEL\ PARKINSON}; Creator
  \texttt{Arbortext\ Advanced\ Print\ Publisher\ 9.1.520/W}; Producer
  \texttt{Acrobat\ Distiller\ 8.1.0\ (Windows);\ modified\ using\ iTextSharp\ 5.4.5};
  CreationDate dan ModDate \texttt{Sat\ Apr\ 11\ 01:39:30\ 2015\ WIB}.
  (Locator: TXT baris 7-13.)
\item
  {[}Laporan sumber{]} Metadata PDF selanjutnya mencatat Custom Metadata
  \texttt{no}, Metadata Stream \texttt{yes}, Tagged \texttt{no},
  UserProperties \texttt{no}, Suspects \texttt{no}, Form \texttt{none},
  JavaScript \texttt{no}, Pages \texttt{16}, Encrypted \texttt{no}, Page
  size \texttt{493.228\ x\ 702.992\ pts}, Page rot \texttt{0}, File size
  \texttt{573888\ bytes}, Optimized \texttt{yes}, dan PDF version
  \texttt{1.6}. (Locator: TXT baris 14-27.)
\item
  {[}Laporan sumber{]} Header teks artikel mencatat
  \texttt{European\ Planning\ Studies,\ 2015}, Vol. 23, No.~6, halaman
  1054-1068, dan DOI yang sama. Header juga mencantumkan Michael
  Parkinson, Richard Meegan, dan Jay Karecha, afiliasi European
  Institute for Urban Affairs, Liverpool John Moores University, serta
  status received August 2013 dan accepted December 2013. (Locator: TXT
  baris 29-45.)
\item
  {[}Inferensi pengolahan{]} Ada perbedaan yang dipertahankan apa adanya
  antara metadata PDF yang mencatat Subject tahun 2014 dan Author
  tunggal dengan header artikel yang mencatat tahun 2015 serta tiga
  penulis. Review ini tidak menyelesaikan perbedaan tersebut memakai
  informasi eksternal. (Locator: TXT baris 7-13 dan 29-45.)
\item
  {[}Inferensi pengolahan{]} Label dalam review ini berarti:
  \texttt{Laporan\ sumber} adalah isi yang dilaporkan atau dinyatakan
  TXT; \texttt{Interpretasi\ penulis} adalah argumen atau pembacaan
  kebijakan yang dinyatakan oleh penulis artikel;
  \texttt{Inferensi\ pengolahan} adalah simpulan terbatas dari
  pengolahan TXT ini dan bukan kutipan atau klaim baru dari sumber.
\item
  {[}Inferensi pengolahan{]} Semua locator halaman menggunakan nomor
  halaman cetak yang muncul dalam ekstraksi, sedangkan locator baris
  merujuk pada nomor baris TXT yang dibaca. Tidak ada informasi
  eksternal yang digunakan dan file TXT sumber tidak diubah.
\end{itemize}

\section{Review B13}\label{review-b13}

\textbf{Kode:} B13

\textbf{Sumber yang direview:}
\texttt{source/journal/B13-meijers-burger-2010-spatial-structure-and-productivity-in-us-metropolitan-areas-10.1068-a42151.txt}

\subsection{Summarized Abstract}\label{summarized-abstract-56}

\textbf{Temuan sumber:} Artikel menganalisis pengaruh struktur spasial
yang berbeda, khususnya dimensi monosentrisitas-polisentrisitas dan
sentralisasi-dispersal, terhadap kinerja ekonomi kawasan metropolitan di
Amerika Serikat. Model ordinary least squares (OLS) dan two-stage least
squares (TSLS) untuk menjelaskan produktivitas tenaga kerja menunjukkan
bahwa struktur spasial berpengaruh: polisentrisitas berkaitan dengan
produktivitas tenaga kerja yang lebih tinggi. Sumber menyatakan bahwa
hasil ini tampak mendukung gagasan bahwa disekonomi aglomerasi relatif
terbatas di kawasan yang lebih polisentris dan bahwa eksternalitas
aglomerasi sampai tingkat tertentu dibagi antarkota dalam kawasan
tersebut. Namun, jaringan kota-kota kecil yang berdekatan secara
geografis tidak dapat menggantikan eksternalitas urbanisasi dari satu
kota besar (TXT lines 47-60; printed PDF p.~1383, abstract).

\textbf{Interpretasi penulis:} Temuan tersebut diposisikan sebagai
langkah empiris untuk menguji apakah kota-kota saling ``meminjam'' atau
berbagi ukuran dan apakah struktur metropolitan perlu dimasukkan ke
dalam analisis eksternalitas aglomerasi (TXT lines 117-139; printed PDF
pp.~1384-85).

\textbf{Inferensi pengolahan:} Review ini memperlakukan produktivitas
tenaga kerja sebagai keluaran empiris utama karena itulah outcome yang
dijelaskan dalam abstrak dan model sumber; tidak ada outcome utama lain
yang dilaporkan.

\subsection{Problem Statement}\label{problem-statement-56}

\textbf{Temuan sumber:} Analisis empiris aglomerasi sebelumnya disebut
mengabaikan organisasi spasial yang multisentris serta kemungkinan
``sharing'' atau ``borrowing'' size antarkota. Penelitian aglomerasi
juga sering menerima bentuk sistem perkotaan begitu saja, sementara
definisi kawasan metropolitan umumnya berangkat dari satu kota inti dan
hinterland-nya. Akibatnya, analisis tersebut tidak selaras dengan
perubahan skala geografis tempat aglomerasi bermanifestasi (TXT lines
105-116; printed PDF p.~1384).

\textbf{Interpretasi penulis:} Kekurangan empiris itu dirumuskan melalui
dua pertanyaan: apakah kota-kota benar-benar meminjam ukuran satu sama
lain, dan apakah kumpulan kota yang berdekatan dapat menjadi substitusi
bagi eksternalitas urbanisasi dari satu kota yang lebih besar (TXT lines
117-125; printed PDF p.~1384).

\subsection{Objectives}\label{objectives-56}

\textbf{Temuan sumber:} Tujuan utama artikel adalah memberikan langkah
awal untuk mengatasi kekurangan empiris dalam perdebatan tentang
pembagian eksternalitas aglomerasi di antara kota-kota yang lokasinya
relatif berdekatan. Pendekatan yang dipilih ialah memasukkan struktur
spasial kawasan metropolitan ke dalam analisis empiris eksternalitas
aglomerasi (TXT lines 117-128; printed PDF p.~1384).

\textbf{Temuan sumber:} Penulis tidak membatasi pengukuran pada dimensi
monosentrisitas-polisentrisitas, tetapi juga memasukkan dimensi
sentralisasi-dispersal. Produktivitas tenaga kerja digunakan sebagai
proksi kinerja metropolitan, lalu fungsi produksi yang sudah ada
diperluas dengan variabel spasial selain ukuran atau kepadatan rata-rata
(TXT lines 126-133; printed PDF p.~1384).

\subsection{Literature Survey}\label{literature-survey-56}

\textbf{Temuan sumber (sintesis pustaka penulis):} Literatur yang
dibahas menggambarkan perubahan dari kota dengan satu inti dan
hinterland rural menuju entitas perkotaan regional yang mencakup
beberapa kota yang terpisah secara fisik tetapi terhubung secara
fungsional. Dalam kerangka ini, eksternalitas ekonomi dipandang tidak
selalu terbatas pada satu inti, melainkan dapat dibagi di antara
permukiman yang saling terhubung. Konsep \emph{borrowed size} dari
Alonso digunakan untuk menjelaskan pendapatan yang lebih tinggi pada
kota kecil yang menjadi bagian dari kompleks megalopolitan dibandingkan
kota mandiri berukuran sama (TXT lines 62-104; printed PDF pp.~1383-84).

\textbf{Temuan sumber (sintesis pustaka penulis):} Ukuran dan kepadatan
dikaitkan dengan eksternalitas urbanisasi, seperti kumpulan tenaga kerja
yang besar dan beragam, infrastruktur, fasilitas publik, institusi
pengetahuan, interaksi lintas sektor, serta pasar internal. Penulis juga
membahas disekonomi urbanisasi yang dapat muncul melalui polusi,
kriminalitas, dan harga lahan atau perumahan yang tinggi (TXT lines
162-184; printed PDF p.~1385).

\textbf{Temuan sumber (sintesis pustaka penulis):} Pengaruh
polisentrisitas terhadap kinerja metropolitan disebut belum jelas dan
masih kekurangan penelitian empiris. Argumen yang mendukung
polisentrisitas menyatakan bahwa keuntungan aglomerasi dapat
diregionalisasi atau digantikan oleh jaringan, sedangkan kerugian
seperti kompetisi lahan dan tenaga kerja, kemacetan, serta paparan
polusi lebih banyak berada pada skala kota individual. Argumen tandingan
menekankan biaya perjalanan dan arus informasi yang lebih panjang serta
berkurangnya keuntungan yang bergantung pada kepadatan, kedekatan, dan
kontak tatap muka (TXT lines 185-236; printed PDF pp.~1385-86).

\textbf{Temuan sumber (sintesis pustaka penulis):} Perdebatan tentang
dispersal juga belum terselesaikan. Dispersal dibedakan dari sprawl
ber-kepadatan rendah karena yang diukur adalah apakah penduduk berada di
pusat-pusat perkotaan atau tersebar di luar pusat, bukan tingkat
kepadatannya. Sumber memuat argumen bahwa pembangunan tersebar dapat
meningkatkan kebutuhan lahan, infrastruktur, perjalanan, dan
ketergantungan pada mobil, tetapi juga memuat pandangan bahwa sprawl
dapat berkaitan dengan kualitas hidup. Artikel menguji dugaan bahwa
kawasan dengan proporsi penduduk lebih besar di tempat perkotaan akan
memiliki produktivitas tenaga kerja lebih tinggi (TXT lines 237-261;
printed PDF pp.~1386-87).

\textbf{Interpretasi penulis:} Dari literatur tersebut dibangun dugaan
bahwa ukuran metropolitan berhubungan positif dengan kinerja, struktur
yang lebih polisentris berdampak positif langsung, polisentrisitas
mengurangi efek ukuran terhadap produktivitas, dan dispersal yang lebih
rendah akan berkaitan dengan kinerja yang lebih baik. Sumber menyebut
pembahasan ini menghasilkan tiga asumsi yang dapat diuji, dengan
hubungan ukuran, monosentrisitas-polisentrisitas, dan
sentralisasi-dispersal sebagai kerangka utamanya (TXT lines 155-161,
178-184, 215-236, 259-261; printed PDF pp.~1385-87).

\subsection{Method Used}\label{method-used-56}

\textbf{Temuan sumber:} Penulis menggunakan fungsi produksi Cobb-Douglas
yang menghubungkan output metropolitan dengan modal, tenaga kerja, modal
manusia, input antara, dan lahan. Setelah dibagi dengan tenaga kerja dan
dilogaritmakan, produktivitas tenaga kerja dijelaskan oleh rasio input
tersebut, ukuran populasi metropolitan, polisentrisitas, dispersal,
serta dummy wilayah sensus. Asumsi constant returns to scale digunakan
dalam penurunan model (TXT lines 411-490; printed PDF pp.~1389-91).

\textbf{Temuan sumber:} Polisentrisitas diukur secara morfologis melalui
kemiringan garis regresi distribusi rank-size kota-kota yang berstatus
\emph{incorporated place}. Kemiringan dihitung untuk 2, 3, dan 4 tempat
terbesar di setiap kawasan metropolitan, kemudian dirata-ratakan.
Kemiringan yang lebih datar berarti kawasan lebih polisentris; tempat
dengan sedikitnya 5.000 penduduk pada 2006 dipertimbangkan, dengan
pengecualian pengukuran untuk tempat terbesar kedua yang tidak memenuhi
ambang tersebut. Kawasan yang dua kota terbesarnya berada dalam
\emph{urban area} yang sama dikeluarkan sebagai polisentris konurbasi
(TXT lines 294-393; printed PDF pp.~1387-89; Figure 2; footnote 1).

\textbf{Temuan sumber:} Dispersal diukur berdasarkan bagian populasi
metropolitan yang tidak tinggal di pusat perkotaan dengan sedikitnya
25.000 penduduk pada 2006. Posisi pada sumbu sentralisasi-dispersal
bergantung pada bagian populasi tersebut; bagian yang lebih besar
berarti lebih dispersal (TXT lines 394-403; printed PDF p.~1389).

\textbf{Temuan sumber:} Estimasi dilakukan dengan OLS dan TSLS. Karena
ukuran dan struktur metropolitan berpotensi bersifat simultan dengan
produktivitas, TSLS menggunakan lima variabel historis sebagai
instrumen: populasi 1950, polisentrisitas 1950, dispersal 1950,
keberadaan rel kereta api pada 1860, dan penggunaan lahan pertanian yang
diukur sebagai kepadatan pekerjaan pertanian di luar kawasan urban. Uji
Anderson, Cragg-Donald, dan Shea digunakan untuk relevansi; uji Sargan
dan Basmann untuk validitas; serta uji Wu-Hausman dan Durbin-Wu-Hausman
untuk eksogenitas (TXT lines 572-630, 632-711; printed PDF pp.~1392-94;
Table 3).

\textbf{Interpretasi penulis:} Karena uji pada spesifikasi dengan ukuran
dan struktur sebagai regressor tidak menunjukkan bukti hubungan endogen,
OLS dipilih sebagai estimasi yang lebih efisien. Perbandingan dengan
TSLS yang memperlakukan ukuran metropolitan sebagai endogen digunakan
sebagai uji ketahanan (TXT lines 717-732; printed PDF p.~1395).

\textbf{Inferensi pengolahan:} Karena pengukuran utama dan observasi
berpusat pada 2006, desain ini diklasifikasikan di sini sebagai analisis
kuantitatif satu-periode; label desain tersebut adalah inferensi
pengolahan, bukan istilah eksplisit yang digunakan sumber.

\subsection{Dataset}\label{dataset-56}

\textbf{Temuan sumber:} Data dikumpulkan untuk metropolitan statistical
areas (MSAs) dan, bila menjadi bagian dari combined statistical area
(CSA), digunakan definisi CSA yang paling luas. Cakupan utama adalah
kawasan metropolitan di Amerika Serikat kontinental dengan populasi
lebih dari 250.000 penduduk pada 2006. Tabel statistik deskriptif
melaporkan 113 observasi (TXT lines 491-519; printed PDF p.~1391; Table
2).

\textbf{Temuan sumber:} Produktivitas tenaga kerja dihitung sebagai GDP
kawasan metropolitan dalam dolar AS riil pada 2006 dibagi jumlah
pekerjaan pada sektor-sektor yang disertakan. Pertanian, perikanan,
perburuan, pertambangan, dan administrasi publik dikeluarkan; angka
untuk pekerja mandiri tidak tersedia. GDP menurut kawasan dan sektor
berasal dari Bureau of Economic Analysis, jumlah pekerjaan dari Bureau
of Labor Statistics, dan data sektor individual dari American Community
Survey 2006 (TXT lines 503-518; printed PDF p.~1391).

\textbf{Temuan sumber:} Variabel input dan kontrol mencakup rasio
modal-tenaga kerja, rasio lahan-tenaga kerja, pendidikan per pekerja,
ukuran metropolitan, polisentrisitas, dispersal, dan dummy sembilan
divisi sensus. Rasio modal-tenaga kerja dibentuk sebagai rata-rata
tertimbang rasio sektor dari Annual Survey of Manufacturers berdasarkan
komposisi industri 15 sektor. Rasio lahan-tenaga kerja adalah rata-rata
acre per pekerja. Pendidikan per pekerja adalah persentase penduduk
berusia setidaknya 25 tahun yang memiliki gelar sarjana atau lebih
tinggi, dari ACS 2006 (TXT lines 545-571; printed PDF p.~1392).

\textbf{Temuan sumber:} Semua variabel non-dummy dalam analisis empiris
ditransformasikan ke logaritma. GDP tidak dipasangkan dengan input
antara karena GDP didefinisikan sebagai total produksi dikurangi input
antara (TXT lines 517-518, 557-559; printed PDF pp.~1391-92). Variabel
historis untuk TSLS meliputi data populasi, polisentrisitas, dan
dispersal 1950, rel kereta api 1860, serta penggunaan lahan pertanian di
luar kawasan urban (TXT lines 607-630; printed PDF p.~1393).

\subsection{Population / Sample}\label{population-sample-56}

\textbf{Temuan sumber:} Unit analisisnya adalah kawasan metropolitan di
Amerika Serikat kontinental. Data dikumpulkan untuk metropolitan
statistical areas (MSAs), tetapi jika sebuah MSA merupakan bagian dari
combined statistical area (CSA), definisi CSA terluas digunakan agar
kota dianalisis dalam konteks spasial yang lebih luas. MSA yang tidak
menjadi bagian CSA tetap digunakan. Populasi kawasan yang disyaratkan
lebih dari 250.000 penduduk pada 2006, dan statistik deskriptif
melaporkan 113 observasi (TXT lines 491-502, 519-539; printed PDF
p.~1391; Table 2).

\textbf{Temuan sumber:} Dalam lampiran, MSA dijelaskan sebagai kawasan
dengan inti urban sedikitnya 50.000 penduduk dan kabupaten terkait yang
memenuhi kriteria hubungan kerja; MSA yang berdekatan dapat digabung
menjadi CSA berdasarkan ukuran pertukaran pekerjaan. Kawasan
mikropolitan umumnya tidak dipertimbangkan, kecuali menjadi bagian CSA
yang juga mencakup sedikitnya satu MSA (TXT lines 1108-1123; printed PDF
p.~1402; Appendix).

\textbf{Temuan sumber:} Beberapa kasus dikeluarkan karena data kawasan
mikropolitan tidak tersedia, data New Orleans dinilai bias akibat
Hurricane Katrina, dan kawasan yang diklasifikasikan sebagai
\emph{polycentric conurbations} tidak dimasukkan. Daftar conurbation
yang dikeluarkan tercantum di lampiran (TXT lines 1124-1149; printed PDF
p.~1402; Appendix). Sumber menyatakan ada dua CSA dengan lebih dari 20\%
penduduk berada di kawasan mikropolitan, tetapi kemudian mencantumkan
tiga nama kawasan; ketidaksesuaian ini dipertahankan sebagai catatan
provenance dan tidak diselesaikan dengan dugaan.

\subsection{Dependent Variables}\label{dependent-variables-56}

\textbf{Temuan sumber:} Variabel dependen tunggal adalah produktivitas
tenaga kerja metropolitan pada 2006. Produktivitas dihitung sebagai GDP
kawasan metropolitan dalam dolar AS riil dibagi jumlah pekerjaan pada
sektor-sektor yang disertakan. Pertanian, perikanan, perburuan,
pertambangan, dan administrasi publik dikeluarkan; angka pekerja mandiri
tidak tersedia. Dalam persamaan empiris, variabel ini ditransformasikan
ke logaritma (TXT lines 131-133, 464-490, 503-518; printed PDF pp.~1384,
1390-91).

\textbf{Temuan sumber:} Produktivitas tenaga kerja digunakan sebagai
proksi kinerja metropolitan dan sebagai ukuran empiris eksternalitas
yang dianalisis. Tidak ada variabel dependen lain yang dilaporkan dalam
file (TXT lines 131-133, 865-869; printed PDF pp.~1384, 1397).

\subsection{Results}\label{results-56}

\textbf{Temuan sumber - pengujian endogenitas:} Instrumen untuk ukuran
metropolitan, polisentrisitas, dan dispersal dinyatakan relevan
berdasarkan Anderson canonical correlation, Cragg-Donald F-test, dan
Shea partial R2; uji Sargan dan Basmann tidak menolak validitas
instrumen. Uji Wu-Hausman dan Durbin-Wu-Hausman menunjukkan struktur
spasial dapat diperlakukan sebagai eksogen. Untuk spesifikasi yang hanya
menginstrumentasikan ukuran metropolitan, Durbin-Wu-Hausman menolak
eksogenitas pada tingkat 5\% dengan p = 0,021; untuk spesifikasi dengan
ukuran dan struktur secara simultan, tidak ada bukti hubungan endogen
(TXT lines 632-711, 717-732; printed PDF pp.~1393-95; Table 3).

\textbf{Interpretasi penulis:} Karena instrumen dinilai relevan dan
valid, serta spesifikasi simultan tidak menunjukkan hubungan endogen
antara produktivitas dan ukuran maupun struktur metropolitan, OLS
dipilih sebagai estimator yang lebih efisien. Model TSLS yang
memperlakukan ukuran metropolitan sebagai endogen digunakan sebagai
pemeriksaan ketahanan (TXT lines 717-732; printed PDF p.~1395).

\textbf{Temuan sumber - model dasar:} Model OLS pada Table 4 menjelaskan
sekitar dua pertiga variasi produktivitas tenaga kerja, dengan R2 = 0,66
dan 113 observasi. Koefisien rasio modal-tenaga kerja adalah 0,754 dan
signifikan pada p \textless{} 0,01; rasio lahan-tenaga kerja adalah
0,013 dan pendidikan per pekerja adalah -0,007, keduanya tidak
signifikan. Ukuran metropolitan memiliki koefisien 0,107 dan signifikan
pada p \textless{} 0,01; polisentrisitas memiliki koefisien 0,055 dan
signifikan pada p \textless{} 0,05; dispersal memiliki koefisien 0,024
dan tidak signifikan (TXT lines 733-763; printed PDF pp.~1395-96; Table
4).

\textbf{Temuan sumber:} Penulis melaporkan bahwa penggandaan ukuran
metropolitan meningkatkan produktivitas tenaga kerja lebih dari 10\%,
sedangkan penggandaan derajat polisentrisitas meningkatkan produktivitas
tenaga kerja sebesar 5,5\%. Tidak ditemukan pengaruh derajat dispersal
terhadap produktivitas tenaga kerja (TXT lines 774-790; printed PDF
p.~1396).

\textbf{Temuan sumber - pemeriksaan TSLS:} Pada model TSLS di Table 4,
koefisien ukuran metropolitan adalah 0,078 dan rasio modal-tenaga kerja
0,843, keduanya signifikan pada p \textless{} 0,01. Koefisien
polisentrisitas adalah 0,051 dan signifikan pada p \textless{} 0,05;
dispersal adalah 0,030 dan tidak signifikan; R2 terpusat yang dilaporkan
adalah 0,65. Efek eksternalitas urbanisasi yang diestimasi TSLS lebih
dari 25\% lebih kecil daripada OLS, sedangkan estimasi polisentrisitas
dan dispersal tidak berbeda signifikan antar-model (TXT lines 746-768,
795-798; printed PDF pp.~1395-96; Table 4).

\textbf{Temuan sumber - interaksi:} Pada spesifikasi interaksi di Table
5, interaksi ukuran metropolitan dan polisentrisitas bernilai -0,059
pada OLS dan -0,063 pada TSLS, keduanya signifikan pada p \textless{}
0,01. Interaksi ukuran metropolitan dan dispersal bernilai -0,008 pada
OLS dan -0,016 pada TSLS, keduanya tidak signifikan. R2 OLS dan R2
terpusat TSLS sama-sama dilaporkan 0,68 (TXT lines 799-827; printed PDF
pp.~1396-97; Table 5).

\textbf{Interpretasi penulis:} Interaksi negatif dan signifikan antara
ukuran metropolitan dan polisentrisitas ditafsirkan sebagai bukti bahwa
eksternalitas urbanisasi lebih besar di kota atau kawasan yang lebih
monosentris, sehingga polisentrisitas mengurangi pengaruh ukuran
terhadap produktivitas. Penulis juga menyatakan bahwa manfaat
polisentrisitas mungkin lebih kuat pada kawasan metropolitan yang lebih
kecil (TXT lines 833-840; printed PDF p.~1397).

\subsection{Findings}\label{findings-56}

\textbf{Temuan sumber:} Pada bagian kesimpulan, penulis merangkum empat
temuan utama:

\begin{itemize}
\tightlist
\item
  Kawasan metropolitan dengan dispersal yang lebih besar tidak terbukti
  memiliki kinerja produktivitas tenaga kerja yang lebih buruk. Proporsi
  penduduk yang tinggal di pusat urban atau di luar pusat tidak
  ditemukan memengaruhi produktivitas (TXT lines 879-892; printed PDF
  p.~1397-98).
\item
  Kawasan metropolitan yang lebih polisentris memiliki produktivitas
  tenaga kerja yang lebih tinggi. Semakin merata penduduk tersebar di
  antara beberapa tempat dalam kawasan metropolitan, semakin tinggi
  produktivitas yang dilaporkan (TXT lines 893-910; printed PDF
  p.~1398).
\item
  Pengaruh ukuran metropolitan menurun ketika kawasan menjadi lebih
  polisentris. Kumpulan kota tidak menyediakan pengganti bagi
  eksternalitas urbanisasi dari satu kota besar, meskipun jumlah
  penduduk metropolitan secara keseluruhan dapat serupa (TXT lines
  911-927; printed PDF p.~1398).
\item
  Polisentrisitas tampak lebih menguntungkan kawasan metropolitan yang
  lebih kecil. Penulis menawarkan kemungkinan bahwa kota-kota dalam
  kawasan polisentris yang lebih kecil lebih terkait secara fungsional,
  tetapi menyatakan kemungkinan ini masih perlu diuji (TXT lines
  928-940; printed PDF pp.~1398-99).
\end{itemize}

\textbf{Interpretasi penulis:} Produktivitas yang lebih tinggi pada
kawasan polisentris dipandang mencerminkan keseimbangan yang lebih baik
antara ekonomi dan disekonomi aglomerasi: keuntungan urbanisasi dianggap
telah mengalami regionalisasi sampai tingkat tertentu, sedangkan
kerugian urbanisasi lebih banyak berada pada batas kota. Penulis
mengaitkan hasil ini dengan kebutuhan untuk mengonseptualisasikan
ekonomi aglomerasi secara relasional, bukan hanya pada inti urban
tunggal (TXT lines 897-910; printed PDF p.~1398).

\subsection{Conclusions}\label{conclusions-56}

\textbf{Interpretasi penulis:} Struktur spasial metropolitan dinilai
berpengaruh terhadap eksternalitas aglomerasi dan produktivitas tenaga
kerja. Polisentrisitas memiliki efek langsung yang positif, sedangkan
dispersal tidak menunjukkan efek; ukuran metropolitan tetap berhubungan
positif dengan produktivitas, tetapi keuntungan ukuran berkurang ketika
polisentrisitas meningkat. Kesimpulan ini mendukung adanya pembagian
eksternalitas di antara permukiman yang terhubung, tetapi sekaligus
membatasi klaim tersebut: jaringan beberapa kota kecil yang berdekatan
tidak sepenuhnya menggantikan eksternalitas urbanisasi satu kota besar
(TXT lines 865-881, 887-927; printed PDF pp.~1397-98).

\subsection{Contributions}\label{contributions-56}

\textbf{Inferensi pengolahan berbasis isi sumber:} Kontribusi empiris
artikel dapat diringkas sebagai berikut:

\begin{itemize}
\tightlist
\item
  Memasukkan struktur spasial metropolitan ke dalam analisis
  eksternalitas aglomerasi sebagai langkah awal untuk mengatasi
  kekurangan empiris yang diidentifikasi penulis (TXT lines 117-139;
  printed PDF p.~1384).
\item
  Menguji dua dimensi struktur, yaitu monosentrisitas-polisentrisitas
  dan sentralisasi-dispersal, alih-alih hanya menggunakan ukuran atau
  kepadatan rata-rata (TXT lines 126-133, 155-161, 263-275; printed PDF
  pp.~1384, 1385, 1387).
\item
  Memberikan bukti kuantitatif untuk 113 kawasan metropolitan AS pada
  2006 tentang efek langsung polisentrisitas, ketiadaan efek dispersal,
  dan interaksi ukuran-polisentrisitas (TXT lines 501-518, 733-840;
  printed PDF pp.~1391, 1395-97).
\item
  Menguji ketahanan hasil terhadap potensi simultanitas dengan TSLS dan
  instrumen historis, kemudian membandingkannya dengan OLS (TXT lines
  572-630, 717-732; printed PDF pp.~1392-95).
\end{itemize}

\subsection{Research Gap}\label{research-gap-56}

\textbf{Temuan sumber:} Sumber menyatakan bahwa analisis empiris
aglomerasi sebelumnya banyak mengabaikan organisasi spasial yang
multisentris dan kemungkinan kota saling berbagi atau meminjam ukuran.
Bentuk sistem urban sering dianggap begitu saja, sementara definisi
kawasan metropolitan yang berangkat dari satu kota inti tidak
menyesuaikan perubahan skala geografis tempat aglomerasi berlangsung.
Sumber juga menyatakan bahwa pengaruh polisentrisitas terhadap kinerja
metropolitan masih tidak jelas dan kekurangan penelitian empiris (TXT
lines 105-125, 185-214; printed PDF pp.~1384-86).

\textbf{Temuan sumber:} Artikel menyebut relasi antara kekuatan dan
jangkauan hubungan fungsional antarkota dengan kinerja sistem urban
regional, hubungan eksternal antarkawasan metropolitan yang lebih jauh,
perubahan kesenjangan eksternalitas dari waktu ke waktu, tingkat
disekonomi aglomerasi, dan perbedaan lintas sektor sebagai persoalan
yang belum terselesaikan (TXT lines 941-956; printed PDF p.~1399;
section 6.2).

\textbf{Inferensi pengolahan:} File tidak memakai heading khusus bernama
``Research Gap''; uraian di atas diambil dari pernyataan \emph{empirical
deficit}, kurangnya bukti, dan agenda riset dalam artikel, bukan dari
sumber eksternal.

\subsection{Limitations}\label{limitations-56}

\textbf{Temuan sumber:} Sumber tidak menyediakan bagian khusus berjudul
``Limitations'', tetapi menyatakan batasan operasional dan cakupan
berikut:

\begin{itemize}
\tightlist
\item
  Polisentrisitas diukur dari perspektif morfologis yang paling tidak
  restriktif, bukan dari hubungan fungsional. Penulis mencatat bahwa
  kategorisasi definitif memerlukan pengetahuan lebih lanjut tentang
  tingkat spesialisasi dan interaksi (TXT lines 294-307, 422-425;
  printed PDF pp.~1387-90; Figure 2 and footnote 2).
\item
  Ukuran polisentrisitas menggunakan kemiringan distribusi rank-size
  untuk dua, tiga, dan empat \emph{incorporated places}, lalu
  merata-ratakannya. Kawasan dengan dua kota terbesar berada dalam
  \emph{urban area} yang sama dikeluarkan karena dapat lebih menyerupai
  kawasan monosentris (TXT lines 313-326, 382-393; printed PDF
  pp.~1388-89; Figure 2 and footnote 1).
\item
  Rasio modal-tenaga kerja merupakan rata-rata tertimbang rasio sektoral
  dan tidak menangkap perbedaan regional di dalam sektor, seperti
  intensitas modal, tingkat teknologi, dan spesialisasi subsektor.
  Sumber menyebutnya sebagai proksi yang wajar (TXT lines 545-552;
  printed PDF p.~1392).
\item
  Produktivitas tidak mencakup sektor pertanian, perikanan, perburuan,
  pertambangan, dan administrasi publik, serta tidak mencakup pekerja
  mandiri karena angkanya tidak tersedia (TXT lines 503-509; printed PDF
  p.~1391).
\item
  Cakupan sampel dibatasi oleh ketersediaan data kawasan mikropolitan,
  pengecualian New Orleans karena bias data akibat Katrina, dan
  pengecualian \emph{polycentric conurbations} (TXT lines 1124-1149;
  printed PDF p.~1402; Appendix).
\end{itemize}

\textbf{Inferensi pengolahan:} Karena pengukuran hasil dan struktur
utama berpusat pada 2006, file ini tidak menyediakan estimasi perubahan
produktivitas atau perubahan efek struktur sepanjang waktu. Keterbatasan
temporal ini merupakan inferensi dari cakupan data, bukan pernyataan
eksplisit penulis; agenda riset penulis sendiri mengusulkan pengujian
perubahan kesenjangan dari waktu ke waktu (TXT lines 492-502, 503-518,
952-956; printed PDF pp.~1391, 1399).

\subsection{Challenges}\label{challenges-56}

\textbf{Temuan sumber:} Tantangan empiris utama adalah
mengoperasionalkan gagasan bahwa eksternalitas aglomerasi dapat dibagi
di antara kota-kota yang berdekatan, sambil menghindari perspektif satu
inti yang melekat pada banyak definisi metropolitan. Penelitian juga
menghadapi ambiguitas kausal: aglomerasi dapat meningkatkan
produktivitas, tetapi perusahaan dan penduduk juga dapat tertarik ke
tempat yang produktivitasnya sudah tinggi; produktivitas yang tinggi dan
harga lahan dapat pula mendorong penyebaran struktur spasial (TXT lines
105-128, 572-590; printed PDF pp.~1384, 1392-93).

\textbf{Temuan sumber:} Tantangan metodologis lain meliputi pemenuhan
relevansi dan eksogenitas instrumen historis, pengukuran polisentrisitas
tanpa menyamakan keberadaan banyak kota dengan keterhubungan fungsional,
serta pembedaan dispersal dari sprawl berkepadatan rendah (TXT lines
237-261, 294-307, 600-630; printed PDF pp.~1386-89, 1392-93).

\subsection{Practical Implications}\label{practical-implications-56}

\textbf{Interpretasi penulis:} Analisis dan perdebatan kebijakan tentang
skala aglomerasi perlu memperhitungkan struktur metropolitan regional,
bukan hanya ukuran, kepadatan, atau satu kota inti. Hasil artikel
mendukung perhatian pada kawasan polisentris sebagai bentuk yang dapat
mempertahankan produktivitas dengan disekonomi yang relatif terbatas,
tetapi ukuran kawasan tetap penting karena manfaat polisentrisitas dan
eksternalitas urbanisasi tidak setara dengan manfaat satu kota besar
(TXT lines 145-161, 215-236, 865-927; printed PDF pp.~1385-88, 1397-98).

\textbf{Temuan sumber:} Artikel menempatkan pertanyaan tersebut dalam
debat kebijakan mengenai megaregion di Amerika Serikat dan Asia serta
kawasan urban polisentris di Eropa; artikel tidak menyatakan bahwa
hasilnya merupakan evaluasi atas satu kebijakan tertentu (TXT lines
117-125, 145-150; printed PDF pp.~1384-85).

\subsection{Applications}\label{applications-56}

\textbf{Temuan sumber:} Aplikasi implementatif atau evaluasi kebijakan
spesifik tidak disebutkan dalam file. Sumber menyebut \emph{new
urbanism}, \emph{smart growth}, dan debat perencanaan sebagai konteks
pembahasan struktur spasial, tetapi penelitian yang dilaporkan
mengestimasi produktivitas tenaga kerja dan bukan menguji pelaksanaan
program tertentu (TXT lines 145-160, 237-261; printed PDF pp.~1385-87).

\subsection{Future Research
(Optional)}\label{future-research-optional-56}

\textbf{Temuan sumber:} Pada bagian ``Research agenda'', penulis
mengusulkan penelitian tentang:

\begin{itemize}
\tightlist
\item
  hubungan antara kekuatan dan jangkauan hubungan fungsional antarkota
  dengan kinerja sistem urban regional;
\item
  hubungan eksternal antarkawasan metropolitan yang lebih jauh dan
  kemungkinan eksternalitas aglomerasi dibagi di antara kota-kota yang
  lebih jauh;
\item
  apakah kesenjangan eksternalitas urbanisasi antara kawasan lebih
  polisentris dan lebih monosentris meningkat atau menurun dari waktu ke
  waktu;
\item
  apakah disekonomi aglomerasi yang lazim memang kurang parah di kawasan
  polisentris; dan
\item
  apakah efek struktur spasial terhadap tingkat produktivitas berbeda
  antar-sektor (TXT lines 941-956; printed PDF p.~1399; section 6.2).
\end{itemize}

\textbf{Temuan sumber:} Penulis juga menyatakan bahwa alasan mengapa
polisentrisitas mungkin lebih menguntungkan kawasan kecil perlu diuji
(TXT lines 928-940; printed PDF pp.~1398-99).

\subsection{Provenance Notes}\label{provenance-notes-56}

\begin{itemize}
\tightlist
\item
  \textbf{Kode dan filename aktual:} B13 dicocokkan dengan satu file TXT
  aktual:
  \texttt{source/journal/B13-meijers-burger-2010-spatial-structure-and-productivity-in-us-metropolitan-areas-10.1068-a42151.txt}.
  File merujuk pada PDF
  \texttt{journal/B13-meijers-burger-2010-spatial-structure-and-productivity-in-us-metropolitan-areas-10.1068-a42151.pdf}
  (TXT lines 1-3).
\item
  \textbf{Identitas yang tersedia dalam TXT:} ``Spatial structure and
  productivity in US metropolitan areas'', Evert J Meijers dan Martijn J
  Burger, \emph{Environment and Planning A} 2010, volume 42, pages
  1383-1402, DOI \texttt{10.1068/a42151} (TXT lines 27-43).
\item
  \textbf{Metadata ekstraksi yang dibaca:} TXT menyatakan diekstrak pada
  2026-08-31 dengan \texttt{pdftotext\ -layout}; PDF memiliki 21
  halaman, ukuran 420.8 x 686 pt, ukuran file 365323 bytes, versi PDF
  1.4, dan tercatat terenkripsi dengan izin cetak/salin tertentu
  sebagaimana blok metadata (TXT lines 1-25). Metadata tersebut
  diperlakukan sebagai provenance ekstraksi, bukan sebagai temuan
  substantif artikel.
\item
  \textbf{Cakupan pembacaan:} Seluruh TXT dibaca, termasuk blok
  \texttt{{[}PDF\ Metadata{]}}, teks artikel, referensi, lampiran, dan
  conditions of use sampai TXT line 1158. Locator
  \texttt{TXT\ lines\ ...} merujuk pada nomor baris hasil pembacaan file
  ini; locator \texttt{printed\ PDF\ p.\ ...} merujuk pada nomor halaman
  yang tercetak di dalam hasil ekstraksi.
\item
  \textbf{Catatan integritas sumber:} Ekstraksi memuat artefak karakter
  dan pemenggalan teks PDF. Selain itu, Appendix menyatakan ada dua CSA
  dengan lebih dari 20\% penduduk di kawasan mikropolitan, tetapi
  mencantumkan tiga nama kawasan (TXT lines 1128-1131). Ketidaksesuaian
  tersebut tidak ditafsirkan atau diperbaiki. Isi substantif tetap dapat
  dibaca dan direview; angka atau klasifikasi yang ambigu tidak diisi
  dengan dugaan.
\item
  \textbf{Ambiguitas internal hasil ekonometrik:} Pada TXT lines
  721-724, sumber menyatakan uji Durbin-Wu-Hausman untuk spesifikasi
  yang menginstrumentasikan ukuran metropolitan menolak hipotesis pada p
  = 0,021, lalu memuat kalimat tentang konsistensi OLS yang tidak
  sepenuhnya selaras dengan pembacaan uji eksogenitas. Review ini
  mempertahankan hasil uji dan keputusan estimator yang dilaporkan tanpa
  menyelesaikan ambiguitas kalimat tersebut. Pada TXT lines 799-806,
  prosa menyebut model interaksi 3 dan 4, sedangkan heading Table 5
  menampilkan Model 3 OLS dan Model 5 TSLS; perbedaan label ini juga
  tidak diperbaiki.
\item
  \textbf{Batas evidensi:} Review ini hanya menggunakan TXT B13 dan
  tidak menambahkan informasi dari referensi yang dikutip, PDF lain,
  atau sumber eksternal. Pernyataan bertanda \textbf{Temuan sumber}
  merangkum isi yang dilaporkan atau disediakan artikel;
  \textbf{Interpretasi penulis} mempertahankan penafsiran artikel;
  \textbf{Inferensi pengolahan} adalah klasifikasi atau sintesis review
  yang diberi label.
\end{itemize}

\section{Review B14}\label{review-b14}

\textbf{Review:} B14

\textbf{Sumber yang direview:}
\texttt{source/journal/B14-li-sun-zhang-zhang-2022-polycentric-spatial-structure-and-its-economic-performance-evidence-from-meta-analysis-10.1080-00343404.2021.2012142.txt}

\textbf{Penanda epistemik:} \textbf{Laporan sumber} berarti isi yang
dinyatakan atau dihitung oleh artikel; \textbf{Interpretasi penulis}
berarti penjelasan atau posisi analitis yang dibangun penulis artikel;
\textbf{Inferensi pengolahan} berarti pembacaan yang ditarik dari
struktur teks dan hasil pengolahan review ini, bukan klaim eksplisit
artikel.

\subsection{Summarized Abstract}\label{summarized-abstract-57}

\begin{itemize}
\tightlist
\item
  \textbf{Laporan sumber:} Artikel menyatakan bahwa harapan umum
  mengenai polisentrisitas sebagai pendorong daya saing ekonomi belum
  didukung secara konsisten oleh temuan empiris. Sebagai meta-analisis
  pertama mengenai hubungan polisentrisitas dan kinerja ekonomi, artikel
  menggabungkan 139 estimasi dan menemukan tidak ada bukti meyakinkan
  bahwa polisentrisitas secara umum menguntungkan atau merugikan
  pembangunan ekonomi (TXT lines 89-96; PDF p.~1888).
\item
  \textbf{Laporan sumber:} Meta-regresi logit menunjukkan bahwa cara
  mengukur polisentrisitas dan apakah ekonomi aglomerasi dikendalikan
  memengaruhi efek polisentrisitas-produktivitas. Sebaliknya, pengukuran
  polisentrisitas dari perspektif yang berbeda serta perlakuan terhadap
  potensi endogenitas tidak tampak berpengaruh signifikan terhadap efek
  yang diestimasi (TXT lines 89-96; PDF p.~1888).
\item
  \textbf{Interpretasi penulis:} Artikel memandang meta-analisis sebagai
  cara yang objektif untuk merangkum hasil yang bertentangan, menyaring
  kemungkinan bias seleksi publikasi, dan menjelaskan mengapa hasil
  studi primer bervariasi (TXT lines 148-187; PDF pp.~1889-90).
\item
  \textbf{Inferensi pengolahan:} Unit bukti artikel ini adalah estimasi
  dari studi terdahulu, bukan pengamatan primer baru atas kota atau
  wilayah. Karena itu, kesimpulan agregatnya tidak dengan sendirinya
  merupakan estimasi kausal untuk satu wilayah tertentu (TXT lines
  280-349, 449-498; PDF pp.~1890-94).
\end{itemize}

\subsection{Problem Statement}\label{problem-statement-57}

\begin{itemize}
\tightlist
\item
  \textbf{Laporan sumber:} Polisentrisitas telah menjadi pola spasial
  yang banyak didukung dalam perencanaan, tetapi bukti tentang dampaknya
  terhadap kinerja ekonomi tidak konklusif. Sebagian studi melaporkan
  pengaruh yang kecil atau tidak ada, sedangkan studi lain menemukan
  arah yang bertentangan: baik polisentrisitas maupun monosentrisitas
  dapat dikaitkan dengan kinerja yang lebih baik (TXT lines 104-124; PDF
  p.~1888).
\item
  \textbf{Laporan sumber:} Artikel menyebut beberapa kemungkinan
  penyebab ketidakkonsistenan tersebut, yaitu perbedaan definisi
  konseptual dan skala geografis, pilihan pengendalian ekonomi
  aglomerasi, serta keputusan untuk menangani atau mengabaikan
  endogenitas. Bentuk data, lokasi dan ukuran sampel, serta
  karakteristik publikasi juga dapat memengaruhi variasi hasil (TXT
  lines 107-124, 230-287; PDF pp.~1888-90).
\item
  \textbf{Interpretasi penulis:} Masalah penelitian dirumuskan bukan
  hanya sebagai pertanyaan apakah polisentrisitas berpengaruh, melainkan
  juga mengapa studi yang menilai hubungan polisentrisitas-produktivitas
  menghasilkan kesimpulan yang berbeda walaupun sebagian studi memakai
  skala dan definisi yang serupa (TXT lines 148-187; PDF p.~1889).
\item
  \textbf{Inferensi pengolahan:} Pertanyaan utama B14 adalah sintesis
  lintas-studi dan pengujian moderator spesifikasi. File tidak
  menunjukkan bahwa artikel menguji suatu intervensi kebijakan
  polisentrisitas secara langsung (TXT lines 280-349, 449-498; PDF
  pp.~1890-94).
\end{itemize}

\subsection{Objectives}\label{objectives-57}

\begin{itemize}
\tightlist
\item
  \textbf{Laporan sumber:} Tujuan artikel adalah memperjelas hubungan
  antara polisentrisitas dan kinerja ekonomi melalui meta-analisis
  terhadap bukti empiris yang sudah tersedia (TXT lines 148-161,
  774-781; PDF pp.~1889, 1899).
\item
  \textbf{Laporan sumber:} Penulis berupaya memberikan ringkasan
  objektif dan komprehensif tentang efek polisentrisitas-produktivitas
  serta menyaring kemungkinan bias seleksi publikasi. Artikel juga
  mengumpulkan 139 estimasi dari 31 studi untuk menilai mengapa efek
  yang dilaporkan bervariasi (TXT lines 161-187; PDF p.~1889).
\item
  \textbf{Laporan sumber:} Tiga penjelasan utama yang diuji adalah
  pengukuran polisentrisitas, keberadaan kontrol untuk ekonomi
  aglomerasi, dan perlakuan terhadap potensi endogenitas. Karakteristik
  bentuk dan periode data, ukuran serta lokasi sampel, dan ciri
  publikasi dimasukkan sebagai moderator tambahan (TXT lines 230-287,
  618-645; PDF pp.~1890, 1897).
\item
  \textbf{Interpretasi penulis:} Tujuan tersebut dimaksudkan untuk
  mengubah perdebatan yang terpecah antarstudi menjadi penilaian
  kuantitatif yang dapat menunjukkan apakah ketidakkonsistenan berasal
  dari perbedaan spesifikasi atau dari efek ekonomi yang memang saling
  mengimbangi (TXT lines 148-187, 496-500; PDF pp.~1889, 1894).
\end{itemize}

\subsection{Literature Survey}\label{literature-survey-57}

\begin{itemize}
\tightlist
\item
  \textbf{Laporan sumber:} Literatur yang diringkas memberi dua alasan
  teoritis mengapa wilayah polisentris diperkirakan berkinerja lebih
  baik. Kota-kota di dalamnya cenderung lebih kecil sehingga dianggap
  lebih sedikit mengalami kemacetan, harga hunian tinggi, polusi, dan
  kriminalitas; selain itu, kota-kota tersebut dapat meminjam ukuran
  atau fungsi satu sama lain melalui jaringan keterkaitan (TXT lines
  194-221; PDF p.~1889).
\item
  \textbf{Laporan sumber:} Argumen tandingannya adalah bahwa struktur
  monesentris dapat menghasilkan ekonomi aglomerasi yang lebih kuat
  melalui mekanisme \emph{sharing}, \emph{matching}, dan
  \emph{learning}. Sekumpulan kota yang berdekatan juga belum tentu
  dapat menggantikan ekonomi aglomerasi dari satu kota besar (TXT lines
  208-229; PDF pp.~1889-90).
\item
  \textbf{Laporan sumber:} Tinjauan membedakan pengukuran
  polisentrisitas secara morfologis dan fungsional. Artikel
  mengelompokkan indikator menjadi indikator berbasis proporsi, berbasis
  rank-size, berbasis simpangan baku, dan indikator lain seperti jumlah
  pusat; pengukuran juga dapat memakai jumlah pusat yang tetap atau
  seluruh pusat. Polisentrisitas dipandang bergantung pada skala mikro,
  meso, atau makro (TXT lines 230-263; PDF p.~1890).
\item
  \textbf{Laporan sumber:} Ekonomi aglomerasi diposisikan sebagai
  variabel kontrol penting karena pola polisentris dapat muncul ketika
  disekonomi aglomerasi melampaui ekonominya, meskipun kegiatan ekonomi
  tetap cenderung terkonsentrasi di pusat-pusat agar memperoleh manfaat
  aglomerasi. Populasi atau kepadatan sering digunakan sebagai proksi
  (TXT lines 265-287; PDF p.~1890).
\item
  \textbf{Laporan sumber:} Endogenitas dibahas terutama sebagai
  kemungkinan kausalitas terbalik: produktivitas yang lebih tinggi dapat
  menarik perusahaan dan pekerja ke lokasi tertentu atau mendorong
  penyebaran ke kota yang lebih kecil dengan harga lahan dan properti
  lebih rendah. Instrumen historis dianggap relevan tetapi mungkin
  kurang eksogen, sedangkan instrumen geologis dianggap eksogen tetapi
  dapat kurang kuat (TXT lines 233-259; PDF p.~1890).
\item
  \textbf{Interpretasi penulis:} Ketiga perbedaan spesifikasi tersebut
  mencakup pengukuran variabel utama, pemilihan kontrol, dan strategi
  estimasi, sehingga dipilih sebagai penjelasan utama atas hasil yang
  berkonflik. Literatur artikel juga menempatkan bentuk data, periode,
  ukuran dan lokasi sampel, serta fitur publikasi sebagai sumber variasi
  tambahan (TXT lines 260-287; PDF p.~1890).
\end{itemize}

\subsection{Method Used}\label{method-used-57}

\begin{itemize}
\tightlist
\item
  \textbf{Laporan sumber:} Artikel menggunakan meta-analisis untuk
  menggabungkan temuan studi empiris terdahulu. Studi diidentifikasi
  melalui daftar referensi dan sitasi studi empiris yang dianggap
  penting, kemudian melalui Web of Science, Google Scholar, dan CNKI
  dengan kombinasi kata kunci \texttt{polycentric*},
  \texttt{mono-centric*}, atau \texttt{spatial\ structure} dengan
  \texttt{economic\ performance}, \texttt{economic\ growth}, atau
  \texttt{growth} (TXT lines 280-299; PDF pp.~1890-91).
\item
  \textbf{Laporan sumber:} Studi berbahasa Inggris atau China
  dimasukkan. Penulis juga memasukkan beberapa estimasi dari satu studi
  yang sama, mengecualikan Lee dan Gordon (2011) karena hanya melaporkan
  suku interaksi polisentrisitas, serta tidak memasukkan studi
  konsentrasi perkotaan yang menggunakan bentuk linear dan kuadrat untuk
  menangkap efek nonlinier (TXT lines 295-313; PDF p.~1891).
\item
  \textbf{Laporan sumber:} CiteSpace digunakan untuk analisis tematik
  atas kata kunci 31 studi terpilih. Jaringan tersebut memiliki 76 node
  dan 199 link; lima studi yang tidak memiliki kata kunci diproses
  menggunakan judul dan abstrak. Pengaturan yang disebutkan adalah
  irisan waktu satu tahun, tipe node kata kunci, dan kriteria seleksi
  g-index dengan k = 25 (TXT lines 312-320, 368-373, 805-824; PDF
  pp.~1891-92, 1900).
\item
  \textbf{Laporan sumber:} Tiga jenis ukuran efek diekstraksi, yaitu
  elastisitas, statistik-t, dan hasil kategorikal. Elastisitas hanya
  tersedia untuk model log-log sehingga terkumpul 98 elastisitas dari 20
  studi. Hasil kategorikal mengelompokkan temuan menjadi signifikan
  negatif, tidak signifikan, dan signifikan positif agar lebih banyak
  studi dapat dimasukkan (TXT lines 322-349; PDF p.~1891).
\item
  \textbf{Laporan sumber:} Analisis dilakukan melalui \emph{vote count},
  funnel plot, dan \emph{precision-effect test} untuk menilai pola efek
  serta bias seleksi publikasi. Selanjutnya, meta-regresi multinomial
  logit menggunakan hasil kategorikal sebagai variabel yang dijelaskan
  dan memasukkan moderator pengukuran, skala, kontrol ekonomi
  aglomerasi, endogenitas, karakteristik data, sampel, serta publikasi
  (TXT lines 353-397, 405-498, 542-592; PDF pp.~1891-94, 1897).
\item
  \textbf{Laporan sumber:} Meta-logit diestimasi dalam enam spesifikasi:
  tiga kelompok karakteristik utama, seluruh moderator, regresi
  \emph{backward stepwise}, dan subsampel berbahasa Inggris. Tabel 5
  melaporkan 139 observasi dari 31 studi pada lima spesifikasi pertama
  dan 80 observasi dari 17 studi pada subsampel Inggris (TXT lines
  646-767; PDF pp.~1897-99).
\item
  \textbf{Inferensi pengolahan:} B14 adalah sintesis kuantitatif
  lintas-studi dengan beberapa tingkat analisis, bukan satu model
  regresi atas populasi kota. Perbedaan bentuk ukuran efek dan
  pengelompokan hasil berarti angka-angka pada bagian berikut tidak
  seluruhnya berasal dari sampel yang sama (TXT lines 322-349, 377-397,
  432-434; PDF pp.~1891-93).
\end{itemize}

\subsection{Dataset}\label{dataset-57}

\begin{itemize}
\tightlist
\item
  \textbf{Laporan sumber:} Dataset meta-analisis terdiri atas 139
  estimasi dari 31 studi primer yang dipilih melalui penelusuran
  referensi, sitasi, Web of Science, Google Scholar, dan CNKI. Artikel
  menyebut kinerja ekonomi secara luas, meliputi produktivitas tenaga
  kerja, GDP per kapita, produktivitas faktor total, pertumbuhan
  penduduk atau pekerjaan, serta ukuran seperti jumlah toko ritel,
  universitas, apotek, atau tempat tidur rumah sakit (TXT lines 280-320;
  PDF pp.~1890-91).
\item
  \textbf{Laporan sumber:} Untuk setiap temuan, penulis mengodekan
  elastisitas, statistik-t, atau kategori hasil. Kategori hasil adalah
  signifikan negatif, tidak signifikan, dan signifikan positif. Dari
  keseluruhan korpus, 98 elastisitas berasal dari 20 studi karena
  elastisitas hanya dapat dikumpulkan dari model log-log (TXT lines
  322-349; PDF p.~1891).
\item
  \textbf{Laporan sumber:} Moderator mencakup cara ukur polisentrisitas,
  metode indikator, jumlah pusat yang digunakan, skala dan konteks
  geografis, kontrol ekonomi aglomerasi, penanganan endogenitas, bentuk
  panel atau penampang lintang, tahun data, ukuran sampel studi primer,
  status publikasi, tahun publikasi, bahasa, dan jenis variabel kinerja
  yang dijelaskan (TXT lines 542-616; PDF p.~1897).
\item
  \textbf{Laporan sumber:} Cakupan geografis moderator dibagi menjadi
  Asia, Eropa, dan Amerika Utara, sedangkan skala dibedakan menjadi
  mikro, meso, dan makro. Nilai rata-rata yang dilaporkan antara lain
  proporsi studi panel 0,280, tahun rata-rata data 2006,151, ukuran
  sampel studi primer 372,381, proporsi estimasi dari jurnal akademik
  0,726, dan proporsi studi berbahasa Inggris 0,580 (TXT lines 604-616;
  PDF p.~1897).
\item
  \textbf{Laporan sumber:} Pada variabel spesifikasi, 87\% estimasi
  menggunakan pengukuran polisentrisitas morfologis. Proporsi estimasi
  yang mengendalikan ekonomi aglomerasi adalah 0,806, sedangkan yang
  menangani potensi endogenitas adalah 0,360 (TXT lines 637-645,
  553-591; PDF p.~1897).
\item
  \textbf{Inferensi pengolahan:} TXT tidak menyediakan baris data
  mentah, berkas kode, atau rekaman ekstraksi untuk setiap estimasi.
  Karena itu, review ini dapat menilai definisi dan angka agregat yang
  dilaporkan, tetapi tidak dapat merekonstruksi seluruh dataset primer
  dari file tersebut (TXT lines 280-349, 542-616; PDF pp.~1890-91,
  1897).
\end{itemize}

\subsection{Population / Sample}\label{population-sample-57}

\begin{itemize}
\tightlist
\item
  \textbf{Laporan sumber:} Populasi substantif yang dituju adalah studi
  empiris terdahulu tentang struktur spasial polisentris dan kinerja
  ekonomi. Sampel yang dianalisis adalah 31 studi dengan 139 estimasi;
  daftar studi yang dikodekan disajikan dalam Table 1 (TXT lines
  295-299, 322-357; PDF p.~1891).
\item
  \textbf{Laporan sumber:} Studi berbahasa Inggris dan China dimasukkan
  karena bidangnya masih kecil dan berkembang. Artikel menerapkan
  \emph{multiple sampling}, sehingga satu studi dapat menyumbang
  beberapa estimasi. Jumlah studi juga dapat dihitung lebih dari sekali
  dalam kategori hasil bila studi yang sama melaporkan hasil campuran
  (TXT lines 300-313, 395-397; PDF pp.~1891-92).
\item
  \textbf{Laporan sumber:} Artikel mengecualikan Lee dan Gordon (2011)
  karena hanya melaporkan suku interaksi polisentrisitas, serta
  mengecualikan studi konsentrasi perkotaan yang memakai variabel linear
  dan kuadrat untuk menangkap nonlinieritas. Prosedur sampling
  probabilistik: \textbf{Tidak disebutkan dalam file} (TXT lines
  302-313; PDF p.~1891).
\item
  \textbf{Laporan sumber:} Untuk analisis elastisitas, sampelnya
  menyusut menjadi 98 estimasi dari 20 studi. Meta-logit subsampel
  Inggris menggunakan 80 observasi dari 17 studi; model dengan lima
  spesifikasi lainnya menggunakan 139 observasi dari 31 studi (TXT lines
  322-349, 761-767; PDF pp.~1891, 1899).
\item
  \textbf{Laporan sumber:} Studi primer ditempatkan dalam konteks Asia,
  Eropa, atau Amerika Utara dan pada skala mikro, meso, atau makro. Pada
  klasifikasi yang digunakan, Table 4 mencantumkan tidak adanya sampel
  Amerika Utara pada skala mikro dan makro (TXT lines 479-492, 567-586;
  PDF pp.~1894, 1897).
\item
  \textbf{Inferensi pengolahan:} Sampel B14 bukan sensus atas semua
  studi atau wilayah yang mungkin relevan. Ketergantungan pada studi
  yang teridentifikasi melalui strategi pencarian artikel dan pada
  beberapa estimasi dari studi yang sama membatasi generalisasi
  universal dan membuat unit studi berbeda dari unit estimasi (TXT lines
  280-313, 395-397; PDF p.~1891-92).
\end{itemize}

\subsection{Dependent Variables}\label{dependent-variables-57}

\begin{itemize}
\tightlist
\item
  \textbf{Laporan sumber:} Pada studi primer, kinerja ekonomi dipakai
  dalam arti umum. Outcome yang dicantumkan artikel mencakup
  produktivitas tenaga kerja, GDP per kapita, produktivitas faktor
  total, pertumbuhan penduduk atau pekerjaan, dan beberapa ukuran fungsi
  ekonomi atau layanan seperti toko ritel, universitas, apotek, dan
  tempat tidur rumah sakit (TXT lines 314-320; PDF p.~1891).
\item
  \textbf{Laporan sumber:} Pada meta-regresi multinomial logit, variabel
  yang dijelaskan adalah \texttt{Poly\_result}: nilai 1 untuk hasil
  polisentrisitas yang signifikan negatif, nilai 2 untuk hasil tidak
  signifikan sebagai kategori dasar, dan nilai 3 untuk hasil signifikan
  positif (TXT lines 449-465, 542-552; PDF pp.~1894, 1897).
\item
  \textbf{Laporan sumber:} Pada \emph{precision-effect test},
  statistik-t dari elastisitas polisentrisitas-produktivitas, yaitu
  \texttt{Tij}, menjadi variabel yang dijelaskan dan dibandingkan dengan
  kebalikan standar deviasinya, \texttt{1/SEij}. Elastisitas sendiri
  berfungsi sebagai ukuran efek ekonomi pada subset model log-log, bukan
  sebagai satu variabel dependen yang seragam untuk seluruh 139 estimasi
  (TXT lines 405-424, 458-465; PDF pp.~1893-94).
\item
  \textbf{Laporan sumber:} Moderator \texttt{Control\_y} membedakan
  studi yang variabel hasilnya berupa GDP per kapita atau GDP per tenaga
  kerja dari ukuran hasil lainnya. File tidak menetapkan satu outcome
  substantif yang sama untuk seluruh studi primer; satu variabel
  dependen substantif yang seragam adalah \textbf{Tidak berlaku} (TXT
  lines 604-606, 618-625; PDF p.~1897).
\item
  \textbf{Inferensi pengolahan:} Istilah variabel dependen dalam B14
  perlu dibaca pada dua tingkat: outcome kinerja ekonomi pada studi
  primer dan kategori hasil yang menjadi outcome meta-logit. Keduanya
  tidak boleh diperlakukan sebagai pengukuran langsung yang sama atas
  produktivitas wilayah (TXT lines 314-320, 542-552; PDF pp.~1891,
  1897).
\end{itemize}

\subsection{Results}\label{results-57}

\begin{itemize}
\tightlist
\item
  \textbf{Laporan sumber:} \emph{Vote count} atas 139 estimasi menemukan
  64 hasil signifikan negatif dari 16 studi, 31 hasil tidak signifikan
  dari 13 studi, dan 44 hasil signifikan positif dari 14 studi. Jumlah
  studi dapat terhitung berulang karena satu studi boleh melaporkan
  hasil campuran (TXT lines 353-397; PDF pp.~1891-92).
\item
  \textbf{Laporan sumber:} Hasil keseluruhan menunjukkan kecenderungan
  negatif yang lemah, tetapi tidak ada kesepakatan definitif bahwa
  polisentrisitas efektif untuk mendorong pertumbuhan ekonomi. Dalam
  hitungan estimasi, hasil signifikan negatif mencakup sekitar 78\%
  pengukuran fungsional, 71\% pengukuran berbasis proporsi, dan 61\%
  pengukuran pada skala meso (TXT lines 405-424; PDF p.~1893).
\item
  \textbf{Laporan sumber:} Funnel plot untuk 98 estimasi elastisitas
  dari 20 studi memperlihatkan banyak observasi berkumpul di sekitar nol
  dan sebagian besar elastisitas berada antara -1 dan 1. Elastisitas
  ekonomi rata-rata adalah -0,065, sedangkan estimasi di bagian atas
  funnel yang dianggap paling presisi adalah -0,042. Polanya menunjukkan
  kecenderungan negatif lemah, tanpa tanda kuat bahwa polisentrisitas
  secara umum mendorong atau menghambat pembangunan ekonomi (TXT lines
  405-447; PDF pp.~1893-94).
\item
  \textbf{Laporan sumber:} \emph{Precision-effect test} untuk 98
  estimasi elastisitas menghasilkan koefisien kemiringan -0,043 dengan
  statistik-t -1,59 dan intersep 0,203 dengan statistik-t 0,16. Keduanya
  tidak signifikan, sehingga penulis melaporkan tidak ada bukti
  statistik untuk efek empiris umum setelah koreksi potensial distorsi
  dan tidak ada bukti bias publikasi pada keseluruhan studi yang diuji
  dalam test tersebut (TXT lines 449-498, 503-529; PDF pp.~1894-95).
\item
  \textbf{Laporan sumber:} Beberapa subkelompok menunjukkan sinyal yang
  lebih lemah dan tidak stabil. Tabel 3 mencatat kemiringan -0,061 pada
  skala mikro dan -0,076 pada estimasi yang tidak menangani endogenitas
  dengan tanda signifikansi; narasi artikel juga menyebut kemungkinan
  bias publikasi pada skala mikro dan meso. Penulis meminta validasi
  lebih lanjut karena ukuran subkelompok terbatas (TXT lines 482-498,
  511-529; PDF pp.~1894-95).
\item
  \textbf{Laporan sumber:} Meta-logit menunjukkan pseudo-R2 sebesar 0,40
  untuk tiga karakteristik spesifikasi utama, 0,20 untuk karakteristik
  studi-invarian, 0,12 untuk karakteristik publikasi, 0,70 untuk seluruh
  moderator, 0,58 untuk \emph{stepwise}, dan 0,71 untuk subsampel
  Inggris. Artikel merangkum bahwa kelompok karakteristik tersebut
  masing-masing menjelaskan sekitar 40\%, 20\%, dan 12\% variasi hasil
  dalam spesifikasinya (TXT lines 618-645, 761-767; PDF pp.~1897, 1899).
\item
  \textbf{Laporan sumber:} Metode pengukuran polisentrisitas dan
  pengendalian ekonomi aglomerasi memengaruhi peluang hasil yang
  dilaporkan. Dibandingkan indikator rank-size, indikator berbasis
  simpangan baku lebih mungkin menghasilkan hasil signifikan, baik
  positif maupun negatif. Pengukuran pada konteks Asia mikro dan meso
  juga menunjukkan kecenderungan tertentu dalam meta-logit (TXT lines
  626-645; PDF p.~1897).
\item
  \textbf{Laporan sumber:} Pengukuran morfologis versus perspektif lain
  dan perlakuan terhadap potensi endogenitas tidak menunjukkan pengaruh
  yang terlihat terhadap kinerja ekonomi yang diestimasi. Penulis
  mengingatkan bahwa 87\% estimasi berbasis polisentrisitas morfologis,
  sehingga ketidakberartian moderator morfologi dapat berkaitan dengan
  keterbatasan sampel fungsional (TXT lines 637-645; PDF p.~1897).
\item
  \textbf{Laporan sumber:} Studi dengan data panel atau data yang lebih
  baru dan studi yang dipublikasikan lebih awal memiliki probabilitas
  lebih tinggi untuk menemukan hasil signifikan. Ukuran sampel yang
  besar meningkatkan odds hasil signifikan negatif, sedangkan publikasi
  yang lebih awal meningkatkan odds hasil signifikan positif.
  \emph{Precision-effect test} tidak konvergen ketika hanya estimasi
  tanpa kontrol ekonomi aglomerasi yang digunakan (TXT lines 646-659,
  693-706, 530-531; PDF pp.~1895, 1897).
\end{itemize}

\subsection{Findings}\label{findings-57}

\begin{itemize}
\tightlist
\item
  \textbf{Interpretasi penulis:} Tidak adanya bukti agregat yang
  meyakinkan dipahami sebagai hasil saling mengimbangi antara efek
  positif dan negatif, bukan sebagai bukti bahwa pengaruh
  polisentrisitas tidak pernah ada. Artikel tetap membuka kemungkinan
  efek yang signifikan pada kasus kota atau wilayah tertentu (TXT lines
  496-500, 774-824; PDF pp.~1894, 1899-1900).
\item
  \textbf{Interpretasi penulis:} Polisentrisitas tidak otomatis menjadi
  alat untuk meningkatkan pembangunan ekonomi. Temuan tersebut juga
  tidak membatalkan kegunaan polisentrisitas untuk tujuan lain, seperti
  mengurangi kesenjangan atau mendorong keberlanjutan lingkungan, tetapi
  tujuan-tujuan tersebut bukan outcome utama yang diuji dalam
  meta-analisis ini (TXT lines 800-824; PDF p.~1900).
\item
  \textbf{Interpretasi penulis:} Perbedaan cara mengukur struktur
  spasial dan keputusan untuk mengendalikan ekonomi aglomerasi merupakan
  penjelasan penting atas variasi temuan primer. Sebaliknya, endogenitas
  tidak tampak sebagai pembeda utama dalam kumpulan studi yang
  dianalisis (TXT lines 626-645, 834-847; PDF pp.~1897, 1900-01).
\item
  \textbf{Inferensi pengolahan:} Bukti terkuat B14 mendukung
  kehati-hatian terhadap klaim universal, bukan pilihan kebijakan
  tertentu. Hasilnya menjelaskan pola dalam literatur yang tersedia,
  tetapi tidak mengidentifikasi mekanisme lokal atau dampak kausal dari
  satu konfigurasi spasial terhadap satu wilayah (TXT lines 449-498,
  774-824; PDF pp.~1894, 1899-1900).
\end{itemize}

\subsection{Conclusions}\label{conclusions-57}

\begin{itemize}
\tightlist
\item
  \textbf{Laporan sumber:} Artikel menyimpulkan bahwa, dari 139 estimasi
  dalam 31 studi, tidak ada bukti meyakinkan bahwa polisentrisitas
  menguntungkan atau merugikan pembangunan ekonomi. Artikel juga
  melaporkan tidak adanya bukti bias seleksi publikasi secara
  keseluruhan (TXT lines 774-799; PDF pp.~1899-1900).
\item
  \textbf{Laporan sumber:} Tiga perbedaan spesifikasi berperan penting
  dalam menjelaskan hasil yang tidak konklusif. Metode pengukuran
  polisentrisitas dan penggunaan kontrol ekonomi aglomerasi memengaruhi
  efek polisentrisitas-produktivitas, sedangkan perspektif pengukuran
  yang berbeda dan penanganan potensi endogenitas tidak menunjukkan
  pengaruh yang terlihat (TXT lines 787-799; PDF p.~1900).
\item
  \textbf{Interpretasi penulis:} Kesimpulan nol perlu dibaca sebagai
  hasil agregat dari efek positif dan negatif yang saling mengimbangi.
  Polisentrisitas mungkin tetap relevan dalam kasus tertentu atau untuk
  tujuan selain pembangunan ekonomi, tetapi artikel tidak mendukung
  penggunaannya sebagai resep umum untuk meningkatkan kinerja ekonomi
  (TXT lines 800-824; PDF p.~1900).
\item
  \textbf{Inferensi pengolahan:} Kesimpulan B14 berlaku pada korpus
  studi yang berhasil diidentifikasi dan dikodekan, bukan pada semua
  kemungkinan bentuk polisentrisitas. Batas ini penting karena artikel
  sendiri menyebut bidang tersebut masih hanya memiliki sejumlah kecil
  studi empiris (TXT lines 815-824; PDF p.~1900).
\end{itemize}

\subsection{Contributions}\label{contributions-57}

\begin{itemize}
\tightlist
\item
  \textbf{Laporan sumber:} Artikel mengklaim sebagai meta-analisis
  pertama tentang efek struktur spasial polisentris terhadap kinerja
  ekonomi. Kontribusinya adalah menyediakan ringkasan objektif dan
  komprehensif, menyaring kemungkinan bias seleksi publikasi, serta
  memberi penilaian kuantitatif atas variasi hasil studi primer (TXT
  lines 148-187; PDF p.~1889).
\item
  \textbf{Laporan sumber:} Artikel mengumpulkan 139 estimasi dari 31
  studi dan memakai elastisitas, statistik-t, serta hasil kategorikal
  agar bukti yang tersedia dapat dianalisis dari sisi ekonomi dan
  statistik. Analisis juga menguji moderator pengukuran, kontrol,
  strategi estimasi, karakteristik data, sampel, serta publikasi (TXT
  lines 176-187, 322-349, 542-616; PDF pp.~1889, 1891, 1897).
\item
  \textbf{Laporan sumber:} Penggunaan \emph{vote count}, funnel plot,
  \emph{precision-effect test}, dan meta-logit memungkinkan artikel
  memisahkan pertanyaan tentang arah efek agregat dari pertanyaan
  tentang sumber variasi antarstudi. CiteSpace menambahkan pemetaan
  tematik terhadap korpus yang dipilih (TXT lines 312-320, 353-397,
  405-498; PDF pp.~1891-95).
\item
  \textbf{Interpretasi penulis:} Artikel menempatkan temuan
  meta-analisis sebagai dasar untuk menilai kembali keyakinan bahwa
  polisentrisitas selalu meningkatkan daya saing ekonomi dan untuk
  mengarahkan perhatian pada definisi, skala, serta kontrol ekonomi
  aglomerasi (TXT lines 774-847; PDF pp.~1899-1901).
\item
  \textbf{Inferensi pengolahan:} Kontribusi B14 terutama berupa sintesis
  dan diagnosis heterogenitas bukti. Artikel tidak menyumbang dataset
  kota primer atau estimasi kebijakan baru yang dapat menggantikan studi
  kontekstual (TXT lines 280-349, 542-616; PDF pp.~1890-91, 1897).
\end{itemize}

\subsection{Research Gap}\label{research-gap-57}

\begin{itemize}
\tightlist
\item
  \textbf{Inferensi pengolahan:} File tidak menyediakan heading khusus
  \texttt{Research\ Gap}. Bagian ini dirangkum dari pertanyaan
  penelitian, keterbatasan korpus, dan agenda yang dinyatakan langsung
  dalam pendahuluan serta kesimpulan artikel.
\item
  \textbf{Laporan sumber:} Artikel menyatakan bahwa penelitian
  sebelumnya menghasilkan temuan yang sangat tidak konklusif dan masih
  sedikit pekerjaan yang menggeneralisasi hasil lintas studi. Karena
  itu, hubungan polisentrisitas-kinerja ekonomi belum memiliki jawaban
  umum yang singkat dan jelas (TXT lines 104-124, 148-187, 774-807; PDF
  pp.~1888-89, 1899-1900).
\item
  \textbf{Laporan sumber:} Definisi dan identifikasi polisentrisitas
  yang diterima secara universal belum tercapai. Perbedaan pengukuran
  dan skala mengurangi keterbandingan antarstudi, sementara pengukuran
  fungsional masih memerlukan lebih banyak bukti empiris (TXT lines
  230-263, 637-645, 825-834; PDF pp.~1890, 1897, 1900).
\item
  \textbf{Laporan sumber:} Artikel menyebut perlunya lebih banyak studi
  empiris agar hasil meta menjadi lebih kredibel, terutama studi dari
  perspektif polisentrisitas fungsional. Artikel juga menegaskan bahwa
  pengaruh ekonomi dapat tetap signifikan dalam kasus-kasus tertentu
  meskipun efek keseluruhannya tidak meyakinkan (TXT lines 800-824; PDF
  p.~1900).
\item
  \textbf{Inferensi pengolahan:} Korpus yang 87\% berbasis pengukuran
  morfologis dan hanya menyediakan 98 elastisitas dari 20 studi
  menunjukkan ketimpangan bukti pada definisi fungsional serta ukuran
  efek yang bermakna secara ekonomi. Ini adalah inferensi dari
  distribusi data yang dilaporkan, bukan klaim tambahan dari sumber (TXT
  lines 322-349, 637-645; PDF pp.~1891, 1897).
\end{itemize}

\subsection{Limitations}\label{limitations-57}

\begin{itemize}
\tightlist
\item
  \textbf{Laporan sumber:} Meta-analisis ini dibangun dari bidang yang
  masih kecil: 139 estimasi dari 31 studi, dengan hanya 98 elastisitas
  dari 20 studi. Penulis menyatakan bahwa hasil meta akan lebih kredibel
  jika didukung oleh jumlah studi empiris yang lebih banyak (TXT lines
  176-187, 322-349, 815-824; PDF pp.~1889, 1891, 1900).
\item
  \textbf{Laporan sumber:} Ukuran efek tidak seragam. Elastisitas hanya
  tersedia untuk model log-log, statistik-t memiliki keterbatasan dalam
  makna ekonomi, dan hasil kategorikal hanya menunjukkan apakah hasil
  signifikan negatif, tidak signifikan, atau signifikan positif (TXT
  lines 322-349; PDF p.~1891).
\item
  \textbf{Laporan sumber:} Ukuran subkelompok untuk sebagian uji
  terbatas. Artikel meminta validasi lebih lanjut untuk perbedaan
  spesifikasi utama, mencatat kemungkinan bias publikasi pada skala
  mikro dan meso, serta melaporkan bahwa model tanpa kontrol ekonomi
  aglomerasi tidak konvergen (TXT lines 482-498, 515-531; PDF
  pp.~1894-95).
\item
  \textbf{Laporan sumber:} 87\% estimasi menggunakan polisentrisitas
  morfologis. Penulis menyatakan bahwa ketidakberartian pengukuran
  morfologis dapat sebagian disebabkan oleh keterbatasan sampel
  pengukuran fungsional (TXT lines 637-645; PDF p.~1897).
\item
  \textbf{Laporan sumber:} Kesimpulan tidak meniadakan pengaruh
  polisentrisitas pada kasus tertentu karena hasil keseluruhan merupakan
  agregasi efek positif dan negatif. Tujuan lain seperti pengurangan
  kesenjangan dan keberlanjutan lingkungan disebut, tetapi bukan outcome
  utama yang diuji (TXT lines 800-824; PDF p.~1900).
\item
  \textbf{Inferensi pengolahan:} \emph{Multiple sampling} membuat
  beberapa estimasi berasal dari studi yang sama, sehingga jumlah
  estimasi tidak sama dengan jumlah studi independen. TXT menyebut
  praktik tersebut, tetapi tidak menyajikan analisis sensitivitas yang
  memisahkan seluruh estimasi menurut studi dalam ringkasan hasilnya
  (TXT lines 295-313, 395-397; PDF p.~1891-92).
\item
  \textbf{Inferensi pengolahan:} TXT tidak melaporkan alur penyaringan
  lengkap, tanggal pencarian, penilaian kualitas setiap studi primer,
  atau dataset baris yang memungkinkan audit ulang seluruh keputusan
  inklusi dan pengodean. Ketiadaan ini membatasi reproduksibilitas
  review dari file yang tersedia (TXT lines 280-349, 542-616; PDF
  pp.~1890-91, 1897).
\end{itemize}

\subsection{Challenges}\label{challenges-57}

\begin{itemize}
\tightlist
\item
  \textbf{Laporan sumber:} Tantangan awal meta-analisis adalah menemukan
  studi yang relevan dalam topik yang masih baru. Penulis menggabungkan
  penelusuran referensi dan sitasi dengan tiga basis pencarian serta
  memasukkan publikasi berbahasa Inggris dan China untuk memperluas
  korpus (TXT lines 280-313; PDF pp.~1890-91).
\item
  \textbf{Laporan sumber:} Heterogenitas definisi, indikator, pusat yang
  dihitung, skala geografis, dan fungsi morfologis atau fungsional
  menyulitkan perbandingan hasil. Ketersediaan data interaksi yang
  terbatas membuat studi primer lebih sering memakai pengukuran
  morfologis (TXT lines 205-263, 637-645; PDF pp.~1889-90, 1897).
\item
  \textbf{Laporan sumber:} Pengodean ukuran efek juga menimbulkan
  tantangan. Elastisitas tidak tersedia untuk semua model, statistik-t
  kurang langsung bermakna secara ekonomi, dan kategori hasil diperlukan
  untuk mempertahankan studi yang tidak dapat menyumbang elastisitas
  (TXT lines 322-349; PDF p.~1891).
\item
  \textbf{Laporan sumber:} Endogenitas dan kausalitas terbalik menjadi
  masalah metodologis karena produktivitas dapat membentuk kembali pola
  spasial. Instrumen historis menghadapi pertanyaan tentang eksogenitas,
  sedangkan instrumen geologis dapat kurang kuat untuk menjelaskan
  variabel yang diinstrumentasi (TXT lines 233-259; PDF p.~1890).
\item
  \textbf{Laporan sumber:} Analisis harus menangani kemungkinan bias
  publikasi, observasi ekstrem, perbedaan periode dan bentuk data, serta
  ketergantungan pada beberapa estimasi dari satu studi. Keterbatasan
  ukuran subkelompok juga memengaruhi validasi moderator (TXT lines
  405-447, 482-498, 530-531; PDF pp.~1893-95).
\item
  \textbf{Inferensi pengolahan:} Keterangan bahwa regresi tanpa kontrol
  ekonomi aglomerasi tidak konvergen dan bahwa meta-logit memakai
  subsampel Inggris karena algoritme likelihood dapat gagal konvergen
  menunjukkan bahwa tidak semua spesifikasi yang substantif menarik
  dapat dibandingkan dengan basis estimasi yang sama (TXT lines 530-531,
  646-655, 761-767; PDF pp.~1895, 1897, 1899).
\end{itemize}

\subsection{Practical Implications}\label{practical-implications-57}

\begin{itemize}
\tightlist
\item
  \textbf{Laporan sumber:} Artikel menekankan bahwa polisentrisitas
  berkaitan erat dengan praktik perencanaan, sehingga ketidakkonsistenan
  bukti ekonomi perlu ditanggapi secara kritis. Hasil keseluruhan tidak
  mendukung anggapan bahwa polisentrisitas selalu merupakan alat untuk
  meningkatkan daya saing ekonomi (TXT lines 104-124, 774-824; PDF
  pp.~1888, 1899-1900).
\item
  \textbf{Interpretasi penulis:} Perencana dan pembuat kebijakan
  sebaiknya tidak mengadopsi polisentrisitas sebagai resep universal
  tanpa memperhatikan cara struktur tersebut diukur, skala analisis, dan
  ekonomi aglomerasi yang mendasarinya. Perbedaan faktor-faktor tersebut
  terbukti berkaitan dengan perbedaan hasil dalam literatur (TXT lines
  618-645, 825-840; PDF pp.~1897, 1900).
\item
  \textbf{Laporan sumber:} Artikel tetap mengakui kemungkinan manfaat
  polisentrisitas pada kasus tertentu dan menyebut tujuan lain seperti
  pengurangan kesenjangan atau keberlanjutan lingkungan sebagai alasan
  yang mungkin, tetapi tujuan tersebut tidak diuji sebagai outcome utama
  dalam meta-analisis (TXT lines 800-824; PDF p.~1900).
\item
  \textbf{Inferensi pengolahan:} Implikasi praktis yang aman dari file
  ini adalah kewajiban kehati-hatian dalam membaca bukti, bukan
  rekomendasi untuk memilih struktur monosentris atau polisentris. File
  tidak menyediakan evaluasi kebijakan yang dapat menetapkan dampak
  implementasi secara langsung (TXT lines 449-498, 774-824; PDF
  pp.~1894, 1899-1900).
\end{itemize}

\subsection{Applications}\label{applications-57}

\begin{itemize}
\tightlist
\item
  \textbf{Laporan sumber:} Tidak disebutkan dalam file adanya
  implementasi kebijakan, eksperimen perencanaan, atau evaluasi program
  yang memakai hasil meta-analisis ini pada suatu wilayah tertentu.
  Artikel membahas relevansinya bagi praktik perencanaan dan menyatakan
  bahwa temuan tersebut memberi panduan bagi penelitian berikutnya (TXT
  lines 774-824; PDF p.~1900).
\item
  \textbf{Inferensi pengolahan:} Kerangka pengodean B14 dapat diterapkan
  sebagai alat diagnosis untuk studi baru: catat definisi morfologis
  atau fungsional, indikator dan jumlah pusat, skala geografis, outcome
  kinerja, kontrol ekonomi aglomerasi, perlakuan endogenitas, bentuk
  data, dan karakteristik publikasi. Penerapan ini adalah inferensi dari
  moderator artikel, bukan aplikasi yang diuji atau dilaporkan oleh
  sumber (TXT lines 542-616, 618-767; PDF pp.~1897-99).
\item
  \textbf{Inferensi pengolahan:} Penerapan hasil B14 untuk memilih
  kebijakan wilayah tertentu: \textbf{Tidak berlaku} tanpa bukti lokal
  tambahan. File hanya memungkinkan perbandingan pola hasil
  lintas-studi, bukan penetapan bahwa satu wilayah akan mengalami efek
  yang sama (TXT lines 496-500, 774-824; PDF pp.~1894, 1899-1900).
\end{itemize}

\subsection{Future Research
(Optional)}\label{future-research-optional-57}

\begin{itemize}
\tightlist
\item
  \textbf{Laporan sumber:} Penulis menyerukan tercapainya definisi dan
  identifikasi polisentrisitas yang lebih dapat diterima secara
  universal agar keterbandingan studi meningkat. Mereka juga menyatakan
  bahwa pengaruh skala polisentrisitas perlu terus diperhatikan (TXT
  lines 825-834; PDF p.~1900).
\item
  \textbf{Laporan sumber:} Penelitian empiris tambahan diperlukan karena
  korpus yang tersedia masih terbatas, terutama penelitian dari
  perspektif polisentrisitas fungsional. Penulis menyebut bahwa
  ketidakberartian pengukuran morfologis mungkin berasal dari
  keterbatasan sampel, bukan penyelesaian umum atas pertanyaan tersebut
  (TXT lines 637-645, 815-824, 825-834; PDF pp.~1897, 1900).
\item
  \textbf{Laporan sumber:} Penelitian berikutnya perlu tetap
  mengendalikan ekonomi aglomerasi ketika menilai hubungan
  polisentrisitas-produktivitas. Penulis juga menyatakan bahwa hasil
  mereka tidak menunjukkan perlunya kekhawatiran berlebihan terhadap
  kausalitas terbalik atau bias endogenitas, meskipun hubungan tersebut
  tetap dibahas sebagai isu metodologis (TXT lines 834-847; PDF
  pp.~1900-01).
\end{itemize}

\subsection{Provenance Notes}\label{provenance-notes-57}

\begin{itemize}
\tightlist
\item
  \textbf{Laporan sumber:} Verifikasi prefiks B14 menemukan tepat satu
  TXT aktual di bawah \texttt{source/}, yaitu
  \texttt{source/journal/B14-li-sun-zhang-zhang-2022-polycentric-spatial-structure-and-its-economic-performance-evidence-from-meta-analysis-10.1080-00343404.2021.2012142.txt}.
  Judul artikel, penulis, jurnal \emph{Regional Studies}, volume 56
  nomor 11, halaman 1888-1902, dan DOI
  \texttt{10.1080/00343404.2021.2012142} tercantum di dalam TXT (TXT
  lines 27-46, 78-101).
\item
  \textbf{Laporan sumber:} Blok metadata menyatakan PDF sumber
  \texttt{journal/B14-li-sun-zhang-zhang-2022-polycentric-spatial-structure-and-its-economic-performance-evidence-from-meta-analysis-10.1080-00343404.2021.2012142.pdf},
  diekstraksi pada 2026-08-31 dengan \texttt{pdftotext\ -layout}.
  Metadata PDF melaporkan 16 halaman, PDF versi 1.4, tidak terenkripsi,
  tidak bertag, tanpa JavaScript, tanpa custom metadata, creator
  \texttt{Servigistics\ Arbortext\ Advanced\ Print\ Publisher\ 11.1.4546/W-x64},
  dan producer
  \texttt{Acrobat\ Distiller\ 8.1.0\ (Windows);\ modified\ using\ iTextSharp\ 5.5.8}
  (TXT lines 1-25).
\item
  \textbf{Inferensi pengolahan:} Seluruh 979 baris TXT dibaca, termasuk
  metadata ekstraksi, teks artikel, tabel dan caption yang berhasil
  diekstrak, catatan kaki, disclosure statement, funding, ORCID, serta
  daftar referensi. Tidak ada dokumen eksternal yang digunakan dan TXT
  sumber tidak diubah dalam pengerjaan ini (TXT lines 1-979).
\item
  \textbf{Inferensi pengolahan:} Locator \texttt{PDF\ p.} dalam review
  ini mengikuti nomor halaman tercetak dan penanda halaman yang muncul
  di TXT hasil ekstraksi; PDF tidak dibaca sebagai sumber tambahan
  secara terpisah. Baris tabel dua kolom, pemenggalan kata, dan karakter
  ligatur dapat menghasilkan artefak tata letak, sehingga angka dan
  hubungan yang penting dipertahankan dengan locator baris TXT (TXT
  lines 78-979).
\item
  \textbf{Inferensi pengolahan:} Teks tampak tersedia dari metadata dan
  halaman awal artikel sampai daftar referensi, tetapi file tidak
  menyediakan data mentah, kode analisis, atau visual lengkap yang dapat
  dipakai untuk merekonstruksi semua estimasi. Karena itu, keterbatasan
  reproduksibilitas dicatat dan tidak diisi dengan dugaan (TXT lines
  280-349, 542-616, 849-978; PDF pp.~1890-1902).
\item
  \textbf{Inferensi pengolahan:} Semua klaim review diberi label sebagai
  laporan sumber, interpretasi penulis, atau inferensi pengolahan.
  Daftar referensi yang tercantum dalam artikel diperlakukan hanya
  sebagai bagian dari teks B14 dan tidak ditelusuri atau diverifikasi
  secara eksternal (TXT lines 836-978).
\end{itemize}

\section{Review B15}\label{review-b15}

\textbf{Kode:} B15

\textbf{Sumber yang direview:}
\texttt{source/journal/B15-ouwehand-van-oort-cortinovis-2022-spatial-structure-and-productivity-in-european-regions-10.1080-00343404.2021.1950912.txt}

\subsection{Summarized Abstract}\label{summarized-abstract-58}

\textbf{Temuan sumber:} Artikel menganalisis total factor productivity
(TFP) pada kawasan regional Eropa dengan strategi identifikasi
ekonometrik. Ukuran urban yang lebih besar berpengaruh positif terhadap
produktivitas regional. Struktur urban polisentris tidak memiliki dampak
langsung yang teridentifikasi terhadap produktivitas. Interaksi antara
ukuran urban dan polisentrisitas berpengaruh negatif, yang menurut
sumber menunjukkan bahwa kawasan polisentris tidak dapat menggantikan
eksternalitas urbanisasi ekonomi yang terkait dengan satu kota besar.
Sumber menyatakan bahwa temuan ini berimplikasi pada agenda kebijakan
Uni Eropa mengenai pembangunan urban dan produktivitas regional (TXT
lines 80-91; printed PDF p.~48).

\textbf{Interpretasi penulis:} Bukti tersebut menempatkan ukuran urban
sebagai sumber keuntungan produktivitas yang lebih kuat daripada yang
dapat disediakan oleh struktur polisentris, terutama ketika
eksternalitas urbanisasi berukuran besar (TXT lines 83-86, 717-730;
printed PDF pp.~48, 57-59).

\textbf{Inferensi pengolahan:} Outcome utama review ini adalah TFP
regional karena TFP secara eksplisit digunakan sebagai dependent
variable pada model utama; hubungan yang dibahas adalah hubungan pada
tingkat region, bukan pada tingkat individu atau perusahaan (TXT lines
401-424, 476-496; printed PDF pp.~53-55).

\subsection{Problem Statement}\label{problem-statement-58}

\textbf{Temuan sumber:} Dalam perdebatan tentang kinerja ekonomi urban
Eropa, bukti mengenai struktur spasial yang paling berkaitan dengan
kinerja disebut terfragmentasi dan tidak selalu meyakinkan. Belum ada
konsensus mengenai keuntungan polisentrisitas, endogenitas tetap
mungkin, dan hasil dapat berubah menurut ukuran, definisi struktur
spasial, serta skala analisis. Analisis pan-Eropa juga disebut belum
tersedia, meskipun sudah ada studi khusus negara (TXT lines 141-174;
printed PDF pp.~49-50).

\textbf{Interpretasi penulis:} Masalah substantifnya dirumuskan sebagai
pertanyaan apakah massa ekonomi kritis hanya menguntungkan perusahaan
dan penduduk di kota besar, atau apakah kawasan polisentris yang terdiri
atas beberapa kota dapat menghasilkan keuntungan ekonomi yang serupa.
Pertanyaan terkaitnya adalah apakah ukuran kota-kota yang dibundel dapat
menggantikan keuntungan aglomerasi dan sekaligus mengurangi kerugian
dari ketergantungan pada satu kota besar (TXT lines 94-118, 141-174;
printed PDF pp.~48-50).

\subsection{Objectives}\label{objectives-58}

\textbf{Temuan sumber:} Artikel memperkenalkan analisis identifikasi
spasial-struktural pada region Eropa untuk periode 2010-12. Empat
hubungan yang diuji adalah: (1) ukuran populasi urban terhadap TFP
regional; (2) polisentrisitas terhadap TFP regional; (3) interaksi
ukuran dan polisentrisitas terhadap TFP regional; dan (4) dispersi
terhadap TFP regional (TXT lines 170-188, 804-812; printed PDF
pp.~49-50, 58-59).

\textbf{Hipotesis yang dinyatakan sumber:}

\begin{itemize}
\tightlist
\item
  Hipotesis 1: eksternalitas urbanisasi dalam bentuk populasi urban yang
  lebih besar berpengaruh positif terhadap produktivitas regional (TXT
  lines 206-215; printed PDF p.~50).
\item
  Hipotesis 2: derajat polisentrisitas yang lebih tinggi berpengaruh
  positif terhadap produktivitas regional (TXT lines 300-308; printed
  PDF p.~51).
\item
  Hipotesis 3: pengaruh positif polisentrisitas terhadap produktivitas
  regional berkurang ketika eksternalitas urbanisasi meningkat (TXT
  lines 271-272; printed PDF p.~51).
\item
  Hipotesis 4: derajat dispersi yang lebih tinggi berpengaruh negatif
  terhadap produktivitas regional, terlepas dari ukuran urban (TXT lines
  274-308; printed PDF pp.~51-52).
\end{itemize}

\textbf{Interpretasi penulis:} Tiga fitur metodologis dipakai untuk
mendukung tujuan tersebut: TFP sebagai ukuran kinerja ekonomi,
instrumental variables (IV) untuk menangani potensi endogenitas dalam
hubungan struktur-produktivitas, dan setting pan-Eropa dengan pengukuran
polisentrisitas morfologis (TXT lines 170-188; printed PDF p.~49).

\subsection{Literature Survey}\label{literature-survey-58}

\textbf{Temuan sumber (sintesis pustaka penulis):} Literatur ekonomi
urban yang dibahas memandang skala dan kepadatan sebagai sumber ekonomi
aglomerasi. Ekonomi aglomerasi dapat berupa return internal dalam
perusahaan atau return eksternal antarperusahaan; ekonomi urbanisasi
menyebarkan eksternalitas kepada perusahaan dalam region dan berkaitan
dengan ukuran atau kepadatan kawasan urban. Mekanisme yang disebut
meliputi tenaga kerja bersama, infrastruktur, fasilitas R\&D,
universitas, institusi lain, serta interaksi lintas sektor. Disekonomi
aglomerasi berasal dari persaingan atas sumber daya langka dan dapat
berupa harga lahan, upah, kemacetan, polusi, dan kebisingan yang lebih
tinggi (TXT lines 162-200; printed PDF pp.~49-50).

\textbf{Temuan sumber (sintesis pustaka penulis):} Literatur tentang
kota-kota yang saling terhubung mengusulkan bahwa kota kecil dan kota
lapis kedua dapat memperoleh ``borrowed size'' melalui limpahan manfaat
dan aliran pasokan dari kota tetangga. Jaringan kota juga diperkirakan
dapat membagi manfaat dan mengurangi sebagian kerugian konsentrasi.
Namun, sumber mencatat argumen tandingan bahwa eksternalitas jaringan
urban hanya dapat menggantikan eksternalitas urbanisasi sederhana secara
terbatas karena jarak antarkota dapat menimbulkan hambatan logistik dan
mengurangi keuntungan interaksi yang melekat pada lingkungan kota
tunggal yang padat (TXT lines 94-115, 206-237, 310-329; printed PDF
pp.~48-52).

\textbf{Temuan sumber (sintesis pustaka penulis):} Polisentrisitas
digunakan dalam beberapa skala: mikro atau intraurban, meso atau
regional/interurban, serta makro atau sistem urban nasional hingga
transnasional. Perspektif morfologis menilai distribusi dan ukuran pusat
urban, sedangkan perspektif fungsional menilai cara pusat-pusat tersebut
mengorganisasi wilayah melalui fungsi dan hubungan yang
multidireksional. Sumber menyebutkan bahwa banyak metode pengukuran
telah menghasilkan temuan yang berbeda dan belum ada konsensus
harmonisasi. Artikel ini memilih perspektif morfologis pada skala
regional Eropa (TXT lines 217-265; printed PDF pp.~50-51).

\textbf{Temuan sumber (sintesis pustaka penulis):} Konteks Eropa
dipandang berbeda dari Amerika Serikat karena populasi urban Eropa
terutama tinggal di kota kecil dan menengah serta kawasan urban Eropa
tertanam dalam region urban dan negara-kota yang terbentuk secara
historis. Sumber juga meninjau argumen bahwa dispersi penduduk di luar
inti urban dapat mengurangi aglomerasi dan membuat faktor produksi
terpusat, seperti modal manusia, kurang dimanfaatkan. Polisentrisitas
dan dispersi diperlakukan sebagai dua dimensi yang berkorelasi rendah
dan dapat membentuk kombinasi struktur spasial yang berbeda (TXT lines
279-308; printed PDF pp.~51-52).

\textbf{Interpretasi penulis:} Tinjauan tersebut mendukung pengujian
terpisah atas ukuran, polisentrisitas, interaksi ukuran-polisentrisitas,
dan dispersi. Fokus morfologis dipilih agar dapat dibandingkan dengan
studi terdahulu dan karena IV yang dibangun dari data historis akan
lebih sulit jika harus memasukkan hubungan fungsional yang kompleks (TXT
lines 238-265, 310-329, 352-360; printed PDF pp.~50-52).

\subsection{Method Used}\label{method-used-58}

\textbf{Temuan sumber:} Analisis utama menggunakan cross-section pada
rata-rata periode 2010-12. Model menjelaskan TFP regional dengan
pendidikan, ukuran populasi urban, polisentrisitas, dispersi,
konektivitas infrastruktur, spesialisasi pekerjaan bernilai tambah
tinggi, dan dummy makrogeografis. Semua variabel yang bukan dummy
ditransformasikan ke logaritma (TXT lines 401-471, 476-496, 516-532;
printed PDF pp.~53-55).

\textbf{Temuan sumber:} TFP dihitung dari fungsi produksi yang
meregresikan gross value added (GVA) terhadap modal dan tenaga kerja,
dengan efek tetap region dan tahun. Residual regresi tersebut menjadi
TFP untuk model utama. TFP dihitung pada tingkat NUTS-2 dan NUTS-1,
mencakup seluruh sektor di region, bukan hanya industri tertentu (TXT
lines 401-446; printed PDF pp.~53-54).

\textbf{Temuan sumber:} Tiga dimensi struktur spasial diukur secara
morfologis. Ukuran urban adalah total populasi yang tinggal di seluruh
city cores dalam region TL2. Polisentrisitas dihitung dari koefisien
rank-size distribution untuk dua, tiga, dan empat city cores terbesar;
koefisien tersebut dirata-ratakan, lalu nilai negatifnya diabsolutkan
dan diinversi sehingga koefisien yang lebih besar merepresentasikan
polisentrisitas yang lebih tinggi. Garis rank-size yang lebih curam
menunjukkan struktur lebih monosentris, sedangkan garis yang lebih datar
menunjukkan distribusi lebih merata dan lebih polisentris. Dispersi
adalah bagian populasi TL2 yang tidak tinggal di city cores, yaitu satu
dikurangi bagian populasi yang tinggal di city cores (TXT lines 310-400,
474-500, 561-574; printed PDF pp.~51-55).

\textbf{Temuan sumber:} Hanya city cores dengan sedikitnya 40.000
penduduk yang dimasukkan. Data populasi LAU-2 dikumpulkan dari Eurostat
Urban Audit, lalu diagregasikan ke TL2 tempat mayoritas populasi core
berada. Untuk France, Greece, Sweden, dan the UK, ketiadaan kode
penghubung memaksa penggunaan data agregat non-modular dari Urban Audit
(TXT lines 448-496; printed PDF pp.~54-55).

\textbf{Temuan sumber:} Estimasi OLS disajikan dalam tujuh spesifikasi
yang memasukkan satu, dua, atau tiga dimensi spasial, lalu dalam model
interaksi ukuran-polisentrisitas dan ukuran-dispersi. Robust standard
errors digunakan. Untuk menguji potensi bidirectional causality,
estimasi TSLS menggunakan ukuran urban dan polisentrisitas historis
tahun 1850 dari Bairoch et al.~serta luas lahan kontemporer untuk
instrumen dispersi. Instrumen interaksi dibuat dari nilai prediksi
first-stage yang kemudian dikalikan (TXT lines 425-471, 557-587,
592-620; printed PDF pp.~53-56).

\textbf{Temuan sumber:} Relevansi dan kekuatan instrumen diperiksa
dengan first-stage F-statistics dan Sanderson-Windmeijer F-statistics;
catatan tabel menamai statistik kedua sebagai SW. Model TSLS dilaporkan
dalam tiga spesifikasi: semua tiga dimensi, interaksi
ukuran-polisentrisitas, dan interaksi ukuran-dispersi (TXT lines
662-674, 818-860; printed PDF pp.~56-59).

\textbf{Inferensi pengolahan:} Karena model utama memakai observasi
region yang dirata-ratakan untuk satu periode cross-section, hasil OLS
dapat dibaca sebagai asosiasi regional bersyarat pada kontrol. TSLS
digunakan sumber untuk menangani potensi endogenitas dan menguji arah
kausalitas, tetapi desain ini tidak dengan sendirinya mengidentifikasi
efek pada individu atau perusahaan (TXT lines 476-496, 557-587, 662-667;
printed PDF pp.~54-56).

\subsection{Dataset}\label{dataset-58}

\textbf{Temuan sumber:} Dataset utama mencakup 168 region di 16 negara,
terutama negara dengan keanggotaan EU yang telah lama. Region
didefinisikan pada tingkat OECD TL2, yang menggunakan tingkat NUTS-2 dan
NUTS-1 Eurostat dan dimaksudkan agar region cukup besar serta kota tidak
melintasi batas interregional. Descriptive statistics pada Table 1
menggunakan N = 168; catatan tabel menyatakan seluruh variabel adalah
rata-rata 2010-12 kecuali jika disebutkan lain (TXT lines 474-496,
516-532; printed PDF pp.~54-55).

\textbf{Temuan sumber:} GVA regional, estimasi modal, dan tenaga kerja
berasal dari Cambridge Econometrics. Modal diestimasi dengan perpetual
inventory method atas gross fixed capital formation, sedangkan tenaga
kerja diberikan sebagai jumlah jam kerja penduduk regional. Populasi
city core dan municipality berasal dari Eurostat Urban Audit; total
populasi region berasal dari data NUTS-2 dan NUTS-1 Eurostat.
Konektivitas jalan, bandara, dan pelabuhan menggunakan indikator PBL
pada skala 1-10 (TXT lines 497-510, 536-544; printed PDF pp.~54-55).

\textbf{Temuan sumber:} Data historis Bairoch et al.~menyediakan
populasi permukiman urban Eropa dari 800 sampai 1850; artikel
menggunakan 1.732 permukiman pada 1850 untuk membentuk instrumen ukuran
dan polisentrisitas. Luas lahan dalam kilometer persegi digunakan
sebagai instrumen dispersi karena tidak tersedia estimasi historis total
populasi TL2 yang dianggap relevan untuk tujuan tersebut (TXT lines
557-574, 536-544; printed PDF pp.~55-56).

\textbf{Temuan sumber:} Table 1 melaporkan mean TFP log 0.0138 dengan
standard deviation 0.552 dan rentang -1.330 sampai 1.330; mean ukuran
populasi urban log 13.16; mean polisentrisitas log 0.284; dan mean
dispersi log -0.484. Connectivity memiliki mean log 1.615. Dummy
ditandai dengan \texttt{d}, sementara variabel log ditandai dengan
\texttt{ln} (TXT lines 516-532; printed PDF p.~55).

\textbf{Batas data yang disebutkan sumber:} Sebagian region tidak
dimasukkan karena data pada variabel atau tingkat analisis tidak
mencukupi. Format dan lokasi file data mentah tidak disebutkan dalam
file (TXT lines 476-496, 880-896; printed PDF pp.~54, 60).

\subsection{Population / Sample}\label{population-sample-58}

\textbf{Temuan sumber:} Unit analisis utama adalah 168 region OECD TL2
di 16 negara Eropa. Jumlah observasi pada model tanpa dispersi adalah
168, sedangkan model yang memasukkan dispersi dan model interaksi
menggunakan 166 observasi; sumber tidak menyatakan alasan spesifik untuk
selisih dua observasi tersebut (TXT lines 476-496, 704-705, 783-786,
856-857; printed PDF pp.~54-55, 57-59).

\textbf{Temuan sumber:} Untuk membangun ukuran struktur urban, sumber
mengolah 104.179 municipality LAU-2 dan memperoleh 730 city cores
setelah menerapkan ambang minimum 40.000 penduduk. City cores yang cukup
digunakan untuk rank-size distribution dua, tiga, dan empat core
terbesar dalam setiap region (TXT lines 452-472, 474-496; printed PDF
p.~54).

\textbf{Tidak berlaku:} Sumber tidak menggunakan responden manusia atau
sampel individu; populasi/sampelnya berupa region, municipality, dan
city cores sebagai unit agregat analisis (TXT lines 452-496; printed PDF
p.~54).

\subsection{Dependent Variables}\label{dependent-variables-58}

\textbf{Temuan sumber:} Dependent variable utama adalah log TFP regional
untuk rata-rata 2010-12. TFP merupakan residual dari fungsi produksi GVA
yang memperhitungkan modal, tenaga kerja, efek tetap region, dan efek
tetap tahun. Karena GVA mencakup seluruh aktivitas di region, TFP juga
dimaksudkan mencakup seluruh sektor yang hadir di region (TXT lines
401-446, 497-506; printed PDF pp.~53-55).

\textbf{Temuan sumber:} Dalam first-stage TSLS, variabel yang
diperlakukan sebagai endogenous dan diprediksi meliputi ukuran,
polisentrisitas, dispersi, serta interaksi ukuran dengan polisentrisitas
atau dispersi. Variabel-variabel tersebut adalah dependent variable
first-stage, bukan outcome substantif tambahan (TXT lines 557-587,
846-855; printed PDF pp.~55-56, 59).

\textbf{Tidak disebutkan dalam file:} Outcome utama lain di luar TFP
regional tidak dilaporkan.

\subsection{Results}\label{results-58}

\textbf{Temuan sumber - OLS:} Ukuran populasi urban berasosiasi positif
kuat dengan TFP dan tetap signifikan pada level 1\% ketika dimensi
spasial lain ditambahkan. Koefisien ukuran adalah 0.15 pada model ukuran
saja, 0.15 pada model ukuran-dispersi, dan 0.18 pada model penuh. Pada
spesifikasi yang memuatnya, polisentrisitas memiliki koefisien negatif;
koefisiennya tidak signifikan ketika berdiri sendiri (-0.04), tetapi
signifikan negatif ketika ukuran dikontrol (-0.07 pada model
ukuran-polisentrisitas dan model penuh). Dispersi berkoefisien negatif
tetapi tidak signifikan, dan menjadi 0.00 pada model penuh (TXT lines
592-661, 680-708; printed PDF pp.~56-57).

\textbf{Temuan sumber - OLS:} Connectivity berkoefisien positif dan
signifikan pada level 1\% pada semua model, dengan nilai antara 0.46 dan
0.55. Dummy Eastern Europe berkoefisien negatif dan signifikan pada
semua spesifikasi yang ditampilkan. Dalam spesifikasi Table 2,
pendidikan dan high-value-added employment LQ tidak signifikan. Jumlah
observasi berkisar 166-168 dan R2 berkisar 0.68-0.74 (TXT lines 662-708;
printed PDF p.~57).

\textbf{Temuan sumber - interaksi OLS:} Dalam Table 3, interaksi
ukuran-polisentrisitas berkoefisien -0.06 dan signifikan pada level 1\%
baik ketika dimasukkan sendiri maupun bersama interaksi ukuran-dispersi.
Interaksi ukuran-dispersi berkoefisien -0.06 tetapi tidak signifikan.
Model dengan kedua interaksi melaporkan koefisien polisentrisitas -0.10,
ukuran 0.18, dan dispersi 0.07 yang signifikan pada level 5\%; sumber
mengaitkan kemunculan koefisien positif dispersi ini dengan collinearity
tinggi antara dispersi dan interaksinya dengan ukuran (TXT lines
621-659, 743-792; printed PDF pp.~56, 58).

\textbf{Temuan sumber - TSLS:} Pada model TSLS semua tiga dimensi,
ukuran berkoefisien 0.34 dan signifikan pada level 1\%, polisentrisitas
-0.16 dan signifikan pada level 5\%, sedangkan dispersi 0.05 dan tidak
signifikan. Pada model dengan interaksi ukuran-polisentrisitas, ukuran
0.27 signifikan pada level 1\%, polisentrisitas -0.11 tidak signifikan,
dan interaksi ukuran-polisentrisitas -0.16 signifikan pada level 1\%.
Pada model dengan interaksi ukuran-dispersi, ukuran 0.35 signifikan pada
level 1\%, polisentrisitas -0.18 signifikan pada level 5\%, dispersi
0.14 tidak signifikan, dan interaksi ukuran-dispersi -0.11 tidak
signifikan (TXT lines 713-737, 818-860; printed PDF pp.~57-59).

\textbf{Temuan sumber - TSLS:} Penulis menyatakan hasil TSLS robust
terhadap estimasi OLS: ukuran populasi urban tetap menjadi prediktor
signifikan, polisentrisitas memiliki efek negatif signifikan dalam
spesifikasi tertentu, dan dispersi tidak menunjukkan efek signifikan.
Dalam spesifikasi IV yang memasukkan interaksi ukuran-polisentrisitas,
efek langsung polisentrisitas tidak lagi signifikan, sedangkan
interaksinya tetap negatif signifikan (TXT lines 711-737; printed PDF
pp.~57-59).

\textbf{Temuan sumber - first stage:} Table 4 melaporkan first-stage F
untuk ukuran 23.57-29.58, polisentrisitas 18.70-24.92, dan dispersi
5.27-8.38; SW F untuk ukuran 42.81-54.00, polisentrisitas 15.39-19.71,
dan dispersi 7.58-24.15. F untuk interaksi ukuran-polisentrisitas adalah
8.51 dengan SW F 30.35; F untuk interaksi ukuran-dispersi adalah 6.19
dengan SW F 8.80 (TXT lines 846-855; printed PDF p.~59).

\textbf{Inferensi pengolahan - audit konsistensi internal:} Narasi
sumber menyatakan bahwa semua F-statistics pada Table 4 berada di atas
rule-of-thumb 10 dan mengonfirmasi validitas pendekatan, tetapi Table 4
sendiri menampilkan beberapa first-stage F di bawah 10, termasuk 5.27,
6.19, dan 8.51. Review ini mencatat perbedaan antara narasi dan angka
tabel tanpa memilih interpretasi yang tidak dinyatakan sumber (TXT lines
711-713, 846-855; printed PDF pp.~57, 59).

\subsection{Findings}\label{findings-58}

\textbf{Interpretasi penulis:} Sumber merangkum empat temuan utama.
Pertama, ukuran populasi urban merupakan prediktor yang sangat
signifikan bagi TFP regional; Hipotesis 1 diterima. Kedua, tidak
ditemukan efek langsung positif polisentrisitas yang signifikan;
Hipotesis 2 ditolak. Ketiga, polisentrisitas berdampak negatif pada
produktivitas ketika ukuran urban meningkat, ditunjukkan oleh interaksi
ukuran-polisentrisitas yang negatif signifikan. Keempat, pengaruh
dispersi tidak signifikan atau ambigu dalam berbagai estimasi (TXT lines
704-708, 743-800, 804-812; printed PDF pp.~57-59).

\textbf{Interpretasi penulis:} Dalam spesifikasi IV dengan interaksi
ukuran-polisentrisitas yang termasuk model paling rinci, efek langsung
polisentrisitas menjadi tidak signifikan ketika interaksi dengan ukuran
dimasukkan, sedangkan interaksi negatif menunjukkan bahwa
polisentrisitas khususnya tidak bekerja pada region dengan populasi
urban relatif lebih besar. Penulis menyimpulkan bahwa, untuk region
dengan populasi urban yang sama, eksternalitas urbanisasi dari beberapa
city cores tidak menggantikan eksternalitas yang dicapai oleh struktur
dengan satu kota yang lebih besar (TXT lines 717-730, 775-792; printed
PDF pp.~57-59).

\textbf{Interpretasi penulis:} Hasil ini berbeda dari studi Amerika
Serikat yang dirujuk, yang menemukan efek langsung positif
polisentrisitas pada labor productivity, tetapi memiliki kemiripan pada
interaksi negatif ukuran-polisentrisitas. Sumber juga mengaitkan
hasilnya dengan studi OECD yang memberikan arah berbeda menurut skala
analisis (TXT lines 757-780; printed PDF p.~58).

\textbf{Interpretasi penulis:} Sumber menyatakan bahwa yang penting
adalah cara populasi didistribusikan di antara city cores, bukan cara
populasi tersebar antara city cores dan wilayah regional lainnya. Namun,
kesimpulan ini disertai caveat bahwa pengukuran dispersi yang digunakan
relatif kasar dan hasilnya sedikit (TXT lines 793-800, 865-872; printed
PDF pp.~58-60).

\subsection{Conclusions}\label{conclusions-58}

\textbf{Interpretasi penulis:} Kesimpulan utama artikel adalah bahwa
urbanization externalities memiliki peran yang lebih besar terhadap
produktivitas regional daripada yang disarankan oleh perdebatan
kebijakan pan-Eropa mengenai polisentrisitas dan dispersi. Ukuran urban
tetap menunjukkan hubungan positif yang kuat, sedangkan polisentrisitas
tidak menunjukkan keuntungan langsung positif dan menjadi negatif
melalui interaksinya dengan ukuran; dispersi tidak menghasilkan efek
yang konsisten (TXT lines 897-905; printed PDF p.~60).

\textbf{Interpretasi penulis:} Hasil tersebut mendukung place-based
regional development strategies yang lebih terintegrasi dengan agenda
ekonomi urban, dengan ketergantungan yang lebih hati-hati pada
polisentrisitas sebagai solusi umum untuk mengurangi ketimpangan
regional (TXT lines 897-905; printed PDF p.~60).

\textbf{Inferensi pengolahan:} Kesimpulan sumber tidak sama dengan
pernyataan bahwa setiap struktur polisentris merugikan produktivitas.
Bukti yang lebih spesifik adalah tidak adanya efek langsung positif yang
konsisten dan adanya interaksi negatif yang bergantung pada ukuran
urban, sementara efek dispersi tidak jelas (TXT lines 717-730, 793-800;
printed PDF pp.~57-59).

\subsection{Contributions}\label{contributions-58}

\textbf{Interpretasi penulis:} Artikel menyatakan tiga aspek baru.
Pertama, artikel menggunakan TFP sebagai ukuran kinerja ekonomi regional
yang dianggap lebih baik dan secara inheren mengontrol endowment modal
serta tenaga kerja. Kedua, artikel menggunakan IV untuk menangkap
potensi endogenitas dalam hubungan struktur-produktivitas dan
mengembangkan kontrol serta IV yang spesifik untuk konteks Eropa.
Ketiga, analisis dilakukan pada setting pan-Eropa dengan pengukuran
polisentrisitas morfologis yang mengikuti delineasi Brezzi dan Veneri
(TXT lines 170-188, 401-446, 731-738; printed PDF pp.~49, 53-54, 57-58).

\textbf{Interpretasi penulis:} Penulis memosisikan penelitian ini
sebagai eksplorasi pertama atas hubungan produktivitas dan struktur
spasial urban dalam setting regional pan-Eropa yang kaya kebijakan.
Kontribusi empirisnya adalah memperluas penelitian morfologis sebelumnya
di Amerika Serikat ke region Eropa dengan tiga dimensi spasial: ukuran,
polisentrisitas, dan dispersi (TXT lines 795-803; printed PDF
pp.~58-59).

\subsection{Research Gap}\label{research-gap-58}

\textbf{Temuan sumber:} Sebelum artikel ini, bukti tentang hubungan
struktur spasial dan kinerja ekonomi di Eropa disebut terfragmentasi,
studi khusus negara lebih tersedia daripada analisis pan-Eropa, dan
konsensus mengenai manfaat polisentrisitas belum tercapai. Potensi
endogenitas serta perbedaan skala dan definisi juga menghambat
identifikasi (TXT lines 154-174; printed PDF p.~49).

\textbf{Temuan sumber:} Artikel menyisakan gap pada hubungan fungsional
antarpusat, ketergantungan skala interurban dan intraurban, serta
ketergantungan waktu. Sumber menyebut bahwa pengukuran morfologis tidak
menangkap seluruh functional linkages, pendekatan fungsional memerlukan
data Eropa antarrregion yang belum tersedia, dan skala TL2 mungkin tidak
sesuai dengan proses interlokal yang berlangsung hanya di bagian region
(TXT lines 141-150, 238-265, 865-872; printed PDF pp.~49-51, 59-60).

\textbf{Interpretasi penulis:} Agenda gap lanjutan yang tersisa mencakup
indikator dispersi yang lebih masuk akal, data panel dengan variasi
waktu kontrol yang lebih besar, transparansi dan konsistensi pengumpulan
Urban Audit, serta pemahaman yang lebih mendalam mengenai hasil variabel
pendidikan dan konektivitas (TXT lines 580-587, 865-892; printed PDF
pp.~56, 59-60).

\subsection{Limitations}\label{limitations-58}

\textbf{Temuan sumber:} Penulis menyatakan bahwa pilihan dan asumsi
dalam rancangan penelitian pada umumnya bersifat arbitrer meskipun
berakar pada literatur. Penekanan pada polisentrisitas morfologis
berarti pendekatan fungsional dapat menghasilkan outcome berbeda, tetapi
data Eropa antarrregion yang canggih untuk pendekatan tersebut belum
tersedia (TXT lines 800-812, 865-872; printed PDF pp.~58-60).

\textbf{Temuan sumber:} Skala administratif TL2 dapat bersifat arbitrer
karena proses hubungan interlokal yang sebenarnya bisa berlangsung hanya
dalam bagian tertentu dari region. Sensitivitas ini tidak dapat diuji
karena tidak ada data regional yang harmonis. Pengukuran dispersi juga
relatif kasar karena dipilih untuk kebutuhan perbandingan dan kesesuaian
dengan IV, sehingga hanya menghasilkan sedikit temuan pada indikator
tersebut (TXT lines 865-872; printed PDF pp.~59-60).

\textbf{Temuan sumber:} Beberapa variabel sensitif terhadap pengukuran,
terutama accessibility yang terutama mengukur konektivitas eksternal
regional. Dalam uraian keterbatasan, penulis menyatakan bahwa kontrol
spesialisasi dan pendidikan tidak menunjukkan hubungan signifikan dengan
produktivitas dan bahwa alasannya masih tidak pasti. Data Urban Audit
mengalami perubahan yang sering dalam definisi, cakupan, dan cara
pengukuran (TXT lines 865-888; printed PDF pp.~59-60).

\textbf{Inferensi pengolahan:} Pernyataan naratif tersebut tidak
sepenuhnya seragam dengan Table 4, yang melaporkan koefisien pendidikan
negatif dan signifikan pada level 1\% dalam ketiga model TSLS, dengan
nilai -0.45, -0.43, dan -0.42. Review mempertahankan kedua informasi
dari TXT ini dan tidak memutuskan mana yang harus menggantikan yang lain
(TXT lines 818-837; printed PDF p.~59).

\textbf{Temuan sumber:} Seluruh variabel berpotensi mendapat manfaat
dari data longitudinal atau panel. Ukuran, polisentrisitas, dan dispersi
tidak berubah cepat, tetapi variasi waktu yang lebih besar pada variabel
kontrol dianggap penting. Periode penelitian 2010-12 berada di tengah
kemunduran ekonomi, sehingga variasi regional dalam economic resilience
dapat berinteraksi dengan hasil (TXT lines 888-892; printed PDF p.~60).

\textbf{Temuan sumber:} Pemilihan aggregate GVA dan tetap memasukkan
aktivitas non-market serta public administration dapat menghasilkan
estimasi TFP yang lebih rendah jika efeknya tidak sepenuhnya dikontrol.
Penulis menyebut temuan tersebut sebagai lower-bound estimates yang
kecil kemungkinan bias ke atas; menghapus industri tertentu juga akan
mengurangi sifat umum temuan dan dapat memasukkan bias atau arbitrari
dalam pemilihan industri (TXT lines 920-938; printed PDF p.~60).

\textbf{Temuan sumber:} Sampel tidak mencakup semua region yang datanya
tersedia. Negara dengan data Eurostat atau Urban Audit yang tidak
lengkap yang disebutkan adalah Bulgaria, Croatia, Cyprus, Denmark,
Estonia, Finland, Iceland, Latvia, Macedonia, Norway, Romania, dan
Switzerland. Contoh region yang dikeluarkan karena tidak ada cukup city
cores untuk menghitung polisentrisitas pada 2010-12 adalah Valle
d'Aosta, Alentejo, Norra Mellansverige, dan Highlands \& Islands;
beberapa region lain juga tidak memiliki perhitungan polisentrisitas
1850 (TXT lines 939-942, 880-896; printed PDF pp.~60-61).

\subsection{Challenges}\label{challenges-58}

\textbf{Temuan sumber:} Tantangan identifikasi utama adalah
bidirectionality: produktivitas regional dapat memengaruhi struktur
urban, bukan hanya sebaliknya. Hal ini mendorong penggunaan IV,
instrumen historis, first-stage regressions, dan instrumen khusus untuk
interaction terms (TXT lines 425-429, 557-587, 662-674; printed PDF
pp.~53-56).

\textbf{Temuan sumber:} Konstruksi ukuran spasial memerlukan proses
multi-langkah untuk menghindari kota melintasi batas region.
Tantangannya meliputi pemetaan LAU-2 ke city cores dan TL2, ambang
minimum 40.000 penduduk, penggunaan agregasi non-modular untuk France,
Greece, Sweden, dan the UK, serta pembentukan rank-size distribution
yang sebanding lintas region (TXT lines 448-496; printed PDF pp.~54-55).

\textbf{Temuan sumber:} Keterbatasan data historis membuat instrumen
dispersi tidak dibentuk dari populasi TL2 historis, melainkan dari luas
lahan kontemporer. Instrumen interaksi juga harus dibentuk melalui nilai
prediksi first-stage sebelum interaksi dapat dihitung (TXT lines
536-587; printed PDF pp.~55-56).

\textbf{Temuan sumber:} Ketersediaan data yang tidak seragam menyebabkan
sebagian negara dan region dikeluarkan, sementara perubahan definisi dan
pengukuran Urban Audit menyulitkan konsistensi. Kekurangan data
fungsional yang harmonis juga membatasi kemampuan untuk menguji hubungan
antarkota yang sebenarnya (TXT lines 476-496, 865-888; printed PDF
pp.~54-55, 59-60).

\textbf{Inferensi pengolahan:} Perbedaan antara klaim naratif mengenai
kekuatan first-stage dan angka F yang tercetak di Table 4 merupakan
tantangan audit interpretasi hasil IV, bukan dasar untuk menambahkan
validitas yang tidak ditunjukkan secara konsisten oleh TXT (TXT lines
711-713, 846-860; printed PDF pp.~57, 59).

\subsection{Practical Implications}\label{practical-implications-58}

\textbf{Interpretasi penulis:} Strategi place-based untuk pembangunan
regional dapat memperoleh manfaat dari integrasi yang lebih besar dengan
agenda ekonomi urban, dengan ketergantungan yang tidak terlalu yakin
pada polisentrisitas sebagai panacea untuk mengurangi ketimpangan
regional. Temuan ukuran urban yang kuat menjadi dasar untuk menilai
kembali bobot kebijakan yang diberikan pada eksternalitas urbanisasi,
polisentrisitas, dan dispersi (TXT lines 897-905; printed PDF p.~60).

\textbf{Temuan sumber:} Artikel menyatakan bahwa temuan mempunyai
implikasi bagi agenda kebijakan Uni Eropa mengenai pembangunan urban dan
produktivitas regional (TXT lines 80-86; printed PDF p.~48).

\subsection{Applications}\label{applications-58}

\textbf{Temuan sumber:} Penerapan yang disebut secara eksplisit adalah
penggunaan hasil untuk agenda kebijakan Uni Eropa mengenai pembangunan
urban dan produktivitas regional serta untuk place-based regional
development strategies (TXT lines 80-86, 897-905; printed PDF pp.~48,
60).

\textbf{Tidak disebutkan dalam file:} Nama program kebijakan tertentu,
prosedur implementasi, atau aplikasi sektoral yang lebih spesifik.

\subsection{Future Research
(Optional)}\label{future-research-optional-58}

\textbf{Temuan sumber:} Sumber menyebut beberapa arah lanjutan:
penelitian yang lebih mendalam mengenai pendidikan dan konektivitas;
analisis polisentrisitas dari perspektif fungsional jika data
antarrregion Eropa yang memadai tersedia; pengembangan indikator
dispersi yang lebih masuk akal; dan penggunaan data longitudinal atau
panel dengan variasi waktu yang lebih besar pada variabel kontrol (TXT
lines 580-587, 807-812, 865-892; printed PDF pp.~56, 59-60).

\textbf{Temuan sumber:} Sumber juga mendorong transparansi yang lebih
besar dan konsistensi dalam cara variabel Urban Audit didefinisikan,
dimasukkan, dan diukur (TXT lines 880-888; printed PDF p.~60).

\subsection{Provenance Notes}\label{provenance-notes-58}

\begin{itemize}
\tightlist
\item
  \textbf{Identifikasi sumber:} Pencarian di bawah \texttt{source/}
  menemukan tepat satu TXT dengan prefiks B15, yaitu file sumber yang
  tercantum di awal review. Tidak ada file TXT sumber lain berprefiks
  B15 yang dipakai.
\item
  \textbf{Cakupan pembacaan:} Seluruh teks yang tersedia dibaca dari
  line 1 sampai line 1074, termasuk blok \texttt{{[}PDF\ Metadata{]}}
  pada TXT lines 1-25 dan blok \texttt{{[}Extracted\ Text{]}} mulai TXT
  line 26. TXT dan PDF sumber tidak diedit.
\item
  \textbf{Metadata ekstraksi:} TXT menyatakan sumber PDF relatifnya,
  tanggal ekstraksi 2026-08-31, metode \texttt{pdftotext\ -layout}, dan
  metadata PDF berupa 16 pages, tidak terenkripsi, serta ukuran halaman
  634.32 x 833.04 pts (TXT lines 1-25).
\item
  \textbf{Metadata bibliografis:} Judulnya \emph{Spatial structure and
  productivity in European regions}, penulis Wessel M. Ouwehand, Frank
  G. van Oort, dan Nicola Cortinovis; terbit di \emph{Regional Studies},
  volume 56 nomor 1, halaman 48-62, DOI
  \texttt{10.1080/00343404.2021.1950912}; published online 17 August
  2021 (TXT lines 37-55, 70-92).
\item
  \textbf{Locator:} \texttt{TXT\ lines} merujuk pada nomor baris hasil
  pembacaan file ekstraksi; \texttt{printed\ PDF\ p./pp.} merujuk pada
  nomor halaman artikel yang tercetak di teks. Table 1-4 dan Appendix
  Tables yang disebut sumber diperlakukan sebagai bagian dari isi TXT.
\item
  \textbf{Keterbatasan akses provenance:} Artikel menyebut supplemental
  data online dan Appendix Tables A1-A3, tetapi review ini hanya
  menggunakan isi TXT yang tersedia dan tidak mengklaim telah membaca
  materi suplementer di luar file (TXT lines 131, 299, 665-674).
\item
  \textbf{Audit ekstraksi:} Teks dapat dibaca lengkap untuk review ini.
  Format dua kolom dan urutan beberapa baris tabel/teks tetap dapat
  menimbulkan interleaving; locator halaman dan label tabel
  dipertahankan. Anomali antara pernyataan naratif tentang first-stage
  F-statistics dan angka Table 4 dicatat pada bagian Results dan
  Challenges, bukan diselesaikan dengan informasi eksternal (TXT lines
  711-713, 818-860).
\item
  \textbf{Hash sumber saat verifikasi:} SHA-256
  \texttt{c9fe72ffb703d25d8fb9531a8e4c2aa16aa1c3bba6b570a694a83d825138b0ea}.
\item
  \textbf{Batas klaim:} Literatur yang diringkas di bagian Literature
  Survey adalah sintesis yang ditulis artikel ini; karya-karya dalam
  daftar referensi tidak dibaca sebagai sumber terpisah. Tidak ada
  informasi dari luar TXT yang digunakan.
\end{itemize}

\section{Review B16}\label{review-b16}

Kode: B16

\subsection{Summarized Abstract}\label{summarized-abstract-59}

\textbf{{[}Temuan dilaporkan sumber{]}} Artikel menyatakan bahwa riset
tentang dampak ekonomi urban polycentricity masih menghasilkan temuan
yang tidak konklusif. Artikel mengembangkan kerangka longitudinal yang
menghubungkan perubahan polycentricity di urban regions (URs) China
dengan perubahan total factor productivity (TFP). Tidak ditemukan bukti
bahwa urban polycentricity secara umum mendorong pertumbuhan ekonomi.
Namun, hubungan tersebut bergantung pada ukuran populasi dan interaksi
antarkota. Artikel juga melaporkan bahwa kota-kota di UR besar dapat
meminjam ukuran dari kota-kota di dekatnya, dan hal itu berkontribusi
pada pertumbuhan ekonomi regional (PDF p.~1018, abstract; TXT lines
95-101).

\textbf{{[}Interpretasi penulis{]}} Temuan digunakan untuk merefleksikan
strategi pembangunan regional dan perkotaan China (PDF p.~1018,
abstract; TXT lines 99-101).

\textbf{{[}Inferensi pengolahan{]}} Unit bukti utama artikel adalah UR
di China dan TFP, sehingga ringkasan ini tidak memperlakukan hasilnya
sebagai bukti universal untuk semua negara atau semua ukuran kinerja
ekonomi.

\subsection{Problem Statement}\label{problem-statement-59}

\textbf{{[}Temuan dilaporkan sumber{]}} Literatur yang ditinjau
menghasilkan tiga pola: polycentricity tidak berdampak langsung pada
kinerja regional; polycentricity berkaitan dengan produktivitas regional
yang lebih tinggi; atau polycentricity berkaitan dengan kinerja yang
lebih rendah. Hasil yang beragam tersebut menciptakan kebutuhan akan
bukti yang lebih kuat mengenai dampak ekonomi pembangunan polycentric
urban regions (PURs) (PDF pp.~1018-1019, Section 1; TXT lines 109-134).

\textbf{{[}Interpretasi penulis{]}} Perbedaan hasil mungkin berasal dari
perbedaan definisi dan operasionalisasi polycentricity, pilihan
indikator kinerja ekonomi, cara menangani endogeneity, serta penggunaan
data sectional atau panel. Artikel menilai bahwa aspek pemilihan
indikator, endogeneity, dan desain longitudinal masih kurang
diperhatikan (PDF pp.~1018-1019, Section 1; TXT lines 126-189).

\textbf{{[}Inferensi pengolahan{]}} Masalah penelitian dapat dirumuskan
sebagai pengujian hubungan yang bersyarat antara struktur ruang
polycentric dan kinerja ekonomi regional, bukan pengujian sederhana
apakah polycentricity selalu baik atau buruk.

\subsection{Objectives}\label{objectives-59}

\textbf{{[}Temuan dilaporkan sumber{]}} Tujuan eksplisit artikel adalah
mengkaji hubungan yang sedang berkembang antara struktur ruang
polycentric dan kinerja ekonomi di UR China pada 2005-2017. Artikel
menggunakan TFP sebagai ukuran kinerja, mengembangkan desain
longitudinal, dan menerapkan time-varying instrumental variables (IV)
untuk menangani potensi endogeneity (PDF pp.~1019, Section 1; TXT lines
185-210).

\textbf{{[}Temuan dilaporkan sumber{]}} Lima hipotesis yang diuji
adalah: (1) tingkat urban polycentricity yang lebih tinggi kondusif bagi
pertumbuhan ekonomi regional; (2) pengaruhnya bergantung pada
keseimbangan relatif borrowed size dan agglomeration shadow; (3)
dampaknya berkurang ketika ukuran populasi urban meningkat; (4) posisi
UR dalam hierarki administratif China memengaruhi dampaknya; dan (5) UR
pesisir dan pedalaman menunjukkan pola yang berbeda (PDF pp.~1020-1021,
Sections 2.2-2.3; TXT lines 250-263 and 303-313).

\subsection{Literature Survey}\label{literature-survey-59}

\textbf{{[}Temuan dilaporkan sumber{]}} Sumber mengelompokkan pengukuran
pertumbuhan ekonomi regional ke dalam pola upah, indikator pertumbuhan
kota seperti populasi dan pekerjaan, serta tingkat produktivitas. Labour
productivity (LP) lebih mudah dihitung dan ditafsirkan, tetapi hanya
memasukkan tenaga kerja. TFP dipaparkan sebagai ukuran yang juga
menangkap faktor produksi lain seperti modal fisik, modal manusia, dan
kemajuan teknologi (PDF pp.~1019-1020, Section 2.1; TXT lines 167-231).

\textbf{{[}Temuan dilaporkan sumber{]}} Kerangka teoritis yang dibahas
menghubungkan PUR dengan regional externalities, agglomeration
economies, network externalities, borrowed size, dan agglomeration
shadow. Borrowed size menggambarkan manfaat fungsi yang dipinjam kota
kecil atau menengah dari kota besar di dekatnya; agglomeration shadow
menggambarkan kemungkinan terhambatnya pertumbuhan kota kecil karena
kompetisi dengan kota besar. Keduanya diperlakukan sebagai dua efek yang
dapat membentuk dampak bersih interaksi antarkota (PDF pp.~1020-1021,
Section 2.2; TXT lines 232-276).

\textbf{{[}Temuan dilaporkan sumber{]}} Tinjauan juga menunjukkan bahwa
pengaruh polycentricity dapat bergantung pada ukuran UR, jarak
antarkota, tingkat konektivitas, serta konteks kelembagaan dan ekonomi.
Dalam konteks China, rencana polycentric dapat dibentuk melalui
interaksi proses top-down dan bottom-up, sementara kepentingan
pemerintah lokal dapat membatasi kerja sama regional (PDF pp.~1020-1021,
Section 2.2-2.3; TXT lines 284-317).

\textbf{{[}Interpretasi penulis{]}} Bukti terdahulu belum konsisten
karena pendekatan pengukuran, indikator ekonomi, penanganan endogeneity,
dan desain data berbeda-beda. Perbandingan China dengan konteks Eropa
atau Amerika Serikat karena itu perlu memperhatikan kondisi
institusional dan ekonomi yang berbeda (PDF pp.~1018-1021, Sections
1-2.3; TXT lines 122-189 and 267-317).

\subsection{Method Used}\label{method-used-59}

\textbf{{[}Temuan dilaporkan sumber{]}} Penelitian menggunakan desain
longitudinal panel pada 19 UR China selama periode 2005-2017. Unit
keputusan dalam pengukuran TFP adalah setiap UR pada setiap tahun. Model
utama adalah fixed-effects panel data model dengan variabel dependen
\texttt{LN(TFP\_it)} (PDF pp.~1021-1023, Sections 3.1-3.3.2; TXT lines
321-424).

\textbf{{[}Temuan dilaporkan sumber{]}} Polycentricity morfologis
dibangun dalam dua tahap. Pertama, LandScan digunakan untuk memilih 5\%
grid berpenduduk paling padat, yaitu grid di atas persentil ke-95
populasi kota. Grid yang berdekatan dalam delapan arah digabungkan
menjadi cluster; cluster dipertahankan jika mencakup sedikitnya 3 km2
dan memuat lebih dari 100.000 penduduk. Kedua, indeks dihitung sebagai
\texttt{1\ -\ s\_obs/s\_max}, dengan rentang 0 untuk tidak polycentric
sampai 1 untuk kondisi ideal ketika semua kota sama besar (PDF
pp.~1022-1023, Section 3.3.1; TXT lines 356-378 and 799-837).

\textbf{{[}Temuan dilaporkan sumber{]}} TFP dihitung menggunakan
Malmquist index non-parametrik dalam data envelopment analysis (DEA).
Setiap UR-tahun diperlakukan sebagai decision-making unit; inputnya
adalah stok modal tetap dan tenaga kerja, sedangkan outputnya adalah GDP
UR. Nilai TFP lebih besar, sama dengan, atau lebih kecil dari 1
masing-masing menunjukkan peningkatan, keadaan tetap, atau penurunan
efisiensi (PDF pp.~1023, Section 3.3.1; TXT lines 394-424).

\textbf{{[}Temuan dilaporkan sumber{]}} Model fixed effects menguji
hipotesis secara bertahap: model polycentricity dasar untuk Hipotesis 1;
penambahan \texttt{BOR} dan \texttt{BOR2} untuk Hipotesis 2; serta
interaksi \texttt{POP*POLY} untuk Hipotesis 3. Analisis subsampel
dilakukan berdasarkan hierarki administratif dan lokasi pesisir atau
pedalaman untuk menguji Hipotesis 4 dan 5. Variabel kontrol meliputi
human capital, informatisation, ukuran relatif industri tersier, FDI,
dan intervensi pemerintah. Variabel ditransformasi ke bentuk logaritmik
untuk memperoleh bentuk linear (PDF pp.~1023-1024, Section 3.3.2; TXT
lines 409-477).

\textbf{{[}Temuan dilaporkan sumber{]}} BOR dihitung dari populasi
kota-kota tetangga yang diberi bobot jarak lalu diagregasikan ke tingkat
regional. Untuk menangani potensi hubungan dua arah antara
polycentricity dan kinerja ekonomi, artikel menggunakan FE-IV. IV
time-varying dibangun dari bentuk kota historis dan jalur ekspansi kota
yang diprediksi secara mekanis, dengan empat indeks bentuk: cohesion,
proximity, spin, dan range. Prosedurnya memakai urban area tahun 1992,
lahan yang dapat dikembangkan, proyeksi populasi 2005-2017, dan bentuk
ekspansi melingkar tanpa kendala geografis (PDF pp.~1024-1026, Section
3.3.2; TXT lines 482-550).

\textbf{{[}Temuan dilaporkan sumber{]}} Sebelum hasil utama, sumber
melaporkan uji panel unit, multicollinearity, heteroscedasticity, dan
autocorrelation. Hasil VIF, Breusch-Pagan, dan Wooldridge dipakai untuk
menyatakan bahwa masalah-masalah tersebut tidak relevan bagi spesifikasi
yang digunakan (PDF p.~1026, Section 3.3.2; TXT lines 553-563).

\subsection{Dataset}\label{dataset-59}

\textbf{{[}Temuan dilaporkan sumber{]}} Dataset yang disebutkan terdiri
atas:

\begin{itemize}
\tightlist
\item
  LandScan High-Resolution Global Population Dataset untuk 2005-2017,
  berupa distribusi populasi pada grid 1 km2.
\item
  Data sosial-ekonomi tingkat prefektur, termasuk GDP dan FDI, dari
  China City Statistical Yearbooks, China Urban database, dan China
  Provincial Statistical Yearbooks untuk 2005-2017.
\item
  Statistical Bulletins untuk 12 autonomous prefectures yang tidak
  terdokumentasi dalam yearbook tersebut.
\item
  DMSP/OLS Night-time Lights sebagai dataset alternatif untuk membangun
  IV; data ini berupa rekaman intensitas cahaya tahunan dengan nilai
  0-63 dan resolusi sekitar 1 km2.
\end{itemize}

\textbf{{[}Temuan dilaporkan sumber{]}} Sumber juga menyebut stok modal
tetap yang diestimasi dengan perpetual inventory method dan jumlah
pekerja yang bekerja di seluruh sektor sebagai input TFP (PDF
pp.~1021-1025, Sections 3.1-3.3.2; TXT lines 333-345, 394-424, and
477-485).

\textbf{{[}Temuan dilaporkan sumber{]}} Karena masalah definisi
statistik, data empat kota hilang untuk beberapa tahun. Penulis
menggunakan linear interpolation untuk mengatasinya. Data supplemental
yang dirujuk artikel tidak tersedia dalam TXT yang diperiksa ini (PDF
pp.~1021-1022, Section 3.2; TXT lines 333-356).

\subsection{Population / Sample}\label{population-sample-59}

\textbf{{[}Temuan dilaporkan sumber{]}} Populasi analisis berupa 19 UR
yang diidentifikasi dalam China's 14th Five-Year Plan (2021-2025): 5
national URs, 8 regional URs, dan 6 subregional URs. Enam UR berada di
China pesisir dan 13 berada di China pedalaman. Pada 2017, 19 UR
tersebut secara bersama-sama mencakup 29\% luas daratan nasional, 68\%
populasi, dan 83\% GDP. BBW UR di Guangxi tidak dimasukkan sebagai UR
pesisir karena tidak seluruh batasnya berada di pesisir dan UR tersebut
berbatasan dengan Vietnam (PDF pp.~1021, 1030, Sections 3.1 and Notes;
TXT lines 321-331 and 787-791).

\textbf{{[}Temuan dilaporkan sumber{]}} Unit analisis ekonometrik adalah
UR-tahun. Tabel 1 melaporkan 228 observasi untuk model sampel penuh,
sedangkan model subsampel melaporkan 60 observasi untuk national URs, 96
untuk regional URs, 72 untuk subregional URs, dan 72 untuk coastal URs
serta 156 untuk inland URs (PDF pp.~1027-1028, Tables 1-2; TXT lines
586-616 and 649-685).

\textbf{{[}Inferensi pengolahan{]}} Angka observasi pada tabel tidak
dijadikan dasar untuk menghitung ulang jumlah tahun atau mengoreksi
periode penelitian, karena artikel menyatakan periode 2005-2017
sekaligus melaporkan data yang hilang dan interpolasi.

\subsection{Dependent Variables}\label{dependent-variables-59}

\textbf{{[}Temuan dilaporkan sumber{]}} Variabel dependen dalam
persamaan regresi adalah \texttt{LN(TFP\_it)}, yaitu log total factor
productivity UR \texttt{i} pada tahun \texttt{t}. TFP diperoleh dari
Malmquist index DEA dengan GDP sebagai output dan stok modal tetap serta
tenaga kerja sebagai input. Sumber menjelaskan bahwa indeks di atas 1
menunjukkan kenaikan efisiensi, sama dengan 1 tidak berubah, dan di
bawah 1 penurunan (PDF pp.~1023-1024, Section 3.3; TXT lines 394-424).

\subsection{Results}\label{results-59}

\textbf{{[}Temuan dilaporkan sumber{]}} Pada model baseline, penambahan
POLY menghasilkan koefisien negatif dan signifikan dalam model 2.
Koefisien tersebut tetap negatif dan signifikan ketika ukuran populasi
serta interaksi POLY dengan POP dimasukkan dalam model 3. Hipotesis 1
ditolak. Sumber melaporkan bahwa tanda variabel kontrol konsisten di
berbagai spesifikasi: HUM, INF, dan FDI berasosiasi positif dengan TFP,
sedangkan TIN dan GOV berasosiasi negatif dalam uraian model baseline
(PDF pp.~1026-1027, Section 4.1; TXT lines 566-579 and 586-616).

\textbf{{[}Temuan dilaporkan sumber{]}} Interaksi \texttt{POLY*POP}
bernilai positif dan signifikan pada model 3. Sumber menghitung bahwa
pengaruh polycentricity diperkirakan mulai menjadi positif pada UR
dengan populasi lebih dari 20,94 juta; sekitar sepertiga UR China
disebut terlalu kecil untuk menangkap sepenuhnya manfaat ekonomi yang
diasosiasikan dengan polycentricity. Hasil ini berlawanan dengan
Hipotesis 3, yang memprediksi pelemahan dampak ketika ukuran populasi
meningkat (PDF p.~1026, Section 4.1; TXT lines 552-565).

\textbf{{[}Temuan dilaporkan sumber{]}} Analisis subsampel menunjukkan
heterogenitas berdasarkan hierarki administratif dan lokasi. Dampak
negatif POLY terhadap TFP paling kuat pada subregional URs, lalu
national URs dan regional URs. Coastal URs menunjukkan pengaruh POLY
yang signifikan dan positif, sedangkan inland URs menunjukkan pengaruh
yang signifikan dan negatif. Hipotesis 4 dan 5 dinyatakan diterima (PDF
pp.~1026-1027, Section 4.1; TXT lines 566-579; Table 1, TXT lines
586-616).

\textbf{{[}Temuan dilaporkan sumber{]}} Untuk mekanisme borrowed size
dan agglomeration shadow, model yang hanya memasukkan BOR dan
\texttt{BOR2} tidak menunjukkan pengaruh yang tampak pada TFP, sehingga
Hipotesis 2 ditolak dalam spesifikasi tersebut. Ketika BOR
diinteraksikan dengan POP, efek bersih menjadi positif setelah ukuran UR
melampaui 18,3 juta. Interaksi \texttt{BOR*POLY} tidak signifikan pada
model yang mengujinya (PDF pp.~1027-1028, Section 4.2; TXT lines 620-643
and 649-685).

\textbf{{[}Temuan dilaporkan sumber{]}} Hasil subsampel untuk BOR
berbeda antarhierarki: uraian penulis menyatakan bahwa national URs
memiliki koefisien BOR positif dan koefisien kuadrat negatif; regional
URs menunjukkan hubungan non-linear negatif sebagaimana dijelaskan dalam
teks; dan direct borrowed-size effect pada subregional URs dinyatakan
signifikan negatif. Sumber kemudian mengaitkan pola tersebut dengan
kemungkinan mitigasi externalities negatif melalui sistem yang lebih
terdesentralisasi dan koordinasi sukarela antarkota (PDF p.~1027,
Section 4.2; TXT lines 628-643).

\textbf{{[}Temuan dilaporkan sumber{]}} Model FE-IV menunjukkan bahwa IV
berkorelasi dengan POLY tetapi tidak dengan TFP. Estimasi IV
mengonfirmasi temuan utama bahwa tingkat polycentricity yang lebih
tinggi berkorelasi dengan TFP yang lebih rendah, dan uji endogeneity
tidak menunjukkan masalah endogeneity dalam spesifikasi yang digunakan
(PDF p.~1028, Section 4.3; TXT lines 689-705).

\textbf{{[}Inferensi pengolahan{]}} Hasil paling konsisten di seluruh
uraian adalah hubungan POLY-TFP yang negatif pada sampel umum, tetapi
efeknya berubah menurut ukuran UR, hierarki administratif, lokasi
pesisir-pedalaman, dan interaksi antarkota. Kesimpulan ini tidak
menghapus perbedaan antarspesifikasi atau keterbatasan pengukuran yang
dicatat sumber.

\subsection{Findings}\label{findings-59}

\textbf{{[}Temuan dilaporkan sumber{]}} Pada bagian diskusi, artikel
merangkum lima temuan utama: UR China yang lebih polycentric memiliki
TFP lebih rendah; pengaruh polycentricity terhadap TFP mulai positif
ketika ukuran populasi meningkat; kota dapat meminjam ukuran dari kota
di dekatnya ketika UR cukup besar; dampak negatif polycentricity lebih
kuat pada UR dengan posisi lebih rendah dalam hierarki administratif;
dan UR pesisir menunjukkan pola pengaruh yang berbeda secara nyata
dibanding UR pedalaman (PDF pp.~1029, Section 5; TXT lines 725-740).

\textbf{{[}Interpretasi penulis{]}} Penulis menilai bahwa agglomeration
economies, ketika diproksikan dengan TFP, relatif lebih sedikit hadir
pada Chinese URs yang lebih polycentric. Perbedaan dengan hasil yang
menggunakan LP ditafsirkan mungkin berkaitan dengan perbedaan jenis
keunggulan urban: faktor tidak berwujud seperti teknologi, inovasi, dan
kemampuan manajerial dapat bergantung pada peran kota pusat, sedangkan
faktor produksi berwujud dapat menghadapi kompetisi dalam aglomerasi
besar. Penulis secara eksplisit menyatakan bahwa gagasan penjelas ini
masih memerlukan pemeriksaan lebih lanjut (PDF pp.~1029, Section 5; TXT
lines 742-766).

\textbf{{[}Interpretasi penulis{]}} Penulis semula mengharapkan
polycentricity lebih menguntungkan UR kecil, tetapi menemukan pola
sebaliknya. Mereka mengaitkan perbedaan tersebut dengan skala teritorial
dan populasi UR China, jarak antarkota, konektivitas, serta konteks
kelembagaan China. Polycentric pattern yang dibentuk secara top-down
tanpa kolaborasi bottom-up yang cukup dinilai dapat menimbulkan
kehilangan efisiensi (PDF pp.~1029, Section 5; TXT lines 767-770 and
712-724).

\textbf{{[}Inferensi pengolahan{]}} Temuan B16 lebih tepat dipakai
sebagai bukti bersyarat tentang kapan struktur polycentric berkaitan
dengan produktivitas dalam kasus yang diteliti, bukan sebagai peringkat
umum antara struktur monocentric dan polycentric.

\subsection{Conclusions}\label{conclusions-59}

\textbf{{[}Interpretasi penulis{]}} Analisis tidak menghasilkan bukti
bahwa pembangunan urban polycentric terkait dengan produktivitas ekonomi
regional yang lebih tinggi di China secara umum. Penulis menyerukan
kehati-hatian dalam menetapkan manfaat ekonomi regional sebagai alasan
untuk merencanakan pola polycentric secara universal (PDF p.~1029,
Section 5; TXT lines 753-770).

\textbf{{[}Interpretasi penulis{]}} Untuk UR terbesar seperti Yangtze
River Delta dan Pearl River Delta, pemerintah lokal disarankan
mempertimbangkan konektivitas antarkota dan jaringan jaminan sosial
terintegrasi untuk memfasilitasi mobilitas tenaga kerja. Untuk UR yang
lebih kecil seperti Ningxia-Yellow River dan Central Shanxi, penulis
menyatakan bahwa agglomeration economies kota pusat belum sepenuhnya
dimanfaatkan dan pembangunan multi-pusat dapat menurunkan produktivitas.
Penulis juga menyarankan mempertahankan kesenjangan ukuran antarkota
yang moderat agar peran kota pusat tidak melemah karena polycentricity
berlebihan (PDF pp.~1029-1030, Section 5; TXT lines 760-780).

\subsection{Contributions}\label{contributions-59}

\textbf{{[}Temuan dilaporkan sumber{]}} Artikel menyatakan tiga
kontribusi utama. Pertama, artikel menggunakan TFP untuk mengukur
kinerja ekonomi regional, yang diposisikan sebagai ukuran yang lebih
komprehensif daripada LP karena mempertimbangkan lebih dari input tenaga
kerja. Kedua, artikel membangun kerangka longitudinal yang menghubungkan
pembentukan PUR dan dampak ekonomi melintasi ruang dan waktu pada UR
China. Ketiga, artikel menerapkan time-varying IV yang memanfaatkan
bentuk kota dan pertumbuhan kota yang diprediksi untuk menangani
endogeneity (PDF pp.~1019 and 1028-1029, Sections 1 and 5; TXT lines
190-210 and 712-724).

\textbf{{[}Temuan dilaporkan sumber{]}} Sumber juga menyatakan bahwa
metodologinya menggabungkan unsur dari studi tentang pengukuran
polycentricity, pengukuran produktivitas dengan TFP, kinerja PUR China
berbasis data cross-sectional, dan IV bentuk kota berbasis hambatan
geografis (PDF p.~1029, Section 5; TXT lines 712-724).

\subsection{Research Gap}\label{research-gap-59}

\textbf{{[}Temuan dilaporkan sumber{]}} Gap yang hendak diisi adalah
kurangnya bukti longitudinal yang lebih robust tentang hubungan struktur
polycentric dan kinerja ekonomi, terutama dengan perhatian yang memadai
pada pilihan indikator ekonomi dan potensi endogeneity. Artikel juga
menempatkan kurangnya analisis komparatif lintas UR China sebagai
konteks bagi penggunaan 19 UR, karena banyak studi sebelumnya berfokus
pada satu UR atau sejumlah kecil UR (PDF pp.~1018-1019 and 1021,
Sections 1 and 2.3; TXT lines 151-189 and 284-301).

\textbf{{[}Temuan dilaporkan sumber{]}} Gap yang masih tersisa menurut
sumber meliputi sensitivitas hasil terhadap ukuran efisiensi ekonomi,
basis pengetahuan TFP yang masih terbatas dalam analisis struktur ruang,
kebutuhan untuk memisahkan borrowed size dari agglomeration shadow,
serta kebutuhan akan analisis fungsional dan multiskalar (PDF pp.~1030,
Notes and Section 5; TXT lines 820-874).

\textbf{{[}Inferensi pengolahan{]}} Dengan demikian, B16 mengisi
sebagian gap desain empiris dan pengendalian endogeneity, tetapi tidak
menyelesaikan persoalan identifikasi efek individual borrowed size
versus agglomeration shadow atau persoalan skala analisis.

\subsection{Limitations}\label{limitations-59}

\textbf{{[}Temuan dilaporkan sumber{]}} Keterbatasan yang disebut atau
diakui dalam file meliputi:

\begin{itemize}
\tightlist
\item
  BOR yang digunakan adalah indikator terpadu sehingga borrowed size dan
  agglomeration shadow terkolaps menjadi satu istilah; akibatnya efek
  masing-masing tidak dapat dinilai secara terpisah.
\item
  Proksi interaksi antarkota dibangun dari populasi dan jarak. Sumber
  menyatakan bahwa data lain akan lebih tepat jika fokus risetnya adalah
  fungsi tingkat tinggi di kota tetangga.
\item
  Analisis hanya berfokus pada skala regional karena di situlah
  inisiatif kebijakan PUR di China beroperasi; sumber menyatakan
  perspektif multiskalar dapat memberi pemahaman lebih baik.
\item
  Penggunaan projected historical population growth dalam IV dapat
  terpengaruh tren spesifik kota yang berasal dari kondisi awal 1992.
  Penulis menyatakan prosedur pengendalian tersebut hanya menangani
  kekhawatiran itu secara sebagian.
\item
  Ambang identifikasi population centres disesuaikan dengan konteks
  China. Sumber mengakui bahwa kota baru yang direncanakan dapat
  memiliki kepadatan penduduk dan pengaruh sekitar yang modest sehingga
  tersaring.
\item
  Data empat kota hilang pada beberapa tahun dan ditangani dengan linear
  interpolation; 12 autonomous prefectures harus dilengkapi dari
  Statistical Bulletins karena tidak tercatat dalam yearbook yang
  digunakan.
\end{itemize}

Locator: PDF pp.~1021-1022, 1024-1026, and 1030; TXT lines 333-356,
482-550, and 787-874.

\textbf{{[}Inferensi pengolahan{]}} Karena supplemental data, Table C1,
Table E1, Table F1-F2, dan Table G1 hanya dirujuk tetapi tidak hadir
sebagai isi penuh dalam TXT, review ini tidak memverifikasi ulang
statistik rinci, konstruksi modal, efektivitas setiap IV, atau hasil uji
tambahan di luar uraian yang tersedia.

\subsection{Challenges}\label{challenges-59}

\textbf{{[}Temuan dilaporkan sumber{]}} Tantangan empiris yang
dijelaskan sumber mencakup ketidaklengkapan data, perbedaan definisi
statistik, kebutuhan menggabungkan sumber data pada tingkat prefektur,
serta kebutuhan melakukan interpolasi. Tantangan operasional lainnya
adalah menentukan population centres dalam kondisi kota China yang
memiliki kota baru terencana, kepadatan sekitar rendah, dan
infrastruktur lemah (PDF pp.~1021-1023, Sections 3.1-3.3.1 and Note 3;
TXT lines 333-378 and 799-837).

\textbf{{[}Temuan dilaporkan sumber{]}} Tantangan identifikasi kausal
muncul karena struktur spasial dapat merupakan akibat, bukan penyebab,
dari kinerja ekonomi. Artikel meresponsnya dengan IV time-varying
berbasis proyeksi bentuk dan pertumbuhan kota, tetapi pembangunan
instrumen itu memerlukan asumsi tentang ekspansi kota melingkar tanpa
kendala geografis dan penggunaan data historis (PDF pp.~1024-1026,
Section 3.3.2 and Note 6; TXT lines 482-550 and 858-873).

\textbf{{[}Interpretasi penulis{]}} Dalam konteks kebijakan, koordinasi
regional juga sulit karena proses top-down dan bottom-up berinteraksi,
pemerintah lokal dapat memprioritaskan kepentingannya sendiri, dan kerja
sama antarkota disebut masih terutama terbatas pada berbagi
infrastruktur transportasi (PDF pp.~1021 and 1029, Sections 2.3 and 5;
TXT lines 297-317 and 735-752).

\subsection{Practical Implications}\label{practical-implications-59}

\textbf{{[}Interpretasi penulis{]}} Hasilnya menyiratkan bahwa pembuat
kebijakan sebaiknya tidak menganggap polycentric development sebagai
resep yang menghasilkan pertumbuhan ekonomi regional pada semua UR.
Pertimbangan ukuran populasi, hierarki administratif, lokasi
pesisir-pedalaman, dan kapasitas kerja sama antarkota perlu mendahului
klaim manfaat ekonomi (PDF pp.~1026-1029, Sections 4-5; TXT lines
552-579 and 742-770).

\textbf{{[}Interpretasi penulis{]}} Pada UR besar, penguatan
konektivitas antarkota dan jaringan jaminan sosial terintegrasi
diposisikan sebagai cara untuk memperoleh manfaat bersih borrowed size
dan memfasilitasi mobilitas tenaga kerja. Pada UR kecil, memperkuat
manfaat agglomeration economies kota pusat dapat lebih sesuai daripada
memaksakan multi-centre development (PDF pp.~1029-1030, Section 5; TXT
lines 760-780).

\textbf{{[}Inferensi pengolahan{]}} Implikasi praktis yang dapat ditarik
dari file terbatas pada penyesuaian kebijakan terhadap kondisi UR yang
diteliti; file tidak menyediakan aturan keputusan atau ambang kebijakan
yang dapat digeneralisasi di luar konteks analisis.

\subsection{Applications}\label{applications-59}

\textbf{{[}Temuan dilaporkan sumber{]}} Dalam studi ini, kerangka
tersebut diterapkan untuk mengukur polycentricity morfologis, TFP, dan
net effect interaksi populasi antarkota pada 19 UR China, lalu
membandingkannya menurut hierarki administratif dan lokasi
pesisir-pedalaman (PDF pp.~1021-1028, Sections 3-4; TXT lines 321-705).

\textbf{{[}Temuan dilaporkan sumber{]}} Hasil penelitian digunakan untuk
merefleksikan strategi pembangunan regional dan perkotaan China,
khususnya pilihan konektivitas dan koordinasi antarkota pada UR besar
serta kehati-hatian terhadap multi-centre development pada UR kecil (PDF
pp.~1018 and 1029-1030; TXT lines 95-101 and 760-780).

\textbf{{[}Informasi cakupan{]}} Aplikasi di luar studi ini: Tidak
disebutkan dalam file.

\subsection{Future Research
(Optional)}\label{future-research-optional-59}

\textbf{{[}Saran eksplisit sumber{]}} Artikel menyebut lima arah
penelitian lanjutan:

\begin{itemize}
\tightlist
\item
  Memeriksa sensitivitas dampak polycentricity terhadap pilihan ukuran
  efisiensi ekonomi.
\item
  Menguji kesesuaian TFP dalam konteks empiris yang lebih beragam,
  memperpanjang kerangka longitudinal, dan memperbaikinya dengan data
  sektoral.
\item
  Menggunakan perspektif fungsional untuk mengidentifikasi struktur
  ruang polycentric, terutama karena konstruk seperti TFP dapat
  berkaitan dengan sifat dan kekuatan hubungan antarkota.
\item
  Memisahkan efek borrowed size dan agglomeration shadow agar kedua
  kekuatan tersebut dapat dipahami dengan lebih terperinci.
\item
  Mengembangkan perspektif multiskalar yang mengintegrasikan beberapa
  skala untuk memahami interaksi produktivitas-struktur
  ruang-polycentricity dan memilih skala intervensi kebijakan yang
  sesuai.
\end{itemize}

Locator: PDF pp.~1030, Section 5; TXT lines 820-874.

\subsection{Provenance Notes}\label{provenance-notes-59}

\textbf{\hyperref[identitas-sumber]{Identitas sumber}} Review ini dibuat
dari file TXT aktual dengan prefiks B16:
\texttt{source/journal/B16-yang-caset-derudder-2024-does-urban-polycentricity-contribute-to-regional-economic-growth-empirical-evidence-from-a-panel-of-chinese-urban-regions-10.1080-00343404.2023.2255623.txt}.
File tersebut menunjuk sumber PDF
\texttt{journal/B16-yang-caset-derudder-2024-does-urban-polycentricity-contribute-to-regional-economic-growth-empirical-evidence-from-a-panel-of-chinese-urban-regions-10.1080-00343404.2023.2255623.pdf}
(TXT lines 1-3).

\textbf{{[}Metadata ekstraksi{]}} Blok metadata menyatakan ekstraksi
dilakukan pada 2026-08-31 dengan \texttt{pdftotext\ -layout}. Metadata
PDF mencatat 16 halaman, PDF versi 1.5, tidak terenkripsi, ukuran
halaman 634.32 x 833.04 points, serta ukuran file 2.039.910 bytes (TXT
lines 1-25). Metadata bibliografis dan isi artikel menyebut Yuting Yang,
Freke Caset, dan Ben Derudder; \emph{Regional Studies}, 2024, volume 58
nomor 5, halaman 1018-1032; DOI 10.1080/00343404.2023.2255623 (TXT lines
27-47 and 83-93).

\textbf{{[}Cakupan pembacaan{]}} Seluruh 1.031 baris TXT dibaca,
termasuk blok \texttt{{[}PDF\ Metadata{]}}, teks artikel, Table 1, Table
2, catatan, dan daftar referensi. Locator dalam review menggunakan nomor
halaman tercetak artikel dan rentang baris TXT hasil ekstraksi. Tidak
ada informasi eksternal yang ditambahkan; supplemental data yang hanya
ditautkan dalam artikel tidak dibaca.

\textbf{{[}Batasan provenance{]}} Ekstraksi \texttt{-layout}
mempertahankan artefak pemenggalan halaman dan tata letak kolom. Pada
Table 2, beberapa angka dan tanda statistik hasil ekstraksi tidak mudah
direkonstruksi dari susunan kolom dan tidak seluruhnya selaras dengan
uraian naratif Section 4.2. Review mempertahankan uraian sumber,
menandai keterbatasan tersebut, dan tidak mengarang atau menghitung
ulang koefisien. Table C1, E1, F1-F2, dan G1 hanya tersedia sebagai
rujukan dalam TXT, bukan sebagai isi tabel lengkap.

\textbf{{[}Status pengolahan{]}} Sumber B16 berhasil dibaca dan review
ini memuat seluruh bagian wajib. Sumber TXT dan PDF tidak diubah. Tidak
ada kode lain yang diproses.

\section{Review B17}\label{review-b17}

\subsection{Summarized Abstract}\label{summarized-abstract-60}

\textbf{Temuan sumber:} Artikel menguji kepentingan relatif
agglomeration externalities dan network externalities dalam menjelaskan
produktivitas regional melalui kerangka pembangunan polisentris di
Tiongkok. Dengan data investasi perusahaan, model ekonometrika spasial,
dan strategi instrumental variables (IV), penulis melaporkan bahwa kedua
jenis eksternalitas penting bagi produktivitas regional. Tidak ditemukan
efek interaksi antara keduanya. Network externalities lebih menonjol
pada urban regions (URs) dengan tingkat polisentrisitas lebih tinggi
dan, berbeda dari agglomeration externalities yang terikat secara
geografis, dapat menghasilkan spillover spasial. (Abstract, hlm. 175)

\textbf{Interpretasi penulis:} Temuan abstrak diposisikan sebagai
dukungan bagi pentingnya dua logika sekaligus, bukan pilihan eksklusif
antara aglomerasi dan jaringan. (Abstract, hlm. 175; Section 1, hlm.
176)

\textbf{Inferensi pengolahan:} Istilah ``lebih menonjol'' dipertahankan
sebagai klaim komparatif penulis; TXT tidak menyediakan dasar untuk
menyetarakan besaran kedua eksternalitas di luar spesifikasi model yang
digunakan.

\subsection{Problem Statement}\label{problem-statement-60}

\textbf{Temuan sumber:} Artikel berangkat dari ketegangan antara
pendekatan teritorial, yang menekankan manfaat ekonomi dari ko-lokasi
agen, dan pendekatan relasional, yang menekankan manfaat hubungan
fungsional lintas ruang. Analisis komparatif tentang seberapa penting
kedua logika tersebut masih dinyatakan belum konklusif. Pertanyaan
dasarnya adalah apakah unit geografis berupa kota, wilayah, atau
aglomerasi merupakan unit utama untuk memahami pembangunan ekonomi, atau
apakah perhatian seharusnya diarahkan pada interaksi antarsatuan
tersebut melalui jaringan. Perdebatan ini juga berimplikasi pada tata
kelola dan cara mengakomodasi urbanisasi masa depan. (Section 1, hlm.
175--176)

\textbf{Temuan sumber:} Dalam konteks UR Tiongkok, penelitian sebelumnya
umumnya berfokus pada dampak ekonomi tingkat kota. Studi yang ada belum
secara eksplisit menghubungkan agglomeration externalities dan network
externalities dalam satu kerangka analitis, belum menguji bagaimana
keduanya tersusun pada UR dengan struktur spasial yang berbeda, dan
belum menangani potensi endogenitas karena keterbatasan data untuk
menguji eksogenitas hubungan tersebut. (Section 1, hlm. 176--177)

\textbf{Inferensi pengolahan:} Masalah penelitian B17 dapat diringkas
sebagai kebutuhan untuk memisahkan, membandingkan, dan menguji
kemungkinan hubungan antara manfaat konsentrasi dan manfaat konektivitas
pada unit analisis UR, dengan perhatian pada struktur polisentris dan
endogenitas.

\subsection{Objectives}\label{objectives-60}

\textbf{Temuan sumber:} Tujuan utama artikel adalah mengurai efek
agglomeration externalities dan network externalities serta potensi
interaksinya dalam konteks UR Tiongkok, khususnya dalam perkembangan
polisentris. Pertanyaan penelitiannya adalah: (1) sejauh mana network
externalities dan agglomeration externalities menjelaskan produktivitas
ekonomi regional di UR Tiongkok; dan (2) bagaimana efek tersebut berbeda
menurut tingkat polisentrisitas. (Section 1, hlm. 176--177)

\textbf{Temuan sumber:} Enam hipotesis yang diuji adalah: H1,
agglomeration externalities berdampak positif pada produktivitas
regional; H2, dampaknya lebih menonjol pada UR yang lebih polisentris;
H3, network externalities berhubungan positif dengan produktivitas
regional; H4, network externalities lebih besar pada UR yang lebih
polisentris; H5, pengaruh positif konektivitas jaringan lebih kuat
ketika jarak rata-rata antarkota lebih kecil; dan H6, terdapat efek
interaksi network externalities dan agglomeration externalities terhadap
produktivitas regional. (Section 2.1, hlm. 177--178)

\textbf{Interpretasi penulis:} Kerangka tujuan tersebut dimaksudkan
untuk menilai apakah penguatan jaringan antarkota atau pertumbuhan kota
besar dan padat lebih bermanfaat, tanpa mengasumsikan sejak awal bahwa
salah satu logika harus menggantikan yang lain. (Section 1, hlm. 176;
Section 2.1, hlm. 177--178)

\subsection{Literature Survey}\label{literature-survey-60}

\textbf{Temuan sumber:} Literatur yang ditinjau membedakan agglomeration
externalities sebagai manfaat ekonomi dari kedekatan geografis dan
ko-lokasi, yang bekerja melalui sharing, matching, dan learning, dari
network externalities sebagai manfaat ekonomi dari hubungan fungsional
lintas ruang. Artikel juga mengaitkan network externalities dengan
borrowed size sebagai spillover positif dan agglomeration shadows
sebagai spillover negatif dalam jaringan kota. (Section 1, hlm.
175--176; Section 2.1, hlm. 177)

\textbf{Temuan sumber:} Tinjauan sebelumnya menunjukkan bahwa
agglomeration externalities dapat melampaui inti kota dan
``regionalize'' dalam sistem polisentris. Network externalities
diperkirakan menguat melalui jaringan perkotaan yang erat, dan
perkembangan polisentris dapat memperkuatnya. Namun, pengaruh jarak
terhadap jaringan tidak seragam: teori dapat menganggap jarak kurang
penting, sedangkan sejumlah studi tentang jaringan inovasi, ekonomi,
transportasi, dan virtual di Tiongkok menemukan penurunan konektivitas
seiring jarak. Literatur juga memunculkan dua kemungkinan hubungan:
aglomerasi dapat membantu menghasilkan jaringan melalui pengurangan
biaya transaksi/transportasi, sedangkan jaringan dapat menggantikan
sebagian manfaat aglomerasi melalui berbagi sumber daya dan pengetahuan.
(Section 2.1, hlm. 177--178)

\textbf{Temuan sumber:} Dalam kerangka analitis terdahulu, jaringan
diukur melalui pendekatan fisik seperti penerbangan, jalan, dan kereta
api, atau pendekatan imaterial seperti jaringan politik, kerja sama,
inovasi, hubungan antarperusahaan, dan hubungan intraperusahaan.
Aglomerasi diukur dengan konsentrasi penduduk total atau ukuran berbasis
kepadatan. Keuntungan ekonomi dianalisis melalui regresi nonspasial,
model count, atau model ekonometrika spasial. Endogenitas telah menjadi
kekhawatiran berulang; IV lebih umum dalam model aglomerasi, sedangkan
pembentukan IV untuk jaringan lebih sulit karena data arus historis
sering tidak tersedia. (Section 2.2, hlm. 178--179)

\textbf{Interpretasi penulis:} Artikel menempatkan enterprise investment
sebagai proksi hubungan antarkota agar kedua logika eksternalitas dapat
dianalisis dalam satu kerangka, sekaligus menguji spillover spasial dan
endogenitas yang belum ditangani oleh sebagian penelitian sebelumnya.
(Section 2.2, hlm. 178--179; Section 1, hlm. 176--177)

\subsection{Method Used}\label{method-used-60}

\textbf{Temuan sumber:} Studi menganalisis 19 Chinese urban regions,
yang didefinisikan sebagai wilayah administratif dalam National Plan
ke-14 dan dianggap memiliki aktivitas ekonomi antarkota yang
terintegrasi. Polisentrisitas morfologis diukur melalui keseimbangan
distribusi ukuran kota menggunakan ukuran berbasis deviasi standar:

\texttt{Poly\ =\ 1\ −\ s\_obs\ /\ s\_max}

\texttt{Poly} bernilai 0 untuk tidak ada polisentrisitas dan 1 untuk UR
ideal-tipikal yang semua kotanya sama besar. Kota atau urban centers
diidentifikasi menggunakan dataset populasi LandScanTM; catatan 1
menjelaskan ambang persentil ke-95, pengelompokan grid yang
bersebelahan, serta syarat lebih dari 100.000 penduduk dan sedikitnya 3
km². (Section 3.1--3.2, hlm. 179--180; Note 1, hlm. 191)

\textbf{Temuan sumber:} Aglomerasi dioperasionalkan sebagai kepadatan
penduduk UR. Jaringan perkotaan dibangun dari matriks kota-ke-kota
berdasarkan arus investasi perusahaan dalam UR yang sama. Tiga indikator
jaringan digunakan: weighted degree dan weighted clustering coefficient
pada tingkat kota, serta global weighted clustering coefficient pada
tingkat UR. Ukuran jaringan regional \texttt{Gr} merupakan nilai
rata-rata clustering coefficient kota-kota dalam UR. (Section 3.3, hlm.
181)

\textbf{Temuan sumber:} Produktivitas regional diproksikan dengan labour
productivity (LP), yaitu GDP sektor sekunder dan tersier dibagi ukuran
angkatan kerja. Setelah Moran's I menunjukkan efek spasial yang
signifikan, penulis menggunakan spatial autoregressive model (SAR)
dengan spatial fixed effects dan time fixed effects. Matriks bobot
spasial menggunakan kebalikan jarak geografis rata-rata antarkota pada
UR yang berbeda, dihitung dari average driving distance. Model
memasukkan interaksi \texttt{AGG*POLY}, \texttt{NET*POLY},
\texttt{NET*DIS}, dan \texttt{NET*AGG}. (Section 3.3--3.4.1, hlm.
181--182; Note 3--5, hlm. 191)

\textbf{Temuan sumber:} Variabel kontrol mencakup capital stock per
tenaga kerja (\texttt{CLR}), human capital (\texttt{HUM}), information
level (\texttt{INF}), industrial level (\texttt{TIN}), foreign direct
investment (\texttt{FDI}), dan government intervention (\texttt{GOV}).
Untuk menangani endogenitas, prosedur IV menggunakan variabel historis
untuk \texttt{AGG} dan \texttt{NET}, serta predicted polycentricity
untuk \texttt{POLY}: kepadatan penduduk 1978 dan 1992 (\texttt{AGG\_H}),
kepadatan jaringan historis 1978 dan 1992 (\texttt{NET\_H}), serta
derajat polisentrisitas terprediksi 2010 dan 2019 (\texttt{POLY\_P}).
Estimasi IV dijelaskan sebagai prosedur dua tahap. (Section
3.4.1--3.4.2, hlm. 181--182; Note 6, hlm. 191)

\textbf{Temuan sumber:} Pemeriksaan tambahan meliputi uji korelasi,
tail-dependence, dan Breusch--Pagan untuk memvalidasi efek interaksi,
serta pemeriksaan kausalitas dengan IV. Robustness check juga mengganti
ukuran \texttt{POLY\_S} berbasis deviasi standar dengan \texttt{POLY\_P}
berbasis primacy, \texttt{POLY\_H} berbasis Herfindahl, dan
\texttt{POLY\_R} berbasis rank-size. (Note 7, hlm. 191--192; Section
4.3, hlm. 187--188)

\subsection{Dataset}\label{dataset-60}

\textbf{Temuan sumber:} Data jaringan perusahaan berasal dari Qichacha.
Penulis mengekstrak struktur kepemilikan seluruh perusahaan terdaftar,
yaitu distribusi controlling interest di pasar saham, untuk
menggambarkan hubungan investasi antarkota yang terarah. Per Juni 2020,
basis data yang dibangun memuat 249 juta catatan perusahaan di 226 kota
terpilih untuk periode 1978--2019. Catatan tersebut diagregasikan
menjadi jaringan untuk 1978 (54 links), 1992 (648 links), 2001 (1.511
links), 2010 (2.195 links), dan 2019 (2.709 links). Data populasi untuk
mengidentifikasi urban centers berasal dari LandScanTM. Sumber dan
statistik deskriptif seluruh variabel disebut tersedia di Appendix B
pada supplemental data online, tetapi isi appendix tersebut tidak
terdapat dalam TXT. (Section 3.2, hlm. 179--180; Section 3.5, hlm. 183;
Appendix B dirujuk pada hlm. 183)

\textbf{Interpretasi penulis:} Lima titik waktu jaringan dianggap cukup
untuk menggambarkan dinamika struktur spasial regional, meskipun penulis
menyatakan panel tersebut ``not ideal''. (Section 3.5, hlm. 183)

\subsection{Population / Sample}\label{population-sample-60}

\textbf{Temuan sumber:} Populasi analisis adalah 19 Chinese urban
regions (URs), yaitu wilayah administratif yang digariskan dalam
National Plan ke-14 dan dianggap memiliki integrasi aktivitas ekonomi
antarkota. Pada 2019, 19 UR tersebut secara kolektif mencakup 70\%
prefecture-level municipalities di Tiongkok, 29\% sumber daya lahan,
73\% populasi total, dan sekitar 80\% GDP. Unit kota yang digunakan
untuk jaringan mencakup 226 kota. Tabel SAR memuat 57 observasi,
sedangkan estimasi IV memuat 38 observasi karena menggunakan tahun 2010
dan 2019. (Section 3.1, hlm. 179; Section 3.5, hlm. 183; Table 1, hlm.
185; Table 3, hlm. 188; Note 6, hlm. 191)

\textbf{Inferensi pengolahan:} Ini merupakan analisis panel wilayah
administratif yang dipilih berdasarkan kerangka National Plan ke-14,
bukan sampel survei individual. Kriteria pemilihan di luar uraian
tersebut tidak disebutkan dalam file.

\subsection{Dependent Variables}\label{dependent-variables-60}

\textbf{Temuan sumber:} Variabel dependen substantif adalah labour
productivity (\texttt{LP}) pada UR dan tahun tertentu. \texttt{LP}
dioperasionalkan sebagai GDP sektor sekunder dan tersier dibagi ukuran
angkatan kerja. Dalam model SAR, \texttt{LP} menjadi variabel dependen
dan produktivitas UR tetangga masuk melalui komponen spatial
autoregressive. (Section 3.3--3.4.1, hlm. 181--182)

\textbf{Temuan sumber:} Pada tahap pertama prosedur IV, teks menyatakan
bahwa variabel endogen diregresikan sebagai variabel dependen untuk
menghasilkan nilai prediksi yang kemudian digunakan dalam tahap
berikutnya. Ini merupakan tahap instrumentasi; outcome utama penelitian
tetap \texttt{LP}. (Section 3.4.2, hlm. 182)

\subsection{Results}\label{results-60}

\textbf{Temuan sumber:} Deskripsi jaringan pada Figure 2 menunjukkan
bahwa jaringan urban terkuat terkonsentrasi di pesisir, terutama
Beijing--Tianjin--Hebei (BTH) dan Pearl River Delta (PRD), disusul
Yangtze River Delta (YRD) dan Shandong Peninsula (SDP). Semua UR
mengalami kenaikan konektivitas, tetapi tidak merata; Yangtze River
Middle-Reach (YRM) dan Chengdu--Chongqing (CHC) menjadi lebih terhubung,
sedangkan UR di timur laut dan barat daya tetap relatif kurang
terhubung. UR dengan konektivitas antarperkotaan lebih kuat umumnya juga
lebih polisentris, dengan polisentrisitas terkuat di Tiongkok timur.
(Section 4.1, hlm. 183--184; Figure 2)

\textbf{Temuan sumber:} Dalam sembilan model SAR, koefisien autoregresif
spasial \texttt{r} signifikan dan berada pada kisaran 0,65--0,75. Nilai
R² model berada pada kisaran 0,67--0,73, dengan 57 observasi pada setiap
model. Pada model kontrol, \texttt{CLR}, \texttt{HUM}, dan \texttt{TIN}
berasosiasi positif dengan produktivitas, sedangkan \texttt{GOV}
berasosiasi negatif; koefisien \texttt{INF} dan \texttt{FDI} bertanda
negatif. (Section 4.2, hlm. 184--185; Table 1, hlm. 185)

\textbf{Interpretasi penulis:} Tanda negatif \texttt{INF} dikaitkan
dengan kemungkinan biaya regulasi dari investasi pemerintah pada layanan
komunikasi, sedangkan tanda negatif \texttt{FDI} dikaitkan dengan
kemungkinan crowding effect terhadap ekonomi lokal. Ini adalah
penjelasan penulis, bukan pengukuran mekanisme tambahan dalam model.
(Section 4.2, hlm. 184--185)

\textbf{Temuan sumber:} \texttt{POLY} berpengaruh positif dan signifikan
terhadap \texttt{LP} pada Model 2. Untuk agglomeration externalities,
koefisien \texttt{AGG} positif; penulis menyatakan bahwa kenaikan
kepadatan penduduk 10\% menghasilkan kenaikan produktivitas 2,6\%.
Interaksi \texttt{AGG*POLY} tidak signifikan, sehingga H1 diterima dan
H2 ditolak. (Section 4.2, hlm. 185; Table 1)

\textbf{Temuan sumber:} Koefisien \texttt{NET} dan \texttt{NET*POLY}
signifikan. Hasil ini dipakai untuk menyatakan bahwa network
externalities bekerja pada tingkat UR dan lebih menonjol pada UR dengan
\texttt{POLY} lebih tinggi; H3 dan H4 diterima. Interaksi
\texttt{NET*DIS} tidak memberikan bukti bahwa manfaat jaringan
bergantung pada jarak antarkota, sehingga H5 ditolak. (Section 4.2, hlm.
185; Table 1)

\textbf{Temuan sumber:} Pada Model 8, \texttt{NET*AGG} tidak signifikan
(koefisien −0,03 dengan standard error 0,05). Penulis tidak menemukan
efek interaksi yang mendukung komplementaritas jaringan dan aglomerasi,
dan membahas hasil tersebut sebagai tidak adanya substitusi jaringan
terhadap manfaat aglomerasi atau sebagai indikasi bahwa manfaat
aglomerasi tidak bergantung pada interaksi antarkota. Pada Model 9, yang
memasukkan seluruh interaksi sekaligus, \texttt{NET*AGG} menjadi
signifikan (−0,09 dengan standard error 0,04). Penulis mengaitkan
perubahan ini dengan kemungkinan confounding dari interaksi lain dan
menetapkan hasil Model 1--8 sebagai temuan final. (Section 4.2, hlm.
185--186; Table 1)

\textbf{Temuan sumber:} Dekomposisi efek spasial dari Model 8 menemukan
spillover spasial \texttt{POLY}. Efek tidak langsung \texttt{AGG} tidak
signifikan, sedangkan efek tidak langsung \texttt{NET} positif dan
signifikan. Karena itu, penulis menyatakan bahwa \texttt{AGG} terbatas
pada skala regional, sementara \texttt{NET} memiliki cakupan spasial
lebih luas. (Section 4.2, hlm. 186; Table 2, hlm. 187)

\textbf{Temuan sumber:} Estimasi IV tahap kedua mengonfirmasi pengaruh
positif \texttt{POLY}, \texttt{AGG}, dan \texttt{NET}, dengan ukuran
koefisien yang lebih besar tetapi sebanding dengan Table 1. Uji
relevansi instrumen dan identifikasi dilaporkan mendukung validitas
instrumen; uji Durbin--Wu--Hausman dan Wu--Hausman tidak signifikan pada
taraf 10\%. Penulis menyimpulkan bahwa \texttt{POLY}, \texttt{NET}, dan
\texttt{AGG} dapat diperlakukan sebagai eksogen dalam estimasi tersebut.
Estimasi IV menggunakan 38 observasi. (Section 4.3.1, hlm. 187--188;
Table 3, hlm. 188; Notes 8--9, hlm. 191--192)

\textbf{Temuan sumber:} Robustness check dengan ukuran polisentrisitas
alternatif menghasilkan koefisien positif untuk \texttt{POLY\_S} dan
\texttt{POLY\_H}, sedangkan \texttt{POLY\_P} dan \texttt{POLY\_R} tidak
signifikan. Penulis tetap menggunakan \texttt{POLY\_S} sebagai
kesimpulan karena ukuran tersebut paling sesuai dengan fokus pada
pemerataan distribusi penduduk dan berkorelasi kuat dengan
\texttt{POLY\_R}; TXT melaporkan koefisien korelasi 0,67 untuk keduanya.
(Section 4.3.2, hlm. 188; Note 10, hlm. 192)

\subsection{Findings}\label{findings-60}

\textbf{Temuan sumber:} Penulis merumuskan enam temuan utama berikut.
(Section 5, hlm. 189--190)

\begin{enumerate}
\def\labelenumi{\arabic{enumi}.}
\tightlist
\item
  Network externalities dan agglomeration externalities sama-sama
  berdampak positif pada produktivitas regional.
\item
  Tidak ada bukti bahwa network externalities melemah seiring jarak
  antarkota di dalam UR.
\item
  Network externalities tidak menggantikan agglomeration externalities.
\item
  Agglomeration externalities tetap konsisten pada UR dengan tingkat
  polisentrisitas yang berbeda.
\item
  Network externalities lebih menonjol pada UR yang lebih polisentris.
\item
  Dibandingkan agglomeration externalities, network externalities dapat
  menyebar pada skala yang lebih luas.
\end{enumerate}

\textbf{Interpretasi penulis:} Jawaban atas pertanyaan apakah jaringan
perlu diperkuat atau kota besar dan padat perlu dikembangkan tidak dapat
ditetapkan secara definitif hanya dari hasil statistik karena kedua
jenis eksternalitas dibangun secara analitis dengan cara yang berbeda.
Meskipun demikian, penulis menilai pentingnya network externalities di
UR Tiongkok meningkat karena konektivitas bertumbuh, spillover-nya
melampaui skala regional, dan UR yang lebih polisentris tampak lebih
kondusif bagi perkembangannya. (Section 5, hlm. 189--190)

\subsection{Conclusions}\label{conclusions-60}

\textbf{Interpretasi penulis:} Artikel menyimpulkan bahwa kedua logika,
yaitu manfaat kepadatan dan manfaat konektivitas, relevan bagi
produktivitas regional, bukan pilihan ``either-or''. Agglomeration
externalities konsisten pada tingkat polisentrisitas yang berbeda,
sedangkan network externalities lebih menonjol pada UR polisentris dan
memiliki spillover yang lebih luas. Network externalities tidak terbukti
menggantikan manfaat aglomerasi dan tidak menunjukkan distance-decay
dalam pengujian ini. (Abstract, hlm. 175; Section 5, hlm. 189--190)

\textbf{Interpretasi penulis:} Pada tingkat UR, network externalities
diperlakukan sekaligus sebagai fenomena lokal dan regional: kepadatan
jaringan dimasukkan ke regresi untuk menangkap pengaruh lokal, sedangkan
matriks jarak antarsesama UR digunakan untuk menguji pengaruh jaringan
satu UR terhadap produktivitas UR tetangga. Penulis menyatakan bahwa
perspektif ini melengkapi konseptualisasi network externalities sebagai
``club'' goods maupun ``public'' goods. (Section 5, hlm. 190)

\textbf{Inferensi pengolahan:} Kesimpulan artikel berlaku pada
operasionalisasi network berbasis investasi perusahaan, ukuran
aglomerasi berbasis kepadatan penduduk, dan unit 19 UR Tiongkok yang
diuji; TXT tidak mendukung generalisasi langsung ke semua jenis jaringan
atau wilayah lain.

\subsection{Contributions}\label{contributions-60}

\textbf{Temuan sumber:} Penulis menyatakan tiga kontribusi utama.
Pertama, artikel mengonsolidasikan network externalities dan
agglomeration externalities dalam satu kerangka. Kedua, artikel
menggunakan metode ekonometrika spasial untuk menguji apakah hubungan
jaringan menghasilkan spatial spillover. Ketiga, artikel menggunakan
strategi IV untuk memperoleh bukti yang disebut penulis lebih kuat dalam
memverifikasi efek kausal yang dihipotesiskan. (Section 5, hlm. 189)

\subsection{Research Gap}\label{research-gap-60}

\textbf{Temuan sumber:} Sebelum artikel ini, sebagian besar penelitian
tentang dampak ekonomi UR Tiongkok berfokus pada tingkat kota. Studi
terdahulu belum secara eksplisit menghubungkan agglomeration
externalities dan network externalities dalam satu kerangka analitis,
belum menguji bagaimana keduanya tersusun pada UR dengan struktur
spasial berbeda, dan belum menangani potensi endogenitas karena
keterbatasan data untuk menguji eksogenitas hubungan tersebut. (Section
1, hlm. 176--177)

\textbf{Temuan sumber:} Artikel juga menunjukkan gap yang masih tersisa:
hasil berbeda antarukuran polisentrisitas, keterbatasan pemahaman
tentang jenis jaringan yang dapat melengkapi atau menggantikan efek
aglomerasi, kemungkinan perbedaan pola pada skala spasial lain,
pentingnya kedekatan organisasional atau komunitas selain kedekatan
geografis, serta perlunya metode yang konsisten untuk mengoperasionalkan
kedua jenis eksternalitas. (Section 5, hlm. 190--191)

\subsection{Limitations}\label{limitations-60}

\textbf{Temuan sumber:} Penulis menyatakan bahwa panel lima titik waktu
yang dipakai untuk menggambarkan dinamika struktur spasial ``not
ideal'', walaupun dianggap cukup. Data interaksi pada tingkat antar-UR
juga terbatas; karena itu, jaringan dibangun dari hubungan investasi
perusahaan antarkota di dalam UR. (Section 3.5, hlm. 183; Section 5,
hlm. 190)

\textbf{Temuan sumber:} Pendekatan jaringan menyederhanakan realitas
urban dan regional. Catatan 2 menjelaskan bahwa county kecil dapat
berada dalam lingkup pengaruh kota besar, sehingga kota-kota tidak
selalu merupakan node yang sepenuhnya terpisah dan dapat menjadi bagian
dari aglomerasi yang lebih besar bersama node yang tidak dimasukkan.
Penulis tetap memakai pendekatan tersebut karena sesuai dengan gagasan
konseptual tentang keterhubungan kota. (Note 2, hlm. 191)

\textbf{Temuan sumber:} Kesimpulan tentang polisentrisitas sensitif
terhadap cara pengukuran. \texttt{POLY\_S} dan \texttt{POLY\_H}
menghasilkan hasil positif, tetapi \texttt{POLY\_P} dan \texttt{POLY\_R}
tidak signifikan. Penulis menilai perbedaan tersebut mencerminkan aspek
polisentrisitas yang ditangkap oleh masing-masing indeks, bukan otomatis
berarti temuan tidak robust. Perbandingan langsung agglomeration dan
network externalities juga dibatasi oleh konstruksi analitis kedua
variabel yang berbeda. (Section 4.3.2, hlm. 188; Section 5, hlm. 189;
Note 10, hlm. 192)

\textbf{Inferensi pengolahan:} Batasan sumber data, pilihan indeks,
skala UR, dan proksi enterprise investment membatasi ruang lingkup
generalisasi review ini. Tidak disebutkan dalam file adanya validasi
dengan jenis jaringan lain atau evaluasi dampak kebijakan langsung.

\subsection{Challenges}\label{challenges-60}

\textbf{Temuan sumber:} Tantangan metodologis utama adalah mengukur
network externalities pada tingkat inter-urban region. Data interaksi
antarsatuan regional yang terintegrasi lebih terbatas dibandingkan data
tingkat kota tentang transportasi, kolaborasi pengetahuan, dan arus
keuangan. Penulis mengatasinya dengan menggabungkan variabel kepadatan
jaringan dalam regresi dan matriks bobot spasial berbasis jarak
berkendara rata-rata antarkota di UR yang berbeda. (Section 5, hlm. 190)

\textbf{Temuan sumber:} Pembangunan IV untuk network externalities lebih
sulit karena membutuhkan data arus historis yang biasanya tidak
tersedia. Dalam studi ini, data jaringan historis yang lebih jarang
tersedia untuk 1978 dan 1992 digunakan bersama kepadatan penduduk
historis dan polisentrisitas terprediksi. (Section 2.2, hlm. 178--179;
Section 3.4.2, hlm. 182; Note 6, hlm. 191)

\textbf{Temuan sumber:} Penulis juga menghadapi tantangan dalam
menyimpulkan efek interaksi. Signifikansi statistik saja dinilai dapat
menyesatkan karena kemungkinan korelasi nonlinear, tail dependence, dan
heterogeneity; karena itu dilakukan uji korelasi, tail-dependence, dan
Breusch--Pagan. Sensitivitas indeks polisentrisitas terhadap pilihan
metodologis menjadi tantangan tambahan. (Note 7, hlm. 191--192; Section
4.3.2, hlm. 188)

\subsection{Practical Implications}\label{practical-implications-60}

\textbf{Interpretasi penulis:} Karena kedua jenis eksternalitas
berasosiasi positif dengan produktivitas, perdebatan perencanaan tidak
seharusnya diperlakukan sebagai pilihan eksklusif antara memperkuat
jaringan urban dan menumbuhkan kota yang besar serta padat. Namun,
penulis menegaskan bahwa hasil statistik saja tidak cukup untuk
menetapkan pilihan definitif karena konstruksi analitis kedua
eksternalitas berbeda. (Section 5, hlm. 189--190)

\textbf{Interpretasi penulis:} Untuk memahami dinamika urban, perspektif
jaringan perlu diintegrasikan dengan ukuran polisentrisitas tradisional.
Perspektif ini menempatkan posisi kota dalam arus antarkota orang,
informasi, dan barang sebagai pelengkap ukuran morfologis. (Section 5,
hlm. 190--191)

\textbf{Inferensi pengolahan:} Implikasi praktis yang paling aman dari
TXT adalah mengevaluasi kepadatan lokal, konektivitas antarkota, dan
kemungkinan spillover lintas-UR secara bersamaan; ini merupakan
inferensi dari desain dan hasil studi, bukan rekomendasi kebijakan
terukur yang dinyatakan langsung oleh penulis.

\subsection{Applications}\label{applications-60}

\textbf{Temuan sumber:} Artikel menerapkan kerangka gabungan
agglomeration/network externalities pada 19 UR Tiongkok untuk
menjelaskan \texttt{LP}. Network urban dibentuk dari investasi
perusahaan antarkota, lalu diuji bersama kepadatan penduduk,
polisentrisitas, jarak, dan efek spasial antarsesama UR. (Sections
3.1--3.4, hlm. 179--182)

\textbf{Temuan sumber:} Aplikasi empiris tersebut dipakai untuk menilai
apakah network externalities menghasilkan spillover di luar UR dan
apakah efeknya berbeda menurut polisentrisitas. Hasil yang dilaporkan
berasal dari data dan estimasi ekonometrika yang dijelaskan dalam TXT.
(Sections 3--4, hlm. 179--188)

\textbf{Inferensi pengolahan:} Tidak disebutkan dalam file adanya
penerapan kebijakan atau intervensi lapangan langsung.

\subsection{Future Research
(Optional)}\label{future-research-optional-60}

\textbf{Temuan sumber:} Penulis mengusulkan beberapa arah riset
lanjutan: menggabungkan perspektif jaringan dengan ukuran
polisentrisitas tradisional; memasukkan berbagai jenis jaringan, seperti
arus penumpang udara atau logistik, dan memakai data multisumber untuk
mengidentifikasi multiplex urban network; menguji kedekatan
organisasional dan komunitas selain kedekatan geografis; membandingkan
pola pada skala spasial yang berbeda karena arus investasi dapat
bergeser dari intra-region ke inter-region; serta menerapkan metode yang
konsisten untuk mengoperasionalkan agglomeration dan network
externalities. (Section 5, hlm. 190--191)

\subsection{Provenance Notes}\label{provenance-notes-60}

\textbf{Temuan sumber:} Satu TXT aktual berprefiks \texttt{B17}
ditemukan dan diverifikasi di
\texttt{source/journal/B17-yang-lu-caset-derudder-2025-agglomeration-externalities-or-network-externalities-explaining-productivity-in-chinese-urban-regions-10.1080-17421772.2024.2414960.txt}.
Blok \texttt{{[}PDF\ Metadata{]}} menyatakan PDF sumber, tanggal
ekstraksi 2026-08-31, metode \texttt{pdftotext\ -layout}, dan 23 halaman
(TXT, baris 1--26). Teks artikel memuat judul, penulis, jurnal, volume,
halaman 175--196, DOI, dan tanggal publikasi daring (TXT, baris 38--58,
84--110). Seluruh TXT dibaca, termasuk bagian metadata, artikel,
catatan, disclosure, funding, dan referensi hingga baris 1.227.

\textbf{Temuan sumber:} Artikel menyatakan tidak ada potential conflict
of interest yang dilaporkan. Sumber pendanaan yang dicantumkan adalah
Bijzonder Onderzoeksfonds UGent, China Scholarship Council, Research
Foundation--Flanders, KU Leuven research grant, dan National Science
Centre (Poland) funds. (Disclosure Statement dan Funding, hlm. 191; TXT,
baris 910--919)

\textbf{Inferensi pengolahan:} Review ini hanya menggunakan isi TXT B17
tersebut; tidak ada informasi dari luar file yang ditambahkan. Locator
\texttt{hlm.} merujuk pada halaman artikel/PDF yang tercetak di TXT,
sedangkan penanda \texttt{TXT,\ baris\ ...} merujuk pada posisi metadata
atau end matter dalam file ekstraksi. TXT sumber tidak diubah.

\section{Review B18}\label{review-b18}

\textbf{Sumber yang direview:} Chris Jacobs-Crisioni, Mert Kompil, dan
Lewis Dijkstra (2023), \emph{Big in the neighbourhood: Identifying local
and regional centres through their network position}, \emph{Papers in
Regional Science}, 102, 421-457, DOI 10.1111/pirs.12727.

\subsection{Summarized Abstract}\label{summarized-abstract-61}

\textbf{Laporan sumber:} Artikel menyatakan bahwa peringkat kota pada
tingkat nasional tidak cukup menggambarkan relevansi lokal dan regional
suatu settlement. Penulis memperkenalkan metode untuk memeringkat
centrality settlement berdasarkan waktu tempuh dan ukuran settlement,
kemudian menguji hubungan antara ukuran settlement, peran sebagai pusat
lokal atau regional, dan endowment layanan. Untuk EU27 dan UK,
pendekatan tersebut dilaporkan dapat menjelaskan endowment layanan per
kapita dengan baik dan menambahkan nuansa yang tidak ditangkap peringkat
nasional. Artikel juga mengidentifikasi preferensi skala spasial
beberapa jenis layanan. Hasilnya menyiratkan bahwa endowment layanan
dapat terkonsentrasi ulang karena koneksi jalan yang lebih cepat,
terutama di area yang lebih terkoneksi. {[}Abstract, pp.~421-422{]}

\textbf{Interpretasi penulis:} Penulis memosisikan metode tersebut
sebagai heuristic pragmatis untuk menangkap hierarki Christallerian
pusat-pusat layanan, bukan sebagai pengganti pengamatan lapangan atau
definisi functional urban area. {[}Section 1, pp.~422-423; Section 5,
pp.~445-446{]}

\subsection{Problem Statement}\label{problem-statement-61}

\textbf{Laporan sumber:} Sebagian besar kajian urbanisme dan geografi
ekonomi berfokus pada kota terbesar, padahal artikel melaporkan bahwa
63\% penduduk EU dan UK tidak tinggal di urban centre dan sekitar 38\%
penduduk EU tinggal di wilayah di luar functional urban areas.
Settlement kecil, town, dan village dapat menjadi hub layanan bagi
hinterland meskipun perannya dalam pasar tenaga kerja terbatas.
Penyediaan layanan sosial dan ekonomi menghadapi biaya lebih tinggi di
wilayah berpenduduk kecil dan tersebar, sementara informasi yang
komprehensif dan harmonis tentang layanan publik serta privat di seluruh
Eropa belum tersedia. {[}Section 1, pp.~421-422{]}

\textbf{Laporan sumber:} Metode berbasis ukuran nasional, distribusi
ukuran kota, jaringan kota global, atau ketergantungan pasar tenaga
kerja tidak memadai untuk menjelaskan fungsi penyedia layanan bagi
hinterland terdekat, khususnya pada settlement kecil. Data interaksi dan
point of interest yang tersedia juga memiliki keterbatasan cakupan,
variasi kualitas, serta potensi bias klasifikasi, commission, dan
omission. {[}Section 2, pp.~423-425{]}

\textbf{Interpretasi penulis:} Masalah inti yang hendak diatasi adalah
ketidakselarasan antara ukuran penduduk sebuah settlement dan peran
layanan yang sebenarnya dijalankannya dalam konteks jaringan tetangga
serta waktu tempuh.

\textbf{Inferensi pengolahan:} Unit persoalan artikel bukan sekadar
lokasi layanan individual, melainkan cara mengidentifikasi settlement
yang berfungsi sebagai pusat pada skala waktu tertentu ketika ukuran
nasional saja menyembunyikan peran lokal atau regional.

\subsection{Objectives}\label{objectives-61}

\textbf{Laporan sumber:} Artikel menetapkan tiga tujuan utama: (1)
memperkenalkan metode berbasis waktu tempuh jalan dan jumlah penduduk
settlement untuk menginferensikan klasifikasi hierarki centrality serta
menyajikan hasilnya untuk EU27 dan UK; (2) memvalidasi kegunaan hierarki
dengan menguji seberapa baik hierarki tersebut menjelaskan endowment
lokasi layanan; dan (3) memeriksa apakah perbedaan efisiensi jaringan
transportasi berhubungan dengan peringkat centrality dan endowment
layanan. Layanan yang dipertimbangkan mencakup layanan ekonomi dan
sosial untuk kepentingan umum. {[}Section 1, pp.~422-423{]}

\textbf{Interpretasi penulis:} Pendekatan dimaksudkan sebagai heuristic
yang hanya memerlukan ukuran settlement dan konektivitas jaringan jalan,
sehingga menurut penulis pada prinsipnya dapat diterapkan secara global,
walaupun validasi empiris artikel dilakukan pada EU27 dan UK. {[}Section
1, p.~423{]}

\subsection{Literature Survey}\label{literature-survey-61}

\textbf{Laporan sumber:} Artikel menempatkan kajian dalam Central Place
Theory (CPT) Christaller dan Losch. CPT dipakai untuk memahami sistem
berulang dari cluster retail dan layanan dengan tingkat kompleksitas,
spesialisasi, dan frekuensi spasial yang berbeda. Artikel membahas tiga
mekanisme: preferensi konsumen atas jarak yang lebih dekat versus
pilihan yang lebih kaya; keputusan penyedia layanan yang menimbang
kedekatan pengguna, pesaing, critical mass, economies of scale, dan
agglomeration benefits; serta pertimbangan equity dan akses dalam
layanan sosial. {[}Section 2, pp.~423-424{]}

\textbf{Laporan sumber:} Kajian terdahulu yang diringkas mencakup
hierarki berdasarkan ukuran kota nasional atau global, jaringan kota
kontinental, data commuting untuk mengidentifikasi ketergantungan
fungsional, pemodelan dinamis kemunculan hierarki, dan penggunaan ukuran
penduduk sebagai proksi ukuran pasar. Artikel menyatakan bahwa commuting
tidak sepenuhnya merepresentasikan seluruh fungsi layanan, data layanan
yang harmonis sulit diperoleh, model hierarki sektoral sangat kompleks,
dan ukuran pasar statis mengabaikan dinamika agglomeration serta biaya
perjalanan. Konteks lingkungan juga membuat settlement berukuran sama
dapat mempunyai peran berbeda; artikel menghubungkannya dengan gagasan
bahwa kota dapat meminjam ukuran dari kota yang terkoneksi. {[}Sections
2.1-2.4, pp.~423-425{]}

\textbf{Interpretasi penulis:} Literatur tersebut mendukung pemakaian
kombinasi ukuran lokal dan konteks jaringan, bukan peringkat nasional
tunggal, untuk membaca hierarki pusat layanan sampai ke settlement
terkecil. Metode artikel diposisikan sebagai pelengkap definisi berbasis
aglomerasi pasar tenaga kerja. {[}Sections 2.4 and 5, pp.~425 and
445-446{]}

\subsection{Method Used}\label{method-used-61}

\textbf{Laporan sumber:} Centrality dihitung dari ukuran relatif
settlement dalam radius waktu tempuh yang telah ditentukan. Untuk
settlement asal i dan ambang waktu t, nilai \texttt{L(t)\_i} menjadi 1
jika settlement i memiliki penduduk setidaknya sebesar setiap settlement
tujuan j yang dapat dicapai dalam \texttt{T\_ij\ \textless{}\ t}; jika
ada settlement yang lebih besar dalam radius tersebut, nilainya 0.
Ambang diuji dari 5 sampai 120 menit dengan kenaikan 5 menit. Hierarki
visual settlement ditentukan dari ambang tertinggi ketika kriteria
tersebut masih benar. {[}Section 3, pp.~425-426; Equation (1); Figures 1
and 3; Tables 1-3{]}

\textbf{Laporan sumber:} Settlement dibentuk dari cluster populasi
berbasis grid 1 x 1 km dan digeneralisasi menjadi polygon. Matriks waktu
tempuh dihitung terpisah untuk setiap negara, menggunakan topologi
jaringan dan kecepatan maksimum untuk navigasi mobil, dari centroid
berbobot penduduk ke centroid berbobot penduduk settlement lain yang
berada dalam jangkauan dua jam berkendara. Ukuran penduduk yang dipakai
dalam analisis berasal dari agregasi grid JRC-GeoStat 2018. {[}Section
3.1, pp.~426-427{]}

\textbf{Laporan sumber:} Lokasi layanan diatribusikan ke settlement
melalui jarak Euclidean ke grid cell terdekat yang berada di dalam
settlement. Lokasi yang berjarak lebih dari 2,5 km dari settlement
terdekat tidak dihitung. Endowment diuji dengan ANOVA dan perbandingan
Tukey pada tingkat kepercayaan 95\% menurut tipe settlement dan kelas
ukuran. {[}Sections 4 and 4.1-4.2, pp.~430-435; Tables 4-5 and Appendix
B{]}

\textbf{Laporan sumber:} Regresi OLS pada Equation (2) dijalankan
terpisah untuk tiap jenis layanan. Variabel terikatnya adalah
\texttt{ln(Y\_i\ +\ 0.01)}, dengan \texttt{Y\_i} jumlah lokasi layanan;
prediktornya mencakup dummy negara, \texttt{ln(P\_i\ +\ 0.01)}, dummy
settlement dalam 10 km dari perbatasan nasional, dummy jaringan
terisolasi, dan dummy \texttt{L(t)} pada beberapa ambang waktu secara
simultan. Equation (3) menambahkan normalized effective speed dan jumlah
destinasi yang dapat dicapai dalam 120 menit. Settlement dengan
effective speed sedikitnya 11\% di atas atau di bawah rata-rata nasional
juga dibandingkan. {[}Sections 4.3-4.4, pp.~436-445; Equations (2)-(3);
Tables 7-9{]}

\textbf{Laporan sumber:} Klasifikasi dihitung dengan GeoDMS dan
algoritme shortest-path; analisis statistik dilakukan dengan SPSS dan R.
{[}Section 1, p.~423{]}

\textbf{Inferensi pengolahan:} Desain ini merupakan analisis
observasional berbasis perbandingan spasial, bukan eksperimen atau
evaluasi before-after. \texttt{L(t)} adalah proksi centrality yang
diturunkan dari populasi dan waktu tempuh, bukan observasi langsung atas
pilihan konsumen atau keputusan setiap penyedia layanan.

\subsection{Dataset}\label{dataset-61}

\textbf{Laporan sumber:} Dataset settlement menggunakan grid populasi
sensus EU 2011 untuk delineasi derajat urbanisasi dan grid populasi
JRC-GeoStat 2018 untuk jumlah penduduk settlement. Settlement yang masuk
adalah cluster city, town, dan village yang terpisah secara fisik,
kontinu secara spasial, serta memenuhi ambang massa dan kepadatan
penduduk; area yang terlalu jarang atau tidak berpenghuni tidak
dimasukkan. Data jaringan jalan dan kecepatan berasal dari database
TomTom 2018. {[}Section 3.1, pp.~426-427{]}

\textbf{Laporan sumber:} Sepuluh jenis layanan dan jumlah lokasi yang
tersedia serta berhasil dihubungkan ke settlement adalah: retail 189.836
dan 164.459 (86,6\%); primary schools 150.633 dan 125.455 (83,3\%);
banks 89.546 dan 82.604 (92,2\%); pharmacies 86.789 dan 79.985 (92,2\%);
doctors 47.822 dan 39.872 (83,4\%); secondary schools 37.464 dan 35.168
(93,9\%); train stations 31.156 dan 23.524 (75,5\%); hospitals 10.032
dan 9.496 (94,7\%); cinemas 7.800 dan 7.375 (94,6\%); universities 3.958
dan 3.677 (92,9\%). Angka pertama adalah total lokasi di wilayah
analisis dan angka kedua adalah lokasi yang diperhitungkan setelah
linkage. {[}Table 4, p.~432{]}

\textbf{Laporan sumber:} Sumber utama lokasi layanan adalah
PROFECY/ESPON dan sebagian besar datanya berasal dari OpenStreetMap.
Data rumah sakit ditambah dari negara anggota melalui EUROSTAT-GISCO;
stasiun kereta berasal dari basis data Poelman dan Dijkstra yang
diperbarui; lokasi universitas berasal dari European Tertiary Education
Register. {[}Section 4, pp.~430-432{]}

\subsection{Population / Sample}\label{population-sample-61}

\textbf{Laporan sumber:} Populasi analisis mencakup 49.470 settlement di
EU27 dan UK. Tabel 2 melaporkan klasifikasi per negara; total settlement
yang menjadi pusat pada 15 menit adalah 8.796, pada 45 menit 840, dan
pada 90 menit 254. Pada Appendix B, komposisi tipe settlement yang
dianalisis adalah 823 urban centres, 8.817 towns, dan 39.790 villages,
dengan total 49.430 observasi untuk analisis layanan per 100.000
penduduk. {[}Table 2, p.~428; Table 3, p.~430; Appendix B, Table A1,
pp.~451-452{]}

\textbf{Laporan sumber:} Regresi Equation (2) menggunakan
\texttt{n\ =\ 49.470}; Equation (3) menggunakan \texttt{n\ =\ 49.459}
karena effective speed tidak dapat dihitung untuk seluruh settlement.
Analisis kelas ukuran hanya mempertahankan settlement dengan sedikitnya
25.000 penduduk dan membaginya menjadi 25.000-50.000, 50.000-100.000,
100.000-250.000, dan lebih dari 250.000 penduduk, dengan total
\texttt{n\ =\ 1.675}. Perbandingan effective speed memakai 7.269
settlement berkecepatan relatif rendah dan 6.647 berkecepatan relatif
tinggi. {[}Table 7, p.~439; Tables 8-9, pp.~443-444; Appendix B,
pp.~452-454{]}

\textbf{Inferensi pengolahan:} Cakupan yang dilaporkan lebih menyerupai
universe settlement yang berhasil diidentifikasi melalui prosedur
harmonisasi daripada sampel survei. Rancangan sampling probabilistik dan
tingkat respons: Tidak disebutkan dalam file.

\subsection{Dependent Variables}\label{dependent-variables-61}

\textbf{Laporan sumber:} Pada Equation (2) dan Equation (3), variabel
terikat \texttt{Y\_i} adalah jumlah lokasi dari satu jenis layanan yang
diatribusikan ke settlement i. Regresi dijalankan terpisah untuk retail,
primary schools, pharmacies, banks, train stations, doctors, secondary
schools, hospitals, cinemas, dan universities. Jumlah tersebut
ditransformasi menjadi \texttt{ln(Y\_i\ +\ 0.01)} untuk mengurangi
skewness dan menghindari log dari nol. {[}Sections 4.3-4.4, pp.~436-445;
Tables 7-8{]}

\textbf{Laporan sumber:} Dalam pengujian ANOVA, variabel terikat adalah
jumlah lokasi layanan per 100.000 penduduk. Tipe settlement dan kelas
ukuran dipakai sebagai kelompok pembanding. \texttt{L(t)} bukan variabel
terikat pada regresi; ia adalah variabel penjelas berbentuk dummy yang
merepresentasikan centrality pada beberapa ambang waktu. {[}Section 4.2,
pp.~434-435; Appendix B, Tables A1 and A3{]}

\textbf{Inferensi pengolahan:} Artikel tidak mengukur outcome
kesejahteraan, akses aktual pengguna, atau performa ekonomi secara
langsung. Outcome empiris utamanya adalah endowment atau jumlah lokasi
layanan yang teramati dalam data POI.

\subsection{Results}\label{results-61}

\textbf{Laporan sumber:} Klasifikasi menghasilkan 49.470 settlement.
Sebanyak 8.796 settlement diklasifikasikan sebagai centre pada ambang 15
menit, 840 pada 45 menit, dan 254 pada 90 menit, masing-masing sekitar
17,8\%, 1,7\%, dan 0,5\% dari total. Settlement berukuran besar lebih
sering menjadi settlement terbesar dalam neighbourhood. Artikel menyebut
lima settlement dengan penduduk kurang dari 1.000 tetap menjadi yang
terbesar dalam jangkauan 120 menit; semuanya merupakan komunitas pulau
yang terisolasi atau memiliki koneksi daratan yang lambat. {[}Section
3.2, pp.~427-430; Tables 2-3{]}

\textbf{Inferensi pengolahan:} Table 3 yang diekstrak menampilkan angka
6 pada baris \texttt{L(120)} untuk kelas \texttt{\textless{}1k},
sedangkan narasi Section 3.2 menyebut lima. Review mempertahankan kedua
informasi tersebut sebagai ketidakkonsistenan internal dan tidak memilih
salah satunya tanpa sumber tambahan.

\textbf{Laporan sumber:} Contoh Belgia memperlihatkan perbedaan antara
ukuran nasional dan posisi jaringan. Antwerp berada dalam 30 menit
berkendara dari Brussels yang jauh lebih besar, sedangkan Arlon yang
kecil mempunyai peran layanan hinterland yang diinferensikan lebih
besar. Dalam Table 6, nilai \texttt{Max\ L(t)} adalah 70 untuk Arlon
town cluster, 120 untuk Brussels, dan 25 untuk Antwerp. Per kapita
endowment Arlon town cluster lebih tinggi daripada Brussels dan Antwerp
untuk banyak layanan, tetapi penulis memperingatkan bahwa hitungan
penduduk Arlon town cluster hanya mencakup 12.714 penduduk, dibanding
32.093 untuk munisipalitas Arlon. {[}Section 3.2 and Section 4.3,
pp.~428 and 436; Figure 3; Table 6{]}

\textbf{Laporan sumber:} Kehadiran setidaknya satu layanan meningkat
seiring ukuran settlement. Sebagai contoh, primary school terdapat pada
37,2\% settlement berpenduduk kurang dari 1.000 dan 99,3\% settlement
berpenduduk lebih dari 250.000; university masing-masing terdapat pada
0,3\% dan 97,8\% kelompok tersebut. {[}Section 4.1, pp.~434-435; Table
5{]}

\textbf{Laporan sumber:} ANOVA menunjukkan perbedaan yang signifikan
pada endowment per 100.000 penduduk antara urban centres, towns, dan
villages untuk jenis layanan yang diuji. Rata-rata per kapita lebih
tinggi di villages untuk retail, banks, pharmacies, primary schools, dan
doctors; towns lebih tinggi untuk secondary schools, hospitals, dan
cinemas; urban centres hanya lebih tinggi untuk universities. Tukey HSD
menunjukkan bahwa sebagian besar pasangan kelompok berbeda signifikan,
dengan beberapa pengecualian yang dirinci di Table A2. {[}Appendix B,
pp.~451-453; Tables A1-A2{]}

\textbf{Laporan sumber:} Untuk perbandingan kelas ukuran settlement
dengan sedikitnya 25.000 penduduk, teks Appendix B menyatakan perbedaan
signifikan untuk retail, banks, pharmacies, primary schools, secondary
schools, hospitals, dan cinemas, sementara doctors dan universities
tidak berbeda signifikan. {[}Appendix B, pp.~452-454{]}

\textbf{Inferensi pengolahan:} Ringkasan teks Appendix B tidak
sepenuhnya cocok dengan Table A3: nilai \texttt{Sig.} banks yang terbaca
adalah 0,07, train stations memiliki 0,000 tetapi tidak disebut dalam
ringkasan teks, dan doctors 0,081 serta universities 0,273. Perbedaan
tersebut dibiarkan terlihat, bukan direkonsiliasi secara spekulatif.

\textbf{Laporan sumber:} Pada regresi Equation (2), ukuran settlement
tetap positif dan signifikan untuk semua sepuluh layanan setelah kontrol
lain dimasukkan. Centrality \texttt{L(t)} menambah daya jelaskan di luar
ukuran penduduk. Primary schools, retail, dan pharmacies terutama
berasosiasi dengan settlement yang dominan pada ambang waktu rendah;
secondary schools, banks, dan hospitals menunjukkan preferensi
supralokal; cinemas dan universities lebih terkait dengan ambang tinggi.
Adjusted R2 model berkisar dari 0,56 untuk primary schools sampai 0,98
untuk universities. {[}Section 4.3, pp.~436-441; Table 7; Figure 8{]}

\textbf{Laporan sumber:} Dampak perbatasan berbeda menurut layanan.
Retail dan banks lebih banyak hadir dekat perbatasan, sedangkan beberapa
layanan sosial lebih sedikit hadir di sana. Ketika settlement dekat
perbatasan dikeluarkan, hasil variabel \texttt{L(t)} disebut tetap
serupa, sehingga perbatasan memengaruhi jumlah layanan tetapi tidak
tampak mengubah hierarki yang diinferensikan. {[}Section 4.3,
pp.~438-439; Table 7 and footnote 2{]}

\textbf{Laporan sumber:} Train stations hampir tidak dipengaruhi
peringkat lokal. Korelasi Pearson antara jumlah train stations dan
layanan lain hanya 0,24-0,32. Penulis mengaitkan hal ini dengan
pertimbangan jalur jaringan dan permintaan perjalanan, bukan semata-mata
ukuran atau centrality settlement. {[}Section 4.3, pp.~438-439{]}

\textbf{Laporan sumber:} Pada Equation (3), effective speed yang lebih
tinggi umumnya berasosiasi dengan lebih sedikit lokasi layanan, dengan
pengecualian utama primary schools, train stations, dan universities
yang menunjukkan koefisien positif signifikan. Efek negatif signifikan
terlihat untuk banks, doctors, secondary schools, hospitals, dan
cinemas; retail serta pharmacies tidak menunjukkan koefisien effective
speed yang signifikan dalam Table 8. Jumlah destinasi yang terjangkau
dalam 120 menit berasosiasi negatif dengan sebagian besar layanan dan
positif dengan primary schools. {[}Section 4.4, pp.~441-445; Table 8{]}

\textbf{Laporan sumber:} Probabilitas menjadi local centre pada ambang
10 atau 20 menit menurun ketika effective speed lebih tinggi.
Perbandingan settlement dengan effective speed sedikitnya 11\% di bawah
atau di atas rata-rata nasional menghasilkan pola yang campuran:
secondary schools dan cinemas mempunyai ambang waktu yang relatif sama,
sedangkan primary schools, retail, dan pharmacies lebih kurang sensitif
terhadap ambang waktu dan tampak lebih mungkin terorganisasi menurut
jarak. {[}Section 4.4, pp.~443-445; Table 9{]}

\textbf{Laporan sumber:} Pada subset settlement berpenduduk kurang dari
25.000, sebagian layanan berorde lebih tinggi, termasuk secondary
schools dan cinemas, tidak lagi menunjukkan efek signifikan pada
\texttt{L(t)} yang lebih tinggi. Artikel menyebut ukuran subset yang
terbatas sebagai kemungkinan penjelasan. Peringkat populasi nasional
sangat berkorelasi dengan ukuran penduduk dan tidak memberi daya jelas
tambahan dalam model utama. {[}Sections 4.3-4.4, pp.~440-445; Figure
9{]}

\textbf{Inferensi pengolahan:} Hasil regresi mendukung hubungan
asosiasional antara posisi jaringan dan jumlah layanan setelah ukuran
penduduk dikontrol, tetapi tidak dengan sendirinya mengidentifikasi arah
kausal antara peningkatan kecepatan jalan dan perpindahan layanan.

\subsection{Findings}\label{findings-61}

\textbf{Laporan sumber:} Penulis menyimpulkan bahwa settlement yang
diinferensikan memiliki peran pusat lokal atau regional memiliki lokasi
layanan yang secara signifikan lebih banyak daripada settlement lain
dengan ukuran sama. Hierarki yang diusulkan dinilai bermakna untuk
menjelaskan service provision dan memperluas gagasan bahwa kinerja
settlement bergantung pada ukuran lokal serta neighbourhood yang
terkoneksi hingga settlement terkecil di EU27 dan UK. {[}Section 5,
pp.~445-446{]}

\textbf{Interpretasi penulis:} Layanan beroperasi pada skala spasial
yang berbeda. Dalam ringkasan Section 4.3, layanan lokal terutama
mencakup primary schools, retail, dan pharmacies; layanan subregional
mencakup secondary schools, banks, dan hospitals; layanan regional
mencakup cinemas dan universities. Temuan ini dibaca sebagai dukungan
terhadap ekspektasi CPT tentang hierarki pusat layanan. {[}Section 4.3,
pp.~438-439; Section 5, p.~446{]}

\textbf{Interpretasi penulis:} Jaringan yang lebih cepat dapat
menyebabkan hierarki neighbourhood menjadi lebih jarang dan layanan
terkonsentrasi di area yang lebih terkoneksi. Namun, hasil perbandingan
settlement berkecepatan rendah dan tinggi tidak seragam untuk semua
layanan, sehingga penulis tidak memperlakukan mekanisme tersebut sebagai
pola universal. {[}Section 4.4, pp.~443-445; Section 5, pp.~446{]}

\textbf{Inferensi pengolahan:} Centrality dalam artikel paling tepat
dibaca sebagai kapasitas relatif suatu settlement untuk menjadi pusat
layanan pada skala waktu tertentu. Ia tidak identik dengan ukuran kota,
status administratif, volume perjalanan, atau akses aktual seluruh
penduduk.

\subsection{Conclusions}\label{conclusions-61}

\textbf{Interpretasi penulis:} Metode yang diperkenalkan dipandang
komplementer terhadap definisi fungsional yang bertumpu pada aglomerasi
pasar tenaga kerja. Metode tersebut membantu menemukan settlement kecil
yang ``punch above their weight'' dalam penyediaan layanan dan penting
bagi konteks rural. {[}Section 5, pp.~445-446{]}

\textbf{Interpretasi penulis:} Penulis menyatakan bahwa peningkatan
jaringan jalan berpotensi memperkuat konsentrasi layanan apabila layanan
memang diatur oleh waktu tempuh, bahkan dapat memperkuat ketergantungan
pada mobil pribadi. Kesimpulan ini tetap disertai kualifikasi bahwa
sebagian layanan lebih dekat dengan logika jarak dan bahwa faktor
struktur spasial yang tidak terukur mungkin menjelaskan pola campuran.
{[}Section 5, pp.~446{]}

\textbf{Laporan sumber:} Artikel menegaskan bahwa klasifikasi ini adalah
heuristic, bergantung pada data sekunder, dan hasilnya tidak boleh
dianggap sebagai pengganti field work lokal. {[}Section 5, p.~446{]}

\subsection{Contributions}\label{contributions-61}

\textbf{Laporan sumber:} Kontribusi metodologis artikel adalah
klasifikasi hierarki centrality settlement yang menggunakan ukuran
penduduk relatif dalam neighbourhood berbasis waktu tempuh jalan.
Klasifikasi tersebut dihitung untuk hampir 50.000 settlement di EU27 dan
UK, bukan hanya kota besar. {[}Sections 1 and 3, pp.~422-430{]}

\textbf{Laporan sumber:} Artikel memvalidasi klasifikasi dengan sepuluh
kategori layanan dan menunjukkan bahwa posisi jaringan menjelaskan
variasi endowment setelah ukuran penduduk dikontrol. Artikel juga
memberi karakterisasi empiris tentang skala spasial layanan serta
menguji kaitan endowment dengan effective speed jaringan. {[}Sections
4.2-4.4, pp.~434-445{]}

\textbf{Interpretasi penulis:} Kontribusi substantifnya adalah
memperluas pembacaan tentang ketergantungan settlement pada
neighbourhood terkoneksi sampai ke town dan village, serta menyediakan
informasi pelengkap bagi peringkat populasi nasional dan definisi
functional urban area. Hasil klasifikasi disebut tersedia secara publik
melalui tautan data yang dicantumkan pada artikel. {[}Sections 1 and 5,
pp.~422-423 and 445-446{]}

\textbf{Inferensi pengolahan:} Nilai tambah utama artikel berada pada
operasionalisasi peran pusat layanan secara harmonis pada banyak
settlement, bukan pada pengukuran langsung kualitas, keterjangkauan,
atau penggunaan layanan.

\subsection{Research Gap}\label{research-gap-61}

\textbf{Laporan sumber:} Gap yang dinyatakan adalah belum tersedianya
cara universal yang memuaskan untuk menangkap fungsi penyedia layanan
settlement bagi hinterland langsung. Informasi yang komprehensif,
harmonis, dan cukup rinci tentang semua layanan publik serta privat di
seluruh wilayah Eropa juga dinyatakan belum tersedia. Karena itu,
settlement kecil yang penting bagi service provision dapat terlewat jika
klasifikasi hanya memakai ukuran penduduk nasional atau ketergantungan
pasar tenaga kerja. {[}Sections 1 and 2.4, pp.~422 and 425{]}

\textbf{Laporan sumber:} Artikel masih menyisakan kebutuhan untuk
memahami peran struktur spasial dan keputusan lokasi layanan dengan
lebih baik, terutama karena respons terhadap travel time berbeda antara
layanan dan antara settlement dengan effective speed rendah atau tinggi.
{[}Section 5, p.~446{]}

\textbf{Inferensi pengolahan:} Karena analisis yang dilaporkan
menggunakan perbandingan spasial dan cross-sectional differences, dampak
kausal perubahan jaringan jalan terhadap relokasi layanan belum
ditetapkan oleh artikel. Ini merupakan batas evidence untuk membaca
dugaan ``reshuffling'', bukan temuan kausal baru.

\subsection{Limitations}\label{limitations-61}

\textbf{Laporan sumber:} Penulis menyebut metode sebagai heuristic dan
mengakui bahwa metode tersebut tidak memasukkan banyak faktor relevan,
termasuk historical inertia. Ukuran penduduk dan waktu tempuh jalan
berasal dari data sekunder dan dapat menyimpang dari kondisi praktik.
Validasi terutama bergantung pada data crowdsourced yang mungkin
mengandung commission error, omission error, perbedaan klasifikasi, dan
perbedaan eligibility antarnegara, zona bahasa, atau wilayah. {[}Section
4, pp.~430-432; Section 5, p.~446{]}

\textbf{Laporan sumber:} Settlement dianalisis sebagai cluster polygon,
sehingga ukuran settlement tidak selalu sama dengan ukuran pasar atau
populasi munisipalitas. Contoh Arlon menunjukkan bahwa 12.714 penduduk
pada town cluster kurang dari setengah 32.093 penduduk munisipalitasnya,
sedangkan batas Brussels dan Antwerp mencakup beberapa munisipalitas.
{[}Section 4.3, p.~436; Table 6{]}

\textbf{Laporan sumber:} Lokasi layanan yang lebih dari 2,5 km dari
settlement terdekat tidak diperhitungkan. Hanya 75,5\% train stations
yang berhasil dihubungkan menurut Table 4; artikel menjelaskan bahwa
sebagian rail stops di lokasi terpencil mungkin melayani wilayah yang
lebih luas. Data transportasi publik komprehensif berskala Eropa juga
tidak tersedia dalam studi ini, sehingga train station hanya menjadi
indikator kasar transportasi publik. {[}Sections 4 and 4.1-4.3,
pp.~430-441; Table 4{]}

\textbf{Laporan sumber:} Matriks travel time dibuat terpisah per negara.
Dengan demikian, peran settlement pusat terhadap service endowment
lintas batas tidak dimasukkan ke dalam matriks utama, walaupun kedekatan
ke perbatasan diuji sebagai variabel kontrol. {[}Section 3.1,
pp.~426-427; Section 4.3, pp.~438-439{]}

\textbf{Laporan sumber:} Hasil pada subset settlement kecil kehilangan
signifikansi untuk sebagian layanan berorde lebih tinggi; artikel
mengaitkannya secara kemungkinan dengan ukuran subset yang terbatas.
Hasil perbandingan effective speed juga campuran, sehingga tidak semua
layanan dapat diperlakukan mengikuti travel time secara seragam.
{[}Section 4.3-4.4, pp.~440-445{]}

\textbf{Inferensi pengolahan:} Outcome berbasis jumlah lokasi POI tidak
cukup untuk menyimpulkan bahwa penduduk benar-benar memakai layanan
tersebut, memperoleh akses yang setara, atau mengalami peningkatan
kesejahteraan. Artikel juga tidak melaporkan observasi longitudinal
before-after untuk perubahan jaringan jalan.

\subsection{Challenges}\label{challenges-61}

\textbf{Laporan sumber:} Tantangan identifikasi adalah menentukan
settlement sebagai node tempat penyedia layanan secara masuk akal dapat
berlokasi, ketika fasilitas sering berada di luar batas settlement.
Penulis harus melakukan linkage ke grid cell terdekat, menetapkan batas
2,5 km, dan menangani lokasi terpencil seperti rail stops yang mungkin
melayani hinterland lebih luas. {[}Sections 3.1 and 4.1, pp.~426-427 and
432-434{]}

\textbf{Laporan sumber:} Ketiadaan data konsumsi barang dan layanan yang
komprehensif serta harmonis membuat fungsi settlement tidak dapat
diidentifikasi langsung dari interaksi aktual. Pemodelan dinamis
hierarki layanan juga dipandang terlalu kompleks untuk tujuan
identifikasi yang mencakup banyak sektor dan distribusi penduduk Eropa
yang tidak seragam. {[}Section 2.4, pp.~424-425{]}

\textbf{Laporan sumber:} Penulis menghadapi variasi lintas negara dalam
kualitas dan cakupan POI, persoalan interaksi lintas batas, serta
korelasi tinggi antara jarak dan waktu tempuh yang menyulitkan
pembandingan kedua prinsip pengorganisasian layanan. Effective speed
juga tidak dapat dihitung untuk seluruh settlement. {[}Sections 3.1 and
4.4, pp.~426-427 and 441-445{]}

\textbf{Inferensi pengolahan:} Tantangan analitis terbesar adalah
memisahkan pengaruh ukuran penduduk, posisi jaringan, jarak, dan
kualitas data POI ketika seluruhnya saling berkaitan secara spasial.
Model kontrol mengurangi masalah tersebut, tetapi tidak menghapusnya.

\subsection{Practical Implications}\label{practical-implications-61}

\textbf{Interpretasi penulis:} Hasil dapat membantu mengarahkan
investasi untuk territorial cohesion dan memprioritaskan settlement
dalam rencana tata ruang strategis nasional maupun regional. Metode ini
juga disebut dapat membantu zoning lokal dan regulasi lalu lintas dengan
mengidentifikasi peran teritorial settlement, termasuk settlement kecil
yang menjadi pusat layanan rural. {[}Section 5, pp.~445-446{]}

\textbf{Laporan sumber:} Penulis memperingatkan bahwa peningkatan
kecepatan jaringan dapat mengurangi peran central settlement pada skala
waktu rendah dan berpotensi memusatkan layanan di area yang lebih
terkoneksi. Jika layanan memang terorganisasi menurut travel time,
kebijakan mempercepat jaringan dapat memperkuat ketergantungan pada
mobil pribadi; karena itu hasil tidak boleh dibaca sebagai alasan
otomatis untuk mempercepat semua koneksi. {[}Sections 4.4 and 5,
pp.~443-446{]}

\textbf{Inferensi pengolahan:} Untuk keputusan praktis, \texttt{L(t)}
lebih aman digunakan sebagai alat screening awal untuk calon pusat
layanan, lalu diverifikasi dengan kondisi lapangan, data layanan lokal,
dan moda transportasi selain mobil.

\subsection{Applications}\label{applications-61}

\textbf{Laporan sumber:} Aplikasi yang ditunjukkan artikel adalah
klasifikasi centre lokal dan regional; pemetaan peran service provider
settlement; perbandingan endowment layanan menurut tipe, ukuran, dan
posisi jaringan; serta analisis hubungan service provision dengan
effective road speed. Penulis menyatakan pendekatan pada prinsipnya
dapat diterapkan secara global, tetapi hasil empiris yang disajikan
hanya untuk EU27 dan UK. {[}Sections 1, 3, and 4, pp.~422-445{]}

\textbf{Interpretasi penulis:} Klasifikasi dapat melengkapi functional
urban area dalam penentuan prioritas investasi territorial cohesion,
spatial planning, zoning, dan traffic regulation, terutama di luar
aglomerasi metropolitan. {[}Section 5, pp.~445-446{]}

\textbf{Inferensi pengolahan:} Penggunaan untuk memproyeksikan
konsekuensi suatu proyek jalan atau untuk menetapkan lokasi layanan
secara final memerlukan data temporal dan validasi lokal yang tidak
disediakan artikel; aplikasi semacam itu tidak dapat diperlakukan
sebagai hasil langsung studi.

\subsection{Future Research
(Optional)}\label{future-research-optional-61}

\textbf{Laporan sumber:} Artikel menyebut perlunya case studies untuk
menganalisis akses spasial publik terhadap layanan secara penuh karena
data public transport komprehensif berskala Eropa belum tersedia.
Artikel juga menyatakan bahwa diperlukan lebih banyak penelitian untuk
memahami peran struktur spasial dan keputusan lokasi layanan ketika
jaringan berubah. {[}Section 4.3, pp.~440-441; Section 5, p.~446{]}

\textbf{Inferensi pengolahan:} Tidak ada agenda future research lain
yang ditambahkan di luar dua arah yang disebutkan sumber.

\subsection{Provenance Notes}\label{provenance-notes-61}

\textbf{Laporan sumber:} File yang diverifikasi adalah tepat satu TXT
dengan prefiks \texttt{B18} di bawah \texttt{source/}:
\texttt{source/journal/B18-jacobs-crisioni-kompil-dijkstra-2023-big-in-the-neighbourhood-identifying-local-and-regional-centres-through-their-network-position-10.1111-pirs.12727.txt}.
Header file menunjuk ke PDF yang bersesuaian, mencatat ekstraksi pada
2026-08-31 dengan \texttt{pdftotext\ -layout}, dan mencatat PDF 38
halaman. {[}PDF Metadata, lines 1-28{]}

\textbf{Laporan sumber:} Seluruh rentang teks bernomor 1-1572 dibaca,
termasuk metadata PDF, artikel utama, references, Appendix A-C, Table
A1-A5, catatan kaki, dan abstrak bahasa Spanyol serta Jepang. Artikel
yang diidentifikasi dari byline adalah Jacobs-Crisioni, Kompil, dan
Dijkstra; DOI, tanggal penerimaan, jurnal, volume, dan halaman diambil
dari file TXT. {[}lines 30-32, 41-73, 1017-1162, 1167-1543, and
1548-1572{]}

\textbf{Interpretasi penulis:} Batas evidence yang dinyatakan artikel
sendiri dipertahankan di review ini. Tidak ada sumber luar yang
digunakan atau diakses; nama karya lain hanya disebut ketika artikel TXT
menggunakannya dalam literature survey atau referensi.

\textbf{Inferensi pengolahan:} \texttt{wc\ -l} melaporkan 1.571 newline,
sedangkan pembacaan file menandai 1.572 baris karena baris terakhir
tidak menjadi newline terpisah; tidak ada rentang isi yang dilewati
dalam pembacaan. Artefak ekstraksi meliputi mojibake pada metadata nama
penulis (\texttt{Chris\ Jacobsâ€\ Crisioni}), karakter kontrol pada
sebagian tanda minus di tabel, form-feed penanda halaman, dan beberapa
tata letak tabel yang terpecah. Nama penulis dalam review mengikuti
byline artikel yang terbaca sebagai \texttt{Chris\ Jacobs-Crisioni}.

\textbf{Inferensi pengolahan:} Artikel menyebut sepuluh service types
pada bagian utama dan Tables 4, 7, serta 8, sedangkan Appendix B pada
satu kalimat menyebut ``nine services'' meskipun tabelnya memuat sepuluh
kategori. Review melaporkan sepuluh kategori dan menandai
ketidakkonsistenan tersebut, bukan mengoreksinya dengan sumber luar.
Ringkasan skala layanan juga memakai ``up to 30 minutes'' pada Section
4.3 dan ``20-minute range'' pada Section 5; kedua locator dipertahankan
pada bagian Results tanpa dipaksakan menjadi satu angka.

\textbf{Inferensi pengolahan:} Untuk perbandingan tipe settlement,
narasi Appendix B menyebut rata-rata per kapita banks lebih tinggi di
villages, sedangkan Table A1 yang terbaca menunjukkan mean towns 30,10
dan villages 29,74; Table A2 juga menunjukkan perbedaan towns-villages
tidak signifikan. Pernyataan narasi dan angka tabel dipertahankan
sebagai dua bukti yang tidak sepenuhnya konsisten.

\textbf{Status review:} Lengkap secara struktur dan evidence untuk TXT
B18 yang tersedia. Review ini tidak mengubah TXT sumber maupun PDF,
tidak menangani kode lain, dan tidak membuat dokumen induk.

\section{Review B19}\label{review-b19}

\textbf{Review kode B19}

\textbf{Identitas sumber:} Artikel jurnal berjudul \emph{Territorial
Arrangements of Small and Medium-Sized Towns from a Functional-Spatial
Perspective}.

\textbf{Penulis:} Luděk Sýkora dan Ondřej Mulíček. Nama pada baris judul
mengalami artefak encoding; entri referensi di dalam TXT menampilkan
Sýkora, L. dan Mulíček, O. {[}TXT, baris 28-36 dan 876-884{]}.

\textbf{Jurnal dan identifikasi:} \emph{Tijdschrift voor Economische en
Sociale Geografie}, 2017, DOI \texttt{10.1111/tesg.12249}. Metadata yang
terbaca mencantumkan Vol. 00, No.~00, pp.~00-00 {[}TXT, p.~1, baris
75-76{]}. Naskah dicatat diterima pada Februari 2015 dan disetujui pada
November 2016 {[}TXT, p.~1, baris 38{]}.

\textbf{Bahan review:} Hanya file
\texttt{source/journal/B19-sykora-mulicek-2017-territorial-arrangements-of-small-and-medium-sized-towns-from-a-functional-spatial-perspective-10.1111-tesg.12249.txt}
yang digunakan.

\subsection{Summarized Abstract}\label{summarized-abstract-62}

{[}Laporan sumber{]} Artikel mengembangkan dan menguji metodologi
analisis fungsional-spasial untuk mengidentifikasi kota atau pusat
urban, membedakan kota besar dari kota kecil dan menengah, serta
menentukan posisi kota kecil dan menengah dalam susunan teritorial.
Posisi tersebut dibedakan menjadi otonom, berjaringan dengan kota atau
kota kecil-menengah lain, atau teraglomerasi dengan kota besar.
Metodologi kemudian digunakan untuk membandingkan sistem permukiman di
Catalonia, Czechia, Central France, dan Slovenia berdasarkan keberadaan
serta sifat jaringan antarkota, juga kinerja populasi dan pekerjaan
menurut jenis susunan teritorial {[}TXT, Abstrak, p.~1, baris 40-51{]}.

{[}Interpretasi penulis{]} Fokus artikel bukan ukuran populasi semata,
melainkan fungsi pusat urban dan keterkaitan fungsional antarpemukiman.
Kota kecil dan menengah ditempatkan dalam sistem permukiman yang lebih
luas, lalu dibandingkan menurut apakah mereka berdiri sendiri,
berhubungan secara hierarkis dengan kota besar, atau menjadi bagian dari
jaringan yang lebih polisentris {[}TXT, p.~2-4, baris 82-98 dan
198-250{]}.

{[}Inferensi review{]} Artikel menghasilkan tipologi dan perbandingan
deskriptif, bukan estimasi kausal mengenai pengaruh suatu susunan
teritorial terhadap pertumbuhan. Penulis sendiri menyebut analisis
kinerja tersebut sebagai analisis sederhana dan observasi awal {[}TXT,
p.~14-15, baris 710-743{]}.

\subsection{Problem Statement}\label{problem-statement-62}

{[}Laporan sumber{]} Kota kecil dan menengah merupakan pusat lokal yang
menyediakan pekerjaan dan layanan, tetapi sebagian menghadapi kehilangan
pekerjaan, layanan, dan fungsi sentralitas. Keterkaitan antarkota
berpotensi memberi manfaat, sehingga penting untuk mengetahui apakah
kota-kota tersebut berdiri sendiri, terhubung dalam jaringan, atau
berada dalam aglomerasi kota besar, serta apakah posisi tersebut
berkaitan dengan kinerja mereka {[}TXT, p.~1, baris 55-74{]}.

{[}Laporan sumber{]} Literatur yang lebih baru mulai memberi perhatian
kepada kota lapis kedua, tetapi posisi kota kecil dan menengah dalam
sistem urban masih kurang diketahui. Artikel merujuk pada dua persoalan
yang diajukan dalam literatur: sejauh mana sistem urban menjadi lebih
polisentris dan terintegrasi secara spasial, serta apakah sistem yang
polisentris dan terintegrasi lebih efisien secara ekonomi daripada
sistem monosentris atau tidak terintegrasi {[}TXT, p.~1, baris 60-74{]}.

{[}Laporan sumber{]} Kajian tentang \emph{borrowed size} telah menelaah
fungsi metropolitan tingkat tinggi dan amenitas budaya kelas atas,
tetapi hubungan kota kecil dengan pasar kerja lokal serta penyediaan
layanan lokal disebut masih belum diteliti {[}TXT, p.~2-3, baris 101-121
dan 138-146{]}.

{[}Interpretasi penulis{]} Masalah penelitian memiliki dua lapisan:
kebutuhan untuk mengidentifikasi pusat urban dan kebutuhan untuk membaca
posisi pusat tersebut dalam susunan teritorial yang dapat menggabungkan
hierarki dan jaringan. Analisis yang hanya berfokus pada kota besar atau
dikotomi urban-rural dianggap mengabaikan mikro-wilayah tempat aktivitas
harian penduduk kota kecil dan menengah berlangsung {[}TXT, p.~3-4,
baris 147-191 dan 198-222{]}.

{[}Inferensi review{]} Dengan demikian, celah yang hendak diisi bukan
hanya kekurangan data mengenai kota kecil, melainkan juga kekurangan
cara operasional untuk membedakan hubungan aglomeratif dan hubungan
jaringan berdasarkan relasi fungsional yang terukur.

\subsection{Objectives}\label{objectives-62}

{[}Laporan sumber{]} Tujuan artikel adalah mengembangkan dan menguji
metode untuk mengidentifikasi kota kecil dan menengah serta susunan
teritorialnya dengan perspektif fungsional-spasial {[}TXT, Abstrak,
p.~1, baris 40-48{]}.

{[}Laporan sumber{]} Tujuan operasional yang dinyatakan atau dijalankan
dalam artikel meliputi:

\begin{itemize}
\tightlist
\item
  mengidentifikasi permukiman yang berfungsi sebagai pusat urban
  berdasarkan pekerjaan dan relasi komuter;
\item
  membedakan pusat urban berperingkat lebih rendah, yang ditafsirkan
  sebagai kota kecil dan menengah, dari pusat urban berperingkat lebih
  tinggi, yang ditafsirkan sebagai kota besar;
\item
  mengklasifikasikan kota kecil dan menengah sebagai otonom,
  teraglomerasi dengan kota besar atau kota kecil-menengah, atau
  berjaringan dengan kota besar atau kota kecil-menengah;
\item
  membandingkan susunan permukiman di Catalonia, Czechia, Central
  France, dan Slovenia;
\item
  memberikan gambaran awal tentang pertumbuhan atau penurunan populasi
  dan pekerjaan menurut posisi kota dalam susunan teritorial {[}TXT,
  p.~2-4, baris 82-98 dan 245-257{]}.
\end{itemize}

{[}Interpretasi penulis{]} Contoh Czechia dipakai untuk
mendokumentasikan langkah-langkah metodologi secara rinci, sedangkan
empat wilayah atau negara digunakan sebagai ilustrasi komparatif dan
bukan sebagai pengujian menyeluruh atas seluruh sistem urban Eropa
{[}TXT, p.~2, baris 82-98; p.~9, baris 490-495{]}.

\subsection{Literature Survey}\label{literature-survey-62}

{[}Laporan sumber{]} Artikel membangun tinjauan literatur di sekitar
beberapa gagasan berikut:

\begin{itemize}
\tightlist
\item
  Kota dipahami sebagai bagian dari sistem kota dan wilayah; nilai kota
  kecil terletak pada karakter fungsional, pengaruh, jangkauan, dan
  relasinya, bukan hanya jumlah penduduk atau kepadatan {[}TXT, p.~2,
  baris 101-119{]}.
\item
  Teori jaringan kota menekankan sinergi melalui aktivitas dan hubungan
  yang komplementer, resiprokal, dan kooperatif. Sistem urban
  polisentris karena itu sering diasosiasikan dengan kondisi ekonomi
  yang menguntungkan, walaupun bukti apakah wilayah polisentris dan
  berjaringan berkinerja lebih baik daripada wilayah monosentris atau
  hierarkis masih belum jelas {[}TXT, p.~2, baris 120-141{]}.
\item
  Konsep \emph{borrowed size} digunakan untuk menggambarkan kota yang
  lebih kecil di kawasan metropolitan dan dapat berkinerja lebih baik
  karena kedekatan geografis serta keterhubungan fungsional dengan kota
  besar. Konsep \emph{agglomeration shadow} dipakai untuk menjelaskan
  kemungkinan bahwa persaingan dengan kota besar justru melemahkan
  fungsi sentralitas kota kecil dan mengubahnya menjadi kawasan hunian
  yang bergantung pada kota besar {[}TXT, p.~2-3, baris 124-146{]}.
\item
  Kajian sebelumnya mengenai jaringan kota, wilayah urban polisentris,
  dan \emph{borrowed size} lebih banyak berfokus pada kota besar, fungsi
  metropolitan tingkat tinggi, atau amenitas budaya kelas atas. Pasar
  kerja lokal dan layanan lokal kota kecil disebut sebagai wilayah yang
  belum dieksplorasi {[}TXT, p.~2, baris 110-121{]}.
\item
  Literatur mengenai hierarki pusat menurut Christaller, Lösch, Berry,
  dan tradisi analisis urban lainnya dipertemukan dengan literatur
  jaringan kota dan wilayah urban polisentris. Artikel menekankan bahwa
  sistem urban dapat memadukan ciri hierarki dan jaringan, sehingga satu
  konsep tunggal tidak selalu cukup {[}TXT, p.~3, baris 163-191{]}.
\end{itemize}

{[}Interpretasi penulis{]} Posisi teoretis artikel adalah
fungsional-spasial: pusat urban dikenali melalui konsentrasi fungsi
sentralitas dan hubungan dengan permukiman lain, sedangkan tingkat kota
besar atau kota kecil-menengah ditentukan melalui peringkat fungsional
dalam sistem urban. Perspektif ini dimaksudkan untuk menghindari
analisis yang menyamakan status urban dengan batas administratif atau
ukuran populasi saja {[}TXT, p.~3-4, baris 198-244{]}.

{[}Inferensi review{]} Tinjauan tersebut memberi dasar konseptual yang
relevan untuk studi jaringan dan aglomerasi, tetapi artikel tidak
melakukan uji sistematis atas seluruh teori yang dirujuk. Literatur
berfungsi terutama sebagai landasan definisi, klasifikasi, dan perumusan
celah penelitian.

\subsection{Method Used}\label{method-used-62}

{[}Laporan sumber{]} Penelitian menggunakan analisis fungsional-spasial
komparatif. Dua kelompok informasi dipakai: ukuran permukiman berupa
jumlah pekerjaan dan populasi, serta relasi antarpemukiman yang
direpresentasikan oleh arus komuter dari tempat tinggal ke tempat kerja.
Penulis memilih komuter kerja sebagai satu relasi utama karena dianggap
paling sering, intensif, dan penting dalam kehidupan harian serta
operasi mikro-wilayah {[}TXT, p.~4, baris 198-250{]}.

{[}Laporan sumber{]} Tahapan metode yang dijelaskan adalah sebagai
berikut:

\begin{itemize}
\tightlist
\item
  Pusat pekerjaan diidentifikasi dari unit LAU2 yang memiliki pekerjaan
  di atas ambang minimum dan menjadi tujuan arus komuter terkuat dari
  setidaknya satu munisipalitas lain. Untuk Czechia, ambang yang dipilih
  adalah 1.000 pekerjaan, berdasarkan distribusi frekuensi munisipalitas
  yang menjadi tujuan arus komuter utama {[}TXT, p.~5-6, baris
  268-323{]}.
\item
  Setiap permukiman sumber dihubungkan ke satu tujuan berdasarkan arus
  komuter terkuat. Jika tujuan terkuat bukan pusat pekerjaan, hubungan
  diteruskan melalui permukiman perantara ke pusat pekerjaan. Hasil
  awalnya adalah \emph{proto-microregions} {[}TXT, p.~6, baris
  325-348{]}.
\item
  \emph{Proto-microregions} yang sangat kecil dihilangkan menggunakan
  ambang populasi 6.000 penduduk di Czechia, lalu wilayah yang
  terfragmentasi dikonsolidasikan menjadi mikro-wilayah yang kontigu.
  Hasilnya adalah 260 pusat urban mikro-regional dari 367 pusat
  pekerjaan awal {[}TXT, p.~6, baris 290-327{]}.
\item
  Arus signifikan antarpusat mikro-regional ditentukan melalui algoritme
  yang mengambil lima arus keluar tertinggi, menghitung proporsinya,
  lalu mengorelasikan distribusi aktual dengan lima distribusi ideal:
  satu tujuan, dua tujuan, tiga tujuan, empat tujuan, atau lima tujuan
  dengan bobot relatif setara. Jumlah arus signifikan ditentukan dari
  korelasi tertinggi {[}TXT, Catatan 1, p.~15-16, baris 756-783{]}.
\item
  Dari arus signifikan tersebut, arus yang penting bagi pasar kerja asal
  dipertahankan jika mencakup sedikitnya 5 persen penduduk aktif secara
  ekonomi di pusat asal. Di Czechia, 205 arus memenuhi kriteria ini.
  Arus lalu dinilai penting bagi pusat tujuan jika komuter masuk
  mencakup sedikitnya 1 persen pekerjaan di pusat tujuan; 116 arus
  memenuhi kriteria tujuan dan 89 hanya signifikan bagi pusat asal
  {[}TXT, p.~7, baris 351-381{]}.
\item
  Hierarki urban dibentuk dari jumlah arus masuk signifikan, ditambah
  proporsi arus keluar ketika sebuah pusat memiliki lebih dari satu arus
  keluar signifikan. Bobot yang digunakan adalah 0,5; 0,33; 0,25; atau
  0,2. Di Czechia, pusat berpenduduk lebih dari 50.000 dengan posisi
  fungsional sedikitnya 2,5 diklasifikasikan sebagai kota besar {[}TXT,
  p.~7-8, baris 411-440{]}.
\item
  Posisi dan susunan teritorial diklasifikasikan ke dalam pusat kota
  besar, pusat otonom, pusat aglomeratif, dan pusat berjaringan.
  Subkategori membedakan hubungan dengan kota besar atau dengan kota
  kecil-menengah, serta apakah arus terutama signifikan bagi asal,
  tujuan, atau keduanya {[}TXT, p.~8, baris 441-464{]}.
\item
  Kinerja dibandingkan secara deskriptif melalui perubahan populasi dan
  pekerjaan. Periode yang digunakan umumnya 2001-2011, sedangkan Central
  France menggunakan 2001-2009 dan data Catalonia tidak tersedia untuk
  perbandingan tersebut {[}TXT, p.~12 dan 14, baris 620-625 dan
  687-704{]}.
\end{itemize}

{[}Interpretasi penulis{]} Metode sengaja tidak diarahkan untuk
mendeliniasi wilayah perjalanan kerja yang sepenuhnya mandiri. Penulis
mengabstraksikan tingkat \emph{self-containment} karena pengukurannya
dinilai dapat menguntungkan kawasan besar yang saling terhubung dan
menyembunyikan struktur mikro-wilayah kota kecil yang berjaringan
{[}TXT, p.~6-7, baris 329-363{]}.

{[}Inferensi review{]} Tidak ada model regresi, desain eksperimen, atau
estimasi efek yang dilaporkan. Metode menggabungkan klasifikasi berbasis
ambang, analisis jaringan komuter, dan perbandingan pertumbuhan
antarkategori.

\subsection{Dataset}\label{dataset-62}

{[}Laporan sumber{]} Dataset terdiri atas jumlah populasi, jumlah
pekerjaan, dan matriks total arus perjalanan ke tempat kerja antarunit
permukiman dalam wilayah yang dikaji. Penulis menyebut bahwa data relasi
idealnya juga mencakup arus lintas batas wilayah, terutama karena
wilayah kajian tidak selalu merupakan kawasan perjalanan kerja yang
sepenuhnya mandiri {[}TXT, p.~4, baris 198-250{]}.

{[}Laporan sumber{]} Data dan unit untuk empat wilayah diringkas sebagai
berikut dalam Tabel 1:

\begin{itemize}
\tightlist
\item
  Catalonia, Spain (ES51): populasi total 6.343.110, 946 unit, 66
  mikro-wilayah atau pusat mikro-regional, 6,98 persen unit menjadi MRC,
  rata-rata populasi mikro-wilayah 96.107, dan rata-rata populasi pusat
  MRC 48.827 atau 22.878 jika kota besar dikeluarkan.
\item
  Central France region (FR24): populasi total 1.565.920, 116 unit, 19
  mikro-wilayah atau MRC, 16,38 persen unit menjadi MRC, rata-rata
  populasi mikro-wilayah 104.395, dan rata-rata populasi pusat MRC
  62.673 atau 24.292 tanpa kota besar.
\item
  Czech Republic (CZ0): populasi total 10.230.060, 6.258 unit, 260
  mikro-wilayah atau MRC, 4,15 persen unit menjadi MRC, rata-rata
  populasi mikro-wilayah 39.346, dan rata-rata populasi pusat MRC 24.193
  atau 14.030 tanpa Praha.
\item
  Slovenia (SI0): populasi total 1.994.026, 192 unit, 51 mikro-wilayah
  atau MRC, 26,56 persen unit menjadi MRC, rata-rata populasi
  mikro-wilayah 39.099, dan rata-rata populasi pusat MRC 26.348 atau
  18.265 tanpa kota besar {[}TXT, Tabel 1, p.~10, baris 509-525{]}.
\end{itemize}

{[}Laporan sumber{]} Unit dasar umumnya LAU2 atau munisipalitas yang
mencakup pusat permukiman dan permukiman kecil di sekitarnya. Di France,
beberapa munisipalitas digabung menjadi \emph{unités urbaines}. Ukuran
unit tidak seragam: rata-rata populasi unit di Central France dan
Slovenia melebihi 10.000, Catalonia sekitar 7.000, sedangkan Czechia di
bawah 2.000 {[}TXT, p.~9, baris 473-483{]}.

{[}Laporan sumber{]} Ambang identifikasi juga berbeda. Ambang 1.000
pekerjaan digunakan di semua wilayah kecuali Slovenia, yang menggunakan
semua arus signifikan, bukan hanya arus maksimum. Ambang populasi untuk
delimitasi mikro-wilayah adalah 5.000 di Catalonia dan 10.000 di Central
France serta Slovenia; Central France juga memakai ambang pekerjaan
minimum 4.000 {[}TXT, p.~9, baris 480-492{]}.

{[}Batas informasi{]} Nama basis data statistik asli, definisi teknis
setiap variabel, prosedur pembersihan data, dan prosedur penanganan data
hilang tidak disebutkan dalam file. Keterlibatan ESPON TOWN disebut
sebagai konteks pengumpulan dan analisis awal data, tetapi sumber asli
setiap seri data tidak dirinci {[}TXT, p.~9, baris 490-495; p.~16, baris
826-830{]}.

\subsection{Population / Sample}\label{population-sample-62}

{[}Laporan sumber{]} Populasi analisis berupa unit permukiman terkecil
yang tersedia dalam sistem wilayah terpilih, bukan responden survei.
Unit tersebut umumnya munisipalitas LAU2, sedangkan France menggunakan
\emph{unités urbaines}. Dari unit-unit tersebut dibentuk pusat
pekerjaan, \emph{proto-microregions}, mikro-wilayah, dan akhirnya pusat
urban mikro-regional {[}TXT, p.~4-6, baris 251-257 dan 268-348{]}.

{[}Laporan sumber{]} Czechia dipakai sebagai contoh rinci dengan 6.258
munisipalitas; 493 memiliki sedikitnya 1.000 pekerjaan, 645 menjadi
tujuan arus komuter maksimum dari munisipalitas lain, 367 memenuhi
irisan kedua kriteria, dan 260 menjadi pusat urban mikro-regional
setelah proses delimitasi {[}TXT, p.~6, baris 290-327{]}. Empat kasus
komparatif menghasilkan 66 MRC di Catalonia, 19 di Central France, 260
di Czechia, dan 51 di Slovenia {[}TXT, Tabel 1, p.~10, baris 515-525{]}.

{[}Interpretasi penulis{]} Artikel memproses unit yang tersedia di
wilayah kajian, dengan penyaringan berdasarkan pekerjaan, arus komuter,
dan populasi. Tidak ada pengambilan sampel probabilistik atau survei
individu yang dilaporkan.

{[}Batas informasi{]} Kriteria pemilihan empat wilayah dari sepuluh
wilayah atau negara yang disebut berada dalam proyek ESPON TOWN tidak
dijelaskan lebih lanjut dalam file {[}TXT, p.~9, baris 490-495{]}.
Ukuran sampel statistik formal, tingkat respons, dan galat sampling
tidak berlaku atau tidak disebutkan dalam file.

\subsection{Dependent Variables}\label{dependent-variables-62}

{[}Laporan sumber{]} Artikel tidak menamai variabel dependen dalam
bentuk model statistik formal. Ukuran kinerja yang dibandingkan adalah
perubahan atau pertumbuhan populasi dan pekerjaan menurut kategori
fungsional kota. Perbandingan ditampilkan pada Figure 7 untuk Czechia,
Slovenia, dan Central France; data Catalonia tidak tersedia untuk
analisis kinerja tersebut {[}TXT, Figure 7, p.~14, baris 650-704{]}.

{[}Interpretasi penulis{]} Dalam pembacaan analitis, populasi dan
pekerjaan merupakan luaran kinerja, sedangkan posisi fungsional dan
jenis susunan teritorial merupakan dimensi klasifikasi yang dipakai
untuk membandingkan luaran. Artikel tidak memberikan rumus perubahan,
model estimasi, koefisien, atau uji signifikansi untuk hubungan
tersebut.

{[}Inferensi review{]} Jumlah pekerjaan memiliki dua fungsi sekaligus:
dipakai untuk mengidentifikasi pusat pekerjaan dan dipakai sebagai salah
satu ukuran kinerja. Karena itu, perbedaan pekerjaan antarkategori tidak
dapat dibaca sebagai bukti mandiri bahwa kategori tersebut menyebabkan
pertumbuhan atau penurunan pekerjaan. TXT tidak membahas implikasi
analitis dari penggunaan ganda tersebut.

\subsection{Results}\label{results-62}

{[}Laporan sumber{]} Hasil prosedur Czechia menunjukkan 367 pusat
pekerjaan dari 6.258 munisipalitas, 260 pusat urban mikro-regional, dan
14 kota besar dari 260 pusat tersebut. Analisis arus menghasilkan 415
arus signifikan antarpusat, 205 arus yang signifikan bagi pasar kerja
asal, serta 116 arus yang juga signifikan bagi pasar kerja tujuan
{[}TXT, p.~6-8, baris 306-327, 367-381, dan 425-440{]}.

{[}Laporan sumber{]} Tabel 2 melaporkan jumlah dan persentase pusat
urban menurut susunan teritorial. Kategori \texttt{LC} adalah kota
besar; \texttt{NETW-LC} berjaringan dengan kota besar; \texttt{AGLO-LC}
teraglomerasi dengan kota besar; \texttt{NETW-SMSC-D} berjaringan dengan
kota kecil-menengah sebagai tujuan; \texttt{NETW-SMSC-S} berjaringan
dengan kota kecil-menengah sebagai asal; \texttt{AGLO-SMSC}
teraglomerasi dengan kota kecil-menengah; dan \texttt{AUTO} otonom
{[}TXT, p.~8 dan Tabel 2, baris 441-464 dan 536-544{]}. Rinciannya
adalah:

\begin{itemize}
\tightlist
\item
  Catalonia: 4 LC atau 6 persen, 10 NETW-LC atau 15 persen, 19 AGLO-LC
  atau 29 persen, 10 NETW-SMSC-D atau 15 persen, 14 NETW-SMSC-S atau 21
  persen, tidak ada AGLO-SMSC, dan 9 AUTO atau 14 persen.
\item
  Central France: 5 LC atau 26 persen, 3 NETW-LC atau 16 persen, 8
  AGLO-LC atau 42 persen, 1 NETW-SMSC-D atau 5 persen, 2 NETW-SMSC-S
  atau 11 persen, tidak ada AGLO-SMSC dan tidak ada AUTO.
\item
  Czech Republic: 14 LC atau 5 persen, 19 NETW-LC atau 7 persen, 66
  AGLO-LC atau 26 persen, 49 NETW-SMSC-D atau 19 persen, 67 NETW-SMSC-S
  atau 26 persen, 6 AGLO-SMSC atau 2 persen, dan 39 AUTO atau 15 persen.
\item
  Slovenia: 4 LC atau 8 persen, 10 NETW-LC atau 20 persen, 9 AGLO-LC
  atau 17 persen, 11 NETW-SMSC-D atau 22 persen, 16 NETW-SMSC-S atau 31
  persen, 1 AGLO-SMSC atau 2 persen, dan tidak ada AUTO {[}TXT, Tabel 2,
  p.~11, baris 536-544{]}.
\end{itemize}

{[}Laporan sumber{]} Tabel 3 menunjukkan proporsi populasi regional yang
berada pada tiap kategori. Di Catalonia, proporsinya adalah LC 28,44
persen, NETW-LC 6,03 persen, AGLO-LC 10,31 persen, NETW-SMSC-D 2,71
persen, NETW-SMSC-S 2,04 persen, AGLO-SMSC 0 persen, AUTO 1,28 persen,
dan non-centres 49,20 persen. Di Central France, proporsinya adalah LC
54,33 persen, NETW-LC 6,48 persen, AGLO-LC 9,26 persen, NETW-SMSC-D 4,09
persen, NETW-SMSC-S 1,89 persen, AGLO-SMSC 0 persen, AUTO 0 persen, dan
non-centres 23,96 persen. Di Czech Republic, proporsinya adalah LC 27,75
persen, NETW-LC 4,16 persen, AGLO-LC 6,84 persen, NETW-SMSC-D 10,89
persen, NETW-SMSC-S 6,74 persen, AGLO-SMSC 0,28 persen, AUTO 4,83
persen, dan non-centres 38,51 persen. Di Slovenia, proporsinya adalah LC
24,34 persen, NETW-LC 9,89 persen, AGLO-LC 6,33 persen, NETW-SMSC-D
14,76 persen, NETW-SMSC-S 11,93 persen, AGLO-SMSC 0,14 persen, AUTO 0
persen, dan non-centres 32,61 persen {[}TXT, Tabel 3, p.~11, baris
548-557{]}.

{[}Laporan sumber{]} Susunan populasi berbeda antarkasus. Central France
memiliki sekitar separuh populasi di kota besar, sekitar 22 persen di
pusat kecil-menengah, dan sisanya di luar pusat urban. Czechia dan
Slovenia lebih merata antara kota besar, pusat kecil-menengah, dan
non-centres. Catalonia memiliki hampir separuh populasi di luar pusat
urban dan sisanya lebih condong ke pusat besar {[}TXT, p.~11-12, baris
560-570 dan Figure 3, baris 572-596{]}.

{[}Laporan sumber{]} AGLO-LC paling lazim di Catalonia dan Central
France. Slovenia memiliki banyak keterkaitan dengan kota besar yang
lebih resiprokal dan kurang hierarkis, tercermin dalam NETW-LC. Czechia
sering memiliki kota yang teraglomerasi dengan kota besar, tetapi juga
dicirikan oleh jaringan padat antarkota kecil-menengah. Kota otonom
lazim di Czechia dan Catalonia; tidak ada SMSC otonom di Central France
atau Slovenia {[}TXT, p.~12, baris 609-630{]}.

{[}Laporan sumber{]} Untuk kinerja pekerjaan, pekerjaan di pusat
kecil-menengah Czechia turun 12 persen sementara pekerjaan di pusat
besar naik 11 persen. Di Slovenia, pekerjaan di kota besar naik 13
persen, di kota berjaringan naik 2 persen, dan di kota teraglomerasi
turun 4 persen. Di Central France, pekerjaan di kota besar naik 7
persen, di kota berjaringan naik 1 persen, dan di kota teraglomerasi
turun 5 persen. Distribusi populasi penduduk di Central France relatif
stabil; penurunan populasi kota berjaringan tampak lebih kuat daripada
kota teraglomerasi di Czechia dan Central France {[}TXT, Figure 7,
p.~14, baris 650-709{]}.

{[}Laporan sumber{]} Tidak ada perbedaan signifikan antara kategori
fungsional SMSC yang lebih rinci. Di tingkat pusat urban individual
terdapat variasi kinerja yang besar, yang oleh penulis dikaitkan dengan
beragam faktor kontekstual {[}TXT, p.~14, baris 689-704{]}.

\subsection{Findings}\label{findings-62}

{[}Laporan sumber{]} Temuan utama artikel adalah bahwa kota yang relatif
kecil dapat berfungsi sebagai pusat lokal, menyediakan pekerjaan,
menarik arus komuter dari desa atau permukiman yang lebih kecil, dan
membentuk mikro-wilayah yang relevan bagi kehidupan sehari-hari penduduk
{[}TXT, p.~15, baris 720-730{]}.

{[}Laporan sumber{]} Susunan teritorial kota kecil-menengah bervariasi
antarwilayah. Ada sistem yang didominasi aglomerasi kota besar, sistem
berjaringan dengan kota besar, jaringan yang terutama terdiri atas kota
kecil-menengah, serta pusat otonom di wilayah yang lebih periferal
{[}TXT, p.~8-12, baris 441-464 dan 609-630{]}.

{[}Interpretasi penulis{]} Penurunan pekerjaan di kota teraglomerasi di
Slovenia dan Central France dapat dijelaskan oleh persaingan dengan kota
besar dan transformasi kota teraglomerasi dari tempat produksi serta
layanan menjadi simpul suburban yang lebih residensial. Penjelasan
tentang tidak adanya kota otonom juga dikaitkan dengan lanskap yang
sangat terurbanisasi di Central France dan pola nasional yang
polisentris di Slovenia {[}TXT, p.~12 dan 14, baris 609-630 dan
689-704{]}.

{[}Interpretasi penulis{]} Kinerja kota besar berbeda secara nyata dari
kota kecil-menengah, tetapi tipe fungsional di antara kota
kecil-menengah tidak menunjukkan perbedaan besar. Artikel menekankan
bahwa hasil ini merupakan observasi awal, bukan penjelasan lengkap atas
kinerja kota {[}TXT, p.~14-15, baris 710-743{]}.

{[}Inferensi review{]} Hasil terutama menunjukkan asosiasi deskriptif
antara posisi berbasis arus komuter dan kinerja populasi atau pekerjaan.
Hasil belum membuktikan mekanisme \emph{borrowed size},
\emph{agglomeration shadow}, atau efek kausal jaringan, karena relasi
lain, faktor kontekstual, dan urutan sebab-akibat tidak diestimasi
secara formal dalam TXT.

\subsection{Conclusions}\label{conclusions-62}

{[}Laporan sumber{]} Penulis menyimpulkan bahwa kota kecil dan menengah
beserta konstelasi spasialnya harus menjadi bagian integral dari
pemahaman tentang sistem urban dan regional. Kota-kota tersebut mencakup
proporsi populasi yang substansial, yaitu 20-45 persen di wilayah yang
diteliti, tetapi sering tidak masuk dalam analisis sistem urban {[}TXT,
p.~15, baris 709-727{]}.

{[}Laporan sumber{]} Pada blok kesimpulan yang terbaca dengan urutan dua
kolom, klausa tentang \emph{small and medium-sized centres} mencantumkan
27 persen populasi Slovenia dan 18 persen populasi Czechia, tetapi hanya
6 persen di Central France dan 5 persen di Catalonia. TXT juga menyebut
bahwa di Catalonia hampir separuh populasi tinggal di permukiman
non-urban dan 45 persen berada di kota besar atau SMSC yang terkait
dengan kota besar {[}TXT, p.~15, baris 720-727{]}. Awal kalimat angka
tersebut terfragmentasi oleh ekstraksi kolom, sehingga review tidak
menambahkan subjek atau makna yang tidak terbaca.

{[}Laporan sumber{]} Metodologi artikel dipandang sebagai kontribusi
utama karena menyediakan tipologi kota berdasarkan posisi dalam sistem
urban: otonom, berjaringan dengan kota lain, atau teraglomerasi dengan
kota besar. Metode itu dimaksudkan bukan hanya untuk mempelajari
struktur dan komposisi sistem permukiman, tetapi juga sebagai dasar
perbandingan kinerja kota menurut keanggotaan dalam konstelasi yang
berbeda {[}TXT, p.~15, baris 731-755{]}.

{[}Laporan sumber{]} Perbandingan Catalonia, Czechia, Central France,
dan Slovenia menunjukkan variasi ukuran pusat urban, ukuran
mikro-wilayah, dan jenis susunan teritorial. Analisis kinerja menemukan
perbedaan besar antara kota besar dan kota kecil-menengah, tetapi tidak
menemukan perbedaan besar antartipe fungsional SMSC. Penulis menegaskan
bahwa analisis yang lebih maju belum dilakukan {[}TXT, p.~15, baris
756-775 dan 731-743{]}.

\subsection{Contributions}\label{contributions-62}

{[}Laporan sumber{]} Kontribusi yang diklaim artikel adalah:

\begin{itemize}
\tightlist
\item
  kontribusi konseptual berupa penempatan kota kecil dan menengah dalam
  sistem urban yang dibaca melalui fungsi dan relasi, bukan ukuran
  populasi saja;
\item
  kontribusi metodologis berupa prosedur untuk mengidentifikasi pusat
  urban, menyusun mikro-wilayah, menentukan arus signifikan, mengukur
  posisi fungsional, dan membedakan susunan otonom, aglomeratif, serta
  berjaringan;
\item
  kontribusi empiris berupa penerapan dan perbandingan prosedur pada
  empat wilayah atau negara Eropa;
\item
  kontribusi analitis berupa dasar untuk membandingkan kinerja populasi
  dan pekerjaan menurut konstelasi kota dan pusat urban {[}TXT, p.~3-4
  dan p.~15, baris 198-257 dan 742-755{]}.
\end{itemize}

{[}Interpretasi penulis{]} Nilai tambah metode terletak pada upaya
menangkap sistem yang dapat sekaligus memiliki hubungan hierarkis dengan
kota besar dan hubungan jaringan antarkota kecil-menengah. Peta,
klasifikasi, dan ringkasan regional menjadi bentuk keluaran untuk
membaca variasi tersebut {[}TXT, Figure 1-6, p.~5, 9, 12-13, baris
260-265, 465-472, dan 595-648{]}.

{[}Inferensi review{]} Kontribusi ini paling kuat sebagai perangkat
pemetaan dan kategorisasi komparatif. Replikasinya tetap memerlukan
keputusan kontekstual mengenai ambang pekerjaan, populasi, dan arus
signifikan, sehingga tipologi tidak dapat dianggap sepenuhnya bebas dari
diskresi peneliti.

\subsection{Research Gap}\label{research-gap-62}

{[}Laporan sumber{]} Celah awal yang diidentifikasi adalah kurangnya
pengetahuan tentang posisi kota kecil dan menengah dalam sistem urban,
khususnya apakah mereka otonom, berjaringan dengan kota lain, atau
teraglomerasi dalam sistem hierarkis yang didominasi kota besar {[}TXT,
p.~1-3, baris 55-74 dan 138-191{]}. Literatur \emph{borrowed size} juga
disebut belum mengeksplorasi fungsi dan kinerja kota kecil dalam
kaitannya dengan pasar kerja lokal dan penyediaan layanan lokal {[}TXT,
p.~2, baris 101-121{]}.

{[}Laporan sumber{]} Artikel menutup sebagian celah tersebut melalui
tipologi fungsional-spasial dan perbandingan empat sistem permukiman.
Namun, penulis menyatakan bahwa masih diperlukan analisis yang lebih
maju dan bernuansa mengenai kinerja kota menurut posisi dalam susunan
teritorial serta faktor ekonomi regional dan lokal, budaya lokal,
pengaturan kelembagaan, dan modal manusia serta sosial {[}TXT, p.~14-15,
baris 689-704 dan 731-743{]}.

{[}Inferensi review{]} Celah yang masih terbuka mencakup pengujian
kausal atas apakah jaringan atau kedekatan dengan kota besar benar-benar
meningkatkan kinerja, pengukuran relasi selain komuter, dan harmonisasi
unit serta ambang lintas negara. Inferensi ini didasarkan pada sifat
deskriptif analisis, penggunaan satu jenis relasi utama, dan perbedaan
unit serta ambang yang dilaporkan dalam TXT; artikel tidak menyajikannya
sebagai uji tersendiri.

\subsection{Limitations}\label{limitations-62}

{[}Laporan sumber{]} Keterbatasan dan pembatasan analisis yang disebut
atau terlihat jelas dalam artikel adalah:

\begin{itemize}
\tightlist
\item
  Populasi dan pekerjaan tidak menggambarkan keseluruhan fungsi
  permukiman. Penulis mengakui bahwa pekerjaan dan populasi hanya
  memberi perspektif awal dan gambaran luas tentang sistem permukiman
  {[}TXT, p.~4, baris 198-209{]}.
\item
  Analisis relasi hanya menggunakan komuter kerja, sehingga
  mengabstraksikan keragaman hubungan transportasi, telekomunikasi,
  pendidikan, belanja, rekreasi, informasi, dan barang. Penulis
  berasumsi bahwa sebagian relasi lain berkaitan dengan komuter, tetapi
  mengakui adanya multiplexity {[}TXT, p.~4, baris 210-244{]}.
\item
  Ambang bersifat diskresioner atau arbitrer dan dapat berbeda menurut
  konteks nasional: 1.000 pekerjaan, 6.000 penduduk, 5 persen penduduk
  aktif secara ekonomi, dan 1 persen pekerjaan tujuan untuk Czechia,
  serta ambang berbeda di wilayah lain {[}TXT, p.~6-9, baris 290-311,
  351-381, dan 480-492{]}.
\item
  Unit geografis berbeda dalam ukuran spasial dan populasi. LAU2
  digunakan di sebagian besar kasus, sedangkan \emph{unités urbaines}
  digunakan di France, sehingga perbandingan langsung harus
  memperhatikan perbedaan skala unit {[}TXT, p.~9-10, baris 473-525{]}.
\item
  Prosedur menghubungkan setiap sumber ke satu tujuan arus terkuat dan
  tidak ditujukan untuk mengidentifikasi wilayah perjalanan kerja yang
  sepenuhnya mandiri. Penulis sengaja mengabaikan
  \emph{self-containment} tertentu untuk menghindari bias terhadap
  kawasan besar yang saling terhubung {[}TXT, p.~6-7, baris 329-363{]}.
\item
  Analisis kinerja bersifat sederhana, menggunakan periode yang tidak
  sepenuhnya sama antarwilayah, dan tidak dapat memasukkan Catalonia
  karena data kinerja tidak tersedia {[}TXT, p.~12 dan 14, baris 620-625
  dan 689-704{]}.
\item
  Artikel hanya menampilkan empat dari sepuluh wilayah atau negara yang
  dianalisis dalam kerangka ESPON TOWN; alasan pemilihan empat kasus
  tidak disebutkan dalam file {[}TXT, p.~9, baris 490-495{]}.
\end{itemize}

{[}Inferensi review{]} Algoritme arus signifikan yang hanya menilai lima
arus keluar tertinggi dan klasifikasi yang bergantung pada ambang dapat
menyederhanakan jaringan dengan banyak relasi lemah atau menengah.
Sejauh yang tersedia dalam TXT, penulis tidak menguji sensitivitas hasil
terhadap perubahan ambang tersebut.

\subsection{Challenges}\label{challenges-62}

{[}Laporan sumber{]} Tantangan metodologis yang dihadapi penulis
meliputi penentuan jumlah pekerjaan minimum agar sebuah permukiman layak
disebut pusat pekerjaan, identifikasi tujuan tidak langsung, penghapusan
\emph{proto-microregions} yang terlalu kecil, dan konsolidasi wilayah
yang terfragmentasi {[}TXT, p.~5-6, baris 290-327{]}.

{[}Laporan sumber{]} Tidak ada garis pemisah yang jelas untuk menentukan
apakah arus penting bagi pasar kerja tujuan. Karena itu, ambang 1 persen
dipilih secara arbitrer untuk Czechia, sementara ambang di negara lain
dapat berbeda {[}TXT, p.~7, baris 369-381{]}.

{[}Laporan sumber{]} Klasifikasi juga menghadapi kasus ketika ukuran
populasi dan posisi fungsional tidak sejalan. Di Czechia terdapat pusat
berpenduduk di bawah 50.000 dengan posisi fungsional tinggi dan pusat
berpenduduk di atas 50.000 dengan posisi rendah. Penulis menilai
kasus-kasus tersebut secara individual berdasarkan konteks spasial dan
historis {[}TXT, p.~8, baris 425-440{]}.

{[}Laporan sumber{]} Dua pengecualian dijelaskan: Náchod yang
berpenduduk sekitar 20.000 tidak dimasukkan sebagai kota besar meskipun
posisi fungsionalnya tinggi, sedangkan Mladá Boleslav yang berpenduduk
sekitar 44.000 dimasukkan karena merupakan pusat pekerjaan yang kuat.
Ústí nad Labem yang berpenduduk hampir 100.000 diklasifikasikan sebagai
kota besar otonom karena tidak menjadi tujuan arus signifikan dari pusat
mikro-regional lain, dengan konteks pegunungan, perbatasan Jerman,
aksesibilitas yang rendah, dan otonomi pasar kerja kota-kota sekitar
sebagai penjelasan {[}TXT, Catatan 2, p.~16, baris 784-825{]}.

{[}Interpretasi penulis{]} Tantangan substantifnya adalah bahwa kinerja
kota dipengaruhi banyak faktor dan sangat bervariasi pada tingkat pusat
urban individual. Posisi dalam sistem urban hanya salah satu kondisi
struktural yang mungkin berpengaruh {[}TXT, p.~14, baris 689-704{]}.

\subsection{Practical Implications}\label{practical-implications-62}

{[}Laporan sumber{]} Artikel menyatakan bahwa metodologi dapat digunakan
untuk mempelajari struktur dan komposisi berbagai sistem permukiman
serta menjadi dasar perbandingan kinerja kota menurut keanggotaan dalam
konstelasi kota dan kota besar {[}TXT, p.~15, baris 742-755{]}.

{[}Inferensi review{]} Bagi analisis perencanaan wilayah, pendekatan ini
dapat mendorong pemetaan pusat lokal berdasarkan pekerjaan dan arus
komuter, bukan hanya ukuran penduduk atau batas administratif. Hasilnya
juga dapat membantu membedakan konteks kota otonom, kota yang
teraglomerasi, dan kota berjaringan sebelum merumuskan intervensi
wilayah.

{[}Inferensi review{]} Ambang tidak sebaiknya dipindahkan secara mekanis
antarnegara atau antarwilayah. Sumber sendiri menekankan bahwa ambang
pekerjaan, populasi, dan signifikansi arus dapat berbeda sesuai konteks
nasional {[}TXT, p.~6-9, baris 300-305, 369-381, dan 480-492{]}.

{[}Batas informasi{]} Rekomendasi kebijakan yang konkret, seperti
program investasi, bentuk kelembagaan, atau skema pembiayaan untuk tipe
kota tertentu, tidak disebutkan dalam file.

\subsection{Applications}\label{applications-62}

{[}Laporan sumber{]} Aplikasi empiris utama adalah penerapan rinci pada
sistem urban Czechia dan penerapan komparatif pada Catalonia, Czechia,
Central France, dan Slovenia dalam proyek ESPON TOWN. Keluaran aplikasi
meliputi identifikasi MRC, klasifikasi pusat, peta susunan teritorial,
serta perbandingan populasi dan pekerjaan {[}TXT, p.~2, baris 82-98;
p.~9-14, baris 490-704{]}.

{[}Laporan sumber{]} Penulis menyatakan bahwa alat tersebut dapat
diterapkan untuk mempelajari struktur berbagai sistem permukiman dan
membandingkan performa kota menurut susunan teritorialnya {[}TXT, p.~15,
baris 742-755{]}.

{[}Inferensi review{]} Dalam batas TXT, aplikasi yang benar-benar
terdokumentasi bersifat analitis dan komparatif. Implementasi perangkat
lunak, protokol operasional yang dapat langsung dijalankan, atau
penggunaan untuk wilayah di luar empat kasus tidak disebutkan secara
rinci dalam file.

\subsection{Future Research
(Optional)}\label{future-research-optional-62}

{[}Laporan sumber{]} Penulis menyarankan analisis multivariat untuk
mengidentifikasi kondisi struktural utama yang memengaruhi kinerja kota,
karena posisi dalam sistem urban hanya salah satu faktor. Penulis juga
menyatakan perlunya studi kasus kualitatif pada kota individual untuk
melengkapi perspektif kuantitatif {[}TXT, p.~14, baris 689-704{]}.

{[}Laporan sumber{]} Analisis lanjutan perlu mempertimbangkan kinerja
kota dalam kaitannya dengan posisi di susunan teritorial serta ekonomi
regional dan lokal, budaya lokal, pengaturan kelembagaan, dan modal
manusia serta sosial. Artikel tidak memberikan desain penelitian, ukuran
sampel, atau jadwal untuk agenda tersebut {[}TXT, p.~15, baris
731-743{]}.

\subsection{Provenance Notes}\label{provenance-notes-62}

\begin{itemize}
\tightlist
\item
  Sumber tunggal yang digunakan adalah file TXT B19 pada direktori
  \texttt{source/journal}; tidak ada akses langsung atau informasi
  tambahan dari PDF, web, artikel lain, atau review kode lain yang
  ditambahkan. Metadata PDF yang disebut dalam review hanya disalin dari
  blok metadata yang sudah ada di TXT.
\item
  Seluruh teks yang tersedia dalam TXT dibaca, dari baris 1 sampai 910,
  termasuk blok metadata ekstraksi, abstrak, isi artikel, catatan kaki,
  acknowledgments, dan daftar referensi.
\item
  Blok metadata menyatakan bahwa sumber PDF memiliki 18 halaman,
  diekstrak pada 2026-08-31 dengan \texttt{pdftotext\ -layout}; metadata
  juga mencantumkan judul, DOI pada badan artikel, dan ukuran file
  {[}TXT, baris 1-26 dan 75-76{]}.
\item
  Locator dalam review menggunakan nomor halaman artikel yang tercetak
  atau tersirat pada pemisah halaman TXT, serta nomor baris TXT. Locator
  \texttt{Tabel\ 1}, \texttt{Tabel\ 2}, \texttt{Tabel\ 3},
  \texttt{Figure\ 1-7}, dan \texttt{Catatan\ 1-2} mengikuti label yang
  terbaca di dalam file.
\item
  Ekstraksi mempertahankan artefak layout dua kolom, karakter encoding
  pada nama penulis dan beberapa nama tempat, serta representasi tabel
  dan gambar yang tidak selalu utuh. Angka atau hubungan yang tidak
  terbaca tidak direkonstruksi dari sumber lain.
\item
  Pada beberapa bagian, terutama kesimpulan halaman 15, urutan teks dua
  kolom membuat klausa terpisah. Review menandai fragmentasi tersebut
  dan tidak mengisi bagian yang tidak dapat dipastikan.
\item
  Label \texttt{{[}Laporan\ sumber{]}} merangkum apa yang dinyatakan
  atau ditampilkan artikel; \texttt{{[}Interpretasi\ penulis{]}}
  merangkum penjelasan atau posisi analitis penulis;
  \texttt{{[}Inferensi\ review{]}} adalah pembacaan kritis yang
  diturunkan dari TXT dan bukan klaim eksplisit artikel.
\item
  Informasi yang tidak tersedia atau tidak berlaku telah dinyatakan
  secara eksplisit sebagai \texttt{Tidak\ disebutkan\ dalam\ file} atau
  \texttt{Tidak\ berlaku}, sesuai batas sumber.
\end{itemize}

\section{Review B20}\label{review-b20}

\textbf{Identitas sumber}

\begin{itemize}
\tightlist
\item
  Kode: B20.
\item
  Judul: \emph{Grid and shake: spatial aggregation and the robustness of
  regionally estimated elasticities}.
\item
  Penulis: Gábor Békés dan Péter Harasztosi.
\item
  Jenis yang tercantum dalam TXT: \texttt{ORIGINAL\ PAPER}. Status
  peer-review tidak diverifikasi secara terpisah.
\item
  Jurnal: \emph{The Annals of Regional Science} atau singkatan yang
  tercetak, \emph{Ann Reg Sci}, volume 60, halaman 143-170, tahun 2018.
  Nomor isu tidak disebutkan dalam file TXT.
\item
  DOI: \texttt{10.1007/s00168-017-0849-y}; URL yang tercantum:
  \texttt{https://doi.org/10.1007/s00168-017-0849-y}.
\item
  Riwayat naskah: diterima 29 Juli 2016, diterima untuk publikasi 26
  September 2017, dan dipublikasikan secara daring 7 Oktober 2017.
\item
  Klasifikasi JEL: R12, R30, C15.
\item
  Afiliasi yang tercantum: Institute of Economics of CERS-HAS and CEPR,
  Central European University, Budapest, untuk Békés; dan Joint Research
  Centre, Ispra, Italy, untuk Harasztosi.
\item
  Bahan yang ditinjau adalah file
  \texttt{/Users/mac/Documents/Mac/{[}2{]}\ Obsidian\ Vault/zettelkasten/source/journal/B20-bekes-harastosi-2018-grid-and-shake-spatial-aggregation-and-the-robustness-of-regionally-estimated-elasticities-10.1007-s00168-017-0849-y.txt}.
  Metadata TXT menyatakan sumber PDF terdiri atas 28 halaman dan teks
  diekstraksi dengan \texttt{pdftotext\ -layout} pada 31 Agustus 2026.
\end{itemize}

\subsection{Summarized Abstract}\label{summarized-abstract-63}

\textbf{Laporan sumber.} Artikel mengusulkan metode yang sederhana dan
transparan untuk mengukur robustnes spasial koefisien yang diestimasi
secara regional. Fokusnya adalah peran distrik administratif dan ukuran
wilayah, terutama ketika variabel dibandingkan pada unit spasial yang
telah diagregasi. Penulis menyatakan bahwa prosedur tersebut memperbaiki
metode yang ada, khususnya ketika unit spasial bersifat heterogen.
Ilustrasi empiris menggunakan data Hungaria untuk membandingkan estimasi
eksternalitas aglomerasi pada beberapa tingkat agregasi. Kesimpulan
abstraknya adalah bahwa cara agregasi spasial tampak sama pentingnya
dengan spesifikasi model ekonometrik (TXT, hlm. 143, abstrak, baris
50-60).

\textbf{Interpretasi penulis.} Kontribusi utama yang diklaim bukan
estimator baru untuk elastisitas aglomerasi, melainkan prosedur untuk
menghasilkan distribusi koefisien dari berbagai realisasi unit spasial
buatan. Distribusi tersebut memungkinkan perbandingan antara unit
administratif, unit buatan dengan ukuran serupa, dan ukuran grid yang
berbeda (TXT, hlm. 143-145, baris 93-146 dan 158-163).

\textbf{Inferensi review.} Abstrak mendukung pembacaan bahwa objek utama
artikel adalah sensitivitas hasil terhadap pilihan unit spasial. Ia
tidak dengan sendirinya membuktikan besar atau arah kausal eksternalitas
aglomerasi; angka elastisitas tetap bergantung pada data Hungaria, model
yang dipakai, dan aturan pembentukan grid yang dijelaskan dalam artikel.

\subsection{Problem Statement}\label{problem-statement-63}

\textbf{Laporan sumber.} Penulis memulai dari masalah bahwa studi
regional sering memilih kota, kabupaten, atau wilayah secara arbitrer.
Estimasi kemudian dapat berubah karena karakteristik unit spasial,
ukuran wilayah, letak batas administratif, serta jarak tempat
berlangsungnya eksternalitas. Ketidakcocokan antara unit observasi dan
jangkauan spasial fenomena dapat menghasilkan kesalahan pengukuran
spasial dan korelasi spasial pada galat di lokasi yang berdekatan (TXT,
hlm. 144-147, baris 88-121 dan 224-248).

Penulis juga menekankan bahwa ukuran unit tidak hanya tidak seragam,
tetapi distribusinya tidak acak. Unit kecil atau besar dapat mengelompok
karena geografi fisik, sejarah, urbanisasi, pertanian, akses pasar, dan
institusi. Karena faktor-faktor tersebut juga berhubungan dengan
kepadatan, upah, atau produktivitas, hasil estimasi dapat mencerminkan
pola spasial unit administratif, bukan hanya hubungan substantif yang
hendak diukur (TXT, hlm. 147-148, baris 255-312).

Batas administratif memiliki persoalan tambahan karena dapat menampung
kebijakan pasar kerja, pajak, subsidi, dan regulasi. Karena itu,
perbedaan antara estimasi pada wilayah administratif dan wilayah buatan
dengan ukuran rata-rata yang sama dapat digunakan untuk memeriksa apakah
batas administratif berkaitan dengan hasil estimasi (TXT, hlm. 148-149,
baris 315-348).

\textbf{Interpretasi penulis.} Masalah tersebut dirumuskan sebagai
kebutuhan untuk membandingkan estimasi antarunit spasial secara
statistik, bukan sekadar menjalankan satu uji robustnes dengan unit
alternatif. Prosedur yang dibutuhkan harus dapat diskalakan menurut
ukuran unit, memperhitungkan pengelompokan unit kecil dan besar,
membandingkan unit administratif dengan unit buatan, serta menghasilkan
dasar untuk menguji signifikansi perbedaan elastisitas (TXT, hlm.
146-149, baris 249-252, 310-312, 340-364).

\textbf{Inferensi review.} Problem statement artikel terutama merupakan
problem komparabilitas dan sensitivitas pengukuran. Ia belum sama dengan
problem identifikasi kausal. OLS dan Poisson dalam artikel digunakan
untuk melihat bagaimana koefisien berubah ketika unit spasial diubah;
perubahan tersebut tidak secara otomatis dapat ditafsirkan sebagai
perubahan kausal pada eksternalitas atau sebagai efek kebijakan.

\subsection{Objectives}\label{objectives-63}

\textbf{Laporan sumber.} Tujuan yang dinyatakan atau dijalankan dalam
artikel meliputi:

\begin{itemize}
\tightlist
\item
  Memperkenalkan prosedur ``grid and shake'' yang menempatkan grid pada
  peta, menggeser posisinya secara acak, menggabungkan unit dasar ke
  dalam kotak, lalu mengulang regresi untuk memperoleh distribusi
  koefisien (TXT, hlm. 145 dan 149-153, baris 122-146 dan 367-380).
\item
  Menguji bagaimana ukuran unit spasial memengaruhi elastisitas
  aglomerasi pada beberapa ukuran grid buatan (15 x 15 km, 26 x 26 km,
  dan 39 x 39 km) (TXT, hlm. 150-152, baris 425-470).
\item
  Membandingkan unit administratif Hungaria, yaitu munisipalitas dan
  mikro-wilayah, dengan unit buatan yang memiliki ukuran atau jumlah
  unit yang dapat dibandingkan (TXT, hlm. 145-146 dan 157-163, baris
  164-183 dan 751-755).
\item
  Menggunakan dua contoh empiris, yaitu hubungan kepadatan dengan upah
  dan hubungan kepadatan dengan pembentukan perusahaan baru (TXT, hlm.
  153-157, baris 538-547 dan 583-721).
\item
  Menilai apakah perbedaan estimasi antarukuran grid, serta antara unit
  administratif dan buatan, berbeda secara statistik (TXT, hlm. 151-153
  dan 157-161, baris 499-535 dan 806-931).
\item
  Membandingkan prosedur tersebut dengan agregasi berbasis lingkaran,
  spatial atau distributed lags, dan metode dots-to-boxes Briant et
  al.~(2010) (TXT, hlm. 162-163, baris 934-1020).
\end{itemize}

\textbf{Interpretasi penulis.} Tujuan akhirnya adalah menyediakan cara
yang transparan untuk mengevaluasi apakah hasil regional merupakan
konsekuensi dari ukuran unit, realisasi batas, atau peran batas
administratif, sambil tetap memungkinkan perbandingan lintas dataset dan
lintas negara (TXT, hlm. 144-146 dan 163-164, baris 93-98, 175-183, dan
1070-1076).

\subsection{Literature Survey}\label{literature-survey-63}

\textbf{Laporan sumber.} Literatur yang dipakai penulis berangkat dari
ekonomi geografi dan eksternalitas aglomerasi. Artikel merujuk pada
Marshall untuk eksternalitas, Ciccone dan Hall untuk hubungan kepadatan
dengan produktivitas, Ács dan Armington untuk pembentukan perusahaan,
serta Melo et al.~untuk meta-analisis estimasi ekonomi aglomerasi.
Penulis menyatakan bahwa literatur tersebut umumnya menemukan
produktivitas yang lebih tinggi pada wilayah yang lebih padat dan
menghubungkannya dengan upah lokal (TXT, hlm. 153-154, baris 583-596;
daftar referensi baris 1187-1245).

Argumen tentang skala dan jarak dibangun dengan merujuk pada
Fotheringham dan Wong mengenai \emph{modifiable areal unit problem} atau
MAUP, Dewhurst dan McCann mengenai ukuran wilayah dan spesialisasi,
Burger et al.~mengenai tingkat geografis eksternalitas, serta Larsson
mengenai peluruhan elastisitas kepadatan-upah menurut resolusi spasial.
Penulis juga memberi ilustrasi korelasi kepadatan penduduk dan PDB per
kapita di sejumlah negara Eropa: 0,74 pada NUTS1 dan 0,54 pada NUTS3.
Data ilustratif itu disebut berasal dari Eurostat, dengan PDB per kapita
tahun 2007 dan kepadatan penduduk rata-rata 2006-2008 pada delapan
negara (TXT, hlm. 146-147, baris 238-276).

Untuk pendekatan metodologis, artikel membahas regresi spasial dan
spatial lag, agregasi berbasis seed dan tetangga oleh Amrhein dan
Flowerdew, dots-to-boxes oleh Briant et al., simulasi Menon, agregasi
berbasis lingkaran, serta spatial dynamic atau distributed lags.
Masing-masing dibahas menurut kemampuan menangani jarak, ukuran unit,
bentuk wilayah, pengelompokan unit dasar, dan pengujian perbedaan
koefisien (TXT, hlm. 145 dan 162-163, baris 147-157 dan 949-1020).

\textbf{Interpretasi penulis.} Penulis menempatkan grid and shake
sebagai pelengkap atau alternatif bagi pendekatan tersebut. Klaim
pembeda utamanya adalah kemampuan membentuk unit dengan ukuran yang
dikendalikan, mengulang realisasi batas, menangani unit dasar yang
heterogen dan tidak acak, menghitung semua unit dasar satu kali, serta
membangun distribusi untuk inferensi (TXT, hlm. 145, 162-164, baris
158-163 dan 1023-1069).

\textbf{Inferensi review.} Survei ini berbentuk tinjauan naratif yang
dipilih untuk membangun motivasi metode, bukan tinjauan sistematis. TXT
tidak memuat protokol pencarian, kriteria inklusi, basis data pencarian,
atau penilaian kualitas korpus literatur. Oleh karena itu, literatur
yang disebut dapat dibaca sebagai landasan yang dipilih penulis, bukan
sebagai inventaris lengkap penelitian tentang MAUP, agregasi spasial,
atau elastisitas aglomerasi.

\subsection{Method Used}\label{method-used-63}

\textbf{Laporan sumber: pembentukan unit spasial.} Penulis mendigitalkan
peta Hungaria dengan marker yang disusun seperti raster pada unit dasar
munisipalitas. Marker ditempatkan pada jarak sekitar 800 meter; 153.113
marker digunakan untuk wilayah Hungaria seluas 93.030 km2. Marker
dicocokkan dengan poligon munisipalitas sehingga dapat menjadi dasar
pengelompokan ke unit yang lebih besar. Sistem koordinat yang disebut
adalah proyeksi Mercator WGS84 EPSG:41001 (TXT, hlm. 150, baris 391-402
dan catatan kaki 10-11).

Grid persegi berukuran sisi L diletakkan di atas marker. Setiap marker
diberi indeks kotak dengan \texttt{gx\ =\ int(x/L)\ +\ 1} dan
\texttt{gy\ =\ int(y/L)\ +\ 1}. Jika satu munisipalitas memotong
beberapa kotak, munisipalitas dimasukkan ke kotak yang memuat mayoritas
markernya, yang dianggap mewakili bagian wilayah terbesar (TXT, hlm.
150, baris 403-418).

Tiga ukuran grid digunakan. Grid 15 x 15 km menghasilkan sekitar 440
unit dan dimaksudkan meniru unit yang seukuran kota besar. Grid 26 x 26
km dipilih untuk mendekati 150 mikro-wilayah, sedangkan grid 39 x 39 km
menghasilkan sekitar 75 unit dan dimaksudkan untuk mensimulasikan ukuran
rata-rata lima mikro-wilayah terbesar (TXT, hlm. 150-151, baris 425-458;
ringkasan area pada Tabel 3, hlm. 164, baris 1083-1094).

Untuk mengurangi efek batas negara yang bergerigi, unit hasil grid yang
luasnya kurang dari ambang 10 persen dari luas L x L km digabungkan
dengan wilayah tetangga yang dipilih secara acak. Penulis mengakui bahwa
persoalan ini lebih penting pada negara kecil, ukuran grid besar, dan
garis batas yang tidak horizontal atau vertikal (TXT, hlm. 152, baris
489-496).

\textbf{Laporan sumber: pengacakan dan inferensi.} Titik awal grid
digeser dengan menambahkan dua angka acak independen, xi dan zeta, yang
masing-masing seragam pada interval 0 sampai L. Setiap posisi
menghasilkan pengelompokan munisipalitas yang dapat berbeda. Penulis
menjelaskan bahwa koefisien regresi dihimpun dari realisasi tersebut
untuk membentuk distribusi dan dibandingkan dengan koefisien unit
administratif atau distribusi dari ukuran lain (TXT, hlm. 152-153, baris
499-535).

Jumlah replikasi yang muncul secara konsisten dalam uraian metode,
gambar, Tabel 1, Tabel 2, dan penjelasan aplikasi adalah 300. Namun,
langkah prosedur pada bagian 3.4 menyebut bahwa penggeseran grid diulang
``a thousand times''. TXT tidak menjelaskan apakah ini perubahan desain,
kesalahan penyalinan, atau langkah berbeda. Review ini mempertahankan
kedua informasi tersebut dan menggunakan angka 300 ketika merujuk pada
tabel dan hasil utama (TXT, hlm. 145, 152-153, 154-159, baris 141-146,
506-535, 643-651, 699-707, dan 724-735).

Untuk pemeriksaan distribusi yang memasukkan ketidakpastian estimasi,
penulis mengambil 50 angka acak dari distribusi normal di sekitar setiap
koefisien dan standard error pada masing-masing dari 300 realisasi.
Distribusi perluasan tersebut terdiri atas 50 x 300 angka dan
dibandingkan dengan distribusi yang terkait dengan batas administratif
menggunakan uji Kolmogorov-Smirnov dua sisi (TXT, hlm. 161, baris
896-918; Appendix, Gambar 12, baris 1133-1137).

\textbf{Laporan sumber: model empiris.} Model pertama adalah regresi OLS
penampang lintang tahun 1999:

\texttt{Wage\_i\ =\ beta\_1\ DENSITY\_i\ +\ beta\_2\ EDUC\_i\ +\ nu\_i}

Wage adalah upah rata-rata pekerja pada unit spasial; DENSITY adalah
kepadatan penduduk usia kerja; EDUC adalah proksi pendidikan. Teks
menyatakan bahwa seluruh variabel digunakan dalam logaritma. Standard
error yang ditampilkan adalah robust. Variabel dan koefisien dihitung
ulang pada setiap tingkat agregasi; rata-rata tahun sekolah diagregasi
dengan pembobot frekuensi berdasarkan jumlah individu (TXT, hlm.
153-155, baris 597-625 dan Tabel 1 baris 631-651).

Model kedua menggunakan Poisson untuk jumlah perusahaan baru, dengan
kepadatan penduduk dan upah sebagai penjelas:

\texttt{Pr(y\_i\ =\ N)\ =\ f(lambda\_i)}

\texttt{lambda\_i\ =\ beta\_1\ DENSITY\_i\ +\ beta\_2\ Wage\_i}

Artikel tidak menjelaskan lebih lanjut dalam TXT bentuk lengkap fungsi
link, prosedur penanganan dispersi, atau definisi operasional rinci
jumlah perusahaan baru. Tabel 2 menyatakan semua variabel dalam
logaritma dan menampilkan robust standard error untuk unit administratif
serta rata-rata standard error dan simpangan baku koefisien untuk 300
replikasi grid (TXT, hlm. 155-156, baris 660-721 dan Tabel 2 baris
687-707).

\textbf{Inferensi review.} Randomisasi artikel mengubah pembagian ruang
dengan mempertahankan data dasar dan model regresi. Dengan demikian,
distribusi koefisien terutama menggambarkan sensitivitas terhadap
realisasi agregasi yang dihasilkan oleh aturan grid, bukan
ketidakpastian sampling baru atas perusahaan atau munisipalitas.
Kesimpulan robustnes karenanya bersyarat pada proyeksi, kerapatan
marker, ukuran grid, ambang koreksi batas, jumlah replikasi, dan
spesifikasi model yang dipilih.

\subsection{Dataset}\label{dataset-63}

\textbf{Laporan sumber.} Data utama berasal dari data neraca perusahaan
tingkat perusahaan milik Central Statistical Office Hungaria. Cakupannya
adalah perusahaan manufaktur dengan pembukuan ganda dan sedikitnya lima
karyawan. Informasi tahun 1999 digunakan untuk penampang lintang OLS,
sedangkan informasi tahun 1998 juga tersedia untuk membentuk variabel
pembentukan perusahaan baru. Data keuangan perusahaan dicocokkan dengan
lokasi pada tingkat munisipalitas, kemudian dipasangkan dengan informasi
ukuran dan penduduk munisipalitas dari basis data T-STAR (TXT, hlm. 153,
baris 550-558).

Upah dihitung dari total biaya upah dan informasi pekerja penuh waktu.
Proksi modal manusia berasal dari Hungarian Labor Force Survey, yang
memuat sampel pekerja sektor manufaktur dan tingkat pendidikan
tertinggi. Data survei dapat dihubungkan ke munisipalitas, tetapi
pemberi kerja aktual tidak diidentifikasi, sehingga informasi pekerja
tidak dapat dihubungkan langsung ke data perusahaan. Rata-rata tahun
sekolah pada tingkat munisipalitas digunakan sebagai proksi kualitas
tenaga kerja dan disebut bervariasi antara 7 dan 14 tahun (TXT, hlm.
153-154, baris 559-580).

Unit spasial dasar terdiri atas 3.125 munisipalitas. Struktur
administratif yang disebut juga mencakup 150 mikro-wilayah atau
kistérség, 20 wilayah NUTS3, dan tujuh wilayah NUTS2. Untuk analisis,
grid buatan memiliki sekitar 440 unit pada ukuran 15 km, 150 unit pada
ukuran 26 km, dan 75 unit pada ukuran 39 km. Angka luas minimum,
rata-rata, dan maksimum tiap jenis unit dirangkum pada Tabel 3; angka
untuk grid adalah rata-rata 300 replikasi (TXT, hlm. 145 dan 164, baris
164-173 dan 1083-1094).

Artikel juga menyebut data Eurostat untuk ilustrasi hubungan kepadatan
dan PDB per kapita pada delapan negara Eropa. Data tersebut berfungsi
sebagai motivasi mengenai perbedaan korelasi menurut NUTS1 dan NUTS3,
bukan sebagai dataset utama untuk estimasi grid and shake dalam dua
aplikasi Hungaria (TXT, hlm. 146-147, baris 238-270).

\textbf{Informasi yang tidak tersedia.} TXT tidak menyebut jumlah
perusahaan yang masuk sebelum agregasi, ukuran sampel numerik Hungarian
Labor Force Survey, skema sampling survei secara rinci, nama berkas data
mentah, atau tautan akses data. TXT juga tidak memberi definisi lengkap
pembentukan perusahaan baru selain keterkaitannya dengan informasi tahun
1998 dan 1999.

\subsection{Population / Sample}\label{population-sample-63}

\textbf{Laporan sumber.} Populasi empiris yang didefinisikan penulis
adalah perusahaan manufaktur di Hungaria yang melakukan pembukuan ganda
dan mempunyai sedikitnya lima karyawan. Namun, unit observasi regresi
bukan lagi perusahaan individual, melainkan unit spasial yang berisi
informasi perusahaan, penduduk, luas wilayah, upah, dan pendidikan
setelah agregasi (TXT, hlm. 153-155, baris 552-568 dan 597-606).

Tabel 1 dan Tabel 2 melaporkan 2.872 observasi pada tingkat
munisipalitas, 149 pada tingkat mikro-wilayah, 437 pada grid 15 km, 150
pada grid 26 km, dan 73 pada grid 39 km. Jumlah unit administratif awal
yang disebut adalah 3.125 munisipalitas dan 150 mikro-wilayah. Penulis
menjelaskan bahwa sebagian munisipalitas tidak memiliki aktivitas
manufaktur yang dapat diidentifikasi, sehingga tidak semuanya masuk
dalam regresi (TXT, hlm. 155-156, Tabel 1 dan Tabel 2, baris 631-651 dan
687-707).

Terdapat perbedaan internal pada jumlah observasi munisipalitas. Tabel 1
dan Tabel 2 menulis 2.872, tetapi paragraf penjelas setelah Tabel 1
menyatakan 2.873. TXT tidak memberikan rekonsiliasi. Sampel pekerja dari
Hungarian Labor Force Survey disebutkan secara substantif, tetapi ukuran
numerik dan rincian pemilihannya tidak disebutkan dalam file (TXT, hlm.
155, baris 655-659; hlm. 153-154, baris 560-568).

\textbf{Inferensi review.} Angka 2.872 atau 2.873 harus dibaca sebagai
jumlah unit spasial yang tersedia untuk regresi tingkat munisipalitas,
bukan jumlah perusahaan atau jumlah individu. Replikasi 300 pada grid
juga bukan 300 sampel independen; itu adalah 300 realisasi pembagian
wilayah yang diterapkan pada basis data yang sama. Perbedaan jumlah
observasi dan tidak tersedianya ukuran sampel HLFS membatasi audit
eksternal atas komposisi sampel.

\subsection{Dependent Variables}\label{dependent-variables-63}

\textbf{Laporan sumber.} Pada Model 1, variabel dependen adalah upah
rata-rata pekerja pada unit spasial. Upah dihitung dari total biaya upah
dan informasi pekerja penuh waktu, lalu digunakan dalam logaritma.
Koefisien kepadatan ditafsirkan penulis sebagai elastisitas aglomerasi
upah. Kepadatan penduduk usia kerja per km2 adalah variabel penjelas
utama, sedangkan rata-rata tahun sekolah adalah kontrol (TXT, hlm.
153-155, baris 559-580 dan 597-606).

Pada Model 2, variabel dependen adalah jumlah perusahaan baru atau
pembentukan perusahaan baru di unit spasial. Modelnya adalah Poisson;
kepadatan menjadi variabel penjelas utama dan upah menjadi proksi
kondisi pasar tenaga kerja. Data tahun 1998 digunakan bersama data tahun
1999 untuk membangun variabel pembentukan perusahaan baru, tetapi rumus
atau aturan klasifikasi perusahaan baru tidak dijelaskan lebih rinci
(TXT, hlm. 153 dan 155-156, baris 552-558 dan 660-721).

\textbf{Informasi yang tidak tersedia.} TXT tidak menyebut mata uang
atau periode pembayaran upah secara rinci, tidak memberikan definisi
formal jumlah perusahaan baru, dan tidak menjelaskan penanganan nilai
nol atau transformasi variabel dependen pada Model 2. Catatan tabel
menyatakan semua variabel dalam logaritma, tetapi rincian implementasi
logaritma untuk jumlah perusahaan pada model Poisson tidak dijabarkan.
Karena itu, review tidak mengasumsikan bentuk transformasi di luar
pernyataan tersebut.

\subsection{Results}\label{results-63}

\textbf{Laporan sumber: regresi upah.} Tabel 1 melaporkan koefisien
kepadatan pada model upah sebesar 0,055 dengan standard error 0,007
untuk munisipalitas dan 0,128 dengan standard error 0,021 untuk
mikro-wilayah. Untuk grid buatan, rata-rata koefisien dari 300 replikasi
adalah 0,082 pada grid 15 km, 0,116 pada grid 26 km, dan 0,137 pada grid
39 km. Rata-rata standard error masing-masing adalah 0,011, 0,019, dan
0,026, sedangkan simpangan baku koefisien lintas replikasi masing-masing
adalah 0,006, 0,008, dan 0,015 (TXT, hlm. 154-155, Tabel 1, baris
631-651).

Koefisien kontrol tahun sekolah pada kolom munisipalitas dan
mikro-wilayah adalah 0,966 dan 0,820. Pada grid 15 km, 26 km, dan 39 km,
rata-ratanya adalah 1,098, 1,123, dan 0,979; simpangan baku koefisiennya
adalah 0,121, 0,219, dan 0,510. R-squared berturut-turut untuk
munisipalitas, mikro-wilayah, serta grid 15, 26, dan 39 km adalah 0,244,
0,383, 0,240, 0,361, dan 0,393. Tabel melaporkan jumlah observasi 2.872,
149, 437, 150, dan 73 (TXT, hlm. 154-155, Tabel 1, baris 637-651).

Narasi penulis menyebut urutan koefisien kepadatan sebagai 5,5 persen
pada munisipalitas, 7,7 persen pada unit buatan kecil, 11,8 persen pada
mikro-wilayah, dan 12,7 persen pada unit buatan sedang. Angka ini tidak
sepenuhnya sama dengan angka Tabel 1, yang masing-masing menunjukkan
0,055, 0,082, 0,128, dan 0,116. TXT tidak menjelaskan perbedaan antara
angka narasi dan tabel; keduanya dicatat terpisah di sini, bukan
diselaraskan secara dugaan (TXT, hlm. 154, baris 612-619; Tabel 1, baris
637-645).

\textbf{Laporan sumber: regresi pembentukan perusahaan.} Pada Tabel 2,
koefisien kepadatan adalah 1,377 dengan standard error 0,079 pada
munisipalitas dan 0,917 dengan standard error 0,172 pada mikro-wilayah.
Rata-rata koefisien dari grid 15, 26, dan 39 km adalah 1,321, 1,257, dan
1,146; rata-rata standard errornya adalah 0,074, 0,096, dan 0,219;
simpangan baku koefisien lintas 300 replikasi adalah 0,027, 0,039, dan
0,085. Koefisien upah adalah 0,608 dengan standard error 0,200 pada
munisipalitas, 0,554 dengan standard error 0,716 pada mikro-wilayah,
serta rata-rata 0,526, 0,392, dan 1,664 pada grid 15, 26, dan 39 km.
Simpangan baku koefisien upah pada grid tersebut adalah 0,110, 0,152,
dan 0,358 (TXT, hlm. 155-156, Tabel 2, baris 687-707).

Jumlah observasi pada Tabel 2 sama dengan Tabel 1. Pseudo-R-squared yang
tercantum adalah 0,467 dan 0,379 untuk dua unit administratif, serta
0,605, 0,615, dan 0,614 untuk grid 15, 26, dan 39 km. Penulis menyatakan
bahwa koefisien kepadatan lebih besar dari satu pada seluruh estimasi
dan menganggapnya sebagai hubungan bersih yang positif antara penduduk
dan pembentukan perusahaan. Pernyataan naratif ini tidak sepenuhnya
selaras dengan koefisien kepadatan 0,917 pada baris mikro-wilayah di
Tabel 2; TXT tidak memberi penjelasan atas perbedaannya (TXT, hlm.
155-156, baris 711-721 dan Tabel 2, baris 693-707).

\textbf{Laporan sumber: perbandingan fragmentasi.} Untuk OLS upah dan
Poisson pembentukan perusahaan, penulis menyatakan bahwa elastisitas
pada tingkat munisipalitas lebih rendah daripada 95 persen elastisitas
yang dihitung dari unit buatan. Untuk regresi upah, teks menyebut
koefisien munisipalitas sebesar 5,2 persen, lebih rendah daripada batas
bawah 95 persen sebesar 6,6 persen pada grid 15 km. Nilai tersebut
divisualisasikan pada Gambar 4 dan Gambar 6 (TXT, hlm. 157-159, baris
751-802).

\textbf{Laporan sumber: perbandingan ukuran grid.} Untuk upah, penulis
menemukan perbedaan elastisitas kepadatan antara grid 15 km dan 26 km,
tetapi tidak menemukan perbedaan lanjutan yang signifikan di luar
perbandingan tersebut. Hubungan dengan tahun sekolah disebut berubah
antara grid 26 km dan 39 km. Untuk pembentukan perusahaan, estimasi
titik pada grid 39 km tampak berbeda dari dua grid lain, tetapi
perbedaannya tidak signifikan secara statistik karena variasi
estimasinya lebih besar. Penulis juga menyatakan bahwa elastisitas upah
terhadap pembentukan perusahaan pada grid besar berbeda secara statistik
(TXT, hlm. 158-160, baris 806-837).

\textbf{Laporan sumber: perbandingan batas administratif dan unit
buatan.} Pada Model 1, koefisien mikro-wilayah sebesar 12,7 persen
dibandingkan dengan rata-rata 11,8 persen pada grid 26 km. Berdasarkan
distribusi pada Gambar 4, perbedaan itu tidak signifikan; hipotesis nol
bahwa efek unit administratif mikro-wilayah relatif terhadap unit buatan
berukuran rata-rata sama dengan nol tidak dapat ditolak. Pada Model 2,
penulis menyatakan bahwa koefisien kepadatan yang dibandingkan lebih
besar dari satu dan kesamaan koefisien ditolak pada tingkat kepercayaan
5 persen, yang berarti terdapat perbedaan kecil tetapi signifikan antara
unit administratif dan unit buatan. Angka 0,917 pada Tabel 2 tetap
merupakan pengecualian yang tidak direkonsiliasi dalam TXT (TXT, hlm.
159-160, baris 839-859; Tabel 2, baris 693-695).

\textbf{Laporan sumber: distribusi perluasan dan validitas luar.} Untuk
menguji perbandingan distribusi secara lebih luas, penulis menggabungkan
300 replikasi dengan 50 pengambilan acak dari distribusi normal di
sekitar setiap estimasi. Uji Kolmogorov-Smirnov dua sisi tidak menolak
hipotesis nol kesamaan antara distribusi yang terkait dengan batas
administratif dan distribusi yang mencakup seluruh permutasi. Dengan
distribusi perluasan tersebut, estimasi administratif dan buatan tidak
dapat dibedakan baik pada mikro-wilayah maupun munisipalitas, dan
perbedaan antara dua ukuran grid buatan juga tidak signifikan. Penulis
menyimpulkan bahwa temuan khusus mengenai ukuran wilayah bersifat
indikatif bila dibawa ke negara lain. Nilai statistik KS tidak
dicantumkan dan disebut tersedia atas permintaan (TXT, hlm. 161, baris
896-931 dan catatan kaki 18 baris 941).

\textbf{Laporan sumber: perbandingan dengan metode lain.} Agregasi
berbasis lingkaran dengan radius 14 km dipilih agar luasnya mendekati
rata-rata mikro-wilayah dan grid 26 km. Di Hungaria, munisipalitas yang
lebih kaya cenderung dihitung lebih sedikit oleh lingkaran, sedangkan
munisipalitas yang lebih miskin dihitung lebih sering; penulis
menganggap hal ini dapat membiasakan estimasi. Unit grid yang saling
terpisah tidak mengalami penghitungan berulang tersebut (TXT, hlm. 162,
baris 949-967).

Pada metode dots-to-boxes dengan 150 seed, artikel melaporkan bahwa unit
buatan lebih kecil di barat daya dan lebih besar di tengah-timur karena
heterogenitas ukuran munisipalitas. Grid menghasilkan 93 unit di sebelah
barat Danube, sedangkan dots-to-boxes menghasilkan 57 unit. Distribusi
koefisien dots-to-boxes memiliki median yang agak lebih rendah dan
simpangan baku yang lebih tinggi, yaitu 0,7 dibandingkan 1,4 persen.
Estimasi dari metode lingkaran lebih kecil daripada estimasi titik
mikro-wilayah, tetapi keduanya berada dalam interval kepercayaan
estimasi inti (TXT, hlm. 162-163, baris 974-1020; Appendix, Gambar
14-15, baris 1163-1182).

\subsection{Findings}\label{findings-63}

\textbf{Interpretasi penulis.} Temuan substantif yang ditarik penulis
adalah sebagai berikut:

\begin{itemize}
\tightlist
\item
  Ukuran unit spasial dapat mengubah elastisitas aglomerasi. Pada contoh
  upah, koefisien meningkat dari unit yang lebih kecil ke unit yang
  lebih besar, tetapi bukti perbedaan statistik terutama terlihat antara
  grid 15 km dan 26 km.
\item
  Fragmentasi munisipalitas Hungaria berkaitan dengan elastisitas
  kepadatan-upah yang lebih rendah daripada estimasi pada unit buatan
  yang kurang terfragmentasi dan berukuran serupa. Penulis
  memperingatkan bahwa perbandingan lintas negara pada nomenklatur
  administratif yang sama dapat membandingkan unit dengan ukuran nyata
  yang berbeda.
\item
  Batas administratif mikro-wilayah tidak menunjukkan perbedaan yang
  signifikan untuk elastisitas upah ketika dibandingkan dengan grid 26
  km yang berukuran rata-rata sama. Sebaliknya, perbedaan pada model
  pembentukan perusahaan ditolak pada tingkat 5 persen.
\item
  Pengujian dengan distribusi perluasan melemahkan pembacaan bahwa
  perbedaan ukuran grid atau batas administratif dapat digeneralisasi
  dengan mudah ke negara lain.
\item
  Grid and shake dipandang lebih sesuai daripada dots-to-boxes ketika
  unit dasar sangat heterogen dan ukuran unit berkorelasi secara nonacak
  dengan kondisi sosial-ekonomi. Dibandingkan metode lingkaran,
  keunggulan yang ditekankan adalah setiap unit dasar dihitung satu
  kali.
\end{itemize}

\textbf{Inferensi review.} Hasil paling kuat dalam artikel adalah
demonstrasi bahwa satu hubungan dapat berubah menurut agregasi dan bahwa
variasi tersebut dapat divisualisasikan serta diuji dengan distribusi
koefisien. Namun, hasil antaroutcome tidak seragam: perbedaan batas
administratif tidak signifikan pada model upah tetapi signifikan pada
model pembentukan perusahaan. Karena itu, temuan tidak mendukung
pernyataan universal bahwa batas administratif selalu penting atau
selalu tidak penting. Pernyataan penulis bahwa mungkin ada kebijakan
yang berperan pada model pembentukan perusahaan adalah interpretasi yang
masuk akal dari perbedaan tersebut, tetapi bukan identifikasi efek
kebijakan.

\subsection{Conclusions}\label{conclusions-63}

\textbf{Laporan sumber.} Penulis menyimpulkan bahwa grid and shake
adalah metode yang fleksibel, transparan, sederhana, dan relatif cepat
untuk menguji robustnes penelitian empiris pada data yang diagregasi
secara spasial. Dalam menghasilkan distribusi parameter, metode ini
diposisikan sebagai alternatif bagi dots-to-boxes; dalam memperhitungkan
jarak, ia diposisikan sebagai alternatif bagi spatial lags berbobot
(TXT, hlm. 163-164, baris 1023-1029).

Metode dinilai sesuai bagi peneliti dan praktisi yang memiliki data pada
tingkat administratif rendah seperti munisipalitas atau kode pos.
Penulis menyebut empat kondisi penggunaan: variabel bermakna pada unit
regional; proses yang dikaji tidak dibatasi pada satu unit dan memiliki
spillover; wilayah tidak dipisahkan dari aliran transportasi atau
pengetahuan oleh geografi pertama seperti laut atau gunung; serta unit
dasar heterogen dan ukuran rata-ratanya berkorelasi dengan fitur tak
teramati (TXT, hlm. 163-164, baris 1030-1057).

Penulis menyatakan bahwa unit buatan yang berukuran sama membantu
menghindari bias dari struktur spasial yang endogen, menghitung semua
unit dasar sekali, dan menggabungkan keunggulan sejumlah metode
terdahulu. Keterbatasan yang diakui dalam perbandingan itu adalah
fleksibilitas yang lebih rendah untuk memodelkan hubungan nonlinier di
dalam unit. Dalam aplikasi Hungaria, penulis menyatakan bahwa ukuran
grid memengaruhi pengukuran premi aglomerasi dan bahwa batas regional
tidak berdampak pada elastisitas upah terhadap kepadatan (TXT, hlm. 164,
baris 1058-1076).

\textbf{Inferensi review.} Kesimpulan tersebut sebaiknya dibaca sebagai
klaim metodologis yang bersyarat, bukan validasi umum untuk seluruh
jenis data regional. Uji validitas luar tidak menolak perbedaan
distribusi, tetapi itu tidak membuktikan bahwa distribusi benar-benar
identik. Selain itu, kesimpulan tentang kebijakan tidak dapat dinaikkan
menjadi efek kausal karena desain yang dilaporkan berupa penampang
lintang dan perbandingan agregasi.

\subsection{Contributions}\label{contributions-63}

\textbf{Laporan sumber.} Kontribusi metodologis yang diklaim artikel
adalah prosedur yang dapat: membandingkan estimasi pada beberapa tingkat
agregasi; memperlakukan distribusi ukuran unit yang tidak seragam dan
tidak acak; membandingkan unit administratif dengan unit buatan; dan
mengukur signifikansi statistik perbedaan (TXT, hlm. 145, baris
158-163).

Kontribusi empirisnya adalah penerapan prosedur pada dua hubungan di
Hungaria, yaitu kepadatan dengan upah dan kepadatan dengan pembentukan
perusahaan baru, serta pembandingan tiga ukuran grid dengan dua tingkat
administratif. Artikel juga menempatkan hasil itu berdampingan dengan
metode lingkaran, spatial lags, dan dots-to-boxes (TXT, hlm. 153-163,
baris 538-547 dan 751-1020).

\textbf{Inferensi review.} Sumbangan utama artikel adalah membuat
pilihan agregasi menjadi objek robustnes yang dapat diamati sebagai
distribusi, bukan sekadar keputusan awal peneliti. Sumbangan itu
bersifat prosedural dan diagnostik. TXT tidak menunjukkan bahwa artikel
menyediakan dataset publik, kode replikasi, atau estimator kausal baru;
keberadaan atau ketiadaan bahan tersebut di luar TXT tidak dapat
dipastikan dari sumber yang digunakan.

\subsection{Research Gap}\label{research-gap-63}

\textbf{Laporan sumber.} Artikel mengidentifikasi celah pada praktik
estimasi regional yang menggunakan unit administratif tertentu tanpa
cara statistik yang memadai untuk mengetahui apakah koefisien sensitif
terhadap batas dan ukuran unit. Metode terdahulu dinilai hanya memenuhi
sebagian kebutuhan: spatial lag menangani aspek spasial tetapi jumlah
tetangga tidak ditentukan oleh ukuran unit; lingkaran sederhana dan
fleksibel tetapi dapat mengulang atau mengoversampling wilayah padat;
dots-to-boxes dapat menghasilkan ukuran dan bentuk yang tidak seragam
ketika unit dasar heterogen; dan robustness check konvensional tidak
selalu memberi distribusi pembanding atau uji signifikansi (TXT, hlm.
147-149 dan 162-163, baris 351-364 dan 949-1020).

Penulis terutama menempatkan konteks Hungaria sebagai kasus yang tidak
cocok untuk metode yang mengandaikan unit dasar relatif sama besar,
karena ukuran munisipalitas berbeda secara besar dan mengelompok secara
spasial. Grid and shake diajukan untuk mengisi celah tersebut dengan
menyamakan ukuran rata-rata unit buatan sambil mengacak letak batasnya
(TXT, hlm. 147-148 dan 162-163, baris 285-312 dan 989-1008).

\textbf{Inferensi review.} Celah yang tersisa adalah pengujian lintas
negara dan lintas waktu yang lebih luas, termasuk negara dengan
geografi, sistem administratif, dan distribusi unit dasar yang berbeda.
TXT juga belum memberi uji sensitivitas menyeluruh terhadap ukuran
marker, ambang koreksi batas, jumlah replikasi, pilihan ukuran L, dan
bentuk ketergantungan spasial. Ini adalah celah yang diinferensikan dari
batas cakupan artikel, bukan agenda eksplisit yang dinyatakan penulis.

\subsection{Limitations}\label{limitations-63}

\textbf{Laporan sumber.} Artikel menyebut beberapa keterbatasan data dan
prosedur secara langsung atau tersebar di dalam pembahasan:

\begin{itemize}
\tightlist
\item
  Data lokasi perusahaan hanya menunjukkan kantor pusat. Penulis merujuk
  bukti bahwa 7 persen perusahaan memiliki lebih dari satu lokasi dan
  rata-rata perusahaan tersebut memiliki 1,15 pabrik, lalu menilai
  potensi biasnya relatif kecil (TXT, hlm. 153, catatan kaki 14, baris
  556-572).
\item
  Pemberi kerja pekerja pada Hungarian Labor Force Survey tidak
  diketahui. Karena itu, informasi pekerja tidak dapat ditautkan
  langsung ke perusahaan dan hanya diringkas menjadi rata-rata
  pendidikan pada tingkat munisipalitas (TXT, hlm. 153-154, baris
  559-580).
\item
  Regresi munisipalitas tidak mencakup seluruh 3.125 munisipalitas
  karena ada unit tanpa aktivitas manufaktur yang dapat diidentifikasi.
  Perbandingan dengan grid 15 km juga dapat dipengaruhi oleh wilayah
  tanpa pekerjaan yang masuk ke penyebut kepadatan; penulis menyatakan
  arah biasnya ambigu dan persoalan itu menonjol di Hungaria bagian
  barat. Catatan kaki menyatakan bahwa persoalan tersebut tidak
  memengaruhi grid yang lebih besar, tetapi menulis ukuran ``26 and 36
  km'', sedangkan bagian lain menggunakan grid 39 km. Perbedaan ini
  tidak direkonsiliasi dalam TXT (TXT, hlm. 155, catatan kaki 16, baris
  655-680).
\item
  Koreksi batas hanya menangani unit yang lebih kecil dari 10 persen
  luas kotak dengan menggabungkannya ke tetangga. Penulis menyebutnya
  sebagai pengimbangan sebagian, bukan penghapusan seluruh efek batas
  (TXT, hlm. 152, baris 489-496).
\item
  Penggunaan metode dibatasi oleh empat kondisi yang dicantumkan
  penulis, termasuk adanya spillover lintas unit, tidak adanya
  penghalang geografis yang relevan, dan korelasi antara heterogenitas
  ukuran unit dengan fitur tak teramati (TXT, hlm. 163-164, baris
  1030-1057).
\end{itemize}

\textbf{Inferensi review.} Bukti empiris dibatasi pada satu negara,
periode utama 1999 untuk upah dan informasi 1998-1999 untuk pembentukan
perusahaan, serta perusahaan manufaktur yang memenuhi kriteria tertentu.
Generalisasi ke wilayah lain tidak dapat diasumsikan. Pengacakan juga
hanya menyapu pergeseran grid dalam aturan yang dipilih, bukan seluruh
kemungkinan partisi spasial. Pemakaian distribusi normal untuk
menggambar ketidakpastian koefisien dan penggunaan 300 replikasi
merupakan keputusan pemodelan yang dapat memengaruhi interval dan uji
perbandingan.

Desain penampang lintang yang dilaporkan tidak mengatasi endogenitas
kepadatan, seleksi lokasi perusahaan, atau kemungkinan bahwa faktor tak
teramati memengaruhi upah dan pembentukan perusahaan sekaligus. Dengan
demikian, istilah ``elasticity'' dalam hasil tidak boleh dibaca sebagai
bukti kausal tanpa asumsi tambahan yang tidak disediakan TXT. Detail
yang tidak tersedia mengenai jumlah perusahaan, ukuran sampel survei,
konstruksi perusahaan baru, kode, dan statistik KS juga membatasi
pemeriksaan replikasi.

\subsection{Challenges}\label{challenges-63}

\textbf{Laporan sumber.} Tantangan implementasi yang dihadapi metode
meliputi digitasi peta, penyediaan marker yang cukup rapat untuk
mengenali wilayah munisipalitas kecil, pencocokan marker dengan poligon,
pengelompokan munisipalitas yang memotong lebih dari satu kotak, serta
pengurangan unit kecil di perbatasan (TXT, hlm. 150-152, baris 391-418
dan 489-519).

Penulis juga menghadapi tantangan substantif berupa distribusi ukuran
munisipalitas yang tidak seimbang dan berkelompok. Tantangan analitis
berikutnya adalah membangun distribusi koefisien dari banyak realisasi,
membandingkannya dengan satu koefisien administratif, serta memasukkan
standard error masing-masing realisasi ke dalam distribusi perluasan
(TXT, hlm. 147-153 dan 161, baris 255-312, 499-535, dan 908-918).

Pencocokan data juga memiliki tantangan: data neraca perusahaan berada
pada tingkat perusahaan, data populasi dan luas pada tingkat
munisipalitas, dan data pendidikan berasal dari survei yang tidak
mengidentifikasi pemberi kerja. Hasil akhirnya harus dibangun ulang pada
setiap agregasi, termasuk dengan pembobotan frekuensi untuk rata-rata
tahun sekolah (TXT, hlm. 153-155, baris 552-580 dan 621-625).

\textbf{Catatan audit review.} TXT memuat beberapa ketidakkonsistenan
yang perlu diperlakukan sebagai tantangan dokumentasi, bukan dikoreksi
dengan dugaan:

\begin{itemize}
\tightlist
\item
  Bagian 3.4 menyebut pengulangan seribu kali, sedangkan sebagian besar
  keterangan metode, tabel, dan gambar menyebut 300 replikasi (TXT, hlm.
  145 dan 152-159, baris 141-146, 506-535, 643-651, 699-707, dan
  724-735).
\item
  Tabel 1 dan Tabel 2 mencantumkan 2.872 observasi munisipalitas,
  sedangkan paragraf sesudah Tabel 1 menyebut 2.873 (TXT, hlm. 155-156,
  baris 644-655 dan 699-701).
\item
  Narasi hasil upah menyebut 7,7 persen untuk grid kecil, 11,8 persen
  untuk mikro-wilayah, dan 12,7 persen untuk grid sedang, sementara
  Tabel 1 mencantumkan 0,082, 0,128, dan 0,116 untuk kolom yang
  bersesuaian (TXT, hlm. 154-155, baris 612-619 dan 637-645).
\item
  Narasi Model 2 menyatakan bahwa semua koefisien kepadatan lebih besar
  dari satu, tetapi Tabel 2 mencantumkan 0,917 pada mikro-wilayah (TXT,
  hlm. 155-156, baris 711-721 dan 693-695).
\item
  Catatan kaki 16 menyebut grid 36 km, sedangkan metode dan tabel
  menggunakan grid 39 km. Tabel 2 juga mewarisi catatan ``years of
  schooling'' meskipun baris kovariatnya adalah Wage. TXT tidak memberi
  penjelasan atas perbedaan penulisan ini (TXT, hlm. 155-156, baris
  676-680 dan 703-707).
\end{itemize}

\textbf{Inferensi review.} Ketidakkonsistenan tersebut tidak cukup untuk
menyatakan hasil salah, tetapi cukup untuk mengharuskan pembaca memilih
angka secara eksplisit dan menahan diri dari rekonstruksi yang tidak
didukung. Audit metodologis atau replikasi memerlukan klarifikasi
penulis, terutama mengenai jumlah replikasi dan jumlah observasi dasar.

\subsection{Practical Implications}\label{practical-implications-63}

\textbf{Laporan sumber.} Penulis menyarankan agar peneliti tidak
membandingkan koefisien regional hanya berdasarkan nomenklatur
administratif. Ukuran aktual unit, distribusi ukuran, dan kemungkinan
pengelompokan geografis perlu diperhitungkan. Grid and shake dapat
digunakan untuk membandingkan estimasi pada unit berukuran rata-rata
sama, menguji apakah unit administratif merupakan realisasi spasial yang
khusus, dan membangun interval atau distribusi perbandingan (TXT, hlm.
144-149 dan 163-164, baris 118-121, 340-364, dan 1070-1076).

Untuk evaluasi kebijakan regional, penulis menyarankan perbandingan
wilayah administratif dengan unit pseudo-administratif yang ukurannya
sebanding. Jika perbedaannya terlihat, prosedur dapat membantu menilai
apakah batas administratif berkaitan dengan hubungan yang diestimasi.
Namun, hasil model upah dalam artikel tidak menunjukkan dampak batas
regional, sementara model pembentukan perusahaan menunjukkan perbedaan
kecil yang signifikan (TXT, hlm. 148-149 dan 159-164, baris 315-348,
839-890, dan 1070-1076).

\textbf{Inferensi review.} Implikasi praktis yang paling aman adalah
menjadikan sensitivitas unit spasial sebagai bagian dari pelaporan
utama, bukan catatan tambahan. Prosedur ini dapat memperingatkan pembaca
ketika perbandingan lintas wilayah lebih banyak mencerminkan ukuran unit
daripada perbedaan hubungan substantif. Ia tetap tidak menggantikan
desain identifikasi kebijakan, pengendalian endogenitas, atau penilaian
kualitas data.

\subsection{Applications}\label{applications-63}

\textbf{Laporan sumber: aplikasi yang benar-benar dilakukan.} Artikel
menerapkan metode pada dua model berbasis data Hungaria. Aplikasi
pertama mengestimasi premi aglomerasi upah dari kepadatan dan
pendidikan. Aplikasi kedua mengestimasi pembentukan perusahaan baru dari
kepadatan dan upah. Keduanya dibandingkan pada tingkat munisipalitas,
mikro-wilayah, serta grid 15 km, 26 km, dan 39 km (TXT, hlm. 153-160,
baris 538-547, 583-721, dan 751-859).

Tiga pertanyaan aplikasi yang lebih spesifik adalah pengaruh fragmentasi
munisipalitas, pengaruh ukuran unit spasial, dan peran pembagian
politik-administratif. Artikel juga menerapkan atau mereplikasi
perbandingan dengan agregasi lingkaran dan dots-to-boxes untuk melihat
dampak penghitungan berulang serta heterogenitas ukuran dan bentuk unit
(TXT, hlm. 157-163, baris 751-755 dan 934-1020).

\textbf{Laporan sumber: aplikasi yang disarankan.} Dalam kesimpulan,
penulis menyebut contoh variabel yang dapat bermakna pada unit regional,
seperti jumlah tempat tidur rumah sakit, dokter, dan kelas universitas,
serta proses spillover seperti aliran pengetahuan. Kondisi itu
disampaikan sebagai ruang penggunaan metode, bukan sebagai aplikasi
empiris tambahan dalam artikel (TXT, hlm. 163-164, baris 1030-1057).

\textbf{Inferensi review.} Di luar Hungaria, TXT hanya memberi ilustrasi
korelasi regional Eropa dan tidak melaporkan penerapan grid and shake
pada negara lain. Karena itu, penggunaan artikel untuk konteks lain
harus diperlakukan sebagai transfer metode yang perlu diuji, bukan
sebagai bukti bahwa pola koefisiennya berlaku di tempat lain.

\subsection{Future Research
(Optional)}\label{future-research-optional-63}

\textbf{Laporan sumber.} Artikel tidak memiliki bagian agenda penelitian
masa depan yang eksplisit. Penulis menyatakan bahwa metode dapat
digunakan dalam pengaturan spasial lain dan berguna untuk membandingkan
dataset dari berbagai bagian dunia, tetapi tidak merumuskan daftar studi
lanjutan, desain baru, atau target data tertentu (TXT, hlm. 143-146 dan
163-164, baris 93-98, 164-183, dan 1070-1076).

\textbf{Inferensi review, bukan agenda yang dinyatakan penulis.}
Penelitian lanjutan yang logis dapat menguji metode pada lebih banyak
negara dan periode, membandingkan jenis unit dasar yang berbeda, serta
melakukan analisis sensitivitas formal terhadap ukuran grid, kerapatan
marker, ambang koreksi batas, jumlah replikasi, dan aturan tetangga.
Penggabungan dengan panel, desain kuasi-eksperimental, atau model yang
secara eksplisit menangani ketergantungan spasial juga dapat memisahkan
sensitivitas agregasi dari masalah identifikasi kausal. Usulan ini tidak
dianggap sebagai isi atau rekomendasi eksplisit artikel.

\subsection{Provenance Notes}\label{provenance-notes-63}

\begin{itemize}
\tightlist
\item
  Seluruh review ini hanya menggunakan file TXT B20 yang ditentukan
  pengguna. Tidak ada informasi dari DOI, laman jurnal, artikel lain,
  atau pencarian eksternal yang ditambahkan.
\item
  Seluruh teks yang tersedia dibaca, termasuk metadata ekstraksi pada
  baris 1-28, artikel pada baris 30-1182, appendix, dan daftar referensi
  pada baris 1187-1251.
\item
  Locator utama menggunakan halaman cetak artikel, bagian, tabel,
  gambar, catatan kaki, dan rentang baris TXT. Nomor halaman PDF dan
  nomor halaman artikel dapat berdekatan tetapi tidak diasumsikan
  identik; halaman yang digunakan di review mengikuti nomor halaman
  artikel yang tercetak, seperti 143-170.
\item
  Pernyataan dengan label \textbf{Laporan sumber} merangkum apa yang
  dinyatakan, ditampilkan, atau dilaporkan dalam TXT. Pernyataan dengan
  label \textbf{Interpretasi penulis} merangkum cara penulis memaknai
  hasil atau keunggulan metodenya. Pernyataan dengan label
  \textbf{Inferensi review} adalah analisis pengolahan ini dan bukan
  klaim yang dinyatakan penulis.
\item
  Informasi yang tidak tersedia atau tidak dapat diverifikasi dari TXT
  meliputi nomor isu jurnal, jumlah perusahaan dan individu sebelum
  agregasi, ukuran numerik serta desain sampling HLFS, definisi rinci
  perusahaan baru, nilai statistik KS, kode replikasi, dan akses data
  mentah.
\item
  Ketidakkonsistenan internal mengenai 300 versus seribu replikasi,
  2.872 versus 2.873 observasi, angka narasi versus Tabel 1, klaim
  koefisien kepadatan Model 2 versus angka 0,917, serta penyebutan 36
  versus 39 km dipertahankan sebagai catatan audit. Review tidak memilih
  salah satu angka dengan bantuan sumber di luar TXT.
\item
  Tidak ada tabel yang digunakan dalam review ini.
\end{itemize}

\section{Review B21}\label{review-b21}

\subsection{Identitas Sumber}\label{identitas-sumber-14}

\begin{itemize}
\tightlist
\item
  Kode kajian: B21.
\item
  Judul: ``Do New Economic Geography agglomeration shadows underlie
  current population dynamics across the urban hierarchy?''
\item
  Penulis: Mark D. Partridge, Dan S. Rickman, Kamar Ali, dan M. Rose
  Olfert.
\item
  Jurnal: \emph{Papers in Regional Science}, volume 88, nomor 2, Juni
  2009, halaman 445-467.
\item
  DOI yang tercantum pada metadata PDF dan artikel berbahasa Inggris:
  \texttt{10.1111/j.1435-5957.2008.00211.x}.
\item
  Waktu editorial yang tercantum: diterima 21 September 2007 dan
  disetujui 17 Juli 2008.
\item
  Jenis sumber: artikel jurnal. Status peer-review tidak diverifikasi di
  dalam TXT.
\item
  File yang diperiksa:
  \texttt{source/journal/B21-partridge-rickman-ali-olfert-2009-do-new-economic-geography-agglomeration-shadows-underlie-current-population-dynamics-across-the-urban-hierarchy-10.1111-j.1435-5957.2008.00211.x.txt}.
\item
  Isi file: metadata PDF, teks artikel berbahasa Inggris, tabel utama
  dan lampiran, daftar referensi, serta abstrak berbahasa Spanyol.
\item
  Basis pembacaan: seluruh TXT dibaca, termasuk blok metadata hasil
  ekstraksi PDF. Metadata menyatakan ekstraksi menggunakan
  \texttt{pdftotext\ -layout} pada 31 Agustus 2026 dan PDF memiliki 23
  halaman.
\item
  Locator identitas: TXT baris 1-28 dan 35-63.
\end{itemize}

\subsection{Summarized Abstract}\label{summarized-abstract-64}

\textbf{Laporan sumber:} Artikel menguji apakah kedekatan dengan pusat
urban yang berukuran sama atau lebih tinggi memengaruhi pertumbuhan
penduduk county di Amerika Serikat pada 1990-2006. Pertanyaan ini
ditempatkan dalam perdebatan New Economic Geography (NEG), khususnya
prediksi bahwa pusat urban yang berdekatan dapat menimbulkan
\emph{agglomeration shadow} melalui kompetisi harga spasial. Hasil umum
yang dilaporkan tidak mendukung bayangan pertumbuhan sebagai pola
dominan: pusat urban yang lebih besar umumnya memberi efek pertumbuhan
positif bagi tempat-tempat berpenduduk kurang dari 250.000 yang lebih
dekat. Namun, terdapat bukti bahwa kawasan urban terbesar dapat
membentuk bayangan pertumbuhan atas area metropolitan berukuran menengah
di sekitarnya dan bahwa terdapat kompetisi spasial antararea
metropolitan kecil.

\textbf{Interpretasi penulis:} Abstrak menempatkan NEG sebagai
penjelasan yang masih sebagian berguna, tetapi tidak memadai untuk
menjelaskan seluruh dinamika penduduk dalam sistem urban Amerika Serikat
yang sudah matang. Efek kedekatan positif dan bayangan pertumbuhan hidup
berdampingan menurut tingkat urban, bukan menghasilkan satu pola yang
seragam.

\textbf{Inferensi review:} Ringkasan ini menunjukkan bahwa kesimpulan
artikel bersifat kondisional dan bertingkat. Artikel tidak menyatakan
bahwa seluruh prediksi NEG salah; yang ditolak adalah penerapan
sederhana bayangan aglomerasi sebagai mekanisme umum bagi semua
tingkatan urban.

\textbf{Locator:} TXT baris 50-63, abstrak, artikel halaman 445.

\subsection{Problem Statement}\label{problem-statement-64}

\textbf{Laporan sumber:} Saling keterkaitan antara pusat urban, daerah
pinggiran rural, keputusan lokasi rumah tangga, dan keputusan lokasi
perusahaan membentuk arus penduduk serta evolusi hierarki urban. Dalam
teori tempat sentral dan banyak formulasi NEG, pusat dengan tingkat
lebih tinggi menyediakan barang dan jasa berurutan lebih tinggi.
Kedekatan dengan pusat yang sama besar atau lebih besar dapat sekaligus
memberikan manfaat aglomerasi dan menimbulkan kompetisi harga spasial.
Efek bersih terhadap pertumbuhan penduduk karena itu belum jelas.

Artikel menyatakan bahwa NEG telah banyak digunakan untuk menjelaskan
kemunculan sistem urban Amerika Serikat, tetapi hanya sedikit studi yang
menguji kemampuannya menjelaskan dinamika penduduk kontemporer dalam
sistem urban yang telah mapan. Literatur terkait terutama berfokus pada
pusat urban besar dan area metropolitan. Studi sebelumnya yang menelaah
area nonmetropolitan dan metropolitan kecil untuk 1950-2000 belum
menguji interaksi antara metropolitan menengah dan besar, serta belum
berfokus pada pengalaman pertumbuhan terbaru. Perubahan biaya
transportasi, teknologi, biaya transaksi, dan struktur ekonomi membuat
relevansi pola NEG historis perlu diuji kembali.

\textbf{Interpretasi penulis:} Masalah utama bukan sekadar apakah
aglomerasi ada, tetapi apakah jarak dari berbagai tingkat pusat urban
memperlambat atau justru mendukung pertumbuhan county. NEG/CPT standar
memprediksi hubungan jarak-pertumbuhan yang positif karena perlindungan
dari kompetisi, sedangkan perluasan model yang memperhitungkan
perjalanan komuter, hubungan input-output di luar batas kota, akses
terhadap layanan, amenitas, dan biaya kemacetan dapat menghasilkan
hubungan negatif.

\textbf{Inferensi review:} Problem statement artikel dapat dirumuskan
sebagai masalah identifikasi pola spasial: pada tingkat urban mana efek
akses dan limpahan aglomerasi mengalahkan bayangan kompetisi, dan pada
tingkat mana bayangan tetap terlihat? Ini juga menjelaskan mengapa satu
koefisien jarak untuk seluruh sistem urban tidak cukup.

\textbf{Locator:} TXT baris 72-130 untuk konteks masalah; baris 148-169
untuk kekosongan studi dan pertanyaan penelitian; artikel halaman
445-447.

\subsection{Objectives}\label{objectives-64}

\textbf{Laporan sumber:} Tujuan yang dinyatakan atau langsung tersirat
dalam desain artikel adalah:

\begin{itemize}
\tightlist
\item
  Menguji pengaruh lokasi geografis county dalam hierarki urban terhadap
  pertumbuhan penduduk.
\item
  Menguji apakah kedekatan dengan pusat urban yang sama tingkatnya atau
  lebih tinggi menghasilkan manfaat akses aglomerasi atau bayangan
  pertumbuhan.
\item
  Mencakup interaksi dari tingkat terendah, yaitu county rural dan area
  mikropolitan, sampai area metropolitan terbesar.
\item
  Membandingkan pertumbuhan 1990-2000 sebagai analisis utama dengan
  pertumbuhan 1990-2006 sebagai perluasan dan pemeriksaan sensitivitas
  terhadap periode yang lebih baru.
\item
  Memisahkan jarak ke pusat urban terdekat dari jarak inkremental untuk
  mencapai tingkat metropolitan yang lebih tinggi, sehingga efek jarak
  dapat berbeda pada setiap segmen hierarki.
\item
  Menguji peran jarak ke pusat urban lain dalam tingkat yang sama
  sebagai kemungkinan kompetisi spasial atau penguatan aglomerasi
  regional.
\item
  Menilai apakah hubungan komuter membantu menjelaskan efek jarak dengan
  memasukkan pertumbuhan pekerjaan berbasis \emph{industry mix} dan
  interaksinya dengan jarak.
\item
  Menilai ketahanan hasil terhadap penambahan variabel ekonomi,
  demografis, amenitas, wilayah sekitar, efek tetap negara bagian,
  definisi batas metropolitan, serta ambang ukuran metropolitan.
\end{itemize}

\textbf{Interpretasi review:} Tujuan empiris artikel lebih luas daripada
menguji satu hipotesis bayangan. Artikel membangun perbandingan
mekanisme menurut tingkat urban dan menguji apakah hasil yang tampak
sebagai limpahan sebenarnya berjalan melalui kesempatan komuter.

TXT tidak merumuskan hipotesis formal bernomor, evaluasi program
kebijakan, atau tujuan untuk menghasilkan prediksi kebijakan spesifik.
Prediksi yang digunakan berbentuk arah hubungan yang bersaing antara
kedekatan, limpahan aglomerasi, dan kompetisi.

\textbf{Locator:} TXT baris 153-169, 333-350, dan 789-800; artikel
halaman 447-450 dan 458.

\subsection{Literature Survey}\label{literature-survey-64}

\textbf{Laporan sumber:} Artikel memulai dari teori tempat sentral (CPT)
Christaller, yang menggambarkan hierarki tempat berdasarkan pasar
potensial dan ketersediaan barang serta jasa dengan tingkat pesanan
berbeda. NEG, terutama Krugman dan pengembangan Fujita, menggabungkan
persaingan tidak sempurna, skala hasil meningkat, faktor yang bergerak,
biaya transportasi, dan keterkaitan pemasok-pelanggan untuk menjelaskan
aglomerasi secara endogen. Dalam kerangka ini, pusat besar melayani
pasar yang lebih luas, sedangkan pusat kecil melayani pasar lokal dan
dapat berada dalam bayangan pertumbuhan pusat yang lebih tinggi akibat
kompetisi.

Teori neoklasik yang dibahas menjelaskan konsentrasi ekonomi melalui
perbedaan geografi, endowment, dan teknologi eksogen. Artikel kemudian
memakai pendekatan \emph{compensating differential} sebagai dasar untuk
mempertimbangkan amenitas konsumsi dan produksi, utilitas rumah tangga,
keuntungan perusahaan, biaya tempat tinggal, kemacetan, kejahatan,
polusi, harga tanah, dan akses terhadap layanan urban. Model Lucas dan
Rossi-Hansberg dipakai untuk menunjukkan bahwa perusahaan dapat
menimbang manfaat konsentrasi input-output terhadap penghematan biaya
komuter jika kegiatan tersebar ke area hunian.

Artikel juga menempatkan perubahan teknologi transportasi dan
komunikasi, pergeseran perdagangan, restrukturisasi industri dari barang
ke jasa, serta penurunan biaya transportasi sebagai alasan mengapa pola
aglomerasi modern dapat berbeda dari pola abad ke-19 dan awal abad
ke-20. Amenitas urban dengan keragaman konsumsi dan akses terhadap
pekerjaan dipaparkan sebagai mekanisme yang dapat membuat daerah dekat
pusat besar lebih menarik, sekalipun daerah tersebut tidak bersaing
langsung untuk layanan tingkat tertinggi.

Untuk bukti empiris, artikel menyebut bahwa banyak penelitian sebelumnya
berfokus pada kota besar dan area metropolitan. Stelder menjadi
pengecualian yang disebut untuk simulasi hierarki kota Eropa, sedangkan
Partridge et al.~(2008b) disebut telah mempelajari county
nonmetropolitan dan area metropolitan kecil pada 1950-2000 tetapi tidak
interaksi metropolitan menengah-besar atau pertumbuhan paling baru.
Artikel juga menyebut bukti tentang stabilitas pusat teratas dalam
hierarki dan lebih banyak perubahan pada tingkat bawah.

\textbf{Interpretasi penulis:} Tinjauan tersebut mengarah pada sintesis
bahwa NEG/CPT memberi prediksi penting mengenai bayangan kompetisi,
tetapi perlu diperluas dengan limpahan komuter, hubungan input-output
yang melampaui batas kota, amenitas, dan gaya sentrifugal seperti
suburbanisasi. Karena setiap tingkat menyediakan layanan yang berbeda,
jarak inkremental ke tiap tingkat dianggap lebih informatif daripada
jarak total tunggal.

\textbf{Inferensi review:} Literatur digunakan secara teoritis dan
selektif untuk membangun model empiris, bukan sebagai tinjauan
sistematis. TXT tidak menyebutkan protokol pencarian, kriteria inklusi,
jumlah korpus literatur, atau prosedur penilaian kualitas studi. Karena
itu, bagian ini tidak dapat diperlakukan sebagai \emph{systematic
literature review}.

\textbf{Locator:} TXT baris 91-141 untuk CPT, NEG, dan studi terdahulu;
baris 172-270 untuk teori alternatif; baris 273-329 untuk sintesis
hierarki dan jarak; artikel halaman 446-450.

\subsection{Method Used}\label{method-used-64}

\textbf{Laporan sumber:} Artikel menggunakan regresi pertumbuhan
penduduk county sebagai fungsi dari kondisi awal periode 1990. Bentuk
paling lengkapnya ditulis sebagai
\texttt{\%DeltaPOP\_is(t-0)\ =\ alpha\ +\ delta\ POPDEN\_is0\ +\ j\ GEOG\_is0\ +\ theta\ DEMOG\_is0\ +\ gamma\ ECON\_is0\ +\ g\ AMENITY\_is0\ +\ sigma\_s\ +\ epsilon\_is(t-0)}.
Variabel dependen adalah perubahan persentase penduduk, sedangkan vektor
penjelas mencakup kepadatan penduduk awal, geografi, demografi, ekonomi,
amenitas, dan efek tetap negara bagian.

Analisis utama memakai pertumbuhan 1990-2000. Analisis 1990-2006
digunakan untuk melihat apakah pola tetap berlaku dalam periode yang
lebih panjang dan lebih baru. Kondisi awal dipakai untuk membantu
mengurangi masalah endogenitas statistik, sedangkan jarak dianggap telah
ditentukan sebelumnya jika hierarki urban Amerika Serikat diasumsikan
sudah mapan pada 1990.

Jarak dioperasionalkan secara bertingkat. Artikel mengukur jarak ke
pusat urban terdekat atau pusat urban tempat county berada, jarak
inkremental ke metropolitan mana pun, metropolitan berpenduduk
sekurang-kurangnya 250.000, sekurang-kurangnya 500.000, dan lebih dari
1,5 juta. Jarak inkremental adalah selisih jarak ke tingkat berikutnya
dengan jarak ke tingkat sebelumnya. Nilainya nol ketika county sudah
berada di tingkat yang relevan atau pusat terdekatnya sudah termasuk
tingkat yang sama atau lebih tinggi. Untuk beberapa subsampel, jarak ke
pusat urban lain dalam tingkat yang sama juga dimasukkan.

Spesifikasi dibangun bertahap: model awal hanya memasukkan variabel yang
jelas eksogen atau telah ditentukan sebelumnya, kemudian menambahkan
populasi pusat urban, variabel ekonomi dan demografi, variabel wilayah
sekitar, serta efek tetap negara bagian. GMM digunakan untuk memperoleh
statistik-t yang tahan terhadap spillover lintas penampang dan korelasi
spasial residual. Korelasi residual diasumsikan menurun terhadap jarak
dan menjadi nol setelah 200 km, dengan pembobotan yang menurun secara
kuadratik.

Untuk menguji mekanisme komuter, artikel memasukkan pertumbuhan
pekerjaan berbasis \emph{industry mix} pada pusat urban terdekat dalam
kategori ukuran yang sesuai dan interaksi antara pertumbuhan tersebut
dengan jarak. Pertumbuhan pekerjaan pusat urban ditetapkan nol bila
jaraknya lebih dari 160 km karena penulis mengasumsikan efek komuter
kemungkinan melemah setelah sekitar 100 mil. Karena pertumbuhan penduduk
dan pekerjaan dapat ditentukan bersama, pertumbuhan \emph{industry mix}
digunakan sebagai proksi eksogen. Uji tambahan memasukkan persentase
tenaga kerja county yang bekerja di county tempat tinggal pada 1990
sebagai ukuran kebalikan dari komuter.

\textbf{Interpretasi penulis:} Rangkaian spesifikasi dan GMM dimaksudkan
untuk memeriksa apakah hasil jarak merupakan artefak dari variabel
terlewat, multikolinearitas, endogenitas, proses autoregresif yang
terlewat, atau spillover spasial. Artikel menafsirkan koefisien jarak
yang tersisa setelah kontrol komuter sebagai indikasi mekanisme lain,
seperti akses layanan, amenitas, dan externalities backward-forward.

\textbf{Inferensi review:} Ini adalah desain observasional reduced-form
berbasis perbandingan antarcounty dan horizon pertumbuhan, bukan
eksperimen atau model struktural rumah tangga/perusahaan. TXT tidak
memberikan strategi identifikasi kausal terpisah untuk setiap mekanisme;
interpretasi mekanisme harus dibaca sebagai penjelasan yang konsisten
dengan pola koefisien, bukan pengukuran langsung seluruh jalur perilaku.

\textbf{Locator:} TXT baris 333-407 untuk persamaan, asumsi, dan GMM;
baris 409-508 untuk pengukuran jarak dan populasi; baris 509-548 untuk
kontrol; baris 826-966 untuk uji komuter; artikel halaman 450-461.

\subsection{Dataset}\label{dataset-64}

\textbf{Laporan sumber:} Dataset berbentuk data agregat tingkat county
untuk wilayah 48 negara bagian daratan Amerika Serikat dan District of
Columbia. Appendix Table 1 mencantumkan sumber dan konstruksi berikut:

\begin{itemize}
\tightlist
\item
  Perubahan penduduk, populasi awal, kepadatan penduduk, demografi, dan
  sebagian besar variabel ekonomi berasal dari Census 1990 dan Census
  2000.
\item
  Pertumbuhan \emph{industry mix} menggunakan data ketenagakerjaan BEA
  1990 dan 2000 serta perhitungan penulis. Ukurannya mengalikan
  pertumbuhan pekerjaan nasional menurut industri dengan pangsa
  pekerjaan awal county.
\item
  Jarak ke pusat urban, jarak inkremental ke tingkat urban, populasi
  inkremental, dan beberapa variabel pusat urban merupakan estimasi
  penulis berdasarkan centroid berbobot penduduk serta definisi area
  urban.
\item
  Variabel amenitas cuaca dan alam berasal dari Economic Research
  Service, USDA. Vektornya mencakup jam sinar matahari Januari, suhu
  Januari, kelembapan relatif Juli, skor topografi, peringkat amenitas
  alam, dan persentase wilayah berair.
\item
  Variabel wilayah sekitar dibangun sebagai rata-rata berbobot county
  lain dalam 177 wilayah ekonomi fungsional BEA, dengan county yang
  sedang diamati dikeluarkan dari rata-rata.
\item
  Pendapatan median rumah tangga 1989, tingkat pengangguran 1990, pangsa
  pekerjaan pertanian dan barang, distribusi umur, pendidikan,
  ras/etnis, dan imigrasi 1985-1990 digunakan sebagai kontrol awal.
\end{itemize}

Appendix Table 1 melaporkan 3.029 county. Beberapa statistik deskriptif
untuk variabel amenitas, pangsa umur, pendidikan, dan ras/etnis ditandai
\texttt{n.a.}; TXT tidak memberi nilai rata-rata dan simpangan bakunya
untuk variabel-variabel tersebut, meskipun variabelnya dipakai dalam
model.

Untuk pertumbuhan 1990-2006, TXT melaporkan hasil pada Appendix Table 3,
tetapi tidak menguraikan secara terpisah pada tabel tersebut sumber data
penduduk 2006 atau prosedur pembentukan seluruh variabel periode
panjang. Karena itu, sumber Census/BEA yang tercantum untuk konstruksi
utama tidak boleh diperluas tanpa kualifikasi menjadi dokumentasi
lengkap dataset 1990-2006.

\textbf{Inferensi review:} Dataset memungkinkan analisis hierarki urban
yang rinci dan pengendalian karakteristik awal, tetapi unit analisisnya
tetap county agregat. TXT tidak menyebutkan adanya data mikro rumah
tangga, data perusahaan individual, panel tahunan lengkap, kode
replikasi, atau tautan unduhan raw data.

\textbf{Locator:} TXT baris 1014-1124 untuk definisi dan sumber
variabel; baris 1128-1150 untuk statistik subsampel; baris 1157-1211
untuk hasil 1990-2006; artikel halaman 462-464.

\subsection{Population / Sample}\label{population-sample-64}

\textbf{Laporan sumber:} Unit observasi adalah county di 48 negara
bagian daratan dan District of Columbia. Empat kelompok utama dibentuk
menurut posisi dalam hierarki urban:

\begin{itemize}
\tightlist
\item
  Noncore rural: county yang bukan bagian dari micropolitan area maupun
  metropolitan area, sebanyak 1.300 county.
\item
  Micropolitan area (MICRO): pusat urban kecil berpenduduk 10.000-50.000
  beserta county yang memiliki hubungan komuter erat dengannya, sebanyak
  672 county.
\item
  Small metropolitan area: county dalam metropolitan area dengan
  penduduk tidak lebih dari 250.000, sebanyak 416 county.
\item
  Large metropolitan area: county dalam metropolitan area dengan
  penduduk lebih dari 250.000, sebanyak 641 county.
\end{itemize}

Jumlah empat kelompok tersebut adalah 3.029 county. Ambang 250.000
membagi sampel metropolitan menjadi dua kelompok yang kira-kira
seimbang. Untuk beberapa analisis, metropolitan besar juga dibaca
melalui kategori 250.000-500.000, 500.000-1,5 juta, dan lebih dari 1,5
juta.

Empat puluh tiga county, sebagian besar county rural kecil, dikeluarkan
karena tidak memiliki data ekonomi. Definisi MA dan MICRO umumnya
menggunakan definisi tahun 2003. Analisis sensitivitas memakai definisi
tahun 1999; untuk batas metropolitan tertentu artikel juga menguji batas
tahun 1993 dan 1999. MICRO baru didefinisikan pada 2003, sehingga
perbandingan historis untuk kategori ini memiliki batas definisi yang
dinyatakan oleh penulis.

\textbf{Interpretasi penulis:} Pembagian tersebut dirancang untuk
menangkap mekanisme berbeda: akses rural ke pusat urban, interaksi dalam
dan antararea metropolitan kecil, serta hubungan antara metropolitan
menengah dan pusat terbesar. Penggunaan batas yang lebih inklusif
dimaksudkan untuk menangkap perubahan pola komuter, bukan hanya
interaksi di dalam batas urban lama.

\textbf{Inferensi review:} Sampel adalah kumpulan observasi county yang
tersedia setelah penghapusan 43 county, bukan sampel responden yang
diambil melalui survei. Generalisasi langsung berlaku pada cakupan
geografis dan definisi hierarki yang digunakan, bukan otomatis pada
Alaska, Hawaii, wilayah teritori, atau unit yang tidak termasuk dalam
TXT.

\textbf{Locator:} TXT baris 351-378 untuk definisi kelompok dan
penghapusan county; baris 563-567 dan 646 untuk jumlah subsampel; baris
775-780 untuk sensitivitas batas; Appendix Table 1 baris 1116-1124.

\subsection{Dependent Variables}\label{dependent-variables-64}

\textbf{Laporan sumber:} Variabel dependen utama adalah persentase
perubahan penduduk county pada 1990-2000, ditulis sebagai
\texttt{\%DeltaPOP\_is(t-0)}. Appendix Table 3 menggunakan persentase
perubahan penduduk 1990-2006. Analisis dilakukan terhadap pertumbuhan
dalam masing-masing kelompok rural, MICRO, small MA, dan large MA.

Pertumbuhan \emph{industry mix}, persentase tenaga kerja yang bekerja di
county tempat tinggal, populasi pusat urban, jarak, kepadatan, amenitas,
dan karakteristik ekonomi-demografi adalah variabel penjelas atau
kontrol, bukan variabel dependen dalam artikel ini.

TXT tidak memberikan formula lebih rinci untuk menghitung persentase
perubahan penduduk selain definisi tersebut, dan tidak memisahkan
perubahan penduduk menjadi migrasi, kelahiran, atau kematian. Dengan
demikian, interpretasi review tidak memperlakukan variabel dependen
sebagai ukuran migrasi murni.

\textbf{Locator:} TXT baris 339-350 untuk rancangan variabel dependen;
baris 1016-1027 untuk definisi perubahan penduduk; baris 1157-1164 untuk
versi 1990-2006.

\subsection{Results}\label{results-64}

\textbf{Laporan sumber, hasil dasar 1990-2000:} Tabel 1 menggunakan
statistik-t robust terhadap korelasi spasial. Dua tanda bintang pada
tabel berarti signifikan pada tingkat paling tinggi 5 persen dan satu
tanda bintang berarti signifikan pada tingkat 10 persen. Hasil utama
menurut kelompok adalah sebagai berikut.

\begin{itemize}
\tightlist
\item
  Untuk county rural, koefisien jarak ke pusat urban terdekat pada model
  lengkap adalah -0,102 (t = -7,62), signifikan pada 5 persen. Jarak
  inkremental ke MA bernilai -0,038 (t = -4,79), ke MA berpenduduk lebih
  dari 250.000 bernilai -0,030 (t = -5,17), ke MA lebih dari 500.000
  bernilai -0,018 (t = -2,42), dan ke MA lebih dari 1,5 juta bernilai
  -0,011 (t = -2,50). Semua tanda negatif yang signifikan menunjukkan
  pertumbuhan lebih rendah ketika county lebih jauh dari tingkat urban
  yang relevan.
\item
  Untuk county rural pada jarak rata-rata, penulis menghitung bahwa
  jarak ke pusat urban terdekat berkaitan dengan sekitar 6 persen lebih
  sedikit pertumbuhan dibanding county yang berdekatan dengan inti
  pusat. Jika seluruh jarak inkremental digabungkan pada nilai
  rata-ratanya, penalti pertumbuhan yang dihitung adalah 11,9 persen.
  Satu simpangan baku tambahan pada jarak ke pusat terdekat dikaitkan
  dengan penurunan 3,1 persen, sedangkan satu simpangan baku pada
  seluruh jarak inkremental dikaitkan dengan pertumbuhan 10,5 persen
  lebih rendah.
\item
  Untuk MICRO, jarak ke pusat urban aktual atau terdekat pada model
  lengkap bernilai positif 0,137 (t = 2,14), signifikan pada 5 persen.
  Jarak inkremental ke MA bernilai -0,026 (t = -2,02) dan ke MA lebih
  dari 250.000 bernilai -0,029 (t = -3,24), keduanya signifikan pada 5
  persen. Jarak inkremental ke MA lebih dari 500.000 dan lebih dari 1,5
  juta tidak signifikan dalam model lengkap. Jarak ke tingkat yang sama
  juga tidak signifikan.
\item
  Untuk small MA, jarak inkremental ke MA lebih dari 250.000 bernilai
  -0,057 (t = -5,53), ke MA lebih dari 500.000 bernilai -0,027 (t =
  -2,23), dan ke MA lebih dari 1,5 juta bernilai -0,026 (t = -2,66),
  semuanya signifikan pada 5 persen dalam model lengkap. Jarak ke pusat
  urban aktual tidak signifikan dalam model lengkap, tetapi koefisien
  jarak ke pusat tingkat yang sama bernilai positif 0,044 (t = 3,14),
  signifikan pada 5 persen.
\item
  Untuk small MA, penulis menghitung bahwa pada jarak rata-rata
  antar-small MA, pertumbuhan sekitar 2 persen lebih lambat; peningkatan
  satu simpangan baku jarak meningkatkan pertumbuhan sekitar 2 persen.
  Ini adalah hasil yang paling jelas konsisten dengan bayangan
  pertumbuhan karena kompetisi antar pusat dalam tingkat yang sama.
\item
  Untuk large MA, jarak ke pusat urban aktual bernilai positif 0,188 (t
  = 4,26), signifikan pada 5 persen, sedangkan jarak inkremental ke MA
  lebih dari 500.000, ke MA lebih dari 1,5 juta, dan jarak ke tingkat
  yang sama tidak signifikan. Pola positif jarak internal ditafsirkan
  dalam artikel sebagai suburbanisasi atau sprawl di dalam area
  metropolitan.
\item
  Pada large MA, kepadatan penduduk awal bernilai positif 0,00042 dan
  signifikan pada 10 persen, sementara populasi county sekitar bernilai
  positif 4,0 x 10\^{}-7 dan signifikan pada 10 persen. Populasi
  inkremental MA lebih dari 1,5 juta bernilai negatif -5,0 x 10\^{}-7
  dan signifikan pada 5 persen. Penulis menafsirkan hasil terakhir
  sebagai bukti pengaruh buruk pusat terbesar terhadap MA berukuran
  250.000-1,5 juta, meskipun jarak ke pusat tersebut tidak signifikan.
\end{itemize}

Uji bersama pada Tabel 1 mendukung signifikansi jarak inkremental untuk
rural, MICRO, dan small MA, tetapi tidak untuk large MA. Nilai R2 model
lengkap adalah 0,60 untuk rural, 0,63 untuk MICRO, 0,64 untuk small MA,
dan 0,63 untuk large MA. Artikel menyatakan bahwa model-model dengan
spesifikasi bertahap memberikan hasil jarak yang serupa untuk tiga
kelompok pertama.

\textbf{Laporan sumber, pencarian tambahan atas growth shadow:}
Interaksi antara variabel jarak dan populasi pusat urban secara
bersama-sama tidak signifikan dalam model MICRO dan small MA, tetapi
signifikan pada tingkat 5 persen untuk rural dan large MA. Pada rural,
koefisien interaksi jarak-populasi umumnya negatif, yaitu county yang
lebih jauh dari pusat dengan ukuran tertentu tumbuh lebih lambat,
berlawanan dengan bayangan aglomerasi. Pada MA berukuran 250.000-1,5
juta, interaksi positif dengan pusat berpenduduk lebih dari 1,5 juta
menunjukkan bahwa jarak yang lebih besar mengurangi dampak penduduk yang
merugikan, konsisten dengan bayangan pertumbuhan.

\textbf{Laporan sumber, periode 1990-2006:} Appendix Table 3
mempertahankan pola umum. Untuk rural, koefisien jarak ke pusat terdekat
adalah -0,150 dan seluruh jarak inkremental ke MA, MA lebih dari
250.000, lebih dari 500.000, dan lebih dari 1,5 juta tetap negatif serta
signifikan pada 5 persen. Untuk MICRO, koefisien jarak ke pusat terdekat
adalah 0,200 dan signifikan pada 5 persen, sedangkan jarak inkremental
ke MA adalah -0,039 dan ke MA lebih dari 250.000 adalah -0,037, keduanya
signifikan pada 5 persen. Untuk small MA, jarak inkremental ke MA lebih
dari 250.000 adalah -0,095, ke lebih dari 500.000 adalah -0,050, dan ke
lebih dari 1,5 juta adalah -0,060, semuanya signifikan pada 5 persen;
jarak ke tingkat yang sama adalah 0,065 dan signifikan pada 5 persen.
Untuk large MA, jarak ke pusat aktual adalah 0,303 dan signifikan pada 5
persen, sedangkan jarak ke pusat tingkat lebih tinggi umumnya tidak
signifikan. Penulis menyatakan bahwa penalti jarak untuk small MA
menjadi lebih negatif daripada proporsi tambahan panjang periode,
sehingga tampak menguat setelah 2000.

\textbf{Laporan sumber, uji komuter:} Dalam model rural pada Tabel 2,
ukuran pertumbuhan pekerjaan pusat urban berbasis \emph{industry mix}
positif dan secara bersama-sama signifikan pada 5 persen. Interaksi
pertumbuhan pekerjaan dengan jarak negatif dan secara bersama-sama
signifikan pada 5 persen, sehingga manfaat pertumbuhan pekerjaan pusat
urban melemah dengan jarak. Koefisien jarak inkremental tetap secara
bersama-sama signifikan, walaupun besarnya hanya sedikit berkurang.

Untuk sampel urban, ukuran pertumbuhan pekerjaan dan interaksi jaraknya
hampir seluruhnya tidak signifikan menurut uraian penulis. Tabel 2 tetap
menunjukkan bahwa jarak inkremental ke pusat tingkat tinggi umumnya
negatif untuk MICRO dan small MA, sedangkan uji bersama jarak
inkremental tidak mendukung interaksi spasial yang kuat untuk large MA.
Penambahan persentase pekerja yang bekerja di county tempat tinggal
sedikit mengecilkan beberapa koefisien jarak, tetapi efek jarak tetap
ada. Setiap kenaikan satu poin persentase dalam proporsi pekerja yang
bekerja di county sendiri dikaitkan dengan 0,1-0,2 poin persentase lebih
rendah pada pertumbuhan penduduk berikutnya dan signifikan pada 5
persen.

\textbf{Interpretasi penulis:} Pertumbuhan pekerjaan pusat urban
menciptakan kesempatan komuter yang membantu pertumbuhan rural. Karena
jarak inkremental tetap berpengaruh setelah ukuran komuter dimasukkan,
penulis menyimpulkan bahwa hubungan backward-forward, akses layanan,
atau limpahan lain juga penting. Ketidakberartian umum ukuran komuter
untuk interaksi antararea urban digunakan untuk menilai bahwa hubungan
antar pusat urban bukan terutama hubungan komuter.

\textbf{Inferensi review:} Hasil yang paling konsisten adalah hubungan
jarak yang nonlinear dan bergantung tingkat. Akses ke pusat berikutnya
penting bagi tingkat bawah; kompetisi paling terlihat dalam relasi yang
lebih sempit, yaitu small MA dengan small MA lain atau MA menengah
dengan pusat terbesar. Besarnya koefisien tidak dengan sendirinya
memisahkan komuter, amenitas, input-output, dan seleksi lokasi.

\textbf{Locator:} Tabel 1 TXT baris 602-663 dan uraian baris 667-780,
artikel halaman 454-457; uji interaksi baris 787-800, halaman 458;
ringkasan periode baris 803-823; Tabel 2 baris 900-954 dan uraian
komuter baris 826-896 serta 958-966, artikel halaman 458-461; Appendix
Table 3 baris 1157-1211.

\subsection{Findings}\label{findings-64}

\textbf{Interpretasi penulis:} Penulis merangkum bahwa pertumbuhan
penduduk county rural, MICRO, dan small MA umumnya berhubungan positif
dengan kedekatan terhadap pusat urban tingkat lebih tinggi. Pola ini
tidak sesuai dengan bayangan pertumbuhan yang diprediksi oleh CPT dan
NEG standar. Koefisien kepadatan sendiri dan populasi county tetangga
sering memiliki tanda yang tidak sesuai atau tidak signifikan, sehingga
potensi pasar regional bukan penjelasan yang kuat untuk seluruh
perubahan penduduk. Penulis mengusulkan bahwa ekspor dari pusat kecil
dan hinterland mungkin lebih diarahkan ke pasar nasional dan
internasional daripada hanya pasar terdekat.

Untuk large MA, penulis menemukan sedikit bukti bahwa kedekatan dengan
pusat tingkat lebih tinggi memengaruhi pertumbuhan. Ada bukti bayangan
pertumbuhan antar-small MA dalam tingkat yang sama dan bukti tambahan
bahwa MA berukuran 250.000-1,5 juta mengalami dampak persaingan dari MA
terbesar. Sprawl tetap terlihat pada fringe large MA. Periode 1990-2006
mengonfirmasi efek negatif jarak dari pusat tingkat lebih tinggi,
terutama karena efek untuk small MA tampak semakin kuat setelah 2000.

\textbf{Laporan sumber mengenai mekanisme:} Pertumbuhan pekerjaan pusat
urban menaikkan pertumbuhan rural sebagian melalui kesempatan komuter,
tetapi tidak menjelaskan sepenuhnya efek jarak. Untuk interaksi
antarpusat urban, commuting ties tidak tampak sebagai mekanisme utama
dalam hasil yang dilaporkan. Penulis menganggap hubungan
backward-forward yang menjangkau jauh serta akses terhadap amenitas dan
layanan urban lebih sesuai dengan pola yang tersisa.

\textbf{Inferensi review:} Temuan artikel lebih tepat dibaca sebagai
penolakan terhadap versi homogen dari teori bayangan, bukan sebagai
penggantian NEG dengan satu teori tunggal. Kekuatan dan arah hubungan
bergantung pada posisi county, ukuran pusat, dan apakah jarak berada di
dalam pusat urban, antar tingkat, atau antar pusat dalam tingkat yang
sama.

\textbf{Locator:} TXT baris 803-823 untuk ringkasan temuan; baris
866-885 dan 958-966 untuk pembacaan mekanisme komuter; artikel halaman
458-461.

\subsection{Conclusions}\label{conclusions-64}

\textbf{Laporan sumber:} Kesimpulan artikel adalah bahwa county rural
dan pusat urban yang lebih kecil memiliki interaksi positif yang
signifikan dengan pusat urban terdekat yang tingkatnya lebih tinggi.
Jarak yang lebih jauh dari setiap tingkat pusat urban yang lebih tinggi
secara berturut-turut berkaitan dengan pertumbuhan yang lebih rendah.
Pertumbuhan pekerjaan urban mendukung pertumbuhan penduduk rural
sebagian melalui komuter, dan pola tersebut menunjukkan kekuatan yang
mendukung sprawl urban.

Bukti yang konsisten dengan bayangan pertumbuhan NEG relatif sedikit.
Pengecualian utamanya adalah kompetisi spasial antar-small MA dan
bayangan yang dipancarkan MA paling besar terhadap MA terdekat berukuran
250.000-1,5 juta. Pada large MA, interaksi dengan pusat lebih tinggi
umumnya lemah; pola yang paling jelas adalah dekonsentrasi dan sprawl di
dalam area metropolitan.

\textbf{Interpretasi penulis:} Prediksi tertentu dari NEG dan CPT
dianggap kurang memadai untuk menggambarkan evolusi berkelanjutan sistem
urban Amerika Serikat. Penulis menyatakan bahwa tempat berorde lebih
rendah lebih banyak memperoleh manfaat dari akses yang dekat, hubungan
backward-forward, dan amenitas urban daripada mengalami bayangan
kompetisi. Elemen NEG tetap berguna, tetapi perlu dilengkapi perspektif
dan model lain.

\textbf{Inferensi review:} Kesimpulan ini memiliki cakupan empiris yang
jelas: sistem urban Amerika Serikat yang dianalisis pada 1990-2006,
dengan unit county dan klasifikasi yang digunakan artikel. Kesimpulan
tersebut tidak membuktikan bahwa akses selalu menyebabkan pertumbuhan
atau bahwa semua sistem urban matang memiliki pola yang sama.

\textbf{Locator:} TXT baris 970-996, bagian ``5 Conclusion'', artikel
halaman 461.

\subsection{Contributions}\label{contributions-64}

\textbf{Laporan sumber:} Penulis menyatakan bahwa studi ini mengisi
kekosongan dengan menggunakan pertumbuhan penduduk county Amerika
Serikat 1990-2000 dan 1990-2006, sehingga interaksi dapat
dipertimbangkan dari pusat urban tingkat terendah sampai tertinggi.
Mereka juga secara eksplisit memperluas titik berangkat konseptual
dengan memasukkan kemungkinan komuter, externalities input-output di
luar batas kota, serta akses terhadap amenitas dan layanan urban.

\textbf{Inferensi review:} Kontribusi yang dapat ditelusuri dari TXT
adalah sebagai berikut:

\begin{itemize}
\tightlist
\item
  Memindahkan pengujian NEG dari fokus dominan pada kemunculan sistem
  urban atau pusat besar ke dinamika penduduk terkini dalam sistem urban
  yang telah mapan.
\item
  Menguji seluruh rentang hierarki urban secara berlapis, termasuk
  county rural, MICRO, small MA, large MA, dan relasi dengan pusat lebih
  dari 1,5 juta penduduk.
\item
  Mengoperasionalkan jarak inkremental antar tingkat, sehingga biaya
  atau manfaat akses ke layanan tingkat lebih tinggi tidak dipaksa
  mengikuti satu hubungan jarak linear.
\item
  Memisahkan kompetisi antar pusat dalam tingkat yang sama dari
  kedekatan terhadap pusat dengan tingkat lebih tinggi.
\item
  Menguji ketahanan hasil dengan kontrol awal, amenitas,
  ekonomi-demografi, wilayah sekitar, efek tetap negara bagian, korelasi
  spasial, definisi batas, dan dua horizon waktu.
\item
  Menghubungkan hasil pertumbuhan penduduk dengan mekanisme komuter
  melalui proksi pertumbuhan pekerjaan dan interaksi jarak.
\item
  Memberi kualifikasi empiris terhadap NEG: bayangan pertumbuhan muncul
  secara selektif, sementara akses dan limpahan positif dominan pada
  banyak county tingkat bawah.
\end{itemize}

TXT tidak menyatakan bahwa artikel memperkenalkan model formal NEG baru,
dataset mikro baru, atau instrumen kebijakan yang telah diuji.
Kontribusi yang paling aman untuk disebut adalah perluasan empiris dan
penilaian ulang mekanisme spasial.

\textbf{Locator:} TXT baris 110-122 dan 148-169 untuk kekosongan serta
tujuan perluasan; baris 333-348 dan 789-800 untuk desain pengujian;
artikel halaman 446-450 dan 458.

\subsection{Research Gap}\label{research-gap-64}

\textbf{Kesenjangan yang dinyatakan sumber:} Artikel mengidentifikasi
beberapa kekosongan sebelum penelitiannya dilakukan:

\begin{itemize}
\tightlist
\item
  Sedikit studi menguji NEG untuk dinamika penduduk saat ini dalam
  sistem urban yang lebih mapan.
\item
  Penelitian terdahulu lebih banyak memeriksa pusat urban besar dan area
  metropolitan, dengan sedikit perhatian pada interaksi di luar pusat
  utama.
\item
  Studi Partridge et al.~(2008b) yang disebut belum menguji interaksi
  antara metropolitan menengah dan besar, serta belum memusatkan
  perhatian pada pertumbuhan terbaru.
\item
  Hubungan bersih antara limpahan aglomerasi positif dan bayangan
  kompetisi negatif di sekitar pusat urban belum dipahami dengan baik.
\item
  Peran hubungan komuter, limpahan input-output, akses layanan, dan
  amenitas belum dapat dibedakan hanya dari model hierarki standar.
\end{itemize}

\textbf{Celah yang dijawab artikel:} Artikel menjawab sebagian celah
tersebut dengan memasukkan county di semua tingkat, jarak ke beberapa
ambang ukuran, jarak antar pusat dalam tingkat yang sama, periode
1990-2000 dan 1990-2006, serta uji pertumbuhan pekerjaan untuk mekanisme
komuter.

\textbf{Celah tersisa, inferensi review:} TXT masih menyisakan kebutuhan
untuk mengidentifikasi jalur rumah tangga dan perusahaan secara lebih
langsung, menguji perubahan efek jarak dari waktu ke waktu, dan
membedakan komponen perubahan penduduk. Dua kebutuhan pertama juga
disebutkan secara eksplisit sebagai agenda lanjutan oleh penulis;
formulasi mengenai dekomposisi komponen penduduk adalah inferensi review
dari ukuran dependen yang digunakan.

\textbf{Locator:} TXT baris 101-122 dan 123-169 untuk gap awal; baris
826-896 untuk gap mekanisme komuter; baris 997-1005 untuk gap yang
diarahkan ke penelitian lanjut; artikel halaman 446-447, 458-461.

\subsection{Limitations}\label{limitations-64}

TXT tidak memiliki bagian terpisah berjudul ``Limitations''. Batasan
berikut dipisahkan antara batasan atau asumsi yang dinyatakan sumber dan
batasan yang merupakan inferensi review dari desainnya.

\textbf{Batasan atau asumsi yang dinyatakan sumber:}

\begin{itemize}
\tightlist
\item
  Penulis menyatakan bahwa endogenitas statistik pada variabel selain
  jarak dapat membiaskan hasil jarak. Penggunaan karakteristik awal dan
  spesifikasi bertahap dimaksudkan untuk mengurangi kekhawatiran
  tersebut, bukan menghapusnya secara pasti.
\item
  Analisis mengasumsikan hierarki urban Amerika Serikat telah mapan pada
  1990, sehingga jarak ke pusat tingkat lebih tinggi dianggap telah
  ditentukan sebelumnya.
\item
  Empat puluh tiga county, sebagian besar county rural kecil,
  dihilangkan karena tidak memiliki data ekonomi. Dampak penghilangan
  tersebut terhadap keterwakilan county yang tersisa tidak dihitung
  dalam TXT.
\item
  Model large MA yang juga dinyatakan robust tidak seluruhnya dilaporkan
  demi keringkasan.
\item
  Cakupan geografis hanya 48 negara bagian daratan dan District of
  Columbia.
\item
  Pengukuran komuter menggunakan batas 160 km dan pengukuran korelasi
  spasial menggunakan bandwidth 200 km. Keduanya adalah keputusan desain
  yang dinyatakan, bukan batas alamiah yang divalidasi secara independen
  dalam TXT.
\item
  Definisi MA/MICRO yang digunakan terutama berasal dari 2003, sedangkan
  definisi yang lebih lama hanya dipakai dalam analisis sensitivitas.
  Perubahan batas dapat memengaruhi pengelompokan dan ukuran jarak.
\end{itemize}

\textbf{Batasan desain yang merupakan inferensi review:}

\begin{itemize}
\tightlist
\item
  Regresi pertumbuhan county dengan kondisi awal mendukung analisis
  asosiasi di bawah asumsi yang dipakai, tetapi TXT tidak menunjukkan
  eksperimen, kejutan eksogen, atau strategi identifikasi terpisah yang
  dapat memastikan hubungan kausal setiap koefisien jarak.
\item
  Variabel dependen adalah perubahan total penduduk. TXT tidak
  menyediakan dekomposisi migrasi bersih, kelahiran, dan kematian,
  sehingga jalur demografis yang menghasilkan pertumbuhan tidak dapat
  dipisahkan.
\item
  \emph{Industry mix} adalah proksi eksogen untuk pertumbuhan pekerjaan,
  bukan pengukuran langsung seluruh kesempatan kerja atau intensitas
  komuter. Karena itu, ketidakberartian atau keberartian proksi tidak
  identik dengan bukti lengkap mengenai tidak ada atau adanya komuter.
\item
  Proksi populasi dan jarak ke pusat urban dapat menangkap beberapa
  mekanisme sekaligus. Koefisien yang tersisa setelah kontrol belum
  memisahkan secara sempurna akses layanan, amenitas, input-output,
  harga lahan, dan seleksi lokasi.
\item
  Horizon utama 1990-2000 dan perluasan 1990-2006 tidak memberi deret
  waktu tahunan yang dapat memperlihatkan kapan perubahan mekanisme
  terjadi.
\end{itemize}

\textbf{Informasi yang tidak tersedia dalam TXT:} Tidak disebutkan
prosedur lengkap untuk menangani seluruh data hilang, pembobotan
observasi, kode replikasi, arsip raw data, validasi eksternal ukuran
jarak, atau formula rinci perubahan penduduk. Ketiadaan informasi ini
dicatat sebagai keterbatasan provenance, bukan sebagai bukti bahwa
prosedurnya tidak dilakukan.

\textbf{Locator:} TXT baris 339-349, 400-408, dan 421-441 untuk asumsi
serta risiko endogenitas; baris 369-378 dan 775-780 untuk cakupan dan
definisi; baris 598 dan 854-860 untuk pelaporan model dan batas komuter;
Appendix Table 1 baris 1116-1124.

\subsection{Challenges}\label{challenges-64}

\textbf{Laporan sumber:} Tantangan analitis yang dihadapi dan ditangani
artikel meliputi:

\begin{itemize}
\tightlist
\item
  Membedakan manfaat aglomerasi dari penalti kompetisi ketika keduanya
  berubah bersama jarak ke pusat urban.
\item
  Merepresentasikan hierarki yang memiliki beberapa tingkat layanan
  tanpa memaksakan satu efek jarak linear. Artikel mengatasi ini dengan
  jarak inkremental ke MA dan ambang ukuran yang makin tinggi.
\item
  Menentukan apakah jarak ke pusat dalam tingkat yang sama menunjukkan
  kompetisi spasial atau penguatan aglomerasi regional, karena kedua
  mekanisme memiliki prediksi arah yang berlawanan.
\item
  Menghadapi korelasi spasial residual antarcounty yang berdekatan; GMM
  Conley dengan pembobotan hingga 200 km dipakai untuk statistik-t yang
  robust.
\item
  Mengurangi multikolinearitas, endogenitas, variabel terlewat, dan
  proses autoregresif yang mungkin memengaruhi hasil jarak melalui model
  yang ditambah secara bertahap.
\item
  Membedakan kesempatan komuter dari externalities input-output dan
  amenitas. Pertumbuhan penduduk dan pekerjaan yang dapat ditentukan
  bersama diatasi secara pendekatan dengan proksi \emph{industry mix}.
\item
  Menangani perubahan definisi area metropolitan dan memasukkan
  kemungkinan suburbanisasi ke dalam jarak ke pusat aktual.
\end{itemize}

\textbf{Interpretasi penulis:} Tantangan terbesar bukan hanya teknis,
tetapi substantif: hubungan rural-pusat urban dapat ditopang komuter,
sedangkan hubungan antarpusat urban mungkin lebih terkait akses,
hubungan input-output, atau layanan. Persistennya efek jarak setelah uji
komuter dipakai penulis untuk menyatakan bahwa satu jalur tidak
menjelaskan seluruh pola.

\textbf{Inferensi review:} Tantangan-tantangan tersebut juga membatasi
ketegasan atribusi mekanisme. Artikel berhasil menguji beberapa prediksi
yang bersaing, tetapi tidak mengubah setiap mekanisme menjadi ukuran
langsung yang berdiri sendiri.

\textbf{Locator:} TXT baris 273-329 untuk konstruksi hierarki; baris
388-408 untuk masalah spesifikasi dan spasial; baris 451-462 untuk jarak
tingkat yang sama; baris 826-896 untuk tantangan mekanisme; artikel
halaman 449-450 dan 458-460.

\subsection{Practical Implications}\label{practical-implications-64}

\textbf{Laporan sumber:} Penulis menyatakan bahwa hasilnya relevan untuk
memahami pola tata guna lahan dan permukiman. Kedekatan dan akses ke
pusat urban yang lebih besar dianggap penting dalam membentuk ukuran dan
pertumbuhan county rural serta area urban dengan penduduk kurang dari
1,5 juta. Efek kedekatan dinilai besar dan dalam beberapa hasil tampak
mengungguli ukuran sendiri serta amenitas. Penulis menyarankan agar
perencana mempertimbangkan interaksi regional yang lebih luas dalam
strategi pembangunan, serta menilai bahwa tren tersebut belum
menunjukkan tanda mereda dalam periode terbaru yang dianalisis.

\textbf{Interpretasi penulis:} Implikasi kebijakan bukan untuk
melindungi semua pusat kecil dari pusat besar, karena pusat besar dapat
mendukung pertumbuhan tempat yang lebih kecil melalui akses dan
pekerjaan. Namun, pusat terbesar dapat tetap menimbulkan tekanan
persaingan pada MA menengah, sementara small MA dapat bersaing satu sama
lain. Strategi pembangunan perlu membaca efek tersebut menurut posisi
dalam hierarki.

\textbf{Inferensi review:} Temuan ini dapat membantu menghindari asumsi
perencanaan yang menyamakan kedekatan dengan ancaman bagi county kecil.
Akan tetapi, karena desain tidak menguji intervensi kebijakan dan
menyisakan risiko endogenitas, implikasi tersebut sebaiknya dipakai
sebagai pertimbangan analitis, bukan sebagai bukti bahwa pembangunan
infrastruktur atau kebijakan kedekatan tertentu pasti meningkatkan
penduduk.

TXT tidak menyebutkan instrumen kebijakan tertentu, skema pembiayaan,
target kepadatan, atau prosedur implementasi tata guna lahan. Informasi
spesifik tersebut tidak tersedia dalam sumber.

\textbf{Locator:} TXT baris 979-1005, bagian kesimpulan, artikel halaman
461; baris 747-780 untuk sprawl dan interaksi large MA.

\subsection{Applications}\label{applications-64}

\textbf{Laporan sumber:} Aplikasi yang dinyatakan secara langsung adalah
penggunaan temuan untuk strategi pembangunan, tata guna lahan, dan pola
permukiman yang mempertimbangkan interaksi regional antara county dan
pusat urban pada tingkat berbeda. Contoh Baltimore-Washington DC
digunakan untuk mengilustrasikan kemungkinan dua area metropolitan yang
berdekatan tumbuh atau melakukan sprawl menjadi satu area metropolitan
lebih besar; TXT menyajikannya sebagai contoh pola, bukan evaluasi
kebijakan.

\textbf{Inferensi review:} Kerangka jarak inkremental dapat dipakai
sebagai alat diagnosis untuk membandingkan akses county terhadap tingkat
layanan urban yang berbeda dan untuk memeriksa apakah suatu lokasi
berada dalam pola akses atau kompetisi. Penggunaan seperti itu adalah
inferensi dari desain penelitian, bukan aplikasi yang diuji atau
perangkat operasional yang dilaporkan artikel.

Tidak disebutkan penerapan pada satu kota tertentu sebagai studi kasus
kebijakan, implementasi oleh lembaga perencana, simulasi keputusan
lokasi, atau evaluasi hasil program. Tidak ada dasar dalam TXT untuk
menyatakan bahwa model telah digunakan di luar analisis empiris artikel.

\textbf{Locator:} TXT baris 749-756 untuk contoh Baltimore-Washington;
baris 997-1005 untuk aplikasi perencanaan; baris 294-329 dan 409-462
untuk struktur ukuran jarak.

\subsection{Future Research
(Optional)}\label{future-research-optional-64}

\textbf{Laporan sumber:} Penulis menyarankan agar penelitian berikutnya:

\begin{itemize}
\tightlist
\item
  Menilai bagaimana efek jarak berubah dari waktu ke waktu.
\item
  Menelusuri lebih dekat jalur pada sisi rumah tangga dan sisi
  perusahaan yang mentransmisikan efek jarak dan akses.
\end{itemize}

Penulis juga menyatakan bahwa perspektif dan model lain masih diperlukan
untuk menjelaskan evolusi serta interaksi dalam sistem urban yang telah
mapan.

\textbf{Inferensi review:} Agenda tersebut dapat diarahkan pada
pemisahan mekanisme komuter, amenitas, input-output, dan kompetisi,
tetapi TXT tidak memberikan desain, data, atau horizon waktu khusus
untuk penelitian lanjutan. Tidak ada agenda lanjutan lain yang
ditambahkan tanpa label inferensi.

\textbf{Locator:} TXT baris 997-1005, artikel halaman 461; kebutuhan
akan perspektif/model lain juga dinyatakan pada baris 163-169.

\subsection{Provenance Notes}\label{provenance-notes-64}

\begin{itemize}
\tightlist
\item
  Review ini hanya menggunakan TXT B21 yang disebutkan pada bagian
  Identitas Sumber. Tidak ada informasi dari PDF asli secara terpisah,
  web, sumber bibliografis lain, atau artikel lain yang ditambahkan
  untuk melengkapi isi.
\item
  Seluruh TXT dibaca dari baris 1 sampai 1.355. Pembacaan mencakup
  metadata PDF pada baris 1-28, artikel Inggris pada baris 30-1.314, dan
  abstrak Spanyol pada baris 1.320-1.354.
\item
  Locator utama memakai nomor baris TXT dan, jika terlihat di ekstraksi,
  halaman artikel. Nomor halaman yang disebut mengikuti penanda halaman
  pada TXT, sedangkan nomor baris merujuk pada posisi dalam file
  ekstraksi.
\item
  \texttt{Laporan\ sumber} berarti fakta, angka, definisi, atau
  pernyataan yang dilaporkan di dalam TXT.
  \texttt{Interpretasi\ penulis} berarti penjelasan mekanisme atau
  pembacaan hasil yang dikaitkan dengan penulis artikel.
  \texttt{Inferensi\ review} berarti pembacaan analitis yang dibuat
  dalam review ini dan tidak ditulis sebagai klaim langsung dari sumber.
\item
  Tabel 1, Tabel 2, Appendix Table 1, Appendix Table 2, dan Appendix
  Table 3 dirujuk sebagai locator sumber. Review ini tidak mereproduksi
  data dalam format tabel.
\item
  Metadata PDF pada TXT baris 8 dan artikel berbahasa Inggris pada baris
  30 mencantumkan DOI \texttt{10.1111/j.1435-5957.2008.00211.x}, tetapi
  baris 1.320 pada abstrak Spanyol menampilkan
  \texttt{10.1111/j.1435-5957.2009.00211.x}. Perbedaan tersebut tidak
  diselesaikan dengan sumber luar; DOI pertama dipakai karena muncul
  pada metadata PDF dan bagian artikel Inggris, sementara
  ketidakkonsistenan dicatat.
\item
  Metadata PDF hanya mencantumkan Mark D. Partridge pada field Author,
  sedangkan byline artikel mencantumkan empat penulis. Review memakai
  byline artikel dan mempertahankan perbedaan metadata sebagai catatan
  provenance.
\item
  Metadata menyebut rentang halaman 445-467, sementara penanda halaman
  yang terlihat pada teks Inggris berakhir pada halaman 466 sebelum
  materi abstrak Spanyol. Review mempertahankan rentang metadata dan
  tidak mengoreksi pagination berdasarkan sumber luar.
\item
  Beberapa karakter hasil ekstraksi mengalami kerusakan encoding,
  misalnya \texttt{GeografÃ\-a} pada kata kunci Spanyol. Istilah teknis,
  angka, dan keterangan tabel dipertahankan sesuai TXT; kata yang tampak
  seperti \texttt{typography} tidak dikoreksi menjadi istilah lain
  karena review tidak menggunakan sumber di luar TXT.
\item
  \texttt{n.a.} pada tabel sumber diperlakukan sebagai informasi yang
  tidak tersedia atau tidak berlaku untuk kolom tersebut, bukan sebagai
  nilai nol. Jika suatu informasi tidak ditampilkan dalam TXT, review
  menyatakannya sebagai tidak tersedia.
\end{itemize}

\section{Review Kode B22}\label{review-kode-b22}

\textbf{Identitas sumber}

\begin{itemize}
\tightlist
\item
  Kode kajian: B22.
\item
  Judul: ``Urban growth shadows''.
\item
  Penulis: David Cuberes, Klaus Desmet, dan Jordan Rappaport.
\item
  Publikasi: \emph{Journal of Urban Economics}, volume 123, tahun 2021,
  artikel 103334.
\item
  DOI: \texttt{10.1016/j.jue.2021.103334}.
\item
  Afiliasi yang tercantum: Department of Economics, Clark University;
  Department of Economics and Cox School of Business, Southern Methodist
  University; Federal Reserve Bank of Kansas City; NBER; dan CEPR.
\item
  Informasi waktu publikasi: diterima 2 Oktober 2019, diterima setelah
  revisi 9 Februari 2021, dan tersedia daring 1 Maret 2021.
\item
  Klasifikasi JEL: R11, R12, N91, dan N92. Kata kunci: urban shadows,
  urban access, commuting, spatial economics, urban systems, city
  growth, United States, dan 1840-2017.
\item
  Issue dan rentang halaman jurnal tidak disebutkan dalam file; file
  hanya mencantumkan volume dan nomor artikel. Status peer review tidak
  dapat diverifikasi dari TXT.
\end{itemize}

\subsection{Summarized Abstract}\label{summarized-abstract-65}

\begin{itemize}
\tightlist
\item
  \textbf{Laporan sumber:} Artikel menanyakan apakah kedekatan geografis
  suatu lokasi dengan pusat ekonomi besar meningkatkan atau menghambat
  pertumbuhan. Kedekatan dapat menimbulkan kompetisi yang merugikan,
  tetapi juga dapat memperbaiki akses pasar. Dengan data county dan
  wilayah metropolitan Amerika Serikat selama 1840-2017, penulis
  melaporkan bahwa kedekatan dengan pusat urban besar berkorelasi
  negatif dengan pertumbuhan pada 1840-1920 dan positif setelah 1920.
  {[}TXT baris 50-67{]}
\item
  \textbf{Laporan sumber:} Model spasial dua kota digunakan untuk
  menunjukkan bahwa perubahan jangka panjang pada biaya commuting
  antarkota dan intrakota dapat menghasilkan pergantian dari urban
  shadows ke urban access. Penurunan biaya pengiriman antarkota yang
  lebih cepat daripada penurunan biaya commuting juga diajukan sebagai
  penjelasan alternatif. {[}TXT baris 111-161{]}
\item
  \textbf{Interpretasi penulis:} Lima fakta bergaya yang ditekankan
  adalah pergantian rezim sekitar 1920, melemahnya urban access sekitar
  2000, meluasnya jangkauan geografis pengaruh pusat besar, menguatnya
  pengaruh ketika pusat tetangga lebih besar, dan menguatnya urban
  shadow ketika infrastruktur streetcar pusat besar membaik pada awal
  abad ke-20. {[}TXT baris 409-417, 461-483, 604-612, 888-896{]}
\item
  \textbf{Inferensi review:} Artikel menggabungkan dua lapisan bukti
  yang tidak setara: regresi empiris terutama mengidentifikasi korelasi
  pertumbuhan populasi dengan kedekatan, sedangkan model teoritis
  memberi mekanisme yang konsisten dengan pola tersebut. Karena penulis
  sendiri memperingatkan bahwa hasilnya tidak dapat ditafsirkan sebagai
  kausal, ringkasan ini tidak menyatakan bahwa kedekatan atau
  infrastruktur pasti menyebabkan pertumbuhan. {[}TXT baris 977-982{]}
\end{itemize}

\subsection{Problem Statement}\label{problem-statement-65}

\begin{itemize}
\tightlist
\item
  \textbf{Laporan sumber:} Masalah utama adalah ketegangan antara dua
  proses spasial. Pusat besar dapat menarik migran dan meningkatkan
  kompetisi atas sumber daya sehingga lokasi kecil di sekitarnya berada
  dalam ``urban shadow''; sebaliknya, klaster aktivitas ekonomi di
  dekatnya dapat memperbaiki akses pasar dan menguntungkan pertumbuhan
  lokasi yang lebih kecil. {[}TXT baris 72-85{]}
\item
  \textbf{Laporan sumber:} Literatur yang dirujuk menunjukkan bahwa
  tanda hubungan tersebut dapat berbeda menurut periode dan konteks.
  Artikel mempersoalkan apakah kedekatan dengan pusat besar
  menguntungkan atau merugikan pertumbuhan lokal, dengan fokus pada
  pertumbuhan populasi lokal Amerika Serikat selama hampir dua abad.
  {[}TXT baris 75-85{]}
\item
  \textbf{Interpretasi penulis:} Perubahan relatif antara shadow dan
  access diduga berkaitan dengan evolusi biaya commuting intrakota,
  commuting antarkota, dan pengiriman antarkota. Penulis menilai bahwa
  penjelasan tersebut perlu diuji bersama bukti historis transportasi
  dan model spasial. {[}TXT baris 133-166{]}
\item
  \textbf{Inferensi review:} Problem statement bukan sekadar mengukur
  efek jarak saat ini, melainkan menjelaskan mengapa tanda hubungan
  kedekatan berubah dari negatif menjadi positif dan mengapa jangkauan
  serta kekuatan hubungan juga berubah sepanjang waktu.
\end{itemize}

\subsection{Objectives}\label{objectives-65}

\begin{itemize}
\tightlist
\item
  \textbf{Laporan sumber:} Tujuan langsungnya adalah mendokumentasikan
  bagaimana korelasi pertumbuhan populasi lokal dengan keberadaan lokasi
  besar di sekitarnya berubah sepanjang 1840-2017, lalu mengidentifikasi
  periode dominasi urban shadows dan urban access. {[}TXT baris
  235-241{]}
\item
  \textbf{Laporan sumber:} Penulis juga bertujuan mengidentifikasi fakta
  bergaya tambahan: pelemahan urban access pada dekade terakhir,
  perubahan jangkauan geografis, pengaruh ukuran lokasi besar, serta
  hubungan infrastruktur commuting lokal dengan urban shadow pada awal
  abad ke-20. {[}TXT baris 118-149, 468-483, 484-612{]}
\item
  \textbf{Laporan sumber:} Tujuan teoretisnya adalah membangun model dua
  kota yang menghubungkan friksi spasial dengan pilihan tinggal, pindah,
  commuting, dan perdagangan, kemudian memeriksa apakah prediksi model
  konsisten dengan pola empiris. {[}TXT baris 154-172, 902-939{]}
\item
  \textbf{Inferensi review:} Artikel berfungsi sebagai studi
  historis-deskriptif dengan elaborasi mekanisme, bukan sebagai estimasi
  dampak kausal kebijakan transportasi atau evaluasi kesejahteraan
  individu.
\end{itemize}

\subsection{Literature Survey}\label{literature-survey-65}

\begin{itemize}
\tightlist
\item
  \textbf{Laporan sumber:} Artikel menempatkan dirinya dalam tradisi
  central place theory, terutama gagasan Christaller dan Losch bahwa
  skala ekonomi serta biaya transportasi dapat membentuk hierarki
  lokasi. Dalam kerangka tersebut, pusat besar dapat memberi spillover
  positif sekaligus membatasi pertumbuhan melalui kompetisi. {[}TXT
  baris 162-172{]}
\item
  \textbf{Laporan sumber:} Penulis menyebut bahwa urban economics hingga
  relatif baru banyak mengabaikan distribusi spasial antarkota. Model
  geografi ekonomi kuantitatif yang memasukkan ruang teratur disebut
  sebagai perkembangan yang lebih baru. {[}TXT baris 162-166, 226-233{]}
\item
  \textbf{Laporan sumber:} Studi empiris yang berfokus pada abad ke-20
  umumnya menemukan pengaruh positif kedekatan. Contoh yang disebut
  mencakup Partridge et al.~untuk 1990-2006, Rappaport untuk periode
  pascaperang, Dobkins dan Ioannides untuk interaksi kota berbeda
  ukuran, serta Liu et al.~untuk kota-kota di China. {[}TXT baris
  186-195{]}
\item
  \textbf{Laporan sumber:} Bukti untuk periode lebih awal cenderung
  menunjukkan urban growth shadows: Bosker dan Buringh untuk Eropa
  praindustri, Rauch untuk hinterland kota-kota Eropa, dan Beltran et
  al.~untuk munisipalitas Spanyol 1800-2000 yang berubah dari negatif
  sebelum 1950 menjadi semakin positif sesudahnya. {[}TXT baris
  196-211{]}
\item
  \textbf{Laporan sumber:} Alternatif structural transformation dibahas
  melalui Eckert dan Peters. Penulis menyatakan bahwa studi tersebut
  menemukan sebagian besar perubahan struktural terjadi di dalam county,
  sehingga perpindahan tenaga kerja dari pertanian ke nonpertanian
  mungkin bukan pendorong urutan pertama bagi fakta bergaya mereka.
  {[}TXT baris 213-225{]}
\item
  \textbf{Laporan sumber:} Literatur transportasi digunakan untuk
  menjelaskan penurunan biaya pengiriman dan commuting, termasuk
  railroad, kanal, highway, containerization, streetcar, automobile,
  suburban rail, congestion, dan opportunity cost of time. Rujukan yang
  dipakai antara lain Glaeser dan Kahn, Baum-Snow, LeRoy dan Sonstelie,
  Warner, Jackson, serta Duranton et al.~{[}TXT baris 575-605,
  607-793{]}
\item
  \textbf{Inferensi review:} Survei literatur bersifat naratif dan
  dipakai untuk membangun hipotesis serta konteks historis. TXT tidak
  memuat protokol pencarian, kriteria inklusi-eksklusi, atau corpus
  review sistematis; karena itu artikel tidak dapat diklasifikasikan
  sebagai systematic literature review berdasarkan informasi yang
  tersedia.
\end{itemize}

\subsection{Method Used}\label{method-used-65}

\begin{itemize}
\tightlist
\item
  \textbf{Laporan sumber:} Bagian empiris menggunakan regresi
  pertumbuhan populasi rata-rata tahunan selama interval dua puluh
  tahun, dengan interval terakhir 2000-2017. Variabel utama berupa
  indikator lokasi tetangga terdekat yang termasuk kategori besar dan
  berada pada bin jarak tertentu. {[}TXT baris 265-317{]}
\item
  \textbf{Laporan sumber:} Ukuran ``besar'' dibuat relatif terhadap
  distribusi populasi pada setiap tahun: minimal persentil ke-95 untuk
  moderately large dan persentil ke-99 untuk very large. Lokasi yang
  dianalisis pada baseline berada pada atau di bawah persentil ke-80,
  yaitu kuintil pertama sampai keempat. {[}TXT baris 318-336{]}
\item
  \textbf{Laporan sumber:} Persamaan baseline memasukkan indikator
  jarak, polinomial orde kelima dari log populasi awal, serta 47 kontrol
  geografis, cuaca, pesisir/pelabuhan, sungai, dan topografi. Jarak
  dihitung sebagai jarak garis lurus antara centroid geografis. Batas
  jarak maksimum dipilih agar sedikitnya 50 persen observasi berada
  dalam kategori pengecualian. {[}TXT baris 267-317{]}
\item
  \textbf{Laporan sumber:} Kesalahan baku dibuat robust terhadap
  korelasi spasial dengan metode Conley. Tabel juga melaporkan
  incremental R2, yaitu perbedaan R2 regresi lengkap dan R2 regresi yang
  hanya memakai kontrol tambahan. {[}TXT baris 368-373{]}
\item
  \textbf{Laporan sumber:} Robustness check menggunakan batas wilayah
  metropolitan yang ditetapkan berdasarkan sensus 1960 untuk periode
  1960-2017. Analisis lain mengganti ambang tetangga dari persentil
  ke-95 ke persentil ke-99, membandingkan ambang persentil ke-80, 90,
  95, dan 99, serta membedakan subsampel menurut ukuran lokasi. {[}TXT
  baris 430-458, 468-483, 484-512{]}
\item
  \textbf{Laporan sumber:} Uji streetcar memperluas regresi baseline
  dengan interaksi indikator tetangga besar dan perubahan keberadaan
  streetcar antara 1902 dan 1907. Spesifikasi alternatif memakai
  perubahan log dari satu ditambah jarak tempuh streetcar. Perubahan
  streetcar di lokasi sendiri dikontrol. {[}TXT baris 835-866{]}
\item
  \textbf{Laporan sumber:} Lapisan teoritis berupa model spasial dua
  kota dengan moving cost, biaya commuting intrakota, biaya commuting
  antarkota, dan biaya perdagangan antarkota. Kasus pertama menganggap
  dua barang sebagai substitusi sempurna, menyamakan biaya commuting
  intrakota dan antarkota, mengabaikan perdagangan, serta mengizinkan
  tinggal, pindah, atau commuting; kasus kedua mempertahankan preferensi
  CES, mengabaikan commuting antarkota, dan mengizinkan perdagangan.
  Model memuat kondisi ekuilibrium, tiga hasil formal, bukti di Appendix
  A, dan satu contoh numerik. {[}TXT baris 902-950, 1059-1143,
  1225-1360{]}
\item
  \textbf{Batas ketersediaan:} TXT tidak menyebut estimator perangkat
  lunak untuk regresi, kode replikasi, atau prosedur inferensi selain
  penggunaan standard error Conley. TXT hanya menyatakan bahwa program
  Python dipakai untuk mencocokkan jalur streetcar dengan county dan
  bahwa sisa kesalahan dicocokkan secara manual. {[}TXT baris 846-877{]}
\end{itemize}

\subsection{Dataset}\label{dataset-65}

\begin{itemize}
\tightlist
\item
  \textbf{Laporan sumber:} Dataset utama adalah data populasi county
  dari U.S. Census Bureau untuk 1840-2017. Intervalnya adalah 1840-1860,
  1860-1880, 1880-1900, 1900-1920, 1920-1940, 1940-1960, 1960-1980,
  1980-2000, dan 2000-2017. {[}TXT baris 244-255{]}
\item
  \textbf{Laporan sumber:} Unit geografis dibentuk sebagai gabungan
  county yang konsisten secara geografis dan wilayah metropolitan jika
  delineasinya tersedia. Untuk 1940 dan sebelumnya, county digabung
  menggunakan kriteria OMB 1950 yang diterapkan pada kondisi awal
  periode; mulai 1960, delineasi OMB setelah sensus digunakan. Analisis
  akhirnya merupakan hybrid metropolitan areas dan non-metropolitan
  counties. {[}TXT baris 186-220{]}
\item
  \textbf{Laporan sumber:} County longitudinal template dan map guide
  digunakan untuk menangani county yang terpecah atau bergabung. Batas
  pada tahun awal periode dipakai untuk menghitung pertumbuhan, dengan
  penyesuaian yang dijelaskan untuk kasus split dan merge. {[}TXT baris
  186-212, 252-255{]}
\item
  \textbf{Laporan sumber:} Dataset streetcar berasal dari dua edisi awal
  \emph{U.S. Census Special Report on Street and Electric Railways},
  yaitu data 1902 dan 1907. Data tersebut memuat kota utama operasi
  setiap jalur dan panjang jalur. Pada 1902 tercatat 521 dari 2.641
  county memiliki streetcar dengan hampir 20.000 mil total; pada 1907
  tercatat 693 dari 2.797 county dengan lebih dari 30.000 mil. {[}TXT
  baris 846-857{]}
\item
  \textbf{Laporan sumber:} Bila satu jalur melintasi beberapa county,
  total panjangnya dialokasikan menurut populasi county. Untuk analisis
  tambahan biaya commuting, artikel juga membahas Travel Time Index dari
  Texas A\&M Transportation Institute dan kecepatan perjalanan dari
  National Household Travel Survey, tetapi keduanya disajikan sebagai
  bukti kontekstual, bukan sebagai outcome utama regresi pertumbuhan.
  {[}TXT baris 715-741{]}
\item
  \textbf{Batas ketersediaan:} TXT tidak memberikan berkas data mentah,
  URL dataset yang dapat diunduh, kamus variabel, atau rincian observasi
  per county. Isi online appendix yang dirujuk juga tidak tersedia di
  dalam TXT.
\end{itemize}

\subsection{Population / Sample}\label{population-sample-65}

\begin{itemize}
\tightlist
\item
  \textbf{Laporan sumber:} Populasi analisis terdiri atas lokasi-lokasi
  Amerika Serikat yang direpresentasikan oleh county, metro area, atau
  county nonmetropolitan setelah harmonisasi geografis. Jumlah lokasi
  dalam dataset meningkat dari 862 pada 1840-1860 menjadi 2.370 pada
  1880-1900, mencapai maksimum 2.982 pada 1940-1960, lalu menurun karena
  lebih banyak county terserap ke metro area; periode 2000-2017 memiliki
  2.369 lokasi. {[}TXT baris 213-220{]}
\item
  \textbf{Laporan sumber:} Sampel regresi baseline dibatasi pada lokasi
  dengan populasi pada atau di bawah persentil ke-80 dan membandingkan
  keberadaan tetangga besar menurut jarak. Jumlah observasi pada Tabel 1
  berturut-turut adalah 691, 1.328, 1.844, 2.110, 2.357, 2.387, 2.283,
  2.104, dan 1.895 untuk sembilan periode. {[}TXT baris 323-336,
  348-365{]}
\item
  \textbf{Laporan sumber:} Tabel 2 dengan batas metro 1960 menggunakan
  2.283, 2.282, dan 2.283 observasi. Tabel 5 untuk uji streetcar
  menggunakan 2.238 observasi pada spesifikasi keberadaan dan 2.224 pada
  spesifikasi jarak tempuh. {[}TXT baris 430-458, 800-833{]}
\item
  \textbf{Interpretasi penulis:} Penyaringan lokasi besar diperlukan
  agar pertumbuhan lokasi kecil dan menengah dibandingkan dengan kondisi
  tetangga besar, bukan dengan pertumbuhan pusat terbesar itu sendiri.
  {[}TXT baris 318-336{]}
\item
  \textbf{Batas ketersediaan:} Tidak ada sampel individu, survei
  responden, bobot sampling, atau kriteria nonresponse karena objek
  empirisnya adalah lokasi geografis dan populasinya, bukan manusia yang
  diwawancarai. Prosedur pembentukan unit yang dijelaskan tidak
  menyebutkan sampling acak.
\end{itemize}

\subsection{Dependent Variables}\label{dependent-variables-65}

\begin{itemize}
\tightlist
\item
  \textbf{Laporan sumber:} Variabel dependen utama adalah pertumbuhan
  populasi rata-rata tahunan suatu lokasi, \texttt{g\_l}, selama setiap
  interval yang dianalisis. Regresi menghubungkan outcome ini dengan
  kategori jarak ke tetangga besar, sambil mengontrol populasi awal dan
  atribut geografis, cuaca, serta topografi. {[}TXT baris 301-317,
  368-373{]}
\item
  \textbf{Laporan sumber:} Uji streetcar tetap memakai pertumbuhan
  populasi rata-rata tahunan 1910-1920 sebagai variabel dependen;
  perubahan keberadaan atau panjang streetcar di tetangga dan lokasi
  sendiri adalah variabel penjelas/interaksi. {[}TXT baris 800-833,
  858-866{]}
\item
  \textbf{Interpretasi penulis:} Istilah urban shadow dan urban access
  bukan variabel dependen terpisah yang diukur langsung. Keduanya adalah
  pembacaan substantif atas korelasi negatif atau positif antara
  pertumbuhan lokasi kecil-menengah dan kedekatan dengan lokasi besar.
  {[}TXT baris 409-417{]}
\item
  \textbf{Laporan sumber:} Dalam model teori, keluaran yang dianalisis
  adalah pilihan tempat tinggal dan kerja, perubahan ukuran relatif dua
  kota, serta perpindahan penduduk dari kota kecil ke kota besar atau
  sebaliknya. Utility, produktivitas, harga, dan sewa tanah menjadi
  bagian dari kondisi model, bukan outcome observasional. {[}TXT baris
  941-1058, 1125-1151, 1250-1345{]}
\item
  \textbf{Batas ketersediaan:} Tidak ada outcome langsung untuk
  pendapatan, pekerjaan, produktivitas teramati, kesejahteraan, volume
  perdagangan, atau waktu commuting pada regresi utama.
\end{itemize}

\subsection{Results}\label{results-65}

\begin{itemize}
\tightlist
\item
  \textbf{Laporan sumber:} Tabel 1 menunjukkan dua rezim. Pada
  1840-1920, lokasi kecil dan menengah dengan tetangga moderately large
  cenderung memiliki pertumbuhan yang lebih rendah. Contoh 1840-1860
  adalah selisih -0,63 poin persentase per tahun untuk tetangga dalam
  1-50 km, -0,67 untuk 50-100 km, dan -0,30 untuk 100-150 km. Pada
  1880-1900, koefisien dalam tiga rentang tersebut masing-masing -1,10,
  -0,86, dan -0,63. {[}TXT baris 318-340, 348-373{]}
\item
  \textbf{Laporan sumber:} Pada 1920-2017, hubungan menjadi positif
  terutama untuk tetangga dalam 1-50 km. Koefisiennya 0,28 pada
  1920-1940, 1,11 pada 1940-1960, 1,19 pada 1960-1980, 0,50 pada
  1980-2000, dan 0,41 pada 2000-2017. Penulis mencatat bahwa kontribusi
  marginal indikator tetangga hanya 0,2-0,8 poin persentase R2 pada
  rezim negatif dan sekitar 3 poin persentase pada periode yang dimulai
  1940 dan 1960. {[}TXT baris 376-408{]}
\item
  \textbf{Laporan sumber:} Dengan batas metro 1960, Tabel 2 menunjukkan
  koefisien 1-50 km sebesar 1,19 pada 1960-1980, 1,43 pada 1980-2000,
  dan 0,60 pada 2000-2017. Hasil ini menggeser puncak hubungan positif
  ke 1980-2000 dan tetap menunjukkan pelemahan mulai sekitar 2000,
  sehingga pelemahan tidak sepenuhnya dijelaskan oleh perubahan
  delineasi metro. {[}TXT baris 386-408, 430-458{]}
\item
  \textbf{Laporan sumber:} Tabel 3 dengan ambang very large persentil
  ke-99 mempertahankan pemisahan rezim sekitar 1920. Pada rezim positif,
  jangkauan hubungan positif awalnya sangat lokal pada 1920-1940,
  kemudian bertambah sekitar 50 km pada tiap periode dua puluh tahun dan
  mencapai sekitar 250 km pada 2000-2017. Dalam uraian hasil, penulis
  juga menyatakan bahwa boost yang signifikan pada 2000-2017 dapat
  menjangkau tetangga sejauh 300 km, dengan besaran yang makin lemah
  ketika jarak meningkat. {[}TXT baris 468-483, 520-552, 578-602{]}
\item
  \textbf{Laporan sumber:} Untuk 1840-1920, bukti perluasan jangkauan
  urban shadow dinyatakan campuran, meskipun ada bukti lemah bahwa
  jangkauan meningkat antara akhir abad ke-19 dan awal abad ke-20.
  {[}TXT baris 470-483{]}
\item
  \textbf{Laporan sumber:} Tabel 4 menunjukkan bahwa kekuatan hubungan
  bergantung pada ukuran tetangga. Contoh 1880-1900: tetangga minimal
  persentil ke-80 dalam 50 km berkaitan dengan pertumbuhan 0,21 poin
  persentase lebih rendah, lalu ambang persentil ke-90 dan ke-95
  menambah penurunan 0,37 dan 0,74 poin persentase dalam contoh penulis.
  Pada 1960-1980, tetangga persentil ke-99 dalam 50 km memberi kenaikan
  marginal 2,43 poin persentase dibanding kategori persentil ke-95
  sampai di bawah persentil ke-99. {[}TXT baris 556-576, 619-668{]}
\item
  \textbf{Laporan sumber:} Pada uji streetcar, tetangga besar yang
  memperoleh streetcar antara 1902 dan 1907 berkorelasi dengan
  pertumbuhan 1910-1920 yang lebih rendah bagi lokasi kecil dan
  menengah. Untuk tetangga persentil ke-95 dalam 1-50 km, koefisien
  interaksi perubahan keberadaan streetcar adalah -1,02 dan interaksi
  perubahan log jarak tempuh adalah -0,28; perubahan streetcar di lokasi
  sendiri berkoefisien positif, masing-masing 0,05 dan 0,14 pada dua
  spesifikasi. Hasil dengan ambang persentil ke-90 serupa arahnya,
  tetapi kebanyakan tidak signifikan secara statistik. {[}TXT baris
  800-833, 835-867{]}
\item
  \textbf{Interpretasi penulis:} Dalam kasus pertama model, penurunan
  bertahap biaya commuting memindahkan ekuilibrium dari staying ke
  inter-city moving, kemudian inter-city moving and commuting, dan
  akhirnya inter-city commuting. Kota kecil mula-mula kehilangan
  penduduk karena migrasi ke kota besar, lalu dapat memperoleh penduduk
  ketika commuting menjadi lebih menarik daripada pindah. {[}TXT baris
  1125-1151{]}
\item
  \textbf{Interpretasi penulis:} Result 2 menyatakan bahwa kota besar
  yang lebih jauh memerlukan biaya commuting yang lebih rendah sebelum
  perpindahan rezim terjadi. Dengan population-weighted productivity
  tetap, Result 3 menyatakan bahwa ambang biaya untuk mulai menarik
  penduduk meningkat bersama ukuran relatif kota besar, sehingga pusat
  yang lebih besar mulai menimbulkan shadow pada biaya commuting yang
  lebih tinggi. {[}TXT baris 1198-1223, 1225-1243{]}
\item
  \textbf{Laporan sumber:} Kasus alternatif dengan perdagangan antarkota
  menunjukkan bahwa penurunan biaya pengiriman yang lebih cepat daripada
  penurunan biaya commuting dapat menghasilkan urutan naik lalu turunnya
  urban shadow. Contoh numerik memakai produktivitas 1,0 dan 1,5,
  elastisitas substitusi 3, total populasi 6, jarak antarkota 3, moving
  cost 0,001, serta biaya awal commuting dan perdagangan 0,25. {[}TXT
  baris 1245-1360{]}
\end{itemize}

\subsection{Findings}\label{findings-65}

\begin{itemize}
\tightlist
\item
  \textbf{Interpretasi penulis, Stylized Fact 1:} Urban shadows
  mendominasi Amerika Serikat pada 1840-1920, sedangkan urban access
  mendominasi 1920-2017. Ini didasarkan pada perubahan tanda korelasi,
  bukan pada estimasi dampak kausal. {[}TXT baris 376-417{]}
\item
  \textbf{Interpretasi penulis, Stylized Fact 2:} Urban access melemah
  sekitar 2000. Penulis menghubungkan hasil ini dengan perlambatan
  penurunan biaya commuting, tetapi juga mengakui bahwa perubahan batas
  metro mempersulit perbandingan antarperiode. {[}TXT baris 386-408,
  461-466, 756-793{]}
\item
  \textbf{Interpretasi penulis, Stylized Fact 3:} Jangkauan access
  meluas dari radius yang sangat lokal pada 1920-1940 ke radius yang
  jauh lebih besar pada 2000-2017. Jangkauan shadow pada periode awal
  tidak sekuat dan seseragam pola access pada periode akhir. {[}TXT
  baris 468-483{]}
\item
  \textbf{Interpretasi penulis, Stylized Fact 4:} Shadow dan access
  sama-sama cenderung lebih kuat ketika tetangga besar relatif lebih
  besar. {[}TXT baris 484-512, 556-612{]}
\item
  \textbf{Interpretasi penulis, Stylized Fact 5:} Pada awal abad ke-20,
  peningkatan infrastruktur commuting berupa streetcar di lokasi besar
  berkaitan dengan urban shadow yang lebih kuat terhadap lokasi kecil di
  dekatnya. {[}TXT baris 835-896{]}
\item
  \textbf{Interpretasi penulis:} Evolusi biaya commuting tunggal dapat
  menangkap banyak fakta bergaya, terutama kebangkitan dan penurunan
  urban shadows. Model perdagangan memberi mekanisme alternatif yang
  juga konsisten dengan hasil utama. {[}TXT baris 1158-1223,
  1322-1359{]}
\item
  \textbf{Inferensi review:} Temuan paling kuat secara deskriptif adalah
  perubahan tanda yang bertahan di berbagai ambang ukuran dan
  spesifikasi jarak. Namun, besaran incremental R2 yang kecil pada
  banyak periode awal, sifat relatif terhadap kategori pengecualian, dan
  peringatan nonkausal membatasi pembacaan temuan sebagai mekanisme
  pertumbuhan yang telah terbukti.
\end{itemize}

\subsection{Conclusions}\label{conclusions-65}

\begin{itemize}
\tightlist
\item
  \textbf{Laporan sumber:} Penulis menyimpulkan bahwa kedekatan dengan
  klaster urban besar berkorelasi negatif dengan pertumbuhan lokasi pada
  1840-1920 dan positif pada 1920-2017, dengan sebagian pelemahan
  korelasi positif dalam beberapa dekade terakhir. {[}TXT baris
  1331-1342{]}
\item
  \textbf{Interpretasi penulis:} Penurunan biaya commuting mula-mula
  mendorong penduduk dari lokasi kecil dekat pusat ke pusat besar,
  tetapi penurunan lebih lanjut memungkinkan penduduk tetap tinggal di
  lokasi kecil dan commuting ke pusat besar. Dengan demikian, biaya
  commuting yang menurun dapat terlebih dahulu merugikan lalu membantu
  pertumbuhan lokasi kecil di sekitar pusat besar. {[}TXT baris
  1343-1352{]}
\item
  \textbf{Interpretasi penulis:} Perdagangan antarkota merupakan
  penjelasan alternatif: penurunan biaya pengiriman memperbaiki akses
  pasar kota kecil dan mengurangi insentif pindah, sehingga dapat
  melemahkan urban shadow relatif terhadap penurunan biaya commuting.
  {[}TXT baris 1346-1359{]}
\item
  \textbf{Inferensi review:} Kesimpulan berlaku pada pola populasi
  lokasi Amerika Serikat yang dibangun dalam unit county/metro dan pada
  rentang historis yang dikaji. Kesimpulan tidak cukup untuk menyatakan
  bahwa satu jenis investasi transportasi akan menghasilkan efek yang
  sama di semua tempat atau periode.
\end{itemize}

\subsection{Contributions}\label{contributions-65}

\begin{itemize}
\tightlist
\item
  \textbf{Interpretasi penulis:} Kontribusi empiris utama adalah
  memperpanjang analisis hubungan pusat besar dan lokasi tetangga hingga
  hampir dua abad, bukan hanya periode abad ke-20 yang dominan dalam
  literatur yang mereka bahas. {[}TXT baris 186-211{]}
\item
  \textbf{Interpretasi penulis:} Artikel menyatukan perubahan tanda,
  perubahan jangkauan, ukuran tetangga, dan bukti streetcar dalam
  serangkaian fakta bergaya yang konsisten secara temporal. {[}TXT baris
  118-149, 235-241, 468-612, 835-896{]}
\item
  \textbf{Interpretasi penulis:} Kontribusi teoritisnya adalah model dua
  kota yang memisahkan moving cost, commuting intrakota, commuting
  antarkota, dan biaya perdagangan, lalu menunjukkan bahwa dua
  konfigurasi friksi yang berbeda dapat menghasilkan pola shadow-access.
  {[}TXT baris 918-939{]}
\item
  \textbf{Inferensi review:} Nilai tambah konseptual artikel terletak
  pada penjelasan non-monotonik: akses terhadap produktivitas pusat
  dapat meningkatkan lokasi kecil tanpa mengharuskan penduduk pindah.
  Namun, kontribusi ini adalah kecocokan teori dengan fakta bergaya,
  bukan identifikasi empiris bahwa commuting, bukan perdagangan atau
  spillover lain, adalah penyebab tunggal.
\end{itemize}

\subsection{Research Gap}\label{research-gap-65}

\begin{itemize}
\tightlist
\item
  \textbf{Laporan sumber:} Artikel mengidentifikasi bahwa urban
  economics sebelumnya relatif mengabaikan distribusi spasial antarkota
  dan bahwa banyak studi empiris tentang efek pusat besar berfokus pada
  abad ke-20. {[}TXT baris 162-166, 186-195{]}
\item
  \textbf{Laporan sumber:} Bukti historis yang dirujuk sebelum artikel
  ini terutama berasal dari Eropa praindustri, kota-kota Eropa, dan
  Spanyol. Artikel mengisi konteks Amerika Serikat, yang menurut penulis
  memiliki jalur settlement, mobilitas, dan adopsi automobile yang
  berbeda dari Spanyol. {[}TXT baris 196-212{]}
\item
  \textbf{Laporan sumber:} Penulis mengakui bahwa kemungkinan structural
  transformation sebagai penjelasan belum dapat ditentukan secara kuat
  oleh data mereka. {[}TXT baris 213-225{]}
\item
  \textbf{Inferensi review:} Gap yang tersisa adalah membedakan secara
  empiris mekanisme commuting, perdagangan, spillover teknologi, dan
  seleksi lokasi. Artikel menunjukkan bahwa beberapa mekanisme dapat
  menghasilkan pola yang sama, sehingga hasilnya belum memilih satu
  mekanisme secara konklusif.
\item
  \textbf{Inferensi review:} Gap lain adalah hubungan antara unit
  agregat county/metro dan perilaku individu. Model memuat pilihan
  individu, tetapi regresi utama tidak mengamati keputusan pindah,
  commuting, produktivitas, atau perdagangan pada tingkat individu.
\item
  \textbf{Batas ketersediaan:} TXT tidak menyatakan gap dalam heading
  tersendiri atau memberikan agenda pengisian gap dengan desain
  penelitian tertentu.
\end{itemize}

\subsection{Limitations}\label{limitations-65}

\begin{itemize}
\tightlist
\item
  \textbf{Laporan sumber:} Penulis menyatakan bahwa hasil tidak dapat
  ditafsirkan sebagai kausal. Hubungan infrastruktur commuting dan
  pertumbuhan dapat mencerminkan arah sebaliknya, yaitu ekspektasi
  pertumbuhan mendorong pembangunan infrastruktur di lokasi yang
  diperkirakan tumbuh. {[}TXT baris 977-982{]}
\item
  \textbf{Laporan sumber:} Perubahan delineasi metro dapat membuat
  county yang cepat tumbuh terserap ke metro pada periode berikutnya.
  Re-delineasi juga mengubah centroid dan jarak ke tetangga serta
  memilih county mana yang diserap, sehingga menimbulkan masalah
  perbandingan dan kemungkinan selection bias. {[}TXT baris 386-423{]}
\item
  \textbf{Laporan sumber:} Data tidak memungkinkan penulis mengambil
  posisi kuat tentang peran structural transformation. Perbedaan ukuran
  county, ekspansi frontier ke arah barat, dan variasi regional juga
  menjadi alasan penulis memasukkan kontrol geografis serta menjalankan
  pemeriksaan tambahan yang dirujuk dalam online appendix. {[}TXT baris
  213-225, 226-255, 410-414{]}
\item
  \textbf{Laporan sumber:} Hasil streetcar bersifat sugestif; koefisien
  negatif tidak selalu signifikan. Pada ambang tetangga persentil ke-90,
  hasil kebanyakan tidak signifikan. {[}TXT baris 977-982, 888-896{]}
\item
  \textbf{Laporan sumber:} Jarak diukur antarcentroid. Jangkauan 200 km
  pada periode akhir dapat melampaui commuting antarkota yang lazim,
  walaupun penulis memberi alasan bahwa jarak ke tepi metro dan wilayah
  interaksi yang tumpang tindih dapat membuat interpretasi tersebut
  tetap mungkin. {[}TXT baris 1171-1223{]}
\item
  \textbf{Laporan sumber:} Korelasi selalu relatif terhadap kategori
  pengecualian. Jika lokasi yang terisolasi tumbuh sangat lambat, semua
  lokasi yang memiliki tetangga dalam jarak jauh dapat tampak tumbuh
  lebih cepat secara relatif. {[}TXT baris 1191-1208{]}
\item
  \textbf{Laporan sumber:} Perbandingan bukti biaya commuting juga perlu
  hati-hati. Penulis menyebut perbandingan data National Household
  Travel Survey lintas waktu sebagai latihan yang rumit dan meminta
  interpretasi yang berhati-hati. Tahun pemisah subperiode juga bukan
  breakpoint yang presisi. {[}TXT baris 781-793{]}
\item
  \textbf{Laporan sumber:} Model teoritis disederhanakan menjadi dua
  kota dengan titik produksi yang eksogen, produktivitas eksogen, bentuk
  kota simetris, serta kasus khusus yang masing-masing mengabaikan
  perdagangan atau commuting antarkota. Model tersebut tidak memasukkan
  agglomeration economies standar sebagai sumber hubungan
  produktivitas-ukuran. {[}TXT baris 949-977, 1225-1243{]}
\item
  \textbf{Inferensi review:} Karena data utama adalah pertumbuhan
  populasi, hasil tidak langsung mengukur apakah access meningkatkan
  produktivitas, pendapatan riil, kesejahteraan, atau kualitas
  pekerjaan. Tidak adanya kode, data mentah, dan isi online appendix di
  TXT juga membatasi penilaian replikasi dari file ini saja.
\end{itemize}

\subsection{Challenges}\label{challenges-65}

\begin{itemize}
\tightlist
\item
  \textbf{Laporan sumber:} Tantangan konstruksi data pertama adalah
  perubahan batas county melalui split dan merge. Penulis membuat unit
  geografis yang konsisten pada setiap interval dua puluh tahun dengan
  menggabungkan county bila diperlukan. {[}TXT baris 186-212, 252-255{]}
\item
  \textbf{Laporan sumber:} Tantangan kedua adalah perubahan cara
  delineasi metro area. Penulis harus menggabungkan county menjadi metro
  area bila dapat dilakukan, tetapi menggunakan kriteria berbeda untuk
  1940 dan sebelumnya dibandingkan periode setelah 1960. {[}TXT baris
  202-220{]}
\item
  \textbf{Laporan sumber:} Tantangan desain regresi adalah
  menyeimbangkan jarak maksimum. Jarak terlalu tinggi menyisakan terlalu
  sedikit observasi dalam kategori pengecualian, sedangkan jarak terlalu
  rendah gagal menangkap pengaruh pada rentang 50-150 km. {[}TXT baris
  262-313{]}
\item
  \textbf{Laporan sumber:} Tantangan pengukuran tetangga adalah
  perubahan distribusi ukuran lokasi dari waktu ke waktu. Karena itu,
  penulis memakai persentil tahunan dan mengontrol polinomial populasi
  awal. {[}TXT baris 301-336{]}
\item
  \textbf{Laporan sumber:} Tantangan data streetcar adalah mencocokkan
  jalur dengan county ketika jalur melewati beberapa county. Program
  Python dipakai untuk matching, hasil yang tersisa diperiksa manual,
  dan panjang jalur dialokasikan menurut populasi. {[}TXT baris
  846-877{]}
\item
  \textbf{Inferensi review:} Tantangan substantif terbesar adalah
  observational equivalence: pola empiris yang sama dapat dijelaskan
  oleh commuting antarkota, perdagangan, spillover teknologi, seleksi
  batas metro, atau perubahan struktural. Artikel memodelkan beberapa
  alternatif, tetapi tidak memiliki desain empiris yang memisahkan
  semuanya.
\end{itemize}

\subsection{Practical Implications}\label{practical-implications-65}

\begin{itemize}
\tightlist
\item
  \textbf{Laporan sumber:} TXT tidak menyajikan rekomendasi kebijakan
  eksplisit, analisis biaya-manfaat, atau aturan perencanaan
  transportasi. Fokusnya adalah korelasi historis, mekanisme teoritis,
  dan fakta bergaya.
\item
  \textbf{Inferensi review:} Implikasi praktis yang berhati-hati adalah
  bahwa kedekatan dengan pusat besar tidak seharusnya diperlakukan
  sebagai keuntungan atau kerugian yang selalu searah. Dampaknya dapat
  bergantung pada biaya commuting, biaya perdagangan, ukuran pusat,
  jarak, dan periode perkembangan transportasi.
\item
  \textbf{Inferensi review:} Infrastruktur commuting di pusat besar
  dapat berkaitan dengan penguatan shadow pada konteks awal yang
  dianalisis, sementara biaya commuting yang lebih rendah pada tahap
  lanjut dapat memungkinkan access bagi lokasi kecil. Ini bukan estimasi
  dampak kebijakan; arah tersebut tidak boleh dipakai sebagai klaim
  kausal tanpa desain identifikasi tambahan.
\item
  \textbf{Inferensi review:} Untuk membaca pertumbuhan hinterland,
  indikator kedekatan perlu dipasangkan dengan pemeriksaan batas
  metropolitan dan kategori pembanding, karena perubahan unit dan
  pertumbuhan lokasi terisolasi dapat mengubah koefisien relatif. {[}TXT
  baris 391-423, 1191-1208{]}
\end{itemize}

\subsection{Applications}\label{applications-65}

\begin{itemize}
\tightlist
\item
  \textbf{Laporan sumber:} Tidak ada aplikasi kebijakan, studi kasus
  perencanaan, atau validasi prediktif di luar lokasi dan periode
  Amerika Serikat yang dianalisis dalam TXT.
\item
  \textbf{Inferensi review:} Kerangka artikel dapat dipakai sebagai
  perangkat konseptual untuk membandingkan apakah hubungan
  pusat-hinterland pada periode tertentu lebih menyerupai shadow atau
  access, dengan syarat unit geografis, jarak, ukuran pusat, dan
  kategori pengecualian dijelaskan secara eksplisit.
\item
  \textbf{Inferensi review:} Model dua kota juga dapat berfungsi sebagai
  alat skenario untuk memahami urutan pindah lalu commuting atau
  membandingkan mekanisme commuting dengan perdagangan. Namun, TXT tidak
  menunjukkan bahwa model tersebut telah dikalibrasi untuk prediksi
  suatu kota tertentu atau digunakan sebagai alat keputusan.
\item
  \textbf{Batas ketersediaan:} Tidak ada bukti dalam file bahwa penulis
  menerapkan hasilnya pada negara lain, membuat indeks operasional urban
  shadow/access, atau menguji manfaat praktis untuk kebijakan tertentu.
\end{itemize}

\subsection{Future Research
(Optional)}\label{future-research-optional-65}

\begin{itemize}
\tightlist
\item
  \textbf{Laporan sumber:} Tidak ada bagian agenda penelitian masa depan
  yang eksplisit dalam TXT.
\item
  \textbf{Inferensi review, bukan agenda yang dinyatakan penulis:}
  Penelitian lanjutan dapat menguji kausalitas pembangunan commuting
  dengan strategi yang mengatasi kemungkinan bahwa ekspektasi
  pertumbuhan mendahului infrastruktur. Dasarnya adalah pengakuan
  penulis bahwa hasil streetcar bersifat korelasional dan dapat
  menghadapi reverse causality. {[}TXT baris 977-982{]}
\item
  \textbf{Inferensi review, bukan agenda yang dinyatakan penulis:}
  Penelitian lanjutan dapat membedakan commuting, perdagangan, dan
  spillover teknologi sebagai mekanisme access, karena penulis
  menunjukkan bahwa penurunan biaya pengiriman dan penurunan biaya
  commuting sama-sama dapat menghasilkan pola urban shadow lalu urban
  access. {[}TXT baris 1207-1223, 1322-1359{]}
\item
  \textbf{Inferensi review, bukan agenda yang dinyatakan penulis:}
  Penelitian lanjutan dapat menguji sensitivitas hasil terhadap batas
  metro, jarak tepi metro, dan ukuran geografis aktual selain jarak
  antarcentroid, serta mengaitkan pertumbuhan populasi dengan outcome
  ekonomi yang tidak diukur di artikel ini. {[}TXT baris 391-423,
  1171-1208{]}
\item
  \textbf{Inferensi review, bukan agenda yang dinyatakan penulis:} Peran
  structural transformation dan heterogenitas lintas wilayah juga tetap
  perlu diuji lebih langsung karena artikel menyatakan bahwa data mereka
  belum memungkinkan posisi yang kuat mengenai faktor tersebut. {[}TXT
  baris 213-225{]}
\end{itemize}

\subsection{Provenance Notes}\label{provenance-notes-65}

\begin{itemize}
\tightlist
\item
  \textbf{Bahan yang digunakan:} Hanya file
  \texttt{/Users/mac/Documents/Mac/{[}2{]}\ Obsidian\ Vault/zettelkasten/source/journal/B22-cuberes-desmet-rappaport-2021-urban-growth-shadows-10.1016-j.jue.2021.103334.txt}
  yang digunakan untuk isi review ini. Tidak ada informasi dari luar TXT
  yang ditambahkan.
\item
  \textbf{Cakupan pembacaan:} Seluruh file, 1.486 baris, dibaca. File
  memuat blok \texttt{{[}PDF\ Metadata{]}} dan
  \texttt{{[}Extracted\ Text{]}}, serta teks artikel sampai Appendix A
  dan daftar referensi. {[}TXT baris 1-21, 1361-1486{]}
\item
  \textbf{Metadata ekstraksi:} TXT menyatakan sumber PDF berada pada
  \texttt{journal/B22-cuberes-desmet-rappaport-2021-urban-growth-shadows-10.1016-j.jue.2021.103334.pdf},
  diekstrak pada 2026-08-31 dengan \texttt{pdftotext\ -layout}, dan
  memiliki 17 halaman. {[}TXT baris 1-20{]}
\item
  \textbf{Konvensi locator:} Rujukan \texttt{{[}TXT\ baris\ x-y{]}}
  menunjuk langsung ke nomor baris file sumber. Rujukan bagian dan tabel
  dipertahankan bila membantu, termasuk bagian 2.1-2.6, 3.1-3.2,
  4.1-4.3, Tabel 1-5, dan Appendix A.
\item
  \textbf{Integritas isi:} Formula dan beberapa kolom tabel mengalami
  gangguan tata letak akibat ekstraksi teks, sehingga angka yang dipakai
  di review diutamakan dari uraian prose, label tabel, dan nilai yang
  terbaca jelas. Tidak ada kutipan langsung yang direkonstruksi dari
  bagian yang formatnya rusak.
\item
  \textbf{Informasi yang tidak tersedia dalam TXT:} Issue jurnal,
  rentang halaman jurnal selain nomor artikel, tanggal akses agen, data
  mentah, kode replikasi, kamus variabel, protokol pencarian literatur,
  isi lengkap online appendix, serta verifikasi status peer review tidak
  disebutkan atau tidak tersedia dalam file.
\item
  \textbf{Batas provenance:} Daftar referensi di akhir TXT dilaporkan
  sebagai bagian dari sumber, tetapi tidak diverifikasi terhadap
  publikasi eksternal. Pernyataan \texttt{Laporan\ sumber} merangkum
  atau memparafrasekan isi TXT; \texttt{Interpretasi\ penulis} merujuk
  pada pembacaan dan klaim yang dibuat artikel;
  \texttt{Inferensi\ review} adalah analisis pengolahan yang tidak boleh
  dibaca sebagai pernyataan penulis.
\end{itemize}

\section{Review B23: Cities, hinterlands and agglomeration
shadows}\label{review-b23-cities-hinterlands-and-agglomeration-shadows}

\textbf{Identitas sumber}

\begin{itemize}
\tightlist
\item
  Kode kajian: B23.
\item
  Penulis: Hannu Tervo, School of Business and Economics, University of
  Jyvaskyla, Finland.
\item
  Judul: \emph{Cities, hinterlands and agglomeration shadows: Spatial
  developments in Finland during 1880-2004}.
\item
  Jurnal: \emph{Explorations in Economic History}, volume 47 (2010),
  halaman 476-486.
\item
  DOI: \texttt{10.1016/j.eeh.2010.05.002}.
\item
  ISSN yang tercantum: \texttt{0014-4983}.
\item
  Riwayat artikel yang tersedia: diterima 9 Desember 2008. Tanggal
  publikasi yang lebih spesifik dan nomor isu tidak disebutkan dalam
  TXT.
\item
  Jenis sumber: artikel empiris jurnal. Status peer-review tidak
  diverifikasi dari TXT.
\item
  Basis pembacaan: seluruh TXT B23, termasuk blok metadata hasil
  ekstraksi PDF dan teks artikel pada baris 1-705.
\end{itemize}

\textbf{Konvensi pembacaan:} ``Laporan sumber'' berarti hal yang secara
eksplisit dilaporkan dalam artikel; ``interpretasi penulis'' berarti
penjelasan atau pemaknaan yang diberikan Tervo; dan ``inferensi review''
berarti penilaian analitis yang diturunkan dari TXT ini, bukan klaim
yang dinyatakan penulis. Locator utama menggunakan nomor baris TXT dan,
bila penanda halaman tersedia, halaman PDF tercetak.

\subsection{Summarized Abstract}\label{summarized-abstract-66}

\textbf{Laporan sumber:} Artikel menganalisis perkembangan spasial
jangka panjang di Finlandia dengan menguji dua prediksi \emph{new
economic geography} (NEG), yaitu semakin persistennya pola lokasi dan
semakin dominannya pusat-pusat pertumbuhan. Analisis menggunakan data
populasi regional 1880-2004. Hasil yang dilaporkan mendukung hipotesis
tersebut: perubahan peringkat dan distribusi rank-size selama
industrialisasi serta urbanisasi menunjukkan meningkatnya persistensi
struktur regional; analisis proses kausal antara pusat populasi dan
daerah hinterland menunjukkan pertumbuhan bersama pada periode sebelum
perang, sedangkan bayang-bayang aglomerasi mulai muncul setelah perang.
{[}Locator: TXT baris 45-56; PDF hlm. 476.{]}

\textbf{Interpretasi penulis:} Finlandia bergerak dari ekonomi yang
didominasi pertanian menuju ekonomi industri, jasa, dan teknologi. Dalam
pembacaan NEG, perubahan itu diharapkan meningkatkan konsentrasi
populasi, mengunci pola lokasi tertentu, dan memperkuat dominasi pusat
atas hinterland. {[}Locator: TXT baris 61-74, 127-151; PDF hlm.
476-478.{]}

\textbf{Inferensi review:} Abstrak menyatakan ``agglomeration shadow'',
tetapi bukti empirisnya dioperasionalkan terutama melalui populasi,
peringkat subwilayah, dan hubungan prediktif pertumbuhan
pusat-hinterland. Jadi, artikel tidak mengukur bayang-bayang aglomerasi
secara langsung melalui arus komuter, migrasi, perusahaan, pekerjaan,
atau output industri. {[}Locator: TXT baris 93-104, 208-214, 422-449.{]}

\subsection{Problem Statement}\label{problem-statement-66}

\textbf{Laporan sumber:} Persoalan penelitian adalah bagaimana struktur
regional dan hubungan antara kota dengan hinterland berubah dalam jangka
panjang ketika Finlandia mengalami industrialisasi dan urbanisasi. NEG
memprediksi bahwa pola lokasi pada fase awal belum persisten, kemudian
menjadi lebih stabil; teori yang sama memprediksi bahwa pusat yang
berkembang mula-mula dapat tumbuh bersama hinterland, tetapi setelah
titik balik ekonomi tercapai, pusat dapat menarik penduduk dan kegiatan
dari daerah sekitarnya sehingga terbentuk \emph{agglomeration shadow}.
{[}Locator: TXT baris 65-74, 147-158; PDF hlm. 476-478.{]}

Artikel juga menempatkan masalah ini dalam kekurangan bukti empiris.
Penulis menyebut kajian mengenai persistensi distribusi aktivitas
ekonomi sudah cukup banyak, tetapi kajian empiris tentang peran pusat
pertumbuhan masih terbatas. Untuk Finlandia, kajian spasial jangka
panjang juga disebut masih sedikit; karya Ottaviano dan Pinelli serta
Tervo sebelumnya terutama mencakup beberapa dekade terakhir. {[}Locator:
TXT baris 105-117; PDF hlm. 477.{]}

\textbf{Interpretasi penulis:} Perang Dunia II dipakai sebagai penanda
kasar perubahan dari Finlandia berbasis pertanian dengan petani yang
relatif tidak bergerak dan biaya transportasi tinggi menuju ekonomi
pascaperang dengan biaya transportasi menurun serta produksi industri
yang lebih mudah berpindah. Penulis mengakui bahwa waktu tepat transisi
tersebut sulit ditentukan. {[}Locator: TXT baris 93-100, 178-189; PDF
hlm. 477-478.{]}

\textbf{Inferensi review:} Problem statement artikel bukan sekadar
apakah kota membesar, melainkan apakah urutan industrialisasi-urbanisasi
mengubah dua hal sekaligus: stabilitas hierarki spasial dan arah
hubungan pertumbuhan pusat-periferi. Namun, TXT tidak menunjukkan
formulasi eksplisit mengenai identifikasi kausal industrialisasi atau
kebijakan tertentu; yang diuji adalah kesesuaian pola populasi dengan
prediksi NEG.

\subsection{Objectives}\label{objectives-66}

\textbf{Laporan sumber:} Tujuan utama yang dinyatakan adalah menguji dua
prediksi NEG melalui data populasi regional Finlandia 1880-2004.
{[}Locator: TXT baris 75-76, 578-581; PDF hlm. 476, 484.{]}

\begin{itemize}
\tightlist
\item
  Menguji apakah struktur regional yang masih berkembang pada fase
  pra-industrialisasi menjadi semakin persisten selama industrialisasi.
  Secara empiris, hal ini dilakukan dengan menilai perubahan peringkat
  subwilayah dari waktu ke waktu dan distribusi rank-size pada berbagai
  tahun. {[}Locator: TXT baris 147-151, 251-264, 322-324.{]}
\item
  Menguji apakah pusat-pusat pertumbuhan menjadi semakin dominan
  terhadap hinterland lokalnya. Secara empiris, hal ini dilakukan dengan
  membandingkan pangsa populasi pusat serta menguji hubungan pertumbuhan
  pusat dan hinterland sebelum dan sesudah perang. {[}Locator: TXT baris
  75-76, 100-104, 373-421.{]}
\item
  Menguji apakah hubungan pusat-hinterland bersifat homogen atau
  heterogen dan apakah arahnya positif atau negatif di tiap wilayah pada
  periode pascaperang. {[}Locator: TXT baris 440-449, 547-574; PDF hlm.
  483-484.{]}
\end{itemize}

Hipotesis substantif penulis adalah bahwa struktur regional makin
persisten selama industrialisasi, dan bahwa pusat serta hinterland
mula-mula saling mendukung tetapi kemudian pusat yang ``beruntung''
dapat menarik mayoritas perusahaan dan penduduk sehingga menghasilkan
bayang-bayang aglomerasi. {[}Locator: TXT baris 147-158; PDF hlm.
478.{]}

\textbf{Informasi yang tidak tersedia:} TXT tidak menyatakan tujuan
terpisah untuk mengevaluasi efektivitas setiap instrumen kebijakan
regional, mengestimasi dampak industri tertentu, atau menguji mekanisme
migrasi secara langsung.

\subsection{Literature Survey}\label{literature-survey-66}

\textbf{Laporan sumber:} Kerangka teorinya bertumpu pada NEG dan
mekanisme \emph{cumulative causation}. Dalam model core-periphery
Krugman, perubahan dari produksi primer ke manufaktur dan jasa dapat
membuat permintaan serta kegiatan produksi terkonsentrasi. Pada model
yang diperluas, hubungan ke belakang dan ke depan menjadi kekuatan
sentripetal, sedangkan tanah yang tidak bergerak menjadi kekuatan
sentrifugal; kombinasi skala ekonomi dan biaya transportasi dapat
membentuk hierarki kota. {[}Locator: TXT baris 125-146; PDF hlm.
477-478.{]}

Artikel juga menghubungkan NEG dengan perdebatan
\emph{backwash/polarization} dan \emph{spread/trickling down}.
Sumber-sumber yang dirujuk membedakan penarikan tenaga kerja dan sumber
daya menuju pusat dari penyebaran pabrik, desentralisasi penduduk,
inovasi, investasi, dan sikap pertumbuhan menuju hinterland. {[}Locator:
TXT baris 151-158; PDF hlm. 478.{]}

Untuk persistensi dan ukuran kota, penulis membahas teori pertumbuhan
acak, teori fundamental lokasi, hukum Zipf atau aturan rank-size, dan
model hierarki perkotaan NEG. Pertumbuhan acak dan fundamental lokasi
dapat menghasilkan pola rank-size, sedangkan model NEG yang ditambah
kekuatan penyebaran berupa kemacetan dapat mereplikasi perubahan
distribusi tersebut. {[}Locator: TXT baris 265-278; PDF hlm. 479-480.{]}

Literatur empiris yang dirangkum mencakup perkembangan kota di Amerika
Serikat, Jepang, Jerman Barat dan Timur, serta studi mengenai distribusi
ukuran kota dan multiple equilibria. Untuk Finlandia, penulis merujuk
pada Ottaviano dan Pinelli mengenai prediksi NEG serta Tervo (2009)
mengenai hubungan pusat-hinterland pada 1970-2004. {[}Locator: TXT baris
105-117, 272-278, 653-704.{]}

\textbf{Interpretasi penulis:} Literatur tersebut menyediakan dua
ekspektasi yang dapat diuji secara historis: pola lokasi akan makin
stabil, dan hubungan pusat-hinterland akan bergeser dari pertumbuhan
bersama menuju dominasi pusat ketika industrialisasi dan urbanisasi
menguat.

\textbf{Inferensi review:} Survei literatur ini bersifat naratif dan
teoritis, bukan tinjauan sistematis. TXT tidak memuat protokol
pencarian, kriteria inklusi, atau batas korpus literatur. Daftar studi
di atas karena itu harus dibaca sebagai literatur yang dipilih dan
dilaporkan oleh penulis, bukan sebagai verifikasi independen atas
seluruh keadaan literatur.

\subsection{Method Used}\label{method-used-66}

\textbf{Laporan sumber:} Desain penelitian adalah analisis empiris
longitudinal dengan dua periode, yaitu 1880-1940 dan 1950-2004.
Pemisahan ini dimaksudkan untuk membandingkan fase sebelum dan sesudah
perubahan struktural yang diperkirakan berhubungan dengan Perang Dunia
II. {[}Locator: TXT baris 93-104, 206-214; PDF hlm. 477-479.{]}

\begin{itemize}
\tightlist
\item
  Untuk persistensi struktur, penulis membandingkan peringkat setiap
  subwilayah pada tanggal berturut-turut. Untuk setiap tanggal,
  penyimpangan absolut peringkat dijumlahkan; rata-rata jumlah
  subperiode menjadi statistik perubahan struktur regional. Perbedaan
  antara periode diuji dengan \emph{paired samples t test}. {[}Locator:
  TXT baris 251-264; PDF hlm. 479.{]}
\item
  Untuk distribusi rank-size, penulis mengestimasi persamaan
  \texttt{log(size)\ =\ log(constant)\ -\ alpha\ log(rank\ -\ 1/2)}
  dengan OLS. Pengurangan \texttt{1/2} mengikuti koreksi
  Gabaix-Ibragimov terhadap bias OLS pada sampel kecil. \texttt{R2}
  dipakai sebagai ukuran kecocokan dan \texttt{alpha} sebagai eksponen
  Pareto; \texttt{alpha\ =\ 1} menjadi kondisi yang dikaitkan dengan
  hukum Zipf. {[}Locator: TXT baris 300-321, 326-327; PDF hlm.
  480-481.{]}
\item
  Persamaan rank-size diestimasi pada 14 tahun: 1880, 1890, 1900, 1910,
  1920, 1930, 1940, 1950, 1960, 1970, 1980, 1990, 2000, dan 2004.
  {[}Locator: TXT baris 322-324, 355-367; PDF hlm. 480-481.{]}
\item
  Untuk hubungan pusat-hinterland, penulis memakai perluasan uji
  kausalitas Granger dalam kerangka panel yang memungkinkan
  heterogenitas antarwilayah. Model panel VAR memakai pertumbuhan
  populasi pusat dan hinterland, efek individual, koefisien autoregresif
  yang sama antarunit, dan koefisien hubungan pusat-hinterland yang
  boleh berbeda antarunit. {[}Locator: TXT baris 416-439; PDF hlm.
  482-483.{]}
\item
  Uji dilakukan untuk arah pusat ke hinterland dan arah sebaliknya, pada
  lag \texttt{t} dan \texttt{t-1}. Kasus instantaneous atau lag
  \texttt{t} diizinkan karena interval data sepuluh tahun; lag yang
  lebih panjang tidak dapat diuji karena jumlah periode sedikit.
  {[}Locator: TXT baris 440-445; PDF hlm. 483.{]}
\item
  Variabel diubah ke logaritma natural dan diferensiasi untuk
  menghilangkan kemungkinan unit root serta mencapai stationarity.
  Pengujian menggunakan statistik Wald, jumlah kuadrat residual model
  tidak terbatas dan model terbatas, estimasi fixed effects yang
  dipandang ekuivalen dengan MLE dalam konteks tersebut, serta
  constrained regression. {[}Locator: TXT baris 422-449; PDF hlm.
  482-483.{]}
\item
  Prosedur Granger bersarang terdiri atas pengujian homogeneous
  instantaneous non-causality (HINC), homogeneous causality (HC), dan
  heterogeneous non-causality (HENC). HENC dipakai untuk
  mengidentifikasi wilayah individual yang menyumbang hubungan pada
  periode pascaperang. {[}Locator: TXT baris 416-421, 450-489, 547-567;
  PDF hlm. 482-484.{]}
\end{itemize}

\textbf{Inferensi review:} Dalam pengertian ekonometrik yang dijelaskan
TXT, ``kausalitas'' Granger berarti tambahan informasi masa kini atau
masa lalu meningkatkan kemampuan memprediksi variabel lain. Itu bukan
dengan sendirinya identifikasi dampak struktural atau bukti bahwa
pertumbuhan pusat secara eksogen menyebabkan perubahan hinterland.
{[}Locator: TXT baris 422-426.{]} Kesimpulan tentang bayang-bayang
aglomerasi perlu dibaca dengan batasan ini.

\textbf{Informasi yang tidak tersedia:} TXT tidak menyebut perangkat
lunak, kode analisis, kriteria pemilihan jumlah lag selain keterbatasan
periode, atau prosedur robustness check di luar rangkaian pengujian yang
dijelaskan.

\subsection{Dataset}\label{dataset-66}

\textbf{Laporan sumber:} Dataset berupa data populasi berbasis
munisipalitas yang diproduksi Statistics Finland dan diolah ulang agar
cocok dengan pembagian wilayah tahun 2005. Rentangnya 1880-2004 dengan
interval satu dekade, kecuali interval terakhir yang panjangnya empat
tahun. Data 1880-1940 berasal dari register gereja, sedangkan data mulai
1950 berasal dari sensus penduduk. {[}Locator: TXT baris 206-214; PDF
hlm. 478-479.{]}

Analisis menggunakan pembagian wilayah kontemporer tahun 2005: NUTS
level 3 untuk region dan NUTS level 4 untuk subregion. Pada 2005
terdapat 20 region dan 77 subregion; Åland dikeluarkan sehingga analisis
memakai 19 region dan 74 subregion. Subregion terdiri dari dua sampai
tiga belas munisipalitas, sedangkan region terdiri dari dua sampai tujuh
subregion. {[}Locator: TXT baris 213-230; PDF hlm. 478-479.{]}

Penulis merekonstruksi sejarah populasi ke pembagian 2005. Jumlah
munisipalitas adalah 432 pada 2005 dan sebelumnya lebih dari 600. Jika
satu munisipalitas digabungkan seluruhnya, penyesuaian relatif
sederhana. Jika munisipalitas dibagi ke beberapa munisipalitas baru,
populasi historisnya ditambahkan kepada munisipalitas penerima mayoritas
penduduk. Penulis menyatakan persoalan ini kecil karena dalam banyak
kasus munisipalitas yang terdampak tetap berada dalam subregion yang
sama. {[}Locator: TXT baris 231-238; PDF hlm. 479.{]}

Wilayah yang diserahkan Finlandia kepada Uni Soviet setelah Perang Dunia
II tidak muncul dalam pembagian wilayah sekarang. Penulis menyatakan
ketiadaan itu tidak terlalu bermasalah untuk analisis, tetapi mengakui
hasil pra-perang sampai tingkat tertentu menjadi bias. Pemukiman sekitar
11\% penduduk yang dievakuasi juga menyebabkan kenaikan mendadak pada
populasi banyak munisipalitas; pemisahan periode dimaksudkan untuk
mengurangi masalah ini. {[}Locator: TXT baris 239-242, 355-356; PDF hlm.
479-481.{]}

Sebagai gambaran deskriptif, rata-rata populasi pusat pada 19 region
meningkat dari 37.975 pada 1880 menjadi 182.834 pada 2004, sedangkan
rata-rata populasi daerah periphery meningkat dari 58.169 menjadi
91.381. Nilai tahunan lengkap beserta standard error dilaporkan dalam
Table 1. {[}Locator: TXT baris 331-351; PDF hlm. 480.{]}

\textbf{Informasi yang tidak tersedia:} TXT tidak menyediakan berkas
data mentah, daftar seluruh munisipalitas hasil harmonisasi, kamus
variabel, prosedur akses data Statistics Finland, atau informasi rinci
tentang pembobotan dan penanganan observasi yang hilang.

\subsection{Population / Sample}\label{population-sample-66}

\textbf{Laporan sumber:} Tidak ada sampel individu atau responden. Unit
analisis utama adalah seluruh 19 region Finlandia yang dipakai dalam
studi dan 74 subregion di dalamnya. NUTS 4 diperlakukan sebagai
pendekatan yang cukup masuk akal untuk wilayah pasar kerja lokal,
walaupun munisipalitas mungkin lebih dekat dengan wilayah fungsional
pada awal periode dan wilayah fungsional kemudian membesar. {[}Locator:
TXT baris 222-230; PDF hlm. 479.{]}

Untuk analisis pusat-hinterland, setiap region memiliki satu pusat:
subregion dengan populasi terbesar pada 2004. Subregion lain dalam
region yang sama disebut periphery atau hinterland. Dalam semua kasus,
subregion tersebut juga merupakan pusat administratif region. Penulis
mengakui bahwa klasifikasi kurang jelas di beberapa region yang
mempunyai lebih dari satu subregion berukuran cukup besar. {[}Locator:
TXT baris 243-248; PDF hlm. 479.{]}

Catatan historis menunjukkan bahwa 15 dari 19 pusat yang ditetapkan
berdasarkan 2004 tetap menjadi subregion terbesar sepanjang periode.
Empat pengecualian adalah Kymenlaakso, South Karelia, Northern Savo, dan
Lapland; keempat subregion tersebut tetap merupakan pusat administratif
serta menjadi yang terbesar pada 2004. {[}Locator: TXT baris 280-283;
PDF hlm. 479.{]}

\textbf{Inferensi review:} Pemilihan unit adalah cakupan geografis yang
telah ditentukan, bukan pengambilan sampel probabilistik. Penggunaan
pusat yang ditetapkan dari kondisi 2004 membuat kategori
pusat-hinterland konsisten untuk analisis, tetapi tidak sepenuhnya
merepresentasikan pusat terbesar pada setiap tanggal historis, terutama
pada empat region yang mengalami perubahan peringkat.

\textbf{Informasi yang tidak tersedia:} TXT tidak melaporkan ukuran
populasi total yang dianalisis untuk setiap tahun, jumlah penduduk
menurut karakteristik sosial-demografis, atau prosedur sampling karena
penelitian ini tidak menggunakan sampel manusia.

\subsection{Dependent Variables}\label{dependent-variables-66}

Artikel tidak memakai satu variabel dependen tunggal; variabel keluaran
berubah sesuai pengujian.

\begin{itemize}
\tightlist
\item
  Dalam analisis persistensi, objek keluaran adalah peringkat relatif
  subregion dan perubahan peringkat antartanggal. Statistiknya berupa
  jumlah penyimpangan absolut peringkat yang dirata-ratakan untuk setiap
  subperiode. {[}Locator: TXT baris 251-264.{]}
\item
  Dalam analisis rank-size, \texttt{log(size)} atau logaritma populasi
  subregion adalah variabel dependen dalam persamaan OLS.
  \texttt{log(rank\ -\ 1/2)} menjadi variabel penjelas, sedangkan
  \texttt{alpha} adalah parameter distribusi yang diestimasi, bukan
  variabel populasi langsung. \texttt{R2} mengukur kecocokan distribusi.
  {[}Locator: TXT baris 300-321, 379-401.{]}
\item
  Dalam analisis panel Granger, pertumbuhan populasi pusat atau
  pertumbuhan populasi hinterland bergantian menjadi \texttt{y}, yaitu
  variabel yang diprediksi. Pertumbuhan wilayah yang lain menjadi
  \texttt{x}; arah pusat ke hinterland dan arah hinterland ke pusat
  diuji secara terpisah. {[}Locator: TXT baris 422-449.{]}
\item
  Pangsa penduduk yang tinggal di pusat juga dipakai sebagai ukuran
  deskriptif dominasi pusat, meningkat dari 39\% pada 1880 menjadi 67\%
  pada 2004. Pangsa ini bukan outcome dari persamaan kausalitas yang
  dijelaskan. {[}Locator: TXT baris 373-415; PDF hlm. 482.{]}
\end{itemize}

\textbf{Inferensi review:} Operasionalisasi dependent variable berfokus
pada perubahan ukuran dan kemampuan prediksi antarwilayah. Ia tidak
langsung mengukur kesejahteraan, pendapatan, produktivitas, pekerjaan,
arus migrasi, arus komuter, atau performa perusahaan. Karena itu,
``dominasi'' dalam artikel terutama berarti dominasi demografis
berdasarkan populasi.

\subsection{Results}\label{results-66}

\textbf{Persistensi peringkat:} Statistik perubahan struktur bernilai
2,26 pada periode pra-perang dan 1,70 pada periode pascaperang. Uji
\emph{paired samples t} menghasilkan \texttt{t\ =\ 3,43} dengan 73
derajat kebebasan dan \texttt{p\ =\ 0,001}, sehingga perbedaan fluktuasi
peringkat dinyatakan signifikan. {[}Locator: TXT baris 253-264; PDF hlm.
479.{]}

\textbf{Distribusi rank-size:} Kecocokan distribusi meningkat secara
umum. \texttt{R2} bernilai 0,69 pada 1880, turun ke 0,65 pada 1900, dan
mencapai 0,95 pada 2004; rata-rata \texttt{R2} adalah 0,78 pada periode
pra-perang dan 0,93 pada periode pascaperang. Nilai \texttt{alpha}
meningkat dari 0,55 pada 1880 menjadi 0,88 pada 2004. Nilai per tahun
yang dilaporkan dalam Table 2 adalah: 1880 \texttt{alpha\ =\ 0,55},
\texttt{R2\ =\ 0,69}; 1890 \texttt{0,55}, \texttt{0,71}; 1900
\texttt{0,60}, \texttt{0,65}; 1910 \texttt{0,57}, \texttt{0,82}; 1920
\texttt{0,58}, \texttt{0,84}; 1930 \texttt{0,58}, \texttt{0,87}; 1940
\texttt{0,61}, \texttt{0,89}; 1950 \texttt{0,62}, \texttt{0,91}; 1960
\texttt{0,67}, \texttt{0,93}; 1970 \texttt{0,74}, \texttt{0,94}; 1980
\texttt{0,79}, \texttt{0,93}; 1990 \texttt{0,81}, \texttt{0,94}; 2000
\texttt{0,86}, \texttt{0,94}; dan 2004 \texttt{0,88}, \texttt{0,95}.
{[}Locator: TXT baris 355-367, 379-402; PDF hlm. 480-481.{]}

Dalam pengujian terhadap \texttt{alpha\ =\ 1}, hasil Table 2 menolak
hipotesis tersebut pada tingkat 1\% untuk tahun 1880-1960 dan pada
tingkat 5\% untuk 1970; tidak ada tanda penolakan untuk 1980-2004.
Penulis karena itu menyatakan hukum Zipf belum berlaku pada periode
pra-perang dan secara bertahap menjadi valid mulai 1980. Sumber juga
mencatat bahwa pada 1990 Helsinki berpenduduk 1.030.200, sedangkan
Kouvola sebagai subregion peringkat ke-10 berpenduduk 102.200;
penyimpangan terbesar berada antara Helsinki dan subregion-subregion
berikutnya, termasuk Tampere yang lebih kecil daripada prediksi aturan
rank-size. {[}Locator: TXT baris 292-299, 322-367, 400-402; PDF hlm.
480-481.{]}

\textbf{Dominasi pusat:} Pangsa populasi 19 pusat meningkat dari 39\%
pada 1880 menjadi 67\% pada 2004, dan sebagian besar kenaikan terjadi
sesudah perang. Pangsa penduduk di pusat meningkat di semua region pada
periode pascaperang, sedangkan pada periode pra-perang menurun di enam
region. {[}Locator: TXT baris 373-415; PDF hlm. 482.{]}

\textbf{Uji HINC dan HC:} Pada lag \texttt{t}, statistik HINC untuk
kausalitas pusat ke hinterland adalah 3,30*** pada 1880-1940 dan 5,59***
pada 1950-2004. Untuk arah hinterland ke pusat, nilainya 4,73*** dan
5,02***. Pada lag \texttt{t-1}, masing-masing hubungan tidak signifikan:
pusat ke hinterland bernilai 0,88 dan 0,40, sedangkan hinterland ke
pusat bernilai 0,74 dan 1,21. Dengan demikian, HINC ditolak pada lag
\texttt{t}, yang berarti setidaknya satu region memiliki hubungan
Granger pada masing-masing arah; hasil ini tidak membuktikan bahwa semua
region memiliki hubungan tersebut. {[}Locator: TXT baris 440-468,
491-506; PDF hlm. 483.{]}

Untuk HC, statistik pada periode pra-perang tidak signifikan pada kedua
arah, masing-masing 1,25. Pada periode pascaperang, HC bernilai 3,24***
untuk pusat ke hinterland dan 2,79*** untuk hinterland ke pusat. Hasil
ini menunjukkan bahwa koefisien kausalitas tidak dapat diperlakukan
homogen antarregion pada periode pascaperang. {[}Locator: TXT baris
469-489, 491-506; PDF hlm. 483.{]}

\textbf{Uji HENC per region:} HENC hanya dihitung untuk 1950-2004 dengan
lag \texttt{t}. Tanda \texttt{*}, \texttt{**}, dan \texttt{***}
masing-masing berarti penolakan hipotesis nol pada tingkat 10\%, 5\%,
dan 1\%. Nilai dan tanda efek yang dilaporkan, ditulis sebagai pusat ke
hinterland lalu hinterland ke pusat, adalah sebagai berikut: Uusimaa
\texttt{0,71\ (+)} dan \texttt{3,70*\ (+)}; Itä-Uusimaa
\texttt{0,19\ (+)} dan \texttt{0,11\ (+)}; Southwest Finland
\texttt{2,28\ (-)} dan \texttt{4,64**\ (-)}; Satakunta
\texttt{1,43\ (+)} dan \texttt{2,30\ (+)}; Häme \texttt{0,53\ (+)} dan
\texttt{0,55\ (+)}; Tampere Region \texttt{0,03\ (-)} dan
\texttt{0,11\ (-)}; Päijät-Häme \texttt{0,66\ (+)} dan
\texttt{1,96\ (+)}; Kymenlaakso \texttt{5,66**\ (+)} dan
\texttt{4,55**\ (+)}; South Karelia \texttt{0,39\ (+)} dan
\texttt{0,74\ (+)}; Etelä-Savo \texttt{3,46*\ (+)} dan
\texttt{0,78\ (+)}; Northern Savo \texttt{0,04\ (+)} dan
\texttt{0,02\ (+)}; Northern Karelia \texttt{10,27***\ (+)} dan
\texttt{3,08*\ (+)}; Central Finland \texttt{2,99*\ (+)} dan
\texttt{0,49\ (+)}; Southern Ostrobothnia \texttt{0,74\ (+)} dan
\texttt{0,18\ (+)}; Ostrobothnia \texttt{0,01\ (-)} dan
\texttt{0,00\ (-)}; Central Ostrobothnia \texttt{1,81\ (-)} dan
\texttt{0,67\ (-)}; Northern Ostrobothnia \texttt{0,35\ (+)} dan
\texttt{0,38\ (+)}; Kainuu \texttt{34,60***\ (+)} dan
\texttt{16,93***\ (+)}; serta Lapland \texttt{16,65***\ (+)} dan
\texttt{32,59***\ (+)}. {[}Locator: TXT baris 510-541; PDF hlm. 484.{]}

Secara khusus, penulis menyoroti Joensuu, Kajaani, dan Rovaniemi sebagai
pusat di wilayah utara dan timur yang masih menunjukkan hubungan positif
kuat dari pusat ke periphery. Ketiganya berada di wilayah yang disebut
sebagai yang paling kurang berkembang dan paling miskin di Finlandia.
Penulis juga membandingkan hasil ini dengan Tervo (2009): pada perubahan
tahunan 1970-2004, pusat besar yang tumbuh cepat dilaporkan berdampak
negatif pada hinterland, sedangkan pusat yang tumbuh lambat atau lemah
berdampak positif. {[}Locator: TXT baris 545-574; PDF hlm. 484.{]}

\textbf{Catatan atas notasi:} Teks ekstraksi menampilkan karakter alpha
sebagai \texttt{â} pada beberapa baris persamaan dan tabel. Review ini
menuliskannya sebagai \texttt{alpha} berdasarkan konteks persamaan dan
penjelasan artikel, tanpa mengubah angka hasil yang dilaporkan.
{[}Locator: TXT baris 302-321, 379-401.{]}

\subsection{Findings}\label{findings-66}

\textbf{Temuan yang dilaporkan sumber:} Persistensi pola lokasi
meningkat setelah industrialisasi dan urbanisasi. Hal itu didukung oleh
penurunan statistik perubahan peringkat dan peningkatan kecocokan
distribusi rank-size. Struktur ukuran subregion Finlandia modern lebih
dekat dengan distribusi Pareto daripada struktur awal periode observasi,
dan ukuran \texttt{alpha} yang meningkat menunjukkan distribusi populasi
menjadi kurang merata. {[}Locator: TXT baris 357-370, 582-589; PDF hlm.
481, 484-485.{]}

\textbf{Interpretasi penulis:} Pada masa pra-perang, ketika produksi
primer masih dominan dan migrasi internal belum luas, pusat dan
hinterland tumbuh bersama secara homogen dan positif. Pada masa
pascaperang, industrialisasi, urbanisasi, dan migrasi dari desa menuju
kota menghasilkan pembangunan regional yang lebih tidak seimbang.
Hubungan pusat-hinterland tetap ada, tetapi menjadi heterogen; pola ini
dibaca sebagai bukti bahwa bayang-bayang aglomerasi mulai muncul.
{[}Locator: TXT baris 590-599; PDF hlm. 485.{]}

\textbf{Inferensi review:} Hasilnya lebih tepat dibaca sebagai
bayang-bayang aglomerasi yang bersifat bersyarat dan berbeda
antarwilayah, bukan sebagai efek negatif yang seragam di semua
hinterland. Tabel HENC memuat hubungan positif maupun negatif dan banyak
nilai yang tidak signifikan. Bahkan, pada periode pascaperang rata-rata
populasi periphery dalam Table 1 tetap lebih tinggi pada 2004 daripada
1880, meskipun turun dari puncaknya pada 1960; ini menunjukkan dominasi
pusat tidak identik dengan pengosongan total seluruh hinterland.
{[}Locator: TXT baris 331-351, 510-541, 568-574.{]}

\subsection{Conclusions}\label{conclusions-66}

\textbf{Laporan sumber:} Penulis menyimpulkan bahwa dua prediksi NEG
yang diuji mendapat dukungan. Distribusi rank-size lebih cocok untuk
periode setelah perang, hukum Zipf menjadi valid secara bertahap selama
industrialisasi dan urbanisasi, serta perubahan peringkat subregion
menunjukkan struktur regional makin stabil. {[}Locator: TXT baris
576-589; PDF hlm. 484-485.{]}

Untuk hubungan spasial, hasil pengujian menunjukkan hubungan yang
homogen dan positif pada 1880-1940, yang dipahami sebagai pertumbuhan
pusat dan hinterland secara bersama. Pada 1950-2004, hipotesis
non-kausalitas homogen dan kausalitas homogen ditolak; hubungan tetap
ada tetapi bersifat heterogen. Penulis mengaitkan heterogenitas ini
dengan bayang-bayang aglomerasi dan menyatakan bahwa kekuatan pasar
lebih kuat daripada upaya intervensi pemerintah untuk mengurangi
pertumbuhan regional yang tidak seimbang. {[}Locator: TXT baris 590-599;
PDF hlm. 485.{]}

\textbf{Inferensi review:} Kesimpulan artikel berlaku pada pola populasi
19 region dan 74 subregion Finlandia yang direkonstruksi ke batas
wilayah 2005. Ia tidak seharusnya diperluas menjadi bukti bahwa seluruh
mekanisme NEG, seluruh kota, atau seluruh negara akan menunjukkan pola
yang sama. Selain itu, istilah kausalitas perlu dibaca dalam arti
Granger yang didefinisikan artikel, yaitu hubungan prediktif dalam
panel.

\subsection{Contributions}\label{contributions-66}

\textbf{Laporan sumber:} Penulis menyebut artikel ini sebagai salah satu
studi pertama yang menganalisis pertumbuhan regional jangka panjang
suatu negara dari perspektif ini. {[}Locator: TXT baris 578-581; PDF
hlm. 484.{]}

Kontribusi empirisnya adalah menyatukan rentang historis 1880-2004 untuk
menilai perubahan struktur regional Finlandia lintas fase
industrialisasi, bukan hanya memeriksa beberapa dekade terakhir.
Kontribusi substantifnya adalah menghubungkan persistensi hierarki
wilayah dengan perubahan hubungan pusat-hinterland sebelum dan sesudah
perang. {[}Locator: TXT baris 93-104, 118-123, 578-599.{]}

Kontribusi metodologis yang dinyatakan adalah penerapan prosedur panel
Granger yang memungkinkan heterogenitas antardaerah, termasuk rangkaian
HINC, HC, dan HENC. Artikel juga menggabungkan statistik perubahan
peringkat, estimasi rank-size yang dikoreksi bias, dan pengujian panel
untuk mendekati dua prediksi NEG dari sisi yang berbeda. {[}Locator: TXT
baris 300-321, 416-421, 547-567.{]}

\textbf{Inferensi review:} Nilai tambah utama artikel adalah menunjukkan
bahwa ``dominasi pusat'' tidak cukup disimpulkan dari membesarnya kota;
arah dan homogenitas hubungan dengan hinterland perlu diperiksa. Namun,
klaim bahwa prosedur tersebut benar-benar novel dibandingkan seluruh
literatur tidak diverifikasi secara terpisah karena review ini hanya
memakai TXT B23.

\subsection{Research Gap}\label{research-gap-66}

\textbf{Kesenjangan yang dilaporkan sumber:} Literatur empiris tentang
persistensi distribusi aktivitas telah berkembang, tetapi peran pusat
pertumbuhan masih kurang diteliti. Untuk Finlandia, studi regional
jangka panjang dari perspektif ekonomi spasial disebut terbatas.
Ottaviano dan Pinelli membahas kekuatan yang membentuk lanskap ekonomi
Finlandia, sedangkan Tervo (2009) menganalisis hubungan pusat-hinterland
hanya pada 1970-2004. {[}Locator: TXT baris 105-117; PDF hlm. 477.{]}

\textbf{Kesenjangan yang dijawab:} Artikel memperluas horizon menjadi
1880-2004, membandingkan fase pra-perang dan pascaperang, serta menguji
perubahan baik pada persistensi regional maupun hubungan
pusat-hinterland. {[}Locator: TXT baris 93-104, 578-581.{]}

\textbf{Kesenjangan tersisa, sebagai inferensi review:} Karena dataset
hanya berisi populasi agregat, artikel belum mengidentifikasi mekanisme
melalui sektor industri, pekerjaan, pendapatan, migrasi, komuter,
perusahaan, atau arus investasi. Penggunaan batas wilayah 2005 juga
menyisakan pertanyaan tentang bagaimana hubungan fungsional yang berubah
dari waktu ke waktu memengaruhi hasil. Heterogenitas pascaperang
ditemukan, tetapi penyebab perbedaan antarregion tidak diestimasi secara
eksplisit.

\textbf{Informasi yang tidak tersedia:} TXT tidak memiliki bagian yang
secara khusus diberi label ``Research Gap'' dan tidak merumuskan daftar
kesenjangan penelitian yang belum terjawab setelah kesimpulan.

\subsection{Limitations}\label{limitations-66}

\textbf{Keterbatasan yang dinyatakan atau diakui sumber:}

\begin{itemize}
\tightlist
\item
  Data populasi dipilih karena merupakan satu-satunya kemungkinan yang
  dianggap cukup baik untuk analisis regional jangka panjang, sehingga
  jenis bukti yang tersedia memang terbatas pada perkembangan penduduk.
  {[}Locator: TXT baris 93-95.{]}
\item
  Perang Dunia II dipakai sebagai batas pra-perang dan pascaperang hanya
  sebagai pendekatan kasar; waktu tepat peralihan ekonomi sulit
  ditentukan. {[}Locator: TXT baris 93-100, 206-212.{]}
\item
  Wilayah historis dipetakan ke pembagian NUTS 2005. Munisipalitas,
  subregion, dan wilayah fungsional tidak memiliki ukuran atau batas
  yang identik sepanjang 124 tahun. {[}Locator: TXT baris 213-238.{]}
\item
  Åland dikeluarkan karena otonomi politik dan karakteristiknya berbeda
  dari Finlandia kontinental. Wilayah yang diserahkan setelah perang
  juga tidak muncul dalam data; penulis mengakui hasil pra-perang
  sedikit terdistorsi. {[}Locator: TXT baris 227-242, 355-356.{]}
\item
  Pemukiman pengungsi setelah perang meningkatkan populasi banyak
  munisipalitas secara mendadak. Pemisahan dua periode mengurangi,
  tetapi tidak dinyatakan menghapus, persoalan ini. {[}Locator: TXT
  baris 239-242.{]}
\item
  Definisi pusat berdasarkan subregion terbesar pada 2004 tidak selalu
  cocok dengan subregion terbesar pada semua tahun; empat dari 19 pusat
  pernah mengalami perubahan posisi tersebut. Beberapa region juga
  mempunyai lebih dari satu subregion besar sehingga klasifikasi
  pusat-periphery tidak sepenuhnya jelas. {[}Locator: TXT baris 243-248,
  280-283.{]}
\item
  Uji Granger dengan lag lebih panjang dari \texttt{t-1} tidak dapat
  dilakukan karena jumlah periode sedikit. {[}Locator: TXT baris
  440-445.{]}
\end{itemize}

\textbf{Keterbatasan yang diturunkan review:}

\begin{itemize}
\tightlist
\item
  Perbedaan sumber historis antara register gereja dan sensus penduduk
  berpotensi menjadi isu keterbandingan, tetapi TXT tidak melaporkan
  pengujian khusus untuk besarnya dampak perbedaan sumber tersebut.
  {[}Locator: TXT baris 208-212.{]}
\item
  Data agregat regional membatasi inferensi pada hubungan ekologis
  antarwilayah; data tersebut tidak cukup untuk menyimpulkan perilaku
  rumah tangga, perusahaan, atau individu.
\item
  Uji Granger menunjukkan urutan informasi dan kemampuan prediksi. Tanpa
  identifikasi tambahan, hasilnya tidak secara otomatis memisahkan
  pengaruh pusat dari pengaruh bersama seperti kebijakan, perubahan
  transportasi, atau guncangan ekonomi.
\item
  Karena hanya satu negara yang dianalisis dan pembagian wilayah
  kontemporer dipakai untuk sejarah panjang, generalisasi ke negara lain
  atau ke batas fungsional lain tidak dapat diasumsikan.
\item
  TXT tidak melaporkan robustness check alternatif untuk definisi pusat,
  batas wilayah, pembagian periode, atau bentuk hubungan nonlinier.
\end{itemize}

\subsection{Challenges}\label{challenges-66}

\textbf{Laporan sumber:} Tantangan pengolahan data terutama berasal dari
perubahan jumlah dan bentuk munisipalitas, termasuk kasus pemekaran atau
pembagian munisipalitas. Penulis harus mengalokasikan populasi historis
ke unit penerima mayoritas agar data dapat diagregasikan ke subregion
yang konsisten. {[}Locator: TXT baris 231-238.{]}

Tantangan historis lainnya adalah menggabungkan data dengan batas
wilayah masa kini ketika Finlandia kehilangan wilayah setelah perang dan
menerima perpindahan penduduk dalam jumlah besar. Pemilihan Perang Dunia
II sebagai titik pemisah juga harus menangani fakta bahwa transisi
menuju ekonomi industri tidak mempunyai tanggal tunggal yang jelas.
{[}Locator: TXT baris 93-100, 178-189, 239-242.{]}

Pada sisi estimasi, penulis menghadapi bias OLS pada sampel kecil untuk
hukum rank-size, keterbatasan sifat finite-sample estimator Hill, serta
kebutuhan memakai koreksi \texttt{rank\ -\ 1/2}. Untuk panel Granger,
penulis perlu mengakomodasi kemungkinan heterogenitas, memastikan
transformasi menuju stationarity, memilih lag yang masih mungkin dengan
sedikit periode, dan menguji beberapa hipotesis dalam prosedur
bersarang. {[}Locator: TXT baris 312-321, 416-449, 547-567.{]}

\textbf{Inferensi review:} Tantangan-tantangan tersebut bukan hanya
persoalan teknis. Rekonstruksi ke batas 2005 menentukan cara pusat dan
hinterland dibaca, sedangkan interval satu dekade mengaburkan urutan
perubahan yang lebih pendek. Karena itu, hasil sensitif terhadap unit
spasial dan resolusi waktu yang dipilih, meskipun sensitivitas tersebut
tidak diuji secara terpisah dalam TXT.

\textbf{Informasi yang tidak tersedia:} Artikel tidak memberi catatan
rinci tentang waktu kerja, prosedur audit data, atau kesulitan
administratif dalam memperoleh data Statistics Finland.

\subsection{Practical Implications}\label{practical-implications-66}

\textbf{Laporan sumber:} Artikel menempatkan hasilnya dalam sejarah
kebijakan regional Finlandia. Kebijakan infrastruktur pada 1950-an dan
1960-an menjadi dasar pembangunan wilayah berpenduduk jarang; kebijakan
awal juga berupaya mendesentralisasikan manufaktur; kebijakan
kesejahteraan dan regionalisasi pendidikan universitas disebut sebagai
orientasi penting. Penulis menyatakan bahwa kebijakan regional tidak
sepenuhnya melawan kekuatan pasar yang mengarahkan pembangunan menuju
struktur terpusat. {[}Locator: TXT baris 190-204; PDF hlm. 478.{]}

Dalam kesimpulan, penulis menyatakan bahwa intervensi pemerintah
berusaha mengurangi pertumbuhan regional yang tidak seimbang, tetapi
kekuatan pasar lebih kuat. {[}Locator: TXT baris 596-599; PDF hlm.
485.{]}

\textbf{Inferensi review:} Implikasi praktis yang paling hati-hati
adalah bahwa pertumbuhan pusat tidak boleh diasumsikan otomatis
menghasilkan \emph{spread effect} yang sama bagi semua hinterland.
Hubungan pascaperang berbeda antarregion: Kainuu, Lapland, dan Northern
Karelia menunjukkan hubungan positif kuat dari pusat ke hinterland,
sementara beberapa hubungan lain negatif atau tidak signifikan.
Kebijakan berbasis pusat pertumbuhan karena itu perlu memperhatikan
kondisi regional dan arah hubungan yang benar-benar teramati, bukan
hanya ukuran kota.

Implikasi lain adalah perlunya memantau pusat dan hinterland secara
bersamaan dalam evaluasi pembangunan regional. Kenaikan pangsa penduduk
pusat dapat berjalan bersamaan dengan pertumbuhan sebagian hinterland,
sehingga indikator konsentrasi nasional saja tidak cukup untuk menilai
apakah manfaat pembangunan menyebar secara lokal.

\textbf{Batas implikasi:} TXT tidak menyajikan rekomendasi kebijakan
operasional, estimasi biaya-manfaat, atau evaluasi kausal atas empat
orientasi kebijakan tersebut. Pernyataan pada dua paragraf terakhir
adalah inferensi review dari hasil artikel, bukan rekomendasi eksplisit
penulis.

\subsection{Applications}\label{applications-66}

\textbf{Aplikasi yang tersirat dalam sumber:} Kerangka dan prosedur
artikel dapat digunakan untuk membaca perubahan hierarki kota,
membandingkan stabilitas peringkat subregion, dan menilai hubungan
pertumbuhan antara pusat administratif dengan hinterland dalam sejarah
regional. Artikel juga menyediakan cara membedakan hubungan yang homogen
pada seluruh wilayah dari hubungan yang hanya muncul pada sebagian
wilayah melalui HINC, HC, dan HENC. {[}Locator: TXT baris 251-264,
416-421, 547-567.{]}

\textbf{Inferensi review:} Dalam konteks kajian regional, pendekatan ini
dapat diterapkan untuk evaluasi historis konsentrasi penduduk,
pemeriksaan apakah pusat pertumbuhan berkaitan dengan penguatan atau
pelemahan hinterland, serta perbandingan antarperiode sebelum dan
sesudah perubahan struktural. Aplikasi tersebut tetap memerlukan data
dan definisi unit spasial yang sebanding; artikel sendiri hanya menguji
Finlandia dan tidak melaporkan penerapan pada negara atau kebijakan
lain.

\textbf{Informasi yang tidak tersedia:} TXT tidak memberikan studi
penerapan lanjutan, panduan replikasi, kode, atau contoh penggunaan
hasil oleh lembaga perencanaan.

\subsection{Future Research
(Optional)}\label{future-research-optional-66}

\textbf{Laporan sumber:} Tidak ada bagian khusus atau agenda eksplisit
\emph{future research} dalam TXT. Penulis menutup artikel dengan
kesimpulan, ucapan terima kasih, lampiran, dan referensi. {[}Locator:
TXT baris 576-705; PDF hlm. 484-486.{]}

\textbf{Inferensi review, bukan agenda penulis:}

\begin{itemize}
\tightlist
\item
  Menguji definisi pusat yang berubah menurut waktu dan menggunakan
  wilayah fungsional historis, bukan hanya subregion terbesar pada 2004.
\item
  Menggabungkan populasi dengan data migrasi, pekerjaan, sektor
  industri, pendapatan, perusahaan, komuter, atau output agar mekanisme
  bayang-bayang aglomerasi dapat dibedakan dari sekadar perubahan ukuran
  penduduk.
\item
  Menguji sensitivitas terhadap batas periode, pengaruh pemukiman
  pengungsi, wilayah yang diserahkan, dan perbedaan register
  gereja-sensus.
\item
  Mengembangkan model yang dapat mengakomodasi lag lebih panjang,
  heterogenitas, dan kemungkinan hubungan nonlinier dengan rentang waktu
  yang lebih kaya.
\item
  Membandingkan hasil Finlandia dengan negara lain atau wilayah yang
  mengalami transisi industrialisasi berbeda, sambil mempertahankan
  kehati-hatian terhadap perbedaan institusi dan batas administratif.
\end{itemize}

\subsection{Provenance Notes}\label{provenance-notes-66}

\begin{itemize}
\tightlist
\item
  Sumber yang digunakan hanya file
  \texttt{/Users/mac/Documents/Mac/{[}2{]}\ Obsidian\ Vault/zettelkasten/source/journal/B23-tervo-2010-cities-hinterlands-and-agglomeration-shadows-spatial-developments-in-finland-during-1880-2004-10.1016-j.eeh.2010.05.002.txt}.
\item
  TXT memuat metadata PDF pada baris 1-20: sumber PDF B23, tanggal
  ekstraksi 2026-08-31, metode \texttt{pdftotext\ -layout}, dan jumlah
  11 halaman. Teks artikel dimulai pada baris 22 dan berakhir pada baris
  705.
\item
  Identitas bibliografis diambil dari judul, nama penulis, jurnal,
  volume, halaman, ISSN, DOI, dan riwayat artikel yang tercantum pada
  baris 22-56 dan 88-90. Nomor isu, tanggal publikasi yang lebih
  spesifik, status peer-review, serta metadata bibliografis lain yang
  tidak tercantum dinyatakan tidak tersedia.
\item
  Locator numerik dalam review merujuk pada baris TXT yang dibaca;
  locator halaman PDF mengikuti nomor halaman tercetak yang muncul di
  dalam ekstraksi, bukan nomor halaman penampil teks.
\item
  Laporan statistik dan isi Table 1-4, lampiran, serta kualifikasi data
  diringkas dari TXT. Referensi yang tercantum pada TXT hanya dipakai
  untuk menjelaskan literatur yang dilaporkan artikel; karya-karya
  tersebut tidak dibaca atau diverifikasi secara independen.
\item
  Pemisahan antara laporan sumber, interpretasi penulis, dan inferensi
  review dibuat untuk mencegah pembacaan inferensi metodologis sebagai
  temuan empiris penulis. Khususnya, penilaian bahwa Granger causality
  bersifat prediktif dan bahwa hasil tidak membuktikan mekanisme
  struktural adalah inferensi review yang didasarkan pada definisi model
  di TXT.
\item
  Ekstraksi \texttt{pdftotext\ -layout} mengandung artefak karakter pada
  persamaan dan tabel, terutama simbol \texttt{alpha} yang muncul
  sebagai \texttt{â}, serta pemisahan baris dan nomor halaman. Angka,
  tanda signifikansi, dan tanda arah efek dipertahankan sesuai teks yang
  tersedia; simbol dijelaskan ulang hanya jika konteksnya eksplisit.
\item
  Review ini tidak menambahkan data dari luar TXT, tidak mengubah file
  sumber, dan tidak mengubah file lain di vault.
\end{itemize}

\section{Review B24}\label{review-b24}

\textbf{Identitas sumber}

\begin{itemize}
\tightlist
\item
  Kode kajian: B24.
\item
  Judul: \emph{Identifying Agglomeration Shadows: Long-Run Evidence from
  Ancient Ports}.
\item
  Penulis: Richard Hornbeck, Guy Michaels, dan Ferdinand Rauch.
\item
  Seri dan waktu: CESifo Working Paper No.~11188, June 2024. TXT tidak
  memverifikasi bahwa naskah ini telah melalui peer review; karena itu
  statusnya diperlakukan sebagai working paper, bukan publikasi jurnal
  final.
\item
  Identifier dan akses yang tercantum: SSRN abstract 4889236, dengan URL
  \texttt{https://ssrn.com/abstract=4889236}.
\item
  Bahan yang direview: file TXT B24 yang memuat metadata PDF, teks
  utama, tabel, gambar, online appendix, dan appendix references.
  Metadata menyebut PDF 69 halaman, diekstrak pada 2026-08-31 dengan
  \texttt{pdftotext\ -layout} (TXT baris 1-30).
\item
  Konvensi locator: rujukan utama menggunakan nomor baris TXT dan, jika
  tersedia, bagian, tabel, atau gambar dalam naskah. Label
  \texttt{{[}Laporan\ sumber{]}} berarti isi yang dilaporkan naskah;
  \texttt{{[}Interpretasi\ penulis{]}} berarti makna yang diberikan oleh
  penulis; \texttt{{[}Inferensi\ review{]}} berarti penilaian yang
  diturunkan dalam review ini.
\end{itemize}

\subsection{Summarized Abstract}\label{summarized-abstract-67}

{[}Laporan sumber{]} Naskah mengkaji \emph{agglomeration shadows}, yaitu
pengaruh ketika kota besar mendorong konsentrasi kegiatan di lokasinya
tetapi menghambat sebagian kegiatan ekonomi di area sekitarnya. Penulis
menekankan dua hambatan identifikasi: pembentukan kota bersifat endogen
dan penggabungan banyak jaringan kota dapat menghasilkan \emph{wave
interference} yang menutupi pola spasial. Untuk memperoleh titik awal
yang lebih terarah, naskah menggunakan 4.263 lokasi pelabuhan kuno di
sekitar Mediterania, Laut Hitam, Laut Merah, dan wilayah pesisir terkait
sebagai lokasi yang men-seed kota modern. Hasil utama yang dilaporkan
adalah munculnya bayangan aglomerasi yang terdeteksi di sekitar kota
besar yang berakar dari pelabuhan kuno; pola tersebut juga diperluas ke
lokasi kota modern secara umum. Penulis mengaitkannya dengan konsekuensi
bahwa pertumbuhan di suatu tempat dapat menghambat pertumbuhan tempat
yang berdekatan (TXT baris 78-87, Abstract).

{[}Interpretasi penulis{]} Pertumbuhan kota dapat menaikkan kepadatan
penduduk di sekitar pusat karena ekspansi spasial, tetapi keuntungan
ekonomi yang lebih luas dapat tertahan ketika area sekitar menjadi
tingkat kedua dalam hierarki perkotaan. Dengan perbedaan produktivitas,
upah, dan biaya perumahan antara kawasan urban dan rural, mekanisme ini
berpotensi memperbesar ketimpangan spasial (TXT baris 117-123). Temuan
empiris yang dibangun kemudian dibaca sebagai bukti path dependence
jangka panjang, bukan hanya di titik pelabuhan/kota awal tetapi juga
dalam hubungan spasial dengan area sekitarnya (TXT baris 178-184,
228-250).

{[}Inferensi review{]} Abstrak menyampaikan klaim kausal yang kuat,
tetapi klaim tersebut bergantung pada asumsi bahwa lokasi pelabuhan kuno
tidak berkaitan dengan pembentukan kota modern di area sekitar melalui
jalur lain selain pertumbuhan di lokasi pelabuhan itu sendiri. Jadi,
kesimpulan paling aman dari TXT adalah adanya pola yang konsisten dengan
bayangan aglomerasi dalam sampel pesisir dan di bawah strategi
identifikasi tersebut, bukan bukti bahwa setiap kota besar selalu
menimbulkan efek yang sama.

\subsection{Problem Statement}\label{problem-statement-67}

{[}Laporan sumber{]} Masalah substantifnya adalah ketegangan antara dua
kekuatan. Kota memperoleh \emph{agglomeration spillovers} yang
meningkatkan produktivitas ketika kegiatan ekonomi berdekatan, tetapi
konsentrasi itu dapat membuat pusat pesaing di sekitar tidak layak atau
hanya mampu menyediakan kegiatan tingkat lebih rendah. Dalam hierarki
urban, kota besar ditempatkan lebih berjauhan dan kota kecil mengisi
ruang antara kota-kota tersebut; model New Economic Geography memberikan
dasar mikro bagi proses tersebut (TXT baris 124-133, 291-305).

{[}Laporan sumber{]} Masalah empirisnya terdiri dari dua bagian.
Pertama, kota tumbuh secara endogen di tempat yang sekaligus
mencerminkan dan memengaruhi potensi pertumbuhan kota sekitar, sehingga
tidak ada titik awal pusat yang secara umum eksogen. Kedua, ketika
distribusi kota dari banyak jaringan yang sedikit berbeda
dirata-ratakan, gelombang lokasi kota menjadi tidak sinkron. Puncak dan
lembah jarak antar-kota dapat saling meniadakan, terutama bagi kota
kecil yang jaraknya lebih rapat, sehingga bayangan yang sebenarnya ada
dapat tampak sebagai penurunan monoton atau hilang (TXT baris 134-143,
431-437).

{[}Interpretasi penulis{]} Karena itu, mendeteksi hubungan negatif di
sekitar kota tidak cukup dilakukan dengan sekadar menghubungkan
kepadatan dengan jarak ke kota yang sudah ada. Diperlukan lokasi awal
yang memiliki alasan historis dan geografis yang cukup lokal untuk
men-seed kota, serta cara membedakan bayangan kota dari pengaruh
langsung pelabuhan atau ketidaksesuaian geografis area sekitar (TXT
baris 144-152, 451-466).

{[}Inferensi review{]} Formulasi masalah naskah kuat karena
menghubungkan teori pembentukan hierarki dengan persoalan observasional
yang spesifik. Namun, strategi tersebut memindahkan beban identifikasi
ke dua hal yang tidak dapat diamati sepenuhnya dari hasil akhir: apakah
pelabuhan kuno hanya memengaruhi area sekitar melalui kota yang tumbuh
di titik itu, dan apakah hilang/bertahannya pelabuhan alami benar-benar
terutama mengubah ketepatan lokasi kota yang di-seed. TXT melaporkan uji
kontrol dan robustness, tetapi tidak menghilangkan kebutuhan untuk
menilai asumsi ini.

\subsection{Objectives}\label{objectives-67}

{[}Laporan sumber{]} Tujuan operasional yang dapat diidentifikasi dari
naskah adalah:

\begin{itemize}
\tightlist
\item
  Menggunakan simulasi model Fujita, Krugman, dan Mori (1999), disebut
  FKM, untuk menunjukkan bagaimana bayangan aglomerasi dan \emph{wave
  interference} terlihat ketika dampak beberapa titik awal
  dirata-ratakan (TXT baris 306-429).
\item
  Menggunakan lokasi pelabuhan kuno sebagai titik awal historis untuk
  mengestimasi pengaruh jangka panjangnya terhadap kepadatan penduduk,
  aktivitas urban, dan aktivitas kota di area sekitarnya (TXT baris
  438-466, 502-559).
\item
  Memakai keberadaan atau hilangnya pelabuhan alami modern untuk
  membandingkan kasus ketika lokasi kota yang di-seed lebih tepat
  terikat pada lokasi pelabuhan dengan kasus ketika lokasi awal lebih
  berisik secara spasial (TXT baris 160-164, 396-420).
\item
  Menguji apakah pola penurunan dan kenaikan probabilitas kepadatan kota
  pada jarak menengah sesuai dengan bayangan aglomerasi, bukan hanya
  dengan penurunan kepadatan monoton (TXT baris 675-715).
\item
  Membedakan persaingan antarkota dari persaingan langsung antarstruktur
  pelabuhan melalui hubungan antara kehilangan pelabuhan alami di satu
  lokasi dan kondisi lokasi pelabuhan di sekitarnya (TXT baris 784-821).
\item
  Membandingkan jarak kota modern yang terealisasi dengan distribusi
  jarak dari lokasi kota acak yang diberi bobot berdasarkan
  karakteristik geografis, untuk melihat apakah bayangan juga tampak
  pada hierarki urban modern secara umum (TXT baris 822-868).
\end{itemize}

{[}Inferensi review{]} Tujuan tersebut membentuk desain berlapis:
teori/simulasi untuk memprediksi bentuk pola, desain pelabuhan kuno
untuk identifikasi historis, dan uji spacing modern sebagai pemeriksaan
eksternal. Uji-uji ini saling melengkapi, tetapi tidak semuanya menjawab
pertanyaan kausal yang sama; analisis spacing modern terutama bersifat
perbandingan cross-section dengan null simulasi, sedangkan analisis
pelabuhan berusaha menelusuri dampak dari titik awal historis.

\subsection{Literature Survey}\label{literature-survey-67}

{[}Laporan sumber{]} Naskah menempatkan pertanyaan dalam beberapa rumpun
literatur. Literatur hierarki kota dimulai dari gagasan kota primat dan
\emph{central places} yang tersusun dalam hierarki, dengan kota besar
berjauhan dan kota lebih kecil mengisi area antara pusat utama.
Literatur New Economic Geography kemudian memberi mekanisme berupa
increasing returns, biaya transportasi, mobilitas tenaga kerja, dan
kekuatan dispersi yang bersama-sama menghasilkan kota dengan ukuran
serta keragaman kegiatan berbeda (TXT baris 265-305).

{[}Laporan sumber{]} Literatur \emph{agglomeration spillovers}
mendokumentasikan konsentrasi industri yang melebihi apa yang dijelaskan
oleh keunggulan alam atau kebetulan, serta kenaikan produktivitas
perusahaan sekitar setelah pembukaan pabrik besar. Naskah membandingkan
fokus bayangan negatifnya dengan literatur spillover positif yang
menggunakan pembukaan establishment sebagai sumber variasi (TXT baris
267-276, 228-236). Literatur kota monocentric juga digunakan sebagai
pembanding: kepadatan dan harga lahan biasanya menurun dari pusat,
sedangkan model polycentric menjelaskan penyimpangan dari pola itu (TXT
baris 282-290).

{[}Laporan sumber{]} Untuk bukti empiris bayangan, naskah merujuk studi
yang mengaitkan jarak dari kota dengan dinamika populasi serta studi
pembentukan kota Eropa pada 800-1800 yang menemukan lebih sedikit kota
baru dalam 20 km dari kota yang ada dan lebih banyak pada jarak 20-100
km. Literatur konektivitas dan jaringan transportasi dipakai untuk
menunjukkan bahwa hubungan antarlokasi dapat berdampak positif maupun
negatif (TXT baris 430-450). Literatur path dependence, termasuk lokasi
portage dan kota pelabuhan, menjadi dasar argumen bahwa keunggulan
historis dapat mengunci aktivitas di suatu lokasi setelah keunggulan
awalnya melemah (TXT baris 458-487).

{[}Laporan sumber{]} Naskah juga membedakan fokusnya dari literatur
jaringan perdagangan lintas kota: objek utamanya bukan evolusi jaringan
perdagangan, melainkan geografi ekonomi lokal di sekitar pelabuhan kuno
dan kota yang di-seed. Pada analisis spacing, penulis mengacu pada
pendekatan \emph{dartboard} dan analisis pola titik untuk membandingkan
distribusi aktual dengan distribusi acak yang mempertahankan
kecenderungan lokasi berdasarkan karakteristik geografis (TXT baris
479-501, 842-862).

{[}Inferensi review{]} Survei literatur berhasil memberi rantai teori
yang koheren dari spillover, hierarki, dan path dependence menuju
prediksi bayangan berbentuk gelombang. Batasnya, seluruh cakupan
literatur di bagian ini adalah sebagaimana dipaparkan dan dirujuk dalam
TXT; review ini tidak memverifikasi status, kelengkapan, atau
perkembangan setiap referensi. Naskah sendiri menyebut bayangan
aglomerasi relatif kurang diperhatikan dibanding spillover positif,
tetapi pernyataan itu diperlakukan di sini sebagai framing penulis,
bukan hasil audit bibliometrik independen.

\subsection{Method Used}\label{method-used-67}

{[}Laporan sumber{]} Metode teori dan simulasi memakai FKM, sebuah model
ekonomi satu dimensi dengan satu kota awal tetap, sektor manufaktur dan
pertanian, tiga industri manufaktur dengan increasing returns, barang
terdiferensiasi, biaya transportasi iceberg, preferensi
Cobb-Douglas/CES, tenaga kerja yang mobile, dan kualitas lahan homogen.
Ketika populasi agregat meningkat, kota-kota baru muncul secara
bertahap; kota orde lebih tinggi memiliki lebih banyak industri dan
biasanya berpopulasi lebih besar (TXT baris 307-340).

{[}Laporan sumber{]} Penulis mengambil 16 nilai populasi agregat yang
dilaporkan pada Figure 7 FKM dan, pada titik bifurkasi, memilih sisi
dengan populasi lebih tinggi. Mereka menghitung pangsa simulasi yang
memiliki kota pada setiap jarak dari kota awal. Untuk memodelkan
ketidakpastian lokasi kota yang di-seed, lokasi awal digeser dengan
noise normal ber-mean nol. Model disimulasikan 2.000 kali untuk setiap
tingkat populasi; noise rendah, sedang, dan tinggi masing-masing
menggunakan simpangan baku 0,25, 0,5, dan 1,0. Figures 1-3 dan Appendix
Figures A.1-A.3 digunakan untuk melihat bagaimana noise meredam
gelombang kota kecil dan menyisakan pola pada kota besar (TXT baris
341-420, 1270-1340, 1848-1910).

{[}Laporan sumber{]} Metode empiris pertama mengestimasi persamaan (1),
yaitu perbedaan outcome untuk sel grid yang memiliki pelabuhan kuno
dibanding sel lain dengan karakteristik geografis serupa. Outcome
mencakup log kepadatan penduduk, indikator
\texttt{ln\ population\ density\ \textgreater{}\ 6} untuk aktivitas
urban, dan indikator \texttt{ln\ population\ density\ \textgreater{}\ 9}
untuk aktivitas kota. Kontrol baseline meliputi log jarak ke pantai, log
jarak ke sungai besar terdekat, lintang, bujur, ruggedness, serta suhu
dan presipitasi rata-rata Januari dan Juli (TXT baris 602-612).

{[}Laporan sumber{]} Persamaan (2) mengganti indikator pelabuhan tunggal
dengan bin jarak 1 km dari 0 sampai 50 km ke pelabuhan kuno terdekat,
dipisahkan antara lokasi dengan dan tanpa pelabuhan alami modern.
Koefisien kedua kelompok diestimasi bersama karena kedekatan dengan satu
tipe pelabuhan berkorelasi dengan kedekatan ke tipe lain. Kategori
pembanding untuk nilai absolutnya adalah jarak lebih dari 50 km untuk
tiap tipe; fokus interpretasi penulis adalah perbedaan bentuk pola dan
perbandingan antartipe, bukan sekadar level absolut setiap koefisien
(TXT baris 613-633, catatan kaki 10 pada baris 651-653).

{[}Laporan sumber{]} Persamaan (3) membatasi sampel pada sel yang
memiliki pelabuhan kuno. Outcome-nya adalah struktur pelabuhan modern,
log kepadatan penduduk, aktivitas urban, atau aktivitas kota. Variabel
penjelasnya adalah kehilangan pelabuhan alami di sel itu sendiri dan
proporsi sel pelabuhan kuno dalam jarak 5-50 km yang juga kehilangan
pelabuhan alami. Spesifikasi mengontrol jumlah sel pelabuhan lain pada
jarak tersebut dan kontrol geografis baseline (TXT baris 784-809).

{[}Laporan sumber{]} Untuk analisis tambahan, penulis menghitung jarak
setiap kota modern ke kota modern terdekat menggunakan GHSL-UCDB 2019
dan membandingkannya dengan 1.000 simulasi lokasi acak berbobot. Bobot
berasal dari nilai fitted probit yang memprediksi kemungkinan kota
berdasarkan fixed effects negara dan kontrol geografis baseline.
Analisis dipisahkan untuk kota berpenduduk di atas 500.000, 250.000, dan
100.000. Penulis juga menggambar ulang kota dengan penyesuaian area,
menolak lokasi yang masuk lingkaran footprint kota yang telah
ditempatkan (TXT baris 842-876, 897-923, 1446-1475).

{[}Laporan sumber{]} Inferensi statistik baseline mengelompokkan
standard error pada grid sekitar 1/12 derajat x 1/12 derajat, kira-kira
8 km x 8 km, untuk memperhitungkan korelasi spasial akibat satu kota
tercermin pada beberapa sel. Robustness menggunakan two-way shifted
clusters, kelompok 25 km x 25 km, serta Conley standard errors dengan
cutoff 4 km dan 8 km (TXT baris 634-638, 772-783, Figures A.18-A.21 pada
baris 2251-2334). Penulis menyatakan simulasi dipakai untuk menunjukkan
tampilan bayangan, bukan untuk mengestimasi parameter struktural FKM
(TXT baris 421-429).

{[}Inferensi review{]} Metode ini konsisten dengan tujuan untuk
mengenali bentuk pola, khususnya kontras antara penurunan monoton
aktivitas urban dan pola turun-naik aktivitas kota. Akan tetapi, desain
observasional tetap membutuhkan asumsi eksogenitas relatif lokasi
pelabuhan serta interpretasi kehilangan/bertahannya pelabuhan alami
sebagai sumber variasi ketepatan lokasi. Robustness terhadap kontrol dan
standard error memperkuat kestabilan pola yang dilaporkan, tetapi tidak
dengan sendirinya membuktikan bahwa semua jalur historis yang relevan
telah terkontrol.

\subsection{Dataset}\label{dataset-67}

{[}Laporan sumber{]} Dataset utama dibangun dari beberapa sumber. Lokasi
pelabuhan kuno berasal dari database de Graauw (2019), yang memuat
koordinat lokasi yang disebut dalam teks dan sejarah kuno. Penulis
kemudian menggunakan citra Google Earth untuk melakukan hand-code
kondisi modern setiap lokasi, terutama keberadaan pelabuhan alami dan
struktur pelabuhan modern. Pelabuhan alami dinilai berada di sekitar
lokasi, kira-kira dalam jarak 1 km. Dari lokasi utama, 80\% masih
memiliki pelabuhan alami dan 20\% tidak; 46\% memiliki struktur
pelabuhan modern apa pun dan 6\% memiliki struktur komersial (TXT baris
502-529, 1617-1644).

{[}Laporan sumber{]} Outcome kependudukan utama adalah GPWv4 (CIESIN,
2018) untuk tahun 2015. Dataset ini mengalokasikan data penduduk rinci
ke sel melalui areal weighting, juga memakai data dari tahun lain dan
estimasi laju pertumbuhan. Pemeriksaan robustness menggunakan GRUMP
untuk tahun 2000 dan GHS population grid/GHSL-POP untuk tahun 2015.
Untuk analisis kota modern, penulis memakai GHS Urban Centre Database
yang memuat koordinat pusat kota, luas geografis, dan total penduduk
tahun 2015 (TXT baris 536-559, 1645-1665).

{[}Laporan sumber{]} Data geografis mencakup garis pantai, sungai, dan
batas negara dari ESRI; ruggedness dari simpangan baku elevasi; suhu dan
presipitasi Januari serta Juli dari WorldClim 2; kesesuaian enam tanaman
dari FAO-GAEZ; kesesuaian perairan untuk 15 spesies ikan dari AquaMaps;
indikator gurun, elevasi, lokasi kota kuno dari Barrington Atlas/Hanson,
dan lokasi jalan Romawi. Untuk sebagian kecil sel yang kekurangan data
iklim, suhu dan presipitasi diinterpolasi dengan inverse distance
weighting (TXT baris 560-600, 1666-1686).

{[}Laporan sumber{]} Seluruh data digabungkan ke dalam grid equal-area
sel 1 km2 dengan proyeksi Lambert Azimuthal Equal Area dan titik
referensi sekitar 39N, 18.5E. Jarak dihitung dari centroid sel
menggunakan jarak geodesik. Lokasi pelabuhan dan kota Barrington
dicocokkan ke centroid grid terdekat, sampai sekitar 1,5 km; bila sel
yang memiliki pelabuhan/kota tidak memiliki kepadatan penduduk, nilai
diganti dengan rata-rata \emph{king's neighbors} yang memiliki nilai
(TXT baris 1698-1727).

{[}Inferensi review{]} Dataset memungkinkan kombinasi variasi spasial
rinci, jejak historis, dan pengukuran modern yang berbeda. Namun, TXT
tidak menyediakan data mentah, kode estimasi, atau informasi tentang
akses replikasi. Ketersediaan bahan-bahan tersebut tidak disebutkan
dalam file; karenanya review ini hanya menilai konstruksi dan hasil
sebagaimana dilaporkan, bukan mereplikasi perhitungan.

\subsection{Population / Sample}\label{population-sample-67}

{[}Laporan sumber{]} Unit analisis utama bukan individu, rumah tangga,
atau kota administratif, melainkan sel grid 1 km x 1 km. Sampel
geografis mencakup sel dalam jarak 50 km dari pantai dan dalam jarak 200
km dari pelabuhan kuno terdekat, pada negara modern yang memiliki
setidaknya satu pelabuhan kuno dalam sampel. Wilayah mencakup pesisir
Mediterania dan sekitarnya, Laut Hitam, Laut Merah, pesisir
Eropa/Selatan Eropa, Afrika Utara, Asia Barat, serta wilayah pesisir
terkait yang dimasukkan dalam database (TXT baris 504-508, 536-545,
1698-1753).

{[}Laporan sumber{]} Sampel utama mencakup 2,3 juta km2 secara
deskriptif dan Table 1 melaporkan 2,278,861 observasi sel. Sebanyak
4,263 sel memiliki setidaknya satu pelabuhan kuno menurut keterangan
data utama. Appendix menjelaskan bahwa 4,333 lokasi pelabuhan berada di
sel yang memiliki data kepadatan dan bahwa sel-sel itu menghasilkan
4,263 sel unik dengan pelabuhan; sebagian kecil sel memuat lebih dari
satu pelabuhan, dengan rincian 4,204 sel berisi satu pelabuhan, 51
berisi dua, 6 berisi tiga, 1 berisi empat, dan 1 berisi lima (TXT baris
540-545, 1718-1738).

{[}Laporan sumber{]} Database awal memuat 4,561 pelabuhan. Naskah utama
menyatakan bahwa setelah mengecualikan lokasi lebih dari 50 km dari
pantai dan area remote dengan cakupan kurang lengkap, tersisa 4,263
pelabuhan (TXT baris 504-508). Appendix lebih rinci menyatakan bahwa 134
lokasi dikeluarkan melalui tiga kriteria geografis, lalu 4,333 lokasi
pelabuhan berada di sel dengan data kepadatan dan 4,263 sel unik
dipakai. Dua cara penyebutan ini tidak sepenuhnya konsisten antara
jumlah record pelabuhan dan jumlah sel pelabuhan; TXT tidak memberi
rekonsiliasi eksplisitnya (TXT baris 1708-1725).

{[}Laporan sumber{]} Untuk analisis mekanisme pada persamaan (3), sampel
dibatasi pada sel pelabuhan kuno yang memiliki setidaknya satu sel
pelabuhan lain dalam jarak 5-50 km. Table 3 melaporkan 4,198 observasi;
catatan kaki menyebut 71 sel tanpa pelabuhan lain dalam rentang tersebut
dikeluarkan (TXT baris 827-829, 1556-1587). Untuk pembanding ancient
city di Table 1, jumlah yang berada dalam sampel terbatas adalah 10
Category 1, 101 Category 2, dan 480 Category 3, berbeda dari jumlah
kategori pada cakupan database ancient city yang lebih luas (TXT baris
593-599, 1512-1515).

{[}Inferensi review{]} Sampel tidak merepresentasikan seluruh sistem
kota dunia atau seluruh kawasan yang pernah memiliki pelabuhan kuno.
Pembatasan pesisir, negara, jarak, cakupan database, dan penghapusan
lokasi yang kurang pasti menentukan populasi target substantifnya:
sel-sel pesisir dalam jaringan lokasi pelabuhan kuno yang teramati.
Exclusion ini masuk akal untuk desain yang berfokus pada pelabuhan dan
mengurangi cakupan yang tidak lengkap, tetapi membatasi generalisasi ke
wilayah inland, kawasan tanpa pelabuhan dalam database, dan sistem urban
di luar kawasan tersebut.

\subsection{Dependent Variables}\label{dependent-variables-67}

{[}Laporan sumber{]} Outcome utama dalam persamaan (1) dan (2) adalah:

\begin{itemize}
\tightlist
\item
  \texttt{Ln\ population\ density}, yaitu log natural kepadatan penduduk
  per km2.
\item
  \texttt{Urban\ Density}, indikator bernilai satu jika log kepadatan
  penduduk lebih besar dari 6, kira-kira lebih dari 400 orang per km2.
\item
  \texttt{City\ Density}, indikator bernilai satu jika log kepadatan
  penduduk lebih besar dari 9, kira-kira lebih dari 8.000 orang per km2.
\end{itemize}

Definisi log sebagai log natural disebutkan dalam catatan kaki; outcome
biner merepresentasikan aktivitas urban dan aktivitas kota pada tingkat
kepadatan yang berbeda (TXT baris 546-549, 564-576, 607-612; Table 1,
TXT baris 1483-1520).

{[}Laporan sumber{]} Persamaan (3) memakai empat outcome pada sel
pelabuhan: indikator ada/tidaknya \texttt{Modern\ Port\ Structure}, log
kepadatan penduduk, \texttt{Urban\ Density}, dan \texttt{City\ Density}.
Struktur pelabuhan modern berarti struktur buatan manusia yang terlihat
dari satelit untuk membantu bongkar-muat kapal, termasuk dermaga dasar.
Variabel \texttt{Surrounding\ Share\ with\ No\ Harbor} bukan dependent
variable; variabel itu adalah ukuran paparan sel terhadap proporsi
pelabuhan sekitar yang kehilangan pelabuhan alami (TXT baris 796-809,
1556-1587).

{[}Laporan sumber{]} Analisis spacing modern memakai distribusi jarak
minimum setiap kota ke kota terdekat dan cumulative distribution
function, bukan outcome kepadatan grid pada regresi utama. Analisis
alternatif memakai data GHSL untuk mendefinisikan sel yang berada dalam
radius kota berpenduduk lebih dari 500.000, dengan radius dihitung dari
luas kota (TXT baris 863-868, 2226-2243).

{[}Laporan sumber{]} Robustness memvariasikan ambang log kepadatan dari
5 sampai 10, memakai penduduk 2000 dari GRUMP, kepadatan 2015 dari
GHSL-POP, dan outcome berbasis lokasi kota GHSL. Outcome seperti
produktivitas, upah, harga rumah, kesejahteraan, kualitas layanan, atau
pertumbuhan penduduk antarperiode tidak diestimasi secara langsung dalam
TXT (TXT baris 763-771, 2130-2243).

{[}Inferensi review{]} Pemisahan \texttt{Urban\ Density} dan
\texttt{City\ Density} merupakan inti pengujian: penurunan monoton pada
ambang urban tidak boleh langsung dibaca sebagai bayangan, karena dapat
mencerminkan gradien monocentric biasa; pola turun lalu naik pada ambang
kota lebih dekat dengan prediksi bayangan dalam simulasi. Ketergantungan
pada ambang kepadatan dan perbedaan antara GPW serta GHSL tetap menjadi
pertimbangan pengukuran, meskipun naskah melaporkan hasil robustness.

\subsection{Results}\label{results-67}

{[}Laporan sumber{]} Hasil lokal awal pada Table 1 menunjukkan bahwa,
setelah kontrol geografis baseline, sel dengan pelabuhan kuno memiliki
koefisien 0,589 untuk log kepadatan penduduk, 0,116 untuk probabilitas
\texttt{Urban\ Density}, dan 0,017 untuk probabilitas
\texttt{City\ Density}. Standard error berturut-turut adalah 0,044,
0,008, dan 0,004. Dengan rata-rata outcome 3,290, 0,071, dan 0,002,
penulis merangkum hubungan itu sebagai kepadatan rata-rata sekitar 60\%
lebih tinggi, probabilitas urban 12 poin persentase lebih tinggi, dan
probabilitas kota 1,7 poin persentase lebih tinggi pada sel pelabuhan.
Table 1 memiliki 2,278,861 observasi (TXT baris 641-650, 1483-1520).

{[}Laporan sumber{]} Figure 5 menunjukkan dampak terhadap log kepadatan
penduduk yang menurun secara monoton menurut jarak dari pelabuhan kuno
dan masih positif pada 50 km relatif terhadap sel yang lebih jauh.
Dampaknya lebih besar di sekitar lokasi yang pelabuhan alaminya
bertahan, tetapi tetap substansial di sekitar lokasi yang kehilangan
pelabuhan alami. Ringkasan abstrak menyatakan bahwa dampak pada 50 km
kira-kira seperempat dari dampak di lokasi pelabuhan (TXT baris 670-674,
1373-1392; ringkasan pada baris 178-184).

{[}Laporan sumber{]} Figure 6 menemukan pola yang berbeda untuk dua
ambang. Probabilitas \texttt{Urban\ Density} menurun lebih monoton dan
terus menurun dengan jarak. Untuk \texttt{City\ Density} di sekitar
pelabuhan dengan natural harbor, probabilitasnya tinggi di lokasi
pelabuhan, turun cepat hingga titik terendah sekitar 20 km, lalu naik
moderat sampai sekitar 40 km. Di sekitar lokasi tanpa natural harbor,
ketidakpastian lokasi awal lebih besar sampai sekitar 20 km sebelum
probabilitas kota menurun. Perbedaan antara estimasi dengan harbor dan
tanpa harbor memperlihatkan dip relatif pada city density dan penurunan
urban activity yang lebih mendatar (TXT baris 675-689, 1393-1415).

{[}Laporan sumber{]} Figure 7 menunjukkan bahwa penurunan urban activity
kira-kira mengikuti hubungan logaritmik terhadap jarak, sedangkan pola
city activity lebih menyimpang dari peluruhan logaritmik: penurunan
lebih cepat pada 0-20 km dan berbalik pada 20-40 km. Table 2 merangkum
perbandingan tersebut. Pada perbandingan dengan lokasi tanpa harbor,
perbedaan probabilitas city density adalah -0,0042 pada 20 km dan 0,0012
pada 40 km, dengan selisih 0,0054 dan standard error 0,0013; penulis
menyatakan selisih ini signifikan secara statistik. Terhadap fitted log
relationship, deviasi pada 20 km adalah -0,0025 dan pada 40 km 0,0016,
dengan selisih 0,0040 dan standard error 0,0005. Rata-rata outcome pada
tabel adalah 0,0020 dan observasinya 2,278,861 (TXT baris 690-715,
1423-1445, 1528-1555).

{[}Laporan sumber{]} Hasil robustness dinyatakan tetap memperlihatkan
bayangan ketika penulis menambahkan karakteristik geografis seperti
kesesuaian tanaman, fishing suitability, jarak ke muara, elevasi, gurun,
pulau, dan kedekatan dengan pantai/sungai; mengecualikan pulau kecil,
216 lokasi yang kurang pasti, 22 open-water ports, atau area sekitar 14
ancient cities besar; serta menambahkan fixed effects 2 derajat dan
negara. Pola juga dilaporkan serupa dengan threshold 5-10, data GRUMP
2000, GHSL-POP 2015, outcome lokasi kota GHSL, dan beberapa cara
menangani korelasi spasial. TXT terutama memberi deskripsi gambar
robustness, bukan angka lengkap untuk setiap spesifikasi (TXT baris
721-783, 1944-2121, 2129-2334).

{[}Laporan sumber{]} Pada uji persaingan menggunakan Table 3, kehilangan
harbor alami sendiri berkaitan dengan penurunan probabilitas struktur
pelabuhan modern sebesar 0,383, log kepadatan sebesar 0,269,
probabilitas urban sebesar 0,050, dan probabilitas city sebesar 0,017
dalam nilai koefisien negatif. Sebaliknya, proporsi pelabuhan sekitar
yang kehilangan harbor memiliki koefisien -0,039 untuk struktur
pelabuhan modern, tetapi 0,973 untuk log kepadatan, 0,134 untuk urban,
dan 0,041 untuk city. Standard error masing-masing dilaporkan dalam
tabel. Sampel terdiri dari 4,198 observasi (TXT baris 784-821,
1556-1588).

{[}Laporan sumber{]} Pada Figure 8, terdapat 61 kota modern dengan
populasi lebih dari 500.000. Dibanding distribusi random berbobot,
kota-kota besar yang terealisasi lebih jarang memiliki kota besar
terdekat dalam jarak 40 km dan lebih banyak memiliki tetangga pada 40-60
km. Pola ini lebih lemah ketika ambang kota diturunkan ke 250.000 dan
100.000. Setelah area adjustment untuk mencegah kota simulasi masuk
footprint kota yang sudah ditempatkan, pola masih terlihat untuk kota di
atas 500.000, tetapi tidak lagi terdeteksi untuk ambang 250.000 atau
100.000 (TXT baris 869-906, 1446-1475).

{[}Inferensi review{]} Bukti paling tajam dalam TXT bukan sekadar bahwa
kepadatan dekat pelabuhan lebih tinggi, melainkan bahwa city density di
sekitar pelabuhan dengan lokasi kota yang lebih terikat menunjukkan
bentuk turun-naik pada jarak menengah. Pada saat yang sama, hilangnya
pola untuk kota lebih kecil setelah area adjustment mengurangi kekuatan
klaim bahwa bayangan merupakan sifat umum seluruh skala kota. Hasil
utama juga lebih tepat dibaca sebagai pola probabilitas dan kepadatan
2015 di sekitar titik historis daripada pengamatan langsung atas seluruh
proses pembentukan kota sepanjang 1.500-3.500 tahun.

\subsection{Findings}\label{findings-67}

{[}Interpretasi penulis{]} Temuan substantif utama adalah bahwa kota
yang di-seed di lokasi pelabuhan kuno tampak menghasilkan bayangan
aglomerasi pada pembentukan kota di sekitarnya. Pada jarak 10-30 km,
probabilitas city density lebih rendah daripada pola yang diharapkan
dari penurunan halus, lalu meningkat pada jarak yang lebih jauh,
terutama ketika natural harbor masih bertahan dan lokasi kota lebih
presisi. Penulis menghubungkan rentang 10-30 km itu dengan jarak
perjalanan satu hari bagi hewan pengangkut atau kereta dalam dunia kuno
(TXT baris 185-190, 700-715).

{[}Interpretasi penulis{]} Temuan kedua adalah perbedaan antara jenis
aktivitas. Kepadatan penduduk rata-rata dan aktivitas urban menurun
secara umum ketika menjauh dari pelabuhan, sedangkan aktivitas kota
menunjukkan dip dan pemulihan pada jarak menengah. Penulis membaca
penurunan umum ini sebagai konsisten dengan area sekitar yang menjadi
lokasi tingkat kedua dalam hierarki urban. Pola turun-naik pada ambang
kota diperlakukan sebagai tanda yang lebih khas dari bayangan daripada
gradien monocentric biasa (TXT baris 178-184, 202-209, 947-955).

{[}Interpretasi penulis{]} Temuan ketiga menyangkut mekanisme.
Kehilangan natural harbor menurunkan aktivitas di lokasi pelabuhan itu
sendiri dan mengurangi peluang adanya struktur pelabuhan modern, tetapi
kehilangan harbor di pelabuhan sekitar tidak membuat struktur pelabuhan
di lokasi target lebih mungkin muncul. Sebaliknya, aktivitas penduduk,
urban, dan kota di lokasi target lebih tinggi ketika pelabuhan sekitar
kehilangan harbor. Penulis menyimpulkan bahwa pola lebih mencerminkan
crowd-out antarkota untuk kegiatan ekonomi umum daripada kompetisi
langsung antarport (TXT baris 210-217, 784-821).

{[}Interpretasi penulis{]} Temuan keempat berasal dari cross-section
kota modern. Kota besar di atas 500.000 penduduk memiliki lebih sedikit
tetangga kota besar dalam 40 km dan lebih banyak tetangga pada 40-60 km
dibanding null random berbobot. Penulis menilai hal tersebut konsisten
dengan bayangan di dalam ekuilibrium spasial modern, tetapi efeknya
kurang menonjol untuk ambang populasi yang lebih rendah dan hilang untuk
kota di atas 250.000 serta 100.000 setelah penyesuaian footprint (TXT
baris 883-910).

{[}Inferensi review{]} Secara gabungan, hasil mendukung mekanisme
hierarki yang berskala: pusat besar menarik atau mempertahankan
aktivitas, area dekat menjadi lebih padat secara agregat, tetapi kota
pesaing pada jarak tertentu lebih sulit muncul. Namun, temuan yang
paling diagnostik bergantung pada natural harbor sebagai kondisi yang
mengurangi \emph{wave interference}, dan temuan spacing modern
bergantung pada cara menetapkan batas kota serta null distribution.
Dengan demikian, label bayangan merupakan interpretasi mekanistik yang
didukung oleh bentuk pola dan simulasi, bukan pengukuran langsung atas
niat bersaing atau arus migrasi antarkota.

\subsection{Conclusions}\label{conclusions-67}

{[}Laporan sumber{]} Naskah menyimpulkan bahwa lokasi pelabuhan kuno
yang dibentuk oleh fitur pantai lokal menyediakan pijakan empiris untuk
menelusuri dampak jangka panjang terhadap tempat sekitar. Lokasi
pelabuhan itu sendiri boleh jadi terus memperoleh manfaat dari geografi
lokal, tetapi penulis berpendapat bahwa geografi tersebut secara masuk
akal tidak berkorelasi dengan pola turun-naik karakteristik pada 20 km
dan 40 km. Bertahannya natural harbor membuat lokasi kota modern lebih
tepat terikat pada titik pelabuhan dan membantu mengurangi \emph{wave
interference}; ketika harbor hilang, pola dapat ter-ratakan menjadi
penurunan halus (TXT baris 930-946).

{[}Laporan sumber{]} Kesimpulan berikutnya adalah bahwa kota besar
kurang mungkin terbentuk pada jarak menengah dari lokasi pelabuhan kuno,
tetapi lebih mungkin pada jarak yang lebih jauh, sedangkan populasi dan
aktivitas urban menurun secara umum menurut jarak. Analisis spacing
modern juga menunjukkan jarak kota besar yang berbeda secara statistik
dari random setelah karakteristik lokasi diperhitungkan. Naskah menutup
dengan klaim bahwa investasi awal pada aktivitas lokal dapat
menghasilkan pertumbuhan lokal yang bertahan sekaligus mengatur
organisasi spasial hierarki urban di sekitarnya (TXT baris 947-970).

{[}Inferensi review{]} Kesimpulan paling terukur adalah adanya bukti
yang konsisten dengan bayangan aglomerasi pada outcome kepadatan dan
lokasi kota dalam sampel pesisir yang dipilih. Kata ``causal'' yang
dipakai penulis perlu dibaca bersama identifying assumption,
keterbatasan data antara masa kuno dan 2015, serta sensitivitas analisis
kota kecil terhadap area adjustment. TXT tidak membuktikan dampak pada
produktivitas, upah, harga perumahan, atau kesejahteraan secara
langsung; implikasi tersebut tetap berada pada tingkat konseptual dalam
naskah.

\subsection{Contributions}\label{contributions-67}

{[}Laporan sumber{]} Naskah mengklaim kontribusi berikut:

\begin{itemize}
\tightlist
\item
  Menunjukkan bukti kausal negatif tentang \emph{agglomeration shadows}
  dengan menggunakan pengaruh historis lokasi pelabuhan kuno sebagai
  penentu lokasi awal kota, sebagai pelengkap literatur yang lebih
  banyak mengidentifikasi spillover aglomerasi positif (TXT baris
  228-236).
\item
  Memperluas literatur path dependence dari persistensi pada satu lokasi
  menjadi persistensi dalam hubungan spasial antara lokasi awal dan area
  sekitarnya (TXT baris 234-236, 458-466).
\item
  Menghubungkan pengujian empiris lebih dekat dengan prediksi teoritis
  melalui simulasi FKM, termasuk dokumentasi bahwa \emph{wave
  interference} dapat menyembunyikan bayangan, terutama untuk kota kecil
  (TXT baris 242-250, 396-429).
\item
  Memanfaatkan keberadaan/hilangnya natural harbor untuk membedakan
  lokasi yang lebih tepat mengikat kota awal dari lokasi yang mengandung
  lebih banyak noise spasial (TXT baris 160-164, 369-420).
\item
  Menawarkan uji tambahan berbasis spacing seluruh kota modern dengan
  null random yang mempertahankan kecenderungan geografis pembentukan
  kota (TXT baris 842-862).
\end{itemize}

{[}Inferensi review{]} Nilai tambah desain terutama terletak pada
triangulasi antara prediksi bentuk pola, anchor historis, dan
pemeriksaan modern, bukan pada satu koefisien tunggal. Kontribusi ini
bersifat empiris-metodologis dan tidak sama dengan mengidentifikasi
seluruh mekanisme mikro. Naskah sendiri menyatakan bahwa FKM hanya satu
microstructure yang dapat menghasilkan pola tersebut, bahwa mekanisme
lain mungkin memberi implikasi kesejahteraan berbeda, dan bahwa simulasi
tidak ditujukan untuk estimasi parameter (TXT baris 421-429).

\subsection{Research Gap}\label{research-gap-67}

{[}Laporan sumber{]} Gap yang dikemukakan penulis adalah bahwa bayangan
aglomerasi telah lama diprediksi oleh teori hierarki urban dan New
Economic Geography, tetapi lebih relatif terabaikan dalam fokus empiris
mutakhir pada spillover aglomerasi positif. Pengidentifikasiannya sulit
karena kota terbentuk endogen dan pola wave dapat hilang ketika banyak
jaringan kota dirata-ratakan. Studi empiris terkait juga sering
mengambil lokasi kota sebagai given dan mengestimasi dinamika populasi
menurut jarak, bukan mengidentifikasi bagaimana sebuah kota awal
membentuk hierarki sekitarnya secara jangka panjang (TXT baris 291-305,
431-457, 251-258).

{[}Interpretasi penulis{]} Naskah mengisi gap itu dengan dua langkah:
menggunakan sejarah lokasi pelabuhan sebagai sumber variasi untuk titik
awal kota dan menyesuaikan bentuk uji empiris dengan pola yang
dihasilkan simulasi FKM. Penggunaan ancient ports, bukan semua ancient
cities, dimaksudkan untuk lebih menekankan fitur pantai yang sangat
lokal dan mengurangi hubungan langsung dengan kesesuaian area sekitar;
kontrol terhadap jalan Romawi dan kota kuno besar kemudian dipakai
sebagai pemeriksaan tambahan (TXT baris 451-501, 740-758).

{[}Inferensi review{]} Gap yang terisi terutama adalah gap identifikasi
dan diagnosis bentuk pola. Gap yang belum terisi adalah observasi
langsung atas lintasan kota pada periode antara masa kuno dan 2015,
bukti sektoral tentang kegiatan yang benar-benar terdorong keluar, serta
pengukuran kesejahteraan/distribusi yang dapat menguji konsekuensi
normatif dari bayangan. Tiga hal ini tidak boleh dibaca sebagai klaim
bahwa naskah gagal total; semuanya berada di luar outcome dan cakupan
yang dilaporkan TXT.

\subsection{Limitations}\label{limitations-67}

{[}Batasan yang dinyatakan sumber{]} Cakupan empiris dibatasi pada area
dalam 50 km dari pantai, dalam 200 km dari pelabuhan kuno, dan negara
yang memiliki pelabuhan dalam sampel. Penulis juga mengakui bahwa
analisis tidak memodelkan atau mengestimasi aktivitas yang sangat
tradable pada jarak jauh maupun interaksi antar banyak titik awal; hal
itu dinyatakan berada di luar scope (TXT baris 716-720). Data populasi
dengan resolusi spasial cukup rinci untuk periode antara era kuno dan
masa kini tidak tersedia, sehingga lintasan peralihan tidak dapat
diamati secara komprehensif (TXT baris 550-559).

{[}Batasan yang dinyatakan sumber{]} FKM memakai ekonomi satu dimensi
dan featureless plain, sementara analisis empiris harus mengendalikan
geografi yang heterogen. Simulasi menambahkan noise lokasi karena alasan
analitis dan sebagai analogi ketepatan lokasi kota, tetapi penulis
mengakui bahwa variasi biaya transportasi, ruang kota, parameter model,
dan mekanisme mikro lain juga dapat menghasilkan wave interference.
Simulasi dimaksudkan untuk memperlihatkan bagaimana bayangan tampak
dalam data, bukan untuk mengestimasi parameter model (TXT baris 307-331,
411-429).

{[}Batasan yang dinyatakan sumber{]} Identifikasi memerlukan asumsi
bahwa ancient ports tidak terkait dengan pembentukan kota di area
sekitar melalui jalur lain selain mendorong pertumbuhan kota di lokasi
pelabuhan. Pelabuhan kuno dapat mempunyai dampak langsung pada aktivitas
modern di titiknya sendiri; fokus penelitian adalah area sekitar.
Penulis menanggapi kekhawatiran itu dengan kontrol jarak ke pantai,
sungai, geografi lain, jalan Romawi, kota kuno, fixed effects, dan
exclusion tests, tetapi asumsi dasar tetap bagian dari interpretasi
kausal (TXT baris 144-152, 158-171, 740-758).

{[}Batasan yang dinyatakan sumber{]} Pengukuran lokasi dan outcome
mengandung keputusan klasifikasi. Keberadaan natural harbor dikode dari
citra satelit modern dan tidak memasukkan pengaruh struktur buatan
manusia dalam definisi harbor; 22 open-water ports yang lebih bergantung
pada perlindungan buatan kadang dikeluarkan. Sebanyak 216 lokasi yang
kurang diyakini sebagai pelabuhan juga diuji lewat exclusion. GPWv4
adalah estimasi kepadatan yang dialokasikan ke grid, sedangkan GHSL
bergantung lebih besar pada area/batas kota yang ditentukan. Penulis
sendiri mengidentifikasi risiko bahwa area tepat di luar batas kota
dapat secara mekanis tampak kurang padat (TXT baris 514-529, 763-771,
897-923, 1628-1644).

{[}Batasan yang dinyatakan sumber{]} Analisis spacing modern tidak
mengamati proses pembentukan kota secara historis. Lokasi dan ukuran
kota telah berevolusi bersama lingkungannya, dan null random dapat
terlalu berjauhan bila karakteristik tak teramati yang mendukung kota
juga berkorelasi spasial positif. Setelah area adjustment, bayangan
tidak lagi terdeteksi untuk ambang 250.000 dan 100.000, sehingga
hasilnya sensitif pada cara menangani footprint dan definisi kota (TXT
baris 838-862, 897-910).

{[}Inferensi review{]} Ada dua keterbatasan interpretatif tambahan.
Pertama, perubahan natural harbor dapat merefleksikan proses pesisir dan
sejarah lokal yang juga berhubungan dengan ekonomi, sehingga status
harbor adalah proxy untuk ketepatan lokasi, bukan eksperimen murni.
Kedua, outcome 2015 menggabungkan akumulasi proses berabad-abad;
koefisien jarak tidak mengungkap kapan, melalui sektor apa, atau melalui
migrasi siapa bayangan terbentuk. Uji robustness memperkecil beberapa
penjelasan alternatif, tetapi TXT tidak melaporkan desain yang secara
langsung mengisolasi semua jalur tersebut.

{[}Inferensi review{]} Terdapat pula isu rekonsiliasi angka yang perlu
diperjelas bila naskah hendak direplikasi: teks utama menyebut 4.263
ancient ports setelah eksklusi, sedangkan appendix membedakan jumlah
record pelabuhan, jumlah lokasi dengan data kepadatan, dan 4.263 sel
unik; Table 3 melaporkan 4.198 observasi sementara catatan kaki menyebut
71 sel dikeluarkan dari basis sel pelabuhan. Ini tidak otomatis
membatalkan hasil, tetapi denominator dan urutan filter perlu
didokumentasikan secara eksplisit. TXT tidak memberi penjelasan tambahan
(TXT baris 504-508, 827-829, 1556-1578, 1708-1738).

\subsection{Challenges}\label{challenges-67}

{[}Laporan sumber{]} Tantangan identifikasi pertama adalah
\emph{reflective endogeneity}: tempat kota tumbuh dipengaruhi kondisi
lokal dan sekaligus memengaruhi peluang kota di sekitarnya. Tantangan
kedua adalah \emph{wave interference}, ketika perbedaan kecil dalam
jaringan urban, parameter, atau lokasi kota membuat pola gelombang tidak
sinkron saat dirata-ratakan. Naskah menekankan bahwa masalah kedua ini
dapat mengubah bayangan kota kecil menjadi penurunan monoton (TXT baris
134-143, 369-404).

{[}Laporan sumber{]} Tantangan pengukuran dan konstruksi data meliputi
pelacakan koordinat pelabuhan kuno, penentuan apakah natural harbor
masih bertahan melalui citra satelit, pembedaan harbor alami dari
perlindungan buatan, pencocokan titik ke grid, dan penanganan sel yang
memiliki beberapa pelabuhan atau kepadatan yang hilang. Cakupan de
Graauw juga tidak dianggap lengkap di semua area, sehingga lokasi remote
dan 216 lokasi yang kurang pasti diperlakukan melalui eksklusi atau
robustness (TXT baris 502-529, 738-739, 1617-1644, 1708-1738).

{[}Laporan sumber{]} Tantangan inferensi statistik muncul karena satu
kota dapat membuat beberapa sel terdekat tampak memiliki aktivitas kota
yang sama. Penulis harus memperhitungkan korelasi spasial melalui
cluster sekitar 8 km, cluster yang digeser, cluster sekitar 25 km, dan
Conley standard errors. Tantangan konseptual lainnya adalah memisahkan
kompetisi antarport dari crowd-out kota yang lebih umum, yang menjadi
alasan untuk mengestimasi persamaan (3) pada sel pelabuhan (TXT baris
634-663, 772-783, 784-821).

{[}Laporan sumber{]} Pada analisis kota modern, definisi administratif
dapat menggabungkan area padat yang berdekatan ke dalam satu kota dan
membuat area di luar batas tampak kurang padat secara mekanis. Karena
itu penulis membandingkan null random berbobot dengan dan tanpa area
adjustment. Hasil yang berubah setelah adjustment menunjukkan bahwa
definisi unit kota merupakan tantangan substantif, bukan hanya pilihan
visualisasi (TXT baris 863-868, 897-923).

{[}Inferensi review{]} Tantangan terbesar yang tersisa adalah
menerjemahkan pola dari model satu dimensi ke ruang geografis dua
dimensi yang heterogen tanpa menyamakan kesesuaian pola dengan bukti
mekanisme. Sumber mengatasi sebagian tantangan itu dengan prediksi
bentuk dan banyak robustness, tetapi tidak menyediakan observasi
langsung atas semua faktor historis yang dapat memengaruhi cincin 10-50
km.

\subsection{Practical Implications}\label{practical-implications-67}

{[}Interpretasi penulis{]} Naskah menilai bahwa temuan bayangan
memperluas konsekuensi kebijakan berbasis tempat. Keberhasilan
pertumbuhan lokal dapat mengurangi pertumbuhan serupa di tempat yang
berdekatan bila pusat-pusat ekonomi saling bersaing; karena mobilitas
manusia tidak sempurna, lokasi aktivitas ekonomi juga berhubungan dengan
kesejahteraan individu, pemilik lahan, produktivitas, upah, dan biaya
perumahan. Dengan demikian, kebijakan yang hanya menilai efek di lokasi
sasaran dapat mengabaikan dampak pada hierarki urban di sekitarnya (TXT
baris 117-123, 251-258, 930-936).

{[}Inferensi review{]} Implikasi praktis yang paling langsung adalah
perlunya evaluasi spasial di luar batas administratif atau lokasi
penerima kebijakan: ukur apakah aktivitas di sekitar menurun, bergeser
ke jarak yang lebih jauh, atau berubah komposisinya. Jarak 10-30 km
dalam hasil ini dapat dipakai sebagai hipotesis tentang zona kompetisi
untuk konteks yang sebanding, bukan sebagai ambang kebijakan universal.
Kebijakan koordinasi antarkota dan penilaian distribusional juga relevan
secara konseptual, tetapi TXT tidak mengestimasi biaya, manfaat, atau
efektivitas intervensi tertentu.

{[}Batas bukti{]} Sumber tidak menguji kebijakan place-based aktual,
tidak mengukur kesejahteraan secara langsung, dan tidak menunjukkan
bahwa setiap kenaikan aktivitas di pusat menyebabkan kerugian neto di
sekitar. Karena itu, practical implications di atas adalah pembacaan
kebijakan dari mekanisme dan hasil yang dilaporkan, bukan rekomendasi
yang telah divalidasi melalui evaluasi kebijakan.

\subsection{Applications}\label{applications-67}

{[}Laporan sumber{]} Penulis menyebut bahwa bentuk bayangan yang serupa
dapat muncul pada akses spasial terhadap layanan tertentu, dengan contoh
perbedaan akses layanan kesehatan dan \emph{food deserts}. Contoh ini
disampaikan sebagai kemungkinan penerapan gagasan bahwa konsentrasi
akses di satu lokasi dapat mengurangi aktivitas atau layanan di lokasi
lain; naskah tidak mengestimasi kedua kasus tersebut (TXT baris
251-256).

{[}Inferensi review{]} Kerangka ini juga dapat diterapkan sebagai desain
analitis untuk kebijakan pertumbuhan kota dan hubungan antarpusat:
tentukan anchor atau shock lokasi, bedakan outcome kepadatan umum dari
kemunculan pusat besar, lalu uji apakah ada pola dip dan pemulihan pada
jarak tertentu. Namun, hasil pelabuhan kuno tidak boleh dipindahkan
langsung ke layanan kesehatan, food deserts, atau wilayah nonpesisir.
Setiap aplikasi memerlukan anchor, outcome, skala perjalanan, dan asumsi
identifikasi yang baru. Tidak ada hasil empiris untuk aplikasi-aplikasi
tersebut dalam TXT.

\subsection{Future Research
(Optional)}\label{future-research-optional-67}

{[}Laporan sumber{]} TXT tidak memiliki subbagian agenda future research
yang eksplisit. Beberapa pernyataan yang dapat menjadi titik berangkat
adalah bahwa data kependudukan resolusi tinggi untuk periode antara era
kuno dan masa kini akan berguna tetapi tidak tersedia, serta bahwa
pemodelan kegiatan yang sangat tradable pada jarak lebih jauh dan
interaksi banyak kota berada di luar scope (TXT baris 550-559, 716-720).
Penulis juga menyatakan bahwa microstructure selain FKM dapat
menghasilkan bayangan yang sama dan mungkin mempunyai implikasi
kesejahteraan berbeda (TXT baris 421-429).

{[}Inferensi review{]} Agenda lanjutan yang dapat diturunkan secara
eksplisit dari keterbatasan itu adalah:

\begin{itemize}
\tightlist
\item
  Menyusun rangkaian data historis antara masa kuno dan 2015 agar waktu
  munculnya, menguatnya, atau menghilangnya bayangan dapat dibedakan.
\item
  Menguji interaksi beberapa kota awal dan outcome kegiatan yang lebih
  tradable pada jarak di luar 50 km.
\item
  Membandingkan mekanisme increasing returns FKM dengan microstructure
  alternatif, lalu mengukur implikasi kesejahteraan atau
  distribusionalnya, bukan hanya pola kepadatan.
\item
  Menguji definisi kota, batas administratif, dan null spatial benchmark
  yang berbeda, terutama karena area adjustment mengubah hasil untuk
  kota di bawah ambang 500.000.
\item
  Memeriksa apakah pola serupa bertahan pada wilayah dan jenis anchor
  yang tidak termasuk dalam sampel pesisir ini, dengan strategi
  identifikasi yang sepadan.
\end{itemize}

Daftar tersebut adalah inferensi review, bukan klaim bahwa TXT telah
menetapkannya sebagai rencana penelitian penulis.

\subsection{Provenance Notes}\label{provenance-notes-67}

\begin{itemize}
\tightlist
\item
  Review ini dibuat hanya dari TXT B24 yang diminta. Seluruh 2.342 baris
  yang tersedia dibaca, termasuk \texttt{{[}PDF\ Metadata{]}}, teks
  utama, Abstract, bagian I-VI, References, Online Appendix, Data
  Appendix, tabel, serta caption Figures 1-8 dan Appendix Figures
  A.1-A.21.
\item
  TXT menyatakan sumber PDF berada pada
  \texttt{working-paper/B24-hornbeck-michaels-rauch-2024-identifying-agglomeration-shadows-long-run-evidence-from-ancient-ports-10.2139-ssrn.4889236.pdf},
  diekstrak pada 2026-08-31 dengan \texttt{pdftotext\ -layout}; metadata
  PDF menyebut 69 halaman (TXT baris 1-30). Review menggunakan teks
  hasil ekstraksi dan caption yang tersedia, bukan pemeriksaan visual
  terpisah atas raster/plot PDF.
\item
  Identitas yang eksplisit di dalam TXT adalah judul, tiga penulis,
  CESifo Working Paper No.~11188, June 2024, JEL R120 dan N900,
  keywords, serta URL SSRN abstract 4889236 (TXT baris 31-111). Status
  peer-review tidak disebutkan dalam file. Versi publikasi final atau
  riwayat versi lanjutan juga tidak disebutkan dalam file dan tidak
  ditambahkan dari luar.
\item
  \texttt{{[}Laporan\ sumber{]}} dipakai untuk hasil, metode, data, dan
  pernyataan yang dikaitkan langsung dengan naskah.
  \texttt{{[}Interpretasi\ penulis{]}} dipakai ketika naskah memberi
  makna mekanistik atau implikasi atas hasilnya.
  \texttt{{[}Inferensi\ review{]}} dipakai untuk evaluasi, kehati-hatian
  kausal, generalisasi, dan agenda lanjutan. Tidak ada kutipan yang
  diperlakukan sebagai bukti di luar konteks locator-nya.
\item
  Locator mengutamakan baris TXT dan dilengkapi nama bagian, tabel, atau
  gambar. Angka dan definisi dipertahankan sebagaimana tertulis dalam
  TXT; khususnya, isu perbedaan penyebutan jumlah 4.263 pelabuhan versus
  sel pelabuhan dan denominator Table 3 dicatat sebagai hal yang perlu
  direkonsiliasi, bukan diperbaiki secara spekulatif.
\item
  Kode estimasi, data mentah, hasil lengkap semua spesifikasi
  robustness, status peer-review, dan publikasi final: Tidak disebutkan
  dalam file. Review tidak mengklaim telah mereplikasi hasil,
  memverifikasi referensi, atau mengakses sumber di luar TXT.
\end{itemize}

\section{Review B25}\label{review-b25}

\textbf{Identitas sumber}

\begin{itemize}
\tightlist
\item
  \textbf{Kode:} B25.
\item
  \textbf{Judul:} \emph{The Identification and Dynamics of Urban Shadow
  Areas from the Perspective of People Flows---A Case Study of Nanjing}.
\item
  \textbf{Penulis:} Weiting Xiong dan Junyan Yang. Afiliasi yang
  tercantum adalah Department of Urban Planning, College of Landscape
  Architecture, Nanjing Forestry University, dan Department of Urban
  Planning, School of Architecture, Southeast University.
\item
  \textbf{Jurnal dan bibliografi:} \emph{Buildings}, 2023, volume 13,
  artikel 2934. Nomor isu tidak disebutkan dalam TXT; nomor 2934 adalah
  nomor artikel, bukan rentang halaman konvensional.
\item
  \textbf{DOI:} \texttt{10.3390/buildings13122934}. DOI ditampilkan
  sebagai \texttt{https://doi.org/10.3390/buildings13122934}.
\item
  \textbf{Tanggal publikasi:} 24 November 2023. TXT juga mencatat
  diterima 21 November 2023, setelah direvisi 13 November 2023 dan
  diterima awal 15 Oktober 2023.
\item
  \textbf{Kata kunci:} urban shadow area, people flow, cellular
  signaling data, urban morphology, Nanjing.
\item
  \textbf{Jenis dan basis sumber:} artikel jurnal; review ini hanya
  menggunakan file TXT B25 yang memuat metadata PDF dan seluruh teks
  hasil ekstraksi \texttt{pdftotext\ -layout}. Status peer-review tidak
  diverifikasi dalam TXT.
\end{itemize}

\subsection{Summarized Abstract}\label{summarized-abstract-68}

\textbf{Laporan sumber:} Artikel membahas area bayangan perkotaan, yaitu
ruang yang terbentuk oleh pembangunan yang tidak seimbang dan tidak
memadai dalam proses urbanisasi cepat. Penulis menyatakan bahwa
perhatian terhadap identifikasi dan karakterisasi area tersebut masih
relatif terbatas. Dengan menggabungkan data morfologi perkotaan dan data
sinyal seluler, studi kasus Nanjing mengusulkan metode identifikasi dari
perspektif arus manusia. Hasil empiris mengidentifikasi 19 area bayangan
di kawasan pusat Nanjing; 11 berada di pusat kota lama dan sisanya
tersebar di bagian periferi. Area-area tersebut heterogen dalam
kepadatan bangunan dan intensitas pembangunan, serta tidak terisolasi
karena terhubung erat dengan bagian lain kota. Berdasarkan pola koneksi
spatio-temporal, 19 area dibagi menjadi empat tipe dan tiap tipe
dianalisis melalui satu area perwakilan. Penulis menyarankan agar
kebijakan memperhatikan heterogenitas spasial, waktu koneksi eksternal,
dan fungsi dominan setiap area (TXT baris 71--90; abstrak, artikel hlm.
1).

\textbf{Interpretasi penulis:} Area bayangan tidak diperlakukan hanya
sebagai ruang negatif di sekitar pusat kota. Arus manusia digunakan
untuk menambahkan dimensi dinamis pada identifikasi morfologis, sehingga
kebijakan pembaruan dapat disesuaikan dengan cara area tersebut
terhubung dan digunakan pada waktu yang berbeda (TXT baris 147--166;
782--793).

\textbf{Inferensi review:} Kontribusi empiris artikel terutama bersifat
pemetaan, pengukuran deskriptif, dan tipologi. TXT tidak menunjukkan
pengujian kausal bahwa karakteristik morfologi tertentu menyebabkan pola
arus tertentu. Karena itu, abstrak mendukung klaim tentang keberadaan,
heterogenitas, dan konektivitas 19 area di Nanjing, tetapi tidak cukup
untuk menggeneralisasi mekanisme pembentukan area bayangan ke kota lain.

\subsection{Problem Statement}\label{problem-statement-68}

\textbf{Laporan sumber:} Penulis memulai dari fenomena
``dark-under-the-lights'', yaitu ruang yang berada dekat pusat kota
tetapi tertinggal secara pembangunan. Dalam konteks kota-kota Tiongkok
yang mengalami urbanisasi cepat, area bayangan digambarkan berkaitan
dengan bangunan berkepadatan rendah, jenis usaha dan fungsi yang lebih
rendah, serta kesenjangan terhadap kawasan pusat. Penelitian terdahulu
lebih banyak mengkaji area bayangan pada skala regional, sedangkan
fenomena intra-perkotaan relatif kurang diteliti (TXT baris 97--116;
175--219).

\textbf{Laporan sumber:} Penulis juga menyatakan bahwa penelitian yang
hanya memakai data morfologi, seperti intensitas pembangunan dan
kepadatan bangunan, belum sepenuhnya menangkap pola perkembangan dinamis
atau mekanisme pembentukan area bayangan. Area-area tersebut dapat
menjadi tempat limpahan fungsi nonpusat, seperti hunian, restoran,
pengantaran, ritel, dan layanan lain. Limpahan itu membentuk arus
manusia, logistik, dan informasi dengan kawasan pusat; arus aktivitas
harian manusia dipandang paling langsung mencerminkan kebutuhan berbagai
kelompok (TXT baris 140--159).

\textbf{Interpretasi penulis:} Masalah penelitian dirumuskan sebagai
kebutuhan untuk mengidentifikasi batas spasial area bayangan secara
lebih kuantitatif sekaligus memahami koneksi dinamisnya dengan bagian
kota lain. Nanjing dipilih sebagai kasus untuk mengembangkan metode dan
merumuskan implikasi kebijakan berbasis perspektif statis dan dinamis
(TXT baris 160--166).

\textbf{Inferensi review:} Masalah yang benar-benar dioperasionalkan
adalah masalah klasifikasi spasial dan deskripsi pola koneksi, bukan
evaluasi dampak area bayangan terhadap kesejahteraan, harga lahan,
produktivitas, atau hasil pembaruan. TXT tidak menyatakan pertanyaan
penelitian formal atau hipotesis eksplisit; tujuan dan masalah di atas
diringkas dari uraian pendahuluan.

\subsection{Objectives}\label{objectives-68}

\textbf{Laporan sumber:} Tujuan yang dinyatakan penulis adalah
mengembangkan metode untuk mengidentifikasi area bayangan perkotaan di
Nanjing, menganalisis bagaimana area tersebut terhubung secara dinamis
dengan bagian lain kota melalui data morfologi dan arus manusia, serta
menyusun implikasi kebijakan untuk mengurangi efek negatifnya (TXT baris
160--166).

\textbf{Laporan sumber:} Secara operasional, artikel bertujuan untuk:
(1) menetapkan batas area bayangan pada unit blok; (2) menggambarkan
karakteristik morfologi area yang teridentifikasi; (3) mengukur dinamika
koneksi melalui volatilitas arus manusia dan rasio siang--malam; (4)
membagi area menjadi empat tipe berdasarkan kedua indikator tersebut;
dan (5) membahas karakteristik tiap tipe melalui area perwakilan (TXT
baris 336--468; 470--582).

\textbf{Informasi yang tidak tersedia:} TXT tidak memberikan daftar
pertanyaan penelitian bernomor, hipotesis statistik, atau target
validasi kuantitatif yang terpisah dari tujuan tersebut.

\subsection{Literature Survey}\label{literature-survey-68}

\textbf{Laporan sumber:} Kajian pustaka menelusuri fenomena area
bayangan dari teori geografi ekonomi baru, khususnya efek
``agglomeration shadow'' dan konsep ``borrowed size''. Pada skala
regional, kota kecil yang dekat kota besar dapat terhambat ketika efek
kompetisi lebih kuat daripada manfaat ukuran yang dipinjam. Literatur
juga memakai konsep ``interplaces between cities'', ``distressed
inner-city neighborhoods'', dan ``inner-city suburbanization'' untuk
fenomena yang berdekatan tetapi tidak sepenuhnya sama (TXT baris
175--201).

\textbf{Laporan sumber:} Dalam konteks Tiongkok, penulis menyebut konsep
metropolitan shadow area, development shadow areas, metropolitan
multiple shadow areas, city center shadow area, dan city shadow area.
Walaupun istilahnya berbeda, literatur yang dirangkum penulis mengaitkan
area bayangan dengan kesenjangan pembangunan, lokasi yang tersebar dekat
pusat kota, fungsi layanan tingkat rendah, bangunan lama, penurunan
kepadatan dan ketinggian fasilitas publik, efek negatif pusat terhadap
area sekitarnya, serta koneksi yang lemah dengan industri pusat (TXT
baris 202--219).

\textbf{Laporan sumber:} Untuk identifikasi, penelitian regional
menggunakan indikator dan metode seperti koefisien sensitivitas dan
pengaruh, break points, comprehensive connection strength, model Huff,
dan autokorelasi spasial. Penelitian intra-perkotaan menggunakan
indikator sosial-ekonomi seperti harga rumah, tingkat kriminalitas, dan
tingkat pekerjaan, serta indikator morfologi seperti intensitas
pembangunan dan kepadatan bangunan. Data besar multi-sumber, termasuk
morfologi, sinyal seluler, dan POI, dipakai oleh studi lain untuk
memperluas analisis ke cakupan kota (TXT baris 221--238).

\textbf{Laporan sumber:} Untuk mekanisme pembentukan dan evolusi,
penulis merangkum teori agglomeration--diffusion, preferensi hunian
keluarga kelas menengah, lokasi dan aksesibilitas, kebijakan
perencanaan, struktur ruang, teori core--edge dan diffusion--echo,
pembatasan kebijakan, pembagian fungsi, struktur industri, aktivitas
harian, strategi integrasi regional, dan eksternalitas jaringan (TXT
baris 244--267).

\textbf{Interpretasi penulis:} Artikel menempatkan dirinya di antara
literatur morfologi dan literatur jaringan perkotaan: identifikasi perlu
menggabungkan kesenjangan fisik dengan aktivitas dan koneksi arus
manusia. Ini merupakan alasan teoritis dan metodologis untuk memakai
perspektif statis sekaligus dinamis (TXT baris 140--166).

\textbf{Inferensi review:} TXT tidak menjelaskan prosedur pencarian
pustaka yang sistematis, rentang waktu pencarian, kriteria
inklusi-eksklusi, jumlah korpus, atau penilaian mutu studi. Dengan
demikian, bagian ini berfungsi sebagai tinjauan konseptual untuk
membangun posisi penelitian, bukan systematic literature review yang
metodenya dapat direplikasi.

\subsection{Method Used}\label{method-used-68}

\textbf{Laporan sumber:} Unit analisis dasar adalah blok. Penulis
memilih blok karena dibatasi jalan dengan batas spasial yang jelas,
relatif utuh, memiliki makna perencanaan dan kognitif, serta dianggap
lebih stabil daripada petak lahan yang dapat menunjukkan karakteristik
ekstrem akibat intensitas pembangunan lokal. Grid dinilai sebagai
pembagian ruang yang acak dan kurang memiliki makna perencanaan (TXT
baris 336--352).

\textbf{Laporan sumber:} Area bayangan didefinisikan sebagai area yang
berdekatan dengan pusat kota, tetapi memiliki intensitas konstruksi
lebih rendah, fasilitas pelayanan publik lebih sedikit, dan vitalitas
penduduk yang tidak memadai. Tiga indikator blok dipakai untuk
identifikasi: construction intensity, Facility, dan Popu atau vitalitas
penduduk (TXT baris 354--381).

\textbf{Laporan sumber:} \texttt{Facility\_i} adalah jumlah fasilitas
pelayanan publik berdasarkan tipe fasilitas. \texttt{Popu\_i} adalah
jumlah pengguna telepon seluler yang dicatat oleh stasiun basis di dalam
blok. Untuk \texttt{Intensity\_i}, TXT menjelaskan komponen
\texttt{floor\_area}, \texttt{height}, dan \texttt{block\_area};
persamaan (1) pada hasil ekstraksi terpecah secara tata letak di sekitar
TXT baris 368--381. Karena itu, review ini hanya menyatakan bahwa
indikator tersebut mengagregasikan luas lantai dan tinggi bangunan
terhadap luas blok, tanpa mengklaim bentuk visual persamaan yang lebih
presisi daripada yang tersedia dalam TXT.

\textbf{Laporan sumber:} Karena area bayangan dipandang sebagai konsep
relatif, penulis menggunakan exploratory spatial data analysis dan local
Moran's I untuk mengidentifikasi autokorelasi lokal setiap indikator.
Blok yang signifikan diklasifikasikan sebagai high--high, low--low,
high--low, atau low--high. Blok kategori low--high pada sedikitnya dua
dari tiga indikator dijadikan calon area bayangan. Kriteria tiga dari
tiga pernah dicoba, tetapi dianggap terlalu ketat karena menghasilkan
terlalu sedikit area dan tidak sesuai dengan observasi survei lapangan
penulis. Blok low--high yang bersebelahan kemudian digabungkan memakai
kriteria rook continuity (TXT baris 382--414).

\textbf{Laporan sumber:} Hasil statistik kemudian dikoreksi secara
substantif. Area yang terutama berisi gunung, badan air, lahan
pertanian, atau ruang hijau perlindungan dikeluarkan. Area yang terutama
berisi museum, perpustakaan, sekolah, dan bangunan publik lain juga
dikeluarkan karena dapat memiliki intensitas konstruksi rendah tetapi
tetap aktif dan ramai (TXT baris 416--437).

\textbf{Laporan sumber:} Dinamika diukur melalui dua indikator.
Volatilitas adalah simpangan baku arus manusia keseluruhan sepanjang 12
interval waktu dua jam, dengan arus keseluruhan didefinisikan sebagai
jumlah arus masuk dan arus keluar. Rasio siang--malam membandingkan arus
pada 08.00--10.00 dengan arus pada 02.00--04.00. Nilai arus per area dan
interval diagregasikan pada tingkat stasiun basis (TXT baris 439--468).

\textbf{Laporan sumber:} Nilai rata-rata volatilitas dan rasio
siang--malam digunakan sebagai ambang untuk membagi 19 area menjadi
empat tipe: rendah--rendah, rendah volatilitas--tinggi rasio, tinggi
volatilitas--rendah rasio, dan tinggi--tinggi. Satu area perwakilan
dipilih untuk tiap tipe dan dianalisis secara rinci melalui pola koneksi
pada beberapa waktu (TXT baris 569--582).

\textbf{Informasi metodologis yang tidak tersedia:} TXT tidak
menyebutkan tingkat signifikansi local Moran's I, matriks bobot spasial,
prosedur standardisasi, jumlah total blok sebelum seleksi, kriteria
pasti untuk memilih area perwakilan, atau prosedur validasi kuantitatif
terhadap hasil klasifikasi.

\subsection{Dataset}\label{dataset-68}

\textbf{Laporan sumber:} Dataset morfologi perkotaan dibangun dari peta
topografi vektor kawasan pusat Nanjing yang disebut lengkap dan akurat.
Peta tersebut dikalibrasi, dikoreksi, dan diperbarui setelah survei
lapangan oleh kelompok riset penulis. Isinya meliputi bangunan, petak
lahan, blok, jalan, dan lingkungan lanskap alami. Seluruh data
diproyeksikan ke sistem koordinat WGS\_1984 dan indikator dihitung
dengan ArcGIS 10.3 (TXT baris 288--301).

\textbf{Laporan sumber:} Dataset sinyal seluler diperoleh dari China
Mobile, yang menurut penulis mencakup sekitar 70\% pangsa pasar telepon
seluler di Nanjing. Wilayah studi memiliki 104.247 stasiun basis. Data
mentah mencakup 14 hari pada 2015, terdiri atas 10 hari kerja dan 4 hari
akhir pekan. Setelah pembersihan dan pertimbangan kondisi cuaca, penulis
memilih data 11 November sebagai hari kerja dan 21 November sebagai
akhir pekan karena distribusinya normal. Data terpilih mencakup lebih
dari 2 juta pengguna telepon seluler per hari (TXT baris 303--318).

\textbf{Laporan sumber:} Data stasiun basis memuat kode identifikasi dan
koordinat geografis stasiun serta jumlah pengguna dalam tiap interval
dua jam. Penulis membuat poligon yang disebut ``Tyson polygons'' dalam
TXT berdasarkan lokasi stasiun, mengalokasikan stasiun ke blok, lalu
memperoleh jumlah orang per blok dan interval. Data pensinyalan memuat
ID terminal yang dienkripsi secara anonim, tipe sinyal, dan waktu
koneksi ke stasiun basis. Lokasi pengguna diproyeksikan secara acak ke
blok yang berdekatan dengan stasiun terkait, kemudian titik tinggal
pengguna dipetakan menjadi trajektori harian (TXT baris 319--334).

\textbf{Laporan sumber:} Pernyataan ketersediaan data menyebutkan bahwa
data dapat diminta kepada penulis korespondensi, tetapi tidak tersedia
untuk publik karena alasan privasi (TXT baris 879--880).

\textbf{Informasi yang tidak tersedia:} TXT tidak memberikan ukuran
berkas, tingkat data hilang setelah pembersihan, daftar seluruh variabel
morfologi beserta satuannya, distribusi pengguna per stasiun, atau
prosedur anonimisasi yang lebih rinci daripada keterangan ID terenkripsi
anonim.

\subsection{Population / Sample}\label{population-sample-68}

\textbf{Laporan sumber:} Populasi spasial studi adalah kawasan pusat
Nanjing yang mencakup Jiangbei, Xianlin, Dongshan, dan main city, dengan
luas 846 km² dan penduduk sekitar 6,7 juta. Kawasan ini mencakup 12,8\%
wilayah administratif Nanjing tetapi menampung 63,2\% penduduk kota (TXT
baris 269--282).

\textbf{Laporan sumber:} Populasi observasi arus manusia berasal dari
pengguna China Mobile yang tercakup sinyal di kawasan studi, yaitu lebih
dari 2 juta pengguna per hari pada data terpilih. ID terminal dienkripsi
secara anonim. Artikel tidak melaporkan jumlah pengguna unik sepanjang
dua hari, karakteristik demografis, atau aturan inklusi individu selain
keberadaan data sinyal (TXT baris 303--329).

\textbf{Laporan sumber:} Sampel waktu awal terdiri atas 14 hari pada
2015, lalu dipersempit menjadi 11 November dan 21 November setelah
pembersihan serta pertimbangan cuaca. Unit analisis akhir bukan
individu, melainkan blok; hasil identifikasi mencakup 19 area bayangan
yang terdiri atas 97 blok (TXT baris 303--318; 338--352; 470--481).

\textbf{Inferensi review:} Angka 19 area dan 97 blok adalah hasil
klasifikasi, bukan sampel probabilitas dari seluruh kota. Total blok di
kawasan pusat, proporsi blok yang masuk hasil, dan desain pengambilan
sampel individu tidak tersedia. Karena itu, representativitas arus
pengguna terhadap seluruh penduduk dan representativitas dua hari
terhadap pola tahunan tidak dapat ditentukan dari TXT.

\subsection{Dependent Variables}\label{dependent-variables-68}

\textbf{Inferensi review:} TXT tidak memuat model kausal atau regresi
dengan variabel dependen dan independen konvensional. Tiga indikator
yang dipakai untuk membatasi area adalah intensitas konstruksi, jumlah
fasilitas publik, dan vitalitas penduduk. Dua ukuran keluaran untuk
dinamika adalah volatilitas arus manusia dan rasio siang--malam (TXT
baris 354--381; 439--468).

\textbf{Interpretasi penulis:} Volatilitas merepresentasikan kestabilan
atau perubahan arus keseluruhan selama 12 interval, sedangkan rasio
siang--malam merepresentasikan perbedaan arus pada dua interval yang
dianggap mewakili perilaku siang dan malam. Arus keseluruhan
menggabungkan arus masuk dan arus keluar, bukan hanya perpindahan bersih
(TXT baris 440--467).

\textbf{Inferensi review:} Jika dipaksa dalam kerangka variabel
dependen, keluaran yang paling dekat adalah status area bayangan dan
tipe dinamisnya. Namun, keduanya merupakan hasil aturan klasifikasi
deskriptif, bukan respons yang diuji terhadap prediktor. \texttt{Popu}
juga merupakan proksi vitalitas berbasis catatan pengguna stasiun basis,
bukan ukuran populasi sensus; hal ini perlu dipertahankan ketika
menafsirkan hasil.

\textbf{Informasi yang tidak tersedia:} TXT tidak memberikan satuan
volatilitas, metode normalisasi rasio, ambang signifikansi untuk
autokorelasi, atau model statistik untuk mengestimasi pengaruh
intensitas, fasilitas, dan vitalitas terhadap arus manusia.

\subsection{Results}\label{results-68}

\textbf{Laporan sumber:} Identifikasi menghasilkan 19 area bayangan di
kawasan studi. Sebelas berada di pusat kota lama dan delapan lainnya
tersebar di periferi. Secara keseluruhan, 19 area mencakup 97 blok dan
luas sekitar 13,53 km², dengan luas lantai bangunan sekitar 7,21 km² dan
luas tapak bangunan sekitar 2,88 km². Penulis melaporkan rasio plot
rata-rata sekitar 0,53 dan kepadatan bangunan rata-rata 0,21. Luas
individual berkisar 4,19--323,59 hektare dan jumlah blok per area
berkisar 1--14; seluruhnya dekat sedikitnya satu pusat kota (TXT baris
470--483; Tabel 1, baris 489--512).

\textbf{Laporan sumber:} Nama 19 area sebagaimana tercantum dalam Tabel
1 adalah Wutang Parcel, Qingliangmen Avenue, Liuhe Parcel, Baijiahu
Parcel, Andemen Avenue, Kening Road, Fengtainan Road, Shengzhou Road,
Huowaxiang Parcel, Youfuxi Street, Hunan Road, Pukou Parcel, Xianlin
Parcel, Yihe Road, Cibeishe Parcel, Yunnanbei Road, Qingshi Street,
Guyilong Parcel, dan Jiankang Road (TXT baris 489--512).

\textbf{Laporan sumber:} Pola spasial disebut sebagai aglomerasi
internal dan dispersi eksternal. Area bayangan terkonsentrasi di sekitar
Xinjiekou, Hunan Road, dan Confucius Temple pada pusat kota lama,
termasuk koridor Zhongshanbei Road, Zhongshan Road, Zhongshandong Road,
serta Zhonghua Road hingga Zhonghuamen Gate. Delapan area di luar pusat
lama lebih tersebar, antara lain di sekitar Hexi, Dongshan, Jiangbei,
dan Xianlin (TXT baris 515--540).

\textbf{Laporan sumber:} Volatilitas arus berkisar 1,80--33,86 dengan
rata-rata 15,27. Wutang Parcel memiliki nilai terendah dan Yihe Road
tertinggi. Rasio siang--malam memiliki rata-rata 5,04; Xianlin Parcel
terendah dengan 2,14 dan Jiankang Road tertinggi dengan 9,70 (TXT baris
548--567).

\textbf{Laporan sumber:} Pembagian berdasarkan rata-rata kedua indikator
ditulis sebagai berikut. Tipe rendah--rendah mencakup Wutang Parcel,
Kening Road, Qingliangmen Avenue, Guyilang Parcel, Hunan Road, Xianlin
Parcel, dan Liuhe Parcel. Tipe rendah volatilitas--tinggi rasio mencakup
Baijiahu Parcel, Andemen Avenue, Fengtainan Road, dan Shengzhou Road.
Tipe tinggi volatilitas--rendah rasio mencakup Huowaxiang Parcel,
Yunanbei Road, Pukou Parcel, dan Youfuxi Street. Tipe tinggi--tinggi
mencakup Jiankang Road, Qingshi Street, Cibeishe Parcel, dan Yihe Road
(TXT baris 569--582).

\textbf{Laporan sumber:} Empat contoh rinci menghasilkan gambaran
berikut. Wutang Parcel terutama terhubung dengan area sekitarnya dengan
kekuatan dan fluktuasi rendah; Baijiahu Parcel menunjukkan arus siang
yang lebih kuat daripada malam dan koneksi yang relatif stabil; Hunan
Road digambarkan sebagai area dengan fluktuasi tinggi dan perbedaan
siang--malam rendah; Cibeishe Parcel memiliki fluktuasi dan rasio
siang--malam tinggi (TXT baris 585--615; 634--667; 684--717; 733--760).

\textbf{Catatan konsistensi hasil:} Ada konflik internal yang harus
dibaca sebagai masalah pelaporan sumber. Hunan Road dimasukkan ke tipe
rendah--rendah dalam daftar pada TXT baris 571--574, tetapi Bagian 5.3
secara eksplisit menjadikannya contoh tipe tinggi volatilitas--rendah
rasio pada baris 684--695. Nama \texttt{Guyilong\ Parcel} pada Tabel 1
muncul sebagai \texttt{Guyilang\ Parcel} dalam uraian distribusi dan
daftar tipe; \texttt{Yunnanbei\ Road} pada Tabel 1 muncul sebagai
\texttt{Yunanbei\ Road} dalam daftar tipe; \texttt{Huaihai\ Road}
disebut dalam uraian pusat Xinjiekou tetapi tidak tercantum dalam daftar
19 area di Tabel 1 (TXT baris 489--512; 515--526; 569--580). Review
tidak memilih versi yang benar dan tidak mengoreksi data tersebut.

\subsection{Findings}\label{findings-68}

\textbf{Interpretasi penulis:} Temuan utama penulis adalah bahwa area
bayangan perkotaan ada di dalam kota dan tidak dapat dipahami sebagai
ruang terisolasi atau semata-mata negatif. Area tersebut membentuk
koneksi dinamis dengan pusat di sekitarnya dan bagian lain kota,
sehingga dapat dipandang sebagai simpul dalam struktur jaringan
perkotaan (TXT baris 774--793).

\textbf{Interpretasi penulis:} Heterogenitas muncul pada tiga dimensi.
Pertama, distribusinya berbeda antara pusat kota lama yang beraglomerasi
dan periferi yang lebih tersebar. Kedua, kekuatan dan volatilitas
koneksi berbeda menurut waktu dalam satu hari. Ketiga, fungsi dominan,
seperti hunian, manufaktur, perdagangan, layanan, atau ruang dalam
pembaruan, berkaitan dengan pola arus yang berbeda (TXT baris 806--842).

\textbf{Interpretasi penulis:} Keempat tipe dinamis dijelaskan secara
kualitatif. Tipe rendah--rendah dikaitkan dengan area periferi yang
terfragmentasi, lahan kosong atau pertanian/kehutanan, aksesibilitas
lemah, dan daya tarik rendah. Tipe rendah volatilitas--tinggi rasio
dikaitkan dengan fungsi produksi/manufaktur, kesempatan kerja relatif
stabil, arus pekerja siang hari, serta basis penduduk tinggal yang
kecil. Tipe tinggi volatilitas--rendah rasio dikaitkan dengan campuran
lahan, fasilitas komersial, ruang terbuka atau konstruksi, aktivitas
beragam, tetapi fungsi hunian dan pekerjaan yang lemah. Tipe
tinggi--tinggi dikaitkan dengan hunian lama dekat pusat, fasilitas
publik dan kesempatan kerja, biaya lahan serta biaya hidup tinggi, dan
fenomena arus pasang-surut (TXT baris 585--760).

\textbf{Interpretasi penulis:} Penjelasan mekanisme pada empat contoh
bersifat penalaran berbasis morfologi, fungsi, aksesibilitas, dan pola
aktivitas. Misalnya, Baijiahu dikaitkan dengan industri dan akses yang
baik, Hunan Road dengan penurunan daya tarik dan
pembongkaran-rekonstruksi, serta Cibeishe dengan hunian lama dekat pusat
komersial. Penulis tidak menguji mekanisme tersebut dengan model
kuantitatif dalam artikel ini (TXT baris 601--615; 647--667; 694--717;
741--760; 863--869).

\textbf{Inferensi review:} Tipologi berhasil mengorganisasi variasi
deskriptif dari 19 kasus, tetapi kekuatan penjelasnya dibatasi oleh
ambang rata-rata dan empat studi contoh. Konflik klasifikasi Hunan Road
juga membuat daftar tipe dalam TXT tidak sepenuhnya konsisten. Karena
itu, tipologi dapat dibaca sebagai hasil eksploratif kasus Nanjing,
bukan klasifikasi tervalidasi yang bebas ambiguitas.

\subsection{Conclusions}\label{conclusions-68}

\textbf{Laporan sumber:} Kesimpulan artikel menyatakan bahwa area
bayangan merupakan fenomena umum selama urbanisasi cepat, tetapi upaya
untuk menggambarkan dan mengkarakterisasinya secara kuantitatif masih
terbatas. Dengan data morfologi perkotaan dan sinyal seluler, penulis
mengidentifikasi serta menelaah koneksi dinamis 19 area di kawasan pusat
Nanjing. Area tersebut berbeda dalam kepadatan bangunan dan intensitas
pembangunan; 11 berada di pusat kota lama dan delapan tersebar di
periferi (TXT baris 844--855).

\textbf{Laporan sumber:} Penulis menyimpulkan bahwa 19 area tidak
terisolasi, melainkan terhubung dengan bagian lain Nanjing melalui pola
koneksi spatio-temporal yang berbeda. Karena itu, area tersebut dibagi
menjadi empat tipe berdasarkan pola koneksi, masing-masing dengan satu
area perwakilan untuk telaah lebih rinci. Kebijakan, menurut penulis,
perlu memperhatikan distribusi spasial, waktu koneksi eksternal, dan
fungsi dominan (TXT baris 855--862).

\textbf{Laporan sumber:} Penulis juga menyebut dua batasan dan arah
lanjutan: dinamika hanya diamati pada tanggal tertentu karena
keterbatasan data, sehingga penelitian berikutnya dapat membandingkan
koneksi dari perspektif evolusioner; serta mekanisme pembentukan tiap
tipe baru dibahas secara kualitatif, sehingga perlu diteliti secara
lebih kuantitatif (TXT baris 863--869).

\subsection{Contributions}\label{contributions-68}

\textbf{Interpretasi penulis:} Penulis menyatakan kontribusi pada
literatur dengan mengidentifikasi dan mengkarakterisasi area bayangan
dari perspektif statis dan dinamis. Secara operasional, kontribusinya
berupa penggabungan indikator morfologi, fasilitas, vitalitas penduduk,
dan arus manusia dalam prosedur berbasis blok; penggunaan local Moran's
I untuk batas relatif; serta tipologi koneksi spatio-temporal untuk 19
area di Nanjing (TXT baris 160--166; 354--414; 844--862).

\textbf{Laporan sumber:} Kontribusi terapan yang diklaim adalah
menyediakan dasar bagi kebijakan pembaruan area bayangan yang peka
terhadap heterogenitas lokasi, waktu, dan fungsi dominan, bukan
kebijakan tunggal untuk semua area (TXT baris 794--842).

\textbf{Inferensi review:} Dari TXT saja, kontribusi yang dapat
dipertahankan adalah kerangka operasional eksploratif dan demonstrasi
kasus. Tidak ada dasar untuk menyebut artikel ini memberi estimasi
kausal, validasi lintas kota, atau teori mekanisme yang telah
dikonfirmasi, karena penulis sendiri menyatakan mekanisme masih dibahas
secara kualitatif dan data dinamis hanya berasal dari tanggal tertentu.

\subsection{Research Gap}\label{research-gap-68}

\textbf{Gap yang dinyatakan penulis:} Literatur sebelumnya lebih banyak
membahas shadow area pada skala regional, sedangkan area bayangan
intra-perkotaan relatif kurang mendapat perhatian. Studi intra-perkotaan
menggunakan indikator sosial-ekonomi dan morfologi; secara khusus,
penulis menilai pendekatan yang hanya berbasis morfologi belum menangkap
pola dinamis dan mekanisme pembentukan. Selain itu, upaya delineasi dan
karakterisasi secara kuantitatif masih relatif sedikit (TXT baris
97--116; 140--166; 221--238; 844--850).

\textbf{Gap yang dinyatakan penulis:} Artikel secara eksplisit
menyisakan kebutuhan untuk membandingkan koneksi area bayangan dari
perspektif evolusioner dan untuk mengeksplorasi mekanisme pembentukan
tipe yang berbeda secara kuantitatif (TXT baris 863--869).

\textbf{Inferensi review:} Berdasarkan desain yang dilaporkan, celah
yang belum terisi mencakup validasi klasifikasi terhadap observasi
independen, pengujian sensitivitas terhadap pilihan indikator dan
ambang, serta pengukuran representativitas data China Mobile dan dua
hari observasi. Ini adalah inferensi review dari informasi yang
tersedia, bukan daftar gap tambahan yang dinyatakan eksplisit oleh
penulis. TXT juga tidak menyediakan protokol tinjauan pustaka yang
memungkinkan gap literatur diverifikasi secara sistematis.

\subsection{Limitations}\label{limitations-68}

\textbf{Laporan sumber:} Keterbatasan eksplisit pertama adalah observasi
koneksi dinamis hanya dilakukan pada satu tanggal yang dipilih dari data
2015, meskipun dataset awal mencakup 14 hari. Keterbatasan kedua adalah
mekanisme pembentukan berbagai tipe hanya dibahas secara kualitatif.
Keduanya disebut sebagai alasan untuk penelitian lanjutan (TXT baris
863--869).

\textbf{Laporan sumber:} Data sinyal tidak tersedia secara publik karena
privasi dan hanya dapat diminta kepada penulis korespondensi. Cakupan
operator juga sekitar 70\% pangsa pasar telepon seluler Nanjing,
sehingga TXT tidak menyatakan bahwa data mewakili seluruh penduduk (TXT
baris 303--305; 879--880).

\textbf{Inferensi review:} Keterbatasan metodologis lain yang terlihat
dari TXT adalah penggunaan pengguna stasiun basis sebagai proksi
vitalitas, pemetaan lokasi pengguna secara acak ke blok yang berdekatan,
dan penggunaan dua interval dua jam untuk rasio siang--malam.
Keterbatasan ini dapat memengaruhi interpretasi arus dan batas area,
tetapi besaran pengaruhnya tidak dapat dihitung karena tingkat
kesalahan, data hilang, dan validasi tidak dilaporkan.

\textbf{Informasi yang tidak tersedia:} TXT tidak melaporkan total blok
awal, tingkat signifikansi dan spesifikasi bobot spasial, uji
robustness, kriteria kuantitatif pengeluaran lanskap/bangunan publik,
atau validasi eksternal dan replikasi di kota lain.

\subsection{Challenges}\label{challenges-68}

\textbf{Laporan sumber:} Penulis menyebut bahwa gangguan jaringan,
keterlambatan transmisi, interferensi lingkungan eksternal, kegagalan
sistem, dan kondisi tidak ideal lain dapat menghasilkan data tidak
valid. Data harus dibersihkan dan dipertimbangkan bersama kondisi cuaca
sebelum dua hari dipilih (TXT baris 306--318).

\textbf{Laporan sumber:} Tantangan konseptualnya adalah area bayangan
merupakan konsep relatif, sehingga batasnya tidak cukup ditentukan oleh
nilai rendah absolut. Penulis memakai low--high local Moran's I dan
mensyaratkan sedikitnya dua dari tiga indikator. Kriteria tiga indikator
dinilai terlalu ketat; hasil statistik juga harus dikoreksi agar lanskap
alami dan bangunan publik tidak salah diklasifikasikan sebagai area
bayangan (TXT baris 382--437).

\textbf{Laporan sumber:} Tantangan interpretasi muncul karena pola arus
harus dibaca bersama struktur jalan, aksesibilitas, fungsi lahan,
bangunan, fasilitas sekitar, dan pusat kota yang berdekatan. Empat
contoh kasus dipakai untuk menelusuri penjelasan tersebut, tetapi
mekanisme dasarnya belum dihitung secara kuantitatif (TXT baris
585--760; 863--869).

\textbf{Inferensi review:} Penggabungan data operator, stasiun basis,
poligon spasial, titik tinggal, dan blok berpotensi menimbulkan
ketidakpastian pemetaan serta ketergantungan pada cakupan operator. TXT
mendeskripsikan prosedur agregasi, tetapi tidak mengukur ketidakpastian
itu. Konflik klasifikasi Hunan Road dan variasi ejaan nama area juga
merupakan tantangan reproduksibilitas pada tingkat pelaporan.

\subsection{Practical Implications}\label{practical-implications-68}

\textbf{Laporan sumber:} Penulis menilai area bayangan hadir pada
berbagai tahap perkembangan kota dan efek negatifnya tidak dapat dihapus
secara permanen hanya melalui perencanaan dan desain karena pembangunan
yang tidak seimbang tetap ada. Sumber daya khas, termasuk peninggalan
sejarah, dapat diaktifkan melalui organisasi ruang, ruang publik, dan
jalur aktivitas agar terhubung dengan konteks budaya kota (TXT baris
794--805).

\textbf{Laporan sumber:} Implikasi pertama adalah membedakan kebijakan
berdasarkan lokasi. Area di pusat kota lama yang memiliki lokasi
menguntungkan dan arus lebih sering dapat memanfaatkan kedekatan dengan
pusat kota untuk memperkuat vitalitas. Area periferi perlu memperkuat
keunggulan komparatif dan integrasinya ke struktur jaringan kota (TXT
baris 806--814).

\textbf{Laporan sumber:} Implikasi kedua adalah memperhitungkan
heterogenitas waktu. Volatilitas dan kekuatan arus berubah menurut
interval harian serta dipengaruhi area bayangan dan pusat di sekitarnya.
Penulis memberi contoh pasar malam sebagai cara memanfaatkan ruang
terbuka dan karakter temporal, terutama di pusat kota lama, dengan tetap
menyatakannya sebagai kemungkinan kebijakan (TXT baris 815--827).

\textbf{Laporan sumber:} Implikasi ketiga adalah menyesuaikan intervensi
dengan fungsi dominan. Area yang didominasi hunian dapat diarahkan pada
diversifikasi tipe bangunan dan fungsi untuk membentuk komunitas yang
lebih beragam dan inklusif. Area manufaktur dapat mempertimbangkan
pembaruan pabrik tua atau kosong dan pembangunan kompleks industri untuk
menarik orang serta mencampurkan penggunaan lahan (TXT baris 828--842).

\textbf{Inferensi review:} Implikasi tersebut paling kuat sebagai
prinsip diagnosis dan diferensiasi kebijakan. TXT tidak menguji
efektivitas pasar malam, diversifikasi fungsi, atau pembaruan industri,
sehingga rekomendasi tidak boleh dibaca sebagai dampak kebijakan yang
sudah terbukti.

\subsection{Applications}\label{applications-68}

\textbf{Aplikasi yang didukung sumber:} Prosedur artikel dapat digunakan
untuk menyaring blok yang berdekatan dengan pusat kota berdasarkan
kombinasi intensitas konstruksi, jumlah fasilitas publik, dan vitalitas
berbasis arus manusia. Hasilnya dapat dipakai untuk memetakan lokasi
prioritas kajian pembaruan dan mengelompokkan area menurut pola
volatilitas serta perbedaan siang--malam (TXT baris 354--468).

\textbf{Aplikasi yang didukung sumber:} Pemerintah atau perencana dapat
memakai empat tipe sebagai perangkat diagnosis awal: area periferi
dengan koneksi lemah dan stabil, area manufaktur dengan arus siang
dominan, area campuran atau dalam transisi pembaruan, serta area hunian
lama dekat pusat dengan arus pasang-surut. Intervensi kemudian diarahkan
menurut lokasi, waktu, dan fungsi dominan sebagaimana dibahas penulis
(TXT baris 585--760; 794--842).

\textbf{Inferensi review:} Aplikasi paling aman adalah sebagai alat
eksplorasi kasus dan penyusunan hipotesis lapangan, bukan sebagai aturan
otomatis penetapan status area bayangan. Penerapan di luar Nanjing, di
luar cakupan China Mobile, atau pada periode lain belum ditunjukkan
dalam TXT dan memerlukan kalibrasi serta validasi tersendiri.

\subsection{Future Research
(Optional)}\label{future-research-optional-68}

\textbf{Laporan sumber:} Penulis menyebut dua agenda penelitian masa
depan. Pertama, penelitian perlu membandingkan koneksi area bayangan
dari perspektif evolusioner, karena artikel hanya mengamati tanggal
tertentu. Kedua, penelitian perlu mengeksplorasi mekanisme pembentukan
tipe area bayangan secara lebih kuantitatif, karena pembahasan artikel
masih kualitatif (TXT baris 863--869).

\textbf{Batas isi:} Tidak ada agenda masa depan lain yang ditambahkan
tanpa label inferensi. TXT tidak menyebut desain longitudinal tertentu,
model kuantitatif tertentu, atau kota pembanding tertentu.

\subsection{Provenance Notes}\label{provenance-notes-68}

\begin{itemize}
\tightlist
\item
  \textbf{Bahan yang digunakan:} Hanya file
  \texttt{source/journal/B25-xiong-yang-2023-the-identification-and-dynamics-of-urban-shadow-areas-from-the-perspective-of-people-flows-a-case-study-of-nanjing-10.3390-buildings13122934.txt}.
  Seluruh teks yang tersedia telah dibaca, dari baris 1 sampai 960,
  termasuk blok \texttt{{[}PDF\ Metadata{]}}, teks artikel, kesimpulan,
  pernyataan ketersediaan data, dan referensi.
\item
  \textbf{Status akses:} TXT menyatakan bahwa teks berasal dari PDF dan
  diekstrak pada 31 Agustus 2026 dengan \texttt{pdftotext\ -layout}.
  Review ini memperlakukan teks yang tersedia sebagai full-text hasil
  ekstraksi, bukan sebagai hasil pemeriksaan web atau sumber eksternal.
\item
  \textbf{Locator:} Rujukan utama memakai nomor baris TXT, ditambah
  bagian, halaman artikel, tabel, atau gambar jika disebutkan dalam
  teks. Penanda halaman PDF menunjukkan artikel pada halaman tercetak
  1--19, sedangkan metadata PDF menyebut 20 halaman; perbedaan ini tidak
  ditafsirkan lebih lanjut.
\item
  \textbf{Pemisahan klaim:} \texttt{Laporan\ sumber} adalah isi atau
  hasil yang dilaporkan artikel; \texttt{Interpretasi\ penulis} adalah
  penjelasan dan pembacaan mekanisme/kebijakan yang dinyatakan penulis;
  \texttt{Inferensi\ review} adalah evaluasi yang ditarik dari TXT dan
  tidak dipresentasikan sebagai temuan penulis.
\item
  \textbf{Artefak ekstraksi:} TXT mengandung karakter atau fragmen tata
  letak seperti \texttt{tt} dan \texttt{ff}, persamaan yang terpecah,
  serta variasi nama area. Persamaan intensitas tidak direkonstruksi
  dari luar TXT. Nama \texttt{Guyilong}/\texttt{Guyilang},
  \texttt{Yunnanbei}/\texttt{Yunanbei}, dan konflik klasifikasi Hunan
  Road dipertahankan sebagai ketidakkonsistenan sumber.
\item
  \textbf{Metadata yang tidak tersedia atau tidak diverifikasi:} Nomor
  isu, rentang halaman jurnal konvensional, tingkat peer-review, jumlah
  total blok kawasan studi, parameter local Moran's I, detail bobot
  spasial, kriteria pemilihan area perwakilan, tingkat data hilang, dan
  validasi eksternal tidak tersedia dalam TXT atau tidak diverifikasi
  melalui sumber lain.
\item
  \textbf{Batas provenance data:} Data morfologi disebut telah
  diperbarui melalui survei lapangan kelompok penulis, tetapi TXT tidak
  menyediakan instrumen, tanggal, cakupan, atau hasil validasi survei
  tersebut. Data sinyal berasal dari China Mobile dan tidak tersedia
  publik karena privasi; review tidak mengklaim akses ke data mentah.
\item
  \textbf{Batas interpretasi:} Tidak ada informasi dari luar TXT yang
  dipakai untuk melengkapi bibliografi, membetulkan nama area,
  menentukan formula yang terpotong, atau menyelesaikan konflik
  klasifikasi. Tidak ada perubahan pada file TXT sumber.
\end{itemize}

\section{Review C02: Urban Cannibalism: Expanding Agglomeration Shadow
to Function Extraction in Arrested
Transition}\label{review-c02-urban-cannibalism-expanding-agglomeration-shadow-to-function-extraction-in-arrested-transition}

\textbf{Identitas sumber}

\begin{itemize}
\tightlist
\item
  Kode kajian: C02.
\item
  Judul: \emph{Urban Cannibalism: Expanding Agglomeration Shadow to
  Function Extraction in Arrested Transition}.
\item
  Penulis: Mahastri Laksita Sam Widita (pada halaman judul ditulis
  Mahastri L.S. Widita) dan Raisal Ahmad.
\item
  Tahun dan tanggal: 2026; metadata dan halaman akhir TXT menyebut
  preprint SocArXiv bertanggal 10 Juni 2026.
\item
  Publikasi/repository: SocArXiv Preprint, Version 1, melalui OSF; bukan
  artikel jurnal yang dinyatakan telah melalui peer review.
\item
  DOI: \texttt{10.31235/osf.io/brxn3\_v1}.
\item
  Afiliasi: keduanya tercantum sebagai Independent Researcher,
  Yogyakarta, Indonesia.
\item
  Kata kunci: agglomeration shadow, urban cannibalism, arrested
  transition, functional polycentricity, cannibalism index,
  Gerbangkertosusila.
\item
  Klasifikasi vault: \texttt{source/article}, mengikuti lokasi dan
  instruksi folder. Sifat dokumen yang dinyatakan dalam TXT tetap
  preprint; status peer-reviewed tidak disebutkan dan tidak diasumsikan.
\end{itemize}

\subsection{Summarized Abstract}\label{summarized-abstract-69}

\textbf{Laporan sumber:} Paper berangkat dari paradoks bahwa
konektivitas metropolitan Indonesia yang dimaksudkan untuk pemerataan
dapat bekerja terbalik. Di Gerbangkertosusila, Bangkalan sebagai
kabupaten yang paling dekat dengan Surabaya melalui Jembatan Suramadu
dilaporkan tumbuh lebih lambat daripada Pamekasan dan Sumenep. Penulis
memperluas konsep \emph{agglomeration shadow} dengan istilah
\emph{cannibalism}, yaitu mode aktif-ekstraktif ketika wilayah
metropolitan terjebak dalam \emph{arrested transition}. Konsep tersebut
dioperasionalkan melalui \emph{Cannibalism Index} (CI), rasio
\emph{external pull} terhadap \emph{local strength}, pada 149 heksagon
H3 menggunakan model gravitasi Newton-Reilly, Point of Interest (POI)
OpenStreetMap, dan kerangka \emph{urbanization economies} O'Sullivan.
Analisis dilengkapi Location Quotient (LQ) dan Shift-Share atas PDRB
sektoral.

\textbf{Laporan sumber:} Hasil utama yang diringkas abstrak adalah tujuh
heksagon \emph{Agglomeration Core}, semuanya di Kota Surabaya; tidak ada
\emph{Independent Hub} di luar inti; dan 16 heksagon \emph{Shadow Zone}.
Perbandingan Bangkalan, Mojokerto, dan Lamongan menyatakan bahwa jarak
mengurangi tarikan inti, tetapi tidak membangun kekuatan fungsional
internal. Studi kasus Mojokerto dilaporkan memperlihatkan
\emph{function-extractive split}: industri berada di Kabupaten Mojokerto
melalui Ngoro Industrial Park, sedangkan jasa profesional tetap di
Surabaya. Berdasarkan rangkaian tersebut, penulis mengklasifikasikan
Gerbangkertosusila sebagai \emph{polynuclear monocentric region in
arrested transition}, bukan sistem polisentris matang.

\textbf{Interpretasi penulis:} Implikasi perencanaan yang diajukan
adalah bahwa konektivitas tanpa dekonsentrasi fungsional bukan
pemberdayaan pinggiran, melainkan mekanisme ekstraksi yang difasilitasi
negara. \textbf{Inferensi review:} abstrak sudah mengandung diagnosis
konseptual dan rekomendasi normatif, tetapi CI sendiri dinyatakan
sebagai ukuran tekanan gravitasi atau sistem peringatan dini, bukan
bukti arus ekstraksi aktual.

\emph{(Lokator: metadata PDF, baris 1-28; Abstrak, hlm. 1-2; Bagian 4-8,
hlm. 12-25.)}

\subsection{Problem Statement}\label{problem-statement-69}

\textbf{Laporan sumber:} Masalah empiris dibingkai melalui peresmian
Jembatan Suramadu pada 2009. Pemerintah pusat berharap jembatan tersebut
membuka isolasi Madura, menjadikan Bangkalan pintu gerbang ekonomi baru,
dan mendistribusikan pertumbuhan dari Surabaya. Namun, angka pertumbuhan
PDRB yang dikutip paper menempatkan Bangkalan pada 6,29\% per tahun,
lebih rendah daripada Pamekasan 6,69\% dan Sumenep 6,36\%. Dengan
demikian, kabupaten yang lebih dekat ke inti dan terhubung infrastruktur
dilaporkan tumbuh lebih lambat daripada kabupaten yang lebih jauh. Paper
juga mengutip Efendi (2013) tentang kenaikan IMB dan munculnya delapan
pengembang real estat di Bangkalan, sementara jasa profesional,
perbankan korporat, dan layanan tingkat tinggi tetap terpusat di
Surabaya.

\textbf{Laporan sumber:} Masalah tersebut diperluas ke seluruh
Gerbangkertosusila. Industri telah dipindahkan ke kawasan pinggiran
seperti Ngoro Industrial Park di Mojokerto, JIIPE di Gresik, dan kawasan
industri Sidoarjo. Secara morfologis, keadaan ini terlihat polisentris,
tetapi fungsi yang menurut penulis menentukan kemandirian kota, yaitu
jasa profesional, kantor pusat, konsultan teknis, dan layanan tingkat
tinggi, tetap terkonsentrasi di Surabaya. Paper menilai literatur
\emph{agglomeration shadow} belum membedakan bayangan yang transisional
dari bayangan yang struktural dan ekstraktif. Model evolusi spasial
Friedmann juga diposisikan belum memadai untuk kasus ketika perpindahan
antarfase berhenti.

\textbf{Interpretasi penulis:} Ketidaksesuaian antara penyebaran
industri dan konsentrasi fungsi pengambilan keputusan disebut sebagai
kontradiksi utama. Infrastruktur yang seharusnya memperkuat pinggiran
dipahami sebagai saluran penyerapan tenaga kerja, konsumsi, dan
investasi ke inti. \textbf{Inferensi review:} problem statement
menggabungkan anomali pertumbuhan, ketimpangan fungsi, dan kekurangan
konsep, tetapi sebagian bukti awalnya berasal dari studi eksternal yang
dikutip, bukan seluruhnya dari pengukuran asli paper ini.

\emph{(Lokator: Bagian 1, hlm. 3-4; Gambar 1-2, hlm. 4; Bagian 2, hlm.
5-8.)}

\subsection{Objectives}\label{objectives-69}

Sumber tidak menyediakan subbagian terpisah berjudul Objectives. Tujuan
berikut direkonstruksi dari tujuan, kontribusi, dan susunan analisis
yang dinyatakan penulis, sehingga bukan kutipan tujuan formal yang
berdiri sendiri.

\begin{itemize}
\tightlist
\item
  \textbf{Laporan tujuan penulis:} memperluas literatur
  \emph{agglomeration shadow} dengan memberi nama pada dimensi aktif dan
  patologis dari ekstraksi fungsional ketika transisi metropolitan
  terhenti.
\item
  \textbf{Laporan tujuan penulis:} mengoperasionalkan konsep tersebut
  melalui CI yang dapat dihitung pada unit intra-metropolitan H3, dengan
  membandingkan tarikan eksternal dan kekuatan lokal.
\item
  \textbf{Laporan tujuan penulis:} memetakan struktur spasial
  Gerbangkertosusila dan mengidentifikasi \emph{Agglomeration Core},
  \emph{Independent Hub}, \emph{Shadow Zone}, serta tipologi lain yang
  dihasilkan oleh matriks 3x3.
\item
  \textbf{Laporan tujuan penulis:} menguji peran jarak dan konektivitas
  melalui perbandingan tiga arah Bangkalan, Mojokerto, dan Lamongan.
\item
  \textbf{Laporan tujuan penulis:} menelusuri mekanisme dekonsentrasi
  industri tanpa dekonsentrasi fungsional melalui studi kasus Mojokerto,
  lalu memvalidasi paradoks konektivitas melalui Bangkalan.
\item
  \textbf{Laporan tujuan penulis:} mengusulkan diagnosis dan implikasi
  perencanaan untuk mengklasifikasikan Gerbangkertosusila secara lebih
  tepat daripada klasifikasi polisentris yang hanya berbasis pertumbuhan
  morfologis.
\end{itemize}

\emph{(Lokator: Abstrak, hlm. 1-2; Bagian 1, hlm. 3-4; Bagian 5-8, hlm.
15-25.)}

\subsection{Literature Survey}\label{literature-survey-69}

\textbf{Laporan sumber:} Kerangka awal paper adalah \emph{agglomeration
shadow} Meijers dan Burger (2010), yang menjelaskan kecenderungan kota
kecil dekat metropolitan besar tumbuh lebih lambat. Penulis
menghubungkannya dengan logika Christaller dan Losch tentang tarikan
pasar serta ekonomi aglomerasi. Myrdal (1957) digunakan untuk konsep
\emph{backwash effect}, yaitu penyerapan tenaga kerja, modal, dan
kapasitas inovasi dari pinggiran ke pusat. Friedmann (1966) dipaparkan
melalui empat fase: \emph{Pre-industrial}, \emph{Transitional/Periphery
Dependence}, \emph{Industrial Dispersion}, dan \emph{Functional
Integration}.

\textbf{Laporan sumber:} Hall dan Pain (2006) serta Meijers (2008)
menjadi dasar pembedaan antara polisentrisme morfologis dan fungsional
serta antara \emph{polynuclear monocentric region} dan \emph{polycentric
urban region}. Menurut paparan paper, pusat industri, bangunan, atau
infrastruktur di pinggiran tidak otomatis menciptakan subpusat
fungsional apabila pengambilan keputusan dan jasa tingkat tinggi tetap
di inti.

\textbf{Laporan sumber:} Penulis membandingkan kerangka tersebut dengan
Sassen (2015), Florida (2005), dan Beaverstock, Smith, dan Taylor
(1999), yang dinilai berguna untuk memahami hierarki fungsional pada
skala kota global atau regional makro, tetapi tidak menyediakan
mekanisme mikro untuk unit subkota. O'Sullivan (1996) dipilih karena
konsep \emph{input sharing}, \emph{labor pooling}, \emph{labor
matching}, dan \emph{knowledge spillovers} dapat dihubungkan dengan
keberadaan POI jasa profesional.

\textbf{Laporan sumber:} Untuk model spasial, paper menyebut
Newton-Reilly (Reilly 1931) dan membandingkannya dengan model Huff
Probabilistic (Huff 1964). Huff dinilai lebih ideal untuk kompetisi
destinasi dan probabilitas pangsa pasar, tetapi memerlukan data arus
asal-tujuan yang tidak tersedia publik. Konsep polarisasi Hirschman
(1958) dinilai tidak menyediakan kerangka spasial intra-metropolitan
yang dapat langsung dioperasionalkan.

\textbf{Interpretasi penulis:} Literatur yang ada diposisikan belum
memberi istilah bagi mode \emph{active-extractive} yang terjadi ketika
transisi antara dekonsentrasi spasial dan dekonsentrasi fungsional
macet. Istilah \emph{cannibalism} diajukan sebagai kontribusi konseptual
dan secara eksplisit membawa konotasi normatif. \textbf{Inferensi
review:} survei literatur dalam paper terutama berfungsi membangun
definisi dan pemilihan indikator; TXT tidak menunjukkan tinjauan
sistematis, kriteria pencarian, atau perbandingan menyeluruh terhadap
seluruh literatur terkait.

\emph{(Lokator: Bagian 1-3, hlm. 3-12; Gambar 3-5, hlm. 6-8; Daftar
Pustaka, hlm. 26-27.)}

\subsection{Method Used}\label{method-used-69}

\textbf{Laporan sumber:} Wilayah studi adalah Gerbangkertosusila, yang
dalam paper mencakup Kota Surabaya, Kabupaten Sidoarjo, Kabupaten
Gresik, Kabupaten Lamongan, Kabupaten Mojokerto, Kota Mojokerto, dan
Kabupaten Bangkalan. Unit analisisnya adalah heksagon H3 resolusi 7
dengan luas sekitar 5 km2. Heksagon yang memotong batas administratif
diatribusikan kepada wilayah dengan luas dominan. Setelah hanya
mempertahankan heksagon yang memiliki minimal satu POI relevan, tersisa
149 heksagon.

\textbf{Laporan sumber:} Setiap heksagon diberi dua massa.
\texttt{M\_econ} adalah jumlah POI primer berbobot yang mewakili jasa
profesional: pengacara, konsultan, dan akuntan berbobot 3; IT,
telekomunikasi, dan periklanan berbobot 2; bank dan asuransi berbobot
1,5. Nilai tersebut dikalikan \texttt{1\ +\ E\_i}, dengan \texttt{E\_i}
berupa Shannon Entropy untuk memperhitungkan diversitas kategori.
\texttt{M\_social} adalah jumlah berbobot fasilitas sekunder: rumah
sakit berbobot 4, universitas 3, serta mal dan department store 2.

\textbf{Laporan sumber:} \emph{External pull} dihitung dengan model
gravitasi Newton-Reilly:
\texttt{P\_ext,i\ =\ jumlah\ M\_j\ /\ d\_ij\^{}beta} untuk semua
heksagon lain. Massa total \texttt{M\_j} adalah
\texttt{M\_econ,j\ +\ M\_social,j}; \texttt{d\_ij} adalah waktu tempuh
dalam menit dari jaringan jalan OpenStreetMap melalui OSMnx; dan
\texttt{beta\ =\ 2}. \emph{Local strength} ditetapkan sama dengan
\texttt{M\_econ}, bukan massa sosial, karena kemandirian fungsional
menurut kerangka penulis ditentukan oleh ekosistem jasa profesional. CI
kemudian dihitung sebagai \texttt{CI\_i\ =\ P\_ext,i\ /\ (L\_i\ +\ 1)},
dengan penambahan 1 untuk menghindari pembagian dengan nol.

\textbf{Laporan sumber:} Nilai \texttt{L\_i} dan \texttt{P\_ext}
dinormalisasi ke tiga kelas rendah, sedang, dan tinggi menggunakan Jenks
natural breaks. Karena \texttt{P\_ext} sangat menceng ke kanan, nilainya
terlebih dahulu ditransformasi logaritmik. Matriks 3x3 menghasilkan
sembilan tipologi: \emph{Agglomeration Core}, \emph{Semi AC},
\emph{Independent Hub}, \emph{Semi SZ}, \emph{Neutral}, \emph{Semi IH},
\emph{Shadow Zone}, \emph{Semi T}, dan \emph{Terpencil}.

\textbf{Laporan sumber:} Analisis empiris terdiri atas pemetaan tipologi
regional; cross-validation dengan pemetaan dasymetric kepadatan
penduduk; perbandingan \emph{distance-decay} tiga arah; studi kasus
sektoral Mojokerto memakai SLQ dan Shift-Share; serta validasi Bangkalan
melalui data pertumbuhan PDRB, transformasi pasca-Suramadu, dan data
APBD. Paper menetapkan sistem tertutup pada batas Gerbangkertosusila dan
menggunakan data utama yang disebut cross-sectional 2025-2026.

\textbf{Inferensi review:} Metode ini adalah kombinasi indeks spasial
deskriptif, klasifikasi, perbandingan kasus, dan analisis sektoral. TXT
tidak menjelaskan uji signifikansi, interval ketidakpastian, prosedur
estimasi parameter, kode analisis, atau bentuk statistik formal untuk
istilah ``pengujian distance-decay''; yang terlihat adalah perbandingan
tipologi, nilai, dan scatter sebagaimana dijelaskan di Bagian 4-6.

\emph{(Lokator: Bagian 3, hlm. 8-12; Tabel 1, hlm. 9; Persamaan 1-7,
hlm. 9-11; Gambar 6, hlm. 11; Bagian 4-7, hlm. 12-23.)}

\subsection{Dataset}\label{dataset-69}

\textbf{Laporan sumber:} Dataset dan bahan analisis yang disebut dalam
TXT meliputi:

\begin{itemize}
\tightlist
\item
  POI OpenStreetMap untuk jasa profesional dan fasilitas sosial, dengan
  bobot kategori pada Tabel 1.
\item
  Jaringan jalan OpenStreetMap dan library OSMnx untuk memperoleh waktu
  tempuh antarsentroid heksagon.
\item
  Geometri spasial Uber H3 resolusi 7 dan batas administratif tujuh
  kabupaten/kota.
\item
  Kepadatan penduduk untuk cross-validation dasymetric, yang disebut
  menggabungkan BPS 2025, WorldPop yang dirujuk melalui Tatem (2017),
  dan penggunaan lahan.
\item
  PDRB sektoral Kota Mojokerto dan Kabupaten Mojokerto untuk 2018-2023,
  dirujuk dari BPS Kota Mojokerto (2024) dan BPS Kabupaten Mojokerto
  (2024), yang diringkas dalam SLQ dan Shift-Share.
\item
  Angka pertumbuhan PDRB tiga kabupaten Madura, kenaikan IMB, delapan
  pengembang real estat, dan status rencana industri pasca-Suramadu dari
  studi Efendi (2013). Paper secara eksplisit menyebut data pertumbuhan
  tersebut sebagai data studi eksternal, bukan analisis asli paper.
\item
  Data APBD dari Kementerian Keuangan RI, DJPK (2026), termasuk proporsi
  PAD dan belanja modal yang dilaporkan untuk Bangkalan serta pembanding
  Sidoarjo dan Gresik.
\end{itemize}

\textbf{Inferensi review:} TXT tidak menyediakan jumlah POI per
kategori, tanggal pengambilan atau versi data OpenStreetMap, daftar
koordinat heksagon, data mentah PDRB/APBD, atau kode pengolahan. Karena
itu, replikasi numerik penuh tidak dapat dipastikan hanya dari TXT
meskipun penulis menyatakan CI kompatibel dengan data POI dan dapat
direplikasi pada wilayah lain.

\emph{(Lokator: Bagian 3, hlm. 8-12; Gambar 9, hlm. 15; Bagian 6, hlm.
18-21; Bagian 7, hlm. 21-23; Daftar Pustaka, hlm. 26-27.)}

\subsection{Population / Sample}\label{population-sample-69}

\textbf{Laporan sumber:} Populasi analisis utama bukan individu atau
rumah tangga, melainkan ruang metropolitan yang dibagi menjadi 149
heksagon H3 valid di tujuh kabupaten/kota. Validitas operasional
ditentukan oleh adanya minimal satu POI dari kategori relevan. Tabel 3
mencatat total per wilayah: Surabaya 52 heksagon, Sidoarjo 35, Gresik
18, Lamongan 18, Kabupaten Mojokerto 16, Bangkalan 9, dan Kota Mojokerto
1.

\textbf{Laporan sumber:} Perbandingan \emph{distance-decay} menggunakan
subset Bangkalan sebanyak 9 heksagon, gabungan Kota dan Kabupaten
Mojokerto sebanyak 17 heksagon, dan Lamongan sebanyak 18 heksagon. Studi
kasus sektoral hanya mencakup Kota Mojokerto dan Kabupaten Mojokerto.
Validasi pertumbuhan historis Bangkalan membandingkan Bangkalan dengan
Pamekasan dan Sumenep.

\textbf{Laporan sumber:} Empat heksagon dengan centroid tepat di tepi
pantai diatribusikan kepada kabupaten terdekat dalam distribusi
administratif Tabel 3. \textbf{Inferensi review:} ini merupakan cakupan
spasial yang disaring dan ditetapkan secara teknis, bukan sampel
probabilitas. TXT tidak menyebut populasi awal seluruh kandidat
heksagon, jumlah observasi POI sebelum penyaringan, ukuran sampel
manusia, survei responden, atau tingkat respons.

\emph{(Lokator: Bagian 3, hlm. 8-9; Tabel 2-4, hlm. 12-17; catatan kaki
Tabel 3, hlm. 13; Bagian 6, hlm. 18-21.)}

\subsection{Dependent Variables}\label{dependent-variables-69}

Sumber tidak menyatakan variabel dependen dan independen dalam
terminologi desain penelitian konvensional. Pemetaan berikut adalah
klasifikasi review atas keluaran dan indikator yang digunakan, bukan
penamaan formal dari penulis.

\begin{itemize}
\tightlist
\item
  \textbf{Keluaran utama:} \texttt{Cannibalism\ Index\ (CI\_i)}, yaitu
  rasio \texttt{P\_ext} terhadap \texttt{L\_i\ +\ 1}. CI tinggi dimaknai
  sebagai tekanan eksternal yang besar relatif terhadap kapasitas
  internal.
\item
  \textbf{Keluaran klasifikasi:} tipologi sembilan kelas dari kombinasi
  \texttt{L\_i} dan \texttt{P\_ext}, terutama \emph{Agglomeration Core},
  \emph{Independent Hub}, \emph{Shadow Zone}, dan \emph{Terpencil}.
  \emph{Shadow Zone} adalah konfigurasi \texttt{L\_i} rendah dan
  \texttt{P\_ext} tinggi.
\item
  \textbf{Komponen penjelas langsung:} \texttt{P\_ext} sebagai tarikan
  dari massa heksagon lain dan \texttt{L\_i} sebagai kekuatan jasa
  profesional lokal. Waktu tempuh, konektivitas, komposisi POI, dan
  posisi relatif terhadap Surabaya membentuk kondisi yang dibaca melalui
  kedua komponen tersebut.
\item
  \textbf{Indikator sektoral dan kontekstual:} SLQ, komponen
  Differential Shift, pertumbuhan PDRB, PAD, dan belanja modal digunakan
  untuk membaca dekonsentrasi industri, kekuatan jasa, dinamika daya
  saing, dan kerentanan fiskal. Paper tidak menyebut salah satunya
  sebagai dependent variable formal.
\end{itemize}

\textbf{Inferensi review:} CI lebih tepat dibaca sebagai indikator
diagnostik atau outcome tipologis daripada ukuran langsung
``ekstraksi'', karena penulis sendiri membedakan \emph{pressure} dari
\emph{flow}. Tidak ada variabel arus aktual, seperti arus komuter,
migrasi, transaksi, atau aliran investasi, yang menjadi outcome terukur
dalam model ini.

\emph{(Lokator: Bagian 3, hlm. 9-12; Persamaan 4-7, hlm. 10-11; Bagian
5-7, hlm. 15-23.)}

\subsection{Results}\label{results-69}

\textbf{Laporan sumber, pola regional:} Dari 149 heksagon,
\emph{Agglomeration Core} berjumlah 7 (4,7\%), \emph{Semi AC} 0,
\emph{Independent Hub} 0, \emph{Semi SZ} 7 (4,7\%), \emph{Neutral} 9
(6,0\%), \emph{Semi IH} 0, \emph{Shadow Zone} 16 (10,7\%), \emph{Semi T}
57 (38,3\%), dan \emph{Terpencil} 53 (35,6\%). \emph{Terpencil} dan
\emph{Semi T} bersama-sama mencakup 110 heksagon atau 73,9\%. Tujuh
\emph{Agglomeration Core} seluruhnya berada di Surabaya. Tidak ada
\emph{Independent Hub} pada distribusi 149 heksagon.

\textbf{Laporan sumber, distribusi administratif:} Surabaya memiliki 7
\emph{Agglomeration Core}, 11 \emph{Shadow Zone}, 7 \emph{Semi SZ}, 2
\emph{Terpencil}, 17 \emph{Semi T}, dan 8 \emph{Neutral}. Sidoarjo
memiliki 1 \emph{Shadow Zone}, 7 \emph{Terpencil}, 26 \emph{Semi T}, dan
1 \emph{Neutral}. Gresik memiliki 8 \emph{Terpencil} dan 10 \emph{Semi
T}. Lamongan seluruhnya 18 \emph{Terpencil}; Kabupaten Mojokerto
seluruhnya 16 \emph{Terpencil}; Bangkalan memiliki 4 \emph{Shadow Zone},
1 \emph{Terpencil}, dan 4 \emph{Semi T}; Kota Mojokerto memiliki 1
\emph{Terpencil}. Dengan demikian, empat \emph{Shadow Zone} Bangkalan
berada pada sisi koridor Suramadu, sedangkan satu heksagon Sidoarjo dan
11 heksagon Surabaya membentuk lokasi lainnya.

\textbf{Laporan sumber, cross-validation spasial:} Keenam belas
\emph{Shadow Zone} dikatakan seluruhnya tumpang-tindih dengan hotspot
kepadatan penduduk pada analisis dasymetric, bukan coldspot atau wilayah
pinggiran yang jarang penduduk. Penulis menggunakan hasil ini untuk
menolak pembacaan bahwa \emph{Shadow Zone} sekadar artefak penyaringan
POI atau efek tepi metodologis.

\textbf{Laporan sumber, perbandingan jarak:} Bangkalan memiliki 4 dari 9
heksagon (44\%) dalam kategori \emph{Shadow Zone} dan tidak memiliki
heksagon kuat. Lima sisanya terdiri atas 1 \emph{Terpencil} dan 4
\emph{Semi T}. Gabungan Kota dan Kabupaten Mojokerto memiliki 17
heksagon dan seluruhnya \emph{Terpencil}, termasuk heksagon yang memuat
Ngoro Industrial Park. Lamongan memiliki 18 heksagon dan seluruhnya
\emph{Terpencil}. Tidak ada \emph{Agglomeration Core}, \emph{Semi AC},
atau \emph{Independent Hub} pada ketiga wilayah pembanding. Tabel 4
merangkum 4 heksagon lemah-tertarik di Bangkalan, 0 di Mojokerto, dan 0
di Lamongan; heksagon terlindung-lemah masing-masing 5, 17, dan 18.

\textbf{Laporan sumber, hasil sektoral Mojokerto:} Rata-rata SLQ
2018-2023 menunjukkan Kota Mojokerto kuat pada jasa keuangan (sekitar
2,85), jasa lainnya (2,58), komunikasi (2,50), pemerintahan (2,04),
kesehatan (1,76), pendidikan (1,62), real estat (1,58), perdagangan
(1,56), serta makan minum (1,19), tetapi lemah pada industri pengolahan
(0,33) dan pertanian (0,05). Kabupaten Mojokerto kuat pada industri
pengolahan (1,85), sementara komunikasi 1,09 dan pemerintahan 1,00
berada sekitar representasi rata-rata; sektor lain yang tercantum lebih
rendah, termasuk keuangan 0,56, perdagangan 0,56, pendidikan 0,48,
kesehatan 0,58, dan jasa perusahaan 0,19. Jasa perusahaan kategori M,N
bernilai sekitar 0,92 di kota dan 0,19 di kabupaten.

\textbf{Laporan sumber, Shift-Share Mojokerto:} Differential Shift
industri pengolahan Kota Mojokerto berubah dari +4.954 sebelum COVID
menjadi +214 selama COVID dan -41.728 setelah COVID. Pada Kabupaten
Mojokerto, nilainya tetap positif tetapi melemah dari +572.390 menjadi
+562.392 lalu +327.437. Differential Shift jasa perusahaan Kabupaten
Mojokerto bernilai -2.071 sebelum COVID, -959 selama COVID, dan +220
setelah COVID, sementara SLQ absolutnya tetap 0,19.

\textbf{Laporan sumber, Bangkalan dan fiskal:} Data eksternal yang
dikutip untuk Bangkalan menunjukkan pertumbuhan 6,29\%, dibandingkan
Pamekasan 6,69\% dan Sumenep 6,36\%. Paper juga melaporkan peningkatan
IMB, delapan pengembang real estat, rencana industri dan Pelabuhan
Internasional Klampis yang belum terealisasi, serta tetap terpusatnya
jasa profesional, teknologi, dan layanan tingkat tinggi di Surabaya.
Data DJPK yang dikutip menunjukkan PAD Bangkalan 22,9\% dari total
pendapatan terealisasi dan belanja modal 32\%; sebagai pembanding, PAD
Sidoarjo 44,1\% dan belanja modal Gresik disebut 66\%.

\emph{(Lokator: Bagian 4, hlm. 12-15; Tabel 2-3; Gambar 7-9; Bagian 5,
hlm. 15-18; Tabel 4; Gambar 10-11; Bagian 6, hlm. 18-21; Tabel 5; Gambar
12-13; Bagian 7, hlm. 21-23.)}

\subsection{Findings}\label{findings-69}

\textbf{Interpretasi penulis:} Empat lapisan bukti diperlakukan sebagai
konvergen: konsentrasi kekuatan fungsional di Surabaya; perlindungan
jarak tanpa pembentukan kekuatan internal; \emph{function-extractive
split} Mojokerto; dan paradoks konektivitas Bangkalan. Penulis
menyimpulkan bahwa dekonsentrasi industri tidak sama dengan
dekonsentrasi fungsional, sehingga struktur Gerbangkertosusila secara
morfologis dapat tampak memiliki banyak pusat tetapi secara fungsional
tetap monocentric.

\textbf{Interpretasi penulis:} Penulis membedakan dua bentuk
\emph{cannibalism}. Mojokerto adalah bentuk struktural-historis karena
dekonsentrasi industri telah terjadi tanpa jasa profesional pendamping.
Bangkalan adalah bentuk \emph{connectivity-induced} karena infrastruktur
Suramadu memperdalam tarikan pada lokasi yang sebelumnya dilindungi oleh
\emph{friction of distance}. Tiga faktor yang diajukan adalah
konektivitas yang mengurangi jarak efektif, ketiadaan kekuatan lokal,
dan absennya dekonsentrasi fungsional.

\textbf{Inferensi review:} Temuan terkuat dari bahan TXT adalah pola
klasifikasi dan ketidakseimbangan sektoral, bukan pembuktian arus
ekstraksi. Absennya \texttt{L\_i} di sekitar Ngoro juga sebagian
merupakan konsekuensi definisi metode, karena \texttt{L\_i} hanya
menghitung jasa profesional dan tidak memasukkan volume industri
manufaktur. Hubungan fiskal Bangkalan dengan kerentanan fungsional
disajikan sebagai mekanisme yang saling memperkuat oleh penulis, tetapi
TXT tidak menunjukkan model kausal atau uji yang mengisolasi pengaruh
PAD dan belanja modal.

\emph{(Lokator: Bagian 5, hlm. 17-18; Bagian 6, hlm. 18-21; Bagian 7,
hlm. 21-23; Bagian 8, hlm. 23-25.)}

\subsection{Conclusions}\label{conclusions-69}

\textbf{Laporan kesimpulan sumber:} Gerbangkertosusila dinilai bukan
sistem polisentris matang, melainkan \emph{polynuclear monocentric
region in arrested transition}. Industri yang tersebar di Ngoro, JIIPE,
dan kawasan industri Sidoarjo belum diikuti pertumbuhan ekosistem jasa
profesional yang dapat menjadikan wilayah-wilayah tersebut subpusat
mandiri. CI memetakan 16 \emph{Shadow Zone} sebagai zona prioritas,
tetapi penulis menegaskan bahwa indeks itu mengukur potensi tekanan dan
bukan arus aktual.

\textbf{Interpretasi penulis:} Kebijakan dekonsentrasi industri tanpa
ekosistem jasa profesional menghasilkan struktur pemasok-konsumen yang
permanen. Infrastruktur konektivitas seperti Suramadu, Tol Trans Jawa,
jaringan kereta Surabaya-Madiun, dan rencana KA Lyn Surabaya-Bangkalan
diperingatkan dapat berfungsi sebagai infrastruktur ekstraksi selama
kondisi \emph{arrested transition} belum berubah. Solusi yang diajukan
adalah memindahkan industri sekaligus menumbuhkan jasa, serta membangun
jalan sekaligus institusi subpusat.

\textbf{Inferensi review:} Kesimpulan tersebut konsisten dengan definisi
konsep dan indikator yang dipilih, tetapi cakupannya harus dibaca
bersama batasan sumber: data utama bersifat cross-sectional, sebagian
validasi berasal dari sumber eksternal, dan hubungan sebab-akibat
konektivitas dengan pertumbuhan tidak diestimasi secara longitudinal.

\emph{(Lokator: Bagian 8, hlm. 23-25; Gambar 14, hlm. 24; Bagian 9, hlm.
25-26.)}

\subsection{Contributions}\label{contributions-69}

\textbf{Klaim kontribusi teoretis penulis:} Paper memperkenalkan
\emph{cannibalism} untuk membedakan \emph{agglomeration shadow} yang
dianggap netral atau transisional dari mode aktif, struktural, dan
ekstraktif ketika dekonsentrasi fungsional tidak mengikuti dekonsentrasi
spasial. Penulis menyebut pembedaan \emph{shadow} transisional dan
\emph{shadow} struktural-patologis sebagai sesuatu yang belum tersedia
dalam literatur yang mereka bahas.

\textbf{Klaim kontribusi metodologis penulis:} CI mengoperasionalkan
rasio tekanan eksternal dan kekuatan lokal pada skala
intra-metropolitan, menggunakan H3 dan model gravitasi yang dibangun di
atas kerangka O'Sullivan serta kompatibel dengan POI OpenStreetMap.
Penulis menyatakan metode ini dapat dipetakan dan direplikasi di wilayah
metropolitan lain di Indonesia.

\textbf{Klaim kontribusi empiris penulis:} Paper memetakan
Gerbangkertosusila dengan tujuh \emph{Agglomeration Core}, 16
\emph{Shadow Zone}, tanpa \emph{Independent Hub}, serta
mendemonstrasikan dua mekanisme melalui Mojokerto dan Bangkalan.
Kontribusi empiris ini dipakai untuk menantang klasifikasi fase 3 yang
dikaitkan paper dengan Putra dan Salim (2022).

\textbf{Inferensi review:} Kontribusi yang benar-benar dapat dinilai
dari TXT adalah usulan konsep, konstruksi indikator, dan penerapan pada
satu wilayah. Klaim bahwa pembedaan tersebut belum ada atau bahwa indeks
dapat direplikasi luas tetap merupakan klaim penulis; TXT tidak memuat
validasi lintas metropolitan atau pembuktian bahwa terminologi tersebut
telah diterima disiplin.

\emph{(Lokator: Bagian 1, hlm. 3-4; Bagian 2, hlm. 5-8; Bagian 8, hlm.
23-25.)}

\subsection{Research Gap}\label{research-gap-69}

\textbf{Gap yang ditargetkan penulis:}

\begin{itemize}
\tightlist
\item
  Literatur \emph{agglomeration shadow} yang dibahas paper menjelaskan
  pertumbuhan lambat di pinggiran, tetapi tidak membedakan kondisi
  transisional dari kondisi struktural-ekstraktif.
\item
  Model empat fase Friedmann dipandang mengandaikan transisi yang
  berjalan, sedangkan kasus Gerbangkertosusila digunakan untuk
  menunjukkan transisi yang terhenti antara ketergantungan pinggiran dan
  dispersi industri.
\item
  Kerangka Sassen, Florida, dan GaWC disebut beroperasi pada skala kota
  global atau regional makro, bukan mekanisme mikro pada unit subkota
  dalam satu metropolitan.
\item
  Studi terdahulu yang dikritik paper dinilai membaca polisentrisme dari
  penyebaran pusat industri, tanpa menguji apakah jasa profesional dan
  fungsi pengambilan keputusan ikut terdesentralisasi.
\end{itemize}

\textbf{Gap yang masih terbuka menurut teks, sebagai inferensi review:}

\begin{itemize}
\tightlist
\item
  Belum ada pengukuran arus komuter, migrasi, telekomunikasi, konsumsi,
  atau investasi yang dapat memverifikasi bahwa tekanan CI benar-benar
  menjadi ekstraksi.
\item
  Belum ada lintasan historis CI atau desain sebelum-sesudah yang dapat
  menguji klaim sebab-akibat mengenai Suramadu.
\item
  Bukti sektoral langsung hanya tersedia untuk Mojokerto; penerapan
  diagnosis pada Bangkalan, Sidoarjo, Gresik, dan Lamongan masih
  memerlukan validasi sektoral per kasus.
\item
  Belum dijelaskan sensitivitas hasil terhadap \texttt{beta}, bobot POI,
  resolusi H3, batas sistem, serta definisi kategori jasa profesional.
\item
  Belum tersedia pengujian eksternal pada metropolitan lain dan belum
  ada informasi peer review dalam TXT yang dapat menilai penerimaan
  akademik atas konsep baru tersebut.
\end{itemize}

\emph{(Lokator: Bagian 1-3, hlm. 3-12; Bagian 8-9, hlm. 23-26.)}

\subsection{Limitations}\label{limitations-69}

Sumber memiliki bagian khusus Keterbatasan dan mengakui batas-batas
berikut.

\begin{enumerate}
\def\labelenumi{\arabic{enumi}.}
\tightlist
\item
  CI mengukur potensi tekanan gravitasi, bukan arus ekstraksi aktual.
  Validasi arus memerlukan data komuter Sakernas, migrasi SP2020, atau
  mobilitas telco yang menurut paper tidak diakses karena berbayar atau
  tidak tersedia publik.
\item
  Data utama disebut cross-sectional 2025-2026 dan tidak menelusuri
  trajektori temporal. Klaim kausal historis, termasuk bahwa Suramadu
  menyebabkan \emph{cannibalism} Bangkalan, bertumpu pada teori dan data
  eksternal, bukan analisis multi-periode asli.
\item
  Analisis sektoral hanya dilakukan untuk Kota dan Kabupaten Mojokerto.
  Klaim tentang Bangkalan, Sidoarjo, Gresik, dan Lamongan dinyatakan
  sebagai ekstrapolasi teoretis dari pola Mojokerto dan struktur
  gravitasi.
\item
  \texttt{Beta\ =\ 2} mengikuti benchmark Newtonian dan tidak
  dikalibrasi dari perilaku perjalanan lokal; analisis sensitivitas
  tidak dilakukan. Model yang dipakai adalah gravitasi klasik, bukan
  Huff Probabilistic. Sistem tertutup mengabaikan tarikan Jakarta,
  Malang, dan Bali serta dapat menimbulkan bias batas pada Lamongan
  utara dan Bangkalan timur.
\item
  POI OpenStreetMap cenderung kurang menangkap ekonomi informal dan
  tidak membedakan kantor pusat dengan cabang. MAUP tetap menjadi risiko
  pada heksagon sekitar 5 km2, sedangkan bobot POI ditetapkan
  berdasarkan pertimbangan teoretis dan bukan kalibrasi empiris.
\item
  Istilah \emph{cannibalism} adalah usulan paper, bukan terminologi
  mapan dalam uraian sumber. Penyebutan \emph{arrested transition}
  sebagai kondisi patologis dan penggunaan \texttt{M\_econ} sebagai
  ukuran kekuatan lokal bersifat theory-laden serta dapat diperdebatkan
  oleh kerangka lain.
\end{enumerate}

\textbf{Inferensi review:} Selain batasan eksplisit tersebut, TXT tidak
menjelaskan ketidakpastian klasifikasi, validasi prediktif, atau
harmonisasi penuh antara data cross-sectional 2025-2026, PDRB 2018-2023,
dan bukti historis Efendi (2013). Hal ini tidak membatalkan pola
deskriptif yang dilaporkan, tetapi membatasi klaim kausal dan
generalisasi.

\emph{(Lokator: Bagian 3, hlm. 11-12; Bagian 8, hlm. 24-25; Bagian 9,
hlm. 25-26.)}

\subsection{Challenges}\label{challenges-69}

Sumber tidak memiliki subbagian terpisah bernama Challenges. Tantangan
penelitian yang dinyatakan atau dapat diidentifikasi langsung dari
uraian metode dan keterbatasannya adalah:

\begin{itemize}
\tightlist
\item
  Ketiadaan data asal-tujuan publik membuat penulis memilih
  Newton-Reilly, bukan Huff Probabilistic, sehingga tarikan dihitung
  sebagai tekanan model dan bukan probabilitas pangsa atau arus yang
  diamati.
\item
  Kalibrasi parameter \texttt{beta} melalui perilaku perjalanan lokal
  tidak tersedia, dan TXT menyebut tidak ada analisis sensitivitas
  terhadap nilai tersebut.
\item
  Konstruksi massa membutuhkan klasifikasi POI, bobot ambang
  Christaller, dan Shannon Entropy; keputusan ini mengubah arti
  ``kekuatan'' sesuai teori O'Sullivan dan tidak dikalibrasi empiris.
\item
  Penggunaan sistem tertutup menimbulkan tantangan batas karena pusat di
  luar Gerbangkertosusila diabaikan dan sel tepi memiliki lebih sedikit
  tetangga.
\item
  Penggabungan heksagon H3 dengan wilayah administratif membutuhkan
  aturan atribusi, termasuk empat centroid di tepi pantai yang
  ditempatkan ke kabupaten terdekat.
\item
  Keterbatasan POI dalam menangkap ekonomi informal dan membedakan
  kantor pusat/cabang menyulitkan pengukuran fungsi profesional.
\item
  Keterbatasan cakupan data sektoral membuat mekanisme yang
  didemonstrasikan di Mojokerto harus diekstrapolasi ketika membaca
  wilayah lain.
\end{itemize}

\textbf{Inferensi review:} Perbedaan periode dan asal data juga
merupakan tantangan interpretasi yang perlu diperhatikan ketika hasil
indeks masa kini dibandingkan dengan PDRB 2018-2023 atau bukti
pasca-Suramadu dari studi eksternal. TXT tidak menjelaskan prosedur
penyesuaian temporalnya.

\emph{(Lokator: Bagian 2-3, hlm. 7-12; catatan kaki Tabel 3, hlm. 13;
Bagian 5-7, hlm. 15-23; Bagian 9, hlm. 25-26.)}

\subsection{Practical Implications}\label{practical-implications-69}

\textbf{Interpretasi penulis:} Kerangka koordinasi antaryurisdiksi yang
mengasumsikan subpusat telah mandiri, termasuk pendekatan BKSP yang
dicontohkan paper, dinilai belum cukup untuk Gerbangkertosusila karena
kekuatan fungsional subpusat belum terbentuk. Koordinasi perlu
disesuaikan dengan ketimpangan kapasitas, bukan menganggap semua wilayah
sebagai mitra setara.

\textbf{Interpretasi penulis:} Kebijakan pemindahan industri ke Ngoro,
JIIPE, dan kawasan industri Sidoarjo perlu disertai kebijakan yang
memungkinkan tumbuhnya jasa profesional, kantor pusat, konsultan,
layanan teknologi, dan fungsi pengambilan keputusan di subpusat. Jika
hanya industri yang dipindahkan, hubungan pemasok-konsumen dapat menetap
tanpa menghasilkan polisentrisme fungsional.

\textbf{Interpretasi penulis:} Investasi konektivitas seperti Suramadu
perlu dinilai bersama dekonsentrasi fungsional. Dalam pembacaan paper,
menghilangkan \emph{friction of distance} tanpa kapasitas lokal dapat
memperbesar ekstraksi. Intervensi yang diajukan harus menggeser industri
sekaligus menumbuhkan jasa, serta membangun infrastruktur fisik
sekaligus institusi subpusat.

\textbf{Inferensi review:} Rekomendasi ini merupakan implikasi normatif
dan perencanaan dari model, bukan evaluasi atas kebijakan yang sudah
dilaksanakan. TXT tidak mengukur dampak atau biaya implementasi
rekomendasi, dan tidak menunjukkan instrumen kebijakan rinci, aktor
pelaksana, atau urutan intervensi.

\emph{(Lokator: Bagian 1, hlm. 3-4; Bagian 8, hlm. 23-25.)}

\subsection{Applications}\label{applications-69}

\textbf{Laporan sumber:} CI dan matriks tipologi dirancang sebagai
\emph{early warning system} bagi perencana. Dalam kasus ini, 16
\emph{Shadow Zone} diposisikan sebagai zona prioritas intervensi. Peta
CI, distribusi administratif, dan cross-validation kepadatan penduduk
dapat digunakan untuk mengidentifikasi lokasi yang kuat ditarik inti
tetapi memiliki kekuatan profesional lokal rendah.

\textbf{Laporan sumber:} Kerangka H3, POI, waktu tempuh, massa
ekonomi-sosial, LQ, dan Shift-Share dapat diterapkan untuk membaca
wilayah metropolitan lain di Indonesia, menurut klaim metodologis
penulis. Analisis sektoral Mojokerto menyediakan contoh penggunaan data
PDRB untuk memeriksa apakah dekonsentrasi industri disertai jasa
profesional.

\textbf{Inferensi review:} Aplikasi paling aman dari bukti TXT adalah
penyaringan awal lokasi dan perumusan hipotesis kebijakan, bukan
penetapan bahwa ekstraksi telah terbukti atau keputusan investasi
tunggal. Setiap penggunaan di luar Gerbangkertosusila perlu menguji
kelengkapan POI, memilih batas sistem yang tepat, mengkalibrasi
parameter, dan memverifikasi arus aktual.

\emph{(Lokator: Bagian 3-5, hlm. 8-18; Bagian 8, hlm. 23-25; Gambar
7-11.)}

\subsection{Future Research
(Optional)}\label{future-research-optional-69}

\textbf{Agenda yang disebutkan atau dituntut langsung oleh sumber:}

\begin{itemize}
\tightlist
\item
  Menelusuri trajektori historis CI melalui data multi-periode untuk
  menguji klaim perubahan dan hubungan Suramadu dengan Bangkalan.
\item
  Memvalidasi tekanan model dengan arus aktual, termasuk data komuter
  Sakernas, migrasi SP2020, atau mobilitas telco apabila tersedia.
\item
  Memvalidasi klaim \emph{function-extractive cannibalism} secara
  sektoral pada Bangkalan, Sidoarjo, Gresik, dan Lamongan, bukan hanya
  melalui ekstrapolasi dari Mojokerto.
\item
  Mengkalibrasi \texttt{beta} dari perilaku perjalanan lokal dan
  melakukan analisis sensitivitas karena klasifikasi dapat berubah pada
  nilai parameter berbeda.
\item
  Memeriksa kembali batas sistem, efek tepi, POI informal, hierarki
  kantor pusat dan cabang, MAUP, serta bobot POI.
\end{itemize}

\textbf{Inferensi review yang diberi label:} Studi komparatif pada
metropolitan lain dapat menguji apakah konsep \emph{polynuclear
monocentric region in arrested transition} dan istilah
\emph{cannibalism} berlaku di luar Gerbangkertosusila. TXT menyatakan
kemungkinan replikasi, tetapi belum melaporkan pelaksanaan atau hasil
replikasi tersebut.

\emph{(Lokator: Bagian 3, hlm. 11-12; Bagian 8, hlm. 23-25; Bagian 9,
hlm. 25-26.)}

\subsection{Provenance Notes}\label{provenance-notes-69}

\begin{itemize}
\tightlist
\item
  Review ini hanya menggunakan file TXT C02 yang diminta, dibaca
  seluruhnya dari baris 1 sampai 1.114. Tidak ada sumber eksternal,
  penelusuran web, atau verifikasi terpisah atas daftar pustaka yang
  digunakan.
\item
  Blok metadata TXT menyebut sumber PDF
  \texttt{working-paper/C02-widita-ahmad-2026-urban-cannibalism-expanding-agglomeration-shadow-to-function-extraction-in-arrested-transition-10.31235-osf.io-brxn3\_v1.pdf},
  diekstrak pada 2026-08-31 dengan \texttt{pdftotext\ -layout}.
\item
  Metadata PDF menyebut 27 halaman, ukuran A4, tidak terenkripsi, judul
  dan DOI sebagaimana tercantum di bagian identitas, serta tanggal
  pembuatan dan modifikasi 10 Juni 2026. Metadata tersebut diperlakukan
  sebagai isi sumber, bukan sebagai hasil verifikasi independen.
\item
  TXT memuat abstrak bahasa Inggris dan bahasa Indonesia, halaman isi,
  Tabel 1-5, Gambar 1-14, daftar pustaka, afiliasi, dan informasi
  SocArXiv pada halaman akhir. Locator dalam review mengikuti nomor
  halaman PDF/bagian/tabel/gambar yang tampak pada hasil ekstraksi;
  baris TXT digunakan untuk memastikan isi metadata dan teks.
\item
  Label \textbf{Laporan sumber} merujuk pada data, metode, hasil, atau
  pernyataan yang dilaporkan paper. Label \textbf{Interpretasi penulis}
  merujuk pada makna konseptual, diagnosis, atau rekomendasi yang
  dibangun penulis. Label \textbf{Inferensi review} merujuk pada
  pembacaan analitis yang diturunkan dari TXT dan bukan klaim penulis.
\item
  URL arsip permanen selain DOI tidak disebutkan dalam TXT. Status peer
  review juga tidak disebutkan; karena dokumen secara eksplisit
  berstatus SocArXiv preprint Version 1, review ini tidak
  memperlakukannya sebagai artikel peer-reviewed.
\item
  Informasi yang tidak tersedia dalam TXT dan tidak diisi dengan dugaan
  mencakup raw data, kode, jumlah POI per kategori, tanggal unduh OSM,
  prosedur rinci cross-validation, statistik ketidakpastian, uji
  sensitivitas, dan data arus aktual.
\end{itemize}

\end{document}
